\documentclass[12pt,a4paper,oneside]{report}

% --- Packages de base ---
\usepackage[utf8]{inputenc}
\usepackage[T1]{fontenc}
\usepackage[french]{babel}
\usepackage{graphicx}
\usepackage{geometry}
\usepackage{amsmath}
\usepackage{amsfonts}
\usepackage{amssymb}
\geometry{hmargin=2.5cm,vmargin=2.5cm, bindingoffset=0.5cm}
\usepackage{setspace}
\onehalfspacing % Interligne 1.5 obligatoire pour les PFE

% --- Optimisation de la Typographie (Style Arial/Helvetica) ---
\usepackage[scaled]{helvet}    % Charge Helvetica (quasi identique à Arial)
\renewcommand{\familydefault}{\sfdefault} % Définit le Sans-Serif par défaut pour tout le document
\usepackage{microtype}         % Améliore l'espacement et l'alignement des blocs
\usepackage{parskip}           % Ajoute un léger espace entre les paragraphes (plus moderne)

% --- Packages pour les graphiques et schémas ---
\usepackage{tikz}
\usepackage{pgfgantt}
\usepackage{circuitikz}
\usetikzlibrary{shapes.geometric, arrows.meta, positioning, calc, fit, backgrounds, patterns}

% --- Packages de mise en forme professionnelle ---
\usepackage{longtable}          % Charger longtable AVANT xcolor[table]
\usepackage[table,xcdraw]{xcolor}
\usepackage{fontawesome5} % Icônes professionnelles
\usepackage{tcolorbox}
\tcbuselibrary{most} % Bibliothèques avancées pour tcolorbox
\usepackage{enumitem}
\usepackage{titlesec}
\usepackage{tocloft}
\usepackage[colorlinks=true, allcolors=macouleur]{hyperref} % Liens cliquables professionnels

% --- Packages additionnels (enrichissement PFE) ---
\usepackage{colortbl}           % Couleurs dans les tableaux (rowcolor)
\usepackage{booktabs}           % Traits de tableaux professionnels
\usepackage{multirow}           % Fusion de cellules dans les tableaux
\usepackage{array}              % Options de colonnes dans les tableaux
\usepackage{pgfplots}           % Graphiques et histogrammes PGF
\pgfplotsset{compat=1.18}
\usepackage{pgf-umlsd}          % Diagrammes de séquence UML
\usepackage{siunitx}            % Unités SI normalisées
\usepackage{fancyhdr}           % En-têtes et pieds de page personnalisés
\setlength{\headheight}{40pt}   % Fixe le warning fancyhdr (titres longs sur 2 lignes)
\addtolength{\topmargin}{-25pt} % Compense l'augmentation de headheight
\sisetup{locale=FR, output-decimal-marker={,}}
\usepackage{listings}           % Affichage du code source (C++, Dart)
\lstset{
    basicstyle=\ttfamily\footnotesize,
    breaklines=true,
    keywordstyle=\color{blue},
    commentstyle=\color{gray},
    stringstyle=\color{orange!70!black},
    frame=single, 
    rulecolor=\color{macouleur!30},
    numbers=left,
    numberstyle=\tiny\color{gray}
}
\usepackage{caption}            % Mise en forme des légendes de figures/tableaux
\usepackage{subcaption}         % Sous-figures (a), (b), etc.
\usepackage{float}              % Placement H absolu des flottants
\usepackage{mdframed}           % Encadrés améliorés
\usepackage{verbatim}           % Blocs de code verbatim
\usepackage[T1]{fontenc}
\usepackage{lipsum}             % Texte de test
\tcbuselibrary{skins, breakable}% TColorBox étendu

% Définition de la couleur bleue pro (similaire aux exemples)
\definecolor{macouleur}{RGB}{31, 78, 121}

% Personnalisation des titres (Bleu pour s'aligner sur le standard professionnel)
\titleformat{\chapter}[display]
  {\normalfont\huge\bfseries\color{macouleur}}{\chaptertitlename\ \thechapter}{20pt}{\Huge}
\titleformat{\section}
  {\normalfont\Large\bfseries\color{macouleur}}{\thesection}{1em}{}
\titleformat{\subsection}
  {\normalfont\large\bfseries\color{macouleur}}{\thesubsection}{1em}{}

% Rangement des listes de figures et tableaux pour qu'ils soient bleus aussi
\renewcommand{\cfttoctitlefont}{\hfill\Huge\bfseries\color{macouleur}}
\renewcommand{\cftloftitlefont}{\hfill\Huge\bfseries\color{macouleur}}
\renewcommand{\cftlottitlefont}{\hfill\Huge\bfseries\color{macouleur}}

% Configuration des en-têtes et pieds de page pro (Fancyhdr)
\pagestyle{fancy}
\fancyhf{} % Nettoie les en-têtes et pieds de page
\fancyhead[L]{\parbox[b]{0.45\textwidth}{\small\bfseries\color{macouleur}\nouppercase{\leftmark}}} % Chapitre à gauche, largeur limitée
\fancyhead[R]{\parbox[b]{0.45\textwidth}{\raggedleft\small\bfseries\color{macouleur}\nouppercase{\rightmark}}} % Section à droite, largeur limitée
\fancyfoot[C]{\textbf{\thepage}} % Numéro de page centré
\fancyfoot[R]{\footnotesize\color{gray}ENI-ABT -- PFE CNC 5-Axes} % Texte en pied de page
\renewcommand{\headrulewidth}{0.8pt} % Épaisseur de la ligne d'en-tête
\renewcommand{\footrulewidth}{0.4pt} % Épaisseur de la ligne de pied
\renewcommand{\headrule}{\color{macouleur}\hrule width \headwidth height \headrulewidth \vskip-\headrulewidth} % Ligne d'en-tête bleue
\renewcommand{\footrule}{\color{macouleur}\hrule width \headwidth height \footrulewidth \vskip-\footrulewidth} % Ligne de pied bleue

% --- Début du document ---
\begin{document}

% !TEX root = main.tex

% --- Page de Résumé / Abstract ---
\chapter*{Résumé}
\addcontentsline{toc}{chapter}{Résumé}

Le présent Projet de Fin d'Études décrit de manière holistique la conception, le dimensionnement, et la programmation d'une micro-fraiseuse à commande numérique (CNC) 5 axes dotée d'une table trunnion basculante rotative. L'enjeu de cette étude est de démocratiser l'accès à l'usinage tridimensionnel complexe, traditionnellement réservé aux industries lourdes, en développant une solution "Desktop" abordable et open-source.

Sur le plan mécanique, une analyse conceptuelle (SysML, RDM) a permis de sélectionner des composants rigides capables d'absorber le couple de renversement. Sur le plan mécatronique et logiciel, le défi de la cinématique inverse a été relevé grâce à une architecture logicielle distribuée innovante. Un microcontrôleur ESP32 gère en temps réel le calcul matriciel complexe et le cadençage "Look-Ahead" des moteurs via FreeRTOS, tandis qu'une application multiplateforme développée en Flutter (Dart) assure la supervision et la traduction du G-Code, le tout de manière fluide grâce au système d'état réactif Riverpod. Les tests pratiques (fraisage simultané 5 axes) confirment la capacité de la machine à produire des géométries d'avant-garde.

\vspace{0.5cm}
\noindent\textbf{Mots-clés :} CNC 5 axes, Fraisage, Trunnion, ESP32, Flutter, Cinématique Inverse, Mécatronique, G-Code, AMDEC.

\vspace{1.5cm}

% --- Abstract ---
\begin{center}
    \textbf{\Large Abstract}
\end{center}
\addcontentsline{toc}{chapter}{Abstract}

This final year project holistically chronicles the design, sizing, and programming of a 5-axis computer numerical control (CNC) micro-milling machine equipped with a rotary trunnion table. The challenge of this thesis is to democratize access to complex three-dimensional machining—traditionally reserved for heavy industry—by developing an affordable and open-source "Desktop" solution.

On the mechanical side, a conceptual analysis (SysML, strength of materials) led to the selection of rigid components capable of absorbing overturning torque. On the mechatronic and software side, the challenge of inverse kinematics was met with an innovative distributed software architecture. An ESP32 microcontroller manages complex matrix calculations and motor "Look-Ahead" sequencing in real-time via FreeRTOS, while a cross-platform application developed in Flutter (Dart) handles supervision and G-Code translation smoothly using the Riverpod reactive state system. Practical tests (simultaneous 5-axis milling) validate the machine's ability to produce avant-garde geometries.

\vspace{0.5cm}
\noindent\textbf{Keywords :} 5-axis CNC, Milling, Trunnion, ESP32, Flutter, Inverse Kinematics, Mechatronics, G-Code, FMEA.

% !TEX root = main.tex

% --- Dédicaces ---
\chapter*{Dédicaces}
\addcontentsline{toc}{chapter}{Dédicaces}

\vspace{2cm}
\begin{center}
    \begin{minipage}{0.8\textwidth}
        \large\itshape
        À mes très chers parents,\\[0.5cm]
        Aucun mot ne saurait exprimer mon profond respect et ma gratitude pour les sacrifices que vous avez consentis pour mon éducation et mon bien-être. Ce travail est le fruit de vos prières, de votre patience et de votre soutien inconditionnel.\\[1cm]
        À mes frères et sœurs,\\[0.5cm]
        Pour votre affection, vos encouragements et votre présence constante à mes côtés.\\[1cm]
        À tous ceux qui me sont chers.
    \end{minipage}
\end{center}

\newpage
% --- Remerciements ---
\chapter*{Remerciements}
\addcontentsline{toc}{chapter}{Remerciements}

\vspace{1cm}
Je tiens à exprimer ma profonde gratitude et mes sincères remerciements à toutes les personnes qui ont contribué, de près ou de loin, à l'aboutissement de ce Projet de Fin d'Études.

\vspace{0.5cm}
Mes remerciements s'adressent tout d'abord à la direction de l'\textbf{École Nationale d'Ingénieurs - Abderhamane Baba Touré (ENI-ABT)}, ainsi qu'à l'ensemble du corps professoral pour la qualité de l'enseignement dispensé durant mon cursus.

\vspace{0.5cm}
Je remercie chaleureusement mon promoteur/encadreur, \textbf{M. [Nom de l'encadreur]}, pour sa disponibilité, ses conseils judicieux, sa rigueur scientifique et la confiance qu'il m'a accordée tout au long de ce projet.

\vspace{0.5cm}
Mes remerciements vont également aux membres du jury qui m'ont fait l'honneur d'évaluer ce travail.

\vspace{0.5cm}
Enfin, je n'oublie pas mes amis et camarades de promotion pour leur entraide, les échanges fructueux et les moments inoubliables partagés durant ces années d'études.

% !TEX root = main.tex

% --- Présentation de l'ENI-ABT ---
\chapter*{Présentation de l'ENI-ABT}
\addcontentsline{toc}{chapter}{Présentation de l'ENI-ABT}

\begin{tcolorbox}[colback=white, colframe=black, sharp corners, halign=center, boxrule=1.5pt]
    \textbf{\Large\textcolor{macouleur}{Présentation de L'ENI-ABT}}
\end{tcolorbox}

\vspace{0.5cm}
Créée le 14 avril 1939 l'\textbf{Ecole Nationale d'Ingénieurs – Abderhamane Baba Touré} (ENI-ABT) a subi plusieurs mutations allant de l'\textbf{Ecole des Travaux Publics} formant des dessinateurs jusqu'à la formation des ingénieurs de conception dans plusieurs domaines. Elle commença comme l'\textbf{Ecole des Travaux Publics} (ETP) de l'A.O.P destinée à former des adjoints techniques en Travaux Publics. La première réforme de l'enseignement érigea l'ETP en \textbf{Ecole Nationale d'Ingénieurs} (ENI) en 1962, en Septembre 1973 le recrutement à l'ENI fut conditionné au Baccalauréat, pour la formation des ingénieurs des sciences appliquées en industrie, génie civil, géologie et en topographie. En 2002 l'ENI, à l'instar des autres grandes écoles (IPR/IFRA, ENSUP), a été directement rattachée à la Direction Nationale de l'Enseignement Supérieur et de la Recherche Scientifique (DNESRS) suivant l'Ordonnance N° 02-054/P-RM du 04 juin 2002. Elle a alors été régie dans son organisation et son fonctionnement par le décret N° 96-378/P-RM du 31 décembre 1996. En juin 2006 l'ENI fut baptisée au nom de son premier Directeur Malien Abderhamane Baba Touré (ENI\textendash ABT).

\vspace{0.5cm}
L'ENI-ABT occupe une place centrale dans la politique industrielle du Mali. En formant des ingénieurs capables de concevoir, de fabriquer et de maintenir des systèmes mécaniques complexes, elle répond directement aux besoins d'une économie en pleine mutation vers l'industrialisation. Le présent Projet de Fin d'Études (PFE) s'inscrit dans cette vision : démontrer qu'il est possible de concevoir et de réaliser localement, à l'ENI-ABT, une machine-outil de précision à 5 axes, en exploitant des technologies de l'Industrie 4.0, sans dépendance aux systèmes propriétaires étrangers. L'ENI-ABT dispose à ce jour d'un atelier mécanique équipé de tours parallèles, de fraiseuses conventionnelles et d'une imprimante 3D FDM, ce qui constitue un socle technique sur lequel ce projet vient greffer une dimension CNC numérique avancée.

\vspace{1cm}
\begin{figure}[htbp]
    \centering
    \caption{Vues des bâtiments de l'ENI-ABT, Bamako, Mali}
    \label{fig:eni_abt}
\end{figure}

\newpage
% --- Structuration pédagogique ---
\section*{\textcolor{macouleur}{Structuration pédagogique}}

Depuis la rentrée 2014-2015, l'ENI-ABT applique le système LMD à travers une nouvelle structuration comportant six DER pour six filières de départ :

\vspace{0.5cm}
\noindent
\begin{minipage}[t]{0.48\textwidth}
    \begin{tcolorbox}[colback=white, colframe=gray, boxrule=1pt, sharp corners, halign=center]
        \textbf{L'ENI-ABT avec le système}\\
        \textbf{«CLASSIQUE»}
    \end{tcolorbox}
    \vspace{0.3cm}
    \textbf{Deux niveaux de formations :}
    \begin{itemize}[label=\raisebox{0.25ex}{\tiny$\blacktriangleright$}, itemsep=2pt]
        \item Cycle Ingénieur ;
        \item Cycle Technicien Supérieur.
    \end{itemize}
    \vspace{0.5cm}
    \textbf{Quatre DER :}
    \begin{itemize}[label=\raisebox{0.25ex}{\tiny$\blacktriangleright$}, itemsep=2pt]
        \item Génie civil ;
        \item Génie Industriel ;
        \item Géodésie ;
        \item Géologie.
    \end{itemize}
\end{minipage}
\hfill
\vrule
\hfill
\begin{minipage}[t]{0.48\textwidth}
    \begin{tcolorbox}[colback=gray!10, colframe=gray, boxrule=1pt, sharp corners, halign=center]
        \textbf{L'ENI-ABT avec le système}\\
        \textbf{«LMD»}
    \end{tcolorbox}
    \vspace{0.3cm}
    \textbf{Trois niveaux de formations :}
    \begin{itemize}[label=\raisebox{0.25ex}{\tiny$\blacktriangleright$}, itemsep=2pt]
        \item Licence ;
        \item Master ;
        \item Doctorat.
    \end{itemize}
    \vspace{0.5cm}
    \textbf{Six DER :}
    \begin{itemize}[label=\raisebox{0.25ex}{\tiny$\blacktriangleright$}, itemsep=2pt]
        \item Génie civil ;
        \item Génie Electrique/Industrie ;
        \item Génie Mécanique et Energie ;
        \item Génie Informatique et Télécommunication ;
        \item Géodésie/Sciences Géographiques.
    \end{itemize}
\end{minipage}

\vspace{1cm}
\noindent L'ENI-ABT mène également des activités :
\begin{itemize}[label=\raisebox{0.25ex}{\tiny$\blacktriangleright$}]
    \item De formation continue diplômante et qualifiante à travers son Service de la Formation Continue (SFC) ;
    \item D'expertise et de production à travers son Unité d'Expertise et de Production (UEP).
\end{itemize}

\newpage
% ===================================================
% --- Sigles et Abréviations ---
\chapter*{Sigles et Abréviations}
\addcontentsline{toc}{chapter}{Sigles et Abréviations}

\renewcommand{\arraystretch}{1.25}
\begin{longtable}{|l|p{10.5cm}|}
\hline
\rowcolor{macouleur!20}
\textbf{Sigle / Abréviation} & \textbf{Signification} \\
\hline
\endfirsthead
\hline
\rowcolor{macouleur!20}
\textbf{Sigle / Abréviation} & \textbf{Signification} \\
\hline
\endhead
\hline
\endfoot
\textbf{AMDEC} & Analyse des Modes de Défaillance, de leurs Effets et de leur Criticité \\
\hline
\textbf{AOP} & Afrique Occidentale Protectorat \\
\hline
\textbf{API} & Application Programming Interface \\
\hline
\textbf{BOM} & Bill of Materials (Nomenclature financière de projet) \\
\hline
\textbf{CAO} & Conception Assistée par Ordinateur \\
\hline
\textbf{CEM} & Compatibilité ÉlectroMagnétique (EMC en anglais) \\
\hline
\textbf{CNC} & Computer Numerical Control (Commande Numérique par Calculateur) \\
\hline
\textbf{D-H} & Denavit-Hartenberg (Paramètres conventionnels de cinématique des robots) \\
\hline
\textbf{DDA} & Digital Differential Analyzer (Algorithme de génération de trajectoire) \\
\hline
\textbf{DER} & Département d'Enseignement et de Recherche \\
\hline
\textbf{DPU} & Data Processing Unit (Unité de Traitement des Données) \\
\hline
\textbf{EEPROM} & Electrically Erasable Programmable Read-Only Memory \\
\hline
\textbf{EMI} & ElectroMagnetic Interference (Perturbation électromagnétique) \\
\hline
\textbf{ENI-ABT} & École Nationale d'Ingénieurs – Abderhamane Baba Touré \\
\hline
\textbf{E-STOP} & Emergency Stop (Arrêt d'Urgence) \\
\hline
\textbf{FAO} & Fabrication Assistée par Ordinateur \\
\hline
\textbf{FAST} & Function Analysis System Technique (Diagramme d'analyse fonctionnelle) \\
\hline
\textbf{FCFA} & Franc de la Communauté Financière Africaine \\
\hline
\textbf{FEA} & Finite Element Analysis (Analyse par Éléments Finis) \\
\hline
\textbf{FFF} & Fused Filament Fabrication (impression 3D par dépôt de matière) \\
\hline
\textbf{FreeRTOS} & Free Real-Time Operating System (OS temps réel open-source) \\
\hline
\textbf{FSM} & Finite State Machine (Machine à États Finis) \\
\hline
\textbf{GPIO} & General Purpose Input/Output (Broches d'entrée/sortie) \\
\hline
\textbf{GRBL} & G-code Real-time Bootstrap Loader (firmware CNC open-source) \\
\hline
\textbf{HB} & Dureté Brinell \\
\hline
\textbf{IDE} & Integrated Development Environment (Environnement de développement) \\
\hline
\textbf{IHM} & Interface Homme-Machine \\
\hline
\textbf{IoT} & Internet of Things (Internet des Objets) \\
\hline
\textbf{ISR} & Interrupt Service Routine (Routine d'interruption matérielle) \\
\hline
\textbf{LDO} & Low Drop-Out (Régulateur de tension linéaire) \\
\hline
\textbf{LUT} & Look-Up Table (Table de correspondance pré-calculée) \\
\hline
\textbf{MISRA} & Motor Industry Software Reliability Association \\
\hline
\textbf{NEMA} & National Electrical Manufacturers Association \\
\hline
\textbf{OEE} & Overall Equipment Effectiveness (Taux de Rendement Synthétique) \\
\hline
\textbf{PCB} & Printed Circuit Board (Circuit Imprimé) \\
\hline
\textbf{PFE} & Projet de Fin d'Études \\
\hline
\textbf{PME} & Petites et Moyennes Entreprises \\
\hline
\textbf{PSU} & Power Supply Unit (Alimentation à Découpage) \\
\hline
\textbf{PWM} & Pulse Width Modulation (Modulation de Largeur d'Impulsion) \\
\hline
\textbf{RDM} & Résistance des Matériaux \\
\hline
\textbf{ROI} & Return on Investment (Retour sur Investissement) \\
\hline
\textbf{RTOS} & Real-Time Operating System \\
\hline
\textbf{SMPS} & Switch-Mode Power Supply (Alimentation à découpage industrielle) \\
\hline
\textbf{SoC} & System on Chip (Système sur puce intégrée) \\
\hline
\textbf{TCP} & Tool Center Point (Point de contact de l'outil) \\
\hline
\textbf{TTL} & Transistor-Transistor Logic (Logique $5\text{V}$) \\
\hline
\textbf{UDP} & User Datagram Protocol (Protocole réseau sans connexion) \\
\hline
\textbf{UML} & Unified Modeling Language \\
\hline
\textbf{USB} & Universal Serial Bus \\
\hline
\textbf{VFD} & Variable Frequency Drive (Variateur de fréquence pour broche) \\
\hline
\textbf{WBS} & Work Breakdown Structure (Structure de découpage du travail) \\
\hline
\end{longtable}
\newpage
% ===================================================
% --- NOTATIONS ET SYMBOLES MATHÉMATIQUES ---
\chapter*{Notations et Symboles Mathématiques}
\addcontentsline{toc}{chapter}{Notations et Symboles Mathématiques}

Le tableau suivant recense les principaux symboles mathématiques utilisés dans les chapitres de calcul, accompagnés de leur signification physique et de leur unité dans le Système International (SI).

\vspace{0.3cm}
\renewcommand{\arraystretch}{1.3}
\begin{center}
\begin{tabular}{|c|p{8cm}|c|}
\hline
\rowcolor{macouleur!20}
\textbf{Symbole} & \textbf{Signification} & \textbf{Unité} \\
\hline
$F_c$ & Force de coupe tangentielle & N \\
\hline
$F_f$ & Force d'avance (axiale) & N \\
\hline
$F_p$ & Force de pénétration (radiale) & N \\
\hline
$F_{ax}$ & Force axiale sur la vis à billes & N \\
\hline
$\vec{R}$ & Résultante des forces de coupe (torseur) & N \\
\hline
$M_f$ & Moment de flexion & N·mm \\
\hline
$M_{renversement}$ & Moment de renversement du chariot Y (Table Trunnion) & N·m \\
\hline
$\sigma_{eq}$ & Contrainte équivalente de Von Mises & MPa \\
\hline
$\sigma_{max}$ & Contrainte normale maximale en flexion & MPa \\
\hline
$\tau$ & Contrainte de cisaillement & MPa \\
\hline
$R_e$ & Limite d'élasticité du matériau & MPa \\
\hline
$R_m$ & Résistance maximale à la rupture & MPa \\
\hline
$E$ & Module de Young (rigidité) & MPa \\
\hline
$I_{yy}$ & Moment quadratique de la section transversale & mm$^4$ \\
\hline
$f_{max}$ & Flèche maximale (déflexion statique) & mm \\
\hline
$F_k$ & Charge critique de flambage (critère d'Euler) & N \\
\hline
$k_L$ & Facteur d'appui selon le montage de la vis & --- \\
\hline
$s$ & Coefficient de sécurité statique & --- \\
\hline
$C_r$ & Couple résistant continu de la vis à billes & N·m \\
\hline
$C_J$ & Couple d'inertie lors des accélérations & N·m \\
\hline
$P_h$ & Pas de la vis à billes & mm/tour \\
\hline
$\eta$ & Rendement de la transmission par vis à billes & \% \\
\hline
$\Delta\theta$ & Erreur angulaire du moteur (Backlash) & ° ou rad \\
\hline
$\Delta L$ & Erreur linéaire tangentielle (effet levier Trunnion) & mm \\
\hline
$R$ & Rayon du plateau Trunnion & mm \\
\hline
$H$ & Hauteur du berceau Trunnion & mm \\
\hline
$L_y, L_z$ & Offsets géométriques de Denavit-Hartenberg & mm \\
\hline
$\theta_A, \theta_B$ & Angles des articulations rotatives Axe A et Axe B & ° \\
\hline
$\,^0T_n$ & Matrice de transformation homogène globale (produit D-H) & --- \\
\hline
$f_{step}$ & Fréquence des impulsions STEP vers les drivers & Hz \\
\hline
$k_c$ & Pression spécifique de coupe du matériau & N/mm$^2$ \\
\hline
$V_c$ & Vitesse de coupe outil & m/min \\
\hline
$f_z$ & Avance par dent de la fraise & mm/dent \\
\hline
$a_p$ & Profondeur de passe axiale & mm \\
\hline
$a_e$ & Profondeur de passe radiale (engagement) & mm \\
\hline
$D_c$ & Diamètre de l'outil coupant (fraise) & mm \\
\hline
$P_{PSU}$ & Puissance nominale de l'alimentation à découpage & W \\
\hline
$C_{bulk}$ & Condensateur de découplage (anti Brown-out) & F (µF) \\
\hline
$L_{10}$ & Durée de vie nominale des roulements (ISO 281) & heures (h) \\
\hline
$C_{dyn}$ & Charge dynamique de base des roulements & kN \\
\hline
$P_{equiv}$ & Charge équivalente sur le roulement & N \\
\hline
$Ra$ & Rugosité moyenne arithmétique de surface & µm \\
\hline
$C_{pk}$ & Indice de capabilité machine (ISO 22514) & --- \\
\hline
\end{tabular}
\end{center}


\newpage
\tableofcontents % Sommaire automatique

\newpage
\renewcommand{\listfigurename}{Liste des figures}
\listoffigures
\addcontentsline{toc}{chapter}{Liste des figures}

\newpage
\renewcommand{\listtablename}{Liste des tableaux}
\listoftables
\addcontentsline{toc}{chapter}{Liste des tableaux}

\newpage
% !TEX root = ../main.tex
\chapter*{Introduction Générale}
\addcontentsline{toc}{chapter}{Introduction Générale}
\label{chap:intro}

\section*{Contexte et Justification du Projet}
L'industrie manufacturière mondiale traverse une phase de mutation sans précédent, portée par des exigences de précision millimétrique, de flexibilité accrue et de réduction des temps de cycle. Au centre de cette dynamique, la Machine-Outil à Commande Numérique (MOCN) s'est imposée comme l'équipement indispensable à la production moderne. Si les architectures à trois axes ont longtemps dominé le marché pour la majorité des opérations d'usinage, l'émergence de besoins de plus en plus complexes - notamment dans les secteurs de l'aéronautique, du médical, et de la mécanique de précision - a propulsé la technologie d'usinage 5 axes au premier plan.

Cependant, aborder la conception d'une CNC 5 axes constitue un défi d'intégration mécatronique majeur. Il ne s'agit plus seulement de contrôler des translations linéaires, mais d'assurer la synchronisation parfaite de cinq degrés de liberté (trois translations et deux rotations) via un logiciel de commande capable de traiter des calculs de cinématique inverse en temps réel.

\section*{Évolution du Concept d'Industrie et Enjeux Locaux}
Le terme « industrie » a connu plusieurs ruptures technologiques. Comprendre ces évolutions permet de situer la pertinence de ce Projet de Fin d'Études (PFE) :
\begin{itemize}
    \item \textbf{L'Industrie 1.0 et 2.0 :} Fondées sur l'énergie thermique puis électrique, elles ont permis la naissance de la machine-outil, mais celle-ci restait asservie à la main de l'homme.
    \item \textbf{L'Industrie 3.0 :} Elle a introduit l'informatique de gestion et l'automatisation. C'est l'époque des CNC "boîtes noires" où l'utilisateur est totalement dépendant des constructeurs de contrôleurs propriétaires.
    \item \textbf{L'Industrie 4.0 :} Elle marque l'avènement des systèmes cyber-physiques. La machine devient un nœud communicant. C'est dans ce paradigme que s'inscrit notre projet : créer une machine intelligente, dont le code source est ouvert, adaptable et pilotable via des interfaces modernes (Wi-Fi, IHM sur tablette).
\end{itemize}

L'usinage multiaxes est le bras armé de cette nouvelle industrie. En permettant l'orientation spatiale de l'outil, la CNC 5 axes supprime de multiples montages/démontages de la pièce, garantissant une meilleure précision géométrique (tolérances de forme restrictives).

\begin{figure}[htbp]
    \centering
    % Placeholder pour une image de CNC 5 axes
    % Remplacer "placeholder_cnc5axes.png" par l'image réelle
%    \includegraphics[width=0.7\textwidth, draft]{placeholder_cnc5axes.png} 
    \caption{Concept d'usinage sur un centre CNC 5 axes (Table Trunnion)}
    \label{fig:concept_5axes}
\end{figure}

\section*{Problématique}
L'acquisition de centres d'usinage 5 axes reste un luxe technologique, souvent inaccessible pour beaucoup de structures de formation, de FabLabs ou de petites unités de production, particulièrement dans les pays en développement comme le Mali. Cette barrière d'entrée n'est pas seulement financière (coût d'achat et d'entretien des machines industrielles), mais aussi logicielle. 
Les systèmes propriétaires (Fanuc, Heidenhain, Siemens) interdisent toute modification algorithmique de bas niveau, limitant l'apprentissage profond de la mécatronique de commande pour les étudiants et ingénieurs.

La problématique qui sous-tend ce mémoire est donc la suivante : \textbf{Comment concevoir et réaliser localement, à moindre coût, une mini-fraiseuse CNC 5 axes et développer l'intégralité de son environnement logiciel (de la cinématique inverse sur microcontrôleur 32 bits jusqu'à l'Interface Homme-Machine multiplateforme) pour obtenir un système d'usinage didactique, performant et ouvert ?}

\section*{Souveraineté Technologique et Enjeux au Mali}
Dans le contexte de la transition industrielle du Mali, l'acquisition de technologies de pointe est le pivot de la croissance économique. Ce projet dépasse le cadre d'un simple exercice académique pour s'inscrire dans une vision de \textbf{Souveraineté Technologique}. L'École Nationale d'Ingénieurs Abderhamane Baba Touré (ENI-ABT) a pour mission de former des cadres capables de s'affranchir de la dépendance externe.

La maîtrise locale du code source et de la conception mécanique d'un système 5 axes permet :
\begin{itemize}
    \item \textbf{La Maintenance Autonome :} Les pannes sur les CNC industrielles étrangères paralysent souvent les entreprises maliennes faute de pièces ou de techniciens agréés. Une machine "open-design" est réparable localement.
    \item \textbf{L'Industrialisation des PME :} Permettre à de petits ateliers de fabriquer des pièces de haute précision (prothèses médicales, pièces aéronautiques de rechange, moules complexes) sans investissements de plusieurs centaines de millions de FCFA.
    \item \textbf{L'Innovation Continue :} La possibilité de modifier les algorithmes de la machine pour l'adapter à des besoins spécifiques (impression 3D béton, découpe laser multiaxes, etc.).
\end{itemize}
C’est une contribution directe à l'édifice d'une industrie nationale résiliente et auto-suffisante.

\section*{Objectifs du Projet}
L'objectif principal est la réalisation de A à Z d'un prototype fonctionnel de CNC 5 axes piloté par un logiciel propriétaire (développé "from scratch"). Les objectifs spécifiques se déclinent comme suit :
\begin{enumerate}
    \item \textbf{Conception Mécanique :} Dimensionnement d'une structure rigide de type "Table basculante et rotative" (Trunnion Table) permettant d'abriter les axes A et B sans perte de précision.
    \item \textbf{Architecture Électronique :} Intégration d'un microcontrôleur hautement performant (ESP32 à double cœur) capable de soutenir les calculs matriciels en temps réel, couplé à des composants fiables (Vis à billes, Moteurs pas-à-pas NEMA de précision).
    \item \textbf{Développement Logiciel Embarqué (Firmware) :} Implémentation algorithmique en C/C++ de la cinématique inverse et de la génération d'impulsions à très haute fréquence via RTOS.
    \item \textbf{Développement Logiciel IHM :} Création d'une Interface de commande (Sender) en langage Dart/Flutter, permettant une communication fluide, asynchrone (JSON) et multiplateforme (Windows/Android) avec la machine.
\end{enumerate}

\section*{Planification du Projet}
Afin de mener à bien ce projet ambitieux, une planification rigoureuse a été établie. Le diagramme de Gantt ci-dessous illustre les différentes phases clés de la réalisation : de l'étude bibliographique à la validation finale des performances d'usinage.

\begin{figure}[htbp]
    \centering
    \resizebox{\textwidth}{!}{
    \begin{ganttchart}[
        y unit title=0.6cm,
        y unit chart=0.7cm,
        x unit=0.8cm,
        vgrid,
        time slot format=isodate-yearmonth,
        time slot unit=month,
        title/.append style={fill=blue!10},
        title label font=\bfseries\footnotesize,
        bar/.append style={fill=blue!50, rounded corners=2pt},
        bar label font=\footnotesize,
        group/.append style={fill=black},
        group label font=\bfseries\footnotesize
    ]{2023-09}{2024-06}
        \gantttitlecalendar{year, month=shortname} \\
        \ganttgroup{Phase 1: Étude \& État de l'art}{2023-09}{2023-11} \\
        \ganttbar{Recherches bibliographiques}{2023-09}{2023-10} \\
        \ganttbar{Choix de l'architecture}{2023-10}{2023-11} \\
        \ganttgroup{Phase 2: Conception (CAO \& Ajustements)}{2023-11}{2024-02} \\
        \ganttbar{Modélisation Trunnion \& Gantry}{2023-11}{2024-01} \\
        \ganttbar{Dimensionnement mécanique}{2024-01}{2024-02} \\
        \ganttgroup{Phase 3: Développement Logiciel}{2023-12}{2024-04} \\
        \ganttbar{Firmware ESP32 (Cinématique)}{2023-12}{2024-03} \\
        \ganttbar{Interface IHM (Flutter)}{2024-02}{2024-04} \\
        \ganttgroup{Phase 4: Réalisation \& Tests}{2024-03}{2024-06} \\
        \ganttbar{Câblage électrique}{2024-03}{2024-04} \\
        \ganttbar{Assemblage mécanique}{2024-04}{2024-05} \\
        \ganttbar{Tests d'usinage \& Calibrations}{2024-05}{2024-06}
    \end{ganttchart}
    }
    \caption{Diagramme de Gantt prévisionnel des grandes phases du projet PFE}
    \label{fig:gantt_projet}
\end{figure}

\section*{Structure du Mémoire}
Pour exposer rigoureusement notre démarche, ce mémoire est structuré en plusieurs chapitres consécutifs :
\begin{itemize}
    \item \textbf{Le Chapitre 1} présente la revue de la littérature sur la technologie d'usinage 5 axes et l'état de l'art des systèmes de commande cyber-physiques.
    \item \textbf{Le Chapitre 2} détaille la conception et les lourds calculs de dimensionnement mécanique (structure, transmission, RDM).
    \item \textbf{Le Chapitre 3} se focalise sur l'architecture électronique et le développement logiciel complet sous Flutter et la cinématique inverse sous ESP32.
    \item \textbf{Le Chapitre 4}, enfin, présente la réalisation pratique, le paramétrage, les tests d'usinage et la validation des performances de notre mini-fraiseuse.
\end{itemize}            % Introduction et Contexte du Projet
% !TEX root = ../main.tex
\chapter{État de l'art et Revue de la Littérature}
\label{chap:etat_art}

\section{Généralités sur la Commande Numérique par Calculateur (CNC)}
La Commande Numérique par Calculateur (CNC) se définit comme le pilotage automatisé des mouvements et des fonctions auxiliaires d'une machine-outil par microprocesseur, via l'interprétation d'un programme d'instructions textuelles (G-Code). Depuis les travaux pionniers de John T. Parsons (qui a utilisé des cartes perforées pour usiner des pales d'hélicoptère dans les années 1940), la CNC a évolué de systèmes câblés immenses et rigides vers des architectures logicielles distribuées, flexibles et miniaturisées.

\subsection{Évolution historique des systèmes CNC}
L'histoire de la commande numérique est intrinsèquement liée à l'évolution de l'informatique industrielle. Le tableau~\ref{tab:histoire_cnc} retrace les grandes ruptures technologiques qui ont conduit aux systèmes modernes ouverts.

\begin{table}[htbp]
\centering
\renewcommand{\arraystretch}{1.3}
\begin{tabular}{|c|p{2.5cm}|p{6cm}|p{4.5cm}|}
\hline
\rowcolor{macouleur!20}
\textbf{Décennie} & \textbf{Technologie} & \textbf{Innovation principale} & \textbf{Limites identifiées} \\
\hline
1940--1950 & Cartes perforées & Travaux de Parsons / MIT : premier usinage de pale d'hélicoptère & Pas de rétroaction, séquentiel pur \\
\hline
1960--1970 & NC câblé & Contrôleurs à relais, tubes à vide & Rigides, non-reprogrammables \\
\hline
1970--1980 & CNC 8 bits & Microprocesseur Intel 8080 / Z80, DNC & Un seul cœur, pas de multi-axes \\
\hline
1980--1990 & Propriétaire & FANUC, Siemens, Heidenhain — OS fermés & Coût prohibitif, dépendance \\
\hline
1990--2010 & PC-Based (Linux) & Interface PC + carte FPGA (Mesa Cards) & Configuration complexe, coût \\
\hline
2015--2024 & RTOS 32-bits open & ESP32, STM32 + FreeRTOS, GRBL-ESP32 & Notre domaine d'innovation \\
\hline
\end{tabular}
\caption{Évolution historique des architectures de commande numérique}
\label{tab:histoire_cnc}
\end{table}

\subsection{Principe d'asservissement d'un système CNC}
Un système CNC moderne s'articule autour d'une chaîne de commande comportant trois unités principales agissant en synergie :
\begin{enumerate}
    \item \textbf{L'Unité de Traitement des Données (Data Processing Unit -- DPU) :} Elle correspond au logiciel IHM. Elle lit le fichier de commande numérique, l'analyse, simule la trajectoire et l'envoie de manière séquencée.
    \item \textbf{L'Unité de Contrôle (Control Loop Unit -- CLU) :} C'est le « cerveau interne » (le microcontrôleur sur la machine). Il reçoit les données du DPU, calcule la cinématique, gère les accélérations (planner) et génère les signaux électriques logiques (Step/Direction ou signaux PWM).
    \item \textbf{L'Unité de Puissance et d'Actionnement :} Composée de drivers (variateurs) qui amplifient les signaux logiques de la CLU, et de moteurs électriques transmettant les couples et mouvements aux axes physiques.
\end{enumerate}

\section{Classification Architecturale des Machines CNC 5 Axes}
Le passage d'une machine 3 axes à une machine 5 axes introduit une complexité majeure dans la chaîne cinématique. Pour qu'une fraise puisse atteindre n'importe quel point d'une pièce sous n'importe quel angle, deux axes de rotation (généralement nommés A, B ou C) doivent s'ajouter aux trois translations (X, Y, Z). Il existe trois grandes familles d'architectures pour les machines 5 axes, classées selon l'emplacement de ces deux axes rotatifs.

\begin{figure}[htbp]
    \centering
    \includegraphics[width=0.6\textwidth]{images/trunnion_diagram.png}
    \caption{Schéma cinématique d'une machine CNC 5 axes de type Trunnion (Table-Table)}
    \label{fig:trunnion_schema}
\end{figure}

\begin{figure}[htbp]
    \centering
    \begin{tikzpicture}[
        node distance=2.5cm and 1.8cm,
        box/.style={rectangle, draw, rounded corners=4pt, align=center, fill=blue!5, minimum width=3.8cm, minimum height=1.8cm, thick},
        title/.style={font=\bfseries\small},
        desc/.style={font=\footnotesize, text width=3.8cm, align=center},
        arrow/.style={-{Stealth[scale=1.2]}, thick, color=gray!80}
    ]
        \node[box] (table_table) {
            \textbf{Table-Table}\\(Trunnion)
        };
        \node[below=0.2cm of table_table, desc, align=center] {Rotations A et C sur la table.\\ Meilleure rigidité.};
        
        \node[box, right=1cm of table_table] (head_head) {
            \textbf{Head-Head}\\(Tête Pivotante)
        };
        \node[below=0.2cm of head_head, desc, align=center] {Rotations A et C sur la broche.\\ Pièces volumineuses.};
        
        \node[box, right=1.5cm of head_head] (head_table) {
            \textbf{Head-Table}\\(Hybride)
        };
        \node[below=0.2cm of head_table, desc] {Répartition des rotations\\(mixte Tête / Table).};
        
        \draw[arrow] (table_table.east) -- (head_head.west) node[midway, above, font=\scriptsize, color=black] {Cinématique};
        \draw[arrow] (head_head.east) -- (head_table.west) node[midway, above, font=\scriptsize, color=black] {Compromis};
    \end{tikzpicture}
    \caption{Comparaison schématique des topologies CNC 5 axes}
    \label{fig:topo_5axes}
\end{figure}

\subsection{Configuration «Table-Table» (Trunnion Table)}
Dans cette architecture, les deux axes rotatifs (souvent l'axe A pour le basculement et l'axe B pour la rotation continue) sont emboîtés l'un dans l'autre et supportent la table de travail.
\begin{itemize}
    \item \textbf{Cinématique :} L'outil se déplace linéairement en X, Y et Z. C'est la pièce à usiner qui tourne (Pan) et bascule (Tilt) face à l'outil.
    \item \textbf{Avantages :} Meilleure rigidité globale et excellente précision sur les petites pièces compactes.
    \item \textbf{Limites :} Le volume de la pièce et son poids sont strictement limités par le diamètre du plateau rotatif.
\end{itemize}
C'est l'architecture que nous avons retenue pour notre projet didactique.

\subsection{Configuration «Head-Head» (Tête Pivotante Bi-Rotative)}
Ici, la table supportant la pièce est fixe. Les deux axes rotatifs sont intégrés dans la tête qui supporte la broche.
\begin{itemize}
    \item \textbf{Avantages :} Parfaitement adaptée à l'usinage de pièces de très grandes dimensions et très lourdes.
    \item \textbf{Limites :} Rigidité angulaire diminuée. Le bras de levier important amplifie considérablement la moindre erreur angulaire.
\end{itemize}

\subsection{Configuration «Head-Table» (Hybride)}
Cette architecture répartit les rotations. Généralement, l'axe de basculement est fixé sur la tête de la broche, et l'axe de rotation continue est placé sur une table tournante. Cette solution hybride permet un compromis intéressant pour les machines de taille moyenne.

\section{Comparatif des Machines CNC 5 Axes Accessibles}
Afin de situer notre projet dans le paysage industriel actuel, le tableau~\ref{tab:comparatif_cnc} compare les caractéristiques de notre prototype avec celles de machines de référence.

\begin{table}[htbp]
\centering
\renewcommand{\arraystretch}{1.3}
\begin{tabular}{|p{2.5cm}|p{2cm}|p{2cm}|p{2.5cm}|p{2.5cm}|}
\hline
\rowcolor{macouleur!20}
\textbf{Critère} & \textbf{PFE (Nous)} & \textbf{PocketNC} & \textbf{Haas UMC-500} & \textbf{Candle+GRBL} \\
\hline
\textbf{Architecture} & Trunnion A+B & Trunnion A+B & Trunnion A+C & 3 axes \\
\hline
\textbf{Volume} & 20cm cube & 12cm cube & 50cm cube & Variable \\
\hline
\textbf{Contrôleur} & ESP32 RTOS & Propriétaire & Haas CNC & Arduino \\
\hline
\textbf{Précision} & $\pm 0{,}05$ mm & $\pm 0{,}025$ mm & $\pm 0{,}005$ mm & $\pm 0{,}1$ mm \\
\hline
\textbf{Open-Source} & Oui (100\%) & Non & Non & Oui (firm) \\
\hline
\textbf{Coût} & $\approx 0{,}5 M$ & $\approx 7 M$ & $\approx 300 M$ & $\approx 0{,}1 M$ \\
\hline
\textbf{IHM} & Flutter App & Kinetic & Tactile HAAS & PC Candle \\
\hline
\end{tabular}
\caption{Tableau comparatif des machines CNC 5 axes du marché vs. notre prototype}
\label{tab:comparatif_cnc}
\end{table}

L'analyse de ce tableau démontre que notre prototype, bien qu'inférieur en précision et en volume aux machines industrielles, se positionne comme la \textbf{seule solution entièrement open-source} dotée d'une IHM multi-plateforme moderne (Flutter) et d'une connectivité sans fil. Le rapport \textbf{performance/coût} est sans équivalent dans le segment éducatif africain.

\section{Problématiques Avancées : Singularités et TCP}
L'usinage 5 axes introduit des défis théoriques absents des machines cartésiennes classiques : la gestion des singularités et le contrôle du point de contact outil.

\subsection{Les Singularités Cinématiques (Gimbal Lock)}
Une singularité cinématique survient lorsque deux axes de rotation s'alignent, entraînant la perte d'un degré de liberté. Pour une table Trunnion (Axe A et B), cela se produit généralement lorsque l'Axe B est vertical et que l'angle de l'Axe A atteint une valeur critique (souvent $0^\circ$).
À ce point précis, une infinité de combinaisons d'angles peut donner la même position spatiale, ce qui peut provoquer des accélérations infinies théoriques des moteurs et des marques de vibration sur la pièce (\textit{chatter}). Le logiciel de contrôle doit intégrer des algorithmes d'évitement ou de lissage de trajectoire aux abords de ces zones critiques.

Il existe trois types principaux de singularités rencontrées en usinage 5 axes :
\begin{enumerate}
    \item \textbf{Singularité de configuration (Gimbal Lock) :} Perte d'un degré de liberté par alignement des axes de rotation. Se manifeste par des vitesses angulaires théoriquement infinies.
    \item \textbf{Singularité de frontière :} Le vecteur outil atteint la limite physique d'une articulation (débattement angulaire maximal). Un interverrouillage logiciel est nécessaire.
    \item \textbf{Singularité de chemin :} Le parcours CAO génère une trajectoire passant exactement par un point singulier. Le post-processeur FAO doit détecter ces cas et déranger légèrement la trajectoire.
\end{enumerate}

\subsection{Le Concept de TCP (Tool Center Point Control)}
Dans un système 3 axes, la position de l'outil est directement liée aux coordonnées X, Y, Z. En 5 axes, toute rotation de la table (Axe A ou B) déplace physiquement la pièce par rapport à la pointe de l'outil.
Le \textbf{TCP Control} est la capacité du contrôleur à compenser automatiquement ces décalages géométriques. Au lieu de programmer les positions des moteurs, le programmeur définit la trajectoire de la \textbf{pointe de l'outil} par rapport à la pièce. C'est l'intelligence embarquée (le microcontrôleur) qui recalcule en temps réel les compensations linéaires $\Delta X, \Delta Y, \Delta Z$ nécessaires pour maintenir l'outil sur le point de contact nominal avec la surface.

\section{L'Écosystème Logiciel et le Langage G-Code 5 Axes}

\subsection{La Chaîne Numérique 5 Axes}
Le flux de données d'une pièce usinée numériquement passe par trois phases :
\begin{itemize}
    \item \textbf{CAO (Conception Assistée par Ordinateur) :} Le concepteur crée le modèle 3D de la pièce.
    \item \textbf{FAO (Fabrication Assistée par Ordinateur) :} Le programmeur génère les parcours d'outil. En 5 axes, on distingue deux méthodes : le \textit{3+2 axes} (Usinage Positionnel) et le \textit{5 axes continus} (simultané).
    \item \textbf{Le Post-Processeur :} Il traduit la géométrie (vecteurs de l'outil I, J, K) en positions angulaires réelles (A, B) des moteurs de la machine spécifique.
\end{itemize}

\subsection{Le Langage G-Code (Norme ISO 6983)}
C'est le langage universel de la commande numérique. Un programme est une suite alphabétique de blocs d'instructions. En multi-axes, un bloc ressemble à ceci :
\texttt{G01 X15.5 Y0.0 Z-5.0 A45.0 B90.0 F800}

Le tableau~\ref{tab:gcode_5axes} présente les codes G et M principaux utilisés dans notre logiciel :

\begin{table}[htbp]
\centering
\renewcommand{\arraystretch}{1.3}
\begin{tabular}{|l|p{6cm}|l|}
\hline
\rowcolor{macouleur!20}
\textbf{Code} & \textbf{Fonction} & \textbf{Utilisation} \\
\hline
G00 & Déplacement rapide (Rapid) en XYZ & Positionnement \\
\hline
G01 & Interpolation linéaire à vitesse contrôlée & Usinage rectiligne \\
\hline
G02 & Interpolation circulaire sens horaire & Congé, rayon \\
\hline
G03 & Interpolation circulaire sens anti-horaire & Congé, rayon \\
\hline
G17/18/19 & Sélection du plan de travail (XY, XZ, YZ) & Fraisage de poches \\
\hline
G28 & Retour à l'origine machine (Homing) & Initialisation \\
\hline
G54--G59 & Systèmes de coordonnées pièce & Multi-fixtures \\
\hline
G91 & Mode de programmation incrémentale & Jog relatif \\
\hline
G92 & Définition de l'origine pièce courante & Palpage TCP \\
\hline
M03 / M04 & Mise en marche de la broche (sens direct / inverse) & Fraisage \\
\hline
M05 & Arrêt de la broche & Fin pass \\
\hline
M08 / M09 & Arrosage ON / OFF & Refroidissement \\
\hline
M30 & Fin de programme + retour au début & Fin cycle \\
\hline
\$H & Homing ESP32 (FluidNC/GRBL étendu) & Calibration \\
\hline
\$X & Déverrouillage de l'état ALARM & Reprise \\
\hline
\end{tabular}
\caption{Tableau des principaux codes G et M supportés par notre contrôleur ESP32}
\label{tab:gcode_5axes}
\end{table}

\section{Analyse Comparative des Protocoles de Communication Industrielle}

Le choix du protocole de communication entre l'IHM Flutter et le microcontrôleur ESP32 est central pour les performances temps réel du système. Le tableau~\ref{tab:protocoles} compare les principales solutions envisageables.

\begin{table}[htbp]
\centering
\renewcommand{\arraystretch}{1.3}
\resizebox{\textwidth}{!}{
\begin{tabular}{|l|c|c|c|c|p{4cm}|}
\hline
\rowcolor{macouleur!20}
\textbf{Protocole} & \textbf{Latence} & \textbf{Débit max} & \textbf{Connectivité} & \textbf{Overhead} & \textbf{Adéquation CNC} \\
\hline
\textbf{USB-Serial (GRBL)} & $< 10$ ms & 115 200 baud & Filaire & Faible & Limité (câble obligatoire) \\
\hline
\textbf{TCP (WebSocket)} & $5$--$50$ ms & 100 Mbps & Wi-Fi/Eth & Moyen & Bon (handshake) \\
\hline
\textbf{UDP (choix retenu)} & $1$--$5$ ms & 100 Mbps & Wi-Fi & Très faible & \textbf{Optimal} (pas d'ACK) \\
\hline
\textbf{Modbus RTU} & $10$--$100$ ms & 1 Mbps & RS-485 & Moyen & Industriel, trop lent \\
\hline
\textbf{EtherCAT} & $< 0{,}1$ ms & 100 Mbps & Ethernet & Faible & Professionnel, complexe \\
\hline
\textbf{Profibus} & $1$--$10$ ms & 12 Mbps & RS-485 & Élevé & Obsolète, propriétaire \\
\hline
\end{tabular}
}
\caption{Tableau comparatif des protocoles de communication pour systèmes CNC}
\label{tab:protocoles}
\end{table}

\textbf{Justification du choix UDP :} La bande passante théorique nécessaire pour notre application est calculée comme suit. En supposant un taux de rafraîchissement de $f_{refresh} = 50$ trames/seconde et une taille de trame JSON compressée de $T_{trame} = 128$ octets :
\begin{equation}
    BW_{required} = N_{trames} \times T_{trame} \times 8 = 50 \times 128 \times 8 = \boxed{51~200 \text{ bps} \approx 51 \text{ kbps}}
\end{equation}
Ce débit est négligeable face aux 100 Mbps disponibles en Wi-Fi 802.11n. L'UDP est donc choisi pour sa latence minimale ($< 2$ ms) et l'absence d'overhead d'accusé de réception (ACK), critique lors de l'envoi en rafale des blocs G-Code.

\section{État de l'Art des Firmwares CNC Open-Source}

L'intelligence de la machine réside dans son micrologiciel (firmware). Pour opérer une machine 5 axes open-source, plusieurs grands firmwares existent. Le tableau~\ref{tab:firmwares} les compare.

\begin{table}[htbp]
\centering
\renewcommand{\arraystretch}{1.3}
\resizebox{\textwidth}{!}{
\begin{tabular}{|l|c|c|c|p{5cm}|}
\hline
\rowcolor{macouleur!20}
\textbf{Firmware} & \textbf{Contrôleur} & \textbf{Axes max} & \textbf{5 axes continus} & \textbf{Limitation principale} \\
\hline
\textbf{GRBL (classique)} & Arduino Uno (AVR 8-bit) & 3 & Non & Insufficient Timers. Inadapté. \\
\hline
\textbf{Marlin 2.x} & AVR 32-bit / ARM & 6 & Limité & Conçu pour impression 3D, pas pour fraisage. \\
\hline
\textbf{Klipper} & Raspberry Pi + MCU & 6 & Partiel & Complexe, nécessite Linux et Python. \\
\hline
\textbf{FluidNC / GRBL-ESP32} & ESP32 & 6 & Oui (avec plugin) & Configuration YAML ; pas de cinématique Trunnion native. \\
\hline
\textbf{LinuxCNC} & PC + FPGA Mesa & Illimité & Oui & Très complexe, onéreux, nécessite un PC dédié sous Linux RT. \\
\hline
\textbf{Notre firmware (Custom)} & ESP32 Dual-Core & 5 & Oui & Spécifique à l'architecture Trunnion de ce PFE (avantage). \\
\hline
\end{tabular}
}
\caption{Comparatif des firmwares CNC open-source pour machine 5 axes}
\label{tab:firmwares}
\end{table}

Ce tableau confirme la nécessité d'un développement ``from scratch'' du firmware ESP32. Aucun des firmwares existants ne gère nativement la cinématique inverse d'une table Trunnion (axes A+B emboîtés) avec compensation TCP en temps réel.

\section{Matériaux Usinables et Usinabilité}

Le choix des matériaux d'essai est guidé par leur usinabilité et leur pertinence industrielle dans le contexte malien. Le tableau~\ref{tab:usinabilite} présente les paramètres de coupe recommandés pour notre broche de 500\,W et notre fraise de $6\,\text{mm}$.

\begin{table}[htbp]
\centering
\renewcommand{\arraystretch}{1.3}
\resizebox{\textwidth}{!}{
\begin{tabular}{|l|c|c|c|c|c|c|}
\hline
\rowcolor{macouleur!20}
\textbf{Matériau} & \textbf{Dureté (HB)} & \textbf{$k_c$ (N/mm²)} & \textbf{$V_c$ (m/min)} & \textbf{$f_z$ (mm/dent)} & \textbf{$a_p$ max (mm)} & \textbf{Indice Usinab.} \\
\hline
Résine Epoxy (Cibatool) & 15--25 & 200 & 300--500 & 0,10 & 3,0 & Très Facile \\
\hline
Bois dur (Chêne, Iroko) & 25--35 & 350 & 200--400 & 0,12 & 4,0 & Facile \\
\hline
HDPE (Polyéthylène HD) & 60--70 & 400 & 200--350 & 0,08 & 3,0 & Facile \\
\hline
Aluminium AW-2017A & 95--105 & 700 & 100--200 & 0,05 & 1,0 & Moyen \\
\hline
Aluminium AW-5052 & 60 & 600 & 150--250 & 0,06 & 1,5 & Moyen \\
\hline
Laiton CuZn37 & 70--80 & 800 & 80--150 & 0,04 & 0,8 & Moyen \\
\hline
Cuivre pur & 45--55 & 600 & 80--120 & 0,05 & 0,8 & Difficile (adhérent) \\
\hline
Acier doux (S235) & 120 & 1\,500 & 20--50 & 0,02 & 0,3 & Hors capacité broche \\
\hline
\end{tabular}
}
\caption{Paramètres d'usinabilité et vitesses de coupe pour les matériaux testables avec notre CNC}
\label{tab:usinabilite}
\end{table}

\section{Conclusion du chapitre}

Ce chapitre a établi le socle théorique et comparatif sur lequel repose notre démarche de conception. Le passage de 3 à 5 axes exige une refonte totale de la géométrie de la machine (table Trunnion), des solutions de transmissions strictes (Vis à billes), et surtout, d'une rupture architecturale en matière de puissance de calcul embarquée (ESP32 Dual-Core). L'analyse comparative des protocoles de communication, des firmwares open-source et des matériaux usinables a permis de valider et justifier chaque choix technologique majeur du projet. Le chapitre suivant s'attachera à formaliser précisément ces besoins et à déterminer les calculs de sollicitations mécaniques de la future CNC 5 axes.
       % État de l'art et Revue de Littérature
\checkmark$\checkmark^\checkmark\\checkmarkc\checkmarki\checkmarkr\checkmarkc\checkmark$\checkmark%\checkmark$\checkmark^\checkmark\\checkmarkc\checkmarki\checkmarkr\checkmarkc\checkmark$\checkmark \checkmark$\checkmark^\checkmark\\checkmarkc\checkmarki\checkmarkr\checkmarkc\checkmark$\checkmark!\checkmark$\checkmark^\checkmark\\checkmarkc\checkmarki\checkmarkr\checkmarkc\checkmark$\checkmarkT\checkmark$\checkmark^\checkmark\\checkmarkc\checkmarki\checkmarkr\checkmarkc\checkmark$\checkmarkE\checkmark$\checkmark^\checkmark\\checkmarkc\checkmarki\checkmarkr\checkmarkc\checkmark$\checkmarkX\checkmark$\checkmark^\checkmark\\checkmarkc\checkmarki\checkmarkr\checkmarkc\checkmark$\checkmark \checkmark$\checkmark^\checkmark\\checkmarkc\checkmarki\checkmarkr\checkmarkc\checkmark$\checkmarkr\checkmark$\checkmark^\checkmark\\checkmarkc\checkmarki\checkmarkr\checkmarkc\checkmark$\checkmarko\checkmark$\checkmark^\checkmark\\checkmarkc\checkmarki\checkmarkr\checkmarkc\checkmark$\checkmarko\checkmark$\checkmark^\checkmark\\checkmarkc\checkmarki\checkmarkr\checkmarkc\checkmark$\checkmarkt\checkmark$\checkmark^\checkmark\\checkmarkc\checkmarki\checkmarkr\checkmarkc\checkmark$\checkmark \checkmark$\checkmark^\checkmark\\checkmarkc\checkmarki\checkmarkr\checkmarkc\checkmark$\checkmark=\checkmark$\checkmark^\checkmark\\checkmarkc\checkmarki\checkmarkr\checkmarkc\checkmark$\checkmark \checkmark$\checkmark^\checkmark\\checkmarkc\checkmarki\checkmarkr\checkmarkc\checkmark$\checkmarkm\checkmark$\checkmark^\checkmark\\checkmarkc\checkmarki\checkmarkr\checkmarkc\checkmark$\checkmarka\checkmark$\checkmark^\checkmark\\checkmarkc\checkmarki\checkmarkr\checkmarkc\checkmark$\checkmarki\checkmark$\checkmark^\checkmark\\checkmarkc\checkmarki\checkmarkr\checkmarkc\checkmark$\checkmarkn\checkmark$\checkmark^\checkmark\\checkmarkc\checkmarki\checkmarkr\checkmarkc\checkmark$\checkmark.\checkmark$\checkmark^\checkmark\\checkmarkc\checkmarki\checkmarkr\checkmarkc\checkmark$\checkmarkt\checkmark$\checkmark^\checkmark\\checkmarkc\checkmarki\checkmarkr\checkmarkc\checkmark$\checkmarke\checkmark$\checkmark^\checkmark\\checkmarkc\checkmarki\checkmarkr\checkmarkc\checkmark$\checkmarkx\checkmark$\checkmark^\checkmark\\checkmarkc\checkmarki\checkmarkr\checkmarkc\checkmark$\checkmark
\checkmark$\checkmark^\checkmark\\checkmarkc\checkmarki\checkmarkr\checkmarkc\checkmark$\checkmark\\checkmark$\checkmark^\checkmark\\checkmarkc\checkmarki\checkmarkr\checkmarkc\checkmark$\checkmarkc\checkmark$\checkmark^\checkmark\\checkmarkc\checkmarki\checkmarkr\checkmarkc\checkmark$\checkmarkh\checkmark$\checkmark^\checkmark\\checkmarkc\checkmarki\checkmarkr\checkmarkc\checkmark$\checkmarka\checkmark$\checkmark^\checkmark\\checkmarkc\checkmarki\checkmarkr\checkmarkc\checkmark$\checkmarkp\checkmark$\checkmark^\checkmark\\checkmarkc\checkmarki\checkmarkr\checkmarkc\checkmark$\checkmarkt\checkmark$\checkmark^\checkmark\\checkmarkc\checkmarki\checkmarkr\checkmarkc\checkmark$\checkmarke\checkmark$\checkmark^\checkmark\\checkmarkc\checkmarki\checkmarkr\checkmarkc\checkmark$\checkmarkr\checkmark$\checkmark^\checkmark\\checkmarkc\checkmarki\checkmarkr\checkmarkc\checkmark$\checkmark{\checkmark$\checkmark^\checkmark\\checkmarkc\checkmarki\checkmarkr\checkmarkc\checkmark$\checkmarkÉ\checkmark$\checkmark^\checkmark\\checkmarkc\checkmarki\checkmarkr\checkmarkc\checkmark$\checkmarkt\checkmark$\checkmark^\checkmark\\checkmarkc\checkmarki\checkmarkr\checkmarkc\checkmark$\checkmarku\checkmark$\checkmark^\checkmark\\checkmarkc\checkmarki\checkmarkr\checkmarkc\checkmark$\checkmarkd\checkmark$\checkmark^\checkmark\\checkmarkc\checkmarki\checkmarkr\checkmarkc\checkmark$\checkmarke\checkmark$\checkmark^\checkmark\\checkmarkc\checkmarki\checkmarkr\checkmarkc\checkmark$\checkmark \checkmark$\checkmark^\checkmark\\checkmarkc\checkmarki\checkmarkr\checkmarkc\checkmark$\checkmarkT\checkmark$\checkmark^\checkmark\\checkmarkc\checkmarki\checkmarkr\checkmarkc\checkmark$\checkmarke\checkmark$\checkmark^\checkmark\\checkmarkc\checkmarki\checkmarkr\checkmarkc\checkmark$\checkmarkc\checkmark$\checkmark^\checkmark\\checkmarkc\checkmarki\checkmarkr\checkmarkc\checkmark$\checkmarkh\checkmark$\checkmark^\checkmark\\checkmarkc\checkmarki\checkmarkr\checkmarkc\checkmark$\checkmarkn\checkmark$\checkmark^\checkmark\\checkmarkc\checkmarki\checkmarkr\checkmarkc\checkmark$\checkmarki\checkmark$\checkmark^\checkmark\\checkmarkc\checkmarki\checkmarkr\checkmarkc\checkmark$\checkmarkq\checkmark$\checkmark^\checkmark\\checkmarkc\checkmarki\checkmarkr\checkmarkc\checkmark$\checkmarku\checkmark$\checkmark^\checkmark\\checkmarkc\checkmarki\checkmarkr\checkmarkc\checkmark$\checkmarke\checkmark$\checkmark^\checkmark\\checkmarkc\checkmarki\checkmarkr\checkmarkc\checkmark$\checkmark \checkmark$\checkmark^\checkmark\\checkmarkc\checkmarki\checkmarkr\checkmarkc\checkmark$\checkmarke\checkmark$\checkmark^\checkmark\\checkmarkc\checkmarki\checkmarkr\checkmarkc\checkmark$\checkmarkt\checkmark$\checkmark^\checkmark\\checkmarkc\checkmarki\checkmarkr\checkmarkc\checkmark$\checkmark \checkmark$\checkmark^\checkmark\\checkmarkc\checkmarki\checkmarkr\checkmarkc\checkmark$\checkmarkD\checkmark$\checkmark^\checkmark\\checkmarkc\checkmarki\checkmarkr\checkmarkc\checkmark$\checkmarki\checkmark$\checkmark^\checkmark\\checkmarkc\checkmarki\checkmarkr\checkmarkc\checkmark$\checkmarkm\checkmark$\checkmark^\checkmark\\checkmarkc\checkmarki\checkmarkr\checkmarkc\checkmark$\checkmarke\checkmark$\checkmark^\checkmark\\checkmarkc\checkmarki\checkmarkr\checkmarkc\checkmark$\checkmarkn\checkmark$\checkmark^\checkmark\\checkmarkc\checkmarki\checkmarkr\checkmarkc\checkmark$\checkmarks\checkmark$\checkmark^\checkmark\\checkmarkc\checkmarki\checkmarkr\checkmarkc\checkmark$\checkmarki\checkmark$\checkmark^\checkmark\\checkmarkc\checkmarki\checkmarkr\checkmarkc\checkmark$\checkmarko\checkmark$\checkmark^\checkmark\\checkmarkc\checkmarki\checkmarkr\checkmarkc\checkmark$\checkmarkn\checkmark$\checkmark^\checkmark\\checkmarkc\checkmarki\checkmarkr\checkmarkc\checkmark$\checkmarkn\checkmark$\checkmark^\checkmark\\checkmarkc\checkmarki\checkmarkr\checkmarkc\checkmark$\checkmarke\checkmark$\checkmark^\checkmark\\checkmarkc\checkmarki\checkmarkr\checkmarkc\checkmark$\checkmarkm\checkmark$\checkmark^\checkmark\\checkmarkc\checkmarki\checkmarkr\checkmarkc\checkmark$\checkmarke\checkmark$\checkmark^\checkmark\\checkmarkc\checkmarki\checkmarkr\checkmarkc\checkmark$\checkmarkn\checkmark$\checkmark^\checkmark\\checkmarkc\checkmarki\checkmarkr\checkmarkc\checkmark$\checkmarkt\checkmark$\checkmark^\checkmark\\checkmarkc\checkmarki\checkmarkr\checkmarkc\checkmark$\checkmark \checkmark$\checkmark^\checkmark\\checkmarkc\checkmarki\checkmarkr\checkmarkc\checkmark$\checkmarkM\checkmark$\checkmark^\checkmark\\checkmarkc\checkmarki\checkmarkr\checkmarkc\checkmark$\checkmarké\checkmark$\checkmark^\checkmark\\checkmarkc\checkmarki\checkmarkr\checkmarkc\checkmark$\checkmarkc\checkmark$\checkmark^\checkmark\\checkmarkc\checkmarki\checkmarkr\checkmarkc\checkmark$\checkmarka\checkmark$\checkmark^\checkmark\\checkmarkc\checkmarki\checkmarkr\checkmarkc\checkmark$\checkmarkn\checkmark$\checkmark^\checkmark\\checkmarkc\checkmarki\checkmarkr\checkmarkc\checkmark$\checkmarki\checkmark$\checkmark^\checkmark\\checkmarkc\checkmarki\checkmarkr\checkmarkc\checkmark$\checkmarkq\checkmark$\checkmark^\checkmark\\checkmarkc\checkmarki\checkmarkr\checkmarkc\checkmark$\checkmarku\checkmark$\checkmark^\checkmark\\checkmarkc\checkmarki\checkmarkr\checkmarkc\checkmark$\checkmarke\checkmark$\checkmark^\checkmark\\checkmarkc\checkmarki\checkmarkr\checkmarkc\checkmark$\checkmark}\checkmark$\checkmark^\checkmark\\checkmarkc\checkmarki\checkmarkr\checkmarkc\checkmark$\checkmark
\checkmark$\checkmark^\checkmark\\checkmarkc\checkmarki\checkmarkr\checkmarkc\checkmark$\checkmark\\checkmark$\checkmark^\checkmark\\checkmarkc\checkmarki\checkmarkr\checkmarkc\checkmark$\checkmarkl\checkmark$\checkmark^\checkmark\\checkmarkc\checkmarki\checkmarkr\checkmarkc\checkmark$\checkmarka\checkmark$\checkmark^\checkmark\\checkmarkc\checkmarki\checkmarkr\checkmarkc\checkmark$\checkmarkb\checkmark$\checkmark^\checkmark\\checkmarkc\checkmarki\checkmarkr\checkmarkc\checkmark$\checkmarke\checkmark$\checkmark^\checkmark\\checkmarkc\checkmarki\checkmarkr\checkmarkc\checkmark$\checkmarkl\checkmark$\checkmark^\checkmark\\checkmarkc\checkmarki\checkmarkr\checkmarkc\checkmark$\checkmark{\checkmark$\checkmark^\checkmark\\checkmarkc\checkmarki\checkmarkr\checkmarkc\checkmark$\checkmarkc\checkmark$\checkmark^\checkmark\\checkmarkc\checkmarki\checkmarkr\checkmarkc\checkmark$\checkmarkh\checkmark$\checkmark^\checkmark\\checkmarkc\checkmarki\checkmarkr\checkmarkc\checkmark$\checkmarka\checkmark$\checkmark^\checkmark\\checkmarkc\checkmarki\checkmarkr\checkmarkc\checkmark$\checkmarkp\checkmark$\checkmark^\checkmark\\checkmarkc\checkmarki\checkmarkr\checkmarkc\checkmark$\checkmark:\checkmark$\checkmark^\checkmark\\checkmarkc\checkmarki\checkmarkr\checkmarkc\checkmark$\checkmarkc\checkmark$\checkmark^\checkmark\\checkmarkc\checkmarki\checkmarkr\checkmarkc\checkmark$\checkmarko\checkmark$\checkmark^\checkmark\\checkmarkc\checkmarki\checkmarkr\checkmarkc\checkmark$\checkmarkn\checkmark$\checkmark^\checkmark\\checkmarkc\checkmarki\checkmarkr\checkmarkc\checkmark$\checkmarkc\checkmark$\checkmark^\checkmark\\checkmarkc\checkmarki\checkmarkr\checkmarkc\checkmark$\checkmarke\checkmark$\checkmark^\checkmark\\checkmarkc\checkmarki\checkmarkr\checkmarkc\checkmark$\checkmarkp\checkmark$\checkmark^\checkmark\\checkmarkc\checkmarki\checkmarkr\checkmarkc\checkmark$\checkmarkt\checkmark$\checkmark^\checkmark\\checkmarkc\checkmarki\checkmarkr\checkmarkc\checkmark$\checkmarki\checkmark$\checkmark^\checkmark\\checkmarkc\checkmarki\checkmarkr\checkmarkc\checkmark$\checkmarko\checkmark$\checkmark^\checkmark\\checkmarkc\checkmarki\checkmarkr\checkmarkc\checkmark$\checkmarkn\checkmark$\checkmark^\checkmark\\checkmarkc\checkmarki\checkmarkr\checkmarkc\checkmark$\checkmark_\checkmark$\checkmark^\checkmark\\checkmarkc\checkmarki\checkmarkr\checkmarkc\checkmark$\checkmarkm\checkmark$\checkmark^\checkmark\\checkmarkc\checkmarki\checkmarkr\checkmarkc\checkmark$\checkmarke\checkmark$\checkmark^\checkmark\\checkmarkc\checkmarki\checkmarkr\checkmarkc\checkmark$\checkmarkc\checkmark$\checkmark^\checkmark\\checkmarkc\checkmarki\checkmarkr\checkmarkc\checkmark$\checkmark}\checkmark$\checkmark^\checkmark\\checkmarkc\checkmarki\checkmarkr\checkmarkc\checkmark$\checkmark
\checkmark$\checkmark^\checkmark\\checkmarkc\checkmarki\checkmarkr\checkmarkc\checkmark$\checkmark
\checkmark$\checkmark^\checkmark\\checkmarkc\checkmarki\checkmarkr\checkmarkc\checkmark$\checkmark\\checkmark$\checkmark^\checkmark\\checkmarkc\checkmarki\checkmarkr\checkmarkc\checkmark$\checkmarks\checkmark$\checkmark^\checkmark\\checkmarkc\checkmarki\checkmarkr\checkmarkc\checkmark$\checkmarke\checkmark$\checkmark^\checkmark\\checkmarkc\checkmarki\checkmarkr\checkmarkc\checkmark$\checkmarkc\checkmark$\checkmark^\checkmark\\checkmarkc\checkmarki\checkmarkr\checkmarkc\checkmark$\checkmarkt\checkmark$\checkmark^\checkmark\\checkmarkc\checkmarki\checkmarkr\checkmarkc\checkmark$\checkmarki\checkmark$\checkmark^\checkmark\\checkmarkc\checkmarki\checkmarkr\checkmarkc\checkmark$\checkmarko\checkmark$\checkmark^\checkmark\\checkmarkc\checkmarki\checkmarkr\checkmarkc\checkmark$\checkmarkn\checkmark$\checkmark^\checkmark\\checkmarkc\checkmarki\checkmarkr\checkmarkc\checkmark$\checkmark{\checkmark$\checkmark^\checkmark\\checkmarkc\checkmarki\checkmarkr\checkmarkc\checkmark$\checkmarkI\checkmark$\checkmark^\checkmark\\checkmarkc\checkmarki\checkmarkr\checkmarkc\checkmark$\checkmarkn\checkmark$\checkmark^\checkmark\\checkmarkc\checkmarki\checkmarkr\checkmarkc\checkmark$\checkmarkt\checkmark$\checkmark^\checkmark\\checkmarkc\checkmarki\checkmarkr\checkmarkc\checkmark$\checkmarkr\checkmark$\checkmark^\checkmark\\checkmarkc\checkmarki\checkmarkr\checkmarkc\checkmark$\checkmarko\checkmark$\checkmark^\checkmark\\checkmarkc\checkmarki\checkmarkr\checkmarkc\checkmark$\checkmarkd\checkmark$\checkmark^\checkmark\\checkmarkc\checkmarki\checkmarkr\checkmarkc\checkmark$\checkmarku\checkmark$\checkmark^\checkmark\\checkmarkc\checkmarki\checkmarkr\checkmarkc\checkmark$\checkmarkc\checkmark$\checkmark^\checkmark\\checkmarkc\checkmarki\checkmarkr\checkmarkc\checkmark$\checkmarkt\checkmark$\checkmark^\checkmark\\checkmarkc\checkmarki\checkmarkr\checkmarkc\checkmark$\checkmarki\checkmark$\checkmark^\checkmark\\checkmarkc\checkmarki\checkmarkr\checkmarkc\checkmark$\checkmarko\checkmark$\checkmark^\checkmark\\checkmarkc\checkmarki\checkmarkr\checkmarkc\checkmark$\checkmarkn\checkmark$\checkmark^\checkmark\\checkmarkc\checkmarki\checkmarkr\checkmarkc\checkmark$\checkmark}\checkmark$\checkmark^\checkmark\\checkmarkc\checkmarki\checkmarkr\checkmarkc\checkmark$\checkmark
\checkmark$\checkmark^\checkmark\\checkmarkc\checkmarki\checkmarkr\checkmarkc\checkmark$\checkmarkL\checkmark$\checkmark^\checkmark\\checkmarkc\checkmarki\checkmarkr\checkmarkc\checkmark$\checkmarke\checkmark$\checkmark^\checkmark\\checkmarkc\checkmarki\checkmarkr\checkmarkc\checkmark$\checkmark \checkmark$\checkmark^\checkmark\\checkmarkc\checkmarki\checkmarkr\checkmarkc\checkmark$\checkmarkp\checkmark$\checkmark^\checkmark\\checkmarkc\checkmarki\checkmarkr\checkmarkc\checkmark$\checkmarkr\checkmark$\checkmark^\checkmark\\checkmarkc\checkmarki\checkmarkr\checkmarkc\checkmark$\checkmarké\checkmark$\checkmark^\checkmark\\checkmarkc\checkmarki\checkmarkr\checkmarkc\checkmark$\checkmarks\checkmark$\checkmark^\checkmark\\checkmarkc\checkmarki\checkmarkr\checkmarkc\checkmark$\checkmarke\checkmark$\checkmark^\checkmark\\checkmarkc\checkmarki\checkmarkr\checkmarkc\checkmark$\checkmarkn\checkmark$\checkmark^\checkmark\\checkmarkc\checkmarki\checkmarkr\checkmarkc\checkmark$\checkmarkt\checkmark$\checkmark^\checkmark\\checkmarkc\checkmarki\checkmarkr\checkmarkc\checkmark$\checkmark \checkmark$\checkmark^\checkmark\\checkmarkc\checkmarki\checkmarkr\checkmarkc\checkmark$\checkmarkc\checkmark$\checkmark^\checkmark\\checkmarkc\checkmarki\checkmarkr\checkmarkc\checkmark$\checkmarkh\checkmark$\checkmark^\checkmark\\checkmarkc\checkmarki\checkmarkr\checkmarkc\checkmark$\checkmarka\checkmark$\checkmark^\checkmark\\checkmarkc\checkmarki\checkmarkr\checkmarkc\checkmark$\checkmarkp\checkmark$\checkmark^\checkmark\\checkmarkc\checkmarki\checkmarkr\checkmarkc\checkmark$\checkmarki\checkmark$\checkmark^\checkmark\\checkmarkc\checkmarki\checkmarkr\checkmarkc\checkmark$\checkmarkt\checkmark$\checkmark^\checkmark\\checkmarkc\checkmarki\checkmarkr\checkmarkc\checkmark$\checkmarkr\checkmark$\checkmark^\checkmark\\checkmarkc\checkmarki\checkmarkr\checkmarkc\checkmark$\checkmarke\checkmark$\checkmark^\checkmark\\checkmarkc\checkmarki\checkmarkr\checkmarkc\checkmark$\checkmark \checkmark$\checkmark^\checkmark\\checkmarkc\checkmarki\checkmarkr\checkmarkc\checkmark$\checkmarkc\checkmark$\checkmark^\checkmark\\checkmarkc\checkmarki\checkmarkr\checkmarkc\checkmark$\checkmarko\checkmark$\checkmark^\checkmark\\checkmarkc\checkmarki\checkmarkr\checkmarkc\checkmark$\checkmarkn\checkmark$\checkmark^\checkmark\\checkmarkc\checkmarki\checkmarkr\checkmarkc\checkmark$\checkmarks\checkmark$\checkmark^\checkmark\\checkmarkc\checkmarki\checkmarkr\checkmarkc\checkmark$\checkmarkt\checkmark$\checkmark^\checkmark\\checkmarkc\checkmarki\checkmarkr\checkmarkc\checkmark$\checkmarki\checkmark$\checkmark^\checkmark\\checkmarkc\checkmarki\checkmarkr\checkmarkc\checkmark$\checkmarkt\checkmark$\checkmark^\checkmark\\checkmarkc\checkmarki\checkmarkr\checkmarkc\checkmark$\checkmarku\checkmark$\checkmark^\checkmark\\checkmarkc\checkmarki\checkmarkr\checkmarkc\checkmark$\checkmarke\checkmark$\checkmark^\checkmark\\checkmarkc\checkmarki\checkmarkr\checkmarkc\checkmark$\checkmark \checkmark$\checkmark^\checkmark\\checkmarkc\checkmarki\checkmarkr\checkmarkc\checkmark$\checkmarkl\checkmark$\checkmark^\checkmark\\checkmarkc\checkmarki\checkmarkr\checkmarkc\checkmark$\checkmarke\checkmark$\checkmark^\checkmark\\checkmarkc\checkmarki\checkmarkr\checkmarkc\checkmark$\checkmark \checkmark$\checkmark^\checkmark\\checkmarkc\checkmarki\checkmarkr\checkmarkc\checkmark$\checkmarkc\checkmark$\checkmark^\checkmark\\checkmarkc\checkmarki\checkmarkr\checkmarkc\checkmark$\checkmarkœ\checkmark$\checkmark^\checkmark\\checkmarkc\checkmarki\checkmarkr\checkmarkc\checkmark$\checkmarku\checkmark$\checkmark^\checkmark\\checkmarkc\checkmarki\checkmarkr\checkmarkc\checkmark$\checkmarkr\checkmark$\checkmark^\checkmark\\checkmarkc\checkmarki\checkmarkr\checkmarkc\checkmark$\checkmark \checkmark$\checkmark^\checkmark\\checkmarkc\checkmarki\checkmarkr\checkmarkc\checkmark$\checkmarkd\checkmark$\checkmark^\checkmark\\checkmarkc\checkmarki\checkmarkr\checkmarkc\checkmark$\checkmarke\checkmark$\checkmark^\checkmark\\checkmarkc\checkmarki\checkmarkr\checkmarkc\checkmark$\checkmark \checkmark$\checkmark^\checkmark\\checkmarkc\checkmarki\checkmarkr\checkmarkc\checkmark$\checkmarkl\checkmark$\checkmark^\checkmark\\checkmarkc\checkmarki\checkmarkr\checkmarkc\checkmark$\checkmarka\checkmark$\checkmark^\checkmark\\checkmarkc\checkmarki\checkmarkr\checkmarkc\checkmark$\checkmark \checkmark$\checkmark^\checkmark\\checkmarkc\checkmarki\checkmarkr\checkmarkc\checkmark$\checkmarkj\checkmark$\checkmark^\checkmark\\checkmarkc\checkmarki\checkmarkr\checkmarkc\checkmark$\checkmarku\checkmark$\checkmark^\checkmark\\checkmarkc\checkmarki\checkmarkr\checkmarkc\checkmark$\checkmarks\checkmark$\checkmark^\checkmark\\checkmarkc\checkmarki\checkmarkr\checkmarkc\checkmark$\checkmarkt\checkmark$\checkmark^\checkmark\\checkmarkc\checkmarki\checkmarkr\checkmarkc\checkmark$\checkmarki\checkmark$\checkmark^\checkmark\\checkmarkc\checkmarki\checkmarkr\checkmarkc\checkmark$\checkmarkf\checkmark$\checkmark^\checkmark\\checkmarkc\checkmarki\checkmarkr\checkmarkc\checkmark$\checkmarki\checkmark$\checkmark^\checkmark\\checkmarkc\checkmarki\checkmarkr\checkmarkc\checkmark$\checkmarkc\checkmark$\checkmark^\checkmark\\checkmarkc\checkmarki\checkmarkr\checkmarkc\checkmark$\checkmarka\checkmark$\checkmark^\checkmark\\checkmarkc\checkmarki\checkmarkr\checkmarkc\checkmark$\checkmarkt\checkmark$\checkmark^\checkmark\\checkmarkc\checkmarki\checkmarkr\checkmarkc\checkmark$\checkmarki\checkmark$\checkmark^\checkmark\\checkmarkc\checkmarki\checkmarkr\checkmarkc\checkmark$\checkmarko\checkmark$\checkmark^\checkmark\\checkmarkc\checkmarki\checkmarkr\checkmarkc\checkmark$\checkmarkn\checkmark$\checkmark^\checkmark\\checkmarkc\checkmarki\checkmarkr\checkmarkc\checkmark$\checkmark \checkmark$\checkmark^\checkmark\\checkmarkc\checkmarki\checkmarkr\checkmarkc\checkmark$\checkmarkp\checkmark$\checkmark^\checkmark\\checkmarkc\checkmarki\checkmarkr\checkmarkc\checkmark$\checkmarkh\checkmark$\checkmark^\checkmark\\checkmarkc\checkmarki\checkmarkr\checkmarkc\checkmark$\checkmarky\checkmark$\checkmark^\checkmark\\checkmarkc\checkmarki\checkmarkr\checkmarkc\checkmark$\checkmarks\checkmark$\checkmark^\checkmark\\checkmarkc\checkmarki\checkmarkr\checkmarkc\checkmark$\checkmarki\checkmark$\checkmark^\checkmark\\checkmarkc\checkmarki\checkmarkr\checkmarkc\checkmark$\checkmarkq\checkmark$\checkmark^\checkmark\\checkmarkc\checkmarki\checkmarkr\checkmarkc\checkmark$\checkmarku\checkmark$\checkmark^\checkmark\\checkmarkc\checkmarki\checkmarkr\checkmarkc\checkmark$\checkmarke\checkmark$\checkmark^\checkmark\\checkmarkc\checkmarki\checkmarkr\checkmarkc\checkmark$\checkmark \checkmark$\checkmark^\checkmark\\checkmarkc\checkmarki\checkmarkr\checkmarkc\checkmark$\checkmarke\checkmark$\checkmark^\checkmark\\checkmarkc\checkmarki\checkmarkr\checkmarkc\checkmark$\checkmarkt\checkmark$\checkmark^\checkmark\\checkmarkc\checkmarki\checkmarkr\checkmarkc\checkmark$\checkmark \checkmark$\checkmark^\checkmark\\checkmarkc\checkmarki\checkmarkr\checkmarkc\checkmark$\checkmarkt\checkmark$\checkmark^\checkmark\\checkmarkc\checkmarki\checkmarkr\checkmarkc\checkmark$\checkmarke\checkmark$\checkmark^\checkmark\\checkmarkc\checkmarki\checkmarkr\checkmarkc\checkmark$\checkmarkc\checkmark$\checkmark^\checkmark\\checkmarkc\checkmarki\checkmarkr\checkmarkc\checkmark$\checkmarkh\checkmark$\checkmark^\checkmark\\checkmarkc\checkmarki\checkmarkr\checkmarkc\checkmark$\checkmarkn\checkmark$\checkmark^\checkmark\\checkmarkc\checkmarki\checkmarkr\checkmarkc\checkmark$\checkmarki\checkmark$\checkmark^\checkmark\\checkmarkc\checkmarki\checkmarkr\checkmarkc\checkmark$\checkmarkq\checkmark$\checkmark^\checkmark\\checkmarkc\checkmarki\checkmarkr\checkmarkc\checkmark$\checkmarku\checkmark$\checkmark^\checkmark\\checkmarkc\checkmarki\checkmarkr\checkmarkc\checkmark$\checkmarke\checkmark$\checkmark^\checkmark\\checkmarkc\checkmarki\checkmarkr\checkmarkc\checkmark$\checkmark \checkmark$\checkmark^\checkmark\\checkmarkc\checkmarki\checkmarkr\checkmarkc\checkmark$\checkmarkd\checkmark$\checkmark^\checkmark\\checkmarkc\checkmarki\checkmarkr\checkmarkc\checkmark$\checkmarku\checkmark$\checkmark^\checkmark\\checkmarkc\checkmarki\checkmarkr\checkmarkc\checkmark$\checkmark \checkmark$\checkmark^\checkmark\\checkmarkc\checkmarki\checkmarkr\checkmarkc\checkmark$\checkmarkp\checkmark$\checkmark^\checkmark\\checkmarkc\checkmarki\checkmarkr\checkmarkc\checkmark$\checkmarkr\checkmark$\checkmark^\checkmark\\checkmarkc\checkmarki\checkmarkr\checkmarkc\checkmark$\checkmarko\checkmark$\checkmark^\checkmark\\checkmarkc\checkmarki\checkmarkr\checkmarkc\checkmark$\checkmarkj\checkmark$\checkmark^\checkmark\\checkmarkc\checkmarki\checkmarkr\checkmarkc\checkmark$\checkmarke\checkmark$\checkmark^\checkmark\\checkmarkc\checkmarki\checkmarkr\checkmarkc\checkmark$\checkmarkt\checkmark$\checkmark^\checkmark\\checkmarkc\checkmarki\checkmarkr\checkmarkc\checkmark$\checkmark \checkmark$\checkmark^\checkmark\\checkmarkc\checkmarki\checkmarkr\checkmarkc\checkmark$\checkmarkd\checkmark$\checkmark^\checkmark\\checkmarkc\checkmarki\checkmarkr\checkmarkc\checkmark$\checkmarke\checkmark$\checkmark^\checkmark\\checkmarkc\checkmarki\checkmarkr\checkmarkc\checkmark$\checkmark \checkmark$\checkmark^\checkmark\\checkmarkc\checkmarki\checkmarkr\checkmarkc\checkmark$\checkmarkc\checkmark$\checkmark^\checkmark\\checkmarkc\checkmarki\checkmarkr\checkmarkc\checkmark$\checkmarko\checkmark$\checkmark^\checkmark\\checkmarkc\checkmarki\checkmarkr\checkmarkc\checkmark$\checkmarkn\checkmark$\checkmark^\checkmark\\checkmarkc\checkmarki\checkmarkr\checkmarkc\checkmark$\checkmarkc\checkmark$\checkmark^\checkmark\\checkmarkc\checkmarki\checkmarkr\checkmarkc\checkmark$\checkmarke\checkmark$\checkmark^\checkmark\\checkmarkc\checkmarki\checkmarkr\checkmarkc\checkmark$\checkmarkp\checkmark$\checkmark^\checkmark\\checkmarkc\checkmarki\checkmarkr\checkmarkc\checkmark$\checkmarkt\checkmark$\checkmark^\checkmark\\checkmarkc\checkmarki\checkmarkr\checkmarkc\checkmark$\checkmarki\checkmark$\checkmark^\checkmark\\checkmarkc\checkmarki\checkmarkr\checkmarkc\checkmark$\checkmarko\checkmark$\checkmark^\checkmark\\checkmarkc\checkmarki\checkmarkr\checkmarkc\checkmark$\checkmarkn\checkmark$\checkmark^\checkmark\\checkmarkc\checkmarki\checkmarkr\checkmarkc\checkmark$\checkmark \checkmark$\checkmark^\checkmark\\checkmarkc\checkmarki\checkmarkr\checkmarkc\checkmark$\checkmarkd\checkmark$\checkmark^\checkmark\\checkmarkc\checkmarki\checkmarkr\checkmarkc\checkmark$\checkmark'\checkmark$\checkmark^\checkmark\\checkmarkc\checkmarki\checkmarkr\checkmarkc\checkmark$\checkmarku\checkmark$\checkmark^\checkmark\\checkmarkc\checkmarki\checkmarkr\checkmarkc\checkmark$\checkmarkn\checkmark$\checkmark^\checkmark\\checkmarkc\checkmarki\checkmarkr\checkmarkc\checkmark$\checkmarke\checkmark$\checkmark^\checkmark\\checkmarkc\checkmarki\checkmarkr\checkmarkc\checkmark$\checkmark \checkmark$\checkmark^\checkmark\\checkmarkc\checkmarki\checkmarkr\checkmarkc\checkmark$\checkmarkm\checkmark$\checkmark^\checkmark\\checkmarkc\checkmarki\checkmarkr\checkmarkc\checkmark$\checkmarki\checkmark$\checkmark^\checkmark\\checkmarkc\checkmarki\checkmarkr\checkmarkc\checkmark$\checkmarkn\checkmark$\checkmark^\checkmark\\checkmarkc\checkmarki\checkmarkr\checkmarkc\checkmark$\checkmarki\checkmark$\checkmark^\checkmark\\checkmarkc\checkmarki\checkmarkr\checkmarkc\checkmark$\checkmark-\checkmark$\checkmark^\checkmark\\checkmarkc\checkmarki\checkmarkr\checkmarkc\checkmark$\checkmarkf\checkmark$\checkmark^\checkmark\\checkmarkc\checkmarki\checkmarkr\checkmarkc\checkmark$\checkmarkr\checkmark$\checkmark^\checkmark\\checkmarkc\checkmarki\checkmarkr\checkmarkc\checkmark$\checkmarka\checkmark$\checkmark^\checkmark\\checkmarkc\checkmarki\checkmarkr\checkmarkc\checkmark$\checkmarki\checkmark$\checkmark^\checkmark\\checkmarkc\checkmarki\checkmarkr\checkmarkc\checkmark$\checkmarks\checkmark$\checkmark^\checkmark\\checkmarkc\checkmarki\checkmarkr\checkmarkc\checkmark$\checkmarke\checkmark$\checkmark^\checkmark\\checkmarkc\checkmarki\checkmarkr\checkmarkc\checkmark$\checkmarku\checkmark$\checkmark^\checkmark\\checkmarkc\checkmarki\checkmarkr\checkmarkc\checkmark$\checkmarks\checkmark$\checkmark^\checkmark\\checkmarkc\checkmarki\checkmarkr\checkmarkc\checkmark$\checkmarke\checkmark$\checkmark^\checkmark\\checkmarkc\checkmarki\checkmarkr\checkmarkc\checkmark$\checkmark \checkmark$\checkmark^\checkmark\\checkmarkc\checkmarki\checkmarkr\checkmarkc\checkmark$\checkmarkC\checkmark$\checkmark^\checkmark\\checkmarkc\checkmarki\checkmarkr\checkmarkc\checkmark$\checkmarkN\checkmark$\checkmark^\checkmark\\checkmarkc\checkmarki\checkmarkr\checkmarkc\checkmark$\checkmarkC\checkmark$\checkmark^\checkmark\\checkmarkc\checkmarki\checkmarkr\checkmarkc\checkmark$\checkmark \checkmark$\checkmark^\checkmark\\checkmarkc\checkmarki\checkmarkr\checkmarkc\checkmark$\checkmark5\checkmark$\checkmark^\checkmark\\checkmarkc\checkmarki\checkmarkr\checkmarkc\checkmark$\checkmark \checkmark$\checkmark^\checkmark\\checkmarkc\checkmarki\checkmarkr\checkmarkc\checkmark$\checkmarka\checkmark$\checkmark^\checkmark\\checkmarkc\checkmarki\checkmarkr\checkmarkc\checkmark$\checkmarkx\checkmark$\checkmark^\checkmark\\checkmarkc\checkmarki\checkmarkr\checkmarkc\checkmark$\checkmarke\checkmark$\checkmark^\checkmark\\checkmarkc\checkmarki\checkmarkr\checkmarkc\checkmark$\checkmarks\checkmark$\checkmark^\checkmark\\checkmarkc\checkmarki\checkmarkr\checkmarkc\checkmark$\checkmark.\checkmark$\checkmark^\checkmark\\checkmarkc\checkmarki\checkmarkr\checkmarkc\checkmark$\checkmark \checkmark$\checkmark^\checkmark\\checkmarkc\checkmarki\checkmarkr\checkmarkc\checkmark$\checkmarkL\checkmark$\checkmark^\checkmark\\checkmarkc\checkmarki\checkmarkr\checkmarkc\checkmark$\checkmark'\checkmark$\checkmark^\checkmark\\checkmarkc\checkmarki\checkmarkr\checkmarkc\checkmark$\checkmarké\checkmark$\checkmark^\checkmark\\checkmarkc\checkmarki\checkmarkr\checkmarkc\checkmark$\checkmarkt\checkmark$\checkmark^\checkmark\\checkmarkc\checkmarki\checkmarkr\checkmarkc\checkmark$\checkmarku\checkmark$\checkmark^\checkmark\\checkmarkc\checkmarki\checkmarkr\checkmarkc\checkmark$\checkmarkd\checkmark$\checkmark^\checkmark\\checkmarkc\checkmarki\checkmarkr\checkmarkc\checkmark$\checkmarke\checkmark$\checkmark^\checkmark\\checkmarkc\checkmarki\checkmarkr\checkmarkc\checkmark$\checkmark \checkmark$\checkmark^\checkmark\\checkmarkc\checkmarki\checkmarkr\checkmarkc\checkmark$\checkmarks\checkmark$\checkmark^\checkmark\\checkmarkc\checkmarki\checkmarkr\checkmarkc\checkmark$\checkmark'\checkmark$\checkmark^\checkmark\\checkmarkc\checkmarki\checkmarkr\checkmarkc\checkmark$\checkmarka\checkmark$\checkmark^\checkmark\\checkmarkc\checkmarki\checkmarkr\checkmarkc\checkmark$\checkmarkr\checkmark$\checkmark^\checkmark\\checkmarkc\checkmarki\checkmarkr\checkmarkc\checkmark$\checkmarkt\checkmark$\checkmark^\checkmark\\checkmarkc\checkmarki\checkmarkr\checkmarkc\checkmark$\checkmarki\checkmark$\checkmark^\checkmark\\checkmarkc\checkmarki\checkmarkr\checkmarkc\checkmark$\checkmarkc\checkmark$\checkmark^\checkmark\\checkmarkc\checkmarki\checkmarkr\checkmarkc\checkmark$\checkmarku\checkmark$\checkmark^\checkmark\\checkmarkc\checkmarki\checkmarkr\checkmarkc\checkmark$\checkmarkl\checkmark$\checkmark^\checkmark\\checkmarkc\checkmarki\checkmarkr\checkmarkc\checkmark$\checkmarke\checkmark$\checkmark^\checkmark\\checkmarkc\checkmarki\checkmarkr\checkmarkc\checkmark$\checkmark \checkmark$\checkmark^\checkmark\\checkmarkc\checkmarki\checkmarkr\checkmarkc\checkmark$\checkmarka\checkmark$\checkmark^\checkmark\\checkmarkc\checkmarki\checkmarkr\checkmarkc\checkmark$\checkmarku\checkmark$\checkmark^\checkmark\\checkmarkc\checkmarki\checkmarkr\checkmarkc\checkmark$\checkmarkt\checkmark$\checkmark^\checkmark\\checkmarkc\checkmarki\checkmarkr\checkmarkc\checkmark$\checkmarko\checkmark$\checkmark^\checkmark\\checkmarkc\checkmarki\checkmarkr\checkmarkc\checkmark$\checkmarku\checkmark$\checkmark^\checkmark\\checkmarkc\checkmarki\checkmarkr\checkmarkc\checkmark$\checkmarkr\checkmark$\checkmark^\checkmark\\checkmarkc\checkmarki\checkmarkr\checkmarkc\checkmark$\checkmark \checkmark$\checkmark^\checkmark\\checkmarkc\checkmarki\checkmarkr\checkmarkc\checkmark$\checkmarkd\checkmark$\checkmark^\checkmark\\checkmarkc\checkmarki\checkmarkr\checkmarkc\checkmark$\checkmarke\checkmark$\checkmark^\checkmark\\checkmarkc\checkmarki\checkmarkr\checkmarkc\checkmark$\checkmarks\checkmark$\checkmark^\checkmark\\checkmarkc\checkmarki\checkmarkr\checkmarkc\checkmark$\checkmark \checkmark$\checkmark^\checkmark\\checkmarkc\checkmarki\checkmarkr\checkmarkc\checkmark$\checkmarka\checkmark$\checkmark^\checkmark\\checkmarkc\checkmarki\checkmarkr\checkmarkc\checkmark$\checkmarkx\checkmark$\checkmark^\checkmark\\checkmarkc\checkmarki\checkmarkr\checkmarkc\checkmark$\checkmarke\checkmark$\checkmark^\checkmark\\checkmarkc\checkmarki\checkmarkr\checkmarkc\checkmark$\checkmarks\checkmark$\checkmark^\checkmark\\checkmarkc\checkmarki\checkmarkr\checkmarkc\checkmark$\checkmark \checkmark$\checkmark^\checkmark\\checkmarkc\checkmarki\checkmarkr\checkmarkc\checkmark$\checkmarks\checkmark$\checkmark^\checkmark\\checkmarkc\checkmarki\checkmarkr\checkmarkc\checkmark$\checkmarku\checkmark$\checkmark^\checkmark\\checkmarkc\checkmarki\checkmarkr\checkmarkc\checkmark$\checkmarki\checkmark$\checkmark^\checkmark\\checkmarkc\checkmarki\checkmarkr\checkmarkc\checkmark$\checkmarkv\checkmark$\checkmark^\checkmark\\checkmarkc\checkmarki\checkmarkr\checkmarkc\checkmark$\checkmarka\checkmark$\checkmark^\checkmark\\checkmarkc\checkmarki\checkmarkr\checkmarkc\checkmark$\checkmarkn\checkmark$\checkmark^\checkmark\\checkmarkc\checkmarki\checkmarkr\checkmarkc\checkmark$\checkmarkt\checkmark$\checkmark^\checkmark\\checkmarkc\checkmarki\checkmarkr\checkmarkc\checkmark$\checkmarks\checkmark$\checkmark^\checkmark\\checkmarkc\checkmarki\checkmarkr\checkmarkc\checkmark$\checkmark \checkmark$\checkmark^\checkmark\\checkmarkc\checkmarki\checkmarkr\checkmarkc\checkmark$\checkmark:\checkmark$\checkmark^\checkmark\\checkmarkc\checkmarki\checkmarkr\checkmarkc\checkmark$\checkmark
\checkmark$\checkmark^\checkmark\\checkmarkc\checkmarki\checkmarkr\checkmarkc\checkmark$\checkmark\\checkmark$\checkmark^\checkmark\\checkmarkc\checkmarki\checkmarkr\checkmarkc\checkmark$\checkmarkb\checkmark$\checkmark^\checkmark\\checkmarkc\checkmarki\checkmarkr\checkmarkc\checkmark$\checkmarke\checkmark$\checkmark^\checkmark\\checkmarkc\checkmarki\checkmarkr\checkmarkc\checkmark$\checkmarkg\checkmark$\checkmark^\checkmark\\checkmarkc\checkmarki\checkmarkr\checkmarkc\checkmark$\checkmarki\checkmark$\checkmark^\checkmark\\checkmarkc\checkmarki\checkmarkr\checkmarkc\checkmark$\checkmarkn\checkmark$\checkmark^\checkmark\\checkmarkc\checkmarki\checkmarkr\checkmarkc\checkmark$\checkmark{\checkmark$\checkmark^\checkmark\\checkmarkc\checkmarki\checkmarkr\checkmarkc\checkmark$\checkmarki\checkmark$\checkmark^\checkmark\\checkmarkc\checkmarki\checkmarkr\checkmarkc\checkmark$\checkmarkt\checkmark$\checkmark^\checkmark\\checkmarkc\checkmarki\checkmarkr\checkmarkc\checkmark$\checkmarke\checkmark$\checkmark^\checkmark\\checkmarkc\checkmarki\checkmarkr\checkmarkc\checkmark$\checkmarkm\checkmark$\checkmark^\checkmark\\checkmarkc\checkmarki\checkmarkr\checkmarkc\checkmark$\checkmarki\checkmark$\checkmark^\checkmark\\checkmarkc\checkmarki\checkmarkr\checkmarkc\checkmark$\checkmarkz\checkmark$\checkmark^\checkmark\\checkmarkc\checkmarki\checkmarkr\checkmarkc\checkmark$\checkmarke\checkmark$\checkmark^\checkmark\\checkmarkc\checkmarki\checkmarkr\checkmarkc\checkmark$\checkmark}\checkmark$\checkmark^\checkmark\\checkmarkc\checkmarki\checkmarkr\checkmarkc\checkmark$\checkmark
\checkmark$\checkmark^\checkmark\\checkmarkc\checkmarki\checkmarkr\checkmarkc\checkmark$\checkmark \checkmark$\checkmark^\checkmark\\checkmarkc\checkmarki\checkmarkr\checkmarkc\checkmark$\checkmark \checkmark$\checkmark^\checkmark\\checkmarkc\checkmarki\checkmarkr\checkmarkc\checkmark$\checkmark \checkmark$\checkmark^\checkmark\\checkmarkc\checkmarki\checkmarkr\checkmarkc\checkmark$\checkmark \checkmark$\checkmark^\checkmark\\checkmarkc\checkmarki\checkmarkr\checkmarkc\checkmark$\checkmark\\checkmark$\checkmark^\checkmark\\checkmarkc\checkmarki\checkmarkr\checkmarkc\checkmark$\checkmarki\checkmark$\checkmark^\checkmark\\checkmarkc\checkmarki\checkmarkr\checkmarkc\checkmark$\checkmarkt\checkmark$\checkmark^\checkmark\\checkmarkc\checkmarki\checkmarkr\checkmarkc\checkmark$\checkmarke\checkmark$\checkmark^\checkmark\\checkmarkc\checkmarki\checkmarkr\checkmarkc\checkmark$\checkmarkm\checkmark$\checkmark^\checkmark\\checkmarkc\checkmarki\checkmarkr\checkmarkc\checkmark$\checkmark \checkmark$\checkmark^\checkmark\\checkmarkc\checkmarki\checkmarkr\checkmarkc\checkmark$\checkmarkL\checkmark$\checkmark^\checkmark\\checkmarkc\checkmarki\checkmarkr\checkmarkc\checkmark$\checkmark'\checkmark$\checkmark^\checkmark\\checkmarkc\checkmarki\checkmarkr\checkmarkc\checkmark$\checkmarka\checkmark$\checkmark^\checkmark\\checkmarkc\checkmarki\checkmarkr\checkmarkc\checkmark$\checkmarkn\checkmark$\checkmark^\checkmark\\checkmarkc\checkmarki\checkmarkr\checkmarkc\checkmark$\checkmarka\checkmark$\checkmark^\checkmark\\checkmarkc\checkmarki\checkmarkr\checkmarkc\checkmark$\checkmarkl\checkmark$\checkmark^\checkmark\\checkmarkc\checkmarki\checkmarkr\checkmarkc\checkmark$\checkmarky\checkmark$\checkmark^\checkmark\\checkmarkc\checkmarki\checkmarkr\checkmarkc\checkmark$\checkmarks\checkmark$\checkmark^\checkmark\\checkmarkc\checkmarki\checkmarkr\checkmarkc\checkmark$\checkmarke\checkmark$\checkmark^\checkmark\\checkmarkc\checkmarki\checkmarkr\checkmarkc\checkmark$\checkmark \checkmark$\checkmark^\checkmark\\checkmarkc\checkmarki\checkmarkr\checkmarkc\checkmark$\checkmarkf\checkmark$\checkmark^\checkmark\\checkmarkc\checkmarki\checkmarkr\checkmarkc\checkmark$\checkmarko\checkmark$\checkmark^\checkmark\\checkmarkc\checkmarki\checkmarkr\checkmarkc\checkmark$\checkmarkn\checkmark$\checkmark^\checkmark\\checkmarkc\checkmarki\checkmarkr\checkmarkc\checkmark$\checkmarkc\checkmark$\checkmark^\checkmark\\checkmarkc\checkmarki\checkmarkr\checkmarkc\checkmark$\checkmarkt\checkmark$\checkmark^\checkmark\\checkmarkc\checkmarki\checkmarkr\checkmarkc\checkmark$\checkmarki\checkmark$\checkmark^\checkmark\\checkmarkc\checkmarki\checkmarkr\checkmarkc\checkmark$\checkmarko\checkmark$\checkmark^\checkmark\\checkmarkc\checkmarki\checkmarkr\checkmarkc\checkmark$\checkmarkn\checkmark$\checkmark^\checkmark\\checkmarkc\checkmarki\checkmarkr\checkmarkc\checkmark$\checkmarkn\checkmark$\checkmark^\checkmark\\checkmarkc\checkmarki\checkmarkr\checkmarkc\checkmark$\checkmarke\checkmark$\checkmark^\checkmark\\checkmarkc\checkmarki\checkmarkr\checkmarkc\checkmark$\checkmarkl\checkmark$\checkmark^\checkmark\\checkmarkc\checkmarki\checkmarkr\checkmarkc\checkmark$\checkmarkl\checkmark$\checkmark^\checkmark\\checkmarkc\checkmarki\checkmarkr\checkmarkc\checkmark$\checkmarke\checkmark$\checkmark^\checkmark\\checkmarkc\checkmarki\checkmarkr\checkmarkc\checkmark$\checkmark \checkmark$\checkmark^\checkmark\\checkmarkc\checkmarki\checkmarkr\checkmarkc\checkmark$\checkmarkd\checkmark$\checkmark^\checkmark\\checkmarkc\checkmarki\checkmarkr\checkmarkc\checkmark$\checkmarku\checkmark$\checkmark^\checkmark\\checkmarkc\checkmarki\checkmarkr\checkmarkc\checkmark$\checkmark \checkmark$\checkmark^\checkmark\\checkmarkc\checkmarki\checkmarkr\checkmarkc\checkmark$\checkmarkb\checkmark$\checkmark^\checkmark\\checkmarkc\checkmarki\checkmarkr\checkmarkc\checkmark$\checkmarke\checkmark$\checkmark^\checkmark\\checkmarkc\checkmarki\checkmarkr\checkmarkc\checkmark$\checkmarks\checkmark$\checkmark^\checkmark\\checkmarkc\checkmarki\checkmarkr\checkmarkc\checkmark$\checkmarko\checkmark$\checkmark^\checkmark\\checkmarkc\checkmarki\checkmarkr\checkmarkc\checkmark$\checkmarki\checkmark$\checkmark^\checkmark\\checkmarkc\checkmarki\checkmarkr\checkmarkc\checkmark$\checkmarkn\checkmark$\checkmark^\checkmark\\checkmarkc\checkmarki\checkmarkr\checkmarkc\checkmark$\checkmark \checkmark$\checkmark^\checkmark\\checkmarkc\checkmarki\checkmarkr\checkmarkc\checkmark$\checkmark(\checkmark$\checkmark^\checkmark\\checkmarkc\checkmarki\checkmarkr\checkmarkc\checkmark$\checkmarkA\checkmark$\checkmark^\checkmark\\checkmarkc\checkmarki\checkmarkr\checkmarkc\checkmark$\checkmarkF\checkmark$\checkmark^\checkmark\\checkmarkc\checkmarki\checkmarkr\checkmarkc\checkmark$\checkmarkB\checkmark$\checkmark^\checkmark\\checkmarkc\checkmarki\checkmarkr\checkmarkc\checkmark$\checkmark)\checkmark$\checkmark^\checkmark\\checkmarkc\checkmarki\checkmarkr\checkmarkc\checkmark$\checkmark \checkmark$\checkmark^\checkmark\\checkmarkc\checkmarki\checkmarkr\checkmarkc\checkmark$\checkmarka\checkmark$\checkmark^\checkmark\\checkmarkc\checkmarki\checkmarkr\checkmarkc\checkmark$\checkmarkv\checkmark$\checkmark^\checkmark\\checkmarkc\checkmarki\checkmarkr\checkmarkc\checkmark$\checkmarke\checkmark$\checkmark^\checkmark\\checkmarkc\checkmarki\checkmarkr\checkmarkc\checkmark$\checkmarkc\checkmark$\checkmark^\checkmark\\checkmarkc\checkmarki\checkmarkr\checkmarkc\checkmark$\checkmark \checkmark$\checkmark^\checkmark\\checkmarkc\checkmarki\checkmarkr\checkmarkc\checkmark$\checkmarkl\checkmark$\checkmark^\checkmark\\checkmarkc\checkmarki\checkmarkr\checkmarkc\checkmark$\checkmarke\checkmark$\checkmark^\checkmark\\checkmarkc\checkmarki\checkmarkr\checkmarkc\checkmark$\checkmarks\checkmark$\checkmark^\checkmark\\checkmarkc\checkmarki\checkmarkr\checkmarkc\checkmark$\checkmark \checkmark$\checkmark^\checkmark\\checkmarkc\checkmarki\checkmarkr\checkmarkc\checkmark$\checkmarko\checkmark$\checkmark^\checkmark\\checkmarkc\checkmarki\checkmarkr\checkmarkc\checkmark$\checkmarku\checkmark$\checkmark^\checkmark\\checkmarkc\checkmarki\checkmarkr\checkmarkc\checkmark$\checkmarkt\checkmark$\checkmark^\checkmark\\checkmarkc\checkmarki\checkmarkr\checkmarkc\checkmark$\checkmarki\checkmark$\checkmark^\checkmark\\checkmarkc\checkmarki\checkmarkr\checkmarkc\checkmark$\checkmarkl\checkmark$\checkmark^\checkmark\\checkmarkc\checkmarki\checkmarkr\checkmarkc\checkmark$\checkmarks\checkmark$\checkmark^\checkmark\\checkmarkc\checkmarki\checkmarkr\checkmarkc\checkmark$\checkmark \checkmark$\checkmark^\checkmark\\checkmarkc\checkmarki\checkmarkr\checkmarkc\checkmark$\checkmark:\checkmark$\checkmark^\checkmark\\checkmarkc\checkmarki\checkmarkr\checkmarkc\checkmark$\checkmark \checkmark$\checkmark^\checkmark\\checkmarkc\checkmarki\checkmarkr\checkmarkc\checkmark$\checkmarkb\checkmark$\checkmark^\checkmark\\checkmarkc\checkmarki\checkmarkr\checkmarkc\checkmark$\checkmarkê\checkmark$\checkmark^\checkmark\\checkmarkc\checkmarki\checkmarkr\checkmarkc\checkmark$\checkmarkt\checkmark$\checkmark^\checkmark\\checkmarkc\checkmarki\checkmarkr\checkmarkc\checkmark$\checkmarke\checkmark$\checkmark^\checkmark\\checkmarkc\checkmarki\checkmarkr\checkmarkc\checkmark$\checkmark \checkmark$\checkmark^\checkmark\\checkmarkc\checkmarki\checkmarkr\checkmarkc\checkmark$\checkmarkà\checkmark$\checkmark^\checkmark\\checkmarkc\checkmarki\checkmarkr\checkmarkc\checkmark$\checkmark \checkmark$\checkmark^\checkmark\\checkmarkc\checkmarki\checkmarkr\checkmarkc\checkmark$\checkmarkc\checkmark$\checkmark^\checkmark\\checkmarkc\checkmarki\checkmarkr\checkmarkc\checkmark$\checkmarko\checkmark$\checkmark^\checkmark\\checkmarkc\checkmarki\checkmarkr\checkmarkc\checkmark$\checkmarkr\checkmark$\checkmark^\checkmark\\checkmarkc\checkmarki\checkmarkr\checkmarkc\checkmark$\checkmarkn\checkmark$\checkmark^\checkmark\\checkmarkc\checkmarki\checkmarkr\checkmarkc\checkmark$\checkmarke\checkmark$\checkmark^\checkmark\\checkmarkc\checkmarki\checkmarkr\checkmarkc\checkmark$\checkmarks\checkmark$\checkmark^\checkmark\\checkmarkc\checkmarki\checkmarkr\checkmarkc\checkmark$\checkmark,\checkmark$\checkmark^\checkmark\\checkmarkc\checkmarki\checkmarkr\checkmarkc\checkmark$\checkmark \checkmark$\checkmark^\checkmark\\checkmarkc\checkmarki\checkmarkr\checkmarkc\checkmark$\checkmarkd\checkmark$\checkmark^\checkmark\\checkmarkc\checkmarki\checkmarkr\checkmarkc\checkmark$\checkmarki\checkmark$\checkmark^\checkmark\\checkmarkc\checkmarki\checkmarkr\checkmarkc\checkmark$\checkmarka\checkmark$\checkmark^\checkmark\\checkmarkc\checkmarki\checkmarkr\checkmarkc\checkmark$\checkmarkg\checkmark$\checkmark^\checkmark\\checkmarkc\checkmarki\checkmarkr\checkmarkc\checkmark$\checkmarkr\checkmark$\checkmark^\checkmark\\checkmarkc\checkmarki\checkmarkr\checkmarkc\checkmark$\checkmarka\checkmark$\checkmark^\checkmark\\checkmarkc\checkmarki\checkmarkr\checkmarkc\checkmark$\checkmarkm\checkmark$\checkmark^\checkmark\\checkmarkc\checkmarki\checkmarkr\checkmarkc\checkmark$\checkmarkm\checkmark$\checkmark^\checkmark\\checkmarkc\checkmarki\checkmarkr\checkmarkc\checkmark$\checkmarke\checkmark$\checkmark^\checkmark\\checkmarkc\checkmarki\checkmarkr\checkmarkc\checkmark$\checkmark \checkmark$\checkmark^\checkmark\\checkmarkc\checkmarki\checkmarkr\checkmarkc\checkmark$\checkmarkp\checkmark$\checkmark^\checkmark\\checkmarkc\checkmarki\checkmarkr\checkmarkc\checkmark$\checkmarki\checkmark$\checkmark^\checkmark\\checkmarkc\checkmarki\checkmarkr\checkmarkc\checkmark$\checkmarke\checkmark$\checkmark^\checkmark\\checkmarkc\checkmarki\checkmarkr\checkmarkc\checkmark$\checkmarku\checkmark$\checkmark^\checkmark\\checkmarkc\checkmarki\checkmarkr\checkmarkc\checkmark$\checkmarkv\checkmark$\checkmark^\checkmark\\checkmarkc\checkmarki\checkmarkr\checkmarkc\checkmark$\checkmarkr\checkmark$\checkmark^\checkmark\\checkmarkc\checkmarki\checkmarkr\checkmarkc\checkmark$\checkmarke\checkmark$\checkmark^\checkmark\\checkmarkc\checkmarki\checkmarkr\checkmarkc\checkmark$\checkmark \checkmark$\checkmark^\checkmark\\checkmarkc\checkmarki\checkmarkr\checkmarkc\checkmark$\checkmarke\checkmark$\checkmark^\checkmark\\checkmarkc\checkmarki\checkmarkr\checkmarkc\checkmark$\checkmarkt\checkmark$\checkmark^\checkmark\\checkmarkc\checkmarki\checkmarkr\checkmarkc\checkmark$\checkmark \checkmark$\checkmark^\checkmark\\checkmarkc\checkmarki\checkmarkr\checkmarkc\checkmark$\checkmarkF\checkmark$\checkmark^\checkmark\\checkmarkc\checkmarki\checkmarkr\checkmarkc\checkmark$\checkmarkA\checkmark$\checkmark^\checkmark\\checkmarkc\checkmarki\checkmarkr\checkmarkc\checkmark$\checkmarkS\checkmark$\checkmark^\checkmark\\checkmarkc\checkmarki\checkmarkr\checkmarkc\checkmark$\checkmarkT\checkmark$\checkmark^\checkmark\\checkmarkc\checkmarki\checkmarkr\checkmarkc\checkmark$\checkmark \checkmark$\checkmark^\checkmark\\checkmarkc\checkmarki\checkmarkr\checkmarkc\checkmark$\checkmark;\checkmark$\checkmark^\checkmark\\checkmarkc\checkmarki\checkmarkr\checkmarkc\checkmark$\checkmark
\checkmark$\checkmark^\checkmark\\checkmarkc\checkmarki\checkmarkr\checkmarkc\checkmark$\checkmark \checkmark$\checkmark^\checkmark\\checkmarkc\checkmarki\checkmarkr\checkmarkc\checkmark$\checkmark \checkmark$\checkmark^\checkmark\\checkmarkc\checkmarki\checkmarkr\checkmarkc\checkmark$\checkmark \checkmark$\checkmark^\checkmark\\checkmarkc\checkmarki\checkmarkr\checkmarkc\checkmark$\checkmark \checkmark$\checkmark^\checkmark\\checkmarkc\checkmarki\checkmarkr\checkmarkc\checkmark$\checkmark\\checkmark$\checkmark^\checkmark\\checkmarkc\checkmarki\checkmarkr\checkmarkc\checkmark$\checkmarki\checkmark$\checkmark^\checkmark\\checkmarkc\checkmarki\checkmarkr\checkmarkc\checkmark$\checkmarkt\checkmark$\checkmark^\checkmark\\checkmarkc\checkmarki\checkmarkr\checkmarkc\checkmark$\checkmarke\checkmark$\checkmark^\checkmark\\checkmarkc\checkmarki\checkmarkr\checkmarkc\checkmark$\checkmarkm\checkmark$\checkmark^\checkmark\\checkmarkc\checkmarki\checkmarkr\checkmarkc\checkmark$\checkmark \checkmark$\checkmark^\checkmark\\checkmarkc\checkmarki\checkmarkr\checkmarkc\checkmark$\checkmarkL\checkmark$\checkmark^\checkmark\\checkmarkc\checkmarki\checkmarkr\checkmarkc\checkmark$\checkmarke\checkmark$\checkmark^\checkmark\\checkmarkc\checkmarki\checkmarkr\checkmarkc\checkmark$\checkmark \checkmark$\checkmark^\checkmark\\checkmarkc\checkmarki\checkmarkr\checkmarkc\checkmark$\checkmarkc\checkmark$\checkmark^\checkmark\\checkmarkc\checkmarki\checkmarkr\checkmarkc\checkmark$\checkmarkh\checkmark$\checkmark^\checkmark\\checkmarkc\checkmarki\checkmarkr\checkmarkc\checkmark$\checkmarko\checkmark$\checkmark^\checkmark\\checkmarkc\checkmarki\checkmarkr\checkmarkc\checkmark$\checkmarki\checkmark$\checkmark^\checkmark\\checkmarkc\checkmarki\checkmarkr\checkmarkc\checkmark$\checkmarkx\checkmark$\checkmark^\checkmark\\checkmarkc\checkmarki\checkmarkr\checkmarkc\checkmark$\checkmark \checkmark$\checkmark^\checkmark\\checkmarkc\checkmarki\checkmarkr\checkmarkc\checkmark$\checkmarkr\checkmark$\checkmark^\checkmark\\checkmarkc\checkmarki\checkmarkr\checkmarkc\checkmark$\checkmarka\checkmark$\checkmark^\checkmark\\checkmarkc\checkmarki\checkmarkr\checkmarkc\checkmark$\checkmarki\checkmark$\checkmark^\checkmark\\checkmarkc\checkmarki\checkmarkr\checkmarkc\checkmark$\checkmarks\checkmark$\checkmark^\checkmark\\checkmarkc\checkmarki\checkmarkr\checkmarkc\checkmark$\checkmarko\checkmark$\checkmark^\checkmark\\checkmarkc\checkmarki\checkmarkr\checkmarkc\checkmark$\checkmarkn\checkmark$\checkmark^\checkmark\\checkmarkc\checkmarki\checkmarkr\checkmarkc\checkmark$\checkmarkn\checkmark$\checkmark^\checkmark\\checkmarkc\checkmarki\checkmarkr\checkmarkc\checkmark$\checkmarké\checkmark$\checkmark^\checkmark\\checkmarkc\checkmarki\checkmarkr\checkmarkc\checkmark$\checkmark \checkmark$\checkmark^\checkmark\\checkmarkc\checkmarki\checkmarkr\checkmarkc\checkmark$\checkmarke\checkmark$\checkmark^\checkmark\\checkmarkc\checkmarki\checkmarkr\checkmarkc\checkmark$\checkmarkt\checkmark$\checkmark^\checkmark\\checkmarkc\checkmarki\checkmarkr\checkmarkc\checkmark$\checkmark \checkmark$\checkmark^\checkmark\\checkmarkc\checkmarki\checkmarkr\checkmarkc\checkmark$\checkmarkc\checkmark$\checkmark^\checkmark\\checkmarkc\checkmarki\checkmarkr\checkmarkc\checkmark$\checkmarko\checkmark$\checkmark^\checkmark\\checkmarkc\checkmarki\checkmarkr\checkmarkc\checkmark$\checkmarkm\checkmark$\checkmark^\checkmark\\checkmarkc\checkmarki\checkmarkr\checkmarkc\checkmark$\checkmarkp\checkmark$\checkmark^\checkmark\\checkmarkc\checkmarki\checkmarkr\checkmarkc\checkmark$\checkmarka\checkmark$\checkmark^\checkmark\\checkmarkc\checkmarki\checkmarkr\checkmarkc\checkmark$\checkmarkr\checkmark$\checkmark^\checkmark\\checkmarkc\checkmarki\checkmarkr\checkmarkc\checkmark$\checkmarka\checkmark$\checkmark^\checkmark\\checkmarkc\checkmarki\checkmarkr\checkmarkc\checkmark$\checkmarkt\checkmark$\checkmark^\checkmark\\checkmarkc\checkmarki\checkmarkr\checkmarkc\checkmark$\checkmarki\checkmark$\checkmark^\checkmark\\checkmarkc\checkmarki\checkmarkr\checkmarkc\checkmark$\checkmarkf\checkmark$\checkmark^\checkmark\\checkmarkc\checkmarki\checkmarkr\checkmarkc\checkmark$\checkmark \checkmark$\checkmark^\checkmark\\checkmarkc\checkmarki\checkmarkr\checkmarkc\checkmark$\checkmarkd\checkmark$\checkmark^\checkmark\\checkmarkc\checkmarki\checkmarkr\checkmarkc\checkmark$\checkmarke\checkmark$\checkmark^\checkmark\\checkmarkc\checkmarki\checkmarkr\checkmarkc\checkmark$\checkmarks\checkmark$\checkmark^\checkmark\\checkmarkc\checkmarki\checkmarkr\checkmarkc\checkmark$\checkmark \checkmark$\checkmark^\checkmark\\checkmarkc\checkmarki\checkmarkr\checkmarkc\checkmark$\checkmarko\checkmark$\checkmark^\checkmark\\checkmarkc\checkmarki\checkmarkr\checkmarkc\checkmark$\checkmarkr\checkmark$\checkmark^\checkmark\\checkmarkc\checkmarki\checkmarkr\checkmarkc\checkmark$\checkmarkg\checkmark$\checkmark^\checkmark\\checkmarkc\checkmarki\checkmarkr\checkmarkc\checkmark$\checkmarka\checkmark$\checkmark^\checkmark\\checkmarkc\checkmarki\checkmarkr\checkmarkc\checkmark$\checkmarkn\checkmark$\checkmark^\checkmark\\checkmarkc\checkmarki\checkmarkr\checkmarkc\checkmark$\checkmarke\checkmark$\checkmark^\checkmark\\checkmarkc\checkmarki\checkmarkr\checkmarkc\checkmark$\checkmarks\checkmark$\checkmark^\checkmark\\checkmarkc\checkmarki\checkmarkr\checkmarkc\checkmark$\checkmark \checkmark$\checkmark^\checkmark\\checkmarkc\checkmarki\checkmarkr\checkmarkc\checkmark$\checkmarkm\checkmark$\checkmark^\checkmark\\checkmarkc\checkmarki\checkmarkr\checkmarkc\checkmark$\checkmarké\checkmark$\checkmark^\checkmark\\checkmarkc\checkmarki\checkmarkr\checkmarkc\checkmark$\checkmarkc\checkmark$\checkmark^\checkmark\\checkmarkc\checkmarki\checkmarkr\checkmarkc\checkmark$\checkmarka\checkmark$\checkmark^\checkmark\\checkmarkc\checkmarki\checkmarkr\checkmarkc\checkmark$\checkmarkn\checkmark$\checkmark^\checkmark\\checkmarkc\checkmarki\checkmarkr\checkmarkc\checkmark$\checkmarki\checkmark$\checkmark^\checkmark\\checkmarkc\checkmarki\checkmarkr\checkmarkc\checkmark$\checkmarkq\checkmark$\checkmark^\checkmark\\checkmarkc\checkmarki\checkmarkr\checkmarkc\checkmark$\checkmarku\checkmark$\checkmark^\checkmark\\checkmarkc\checkmarki\checkmarkr\checkmarkc\checkmark$\checkmarke\checkmark$\checkmark^\checkmark\\checkmarkc\checkmarki\checkmarkr\checkmarkc\checkmark$\checkmarks\checkmark$\checkmark^\checkmark\\checkmarkc\checkmarki\checkmarkr\checkmarkc\checkmark$\checkmark \checkmark$\checkmark^\checkmark\\checkmarkc\checkmarki\checkmarkr\checkmarkc\checkmark$\checkmark;\checkmark$\checkmark^\checkmark\\checkmarkc\checkmarki\checkmarkr\checkmarkc\checkmark$\checkmark
\checkmark$\checkmark^\checkmark\\checkmarkc\checkmarki\checkmarkr\checkmarkc\checkmark$\checkmark \checkmark$\checkmark^\checkmark\\checkmarkc\checkmarki\checkmarkr\checkmarkc\checkmark$\checkmark \checkmark$\checkmark^\checkmark\\checkmarkc\checkmarki\checkmarkr\checkmarkc\checkmark$\checkmark \checkmark$\checkmark^\checkmark\\checkmarkc\checkmarki\checkmarkr\checkmarkc\checkmark$\checkmark \checkmark$\checkmark^\checkmark\\checkmarkc\checkmarki\checkmarkr\checkmarkc\checkmark$\checkmark\\checkmark$\checkmark^\checkmark\\checkmarkc\checkmarki\checkmarkr\checkmarkc\checkmark$\checkmarki\checkmark$\checkmark^\checkmark\\checkmarkc\checkmarki\checkmarkr\checkmarkc\checkmark$\checkmarkt\checkmark$\checkmark^\checkmark\\checkmarkc\checkmarki\checkmarkr\checkmarkc\checkmark$\checkmarke\checkmark$\checkmark^\checkmark\\checkmarkc\checkmarki\checkmarkr\checkmarkc\checkmark$\checkmarkm\checkmark$\checkmark^\checkmark\\checkmarkc\checkmarki\checkmarkr\checkmarkc\checkmark$\checkmark \checkmark$\checkmark^\checkmark\\checkmarkc\checkmarki\checkmarkr\checkmarkc\checkmark$\checkmarkL\checkmark$\checkmark^\checkmark\\checkmarkc\checkmarki\checkmarkr\checkmarkc\checkmark$\checkmarke\checkmark$\checkmark^\checkmark\\checkmarkc\checkmarki\checkmarkr\checkmarkc\checkmark$\checkmark \checkmark$\checkmark^\checkmark\\checkmarkc\checkmarki\checkmarkr\checkmarkc\checkmark$\checkmarkd\checkmark$\checkmark^\checkmark\\checkmarkc\checkmarki\checkmarkr\checkmarkc\checkmark$\checkmarki\checkmark$\checkmark^\checkmark\\checkmarkc\checkmarki\checkmarkr\checkmarkc\checkmark$\checkmarkm\checkmark$\checkmark^\checkmark\\checkmarkc\checkmarki\checkmarkr\checkmarkc\checkmark$\checkmarke\checkmark$\checkmark^\checkmark\\checkmarkc\checkmarki\checkmarkr\checkmarkc\checkmark$\checkmarkn\checkmark$\checkmark^\checkmark\\checkmarkc\checkmarki\checkmarkr\checkmarkc\checkmark$\checkmarks\checkmark$\checkmark^\checkmark\\checkmarkc\checkmarki\checkmarkr\checkmarkc\checkmark$\checkmarki\checkmark$\checkmark^\checkmark\\checkmarkc\checkmarki\checkmarkr\checkmarkc\checkmark$\checkmarko\checkmark$\checkmark^\checkmark\\checkmarkc\checkmarki\checkmarkr\checkmarkc\checkmark$\checkmarkn\checkmark$\checkmark^\checkmark\\checkmarkc\checkmarki\checkmarkr\checkmarkc\checkmark$\checkmarkn\checkmark$\checkmark^\checkmark\\checkmarkc\checkmarki\checkmarkr\checkmarkc\checkmark$\checkmarke\checkmark$\checkmark^\checkmark\\checkmarkc\checkmarki\checkmarkr\checkmarkc\checkmark$\checkmarkm\checkmark$\checkmark^\checkmark\\checkmarkc\checkmarki\checkmarkr\checkmarkc\checkmark$\checkmarke\checkmark$\checkmark^\checkmark\\checkmarkc\checkmarki\checkmarkr\checkmarkc\checkmark$\checkmarkn\checkmark$\checkmark^\checkmark\\checkmarkc\checkmarki\checkmarkr\checkmarkc\checkmark$\checkmarkt\checkmark$\checkmark^\checkmark\\checkmarkc\checkmarki\checkmarkr\checkmarkc\checkmark$\checkmark \checkmark$\checkmark^\checkmark\\checkmarkc\checkmarki\checkmarkr\checkmarkc\checkmark$\checkmarka\checkmark$\checkmark^\checkmark\\checkmarkc\checkmarki\checkmarkr\checkmarkc\checkmark$\checkmarkn\checkmark$\checkmark^\checkmark\\checkmarkc\checkmarki\checkmarkr\checkmarkc\checkmark$\checkmarka\checkmark$\checkmark^\checkmark\\checkmarkc\checkmarki\checkmarkr\checkmarkc\checkmark$\checkmarkl\checkmark$\checkmark^\checkmark\\checkmarkc\checkmarki\checkmarkr\checkmarkc\checkmark$\checkmarky\checkmark$\checkmark^\checkmark\\checkmarkc\checkmarki\checkmarkr\checkmarkc\checkmark$\checkmarkt\checkmark$\checkmark^\checkmark\\checkmarkc\checkmarki\checkmarkr\checkmarkc\checkmark$\checkmarki\checkmark$\checkmark^\checkmark\\checkmarkc\checkmarki\checkmarkr\checkmarkc\checkmark$\checkmarkq\checkmark$\checkmark^\checkmark\\checkmarkc\checkmarki\checkmarkr\checkmarkc\checkmark$\checkmarku\checkmark$\checkmark^\checkmark\\checkmarkc\checkmarki\checkmarkr\checkmarkc\checkmark$\checkmarke\checkmark$\checkmark^\checkmark\\checkmarkc\checkmarki\checkmarkr\checkmarkc\checkmark$\checkmark \checkmark$\checkmark^\checkmark\\checkmarkc\checkmarki\checkmarkr\checkmarkc\checkmark$\checkmarkc\checkmark$\checkmark^\checkmark\\checkmarkc\checkmarki\checkmarkr\checkmarkc\checkmark$\checkmarko\checkmark$\checkmark^\checkmark\\checkmarkc\checkmarki\checkmarkr\checkmarkc\checkmark$\checkmarkm\checkmark$\checkmark^\checkmark\\checkmarkc\checkmarki\checkmarkr\checkmarkc\checkmark$\checkmarkp\checkmark$\checkmark^\checkmark\\checkmarkc\checkmarki\checkmarkr\checkmarkc\checkmark$\checkmarkl\checkmark$\checkmark^\checkmark\\checkmarkc\checkmarki\checkmarkr\checkmarkc\checkmark$\checkmarke\checkmark$\checkmark^\checkmark\\checkmarkc\checkmarki\checkmarkr\checkmarkc\checkmark$\checkmarkt\checkmark$\checkmark^\checkmark\\checkmarkc\checkmarki\checkmarkr\checkmarkc\checkmark$\checkmark \checkmark$\checkmark^\checkmark\\checkmarkc\checkmarki\checkmarkr\checkmarkc\checkmark$\checkmarkd\checkmark$\checkmark^\checkmark\\checkmarkc\checkmarki\checkmarkr\checkmarkc\checkmark$\checkmarke\checkmark$\checkmark^\checkmark\\checkmarkc\checkmarki\checkmarkr\checkmarkc\checkmark$\checkmarks\checkmark$\checkmark^\checkmark\\checkmarkc\checkmarki\checkmarkr\checkmarkc\checkmark$\checkmark \checkmark$\checkmark^\checkmark\\checkmarkc\checkmarki\checkmarkr\checkmarkc\checkmark$\checkmarko\checkmark$\checkmark^\checkmark\\checkmarkc\checkmarki\checkmarkr\checkmarkc\checkmark$\checkmarkr\checkmark$\checkmark^\checkmark\\checkmarkc\checkmarki\checkmarkr\checkmarkc\checkmark$\checkmarkg\checkmark$\checkmark^\checkmark\\checkmarkc\checkmarki\checkmarkr\checkmarkc\checkmark$\checkmarka\checkmark$\checkmark^\checkmark\\checkmarkc\checkmarki\checkmarkr\checkmarkc\checkmark$\checkmarkn\checkmark$\checkmark^\checkmark\\checkmarkc\checkmarki\checkmarkr\checkmarkc\checkmark$\checkmarke\checkmark$\checkmark^\checkmark\\checkmarkc\checkmarki\checkmarkr\checkmarkc\checkmark$\checkmarks\checkmark$\checkmark^\checkmark\\checkmarkc\checkmarki\checkmarkr\checkmarkc\checkmark$\checkmark \checkmark$\checkmark^\checkmark\\checkmarkc\checkmarki\checkmarkr\checkmarkc\checkmark$\checkmarkd\checkmark$\checkmark^\checkmark\\checkmarkc\checkmarki\checkmarkr\checkmarkc\checkmark$\checkmarke\checkmark$\checkmark^\checkmark\\checkmarkc\checkmarki\checkmarkr\checkmarkc\checkmark$\checkmark \checkmark$\checkmark^\checkmark\\checkmarkc\checkmarki\checkmarkr\checkmarkc\checkmark$\checkmarkt\checkmark$\checkmark^\checkmark\\checkmarkc\checkmarki\checkmarkr\checkmarkc\checkmark$\checkmarkr\checkmark$\checkmark^\checkmark\\checkmarkc\checkmarki\checkmarkr\checkmarkc\checkmark$\checkmarka\checkmark$\checkmark^\checkmark\\checkmarkc\checkmarki\checkmarkr\checkmarkc\checkmark$\checkmarkn\checkmark$\checkmark^\checkmark\\checkmarkc\checkmarki\checkmarkr\checkmarkc\checkmark$\checkmarks\checkmark$\checkmark^\checkmark\\checkmarkc\checkmarki\checkmarkr\checkmarkc\checkmark$\checkmarkl\checkmark$\checkmark^\checkmark\\checkmarkc\checkmarki\checkmarkr\checkmarkc\checkmark$\checkmarka\checkmark$\checkmark^\checkmark\\checkmarkc\checkmarki\checkmarkr\checkmarkc\checkmark$\checkmarkt\checkmark$\checkmark^\checkmark\\checkmarkc\checkmarki\checkmarkr\checkmarkc\checkmark$\checkmarki\checkmark$\checkmark^\checkmark\\checkmarkc\checkmarki\checkmarkr\checkmarkc\checkmark$\checkmarko\checkmark$\checkmark^\checkmark\\checkmarkc\checkmarki\checkmarkr\checkmarkc\checkmark$\checkmarkn\checkmark$\checkmark^\checkmark\\checkmarkc\checkmarki\checkmarkr\checkmarkc\checkmark$\checkmark \checkmark$\checkmark^\checkmark\\checkmarkc\checkmarki\checkmarkr\checkmarkc\checkmark$\checkmark(\checkmark$\checkmark^\checkmark\\checkmarkc\checkmarki\checkmarkr\checkmarkc\checkmark$\checkmarkv\checkmark$\checkmark^\checkmark\\checkmarkc\checkmarki\checkmarkr\checkmarkc\checkmark$\checkmarki\checkmark$\checkmark^\checkmark\\checkmarkc\checkmarki\checkmarkr\checkmarkc\checkmark$\checkmarks\checkmark$\checkmark^\checkmark\\checkmarkc\checkmarki\checkmarkr\checkmarkc\checkmark$\checkmark \checkmark$\checkmark^\checkmark\\checkmarkc\checkmarki\checkmarkr\checkmarkc\checkmark$\checkmarkà\checkmark$\checkmark^\checkmark\\checkmarkc\checkmarki\checkmarkr\checkmarkc\checkmark$\checkmark \checkmark$\checkmark^\checkmark\\checkmarkc\checkmarki\checkmarkr\checkmarkc\checkmark$\checkmarkb\checkmark$\checkmark^\checkmark\\checkmarkc\checkmarki\checkmarkr\checkmarkc\checkmark$\checkmarki\checkmark$\checkmark^\checkmark\\checkmarkc\checkmarki\checkmarkr\checkmarkc\checkmark$\checkmarkl\checkmark$\checkmark^\checkmark\\checkmarkc\checkmarki\checkmarkr\checkmarkc\checkmark$\checkmarkl\checkmark$\checkmark^\checkmark\\checkmarkc\checkmarki\checkmarkr\checkmarkc\checkmark$\checkmarke\checkmark$\checkmark^\checkmark\\checkmarkc\checkmarki\checkmarkr\checkmarkc\checkmark$\checkmarks\checkmark$\checkmark^\checkmark\\checkmarkc\checkmarki\checkmarkr\checkmarkc\checkmark$\checkmark)\checkmark$\checkmark^\checkmark\\checkmarkc\checkmarki\checkmarkr\checkmarkc\checkmark$\checkmark \checkmark$\checkmark^\checkmark\\checkmarkc\checkmarki\checkmarkr\checkmarkc\checkmark$\checkmarke\checkmark$\checkmark^\checkmark\\checkmarkc\checkmarki\checkmarkr\checkmarkc\checkmark$\checkmarkt\checkmark$\checkmark^\checkmark\\checkmarkc\checkmarki\checkmarkr\checkmarkc\checkmark$\checkmark \checkmark$\checkmark^\checkmark\\checkmarkc\checkmarki\checkmarkr\checkmarkc\checkmark$\checkmarkd\checkmark$\checkmark^\checkmark\\checkmarkc\checkmarki\checkmarkr\checkmarkc\checkmark$\checkmarke\checkmark$\checkmark^\checkmark\\checkmarkc\checkmarki\checkmarkr\checkmarkc\checkmark$\checkmark \checkmark$\checkmark^\checkmark\\checkmarkc\checkmarki\checkmarkr\checkmarkc\checkmark$\checkmarkr\checkmark$\checkmark^\checkmark\\checkmarkc\checkmarki\checkmarkr\checkmarkc\checkmark$\checkmarko\checkmark$\checkmark^\checkmark\\checkmarkc\checkmarki\checkmarkr\checkmarkc\checkmark$\checkmarkt\checkmark$\checkmark^\checkmark\\checkmarkc\checkmarki\checkmarkr\checkmarkc\checkmark$\checkmarka\checkmark$\checkmark^\checkmark\\checkmarkc\checkmarki\checkmarkr\checkmarkc\checkmark$\checkmarkt\checkmark$\checkmark^\checkmark\\checkmarkc\checkmarki\checkmarkr\checkmarkc\checkmark$\checkmarki\checkmark$\checkmark^\checkmark\\checkmarkc\checkmarki\checkmarkr\checkmarkc\checkmark$\checkmarko\checkmark$\checkmark^\checkmark\\checkmarkc\checkmarki\checkmarkr\checkmarkc\checkmark$\checkmarkn\checkmark$\checkmark^\checkmark\\checkmarkc\checkmarki\checkmarkr\checkmarkc\checkmark$\checkmark \checkmark$\checkmark^\checkmark\\checkmarkc\checkmarki\checkmarkr\checkmarkc\checkmark$\checkmark(\checkmark$\checkmark^\checkmark\\checkmarkc\checkmarki\checkmarkr\checkmarkc\checkmark$\checkmarka\checkmark$\checkmark^\checkmark\\checkmarkc\checkmarki\checkmarkr\checkmarkc\checkmark$\checkmarkr\checkmark$\checkmark^\checkmark\\checkmarkc\checkmarki\checkmarkr\checkmarkc\checkmark$\checkmarkb\checkmark$\checkmark^\checkmark\\checkmarkc\checkmarki\checkmarkr\checkmarkc\checkmark$\checkmarkr\checkmark$\checkmark^\checkmark\\checkmarkc\checkmarki\checkmarkr\checkmarkc\checkmark$\checkmarke\checkmark$\checkmark^\checkmark\\checkmarkc\checkmarki\checkmarkr\checkmarkc\checkmark$\checkmark \checkmark$\checkmark^\checkmark\\checkmarkc\checkmarki\checkmarkr\checkmarkc\checkmark$\checkmarkd\checkmark$\checkmark^\checkmark\\checkmarkc\checkmarki\checkmarkr\checkmarkc\checkmark$\checkmarke\checkmark$\checkmark^\checkmark\\checkmarkc\checkmarki\checkmarkr\checkmarkc\checkmark$\checkmark \checkmark$\checkmark^\checkmark\\checkmarkc\checkmarki\checkmarkr\checkmarkc\checkmark$\checkmarkp\checkmark$\checkmark^\checkmark\\checkmarkc\checkmarki\checkmarkr\checkmarkc\checkmark$\checkmarki\checkmark$\checkmark^\checkmark\\checkmarkc\checkmarki\checkmarkr\checkmarkc\checkmark$\checkmarkv\checkmark$\checkmark^\checkmark\\checkmarkc\checkmarki\checkmarkr\checkmarkc\checkmark$\checkmarko\checkmark$\checkmark^\checkmark\\checkmarkc\checkmarki\checkmarkr\checkmarkc\checkmark$\checkmarkt\checkmark$\checkmark^\checkmark\\checkmarkc\checkmarki\checkmarkr\checkmarkc\checkmark$\checkmarke\checkmark$\checkmark^\checkmark\\checkmarkc\checkmarki\checkmarkr\checkmarkc\checkmark$\checkmarkm\checkmark$\checkmark^\checkmark\\checkmarkc\checkmarki\checkmarkr\checkmarkc\checkmark$\checkmarke\checkmark$\checkmark^\checkmark\\checkmarkc\checkmarki\checkmarkr\checkmarkc\checkmark$\checkmarkn\checkmark$\checkmark^\checkmark\\checkmarkc\checkmarki\checkmarkr\checkmarkc\checkmark$\checkmarkt\checkmark$\checkmark^\checkmark\\checkmarkc\checkmarki\checkmarkr\checkmarkc\checkmark$\checkmark,\checkmark$\checkmark^\checkmark\\checkmarkc\checkmarki\checkmarkr\checkmarkc\checkmark$\checkmark \checkmark$\checkmark^\checkmark\\checkmarkc\checkmarki\checkmarkr\checkmarkc\checkmark$\checkmarkt\checkmark$\checkmark^\checkmark\\checkmarkc\checkmarki\checkmarkr\checkmarkc\checkmark$\checkmarkr\checkmark$\checkmark^\checkmark\\checkmarkc\checkmarki\checkmarkr\checkmarkc\checkmark$\checkmarka\checkmark$\checkmark^\checkmark\\checkmarkc\checkmarki\checkmarkr\checkmarkc\checkmark$\checkmarkn\checkmark$\checkmark^\checkmark\\checkmarkc\checkmarki\checkmarkr\checkmarkc\checkmark$\checkmarks\checkmark$\checkmark^\checkmark\\checkmarkc\checkmarki\checkmarkr\checkmarkc\checkmark$\checkmarkm\checkmark$\checkmark^\checkmark\\checkmarkc\checkmarki\checkmarkr\checkmarkc\checkmark$\checkmarki\checkmark$\checkmark^\checkmark\\checkmarkc\checkmarki\checkmarkr\checkmarkc\checkmark$\checkmarks\checkmark$\checkmark^\checkmark\\checkmarkc\checkmarki\checkmarkr\checkmarkc\checkmark$\checkmarks\checkmark$\checkmark^\checkmark\\checkmarkc\checkmarki\checkmarkr\checkmarkc\checkmark$\checkmarki\checkmark$\checkmark^\checkmark\\checkmarkc\checkmarki\checkmarkr\checkmarkc\checkmark$\checkmarko\checkmark$\checkmark^\checkmark\\checkmarkc\checkmarki\checkmarkr\checkmarkc\checkmark$\checkmarkn\checkmark$\checkmark^\checkmark\\checkmarkc\checkmarki\checkmarkr\checkmarkc\checkmark$\checkmarks\checkmark$\checkmark^\checkmark\\checkmarkc\checkmarki\checkmarkr\checkmarkc\checkmark$\checkmark)\checkmark$\checkmark^\checkmark\\checkmarkc\checkmarki\checkmarkr\checkmarkc\checkmark$\checkmark \checkmark$\checkmark^\checkmark\\checkmarkc\checkmarki\checkmarkr\checkmarkc\checkmark$\checkmark;\checkmark$\checkmark^\checkmark\\checkmarkc\checkmarki\checkmarkr\checkmarkc\checkmark$\checkmark
\checkmark$\checkmark^\checkmark\\checkmarkc\checkmarki\checkmarkr\checkmarkc\checkmark$\checkmark \checkmark$\checkmark^\checkmark\\checkmarkc\checkmarki\checkmarkr\checkmarkc\checkmark$\checkmark \checkmark$\checkmark^\checkmark\\checkmarkc\checkmarki\checkmarkr\checkmarkc\checkmark$\checkmark \checkmark$\checkmark^\checkmark\\checkmarkc\checkmarki\checkmarkr\checkmarkc\checkmark$\checkmark \checkmark$\checkmark^\checkmark\\checkmarkc\checkmarki\checkmarkr\checkmarkc\checkmark$\checkmark\\checkmark$\checkmark^\checkmark\\checkmarkc\checkmarki\checkmarkr\checkmarkc\checkmark$\checkmarki\checkmark$\checkmark^\checkmark\\checkmarkc\checkmarki\checkmarkr\checkmarkc\checkmark$\checkmarkt\checkmark$\checkmark^\checkmark\\checkmarkc\checkmarki\checkmarkr\checkmarkc\checkmark$\checkmarke\checkmark$\checkmark^\checkmark\\checkmarkc\checkmarki\checkmarkr\checkmarkc\checkmark$\checkmarkm\checkmark$\checkmark^\checkmark\\checkmarkc\checkmarki\checkmarkr\checkmarkc\checkmark$\checkmark \checkmark$\checkmark^\checkmark\\checkmarkc\checkmarki\checkmarkr\checkmarkc\checkmark$\checkmarkL\checkmark$\checkmark^\checkmark\\checkmarkc\checkmarki\checkmarkr\checkmarkc\checkmark$\checkmarka\checkmark$\checkmark^\checkmark\\checkmarkc\checkmarki\checkmarkr\checkmarkc\checkmark$\checkmark \checkmark$\checkmark^\checkmark\\checkmarkc\checkmarki\checkmarkr\checkmarkc\checkmark$\checkmarkv\checkmark$\checkmark^\checkmark\\checkmarkc\checkmarki\checkmarkr\checkmarkc\checkmark$\checkmarké\checkmark$\checkmark^\checkmark\\checkmarkc\checkmarki\checkmarkr\checkmarkc\checkmark$\checkmarkr\checkmark$\checkmark^\checkmark\\checkmarkc\checkmarki\checkmarkr\checkmarkc\checkmark$\checkmarki\checkmark$\checkmark^\checkmark\\checkmarkc\checkmarki\checkmarkr\checkmarkc\checkmark$\checkmarkf\checkmark$\checkmark^\checkmark\\checkmarkc\checkmarki\checkmarkr\checkmarkc\checkmark$\checkmarki\checkmark$\checkmark^\checkmark\\checkmarkc\checkmarki\checkmarkr\checkmarkc\checkmark$\checkmarkc\checkmark$\checkmark^\checkmark\\checkmarkc\checkmarki\checkmarkr\checkmarkc\checkmark$\checkmarka\checkmark$\checkmark^\checkmark\\checkmarkc\checkmarki\checkmarkr\checkmarkc\checkmark$\checkmarkt\checkmark$\checkmark^\checkmark\\checkmarkc\checkmarki\checkmarkr\checkmarkc\checkmark$\checkmarki\checkmark$\checkmark^\checkmark\\checkmarkc\checkmarki\checkmarkr\checkmarkc\checkmark$\checkmarko\checkmark$\checkmark^\checkmark\\checkmarkc\checkmarki\checkmarkr\checkmarkc\checkmark$\checkmarkn\checkmark$\checkmark^\checkmark\\checkmarkc\checkmarki\checkmarkr\checkmarkc\checkmark$\checkmark \checkmark$\checkmark^\checkmark\\checkmarkc\checkmarki\checkmarkr\checkmarkc\checkmark$\checkmarkd\checkmark$\checkmark^\checkmark\\checkmarkc\checkmarki\checkmarkr\checkmarkc\checkmark$\checkmarke\checkmark$\checkmark^\checkmark\\checkmarkc\checkmarki\checkmarkr\checkmarkc\checkmark$\checkmark \checkmark$\checkmark^\checkmark\\checkmarkc\checkmarki\checkmarkr\checkmarkc\checkmark$\checkmarkl\checkmark$\checkmark^\checkmark\\checkmarkc\checkmarki\checkmarkr\checkmarkc\checkmark$\checkmarka\checkmark$\checkmark^\checkmark\\checkmarkc\checkmarki\checkmarkr\checkmarkc\checkmark$\checkmark \checkmark$\checkmark^\checkmark\\checkmarkc\checkmarki\checkmarkr\checkmarkc\checkmark$\checkmarkr\checkmark$\checkmark^\checkmark\\checkmarkc\checkmarki\checkmarkr\checkmarkc\checkmark$\checkmarké\checkmark$\checkmark^\checkmark\\checkmarkc\checkmarki\checkmarkr\checkmarkc\checkmark$\checkmarks\checkmark$\checkmark^\checkmark\\checkmarkc\checkmarki\checkmarkr\checkmarkc\checkmark$\checkmarki\checkmark$\checkmark^\checkmark\\checkmarkc\checkmarki\checkmarkr\checkmarkc\checkmark$\checkmarks\checkmark$\checkmark^\checkmark\\checkmarkc\checkmarki\checkmarkr\checkmarkc\checkmark$\checkmarkt\checkmark$\checkmark^\checkmark\\checkmarkc\checkmarki\checkmarkr\checkmarkc\checkmark$\checkmarka\checkmark$\checkmark^\checkmark\\checkmarkc\checkmarki\checkmarkr\checkmarkc\checkmark$\checkmarkn\checkmark$\checkmark^\checkmark\\checkmarkc\checkmarki\checkmarkr\checkmarkc\checkmark$\checkmarkc\checkmark$\checkmark^\checkmark\\checkmarkc\checkmarki\checkmarkr\checkmarkc\checkmark$\checkmarke\checkmark$\checkmark^\checkmark\\checkmarkc\checkmarki\checkmarkr\checkmarkc\checkmark$\checkmark \checkmark$\checkmark^\checkmark\\checkmarkc\checkmarki\checkmarkr\checkmarkc\checkmark$\checkmarkd\checkmark$\checkmark^\checkmark\\checkmarkc\checkmarki\checkmarkr\checkmarkc\checkmark$\checkmarke\checkmark$\checkmark^\checkmark\\checkmarkc\checkmarki\checkmarkr\checkmarkc\checkmark$\checkmarks\checkmark$\checkmark^\checkmark\\checkmarkc\checkmarki\checkmarkr\checkmarkc\checkmark$\checkmark \checkmark$\checkmark^\checkmark\\checkmarkc\checkmarki\checkmarkr\checkmarkc\checkmark$\checkmarkm\checkmark$\checkmark^\checkmark\\checkmarkc\checkmarki\checkmarkr\checkmarkc\checkmark$\checkmarka\checkmark$\checkmark^\checkmark\\checkmarkc\checkmarki\checkmarkr\checkmarkc\checkmark$\checkmarkt\checkmark$\checkmark^\checkmark\\checkmarkc\checkmarki\checkmarkr\checkmarkc\checkmark$\checkmarké\checkmark$\checkmark^\checkmark\\checkmarkc\checkmarki\checkmarkr\checkmarkc\checkmark$\checkmarkr\checkmark$\checkmark^\checkmark\\checkmarkc\checkmarki\checkmarkr\checkmarkc\checkmark$\checkmarki\checkmark$\checkmark^\checkmark\\checkmarkc\checkmarki\checkmarkr\checkmarkc\checkmark$\checkmarka\checkmark$\checkmark^\checkmark\\checkmarkc\checkmarki\checkmarkr\checkmarkc\checkmark$\checkmarku\checkmark$\checkmark^\checkmark\\checkmarkc\checkmarki\checkmarkr\checkmarkc\checkmark$\checkmarkx\checkmark$\checkmark^\checkmark\\checkmarkc\checkmarki\checkmarkr\checkmarkc\checkmark$\checkmark \checkmark$\checkmark^\checkmark\\checkmarkc\checkmarki\checkmarkr\checkmarkc\checkmark$\checkmark(\checkmark$\checkmark^\checkmark\\checkmarkc\checkmarki\checkmarkr\checkmarkc\checkmark$\checkmarkR\checkmark$\checkmark^\checkmark\\checkmarkc\checkmarki\checkmarkr\checkmarkc\checkmark$\checkmarkD\checkmark$\checkmark^\checkmark\\checkmarkc\checkmarki\checkmarkr\checkmarkc\checkmark$\checkmarkM\checkmark$\checkmark^\checkmark\\checkmarkc\checkmarki\checkmarkr\checkmarkc\checkmark$\checkmark)\checkmark$\checkmark^\checkmark\\checkmarkc\checkmarki\checkmarkr\checkmarkc\checkmark$\checkmark \checkmark$\checkmark^\checkmark\\checkmarkc\checkmarki\checkmarkr\checkmarkc\checkmark$\checkmarkp\checkmark$\checkmark^\checkmark\\checkmarkc\checkmarki\checkmarkr\checkmarkc\checkmark$\checkmarka\checkmark$\checkmark^\checkmark\\checkmarkc\checkmarki\checkmarkr\checkmarkc\checkmark$\checkmarkr\checkmark$\checkmark^\checkmark\\checkmarkc\checkmarki\checkmarkr\checkmarkc\checkmark$\checkmark \checkmark$\checkmark^\checkmark\\checkmarkc\checkmarki\checkmarkr\checkmarkc\checkmark$\checkmarkl\checkmark$\checkmark^\checkmark\\checkmarkc\checkmarki\checkmarkr\checkmarkc\checkmark$\checkmarke\checkmark$\checkmark^\checkmark\\checkmarkc\checkmarki\checkmarkr\checkmarkc\checkmark$\checkmarks\checkmark$\checkmark^\checkmark\\checkmarkc\checkmarki\checkmarkr\checkmarkc\checkmark$\checkmark \checkmark$\checkmark^\checkmark\\checkmarkc\checkmarki\checkmarkr\checkmarkc\checkmark$\checkmarkc\checkmark$\checkmark^\checkmark\\checkmarkc\checkmarki\checkmarkr\checkmarkc\checkmark$\checkmarkr\checkmark$\checkmark^\checkmark\\checkmarkc\checkmarki\checkmarkr\checkmarkc\checkmark$\checkmarki\checkmark$\checkmark^\checkmark\\checkmarkc\checkmarki\checkmarkr\checkmarkc\checkmark$\checkmarkt\checkmark$\checkmark^\checkmark\\checkmarkc\checkmarki\checkmarkr\checkmarkc\checkmark$\checkmarkè\checkmark$\checkmark^\checkmark\\checkmarkc\checkmarki\checkmarkr\checkmarkc\checkmark$\checkmarkr\checkmark$\checkmark^\checkmark\\checkmarkc\checkmarki\checkmarkr\checkmarkc\checkmark$\checkmarke\checkmark$\checkmark^\checkmark\\checkmarkc\checkmarki\checkmarkr\checkmarkc\checkmark$\checkmarks\checkmark$\checkmark^\checkmark\\checkmarkc\checkmarki\checkmarkr\checkmarkc\checkmark$\checkmark \checkmark$\checkmark^\checkmark\\checkmarkc\checkmarki\checkmarkr\checkmarkc\checkmark$\checkmarkd\checkmark$\checkmark^\checkmark\\checkmarkc\checkmarki\checkmarkr\checkmarkc\checkmark$\checkmarke\checkmark$\checkmark^\checkmark\\checkmarkc\checkmarki\checkmarkr\checkmarkc\checkmark$\checkmark \checkmark$\checkmark^\checkmark\\checkmarkc\checkmarki\checkmarkr\checkmarkc\checkmark$\checkmarkV\checkmark$\checkmark^\checkmark\\checkmarkc\checkmarki\checkmarkr\checkmarkc\checkmark$\checkmarko\checkmark$\checkmark^\checkmark\\checkmarkc\checkmarki\checkmarkr\checkmarkc\checkmark$\checkmarkn\checkmark$\checkmark^\checkmark\\checkmarkc\checkmarki\checkmarkr\checkmarkc\checkmark$\checkmark \checkmark$\checkmark^\checkmark\\checkmarkc\checkmarki\checkmarkr\checkmarkc\checkmark$\checkmarkM\checkmark$\checkmark^\checkmark\\checkmarkc\checkmarki\checkmarkr\checkmarkc\checkmark$\checkmarki\checkmark$\checkmark^\checkmark\\checkmarkc\checkmarki\checkmarkr\checkmarkc\checkmark$\checkmarks\checkmark$\checkmark^\checkmark\\checkmarkc\checkmarki\checkmarkr\checkmarkc\checkmark$\checkmarke\checkmark$\checkmark^\checkmark\\checkmarkc\checkmarki\checkmarkr\checkmarkc\checkmark$\checkmarks\checkmark$\checkmark^\checkmark\\checkmarkc\checkmarki\checkmarkr\checkmarkc\checkmark$\checkmark,\checkmark$\checkmark^\checkmark\\checkmarkc\checkmarki\checkmarkr\checkmarkc\checkmark$\checkmark \checkmark$\checkmark^\checkmark\\checkmarkc\checkmarki\checkmarkr\checkmarkc\checkmark$\checkmarkd\checkmark$\checkmark^\checkmark\\checkmarkc\checkmarki\checkmarkr\checkmarkc\checkmark$\checkmark'\checkmark$\checkmark^\checkmark\\checkmarkc\checkmarki\checkmarkr\checkmarkc\checkmark$\checkmarkE\checkmark$\checkmark^\checkmark\\checkmarkc\checkmarki\checkmarkr\checkmarkc\checkmark$\checkmarku\checkmark$\checkmark^\checkmark\\checkmarkc\checkmarki\checkmarkr\checkmarkc\checkmark$\checkmarkl\checkmark$\checkmark^\checkmark\\checkmarkc\checkmarki\checkmarkr\checkmarkc\checkmark$\checkmarke\checkmark$\checkmark^\checkmark\\checkmarkc\checkmarki\checkmarkr\checkmarkc\checkmark$\checkmarkr\checkmark$\checkmark^\checkmark\\checkmarkc\checkmarki\checkmarkr\checkmarkc\checkmark$\checkmark \checkmark$\checkmark^\checkmark\\checkmarkc\checkmarki\checkmarkr\checkmarkc\checkmark$\checkmarke\checkmark$\checkmark^\checkmark\\checkmarkc\checkmarki\checkmarkr\checkmarkc\checkmark$\checkmarkt\checkmark$\checkmark^\checkmark\\checkmarkc\checkmarki\checkmarkr\checkmarkc\checkmark$\checkmark \checkmark$\checkmark^\checkmark\\checkmarkc\checkmarki\checkmarkr\checkmarkc\checkmark$\checkmarkI\checkmark$\checkmark^\checkmark\\checkmarkc\checkmarki\checkmarkr\checkmarkc\checkmark$\checkmarkS\checkmark$\checkmark^\checkmark\\checkmarkc\checkmarki\checkmarkr\checkmarkc\checkmark$\checkmarkO\checkmark$\checkmark^\checkmark\\checkmarkc\checkmarki\checkmarkr\checkmarkc\checkmark$\checkmark \checkmark$\checkmark^\checkmark\\checkmarkc\checkmarki\checkmarkr\checkmarkc\checkmark$\checkmark2\checkmark$\checkmark^\checkmark\\checkmarkc\checkmarki\checkmarkr\checkmarkc\checkmark$\checkmark8\checkmark$\checkmark^\checkmark\\checkmarkc\checkmarki\checkmarkr\checkmarkc\checkmark$\checkmark1\checkmark$\checkmark^\checkmark\\checkmarkc\checkmarki\checkmarkr\checkmarkc\checkmark$\checkmark.\checkmark$\checkmark^\checkmark\\checkmarkc\checkmarki\checkmarkr\checkmarkc\checkmark$\checkmark
\checkmark$\checkmark^\checkmark\\checkmarkc\checkmarki\checkmarkr\checkmarkc\checkmark$\checkmark\\checkmark$\checkmark^\checkmark\\checkmarkc\checkmarki\checkmarkr\checkmarkc\checkmark$\checkmarke\checkmark$\checkmark^\checkmark\\checkmarkc\checkmarki\checkmarkr\checkmarkc\checkmark$\checkmarkn\checkmark$\checkmark^\checkmark\\checkmarkc\checkmarki\checkmarkr\checkmarkc\checkmark$\checkmarkd\checkmark$\checkmark^\checkmark\\checkmarkc\checkmarki\checkmarkr\checkmarkc\checkmark$\checkmark{\checkmark$\checkmark^\checkmark\\checkmarkc\checkmarki\checkmarkr\checkmarkc\checkmark$\checkmarki\checkmark$\checkmark^\checkmark\\checkmarkc\checkmarki\checkmarkr\checkmarkc\checkmark$\checkmarkt\checkmark$\checkmark^\checkmark\\checkmarkc\checkmarki\checkmarkr\checkmarkc\checkmark$\checkmarke\checkmark$\checkmark^\checkmark\\checkmarkc\checkmarki\checkmarkr\checkmarkc\checkmark$\checkmarkm\checkmark$\checkmark^\checkmark\\checkmarkc\checkmarki\checkmarkr\checkmarkc\checkmark$\checkmarki\checkmark$\checkmark^\checkmark\\checkmarkc\checkmarki\checkmarkr\checkmarkc\checkmark$\checkmarkz\checkmark$\checkmark^\checkmark\\checkmarkc\checkmarki\checkmarkr\checkmarkc\checkmark$\checkmarke\checkmark$\checkmark^\checkmark\\checkmarkc\checkmarki\checkmarkr\checkmarkc\checkmark$\checkmark}\checkmark$\checkmark^\checkmark\\checkmarkc\checkmarki\checkmarkr\checkmarkc\checkmark$\checkmark
\checkmark$\checkmark^\checkmark\\checkmarkc\checkmarki\checkmarkr\checkmarkc\checkmark$\checkmark
\checkmark$\checkmark^\checkmark\\checkmarkc\checkmarki\checkmarkr\checkmarkc\checkmark$\checkmarkL\checkmark$\checkmark^\checkmark\\checkmarkc\checkmarki\checkmarkr\checkmarkc\checkmark$\checkmarke\checkmark$\checkmark^\checkmark\\checkmarkc\checkmarki\checkmarkr\checkmarkc\checkmark$\checkmarks\checkmark$\checkmark^\checkmark\\checkmarkc\checkmarki\checkmarkr\checkmarkc\checkmark$\checkmark \checkmark$\checkmark^\checkmark\\checkmarkc\checkmarki\checkmarkr\checkmarkc\checkmark$\checkmarke\checkmark$\checkmark^\checkmark\\checkmarkc\checkmarki\checkmarkr\checkmarkc\checkmark$\checkmarkf\checkmark$\checkmark^\checkmark\\checkmarkc\checkmarki\checkmarkr\checkmarkc\checkmark$\checkmarkf\checkmark$\checkmark^\checkmark\\checkmarkc\checkmarki\checkmarkr\checkmarkc\checkmark$\checkmarko\checkmark$\checkmark^\checkmark\\checkmarkc\checkmarki\checkmarkr\checkmarkc\checkmark$\checkmarkr\checkmark$\checkmark^\checkmark\\checkmarkc\checkmarki\checkmarkr\checkmarkc\checkmark$\checkmarkt\checkmark$\checkmark^\checkmark\\checkmarkc\checkmarki\checkmarkr\checkmarkc\checkmark$\checkmarks\checkmark$\checkmark^\checkmark\\checkmarkc\checkmarki\checkmarkr\checkmarkc\checkmark$\checkmark \checkmark$\checkmark^\checkmark\\checkmarkc\checkmarki\checkmarkr\checkmarkc\checkmark$\checkmarkd\checkmark$\checkmark^\checkmark\\checkmarkc\checkmarki\checkmarkr\checkmarkc\checkmark$\checkmarke\checkmark$\checkmark^\checkmark\\checkmarkc\checkmarki\checkmarkr\checkmarkc\checkmark$\checkmark \checkmark$\checkmark^\checkmark\\checkmarkc\checkmarki\checkmarkr\checkmarkc\checkmark$\checkmarkc\checkmark$\checkmark^\checkmark\\checkmarkc\checkmarki\checkmarkr\checkmarkc\checkmark$\checkmarko\checkmark$\checkmark^\checkmark\\checkmarkc\checkmarki\checkmarkr\checkmarkc\checkmark$\checkmarku\checkmark$\checkmark^\checkmark\\checkmarkc\checkmarki\checkmarkr\checkmarkc\checkmark$\checkmarkp\checkmark$\checkmark^\checkmark\\checkmarkc\checkmarki\checkmarkr\checkmarkc\checkmark$\checkmarke\checkmark$\checkmark^\checkmark\\checkmarkc\checkmarki\checkmarkr\checkmarkc\checkmark$\checkmark \checkmark$\checkmark^\checkmark\\checkmarkc\checkmarki\checkmarkr\checkmarkc\checkmark$\checkmarkc\checkmark$\checkmark^\checkmark\\checkmarkc\checkmarki\checkmarkr\checkmarkc\checkmark$\checkmarko\checkmark$\checkmark^\checkmark\\checkmarkc\checkmarki\checkmarkr\checkmarkc\checkmark$\checkmarkn\checkmark$\checkmark^\checkmark\\checkmarkc\checkmarki\checkmarkr\checkmarkc\checkmark$\checkmarks\checkmark$\checkmark^\checkmark\\checkmarkc\checkmarki\checkmarkr\checkmarkc\checkmark$\checkmarkt\checkmark$\checkmark^\checkmark\\checkmarkc\checkmarki\checkmarkr\checkmarkc\checkmark$\checkmarki\checkmark$\checkmark^\checkmark\\checkmarkc\checkmarki\checkmarkr\checkmarkc\checkmark$\checkmarkt\checkmark$\checkmark^\checkmark\\checkmarkc\checkmarki\checkmarkr\checkmarkc\checkmark$\checkmarku\checkmark$\checkmark^\checkmark\\checkmarkc\checkmarki\checkmarkr\checkmarkc\checkmark$\checkmarke\checkmark$\checkmark^\checkmark\\checkmarkc\checkmarki\checkmarkr\checkmarkc\checkmark$\checkmarkn\checkmark$\checkmark^\checkmark\\checkmarkc\checkmarki\checkmarkr\checkmarkc\checkmark$\checkmarkt\checkmark$\checkmark^\checkmark\\checkmarkc\checkmarki\checkmarkr\checkmarkc\checkmark$\checkmark \checkmark$\checkmark^\checkmark\\checkmarkc\checkmarki\checkmarkr\checkmarkc\checkmark$\checkmarkl\checkmark$\checkmark^\checkmark\\checkmarkc\checkmarki\checkmarkr\checkmarkc\checkmark$\checkmarke\checkmark$\checkmark^\checkmark\\checkmarkc\checkmarki\checkmarkr\checkmarkc\checkmark$\checkmarks\checkmark$\checkmark^\checkmark\\checkmarkc\checkmarki\checkmarkr\checkmarkc\checkmark$\checkmark \checkmark$\checkmark^\checkmark\\checkmarkc\checkmarki\checkmarkr\checkmarkc\checkmark$\checkmarkd\checkmark$\checkmark^\checkmark\\checkmarkc\checkmarki\checkmarkr\checkmarkc\checkmark$\checkmarko\checkmark$\checkmark^\checkmark\\checkmarkc\checkmarki\checkmarkr\checkmarkc\checkmark$\checkmarkn\checkmark$\checkmark^\checkmark\\checkmarkc\checkmarki\checkmarkr\checkmarkc\checkmark$\checkmarkn\checkmark$\checkmark^\checkmark\\checkmarkc\checkmarki\checkmarkr\checkmarkc\checkmark$\checkmarké\checkmark$\checkmark^\checkmark\\checkmarkc\checkmarki\checkmarkr\checkmarkc\checkmark$\checkmarke\checkmark$\checkmark^\checkmark\\checkmarkc\checkmarki\checkmarkr\checkmarkc\checkmark$\checkmarks\checkmark$\checkmark^\checkmark\\checkmarkc\checkmarki\checkmarkr\checkmarkc\checkmark$\checkmark \checkmark$\checkmark^\checkmark\\checkmarkc\checkmarki\checkmarkr\checkmarkc\checkmark$\checkmarkd\checkmark$\checkmark^\checkmark\\checkmarkc\checkmarki\checkmarkr\checkmarkc\checkmark$\checkmark'\checkmark$\checkmark^\checkmark\\checkmarkc\checkmarki\checkmarkr\checkmarkc\checkmark$\checkmarke\checkmark$\checkmark^\checkmark\\checkmarkc\checkmarki\checkmarkr\checkmarkc\checkmark$\checkmarkn\checkmark$\checkmark^\checkmark\\checkmarkc\checkmarki\checkmarkr\checkmarkc\checkmark$\checkmarkt\checkmark$\checkmark^\checkmark\\checkmarkc\checkmarki\checkmarkr\checkmarkc\checkmark$\checkmarkr\checkmark$\checkmark^\checkmark\\checkmarkc\checkmarki\checkmarkr\checkmarkc\checkmark$\checkmarké\checkmark$\checkmark^\checkmark\\checkmarkc\checkmarki\checkmarkr\checkmarkc\checkmark$\checkmarke\checkmark$\checkmark^\checkmark\\checkmarkc\checkmarki\checkmarkr\checkmarkc\checkmark$\checkmark \checkmark$\checkmark^\checkmark\\checkmarkc\checkmarki\checkmarkr\checkmarkc\checkmark$\checkmarkf\checkmark$\checkmark^\checkmark\\checkmarkc\checkmarki\checkmarkr\checkmarkc\checkmark$\checkmarko\checkmark$\checkmark^\checkmark\\checkmarkc\checkmarki\checkmarkr\checkmarkc\checkmark$\checkmarkn\checkmark$\checkmark^\checkmark\\checkmarkc\checkmarki\checkmarkr\checkmarkc\checkmark$\checkmarkd\checkmark$\checkmark^\checkmark\\checkmarkc\checkmarki\checkmarkr\checkmarkc\checkmark$\checkmarka\checkmark$\checkmark^\checkmark\\checkmarkc\checkmarki\checkmarkr\checkmarkc\checkmark$\checkmarkm\checkmark$\checkmark^\checkmark\\checkmarkc\checkmarki\checkmarkr\checkmarkc\checkmark$\checkmarke\checkmark$\checkmark^\checkmark\\checkmarkc\checkmarki\checkmarkr\checkmarkc\checkmark$\checkmarkn\checkmark$\checkmark^\checkmark\\checkmarkc\checkmarki\checkmarkr\checkmarkc\checkmark$\checkmarkt\checkmark$\checkmark^\checkmark\\checkmarkc\checkmarki\checkmarkr\checkmarkc\checkmark$\checkmarka\checkmark$\checkmark^\checkmark\\checkmarkc\checkmarki\checkmarkr\checkmarkc\checkmark$\checkmarkl\checkmark$\checkmark^\checkmark\\checkmarkc\checkmarki\checkmarkr\checkmarkc\checkmark$\checkmarke\checkmark$\checkmark^\checkmark\\checkmarkc\checkmarki\checkmarkr\checkmarkc\checkmark$\checkmarks\checkmark$\checkmark^\checkmark\\checkmarkc\checkmarki\checkmarkr\checkmarkc\checkmark$\checkmark \checkmark$\checkmark^\checkmark\\checkmarkc\checkmarki\checkmarkr\checkmarkc\checkmark$\checkmarkd\checkmark$\checkmark^\checkmark\\checkmarkc\checkmarki\checkmarkr\checkmarkc\checkmark$\checkmarke\checkmark$\checkmark^\checkmark\\checkmarkc\checkmarki\checkmarkr\checkmarkc\checkmark$\checkmark \checkmark$\checkmark^\checkmark\\checkmarkc\checkmarki\checkmarkr\checkmarkc\checkmark$\checkmarkt\checkmark$\checkmark^\checkmark\\checkmarkc\checkmarki\checkmarkr\checkmarkc\checkmark$\checkmarko\checkmark$\checkmark^\checkmark\\checkmarkc\checkmarki\checkmarkr\checkmarkc\checkmark$\checkmarku\checkmark$\checkmark^\checkmark\\checkmarkc\checkmarki\checkmarkr\checkmarkc\checkmark$\checkmarkt\checkmark$\checkmark^\checkmark\\checkmarkc\checkmarki\checkmarkr\checkmarkc\checkmark$\checkmark \checkmark$\checkmark^\checkmark\\checkmarkc\checkmarki\checkmarkr\checkmarkc\checkmark$\checkmarkl\checkmark$\checkmark^\checkmark\\checkmarkc\checkmarki\checkmarkr\checkmarkc\checkmark$\checkmarke\checkmark$\checkmark^\checkmark\\checkmarkc\checkmarki\checkmarkr\checkmarkc\checkmark$\checkmark \checkmark$\checkmark^\checkmark\\checkmarkc\checkmarki\checkmarkr\checkmarkc\checkmark$\checkmarkd\checkmark$\checkmark^\checkmark\\checkmarkc\checkmarki\checkmarkr\checkmarkc\checkmark$\checkmarki\checkmark$\checkmark^\checkmark\\checkmarkc\checkmarki\checkmarkr\checkmarkc\checkmark$\checkmarkm\checkmark$\checkmark^\checkmark\\checkmarkc\checkmarki\checkmarkr\checkmarkc\checkmark$\checkmarke\checkmark$\checkmark^\checkmark\\checkmarkc\checkmarki\checkmarkr\checkmarkc\checkmark$\checkmarkn\checkmark$\checkmark^\checkmark\\checkmarkc\checkmarki\checkmarkr\checkmarkc\checkmark$\checkmarks\checkmark$\checkmark^\checkmark\\checkmarkc\checkmarki\checkmarkr\checkmarkc\checkmark$\checkmarki\checkmark$\checkmark^\checkmark\\checkmarkc\checkmarki\checkmarkr\checkmarkc\checkmark$\checkmarko\checkmark$\checkmark^\checkmark\\checkmarkc\checkmarki\checkmarkr\checkmarkc\checkmark$\checkmarkn\checkmark$\checkmark^\checkmark\\checkmarkc\checkmarki\checkmarkr\checkmarkc\checkmark$\checkmarkn\checkmark$\checkmark^\checkmark\\checkmarkc\checkmarki\checkmarkr\checkmarkc\checkmark$\checkmarke\checkmark$\checkmark^\checkmark\\checkmarkc\checkmarki\checkmarkr\checkmarkc\checkmark$\checkmarkm\checkmark$\checkmark^\checkmark\\checkmarkc\checkmarki\checkmarkr\checkmarkc\checkmark$\checkmarke\checkmark$\checkmark^\checkmark\\checkmarkc\checkmarki\checkmarkr\checkmarkc\checkmark$\checkmarkn\checkmark$\checkmark^\checkmark\\checkmarkc\checkmarki\checkmarkr\checkmarkc\checkmark$\checkmarkt\checkmark$\checkmark^\checkmark\\checkmarkc\checkmarki\checkmarkr\checkmarkc\checkmark$\checkmark.\checkmark$\checkmark^\checkmark\\checkmarkc\checkmarki\checkmarkr\checkmarkc\checkmark$\checkmark \checkmark$\checkmark^\checkmark\\checkmarkc\checkmarki\checkmarkr\checkmarkc\checkmark$\checkmarkL\checkmark$\checkmark^\checkmark\\checkmarkc\checkmarki\checkmarkr\checkmarkc\checkmark$\checkmarke\checkmark$\checkmark^\checkmark\\checkmarkc\checkmarki\checkmarkr\checkmarkc\checkmark$\checkmarks\checkmark$\checkmark^\checkmark\\checkmarkc\checkmarki\checkmarkr\checkmarkc\checkmark$\checkmark \checkmark$\checkmark^\checkmark\\checkmarkc\checkmarki\checkmarkr\checkmarkc\checkmark$\checkmarkc\checkmark$\checkmark^\checkmark\\checkmarkc\checkmarki\checkmarkr\checkmarkc\checkmark$\checkmarko\checkmark$\checkmark^\checkmark\\checkmarkc\checkmarki\checkmarkr\checkmarkc\checkmark$\checkmarke\checkmark$\checkmark^\checkmark\\checkmarkc\checkmarki\checkmarkr\checkmarkc\checkmark$\checkmarkf\checkmark$\checkmark^\checkmark\\checkmarkc\checkmarki\checkmarkr\checkmarkc\checkmark$\checkmarkf\checkmark$\checkmark^\checkmark\\checkmarkc\checkmarki\checkmarkr\checkmarkc\checkmark$\checkmarki\checkmark$\checkmark^\checkmark\\checkmarkc\checkmarki\checkmarkr\checkmarkc\checkmark$\checkmarkc\checkmark$\checkmark^\checkmark\\checkmarkc\checkmarki\checkmarkr\checkmarkc\checkmark$\checkmarki\checkmark$\checkmark^\checkmark\\checkmarkc\checkmarki\checkmarkr\checkmarkc\checkmark$\checkmarke\checkmark$\checkmark^\checkmark\\checkmarkc\checkmarki\checkmarkr\checkmarkc\checkmark$\checkmarkn\checkmark$\checkmark^\checkmark\\checkmarkc\checkmarki\checkmarkr\checkmarkc\checkmark$\checkmarkt\checkmark$\checkmark^\checkmark\\checkmarkc\checkmarki\checkmarkr\checkmarkc\checkmark$\checkmarks\checkmark$\checkmark^\checkmark\\checkmarkc\checkmarki\checkmarkr\checkmarkc\checkmark$\checkmark \checkmark$\checkmark^\checkmark\\checkmarkc\checkmarki\checkmarkr\checkmarkc\checkmark$\checkmarkd\checkmark$\checkmark^\checkmark\\checkmarkc\checkmarki\checkmarkr\checkmarkc\checkmark$\checkmarke\checkmark$\checkmark^\checkmark\\checkmarkc\checkmarki\checkmarkr\checkmarkc\checkmark$\checkmark \checkmark$\checkmark^\checkmark\\checkmarkc\checkmarki\checkmarkr\checkmarkc\checkmark$\checkmarks\checkmark$\checkmark^\checkmark\\checkmarkc\checkmarki\checkmarkr\checkmarkc\checkmark$\checkmarké\checkmark$\checkmark^\checkmark\\checkmarkc\checkmarki\checkmarkr\checkmarkc\checkmark$\checkmarkc\checkmark$\checkmark^\checkmark\\checkmarkc\checkmarki\checkmarkr\checkmarkc\checkmark$\checkmarku\checkmark$\checkmark^\checkmark\\checkmarkc\checkmarki\checkmarkr\checkmarkc\checkmark$\checkmarkr\checkmark$\checkmark^\checkmark\\checkmarkc\checkmarki\checkmarkr\checkmarkc\checkmark$\checkmarki\checkmark$\checkmark^\checkmark\\checkmarkc\checkmarki\checkmarkr\checkmarkc\checkmark$\checkmarkt\checkmark$\checkmark^\checkmark\\checkmarkc\checkmarki\checkmarkr\checkmarkc\checkmark$\checkmarké\checkmark$\checkmark^\checkmark\\checkmarkc\checkmarki\checkmarkr\checkmarkc\checkmark$\checkmark \checkmark$\checkmark^\checkmark\\checkmarkc\checkmarki\checkmarkr\checkmarkc\checkmark$\checkmarka\checkmark$\checkmark^\checkmark\\checkmarkc\checkmarki\checkmarkr\checkmarkc\checkmark$\checkmarkd\checkmark$\checkmark^\checkmark\\checkmarkc\checkmarki\checkmarkr\checkmarkc\checkmark$\checkmarko\checkmark$\checkmark^\checkmark\\checkmarkc\checkmarki\checkmarkr\checkmarkc\checkmark$\checkmarkp\checkmark$\checkmark^\checkmark\\checkmarkc\checkmarki\checkmarkr\checkmarkc\checkmark$\checkmarkt\checkmark$\checkmark^\checkmark\\checkmarkc\checkmarki\checkmarkr\checkmarkc\checkmark$\checkmarké\checkmark$\checkmark^\checkmark\\checkmarkc\checkmarki\checkmarkr\checkmarkc\checkmark$\checkmarks\checkmark$\checkmark^\checkmark\\checkmarkc\checkmarki\checkmarkr\checkmarkc\checkmark$\checkmark \checkmark$\checkmark^\checkmark\\checkmarkc\checkmarki\checkmarkr\checkmarkc\checkmark$\checkmarks\checkmark$\checkmark^\checkmark\\checkmarkc\checkmarki\checkmarkr\checkmarkc\checkmark$\checkmarko\checkmark$\checkmark^\checkmark\\checkmarkc\checkmarki\checkmarkr\checkmarkc\checkmark$\checkmarkn\checkmark$\checkmark^\checkmark\\checkmarkc\checkmarki\checkmarkr\checkmarkc\checkmark$\checkmarkt\checkmark$\checkmark^\checkmark\\checkmarkc\checkmarki\checkmarkr\checkmarkc\checkmark$\checkmark \checkmark$\checkmark^\checkmark\\checkmarkc\checkmarki\checkmarkr\checkmarkc\checkmark$\checkmarkc\checkmark$\checkmark^\checkmark\\checkmarkc\checkmarki\checkmarkr\checkmarkc\checkmark$\checkmarko\checkmark$\checkmark^\checkmark\\checkmarkc\checkmarki\checkmarkr\checkmarkc\checkmark$\checkmarkn\checkmark$\checkmark^\checkmark\\checkmarkc\checkmarki\checkmarkr\checkmarkc\checkmark$\checkmarkf\checkmark$\checkmark^\checkmark\\checkmarkc\checkmarki\checkmarkr\checkmarkc\checkmark$\checkmarko\checkmark$\checkmark^\checkmark\\checkmarkc\checkmarki\checkmarkr\checkmarkc\checkmark$\checkmarkr\checkmark$\checkmark^\checkmark\\checkmarkc\checkmarki\checkmarkr\checkmarkc\checkmark$\checkmarkm\checkmark$\checkmark^\checkmark\\checkmarkc\checkmarki\checkmarkr\checkmarkc\checkmark$\checkmarke\checkmark$\checkmark^\checkmark\\checkmarkc\checkmarki\checkmarkr\checkmarkc\checkmark$\checkmarks\checkmark$\checkmark^\checkmark\\checkmarkc\checkmarki\checkmarkr\checkmarkc\checkmark$\checkmark \checkmark$\checkmark^\checkmark\\checkmarkc\checkmarki\checkmarkr\checkmarkc\checkmark$\checkmarka\checkmark$\checkmark^\checkmark\\checkmarkc\checkmarki\checkmarkr\checkmarkc\checkmark$\checkmarku\checkmark$\checkmark^\checkmark\\checkmarkc\checkmarki\checkmarkr\checkmarkc\checkmark$\checkmarkx\checkmark$\checkmark^\checkmark\\checkmarkc\checkmarki\checkmarkr\checkmarkc\checkmark$\checkmark \checkmark$\checkmark^\checkmark\\checkmarkc\checkmarki\checkmarkr\checkmarkc\checkmark$\checkmarkr\checkmark$\checkmark^\checkmark\\checkmarkc\checkmarki\checkmarkr\checkmarkc\checkmark$\checkmarke\checkmark$\checkmark^\checkmark\\checkmarkc\checkmarki\checkmarkr\checkmarkc\checkmark$\checkmarkc\checkmark$\checkmark^\checkmark\\checkmarkc\checkmarki\checkmarkr\checkmarkc\checkmark$\checkmarko\checkmark$\checkmark^\checkmark\\checkmarkc\checkmarki\checkmarkr\checkmarkc\checkmark$\checkmarkm\checkmark$\checkmark^\checkmark\\checkmarkc\checkmarki\checkmarkr\checkmarkc\checkmark$\checkmarkm\checkmark$\checkmark^\checkmark\\checkmarkc\checkmarki\checkmarkr\checkmarkc\checkmark$\checkmarka\checkmark$\checkmark^\checkmark\\checkmarkc\checkmarki\checkmarkr\checkmarkc\checkmark$\checkmarkn\checkmark$\checkmark^\checkmark\\checkmarkc\checkmarki\checkmarkr\checkmarkc\checkmark$\checkmarkd\checkmark$\checkmark^\checkmark\\checkmarkc\checkmarki\checkmarkr\checkmarkc\checkmark$\checkmarka\checkmark$\checkmark^\checkmark\\checkmarkc\checkmarki\checkmarkr\checkmarkc\checkmark$\checkmarkt\checkmark$\checkmark^\checkmark\\checkmarkc\checkmarki\checkmarkr\checkmarkc\checkmark$\checkmarki\checkmark$\checkmark^\checkmark\\checkmarkc\checkmarki\checkmarkr\checkmarkc\checkmark$\checkmarko\checkmark$\checkmark^\checkmark\\checkmarkc\checkmarki\checkmarkr\checkmarkc\checkmark$\checkmarkn\checkmark$\checkmark^\checkmark\\checkmarkc\checkmarki\checkmarkr\checkmarkc\checkmark$\checkmarks\checkmark$\checkmark^\checkmark\\checkmarkc\checkmarki\checkmarkr\checkmarkc\checkmark$\checkmark \checkmark$\checkmark^\checkmark\\checkmarkc\checkmarki\checkmarkr\checkmarkc\checkmark$\checkmarkd\checkmark$\checkmark^\checkmark\\checkmarkc\checkmarki\checkmarkr\checkmarkc\checkmark$\checkmarku\checkmark$\checkmark^\checkmark\\checkmarkc\checkmarki\checkmarkr\checkmarkc\checkmark$\checkmark \checkmark$\checkmark^\checkmark\\checkmarkc\checkmarki\checkmarkr\checkmarkc\checkmark$\checkmarkc\checkmark$\checkmark^\checkmark\\checkmarkc\checkmarki\checkmarkr\checkmarkc\checkmark$\checkmarko\checkmark$\checkmark^\checkmark\\checkmarkc\checkmarki\checkmarkr\checkmarkc\checkmark$\checkmarku\checkmark$\checkmark^\checkmark\\checkmarkc\checkmarki\checkmarkr\checkmarkc\checkmark$\checkmarkr\checkmark$\checkmark^\checkmark\\checkmarkc\checkmarki\checkmarkr\checkmarkc\checkmark$\checkmarks\checkmark$\checkmark^\checkmark\\checkmarkc\checkmarki\checkmarkr\checkmarkc\checkmark$\checkmark \checkmark$\checkmark^\checkmark\\checkmarkc\checkmarki\checkmarkr\checkmarkc\checkmark$\checkmarkd\checkmark$\checkmark^\checkmark\\checkmarkc\checkmarki\checkmarkr\checkmarkc\checkmark$\checkmarke\checkmark$\checkmark^\checkmark\\checkmarkc\checkmarki\checkmarkr\checkmarkc\checkmark$\checkmark \checkmark$\checkmark^\checkmark\\checkmarkc\checkmarki\checkmarkr\checkmarkc\checkmark$\checkmarkc\checkmark$\checkmark^\checkmark\\checkmarkc\checkmarki\checkmarkr\checkmarkc\checkmark$\checkmarko\checkmark$\checkmark^\checkmark\\checkmarkc\checkmarki\checkmarkr\checkmarkc\checkmark$\checkmarkn\checkmark$\checkmark^\checkmark\\checkmarkc\checkmarki\checkmarkr\checkmarkc\checkmark$\checkmarkc\checkmark$\checkmark^\checkmark\\checkmarkc\checkmarki\checkmarkr\checkmarkc\checkmark$\checkmarke\checkmark$\checkmark^\checkmark\\checkmarkc\checkmarki\checkmarkr\checkmarkc\checkmark$\checkmarkp\checkmark$\checkmark^\checkmark\\checkmarkc\checkmarki\checkmarkr\checkmarkc\checkmark$\checkmarkt\checkmark$\checkmark^\checkmark\\checkmarkc\checkmarki\checkmarkr\checkmarkc\checkmark$\checkmarki\checkmark$\checkmark^\checkmark\\checkmarkc\checkmarki\checkmarkr\checkmarkc\checkmark$\checkmarko\checkmark$\checkmark^\checkmark\\checkmarkc\checkmarki\checkmarkr\checkmarkc\checkmark$\checkmarkn\checkmark$\checkmark^\checkmark\\checkmarkc\checkmarki\checkmarkr\checkmarkc\checkmark$\checkmark \checkmark$\checkmark^\checkmark\\checkmarkc\checkmarki\checkmarkr\checkmarkc\checkmark$\checkmarkd\checkmark$\checkmark^\checkmark\\checkmarkc\checkmarki\checkmarkr\checkmarkc\checkmark$\checkmarke\checkmark$\checkmark^\checkmark\\checkmarkc\checkmarki\checkmarkr\checkmarkc\checkmark$\checkmarks\checkmark$\checkmark^\checkmark\\checkmarkc\checkmarki\checkmarkr\checkmarkc\checkmark$\checkmark \checkmark$\checkmark^\checkmark\\checkmarkc\checkmarki\checkmarkr\checkmarkc\checkmark$\checkmarks\checkmark$\checkmark^\checkmark\\checkmarkc\checkmarki\checkmarkr\checkmarkc\checkmark$\checkmarky\checkmark$\checkmark^\checkmark\\checkmarkc\checkmarki\checkmarkr\checkmarkc\checkmark$\checkmarks\checkmark$\checkmark^\checkmark\\checkmarkc\checkmarki\checkmarkr\checkmarkc\checkmark$\checkmarkt\checkmark$\checkmark^\checkmark\\checkmarkc\checkmarki\checkmarkr\checkmarkc\checkmark$\checkmarkè\checkmark$\checkmark^\checkmark\\checkmarkc\checkmarki\checkmarkr\checkmarkc\checkmark$\checkmarkm\checkmark$\checkmark^\checkmark\\checkmarkc\checkmarki\checkmarkr\checkmarkc\checkmark$\checkmarke\checkmark$\checkmark^\checkmark\\checkmarkc\checkmarki\checkmarkr\checkmarkc\checkmark$\checkmarks\checkmark$\checkmark^\checkmark\\checkmarkc\checkmarki\checkmarkr\checkmarkc\checkmark$\checkmark \checkmark$\checkmark^\checkmark\\checkmarkc\checkmarki\checkmarkr\checkmarkc\checkmark$\checkmarkm\checkmark$\checkmark^\checkmark\\checkmarkc\checkmarki\checkmarkr\checkmarkc\checkmark$\checkmarké\checkmark$\checkmark^\checkmark\\checkmarkc\checkmarki\checkmarkr\checkmarkc\checkmark$\checkmarkc\checkmark$\checkmark^\checkmark\\checkmarkc\checkmarki\checkmarkr\checkmarkc\checkmark$\checkmarka\checkmark$\checkmark^\checkmark\\checkmarkc\checkmarki\checkmarkr\checkmarkc\checkmark$\checkmarkn\checkmark$\checkmark^\checkmark\\checkmarkc\checkmarki\checkmarkr\checkmarkc\checkmark$\checkmarki\checkmark$\checkmark^\checkmark\\checkmarkc\checkmarki\checkmarkr\checkmarkc\checkmark$\checkmarkq\checkmark$\checkmark^\checkmark\\checkmarkc\checkmarki\checkmarkr\checkmarkc\checkmark$\checkmarku\checkmark$\checkmark^\checkmark\\checkmarkc\checkmarki\checkmarkr\checkmarkc\checkmark$\checkmarke\checkmark$\checkmark^\checkmark\\checkmarkc\checkmarki\checkmarkr\checkmarkc\checkmark$\checkmarks\checkmark$\checkmark^\checkmark\\checkmarkc\checkmarki\checkmarkr\checkmarkc\checkmark$\checkmark \checkmark$\checkmark^\checkmark\\checkmarkc\checkmarki\checkmarkr\checkmarkc\checkmark$\checkmarkd\checkmark$\checkmark^\checkmark\\checkmarkc\checkmarki\checkmarkr\checkmarkc\checkmark$\checkmarke\checkmark$\checkmark^\checkmark\\checkmarkc\checkmarki\checkmarkr\checkmarkc\checkmark$\checkmark \checkmark$\checkmark^\checkmark\\checkmarkc\checkmarki\checkmarkr\checkmarkc\checkmark$\checkmarkl\checkmark$\checkmark^\checkmark\\checkmarkc\checkmarki\checkmarkr\checkmarkc\checkmark$\checkmark'\checkmark$\checkmark^\checkmark\\checkmarkc\checkmarki\checkmarkr\checkmarkc\checkmark$\checkmarkE\checkmark$\checkmark^\checkmark\\checkmarkc\checkmarki\checkmarkr\checkmarkc\checkmark$\checkmarkN\checkmark$\checkmark^\checkmark\\checkmarkc\checkmarki\checkmarkr\checkmarkc\checkmark$\checkmarkI\checkmark$\checkmark^\checkmark\\checkmarkc\checkmarki\checkmarkr\checkmarkc\checkmark$\checkmark-\checkmark$\checkmark^\checkmark\\checkmarkc\checkmarki\checkmarkr\checkmarkc\checkmark$\checkmarkA\checkmark$\checkmark^\checkmark\\checkmarkc\checkmarki\checkmarkr\checkmarkc\checkmark$\checkmarkB\checkmark$\checkmark^\checkmark\\checkmarkc\checkmarki\checkmarkr\checkmarkc\checkmark$\checkmarkT\checkmark$\checkmark^\checkmark\\checkmarkc\checkmarki\checkmarkr\checkmarkc\checkmark$\checkmark \checkmark$\checkmark^\checkmark\\checkmarkc\checkmarki\checkmarkr\checkmarkc\checkmark$\checkmark(\checkmark$\checkmark^\checkmark\\checkmarkc\checkmarki\checkmarkr\checkmarkc\checkmark$\checkmark$\checkmark$\checkmark^\checkmark\\checkmarkc\checkmarki\checkmarkr\checkmarkc\checkmark$\checkmarks\checkmark$\checkmark^\checkmark\\checkmarkc\checkmarki\checkmarkr\checkmarkc\checkmark$\checkmark \checkmark$\checkmark^\checkmark\\checkmarkc\checkmarki\checkmarkr\checkmarkc\checkmark$\checkmark\\checkmark$\checkmark^\checkmark\\checkmarkc\checkmarki\checkmarkr\checkmarkc\checkmark$\checkmarkg\checkmark$\checkmark^\checkmark\\checkmarkc\checkmarki\checkmarkr\checkmarkc\checkmark$\checkmarke\checkmark$\checkmark^\checkmark\\checkmarkc\checkmarki\checkmarkr\checkmarkc\checkmark$\checkmarkq\checkmark$\checkmark^\checkmark\\checkmarkc\checkmarki\checkmarkr\checkmarkc\checkmark$\checkmark \checkmark$\checkmark^\checkmark\\checkmarkc\checkmarki\checkmarkr\checkmarkc\checkmark$\checkmark1\checkmark$\checkmark^\checkmark\\checkmarkc\checkmarki\checkmarkr\checkmarkc\checkmark$\checkmark{\checkmark$\checkmark^\checkmark\\checkmarkc\checkmarki\checkmarkr\checkmarkc\checkmark$\checkmark,\checkmark$\checkmark^\checkmark\\checkmarkc\checkmarki\checkmarkr\checkmarkc\checkmark$\checkmark}\checkmark$\checkmark^\checkmark\\checkmarkc\checkmarki\checkmarkr\checkmarkc\checkmark$\checkmark5\checkmark$\checkmark^\checkmark\\checkmarkc\checkmarki\checkmarkr\checkmarkc\checkmark$\checkmark$\checkmark$\checkmark^\checkmark\\checkmarkc\checkmarki\checkmarkr\checkmarkc\checkmark$\checkmark \checkmark$\checkmark^\checkmark\\checkmarkc\checkmarki\checkmarkr\checkmarkc\checkmark$\checkmarkp\checkmark$\checkmark^\checkmark\\checkmarkc\checkmarki\checkmarkr\checkmarkc\checkmark$\checkmarko\checkmark$\checkmark^\checkmark\\checkmarkc\checkmarki\checkmarkr\checkmarkc\checkmark$\checkmarku\checkmark$\checkmark^\checkmark\\checkmarkc\checkmarki\checkmarkr\checkmarkc\checkmark$\checkmarkr\checkmark$\checkmark^\checkmark\\checkmarkc\checkmarki\checkmarkr\checkmarkc\checkmark$\checkmark \checkmark$\checkmark^\checkmark\\checkmarkc\checkmarki\checkmarkr\checkmarkc\checkmark$\checkmarkl\checkmark$\checkmark^\checkmark\\checkmarkc\checkmarki\checkmarkr\checkmarkc\checkmark$\checkmarka\checkmark$\checkmark^\checkmark\\checkmarkc\checkmarki\checkmarkr\checkmarkc\checkmark$\checkmark \checkmark$\checkmark^\checkmark\\checkmarkc\checkmarki\checkmarkr\checkmarkc\checkmark$\checkmarks\checkmark$\checkmark^\checkmark\\checkmarkc\checkmarki\checkmarkr\checkmarkc\checkmark$\checkmarkt\checkmark$\checkmark^\checkmark\\checkmarkc\checkmarki\checkmarkr\checkmarkc\checkmark$\checkmarka\checkmark$\checkmark^\checkmark\\checkmarkc\checkmarki\checkmarkr\checkmarkc\checkmark$\checkmarkt\checkmark$\checkmark^\checkmark\\checkmarkc\checkmarki\checkmarkr\checkmarkc\checkmark$\checkmarki\checkmark$\checkmark^\checkmark\\checkmarkc\checkmarki\checkmarkr\checkmarkc\checkmark$\checkmarkq\checkmark$\checkmark^\checkmark\\checkmarkc\checkmarki\checkmarkr\checkmarkc\checkmark$\checkmarku\checkmark$\checkmark^\checkmark\\checkmarkc\checkmarki\checkmarkr\checkmarkc\checkmark$\checkmarke\checkmark$\checkmark^\checkmark\\checkmarkc\checkmarki\checkmarkr\checkmarkc\checkmark$\checkmark,\checkmark$\checkmark^\checkmark\\checkmarkc\checkmarki\checkmarkr\checkmarkc\checkmark$\checkmark \checkmark$\checkmark^\checkmark\\checkmarkc\checkmarki\checkmarkr\checkmarkc\checkmark$\checkmark$\checkmark$\checkmark^\checkmark\\checkmarkc\checkmarki\checkmarkr\checkmarkc\checkmark$\checkmarks\checkmark$\checkmark^\checkmark\\checkmarkc\checkmarki\checkmarkr\checkmarkc\checkmark$\checkmark \checkmark$\checkmark^\checkmark\\checkmarkc\checkmarki\checkmarkr\checkmarkc\checkmark$\checkmark\\checkmark$\checkmark^\checkmark\\checkmarkc\checkmarki\checkmarkr\checkmarkc\checkmark$\checkmarkg\checkmark$\checkmark^\checkmark\\checkmarkc\checkmarki\checkmarkr\checkmarkc\checkmark$\checkmarke\checkmark$\checkmark^\checkmark\\checkmarkc\checkmarki\checkmarkr\checkmarkc\checkmark$\checkmarkq\checkmark$\checkmark^\checkmark\\checkmarkc\checkmarki\checkmarkr\checkmarkc\checkmark$\checkmark \checkmark$\checkmark^\checkmark\\checkmarkc\checkmarki\checkmarkr\checkmarkc\checkmark$\checkmark2\checkmark$\checkmark^\checkmark\\checkmarkc\checkmarki\checkmarkr\checkmarkc\checkmark$\checkmark{\checkmark$\checkmark^\checkmark\\checkmarkc\checkmarki\checkmarkr\checkmarkc\checkmark$\checkmark,\checkmark$\checkmark^\checkmark\\checkmarkc\checkmarki\checkmarkr\checkmarkc\checkmark$\checkmark}\checkmark$\checkmark^\checkmark\\checkmarkc\checkmarki\checkmarkr\checkmarkc\checkmark$\checkmark5\checkmark$\checkmark^\checkmark\\checkmarkc\checkmarki\checkmarkr\checkmarkc\checkmark$\checkmark$\checkmark$\checkmark^\checkmark\\checkmarkc\checkmarki\checkmarkr\checkmarkc\checkmark$\checkmark \checkmark$\checkmark^\checkmark\\checkmarkc\checkmarki\checkmarkr\checkmarkc\checkmark$\checkmarkp\checkmark$\checkmark^\checkmark\\checkmarkc\checkmarki\checkmarkr\checkmarkc\checkmark$\checkmarko\checkmark$\checkmark^\checkmark\\checkmarkc\checkmarki\checkmarkr\checkmarkc\checkmark$\checkmarku\checkmark$\checkmark^\checkmark\\checkmarkc\checkmarki\checkmarkr\checkmarkc\checkmark$\checkmarkr\checkmark$\checkmark^\checkmark\\checkmarkc\checkmarki\checkmarkr\checkmarkc\checkmark$\checkmark \checkmark$\checkmark^\checkmark\\checkmarkc\checkmarki\checkmarkr\checkmarkc\checkmark$\checkmarkl\checkmark$\checkmark^\checkmark\\checkmarkc\checkmarki\checkmarkr\checkmarkc\checkmark$\checkmarka\checkmark$\checkmark^\checkmark\\checkmarkc\checkmarki\checkmarkr\checkmarkc\checkmark$\checkmark \checkmark$\checkmark^\checkmark\\checkmarkc\checkmarki\checkmarkr\checkmarkc\checkmark$\checkmarkf\checkmark$\checkmark^\checkmark\\checkmarkc\checkmarki\checkmarkr\checkmarkc\checkmark$\checkmarka\checkmark$\checkmark^\checkmark\\checkmarkc\checkmarki\checkmarkr\checkmarkc\checkmark$\checkmarkt\checkmark$\checkmark^\checkmark\\checkmarkc\checkmarki\checkmarkr\checkmarkc\checkmark$\checkmarki\checkmark$\checkmark^\checkmark\\checkmarkc\checkmarki\checkmarkr\checkmarkc\checkmark$\checkmarkg\checkmark$\checkmark^\checkmark\\checkmarkc\checkmarki\checkmarkr\checkmarkc\checkmark$\checkmarku\checkmark$\checkmark^\checkmark\\checkmarkc\checkmarki\checkmarkr\checkmarkc\checkmark$\checkmarke\checkmark$\checkmark^\checkmark\\checkmarkc\checkmarki\checkmarkr\checkmarkc\checkmark$\checkmark \checkmark$\checkmark^\checkmark\\checkmarkc\checkmarki\checkmarkr\checkmarkc\checkmark$\checkmarke\checkmark$\checkmark^\checkmark\\checkmarkc\checkmarki\checkmarkr\checkmarkc\checkmark$\checkmarkt\checkmark$\checkmark^\checkmark\\checkmarkc\checkmarki\checkmarkr\checkmarkc\checkmark$\checkmark \checkmark$\checkmark^\checkmark\\checkmarkc\checkmarki\checkmarkr\checkmarkc\checkmark$\checkmarkl\checkmark$\checkmark^\checkmark\\checkmarkc\checkmarki\checkmarkr\checkmarkc\checkmark$\checkmarka\checkmark$\checkmark^\checkmark\\checkmarkc\checkmarki\checkmarkr\checkmarkc\checkmark$\checkmark \checkmark$\checkmark^\checkmark\\checkmarkc\checkmarki\checkmarkr\checkmarkc\checkmark$\checkmarkd\checkmark$\checkmark^\checkmark\\checkmarkc\checkmarki\checkmarkr\checkmarkc\checkmark$\checkmarky\checkmark$\checkmark^\checkmark\\checkmarkc\checkmarki\checkmarkr\checkmarkc\checkmark$\checkmarkn\checkmark$\checkmark^\checkmark\\checkmarkc\checkmarki\checkmarkr\checkmarkc\checkmark$\checkmarka\checkmark$\checkmark^\checkmark\\checkmarkc\checkmarki\checkmarkr\checkmarkc\checkmark$\checkmarkm\checkmark$\checkmark^\checkmark\\checkmarkc\checkmarki\checkmarkr\checkmarkc\checkmark$\checkmarki\checkmark$\checkmark^\checkmark\\checkmarkc\checkmarki\checkmarkr\checkmarkc\checkmark$\checkmarkq\checkmark$\checkmark^\checkmark\\checkmarkc\checkmarki\checkmarkr\checkmarkc\checkmark$\checkmarku\checkmark$\checkmark^\checkmark\\checkmarkc\checkmarki\checkmarkr\checkmarkc\checkmark$\checkmarke\checkmark$\checkmark^\checkmark\\checkmarkc\checkmarki\checkmarkr\checkmarkc\checkmark$\checkmark)\checkmark$\checkmark^\checkmark\\checkmarkc\checkmarki\checkmarkr\checkmarkc\checkmark$\checkmark.\checkmark$\checkmark^\checkmark\\checkmarkc\checkmarki\checkmarkr\checkmarkc\checkmark$\checkmark
\checkmark$\checkmark^\checkmark\\checkmarkc\checkmarki\checkmarkr\checkmarkc\checkmark$\checkmark
\checkmark$\checkmark^\checkmark\\checkmarkc\checkmarki\checkmarkr\checkmarkc\checkmark$\checkmark%\checkmark$\checkmark^\checkmark\\checkmarkc\checkmarki\checkmarkr\checkmarkc\checkmark$\checkmark \checkmark$\checkmark^\checkmark\\checkmarkc\checkmarki\checkmarkr\checkmarkc\checkmark$\checkmark=\checkmark$\checkmark^\checkmark\\checkmarkc\checkmarki\checkmarkr\checkmarkc\checkmark$\checkmark=\checkmark$\checkmark^\checkmark\\checkmarkc\checkmarki\checkmarkr\checkmarkc\checkmark$\checkmark=\checkmark$\checkmark^\checkmark\\checkmarkc\checkmarki\checkmarkr\checkmarkc\checkmark$\checkmark=\checkmark$\checkmark^\checkmark\\checkmarkc\checkmarki\checkmarkr\checkmarkc\checkmark$\checkmark=\checkmark$\checkmark^\checkmark\\checkmarkc\checkmarki\checkmarkr\checkmarkc\checkmark$\checkmark=\checkmark$\checkmark^\checkmark\\checkmarkc\checkmarki\checkmarkr\checkmarkc\checkmark$\checkmark=\checkmark$\checkmark^\checkmark\\checkmarkc\checkmarki\checkmarkr\checkmarkc\checkmark$\checkmark=\checkmark$\checkmark^\checkmark\\checkmarkc\checkmarki\checkmarkr\checkmarkc\checkmark$\checkmark=\checkmark$\checkmark^\checkmark\\checkmarkc\checkmarki\checkmarkr\checkmarkc\checkmark$\checkmark=\checkmark$\checkmark^\checkmark\\checkmarkc\checkmarki\checkmarkr\checkmarkc\checkmark$\checkmark=\checkmark$\checkmark^\checkmark\\checkmarkc\checkmarki\checkmarkr\checkmarkc\checkmark$\checkmark=\checkmark$\checkmark^\checkmark\\checkmarkc\checkmarki\checkmarkr\checkmarkc\checkmark$\checkmark=\checkmark$\checkmark^\checkmark\\checkmarkc\checkmarki\checkmarkr\checkmarkc\checkmark$\checkmark=\checkmark$\checkmark^\checkmark\\checkmarkc\checkmarki\checkmarkr\checkmarkc\checkmark$\checkmark=\checkmark$\checkmark^\checkmark\\checkmarkc\checkmarki\checkmarkr\checkmarkc\checkmark$\checkmark=\checkmark$\checkmark^\checkmark\\checkmarkc\checkmarki\checkmarkr\checkmarkc\checkmark$\checkmark=\checkmark$\checkmark^\checkmark\\checkmarkc\checkmarki\checkmarkr\checkmarkc\checkmark$\checkmark=\checkmark$\checkmark^\checkmark\\checkmarkc\checkmarki\checkmarkr\checkmarkc\checkmark$\checkmark=\checkmark$\checkmark^\checkmark\\checkmarkc\checkmarki\checkmarkr\checkmarkc\checkmark$\checkmark=\checkmark$\checkmark^\checkmark\\checkmarkc\checkmarki\checkmarkr\checkmarkc\checkmark$\checkmark=\checkmark$\checkmark^\checkmark\\checkmarkc\checkmarki\checkmarkr\checkmarkc\checkmark$\checkmark=\checkmark$\checkmark^\checkmark\\checkmarkc\checkmarki\checkmarkr\checkmarkc\checkmark$\checkmark=\checkmark$\checkmark^\checkmark\\checkmarkc\checkmarki\checkmarkr\checkmarkc\checkmark$\checkmark=\checkmark$\checkmark^\checkmark\\checkmarkc\checkmarki\checkmarkr\checkmarkc\checkmark$\checkmark=\checkmark$\checkmark^\checkmark\\checkmarkc\checkmarki\checkmarkr\checkmarkc\checkmark$\checkmark=\checkmark$\checkmark^\checkmark\\checkmarkc\checkmarki\checkmarkr\checkmarkc\checkmark$\checkmark=\checkmark$\checkmark^\checkmark\\checkmarkc\checkmarki\checkmarkr\checkmarkc\checkmark$\checkmark=\checkmark$\checkmark^\checkmark\\checkmarkc\checkmarki\checkmarkr\checkmarkc\checkmark$\checkmark=\checkmark$\checkmark^\checkmark\\checkmarkc\checkmarki\checkmarkr\checkmarkc\checkmark$\checkmark=\checkmark$\checkmark^\checkmark\\checkmarkc\checkmarki\checkmarkr\checkmarkc\checkmark$\checkmark=\checkmark$\checkmark^\checkmark\\checkmarkc\checkmarki\checkmarkr\checkmarkc\checkmark$\checkmark=\checkmark$\checkmark^\checkmark\\checkmarkc\checkmarki\checkmarkr\checkmarkc\checkmark$\checkmark=\checkmark$\checkmark^\checkmark\\checkmarkc\checkmarki\checkmarkr\checkmarkc\checkmark$\checkmark=\checkmark$\checkmark^\checkmark\\checkmarkc\checkmarki\checkmarkr\checkmarkc\checkmark$\checkmark=\checkmark$\checkmark^\checkmark\\checkmarkc\checkmarki\checkmarkr\checkmarkc\checkmark$\checkmark=\checkmark$\checkmark^\checkmark\\checkmarkc\checkmarki\checkmarkr\checkmarkc\checkmark$\checkmark=\checkmark$\checkmark^\checkmark\\checkmarkc\checkmarki\checkmarkr\checkmarkc\checkmark$\checkmark=\checkmark$\checkmark^\checkmark\\checkmarkc\checkmarki\checkmarkr\checkmarkc\checkmark$\checkmark=\checkmark$\checkmark^\checkmark\\checkmarkc\checkmarki\checkmarkr\checkmarkc\checkmark$\checkmark=\checkmark$\checkmark^\checkmark\\checkmarkc\checkmarki\checkmarkr\checkmarkc\checkmark$\checkmark=\checkmark$\checkmark^\checkmark\\checkmarkc\checkmarki\checkmarkr\checkmarkc\checkmark$\checkmark=\checkmark$\checkmark^\checkmark\\checkmarkc\checkmarki\checkmarkr\checkmarkc\checkmark$\checkmark=\checkmark$\checkmark^\checkmark\\checkmarkc\checkmarki\checkmarkr\checkmarkc\checkmark$\checkmark=\checkmark$\checkmark^\checkmark\\checkmarkc\checkmarki\checkmarkr\checkmarkc\checkmark$\checkmark=\checkmark$\checkmark^\checkmark\\checkmarkc\checkmarki\checkmarkr\checkmarkc\checkmark$\checkmark=\checkmark$\checkmark^\checkmark\\checkmarkc\checkmarki\checkmarkr\checkmarkc\checkmark$\checkmark=\checkmark$\checkmark^\checkmark\\checkmarkc\checkmarki\checkmarkr\checkmarkc\checkmark$\checkmark=\checkmark$\checkmark^\checkmark\\checkmarkc\checkmarki\checkmarkr\checkmarkc\checkmark$\checkmark=\checkmark$\checkmark^\checkmark\\checkmarkc\checkmarki\checkmarkr\checkmarkc\checkmark$\checkmark=\checkmark$\checkmark^\checkmark\\checkmarkc\checkmarki\checkmarkr\checkmarkc\checkmark$\checkmark=\checkmark$\checkmark^\checkmark\\checkmarkc\checkmarki\checkmarkr\checkmarkc\checkmark$\checkmark=\checkmark$\checkmark^\checkmark\\checkmarkc\checkmarki\checkmarkr\checkmarkc\checkmark$\checkmark=\checkmark$\checkmark^\checkmark\\checkmarkc\checkmarki\checkmarkr\checkmarkc\checkmark$\checkmark=\checkmark$\checkmark^\checkmark\\checkmarkc\checkmarki\checkmarkr\checkmarkc\checkmark$\checkmark=\checkmark$\checkmark^\checkmark\\checkmarkc\checkmarki\checkmarkr\checkmarkc\checkmark$\checkmark=\checkmark$\checkmark^\checkmark\\checkmarkc\checkmarki\checkmarkr\checkmarkc\checkmark$\checkmark=\checkmark$\checkmark^\checkmark\\checkmarkc\checkmarki\checkmarkr\checkmarkc\checkmark$\checkmark=\checkmark$\checkmark^\checkmark\\checkmarkc\checkmarki\checkmarkr\checkmarkc\checkmark$\checkmark=\checkmark$\checkmark^\checkmark\\checkmarkc\checkmarki\checkmarkr\checkmarkc\checkmark$\checkmark=\checkmark$\checkmark^\checkmark\\checkmarkc\checkmarki\checkmarkr\checkmarkc\checkmark$\checkmark=\checkmark$\checkmark^\checkmark\\checkmarkc\checkmarki\checkmarkr\checkmarkc\checkmark$\checkmark=\checkmark$\checkmark^\checkmark\\checkmarkc\checkmarki\checkmarkr\checkmarkc\checkmark$\checkmark
\checkmark$\checkmark^\checkmark\\checkmarkc\checkmarki\checkmarkr\checkmarkc\checkmark$\checkmark\\checkmark$\checkmark^\checkmark\\checkmarkc\checkmarki\checkmarkr\checkmarkc\checkmark$\checkmarks\checkmark$\checkmark^\checkmark\\checkmarkc\checkmarki\checkmarkr\checkmarkc\checkmark$\checkmarke\checkmark$\checkmark^\checkmark\\checkmarkc\checkmarki\checkmarkr\checkmarkc\checkmark$\checkmarkc\checkmark$\checkmark^\checkmark\\checkmarkc\checkmarki\checkmarkr\checkmarkc\checkmark$\checkmarkt\checkmark$\checkmark^\checkmark\\checkmarkc\checkmarki\checkmarkr\checkmarkc\checkmark$\checkmarki\checkmark$\checkmark^\checkmark\\checkmarkc\checkmarki\checkmarkr\checkmarkc\checkmark$\checkmarko\checkmark$\checkmark^\checkmark\\checkmarkc\checkmarki\checkmarkr\checkmarkc\checkmark$\checkmarkn\checkmark$\checkmark^\checkmark\\checkmarkc\checkmarki\checkmarkr\checkmarkc\checkmark$\checkmark{\checkmark$\checkmark^\checkmark\\checkmarkc\checkmarki\checkmarkr\checkmarkc\checkmark$\checkmarkA\checkmark$\checkmark^\checkmark\\checkmarkc\checkmarki\checkmarkr\checkmarkc\checkmark$\checkmarkn\checkmark$\checkmark^\checkmark\\checkmarkc\checkmarki\checkmarkr\checkmarkc\checkmark$\checkmarka\checkmark$\checkmark^\checkmark\\checkmarkc\checkmarki\checkmarkr\checkmarkc\checkmark$\checkmarkl\checkmark$\checkmark^\checkmark\\checkmarkc\checkmarki\checkmarkr\checkmarkc\checkmark$\checkmarky\checkmark$\checkmark^\checkmark\\checkmarkc\checkmarki\checkmarkr\checkmarkc\checkmark$\checkmarks\checkmark$\checkmark^\checkmark\\checkmarkc\checkmarki\checkmarkr\checkmarkc\checkmark$\checkmarke\checkmark$\checkmark^\checkmark\\checkmarkc\checkmarki\checkmarkr\checkmarkc\checkmark$\checkmark \checkmark$\checkmark^\checkmark\\checkmarkc\checkmarki\checkmarkr\checkmarkc\checkmark$\checkmarkF\checkmark$\checkmark^\checkmark\\checkmarkc\checkmarki\checkmarkr\checkmarkc\checkmark$\checkmarko\checkmark$\checkmark^\checkmark\\checkmarkc\checkmarki\checkmarkr\checkmarkc\checkmark$\checkmarkn\checkmark$\checkmark^\checkmark\\checkmarkc\checkmarki\checkmarkr\checkmarkc\checkmark$\checkmarkc\checkmark$\checkmark^\checkmark\\checkmarkc\checkmarki\checkmarkr\checkmarkc\checkmark$\checkmarkt\checkmark$\checkmark^\checkmark\\checkmarkc\checkmarki\checkmarkr\checkmarkc\checkmark$\checkmarki\checkmark$\checkmark^\checkmark\\checkmarkc\checkmarki\checkmarkr\checkmarkc\checkmark$\checkmarko\checkmark$\checkmark^\checkmark\\checkmarkc\checkmarki\checkmarkr\checkmarkc\checkmark$\checkmarkn\checkmark$\checkmark^\checkmark\\checkmarkc\checkmarki\checkmarkr\checkmarkc\checkmark$\checkmarkn\checkmark$\checkmark^\checkmark\\checkmarkc\checkmarki\checkmarkr\checkmarkc\checkmark$\checkmarke\checkmark$\checkmark^\checkmark\\checkmarkc\checkmarki\checkmarkr\checkmarkc\checkmark$\checkmarkl\checkmark$\checkmark^\checkmark\\checkmarkc\checkmarki\checkmarkr\checkmarkc\checkmark$\checkmarkl\checkmark$\checkmark^\checkmark\\checkmarkc\checkmarki\checkmarkr\checkmarkc\checkmark$\checkmarke\checkmark$\checkmark^\checkmark\\checkmarkc\checkmarki\checkmarkr\checkmarkc\checkmark$\checkmark \checkmark$\checkmark^\checkmark\\checkmarkc\checkmarki\checkmarkr\checkmarkc\checkmark$\checkmarkd\checkmark$\checkmark^\checkmark\\checkmarkc\checkmarki\checkmarkr\checkmarkc\checkmark$\checkmarku\checkmark$\checkmark^\checkmark\\checkmarkc\checkmarki\checkmarkr\checkmarkc\checkmark$\checkmark \checkmark$\checkmark^\checkmark\\checkmarkc\checkmarki\checkmarkr\checkmarkc\checkmark$\checkmarkB\checkmark$\checkmark^\checkmark\\checkmarkc\checkmarki\checkmarkr\checkmarkc\checkmark$\checkmarke\checkmark$\checkmark^\checkmark\\checkmarkc\checkmarki\checkmarkr\checkmarkc\checkmark$\checkmarks\checkmark$\checkmark^\checkmark\\checkmarkc\checkmarki\checkmarkr\checkmarkc\checkmark$\checkmarko\checkmark$\checkmark^\checkmark\\checkmarkc\checkmarki\checkmarkr\checkmarkc\checkmark$\checkmarki\checkmark$\checkmark^\checkmark\\checkmarkc\checkmarki\checkmarkr\checkmarkc\checkmark$\checkmarkn\checkmark$\checkmark^\checkmark\\checkmarkc\checkmarki\checkmarkr\checkmarkc\checkmark$\checkmark \checkmark$\checkmark^\checkmark\\checkmarkc\checkmarki\checkmarkr\checkmarkc\checkmark$\checkmark(\checkmark$\checkmark^\checkmark\\checkmarkc\checkmarki\checkmarkr\checkmarkc\checkmark$\checkmarkA\checkmark$\checkmark^\checkmark\\checkmarkc\checkmarki\checkmarkr\checkmarkc\checkmark$\checkmarkF\checkmark$\checkmark^\checkmark\\checkmarkc\checkmarki\checkmarkr\checkmarkc\checkmark$\checkmarkB\checkmark$\checkmark^\checkmark\\checkmarkc\checkmarki\checkmarkr\checkmarkc\checkmark$\checkmark)\checkmark$\checkmark^\checkmark\\checkmarkc\checkmarki\checkmarkr\checkmarkc\checkmark$\checkmark}\checkmark$\checkmark^\checkmark\\checkmarkc\checkmarki\checkmarkr\checkmarkc\checkmark$\checkmark
\checkmark$\checkmark^\checkmark\\checkmarkc\checkmarki\checkmarkr\checkmarkc\checkmark$\checkmark
\checkmark$\checkmark^\checkmark\\checkmarkc\checkmarki\checkmarkr\checkmarkc\checkmark$\checkmark\\checkmark$\checkmark^\checkmark\\checkmarkc\checkmarki\checkmarkr\checkmarkc\checkmark$\checkmarks\checkmark$\checkmark^\checkmark\\checkmarkc\checkmarki\checkmarkr\checkmarkc\checkmark$\checkmarku\checkmark$\checkmark^\checkmark\\checkmarkc\checkmarki\checkmarkr\checkmarkc\checkmark$\checkmarkb\checkmark$\checkmark^\checkmark\\checkmarkc\checkmarki\checkmarkr\checkmarkc\checkmark$\checkmarks\checkmark$\checkmark^\checkmark\\checkmarkc\checkmarki\checkmarkr\checkmarkc\checkmark$\checkmarke\checkmark$\checkmark^\checkmark\\checkmarkc\checkmarki\checkmarkr\checkmarkc\checkmark$\checkmarkc\checkmark$\checkmark^\checkmark\\checkmarkc\checkmarki\checkmarkr\checkmarkc\checkmark$\checkmarkt\checkmark$\checkmark^\checkmark\\checkmarkc\checkmarki\checkmarkr\checkmarkc\checkmark$\checkmarki\checkmark$\checkmark^\checkmark\\checkmarkc\checkmarki\checkmarkr\checkmarkc\checkmark$\checkmarko\checkmark$\checkmark^\checkmark\\checkmarkc\checkmarki\checkmarkr\checkmarkc\checkmark$\checkmarkn\checkmark$\checkmark^\checkmark\\checkmarkc\checkmarki\checkmarkr\checkmarkc\checkmark$\checkmark{\checkmark$\checkmark^\checkmark\\checkmarkc\checkmarki\checkmarkr\checkmarkc\checkmark$\checkmarkE\checkmark$\checkmark^\checkmark\\checkmarkc\checkmarki\checkmarkr\checkmarkc\checkmark$\checkmarkx\checkmark$\checkmark^\checkmark\\checkmarkc\checkmarki\checkmarkr\checkmarkc\checkmark$\checkmarkp\checkmark$\checkmark^\checkmark\\checkmarkc\checkmarki\checkmarkr\checkmarkc\checkmark$\checkmarkr\checkmark$\checkmark^\checkmark\\checkmarkc\checkmarki\checkmarkr\checkmarkc\checkmark$\checkmarke\checkmark$\checkmark^\checkmark\\checkmarkc\checkmarki\checkmarkr\checkmarkc\checkmark$\checkmarks\checkmark$\checkmark^\checkmark\\checkmarkc\checkmarki\checkmarkr\checkmarkc\checkmark$\checkmarks\checkmark$\checkmark^\checkmark\\checkmarkc\checkmarki\checkmarkr\checkmarkc\checkmark$\checkmarki\checkmark$\checkmark^\checkmark\\checkmarkc\checkmarki\checkmarkr\checkmarkc\checkmark$\checkmarko\checkmark$\checkmark^\checkmark\\checkmarkc\checkmarki\checkmarkr\checkmarkc\checkmark$\checkmarkn\checkmark$\checkmark^\checkmark\\checkmarkc\checkmarki\checkmarkr\checkmarkc\checkmark$\checkmark \checkmark$\checkmark^\checkmark\\checkmarkc\checkmarki\checkmarkr\checkmarkc\checkmark$\checkmarkd\checkmark$\checkmark^\checkmark\\checkmarkc\checkmarki\checkmarkr\checkmarkc\checkmark$\checkmarku\checkmark$\checkmark^\checkmark\\checkmarkc\checkmarki\checkmarkr\checkmarkc\checkmark$\checkmark \checkmark$\checkmark^\checkmark\\checkmarkc\checkmarki\checkmarkr\checkmarkc\checkmark$\checkmarkb\checkmark$\checkmark^\checkmark\\checkmarkc\checkmarki\checkmarkr\checkmarkc\checkmark$\checkmarke\checkmark$\checkmark^\checkmark\\checkmarkc\checkmarki\checkmarkr\checkmarkc\checkmark$\checkmarks\checkmark$\checkmark^\checkmark\\checkmarkc\checkmarki\checkmarkr\checkmarkc\checkmark$\checkmarko\checkmark$\checkmark^\checkmark\\checkmarkc\checkmarki\checkmarkr\checkmarkc\checkmark$\checkmarki\checkmark$\checkmark^\checkmark\\checkmarkc\checkmarki\checkmarkr\checkmarkc\checkmark$\checkmarkn\checkmark$\checkmark^\checkmark\\checkmarkc\checkmarki\checkmarkr\checkmarkc\checkmark$\checkmark \checkmark$\checkmark^\checkmark\\checkmarkc\checkmarki\checkmarkr\checkmarkc\checkmark$\checkmark:\checkmark$\checkmark^\checkmark\\checkmarkc\checkmarki\checkmarkr\checkmarkc\checkmark$\checkmark \checkmark$\checkmark^\checkmark\\checkmarkc\checkmarki\checkmarkr\checkmarkc\checkmark$\checkmarkD\checkmark$\checkmark^\checkmark\\checkmarkc\checkmarki\checkmarkr\checkmarkc\checkmark$\checkmarki\checkmark$\checkmark^\checkmark\\checkmarkc\checkmarki\checkmarkr\checkmarkc\checkmark$\checkmarka\checkmark$\checkmark^\checkmark\\checkmarkc\checkmarki\checkmarkr\checkmarkc\checkmark$\checkmarkg\checkmark$\checkmark^\checkmark\\checkmarkc\checkmarki\checkmarkr\checkmarkc\checkmark$\checkmarkr\checkmark$\checkmark^\checkmark\\checkmarkc\checkmarki\checkmarkr\checkmarkc\checkmark$\checkmarka\checkmark$\checkmark^\checkmark\\checkmarkc\checkmarki\checkmarkr\checkmarkc\checkmark$\checkmarkm\checkmark$\checkmark^\checkmark\\checkmarkc\checkmarki\checkmarkr\checkmarkc\checkmark$\checkmarkm\checkmark$\checkmark^\checkmark\\checkmarkc\checkmarki\checkmarkr\checkmarkc\checkmark$\checkmarke\checkmark$\checkmark^\checkmark\\checkmarkc\checkmarki\checkmarkr\checkmarkc\checkmark$\checkmark \checkmark$\checkmark^\checkmark\\checkmarkc\checkmarki\checkmarkr\checkmarkc\checkmark$\checkmark«\checkmark$\checkmark^\checkmark\\checkmarkc\checkmarki\checkmarkr\checkmarkc\checkmark$\checkmarkB\checkmark$\checkmark^\checkmark\\checkmarkc\checkmarki\checkmarkr\checkmarkc\checkmark$\checkmarkê\checkmark$\checkmark^\checkmark\\checkmarkc\checkmarki\checkmarkr\checkmarkc\checkmark$\checkmarkt\checkmark$\checkmark^\checkmark\\checkmarkc\checkmarki\checkmarkr\checkmarkc\checkmark$\checkmarke\checkmark$\checkmark^\checkmark\\checkmarkc\checkmarki\checkmarkr\checkmarkc\checkmark$\checkmark \checkmark$\checkmark^\checkmark\\checkmarkc\checkmarki\checkmarkr\checkmarkc\checkmark$\checkmarkà\checkmark$\checkmark^\checkmark\\checkmarkc\checkmarki\checkmarkr\checkmarkc\checkmark$\checkmark \checkmark$\checkmark^\checkmark\\checkmarkc\checkmarki\checkmarkr\checkmarkc\checkmark$\checkmarkC\checkmark$\checkmark^\checkmark\\checkmarkc\checkmarki\checkmarkr\checkmarkc\checkmark$\checkmarko\checkmark$\checkmark^\checkmark\\checkmarkc\checkmarki\checkmarkr\checkmarkc\checkmark$\checkmarkr\checkmark$\checkmark^\checkmark\\checkmarkc\checkmarki\checkmarkr\checkmarkc\checkmark$\checkmarkn\checkmark$\checkmark^\checkmark\\checkmarkc\checkmarki\checkmarkr\checkmarkc\checkmark$\checkmarke\checkmark$\checkmark^\checkmark\\checkmarkc\checkmarki\checkmarkr\checkmarkc\checkmark$\checkmarks\checkmark$\checkmark^\checkmark\\checkmarkc\checkmarki\checkmarkr\checkmarkc\checkmark$\checkmark»\checkmark$\checkmark^\checkmark\\checkmarkc\checkmarki\checkmarkr\checkmarkc\checkmark$\checkmark}\checkmark$\checkmark^\checkmark\\checkmarkc\checkmarki\checkmarkr\checkmarkc\checkmark$\checkmark
\checkmark$\checkmark^\checkmark\\checkmarkc\checkmarki\checkmarkr\checkmarkc\checkmark$\checkmarkL\checkmark$\checkmark^\checkmark\\checkmarkc\checkmarki\checkmarkr\checkmarkc\checkmark$\checkmarka\checkmark$\checkmark^\checkmark\\checkmarkc\checkmarki\checkmarkr\checkmarkc\checkmark$\checkmark \checkmark$\checkmark^\checkmark\\checkmarkc\checkmarki\checkmarkr\checkmarkc\checkmark$\checkmarkr\checkmark$\checkmark^\checkmark\\checkmarkc\checkmarki\checkmarkr\checkmarkc\checkmark$\checkmarké\checkmark$\checkmark^\checkmark\\checkmarkc\checkmarki\checkmarkr\checkmarkc\checkmark$\checkmarkp\checkmark$\checkmark^\checkmark\\checkmarkc\checkmarki\checkmarkr\checkmarkc\checkmark$\checkmarko\checkmark$\checkmark^\checkmark\\checkmarkc\checkmarki\checkmarkr\checkmarkc\checkmark$\checkmarkn\checkmark$\checkmark^\checkmark\\checkmarkc\checkmarki\checkmarkr\checkmarkc\checkmark$\checkmarks\checkmark$\checkmark^\checkmark\\checkmarkc\checkmarki\checkmarkr\checkmarkc\checkmark$\checkmarke\checkmark$\checkmark^\checkmark\\checkmarkc\checkmarki\checkmarkr\checkmarkc\checkmark$\checkmark \checkmark$\checkmark^\checkmark\\checkmarkc\checkmarki\checkmarkr\checkmarkc\checkmark$\checkmarka\checkmark$\checkmark^\checkmark\\checkmarkc\checkmarki\checkmarkr\checkmarkc\checkmark$\checkmarku\checkmark$\checkmark^\checkmark\\checkmarkc\checkmarki\checkmarkr\checkmarkc\checkmark$\checkmarkx\checkmark$\checkmark^\checkmark\\checkmarkc\checkmarki\checkmarkr\checkmarkc\checkmark$\checkmark \checkmark$\checkmark^\checkmark\\checkmarkc\checkmarki\checkmarkr\checkmarkc\checkmark$\checkmarkt\checkmark$\checkmark^\checkmark\\checkmarkc\checkmarki\checkmarkr\checkmarkc\checkmark$\checkmarkr\checkmark$\checkmark^\checkmark\\checkmarkc\checkmarki\checkmarkr\checkmarkc\checkmark$\checkmarko\checkmark$\checkmark^\checkmark\\checkmarkc\checkmarki\checkmarkr\checkmarkc\checkmark$\checkmarki\checkmark$\checkmark^\checkmark\\checkmarkc\checkmarki\checkmarkr\checkmarkc\checkmark$\checkmarks\checkmark$\checkmark^\checkmark\\checkmarkc\checkmarki\checkmarkr\checkmarkc\checkmark$\checkmark \checkmark$\checkmark^\checkmark\\checkmarkc\checkmarki\checkmarkr\checkmarkc\checkmark$\checkmarkq\checkmark$\checkmark^\checkmark\\checkmarkc\checkmarki\checkmarkr\checkmarkc\checkmark$\checkmarku\checkmark$\checkmark^\checkmark\\checkmarkc\checkmarki\checkmarkr\checkmarkc\checkmark$\checkmarke\checkmark$\checkmark^\checkmark\\checkmarkc\checkmarki\checkmarkr\checkmarkc\checkmark$\checkmarks\checkmark$\checkmark^\checkmark\\checkmarkc\checkmarki\checkmarkr\checkmarkc\checkmark$\checkmarkt\checkmark$\checkmark^\checkmark\\checkmarkc\checkmarki\checkmarkr\checkmarkc\checkmark$\checkmarki\checkmark$\checkmark^\checkmark\\checkmarkc\checkmarki\checkmarkr\checkmarkc\checkmark$\checkmarko\checkmark$\checkmark^\checkmark\\checkmarkc\checkmarki\checkmarkr\checkmarkc\checkmark$\checkmarkn\checkmark$\checkmark^\checkmark\\checkmarkc\checkmarki\checkmarkr\checkmarkc\checkmark$\checkmarks\checkmark$\checkmark^\checkmark\\checkmarkc\checkmarki\checkmarkr\checkmarkc\checkmark$\checkmark \checkmark$\checkmark^\checkmark\\checkmarkc\checkmarki\checkmarkr\checkmarkc\checkmark$\checkmarkf\checkmark$\checkmark^\checkmark\\checkmarkc\checkmarki\checkmarkr\checkmarkc\checkmark$\checkmarko\checkmark$\checkmark^\checkmark\\checkmarkc\checkmarki\checkmarkr\checkmarkc\checkmark$\checkmarkn\checkmark$\checkmark^\checkmark\\checkmarkc\checkmarki\checkmarkr\checkmarkc\checkmark$\checkmarkd\checkmark$\checkmark^\checkmark\\checkmarkc\checkmarki\checkmarkr\checkmarkc\checkmark$\checkmarka\checkmark$\checkmark^\checkmark\\checkmarkc\checkmarki\checkmarkr\checkmarkc\checkmark$\checkmarkm\checkmark$\checkmark^\checkmark\\checkmarkc\checkmarki\checkmarkr\checkmarkc\checkmark$\checkmarke\checkmark$\checkmark^\checkmark\\checkmarkc\checkmarki\checkmarkr\checkmarkc\checkmark$\checkmarkn\checkmark$\checkmark^\checkmark\\checkmarkc\checkmarki\checkmarkr\checkmarkc\checkmark$\checkmarkt\checkmark$\checkmark^\checkmark\\checkmarkc\checkmarki\checkmarkr\checkmarkc\checkmark$\checkmarka\checkmark$\checkmark^\checkmark\\checkmarkc\checkmarki\checkmarkr\checkmarkc\checkmark$\checkmarkl\checkmark$\checkmark^\checkmark\\checkmarkc\checkmarki\checkmarkr\checkmarkc\checkmark$\checkmarke\checkmark$\checkmark^\checkmark\\checkmarkc\checkmarki\checkmarkr\checkmarkc\checkmark$\checkmarks\checkmark$\checkmark^\checkmark\\checkmarkc\checkmarki\checkmarkr\checkmarkc\checkmark$\checkmark \checkmark$\checkmark^\checkmark\\checkmarkc\checkmarki\checkmarkr\checkmarkc\checkmark$\checkmarkn\checkmark$\checkmark^\checkmark\\checkmarkc\checkmarki\checkmarkr\checkmarkc\checkmark$\checkmarko\checkmark$\checkmark^\checkmark\\checkmarkc\checkmarki\checkmarkr\checkmarkc\checkmark$\checkmarku\checkmark$\checkmark^\checkmark\\checkmarkc\checkmarki\checkmarkr\checkmarkc\checkmark$\checkmarks\checkmark$\checkmark^\checkmark\\checkmarkc\checkmarki\checkmarkr\checkmarkc\checkmark$\checkmark \checkmark$\checkmark^\checkmark\\checkmarkc\checkmarki\checkmarkr\checkmarkc\checkmark$\checkmarkp\checkmark$\checkmark^\checkmark\\checkmarkc\checkmarki\checkmarkr\checkmarkc\checkmark$\checkmarke\checkmark$\checkmark^\checkmark\\checkmarkc\checkmarki\checkmarkr\checkmarkc\checkmark$\checkmarkr\checkmark$\checkmark^\checkmark\\checkmarkc\checkmarki\checkmarkr\checkmarkc\checkmark$\checkmarkm\checkmark$\checkmark^\checkmark\\checkmarkc\checkmarki\checkmarkr\checkmarkc\checkmark$\checkmarke\checkmark$\checkmark^\checkmark\\checkmarkc\checkmarki\checkmarkr\checkmarkc\checkmark$\checkmarkt\checkmark$\checkmark^\checkmark\\checkmarkc\checkmarki\checkmarkr\checkmarkc\checkmark$\checkmark \checkmark$\checkmark^\checkmark\\checkmarkc\checkmarki\checkmarkr\checkmarkc\checkmark$\checkmarkd\checkmark$\checkmark^\checkmark\\checkmarkc\checkmarki\checkmarkr\checkmarkc\checkmark$\checkmark'\checkmark$\checkmark^\checkmark\\checkmarkc\checkmarki\checkmarkr\checkmarkc\checkmark$\checkmarke\checkmark$\checkmark^\checkmark\\checkmarkc\checkmarki\checkmarkr\checkmarkc\checkmark$\checkmarkx\checkmark$\checkmark^\checkmark\\checkmarkc\checkmarki\checkmarkr\checkmarkc\checkmark$\checkmarkp\checkmark$\checkmark^\checkmark\\checkmarkc\checkmarki\checkmarkr\checkmarkc\checkmark$\checkmarkr\checkmark$\checkmark^\checkmark\\checkmarkc\checkmarki\checkmarkr\checkmarkc\checkmark$\checkmarki\checkmark$\checkmark^\checkmark\\checkmarkc\checkmarki\checkmarkr\checkmarkc\checkmark$\checkmarkm\checkmark$\checkmark^\checkmark\\checkmarkc\checkmarki\checkmarkr\checkmarkc\checkmark$\checkmarke\checkmark$\checkmark^\checkmark\\checkmarkc\checkmarki\checkmarkr\checkmarkc\checkmark$\checkmarkr\checkmark$\checkmark^\checkmark\\checkmarkc\checkmarki\checkmarkr\checkmarkc\checkmark$\checkmark \checkmark$\checkmark^\checkmark\\checkmarkc\checkmarki\checkmarkr\checkmarkc\checkmark$\checkmarkl\checkmark$\checkmark^\checkmark\\checkmarkc\checkmarki\checkmarkr\checkmarkc\checkmark$\checkmarke\checkmark$\checkmark^\checkmark\\checkmarkc\checkmarki\checkmarkr\checkmarkc\checkmark$\checkmark \checkmark$\checkmark^\checkmark\\checkmarkc\checkmarki\checkmarkr\checkmarkc\checkmark$\checkmarkb\checkmark$\checkmark^\checkmark\\checkmarkc\checkmarki\checkmarkr\checkmarkc\checkmark$\checkmarke\checkmark$\checkmark^\checkmark\\checkmarkc\checkmarki\checkmarkr\checkmarkc\checkmark$\checkmarks\checkmark$\checkmark^\checkmark\\checkmarkc\checkmarki\checkmarkr\checkmarkc\checkmark$\checkmarko\checkmark$\checkmark^\checkmark\\checkmarkc\checkmarki\checkmarkr\checkmarkc\checkmark$\checkmarki\checkmark$\checkmark^\checkmark\\checkmarkc\checkmarki\checkmarkr\checkmarkc\checkmark$\checkmarkn\checkmark$\checkmark^\checkmark\\checkmarkc\checkmarki\checkmarkr\checkmarkc\checkmark$\checkmark \checkmark$\checkmark^\checkmark\\checkmarkc\checkmarki\checkmarkr\checkmarkc\checkmark$\checkmark:\checkmark$\checkmark^\checkmark\\checkmarkc\checkmarki\checkmarkr\checkmarkc\checkmark$\checkmark
\checkmark$\checkmark^\checkmark\\checkmarkc\checkmarki\checkmarkr\checkmarkc\checkmark$\checkmark\\checkmark$\checkmark^\checkmark\\checkmarkc\checkmarki\checkmarkr\checkmarkc\checkmark$\checkmarkb\checkmark$\checkmark^\checkmark\\checkmarkc\checkmarki\checkmarkr\checkmarkc\checkmark$\checkmarke\checkmark$\checkmark^\checkmark\\checkmarkc\checkmarki\checkmarkr\checkmarkc\checkmark$\checkmarkg\checkmark$\checkmark^\checkmark\\checkmarkc\checkmarki\checkmarkr\checkmarkc\checkmark$\checkmarki\checkmark$\checkmark^\checkmark\\checkmarkc\checkmarki\checkmarkr\checkmarkc\checkmark$\checkmarkn\checkmark$\checkmark^\checkmark\\checkmarkc\checkmarki\checkmarkr\checkmarkc\checkmark$\checkmark{\checkmark$\checkmark^\checkmark\\checkmarkc\checkmarki\checkmarkr\checkmarkc\checkmark$\checkmarki\checkmark$\checkmark^\checkmark\\checkmarkc\checkmarki\checkmarkr\checkmarkc\checkmark$\checkmarkt\checkmark$\checkmark^\checkmark\\checkmarkc\checkmarki\checkmarkr\checkmarkc\checkmark$\checkmarke\checkmark$\checkmark^\checkmark\\checkmarkc\checkmarki\checkmarkr\checkmarkc\checkmark$\checkmarkm\checkmark$\checkmark^\checkmark\\checkmarkc\checkmarki\checkmarkr\checkmarkc\checkmark$\checkmarki\checkmark$\checkmark^\checkmark\\checkmarkc\checkmarki\checkmarkr\checkmarkc\checkmark$\checkmarkz\checkmark$\checkmark^\checkmark\\checkmarkc\checkmarki\checkmarkr\checkmarkc\checkmark$\checkmarke\checkmark$\checkmark^\checkmark\\checkmarkc\checkmarki\checkmarkr\checkmarkc\checkmark$\checkmark}\checkmark$\checkmark^\checkmark\\checkmarkc\checkmarki\checkmarkr\checkmarkc\checkmark$\checkmark
\checkmark$\checkmark^\checkmark\\checkmarkc\checkmarki\checkmarkr\checkmarkc\checkmark$\checkmark \checkmark$\checkmark^\checkmark\\checkmarkc\checkmarki\checkmarkr\checkmarkc\checkmark$\checkmark \checkmark$\checkmark^\checkmark\\checkmarkc\checkmarki\checkmarkr\checkmarkc\checkmark$\checkmark \checkmark$\checkmark^\checkmark\\checkmarkc\checkmarki\checkmarkr\checkmarkc\checkmark$\checkmark \checkmark$\checkmark^\checkmark\\checkmarkc\checkmarki\checkmarkr\checkmarkc\checkmark$\checkmark\\checkmark$\checkmark^\checkmark\\checkmarkc\checkmarki\checkmarkr\checkmarkc\checkmark$\checkmarki\checkmark$\checkmark^\checkmark\\checkmarkc\checkmarki\checkmarkr\checkmarkc\checkmark$\checkmarkt\checkmark$\checkmark^\checkmark\\checkmarkc\checkmarki\checkmarkr\checkmarkc\checkmark$\checkmarke\checkmark$\checkmark^\checkmark\\checkmarkc\checkmarki\checkmarkr\checkmarkc\checkmark$\checkmarkm\checkmark$\checkmark^\checkmark\\checkmarkc\checkmarki\checkmarkr\checkmarkc\checkmark$\checkmark \checkmark$\checkmark^\checkmark\\checkmarkc\checkmarki\checkmarkr\checkmarkc\checkmark$\checkmark\\checkmark$\checkmark^\checkmark\\checkmarkc\checkmarki\checkmarkr\checkmarkc\checkmark$\checkmarkt\checkmark$\checkmark^\checkmark\\checkmarkc\checkmarki\checkmarkr\checkmarkc\checkmark$\checkmarke\checkmark$\checkmark^\checkmark\\checkmarkc\checkmarki\checkmarkr\checkmarkc\checkmark$\checkmarkx\checkmark$\checkmark^\checkmark\\checkmarkc\checkmarki\checkmarkr\checkmarkc\checkmark$\checkmarkt\checkmark$\checkmark^\checkmark\\checkmarkc\checkmarki\checkmarkr\checkmarkc\checkmark$\checkmarkb\checkmark$\checkmark^\checkmark\\checkmarkc\checkmarki\checkmarkr\checkmarkc\checkmark$\checkmarkf\checkmark$\checkmark^\checkmark\\checkmarkc\checkmarki\checkmarkr\checkmarkc\checkmark$\checkmark{\checkmark$\checkmark^\checkmark\\checkmarkc\checkmarki\checkmarkr\checkmarkc\checkmark$\checkmarkÀ\checkmark$\checkmark^\checkmark\\checkmarkc\checkmarki\checkmarkr\checkmarkc\checkmark$\checkmark \checkmark$\checkmark^\checkmark\\checkmarkc\checkmarki\checkmarkr\checkmarkc\checkmark$\checkmarkq\checkmark$\checkmark^\checkmark\\checkmarkc\checkmarki\checkmarkr\checkmarkc\checkmark$\checkmarku\checkmark$\checkmark^\checkmark\\checkmarkc\checkmarki\checkmarkr\checkmarkc\checkmark$\checkmarki\checkmark$\checkmark^\checkmark\\checkmarkc\checkmarki\checkmarkr\checkmarkc\checkmark$\checkmark \checkmark$\checkmark^\checkmark\\checkmarkc\checkmarki\checkmarkr\checkmarkc\checkmark$\checkmarkr\checkmark$\checkmark^\checkmark\\checkmarkc\checkmarki\checkmarkr\checkmarkc\checkmark$\checkmarke\checkmark$\checkmark^\checkmark\\checkmarkc\checkmarki\checkmarkr\checkmarkc\checkmark$\checkmarkn\checkmark$\checkmark^\checkmark\\checkmarkc\checkmarki\checkmarkr\checkmarkc\checkmark$\checkmarkd\checkmark$\checkmark^\checkmark\\checkmarkc\checkmarki\checkmarkr\checkmarkc\checkmark$\checkmark-\checkmark$\checkmark^\checkmark\\checkmarkc\checkmarki\checkmarkr\checkmarkc\checkmark$\checkmarki\checkmark$\checkmark^\checkmark\\checkmarkc\checkmarki\checkmarkr\checkmarkc\checkmark$\checkmarkl\checkmark$\checkmark^\checkmark\\checkmarkc\checkmarki\checkmarkr\checkmarkc\checkmark$\checkmark \checkmark$\checkmark^\checkmark\\checkmarkc\checkmarki\checkmarkr\checkmarkc\checkmark$\checkmarks\checkmark$\checkmark^\checkmark\\checkmarkc\checkmarki\checkmarkr\checkmarkc\checkmark$\checkmarke\checkmark$\checkmark^\checkmark\\checkmarkc\checkmarki\checkmarkr\checkmarkc\checkmark$\checkmarkr\checkmark$\checkmark^\checkmark\\checkmarkc\checkmarki\checkmarkr\checkmarkc\checkmark$\checkmarkv\checkmark$\checkmark^\checkmark\\checkmarkc\checkmarki\checkmarkr\checkmarkc\checkmark$\checkmarki\checkmark$\checkmark^\checkmark\\checkmarkc\checkmarki\checkmarkr\checkmarkc\checkmark$\checkmarkc\checkmark$\checkmark^\checkmark\\checkmarkc\checkmarki\checkmarkr\checkmarkc\checkmark$\checkmarke\checkmark$\checkmark^\checkmark\\checkmarkc\checkmarki\checkmarkr\checkmarkc\checkmark$\checkmark \checkmark$\checkmark^\checkmark\\checkmarkc\checkmarki\checkmarkr\checkmarkc\checkmark$\checkmark?\checkmark$\checkmark^\checkmark\\checkmarkc\checkmarki\checkmarkr\checkmarkc\checkmark$\checkmark}\checkmark$\checkmark^\checkmark\\checkmarkc\checkmarki\checkmarkr\checkmarkc\checkmark$\checkmark \checkmark$\checkmark^\checkmark\\checkmarkc\checkmarki\checkmarkr\checkmarkc\checkmark$\checkmarkÀ\checkmark$\checkmark^\checkmark\\checkmarkc\checkmarki\checkmarkr\checkmarkc\checkmark$\checkmark \checkmark$\checkmark^\checkmark\\checkmarkc\checkmarki\checkmarkr\checkmarkc\checkmark$\checkmarkl\checkmark$\checkmark^\checkmark\\checkmarkc\checkmarki\checkmarkr\checkmarkc\checkmark$\checkmark'\checkmark$\checkmark^\checkmark\\checkmarkc\checkmarki\checkmarkr\checkmarkc\checkmark$\checkmarki\checkmark$\checkmark^\checkmark\\checkmarkc\checkmarki\checkmarkr\checkmarkc\checkmark$\checkmarkn\checkmark$\checkmark^\checkmark\\checkmarkc\checkmarki\checkmarkr\checkmarkc\checkmark$\checkmarkg\checkmark$\checkmark^\checkmark\\checkmarkc\checkmarki\checkmarkr\checkmarkc\checkmark$\checkmarké\checkmark$\checkmark^\checkmark\\checkmarkc\checkmarki\checkmarkr\checkmarkc\checkmark$\checkmarkn\checkmark$\checkmark^\checkmark\\checkmarkc\checkmarki\checkmarkr\checkmarkc\checkmark$\checkmarki\checkmark$\checkmark^\checkmark\\checkmarkc\checkmarki\checkmarkr\checkmarkc\checkmark$\checkmarke\checkmark$\checkmark^\checkmark\\checkmarkc\checkmarki\checkmarkr\checkmarkc\checkmark$\checkmarku\checkmark$\checkmark^\checkmark\\checkmarkc\checkmarki\checkmarkr\checkmarkc\checkmark$\checkmarkr\checkmark$\checkmark^\checkmark\\checkmarkc\checkmarki\checkmarkr\checkmarkc\checkmark$\checkmark-\checkmark$\checkmark^\checkmark\\checkmarkc\checkmarki\checkmarkr\checkmarkc\checkmark$\checkmarkc\checkmark$\checkmark^\checkmark\\checkmarkc\checkmarki\checkmarkr\checkmarkc\checkmark$\checkmarko\checkmark$\checkmark^\checkmark\\checkmarkc\checkmarki\checkmarkr\checkmarkc\checkmark$\checkmarkn\checkmark$\checkmark^\checkmark\\checkmarkc\checkmarki\checkmarkr\checkmarkc\checkmark$\checkmarkc\checkmark$\checkmark^\checkmark\\checkmarkc\checkmarki\checkmarkr\checkmarkc\checkmark$\checkmarke\checkmark$\checkmark^\checkmark\\checkmarkc\checkmarki\checkmarkr\checkmarkc\checkmark$\checkmarkp\checkmark$\checkmark^\checkmark\\checkmarkc\checkmarki\checkmarkr\checkmarkc\checkmark$\checkmarkt\checkmark$\checkmark^\checkmark\\checkmarkc\checkmarki\checkmarkr\checkmarkc\checkmark$\checkmarke\checkmark$\checkmark^\checkmark\\checkmarkc\checkmarki\checkmarkr\checkmarkc\checkmark$\checkmarku\checkmark$\checkmark^\checkmark\\checkmarkc\checkmarki\checkmarkr\checkmarkc\checkmark$\checkmarkr\checkmark$\checkmark^\checkmark\\checkmarkc\checkmarki\checkmarkr\checkmarkc\checkmark$\checkmark \checkmark$\checkmark^\checkmark\\checkmarkc\checkmarki\checkmarkr\checkmarkc\checkmark$\checkmarke\checkmark$\checkmark^\checkmark\\checkmarkc\checkmarki\checkmarkr\checkmarkc\checkmark$\checkmarkt\checkmark$\checkmark^\checkmark\\checkmarkc\checkmarki\checkmarkr\checkmarkc\checkmark$\checkmark \checkmark$\checkmark^\checkmark\\checkmarkc\checkmarki\checkmarkr\checkmarkc\checkmark$\checkmarka\checkmark$\checkmark^\checkmark\\checkmarkc\checkmarki\checkmarkr\checkmarkc\checkmark$\checkmarku\checkmark$\checkmark^\checkmark\\checkmarkc\checkmarki\checkmarkr\checkmarkc\checkmark$\checkmark \checkmark$\checkmark^\checkmark\\checkmarkc\checkmarki\checkmarkr\checkmarkc\checkmark$\checkmarkt\checkmark$\checkmark^\checkmark\\checkmarkc\checkmarki\checkmarkr\checkmarkc\checkmark$\checkmarke\checkmark$\checkmark^\checkmark\\checkmarkc\checkmarki\checkmarkr\checkmarkc\checkmark$\checkmarkc\checkmark$\checkmark^\checkmark\\checkmarkc\checkmarki\checkmarkr\checkmarkc\checkmark$\checkmarkh\checkmark$\checkmark^\checkmark\\checkmarkc\checkmarki\checkmarkr\checkmarkc\checkmark$\checkmarkn\checkmark$\checkmark^\checkmark\\checkmarkc\checkmarki\checkmarkr\checkmarkc\checkmark$\checkmarki\checkmark$\checkmark^\checkmark\\checkmarkc\checkmarki\checkmarkr\checkmarkc\checkmark$\checkmarkc\checkmark$\checkmark^\checkmark\\checkmarkc\checkmarki\checkmarkr\checkmarkc\checkmark$\checkmarki\checkmark$\checkmark^\checkmark\\checkmarkc\checkmarki\checkmarkr\checkmarkc\checkmark$\checkmarke\checkmark$\checkmark^\checkmark\\checkmarkc\checkmarki\checkmarkr\checkmarkc\checkmark$\checkmarkn\checkmark$\checkmark^\checkmark\\checkmarkc\checkmarki\checkmarkr\checkmarkc\checkmark$\checkmark \checkmark$\checkmark^\checkmark\\checkmarkc\checkmarki\checkmarkr\checkmarkc\checkmark$\checkmarkd\checkmark$\checkmark^\checkmark\\checkmarkc\checkmarki\checkmarkr\checkmarkc\checkmark$\checkmark'\checkmark$\checkmark^\checkmark\\checkmarkc\checkmarki\checkmarkr\checkmarkc\checkmark$\checkmarka\checkmark$\checkmark^\checkmark\\checkmarkc\checkmarki\checkmarkr\checkmarkc\checkmark$\checkmarkt\checkmark$\checkmark^\checkmark\\checkmarkc\checkmarki\checkmarkr\checkmarkc\checkmark$\checkmarke\checkmark$\checkmark^\checkmark\\checkmarkc\checkmarki\checkmarkr\checkmarkc\checkmark$\checkmarkl\checkmark$\checkmark^\checkmark\\checkmarkc\checkmarki\checkmarkr\checkmarkc\checkmark$\checkmarki\checkmark$\checkmark^\checkmark\\checkmarkc\checkmarki\checkmarkr\checkmarkc\checkmark$\checkmarke\checkmark$\checkmark^\checkmark\\checkmarkc\checkmarki\checkmarkr\checkmarkc\checkmark$\checkmarkr\checkmark$\checkmark^\checkmark\\checkmarkc\checkmarki\checkmarkr\checkmarkc\checkmark$\checkmark.\checkmark$\checkmark^\checkmark\\checkmarkc\checkmarki\checkmarkr\checkmarkc\checkmark$\checkmark
\checkmark$\checkmark^\checkmark\\checkmarkc\checkmarki\checkmarkr\checkmarkc\checkmark$\checkmark \checkmark$\checkmark^\checkmark\\checkmarkc\checkmarki\checkmarkr\checkmarkc\checkmark$\checkmark \checkmark$\checkmark^\checkmark\\checkmarkc\checkmarki\checkmarkr\checkmarkc\checkmark$\checkmark \checkmark$\checkmark^\checkmark\\checkmarkc\checkmarki\checkmarkr\checkmarkc\checkmark$\checkmark \checkmark$\checkmark^\checkmark\\checkmarkc\checkmarki\checkmarkr\checkmarkc\checkmark$\checkmark\\checkmark$\checkmark^\checkmark\\checkmarkc\checkmarki\checkmarkr\checkmarkc\checkmark$\checkmarki\checkmark$\checkmark^\checkmark\\checkmarkc\checkmarki\checkmarkr\checkmarkc\checkmark$\checkmarkt\checkmark$\checkmark^\checkmark\\checkmarkc\checkmarki\checkmarkr\checkmarkc\checkmark$\checkmarke\checkmark$\checkmark^\checkmark\\checkmarkc\checkmarki\checkmarkr\checkmarkc\checkmark$\checkmarkm\checkmark$\checkmark^\checkmark\\checkmarkc\checkmarki\checkmarkr\checkmarkc\checkmark$\checkmark \checkmark$\checkmark^\checkmark\\checkmarkc\checkmarki\checkmarkr\checkmarkc\checkmark$\checkmark\\checkmark$\checkmark^\checkmark\\checkmarkc\checkmarki\checkmarkr\checkmarkc\checkmark$\checkmarkt\checkmark$\checkmark^\checkmark\\checkmarkc\checkmarki\checkmarkr\checkmarkc\checkmark$\checkmarke\checkmark$\checkmark^\checkmark\\checkmarkc\checkmarki\checkmarkr\checkmarkc\checkmark$\checkmarkx\checkmark$\checkmark^\checkmark\\checkmarkc\checkmarki\checkmarkr\checkmarkc\checkmark$\checkmarkt\checkmark$\checkmark^\checkmark\\checkmarkc\checkmarki\checkmarkr\checkmarkc\checkmark$\checkmarkb\checkmark$\checkmark^\checkmark\\checkmarkc\checkmarki\checkmarkr\checkmarkc\checkmark$\checkmarkf\checkmark$\checkmark^\checkmark\\checkmarkc\checkmarki\checkmarkr\checkmarkc\checkmark$\checkmark{\checkmark$\checkmark^\checkmark\\checkmarkc\checkmarki\checkmarkr\checkmarkc\checkmark$\checkmarkS\checkmark$\checkmark^\checkmark\\checkmarkc\checkmarki\checkmarkr\checkmarkc\checkmark$\checkmarku\checkmark$\checkmark^\checkmark\\checkmarkc\checkmarki\checkmarkr\checkmarkc\checkmark$\checkmarkr\checkmark$\checkmark^\checkmark\\checkmarkc\checkmarki\checkmarkr\checkmarkc\checkmark$\checkmark \checkmark$\checkmark^\checkmark\\checkmarkc\checkmarki\checkmarkr\checkmarkc\checkmark$\checkmarkq\checkmark$\checkmark^\checkmark\\checkmarkc\checkmarki\checkmarkr\checkmarkc\checkmark$\checkmarku\checkmark$\checkmark^\checkmark\\checkmarkc\checkmarki\checkmarkr\checkmarkc\checkmark$\checkmarko\checkmark$\checkmark^\checkmark\\checkmarkc\checkmarki\checkmarkr\checkmarkc\checkmark$\checkmarki\checkmark$\checkmark^\checkmark\\checkmarkc\checkmarki\checkmarkr\checkmarkc\checkmark$\checkmark \checkmark$\checkmark^\checkmark\\checkmarkc\checkmarki\checkmarkr\checkmarkc\checkmark$\checkmarka\checkmark$\checkmark^\checkmark\\checkmarkc\checkmarki\checkmarkr\checkmarkc\checkmark$\checkmarkg\checkmark$\checkmark^\checkmark\\checkmarkc\checkmarki\checkmarkr\checkmarkc\checkmark$\checkmarki\checkmark$\checkmark^\checkmark\\checkmarkc\checkmarki\checkmarkr\checkmarkc\checkmark$\checkmarkt\checkmark$\checkmark^\checkmark\\checkmarkc\checkmarki\checkmarkr\checkmarkc\checkmark$\checkmark-\checkmark$\checkmark^\checkmark\\checkmarkc\checkmarki\checkmarkr\checkmarkc\checkmark$\checkmarki\checkmark$\checkmark^\checkmark\\checkmarkc\checkmarki\checkmarkr\checkmarkc\checkmark$\checkmarkl\checkmark$\checkmark^\checkmark\\checkmarkc\checkmarki\checkmarkr\checkmarkc\checkmark$\checkmark \checkmark$\checkmark^\checkmark\\checkmarkc\checkmarki\checkmarkr\checkmarkc\checkmark$\checkmark?\checkmark$\checkmark^\checkmark\\checkmarkc\checkmarki\checkmarkr\checkmarkc\checkmark$\checkmark}\checkmark$\checkmark^\checkmark\\checkmarkc\checkmarki\checkmarkr\checkmarkc\checkmark$\checkmark \checkmark$\checkmark^\checkmark\\checkmarkc\checkmarki\checkmarkr\checkmarkc\checkmark$\checkmarkS\checkmark$\checkmark^\checkmark\\checkmarkc\checkmarki\checkmarkr\checkmarkc\checkmark$\checkmarku\checkmark$\checkmark^\checkmark\\checkmarkc\checkmarki\checkmarkr\checkmarkc\checkmark$\checkmarkr\checkmark$\checkmark^\checkmark\\checkmarkc\checkmarki\checkmarkr\checkmarkc\checkmark$\checkmark \checkmark$\checkmark^\checkmark\\checkmarkc\checkmarki\checkmarkr\checkmarkc\checkmark$\checkmarkl\checkmark$\checkmark^\checkmark\\checkmarkc\checkmarki\checkmarkr\checkmarkc\checkmark$\checkmarka\checkmark$\checkmark^\checkmark\\checkmarkc\checkmarki\checkmarkr\checkmarkc\checkmark$\checkmark \checkmark$\checkmark^\checkmark\\checkmarkc\checkmarki\checkmarkr\checkmarkc\checkmark$\checkmarkp\checkmark$\checkmark^\checkmark\\checkmarkc\checkmarki\checkmarkr\checkmarkc\checkmark$\checkmarki\checkmark$\checkmark^\checkmark\\checkmarkc\checkmarki\checkmarkr\checkmarkc\checkmark$\checkmarkè\checkmark$\checkmark^\checkmark\\checkmarkc\checkmarki\checkmarkr\checkmarkc\checkmark$\checkmarkc\checkmark$\checkmark^\checkmark\\checkmarkc\checkmarki\checkmarkr\checkmarkc\checkmark$\checkmarke\checkmark$\checkmark^\checkmark\\checkmarkc\checkmarki\checkmarkr\checkmarkc\checkmark$\checkmark \checkmark$\checkmark^\checkmark\\checkmarkc\checkmarki\checkmarkr\checkmarkc\checkmark$\checkmarkb\checkmark$\checkmark^\checkmark\\checkmarkc\checkmarki\checkmarkr\checkmarkc\checkmark$\checkmarkr\checkmark$\checkmark^\checkmark\\checkmarkc\checkmarki\checkmarkr\checkmarkc\checkmark$\checkmarku\checkmark$\checkmark^\checkmark\\checkmarkc\checkmarki\checkmarkr\checkmarkc\checkmark$\checkmarkt\checkmark$\checkmark^\checkmark\\checkmarkc\checkmarki\checkmarkr\checkmarkc\checkmark$\checkmarke\checkmark$\checkmark^\checkmark\\checkmarkc\checkmarki\checkmarkr\checkmarkc\checkmark$\checkmark \checkmark$\checkmark^\checkmark\\checkmarkc\checkmarki\checkmarkr\checkmarkc\checkmark$\checkmarkà\checkmark$\checkmark^\checkmark\\checkmarkc\checkmarki\checkmarkr\checkmarkc\checkmark$\checkmark \checkmark$\checkmark^\checkmark\\checkmarkc\checkmarki\checkmarkr\checkmarkc\checkmark$\checkmarku\checkmark$\checkmark^\checkmark\\checkmarkc\checkmarki\checkmarkr\checkmarkc\checkmark$\checkmarks\checkmark$\checkmark^\checkmark\\checkmarkc\checkmarki\checkmarkr\checkmarkc\checkmark$\checkmarki\checkmark$\checkmark^\checkmark\\checkmarkc\checkmarki\checkmarkr\checkmarkc\checkmark$\checkmarkn\checkmark$\checkmark^\checkmark\\checkmarkc\checkmarki\checkmarkr\checkmarkc\checkmark$\checkmarke\checkmark$\checkmark^\checkmark\\checkmarkc\checkmarki\checkmarkr\checkmarkc\checkmark$\checkmarkr\checkmark$\checkmark^\checkmark\\checkmarkc\checkmarki\checkmarkr\checkmarkc\checkmark$\checkmark \checkmark$\checkmark^\checkmark\\checkmarkc\checkmarki\checkmarkr\checkmarkc\checkmark$\checkmark(\checkmark$\checkmark^\checkmark\\checkmarkc\checkmarki\checkmarkr\checkmarkc\checkmark$\checkmarkr\checkmark$\checkmark^\checkmark\\checkmarkc\checkmarki\checkmarkr\checkmarkc\checkmark$\checkmarké\checkmark$\checkmark^\checkmark\\checkmarkc\checkmarki\checkmarkr\checkmarkc\checkmark$\checkmarks\checkmark$\checkmark^\checkmark\\checkmarkc\checkmarki\checkmarkr\checkmarkc\checkmark$\checkmarki\checkmark$\checkmark^\checkmark\\checkmarkc\checkmarki\checkmarkr\checkmarkc\checkmark$\checkmarkn\checkmark$\checkmark^\checkmark\\checkmarkc\checkmarki\checkmarkr\checkmarkc\checkmark$\checkmarke\checkmark$\checkmark^\checkmark\\checkmarkc\checkmarki\checkmarkr\checkmarkc\checkmark$\checkmark,\checkmark$\checkmark^\checkmark\\checkmarkc\checkmarki\checkmarkr\checkmarkc\checkmark$\checkmark \checkmark$\checkmark^\checkmark\\checkmarkc\checkmarki\checkmarkr\checkmarkc\checkmark$\checkmarkb\checkmark$\checkmark^\checkmark\\checkmarkc\checkmarki\checkmarkr\checkmarkc\checkmark$\checkmarko\checkmark$\checkmark^\checkmark\\checkmarkc\checkmarki\checkmarkr\checkmarkc\checkmark$\checkmarki\checkmark$\checkmark^\checkmark\\checkmarkc\checkmarki\checkmarkr\checkmarkc\checkmark$\checkmarks\checkmark$\checkmark^\checkmark\\checkmarkc\checkmarki\checkmarkr\checkmarkc\checkmark$\checkmark \checkmark$\checkmark^\checkmark\\checkmarkc\checkmarki\checkmarkr\checkmarkc\checkmark$\checkmarkd\checkmark$\checkmark^\checkmark\\checkmarkc\checkmarki\checkmarkr\checkmarkc\checkmark$\checkmarku\checkmark$\checkmark^\checkmark\\checkmarkc\checkmarki\checkmarkr\checkmarkc\checkmark$\checkmarkr\checkmark$\checkmark^\checkmark\\checkmarkc\checkmarki\checkmarkr\checkmarkc\checkmark$\checkmark,\checkmark$\checkmark^\checkmark\\checkmarkc\checkmarki\checkmarkr\checkmarkc\checkmark$\checkmark \checkmark$\checkmark^\checkmark\\checkmarkc\checkmarki\checkmarkr\checkmarkc\checkmark$\checkmarka\checkmark$\checkmark^\checkmark\\checkmarkc\checkmarki\checkmarkr\checkmarkc\checkmark$\checkmarkl\checkmark$\checkmark^\checkmark\\checkmarkc\checkmarki\checkmarkr\checkmarkc\checkmark$\checkmarku\checkmark$\checkmark^\checkmark\\checkmarkc\checkmarki\checkmarkr\checkmarkc\checkmark$\checkmarkm\checkmark$\checkmark^\checkmark\\checkmarkc\checkmarki\checkmarkr\checkmarkc\checkmark$\checkmarki\checkmark$\checkmark^\checkmark\\checkmarkc\checkmarki\checkmarkr\checkmarkc\checkmark$\checkmarkn\checkmark$\checkmark^\checkmark\\checkmarkc\checkmarki\checkmarkr\checkmarkc\checkmark$\checkmarki\checkmark$\checkmark^\checkmark\\checkmarkc\checkmarki\checkmarkr\checkmarkc\checkmark$\checkmarku\checkmark$\checkmark^\checkmark\\checkmarkc\checkmarki\checkmarkr\checkmarkc\checkmark$\checkmarkm\checkmark$\checkmark^\checkmark\\checkmarkc\checkmarki\checkmarkr\checkmarkc\checkmark$\checkmark)\checkmark$\checkmark^\checkmark\\checkmarkc\checkmarki\checkmarkr\checkmarkc\checkmark$\checkmark.\checkmark$\checkmark^\checkmark\\checkmarkc\checkmarki\checkmarkr\checkmarkc\checkmark$\checkmark
\checkmark$\checkmark^\checkmark\\checkmarkc\checkmarki\checkmarkr\checkmarkc\checkmark$\checkmark \checkmark$\checkmark^\checkmark\\checkmarkc\checkmarki\checkmarkr\checkmarkc\checkmark$\checkmark \checkmark$\checkmark^\checkmark\\checkmarkc\checkmarki\checkmarkr\checkmarkc\checkmark$\checkmark \checkmark$\checkmark^\checkmark\\checkmarkc\checkmarki\checkmarkr\checkmarkc\checkmark$\checkmark \checkmark$\checkmark^\checkmark\\checkmarkc\checkmarki\checkmarkr\checkmarkc\checkmark$\checkmark\\checkmark$\checkmark^\checkmark\\checkmarkc\checkmarki\checkmarkr\checkmarkc\checkmark$\checkmarki\checkmark$\checkmark^\checkmark\\checkmarkc\checkmarki\checkmarkr\checkmarkc\checkmark$\checkmarkt\checkmark$\checkmark^\checkmark\\checkmarkc\checkmarki\checkmarkr\checkmarkc\checkmark$\checkmarke\checkmark$\checkmark^\checkmark\\checkmarkc\checkmarki\checkmarkr\checkmarkc\checkmark$\checkmarkm\checkmark$\checkmark^\checkmark\\checkmarkc\checkmarki\checkmarkr\checkmarkc\checkmark$\checkmark \checkmark$\checkmark^\checkmark\\checkmarkc\checkmarki\checkmarkr\checkmarkc\checkmark$\checkmark\\checkmark$\checkmark^\checkmark\\checkmarkc\checkmarki\checkmarkr\checkmarkc\checkmark$\checkmarkt\checkmark$\checkmark^\checkmark\\checkmarkc\checkmarki\checkmarkr\checkmarkc\checkmark$\checkmarke\checkmark$\checkmark^\checkmark\\checkmarkc\checkmarki\checkmarkr\checkmarkc\checkmark$\checkmarkx\checkmark$\checkmark^\checkmark\\checkmarkc\checkmarki\checkmarkr\checkmarkc\checkmark$\checkmarkt\checkmark$\checkmark^\checkmark\\checkmarkc\checkmarki\checkmarkr\checkmarkc\checkmark$\checkmarkb\checkmark$\checkmark^\checkmark\\checkmarkc\checkmarki\checkmarkr\checkmarkc\checkmark$\checkmarkf\checkmark$\checkmark^\checkmark\\checkmarkc\checkmarki\checkmarkr\checkmarkc\checkmark$\checkmark{\checkmark$\checkmark^\checkmark\\checkmarkc\checkmarki\checkmarkr\checkmarkc\checkmark$\checkmarkD\checkmark$\checkmark^\checkmark\\checkmarkc\checkmarki\checkmarkr\checkmarkc\checkmark$\checkmarka\checkmark$\checkmark^\checkmark\\checkmarkc\checkmarki\checkmarkr\checkmarkc\checkmark$\checkmarkn\checkmark$\checkmark^\checkmark\\checkmarkc\checkmarki\checkmarkr\checkmarkc\checkmark$\checkmarks\checkmark$\checkmark^\checkmark\\checkmarkc\checkmarki\checkmarkr\checkmarkc\checkmark$\checkmark \checkmark$\checkmark^\checkmark\\checkmarkc\checkmarki\checkmarkr\checkmarkc\checkmark$\checkmarkq\checkmark$\checkmark^\checkmark\\checkmarkc\checkmarki\checkmarkr\checkmarkc\checkmark$\checkmarku\checkmark$\checkmark^\checkmark\\checkmarkc\checkmarki\checkmarkr\checkmarkc\checkmark$\checkmarke\checkmark$\checkmark^\checkmark\\checkmarkc\checkmarki\checkmarkr\checkmarkc\checkmark$\checkmarkl\checkmark$\checkmark^\checkmark\\checkmarkc\checkmarki\checkmarkr\checkmarkc\checkmark$\checkmark \checkmark$\checkmark^\checkmark\\checkmarkc\checkmarki\checkmarkr\checkmarkc\checkmark$\checkmarkb\checkmark$\checkmark^\checkmark\\checkmarkc\checkmarki\checkmarkr\checkmarkc\checkmark$\checkmarku\checkmark$\checkmark^\checkmark\\checkmarkc\checkmarki\checkmarkr\checkmarkc\checkmark$\checkmarkt\checkmark$\checkmark^\checkmark\\checkmarkc\checkmarki\checkmarkr\checkmarkc\checkmark$\checkmark \checkmark$\checkmark^\checkmark\\checkmarkc\checkmarki\checkmarkr\checkmarkc\checkmark$\checkmark?\checkmark$\checkmark^\checkmark\\checkmarkc\checkmarki\checkmarkr\checkmarkc\checkmark$\checkmark}\checkmark$\checkmark^\checkmark\\checkmarkc\checkmarki\checkmarkr\checkmarkc\checkmark$\checkmark \checkmark$\checkmark^\checkmark\\checkmarkc\checkmarki\checkmarkr\checkmarkc\checkmark$\checkmarkP\checkmark$\checkmark^\checkmark\\checkmarkc\checkmarki\checkmarkr\checkmarkc\checkmark$\checkmarke\checkmark$\checkmark^\checkmark\\checkmarkc\checkmarki\checkmarkr\checkmarkc\checkmark$\checkmarkr\checkmark$\checkmark^\checkmark\\checkmarkc\checkmarki\checkmarkr\checkmarkc\checkmark$\checkmarkm\checkmark$\checkmark^\checkmark\\checkmarkc\checkmarki\checkmarkr\checkmarkc\checkmark$\checkmarke\checkmark$\checkmark^\checkmark\\checkmarkc\checkmarki\checkmarkr\checkmarkc\checkmark$\checkmarkt\checkmark$\checkmark^\checkmark\\checkmarkc\checkmarki\checkmarkr\checkmarkc\checkmark$\checkmarkt\checkmark$\checkmark^\checkmark\\checkmarkc\checkmarki\checkmarkr\checkmarkc\checkmark$\checkmarkr\checkmark$\checkmark^\checkmark\\checkmarkc\checkmarki\checkmarkr\checkmarkc\checkmark$\checkmarke\checkmark$\checkmark^\checkmark\\checkmarkc\checkmarki\checkmarkr\checkmarkc\checkmark$\checkmark \checkmark$\checkmark^\checkmark\\checkmarkc\checkmarki\checkmarkr\checkmarkc\checkmark$\checkmarkl\checkmark$\checkmark^\checkmark\\checkmarkc\checkmarki\checkmarkr\checkmarkc\checkmark$\checkmark'\checkmark$\checkmark^\checkmark\\checkmarkc\checkmarki\checkmarkr\checkmarkc\checkmark$\checkmarku\checkmark$\checkmark^\checkmark\\checkmarkc\checkmarki\checkmarkr\checkmarkc\checkmark$\checkmarks\checkmark$\checkmark^\checkmark\\checkmarkc\checkmarki\checkmarkr\checkmarkc\checkmark$\checkmarki\checkmark$\checkmark^\checkmark\\checkmarkc\checkmarki\checkmarkr\checkmarkc\checkmark$\checkmarkn\checkmark$\checkmark^\checkmark\\checkmarkc\checkmarki\checkmarkr\checkmarkc\checkmark$\checkmarka\checkmark$\checkmark^\checkmark\\checkmarkc\checkmarki\checkmarkr\checkmarkc\checkmark$\checkmarkg\checkmark$\checkmark^\checkmark\\checkmarkc\checkmarki\checkmarkr\checkmarkc\checkmark$\checkmarke\checkmark$\checkmark^\checkmark\\checkmarkc\checkmarki\checkmarkr\checkmarkc\checkmark$\checkmark \checkmark$\checkmark^\checkmark\\checkmarkc\checkmarki\checkmarkr\checkmarkc\checkmark$\checkmarka\checkmark$\checkmark^\checkmark\\checkmarkc\checkmarki\checkmarkr\checkmarkc\checkmark$\checkmarku\checkmark$\checkmark^\checkmark\\checkmarkc\checkmarki\checkmarkr\checkmarkc\checkmark$\checkmarkt\checkmark$\checkmark^\checkmark\\checkmarkc\checkmarki\checkmarkr\checkmarkc\checkmark$\checkmarko\checkmark$\checkmark^\checkmark\\checkmarkc\checkmarki\checkmarkr\checkmarkc\checkmark$\checkmarkm\checkmark$\checkmark^\checkmark\\checkmarkc\checkmarki\checkmarkr\checkmarkc\checkmark$\checkmarka\checkmark$\checkmark^\checkmark\\checkmarkc\checkmarki\checkmarkr\checkmarkc\checkmark$\checkmarkt\checkmark$\checkmark^\checkmark\\checkmarkc\checkmarki\checkmarkr\checkmarkc\checkmark$\checkmarki\checkmark$\checkmark^\checkmark\\checkmarkc\checkmarki\checkmarkr\checkmarkc\checkmark$\checkmarkq\checkmark$\checkmark^\checkmark\\checkmarkc\checkmarki\checkmarkr\checkmarkc\checkmark$\checkmarku\checkmark$\checkmark^\checkmark\\checkmarkc\checkmarki\checkmarkr\checkmarkc\checkmark$\checkmarke\checkmark$\checkmark^\checkmark\\checkmarkc\checkmarki\checkmarkr\checkmarkc\checkmark$\checkmark \checkmark$\checkmark^\checkmark\\checkmarkc\checkmarki\checkmarkr\checkmarkc\checkmark$\checkmarkd\checkmark$\checkmark^\checkmark\\checkmarkc\checkmarki\checkmarkr\checkmarkc\checkmark$\checkmarke\checkmark$\checkmark^\checkmark\\checkmarkc\checkmarki\checkmarkr\checkmarkc\checkmark$\checkmark \checkmark$\checkmark^\checkmark\\checkmarkc\checkmarki\checkmarkr\checkmarkc\checkmark$\checkmarkf\checkmark$\checkmark^\checkmark\\checkmarkc\checkmarki\checkmarkr\checkmarkc\checkmark$\checkmarko\checkmark$\checkmark^\checkmark\\checkmarkc\checkmarki\checkmarkr\checkmarkc\checkmark$\checkmarkr\checkmark$\checkmark^\checkmark\\checkmarkc\checkmarki\checkmarkr\checkmarkc\checkmark$\checkmarkm\checkmark$\checkmark^\checkmark\\checkmarkc\checkmarki\checkmarkr\checkmarkc\checkmark$\checkmarke\checkmark$\checkmark^\checkmark\\checkmarkc\checkmarki\checkmarkr\checkmarkc\checkmark$\checkmarks\checkmark$\checkmark^\checkmark\\checkmarkc\checkmarki\checkmarkr\checkmarkc\checkmark$\checkmark \checkmark$\checkmark^\checkmark\\checkmarkc\checkmarki\checkmarkr\checkmarkc\checkmark$\checkmarkt\checkmark$\checkmark^\checkmark\\checkmarkc\checkmarki\checkmarkr\checkmarkc\checkmark$\checkmarkr\checkmark$\checkmark^\checkmark\\checkmarkc\checkmarki\checkmarkr\checkmarkc\checkmark$\checkmarki\checkmark$\checkmark^\checkmark\\checkmarkc\checkmarki\checkmarkr\checkmarkc\checkmark$\checkmarkd\checkmark$\checkmark^\checkmark\\checkmarkc\checkmarki\checkmarkr\checkmarkc\checkmark$\checkmarki\checkmark$\checkmark^\checkmark\\checkmarkc\checkmarki\checkmarkr\checkmarkc\checkmark$\checkmarkm\checkmark$\checkmark^\checkmark\\checkmarkc\checkmarki\checkmarkr\checkmarkc\checkmark$\checkmarke\checkmark$\checkmark^\checkmark\\checkmarkc\checkmarki\checkmarkr\checkmarkc\checkmark$\checkmarkn\checkmark$\checkmark^\checkmark\\checkmarkc\checkmarki\checkmarkr\checkmarkc\checkmark$\checkmarks\checkmark$\checkmark^\checkmark\\checkmarkc\checkmarki\checkmarkr\checkmarkc\checkmark$\checkmarki\checkmark$\checkmark^\checkmark\\checkmarkc\checkmarki\checkmarkr\checkmarkc\checkmark$\checkmarko\checkmark$\checkmark^\checkmark\\checkmarkc\checkmarki\checkmarkr\checkmarkc\checkmark$\checkmarkn\checkmark$\checkmark^\checkmark\\checkmarkc\checkmarki\checkmarkr\checkmarkc\checkmark$\checkmarkn\checkmark$\checkmark^\checkmark\\checkmarkc\checkmarki\checkmarkr\checkmarkc\checkmark$\checkmarke\checkmark$\checkmark^\checkmark\\checkmarkc\checkmarki\checkmarkr\checkmarkc\checkmark$\checkmarkl\checkmark$\checkmark^\checkmark\\checkmarkc\checkmarki\checkmarkr\checkmarkc\checkmark$\checkmarkl\checkmark$\checkmark^\checkmark\\checkmarkc\checkmarki\checkmarkr\checkmarkc\checkmark$\checkmarke\checkmark$\checkmark^\checkmark\\checkmarkc\checkmarki\checkmarkr\checkmarkc\checkmark$\checkmarks\checkmark$\checkmark^\checkmark\\checkmarkc\checkmarki\checkmarkr\checkmarkc\checkmark$\checkmark \checkmark$\checkmark^\checkmark\\checkmarkc\checkmarki\checkmarkr\checkmarkc\checkmark$\checkmarkc\checkmark$\checkmark^\checkmark\\checkmarkc\checkmarki\checkmarkr\checkmarkc\checkmark$\checkmarko\checkmark$\checkmark^\checkmark\\checkmarkc\checkmarki\checkmarkr\checkmarkc\checkmark$\checkmarkm\checkmark$\checkmark^\checkmark\\checkmarkc\checkmarki\checkmarkr\checkmarkc\checkmark$\checkmarkp\checkmark$\checkmark^\checkmark\\checkmarkc\checkmarki\checkmarkr\checkmarkc\checkmark$\checkmarkl\checkmark$\checkmark^\checkmark\\checkmarkc\checkmarki\checkmarkr\checkmarkc\checkmark$\checkmarke\checkmark$\checkmark^\checkmark\\checkmarkc\checkmarki\checkmarkr\checkmarkc\checkmark$\checkmarkx\checkmark$\checkmark^\checkmark\\checkmarkc\checkmarki\checkmarkr\checkmarkc\checkmark$\checkmarke\checkmark$\checkmark^\checkmark\\checkmarkc\checkmarki\checkmarkr\checkmarkc\checkmark$\checkmarks\checkmark$\checkmark^\checkmark\\checkmarkc\checkmarki\checkmarkr\checkmarkc\checkmark$\checkmark \checkmark$\checkmark^\checkmark\\checkmarkc\checkmarki\checkmarkr\checkmarkc\checkmark$\checkmarke\checkmark$\checkmark^\checkmark\\checkmarkc\checkmarki\checkmarkr\checkmarkc\checkmark$\checkmarkn\checkmark$\checkmark^\checkmark\\checkmarkc\checkmarki\checkmarkr\checkmarkc\checkmark$\checkmark \checkmark$\checkmark^\checkmark\\checkmarkc\checkmarki\checkmarkr\checkmarkc\checkmark$\checkmarks\checkmark$\checkmark^\checkmark\\checkmarkc\checkmarki\checkmarkr\checkmarkc\checkmark$\checkmarky\checkmark$\checkmark^\checkmark\\checkmarkc\checkmarki\checkmarkr\checkmarkc\checkmark$\checkmarkn\checkmark$\checkmark^\checkmark\\checkmarkc\checkmarki\checkmarkr\checkmarkc\checkmark$\checkmarkc\checkmark$\checkmark^\checkmark\\checkmarkc\checkmarki\checkmarkr\checkmarkc\checkmark$\checkmarkh\checkmark$\checkmark^\checkmark\\checkmarkc\checkmarki\checkmarkr\checkmarkc\checkmark$\checkmarkr\checkmark$\checkmark^\checkmark\\checkmarkc\checkmarki\checkmarkr\checkmarkc\checkmark$\checkmarko\checkmark$\checkmark^\checkmark\\checkmarkc\checkmarki\checkmarkr\checkmarkc\checkmark$\checkmarkn\checkmark$\checkmark^\checkmark\\checkmarkc\checkmarki\checkmarkr\checkmarkc\checkmark$\checkmarki\checkmark$\checkmark^\checkmark\\checkmarkc\checkmarki\checkmarkr\checkmarkc\checkmark$\checkmarks\checkmark$\checkmark^\checkmark\\checkmarkc\checkmarki\checkmarkr\checkmarkc\checkmark$\checkmarka\checkmark$\checkmark^\checkmark\\checkmarkc\checkmarki\checkmarkr\checkmarkc\checkmark$\checkmarkn\checkmark$\checkmark^\checkmark\\checkmarkc\checkmarki\checkmarkr\checkmarkc\checkmark$\checkmarkt\checkmark$\checkmark^\checkmark\\checkmarkc\checkmarki\checkmarkr\checkmarkc\checkmark$\checkmark \checkmark$\checkmark^\checkmark\\checkmarkc\checkmarki\checkmarkr\checkmarkc\checkmark$\checkmark5\checkmark$\checkmark^\checkmark\\checkmarkc\checkmarki\checkmarkr\checkmarkc\checkmark$\checkmark \checkmark$\checkmark^\checkmark\\checkmarkc\checkmarki\checkmarkr\checkmarkc\checkmark$\checkmarka\checkmark$\checkmark^\checkmark\\checkmarkc\checkmarki\checkmarkr\checkmarkc\checkmark$\checkmarkx\checkmark$\checkmark^\checkmark\\checkmarkc\checkmarki\checkmarkr\checkmarkc\checkmark$\checkmarke\checkmark$\checkmark^\checkmark\\checkmarkc\checkmarki\checkmarkr\checkmarkc\checkmark$\checkmarks\checkmark$\checkmark^\checkmark\\checkmarkc\checkmarki\checkmarkr\checkmarkc\checkmark$\checkmark \checkmark$\checkmark^\checkmark\\checkmarkc\checkmarki\checkmarkr\checkmarkc\checkmark$\checkmarks\checkmark$\checkmark^\checkmark\\checkmarkc\checkmarki\checkmarkr\checkmarkc\checkmark$\checkmarki\checkmark$\checkmark^\checkmark\\checkmarkc\checkmarki\checkmarkr\checkmarkc\checkmark$\checkmarkm\checkmark$\checkmark^\checkmark\\checkmarkc\checkmarki\checkmarkr\checkmarkc\checkmark$\checkmarku\checkmark$\checkmark^\checkmark\\checkmarkc\checkmarki\checkmarkr\checkmarkc\checkmark$\checkmarkl\checkmark$\checkmark^\checkmark\\checkmarkc\checkmarki\checkmarkr\checkmarkc\checkmark$\checkmarkt\checkmark$\checkmark^\checkmark\\checkmarkc\checkmarki\checkmarkr\checkmarkc\checkmark$\checkmarka\checkmark$\checkmark^\checkmark\\checkmarkc\checkmarki\checkmarkr\checkmarkc\checkmark$\checkmarkn\checkmark$\checkmark^\checkmark\\checkmarkc\checkmarki\checkmarkr\checkmarkc\checkmark$\checkmarké\checkmark$\checkmark^\checkmark\\checkmarkc\checkmarki\checkmarkr\checkmarkc\checkmark$\checkmarkm\checkmark$\checkmark^\checkmark\\checkmarkc\checkmarki\checkmarkr\checkmarkc\checkmark$\checkmarke\checkmark$\checkmark^\checkmark\\checkmarkc\checkmarki\checkmarkr\checkmarkc\checkmark$\checkmarkn\checkmark$\checkmark^\checkmark\\checkmarkc\checkmarki\checkmarkr\checkmarkc\checkmark$\checkmarkt\checkmark$\checkmark^\checkmark\\checkmarkc\checkmarki\checkmarkr\checkmarkc\checkmark$\checkmark \checkmark$\checkmark^\checkmark\\checkmarkc\checkmarki\checkmarkr\checkmarkc\checkmark$\checkmarkd\checkmark$\checkmark^\checkmark\\checkmarkc\checkmarki\checkmarkr\checkmarkc\checkmark$\checkmarke\checkmark$\checkmark^\checkmark\\checkmarkc\checkmarki\checkmarkr\checkmarkc\checkmark$\checkmarkp\checkmark$\checkmark^\checkmark\\checkmarkc\checkmarki\checkmarkr\checkmarkc\checkmark$\checkmarku\checkmark$\checkmark^\checkmark\\checkmarkc\checkmarki\checkmarkr\checkmarkc\checkmark$\checkmarki\checkmark$\checkmark^\checkmark\\checkmarkc\checkmarki\checkmarkr\checkmarkc\checkmark$\checkmarks\checkmark$\checkmark^\checkmark\\checkmarkc\checkmarki\checkmarkr\checkmarkc\checkmark$\checkmark \checkmark$\checkmark^\checkmark\\checkmarkc\checkmarki\checkmarkr\checkmarkc\checkmark$\checkmarku\checkmark$\checkmark^\checkmark\\checkmarkc\checkmarki\checkmarkr\checkmarkc\checkmark$\checkmarkn\checkmark$\checkmark^\checkmark\\checkmarkc\checkmarki\checkmarkr\checkmarkc\checkmark$\checkmarke\checkmark$\checkmark^\checkmark\\checkmarkc\checkmarki\checkmarkr\checkmarkc\checkmark$\checkmark \checkmark$\checkmark^\checkmark\\checkmarkc\checkmarki\checkmarkr\checkmarkc\checkmark$\checkmarki\checkmark$\checkmark^\checkmark\\checkmarkc\checkmarki\checkmarkr\checkmarkc\checkmark$\checkmarkn\checkmark$\checkmark^\checkmark\\checkmarkc\checkmarki\checkmarkr\checkmarkc\checkmark$\checkmarkt\checkmark$\checkmark^\checkmark\\checkmarkc\checkmarki\checkmarkr\checkmarkc\checkmark$\checkmarke\checkmark$\checkmark^\checkmark\\checkmarkc\checkmarki\checkmarkr\checkmarkc\checkmark$\checkmarkr\checkmark$\checkmark^\checkmark\\checkmarkc\checkmarki\checkmarkr\checkmarkc\checkmark$\checkmarkf\checkmark$\checkmark^\checkmark\\checkmarkc\checkmarki\checkmarkr\checkmarkc\checkmark$\checkmarka\checkmark$\checkmark^\checkmark\\checkmarkc\checkmarki\checkmarkr\checkmarkc\checkmark$\checkmarkc\checkmark$\checkmark^\checkmark\\checkmarkc\checkmarki\checkmarkr\checkmarkc\checkmark$\checkmarke\checkmark$\checkmark^\checkmark\\checkmarkc\checkmarki\checkmarkr\checkmarkc\checkmark$\checkmark \checkmark$\checkmark^\checkmark\\checkmarkc\checkmarki\checkmarkr\checkmarkc\checkmark$\checkmarkn\checkmark$\checkmark^\checkmark\\checkmarkc\checkmarki\checkmarkr\checkmarkc\checkmark$\checkmarku\checkmark$\checkmark^\checkmark\\checkmarkc\checkmarki\checkmarkr\checkmarkc\checkmark$\checkmarkm\checkmark$\checkmark^\checkmark\\checkmarkc\checkmarki\checkmarkr\checkmarkc\checkmark$\checkmarké\checkmark$\checkmark^\checkmark\\checkmarkc\checkmarki\checkmarkr\checkmarkc\checkmark$\checkmarkr\checkmark$\checkmark^\checkmark\\checkmarkc\checkmarki\checkmarkr\checkmarkc\checkmark$\checkmarki\checkmark$\checkmark^\checkmark\\checkmarkc\checkmarki\checkmarkr\checkmarkc\checkmark$\checkmarkq\checkmark$\checkmark^\checkmark\\checkmarkc\checkmarki\checkmarkr\checkmarkc\checkmark$\checkmarku\checkmark$\checkmark^\checkmark\\checkmarkc\checkmarki\checkmarkr\checkmarkc\checkmark$\checkmarke\checkmark$\checkmark^\checkmark\\checkmarkc\checkmarki\checkmarkr\checkmarkc\checkmark$\checkmark.\checkmark$\checkmark^\checkmark\\checkmarkc\checkmarki\checkmarkr\checkmarkc\checkmark$\checkmark
\checkmark$\checkmark^\checkmark\\checkmarkc\checkmarki\checkmarkr\checkmarkc\checkmark$\checkmark\\checkmark$\checkmark^\checkmark\\checkmarkc\checkmarki\checkmarkr\checkmarkc\checkmark$\checkmarke\checkmark$\checkmark^\checkmark\\checkmarkc\checkmarki\checkmarkr\checkmarkc\checkmark$\checkmarkn\checkmark$\checkmark^\checkmark\\checkmarkc\checkmarki\checkmarkr\checkmarkc\checkmark$\checkmarkd\checkmark$\checkmark^\checkmark\\checkmarkc\checkmarki\checkmarkr\checkmarkc\checkmark$\checkmark{\checkmark$\checkmark^\checkmark\\checkmarkc\checkmarki\checkmarkr\checkmarkc\checkmark$\checkmarki\checkmark$\checkmark^\checkmark\\checkmarkc\checkmarki\checkmarkr\checkmarkc\checkmark$\checkmarkt\checkmark$\checkmark^\checkmark\\checkmarkc\checkmarki\checkmarkr\checkmarkc\checkmark$\checkmarke\checkmark$\checkmark^\checkmark\\checkmarkc\checkmarki\checkmarkr\checkmarkc\checkmark$\checkmarkm\checkmark$\checkmark^\checkmark\\checkmarkc\checkmarki\checkmarkr\checkmarkc\checkmark$\checkmarki\checkmark$\checkmark^\checkmark\\checkmarkc\checkmarki\checkmarkr\checkmarkc\checkmark$\checkmarkz\checkmark$\checkmark^\checkmark\\checkmarkc\checkmarki\checkmarkr\checkmarkc\checkmark$\checkmarke\checkmark$\checkmark^\checkmark\\checkmarkc\checkmarki\checkmarkr\checkmarkc\checkmark$\checkmark}\checkmark$\checkmark^\checkmark\\checkmarkc\checkmarki\checkmarkr\checkmarkc\checkmark$\checkmark
\checkmark$\checkmark^\checkmark\\checkmarkc\checkmarki\checkmarkr\checkmarkc\checkmark$\checkmark
\checkmark$\checkmark^\checkmark\\checkmarkc\checkmarki\checkmarkr\checkmarkc\checkmark$\checkmark\\checkmark$\checkmark^\checkmark\\checkmarkc\checkmarki\checkmarkr\checkmarkc\checkmark$\checkmarkb\checkmark$\checkmark^\checkmark\\checkmarkc\checkmarki\checkmarkr\checkmarkc\checkmark$\checkmarke\checkmark$\checkmark^\checkmark\\checkmarkc\checkmarki\checkmarkr\checkmarkc\checkmark$\checkmarkg\checkmark$\checkmark^\checkmark\\checkmarkc\checkmarki\checkmarkr\checkmarkc\checkmark$\checkmarki\checkmark$\checkmark^\checkmark\\checkmarkc\checkmarki\checkmarkr\checkmarkc\checkmark$\checkmarkn\checkmark$\checkmark^\checkmark\\checkmarkc\checkmarki\checkmarkr\checkmarkc\checkmark$\checkmark{\checkmark$\checkmark^\checkmark\\checkmarkc\checkmarki\checkmarkr\checkmarkc\checkmark$\checkmarkf\checkmark$\checkmark^\checkmark\\checkmarkc\checkmarki\checkmarkr\checkmarkc\checkmark$\checkmarki\checkmark$\checkmark^\checkmark\\checkmarkc\checkmarki\checkmarkr\checkmarkc\checkmark$\checkmarkg\checkmark$\checkmark^\checkmark\\checkmarkc\checkmarki\checkmarkr\checkmarkc\checkmark$\checkmarku\checkmark$\checkmark^\checkmark\\checkmarkc\checkmarki\checkmarkr\checkmarkc\checkmark$\checkmarkr\checkmark$\checkmark^\checkmark\\checkmarkc\checkmarki\checkmarkr\checkmarkc\checkmark$\checkmarke\checkmark$\checkmark^\checkmark\\checkmarkc\checkmarki\checkmarkr\checkmarkc\checkmark$\checkmark}\checkmark$\checkmark^\checkmark\\checkmarkc\checkmarki\checkmarkr\checkmarkc\checkmark$\checkmark[\checkmark$\checkmark^\checkmark\\checkmarkc\checkmarki\checkmarkr\checkmarkc\checkmark$\checkmarkh\checkmark$\checkmark^\checkmark\\checkmarkc\checkmarki\checkmarkr\checkmarkc\checkmark$\checkmarkt\checkmark$\checkmark^\checkmark\\checkmarkc\checkmarki\checkmarkr\checkmarkc\checkmark$\checkmarkb\checkmark$\checkmark^\checkmark\\checkmarkc\checkmarki\checkmarkr\checkmarkc\checkmark$\checkmarkp\checkmark$\checkmark^\checkmark\\checkmarkc\checkmarki\checkmarkr\checkmarkc\checkmark$\checkmark]\checkmark$\checkmark^\checkmark\\checkmarkc\checkmarki\checkmarkr\checkmarkc\checkmark$\checkmark
\checkmark$\checkmark^\checkmark\\checkmarkc\checkmarki\checkmarkr\checkmarkc\checkmark$\checkmark \checkmark$\checkmark^\checkmark\\checkmarkc\checkmarki\checkmarkr\checkmarkc\checkmark$\checkmark \checkmark$\checkmark^\checkmark\\checkmarkc\checkmarki\checkmarkr\checkmarkc\checkmark$\checkmark \checkmark$\checkmark^\checkmark\\checkmarkc\checkmarki\checkmarkr\checkmarkc\checkmark$\checkmark \checkmark$\checkmark^\checkmark\\checkmarkc\checkmarki\checkmarkr\checkmarkc\checkmark$\checkmark\\checkmark$\checkmark^\checkmark\\checkmarkc\checkmarki\checkmarkr\checkmarkc\checkmark$\checkmarkc\checkmark$\checkmark^\checkmark\\checkmarkc\checkmarki\checkmarkr\checkmarkc\checkmark$\checkmarke\checkmark$\checkmark^\checkmark\\checkmarkc\checkmarki\checkmarkr\checkmarkc\checkmark$\checkmarkn\checkmark$\checkmark^\checkmark\\checkmarkc\checkmarki\checkmarkr\checkmarkc\checkmark$\checkmarkt\checkmark$\checkmark^\checkmark\\checkmarkc\checkmarki\checkmarkr\checkmarkc\checkmark$\checkmarke\checkmark$\checkmark^\checkmark\\checkmarkc\checkmarki\checkmarkr\checkmarkc\checkmark$\checkmarkr\checkmark$\checkmark^\checkmark\\checkmarkc\checkmarki\checkmarkr\checkmarkc\checkmark$\checkmarki\checkmark$\checkmark^\checkmark\\checkmarkc\checkmarki\checkmarkr\checkmarkc\checkmark$\checkmarkn\checkmark$\checkmark^\checkmark\\checkmarkc\checkmarki\checkmarkr\checkmarkc\checkmark$\checkmarkg\checkmark$\checkmark^\checkmark\\checkmarkc\checkmarki\checkmarkr\checkmarkc\checkmark$\checkmark
\checkmark$\checkmark^\checkmark\\checkmarkc\checkmarki\checkmarkr\checkmarkc\checkmark$\checkmark \checkmark$\checkmark^\checkmark\\checkmarkc\checkmarki\checkmarkr\checkmarkc\checkmark$\checkmark \checkmark$\checkmark^\checkmark\\checkmarkc\checkmarki\checkmarkr\checkmarkc\checkmark$\checkmark \checkmark$\checkmark^\checkmark\\checkmarkc\checkmarki\checkmarkr\checkmarkc\checkmark$\checkmark \checkmark$\checkmark^\checkmark\\checkmarkc\checkmarki\checkmarkr\checkmarkc\checkmark$\checkmark\\checkmark$\checkmark^\checkmark\\checkmarkc\checkmarki\checkmarkr\checkmarkc\checkmark$\checkmarkb\checkmark$\checkmark^\checkmark\\checkmarkc\checkmarki\checkmarkr\checkmarkc\checkmark$\checkmarke\checkmark$\checkmark^\checkmark\\checkmarkc\checkmarki\checkmarkr\checkmarkc\checkmark$\checkmarkg\checkmark$\checkmark^\checkmark\\checkmarkc\checkmarki\checkmarkr\checkmarkc\checkmark$\checkmarki\checkmark$\checkmark^\checkmark\\checkmarkc\checkmarki\checkmarkr\checkmarkc\checkmark$\checkmarkn\checkmark$\checkmark^\checkmark\\checkmarkc\checkmarki\checkmarkr\checkmarkc\checkmark$\checkmark{\checkmark$\checkmark^\checkmark\\checkmarkc\checkmarki\checkmarkr\checkmarkc\checkmark$\checkmarkt\checkmark$\checkmark^\checkmark\\checkmarkc\checkmarki\checkmarkr\checkmarkc\checkmark$\checkmarki\checkmark$\checkmark^\checkmark\\checkmarkc\checkmarki\checkmarkr\checkmarkc\checkmark$\checkmarkk\checkmark$\checkmark^\checkmark\\checkmarkc\checkmarki\checkmarkr\checkmarkc\checkmark$\checkmarkz\checkmark$\checkmark^\checkmark\\checkmarkc\checkmarki\checkmarkr\checkmarkc\checkmark$\checkmarkp\checkmark$\checkmark^\checkmark\\checkmarkc\checkmarki\checkmarkr\checkmarkc\checkmark$\checkmarki\checkmark$\checkmark^\checkmark\\checkmarkc\checkmarki\checkmarkr\checkmarkc\checkmark$\checkmarkc\checkmark$\checkmark^\checkmark\\checkmarkc\checkmarki\checkmarkr\checkmarkc\checkmark$\checkmarkt\checkmark$\checkmark^\checkmark\\checkmarkc\checkmarki\checkmarkr\checkmarkc\checkmark$\checkmarku\checkmark$\checkmark^\checkmark\\checkmarkc\checkmarki\checkmarkr\checkmarkc\checkmark$\checkmarkr\checkmark$\checkmark^\checkmark\\checkmarkc\checkmarki\checkmarkr\checkmarkc\checkmark$\checkmarke\checkmark$\checkmark^\checkmark\\checkmarkc\checkmarki\checkmarkr\checkmarkc\checkmark$\checkmark}\checkmark$\checkmark^\checkmark\\checkmarkc\checkmarki\checkmarkr\checkmarkc\checkmark$\checkmark[\checkmark$\checkmark^\checkmark\\checkmarkc\checkmarki\checkmarkr\checkmarkc\checkmark$\checkmark
\checkmark$\checkmark^\checkmark\\checkmarkc\checkmarki\checkmarkr\checkmarkc\checkmark$\checkmark \checkmark$\checkmark^\checkmark\\checkmarkc\checkmarki\checkmarkr\checkmarkc\checkmark$\checkmark \checkmark$\checkmark^\checkmark\\checkmarkc\checkmarki\checkmarkr\checkmarkc\checkmark$\checkmark \checkmark$\checkmark^\checkmark\\checkmarkc\checkmarki\checkmarkr\checkmarkc\checkmark$\checkmark \checkmark$\checkmark^\checkmark\\checkmarkc\checkmarki\checkmarkr\checkmarkc\checkmark$\checkmark \checkmark$\checkmark^\checkmark\\checkmarkc\checkmarki\checkmarkr\checkmarkc\checkmark$\checkmark \checkmark$\checkmark^\checkmark\\checkmarkc\checkmarki\checkmarkr\checkmarkc\checkmark$\checkmark \checkmark$\checkmark^\checkmark\\checkmarkc\checkmarki\checkmarkr\checkmarkc\checkmark$\checkmark \checkmark$\checkmark^\checkmark\\checkmarkc\checkmarki\checkmarkr\checkmarkc\checkmark$\checkmarkb\checkmark$\checkmark^\checkmark\\checkmarkc\checkmarki\checkmarkr\checkmarkc\checkmark$\checkmarku\checkmark$\checkmark^\checkmark\\checkmarkc\checkmarki\checkmarkr\checkmarkc\checkmark$\checkmarkb\checkmark$\checkmark^\checkmark\\checkmarkc\checkmarki\checkmarkr\checkmarkc\checkmark$\checkmarkb\checkmark$\checkmark^\checkmark\\checkmarkc\checkmarki\checkmarkr\checkmarkc\checkmark$\checkmarkl\checkmark$\checkmark^\checkmark\\checkmarkc\checkmarki\checkmarkr\checkmarkc\checkmark$\checkmarke\checkmark$\checkmark^\checkmark\\checkmarkc\checkmarki\checkmarkr\checkmarkc\checkmark$\checkmark/\checkmark$\checkmark^\checkmark\\checkmarkc\checkmarki\checkmarkr\checkmarkc\checkmark$\checkmark.\checkmark$\checkmark^\checkmark\\checkmarkc\checkmarki\checkmarkr\checkmarkc\checkmark$\checkmarks\checkmark$\checkmark^\checkmark\\checkmarkc\checkmarki\checkmarkr\checkmarkc\checkmark$\checkmarkt\checkmark$\checkmark^\checkmark\\checkmarkc\checkmarki\checkmarkr\checkmarkc\checkmark$\checkmarky\checkmark$\checkmark^\checkmark\\checkmarkc\checkmarki\checkmarkr\checkmarkc\checkmark$\checkmarkl\checkmark$\checkmark^\checkmark\\checkmarkc\checkmarki\checkmarkr\checkmarkc\checkmark$\checkmarke\checkmark$\checkmark^\checkmark\\checkmarkc\checkmarki\checkmarkr\checkmarkc\checkmark$\checkmark=\checkmark$\checkmark^\checkmark\\checkmarkc\checkmarki\checkmarkr\checkmarkc\checkmark$\checkmark{\checkmark$\checkmark^\checkmark\\checkmarkc\checkmarki\checkmarkr\checkmarkc\checkmark$\checkmarke\checkmark$\checkmark^\checkmark\\checkmarkc\checkmarki\checkmarkr\checkmarkc\checkmark$\checkmarkl\checkmark$\checkmark^\checkmark\\checkmarkc\checkmarki\checkmarkr\checkmarkc\checkmark$\checkmarkl\checkmark$\checkmark^\checkmark\\checkmarkc\checkmarki\checkmarkr\checkmarkc\checkmark$\checkmarki\checkmark$\checkmark^\checkmark\\checkmarkc\checkmarki\checkmarkr\checkmarkc\checkmark$\checkmarkp\checkmark$\checkmark^\checkmark\\checkmarkc\checkmarki\checkmarkr\checkmarkc\checkmark$\checkmarks\checkmark$\checkmark^\checkmark\\checkmarkc\checkmarki\checkmarkr\checkmarkc\checkmark$\checkmarke\checkmark$\checkmark^\checkmark\\checkmarkc\checkmarki\checkmarkr\checkmarkc\checkmark$\checkmark,\checkmark$\checkmark^\checkmark\\checkmarkc\checkmarki\checkmarkr\checkmarkc\checkmark$\checkmark \checkmark$\checkmark^\checkmark\\checkmarkc\checkmarki\checkmarkr\checkmarkc\checkmark$\checkmarkd\checkmark$\checkmark^\checkmark\\checkmarkc\checkmarki\checkmarkr\checkmarkc\checkmark$\checkmarkr\checkmark$\checkmark^\checkmark\\checkmarkc\checkmarki\checkmarkr\checkmarkc\checkmark$\checkmarka\checkmark$\checkmark^\checkmark\\checkmarkc\checkmarki\checkmarkr\checkmarkc\checkmark$\checkmarkw\checkmark$\checkmark^\checkmark\\checkmarkc\checkmarki\checkmarkr\checkmarkc\checkmark$\checkmark=\checkmark$\checkmark^\checkmark\\checkmarkc\checkmarki\checkmarkr\checkmarkc\checkmark$\checkmarkm\checkmark$\checkmark^\checkmark\\checkmarkc\checkmarki\checkmarkr\checkmarkc\checkmark$\checkmarka\checkmark$\checkmark^\checkmark\\checkmarkc\checkmarki\checkmarkr\checkmarkc\checkmark$\checkmarkc\checkmark$\checkmark^\checkmark\\checkmarkc\checkmarki\checkmarkr\checkmarkc\checkmark$\checkmarko\checkmark$\checkmark^\checkmark\\checkmarkc\checkmarki\checkmarkr\checkmarkc\checkmark$\checkmarku\checkmark$\checkmark^\checkmark\\checkmarkc\checkmarki\checkmarkr\checkmarkc\checkmark$\checkmarkl\checkmark$\checkmark^\checkmark\\checkmarkc\checkmarki\checkmarkr\checkmarkc\checkmark$\checkmarke\checkmark$\checkmark^\checkmark\\checkmarkc\checkmarki\checkmarkr\checkmarkc\checkmark$\checkmarku\checkmark$\checkmark^\checkmark\\checkmarkc\checkmarki\checkmarkr\checkmarkc\checkmark$\checkmarkr\checkmark$\checkmark^\checkmark\\checkmarkc\checkmarki\checkmarkr\checkmarkc\checkmark$\checkmark,\checkmark$\checkmark^\checkmark\\checkmarkc\checkmarki\checkmarkr\checkmarkc\checkmark$\checkmark \checkmark$\checkmark^\checkmark\\checkmarkc\checkmarki\checkmarkr\checkmarkc\checkmark$\checkmarkt\checkmark$\checkmark^\checkmark\\checkmarkc\checkmarki\checkmarkr\checkmarkc\checkmark$\checkmarkh\checkmark$\checkmark^\checkmark\\checkmarkc\checkmarki\checkmarkr\checkmarkc\checkmark$\checkmarki\checkmark$\checkmark^\checkmark\\checkmarkc\checkmarki\checkmarkr\checkmarkc\checkmark$\checkmarkc\checkmark$\checkmark^\checkmark\\checkmarkc\checkmarki\checkmarkr\checkmarkc\checkmark$\checkmarkk\checkmark$\checkmark^\checkmark\\checkmarkc\checkmarki\checkmarkr\checkmarkc\checkmark$\checkmark,\checkmark$\checkmark^\checkmark\\checkmarkc\checkmarki\checkmarkr\checkmarkc\checkmark$\checkmark \checkmark$\checkmark^\checkmark\\checkmarkc\checkmarki\checkmarkr\checkmarkc\checkmark$\checkmarkf\checkmark$\checkmark^\checkmark\\checkmarkc\checkmarki\checkmarkr\checkmarkc\checkmark$\checkmarki\checkmark$\checkmark^\checkmark\\checkmarkc\checkmarki\checkmarkr\checkmarkc\checkmark$\checkmarkl\checkmark$\checkmark^\checkmark\\checkmarkc\checkmarki\checkmarkr\checkmarkc\checkmark$\checkmarkl\checkmark$\checkmark^\checkmark\\checkmarkc\checkmarki\checkmarkr\checkmarkc\checkmark$\checkmark=\checkmark$\checkmark^\checkmark\\checkmarkc\checkmarki\checkmarkr\checkmarkc\checkmark$\checkmarkb\checkmark$\checkmark^\checkmark\\checkmarkc\checkmarki\checkmarkr\checkmarkc\checkmark$\checkmarkl\checkmark$\checkmark^\checkmark\\checkmarkc\checkmarki\checkmarkr\checkmarkc\checkmark$\checkmarku\checkmark$\checkmark^\checkmark\\checkmarkc\checkmarki\checkmarkr\checkmarkc\checkmark$\checkmarke\checkmark$\checkmark^\checkmark\\checkmarkc\checkmarki\checkmarkr\checkmarkc\checkmark$\checkmark!\checkmark$\checkmark^\checkmark\\checkmarkc\checkmarki\checkmarkr\checkmarkc\checkmark$\checkmark8\checkmark$\checkmark^\checkmark\\checkmarkc\checkmarki\checkmarkr\checkmarkc\checkmark$\checkmark,\checkmark$\checkmark^\checkmark\\checkmarkc\checkmarki\checkmarkr\checkmarkc\checkmark$\checkmark \checkmark$\checkmark^\checkmark\\checkmarkc\checkmarki\checkmarkr\checkmarkc\checkmark$\checkmarkt\checkmark$\checkmark^\checkmark\\checkmarkc\checkmarki\checkmarkr\checkmarkc\checkmark$\checkmarke\checkmark$\checkmark^\checkmark\\checkmarkc\checkmarki\checkmarkr\checkmarkc\checkmark$\checkmarkx\checkmark$\checkmark^\checkmark\\checkmarkc\checkmarki\checkmarkr\checkmarkc\checkmark$\checkmarkt\checkmark$\checkmark^\checkmark\\checkmarkc\checkmarki\checkmarkr\checkmarkc\checkmark$\checkmark \checkmark$\checkmark^\checkmark\\checkmarkc\checkmarki\checkmarkr\checkmarkc\checkmark$\checkmarkw\checkmark$\checkmark^\checkmark\\checkmarkc\checkmarki\checkmarkr\checkmarkc\checkmark$\checkmarki\checkmark$\checkmark^\checkmark\\checkmarkc\checkmarki\checkmarkr\checkmarkc\checkmark$\checkmarkd\checkmark$\checkmark^\checkmark\\checkmarkc\checkmarki\checkmarkr\checkmarkc\checkmark$\checkmarkt\checkmark$\checkmark^\checkmark\\checkmarkc\checkmarki\checkmarkr\checkmarkc\checkmark$\checkmarkh\checkmark$\checkmark^\checkmark\\checkmarkc\checkmarki\checkmarkr\checkmarkc\checkmark$\checkmark=\checkmark$\checkmark^\checkmark\\checkmarkc\checkmarki\checkmarkr\checkmarkc\checkmark$\checkmark3\checkmark$\checkmark^\checkmark\\checkmarkc\checkmarki\checkmarkr\checkmarkc\checkmark$\checkmarkc\checkmark$\checkmark^\checkmark\\checkmarkc\checkmarki\checkmarkr\checkmarkc\checkmark$\checkmarkm\checkmark$\checkmark^\checkmark\\checkmarkc\checkmarki\checkmarkr\checkmarkc\checkmark$\checkmark,\checkmark$\checkmark^\checkmark\\checkmarkc\checkmarki\checkmarkr\checkmarkc\checkmark$\checkmark \checkmark$\checkmark^\checkmark\\checkmarkc\checkmarki\checkmarkr\checkmarkc\checkmark$\checkmarka\checkmark$\checkmark^\checkmark\\checkmarkc\checkmarki\checkmarkr\checkmarkc\checkmark$\checkmarkl\checkmark$\checkmark^\checkmark\\checkmarkc\checkmarki\checkmarkr\checkmarkc\checkmark$\checkmarki\checkmark$\checkmark^\checkmark\\checkmarkc\checkmarki\checkmarkr\checkmarkc\checkmark$\checkmarkg\checkmark$\checkmark^\checkmark\\checkmarkc\checkmarki\checkmarkr\checkmarkc\checkmark$\checkmarkn\checkmark$\checkmark^\checkmark\\checkmarkc\checkmarki\checkmarkr\checkmarkc\checkmark$\checkmark=\checkmark$\checkmark^\checkmark\\checkmarkc\checkmarki\checkmarkr\checkmarkc\checkmark$\checkmarkc\checkmark$\checkmark^\checkmark\\checkmarkc\checkmarki\checkmarkr\checkmarkc\checkmark$\checkmarke\checkmark$\checkmark^\checkmark\\checkmarkc\checkmarki\checkmarkr\checkmarkc\checkmark$\checkmarkn\checkmark$\checkmark^\checkmark\\checkmarkc\checkmarki\checkmarkr\checkmarkc\checkmark$\checkmarkt\checkmark$\checkmark^\checkmark\\checkmarkc\checkmarki\checkmarkr\checkmarkc\checkmark$\checkmarke\checkmark$\checkmark^\checkmark\\checkmarkc\checkmarki\checkmarkr\checkmarkc\checkmark$\checkmarkr\checkmark$\checkmark^\checkmark\\checkmarkc\checkmarki\checkmarkr\checkmarkc\checkmark$\checkmark,\checkmark$\checkmark^\checkmark\\checkmarkc\checkmarki\checkmarkr\checkmarkc\checkmark$\checkmark \checkmark$\checkmark^\checkmark\\checkmarkc\checkmarki\checkmarkr\checkmarkc\checkmark$\checkmarkm\checkmark$\checkmark^\checkmark\\checkmarkc\checkmarki\checkmarkr\checkmarkc\checkmark$\checkmarki\checkmark$\checkmark^\checkmark\\checkmarkc\checkmarki\checkmarkr\checkmarkc\checkmark$\checkmarkn\checkmark$\checkmark^\checkmark\\checkmarkc\checkmarki\checkmarkr\checkmarkc\checkmark$\checkmarki\checkmark$\checkmark^\checkmark\\checkmarkc\checkmarki\checkmarkr\checkmarkc\checkmark$\checkmarkm\checkmark$\checkmark^\checkmark\\checkmarkc\checkmarki\checkmarkr\checkmarkc\checkmark$\checkmarku\checkmark$\checkmark^\checkmark\\checkmarkc\checkmarki\checkmarkr\checkmarkc\checkmark$\checkmarkm\checkmark$\checkmark^\checkmark\\checkmarkc\checkmarki\checkmarkr\checkmarkc\checkmark$\checkmark \checkmark$\checkmark^\checkmark\\checkmarkc\checkmarki\checkmarkr\checkmarkc\checkmark$\checkmarkh\checkmark$\checkmark^\checkmark\\checkmarkc\checkmarki\checkmarkr\checkmarkc\checkmark$\checkmarke\checkmark$\checkmark^\checkmark\\checkmarkc\checkmarki\checkmarkr\checkmarkc\checkmark$\checkmarki\checkmark$\checkmark^\checkmark\\checkmarkc\checkmarki\checkmarkr\checkmarkc\checkmark$\checkmarkg\checkmark$\checkmark^\checkmark\\checkmarkc\checkmarki\checkmarkr\checkmarkc\checkmark$\checkmarkh\checkmark$\checkmark^\checkmark\\checkmarkc\checkmarki\checkmarkr\checkmarkc\checkmark$\checkmarkt\checkmark$\checkmark^\checkmark\\checkmarkc\checkmarki\checkmarkr\checkmarkc\checkmark$\checkmark=\checkmark$\checkmark^\checkmark\\checkmarkc\checkmarki\checkmarkr\checkmarkc\checkmark$\checkmark1\checkmark$\checkmark^\checkmark\\checkmarkc\checkmarki\checkmarkr\checkmarkc\checkmark$\checkmark.\checkmark$\checkmark^\checkmark\\checkmarkc\checkmarki\checkmarkr\checkmarkc\checkmark$\checkmark4\checkmark$\checkmark^\checkmark\\checkmarkc\checkmarki\checkmarkr\checkmarkc\checkmark$\checkmarkc\checkmark$\checkmark^\checkmark\\checkmarkc\checkmarki\checkmarkr\checkmarkc\checkmark$\checkmarkm\checkmark$\checkmark^\checkmark\\checkmarkc\checkmarki\checkmarkr\checkmarkc\checkmark$\checkmark}\checkmark$\checkmark^\checkmark\\checkmarkc\checkmarki\checkmarkr\checkmarkc\checkmark$\checkmark,\checkmark$\checkmark^\checkmark\\checkmarkc\checkmarki\checkmarkr\checkmarkc\checkmark$\checkmark
\checkmark$\checkmark^\checkmark\\checkmarkc\checkmarki\checkmarkr\checkmarkc\checkmark$\checkmark \checkmark$\checkmark^\checkmark\\checkmarkc\checkmarki\checkmarkr\checkmarkc\checkmark$\checkmark \checkmark$\checkmark^\checkmark\\checkmarkc\checkmarki\checkmarkr\checkmarkc\checkmark$\checkmark \checkmark$\checkmark^\checkmark\\checkmarkc\checkmarki\checkmarkr\checkmarkc\checkmark$\checkmark \checkmark$\checkmark^\checkmark\\checkmarkc\checkmarki\checkmarkr\checkmarkc\checkmark$\checkmark \checkmark$\checkmark^\checkmark\\checkmarkc\checkmarki\checkmarkr\checkmarkc\checkmark$\checkmark \checkmark$\checkmark^\checkmark\\checkmarkc\checkmarki\checkmarkr\checkmarkc\checkmark$\checkmark \checkmark$\checkmark^\checkmark\\checkmarkc\checkmarki\checkmarkr\checkmarkc\checkmark$\checkmark \checkmark$\checkmark^\checkmark\\checkmarkc\checkmarki\checkmarkr\checkmarkc\checkmark$\checkmarkb\checkmark$\checkmark^\checkmark\\checkmarkc\checkmarki\checkmarkr\checkmarkc\checkmark$\checkmarko\checkmark$\checkmark^\checkmark\\checkmarkc\checkmarki\checkmarkr\checkmarkc\checkmark$\checkmarkx\checkmark$\checkmark^\checkmark\\checkmarkc\checkmarki\checkmarkr\checkmarkc\checkmark$\checkmark/\checkmark$\checkmark^\checkmark\\checkmarkc\checkmarki\checkmarkr\checkmarkc\checkmark$\checkmark.\checkmark$\checkmark^\checkmark\\checkmarkc\checkmarki\checkmarkr\checkmarkc\checkmark$\checkmarks\checkmark$\checkmark^\checkmark\\checkmarkc\checkmarki\checkmarkr\checkmarkc\checkmark$\checkmarkt\checkmark$\checkmark^\checkmark\\checkmarkc\checkmarki\checkmarkr\checkmarkc\checkmark$\checkmarky\checkmark$\checkmark^\checkmark\\checkmarkc\checkmarki\checkmarkr\checkmarkc\checkmark$\checkmarkl\checkmark$\checkmark^\checkmark\\checkmarkc\checkmarki\checkmarkr\checkmarkc\checkmark$\checkmarke\checkmark$\checkmark^\checkmark\\checkmarkc\checkmarki\checkmarkr\checkmarkc\checkmark$\checkmark=\checkmark$\checkmark^\checkmark\\checkmarkc\checkmarki\checkmarkr\checkmarkc\checkmark$\checkmark{\checkmark$\checkmark^\checkmark\\checkmarkc\checkmarki\checkmarkr\checkmarkc\checkmark$\checkmarkr\checkmark$\checkmark^\checkmark\\checkmarkc\checkmarki\checkmarkr\checkmarkc\checkmark$\checkmarke\checkmark$\checkmark^\checkmark\\checkmarkc\checkmarki\checkmarkr\checkmarkc\checkmark$\checkmarkc\checkmark$\checkmark^\checkmark\\checkmarkc\checkmarki\checkmarkr\checkmarkc\checkmark$\checkmarkt\checkmark$\checkmark^\checkmark\\checkmarkc\checkmarki\checkmarkr\checkmarkc\checkmark$\checkmarka\checkmark$\checkmark^\checkmark\\checkmarkc\checkmarki\checkmarkr\checkmarkc\checkmark$\checkmarkn\checkmark$\checkmark^\checkmark\\checkmarkc\checkmarki\checkmarkr\checkmarkc\checkmark$\checkmarkg\checkmark$\checkmark^\checkmark\\checkmarkc\checkmarki\checkmarkr\checkmarkc\checkmark$\checkmarkl\checkmark$\checkmark^\checkmark\\checkmarkc\checkmarki\checkmarkr\checkmarkc\checkmark$\checkmarke\checkmark$\checkmark^\checkmark\\checkmarkc\checkmarki\checkmarkr\checkmarkc\checkmark$\checkmark,\checkmark$\checkmark^\checkmark\\checkmarkc\checkmarki\checkmarkr\checkmarkc\checkmark$\checkmark \checkmark$\checkmark^\checkmark\\checkmarkc\checkmarki\checkmarkr\checkmarkc\checkmark$\checkmarkd\checkmark$\checkmark^\checkmark\\checkmarkc\checkmarki\checkmarkr\checkmarkc\checkmark$\checkmarkr\checkmark$\checkmark^\checkmark\\checkmarkc\checkmarki\checkmarkr\checkmarkc\checkmark$\checkmarka\checkmark$\checkmark^\checkmark\\checkmarkc\checkmarki\checkmarkr\checkmarkc\checkmark$\checkmarkw\checkmark$\checkmark^\checkmark\\checkmarkc\checkmarki\checkmarkr\checkmarkc\checkmark$\checkmark=\checkmark$\checkmark^\checkmark\\checkmarkc\checkmarki\checkmarkr\checkmarkc\checkmark$\checkmarkm\checkmark$\checkmark^\checkmark\\checkmarkc\checkmarki\checkmarkr\checkmarkc\checkmark$\checkmarka\checkmark$\checkmark^\checkmark\\checkmarkc\checkmarki\checkmarkr\checkmarkc\checkmark$\checkmarkc\checkmark$\checkmark^\checkmark\\checkmarkc\checkmarki\checkmarkr\checkmarkc\checkmark$\checkmarko\checkmark$\checkmark^\checkmark\\checkmarkc\checkmarki\checkmarkr\checkmarkc\checkmark$\checkmarku\checkmark$\checkmark^\checkmark\\checkmarkc\checkmarki\checkmarkr\checkmarkc\checkmark$\checkmarkl\checkmark$\checkmark^\checkmark\\checkmarkc\checkmarki\checkmarkr\checkmarkc\checkmark$\checkmarke\checkmark$\checkmark^\checkmark\\checkmarkc\checkmarki\checkmarkr\checkmarkc\checkmark$\checkmarku\checkmark$\checkmark^\checkmark\\checkmarkc\checkmarki\checkmarkr\checkmarkc\checkmark$\checkmarkr\checkmark$\checkmark^\checkmark\\checkmarkc\checkmarki\checkmarkr\checkmarkc\checkmark$\checkmark,\checkmark$\checkmark^\checkmark\\checkmarkc\checkmarki\checkmarkr\checkmarkc\checkmark$\checkmark \checkmark$\checkmark^\checkmark\\checkmarkc\checkmarki\checkmarkr\checkmarkc\checkmark$\checkmarkr\checkmark$\checkmark^\checkmark\\checkmarkc\checkmarki\checkmarkr\checkmarkc\checkmark$\checkmarko\checkmark$\checkmark^\checkmark\\checkmarkc\checkmarki\checkmarkr\checkmarkc\checkmark$\checkmarku\checkmark$\checkmark^\checkmark\\checkmarkc\checkmarki\checkmarkr\checkmarkc\checkmark$\checkmarkn\checkmark$\checkmark^\checkmark\\checkmarkc\checkmarki\checkmarkr\checkmarkc\checkmark$\checkmarkd\checkmark$\checkmark^\checkmark\\checkmarkc\checkmarki\checkmarkr\checkmarkc\checkmark$\checkmarke\checkmark$\checkmark^\checkmark\\checkmarkc\checkmarki\checkmarkr\checkmarkc\checkmark$\checkmarkd\checkmark$\checkmark^\checkmark\\checkmarkc\checkmarki\checkmarkr\checkmarkc\checkmark$\checkmark \checkmark$\checkmark^\checkmark\\checkmarkc\checkmarki\checkmarkr\checkmarkc\checkmark$\checkmarkc\checkmark$\checkmark^\checkmark\\checkmarkc\checkmarki\checkmarkr\checkmarkc\checkmark$\checkmarko\checkmark$\checkmark^\checkmark\\checkmarkc\checkmarki\checkmarkr\checkmarkc\checkmark$\checkmarkr\checkmark$\checkmark^\checkmark\\checkmarkc\checkmarki\checkmarkr\checkmarkc\checkmark$\checkmarkn\checkmark$\checkmark^\checkmark\\checkmarkc\checkmarki\checkmarkr\checkmarkc\checkmark$\checkmarke\checkmark$\checkmark^\checkmark\\checkmarkc\checkmarki\checkmarkr\checkmarkc\checkmark$\checkmarkr\checkmark$\checkmark^\checkmark\\checkmarkc\checkmarki\checkmarkr\checkmarkc\checkmark$\checkmarks\checkmark$\checkmark^\checkmark\\checkmarkc\checkmarki\checkmarkr\checkmarkc\checkmark$\checkmark,\checkmark$\checkmark^\checkmark\\checkmarkc\checkmarki\checkmarkr\checkmarkc\checkmark$\checkmark \checkmark$\checkmark^\checkmark\\checkmarkc\checkmarki\checkmarkr\checkmarkc\checkmark$\checkmarkm\checkmark$\checkmark^\checkmark\\checkmarkc\checkmarki\checkmarkr\checkmarkc\checkmark$\checkmarki\checkmark$\checkmark^\checkmark\\checkmarkc\checkmarki\checkmarkr\checkmarkc\checkmark$\checkmarkn\checkmark$\checkmark^\checkmark\\checkmarkc\checkmarki\checkmarkr\checkmarkc\checkmark$\checkmarki\checkmark$\checkmark^\checkmark\\checkmarkc\checkmarki\checkmarkr\checkmarkc\checkmark$\checkmarkm\checkmark$\checkmark^\checkmark\\checkmarkc\checkmarki\checkmarkr\checkmarkc\checkmark$\checkmarku\checkmark$\checkmark^\checkmark\\checkmarkc\checkmarki\checkmarkr\checkmarkc\checkmark$\checkmarkm\checkmark$\checkmark^\checkmark\\checkmarkc\checkmarki\checkmarkr\checkmarkc\checkmark$\checkmark \checkmark$\checkmark^\checkmark\\checkmarkc\checkmarki\checkmarkr\checkmarkc\checkmark$\checkmarkh\checkmark$\checkmark^\checkmark\\checkmarkc\checkmarki\checkmarkr\checkmarkc\checkmark$\checkmarke\checkmark$\checkmark^\checkmark\\checkmarkc\checkmarki\checkmarkr\checkmarkc\checkmark$\checkmarki\checkmark$\checkmark^\checkmark\\checkmarkc\checkmarki\checkmarkr\checkmarkc\checkmark$\checkmarkg\checkmark$\checkmark^\checkmark\\checkmarkc\checkmarki\checkmarkr\checkmarkc\checkmark$\checkmarkh\checkmark$\checkmark^\checkmark\\checkmarkc\checkmarki\checkmarkr\checkmarkc\checkmark$\checkmarkt\checkmark$\checkmark^\checkmark\\checkmarkc\checkmarki\checkmarkr\checkmarkc\checkmark$\checkmark=\checkmark$\checkmark^\checkmark\\checkmarkc\checkmarki\checkmarkr\checkmarkc\checkmark$\checkmark1\checkmark$\checkmark^\checkmark\\checkmarkc\checkmarki\checkmarkr\checkmarkc\checkmark$\checkmark.\checkmark$\checkmark^\checkmark\\checkmarkc\checkmarki\checkmarkr\checkmarkc\checkmark$\checkmark2\checkmark$\checkmark^\checkmark\\checkmarkc\checkmarki\checkmarkr\checkmarkc\checkmark$\checkmarkc\checkmark$\checkmark^\checkmark\\checkmarkc\checkmarki\checkmarkr\checkmarkc\checkmark$\checkmarkm\checkmark$\checkmark^\checkmark\\checkmarkc\checkmarki\checkmarkr\checkmarkc\checkmark$\checkmark,\checkmark$\checkmark^\checkmark\\checkmarkc\checkmarki\checkmarkr\checkmarkc\checkmark$\checkmark \checkmark$\checkmark^\checkmark\\checkmarkc\checkmarki\checkmarkr\checkmarkc\checkmark$\checkmarkt\checkmark$\checkmark^\checkmark\\checkmarkc\checkmarki\checkmarkr\checkmarkc\checkmark$\checkmarke\checkmark$\checkmark^\checkmark\\checkmarkc\checkmarki\checkmarkr\checkmarkc\checkmark$\checkmarkx\checkmark$\checkmark^\checkmark\\checkmarkc\checkmarki\checkmarkr\checkmarkc\checkmark$\checkmarkt\checkmark$\checkmark^\checkmark\\checkmarkc\checkmarki\checkmarkr\checkmarkc\checkmark$\checkmark \checkmark$\checkmark^\checkmark\\checkmarkc\checkmarki\checkmarkr\checkmarkc\checkmark$\checkmarkw\checkmark$\checkmark^\checkmark\\checkmarkc\checkmarki\checkmarkr\checkmarkc\checkmark$\checkmarki\checkmark$\checkmark^\checkmark\\checkmarkc\checkmarki\checkmarkr\checkmarkc\checkmark$\checkmarkd\checkmark$\checkmark^\checkmark\\checkmarkc\checkmarki\checkmarkr\checkmarkc\checkmark$\checkmarkt\checkmark$\checkmark^\checkmark\\checkmarkc\checkmarki\checkmarkr\checkmarkc\checkmark$\checkmarkh\checkmark$\checkmark^\checkmark\\checkmarkc\checkmarki\checkmarkr\checkmarkc\checkmark$\checkmark=\checkmark$\checkmark^\checkmark\\checkmarkc\checkmarki\checkmarkr\checkmarkc\checkmark$\checkmark5\checkmark$\checkmark^\checkmark\\checkmarkc\checkmarki\checkmarkr\checkmarkc\checkmark$\checkmarkc\checkmark$\checkmark^\checkmark\\checkmarkc\checkmarki\checkmarkr\checkmarkc\checkmark$\checkmarkm\checkmark$\checkmark^\checkmark\\checkmarkc\checkmarki\checkmarkr\checkmarkc\checkmark$\checkmark,\checkmark$\checkmark^\checkmark\\checkmarkc\checkmarki\checkmarkr\checkmarkc\checkmark$\checkmark \checkmark$\checkmark^\checkmark\\checkmarkc\checkmarki\checkmarkr\checkmarkc\checkmark$\checkmarka\checkmark$\checkmark^\checkmark\\checkmarkc\checkmarki\checkmarkr\checkmarkc\checkmark$\checkmarkl\checkmark$\checkmark^\checkmark\\checkmarkc\checkmarki\checkmarkr\checkmarkc\checkmark$\checkmarki\checkmark$\checkmark^\checkmark\\checkmarkc\checkmarki\checkmarkr\checkmarkc\checkmark$\checkmarkg\checkmark$\checkmark^\checkmark\\checkmarkc\checkmarki\checkmarkr\checkmarkc\checkmark$\checkmarkn\checkmark$\checkmark^\checkmark\\checkmarkc\checkmarki\checkmarkr\checkmarkc\checkmark$\checkmark=\checkmark$\checkmark^\checkmark\\checkmarkc\checkmarki\checkmarkr\checkmarkc\checkmark$\checkmarkc\checkmark$\checkmark^\checkmark\\checkmarkc\checkmarki\checkmarkr\checkmarkc\checkmark$\checkmarke\checkmark$\checkmark^\checkmark\\checkmarkc\checkmarki\checkmarkr\checkmarkc\checkmark$\checkmarkn\checkmark$\checkmark^\checkmark\\checkmarkc\checkmarki\checkmarkr\checkmarkc\checkmark$\checkmarkt\checkmark$\checkmark^\checkmark\\checkmarkc\checkmarki\checkmarkr\checkmarkc\checkmark$\checkmarke\checkmark$\checkmark^\checkmark\\checkmarkc\checkmarki\checkmarkr\checkmarkc\checkmark$\checkmarkr\checkmark$\checkmark^\checkmark\\checkmarkc\checkmarki\checkmarkr\checkmarkc\checkmark$\checkmark,\checkmark$\checkmark^\checkmark\\checkmarkc\checkmarki\checkmarkr\checkmarkc\checkmark$\checkmark \checkmark$\checkmark^\checkmark\\checkmarkc\checkmarki\checkmarkr\checkmarkc\checkmark$\checkmarkt\checkmark$\checkmark^\checkmark\\checkmarkc\checkmarki\checkmarkr\checkmarkc\checkmark$\checkmarkh\checkmark$\checkmark^\checkmark\\checkmarkc\checkmarki\checkmarkr\checkmarkc\checkmark$\checkmarki\checkmark$\checkmark^\checkmark\\checkmarkc\checkmarki\checkmarkr\checkmarkc\checkmark$\checkmarkc\checkmark$\checkmark^\checkmark\\checkmarkc\checkmarki\checkmarkr\checkmarkc\checkmark$\checkmarkk\checkmark$\checkmark^\checkmark\\checkmarkc\checkmarki\checkmarkr\checkmarkc\checkmark$\checkmark,\checkmark$\checkmark^\checkmark\\checkmarkc\checkmarki\checkmarkr\checkmarkc\checkmark$\checkmark \checkmark$\checkmark^\checkmark\\checkmarkc\checkmarki\checkmarkr\checkmarkc\checkmark$\checkmarkf\checkmark$\checkmark^\checkmark\\checkmarkc\checkmarki\checkmarkr\checkmarkc\checkmark$\checkmarki\checkmark$\checkmark^\checkmark\\checkmarkc\checkmarki\checkmarkr\checkmarkc\checkmark$\checkmarkl\checkmark$\checkmark^\checkmark\\checkmarkc\checkmarki\checkmarkr\checkmarkc\checkmark$\checkmarkl\checkmark$\checkmark^\checkmark\\checkmarkc\checkmarki\checkmarkr\checkmarkc\checkmark$\checkmark=\checkmark$\checkmark^\checkmark\\checkmarkc\checkmarki\checkmarkr\checkmarkc\checkmark$\checkmarkb\checkmark$\checkmark^\checkmark\\checkmarkc\checkmarki\checkmarkr\checkmarkc\checkmark$\checkmarkl\checkmark$\checkmark^\checkmark\\checkmarkc\checkmarki\checkmarkr\checkmarkc\checkmark$\checkmarku\checkmark$\checkmark^\checkmark\\checkmarkc\checkmarki\checkmarkr\checkmarkc\checkmark$\checkmarke\checkmark$\checkmark^\checkmark\\checkmarkc\checkmarki\checkmarkr\checkmarkc\checkmark$\checkmark!\checkmark$\checkmark^\checkmark\\checkmarkc\checkmarki\checkmarkr\checkmarkc\checkmark$\checkmark5\checkmark$\checkmark^\checkmark\\checkmarkc\checkmarki\checkmarkr\checkmarkc\checkmark$\checkmark}\checkmark$\checkmark^\checkmark\\checkmarkc\checkmarki\checkmarkr\checkmarkc\checkmark$\checkmark
\checkmark$\checkmark^\checkmark\\checkmarkc\checkmarki\checkmarkr\checkmarkc\checkmark$\checkmark \checkmark$\checkmark^\checkmark\\checkmarkc\checkmarki\checkmarkr\checkmarkc\checkmark$\checkmark \checkmark$\checkmark^\checkmark\\checkmarkc\checkmarki\checkmarkr\checkmarkc\checkmark$\checkmark \checkmark$\checkmark^\checkmark\\checkmarkc\checkmarki\checkmarkr\checkmarkc\checkmark$\checkmark \checkmark$\checkmark^\checkmark\\checkmarkc\checkmarki\checkmarkr\checkmarkc\checkmark$\checkmark]\checkmark$\checkmark^\checkmark\\checkmarkc\checkmarki\checkmarkr\checkmarkc\checkmark$\checkmark
\checkmark$\checkmark^\checkmark\\checkmarkc\checkmarki\checkmarkr\checkmarkc\checkmark$\checkmark \checkmark$\checkmark^\checkmark\\checkmarkc\checkmarki\checkmarkr\checkmarkc\checkmark$\checkmark \checkmark$\checkmark^\checkmark\\checkmarkc\checkmarki\checkmarkr\checkmarkc\checkmark$\checkmark \checkmark$\checkmark^\checkmark\\checkmarkc\checkmarki\checkmarkr\checkmarkc\checkmark$\checkmark \checkmark$\checkmark^\checkmark\\checkmarkc\checkmarki\checkmarkr\checkmarkc\checkmark$\checkmark \checkmark$\checkmark^\checkmark\\checkmarkc\checkmarki\checkmarkr\checkmarkc\checkmark$\checkmark \checkmark$\checkmark^\checkmark\\checkmarkc\checkmarki\checkmarkr\checkmarkc\checkmark$\checkmark \checkmark$\checkmark^\checkmark\\checkmarkc\checkmarki\checkmarkr\checkmarkc\checkmark$\checkmark \checkmark$\checkmark^\checkmark\\checkmarkc\checkmarki\checkmarkr\checkmarkc\checkmark$\checkmark\\checkmark$\checkmark^\checkmark\\checkmarkc\checkmarki\checkmarkr\checkmarkc\checkmark$\checkmarkn\checkmark$\checkmark^\checkmark\\checkmarkc\checkmarki\checkmarkr\checkmarkc\checkmark$\checkmarko\checkmark$\checkmark^\checkmark\\checkmarkc\checkmarki\checkmarkr\checkmarkc\checkmark$\checkmarkd\checkmark$\checkmark^\checkmark\\checkmarkc\checkmarki\checkmarkr\checkmarkc\checkmark$\checkmarke\checkmark$\checkmark^\checkmark\\checkmarkc\checkmarki\checkmarkr\checkmarkc\checkmark$\checkmark[\checkmark$\checkmark^\checkmark\\checkmarkc\checkmarki\checkmarkr\checkmarkc\checkmark$\checkmarkb\checkmark$\checkmark^\checkmark\\checkmarkc\checkmarki\checkmarkr\checkmarkc\checkmark$\checkmarku\checkmark$\checkmark^\checkmark\\checkmarkc\checkmarki\checkmarkr\checkmarkc\checkmark$\checkmarkb\checkmark$\checkmark^\checkmark\\checkmarkc\checkmarki\checkmarkr\checkmarkc\checkmark$\checkmarkb\checkmark$\checkmark^\checkmark\\checkmarkc\checkmarki\checkmarkr\checkmarkc\checkmark$\checkmarkl\checkmark$\checkmark^\checkmark\\checkmarkc\checkmarki\checkmarkr\checkmarkc\checkmark$\checkmarke\checkmark$\checkmark^\checkmark\\checkmarkc\checkmarki\checkmarkr\checkmarkc\checkmark$\checkmark]\checkmark$\checkmark^\checkmark\\checkmarkc\checkmarki\checkmarkr\checkmarkc\checkmark$\checkmark \checkmark$\checkmark^\checkmark\\checkmarkc\checkmarki\checkmarkr\checkmarkc\checkmark$\checkmark(\checkmark$\checkmark^\checkmark\\checkmarkc\checkmarki\checkmarkr\checkmarkc\checkmark$\checkmarku\checkmark$\checkmark^\checkmark\\checkmarkc\checkmarki\checkmarkr\checkmarkc\checkmark$\checkmarks\checkmark$\checkmark^\checkmark\\checkmarkc\checkmarki\checkmarkr\checkmarkc\checkmark$\checkmarke\checkmark$\checkmark^\checkmark\\checkmarkc\checkmarki\checkmarkr\checkmarkc\checkmark$\checkmarkr\checkmark$\checkmark^\checkmark\\checkmarkc\checkmarki\checkmarkr\checkmarkc\checkmark$\checkmark)\checkmark$\checkmark^\checkmark\\checkmarkc\checkmarki\checkmarkr\checkmarkc\checkmark$\checkmark \checkmark$\checkmark^\checkmark\\checkmarkc\checkmarki\checkmarkr\checkmarkc\checkmark$\checkmark{\checkmark$\checkmark^\checkmark\\checkmarkc\checkmarki\checkmarkr\checkmarkc\checkmark$\checkmarkI\checkmark$\checkmark^\checkmark\\checkmarkc\checkmarki\checkmarkr\checkmarkc\checkmark$\checkmarkn\checkmark$\checkmark^\checkmark\\checkmarkc\checkmarki\checkmarkr\checkmarkc\checkmark$\checkmarkg\checkmark$\checkmark^\checkmark\\checkmarkc\checkmarki\checkmarkr\checkmarkc\checkmark$\checkmarké\checkmark$\checkmark^\checkmark\\checkmarkc\checkmarki\checkmarkr\checkmarkc\checkmark$\checkmarkn\checkmark$\checkmark^\checkmark\\checkmarkc\checkmarki\checkmarkr\checkmarkc\checkmark$\checkmarki\checkmark$\checkmark^\checkmark\\checkmarkc\checkmarki\checkmarkr\checkmarkc\checkmark$\checkmarke\checkmark$\checkmark^\checkmark\\checkmarkc\checkmarki\checkmarkr\checkmarkc\checkmark$\checkmarku\checkmark$\checkmark^\checkmark\\checkmarkc\checkmarki\checkmarkr\checkmarkc\checkmark$\checkmarkr\checkmark$\checkmark^\checkmark\\checkmarkc\checkmarki\checkmarkr\checkmarkc\checkmark$\checkmark \checkmark$\checkmark^\checkmark\\checkmarkc\checkmarki\checkmarkr\checkmarkc\checkmark$\checkmark/\checkmark$\checkmark^\checkmark\\checkmarkc\checkmarki\checkmarkr\checkmarkc\checkmark$\checkmark\\checkmark$\checkmark^\checkmark\\checkmarkc\checkmarki\checkmarkr\checkmarkc\checkmark$\checkmark\\checkmark$\checkmark^\checkmark\\checkmarkc\checkmarki\checkmarkr\checkmarkc\checkmark$\checkmark \checkmark$\checkmark^\checkmark\\checkmarkc\checkmarki\checkmarkr\checkmarkc\checkmark$\checkmarkT\checkmark$\checkmark^\checkmark\\checkmarkc\checkmarki\checkmarkr\checkmarkc\checkmark$\checkmarke\checkmark$\checkmark^\checkmark\\checkmarkc\checkmarki\checkmarkr\checkmarkc\checkmark$\checkmarkc\checkmark$\checkmark^\checkmark\\checkmarkc\checkmarki\checkmarkr\checkmarkc\checkmark$\checkmarkh\checkmark$\checkmark^\checkmark\\checkmarkc\checkmarki\checkmarkr\checkmarkc\checkmark$\checkmarkn\checkmark$\checkmark^\checkmark\\checkmarkc\checkmarki\checkmarkr\checkmarkc\checkmark$\checkmarki\checkmark$\checkmark^\checkmark\\checkmarkc\checkmarki\checkmarkr\checkmarkc\checkmark$\checkmarkc\checkmark$\checkmark^\checkmark\\checkmarkc\checkmarki\checkmarkr\checkmarkc\checkmark$\checkmarki\checkmark$\checkmark^\checkmark\\checkmarkc\checkmarki\checkmarkr\checkmarkc\checkmark$\checkmarke\checkmark$\checkmark^\checkmark\\checkmarkc\checkmarki\checkmarkr\checkmarkc\checkmark$\checkmarkn\checkmark$\checkmark^\checkmark\\checkmarkc\checkmarki\checkmarkr\checkmarkc\checkmark$\checkmark}\checkmark$\checkmark^\checkmark\\checkmarkc\checkmarki\checkmarkr\checkmarkc\checkmark$\checkmark;\checkmark$\checkmark^\checkmark\\checkmarkc\checkmarki\checkmarkr\checkmarkc\checkmark$\checkmark
\checkmark$\checkmark^\checkmark\\checkmarkc\checkmarki\checkmarkr\checkmarkc\checkmark$\checkmark \checkmark$\checkmark^\checkmark\\checkmarkc\checkmarki\checkmarkr\checkmarkc\checkmark$\checkmark \checkmark$\checkmark^\checkmark\\checkmarkc\checkmarki\checkmarkr\checkmarkc\checkmark$\checkmark \checkmark$\checkmark^\checkmark\\checkmarkc\checkmarki\checkmarkr\checkmarkc\checkmark$\checkmark \checkmark$\checkmark^\checkmark\\checkmarkc\checkmarki\checkmarkr\checkmarkc\checkmark$\checkmark \checkmark$\checkmark^\checkmark\\checkmarkc\checkmarki\checkmarkr\checkmarkc\checkmark$\checkmark \checkmark$\checkmark^\checkmark\\checkmarkc\checkmarki\checkmarkr\checkmarkc\checkmark$\checkmark \checkmark$\checkmark^\checkmark\\checkmarkc\checkmarki\checkmarkr\checkmarkc\checkmark$\checkmark \checkmark$\checkmark^\checkmark\\checkmarkc\checkmarki\checkmarkr\checkmarkc\checkmark$\checkmark\\checkmark$\checkmark^\checkmark\\checkmarkc\checkmarki\checkmarkr\checkmarkc\checkmark$\checkmarkn\checkmark$\checkmark^\checkmark\\checkmarkc\checkmarki\checkmarkr\checkmarkc\checkmark$\checkmarko\checkmark$\checkmark^\checkmark\\checkmarkc\checkmarki\checkmarkr\checkmarkc\checkmark$\checkmarkd\checkmark$\checkmark^\checkmark\\checkmarkc\checkmarki\checkmarkr\checkmarkc\checkmark$\checkmarke\checkmark$\checkmark^\checkmark\\checkmarkc\checkmarki\checkmarkr\checkmarkc\checkmark$\checkmark[\checkmark$\checkmark^\checkmark\\checkmarkc\checkmarki\checkmarkr\checkmarkc\checkmark$\checkmarkb\checkmark$\checkmark^\checkmark\\checkmarkc\checkmarki\checkmarkr\checkmarkc\checkmark$\checkmarku\checkmark$\checkmark^\checkmark\\checkmarkc\checkmarki\checkmarkr\checkmarkc\checkmark$\checkmarkb\checkmark$\checkmark^\checkmark\\checkmarkc\checkmarki\checkmarkr\checkmarkc\checkmark$\checkmarkb\checkmark$\checkmark^\checkmark\\checkmarkc\checkmarki\checkmarkr\checkmarkc\checkmark$\checkmarkl\checkmark$\checkmark^\checkmark\\checkmarkc\checkmarki\checkmarkr\checkmarkc\checkmark$\checkmarke\checkmark$\checkmark^\checkmark\\checkmarkc\checkmarki\checkmarkr\checkmarkc\checkmark$\checkmark,\checkmark$\checkmark^\checkmark\\checkmarkc\checkmarki\checkmarkr\checkmarkc\checkmark$\checkmark \checkmark$\checkmark^\checkmark\\checkmarkc\checkmarki\checkmarkr\checkmarkc\checkmark$\checkmarkr\checkmark$\checkmark^\checkmark\\checkmarkc\checkmarki\checkmarkr\checkmarkc\checkmark$\checkmarki\checkmark$\checkmark^\checkmark\\checkmarkc\checkmarki\checkmarkr\checkmarkc\checkmark$\checkmarkg\checkmark$\checkmark^\checkmark\\checkmarkc\checkmarki\checkmarkr\checkmarkc\checkmark$\checkmarkh\checkmark$\checkmark^\checkmark\\checkmarkc\checkmarki\checkmarkr\checkmarkc\checkmark$\checkmarkt\checkmark$\checkmark^\checkmark\\checkmarkc\checkmarki\checkmarkr\checkmarkc\checkmark$\checkmark=\checkmark$\checkmark^\checkmark\\checkmarkc\checkmarki\checkmarkr\checkmarkc\checkmark$\checkmark5\checkmark$\checkmark^\checkmark\\checkmarkc\checkmarki\checkmarkr\checkmarkc\checkmark$\checkmarkc\checkmark$\checkmark^\checkmark\\checkmarkc\checkmarki\checkmarkr\checkmarkc\checkmark$\checkmarkm\checkmark$\checkmark^\checkmark\\checkmarkc\checkmarki\checkmarkr\checkmarkc\checkmark$\checkmark \checkmark$\checkmark^\checkmark\\checkmarkc\checkmarki\checkmarkr\checkmarkc\checkmark$\checkmarko\checkmark$\checkmark^\checkmark\\checkmarkc\checkmarki\checkmarkr\checkmarkc\checkmark$\checkmarkf\checkmark$\checkmark^\checkmark\\checkmarkc\checkmarki\checkmarkr\checkmarkc\checkmark$\checkmark \checkmark$\checkmark^\checkmark\\checkmarkc\checkmarki\checkmarkr\checkmarkc\checkmark$\checkmarku\checkmark$\checkmark^\checkmark\\checkmarkc\checkmarki\checkmarkr\checkmarkc\checkmark$\checkmarks\checkmark$\checkmark^\checkmark\\checkmarkc\checkmarki\checkmarkr\checkmarkc\checkmark$\checkmarke\checkmark$\checkmark^\checkmark\\checkmarkc\checkmarki\checkmarkr\checkmarkc\checkmark$\checkmarkr\checkmark$\checkmark^\checkmark\\checkmarkc\checkmarki\checkmarkr\checkmarkc\checkmark$\checkmark]\checkmark$\checkmark^\checkmark\\checkmarkc\checkmarki\checkmarkr\checkmarkc\checkmark$\checkmark \checkmark$\checkmark^\checkmark\\checkmarkc\checkmarki\checkmarkr\checkmarkc\checkmark$\checkmark(\checkmark$\checkmark^\checkmark\\checkmarkc\checkmarki\checkmarkr\checkmarkc\checkmark$\checkmarkm\checkmark$\checkmark^\checkmark\\checkmarkc\checkmarki\checkmarkr\checkmarkc\checkmark$\checkmarka\checkmark$\checkmark^\checkmark\\checkmarkc\checkmarki\checkmarkr\checkmarkc\checkmark$\checkmarkt\checkmark$\checkmark^\checkmark\\checkmarkc\checkmarki\checkmarkr\checkmarkc\checkmark$\checkmarki\checkmark$\checkmark^\checkmark\\checkmarkc\checkmarki\checkmarkr\checkmarkc\checkmark$\checkmarke\checkmark$\checkmark^\checkmark\\checkmarkc\checkmarki\checkmarkr\checkmarkc\checkmark$\checkmarkr\checkmark$\checkmark^\checkmark\\checkmarkc\checkmarki\checkmarkr\checkmarkc\checkmark$\checkmarke\checkmark$\checkmark^\checkmark\\checkmarkc\checkmarki\checkmarkr\checkmarkc\checkmark$\checkmark)\checkmark$\checkmark^\checkmark\\checkmarkc\checkmarki\checkmarkr\checkmarkc\checkmark$\checkmark \checkmark$\checkmark^\checkmark\\checkmarkc\checkmarki\checkmarkr\checkmarkc\checkmark$\checkmark{\checkmark$\checkmark^\checkmark\\checkmarkc\checkmarki\checkmarkr\checkmarkc\checkmark$\checkmarkP\checkmark$\checkmark^\checkmark\\checkmarkc\checkmarki\checkmarkr\checkmarkc\checkmark$\checkmarki\checkmark$\checkmark^\checkmark\\checkmarkc\checkmarki\checkmarkr\checkmarkc\checkmark$\checkmarkè\checkmark$\checkmark^\checkmark\\checkmarkc\checkmarki\checkmarkr\checkmarkc\checkmark$\checkmarkc\checkmark$\checkmark^\checkmark\\checkmarkc\checkmarki\checkmarkr\checkmarkc\checkmark$\checkmarke\checkmark$\checkmark^\checkmark\\checkmarkc\checkmarki\checkmarkr\checkmarkc\checkmark$\checkmark \checkmark$\checkmark^\checkmark\\checkmarkc\checkmarki\checkmarkr\checkmarkc\checkmark$\checkmarkb\checkmark$\checkmark^\checkmark\\checkmarkc\checkmarki\checkmarkr\checkmarkc\checkmark$\checkmarkr\checkmark$\checkmark^\checkmark\\checkmarkc\checkmarki\checkmarkr\checkmarkc\checkmark$\checkmarku\checkmark$\checkmark^\checkmark\\checkmarkc\checkmarki\checkmarkr\checkmarkc\checkmark$\checkmarkt\checkmark$\checkmark^\checkmark\\checkmarkc\checkmarki\checkmarkr\checkmarkc\checkmark$\checkmarke\checkmark$\checkmark^\checkmark\\checkmarkc\checkmarki\checkmarkr\checkmarkc\checkmark$\checkmark\\checkmark$\checkmark^\checkmark\\checkmarkc\checkmarki\checkmarkr\checkmarkc\checkmark$\checkmark\\checkmark$\checkmark^\checkmark\\checkmarkc\checkmarki\checkmarkr\checkmarkc\checkmark$\checkmark \checkmark$\checkmark^\checkmark\\checkmarkc\checkmarki\checkmarkr\checkmarkc\checkmark$\checkmark(\checkmark$\checkmark^\checkmark\\checkmarkc\checkmarki\checkmarkr\checkmarkc\checkmark$\checkmarkA\checkmark$\checkmark^\checkmark\\checkmarkc\checkmarki\checkmarkr\checkmarkc\checkmark$\checkmarkl\checkmark$\checkmark^\checkmark\\checkmarkc\checkmarki\checkmarkr\checkmarkc\checkmark$\checkmark,\checkmark$\checkmark^\checkmark\\checkmarkc\checkmarki\checkmarkr\checkmarkc\checkmark$\checkmark \checkmark$\checkmark^\checkmark\\checkmarkc\checkmarki\checkmarkr\checkmarkc\checkmark$\checkmarkB\checkmark$\checkmark^\checkmark\\checkmarkc\checkmarki\checkmarkr\checkmarkc\checkmark$\checkmarko\checkmark$\checkmark^\checkmark\\checkmarkc\checkmarki\checkmarkr\checkmarkc\checkmark$\checkmarki\checkmark$\checkmark^\checkmark\\checkmarkc\checkmarki\checkmarkr\checkmarkc\checkmark$\checkmarks\checkmark$\checkmark^\checkmark\\checkmarkc\checkmarki\checkmarkr\checkmarkc\checkmark$\checkmark,\checkmark$\checkmark^\checkmark\\checkmarkc\checkmarki\checkmarkr\checkmarkc\checkmark$\checkmark \checkmark$\checkmark^\checkmark\\checkmarkc\checkmarki\checkmarkr\checkmarkc\checkmark$\checkmarkR\checkmark$\checkmark^\checkmark\\checkmarkc\checkmarki\checkmarkr\checkmarkc\checkmark$\checkmarké\checkmark$\checkmark^\checkmark\\checkmarkc\checkmarki\checkmarkr\checkmarkc\checkmark$\checkmarks\checkmark$\checkmark^\checkmark\\checkmarkc\checkmarki\checkmarkr\checkmarkc\checkmark$\checkmarki\checkmark$\checkmark^\checkmark\\checkmarkc\checkmarki\checkmarkr\checkmarkc\checkmark$\checkmarkn\checkmark$\checkmark^\checkmark\\checkmarkc\checkmarki\checkmarkr\checkmarkc\checkmark$\checkmarke\checkmark$\checkmark^\checkmark\\checkmarkc\checkmarki\checkmarkr\checkmarkc\checkmark$\checkmark)\checkmark$\checkmark^\checkmark\\checkmarkc\checkmarki\checkmarkr\checkmarkc\checkmark$\checkmark}\checkmark$\checkmark^\checkmark\\checkmarkc\checkmarki\checkmarkr\checkmarkc\checkmark$\checkmark;\checkmark$\checkmark^\checkmark\\checkmarkc\checkmarki\checkmarkr\checkmarkc\checkmark$\checkmark
\checkmark$\checkmark^\checkmark\\checkmarkc\checkmarki\checkmarkr\checkmarkc\checkmark$\checkmark \checkmark$\checkmark^\checkmark\\checkmarkc\checkmarki\checkmarkr\checkmarkc\checkmark$\checkmark \checkmark$\checkmark^\checkmark\\checkmarkc\checkmarki\checkmarkr\checkmarkc\checkmark$\checkmark \checkmark$\checkmark^\checkmark\\checkmarkc\checkmarki\checkmarkr\checkmarkc\checkmark$\checkmark \checkmark$\checkmark^\checkmark\\checkmarkc\checkmarki\checkmarkr\checkmarkc\checkmark$\checkmark \checkmark$\checkmark^\checkmark\\checkmarkc\checkmarki\checkmarkr\checkmarkc\checkmark$\checkmark \checkmark$\checkmark^\checkmark\\checkmarkc\checkmarki\checkmarkr\checkmarkc\checkmark$\checkmark \checkmark$\checkmark^\checkmark\\checkmarkc\checkmarki\checkmarkr\checkmarkc\checkmark$\checkmark \checkmark$\checkmark^\checkmark\\checkmarkc\checkmarki\checkmarkr\checkmarkc\checkmark$\checkmark\\checkmark$\checkmark^\checkmark\\checkmarkc\checkmarki\checkmarkr\checkmarkc\checkmark$\checkmarkn\checkmark$\checkmark^\checkmark\\checkmarkc\checkmarki\checkmarkr\checkmarkc\checkmark$\checkmarko\checkmark$\checkmark^\checkmark\\checkmarkc\checkmarki\checkmarkr\checkmarkc\checkmark$\checkmarkd\checkmark$\checkmark^\checkmark\\checkmarkc\checkmarki\checkmarkr\checkmarkc\checkmark$\checkmarke\checkmark$\checkmark^\checkmark\\checkmarkc\checkmarki\checkmarkr\checkmarkc\checkmark$\checkmark[\checkmark$\checkmark^\checkmark\\checkmarkc\checkmarki\checkmarkr\checkmarkc\checkmark$\checkmarkb\checkmark$\checkmark^\checkmark\\checkmarkc\checkmarki\checkmarkr\checkmarkc\checkmark$\checkmarko\checkmark$\checkmark^\checkmark\\checkmarkc\checkmarki\checkmarkr\checkmarkc\checkmark$\checkmarkx\checkmark$\checkmark^\checkmark\\checkmarkc\checkmarki\checkmarkr\checkmarkc\checkmark$\checkmark,\checkmark$\checkmark^\checkmark\\checkmarkc\checkmarki\checkmarkr\checkmarkc\checkmark$\checkmark \checkmark$\checkmark^\checkmark\\checkmarkc\checkmarki\checkmarkr\checkmarkc\checkmark$\checkmarkb\checkmark$\checkmark^\checkmark\\checkmarkc\checkmarki\checkmarkr\checkmarkc\checkmark$\checkmarke\checkmark$\checkmark^\checkmark\\checkmarkc\checkmarki\checkmarkr\checkmarkc\checkmark$\checkmarkl\checkmark$\checkmark^\checkmark\\checkmarkc\checkmarki\checkmarkr\checkmarkc\checkmark$\checkmarko\checkmark$\checkmark^\checkmark\\checkmarkc\checkmarki\checkmarkr\checkmarkc\checkmark$\checkmarkw\checkmark$\checkmark^\checkmark\\checkmarkc\checkmarki\checkmarkr\checkmarkc\checkmark$\checkmark=\checkmark$\checkmark^\checkmark\\checkmarkc\checkmarki\checkmarkr\checkmarkc\checkmark$\checkmark2\checkmark$\checkmark^\checkmark\\checkmarkc\checkmarki\checkmarkr\checkmarkc\checkmark$\checkmarkc\checkmark$\checkmark^\checkmark\\checkmarkc\checkmarki\checkmarkr\checkmarkc\checkmark$\checkmarkm\checkmark$\checkmark^\checkmark\\checkmarkc\checkmarki\checkmarkr\checkmarkc\checkmark$\checkmark \checkmark$\checkmark^\checkmark\\checkmarkc\checkmarki\checkmarkr\checkmarkc\checkmark$\checkmarko\checkmark$\checkmark^\checkmark\\checkmarkc\checkmarki\checkmarkr\checkmarkc\checkmark$\checkmarkf\checkmark$\checkmark^\checkmark\\checkmarkc\checkmarki\checkmarkr\checkmarkc\checkmark$\checkmark \checkmark$\checkmark^\checkmark\\checkmarkc\checkmarki\checkmarkr\checkmarkc\checkmark$\checkmarku\checkmark$\checkmark^\checkmark\\checkmarkc\checkmarki\checkmarkr\checkmarkc\checkmark$\checkmarks\checkmark$\checkmark^\checkmark\\checkmarkc\checkmarki\checkmarkr\checkmarkc\checkmark$\checkmarke\checkmark$\checkmark^\checkmark\\checkmarkc\checkmarki\checkmarkr\checkmarkc\checkmark$\checkmarkr\checkmark$\checkmark^\checkmark\\checkmarkc\checkmarki\checkmarkr\checkmarkc\checkmark$\checkmark,\checkmark$\checkmark^\checkmark\\checkmarkc\checkmarki\checkmarkr\checkmarkc\checkmark$\checkmark \checkmark$\checkmark^\checkmark\\checkmarkc\checkmarki\checkmarkr\checkmarkc\checkmark$\checkmarkx\checkmark$\checkmark^\checkmark\\checkmarkc\checkmarki\checkmarkr\checkmarkc\checkmark$\checkmarks\checkmark$\checkmark^\checkmark\\checkmarkc\checkmarki\checkmarkr\checkmarkc\checkmark$\checkmarkh\checkmark$\checkmark^\checkmark\\checkmarkc\checkmarki\checkmarkr\checkmarkc\checkmark$\checkmarki\checkmark$\checkmark^\checkmark\\checkmarkc\checkmarki\checkmarkr\checkmarkc\checkmark$\checkmarkf\checkmark$\checkmark^\checkmark\\checkmarkc\checkmarki\checkmarkr\checkmarkc\checkmark$\checkmarkt\checkmark$\checkmark^\checkmark\\checkmarkc\checkmarki\checkmarkr\checkmarkc\checkmark$\checkmark=\checkmark$\checkmark^\checkmark\\checkmarkc\checkmarki\checkmarkr\checkmarkc\checkmark$\checkmark4\checkmark$\checkmark^\checkmark\\checkmarkc\checkmarki\checkmarkr\checkmarkc\checkmark$\checkmarkc\checkmark$\checkmark^\checkmark\\checkmarkc\checkmarki\checkmarkr\checkmarkc\checkmark$\checkmarkm\checkmark$\checkmark^\checkmark\\checkmarkc\checkmarki\checkmarkr\checkmarkc\checkmark$\checkmark]\checkmark$\checkmark^\checkmark\\checkmarkc\checkmarki\checkmarkr\checkmarkc\checkmark$\checkmark \checkmark$\checkmark^\checkmark\\checkmarkc\checkmarki\checkmarkr\checkmarkc\checkmark$\checkmark(\checkmark$\checkmark^\checkmark\\checkmarkc\checkmarki\checkmarkr\checkmarkc\checkmark$\checkmarks\checkmark$\checkmark^\checkmark\\checkmarkc\checkmarki\checkmarkr\checkmarkc\checkmark$\checkmarky\checkmark$\checkmark^\checkmark\\checkmarkc\checkmarki\checkmarkr\checkmarkc\checkmark$\checkmarks\checkmark$\checkmark^\checkmark\\checkmarkc\checkmarki\checkmarkr\checkmarkc\checkmark$\checkmark)\checkmark$\checkmark^\checkmark\\checkmarkc\checkmarki\checkmarkr\checkmarkc\checkmark$\checkmark \checkmark$\checkmark^\checkmark\\checkmarkc\checkmarki\checkmarkr\checkmarkc\checkmark$\checkmark{\checkmark$\checkmark^\checkmark\\checkmarkc\checkmarki\checkmarkr\checkmarkc\checkmark$\checkmark\\checkmark$\checkmark^\checkmark\\checkmarkc\checkmarki\checkmarkr\checkmarkc\checkmark$\checkmarkt\checkmark$\checkmark^\checkmark\\checkmarkc\checkmarki\checkmarkr\checkmarkc\checkmark$\checkmarke\checkmark$\checkmark^\checkmark\\checkmarkc\checkmarki\checkmarkr\checkmarkc\checkmark$\checkmarkx\checkmark$\checkmark^\checkmark\\checkmarkc\checkmarki\checkmarkr\checkmarkc\checkmark$\checkmarkt\checkmark$\checkmark^\checkmark\\checkmarkc\checkmarki\checkmarkr\checkmarkc\checkmark$\checkmarkb\checkmark$\checkmark^\checkmark\\checkmarkc\checkmarki\checkmarkr\checkmarkc\checkmark$\checkmarkf\checkmark$\checkmark^\checkmark\\checkmarkc\checkmarki\checkmarkr\checkmarkc\checkmark$\checkmark{\checkmark$\checkmark^\checkmark\\checkmarkc\checkmarki\checkmarkr\checkmarkc\checkmark$\checkmarkM\checkmark$\checkmark^\checkmark\\checkmarkc\checkmarki\checkmarkr\checkmarkc\checkmark$\checkmarki\checkmark$\checkmark^\checkmark\\checkmarkc\checkmarki\checkmarkr\checkmarkc\checkmark$\checkmarkn\checkmark$\checkmark^\checkmark\\checkmarkc\checkmarki\checkmarkr\checkmarkc\checkmark$\checkmarki\checkmark$\checkmark^\checkmark\\checkmarkc\checkmarki\checkmarkr\checkmarkc\checkmark$\checkmark-\checkmark$\checkmark^\checkmark\\checkmarkc\checkmarki\checkmarkr\checkmarkc\checkmark$\checkmarkF\checkmark$\checkmark^\checkmark\\checkmarkc\checkmarki\checkmarkr\checkmarkc\checkmark$\checkmarkr\checkmark$\checkmark^\checkmark\\checkmarkc\checkmarki\checkmarkr\checkmarkc\checkmark$\checkmarka\checkmark$\checkmark^\checkmark\\checkmarkc\checkmarki\checkmarkr\checkmarkc\checkmark$\checkmarki\checkmark$\checkmark^\checkmark\\checkmarkc\checkmarki\checkmarkr\checkmarkc\checkmark$\checkmarks\checkmark$\checkmark^\checkmark\\checkmarkc\checkmarki\checkmarkr\checkmarkc\checkmark$\checkmarke\checkmark$\checkmark^\checkmark\\checkmarkc\checkmarki\checkmarkr\checkmarkc\checkmark$\checkmarku\checkmark$\checkmark^\checkmark\\checkmarkc\checkmarki\checkmarkr\checkmarkc\checkmark$\checkmarks\checkmark$\checkmark^\checkmark\\checkmarkc\checkmarki\checkmarkr\checkmarkc\checkmark$\checkmarke\checkmark$\checkmark^\checkmark\\checkmarkc\checkmarki\checkmarkr\checkmarkc\checkmark$\checkmark \checkmark$\checkmark^\checkmark\\checkmarkc\checkmarki\checkmarkr\checkmarkc\checkmark$\checkmarkC\checkmark$\checkmark^\checkmark\\checkmarkc\checkmarki\checkmarkr\checkmarkc\checkmark$\checkmarkN\checkmark$\checkmark^\checkmark\\checkmarkc\checkmarki\checkmarkr\checkmarkc\checkmark$\checkmarkC\checkmark$\checkmark^\checkmark\\checkmarkc\checkmarki\checkmarkr\checkmarkc\checkmark$\checkmark\\checkmark$\checkmark^\checkmark\\checkmarkc\checkmarki\checkmarkr\checkmarkc\checkmark$\checkmark\\checkmark$\checkmark^\checkmark\\checkmarkc\checkmarki\checkmarkr\checkmarkc\checkmark$\checkmark \checkmark$\checkmark^\checkmark\\checkmarkc\checkmarki\checkmarkr\checkmarkc\checkmark$\checkmark5\checkmark$\checkmark^\checkmark\\checkmarkc\checkmarki\checkmarkr\checkmarkc\checkmark$\checkmark \checkmark$\checkmark^\checkmark\\checkmarkc\checkmarki\checkmarkr\checkmarkc\checkmark$\checkmarkA\checkmark$\checkmark^\checkmark\\checkmarkc\checkmarki\checkmarkr\checkmarkc\checkmark$\checkmarkx\checkmark$\checkmark^\checkmark\\checkmarkc\checkmarki\checkmarkr\checkmarkc\checkmark$\checkmarke\checkmark$\checkmark^\checkmark\\checkmarkc\checkmarki\checkmarkr\checkmarkc\checkmark$\checkmarks\checkmark$\checkmark^\checkmark\\checkmarkc\checkmarki\checkmarkr\checkmarkc\checkmark$\checkmark \checkmark$\checkmark^\checkmark\\checkmarkc\checkmarki\checkmarkr\checkmarkc\checkmark$\checkmark«\checkmark$\checkmark^\checkmark\\checkmarkc\checkmarki\checkmarkr\checkmarkc\checkmark$\checkmarkF\checkmark$\checkmark^\checkmark\\checkmarkc\checkmarki\checkmarkr\checkmarkc\checkmark$\checkmarkO\checkmark$\checkmark^\checkmark\\checkmarkc\checkmarki\checkmarkr\checkmarkc\checkmark$\checkmarkR\checkmark$\checkmark^\checkmark\\checkmarkc\checkmarki\checkmarkr\checkmarkc\checkmark$\checkmarkG\checkmark$\checkmark^\checkmark\\checkmarkc\checkmarki\checkmarkr\checkmarkc\checkmark$\checkmarkE\checkmark$\checkmark^\checkmark\\checkmarkc\checkmarki\checkmarkr\checkmarkc\checkmark$\checkmarkR\checkmark$\checkmark^\checkmark\\checkmarkc\checkmarki\checkmarkr\checkmarkc\checkmark$\checkmarkO\checkmark$\checkmark^\checkmark\\checkmarkc\checkmarki\checkmarkr\checkmarkc\checkmark$\checkmarkN\checkmark$\checkmark^\checkmark\\checkmarkc\checkmarki\checkmarkr\checkmarkc\checkmark$\checkmark»\checkmark$\checkmark^\checkmark\\checkmarkc\checkmarki\checkmarkr\checkmarkc\checkmark$\checkmark}\checkmark$\checkmark^\checkmark\\checkmarkc\checkmarki\checkmarkr\checkmarkc\checkmark$\checkmark}\checkmark$\checkmark^\checkmark\\checkmarkc\checkmarki\checkmarkr\checkmarkc\checkmark$\checkmark;\checkmark$\checkmark^\checkmark\\checkmarkc\checkmarki\checkmarkr\checkmarkc\checkmark$\checkmark
\checkmark$\checkmark^\checkmark\\checkmarkc\checkmarki\checkmarkr\checkmarkc\checkmark$\checkmark \checkmark$\checkmark^\checkmark\\checkmarkc\checkmarki\checkmarkr\checkmarkc\checkmark$\checkmark \checkmark$\checkmark^\checkmark\\checkmarkc\checkmarki\checkmarkr\checkmarkc\checkmark$\checkmark \checkmark$\checkmark^\checkmark\\checkmarkc\checkmarki\checkmarkr\checkmarkc\checkmark$\checkmark \checkmark$\checkmark^\checkmark\\checkmarkc\checkmarki\checkmarkr\checkmarkc\checkmark$\checkmark \checkmark$\checkmark^\checkmark\\checkmarkc\checkmarki\checkmarkr\checkmarkc\checkmark$\checkmark \checkmark$\checkmark^\checkmark\\checkmarkc\checkmarki\checkmarkr\checkmarkc\checkmark$\checkmark \checkmark$\checkmark^\checkmark\\checkmarkc\checkmarki\checkmarkr\checkmarkc\checkmark$\checkmark \checkmark$\checkmark^\checkmark\\checkmarkc\checkmarki\checkmarkr\checkmarkc\checkmark$\checkmark\\checkmark$\checkmark^\checkmark\\checkmarkc\checkmarki\checkmarkr\checkmarkc\checkmark$\checkmarkd\checkmark$\checkmark^\checkmark\\checkmarkc\checkmarki\checkmarkr\checkmarkc\checkmark$\checkmarkr\checkmark$\checkmark^\checkmark\\checkmarkc\checkmarki\checkmarkr\checkmarkc\checkmark$\checkmarka\checkmark$\checkmark^\checkmark\\checkmarkc\checkmarki\checkmarkr\checkmarkc\checkmark$\checkmarkw\checkmark$\checkmark^\checkmark\\checkmarkc\checkmarki\checkmarkr\checkmarkc\checkmark$\checkmark[\checkmark$\checkmark^\checkmark\\checkmarkc\checkmarki\checkmarkr\checkmarkc\checkmark$\checkmarkt\checkmark$\checkmark^\checkmark\\checkmarkc\checkmarki\checkmarkr\checkmarkc\checkmark$\checkmarkh\checkmark$\checkmark^\checkmark\\checkmarkc\checkmarki\checkmarkr\checkmarkc\checkmark$\checkmarki\checkmark$\checkmark^\checkmark\\checkmarkc\checkmarki\checkmarkr\checkmarkc\checkmark$\checkmarkc\checkmark$\checkmark^\checkmark\\checkmarkc\checkmarki\checkmarkr\checkmarkc\checkmark$\checkmarkk\checkmark$\checkmark^\checkmark\\checkmarkc\checkmarki\checkmarkr\checkmarkc\checkmark$\checkmark,\checkmark$\checkmark^\checkmark\\checkmarkc\checkmarki\checkmarkr\checkmarkc\checkmark$\checkmark \checkmark$\checkmark^\checkmark\\checkmarkc\checkmarki\checkmarkr\checkmarkc\checkmark$\checkmark-\checkmark$\checkmark^\checkmark\\checkmarkc\checkmarki\checkmarkr\checkmarkc\checkmark$\checkmark{\checkmark$\checkmark^\checkmark\\checkmarkc\checkmarki\checkmarkr\checkmarkc\checkmark$\checkmarkS\checkmark$\checkmark^\checkmark\\checkmarkc\checkmarki\checkmarkr\checkmarkc\checkmark$\checkmarkt\checkmark$\checkmark^\checkmark\\checkmarkc\checkmarki\checkmarkr\checkmarkc\checkmark$\checkmarke\checkmark$\checkmark^\checkmark\\checkmarkc\checkmarki\checkmarkr\checkmarkc\checkmark$\checkmarka\checkmark$\checkmark^\checkmark\\checkmarkc\checkmarki\checkmarkr\checkmarkc\checkmark$\checkmarkl\checkmark$\checkmark^\checkmark\\checkmarkc\checkmarki\checkmarkr\checkmarkc\checkmark$\checkmarkt\checkmark$\checkmark^\checkmark\\checkmarkc\checkmarki\checkmarkr\checkmarkc\checkmark$\checkmarkh\checkmark$\checkmark^\checkmark\\checkmarkc\checkmarki\checkmarkr\checkmarkc\checkmark$\checkmark[\checkmark$\checkmark^\checkmark\\checkmarkc\checkmarki\checkmarkr\checkmarkc\checkmark$\checkmarks\checkmark$\checkmark^\checkmark\\checkmarkc\checkmarki\checkmarkr\checkmarkc\checkmark$\checkmarkc\checkmark$\checkmark^\checkmark\\checkmarkc\checkmarki\checkmarkr\checkmarkc\checkmark$\checkmarka\checkmark$\checkmark^\checkmark\\checkmarkc\checkmarki\checkmarkr\checkmarkc\checkmark$\checkmarkl\checkmark$\checkmark^\checkmark\\checkmarkc\checkmarki\checkmarkr\checkmarkc\checkmark$\checkmarke\checkmark$\checkmark^\checkmark\\checkmarkc\checkmarki\checkmarkr\checkmarkc\checkmark$\checkmark=\checkmark$\checkmark^\checkmark\\checkmarkc\checkmarki\checkmarkr\checkmarkc\checkmark$\checkmark1\checkmark$\checkmark^\checkmark\\checkmarkc\checkmarki\checkmarkr\checkmarkc\checkmark$\checkmark.\checkmark$\checkmark^\checkmark\\checkmarkc\checkmarki\checkmarkr\checkmarkc\checkmark$\checkmark2\checkmark$\checkmark^\checkmark\\checkmarkc\checkmarki\checkmarkr\checkmarkc\checkmark$\checkmark]\checkmark$\checkmark^\checkmark\\checkmarkc\checkmarki\checkmarkr\checkmarkc\checkmark$\checkmark}\checkmark$\checkmark^\checkmark\\checkmarkc\checkmarki\checkmarkr\checkmarkc\checkmark$\checkmark]\checkmark$\checkmark^\checkmark\\checkmarkc\checkmarki\checkmarkr\checkmarkc\checkmark$\checkmark \checkmark$\checkmark^\checkmark\\checkmarkc\checkmarki\checkmarkr\checkmarkc\checkmark$\checkmark(\checkmark$\checkmark^\checkmark\\checkmarkc\checkmarki\checkmarkr\checkmarkc\checkmark$\checkmarku\checkmark$\checkmark^\checkmark\\checkmarkc\checkmarki\checkmarkr\checkmarkc\checkmark$\checkmarks\checkmark$\checkmark^\checkmark\\checkmarkc\checkmarki\checkmarkr\checkmarkc\checkmark$\checkmarke\checkmark$\checkmark^\checkmark\\checkmarkc\checkmarki\checkmarkr\checkmarkc\checkmark$\checkmarkr\checkmark$\checkmark^\checkmark\\checkmarkc\checkmarki\checkmarkr\checkmarkc\checkmark$\checkmark)\checkmark$\checkmark^\checkmark\\checkmarkc\checkmarki\checkmarkr\checkmarkc\checkmark$\checkmark \checkmark$\checkmark^\checkmark\\checkmarkc\checkmarki\checkmarkr\checkmarkc\checkmark$\checkmark-\checkmark$\checkmark^\checkmark\\checkmarkc\checkmarki\checkmarkr\checkmarkc\checkmark$\checkmark-\checkmark$\checkmark^\checkmark\\checkmarkc\checkmarki\checkmarkr\checkmarkc\checkmark$\checkmark \checkmark$\checkmark^\checkmark\\checkmarkc\checkmarki\checkmarkr\checkmarkc\checkmark$\checkmark(\checkmark$\checkmark^\checkmark\\checkmarkc\checkmarki\checkmarkr\checkmarkc\checkmark$\checkmarks\checkmark$\checkmark^\checkmark\\checkmarkc\checkmarki\checkmarkr\checkmarkc\checkmark$\checkmarky\checkmark$\checkmark^\checkmark\\checkmarkc\checkmarki\checkmarkr\checkmarkc\checkmark$\checkmarks\checkmark$\checkmark^\checkmark\\checkmarkc\checkmarki\checkmarkr\checkmarkc\checkmark$\checkmark)\checkmark$\checkmark^\checkmark\\checkmarkc\checkmarki\checkmarkr\checkmarkc\checkmark$\checkmark \checkmark$\checkmark^\checkmark\\checkmarkc\checkmarki\checkmarkr\checkmarkc\checkmark$\checkmarkn\checkmark$\checkmark^\checkmark\\checkmarkc\checkmarki\checkmarkr\checkmarkc\checkmark$\checkmarko\checkmark$\checkmark^\checkmark\\checkmarkc\checkmarki\checkmarkr\checkmarkc\checkmark$\checkmarkd\checkmark$\checkmark^\checkmark\\checkmarkc\checkmarki\checkmarkr\checkmarkc\checkmark$\checkmarke\checkmark$\checkmark^\checkmark\\checkmarkc\checkmarki\checkmarkr\checkmarkc\checkmark$\checkmark[\checkmark$\checkmark^\checkmark\\checkmarkc\checkmarki\checkmarkr\checkmarkc\checkmark$\checkmarkm\checkmark$\checkmark^\checkmark\\checkmarkc\checkmarki\checkmarkr\checkmarkc\checkmark$\checkmarki\checkmark$\checkmark^\checkmark\\checkmarkc\checkmarki\checkmarkr\checkmarkc\checkmark$\checkmarkd\checkmark$\checkmark^\checkmark\\checkmarkc\checkmarki\checkmarkr\checkmarkc\checkmark$\checkmarkw\checkmark$\checkmark^\checkmark\\checkmarkc\checkmarki\checkmarkr\checkmarkc\checkmark$\checkmarka\checkmark$\checkmark^\checkmark\\checkmarkc\checkmarki\checkmarkr\checkmarkc\checkmark$\checkmarky\checkmark$\checkmark^\checkmark\\checkmarkc\checkmarki\checkmarkr\checkmarkc\checkmark$\checkmark,\checkmark$\checkmark^\checkmark\\checkmarkc\checkmarki\checkmarkr\checkmarkc\checkmark$\checkmarkl\checkmark$\checkmark^\checkmark\\checkmarkc\checkmarki\checkmarkr\checkmarkc\checkmark$\checkmarke\checkmark$\checkmark^\checkmark\\checkmarkc\checkmarki\checkmarkr\checkmarkc\checkmark$\checkmarkf\checkmark$\checkmark^\checkmark\\checkmarkc\checkmarki\checkmarkr\checkmarkc\checkmark$\checkmarkt\checkmark$\checkmark^\checkmark\\checkmarkc\checkmarki\checkmarkr\checkmarkc\checkmark$\checkmark]\checkmark$\checkmark^\checkmark\\checkmarkc\checkmarki\checkmarkr\checkmarkc\checkmark$\checkmark \checkmark$\checkmark^\checkmark\\checkmarkc\checkmarki\checkmarkr\checkmarkc\checkmark$\checkmark{\checkmark$\checkmark^\checkmark\\checkmarkc\checkmarki\checkmarkr\checkmarkc\checkmark$\checkmark\\checkmark$\checkmark^\checkmark\\checkmarkc\checkmarki\checkmarkr\checkmarkc\checkmark$\checkmarks\checkmark$\checkmark^\checkmark\\checkmarkc\checkmarki\checkmarkr\checkmarkc\checkmark$\checkmarkm\checkmark$\checkmark^\checkmark\\checkmarkc\checkmarki\checkmarkr\checkmarkc\checkmark$\checkmarka\checkmark$\checkmark^\checkmark\\checkmarkc\checkmarki\checkmarkr\checkmarkc\checkmark$\checkmarkl\checkmark$\checkmark^\checkmark\\checkmarkc\checkmarki\checkmarkr\checkmarkc\checkmark$\checkmarkl\checkmark$\checkmark^\checkmark\\checkmarkc\checkmarki\checkmarkr\checkmarkc\checkmark$\checkmark \checkmark$\checkmark^\checkmark\\checkmarkc\checkmarki\checkmarkr\checkmarkc\checkmark$\checkmarkÀ\checkmark$\checkmark^\checkmark\\checkmarkc\checkmarki\checkmarkr\checkmarkc\checkmark$\checkmark \checkmark$\checkmark^\checkmark\\checkmarkc\checkmarki\checkmarkr\checkmarkc\checkmark$\checkmarkq\checkmark$\checkmark^\checkmark\\checkmarkc\checkmarki\checkmarkr\checkmarkc\checkmark$\checkmarku\checkmark$\checkmark^\checkmark\\checkmarkc\checkmarki\checkmarkr\checkmarkc\checkmark$\checkmarki\checkmark$\checkmark^\checkmark\\checkmarkc\checkmarki\checkmarkr\checkmarkc\checkmark$\checkmark \checkmark$\checkmark^\checkmark\\checkmarkc\checkmarki\checkmarkr\checkmarkc\checkmark$\checkmark?\checkmark$\checkmark^\checkmark\\checkmarkc\checkmarki\checkmarkr\checkmarkc\checkmark$\checkmark}\checkmark$\checkmark^\checkmark\\checkmarkc\checkmarki\checkmarkr\checkmarkc\checkmark$\checkmark;\checkmark$\checkmark^\checkmark\\checkmarkc\checkmarki\checkmarkr\checkmarkc\checkmark$\checkmark
\checkmark$\checkmark^\checkmark\\checkmarkc\checkmarki\checkmarkr\checkmarkc\checkmark$\checkmark \checkmark$\checkmark^\checkmark\\checkmarkc\checkmarki\checkmarkr\checkmarkc\checkmark$\checkmark \checkmark$\checkmark^\checkmark\\checkmarkc\checkmarki\checkmarkr\checkmarkc\checkmark$\checkmark \checkmark$\checkmark^\checkmark\\checkmarkc\checkmarki\checkmarkr\checkmarkc\checkmark$\checkmark \checkmark$\checkmark^\checkmark\\checkmarkc\checkmarki\checkmarkr\checkmarkc\checkmark$\checkmark \checkmark$\checkmark^\checkmark\\checkmarkc\checkmarki\checkmarkr\checkmarkc\checkmark$\checkmark \checkmark$\checkmark^\checkmark\\checkmarkc\checkmarki\checkmarkr\checkmarkc\checkmark$\checkmark \checkmark$\checkmark^\checkmark\\checkmarkc\checkmarki\checkmarkr\checkmarkc\checkmark$\checkmark \checkmark$\checkmark^\checkmark\\checkmarkc\checkmarki\checkmarkr\checkmarkc\checkmark$\checkmark\\checkmark$\checkmark^\checkmark\\checkmarkc\checkmarki\checkmarkr\checkmarkc\checkmark$\checkmarkd\checkmark$\checkmark^\checkmark\\checkmarkc\checkmarki\checkmarkr\checkmarkc\checkmark$\checkmarkr\checkmark$\checkmark^\checkmark\\checkmarkc\checkmarki\checkmarkr\checkmarkc\checkmark$\checkmarka\checkmark$\checkmark^\checkmark\\checkmarkc\checkmarki\checkmarkr\checkmarkc\checkmark$\checkmarkw\checkmark$\checkmark^\checkmark\\checkmarkc\checkmarki\checkmarkr\checkmarkc\checkmark$\checkmark[\checkmark$\checkmark^\checkmark\\checkmarkc\checkmarki\checkmarkr\checkmarkc\checkmark$\checkmarkt\checkmark$\checkmark^\checkmark\\checkmarkc\checkmarki\checkmarkr\checkmarkc\checkmark$\checkmarkh\checkmark$\checkmark^\checkmark\\checkmarkc\checkmarki\checkmarkr\checkmarkc\checkmark$\checkmarki\checkmark$\checkmark^\checkmark\\checkmarkc\checkmarki\checkmarkr\checkmarkc\checkmark$\checkmarkc\checkmark$\checkmark^\checkmark\\checkmarkc\checkmarki\checkmarkr\checkmarkc\checkmark$\checkmarkk\checkmark$\checkmark^\checkmark\\checkmarkc\checkmarki\checkmarkr\checkmarkc\checkmark$\checkmark,\checkmark$\checkmark^\checkmark\\checkmarkc\checkmarki\checkmarkr\checkmarkc\checkmark$\checkmark \checkmark$\checkmark^\checkmark\\checkmarkc\checkmarki\checkmarkr\checkmarkc\checkmark$\checkmark-\checkmark$\checkmark^\checkmark\\checkmarkc\checkmarki\checkmarkr\checkmarkc\checkmark$\checkmark{\checkmark$\checkmark^\checkmark\\checkmarkc\checkmarki\checkmarkr\checkmarkc\checkmark$\checkmarkS\checkmark$\checkmark^\checkmark\\checkmarkc\checkmarki\checkmarkr\checkmarkc\checkmark$\checkmarkt\checkmark$\checkmark^\checkmark\\checkmarkc\checkmarki\checkmarkr\checkmarkc\checkmark$\checkmarke\checkmark$\checkmark^\checkmark\\checkmarkc\checkmarki\checkmarkr\checkmarkc\checkmark$\checkmarka\checkmark$\checkmark^\checkmark\\checkmarkc\checkmarki\checkmarkr\checkmarkc\checkmark$\checkmarkl\checkmark$\checkmark^\checkmark\\checkmarkc\checkmarki\checkmarkr\checkmarkc\checkmark$\checkmarkt\checkmark$\checkmark^\checkmark\\checkmarkc\checkmarki\checkmarkr\checkmarkc\checkmark$\checkmarkh\checkmark$\checkmark^\checkmark\\checkmarkc\checkmarki\checkmarkr\checkmarkc\checkmark$\checkmark[\checkmark$\checkmark^\checkmark\\checkmarkc\checkmarki\checkmarkr\checkmarkc\checkmark$\checkmarks\checkmark$\checkmark^\checkmark\\checkmarkc\checkmarki\checkmarkr\checkmarkc\checkmark$\checkmarkc\checkmark$\checkmark^\checkmark\\checkmarkc\checkmarki\checkmarkr\checkmarkc\checkmark$\checkmarka\checkmark$\checkmark^\checkmark\\checkmarkc\checkmarki\checkmarkr\checkmarkc\checkmark$\checkmarkl\checkmark$\checkmark^\checkmark\\checkmarkc\checkmarki\checkmarkr\checkmarkc\checkmark$\checkmarke\checkmark$\checkmark^\checkmark\\checkmarkc\checkmarki\checkmarkr\checkmarkc\checkmark$\checkmark=\checkmark$\checkmark^\checkmark\\checkmarkc\checkmarki\checkmarkr\checkmarkc\checkmark$\checkmark1\checkmark$\checkmark^\checkmark\\checkmarkc\checkmarki\checkmarkr\checkmarkc\checkmark$\checkmark.\checkmark$\checkmark^\checkmark\\checkmarkc\checkmarki\checkmarkr\checkmarkc\checkmark$\checkmark2\checkmark$\checkmark^\checkmark\\checkmarkc\checkmarki\checkmarkr\checkmarkc\checkmark$\checkmark]\checkmark$\checkmark^\checkmark\\checkmarkc\checkmarki\checkmarkr\checkmarkc\checkmark$\checkmark}\checkmark$\checkmark^\checkmark\\checkmarkc\checkmarki\checkmarkr\checkmarkc\checkmark$\checkmark]\checkmark$\checkmark^\checkmark\\checkmarkc\checkmarki\checkmarkr\checkmarkc\checkmark$\checkmark \checkmark$\checkmark^\checkmark\\checkmarkc\checkmarki\checkmarkr\checkmarkc\checkmark$\checkmark(\checkmark$\checkmark^\checkmark\\checkmarkc\checkmarki\checkmarkr\checkmarkc\checkmark$\checkmarkm\checkmark$\checkmark^\checkmark\\checkmarkc\checkmarki\checkmarkr\checkmarkc\checkmark$\checkmarka\checkmark$\checkmark^\checkmark\\checkmarkc\checkmarki\checkmarkr\checkmarkc\checkmark$\checkmarkt\checkmark$\checkmark^\checkmark\\checkmarkc\checkmarki\checkmarkr\checkmarkc\checkmark$\checkmarki\checkmark$\checkmark^\checkmark\\checkmarkc\checkmarki\checkmarkr\checkmarkc\checkmark$\checkmarke\checkmark$\checkmark^\checkmark\\checkmarkc\checkmarki\checkmarkr\checkmarkc\checkmark$\checkmarkr\checkmark$\checkmark^\checkmark\\checkmarkc\checkmarki\checkmarkr\checkmarkc\checkmark$\checkmarke\checkmark$\checkmark^\checkmark\\checkmarkc\checkmarki\checkmarkr\checkmarkc\checkmark$\checkmark)\checkmark$\checkmark^\checkmark\\checkmarkc\checkmarki\checkmarkr\checkmarkc\checkmark$\checkmark \checkmark$\checkmark^\checkmark\\checkmarkc\checkmarki\checkmarkr\checkmarkc\checkmark$\checkmark-\checkmark$\checkmark^\checkmark\\checkmarkc\checkmarki\checkmarkr\checkmarkc\checkmark$\checkmark-\checkmark$\checkmark^\checkmark\\checkmarkc\checkmarki\checkmarkr\checkmarkc\checkmark$\checkmark \checkmark$\checkmark^\checkmark\\checkmarkc\checkmarki\checkmarkr\checkmarkc\checkmark$\checkmark(\checkmark$\checkmark^\checkmark\\checkmarkc\checkmarki\checkmarkr\checkmarkc\checkmark$\checkmarks\checkmark$\checkmark^\checkmark\\checkmarkc\checkmarki\checkmarkr\checkmarkc\checkmark$\checkmarky\checkmark$\checkmark^\checkmark\\checkmarkc\checkmarki\checkmarkr\checkmarkc\checkmark$\checkmarks\checkmark$\checkmark^\checkmark\\checkmarkc\checkmarki\checkmarkr\checkmarkc\checkmark$\checkmark)\checkmark$\checkmark^\checkmark\\checkmarkc\checkmarki\checkmarkr\checkmarkc\checkmark$\checkmark \checkmark$\checkmark^\checkmark\\checkmarkc\checkmarki\checkmarkr\checkmarkc\checkmark$\checkmarkn\checkmark$\checkmark^\checkmark\\checkmarkc\checkmarki\checkmarkr\checkmarkc\checkmark$\checkmarko\checkmark$\checkmark^\checkmark\\checkmarkc\checkmarki\checkmarkr\checkmarkc\checkmark$\checkmarkd\checkmark$\checkmark^\checkmark\\checkmarkc\checkmarki\checkmarkr\checkmarkc\checkmark$\checkmarke\checkmark$\checkmark^\checkmark\\checkmarkc\checkmarki\checkmarkr\checkmarkc\checkmark$\checkmark[\checkmark$\checkmark^\checkmark\\checkmarkc\checkmarki\checkmarkr\checkmarkc\checkmark$\checkmarkm\checkmark$\checkmark^\checkmark\\checkmarkc\checkmarki\checkmarkr\checkmarkc\checkmark$\checkmarki\checkmark$\checkmark^\checkmark\\checkmarkc\checkmarki\checkmarkr\checkmarkc\checkmark$\checkmarkd\checkmark$\checkmark^\checkmark\\checkmarkc\checkmarki\checkmarkr\checkmarkc\checkmark$\checkmarkw\checkmark$\checkmark^\checkmark\\checkmarkc\checkmarki\checkmarkr\checkmarkc\checkmark$\checkmarka\checkmark$\checkmark^\checkmark\\checkmarkc\checkmarki\checkmarkr\checkmarkc\checkmark$\checkmarky\checkmark$\checkmark^\checkmark\\checkmarkc\checkmarki\checkmarkr\checkmarkc\checkmark$\checkmark,\checkmark$\checkmark^\checkmark\\checkmarkc\checkmarki\checkmarkr\checkmarkc\checkmark$\checkmarkr\checkmark$\checkmark^\checkmark\\checkmarkc\checkmarki\checkmarkr\checkmarkc\checkmark$\checkmarki\checkmark$\checkmark^\checkmark\\checkmarkc\checkmarki\checkmarkr\checkmarkc\checkmark$\checkmarkg\checkmark$\checkmark^\checkmark\\checkmarkc\checkmarki\checkmarkr\checkmarkc\checkmark$\checkmarkh\checkmark$\checkmark^\checkmark\\checkmarkc\checkmarki\checkmarkr\checkmarkc\checkmark$\checkmarkt\checkmark$\checkmark^\checkmark\\checkmarkc\checkmarki\checkmarkr\checkmarkc\checkmark$\checkmark]\checkmark$\checkmark^\checkmark\\checkmarkc\checkmarki\checkmarkr\checkmarkc\checkmark$\checkmark \checkmark$\checkmark^\checkmark\\checkmarkc\checkmarki\checkmarkr\checkmarkc\checkmark$\checkmark{\checkmark$\checkmark^\checkmark\\checkmarkc\checkmarki\checkmarkr\checkmarkc\checkmark$\checkmark\\checkmark$\checkmark^\checkmark\\checkmarkc\checkmarki\checkmarkr\checkmarkc\checkmark$\checkmarks\checkmark$\checkmark^\checkmark\\checkmarkc\checkmarki\checkmarkr\checkmarkc\checkmark$\checkmarkm\checkmark$\checkmark^\checkmark\\checkmarkc\checkmarki\checkmarkr\checkmarkc\checkmark$\checkmarka\checkmark$\checkmark^\checkmark\\checkmarkc\checkmarki\checkmarkr\checkmarkc\checkmark$\checkmarkl\checkmark$\checkmark^\checkmark\\checkmarkc\checkmarki\checkmarkr\checkmarkc\checkmark$\checkmarkl\checkmark$\checkmark^\checkmark\\checkmarkc\checkmarki\checkmarkr\checkmarkc\checkmark$\checkmark \checkmark$\checkmark^\checkmark\\checkmarkc\checkmarki\checkmarkr\checkmarkc\checkmark$\checkmarkS\checkmark$\checkmark^\checkmark\\checkmarkc\checkmarki\checkmarkr\checkmarkc\checkmark$\checkmarku\checkmark$\checkmark^\checkmark\\checkmarkc\checkmarki\checkmarkr\checkmarkc\checkmark$\checkmarkr\checkmark$\checkmark^\checkmark\\checkmarkc\checkmarki\checkmarkr\checkmarkc\checkmark$\checkmark \checkmark$\checkmark^\checkmark\\checkmarkc\checkmarki\checkmarkr\checkmarkc\checkmark$\checkmarkq\checkmark$\checkmark^\checkmark\\checkmarkc\checkmarki\checkmarkr\checkmarkc\checkmark$\checkmarku\checkmark$\checkmark^\checkmark\\checkmarkc\checkmarki\checkmarkr\checkmarkc\checkmark$\checkmarko\checkmark$\checkmark^\checkmark\\checkmarkc\checkmarki\checkmarkr\checkmarkc\checkmark$\checkmarki\checkmark$\checkmark^\checkmark\\checkmarkc\checkmarki\checkmarkr\checkmarkc\checkmark$\checkmark \checkmark$\checkmark^\checkmark\\checkmarkc\checkmarki\checkmarkr\checkmarkc\checkmark$\checkmark?\checkmark$\checkmark^\checkmark\\checkmarkc\checkmarki\checkmarkr\checkmarkc\checkmark$\checkmark}\checkmark$\checkmark^\checkmark\\checkmarkc\checkmarki\checkmarkr\checkmarkc\checkmark$\checkmark;\checkmark$\checkmark^\checkmark\\checkmarkc\checkmarki\checkmarkr\checkmarkc\checkmark$\checkmark
\checkmark$\checkmark^\checkmark\\checkmarkc\checkmarki\checkmarkr\checkmarkc\checkmark$\checkmark \checkmark$\checkmark^\checkmark\\checkmarkc\checkmarki\checkmarkr\checkmarkc\checkmark$\checkmark \checkmark$\checkmark^\checkmark\\checkmarkc\checkmarki\checkmarkr\checkmarkc\checkmark$\checkmark \checkmark$\checkmark^\checkmark\\checkmarkc\checkmarki\checkmarkr\checkmarkc\checkmark$\checkmark \checkmark$\checkmark^\checkmark\\checkmarkc\checkmarki\checkmarkr\checkmarkc\checkmark$\checkmark \checkmark$\checkmark^\checkmark\\checkmarkc\checkmarki\checkmarkr\checkmarkc\checkmark$\checkmark \checkmark$\checkmark^\checkmark\\checkmarkc\checkmarki\checkmarkr\checkmarkc\checkmark$\checkmark \checkmark$\checkmark^\checkmark\\checkmarkc\checkmarki\checkmarkr\checkmarkc\checkmark$\checkmark \checkmark$\checkmark^\checkmark\\checkmarkc\checkmarki\checkmarkr\checkmarkc\checkmark$\checkmark\\checkmark$\checkmark^\checkmark\\checkmarkc\checkmarki\checkmarkr\checkmarkc\checkmark$\checkmarkd\checkmark$\checkmark^\checkmark\\checkmarkc\checkmarki\checkmarkr\checkmarkc\checkmark$\checkmarkr\checkmark$\checkmark^\checkmark\\checkmarkc\checkmarki\checkmarkr\checkmarkc\checkmark$\checkmarka\checkmark$\checkmark^\checkmark\\checkmarkc\checkmarki\checkmarkr\checkmarkc\checkmark$\checkmarkw\checkmark$\checkmark^\checkmark\\checkmarkc\checkmarki\checkmarkr\checkmarkc\checkmark$\checkmark[\checkmark$\checkmark^\checkmark\\checkmarkc\checkmarki\checkmarkr\checkmarkc\checkmark$\checkmarkt\checkmark$\checkmark^\checkmark\\checkmarkc\checkmarki\checkmarkr\checkmarkc\checkmark$\checkmarkh\checkmark$\checkmark^\checkmark\\checkmarkc\checkmarki\checkmarkr\checkmarkc\checkmark$\checkmarki\checkmark$\checkmark^\checkmark\\checkmarkc\checkmarki\checkmarkr\checkmarkc\checkmark$\checkmarkc\checkmark$\checkmark^\checkmark\\checkmarkc\checkmarki\checkmarkr\checkmarkc\checkmark$\checkmarkk\checkmark$\checkmark^\checkmark\\checkmarkc\checkmarki\checkmarkr\checkmarkc\checkmark$\checkmark,\checkmark$\checkmark^\checkmark\\checkmarkc\checkmarki\checkmarkr\checkmarkc\checkmark$\checkmark \checkmark$\checkmark^\checkmark\\checkmarkc\checkmarki\checkmarkr\checkmarkc\checkmark$\checkmark-\checkmark$\checkmark^\checkmark\\checkmarkc\checkmarki\checkmarkr\checkmarkc\checkmark$\checkmark{\checkmark$\checkmark^\checkmark\\checkmarkc\checkmarki\checkmarkr\checkmarkc\checkmark$\checkmarkS\checkmark$\checkmark^\checkmark\\checkmarkc\checkmarki\checkmarkr\checkmarkc\checkmark$\checkmarkt\checkmark$\checkmark^\checkmark\\checkmarkc\checkmarki\checkmarkr\checkmarkc\checkmark$\checkmarke\checkmark$\checkmark^\checkmark\\checkmarkc\checkmarki\checkmarkr\checkmarkc\checkmark$\checkmarka\checkmark$\checkmark^\checkmark\\checkmarkc\checkmarki\checkmarkr\checkmarkc\checkmark$\checkmarkl\checkmark$\checkmark^\checkmark\\checkmarkc\checkmarki\checkmarkr\checkmarkc\checkmark$\checkmarkt\checkmark$\checkmark^\checkmark\\checkmarkc\checkmarki\checkmarkr\checkmarkc\checkmark$\checkmarkh\checkmark$\checkmark^\checkmark\\checkmarkc\checkmarki\checkmarkr\checkmarkc\checkmark$\checkmark[\checkmark$\checkmark^\checkmark\\checkmarkc\checkmarki\checkmarkr\checkmarkc\checkmark$\checkmarks\checkmark$\checkmark^\checkmark\\checkmarkc\checkmarki\checkmarkr\checkmarkc\checkmark$\checkmarkc\checkmark$\checkmark^\checkmark\\checkmarkc\checkmarki\checkmarkr\checkmarkc\checkmark$\checkmarka\checkmark$\checkmark^\checkmark\\checkmarkc\checkmarki\checkmarkr\checkmarkc\checkmark$\checkmarkl\checkmark$\checkmark^\checkmark\\checkmarkc\checkmarki\checkmarkr\checkmarkc\checkmark$\checkmarke\checkmark$\checkmark^\checkmark\\checkmarkc\checkmarki\checkmarkr\checkmarkc\checkmark$\checkmark=\checkmark$\checkmark^\checkmark\\checkmarkc\checkmarki\checkmarkr\checkmarkc\checkmark$\checkmark1\checkmark$\checkmark^\checkmark\\checkmarkc\checkmarki\checkmarkr\checkmarkc\checkmark$\checkmark.\checkmark$\checkmark^\checkmark\\checkmarkc\checkmarki\checkmarkr\checkmarkc\checkmark$\checkmark2\checkmark$\checkmark^\checkmark\\checkmarkc\checkmarki\checkmarkr\checkmarkc\checkmark$\checkmark]\checkmark$\checkmark^\checkmark\\checkmarkc\checkmarki\checkmarkr\checkmarkc\checkmark$\checkmark}\checkmark$\checkmark^\checkmark\\checkmarkc\checkmarki\checkmarkr\checkmarkc\checkmark$\checkmark,\checkmark$\checkmark^\checkmark\\checkmarkc\checkmarki\checkmarkr\checkmarkc\checkmark$\checkmark \checkmark$\checkmark^\checkmark\\checkmarkc\checkmarki\checkmarkr\checkmarkc\checkmark$\checkmarkr\checkmark$\checkmark^\checkmark\\checkmarkc\checkmarki\checkmarkr\checkmarkc\checkmark$\checkmarke\checkmark$\checkmark^\checkmark\\checkmarkc\checkmarki\checkmarkr\checkmarkc\checkmark$\checkmarkd\checkmark$\checkmark^\checkmark\\checkmarkc\checkmarki\checkmarkr\checkmarkc\checkmark$\checkmark]\checkmark$\checkmark^\checkmark\\checkmarkc\checkmarki\checkmarkr\checkmarkc\checkmark$\checkmark \checkmark$\checkmark^\checkmark\\checkmarkc\checkmarki\checkmarkr\checkmarkc\checkmark$\checkmark(\checkmark$\checkmark^\checkmark\\checkmarkc\checkmarki\checkmarkr\checkmarkc\checkmark$\checkmarks\checkmark$\checkmark^\checkmark\\checkmarkc\checkmarki\checkmarkr\checkmarkc\checkmark$\checkmarky\checkmark$\checkmark^\checkmark\\checkmarkc\checkmarki\checkmarkr\checkmarkc\checkmark$\checkmarks\checkmark$\checkmark^\checkmark\\checkmarkc\checkmarki\checkmarkr\checkmarkc\checkmark$\checkmark.\checkmark$\checkmark^\checkmark\\checkmarkc\checkmarki\checkmarkr\checkmarkc\checkmark$\checkmarks\checkmark$\checkmark^\checkmark\\checkmarkc\checkmarki\checkmarkr\checkmarkc\checkmark$\checkmarko\checkmark$\checkmark^\checkmark\\checkmarkc\checkmarki\checkmarkr\checkmarkc\checkmark$\checkmarku\checkmark$\checkmark^\checkmark\\checkmarkc\checkmarki\checkmarkr\checkmarkc\checkmark$\checkmarkt\checkmark$\checkmark^\checkmark\\checkmarkc\checkmarki\checkmarkr\checkmarkc\checkmark$\checkmarkh\checkmark$\checkmark^\checkmark\\checkmarkc\checkmarki\checkmarkr\checkmarkc\checkmark$\checkmark)\checkmark$\checkmark^\checkmark\\checkmarkc\checkmarki\checkmarkr\checkmarkc\checkmark$\checkmark \checkmark$\checkmark^\checkmark\\checkmarkc\checkmarki\checkmarkr\checkmarkc\checkmark$\checkmark-\checkmark$\checkmark^\checkmark\\checkmarkc\checkmarki\checkmarkr\checkmarkc\checkmark$\checkmark-\checkmark$\checkmark^\checkmark\\checkmarkc\checkmarki\checkmarkr\checkmarkc\checkmark$\checkmark \checkmark$\checkmark^\checkmark\\checkmarkc\checkmarki\checkmarkr\checkmarkc\checkmark$\checkmark+\checkmark$\checkmark^\checkmark\\checkmarkc\checkmarki\checkmarkr\checkmarkc\checkmark$\checkmark+\checkmark$\checkmark^\checkmark\\checkmarkc\checkmarki\checkmarkr\checkmarkc\checkmark$\checkmark(\checkmark$\checkmark^\checkmark\\checkmarkc\checkmarki\checkmarkr\checkmarkc\checkmark$\checkmark0\checkmark$\checkmark^\checkmark\\checkmarkc\checkmarki\checkmarkr\checkmarkc\checkmark$\checkmark,\checkmark$\checkmark^\checkmark\\checkmarkc\checkmarki\checkmarkr\checkmarkc\checkmark$\checkmark-\checkmark$\checkmark^\checkmark\\checkmarkc\checkmarki\checkmarkr\checkmarkc\checkmark$\checkmark1\checkmark$\checkmark^\checkmark\\checkmarkc\checkmarki\checkmarkr\checkmarkc\checkmark$\checkmark.\checkmark$\checkmark^\checkmark\\checkmarkc\checkmarki\checkmarkr\checkmarkc\checkmark$\checkmark2\checkmark$\checkmark^\checkmark\\checkmarkc\checkmarki\checkmarkr\checkmarkc\checkmark$\checkmark)\checkmark$\checkmark^\checkmark\\checkmarkc\checkmarki\checkmarkr\checkmarkc\checkmark$\checkmark \checkmark$\checkmark^\checkmark\\checkmarkc\checkmarki\checkmarkr\checkmarkc\checkmark$\checkmarkn\checkmark$\checkmark^\checkmark\\checkmarkc\checkmarki\checkmarkr\checkmarkc\checkmark$\checkmarko\checkmark$\checkmark^\checkmark\\checkmarkc\checkmarki\checkmarkr\checkmarkc\checkmark$\checkmarkd\checkmark$\checkmark^\checkmark\\checkmarkc\checkmarki\checkmarkr\checkmarkc\checkmark$\checkmarke\checkmark$\checkmark^\checkmark\\checkmarkc\checkmarki\checkmarkr\checkmarkc\checkmark$\checkmark[\checkmark$\checkmark^\checkmark\\checkmarkc\checkmarki\checkmarkr\checkmarkc\checkmark$\checkmarkb\checkmark$\checkmark^\checkmark\\checkmarkc\checkmarki\checkmarkr\checkmarkc\checkmark$\checkmarke\checkmark$\checkmark^\checkmark\\checkmarkc\checkmarki\checkmarkr\checkmarkc\checkmark$\checkmarkl\checkmark$\checkmark^\checkmark\\checkmarkc\checkmarki\checkmarkr\checkmarkc\checkmark$\checkmarko\checkmark$\checkmark^\checkmark\\checkmarkc\checkmarki\checkmarkr\checkmarkc\checkmark$\checkmarkw\checkmark$\checkmark^\checkmark\\checkmarkc\checkmarki\checkmarkr\checkmarkc\checkmark$\checkmark,\checkmark$\checkmark^\checkmark\\checkmarkc\checkmarki\checkmarkr\checkmarkc\checkmark$\checkmark \checkmark$\checkmark^\checkmark\\checkmarkc\checkmarki\checkmarkr\checkmarkc\checkmark$\checkmarka\checkmark$\checkmark^\checkmark\\checkmarkc\checkmarki\checkmarkr\checkmarkc\checkmark$\checkmarkl\checkmark$\checkmark^\checkmark\\checkmarkc\checkmarki\checkmarkr\checkmarkc\checkmark$\checkmarki\checkmark$\checkmark^\checkmark\\checkmarkc\checkmarki\checkmarkr\checkmarkc\checkmark$\checkmarkg\checkmark$\checkmark^\checkmark\\checkmarkc\checkmarki\checkmarkr\checkmarkc\checkmark$\checkmarkn\checkmark$\checkmark^\checkmark\\checkmarkc\checkmarki\checkmarkr\checkmarkc\checkmark$\checkmark=\checkmark$\checkmark^\checkmark\\checkmarkc\checkmarki\checkmarkr\checkmarkc\checkmark$\checkmarkc\checkmark$\checkmark^\checkmark\\checkmarkc\checkmarki\checkmarkr\checkmarkc\checkmark$\checkmarke\checkmark$\checkmark^\checkmark\\checkmarkc\checkmarki\checkmarkr\checkmarkc\checkmark$\checkmarkn\checkmark$\checkmark^\checkmark\\checkmarkc\checkmarki\checkmarkr\checkmarkc\checkmark$\checkmarkt\checkmark$\checkmark^\checkmark\\checkmarkc\checkmarki\checkmarkr\checkmarkc\checkmark$\checkmarke\checkmark$\checkmark^\checkmark\\checkmarkc\checkmarki\checkmarkr\checkmarkc\checkmark$\checkmarkr\checkmark$\checkmark^\checkmark\\checkmarkc\checkmarki\checkmarkr\checkmarkc\checkmark$\checkmark,\checkmark$\checkmark^\checkmark\\checkmarkc\checkmarki\checkmarkr\checkmarkc\checkmark$\checkmark \checkmark$\checkmark^\checkmark\\checkmarkc\checkmarki\checkmarkr\checkmarkc\checkmark$\checkmarkt\checkmark$\checkmark^\checkmark\\checkmarkc\checkmarki\checkmarkr\checkmarkc\checkmark$\checkmarke\checkmark$\checkmark^\checkmark\\checkmarkc\checkmarki\checkmarkr\checkmarkc\checkmark$\checkmarkx\checkmark$\checkmark^\checkmark\\checkmarkc\checkmarki\checkmarkr\checkmarkc\checkmark$\checkmarkt\checkmark$\checkmark^\checkmark\\checkmarkc\checkmarki\checkmarkr\checkmarkc\checkmark$\checkmark \checkmark$\checkmark^\checkmark\\checkmarkc\checkmarki\checkmarkr\checkmarkc\checkmark$\checkmarkw\checkmark$\checkmark^\checkmark\\checkmarkc\checkmarki\checkmarkr\checkmarkc\checkmark$\checkmarki\checkmark$\checkmark^\checkmark\\checkmarkc\checkmarki\checkmarkr\checkmarkc\checkmark$\checkmarkd\checkmark$\checkmark^\checkmark\\checkmarkc\checkmarki\checkmarkr\checkmarkc\checkmark$\checkmarkt\checkmark$\checkmark^\checkmark\\checkmarkc\checkmarki\checkmarkr\checkmarkc\checkmark$\checkmarkh\checkmark$\checkmark^\checkmark\\checkmarkc\checkmarki\checkmarkr\checkmarkc\checkmark$\checkmark=\checkmark$\checkmark^\checkmark\\checkmarkc\checkmarki\checkmarkr\checkmarkc\checkmark$\checkmark9\checkmark$\checkmark^\checkmark\\checkmarkc\checkmarki\checkmarkr\checkmarkc\checkmark$\checkmarkc\checkmark$\checkmark^\checkmark\\checkmarkc\checkmarki\checkmarkr\checkmarkc\checkmark$\checkmarkm\checkmark$\checkmark^\checkmark\\checkmarkc\checkmarki\checkmarkr\checkmarkc\checkmark$\checkmark]\checkmark$\checkmark^\checkmark\\checkmarkc\checkmarki\checkmarkr\checkmarkc\checkmark$\checkmark \checkmark$\checkmark^\checkmark\\checkmarkc\checkmarki\checkmarkr\checkmarkc\checkmark$\checkmark{\checkmark$\checkmark^\checkmark\\checkmarkc\checkmarki\checkmarkr\checkmarkc\checkmark$\checkmark\\checkmark$\checkmark^\checkmark\\checkmarkc\checkmarki\checkmarkr\checkmarkc\checkmark$\checkmarkt\checkmark$\checkmark^\checkmark\\checkmarkc\checkmarki\checkmarkr\checkmarkc\checkmark$\checkmarke\checkmark$\checkmark^\checkmark\\checkmarkc\checkmarki\checkmarkr\checkmarkc\checkmark$\checkmarkx\checkmark$\checkmark^\checkmark\\checkmarkc\checkmarki\checkmarkr\checkmarkc\checkmark$\checkmarkt\checkmark$\checkmark^\checkmark\\checkmarkc\checkmarki\checkmarkr\checkmarkc\checkmark$\checkmarki\checkmark$\checkmark^\checkmark\\checkmarkc\checkmarki\checkmarkr\checkmarkc\checkmark$\checkmarkt\checkmark$\checkmark^\checkmark\\checkmarkc\checkmarki\checkmarkr\checkmarkc\checkmark$\checkmark{\checkmark$\checkmark^\checkmark\\checkmarkc\checkmarki\checkmarkr\checkmarkc\checkmark$\checkmarkP\checkmark$\checkmark^\checkmark\\checkmarkc\checkmarki\checkmarkr\checkmarkc\checkmark$\checkmarke\checkmark$\checkmark^\checkmark\\checkmarkc\checkmarki\checkmarkr\checkmarkc\checkmark$\checkmarkr\checkmark$\checkmark^\checkmark\\checkmarkc\checkmarki\checkmarkr\checkmarkc\checkmark$\checkmarkm\checkmark$\checkmark^\checkmark\\checkmarkc\checkmarki\checkmarkr\checkmarkc\checkmark$\checkmarke\checkmark$\checkmark^\checkmark\\checkmarkc\checkmarki\checkmarkr\checkmarkc\checkmark$\checkmarkt\checkmark$\checkmark^\checkmark\\checkmarkc\checkmarki\checkmarkr\checkmarkc\checkmark$\checkmarkt\checkmark$\checkmark^\checkmark\\checkmarkc\checkmarki\checkmarkr\checkmarkc\checkmark$\checkmarkr\checkmark$\checkmark^\checkmark\\checkmarkc\checkmarki\checkmarkr\checkmarkc\checkmark$\checkmarke\checkmark$\checkmark^\checkmark\\checkmarkc\checkmarki\checkmarkr\checkmarkc\checkmark$\checkmark \checkmark$\checkmark^\checkmark\\checkmarkc\checkmarki\checkmarkr\checkmarkc\checkmark$\checkmarkl\checkmark$\checkmark^\checkmark\\checkmarkc\checkmarki\checkmarkr\checkmarkc\checkmark$\checkmark'\checkmark$\checkmark^\checkmark\\checkmarkc\checkmarki\checkmarkr\checkmarkc\checkmark$\checkmarku\checkmark$\checkmark^\checkmark\\checkmarkc\checkmarki\checkmarkr\checkmarkc\checkmark$\checkmarks\checkmark$\checkmark^\checkmark\\checkmarkc\checkmarki\checkmarkr\checkmarkc\checkmark$\checkmarki\checkmark$\checkmark^\checkmark\\checkmarkc\checkmarki\checkmarkr\checkmarkc\checkmark$\checkmarkn\checkmark$\checkmark^\checkmark\\checkmarkc\checkmarki\checkmarkr\checkmarkc\checkmark$\checkmarka\checkmark$\checkmark^\checkmark\\checkmarkc\checkmarki\checkmarkr\checkmarkc\checkmark$\checkmarkg\checkmark$\checkmark^\checkmark\\checkmarkc\checkmarki\checkmarkr\checkmarkc\checkmark$\checkmarke\checkmark$\checkmark^\checkmark\\checkmarkc\checkmarki\checkmarkr\checkmarkc\checkmark$\checkmark \checkmark$\checkmark^\checkmark\\checkmarkc\checkmarki\checkmarkr\checkmarkc\checkmark$\checkmarkm\checkmark$\checkmark^\checkmark\\checkmarkc\checkmarki\checkmarkr\checkmarkc\checkmark$\checkmarku\checkmark$\checkmark^\checkmark\\checkmarkc\checkmarki\checkmarkr\checkmarkc\checkmark$\checkmarkl\checkmark$\checkmark^\checkmark\\checkmarkc\checkmarki\checkmarkr\checkmarkc\checkmark$\checkmarkt\checkmark$\checkmark^\checkmark\\checkmarkc\checkmarki\checkmarkr\checkmarkc\checkmark$\checkmarki\checkmark$\checkmark^\checkmark\\checkmarkc\checkmarki\checkmarkr\checkmarkc\checkmark$\checkmarka\checkmark$\checkmark^\checkmark\\checkmarkc\checkmarki\checkmarkr\checkmarkc\checkmark$\checkmarkx\checkmark$\checkmark^\checkmark\\checkmarkc\checkmarki\checkmarkr\checkmarkc\checkmark$\checkmarke\checkmark$\checkmark^\checkmark\\checkmarkc\checkmarki\checkmarkr\checkmarkc\checkmark$\checkmarks\checkmark$\checkmark^\checkmark\\checkmarkc\checkmarki\checkmarkr\checkmarkc\checkmark$\checkmark \checkmark$\checkmark^\checkmark\\checkmarkc\checkmarki\checkmarkr\checkmarkc\checkmark$\checkmarka\checkmark$\checkmark^\checkmark\\checkmarkc\checkmarki\checkmarkr\checkmarkc\checkmark$\checkmarku\checkmark$\checkmark^\checkmark\\checkmarkc\checkmarki\checkmarkr\checkmarkc\checkmark$\checkmarkt\checkmark$\checkmark^\checkmark\\checkmarkc\checkmarki\checkmarkr\checkmarkc\checkmark$\checkmarko\checkmark$\checkmark^\checkmark\\checkmarkc\checkmarki\checkmarkr\checkmarkc\checkmark$\checkmarkm\checkmark$\checkmark^\checkmark\\checkmarkc\checkmarki\checkmarkr\checkmarkc\checkmark$\checkmarka\checkmark$\checkmark^\checkmark\\checkmarkc\checkmarki\checkmarkr\checkmarkc\checkmark$\checkmarkt\checkmark$\checkmark^\checkmark\\checkmarkc\checkmarki\checkmarkr\checkmarkc\checkmark$\checkmarki\checkmark$\checkmark^\checkmark\\checkmarkc\checkmarki\checkmarkr\checkmarkc\checkmark$\checkmarkq\checkmark$\checkmark^\checkmark\\checkmarkc\checkmarki\checkmarkr\checkmarkc\checkmark$\checkmarku\checkmark$\checkmark^\checkmark\\checkmarkc\checkmarki\checkmarkr\checkmarkc\checkmark$\checkmarke\checkmark$\checkmark^\checkmark\\checkmarkc\checkmarki\checkmarkr\checkmarkc\checkmark$\checkmark \checkmark$\checkmark^\checkmark\\checkmarkc\checkmarki\checkmarkr\checkmarkc\checkmark$\checkmarke\checkmark$\checkmark^\checkmark\\checkmarkc\checkmarki\checkmarkr\checkmarkc\checkmark$\checkmarkt\checkmark$\checkmark^\checkmark\\checkmarkc\checkmarki\checkmarkr\checkmarkc\checkmark$\checkmark \checkmark$\checkmark^\checkmark\\checkmarkc\checkmarki\checkmarkr\checkmarkc\checkmark$\checkmarkp\checkmark$\checkmark^\checkmark\\checkmarkc\checkmarki\checkmarkr\checkmarkc\checkmark$\checkmarkr\checkmark$\checkmark^\checkmark\\checkmarkc\checkmarki\checkmarkr\checkmarkc\checkmark$\checkmarké\checkmark$\checkmark^\checkmark\\checkmarkc\checkmarki\checkmarkr\checkmarkc\checkmark$\checkmarkc\checkmark$\checkmark^\checkmark\\checkmarkc\checkmarki\checkmarkr\checkmarkc\checkmark$\checkmarki\checkmark$\checkmark^\checkmark\\checkmarkc\checkmarki\checkmarkr\checkmarkc\checkmark$\checkmarks\checkmark$\checkmark^\checkmark\\checkmarkc\checkmarki\checkmarkr\checkmarkc\checkmark$\checkmark \checkmark$\checkmark^\checkmark\\checkmarkc\checkmarki\checkmarkr\checkmarkc\checkmark$\checkmarkd\checkmark$\checkmark^\checkmark\\checkmarkc\checkmarki\checkmarkr\checkmarkc\checkmark$\checkmarke\checkmark$\checkmark^\checkmark\\checkmarkc\checkmarki\checkmarkr\checkmarkc\checkmark$\checkmark \checkmark$\checkmark^\checkmark\\checkmarkc\checkmarki\checkmarkr\checkmarkc\checkmark$\checkmarkp\checkmark$\checkmark^\checkmark\\checkmarkc\checkmarki\checkmarkr\checkmarkc\checkmark$\checkmarki\checkmark$\checkmark^\checkmark\\checkmarkc\checkmarki\checkmarkr\checkmarkc\checkmark$\checkmarkè\checkmark$\checkmark^\checkmark\\checkmarkc\checkmarki\checkmarkr\checkmarkc\checkmark$\checkmarkc\checkmark$\checkmark^\checkmark\\checkmarkc\checkmarki\checkmarkr\checkmarkc\checkmark$\checkmarke\checkmark$\checkmark^\checkmark\\checkmarkc\checkmarki\checkmarkr\checkmarkc\checkmark$\checkmarks\checkmark$\checkmark^\checkmark\\checkmarkc\checkmarki\checkmarkr\checkmarkc\checkmark$\checkmark \checkmark$\checkmark^\checkmark\\checkmarkc\checkmarki\checkmarkr\checkmarkc\checkmark$\checkmarkc\checkmark$\checkmark^\checkmark\\checkmarkc\checkmarki\checkmarkr\checkmarkc\checkmark$\checkmarko\checkmark$\checkmark^\checkmark\\checkmarkc\checkmarki\checkmarkr\checkmarkc\checkmark$\checkmarkm\checkmark$\checkmark^\checkmark\\checkmarkc\checkmarki\checkmarkr\checkmarkc\checkmark$\checkmarkp\checkmark$\checkmark^\checkmark\\checkmarkc\checkmarki\checkmarkr\checkmarkc\checkmark$\checkmarkl\checkmark$\checkmark^\checkmark\\checkmarkc\checkmarki\checkmarkr\checkmarkc\checkmark$\checkmarke\checkmark$\checkmark^\checkmark\\checkmarkc\checkmarki\checkmarkr\checkmarkc\checkmark$\checkmarkx\checkmark$\checkmark^\checkmark\\checkmarkc\checkmarki\checkmarkr\checkmarkc\checkmark$\checkmarke\checkmark$\checkmark^\checkmark\\checkmarkc\checkmarki\checkmarkr\checkmarkc\checkmark$\checkmarks\checkmark$\checkmark^\checkmark\\checkmarkc\checkmarki\checkmarkr\checkmarkc\checkmark$\checkmark \checkmark$\checkmark^\checkmark\\checkmarkc\checkmarki\checkmarkr\checkmarkc\checkmark$\checkmarks\checkmark$\checkmark^\checkmark\\checkmarkc\checkmarki\checkmarkr\checkmarkc\checkmark$\checkmarka\checkmark$\checkmark^\checkmark\\checkmarkc\checkmarki\checkmarkr\checkmarkc\checkmark$\checkmarkn\checkmark$\checkmark^\checkmark\\checkmarkc\checkmarki\checkmarkr\checkmarkc\checkmark$\checkmarks\checkmark$\checkmark^\checkmark\\checkmarkc\checkmarki\checkmarkr\checkmarkc\checkmark$\checkmark \checkmark$\checkmark^\checkmark\\checkmarkc\checkmarki\checkmarkr\checkmarkc\checkmark$\checkmarkd\checkmark$\checkmark^\checkmark\\checkmarkc\checkmarki\checkmarkr\checkmarkc\checkmark$\checkmarké\checkmark$\checkmark^\checkmark\\checkmarkc\checkmarki\checkmarkr\checkmarkc\checkmark$\checkmarkm\checkmark$\checkmark^\checkmark\\checkmarkc\checkmarki\checkmarkr\checkmarkc\checkmark$\checkmarko\checkmark$\checkmark^\checkmark\\checkmarkc\checkmarki\checkmarkr\checkmarkc\checkmark$\checkmarkn\checkmark$\checkmark^\checkmark\\checkmarkc\checkmarki\checkmarkr\checkmarkc\checkmark$\checkmarkt\checkmark$\checkmark^\checkmark\\checkmarkc\checkmarki\checkmarkr\checkmarkc\checkmark$\checkmarka\checkmark$\checkmark^\checkmark\\checkmarkc\checkmarki\checkmarkr\checkmarkc\checkmark$\checkmarkg\checkmark$\checkmark^\checkmark\\checkmarkc\checkmarki\checkmarkr\checkmarkc\checkmark$\checkmarke\checkmark$\checkmark^\checkmark\\checkmarkc\checkmarki\checkmarkr\checkmarkc\checkmark$\checkmark.\checkmark$\checkmark^\checkmark\\checkmarkc\checkmarki\checkmarkr\checkmarkc\checkmark$\checkmark}\checkmark$\checkmark^\checkmark\\checkmarkc\checkmarki\checkmarkr\checkmarkc\checkmark$\checkmark}\checkmark$\checkmark^\checkmark\\checkmarkc\checkmarki\checkmarkr\checkmarkc\checkmark$\checkmark;\checkmark$\checkmark^\checkmark\\checkmarkc\checkmarki\checkmarkr\checkmarkc\checkmark$\checkmark
\checkmark$\checkmark^\checkmark\\checkmarkc\checkmarki\checkmarkr\checkmarkc\checkmark$\checkmark \checkmark$\checkmark^\checkmark\\checkmarkc\checkmarki\checkmarkr\checkmarkc\checkmark$\checkmark \checkmark$\checkmark^\checkmark\\checkmarkc\checkmarki\checkmarkr\checkmarkc\checkmark$\checkmark \checkmark$\checkmark^\checkmark\\checkmarkc\checkmarki\checkmarkr\checkmarkc\checkmark$\checkmark \checkmark$\checkmark^\checkmark\\checkmarkc\checkmarki\checkmarkr\checkmarkc\checkmark$\checkmark\\checkmark$\checkmark^\checkmark\\checkmarkc\checkmarki\checkmarkr\checkmarkc\checkmark$\checkmarke\checkmark$\checkmark^\checkmark\\checkmarkc\checkmarki\checkmarkr\checkmarkc\checkmark$\checkmarkn\checkmark$\checkmark^\checkmark\\checkmarkc\checkmarki\checkmarkr\checkmarkc\checkmark$\checkmarkd\checkmark$\checkmark^\checkmark\\checkmarkc\checkmarki\checkmarkr\checkmarkc\checkmark$\checkmark{\checkmark$\checkmark^\checkmark\\checkmarkc\checkmarki\checkmarkr\checkmarkc\checkmark$\checkmarkt\checkmark$\checkmark^\checkmark\\checkmarkc\checkmarki\checkmarkr\checkmarkc\checkmark$\checkmarki\checkmark$\checkmark^\checkmark\\checkmarkc\checkmarki\checkmarkr\checkmarkc\checkmark$\checkmarkk\checkmark$\checkmark^\checkmark\\checkmarkc\checkmarki\checkmarkr\checkmarkc\checkmark$\checkmarkz\checkmark$\checkmark^\checkmark\\checkmarkc\checkmarki\checkmarkr\checkmarkc\checkmark$\checkmarkp\checkmark$\checkmark^\checkmark\\checkmarkc\checkmarki\checkmarkr\checkmarkc\checkmark$\checkmarki\checkmark$\checkmark^\checkmark\\checkmarkc\checkmarki\checkmarkr\checkmarkc\checkmark$\checkmarkc\checkmark$\checkmark^\checkmark\\checkmarkc\checkmarki\checkmarkr\checkmarkc\checkmark$\checkmarkt\checkmark$\checkmark^\checkmark\\checkmarkc\checkmarki\checkmarkr\checkmarkc\checkmark$\checkmarku\checkmark$\checkmark^\checkmark\\checkmarkc\checkmarki\checkmarkr\checkmarkc\checkmark$\checkmarkr\checkmark$\checkmark^\checkmark\\checkmarkc\checkmarki\checkmarkr\checkmarkc\checkmark$\checkmarke\checkmark$\checkmark^\checkmark\\checkmarkc\checkmarki\checkmarkr\checkmarkc\checkmark$\checkmark}\checkmark$\checkmark^\checkmark\\checkmarkc\checkmarki\checkmarkr\checkmarkc\checkmark$\checkmark
\checkmark$\checkmark^\checkmark\\checkmarkc\checkmarki\checkmarkr\checkmarkc\checkmark$\checkmark \checkmark$\checkmark^\checkmark\\checkmarkc\checkmarki\checkmarkr\checkmarkc\checkmark$\checkmark \checkmark$\checkmark^\checkmark\\checkmarkc\checkmarki\checkmarkr\checkmarkc\checkmark$\checkmark \checkmark$\checkmark^\checkmark\\checkmarkc\checkmarki\checkmarkr\checkmarkc\checkmark$\checkmark \checkmark$\checkmark^\checkmark\\checkmarkc\checkmarki\checkmarkr\checkmarkc\checkmark$\checkmark\\checkmark$\checkmark^\checkmark\\checkmarkc\checkmarki\checkmarkr\checkmarkc\checkmark$\checkmarkc\checkmark$\checkmark^\checkmark\\checkmarkc\checkmarki\checkmarkr\checkmarkc\checkmark$\checkmarka\checkmark$\checkmark^\checkmark\\checkmarkc\checkmarki\checkmarkr\checkmarkc\checkmark$\checkmarkp\checkmark$\checkmark^\checkmark\\checkmarkc\checkmarki\checkmarkr\checkmarkc\checkmark$\checkmarkt\checkmark$\checkmark^\checkmark\\checkmarkc\checkmarki\checkmarkr\checkmarkc\checkmark$\checkmarki\checkmark$\checkmark^\checkmark\\checkmarkc\checkmarki\checkmarkr\checkmarkc\checkmark$\checkmarko\checkmark$\checkmark^\checkmark\\checkmarkc\checkmarki\checkmarkr\checkmarkc\checkmark$\checkmarkn\checkmark$\checkmark^\checkmark\\checkmarkc\checkmarki\checkmarkr\checkmarkc\checkmark$\checkmark{\checkmark$\checkmark^\checkmark\\checkmarkc\checkmarki\checkmarkr\checkmarkc\checkmark$\checkmarkD\checkmark$\checkmark^\checkmark\\checkmarkc\checkmarki\checkmarkr\checkmarkc\checkmark$\checkmarki\checkmark$\checkmark^\checkmark\\checkmarkc\checkmarki\checkmarkr\checkmarkc\checkmark$\checkmarka\checkmark$\checkmark^\checkmark\\checkmarkc\checkmarki\checkmarkr\checkmarkc\checkmark$\checkmarkg\checkmark$\checkmark^\checkmark\\checkmarkc\checkmarki\checkmarkr\checkmarkc\checkmark$\checkmarkr\checkmark$\checkmark^\checkmark\\checkmarkc\checkmarki\checkmarkr\checkmarkc\checkmark$\checkmarka\checkmark$\checkmark^\checkmark\\checkmarkc\checkmarki\checkmarkr\checkmarkc\checkmark$\checkmarkm\checkmark$\checkmark^\checkmark\\checkmarkc\checkmarki\checkmarkr\checkmarkc\checkmark$\checkmarkm\checkmark$\checkmark^\checkmark\\checkmarkc\checkmarki\checkmarkr\checkmarkc\checkmark$\checkmarke\checkmark$\checkmark^\checkmark\\checkmarkc\checkmarki\checkmarkr\checkmarkc\checkmark$\checkmark \checkmark$\checkmark^\checkmark\\checkmarkc\checkmarki\checkmarkr\checkmarkc\checkmark$\checkmark«\checkmark$\checkmark^\checkmark\\checkmarkc\checkmarki\checkmarkr\checkmarkc\checkmark$\checkmarkB\checkmark$\checkmark^\checkmark\\checkmarkc\checkmarki\checkmarkr\checkmarkc\checkmark$\checkmarkê\checkmark$\checkmark^\checkmark\\checkmarkc\checkmarki\checkmarkr\checkmarkc\checkmark$\checkmarkt\checkmark$\checkmark^\checkmark\\checkmarkc\checkmarki\checkmarkr\checkmarkc\checkmark$\checkmarke\checkmark$\checkmark^\checkmark\\checkmarkc\checkmarki\checkmarkr\checkmarkc\checkmark$\checkmark \checkmark$\checkmark^\checkmark\\checkmarkc\checkmarki\checkmarkr\checkmarkc\checkmark$\checkmarkà\checkmark$\checkmark^\checkmark\\checkmarkc\checkmarki\checkmarkr\checkmarkc\checkmark$\checkmark \checkmark$\checkmark^\checkmark\\checkmarkc\checkmarki\checkmarkr\checkmarkc\checkmark$\checkmarkc\checkmark$\checkmark^\checkmark\\checkmarkc\checkmarki\checkmarkr\checkmarkc\checkmark$\checkmarko\checkmark$\checkmark^\checkmark\\checkmarkc\checkmarki\checkmarkr\checkmarkc\checkmark$\checkmarkr\checkmark$\checkmark^\checkmark\\checkmarkc\checkmarki\checkmarkr\checkmarkc\checkmark$\checkmarkn\checkmark$\checkmark^\checkmark\\checkmarkc\checkmarki\checkmarkr\checkmarkc\checkmark$\checkmarke\checkmark$\checkmark^\checkmark\\checkmarkc\checkmarki\checkmarkr\checkmarkc\checkmark$\checkmarks\checkmark$\checkmark^\checkmark\\checkmarkc\checkmarki\checkmarkr\checkmarkc\checkmark$\checkmark»\checkmark$\checkmark^\checkmark\\checkmarkc\checkmarki\checkmarkr\checkmarkc\checkmark$\checkmark \checkmark$\checkmark^\checkmark\\checkmarkc\checkmarki\checkmarkr\checkmarkc\checkmark$\checkmark—\checkmark$\checkmark^\checkmark\\checkmarkc\checkmarki\checkmarkr\checkmarkc\checkmark$\checkmark \checkmark$\checkmark^\checkmark\\checkmarkc\checkmarki\checkmarkr\checkmarkc\checkmark$\checkmarkE\checkmark$\checkmark^\checkmark\\checkmarkc\checkmarki\checkmarkr\checkmarkc\checkmark$\checkmarkx\checkmark$\checkmark^\checkmark\\checkmarkc\checkmarki\checkmarkr\checkmarkc\checkmark$\checkmarkp\checkmark$\checkmark^\checkmark\\checkmarkc\checkmarki\checkmarkr\checkmarkc\checkmark$\checkmarkr\checkmark$\checkmark^\checkmark\\checkmarkc\checkmarki\checkmarkr\checkmarkc\checkmark$\checkmarke\checkmark$\checkmark^\checkmark\\checkmarkc\checkmarki\checkmarkr\checkmarkc\checkmark$\checkmarks\checkmark$\checkmark^\checkmark\\checkmarkc\checkmarki\checkmarkr\checkmarkc\checkmark$\checkmarks\checkmark$\checkmark^\checkmark\\checkmarkc\checkmarki\checkmarkr\checkmarkc\checkmark$\checkmarki\checkmark$\checkmark^\checkmark\\checkmarkc\checkmarki\checkmarkr\checkmarkc\checkmark$\checkmarko\checkmark$\checkmark^\checkmark\\checkmarkc\checkmarki\checkmarkr\checkmarkc\checkmark$\checkmarkn\checkmark$\checkmark^\checkmark\\checkmarkc\checkmarki\checkmarkr\checkmarkc\checkmark$\checkmark \checkmark$\checkmark^\checkmark\\checkmarkc\checkmarki\checkmarkr\checkmarkc\checkmark$\checkmarkd\checkmark$\checkmark^\checkmark\\checkmarkc\checkmarki\checkmarkr\checkmarkc\checkmark$\checkmarku\checkmark$\checkmark^\checkmark\\checkmarkc\checkmarki\checkmarkr\checkmarkc\checkmark$\checkmark \checkmark$\checkmark^\checkmark\\checkmarkc\checkmarki\checkmarkr\checkmarkc\checkmark$\checkmarkb\checkmark$\checkmark^\checkmark\\checkmarkc\checkmarki\checkmarkr\checkmarkc\checkmark$\checkmarke\checkmark$\checkmark^\checkmark\\checkmarkc\checkmarki\checkmarkr\checkmarkc\checkmark$\checkmarks\checkmark$\checkmark^\checkmark\\checkmarkc\checkmarki\checkmarkr\checkmarkc\checkmark$\checkmarko\checkmark$\checkmark^\checkmark\\checkmarkc\checkmarki\checkmarkr\checkmarkc\checkmark$\checkmarki\checkmark$\checkmark^\checkmark\\checkmarkc\checkmarki\checkmarkr\checkmarkc\checkmark$\checkmarkn\checkmark$\checkmark^\checkmark\\checkmarkc\checkmarki\checkmarkr\checkmarkc\checkmark$\checkmark \checkmark$\checkmark^\checkmark\\checkmarkc\checkmarki\checkmarkr\checkmarkc\checkmark$\checkmarkf\checkmark$\checkmark^\checkmark\\checkmarkc\checkmarki\checkmarkr\checkmarkc\checkmark$\checkmarko\checkmark$\checkmark^\checkmark\\checkmarkc\checkmarki\checkmarkr\checkmarkc\checkmark$\checkmarkn\checkmark$\checkmark^\checkmark\\checkmarkc\checkmarki\checkmarkr\checkmarkc\checkmark$\checkmarkd\checkmark$\checkmark^\checkmark\\checkmarkc\checkmarki\checkmarkr\checkmarkc\checkmark$\checkmarka\checkmark$\checkmark^\checkmark\\checkmarkc\checkmarki\checkmarkr\checkmarkc\checkmark$\checkmarkm\checkmark$\checkmark^\checkmark\\checkmarkc\checkmarki\checkmarkr\checkmarkc\checkmark$\checkmarke\checkmark$\checkmark^\checkmark\\checkmarkc\checkmarki\checkmarkr\checkmarkc\checkmark$\checkmarkn\checkmark$\checkmark^\checkmark\\checkmarkc\checkmarki\checkmarkr\checkmarkc\checkmark$\checkmarkt\checkmark$\checkmark^\checkmark\\checkmarkc\checkmarki\checkmarkr\checkmarkc\checkmark$\checkmarka\checkmark$\checkmark^\checkmark\\checkmarkc\checkmarki\checkmarkr\checkmarkc\checkmark$\checkmarkl\checkmark$\checkmark^\checkmark\\checkmarkc\checkmarki\checkmarkr\checkmarkc\checkmark$\checkmark \checkmark$\checkmark^\checkmark\\checkmarkc\checkmarki\checkmarkr\checkmarkc\checkmark$\checkmarkd\checkmark$\checkmark^\checkmark\\checkmarkc\checkmarki\checkmarkr\checkmarkc\checkmark$\checkmarke\checkmark$\checkmark^\checkmark\\checkmarkc\checkmarki\checkmarkr\checkmarkc\checkmark$\checkmark \checkmark$\checkmark^\checkmark\\checkmarkc\checkmarki\checkmarkr\checkmarkc\checkmark$\checkmarkl\checkmark$\checkmark^\checkmark\\checkmarkc\checkmarki\checkmarkr\checkmarkc\checkmark$\checkmarka\checkmark$\checkmark^\checkmark\\checkmarkc\checkmarki\checkmarkr\checkmarkc\checkmark$\checkmark \checkmark$\checkmark^\checkmark\\checkmarkc\checkmarki\checkmarkr\checkmarkc\checkmark$\checkmarkC\checkmark$\checkmark^\checkmark\\checkmarkc\checkmarki\checkmarkr\checkmarkc\checkmark$\checkmarkN\checkmark$\checkmark^\checkmark\\checkmarkc\checkmarki\checkmarkr\checkmarkc\checkmark$\checkmarkC\checkmark$\checkmark^\checkmark\\checkmarkc\checkmarki\checkmarkr\checkmarkc\checkmark$\checkmark \checkmark$\checkmark^\checkmark\\checkmarkc\checkmarki\checkmarkr\checkmarkc\checkmark$\checkmark5\checkmark$\checkmark^\checkmark\\checkmarkc\checkmarki\checkmarkr\checkmarkc\checkmark$\checkmark \checkmark$\checkmark^\checkmark\\checkmarkc\checkmarki\checkmarkr\checkmarkc\checkmark$\checkmarka\checkmark$\checkmark^\checkmark\\checkmarkc\checkmarki\checkmarkr\checkmarkc\checkmark$\checkmarkx\checkmark$\checkmark^\checkmark\\checkmarkc\checkmarki\checkmarkr\checkmarkc\checkmark$\checkmarke\checkmark$\checkmark^\checkmark\\checkmarkc\checkmarki\checkmarkr\checkmarkc\checkmark$\checkmarks\checkmark$\checkmark^\checkmark\\checkmarkc\checkmarki\checkmarkr\checkmarkc\checkmark$\checkmark}\checkmark$\checkmark^\checkmark\\checkmarkc\checkmarki\checkmarkr\checkmarkc\checkmark$\checkmark
\checkmark$\checkmark^\checkmark\\checkmarkc\checkmarki\checkmarkr\checkmarkc\checkmark$\checkmark \checkmark$\checkmark^\checkmark\\checkmarkc\checkmarki\checkmarkr\checkmarkc\checkmark$\checkmark \checkmark$\checkmark^\checkmark\\checkmarkc\checkmarki\checkmarkr\checkmarkc\checkmark$\checkmark \checkmark$\checkmark^\checkmark\\checkmarkc\checkmarki\checkmarkr\checkmarkc\checkmark$\checkmark \checkmark$\checkmark^\checkmark\\checkmarkc\checkmarki\checkmarkr\checkmarkc\checkmark$\checkmark\\checkmark$\checkmark^\checkmark\\checkmarkc\checkmarki\checkmarkr\checkmarkc\checkmark$\checkmarkl\checkmark$\checkmark^\checkmark\\checkmarkc\checkmarki\checkmarkr\checkmarkc\checkmark$\checkmarka\checkmark$\checkmark^\checkmark\\checkmarkc\checkmarki\checkmarkr\checkmarkc\checkmark$\checkmarkb\checkmark$\checkmark^\checkmark\\checkmarkc\checkmarki\checkmarkr\checkmarkc\checkmark$\checkmarke\checkmark$\checkmark^\checkmark\\checkmarkc\checkmarki\checkmarkr\checkmarkc\checkmark$\checkmarkl\checkmark$\checkmark^\checkmark\\checkmarkc\checkmarki\checkmarkr\checkmarkc\checkmark$\checkmark{\checkmark$\checkmark^\checkmark\\checkmarkc\checkmarki\checkmarkr\checkmarkc\checkmark$\checkmarkf\checkmark$\checkmark^\checkmark\\checkmarkc\checkmarki\checkmarkr\checkmarkc\checkmark$\checkmarki\checkmark$\checkmark^\checkmark\\checkmarkc\checkmarki\checkmarkr\checkmarkc\checkmark$\checkmarkg\checkmark$\checkmark^\checkmark\\checkmarkc\checkmarki\checkmarkr\checkmarkc\checkmark$\checkmark:\checkmark$\checkmark^\checkmark\\checkmarkc\checkmarki\checkmarkr\checkmarkc\checkmark$\checkmarkb\checkmark$\checkmark^\checkmark\\checkmarkc\checkmarki\checkmarkr\checkmarkc\checkmark$\checkmarke\checkmark$\checkmark^\checkmark\\checkmarkc\checkmarki\checkmarkr\checkmarkc\checkmark$\checkmarkt\checkmark$\checkmark^\checkmark\\checkmarkc\checkmarki\checkmarkr\checkmarkc\checkmark$\checkmarke\checkmark$\checkmark^\checkmark\\checkmarkc\checkmarki\checkmarkr\checkmarkc\checkmark$\checkmark_\checkmark$\checkmark^\checkmark\\checkmarkc\checkmarki\checkmarkr\checkmarkc\checkmark$\checkmarka\checkmark$\checkmark^\checkmark\\checkmarkc\checkmarki\checkmarkr\checkmarkc\checkmark$\checkmark_\checkmark$\checkmark^\checkmark\\checkmarkc\checkmarki\checkmarkr\checkmarkc\checkmark$\checkmarkc\checkmark$\checkmark^\checkmark\\checkmarkc\checkmarki\checkmarkr\checkmarkc\checkmark$\checkmarko\checkmark$\checkmark^\checkmark\\checkmarkc\checkmarki\checkmarkr\checkmarkc\checkmark$\checkmarkr\checkmark$\checkmark^\checkmark\\checkmarkc\checkmarki\checkmarkr\checkmarkc\checkmark$\checkmarkn\checkmark$\checkmark^\checkmark\\checkmarkc\checkmarki\checkmarkr\checkmarkc\checkmark$\checkmarke\checkmark$\checkmark^\checkmark\\checkmarkc\checkmarki\checkmarkr\checkmarkc\checkmark$\checkmarks\checkmark$\checkmark^\checkmark\\checkmarkc\checkmarki\checkmarkr\checkmarkc\checkmark$\checkmark}\checkmark$\checkmark^\checkmark\\checkmarkc\checkmarki\checkmarkr\checkmarkc\checkmark$\checkmark
\checkmark$\checkmark^\checkmark\\checkmarkc\checkmarki\checkmarkr\checkmarkc\checkmark$\checkmark\\checkmark$\checkmark^\checkmark\\checkmarkc\checkmarki\checkmarkr\checkmarkc\checkmark$\checkmarke\checkmark$\checkmark^\checkmark\\checkmarkc\checkmarki\checkmarkr\checkmarkc\checkmark$\checkmarkn\checkmark$\checkmark^\checkmark\\checkmarkc\checkmarki\checkmarkr\checkmarkc\checkmark$\checkmarkd\checkmark$\checkmark^\checkmark\\checkmarkc\checkmarki\checkmarkr\checkmarkc\checkmark$\checkmark{\checkmark$\checkmark^\checkmark\\checkmarkc\checkmarki\checkmarkr\checkmarkc\checkmark$\checkmarkf\checkmark$\checkmark^\checkmark\\checkmarkc\checkmarki\checkmarkr\checkmarkc\checkmark$\checkmarki\checkmark$\checkmark^\checkmark\\checkmarkc\checkmarki\checkmarkr\checkmarkc\checkmark$\checkmarkg\checkmark$\checkmark^\checkmark\\checkmarkc\checkmarki\checkmarkr\checkmarkc\checkmark$\checkmarku\checkmark$\checkmark^\checkmark\\checkmarkc\checkmarki\checkmarkr\checkmarkc\checkmark$\checkmarkr\checkmark$\checkmark^\checkmark\\checkmarkc\checkmarki\checkmarkr\checkmarkc\checkmark$\checkmarke\checkmark$\checkmark^\checkmark\\checkmarkc\checkmarki\checkmarkr\checkmarkc\checkmark$\checkmark}\checkmark$\checkmark^\checkmark\\checkmarkc\checkmarki\checkmarkr\checkmarkc\checkmark$\checkmark
\checkmark$\checkmark^\checkmark\\checkmarkc\checkmarki\checkmarkr\checkmarkc\checkmark$\checkmark
\checkmark$\checkmark^\checkmark\\checkmarkc\checkmarki\checkmarkr\checkmarkc\checkmark$\checkmark\\checkmark$\checkmark^\checkmark\\checkmarkc\checkmarki\checkmarkr\checkmarkc\checkmark$\checkmarks\checkmark$\checkmark^\checkmark\\checkmarkc\checkmarki\checkmarkr\checkmarkc\checkmark$\checkmarku\checkmark$\checkmark^\checkmark\\checkmarkc\checkmarki\checkmarkr\checkmarkc\checkmark$\checkmarkb\checkmark$\checkmark^\checkmark\\checkmarkc\checkmarki\checkmarkr\checkmarkc\checkmark$\checkmarks\checkmark$\checkmark^\checkmark\\checkmarkc\checkmarki\checkmarkr\checkmarkc\checkmark$\checkmarke\checkmark$\checkmark^\checkmark\\checkmarkc\checkmarki\checkmarkr\checkmarkc\checkmark$\checkmarkc\checkmark$\checkmark^\checkmark\\checkmarkc\checkmarki\checkmarkr\checkmarkc\checkmark$\checkmarkt\checkmark$\checkmark^\checkmark\\checkmarkc\checkmarki\checkmarkr\checkmarkc\checkmark$\checkmarki\checkmark$\checkmark^\checkmark\\checkmarkc\checkmarki\checkmarkr\checkmarkc\checkmark$\checkmarko\checkmark$\checkmark^\checkmark\\checkmarkc\checkmarki\checkmarkr\checkmarkc\checkmark$\checkmarkn\checkmark$\checkmark^\checkmark\\checkmarkc\checkmarki\checkmarkr\checkmarkc\checkmark$\checkmark{\checkmark$\checkmark^\checkmark\\checkmarkc\checkmarki\checkmarkr\checkmarkc\checkmark$\checkmarkD\checkmark$\checkmark^\checkmark\\checkmarkc\checkmarki\checkmarkr\checkmarkc\checkmark$\checkmarki\checkmark$\checkmark^\checkmark\\checkmarkc\checkmarki\checkmarkr\checkmarkc\checkmark$\checkmarka\checkmark$\checkmark^\checkmark\\checkmarkc\checkmarki\checkmarkr\checkmarkc\checkmark$\checkmarkg\checkmark$\checkmark^\checkmark\\checkmarkc\checkmarki\checkmarkr\checkmarkc\checkmark$\checkmarkr\checkmark$\checkmark^\checkmark\\checkmarkc\checkmarki\checkmarkr\checkmarkc\checkmark$\checkmarka\checkmark$\checkmark^\checkmark\\checkmarkc\checkmarki\checkmarkr\checkmarkc\checkmark$\checkmarkm\checkmark$\checkmark^\checkmark\\checkmarkc\checkmarki\checkmarkr\checkmarkc\checkmark$\checkmarkm\checkmark$\checkmark^\checkmark\\checkmarkc\checkmarki\checkmarkr\checkmarkc\checkmark$\checkmarke\checkmark$\checkmark^\checkmark\\checkmarkc\checkmarki\checkmarkr\checkmarkc\checkmark$\checkmark \checkmark$\checkmark^\checkmark\\checkmarkc\checkmarki\checkmarkr\checkmarkc\checkmark$\checkmarkd\checkmark$\checkmark^\checkmark\\checkmarkc\checkmarki\checkmarkr\checkmarkc\checkmark$\checkmarke\checkmark$\checkmark^\checkmark\\checkmarkc\checkmarki\checkmarkr\checkmarkc\checkmark$\checkmarks\checkmark$\checkmark^\checkmark\\checkmarkc\checkmarki\checkmarkr\checkmarkc\checkmark$\checkmark \checkmark$\checkmark^\checkmark\\checkmarkc\checkmarki\checkmarkr\checkmarkc\checkmark$\checkmarki\checkmark$\checkmark^\checkmark\\checkmarkc\checkmarki\checkmarkr\checkmarkc\checkmark$\checkmarkn\checkmark$\checkmark^\checkmark\\checkmarkc\checkmarki\checkmarkr\checkmarkc\checkmark$\checkmarkt\checkmark$\checkmark^\checkmark\\checkmarkc\checkmarki\checkmarkr\checkmarkc\checkmark$\checkmarke\checkmark$\checkmark^\checkmark\\checkmarkc\checkmarki\checkmarkr\checkmarkc\checkmark$\checkmarkr\checkmark$\checkmark^\checkmark\\checkmarkc\checkmarki\checkmarkr\checkmarkc\checkmark$\checkmarka\checkmark$\checkmark^\checkmark\\checkmarkc\checkmarki\checkmarkr\checkmarkc\checkmark$\checkmarkc\checkmark$\checkmark^\checkmark\\checkmarkc\checkmarki\checkmarkr\checkmarkc\checkmark$\checkmarkt\checkmark$\checkmark^\checkmark\\checkmarkc\checkmarki\checkmarkr\checkmarkc\checkmark$\checkmarke\checkmark$\checkmark^\checkmark\\checkmarkc\checkmarki\checkmarkr\checkmarkc\checkmark$\checkmarku\checkmark$\checkmark^\checkmark\\checkmarkc\checkmarki\checkmarkr\checkmarkc\checkmark$\checkmarkr\checkmark$\checkmark^\checkmark\\checkmarkc\checkmarki\checkmarkr\checkmarkc\checkmark$\checkmarks\checkmark$\checkmark^\checkmark\\checkmarkc\checkmarki\checkmarkr\checkmarkc\checkmark$\checkmark \checkmark$\checkmark^\checkmark\\checkmarkc\checkmarki\checkmarkr\checkmarkc\checkmark$\checkmark«\checkmark$\checkmark^\checkmark\\checkmarkc\checkmarki\checkmarkr\checkmarkc\checkmark$\checkmarkL\checkmark$\checkmark^\checkmark\\checkmarkc\checkmarki\checkmarkr\checkmarkc\checkmark$\checkmarka\checkmark$\checkmark^\checkmark\\checkmarkc\checkmarki\checkmarkr\checkmarkc\checkmark$\checkmark \checkmark$\checkmark^\checkmark\\checkmarkc\checkmarki\checkmarkr\checkmarkc\checkmark$\checkmarkP\checkmark$\checkmark^\checkmark\\checkmarkc\checkmarki\checkmarkr\checkmarkc\checkmark$\checkmarki\checkmark$\checkmark^\checkmark\\checkmarkc\checkmarki\checkmarkr\checkmarkc\checkmark$\checkmarke\checkmark$\checkmark^\checkmark\\checkmarkc\checkmarki\checkmarkr\checkmarkc\checkmark$\checkmarku\checkmark$\checkmark^\checkmark\\checkmarkc\checkmarki\checkmarkr\checkmarkc\checkmark$\checkmarkv\checkmark$\checkmark^\checkmark\\checkmarkc\checkmarki\checkmarkr\checkmarkc\checkmark$\checkmarkr\checkmark$\checkmark^\checkmark\\checkmarkc\checkmarki\checkmarkr\checkmarkc\checkmark$\checkmarke\checkmark$\checkmark^\checkmark\\checkmarkc\checkmarki\checkmarkr\checkmarkc\checkmark$\checkmark»\checkmark$\checkmark^\checkmark\\checkmarkc\checkmarki\checkmarkr\checkmarkc\checkmark$\checkmark}\checkmark$\checkmark^\checkmark\\checkmarkc\checkmarki\checkmarkr\checkmarkc\checkmark$\checkmark
\checkmark$\checkmark^\checkmark\\checkmarkc\checkmarki\checkmarkr\checkmarkc\checkmark$\checkmarkC\checkmark$\checkmark^\checkmark\\checkmarkc\checkmarki\checkmarkr\checkmarkc\checkmark$\checkmarke\checkmark$\checkmark^\checkmark\\checkmarkc\checkmarki\checkmarkr\checkmarkc\checkmark$\checkmark \checkmark$\checkmark^\checkmark\\checkmarkc\checkmarki\checkmarkr\checkmarkc\checkmark$\checkmarkd\checkmark$\checkmark^\checkmark\\checkmarkc\checkmarki\checkmarkr\checkmarkc\checkmark$\checkmarki\checkmark$\checkmark^\checkmark\\checkmarkc\checkmarki\checkmarkr\checkmarkc\checkmark$\checkmarka\checkmark$\checkmark^\checkmark\\checkmarkc\checkmarki\checkmarkr\checkmarkc\checkmark$\checkmarkg\checkmark$\checkmark^\checkmark\\checkmarkc\checkmarki\checkmarkr\checkmarkc\checkmark$\checkmarkr\checkmark$\checkmark^\checkmark\\checkmarkc\checkmarki\checkmarkr\checkmarkc\checkmark$\checkmarka\checkmark$\checkmark^\checkmark\\checkmarkc\checkmarki\checkmarkr\checkmarkc\checkmark$\checkmarkm\checkmark$\checkmark^\checkmark\\checkmarkc\checkmarki\checkmarkr\checkmarkc\checkmark$\checkmarkm\checkmark$\checkmark^\checkmark\\checkmarkc\checkmarki\checkmarkr\checkmarkc\checkmark$\checkmarke\checkmark$\checkmark^\checkmark\\checkmarkc\checkmarki\checkmarkr\checkmarkc\checkmark$\checkmark \checkmark$\checkmark^\checkmark\\checkmarkc\checkmarki\checkmarkr\checkmarkc\checkmark$\checkmarkm\checkmark$\checkmark^\checkmark\\checkmarkc\checkmarki\checkmarkr\checkmarkc\checkmark$\checkmarke\checkmark$\checkmark^\checkmark\\checkmarkc\checkmarki\checkmarkr\checkmarkc\checkmark$\checkmarkt\checkmark$\checkmark^\checkmark\\checkmarkc\checkmarki\checkmarkr\checkmarkc\checkmark$\checkmark \checkmark$\checkmark^\checkmark\\checkmarkc\checkmarki\checkmarkr\checkmarkc\checkmark$\checkmarke\checkmark$\checkmark^\checkmark\\checkmarkc\checkmarki\checkmarkr\checkmarkc\checkmark$\checkmarkn\checkmark$\checkmark^\checkmark\\checkmarkc\checkmarki\checkmarkr\checkmarkc\checkmark$\checkmark \checkmark$\checkmark^\checkmark\\checkmarkc\checkmarki\checkmarkr\checkmarkc\checkmark$\checkmarké\checkmark$\checkmark^\checkmark\\checkmarkc\checkmarki\checkmarkr\checkmarkc\checkmark$\checkmarkv\checkmark$\checkmark^\checkmark\\checkmarkc\checkmarki\checkmarkr\checkmarkc\checkmark$\checkmarki\checkmark$\checkmark^\checkmark\\checkmarkc\checkmarki\checkmarkr\checkmarkc\checkmark$\checkmarkd\checkmark$\checkmark^\checkmark\\checkmarkc\checkmarki\checkmarkr\checkmarkc\checkmark$\checkmarke\checkmark$\checkmark^\checkmark\\checkmarkc\checkmarki\checkmarkr\checkmarkc\checkmark$\checkmarkn\checkmark$\checkmark^\checkmark\\checkmarkc\checkmarki\checkmarkr\checkmarkc\checkmark$\checkmarkc\checkmark$\checkmark^\checkmark\\checkmarkc\checkmarki\checkmarkr\checkmarkc\checkmark$\checkmarke\checkmark$\checkmark^\checkmark\\checkmarkc\checkmarki\checkmarkr\checkmarkc\checkmark$\checkmark \checkmark$\checkmark^\checkmark\\checkmarkc\checkmarki\checkmarkr\checkmarkc\checkmark$\checkmarkl\checkmark$\checkmark^\checkmark\\checkmarkc\checkmarki\checkmarkr\checkmarkc\checkmark$\checkmarka\checkmark$\checkmark^\checkmark\\checkmarkc\checkmarki\checkmarkr\checkmarkc\checkmark$\checkmark \checkmark$\checkmark^\checkmark\\checkmarkc\checkmarki\checkmarkr\checkmarkc\checkmark$\checkmark\\checkmark$\checkmark^\checkmark\\checkmarkc\checkmarki\checkmarkr\checkmarkc\checkmark$\checkmarkt\checkmark$\checkmark^\checkmark\\checkmarkc\checkmarki\checkmarkr\checkmarkc\checkmark$\checkmarke\checkmark$\checkmark^\checkmark\\checkmarkc\checkmarki\checkmarkr\checkmarkc\checkmark$\checkmarkx\checkmark$\checkmark^\checkmark\\checkmarkc\checkmarki\checkmarkr\checkmarkc\checkmark$\checkmarkt\checkmark$\checkmark^\checkmark\\checkmarkc\checkmarki\checkmarkr\checkmarkc\checkmark$\checkmarkb\checkmark$\checkmark^\checkmark\\checkmarkc\checkmarki\checkmarkr\checkmarkc\checkmark$\checkmarkf\checkmark$\checkmark^\checkmark\\checkmarkc\checkmarki\checkmarkr\checkmarkc\checkmark$\checkmark{\checkmark$\checkmark^\checkmark\\checkmarkc\checkmarki\checkmarkr\checkmarkc\checkmark$\checkmarkF\checkmark$\checkmark^\checkmark\\checkmarkc\checkmarki\checkmarkr\checkmarkc\checkmark$\checkmarko\checkmark$\checkmark^\checkmark\\checkmarkc\checkmarki\checkmarkr\checkmarkc\checkmark$\checkmarkn\checkmark$\checkmark^\checkmark\\checkmarkc\checkmarki\checkmarkr\checkmarkc\checkmark$\checkmarkc\checkmark$\checkmark^\checkmark\\checkmarkc\checkmarki\checkmarkr\checkmarkc\checkmark$\checkmarkt\checkmark$\checkmark^\checkmark\\checkmarkc\checkmarki\checkmarkr\checkmarkc\checkmark$\checkmarki\checkmark$\checkmark^\checkmark\\checkmarkc\checkmarki\checkmarkr\checkmarkc\checkmark$\checkmarko\checkmark$\checkmark^\checkmark\\checkmarkc\checkmarki\checkmarkr\checkmarkc\checkmark$\checkmarkn\checkmark$\checkmark^\checkmark\\checkmarkc\checkmarki\checkmarkr\checkmarkc\checkmark$\checkmark \checkmark$\checkmark^\checkmark\\checkmarkc\checkmarki\checkmarkr\checkmarkc\checkmark$\checkmarkP\checkmark$\checkmark^\checkmark\\checkmarkc\checkmarki\checkmarkr\checkmarkc\checkmark$\checkmarkr\checkmark$\checkmark^\checkmark\\checkmarkc\checkmarki\checkmarkr\checkmarkc\checkmark$\checkmarki\checkmark$\checkmark^\checkmark\\checkmarkc\checkmarki\checkmarkr\checkmarkc\checkmark$\checkmarkn\checkmark$\checkmark^\checkmark\\checkmarkc\checkmarki\checkmarkr\checkmarkc\checkmark$\checkmarkc\checkmark$\checkmark^\checkmark\\checkmarkc\checkmarki\checkmarkr\checkmarkc\checkmark$\checkmarki\checkmark$\checkmark^\checkmark\\checkmarkc\checkmarki\checkmarkr\checkmarkc\checkmark$\checkmarkp\checkmark$\checkmark^\checkmark\\checkmarkc\checkmarki\checkmarkr\checkmarkc\checkmark$\checkmarka\checkmark$\checkmark^\checkmark\\checkmarkc\checkmarki\checkmarkr\checkmarkc\checkmark$\checkmarkl\checkmark$\checkmark^\checkmark\\checkmarkc\checkmarki\checkmarkr\checkmarkc\checkmark$\checkmarke\checkmark$\checkmark^\checkmark\\checkmarkc\checkmarki\checkmarkr\checkmarkc\checkmark$\checkmark \checkmark$\checkmark^\checkmark\\checkmarkc\checkmarki\checkmarkr\checkmarkc\checkmark$\checkmark(\checkmark$\checkmark^\checkmark\\checkmarkc\checkmarki\checkmarkr\checkmarkc\checkmark$\checkmarkF\checkmark$\checkmark^\checkmark\\checkmarkc\checkmarki\checkmarkr\checkmarkc\checkmark$\checkmarkP\checkmark$\checkmark^\checkmark\\checkmarkc\checkmarki\checkmarkr\checkmarkc\checkmark$\checkmark)\checkmark$\checkmark^\checkmark\\checkmarkc\checkmarki\checkmarkr\checkmarkc\checkmark$\checkmark}\checkmark$\checkmark^\checkmark\\checkmarkc\checkmarki\checkmarkr\checkmarkc\checkmark$\checkmark \checkmark$\checkmark^\checkmark\\checkmarkc\checkmarki\checkmarkr\checkmarkc\checkmark$\checkmarke\checkmark$\checkmark^\checkmark\\checkmarkc\checkmarki\checkmarkr\checkmarkc\checkmark$\checkmarkt\checkmark$\checkmark^\checkmark\\checkmarkc\checkmarki\checkmarkr\checkmarkc\checkmark$\checkmark \checkmark$\checkmark^\checkmark\\checkmarkc\checkmarki\checkmarkr\checkmarkc\checkmark$\checkmarkl\checkmark$\checkmark^\checkmark\\checkmarkc\checkmarki\checkmarkr\checkmarkc\checkmark$\checkmarke\checkmark$\checkmark^\checkmark\\checkmarkc\checkmarki\checkmarkr\checkmarkc\checkmark$\checkmarks\checkmark$\checkmark^\checkmark\\checkmarkc\checkmarki\checkmarkr\checkmarkc\checkmark$\checkmark \checkmark$\checkmark^\checkmark\\checkmarkc\checkmarki\checkmarkr\checkmarkc\checkmark$\checkmark\\checkmark$\checkmark^\checkmark\\checkmarkc\checkmarki\checkmarkr\checkmarkc\checkmark$\checkmarkt\checkmark$\checkmark^\checkmark\\checkmarkc\checkmarki\checkmarkr\checkmarkc\checkmark$\checkmarke\checkmark$\checkmark^\checkmark\\checkmarkc\checkmarki\checkmarkr\checkmarkc\checkmark$\checkmarkx\checkmark$\checkmark^\checkmark\\checkmarkc\checkmarki\checkmarkr\checkmarkc\checkmark$\checkmarkt\checkmark$\checkmark^\checkmark\\checkmarkc\checkmarki\checkmarkr\checkmarkc\checkmark$\checkmarkb\checkmark$\checkmark^\checkmark\\checkmarkc\checkmarki\checkmarkr\checkmarkc\checkmark$\checkmarkf\checkmark$\checkmark^\checkmark\\checkmarkc\checkmarki\checkmarkr\checkmarkc\checkmark$\checkmark{\checkmark$\checkmark^\checkmark\\checkmarkc\checkmarki\checkmarkr\checkmarkc\checkmark$\checkmarkF\checkmark$\checkmark^\checkmark\\checkmarkc\checkmarki\checkmarkr\checkmarkc\checkmark$\checkmarko\checkmark$\checkmark^\checkmark\\checkmarkc\checkmarki\checkmarkr\checkmarkc\checkmark$\checkmarkn\checkmark$\checkmark^\checkmark\\checkmarkc\checkmarki\checkmarkr\checkmarkc\checkmark$\checkmarkc\checkmark$\checkmark^\checkmark\\checkmarkc\checkmarki\checkmarkr\checkmarkc\checkmark$\checkmarkt\checkmark$\checkmark^\checkmark\\checkmarkc\checkmarki\checkmarkr\checkmarkc\checkmark$\checkmarki\checkmark$\checkmark^\checkmark\\checkmarkc\checkmarki\checkmarkr\checkmarkc\checkmark$\checkmarko\checkmark$\checkmark^\checkmark\\checkmarkc\checkmarki\checkmarkr\checkmarkc\checkmark$\checkmarkn\checkmark$\checkmark^\checkmark\\checkmarkc\checkmarki\checkmarkr\checkmarkc\checkmark$\checkmarks\checkmark$\checkmark^\checkmark\\checkmarkc\checkmarki\checkmarkr\checkmarkc\checkmark$\checkmark \checkmark$\checkmark^\checkmark\\checkmarkc\checkmarki\checkmarkr\checkmarkc\checkmark$\checkmarkC\checkmark$\checkmark^\checkmark\\checkmarkc\checkmarki\checkmarkr\checkmarkc\checkmark$\checkmarko\checkmark$\checkmark^\checkmark\\checkmarkc\checkmarki\checkmarkr\checkmarkc\checkmark$\checkmarkn\checkmark$\checkmark^\checkmark\\checkmarkc\checkmarki\checkmarkr\checkmarkc\checkmark$\checkmarkt\checkmark$\checkmark^\checkmark\\checkmarkc\checkmarki\checkmarkr\checkmarkc\checkmark$\checkmarkr\checkmark$\checkmark^\checkmark\\checkmarkc\checkmarki\checkmarkr\checkmarkc\checkmark$\checkmarka\checkmark$\checkmark^\checkmark\\checkmarkc\checkmarki\checkmarkr\checkmarkc\checkmark$\checkmarki\checkmark$\checkmark^\checkmark\\checkmarkc\checkmarki\checkmarkr\checkmarkc\checkmark$\checkmarkn\checkmark$\checkmark^\checkmark\\checkmarkc\checkmarki\checkmarkr\checkmarkc\checkmark$\checkmarkt\checkmark$\checkmark^\checkmark\\checkmarkc\checkmarki\checkmarkr\checkmarkc\checkmark$\checkmarke\checkmark$\checkmark^\checkmark\\checkmarkc\checkmarki\checkmarkr\checkmarkc\checkmark$\checkmarks\checkmark$\checkmark^\checkmark\\checkmarkc\checkmarki\checkmarkr\checkmarkc\checkmark$\checkmark \checkmark$\checkmark^\checkmark\\checkmarkc\checkmarki\checkmarkr\checkmarkc\checkmark$\checkmark(\checkmark$\checkmark^\checkmark\\checkmarkc\checkmarki\checkmarkr\checkmarkc\checkmark$\checkmarkF\checkmark$\checkmark^\checkmark\\checkmarkc\checkmarki\checkmarkr\checkmarkc\checkmark$\checkmarkC\checkmark$\checkmark^\checkmark\\checkmarkc\checkmarki\checkmarkr\checkmarkc\checkmark$\checkmark)\checkmark$\checkmark^\checkmark\\checkmarkc\checkmarki\checkmarkr\checkmarkc\checkmark$\checkmark}\checkmark$\checkmark^\checkmark\\checkmarkc\checkmarki\checkmarkr\checkmarkc\checkmark$\checkmark \checkmark$\checkmark^\checkmark\\checkmarkc\checkmarki\checkmarkr\checkmarkc\checkmark$\checkmarka\checkmark$\checkmark^\checkmark\\checkmarkc\checkmarki\checkmarkr\checkmarkc\checkmark$\checkmarku\checkmark$\checkmark^\checkmark\\checkmarkc\checkmarki\checkmarkr\checkmarkc\checkmark$\checkmarkx\checkmark$\checkmark^\checkmark\\checkmarkc\checkmarki\checkmarkr\checkmarkc\checkmark$\checkmarkq\checkmark$\checkmark^\checkmark\\checkmarkc\checkmarki\checkmarkr\checkmarkc\checkmark$\checkmarku\checkmark$\checkmark^\checkmark\\checkmarkc\checkmarki\checkmarkr\checkmarkc\checkmark$\checkmarke\checkmark$\checkmark^\checkmark\\checkmarkc\checkmarki\checkmarkr\checkmarkc\checkmark$\checkmarkl\checkmark$\checkmark^\checkmark\\checkmarkc\checkmarki\checkmarkr\checkmarkc\checkmark$\checkmarkl\checkmark$\checkmark^\checkmark\\checkmarkc\checkmarki\checkmarkr\checkmarkc\checkmark$\checkmarke\checkmark$\checkmark^\checkmark\\checkmarkc\checkmarki\checkmarkr\checkmarkc\checkmark$\checkmarks\checkmark$\checkmark^\checkmark\\checkmarkc\checkmarki\checkmarkr\checkmarkc\checkmark$\checkmark \checkmark$\checkmark^\checkmark\\checkmarkc\checkmarki\checkmarkr\checkmarkc\checkmark$\checkmarkl\checkmark$\checkmark^\checkmark\\checkmarkc\checkmarki\checkmarkr\checkmarkc\checkmark$\checkmarka\checkmark$\checkmark^\checkmark\\checkmarkc\checkmarki\checkmarkr\checkmarkc\checkmark$\checkmark \checkmark$\checkmark^\checkmark\\checkmarkc\checkmarki\checkmarkr\checkmarkc\checkmark$\checkmarkf\checkmark$\checkmark^\checkmark\\checkmarkc\checkmarki\checkmarkr\checkmarkc\checkmark$\checkmarkr\checkmark$\checkmark^\checkmark\\checkmarkc\checkmarki\checkmarkr\checkmarkc\checkmark$\checkmarka\checkmark$\checkmark^\checkmark\\checkmarkc\checkmarki\checkmarkr\checkmarkc\checkmark$\checkmarki\checkmark$\checkmark^\checkmark\\checkmarkc\checkmarki\checkmarkr\checkmarkc\checkmark$\checkmarks\checkmark$\checkmark^\checkmark\\checkmarkc\checkmarki\checkmarkr\checkmarkc\checkmark$\checkmarke\checkmark$\checkmark^\checkmark\\checkmarkc\checkmarki\checkmarkr\checkmarkc\checkmark$\checkmarku\checkmark$\checkmark^\checkmark\\checkmarkc\checkmarki\checkmarkr\checkmarkc\checkmark$\checkmarks\checkmark$\checkmark^\checkmark\\checkmarkc\checkmarki\checkmarkr\checkmarkc\checkmark$\checkmarke\checkmark$\checkmark^\checkmark\\checkmarkc\checkmarki\checkmarkr\checkmarkc\checkmark$\checkmark \checkmark$\checkmark^\checkmark\\checkmarkc\checkmarki\checkmarkr\checkmarkc\checkmark$\checkmarkd\checkmark$\checkmark^\checkmark\\checkmarkc\checkmarki\checkmarkr\checkmarkc\checkmark$\checkmarko\checkmark$\checkmark^\checkmark\\checkmarkc\checkmarki\checkmarkr\checkmarkc\checkmark$\checkmarki\checkmark$\checkmark^\checkmark\\checkmarkc\checkmarki\checkmarkr\checkmarkc\checkmark$\checkmarkt\checkmark$\checkmark^\checkmark\\checkmarkc\checkmarki\checkmarkr\checkmarkc\checkmark$\checkmark \checkmark$\checkmark^\checkmark\\checkmarkc\checkmarki\checkmarkr\checkmarkc\checkmark$\checkmarks\checkmark$\checkmark^\checkmark\\checkmarkc\checkmarki\checkmarkr\checkmarkc\checkmark$\checkmarka\checkmark$\checkmark^\checkmark\\checkmarkc\checkmarki\checkmarkr\checkmarkc\checkmark$\checkmarkt\checkmark$\checkmark^\checkmark\\checkmarkc\checkmarki\checkmarkr\checkmarkc\checkmark$\checkmarki\checkmark$\checkmark^\checkmark\\checkmarkc\checkmarki\checkmarkr\checkmarkc\checkmark$\checkmarks\checkmark$\checkmark^\checkmark\\checkmarkc\checkmarki\checkmarkr\checkmarkc\checkmark$\checkmarkf\checkmark$\checkmark^\checkmark\\checkmarkc\checkmarki\checkmarkr\checkmarkc\checkmark$\checkmarka\checkmark$\checkmark^\checkmark\\checkmarkc\checkmarki\checkmarkr\checkmarkc\checkmark$\checkmarki\checkmark$\checkmark^\checkmark\\checkmarkc\checkmarki\checkmarkr\checkmarkc\checkmark$\checkmarkr\checkmark$\checkmark^\checkmark\\checkmarkc\checkmarki\checkmarkr\checkmarkc\checkmark$\checkmarke\checkmark$\checkmark^\checkmark\\checkmarkc\checkmarki\checkmarkr\checkmarkc\checkmark$\checkmark.\checkmark$\checkmark^\checkmark\\checkmarkc\checkmarki\checkmarkr\checkmarkc\checkmark$\checkmark
\checkmark$\checkmark^\checkmark\\checkmarkc\checkmarki\checkmarkr\checkmarkc\checkmark$\checkmark
\checkmark$\checkmark^\checkmark\\checkmarkc\checkmarki\checkmarkr\checkmarkc\checkmark$\checkmark\\checkmark$\checkmark^\checkmark\\checkmarkc\checkmarki\checkmarkr\checkmarkc\checkmark$\checkmarkb\checkmark$\checkmark^\checkmark\\checkmarkc\checkmarki\checkmarkr\checkmarkc\checkmark$\checkmarke\checkmark$\checkmark^\checkmark\\checkmarkc\checkmarki\checkmarkr\checkmarkc\checkmark$\checkmarkg\checkmark$\checkmark^\checkmark\\checkmarkc\checkmarki\checkmarkr\checkmarkc\checkmark$\checkmarki\checkmark$\checkmark^\checkmark\\checkmarkc\checkmarki\checkmarkr\checkmarkc\checkmark$\checkmarkn\checkmark$\checkmark^\checkmark\\checkmarkc\checkmarki\checkmarkr\checkmarkc\checkmark$\checkmark{\checkmark$\checkmark^\checkmark\\checkmarkc\checkmarki\checkmarkr\checkmarkc\checkmark$\checkmarki\checkmark$\checkmark^\checkmark\\checkmarkc\checkmarki\checkmarkr\checkmarkc\checkmark$\checkmarkt\checkmark$\checkmark^\checkmark\\checkmarkc\checkmarki\checkmarkr\checkmarkc\checkmark$\checkmarke\checkmark$\checkmark^\checkmark\\checkmarkc\checkmarki\checkmarkr\checkmarkc\checkmark$\checkmarkm\checkmark$\checkmark^\checkmark\\checkmarkc\checkmarki\checkmarkr\checkmarkc\checkmark$\checkmarki\checkmark$\checkmark^\checkmark\\checkmarkc\checkmarki\checkmarkr\checkmarkc\checkmark$\checkmarkz\checkmark$\checkmark^\checkmark\\checkmarkc\checkmarki\checkmarkr\checkmarkc\checkmark$\checkmarke\checkmark$\checkmark^\checkmark\\checkmarkc\checkmarki\checkmarkr\checkmarkc\checkmark$\checkmark}\checkmark$\checkmark^\checkmark\\checkmarkc\checkmarki\checkmarkr\checkmarkc\checkmark$\checkmark
\checkmark$\checkmark^\checkmark\\checkmarkc\checkmarki\checkmarkr\checkmarkc\checkmark$\checkmark \checkmark$\checkmark^\checkmark\\checkmarkc\checkmarki\checkmarkr\checkmarkc\checkmark$\checkmark \checkmark$\checkmark^\checkmark\\checkmarkc\checkmarki\checkmarkr\checkmarkc\checkmark$\checkmark \checkmark$\checkmark^\checkmark\\checkmarkc\checkmarki\checkmarkr\checkmarkc\checkmark$\checkmark \checkmark$\checkmark^\checkmark\\checkmarkc\checkmarki\checkmarkr\checkmarkc\checkmark$\checkmark\\checkmark$\checkmark^\checkmark\\checkmarkc\checkmarki\checkmarkr\checkmarkc\checkmark$\checkmarki\checkmark$\checkmark^\checkmark\\checkmarkc\checkmarki\checkmarkr\checkmarkc\checkmark$\checkmarkt\checkmark$\checkmark^\checkmark\\checkmarkc\checkmarki\checkmarkr\checkmarkc\checkmark$\checkmarke\checkmark$\checkmark^\checkmark\\checkmarkc\checkmarki\checkmarkr\checkmarkc\checkmark$\checkmarkm\checkmark$\checkmark^\checkmark\\checkmarkc\checkmarki\checkmarkr\checkmarkc\checkmark$\checkmark \checkmark$\checkmark^\checkmark\\checkmarkc\checkmarki\checkmarkr\checkmarkc\checkmark$\checkmark\\checkmark$\checkmark^\checkmark\\checkmarkc\checkmarki\checkmarkr\checkmarkc\checkmark$\checkmarkt\checkmark$\checkmark^\checkmark\\checkmarkc\checkmarki\checkmarkr\checkmarkc\checkmark$\checkmarke\checkmark$\checkmark^\checkmark\\checkmarkc\checkmarki\checkmarkr\checkmarkc\checkmark$\checkmarkx\checkmark$\checkmark^\checkmark\\checkmarkc\checkmarki\checkmarkr\checkmarkc\checkmark$\checkmarkt\checkmark$\checkmark^\checkmark\\checkmarkc\checkmarki\checkmarkr\checkmarkc\checkmark$\checkmarkb\checkmark$\checkmark^\checkmark\\checkmarkc\checkmarki\checkmarkr\checkmarkc\checkmark$\checkmarkf\checkmark$\checkmark^\checkmark\\checkmarkc\checkmarki\checkmarkr\checkmarkc\checkmark$\checkmark{\checkmark$\checkmark^\checkmark\\checkmarkc\checkmarki\checkmarkr\checkmarkc\checkmark$\checkmarkF\checkmark$\checkmark^\checkmark\\checkmarkc\checkmarki\checkmarkr\checkmarkc\checkmark$\checkmarkP\checkmark$\checkmark^\checkmark\\checkmarkc\checkmarki\checkmarkr\checkmarkc\checkmark$\checkmark \checkmark$\checkmark^\checkmark\\checkmarkc\checkmarki\checkmarkr\checkmarkc\checkmark$\checkmark:\checkmark$\checkmark^\checkmark\\checkmarkc\checkmarki\checkmarkr\checkmarkc\checkmark$\checkmark}\checkmark$\checkmark^\checkmark\\checkmarkc\checkmarki\checkmarkr\checkmarkc\checkmark$\checkmark \checkmark$\checkmark^\checkmark\\checkmarkc\checkmarki\checkmarkr\checkmarkc\checkmark$\checkmarkA\checkmark$\checkmark^\checkmark\\checkmarkc\checkmarki\checkmarkr\checkmarkc\checkmark$\checkmarks\checkmark$\checkmark^\checkmark\\checkmarkc\checkmarki\checkmarkr\checkmarkc\checkmark$\checkmarks\checkmark$\checkmark^\checkmark\\checkmarkc\checkmarki\checkmarkr\checkmarkc\checkmark$\checkmarku\checkmark$\checkmark^\checkmark\\checkmarkc\checkmarki\checkmarkr\checkmarkc\checkmark$\checkmarkr\checkmark$\checkmark^\checkmark\\checkmarkc\checkmarki\checkmarkr\checkmarkc\checkmark$\checkmarke\checkmark$\checkmark^\checkmark\\checkmarkc\checkmarki\checkmarkr\checkmarkc\checkmark$\checkmarkr\checkmark$\checkmark^\checkmark\\checkmarkc\checkmarki\checkmarkr\checkmarkc\checkmark$\checkmark \checkmark$\checkmark^\checkmark\\checkmarkc\checkmarki\checkmarkr\checkmarkc\checkmark$\checkmarkl\checkmark$\checkmark^\checkmark\\checkmarkc\checkmarki\checkmarkr\checkmarkc\checkmark$\checkmark'\checkmark$\checkmark^\checkmark\\checkmarkc\checkmarki\checkmarkr\checkmarkc\checkmark$\checkmarku\checkmark$\checkmark^\checkmark\\checkmarkc\checkmarki\checkmarkr\checkmarkc\checkmark$\checkmarks\checkmark$\checkmark^\checkmark\\checkmarkc\checkmarki\checkmarkr\checkmarkc\checkmark$\checkmarki\checkmark$\checkmark^\checkmark\\checkmarkc\checkmarki\checkmarkr\checkmarkc\checkmark$\checkmarkn\checkmark$\checkmark^\checkmark\\checkmarkc\checkmarki\checkmarkr\checkmarkc\checkmark$\checkmarka\checkmark$\checkmark^\checkmark\\checkmarkc\checkmarki\checkmarkr\checkmarkc\checkmark$\checkmarkg\checkmark$\checkmark^\checkmark\\checkmarkc\checkmarki\checkmarkr\checkmarkc\checkmark$\checkmarke\checkmark$\checkmark^\checkmark\\checkmarkc\checkmarki\checkmarkr\checkmarkc\checkmark$\checkmark \checkmark$\checkmark^\checkmark\\checkmarkc\checkmarki\checkmarkr\checkmarkc\checkmark$\checkmarka\checkmark$\checkmark^\checkmark\\checkmarkc\checkmarki\checkmarkr\checkmarkc\checkmark$\checkmarku\checkmark$\checkmark^\checkmark\\checkmarkc\checkmarki\checkmarkr\checkmarkc\checkmark$\checkmarkt\checkmark$\checkmark^\checkmark\\checkmarkc\checkmarki\checkmarkr\checkmarkc\checkmark$\checkmarko\checkmark$\checkmark^\checkmark\\checkmarkc\checkmarki\checkmarkr\checkmarkc\checkmark$\checkmarkm\checkmark$\checkmark^\checkmark\\checkmarkc\checkmarki\checkmarkr\checkmarkc\checkmark$\checkmarka\checkmark$\checkmark^\checkmark\\checkmarkc\checkmarki\checkmarkr\checkmarkc\checkmark$\checkmarkt\checkmark$\checkmark^\checkmark\\checkmarkc\checkmarki\checkmarkr\checkmarkc\checkmark$\checkmarki\checkmark$\checkmark^\checkmark\\checkmarkc\checkmarki\checkmarkr\checkmarkc\checkmark$\checkmarkq\checkmark$\checkmark^\checkmark\\checkmarkc\checkmarki\checkmarkr\checkmarkc\checkmark$\checkmarku\checkmark$\checkmark^\checkmark\\checkmarkc\checkmarki\checkmarkr\checkmarkc\checkmark$\checkmarke\checkmark$\checkmark^\checkmark\\checkmarkc\checkmarki\checkmarkr\checkmarkc\checkmark$\checkmark \checkmark$\checkmark^\checkmark\\checkmarkc\checkmarki\checkmarkr\checkmarkc\checkmark$\checkmarke\checkmark$\checkmark^\checkmark\\checkmarkc\checkmarki\checkmarkr\checkmarkc\checkmark$\checkmarkt\checkmark$\checkmark^\checkmark\\checkmarkc\checkmarki\checkmarkr\checkmarkc\checkmark$\checkmark \checkmark$\checkmark^\checkmark\\checkmarkc\checkmarki\checkmarkr\checkmarkc\checkmark$\checkmarkp\checkmark$\checkmark^\checkmark\\checkmarkc\checkmarki\checkmarkr\checkmarkc\checkmark$\checkmarkr\checkmark$\checkmark^\checkmark\\checkmarkc\checkmarki\checkmarkr\checkmarkc\checkmark$\checkmarké\checkmark$\checkmark^\checkmark\\checkmarkc\checkmarki\checkmarkr\checkmarkc\checkmark$\checkmarkc\checkmark$\checkmark^\checkmark\\checkmarkc\checkmarki\checkmarkr\checkmarkc\checkmark$\checkmarki\checkmark$\checkmark^\checkmark\\checkmarkc\checkmarki\checkmarkr\checkmarkc\checkmark$\checkmarks\checkmark$\checkmark^\checkmark\\checkmarkc\checkmarki\checkmarkr\checkmarkc\checkmark$\checkmark \checkmark$\checkmark^\checkmark\\checkmarkc\checkmarki\checkmarkr\checkmarkc\checkmark$\checkmark(\checkmark$\checkmark^\checkmark\\checkmarkc\checkmarki\checkmarkr\checkmarkc\checkmark$\checkmark$\checkmark$\checkmark^\checkmark\\checkmarkc\checkmarki\checkmarkr\checkmarkc\checkmark$\checkmark\\checkmark$\checkmark^\checkmark\\checkmarkc\checkmarki\checkmarkr\checkmarkc\checkmark$\checkmarkp\checkmark$\checkmark^\checkmark\\checkmarkc\checkmarki\checkmarkr\checkmarkc\checkmark$\checkmarkm\checkmark$\checkmark^\checkmark\\checkmarkc\checkmarki\checkmarkr\checkmarkc\checkmark$\checkmark \checkmark$\checkmark^\checkmark\\checkmarkc\checkmarki\checkmarkr\checkmarkc\checkmark$\checkmark0\checkmark$\checkmark^\checkmark\\checkmarkc\checkmarki\checkmarkr\checkmarkc\checkmark$\checkmark{\checkmark$\checkmark^\checkmark\\checkmarkc\checkmarki\checkmarkr\checkmarkc\checkmark$\checkmark,\checkmark$\checkmark^\checkmark\\checkmarkc\checkmarki\checkmarkr\checkmarkc\checkmark$\checkmark}\checkmark$\checkmark^\checkmark\\checkmarkc\checkmarki\checkmarkr\checkmarkc\checkmark$\checkmark0\checkmark$\checkmark^\checkmark\\checkmarkc\checkmarki\checkmarkr\checkmarkc\checkmark$\checkmark5\checkmark$\checkmark^\checkmark\\checkmarkc\checkmarki\checkmarkr\checkmarkc\checkmark$\checkmark$\checkmark$\checkmark^\checkmark\\checkmarkc\checkmarki\checkmarkr\checkmarkc\checkmark$\checkmark \checkmark$\checkmark^\checkmark\\checkmarkc\checkmarki\checkmarkr\checkmarkc\checkmark$\checkmarkm\checkmark$\checkmark^\checkmark\\checkmarkc\checkmarki\checkmarkr\checkmarkc\checkmark$\checkmarkm\checkmark$\checkmark^\checkmark\\checkmarkc\checkmarki\checkmarkr\checkmarkc\checkmark$\checkmark)\checkmark$\checkmark^\checkmark\\checkmarkc\checkmarki\checkmarkr\checkmarkc\checkmark$\checkmark \checkmark$\checkmark^\checkmark\\checkmarkc\checkmarki\checkmarkr\checkmarkc\checkmark$\checkmarkd\checkmark$\checkmark^\checkmark\\checkmarkc\checkmarki\checkmarkr\checkmarkc\checkmark$\checkmark'\checkmark$\checkmark^\checkmark\\checkmarkc\checkmarki\checkmarkr\checkmarkc\checkmark$\checkmarku\checkmark$\checkmark^\checkmark\\checkmarkc\checkmarki\checkmarkr\checkmarkc\checkmark$\checkmarkn\checkmark$\checkmark^\checkmark\\checkmarkc\checkmarki\checkmarkr\checkmarkc\checkmark$\checkmarke\checkmark$\checkmark^\checkmark\\checkmarkc\checkmarki\checkmarkr\checkmarkc\checkmark$\checkmark \checkmark$\checkmark^\checkmark\\checkmarkc\checkmarki\checkmarkr\checkmarkc\checkmark$\checkmarkp\checkmark$\checkmark^\checkmark\\checkmarkc\checkmarki\checkmarkr\checkmarkc\checkmark$\checkmarki\checkmark$\checkmark^\checkmark\\checkmarkc\checkmarki\checkmarkr\checkmarkc\checkmark$\checkmarkè\checkmark$\checkmark^\checkmark\\checkmarkc\checkmarki\checkmarkr\checkmarkc\checkmark$\checkmarkc\checkmark$\checkmark^\checkmark\\checkmarkc\checkmarki\checkmarkr\checkmarkc\checkmark$\checkmarke\checkmark$\checkmark^\checkmark\\checkmarkc\checkmarki\checkmarkr\checkmarkc\checkmark$\checkmark \checkmark$\checkmark^\checkmark\\checkmarkc\checkmarki\checkmarkr\checkmarkc\checkmark$\checkmarkb\checkmark$\checkmark^\checkmark\\checkmarkc\checkmarki\checkmarkr\checkmarkc\checkmark$\checkmarkr\checkmark$\checkmark^\checkmark\\checkmarkc\checkmarki\checkmarkr\checkmarkc\checkmark$\checkmarku\checkmark$\checkmark^\checkmark\\checkmarkc\checkmarki\checkmarkr\checkmarkc\checkmark$\checkmarkt\checkmark$\checkmark^\checkmark\\checkmarkc\checkmarki\checkmarkr\checkmarkc\checkmark$\checkmarke\checkmark$\checkmark^\checkmark\\checkmarkc\checkmarki\checkmarkr\checkmarkc\checkmark$\checkmark \checkmark$\checkmark^\checkmark\\checkmarkc\checkmarki\checkmarkr\checkmarkc\checkmark$\checkmarke\checkmark$\checkmark^\checkmark\\checkmarkc\checkmarki\checkmarkr\checkmarkc\checkmark$\checkmarkn\checkmark$\checkmark^\checkmark\\checkmarkc\checkmarki\checkmarkr\checkmarkc\checkmark$\checkmark \checkmark$\checkmark^\checkmark\\checkmarkc\checkmarki\checkmarkr\checkmarkc\checkmark$\checkmark5\checkmark$\checkmark^\checkmark\\checkmarkc\checkmarki\checkmarkr\checkmarkc\checkmark$\checkmark \checkmark$\checkmark^\checkmark\\checkmarkc\checkmarki\checkmarkr\checkmarkc\checkmark$\checkmarka\checkmark$\checkmark^\checkmark\\checkmarkc\checkmarki\checkmarkr\checkmarkc\checkmark$\checkmarkx\checkmark$\checkmark^\checkmark\\checkmarkc\checkmarki\checkmarkr\checkmarkc\checkmark$\checkmarke\checkmark$\checkmark^\checkmark\\checkmarkc\checkmarki\checkmarkr\checkmarkc\checkmark$\checkmarks\checkmark$\checkmark^\checkmark\\checkmarkc\checkmarki\checkmarkr\checkmarkc\checkmark$\checkmark \checkmark$\checkmark^\checkmark\\checkmarkc\checkmarki\checkmarkr\checkmarkc\checkmark$\checkmarks\checkmark$\checkmark^\checkmark\\checkmarkc\checkmarki\checkmarkr\checkmarkc\checkmark$\checkmarki\checkmark$\checkmark^\checkmark\\checkmarkc\checkmarki\checkmarkr\checkmarkc\checkmark$\checkmarkm\checkmark$\checkmark^\checkmark\\checkmarkc\checkmarki\checkmarkr\checkmarkc\checkmark$\checkmarku\checkmark$\checkmark^\checkmark\\checkmarkc\checkmarki\checkmarkr\checkmarkc\checkmark$\checkmarkl\checkmark$\checkmark^\checkmark\\checkmarkc\checkmarki\checkmarkr\checkmarkc\checkmark$\checkmarkt\checkmark$\checkmark^\checkmark\\checkmarkc\checkmarki\checkmarkr\checkmarkc\checkmark$\checkmarka\checkmark$\checkmark^\checkmark\\checkmarkc\checkmarki\checkmarkr\checkmarkc\checkmark$\checkmarkn\checkmark$\checkmark^\checkmark\\checkmarkc\checkmarki\checkmarkr\checkmarkc\checkmark$\checkmarké\checkmark$\checkmark^\checkmark\\checkmarkc\checkmarki\checkmarkr\checkmarkc\checkmark$\checkmarks\checkmark$\checkmark^\checkmark\\checkmarkc\checkmarki\checkmarkr\checkmarkc\checkmark$\checkmark \checkmark$\checkmark^\checkmark\\checkmarkc\checkmarki\checkmarkr\checkmarkc\checkmark$\checkmarkd\checkmark$\checkmark^\checkmark\\checkmarkc\checkmarki\checkmarkr\checkmarkc\checkmark$\checkmarke\checkmark$\checkmark^\checkmark\\checkmarkc\checkmarki\checkmarkr\checkmarkc\checkmark$\checkmarkp\checkmark$\checkmark^\checkmark\\checkmarkc\checkmarki\checkmarkr\checkmarkc\checkmark$\checkmarku\checkmark$\checkmark^\checkmark\\checkmarkc\checkmarki\checkmarkr\checkmarkc\checkmark$\checkmarki\checkmark$\checkmark^\checkmark\\checkmarkc\checkmarki\checkmarkr\checkmarkc\checkmark$\checkmarks\checkmark$\checkmark^\checkmark\\checkmarkc\checkmarki\checkmarkr\checkmarkc\checkmark$\checkmark \checkmark$\checkmark^\checkmark\\checkmarkc\checkmarki\checkmarkr\checkmarkc\checkmark$\checkmarku\checkmark$\checkmark^\checkmark\\checkmarkc\checkmarki\checkmarkr\checkmarkc\checkmark$\checkmarkn\checkmark$\checkmark^\checkmark\\checkmarkc\checkmarki\checkmarkr\checkmarkc\checkmark$\checkmark \checkmark$\checkmark^\checkmark\\checkmarkc\checkmarki\checkmarkr\checkmarkc\checkmark$\checkmarkr\checkmark$\checkmark^\checkmark\\checkmarkc\checkmarki\checkmarkr\checkmarkc\checkmark$\checkmarké\checkmark$\checkmark^\checkmark\\checkmarkc\checkmarki\checkmarkr\checkmarkc\checkmark$\checkmarks\checkmark$\checkmark^\checkmark\\checkmarkc\checkmarki\checkmarkr\checkmarkc\checkmark$\checkmarke\checkmark$\checkmark^\checkmark\\checkmarkc\checkmarki\checkmarkr\checkmarkc\checkmark$\checkmarka\checkmark$\checkmark^\checkmark\\checkmarkc\checkmarki\checkmarkr\checkmarkc\checkmark$\checkmarku\checkmark$\checkmark^\checkmark\\checkmarkc\checkmarki\checkmarkr\checkmarkc\checkmark$\checkmark \checkmark$\checkmark^\checkmark\\checkmarkc\checkmarki\checkmarkr\checkmarkc\checkmark$\checkmarkW\checkmark$\checkmark^\checkmark\\checkmarkc\checkmarki\checkmarkr\checkmarkc\checkmark$\checkmarki\checkmark$\checkmark^\checkmark\\checkmarkc\checkmarki\checkmarkr\checkmarkc\checkmark$\checkmark-\checkmark$\checkmark^\checkmark\\checkmarkc\checkmarki\checkmarkr\checkmarkc\checkmark$\checkmarkF\checkmark$\checkmark^\checkmark\\checkmarkc\checkmarki\checkmarkr\checkmarkc\checkmark$\checkmarki\checkmark$\checkmark^\checkmark\\checkmarkc\checkmarki\checkmarkr\checkmarkc\checkmark$\checkmark.\checkmark$\checkmark^\checkmark\\checkmarkc\checkmarki\checkmarkr\checkmarkc\checkmark$\checkmark
\checkmark$\checkmark^\checkmark\\checkmarkc\checkmarki\checkmarkr\checkmarkc\checkmark$\checkmark \checkmark$\checkmark^\checkmark\\checkmarkc\checkmarki\checkmarkr\checkmarkc\checkmark$\checkmark \checkmark$\checkmark^\checkmark\\checkmarkc\checkmarki\checkmarkr\checkmarkc\checkmark$\checkmark \checkmark$\checkmark^\checkmark\\checkmarkc\checkmarki\checkmarkr\checkmarkc\checkmark$\checkmark \checkmark$\checkmark^\checkmark\\checkmarkc\checkmarki\checkmarkr\checkmarkc\checkmark$\checkmark\\checkmark$\checkmark^\checkmark\\checkmarkc\checkmarki\checkmarkr\checkmarkc\checkmark$\checkmarki\checkmark$\checkmark^\checkmark\\checkmarkc\checkmarki\checkmarkr\checkmarkc\checkmark$\checkmarkt\checkmark$\checkmark^\checkmark\\checkmarkc\checkmarki\checkmarkr\checkmarkc\checkmark$\checkmarke\checkmark$\checkmark^\checkmark\\checkmarkc\checkmarki\checkmarkr\checkmarkc\checkmark$\checkmarkm\checkmark$\checkmark^\checkmark\\checkmarkc\checkmarki\checkmarkr\checkmarkc\checkmark$\checkmark \checkmark$\checkmark^\checkmark\\checkmarkc\checkmarki\checkmarkr\checkmarkc\checkmark$\checkmark\\checkmark$\checkmark^\checkmark\\checkmarkc\checkmarki\checkmarkr\checkmarkc\checkmark$\checkmarkt\checkmark$\checkmark^\checkmark\\checkmarkc\checkmarki\checkmarkr\checkmarkc\checkmark$\checkmarke\checkmark$\checkmark^\checkmark\\checkmarkc\checkmarki\checkmarkr\checkmarkc\checkmark$\checkmarkx\checkmark$\checkmark^\checkmark\\checkmarkc\checkmarki\checkmarkr\checkmarkc\checkmark$\checkmarkt\checkmark$\checkmark^\checkmark\\checkmarkc\checkmarki\checkmarkr\checkmarkc\checkmark$\checkmarkb\checkmark$\checkmark^\checkmark\\checkmarkc\checkmarki\checkmarkr\checkmarkc\checkmark$\checkmarkf\checkmark$\checkmark^\checkmark\\checkmarkc\checkmarki\checkmarkr\checkmarkc\checkmark$\checkmark{\checkmark$\checkmark^\checkmark\\checkmarkc\checkmarki\checkmarkr\checkmarkc\checkmark$\checkmarkF\checkmark$\checkmark^\checkmark\\checkmarkc\checkmarki\checkmarkr\checkmarkc\checkmark$\checkmarkC\checkmark$\checkmark^\checkmark\\checkmarkc\checkmarki\checkmarkr\checkmarkc\checkmark$\checkmark1\checkmark$\checkmark^\checkmark\\checkmarkc\checkmarki\checkmarkr\checkmarkc\checkmark$\checkmark \checkmark$\checkmark^\checkmark\\checkmarkc\checkmarki\checkmarkr\checkmarkc\checkmark$\checkmark:\checkmark$\checkmark^\checkmark\\checkmarkc\checkmarki\checkmarkr\checkmarkc\checkmark$\checkmark}\checkmark$\checkmark^\checkmark\\checkmarkc\checkmarki\checkmarkr\checkmarkc\checkmark$\checkmark \checkmark$\checkmark^\checkmark\\checkmarkc\checkmarki\checkmarkr\checkmarkc\checkmark$\checkmarkL\checkmark$\checkmark^\checkmark\\checkmarkc\checkmarki\checkmarkr\checkmarkc\checkmark$\checkmarke\checkmark$\checkmark^\checkmark\\checkmarkc\checkmarki\checkmarkr\checkmarkc\checkmark$\checkmark \checkmark$\checkmark^\checkmark\\checkmarkc\checkmarki\checkmarkr\checkmarkc\checkmark$\checkmarks\checkmark$\checkmark^\checkmark\\checkmarkc\checkmarki\checkmarkr\checkmarkc\checkmark$\checkmarky\checkmark$\checkmark^\checkmark\\checkmarkc\checkmarki\checkmarkr\checkmarkc\checkmark$\checkmarks\checkmark$\checkmark^\checkmark\\checkmarkc\checkmarki\checkmarkr\checkmarkc\checkmark$\checkmarkt\checkmark$\checkmark^\checkmark\\checkmarkc\checkmarki\checkmarkr\checkmarkc\checkmark$\checkmarkè\checkmark$\checkmark^\checkmark\\checkmarkc\checkmarki\checkmarkr\checkmarkc\checkmark$\checkmarkm\checkmark$\checkmark^\checkmark\\checkmarkc\checkmarki\checkmarkr\checkmarkc\checkmark$\checkmarke\checkmark$\checkmark^\checkmark\\checkmarkc\checkmarki\checkmarkr\checkmarkc\checkmark$\checkmark \checkmark$\checkmark^\checkmark\\checkmarkc\checkmarki\checkmarkr\checkmarkc\checkmark$\checkmarkd\checkmark$\checkmark^\checkmark\\checkmarkc\checkmarki\checkmarkr\checkmarkc\checkmark$\checkmarko\checkmark$\checkmark^\checkmark\\checkmarkc\checkmarki\checkmarkr\checkmarkc\checkmark$\checkmarki\checkmark$\checkmark^\checkmark\\checkmarkc\checkmarki\checkmarkr\checkmarkc\checkmark$\checkmarkt\checkmark$\checkmark^\checkmark\\checkmarkc\checkmarki\checkmarkr\checkmarkc\checkmark$\checkmark \checkmark$\checkmark^\checkmark\\checkmarkc\checkmarki\checkmarkr\checkmarkc\checkmark$\checkmarkr\checkmark$\checkmark^\checkmark\\checkmarkc\checkmarki\checkmarkr\checkmarkc\checkmark$\checkmarke\checkmark$\checkmark^\checkmark\\checkmarkc\checkmarki\checkmarkr\checkmarkc\checkmark$\checkmarks\checkmark$\checkmark^\checkmark\\checkmarkc\checkmarki\checkmarkr\checkmarkc\checkmark$\checkmarkp\checkmark$\checkmark^\checkmark\\checkmarkc\checkmarki\checkmarkr\checkmarkc\checkmark$\checkmarke\checkmark$\checkmark^\checkmark\\checkmarkc\checkmarki\checkmarkr\checkmarkc\checkmark$\checkmarkc\checkmark$\checkmark^\checkmark\\checkmarkc\checkmarki\checkmarkr\checkmarkc\checkmark$\checkmarkt\checkmark$\checkmark^\checkmark\\checkmarkc\checkmarki\checkmarkr\checkmarkc\checkmark$\checkmarke\checkmark$\checkmark^\checkmark\\checkmarkc\checkmarki\checkmarkr\checkmarkc\checkmark$\checkmarkr\checkmark$\checkmark^\checkmark\\checkmarkc\checkmarki\checkmarkr\checkmarkc\checkmark$\checkmark \checkmark$\checkmark^\checkmark\\checkmarkc\checkmarki\checkmarkr\checkmarkc\checkmark$\checkmarkl\checkmark$\checkmark^\checkmark\\checkmarkc\checkmarki\checkmarkr\checkmarkc\checkmark$\checkmarke\checkmark$\checkmark^\checkmark\\checkmarkc\checkmarki\checkmarkr\checkmarkc\checkmark$\checkmarks\checkmark$\checkmark^\checkmark\\checkmarkc\checkmarki\checkmarkr\checkmarkc\checkmark$\checkmark \checkmark$\checkmark^\checkmark\\checkmarkc\checkmarki\checkmarkr\checkmarkc\checkmark$\checkmarkn\checkmark$\checkmark^\checkmark\\checkmarkc\checkmarki\checkmarkr\checkmarkc\checkmark$\checkmarko\checkmark$\checkmark^\checkmark\\checkmarkc\checkmarki\checkmarkr\checkmarkc\checkmark$\checkmarkr\checkmark$\checkmark^\checkmark\\checkmarkc\checkmarki\checkmarkr\checkmarkc\checkmark$\checkmarkm\checkmark$\checkmark^\checkmark\\checkmarkc\checkmarki\checkmarkr\checkmarkc\checkmark$\checkmarke\checkmark$\checkmark^\checkmark\\checkmarkc\checkmarki\checkmarkr\checkmarkc\checkmark$\checkmarks\checkmark$\checkmark^\checkmark\\checkmarkc\checkmarki\checkmarkr\checkmarkc\checkmark$\checkmark \checkmark$\checkmark^\checkmark\\checkmarkc\checkmarki\checkmarkr\checkmarkc\checkmark$\checkmarkd\checkmark$\checkmark^\checkmark\\checkmarkc\checkmarki\checkmarkr\checkmarkc\checkmark$\checkmarke\checkmark$\checkmark^\checkmark\\checkmarkc\checkmarki\checkmarkr\checkmarkc\checkmark$\checkmark \checkmark$\checkmark^\checkmark\\checkmarkc\checkmarki\checkmarkr\checkmarkc\checkmark$\checkmarks\checkmark$\checkmark^\checkmark\\checkmarkc\checkmarki\checkmarkr\checkmarkc\checkmark$\checkmarké\checkmark$\checkmark^\checkmark\\checkmarkc\checkmarki\checkmarkr\checkmarkc\checkmark$\checkmarkc\checkmark$\checkmark^\checkmark\\checkmarkc\checkmarki\checkmarkr\checkmarkc\checkmark$\checkmarku\checkmark$\checkmark^\checkmark\\checkmarkc\checkmarki\checkmarkr\checkmarkc\checkmark$\checkmarkr\checkmark$\checkmark^\checkmark\\checkmarkc\checkmarki\checkmarkr\checkmarkc\checkmark$\checkmarki\checkmark$\checkmark^\checkmark\\checkmarkc\checkmarki\checkmarkr\checkmarkc\checkmark$\checkmarkt\checkmark$\checkmark^\checkmark\\checkmarkc\checkmarki\checkmarkr\checkmarkc\checkmark$\checkmarké\checkmark$\checkmark^\checkmark\\checkmarkc\checkmarki\checkmarkr\checkmarkc\checkmark$\checkmark \checkmark$\checkmark^\checkmark\\checkmarkc\checkmarki\checkmarkr\checkmarkc\checkmark$\checkmarkm\checkmark$\checkmark^\checkmark\\checkmarkc\checkmarki\checkmarkr\checkmarkc\checkmark$\checkmarka\checkmark$\checkmark^\checkmark\\checkmarkc\checkmarki\checkmarkr\checkmarkc\checkmark$\checkmarkc\checkmark$\checkmark^\checkmark\\checkmarkc\checkmarki\checkmarkr\checkmarkc\checkmark$\checkmarkh\checkmark$\checkmark^\checkmark\\checkmarkc\checkmarki\checkmarkr\checkmarkc\checkmark$\checkmarki\checkmark$\checkmark^\checkmark\\checkmarkc\checkmarki\checkmarkr\checkmarkc\checkmark$\checkmarkn\checkmark$\checkmark^\checkmark\\checkmarkc\checkmarki\checkmarkr\checkmarkc\checkmark$\checkmarke\checkmark$\checkmark^\checkmark\\checkmarkc\checkmarki\checkmarkr\checkmarkc\checkmark$\checkmark \checkmark$\checkmark^\checkmark\\checkmarkc\checkmarki\checkmarkr\checkmarkc\checkmark$\checkmark(\checkmark$\checkmark^\checkmark\\checkmarkc\checkmarki\checkmarkr\checkmarkc\checkmark$\checkmarkE\checkmark$\checkmark^\checkmark\\checkmarkc\checkmarki\checkmarkr\checkmarkc\checkmark$\checkmarkN\checkmark$\checkmark^\checkmark\\checkmarkc\checkmarki\checkmarkr\checkmarkc\checkmark$\checkmark \checkmark$\checkmark^\checkmark\\checkmarkc\checkmarki\checkmarkr\checkmarkc\checkmark$\checkmarkI\checkmark$\checkmark^\checkmark\\checkmarkc\checkmarki\checkmarkr\checkmarkc\checkmark$\checkmarkS\checkmark$\checkmark^\checkmark\\checkmarkc\checkmarki\checkmarkr\checkmarkc\checkmark$\checkmarkO\checkmark$\checkmark^\checkmark\\checkmarkc\checkmarki\checkmarkr\checkmarkc\checkmark$\checkmark \checkmark$\checkmark^\checkmark\\checkmarkc\checkmarki\checkmarkr\checkmarkc\checkmark$\checkmark1\checkmark$\checkmark^\checkmark\\checkmarkc\checkmarki\checkmarkr\checkmarkc\checkmark$\checkmark3\checkmark$\checkmark^\checkmark\\checkmarkc\checkmarki\checkmarkr\checkmarkc\checkmark$\checkmark8\checkmark$\checkmark^\checkmark\\checkmarkc\checkmarki\checkmarkr\checkmarkc\checkmark$\checkmark5\checkmark$\checkmark^\checkmark\\checkmarkc\checkmarki\checkmarkr\checkmarkc\checkmark$\checkmark0\checkmark$\checkmark^\checkmark\\checkmarkc\checkmarki\checkmarkr\checkmarkc\checkmark$\checkmark \checkmark$\checkmark^\checkmark\\checkmarkc\checkmarki\checkmarkr\checkmarkc\checkmark$\checkmark—\checkmark$\checkmark^\checkmark\\checkmarkc\checkmarki\checkmarkr\checkmarkc\checkmark$\checkmark \checkmark$\checkmark^\checkmark\\checkmarkc\checkmarki\checkmarkr\checkmarkc\checkmark$\checkmarkE\checkmark$\checkmark^\checkmark\\checkmarkc\checkmarki\checkmarkr\checkmarkc\checkmark$\checkmark-\checkmark$\checkmark^\checkmark\\checkmarkc\checkmarki\checkmarkr\checkmarkc\checkmark$\checkmarkS\checkmark$\checkmark^\checkmark\\checkmarkc\checkmarki\checkmarkr\checkmarkc\checkmark$\checkmarkT\checkmark$\checkmark^\checkmark\\checkmarkc\checkmarki\checkmarkr\checkmarkc\checkmark$\checkmarkO\checkmark$\checkmark^\checkmark\\checkmarkc\checkmarki\checkmarkr\checkmarkc\checkmark$\checkmarkP\checkmark$\checkmark^\checkmark\\checkmarkc\checkmarki\checkmarkr\checkmarkc\checkmark$\checkmark \checkmark$\checkmark^\checkmark\\checkmarkc\checkmarki\checkmarkr\checkmarkc\checkmark$\checkmarkc\checkmark$\checkmark^\checkmark\\checkmarkc\checkmarki\checkmarkr\checkmarkc\checkmark$\checkmarka\checkmark$\checkmark^\checkmark\\checkmarkc\checkmarki\checkmarkr\checkmarkc\checkmark$\checkmarkt\checkmark$\checkmark^\checkmark\\checkmarkc\checkmarki\checkmarkr\checkmarkc\checkmark$\checkmarké\checkmark$\checkmark^\checkmark\\checkmarkc\checkmarki\checkmarkr\checkmarkc\checkmark$\checkmarkg\checkmark$\checkmark^\checkmark\\checkmarkc\checkmarki\checkmarkr\checkmarkc\checkmark$\checkmarko\checkmark$\checkmark^\checkmark\\checkmarkc\checkmarki\checkmarkr\checkmarkc\checkmark$\checkmarkr\checkmark$\checkmark^\checkmark\\checkmarkc\checkmarki\checkmarkr\checkmarkc\checkmark$\checkmarki\checkmark$\checkmark^\checkmark\\checkmarkc\checkmarki\checkmarkr\checkmarkc\checkmark$\checkmarke\checkmark$\checkmark^\checkmark\\checkmarkc\checkmarki\checkmarkr\checkmarkc\checkmark$\checkmark \checkmark$\checkmark^\checkmark\\checkmarkc\checkmarki\checkmarkr\checkmarkc\checkmark$\checkmark0\checkmark$\checkmark^\checkmark\\checkmarkc\checkmarki\checkmarkr\checkmarkc\checkmark$\checkmark)\checkmark$\checkmark^\checkmark\\checkmarkc\checkmarki\checkmarkr\checkmarkc\checkmark$\checkmark.\checkmark$\checkmark^\checkmark\\checkmarkc\checkmarki\checkmarkr\checkmarkc\checkmark$\checkmark
\checkmark$\checkmark^\checkmark\\checkmarkc\checkmarki\checkmarkr\checkmarkc\checkmark$\checkmark \checkmark$\checkmark^\checkmark\\checkmarkc\checkmarki\checkmarkr\checkmarkc\checkmark$\checkmark \checkmark$\checkmark^\checkmark\\checkmarkc\checkmarki\checkmarkr\checkmarkc\checkmark$\checkmark \checkmark$\checkmark^\checkmark\\checkmarkc\checkmarki\checkmarkr\checkmarkc\checkmark$\checkmark \checkmark$\checkmark^\checkmark\\checkmarkc\checkmarki\checkmarkr\checkmarkc\checkmark$\checkmark\\checkmark$\checkmark^\checkmark\\checkmarkc\checkmarki\checkmarkr\checkmarkc\checkmark$\checkmarki\checkmark$\checkmark^\checkmark\\checkmarkc\checkmarki\checkmarkr\checkmarkc\checkmark$\checkmarkt\checkmark$\checkmark^\checkmark\\checkmarkc\checkmarki\checkmarkr\checkmarkc\checkmark$\checkmarke\checkmark$\checkmark^\checkmark\\checkmarkc\checkmarki\checkmarkr\checkmarkc\checkmark$\checkmarkm\checkmark$\checkmark^\checkmark\\checkmarkc\checkmarki\checkmarkr\checkmarkc\checkmark$\checkmark \checkmark$\checkmark^\checkmark\\checkmarkc\checkmarki\checkmarkr\checkmarkc\checkmark$\checkmark\\checkmark$\checkmark^\checkmark\\checkmarkc\checkmarki\checkmarkr\checkmarkc\checkmark$\checkmarkt\checkmark$\checkmark^\checkmark\\checkmarkc\checkmarki\checkmarkr\checkmarkc\checkmark$\checkmarke\checkmark$\checkmark^\checkmark\\checkmarkc\checkmarki\checkmarkr\checkmarkc\checkmark$\checkmarkx\checkmark$\checkmark^\checkmark\\checkmarkc\checkmarki\checkmarkr\checkmarkc\checkmark$\checkmarkt\checkmark$\checkmark^\checkmark\\checkmarkc\checkmarki\checkmarkr\checkmarkc\checkmark$\checkmarkb\checkmark$\checkmark^\checkmark\\checkmarkc\checkmarki\checkmarkr\checkmarkc\checkmark$\checkmarkf\checkmark$\checkmark^\checkmark\\checkmarkc\checkmarki\checkmarkr\checkmarkc\checkmark$\checkmark{\checkmark$\checkmark^\checkmark\\checkmarkc\checkmarki\checkmarkr\checkmarkc\checkmark$\checkmarkF\checkmark$\checkmark^\checkmark\\checkmarkc\checkmarki\checkmarkr\checkmarkc\checkmark$\checkmarkC\checkmark$\checkmark^\checkmark\\checkmarkc\checkmarki\checkmarkr\checkmarkc\checkmark$\checkmark2\checkmark$\checkmark^\checkmark\\checkmarkc\checkmarki\checkmarkr\checkmarkc\checkmark$\checkmark \checkmark$\checkmark^\checkmark\\checkmarkc\checkmarki\checkmarkr\checkmarkc\checkmark$\checkmark:\checkmark$\checkmark^\checkmark\\checkmarkc\checkmarki\checkmarkr\checkmarkc\checkmark$\checkmark}\checkmark$\checkmark^\checkmark\\checkmarkc\checkmarki\checkmarkr\checkmarkc\checkmark$\checkmark \checkmark$\checkmark^\checkmark\\checkmarkc\checkmarki\checkmarkr\checkmarkc\checkmark$\checkmarkL\checkmark$\checkmark^\checkmark\\checkmarkc\checkmarki\checkmarkr\checkmarkc\checkmark$\checkmarke\checkmark$\checkmark^\checkmark\\checkmarkc\checkmarki\checkmarkr\checkmarkc\checkmark$\checkmark \checkmark$\checkmark^\checkmark\\checkmarkc\checkmarki\checkmarkr\checkmarkc\checkmark$\checkmarkc\checkmark$\checkmark^\checkmark\\checkmarkc\checkmarki\checkmarkr\checkmarkc\checkmark$\checkmarko\checkmark$\checkmark^\checkmark\\checkmarkc\checkmarki\checkmarkr\checkmarkc\checkmark$\checkmarkû\checkmark$\checkmark^\checkmark\\checkmarkc\checkmarki\checkmarkr\checkmarkc\checkmark$\checkmarkt\checkmark$\checkmark^\checkmark\\checkmarkc\checkmarki\checkmarkr\checkmarkc\checkmark$\checkmark \checkmark$\checkmark^\checkmark\\checkmarkc\checkmarki\checkmarkr\checkmarkc\checkmark$\checkmarkd\checkmark$\checkmark^\checkmark\\checkmarkc\checkmarki\checkmarkr\checkmarkc\checkmark$\checkmarke\checkmark$\checkmark^\checkmark\\checkmarkc\checkmarki\checkmarkr\checkmarkc\checkmark$\checkmark \checkmark$\checkmark^\checkmark\\checkmarkc\checkmarki\checkmarkr\checkmarkc\checkmark$\checkmarkf\checkmark$\checkmark^\checkmark\\checkmarkc\checkmarki\checkmarkr\checkmarkc\checkmark$\checkmarka\checkmark$\checkmark^\checkmark\\checkmarkc\checkmarki\checkmarkr\checkmarkc\checkmark$\checkmarkb\checkmark$\checkmark^\checkmark\\checkmarkc\checkmarki\checkmarkr\checkmarkc\checkmark$\checkmarkr\checkmark$\checkmark^\checkmark\\checkmarkc\checkmarki\checkmarkr\checkmarkc\checkmark$\checkmarki\checkmark$\checkmark^\checkmark\\checkmarkc\checkmarki\checkmarkr\checkmarkc\checkmark$\checkmarkc\checkmark$\checkmark^\checkmark\\checkmarkc\checkmarki\checkmarkr\checkmarkc\checkmark$\checkmarka\checkmark$\checkmark^\checkmark\\checkmarkc\checkmarki\checkmarkr\checkmarkc\checkmark$\checkmarkt\checkmark$\checkmark^\checkmark\\checkmarkc\checkmarki\checkmarkr\checkmarkc\checkmark$\checkmarki\checkmark$\checkmark^\checkmark\\checkmarkc\checkmarki\checkmarkr\checkmarkc\checkmark$\checkmarko\checkmark$\checkmark^\checkmark\\checkmarkc\checkmarki\checkmarkr\checkmarkc\checkmark$\checkmarkn\checkmark$\checkmark^\checkmark\\checkmarkc\checkmarki\checkmarkr\checkmarkc\checkmark$\checkmark \checkmark$\checkmark^\checkmark\\checkmarkc\checkmarki\checkmarkr\checkmarkc\checkmark$\checkmarkd\checkmark$\checkmark^\checkmark\\checkmarkc\checkmarki\checkmarkr\checkmarkc\checkmark$\checkmarko\checkmark$\checkmark^\checkmark\\checkmarkc\checkmarki\checkmarkr\checkmarkc\checkmark$\checkmarki\checkmark$\checkmark^\checkmark\\checkmarkc\checkmarki\checkmarkr\checkmarkc\checkmark$\checkmarkt\checkmark$\checkmark^\checkmark\\checkmarkc\checkmarki\checkmarkr\checkmarkc\checkmark$\checkmark \checkmark$\checkmark^\checkmark\\checkmarkc\checkmarki\checkmarkr\checkmarkc\checkmark$\checkmarkê\checkmark$\checkmark^\checkmark\\checkmarkc\checkmarki\checkmarkr\checkmarkc\checkmark$\checkmarkt\checkmark$\checkmark^\checkmark\\checkmarkc\checkmarki\checkmarkr\checkmarkc\checkmark$\checkmarkr\checkmark$\checkmark^\checkmark\\checkmarkc\checkmarki\checkmarkr\checkmarkc\checkmark$\checkmarke\checkmark$\checkmark^\checkmark\\checkmarkc\checkmarki\checkmarkr\checkmarkc\checkmark$\checkmark \checkmark$\checkmark^\checkmark\\checkmarkc\checkmarki\checkmarkr\checkmarkc\checkmark$\checkmarkm\checkmark$\checkmark^\checkmark\\checkmarkc\checkmarki\checkmarkr\checkmarkc\checkmark$\checkmarki\checkmark$\checkmark^\checkmark\\checkmarkc\checkmarki\checkmarkr\checkmarkc\checkmark$\checkmarkn\checkmark$\checkmark^\checkmark\\checkmarkc\checkmarki\checkmarkr\checkmarkc\checkmark$\checkmarki\checkmark$\checkmark^\checkmark\\checkmarkc\checkmarki\checkmarkr\checkmarkc\checkmark$\checkmarkm\checkmark$\checkmark^\checkmark\\checkmarkc\checkmarki\checkmarkr\checkmarkc\checkmark$\checkmarki\checkmark$\checkmark^\checkmark\\checkmarkc\checkmarki\checkmarkr\checkmarkc\checkmark$\checkmarks\checkmark$\checkmark^\checkmark\\checkmarkc\checkmarki\checkmarkr\checkmarkc\checkmark$\checkmarké\checkmark$\checkmark^\checkmark\\checkmarkc\checkmarki\checkmarkr\checkmarkc\checkmark$\checkmark \checkmark$\checkmark^\checkmark\\checkmarkc\checkmarki\checkmarkr\checkmarkc\checkmark$\checkmarke\checkmark$\checkmark^\checkmark\\checkmarkc\checkmarki\checkmarkr\checkmarkc\checkmark$\checkmarkn\checkmark$\checkmark^\checkmark\\checkmarkc\checkmarki\checkmarkr\checkmarkc\checkmark$\checkmark \checkmark$\checkmark^\checkmark\\checkmarkc\checkmarki\checkmarkr\checkmarkc\checkmark$\checkmarku\checkmark$\checkmark^\checkmark\\checkmarkc\checkmarki\checkmarkr\checkmarkc\checkmark$\checkmarkt\checkmark$\checkmark^\checkmark\\checkmarkc\checkmarki\checkmarkr\checkmarkc\checkmark$\checkmarki\checkmark$\checkmark^\checkmark\\checkmarkc\checkmarki\checkmarkr\checkmarkc\checkmark$\checkmarkl\checkmark$\checkmark^\checkmark\\checkmarkc\checkmarki\checkmarkr\checkmarkc\checkmark$\checkmarki\checkmark$\checkmark^\checkmark\\checkmarkc\checkmarki\checkmarkr\checkmarkc\checkmark$\checkmarks\checkmark$\checkmark^\checkmark\\checkmarkc\checkmarki\checkmarkr\checkmarkc\checkmark$\checkmarka\checkmark$\checkmark^\checkmark\\checkmarkc\checkmarki\checkmarkr\checkmarkc\checkmark$\checkmarkn\checkmark$\checkmark^\checkmark\\checkmarkc\checkmarki\checkmarkr\checkmarkc\checkmark$\checkmarkt\checkmark$\checkmark^\checkmark\\checkmarkc\checkmarki\checkmarkr\checkmarkc\checkmark$\checkmark \checkmark$\checkmark^\checkmark\\checkmarkc\checkmarki\checkmarkr\checkmarkc\checkmark$\checkmarkd\checkmark$\checkmark^\checkmark\\checkmarkc\checkmarki\checkmarkr\checkmarkc\checkmark$\checkmarke\checkmark$\checkmark^\checkmark\\checkmarkc\checkmarki\checkmarkr\checkmarkc\checkmark$\checkmarks\checkmark$\checkmark^\checkmark\\checkmarkc\checkmarki\checkmarkr\checkmarkc\checkmark$\checkmark \checkmark$\checkmark^\checkmark\\checkmarkc\checkmarki\checkmarkr\checkmarkc\checkmark$\checkmarkc\checkmark$\checkmark^\checkmark\\checkmarkc\checkmarki\checkmarkr\checkmarkc\checkmark$\checkmarko\checkmark$\checkmark^\checkmark\\checkmarkc\checkmarki\checkmarkr\checkmarkc\checkmark$\checkmarkm\checkmark$\checkmark^\checkmark\\checkmarkc\checkmarki\checkmarkr\checkmarkc\checkmark$\checkmarkp\checkmark$\checkmark^\checkmark\\checkmarkc\checkmarki\checkmarkr\checkmarkc\checkmark$\checkmarko\checkmark$\checkmark^\checkmark\\checkmarkc\checkmarki\checkmarkr\checkmarkc\checkmark$\checkmarks\checkmark$\checkmark^\checkmark\\checkmarkc\checkmarki\checkmarkr\checkmarkc\checkmark$\checkmarka\checkmark$\checkmark^\checkmark\\checkmarkc\checkmarki\checkmarkr\checkmarkc\checkmark$\checkmarkn\checkmark$\checkmark^\checkmark\\checkmarkc\checkmarki\checkmarkr\checkmarkc\checkmark$\checkmarkt\checkmark$\checkmark^\checkmark\\checkmarkc\checkmarki\checkmarkr\checkmarkc\checkmark$\checkmarks\checkmark$\checkmark^\checkmark\\checkmarkc\checkmarki\checkmarkr\checkmarkc\checkmark$\checkmark \checkmark$\checkmark^\checkmark\\checkmarkc\checkmarki\checkmarkr\checkmarkc\checkmark$\checkmarkn\checkmark$\checkmark^\checkmark\\checkmarkc\checkmarki\checkmarkr\checkmarkc\checkmark$\checkmarko\checkmark$\checkmark^\checkmark\\checkmarkc\checkmarki\checkmarkr\checkmarkc\checkmark$\checkmarkr\checkmark$\checkmark^\checkmark\\checkmarkc\checkmarki\checkmarkr\checkmarkc\checkmark$\checkmarkm\checkmark$\checkmark^\checkmark\\checkmarkc\checkmarki\checkmarkr\checkmarkc\checkmark$\checkmarka\checkmark$\checkmark^\checkmark\\checkmarkc\checkmarki\checkmarkr\checkmarkc\checkmark$\checkmarkl\checkmark$\checkmark^\checkmark\\checkmarkc\checkmarki\checkmarkr\checkmarkc\checkmark$\checkmarki\checkmark$\checkmark^\checkmark\\checkmarkc\checkmarki\checkmarkr\checkmarkc\checkmark$\checkmarks\checkmark$\checkmark^\checkmark\\checkmarkc\checkmarki\checkmarkr\checkmarkc\checkmark$\checkmarké\checkmark$\checkmark^\checkmark\\checkmarkc\checkmarki\checkmarkr\checkmarkc\checkmark$\checkmarks\checkmark$\checkmark^\checkmark\\checkmarkc\checkmarki\checkmarkr\checkmarkc\checkmark$\checkmark \checkmark$\checkmark^\checkmark\\checkmarkc\checkmarki\checkmarkr\checkmarkc\checkmark$\checkmarke\checkmark$\checkmark^\checkmark\\checkmarkc\checkmarki\checkmarkr\checkmarkc\checkmark$\checkmarkt\checkmark$\checkmark^\checkmark\\checkmarkc\checkmarki\checkmarkr\checkmarkc\checkmark$\checkmark \checkmark$\checkmark^\checkmark\\checkmarkc\checkmarki\checkmarkr\checkmarkc\checkmark$\checkmarkd\checkmark$\checkmark^\checkmark\\checkmarkc\checkmarki\checkmarkr\checkmarkc\checkmark$\checkmarki\checkmark$\checkmark^\checkmark\\checkmarkc\checkmarki\checkmarkr\checkmarkc\checkmark$\checkmarks\checkmark$\checkmark^\checkmark\\checkmarkc\checkmarki\checkmarkr\checkmarkc\checkmark$\checkmarkp\checkmark$\checkmark^\checkmark\\checkmarkc\checkmarki\checkmarkr\checkmarkc\checkmark$\checkmarko\checkmark$\checkmark^\checkmark\\checkmarkc\checkmarki\checkmarkr\checkmarkc\checkmark$\checkmarkn\checkmark$\checkmark^\checkmark\\checkmarkc\checkmarki\checkmarkr\checkmarkc\checkmark$\checkmarki\checkmark$\checkmark^\checkmark\\checkmarkc\checkmarki\checkmarkr\checkmarkc\checkmark$\checkmarkb\checkmark$\checkmark^\checkmark\\checkmarkc\checkmarki\checkmarkr\checkmarkc\checkmark$\checkmarkl\checkmark$\checkmark^\checkmark\\checkmarkc\checkmarki\checkmarkr\checkmarkc\checkmark$\checkmarke\checkmark$\checkmark^\checkmark\\checkmarkc\checkmarki\checkmarkr\checkmarkc\checkmark$\checkmarks\checkmark$\checkmark^\checkmark\\checkmarkc\checkmarki\checkmarkr\checkmarkc\checkmark$\checkmark.\checkmark$\checkmark^\checkmark\\checkmarkc\checkmarki\checkmarkr\checkmarkc\checkmark$\checkmark
\checkmark$\checkmark^\checkmark\\checkmarkc\checkmarki\checkmarkr\checkmarkc\checkmark$\checkmark \checkmark$\checkmark^\checkmark\\checkmarkc\checkmarki\checkmarkr\checkmarkc\checkmark$\checkmark \checkmark$\checkmark^\checkmark\\checkmarkc\checkmarki\checkmarkr\checkmarkc\checkmark$\checkmark \checkmark$\checkmark^\checkmark\\checkmarkc\checkmarki\checkmarkr\checkmarkc\checkmark$\checkmark \checkmark$\checkmark^\checkmark\\checkmarkc\checkmarki\checkmarkr\checkmarkc\checkmark$\checkmark\\checkmark$\checkmark^\checkmark\\checkmarkc\checkmarki\checkmarkr\checkmarkc\checkmark$\checkmarki\checkmark$\checkmark^\checkmark\\checkmarkc\checkmarki\checkmarkr\checkmarkc\checkmark$\checkmarkt\checkmark$\checkmark^\checkmark\\checkmarkc\checkmarki\checkmarkr\checkmarkc\checkmark$\checkmarke\checkmark$\checkmark^\checkmark\\checkmarkc\checkmarki\checkmarkr\checkmarkc\checkmark$\checkmarkm\checkmark$\checkmark^\checkmark\\checkmarkc\checkmarki\checkmarkr\checkmarkc\checkmark$\checkmark \checkmark$\checkmark^\checkmark\\checkmarkc\checkmarki\checkmarkr\checkmarkc\checkmark$\checkmark\\checkmark$\checkmark^\checkmark\\checkmarkc\checkmarki\checkmarkr\checkmarkc\checkmark$\checkmarkt\checkmark$\checkmark^\checkmark\\checkmarkc\checkmarki\checkmarkr\checkmarkc\checkmark$\checkmarke\checkmark$\checkmark^\checkmark\\checkmarkc\checkmarki\checkmarkr\checkmarkc\checkmark$\checkmarkx\checkmark$\checkmark^\checkmark\\checkmarkc\checkmarki\checkmarkr\checkmarkc\checkmark$\checkmarkt\checkmark$\checkmark^\checkmark\\checkmarkc\checkmarki\checkmarkr\checkmarkc\checkmark$\checkmarkb\checkmark$\checkmark^\checkmark\\checkmarkc\checkmarki\checkmarkr\checkmarkc\checkmark$\checkmarkf\checkmark$\checkmark^\checkmark\\checkmarkc\checkmarki\checkmarkr\checkmarkc\checkmark$\checkmark{\checkmark$\checkmark^\checkmark\\checkmarkc\checkmarki\checkmarkr\checkmarkc\checkmark$\checkmarkF\checkmark$\checkmark^\checkmark\\checkmarkc\checkmarki\checkmarkr\checkmarkc\checkmark$\checkmarkC\checkmark$\checkmark^\checkmark\\checkmarkc\checkmarki\checkmarkr\checkmarkc\checkmark$\checkmark3\checkmark$\checkmark^\checkmark\\checkmarkc\checkmarki\checkmarkr\checkmarkc\checkmark$\checkmark \checkmark$\checkmark^\checkmark\\checkmarkc\checkmarki\checkmarkr\checkmarkc\checkmark$\checkmark:\checkmark$\checkmark^\checkmark\\checkmarkc\checkmarki\checkmarkr\checkmarkc\checkmark$\checkmark}\checkmark$\checkmark^\checkmark\\checkmarkc\checkmarki\checkmarkr\checkmarkc\checkmark$\checkmark \checkmark$\checkmark^\checkmark\\checkmarkc\checkmarki\checkmarkr\checkmarkc\checkmark$\checkmarkL\checkmark$\checkmark^\checkmark\\checkmarkc\checkmarki\checkmarkr\checkmarkc\checkmark$\checkmarke\checkmark$\checkmark^\checkmark\\checkmarkc\checkmarki\checkmarkr\checkmarkc\checkmark$\checkmark \checkmark$\checkmark^\checkmark\\checkmarkc\checkmarki\checkmarkr\checkmarkc\checkmark$\checkmarkc\checkmark$\checkmark^\checkmark\\checkmarkc\checkmarki\checkmarkr\checkmarkc\checkmark$\checkmarko\checkmark$\checkmark^\checkmark\\checkmarkc\checkmarki\checkmarkr\checkmarkc\checkmark$\checkmarke\checkmark$\checkmark^\checkmark\\checkmarkc\checkmarki\checkmarkr\checkmarkc\checkmark$\checkmarkf\checkmark$\checkmark^\checkmark\\checkmarkc\checkmarki\checkmarkr\checkmarkc\checkmark$\checkmarkf\checkmark$\checkmark^\checkmark\\checkmarkc\checkmarki\checkmarkr\checkmarkc\checkmark$\checkmarki\checkmark$\checkmark^\checkmark\\checkmarkc\checkmarki\checkmarkr\checkmarkc\checkmark$\checkmarkc\checkmark$\checkmark^\checkmark\\checkmarkc\checkmarki\checkmarkr\checkmarkc\checkmark$\checkmarki\checkmark$\checkmark^\checkmark\\checkmarkc\checkmarki\checkmarkr\checkmarkc\checkmark$\checkmarke\checkmark$\checkmark^\checkmark\\checkmarkc\checkmarki\checkmarkr\checkmarkc\checkmark$\checkmarkn\checkmark$\checkmark^\checkmark\\checkmarkc\checkmarki\checkmarkr\checkmarkc\checkmark$\checkmarkt\checkmark$\checkmark^\checkmark\\checkmarkc\checkmarki\checkmarkr\checkmarkc\checkmark$\checkmark \checkmark$\checkmark^\checkmark\\checkmarkc\checkmarki\checkmarkr\checkmarkc\checkmark$\checkmarkd\checkmark$\checkmark^\checkmark\\checkmarkc\checkmarki\checkmarkr\checkmarkc\checkmark$\checkmarke\checkmark$\checkmark^\checkmark\\checkmarkc\checkmarki\checkmarkr\checkmarkc\checkmark$\checkmark \checkmark$\checkmark^\checkmark\\checkmarkc\checkmarki\checkmarkr\checkmarkc\checkmark$\checkmarks\checkmark$\checkmark^\checkmark\\checkmarkc\checkmarki\checkmarkr\checkmarkc\checkmark$\checkmarké\checkmark$\checkmark^\checkmark\\checkmarkc\checkmarki\checkmarkr\checkmarkc\checkmark$\checkmarkc\checkmark$\checkmark^\checkmark\\checkmarkc\checkmarki\checkmarkr\checkmarkc\checkmark$\checkmarku\checkmark$\checkmark^\checkmark\\checkmarkc\checkmarki\checkmarkr\checkmarkc\checkmark$\checkmarkr\checkmark$\checkmark^\checkmark\\checkmarkc\checkmarki\checkmarkr\checkmarkc\checkmark$\checkmarki\checkmark$\checkmark^\checkmark\\checkmarkc\checkmarki\checkmarkr\checkmarkc\checkmark$\checkmarkt\checkmark$\checkmark^\checkmark\\checkmarkc\checkmarki\checkmarkr\checkmarkc\checkmark$\checkmarké\checkmark$\checkmark^\checkmark\\checkmarkc\checkmarki\checkmarkr\checkmarkc\checkmark$\checkmark \checkmark$\checkmark^\checkmark\\checkmarkc\checkmarki\checkmarkr\checkmarkc\checkmark$\checkmarks\checkmark$\checkmark^\checkmark\\checkmarkc\checkmarki\checkmarkr\checkmarkc\checkmark$\checkmarkt\checkmark$\checkmark^\checkmark\\checkmarkc\checkmarki\checkmarkr\checkmarkc\checkmark$\checkmarkr\checkmark$\checkmark^\checkmark\\checkmarkc\checkmarki\checkmarkr\checkmarkc\checkmark$\checkmarku\checkmark$\checkmark^\checkmark\\checkmarkc\checkmarki\checkmarkr\checkmarkc\checkmark$\checkmarkc\checkmark$\checkmark^\checkmark\\checkmarkc\checkmarki\checkmarkr\checkmarkc\checkmark$\checkmarkt\checkmark$\checkmark^\checkmark\\checkmarkc\checkmarki\checkmarkr\checkmarkc\checkmark$\checkmarku\checkmark$\checkmark^\checkmark\\checkmarkc\checkmarki\checkmarkr\checkmarkc\checkmark$\checkmarkr\checkmark$\checkmark^\checkmark\\checkmarkc\checkmarki\checkmarkr\checkmarkc\checkmark$\checkmarka\checkmark$\checkmark^\checkmark\\checkmarkc\checkmarki\checkmarkr\checkmarkc\checkmark$\checkmarkl\checkmark$\checkmark^\checkmark\\checkmarkc\checkmarki\checkmarkr\checkmarkc\checkmark$\checkmark \checkmark$\checkmark^\checkmark\\checkmarkc\checkmarki\checkmarkr\checkmarkc\checkmark$\checkmarkd\checkmark$\checkmark^\checkmark\\checkmarkc\checkmarki\checkmarkr\checkmarkc\checkmark$\checkmarko\checkmark$\checkmark^\checkmark\\checkmarkc\checkmarki\checkmarkr\checkmarkc\checkmark$\checkmarki\checkmark$\checkmark^\checkmark\\checkmarkc\checkmarki\checkmarkr\checkmarkc\checkmark$\checkmarkt\checkmark$\checkmark^\checkmark\\checkmarkc\checkmarki\checkmarkr\checkmarkc\checkmark$\checkmark \checkmark$\checkmark^\checkmark\\checkmarkc\checkmarki\checkmarkr\checkmarkc\checkmark$\checkmarks\checkmark$\checkmark^\checkmark\\checkmarkc\checkmarki\checkmarkr\checkmarkc\checkmark$\checkmarka\checkmark$\checkmark^\checkmark\\checkmarkc\checkmarki\checkmarkr\checkmarkc\checkmark$\checkmarkt\checkmark$\checkmark^\checkmark\\checkmarkc\checkmarki\checkmarkr\checkmarkc\checkmark$\checkmarki\checkmark$\checkmark^\checkmark\\checkmarkc\checkmarki\checkmarkr\checkmarkc\checkmark$\checkmarks\checkmark$\checkmark^\checkmark\\checkmarkc\checkmarki\checkmarkr\checkmarkc\checkmark$\checkmarkf\checkmark$\checkmark^\checkmark\\checkmarkc\checkmarki\checkmarkr\checkmarkc\checkmark$\checkmarka\checkmark$\checkmark^\checkmark\\checkmarkc\checkmarki\checkmarkr\checkmarkc\checkmark$\checkmarki\checkmark$\checkmark^\checkmark\\checkmarkc\checkmarki\checkmarkr\checkmarkc\checkmark$\checkmarkr\checkmark$\checkmark^\checkmark\\checkmarkc\checkmarki\checkmarkr\checkmarkc\checkmark$\checkmarke\checkmark$\checkmark^\checkmark\\checkmarkc\checkmarki\checkmarkr\checkmarkc\checkmark$\checkmark \checkmark$\checkmark^\checkmark\\checkmarkc\checkmarki\checkmarkr\checkmarkc\checkmark$\checkmark$\checkmark$\checkmark^\checkmark\\checkmarkc\checkmarki\checkmarkr\checkmarkc\checkmark$\checkmarks\checkmark$\checkmark^\checkmark\\checkmarkc\checkmarki\checkmarkr\checkmarkc\checkmark$\checkmark \checkmark$\checkmark^\checkmark\\checkmarkc\checkmarki\checkmarkr\checkmarkc\checkmark$\checkmark\\checkmark$\checkmark^\checkmark\\checkmarkc\checkmarki\checkmarkr\checkmarkc\checkmark$\checkmarkg\checkmark$\checkmark^\checkmark\\checkmarkc\checkmarki\checkmarkr\checkmarkc\checkmark$\checkmarke\checkmark$\checkmark^\checkmark\\checkmarkc\checkmarki\checkmarkr\checkmarkc\checkmark$\checkmarkq\checkmark$\checkmark^\checkmark\\checkmarkc\checkmarki\checkmarkr\checkmarkc\checkmark$\checkmark \checkmark$\checkmark^\checkmark\\checkmarkc\checkmarki\checkmarkr\checkmarkc\checkmark$\checkmark1\checkmark$\checkmark^\checkmark\\checkmarkc\checkmarki\checkmarkr\checkmarkc\checkmark$\checkmark{\checkmark$\checkmark^\checkmark\\checkmarkc\checkmarki\checkmarkr\checkmarkc\checkmark$\checkmark,\checkmark$\checkmark^\checkmark\\checkmarkc\checkmarki\checkmarkr\checkmarkc\checkmark$\checkmark}\checkmark$\checkmark^\checkmark\\checkmarkc\checkmarki\checkmarkr\checkmarkc\checkmark$\checkmark5\checkmark$\checkmark^\checkmark\\checkmarkc\checkmarki\checkmarkr\checkmarkc\checkmark$\checkmark$\checkmark$\checkmark^\checkmark\\checkmarkc\checkmarki\checkmarkr\checkmarkc\checkmark$\checkmark.\checkmark$\checkmark^\checkmark\\checkmarkc\checkmarki\checkmarkr\checkmarkc\checkmark$\checkmark
\checkmark$\checkmark^\checkmark\\checkmarkc\checkmarki\checkmarkr\checkmarkc\checkmark$\checkmark \checkmark$\checkmark^\checkmark\\checkmarkc\checkmarki\checkmarkr\checkmarkc\checkmark$\checkmark \checkmark$\checkmark^\checkmark\\checkmarkc\checkmarki\checkmarkr\checkmarkc\checkmark$\checkmark \checkmark$\checkmark^\checkmark\\checkmarkc\checkmarki\checkmarkr\checkmarkc\checkmark$\checkmark \checkmark$\checkmark^\checkmark\\checkmarkc\checkmarki\checkmarkr\checkmarkc\checkmark$\checkmark\\checkmark$\checkmark^\checkmark\\checkmarkc\checkmarki\checkmarkr\checkmarkc\checkmark$\checkmarki\checkmark$\checkmark^\checkmark\\checkmarkc\checkmarki\checkmarkr\checkmarkc\checkmark$\checkmarkt\checkmark$\checkmark^\checkmark\\checkmarkc\checkmarki\checkmarkr\checkmarkc\checkmark$\checkmarke\checkmark$\checkmark^\checkmark\\checkmarkc\checkmarki\checkmarkr\checkmarkc\checkmark$\checkmarkm\checkmark$\checkmark^\checkmark\\checkmarkc\checkmarki\checkmarkr\checkmarkc\checkmark$\checkmark \checkmark$\checkmark^\checkmark\\checkmarkc\checkmarki\checkmarkr\checkmarkc\checkmark$\checkmark\\checkmark$\checkmark^\checkmark\\checkmarkc\checkmarki\checkmarkr\checkmarkc\checkmark$\checkmarkt\checkmark$\checkmark^\checkmark\\checkmarkc\checkmarki\checkmarkr\checkmarkc\checkmark$\checkmarke\checkmark$\checkmark^\checkmark\\checkmarkc\checkmarki\checkmarkr\checkmarkc\checkmark$\checkmarkx\checkmark$\checkmark^\checkmark\\checkmarkc\checkmarki\checkmarkr\checkmarkc\checkmark$\checkmarkt\checkmark$\checkmark^\checkmark\\checkmarkc\checkmarki\checkmarkr\checkmarkc\checkmark$\checkmarkb\checkmark$\checkmark^\checkmark\\checkmarkc\checkmarki\checkmarkr\checkmarkc\checkmark$\checkmarkf\checkmark$\checkmark^\checkmark\\checkmarkc\checkmarki\checkmarkr\checkmarkc\checkmark$\checkmark{\checkmark$\checkmark^\checkmark\\checkmarkc\checkmarki\checkmarkr\checkmarkc\checkmark$\checkmarkF\checkmark$\checkmark^\checkmark\\checkmarkc\checkmarki\checkmarkr\checkmarkc\checkmark$\checkmarkC\checkmark$\checkmark^\checkmark\\checkmarkc\checkmarki\checkmarkr\checkmarkc\checkmark$\checkmark4\checkmark$\checkmark^\checkmark\\checkmarkc\checkmarki\checkmarkr\checkmarkc\checkmark$\checkmark \checkmark$\checkmark^\checkmark\\checkmarkc\checkmarki\checkmarkr\checkmarkc\checkmark$\checkmark:\checkmark$\checkmark^\checkmark\\checkmarkc\checkmarki\checkmarkr\checkmarkc\checkmark$\checkmark}\checkmark$\checkmark^\checkmark\\checkmarkc\checkmarki\checkmarkr\checkmarkc\checkmark$\checkmark \checkmark$\checkmark^\checkmark\\checkmarkc\checkmarki\checkmarkr\checkmarkc\checkmark$\checkmarkL\checkmark$\checkmark^\checkmark\\checkmarkc\checkmarki\checkmarkr\checkmarkc\checkmark$\checkmark'\checkmark$\checkmark^\checkmark\\checkmarkc\checkmarki\checkmarkr\checkmarkc\checkmark$\checkmarki\checkmark$\checkmark^\checkmark\\checkmarkc\checkmarki\checkmarkr\checkmarkc\checkmark$\checkmarkn\checkmark$\checkmark^\checkmark\\checkmarkc\checkmarki\checkmarkr\checkmarkc\checkmark$\checkmarkt\checkmark$\checkmark^\checkmark\\checkmarkc\checkmarki\checkmarkr\checkmarkc\checkmark$\checkmarke\checkmark$\checkmark^\checkmark\\checkmarkc\checkmarki\checkmarkr\checkmarkc\checkmark$\checkmarkr\checkmark$\checkmark^\checkmark\\checkmarkc\checkmarki\checkmarkr\checkmarkc\checkmark$\checkmarkf\checkmark$\checkmark^\checkmark\\checkmarkc\checkmarki\checkmarkr\checkmarkc\checkmark$\checkmarka\checkmark$\checkmark^\checkmark\\checkmarkc\checkmarki\checkmarkr\checkmarkc\checkmark$\checkmarkc\checkmark$\checkmark^\checkmark\\checkmarkc\checkmarki\checkmarkr\checkmarkc\checkmark$\checkmarke\checkmark$\checkmark^\checkmark\\checkmarkc\checkmarki\checkmarkr\checkmarkc\checkmark$\checkmark \checkmark$\checkmark^\checkmark\\checkmarkc\checkmarki\checkmarkr\checkmarkc\checkmark$\checkmarku\checkmark$\checkmark^\checkmark\\checkmarkc\checkmarki\checkmarkr\checkmarkc\checkmark$\checkmarkt\checkmark$\checkmark^\checkmark\\checkmarkc\checkmarki\checkmarkr\checkmarkc\checkmark$\checkmarki\checkmark$\checkmark^\checkmark\\checkmarkc\checkmarki\checkmarkr\checkmarkc\checkmark$\checkmarkl\checkmark$\checkmark^\checkmark\\checkmarkc\checkmarki\checkmarkr\checkmarkc\checkmark$\checkmarki\checkmark$\checkmark^\checkmark\\checkmarkc\checkmarki\checkmarkr\checkmarkc\checkmark$\checkmarks\checkmark$\checkmark^\checkmark\\checkmarkc\checkmarki\checkmarkr\checkmarkc\checkmark$\checkmarka\checkmark$\checkmark^\checkmark\\checkmarkc\checkmarki\checkmarkr\checkmarkc\checkmark$\checkmarkt\checkmark$\checkmark^\checkmark\\checkmarkc\checkmarki\checkmarkr\checkmarkc\checkmark$\checkmarke\checkmark$\checkmark^\checkmark\\checkmarkc\checkmarki\checkmarkr\checkmarkc\checkmark$\checkmarku\checkmark$\checkmark^\checkmark\\checkmarkc\checkmarki\checkmarkr\checkmarkc\checkmark$\checkmarkr\checkmark$\checkmark^\checkmark\\checkmarkc\checkmarki\checkmarkr\checkmarkc\checkmark$\checkmark \checkmark$\checkmark^\checkmark\\checkmarkc\checkmarki\checkmarkr\checkmarkc\checkmark$\checkmarkd\checkmark$\checkmark^\checkmark\\checkmarkc\checkmarki\checkmarkr\checkmarkc\checkmark$\checkmarko\checkmark$\checkmark^\checkmark\\checkmarkc\checkmarki\checkmarkr\checkmarkc\checkmark$\checkmarki\checkmark$\checkmark^\checkmark\\checkmarkc\checkmarki\checkmarkr\checkmarkc\checkmark$\checkmarkt\checkmark$\checkmark^\checkmark\\checkmarkc\checkmarki\checkmarkr\checkmarkc\checkmark$\checkmark \checkmark$\checkmark^\checkmark\\checkmarkc\checkmarki\checkmarkr\checkmarkc\checkmark$\checkmarkê\checkmark$\checkmark^\checkmark\\checkmarkc\checkmarki\checkmarkr\checkmarkc\checkmark$\checkmarkt\checkmark$\checkmark^\checkmark\\checkmarkc\checkmarki\checkmarkr\checkmarkc\checkmark$\checkmarkr\checkmark$\checkmark^\checkmark\\checkmarkc\checkmarki\checkmarkr\checkmarkc\checkmark$\checkmarke\checkmark$\checkmark^\checkmark\\checkmarkc\checkmarki\checkmarkr\checkmarkc\checkmark$\checkmark \checkmark$\checkmark^\checkmark\\checkmarkc\checkmarki\checkmarkr\checkmarkc\checkmark$\checkmarki\checkmark$\checkmark^\checkmark\\checkmarkc\checkmarki\checkmarkr\checkmarkc\checkmark$\checkmarkn\checkmark$\checkmark^\checkmark\\checkmarkc\checkmarki\checkmarkr\checkmarkc\checkmark$\checkmarkt\checkmark$\checkmark^\checkmark\\checkmarkc\checkmarki\checkmarkr\checkmarkc\checkmark$\checkmarku\checkmark$\checkmark^\checkmark\\checkmarkc\checkmarki\checkmarkr\checkmarkc\checkmark$\checkmarki\checkmark$\checkmark^\checkmark\\checkmarkc\checkmarki\checkmarkr\checkmarkc\checkmark$\checkmarkt\checkmark$\checkmark^\checkmark\\checkmarkc\checkmarki\checkmarkr\checkmarkc\checkmark$\checkmarki\checkmark$\checkmark^\checkmark\\checkmarkc\checkmarki\checkmarkr\checkmarkc\checkmark$\checkmarkv\checkmark$\checkmark^\checkmark\\checkmarkc\checkmarki\checkmarkr\checkmarkc\checkmark$\checkmarke\checkmark$\checkmark^\checkmark\\checkmarkc\checkmarki\checkmarkr\checkmarkc\checkmark$\checkmark \checkmark$\checkmark^\checkmark\\checkmarkc\checkmarki\checkmarkr\checkmarkc\checkmark$\checkmarke\checkmark$\checkmark^\checkmark\\checkmarkc\checkmarki\checkmarkr\checkmarkc\checkmark$\checkmarkt\checkmark$\checkmark^\checkmark\\checkmarkc\checkmarki\checkmarkr\checkmarkc\checkmark$\checkmark \checkmark$\checkmark^\checkmark\\checkmarkc\checkmarki\checkmarkr\checkmarkc\checkmark$\checkmarkm\checkmark$\checkmark^\checkmark\\checkmarkc\checkmarki\checkmarkr\checkmarkc\checkmark$\checkmarku\checkmark$\checkmark^\checkmark\\checkmarkc\checkmarki\checkmarkr\checkmarkc\checkmark$\checkmarkl\checkmark$\checkmark^\checkmark\\checkmarkc\checkmarki\checkmarkr\checkmarkc\checkmark$\checkmarkt\checkmark$\checkmark^\checkmark\\checkmarkc\checkmarki\checkmarkr\checkmarkc\checkmark$\checkmarki\checkmark$\checkmark^\checkmark\\checkmarkc\checkmarki\checkmarkr\checkmarkc\checkmark$\checkmarkp\checkmark$\checkmark^\checkmark\\checkmarkc\checkmarki\checkmarkr\checkmarkc\checkmark$\checkmarkl\checkmark$\checkmark^\checkmark\\checkmarkc\checkmarki\checkmarkr\checkmarkc\checkmark$\checkmarka\checkmark$\checkmark^\checkmark\\checkmarkc\checkmarki\checkmarkr\checkmarkc\checkmark$\checkmarkt\checkmark$\checkmark^\checkmark\\checkmarkc\checkmarki\checkmarkr\checkmarkc\checkmark$\checkmarke\checkmark$\checkmark^\checkmark\\checkmarkc\checkmarki\checkmarkr\checkmarkc\checkmark$\checkmarkf\checkmark$\checkmark^\checkmark\\checkmarkc\checkmarki\checkmarkr\checkmarkc\checkmark$\checkmarko\checkmark$\checkmark^\checkmark\\checkmarkc\checkmarki\checkmarkr\checkmarkc\checkmark$\checkmarkr\checkmark$\checkmark^\checkmark\\checkmarkc\checkmarki\checkmarkr\checkmarkc\checkmark$\checkmarkm\checkmark$\checkmark^\checkmark\\checkmarkc\checkmarki\checkmarkr\checkmarkc\checkmark$\checkmarke\checkmark$\checkmark^\checkmark\\checkmarkc\checkmarki\checkmarkr\checkmarkc\checkmark$\checkmark \checkmark$\checkmark^\checkmark\\checkmarkc\checkmarki\checkmarkr\checkmarkc\checkmark$\checkmark(\checkmark$\checkmark^\checkmark\\checkmarkc\checkmarki\checkmarkr\checkmarkc\checkmark$\checkmarkA\checkmark$\checkmark^\checkmark\\checkmarkc\checkmarki\checkmarkr\checkmarkc\checkmark$\checkmarkn\checkmark$\checkmark^\checkmark\\checkmarkc\checkmarki\checkmarkr\checkmarkc\checkmark$\checkmarkd\checkmark$\checkmark^\checkmark\\checkmarkc\checkmarki\checkmarkr\checkmarkc\checkmark$\checkmarkr\checkmark$\checkmark^\checkmark\\checkmarkc\checkmarki\checkmarkr\checkmarkc\checkmark$\checkmarko\checkmark$\checkmark^\checkmark\\checkmarkc\checkmarki\checkmarkr\checkmarkc\checkmark$\checkmarki\checkmark$\checkmark^\checkmark\\checkmarkc\checkmarki\checkmarkr\checkmarkc\checkmark$\checkmarkd\checkmark$\checkmark^\checkmark\\checkmarkc\checkmarki\checkmarkr\checkmarkc\checkmark$\checkmark \checkmark$\checkmark^\checkmark\\checkmarkc\checkmarki\checkmarkr\checkmarkc\checkmark$\checkmark/\checkmark$\checkmark^\checkmark\\checkmarkc\checkmarki\checkmarkr\checkmarkc\checkmark$\checkmark \checkmark$\checkmark^\checkmark\\checkmarkc\checkmarki\checkmarkr\checkmarkc\checkmark$\checkmarkW\checkmark$\checkmark^\checkmark\\checkmarkc\checkmarki\checkmarkr\checkmarkc\checkmark$\checkmarki\checkmark$\checkmark^\checkmark\\checkmarkc\checkmarki\checkmarkr\checkmarkc\checkmark$\checkmarkn\checkmark$\checkmark^\checkmark\\checkmarkc\checkmarki\checkmarkr\checkmarkc\checkmark$\checkmarkd\checkmark$\checkmark^\checkmark\\checkmarkc\checkmarki\checkmarkr\checkmarkc\checkmark$\checkmarko\checkmark$\checkmark^\checkmark\\checkmarkc\checkmarki\checkmarkr\checkmarkc\checkmark$\checkmarkw\checkmark$\checkmark^\checkmark\\checkmarkc\checkmarki\checkmarkr\checkmarkc\checkmark$\checkmarks\checkmark$\checkmark^\checkmark\\checkmarkc\checkmarki\checkmarkr\checkmarkc\checkmark$\checkmark)\checkmark$\checkmark^\checkmark\\checkmarkc\checkmarki\checkmarkr\checkmarkc\checkmark$\checkmark.\checkmark$\checkmark^\checkmark\\checkmarkc\checkmarki\checkmarkr\checkmarkc\checkmark$\checkmark
\checkmark$\checkmark^\checkmark\\checkmarkc\checkmarki\checkmarkr\checkmarkc\checkmark$\checkmark \checkmark$\checkmark^\checkmark\\checkmarkc\checkmarki\checkmarkr\checkmarkc\checkmark$\checkmark \checkmark$\checkmark^\checkmark\\checkmarkc\checkmarki\checkmarkr\checkmarkc\checkmark$\checkmark \checkmark$\checkmark^\checkmark\\checkmarkc\checkmarki\checkmarkr\checkmarkc\checkmark$\checkmark \checkmark$\checkmark^\checkmark\\checkmarkc\checkmarki\checkmarkr\checkmarkc\checkmark$\checkmark\\checkmark$\checkmark^\checkmark\\checkmarkc\checkmarki\checkmarkr\checkmarkc\checkmark$\checkmarki\checkmark$\checkmark^\checkmark\\checkmarkc\checkmarki\checkmarkr\checkmarkc\checkmark$\checkmarkt\checkmark$\checkmark^\checkmark\\checkmarkc\checkmarki\checkmarkr\checkmarkc\checkmark$\checkmarke\checkmark$\checkmark^\checkmark\\checkmarkc\checkmarki\checkmarkr\checkmarkc\checkmark$\checkmarkm\checkmark$\checkmark^\checkmark\\checkmarkc\checkmarki\checkmarkr\checkmarkc\checkmark$\checkmark \checkmark$\checkmark^\checkmark\\checkmarkc\checkmarki\checkmarkr\checkmarkc\checkmark$\checkmark\\checkmark$\checkmark^\checkmark\\checkmarkc\checkmarki\checkmarkr\checkmarkc\checkmark$\checkmarkt\checkmark$\checkmark^\checkmark\\checkmarkc\checkmarki\checkmarkr\checkmarkc\checkmark$\checkmarke\checkmark$\checkmark^\checkmark\\checkmarkc\checkmarki\checkmarkr\checkmarkc\checkmark$\checkmarkx\checkmark$\checkmark^\checkmark\\checkmarkc\checkmarki\checkmarkr\checkmarkc\checkmark$\checkmarkt\checkmark$\checkmark^\checkmark\\checkmarkc\checkmarki\checkmarkr\checkmarkc\checkmark$\checkmarkb\checkmark$\checkmark^\checkmark\\checkmarkc\checkmarki\checkmarkr\checkmarkc\checkmark$\checkmarkf\checkmark$\checkmark^\checkmark\\checkmarkc\checkmarki\checkmarkr\checkmarkc\checkmark$\checkmark{\checkmark$\checkmark^\checkmark\\checkmarkc\checkmarki\checkmarkr\checkmarkc\checkmark$\checkmarkF\checkmark$\checkmark^\checkmark\\checkmarkc\checkmarki\checkmarkr\checkmarkc\checkmark$\checkmarkC\checkmark$\checkmark^\checkmark\\checkmarkc\checkmarki\checkmarkr\checkmarkc\checkmark$\checkmark5\checkmark$\checkmark^\checkmark\\checkmarkc\checkmarki\checkmarkr\checkmarkc\checkmark$\checkmark \checkmark$\checkmark^\checkmark\\checkmarkc\checkmarki\checkmarkr\checkmarkc\checkmark$\checkmark:\checkmark$\checkmark^\checkmark\\checkmarkc\checkmarki\checkmarkr\checkmarkc\checkmark$\checkmark}\checkmark$\checkmark^\checkmark\\checkmarkc\checkmarki\checkmarkr\checkmarkc\checkmark$\checkmark \checkmark$\checkmark^\checkmark\\checkmarkc\checkmarki\checkmarkr\checkmarkc\checkmark$\checkmarkL\checkmark$\checkmark^\checkmark\\checkmarkc\checkmarki\checkmarkr\checkmarkc\checkmark$\checkmarke\checkmark$\checkmark^\checkmark\\checkmarkc\checkmarki\checkmarkr\checkmarkc\checkmark$\checkmark \checkmark$\checkmark^\checkmark\\checkmarkc\checkmarki\checkmarkr\checkmarkc\checkmark$\checkmarks\checkmark$\checkmark^\checkmark\\checkmarkc\checkmarki\checkmarkr\checkmarkc\checkmark$\checkmarky\checkmark$\checkmark^\checkmark\\checkmarkc\checkmarki\checkmarkr\checkmarkc\checkmark$\checkmarks\checkmark$\checkmark^\checkmark\\checkmarkc\checkmarki\checkmarkr\checkmarkc\checkmark$\checkmarkt\checkmark$\checkmark^\checkmark\\checkmarkc\checkmarki\checkmarkr\checkmarkc\checkmark$\checkmarkè\checkmark$\checkmark^\checkmark\\checkmarkc\checkmarki\checkmarkr\checkmarkc\checkmark$\checkmarkm\checkmark$\checkmark^\checkmark\\checkmarkc\checkmarki\checkmarkr\checkmarkc\checkmark$\checkmarke\checkmark$\checkmark^\checkmark\\checkmarkc\checkmarki\checkmarkr\checkmarkc\checkmark$\checkmark \checkmark$\checkmark^\checkmark\\checkmarkc\checkmarki\checkmarkr\checkmarkc\checkmark$\checkmarkd\checkmark$\checkmark^\checkmark\\checkmarkc\checkmarki\checkmarkr\checkmarkc\checkmark$\checkmarko\checkmark$\checkmark^\checkmark\\checkmarkc\checkmarki\checkmarkr\checkmarkc\checkmark$\checkmarki\checkmark$\checkmark^\checkmark\\checkmarkc\checkmarki\checkmarkr\checkmarkc\checkmark$\checkmarkt\checkmark$\checkmark^\checkmark\\checkmarkc\checkmarki\checkmarkr\checkmarkc\checkmark$\checkmark \checkmark$\checkmark^\checkmark\\checkmarkc\checkmarki\checkmarkr\checkmarkc\checkmark$\checkmarkê\checkmark$\checkmark^\checkmark\\checkmarkc\checkmarki\checkmarkr\checkmarkc\checkmark$\checkmarkt\checkmark$\checkmark^\checkmark\\checkmarkc\checkmarki\checkmarkr\checkmarkc\checkmark$\checkmarkr\checkmark$\checkmark^\checkmark\\checkmarkc\checkmarki\checkmarkr\checkmarkc\checkmark$\checkmarke\checkmark$\checkmark^\checkmark\\checkmarkc\checkmarki\checkmarkr\checkmarkc\checkmark$\checkmark \checkmark$\checkmark^\checkmark\\checkmarkc\checkmarki\checkmarkr\checkmarkc\checkmark$\checkmarkc\checkmark$\checkmark^\checkmark\\checkmarkc\checkmarki\checkmarkr\checkmarkc\checkmark$\checkmarko\checkmark$\checkmark^\checkmark\\checkmarkc\checkmarki\checkmarkr\checkmarkc\checkmark$\checkmarkm\checkmark$\checkmark^\checkmark\\checkmarkc\checkmarki\checkmarkr\checkmarkc\checkmark$\checkmarkp\checkmark$\checkmark^\checkmark\\checkmarkc\checkmarki\checkmarkr\checkmarkc\checkmark$\checkmarka\checkmark$\checkmark^\checkmark\\checkmarkc\checkmarki\checkmarkr\checkmarkc\checkmark$\checkmarkc\checkmark$\checkmark^\checkmark\\checkmarkc\checkmarki\checkmarkr\checkmarkc\checkmark$\checkmarkt\checkmark$\checkmark^\checkmark\\checkmarkc\checkmarki\checkmarkr\checkmarkc\checkmark$\checkmark \checkmark$\checkmark^\checkmark\\checkmarkc\checkmarki\checkmarkr\checkmarkc\checkmark$\checkmarke\checkmark$\checkmark^\checkmark\\checkmarkc\checkmarki\checkmarkr\checkmarkc\checkmark$\checkmarkt\checkmark$\checkmark^\checkmark\\checkmarkc\checkmarki\checkmarkr\checkmarkc\checkmark$\checkmark \checkmark$\checkmark^\checkmark\\checkmarkc\checkmarki\checkmarkr\checkmarkc\checkmark$\checkmarkt\checkmark$\checkmark^\checkmark\\checkmarkc\checkmarki\checkmarkr\checkmarkc\checkmark$\checkmarkr\checkmark$\checkmark^\checkmark\\checkmarkc\checkmarki\checkmarkr\checkmarkc\checkmark$\checkmarka\checkmark$\checkmark^\checkmark\\checkmarkc\checkmarki\checkmarkr\checkmarkc\checkmark$\checkmarkn\checkmark$\checkmark^\checkmark\\checkmarkc\checkmarki\checkmarkr\checkmarkc\checkmark$\checkmarks\checkmark$\checkmark^\checkmark\\checkmarkc\checkmarki\checkmarkr\checkmarkc\checkmark$\checkmarkp\checkmark$\checkmark^\checkmark\\checkmarkc\checkmarki\checkmarkr\checkmarkc\checkmark$\checkmarko\checkmark$\checkmark^\checkmark\\checkmarkc\checkmarki\checkmarkr\checkmarkc\checkmark$\checkmarkr\checkmark$\checkmark^\checkmark\\checkmarkc\checkmarki\checkmarkr\checkmarkc\checkmark$\checkmarkt\checkmark$\checkmark^\checkmark\\checkmarkc\checkmarki\checkmarkr\checkmarkc\checkmark$\checkmarka\checkmark$\checkmark^\checkmark\\checkmarkc\checkmarki\checkmarkr\checkmarkc\checkmark$\checkmarkb\checkmark$\checkmark^\checkmark\\checkmarkc\checkmarki\checkmarkr\checkmarkc\checkmark$\checkmarkl\checkmark$\checkmark^\checkmark\\checkmarkc\checkmarki\checkmarkr\checkmarkc\checkmark$\checkmarke\checkmark$\checkmark^\checkmark\\checkmarkc\checkmarki\checkmarkr\checkmarkc\checkmark$\checkmark \checkmark$\checkmark^\checkmark\\checkmarkc\checkmarki\checkmarkr\checkmarkc\checkmark$\checkmark(\checkmark$\checkmark^\checkmark\\checkmarkc\checkmarki\checkmarkr\checkmarkc\checkmark$\checkmarkv\checkmark$\checkmark^\checkmark\\checkmarkc\checkmarki\checkmarkr\checkmarkc\checkmark$\checkmarko\checkmark$\checkmark^\checkmark\\checkmarkc\checkmarki\checkmarkr\checkmarkc\checkmark$\checkmarkl\checkmark$\checkmark^\checkmark\\checkmarkc\checkmarki\checkmarkr\checkmarkc\checkmark$\checkmarku\checkmark$\checkmark^\checkmark\\checkmarkc\checkmarki\checkmarkr\checkmarkc\checkmark$\checkmarkm\checkmark$\checkmark^\checkmark\\checkmarkc\checkmarki\checkmarkr\checkmarkc\checkmark$\checkmarke\checkmark$\checkmark^\checkmark\\checkmarkc\checkmarki\checkmarkr\checkmarkc\checkmark$\checkmark \checkmark$\checkmark^\checkmark\\checkmarkc\checkmarki\checkmarkr\checkmarkc\checkmark$\checkmark$\checkmark$\checkmark^\checkmark\\checkmarkc\checkmarki\checkmarkr\checkmarkc\checkmark$\checkmark<\checkmark$\checkmark^\checkmark\\checkmarkc\checkmarki\checkmarkr\checkmarkc\checkmark$\checkmark \checkmark$\checkmark^\checkmark\\checkmarkc\checkmarki\checkmarkr\checkmarkc\checkmark$\checkmark6\checkmark$\checkmark^\checkmark\\checkmarkc\checkmarki\checkmarkr\checkmarkc\checkmark$\checkmark0\checkmark$\checkmark^\checkmark\\checkmarkc\checkmarki\checkmarkr\checkmarkc\checkmark$\checkmark0\checkmark$\checkmark^\checkmark\\checkmarkc\checkmarki\checkmarkr\checkmarkc\checkmark$\checkmark \checkmark$\checkmark^\checkmark\\checkmarkc\checkmarki\checkmarkr\checkmarkc\checkmark$\checkmark\\checkmark$\checkmark^\checkmark\\checkmarkc\checkmarki\checkmarkr\checkmarkc\checkmark$\checkmarkt\checkmark$\checkmark^\checkmark\\checkmarkc\checkmarki\checkmarkr\checkmarkc\checkmark$\checkmarki\checkmark$\checkmark^\checkmark\\checkmarkc\checkmarki\checkmarkr\checkmarkc\checkmark$\checkmarkm\checkmark$\checkmark^\checkmark\\checkmarkc\checkmarki\checkmarkr\checkmarkc\checkmark$\checkmarke\checkmark$\checkmark^\checkmark\\checkmarkc\checkmarki\checkmarkr\checkmarkc\checkmark$\checkmarks\checkmark$\checkmark^\checkmark\\checkmarkc\checkmarki\checkmarkr\checkmarkc\checkmark$\checkmark \checkmark$\checkmark^\checkmark\\checkmarkc\checkmarki\checkmarkr\checkmarkc\checkmark$\checkmark6\checkmark$\checkmark^\checkmark\\checkmarkc\checkmarki\checkmarkr\checkmarkc\checkmark$\checkmark0\checkmark$\checkmark^\checkmark\\checkmarkc\checkmarki\checkmarkr\checkmarkc\checkmark$\checkmark0\checkmark$\checkmark^\checkmark\\checkmarkc\checkmarki\checkmarkr\checkmarkc\checkmark$\checkmark \checkmark$\checkmark^\checkmark\\checkmarkc\checkmarki\checkmarkr\checkmarkc\checkmark$\checkmark\\checkmark$\checkmark^\checkmark\\checkmarkc\checkmarki\checkmarkr\checkmarkc\checkmark$\checkmarkt\checkmark$\checkmark^\checkmark\\checkmarkc\checkmarki\checkmarkr\checkmarkc\checkmark$\checkmarki\checkmark$\checkmark^\checkmark\\checkmarkc\checkmarki\checkmarkr\checkmarkc\checkmark$\checkmarkm\checkmark$\checkmark^\checkmark\\checkmarkc\checkmarki\checkmarkr\checkmarkc\checkmark$\checkmarke\checkmark$\checkmark^\checkmark\\checkmarkc\checkmarki\checkmarkr\checkmarkc\checkmark$\checkmarks\checkmark$\checkmark^\checkmark\\checkmarkc\checkmarki\checkmarkr\checkmarkc\checkmark$\checkmark \checkmark$\checkmark^\checkmark\\checkmarkc\checkmarki\checkmarkr\checkmarkc\checkmark$\checkmark6\checkmark$\checkmark^\checkmark\\checkmarkc\checkmarki\checkmarkr\checkmarkc\checkmark$\checkmark0\checkmark$\checkmark^\checkmark\\checkmarkc\checkmarki\checkmarkr\checkmarkc\checkmark$\checkmark0\checkmark$\checkmark^\checkmark\\checkmarkc\checkmarki\checkmarkr\checkmarkc\checkmark$\checkmark$\checkmark$\checkmark^\checkmark\\checkmarkc\checkmarki\checkmarkr\checkmarkc\checkmark$\checkmark \checkmark$\checkmark^\checkmark\\checkmarkc\checkmarki\checkmarkr\checkmarkc\checkmark$\checkmarkm\checkmark$\checkmark^\checkmark\\checkmarkc\checkmarki\checkmarkr\checkmarkc\checkmark$\checkmarkm\checkmark$\checkmark^\checkmark\\checkmarkc\checkmarki\checkmarkr\checkmarkc\checkmark$\checkmark)\checkmark$\checkmark^\checkmark\\checkmarkc\checkmarki\checkmarkr\checkmarkc\checkmark$\checkmark.\checkmark$\checkmark^\checkmark\\checkmarkc\checkmarki\checkmarkr\checkmarkc\checkmark$\checkmark
\checkmark$\checkmark^\checkmark\\checkmarkc\checkmarki\checkmarkr\checkmarkc\checkmark$\checkmark\\checkmark$\checkmark^\checkmark\\checkmarkc\checkmarki\checkmarkr\checkmarkc\checkmark$\checkmarke\checkmark$\checkmark^\checkmark\\checkmarkc\checkmarki\checkmarkr\checkmarkc\checkmark$\checkmarkn\checkmark$\checkmark^\checkmark\\checkmarkc\checkmarki\checkmarkr\checkmarkc\checkmark$\checkmarkd\checkmark$\checkmark^\checkmark\\checkmarkc\checkmarki\checkmarkr\checkmarkc\checkmark$\checkmark{\checkmark$\checkmark^\checkmark\\checkmarkc\checkmarki\checkmarkr\checkmarkc\checkmark$\checkmarki\checkmark$\checkmark^\checkmark\\checkmarkc\checkmarki\checkmarkr\checkmarkc\checkmark$\checkmarkt\checkmark$\checkmark^\checkmark\\checkmarkc\checkmarki\checkmarkr\checkmarkc\checkmark$\checkmarke\checkmark$\checkmark^\checkmark\\checkmarkc\checkmarki\checkmarkr\checkmarkc\checkmark$\checkmarkm\checkmark$\checkmark^\checkmark\\checkmarkc\checkmarki\checkmarkr\checkmarkc\checkmark$\checkmarki\checkmark$\checkmark^\checkmark\\checkmarkc\checkmarki\checkmarkr\checkmarkc\checkmark$\checkmarkz\checkmark$\checkmark^\checkmark\\checkmarkc\checkmarki\checkmarkr\checkmarkc\checkmark$\checkmarke\checkmark$\checkmark^\checkmark\\checkmarkc\checkmarki\checkmarkr\checkmarkc\checkmark$\checkmark}\checkmark$\checkmark^\checkmark\\checkmarkc\checkmarki\checkmarkr\checkmarkc\checkmark$\checkmark
\checkmark$\checkmark^\checkmark\\checkmarkc\checkmarki\checkmarkr\checkmarkc\checkmark$\checkmark
\checkmark$\checkmark^\checkmark\\checkmarkc\checkmarki\checkmarkr\checkmarkc\checkmark$\checkmark\\checkmark$\checkmark^\checkmark\\checkmarkc\checkmarki\checkmarkr\checkmarkc\checkmark$\checkmarks\checkmark$\checkmark^\checkmark\\checkmarkc\checkmarki\checkmarkr\checkmarkc\checkmark$\checkmarku\checkmark$\checkmark^\checkmark\\checkmarkc\checkmarki\checkmarkr\checkmarkc\checkmark$\checkmarkb\checkmark$\checkmark^\checkmark\\checkmarkc\checkmarki\checkmarkr\checkmarkc\checkmark$\checkmarks\checkmark$\checkmark^\checkmark\\checkmarkc\checkmarki\checkmarkr\checkmarkc\checkmark$\checkmarke\checkmark$\checkmark^\checkmark\\checkmarkc\checkmarki\checkmarkr\checkmarkc\checkmark$\checkmarkc\checkmark$\checkmark^\checkmark\\checkmarkc\checkmarki\checkmarkr\checkmarkc\checkmark$\checkmarkt\checkmark$\checkmark^\checkmark\\checkmarkc\checkmarki\checkmarkr\checkmarkc\checkmark$\checkmarki\checkmark$\checkmark^\checkmark\\checkmarkc\checkmarki\checkmarkr\checkmarkc\checkmark$\checkmarko\checkmark$\checkmark^\checkmark\\checkmarkc\checkmarki\checkmarkr\checkmarkc\checkmark$\checkmarkn\checkmark$\checkmark^\checkmark\\checkmarkc\checkmarki\checkmarkr\checkmarkc\checkmark$\checkmark{\checkmark$\checkmark^\checkmark\\checkmarkc\checkmarki\checkmarkr\checkmarkc\checkmark$\checkmarkC\checkmark$\checkmark^\checkmark\\checkmarkc\checkmarki\checkmarkr\checkmarkc\checkmark$\checkmarka\checkmark$\checkmark^\checkmark\\checkmarkc\checkmarki\checkmarkr\checkmarkc\checkmark$\checkmarkh\checkmark$\checkmark^\checkmark\\checkmarkc\checkmarki\checkmarkr\checkmarkc\checkmark$\checkmarki\checkmark$\checkmark^\checkmark\\checkmarkc\checkmarki\checkmarkr\checkmarkc\checkmark$\checkmarke\checkmark$\checkmark^\checkmark\\checkmarkc\checkmarki\checkmarkr\checkmarkc\checkmark$\checkmarkr\checkmark$\checkmark^\checkmark\\checkmarkc\checkmarki\checkmarkr\checkmarkc\checkmark$\checkmark \checkmark$\checkmark^\checkmark\\checkmarkc\checkmarki\checkmarkr\checkmarkc\checkmark$\checkmarkd\checkmark$\checkmark^\checkmark\\checkmarkc\checkmarki\checkmarkr\checkmarkc\checkmark$\checkmarke\checkmark$\checkmark^\checkmark\\checkmarkc\checkmarki\checkmarkr\checkmarkc\checkmark$\checkmarks\checkmark$\checkmark^\checkmark\\checkmarkc\checkmarki\checkmarkr\checkmarkc\checkmark$\checkmark \checkmark$\checkmark^\checkmark\\checkmarkc\checkmarki\checkmarkr\checkmarkc\checkmark$\checkmarkC\checkmark$\checkmark^\checkmark\\checkmarkc\checkmarki\checkmarkr\checkmarkc\checkmark$\checkmarkh\checkmark$\checkmark^\checkmark\\checkmarkc\checkmarki\checkmarkr\checkmarkc\checkmark$\checkmarka\checkmark$\checkmark^\checkmark\\checkmarkc\checkmarki\checkmarkr\checkmarkc\checkmark$\checkmarkr\checkmark$\checkmark^\checkmark\\checkmarkc\checkmarki\checkmarkr\checkmarkc\checkmark$\checkmarkg\checkmark$\checkmark^\checkmark\\checkmarkc\checkmarki\checkmarkr\checkmarkc\checkmark$\checkmarke\checkmark$\checkmark^\checkmark\\checkmarkc\checkmarki\checkmarkr\checkmarkc\checkmark$\checkmarks\checkmark$\checkmark^\checkmark\\checkmarkc\checkmarki\checkmarkr\checkmarkc\checkmark$\checkmark \checkmark$\checkmark^\checkmark\\checkmarkc\checkmarki\checkmarkr\checkmarkc\checkmark$\checkmarkF\checkmark$\checkmark^\checkmark\\checkmarkc\checkmarki\checkmarkr\checkmarkc\checkmark$\checkmarko\checkmark$\checkmark^\checkmark\\checkmarkc\checkmarki\checkmarkr\checkmarkc\checkmark$\checkmarkn\checkmark$\checkmark^\checkmark\\checkmarkc\checkmarki\checkmarkr\checkmarkc\checkmark$\checkmarkc\checkmark$\checkmark^\checkmark\\checkmarkc\checkmarki\checkmarkr\checkmarkc\checkmark$\checkmarkt\checkmark$\checkmark^\checkmark\\checkmarkc\checkmarki\checkmarkr\checkmarkc\checkmark$\checkmarki\checkmark$\checkmark^\checkmark\\checkmarkc\checkmarki\checkmarkr\checkmarkc\checkmark$\checkmarko\checkmark$\checkmark^\checkmark\\checkmarkc\checkmarki\checkmarkr\checkmarkc\checkmark$\checkmarkn\checkmark$\checkmark^\checkmark\\checkmarkc\checkmarki\checkmarkr\checkmarkc\checkmark$\checkmarkn\checkmark$\checkmark^\checkmark\\checkmarkc\checkmarki\checkmarkr\checkmarkc\checkmark$\checkmarke\checkmark$\checkmark^\checkmark\\checkmarkc\checkmarki\checkmarkr\checkmarkc\checkmark$\checkmarkl\checkmark$\checkmark^\checkmark\\checkmarkc\checkmarki\checkmarkr\checkmarkc\checkmark$\checkmark \checkmark$\checkmark^\checkmark\\checkmarkc\checkmarki\checkmarkr\checkmarkc\checkmark$\checkmarkQ\checkmark$\checkmark^\checkmark\\checkmarkc\checkmarki\checkmarkr\checkmarkc\checkmark$\checkmarku\checkmark$\checkmark^\checkmark\\checkmarkc\checkmarki\checkmarkr\checkmarkc\checkmark$\checkmarka\checkmark$\checkmark^\checkmark\\checkmarkc\checkmarki\checkmarkr\checkmarkc\checkmark$\checkmarkn\checkmark$\checkmark^\checkmark\\checkmarkc\checkmarki\checkmarkr\checkmarkc\checkmark$\checkmarkt\checkmark$\checkmark^\checkmark\\checkmarkc\checkmarki\checkmarkr\checkmarkc\checkmark$\checkmarki\checkmark$\checkmark^\checkmark\\checkmarkc\checkmarki\checkmarkr\checkmarkc\checkmark$\checkmarkt\checkmark$\checkmark^\checkmark\\checkmarkc\checkmarki\checkmarkr\checkmarkc\checkmark$\checkmarka\checkmark$\checkmark^\checkmark\\checkmarkc\checkmarki\checkmarkr\checkmarkc\checkmark$\checkmarkt\checkmark$\checkmark^\checkmark\\checkmarkc\checkmarki\checkmarkr\checkmarkc\checkmark$\checkmarki\checkmark$\checkmark^\checkmark\\checkmarkc\checkmarki\checkmarkr\checkmarkc\checkmark$\checkmarkf\checkmark$\checkmark^\checkmark\\checkmarkc\checkmarki\checkmarkr\checkmarkc\checkmark$\checkmark}\checkmark$\checkmark^\checkmark\\checkmarkc\checkmarki\checkmarkr\checkmarkc\checkmark$\checkmark
\checkmark$\checkmark^\checkmark\\checkmarkc\checkmarki\checkmarkr\checkmarkc\checkmark$\checkmark
\checkmark$\checkmark^\checkmark\\checkmarkc\checkmarki\checkmarkr\checkmarkc\checkmark$\checkmark\\checkmark$\checkmark^\checkmark\\checkmarkc\checkmarki\checkmarkr\checkmarkc\checkmark$\checkmarkb\checkmark$\checkmark^\checkmark\\checkmarkc\checkmarki\checkmarkr\checkmarkc\checkmark$\checkmarke\checkmark$\checkmark^\checkmark\\checkmarkc\checkmarki\checkmarkr\checkmarkc\checkmark$\checkmarkg\checkmark$\checkmark^\checkmark\\checkmarkc\checkmarki\checkmarkr\checkmarkc\checkmark$\checkmarki\checkmark$\checkmark^\checkmark\\checkmarkc\checkmarki\checkmarkr\checkmarkc\checkmark$\checkmarkn\checkmark$\checkmark^\checkmark\\checkmarkc\checkmarki\checkmarkr\checkmarkc\checkmark$\checkmark{\checkmark$\checkmark^\checkmark\\checkmarkc\checkmarki\checkmarkr\checkmarkc\checkmark$\checkmarkt\checkmark$\checkmark^\checkmark\\checkmarkc\checkmarki\checkmarkr\checkmarkc\checkmark$\checkmarka\checkmark$\checkmark^\checkmark\\checkmarkc\checkmarki\checkmarkr\checkmarkc\checkmark$\checkmarkb\checkmark$\checkmark^\checkmark\\checkmarkc\checkmarki\checkmarkr\checkmarkc\checkmark$\checkmarkl\checkmark$\checkmark^\checkmark\\checkmarkc\checkmarki\checkmarkr\checkmarkc\checkmark$\checkmarke\checkmark$\checkmark^\checkmark\\checkmarkc\checkmarki\checkmarkr\checkmarkc\checkmark$\checkmark}\checkmark$\checkmark^\checkmark\\checkmarkc\checkmarki\checkmarkr\checkmarkc\checkmark$\checkmark[\checkmark$\checkmark^\checkmark\\checkmarkc\checkmarki\checkmarkr\checkmarkc\checkmark$\checkmarkh\checkmark$\checkmark^\checkmark\\checkmarkc\checkmarki\checkmarkr\checkmarkc\checkmark$\checkmarkt\checkmark$\checkmark^\checkmark\\checkmarkc\checkmarki\checkmarkr\checkmarkc\checkmark$\checkmarkb\checkmark$\checkmark^\checkmark\\checkmarkc\checkmarki\checkmarkr\checkmarkc\checkmark$\checkmarkp\checkmark$\checkmark^\checkmark\\checkmarkc\checkmarki\checkmarkr\checkmarkc\checkmark$\checkmark]\checkmark$\checkmark^\checkmark\\checkmarkc\checkmarki\checkmarkr\checkmarkc\checkmark$\checkmark
\checkmark$\checkmark^\checkmark\\checkmarkc\checkmarki\checkmarkr\checkmarkc\checkmark$\checkmark\\checkmark$\checkmark^\checkmark\\checkmarkc\checkmarki\checkmarkr\checkmarkc\checkmark$\checkmarkc\checkmark$\checkmark^\checkmark\\checkmarkc\checkmarki\checkmarkr\checkmarkc\checkmark$\checkmarke\checkmark$\checkmark^\checkmark\\checkmarkc\checkmarki\checkmarkr\checkmarkc\checkmark$\checkmarkn\checkmark$\checkmark^\checkmark\\checkmarkc\checkmarki\checkmarkr\checkmarkc\checkmark$\checkmarkt\checkmark$\checkmark^\checkmark\\checkmarkc\checkmarki\checkmarkr\checkmarkc\checkmark$\checkmarke\checkmark$\checkmark^\checkmark\\checkmarkc\checkmarki\checkmarkr\checkmarkc\checkmark$\checkmarkr\checkmark$\checkmark^\checkmark\\checkmarkc\checkmarki\checkmarkr\checkmarkc\checkmark$\checkmarki\checkmark$\checkmark^\checkmark\\checkmarkc\checkmarki\checkmarkr\checkmarkc\checkmark$\checkmarkn\checkmark$\checkmark^\checkmark\\checkmarkc\checkmarki\checkmarkr\checkmarkc\checkmark$\checkmarkg\checkmark$\checkmark^\checkmark\\checkmarkc\checkmarki\checkmarkr\checkmarkc\checkmark$\checkmark
\checkmark$\checkmark^\checkmark\\checkmarkc\checkmarki\checkmarkr\checkmarkc\checkmark$\checkmark\\checkmark$\checkmark^\checkmark\\checkmarkc\checkmarki\checkmarkr\checkmarkc\checkmark$\checkmarkr\checkmark$\checkmark^\checkmark\\checkmarkc\checkmarki\checkmarkr\checkmarkc\checkmark$\checkmarke\checkmark$\checkmark^\checkmark\\checkmarkc\checkmarki\checkmarkr\checkmarkc\checkmark$\checkmarks\checkmark$\checkmark^\checkmark\\checkmarkc\checkmarki\checkmarkr\checkmarkc\checkmark$\checkmarki\checkmark$\checkmark^\checkmark\\checkmarkc\checkmarki\checkmarkr\checkmarkc\checkmark$\checkmarkz\checkmark$\checkmark^\checkmark\\checkmarkc\checkmarki\checkmarkr\checkmarkc\checkmark$\checkmarke\checkmark$\checkmark^\checkmark\\checkmarkc\checkmarki\checkmarkr\checkmarkc\checkmark$\checkmarkb\checkmark$\checkmark^\checkmark\\checkmarkc\checkmarki\checkmarkr\checkmarkc\checkmark$\checkmarko\checkmark$\checkmark^\checkmark\\checkmarkc\checkmarki\checkmarkr\checkmarkc\checkmark$\checkmarkx\checkmark$\checkmark^\checkmark\\checkmarkc\checkmarki\checkmarkr\checkmarkc\checkmark$\checkmark{\checkmark$\checkmark^\checkmark\\checkmarkc\checkmarki\checkmarkr\checkmarkc\checkmark$\checkmark\\checkmark$\checkmark^\checkmark\\checkmarkc\checkmarki\checkmarkr\checkmarkc\checkmark$\checkmarkt\checkmark$\checkmark^\checkmark\\checkmarkc\checkmarki\checkmarkr\checkmarkc\checkmark$\checkmarke\checkmark$\checkmark^\checkmark\\checkmarkc\checkmarki\checkmarkr\checkmarkc\checkmark$\checkmarkx\checkmark$\checkmark^\checkmark\\checkmarkc\checkmarki\checkmarkr\checkmarkc\checkmark$\checkmarkt\checkmark$\checkmark^\checkmark\\checkmarkc\checkmarki\checkmarkr\checkmarkc\checkmark$\checkmarkw\checkmark$\checkmark^\checkmark\\checkmarkc\checkmarki\checkmarkr\checkmarkc\checkmark$\checkmarki\checkmark$\checkmark^\checkmark\\checkmarkc\checkmarki\checkmarkr\checkmarkc\checkmark$\checkmarkd\checkmark$\checkmark^\checkmark\\checkmarkc\checkmarki\checkmarkr\checkmarkc\checkmark$\checkmarkt\checkmark$\checkmark^\checkmark\\checkmarkc\checkmarki\checkmarkr\checkmarkc\checkmark$\checkmarkh\checkmark$\checkmark^\checkmark\\checkmarkc\checkmarki\checkmarkr\checkmarkc\checkmark$\checkmark}\checkmark$\checkmark^\checkmark\\checkmarkc\checkmarki\checkmarkr\checkmarkc\checkmark$\checkmark{\checkmark$\checkmark^\checkmark\\checkmarkc\checkmarki\checkmarkr\checkmarkc\checkmark$\checkmark!\checkmark$\checkmark^\checkmark\\checkmarkc\checkmarki\checkmarkr\checkmarkc\checkmark$\checkmark}\checkmark$\checkmark^\checkmark\\checkmarkc\checkmarki\checkmarkr\checkmarkc\checkmark$\checkmark{\checkmark$\checkmark^\checkmark\\checkmarkc\checkmarki\checkmarkr\checkmarkc\checkmark$\checkmark
\checkmark$\checkmark^\checkmark\\checkmarkc\checkmarki\checkmarkr\checkmarkc\checkmark$\checkmark\\checkmark$\checkmark^\checkmark\\checkmarkc\checkmarki\checkmarkr\checkmarkc\checkmark$\checkmarkb\checkmark$\checkmark^\checkmark\\checkmarkc\checkmarki\checkmarkr\checkmarkc\checkmark$\checkmarke\checkmark$\checkmark^\checkmark\\checkmarkc\checkmarki\checkmarkr\checkmarkc\checkmark$\checkmarkg\checkmark$\checkmark^\checkmark\\checkmarkc\checkmarki\checkmarkr\checkmarkc\checkmark$\checkmarki\checkmark$\checkmark^\checkmark\\checkmarkc\checkmarki\checkmarkr\checkmarkc\checkmark$\checkmarkn\checkmark$\checkmark^\checkmark\\checkmarkc\checkmarki\checkmarkr\checkmarkc\checkmark$\checkmark{\checkmark$\checkmark^\checkmark\\checkmarkc\checkmarki\checkmarkr\checkmarkc\checkmark$\checkmarkt\checkmark$\checkmark^\checkmark\\checkmarkc\checkmarki\checkmarkr\checkmarkc\checkmark$\checkmarka\checkmark$\checkmark^\checkmark\\checkmarkc\checkmarki\checkmarkr\checkmarkc\checkmark$\checkmarkb\checkmark$\checkmark^\checkmark\\checkmarkc\checkmarki\checkmarkr\checkmarkc\checkmark$\checkmarku\checkmark$\checkmark^\checkmark\\checkmarkc\checkmarki\checkmarkr\checkmarkc\checkmark$\checkmarkl\checkmark$\checkmark^\checkmark\\checkmarkc\checkmarki\checkmarkr\checkmarkc\checkmark$\checkmarka\checkmark$\checkmark^\checkmark\\checkmarkc\checkmarki\checkmarkr\checkmarkc\checkmark$\checkmarkr\checkmark$\checkmark^\checkmark\\checkmarkc\checkmarki\checkmarkr\checkmarkc\checkmark$\checkmark}\checkmark$\checkmark^\checkmark\\checkmarkc\checkmarki\checkmarkr\checkmarkc\checkmark$\checkmark{\checkmark$\checkmark^\checkmark\\checkmarkc\checkmarki\checkmarkr\checkmarkc\checkmark$\checkmark|\checkmark$\checkmark^\checkmark\\checkmarkc\checkmarki\checkmarkr\checkmarkc\checkmark$\checkmarkl\checkmark$\checkmark^\checkmark\\checkmarkc\checkmarki\checkmarkr\checkmarkc\checkmark$\checkmark|\checkmark$\checkmark^\checkmark\\checkmarkc\checkmarki\checkmarkr\checkmarkc\checkmark$\checkmarkl\checkmark$\checkmark^\checkmark\\checkmarkc\checkmarki\checkmarkr\checkmarkc\checkmark$\checkmark|\checkmark$\checkmark^\checkmark\\checkmarkc\checkmarki\checkmarkr\checkmarkc\checkmark$\checkmarkc\checkmark$\checkmark^\checkmark\\checkmarkc\checkmarki\checkmarkr\checkmarkc\checkmark$\checkmark|\checkmark$\checkmark^\checkmark\\checkmarkc\checkmarki\checkmarkr\checkmarkc\checkmark$\checkmarkl\checkmark$\checkmark^\checkmark\\checkmarkc\checkmarki\checkmarkr\checkmarkc\checkmark$\checkmark|\checkmark$\checkmark^\checkmark\\checkmarkc\checkmarki\checkmarkr\checkmarkc\checkmark$\checkmark}\checkmark$\checkmark^\checkmark\\checkmarkc\checkmarki\checkmarkr\checkmarkc\checkmark$\checkmark
\checkmark$\checkmark^\checkmark\\checkmarkc\checkmarki\checkmarkr\checkmarkc\checkmark$\checkmark\\checkmark$\checkmark^\checkmark\\checkmarkc\checkmarki\checkmarkr\checkmarkc\checkmark$\checkmarkh\checkmark$\checkmark^\checkmark\\checkmarkc\checkmarki\checkmarkr\checkmarkc\checkmark$\checkmarkl\checkmark$\checkmark^\checkmark\\checkmarkc\checkmarki\checkmarkr\checkmarkc\checkmark$\checkmarki\checkmark$\checkmark^\checkmark\\checkmarkc\checkmarki\checkmarkr\checkmarkc\checkmark$\checkmarkn\checkmark$\checkmark^\checkmark\\checkmarkc\checkmarki\checkmarkr\checkmarkc\checkmark$\checkmarke\checkmark$\checkmark^\checkmark\\checkmarkc\checkmarki\checkmarkr\checkmarkc\checkmark$\checkmark
\checkmark$\checkmark^\checkmark\\checkmarkc\checkmarki\checkmarkr\checkmarkc\checkmark$\checkmark\\checkmark$\checkmark^\checkmark\\checkmarkc\checkmarki\checkmarkr\checkmarkc\checkmark$\checkmarkt\checkmark$\checkmark^\checkmark\\checkmarkc\checkmarki\checkmarkr\checkmarkc\checkmark$\checkmarke\checkmark$\checkmark^\checkmark\\checkmarkc\checkmarki\checkmarkr\checkmarkc\checkmark$\checkmarkx\checkmark$\checkmark^\checkmark\\checkmarkc\checkmarki\checkmarkr\checkmarkc\checkmark$\checkmarkt\checkmark$\checkmark^\checkmark\\checkmarkc\checkmarki\checkmarkr\checkmarkc\checkmark$\checkmarkb\checkmark$\checkmark^\checkmark\\checkmarkc\checkmarki\checkmarkr\checkmarkc\checkmark$\checkmarkf\checkmark$\checkmark^\checkmark\\checkmarkc\checkmarki\checkmarkr\checkmarkc\checkmark$\checkmark{\checkmark$\checkmark^\checkmark\\checkmarkc\checkmarki\checkmarkr\checkmarkc\checkmark$\checkmarkF\checkmark$\checkmark^\checkmark\\checkmarkc\checkmarki\checkmarkr\checkmarkc\checkmark$\checkmarko\checkmark$\checkmark^\checkmark\\checkmarkc\checkmarki\checkmarkr\checkmarkc\checkmark$\checkmarkn\checkmark$\checkmark^\checkmark\\checkmarkc\checkmarki\checkmarkr\checkmarkc\checkmark$\checkmarkc\checkmark$\checkmark^\checkmark\\checkmarkc\checkmarki\checkmarkr\checkmarkc\checkmark$\checkmarkt\checkmark$\checkmark^\checkmark\\checkmarkc\checkmarki\checkmarkr\checkmarkc\checkmark$\checkmarki\checkmark$\checkmark^\checkmark\\checkmarkc\checkmarki\checkmarkr\checkmarkc\checkmark$\checkmarko\checkmark$\checkmark^\checkmark\\checkmarkc\checkmarki\checkmarkr\checkmarkc\checkmark$\checkmarkn\checkmark$\checkmark^\checkmark\\checkmarkc\checkmarki\checkmarkr\checkmarkc\checkmark$\checkmark}\checkmark$\checkmark^\checkmark\\checkmarkc\checkmarki\checkmarkr\checkmarkc\checkmark$\checkmark \checkmark$\checkmark^\checkmark\\checkmarkc\checkmarki\checkmarkr\checkmarkc\checkmark$\checkmark&\checkmark$\checkmark^\checkmark\\checkmarkc\checkmarki\checkmarkr\checkmarkc\checkmark$\checkmark \checkmark$\checkmark^\checkmark\\checkmarkc\checkmarki\checkmarkr\checkmarkc\checkmark$\checkmark\\checkmark$\checkmark^\checkmark\\checkmarkc\checkmarki\checkmarkr\checkmarkc\checkmark$\checkmarkt\checkmark$\checkmark^\checkmark\\checkmarkc\checkmarki\checkmarkr\checkmarkc\checkmark$\checkmarke\checkmark$\checkmark^\checkmark\\checkmarkc\checkmarki\checkmarkr\checkmarkc\checkmark$\checkmarkx\checkmark$\checkmark^\checkmark\\checkmarkc\checkmarki\checkmarkr\checkmarkc\checkmark$\checkmarkt\checkmark$\checkmark^\checkmark\\checkmarkc\checkmarki\checkmarkr\checkmarkc\checkmark$\checkmarkb\checkmark$\checkmark^\checkmark\\checkmarkc\checkmarki\checkmarkr\checkmarkc\checkmark$\checkmarkf\checkmark$\checkmark^\checkmark\\checkmarkc\checkmarki\checkmarkr\checkmarkc\checkmark$\checkmark{\checkmark$\checkmark^\checkmark\\checkmarkc\checkmarki\checkmarkr\checkmarkc\checkmark$\checkmarkC\checkmark$\checkmark^\checkmark\\checkmarkc\checkmarki\checkmarkr\checkmarkc\checkmark$\checkmarkr\checkmark$\checkmark^\checkmark\\checkmarkc\checkmarki\checkmarkr\checkmarkc\checkmark$\checkmarki\checkmark$\checkmark^\checkmark\\checkmarkc\checkmarki\checkmarkr\checkmarkc\checkmark$\checkmarkt\checkmark$\checkmark^\checkmark\\checkmarkc\checkmarki\checkmarkr\checkmarkc\checkmark$\checkmarkè\checkmark$\checkmark^\checkmark\\checkmarkc\checkmarki\checkmarkr\checkmarkc\checkmark$\checkmarkr\checkmark$\checkmark^\checkmark\\checkmarkc\checkmarki\checkmarkr\checkmarkc\checkmark$\checkmarke\checkmark$\checkmark^\checkmark\\checkmarkc\checkmarki\checkmarkr\checkmarkc\checkmark$\checkmark \checkmark$\checkmark^\checkmark\\checkmarkc\checkmarki\checkmarkr\checkmarkc\checkmark$\checkmarkd\checkmark$\checkmark^\checkmark\\checkmarkc\checkmarki\checkmarkr\checkmarkc\checkmark$\checkmark'\checkmark$\checkmark^\checkmark\\checkmarkc\checkmarki\checkmarkr\checkmarkc\checkmark$\checkmarké\checkmark$\checkmark^\checkmark\\checkmarkc\checkmarki\checkmarkr\checkmarkc\checkmark$\checkmarkv\checkmark$\checkmark^\checkmark\\checkmarkc\checkmarki\checkmarkr\checkmarkc\checkmark$\checkmarka\checkmark$\checkmark^\checkmark\\checkmarkc\checkmarki\checkmarkr\checkmarkc\checkmark$\checkmarkl\checkmark$\checkmark^\checkmark\\checkmarkc\checkmarki\checkmarkr\checkmarkc\checkmark$\checkmarku\checkmark$\checkmark^\checkmark\\checkmarkc\checkmarki\checkmarkr\checkmarkc\checkmark$\checkmarka\checkmark$\checkmark^\checkmark\\checkmarkc\checkmarki\checkmarkr\checkmarkc\checkmark$\checkmarkt\checkmark$\checkmark^\checkmark\\checkmarkc\checkmarki\checkmarkr\checkmarkc\checkmark$\checkmarki\checkmark$\checkmark^\checkmark\\checkmarkc\checkmarki\checkmarkr\checkmarkc\checkmark$\checkmarko\checkmark$\checkmark^\checkmark\\checkmarkc\checkmarki\checkmarkr\checkmarkc\checkmark$\checkmarkn\checkmark$\checkmark^\checkmark\\checkmarkc\checkmarki\checkmarkr\checkmarkc\checkmark$\checkmark}\checkmark$\checkmark^\checkmark\\checkmarkc\checkmarki\checkmarkr\checkmarkc\checkmark$\checkmark \checkmark$\checkmark^\checkmark\\checkmarkc\checkmarki\checkmarkr\checkmarkc\checkmark$\checkmark&\checkmark$\checkmark^\checkmark\\checkmarkc\checkmarki\checkmarkr\checkmarkc\checkmark$\checkmark \checkmark$\checkmark^\checkmark\\checkmarkc\checkmarki\checkmarkr\checkmarkc\checkmark$\checkmark\\checkmark$\checkmark^\checkmark\\checkmarkc\checkmarki\checkmarkr\checkmarkc\checkmark$\checkmarkt\checkmark$\checkmark^\checkmark\\checkmarkc\checkmarki\checkmarkr\checkmarkc\checkmark$\checkmarke\checkmark$\checkmark^\checkmark\\checkmarkc\checkmarki\checkmarkr\checkmarkc\checkmark$\checkmarkx\checkmark$\checkmark^\checkmark\\checkmarkc\checkmarki\checkmarkr\checkmarkc\checkmark$\checkmarkt\checkmark$\checkmark^\checkmark\\checkmarkc\checkmarki\checkmarkr\checkmarkc\checkmark$\checkmarkb\checkmark$\checkmark^\checkmark\\checkmarkc\checkmarki\checkmarkr\checkmarkc\checkmark$\checkmarkf\checkmark$\checkmark^\checkmark\\checkmarkc\checkmarki\checkmarkr\checkmarkc\checkmark$\checkmark{\checkmark$\checkmark^\checkmark\\checkmarkc\checkmarki\checkmarkr\checkmarkc\checkmark$\checkmarkN\checkmark$\checkmark^\checkmark\\checkmarkc\checkmarki\checkmarkr\checkmarkc\checkmark$\checkmarki\checkmark$\checkmark^\checkmark\\checkmarkc\checkmarki\checkmarkr\checkmarkc\checkmark$\checkmarkv\checkmark$\checkmark^\checkmark\\checkmarkc\checkmarki\checkmarkr\checkmarkc\checkmark$\checkmarke\checkmark$\checkmark^\checkmark\\checkmarkc\checkmarki\checkmarkr\checkmarkc\checkmark$\checkmarka\checkmark$\checkmark^\checkmark\\checkmarkc\checkmarki\checkmarkr\checkmarkc\checkmark$\checkmarku\checkmark$\checkmark^\checkmark\\checkmarkc\checkmarki\checkmarkr\checkmarkc\checkmark$\checkmark \checkmark$\checkmark^\checkmark\\checkmarkc\checkmarki\checkmarkr\checkmarkc\checkmark$\checkmarkd\checkmark$\checkmark^\checkmark\\checkmarkc\checkmarki\checkmarkr\checkmarkc\checkmark$\checkmark'\checkmark$\checkmark^\checkmark\\checkmarkc\checkmarki\checkmarkr\checkmarkc\checkmark$\checkmarke\checkmark$\checkmark^\checkmark\\checkmarkc\checkmarki\checkmarkr\checkmarkc\checkmark$\checkmarkx\checkmark$\checkmark^\checkmark\\checkmarkc\checkmarki\checkmarkr\checkmarkc\checkmark$\checkmarki\checkmark$\checkmark^\checkmark\\checkmarkc\checkmarki\checkmarkr\checkmarkc\checkmark$\checkmarkg\checkmark$\checkmark^\checkmark\\checkmarkc\checkmarki\checkmarkr\checkmarkc\checkmark$\checkmarke\checkmark$\checkmark^\checkmark\\checkmarkc\checkmarki\checkmarkr\checkmarkc\checkmark$\checkmarkn\checkmark$\checkmark^\checkmark\\checkmarkc\checkmarki\checkmarkr\checkmarkc\checkmark$\checkmarkc\checkmark$\checkmark^\checkmark\\checkmarkc\checkmarki\checkmarkr\checkmarkc\checkmark$\checkmarke\checkmark$\checkmark^\checkmark\\checkmarkc\checkmarki\checkmarkr\checkmarkc\checkmark$\checkmark}\checkmark$\checkmark^\checkmark\\checkmarkc\checkmarki\checkmarkr\checkmarkc\checkmark$\checkmark \checkmark$\checkmark^\checkmark\\checkmarkc\checkmarki\checkmarkr\checkmarkc\checkmark$\checkmark&\checkmark$\checkmark^\checkmark\\checkmarkc\checkmarki\checkmarkr\checkmarkc\checkmark$\checkmark \checkmark$\checkmark^\checkmark\\checkmarkc\checkmarki\checkmarkr\checkmarkc\checkmark$\checkmark\\checkmark$\checkmark^\checkmark\\checkmarkc\checkmarki\checkmarkr\checkmarkc\checkmark$\checkmarkt\checkmark$\checkmark^\checkmark\\checkmarkc\checkmarki\checkmarkr\checkmarkc\checkmark$\checkmarke\checkmark$\checkmark^\checkmark\\checkmarkc\checkmarki\checkmarkr\checkmarkc\checkmark$\checkmarkx\checkmark$\checkmark^\checkmark\\checkmarkc\checkmarki\checkmarkr\checkmarkc\checkmark$\checkmarkt\checkmark$\checkmark^\checkmark\\checkmarkc\checkmarki\checkmarkr\checkmarkc\checkmark$\checkmarkb\checkmark$\checkmark^\checkmark\\checkmarkc\checkmarki\checkmarkr\checkmarkc\checkmark$\checkmarkf\checkmark$\checkmark^\checkmark\\checkmarkc\checkmarki\checkmarkr\checkmarkc\checkmark$\checkmark{\checkmark$\checkmark^\checkmark\\checkmarkc\checkmarki\checkmarkr\checkmarkc\checkmark$\checkmarkF\checkmark$\checkmark^\checkmark\\checkmarkc\checkmarki\checkmarkr\checkmarkc\checkmark$\checkmarkl\checkmark$\checkmark^\checkmark\\checkmarkc\checkmarki\checkmarkr\checkmarkc\checkmark$\checkmarke\checkmark$\checkmark^\checkmark\\checkmarkc\checkmarki\checkmarkr\checkmarkc\checkmark$\checkmarkx\checkmark$\checkmark^\checkmark\\checkmarkc\checkmarki\checkmarkr\checkmarkc\checkmark$\checkmarki\checkmark$\checkmark^\checkmark\\checkmarkc\checkmarki\checkmarkr\checkmarkc\checkmark$\checkmarkb\checkmark$\checkmark^\checkmark\\checkmarkc\checkmarki\checkmarkr\checkmarkc\checkmark$\checkmarki\checkmark$\checkmark^\checkmark\\checkmarkc\checkmarki\checkmarkr\checkmarkc\checkmark$\checkmarkl\checkmark$\checkmark^\checkmark\\checkmarkc\checkmarki\checkmarkr\checkmarkc\checkmark$\checkmarki\checkmark$\checkmark^\checkmark\\checkmarkc\checkmarki\checkmarkr\checkmarkc\checkmark$\checkmarkt\checkmark$\checkmark^\checkmark\\checkmarkc\checkmarki\checkmarkr\checkmarkc\checkmark$\checkmarké\checkmark$\checkmark^\checkmark\\checkmarkc\checkmarki\checkmarkr\checkmarkc\checkmark$\checkmark}\checkmark$\checkmark^\checkmark\\checkmarkc\checkmarki\checkmarkr\checkmarkc\checkmark$\checkmark \checkmark$\checkmark^\checkmark\\checkmarkc\checkmarki\checkmarkr\checkmarkc\checkmark$\checkmark\\checkmark$\checkmark^\checkmark\\checkmarkc\checkmarki\checkmarkr\checkmarkc\checkmark$\checkmark\\checkmark$\checkmark^\checkmark\\checkmarkc\checkmarki\checkmarkr\checkmarkc\checkmark$\checkmark \checkmark$\checkmark^\checkmark\\checkmarkc\checkmarki\checkmarkr\checkmarkc\checkmark$\checkmark\\checkmark$\checkmark^\checkmark\\checkmarkc\checkmarki\checkmarkr\checkmarkc\checkmark$\checkmarkh\checkmark$\checkmark^\checkmark\\checkmarkc\checkmarki\checkmarkr\checkmarkc\checkmark$\checkmarkl\checkmark$\checkmark^\checkmark\\checkmarkc\checkmarki\checkmarkr\checkmarkc\checkmark$\checkmarki\checkmark$\checkmark^\checkmark\\checkmarkc\checkmarki\checkmarkr\checkmarkc\checkmark$\checkmarkn\checkmark$\checkmark^\checkmark\\checkmarkc\checkmarki\checkmarkr\checkmarkc\checkmark$\checkmarke\checkmark$\checkmark^\checkmark\\checkmarkc\checkmarki\checkmarkr\checkmarkc\checkmark$\checkmark
\checkmark$\checkmark^\checkmark\\checkmarkc\checkmarki\checkmarkr\checkmarkc\checkmark$\checkmark\\checkmark$\checkmark^\checkmark\\checkmarkc\checkmarki\checkmarkr\checkmarkc\checkmark$\checkmarkm\checkmark$\checkmark^\checkmark\\checkmarkc\checkmarki\checkmarkr\checkmarkc\checkmark$\checkmarku\checkmark$\checkmark^\checkmark\\checkmarkc\checkmarki\checkmarkr\checkmarkc\checkmark$\checkmarkl\checkmark$\checkmark^\checkmark\\checkmarkc\checkmarki\checkmarkr\checkmarkc\checkmark$\checkmarkt\checkmark$\checkmark^\checkmark\\checkmarkc\checkmarki\checkmarkr\checkmarkc\checkmark$\checkmarki\checkmark$\checkmark^\checkmark\\checkmarkc\checkmarki\checkmarkr\checkmarkc\checkmark$\checkmarkr\checkmark$\checkmark^\checkmark\\checkmarkc\checkmarki\checkmarkr\checkmarkc\checkmark$\checkmarko\checkmark$\checkmark^\checkmark\\checkmarkc\checkmarki\checkmarkr\checkmarkc\checkmark$\checkmarkw\checkmark$\checkmark^\checkmark\\checkmarkc\checkmarki\checkmarkr\checkmarkc\checkmark$\checkmark{\checkmark$\checkmark^\checkmark\\checkmarkc\checkmarki\checkmarkr\checkmarkc\checkmark$\checkmark4\checkmark$\checkmark^\checkmark\\checkmarkc\checkmarki\checkmarkr\checkmarkc\checkmark$\checkmark}\checkmark$\checkmark^\checkmark\\checkmarkc\checkmarki\checkmarkr\checkmarkc\checkmark$\checkmark{\checkmark$\checkmark^\checkmark\\checkmarkc\checkmarki\checkmarkr\checkmarkc\checkmark$\checkmark*\checkmark$\checkmark^\checkmark\\checkmarkc\checkmarki\checkmarkr\checkmarkc\checkmark$\checkmark}\checkmark$\checkmark^\checkmark\\checkmarkc\checkmarki\checkmarkr\checkmarkc\checkmark$\checkmark{\checkmark$\checkmark^\checkmark\\checkmarkc\checkmarki\checkmarkr\checkmarkc\checkmark$\checkmark\\checkmark$\checkmark^\checkmark\\checkmarkc\checkmarki\checkmarkr\checkmarkc\checkmark$\checkmarkt\checkmark$\checkmark^\checkmark\\checkmarkc\checkmarki\checkmarkr\checkmarkc\checkmark$\checkmarke\checkmark$\checkmark^\checkmark\\checkmarkc\checkmarki\checkmarkr\checkmarkc\checkmark$\checkmarkx\checkmark$\checkmark^\checkmark\\checkmarkc\checkmarki\checkmarkr\checkmarkc\checkmark$\checkmarkt\checkmark$\checkmark^\checkmark\\checkmarkc\checkmarki\checkmarkr\checkmarkc\checkmark$\checkmarkb\checkmark$\checkmark^\checkmark\\checkmarkc\checkmarki\checkmarkr\checkmarkc\checkmark$\checkmarkf\checkmark$\checkmark^\checkmark\\checkmarkc\checkmarki\checkmarkr\checkmarkc\checkmark$\checkmark{\checkmark$\checkmark^\checkmark\\checkmarkc\checkmarki\checkmarkr\checkmarkc\checkmark$\checkmarkF\checkmark$\checkmark^\checkmark\\checkmarkc\checkmarki\checkmarkr\checkmarkc\checkmark$\checkmarkP\checkmark$\checkmark^\checkmark\\checkmarkc\checkmarki\checkmarkr\checkmarkc\checkmark$\checkmark}\checkmark$\checkmark^\checkmark\\checkmarkc\checkmarki\checkmarkr\checkmarkc\checkmark$\checkmark}\checkmark$\checkmark^\checkmark\\checkmarkc\checkmarki\checkmarkr\checkmarkc\checkmark$\checkmark \checkmark$\checkmark^\checkmark\\checkmarkc\checkmarki\checkmarkr\checkmarkc\checkmark$\checkmark&\checkmark$\checkmark^\checkmark\\checkmarkc\checkmarki\checkmarkr\checkmarkc\checkmark$\checkmark \checkmark$\checkmark^\checkmark\\checkmarkc\checkmarki\checkmarkr\checkmarkc\checkmark$\checkmarkZ\checkmark$\checkmark^\checkmark\\checkmarkc\checkmarki\checkmarkr\checkmarkc\checkmark$\checkmarko\checkmark$\checkmark^\checkmark\\checkmarkc\checkmarki\checkmarkr\checkmarkc\checkmark$\checkmarkn\checkmark$\checkmark^\checkmark\\checkmarkc\checkmarki\checkmarkr\checkmarkc\checkmark$\checkmarke\checkmark$\checkmark^\checkmark\\checkmarkc\checkmarki\checkmarkr\checkmarkc\checkmark$\checkmark \checkmark$\checkmark^\checkmark\\checkmarkc\checkmarki\checkmarkr\checkmarkc\checkmark$\checkmarkd\checkmark$\checkmark^\checkmark\\checkmarkc\checkmarki\checkmarkr\checkmarkc\checkmark$\checkmarke\checkmark$\checkmark^\checkmark\\checkmarkc\checkmarki\checkmarkr\checkmarkc\checkmark$\checkmark \checkmark$\checkmark^\checkmark\\checkmarkc\checkmarki\checkmarkr\checkmarkc\checkmark$\checkmarkt\checkmark$\checkmark^\checkmark\\checkmarkc\checkmarki\checkmarkr\checkmarkc\checkmark$\checkmarkr\checkmark$\checkmark^\checkmark\\checkmarkc\checkmarki\checkmarkr\checkmarkc\checkmark$\checkmarka\checkmark$\checkmark^\checkmark\\checkmarkc\checkmarki\checkmarkr\checkmarkc\checkmark$\checkmarkv\checkmark$\checkmark^\checkmark\\checkmarkc\checkmarki\checkmarkr\checkmarkc\checkmark$\checkmarka\checkmark$\checkmark^\checkmark\\checkmarkc\checkmarki\checkmarkr\checkmarkc\checkmark$\checkmarki\checkmark$\checkmark^\checkmark\\checkmarkc\checkmarki\checkmarkr\checkmarkc\checkmark$\checkmarkl\checkmark$\checkmark^\checkmark\\checkmarkc\checkmarki\checkmarkr\checkmarkc\checkmark$\checkmark \checkmark$\checkmark^\checkmark\\checkmarkc\checkmarki\checkmarkr\checkmarkc\checkmark$\checkmarkX\checkmark$\checkmark^\checkmark\\checkmarkc\checkmarki\checkmarkr\checkmarkc\checkmark$\checkmark \checkmark$\checkmark^\checkmark\\checkmarkc\checkmarki\checkmarkr\checkmarkc\checkmark$\checkmark&\checkmark$\checkmark^\checkmark\\checkmarkc\checkmarki\checkmarkr\checkmarkc\checkmark$\checkmark \checkmark$\checkmark^\checkmark\\checkmarkc\checkmarki\checkmarkr\checkmarkc\checkmark$\checkmark$\checkmark$\checkmark^\checkmark\\checkmarkc\checkmarki\checkmarkr\checkmarkc\checkmark$\checkmark2\checkmark$\checkmark^\checkmark\\checkmarkc\checkmarki\checkmarkr\checkmarkc\checkmark$\checkmark0\checkmark$\checkmark^\checkmark\\checkmarkc\checkmarki\checkmarkr\checkmarkc\checkmark$\checkmark0\checkmark$\checkmark^\checkmark\\checkmarkc\checkmarki\checkmarkr\checkmarkc\checkmark$\checkmark$\checkmark$\checkmark^\checkmark\\checkmarkc\checkmarki\checkmarkr\checkmarkc\checkmark$\checkmark \checkmark$\checkmark^\checkmark\\checkmarkc\checkmarki\checkmarkr\checkmarkc\checkmark$\checkmarkm\checkmark$\checkmark^\checkmark\\checkmarkc\checkmarki\checkmarkr\checkmarkc\checkmark$\checkmarkm\checkmark$\checkmark^\checkmark\\checkmarkc\checkmarki\checkmarkr\checkmarkc\checkmark$\checkmark \checkmark$\checkmark^\checkmark\\checkmarkc\checkmarki\checkmarkr\checkmarkc\checkmark$\checkmark&\checkmark$\checkmark^\checkmark\\checkmarkc\checkmarki\checkmarkr\checkmarkc\checkmark$\checkmark \checkmark$\checkmark^\checkmark\\checkmarkc\checkmarki\checkmarkr\checkmarkc\checkmark$\checkmark$\checkmark$\checkmark^\checkmark\\checkmarkc\checkmarki\checkmarkr\checkmarkc\checkmark$\checkmark\\checkmark$\checkmark^\checkmark\\checkmarkc\checkmarki\checkmarkr\checkmarkc\checkmark$\checkmarkp\checkmark$\checkmark^\checkmark\\checkmarkc\checkmarki\checkmarkr\checkmarkc\checkmark$\checkmarkm\checkmark$\checkmark^\checkmark\\checkmarkc\checkmarki\checkmarkr\checkmarkc\checkmark$\checkmark \checkmark$\checkmark^\checkmark\\checkmarkc\checkmarki\checkmarkr\checkmarkc\checkmark$\checkmark5\checkmark$\checkmark^\checkmark\\checkmarkc\checkmarki\checkmarkr\checkmarkc\checkmark$\checkmark$\checkmark$\checkmark^\checkmark\\checkmarkc\checkmarki\checkmarkr\checkmarkc\checkmark$\checkmark \checkmark$\checkmark^\checkmark\\checkmarkc\checkmarki\checkmarkr\checkmarkc\checkmark$\checkmarkm\checkmark$\checkmark^\checkmark\\checkmarkc\checkmarki\checkmarkr\checkmarkc\checkmark$\checkmarkm\checkmark$\checkmark^\checkmark\\checkmarkc\checkmarki\checkmarkr\checkmarkc\checkmark$\checkmark \checkmark$\checkmark^\checkmark\\checkmarkc\checkmarki\checkmarkr\checkmarkc\checkmark$\checkmark\\checkmark$\checkmark^\checkmark\\checkmarkc\checkmarki\checkmarkr\checkmarkc\checkmark$\checkmark\\checkmark$\checkmark^\checkmark\\checkmarkc\checkmarki\checkmarkr\checkmarkc\checkmark$\checkmark \checkmark$\checkmark^\checkmark\\checkmarkc\checkmarki\checkmarkr\checkmarkc\checkmark$\checkmark\\checkmark$\checkmark^\checkmark\\checkmarkc\checkmarki\checkmarkr\checkmarkc\checkmark$\checkmarkc\checkmark$\checkmark^\checkmark\\checkmarkc\checkmarki\checkmarkr\checkmarkc\checkmark$\checkmarkl\checkmark$\checkmark^\checkmark\\checkmarkc\checkmarki\checkmarkr\checkmarkc\checkmark$\checkmarki\checkmark$\checkmark^\checkmark\\checkmarkc\checkmarki\checkmarkr\checkmarkc\checkmark$\checkmarkn\checkmark$\checkmark^\checkmark\\checkmarkc\checkmarki\checkmarkr\checkmarkc\checkmark$\checkmarke\checkmark$\checkmark^\checkmark\\checkmarkc\checkmarki\checkmarkr\checkmarkc\checkmark$\checkmark{\checkmark$\checkmark^\checkmark\\checkmarkc\checkmarki\checkmarkr\checkmarkc\checkmark$\checkmark2\checkmark$\checkmark^\checkmark\\checkmarkc\checkmarki\checkmarkr\checkmarkc\checkmark$\checkmark-\checkmark$\checkmark^\checkmark\\checkmarkc\checkmarki\checkmarkr\checkmarkc\checkmark$\checkmark4\checkmark$\checkmark^\checkmark\\checkmarkc\checkmarki\checkmarkr\checkmarkc\checkmark$\checkmark}\checkmark$\checkmark^\checkmark\\checkmarkc\checkmarki\checkmarkr\checkmarkc\checkmark$\checkmark
\checkmark$\checkmark^\checkmark\\checkmarkc\checkmarki\checkmarkr\checkmarkc\checkmark$\checkmark \checkmark$\checkmark^\checkmark\\checkmarkc\checkmarki\checkmarkr\checkmarkc\checkmark$\checkmark&\checkmark$\checkmark^\checkmark\\checkmarkc\checkmarki\checkmarkr\checkmarkc\checkmark$\checkmark \checkmark$\checkmark^\checkmark\\checkmarkc\checkmarki\checkmarkr\checkmarkc\checkmark$\checkmarkZ\checkmark$\checkmark^\checkmark\\checkmarkc\checkmarki\checkmarkr\checkmarkc\checkmark$\checkmarko\checkmark$\checkmark^\checkmark\\checkmarkc\checkmarki\checkmarkr\checkmarkc\checkmark$\checkmarkn\checkmark$\checkmark^\checkmark\\checkmarkc\checkmarki\checkmarkr\checkmarkc\checkmark$\checkmarke\checkmark$\checkmark^\checkmark\\checkmarkc\checkmarki\checkmarkr\checkmarkc\checkmark$\checkmark \checkmark$\checkmark^\checkmark\\checkmarkc\checkmarki\checkmarkr\checkmarkc\checkmark$\checkmarkd\checkmark$\checkmark^\checkmark\\checkmarkc\checkmarki\checkmarkr\checkmarkc\checkmark$\checkmarke\checkmark$\checkmark^\checkmark\\checkmarkc\checkmarki\checkmarkr\checkmarkc\checkmark$\checkmark \checkmark$\checkmark^\checkmark\\checkmarkc\checkmarki\checkmarkr\checkmarkc\checkmark$\checkmarkt\checkmark$\checkmark^\checkmark\\checkmarkc\checkmarki\checkmarkr\checkmarkc\checkmark$\checkmarkr\checkmark$\checkmark^\checkmark\\checkmarkc\checkmarki\checkmarkr\checkmarkc\checkmark$\checkmarka\checkmark$\checkmark^\checkmark\\checkmarkc\checkmarki\checkmarkr\checkmarkc\checkmark$\checkmarkv\checkmark$\checkmark^\checkmark\\checkmarkc\checkmarki\checkmarkr\checkmarkc\checkmark$\checkmarka\checkmark$\checkmark^\checkmark\\checkmarkc\checkmarki\checkmarkr\checkmarkc\checkmark$\checkmarki\checkmark$\checkmark^\checkmark\\checkmarkc\checkmarki\checkmarkr\checkmarkc\checkmark$\checkmarkl\checkmark$\checkmark^\checkmark\\checkmarkc\checkmarki\checkmarkr\checkmarkc\checkmark$\checkmark \checkmark$\checkmark^\checkmark\\checkmarkc\checkmarki\checkmarkr\checkmarkc\checkmark$\checkmarkY\checkmark$\checkmark^\checkmark\\checkmarkc\checkmarki\checkmarkr\checkmarkc\checkmark$\checkmark \checkmark$\checkmark^\checkmark\\checkmarkc\checkmarki\checkmarkr\checkmarkc\checkmark$\checkmark&\checkmark$\checkmark^\checkmark\\checkmarkc\checkmarki\checkmarkr\checkmarkc\checkmark$\checkmark \checkmark$\checkmark^\checkmark\\checkmarkc\checkmarki\checkmarkr\checkmarkc\checkmark$\checkmark$\checkmark$\checkmark^\checkmark\\checkmarkc\checkmarki\checkmarkr\checkmarkc\checkmark$\checkmark2\checkmark$\checkmark^\checkmark\\checkmarkc\checkmarki\checkmarkr\checkmarkc\checkmark$\checkmark0\checkmark$\checkmark^\checkmark\\checkmarkc\checkmarki\checkmarkr\checkmarkc\checkmark$\checkmark0\checkmark$\checkmark^\checkmark\\checkmarkc\checkmarki\checkmarkr\checkmarkc\checkmark$\checkmark$\checkmark$\checkmark^\checkmark\\checkmarkc\checkmarki\checkmarkr\checkmarkc\checkmark$\checkmark \checkmark$\checkmark^\checkmark\\checkmarkc\checkmarki\checkmarkr\checkmarkc\checkmark$\checkmarkm\checkmark$\checkmark^\checkmark\\checkmarkc\checkmarki\checkmarkr\checkmarkc\checkmark$\checkmarkm\checkmark$\checkmark^\checkmark\\checkmarkc\checkmarki\checkmarkr\checkmarkc\checkmark$\checkmark \checkmark$\checkmark^\checkmark\\checkmarkc\checkmarki\checkmarkr\checkmarkc\checkmark$\checkmark&\checkmark$\checkmark^\checkmark\\checkmarkc\checkmarki\checkmarkr\checkmarkc\checkmark$\checkmark \checkmark$\checkmark^\checkmark\\checkmarkc\checkmarki\checkmarkr\checkmarkc\checkmark$\checkmark$\checkmark$\checkmark^\checkmark\\checkmarkc\checkmarki\checkmarkr\checkmarkc\checkmark$\checkmark\\checkmark$\checkmark^\checkmark\\checkmarkc\checkmarki\checkmarkr\checkmarkc\checkmark$\checkmarkp\checkmark$\checkmark^\checkmark\\checkmarkc\checkmarki\checkmarkr\checkmarkc\checkmark$\checkmarkm\checkmark$\checkmark^\checkmark\\checkmarkc\checkmarki\checkmarkr\checkmarkc\checkmark$\checkmark \checkmark$\checkmark^\checkmark\\checkmarkc\checkmarki\checkmarkr\checkmarkc\checkmark$\checkmark5\checkmark$\checkmark^\checkmark\\checkmarkc\checkmarki\checkmarkr\checkmarkc\checkmark$\checkmark$\checkmark$\checkmark^\checkmark\\checkmarkc\checkmarki\checkmarkr\checkmarkc\checkmark$\checkmark \checkmark$\checkmark^\checkmark\\checkmarkc\checkmarki\checkmarkr\checkmarkc\checkmark$\checkmarkm\checkmark$\checkmark^\checkmark\\checkmarkc\checkmarki\checkmarkr\checkmarkc\checkmark$\checkmarkm\checkmark$\checkmark^\checkmark\\checkmarkc\checkmarki\checkmarkr\checkmarkc\checkmark$\checkmark \checkmark$\checkmark^\checkmark\\checkmarkc\checkmarki\checkmarkr\checkmarkc\checkmark$\checkmark\\checkmark$\checkmark^\checkmark\\checkmarkc\checkmarki\checkmarkr\checkmarkc\checkmark$\checkmark\\checkmark$\checkmark^\checkmark\\checkmarkc\checkmarki\checkmarkr\checkmarkc\checkmark$\checkmark \checkmark$\checkmark^\checkmark\\checkmarkc\checkmarki\checkmarkr\checkmarkc\checkmark$\checkmark\\checkmark$\checkmark^\checkmark\\checkmarkc\checkmarki\checkmarkr\checkmarkc\checkmark$\checkmarkc\checkmark$\checkmark^\checkmark\\checkmarkc\checkmarki\checkmarkr\checkmarkc\checkmark$\checkmarkl\checkmark$\checkmark^\checkmark\\checkmarkc\checkmarki\checkmarkr\checkmarkc\checkmark$\checkmarki\checkmark$\checkmark^\checkmark\\checkmarkc\checkmarki\checkmarkr\checkmarkc\checkmark$\checkmarkn\checkmark$\checkmark^\checkmark\\checkmarkc\checkmarki\checkmarkr\checkmarkc\checkmark$\checkmarke\checkmark$\checkmark^\checkmark\\checkmarkc\checkmarki\checkmarkr\checkmarkc\checkmark$\checkmark{\checkmark$\checkmark^\checkmark\\checkmarkc\checkmarki\checkmarkr\checkmarkc\checkmark$\checkmark2\checkmark$\checkmark^\checkmark\\checkmarkc\checkmarki\checkmarkr\checkmarkc\checkmark$\checkmark-\checkmark$\checkmark^\checkmark\\checkmarkc\checkmarki\checkmarkr\checkmarkc\checkmark$\checkmark4\checkmark$\checkmark^\checkmark\\checkmarkc\checkmarki\checkmarkr\checkmarkc\checkmark$\checkmark}\checkmark$\checkmark^\checkmark\\checkmarkc\checkmarki\checkmarkr\checkmarkc\checkmark$\checkmark
\checkmark$\checkmark^\checkmark\\checkmarkc\checkmarki\checkmarkr\checkmarkc\checkmark$\checkmark \checkmark$\checkmark^\checkmark\\checkmarkc\checkmarki\checkmarkr\checkmarkc\checkmark$\checkmark&\checkmark$\checkmark^\checkmark\\checkmarkc\checkmarki\checkmarkr\checkmarkc\checkmark$\checkmark \checkmark$\checkmark^\checkmark\\checkmarkc\checkmarki\checkmarkr\checkmarkc\checkmark$\checkmarkZ\checkmark$\checkmark^\checkmark\\checkmarkc\checkmarki\checkmarkr\checkmarkc\checkmark$\checkmarko\checkmark$\checkmark^\checkmark\\checkmarkc\checkmarki\checkmarkr\checkmarkc\checkmark$\checkmarkn\checkmark$\checkmark^\checkmark\\checkmarkc\checkmarki\checkmarkr\checkmarkc\checkmark$\checkmarke\checkmark$\checkmark^\checkmark\\checkmarkc\checkmarki\checkmarkr\checkmarkc\checkmark$\checkmark \checkmark$\checkmark^\checkmark\\checkmarkc\checkmarki\checkmarkr\checkmarkc\checkmark$\checkmarkd\checkmark$\checkmark^\checkmark\\checkmarkc\checkmarki\checkmarkr\checkmarkc\checkmark$\checkmarke\checkmark$\checkmark^\checkmark\\checkmarkc\checkmarki\checkmarkr\checkmarkc\checkmark$\checkmark \checkmark$\checkmark^\checkmark\\checkmarkc\checkmarki\checkmarkr\checkmarkc\checkmark$\checkmarkt\checkmark$\checkmark^\checkmark\\checkmarkc\checkmarki\checkmarkr\checkmarkc\checkmark$\checkmarkr\checkmark$\checkmark^\checkmark\\checkmarkc\checkmarki\checkmarkr\checkmarkc\checkmark$\checkmarka\checkmark$\checkmark^\checkmark\\checkmarkc\checkmarki\checkmarkr\checkmarkc\checkmark$\checkmarkv\checkmark$\checkmark^\checkmark\\checkmarkc\checkmarki\checkmarkr\checkmarkc\checkmark$\checkmarka\checkmark$\checkmark^\checkmark\\checkmarkc\checkmarki\checkmarkr\checkmarkc\checkmark$\checkmarki\checkmark$\checkmark^\checkmark\\checkmarkc\checkmarki\checkmarkr\checkmarkc\checkmark$\checkmarkl\checkmark$\checkmark^\checkmark\\checkmarkc\checkmarki\checkmarkr\checkmarkc\checkmark$\checkmark \checkmark$\checkmark^\checkmark\\checkmarkc\checkmarki\checkmarkr\checkmarkc\checkmark$\checkmarkZ\checkmark$\checkmark^\checkmark\\checkmarkc\checkmarki\checkmarkr\checkmarkc\checkmark$\checkmark \checkmark$\checkmark^\checkmark\\checkmarkc\checkmarki\checkmarkr\checkmarkc\checkmark$\checkmark&\checkmark$\checkmark^\checkmark\\checkmarkc\checkmarki\checkmarkr\checkmarkc\checkmark$\checkmark \checkmark$\checkmark^\checkmark\\checkmarkc\checkmarki\checkmarkr\checkmarkc\checkmark$\checkmark$\checkmark$\checkmark^\checkmark\\checkmarkc\checkmarki\checkmarkr\checkmarkc\checkmark$\checkmark1\checkmark$\checkmark^\checkmark\\checkmarkc\checkmarki\checkmarkr\checkmarkc\checkmark$\checkmark5\checkmark$\checkmark^\checkmark\\checkmarkc\checkmarki\checkmarkr\checkmarkc\checkmark$\checkmark0\checkmark$\checkmark^\checkmark\\checkmarkc\checkmarki\checkmarkr\checkmarkc\checkmark$\checkmark$\checkmark$\checkmark^\checkmark\\checkmarkc\checkmarki\checkmarkr\checkmarkc\checkmark$\checkmark \checkmark$\checkmark^\checkmark\\checkmarkc\checkmarki\checkmarkr\checkmarkc\checkmark$\checkmarkm\checkmark$\checkmark^\checkmark\\checkmarkc\checkmarki\checkmarkr\checkmarkc\checkmark$\checkmarkm\checkmark$\checkmark^\checkmark\\checkmarkc\checkmarki\checkmarkr\checkmarkc\checkmark$\checkmark \checkmark$\checkmark^\checkmark\\checkmarkc\checkmarki\checkmarkr\checkmarkc\checkmark$\checkmark&\checkmark$\checkmark^\checkmark\\checkmarkc\checkmarki\checkmarkr\checkmarkc\checkmark$\checkmark \checkmark$\checkmark^\checkmark\\checkmarkc\checkmarki\checkmarkr\checkmarkc\checkmark$\checkmark$\checkmark$\checkmark^\checkmark\\checkmarkc\checkmarki\checkmarkr\checkmarkc\checkmark$\checkmark\\checkmark$\checkmark^\checkmark\\checkmarkc\checkmarki\checkmarkr\checkmarkc\checkmark$\checkmarkp\checkmark$\checkmark^\checkmark\\checkmarkc\checkmarki\checkmarkr\checkmarkc\checkmark$\checkmarkm\checkmark$\checkmark^\checkmark\\checkmarkc\checkmarki\checkmarkr\checkmarkc\checkmark$\checkmark \checkmark$\checkmark^\checkmark\\checkmarkc\checkmarki\checkmarkr\checkmarkc\checkmark$\checkmark5\checkmark$\checkmark^\checkmark\\checkmarkc\checkmarki\checkmarkr\checkmarkc\checkmark$\checkmark$\checkmark$\checkmark^\checkmark\\checkmarkc\checkmarki\checkmarkr\checkmarkc\checkmark$\checkmark \checkmark$\checkmark^\checkmark\\checkmarkc\checkmarki\checkmarkr\checkmarkc\checkmark$\checkmarkm\checkmark$\checkmark^\checkmark\\checkmarkc\checkmarki\checkmarkr\checkmarkc\checkmark$\checkmarkm\checkmark$\checkmark^\checkmark\\checkmarkc\checkmarki\checkmarkr\checkmarkc\checkmark$\checkmark \checkmark$\checkmark^\checkmark\\checkmarkc\checkmarki\checkmarkr\checkmarkc\checkmark$\checkmark\\checkmark$\checkmark^\checkmark\\checkmarkc\checkmarki\checkmarkr\checkmarkc\checkmark$\checkmark\\checkmark$\checkmark^\checkmark\\checkmarkc\checkmarki\checkmarkr\checkmarkc\checkmark$\checkmark \checkmark$\checkmark^\checkmark\\checkmarkc\checkmarki\checkmarkr\checkmarkc\checkmark$\checkmark\\checkmark$\checkmark^\checkmark\\checkmarkc\checkmarki\checkmarkr\checkmarkc\checkmark$\checkmarkc\checkmark$\checkmark^\checkmark\\checkmarkc\checkmarki\checkmarkr\checkmarkc\checkmark$\checkmarkl\checkmark$\checkmark^\checkmark\\checkmarkc\checkmarki\checkmarkr\checkmarkc\checkmark$\checkmarki\checkmark$\checkmark^\checkmark\\checkmarkc\checkmarki\checkmarkr\checkmarkc\checkmark$\checkmarkn\checkmark$\checkmark^\checkmark\\checkmarkc\checkmarki\checkmarkr\checkmarkc\checkmark$\checkmarke\checkmark$\checkmark^\checkmark\\checkmarkc\checkmarki\checkmarkr\checkmarkc\checkmark$\checkmark{\checkmark$\checkmark^\checkmark\\checkmarkc\checkmarki\checkmarkr\checkmarkc\checkmark$\checkmark2\checkmark$\checkmark^\checkmark\\checkmarkc\checkmarki\checkmarkr\checkmarkc\checkmark$\checkmark-\checkmark$\checkmark^\checkmark\\checkmarkc\checkmarki\checkmarkr\checkmarkc\checkmark$\checkmark4\checkmark$\checkmark^\checkmark\\checkmarkc\checkmarki\checkmarkr\checkmarkc\checkmark$\checkmark}\checkmark$\checkmark^\checkmark\\checkmarkc\checkmarki\checkmarkr\checkmarkc\checkmark$\checkmark
\checkmark$\checkmark^\checkmark\\checkmarkc\checkmarki\checkmarkr\checkmarkc\checkmark$\checkmark \checkmark$\checkmark^\checkmark\\checkmarkc\checkmarki\checkmarkr\checkmarkc\checkmark$\checkmark&\checkmark$\checkmark^\checkmark\\checkmarkc\checkmarki\checkmarkr\checkmarkc\checkmark$\checkmark \checkmark$\checkmark^\checkmark\\checkmarkc\checkmarki\checkmarkr\checkmarkc\checkmark$\checkmarkP\checkmark$\checkmark^\checkmark\\checkmarkc\checkmarki\checkmarkr\checkmarkc\checkmark$\checkmarkr\checkmark$\checkmark^\checkmark\\checkmarkc\checkmarki\checkmarkr\checkmarkc\checkmark$\checkmarké\checkmark$\checkmark^\checkmark\\checkmarkc\checkmarki\checkmarkr\checkmarkc\checkmark$\checkmarkc\checkmark$\checkmark^\checkmark\\checkmarkc\checkmarki\checkmarkr\checkmarkc\checkmark$\checkmarki\checkmark$\checkmark^\checkmark\\checkmarkc\checkmarki\checkmarkr\checkmarkc\checkmark$\checkmarks\checkmark$\checkmark^\checkmark\\checkmarkc\checkmarki\checkmarkr\checkmarkc\checkmark$\checkmarki\checkmark$\checkmark^\checkmark\\checkmarkc\checkmarki\checkmarkr\checkmarkc\checkmark$\checkmarko\checkmark$\checkmark^\checkmark\\checkmarkc\checkmarki\checkmarkr\checkmarkc\checkmark$\checkmarkn\checkmark$\checkmark^\checkmark\\checkmarkc\checkmarki\checkmarkr\checkmarkc\checkmark$\checkmark \checkmark$\checkmark^\checkmark\\checkmarkc\checkmarki\checkmarkr\checkmarkc\checkmark$\checkmarkl\checkmark$\checkmark^\checkmark\\checkmarkc\checkmarki\checkmarkr\checkmarkc\checkmark$\checkmarki\checkmark$\checkmark^\checkmark\\checkmarkc\checkmarki\checkmarkr\checkmarkc\checkmark$\checkmarkn\checkmark$\checkmark^\checkmark\\checkmarkc\checkmarki\checkmarkr\checkmarkc\checkmark$\checkmarké\checkmark$\checkmark^\checkmark\\checkmarkc\checkmarki\checkmarkr\checkmarkc\checkmark$\checkmarka\checkmark$\checkmark^\checkmark\\checkmarkc\checkmarki\checkmarkr\checkmarkc\checkmark$\checkmarki\checkmark$\checkmark^\checkmark\\checkmarkc\checkmarki\checkmarkr\checkmarkc\checkmark$\checkmarkr\checkmark$\checkmark^\checkmark\\checkmarkc\checkmarki\checkmarkr\checkmarkc\checkmark$\checkmarke\checkmark$\checkmark^\checkmark\\checkmarkc\checkmarki\checkmarkr\checkmarkc\checkmark$\checkmark \checkmark$\checkmark^\checkmark\\checkmarkc\checkmarki\checkmarkr\checkmarkc\checkmark$\checkmark&\checkmark$\checkmark^\checkmark\\checkmarkc\checkmarki\checkmarkr\checkmarkc\checkmark$\checkmark \checkmark$\checkmark^\checkmark\\checkmarkc\checkmarki\checkmarkr\checkmarkc\checkmark$\checkmark$\checkmark$\checkmark^\checkmark\\checkmarkc\checkmarki\checkmarkr\checkmarkc\checkmark$\checkmark\\checkmark$\checkmark^\checkmark\\checkmarkc\checkmarki\checkmarkr\checkmarkc\checkmark$\checkmarkp\checkmark$\checkmark^\checkmark\\checkmarkc\checkmarki\checkmarkr\checkmarkc\checkmark$\checkmarkm\checkmark$\checkmark^\checkmark\\checkmarkc\checkmarki\checkmarkr\checkmarkc\checkmark$\checkmark \checkmark$\checkmark^\checkmark\\checkmarkc\checkmarki\checkmarkr\checkmarkc\checkmark$\checkmark0\checkmark$\checkmark^\checkmark\\checkmarkc\checkmarki\checkmarkr\checkmarkc\checkmark$\checkmark{\checkmark$\checkmark^\checkmark\\checkmarkc\checkmarki\checkmarkr\checkmarkc\checkmark$\checkmark,\checkmark$\checkmark^\checkmark\\checkmarkc\checkmarki\checkmarkr\checkmarkc\checkmark$\checkmark}\checkmark$\checkmark^\checkmark\\checkmarkc\checkmarki\checkmarkr\checkmarkc\checkmark$\checkmark0\checkmark$\checkmark^\checkmark\\checkmarkc\checkmarki\checkmarkr\checkmarkc\checkmark$\checkmark5\checkmark$\checkmark^\checkmark\\checkmarkc\checkmarki\checkmarkr\checkmarkc\checkmark$\checkmark$\checkmark$\checkmark^\checkmark\\checkmarkc\checkmarki\checkmarkr\checkmarkc\checkmark$\checkmark \checkmark$\checkmark^\checkmark\\checkmarkc\checkmarki\checkmarkr\checkmarkc\checkmark$\checkmarkm\checkmark$\checkmark^\checkmark\\checkmarkc\checkmarki\checkmarkr\checkmarkc\checkmark$\checkmarkm\checkmark$\checkmark^\checkmark\\checkmarkc\checkmarki\checkmarkr\checkmarkc\checkmark$\checkmark \checkmark$\checkmark^\checkmark\\checkmarkc\checkmarki\checkmarkr\checkmarkc\checkmark$\checkmark&\checkmark$\checkmark^\checkmark\\checkmarkc\checkmarki\checkmarkr\checkmarkc\checkmark$\checkmark \checkmark$\checkmark^\checkmark\\checkmarkc\checkmarki\checkmarkr\checkmarkc\checkmark$\checkmarkN\checkmark$\checkmark^\checkmark\\checkmarkc\checkmarki\checkmarkr\checkmarkc\checkmark$\checkmarko\checkmark$\checkmark^\checkmark\\checkmarkc\checkmarki\checkmarkr\checkmarkc\checkmark$\checkmarkn\checkmark$\checkmark^\checkmark\\checkmarkc\checkmarki\checkmarkr\checkmarkc\checkmark$\checkmark \checkmark$\checkmark^\checkmark\\checkmarkc\checkmarki\checkmarkr\checkmarkc\checkmark$\checkmarkn\checkmark$\checkmark^\checkmark\\checkmarkc\checkmarki\checkmarkr\checkmarkc\checkmark$\checkmarké\checkmark$\checkmark^\checkmark\\checkmarkc\checkmarki\checkmarkr\checkmarkc\checkmark$\checkmarkg\checkmark$\checkmark^\checkmark\\checkmarkc\checkmarki\checkmarkr\checkmarkc\checkmark$\checkmarko\checkmark$\checkmark^\checkmark\\checkmarkc\checkmarki\checkmarkr\checkmarkc\checkmark$\checkmarkc\checkmark$\checkmark^\checkmark\\checkmarkc\checkmarki\checkmarkr\checkmarkc\checkmark$\checkmarki\checkmark$\checkmark^\checkmark\\checkmarkc\checkmarki\checkmarkr\checkmarkc\checkmark$\checkmarka\checkmark$\checkmark^\checkmark\\checkmarkc\checkmarki\checkmarkr\checkmarkc\checkmark$\checkmarkb\checkmark$\checkmark^\checkmark\\checkmarkc\checkmarki\checkmarkr\checkmarkc\checkmark$\checkmarkl\checkmark$\checkmark^\checkmark\\checkmarkc\checkmarki\checkmarkr\checkmarkc\checkmark$\checkmarke\checkmark$\checkmark^\checkmark\\checkmarkc\checkmarki\checkmarkr\checkmarkc\checkmark$\checkmark \checkmark$\checkmark^\checkmark\\checkmarkc\checkmarki\checkmarkr\checkmarkc\checkmark$\checkmark\\checkmark$\checkmark^\checkmark\\checkmarkc\checkmarki\checkmarkr\checkmarkc\checkmark$\checkmark\\checkmark$\checkmark^\checkmark\\checkmarkc\checkmarki\checkmarkr\checkmarkc\checkmark$\checkmark \checkmark$\checkmark^\checkmark\\checkmarkc\checkmarki\checkmarkr\checkmarkc\checkmark$\checkmark\\checkmark$\checkmark^\checkmark\\checkmarkc\checkmarki\checkmarkr\checkmarkc\checkmark$\checkmarkh\checkmark$\checkmark^\checkmark\\checkmarkc\checkmarki\checkmarkr\checkmarkc\checkmark$\checkmarkl\checkmark$\checkmark^\checkmark\\checkmarkc\checkmarki\checkmarkr\checkmarkc\checkmark$\checkmarki\checkmark$\checkmark^\checkmark\\checkmarkc\checkmarki\checkmarkr\checkmarkc\checkmark$\checkmarkn\checkmark$\checkmark^\checkmark\\checkmarkc\checkmarki\checkmarkr\checkmarkc\checkmark$\checkmarke\checkmark$\checkmark^\checkmark\\checkmarkc\checkmarki\checkmarkr\checkmarkc\checkmark$\checkmark
\checkmark$\checkmark^\checkmark\\checkmarkc\checkmarki\checkmarkr\checkmarkc\checkmark$\checkmark\\checkmark$\checkmark^\checkmark\\checkmarkc\checkmarki\checkmarkr\checkmarkc\checkmark$\checkmarkm\checkmark$\checkmark^\checkmark\\checkmarkc\checkmarki\checkmarkr\checkmarkc\checkmark$\checkmarku\checkmark$\checkmark^\checkmark\\checkmarkc\checkmarki\checkmarkr\checkmarkc\checkmark$\checkmarkl\checkmark$\checkmark^\checkmark\\checkmarkc\checkmarki\checkmarkr\checkmarkc\checkmark$\checkmarkt\checkmark$\checkmark^\checkmark\\checkmarkc\checkmarki\checkmarkr\checkmarkc\checkmark$\checkmarki\checkmark$\checkmark^\checkmark\\checkmarkc\checkmarki\checkmarkr\checkmarkc\checkmark$\checkmarkr\checkmark$\checkmark^\checkmark\\checkmarkc\checkmarki\checkmarkr\checkmarkc\checkmark$\checkmarko\checkmark$\checkmark^\checkmark\\checkmarkc\checkmarki\checkmarkr\checkmarkc\checkmark$\checkmarkw\checkmark$\checkmark^\checkmark\\checkmarkc\checkmarki\checkmarkr\checkmarkc\checkmark$\checkmark{\checkmark$\checkmark^\checkmark\\checkmarkc\checkmarki\checkmarkr\checkmarkc\checkmark$\checkmark2\checkmark$\checkmark^\checkmark\\checkmarkc\checkmarki\checkmarkr\checkmarkc\checkmark$\checkmark}\checkmark$\checkmark^\checkmark\\checkmarkc\checkmarki\checkmarkr\checkmarkc\checkmark$\checkmark{\checkmark$\checkmark^\checkmark\\checkmarkc\checkmarki\checkmarkr\checkmarkc\checkmark$\checkmark*\checkmark$\checkmark^\checkmark\\checkmarkc\checkmarki\checkmarkr\checkmarkc\checkmark$\checkmark}\checkmark$\checkmark^\checkmark\\checkmarkc\checkmarki\checkmarkr\checkmarkc\checkmark$\checkmark{\checkmark$\checkmark^\checkmark\\checkmarkc\checkmarki\checkmarkr\checkmarkc\checkmark$\checkmark\\checkmark$\checkmark^\checkmark\\checkmarkc\checkmarki\checkmarkr\checkmarkc\checkmark$\checkmarkt\checkmark$\checkmark^\checkmark\\checkmarkc\checkmarki\checkmarkr\checkmarkc\checkmark$\checkmarke\checkmark$\checkmark^\checkmark\\checkmarkc\checkmarki\checkmarkr\checkmarkc\checkmark$\checkmarkx\checkmark$\checkmark^\checkmark\\checkmarkc\checkmarki\checkmarkr\checkmarkc\checkmark$\checkmarkt\checkmark$\checkmark^\checkmark\\checkmarkc\checkmarki\checkmarkr\checkmarkc\checkmark$\checkmarkb\checkmark$\checkmark^\checkmark\\checkmarkc\checkmarki\checkmarkr\checkmarkc\checkmark$\checkmarkf\checkmark$\checkmark^\checkmark\\checkmarkc\checkmarki\checkmarkr\checkmarkc\checkmark$\checkmark{\checkmark$\checkmark^\checkmark\\checkmarkc\checkmarki\checkmarkr\checkmarkc\checkmark$\checkmarkF\checkmark$\checkmark^\checkmark\\checkmarkc\checkmarki\checkmarkr\checkmarkc\checkmark$\checkmarkC\checkmark$\checkmark^\checkmark\\checkmarkc\checkmarki\checkmarkr\checkmarkc\checkmark$\checkmark1\checkmark$\checkmark^\checkmark\\checkmarkc\checkmarki\checkmarkr\checkmarkc\checkmark$\checkmark \checkmark$\checkmark^\checkmark\\checkmarkc\checkmarki\checkmarkr\checkmarkc\checkmark$\checkmark(\checkmark$\checkmark^\checkmark\\checkmarkc\checkmarki\checkmarkr\checkmarkc\checkmark$\checkmarkC\checkmark$\checkmark^\checkmark\\checkmarkc\checkmarki\checkmarkr\checkmarkc\checkmark$\checkmarki\checkmark$\checkmark^\checkmark\\checkmarkc\checkmarki\checkmarkr\checkmarkc\checkmark$\checkmarkn\checkmark$\checkmark^\checkmark\\checkmarkc\checkmarki\checkmarkr\checkmarkc\checkmark$\checkmarké\checkmark$\checkmark^\checkmark\\checkmarkc\checkmarki\checkmarkr\checkmarkc\checkmark$\checkmarkm\checkmark$\checkmark^\checkmark\\checkmarkc\checkmarki\checkmarkr\checkmarkc\checkmark$\checkmarka\checkmark$\checkmark^\checkmark\\checkmarkc\checkmarki\checkmarkr\checkmarkc\checkmark$\checkmarkt\checkmark$\checkmark^\checkmark\\checkmarkc\checkmarki\checkmarkr\checkmarkc\checkmark$\checkmarki\checkmark$\checkmark^\checkmark\\checkmarkc\checkmarki\checkmarkr\checkmarkc\checkmark$\checkmarkq\checkmark$\checkmark^\checkmark\\checkmarkc\checkmarki\checkmarkr\checkmarkc\checkmark$\checkmarku\checkmark$\checkmark^\checkmark\\checkmarkc\checkmarki\checkmarkr\checkmarkc\checkmark$\checkmarke\checkmark$\checkmark^\checkmark\\checkmarkc\checkmarki\checkmarkr\checkmarkc\checkmark$\checkmark)\checkmark$\checkmark^\checkmark\\checkmarkc\checkmarki\checkmarkr\checkmarkc\checkmark$\checkmark}\checkmark$\checkmark^\checkmark\\checkmarkc\checkmarki\checkmarkr\checkmarkc\checkmark$\checkmark}\checkmark$\checkmark^\checkmark\\checkmarkc\checkmarki\checkmarkr\checkmarkc\checkmark$\checkmark \checkmark$\checkmark^\checkmark\\checkmarkc\checkmarki\checkmarkr\checkmarkc\checkmark$\checkmark&\checkmark$\checkmark^\checkmark\\checkmarkc\checkmarki\checkmarkr\checkmarkc\checkmark$\checkmark \checkmark$\checkmark^\checkmark\\checkmarkc\checkmarki\checkmarkr\checkmarkc\checkmark$\checkmarkC\checkmark$\checkmark^\checkmark\\checkmarkc\checkmarki\checkmarkr\checkmarkc\checkmark$\checkmarko\checkmark$\checkmark^\checkmark\\checkmarkc\checkmarki\checkmarkr\checkmarkc\checkmark$\checkmarku\checkmark$\checkmark^\checkmark\\checkmarkc\checkmarki\checkmarkr\checkmarkc\checkmark$\checkmarkr\checkmark$\checkmark^\checkmark\\checkmarkc\checkmarki\checkmarkr\checkmarkc\checkmark$\checkmarks\checkmark$\checkmark^\checkmark\\checkmarkc\checkmarki\checkmarkr\checkmarkc\checkmark$\checkmarke\checkmark$\checkmark^\checkmark\\checkmarkc\checkmarki\checkmarkr\checkmarkc\checkmark$\checkmark \checkmark$\checkmark^\checkmark\\checkmarkc\checkmarki\checkmarkr\checkmarkc\checkmark$\checkmarka\checkmark$\checkmark^\checkmark\\checkmarkc\checkmarki\checkmarkr\checkmarkc\checkmark$\checkmarkn\checkmark$\checkmark^\checkmark\\checkmarkc\checkmarki\checkmarkr\checkmarkc\checkmark$\checkmarkg\checkmark$\checkmark^\checkmark\\checkmarkc\checkmarki\checkmarkr\checkmarkc\checkmark$\checkmarku\checkmark$\checkmark^\checkmark\\checkmarkc\checkmarki\checkmarkr\checkmarkc\checkmark$\checkmarkl\checkmark$\checkmark^\checkmark\\checkmarkc\checkmarki\checkmarkr\checkmarkc\checkmark$\checkmarka\checkmark$\checkmark^\checkmark\\checkmarkc\checkmarki\checkmarkr\checkmarkc\checkmark$\checkmarki\checkmark$\checkmark^\checkmark\\checkmarkc\checkmarki\checkmarkr\checkmarkc\checkmark$\checkmarkr\checkmark$\checkmark^\checkmark\\checkmarkc\checkmarki\checkmarkr\checkmarkc\checkmark$\checkmarke\checkmark$\checkmark^\checkmark\\checkmarkc\checkmarki\checkmarkr\checkmarkc\checkmark$\checkmark \checkmark$\checkmark^\checkmark\\checkmarkc\checkmarki\checkmarkr\checkmarkc\checkmark$\checkmarkA\checkmark$\checkmark^\checkmark\\checkmarkc\checkmarki\checkmarkr\checkmarkc\checkmark$\checkmarkx\checkmark$\checkmark^\checkmark\\checkmarkc\checkmarki\checkmarkr\checkmarkc\checkmark$\checkmarke\checkmark$\checkmark^\checkmark\\checkmarkc\checkmarki\checkmarkr\checkmarkc\checkmark$\checkmark \checkmark$\checkmark^\checkmark\\checkmarkc\checkmarki\checkmarkr\checkmarkc\checkmark$\checkmarkA\checkmark$\checkmark^\checkmark\\checkmarkc\checkmarki\checkmarkr\checkmarkc\checkmark$\checkmark \checkmark$\checkmark^\checkmark\\checkmarkc\checkmarki\checkmarkr\checkmarkc\checkmark$\checkmark(\checkmark$\checkmark^\checkmark\\checkmarkc\checkmarki\checkmarkr\checkmarkc\checkmark$\checkmarkT\checkmark$\checkmark^\checkmark\\checkmarkc\checkmarki\checkmarkr\checkmarkc\checkmark$\checkmarki\checkmark$\checkmark^\checkmark\\checkmarkc\checkmarki\checkmarkr\checkmarkc\checkmark$\checkmarkl\checkmark$\checkmark^\checkmark\\checkmarkc\checkmarki\checkmarkr\checkmarkc\checkmark$\checkmarkt\checkmark$\checkmark^\checkmark\\checkmarkc\checkmarki\checkmarkr\checkmarkc\checkmark$\checkmark)\checkmark$\checkmark^\checkmark\\checkmarkc\checkmarki\checkmarkr\checkmarkc\checkmark$\checkmark \checkmark$\checkmark^\checkmark\\checkmarkc\checkmarki\checkmarkr\checkmarkc\checkmark$\checkmark&\checkmark$\checkmark^\checkmark\\checkmarkc\checkmarki\checkmarkr\checkmarkc\checkmark$\checkmark \checkmark$\checkmark^\checkmark\\checkmarkc\checkmarki\checkmarkr\checkmarkc\checkmark$\checkmark$\checkmark$\checkmark^\checkmark\\checkmarkc\checkmarki\checkmarkr\checkmarkc\checkmark$\checkmark\\checkmark$\checkmark^\checkmark\\checkmarkc\checkmarki\checkmarkr\checkmarkc\checkmark$\checkmarkp\checkmark$\checkmark^\checkmark\\checkmarkc\checkmarki\checkmarkr\checkmarkc\checkmark$\checkmarkm\checkmark$\checkmark^\checkmark\\checkmarkc\checkmarki\checkmarkr\checkmarkc\checkmark$\checkmark \checkmark$\checkmark^\checkmark\\checkmarkc\checkmarki\checkmarkr\checkmarkc\checkmark$\checkmark9\checkmark$\checkmark^\checkmark\\checkmarkc\checkmarki\checkmarkr\checkmarkc\checkmark$\checkmark0\checkmark$\checkmark^\checkmark\\checkmarkc\checkmarki\checkmarkr\checkmarkc\checkmark$\checkmark°\checkmark$\checkmark^\checkmark\\checkmarkc\checkmarki\checkmarkr\checkmarkc\checkmark$\checkmark$\checkmark$\checkmark^\checkmark\\checkmarkc\checkmarki\checkmarkr\checkmarkc\checkmark$\checkmark \checkmark$\checkmark^\checkmark\\checkmarkc\checkmarki\checkmarkr\checkmarkc\checkmark$\checkmark&\checkmark$\checkmark^\checkmark\\checkmarkc\checkmarki\checkmarkr\checkmarkc\checkmark$\checkmark \checkmark$\checkmark^\checkmark\\checkmarkc\checkmarki\checkmarkr\checkmarkc\checkmark$\checkmark$\checkmark$\checkmark^\checkmark\\checkmarkc\checkmarki\checkmarkr\checkmarkc\checkmark$\checkmark\\checkmark$\checkmark^\checkmark\\checkmarkc\checkmarki\checkmarkr\checkmarkc\checkmark$\checkmarkp\checkmark$\checkmark^\checkmark\\checkmarkc\checkmarki\checkmarkr\checkmarkc\checkmark$\checkmarkm\checkmark$\checkmark^\checkmark\\checkmarkc\checkmarki\checkmarkr\checkmarkc\checkmark$\checkmark \checkmark$\checkmark^\checkmark\\checkmarkc\checkmarki\checkmarkr\checkmarkc\checkmark$\checkmark5\checkmark$\checkmark^\checkmark\\checkmarkc\checkmarki\checkmarkr\checkmarkc\checkmark$\checkmark°\checkmark$\checkmark^\checkmark\\checkmarkc\checkmarki\checkmarkr\checkmarkc\checkmark$\checkmark$\checkmark$\checkmark^\checkmark\\checkmarkc\checkmarki\checkmarkr\checkmarkc\checkmark$\checkmark \checkmark$\checkmark^\checkmark\\checkmarkc\checkmarki\checkmarkr\checkmarkc\checkmark$\checkmark\\checkmark$\checkmark^\checkmark\\checkmarkc\checkmarki\checkmarkr\checkmarkc\checkmark$\checkmark\\checkmark$\checkmark^\checkmark\\checkmarkc\checkmarki\checkmarkr\checkmarkc\checkmark$\checkmark \checkmark$\checkmark^\checkmark\\checkmarkc\checkmarki\checkmarkr\checkmarkc\checkmark$\checkmark\\checkmark$\checkmark^\checkmark\\checkmarkc\checkmarki\checkmarkr\checkmarkc\checkmark$\checkmarkc\checkmark$\checkmark^\checkmark\\checkmarkc\checkmarki\checkmarkr\checkmarkc\checkmark$\checkmarkl\checkmark$\checkmark^\checkmark\\checkmarkc\checkmarki\checkmarkr\checkmarkc\checkmark$\checkmarki\checkmark$\checkmark^\checkmark\\checkmarkc\checkmarki\checkmarkr\checkmarkc\checkmark$\checkmarkn\checkmark$\checkmark^\checkmark\\checkmarkc\checkmarki\checkmarkr\checkmarkc\checkmark$\checkmarke\checkmark$\checkmark^\checkmark\\checkmarkc\checkmarki\checkmarkr\checkmarkc\checkmark$\checkmark{\checkmark$\checkmark^\checkmark\\checkmarkc\checkmarki\checkmarkr\checkmarkc\checkmark$\checkmark2\checkmark$\checkmark^\checkmark\\checkmarkc\checkmarki\checkmarkr\checkmarkc\checkmark$\checkmark-\checkmark$\checkmark^\checkmark\\checkmarkc\checkmarki\checkmarkr\checkmarkc\checkmark$\checkmark4\checkmark$\checkmark^\checkmark\\checkmarkc\checkmarki\checkmarkr\checkmarkc\checkmark$\checkmark}\checkmark$\checkmark^\checkmark\\checkmarkc\checkmarki\checkmarkr\checkmarkc\checkmark$\checkmark
\checkmark$\checkmark^\checkmark\\checkmarkc\checkmarki\checkmarkr\checkmarkc\checkmark$\checkmark \checkmark$\checkmark^\checkmark\\checkmarkc\checkmarki\checkmarkr\checkmarkc\checkmark$\checkmark&\checkmark$\checkmark^\checkmark\\checkmarkc\checkmarki\checkmarkr\checkmarkc\checkmark$\checkmark \checkmark$\checkmark^\checkmark\\checkmarkc\checkmarki\checkmarkr\checkmarkc\checkmark$\checkmarkC\checkmark$\checkmark^\checkmark\\checkmarkc\checkmarki\checkmarkr\checkmarkc\checkmark$\checkmarko\checkmark$\checkmark^\checkmark\\checkmarkc\checkmarki\checkmarkr\checkmarkc\checkmark$\checkmarku\checkmark$\checkmark^\checkmark\\checkmarkc\checkmarki\checkmarkr\checkmarkc\checkmark$\checkmarkr\checkmark$\checkmark^\checkmark\\checkmarkc\checkmarki\checkmarkr\checkmarkc\checkmark$\checkmarks\checkmark$\checkmark^\checkmark\\checkmarkc\checkmarki\checkmarkr\checkmarkc\checkmark$\checkmarke\checkmark$\checkmark^\checkmark\\checkmarkc\checkmarki\checkmarkr\checkmarkc\checkmark$\checkmark \checkmark$\checkmark^\checkmark\\checkmarkc\checkmarki\checkmarkr\checkmarkc\checkmark$\checkmarka\checkmark$\checkmark^\checkmark\\checkmarkc\checkmarki\checkmarkr\checkmarkc\checkmark$\checkmarkn\checkmark$\checkmark^\checkmark\\checkmarkc\checkmarki\checkmarkr\checkmarkc\checkmark$\checkmarkg\checkmark$\checkmark^\checkmark\\checkmarkc\checkmarki\checkmarkr\checkmarkc\checkmark$\checkmarku\checkmark$\checkmark^\checkmark\\checkmarkc\checkmarki\checkmarkr\checkmarkc\checkmark$\checkmarkl\checkmark$\checkmark^\checkmark\\checkmarkc\checkmarki\checkmarkr\checkmarkc\checkmark$\checkmarka\checkmark$\checkmark^\checkmark\\checkmarkc\checkmarki\checkmarkr\checkmarkc\checkmark$\checkmarki\checkmark$\checkmark^\checkmark\\checkmarkc\checkmarki\checkmarkr\checkmarkc\checkmark$\checkmarkr\checkmark$\checkmark^\checkmark\\checkmarkc\checkmarki\checkmarkr\checkmarkc\checkmark$\checkmarke\checkmark$\checkmark^\checkmark\\checkmarkc\checkmarki\checkmarkr\checkmarkc\checkmark$\checkmark \checkmark$\checkmark^\checkmark\\checkmarkc\checkmarki\checkmarkr\checkmarkc\checkmark$\checkmarkA\checkmark$\checkmark^\checkmark\\checkmarkc\checkmarki\checkmarkr\checkmarkc\checkmark$\checkmarkx\checkmark$\checkmark^\checkmark\\checkmarkc\checkmarki\checkmarkr\checkmarkc\checkmark$\checkmarke\checkmark$\checkmark^\checkmark\\checkmarkc\checkmarki\checkmarkr\checkmarkc\checkmark$\checkmark \checkmark$\checkmark^\checkmark\\checkmarkc\checkmarki\checkmarkr\checkmarkc\checkmark$\checkmarkB\checkmark$\checkmark^\checkmark\\checkmarkc\checkmarki\checkmarkr\checkmarkc\checkmark$\checkmark \checkmark$\checkmark^\checkmark\\checkmarkc\checkmarki\checkmarkr\checkmarkc\checkmark$\checkmark(\checkmark$\checkmark^\checkmark\\checkmarkc\checkmarki\checkmarkr\checkmarkc\checkmark$\checkmarkP\checkmark$\checkmark^\checkmark\\checkmarkc\checkmarki\checkmarkr\checkmarkc\checkmark$\checkmarka\checkmark$\checkmark^\checkmark\\checkmarkc\checkmarki\checkmarkr\checkmarkc\checkmark$\checkmarkn\checkmark$\checkmark^\checkmark\\checkmarkc\checkmarki\checkmarkr\checkmarkc\checkmark$\checkmark)\checkmark$\checkmark^\checkmark\\checkmarkc\checkmarki\checkmarkr\checkmarkc\checkmark$\checkmark \checkmark$\checkmark^\checkmark\\checkmarkc\checkmarki\checkmarkr\checkmarkc\checkmark$\checkmark&\checkmark$\checkmark^\checkmark\\checkmarkc\checkmarki\checkmarkr\checkmarkc\checkmark$\checkmark \checkmark$\checkmark^\checkmark\\checkmarkc\checkmarki\checkmarkr\checkmarkc\checkmark$\checkmark$\checkmark$\checkmark^\checkmark\\checkmarkc\checkmarki\checkmarkr\checkmarkc\checkmark$\checkmark3\checkmark$\checkmark^\checkmark\\checkmarkc\checkmarki\checkmarkr\checkmarkc\checkmark$\checkmark6\checkmark$\checkmark^\checkmark\\checkmarkc\checkmarki\checkmarkr\checkmarkc\checkmark$\checkmark0\checkmark$\checkmark^\checkmark\\checkmarkc\checkmarki\checkmarkr\checkmarkc\checkmark$\checkmark°\checkmark$\checkmark^\checkmark\\checkmarkc\checkmarki\checkmarkr\checkmarkc\checkmark$\checkmark$\checkmark$\checkmark^\checkmark\\checkmarkc\checkmarki\checkmarkr\checkmarkc\checkmark$\checkmark \checkmark$\checkmark^\checkmark\\checkmarkc\checkmarki\checkmarkr\checkmarkc\checkmark$\checkmark&\checkmark$\checkmark^\checkmark\\checkmarkc\checkmarki\checkmarkr\checkmarkc\checkmark$\checkmark \checkmark$\checkmark^\checkmark\\checkmarkc\checkmarki\checkmarkr\checkmarkc\checkmark$\checkmark$\checkmark$\checkmark^\checkmark\\checkmarkc\checkmarki\checkmarkr\checkmarkc\checkmark$\checkmark0\checkmark$\checkmark^\checkmark\\checkmarkc\checkmarki\checkmarkr\checkmarkc\checkmark$\checkmark°\checkmark$\checkmark^\checkmark\\checkmarkc\checkmarki\checkmarkr\checkmarkc\checkmark$\checkmark$\checkmark$\checkmark^\checkmark\\checkmarkc\checkmarki\checkmarkr\checkmarkc\checkmark$\checkmark \checkmark$\checkmark^\checkmark\\checkmarkc\checkmarki\checkmarkr\checkmarkc\checkmark$\checkmark\\checkmark$\checkmark^\checkmark\\checkmarkc\checkmarki\checkmarkr\checkmarkc\checkmark$\checkmark\\checkmark$\checkmark^\checkmark\\checkmarkc\checkmarki\checkmarkr\checkmarkc\checkmark$\checkmark \checkmark$\checkmark^\checkmark\\checkmarkc\checkmarki\checkmarkr\checkmarkc\checkmark$\checkmark\\checkmark$\checkmark^\checkmark\\checkmarkc\checkmarki\checkmarkr\checkmarkc\checkmark$\checkmarkh\checkmark$\checkmark^\checkmark\\checkmarkc\checkmarki\checkmarkr\checkmarkc\checkmark$\checkmarkl\checkmark$\checkmark^\checkmark\\checkmarkc\checkmarki\checkmarkr\checkmarkc\checkmark$\checkmarki\checkmark$\checkmark^\checkmark\\checkmarkc\checkmarki\checkmarkr\checkmarkc\checkmark$\checkmarkn\checkmark$\checkmark^\checkmark\\checkmarkc\checkmarki\checkmarkr\checkmarkc\checkmark$\checkmarke\checkmark$\checkmark^\checkmark\\checkmarkc\checkmarki\checkmarkr\checkmarkc\checkmark$\checkmark
\checkmark$\checkmark^\checkmark\\checkmarkc\checkmarki\checkmarkr\checkmarkc\checkmark$\checkmark\\checkmark$\checkmark^\checkmark\\checkmarkc\checkmarki\checkmarkr\checkmarkc\checkmark$\checkmarkt\checkmark$\checkmark^\checkmark\\checkmarkc\checkmarki\checkmarkr\checkmarkc\checkmark$\checkmarke\checkmark$\checkmark^\checkmark\\checkmarkc\checkmarki\checkmarkr\checkmarkc\checkmark$\checkmarkx\checkmark$\checkmark^\checkmark\\checkmarkc\checkmarki\checkmarkr\checkmarkc\checkmark$\checkmarkt\checkmark$\checkmark^\checkmark\\checkmarkc\checkmarki\checkmarkr\checkmarkc\checkmark$\checkmarkb\checkmark$\checkmark^\checkmark\\checkmarkc\checkmarki\checkmarkr\checkmarkc\checkmark$\checkmarkf\checkmark$\checkmark^\checkmark\\checkmarkc\checkmarki\checkmarkr\checkmarkc\checkmark$\checkmark{\checkmark$\checkmark^\checkmark\\checkmarkc\checkmarki\checkmarkr\checkmarkc\checkmark$\checkmarkF\checkmark$\checkmark^\checkmark\\checkmarkc\checkmarki\checkmarkr\checkmarkc\checkmark$\checkmarkC\checkmark$\checkmark^\checkmark\\checkmarkc\checkmarki\checkmarkr\checkmarkc\checkmark$\checkmark2\checkmark$\checkmark^\checkmark\\checkmarkc\checkmarki\checkmarkr\checkmarkc\checkmark$\checkmark \checkmark$\checkmark^\checkmark\\checkmarkc\checkmarki\checkmarkr\checkmarkc\checkmark$\checkmark(\checkmark$\checkmark^\checkmark\\checkmarkc\checkmarki\checkmarkr\checkmarkc\checkmark$\checkmarkM\checkmark$\checkmark^\checkmark\\checkmarkc\checkmarki\checkmarkr\checkmarkc\checkmark$\checkmarka\checkmark$\checkmark^\checkmark\\checkmarkc\checkmarki\checkmarkr\checkmarkc\checkmark$\checkmarkt\checkmark$\checkmark^\checkmark\\checkmarkc\checkmarki\checkmarkr\checkmarkc\checkmark$\checkmarké\checkmark$\checkmark^\checkmark\\checkmarkc\checkmarki\checkmarkr\checkmarkc\checkmark$\checkmarkr\checkmark$\checkmark^\checkmark\\checkmarkc\checkmarki\checkmarkr\checkmarkc\checkmark$\checkmarki\checkmark$\checkmark^\checkmark\\checkmarkc\checkmarki\checkmarkr\checkmarkc\checkmark$\checkmarka\checkmark$\checkmark^\checkmark\\checkmarkc\checkmarki\checkmarkr\checkmarkc\checkmark$\checkmarku\checkmark$\checkmark^\checkmark\\checkmarkc\checkmarki\checkmarkr\checkmarkc\checkmark$\checkmarkx\checkmark$\checkmark^\checkmark\\checkmarkc\checkmarki\checkmarkr\checkmarkc\checkmark$\checkmark)\checkmark$\checkmark^\checkmark\\checkmarkc\checkmarki\checkmarkr\checkmarkc\checkmark$\checkmark}\checkmark$\checkmark^\checkmark\\checkmarkc\checkmarki\checkmarkr\checkmarkc\checkmark$\checkmark \checkmark$\checkmark^\checkmark\\checkmarkc\checkmarki\checkmarkr\checkmarkc\checkmark$\checkmark&\checkmark$\checkmark^\checkmark\\checkmarkc\checkmarki\checkmarkr\checkmarkc\checkmark$\checkmark \checkmark$\checkmark^\checkmark\\checkmarkc\checkmarki\checkmarkr\checkmarkc\checkmark$\checkmarkD\checkmark$\checkmark^\checkmark\\checkmarkc\checkmarki\checkmarkr\checkmarkc\checkmark$\checkmarku\checkmark$\checkmark^\checkmark\\checkmarkc\checkmarki\checkmarkr\checkmarkc\checkmark$\checkmarkr\checkmark$\checkmark^\checkmark\\checkmarkc\checkmarki\checkmarkr\checkmarkc\checkmark$\checkmarke\checkmark$\checkmark^\checkmark\\checkmarkc\checkmarki\checkmarkr\checkmarkc\checkmark$\checkmarkt\checkmark$\checkmark^\checkmark\\checkmarkc\checkmarki\checkmarkr\checkmarkc\checkmark$\checkmarké\checkmark$\checkmark^\checkmark\\checkmarkc\checkmarki\checkmarkr\checkmarkc\checkmark$\checkmark \checkmark$\checkmark^\checkmark\\checkmarkc\checkmarki\checkmarkr\checkmarkc\checkmark$\checkmarkm\checkmark$\checkmark^\checkmark\\checkmarkc\checkmarki\checkmarkr\checkmarkc\checkmark$\checkmarka\checkmark$\checkmark^\checkmark\\checkmarkc\checkmarki\checkmarkr\checkmarkc\checkmark$\checkmarkx\checkmark$\checkmark^\checkmark\\checkmarkc\checkmarki\checkmarkr\checkmarkc\checkmark$\checkmark \checkmark$\checkmark^\checkmark\\checkmarkc\checkmarki\checkmarkr\checkmarkc\checkmark$\checkmarkm\checkmark$\checkmark^\checkmark\\checkmarkc\checkmarki\checkmarkr\checkmarkc\checkmark$\checkmarka\checkmark$\checkmark^\checkmark\\checkmarkc\checkmarki\checkmarkr\checkmarkc\checkmark$\checkmarkt\checkmark$\checkmark^\checkmark\\checkmarkc\checkmarki\checkmarkr\checkmarkc\checkmark$\checkmarki\checkmark$\checkmark^\checkmark\\checkmarkc\checkmarki\checkmarkr\checkmarkc\checkmark$\checkmarkè\checkmark$\checkmark^\checkmark\\checkmarkc\checkmarki\checkmarkr\checkmarkc\checkmark$\checkmarkr\checkmark$\checkmark^\checkmark\\checkmarkc\checkmarki\checkmarkr\checkmarkc\checkmark$\checkmarke\checkmark$\checkmark^\checkmark\\checkmarkc\checkmarki\checkmarkr\checkmarkc\checkmark$\checkmark \checkmark$\checkmark^\checkmark\\checkmarkc\checkmarki\checkmarkr\checkmarkc\checkmark$\checkmarkà\checkmark$\checkmark^\checkmark\\checkmarkc\checkmarki\checkmarkr\checkmarkc\checkmark$\checkmark \checkmark$\checkmark^\checkmark\\checkmarkc\checkmarki\checkmarkr\checkmarkc\checkmark$\checkmarku\checkmark$\checkmark^\checkmark\\checkmarkc\checkmarki\checkmarkr\checkmarkc\checkmark$\checkmarks\checkmark$\checkmark^\checkmark\\checkmarkc\checkmarki\checkmarkr\checkmarkc\checkmark$\checkmarki\checkmark$\checkmark^\checkmark\\checkmarkc\checkmarki\checkmarkr\checkmarkc\checkmark$\checkmarkn\checkmark$\checkmark^\checkmark\\checkmarkc\checkmarki\checkmarkr\checkmarkc\checkmark$\checkmarke\checkmark$\checkmark^\checkmark\\checkmarkc\checkmarki\checkmarkr\checkmarkc\checkmark$\checkmarkr\checkmark$\checkmark^\checkmark\\checkmarkc\checkmarki\checkmarkr\checkmarkc\checkmark$\checkmark \checkmark$\checkmark^\checkmark\\checkmarkc\checkmarki\checkmarkr\checkmarkc\checkmark$\checkmark&\checkmark$\checkmark^\checkmark\\checkmarkc\checkmarki\checkmarkr\checkmarkc\checkmark$\checkmark \checkmark$\checkmark^\checkmark\\checkmarkc\checkmarki\checkmarkr\checkmarkc\checkmark$\checkmark$\checkmark$\checkmark^\checkmark\\checkmarkc\checkmarki\checkmarkr\checkmarkc\checkmark$\checkmark1\checkmark$\checkmark^\checkmark\\checkmarkc\checkmarki\checkmarkr\checkmarkc\checkmark$\checkmark1\checkmark$\checkmark^\checkmark\\checkmarkc\checkmarki\checkmarkr\checkmarkc\checkmark$\checkmark0\checkmark$\checkmark^\checkmark\\checkmarkc\checkmarki\checkmarkr\checkmarkc\checkmark$\checkmark$\checkmark$\checkmark^\checkmark\\checkmarkc\checkmarki\checkmarkr\checkmarkc\checkmark$\checkmark \checkmark$\checkmark^\checkmark\\checkmarkc\checkmarki\checkmarkr\checkmarkc\checkmark$\checkmarkH\checkmark$\checkmark^\checkmark\\checkmarkc\checkmarki\checkmarkr\checkmarkc\checkmark$\checkmarkB\checkmark$\checkmark^\checkmark\\checkmarkc\checkmarki\checkmarkr\checkmarkc\checkmark$\checkmark \checkmark$\checkmark^\checkmark\\checkmarkc\checkmarki\checkmarkr\checkmarkc\checkmark$\checkmark(\checkmark$\checkmark^\checkmark\\checkmarkc\checkmarki\checkmarkr\checkmarkc\checkmark$\checkmarkA\checkmark$\checkmark^\checkmark\\checkmarkc\checkmarki\checkmarkr\checkmarkc\checkmark$\checkmarkl\checkmark$\checkmark^\checkmark\\checkmarkc\checkmarki\checkmarkr\checkmarkc\checkmark$\checkmark \checkmark$\checkmark^\checkmark\\checkmarkc\checkmarki\checkmarkr\checkmarkc\checkmark$\checkmark2\checkmark$\checkmark^\checkmark\\checkmarkc\checkmarki\checkmarkr\checkmarkc\checkmark$\checkmark0\checkmark$\checkmark^\checkmark\\checkmarkc\checkmarki\checkmarkr\checkmarkc\checkmark$\checkmark1\checkmark$\checkmark^\checkmark\\checkmarkc\checkmarki\checkmarkr\checkmarkc\checkmark$\checkmark7\checkmark$\checkmark^\checkmark\\checkmarkc\checkmarki\checkmarkr\checkmarkc\checkmark$\checkmark)\checkmark$\checkmark^\checkmark\\checkmarkc\checkmarki\checkmarkr\checkmarkc\checkmark$\checkmark \checkmark$\checkmark^\checkmark\\checkmarkc\checkmarki\checkmarkr\checkmarkc\checkmark$\checkmark&\checkmark$\checkmark^\checkmark\\checkmarkc\checkmarki\checkmarkr\checkmarkc\checkmark$\checkmark \checkmark$\checkmark^\checkmark\\checkmarkc\checkmarki\checkmarkr\checkmarkc\checkmark$\checkmarkN\checkmark$\checkmark^\checkmark\\checkmarkc\checkmarki\checkmarkr\checkmarkc\checkmark$\checkmarko\checkmark$\checkmark^\checkmark\\checkmarkc\checkmarki\checkmarkr\checkmarkc\checkmark$\checkmarkn\checkmark$\checkmark^\checkmark\\checkmarkc\checkmarki\checkmarkr\checkmarkc\checkmark$\checkmark \checkmark$\checkmark^\checkmark\\checkmarkc\checkmarki\checkmarkr\checkmarkc\checkmark$\checkmarkn\checkmark$\checkmark^\checkmark\\checkmarkc\checkmarki\checkmarkr\checkmarkc\checkmark$\checkmarké\checkmark$\checkmark^\checkmark\\checkmarkc\checkmarki\checkmarkr\checkmarkc\checkmark$\checkmarkg\checkmark$\checkmark^\checkmark\\checkmarkc\checkmarki\checkmarkr\checkmarkc\checkmark$\checkmarko\checkmark$\checkmark^\checkmark\\checkmarkc\checkmarki\checkmarkr\checkmarkc\checkmark$\checkmarkc\checkmark$\checkmark^\checkmark\\checkmarkc\checkmarki\checkmarkr\checkmarkc\checkmark$\checkmarki\checkmark$\checkmark^\checkmark\\checkmarkc\checkmarki\checkmarkr\checkmarkc\checkmark$\checkmarka\checkmark$\checkmark^\checkmark\\checkmarkc\checkmarki\checkmarkr\checkmarkc\checkmark$\checkmarkb\checkmark$\checkmark^\checkmark\\checkmarkc\checkmarki\checkmarkr\checkmarkc\checkmark$\checkmarkl\checkmark$\checkmark^\checkmark\\checkmarkc\checkmarki\checkmarkr\checkmarkc\checkmark$\checkmarke\checkmark$\checkmark^\checkmark\\checkmarkc\checkmarki\checkmarkr\checkmarkc\checkmark$\checkmark \checkmark$\checkmark^\checkmark\\checkmarkc\checkmarki\checkmarkr\checkmarkc\checkmark$\checkmark\\checkmark$\checkmark^\checkmark\\checkmarkc\checkmarki\checkmarkr\checkmarkc\checkmark$\checkmark\\checkmark$\checkmark^\checkmark\\checkmarkc\checkmarki\checkmarkr\checkmarkc\checkmark$\checkmark \checkmark$\checkmark^\checkmark\\checkmarkc\checkmarki\checkmarkr\checkmarkc\checkmark$\checkmark\\checkmark$\checkmark^\checkmark\\checkmarkc\checkmarki\checkmarkr\checkmarkc\checkmark$\checkmarkh\checkmark$\checkmark^\checkmark\\checkmarkc\checkmarki\checkmarkr\checkmarkc\checkmark$\checkmarkl\checkmark$\checkmark^\checkmark\\checkmarkc\checkmarki\checkmarkr\checkmarkc\checkmark$\checkmarki\checkmark$\checkmark^\checkmark\\checkmarkc\checkmarki\checkmarkr\checkmarkc\checkmark$\checkmarkn\checkmark$\checkmark^\checkmark\\checkmarkc\checkmarki\checkmarkr\checkmarkc\checkmark$\checkmarke\checkmark$\checkmark^\checkmark\\checkmarkc\checkmarki\checkmarkr\checkmarkc\checkmark$\checkmark
\checkmark$\checkmark^\checkmark\\checkmarkc\checkmarki\checkmarkr\checkmarkc\checkmark$\checkmark\\checkmark$\checkmark^\checkmark\\checkmarkc\checkmarki\checkmarkr\checkmarkc\checkmark$\checkmarkt\checkmark$\checkmark^\checkmark\\checkmarkc\checkmarki\checkmarkr\checkmarkc\checkmark$\checkmarke\checkmark$\checkmark^\checkmark\\checkmarkc\checkmarki\checkmarkr\checkmarkc\checkmark$\checkmarkx\checkmark$\checkmark^\checkmark\\checkmarkc\checkmarki\checkmarkr\checkmarkc\checkmark$\checkmarkt\checkmark$\checkmark^\checkmark\\checkmarkc\checkmarki\checkmarkr\checkmarkc\checkmark$\checkmarkb\checkmark$\checkmark^\checkmark\\checkmarkc\checkmarki\checkmarkr\checkmarkc\checkmark$\checkmarkf\checkmark$\checkmark^\checkmark\\checkmarkc\checkmarki\checkmarkr\checkmarkc\checkmark$\checkmark{\checkmark$\checkmark^\checkmark\\checkmarkc\checkmarki\checkmarkr\checkmarkc\checkmark$\checkmarkF\checkmark$\checkmark^\checkmark\\checkmarkc\checkmarki\checkmarkr\checkmarkc\checkmark$\checkmarkC\checkmark$\checkmark^\checkmark\\checkmarkc\checkmarki\checkmarkr\checkmarkc\checkmark$\checkmark3\checkmark$\checkmark^\checkmark\\checkmarkc\checkmarki\checkmarkr\checkmarkc\checkmark$\checkmark \checkmark$\checkmark^\checkmark\\checkmarkc\checkmarki\checkmarkr\checkmarkc\checkmark$\checkmark(\checkmark$\checkmark^\checkmark\\checkmarkc\checkmarki\checkmarkr\checkmarkc\checkmark$\checkmarkR\checkmark$\checkmark^\checkmark\\checkmarkc\checkmarki\checkmarkr\checkmarkc\checkmark$\checkmarké\checkmark$\checkmark^\checkmark\\checkmarkc\checkmarki\checkmarkr\checkmarkc\checkmark$\checkmarks\checkmark$\checkmark^\checkmark\\checkmarkc\checkmarki\checkmarkr\checkmarkc\checkmark$\checkmarki\checkmark$\checkmark^\checkmark\\checkmarkc\checkmarki\checkmarkr\checkmarkc\checkmark$\checkmarks\checkmark$\checkmark^\checkmark\\checkmarkc\checkmarki\checkmarkr\checkmarkc\checkmark$\checkmarkt\checkmark$\checkmark^\checkmark\\checkmarkc\checkmarki\checkmarkr\checkmarkc\checkmark$\checkmarka\checkmark$\checkmark^\checkmark\\checkmarkc\checkmarki\checkmarkr\checkmarkc\checkmark$\checkmarkn\checkmark$\checkmark^\checkmark\\checkmarkc\checkmarki\checkmarkr\checkmarkc\checkmark$\checkmarkc\checkmark$\checkmark^\checkmark\\checkmarkc\checkmarki\checkmarkr\checkmarkc\checkmark$\checkmarke\checkmark$\checkmark^\checkmark\\checkmarkc\checkmarki\checkmarkr\checkmarkc\checkmark$\checkmark)\checkmark$\checkmark^\checkmark\\checkmarkc\checkmarki\checkmarkr\checkmarkc\checkmark$\checkmark}\checkmark$\checkmark^\checkmark\\checkmarkc\checkmarki\checkmarkr\checkmarkc\checkmark$\checkmark \checkmark$\checkmark^\checkmark\\checkmarkc\checkmarki\checkmarkr\checkmarkc\checkmark$\checkmark&\checkmark$\checkmark^\checkmark\\checkmarkc\checkmarki\checkmarkr\checkmarkc\checkmark$\checkmark \checkmark$\checkmark^\checkmark\\checkmarkc\checkmarki\checkmarkr\checkmarkc\checkmark$\checkmarkC\checkmark$\checkmark^\checkmark\\checkmarkc\checkmarki\checkmarkr\checkmarkc\checkmark$\checkmarko\checkmark$\checkmark^\checkmark\\checkmarkc\checkmarki\checkmarkr\checkmarkc\checkmark$\checkmarke\checkmark$\checkmark^\checkmark\\checkmarkc\checkmarki\checkmarkr\checkmarkc\checkmark$\checkmarkf\checkmark$\checkmark^\checkmark\\checkmarkc\checkmarki\checkmarkr\checkmarkc\checkmark$\checkmarkf\checkmark$\checkmark^\checkmark\\checkmarkc\checkmarki\checkmarkr\checkmarkc\checkmark$\checkmarki\checkmark$\checkmark^\checkmark\\checkmarkc\checkmarki\checkmarkr\checkmarkc\checkmark$\checkmarkc\checkmark$\checkmark^\checkmark\\checkmarkc\checkmarki\checkmarkr\checkmarkc\checkmark$\checkmarki\checkmark$\checkmark^\checkmark\\checkmarkc\checkmarki\checkmarkr\checkmarkc\checkmark$\checkmarke\checkmark$\checkmark^\checkmark\\checkmarkc\checkmarki\checkmarkr\checkmarkc\checkmark$\checkmarkn\checkmark$\checkmark^\checkmark\\checkmarkc\checkmarki\checkmarkr\checkmarkc\checkmark$\checkmarkt\checkmark$\checkmark^\checkmark\\checkmarkc\checkmarki\checkmarkr\checkmarkc\checkmark$\checkmark \checkmark$\checkmark^\checkmark\\checkmarkc\checkmarki\checkmarkr\checkmarkc\checkmark$\checkmarkd\checkmark$\checkmark^\checkmark\\checkmarkc\checkmarki\checkmarkr\checkmarkc\checkmark$\checkmarke\checkmark$\checkmark^\checkmark\\checkmarkc\checkmarki\checkmarkr\checkmarkc\checkmark$\checkmark \checkmark$\checkmark^\checkmark\\checkmarkc\checkmarki\checkmarkr\checkmarkc\checkmark$\checkmarks\checkmark$\checkmark^\checkmark\\checkmarkc\checkmarki\checkmarkr\checkmarkc\checkmark$\checkmarké\checkmark$\checkmark^\checkmark\\checkmarkc\checkmarki\checkmarkr\checkmarkc\checkmark$\checkmarkc\checkmark$\checkmark^\checkmark\\checkmarkc\checkmarki\checkmarkr\checkmarkc\checkmark$\checkmarku\checkmark$\checkmark^\checkmark\\checkmarkc\checkmarki\checkmarkr\checkmarkc\checkmark$\checkmarkr\checkmark$\checkmark^\checkmark\\checkmarkc\checkmarki\checkmarkr\checkmarkc\checkmark$\checkmarki\checkmark$\checkmark^\checkmark\\checkmarkc\checkmarki\checkmarkr\checkmarkc\checkmark$\checkmarkt\checkmark$\checkmark^\checkmark\\checkmarkc\checkmarki\checkmarkr\checkmarkc\checkmark$\checkmarké\checkmark$\checkmark^\checkmark\\checkmarkc\checkmarki\checkmarkr\checkmarkc\checkmark$\checkmark \checkmark$\checkmark^\checkmark\\checkmarkc\checkmarki\checkmarkr\checkmarkc\checkmark$\checkmarks\checkmark$\checkmark^\checkmark\\checkmarkc\checkmarki\checkmarkr\checkmarkc\checkmark$\checkmarkt\checkmark$\checkmark^\checkmark\\checkmarkc\checkmarki\checkmarkr\checkmarkc\checkmark$\checkmarka\checkmark$\checkmark^\checkmark\\checkmarkc\checkmarki\checkmarkr\checkmarkc\checkmark$\checkmarkt\checkmark$\checkmark^\checkmark\\checkmarkc\checkmarki\checkmarkr\checkmarkc\checkmark$\checkmarki\checkmark$\checkmark^\checkmark\\checkmarkc\checkmarki\checkmarkr\checkmarkc\checkmark$\checkmarkq\checkmark$\checkmark^\checkmark\\checkmarkc\checkmarki\checkmarkr\checkmarkc\checkmark$\checkmarku\checkmark$\checkmark^\checkmark\\checkmarkc\checkmarki\checkmarkr\checkmarkc\checkmark$\checkmarke\checkmark$\checkmark^\checkmark\\checkmarkc\checkmarki\checkmarkr\checkmarkc\checkmark$\checkmark \checkmark$\checkmark^\checkmark\\checkmarkc\checkmarki\checkmarkr\checkmarkc\checkmark$\checkmark&\checkmark$\checkmark^\checkmark\\checkmarkc\checkmarki\checkmarkr\checkmarkc\checkmark$\checkmark \checkmark$\checkmark^\checkmark\\checkmarkc\checkmarki\checkmarkr\checkmarkc\checkmark$\checkmark$\checkmark$\checkmark^\checkmark\\checkmarkc\checkmarki\checkmarkr\checkmarkc\checkmark$\checkmarks\checkmark$\checkmark^\checkmark\\checkmarkc\checkmarki\checkmarkr\checkmarkc\checkmark$\checkmark \checkmark$\checkmark^\checkmark\\checkmarkc\checkmarki\checkmarkr\checkmarkc\checkmark$\checkmark\\checkmark$\checkmark^\checkmark\\checkmarkc\checkmarki\checkmarkr\checkmarkc\checkmark$\checkmarkg\checkmark$\checkmark^\checkmark\\checkmarkc\checkmarki\checkmarkr\checkmarkc\checkmark$\checkmarke\checkmark$\checkmark^\checkmark\\checkmarkc\checkmarki\checkmarkr\checkmarkc\checkmark$\checkmarkq\checkmark$\checkmark^\checkmark\\checkmarkc\checkmarki\checkmarkr\checkmarkc\checkmark$\checkmark \checkmark$\checkmark^\checkmark\\checkmarkc\checkmarki\checkmarkr\checkmarkc\checkmark$\checkmark1\checkmark$\checkmark^\checkmark\\checkmarkc\checkmarki\checkmarkr\checkmarkc\checkmark$\checkmark{\checkmark$\checkmark^\checkmark\\checkmarkc\checkmarki\checkmarkr\checkmarkc\checkmark$\checkmark,\checkmark$\checkmark^\checkmark\\checkmarkc\checkmarki\checkmarkr\checkmarkc\checkmark$\checkmark}\checkmark$\checkmark^\checkmark\\checkmarkc\checkmarki\checkmarkr\checkmarkc\checkmark$\checkmark5\checkmark$\checkmark^\checkmark\\checkmarkc\checkmarki\checkmarkr\checkmarkc\checkmark$\checkmark$\checkmark$\checkmark^\checkmark\\checkmarkc\checkmarki\checkmarkr\checkmarkc\checkmark$\checkmark \checkmark$\checkmark^\checkmark\\checkmarkc\checkmarki\checkmarkr\checkmarkc\checkmark$\checkmark&\checkmark$\checkmark^\checkmark\\checkmarkc\checkmarki\checkmarkr\checkmarkc\checkmark$\checkmark \checkmark$\checkmark^\checkmark\\checkmarkc\checkmarki\checkmarkr\checkmarkc\checkmark$\checkmarkN\checkmark$\checkmark^\checkmark\\checkmarkc\checkmarki\checkmarkr\checkmarkc\checkmark$\checkmarko\checkmark$\checkmark^\checkmark\\checkmarkc\checkmarki\checkmarkr\checkmarkc\checkmark$\checkmarkn\checkmark$\checkmark^\checkmark\\checkmarkc\checkmarki\checkmarkr\checkmarkc\checkmark$\checkmark \checkmark$\checkmark^\checkmark\\checkmarkc\checkmarki\checkmarkr\checkmarkc\checkmark$\checkmarkn\checkmark$\checkmark^\checkmark\\checkmarkc\checkmarki\checkmarkr\checkmarkc\checkmark$\checkmarké\checkmark$\checkmark^\checkmark\\checkmarkc\checkmarki\checkmarkr\checkmarkc\checkmark$\checkmarkg\checkmark$\checkmark^\checkmark\\checkmarkc\checkmarki\checkmarkr\checkmarkc\checkmark$\checkmarko\checkmark$\checkmark^\checkmark\\checkmarkc\checkmarki\checkmarkr\checkmarkc\checkmark$\checkmarkc\checkmark$\checkmark^\checkmark\\checkmarkc\checkmarki\checkmarkr\checkmarkc\checkmark$\checkmarki\checkmark$\checkmark^\checkmark\\checkmarkc\checkmarki\checkmarkr\checkmarkc\checkmark$\checkmarka\checkmark$\checkmark^\checkmark\\checkmarkc\checkmarki\checkmarkr\checkmarkc\checkmark$\checkmarkb\checkmark$\checkmark^\checkmark\\checkmarkc\checkmarki\checkmarkr\checkmarkc\checkmark$\checkmarkl\checkmark$\checkmark^\checkmark\\checkmarkc\checkmarki\checkmarkr\checkmarkc\checkmark$\checkmarke\checkmark$\checkmark^\checkmark\\checkmarkc\checkmarki\checkmarkr\checkmarkc\checkmark$\checkmark \checkmark$\checkmark^\checkmark\\checkmarkc\checkmarki\checkmarkr\checkmarkc\checkmark$\checkmark\\checkmark$\checkmark^\checkmark\\checkmarkc\checkmarki\checkmarkr\checkmarkc\checkmark$\checkmark\\checkmark$\checkmark^\checkmark\\checkmarkc\checkmarki\checkmarkr\checkmarkc\checkmark$\checkmark \checkmark$\checkmark^\checkmark\\checkmarkc\checkmarki\checkmarkr\checkmarkc\checkmark$\checkmark\\checkmark$\checkmark^\checkmark\\checkmarkc\checkmarki\checkmarkr\checkmarkc\checkmark$\checkmarkh\checkmark$\checkmark^\checkmark\\checkmarkc\checkmarki\checkmarkr\checkmarkc\checkmark$\checkmarkl\checkmark$\checkmark^\checkmark\\checkmarkc\checkmarki\checkmarkr\checkmarkc\checkmark$\checkmarki\checkmark$\checkmark^\checkmark\\checkmarkc\checkmarki\checkmarkr\checkmarkc\checkmark$\checkmarkn\checkmark$\checkmark^\checkmark\\checkmarkc\checkmarki\checkmarkr\checkmarkc\checkmark$\checkmarke\checkmark$\checkmark^\checkmark\\checkmarkc\checkmarki\checkmarkr\checkmarkc\checkmark$\checkmark
\checkmark$\checkmark^\checkmark\\checkmarkc\checkmarki\checkmarkr\checkmarkc\checkmark$\checkmark\\checkmark$\checkmark^\checkmark\\checkmarkc\checkmarki\checkmarkr\checkmarkc\checkmark$\checkmarkt\checkmark$\checkmark^\checkmark\\checkmarkc\checkmarki\checkmarkr\checkmarkc\checkmark$\checkmarke\checkmark$\checkmark^\checkmark\\checkmarkc\checkmarki\checkmarkr\checkmarkc\checkmark$\checkmarkx\checkmark$\checkmark^\checkmark\\checkmarkc\checkmarki\checkmarkr\checkmarkc\checkmark$\checkmarkt\checkmark$\checkmark^\checkmark\\checkmarkc\checkmarki\checkmarkr\checkmarkc\checkmark$\checkmarkb\checkmark$\checkmark^\checkmark\\checkmarkc\checkmarki\checkmarkr\checkmarkc\checkmark$\checkmarkf\checkmark$\checkmark^\checkmark\\checkmarkc\checkmarki\checkmarkr\checkmarkc\checkmark$\checkmark{\checkmark$\checkmark^\checkmark\\checkmarkc\checkmarki\checkmarkr\checkmarkc\checkmark$\checkmarkF\checkmark$\checkmark^\checkmark\\checkmarkc\checkmarki\checkmarkr\checkmarkc\checkmark$\checkmarkC\checkmark$\checkmark^\checkmark\\checkmarkc\checkmarki\checkmarkr\checkmarkc\checkmark$\checkmark4\checkmark$\checkmark^\checkmark\\checkmarkc\checkmarki\checkmarkr\checkmarkc\checkmark$\checkmark \checkmark$\checkmark^\checkmark\\checkmarkc\checkmarki\checkmarkr\checkmarkc\checkmark$\checkmark(\checkmark$\checkmark^\checkmark\\checkmarkc\checkmarki\checkmarkr\checkmarkc\checkmark$\checkmarkC\checkmark$\checkmark^\checkmark\\checkmarkc\checkmarki\checkmarkr\checkmarkc\checkmark$\checkmarko\checkmark$\checkmark^\checkmark\\checkmarkc\checkmarki\checkmarkr\checkmarkc\checkmark$\checkmarkm\checkmark$\checkmark^\checkmark\\checkmarkc\checkmarki\checkmarkr\checkmarkc\checkmark$\checkmarkm\checkmark$\checkmark^\checkmark\\checkmarkc\checkmarki\checkmarkr\checkmarkc\checkmark$\checkmarka\checkmark$\checkmark^\checkmark\\checkmarkc\checkmarki\checkmarkr\checkmarkc\checkmark$\checkmarkn\checkmark$\checkmark^\checkmark\\checkmarkc\checkmarki\checkmarkr\checkmarkc\checkmark$\checkmarkd\checkmark$\checkmark^\checkmark\\checkmarkc\checkmarki\checkmarkr\checkmarkc\checkmark$\checkmarke\checkmark$\checkmark^\checkmark\\checkmarkc\checkmarki\checkmarkr\checkmarkc\checkmark$\checkmark)\checkmark$\checkmark^\checkmark\\checkmarkc\checkmarki\checkmarkr\checkmarkc\checkmark$\checkmark}\checkmark$\checkmark^\checkmark\\checkmarkc\checkmarki\checkmarkr\checkmarkc\checkmark$\checkmark \checkmark$\checkmark^\checkmark\\checkmarkc\checkmarki\checkmarkr\checkmarkc\checkmark$\checkmark&\checkmark$\checkmark^\checkmark\\checkmarkc\checkmarki\checkmarkr\checkmarkc\checkmark$\checkmark \checkmark$\checkmark^\checkmark\\checkmarkc\checkmarki\checkmarkr\checkmarkc\checkmark$\checkmarkC\checkmark$\checkmark^\checkmark\\checkmarkc\checkmarki\checkmarkr\checkmarkc\checkmark$\checkmarko\checkmark$\checkmark^\checkmark\\checkmarkc\checkmarki\checkmarkr\checkmarkc\checkmark$\checkmarkn\checkmark$\checkmark^\checkmark\\checkmarkc\checkmarki\checkmarkr\checkmarkc\checkmark$\checkmarkt\checkmark$\checkmark^\checkmark\\checkmarkc\checkmarki\checkmarkr\checkmarkc\checkmark$\checkmarkr\checkmark$\checkmark^\checkmark\\checkmarkc\checkmarki\checkmarkr\checkmarkc\checkmark$\checkmarkô\checkmark$\checkmark^\checkmark\\checkmarkc\checkmarki\checkmarkr\checkmarkc\checkmark$\checkmarkl\checkmark$\checkmark^\checkmark\\checkmarkc\checkmarki\checkmarkr\checkmarkc\checkmark$\checkmarke\checkmark$\checkmark^\checkmark\\checkmarkc\checkmarki\checkmarkr\checkmarkc\checkmark$\checkmarku\checkmark$\checkmark^\checkmark\\checkmarkc\checkmarki\checkmarkr\checkmarkc\checkmark$\checkmarkr\checkmark$\checkmark^\checkmark\\checkmarkc\checkmarki\checkmarkr\checkmarkc\checkmark$\checkmark \checkmark$\checkmark^\checkmark\\checkmarkc\checkmarki\checkmarkr\checkmarkc\checkmark$\checkmark&\checkmark$\checkmark^\checkmark\\checkmarkc\checkmarki\checkmarkr\checkmarkc\checkmark$\checkmark \checkmark$\checkmark^\checkmark\\checkmarkc\checkmarki\checkmarkr\checkmarkc\checkmark$\checkmarkE\checkmark$\checkmark^\checkmark\\checkmarkc\checkmarki\checkmarkr\checkmarkc\checkmark$\checkmarkS\checkmark$\checkmark^\checkmark\\checkmarkc\checkmarki\checkmarkr\checkmarkc\checkmark$\checkmarkP\checkmark$\checkmark^\checkmark\\checkmarkc\checkmarki\checkmarkr\checkmarkc\checkmark$\checkmark3\checkmark$\checkmark^\checkmark\\checkmarkc\checkmarki\checkmarkr\checkmarkc\checkmark$\checkmark2\checkmark$\checkmark^\checkmark\\checkmarkc\checkmarki\checkmarkr\checkmarkc\checkmark$\checkmark \checkmark$\checkmark^\checkmark\\checkmarkc\checkmarki\checkmarkr\checkmarkc\checkmark$\checkmarkD\checkmark$\checkmark^\checkmark\\checkmarkc\checkmarki\checkmarkr\checkmarkc\checkmark$\checkmarku\checkmark$\checkmark^\checkmark\\checkmarkc\checkmarki\checkmarkr\checkmarkc\checkmark$\checkmarka\checkmark$\checkmark^\checkmark\\checkmarkc\checkmarki\checkmarkr\checkmarkc\checkmark$\checkmarkl\checkmark$\checkmark^\checkmark\\checkmarkc\checkmarki\checkmarkr\checkmarkc\checkmark$\checkmark-\checkmark$\checkmark^\checkmark\\checkmarkc\checkmarki\checkmarkr\checkmarkc\checkmark$\checkmarkC\checkmark$\checkmark^\checkmark\\checkmarkc\checkmarki\checkmarkr\checkmarkc\checkmark$\checkmarko\checkmark$\checkmark^\checkmark\\checkmarkc\checkmarki\checkmarkr\checkmarkc\checkmark$\checkmarkr\checkmark$\checkmark^\checkmark\\checkmarkc\checkmarki\checkmarkr\checkmarkc\checkmark$\checkmarke\checkmark$\checkmark^\checkmark\\checkmarkc\checkmarki\checkmarkr\checkmarkc\checkmark$\checkmark \checkmark$\checkmark^\checkmark\\checkmarkc\checkmarki\checkmarkr\checkmarkc\checkmark$\checkmark2\checkmark$\checkmark^\checkmark\\checkmarkc\checkmarki\checkmarkr\checkmarkc\checkmark$\checkmark4\checkmark$\checkmark^\checkmark\\checkmarkc\checkmarki\checkmarkr\checkmarkc\checkmark$\checkmark0\checkmark$\checkmark^\checkmark\\checkmarkc\checkmarki\checkmarkr\checkmarkc\checkmark$\checkmark \checkmark$\checkmark^\checkmark\\checkmarkc\checkmarki\checkmarkr\checkmarkc\checkmark$\checkmarkM\checkmark$\checkmark^\checkmark\\checkmarkc\checkmarki\checkmarkr\checkmarkc\checkmark$\checkmarkH\checkmark$\checkmark^\checkmark\\checkmarkc\checkmarki\checkmarkr\checkmarkc\checkmark$\checkmarkz\checkmark$\checkmark^\checkmark\\checkmarkc\checkmarki\checkmarkr\checkmarkc\checkmark$\checkmark \checkmark$\checkmark^\checkmark\\checkmarkc\checkmarki\checkmarkr\checkmarkc\checkmark$\checkmark&\checkmark$\checkmark^\checkmark\\checkmarkc\checkmarki\checkmarkr\checkmarkc\checkmark$\checkmark \checkmark$\checkmark^\checkmark\\checkmarkc\checkmarki\checkmarkr\checkmarkc\checkmark$\checkmark\\checkmark$\checkmark^\checkmark\\checkmarkc\checkmarki\checkmarkr\checkmarkc\checkmark$\checkmark\\checkmark$\checkmark^\checkmark\\checkmarkc\checkmarki\checkmarkr\checkmarkc\checkmark$\checkmark \checkmark$\checkmark^\checkmark\\checkmarkc\checkmarki\checkmarkr\checkmarkc\checkmark$\checkmark\\checkmark$\checkmark^\checkmark\\checkmarkc\checkmarki\checkmarkr\checkmarkc\checkmark$\checkmarkh\checkmark$\checkmark^\checkmark\\checkmarkc\checkmarki\checkmarkr\checkmarkc\checkmark$\checkmarkl\checkmark$\checkmark^\checkmark\\checkmarkc\checkmarki\checkmarkr\checkmarkc\checkmark$\checkmarki\checkmark$\checkmark^\checkmark\\checkmarkc\checkmarki\checkmarkr\checkmarkc\checkmark$\checkmarkn\checkmark$\checkmark^\checkmark\\checkmarkc\checkmarki\checkmarkr\checkmarkc\checkmark$\checkmarke\checkmark$\checkmark^\checkmark\\checkmarkc\checkmarki\checkmarkr\checkmarkc\checkmark$\checkmark
\checkmark$\checkmark^\checkmark\\checkmarkc\checkmarki\checkmarkr\checkmarkc\checkmark$\checkmark\\checkmark$\checkmark^\checkmark\\checkmarkc\checkmarki\checkmarkr\checkmarkc\checkmark$\checkmarkt\checkmark$\checkmark^\checkmark\\checkmarkc\checkmarki\checkmarkr\checkmarkc\checkmark$\checkmarke\checkmark$\checkmark^\checkmark\\checkmarkc\checkmarki\checkmarkr\checkmarkc\checkmark$\checkmarkx\checkmark$\checkmark^\checkmark\\checkmarkc\checkmarki\checkmarkr\checkmarkc\checkmark$\checkmarkt\checkmark$\checkmark^\checkmark\\checkmarkc\checkmarki\checkmarkr\checkmarkc\checkmark$\checkmarkb\checkmark$\checkmark^\checkmark\\checkmarkc\checkmarki\checkmarkr\checkmarkc\checkmark$\checkmarkf\checkmark$\checkmark^\checkmark\\checkmarkc\checkmarki\checkmarkr\checkmarkc\checkmark$\checkmark{\checkmark$\checkmark^\checkmark\\checkmarkc\checkmarki\checkmarkr\checkmarkc\checkmark$\checkmarkF\checkmark$\checkmark^\checkmark\\checkmarkc\checkmarki\checkmarkr\checkmarkc\checkmark$\checkmarkC\checkmark$\checkmark^\checkmark\\checkmarkc\checkmarki\checkmarkr\checkmarkc\checkmark$\checkmark5\checkmark$\checkmark^\checkmark\\checkmarkc\checkmarki\checkmarkr\checkmarkc\checkmark$\checkmark \checkmark$\checkmark^\checkmark\\checkmarkc\checkmarki\checkmarkr\checkmarkc\checkmark$\checkmark(\checkmark$\checkmark^\checkmark\\checkmarkc\checkmarki\checkmarkr\checkmarkc\checkmark$\checkmarkC\checkmark$\checkmark^\checkmark\\checkmarkc\checkmarki\checkmarkr\checkmarkc\checkmark$\checkmarko\checkmark$\checkmark^\checkmark\\checkmarkc\checkmarki\checkmarkr\checkmarkc\checkmark$\checkmarkm\checkmark$\checkmark^\checkmark\\checkmarkc\checkmarki\checkmarkr\checkmarkc\checkmark$\checkmarkm\checkmark$\checkmark^\checkmark\\checkmarkc\checkmarki\checkmarkr\checkmarkc\checkmark$\checkmarku\checkmark$\checkmark^\checkmark\\checkmarkc\checkmarki\checkmarkr\checkmarkc\checkmark$\checkmarkn\checkmark$\checkmark^\checkmark\\checkmarkc\checkmarki\checkmarkr\checkmarkc\checkmark$\checkmarki\checkmark$\checkmark^\checkmark\\checkmarkc\checkmarki\checkmarkr\checkmarkc\checkmark$\checkmarkc\checkmark$\checkmark^\checkmark\\checkmarkc\checkmarki\checkmarkr\checkmarkc\checkmark$\checkmarka\checkmark$\checkmark^\checkmark\\checkmarkc\checkmarki\checkmarkr\checkmarkc\checkmark$\checkmarkt\checkmark$\checkmark^\checkmark\\checkmarkc\checkmarki\checkmarkr\checkmarkc\checkmark$\checkmarki\checkmark$\checkmark^\checkmark\\checkmarkc\checkmarki\checkmarkr\checkmarkc\checkmark$\checkmarko\checkmark$\checkmark^\checkmark\\checkmarkc\checkmarki\checkmarkr\checkmarkc\checkmark$\checkmarkn\checkmark$\checkmark^\checkmark\\checkmarkc\checkmarki\checkmarkr\checkmarkc\checkmark$\checkmark)\checkmark$\checkmark^\checkmark\\checkmarkc\checkmarki\checkmarkr\checkmarkc\checkmark$\checkmark}\checkmark$\checkmark^\checkmark\\checkmarkc\checkmarki\checkmarkr\checkmarkc\checkmark$\checkmark \checkmark$\checkmark^\checkmark\\checkmarkc\checkmarki\checkmarkr\checkmarkc\checkmark$\checkmark&\checkmark$\checkmark^\checkmark\\checkmarkc\checkmarki\checkmarkr\checkmarkc\checkmark$\checkmark \checkmark$\checkmark^\checkmark\\checkmarkc\checkmarki\checkmarkr\checkmarkc\checkmark$\checkmarkP\checkmark$\checkmark^\checkmark\\checkmarkc\checkmarki\checkmarkr\checkmarkc\checkmark$\checkmarkr\checkmark$\checkmark^\checkmark\\checkmarkc\checkmarki\checkmarkr\checkmarkc\checkmark$\checkmarko\checkmark$\checkmark^\checkmark\\checkmarkc\checkmarki\checkmarkr\checkmarkc\checkmark$\checkmarkt\checkmark$\checkmark^\checkmark\\checkmarkc\checkmarki\checkmarkr\checkmarkc\checkmark$\checkmarko\checkmark$\checkmark^\checkmark\\checkmarkc\checkmarki\checkmarkr\checkmarkc\checkmark$\checkmarkc\checkmark$\checkmark^\checkmark\\checkmarkc\checkmarki\checkmarkr\checkmarkc\checkmark$\checkmarko\checkmark$\checkmark^\checkmark\\checkmarkc\checkmarki\checkmarkr\checkmarkc\checkmark$\checkmarkl\checkmark$\checkmark^\checkmark\\checkmarkc\checkmarki\checkmarkr\checkmarkc\checkmark$\checkmarke\checkmark$\checkmark^\checkmark\\checkmarkc\checkmarki\checkmarkr\checkmarkc\checkmark$\checkmark \checkmark$\checkmark^\checkmark\\checkmarkc\checkmarki\checkmarkr\checkmarkc\checkmark$\checkmark&\checkmark$\checkmark^\checkmark\\checkmarkc\checkmarki\checkmarkr\checkmarkc\checkmark$\checkmark \checkmark$\checkmark^\checkmark\\checkmarkc\checkmarki\checkmarkr\checkmarkc\checkmark$\checkmarkW\checkmark$\checkmark^\checkmark\\checkmarkc\checkmarki\checkmarkr\checkmarkc\checkmark$\checkmarki\checkmark$\checkmark^\checkmark\\checkmarkc\checkmarki\checkmarkr\checkmarkc\checkmark$\checkmark-\checkmark$\checkmark^\checkmark\\checkmarkc\checkmarki\checkmarkr\checkmarkc\checkmark$\checkmarkF\checkmark$\checkmark^\checkmark\\checkmarkc\checkmarki\checkmarkr\checkmarkc\checkmark$\checkmarki\checkmark$\checkmark^\checkmark\\checkmarkc\checkmarki\checkmarkr\checkmarkc\checkmark$\checkmark \checkmark$\checkmark^\checkmark\\checkmarkc\checkmarki\checkmarkr\checkmarkc\checkmark$\checkmarkU\checkmark$\checkmark^\checkmark\\checkmarkc\checkmarki\checkmarkr\checkmarkc\checkmark$\checkmarkD\checkmark$\checkmark^\checkmark\\checkmarkc\checkmarki\checkmarkr\checkmarkc\checkmark$\checkmarkP\checkmark$\checkmark^\checkmark\\checkmarkc\checkmarki\checkmarkr\checkmarkc\checkmark$\checkmark \checkmark$\checkmark^\checkmark\\checkmarkc\checkmarki\checkmarkr\checkmarkc\checkmark$\checkmark/\checkmark$\checkmark^\checkmark\\checkmarkc\checkmarki\checkmarkr\checkmarkc\checkmark$\checkmark \checkmark$\checkmark^\checkmark\\checkmarkc\checkmarki\checkmarkr\checkmarkc\checkmark$\checkmarkL\checkmark$\checkmark^\checkmark\\checkmarkc\checkmarki\checkmarkr\checkmarkc\checkmark$\checkmarka\checkmark$\checkmark^\checkmark\\checkmarkc\checkmarki\checkmarkr\checkmarkc\checkmark$\checkmarkt\checkmark$\checkmark^\checkmark\\checkmarkc\checkmarki\checkmarkr\checkmarkc\checkmark$\checkmarke\checkmark$\checkmark^\checkmark\\checkmarkc\checkmarki\checkmarkr\checkmarkc\checkmark$\checkmarkn\checkmark$\checkmark^\checkmark\\checkmarkc\checkmarki\checkmarkr\checkmarkc\checkmark$\checkmarkc\checkmark$\checkmark^\checkmark\\checkmarkc\checkmarki\checkmarkr\checkmarkc\checkmark$\checkmarke\checkmark$\checkmark^\checkmark\\checkmarkc\checkmarki\checkmarkr\checkmarkc\checkmark$\checkmark \checkmark$\checkmark^\checkmark\\checkmarkc\checkmarki\checkmarkr\checkmarkc\checkmark$\checkmark$\checkmark$\checkmark^\checkmark\\checkmarkc\checkmarki\checkmarkr\checkmarkc\checkmark$\checkmark<\checkmark$\checkmark^\checkmark\\checkmarkc\checkmarki\checkmarkr\checkmarkc\checkmark$\checkmark \checkmark$\checkmark^\checkmark\\checkmarkc\checkmarki\checkmarkr\checkmarkc\checkmark$\checkmark5\checkmark$\checkmark^\checkmark\\checkmarkc\checkmarki\checkmarkr\checkmarkc\checkmark$\checkmark$\checkmark$\checkmark^\checkmark\\checkmarkc\checkmarki\checkmarkr\checkmarkc\checkmark$\checkmark \checkmark$\checkmark^\checkmark\\checkmarkc\checkmarki\checkmarkr\checkmarkc\checkmark$\checkmarkm\checkmark$\checkmark^\checkmark\\checkmarkc\checkmarki\checkmarkr\checkmarkc\checkmark$\checkmarks\checkmark$\checkmark^\checkmark\\checkmarkc\checkmarki\checkmarkr\checkmarkc\checkmark$\checkmark \checkmark$\checkmark^\checkmark\\checkmarkc\checkmarki\checkmarkr\checkmarkc\checkmark$\checkmark&\checkmark$\checkmark^\checkmark\\checkmarkc\checkmarki\checkmarkr\checkmarkc\checkmark$\checkmark \checkmark$\checkmark^\checkmark\\checkmarkc\checkmarki\checkmarkr\checkmarkc\checkmark$\checkmark\\checkmark$\checkmark^\checkmark\\checkmarkc\checkmarki\checkmarkr\checkmarkc\checkmark$\checkmark\\checkmark$\checkmark^\checkmark\\checkmarkc\checkmarki\checkmarkr\checkmarkc\checkmark$\checkmark \checkmark$\checkmark^\checkmark\\checkmarkc\checkmarki\checkmarkr\checkmarkc\checkmark$\checkmark\\checkmark$\checkmark^\checkmark\\checkmarkc\checkmarki\checkmarkr\checkmarkc\checkmark$\checkmarkh\checkmark$\checkmark^\checkmark\\checkmarkc\checkmarki\checkmarkr\checkmarkc\checkmark$\checkmarkl\checkmark$\checkmark^\checkmark\\checkmarkc\checkmarki\checkmarkr\checkmarkc\checkmark$\checkmarki\checkmark$\checkmark^\checkmark\\checkmarkc\checkmarki\checkmarkr\checkmarkc\checkmark$\checkmarkn\checkmark$\checkmark^\checkmark\\checkmarkc\checkmarki\checkmarkr\checkmarkc\checkmark$\checkmarke\checkmark$\checkmark^\checkmark\\checkmarkc\checkmarki\checkmarkr\checkmarkc\checkmark$\checkmark
\checkmark$\checkmark^\checkmark\\checkmarkc\checkmarki\checkmarkr\checkmarkc\checkmark$\checkmark\\checkmark$\checkmark^\checkmark\\checkmarkc\checkmarki\checkmarkr\checkmarkc\checkmark$\checkmarkt\checkmark$\checkmark^\checkmark\\checkmarkc\checkmarki\checkmarkr\checkmarkc\checkmark$\checkmarke\checkmark$\checkmark^\checkmark\\checkmarkc\checkmarki\checkmarkr\checkmarkc\checkmark$\checkmarkx\checkmark$\checkmark^\checkmark\\checkmarkc\checkmarki\checkmarkr\checkmarkc\checkmark$\checkmarkt\checkmark$\checkmark^\checkmark\\checkmarkc\checkmarki\checkmarkr\checkmarkc\checkmark$\checkmarkb\checkmark$\checkmark^\checkmark\\checkmarkc\checkmarki\checkmarkr\checkmarkc\checkmark$\checkmarkf\checkmark$\checkmark^\checkmark\\checkmarkc\checkmarki\checkmarkr\checkmarkc\checkmark$\checkmark{\checkmark$\checkmark^\checkmark\\checkmarkc\checkmarki\checkmarkr\checkmarkc\checkmark$\checkmarkF\checkmark$\checkmark^\checkmark\\checkmarkc\checkmarki\checkmarkr\checkmarkc\checkmark$\checkmarkC\checkmark$\checkmark^\checkmark\\checkmarkc\checkmarki\checkmarkr\checkmarkc\checkmark$\checkmark6\checkmark$\checkmark^\checkmark\\checkmarkc\checkmarki\checkmarkr\checkmarkc\checkmark$\checkmark \checkmark$\checkmark^\checkmark\\checkmarkc\checkmarki\checkmarkr\checkmarkc\checkmark$\checkmark(\checkmark$\checkmark^\checkmark\\checkmarkc\checkmarki\checkmarkr\checkmarkc\checkmark$\checkmarkÉ\checkmark$\checkmark^\checkmark\\checkmarkc\checkmarki\checkmarkr\checkmarkc\checkmark$\checkmarkn\checkmark$\checkmark^\checkmark\\checkmarkc\checkmarki\checkmarkr\checkmarkc\checkmark$\checkmarke\checkmark$\checkmark^\checkmark\\checkmarkc\checkmarki\checkmarkr\checkmarkc\checkmark$\checkmarkr\checkmark$\checkmark^\checkmark\\checkmarkc\checkmarki\checkmarkr\checkmarkc\checkmark$\checkmarkg\checkmark$\checkmark^\checkmark\\checkmarkc\checkmarki\checkmarkr\checkmarkc\checkmark$\checkmarki\checkmark$\checkmark^\checkmark\\checkmarkc\checkmarki\checkmarkr\checkmarkc\checkmark$\checkmarke\checkmark$\checkmark^\checkmark\\checkmarkc\checkmarki\checkmarkr\checkmarkc\checkmark$\checkmark)\checkmark$\checkmark^\checkmark\\checkmarkc\checkmarki\checkmarkr\checkmarkc\checkmark$\checkmark}\checkmark$\checkmark^\checkmark\\checkmarkc\checkmarki\checkmarkr\checkmarkc\checkmark$\checkmark \checkmark$\checkmark^\checkmark\\checkmarkc\checkmarki\checkmarkr\checkmarkc\checkmark$\checkmark&\checkmark$\checkmark^\checkmark\\checkmarkc\checkmarki\checkmarkr\checkmarkc\checkmark$\checkmark \checkmark$\checkmark^\checkmark\\checkmarkc\checkmarki\checkmarkr\checkmarkc\checkmark$\checkmarkA\checkmark$\checkmark^\checkmark\\checkmarkc\checkmarki\checkmarkr\checkmarkc\checkmark$\checkmarkl\checkmark$\checkmark^\checkmark\\checkmarkc\checkmarki\checkmarkr\checkmarkc\checkmark$\checkmarki\checkmark$\checkmark^\checkmark\\checkmarkc\checkmarki\checkmarkr\checkmarkc\checkmark$\checkmarkm\checkmark$\checkmark^\checkmark\\checkmarkc\checkmarki\checkmarkr\checkmarkc\checkmark$\checkmarke\checkmark$\checkmark^\checkmark\\checkmarkc\checkmarki\checkmarkr\checkmarkc\checkmark$\checkmarkn\checkmark$\checkmark^\checkmark\\checkmarkc\checkmarki\checkmarkr\checkmarkc\checkmark$\checkmarkt\checkmark$\checkmark^\checkmark\\checkmarkc\checkmarki\checkmarkr\checkmarkc\checkmark$\checkmarka\checkmark$\checkmark^\checkmark\\checkmarkc\checkmarki\checkmarkr\checkmarkc\checkmark$\checkmarkt\checkmark$\checkmark^\checkmark\\checkmarkc\checkmarki\checkmarkr\checkmarkc\checkmark$\checkmarki\checkmark$\checkmark^\checkmark\\checkmarkc\checkmarki\checkmarkr\checkmarkc\checkmark$\checkmarko\checkmark$\checkmark^\checkmark\\checkmarkc\checkmarki\checkmarkr\checkmarkc\checkmark$\checkmarkn\checkmark$\checkmark^\checkmark\\checkmarkc\checkmarki\checkmarkr\checkmarkc\checkmark$\checkmark \checkmark$\checkmark^\checkmark\\checkmarkc\checkmarki\checkmarkr\checkmarkc\checkmark$\checkmark&\checkmark$\checkmark^\checkmark\\checkmarkc\checkmarki\checkmarkr\checkmarkc\checkmark$\checkmark \checkmark$\checkmark^\checkmark\\checkmarkc\checkmarki\checkmarkr\checkmarkc\checkmark$\checkmark3\checkmark$\checkmark^\checkmark\\checkmarkc\checkmarki\checkmarkr\checkmarkc\checkmark$\checkmark6\checkmark$\checkmark^\checkmark\\checkmarkc\checkmarki\checkmarkr\checkmarkc\checkmark$\checkmark \checkmark$\checkmark^\checkmark\\checkmarkc\checkmarki\checkmarkr\checkmarkc\checkmark$\checkmarkV\checkmark$\checkmark^\checkmark\\checkmarkc\checkmarki\checkmarkr\checkmarkc\checkmark$\checkmark \checkmark$\checkmark^\checkmark\\checkmarkc\checkmarki\checkmarkr\checkmarkc\checkmark$\checkmarkD\checkmark$\checkmark^\checkmark\\checkmarkc\checkmarki\checkmarkr\checkmarkc\checkmark$\checkmarkC\checkmark$\checkmark^\checkmark\\checkmarkc\checkmarki\checkmarkr\checkmarkc\checkmark$\checkmark \checkmark$\checkmark^\checkmark\\checkmarkc\checkmarki\checkmarkr\checkmarkc\checkmark$\checkmark/\checkmark$\checkmark^\checkmark\\checkmarkc\checkmarki\checkmarkr\checkmarkc\checkmark$\checkmark \checkmark$\checkmark^\checkmark\\checkmarkc\checkmarki\checkmarkr\checkmarkc\checkmark$\checkmark3\checkmark$\checkmark^\checkmark\\checkmarkc\checkmarki\checkmarkr\checkmarkc\checkmark$\checkmark5\checkmark$\checkmark^\checkmark\\checkmarkc\checkmarki\checkmarkr\checkmarkc\checkmark$\checkmark0\checkmark$\checkmark^\checkmark\\checkmarkc\checkmarki\checkmarkr\checkmarkc\checkmark$\checkmark \checkmark$\checkmark^\checkmark\\checkmarkc\checkmarki\checkmarkr\checkmarkc\checkmark$\checkmarkW\checkmark$\checkmark^\checkmark\\checkmarkc\checkmarki\checkmarkr\checkmarkc\checkmark$\checkmark \checkmark$\checkmark^\checkmark\\checkmarkc\checkmarki\checkmarkr\checkmarkc\checkmark$\checkmark&\checkmark$\checkmark^\checkmark\\checkmarkc\checkmarki\checkmarkr\checkmarkc\checkmark$\checkmark \checkmark$\checkmark^\checkmark\\checkmarkc\checkmarki\checkmarkr\checkmarkc\checkmark$\checkmarkC\checkmark$\checkmark^\checkmark\\checkmarkc\checkmarki\checkmarkr\checkmarkc\checkmark$\checkmarko\checkmark$\checkmark^\checkmark\\checkmarkc\checkmarki\checkmarkr\checkmarkc\checkmark$\checkmarke\checkmark$\checkmark^\checkmark\\checkmarkc\checkmarki\checkmarkr\checkmarkc\checkmark$\checkmarkf\checkmark$\checkmark^\checkmark\\checkmarkc\checkmarki\checkmarkr\checkmarkc\checkmark$\checkmarkf\checkmark$\checkmark^\checkmark\\checkmarkc\checkmarki\checkmarkr\checkmarkc\checkmark$\checkmarki\checkmark$\checkmark^\checkmark\\checkmarkc\checkmarki\checkmarkr\checkmarkc\checkmark$\checkmarkc\checkmark$\checkmark^\checkmark\\checkmarkc\checkmarki\checkmarkr\checkmarkc\checkmark$\checkmarki\checkmark$\checkmark^\checkmark\\checkmarkc\checkmarki\checkmarkr\checkmarkc\checkmark$\checkmarke\checkmark$\checkmark^\checkmark\\checkmarkc\checkmarki\checkmarkr\checkmarkc\checkmark$\checkmarkn\checkmark$\checkmark^\checkmark\\checkmarkc\checkmarki\checkmarkr\checkmarkc\checkmark$\checkmarkt\checkmark$\checkmark^\checkmark\\checkmarkc\checkmarki\checkmarkr\checkmarkc\checkmark$\checkmark \checkmark$\checkmark^\checkmark\\checkmarkc\checkmarki\checkmarkr\checkmarkc\checkmark$\checkmark$\checkmark$\checkmark^\checkmark\\checkmarkc\checkmarki\checkmarkr\checkmarkc\checkmark$\checkmark\\checkmark$\checkmark^\checkmark\\checkmarkc\checkmarki\checkmarkr\checkmarkc\checkmark$\checkmarkt\checkmark$\checkmark^\checkmark\\checkmarkc\checkmarki\checkmarkr\checkmarkc\checkmark$\checkmarki\checkmark$\checkmark^\checkmark\\checkmarkc\checkmarki\checkmarkr\checkmarkc\checkmark$\checkmarkm\checkmark$\checkmark^\checkmark\\checkmarkc\checkmarki\checkmarkr\checkmarkc\checkmark$\checkmarke\checkmark$\checkmark^\checkmark\\checkmarkc\checkmarki\checkmarkr\checkmarkc\checkmark$\checkmarks\checkmark$\checkmark^\checkmark\\checkmarkc\checkmarki\checkmarkr\checkmarkc\checkmark$\checkmark \checkmark$\checkmark^\checkmark\\checkmarkc\checkmarki\checkmarkr\checkmarkc\checkmark$\checkmark1\checkmark$\checkmark^\checkmark\\checkmarkc\checkmarki\checkmarkr\checkmarkc\checkmark$\checkmark{\checkmark$\checkmark^\checkmark\\checkmarkc\checkmarki\checkmarkr\checkmarkc\checkmark$\checkmark,\checkmark$\checkmark^\checkmark\\checkmarkc\checkmarki\checkmarkr\checkmarkc\checkmark$\checkmark}\checkmark$\checkmark^\checkmark\\checkmarkc\checkmarki\checkmarkr\checkmarkc\checkmark$\checkmark2\checkmark$\checkmark^\checkmark\\checkmarkc\checkmarki\checkmarkr\checkmarkc\checkmark$\checkmark5\checkmark$\checkmark^\checkmark\\checkmarkc\checkmarki\checkmarkr\checkmarkc\checkmark$\checkmark$\checkmark$\checkmark^\checkmark\\checkmarkc\checkmarki\checkmarkr\checkmarkc\checkmark$\checkmark \checkmark$\checkmark^\checkmark\\checkmarkc\checkmarki\checkmarkr\checkmarkc\checkmark$\checkmark\\checkmark$\checkmark^\checkmark\\checkmarkc\checkmarki\checkmarkr\checkmarkc\checkmark$\checkmark\\checkmark$\checkmark^\checkmark\\checkmarkc\checkmarki\checkmarkr\checkmarkc\checkmark$\checkmark \checkmark$\checkmark^\checkmark\\checkmarkc\checkmarki\checkmarkr\checkmarkc\checkmark$\checkmark\\checkmark$\checkmark^\checkmark\\checkmarkc\checkmarki\checkmarkr\checkmarkc\checkmark$\checkmarkh\checkmark$\checkmark^\checkmark\\checkmarkc\checkmarki\checkmarkr\checkmarkc\checkmark$\checkmarkl\checkmark$\checkmark^\checkmark\\checkmarkc\checkmarki\checkmarkr\checkmarkc\checkmark$\checkmarki\checkmark$\checkmark^\checkmark\\checkmarkc\checkmarki\checkmarkr\checkmarkc\checkmark$\checkmarkn\checkmark$\checkmark^\checkmark\\checkmarkc\checkmarki\checkmarkr\checkmarkc\checkmark$\checkmarke\checkmark$\checkmark^\checkmark\\checkmarkc\checkmarki\checkmarkr\checkmarkc\checkmark$\checkmark
\checkmark$\checkmark^\checkmark\\checkmarkc\checkmarki\checkmarkr\checkmarkc\checkmark$\checkmark\\checkmark$\checkmark^\checkmark\\checkmarkc\checkmarki\checkmarkr\checkmarkc\checkmark$\checkmarke\checkmark$\checkmark^\checkmark\\checkmarkc\checkmarki\checkmarkr\checkmarkc\checkmark$\checkmarkn\checkmark$\checkmark^\checkmark\\checkmarkc\checkmarki\checkmarkr\checkmarkc\checkmark$\checkmarkd\checkmark$\checkmark^\checkmark\\checkmarkc\checkmarki\checkmarkr\checkmarkc\checkmark$\checkmark{\checkmark$\checkmark^\checkmark\\checkmarkc\checkmarki\checkmarkr\checkmarkc\checkmark$\checkmarkt\checkmark$\checkmark^\checkmark\\checkmarkc\checkmarki\checkmarkr\checkmarkc\checkmark$\checkmarka\checkmark$\checkmark^\checkmark\\checkmarkc\checkmarki\checkmarkr\checkmarkc\checkmark$\checkmarkb\checkmark$\checkmark^\checkmark\\checkmarkc\checkmarki\checkmarkr\checkmarkc\checkmark$\checkmarku\checkmark$\checkmark^\checkmark\\checkmarkc\checkmarki\checkmarkr\checkmarkc\checkmark$\checkmarkl\checkmark$\checkmark^\checkmark\\checkmarkc\checkmarki\checkmarkr\checkmarkc\checkmark$\checkmarka\checkmark$\checkmark^\checkmark\\checkmarkc\checkmarki\checkmarkr\checkmarkc\checkmark$\checkmarkr\checkmark$\checkmark^\checkmark\\checkmarkc\checkmarki\checkmarkr\checkmarkc\checkmark$\checkmark}\checkmark$\checkmark^\checkmark\\checkmarkc\checkmarki\checkmarkr\checkmarkc\checkmark$\checkmark
\checkmark$\checkmark^\checkmark\\checkmarkc\checkmarki\checkmarkr\checkmarkc\checkmark$\checkmark}\checkmark$\checkmark^\checkmark\\checkmarkc\checkmarki\checkmarkr\checkmarkc\checkmark$\checkmark
\checkmark$\checkmark^\checkmark\\checkmarkc\checkmarki\checkmarkr\checkmarkc\checkmark$\checkmark\\checkmark$\checkmark^\checkmark\\checkmarkc\checkmarki\checkmarkr\checkmarkc\checkmark$\checkmarkc\checkmark$\checkmark^\checkmark\\checkmarkc\checkmarki\checkmarkr\checkmarkc\checkmark$\checkmarka\checkmark$\checkmark^\checkmark\\checkmarkc\checkmarki\checkmarkr\checkmarkc\checkmark$\checkmarkp\checkmark$\checkmark^\checkmark\\checkmarkc\checkmarki\checkmarkr\checkmarkc\checkmark$\checkmarkt\checkmark$\checkmark^\checkmark\\checkmarkc\checkmarki\checkmarkr\checkmarkc\checkmark$\checkmarki\checkmark$\checkmark^\checkmark\\checkmarkc\checkmarki\checkmarkr\checkmarkc\checkmark$\checkmarko\checkmark$\checkmark^\checkmark\\checkmarkc\checkmarki\checkmarkr\checkmarkc\checkmark$\checkmarkn\checkmark$\checkmark^\checkmark\\checkmarkc\checkmarki\checkmarkr\checkmarkc\checkmark$\checkmark{\checkmark$\checkmark^\checkmark\\checkmarkc\checkmarki\checkmarkr\checkmarkc\checkmark$\checkmarkC\checkmark$\checkmark^\checkmark\\checkmarkc\checkmarki\checkmarkr\checkmarkc\checkmark$\checkmarka\checkmark$\checkmark^\checkmark\\checkmarkc\checkmarki\checkmarkr\checkmarkc\checkmark$\checkmarkh\checkmark$\checkmark^\checkmark\\checkmarkc\checkmarki\checkmarkr\checkmarkc\checkmark$\checkmarki\checkmark$\checkmark^\checkmark\\checkmarkc\checkmarki\checkmarkr\checkmarkc\checkmark$\checkmarke\checkmark$\checkmark^\checkmark\\checkmarkc\checkmarki\checkmarkr\checkmarkc\checkmark$\checkmarkr\checkmark$\checkmark^\checkmark\\checkmarkc\checkmarki\checkmarkr\checkmarkc\checkmark$\checkmark \checkmark$\checkmark^\checkmark\\checkmarkc\checkmarki\checkmarkr\checkmarkc\checkmark$\checkmarkd\checkmark$\checkmark^\checkmark\\checkmarkc\checkmarki\checkmarkr\checkmarkc\checkmark$\checkmarke\checkmark$\checkmark^\checkmark\\checkmarkc\checkmarki\checkmarkr\checkmarkc\checkmark$\checkmarks\checkmark$\checkmark^\checkmark\\checkmarkc\checkmarki\checkmarkr\checkmarkc\checkmark$\checkmark \checkmark$\checkmark^\checkmark\\checkmarkc\checkmarki\checkmarkr\checkmarkc\checkmark$\checkmarkc\checkmark$\checkmark^\checkmark\\checkmarkc\checkmarki\checkmarkr\checkmarkc\checkmark$\checkmarkh\checkmark$\checkmark^\checkmark\\checkmarkc\checkmarki\checkmarkr\checkmarkc\checkmark$\checkmarka\checkmark$\checkmark^\checkmark\\checkmarkc\checkmarki\checkmarkr\checkmarkc\checkmark$\checkmarkr\checkmark$\checkmark^\checkmark\\checkmarkc\checkmarki\checkmarkr\checkmarkc\checkmark$\checkmarkg\checkmark$\checkmark^\checkmark\\checkmarkc\checkmarki\checkmarkr\checkmarkc\checkmark$\checkmarke\checkmark$\checkmark^\checkmark\\checkmarkc\checkmarki\checkmarkr\checkmarkc\checkmark$\checkmarks\checkmark$\checkmark^\checkmark\\checkmarkc\checkmarki\checkmarkr\checkmarkc\checkmark$\checkmark \checkmark$\checkmark^\checkmark\\checkmarkc\checkmarki\checkmarkr\checkmarkc\checkmark$\checkmarkf\checkmark$\checkmark^\checkmark\\checkmarkc\checkmarki\checkmarkr\checkmarkc\checkmark$\checkmarko\checkmark$\checkmark^\checkmark\\checkmarkc\checkmarki\checkmarkr\checkmarkc\checkmark$\checkmarkn\checkmark$\checkmark^\checkmark\\checkmarkc\checkmarki\checkmarkr\checkmarkc\checkmark$\checkmarkc\checkmark$\checkmark^\checkmark\\checkmarkc\checkmarki\checkmarkr\checkmarkc\checkmark$\checkmarkt\checkmark$\checkmark^\checkmark\\checkmarkc\checkmarki\checkmarkr\checkmarkc\checkmark$\checkmarki\checkmark$\checkmark^\checkmark\\checkmarkc\checkmarki\checkmarkr\checkmarkc\checkmark$\checkmarko\checkmark$\checkmark^\checkmark\\checkmarkc\checkmarki\checkmarkr\checkmarkc\checkmark$\checkmarkn\checkmark$\checkmark^\checkmark\\checkmarkc\checkmarki\checkmarkr\checkmarkc\checkmark$\checkmarkn\checkmark$\checkmark^\checkmark\\checkmarkc\checkmarki\checkmarkr\checkmarkc\checkmark$\checkmarke\checkmark$\checkmark^\checkmark\\checkmarkc\checkmarki\checkmarkr\checkmarkc\checkmark$\checkmarkl\checkmark$\checkmark^\checkmark\\checkmarkc\checkmarki\checkmarkr\checkmarkc\checkmark$\checkmark \checkmark$\checkmark^\checkmark\\checkmarkc\checkmarki\checkmarkr\checkmarkc\checkmark$\checkmarkq\checkmark$\checkmark^\checkmark\\checkmarkc\checkmarki\checkmarkr\checkmarkc\checkmark$\checkmarku\checkmark$\checkmark^\checkmark\\checkmarkc\checkmarki\checkmarkr\checkmarkc\checkmark$\checkmarka\checkmark$\checkmark^\checkmark\\checkmarkc\checkmarki\checkmarkr\checkmarkc\checkmark$\checkmarkn\checkmark$\checkmark^\checkmark\\checkmarkc\checkmarki\checkmarkr\checkmarkc\checkmark$\checkmarkt\checkmark$\checkmark^\checkmark\\checkmarkc\checkmarki\checkmarkr\checkmarkc\checkmark$\checkmarki\checkmark$\checkmark^\checkmark\\checkmarkc\checkmarki\checkmarkr\checkmarkc\checkmark$\checkmarkt\checkmark$\checkmark^\checkmark\\checkmarkc\checkmarki\checkmarkr\checkmarkc\checkmark$\checkmarka\checkmark$\checkmark^\checkmark\\checkmarkc\checkmarki\checkmarkr\checkmarkc\checkmark$\checkmarkt\checkmark$\checkmark^\checkmark\\checkmarkc\checkmarki\checkmarkr\checkmarkc\checkmark$\checkmarki\checkmark$\checkmark^\checkmark\\checkmarkc\checkmarki\checkmarkr\checkmarkc\checkmark$\checkmarkf\checkmark$\checkmark^\checkmark\\checkmarkc\checkmarki\checkmarkr\checkmarkc\checkmark$\checkmark \checkmark$\checkmark^\checkmark\\checkmarkc\checkmarki\checkmarkr\checkmarkc\checkmark$\checkmarkd\checkmark$\checkmark^\checkmark\\checkmarkc\checkmarki\checkmarkr\checkmarkc\checkmark$\checkmarke\checkmark$\checkmark^\checkmark\\checkmarkc\checkmarki\checkmarkr\checkmarkc\checkmark$\checkmark \checkmark$\checkmark^\checkmark\\checkmarkc\checkmarki\checkmarkr\checkmarkc\checkmark$\checkmarkl\checkmark$\checkmark^\checkmark\\checkmarkc\checkmarki\checkmarkr\checkmarkc\checkmark$\checkmarka\checkmark$\checkmark^\checkmark\\checkmarkc\checkmarki\checkmarkr\checkmarkc\checkmark$\checkmark \checkmark$\checkmark^\checkmark\\checkmarkc\checkmarki\checkmarkr\checkmarkc\checkmark$\checkmarkm\checkmark$\checkmark^\checkmark\\checkmarkc\checkmarki\checkmarkr\checkmarkc\checkmark$\checkmarki\checkmark$\checkmark^\checkmark\\checkmarkc\checkmarki\checkmarkr\checkmarkc\checkmark$\checkmarkn\checkmark$\checkmark^\checkmark\\checkmarkc\checkmarki\checkmarkr\checkmarkc\checkmark$\checkmarki\checkmark$\checkmark^\checkmark\\checkmarkc\checkmarki\checkmarkr\checkmarkc\checkmark$\checkmark-\checkmark$\checkmark^\checkmark\\checkmarkc\checkmarki\checkmarkr\checkmarkc\checkmark$\checkmarkf\checkmark$\checkmark^\checkmark\\checkmarkc\checkmarki\checkmarkr\checkmarkc\checkmark$\checkmarkr\checkmark$\checkmark^\checkmark\\checkmarkc\checkmarki\checkmarkr\checkmarkc\checkmark$\checkmarka\checkmark$\checkmark^\checkmark\\checkmarkc\checkmarki\checkmarkr\checkmarkc\checkmark$\checkmarki\checkmark$\checkmark^\checkmark\\checkmarkc\checkmarki\checkmarkr\checkmarkc\checkmark$\checkmarks\checkmark$\checkmark^\checkmark\\checkmarkc\checkmarki\checkmarkr\checkmarkc\checkmark$\checkmarke\checkmark$\checkmark^\checkmark\\checkmarkc\checkmarki\checkmarkr\checkmarkc\checkmark$\checkmarku\checkmark$\checkmark^\checkmark\\checkmarkc\checkmarki\checkmarkr\checkmarkc\checkmark$\checkmarks\checkmark$\checkmark^\checkmark\\checkmarkc\checkmarki\checkmarkr\checkmarkc\checkmark$\checkmarke\checkmark$\checkmark^\checkmark\\checkmarkc\checkmarki\checkmarkr\checkmarkc\checkmark$\checkmark \checkmark$\checkmark^\checkmark\\checkmarkc\checkmarki\checkmarkr\checkmarkc\checkmark$\checkmarkC\checkmark$\checkmark^\checkmark\\checkmarkc\checkmarki\checkmarkr\checkmarkc\checkmark$\checkmarkN\checkmark$\checkmark^\checkmark\\checkmarkc\checkmarki\checkmarkr\checkmarkc\checkmark$\checkmarkC\checkmark$\checkmark^\checkmark\\checkmarkc\checkmarki\checkmarkr\checkmarkc\checkmark$\checkmark \checkmark$\checkmark^\checkmark\\checkmarkc\checkmarki\checkmarkr\checkmarkc\checkmark$\checkmark5\checkmark$\checkmark^\checkmark\\checkmarkc\checkmarki\checkmarkr\checkmarkc\checkmark$\checkmark \checkmark$\checkmark^\checkmark\\checkmarkc\checkmarki\checkmarkr\checkmarkc\checkmark$\checkmarka\checkmark$\checkmark^\checkmark\\checkmarkc\checkmarki\checkmarkr\checkmarkc\checkmark$\checkmarkx\checkmark$\checkmark^\checkmark\\checkmarkc\checkmarki\checkmarkr\checkmarkc\checkmark$\checkmarke\checkmark$\checkmark^\checkmark\\checkmarkc\checkmarki\checkmarkr\checkmarkc\checkmark$\checkmarks\checkmark$\checkmark^\checkmark\\checkmarkc\checkmarki\checkmarkr\checkmarkc\checkmark$\checkmark}\checkmark$\checkmark^\checkmark\\checkmarkc\checkmarki\checkmarkr\checkmarkc\checkmark$\checkmark
\checkmark$\checkmark^\checkmark\\checkmarkc\checkmarki\checkmarkr\checkmarkc\checkmark$\checkmark\\checkmark$\checkmark^\checkmark\\checkmarkc\checkmarki\checkmarkr\checkmarkc\checkmark$\checkmarkl\checkmark$\checkmark^\checkmark\\checkmarkc\checkmarki\checkmarkr\checkmarkc\checkmark$\checkmarka\checkmark$\checkmark^\checkmark\\checkmarkc\checkmarki\checkmarkr\checkmarkc\checkmark$\checkmarkb\checkmark$\checkmark^\checkmark\\checkmarkc\checkmarki\checkmarkr\checkmarkc\checkmark$\checkmarke\checkmark$\checkmark^\checkmark\\checkmarkc\checkmarki\checkmarkr\checkmarkc\checkmark$\checkmarkl\checkmark$\checkmark^\checkmark\\checkmarkc\checkmarki\checkmarkr\checkmarkc\checkmark$\checkmark{\checkmark$\checkmark^\checkmark\\checkmarkc\checkmarki\checkmarkr\checkmarkc\checkmark$\checkmarkt\checkmark$\checkmark^\checkmark\\checkmarkc\checkmarki\checkmarkr\checkmarkc\checkmark$\checkmarka\checkmark$\checkmark^\checkmark\\checkmarkc\checkmarki\checkmarkr\checkmarkc\checkmark$\checkmarkb\checkmark$\checkmark^\checkmark\\checkmarkc\checkmarki\checkmarkr\checkmarkc\checkmark$\checkmark:\checkmark$\checkmark^\checkmark\\checkmarkc\checkmarki\checkmarkr\checkmarkc\checkmark$\checkmarkc\checkmark$\checkmark^\checkmark\\checkmarkc\checkmarki\checkmarkr\checkmarkc\checkmark$\checkmarkd\checkmark$\checkmark^\checkmark\\checkmarkc\checkmarki\checkmarkr\checkmarkc\checkmark$\checkmarkc\checkmark$\checkmark^\checkmark\\checkmarkc\checkmarki\checkmarkr\checkmarkc\checkmark$\checkmarkf\checkmark$\checkmark^\checkmark\\checkmarkc\checkmarki\checkmarkr\checkmarkc\checkmark$\checkmark}\checkmark$\checkmark^\checkmark\\checkmarkc\checkmarki\checkmarkr\checkmarkc\checkmark$\checkmark
\checkmark$\checkmark^\checkmark\\checkmarkc\checkmarki\checkmarkr\checkmarkc\checkmark$\checkmark\\checkmark$\checkmark^\checkmark\\checkmarkc\checkmarki\checkmarkr\checkmarkc\checkmark$\checkmarke\checkmark$\checkmark^\checkmark\\checkmarkc\checkmarki\checkmarkr\checkmarkc\checkmark$\checkmarkn\checkmark$\checkmark^\checkmark\\checkmarkc\checkmarki\checkmarkr\checkmarkc\checkmark$\checkmarkd\checkmark$\checkmark^\checkmark\\checkmarkc\checkmarki\checkmarkr\checkmarkc\checkmark$\checkmark{\checkmark$\checkmark^\checkmark\\checkmarkc\checkmarki\checkmarkr\checkmarkc\checkmark$\checkmarkt\checkmark$\checkmark^\checkmark\\checkmarkc\checkmarki\checkmarkr\checkmarkc\checkmark$\checkmarka\checkmark$\checkmark^\checkmark\\checkmarkc\checkmarki\checkmarkr\checkmarkc\checkmark$\checkmarkb\checkmark$\checkmark^\checkmark\\checkmarkc\checkmarki\checkmarkr\checkmarkc\checkmark$\checkmarkl\checkmark$\checkmark^\checkmark\\checkmarkc\checkmarki\checkmarkr\checkmarkc\checkmark$\checkmarke\checkmark$\checkmark^\checkmark\\checkmarkc\checkmarki\checkmarkr\checkmarkc\checkmark$\checkmark}\checkmark$\checkmark^\checkmark\\checkmarkc\checkmarki\checkmarkr\checkmarkc\checkmark$\checkmark
\checkmark$\checkmark^\checkmark\\checkmarkc\checkmarki\checkmarkr\checkmarkc\checkmark$\checkmark
\checkmark$\checkmark^\checkmark\\checkmarkc\checkmarki\checkmarkr\checkmarkc\checkmark$\checkmark\\checkmark$\checkmark^\checkmark\\checkmarkc\checkmarki\checkmarkr\checkmarkc\checkmark$\checkmarks\checkmark$\checkmark^\checkmark\\checkmarkc\checkmarki\checkmarkr\checkmarkc\checkmark$\checkmarku\checkmark$\checkmark^\checkmark\\checkmarkc\checkmarki\checkmarkr\checkmarkc\checkmark$\checkmarkb\checkmark$\checkmark^\checkmark\\checkmarkc\checkmarki\checkmarkr\checkmarkc\checkmark$\checkmarks\checkmark$\checkmark^\checkmark\\checkmarkc\checkmarki\checkmarkr\checkmarkc\checkmark$\checkmarke\checkmark$\checkmark^\checkmark\\checkmarkc\checkmarki\checkmarkr\checkmarkc\checkmark$\checkmarkc\checkmark$\checkmark^\checkmark\\checkmarkc\checkmarki\checkmarkr\checkmarkc\checkmark$\checkmarkt\checkmark$\checkmark^\checkmark\\checkmarkc\checkmarki\checkmarkr\checkmarkc\checkmark$\checkmarki\checkmark$\checkmark^\checkmark\\checkmarkc\checkmarki\checkmarkr\checkmarkc\checkmark$\checkmarko\checkmark$\checkmark^\checkmark\\checkmarkc\checkmarki\checkmarkr\checkmarkc\checkmark$\checkmarkn\checkmark$\checkmark^\checkmark\\checkmarkc\checkmarki\checkmarkr\checkmarkc\checkmark$\checkmark{\checkmark$\checkmark^\checkmark\\checkmarkc\checkmarki\checkmarkr\checkmarkc\checkmark$\checkmarkD\checkmark$\checkmark^\checkmark\\checkmarkc\checkmarki\checkmarkr\checkmarkc\checkmark$\checkmarki\checkmark$\checkmark^\checkmark\\checkmarkc\checkmarki\checkmarkr\checkmarkc\checkmark$\checkmarka\checkmark$\checkmark^\checkmark\\checkmarkc\checkmarki\checkmarkr\checkmarkc\checkmark$\checkmarkg\checkmark$\checkmark^\checkmark\\checkmarkc\checkmarki\checkmarkr\checkmarkc\checkmark$\checkmarkr\checkmark$\checkmark^\checkmark\\checkmarkc\checkmarki\checkmarkr\checkmarkc\checkmark$\checkmarka\checkmark$\checkmark^\checkmark\\checkmarkc\checkmarki\checkmarkr\checkmarkc\checkmark$\checkmarkm\checkmark$\checkmark^\checkmark\\checkmarkc\checkmarki\checkmarkr\checkmarkc\checkmark$\checkmarkm\checkmark$\checkmark^\checkmark\\checkmarkc\checkmarki\checkmarkr\checkmarkc\checkmark$\checkmarke\checkmark$\checkmark^\checkmark\\checkmarkc\checkmarki\checkmarkr\checkmarkc\checkmark$\checkmark \checkmark$\checkmark^\checkmark\\checkmarkc\checkmarki\checkmarkr\checkmarkc\checkmark$\checkmarkF\checkmark$\checkmark^\checkmark\\checkmarkc\checkmarki\checkmarkr\checkmarkc\checkmark$\checkmarkA\checkmark$\checkmark^\checkmark\\checkmarkc\checkmarki\checkmarkr\checkmarkc\checkmark$\checkmarkS\checkmark$\checkmark^\checkmark\\checkmarkc\checkmarki\checkmarkr\checkmarkc\checkmark$\checkmarkT\checkmark$\checkmark^\checkmark\\checkmarkc\checkmarki\checkmarkr\checkmarkc\checkmark$\checkmark \checkmark$\checkmark^\checkmark\\checkmarkc\checkmarki\checkmarkr\checkmarkc\checkmark$\checkmark—\checkmark$\checkmark^\checkmark\\checkmarkc\checkmarki\checkmarkr\checkmarkc\checkmark$\checkmark \checkmark$\checkmark^\checkmark\\checkmarkc\checkmarki\checkmarkr\checkmarkc\checkmark$\checkmarkA\checkmark$\checkmark^\checkmark\\checkmarkc\checkmarki\checkmarkr\checkmarkc\checkmark$\checkmarkr\checkmark$\checkmark^\checkmark\\checkmarkc\checkmarki\checkmarkr\checkmarkc\checkmark$\checkmarkc\checkmark$\checkmark^\checkmark\\checkmarkc\checkmarki\checkmarkr\checkmarkc\checkmark$\checkmarkh\checkmark$\checkmark^\checkmark\\checkmarkc\checkmarki\checkmarkr\checkmarkc\checkmark$\checkmarki\checkmark$\checkmark^\checkmark\\checkmarkc\checkmarki\checkmarkr\checkmarkc\checkmark$\checkmarkt\checkmark$\checkmark^\checkmark\\checkmarkc\checkmarki\checkmarkr\checkmarkc\checkmark$\checkmarke\checkmark$\checkmark^\checkmark\\checkmarkc\checkmarki\checkmarkr\checkmarkc\checkmark$\checkmarkc\checkmark$\checkmark^\checkmark\\checkmarkc\checkmarki\checkmarkr\checkmarkc\checkmark$\checkmarkt\checkmark$\checkmark^\checkmark\\checkmarkc\checkmarki\checkmarkr\checkmarkc\checkmark$\checkmarku\checkmark$\checkmark^\checkmark\\checkmarkc\checkmarki\checkmarkr\checkmarkc\checkmark$\checkmarkr\checkmark$\checkmark^\checkmark\\checkmarkc\checkmarki\checkmarkr\checkmarkc\checkmark$\checkmarke\checkmark$\checkmark^\checkmark\\checkmarkc\checkmarki\checkmarkr\checkmarkc\checkmark$\checkmark \checkmark$\checkmark^\checkmark\\checkmarkc\checkmarki\checkmarkr\checkmarkc\checkmark$\checkmarkF\checkmark$\checkmark^\checkmark\\checkmarkc\checkmarki\checkmarkr\checkmarkc\checkmark$\checkmarko\checkmark$\checkmark^\checkmark\\checkmarkc\checkmarki\checkmarkr\checkmarkc\checkmark$\checkmarkn\checkmark$\checkmark^\checkmark\\checkmarkc\checkmarki\checkmarkr\checkmarkc\checkmark$\checkmarkc\checkmark$\checkmark^\checkmark\\checkmarkc\checkmarki\checkmarkr\checkmarkc\checkmark$\checkmarkt\checkmark$\checkmark^\checkmark\\checkmarkc\checkmarki\checkmarkr\checkmarkc\checkmark$\checkmarki\checkmark$\checkmark^\checkmark\\checkmarkc\checkmarki\checkmarkr\checkmarkc\checkmark$\checkmarko\checkmark$\checkmark^\checkmark\\checkmarkc\checkmarki\checkmarkr\checkmarkc\checkmark$\checkmarkn\checkmark$\checkmark^\checkmark\\checkmarkc\checkmarki\checkmarkr\checkmarkc\checkmark$\checkmarkn\checkmark$\checkmark^\checkmark\\checkmarkc\checkmarki\checkmarkr\checkmarkc\checkmark$\checkmarke\checkmark$\checkmark^\checkmark\\checkmarkc\checkmarki\checkmarkr\checkmarkc\checkmark$\checkmarkl\checkmark$\checkmark^\checkmark\\checkmarkc\checkmarki\checkmarkr\checkmarkc\checkmark$\checkmarkl\checkmark$\checkmark^\checkmark\\checkmarkc\checkmarki\checkmarkr\checkmarkc\checkmark$\checkmarke\checkmark$\checkmark^\checkmark\\checkmarkc\checkmarki\checkmarkr\checkmarkc\checkmark$\checkmark}\checkmark$\checkmark^\checkmark\\checkmarkc\checkmarki\checkmarkr\checkmarkc\checkmark$\checkmark
\checkmark$\checkmark^\checkmark\\checkmarkc\checkmarki\checkmarkr\checkmarkc\checkmark$\checkmarkÀ\checkmark$\checkmark^\checkmark\\checkmarkc\checkmarki\checkmarkr\checkmarkc\checkmark$\checkmark \checkmark$\checkmark^\checkmark\\checkmarkc\checkmarki\checkmarkr\checkmarkc\checkmark$\checkmarkp\checkmark$\checkmark^\checkmark\\checkmarkc\checkmarki\checkmarkr\checkmarkc\checkmark$\checkmarka\checkmark$\checkmark^\checkmark\\checkmarkc\checkmarki\checkmarkr\checkmarkc\checkmark$\checkmarkr\checkmark$\checkmark^\checkmark\\checkmarkc\checkmarki\checkmarkr\checkmarkc\checkmark$\checkmarkt\checkmark$\checkmark^\checkmark\\checkmarkc\checkmarki\checkmarkr\checkmarkc\checkmark$\checkmarki\checkmark$\checkmark^\checkmark\\checkmarkc\checkmarki\checkmarkr\checkmarkc\checkmark$\checkmarkr\checkmark$\checkmark^\checkmark\\checkmarkc\checkmarki\checkmarkr\checkmarkc\checkmark$\checkmark \checkmark$\checkmark^\checkmark\\checkmarkc\checkmarki\checkmarkr\checkmarkc\checkmark$\checkmarkd\checkmark$\checkmark^\checkmark\\checkmarkc\checkmarki\checkmarkr\checkmarkc\checkmark$\checkmarke\checkmark$\checkmark^\checkmark\\checkmarkc\checkmarki\checkmarkr\checkmarkc\checkmark$\checkmark \checkmark$\checkmark^\checkmark\\checkmarkc\checkmarki\checkmarkr\checkmarkc\checkmark$\checkmarkl\checkmark$\checkmark^\checkmark\\checkmarkc\checkmarki\checkmarkr\checkmarkc\checkmark$\checkmarka\checkmark$\checkmark^\checkmark\\checkmarkc\checkmarki\checkmarkr\checkmarkc\checkmark$\checkmark \checkmark$\checkmark^\checkmark\\checkmarkc\checkmarki\checkmarkr\checkmarkc\checkmark$\checkmarkf\checkmark$\checkmark^\checkmark\\checkmarkc\checkmarki\checkmarkr\checkmarkc\checkmark$\checkmarko\checkmark$\checkmark^\checkmark\\checkmarkc\checkmarki\checkmarkr\checkmarkc\checkmark$\checkmarkn\checkmark$\checkmark^\checkmark\\checkmarkc\checkmarki\checkmarkr\checkmarkc\checkmark$\checkmarkc\checkmark$\checkmark^\checkmark\\checkmarkc\checkmarki\checkmarkr\checkmarkc\checkmark$\checkmarkt\checkmark$\checkmark^\checkmark\\checkmarkc\checkmarki\checkmarkr\checkmarkc\checkmark$\checkmarki\checkmark$\checkmark^\checkmark\\checkmarkc\checkmarki\checkmarkr\checkmarkc\checkmark$\checkmarko\checkmark$\checkmark^\checkmark\\checkmarkc\checkmarki\checkmarkr\checkmarkc\checkmark$\checkmarkn\checkmark$\checkmark^\checkmark\\checkmarkc\checkmarki\checkmarkr\checkmarkc\checkmark$\checkmark \checkmark$\checkmark^\checkmark\\checkmarkc\checkmarki\checkmarkr\checkmarkc\checkmark$\checkmarkp\checkmark$\checkmark^\checkmark\\checkmarkc\checkmarki\checkmarkr\checkmarkc\checkmark$\checkmarkr\checkmark$\checkmark^\checkmark\\checkmarkc\checkmarki\checkmarkr\checkmarkc\checkmark$\checkmarki\checkmark$\checkmark^\checkmark\\checkmarkc\checkmarki\checkmarkr\checkmarkc\checkmark$\checkmarkn\checkmark$\checkmark^\checkmark\\checkmarkc\checkmarki\checkmarkr\checkmarkc\checkmark$\checkmarkc\checkmark$\checkmark^\checkmark\\checkmarkc\checkmarki\checkmarkr\checkmarkc\checkmark$\checkmarki\checkmark$\checkmark^\checkmark\\checkmarkc\checkmarki\checkmarkr\checkmarkc\checkmark$\checkmarkp\checkmark$\checkmark^\checkmark\\checkmarkc\checkmarki\checkmarkr\checkmarkc\checkmark$\checkmarka\checkmark$\checkmark^\checkmark\\checkmarkc\checkmarki\checkmarkr\checkmarkc\checkmark$\checkmarkl\checkmark$\checkmark^\checkmark\\checkmarkc\checkmarki\checkmarkr\checkmarkc\checkmark$\checkmarke\checkmark$\checkmark^\checkmark\\checkmarkc\checkmarki\checkmarkr\checkmarkc\checkmark$\checkmark \checkmark$\checkmark^\checkmark\\checkmarkc\checkmarki\checkmarkr\checkmarkc\checkmark$\checkmark(\checkmark$\checkmark^\checkmark\\checkmarkc\checkmarki\checkmarkr\checkmarkc\checkmark$\checkmarkF\checkmark$\checkmark^\checkmark\\checkmarkc\checkmarki\checkmarkr\checkmarkc\checkmark$\checkmarkP\checkmark$\checkmark^\checkmark\\checkmarkc\checkmarki\checkmarkr\checkmarkc\checkmark$\checkmark \checkmark$\checkmark^\checkmark\\checkmarkc\checkmarki\checkmarkr\checkmarkc\checkmark$\checkmark:\checkmark$\checkmark^\checkmark\\checkmarkc\checkmarki\checkmarkr\checkmarkc\checkmark$\checkmark \checkmark$\checkmark^\checkmark\\checkmarkc\checkmarki\checkmarkr\checkmarkc\checkmark$\checkmarkU\checkmark$\checkmark^\checkmark\\checkmarkc\checkmarki\checkmarkr\checkmarkc\checkmark$\checkmarks\checkmark$\checkmark^\checkmark\\checkmarkc\checkmarki\checkmarkr\checkmarkc\checkmark$\checkmarki\checkmark$\checkmark^\checkmark\\checkmarkc\checkmarki\checkmarkr\checkmarkc\checkmark$\checkmarkn\checkmark$\checkmark^\checkmark\\checkmarkc\checkmarki\checkmarkr\checkmarkc\checkmark$\checkmarke\checkmark$\checkmark^\checkmark\\checkmarkc\checkmarki\checkmarkr\checkmarkc\checkmark$\checkmarkr\checkmark$\checkmark^\checkmark\\checkmarkc\checkmarki\checkmarkr\checkmarkc\checkmark$\checkmark \checkmark$\checkmark^\checkmark\\checkmarkc\checkmarki\checkmarkr\checkmarkc\checkmark$\checkmark5\checkmark$\checkmark^\checkmark\\checkmarkc\checkmarki\checkmarkr\checkmarkc\checkmark$\checkmark \checkmark$\checkmark^\checkmark\\checkmarkc\checkmarki\checkmarkr\checkmarkc\checkmark$\checkmarka\checkmark$\checkmark^\checkmark\\checkmarkc\checkmarki\checkmarkr\checkmarkc\checkmark$\checkmarkx\checkmark$\checkmark^\checkmark\\checkmarkc\checkmarki\checkmarkr\checkmarkc\checkmark$\checkmarke\checkmark$\checkmark^\checkmark\\checkmarkc\checkmarki\checkmarkr\checkmarkc\checkmark$\checkmarks\checkmark$\checkmark^\checkmark\\checkmarkc\checkmarki\checkmarkr\checkmarkc\checkmark$\checkmark)\checkmark$\checkmark^\checkmark\\checkmarkc\checkmarki\checkmarkr\checkmarkc\checkmark$\checkmark,\checkmark$\checkmark^\checkmark\\checkmarkc\checkmarki\checkmarkr\checkmarkc\checkmark$\checkmark \checkmark$\checkmark^\checkmark\\checkmarkc\checkmarki\checkmarkr\checkmarkc\checkmark$\checkmarkl\checkmark$\checkmark^\checkmark\\checkmarkc\checkmarki\checkmarkr\checkmarkc\checkmark$\checkmarke\checkmark$\checkmark^\checkmark\\checkmarkc\checkmarki\checkmarkr\checkmarkc\checkmark$\checkmark \checkmark$\checkmark^\checkmark\\checkmarkc\checkmarki\checkmarkr\checkmarkc\checkmark$\checkmarkd\checkmark$\checkmark^\checkmark\\checkmarkc\checkmarki\checkmarkr\checkmarkc\checkmark$\checkmarki\checkmark$\checkmark^\checkmark\\checkmarkc\checkmarki\checkmarkr\checkmarkc\checkmark$\checkmarka\checkmark$\checkmark^\checkmark\\checkmarkc\checkmarki\checkmarkr\checkmarkc\checkmark$\checkmarkg\checkmark$\checkmark^\checkmark\\checkmarkc\checkmarki\checkmarkr\checkmarkc\checkmark$\checkmarkr\checkmark$\checkmark^\checkmark\\checkmarkc\checkmarki\checkmarkr\checkmarkc\checkmark$\checkmarka\checkmark$\checkmark^\checkmark\\checkmarkc\checkmarki\checkmarkr\checkmarkc\checkmark$\checkmarkm\checkmark$\checkmark^\checkmark\\checkmarkc\checkmarki\checkmarkr\checkmarkc\checkmark$\checkmarkm\checkmark$\checkmark^\checkmark\\checkmarkc\checkmarki\checkmarkr\checkmarkc\checkmark$\checkmarke\checkmark$\checkmark^\checkmark\\checkmarkc\checkmarki\checkmarkr\checkmarkc\checkmark$\checkmark \checkmark$\checkmark^\checkmark\\checkmarkc\checkmarki\checkmarkr\checkmarkc\checkmark$\checkmarkF\checkmark$\checkmark^\checkmark\\checkmarkc\checkmarki\checkmarkr\checkmarkc\checkmark$\checkmarkA\checkmark$\checkmark^\checkmark\\checkmarkc\checkmarki\checkmarkr\checkmarkc\checkmark$\checkmarkS\checkmark$\checkmark^\checkmark\\checkmarkc\checkmarki\checkmarkr\checkmarkc\checkmark$\checkmarkT\checkmark$\checkmark^\checkmark\\checkmarkc\checkmarki\checkmarkr\checkmarkc\checkmark$\checkmark \checkmark$\checkmark^\checkmark\\checkmarkc\checkmarki\checkmarkr\checkmarkc\checkmark$\checkmarkd\checkmark$\checkmark^\checkmark\\checkmarkc\checkmarki\checkmarkr\checkmarkc\checkmark$\checkmarké\checkmark$\checkmark^\checkmark\\checkmarkc\checkmarki\checkmarkr\checkmarkc\checkmark$\checkmarkc\checkmark$\checkmark^\checkmark\\checkmarkc\checkmarki\checkmarkr\checkmarkc\checkmark$\checkmarko\checkmark$\checkmark^\checkmark\\checkmarkc\checkmarki\checkmarkr\checkmarkc\checkmark$\checkmarkm\checkmark$\checkmark^\checkmark\\checkmarkc\checkmarki\checkmarkr\checkmarkc\checkmark$\checkmarkp\checkmark$\checkmark^\checkmark\\checkmarkc\checkmarki\checkmarkr\checkmarkc\checkmark$\checkmarko\checkmark$\checkmark^\checkmark\\checkmarkc\checkmarki\checkmarkr\checkmarkc\checkmark$\checkmarks\checkmark$\checkmark^\checkmark\\checkmarkc\checkmarki\checkmarkr\checkmarkc\checkmark$\checkmarke\checkmark$\checkmark^\checkmark\\checkmarkc\checkmarki\checkmarkr\checkmarkc\checkmark$\checkmark \checkmark$\checkmark^\checkmark\\checkmarkc\checkmarki\checkmarkr\checkmarkc\checkmark$\checkmarkl\checkmark$\checkmark^\checkmark\\checkmarkc\checkmarki\checkmarkr\checkmarkc\checkmark$\checkmarke\checkmark$\checkmark^\checkmark\\checkmarkc\checkmarki\checkmarkr\checkmarkc\checkmark$\checkmarks\checkmark$\checkmark^\checkmark\\checkmarkc\checkmarki\checkmarkr\checkmarkc\checkmark$\checkmark \checkmark$\checkmark^\checkmark\\checkmarkc\checkmarki\checkmarkr\checkmarkc\checkmark$\checkmarkf\checkmark$\checkmark^\checkmark\\checkmarkc\checkmarki\checkmarkr\checkmarkc\checkmark$\checkmarko\checkmark$\checkmark^\checkmark\\checkmarkc\checkmarki\checkmarkr\checkmarkc\checkmark$\checkmarkn\checkmark$\checkmark^\checkmark\\checkmarkc\checkmarki\checkmarkr\checkmarkc\checkmark$\checkmarkc\checkmark$\checkmark^\checkmark\\checkmarkc\checkmarki\checkmarkr\checkmarkc\checkmark$\checkmarkt\checkmark$\checkmark^\checkmark\\checkmarkc\checkmarki\checkmarkr\checkmarkc\checkmark$\checkmarki\checkmark$\checkmark^\checkmark\\checkmarkc\checkmarki\checkmarkr\checkmarkc\checkmark$\checkmarko\checkmark$\checkmark^\checkmark\\checkmarkc\checkmarki\checkmarkr\checkmarkc\checkmark$\checkmarkn\checkmark$\checkmark^\checkmark\\checkmarkc\checkmarki\checkmarkr\checkmarkc\checkmark$\checkmarks\checkmark$\checkmark^\checkmark\\checkmarkc\checkmarki\checkmarkr\checkmarkc\checkmark$\checkmark \checkmark$\checkmark^\checkmark\\checkmarkc\checkmarki\checkmarkr\checkmarkc\checkmark$\checkmarkt\checkmark$\checkmark^\checkmark\\checkmarkc\checkmarki\checkmarkr\checkmarkc\checkmark$\checkmarke\checkmark$\checkmark^\checkmark\\checkmarkc\checkmarki\checkmarkr\checkmarkc\checkmark$\checkmarkc\checkmark$\checkmark^\checkmark\\checkmarkc\checkmarki\checkmarkr\checkmarkc\checkmark$\checkmarkh\checkmark$\checkmark^\checkmark\\checkmarkc\checkmarki\checkmarkr\checkmarkc\checkmark$\checkmarkn\checkmark$\checkmark^\checkmark\\checkmarkc\checkmarki\checkmarkr\checkmarkc\checkmark$\checkmarki\checkmark$\checkmark^\checkmark\\checkmarkc\checkmarki\checkmarkr\checkmarkc\checkmark$\checkmarkq\checkmark$\checkmark^\checkmark\\checkmarkc\checkmarki\checkmarkr\checkmarkc\checkmark$\checkmarku\checkmark$\checkmark^\checkmark\\checkmarkc\checkmarki\checkmarkr\checkmarkc\checkmark$\checkmarke\checkmark$\checkmark^\checkmark\\checkmarkc\checkmarki\checkmarkr\checkmarkc\checkmark$\checkmarks\checkmark$\checkmark^\checkmark\\checkmarkc\checkmarki\checkmarkr\checkmarkc\checkmark$\checkmark \checkmark$\checkmark^\checkmark\\checkmarkc\checkmarki\checkmarkr\checkmarkc\checkmark$\checkmark(\checkmark$\checkmark^\checkmark\\checkmarkc\checkmarki\checkmarkr\checkmarkc\checkmark$\checkmarkF\checkmark$\checkmark^\checkmark\\checkmarkc\checkmarki\checkmarkr\checkmarkc\checkmark$\checkmarkT\checkmark$\checkmark^\checkmark\\checkmarkc\checkmarki\checkmarkr\checkmarkc\checkmark$\checkmark)\checkmark$\checkmark^\checkmark\\checkmarkc\checkmarki\checkmarkr\checkmarkc\checkmark$\checkmark \checkmark$\checkmark^\checkmark\\checkmarkc\checkmarki\checkmarkr\checkmarkc\checkmark$\checkmarkj\checkmark$\checkmark^\checkmark\\checkmarkc\checkmarki\checkmarkr\checkmarkc\checkmark$\checkmarku\checkmark$\checkmark^\checkmark\\checkmarkc\checkmarki\checkmarkr\checkmarkc\checkmark$\checkmarks\checkmark$\checkmark^\checkmark\\checkmarkc\checkmarki\checkmarkr\checkmarkc\checkmark$\checkmarkq\checkmark$\checkmark^\checkmark\\checkmarkc\checkmarki\checkmarkr\checkmarkc\checkmark$\checkmarku\checkmark$\checkmark^\checkmark\\checkmarkc\checkmarki\checkmarkr\checkmarkc\checkmark$\checkmark'\checkmark$\checkmark^\checkmark\\checkmarkc\checkmarki\checkmarkr\checkmarkc\checkmark$\checkmarka\checkmark$\checkmark^\checkmark\\checkmarkc\checkmarki\checkmarkr\checkmarkc\checkmark$\checkmarku\checkmark$\checkmark^\checkmark\\checkmarkc\checkmarki\checkmarkr\checkmarkc\checkmark$\checkmarkx\checkmark$\checkmark^\checkmark\\checkmarkc\checkmarki\checkmarkr\checkmarkc\checkmark$\checkmark \checkmark$\checkmark^\checkmark\\checkmarkc\checkmarki\checkmarkr\checkmarkc\checkmark$\checkmarks\checkmark$\checkmark^\checkmark\\checkmarkc\checkmarki\checkmarkr\checkmarkc\checkmark$\checkmarko\checkmark$\checkmark^\checkmark\\checkmarkc\checkmarki\checkmarkr\checkmarkc\checkmark$\checkmarkl\checkmark$\checkmark^\checkmark\\checkmarkc\checkmarki\checkmarkr\checkmarkc\checkmark$\checkmarku\checkmark$\checkmark^\checkmark\\checkmarkc\checkmarki\checkmarkr\checkmarkc\checkmark$\checkmarkt\checkmark$\checkmark^\checkmark\\checkmarkc\checkmarki\checkmarkr\checkmarkc\checkmark$\checkmarki\checkmark$\checkmark^\checkmark\\checkmarkc\checkmarki\checkmarkr\checkmarkc\checkmark$\checkmarko\checkmark$\checkmark^\checkmark\\checkmarkc\checkmarki\checkmarkr\checkmarkc\checkmark$\checkmarkn\checkmark$\checkmark^\checkmark\\checkmarkc\checkmarki\checkmarkr\checkmarkc\checkmark$\checkmarks\checkmark$\checkmark^\checkmark\\checkmarkc\checkmarki\checkmarkr\checkmarkc\checkmark$\checkmark \checkmark$\checkmark^\checkmark\\checkmarkc\checkmarki\checkmarkr\checkmarkc\checkmark$\checkmarkc\checkmark$\checkmark^\checkmark\\checkmarkc\checkmarki\checkmarkr\checkmarkc\checkmark$\checkmarko\checkmark$\checkmark^\checkmark\\checkmarkc\checkmarki\checkmarkr\checkmarkc\checkmark$\checkmarkn\checkmark$\checkmark^\checkmark\\checkmarkc\checkmarki\checkmarkr\checkmarkc\checkmark$\checkmarks\checkmark$\checkmark^\checkmark\\checkmarkc\checkmarki\checkmarkr\checkmarkc\checkmark$\checkmarkt\checkmark$\checkmark^\checkmark\\checkmarkc\checkmarki\checkmarkr\checkmarkc\checkmark$\checkmarkr\checkmark$\checkmark^\checkmark\\checkmarkc\checkmarki\checkmarkr\checkmarkc\checkmark$\checkmarku\checkmark$\checkmark^\checkmark\\checkmarkc\checkmarki\checkmarkr\checkmarkc\checkmark$\checkmarkc\checkmark$\checkmark^\checkmark\\checkmarkc\checkmarki\checkmarkr\checkmarkc\checkmark$\checkmarkt\checkmark$\checkmark^\checkmark\\checkmarkc\checkmarki\checkmarkr\checkmarkc\checkmark$\checkmarki\checkmark$\checkmark^\checkmark\\checkmarkc\checkmarki\checkmarkr\checkmarkc\checkmark$\checkmarkv\checkmark$\checkmark^\checkmark\\checkmarkc\checkmarki\checkmarkr\checkmarkc\checkmark$\checkmarke\checkmark$\checkmark^\checkmark\\checkmarkc\checkmarki\checkmarkr\checkmarkc\checkmark$\checkmarks\checkmark$\checkmark^\checkmark\\checkmarkc\checkmarki\checkmarkr\checkmarkc\checkmark$\checkmark \checkmark$\checkmark^\checkmark\\checkmarkc\checkmarki\checkmarkr\checkmarkc\checkmark$\checkmarkr\checkmark$\checkmark^\checkmark\\checkmarkc\checkmarki\checkmarkr\checkmarkc\checkmark$\checkmarke\checkmark$\checkmark^\checkmark\\checkmarkc\checkmarki\checkmarkr\checkmarkc\checkmark$\checkmarkt\checkmark$\checkmark^\checkmark\\checkmarkc\checkmarki\checkmarkr\checkmarkc\checkmark$\checkmarke\checkmark$\checkmark^\checkmark\\checkmarkc\checkmarki\checkmarkr\checkmarkc\checkmark$\checkmarkn\checkmark$\checkmark^\checkmark\\checkmarkc\checkmarki\checkmarkr\checkmarkc\checkmark$\checkmarku\checkmark$\checkmark^\checkmark\\checkmarkc\checkmarki\checkmarkr\checkmarkc\checkmark$\checkmarke\checkmark$\checkmark^\checkmark\\checkmarkc\checkmarki\checkmarkr\checkmarkc\checkmark$\checkmarks\checkmark$\checkmark^\checkmark\\checkmarkc\checkmarki\checkmarkr\checkmarkc\checkmark$\checkmark.\checkmark$\checkmark^\checkmark\\checkmarkc\checkmarki\checkmarkr\checkmarkc\checkmark$\checkmark
\checkmark$\checkmark^\checkmark\\checkmarkc\checkmarki\checkmarkr\checkmarkc\checkmark$\checkmark
\checkmark$\checkmark^\checkmark\\checkmarkc\checkmarki\checkmarkr\checkmarkc\checkmark$\checkmark\\checkmark$\checkmark^\checkmark\\checkmarkc\checkmarki\checkmarkr\checkmarkc\checkmark$\checkmarkb\checkmark$\checkmark^\checkmark\\checkmarkc\checkmarki\checkmarkr\checkmarkc\checkmark$\checkmarke\checkmark$\checkmark^\checkmark\\checkmarkc\checkmarki\checkmarkr\checkmarkc\checkmark$\checkmarkg\checkmark$\checkmark^\checkmark\\checkmarkc\checkmarki\checkmarkr\checkmarkc\checkmark$\checkmarki\checkmark$\checkmark^\checkmark\\checkmarkc\checkmarki\checkmarkr\checkmarkc\checkmark$\checkmarkn\checkmark$\checkmark^\checkmark\\checkmarkc\checkmarki\checkmarkr\checkmarkc\checkmark$\checkmark{\checkmark$\checkmark^\checkmark\\checkmarkc\checkmarki\checkmarkr\checkmarkc\checkmark$\checkmarkt\checkmark$\checkmark^\checkmark\\checkmarkc\checkmarki\checkmarkr\checkmarkc\checkmark$\checkmarka\checkmark$\checkmark^\checkmark\\checkmarkc\checkmarki\checkmarkr\checkmarkc\checkmark$\checkmarkb\checkmark$\checkmark^\checkmark\\checkmarkc\checkmarki\checkmarkr\checkmarkc\checkmark$\checkmarkl\checkmark$\checkmark^\checkmark\\checkmarkc\checkmarki\checkmarkr\checkmarkc\checkmark$\checkmarke\checkmark$\checkmark^\checkmark\\checkmarkc\checkmarki\checkmarkr\checkmarkc\checkmark$\checkmark}\checkmark$\checkmark^\checkmark\\checkmarkc\checkmarki\checkmarkr\checkmarkc\checkmark$\checkmark[\checkmark$\checkmark^\checkmark\\checkmarkc\checkmarki\checkmarkr\checkmarkc\checkmark$\checkmarkh\checkmark$\checkmark^\checkmark\\checkmarkc\checkmarki\checkmarkr\checkmarkc\checkmark$\checkmarkt\checkmark$\checkmark^\checkmark\\checkmarkc\checkmarki\checkmarkr\checkmarkc\checkmark$\checkmarkb\checkmark$\checkmark^\checkmark\\checkmarkc\checkmarki\checkmarkr\checkmarkc\checkmark$\checkmarkp\checkmark$\checkmark^\checkmark\\checkmarkc\checkmarki\checkmarkr\checkmarkc\checkmark$\checkmark]\checkmark$\checkmark^\checkmark\\checkmarkc\checkmarki\checkmarkr\checkmarkc\checkmark$\checkmark
\checkmark$\checkmark^\checkmark\\checkmarkc\checkmarki\checkmarkr\checkmarkc\checkmark$\checkmark\\checkmark$\checkmark^\checkmark\\checkmarkc\checkmarki\checkmarkr\checkmarkc\checkmark$\checkmarkc\checkmark$\checkmark^\checkmark\\checkmarkc\checkmarki\checkmarkr\checkmarkc\checkmark$\checkmarke\checkmark$\checkmark^\checkmark\\checkmarkc\checkmarki\checkmarkr\checkmarkc\checkmark$\checkmarkn\checkmark$\checkmark^\checkmark\\checkmarkc\checkmarki\checkmarkr\checkmarkc\checkmark$\checkmarkt\checkmark$\checkmark^\checkmark\\checkmarkc\checkmarki\checkmarkr\checkmarkc\checkmark$\checkmarke\checkmark$\checkmark^\checkmark\\checkmarkc\checkmarki\checkmarkr\checkmarkc\checkmark$\checkmarkr\checkmark$\checkmark^\checkmark\\checkmarkc\checkmarki\checkmarkr\checkmarkc\checkmark$\checkmarki\checkmark$\checkmark^\checkmark\\checkmarkc\checkmarki\checkmarkr\checkmarkc\checkmark$\checkmarkn\checkmark$\checkmark^\checkmark\\checkmarkc\checkmarki\checkmarkr\checkmarkc\checkmark$\checkmarkg\checkmark$\checkmark^\checkmark\\checkmarkc\checkmarki\checkmarkr\checkmarkc\checkmark$\checkmark
\checkmark$\checkmark^\checkmark\\checkmarkc\checkmarki\checkmarkr\checkmarkc\checkmark$\checkmark\\checkmark$\checkmark^\checkmark\\checkmarkc\checkmarki\checkmarkr\checkmarkc\checkmark$\checkmarkr\checkmark$\checkmark^\checkmark\\checkmarkc\checkmarki\checkmarkr\checkmarkc\checkmark$\checkmarke\checkmark$\checkmark^\checkmark\\checkmarkc\checkmarki\checkmarkr\checkmarkc\checkmark$\checkmarks\checkmark$\checkmark^\checkmark\\checkmarkc\checkmarki\checkmarkr\checkmarkc\checkmark$\checkmarki\checkmark$\checkmark^\checkmark\\checkmarkc\checkmarki\checkmarkr\checkmarkc\checkmark$\checkmarkz\checkmark$\checkmark^\checkmark\\checkmarkc\checkmarki\checkmarkr\checkmarkc\checkmark$\checkmarke\checkmark$\checkmark^\checkmark\\checkmarkc\checkmarki\checkmarkr\checkmarkc\checkmark$\checkmarkb\checkmark$\checkmark^\checkmark\\checkmarkc\checkmarki\checkmarkr\checkmarkc\checkmark$\checkmarko\checkmark$\checkmark^\checkmark\\checkmarkc\checkmarki\checkmarkr\checkmarkc\checkmark$\checkmarkx\checkmark$\checkmark^\checkmark\\checkmarkc\checkmarki\checkmarkr\checkmarkc\checkmark$\checkmark{\checkmark$\checkmark^\checkmark\\checkmarkc\checkmarki\checkmarkr\checkmarkc\checkmark$\checkmark\\checkmark$\checkmark^\checkmark\\checkmarkc\checkmarki\checkmarkr\checkmarkc\checkmark$\checkmarkt\checkmark$\checkmark^\checkmark\\checkmarkc\checkmarki\checkmarkr\checkmarkc\checkmark$\checkmarke\checkmark$\checkmark^\checkmark\\checkmarkc\checkmarki\checkmarkr\checkmarkc\checkmark$\checkmarkx\checkmark$\checkmark^\checkmark\\checkmarkc\checkmarki\checkmarkr\checkmarkc\checkmark$\checkmarkt\checkmark$\checkmark^\checkmark\\checkmarkc\checkmarki\checkmarkr\checkmarkc\checkmark$\checkmarkw\checkmark$\checkmark^\checkmark\\checkmarkc\checkmarki\checkmarkr\checkmarkc\checkmark$\checkmarki\checkmark$\checkmark^\checkmark\\checkmarkc\checkmarki\checkmarkr\checkmarkc\checkmark$\checkmarkd\checkmark$\checkmark^\checkmark\\checkmarkc\checkmarki\checkmarkr\checkmarkc\checkmark$\checkmarkt\checkmark$\checkmark^\checkmark\\checkmarkc\checkmarki\checkmarkr\checkmarkc\checkmark$\checkmarkh\checkmark$\checkmark^\checkmark\\checkmarkc\checkmarki\checkmarkr\checkmarkc\checkmark$\checkmark}\checkmark$\checkmark^\checkmark\\checkmarkc\checkmarki\checkmarkr\checkmarkc\checkmark$\checkmark{\checkmark$\checkmark^\checkmark\\checkmarkc\checkmarki\checkmarkr\checkmarkc\checkmark$\checkmark!\checkmark$\checkmark^\checkmark\\checkmarkc\checkmarki\checkmarkr\checkmarkc\checkmark$\checkmark}\checkmark$\checkmark^\checkmark\\checkmarkc\checkmarki\checkmarkr\checkmarkc\checkmark$\checkmark{\checkmark$\checkmark^\checkmark\\checkmarkc\checkmarki\checkmarkr\checkmarkc\checkmark$\checkmark
\checkmark$\checkmark^\checkmark\\checkmarkc\checkmarki\checkmarkr\checkmarkc\checkmark$\checkmark\\checkmark$\checkmark^\checkmark\\checkmarkc\checkmarki\checkmarkr\checkmarkc\checkmark$\checkmarkb\checkmark$\checkmark^\checkmark\\checkmarkc\checkmarki\checkmarkr\checkmarkc\checkmark$\checkmarke\checkmark$\checkmark^\checkmark\\checkmarkc\checkmarki\checkmarkr\checkmarkc\checkmark$\checkmarkg\checkmark$\checkmark^\checkmark\\checkmarkc\checkmarki\checkmarkr\checkmarkc\checkmark$\checkmarki\checkmark$\checkmark^\checkmark\\checkmarkc\checkmarki\checkmarkr\checkmarkc\checkmark$\checkmarkn\checkmark$\checkmark^\checkmark\\checkmarkc\checkmarki\checkmarkr\checkmarkc\checkmark$\checkmark{\checkmark$\checkmark^\checkmark\\checkmarkc\checkmarki\checkmarkr\checkmarkc\checkmark$\checkmarkt\checkmark$\checkmark^\checkmark\\checkmarkc\checkmarki\checkmarkr\checkmarkc\checkmark$\checkmarka\checkmark$\checkmark^\checkmark\\checkmarkc\checkmarki\checkmarkr\checkmarkc\checkmark$\checkmarkb\checkmark$\checkmark^\checkmark\\checkmarkc\checkmarki\checkmarkr\checkmarkc\checkmark$\checkmarku\checkmark$\checkmark^\checkmark\\checkmarkc\checkmarki\checkmarkr\checkmarkc\checkmark$\checkmarkl\checkmark$\checkmark^\checkmark\\checkmarkc\checkmarki\checkmarkr\checkmarkc\checkmark$\checkmarka\checkmark$\checkmark^\checkmark\\checkmarkc\checkmarki\checkmarkr\checkmarkc\checkmark$\checkmarkr\checkmark$\checkmark^\checkmark\\checkmarkc\checkmarki\checkmarkr\checkmarkc\checkmark$\checkmark}\checkmark$\checkmark^\checkmark\\checkmarkc\checkmarki\checkmarkr\checkmarkc\checkmark$\checkmark{\checkmark$\checkmark^\checkmark\\checkmarkc\checkmarki\checkmarkr\checkmarkc\checkmark$\checkmark|\checkmark$\checkmark^\checkmark\\checkmarkc\checkmarki\checkmarkr\checkmarkc\checkmark$\checkmarkl\checkmark$\checkmark^\checkmark\\checkmarkc\checkmarki\checkmarkr\checkmarkc\checkmark$\checkmark|\checkmark$\checkmark^\checkmark\\checkmarkc\checkmarki\checkmarkr\checkmarkc\checkmark$\checkmarkl\checkmark$\checkmark^\checkmark\\checkmarkc\checkmarki\checkmarkr\checkmarkc\checkmark$\checkmark|\checkmark$\checkmark^\checkmark\\checkmarkc\checkmarki\checkmarkr\checkmarkc\checkmark$\checkmarkl\checkmark$\checkmark^\checkmark\\checkmarkc\checkmarki\checkmarkr\checkmarkc\checkmark$\checkmark|\checkmark$\checkmark^\checkmark\\checkmarkc\checkmarki\checkmarkr\checkmarkc\checkmark$\checkmark}\checkmark$\checkmark^\checkmark\\checkmarkc\checkmarki\checkmarkr\checkmarkc\checkmark$\checkmark
\checkmark$\checkmark^\checkmark\\checkmarkc\checkmarki\checkmarkr\checkmarkc\checkmark$\checkmark\\checkmark$\checkmark^\checkmark\\checkmarkc\checkmarki\checkmarkr\checkmarkc\checkmark$\checkmarkh\checkmark$\checkmark^\checkmark\\checkmarkc\checkmarki\checkmarkr\checkmarkc\checkmark$\checkmarkl\checkmark$\checkmark^\checkmark\\checkmarkc\checkmarki\checkmarkr\checkmarkc\checkmark$\checkmarki\checkmark$\checkmark^\checkmark\\checkmarkc\checkmarki\checkmarkr\checkmarkc\checkmark$\checkmarkn\checkmark$\checkmark^\checkmark\\checkmarkc\checkmarki\checkmarkr\checkmarkc\checkmark$\checkmarke\checkmark$\checkmark^\checkmark\\checkmarkc\checkmarki\checkmarkr\checkmarkc\checkmark$\checkmark
\checkmark$\checkmark^\checkmark\\checkmarkc\checkmarki\checkmarkr\checkmarkc\checkmark$\checkmark\\checkmark$\checkmark^\checkmark\\checkmarkc\checkmarki\checkmarkr\checkmarkc\checkmark$\checkmarkt\checkmark$\checkmark^\checkmark\\checkmarkc\checkmarki\checkmarkr\checkmarkc\checkmark$\checkmarke\checkmark$\checkmark^\checkmark\\checkmarkc\checkmarki\checkmarkr\checkmarkc\checkmark$\checkmarkx\checkmark$\checkmark^\checkmark\\checkmarkc\checkmarki\checkmarkr\checkmarkc\checkmark$\checkmarkt\checkmark$\checkmark^\checkmark\\checkmarkc\checkmarki\checkmarkr\checkmarkc\checkmark$\checkmarkb\checkmark$\checkmark^\checkmark\\checkmarkc\checkmarki\checkmarkr\checkmarkc\checkmark$\checkmarkf\checkmark$\checkmark^\checkmark\\checkmarkc\checkmarki\checkmarkr\checkmarkc\checkmark$\checkmark{\checkmark$\checkmark^\checkmark\\checkmarkc\checkmarki\checkmarkr\checkmarkc\checkmark$\checkmarkF\checkmark$\checkmark^\checkmark\\checkmarkc\checkmarki\checkmarkr\checkmarkc\checkmark$\checkmarko\checkmark$\checkmark^\checkmark\\checkmarkc\checkmarki\checkmarkr\checkmarkc\checkmark$\checkmarkn\checkmark$\checkmark^\checkmark\\checkmarkc\checkmarki\checkmarkr\checkmarkc\checkmark$\checkmarkc\checkmark$\checkmark^\checkmark\\checkmarkc\checkmarki\checkmarkr\checkmarkc\checkmark$\checkmarkt\checkmark$\checkmark^\checkmark\\checkmarkc\checkmarki\checkmarkr\checkmarkc\checkmark$\checkmarki\checkmark$\checkmark^\checkmark\\checkmarkc\checkmarki\checkmarkr\checkmarkc\checkmark$\checkmarko\checkmark$\checkmark^\checkmark\\checkmarkc\checkmarki\checkmarkr\checkmarkc\checkmark$\checkmarkn\checkmark$\checkmark^\checkmark\\checkmarkc\checkmarki\checkmarkr\checkmarkc\checkmark$\checkmark \checkmark$\checkmark^\checkmark\\checkmarkc\checkmarki\checkmarkr\checkmarkc\checkmark$\checkmarkT\checkmark$\checkmark^\checkmark\\checkmarkc\checkmarki\checkmarkr\checkmarkc\checkmark$\checkmarke\checkmark$\checkmark^\checkmark\\checkmarkc\checkmarki\checkmarkr\checkmarkc\checkmark$\checkmarkc\checkmark$\checkmark^\checkmark\\checkmarkc\checkmarki\checkmarkr\checkmarkc\checkmark$\checkmarkh\checkmark$\checkmark^\checkmark\\checkmarkc\checkmarki\checkmarkr\checkmarkc\checkmark$\checkmarkn\checkmark$\checkmark^\checkmark\\checkmarkc\checkmarki\checkmarkr\checkmarkc\checkmark$\checkmarki\checkmark$\checkmark^\checkmark\\checkmarkc\checkmarki\checkmarkr\checkmarkc\checkmark$\checkmarkq\checkmark$\checkmark^\checkmark\\checkmarkc\checkmarki\checkmarkr\checkmarkc\checkmark$\checkmarku\checkmark$\checkmark^\checkmark\\checkmarkc\checkmarki\checkmarkr\checkmarkc\checkmark$\checkmarke\checkmark$\checkmark^\checkmark\\checkmarkc\checkmarki\checkmarkr\checkmarkc\checkmark$\checkmark}\checkmark$\checkmark^\checkmark\\checkmarkc\checkmarki\checkmarkr\checkmarkc\checkmark$\checkmark \checkmark$\checkmark^\checkmark\\checkmarkc\checkmarki\checkmarkr\checkmarkc\checkmark$\checkmark&\checkmark$\checkmark^\checkmark\\checkmarkc\checkmarki\checkmarkr\checkmarkc\checkmark$\checkmark \checkmark$\checkmark^\checkmark\\checkmarkc\checkmarki\checkmarkr\checkmarkc\checkmark$\checkmark\\checkmark$\checkmark^\checkmark\\checkmarkc\checkmarki\checkmarkr\checkmarkc\checkmark$\checkmarkt\checkmark$\checkmark^\checkmark\\checkmarkc\checkmarki\checkmarkr\checkmarkc\checkmark$\checkmarke\checkmark$\checkmark^\checkmark\\checkmarkc\checkmarki\checkmarkr\checkmarkc\checkmark$\checkmarkx\checkmark$\checkmark^\checkmark\\checkmarkc\checkmarki\checkmarkr\checkmarkc\checkmark$\checkmarkt\checkmark$\checkmark^\checkmark\\checkmarkc\checkmarki\checkmarkr\checkmarkc\checkmark$\checkmarkb\checkmark$\checkmark^\checkmark\\checkmarkc\checkmarki\checkmarkr\checkmarkc\checkmark$\checkmarkf\checkmark$\checkmark^\checkmark\\checkmarkc\checkmarki\checkmarkr\checkmarkc\checkmark$\checkmark{\checkmark$\checkmark^\checkmark\\checkmarkc\checkmarki\checkmarkr\checkmarkc\checkmark$\checkmarkS\checkmark$\checkmark^\checkmark\\checkmarkc\checkmarki\checkmarkr\checkmarkc\checkmark$\checkmarko\checkmark$\checkmark^\checkmark\\checkmarkc\checkmarki\checkmarkr\checkmarkc\checkmark$\checkmarku\checkmark$\checkmark^\checkmark\\checkmarkc\checkmarki\checkmarkr\checkmarkc\checkmark$\checkmarks\checkmark$\checkmark^\checkmark\\checkmarkc\checkmarki\checkmarkr\checkmarkc\checkmark$\checkmark-\checkmark$\checkmark^\checkmark\\checkmarkc\checkmarki\checkmarkr\checkmarkc\checkmark$\checkmarkF\checkmark$\checkmark^\checkmark\\checkmarkc\checkmarki\checkmarkr\checkmarkc\checkmark$\checkmarko\checkmark$\checkmark^\checkmark\\checkmarkc\checkmarki\checkmarkr\checkmarkc\checkmark$\checkmarkn\checkmark$\checkmark^\checkmark\\checkmarkc\checkmarki\checkmarkr\checkmarkc\checkmark$\checkmarkc\checkmark$\checkmark^\checkmark\\checkmarkc\checkmarki\checkmarkr\checkmarkc\checkmark$\checkmarkt\checkmark$\checkmark^\checkmark\\checkmarkc\checkmarki\checkmarkr\checkmarkc\checkmark$\checkmarki\checkmark$\checkmark^\checkmark\\checkmarkc\checkmarki\checkmarkr\checkmarkc\checkmark$\checkmarko\checkmark$\checkmark^\checkmark\\checkmarkc\checkmarki\checkmarkr\checkmarkc\checkmark$\checkmarkn\checkmark$\checkmark^\checkmark\\checkmarkc\checkmarki\checkmarkr\checkmarkc\checkmark$\checkmark}\checkmark$\checkmark^\checkmark\\checkmarkc\checkmarki\checkmarkr\checkmarkc\checkmark$\checkmark \checkmark$\checkmark^\checkmark\\checkmarkc\checkmarki\checkmarkr\checkmarkc\checkmark$\checkmark&\checkmark$\checkmark^\checkmark\\checkmarkc\checkmarki\checkmarkr\checkmarkc\checkmark$\checkmark \checkmark$\checkmark^\checkmark\\checkmarkc\checkmarki\checkmarkr\checkmarkc\checkmark$\checkmark\\checkmark$\checkmark^\checkmark\\checkmarkc\checkmarki\checkmarkr\checkmarkc\checkmark$\checkmarkt\checkmark$\checkmark^\checkmark\\checkmarkc\checkmarki\checkmarkr\checkmarkc\checkmark$\checkmarke\checkmark$\checkmark^\checkmark\\checkmarkc\checkmarki\checkmarkr\checkmarkc\checkmark$\checkmarkx\checkmark$\checkmark^\checkmark\\checkmarkc\checkmarki\checkmarkr\checkmarkc\checkmark$\checkmarkt\checkmark$\checkmark^\checkmark\\checkmarkc\checkmarki\checkmarkr\checkmarkc\checkmark$\checkmarkb\checkmark$\checkmark^\checkmark\\checkmarkc\checkmarki\checkmarkr\checkmarkc\checkmark$\checkmarkf\checkmark$\checkmark^\checkmark\\checkmarkc\checkmarki\checkmarkr\checkmarkc\checkmark$\checkmark{\checkmark$\checkmark^\checkmark\\checkmarkc\checkmarki\checkmarkr\checkmarkc\checkmark$\checkmarkS\checkmark$\checkmark^\checkmark\\checkmarkc\checkmarki\checkmarkr\checkmarkc\checkmark$\checkmarko\checkmark$\checkmark^\checkmark\\checkmarkc\checkmarki\checkmarkr\checkmarkc\checkmark$\checkmarkl\checkmark$\checkmark^\checkmark\\checkmarkc\checkmarki\checkmarkr\checkmarkc\checkmark$\checkmarku\checkmark$\checkmark^\checkmark\\checkmarkc\checkmarki\checkmarkr\checkmarkc\checkmark$\checkmarkt\checkmark$\checkmark^\checkmark\\checkmarkc\checkmarki\checkmarkr\checkmarkc\checkmark$\checkmarki\checkmark$\checkmark^\checkmark\\checkmarkc\checkmarki\checkmarkr\checkmarkc\checkmark$\checkmarko\checkmark$\checkmark^\checkmark\\checkmarkc\checkmarki\checkmarkr\checkmarkc\checkmark$\checkmarkn\checkmark$\checkmark^\checkmark\\checkmarkc\checkmarki\checkmarkr\checkmarkc\checkmark$\checkmark \checkmark$\checkmark^\checkmark\\checkmarkc\checkmarki\checkmarkr\checkmarkc\checkmark$\checkmarkC\checkmark$\checkmark^\checkmark\\checkmarkc\checkmarki\checkmarkr\checkmarkc\checkmark$\checkmarko\checkmark$\checkmark^\checkmark\\checkmarkc\checkmarki\checkmarkr\checkmarkc\checkmark$\checkmarkn\checkmark$\checkmark^\checkmark\\checkmarkc\checkmarki\checkmarkr\checkmarkc\checkmark$\checkmarks\checkmark$\checkmark^\checkmark\\checkmarkc\checkmarki\checkmarkr\checkmarkc\checkmark$\checkmarkt\checkmark$\checkmark^\checkmark\\checkmarkc\checkmarki\checkmarkr\checkmarkc\checkmark$\checkmarkr\checkmark$\checkmark^\checkmark\\checkmarkc\checkmarki\checkmarkr\checkmarkc\checkmark$\checkmarku\checkmark$\checkmark^\checkmark\\checkmarkc\checkmarki\checkmarkr\checkmarkc\checkmark$\checkmarkc\checkmark$\checkmark^\checkmark\\checkmarkc\checkmarki\checkmarkr\checkmarkc\checkmark$\checkmarkt\checkmark$\checkmark^\checkmark\\checkmarkc\checkmarki\checkmarkr\checkmarkc\checkmark$\checkmarki\checkmark$\checkmark^\checkmark\\checkmarkc\checkmarki\checkmarkr\checkmarkc\checkmark$\checkmarkv\checkmark$\checkmark^\checkmark\\checkmarkc\checkmarki\checkmarkr\checkmarkc\checkmark$\checkmarke\checkmark$\checkmark^\checkmark\\checkmarkc\checkmarki\checkmarkr\checkmarkc\checkmark$\checkmark}\checkmark$\checkmark^\checkmark\\checkmarkc\checkmarki\checkmarkr\checkmarkc\checkmark$\checkmark \checkmark$\checkmark^\checkmark\\checkmarkc\checkmarki\checkmarkr\checkmarkc\checkmark$\checkmark\\checkmark$\checkmark^\checkmark\\checkmarkc\checkmarki\checkmarkr\checkmarkc\checkmark$\checkmark\\checkmark$\checkmark^\checkmark\\checkmarkc\checkmarki\checkmarkr\checkmarkc\checkmark$\checkmark \checkmark$\checkmark^\checkmark\\checkmarkc\checkmarki\checkmarkr\checkmarkc\checkmark$\checkmark\\checkmark$\checkmark^\checkmark\\checkmarkc\checkmarki\checkmarkr\checkmarkc\checkmark$\checkmarkh\checkmark$\checkmark^\checkmark\\checkmarkc\checkmarki\checkmarkr\checkmarkc\checkmark$\checkmarkl\checkmark$\checkmark^\checkmark\\checkmarkc\checkmarki\checkmarkr\checkmarkc\checkmark$\checkmarki\checkmark$\checkmark^\checkmark\\checkmarkc\checkmarki\checkmarkr\checkmarkc\checkmark$\checkmarkn\checkmark$\checkmark^\checkmark\\checkmarkc\checkmarki\checkmarkr\checkmarkc\checkmark$\checkmarke\checkmark$\checkmark^\checkmark\\checkmarkc\checkmarki\checkmarkr\checkmarkc\checkmark$\checkmark
\checkmark$\checkmark^\checkmark\\checkmarkc\checkmarki\checkmarkr\checkmarkc\checkmark$\checkmark\\checkmark$\checkmark^\checkmark\\checkmarkc\checkmarki\checkmarkr\checkmarkc\checkmark$\checkmarkm\checkmark$\checkmark^\checkmark\\checkmarkc\checkmarki\checkmarkr\checkmarkc\checkmark$\checkmarku\checkmark$\checkmark^\checkmark\\checkmarkc\checkmarki\checkmarkr\checkmarkc\checkmark$\checkmarkl\checkmark$\checkmark^\checkmark\\checkmarkc\checkmarki\checkmarkr\checkmarkc\checkmark$\checkmarkt\checkmark$\checkmark^\checkmark\\checkmarkc\checkmarki\checkmarkr\checkmarkc\checkmark$\checkmarki\checkmark$\checkmark^\checkmark\\checkmarkc\checkmarki\checkmarkr\checkmarkc\checkmark$\checkmarkr\checkmark$\checkmark^\checkmark\\checkmarkc\checkmarki\checkmarkr\checkmarkc\checkmark$\checkmarko\checkmark$\checkmark^\checkmark\\checkmarkc\checkmarki\checkmarkr\checkmarkc\checkmark$\checkmarkw\checkmark$\checkmark^\checkmark\\checkmarkc\checkmarki\checkmarkr\checkmarkc\checkmark$\checkmark{\checkmark$\checkmark^\checkmark\\checkmarkc\checkmarki\checkmarkr\checkmarkc\checkmark$\checkmark2\checkmark$\checkmark^\checkmark\\checkmarkc\checkmarki\checkmarkr\checkmarkc\checkmark$\checkmark}\checkmark$\checkmark^\checkmark\\checkmarkc\checkmarki\checkmarkr\checkmarkc\checkmark$\checkmark{\checkmark$\checkmark^\checkmark\\checkmarkc\checkmarki\checkmarkr\checkmarkc\checkmark$\checkmark*\checkmark$\checkmark^\checkmark\\checkmarkc\checkmarki\checkmarkr\checkmarkc\checkmark$\checkmark}\checkmark$\checkmark^\checkmark\\checkmarkc\checkmarki\checkmarkr\checkmarkc\checkmark$\checkmark{\checkmark$\checkmark^\checkmark\\checkmarkc\checkmarki\checkmarkr\checkmarkc\checkmark$\checkmarkF\checkmark$\checkmark^\checkmark\\checkmarkc\checkmarki\checkmarkr\checkmarkc\checkmark$\checkmarkT\checkmark$\checkmark^\checkmark\\checkmarkc\checkmarki\checkmarkr\checkmarkc\checkmark$\checkmark1\checkmark$\checkmark^\checkmark\\checkmarkc\checkmarki\checkmarkr\checkmarkc\checkmark$\checkmark \checkmark$\checkmark^\checkmark\\checkmarkc\checkmarki\checkmarkr\checkmarkc\checkmark$\checkmark:\checkmark$\checkmark^\checkmark\\checkmarkc\checkmarki\checkmarkr\checkmarkc\checkmark$\checkmark \checkmark$\checkmark^\checkmark\\checkmarkc\checkmarki\checkmarkr\checkmarkc\checkmark$\checkmarkA\checkmark$\checkmark^\checkmark\\checkmarkc\checkmarki\checkmarkr\checkmarkc\checkmark$\checkmarkn\checkmark$\checkmark^\checkmark\\checkmarkc\checkmarki\checkmarkr\checkmarkc\checkmark$\checkmarki\checkmark$\checkmark^\checkmark\\checkmarkc\checkmarki\checkmarkr\checkmarkc\checkmark$\checkmarkm\checkmark$\checkmark^\checkmark\\checkmarkc\checkmarki\checkmarkr\checkmarkc\checkmark$\checkmarke\checkmark$\checkmark^\checkmark\\checkmarkc\checkmarki\checkmarkr\checkmarkc\checkmark$\checkmarkr\checkmark$\checkmark^\checkmark\\checkmarkc\checkmarki\checkmarkr\checkmarkc\checkmark$\checkmark \checkmark$\checkmark^\checkmark\\checkmarkc\checkmarki\checkmarkr\checkmarkc\checkmark$\checkmarkl\checkmark$\checkmark^\checkmark\\checkmarkc\checkmarki\checkmarkr\checkmarkc\checkmark$\checkmark'\checkmark$\checkmark^\checkmark\\checkmarkc\checkmarki\checkmarkr\checkmarkc\checkmark$\checkmarko\checkmark$\checkmark^\checkmark\\checkmarkc\checkmarki\checkmarkr\checkmarkc\checkmark$\checkmarku\checkmark$\checkmark^\checkmark\\checkmarkc\checkmarki\checkmarkr\checkmarkc\checkmark$\checkmarkt\checkmark$\checkmark^\checkmark\\checkmarkc\checkmarki\checkmarkr\checkmarkc\checkmark$\checkmarki\checkmark$\checkmark^\checkmark\\checkmarkc\checkmarki\checkmarkr\checkmarkc\checkmark$\checkmarkl\checkmark$\checkmark^\checkmark\\checkmarkc\checkmarki\checkmarkr\checkmarkc\checkmark$\checkmark \checkmark$\checkmark^\checkmark\\checkmarkc\checkmarki\checkmarkr\checkmarkc\checkmark$\checkmark(\checkmark$\checkmark^\checkmark\\checkmarkc\checkmarki\checkmarkr\checkmarkc\checkmark$\checkmarkB\checkmark$\checkmark^\checkmark\\checkmarkc\checkmarki\checkmarkr\checkmarkc\checkmark$\checkmarkr\checkmark$\checkmark^\checkmark\\checkmarkc\checkmarki\checkmarkr\checkmarkc\checkmark$\checkmarko\checkmark$\checkmark^\checkmark\\checkmarkc\checkmarki\checkmarkr\checkmarkc\checkmark$\checkmarkc\checkmark$\checkmark^\checkmark\\checkmarkc\checkmarki\checkmarkr\checkmarkc\checkmark$\checkmarkh\checkmark$\checkmark^\checkmark\\checkmarkc\checkmarki\checkmarkr\checkmarkc\checkmark$\checkmarke\checkmark$\checkmark^\checkmark\\checkmarkc\checkmarki\checkmarkr\checkmarkc\checkmark$\checkmark)\checkmark$\checkmark^\checkmark\\checkmarkc\checkmarki\checkmarkr\checkmarkc\checkmark$\checkmark}\checkmark$\checkmark^\checkmark\\checkmarkc\checkmarki\checkmarkr\checkmarkc\checkmark$\checkmark \checkmark$\checkmark^\checkmark\\checkmarkc\checkmarki\checkmarkr\checkmarkc\checkmark$\checkmark&\checkmark$\checkmark^\checkmark\\checkmarkc\checkmarki\checkmarkr\checkmarkc\checkmark$\checkmark \checkmark$\checkmark^\checkmark\\checkmarkc\checkmarki\checkmarkr\checkmarkc\checkmark$\checkmarkF\checkmark$\checkmark^\checkmark\\checkmarkc\checkmarki\checkmarkr\checkmarkc\checkmark$\checkmarkT\checkmark$\checkmark^\checkmark\\checkmarkc\checkmarki\checkmarkr\checkmarkc\checkmark$\checkmark1\checkmark$\checkmark^\checkmark\\checkmarkc\checkmarki\checkmarkr\checkmarkc\checkmark$\checkmark1\checkmark$\checkmark^\checkmark\\checkmarkc\checkmarki\checkmarkr\checkmarkc\checkmark$\checkmark \checkmark$\checkmark^\checkmark\\checkmarkc\checkmarki\checkmarkr\checkmarkc\checkmark$\checkmark:\checkmark$\checkmark^\checkmark\\checkmarkc\checkmarki\checkmarkr\checkmarkc\checkmark$\checkmark \checkmark$\checkmark^\checkmark\\checkmarkc\checkmarki\checkmarkr\checkmarkc\checkmark$\checkmarkC\checkmark$\checkmark^\checkmark\\checkmarkc\checkmarki\checkmarkr\checkmarkc\checkmark$\checkmarko\checkmark$\checkmark^\checkmark\\checkmarkc\checkmarki\checkmarkr\checkmarkc\checkmark$\checkmarkn\checkmark$\checkmark^\checkmark\\checkmarkc\checkmarki\checkmarkr\checkmarkc\checkmark$\checkmarkv\checkmark$\checkmark^\checkmark\\checkmarkc\checkmarki\checkmarkr\checkmarkc\checkmark$\checkmarke\checkmark$\checkmark^\checkmark\\checkmarkc\checkmarki\checkmarkr\checkmarkc\checkmark$\checkmarkr\checkmark$\checkmark^\checkmark\\checkmarkc\checkmarki\checkmarkr\checkmarkc\checkmark$\checkmarkt\checkmark$\checkmark^\checkmark\\checkmarkc\checkmarki\checkmarkr\checkmarkc\checkmark$\checkmarki\checkmark$\checkmark^\checkmark\\checkmarkc\checkmarki\checkmarkr\checkmarkc\checkmark$\checkmarkr\checkmark$\checkmark^\checkmark\\checkmarkc\checkmarki\checkmarkr\checkmarkc\checkmark$\checkmark \checkmark$\checkmark^\checkmark\\checkmarkc\checkmarki\checkmarkr\checkmarkc\checkmark$\checkmarkl\checkmark$\checkmark^\checkmark\\checkmarkc\checkmarki\checkmarkr\checkmarkc\checkmark$\checkmark'\checkmark$\checkmark^\checkmark\\checkmarkc\checkmarki\checkmarkr\checkmarkc\checkmark$\checkmarké\checkmark$\checkmark^\checkmark\\checkmarkc\checkmarki\checkmarkr\checkmarkc\checkmark$\checkmarkn\checkmark$\checkmark^\checkmark\\checkmarkc\checkmarki\checkmarkr\checkmarkc\checkmark$\checkmarke\checkmark$\checkmark^\checkmark\\checkmarkc\checkmarki\checkmarkr\checkmarkc\checkmark$\checkmarkr\checkmark$\checkmark^\checkmark\\checkmarkc\checkmarki\checkmarkr\checkmarkc\checkmark$\checkmarkg\checkmark$\checkmark^\checkmark\\checkmarkc\checkmarki\checkmarkr\checkmarkc\checkmark$\checkmarki\checkmark$\checkmark^\checkmark\\checkmarkc\checkmarki\checkmarkr\checkmarkc\checkmark$\checkmarke\checkmark$\checkmark^\checkmark\\checkmarkc\checkmarki\checkmarkr\checkmarkc\checkmark$\checkmark \checkmark$\checkmark^\checkmark\\checkmarkc\checkmarki\checkmarkr\checkmarkc\checkmark$\checkmarké\checkmark$\checkmark^\checkmark\\checkmarkc\checkmarki\checkmarkr\checkmarkc\checkmark$\checkmarkl\checkmark$\checkmark^\checkmark\\checkmarkc\checkmarki\checkmarkr\checkmarkc\checkmark$\checkmarke\checkmark$\checkmark^\checkmark\\checkmarkc\checkmarki\checkmarkr\checkmarkc\checkmark$\checkmarkc\checkmark$\checkmark^\checkmark\\checkmarkc\checkmarki\checkmarkr\checkmarkc\checkmark$\checkmarkt\checkmark$\checkmark^\checkmark\\checkmarkc\checkmarki\checkmarkr\checkmarkc\checkmark$\checkmarkr\checkmark$\checkmark^\checkmark\\checkmarkc\checkmarki\checkmarkr\checkmarkc\checkmark$\checkmarki\checkmark$\checkmark^\checkmark\\checkmarkc\checkmarki\checkmarkr\checkmarkc\checkmark$\checkmarkq\checkmark$\checkmark^\checkmark\\checkmarkc\checkmarki\checkmarkr\checkmarkc\checkmark$\checkmarku\checkmark$\checkmark^\checkmark\\checkmarkc\checkmarki\checkmarkr\checkmarkc\checkmark$\checkmarke\checkmark$\checkmark^\checkmark\\checkmarkc\checkmarki\checkmarkr\checkmarkc\checkmark$\checkmark \checkmark$\checkmark^\checkmark\\checkmarkc\checkmarki\checkmarkr\checkmarkc\checkmark$\checkmark&\checkmark$\checkmark^\checkmark\\checkmarkc\checkmarki\checkmarkr\checkmarkc\checkmark$\checkmark \checkmark$\checkmark^\checkmark\\checkmarkc\checkmarki\checkmarkr\checkmarkc\checkmark$\checkmarkM\checkmark$\checkmark^\checkmark\\checkmarkc\checkmarki\checkmarkr\checkmarkc\checkmark$\checkmarko\checkmark$\checkmark^\checkmark\\checkmarkc\checkmarki\checkmarkr\checkmarkc\checkmark$\checkmarkt\checkmark$\checkmark^\checkmark\\checkmarkc\checkmarki\checkmarkr\checkmarkc\checkmark$\checkmarke\checkmark$\checkmark^\checkmark\\checkmarkc\checkmarki\checkmarkr\checkmarkc\checkmark$\checkmarku\checkmark$\checkmark^\checkmark\\checkmarkc\checkmarki\checkmarkr\checkmarkc\checkmark$\checkmarkr\checkmark$\checkmark^\checkmark\\checkmarkc\checkmarki\checkmarkr\checkmarkc\checkmark$\checkmark \checkmark$\checkmark^\checkmark\\checkmarkc\checkmarki\checkmarkr\checkmarkc\checkmark$\checkmarkB\checkmark$\checkmark^\checkmark\\checkmarkc\checkmarki\checkmarkr\checkmarkc\checkmark$\checkmarkr\checkmark$\checkmark^\checkmark\\checkmarkc\checkmarki\checkmarkr\checkmarkc\checkmark$\checkmarku\checkmark$\checkmark^\checkmark\\checkmarkc\checkmarki\checkmarkr\checkmarkc\checkmark$\checkmarks\checkmark$\checkmark^\checkmark\\checkmarkc\checkmarki\checkmarkr\checkmarkc\checkmark$\checkmarkh\checkmark$\checkmark^\checkmark\\checkmarkc\checkmarki\checkmarkr\checkmarkc\checkmark$\checkmarkl\checkmark$\checkmark^\checkmark\\checkmarkc\checkmarki\checkmarkr\checkmarkc\checkmark$\checkmarke\checkmark$\checkmark^\checkmark\\checkmarkc\checkmarki\checkmarkr\checkmarkc\checkmark$\checkmarks\checkmark$\checkmark^\checkmark\\checkmarkc\checkmarki\checkmarkr\checkmarkc\checkmark$\checkmarks\checkmark$\checkmark^\checkmark\\checkmarkc\checkmarki\checkmarkr\checkmarkc\checkmark$\checkmark \checkmark$\checkmark^\checkmark\\checkmarkc\checkmarki\checkmarkr\checkmarkc\checkmark$\checkmarkD\checkmark$\checkmark^\checkmark\\checkmarkc\checkmarki\checkmarkr\checkmarkc\checkmark$\checkmarkC\checkmark$\checkmark^\checkmark\\checkmarkc\checkmarki\checkmarkr\checkmarkc\checkmark$\checkmark \checkmark$\checkmark^\checkmark\\checkmarkc\checkmarki\checkmarkr\checkmarkc\checkmark$\checkmark3\checkmark$\checkmark^\checkmark\\checkmarkc\checkmarki\checkmarkr\checkmarkc\checkmark$\checkmark0\checkmark$\checkmark^\checkmark\\checkmarkc\checkmarki\checkmarkr\checkmarkc\checkmark$\checkmark0\checkmark$\checkmark^\checkmark\\checkmarkc\checkmarki\checkmarkr\checkmarkc\checkmark$\checkmarkW\checkmark$\checkmark^\checkmark\\checkmarkc\checkmarki\checkmarkr\checkmarkc\checkmark$\checkmark \checkmark$\checkmark^\checkmark\\checkmarkc\checkmarki\checkmarkr\checkmarkc\checkmark$\checkmark\\checkmark$\checkmark^\checkmark\\checkmarkc\checkmarki\checkmarkr\checkmarkc\checkmark$\checkmark\\checkmark$\checkmark^\checkmark\\checkmarkc\checkmarki\checkmarkr\checkmarkc\checkmark$\checkmark \checkmark$\checkmark^\checkmark\\checkmarkc\checkmarki\checkmarkr\checkmarkc\checkmark$\checkmark\\checkmark$\checkmark^\checkmark\\checkmarkc\checkmarki\checkmarkr\checkmarkc\checkmark$\checkmarkc\checkmark$\checkmark^\checkmark\\checkmarkc\checkmarki\checkmarkr\checkmarkc\checkmark$\checkmarkl\checkmark$\checkmark^\checkmark\\checkmarkc\checkmarki\checkmarkr\checkmarkc\checkmark$\checkmarki\checkmark$\checkmark^\checkmark\\checkmarkc\checkmarki\checkmarkr\checkmarkc\checkmark$\checkmarkn\checkmark$\checkmark^\checkmark\\checkmarkc\checkmarki\checkmarkr\checkmarkc\checkmark$\checkmarke\checkmark$\checkmark^\checkmark\\checkmarkc\checkmarki\checkmarkr\checkmarkc\checkmark$\checkmark{\checkmark$\checkmark^\checkmark\\checkmarkc\checkmarki\checkmarkr\checkmarkc\checkmark$\checkmark2\checkmark$\checkmark^\checkmark\\checkmarkc\checkmarki\checkmarkr\checkmarkc\checkmark$\checkmark-\checkmark$\checkmark^\checkmark\\checkmarkc\checkmarki\checkmarkr\checkmarkc\checkmark$\checkmark3\checkmark$\checkmark^\checkmark\\checkmarkc\checkmarki\checkmarkr\checkmarkc\checkmark$\checkmark}\checkmark$\checkmark^\checkmark\\checkmarkc\checkmarki\checkmarkr\checkmarkc\checkmark$\checkmark
\checkmark$\checkmark^\checkmark\\checkmarkc\checkmarki\checkmarkr\checkmarkc\checkmark$\checkmark \checkmark$\checkmark^\checkmark\\checkmarkc\checkmarki\checkmarkr\checkmarkc\checkmark$\checkmark&\checkmark$\checkmark^\checkmark\\checkmarkc\checkmarki\checkmarkr\checkmarkc\checkmark$\checkmark \checkmark$\checkmark^\checkmark\\checkmarkc\checkmarki\checkmarkr\checkmarkc\checkmark$\checkmarkF\checkmark$\checkmark^\checkmark\\checkmarkc\checkmarki\checkmarkr\checkmarkc\checkmark$\checkmarkT\checkmark$\checkmark^\checkmark\\checkmarkc\checkmarki\checkmarkr\checkmarkc\checkmark$\checkmark1\checkmark$\checkmark^\checkmark\\checkmarkc\checkmarki\checkmarkr\checkmarkc\checkmark$\checkmark2\checkmark$\checkmark^\checkmark\\checkmarkc\checkmarki\checkmarkr\checkmarkc\checkmark$\checkmark \checkmark$\checkmark^\checkmark\\checkmarkc\checkmarki\checkmarkr\checkmarkc\checkmark$\checkmark:\checkmark$\checkmark^\checkmark\\checkmarkc\checkmarki\checkmarkr\checkmarkc\checkmark$\checkmark \checkmark$\checkmark^\checkmark\\checkmarkc\checkmarki\checkmarkr\checkmarkc\checkmark$\checkmarkG\checkmark$\checkmark^\checkmark\\checkmarkc\checkmarki\checkmarkr\checkmarkc\checkmark$\checkmarku\checkmark$\checkmark^\checkmark\\checkmarkc\checkmarki\checkmarkr\checkmarkc\checkmark$\checkmarki\checkmark$\checkmark^\checkmark\\checkmarkc\checkmarki\checkmarkr\checkmarkc\checkmark$\checkmarkd\checkmark$\checkmark^\checkmark\\checkmarkc\checkmarki\checkmarkr\checkmarkc\checkmark$\checkmarke\checkmark$\checkmark^\checkmark\\checkmarkc\checkmarki\checkmarkr\checkmarkc\checkmark$\checkmarkr\checkmark$\checkmark^\checkmark\\checkmarkc\checkmarki\checkmarkr\checkmarkc\checkmark$\checkmark \checkmark$\checkmark^\checkmark\\checkmarkc\checkmarki\checkmarkr\checkmarkc\checkmark$\checkmarke\checkmark$\checkmark^\checkmark\\checkmarkc\checkmarki\checkmarkr\checkmarkc\checkmark$\checkmarkn\checkmark$\checkmark^\checkmark\\checkmarkc\checkmarki\checkmarkr\checkmarkc\checkmark$\checkmark \checkmark$\checkmark^\checkmark\\checkmarkc\checkmarki\checkmarkr\checkmarkc\checkmark$\checkmarkr\checkmark$\checkmark^\checkmark\\checkmarkc\checkmarki\checkmarkr\checkmarkc\checkmark$\checkmarko\checkmark$\checkmark^\checkmark\\checkmarkc\checkmarki\checkmarkr\checkmarkc\checkmark$\checkmarkt\checkmark$\checkmark^\checkmark\\checkmarkc\checkmarki\checkmarkr\checkmarkc\checkmark$\checkmarka\checkmark$\checkmark^\checkmark\\checkmarkc\checkmarki\checkmarkr\checkmarkc\checkmark$\checkmarkt\checkmark$\checkmark^\checkmark\\checkmarkc\checkmarki\checkmarkr\checkmarkc\checkmark$\checkmarki\checkmark$\checkmark^\checkmark\\checkmarkc\checkmarki\checkmarkr\checkmarkc\checkmark$\checkmarko\checkmark$\checkmark^\checkmark\\checkmarkc\checkmarki\checkmarkr\checkmarkc\checkmark$\checkmarkn\checkmark$\checkmark^\checkmark\\checkmarkc\checkmarki\checkmarkr\checkmarkc\checkmark$\checkmark \checkmark$\checkmark^\checkmark\\checkmarkc\checkmarki\checkmarkr\checkmarkc\checkmark$\checkmarkà\checkmark$\checkmark^\checkmark\\checkmarkc\checkmarki\checkmarkr\checkmarkc\checkmark$\checkmark \checkmark$\checkmark^\checkmark\\checkmarkc\checkmarki\checkmarkr\checkmarkc\checkmark$\checkmarkh\checkmark$\checkmark^\checkmark\\checkmarkc\checkmarki\checkmarkr\checkmarkc\checkmark$\checkmarka\checkmark$\checkmark^\checkmark\\checkmarkc\checkmarki\checkmarkr\checkmarkc\checkmark$\checkmarku\checkmark$\checkmark^\checkmark\\checkmarkc\checkmarki\checkmarkr\checkmarkc\checkmark$\checkmarkt\checkmark$\checkmark^\checkmark\\checkmarkc\checkmarki\checkmarkr\checkmarkc\checkmark$\checkmarke\checkmark$\checkmark^\checkmark\\checkmarkc\checkmarki\checkmarkr\checkmarkc\checkmark$\checkmark \checkmark$\checkmark^\checkmark\\checkmarkc\checkmarki\checkmarkr\checkmarkc\checkmark$\checkmarkv\checkmark$\checkmark^\checkmark\\checkmarkc\checkmarki\checkmarkr\checkmarkc\checkmark$\checkmarki\checkmark$\checkmark^\checkmark\\checkmarkc\checkmarki\checkmarkr\checkmarkc\checkmark$\checkmarkt\checkmark$\checkmark^\checkmark\\checkmarkc\checkmarki\checkmarkr\checkmarkc\checkmark$\checkmarke\checkmark$\checkmark^\checkmark\\checkmarkc\checkmarki\checkmarkr\checkmarkc\checkmark$\checkmarks\checkmark$\checkmark^\checkmark\\checkmarkc\checkmarki\checkmarkr\checkmarkc\checkmark$\checkmarks\checkmark$\checkmark^\checkmark\\checkmarkc\checkmarki\checkmarkr\checkmarkc\checkmark$\checkmarke\checkmark$\checkmark^\checkmark\\checkmarkc\checkmarki\checkmarkr\checkmarkc\checkmark$\checkmark \checkmark$\checkmark^\checkmark\\checkmarkc\checkmarki\checkmarkr\checkmarkc\checkmark$\checkmark&\checkmark$\checkmark^\checkmark\\checkmarkc\checkmarki\checkmarkr\checkmarkc\checkmark$\checkmark \checkmark$\checkmark^\checkmark\\checkmarkc\checkmarki\checkmarkr\checkmarkc\checkmark$\checkmarkR\checkmark$\checkmark^\checkmark\\checkmarkc\checkmarki\checkmarkr\checkmarkc\checkmark$\checkmarko\checkmark$\checkmark^\checkmark\\checkmarkc\checkmarki\checkmarkr\checkmarkc\checkmark$\checkmarku\checkmark$\checkmark^\checkmark\\checkmarkc\checkmarki\checkmarkr\checkmarkc\checkmark$\checkmarkl\checkmark$\checkmark^\checkmark\\checkmarkc\checkmarki\checkmarkr\checkmarkc\checkmark$\checkmarke\checkmark$\checkmark^\checkmark\\checkmarkc\checkmarki\checkmarkr\checkmarkc\checkmark$\checkmarkm\checkmark$\checkmark^\checkmark\\checkmarkc\checkmarki\checkmarkr\checkmarkc\checkmark$\checkmarke\checkmark$\checkmark^\checkmark\\checkmarkc\checkmarki\checkmarkr\checkmarkc\checkmark$\checkmarkn\checkmark$\checkmark^\checkmark\\checkmarkc\checkmarki\checkmarkr\checkmarkc\checkmark$\checkmarkt\checkmark$\checkmark^\checkmark\\checkmarkc\checkmarki\checkmarkr\checkmarkc\checkmark$\checkmarks\checkmark$\checkmark^\checkmark\\checkmarkc\checkmarki\checkmarkr\checkmarkc\checkmark$\checkmark \checkmark$\checkmark^\checkmark\\checkmarkc\checkmarki\checkmarkr\checkmarkc\checkmark$\checkmarkà\checkmark$\checkmark^\checkmark\\checkmarkc\checkmarki\checkmarkr\checkmarkc\checkmark$\checkmark \checkmark$\checkmark^\checkmark\\checkmarkc\checkmarki\checkmarkr\checkmarkc\checkmark$\checkmarkc\checkmark$\checkmark^\checkmark\\checkmarkc\checkmarki\checkmarkr\checkmarkc\checkmark$\checkmarko\checkmark$\checkmark^\checkmark\\checkmarkc\checkmarki\checkmarkr\checkmarkc\checkmark$\checkmarkn\checkmark$\checkmark^\checkmark\\checkmarkc\checkmarki\checkmarkr\checkmarkc\checkmark$\checkmarkt\checkmark$\checkmark^\checkmark\\checkmarkc\checkmarki\checkmarkr\checkmarkc\checkmark$\checkmarka\checkmark$\checkmark^\checkmark\\checkmarkc\checkmarki\checkmarkr\checkmarkc\checkmark$\checkmarkc\checkmark$\checkmark^\checkmark\\checkmarkc\checkmarki\checkmarkr\checkmarkc\checkmark$\checkmarkt\checkmark$\checkmark^\checkmark\\checkmarkc\checkmarki\checkmarkr\checkmarkc\checkmark$\checkmark \checkmark$\checkmark^\checkmark\\checkmarkc\checkmarki\checkmarkr\checkmarkc\checkmark$\checkmarko\checkmark$\checkmark^\checkmark\\checkmarkc\checkmarki\checkmarkr\checkmarkc\checkmark$\checkmarkb\checkmark$\checkmark^\checkmark\\checkmarkc\checkmarki\checkmarkr\checkmarkc\checkmark$\checkmarkl\checkmark$\checkmark^\checkmark\\checkmarkc\checkmarki\checkmarkr\checkmarkc\checkmark$\checkmarki\checkmark$\checkmark^\checkmark\\checkmarkc\checkmarki\checkmarkr\checkmarkc\checkmark$\checkmarkq\checkmark$\checkmark^\checkmark\\checkmarkc\checkmarki\checkmarkr\checkmarkc\checkmark$\checkmarku\checkmark$\checkmark^\checkmark\\checkmarkc\checkmarki\checkmarkr\checkmarkc\checkmark$\checkmarke\checkmark$\checkmark^\checkmark\\checkmarkc\checkmarki\checkmarkr\checkmarkc\checkmark$\checkmark \checkmark$\checkmark^\checkmark\\checkmarkc\checkmarki\checkmarkr\checkmarkc\checkmark$\checkmark\\checkmark$\checkmark^\checkmark\\checkmarkc\checkmarki\checkmarkr\checkmarkc\checkmark$\checkmark\\checkmark$\checkmark^\checkmark\\checkmarkc\checkmarki\checkmarkr\checkmarkc\checkmark$\checkmark \checkmark$\checkmark^\checkmark\\checkmarkc\checkmarki\checkmarkr\checkmarkc\checkmark$\checkmark\\checkmark$\checkmark^\checkmark\\checkmarkc\checkmarki\checkmarkr\checkmarkc\checkmark$\checkmarkh\checkmark$\checkmark^\checkmark\\checkmarkc\checkmarki\checkmarkr\checkmarkc\checkmark$\checkmarkl\checkmark$\checkmark^\checkmark\\checkmarkc\checkmarki\checkmarkr\checkmarkc\checkmark$\checkmarki\checkmark$\checkmark^\checkmark\\checkmarkc\checkmarki\checkmarkr\checkmarkc\checkmark$\checkmarkn\checkmark$\checkmark^\checkmark\\checkmarkc\checkmarki\checkmarkr\checkmarkc\checkmark$\checkmarke\checkmark$\checkmark^\checkmark\\checkmarkc\checkmarki\checkmarkr\checkmarkc\checkmark$\checkmark
\checkmark$\checkmark^\checkmark\\checkmarkc\checkmarki\checkmarkr\checkmarkc\checkmark$\checkmark\\checkmark$\checkmark^\checkmark\\checkmarkc\checkmarki\checkmarkr\checkmarkc\checkmark$\checkmarkm\checkmark$\checkmark^\checkmark\\checkmarkc\checkmarki\checkmarkr\checkmarkc\checkmark$\checkmarku\checkmark$\checkmark^\checkmark\\checkmarkc\checkmarki\checkmarkr\checkmarkc\checkmark$\checkmarkl\checkmark$\checkmark^\checkmark\\checkmarkc\checkmarki\checkmarkr\checkmarkc\checkmark$\checkmarkt\checkmark$\checkmark^\checkmark\\checkmarkc\checkmarki\checkmarkr\checkmarkc\checkmark$\checkmarki\checkmark$\checkmark^\checkmark\\checkmarkc\checkmarki\checkmarkr\checkmarkc\checkmark$\checkmarkr\checkmark$\checkmark^\checkmark\\checkmarkc\checkmarki\checkmarkr\checkmarkc\checkmark$\checkmarko\checkmark$\checkmark^\checkmark\\checkmarkc\checkmarki\checkmarkr\checkmarkc\checkmark$\checkmarkw\checkmark$\checkmark^\checkmark\\checkmarkc\checkmarki\checkmarkr\checkmarkc\checkmark$\checkmark{\checkmark$\checkmark^\checkmark\\checkmarkc\checkmarki\checkmarkr\checkmarkc\checkmark$\checkmark4\checkmark$\checkmark^\checkmark\\checkmarkc\checkmarki\checkmarkr\checkmarkc\checkmark$\checkmark}\checkmark$\checkmark^\checkmark\\checkmarkc\checkmarki\checkmarkr\checkmarkc\checkmark$\checkmark{\checkmark$\checkmark^\checkmark\\checkmarkc\checkmarki\checkmarkr\checkmarkc\checkmark$\checkmark*\checkmark$\checkmark^\checkmark\\checkmarkc\checkmarki\checkmarkr\checkmarkc\checkmark$\checkmark}\checkmark$\checkmark^\checkmark\\checkmarkc\checkmarki\checkmarkr\checkmarkc\checkmark$\checkmark{\checkmark$\checkmark^\checkmark\\checkmarkc\checkmarki\checkmarkr\checkmarkc\checkmark$\checkmarkF\checkmark$\checkmark^\checkmark\\checkmarkc\checkmarki\checkmarkr\checkmarkc\checkmark$\checkmarkT\checkmark$\checkmark^\checkmark\\checkmarkc\checkmarki\checkmarkr\checkmarkc\checkmark$\checkmark2\checkmark$\checkmark^\checkmark\\checkmarkc\checkmarki\checkmarkr\checkmarkc\checkmark$\checkmark \checkmark$\checkmark^\checkmark\\checkmarkc\checkmarki\checkmarkr\checkmarkc\checkmark$\checkmark:\checkmark$\checkmark^\checkmark\\checkmarkc\checkmarki\checkmarkr\checkmarkc\checkmark$\checkmark \checkmark$\checkmark^\checkmark\\checkmarkc\checkmarki\checkmarkr\checkmarkc\checkmark$\checkmarkD\checkmark$\checkmark^\checkmark\\checkmarkc\checkmarki\checkmarkr\checkmarkc\checkmark$\checkmarké\checkmark$\checkmark^\checkmark\\checkmarkc\checkmarki\checkmarkr\checkmarkc\checkmark$\checkmarkp\checkmark$\checkmark^\checkmark\\checkmarkc\checkmarki\checkmarkr\checkmarkc\checkmark$\checkmarkl\checkmark$\checkmark^\checkmark\\checkmarkc\checkmarki\checkmarkr\checkmarkc\checkmark$\checkmarka\checkmark$\checkmark^\checkmark\\checkmarkc\checkmarki\checkmarkr\checkmarkc\checkmark$\checkmarkc\checkmark$\checkmark^\checkmark\\checkmarkc\checkmarki\checkmarkr\checkmarkc\checkmark$\checkmarke\checkmark$\checkmark^\checkmark\\checkmarkc\checkmarki\checkmarkr\checkmarkc\checkmark$\checkmarkr\checkmark$\checkmark^\checkmark\\checkmarkc\checkmarki\checkmarkr\checkmarkc\checkmark$\checkmark \checkmark$\checkmark^\checkmark\\checkmarkc\checkmarki\checkmarkr\checkmarkc\checkmark$\checkmarkX\checkmark$\checkmark^\checkmark\\checkmarkc\checkmarki\checkmarkr\checkmarkc\checkmark$\checkmark,\checkmark$\checkmark^\checkmark\\checkmarkc\checkmarki\checkmarkr\checkmarkc\checkmark$\checkmark \checkmark$\checkmark^\checkmark\\checkmarkc\checkmarki\checkmarkr\checkmarkc\checkmark$\checkmarkY\checkmark$\checkmark^\checkmark\\checkmarkc\checkmarki\checkmarkr\checkmarkc\checkmark$\checkmark,\checkmark$\checkmark^\checkmark\\checkmarkc\checkmarki\checkmarkr\checkmarkc\checkmark$\checkmark \checkmark$\checkmark^\checkmark\\checkmarkc\checkmarki\checkmarkr\checkmarkc\checkmark$\checkmarkZ\checkmark$\checkmark^\checkmark\\checkmarkc\checkmarki\checkmarkr\checkmarkc\checkmark$\checkmark \checkmark$\checkmark^\checkmark\\checkmarkc\checkmarki\checkmarkr\checkmarkc\checkmark$\checkmark(\checkmark$\checkmark^\checkmark\\checkmarkc\checkmarki\checkmarkr\checkmarkc\checkmark$\checkmarkT\checkmark$\checkmark^\checkmark\\checkmarkc\checkmarki\checkmarkr\checkmarkc\checkmark$\checkmarkr\checkmark$\checkmark^\checkmark\\checkmarkc\checkmarki\checkmarkr\checkmarkc\checkmark$\checkmarka\checkmark$\checkmark^\checkmark\\checkmarkc\checkmarki\checkmarkr\checkmarkc\checkmark$\checkmarkn\checkmark$\checkmark^\checkmark\\checkmarkc\checkmarki\checkmarkr\checkmarkc\checkmark$\checkmarks\checkmark$\checkmark^\checkmark\\checkmarkc\checkmarki\checkmarkr\checkmarkc\checkmark$\checkmarkl\checkmark$\checkmark^\checkmark\\checkmarkc\checkmarki\checkmarkr\checkmarkc\checkmark$\checkmarka\checkmark$\checkmark^\checkmark\\checkmarkc\checkmarki\checkmarkr\checkmarkc\checkmark$\checkmarkt\checkmark$\checkmark^\checkmark\\checkmarkc\checkmarki\checkmarkr\checkmarkc\checkmark$\checkmarki\checkmark$\checkmark^\checkmark\\checkmarkc\checkmarki\checkmarkr\checkmarkc\checkmark$\checkmarko\checkmark$\checkmark^\checkmark\\checkmarkc\checkmarki\checkmarkr\checkmarkc\checkmark$\checkmarkn\checkmark$\checkmark^\checkmark\\checkmarkc\checkmarki\checkmarkr\checkmarkc\checkmark$\checkmark)\checkmark$\checkmark^\checkmark\\checkmarkc\checkmarki\checkmarkr\checkmarkc\checkmark$\checkmark}\checkmark$\checkmark^\checkmark\\checkmarkc\checkmarki\checkmarkr\checkmarkc\checkmark$\checkmark \checkmark$\checkmark^\checkmark\\checkmarkc\checkmarki\checkmarkr\checkmarkc\checkmark$\checkmark&\checkmark$\checkmark^\checkmark\\checkmarkc\checkmarki\checkmarkr\checkmarkc\checkmark$\checkmark \checkmark$\checkmark^\checkmark\\checkmarkc\checkmarki\checkmarkr\checkmarkc\checkmark$\checkmarkF\checkmark$\checkmark^\checkmark\\checkmarkc\checkmarki\checkmarkr\checkmarkc\checkmark$\checkmarkT\checkmark$\checkmark^\checkmark\\checkmarkc\checkmarki\checkmarkr\checkmarkc\checkmark$\checkmark2\checkmark$\checkmark^\checkmark\\checkmarkc\checkmarki\checkmarkr\checkmarkc\checkmark$\checkmark1\checkmark$\checkmark^\checkmark\\checkmarkc\checkmarki\checkmarkr\checkmarkc\checkmark$\checkmark \checkmark$\checkmark^\checkmark\\checkmarkc\checkmarki\checkmarkr\checkmarkc\checkmark$\checkmark:\checkmark$\checkmark^\checkmark\\checkmarkc\checkmarki\checkmarkr\checkmarkc\checkmark$\checkmark \checkmark$\checkmark^\checkmark\\checkmarkc\checkmarki\checkmarkr\checkmarkc\checkmark$\checkmarkG\checkmark$\checkmark^\checkmark\\checkmarkc\checkmarki\checkmarkr\checkmarkc\checkmark$\checkmarku\checkmark$\checkmark^\checkmark\\checkmarkc\checkmarki\checkmarkr\checkmarkc\checkmark$\checkmarki\checkmark$\checkmark^\checkmark\\checkmarkc\checkmarki\checkmarkr\checkmarkc\checkmark$\checkmarkd\checkmark$\checkmark^\checkmark\\checkmarkc\checkmarki\checkmarkr\checkmarkc\checkmark$\checkmarke\checkmark$\checkmark^\checkmark\\checkmarkc\checkmarki\checkmarkr\checkmarkc\checkmark$\checkmarkr\checkmark$\checkmark^\checkmark\\checkmarkc\checkmarki\checkmarkr\checkmarkc\checkmark$\checkmark \checkmark$\checkmark^\checkmark\\checkmarkc\checkmarki\checkmarkr\checkmarkc\checkmark$\checkmarks\checkmark$\checkmark^\checkmark\\checkmarkc\checkmarki\checkmarkr\checkmarkc\checkmark$\checkmarkt\checkmark$\checkmark^\checkmark\\checkmarkc\checkmarki\checkmarkr\checkmarkc\checkmark$\checkmarkr\checkmark$\checkmark^\checkmark\\checkmarkc\checkmarki\checkmarkr\checkmarkc\checkmark$\checkmarki\checkmark$\checkmark^\checkmark\\checkmarkc\checkmarki\checkmarkr\checkmarkc\checkmark$\checkmarkc\checkmark$\checkmark^\checkmark\\checkmarkc\checkmarki\checkmarkr\checkmarkc\checkmark$\checkmarkt\checkmark$\checkmark^\checkmark\\checkmarkc\checkmarki\checkmarkr\checkmarkc\checkmark$\checkmarke\checkmark$\checkmark^\checkmark\\checkmarkc\checkmarki\checkmarkr\checkmarkc\checkmark$\checkmarkm\checkmark$\checkmark^\checkmark\\checkmarkc\checkmarki\checkmarkr\checkmarkc\checkmark$\checkmarke\checkmark$\checkmark^\checkmark\\checkmarkc\checkmarki\checkmarkr\checkmarkc\checkmark$\checkmarkn\checkmark$\checkmark^\checkmark\\checkmarkc\checkmarki\checkmarkr\checkmarkc\checkmark$\checkmarkt\checkmark$\checkmark^\checkmark\\checkmarkc\checkmarki\checkmarkr\checkmarkc\checkmark$\checkmark \checkmark$\checkmark^\checkmark\\checkmarkc\checkmarki\checkmarkr\checkmarkc\checkmark$\checkmarke\checkmark$\checkmark^\checkmark\\checkmarkc\checkmarki\checkmarkr\checkmarkc\checkmark$\checkmarkn\checkmark$\checkmark^\checkmark\\checkmarkc\checkmarki\checkmarkr\checkmarkc\checkmark$\checkmark \checkmark$\checkmark^\checkmark\\checkmarkc\checkmarki\checkmarkr\checkmarkc\checkmark$\checkmarkt\checkmark$\checkmark^\checkmark\\checkmarkc\checkmarki\checkmarkr\checkmarkc\checkmark$\checkmarkr\checkmark$\checkmark^\checkmark\\checkmarkc\checkmarki\checkmarkr\checkmarkc\checkmark$\checkmarka\checkmark$\checkmark^\checkmark\\checkmarkc\checkmarki\checkmarkr\checkmarkc\checkmark$\checkmarkn\checkmark$\checkmark^\checkmark\\checkmarkc\checkmarki\checkmarkr\checkmarkc\checkmark$\checkmarks\checkmark$\checkmark^\checkmark\\checkmarkc\checkmarki\checkmarkr\checkmarkc\checkmark$\checkmarkl\checkmark$\checkmark^\checkmark\\checkmarkc\checkmarki\checkmarkr\checkmarkc\checkmark$\checkmarka\checkmark$\checkmark^\checkmark\\checkmarkc\checkmarki\checkmarkr\checkmarkc\checkmark$\checkmarkt\checkmark$\checkmark^\checkmark\\checkmarkc\checkmarki\checkmarkr\checkmarkc\checkmark$\checkmarki\checkmark$\checkmark^\checkmark\\checkmarkc\checkmarki\checkmarkr\checkmarkc\checkmark$\checkmarko\checkmark$\checkmark^\checkmark\\checkmarkc\checkmarki\checkmarkr\checkmarkc\checkmark$\checkmarkn\checkmark$\checkmark^\checkmark\\checkmarkc\checkmarki\checkmarkr\checkmarkc\checkmark$\checkmark \checkmark$\checkmark^\checkmark\\checkmarkc\checkmarki\checkmarkr\checkmarkc\checkmark$\checkmark&\checkmark$\checkmark^\checkmark\\checkmarkc\checkmarki\checkmarkr\checkmarkc\checkmark$\checkmark \checkmark$\checkmark^\checkmark\\checkmarkc\checkmarki\checkmarkr\checkmarkc\checkmark$\checkmarkR\checkmark$\checkmark^\checkmark\\checkmarkc\checkmarki\checkmarkr\checkmarkc\checkmark$\checkmarka\checkmark$\checkmark^\checkmark\\checkmarkc\checkmarki\checkmarkr\checkmarkc\checkmark$\checkmarki\checkmark$\checkmark^\checkmark\\checkmarkc\checkmarki\checkmarkr\checkmarkc\checkmark$\checkmarkl\checkmark$\checkmark^\checkmark\\checkmarkc\checkmarki\checkmarkr\checkmarkc\checkmark$\checkmarks\checkmark$\checkmark^\checkmark\\checkmarkc\checkmarki\checkmarkr\checkmarkc\checkmark$\checkmark \checkmark$\checkmark^\checkmark\\checkmarkc\checkmarki\checkmarkr\checkmarkc\checkmark$\checkmarkp\checkmark$\checkmark^\checkmark\\checkmarkc\checkmarki\checkmarkr\checkmarkc\checkmark$\checkmarkr\checkmark$\checkmark^\checkmark\\checkmarkc\checkmarki\checkmarkr\checkmarkc\checkmark$\checkmarko\checkmark$\checkmark^\checkmark\\checkmarkc\checkmarki\checkmarkr\checkmarkc\checkmark$\checkmarkf\checkmark$\checkmark^\checkmark\\checkmarkc\checkmarki\checkmarkr\checkmarkc\checkmark$\checkmarki\checkmark$\checkmark^\checkmark\\checkmarkc\checkmarki\checkmarkr\checkmarkc\checkmark$\checkmarkl\checkmark$\checkmark^\checkmark\\checkmarkc\checkmarki\checkmarkr\checkmarkc\checkmark$\checkmarké\checkmark$\checkmark^\checkmark\\checkmarkc\checkmarki\checkmarkr\checkmarkc\checkmark$\checkmarks\checkmark$\checkmark^\checkmark\\checkmarkc\checkmarki\checkmarkr\checkmarkc\checkmark$\checkmark \checkmark$\checkmark^\checkmark\\checkmarkc\checkmarki\checkmarkr\checkmarkc\checkmark$\checkmarkH\checkmark$\checkmark^\checkmark\\checkmarkc\checkmarki\checkmarkr\checkmarkc\checkmark$\checkmarkG\checkmark$\checkmark^\checkmark\\checkmarkc\checkmarki\checkmarkr\checkmarkc\checkmark$\checkmarkR\checkmark$\checkmark^\checkmark\\checkmarkc\checkmarki\checkmarkr\checkmarkc\checkmark$\checkmark1\checkmark$\checkmark^\checkmark\\checkmarkc\checkmarki\checkmarkr\checkmarkc\checkmark$\checkmark5\checkmark$\checkmark^\checkmark\\checkmarkc\checkmarki\checkmarkr\checkmarkc\checkmark$\checkmark \checkmark$\checkmark^\checkmark\\checkmarkc\checkmarki\checkmarkr\checkmarkc\checkmark$\checkmark(\checkmark$\checkmark^\checkmark\\checkmarkc\checkmarki\checkmarkr\checkmarkc\checkmark$\checkmarkH\checkmark$\checkmark^\checkmark\\checkmarkc\checkmarki\checkmarkr\checkmarkc\checkmark$\checkmarkI\checkmark$\checkmark^\checkmark\\checkmarkc\checkmarki\checkmarkr\checkmarkc\checkmark$\checkmarkW\checkmark$\checkmark^\checkmark\\checkmarkc\checkmarki\checkmarkr\checkmarkc\checkmark$\checkmarkI\checkmark$\checkmark^\checkmark\\checkmarkc\checkmarki\checkmarkr\checkmarkc\checkmark$\checkmarkN\checkmark$\checkmark^\checkmark\\checkmarkc\checkmarki\checkmarkr\checkmarkc\checkmark$\checkmark)\checkmark$\checkmark^\checkmark\\checkmarkc\checkmarki\checkmarkr\checkmarkc\checkmark$\checkmark \checkmark$\checkmark^\checkmark\\checkmarkc\checkmarki\checkmarkr\checkmarkc\checkmark$\checkmark\\checkmark$\checkmark^\checkmark\\checkmarkc\checkmarki\checkmarkr\checkmarkc\checkmark$\checkmark\\checkmark$\checkmark^\checkmark\\checkmarkc\checkmarki\checkmarkr\checkmarkc\checkmark$\checkmark \checkmark$\checkmark^\checkmark\\checkmarkc\checkmarki\checkmarkr\checkmarkc\checkmark$\checkmark\\checkmark$\checkmark^\checkmark\\checkmarkc\checkmarki\checkmarkr\checkmarkc\checkmark$\checkmarkc\checkmark$\checkmark^\checkmark\\checkmarkc\checkmarki\checkmarkr\checkmarkc\checkmark$\checkmarkl\checkmark$\checkmark^\checkmark\\checkmarkc\checkmarki\checkmarkr\checkmarkc\checkmark$\checkmarki\checkmark$\checkmark^\checkmark\\checkmarkc\checkmarki\checkmarkr\checkmarkc\checkmark$\checkmarkn\checkmark$\checkmark^\checkmark\\checkmarkc\checkmarki\checkmarkr\checkmarkc\checkmark$\checkmarke\checkmark$\checkmark^\checkmark\\checkmarkc\checkmarki\checkmarkr\checkmarkc\checkmark$\checkmark{\checkmark$\checkmark^\checkmark\\checkmarkc\checkmarki\checkmarkr\checkmarkc\checkmark$\checkmark2\checkmark$\checkmark^\checkmark\\checkmarkc\checkmarki\checkmarkr\checkmarkc\checkmark$\checkmark-\checkmark$\checkmark^\checkmark\\checkmarkc\checkmarki\checkmarkr\checkmarkc\checkmark$\checkmark3\checkmark$\checkmark^\checkmark\\checkmarkc\checkmarki\checkmarkr\checkmarkc\checkmark$\checkmark}\checkmark$\checkmark^\checkmark\\checkmarkc\checkmarki\checkmarkr\checkmarkc\checkmark$\checkmark
\checkmark$\checkmark^\checkmark\\checkmarkc\checkmarki\checkmarkr\checkmarkc\checkmark$\checkmark \checkmark$\checkmark^\checkmark\\checkmarkc\checkmarki\checkmarkr\checkmarkc\checkmark$\checkmark&\checkmark$\checkmark^\checkmark\\checkmarkc\checkmarki\checkmarkr\checkmarkc\checkmark$\checkmark \checkmark$\checkmark^\checkmark\\checkmarkc\checkmarki\checkmarkr\checkmarkc\checkmark$\checkmarkF\checkmark$\checkmark^\checkmark\\checkmarkc\checkmarki\checkmarkr\checkmarkc\checkmark$\checkmarkT\checkmark$\checkmark^\checkmark\\checkmarkc\checkmarki\checkmarkr\checkmarkc\checkmark$\checkmark2\checkmark$\checkmark^\checkmark\\checkmarkc\checkmarki\checkmarkr\checkmarkc\checkmark$\checkmark2\checkmark$\checkmark^\checkmark\\checkmarkc\checkmarki\checkmarkr\checkmarkc\checkmark$\checkmark \checkmark$\checkmark^\checkmark\\checkmarkc\checkmarki\checkmarkr\checkmarkc\checkmark$\checkmark:\checkmark$\checkmark^\checkmark\\checkmarkc\checkmarki\checkmarkr\checkmarkc\checkmark$\checkmark \checkmark$\checkmark^\checkmark\\checkmarkc\checkmarki\checkmarkr\checkmarkc\checkmark$\checkmarkT\checkmark$\checkmark^\checkmark\\checkmarkc\checkmarki\checkmarkr\checkmarkc\checkmark$\checkmarkr\checkmark$\checkmark^\checkmark\\checkmarkc\checkmarki\checkmarkr\checkmarkc\checkmark$\checkmarka\checkmark$\checkmark^\checkmark\\checkmarkc\checkmarki\checkmarkr\checkmarkc\checkmark$\checkmarkn\checkmark$\checkmark^\checkmark\\checkmarkc\checkmarki\checkmarkr\checkmarkc\checkmark$\checkmarks\checkmark$\checkmark^\checkmark\\checkmarkc\checkmarki\checkmarkr\checkmarkc\checkmark$\checkmarkf\checkmark$\checkmark^\checkmark\\checkmarkc\checkmarki\checkmarkr\checkmarkc\checkmark$\checkmarko\checkmark$\checkmark^\checkmark\\checkmarkc\checkmarki\checkmarkr\checkmarkc\checkmark$\checkmarkr\checkmark$\checkmark^\checkmark\\checkmarkc\checkmarki\checkmarkr\checkmarkc\checkmark$\checkmarkm\checkmark$\checkmark^\checkmark\\checkmarkc\checkmarki\checkmarkr\checkmarkc\checkmark$\checkmarke\checkmark$\checkmark^\checkmark\\checkmarkc\checkmarki\checkmarkr\checkmarkc\checkmark$\checkmarkr\checkmark$\checkmark^\checkmark\\checkmarkc\checkmarki\checkmarkr\checkmarkc\checkmark$\checkmark \checkmark$\checkmark^\checkmark\\checkmarkc\checkmarki\checkmarkr\checkmarkc\checkmark$\checkmarkr\checkmark$\checkmark^\checkmark\\checkmarkc\checkmarki\checkmarkr\checkmarkc\checkmark$\checkmarko\checkmark$\checkmark^\checkmark\\checkmarkc\checkmarki\checkmarkr\checkmarkc\checkmark$\checkmarkt\checkmark$\checkmark^\checkmark\\checkmarkc\checkmarki\checkmarkr\checkmarkc\checkmark$\checkmarka\checkmark$\checkmark^\checkmark\\checkmarkc\checkmarki\checkmarkr\checkmarkc\checkmark$\checkmarkt\checkmark$\checkmark^\checkmark\\checkmarkc\checkmarki\checkmarkr\checkmarkc\checkmark$\checkmarki\checkmark$\checkmark^\checkmark\\checkmarkc\checkmarki\checkmarkr\checkmarkc\checkmark$\checkmarko\checkmark$\checkmark^\checkmark\\checkmarkc\checkmarki\checkmarkr\checkmarkc\checkmark$\checkmarkn\checkmark$\checkmark^\checkmark\\checkmarkc\checkmarki\checkmarkr\checkmarkc\checkmark$\checkmark \checkmark$\checkmark^\checkmark\\checkmarkc\checkmarki\checkmarkr\checkmarkc\checkmark$\checkmark→\checkmark$\checkmark^\checkmark\\checkmarkc\checkmarki\checkmarkr\checkmarkc\checkmark$\checkmark \checkmark$\checkmark^\checkmark\\checkmarkc\checkmarki\checkmarkr\checkmarkc\checkmark$\checkmarkt\checkmark$\checkmark^\checkmark\\checkmarkc\checkmarki\checkmarkr\checkmarkc\checkmark$\checkmarkr\checkmark$\checkmark^\checkmark\\checkmarkc\checkmarki\checkmarkr\checkmarkc\checkmark$\checkmarka\checkmark$\checkmark^\checkmark\\checkmarkc\checkmarki\checkmarkr\checkmarkc\checkmark$\checkmarkn\checkmark$\checkmark^\checkmark\\checkmarkc\checkmarki\checkmarkr\checkmarkc\checkmark$\checkmarks\checkmark$\checkmark^\checkmark\\checkmarkc\checkmarki\checkmarkr\checkmarkc\checkmark$\checkmarkl\checkmark$\checkmark^\checkmark\\checkmarkc\checkmarki\checkmarkr\checkmarkc\checkmark$\checkmarka\checkmark$\checkmark^\checkmark\\checkmarkc\checkmarki\checkmarkr\checkmarkc\checkmark$\checkmarkt\checkmark$\checkmark^\checkmark\\checkmarkc\checkmarki\checkmarkr\checkmarkc\checkmark$\checkmarki\checkmark$\checkmark^\checkmark\\checkmarkc\checkmarki\checkmarkr\checkmarkc\checkmark$\checkmarko\checkmark$\checkmark^\checkmark\\checkmarkc\checkmarki\checkmarkr\checkmarkc\checkmark$\checkmarkn\checkmark$\checkmark^\checkmark\\checkmarkc\checkmarki\checkmarkr\checkmarkc\checkmark$\checkmark \checkmark$\checkmark^\checkmark\\checkmarkc\checkmarki\checkmarkr\checkmarkc\checkmark$\checkmark&\checkmark$\checkmark^\checkmark\\checkmarkc\checkmarki\checkmarkr\checkmarkc\checkmark$\checkmark \checkmark$\checkmark^\checkmark\\checkmarkc\checkmarki\checkmarkr\checkmarkc\checkmark$\checkmarkV\checkmark$\checkmark^\checkmark\\checkmarkc\checkmarki\checkmarkr\checkmarkc\checkmark$\checkmarki\checkmark$\checkmark^\checkmark\\checkmarkc\checkmarki\checkmarkr\checkmarkc\checkmark$\checkmarks\checkmark$\checkmark^\checkmark\\checkmarkc\checkmarki\checkmarkr\checkmarkc\checkmark$\checkmark \checkmark$\checkmark^\checkmark\\checkmarkc\checkmarki\checkmarkr\checkmarkc\checkmark$\checkmarkà\checkmark$\checkmark^\checkmark\\checkmarkc\checkmarki\checkmarkr\checkmarkc\checkmark$\checkmark \checkmark$\checkmark^\checkmark\\checkmarkc\checkmarki\checkmarkr\checkmarkc\checkmark$\checkmarkb\checkmark$\checkmark^\checkmark\\checkmarkc\checkmarki\checkmarkr\checkmarkc\checkmark$\checkmarki\checkmark$\checkmark^\checkmark\\checkmarkc\checkmarki\checkmarkr\checkmarkc\checkmark$\checkmarkl\checkmark$\checkmark^\checkmark\\checkmarkc\checkmarki\checkmarkr\checkmarkc\checkmark$\checkmarkl\checkmark$\checkmark^\checkmark\\checkmarkc\checkmarki\checkmarkr\checkmarkc\checkmark$\checkmarke\checkmark$\checkmark^\checkmark\\checkmarkc\checkmarki\checkmarkr\checkmarkc\checkmark$\checkmarks\checkmark$\checkmark^\checkmark\\checkmarkc\checkmarki\checkmarkr\checkmarkc\checkmark$\checkmark \checkmark$\checkmark^\checkmark\\checkmarkc\checkmarki\checkmarkr\checkmarkc\checkmark$\checkmarkS\checkmark$\checkmark^\checkmark\\checkmarkc\checkmarki\checkmarkr\checkmarkc\checkmark$\checkmarkF\checkmark$\checkmark^\checkmark\\checkmarkc\checkmarki\checkmarkr\checkmarkc\checkmark$\checkmarkU\checkmark$\checkmark^\checkmark\\checkmarkc\checkmarki\checkmarkr\checkmarkc\checkmark$\checkmark1\checkmark$\checkmark^\checkmark\\checkmarkc\checkmarki\checkmarkr\checkmarkc\checkmark$\checkmark6\checkmark$\checkmark^\checkmark\\checkmarkc\checkmarki\checkmarkr\checkmarkc\checkmark$\checkmark0\checkmark$\checkmark^\checkmark\\checkmarkc\checkmarki\checkmarkr\checkmarkc\checkmark$\checkmark5\checkmark$\checkmark^\checkmark\\checkmarkc\checkmarki\checkmarkr\checkmarkc\checkmark$\checkmark \checkmark$\checkmark^\checkmark\\checkmarkc\checkmarki\checkmarkr\checkmarkc\checkmark$\checkmark(\checkmark$\checkmark^\checkmark\\checkmarkc\checkmarki\checkmarkr\checkmarkc\checkmark$\checkmark$\checkmark$\checkmark^\checkmark\\checkmarkc\checkmarki\checkmarkr\checkmarkc\checkmark$\checkmarkP\checkmark$\checkmark^\checkmark\\checkmarkc\checkmarki\checkmarkr\checkmarkc\checkmark$\checkmark_\checkmark$\checkmark^\checkmark\\checkmarkc\checkmarki\checkmarkr\checkmarkc\checkmark$\checkmarkh\checkmark$\checkmark^\checkmark\\checkmarkc\checkmarki\checkmarkr\checkmarkc\checkmark$\checkmark=\checkmark$\checkmark^\checkmark\\checkmarkc\checkmarki\checkmarkr\checkmarkc\checkmark$\checkmark5\checkmark$\checkmark^\checkmark\\checkmarkc\checkmarki\checkmarkr\checkmarkc\checkmark$\checkmark$\checkmark$\checkmark^\checkmark\\checkmarkc\checkmarki\checkmarkr\checkmarkc\checkmark$\checkmark \checkmark$\checkmark^\checkmark\\checkmarkc\checkmarki\checkmarkr\checkmarkc\checkmark$\checkmarkm\checkmark$\checkmark^\checkmark\\checkmarkc\checkmarki\checkmarkr\checkmarkc\checkmark$\checkmarkm\checkmark$\checkmark^\checkmark\\checkmarkc\checkmarki\checkmarkr\checkmarkc\checkmark$\checkmark)\checkmark$\checkmark^\checkmark\\checkmarkc\checkmarki\checkmarkr\checkmarkc\checkmark$\checkmark \checkmark$\checkmark^\checkmark\\checkmarkc\checkmarki\checkmarkr\checkmarkc\checkmark$\checkmark\\checkmark$\checkmark^\checkmark\\checkmarkc\checkmarki\checkmarkr\checkmarkc\checkmark$\checkmark\\checkmark$\checkmark^\checkmark\\checkmarkc\checkmarki\checkmarkr\checkmarkc\checkmark$\checkmark \checkmark$\checkmark^\checkmark\\checkmarkc\checkmarki\checkmarkr\checkmarkc\checkmark$\checkmark\\checkmark$\checkmark^\checkmark\\checkmarkc\checkmarki\checkmarkr\checkmarkc\checkmark$\checkmarkc\checkmark$\checkmark^\checkmark\\checkmarkc\checkmarki\checkmarkr\checkmarkc\checkmark$\checkmarkl\checkmark$\checkmark^\checkmark\\checkmarkc\checkmarki\checkmarkr\checkmarkc\checkmark$\checkmarki\checkmark$\checkmark^\checkmark\\checkmarkc\checkmarki\checkmarkr\checkmarkc\checkmark$\checkmarkn\checkmark$\checkmark^\checkmark\\checkmarkc\checkmarki\checkmarkr\checkmarkc\checkmark$\checkmarke\checkmark$\checkmark^\checkmark\\checkmarkc\checkmarki\checkmarkr\checkmarkc\checkmark$\checkmark{\checkmark$\checkmark^\checkmark\\checkmarkc\checkmarki\checkmarkr\checkmarkc\checkmark$\checkmark2\checkmark$\checkmark^\checkmark\\checkmarkc\checkmarki\checkmarkr\checkmarkc\checkmark$\checkmark-\checkmark$\checkmark^\checkmark\\checkmarkc\checkmarki\checkmarkr\checkmarkc\checkmark$\checkmark3\checkmark$\checkmark^\checkmark\\checkmarkc\checkmarki\checkmarkr\checkmarkc\checkmark$\checkmark}\checkmark$\checkmark^\checkmark\\checkmarkc\checkmarki\checkmarkr\checkmarkc\checkmark$\checkmark
\checkmark$\checkmark^\checkmark\\checkmarkc\checkmarki\checkmarkr\checkmarkc\checkmark$\checkmark \checkmark$\checkmark^\checkmark\\checkmarkc\checkmarki\checkmarkr\checkmarkc\checkmark$\checkmark&\checkmark$\checkmark^\checkmark\\checkmarkc\checkmarki\checkmarkr\checkmarkc\checkmark$\checkmark \checkmark$\checkmark^\checkmark\\checkmarkc\checkmarki\checkmarkr\checkmarkc\checkmark$\checkmarkF\checkmark$\checkmark^\checkmark\\checkmarkc\checkmarki\checkmarkr\checkmarkc\checkmark$\checkmarkT\checkmark$\checkmark^\checkmark\\checkmarkc\checkmarki\checkmarkr\checkmarkc\checkmark$\checkmark2\checkmark$\checkmark^\checkmark\\checkmarkc\checkmarki\checkmarkr\checkmarkc\checkmark$\checkmark3\checkmark$\checkmark^\checkmark\\checkmarkc\checkmarki\checkmarkr\checkmarkc\checkmark$\checkmark \checkmark$\checkmark^\checkmark\\checkmarkc\checkmarki\checkmarkr\checkmarkc\checkmark$\checkmark:\checkmark$\checkmark^\checkmark\\checkmarkc\checkmarki\checkmarkr\checkmarkc\checkmark$\checkmark \checkmark$\checkmark^\checkmark\\checkmarkc\checkmarki\checkmarkr\checkmarkc\checkmark$\checkmarkG\checkmark$\checkmark^\checkmark\\checkmarkc\checkmarki\checkmarkr\checkmarkc\checkmark$\checkmarké\checkmark$\checkmark^\checkmark\\checkmarkc\checkmarki\checkmarkr\checkmarkc\checkmark$\checkmarkn\checkmark$\checkmark^\checkmark\\checkmarkc\checkmarki\checkmarkr\checkmarkc\checkmark$\checkmarké\checkmark$\checkmark^\checkmark\\checkmarkc\checkmarki\checkmarkr\checkmarkc\checkmark$\checkmarkr\checkmark$\checkmark^\checkmark\\checkmarkc\checkmarki\checkmarkr\checkmarkc\checkmark$\checkmarke\checkmark$\checkmark^\checkmark\\checkmarkc\checkmarki\checkmarkr\checkmarkc\checkmark$\checkmarkr\checkmark$\checkmark^\checkmark\\checkmarkc\checkmarki\checkmarkr\checkmarkc\checkmark$\checkmark \checkmark$\checkmark^\checkmark\\checkmarkc\checkmarki\checkmarkr\checkmarkc\checkmark$\checkmarkl\checkmark$\checkmark^\checkmark\\checkmarkc\checkmarki\checkmarkr\checkmarkc\checkmark$\checkmarke\checkmark$\checkmark^\checkmark\\checkmarkc\checkmarki\checkmarkr\checkmarkc\checkmark$\checkmark \checkmark$\checkmark^\checkmark\\checkmarkc\checkmarki\checkmarkr\checkmarkc\checkmark$\checkmarkm\checkmark$\checkmark^\checkmark\\checkmarkc\checkmarki\checkmarkr\checkmarkc\checkmark$\checkmarko\checkmark$\checkmark^\checkmark\\checkmarkc\checkmarki\checkmarkr\checkmarkc\checkmark$\checkmarku\checkmark$\checkmark^\checkmark\\checkmarkc\checkmarki\checkmarkr\checkmarkc\checkmark$\checkmarkv\checkmark$\checkmark^\checkmark\\checkmarkc\checkmarki\checkmarkr\checkmarkc\checkmark$\checkmarke\checkmark$\checkmark^\checkmark\\checkmarkc\checkmarki\checkmarkr\checkmarkc\checkmark$\checkmarkm\checkmark$\checkmark^\checkmark\\checkmarkc\checkmarki\checkmarkr\checkmarkc\checkmark$\checkmarke\checkmark$\checkmark^\checkmark\\checkmarkc\checkmarki\checkmarkr\checkmarkc\checkmark$\checkmarkn\checkmark$\checkmark^\checkmark\\checkmarkc\checkmarki\checkmarkr\checkmarkc\checkmark$\checkmarkt\checkmark$\checkmark^\checkmark\\checkmarkc\checkmarki\checkmarkr\checkmarkc\checkmark$\checkmark \checkmark$\checkmark^\checkmark\\checkmarkc\checkmarki\checkmarkr\checkmarkc\checkmark$\checkmark&\checkmark$\checkmark^\checkmark\\checkmarkc\checkmarki\checkmarkr\checkmarkc\checkmark$\checkmark \checkmark$\checkmark^\checkmark\\checkmarkc\checkmarki\checkmarkr\checkmarkc\checkmark$\checkmarkM\checkmark$\checkmark^\checkmark\\checkmarkc\checkmarki\checkmarkr\checkmarkc\checkmark$\checkmarko\checkmark$\checkmark^\checkmark\\checkmarkc\checkmarki\checkmarkr\checkmarkc\checkmark$\checkmarkt\checkmark$\checkmark^\checkmark\\checkmarkc\checkmarki\checkmarkr\checkmarkc\checkmark$\checkmarke\checkmark$\checkmark^\checkmark\\checkmarkc\checkmarki\checkmarkr\checkmarkc\checkmark$\checkmarku\checkmark$\checkmark^\checkmark\\checkmarkc\checkmarki\checkmarkr\checkmarkc\checkmark$\checkmarkr\checkmark$\checkmark^\checkmark\\checkmarkc\checkmarki\checkmarkr\checkmarkc\checkmark$\checkmark \checkmark$\checkmark^\checkmark\\checkmarkc\checkmarki\checkmarkr\checkmarkc\checkmark$\checkmarkp\checkmark$\checkmark^\checkmark\\checkmarkc\checkmarki\checkmarkr\checkmarkc\checkmark$\checkmarka\checkmark$\checkmark^\checkmark\\checkmarkc\checkmarki\checkmarkr\checkmarkc\checkmark$\checkmarks\checkmark$\checkmark^\checkmark\\checkmarkc\checkmarki\checkmarkr\checkmarkc\checkmark$\checkmark \checkmark$\checkmark^\checkmark\\checkmarkc\checkmarki\checkmarkr\checkmarkc\checkmark$\checkmarkà\checkmark$\checkmark^\checkmark\\checkmarkc\checkmarki\checkmarkr\checkmarkc\checkmark$\checkmark \checkmark$\checkmark^\checkmark\\checkmarkc\checkmarki\checkmarkr\checkmarkc\checkmark$\checkmarkp\checkmark$\checkmark^\checkmark\\checkmarkc\checkmarki\checkmarkr\checkmarkc\checkmark$\checkmarka\checkmark$\checkmark^\checkmark\\checkmarkc\checkmarki\checkmarkr\checkmarkc\checkmark$\checkmarks\checkmark$\checkmark^\checkmark\\checkmarkc\checkmarki\checkmarkr\checkmarkc\checkmark$\checkmark \checkmark$\checkmark^\checkmark\\checkmarkc\checkmarki\checkmarkr\checkmarkc\checkmark$\checkmarkN\checkmark$\checkmark^\checkmark\\checkmarkc\checkmarki\checkmarkr\checkmarkc\checkmark$\checkmarkE\checkmark$\checkmark^\checkmark\\checkmarkc\checkmarki\checkmarkr\checkmarkc\checkmark$\checkmarkM\checkmark$\checkmark^\checkmark\\checkmarkc\checkmarki\checkmarkr\checkmarkc\checkmark$\checkmarkA\checkmark$\checkmark^\checkmark\\checkmarkc\checkmarki\checkmarkr\checkmarkc\checkmark$\checkmark \checkmark$\checkmark^\checkmark\\checkmarkc\checkmarki\checkmarkr\checkmarkc\checkmark$\checkmark1\checkmark$\checkmark^\checkmark\\checkmarkc\checkmarki\checkmarkr\checkmarkc\checkmark$\checkmark7\checkmark$\checkmark^\checkmark\\checkmarkc\checkmarki\checkmarkr\checkmarkc\checkmark$\checkmark \checkmark$\checkmark^\checkmark\\checkmarkc\checkmarki\checkmarkr\checkmarkc\checkmark$\checkmark\\checkmark$\checkmark^\checkmark\\checkmarkc\checkmarki\checkmarkr\checkmarkc\checkmark$\checkmark\\checkmark$\checkmark^\checkmark\\checkmarkc\checkmarki\checkmarkr\checkmarkc\checkmark$\checkmark \checkmark$\checkmark^\checkmark\\checkmarkc\checkmarki\checkmarkr\checkmarkc\checkmark$\checkmark\\checkmark$\checkmark^\checkmark\\checkmarkc\checkmarki\checkmarkr\checkmarkc\checkmark$\checkmarkc\checkmark$\checkmark^\checkmark\\checkmarkc\checkmarki\checkmarkr\checkmarkc\checkmark$\checkmarkl\checkmark$\checkmark^\checkmark\\checkmarkc\checkmarki\checkmarkr\checkmarkc\checkmark$\checkmarki\checkmark$\checkmark^\checkmark\\checkmarkc\checkmarki\checkmarkr\checkmarkc\checkmark$\checkmarkn\checkmark$\checkmark^\checkmark\\checkmarkc\checkmarki\checkmarkr\checkmarkc\checkmark$\checkmarke\checkmark$\checkmark^\checkmark\\checkmarkc\checkmarki\checkmarkr\checkmarkc\checkmark$\checkmark{\checkmark$\checkmark^\checkmark\\checkmarkc\checkmarki\checkmarkr\checkmarkc\checkmark$\checkmark2\checkmark$\checkmark^\checkmark\\checkmarkc\checkmarki\checkmarkr\checkmarkc\checkmark$\checkmark-\checkmark$\checkmark^\checkmark\\checkmarkc\checkmarki\checkmarkr\checkmarkc\checkmark$\checkmark3\checkmark$\checkmark^\checkmark\\checkmarkc\checkmarki\checkmarkr\checkmarkc\checkmark$\checkmark}\checkmark$\checkmark^\checkmark\\checkmarkc\checkmarki\checkmarkr\checkmarkc\checkmark$\checkmark
\checkmark$\checkmark^\checkmark\\checkmarkc\checkmarki\checkmarkr\checkmarkc\checkmark$\checkmark \checkmark$\checkmark^\checkmark\\checkmarkc\checkmarki\checkmarkr\checkmarkc\checkmark$\checkmark&\checkmark$\checkmark^\checkmark\\checkmarkc\checkmarki\checkmarkr\checkmarkc\checkmark$\checkmark \checkmark$\checkmark^\checkmark\\checkmarkc\checkmarki\checkmarkr\checkmarkc\checkmark$\checkmarkF\checkmark$\checkmark^\checkmark\\checkmarkc\checkmarki\checkmarkr\checkmarkc\checkmark$\checkmarkT\checkmark$\checkmark^\checkmark\\checkmarkc\checkmarki\checkmarkr\checkmarkc\checkmark$\checkmark2\checkmark$\checkmark^\checkmark\\checkmarkc\checkmarki\checkmarkr\checkmarkc\checkmark$\checkmark4\checkmark$\checkmark^\checkmark\\checkmarkc\checkmarki\checkmarkr\checkmarkc\checkmark$\checkmark \checkmark$\checkmark^\checkmark\\checkmarkc\checkmarki\checkmarkr\checkmarkc\checkmark$\checkmark:\checkmark$\checkmark^\checkmark\\checkmarkc\checkmarki\checkmarkr\checkmarkc\checkmark$\checkmark \checkmark$\checkmark^\checkmark\\checkmarkc\checkmarki\checkmarkr\checkmarkc\checkmark$\checkmarkA\checkmark$\checkmark^\checkmark\\checkmarkc\checkmarki\checkmarkr\checkmarkc\checkmark$\checkmarkl\checkmark$\checkmark^\checkmark\\checkmarkc\checkmarki\checkmarkr\checkmarkc\checkmark$\checkmarki\checkmark$\checkmark^\checkmark\\checkmarkc\checkmarki\checkmarkr\checkmarkc\checkmark$\checkmarkm\checkmark$\checkmark^\checkmark\\checkmarkc\checkmarki\checkmarkr\checkmarkc\checkmark$\checkmarke\checkmark$\checkmark^\checkmark\\checkmarkc\checkmarki\checkmarkr\checkmarkc\checkmark$\checkmarkn\checkmark$\checkmark^\checkmark\\checkmarkc\checkmarki\checkmarkr\checkmarkc\checkmark$\checkmarkt\checkmark$\checkmark^\checkmark\\checkmarkc\checkmarki\checkmarkr\checkmarkc\checkmark$\checkmarke\checkmark$\checkmark^\checkmark\\checkmarkc\checkmarki\checkmarkr\checkmarkc\checkmark$\checkmarkr\checkmark$\checkmark^\checkmark\\checkmarkc\checkmarki\checkmarkr\checkmarkc\checkmark$\checkmark \checkmark$\checkmark^\checkmark\\checkmarkc\checkmarki\checkmarkr\checkmarkc\checkmark$\checkmarke\checkmark$\checkmark^\checkmark\\checkmarkc\checkmarki\checkmarkr\checkmarkc\checkmark$\checkmarkn\checkmark$\checkmark^\checkmark\\checkmarkc\checkmarki\checkmarkr\checkmarkc\checkmark$\checkmark \checkmark$\checkmark^\checkmark\\checkmarkc\checkmarki\checkmarkr\checkmarkc\checkmark$\checkmarké\checkmark$\checkmark^\checkmark\\checkmarkc\checkmarki\checkmarkr\checkmarkc\checkmark$\checkmarkn\checkmark$\checkmark^\checkmark\\checkmarkc\checkmarki\checkmarkr\checkmarkc\checkmark$\checkmarke\checkmark$\checkmark^\checkmark\\checkmarkc\checkmarki\checkmarkr\checkmarkc\checkmark$\checkmarkr\checkmark$\checkmark^\checkmark\\checkmarkc\checkmarki\checkmarkr\checkmarkc\checkmark$\checkmarkg\checkmark$\checkmark^\checkmark\\checkmarkc\checkmarki\checkmarkr\checkmarkc\checkmark$\checkmarki\checkmark$\checkmark^\checkmark\\checkmarkc\checkmarki\checkmarkr\checkmarkc\checkmark$\checkmarke\checkmark$\checkmark^\checkmark\\checkmarkc\checkmarki\checkmarkr\checkmarkc\checkmark$\checkmark \checkmark$\checkmark^\checkmark\\checkmarkc\checkmarki\checkmarkr\checkmarkc\checkmark$\checkmark&\checkmark$\checkmark^\checkmark\\checkmarkc\checkmarki\checkmarkr\checkmarkc\checkmark$\checkmark \checkmark$\checkmark^\checkmark\\checkmarkc\checkmarki\checkmarkr\checkmarkc\checkmark$\checkmarkA\checkmark$\checkmark^\checkmark\\checkmarkc\checkmarki\checkmarkr\checkmarkc\checkmark$\checkmarkl\checkmark$\checkmark^\checkmark\\checkmarkc\checkmarki\checkmarkr\checkmarkc\checkmark$\checkmarki\checkmark$\checkmark^\checkmark\\checkmarkc\checkmarki\checkmarkr\checkmarkc\checkmark$\checkmarkm\checkmark$\checkmark^\checkmark\\checkmarkc\checkmarki\checkmarkr\checkmarkc\checkmark$\checkmarke\checkmark$\checkmark^\checkmark\\checkmarkc\checkmarki\checkmarkr\checkmarkc\checkmark$\checkmarkn\checkmark$\checkmark^\checkmark\\checkmarkc\checkmarki\checkmarkr\checkmarkc\checkmark$\checkmarkt\checkmark$\checkmark^\checkmark\\checkmarkc\checkmarki\checkmarkr\checkmarkc\checkmark$\checkmarka\checkmark$\checkmark^\checkmark\\checkmarkc\checkmarki\checkmarkr\checkmarkc\checkmark$\checkmarkt\checkmark$\checkmark^\checkmark\\checkmarkc\checkmarki\checkmarkr\checkmarkc\checkmark$\checkmarki\checkmark$\checkmark^\checkmark\\checkmarkc\checkmarki\checkmarkr\checkmarkc\checkmark$\checkmarko\checkmark$\checkmark^\checkmark\\checkmarkc\checkmarki\checkmarkr\checkmarkc\checkmark$\checkmarkn\checkmark$\checkmark^\checkmark\\checkmarkc\checkmarki\checkmarkr\checkmarkc\checkmark$\checkmark \checkmark$\checkmark^\checkmark\\checkmarkc\checkmarki\checkmarkr\checkmarkc\checkmark$\checkmark3\checkmark$\checkmark^\checkmark\\checkmarkc\checkmarki\checkmarkr\checkmarkc\checkmark$\checkmark6\checkmark$\checkmark^\checkmark\\checkmarkc\checkmarki\checkmarkr\checkmarkc\checkmark$\checkmarkV\checkmark$\checkmark^\checkmark\\checkmarkc\checkmarki\checkmarkr\checkmarkc\checkmark$\checkmark \checkmark$\checkmark^\checkmark\\checkmarkc\checkmarki\checkmarkr\checkmarkc\checkmark$\checkmark/\checkmark$\checkmark^\checkmark\\checkmarkc\checkmarki\checkmarkr\checkmarkc\checkmark$\checkmark \checkmark$\checkmark^\checkmark\\checkmarkc\checkmarki\checkmarkr\checkmarkc\checkmark$\checkmark3\checkmark$\checkmark^\checkmark\\checkmarkc\checkmarki\checkmarkr\checkmarkc\checkmark$\checkmark5\checkmark$\checkmark^\checkmark\\checkmarkc\checkmarki\checkmarkr\checkmarkc\checkmark$\checkmark0\checkmark$\checkmark^\checkmark\\checkmarkc\checkmarki\checkmarkr\checkmarkc\checkmark$\checkmarkW\checkmark$\checkmark^\checkmark\\checkmarkc\checkmarki\checkmarkr\checkmarkc\checkmark$\checkmark \checkmark$\checkmark^\checkmark\\checkmarkc\checkmarki\checkmarkr\checkmarkc\checkmark$\checkmark\\checkmark$\checkmark^\checkmark\\checkmarkc\checkmarki\checkmarkr\checkmarkc\checkmark$\checkmark\\checkmark$\checkmark^\checkmark\\checkmarkc\checkmarki\checkmarkr\checkmarkc\checkmark$\checkmark \checkmark$\checkmark^\checkmark\\checkmarkc\checkmarki\checkmarkr\checkmarkc\checkmark$\checkmark\\checkmark$\checkmark^\checkmark\\checkmarkc\checkmarki\checkmarkr\checkmarkc\checkmark$\checkmarkh\checkmark$\checkmark^\checkmark\\checkmarkc\checkmarki\checkmarkr\checkmarkc\checkmark$\checkmarkl\checkmark$\checkmark^\checkmark\\checkmarkc\checkmarki\checkmarkr\checkmarkc\checkmark$\checkmarki\checkmark$\checkmark^\checkmark\\checkmarkc\checkmarki\checkmarkr\checkmarkc\checkmark$\checkmarkn\checkmark$\checkmark^\checkmark\\checkmarkc\checkmarki\checkmarkr\checkmarkc\checkmark$\checkmarke\checkmark$\checkmark^\checkmark\\checkmarkc\checkmarki\checkmarkr\checkmarkc\checkmark$\checkmark
\checkmark$\checkmark^\checkmark\\checkmarkc\checkmarki\checkmarkr\checkmarkc\checkmark$\checkmark\\checkmark$\checkmark^\checkmark\\checkmarkc\checkmarki\checkmarkr\checkmarkc\checkmark$\checkmarkm\checkmark$\checkmark^\checkmark\\checkmarkc\checkmarki\checkmarkr\checkmarkc\checkmark$\checkmarku\checkmark$\checkmark^\checkmark\\checkmarkc\checkmarki\checkmarkr\checkmarkc\checkmark$\checkmarkl\checkmark$\checkmark^\checkmark\\checkmarkc\checkmarki\checkmarkr\checkmarkc\checkmark$\checkmarkt\checkmark$\checkmark^\checkmark\\checkmarkc\checkmarki\checkmarkr\checkmarkc\checkmark$\checkmarki\checkmark$\checkmark^\checkmark\\checkmarkc\checkmarki\checkmarkr\checkmarkc\checkmark$\checkmarkr\checkmark$\checkmark^\checkmark\\checkmarkc\checkmarki\checkmarkr\checkmarkc\checkmark$\checkmarko\checkmark$\checkmark^\checkmark\\checkmarkc\checkmarki\checkmarkr\checkmarkc\checkmark$\checkmarkw\checkmark$\checkmark^\checkmark\\checkmarkc\checkmarki\checkmarkr\checkmarkc\checkmark$\checkmark{\checkmark$\checkmark^\checkmark\\checkmarkc\checkmarki\checkmarkr\checkmarkc\checkmark$\checkmark4\checkmark$\checkmark^\checkmark\\checkmarkc\checkmarki\checkmarkr\checkmarkc\checkmark$\checkmark}\checkmark$\checkmark^\checkmark\\checkmarkc\checkmarki\checkmarkr\checkmarkc\checkmark$\checkmark{\checkmark$\checkmark^\checkmark\\checkmarkc\checkmarki\checkmarkr\checkmarkc\checkmark$\checkmark*\checkmark$\checkmark^\checkmark\\checkmarkc\checkmarki\checkmarkr\checkmarkc\checkmark$\checkmark}\checkmark$\checkmark^\checkmark\\checkmarkc\checkmarki\checkmarkr\checkmarkc\checkmark$\checkmark{\checkmark$\checkmark^\checkmark\\checkmarkc\checkmarki\checkmarkr\checkmarkc\checkmark$\checkmarkF\checkmark$\checkmark^\checkmark\\checkmarkc\checkmarki\checkmarkr\checkmarkc\checkmark$\checkmarkT\checkmark$\checkmark^\checkmark\\checkmarkc\checkmarki\checkmarkr\checkmarkc\checkmark$\checkmark3\checkmark$\checkmark^\checkmark\\checkmarkc\checkmarki\checkmarkr\checkmarkc\checkmark$\checkmark \checkmark$\checkmark^\checkmark\\checkmarkc\checkmarki\checkmarkr\checkmarkc\checkmark$\checkmark:\checkmark$\checkmark^\checkmark\\checkmarkc\checkmarki\checkmarkr\checkmarkc\checkmark$\checkmark \checkmark$\checkmark^\checkmark\\checkmarkc\checkmarki\checkmarkr\checkmarkc\checkmark$\checkmarkB\checkmark$\checkmark^\checkmark\\checkmarkc\checkmarki\checkmarkr\checkmarkc\checkmark$\checkmarka\checkmark$\checkmark^\checkmark\\checkmarkc\checkmarki\checkmarkr\checkmarkc\checkmark$\checkmarks\checkmark$\checkmark^\checkmark\\checkmarkc\checkmarki\checkmarkr\checkmarkc\checkmark$\checkmarkc\checkmark$\checkmark^\checkmark\\checkmarkc\checkmarki\checkmarkr\checkmarkc\checkmark$\checkmarku\checkmark$\checkmark^\checkmark\\checkmarkc\checkmarki\checkmarkr\checkmarkc\checkmark$\checkmarkl\checkmark$\checkmark^\checkmark\\checkmarkc\checkmarki\checkmarkr\checkmarkc\checkmark$\checkmarke\checkmark$\checkmark^\checkmark\\checkmarkc\checkmarki\checkmarkr\checkmarkc\checkmark$\checkmarkr\checkmark$\checkmark^\checkmark\\checkmarkc\checkmarki\checkmarkr\checkmarkc\checkmark$\checkmark \checkmark$\checkmark^\checkmark\\checkmarkc\checkmarki\checkmarkr\checkmarkc\checkmark$\checkmarke\checkmark$\checkmark^\checkmark\\checkmarkc\checkmarki\checkmarkr\checkmarkc\checkmark$\checkmarkt\checkmark$\checkmark^\checkmark\\checkmarkc\checkmarki\checkmarkr\checkmarkc\checkmark$\checkmark \checkmark$\checkmark^\checkmark\\checkmarkc\checkmarki\checkmarkr\checkmarkc\checkmark$\checkmarkp\checkmark$\checkmark^\checkmark\\checkmarkc\checkmarki\checkmarkr\checkmarkc\checkmark$\checkmarki\checkmark$\checkmark^\checkmark\\checkmarkc\checkmarki\checkmarkr\checkmarkc\checkmark$\checkmarkv\checkmark$\checkmark^\checkmark\\checkmarkc\checkmarki\checkmarkr\checkmarkc\checkmark$\checkmarko\checkmark$\checkmark^\checkmark\\checkmarkc\checkmarki\checkmarkr\checkmarkc\checkmark$\checkmarkt\checkmark$\checkmark^\checkmark\\checkmarkc\checkmarki\checkmarkr\checkmarkc\checkmark$\checkmarke\checkmark$\checkmark^\checkmark\\checkmarkc\checkmarki\checkmarkr\checkmarkc\checkmark$\checkmarkr\checkmark$\checkmark^\checkmark\\checkmarkc\checkmarki\checkmarkr\checkmarkc\checkmark$\checkmark \checkmark$\checkmark^\checkmark\\checkmarkc\checkmarki\checkmarkr\checkmarkc\checkmark$\checkmarkA\checkmark$\checkmark^\checkmark\\checkmarkc\checkmarki\checkmarkr\checkmarkc\checkmark$\checkmark \checkmark$\checkmark^\checkmark\\checkmarkc\checkmarki\checkmarkr\checkmarkc\checkmark$\checkmarke\checkmark$\checkmark^\checkmark\\checkmarkc\checkmarki\checkmarkr\checkmarkc\checkmark$\checkmarkt\checkmark$\checkmark^\checkmark\\checkmarkc\checkmarki\checkmarkr\checkmarkc\checkmark$\checkmark \checkmark$\checkmark^\checkmark\\checkmarkc\checkmarki\checkmarkr\checkmarkc\checkmark$\checkmarkB\checkmark$\checkmark^\checkmark\\checkmarkc\checkmarki\checkmarkr\checkmarkc\checkmark$\checkmark \checkmark$\checkmark^\checkmark\\checkmarkc\checkmarki\checkmarkr\checkmarkc\checkmark$\checkmark(\checkmark$\checkmark^\checkmark\\checkmarkc\checkmarki\checkmarkr\checkmarkc\checkmark$\checkmarkR\checkmark$\checkmark^\checkmark\\checkmarkc\checkmarki\checkmarkr\checkmarkc\checkmark$\checkmarko\checkmark$\checkmark^\checkmark\\checkmarkc\checkmarki\checkmarkr\checkmarkc\checkmark$\checkmarkt\checkmark$\checkmark^\checkmark\\checkmarkc\checkmarki\checkmarkr\checkmarkc\checkmark$\checkmarka\checkmark$\checkmark^\checkmark\\checkmarkc\checkmarki\checkmarkr\checkmarkc\checkmark$\checkmarkt\checkmark$\checkmark^\checkmark\\checkmarkc\checkmarki\checkmarkr\checkmarkc\checkmark$\checkmarki\checkmark$\checkmark^\checkmark\\checkmarkc\checkmarki\checkmarkr\checkmarkc\checkmark$\checkmarko\checkmark$\checkmark^\checkmark\\checkmarkc\checkmarki\checkmarkr\checkmarkc\checkmark$\checkmarkn\checkmark$\checkmark^\checkmark\\checkmarkc\checkmarki\checkmarkr\checkmarkc\checkmark$\checkmark)\checkmark$\checkmark^\checkmark\\checkmarkc\checkmarki\checkmarkr\checkmarkc\checkmark$\checkmark}\checkmark$\checkmark^\checkmark\\checkmarkc\checkmarki\checkmarkr\checkmarkc\checkmark$\checkmark \checkmark$\checkmark^\checkmark\\checkmarkc\checkmarki\checkmarkr\checkmarkc\checkmark$\checkmark&\checkmark$\checkmark^\checkmark\\checkmarkc\checkmarki\checkmarkr\checkmarkc\checkmark$\checkmark \checkmark$\checkmark^\checkmark\\checkmarkc\checkmarki\checkmarkr\checkmarkc\checkmark$\checkmarkF\checkmark$\checkmark^\checkmark\\checkmarkc\checkmarki\checkmarkr\checkmarkc\checkmark$\checkmarkT\checkmark$\checkmark^\checkmark\\checkmarkc\checkmarki\checkmarkr\checkmarkc\checkmark$\checkmark3\checkmark$\checkmark^\checkmark\\checkmarkc\checkmarki\checkmarkr\checkmarkc\checkmark$\checkmark1\checkmark$\checkmark^\checkmark\\checkmarkc\checkmarki\checkmarkr\checkmarkc\checkmark$\checkmark \checkmark$\checkmark^\checkmark\\checkmarkc\checkmarki\checkmarkr\checkmarkc\checkmark$\checkmark:\checkmark$\checkmark^\checkmark\\checkmarkc\checkmarki\checkmarkr\checkmarkc\checkmark$\checkmark \checkmark$\checkmark^\checkmark\\checkmarkc\checkmarki\checkmarkr\checkmarkc\checkmark$\checkmarkG\checkmark$\checkmark^\checkmark\\checkmarkc\checkmarki\checkmarkr\checkmarkc\checkmark$\checkmarku\checkmark$\checkmark^\checkmark\\checkmarkc\checkmarki\checkmarkr\checkmarkc\checkmark$\checkmarki\checkmark$\checkmark^\checkmark\\checkmarkc\checkmarki\checkmarkr\checkmarkc\checkmark$\checkmarkd\checkmark$\checkmark^\checkmark\\checkmarkc\checkmarki\checkmarkr\checkmarkc\checkmark$\checkmarke\checkmark$\checkmark^\checkmark\\checkmarkc\checkmarki\checkmarkr\checkmarkc\checkmark$\checkmarkr\checkmark$\checkmark^\checkmark\\checkmarkc\checkmarki\checkmarkr\checkmarkc\checkmark$\checkmark \checkmark$\checkmark^\checkmark\\checkmarkc\checkmarki\checkmarkr\checkmarkc\checkmark$\checkmarke\checkmark$\checkmark^\checkmark\\checkmarkc\checkmarki\checkmarkr\checkmarkc\checkmark$\checkmarkn\checkmark$\checkmark^\checkmark\\checkmarkc\checkmarki\checkmarkr\checkmarkc\checkmark$\checkmark \checkmark$\checkmark^\checkmark\\checkmarkc\checkmarki\checkmarkr\checkmarkc\checkmark$\checkmarkr\checkmark$\checkmark^\checkmark\\checkmarkc\checkmarki\checkmarkr\checkmarkc\checkmark$\checkmarko\checkmark$\checkmark^\checkmark\\checkmarkc\checkmarki\checkmarkr\checkmarkc\checkmark$\checkmarkt\checkmark$\checkmark^\checkmark\\checkmarkc\checkmarki\checkmarkr\checkmarkc\checkmark$\checkmarka\checkmark$\checkmark^\checkmark\\checkmarkc\checkmarki\checkmarkr\checkmarkc\checkmark$\checkmarkt\checkmark$\checkmark^\checkmark\\checkmarkc\checkmarki\checkmarkr\checkmarkc\checkmark$\checkmarki\checkmark$\checkmark^\checkmark\\checkmarkc\checkmarki\checkmarkr\checkmarkc\checkmark$\checkmarko\checkmark$\checkmark^\checkmark\\checkmarkc\checkmarki\checkmarkr\checkmarkc\checkmark$\checkmarkn\checkmark$\checkmark^\checkmark\\checkmarkc\checkmarki\checkmarkr\checkmarkc\checkmark$\checkmark \checkmark$\checkmark^\checkmark\\checkmarkc\checkmarki\checkmarkr\checkmarkc\checkmark$\checkmarkà\checkmark$\checkmark^\checkmark\\checkmarkc\checkmarki\checkmarkr\checkmarkc\checkmark$\checkmark \checkmark$\checkmark^\checkmark\\checkmarkc\checkmarki\checkmarkr\checkmarkc\checkmark$\checkmarkf\checkmark$\checkmark^\checkmark\\checkmarkc\checkmarki\checkmarkr\checkmarkc\checkmark$\checkmarko\checkmark$\checkmark^\checkmark\\checkmarkc\checkmarki\checkmarkr\checkmarkc\checkmark$\checkmarkr\checkmark$\checkmark^\checkmark\\checkmarkc\checkmarki\checkmarkr\checkmarkc\checkmark$\checkmarkt\checkmark$\checkmark^\checkmark\\checkmarkc\checkmarki\checkmarkr\checkmarkc\checkmark$\checkmarke\checkmark$\checkmark^\checkmark\\checkmarkc\checkmarki\checkmarkr\checkmarkc\checkmark$\checkmark \checkmark$\checkmark^\checkmark\\checkmarkc\checkmarki\checkmarkr\checkmarkc\checkmark$\checkmarkc\checkmark$\checkmark^\checkmark\\checkmarkc\checkmarki\checkmarkr\checkmarkc\checkmark$\checkmarkh\checkmark$\checkmark^\checkmark\\checkmarkc\checkmarki\checkmarkr\checkmarkc\checkmark$\checkmarka\checkmark$\checkmark^\checkmark\\checkmarkc\checkmarki\checkmarkr\checkmarkc\checkmark$\checkmarkr\checkmark$\checkmark^\checkmark\\checkmarkc\checkmarki\checkmarkr\checkmarkc\checkmark$\checkmarkg\checkmark$\checkmark^\checkmark\\checkmarkc\checkmarki\checkmarkr\checkmarkc\checkmark$\checkmarke\checkmark$\checkmark^\checkmark\\checkmarkc\checkmarki\checkmarkr\checkmarkc\checkmark$\checkmark \checkmark$\checkmark^\checkmark\\checkmarkc\checkmarki\checkmarkr\checkmarkc\checkmark$\checkmark&\checkmark$\checkmark^\checkmark\\checkmarkc\checkmarki\checkmarkr\checkmarkc\checkmark$\checkmark \checkmark$\checkmark^\checkmark\\checkmarkc\checkmarki\checkmarkr\checkmarkc\checkmark$\checkmarkR\checkmark$\checkmark^\checkmark\\checkmarkc\checkmarki\checkmarkr\checkmarkc\checkmark$\checkmarko\checkmark$\checkmark^\checkmark\\checkmarkc\checkmarki\checkmarkr\checkmarkc\checkmark$\checkmarku\checkmark$\checkmark^\checkmark\\checkmarkc\checkmarki\checkmarkr\checkmarkc\checkmark$\checkmarkl\checkmark$\checkmark^\checkmark\\checkmarkc\checkmarki\checkmarkr\checkmarkc\checkmark$\checkmarke\checkmark$\checkmark^\checkmark\\checkmarkc\checkmarki\checkmarkr\checkmarkc\checkmark$\checkmarkm\checkmark$\checkmark^\checkmark\\checkmarkc\checkmarki\checkmarkr\checkmarkc\checkmark$\checkmarke\checkmark$\checkmark^\checkmark\\checkmarkc\checkmarki\checkmarkr\checkmarkc\checkmark$\checkmarkn\checkmark$\checkmark^\checkmark\\checkmarkc\checkmarki\checkmarkr\checkmarkc\checkmark$\checkmarkt\checkmark$\checkmark^\checkmark\\checkmarkc\checkmarki\checkmarkr\checkmarkc\checkmark$\checkmarks\checkmark$\checkmark^\checkmark\\checkmarkc\checkmarki\checkmarkr\checkmarkc\checkmark$\checkmark \checkmark$\checkmark^\checkmark\\checkmarkc\checkmarki\checkmarkr\checkmarkc\checkmark$\checkmarkc\checkmark$\checkmark^\checkmark\\checkmarkc\checkmarki\checkmarkr\checkmarkc\checkmark$\checkmarkr\checkmark$\checkmark^\checkmark\\checkmarkc\checkmarki\checkmarkr\checkmarkc\checkmark$\checkmarko\checkmark$\checkmark^\checkmark\\checkmarkc\checkmarki\checkmarkr\checkmarkc\checkmark$\checkmarki\checkmark$\checkmark^\checkmark\\checkmarkc\checkmarki\checkmarkr\checkmarkc\checkmark$\checkmarks\checkmark$\checkmark^\checkmark\\checkmarkc\checkmarki\checkmarkr\checkmarkc\checkmark$\checkmarké\checkmark$\checkmark^\checkmark\\checkmarkc\checkmarki\checkmarkr\checkmarkc\checkmark$\checkmarks\checkmark$\checkmark^\checkmark\\checkmarkc\checkmarki\checkmarkr\checkmarkc\checkmark$\checkmark \checkmark$\checkmark^\checkmark\\checkmarkc\checkmarki\checkmarkr\checkmarkc\checkmark$\checkmark(\checkmark$\checkmark^\checkmark\\checkmarkc\checkmarki\checkmarkr\checkmarkc\checkmark$\checkmarkA\checkmark$\checkmark^\checkmark\\checkmarkc\checkmarki\checkmarkr\checkmarkc\checkmark$\checkmarkx\checkmark$\checkmark^\checkmark\\checkmarkc\checkmarki\checkmarkr\checkmarkc\checkmark$\checkmarke\checkmark$\checkmark^\checkmark\\checkmarkc\checkmarki\checkmarkr\checkmarkc\checkmark$\checkmark \checkmark$\checkmark^\checkmark\\checkmarkc\checkmarki\checkmarkr\checkmarkc\checkmark$\checkmarkB\checkmark$\checkmark^\checkmark\\checkmarkc\checkmarki\checkmarkr\checkmarkc\checkmark$\checkmark)\checkmark$\checkmark^\checkmark\\checkmarkc\checkmarki\checkmarkr\checkmarkc\checkmark$\checkmark \checkmark$\checkmark^\checkmark\\checkmarkc\checkmarki\checkmarkr\checkmarkc\checkmark$\checkmark\\checkmark$\checkmark^\checkmark\\checkmarkc\checkmarki\checkmarkr\checkmarkc\checkmark$\checkmark\\checkmark$\checkmark^\checkmark\\checkmarkc\checkmarki\checkmarkr\checkmarkc\checkmark$\checkmark \checkmark$\checkmark^\checkmark\\checkmarkc\checkmarki\checkmarkr\checkmarkc\checkmark$\checkmark\\checkmark$\checkmark^\checkmark\\checkmarkc\checkmarki\checkmarkr\checkmarkc\checkmark$\checkmarkc\checkmark$\checkmark^\checkmark\\checkmarkc\checkmarki\checkmarkr\checkmarkc\checkmark$\checkmarkl\checkmark$\checkmark^\checkmark\\checkmarkc\checkmarki\checkmarkr\checkmarkc\checkmark$\checkmarki\checkmark$\checkmark^\checkmark\\checkmarkc\checkmarki\checkmarkr\checkmarkc\checkmark$\checkmarkn\checkmark$\checkmark^\checkmark\\checkmarkc\checkmarki\checkmarkr\checkmarkc\checkmark$\checkmarke\checkmark$\checkmark^\checkmark\\checkmarkc\checkmarki\checkmarkr\checkmarkc\checkmark$\checkmark{\checkmark$\checkmark^\checkmark\\checkmarkc\checkmarki\checkmarkr\checkmarkc\checkmark$\checkmark2\checkmark$\checkmark^\checkmark\\checkmarkc\checkmarki\checkmarkr\checkmarkc\checkmark$\checkmark-\checkmark$\checkmark^\checkmark\\checkmarkc\checkmarki\checkmarkr\checkmarkc\checkmark$\checkmark3\checkmark$\checkmark^\checkmark\\checkmarkc\checkmarki\checkmarkr\checkmarkc\checkmark$\checkmark}\checkmark$\checkmark^\checkmark\\checkmarkc\checkmarki\checkmarkr\checkmarkc\checkmark$\checkmark
\checkmark$\checkmark^\checkmark\\checkmarkc\checkmarki\checkmarkr\checkmarkc\checkmark$\checkmark \checkmark$\checkmark^\checkmark\\checkmarkc\checkmarki\checkmarkr\checkmarkc\checkmark$\checkmark&\checkmark$\checkmark^\checkmark\\checkmarkc\checkmarki\checkmarkr\checkmarkc\checkmark$\checkmark \checkmark$\checkmark^\checkmark\\checkmarkc\checkmarki\checkmarkr\checkmarkc\checkmark$\checkmarkF\checkmark$\checkmark^\checkmark\\checkmarkc\checkmarki\checkmarkr\checkmarkc\checkmark$\checkmarkT\checkmark$\checkmark^\checkmark\\checkmarkc\checkmarki\checkmarkr\checkmarkc\checkmark$\checkmark3\checkmark$\checkmark^\checkmark\\checkmarkc\checkmarki\checkmarkr\checkmarkc\checkmark$\checkmark2\checkmark$\checkmark^\checkmark\\checkmarkc\checkmarki\checkmarkr\checkmarkc\checkmark$\checkmark \checkmark$\checkmark^\checkmark\\checkmarkc\checkmarki\checkmarkr\checkmarkc\checkmark$\checkmark:\checkmark$\checkmark^\checkmark\\checkmarkc\checkmarki\checkmarkr\checkmarkc\checkmark$\checkmark \checkmark$\checkmark^\checkmark\\checkmarkc\checkmarki\checkmarkr\checkmarkc\checkmark$\checkmarkD\checkmark$\checkmark^\checkmark\\checkmarkc\checkmarki\checkmarkr\checkmarkc\checkmark$\checkmarké\checkmark$\checkmark^\checkmark\\checkmarkc\checkmarki\checkmarkr\checkmarkc\checkmark$\checkmarkm\checkmark$\checkmark^\checkmark\\checkmarkc\checkmarki\checkmarkr\checkmarkc\checkmark$\checkmarku\checkmark$\checkmark^\checkmark\\checkmarkc\checkmarki\checkmarkr\checkmarkc\checkmark$\checkmarkl\checkmark$\checkmark^\checkmark\\checkmarkc\checkmarki\checkmarkr\checkmarkc\checkmark$\checkmarkt\checkmark$\checkmark^\checkmark\\checkmarkc\checkmarki\checkmarkr\checkmarkc\checkmark$\checkmarki\checkmark$\checkmark^\checkmark\\checkmarkc\checkmarki\checkmarkr\checkmarkc\checkmark$\checkmarkp\checkmark$\checkmark^\checkmark\\checkmarkc\checkmarki\checkmarkr\checkmarkc\checkmark$\checkmarkl\checkmark$\checkmark^\checkmark\\checkmarkc\checkmarki\checkmarkr\checkmarkc\checkmark$\checkmarki\checkmark$\checkmark^\checkmark\\checkmarkc\checkmarki\checkmarkr\checkmarkc\checkmark$\checkmarke\checkmark$\checkmark^\checkmark\\checkmarkc\checkmarki\checkmarkr\checkmarkc\checkmark$\checkmarkr\checkmark$\checkmark^\checkmark\\checkmarkc\checkmarki\checkmarkr\checkmarkc\checkmark$\checkmark \checkmark$\checkmark^\checkmark\\checkmarkc\checkmarki\checkmarkr\checkmarkc\checkmark$\checkmarkl\checkmark$\checkmark^\checkmark\\checkmarkc\checkmarki\checkmarkr\checkmarkc\checkmark$\checkmarke\checkmark$\checkmark^\checkmark\\checkmarkc\checkmarki\checkmarkr\checkmarkc\checkmark$\checkmark \checkmark$\checkmark^\checkmark\\checkmarkc\checkmarki\checkmarkr\checkmarkc\checkmark$\checkmarkc\checkmark$\checkmark^\checkmark\\checkmarkc\checkmarki\checkmarkr\checkmarkc\checkmark$\checkmarko\checkmark$\checkmark^\checkmark\\checkmarkc\checkmarki\checkmarkr\checkmarkc\checkmark$\checkmarku\checkmark$\checkmark^\checkmark\\checkmarkc\checkmarki\checkmarkr\checkmarkc\checkmark$\checkmarkp\checkmark$\checkmark^\checkmark\\checkmarkc\checkmarki\checkmarkr\checkmarkc\checkmark$\checkmarkl\checkmark$\checkmark^\checkmark\\checkmarkc\checkmarki\checkmarkr\checkmarkc\checkmark$\checkmarke\checkmark$\checkmark^\checkmark\\checkmarkc\checkmarki\checkmarkr\checkmarkc\checkmark$\checkmark \checkmark$\checkmark^\checkmark\\checkmarkc\checkmarki\checkmarkr\checkmarkc\checkmark$\checkmark(\checkmark$\checkmark^\checkmark\\checkmarkc\checkmarki\checkmarkr\checkmarkc\checkmark$\checkmarkA\checkmark$\checkmark^\checkmark\\checkmarkc\checkmarki\checkmarkr\checkmarkc\checkmark$\checkmarkn\checkmark$\checkmark^\checkmark\\checkmarkc\checkmarki\checkmarkr\checkmarkc\checkmark$\checkmarkt\checkmark$\checkmark^\checkmark\\checkmarkc\checkmarki\checkmarkr\checkmarkc\checkmark$\checkmarki\checkmark$\checkmark^\checkmark\\checkmarkc\checkmarki\checkmarkr\checkmarkc\checkmark$\checkmark-\checkmark$\checkmark^\checkmark\\checkmarkc\checkmarki\checkmarkr\checkmarkc\checkmark$\checkmarkB\checkmark$\checkmark^\checkmark\\checkmarkc\checkmarki\checkmarkr\checkmarkc\checkmark$\checkmarka\checkmark$\checkmark^\checkmark\\checkmarkc\checkmarki\checkmarkr\checkmarkc\checkmark$\checkmarkc\checkmark$\checkmark^\checkmark\\checkmarkc\checkmarki\checkmarkr\checkmarkc\checkmark$\checkmarkk\checkmark$\checkmark^\checkmark\\checkmarkc\checkmarki\checkmarkr\checkmarkc\checkmark$\checkmarkl\checkmark$\checkmark^\checkmark\\checkmarkc\checkmarki\checkmarkr\checkmarkc\checkmark$\checkmarka\checkmark$\checkmark^\checkmark\\checkmarkc\checkmarki\checkmarkr\checkmarkc\checkmark$\checkmarks\checkmark$\checkmark^\checkmark\\checkmarkc\checkmarki\checkmarkr\checkmarkc\checkmark$\checkmarkh\checkmark$\checkmark^\checkmark\\checkmarkc\checkmarki\checkmarkr\checkmarkc\checkmark$\checkmark)\checkmark$\checkmark^\checkmark\\checkmarkc\checkmarki\checkmarkr\checkmarkc\checkmark$\checkmark \checkmark$\checkmark^\checkmark\\checkmarkc\checkmarki\checkmarkr\checkmarkc\checkmark$\checkmark&\checkmark$\checkmark^\checkmark\\checkmarkc\checkmarki\checkmarkr\checkmarkc\checkmark$\checkmark \checkmark$\checkmark^\checkmark\\checkmarkc\checkmarki\checkmarkr\checkmarkc\checkmark$\checkmarkC\checkmark$\checkmark^\checkmark\\checkmarkc\checkmarki\checkmarkr\checkmarkc\checkmark$\checkmarko\checkmark$\checkmark^\checkmark\\checkmarkc\checkmarki\checkmarkr\checkmarkc\checkmark$\checkmarku\checkmark$\checkmark^\checkmark\\checkmarkc\checkmarki\checkmarkr\checkmarkc\checkmark$\checkmarkr\checkmark$\checkmark^\checkmark\\checkmarkc\checkmarki\checkmarkr\checkmarkc\checkmark$\checkmarkr\checkmark$\checkmark^\checkmark\\checkmarkc\checkmarki\checkmarkr\checkmarkc\checkmark$\checkmarko\checkmark$\checkmark^\checkmark\\checkmarkc\checkmarki\checkmarkr\checkmarkc\checkmark$\checkmarki\checkmark$\checkmark^\checkmark\\checkmarkc\checkmarki\checkmarkr\checkmarkc\checkmark$\checkmarke\checkmark$\checkmark^\checkmark\\checkmarkc\checkmarki\checkmarkr\checkmarkc\checkmark$\checkmark \checkmark$\checkmark^\checkmark\\checkmarkc\checkmarki\checkmarkr\checkmarkc\checkmark$\checkmarkc\checkmark$\checkmark^\checkmark\\checkmarkc\checkmarki\checkmarkr\checkmarkc\checkmark$\checkmarkr\checkmark$\checkmark^\checkmark\\checkmarkc\checkmarki\checkmarkr\checkmarkc\checkmark$\checkmarka\checkmark$\checkmark^\checkmark\\checkmarkc\checkmarki\checkmarkr\checkmarkc\checkmark$\checkmarkn\checkmark$\checkmark^\checkmark\\checkmarkc\checkmarki\checkmarkr\checkmarkc\checkmark$\checkmarkt\checkmark$\checkmark^\checkmark\\checkmarkc\checkmarki\checkmarkr\checkmarkc\checkmark$\checkmarké\checkmark$\checkmark^\checkmark\\checkmarkc\checkmarki\checkmarkr\checkmarkc\checkmark$\checkmarke\checkmark$\checkmark^\checkmark\\checkmarkc\checkmarki\checkmarkr\checkmarkc\checkmark$\checkmark \checkmark$\checkmark^\checkmark\\checkmarkc\checkmarki\checkmarkr\checkmarkc\checkmark$\checkmarkH\checkmark$\checkmark^\checkmark\\checkmarkc\checkmarki\checkmarkr\checkmarkc\checkmark$\checkmarkT\checkmark$\checkmark^\checkmark\\checkmarkc\checkmarki\checkmarkr\checkmarkc\checkmark$\checkmarkD\checkmark$\checkmark^\checkmark\\checkmarkc\checkmarki\checkmarkr\checkmarkc\checkmark$\checkmark \checkmark$\checkmark^\checkmark\\checkmarkc\checkmarki\checkmarkr\checkmarkc\checkmark$\checkmark5\checkmark$\checkmark^\checkmark\\checkmarkc\checkmarki\checkmarkr\checkmarkc\checkmark$\checkmarkM\checkmark$\checkmark^\checkmark\\checkmarkc\checkmarki\checkmarkr\checkmarkc\checkmark$\checkmark \checkmark$\checkmark^\checkmark\\checkmarkc\checkmarki\checkmarkr\checkmarkc\checkmark$\checkmark(\checkmark$\checkmark^\checkmark\\checkmarkc\checkmarki\checkmarkr\checkmarkc\checkmark$\checkmark1\checkmark$\checkmark^\checkmark\\checkmarkc\checkmarki\checkmarkr\checkmarkc\checkmark$\checkmark:\checkmark$\checkmark^\checkmark\\checkmarkc\checkmarki\checkmarkr\checkmarkc\checkmark$\checkmark5\checkmark$\checkmark^\checkmark\\checkmarkc\checkmarki\checkmarkr\checkmarkc\checkmark$\checkmark0\checkmark$\checkmark^\checkmark\\checkmarkc\checkmarki\checkmarkr\checkmarkc\checkmark$\checkmark)\checkmark$\checkmark^\checkmark\\checkmarkc\checkmarki\checkmarkr\checkmarkc\checkmark$\checkmark \checkmark$\checkmark^\checkmark\\checkmarkc\checkmarki\checkmarkr\checkmarkc\checkmark$\checkmark\\checkmark$\checkmark^\checkmark\\checkmarkc\checkmarki\checkmarkr\checkmarkc\checkmark$\checkmark\\checkmark$\checkmark^\checkmark\\checkmarkc\checkmarki\checkmarkr\checkmarkc\checkmark$\checkmark \checkmark$\checkmark^\checkmark\\checkmarkc\checkmarki\checkmarkr\checkmarkc\checkmark$\checkmark\\checkmark$\checkmark^\checkmark\\checkmarkc\checkmarki\checkmarkr\checkmarkc\checkmark$\checkmarkc\checkmark$\checkmark^\checkmark\\checkmarkc\checkmarki\checkmarkr\checkmarkc\checkmark$\checkmarkl\checkmark$\checkmark^\checkmark\\checkmarkc\checkmarki\checkmarkr\checkmarkc\checkmark$\checkmarki\checkmark$\checkmark^\checkmark\\checkmarkc\checkmarki\checkmarkr\checkmarkc\checkmark$\checkmarkn\checkmark$\checkmark^\checkmark\\checkmarkc\checkmarki\checkmarkr\checkmarkc\checkmark$\checkmarke\checkmark$\checkmark^\checkmark\\checkmarkc\checkmarki\checkmarkr\checkmarkc\checkmark$\checkmark{\checkmark$\checkmark^\checkmark\\checkmarkc\checkmarki\checkmarkr\checkmarkc\checkmark$\checkmark2\checkmark$\checkmark^\checkmark\\checkmarkc\checkmarki\checkmarkr\checkmarkc\checkmark$\checkmark-\checkmark$\checkmark^\checkmark\\checkmarkc\checkmarki\checkmarkr\checkmarkc\checkmark$\checkmark3\checkmark$\checkmark^\checkmark\\checkmarkc\checkmarki\checkmarkr\checkmarkc\checkmark$\checkmark}\checkmark$\checkmark^\checkmark\\checkmarkc\checkmarki\checkmarkr\checkmarkc\checkmark$\checkmark
\checkmark$\checkmark^\checkmark\\checkmarkc\checkmarki\checkmarkr\checkmarkc\checkmark$\checkmark \checkmark$\checkmark^\checkmark\\checkmarkc\checkmarki\checkmarkr\checkmarkc\checkmark$\checkmark&\checkmark$\checkmark^\checkmark\\checkmarkc\checkmarki\checkmarkr\checkmarkc\checkmark$\checkmark \checkmark$\checkmark^\checkmark\\checkmarkc\checkmarki\checkmarkr\checkmarkc\checkmark$\checkmarkF\checkmark$\checkmark^\checkmark\\checkmarkc\checkmarki\checkmarkr\checkmarkc\checkmark$\checkmarkT\checkmark$\checkmark^\checkmark\\checkmarkc\checkmarki\checkmarkr\checkmarkc\checkmark$\checkmark3\checkmark$\checkmark^\checkmark\\checkmarkc\checkmarki\checkmarkr\checkmarkc\checkmark$\checkmark3\checkmark$\checkmark^\checkmark\\checkmarkc\checkmarki\checkmarkr\checkmarkc\checkmark$\checkmark \checkmark$\checkmark^\checkmark\\checkmarkc\checkmarki\checkmarkr\checkmarkc\checkmark$\checkmark:\checkmark$\checkmark^\checkmark\\checkmarkc\checkmarki\checkmarkr\checkmarkc\checkmark$\checkmark \checkmark$\checkmark^\checkmark\\checkmarkc\checkmarki\checkmarkr\checkmarkc\checkmark$\checkmarkT\checkmark$\checkmark^\checkmark\\checkmarkc\checkmarki\checkmarkr\checkmarkc\checkmark$\checkmarkr\checkmark$\checkmark^\checkmark\\checkmarkc\checkmarki\checkmarkr\checkmarkc\checkmark$\checkmarka\checkmark$\checkmark^\checkmark\\checkmarkc\checkmarki\checkmarkr\checkmarkc\checkmark$\checkmarkn\checkmark$\checkmark^\checkmark\\checkmarkc\checkmarki\checkmarkr\checkmarkc\checkmark$\checkmarks\checkmark$\checkmark^\checkmark\\checkmarkc\checkmarki\checkmarkr\checkmarkc\checkmark$\checkmarkm\checkmark$\checkmark^\checkmark\\checkmarkc\checkmarki\checkmarkr\checkmarkc\checkmark$\checkmarke\checkmark$\checkmark^\checkmark\\checkmarkc\checkmarki\checkmarkr\checkmarkc\checkmark$\checkmarkt\checkmark$\checkmark^\checkmark\\checkmarkc\checkmarki\checkmarkr\checkmarkc\checkmark$\checkmarkt\checkmark$\checkmark^\checkmark\\checkmarkc\checkmarki\checkmarkr\checkmarkc\checkmark$\checkmarkr\checkmark$\checkmark^\checkmark\\checkmarkc\checkmarki\checkmarkr\checkmarkc\checkmark$\checkmarke\checkmark$\checkmark^\checkmark\\checkmarkc\checkmarki\checkmarkr\checkmarkc\checkmark$\checkmark \checkmark$\checkmark^\checkmark\\checkmarkc\checkmarki\checkmarkr\checkmarkc\checkmark$\checkmarkl\checkmark$\checkmark^\checkmark\\checkmarkc\checkmarki\checkmarkr\checkmarkc\checkmark$\checkmarke\checkmark$\checkmark^\checkmark\\checkmarkc\checkmarki\checkmarkr\checkmarkc\checkmark$\checkmark \checkmark$\checkmark^\checkmark\\checkmarkc\checkmarki\checkmarkr\checkmarkc\checkmark$\checkmarkc\checkmark$\checkmark^\checkmark\\checkmarkc\checkmarki\checkmarkr\checkmarkc\checkmark$\checkmarko\checkmark$\checkmark^\checkmark\\checkmarkc\checkmarki\checkmarkr\checkmarkc\checkmark$\checkmarku\checkmark$\checkmark^\checkmark\\checkmarkc\checkmarki\checkmarkr\checkmarkc\checkmark$\checkmarkp\checkmark$\checkmark^\checkmark\\checkmarkc\checkmarki\checkmarkr\checkmarkc\checkmark$\checkmarkl\checkmark$\checkmark^\checkmark\\checkmarkc\checkmarki\checkmarkr\checkmarkc\checkmark$\checkmarke\checkmark$\checkmark^\checkmark\\checkmarkc\checkmarki\checkmarkr\checkmarkc\checkmark$\checkmark \checkmark$\checkmark^\checkmark\\checkmarkc\checkmarki\checkmarkr\checkmarkc\checkmark$\checkmark&\checkmark$\checkmark^\checkmark\\checkmarkc\checkmarki\checkmarkr\checkmarkc\checkmark$\checkmark \checkmark$\checkmark^\checkmark\\checkmarkc\checkmarki\checkmarkr\checkmarkc\checkmark$\checkmarkA\checkmark$\checkmark^\checkmark\\checkmarkc\checkmarki\checkmarkr\checkmarkc\checkmark$\checkmarkr\checkmark$\checkmark^\checkmark\\checkmarkc\checkmarki\checkmarkr\checkmarkc\checkmark$\checkmarkb\checkmark$\checkmark^\checkmark\\checkmarkc\checkmarki\checkmarkr\checkmarkc\checkmark$\checkmarkr\checkmark$\checkmark^\checkmark\\checkmarkc\checkmarki\checkmarkr\checkmarkc\checkmark$\checkmarke\checkmark$\checkmark^\checkmark\\checkmarkc\checkmarki\checkmarkr\checkmarkc\checkmark$\checkmark \checkmark$\checkmark^\checkmark\\checkmarkc\checkmarki\checkmarkr\checkmarkc\checkmark$\checkmarke\checkmark$\checkmark^\checkmark\\checkmarkc\checkmarki\checkmarkr\checkmarkc\checkmark$\checkmarkn\checkmark$\checkmark^\checkmark\\checkmarkc\checkmarki\checkmarkr\checkmarkc\checkmark$\checkmark \checkmark$\checkmark^\checkmark\\checkmarkc\checkmarki\checkmarkr\checkmarkc\checkmark$\checkmarkA\checkmark$\checkmark^\checkmark\\checkmarkc\checkmarki\checkmarkr\checkmarkc\checkmark$\checkmarkc\checkmark$\checkmark^\checkmark\\checkmarkc\checkmarki\checkmarkr\checkmarkc\checkmark$\checkmarki\checkmark$\checkmark^\checkmark\\checkmarkc\checkmarki\checkmarkr\checkmarkc\checkmark$\checkmarke\checkmark$\checkmark^\checkmark\\checkmarkc\checkmarki\checkmarkr\checkmarkc\checkmark$\checkmarkr\checkmark$\checkmark^\checkmark\\checkmarkc\checkmarki\checkmarkr\checkmarkc\checkmark$\checkmark \checkmark$\checkmark^\checkmark\\checkmarkc\checkmarki\checkmarkr\checkmarkc\checkmark$\checkmarkC\checkmark$\checkmark^\checkmark\\checkmarkc\checkmarki\checkmarkr\checkmarkc\checkmark$\checkmark4\checkmark$\checkmark^\checkmark\\checkmarkc\checkmarki\checkmarkr\checkmarkc\checkmark$\checkmark5\checkmark$\checkmark^\checkmark\\checkmarkc\checkmarki\checkmarkr\checkmarkc\checkmark$\checkmark \checkmark$\checkmark^\checkmark\\checkmarkc\checkmarki\checkmarkr\checkmarkc\checkmark$\checkmark\\checkmark$\checkmark^\checkmark\\checkmarkc\checkmarki\checkmarkr\checkmarkc\checkmark$\checkmark\\checkmark$\checkmark^\checkmark\\checkmarkc\checkmarki\checkmarkr\checkmarkc\checkmark$\checkmark \checkmark$\checkmark^\checkmark\\checkmarkc\checkmarki\checkmarkr\checkmarkc\checkmark$\checkmark\\checkmark$\checkmark^\checkmark\\checkmarkc\checkmarki\checkmarkr\checkmarkc\checkmark$\checkmarkc\checkmark$\checkmark^\checkmark\\checkmarkc\checkmarki\checkmarkr\checkmarkc\checkmark$\checkmarkl\checkmark$\checkmark^\checkmark\\checkmarkc\checkmarki\checkmarkr\checkmarkc\checkmark$\checkmarki\checkmark$\checkmark^\checkmark\\checkmarkc\checkmarki\checkmarkr\checkmarkc\checkmark$\checkmarkn\checkmark$\checkmark^\checkmark\\checkmarkc\checkmarki\checkmarkr\checkmarkc\checkmark$\checkmarke\checkmark$\checkmark^\checkmark\\checkmarkc\checkmarki\checkmarkr\checkmarkc\checkmark$\checkmark{\checkmark$\checkmark^\checkmark\\checkmarkc\checkmarki\checkmarkr\checkmarkc\checkmark$\checkmark2\checkmark$\checkmark^\checkmark\\checkmarkc\checkmarki\checkmarkr\checkmarkc\checkmark$\checkmark-\checkmark$\checkmark^\checkmark\\checkmarkc\checkmarki\checkmarkr\checkmarkc\checkmark$\checkmark3\checkmark$\checkmark^\checkmark\\checkmarkc\checkmarki\checkmarkr\checkmarkc\checkmark$\checkmark}\checkmark$\checkmark^\checkmark\\checkmarkc\checkmarki\checkmarkr\checkmarkc\checkmark$\checkmark
\checkmark$\checkmark^\checkmark\\checkmarkc\checkmarki\checkmarkr\checkmarkc\checkmark$\checkmark \checkmark$\checkmark^\checkmark\\checkmarkc\checkmarki\checkmarkr\checkmarkc\checkmark$\checkmark&\checkmark$\checkmark^\checkmark\\checkmarkc\checkmarki\checkmarkr\checkmarkc\checkmark$\checkmark \checkmark$\checkmark^\checkmark\\checkmarkc\checkmarki\checkmarkr\checkmarkc\checkmark$\checkmarkF\checkmark$\checkmark^\checkmark\\checkmarkc\checkmarki\checkmarkr\checkmarkc\checkmark$\checkmarkT\checkmark$\checkmark^\checkmark\\checkmarkc\checkmarki\checkmarkr\checkmarkc\checkmark$\checkmark3\checkmark$\checkmark^\checkmark\\checkmarkc\checkmarki\checkmarkr\checkmarkc\checkmark$\checkmark4\checkmark$\checkmark^\checkmark\\checkmarkc\checkmarki\checkmarkr\checkmarkc\checkmark$\checkmark \checkmark$\checkmark^\checkmark\\checkmarkc\checkmarki\checkmarkr\checkmarkc\checkmark$\checkmark:\checkmark$\checkmark^\checkmark\\checkmarkc\checkmarki\checkmarkr\checkmarkc\checkmark$\checkmark \checkmark$\checkmark^\checkmark\\checkmarkc\checkmarki\checkmarkr\checkmarkc\checkmark$\checkmarkG\checkmark$\checkmark^\checkmark\\checkmarkc\checkmarki\checkmarkr\checkmarkc\checkmark$\checkmarké\checkmark$\checkmark^\checkmark\\checkmarkc\checkmarki\checkmarkr\checkmarkc\checkmark$\checkmarkn\checkmark$\checkmark^\checkmark\\checkmarkc\checkmarki\checkmarkr\checkmarkc\checkmark$\checkmarké\checkmark$\checkmark^\checkmark\\checkmarkc\checkmarki\checkmarkr\checkmarkc\checkmark$\checkmarkr\checkmark$\checkmark^\checkmark\\checkmarkc\checkmarki\checkmarkr\checkmarkc\checkmark$\checkmarke\checkmark$\checkmark^\checkmark\\checkmarkc\checkmarki\checkmarkr\checkmarkc\checkmark$\checkmarkr\checkmark$\checkmark^\checkmark\\checkmarkc\checkmarki\checkmarkr\checkmarkc\checkmark$\checkmark \checkmark$\checkmark^\checkmark\\checkmarkc\checkmarki\checkmarkr\checkmarkc\checkmark$\checkmarkl\checkmark$\checkmark^\checkmark\\checkmarkc\checkmarki\checkmarkr\checkmarkc\checkmark$\checkmarke\checkmark$\checkmark^\checkmark\\checkmarkc\checkmarki\checkmarkr\checkmarkc\checkmark$\checkmark \checkmark$\checkmark^\checkmark\\checkmarkc\checkmarki\checkmarkr\checkmarkc\checkmark$\checkmarkc\checkmark$\checkmark^\checkmark\\checkmarkc\checkmarki\checkmarkr\checkmarkc\checkmark$\checkmarko\checkmark$\checkmark^\checkmark\\checkmarkc\checkmarki\checkmarkr\checkmarkc\checkmark$\checkmarku\checkmark$\checkmark^\checkmark\\checkmarkc\checkmarki\checkmarkr\checkmarkc\checkmark$\checkmarkp\checkmark$\checkmark^\checkmark\\checkmarkc\checkmarki\checkmarkr\checkmarkc\checkmark$\checkmarkl\checkmark$\checkmark^\checkmark\\checkmarkc\checkmarki\checkmarkr\checkmarkc\checkmark$\checkmarke\checkmark$\checkmark^\checkmark\\checkmarkc\checkmarki\checkmarkr\checkmarkc\checkmark$\checkmark \checkmark$\checkmark^\checkmark\\checkmarkc\checkmarki\checkmarkr\checkmarkc\checkmark$\checkmarkr\checkmark$\checkmark^\checkmark\\checkmarkc\checkmarki\checkmarkr\checkmarkc\checkmark$\checkmarko\checkmark$\checkmark^\checkmark\\checkmarkc\checkmarki\checkmarkr\checkmarkc\checkmark$\checkmarkt\checkmark$\checkmark^\checkmark\\checkmarkc\checkmarki\checkmarkr\checkmarkc\checkmark$\checkmarka\checkmark$\checkmark^\checkmark\\checkmarkc\checkmarki\checkmarkr\checkmarkc\checkmark$\checkmarkt\checkmark$\checkmark^\checkmark\\checkmarkc\checkmarki\checkmarkr\checkmarkc\checkmark$\checkmarki\checkmark$\checkmark^\checkmark\\checkmarkc\checkmarki\checkmarkr\checkmarkc\checkmark$\checkmarkf\checkmark$\checkmark^\checkmark\\checkmarkc\checkmarki\checkmarkr\checkmarkc\checkmark$\checkmark \checkmark$\checkmark^\checkmark\\checkmarkc\checkmarki\checkmarkr\checkmarkc\checkmark$\checkmark&\checkmark$\checkmark^\checkmark\\checkmarkc\checkmarki\checkmarkr\checkmarkc\checkmark$\checkmark \checkmark$\checkmark^\checkmark\\checkmarkc\checkmarki\checkmarkr\checkmarkc\checkmark$\checkmarkM\checkmark$\checkmark^\checkmark\\checkmarkc\checkmarki\checkmarkr\checkmarkc\checkmark$\checkmarko\checkmark$\checkmark^\checkmark\\checkmarkc\checkmarki\checkmarkr\checkmarkc\checkmark$\checkmarkt\checkmark$\checkmark^\checkmark\\checkmarkc\checkmarki\checkmarkr\checkmarkc\checkmark$\checkmarke\checkmark$\checkmark^\checkmark\\checkmarkc\checkmarki\checkmarkr\checkmarkc\checkmark$\checkmarku\checkmark$\checkmark^\checkmark\\checkmarkc\checkmarki\checkmarkr\checkmarkc\checkmark$\checkmarkr\checkmark$\checkmark^\checkmark\\checkmarkc\checkmarki\checkmarkr\checkmarkc\checkmark$\checkmark \checkmark$\checkmark^\checkmark\\checkmarkc\checkmarki\checkmarkr\checkmarkc\checkmark$\checkmarkN\checkmark$\checkmark^\checkmark\\checkmarkc\checkmarki\checkmarkr\checkmarkc\checkmark$\checkmarkE\checkmark$\checkmark^\checkmark\\checkmarkc\checkmarki\checkmarkr\checkmarkc\checkmark$\checkmarkM\checkmark$\checkmark^\checkmark\\checkmarkc\checkmarki\checkmarkr\checkmarkc\checkmark$\checkmarkA\checkmark$\checkmark^\checkmark\\checkmarkc\checkmarki\checkmarkr\checkmarkc\checkmark$\checkmark \checkmark$\checkmark^\checkmark\\checkmarkc\checkmarki\checkmarkr\checkmarkc\checkmark$\checkmark2\checkmark$\checkmark^\checkmark\\checkmarkc\checkmarki\checkmarkr\checkmarkc\checkmark$\checkmark3\checkmark$\checkmark^\checkmark\\checkmarkc\checkmarki\checkmarkr\checkmarkc\checkmark$\checkmark \checkmark$\checkmark^\checkmark\\checkmarkc\checkmarki\checkmarkr\checkmarkc\checkmark$\checkmark(\checkmark$\checkmark^\checkmark\\checkmarkc\checkmarki\checkmarkr\checkmarkc\checkmark$\checkmarkA\checkmark$\checkmark^\checkmark\\checkmarkc\checkmarki\checkmarkr\checkmarkc\checkmark$\checkmarkx\checkmark$\checkmark^\checkmark\\checkmarkc\checkmarki\checkmarkr\checkmarkc\checkmark$\checkmarke\checkmark$\checkmark^\checkmark\\checkmarkc\checkmarki\checkmarkr\checkmarkc\checkmark$\checkmark \checkmark$\checkmark^\checkmark\\checkmarkc\checkmarki\checkmarkr\checkmarkc\checkmark$\checkmarkA\checkmark$\checkmark^\checkmark\\checkmarkc\checkmarki\checkmarkr\checkmarkc\checkmark$\checkmark)\checkmark$\checkmark^\checkmark\\checkmarkc\checkmarki\checkmarkr\checkmarkc\checkmark$\checkmark \checkmark$\checkmark^\checkmark\\checkmarkc\checkmarki\checkmarkr\checkmarkc\checkmark$\checkmark/\checkmark$\checkmark^\checkmark\\checkmarkc\checkmarki\checkmarkr\checkmarkc\checkmark$\checkmark \checkmark$\checkmark^\checkmark\\checkmarkc\checkmarki\checkmarkr\checkmarkc\checkmark$\checkmarkN\checkmark$\checkmark^\checkmark\\checkmarkc\checkmarki\checkmarkr\checkmarkc\checkmark$\checkmarkE\checkmark$\checkmark^\checkmark\\checkmarkc\checkmarki\checkmarkr\checkmarkc\checkmark$\checkmarkM\checkmark$\checkmark^\checkmark\\checkmarkc\checkmarki\checkmarkr\checkmarkc\checkmark$\checkmarkA\checkmark$\checkmark^\checkmark\\checkmarkc\checkmarki\checkmarkr\checkmarkc\checkmark$\checkmark \checkmark$\checkmark^\checkmark\\checkmarkc\checkmarki\checkmarkr\checkmarkc\checkmark$\checkmark1\checkmark$\checkmark^\checkmark\\checkmarkc\checkmarki\checkmarkr\checkmarkc\checkmark$\checkmark7\checkmark$\checkmark^\checkmark\\checkmarkc\checkmarki\checkmarkr\checkmarkc\checkmark$\checkmark \checkmark$\checkmark^\checkmark\\checkmarkc\checkmarki\checkmarkr\checkmarkc\checkmark$\checkmark(\checkmark$\checkmark^\checkmark\\checkmarkc\checkmarki\checkmarkr\checkmarkc\checkmark$\checkmarkA\checkmark$\checkmark^\checkmark\\checkmarkc\checkmarki\checkmarkr\checkmarkc\checkmark$\checkmarkx\checkmark$\checkmark^\checkmark\\checkmarkc\checkmarki\checkmarkr\checkmarkc\checkmark$\checkmarke\checkmark$\checkmark^\checkmark\\checkmarkc\checkmarki\checkmarkr\checkmarkc\checkmark$\checkmark \checkmark$\checkmark^\checkmark\\checkmarkc\checkmarki\checkmarkr\checkmarkc\checkmark$\checkmarkB\checkmark$\checkmark^\checkmark\\checkmarkc\checkmarki\checkmarkr\checkmarkc\checkmark$\checkmark)\checkmark$\checkmark^\checkmark\\checkmarkc\checkmarki\checkmarkr\checkmarkc\checkmark$\checkmark \checkmark$\checkmark^\checkmark\\checkmarkc\checkmarki\checkmarkr\checkmarkc\checkmark$\checkmark\\checkmark$\checkmark^\checkmark\\checkmarkc\checkmarki\checkmarkr\checkmarkc\checkmark$\checkmark\\checkmark$\checkmark^\checkmark\\checkmarkc\checkmarki\checkmarkr\checkmarkc\checkmark$\checkmark \checkmark$\checkmark^\checkmark\\checkmarkc\checkmarki\checkmarkr\checkmarkc\checkmark$\checkmark\\checkmark$\checkmark^\checkmark\\checkmarkc\checkmarki\checkmarkr\checkmarkc\checkmark$\checkmarkh\checkmark$\checkmark^\checkmark\\checkmarkc\checkmarki\checkmarkr\checkmarkc\checkmark$\checkmarkl\checkmark$\checkmark^\checkmark\\checkmarkc\checkmarki\checkmarkr\checkmarkc\checkmark$\checkmarki\checkmark$\checkmark^\checkmark\\checkmarkc\checkmarki\checkmarkr\checkmarkc\checkmark$\checkmarkn\checkmark$\checkmark^\checkmark\\checkmarkc\checkmarki\checkmarkr\checkmarkc\checkmark$\checkmarke\checkmark$\checkmark^\checkmark\\checkmarkc\checkmarki\checkmarkr\checkmarkc\checkmark$\checkmark
\checkmark$\checkmark^\checkmark\\checkmarkc\checkmarki\checkmarkr\checkmarkc\checkmark$\checkmarkF\checkmark$\checkmark^\checkmark\\checkmarkc\checkmarki\checkmarkr\checkmarkc\checkmark$\checkmarkT\checkmark$\checkmark^\checkmark\\checkmarkc\checkmarki\checkmarkr\checkmarkc\checkmark$\checkmark4\checkmark$\checkmark^\checkmark\\checkmarkc\checkmarki\checkmarkr\checkmarkc\checkmark$\checkmark \checkmark$\checkmark^\checkmark\\checkmarkc\checkmarki\checkmarkr\checkmarkc\checkmark$\checkmark:\checkmark$\checkmark^\checkmark\\checkmarkc\checkmarki\checkmarkr\checkmarkc\checkmark$\checkmark \checkmark$\checkmark^\checkmark\\checkmarkc\checkmarki\checkmarkr\checkmarkc\checkmark$\checkmarkC\checkmark$\checkmark^\checkmark\\checkmarkc\checkmarki\checkmarkr\checkmarkc\checkmark$\checkmarko\checkmark$\checkmark^\checkmark\\checkmarkc\checkmarki\checkmarkr\checkmarkc\checkmark$\checkmarkn\checkmark$\checkmark^\checkmark\\checkmarkc\checkmarki\checkmarkr\checkmarkc\checkmark$\checkmarkt\checkmark$\checkmark^\checkmark\\checkmarkc\checkmarki\checkmarkr\checkmarkc\checkmark$\checkmarkr\checkmark$\checkmark^\checkmark\\checkmarkc\checkmarki\checkmarkr\checkmarkc\checkmark$\checkmarkô\checkmark$\checkmark^\checkmark\\checkmarkc\checkmarki\checkmarkr\checkmarkc\checkmark$\checkmarkl\checkmark$\checkmark^\checkmark\\checkmarkc\checkmarki\checkmarkr\checkmarkc\checkmark$\checkmarke\checkmark$\checkmark^\checkmark\\checkmarkc\checkmarki\checkmarkr\checkmarkc\checkmark$\checkmarkr\checkmark$\checkmark^\checkmark\\checkmarkc\checkmarki\checkmarkr\checkmarkc\checkmark$\checkmark \checkmark$\checkmark^\checkmark\\checkmarkc\checkmarki\checkmarkr\checkmarkc\checkmark$\checkmarkl\checkmark$\checkmark^\checkmark\\checkmarkc\checkmarki\checkmarkr\checkmarkc\checkmark$\checkmarke\checkmark$\checkmark^\checkmark\\checkmarkc\checkmarki\checkmarkr\checkmarkc\checkmark$\checkmarks\checkmark$\checkmark^\checkmark\\checkmarkc\checkmarki\checkmarkr\checkmarkc\checkmark$\checkmark \checkmark$\checkmark^\checkmark\\checkmarkc\checkmarki\checkmarkr\checkmarkc\checkmark$\checkmark5\checkmark$\checkmark^\checkmark\\checkmarkc\checkmarki\checkmarkr\checkmarkc\checkmark$\checkmark \checkmark$\checkmark^\checkmark\\checkmarkc\checkmarki\checkmarkr\checkmarkc\checkmark$\checkmarka\checkmark$\checkmark^\checkmark\\checkmarkc\checkmarki\checkmarkr\checkmarkc\checkmark$\checkmarkx\checkmark$\checkmark^\checkmark\\checkmarkc\checkmarki\checkmarkr\checkmarkc\checkmark$\checkmarke\checkmark$\checkmark^\checkmark\\checkmarkc\checkmarki\checkmarkr\checkmarkc\checkmark$\checkmarks\checkmark$\checkmark^\checkmark\\checkmarkc\checkmarki\checkmarkr\checkmarkc\checkmark$\checkmark \checkmark$\checkmark^\checkmark\\checkmarkc\checkmarki\checkmarkr\checkmarkc\checkmark$\checkmark&\checkmark$\checkmark^\checkmark\\checkmarkc\checkmarki\checkmarkr\checkmarkc\checkmark$\checkmark \checkmark$\checkmark^\checkmark\\checkmarkc\checkmarki\checkmarkr\checkmarkc\checkmark$\checkmarkF\checkmark$\checkmark^\checkmark\\checkmarkc\checkmarki\checkmarkr\checkmarkc\checkmark$\checkmarkT\checkmark$\checkmark^\checkmark\\checkmarkc\checkmarki\checkmarkr\checkmarkc\checkmark$\checkmark4\checkmark$\checkmark^\checkmark\\checkmarkc\checkmarki\checkmarkr\checkmarkc\checkmark$\checkmark1\checkmark$\checkmark^\checkmark\\checkmarkc\checkmarki\checkmarkr\checkmarkc\checkmark$\checkmark \checkmark$\checkmark^\checkmark\\checkmarkc\checkmarki\checkmarkr\checkmarkc\checkmark$\checkmark:\checkmark$\checkmark^\checkmark\\checkmarkc\checkmarki\checkmarkr\checkmarkc\checkmark$\checkmark \checkmark$\checkmark^\checkmark\\checkmarkc\checkmarki\checkmarkr\checkmarkc\checkmark$\checkmarkC\checkmark$\checkmark^\checkmark\\checkmarkc\checkmarki\checkmarkr\checkmarkc\checkmark$\checkmarka\checkmark$\checkmark^\checkmark\\checkmarkc\checkmarki\checkmarkr\checkmarkc\checkmark$\checkmarkl\checkmark$\checkmark^\checkmark\\checkmarkc\checkmarki\checkmarkr\checkmarkc\checkmark$\checkmarkc\checkmark$\checkmark^\checkmark\\checkmarkc\checkmarki\checkmarkr\checkmarkc\checkmark$\checkmarku\checkmark$\checkmark^\checkmark\\checkmarkc\checkmarki\checkmarkr\checkmarkc\checkmark$\checkmarkl\checkmark$\checkmark^\checkmark\\checkmarkc\checkmarki\checkmarkr\checkmarkc\checkmark$\checkmarke\checkmark$\checkmark^\checkmark\\checkmarkc\checkmarki\checkmarkr\checkmarkc\checkmark$\checkmarkr\checkmark$\checkmark^\checkmark\\checkmarkc\checkmarki\checkmarkr\checkmarkc\checkmark$\checkmark \checkmark$\checkmark^\checkmark\\checkmarkc\checkmarki\checkmarkr\checkmarkc\checkmark$\checkmarkl\checkmark$\checkmark^\checkmark\\checkmarkc\checkmarki\checkmarkr\checkmarkc\checkmark$\checkmarke\checkmark$\checkmark^\checkmark\\checkmarkc\checkmarki\checkmarkr\checkmarkc\checkmark$\checkmarks\checkmark$\checkmark^\checkmark\\checkmarkc\checkmarki\checkmarkr\checkmarkc\checkmark$\checkmark \checkmark$\checkmark^\checkmark\\checkmarkc\checkmarki\checkmarkr\checkmarkc\checkmark$\checkmarkt\checkmark$\checkmark^\checkmark\\checkmarkc\checkmarki\checkmarkr\checkmarkc\checkmark$\checkmarkr\checkmark$\checkmark^\checkmark\\checkmarkc\checkmarki\checkmarkr\checkmarkc\checkmark$\checkmarka\checkmark$\checkmark^\checkmark\\checkmarkc\checkmarki\checkmarkr\checkmarkc\checkmark$\checkmarkj\checkmark$\checkmark^\checkmark\\checkmarkc\checkmarki\checkmarkr\checkmarkc\checkmark$\checkmarke\checkmark$\checkmark^\checkmark\\checkmarkc\checkmarki\checkmarkr\checkmarkc\checkmark$\checkmarkc\checkmark$\checkmark^\checkmark\\checkmarkc\checkmarki\checkmarkr\checkmarkc\checkmark$\checkmarkt\checkmark$\checkmark^\checkmark\\checkmarkc\checkmarki\checkmarkr\checkmarkc\checkmark$\checkmarko\checkmark$\checkmark^\checkmark\\checkmarkc\checkmarki\checkmarkr\checkmarkc\checkmark$\checkmarki\checkmark$\checkmark^\checkmark\\checkmarkc\checkmarki\checkmarkr\checkmarkc\checkmark$\checkmarkr\checkmark$\checkmark^\checkmark\\checkmarkc\checkmarki\checkmarkr\checkmarkc\checkmark$\checkmarke\checkmark$\checkmark^\checkmark\\checkmarkc\checkmarki\checkmarkr\checkmarkc\checkmark$\checkmarks\checkmark$\checkmark^\checkmark\\checkmarkc\checkmarki\checkmarkr\checkmarkc\checkmark$\checkmark \checkmark$\checkmark^\checkmark\\checkmarkc\checkmarki\checkmarkr\checkmarkc\checkmark$\checkmark(\checkmark$\checkmark^\checkmark\\checkmarkc\checkmarki\checkmarkr\checkmarkc\checkmark$\checkmarkR\checkmark$\checkmark^\checkmark\\checkmarkc\checkmarki\checkmarkr\checkmarkc\checkmark$\checkmarkT\checkmark$\checkmark^\checkmark\\checkmarkc\checkmarki\checkmarkr\checkmarkc\checkmark$\checkmarkC\checkmark$\checkmark^\checkmark\\checkmarkc\checkmarki\checkmarkr\checkmarkc\checkmark$\checkmarkP\checkmark$\checkmark^\checkmark\\checkmarkc\checkmarki\checkmarkr\checkmarkc\checkmark$\checkmark)\checkmark$\checkmark^\checkmark\\checkmarkc\checkmarki\checkmarkr\checkmarkc\checkmark$\checkmark \checkmark$\checkmark^\checkmark\\checkmarkc\checkmarki\checkmarkr\checkmarkc\checkmark$\checkmark&\checkmark$\checkmark^\checkmark\\checkmarkc\checkmarki\checkmarkr\checkmarkc\checkmark$\checkmark \checkmark$\checkmark^\checkmark\\checkmarkc\checkmarki\checkmarkr\checkmarkc\checkmark$\checkmarkF\checkmark$\checkmark^\checkmark\\checkmarkc\checkmarki\checkmarkr\checkmarkc\checkmark$\checkmarki\checkmark$\checkmark^\checkmark\\checkmarkc\checkmarki\checkmarkr\checkmarkc\checkmark$\checkmarkr\checkmark$\checkmark^\checkmark\\checkmarkc\checkmarki\checkmarkr\checkmarkc\checkmark$\checkmarkm\checkmark$\checkmark^\checkmark\\checkmarkc\checkmarki\checkmarkr\checkmarkc\checkmark$\checkmarkw\checkmark$\checkmark^\checkmark\\checkmarkc\checkmarki\checkmarkr\checkmarkc\checkmark$\checkmarka\checkmark$\checkmark^\checkmark\\checkmarkc\checkmarki\checkmarkr\checkmarkc\checkmark$\checkmarkr\checkmark$\checkmark^\checkmark\\checkmarkc\checkmarki\checkmarkr\checkmarkc\checkmark$\checkmarke\checkmark$\checkmark^\checkmark\\checkmarkc\checkmarki\checkmarkr\checkmarkc\checkmark$\checkmark \checkmark$\checkmark^\checkmark\\checkmarkc\checkmarki\checkmarkr\checkmarkc\checkmark$\checkmarkE\checkmark$\checkmark^\checkmark\\checkmarkc\checkmarki\checkmarkr\checkmarkc\checkmark$\checkmarkS\checkmark$\checkmark^\checkmark\\checkmarkc\checkmarki\checkmarkr\checkmarkc\checkmark$\checkmarkP\checkmark$\checkmark^\checkmark\\checkmarkc\checkmarki\checkmarkr\checkmarkc\checkmark$\checkmark3\checkmark$\checkmark^\checkmark\\checkmarkc\checkmarki\checkmarkr\checkmarkc\checkmark$\checkmark2\checkmark$\checkmark^\checkmark\\checkmarkc\checkmarki\checkmarkr\checkmarkc\checkmark$\checkmark \checkmark$\checkmark^\checkmark\\checkmarkc\checkmarki\checkmarkr\checkmarkc\checkmark$\checkmark(\checkmark$\checkmark^\checkmark\\checkmarkc\checkmarki\checkmarkr\checkmarkc\checkmark$\checkmarkF\checkmark$\checkmark^\checkmark\\checkmarkc\checkmarki\checkmarkr\checkmarkc\checkmark$\checkmarkr\checkmark$\checkmark^\checkmark\\checkmarkc\checkmarki\checkmarkr\checkmarkc\checkmark$\checkmarke\checkmark$\checkmark^\checkmark\\checkmarkc\checkmarki\checkmarkr\checkmarkc\checkmark$\checkmarke\checkmark$\checkmark^\checkmark\\checkmarkc\checkmarki\checkmarkr\checkmarkc\checkmark$\checkmarkR\checkmark$\checkmark^\checkmark\\checkmarkc\checkmarki\checkmarkr\checkmarkc\checkmark$\checkmarkT\checkmark$\checkmark^\checkmark\\checkmarkc\checkmarki\checkmarkr\checkmarkc\checkmark$\checkmarkO\checkmark$\checkmark^\checkmark\\checkmarkc\checkmarki\checkmarkr\checkmarkc\checkmark$\checkmarkS\checkmark$\checkmark^\checkmark\\checkmarkc\checkmarki\checkmarkr\checkmarkc\checkmark$\checkmark)\checkmark$\checkmark^\checkmark\\checkmarkc\checkmarki\checkmarkr\checkmarkc\checkmark$\checkmark \checkmark$\checkmark^\checkmark\\checkmarkc\checkmarki\checkmarkr\checkmarkc\checkmark$\checkmark\\checkmark$\checkmark^\checkmark\\checkmarkc\checkmarki\checkmarkr\checkmarkc\checkmark$\checkmark\\checkmark$\checkmark^\checkmark\\checkmarkc\checkmarki\checkmarkr\checkmarkc\checkmark$\checkmark \checkmark$\checkmark^\checkmark\\checkmarkc\checkmarki\checkmarkr\checkmarkc\checkmark$\checkmark\\checkmark$\checkmark^\checkmark\\checkmarkc\checkmarki\checkmarkr\checkmarkc\checkmark$\checkmarkh\checkmark$\checkmark^\checkmark\\checkmarkc\checkmarki\checkmarkr\checkmarkc\checkmark$\checkmarkl\checkmark$\checkmark^\checkmark\\checkmarkc\checkmarki\checkmarkr\checkmarkc\checkmark$\checkmarki\checkmark$\checkmark^\checkmark\\checkmarkc\checkmarki\checkmarkr\checkmarkc\checkmark$\checkmarkn\checkmark$\checkmark^\checkmark\\checkmarkc\checkmarki\checkmarkr\checkmarkc\checkmark$\checkmarke\checkmark$\checkmark^\checkmark\\checkmarkc\checkmarki\checkmarkr\checkmarkc\checkmark$\checkmark
\checkmark$\checkmark^\checkmark\\checkmarkc\checkmarki\checkmarkr\checkmarkc\checkmark$\checkmarkF\checkmark$\checkmark^\checkmark\\checkmarkc\checkmarki\checkmarkr\checkmarkc\checkmark$\checkmarkT\checkmark$\checkmark^\checkmark\\checkmarkc\checkmarki\checkmarkr\checkmarkc\checkmark$\checkmark5\checkmark$\checkmark^\checkmark\\checkmarkc\checkmarki\checkmarkr\checkmarkc\checkmark$\checkmark \checkmark$\checkmark^\checkmark\\checkmarkc\checkmarki\checkmarkr\checkmarkc\checkmark$\checkmark:\checkmark$\checkmark^\checkmark\\checkmarkc\checkmarki\checkmarkr\checkmarkc\checkmark$\checkmark \checkmark$\checkmark^\checkmark\\checkmarkc\checkmarki\checkmarkr\checkmarkc\checkmark$\checkmarkS\checkmark$\checkmark^\checkmark\\checkmarkc\checkmarki\checkmarkr\checkmarkc\checkmark$\checkmarku\checkmark$\checkmark^\checkmark\\checkmarkc\checkmarki\checkmarkr\checkmarkc\checkmark$\checkmarkp\checkmark$\checkmark^\checkmark\\checkmarkc\checkmarki\checkmarkr\checkmarkc\checkmark$\checkmarke\checkmark$\checkmark^\checkmark\\checkmarkc\checkmarki\checkmarkr\checkmarkc\checkmark$\checkmarkr\checkmark$\checkmark^\checkmark\\checkmarkc\checkmarki\checkmarkr\checkmarkc\checkmark$\checkmarkv\checkmark$\checkmark^\checkmark\\checkmarkc\checkmarki\checkmarkr\checkmarkc\checkmark$\checkmarki\checkmark$\checkmark^\checkmark\\checkmarkc\checkmarki\checkmarkr\checkmarkc\checkmark$\checkmarks\checkmark$\checkmark^\checkmark\\checkmarkc\checkmarki\checkmarkr\checkmarkc\checkmark$\checkmarke\checkmark$\checkmark^\checkmark\\checkmarkc\checkmarki\checkmarkr\checkmarkc\checkmark$\checkmarkr\checkmark$\checkmark^\checkmark\\checkmarkc\checkmarki\checkmarkr\checkmarkc\checkmark$\checkmark \checkmark$\checkmark^\checkmark\\checkmarkc\checkmarki\checkmarkr\checkmarkc\checkmark$\checkmarke\checkmark$\checkmark^\checkmark\\checkmarkc\checkmarki\checkmarkr\checkmarkc\checkmark$\checkmarkn\checkmark$\checkmark^\checkmark\\checkmarkc\checkmarki\checkmarkr\checkmarkc\checkmark$\checkmark \checkmark$\checkmark^\checkmark\\checkmarkc\checkmarki\checkmarkr\checkmarkc\checkmark$\checkmarkt\checkmark$\checkmark^\checkmark\\checkmarkc\checkmarki\checkmarkr\checkmarkc\checkmark$\checkmarke\checkmark$\checkmark^\checkmark\\checkmarkc\checkmarki\checkmarkr\checkmarkc\checkmark$\checkmarkm\checkmark$\checkmark^\checkmark\\checkmarkc\checkmarki\checkmarkr\checkmarkc\checkmark$\checkmarkp\checkmark$\checkmark^\checkmark\\checkmarkc\checkmarki\checkmarkr\checkmarkc\checkmark$\checkmarks\checkmark$\checkmark^\checkmark\\checkmarkc\checkmarki\checkmarkr\checkmarkc\checkmark$\checkmark \checkmark$\checkmark^\checkmark\\checkmarkc\checkmarki\checkmarkr\checkmarkc\checkmark$\checkmarkr\checkmark$\checkmark^\checkmark\\checkmarkc\checkmarki\checkmarkr\checkmarkc\checkmark$\checkmarké\checkmark$\checkmark^\checkmark\\checkmarkc\checkmarki\checkmarkr\checkmarkc\checkmark$\checkmarke\checkmark$\checkmark^\checkmark\\checkmarkc\checkmarki\checkmarkr\checkmarkc\checkmark$\checkmarkl\checkmark$\checkmark^\checkmark\\checkmarkc\checkmarki\checkmarkr\checkmarkc\checkmark$\checkmark \checkmark$\checkmark^\checkmark\\checkmarkc\checkmarki\checkmarkr\checkmarkc\checkmark$\checkmark&\checkmark$\checkmark^\checkmark\\checkmarkc\checkmarki\checkmarkr\checkmarkc\checkmark$\checkmark \checkmark$\checkmark^\checkmark\\checkmarkc\checkmarki\checkmarkr\checkmarkc\checkmark$\checkmarkF\checkmark$\checkmark^\checkmark\\checkmarkc\checkmarki\checkmarkr\checkmarkc\checkmark$\checkmarkT\checkmark$\checkmark^\checkmark\\checkmarkc\checkmarki\checkmarkr\checkmarkc\checkmark$\checkmark5\checkmark$\checkmark^\checkmark\\checkmarkc\checkmarki\checkmarkr\checkmarkc\checkmark$\checkmark1\checkmark$\checkmark^\checkmark\\checkmarkc\checkmarki\checkmarkr\checkmarkc\checkmark$\checkmark \checkmark$\checkmark^\checkmark\\checkmarkc\checkmarki\checkmarkr\checkmarkc\checkmark$\checkmark:\checkmark$\checkmark^\checkmark\\checkmarkc\checkmarki\checkmarkr\checkmarkc\checkmark$\checkmark \checkmark$\checkmark^\checkmark\\checkmarkc\checkmarki\checkmarkr\checkmarkc\checkmark$\checkmarkI\checkmark$\checkmark^\checkmark\\checkmarkc\checkmarki\checkmarkr\checkmarkc\checkmark$\checkmarkn\checkmark$\checkmark^\checkmark\\checkmarkc\checkmarki\checkmarkr\checkmarkc\checkmark$\checkmarkt\checkmark$\checkmark^\checkmark\\checkmarkc\checkmarki\checkmarkr\checkmarkc\checkmark$\checkmarke\checkmark$\checkmark^\checkmark\\checkmarkc\checkmarki\checkmarkr\checkmarkc\checkmark$\checkmarkr\checkmark$\checkmark^\checkmark\\checkmarkc\checkmarki\checkmarkr\checkmarkc\checkmark$\checkmarkf\checkmark$\checkmark^\checkmark\\checkmarkc\checkmarki\checkmarkr\checkmarkc\checkmark$\checkmarka\checkmark$\checkmark^\checkmark\\checkmarkc\checkmarki\checkmarkr\checkmarkc\checkmark$\checkmarkc\checkmark$\checkmark^\checkmark\\checkmarkc\checkmarki\checkmarkr\checkmarkc\checkmark$\checkmarke\checkmark$\checkmark^\checkmark\\checkmarkc\checkmarki\checkmarkr\checkmarkc\checkmark$\checkmarkr\checkmark$\checkmark^\checkmark\\checkmarkc\checkmarki\checkmarkr\checkmarkc\checkmark$\checkmark \checkmark$\checkmark^\checkmark\\checkmarkc\checkmarki\checkmarkr\checkmarkc\checkmark$\checkmarkl\checkmark$\checkmark^\checkmark\\checkmarkc\checkmarki\checkmarkr\checkmarkc\checkmark$\checkmark'\checkmark$\checkmark^\checkmark\\checkmarkc\checkmarki\checkmarkr\checkmarkc\checkmark$\checkmarko\checkmark$\checkmark^\checkmark\\checkmarkc\checkmarki\checkmarkr\checkmarkc\checkmark$\checkmarkp\checkmark$\checkmark^\checkmark\\checkmarkc\checkmarki\checkmarkr\checkmarkc\checkmark$\checkmarké\checkmark$\checkmark^\checkmark\\checkmarkc\checkmarki\checkmarkr\checkmarkc\checkmark$\checkmarkr\checkmark$\checkmark^\checkmark\\checkmarkc\checkmarki\checkmarkr\checkmarkc\checkmark$\checkmarka\checkmark$\checkmark^\checkmark\\checkmarkc\checkmarki\checkmarkr\checkmarkc\checkmark$\checkmarkt\checkmark$\checkmark^\checkmark\\checkmarkc\checkmarki\checkmarkr\checkmarkc\checkmark$\checkmarke\checkmark$\checkmark^\checkmark\\checkmarkc\checkmarki\checkmarkr\checkmarkc\checkmark$\checkmarku\checkmark$\checkmark^\checkmark\\checkmarkc\checkmarki\checkmarkr\checkmarkc\checkmark$\checkmarkr\checkmark$\checkmark^\checkmark\\checkmarkc\checkmarki\checkmarkr\checkmarkc\checkmark$\checkmark \checkmark$\checkmark^\checkmark\\checkmarkc\checkmarki\checkmarkr\checkmarkc\checkmark$\checkmark&\checkmark$\checkmark^\checkmark\\checkmarkc\checkmarki\checkmarkr\checkmarkc\checkmark$\checkmark \checkmark$\checkmark^\checkmark\\checkmarkc\checkmarki\checkmarkr\checkmarkc\checkmark$\checkmarkI\checkmark$\checkmark^\checkmark\\checkmarkc\checkmarki\checkmarkr\checkmarkc\checkmark$\checkmarkH\checkmark$\checkmark^\checkmark\\checkmarkc\checkmarki\checkmarkr\checkmarkc\checkmark$\checkmarkM\checkmark$\checkmark^\checkmark\\checkmarkc\checkmarki\checkmarkr\checkmarkc\checkmark$\checkmark \checkmark$\checkmark^\checkmark\\checkmarkc\checkmarki\checkmarkr\checkmarkc\checkmark$\checkmarkF\checkmark$\checkmark^\checkmark\\checkmarkc\checkmarki\checkmarkr\checkmarkc\checkmark$\checkmarkl\checkmark$\checkmark^\checkmark\\checkmarkc\checkmarki\checkmarkr\checkmarkc\checkmark$\checkmarku\checkmark$\checkmark^\checkmark\\checkmarkc\checkmarki\checkmarkr\checkmarkc\checkmark$\checkmarkt\checkmark$\checkmark^\checkmark\\checkmarkc\checkmarki\checkmarkr\checkmarkc\checkmark$\checkmarkt\checkmark$\checkmark^\checkmark\\checkmarkc\checkmarki\checkmarkr\checkmarkc\checkmark$\checkmarke\checkmark$\checkmark^\checkmark\\checkmarkc\checkmarki\checkmarkr\checkmarkc\checkmark$\checkmarkr\checkmark$\checkmark^\checkmark\\checkmarkc\checkmarki\checkmarkr\checkmarkc\checkmark$\checkmark \checkmark$\checkmark^\checkmark\\checkmarkc\checkmarki\checkmarkr\checkmarkc\checkmark$\checkmark(\checkmark$\checkmark^\checkmark\\checkmarkc\checkmarki\checkmarkr\checkmarkc\checkmark$\checkmarkA\checkmark$\checkmark^\checkmark\\checkmarkc\checkmarki\checkmarkr\checkmarkc\checkmark$\checkmarkn\checkmark$\checkmark^\checkmark\\checkmarkc\checkmarki\checkmarkr\checkmarkc\checkmark$\checkmarkd\checkmark$\checkmark^\checkmark\\checkmarkc\checkmarki\checkmarkr\checkmarkc\checkmark$\checkmarkr\checkmark$\checkmark^\checkmark\\checkmarkc\checkmarki\checkmarkr\checkmarkc\checkmark$\checkmarko\checkmark$\checkmark^\checkmark\\checkmarkc\checkmarki\checkmarkr\checkmarkc\checkmark$\checkmarki\checkmark$\checkmark^\checkmark\\checkmarkc\checkmarki\checkmarkr\checkmarkc\checkmark$\checkmarkd\checkmark$\checkmark^\checkmark\\checkmarkc\checkmarki\checkmarkr\checkmarkc\checkmark$\checkmark/\checkmark$\checkmark^\checkmark\\checkmarkc\checkmarki\checkmarkr\checkmarkc\checkmark$\checkmarkW\checkmark$\checkmark^\checkmark\\checkmarkc\checkmarki\checkmarkr\checkmarkc\checkmark$\checkmarki\checkmark$\checkmark^\checkmark\\checkmarkc\checkmarki\checkmarkr\checkmarkc\checkmark$\checkmarkn\checkmark$\checkmark^\checkmark\\checkmarkc\checkmarki\checkmarkr\checkmarkc\checkmark$\checkmarkd\checkmark$\checkmark^\checkmark\\checkmarkc\checkmarki\checkmarkr\checkmarkc\checkmark$\checkmarko\checkmark$\checkmark^\checkmark\\checkmarkc\checkmarki\checkmarkr\checkmarkc\checkmark$\checkmarkw\checkmark$\checkmark^\checkmark\\checkmarkc\checkmarki\checkmarkr\checkmarkc\checkmark$\checkmarks\checkmark$\checkmark^\checkmark\\checkmarkc\checkmarki\checkmarkr\checkmarkc\checkmark$\checkmark)\checkmark$\checkmark^\checkmark\\checkmarkc\checkmarki\checkmarkr\checkmarkc\checkmark$\checkmark \checkmark$\checkmark^\checkmark\\checkmarkc\checkmarki\checkmarkr\checkmarkc\checkmark$\checkmark\\checkmark$\checkmark^\checkmark\\checkmarkc\checkmarki\checkmarkr\checkmarkc\checkmark$\checkmark\\checkmark$\checkmark^\checkmark\\checkmarkc\checkmarki\checkmarkr\checkmarkc\checkmark$\checkmark \checkmark$\checkmark^\checkmark\\checkmarkc\checkmarki\checkmarkr\checkmarkc\checkmark$\checkmark\\checkmark$\checkmark^\checkmark\\checkmarkc\checkmarki\checkmarkr\checkmarkc\checkmark$\checkmarkh\checkmark$\checkmark^\checkmark\\checkmarkc\checkmarki\checkmarkr\checkmarkc\checkmark$\checkmarkl\checkmark$\checkmark^\checkmark\\checkmarkc\checkmarki\checkmarkr\checkmarkc\checkmark$\checkmarki\checkmark$\checkmark^\checkmark\\checkmarkc\checkmarki\checkmarkr\checkmarkc\checkmark$\checkmarkn\checkmark$\checkmark^\checkmark\\checkmarkc\checkmarki\checkmarkr\checkmarkc\checkmark$\checkmarke\checkmark$\checkmark^\checkmark\\checkmarkc\checkmarki\checkmarkr\checkmarkc\checkmark$\checkmark
\checkmark$\checkmark^\checkmark\\checkmarkc\checkmarki\checkmarkr\checkmarkc\checkmark$\checkmark\\checkmark$\checkmark^\checkmark\\checkmarkc\checkmarki\checkmarkr\checkmarkc\checkmark$\checkmarke\checkmark$\checkmark^\checkmark\\checkmarkc\checkmarki\checkmarkr\checkmarkc\checkmark$\checkmarkn\checkmark$\checkmark^\checkmark\\checkmarkc\checkmarki\checkmarkr\checkmarkc\checkmark$\checkmarkd\checkmark$\checkmark^\checkmark\\checkmarkc\checkmarki\checkmarkr\checkmarkc\checkmark$\checkmark{\checkmark$\checkmark^\checkmark\\checkmarkc\checkmarki\checkmarkr\checkmarkc\checkmark$\checkmarkt\checkmark$\checkmark^\checkmark\\checkmarkc\checkmarki\checkmarkr\checkmarkc\checkmark$\checkmarka\checkmark$\checkmark^\checkmark\\checkmarkc\checkmarki\checkmarkr\checkmarkc\checkmark$\checkmarkb\checkmark$\checkmark^\checkmark\\checkmarkc\checkmarki\checkmarkr\checkmarkc\checkmark$\checkmarku\checkmark$\checkmark^\checkmark\\checkmarkc\checkmarki\checkmarkr\checkmarkc\checkmark$\checkmarkl\checkmark$\checkmark^\checkmark\\checkmarkc\checkmarki\checkmarkr\checkmarkc\checkmark$\checkmarka\checkmark$\checkmark^\checkmark\\checkmarkc\checkmarki\checkmarkr\checkmarkc\checkmark$\checkmarkr\checkmark$\checkmark^\checkmark\\checkmarkc\checkmarki\checkmarkr\checkmarkc\checkmark$\checkmark}\checkmark$\checkmark^\checkmark\\checkmarkc\checkmarki\checkmarkr\checkmarkc\checkmark$\checkmark
\checkmark$\checkmark^\checkmark\\checkmarkc\checkmarki\checkmarkr\checkmarkc\checkmark$\checkmark}\checkmark$\checkmark^\checkmark\\checkmarkc\checkmarki\checkmarkr\checkmarkc\checkmark$\checkmark
\checkmark$\checkmark^\checkmark\\checkmarkc\checkmarki\checkmarkr\checkmarkc\checkmark$\checkmark\\checkmark$\checkmark^\checkmark\\checkmarkc\checkmarki\checkmarkr\checkmarkc\checkmark$\checkmarkc\checkmark$\checkmark^\checkmark\\checkmarkc\checkmarki\checkmarkr\checkmarkc\checkmark$\checkmarka\checkmark$\checkmark^\checkmark\\checkmarkc\checkmarki\checkmarkr\checkmarkc\checkmark$\checkmarkp\checkmark$\checkmark^\checkmark\\checkmarkc\checkmarki\checkmarkr\checkmarkc\checkmark$\checkmarkt\checkmark$\checkmark^\checkmark\\checkmarkc\checkmarki\checkmarkr\checkmarkc\checkmark$\checkmarki\checkmark$\checkmark^\checkmark\\checkmarkc\checkmarki\checkmarkr\checkmarkc\checkmark$\checkmarko\checkmark$\checkmark^\checkmark\\checkmarkc\checkmarki\checkmarkr\checkmarkc\checkmark$\checkmarkn\checkmark$\checkmark^\checkmark\\checkmarkc\checkmarki\checkmarkr\checkmarkc\checkmark$\checkmark{\checkmark$\checkmark^\checkmark\\checkmarkc\checkmarki\checkmarkr\checkmarkc\checkmark$\checkmarkD\checkmark$\checkmark^\checkmark\\checkmarkc\checkmarki\checkmarkr\checkmarkc\checkmark$\checkmarki\checkmark$\checkmark^\checkmark\\checkmarkc\checkmarki\checkmarkr\checkmarkc\checkmark$\checkmarka\checkmark$\checkmark^\checkmark\\checkmarkc\checkmarki\checkmarkr\checkmarkc\checkmark$\checkmarkg\checkmark$\checkmark^\checkmark\\checkmarkc\checkmarki\checkmarkr\checkmarkc\checkmark$\checkmarkr\checkmark$\checkmark^\checkmark\\checkmarkc\checkmarki\checkmarkr\checkmarkc\checkmark$\checkmarka\checkmark$\checkmark^\checkmark\\checkmarkc\checkmarki\checkmarkr\checkmarkc\checkmark$\checkmarkm\checkmark$\checkmark^\checkmark\\checkmarkc\checkmarki\checkmarkr\checkmarkc\checkmark$\checkmarkm\checkmark$\checkmark^\checkmark\\checkmarkc\checkmarki\checkmarkr\checkmarkc\checkmark$\checkmarke\checkmark$\checkmark^\checkmark\\checkmarkc\checkmarki\checkmarkr\checkmarkc\checkmark$\checkmark \checkmark$\checkmark^\checkmark\\checkmarkc\checkmarki\checkmarkr\checkmarkc\checkmark$\checkmarkF\checkmark$\checkmark^\checkmark\\checkmarkc\checkmarki\checkmarkr\checkmarkc\checkmark$\checkmarkA\checkmark$\checkmark^\checkmark\\checkmarkc\checkmarki\checkmarkr\checkmarkc\checkmark$\checkmarkS\checkmark$\checkmark^\checkmark\\checkmarkc\checkmarki\checkmarkr\checkmarkc\checkmark$\checkmarkT\checkmark$\checkmark^\checkmark\\checkmarkc\checkmarki\checkmarkr\checkmarkc\checkmark$\checkmark \checkmark$\checkmark^\checkmark\\checkmarkc\checkmarki\checkmarkr\checkmarkc\checkmark$\checkmarkc\checkmark$\checkmark^\checkmark\\checkmarkc\checkmarki\checkmarkr\checkmarkc\checkmark$\checkmarko\checkmark$\checkmark^\checkmark\\checkmarkc\checkmarki\checkmarkr\checkmarkc\checkmark$\checkmarkn\checkmark$\checkmark^\checkmark\\checkmarkc\checkmarki\checkmarkr\checkmarkc\checkmark$\checkmarkd\checkmark$\checkmark^\checkmark\\checkmarkc\checkmarki\checkmarkr\checkmarkc\checkmark$\checkmarke\checkmark$\checkmark^\checkmark\\checkmarkc\checkmarki\checkmarkr\checkmarkc\checkmark$\checkmarkn\checkmark$\checkmark^\checkmark\\checkmarkc\checkmarki\checkmarkr\checkmarkc\checkmark$\checkmarks\checkmark$\checkmark^\checkmark\\checkmarkc\checkmarki\checkmarkr\checkmarkc\checkmark$\checkmarké\checkmark$\checkmark^\checkmark\\checkmarkc\checkmarki\checkmarkr\checkmarkc\checkmark$\checkmark \checkmark$\checkmark^\checkmark\\checkmarkc\checkmarki\checkmarkr\checkmarkc\checkmark$\checkmark—\checkmark$\checkmark^\checkmark\\checkmarkc\checkmarki\checkmarkr\checkmarkc\checkmark$\checkmark \checkmark$\checkmark^\checkmark\\checkmarkc\checkmarki\checkmarkr\checkmarkc\checkmark$\checkmarkD\checkmark$\checkmark^\checkmark\\checkmarkc\checkmarki\checkmarkr\checkmarkc\checkmark$\checkmarké\checkmark$\checkmark^\checkmark\\checkmarkc\checkmarki\checkmarkr\checkmarkc\checkmark$\checkmarkc\checkmark$\checkmark^\checkmark\\checkmarkc\checkmarki\checkmarkr\checkmarkc\checkmark$\checkmarko\checkmark$\checkmark^\checkmark\\checkmarkc\checkmarki\checkmarkr\checkmarkc\checkmark$\checkmarkm\checkmark$\checkmark^\checkmark\\checkmarkc\checkmarki\checkmarkr\checkmarkc\checkmark$\checkmarkp\checkmark$\checkmark^\checkmark\\checkmarkc\checkmarki\checkmarkr\checkmarkc\checkmark$\checkmarko\checkmark$\checkmark^\checkmark\\checkmarkc\checkmarki\checkmarkr\checkmarkc\checkmark$\checkmarks\checkmark$\checkmark^\checkmark\\checkmarkc\checkmarki\checkmarkr\checkmarkc\checkmark$\checkmarki\checkmark$\checkmark^\checkmark\\checkmarkc\checkmarki\checkmarkr\checkmarkc\checkmark$\checkmarkt\checkmark$\checkmark^\checkmark\\checkmarkc\checkmarki\checkmarkr\checkmarkc\checkmark$\checkmarki\checkmark$\checkmark^\checkmark\\checkmarkc\checkmarki\checkmarkr\checkmarkc\checkmark$\checkmarko\checkmark$\checkmark^\checkmark\\checkmarkc\checkmarki\checkmarkr\checkmarkc\checkmark$\checkmarkn\checkmark$\checkmark^\checkmark\\checkmarkc\checkmarki\checkmarkr\checkmarkc\checkmark$\checkmark \checkmark$\checkmark^\checkmark\\checkmarkc\checkmarki\checkmarkr\checkmarkc\checkmark$\checkmarkf\checkmark$\checkmark^\checkmark\\checkmarkc\checkmarki\checkmarkr\checkmarkc\checkmark$\checkmarko\checkmark$\checkmark^\checkmark\\checkmarkc\checkmarki\checkmarkr\checkmarkc\checkmark$\checkmarkn\checkmark$\checkmark^\checkmark\\checkmarkc\checkmarki\checkmarkr\checkmarkc\checkmark$\checkmarkc\checkmark$\checkmark^\checkmark\\checkmarkc\checkmarki\checkmarkr\checkmarkc\checkmark$\checkmarkt\checkmark$\checkmark^\checkmark\\checkmarkc\checkmarki\checkmarkr\checkmarkc\checkmark$\checkmarki\checkmark$\checkmark^\checkmark\\checkmarkc\checkmarki\checkmarkr\checkmarkc\checkmark$\checkmarko\checkmark$\checkmark^\checkmark\\checkmarkc\checkmarki\checkmarkr\checkmarkc\checkmark$\checkmarkn\checkmark$\checkmark^\checkmark\\checkmarkc\checkmarki\checkmarkr\checkmarkc\checkmark$\checkmarkn\checkmark$\checkmark^\checkmark\\checkmarkc\checkmarki\checkmarkr\checkmarkc\checkmark$\checkmarke\checkmark$\checkmark^\checkmark\\checkmarkc\checkmarki\checkmarkr\checkmarkc\checkmark$\checkmarkl\checkmark$\checkmark^\checkmark\\checkmarkc\checkmarki\checkmarkr\checkmarkc\checkmark$\checkmarkl\checkmark$\checkmark^\checkmark\\checkmarkc\checkmarki\checkmarkr\checkmarkc\checkmark$\checkmarke\checkmark$\checkmark^\checkmark\\checkmarkc\checkmarki\checkmarkr\checkmarkc\checkmark$\checkmark \checkmark$\checkmark^\checkmark\\checkmarkc\checkmarki\checkmarkr\checkmarkc\checkmark$\checkmarke\checkmark$\checkmark^\checkmark\\checkmarkc\checkmarki\checkmarkr\checkmarkc\checkmark$\checkmarkn\checkmark$\checkmark^\checkmark\\checkmarkc\checkmarki\checkmarkr\checkmarkc\checkmark$\checkmark \checkmark$\checkmark^\checkmark\\checkmarkc\checkmarki\checkmarkr\checkmarkc\checkmark$\checkmarks\checkmark$\checkmark^\checkmark\\checkmarkc\checkmarki\checkmarkr\checkmarkc\checkmark$\checkmarko\checkmark$\checkmark^\checkmark\\checkmarkc\checkmarki\checkmarkr\checkmarkc\checkmark$\checkmarkl\checkmark$\checkmark^\checkmark\\checkmarkc\checkmarki\checkmarkr\checkmarkc\checkmark$\checkmarku\checkmark$\checkmark^\checkmark\\checkmarkc\checkmarki\checkmarkr\checkmarkc\checkmark$\checkmarkt\checkmark$\checkmark^\checkmark\\checkmarkc\checkmarki\checkmarkr\checkmarkc\checkmark$\checkmarki\checkmark$\checkmark^\checkmark\\checkmarkc\checkmarki\checkmarkr\checkmarkc\checkmark$\checkmarko\checkmark$\checkmark^\checkmark\\checkmarkc\checkmarki\checkmarkr\checkmarkc\checkmark$\checkmarkn\checkmark$\checkmark^\checkmark\\checkmarkc\checkmarki\checkmarkr\checkmarkc\checkmark$\checkmarks\checkmark$\checkmark^\checkmark\\checkmarkc\checkmarki\checkmarkr\checkmarkc\checkmark$\checkmark \checkmark$\checkmark^\checkmark\\checkmarkc\checkmarki\checkmarkr\checkmarkc\checkmark$\checkmarkc\checkmark$\checkmark^\checkmark\\checkmarkc\checkmarki\checkmarkr\checkmarkc\checkmark$\checkmarko\checkmark$\checkmark^\checkmark\\checkmarkc\checkmarki\checkmarkr\checkmarkc\checkmark$\checkmarkn\checkmark$\checkmark^\checkmark\\checkmarkc\checkmarki\checkmarkr\checkmarkc\checkmark$\checkmarks\checkmark$\checkmark^\checkmark\\checkmarkc\checkmarki\checkmarkr\checkmarkc\checkmark$\checkmarkt\checkmark$\checkmark^\checkmark\\checkmarkc\checkmarki\checkmarkr\checkmarkc\checkmark$\checkmarkr\checkmark$\checkmark^\checkmark\\checkmarkc\checkmarki\checkmarkr\checkmarkc\checkmark$\checkmarku\checkmark$\checkmark^\checkmark\\checkmarkc\checkmarki\checkmarkr\checkmarkc\checkmark$\checkmarkc\checkmark$\checkmark^\checkmark\\checkmarkc\checkmarki\checkmarkr\checkmarkc\checkmark$\checkmarkt\checkmark$\checkmark^\checkmark\\checkmarkc\checkmarki\checkmarkr\checkmarkc\checkmark$\checkmarki\checkmark$\checkmark^\checkmark\\checkmarkc\checkmarki\checkmarkr\checkmarkc\checkmark$\checkmarkv\checkmark$\checkmark^\checkmark\\checkmarkc\checkmarki\checkmarkr\checkmarkc\checkmark$\checkmarke\checkmark$\checkmark^\checkmark\\checkmarkc\checkmarki\checkmarkr\checkmarkc\checkmark$\checkmarks\checkmark$\checkmark^\checkmark\\checkmarkc\checkmarki\checkmarkr\checkmarkc\checkmark$\checkmark}\checkmark$\checkmark^\checkmark\\checkmarkc\checkmarki\checkmarkr\checkmarkc\checkmark$\checkmark
\checkmark$\checkmark^\checkmark\\checkmarkc\checkmarki\checkmarkr\checkmarkc\checkmark$\checkmark\\checkmark$\checkmark^\checkmark\\checkmarkc\checkmarki\checkmarkr\checkmarkc\checkmark$\checkmarkl\checkmark$\checkmark^\checkmark\\checkmarkc\checkmarki\checkmarkr\checkmarkc\checkmark$\checkmarka\checkmark$\checkmark^\checkmark\\checkmarkc\checkmarki\checkmarkr\checkmarkc\checkmark$\checkmarkb\checkmark$\checkmark^\checkmark\\checkmarkc\checkmarki\checkmarkr\checkmarkc\checkmark$\checkmarke\checkmark$\checkmark^\checkmark\\checkmarkc\checkmarki\checkmarkr\checkmarkc\checkmark$\checkmarkl\checkmark$\checkmark^\checkmark\\checkmarkc\checkmarki\checkmarkr\checkmarkc\checkmark$\checkmark{\checkmark$\checkmark^\checkmark\\checkmarkc\checkmarki\checkmarkr\checkmarkc\checkmark$\checkmarkt\checkmark$\checkmark^\checkmark\\checkmarkc\checkmarki\checkmarkr\checkmarkc\checkmark$\checkmarka\checkmark$\checkmark^\checkmark\\checkmarkc\checkmarki\checkmarkr\checkmarkc\checkmark$\checkmarkb\checkmark$\checkmark^\checkmark\\checkmarkc\checkmarki\checkmarkr\checkmarkc\checkmark$\checkmark:\checkmark$\checkmark^\checkmark\\checkmarkc\checkmarki\checkmarkr\checkmarkc\checkmark$\checkmarkf\checkmark$\checkmark^\checkmark\\checkmarkc\checkmarki\checkmarkr\checkmarkc\checkmark$\checkmarka\checkmark$\checkmark^\checkmark\\checkmarkc\checkmarki\checkmarkr\checkmarkc\checkmark$\checkmarks\checkmark$\checkmark^\checkmark\\checkmarkc\checkmarki\checkmarkr\checkmarkc\checkmark$\checkmarkt\checkmark$\checkmark^\checkmark\\checkmarkc\checkmarki\checkmarkr\checkmarkc\checkmark$\checkmark}\checkmark$\checkmark^\checkmark\\checkmarkc\checkmarki\checkmarkr\checkmarkc\checkmark$\checkmark
\checkmark$\checkmark^\checkmark\\checkmarkc\checkmarki\checkmarkr\checkmarkc\checkmark$\checkmark\\checkmark$\checkmark^\checkmark\\checkmarkc\checkmarki\checkmarkr\checkmarkc\checkmark$\checkmarke\checkmark$\checkmark^\checkmark\\checkmarkc\checkmarki\checkmarkr\checkmarkc\checkmark$\checkmarkn\checkmark$\checkmark^\checkmark\\checkmarkc\checkmarki\checkmarkr\checkmarkc\checkmark$\checkmarkd\checkmark$\checkmark^\checkmark\\checkmarkc\checkmarki\checkmarkr\checkmarkc\checkmark$\checkmark{\checkmark$\checkmark^\checkmark\\checkmarkc\checkmarki\checkmarkr\checkmarkc\checkmark$\checkmarkt\checkmark$\checkmark^\checkmark\\checkmarkc\checkmarki\checkmarkr\checkmarkc\checkmark$\checkmarka\checkmark$\checkmark^\checkmark\\checkmarkc\checkmarki\checkmarkr\checkmarkc\checkmark$\checkmarkb\checkmark$\checkmark^\checkmark\\checkmarkc\checkmarki\checkmarkr\checkmarkc\checkmark$\checkmarkl\checkmark$\checkmark^\checkmark\\checkmarkc\checkmarki\checkmarkr\checkmarkc\checkmark$\checkmarke\checkmark$\checkmark^\checkmark\\checkmarkc\checkmarki\checkmarkr\checkmarkc\checkmark$\checkmark}\checkmark$\checkmark^\checkmark\\checkmarkc\checkmarki\checkmarkr\checkmarkc\checkmark$\checkmark
\checkmark$\checkmark^\checkmark\\checkmarkc\checkmarki\checkmarkr\checkmarkc\checkmark$\checkmark
\checkmark$\checkmark^\checkmark\\checkmarkc\checkmarki\checkmarkr\checkmarkc\checkmark$\checkmark%\checkmark$\checkmark^\checkmark\\checkmarkc\checkmarki\checkmarkr\checkmarkc\checkmark$\checkmark \checkmark$\checkmark^\checkmark\\checkmarkc\checkmarki\checkmarkr\checkmarkc\checkmark$\checkmark=\checkmark$\checkmark^\checkmark\\checkmarkc\checkmarki\checkmarkr\checkmarkc\checkmark$\checkmark=\checkmark$\checkmark^\checkmark\\checkmarkc\checkmarki\checkmarkr\checkmarkc\checkmark$\checkmark=\checkmark$\checkmark^\checkmark\\checkmarkc\checkmarki\checkmarkr\checkmarkc\checkmark$\checkmark=\checkmark$\checkmark^\checkmark\\checkmarkc\checkmarki\checkmarkr\checkmarkc\checkmark$\checkmark=\checkmark$\checkmark^\checkmark\\checkmarkc\checkmarki\checkmarkr\checkmarkc\checkmark$\checkmark=\checkmark$\checkmark^\checkmark\\checkmarkc\checkmarki\checkmarkr\checkmarkc\checkmark$\checkmark=\checkmark$\checkmark^\checkmark\\checkmarkc\checkmarki\checkmarkr\checkmarkc\checkmark$\checkmark=\checkmark$\checkmark^\checkmark\\checkmarkc\checkmarki\checkmarkr\checkmarkc\checkmark$\checkmark=\checkmark$\checkmark^\checkmark\\checkmarkc\checkmarki\checkmarkr\checkmarkc\checkmark$\checkmark=\checkmark$\checkmark^\checkmark\\checkmarkc\checkmarki\checkmarkr\checkmarkc\checkmark$\checkmark=\checkmark$\checkmark^\checkmark\\checkmarkc\checkmarki\checkmarkr\checkmarkc\checkmark$\checkmark=\checkmark$\checkmark^\checkmark\\checkmarkc\checkmarki\checkmarkr\checkmarkc\checkmark$\checkmark=\checkmark$\checkmark^\checkmark\\checkmarkc\checkmarki\checkmarkr\checkmarkc\checkmark$\checkmark=\checkmark$\checkmark^\checkmark\\checkmarkc\checkmarki\checkmarkr\checkmarkc\checkmark$\checkmark=\checkmark$\checkmark^\checkmark\\checkmarkc\checkmarki\checkmarkr\checkmarkc\checkmark$\checkmark=\checkmark$\checkmark^\checkmark\\checkmarkc\checkmarki\checkmarkr\checkmarkc\checkmark$\checkmark=\checkmark$\checkmark^\checkmark\\checkmarkc\checkmarki\checkmarkr\checkmarkc\checkmark$\checkmark=\checkmark$\checkmark^\checkmark\\checkmarkc\checkmarki\checkmarkr\checkmarkc\checkmark$\checkmark=\checkmark$\checkmark^\checkmark\\checkmarkc\checkmarki\checkmarkr\checkmarkc\checkmark$\checkmark=\checkmark$\checkmark^\checkmark\\checkmarkc\checkmarki\checkmarkr\checkmarkc\checkmark$\checkmark=\checkmark$\checkmark^\checkmark\\checkmarkc\checkmarki\checkmarkr\checkmarkc\checkmark$\checkmark=\checkmark$\checkmark^\checkmark\\checkmarkc\checkmarki\checkmarkr\checkmarkc\checkmark$\checkmark=\checkmark$\checkmark^\checkmark\\checkmarkc\checkmarki\checkmarkr\checkmarkc\checkmark$\checkmark=\checkmark$\checkmark^\checkmark\\checkmarkc\checkmarki\checkmarkr\checkmarkc\checkmark$\checkmark=\checkmark$\checkmark^\checkmark\\checkmarkc\checkmarki\checkmarkr\checkmarkc\checkmark$\checkmark=\checkmark$\checkmark^\checkmark\\checkmarkc\checkmarki\checkmarkr\checkmarkc\checkmark$\checkmark=\checkmark$\checkmark^\checkmark\\checkmarkc\checkmarki\checkmarkr\checkmarkc\checkmark$\checkmark=\checkmark$\checkmark^\checkmark\\checkmarkc\checkmarki\checkmarkr\checkmarkc\checkmark$\checkmark=\checkmark$\checkmark^\checkmark\\checkmarkc\checkmarki\checkmarkr\checkmarkc\checkmark$\checkmark=\checkmark$\checkmark^\checkmark\\checkmarkc\checkmarki\checkmarkr\checkmarkc\checkmark$\checkmark=\checkmark$\checkmark^\checkmark\\checkmarkc\checkmarki\checkmarkr\checkmarkc\checkmark$\checkmark=\checkmark$\checkmark^\checkmark\\checkmarkc\checkmarki\checkmarkr\checkmarkc\checkmark$\checkmark=\checkmark$\checkmark^\checkmark\\checkmarkc\checkmarki\checkmarkr\checkmarkc\checkmark$\checkmark=\checkmark$\checkmark^\checkmark\\checkmarkc\checkmarki\checkmarkr\checkmarkc\checkmark$\checkmark=\checkmark$\checkmark^\checkmark\\checkmarkc\checkmarki\checkmarkr\checkmarkc\checkmark$\checkmark=\checkmark$\checkmark^\checkmark\\checkmarkc\checkmarki\checkmarkr\checkmarkc\checkmark$\checkmark=\checkmark$\checkmark^\checkmark\\checkmarkc\checkmarki\checkmarkr\checkmarkc\checkmark$\checkmark=\checkmark$\checkmark^\checkmark\\checkmarkc\checkmarki\checkmarkr\checkmarkc\checkmark$\checkmark=\checkmark$\checkmark^\checkmark\\checkmarkc\checkmarki\checkmarkr\checkmarkc\checkmark$\checkmark=\checkmark$\checkmark^\checkmark\\checkmarkc\checkmarki\checkmarkr\checkmarkc\checkmark$\checkmark=\checkmark$\checkmark^\checkmark\\checkmarkc\checkmarki\checkmarkr\checkmarkc\checkmark$\checkmark=\checkmark$\checkmark^\checkmark\\checkmarkc\checkmarki\checkmarkr\checkmarkc\checkmark$\checkmark=\checkmark$\checkmark^\checkmark\\checkmarkc\checkmarki\checkmarkr\checkmarkc\checkmark$\checkmark=\checkmark$\checkmark^\checkmark\\checkmarkc\checkmarki\checkmarkr\checkmarkc\checkmark$\checkmark=\checkmark$\checkmark^\checkmark\\checkmarkc\checkmarki\checkmarkr\checkmarkc\checkmark$\checkmark=\checkmark$\checkmark^\checkmark\\checkmarkc\checkmarki\checkmarkr\checkmarkc\checkmark$\checkmark=\checkmark$\checkmark^\checkmark\\checkmarkc\checkmarki\checkmarkr\checkmarkc\checkmark$\checkmark=\checkmark$\checkmark^\checkmark\\checkmarkc\checkmarki\checkmarkr\checkmarkc\checkmark$\checkmark=\checkmark$\checkmark^\checkmark\\checkmarkc\checkmarki\checkmarkr\checkmarkc\checkmark$\checkmark=\checkmark$\checkmark^\checkmark\\checkmarkc\checkmarki\checkmarkr\checkmarkc\checkmark$\checkmark=\checkmark$\checkmark^\checkmark\\checkmarkc\checkmarki\checkmarkr\checkmarkc\checkmark$\checkmark=\checkmark$\checkmark^\checkmark\\checkmarkc\checkmarki\checkmarkr\checkmarkc\checkmark$\checkmark=\checkmark$\checkmark^\checkmark\\checkmarkc\checkmarki\checkmarkr\checkmarkc\checkmark$\checkmark=\checkmark$\checkmark^\checkmark\\checkmarkc\checkmarki\checkmarkr\checkmarkc\checkmark$\checkmark=\checkmark$\checkmark^\checkmark\\checkmarkc\checkmarki\checkmarkr\checkmarkc\checkmark$\checkmark=\checkmark$\checkmark^\checkmark\\checkmarkc\checkmarki\checkmarkr\checkmarkc\checkmark$\checkmark=\checkmark$\checkmark^\checkmark\\checkmarkc\checkmarki\checkmarkr\checkmarkc\checkmark$\checkmark=\checkmark$\checkmark^\checkmark\\checkmarkc\checkmarki\checkmarkr\checkmarkc\checkmark$\checkmark=\checkmark$\checkmark^\checkmark\\checkmarkc\checkmarki\checkmarkr\checkmarkc\checkmark$\checkmark=\checkmark$\checkmark^\checkmark\\checkmarkc\checkmarki\checkmarkr\checkmarkc\checkmark$\checkmark=\checkmark$\checkmark^\checkmark\\checkmarkc\checkmarki\checkmarkr\checkmarkc\checkmark$\checkmark=\checkmark$\checkmark^\checkmark\\checkmarkc\checkmarki\checkmarkr\checkmarkc\checkmark$\checkmark
\checkmark$\checkmark^\checkmark\\checkmarkc\checkmarki\checkmarkr\checkmarkc\checkmark$\checkmark\\checkmark$\checkmark^\checkmark\\checkmarkc\checkmarki\checkmarkr\checkmarkc\checkmark$\checkmarks\checkmark$\checkmark^\checkmark\\checkmarkc\checkmarki\checkmarkr\checkmarkc\checkmark$\checkmarke\checkmark$\checkmark^\checkmark\\checkmarkc\checkmarki\checkmarkr\checkmarkc\checkmark$\checkmarkc\checkmark$\checkmark^\checkmark\\checkmarkc\checkmarki\checkmarkr\checkmarkc\checkmark$\checkmarkt\checkmark$\checkmark^\checkmark\\checkmarkc\checkmarki\checkmarkr\checkmarkc\checkmark$\checkmarki\checkmark$\checkmark^\checkmark\\checkmarkc\checkmarki\checkmarkr\checkmarkc\checkmark$\checkmarko\checkmark$\checkmark^\checkmark\\checkmarkc\checkmarki\checkmarkr\checkmarkc\checkmark$\checkmarkn\checkmark$\checkmark^\checkmark\\checkmarkc\checkmarki\checkmarkr\checkmarkc\checkmark$\checkmark{\checkmark$\checkmark^\checkmark\\checkmarkc\checkmarki\checkmarkr\checkmarkc\checkmark$\checkmarkA\checkmark$\checkmark^\checkmark\\checkmarkc\checkmarki\checkmarkr\checkmarkc\checkmark$\checkmarkr\checkmark$\checkmark^\checkmark\\checkmarkc\checkmarki\checkmarkr\checkmarkc\checkmark$\checkmarkc\checkmark$\checkmark^\checkmark\\checkmarkc\checkmarki\checkmarkr\checkmarkc\checkmark$\checkmarkh\checkmark$\checkmark^\checkmark\\checkmarkc\checkmarki\checkmarkr\checkmarkc\checkmark$\checkmarki\checkmark$\checkmark^\checkmark\\checkmarkc\checkmarki\checkmarkr\checkmarkc\checkmark$\checkmarkt\checkmark$\checkmark^\checkmark\\checkmarkc\checkmarki\checkmarkr\checkmarkc\checkmark$\checkmarke\checkmark$\checkmark^\checkmark\\checkmarkc\checkmarki\checkmarkr\checkmarkc\checkmark$\checkmarkc\checkmark$\checkmark^\checkmark\\checkmarkc\checkmarki\checkmarkr\checkmarkc\checkmark$\checkmarkt\checkmark$\checkmark^\checkmark\\checkmarkc\checkmarki\checkmarkr\checkmarkc\checkmark$\checkmarku\checkmark$\checkmark^\checkmark\\checkmarkc\checkmarki\checkmarkr\checkmarkc\checkmark$\checkmarkr\checkmark$\checkmark^\checkmark\\checkmarkc\checkmarki\checkmarkr\checkmarkc\checkmark$\checkmarke\checkmark$\checkmark^\checkmark\\checkmarkc\checkmarki\checkmarkr\checkmarkc\checkmark$\checkmark \checkmark$\checkmark^\checkmark\\checkmarkc\checkmarki\checkmarkr\checkmarkc\checkmark$\checkmarkR\checkmark$\checkmark^\checkmark\\checkmarkc\checkmarki\checkmarkr\checkmarkc\checkmark$\checkmarke\checkmark$\checkmark^\checkmark\\checkmarkc\checkmarki\checkmarkr\checkmarkc\checkmark$\checkmarkt\checkmark$\checkmark^\checkmark\\checkmarkc\checkmarki\checkmarkr\checkmarkc\checkmark$\checkmarke\checkmark$\checkmark^\checkmark\\checkmarkc\checkmarki\checkmarkr\checkmarkc\checkmark$\checkmarkn\checkmark$\checkmark^\checkmark\\checkmarkc\checkmarki\checkmarkr\checkmarkc\checkmark$\checkmarku\checkmark$\checkmark^\checkmark\\checkmarkc\checkmarki\checkmarkr\checkmarkc\checkmark$\checkmarke\checkmark$\checkmark^\checkmark\\checkmarkc\checkmarki\checkmarkr\checkmarkc\checkmark$\checkmark \checkmark$\checkmark^\checkmark\\checkmarkc\checkmarki\checkmarkr\checkmarkc\checkmark$\checkmark:\checkmark$\checkmark^\checkmark\\checkmarkc\checkmarki\checkmarkr\checkmarkc\checkmark$\checkmark \checkmark$\checkmark^\checkmark\\checkmarkc\checkmarki\checkmarkr\checkmarkc\checkmark$\checkmarkC\checkmark$\checkmark^\checkmark\\checkmarkc\checkmarki\checkmarkr\checkmarkc\checkmark$\checkmarko\checkmark$\checkmark^\checkmark\\checkmarkc\checkmarki\checkmarkr\checkmarkc\checkmark$\checkmarkn\checkmark$\checkmark^\checkmark\\checkmarkc\checkmarki\checkmarkr\checkmarkc\checkmark$\checkmarkf\checkmark$\checkmark^\checkmark\\checkmarkc\checkmarki\checkmarkr\checkmarkc\checkmark$\checkmarki\checkmark$\checkmark^\checkmark\\checkmarkc\checkmarki\checkmarkr\checkmarkc\checkmark$\checkmarkg\checkmark$\checkmark^\checkmark\\checkmarkc\checkmarki\checkmarkr\checkmarkc\checkmark$\checkmarku\checkmark$\checkmark^\checkmark\\checkmarkc\checkmarki\checkmarkr\checkmarkc\checkmark$\checkmarkr\checkmark$\checkmark^\checkmark\\checkmarkc\checkmarki\checkmarkr\checkmarkc\checkmark$\checkmarka\checkmark$\checkmark^\checkmark\\checkmarkc\checkmarki\checkmarkr\checkmarkc\checkmark$\checkmarkt\checkmark$\checkmark^\checkmark\\checkmarkc\checkmarki\checkmarkr\checkmarkc\checkmark$\checkmarki\checkmark$\checkmark^\checkmark\\checkmarkc\checkmarki\checkmarkr\checkmarkc\checkmark$\checkmarko\checkmark$\checkmark^\checkmark\\checkmarkc\checkmarki\checkmarkr\checkmarkc\checkmark$\checkmarkn\checkmark$\checkmark^\checkmark\\checkmarkc\checkmarki\checkmarkr\checkmarkc\checkmark$\checkmark \checkmark$\checkmark^\checkmark\\checkmarkc\checkmarki\checkmarkr\checkmarkc\checkmark$\checkmarkT\checkmark$\checkmark^\checkmark\\checkmarkc\checkmarki\checkmarkr\checkmarkc\checkmark$\checkmarka\checkmark$\checkmark^\checkmark\\checkmarkc\checkmarki\checkmarkr\checkmarkc\checkmark$\checkmarkb\checkmark$\checkmark^\checkmark\\checkmarkc\checkmarki\checkmarkr\checkmarkc\checkmark$\checkmarkl\checkmark$\checkmark^\checkmark\\checkmarkc\checkmarki\checkmarkr\checkmarkc\checkmark$\checkmarke\checkmark$\checkmark^\checkmark\\checkmarkc\checkmarki\checkmarkr\checkmarkc\checkmark$\checkmark-\checkmark$\checkmark^\checkmark\\checkmarkc\checkmarki\checkmarkr\checkmarkc\checkmark$\checkmarkT\checkmark$\checkmark^\checkmark\\checkmarkc\checkmarki\checkmarkr\checkmarkc\checkmark$\checkmarka\checkmark$\checkmark^\checkmark\\checkmarkc\checkmarki\checkmarkr\checkmarkc\checkmark$\checkmarkb\checkmark$\checkmark^\checkmark\\checkmarkc\checkmarki\checkmarkr\checkmarkc\checkmark$\checkmarkl\checkmark$\checkmark^\checkmark\\checkmarkc\checkmarki\checkmarkr\checkmarkc\checkmark$\checkmarke\checkmark$\checkmark^\checkmark\\checkmarkc\checkmarki\checkmarkr\checkmarkc\checkmark$\checkmark \checkmark$\checkmark^\checkmark\\checkmarkc\checkmarki\checkmarkr\checkmarkc\checkmark$\checkmark(\checkmark$\checkmark^\checkmark\\checkmarkc\checkmarki\checkmarkr\checkmarkc\checkmark$\checkmarkT\checkmark$\checkmark^\checkmark\\checkmarkc\checkmarki\checkmarkr\checkmarkc\checkmark$\checkmarkr\checkmark$\checkmark^\checkmark\\checkmarkc\checkmarki\checkmarkr\checkmarkc\checkmark$\checkmarku\checkmark$\checkmark^\checkmark\\checkmarkc\checkmarki\checkmarkr\checkmarkc\checkmark$\checkmarkn\checkmark$\checkmark^\checkmark\\checkmarkc\checkmarki\checkmarkr\checkmarkc\checkmark$\checkmarkn\checkmark$\checkmark^\checkmark\\checkmarkc\checkmarki\checkmarkr\checkmarkc\checkmark$\checkmarki\checkmark$\checkmark^\checkmark\\checkmarkc\checkmarki\checkmarkr\checkmarkc\checkmark$\checkmarko\checkmark$\checkmark^\checkmark\\checkmarkc\checkmarki\checkmarkr\checkmarkc\checkmark$\checkmarkn\checkmark$\checkmark^\checkmark\\checkmarkc\checkmarki\checkmarkr\checkmarkc\checkmark$\checkmark)\checkmark$\checkmark^\checkmark\\checkmarkc\checkmarki\checkmarkr\checkmarkc\checkmark$\checkmark}\checkmark$\checkmark^\checkmark\\checkmarkc\checkmarki\checkmarkr\checkmarkc\checkmark$\checkmark
\checkmark$\checkmark^\checkmark\\checkmarkc\checkmarki\checkmarkr\checkmarkc\checkmark$\checkmark
\checkmark$\checkmark^\checkmark\\checkmarkc\checkmarki\checkmarkr\checkmarkc\checkmark$\checkmark\\checkmark$\checkmark^\checkmark\\checkmarkc\checkmarki\checkmarkr\checkmarkc\checkmark$\checkmarks\checkmark$\checkmark^\checkmark\\checkmarkc\checkmarki\checkmarkr\checkmarkc\checkmark$\checkmarku\checkmark$\checkmark^\checkmark\\checkmarkc\checkmarki\checkmarkr\checkmarkc\checkmark$\checkmarkb\checkmark$\checkmark^\checkmark\\checkmarkc\checkmarki\checkmarkr\checkmarkc\checkmark$\checkmarks\checkmark$\checkmark^\checkmark\\checkmarkc\checkmarki\checkmarkr\checkmarkc\checkmark$\checkmarke\checkmark$\checkmark^\checkmark\\checkmarkc\checkmarki\checkmarkr\checkmarkc\checkmark$\checkmarkc\checkmark$\checkmark^\checkmark\\checkmarkc\checkmarki\checkmarkr\checkmarkc\checkmark$\checkmarkt\checkmark$\checkmark^\checkmark\\checkmarkc\checkmarki\checkmarkr\checkmarkc\checkmark$\checkmarki\checkmark$\checkmark^\checkmark\\checkmarkc\checkmarki\checkmarkr\checkmarkc\checkmark$\checkmarko\checkmark$\checkmark^\checkmark\\checkmarkc\checkmarki\checkmarkr\checkmarkc\checkmark$\checkmarkn\checkmark$\checkmark^\checkmark\\checkmarkc\checkmarki\checkmarkr\checkmarkc\checkmark$\checkmark{\checkmark$\checkmark^\checkmark\\checkmarkc\checkmarki\checkmarkr\checkmarkc\checkmark$\checkmarkD\checkmark$\checkmark^\checkmark\\checkmarkc\checkmarki\checkmarkr\checkmarkc\checkmark$\checkmarke\checkmark$\checkmark^\checkmark\\checkmarkc\checkmarki\checkmarkr\checkmarkc\checkmark$\checkmarks\checkmark$\checkmark^\checkmark\\checkmarkc\checkmarki\checkmarkr\checkmarkc\checkmark$\checkmarkc\checkmark$\checkmark^\checkmark\\checkmarkc\checkmarki\checkmarkr\checkmarkc\checkmark$\checkmarkr\checkmark$\checkmark^\checkmark\\checkmarkc\checkmarki\checkmarkr\checkmarkc\checkmark$\checkmarki\checkmark$\checkmark^\checkmark\\checkmarkc\checkmarki\checkmarkr\checkmarkc\checkmark$\checkmarkp\checkmark$\checkmark^\checkmark\\checkmarkc\checkmarki\checkmarkr\checkmarkc\checkmark$\checkmarkt\checkmark$\checkmark^\checkmark\\checkmarkc\checkmarki\checkmarkr\checkmarkc\checkmark$\checkmarki\checkmark$\checkmark^\checkmark\\checkmarkc\checkmarki\checkmarkr\checkmarkc\checkmark$\checkmarko\checkmark$\checkmark^\checkmark\\checkmarkc\checkmarki\checkmarkr\checkmarkc\checkmark$\checkmarkn\checkmark$\checkmark^\checkmark\\checkmarkc\checkmarki\checkmarkr\checkmarkc\checkmark$\checkmark \checkmark$\checkmark^\checkmark\\checkmarkc\checkmarki\checkmarkr\checkmarkc\checkmark$\checkmarkd\checkmark$\checkmark^\checkmark\\checkmarkc\checkmarki\checkmarkr\checkmarkc\checkmark$\checkmarke\checkmark$\checkmark^\checkmark\\checkmarkc\checkmarki\checkmarkr\checkmarkc\checkmark$\checkmark \checkmark$\checkmark^\checkmark\\checkmarkc\checkmarki\checkmarkr\checkmarkc\checkmark$\checkmarkl\checkmark$\checkmark^\checkmark\\checkmarkc\checkmarki\checkmarkr\checkmarkc\checkmark$\checkmark'\checkmark$\checkmark^\checkmark\\checkmarkc\checkmarki\checkmarkr\checkmarkc\checkmark$\checkmarka\checkmark$\checkmark^\checkmark\\checkmarkc\checkmarki\checkmarkr\checkmarkc\checkmark$\checkmarkr\checkmark$\checkmark^\checkmark\\checkmarkc\checkmarki\checkmarkr\checkmarkc\checkmark$\checkmarkc\checkmark$\checkmark^\checkmark\\checkmarkc\checkmarki\checkmarkr\checkmarkc\checkmark$\checkmarkh\checkmark$\checkmark^\checkmark\\checkmarkc\checkmarki\checkmarkr\checkmarkc\checkmark$\checkmarki\checkmark$\checkmark^\checkmark\\checkmarkc\checkmarki\checkmarkr\checkmarkc\checkmark$\checkmarkt\checkmark$\checkmark^\checkmark\\checkmarkc\checkmarki\checkmarkr\checkmarkc\checkmark$\checkmarke\checkmark$\checkmark^\checkmark\\checkmarkc\checkmarki\checkmarkr\checkmarkc\checkmark$\checkmarkc\checkmark$\checkmark^\checkmark\\checkmarkc\checkmarki\checkmarkr\checkmarkc\checkmark$\checkmarkt\checkmark$\checkmark^\checkmark\\checkmarkc\checkmarki\checkmarkr\checkmarkc\checkmark$\checkmarku\checkmark$\checkmark^\checkmark\\checkmarkc\checkmarki\checkmarkr\checkmarkc\checkmark$\checkmarkr\checkmark$\checkmark^\checkmark\\checkmarkc\checkmarki\checkmarkr\checkmarkc\checkmark$\checkmarke\checkmark$\checkmark^\checkmark\\checkmarkc\checkmarki\checkmarkr\checkmarkc\checkmark$\checkmark}\checkmark$\checkmark^\checkmark\\checkmarkc\checkmarki\checkmarkr\checkmarkc\checkmark$\checkmark
\checkmark$\checkmark^\checkmark\\checkmarkc\checkmarki\checkmarkr\checkmarkc\checkmark$\checkmarkL\checkmark$\checkmark^\checkmark\\checkmarkc\checkmarki\checkmarkr\checkmarkc\checkmark$\checkmark'\checkmark$\checkmark^\checkmark\\checkmarkc\checkmarki\checkmarkr\checkmarkc\checkmark$\checkmarka\checkmark$\checkmark^\checkmark\\checkmarkc\checkmarki\checkmarkr\checkmarkc\checkmark$\checkmarkr\checkmark$\checkmark^\checkmark\\checkmarkc\checkmarki\checkmarkr\checkmarkc\checkmark$\checkmarkc\checkmark$\checkmark^\checkmark\\checkmarkc\checkmarki\checkmarkr\checkmarkc\checkmark$\checkmarkh\checkmark$\checkmark^\checkmark\\checkmarkc\checkmarki\checkmarkr\checkmarkc\checkmark$\checkmarki\checkmark$\checkmark^\checkmark\\checkmarkc\checkmarki\checkmarkr\checkmarkc\checkmark$\checkmarkt\checkmark$\checkmark^\checkmark\\checkmarkc\checkmarki\checkmarkr\checkmarkc\checkmark$\checkmarke\checkmark$\checkmark^\checkmark\\checkmarkc\checkmarki\checkmarkr\checkmarkc\checkmark$\checkmarkc\checkmark$\checkmark^\checkmark\\checkmarkc\checkmarki\checkmarkr\checkmarkc\checkmark$\checkmarkt\checkmark$\checkmark^\checkmark\\checkmarkc\checkmarki\checkmarkr\checkmarkc\checkmark$\checkmarku\checkmark$\checkmark^\checkmark\\checkmarkc\checkmarki\checkmarkr\checkmarkc\checkmark$\checkmarkr\checkmark$\checkmark^\checkmark\\checkmarkc\checkmarki\checkmarkr\checkmarkc\checkmark$\checkmarke\checkmark$\checkmark^\checkmark\\checkmarkc\checkmarki\checkmarkr\checkmarkc\checkmark$\checkmark \checkmark$\checkmark^\checkmark\\checkmarkc\checkmarki\checkmarkr\checkmarkc\checkmark$\checkmark\\checkmark$\checkmark^\checkmark\\checkmarkc\checkmarki\checkmarkr\checkmarkc\checkmark$\checkmarkt\checkmark$\checkmark^\checkmark\\checkmarkc\checkmarki\checkmarkr\checkmarkc\checkmark$\checkmarke\checkmark$\checkmark^\checkmark\\checkmarkc\checkmarki\checkmarkr\checkmarkc\checkmark$\checkmarkx\checkmark$\checkmark^\checkmark\\checkmarkc\checkmarki\checkmarkr\checkmarkc\checkmark$\checkmarkt\checkmark$\checkmark^\checkmark\\checkmarkc\checkmarki\checkmarkr\checkmarkc\checkmark$\checkmarkb\checkmark$\checkmark^\checkmark\\checkmarkc\checkmarki\checkmarkr\checkmarkc\checkmark$\checkmarkf\checkmark$\checkmark^\checkmark\\checkmarkc\checkmarki\checkmarkr\checkmarkc\checkmark$\checkmark{\checkmark$\checkmark^\checkmark\\checkmarkc\checkmarki\checkmarkr\checkmarkc\checkmark$\checkmark«\checkmark$\checkmark^\checkmark\\checkmarkc\checkmarki\checkmarkr\checkmarkc\checkmark$\checkmarkT\checkmark$\checkmark^\checkmark\\checkmarkc\checkmarki\checkmarkr\checkmarkc\checkmark$\checkmarka\checkmark$\checkmark^\checkmark\\checkmarkc\checkmarki\checkmarkr\checkmarkc\checkmark$\checkmarkb\checkmark$\checkmark^\checkmark\\checkmarkc\checkmarki\checkmarkr\checkmarkc\checkmark$\checkmarkl\checkmark$\checkmark^\checkmark\\checkmarkc\checkmarki\checkmarkr\checkmarkc\checkmark$\checkmarke\checkmark$\checkmark^\checkmark\\checkmarkc\checkmarki\checkmarkr\checkmarkc\checkmark$\checkmark-\checkmark$\checkmark^\checkmark\\checkmarkc\checkmarki\checkmarkr\checkmarkc\checkmark$\checkmarkT\checkmark$\checkmark^\checkmark\\checkmarkc\checkmarki\checkmarkr\checkmarkc\checkmark$\checkmarka\checkmark$\checkmark^\checkmark\\checkmarkc\checkmarki\checkmarkr\checkmarkc\checkmark$\checkmarkb\checkmark$\checkmark^\checkmark\\checkmarkc\checkmarki\checkmarkr\checkmarkc\checkmark$\checkmarkl\checkmark$\checkmark^\checkmark\\checkmarkc\checkmarki\checkmarkr\checkmarkc\checkmark$\checkmarke\checkmark$\checkmark^\checkmark\\checkmarkc\checkmarki\checkmarkr\checkmarkc\checkmark$\checkmark»\checkmark$\checkmark^\checkmark\\checkmarkc\checkmarki\checkmarkr\checkmarkc\checkmark$\checkmark}\checkmark$\checkmark^\checkmark\\checkmarkc\checkmarki\checkmarkr\checkmarkc\checkmark$\checkmark \checkmark$\checkmark^\checkmark\\checkmarkc\checkmarki\checkmarkr\checkmarkc\checkmark$\checkmark(\checkmark$\checkmark^\checkmark\\checkmarkc\checkmarki\checkmarkr\checkmarkc\checkmark$\checkmarko\checkmark$\checkmark^\checkmark\\checkmarkc\checkmarki\checkmarkr\checkmarkc\checkmark$\checkmarku\checkmark$\checkmark^\checkmark\\checkmarkc\checkmarki\checkmarkr\checkmarkc\checkmark$\checkmark \checkmark$\checkmark^\checkmark\\checkmarkc\checkmarki\checkmarkr\checkmarkc\checkmark$\checkmarkT\checkmark$\checkmark^\checkmark\\checkmarkc\checkmarki\checkmarkr\checkmarkc\checkmark$\checkmarkr\checkmark$\checkmark^\checkmark\\checkmarkc\checkmarki\checkmarkr\checkmarkc\checkmark$\checkmarku\checkmark$\checkmark^\checkmark\\checkmarkc\checkmarki\checkmarkr\checkmarkc\checkmark$\checkmarkn\checkmark$\checkmark^\checkmark\\checkmarkc\checkmarki\checkmarkr\checkmarkc\checkmark$\checkmarkn\checkmark$\checkmark^\checkmark\\checkmarkc\checkmarki\checkmarkr\checkmarkc\checkmark$\checkmarki\checkmark$\checkmark^\checkmark\\checkmarkc\checkmarki\checkmarkr\checkmarkc\checkmark$\checkmarko\checkmark$\checkmark^\checkmark\\checkmarkc\checkmarki\checkmarkr\checkmarkc\checkmark$\checkmarkn\checkmark$\checkmark^\checkmark\\checkmarkc\checkmarki\checkmarkr\checkmarkc\checkmark$\checkmark)\checkmark$\checkmark^\checkmark\\checkmarkc\checkmarki\checkmarkr\checkmarkc\checkmark$\checkmark \checkmark$\checkmark^\checkmark\\checkmarkc\checkmarki\checkmarkr\checkmarkc\checkmark$\checkmarke\checkmark$\checkmark^\checkmark\\checkmarkc\checkmarki\checkmarkr\checkmarkc\checkmark$\checkmarks\checkmark$\checkmark^\checkmark\\checkmarkc\checkmarki\checkmarkr\checkmarkc\checkmark$\checkmarkt\checkmark$\checkmark^\checkmark\\checkmarkc\checkmarki\checkmarkr\checkmarkc\checkmark$\checkmark \checkmark$\checkmark^\checkmark\\checkmarkc\checkmarki\checkmarkr\checkmarkc\checkmark$\checkmarkl\checkmark$\checkmark^\checkmark\\checkmarkc\checkmarki\checkmarkr\checkmarkc\checkmark$\checkmark'\checkmark$\checkmark^\checkmark\\checkmarkc\checkmarki\checkmarkr\checkmarkc\checkmark$\checkmarka\checkmark$\checkmark^\checkmark\\checkmarkc\checkmarki\checkmarkr\checkmarkc\checkmark$\checkmarkr\checkmark$\checkmark^\checkmark\\checkmarkc\checkmarki\checkmarkr\checkmarkc\checkmark$\checkmarkc\checkmark$\checkmark^\checkmark\\checkmarkc\checkmarki\checkmarkr\checkmarkc\checkmark$\checkmarkh\checkmark$\checkmark^\checkmark\\checkmarkc\checkmarki\checkmarkr\checkmarkc\checkmark$\checkmarki\checkmark$\checkmark^\checkmark\\checkmarkc\checkmarki\checkmarkr\checkmarkc\checkmark$\checkmarkt\checkmark$\checkmark^\checkmark\\checkmarkc\checkmarki\checkmarkr\checkmarkc\checkmark$\checkmarke\checkmark$\checkmark^\checkmark\\checkmarkc\checkmarki\checkmarkr\checkmarkc\checkmark$\checkmarkc\checkmark$\checkmark^\checkmark\\checkmarkc\checkmarki\checkmarkr\checkmarkc\checkmark$\checkmarkt\checkmark$\checkmark^\checkmark\\checkmarkc\checkmarki\checkmarkr\checkmarkc\checkmark$\checkmarku\checkmark$\checkmark^\checkmark\\checkmarkc\checkmarki\checkmarkr\checkmarkc\checkmark$\checkmarkr\checkmark$\checkmark^\checkmark\\checkmarkc\checkmarki\checkmarkr\checkmarkc\checkmark$\checkmarke\checkmark$\checkmark^\checkmark\\checkmarkc\checkmarki\checkmarkr\checkmarkc\checkmark$\checkmark \checkmark$\checkmark^\checkmark\\checkmarkc\checkmarki\checkmarkr\checkmarkc\checkmark$\checkmarkd\checkmark$\checkmark^\checkmark\\checkmarkc\checkmarki\checkmarkr\checkmarkc\checkmark$\checkmarka\checkmark$\checkmark^\checkmark\\checkmarkc\checkmarki\checkmarkr\checkmarkc\checkmark$\checkmarkn\checkmark$\checkmark^\checkmark\\checkmarkc\checkmarki\checkmarkr\checkmarkc\checkmark$\checkmarks\checkmark$\checkmark^\checkmark\\checkmarkc\checkmarki\checkmarkr\checkmarkc\checkmark$\checkmark \checkmark$\checkmark^\checkmark\\checkmarkc\checkmarki\checkmarkr\checkmarkc\checkmark$\checkmarkl\checkmark$\checkmark^\checkmark\\checkmarkc\checkmarki\checkmarkr\checkmarkc\checkmark$\checkmarka\checkmark$\checkmark^\checkmark\\checkmarkc\checkmarki\checkmarkr\checkmarkc\checkmark$\checkmarkq\checkmark$\checkmark^\checkmark\\checkmarkc\checkmarki\checkmarkr\checkmarkc\checkmark$\checkmarku\checkmark$\checkmark^\checkmark\\checkmarkc\checkmarki\checkmarkr\checkmarkc\checkmark$\checkmarke\checkmark$\checkmark^\checkmark\\checkmarkc\checkmarki\checkmarkr\checkmarkc\checkmark$\checkmarkl\checkmark$\checkmark^\checkmark\\checkmarkc\checkmarki\checkmarkr\checkmarkc\checkmark$\checkmarkl\checkmark$\checkmark^\checkmark\\checkmarkc\checkmarki\checkmarkr\checkmarkc\checkmark$\checkmarke\checkmark$\checkmark^\checkmark\\checkmarkc\checkmarki\checkmarkr\checkmarkc\checkmark$\checkmark \checkmark$\checkmark^\checkmark\\checkmarkc\checkmarki\checkmarkr\checkmarkc\checkmark$\checkmarkl\checkmark$\checkmark^\checkmark\\checkmarkc\checkmarki\checkmarkr\checkmarkc\checkmark$\checkmarke\checkmark$\checkmark^\checkmark\\checkmarkc\checkmarki\checkmarkr\checkmarkc\checkmark$\checkmarks\checkmark$\checkmark^\checkmark\\checkmarkc\checkmarki\checkmarkr\checkmarkc\checkmark$\checkmark \checkmark$\checkmark^\checkmark\\checkmarkc\checkmarki\checkmarkr\checkmarkc\checkmark$\checkmarkd\checkmark$\checkmark^\checkmark\\checkmarkc\checkmarki\checkmarkr\checkmarkc\checkmark$\checkmarke\checkmark$\checkmark^\checkmark\\checkmarkc\checkmarki\checkmarkr\checkmarkc\checkmark$\checkmarku\checkmark$\checkmark^\checkmark\\checkmarkc\checkmarki\checkmarkr\checkmarkc\checkmark$\checkmarkx\checkmark$\checkmark^\checkmark\\checkmarkc\checkmarki\checkmarkr\checkmarkc\checkmark$\checkmark \checkmark$\checkmark^\checkmark\\checkmarkc\checkmarki\checkmarkr\checkmarkc\checkmark$\checkmarka\checkmark$\checkmark^\checkmark\\checkmarkc\checkmarki\checkmarkr\checkmarkc\checkmark$\checkmarkx\checkmark$\checkmark^\checkmark\\checkmarkc\checkmarki\checkmarkr\checkmarkc\checkmark$\checkmarke\checkmark$\checkmark^\checkmark\\checkmarkc\checkmarki\checkmarkr\checkmarkc\checkmark$\checkmarks\checkmark$\checkmark^\checkmark\\checkmarkc\checkmarki\checkmarkr\checkmarkc\checkmark$\checkmark \checkmark$\checkmark^\checkmark\\checkmarkc\checkmarki\checkmarkr\checkmarkc\checkmark$\checkmarkd\checkmark$\checkmark^\checkmark\\checkmarkc\checkmarki\checkmarkr\checkmarkc\checkmark$\checkmarke\checkmark$\checkmark^\checkmark\\checkmarkc\checkmarki\checkmarkr\checkmarkc\checkmark$\checkmark \checkmark$\checkmark^\checkmark\\checkmarkc\checkmarki\checkmarkr\checkmarkc\checkmark$\checkmarkr\checkmark$\checkmark^\checkmark\\checkmarkc\checkmarki\checkmarkr\checkmarkc\checkmark$\checkmarko\checkmark$\checkmark^\checkmark\\checkmarkc\checkmarki\checkmarkr\checkmarkc\checkmark$\checkmarkt\checkmark$\checkmark^\checkmark\\checkmarkc\checkmarki\checkmarkr\checkmarkc\checkmark$\checkmarka\checkmark$\checkmark^\checkmark\\checkmarkc\checkmarki\checkmarkr\checkmarkc\checkmark$\checkmarkt\checkmark$\checkmark^\checkmark\\checkmarkc\checkmarki\checkmarkr\checkmarkc\checkmark$\checkmarki\checkmark$\checkmark^\checkmark\\checkmarkc\checkmarki\checkmarkr\checkmarkc\checkmark$\checkmarko\checkmark$\checkmark^\checkmark\\checkmarkc\checkmarki\checkmarkr\checkmarkc\checkmark$\checkmarkn\checkmark$\checkmark^\checkmark\\checkmarkc\checkmarki\checkmarkr\checkmarkc\checkmark$\checkmark \checkmark$\checkmark^\checkmark\\checkmarkc\checkmarki\checkmarkr\checkmarkc\checkmark$\checkmark(\checkmark$\checkmark^\checkmark\\checkmarkc\checkmarki\checkmarkr\checkmarkc\checkmark$\checkmarkA\checkmark$\checkmark^\checkmark\\checkmarkc\checkmarki\checkmarkr\checkmarkc\checkmark$\checkmark \checkmark$\checkmark^\checkmark\\checkmarkc\checkmarki\checkmarkr\checkmarkc\checkmark$\checkmarke\checkmark$\checkmark^\checkmark\\checkmarkc\checkmarki\checkmarkr\checkmarkc\checkmark$\checkmarkt\checkmark$\checkmark^\checkmark\\checkmarkc\checkmarki\checkmarkr\checkmarkc\checkmark$\checkmark \checkmark$\checkmark^\checkmark\\checkmarkc\checkmarki\checkmarkr\checkmarkc\checkmark$\checkmarkB\checkmark$\checkmark^\checkmark\\checkmarkc\checkmarki\checkmarkr\checkmarkc\checkmark$\checkmark)\checkmark$\checkmark^\checkmark\\checkmarkc\checkmarki\checkmarkr\checkmarkc\checkmark$\checkmark \checkmark$\checkmark^\checkmark\\checkmarkc\checkmarki\checkmarkr\checkmarkc\checkmark$\checkmarkp\checkmark$\checkmark^\checkmark\\checkmarkc\checkmarki\checkmarkr\checkmarkc\checkmark$\checkmarko\checkmark$\checkmark^\checkmark\\checkmarkc\checkmarki\checkmarkr\checkmarkc\checkmark$\checkmarkr\checkmark$\checkmark^\checkmark\\checkmarkc\checkmarki\checkmarkr\checkmarkc\checkmark$\checkmarkt\checkmark$\checkmark^\checkmark\\checkmarkc\checkmarki\checkmarkr\checkmarkc\checkmark$\checkmarke\checkmark$\checkmark^\checkmark\\checkmarkc\checkmarki\checkmarkr\checkmarkc\checkmark$\checkmarkn\checkmark$\checkmark^\checkmark\\checkmarkc\checkmarki\checkmarkr\checkmarkc\checkmark$\checkmarkt\checkmark$\checkmark^\checkmark\\checkmarkc\checkmarki\checkmarkr\checkmarkc\checkmark$\checkmark \checkmark$\checkmark^\checkmark\\checkmarkc\checkmarki\checkmarkr\checkmarkc\checkmark$\checkmarkl\checkmark$\checkmark^\checkmark\\checkmarkc\checkmarki\checkmarkr\checkmarkc\checkmark$\checkmarka\checkmark$\checkmark^\checkmark\\checkmarkc\checkmarki\checkmarkr\checkmarkc\checkmark$\checkmark \checkmark$\checkmark^\checkmark\\checkmarkc\checkmarki\checkmarkr\checkmarkc\checkmark$\checkmarkp\checkmark$\checkmark^\checkmark\\checkmarkc\checkmarki\checkmarkr\checkmarkc\checkmark$\checkmarki\checkmark$\checkmark^\checkmark\\checkmarkc\checkmarki\checkmarkr\checkmarkc\checkmark$\checkmarkè\checkmark$\checkmark^\checkmark\\checkmarkc\checkmarki\checkmarkr\checkmarkc\checkmark$\checkmarkc\checkmark$\checkmark^\checkmark\\checkmarkc\checkmarki\checkmarkr\checkmarkc\checkmark$\checkmarke\checkmark$\checkmark^\checkmark\\checkmarkc\checkmarki\checkmarkr\checkmarkc\checkmark$\checkmark,\checkmark$\checkmark^\checkmark\\checkmarkc\checkmarki\checkmarkr\checkmarkc\checkmark$\checkmark \checkmark$\checkmark^\checkmark\\checkmarkc\checkmarki\checkmarkr\checkmarkc\checkmark$\checkmarkt\checkmark$\checkmark^\checkmark\\checkmarkc\checkmarki\checkmarkr\checkmarkc\checkmark$\checkmarka\checkmark$\checkmark^\checkmark\\checkmarkc\checkmarki\checkmarkr\checkmarkc\checkmark$\checkmarkn\checkmark$\checkmark^\checkmark\\checkmarkc\checkmarki\checkmarkr\checkmarkc\checkmark$\checkmarkd\checkmark$\checkmark^\checkmark\\checkmarkc\checkmarki\checkmarkr\checkmarkc\checkmark$\checkmarki\checkmark$\checkmark^\checkmark\\checkmarkc\checkmarki\checkmarkr\checkmarkc\checkmark$\checkmarks\checkmark$\checkmark^\checkmark\\checkmarkc\checkmarki\checkmarkr\checkmarkc\checkmark$\checkmark \checkmark$\checkmark^\checkmark\\checkmarkc\checkmarki\checkmarkr\checkmarkc\checkmark$\checkmarkq\checkmark$\checkmark^\checkmark\\checkmarkc\checkmarki\checkmarkr\checkmarkc\checkmark$\checkmarku\checkmark$\checkmark^\checkmark\\checkmarkc\checkmarki\checkmarkr\checkmarkc\checkmark$\checkmarke\checkmark$\checkmark^\checkmark\\checkmarkc\checkmarki\checkmarkr\checkmarkc\checkmark$\checkmark \checkmark$\checkmark^\checkmark\\checkmarkc\checkmarki\checkmarkr\checkmarkc\checkmark$\checkmarkl\checkmark$\checkmark^\checkmark\\checkmarkc\checkmarki\checkmarkr\checkmarkc\checkmark$\checkmarke\checkmark$\checkmark^\checkmark\\checkmarkc\checkmarki\checkmarkr\checkmarkc\checkmark$\checkmarks\checkmark$\checkmark^\checkmark\\checkmarkc\checkmarki\checkmarkr\checkmarkc\checkmark$\checkmark \checkmark$\checkmark^\checkmark\\checkmarkc\checkmarki\checkmarkr\checkmarkc\checkmark$\checkmarkt\checkmark$\checkmark^\checkmark\\checkmarkc\checkmarki\checkmarkr\checkmarkc\checkmark$\checkmarkr\checkmark$\checkmark^\checkmark\\checkmarkc\checkmarki\checkmarkr\checkmarkc\checkmark$\checkmarko\checkmark$\checkmark^\checkmark\\checkmarkc\checkmarki\checkmarkr\checkmarkc\checkmark$\checkmarki\checkmark$\checkmark^\checkmark\\checkmarkc\checkmarki\checkmarkr\checkmarkc\checkmark$\checkmarks\checkmark$\checkmark^\checkmark\\checkmarkc\checkmarki\checkmarkr\checkmarkc\checkmark$\checkmark \checkmark$\checkmark^\checkmark\\checkmarkc\checkmarki\checkmarkr\checkmarkc\checkmark$\checkmarka\checkmark$\checkmark^\checkmark\\checkmarkc\checkmarki\checkmarkr\checkmarkc\checkmark$\checkmarkx\checkmark$\checkmark^\checkmark\\checkmarkc\checkmarki\checkmarkr\checkmarkc\checkmark$\checkmarke\checkmark$\checkmark^\checkmark\\checkmarkc\checkmarki\checkmarkr\checkmarkc\checkmark$\checkmarks\checkmark$\checkmark^\checkmark\\checkmarkc\checkmarki\checkmarkr\checkmarkc\checkmark$\checkmark \checkmark$\checkmark^\checkmark\\checkmarkc\checkmarki\checkmarkr\checkmarkc\checkmark$\checkmarkd\checkmark$\checkmark^\checkmark\\checkmarkc\checkmarki\checkmarkr\checkmarkc\checkmark$\checkmarke\checkmark$\checkmark^\checkmark\\checkmarkc\checkmarki\checkmarkr\checkmarkc\checkmark$\checkmark \checkmark$\checkmark^\checkmark\\checkmarkc\checkmarki\checkmarkr\checkmarkc\checkmark$\checkmarkt\checkmark$\checkmark^\checkmark\\checkmarkc\checkmarki\checkmarkr\checkmarkc\checkmark$\checkmarkr\checkmark$\checkmark^\checkmark\\checkmarkc\checkmarki\checkmarkr\checkmarkc\checkmark$\checkmarka\checkmark$\checkmark^\checkmark\\checkmarkc\checkmarki\checkmarkr\checkmarkc\checkmark$\checkmarkn\checkmark$\checkmark^\checkmark\\checkmarkc\checkmarki\checkmarkr\checkmarkc\checkmark$\checkmarks\checkmark$\checkmark^\checkmark\\checkmarkc\checkmarki\checkmarkr\checkmarkc\checkmark$\checkmarkl\checkmark$\checkmark^\checkmark\\checkmarkc\checkmarki\checkmarkr\checkmarkc\checkmark$\checkmarka\checkmark$\checkmark^\checkmark\\checkmarkc\checkmarki\checkmarkr\checkmarkc\checkmark$\checkmarkt\checkmark$\checkmark^\checkmark\\checkmarkc\checkmarki\checkmarkr\checkmarkc\checkmark$\checkmarki\checkmark$\checkmark^\checkmark\\checkmarkc\checkmarki\checkmarkr\checkmarkc\checkmark$\checkmarko\checkmark$\checkmark^\checkmark\\checkmarkc\checkmarki\checkmarkr\checkmarkc\checkmark$\checkmarkn\checkmark$\checkmark^\checkmark\\checkmarkc\checkmarki\checkmarkr\checkmarkc\checkmark$\checkmark \checkmark$\checkmark^\checkmark\\checkmarkc\checkmarki\checkmarkr\checkmarkc\checkmark$\checkmark(\checkmark$\checkmark^\checkmark\\checkmarkc\checkmarki\checkmarkr\checkmarkc\checkmark$\checkmarkX\checkmark$\checkmark^\checkmark\\checkmarkc\checkmarki\checkmarkr\checkmarkc\checkmark$\checkmark,\checkmark$\checkmark^\checkmark\\checkmarkc\checkmarki\checkmarkr\checkmarkc\checkmark$\checkmark \checkmark$\checkmark^\checkmark\\checkmarkc\checkmarki\checkmarkr\checkmarkc\checkmark$\checkmarkY\checkmark$\checkmark^\checkmark\\checkmarkc\checkmarki\checkmarkr\checkmarkc\checkmark$\checkmark,\checkmark$\checkmark^\checkmark\\checkmarkc\checkmarki\checkmarkr\checkmarkc\checkmark$\checkmark \checkmark$\checkmark^\checkmark\\checkmarkc\checkmarki\checkmarkr\checkmarkc\checkmark$\checkmarkZ\checkmark$\checkmark^\checkmark\\checkmarkc\checkmarki\checkmarkr\checkmarkc\checkmark$\checkmark)\checkmark$\checkmark^\checkmark\\checkmarkc\checkmarki\checkmarkr\checkmarkc\checkmark$\checkmark \checkmark$\checkmark^\checkmark\\checkmarkc\checkmarki\checkmarkr\checkmarkc\checkmark$\checkmarkd\checkmark$\checkmark^\checkmark\\checkmarkc\checkmarki\checkmarkr\checkmarkc\checkmark$\checkmarké\checkmark$\checkmark^\checkmark\\checkmarkc\checkmarki\checkmarkr\checkmarkc\checkmark$\checkmarkp\checkmark$\checkmark^\checkmark\\checkmarkc\checkmarki\checkmarkr\checkmarkc\checkmark$\checkmarkl\checkmark$\checkmark^\checkmark\\checkmarkc\checkmarki\checkmarkr\checkmarkc\checkmark$\checkmarka\checkmark$\checkmark^\checkmark\\checkmarkc\checkmarki\checkmarkr\checkmarkc\checkmark$\checkmarkc\checkmark$\checkmark^\checkmark\\checkmarkc\checkmarki\checkmarkr\checkmarkc\checkmark$\checkmarke\checkmark$\checkmark^\checkmark\\checkmarkc\checkmarki\checkmarkr\checkmarkc\checkmark$\checkmarkn\checkmark$\checkmark^\checkmark\\checkmarkc\checkmarki\checkmarkr\checkmarkc\checkmark$\checkmarkt\checkmark$\checkmark^\checkmark\\checkmarkc\checkmarki\checkmarkr\checkmarkc\checkmark$\checkmark \checkmark$\checkmark^\checkmark\\checkmarkc\checkmarki\checkmarkr\checkmarkc\checkmark$\checkmarkl\checkmark$\checkmark^\checkmark\\checkmarkc\checkmarki\checkmarkr\checkmarkc\checkmark$\checkmark'\checkmark$\checkmark^\checkmark\\checkmarkc\checkmarki\checkmarkr\checkmarkc\checkmark$\checkmarko\checkmark$\checkmark^\checkmark\\checkmarkc\checkmarki\checkmarkr\checkmarkc\checkmark$\checkmarku\checkmark$\checkmark^\checkmark\\checkmarkc\checkmarki\checkmarkr\checkmarkc\checkmark$\checkmarkt\checkmark$\checkmark^\checkmark\\checkmarkc\checkmarki\checkmarkr\checkmarkc\checkmark$\checkmarki\checkmark$\checkmark^\checkmark\\checkmarkc\checkmarki\checkmarkr\checkmarkc\checkmark$\checkmarkl\checkmark$\checkmark^\checkmark\\checkmarkc\checkmarki\checkmarkr\checkmarkc\checkmark$\checkmark.\checkmark$\checkmark^\checkmark\\checkmarkc\checkmarki\checkmarkr\checkmarkc\checkmark$\checkmark
\checkmark$\checkmark^\checkmark\\checkmarkc\checkmarki\checkmarkr\checkmarkc\checkmark$\checkmark
\checkmark$\checkmark^\checkmark\\checkmarkc\checkmarki\checkmarkr\checkmarkc\checkmark$\checkmarkL\checkmark$\checkmark^\checkmark\\checkmarkc\checkmarki\checkmarkr\checkmarkc\checkmark$\checkmarke\checkmark$\checkmark^\checkmark\\checkmarkc\checkmarki\checkmarkr\checkmarkc\checkmark$\checkmark \checkmark$\checkmark^\checkmark\\checkmarkc\checkmarki\checkmarkr\checkmarkc\checkmark$\checkmark\\checkmark$\checkmark^\checkmark\\checkmarkc\checkmarki\checkmarkr\checkmarkc\checkmark$\checkmarkt\checkmark$\checkmark^\checkmark\\checkmarkc\checkmarki\checkmarkr\checkmarkc\checkmark$\checkmarke\checkmark$\checkmark^\checkmark\\checkmarkc\checkmarki\checkmarkr\checkmarkc\checkmark$\checkmarkx\checkmark$\checkmark^\checkmark\\checkmarkc\checkmarki\checkmarkr\checkmarkc\checkmark$\checkmarkt\checkmark$\checkmark^\checkmark\\checkmarkc\checkmarki\checkmarkr\checkmarkc\checkmark$\checkmarkb\checkmark$\checkmark^\checkmark\\checkmarkc\checkmarki\checkmarkr\checkmarkc\checkmark$\checkmarkf\checkmark$\checkmark^\checkmark\\checkmarkc\checkmarki\checkmarkr\checkmarkc\checkmark$\checkmark{\checkmark$\checkmark^\checkmark\\checkmarkc\checkmarki\checkmarkr\checkmarkc\checkmark$\checkmarkb\checkmark$\checkmark^\checkmark\\checkmarkc\checkmarki\checkmarkr\checkmarkc\checkmark$\checkmarke\checkmark$\checkmark^\checkmark\\checkmarkc\checkmarki\checkmarkr\checkmarkc\checkmark$\checkmarkr\checkmark$\checkmark^\checkmark\\checkmarkc\checkmarki\checkmarkr\checkmarkc\checkmark$\checkmarkc\checkmark$\checkmark^\checkmark\\checkmarkc\checkmarki\checkmarkr\checkmarkc\checkmark$\checkmarke\checkmark$\checkmark^\checkmark\\checkmarkc\checkmarki\checkmarkr\checkmarkc\checkmark$\checkmarka\checkmark$\checkmark^\checkmark\\checkmarkc\checkmarki\checkmarkr\checkmarkc\checkmark$\checkmarku\checkmark$\checkmark^\checkmark\\checkmarkc\checkmarki\checkmarkr\checkmarkc\checkmark$\checkmark \checkmark$\checkmark^\checkmark\\checkmarkc\checkmarki\checkmarkr\checkmarkc\checkmark$\checkmark(\checkmark$\checkmark^\checkmark\\checkmarkc\checkmarki\checkmarkr\checkmarkc\checkmark$\checkmarkA\checkmark$\checkmark^\checkmark\\checkmarkc\checkmarki\checkmarkr\checkmarkc\checkmark$\checkmarkx\checkmark$\checkmark^\checkmark\\checkmarkc\checkmarki\checkmarkr\checkmarkc\checkmark$\checkmarke\checkmark$\checkmark^\checkmark\\checkmarkc\checkmarki\checkmarkr\checkmarkc\checkmark$\checkmark \checkmark$\checkmark^\checkmark\\checkmarkc\checkmarki\checkmarkr\checkmarkc\checkmark$\checkmarkA\checkmark$\checkmark^\checkmark\\checkmarkc\checkmarki\checkmarkr\checkmarkc\checkmark$\checkmark)\checkmark$\checkmark^\checkmark\\checkmarkc\checkmarki\checkmarkr\checkmarkc\checkmark$\checkmark}\checkmark$\checkmark^\checkmark\\checkmarkc\checkmarki\checkmarkr\checkmarkc\checkmark$\checkmark \checkmark$\checkmark^\checkmark\\checkmarkc\checkmarki\checkmarkr\checkmarkc\checkmark$\checkmark—\checkmark$\checkmark^\checkmark\\checkmarkc\checkmarki\checkmarkr\checkmarkc\checkmark$\checkmark \checkmark$\checkmark^\checkmark\\checkmarkc\checkmarki\checkmarkr\checkmarkc\checkmark$\checkmarkm\checkmark$\checkmark^\checkmark\\checkmarkc\checkmarki\checkmarkr\checkmarkc\checkmark$\checkmarko\checkmark$\checkmark^\checkmark\\checkmarkc\checkmarki\checkmarkr\checkmarkc\checkmark$\checkmarkn\checkmark$\checkmark^\checkmark\\checkmarkc\checkmarki\checkmarkr\checkmarkc\checkmark$\checkmarkt\checkmark$\checkmark^\checkmark\\checkmarkc\checkmarki\checkmarkr\checkmarkc\checkmark$\checkmarké\checkmark$\checkmark^\checkmark\\checkmarkc\checkmarki\checkmarkr\checkmarkc\checkmark$\checkmark \checkmark$\checkmark^\checkmark\\checkmarkc\checkmarki\checkmarkr\checkmarkc\checkmark$\checkmarks\checkmark$\checkmark^\checkmark\\checkmarkc\checkmarki\checkmarkr\checkmarkc\checkmark$\checkmarku\checkmark$\checkmark^\checkmark\\checkmarkc\checkmarki\checkmarkr\checkmarkc\checkmark$\checkmarkr\checkmark$\checkmark^\checkmark\\checkmarkc\checkmarki\checkmarkr\checkmarkc\checkmark$\checkmark \checkmark$\checkmark^\checkmark\\checkmarkc\checkmarki\checkmarkr\checkmarkc\checkmark$\checkmarkl\checkmark$\checkmark^\checkmark\\checkmarkc\checkmarki\checkmarkr\checkmarkc\checkmark$\checkmarke\checkmark$\checkmark^\checkmark\\checkmarkc\checkmarki\checkmarkr\checkmarkc\checkmark$\checkmark \checkmark$\checkmark^\checkmark\\checkmarkc\checkmarki\checkmarkr\checkmarkc\checkmark$\checkmarkc\checkmark$\checkmark^\checkmark\\checkmarkc\checkmarki\checkmarkr\checkmarkc\checkmark$\checkmarkh\checkmark$\checkmark^\checkmark\\checkmarkc\checkmarki\checkmarkr\checkmarkc\checkmark$\checkmarka\checkmark$\checkmark^\checkmark\\checkmarkc\checkmarki\checkmarkr\checkmarkc\checkmark$\checkmarkr\checkmark$\checkmark^\checkmark\\checkmarkc\checkmarki\checkmarkr\checkmarkc\checkmark$\checkmarki\checkmark$\checkmark^\checkmark\\checkmarkc\checkmarki\checkmarkr\checkmarkc\checkmark$\checkmarko\checkmark$\checkmark^\checkmark\\checkmarkc\checkmarki\checkmarkr\checkmarkc\checkmark$\checkmarkt\checkmark$\checkmark^\checkmark\\checkmarkc\checkmarki\checkmarkr\checkmarkc\checkmark$\checkmark \checkmark$\checkmark^\checkmark\\checkmarkc\checkmarki\checkmarkr\checkmarkc\checkmark$\checkmarkY\checkmark$\checkmark^\checkmark\\checkmarkc\checkmarki\checkmarkr\checkmarkc\checkmark$\checkmark \checkmark$\checkmark^\checkmark\\checkmarkc\checkmarki\checkmarkr\checkmarkc\checkmark$\checkmark—\checkmark$\checkmark^\checkmark\\checkmarkc\checkmarki\checkmarkr\checkmarkc\checkmark$\checkmark \checkmark$\checkmark^\checkmark\\checkmarkc\checkmarki\checkmarkr\checkmarkc\checkmark$\checkmarkb\checkmark$\checkmark^\checkmark\\checkmarkc\checkmarki\checkmarkr\checkmarkc\checkmark$\checkmarka\checkmark$\checkmark^\checkmark\\checkmarkc\checkmarki\checkmarkr\checkmarkc\checkmark$\checkmarks\checkmark$\checkmark^\checkmark\\checkmarkc\checkmarki\checkmarkr\checkmarkc\checkmark$\checkmarkc\checkmark$\checkmark^\checkmark\\checkmarkc\checkmarki\checkmarkr\checkmarkc\checkmark$\checkmarku\checkmark$\checkmark^\checkmark\\checkmarkc\checkmarki\checkmarkr\checkmarkc\checkmark$\checkmarkl\checkmark$\checkmark^\checkmark\\checkmarkc\checkmarki\checkmarkr\checkmarkc\checkmark$\checkmarke\checkmark$\checkmark^\checkmark\\checkmarkc\checkmarki\checkmarkr\checkmarkc\checkmark$\checkmark \checkmark$\checkmark^\checkmark\\checkmarkc\checkmarki\checkmarkr\checkmarkc\checkmark$\checkmarkd\checkmark$\checkmark^\checkmark\\checkmarkc\checkmarki\checkmarkr\checkmarkc\checkmark$\checkmarke\checkmark$\checkmark^\checkmark\\checkmarkc\checkmarki\checkmarkr\checkmarkc\checkmark$\checkmark \checkmark$\checkmark^\checkmark\\checkmarkc\checkmarki\checkmarkr\checkmarkc\checkmark$\checkmark$\checkmark$\checkmark^\checkmark\\checkmarkc\checkmarki\checkmarkr\checkmarkc\checkmark$\checkmark-\checkmark$\checkmark^\checkmark\\checkmarkc\checkmarki\checkmarkr\checkmarkc\checkmark$\checkmark9\checkmark$\checkmark^\checkmark\\checkmarkc\checkmarki\checkmarkr\checkmarkc\checkmark$\checkmark0\checkmark$\checkmark^\checkmark\\checkmarkc\checkmarki\checkmarkr\checkmarkc\checkmark$\checkmark°\checkmark$\checkmark^\checkmark\\checkmarkc\checkmarki\checkmarkr\checkmarkc\checkmark$\checkmark$\checkmark$\checkmark^\checkmark\\checkmarkc\checkmarki\checkmarkr\checkmarkc\checkmark$\checkmark \checkmark$\checkmark^\checkmark\\checkmarkc\checkmarki\checkmarkr\checkmarkc\checkmark$\checkmarkà\checkmark$\checkmark^\checkmark\\checkmarkc\checkmarki\checkmarkr\checkmarkc\checkmark$\checkmark \checkmark$\checkmark^\checkmark\\checkmarkc\checkmarki\checkmarkr\checkmarkc\checkmark$\checkmark$\checkmark$\checkmark^\checkmark\\checkmarkc\checkmarki\checkmarkr\checkmarkc\checkmark$\checkmark+\checkmark$\checkmark^\checkmark\\checkmarkc\checkmarki\checkmarkr\checkmarkc\checkmark$\checkmark9\checkmark$\checkmark^\checkmark\\checkmarkc\checkmarki\checkmarkr\checkmarkc\checkmark$\checkmark0\checkmark$\checkmark^\checkmark\\checkmarkc\checkmarki\checkmarkr\checkmarkc\checkmark$\checkmark°\checkmark$\checkmark^\checkmark\\checkmarkc\checkmarki\checkmarkr\checkmarkc\checkmark$\checkmark$\checkmark$\checkmark^\checkmark\\checkmarkc\checkmarki\checkmarkr\checkmarkc\checkmark$\checkmark.\checkmark$\checkmark^\checkmark\\checkmarkc\checkmarki\checkmarkr\checkmarkc\checkmark$\checkmark \checkmark$\checkmark^\checkmark\\checkmarkc\checkmarki\checkmarkr\checkmarkc\checkmark$\checkmarkE\checkmark$\checkmark^\checkmark\\checkmarkc\checkmarki\checkmarkr\checkmarkc\checkmark$\checkmarkn\checkmark$\checkmark^\checkmark\\checkmarkc\checkmarki\checkmarkr\checkmarkc\checkmark$\checkmark \checkmark$\checkmark^\checkmark\\checkmarkc\checkmarki\checkmarkr\checkmarkc\checkmark$\checkmarks\checkmark$\checkmark^\checkmark\\checkmarkc\checkmarki\checkmarkr\checkmarkc\checkmark$\checkmarko\checkmark$\checkmark^\checkmark\\checkmarkc\checkmarki\checkmarkr\checkmarkc\checkmark$\checkmarkn\checkmark$\checkmark^\checkmark\\checkmarkc\checkmarki\checkmarkr\checkmarkc\checkmark$\checkmark \checkmark$\checkmark^\checkmark\\checkmarkc\checkmarki\checkmarkr\checkmarkc\checkmark$\checkmarkc\checkmark$\checkmark^\checkmark\\checkmarkc\checkmarki\checkmarkr\checkmarkc\checkmark$\checkmarke\checkmark$\checkmark^\checkmark\\checkmarkc\checkmarki\checkmarkr\checkmarkc\checkmark$\checkmarkn\checkmark$\checkmark^\checkmark\\checkmarkc\checkmarki\checkmarkr\checkmarkc\checkmark$\checkmarkt\checkmark$\checkmark^\checkmark\\checkmarkc\checkmarki\checkmarkr\checkmarkc\checkmark$\checkmarkr\checkmark$\checkmark^\checkmark\\checkmarkc\checkmarki\checkmarkr\checkmarkc\checkmark$\checkmarke\checkmark$\checkmark^\checkmark\\checkmarkc\checkmarki\checkmarkr\checkmarkc\checkmark$\checkmark \checkmark$\checkmark^\checkmark\\checkmarkc\checkmarki\checkmarkr\checkmarkc\checkmark$\checkmarks\checkmark$\checkmark^\checkmark\\checkmarkc\checkmarki\checkmarkr\checkmarkc\checkmark$\checkmarke\checkmark$\checkmark^\checkmark\\checkmarkc\checkmarki\checkmarkr\checkmarkc\checkmark$\checkmark \checkmark$\checkmark^\checkmark\\checkmarkc\checkmarki\checkmarkr\checkmarkc\checkmark$\checkmarkt\checkmark$\checkmark^\checkmark\\checkmarkc\checkmarki\checkmarkr\checkmarkc\checkmark$\checkmarkr\checkmark$\checkmark^\checkmark\\checkmarkc\checkmarki\checkmarkr\checkmarkc\checkmark$\checkmarko\checkmark$\checkmark^\checkmark\\checkmarkc\checkmarki\checkmarkr\checkmarkc\checkmark$\checkmarku\checkmark$\checkmark^\checkmark\\checkmarkc\checkmarki\checkmarkr\checkmarkc\checkmark$\checkmarkv\checkmark$\checkmark^\checkmark\\checkmarkc\checkmarki\checkmarkr\checkmarkc\checkmark$\checkmarke\checkmark$\checkmark^\checkmark\\checkmarkc\checkmarki\checkmarkr\checkmarkc\checkmark$\checkmark \checkmark$\checkmark^\checkmark\\checkmarkc\checkmarki\checkmarkr\checkmarkc\checkmark$\checkmarkl\checkmark$\checkmark^\checkmark\\checkmarkc\checkmarki\checkmarkr\checkmarkc\checkmark$\checkmarke\checkmark$\checkmark^\checkmark\\checkmarkc\checkmarki\checkmarkr\checkmarkc\checkmark$\checkmark \checkmark$\checkmark^\checkmark\\checkmarkc\checkmarki\checkmarkr\checkmarkc\checkmark$\checkmark\\checkmark$\checkmark^\checkmark\\checkmarkc\checkmarki\checkmarkr\checkmarkc\checkmark$\checkmarkt\checkmark$\checkmark^\checkmark\\checkmarkc\checkmarki\checkmarkr\checkmarkc\checkmark$\checkmarke\checkmark$\checkmark^\checkmark\\checkmarkc\checkmarki\checkmarkr\checkmarkc\checkmark$\checkmarkx\checkmark$\checkmark^\checkmark\\checkmarkc\checkmarki\checkmarkr\checkmarkc\checkmark$\checkmarkt\checkmark$\checkmark^\checkmark\\checkmarkc\checkmarki\checkmarkr\checkmarkc\checkmark$\checkmarkb\checkmark$\checkmark^\checkmark\\checkmarkc\checkmarki\checkmarkr\checkmarkc\checkmark$\checkmarkf\checkmark$\checkmark^\checkmark\\checkmarkc\checkmarki\checkmarkr\checkmarkc\checkmark$\checkmark{\checkmark$\checkmark^\checkmark\\checkmarkc\checkmarki\checkmarkr\checkmarkc\checkmark$\checkmarkp\checkmark$\checkmark^\checkmark\\checkmarkc\checkmarki\checkmarkr\checkmarkc\checkmark$\checkmarkl\checkmark$\checkmark^\checkmark\\checkmarkc\checkmarki\checkmarkr\checkmarkc\checkmark$\checkmarka\checkmark$\checkmark^\checkmark\\checkmarkc\checkmarki\checkmarkr\checkmarkc\checkmark$\checkmarkt\checkmark$\checkmark^\checkmark\\checkmarkc\checkmarki\checkmarkr\checkmarkc\checkmark$\checkmarke\checkmark$\checkmark^\checkmark\\checkmarkc\checkmarki\checkmarkr\checkmarkc\checkmark$\checkmarka\checkmark$\checkmark^\checkmark\\checkmarkc\checkmarki\checkmarkr\checkmarkc\checkmark$\checkmarku\checkmark$\checkmark^\checkmark\\checkmarkc\checkmarki\checkmarkr\checkmarkc\checkmark$\checkmark \checkmark$\checkmark^\checkmark\\checkmarkc\checkmarki\checkmarkr\checkmarkc\checkmark$\checkmarkr\checkmark$\checkmark^\checkmark\\checkmarkc\checkmarki\checkmarkr\checkmarkc\checkmark$\checkmarko\checkmark$\checkmark^\checkmark\\checkmarkc\checkmarki\checkmarkr\checkmarkc\checkmark$\checkmarkt\checkmark$\checkmark^\checkmark\\checkmarkc\checkmarki\checkmarkr\checkmarkc\checkmark$\checkmarka\checkmark$\checkmark^\checkmark\\checkmarkc\checkmarki\checkmarkr\checkmarkc\checkmark$\checkmarkt\checkmark$\checkmark^\checkmark\\checkmarkc\checkmarki\checkmarkr\checkmarkc\checkmark$\checkmarki\checkmark$\checkmark^\checkmark\\checkmarkc\checkmarki\checkmarkr\checkmarkc\checkmark$\checkmarkf\checkmark$\checkmark^\checkmark\\checkmarkc\checkmarki\checkmarkr\checkmarkc\checkmark$\checkmark \checkmark$\checkmark^\checkmark\\checkmarkc\checkmarki\checkmarkr\checkmarkc\checkmark$\checkmark(\checkmark$\checkmark^\checkmark\\checkmarkc\checkmarki\checkmarkr\checkmarkc\checkmark$\checkmarkA\checkmark$\checkmark^\checkmark\\checkmarkc\checkmarki\checkmarkr\checkmarkc\checkmark$\checkmarkx\checkmark$\checkmark^\checkmark\\checkmarkc\checkmarki\checkmarkr\checkmarkc\checkmark$\checkmarke\checkmark$\checkmark^\checkmark\\checkmarkc\checkmarki\checkmarkr\checkmarkc\checkmark$\checkmark \checkmark$\checkmark^\checkmark\\checkmarkc\checkmarki\checkmarkr\checkmarkc\checkmark$\checkmarkB\checkmark$\checkmark^\checkmark\\checkmarkc\checkmarki\checkmarkr\checkmarkc\checkmark$\checkmark)\checkmark$\checkmark^\checkmark\\checkmarkc\checkmarki\checkmarkr\checkmarkc\checkmark$\checkmark}\checkmark$\checkmark^\checkmark\\checkmarkc\checkmarki\checkmarkr\checkmarkc\checkmark$\checkmark \checkmark$\checkmark^\checkmark\\checkmarkc\checkmarki\checkmarkr\checkmarkc\checkmark$\checkmarkq\checkmark$\checkmark^\checkmark\\checkmarkc\checkmarki\checkmarkr\checkmarkc\checkmark$\checkmarku\checkmark$\checkmark^\checkmark\\checkmarkc\checkmarki\checkmarkr\checkmarkc\checkmark$\checkmarki\checkmark$\checkmark^\checkmark\\checkmarkc\checkmarki\checkmarkr\checkmarkc\checkmark$\checkmark \checkmark$\checkmark^\checkmark\\checkmarkc\checkmarki\checkmarkr\checkmarkc\checkmark$\checkmarke\checkmark$\checkmark^\checkmark\\checkmarkc\checkmarki\checkmarkr\checkmarkc\checkmark$\checkmarkf\checkmark$\checkmark^\checkmark\\checkmarkc\checkmarki\checkmarkr\checkmarkc\checkmark$\checkmarkf\checkmark$\checkmark^\checkmark\\checkmarkc\checkmarki\checkmarkr\checkmarkc\checkmark$\checkmarke\checkmark$\checkmark^\checkmark\\checkmarkc\checkmarki\checkmarkr\checkmarkc\checkmark$\checkmarkc\checkmark$\checkmark^\checkmark\\checkmarkc\checkmarki\checkmarkr\checkmarkc\checkmark$\checkmarkt\checkmark$\checkmark^\checkmark\\checkmarkc\checkmarki\checkmarkr\checkmarkc\checkmark$\checkmarku\checkmark$\checkmark^\checkmark\\checkmarkc\checkmarki\checkmarkr\checkmarkc\checkmark$\checkmarke\checkmark$\checkmark^\checkmark\\checkmarkc\checkmarki\checkmarkr\checkmarkc\checkmark$\checkmark \checkmark$\checkmark^\checkmark\\checkmarkc\checkmarki\checkmarkr\checkmarkc\checkmark$\checkmarku\checkmark$\checkmark^\checkmark\\checkmarkc\checkmarki\checkmarkr\checkmarkc\checkmark$\checkmarkn\checkmark$\checkmark^\checkmark\\checkmarkc\checkmarki\checkmarkr\checkmarkc\checkmark$\checkmarke\checkmark$\checkmark^\checkmark\\checkmarkc\checkmarki\checkmarkr\checkmarkc\checkmark$\checkmark \checkmark$\checkmark^\checkmark\\checkmarkc\checkmarki\checkmarkr\checkmarkc\checkmark$\checkmarkr\checkmark$\checkmark^\checkmark\\checkmarkc\checkmarki\checkmarkr\checkmarkc\checkmark$\checkmarko\checkmark$\checkmark^\checkmark\\checkmarkc\checkmarki\checkmarkr\checkmarkc\checkmark$\checkmarkt\checkmark$\checkmark^\checkmark\\checkmarkc\checkmarki\checkmarkr\checkmarkc\checkmark$\checkmarka\checkmark$\checkmark^\checkmark\\checkmarkc\checkmarki\checkmarkr\checkmarkc\checkmark$\checkmarkt\checkmark$\checkmark^\checkmark\\checkmarkc\checkmarki\checkmarkr\checkmarkc\checkmark$\checkmarki\checkmark$\checkmark^\checkmark\\checkmarkc\checkmarki\checkmarkr\checkmarkc\checkmark$\checkmarko\checkmark$\checkmark^\checkmark\\checkmarkc\checkmarki\checkmarkr\checkmarkc\checkmark$\checkmarkn\checkmark$\checkmark^\checkmark\\checkmarkc\checkmarki\checkmarkr\checkmarkc\checkmark$\checkmark \checkmark$\checkmark^\checkmark\\checkmarkc\checkmarki\checkmarkr\checkmarkc\checkmark$\checkmarkc\checkmark$\checkmark^\checkmark\\checkmarkc\checkmarki\checkmarkr\checkmarkc\checkmark$\checkmarko\checkmark$\checkmark^\checkmark\\checkmarkc\checkmarki\checkmarkr\checkmarkc\checkmark$\checkmarkm\checkmark$\checkmark^\checkmark\\checkmarkc\checkmarki\checkmarkr\checkmarkc\checkmark$\checkmarkp\checkmark$\checkmark^\checkmark\\checkmarkc\checkmarki\checkmarkr\checkmarkc\checkmark$\checkmarkl\checkmark$\checkmark^\checkmark\\checkmarkc\checkmarki\checkmarkr\checkmarkc\checkmark$\checkmarkè\checkmark$\checkmark^\checkmark\\checkmarkc\checkmarki\checkmarkr\checkmarkc\checkmark$\checkmarkt\checkmark$\checkmark^\checkmark\\checkmarkc\checkmarki\checkmarkr\checkmarkc\checkmark$\checkmarke\checkmark$\checkmark^\checkmark\\checkmarkc\checkmarki\checkmarkr\checkmarkc\checkmark$\checkmark \checkmark$\checkmark^\checkmark\\checkmarkc\checkmarki\checkmarkr\checkmarkc\checkmark$\checkmarkà\checkmark$\checkmark^\checkmark\\checkmarkc\checkmarki\checkmarkr\checkmarkc\checkmark$\checkmark \checkmark$\checkmark^\checkmark\\checkmarkc\checkmarki\checkmarkr\checkmarkc\checkmark$\checkmark$\checkmark$\checkmark^\checkmark\\checkmarkc\checkmarki\checkmarkr\checkmarkc\checkmark$\checkmark3\checkmark$\checkmark^\checkmark\\checkmarkc\checkmarki\checkmarkr\checkmarkc\checkmark$\checkmark6\checkmark$\checkmark^\checkmark\\checkmarkc\checkmarki\checkmarkr\checkmarkc\checkmark$\checkmark0\checkmark$\checkmark^\checkmark\\checkmarkc\checkmarki\checkmarkr\checkmarkc\checkmark$\checkmark°\checkmark$\checkmark^\checkmark\\checkmarkc\checkmarki\checkmarkr\checkmarkc\checkmark$\checkmark$\checkmark$\checkmark^\checkmark\\checkmarkc\checkmarki\checkmarkr\checkmarkc\checkmark$\checkmark \checkmark$\checkmark^\checkmark\\checkmarkc\checkmarki\checkmarkr\checkmarkc\checkmark$\checkmarke\checkmark$\checkmark^\checkmark\\checkmarkc\checkmarki\checkmarkr\checkmarkc\checkmark$\checkmarkt\checkmark$\checkmark^\checkmark\\checkmarkc\checkmarki\checkmarkr\checkmarkc\checkmark$\checkmark \checkmark$\checkmark^\checkmark\\checkmarkc\checkmarki\checkmarkr\checkmarkc\checkmark$\checkmarkm\checkmark$\checkmark^\checkmark\\checkmarkc\checkmarki\checkmarkr\checkmarkc\checkmark$\checkmarka\checkmark$\checkmark^\checkmark\\checkmarkc\checkmarki\checkmarkr\checkmarkc\checkmark$\checkmarki\checkmark$\checkmark^\checkmark\\checkmarkc\checkmarki\checkmarkr\checkmarkc\checkmark$\checkmarkn\checkmark$\checkmark^\checkmark\\checkmarkc\checkmarki\checkmarkr\checkmarkc\checkmark$\checkmarkt\checkmark$\checkmark^\checkmark\\checkmarkc\checkmarki\checkmarkr\checkmarkc\checkmark$\checkmarki\checkmark$\checkmark^\checkmark\\checkmarkc\checkmarki\checkmarkr\checkmarkc\checkmark$\checkmarke\checkmark$\checkmark^\checkmark\\checkmarkc\checkmarki\checkmarkr\checkmarkc\checkmark$\checkmarkn\checkmark$\checkmark^\checkmark\\checkmarkc\checkmarki\checkmarkr\checkmarkc\checkmark$\checkmarkt\checkmark$\checkmark^\checkmark\\checkmarkc\checkmarki\checkmarkr\checkmarkc\checkmark$\checkmark \checkmark$\checkmark^\checkmark\\checkmarkc\checkmarki\checkmarkr\checkmarkc\checkmark$\checkmarkl\checkmark$\checkmark^\checkmark\\checkmarkc\checkmarki\checkmarkr\checkmarkc\checkmark$\checkmarka\checkmark$\checkmark^\checkmark\\checkmarkc\checkmarki\checkmarkr\checkmarkc\checkmark$\checkmark \checkmark$\checkmark^\checkmark\\checkmarkc\checkmarki\checkmarkr\checkmarkc\checkmark$\checkmarkp\checkmark$\checkmark^\checkmark\\checkmarkc\checkmarki\checkmarkr\checkmarkc\checkmark$\checkmarki\checkmark$\checkmark^\checkmark\\checkmarkc\checkmarki\checkmarkr\checkmarkc\checkmark$\checkmarkè\checkmark$\checkmark^\checkmark\\checkmarkc\checkmarki\checkmarkr\checkmarkc\checkmark$\checkmarkc\checkmark$\checkmark^\checkmark\\checkmarkc\checkmarki\checkmarkr\checkmarkc\checkmark$\checkmarke\checkmark$\checkmark^\checkmark\\checkmarkc\checkmarki\checkmarkr\checkmarkc\checkmark$\checkmark.\checkmark$\checkmark^\checkmark\\checkmarkc\checkmarki\checkmarkr\checkmarkc\checkmark$\checkmark \checkmark$\checkmark^\checkmark\\checkmarkc\checkmarki\checkmarkr\checkmarkc\checkmark$\checkmarkL\checkmark$\checkmark^\checkmark\\checkmarkc\checkmarki\checkmarkr\checkmarkc\checkmark$\checkmarka\checkmark$\checkmark^\checkmark\\checkmarkc\checkmarki\checkmarkr\checkmarkc\checkmark$\checkmark \checkmark$\checkmark^\checkmark\\checkmarkc\checkmarki\checkmarkr\checkmarkc\checkmark$\checkmarkb\checkmark$\checkmark^\checkmark\\checkmarkc\checkmarki\checkmarkr\checkmarkc\checkmark$\checkmarkr\checkmark$\checkmark^\checkmark\\checkmarkc\checkmarki\checkmarkr\checkmarkc\checkmark$\checkmarko\checkmark$\checkmark^\checkmark\\checkmarkc\checkmarki\checkmarkr\checkmarkc\checkmark$\checkmarkc\checkmark$\checkmark^\checkmark\\checkmarkc\checkmarki\checkmarkr\checkmarkc\checkmark$\checkmarkh\checkmark$\checkmark^\checkmark\\checkmarkc\checkmarki\checkmarkr\checkmarkc\checkmark$\checkmarke\checkmark$\checkmark^\checkmark\\checkmarkc\checkmarki\checkmarkr\checkmarkc\checkmark$\checkmark \checkmark$\checkmark^\checkmark\\checkmarkc\checkmarki\checkmarkr\checkmarkc\checkmark$\checkmark(\checkmark$\checkmark^\checkmark\\checkmarkc\checkmarki\checkmarkr\checkmarkc\checkmark$\checkmarko\checkmark$\checkmark^\checkmark\\checkmarkc\checkmarki\checkmarkr\checkmarkc\checkmark$\checkmarku\checkmark$\checkmark^\checkmark\\checkmarkc\checkmarki\checkmarkr\checkmarkc\checkmark$\checkmarkt\checkmark$\checkmark^\checkmark\\checkmarkc\checkmarki\checkmarkr\checkmarkc\checkmark$\checkmarki\checkmark$\checkmark^\checkmark\\checkmarkc\checkmarki\checkmarkr\checkmarkc\checkmark$\checkmarkl\checkmark$\checkmark^\checkmark\\checkmarkc\checkmarki\checkmarkr\checkmarkc\checkmark$\checkmark)\checkmark$\checkmark^\checkmark\\checkmarkc\checkmarki\checkmarkr\checkmarkc\checkmark$\checkmark \checkmark$\checkmark^\checkmark\\checkmarkc\checkmarki\checkmarkr\checkmarkc\checkmark$\checkmarke\checkmark$\checkmark^\checkmark\\checkmarkc\checkmarki\checkmarkr\checkmarkc\checkmark$\checkmarks\checkmark$\checkmark^\checkmark\\checkmarkc\checkmarki\checkmarkr\checkmarkc\checkmark$\checkmarkt\checkmark$\checkmark^\checkmark\\checkmarkc\checkmarki\checkmarkr\checkmarkc\checkmark$\checkmark \checkmark$\checkmark^\checkmark\\checkmarkc\checkmarki\checkmarkr\checkmarkc\checkmark$\checkmarkf\checkmark$\checkmark^\checkmark\\checkmarkc\checkmarki\checkmarkr\checkmarkc\checkmark$\checkmarki\checkmark$\checkmark^\checkmark\\checkmarkc\checkmarki\checkmarkr\checkmarkc\checkmark$\checkmarkx\checkmark$\checkmark^\checkmark\\checkmarkc\checkmarki\checkmarkr\checkmarkc\checkmark$\checkmarké\checkmark$\checkmark^\checkmark\\checkmarkc\checkmarki\checkmarkr\checkmarkc\checkmark$\checkmarke\checkmark$\checkmark^\checkmark\\checkmarkc\checkmarki\checkmarkr\checkmarkc\checkmark$\checkmark \checkmark$\checkmark^\checkmark\\checkmarkc\checkmarki\checkmarkr\checkmarkc\checkmark$\checkmarks\checkmark$\checkmark^\checkmark\\checkmarkc\checkmarki\checkmarkr\checkmarkc\checkmark$\checkmarku\checkmark$\checkmark^\checkmark\\checkmarkc\checkmarki\checkmarkr\checkmarkc\checkmark$\checkmarkr\checkmark$\checkmark^\checkmark\\checkmarkc\checkmarki\checkmarkr\checkmarkc\checkmark$\checkmark \checkmark$\checkmark^\checkmark\\checkmarkc\checkmarki\checkmarkr\checkmarkc\checkmark$\checkmarkl\checkmark$\checkmark^\checkmark\\checkmarkc\checkmarki\checkmarkr\checkmarkc\checkmark$\checkmarke\checkmark$\checkmark^\checkmark\\checkmarkc\checkmarki\checkmarkr\checkmarkc\checkmark$\checkmark \checkmark$\checkmark^\checkmark\\checkmarkc\checkmarki\checkmarkr\checkmarkc\checkmark$\checkmarkp\checkmark$\checkmark^\checkmark\\checkmarkc\checkmarki\checkmarkr\checkmarkc\checkmark$\checkmarko\checkmark$\checkmark^\checkmark\\checkmarkc\checkmarki\checkmarkr\checkmarkc\checkmark$\checkmarkr\checkmark$\checkmark^\checkmark\\checkmarkc\checkmarki\checkmarkr\checkmarkc\checkmark$\checkmarkt\checkmark$\checkmark^\checkmark\\checkmarkc\checkmarki\checkmarkr\checkmarkc\checkmark$\checkmarki\checkmark$\checkmark^\checkmark\\checkmarkc\checkmarki\checkmarkr\checkmarkc\checkmark$\checkmarkq\checkmark$\checkmark^\checkmark\\checkmarkc\checkmarki\checkmarkr\checkmarkc\checkmark$\checkmarku\checkmark$\checkmark^\checkmark\\checkmarkc\checkmarki\checkmarkr\checkmarkc\checkmark$\checkmarke\checkmark$\checkmark^\checkmark\\checkmarkc\checkmarki\checkmarkr\checkmarkc\checkmark$\checkmark \checkmark$\checkmark^\checkmark\\checkmarkc\checkmarki\checkmarkr\checkmarkc\checkmark$\checkmarkX\checkmark$\checkmark^\checkmark\\checkmarkc\checkmarki\checkmarkr\checkmarkc\checkmark$\checkmark-\checkmark$\checkmark^\checkmark\\checkmarkc\checkmarki\checkmarkr\checkmarkc\checkmark$\checkmarkZ\checkmark$\checkmark^\checkmark\\checkmarkc\checkmarki\checkmarkr\checkmarkc\checkmark$\checkmark.\checkmark$\checkmark^\checkmark\\checkmarkc\checkmarki\checkmarkr\checkmarkc\checkmark$\checkmark
\checkmark$\checkmark^\checkmark\\checkmarkc\checkmarki\checkmarkr\checkmarkc\checkmark$\checkmark
\checkmark$\checkmark^\checkmark\\checkmarkc\checkmarki\checkmarkr\checkmarkc\checkmark$\checkmarkL\checkmark$\checkmark^\checkmark\\checkmarkc\checkmarki\checkmarkr\checkmarkc\checkmark$\checkmark'\checkmark$\checkmark^\checkmark\\checkmarkc\checkmarki\checkmarkr\checkmarkc\checkmark$\checkmarka\checkmark$\checkmark^\checkmark\\checkmarkc\checkmarki\checkmarkr\checkmarkc\checkmark$\checkmarkv\checkmark$\checkmark^\checkmark\\checkmarkc\checkmarki\checkmarkr\checkmarkc\checkmark$\checkmarka\checkmark$\checkmark^\checkmark\\checkmarkc\checkmarki\checkmarkr\checkmarkc\checkmark$\checkmarkn\checkmark$\checkmark^\checkmark\\checkmarkc\checkmarki\checkmarkr\checkmarkc\checkmark$\checkmarkt\checkmark$\checkmark^\checkmark\\checkmarkc\checkmarki\checkmarkr\checkmarkc\checkmark$\checkmarka\checkmark$\checkmark^\checkmark\\checkmarkc\checkmarki\checkmarkr\checkmarkc\checkmark$\checkmarkg\checkmark$\checkmark^\checkmark\\checkmarkc\checkmarki\checkmarkr\checkmarkc\checkmark$\checkmarke\checkmark$\checkmark^\checkmark\\checkmarkc\checkmarki\checkmarkr\checkmarkc\checkmark$\checkmark \checkmark$\checkmark^\checkmark\\checkmarkc\checkmarki\checkmarkr\checkmarkc\checkmark$\checkmarkp\checkmark$\checkmark^\checkmark\\checkmarkc\checkmarki\checkmarkr\checkmarkc\checkmark$\checkmarkr\checkmark$\checkmark^\checkmark\\checkmarkc\checkmarki\checkmarkr\checkmarkc\checkmark$\checkmarki\checkmark$\checkmark^\checkmark\\checkmarkc\checkmarki\checkmarkr\checkmarkc\checkmark$\checkmarkn\checkmark$\checkmark^\checkmark\\checkmarkc\checkmarki\checkmarkr\checkmarkc\checkmark$\checkmarkc\checkmark$\checkmark^\checkmark\\checkmarkc\checkmarki\checkmarkr\checkmarkc\checkmark$\checkmarki\checkmark$\checkmark^\checkmark\\checkmarkc\checkmarki\checkmarkr\checkmarkc\checkmark$\checkmarkp\checkmark$\checkmark^\checkmark\\checkmarkc\checkmarki\checkmarkr\checkmarkc\checkmark$\checkmarka\checkmark$\checkmark^\checkmark\\checkmarkc\checkmarki\checkmarkr\checkmarkc\checkmark$\checkmarkl\checkmark$\checkmark^\checkmark\\checkmarkc\checkmarki\checkmarkr\checkmarkc\checkmark$\checkmark \checkmark$\checkmark^\checkmark\\checkmarkc\checkmarki\checkmarkr\checkmarkc\checkmark$\checkmarke\checkmark$\checkmark^\checkmark\\checkmarkc\checkmarki\checkmarkr\checkmarkc\checkmark$\checkmarks\checkmark$\checkmark^\checkmark\\checkmarkc\checkmarki\checkmarkr\checkmarkc\checkmark$\checkmarkt\checkmark$\checkmark^\checkmark\\checkmarkc\checkmarki\checkmarkr\checkmarkc\checkmark$\checkmark \checkmark$\checkmark^\checkmark\\checkmarkc\checkmarki\checkmarkr\checkmarkc\checkmark$\checkmarkl\checkmark$\checkmark^\checkmark\\checkmarkc\checkmarki\checkmarkr\checkmarkc\checkmark$\checkmarka\checkmark$\checkmark^\checkmark\\checkmarkc\checkmarki\checkmarkr\checkmarkc\checkmark$\checkmark \checkmark$\checkmark^\checkmark\\checkmarkc\checkmarki\checkmarkr\checkmarkc\checkmark$\checkmark\\checkmark$\checkmark^\checkmark\\checkmarkc\checkmarki\checkmarkr\checkmarkc\checkmark$\checkmarkt\checkmark$\checkmark^\checkmark\\checkmarkc\checkmarki\checkmarkr\checkmarkc\checkmark$\checkmarke\checkmark$\checkmark^\checkmark\\checkmarkc\checkmarki\checkmarkr\checkmarkc\checkmark$\checkmarkx\checkmark$\checkmark^\checkmark\\checkmarkc\checkmarki\checkmarkr\checkmarkc\checkmark$\checkmarkt\checkmark$\checkmark^\checkmark\\checkmarkc\checkmarki\checkmarkr\checkmarkc\checkmark$\checkmarkb\checkmark$\checkmark^\checkmark\\checkmarkc\checkmarki\checkmarkr\checkmarkc\checkmark$\checkmarkf\checkmark$\checkmark^\checkmark\\checkmarkc\checkmarki\checkmarkr\checkmarkc\checkmark$\checkmark{\checkmark$\checkmark^\checkmark\\checkmarkc\checkmarki\checkmarkr\checkmarkc\checkmark$\checkmarkr\checkmark$\checkmark^\checkmark\\checkmarkc\checkmarki\checkmarkr\checkmarkc\checkmark$\checkmarki\checkmark$\checkmark^\checkmark\\checkmarkc\checkmarki\checkmarkr\checkmarkc\checkmark$\checkmarkg\checkmark$\checkmark^\checkmark\\checkmarkc\checkmarki\checkmarkr\checkmarkc\checkmark$\checkmarki\checkmark$\checkmark^\checkmark\\checkmarkc\checkmarki\checkmarkr\checkmarkc\checkmark$\checkmarkd\checkmark$\checkmark^\checkmark\\checkmarkc\checkmarki\checkmarkr\checkmarkc\checkmark$\checkmarki\checkmark$\checkmark^\checkmark\\checkmarkc\checkmarki\checkmarkr\checkmarkc\checkmark$\checkmarkt\checkmark$\checkmark^\checkmark\\checkmarkc\checkmarki\checkmarkr\checkmarkc\checkmark$\checkmarké\checkmark$\checkmark^\checkmark\\checkmarkc\checkmarki\checkmarkr\checkmarkc\checkmark$\checkmark \checkmark$\checkmark^\checkmark\\checkmarkc\checkmarki\checkmarkr\checkmarkc\checkmark$\checkmarkd\checkmark$\checkmark^\checkmark\\checkmarkc\checkmarki\checkmarkr\checkmarkc\checkmark$\checkmarku\checkmark$\checkmark^\checkmark\\checkmarkc\checkmarki\checkmarkr\checkmarkc\checkmark$\checkmark \checkmark$\checkmark^\checkmark\\checkmarkc\checkmarki\checkmarkr\checkmarkc\checkmark$\checkmarkg\checkmark$\checkmark^\checkmark\\checkmarkc\checkmarki\checkmarkr\checkmarkc\checkmark$\checkmarkr\checkmark$\checkmark^\checkmark\\checkmarkc\checkmarki\checkmarkr\checkmarkc\checkmark$\checkmarko\checkmark$\checkmark^\checkmark\\checkmarkc\checkmarki\checkmarkr\checkmarkc\checkmark$\checkmarku\checkmark$\checkmark^\checkmark\\checkmarkc\checkmarki\checkmarkr\checkmarkc\checkmark$\checkmarkp\checkmark$\checkmark^\checkmark\\checkmarkc\checkmarki\checkmarkr\checkmarkc\checkmark$\checkmarke\checkmark$\checkmark^\checkmark\\checkmarkc\checkmarki\checkmarkr\checkmarkc\checkmark$\checkmark \checkmark$\checkmark^\checkmark\\checkmarkc\checkmarki\checkmarkr\checkmarkc\checkmark$\checkmarkp\checkmark$\checkmark^\checkmark\\checkmarkc\checkmarki\checkmarkr\checkmarkc\checkmark$\checkmarko\checkmark$\checkmark^\checkmark\\checkmarkc\checkmarki\checkmarkr\checkmarkc\checkmark$\checkmarkr\checkmark$\checkmark^\checkmark\\checkmarkc\checkmarki\checkmarkr\checkmarkc\checkmark$\checkmarkt\checkmark$\checkmark^\checkmark\\checkmarkc\checkmarki\checkmarkr\checkmarkc\checkmark$\checkmarke\checkmark$\checkmark^\checkmark\\checkmarkc\checkmarki\checkmarkr\checkmarkc\checkmark$\checkmark-\checkmark$\checkmark^\checkmark\\checkmarkc\checkmarki\checkmarkr\checkmarkc\checkmark$\checkmarko\checkmark$\checkmark^\checkmark\\checkmarkc\checkmarki\checkmarkr\checkmarkc\checkmark$\checkmarku\checkmark$\checkmark^\checkmark\\checkmarkc\checkmarki\checkmarkr\checkmarkc\checkmark$\checkmarkt\checkmark$\checkmark^\checkmark\\checkmarkc\checkmarki\checkmarkr\checkmarkc\checkmark$\checkmarki\checkmark$\checkmark^\checkmark\\checkmarkc\checkmarki\checkmarkr\checkmarkc\checkmark$\checkmarkl\checkmark$\checkmark^\checkmark\\checkmarkc\checkmarki\checkmarkr\checkmarkc\checkmark$\checkmark}\checkmark$\checkmark^\checkmark\\checkmarkc\checkmarki\checkmarkr\checkmarkc\checkmark$\checkmark \checkmark$\checkmark^\checkmark\\checkmarkc\checkmarki\checkmarkr\checkmarkc\checkmark$\checkmark:\checkmark$\checkmark^\checkmark\\checkmarkc\checkmarki\checkmarkr\checkmarkc\checkmark$\checkmark \checkmark$\checkmark^\checkmark\\checkmarkc\checkmarki\checkmarkr\checkmarkc\checkmark$\checkmarkl\checkmark$\checkmark^\checkmark\\checkmarkc\checkmarki\checkmarkr\checkmarkc\checkmark$\checkmarka\checkmark$\checkmark^\checkmark\\checkmarkc\checkmarki\checkmarkr\checkmarkc\checkmark$\checkmark \checkmark$\checkmark^\checkmark\\checkmarkc\checkmarki\checkmarkr\checkmarkc\checkmark$\checkmarkt\checkmark$\checkmark^\checkmark\\checkmarkc\checkmarki\checkmarkr\checkmarkc\checkmark$\checkmarkê\checkmark$\checkmark^\checkmark\\checkmarkc\checkmarki\checkmarkr\checkmarkc\checkmark$\checkmarkt\checkmark$\checkmark^\checkmark\\checkmarkc\checkmarki\checkmarkr\checkmarkc\checkmark$\checkmarke\checkmark$\checkmark^\checkmark\\checkmarkc\checkmarki\checkmarkr\checkmarkc\checkmark$\checkmark \checkmark$\checkmark^\checkmark\\checkmarkc\checkmarki\checkmarkr\checkmarkc\checkmark$\checkmarkd\checkmark$\checkmark^\checkmark\\checkmarkc\checkmarki\checkmarkr\checkmarkc\checkmark$\checkmarke\checkmark$\checkmark^\checkmark\\checkmarkc\checkmarki\checkmarkr\checkmarkc\checkmark$\checkmark \checkmark$\checkmark^\checkmark\\checkmarkc\checkmarki\checkmarkr\checkmarkc\checkmark$\checkmarkc\checkmark$\checkmark^\checkmark\\checkmarkc\checkmarki\checkmarkr\checkmarkc\checkmark$\checkmarko\checkmark$\checkmark^\checkmark\\checkmarkc\checkmarki\checkmarkr\checkmarkc\checkmark$\checkmarku\checkmark$\checkmark^\checkmark\\checkmarkc\checkmarki\checkmarkr\checkmarkc\checkmark$\checkmarkp\checkmark$\checkmark^\checkmark\\checkmarkc\checkmarki\checkmarkr\checkmarkc\checkmark$\checkmarke\checkmark$\checkmark^\checkmark\\checkmarkc\checkmarki\checkmarkr\checkmarkc\checkmark$\checkmark \checkmark$\checkmark^\checkmark\\checkmarkc\checkmarki\checkmarkr\checkmarkc\checkmark$\checkmarke\checkmark$\checkmark^\checkmark\\checkmarkc\checkmarki\checkmarkr\checkmarkc\checkmark$\checkmarks\checkmark$\checkmark^\checkmark\\checkmarkc\checkmarki\checkmarkr\checkmarkc\checkmark$\checkmarkt\checkmark$\checkmark^\checkmark\\checkmarkc\checkmarki\checkmarkr\checkmarkc\checkmark$\checkmark \checkmark$\checkmark^\checkmark\\checkmarkc\checkmarki\checkmarkr\checkmarkc\checkmark$\checkmarkp\checkmark$\checkmark^\checkmark\\checkmarkc\checkmarki\checkmarkr\checkmarkc\checkmark$\checkmarko\checkmark$\checkmark^\checkmark\\checkmarkc\checkmarki\checkmarkr\checkmarkc\checkmark$\checkmarkr\checkmark$\checkmark^\checkmark\\checkmarkc\checkmarki\checkmarkr\checkmarkc\checkmark$\checkmarkt\checkmark$\checkmark^\checkmark\\checkmarkc\checkmarki\checkmarkr\checkmarkc\checkmark$\checkmarké\checkmark$\checkmark^\checkmark\\checkmarkc\checkmarki\checkmarkr\checkmarkc\checkmark$\checkmarke\checkmark$\checkmark^\checkmark\\checkmarkc\checkmarki\checkmarkr\checkmarkc\checkmark$\checkmark \checkmark$\checkmark^\checkmark\\checkmarkc\checkmarki\checkmarkr\checkmarkc\checkmark$\checkmarkp\checkmark$\checkmark^\checkmark\\checkmarkc\checkmarki\checkmarkr\checkmarkc\checkmark$\checkmarka\checkmark$\checkmark^\checkmark\\checkmarkc\checkmarki\checkmarkr\checkmarkc\checkmark$\checkmarkr\checkmark$\checkmark^\checkmark\\checkmarkc\checkmarki\checkmarkr\checkmarkc\checkmark$\checkmark \checkmark$\checkmark^\checkmark\\checkmarkc\checkmarki\checkmarkr\checkmarkc\checkmark$\checkmarku\checkmark$\checkmark^\checkmark\\checkmarkc\checkmarki\checkmarkr\checkmarkc\checkmark$\checkmarkn\checkmark$\checkmark^\checkmark\\checkmarkc\checkmarki\checkmarkr\checkmarkc\checkmark$\checkmarke\checkmark$\checkmark^\checkmark\\checkmarkc\checkmarki\checkmarkr\checkmarkc\checkmark$\checkmark \checkmark$\checkmark^\checkmark\\checkmarkc\checkmarki\checkmarkr\checkmarkc\checkmark$\checkmarks\checkmark$\checkmark^\checkmark\\checkmarkc\checkmarki\checkmarkr\checkmarkc\checkmark$\checkmarkt\checkmark$\checkmark^\checkmark\\checkmarkc\checkmarki\checkmarkr\checkmarkc\checkmark$\checkmarkr\checkmark$\checkmark^\checkmark\\checkmarkc\checkmarki\checkmarkr\checkmarkc\checkmark$\checkmarku\checkmark$\checkmark^\checkmark\\checkmarkc\checkmarki\checkmarkr\checkmarkc\checkmark$\checkmarkc\checkmark$\checkmark^\checkmark\\checkmarkc\checkmarki\checkmarkr\checkmarkc\checkmark$\checkmarkt\checkmark$\checkmark^\checkmark\\checkmarkc\checkmarki\checkmarkr\checkmarkc\checkmark$\checkmarku\checkmark$\checkmark^\checkmark\\checkmarkc\checkmarki\checkmarkr\checkmarkc\checkmark$\checkmarkr\checkmark$\checkmark^\checkmark\\checkmarkc\checkmarki\checkmarkr\checkmarkc\checkmark$\checkmarke\checkmark$\checkmark^\checkmark\\checkmarkc\checkmarki\checkmarkr\checkmarkc\checkmark$\checkmark \checkmark$\checkmark^\checkmark\\checkmarkc\checkmarki\checkmarkr\checkmarkc\checkmark$\checkmarkm\checkmark$\checkmark^\checkmark\\checkmarkc\checkmarki\checkmarkr\checkmarkc\checkmark$\checkmarko\checkmark$\checkmark^\checkmark\\checkmarkc\checkmarki\checkmarkr\checkmarkc\checkmark$\checkmarkn\checkmark$\checkmark^\checkmark\\checkmarkc\checkmarki\checkmarkr\checkmarkc\checkmark$\checkmarko\checkmark$\checkmark^\checkmark\\checkmarkc\checkmarki\checkmarkr\checkmarkc\checkmark$\checkmarkl\checkmark$\checkmark^\checkmark\\checkmarkc\checkmarki\checkmarkr\checkmarkc\checkmark$\checkmarki\checkmark$\checkmark^\checkmark\\checkmarkc\checkmarki\checkmarkr\checkmarkc\checkmark$\checkmarkt\checkmark$\checkmark^\checkmark\\checkmarkc\checkmarki\checkmarkr\checkmarkc\checkmark$\checkmarkh\checkmark$\checkmark^\checkmark\\checkmarkc\checkmarki\checkmarkr\checkmarkc\checkmark$\checkmarki\checkmark$\checkmark^\checkmark\\checkmarkc\checkmarki\checkmarkr\checkmarkc\checkmark$\checkmarkq\checkmark$\checkmark^\checkmark\\checkmarkc\checkmarki\checkmarkr\checkmarkc\checkmark$\checkmarku\checkmark$\checkmark^\checkmark\\checkmarkc\checkmarki\checkmarkr\checkmarkc\checkmark$\checkmarke\checkmark$\checkmark^\checkmark\\checkmarkc\checkmarki\checkmarkr\checkmarkc\checkmark$\checkmark \checkmark$\checkmark^\checkmark\\checkmarkc\checkmarki\checkmarkr\checkmarkc\checkmark$\checkmarke\checkmark$\checkmark^\checkmark\\checkmarkc\checkmarki\checkmarkr\checkmarkc\checkmark$\checkmarkn\checkmark$\checkmark^\checkmark\\checkmarkc\checkmarki\checkmarkr\checkmarkc\checkmark$\checkmark \checkmark$\checkmark^\checkmark\\checkmarkc\checkmarki\checkmarkr\checkmarkc\checkmark$\checkmarkp\checkmark$\checkmark^\checkmark\\checkmarkc\checkmarki\checkmarkr\checkmarkc\checkmark$\checkmarko\checkmark$\checkmark^\checkmark\\checkmarkc\checkmarki\checkmarkr\checkmarkc\checkmark$\checkmarkr\checkmark$\checkmark^\checkmark\\checkmarkc\checkmarki\checkmarkr\checkmarkc\checkmark$\checkmarkt\checkmark$\checkmark^\checkmark\\checkmarkc\checkmarki\checkmarkr\checkmarkc\checkmark$\checkmarki\checkmark$\checkmark^\checkmark\\checkmarkc\checkmarki\checkmarkr\checkmarkc\checkmark$\checkmarkq\checkmark$\checkmark^\checkmark\\checkmarkc\checkmarki\checkmarkr\checkmarkc\checkmark$\checkmarku\checkmark$\checkmark^\checkmark\\checkmarkc\checkmarki\checkmarkr\checkmarkc\checkmark$\checkmarke\checkmark$\checkmark^\checkmark\\checkmarkc\checkmarki\checkmarkr\checkmarkc\checkmark$\checkmark.\checkmark$\checkmark^\checkmark\\checkmarkc\checkmarki\checkmarkr\checkmarkc\checkmark$\checkmark \checkmark$\checkmark^\checkmark\\checkmarkc\checkmarki\checkmarkr\checkmarkc\checkmark$\checkmarkL\checkmark$\checkmark^\checkmark\\checkmarkc\checkmarki\checkmarkr\checkmarkc\checkmark$\checkmark'\checkmark$\checkmark^\checkmark\\checkmarkc\checkmarki\checkmarkr\checkmarkc\checkmark$\checkmarki\checkmark$\checkmark^\checkmark\\checkmarkc\checkmarki\checkmarkr\checkmarkc\checkmark$\checkmarkn\checkmark$\checkmark^\checkmark\\checkmarkc\checkmarki\checkmarkr\checkmarkc\checkmark$\checkmarkc\checkmark$\checkmark^\checkmark\\checkmarkc\checkmarki\checkmarkr\checkmarkc\checkmark$\checkmarko\checkmark$\checkmark^\checkmark\\checkmarkc\checkmarki\checkmarkr\checkmarkc\checkmark$\checkmarkn\checkmark$\checkmark^\checkmark\\checkmarkc\checkmarki\checkmarkr\checkmarkc\checkmark$\checkmarkv\checkmark$\checkmark^\checkmark\\checkmarkc\checkmarki\checkmarkr\checkmarkc\checkmark$\checkmarké\checkmark$\checkmark^\checkmark\\checkmarkc\checkmarki\checkmarkr\checkmarkc\checkmark$\checkmarkn\checkmark$\checkmark^\checkmark\\checkmarkc\checkmarki\checkmarkr\checkmarkc\checkmark$\checkmarki\checkmark$\checkmark^\checkmark\\checkmarkc\checkmarki\checkmarkr\checkmarkc\checkmark$\checkmarke\checkmark$\checkmark^\checkmark\\checkmarkc\checkmarki\checkmarkr\checkmarkc\checkmark$\checkmarkn\checkmark$\checkmark^\checkmark\\checkmarkc\checkmarki\checkmarkr\checkmarkc\checkmark$\checkmarkt\checkmark$\checkmark^\checkmark\\checkmarkc\checkmarki\checkmarkr\checkmarkc\checkmark$\checkmark \checkmark$\checkmark^\checkmark\\checkmarkc\checkmarki\checkmarkr\checkmarkc\checkmark$\checkmarke\checkmark$\checkmark^\checkmark\\checkmarkc\checkmarki\checkmarkr\checkmarkc\checkmark$\checkmarks\checkmark$\checkmark^\checkmark\\checkmarkc\checkmarki\checkmarkr\checkmarkc\checkmark$\checkmarkt\checkmark$\checkmark^\checkmark\\checkmarkc\checkmarki\checkmarkr\checkmarkc\checkmark$\checkmark \checkmark$\checkmark^\checkmark\\checkmarkc\checkmarki\checkmarkr\checkmarkc\checkmark$\checkmarkl\checkmark$\checkmark^\checkmark\\checkmarkc\checkmarki\checkmarkr\checkmarkc\checkmark$\checkmarka\checkmark$\checkmark^\checkmark\\checkmarkc\checkmarki\checkmarkr\checkmarkc\checkmark$\checkmark \checkmark$\checkmark^\checkmark\\checkmarkc\checkmarki\checkmarkr\checkmarkc\checkmark$\checkmarks\checkmark$\checkmark^\checkmark\\checkmarkc\checkmarki\checkmarkr\checkmarkc\checkmark$\checkmarku\checkmark$\checkmark^\checkmark\\checkmarkc\checkmarki\checkmarkr\checkmarkc\checkmark$\checkmarkr\checkmark$\checkmark^\checkmark\\checkmarkc\checkmarki\checkmarkr\checkmarkc\checkmark$\checkmarké\checkmark$\checkmark^\checkmark\\checkmarkc\checkmarki\checkmarkr\checkmarkc\checkmark$\checkmarkl\checkmark$\checkmark^\checkmark\\checkmarkc\checkmarki\checkmarkr\checkmarkc\checkmark$\checkmarké\checkmark$\checkmark^\checkmark\\checkmarkc\checkmarki\checkmarkr\checkmarkc\checkmark$\checkmarkv\checkmark$\checkmark^\checkmark\\checkmarkc\checkmarki\checkmarkr\checkmarkc\checkmark$\checkmarka\checkmark$\checkmark^\checkmark\\checkmarkc\checkmarki\checkmarkr\checkmarkc\checkmark$\checkmarkt\checkmark$\checkmark^\checkmark\\checkmarkc\checkmarki\checkmarkr\checkmarkc\checkmark$\checkmarki\checkmark$\checkmark^\checkmark\\checkmarkc\checkmarki\checkmarkr\checkmarkc\checkmark$\checkmarko\checkmark$\checkmark^\checkmark\\checkmarkc\checkmarki\checkmarkr\checkmarkc\checkmark$\checkmarkn\checkmark$\checkmark^\checkmark\\checkmarkc\checkmarki\checkmarkr\checkmarkc\checkmark$\checkmark \checkmark$\checkmark^\checkmark\\checkmarkc\checkmarki\checkmarkr\checkmarkc\checkmark$\checkmarkd\checkmark$\checkmark^\checkmark\\checkmarkc\checkmarki\checkmarkr\checkmarkc\checkmark$\checkmarku\checkmark$\checkmark^\checkmark\\checkmarkc\checkmarki\checkmarkr\checkmarkc\checkmark$\checkmark \checkmark$\checkmark^\checkmark\\checkmarkc\checkmarki\checkmarkr\checkmarkc\checkmark$\checkmarkc\checkmark$\checkmark^\checkmark\\checkmarkc\checkmarki\checkmarkr\checkmarkc\checkmark$\checkmarke\checkmark$\checkmark^\checkmark\\checkmarkc\checkmarki\checkmarkr\checkmarkc\checkmark$\checkmarkn\checkmark$\checkmark^\checkmark\\checkmarkc\checkmarki\checkmarkr\checkmarkc\checkmark$\checkmarkt\checkmark$\checkmark^\checkmark\\checkmarkc\checkmarki\checkmarkr\checkmarkc\checkmark$\checkmarkr\checkmark$\checkmark^\checkmark\\checkmarkc\checkmarki\checkmarkr\checkmarkc\checkmark$\checkmarke\checkmark$\checkmark^\checkmark\\checkmarkc\checkmarki\checkmarkr\checkmarkc\checkmark$\checkmark \checkmark$\checkmark^\checkmark\\checkmarkc\checkmarki\checkmarkr\checkmarkc\checkmark$\checkmarkd\checkmark$\checkmark^\checkmark\\checkmarkc\checkmarki\checkmarkr\checkmarkc\checkmark$\checkmarke\checkmark$\checkmark^\checkmark\\checkmarkc\checkmarki\checkmarkr\checkmarkc\checkmark$\checkmark \checkmark$\checkmark^\checkmark\\checkmarkc\checkmarki\checkmarkr\checkmarkc\checkmark$\checkmarkg\checkmark$\checkmark^\checkmark\\checkmarkc\checkmarki\checkmarkr\checkmarkc\checkmark$\checkmarkr\checkmark$\checkmark^\checkmark\\checkmarkc\checkmarki\checkmarkr\checkmarkc\checkmark$\checkmarka\checkmark$\checkmark^\checkmark\\checkmarkc\checkmarki\checkmarkr\checkmarkc\checkmark$\checkmarkv\checkmark$\checkmark^\checkmark\\checkmarkc\checkmarki\checkmarkr\checkmarkc\checkmark$\checkmarki\checkmark$\checkmark^\checkmark\\checkmarkc\checkmarki\checkmarkr\checkmarkc\checkmark$\checkmarkt\checkmark$\checkmark^\checkmark\\checkmarkc\checkmarki\checkmarkr\checkmarkc\checkmark$\checkmarké\checkmark$\checkmark^\checkmark\\checkmarkc\checkmarki\checkmarkr\checkmarkc\checkmark$\checkmark \checkmark$\checkmark^\checkmark\\checkmarkc\checkmarki\checkmarkr\checkmarkc\checkmark$\checkmarkd\checkmark$\checkmark^\checkmark\\checkmarkc\checkmarki\checkmarkr\checkmarkc\checkmark$\checkmarke\checkmark$\checkmark^\checkmark\\checkmarkc\checkmarki\checkmarkr\checkmarkc\checkmark$\checkmark \checkmark$\checkmark^\checkmark\\checkmarkc\checkmarki\checkmarkr\checkmarkc\checkmark$\checkmarkl\checkmark$\checkmark^\checkmark\\checkmarkc\checkmarki\checkmarkr\checkmarkc\checkmark$\checkmarka\checkmark$\checkmark^\checkmark\\checkmarkc\checkmarki\checkmarkr\checkmarkc\checkmark$\checkmark \checkmark$\checkmark^\checkmark\\checkmarkc\checkmarki\checkmarkr\checkmarkc\checkmark$\checkmarkp\checkmark$\checkmark^\checkmark\\checkmarkc\checkmarki\checkmarkr\checkmarkc\checkmark$\checkmarki\checkmark$\checkmark^\checkmark\\checkmarkc\checkmarki\checkmarkr\checkmarkc\checkmark$\checkmarkè\checkmark$\checkmark^\checkmark\\checkmarkc\checkmarki\checkmarkr\checkmarkc\checkmark$\checkmarkc\checkmark$\checkmark^\checkmark\\checkmarkc\checkmarki\checkmarkr\checkmarkc\checkmark$\checkmarke\checkmark$\checkmark^\checkmark\\checkmarkc\checkmarki\checkmarkr\checkmarkc\checkmark$\checkmark \checkmark$\checkmark^\checkmark\\checkmarkc\checkmarki\checkmarkr\checkmarkc\checkmark$\checkmarkp\checkmark$\checkmark^\checkmark\\checkmarkc\checkmarki\checkmarkr\checkmarkc\checkmark$\checkmarka\checkmark$\checkmark^\checkmark\\checkmarkc\checkmarki\checkmarkr\checkmarkc\checkmark$\checkmarkr\checkmark$\checkmark^\checkmark\\checkmarkc\checkmarki\checkmarkr\checkmarkc\checkmark$\checkmark \checkmark$\checkmark^\checkmark\\checkmarkc\checkmarki\checkmarkr\checkmarkc\checkmark$\checkmarkr\checkmark$\checkmark^\checkmark\\checkmarkc\checkmarki\checkmarkr\checkmarkc\checkmark$\checkmarka\checkmark$\checkmark^\checkmark\\checkmarkc\checkmarki\checkmarkr\checkmarkc\checkmark$\checkmarkp\checkmark$\checkmark^\checkmark\\checkmarkc\checkmarki\checkmarkr\checkmarkc\checkmark$\checkmarkp\checkmark$\checkmark^\checkmark\\checkmarkc\checkmarki\checkmarkr\checkmarkc\checkmark$\checkmarko\checkmark$\checkmark^\checkmark\\checkmarkc\checkmarki\checkmarkr\checkmarkc\checkmark$\checkmarkr\checkmark$\checkmark^\checkmark\\checkmarkc\checkmarki\checkmarkr\checkmarkc\checkmark$\checkmarkt\checkmark$\checkmark^\checkmark\\checkmarkc\checkmarki\checkmarkr\checkmarkc\checkmark$\checkmark \checkmark$\checkmark^\checkmark\\checkmarkc\checkmarki\checkmarkr\checkmarkc\checkmark$\checkmarka\checkmark$\checkmark^\checkmark\\checkmarkc\checkmarki\checkmarkr\checkmarkc\checkmark$\checkmarku\checkmark$\checkmark^\checkmark\\checkmarkc\checkmarki\checkmarkr\checkmarkc\checkmark$\checkmarkx\checkmark$\checkmark^\checkmark\\checkmarkc\checkmarki\checkmarkr\checkmarkc\checkmark$\checkmark \checkmark$\checkmark^\checkmark\\checkmarkc\checkmarki\checkmarkr\checkmarkc\checkmark$\checkmarkg\checkmark$\checkmark^\checkmark\\checkmarkc\checkmarki\checkmarkr\checkmarkc\checkmark$\checkmarku\checkmark$\checkmark^\checkmark\\checkmarkc\checkmarki\checkmarkr\checkmarkc\checkmark$\checkmarki\checkmark$\checkmark^\checkmark\\checkmarkc\checkmarki\checkmarkr\checkmarkc\checkmark$\checkmarkd\checkmark$\checkmark^\checkmark\\checkmarkc\checkmarki\checkmarkr\checkmarkc\checkmark$\checkmarka\checkmark$\checkmark^\checkmark\\checkmarkc\checkmarki\checkmarkr\checkmarkc\checkmark$\checkmarkg\checkmark$\checkmark^\checkmark\\checkmarkc\checkmarki\checkmarkr\checkmarkc\checkmark$\checkmarke\checkmark$\checkmark^\checkmark\\checkmarkc\checkmarki\checkmarkr\checkmarkc\checkmark$\checkmarks\checkmark$\checkmark^\checkmark\\checkmarkc\checkmarki\checkmarkr\checkmarkc\checkmark$\checkmark \checkmark$\checkmark^\checkmark\\checkmarkc\checkmarki\checkmarkr\checkmarkc\checkmark$\checkmarkY\checkmark$\checkmark^\checkmark\\checkmarkc\checkmarki\checkmarkr\checkmarkc\checkmark$\checkmark,\checkmark$\checkmark^\checkmark\\checkmarkc\checkmarki\checkmarkr\checkmarkc\checkmark$\checkmark \checkmark$\checkmark^\checkmark\\checkmarkc\checkmarki\checkmarkr\checkmarkc\checkmark$\checkmarkq\checkmark$\checkmark^\checkmark\\checkmarkc\checkmarki\checkmarkr\checkmarkc\checkmark$\checkmarku\checkmark$\checkmark^\checkmark\\checkmarkc\checkmarki\checkmarkr\checkmarkc\checkmark$\checkmarki\checkmark$\checkmark^\checkmark\\checkmarkc\checkmarki\checkmarkr\checkmarkc\checkmark$\checkmark \checkmark$\checkmark^\checkmark\\checkmarkc\checkmarki\checkmarkr\checkmarkc\checkmark$\checkmarkg\checkmark$\checkmark^\checkmark\\checkmarkc\checkmarki\checkmarkr\checkmarkc\checkmark$\checkmarké\checkmark$\checkmark^\checkmark\\checkmarkc\checkmarki\checkmarkr\checkmarkc\checkmark$\checkmarkn\checkmark$\checkmark^\checkmark\\checkmarkc\checkmarki\checkmarkr\checkmarkc\checkmark$\checkmarkè\checkmark$\checkmark^\checkmark\\checkmarkc\checkmarki\checkmarkr\checkmarkc\checkmark$\checkmarkr\checkmark$\checkmark^\checkmark\\checkmarkc\checkmarki\checkmarkr\checkmarkc\checkmark$\checkmarke\checkmark$\checkmark^\checkmark\\checkmarkc\checkmarki\checkmarkr\checkmarkc\checkmark$\checkmark \checkmark$\checkmark^\checkmark\\checkmarkc\checkmarki\checkmarkr\checkmarkc\checkmark$\checkmarku\checkmark$\checkmark^\checkmark\\checkmarkc\checkmarki\checkmarkr\checkmarkc\checkmark$\checkmarkn\checkmark$\checkmark^\checkmark\\checkmarkc\checkmarki\checkmarkr\checkmarkc\checkmark$\checkmark \checkmark$\checkmark^\checkmark\\checkmarkc\checkmarki\checkmarkr\checkmarkc\checkmark$\checkmark\\checkmark$\checkmark^\checkmark\\checkmarkc\checkmarki\checkmarkr\checkmarkc\checkmark$\checkmarkt\checkmark$\checkmark^\checkmark\\checkmarkc\checkmarki\checkmarkr\checkmarkc\checkmark$\checkmarke\checkmark$\checkmark^\checkmark\\checkmarkc\checkmarki\checkmarkr\checkmarkc\checkmark$\checkmarkx\checkmark$\checkmark^\checkmark\\checkmarkc\checkmarki\checkmarkr\checkmarkc\checkmark$\checkmarkt\checkmark$\checkmark^\checkmark\\checkmarkc\checkmarki\checkmarkr\checkmarkc\checkmark$\checkmarkb\checkmark$\checkmark^\checkmark\\checkmarkc\checkmarki\checkmarkr\checkmarkc\checkmark$\checkmarkf\checkmark$\checkmark^\checkmark\\checkmarkc\checkmarki\checkmarkr\checkmarkc\checkmark$\checkmark{\checkmark$\checkmark^\checkmark\\checkmarkc\checkmarki\checkmarkr\checkmarkc\checkmark$\checkmarkm\checkmark$\checkmark^\checkmark\\checkmarkc\checkmarki\checkmarkr\checkmarkc\checkmark$\checkmarko\checkmark$\checkmark^\checkmark\\checkmarkc\checkmarki\checkmarkr\checkmarkc\checkmark$\checkmarkm\checkmark$\checkmark^\checkmark\\checkmarkc\checkmarki\checkmarkr\checkmarkc\checkmark$\checkmarke\checkmark$\checkmark^\checkmark\\checkmarkc\checkmarki\checkmarkr\checkmarkc\checkmark$\checkmarkn\checkmark$\checkmark^\checkmark\\checkmarkc\checkmarki\checkmarkr\checkmarkc\checkmark$\checkmarkt\checkmark$\checkmark^\checkmark\\checkmarkc\checkmarki\checkmarkr\checkmarkc\checkmark$\checkmark \checkmark$\checkmark^\checkmark\\checkmarkc\checkmarki\checkmarkr\checkmarkc\checkmark$\checkmarkd\checkmark$\checkmark^\checkmark\\checkmarkc\checkmarki\checkmarkr\checkmarkc\checkmark$\checkmarke\checkmark$\checkmark^\checkmark\\checkmarkc\checkmarki\checkmarkr\checkmarkc\checkmark$\checkmark \checkmark$\checkmark^\checkmark\\checkmarkc\checkmarki\checkmarkr\checkmarkc\checkmark$\checkmarkr\checkmark$\checkmark^\checkmark\\checkmarkc\checkmarki\checkmarkr\checkmarkc\checkmark$\checkmarke\checkmark$\checkmark^\checkmark\\checkmarkc\checkmarki\checkmarkr\checkmarkc\checkmark$\checkmarkn\checkmark$\checkmark^\checkmark\\checkmarkc\checkmarki\checkmarkr\checkmarkc\checkmark$\checkmarkv\checkmark$\checkmark^\checkmark\\checkmarkc\checkmarki\checkmarkr\checkmarkc\checkmark$\checkmarke\checkmark$\checkmark^\checkmark\\checkmarkc\checkmarki\checkmarkr\checkmarkc\checkmark$\checkmarkr\checkmark$\checkmark^\checkmark\\checkmarkc\checkmarki\checkmarkr\checkmarkc\checkmark$\checkmarks\checkmark$\checkmark^\checkmark\\checkmarkc\checkmarki\checkmarkr\checkmarkc\checkmark$\checkmarke\checkmark$\checkmark^\checkmark\\checkmarkc\checkmarki\checkmarkr\checkmarkc\checkmark$\checkmarkm\checkmark$\checkmark^\checkmark\\checkmarkc\checkmarki\checkmarkr\checkmarkc\checkmark$\checkmarke\checkmark$\checkmark^\checkmark\\checkmarkc\checkmarki\checkmarkr\checkmarkc\checkmark$\checkmarkn\checkmark$\checkmark^\checkmark\\checkmarkc\checkmarki\checkmarkr\checkmarkc\checkmark$\checkmarkt\checkmark$\checkmark^\checkmark\\checkmarkc\checkmarki\checkmarkr\checkmarkc\checkmark$\checkmark}\checkmark$\checkmark^\checkmark\\checkmarkc\checkmarki\checkmarkr\checkmarkc\checkmark$\checkmark \checkmark$\checkmark^\checkmark\\checkmarkc\checkmarki\checkmarkr\checkmarkc\checkmark$\checkmarkq\checkmark$\checkmark^\checkmark\\checkmarkc\checkmarki\checkmarkr\checkmarkc\checkmark$\checkmarku\checkmark$\checkmark^\checkmark\\checkmarkc\checkmarki\checkmarkr\checkmarkc\checkmark$\checkmark'\checkmark$\checkmark^\checkmark\\checkmarkc\checkmarki\checkmarkr\checkmarkc\checkmark$\checkmarki\checkmark$\checkmark^\checkmark\\checkmarkc\checkmarki\checkmarkr\checkmarkc\checkmark$\checkmarkl\checkmark$\checkmark^\checkmark\\checkmarkc\checkmarki\checkmarkr\checkmarkc\checkmark$\checkmark \checkmark$\checkmark^\checkmark\\checkmarkc\checkmarki\checkmarkr\checkmarkc\checkmark$\checkmarkc\checkmark$\checkmark^\checkmark\\checkmarkc\checkmarki\checkmarkr\checkmarkc\checkmark$\checkmarko\checkmark$\checkmark^\checkmark\\checkmarkc\checkmarki\checkmarkr\checkmarkc\checkmark$\checkmarkn\checkmark$\checkmark^\checkmark\\checkmarkc\checkmarki\checkmarkr\checkmarkc\checkmark$\checkmarkv\checkmark$\checkmark^\checkmark\\checkmarkc\checkmarki\checkmarkr\checkmarkc\checkmark$\checkmarki\checkmark$\checkmark^\checkmark\\checkmarkc\checkmarki\checkmarkr\checkmarkc\checkmark$\checkmarke\checkmark$\checkmark^\checkmark\\checkmarkc\checkmarki\checkmarkr\checkmarkc\checkmark$\checkmarkn\checkmark$\checkmark^\checkmark\\checkmarkc\checkmarki\checkmarkr\checkmarkc\checkmark$\checkmarkt\checkmark$\checkmark^\checkmark\\checkmarkc\checkmarki\checkmarkr\checkmarkc\checkmark$\checkmark \checkmark$\checkmark^\checkmark\\checkmarkc\checkmarki\checkmarkr\checkmarkc\checkmark$\checkmarkd\checkmark$\checkmark^\checkmark\\checkmarkc\checkmarki\checkmarkr\checkmarkc\checkmark$\checkmarke\checkmark$\checkmark^\checkmark\\checkmarkc\checkmarki\checkmarkr\checkmarkc\checkmark$\checkmark \checkmark$\checkmark^\checkmark\\checkmarkc\checkmarki\checkmarkr\checkmarkc\checkmark$\checkmarkd\checkmark$\checkmark^\checkmark\\checkmarkc\checkmarki\checkmarkr\checkmarkc\checkmark$\checkmarki\checkmark$\checkmark^\checkmark\\checkmarkc\checkmarki\checkmarkr\checkmarkc\checkmark$\checkmarkm\checkmark$\checkmark^\checkmark\\checkmarkc\checkmarki\checkmarkr\checkmarkc\checkmark$\checkmarke\checkmark$\checkmark^\checkmark\\checkmarkc\checkmarki\checkmarkr\checkmarkc\checkmark$\checkmarkn\checkmark$\checkmark^\checkmark\\checkmarkc\checkmarki\checkmarkr\checkmarkc\checkmark$\checkmarks\checkmark$\checkmark^\checkmark\\checkmarkc\checkmarki\checkmarkr\checkmarkc\checkmark$\checkmarki\checkmark$\checkmark^\checkmark\\checkmarkc\checkmarki\checkmarkr\checkmarkc\checkmark$\checkmarko\checkmark$\checkmark^\checkmark\\checkmarkc\checkmarki\checkmarkr\checkmarkc\checkmark$\checkmarkn\checkmark$\checkmark^\checkmark\\checkmarkc\checkmarki\checkmarkr\checkmarkc\checkmark$\checkmarkn\checkmark$\checkmark^\checkmark\\checkmarkc\checkmarki\checkmarkr\checkmarkc\checkmark$\checkmarke\checkmark$\checkmark^\checkmark\\checkmarkc\checkmarki\checkmarkr\checkmarkc\checkmark$\checkmarkr\checkmark$\checkmark^\checkmark\\checkmarkc\checkmarki\checkmarkr\checkmarkc\checkmark$\checkmark \checkmark$\checkmark^\checkmark\\checkmarkc\checkmarki\checkmarkr\checkmarkc\checkmark$\checkmarkr\checkmark$\checkmark^\checkmark\\checkmarkc\checkmarki\checkmarkr\checkmarkc\checkmark$\checkmarki\checkmark$\checkmark^\checkmark\\checkmarkc\checkmarki\checkmarkr\checkmarkc\checkmark$\checkmarkg\checkmark$\checkmark^\checkmark\\checkmarkc\checkmarki\checkmarkr\checkmarkc\checkmark$\checkmarko\checkmark$\checkmark^\checkmark\\checkmarkc\checkmarki\checkmarkr\checkmarkc\checkmark$\checkmarku\checkmark$\checkmark^\checkmark\\checkmarkc\checkmarki\checkmarkr\checkmarkc\checkmark$\checkmarkr\checkmark$\checkmark^\checkmark\\checkmarkc\checkmarki\checkmarkr\checkmarkc\checkmark$\checkmarke\checkmark$\checkmark^\checkmark\\checkmarkc\checkmarki\checkmarkr\checkmarkc\checkmark$\checkmarku\checkmark$\checkmark^\checkmark\\checkmarkc\checkmarki\checkmarkr\checkmarkc\checkmark$\checkmarks\checkmark$\checkmark^\checkmark\\checkmarkc\checkmarki\checkmarkr\checkmarkc\checkmark$\checkmarke\checkmark$\checkmark^\checkmark\\checkmarkc\checkmarki\checkmarkr\checkmarkc\checkmark$\checkmarkm\checkmark$\checkmark^\checkmark\\checkmarkc\checkmarki\checkmarkr\checkmarkc\checkmark$\checkmarke\checkmark$\checkmark^\checkmark\\checkmarkc\checkmarki\checkmarkr\checkmarkc\checkmark$\checkmarkn\checkmark$\checkmark^\checkmark\\checkmarkc\checkmarki\checkmarkr\checkmarkc\checkmark$\checkmarkt\checkmark$\checkmark^\checkmark\\checkmarkc\checkmarki\checkmarkr\checkmarkc\checkmark$\checkmark.\checkmark$\checkmark^\checkmark\\checkmarkc\checkmarki\checkmarkr\checkmarkc\checkmark$\checkmark
\checkmark$\checkmark^\checkmark\\checkmarkc\checkmarki\checkmarkr\checkmarkc\checkmark$\checkmark
\checkmark$\checkmark^\checkmark\\checkmarkc\checkmarki\checkmarkr\checkmarkc\checkmark$\checkmark\\checkmark$\checkmark^\checkmark\\checkmarkc\checkmarki\checkmarkr\checkmarkc\checkmark$\checkmarks\checkmark$\checkmark^\checkmark\\checkmarkc\checkmarki\checkmarkr\checkmarkc\checkmark$\checkmarku\checkmark$\checkmark^\checkmark\\checkmarkc\checkmarki\checkmarkr\checkmarkc\checkmark$\checkmarkb\checkmark$\checkmark^\checkmark\\checkmarkc\checkmarki\checkmarkr\checkmarkc\checkmark$\checkmarks\checkmark$\checkmark^\checkmark\\checkmarkc\checkmarki\checkmarkr\checkmarkc\checkmark$\checkmarke\checkmark$\checkmark^\checkmark\\checkmarkc\checkmarki\checkmarkr\checkmarkc\checkmark$\checkmarkc\checkmark$\checkmark^\checkmark\\checkmarkc\checkmarki\checkmarkr\checkmarkc\checkmark$\checkmarkt\checkmark$\checkmark^\checkmark\\checkmarkc\checkmarki\checkmarkr\checkmarkc\checkmark$\checkmarki\checkmark$\checkmark^\checkmark\\checkmarkc\checkmarki\checkmarkr\checkmarkc\checkmark$\checkmarko\checkmark$\checkmark^\checkmark\\checkmarkc\checkmarki\checkmarkr\checkmarkc\checkmark$\checkmarkn\checkmark$\checkmark^\checkmark\\checkmarkc\checkmarki\checkmarkr\checkmarkc\checkmark$\checkmark{\checkmark$\checkmark^\checkmark\\checkmarkc\checkmarki\checkmarkr\checkmarkc\checkmark$\checkmarkS\checkmark$\checkmark^\checkmark\\checkmarkc\checkmarki\checkmarkr\checkmarkc\checkmark$\checkmarkc\checkmark$\checkmark^\checkmark\\checkmarkc\checkmarki\checkmarkr\checkmarkc\checkmark$\checkmarkh\checkmark$\checkmark^\checkmark\\checkmarkc\checkmarki\checkmarkr\checkmarkc\checkmark$\checkmarké\checkmark$\checkmark^\checkmark\\checkmarkc\checkmarki\checkmarkr\checkmarkc\checkmark$\checkmarkm\checkmark$\checkmark^\checkmark\\checkmarkc\checkmarki\checkmarkr\checkmarkc\checkmark$\checkmarka\checkmark$\checkmark^\checkmark\\checkmarkc\checkmarki\checkmarkr\checkmarkc\checkmark$\checkmark \checkmark$\checkmark^\checkmark\\checkmarkc\checkmarki\checkmarkr\checkmarkc\checkmark$\checkmarkc\checkmark$\checkmark^\checkmark\\checkmarkc\checkmarki\checkmarkr\checkmarkc\checkmark$\checkmarki\checkmark$\checkmark^\checkmark\\checkmarkc\checkmarki\checkmarkr\checkmarkc\checkmark$\checkmarkn\checkmark$\checkmark^\checkmark\\checkmarkc\checkmarki\checkmarkr\checkmarkc\checkmark$\checkmarké\checkmark$\checkmark^\checkmark\\checkmarkc\checkmarki\checkmarkr\checkmarkc\checkmark$\checkmarkm\checkmark$\checkmark^\checkmark\\checkmarkc\checkmarki\checkmarkr\checkmarkc\checkmark$\checkmarka\checkmark$\checkmark^\checkmark\\checkmarkc\checkmarki\checkmarkr\checkmarkc\checkmark$\checkmarkt\checkmark$\checkmark^\checkmark\\checkmarkc\checkmarki\checkmarkr\checkmarkc\checkmark$\checkmarki\checkmark$\checkmark^\checkmark\\checkmarkc\checkmarki\checkmarkr\checkmarkc\checkmark$\checkmarkq\checkmark$\checkmark^\checkmark\\checkmarkc\checkmarki\checkmarkr\checkmarkc\checkmark$\checkmarku\checkmark$\checkmark^\checkmark\\checkmarkc\checkmarki\checkmarkr\checkmarkc\checkmark$\checkmarke\checkmark$\checkmark^\checkmark\\checkmarkc\checkmarki\checkmarkr\checkmarkc\checkmark$\checkmark \checkmark$\checkmark^\checkmark\\checkmarkc\checkmarki\checkmarkr\checkmarkc\checkmark$\checkmarkg\checkmark$\checkmark^\checkmark\\checkmarkc\checkmarki\checkmarkr\checkmarkc\checkmark$\checkmarkl\checkmark$\checkmark^\checkmark\\checkmarkc\checkmarki\checkmarkr\checkmarkc\checkmark$\checkmarko\checkmark$\checkmark^\checkmark\\checkmarkc\checkmarki\checkmarkr\checkmarkc\checkmark$\checkmarkb\checkmark$\checkmark^\checkmark\\checkmarkc\checkmarki\checkmarkr\checkmarkc\checkmark$\checkmarka\checkmark$\checkmark^\checkmark\\checkmarkc\checkmarki\checkmarkr\checkmarkc\checkmark$\checkmarkl\checkmark$\checkmark^\checkmark\\checkmarkc\checkmarki\checkmarkr\checkmarkc\checkmark$\checkmark}\checkmark$\checkmark^\checkmark\\checkmarkc\checkmarki\checkmarkr\checkmarkc\checkmark$\checkmark
\checkmark$\checkmark^\checkmark\\checkmarkc\checkmarki\checkmarkr\checkmarkc\checkmark$\checkmark\\checkmark$\checkmark^\checkmark\\checkmarkc\checkmarki\checkmarkr\checkmarkc\checkmark$\checkmarkb\checkmark$\checkmark^\checkmark\\checkmarkc\checkmarki\checkmarkr\checkmarkc\checkmark$\checkmarke\checkmark$\checkmark^\checkmark\\checkmarkc\checkmarki\checkmarkr\checkmarkc\checkmark$\checkmarkg\checkmark$\checkmark^\checkmark\\checkmarkc\checkmarki\checkmarkr\checkmarkc\checkmark$\checkmarki\checkmark$\checkmark^\checkmark\\checkmarkc\checkmarki\checkmarkr\checkmarkc\checkmark$\checkmarkn\checkmark$\checkmark^\checkmark\\checkmarkc\checkmarki\checkmarkr\checkmarkc\checkmark$\checkmark{\checkmark$\checkmark^\checkmark\\checkmarkc\checkmarki\checkmarkr\checkmarkc\checkmark$\checkmarkf\checkmark$\checkmark^\checkmark\\checkmarkc\checkmarki\checkmarkr\checkmarkc\checkmark$\checkmarki\checkmark$\checkmark^\checkmark\\checkmarkc\checkmarki\checkmarkr\checkmarkc\checkmark$\checkmarkg\checkmark$\checkmark^\checkmark\\checkmarkc\checkmarki\checkmarkr\checkmarkc\checkmark$\checkmarku\checkmark$\checkmark^\checkmark\\checkmarkc\checkmarki\checkmarkr\checkmarkc\checkmark$\checkmarkr\checkmark$\checkmark^\checkmark\\checkmarkc\checkmarki\checkmarkr\checkmarkc\checkmark$\checkmarke\checkmark$\checkmark^\checkmark\\checkmarkc\checkmarki\checkmarkr\checkmarkc\checkmark$\checkmark}\checkmark$\checkmark^\checkmark\\checkmarkc\checkmarki\checkmarkr\checkmarkc\checkmark$\checkmark[\checkmark$\checkmark^\checkmark\\checkmarkc\checkmarki\checkmarkr\checkmarkc\checkmark$\checkmarkh\checkmark$\checkmark^\checkmark\\checkmarkc\checkmarki\checkmarkr\checkmarkc\checkmark$\checkmarkt\checkmark$\checkmark^\checkmark\\checkmarkc\checkmarki\checkmarkr\checkmarkc\checkmark$\checkmarkb\checkmark$\checkmark^\checkmark\\checkmarkc\checkmarki\checkmarkr\checkmarkc\checkmark$\checkmarkp\checkmark$\checkmark^\checkmark\\checkmarkc\checkmarki\checkmarkr\checkmarkc\checkmark$\checkmark]\checkmark$\checkmark^\checkmark\\checkmarkc\checkmarki\checkmarkr\checkmarkc\checkmark$\checkmark
\checkmark$\checkmark^\checkmark\\checkmarkc\checkmarki\checkmarkr\checkmarkc\checkmark$\checkmark \checkmark$\checkmark^\checkmark\\checkmarkc\checkmarki\checkmarkr\checkmarkc\checkmark$\checkmark \checkmark$\checkmark^\checkmark\\checkmarkc\checkmarki\checkmarkr\checkmarkc\checkmark$\checkmark \checkmark$\checkmark^\checkmark\\checkmarkc\checkmarki\checkmarkr\checkmarkc\checkmark$\checkmark \checkmark$\checkmark^\checkmark\\checkmarkc\checkmarki\checkmarkr\checkmarkc\checkmark$\checkmark\\checkmark$\checkmark^\checkmark\\checkmarkc\checkmarki\checkmarkr\checkmarkc\checkmark$\checkmarkc\checkmark$\checkmark^\checkmark\\checkmarkc\checkmarki\checkmarkr\checkmarkc\checkmark$\checkmarke\checkmark$\checkmark^\checkmark\\checkmarkc\checkmarki\checkmarkr\checkmarkc\checkmark$\checkmarkn\checkmark$\checkmark^\checkmark\\checkmarkc\checkmarki\checkmarkr\checkmarkc\checkmark$\checkmarkt\checkmark$\checkmark^\checkmark\\checkmarkc\checkmarki\checkmarkr\checkmarkc\checkmark$\checkmarke\checkmark$\checkmark^\checkmark\\checkmarkc\checkmarki\checkmarkr\checkmarkc\checkmark$\checkmarkr\checkmark$\checkmark^\checkmark\\checkmarkc\checkmarki\checkmarkr\checkmarkc\checkmark$\checkmarki\checkmark$\checkmark^\checkmark\\checkmarkc\checkmarki\checkmarkr\checkmarkc\checkmark$\checkmarkn\checkmark$\checkmark^\checkmark\\checkmarkc\checkmarki\checkmarkr\checkmarkc\checkmark$\checkmarkg\checkmark$\checkmark^\checkmark\\checkmarkc\checkmarki\checkmarkr\checkmarkc\checkmark$\checkmark
\checkmark$\checkmark^\checkmark\\checkmarkc\checkmarki\checkmarkr\checkmarkc\checkmark$\checkmark \checkmark$\checkmark^\checkmark\\checkmarkc\checkmarki\checkmarkr\checkmarkc\checkmark$\checkmark \checkmark$\checkmark^\checkmark\\checkmarkc\checkmarki\checkmarkr\checkmarkc\checkmark$\checkmark \checkmark$\checkmark^\checkmark\\checkmarkc\checkmarki\checkmarkr\checkmarkc\checkmark$\checkmark \checkmark$\checkmark^\checkmark\\checkmarkc\checkmarki\checkmarkr\checkmarkc\checkmark$\checkmark%\checkmark$\checkmark^\checkmark\\checkmarkc\checkmarki\checkmarkr\checkmarkc\checkmark$\checkmark \checkmark$\checkmark^\checkmark\\checkmarkc\checkmarki\checkmarkr\checkmarkc\checkmark$\checkmarkP\checkmark$\checkmark^\checkmark\\checkmarkc\checkmarki\checkmarkr\checkmarkc\checkmark$\checkmarkl\checkmark$\checkmark^\checkmark\\checkmarkc\checkmarki\checkmarkr\checkmarkc\checkmark$\checkmarka\checkmark$\checkmark^\checkmark\\checkmarkc\checkmarki\checkmarkr\checkmarkc\checkmark$\checkmarkc\checkmark$\checkmark^\checkmark\\checkmarkc\checkmarki\checkmarkr\checkmarkc\checkmark$\checkmarke\checkmark$\checkmark^\checkmark\\checkmarkc\checkmarki\checkmarkr\checkmarkc\checkmark$\checkmarkh\checkmark$\checkmark^\checkmark\\checkmarkc\checkmarki\checkmarkr\checkmarkc\checkmark$\checkmarko\checkmark$\checkmark^\checkmark\\checkmarkc\checkmarki\checkmarkr\checkmarkc\checkmark$\checkmarkl\checkmark$\checkmark^\checkmark\\checkmarkc\checkmarki\checkmarkr\checkmarkc\checkmark$\checkmarkd\checkmark$\checkmark^\checkmark\\checkmarkc\checkmarki\checkmarkr\checkmarkc\checkmark$\checkmarke\checkmark$\checkmark^\checkmark\\checkmarkc\checkmarki\checkmarkr\checkmarkc\checkmark$\checkmarkr\checkmark$\checkmark^\checkmark\\checkmarkc\checkmarki\checkmarkr\checkmarkc\checkmark$\checkmark \checkmark$\checkmark^\checkmark\\checkmarkc\checkmarki\checkmarkr\checkmarkc\checkmark$\checkmarkp\checkmark$\checkmark^\checkmark\\checkmarkc\checkmarki\checkmarkr\checkmarkc\checkmark$\checkmarko\checkmark$\checkmark^\checkmark\\checkmarkc\checkmarki\checkmarkr\checkmarkc\checkmark$\checkmarku\checkmark$\checkmark^\checkmark\\checkmarkc\checkmarki\checkmarkr\checkmarkc\checkmark$\checkmarkr\checkmark$\checkmark^\checkmark\\checkmarkc\checkmarki\checkmarkr\checkmarkc\checkmark$\checkmark \checkmark$\checkmark^\checkmark\\checkmarkc\checkmarki\checkmarkr\checkmarkc\checkmark$\checkmarkl\checkmark$\checkmark^\checkmark\\checkmarkc\checkmarki\checkmarkr\checkmarkc\checkmark$\checkmarke\checkmark$\checkmark^\checkmark\\checkmarkc\checkmarki\checkmarkr\checkmarkc\checkmark$\checkmark \checkmark$\checkmark^\checkmark\\checkmarkc\checkmarki\checkmarkr\checkmarkc\checkmark$\checkmarks\checkmark$\checkmark^\checkmark\\checkmarkc\checkmarki\checkmarkr\checkmarkc\checkmark$\checkmarkc\checkmark$\checkmark^\checkmark\\checkmarkc\checkmarki\checkmarkr\checkmarkc\checkmark$\checkmarkh\checkmark$\checkmark^\checkmark\\checkmarkc\checkmarki\checkmarkr\checkmarkc\checkmark$\checkmarké\checkmark$\checkmark^\checkmark\\checkmarkc\checkmarki\checkmarkr\checkmarkc\checkmark$\checkmarkm\checkmark$\checkmark^\checkmark\\checkmarkc\checkmarki\checkmarkr\checkmarkc\checkmark$\checkmarka\checkmark$\checkmark^\checkmark\\checkmarkc\checkmarki\checkmarkr\checkmarkc\checkmark$\checkmark \checkmark$\checkmark^\checkmark\\checkmarkc\checkmarki\checkmarkr\checkmarkc\checkmark$\checkmarkc\checkmark$\checkmark^\checkmark\\checkmarkc\checkmarki\checkmarkr\checkmarkc\checkmark$\checkmarki\checkmark$\checkmark^\checkmark\\checkmarkc\checkmarki\checkmarkr\checkmarkc\checkmark$\checkmarkn\checkmark$\checkmark^\checkmark\\checkmarkc\checkmarki\checkmarkr\checkmarkc\checkmark$\checkmarké\checkmark$\checkmark^\checkmark\\checkmarkc\checkmarki\checkmarkr\checkmarkc\checkmark$\checkmarkm\checkmark$\checkmark^\checkmark\\checkmarkc\checkmarki\checkmarkr\checkmarkc\checkmark$\checkmarka\checkmark$\checkmark^\checkmark\\checkmarkc\checkmarki\checkmarkr\checkmarkc\checkmark$\checkmarkt\checkmark$\checkmark^\checkmark\\checkmarkc\checkmarki\checkmarkr\checkmarkc\checkmark$\checkmarki\checkmark$\checkmark^\checkmark\\checkmarkc\checkmarki\checkmarkr\checkmarkc\checkmark$\checkmarkq\checkmark$\checkmark^\checkmark\\checkmarkc\checkmarki\checkmarkr\checkmarkc\checkmark$\checkmarku\checkmark$\checkmark^\checkmark\\checkmarkc\checkmarki\checkmarkr\checkmarkc\checkmark$\checkmarke\checkmark$\checkmark^\checkmark\\checkmarkc\checkmarki\checkmarkr\checkmarkc\checkmark$\checkmark \checkmark$\checkmark^\checkmark\\checkmarkc\checkmarki\checkmarkr\checkmarkc\checkmark$\checkmarkd\checkmark$\checkmark^\checkmark\\checkmarkc\checkmarki\checkmarkr\checkmarkc\checkmark$\checkmarke\checkmark$\checkmark^\checkmark\\checkmarkc\checkmarki\checkmarkr\checkmarkc\checkmark$\checkmark \checkmark$\checkmark^\checkmark\\checkmarkc\checkmarki\checkmarkr\checkmarkc\checkmark$\checkmarkl\checkmark$\checkmark^\checkmark\\checkmarkc\checkmarki\checkmarkr\checkmarkc\checkmark$\checkmarka\checkmark$\checkmark^\checkmark\\checkmarkc\checkmarki\checkmarkr\checkmarkc\checkmark$\checkmark \checkmark$\checkmark^\checkmark\\checkmarkc\checkmarki\checkmarkr\checkmarkc\checkmark$\checkmarkm\checkmark$\checkmark^\checkmark\\checkmarkc\checkmarki\checkmarkr\checkmarkc\checkmark$\checkmarka\checkmark$\checkmark^\checkmark\\checkmarkc\checkmarki\checkmarkr\checkmarkc\checkmark$\checkmarkc\checkmark$\checkmark^\checkmark\\checkmarkc\checkmarki\checkmarkr\checkmarkc\checkmark$\checkmarkh\checkmark$\checkmark^\checkmark\\checkmarkc\checkmarki\checkmarkr\checkmarkc\checkmark$\checkmarki\checkmark$\checkmark^\checkmark\\checkmarkc\checkmarki\checkmarkr\checkmarkc\checkmark$\checkmarkn\checkmark$\checkmark^\checkmark\\checkmarkc\checkmarki\checkmarkr\checkmarkc\checkmark$\checkmarke\checkmark$\checkmark^\checkmark\\checkmarkc\checkmarki\checkmarkr\checkmarkc\checkmark$\checkmark \checkmark$\checkmark^\checkmark\\checkmarkc\checkmarki\checkmarkr\checkmarkc\checkmark$\checkmark(\checkmark$\checkmark^\checkmark\\checkmarkc\checkmarki\checkmarkr\checkmarkc\checkmark$\checkmarki\checkmark$\checkmark^\checkmark\\checkmarkc\checkmarki\checkmarkr\checkmarkc\checkmark$\checkmarks\checkmark$\checkmark^\checkmark\\checkmarkc\checkmarki\checkmarkr\checkmarkc\checkmark$\checkmarks\checkmark$\checkmark^\checkmark\\checkmarkc\checkmarki\checkmarkr\checkmarkc\checkmark$\checkmarku\checkmark$\checkmark^\checkmark\\checkmarkc\checkmarki\checkmarkr\checkmarkc\checkmark$\checkmark \checkmark$\checkmark^\checkmark\\checkmarkc\checkmarki\checkmarkr\checkmarkc\checkmark$\checkmarkd\checkmark$\checkmark^\checkmark\\checkmarkc\checkmarki\checkmarkr\checkmarkc\checkmark$\checkmarke\checkmark$\checkmark^\checkmark\\checkmarkc\checkmarki\checkmarkr\checkmarkc\checkmark$\checkmark \checkmark$\checkmark^\checkmark\\checkmarkc\checkmarki\checkmarkr\checkmarkc\checkmark$\checkmarkl\checkmark$\checkmark^\checkmark\\checkmarkc\checkmarki\checkmarkr\checkmarkc\checkmark$\checkmarka\checkmark$\checkmark^\checkmark\\checkmarkc\checkmarki\checkmarkr\checkmarkc\checkmark$\checkmark \checkmark$\checkmark^\checkmark\\checkmarkc\checkmarki\checkmarkr\checkmarkc\checkmark$\checkmarkC\checkmark$\checkmark^\checkmark\\checkmarkc\checkmarki\checkmarkr\checkmarkc\checkmark$\checkmarkA\checkmark$\checkmark^\checkmark\\checkmarkc\checkmarki\checkmarkr\checkmarkc\checkmark$\checkmarkO\checkmark$\checkmark^\checkmark\\checkmarkc\checkmarki\checkmarkr\checkmarkc\checkmark$\checkmark \checkmark$\checkmark^\checkmark\\checkmarkc\checkmarki\checkmarkr\checkmarkc\checkmark$\checkmarkS\checkmark$\checkmark^\checkmark\\checkmarkc\checkmarki\checkmarkr\checkmarkc\checkmark$\checkmarko\checkmark$\checkmark^\checkmark\\checkmarkc\checkmarki\checkmarkr\checkmarkc\checkmark$\checkmarkl\checkmark$\checkmark^\checkmark\\checkmarkc\checkmarki\checkmarkr\checkmarkc\checkmark$\checkmarki\checkmark$\checkmark^\checkmark\\checkmarkc\checkmarki\checkmarkr\checkmarkc\checkmark$\checkmarkd\checkmark$\checkmark^\checkmark\\checkmarkc\checkmarki\checkmarkr\checkmarkc\checkmark$\checkmarkW\checkmark$\checkmark^\checkmark\\checkmarkc\checkmarki\checkmarkr\checkmarkc\checkmark$\checkmarko\checkmark$\checkmark^\checkmark\\checkmarkc\checkmarki\checkmarkr\checkmarkc\checkmark$\checkmarkr\checkmark$\checkmark^\checkmark\\checkmarkc\checkmarki\checkmarkr\checkmarkc\checkmark$\checkmarkk\checkmark$\checkmark^\checkmark\\checkmarkc\checkmarki\checkmarkr\checkmarkc\checkmark$\checkmarks\checkmark$\checkmark^\checkmark\\checkmarkc\checkmarki\checkmarkr\checkmarkc\checkmark$\checkmark)\checkmark$\checkmark^\checkmark\\checkmarkc\checkmarki\checkmarkr\checkmarkc\checkmark$\checkmark
\checkmark$\checkmark^\checkmark\\checkmarkc\checkmarki\checkmarkr\checkmarkc\checkmark$\checkmark \checkmark$\checkmark^\checkmark\\checkmarkc\checkmarki\checkmarkr\checkmarkc\checkmark$\checkmark \checkmark$\checkmark^\checkmark\\checkmarkc\checkmarki\checkmarkr\checkmarkc\checkmark$\checkmark \checkmark$\checkmark^\checkmark\\checkmarkc\checkmarki\checkmarkr\checkmarkc\checkmark$\checkmark \checkmark$\checkmark^\checkmark\\checkmarkc\checkmarki\checkmarkr\checkmarkc\checkmark$\checkmark\\checkmark$\checkmark^\checkmark\\checkmarkc\checkmarki\checkmarkr\checkmarkc\checkmark$\checkmarkf\checkmark$\checkmark^\checkmark\\checkmarkc\checkmarki\checkmarkr\checkmarkc\checkmark$\checkmarkb\checkmark$\checkmark^\checkmark\\checkmarkc\checkmarki\checkmarkr\checkmarkc\checkmark$\checkmarko\checkmark$\checkmark^\checkmark\\checkmarkc\checkmarki\checkmarkr\checkmarkc\checkmark$\checkmarkx\checkmark$\checkmark^\checkmark\\checkmarkc\checkmarki\checkmarkr\checkmarkc\checkmark$\checkmark{\checkmark$\checkmark^\checkmark\\checkmarkc\checkmarki\checkmarkr\checkmarkc\checkmark$\checkmark\\checkmark$\checkmark^\checkmark\\checkmarkc\checkmarki\checkmarkr\checkmarkc\checkmark$\checkmarkp\checkmark$\checkmark^\checkmark\\checkmarkc\checkmarki\checkmarkr\checkmarkc\checkmark$\checkmarka\checkmark$\checkmark^\checkmark\\checkmarkc\checkmarki\checkmarkr\checkmarkc\checkmark$\checkmarkr\checkmark$\checkmark^\checkmark\\checkmarkc\checkmarki\checkmarkr\checkmarkc\checkmark$\checkmarkb\checkmark$\checkmark^\checkmark\\checkmarkc\checkmarki\checkmarkr\checkmarkc\checkmark$\checkmarko\checkmark$\checkmark^\checkmark\\checkmarkc\checkmarki\checkmarkr\checkmarkc\checkmark$\checkmarkx\checkmark$\checkmark^\checkmark\\checkmarkc\checkmarki\checkmarkr\checkmarkc\checkmark$\checkmark{\checkmark$\checkmark^\checkmark\\checkmarkc\checkmarki\checkmarkr\checkmarkc\checkmark$\checkmark0\checkmark$\checkmark^\checkmark\\checkmarkc\checkmarki\checkmarkr\checkmarkc\checkmark$\checkmark.\checkmark$\checkmark^\checkmark\\checkmarkc\checkmarki\checkmarkr\checkmarkc\checkmark$\checkmark8\checkmark$\checkmark^\checkmark\\checkmarkc\checkmarki\checkmarkr\checkmarkc\checkmark$\checkmark5\checkmark$\checkmark^\checkmark\\checkmarkc\checkmarki\checkmarkr\checkmarkc\checkmark$\checkmark\\checkmark$\checkmark^\checkmark\\checkmarkc\checkmarki\checkmarkr\checkmarkc\checkmark$\checkmarkt\checkmark$\checkmark^\checkmark\\checkmarkc\checkmarki\checkmarkr\checkmarkc\checkmark$\checkmarke\checkmark$\checkmark^\checkmark\\checkmarkc\checkmarki\checkmarkr\checkmarkc\checkmark$\checkmarkx\checkmark$\checkmark^\checkmark\\checkmarkc\checkmarki\checkmarkr\checkmarkc\checkmark$\checkmarkt\checkmark$\checkmark^\checkmark\\checkmarkc\checkmarki\checkmarkr\checkmarkc\checkmark$\checkmarkw\checkmark$\checkmark^\checkmark\\checkmarkc\checkmarki\checkmarkr\checkmarkc\checkmark$\checkmarki\checkmark$\checkmark^\checkmark\\checkmarkc\checkmarki\checkmarkr\checkmarkc\checkmark$\checkmarkd\checkmark$\checkmark^\checkmark\\checkmarkc\checkmarki\checkmarkr\checkmarkc\checkmark$\checkmarkt\checkmark$\checkmark^\checkmark\\checkmarkc\checkmarki\checkmarkr\checkmarkc\checkmark$\checkmarkh\checkmark$\checkmark^\checkmark\\checkmarkc\checkmarki\checkmarkr\checkmarkc\checkmark$\checkmark}\checkmark$\checkmark^\checkmark\\checkmarkc\checkmarki\checkmarkr\checkmarkc\checkmark$\checkmark{\checkmark$\checkmark^\checkmark\\checkmarkc\checkmarki\checkmarkr\checkmarkc\checkmark$\checkmark\\checkmark$\checkmark^\checkmark\\checkmarkc\checkmarki\checkmarkr\checkmarkc\checkmark$\checkmarkc\checkmark$\checkmark^\checkmark\\checkmarkc\checkmarki\checkmarkr\checkmarkc\checkmark$\checkmarke\checkmark$\checkmark^\checkmark\\checkmarkc\checkmarki\checkmarkr\checkmarkc\checkmark$\checkmarkn\checkmark$\checkmark^\checkmark\\checkmarkc\checkmarki\checkmarkr\checkmarkc\checkmark$\checkmarkt\checkmark$\checkmark^\checkmark\\checkmarkc\checkmarki\checkmarkr\checkmarkc\checkmark$\checkmarke\checkmark$\checkmark^\checkmark\\checkmarkc\checkmarki\checkmarkr\checkmarkc\checkmark$\checkmarkr\checkmark$\checkmark^\checkmark\\checkmarkc\checkmarki\checkmarkr\checkmarkc\checkmark$\checkmarki\checkmark$\checkmark^\checkmark\\checkmarkc\checkmarki\checkmarkr\checkmarkc\checkmark$\checkmarkn\checkmark$\checkmark^\checkmark\\checkmarkc\checkmarki\checkmarkr\checkmarkc\checkmark$\checkmarkg\checkmark$\checkmark^\checkmark\\checkmarkc\checkmarki\checkmarkr\checkmarkc\checkmark$\checkmark \checkmark$\checkmark^\checkmark\\checkmarkc\checkmarki\checkmarkr\checkmarkc\checkmark$\checkmark\\checkmark$\checkmark^\checkmark\\checkmarkc\checkmarki\checkmarkr\checkmarkc\checkmark$\checkmarkv\checkmark$\checkmark^\checkmark\\checkmarkc\checkmarki\checkmarkr\checkmarkc\checkmark$\checkmarks\checkmark$\checkmark^\checkmark\\checkmarkc\checkmarki\checkmarkr\checkmarkc\checkmark$\checkmarkp\checkmark$\checkmark^\checkmark\\checkmarkc\checkmarki\checkmarkr\checkmarkc\checkmark$\checkmarka\checkmark$\checkmark^\checkmark\\checkmarkc\checkmarki\checkmarkr\checkmarkc\checkmark$\checkmarkc\checkmark$\checkmark^\checkmark\\checkmarkc\checkmarki\checkmarkr\checkmarkc\checkmark$\checkmarke\checkmark$\checkmark^\checkmark\\checkmarkc\checkmarki\checkmarkr\checkmarkc\checkmark$\checkmark{\checkmark$\checkmark^\checkmark\\checkmarkc\checkmarki\checkmarkr\checkmarkc\checkmark$\checkmark3\checkmark$\checkmark^\checkmark\\checkmarkc\checkmarki\checkmarkr\checkmarkc\checkmark$\checkmarkc\checkmark$\checkmark^\checkmark\\checkmarkc\checkmarki\checkmarkr\checkmarkc\checkmark$\checkmarkm\checkmark$\checkmark^\checkmark\\checkmarkc\checkmarki\checkmarkr\checkmarkc\checkmark$\checkmark}\checkmark$\checkmark^\checkmark\\checkmarkc\checkmarki\checkmarkr\checkmarkc\checkmark$\checkmark \checkmark$\checkmark^\checkmark\\checkmarkc\checkmarki\checkmarkr\checkmarkc\checkmark$\checkmark\\checkmark$\checkmark^\checkmark\\checkmarkc\checkmarki\checkmarkr\checkmarkc\checkmark$\checkmarkt\checkmark$\checkmark^\checkmark\\checkmarkc\checkmarki\checkmarkr\checkmarkc\checkmark$\checkmarke\checkmark$\checkmark^\checkmark\\checkmarkc\checkmarki\checkmarkr\checkmarkc\checkmark$\checkmarkx\checkmark$\checkmark^\checkmark\\checkmarkc\checkmarki\checkmarkr\checkmarkc\checkmark$\checkmarkt\checkmark$\checkmark^\checkmark\\checkmarkc\checkmarki\checkmarkr\checkmarkc\checkmark$\checkmarki\checkmark$\checkmark^\checkmark\\checkmarkc\checkmarki\checkmarkr\checkmarkc\checkmark$\checkmarkt\checkmark$\checkmark^\checkmark\\checkmarkc\checkmarki\checkmarkr\checkmarkc\checkmark$\checkmark{\checkmark$\checkmark^\checkmark\\checkmarkc\checkmarki\checkmarkr\checkmarkc\checkmark$\checkmark[\checkmark$\checkmark^\checkmark\\checkmarkc\checkmarki\checkmarkr\checkmarkc\checkmark$\checkmarkS\checkmark$\checkmark^\checkmark\\checkmarkc\checkmarki\checkmarkr\checkmarkc\checkmark$\checkmarkc\checkmark$\checkmark^\checkmark\\checkmarkc\checkmarki\checkmarkr\checkmarkc\checkmark$\checkmarkh\checkmark$\checkmark^\checkmark\\checkmarkc\checkmarki\checkmarkr\checkmarkc\checkmark$\checkmarké\checkmark$\checkmark^\checkmark\\checkmarkc\checkmarki\checkmarkr\checkmarkc\checkmark$\checkmarkm\checkmark$\checkmark^\checkmark\\checkmarkc\checkmarki\checkmarkr\checkmarkc\checkmark$\checkmarka\checkmark$\checkmark^\checkmark\\checkmarkc\checkmarki\checkmarkr\checkmarkc\checkmark$\checkmark \checkmark$\checkmark^\checkmark\\checkmarkc\checkmarki\checkmarkr\checkmarkc\checkmark$\checkmarkc\checkmark$\checkmark^\checkmark\\checkmarkc\checkmarki\checkmarkr\checkmarkc\checkmark$\checkmarki\checkmark$\checkmark^\checkmark\\checkmarkc\checkmarki\checkmarkr\checkmarkc\checkmark$\checkmarkn\checkmark$\checkmark^\checkmark\\checkmarkc\checkmarki\checkmarkr\checkmarkc\checkmark$\checkmarké\checkmark$\checkmark^\checkmark\\checkmarkc\checkmarki\checkmarkr\checkmarkc\checkmark$\checkmarkm\checkmark$\checkmark^\checkmark\\checkmarkc\checkmarki\checkmarkr\checkmarkc\checkmark$\checkmarka\checkmark$\checkmark^\checkmark\\checkmarkc\checkmarki\checkmarkr\checkmarkc\checkmark$\checkmarkt\checkmark$\checkmark^\checkmark\\checkmarkc\checkmarki\checkmarkr\checkmarkc\checkmark$\checkmarki\checkmark$\checkmark^\checkmark\\checkmarkc\checkmarki\checkmarkr\checkmarkc\checkmark$\checkmarkq\checkmark$\checkmark^\checkmark\\checkmarkc\checkmarki\checkmarkr\checkmarkc\checkmark$\checkmarku\checkmark$\checkmark^\checkmark\\checkmarkc\checkmarki\checkmarkr\checkmarkc\checkmark$\checkmarke\checkmark$\checkmark^\checkmark\\checkmarkc\checkmarki\checkmarkr\checkmarkc\checkmark$\checkmark \checkmark$\checkmark^\checkmark\\checkmarkc\checkmarki\checkmarkr\checkmarkc\checkmark$\checkmarkd\checkmark$\checkmark^\checkmark\\checkmarkc\checkmarki\checkmarkr\checkmarkc\checkmark$\checkmarke\checkmark$\checkmark^\checkmark\\checkmarkc\checkmarki\checkmarkr\checkmarkc\checkmark$\checkmark \checkmark$\checkmark^\checkmark\\checkmarkc\checkmarki\checkmarkr\checkmarkc\checkmark$\checkmarkl\checkmark$\checkmark^\checkmark\\checkmarkc\checkmarki\checkmarkr\checkmarkc\checkmark$\checkmarka\checkmark$\checkmark^\checkmark\\checkmarkc\checkmarki\checkmarkr\checkmarkc\checkmark$\checkmark \checkmark$\checkmark^\checkmark\\checkmarkc\checkmarki\checkmarkr\checkmarkc\checkmark$\checkmarkC\checkmark$\checkmark^\checkmark\\checkmarkc\checkmarki\checkmarkr\checkmarkc\checkmark$\checkmarkN\checkmark$\checkmark^\checkmark\\checkmarkc\checkmarki\checkmarkr\checkmarkc\checkmark$\checkmarkC\checkmark$\checkmark^\checkmark\\checkmarkc\checkmarki\checkmarkr\checkmarkc\checkmark$\checkmark \checkmark$\checkmark^\checkmark\\checkmarkc\checkmarki\checkmarkr\checkmarkc\checkmark$\checkmark5\checkmark$\checkmark^\checkmark\\checkmarkc\checkmarki\checkmarkr\checkmarkc\checkmark$\checkmark \checkmark$\checkmark^\checkmark\\checkmarkc\checkmarki\checkmarkr\checkmarkc\checkmark$\checkmarka\checkmark$\checkmark^\checkmark\\checkmarkc\checkmarki\checkmarkr\checkmarkc\checkmark$\checkmarkx\checkmark$\checkmark^\checkmark\\checkmarkc\checkmarki\checkmarkr\checkmarkc\checkmark$\checkmarke\checkmark$\checkmark^\checkmark\\checkmarkc\checkmarki\checkmarkr\checkmarkc\checkmark$\checkmarks\checkmark$\checkmark^\checkmark\\checkmarkc\checkmarki\checkmarkr\checkmarkc\checkmark$\checkmark \checkmark$\checkmark^\checkmark\\checkmarkc\checkmarki\checkmarkr\checkmarkc\checkmark$\checkmark—\checkmark$\checkmark^\checkmark\\checkmarkc\checkmarki\checkmarkr\checkmarkc\checkmark$\checkmark \checkmark$\checkmark^\checkmark\\checkmarkc\checkmarki\checkmarkr\checkmarkc\checkmark$\checkmarkF\checkmark$\checkmark^\checkmark\\checkmarkc\checkmarki\checkmarkr\checkmarkc\checkmark$\checkmarki\checkmark$\checkmark^\checkmark\\checkmarkc\checkmarki\checkmarkr\checkmarkc\checkmark$\checkmarkg\checkmark$\checkmark^\checkmark\\checkmarkc\checkmarki\checkmarkr\checkmarkc\checkmark$\checkmarku\checkmark$\checkmark^\checkmark\\checkmarkc\checkmarki\checkmarkr\checkmarkc\checkmark$\checkmarkr\checkmark$\checkmark^\checkmark\\checkmarkc\checkmarki\checkmarkr\checkmarkc\checkmark$\checkmarke\checkmark$\checkmark^\checkmark\\checkmarkc\checkmarki\checkmarkr\checkmarkc\checkmark$\checkmark \checkmark$\checkmark^\checkmark\\checkmarkc\checkmarki\checkmarkr\checkmarkc\checkmark$\checkmarki\checkmark$\checkmark^\checkmark\\checkmarkc\checkmarki\checkmarkr\checkmarkc\checkmark$\checkmarks\checkmark$\checkmark^\checkmark\\checkmarkc\checkmarki\checkmarkr\checkmarkc\checkmark$\checkmarks\checkmark$\checkmark^\checkmark\\checkmarkc\checkmarki\checkmarkr\checkmarkc\checkmark$\checkmarku\checkmark$\checkmark^\checkmark\\checkmarkc\checkmarki\checkmarkr\checkmarkc\checkmark$\checkmarke\checkmark$\checkmark^\checkmark\\checkmarkc\checkmarki\checkmarkr\checkmarkc\checkmark$\checkmark \checkmark$\checkmark^\checkmark\\checkmarkc\checkmarki\checkmarkr\checkmarkc\checkmark$\checkmarkd\checkmark$\checkmark^\checkmark\\checkmarkc\checkmarki\checkmarkr\checkmarkc\checkmark$\checkmarke\checkmark$\checkmark^\checkmark\\checkmarkc\checkmarki\checkmarkr\checkmarkc\checkmark$\checkmark \checkmark$\checkmark^\checkmark\\checkmarkc\checkmarki\checkmarkr\checkmarkc\checkmark$\checkmarkl\checkmark$\checkmark^\checkmark\\checkmarkc\checkmarki\checkmarkr\checkmarkc\checkmark$\checkmarka\checkmark$\checkmark^\checkmark\\checkmarkc\checkmarki\checkmarkr\checkmarkc\checkmark$\checkmark \checkmark$\checkmark^\checkmark\\checkmarkc\checkmarki\checkmarkr\checkmarkc\checkmark$\checkmarkC\checkmark$\checkmark^\checkmark\\checkmarkc\checkmarki\checkmarkr\checkmarkc\checkmark$\checkmarkA\checkmark$\checkmark^\checkmark\\checkmarkc\checkmarki\checkmarkr\checkmarkc\checkmark$\checkmarkO\checkmark$\checkmark^\checkmark\\checkmarkc\checkmarki\checkmarkr\checkmarkc\checkmark$\checkmark]\checkmark$\checkmark^\checkmark\\checkmarkc\checkmarki\checkmarkr\checkmarkc\checkmark$\checkmark}\checkmark$\checkmark^\checkmark\\checkmarkc\checkmarki\checkmarkr\checkmarkc\checkmark$\checkmark \checkmark$\checkmark^\checkmark\\checkmarkc\checkmarki\checkmarkr\checkmarkc\checkmark$\checkmark\\checkmark$\checkmark^\checkmark\\checkmarkc\checkmarki\checkmarkr\checkmarkc\checkmark$\checkmarkv\checkmark$\checkmark^\checkmark\\checkmarkc\checkmarki\checkmarkr\checkmarkc\checkmark$\checkmarks\checkmark$\checkmark^\checkmark\\checkmarkc\checkmarki\checkmarkr\checkmarkc\checkmark$\checkmarkp\checkmark$\checkmark^\checkmark\\checkmarkc\checkmarki\checkmarkr\checkmarkc\checkmark$\checkmarka\checkmark$\checkmark^\checkmark\\checkmarkc\checkmarki\checkmarkr\checkmarkc\checkmark$\checkmarkc\checkmark$\checkmark^\checkmark\\checkmarkc\checkmarki\checkmarkr\checkmarkc\checkmark$\checkmarke\checkmark$\checkmark^\checkmark\\checkmarkc\checkmarki\checkmarkr\checkmarkc\checkmark$\checkmark{\checkmark$\checkmark^\checkmark\\checkmarkc\checkmarki\checkmarkr\checkmarkc\checkmark$\checkmark3\checkmark$\checkmark^\checkmark\\checkmarkc\checkmarki\checkmarkr\checkmarkc\checkmark$\checkmarkc\checkmark$\checkmark^\checkmark\\checkmarkc\checkmarki\checkmarkr\checkmarkc\checkmark$\checkmarkm\checkmark$\checkmark^\checkmark\\checkmarkc\checkmarki\checkmarkr\checkmarkc\checkmark$\checkmark}\checkmark$\checkmark^\checkmark\\checkmarkc\checkmarki\checkmarkr\checkmarkc\checkmark$\checkmark}\checkmark$\checkmark^\checkmark\\checkmarkc\checkmarki\checkmarkr\checkmarkc\checkmark$\checkmark}\checkmark$\checkmark^\checkmark\\checkmarkc\checkmarki\checkmarkr\checkmarkc\checkmark$\checkmark
\checkmark$\checkmark^\checkmark\\checkmarkc\checkmarki\checkmarkr\checkmarkc\checkmark$\checkmark \checkmark$\checkmark^\checkmark\\checkmarkc\checkmarki\checkmarkr\checkmarkc\checkmark$\checkmark \checkmark$\checkmark^\checkmark\\checkmarkc\checkmarki\checkmarkr\checkmarkc\checkmark$\checkmark \checkmark$\checkmark^\checkmark\\checkmarkc\checkmarki\checkmarkr\checkmarkc\checkmark$\checkmark \checkmark$\checkmark^\checkmark\\checkmarkc\checkmarki\checkmarkr\checkmarkc\checkmark$\checkmark\\checkmark$\checkmark^\checkmark\\checkmarkc\checkmarki\checkmarkr\checkmarkc\checkmark$\checkmarkc\checkmark$\checkmark^\checkmark\\checkmarkc\checkmarki\checkmarkr\checkmarkc\checkmark$\checkmarka\checkmark$\checkmark^\checkmark\\checkmarkc\checkmarki\checkmarkr\checkmarkc\checkmark$\checkmarkp\checkmark$\checkmark^\checkmark\\checkmarkc\checkmarki\checkmarkr\checkmarkc\checkmark$\checkmarkt\checkmark$\checkmark^\checkmark\\checkmarkc\checkmarki\checkmarkr\checkmarkc\checkmark$\checkmarki\checkmark$\checkmark^\checkmark\\checkmarkc\checkmarki\checkmarkr\checkmarkc\checkmark$\checkmarko\checkmark$\checkmark^\checkmark\\checkmarkc\checkmarki\checkmarkr\checkmarkc\checkmark$\checkmarkn\checkmark$\checkmark^\checkmark\\checkmarkc\checkmarki\checkmarkr\checkmarkc\checkmark$\checkmark{\checkmark$\checkmark^\checkmark\\checkmarkc\checkmarki\checkmarkr\checkmarkc\checkmark$\checkmarkS\checkmark$\checkmark^\checkmark\\checkmarkc\checkmarki\checkmarkr\checkmarkc\checkmark$\checkmarkc\checkmark$\checkmark^\checkmark\\checkmarkc\checkmarki\checkmarkr\checkmarkc\checkmark$\checkmarkh\checkmark$\checkmark^\checkmark\\checkmarkc\checkmarki\checkmarkr\checkmarkc\checkmark$\checkmarké\checkmark$\checkmark^\checkmark\\checkmarkc\checkmarki\checkmarkr\checkmarkc\checkmark$\checkmarkm\checkmark$\checkmark^\checkmark\\checkmarkc\checkmarki\checkmarkr\checkmarkc\checkmark$\checkmarka\checkmark$\checkmark^\checkmark\\checkmarkc\checkmarki\checkmarkr\checkmarkc\checkmark$\checkmark \checkmark$\checkmark^\checkmark\\checkmarkc\checkmarki\checkmarkr\checkmarkc\checkmark$\checkmarkc\checkmark$\checkmark^\checkmark\\checkmarkc\checkmarki\checkmarkr\checkmarkc\checkmark$\checkmarki\checkmark$\checkmark^\checkmark\\checkmarkc\checkmarki\checkmarkr\checkmarkc\checkmark$\checkmarkn\checkmark$\checkmark^\checkmark\\checkmarkc\checkmarki\checkmarkr\checkmarkc\checkmark$\checkmarké\checkmark$\checkmark^\checkmark\\checkmarkc\checkmarki\checkmarkr\checkmarkc\checkmark$\checkmarkm\checkmark$\checkmark^\checkmark\\checkmarkc\checkmarki\checkmarkr\checkmarkc\checkmark$\checkmarka\checkmark$\checkmark^\checkmark\\checkmarkc\checkmarki\checkmarkr\checkmarkc\checkmark$\checkmarkt\checkmark$\checkmark^\checkmark\\checkmarkc\checkmarki\checkmarkr\checkmarkc\checkmark$\checkmarki\checkmark$\checkmark^\checkmark\\checkmarkc\checkmarki\checkmarkr\checkmarkc\checkmark$\checkmarkq\checkmark$\checkmark^\checkmark\\checkmarkc\checkmarki\checkmarkr\checkmarkc\checkmark$\checkmarku\checkmark$\checkmark^\checkmark\\checkmarkc\checkmarki\checkmarkr\checkmarkc\checkmark$\checkmarke\checkmark$\checkmark^\checkmark\\checkmarkc\checkmarki\checkmarkr\checkmarkc\checkmark$\checkmark \checkmark$\checkmark^\checkmark\\checkmarkc\checkmarki\checkmarkr\checkmarkc\checkmark$\checkmarkd\checkmark$\checkmark^\checkmark\\checkmarkc\checkmarki\checkmarkr\checkmarkc\checkmark$\checkmarke\checkmark$\checkmark^\checkmark\\checkmarkc\checkmarki\checkmarkr\checkmarkc\checkmark$\checkmark \checkmark$\checkmark^\checkmark\\checkmarkc\checkmarki\checkmarkr\checkmarkc\checkmark$\checkmarkl\checkmark$\checkmark^\checkmark\\checkmarkc\checkmarki\checkmarkr\checkmarkc\checkmark$\checkmarka\checkmark$\checkmark^\checkmark\\checkmarkc\checkmarki\checkmarkr\checkmarkc\checkmark$\checkmark \checkmark$\checkmark^\checkmark\\checkmarkc\checkmarki\checkmarkr\checkmarkc\checkmark$\checkmarkm\checkmark$\checkmark^\checkmark\\checkmarkc\checkmarki\checkmarkr\checkmarkc\checkmark$\checkmarki\checkmark$\checkmark^\checkmark\\checkmarkc\checkmarki\checkmarkr\checkmarkc\checkmark$\checkmarkn\checkmark$\checkmark^\checkmark\\checkmarkc\checkmarki\checkmarkr\checkmarkc\checkmark$\checkmarki\checkmark$\checkmark^\checkmark\\checkmarkc\checkmarki\checkmarkr\checkmarkc\checkmark$\checkmark-\checkmark$\checkmark^\checkmark\\checkmarkc\checkmarki\checkmarkr\checkmarkc\checkmark$\checkmarkf\checkmark$\checkmark^\checkmark\\checkmarkc\checkmarki\checkmarkr\checkmarkc\checkmark$\checkmarkr\checkmark$\checkmark^\checkmark\\checkmarkc\checkmarki\checkmarkr\checkmarkc\checkmark$\checkmarka\checkmark$\checkmark^\checkmark\\checkmarkc\checkmarki\checkmarkr\checkmarkc\checkmark$\checkmarki\checkmark$\checkmark^\checkmark\\checkmarkc\checkmarki\checkmarkr\checkmarkc\checkmark$\checkmarks\checkmark$\checkmark^\checkmark\\checkmarkc\checkmarki\checkmarkr\checkmarkc\checkmark$\checkmarke\checkmark$\checkmark^\checkmark\\checkmarkc\checkmarki\checkmarkr\checkmarkc\checkmark$\checkmarku\checkmark$\checkmark^\checkmark\\checkmarkc\checkmarki\checkmarkr\checkmarkc\checkmark$\checkmarks\checkmark$\checkmark^\checkmark\\checkmarkc\checkmarki\checkmarkr\checkmarkc\checkmark$\checkmarke\checkmark$\checkmark^\checkmark\\checkmarkc\checkmarki\checkmarkr\checkmarkc\checkmark$\checkmark \checkmark$\checkmark^\checkmark\\checkmarkc\checkmarki\checkmarkr\checkmarkc\checkmark$\checkmarkC\checkmark$\checkmark^\checkmark\\checkmarkc\checkmarki\checkmarkr\checkmarkc\checkmark$\checkmarkN\checkmark$\checkmark^\checkmark\\checkmarkc\checkmarki\checkmarkr\checkmarkc\checkmark$\checkmarkC\checkmark$\checkmark^\checkmark\\checkmarkc\checkmarki\checkmarkr\checkmarkc\checkmark$\checkmark \checkmark$\checkmark^\checkmark\\checkmarkc\checkmarki\checkmarkr\checkmarkc\checkmark$\checkmark5\checkmark$\checkmark^\checkmark\\checkmarkc\checkmarki\checkmarkr\checkmarkc\checkmark$\checkmark \checkmark$\checkmark^\checkmark\\checkmarkc\checkmarki\checkmarkr\checkmarkc\checkmark$\checkmarka\checkmark$\checkmark^\checkmark\\checkmarkc\checkmarki\checkmarkr\checkmarkc\checkmark$\checkmarkx\checkmark$\checkmark^\checkmark\\checkmarkc\checkmarki\checkmarkr\checkmarkc\checkmark$\checkmarke\checkmark$\checkmark^\checkmark\\checkmarkc\checkmarki\checkmarkr\checkmarkc\checkmark$\checkmarks\checkmark$\checkmark^\checkmark\\checkmarkc\checkmarki\checkmarkr\checkmarkc\checkmark$\checkmark \checkmark$\checkmark^\checkmark\\checkmarkc\checkmarki\checkmarkr\checkmarkc\checkmark$\checkmarkT\checkmark$\checkmark^\checkmark\\checkmarkc\checkmarki\checkmarkr\checkmarkc\checkmark$\checkmarkr\checkmark$\checkmark^\checkmark\\checkmarkc\checkmarki\checkmarkr\checkmarkc\checkmark$\checkmarku\checkmark$\checkmark^\checkmark\\checkmarkc\checkmarki\checkmarkr\checkmarkc\checkmark$\checkmarkn\checkmark$\checkmark^\checkmark\\checkmarkc\checkmarki\checkmarkr\checkmarkc\checkmark$\checkmarkn\checkmark$\checkmark^\checkmark\\checkmarkc\checkmarki\checkmarkr\checkmarkc\checkmark$\checkmarki\checkmark$\checkmark^\checkmark\\checkmarkc\checkmarki\checkmarkr\checkmarkc\checkmark$\checkmarko\checkmark$\checkmark^\checkmark\\checkmarkc\checkmarki\checkmarkr\checkmarkc\checkmark$\checkmarkn\checkmark$\checkmark^\checkmark\\checkmarkc\checkmarki\checkmarkr\checkmarkc\checkmark$\checkmark}\checkmark$\checkmark^\checkmark\\checkmarkc\checkmarki\checkmarkr\checkmarkc\checkmark$\checkmark
\checkmark$\checkmark^\checkmark\\checkmarkc\checkmarki\checkmarkr\checkmarkc\checkmark$\checkmark \checkmark$\checkmark^\checkmark\\checkmarkc\checkmarki\checkmarkr\checkmarkc\checkmark$\checkmark \checkmark$\checkmark^\checkmark\\checkmarkc\checkmarki\checkmarkr\checkmarkc\checkmark$\checkmark \checkmark$\checkmark^\checkmark\\checkmarkc\checkmarki\checkmarkr\checkmarkc\checkmark$\checkmark \checkmark$\checkmark^\checkmark\\checkmarkc\checkmarki\checkmarkr\checkmarkc\checkmark$\checkmark\\checkmark$\checkmark^\checkmark\\checkmarkc\checkmarki\checkmarkr\checkmarkc\checkmark$\checkmarkl\checkmark$\checkmark^\checkmark\\checkmarkc\checkmarki\checkmarkr\checkmarkc\checkmark$\checkmarka\checkmark$\checkmark^\checkmark\\checkmarkc\checkmarki\checkmarkr\checkmarkc\checkmark$\checkmarkb\checkmark$\checkmark^\checkmark\\checkmarkc\checkmarki\checkmarkr\checkmarkc\checkmark$\checkmarke\checkmark$\checkmark^\checkmark\\checkmarkc\checkmarki\checkmarkr\checkmarkc\checkmark$\checkmarkl\checkmark$\checkmark^\checkmark\\checkmarkc\checkmarki\checkmarkr\checkmarkc\checkmark$\checkmark{\checkmark$\checkmark^\checkmark\\checkmarkc\checkmarki\checkmarkr\checkmarkc\checkmark$\checkmarkf\checkmark$\checkmark^\checkmark\\checkmarkc\checkmarki\checkmarkr\checkmarkc\checkmark$\checkmarki\checkmark$\checkmark^\checkmark\\checkmarkc\checkmarki\checkmarkr\checkmarkc\checkmark$\checkmarkg\checkmark$\checkmark^\checkmark\\checkmarkc\checkmarki\checkmarkr\checkmarkc\checkmark$\checkmark:\checkmark$\checkmark^\checkmark\\checkmarkc\checkmarki\checkmarkr\checkmarkc\checkmark$\checkmarks\checkmark$\checkmark^\checkmark\\checkmarkc\checkmarki\checkmarkr\checkmarkc\checkmark$\checkmarkc\checkmark$\checkmark^\checkmark\\checkmarkc\checkmarki\checkmarkr\checkmarkc\checkmark$\checkmarkh\checkmark$\checkmark^\checkmark\\checkmarkc\checkmarki\checkmarkr\checkmarkc\checkmark$\checkmarke\checkmark$\checkmark^\checkmark\\checkmarkc\checkmarki\checkmarkr\checkmarkc\checkmark$\checkmarkm\checkmark$\checkmark^\checkmark\\checkmarkc\checkmarki\checkmarkr\checkmarkc\checkmark$\checkmarka\checkmark$\checkmark^\checkmark\\checkmarkc\checkmarki\checkmarkr\checkmarkc\checkmark$\checkmark_\checkmark$\checkmark^\checkmark\\checkmarkc\checkmarki\checkmarkr\checkmarkc\checkmark$\checkmarkc\checkmark$\checkmark^\checkmark\\checkmarkc\checkmarki\checkmarkr\checkmarkc\checkmark$\checkmarki\checkmark$\checkmark^\checkmark\\checkmarkc\checkmarki\checkmarkr\checkmarkc\checkmark$\checkmarkn\checkmark$\checkmark^\checkmark\\checkmarkc\checkmarki\checkmarkr\checkmarkc\checkmark$\checkmarke\checkmark$\checkmark^\checkmark\\checkmarkc\checkmarki\checkmarkr\checkmarkc\checkmark$\checkmarkm\checkmark$\checkmark^\checkmark\\checkmarkc\checkmarki\checkmarkr\checkmarkc\checkmark$\checkmarka\checkmark$\checkmark^\checkmark\\checkmarkc\checkmarki\checkmarkr\checkmarkc\checkmark$\checkmarkt\checkmark$\checkmark^\checkmark\\checkmarkc\checkmarki\checkmarkr\checkmarkc\checkmark$\checkmarki\checkmark$\checkmark^\checkmark\\checkmarkc\checkmarki\checkmarkr\checkmarkc\checkmark$\checkmarkq\checkmark$\checkmark^\checkmark\\checkmarkc\checkmarki\checkmarkr\checkmarkc\checkmark$\checkmarku\checkmark$\checkmark^\checkmark\\checkmarkc\checkmarki\checkmarkr\checkmarkc\checkmark$\checkmarke\checkmark$\checkmark^\checkmark\\checkmarkc\checkmarki\checkmarkr\checkmarkc\checkmark$\checkmark}\checkmark$\checkmark^\checkmark\\checkmarkc\checkmarki\checkmarkr\checkmarkc\checkmark$\checkmark
\checkmark$\checkmark^\checkmark\\checkmarkc\checkmarki\checkmarkr\checkmarkc\checkmark$\checkmark\\checkmark$\checkmark^\checkmark\\checkmarkc\checkmarki\checkmarkr\checkmarkc\checkmark$\checkmarke\checkmark$\checkmark^\checkmark\\checkmarkc\checkmarki\checkmarkr\checkmarkc\checkmark$\checkmarkn\checkmark$\checkmark^\checkmark\\checkmarkc\checkmarki\checkmarkr\checkmarkc\checkmark$\checkmarkd\checkmark$\checkmark^\checkmark\\checkmarkc\checkmarki\checkmarkr\checkmarkc\checkmark$\checkmark{\checkmark$\checkmark^\checkmark\\checkmarkc\checkmarki\checkmarkr\checkmarkc\checkmark$\checkmarkf\checkmark$\checkmark^\checkmark\\checkmarkc\checkmarki\checkmarkr\checkmarkc\checkmark$\checkmarki\checkmark$\checkmark^\checkmark\\checkmarkc\checkmarki\checkmarkr\checkmarkc\checkmark$\checkmarkg\checkmark$\checkmark^\checkmark\\checkmarkc\checkmarki\checkmarkr\checkmarkc\checkmark$\checkmarku\checkmark$\checkmark^\checkmark\\checkmarkc\checkmarki\checkmarkr\checkmarkc\checkmark$\checkmarkr\checkmark$\checkmark^\checkmark\\checkmarkc\checkmarki\checkmarkr\checkmarkc\checkmark$\checkmarke\checkmark$\checkmark^\checkmark\\checkmarkc\checkmarki\checkmarkr\checkmarkc\checkmark$\checkmark}\checkmark$\checkmark^\checkmark\\checkmarkc\checkmarki\checkmarkr\checkmarkc\checkmark$\checkmark
\checkmark$\checkmark^\checkmark\\checkmarkc\checkmarki\checkmarkr\checkmarkc\checkmark$\checkmark
\checkmark$\checkmark^\checkmark\\checkmarkc\checkmarki\checkmarkr\checkmarkc\checkmark$\checkmark%\checkmark$\checkmark^\checkmark\\checkmarkc\checkmarki\checkmarkr\checkmarkc\checkmark$\checkmark \checkmark$\checkmark^\checkmark\\checkmarkc\checkmarki\checkmarkr\checkmarkc\checkmark$\checkmark=\checkmark$\checkmark^\checkmark\\checkmarkc\checkmarki\checkmarkr\checkmarkc\checkmark$\checkmark=\checkmark$\checkmark^\checkmark\\checkmarkc\checkmarki\checkmarkr\checkmarkc\checkmark$\checkmark=\checkmark$\checkmark^\checkmark\\checkmarkc\checkmarki\checkmarkr\checkmarkc\checkmark$\checkmark=\checkmark$\checkmark^\checkmark\\checkmarkc\checkmarki\checkmarkr\checkmarkc\checkmark$\checkmark=\checkmark$\checkmark^\checkmark\\checkmarkc\checkmarki\checkmarkr\checkmarkc\checkmark$\checkmark=\checkmark$\checkmark^\checkmark\\checkmarkc\checkmarki\checkmarkr\checkmarkc\checkmark$\checkmark=\checkmark$\checkmark^\checkmark\\checkmarkc\checkmarki\checkmarkr\checkmarkc\checkmark$\checkmark=\checkmark$\checkmark^\checkmark\\checkmarkc\checkmarki\checkmarkr\checkmarkc\checkmark$\checkmark=\checkmark$\checkmark^\checkmark\\checkmarkc\checkmarki\checkmarkr\checkmarkc\checkmark$\checkmark=\checkmark$\checkmark^\checkmark\\checkmarkc\checkmarki\checkmarkr\checkmarkc\checkmark$\checkmark=\checkmark$\checkmark^\checkmark\\checkmarkc\checkmarki\checkmarkr\checkmarkc\checkmark$\checkmark=\checkmark$\checkmark^\checkmark\\checkmarkc\checkmarki\checkmarkr\checkmarkc\checkmark$\checkmark=\checkmark$\checkmark^\checkmark\\checkmarkc\checkmarki\checkmarkr\checkmarkc\checkmark$\checkmark=\checkmark$\checkmark^\checkmark\\checkmarkc\checkmarki\checkmarkr\checkmarkc\checkmark$\checkmark=\checkmark$\checkmark^\checkmark\\checkmarkc\checkmarki\checkmarkr\checkmarkc\checkmark$\checkmark=\checkmark$\checkmark^\checkmark\\checkmarkc\checkmarki\checkmarkr\checkmarkc\checkmark$\checkmark=\checkmark$\checkmark^\checkmark\\checkmarkc\checkmarki\checkmarkr\checkmarkc\checkmark$\checkmark=\checkmark$\checkmark^\checkmark\\checkmarkc\checkmarki\checkmarkr\checkmarkc\checkmark$\checkmark=\checkmark$\checkmark^\checkmark\\checkmarkc\checkmarki\checkmarkr\checkmarkc\checkmark$\checkmark=\checkmark$\checkmark^\checkmark\\checkmarkc\checkmarki\checkmarkr\checkmarkc\checkmark$\checkmark=\checkmark$\checkmark^\checkmark\\checkmarkc\checkmarki\checkmarkr\checkmarkc\checkmark$\checkmark=\checkmark$\checkmark^\checkmark\\checkmarkc\checkmarki\checkmarkr\checkmarkc\checkmark$\checkmark=\checkmark$\checkmark^\checkmark\\checkmarkc\checkmarki\checkmarkr\checkmarkc\checkmark$\checkmark=\checkmark$\checkmark^\checkmark\\checkmarkc\checkmarki\checkmarkr\checkmarkc\checkmark$\checkmark=\checkmark$\checkmark^\checkmark\\checkmarkc\checkmarki\checkmarkr\checkmarkc\checkmark$\checkmark=\checkmark$\checkmark^\checkmark\\checkmarkc\checkmarki\checkmarkr\checkmarkc\checkmark$\checkmark=\checkmark$\checkmark^\checkmark\\checkmarkc\checkmarki\checkmarkr\checkmarkc\checkmark$\checkmark=\checkmark$\checkmark^\checkmark\\checkmarkc\checkmarki\checkmarkr\checkmarkc\checkmark$\checkmark=\checkmark$\checkmark^\checkmark\\checkmarkc\checkmarki\checkmarkr\checkmarkc\checkmark$\checkmark=\checkmark$\checkmark^\checkmark\\checkmarkc\checkmarki\checkmarkr\checkmarkc\checkmark$\checkmark=\checkmark$\checkmark^\checkmark\\checkmarkc\checkmarki\checkmarkr\checkmarkc\checkmark$\checkmark=\checkmark$\checkmark^\checkmark\\checkmarkc\checkmarki\checkmarkr\checkmarkc\checkmark$\checkmark=\checkmark$\checkmark^\checkmark\\checkmarkc\checkmarki\checkmarkr\checkmarkc\checkmark$\checkmark=\checkmark$\checkmark^\checkmark\\checkmarkc\checkmarki\checkmarkr\checkmarkc\checkmark$\checkmark=\checkmark$\checkmark^\checkmark\\checkmarkc\checkmarki\checkmarkr\checkmarkc\checkmark$\checkmark=\checkmark$\checkmark^\checkmark\\checkmarkc\checkmarki\checkmarkr\checkmarkc\checkmark$\checkmark=\checkmark$\checkmark^\checkmark\\checkmarkc\checkmarki\checkmarkr\checkmarkc\checkmark$\checkmark=\checkmark$\checkmark^\checkmark\\checkmarkc\checkmarki\checkmarkr\checkmarkc\checkmark$\checkmark=\checkmark$\checkmark^\checkmark\\checkmarkc\checkmarki\checkmarkr\checkmarkc\checkmark$\checkmark=\checkmark$\checkmark^\checkmark\\checkmarkc\checkmarki\checkmarkr\checkmarkc\checkmark$\checkmark=\checkmark$\checkmark^\checkmark\\checkmarkc\checkmarki\checkmarkr\checkmarkc\checkmark$\checkmark=\checkmark$\checkmark^\checkmark\\checkmarkc\checkmarki\checkmarkr\checkmarkc\checkmark$\checkmark=\checkmark$\checkmark^\checkmark\\checkmarkc\checkmarki\checkmarkr\checkmarkc\checkmark$\checkmark=\checkmark$\checkmark^\checkmark\\checkmarkc\checkmarki\checkmarkr\checkmarkc\checkmark$\checkmark=\checkmark$\checkmark^\checkmark\\checkmarkc\checkmarki\checkmarkr\checkmarkc\checkmark$\checkmark=\checkmark$\checkmark^\checkmark\\checkmarkc\checkmarki\checkmarkr\checkmarkc\checkmark$\checkmark=\checkmark$\checkmark^\checkmark\\checkmarkc\checkmarki\checkmarkr\checkmarkc\checkmark$\checkmark=\checkmark$\checkmark^\checkmark\\checkmarkc\checkmarki\checkmarkr\checkmarkc\checkmark$\checkmark=\checkmark$\checkmark^\checkmark\\checkmarkc\checkmarki\checkmarkr\checkmarkc\checkmark$\checkmark=\checkmark$\checkmark^\checkmark\\checkmarkc\checkmarki\checkmarkr\checkmarkc\checkmark$\checkmark=\checkmark$\checkmark^\checkmark\\checkmarkc\checkmarki\checkmarkr\checkmarkc\checkmark$\checkmark=\checkmark$\checkmark^\checkmark\\checkmarkc\checkmarki\checkmarkr\checkmarkc\checkmark$\checkmark=\checkmark$\checkmark^\checkmark\\checkmarkc\checkmarki\checkmarkr\checkmarkc\checkmark$\checkmark=\checkmark$\checkmark^\checkmark\\checkmarkc\checkmarki\checkmarkr\checkmarkc\checkmark$\checkmark=\checkmark$\checkmark^\checkmark\\checkmarkc\checkmarki\checkmarkr\checkmarkc\checkmark$\checkmark=\checkmark$\checkmark^\checkmark\\checkmarkc\checkmarki\checkmarkr\checkmarkc\checkmark$\checkmark=\checkmark$\checkmark^\checkmark\\checkmarkc\checkmarki\checkmarkr\checkmarkc\checkmark$\checkmark=\checkmark$\checkmark^\checkmark\\checkmarkc\checkmarki\checkmarkr\checkmarkc\checkmark$\checkmark=\checkmark$\checkmark^\checkmark\\checkmarkc\checkmarki\checkmarkr\checkmarkc\checkmark$\checkmark=\checkmark$\checkmark^\checkmark\\checkmarkc\checkmarki\checkmarkr\checkmarkc\checkmark$\checkmark=\checkmark$\checkmark^\checkmark\\checkmarkc\checkmarki\checkmarkr\checkmarkc\checkmark$\checkmark=\checkmark$\checkmark^\checkmark\\checkmarkc\checkmarki\checkmarkr\checkmarkc\checkmark$\checkmark
\checkmark$\checkmark^\checkmark\\checkmarkc\checkmarki\checkmarkr\checkmarkc\checkmark$\checkmark\\checkmark$\checkmark^\checkmark\\checkmarkc\checkmarki\checkmarkr\checkmarkc\checkmark$\checkmarks\checkmark$\checkmark^\checkmark\\checkmarkc\checkmarki\checkmarkr\checkmarkc\checkmark$\checkmarke\checkmark$\checkmark^\checkmark\\checkmarkc\checkmarki\checkmarkr\checkmarkc\checkmark$\checkmarkc\checkmark$\checkmark^\checkmark\\checkmarkc\checkmarki\checkmarkr\checkmarkc\checkmark$\checkmarkt\checkmark$\checkmark^\checkmark\\checkmarkc\checkmarki\checkmarkr\checkmarkc\checkmark$\checkmarki\checkmark$\checkmark^\checkmark\\checkmarkc\checkmarki\checkmarkr\checkmarkc\checkmark$\checkmarko\checkmark$\checkmark^\checkmark\\checkmarkc\checkmarki\checkmarkr\checkmarkc\checkmark$\checkmarkn\checkmark$\checkmark^\checkmark\\checkmarkc\checkmarki\checkmarkr\checkmarkc\checkmark$\checkmark{\checkmark$\checkmark^\checkmark\\checkmarkc\checkmarki\checkmarkr\checkmarkc\checkmark$\checkmarkÉ\checkmark$\checkmark^\checkmark\\checkmarkc\checkmarki\checkmarkr\checkmarkc\checkmark$\checkmarkt\checkmark$\checkmark^\checkmark\\checkmarkc\checkmarki\checkmarkr\checkmarkc\checkmark$\checkmarku\checkmark$\checkmark^\checkmark\\checkmarkc\checkmarki\checkmarkr\checkmarkc\checkmark$\checkmarkd\checkmark$\checkmark^\checkmark\\checkmarkc\checkmarki\checkmarkr\checkmarkc\checkmark$\checkmarke\checkmark$\checkmark^\checkmark\\checkmarkc\checkmarki\checkmarkr\checkmarkc\checkmark$\checkmark \checkmark$\checkmark^\checkmark\\checkmarkc\checkmarki\checkmarkr\checkmarkc\checkmark$\checkmarkd\checkmark$\checkmark^\checkmark\\checkmarkc\checkmarki\checkmarkr\checkmarkc\checkmark$\checkmarke\checkmark$\checkmark^\checkmark\\checkmarkc\checkmarki\checkmarkr\checkmarkc\checkmark$\checkmarks\checkmark$\checkmark^\checkmark\\checkmarkc\checkmarki\checkmarkr\checkmarkc\checkmark$\checkmark \checkmark$\checkmark^\checkmark\\checkmarkc\checkmarki\checkmarkr\checkmarkc\checkmark$\checkmarkE\checkmark$\checkmark^\checkmark\\checkmarkc\checkmarki\checkmarkr\checkmarkc\checkmark$\checkmarkf\checkmark$\checkmark^\checkmark\\checkmarkc\checkmarki\checkmarkr\checkmarkc\checkmark$\checkmarkf\checkmark$\checkmark^\checkmark\\checkmarkc\checkmarki\checkmarkr\checkmarkc\checkmark$\checkmarko\checkmark$\checkmark^\checkmark\\checkmarkc\checkmarki\checkmarkr\checkmarkc\checkmark$\checkmarkr\checkmark$\checkmark^\checkmark\\checkmarkc\checkmarki\checkmarkr\checkmarkc\checkmark$\checkmarkt\checkmark$\checkmark^\checkmark\\checkmarkc\checkmarki\checkmarkr\checkmarkc\checkmark$\checkmarks\checkmark$\checkmark^\checkmark\\checkmarkc\checkmarki\checkmarkr\checkmarkc\checkmark$\checkmark \checkmark$\checkmark^\checkmark\\checkmarkc\checkmarki\checkmarkr\checkmarkc\checkmark$\checkmarkd\checkmark$\checkmark^\checkmark\\checkmarkc\checkmarki\checkmarkr\checkmarkc\checkmark$\checkmarke\checkmark$\checkmark^\checkmark\\checkmarkc\checkmarki\checkmarkr\checkmarkc\checkmark$\checkmark \checkmark$\checkmark^\checkmark\\checkmarkc\checkmarki\checkmarkr\checkmarkc\checkmark$\checkmarkC\checkmark$\checkmark^\checkmark\\checkmarkc\checkmarki\checkmarkr\checkmarkc\checkmark$\checkmarko\checkmark$\checkmark^\checkmark\\checkmarkc\checkmarki\checkmarkr\checkmarkc\checkmark$\checkmarku\checkmark$\checkmark^\checkmark\\checkmarkc\checkmarki\checkmarkr\checkmarkc\checkmark$\checkmarkp\checkmark$\checkmark^\checkmark\\checkmarkc\checkmarki\checkmarkr\checkmarkc\checkmark$\checkmarke\checkmark$\checkmark^\checkmark\\checkmarkc\checkmarki\checkmarkr\checkmarkc\checkmark$\checkmark}\checkmark$\checkmark^\checkmark\\checkmarkc\checkmarki\checkmarkr\checkmarkc\checkmark$\checkmark
\checkmark$\checkmark^\checkmark\\checkmarkc\checkmarki\checkmarkr\checkmarkc\checkmark$\checkmark
\checkmark$\checkmark^\checkmark\\checkmarkc\checkmarki\checkmarkr\checkmarkc\checkmark$\checkmark\\checkmark$\checkmark^\checkmark\\checkmarkc\checkmarki\checkmarkr\checkmarkc\checkmark$\checkmarks\checkmark$\checkmark^\checkmark\\checkmarkc\checkmarki\checkmarkr\checkmarkc\checkmark$\checkmarku\checkmark$\checkmark^\checkmark\\checkmarkc\checkmarki\checkmarkr\checkmarkc\checkmark$\checkmarkb\checkmark$\checkmark^\checkmark\\checkmarkc\checkmarki\checkmarkr\checkmarkc\checkmark$\checkmarks\checkmark$\checkmark^\checkmark\\checkmarkc\checkmarki\checkmarkr\checkmarkc\checkmark$\checkmarke\checkmark$\checkmark^\checkmark\\checkmarkc\checkmarki\checkmarkr\checkmarkc\checkmark$\checkmarkc\checkmark$\checkmark^\checkmark\\checkmarkc\checkmarki\checkmarkr\checkmarkc\checkmark$\checkmarkt\checkmark$\checkmark^\checkmark\\checkmarkc\checkmarki\checkmarkr\checkmarkc\checkmark$\checkmarki\checkmark$\checkmark^\checkmark\\checkmarkc\checkmarki\checkmarkr\checkmarkc\checkmark$\checkmarko\checkmark$\checkmark^\checkmark\\checkmarkc\checkmarki\checkmarkr\checkmarkc\checkmark$\checkmarkn\checkmark$\checkmark^\checkmark\\checkmarkc\checkmarki\checkmarkr\checkmarkc\checkmark$\checkmark{\checkmark$\checkmark^\checkmark\\checkmarkc\checkmarki\checkmarkr\checkmarkc\checkmark$\checkmarkI\checkmark$\checkmark^\checkmark\\checkmarkc\checkmarki\checkmarkr\checkmarkc\checkmark$\checkmarkn\checkmark$\checkmark^\checkmark\\checkmarkc\checkmarki\checkmarkr\checkmarkc\checkmark$\checkmarkt\checkmark$\checkmark^\checkmark\\checkmarkc\checkmarki\checkmarkr\checkmarkc\checkmark$\checkmarkr\checkmark$\checkmark^\checkmark\\checkmarkc\checkmarki\checkmarkr\checkmarkc\checkmark$\checkmarko\checkmark$\checkmark^\checkmark\\checkmarkc\checkmarki\checkmarkr\checkmarkc\checkmark$\checkmarkd\checkmark$\checkmark^\checkmark\\checkmarkc\checkmarki\checkmarkr\checkmarkc\checkmark$\checkmarku\checkmark$\checkmark^\checkmark\\checkmarkc\checkmarki\checkmarkr\checkmarkc\checkmark$\checkmarkc\checkmark$\checkmark^\checkmark\\checkmarkc\checkmarki\checkmarkr\checkmarkc\checkmark$\checkmarkt\checkmark$\checkmark^\checkmark\\checkmarkc\checkmarki\checkmarkr\checkmarkc\checkmark$\checkmarki\checkmark$\checkmark^\checkmark\\checkmarkc\checkmarki\checkmarkr\checkmarkc\checkmark$\checkmarko\checkmark$\checkmark^\checkmark\\checkmarkc\checkmarki\checkmarkr\checkmarkc\checkmark$\checkmarkn\checkmark$\checkmark^\checkmark\\checkmarkc\checkmarki\checkmarkr\checkmarkc\checkmark$\checkmark \checkmark$\checkmark^\checkmark\\checkmarkc\checkmarki\checkmarkr\checkmarkc\checkmark$\checkmarke\checkmark$\checkmark^\checkmark\\checkmarkc\checkmarki\checkmarkr\checkmarkc\checkmark$\checkmarkt\checkmark$\checkmark^\checkmark\\checkmarkc\checkmarki\checkmarkr\checkmarkc\checkmark$\checkmark \checkmark$\checkmark^\checkmark\\checkmarkc\checkmarki\checkmarkr\checkmarkc\checkmark$\checkmarkp\checkmark$\checkmark^\checkmark\\checkmarkc\checkmarki\checkmarkr\checkmarkc\checkmark$\checkmarka\checkmark$\checkmark^\checkmark\\checkmarkc\checkmarki\checkmarkr\checkmarkc\checkmark$\checkmarkr\checkmark$\checkmark^\checkmark\\checkmarkc\checkmarki\checkmarkr\checkmarkc\checkmark$\checkmarka\checkmark$\checkmark^\checkmark\\checkmarkc\checkmarki\checkmarkr\checkmarkc\checkmark$\checkmarkm\checkmark$\checkmark^\checkmark\\checkmarkc\checkmarki\checkmarkr\checkmarkc\checkmark$\checkmarkè\checkmark$\checkmark^\checkmark\\checkmarkc\checkmarki\checkmarkr\checkmarkc\checkmark$\checkmarkt\checkmark$\checkmark^\checkmark\\checkmarkc\checkmarki\checkmarkr\checkmarkc\checkmark$\checkmarkr\checkmark$\checkmark^\checkmark\\checkmarkc\checkmarki\checkmarkr\checkmarkc\checkmark$\checkmarke\checkmark$\checkmark^\checkmark\\checkmarkc\checkmarki\checkmarkr\checkmarkc\checkmark$\checkmarks\checkmark$\checkmark^\checkmark\\checkmarkc\checkmarki\checkmarkr\checkmarkc\checkmark$\checkmark \checkmark$\checkmark^\checkmark\\checkmarkc\checkmarki\checkmarkr\checkmarkc\checkmark$\checkmarkd\checkmark$\checkmark^\checkmark\\checkmarkc\checkmarki\checkmarkr\checkmarkc\checkmark$\checkmarke\checkmark$\checkmark^\checkmark\\checkmarkc\checkmarki\checkmarkr\checkmarkc\checkmark$\checkmark \checkmark$\checkmark^\checkmark\\checkmarkc\checkmarki\checkmarkr\checkmarkc\checkmark$\checkmarkc\checkmark$\checkmark^\checkmark\\checkmarkc\checkmarki\checkmarkr\checkmarkc\checkmark$\checkmarko\checkmark$\checkmark^\checkmark\\checkmarkc\checkmarki\checkmarkr\checkmarkc\checkmark$\checkmarku\checkmark$\checkmark^\checkmark\\checkmarkc\checkmarki\checkmarkr\checkmarkc\checkmark$\checkmarkp\checkmark$\checkmark^\checkmark\\checkmarkc\checkmarki\checkmarkr\checkmarkc\checkmark$\checkmarke\checkmark$\checkmark^\checkmark\\checkmarkc\checkmarki\checkmarkr\checkmarkc\checkmark$\checkmark}\checkmark$\checkmark^\checkmark\\checkmarkc\checkmarki\checkmarkr\checkmarkc\checkmark$\checkmark
\checkmark$\checkmark^\checkmark\\checkmarkc\checkmarki\checkmarkr\checkmarkc\checkmark$\checkmarkL\checkmark$\checkmark^\checkmark\\checkmarkc\checkmarki\checkmarkr\checkmarkc\checkmark$\checkmarke\checkmark$\checkmark^\checkmark\\checkmarkc\checkmarki\checkmarkr\checkmarkc\checkmark$\checkmarks\checkmark$\checkmark^\checkmark\\checkmarkc\checkmarki\checkmarkr\checkmarkc\checkmark$\checkmark \checkmark$\checkmark^\checkmark\\checkmarkc\checkmarki\checkmarkr\checkmarkc\checkmark$\checkmarke\checkmark$\checkmark^\checkmark\\checkmarkc\checkmarki\checkmarkr\checkmarkc\checkmark$\checkmarkf\checkmark$\checkmark^\checkmark\\checkmarkc\checkmarki\checkmarkr\checkmarkc\checkmark$\checkmarkf\checkmark$\checkmark^\checkmark\\checkmarkc\checkmarki\checkmarkr\checkmarkc\checkmark$\checkmarko\checkmark$\checkmark^\checkmark\\checkmarkc\checkmarki\checkmarkr\checkmarkc\checkmark$\checkmarkr\checkmark$\checkmark^\checkmark\\checkmarkc\checkmarki\checkmarkr\checkmarkc\checkmark$\checkmarkt\checkmark$\checkmark^\checkmark\\checkmarkc\checkmarki\checkmarkr\checkmarkc\checkmark$\checkmarks\checkmark$\checkmark^\checkmark\\checkmarkc\checkmarki\checkmarkr\checkmarkc\checkmark$\checkmark \checkmark$\checkmark^\checkmark\\checkmarkc\checkmarki\checkmarkr\checkmarkc\checkmark$\checkmarkd\checkmark$\checkmark^\checkmark\\checkmarkc\checkmarki\checkmarkr\checkmarkc\checkmark$\checkmarke\checkmark$\checkmark^\checkmark\\checkmarkc\checkmarki\checkmarkr\checkmarkc\checkmark$\checkmark \checkmark$\checkmark^\checkmark\\checkmarkc\checkmarki\checkmarkr\checkmarkc\checkmark$\checkmarkc\checkmark$\checkmark^\checkmark\\checkmarkc\checkmarki\checkmarkr\checkmarkc\checkmark$\checkmarko\checkmark$\checkmark^\checkmark\\checkmarkc\checkmarki\checkmarkr\checkmarkc\checkmark$\checkmarku\checkmark$\checkmark^\checkmark\\checkmarkc\checkmarki\checkmarkr\checkmarkc\checkmark$\checkmarkp\checkmark$\checkmark^\checkmark\\checkmarkc\checkmarki\checkmarkr\checkmarkc\checkmark$\checkmarke\checkmark$\checkmark^\checkmark\\checkmarkc\checkmarki\checkmarkr\checkmarkc\checkmark$\checkmark \checkmark$\checkmark^\checkmark\\checkmarkc\checkmarki\checkmarkr\checkmarkc\checkmark$\checkmarkc\checkmark$\checkmark^\checkmark\\checkmarkc\checkmarki\checkmarkr\checkmarkc\checkmark$\checkmarko\checkmark$\checkmark^\checkmark\\checkmarkc\checkmarki\checkmarkr\checkmarkc\checkmark$\checkmarkn\checkmark$\checkmark^\checkmark\\checkmarkc\checkmarki\checkmarkr\checkmarkc\checkmark$\checkmarks\checkmark$\checkmark^\checkmark\\checkmarkc\checkmarki\checkmarkr\checkmarkc\checkmark$\checkmarkt\checkmark$\checkmark^\checkmark\\checkmarkc\checkmarki\checkmarkr\checkmarkc\checkmark$\checkmarki\checkmark$\checkmark^\checkmark\\checkmarkc\checkmarki\checkmarkr\checkmarkc\checkmark$\checkmarkt\checkmark$\checkmark^\checkmark\\checkmarkc\checkmarki\checkmarkr\checkmarkc\checkmark$\checkmarku\checkmark$\checkmark^\checkmark\\checkmarkc\checkmarki\checkmarkr\checkmarkc\checkmark$\checkmarke\checkmark$\checkmark^\checkmark\\checkmarkc\checkmarki\checkmarkr\checkmarkc\checkmark$\checkmarkn\checkmark$\checkmark^\checkmark\\checkmarkc\checkmarki\checkmarkr\checkmarkc\checkmark$\checkmarkt\checkmark$\checkmark^\checkmark\\checkmarkc\checkmarki\checkmarkr\checkmarkc\checkmark$\checkmark \checkmark$\checkmark^\checkmark\\checkmarkc\checkmarki\checkmarkr\checkmarkc\checkmark$\checkmarkl\checkmark$\checkmark^\checkmark\\checkmarkc\checkmarki\checkmarkr\checkmarkc\checkmark$\checkmarke\checkmark$\checkmark^\checkmark\\checkmarkc\checkmarki\checkmarkr\checkmarkc\checkmark$\checkmarks\checkmark$\checkmark^\checkmark\\checkmarkc\checkmarki\checkmarkr\checkmarkc\checkmark$\checkmark \checkmark$\checkmark^\checkmark\\checkmarkc\checkmarki\checkmarkr\checkmarkc\checkmark$\checkmark\\checkmark$\checkmark^\checkmark\\checkmarkc\checkmarki\checkmarkr\checkmarkc\checkmark$\checkmarkt\checkmark$\checkmark^\checkmark\\checkmarkc\checkmarki\checkmarkr\checkmarkc\checkmark$\checkmarke\checkmark$\checkmark^\checkmark\\checkmarkc\checkmarki\checkmarkr\checkmarkc\checkmark$\checkmarkx\checkmark$\checkmark^\checkmark\\checkmarkc\checkmarki\checkmarkr\checkmarkc\checkmark$\checkmarkt\checkmark$\checkmark^\checkmark\\checkmarkc\checkmarki\checkmarkr\checkmarkc\checkmark$\checkmarkb\checkmark$\checkmark^\checkmark\\checkmarkc\checkmarki\checkmarkr\checkmarkc\checkmark$\checkmarkf\checkmark$\checkmark^\checkmark\\checkmarkc\checkmarki\checkmarkr\checkmarkc\checkmark$\checkmark{\checkmark$\checkmark^\checkmark\\checkmarkc\checkmarki\checkmarkr\checkmarkc\checkmark$\checkmarkd\checkmark$\checkmark^\checkmark\\checkmarkc\checkmarki\checkmarkr\checkmarkc\checkmark$\checkmarko\checkmark$\checkmark^\checkmark\\checkmarkc\checkmarki\checkmarkr\checkmarkc\checkmark$\checkmarkn\checkmark$\checkmark^\checkmark\\checkmarkc\checkmarki\checkmarkr\checkmarkc\checkmark$\checkmarkn\checkmark$\checkmark^\checkmark\\checkmarkc\checkmarki\checkmarkr\checkmarkc\checkmark$\checkmarké\checkmark$\checkmark^\checkmark\\checkmarkc\checkmarki\checkmarkr\checkmarkc\checkmark$\checkmarke\checkmark$\checkmark^\checkmark\\checkmarkc\checkmarki\checkmarkr\checkmarkc\checkmark$\checkmarks\checkmark$\checkmark^\checkmark\\checkmarkc\checkmarki\checkmarkr\checkmarkc\checkmark$\checkmark \checkmark$\checkmark^\checkmark\\checkmarkc\checkmarki\checkmarkr\checkmarkc\checkmark$\checkmarkd\checkmark$\checkmark^\checkmark\\checkmarkc\checkmarki\checkmarkr\checkmarkc\checkmark$\checkmark'\checkmark$\checkmark^\checkmark\\checkmarkc\checkmarki\checkmarkr\checkmarkc\checkmark$\checkmarke\checkmark$\checkmark^\checkmark\\checkmarkc\checkmarki\checkmarkr\checkmarkc\checkmark$\checkmarkn\checkmark$\checkmark^\checkmark\\checkmarkc\checkmarki\checkmarkr\checkmarkc\checkmark$\checkmarkt\checkmark$\checkmark^\checkmark\\checkmarkc\checkmarki\checkmarkr\checkmarkc\checkmark$\checkmarkr\checkmark$\checkmark^\checkmark\\checkmarkc\checkmarki\checkmarkr\checkmarkc\checkmark$\checkmarké\checkmark$\checkmark^\checkmark\\checkmarkc\checkmarki\checkmarkr\checkmarkc\checkmark$\checkmarke\checkmark$\checkmark^\checkmark\\checkmarkc\checkmarki\checkmarkr\checkmarkc\checkmark$\checkmark \checkmark$\checkmark^\checkmark\\checkmarkc\checkmarki\checkmarkr\checkmarkc\checkmark$\checkmarkf\checkmark$\checkmark^\checkmark\\checkmarkc\checkmarki\checkmarkr\checkmarkc\checkmark$\checkmarko\checkmark$\checkmark^\checkmark\\checkmarkc\checkmarki\checkmarkr\checkmarkc\checkmark$\checkmarkn\checkmark$\checkmark^\checkmark\\checkmarkc\checkmarki\checkmarkr\checkmarkc\checkmark$\checkmarkd\checkmark$\checkmark^\checkmark\\checkmarkc\checkmarki\checkmarkr\checkmarkc\checkmark$\checkmarka\checkmark$\checkmark^\checkmark\\checkmarkc\checkmarki\checkmarkr\checkmarkc\checkmark$\checkmarkm\checkmark$\checkmark^\checkmark\\checkmarkc\checkmarki\checkmarkr\checkmarkc\checkmark$\checkmarke\checkmark$\checkmark^\checkmark\\checkmarkc\checkmarki\checkmarkr\checkmarkc\checkmark$\checkmarkn\checkmark$\checkmark^\checkmark\\checkmarkc\checkmarki\checkmarkr\checkmarkc\checkmark$\checkmarkt\checkmark$\checkmark^\checkmark\\checkmarkc\checkmarki\checkmarkr\checkmarkc\checkmark$\checkmarka\checkmark$\checkmark^\checkmark\\checkmarkc\checkmarki\checkmarkr\checkmarkc\checkmark$\checkmarkl\checkmark$\checkmark^\checkmark\\checkmarkc\checkmarki\checkmarkr\checkmarkc\checkmark$\checkmarke\checkmark$\checkmark^\checkmark\\checkmarkc\checkmarki\checkmarkr\checkmarkc\checkmark$\checkmarks\checkmark$\checkmark^\checkmark\\checkmarkc\checkmarki\checkmarkr\checkmarkc\checkmark$\checkmark}\checkmark$\checkmark^\checkmark\\checkmarkc\checkmarki\checkmarkr\checkmarkc\checkmark$\checkmark \checkmark$\checkmark^\checkmark\\checkmarkc\checkmarki\checkmarkr\checkmarkc\checkmark$\checkmarkd\checkmark$\checkmark^\checkmark\\checkmarkc\checkmarki\checkmarkr\checkmarkc\checkmark$\checkmarke\checkmark$\checkmark^\checkmark\\checkmarkc\checkmarki\checkmarkr\checkmarkc\checkmark$\checkmark \checkmark$\checkmark^\checkmark\\checkmarkc\checkmarki\checkmarkr\checkmarkc\checkmark$\checkmarkt\checkmark$\checkmark^\checkmark\\checkmarkc\checkmarki\checkmarkr\checkmarkc\checkmark$\checkmarko\checkmark$\checkmark^\checkmark\\checkmarkc\checkmarki\checkmarkr\checkmarkc\checkmark$\checkmarku\checkmark$\checkmark^\checkmark\\checkmarkc\checkmarki\checkmarkr\checkmarkc\checkmark$\checkmarkt\checkmark$\checkmark^\checkmark\\checkmarkc\checkmarki\checkmarkr\checkmarkc\checkmark$\checkmark \checkmark$\checkmark^\checkmark\\checkmarkc\checkmarki\checkmarkr\checkmarkc\checkmark$\checkmarkl\checkmark$\checkmark^\checkmark\\checkmarkc\checkmarki\checkmarkr\checkmarkc\checkmark$\checkmarke\checkmark$\checkmark^\checkmark\\checkmarkc\checkmarki\checkmarkr\checkmarkc\checkmark$\checkmark \checkmark$\checkmark^\checkmark\\checkmarkc\checkmarki\checkmarkr\checkmarkc\checkmark$\checkmarkd\checkmark$\checkmark^\checkmark\\checkmarkc\checkmarki\checkmarkr\checkmarkc\checkmark$\checkmarki\checkmark$\checkmark^\checkmark\\checkmarkc\checkmarki\checkmarkr\checkmarkc\checkmark$\checkmarkm\checkmark$\checkmark^\checkmark\\checkmarkc\checkmarki\checkmarkr\checkmarkc\checkmark$\checkmarke\checkmark$\checkmark^\checkmark\\checkmarkc\checkmarki\checkmarkr\checkmarkc\checkmark$\checkmarkn\checkmark$\checkmark^\checkmark\\checkmarkc\checkmarki\checkmarkr\checkmarkc\checkmark$\checkmarks\checkmark$\checkmark^\checkmark\\checkmarkc\checkmarki\checkmarkr\checkmarkc\checkmark$\checkmarki\checkmark$\checkmark^\checkmark\\checkmarkc\checkmarki\checkmarkr\checkmarkc\checkmark$\checkmarko\checkmark$\checkmark^\checkmark\\checkmarkc\checkmarki\checkmarkr\checkmarkc\checkmark$\checkmarkn\checkmark$\checkmark^\checkmark\\checkmarkc\checkmarki\checkmarkr\checkmarkc\checkmark$\checkmarkn\checkmark$\checkmark^\checkmark\\checkmarkc\checkmarki\checkmarkr\checkmarkc\checkmark$\checkmarke\checkmark$\checkmark^\checkmark\\checkmarkc\checkmarki\checkmarkr\checkmarkc\checkmark$\checkmarkm\checkmark$\checkmark^\checkmark\\checkmarkc\checkmarki\checkmarkr\checkmarkc\checkmark$\checkmarke\checkmark$\checkmark^\checkmark\\checkmarkc\checkmarki\checkmarkr\checkmarkc\checkmark$\checkmarkn\checkmark$\checkmark^\checkmark\\checkmarkc\checkmarki\checkmarkr\checkmarkc\checkmark$\checkmarkt\checkmark$\checkmark^\checkmark\\checkmarkc\checkmarki\checkmarkr\checkmarkc\checkmark$\checkmark \checkmark$\checkmark^\checkmark\\checkmarkc\checkmarki\checkmarkr\checkmarkc\checkmark$\checkmarkm\checkmark$\checkmark^\checkmark\\checkmarkc\checkmarki\checkmarkr\checkmarkc\checkmark$\checkmarké\checkmark$\checkmark^\checkmark\\checkmarkc\checkmarki\checkmarkr\checkmarkc\checkmark$\checkmarkc\checkmark$\checkmark^\checkmark\\checkmarkc\checkmarki\checkmarkr\checkmarkc\checkmark$\checkmarka\checkmark$\checkmark^\checkmark\\checkmarkc\checkmarki\checkmarkr\checkmarkc\checkmark$\checkmarkn\checkmark$\checkmark^\checkmark\\checkmarkc\checkmarki\checkmarkr\checkmarkc\checkmark$\checkmarki\checkmark$\checkmark^\checkmark\\checkmarkc\checkmarki\checkmarkr\checkmarkc\checkmark$\checkmarkq\checkmark$\checkmark^\checkmark\\checkmarkc\checkmarki\checkmarkr\checkmarkc\checkmark$\checkmarku\checkmark$\checkmark^\checkmark\\checkmarkc\checkmarki\checkmarkr\checkmarkc\checkmark$\checkmarke\checkmark$\checkmark^\checkmark\\checkmarkc\checkmarki\checkmarkr\checkmarkc\checkmark$\checkmark \checkmark$\checkmark^\checkmark\\checkmarkc\checkmarki\checkmarkr\checkmarkc\checkmark$\checkmarkd\checkmark$\checkmark^\checkmark\\checkmarkc\checkmarki\checkmarkr\checkmarkc\checkmark$\checkmarke\checkmark$\checkmark^\checkmark\\checkmarkc\checkmarki\checkmarkr\checkmarkc\checkmark$\checkmark \checkmark$\checkmark^\checkmark\\checkmarkc\checkmarki\checkmarkr\checkmarkc\checkmark$\checkmarkl\checkmark$\checkmark^\checkmark\\checkmarkc\checkmarki\checkmarkr\checkmarkc\checkmark$\checkmarka\checkmark$\checkmark^\checkmark\\checkmarkc\checkmarki\checkmarkr\checkmarkc\checkmark$\checkmark \checkmark$\checkmark^\checkmark\\checkmarkc\checkmarki\checkmarkr\checkmarkc\checkmark$\checkmarkm\checkmark$\checkmark^\checkmark\\checkmarkc\checkmarki\checkmarkr\checkmarkc\checkmark$\checkmarka\checkmark$\checkmark^\checkmark\\checkmarkc\checkmarki\checkmarkr\checkmarkc\checkmark$\checkmarkc\checkmark$\checkmark^\checkmark\\checkmarkc\checkmarki\checkmarkr\checkmarkc\checkmark$\checkmarkh\checkmark$\checkmark^\checkmark\\checkmarkc\checkmarki\checkmarkr\checkmarkc\checkmark$\checkmarki\checkmark$\checkmark^\checkmark\\checkmarkc\checkmarki\checkmarkr\checkmarkc\checkmark$\checkmarkn\checkmark$\checkmark^\checkmark\\checkmarkc\checkmarki\checkmarkr\checkmarkc\checkmark$\checkmarke\checkmark$\checkmark^\checkmark\\checkmarkc\checkmarki\checkmarkr\checkmarkc\checkmark$\checkmark.\checkmark$\checkmark^\checkmark\\checkmarkc\checkmarki\checkmarkr\checkmarkc\checkmark$\checkmark \checkmark$\checkmark^\checkmark\\checkmarkc\checkmarki\checkmarkr\checkmarkc\checkmark$\checkmarkL\checkmark$\checkmark^\checkmark\\checkmarkc\checkmarki\checkmarkr\checkmarkc\checkmark$\checkmark'\checkmark$\checkmark^\checkmark\\checkmarkc\checkmarki\checkmarkr\checkmarkc\checkmark$\checkmarké\checkmark$\checkmark^\checkmark\\checkmarkc\checkmarki\checkmarkr\checkmarkc\checkmark$\checkmarkt\checkmark$\checkmark^\checkmark\\checkmarkc\checkmarki\checkmarkr\checkmarkc\checkmark$\checkmarku\checkmark$\checkmark^\checkmark\\checkmarkc\checkmarki\checkmarkr\checkmarkc\checkmark$\checkmarkd\checkmark$\checkmark^\checkmark\\checkmarkc\checkmarki\checkmarkr\checkmarkc\checkmark$\checkmarke\checkmark$\checkmark^\checkmark\\checkmarkc\checkmarki\checkmarkr\checkmarkc\checkmark$\checkmark \checkmark$\checkmark^\checkmark\\checkmarkc\checkmarki\checkmarkr\checkmarkc\checkmark$\checkmarkp\checkmark$\checkmark^\checkmark\\checkmarkc\checkmarki\checkmarkr\checkmarkc\checkmark$\checkmarkr\checkmark$\checkmark^\checkmark\\checkmarkc\checkmarki\checkmarkr\checkmarkc\checkmark$\checkmarke\checkmark$\checkmark^\checkmark\\checkmarkc\checkmarki\checkmarkr\checkmarkc\checkmark$\checkmarkn\checkmark$\checkmark^\checkmark\\checkmarkc\checkmarki\checkmarkr\checkmarkc\checkmark$\checkmarkd\checkmark$\checkmark^\checkmark\\checkmarkc\checkmarki\checkmarkr\checkmarkc\checkmark$\checkmark,\checkmark$\checkmark^\checkmark\\checkmarkc\checkmarki\checkmarkr\checkmarkc\checkmark$\checkmark \checkmark$\checkmark^\checkmark\\checkmarkc\checkmarki\checkmarkr\checkmarkc\checkmark$\checkmarkp\checkmark$\checkmark^\checkmark\\checkmarkc\checkmarki\checkmarkr\checkmarkc\checkmark$\checkmarko\checkmark$\checkmark^\checkmark\\checkmarkc\checkmarki\checkmarkr\checkmarkc\checkmark$\checkmarku\checkmark$\checkmark^\checkmark\\checkmarkc\checkmarki\checkmarkr\checkmarkc\checkmark$\checkmarkr\checkmark$\checkmark^\checkmark\\checkmarkc\checkmarki\checkmarkr\checkmarkc\checkmark$\checkmark \checkmark$\checkmark^\checkmark\\checkmarkc\checkmarki\checkmarkr\checkmarkc\checkmark$\checkmarkc\checkmark$\checkmark^\checkmark\\checkmarkc\checkmarki\checkmarkr\checkmarkc\checkmark$\checkmarka\checkmark$\checkmark^\checkmark\\checkmarkc\checkmarki\checkmarkr\checkmarkc\checkmark$\checkmarks\checkmark$\checkmark^\checkmark\\checkmarkc\checkmarki\checkmarkr\checkmarkc\checkmark$\checkmark \checkmark$\checkmark^\checkmark\\checkmarkc\checkmarki\checkmarkr\checkmarkc\checkmark$\checkmarkl\checkmark$\checkmark^\checkmark\\checkmarkc\checkmarki\checkmarkr\checkmarkc\checkmark$\checkmarke\checkmark$\checkmark^\checkmark\\checkmarkc\checkmarki\checkmarkr\checkmarkc\checkmark$\checkmark \checkmark$\checkmark^\checkmark\\checkmarkc\checkmarki\checkmarkr\checkmarkc\checkmark$\checkmarkp\checkmark$\checkmark^\checkmark\\checkmarkc\checkmarki\checkmarkr\checkmarkc\checkmark$\checkmarkl\checkmark$\checkmark^\checkmark\\checkmarkc\checkmarki\checkmarkr\checkmarkc\checkmark$\checkmarku\checkmark$\checkmark^\checkmark\\checkmarkc\checkmarki\checkmarkr\checkmarkc\checkmark$\checkmarks\checkmark$\checkmark^\checkmark\\checkmarkc\checkmarki\checkmarkr\checkmarkc\checkmark$\checkmark \checkmark$\checkmark^\checkmark\\checkmarkc\checkmarki\checkmarkr\checkmarkc\checkmark$\checkmarkd\checkmark$\checkmark^\checkmark\\checkmarkc\checkmarki\checkmarkr\checkmarkc\checkmark$\checkmarké\checkmark$\checkmark^\checkmark\\checkmarkc\checkmarki\checkmarkr\checkmarkc\checkmark$\checkmarkf\checkmark$\checkmark^\checkmark\\checkmarkc\checkmarki\checkmarkr\checkmarkc\checkmark$\checkmarka\checkmark$\checkmark^\checkmark\\checkmarkc\checkmarki\checkmarkr\checkmarkc\checkmark$\checkmarkv\checkmark$\checkmark^\checkmark\\checkmarkc\checkmarki\checkmarkr\checkmarkc\checkmark$\checkmarko\checkmark$\checkmark^\checkmark\\checkmarkc\checkmarki\checkmarkr\checkmarkc\checkmark$\checkmarkr\checkmark$\checkmark^\checkmark\\checkmarkc\checkmarki\checkmarkr\checkmarkc\checkmark$\checkmarka\checkmark$\checkmark^\checkmark\\checkmarkc\checkmarki\checkmarkr\checkmarkc\checkmark$\checkmarkb\checkmark$\checkmark^\checkmark\\checkmarkc\checkmarki\checkmarkr\checkmarkc\checkmark$\checkmarkl\checkmark$\checkmark^\checkmark\\checkmarkc\checkmarki\checkmarkr\checkmarkc\checkmark$\checkmarke\checkmark$\checkmark^\checkmark\\checkmarkc\checkmarki\checkmarkr\checkmarkc\checkmark$\checkmark,\checkmark$\checkmark^\checkmark\\checkmarkc\checkmarki\checkmarkr\checkmarkc\checkmark$\checkmark \checkmark$\checkmark^\checkmark\\checkmarkc\checkmarki\checkmarkr\checkmarkc\checkmark$\checkmarkl\checkmark$\checkmark^\checkmark\\checkmarkc\checkmarki\checkmarkr\checkmarkc\checkmark$\checkmarke\checkmark$\checkmark^\checkmark\\checkmarkc\checkmarki\checkmarkr\checkmarkc\checkmark$\checkmark \checkmark$\checkmark^\checkmark\\checkmarkc\checkmarki\checkmarkr\checkmarkc\checkmark$\checkmarkf\checkmark$\checkmark^\checkmark\\checkmarkc\checkmarki\checkmarkr\checkmarkc\checkmark$\checkmarkr\checkmark$\checkmark^\checkmark\\checkmarkc\checkmarki\checkmarkr\checkmarkc\checkmark$\checkmarka\checkmark$\checkmark^\checkmark\\checkmarkc\checkmarki\checkmarkr\checkmarkc\checkmark$\checkmarki\checkmark$\checkmark^\checkmark\\checkmarkc\checkmarki\checkmarkr\checkmarkc\checkmark$\checkmarks\checkmark$\checkmark^\checkmark\\checkmarkc\checkmarki\checkmarkr\checkmarkc\checkmark$\checkmarka\checkmark$\checkmark^\checkmark\\checkmarkc\checkmarki\checkmarkr\checkmarkc\checkmark$\checkmarkg\checkmark$\checkmark^\checkmark\\checkmarkc\checkmarki\checkmarkr\checkmarkc\checkmark$\checkmarke\checkmark$\checkmark^\checkmark\\checkmarkc\checkmarki\checkmarkr\checkmarkc\checkmark$\checkmark \checkmark$\checkmark^\checkmark\\checkmarkc\checkmarki\checkmarkr\checkmarkc\checkmark$\checkmarke\checkmark$\checkmark^\checkmark\\checkmarkc\checkmarki\checkmarkr\checkmarkc\checkmark$\checkmarkn\checkmark$\checkmark^\checkmark\\checkmarkc\checkmarki\checkmarkr\checkmarkc\checkmark$\checkmark \checkmark$\checkmark^\checkmark\\checkmarkc\checkmarki\checkmarkr\checkmarkc\checkmark$\checkmarkp\checkmark$\checkmark^\checkmark\\checkmarkc\checkmarki\checkmarkr\checkmarkc\checkmark$\checkmarkl\checkmark$\checkmark^\checkmark\\checkmarkc\checkmarki\checkmarkr\checkmarkc\checkmark$\checkmarke\checkmark$\checkmark^\checkmark\\checkmarkc\checkmarki\checkmarkr\checkmarkc\checkmark$\checkmarki\checkmark$\checkmark^\checkmark\\checkmarkc\checkmarki\checkmarkr\checkmarkc\checkmark$\checkmarkn\checkmark$\checkmark^\checkmark\\checkmarkc\checkmarki\checkmarkr\checkmarkc\checkmark$\checkmarke\checkmark$\checkmark^\checkmark\\checkmarkc\checkmarki\checkmarkr\checkmarkc\checkmark$\checkmark \checkmark$\checkmark^\checkmark\\checkmarkc\checkmarki\checkmarkr\checkmarkc\checkmark$\checkmarkm\checkmark$\checkmark^\checkmark\\checkmarkc\checkmarki\checkmarkr\checkmarkc\checkmark$\checkmarka\checkmark$\checkmark^\checkmark\\checkmarkc\checkmarki\checkmarkr\checkmarkc\checkmark$\checkmarkt\checkmark$\checkmark^\checkmark\\checkmarkc\checkmarki\checkmarkr\checkmarkc\checkmark$\checkmarki\checkmark$\checkmark^\checkmark\\checkmarkc\checkmarki\checkmarkr\checkmarkc\checkmark$\checkmarkè\checkmark$\checkmark^\checkmark\\checkmarkc\checkmarki\checkmarkr\checkmarkc\checkmark$\checkmarkr\checkmark$\checkmark^\checkmark\\checkmarkc\checkmarki\checkmarkr\checkmarkc\checkmark$\checkmarke\checkmark$\checkmark^\checkmark\\checkmarkc\checkmarki\checkmarkr\checkmarkc\checkmark$\checkmark \checkmark$\checkmark^\checkmark\\checkmarkc\checkmarki\checkmarkr\checkmarkc\checkmark$\checkmarkd\checkmark$\checkmark^\checkmark\\checkmarkc\checkmarki\checkmarkr\checkmarkc\checkmark$\checkmark'\checkmark$\checkmark^\checkmark\\checkmarkc\checkmarki\checkmarkr\checkmarkc\checkmark$\checkmarku\checkmark$\checkmark^\checkmark\\checkmarkc\checkmarki\checkmarkr\checkmarkc\checkmark$\checkmarkn\checkmark$\checkmark^\checkmark\\checkmarkc\checkmarki\checkmarkr\checkmarkc\checkmark$\checkmarke\checkmark$\checkmark^\checkmark\\checkmarkc\checkmarki\checkmarkr\checkmarkc\checkmark$\checkmark \checkmark$\checkmark^\checkmark\\checkmarkc\checkmarki\checkmarkr\checkmarkc\checkmark$\checkmarkp\checkmark$\checkmark^\checkmark\\checkmarkc\checkmarki\checkmarkr\checkmarkc\checkmark$\checkmarki\checkmark$\checkmark^\checkmark\\checkmarkc\checkmarki\checkmarkr\checkmarkc\checkmark$\checkmarkè\checkmark$\checkmark^\checkmark\\checkmarkc\checkmarki\checkmarkr\checkmarkc\checkmark$\checkmarkc\checkmark$\checkmark^\checkmark\\checkmarkc\checkmarki\checkmarkr\checkmarkc\checkmark$\checkmarke\checkmark$\checkmark^\checkmark\\checkmarkc\checkmarki\checkmarkr\checkmarkc\checkmark$\checkmark \checkmark$\checkmark^\checkmark\\checkmarkc\checkmarki\checkmarkr\checkmarkc\checkmark$\checkmarke\checkmark$\checkmark^\checkmark\\checkmarkc\checkmarki\checkmarkr\checkmarkc\checkmark$\checkmarkn\checkmark$\checkmark^\checkmark\\checkmarkc\checkmarki\checkmarkr\checkmarkc\checkmark$\checkmark \checkmark$\checkmark^\checkmark\\checkmarkc\checkmarki\checkmarkr\checkmarkc\checkmark$\checkmarka\checkmark$\checkmark^\checkmark\\checkmarkc\checkmarki\checkmarkr\checkmarkc\checkmark$\checkmarkl\checkmark$\checkmark^\checkmark\\checkmarkc\checkmarki\checkmarkr\checkmarkc\checkmark$\checkmarkl\checkmark$\checkmark^\checkmark\\checkmarkc\checkmarki\checkmarkr\checkmarkc\checkmark$\checkmarki\checkmark$\checkmark^\checkmark\\checkmarkc\checkmarki\checkmarkr\checkmarkc\checkmark$\checkmarka\checkmark$\checkmark^\checkmark\\checkmarkc\checkmarki\checkmarkr\checkmarkc\checkmark$\checkmarkg\checkmark$\checkmark^\checkmark\\checkmarkc\checkmarki\checkmarkr\checkmarkc\checkmark$\checkmarke\checkmark$\checkmark^\checkmark\\checkmarkc\checkmarki\checkmarkr\checkmarkc\checkmark$\checkmark \checkmark$\checkmark^\checkmark\\checkmarkc\checkmarki\checkmarkr\checkmarkc\checkmark$\checkmarkd\checkmark$\checkmark^\checkmark\\checkmarkc\checkmarki\checkmarkr\checkmarkc\checkmark$\checkmark'\checkmark$\checkmark^\checkmark\\checkmarkc\checkmarki\checkmarkr\checkmarkc\checkmark$\checkmarka\checkmark$\checkmark^\checkmark\\checkmarkc\checkmarki\checkmarkr\checkmarkc\checkmark$\checkmarkl\checkmark$\checkmark^\checkmark\\checkmarkc\checkmarki\checkmarkr\checkmarkc\checkmark$\checkmarku\checkmark$\checkmark^\checkmark\\checkmarkc\checkmarki\checkmarkr\checkmarkc\checkmark$\checkmarkm\checkmark$\checkmark^\checkmark\\checkmarkc\checkmarki\checkmarkr\checkmarkc\checkmark$\checkmarki\checkmark$\checkmark^\checkmark\\checkmarkc\checkmarki\checkmarkr\checkmarkc\checkmark$\checkmarkn\checkmark$\checkmark^\checkmark\\checkmarkc\checkmarki\checkmarkr\checkmarkc\checkmark$\checkmarki\checkmark$\checkmark^\checkmark\\checkmarkc\checkmarki\checkmarkr\checkmarkc\checkmark$\checkmarku\checkmark$\checkmark^\checkmark\\checkmarkc\checkmarki\checkmarkr\checkmarkc\checkmark$\checkmarkm\checkmark$\checkmark^\checkmark\\checkmarkc\checkmarki\checkmarkr\checkmarkc\checkmark$\checkmark \checkmark$\checkmark^\checkmark\\checkmarkc\checkmarki\checkmarkr\checkmarkc\checkmark$\checkmarkA\checkmark$\checkmark^\checkmark\\checkmarkc\checkmarki\checkmarkr\checkmarkc\checkmark$\checkmarkW\checkmark$\checkmark^\checkmark\\checkmarkc\checkmarki\checkmarkr\checkmarkc\checkmark$\checkmark-\checkmark$\checkmark^\checkmark\\checkmarkc\checkmarki\checkmarkr\checkmarkc\checkmark$\checkmark2\checkmark$\checkmark^\checkmark\\checkmarkc\checkmarki\checkmarkr\checkmarkc\checkmark$\checkmark0\checkmark$\checkmark^\checkmark\\checkmarkc\checkmarki\checkmarkr\checkmarkc\checkmark$\checkmark1\checkmark$\checkmark^\checkmark\\checkmarkc\checkmarki\checkmarkr\checkmarkc\checkmark$\checkmark7\checkmark$\checkmark^\checkmark\\checkmarkc\checkmarki\checkmarkr\checkmarkc\checkmark$\checkmarkA\checkmark$\checkmark^\checkmark\\checkmarkc\checkmarki\checkmarkr\checkmarkc\checkmark$\checkmark \checkmark$\checkmark^\checkmark\\checkmarkc\checkmarki\checkmarkr\checkmarkc\checkmark$\checkmark(\checkmark$\checkmark^\checkmark\\checkmarkc\checkmarki\checkmarkr\checkmarkc\checkmark$\checkmarkm\checkmark$\checkmark^\checkmark\\checkmarkc\checkmarki\checkmarkr\checkmarkc\checkmark$\checkmarka\checkmark$\checkmark^\checkmark\\checkmarkc\checkmarki\checkmarkr\checkmarkc\checkmark$\checkmarkt\checkmark$\checkmark^\checkmark\\checkmarkc\checkmarki\checkmarkr\checkmarkc\checkmark$\checkmarké\checkmark$\checkmark^\checkmark\\checkmarkc\checkmarki\checkmarkr\checkmarkc\checkmark$\checkmarkr\checkmark$\checkmark^\checkmark\\checkmarkc\checkmarki\checkmarkr\checkmarkc\checkmark$\checkmarki\checkmark$\checkmark^\checkmark\\checkmarkc\checkmarki\checkmarkr\checkmarkc\checkmark$\checkmarka\checkmark$\checkmark^\checkmark\\checkmarkc\checkmarki\checkmarkr\checkmarkc\checkmark$\checkmarku\checkmark$\checkmark^\checkmark\\checkmarkc\checkmarki\checkmarkr\checkmarkc\checkmark$\checkmark \checkmark$\checkmark^\checkmark\\checkmarkc\checkmarki\checkmarkr\checkmarkc\checkmark$\checkmarkl\checkmark$\checkmark^\checkmark\\checkmarkc\checkmarki\checkmarkr\checkmarkc\checkmark$\checkmarke\checkmark$\checkmark^\checkmark\\checkmarkc\checkmarki\checkmarkr\checkmarkc\checkmark$\checkmark \checkmark$\checkmark^\checkmark\\checkmarkc\checkmarki\checkmarkr\checkmarkc\checkmark$\checkmarkp\checkmark$\checkmark^\checkmark\\checkmarkc\checkmarki\checkmarkr\checkmarkc\checkmark$\checkmarkl\checkmark$\checkmark^\checkmark\\checkmarkc\checkmarki\checkmarkr\checkmarkc\checkmark$\checkmarku\checkmark$\checkmark^\checkmark\\checkmarkc\checkmarki\checkmarkr\checkmarkc\checkmark$\checkmarks\checkmark$\checkmark^\checkmark\\checkmarkc\checkmarki\checkmarkr\checkmarkc\checkmark$\checkmark \checkmark$\checkmark^\checkmark\\checkmarkc\checkmarki\checkmarkr\checkmarkc\checkmark$\checkmarkd\checkmark$\checkmark^\checkmark\\checkmarkc\checkmarki\checkmarkr\checkmarkc\checkmark$\checkmarku\checkmark$\checkmark^\checkmark\\checkmarkc\checkmarki\checkmarkr\checkmarkc\checkmark$\checkmarkr\checkmark$\checkmark^\checkmark\\checkmarkc\checkmarki\checkmarkr\checkmarkc\checkmark$\checkmark \checkmark$\checkmark^\checkmark\\checkmarkc\checkmarki\checkmarkr\checkmarkc\checkmark$\checkmarka\checkmark$\checkmark^\checkmark\\checkmarkc\checkmarki\checkmarkr\checkmarkc\checkmark$\checkmarkd\checkmark$\checkmark^\checkmark\\checkmarkc\checkmarki\checkmarkr\checkmarkc\checkmark$\checkmarkm\checkmark$\checkmark^\checkmark\\checkmarkc\checkmarki\checkmarkr\checkmarkc\checkmark$\checkmarki\checkmark$\checkmark^\checkmark\\checkmarkc\checkmarki\checkmarkr\checkmarkc\checkmark$\checkmarks\checkmark$\checkmark^\checkmark\\checkmarkc\checkmarki\checkmarkr\checkmarkc\checkmark$\checkmark \checkmark$\checkmark^\checkmark\\checkmarkc\checkmarki\checkmarkr\checkmarkc\checkmark$\checkmarkp\checkmark$\checkmark^\checkmark\\checkmarkc\checkmarki\checkmarkr\checkmarkc\checkmark$\checkmarka\checkmark$\checkmark^\checkmark\\checkmarkc\checkmarki\checkmarkr\checkmarkc\checkmark$\checkmarkr\checkmark$\checkmark^\checkmark\\checkmarkc\checkmarki\checkmarkr\checkmarkc\checkmark$\checkmark \checkmark$\checkmark^\checkmark\\checkmarkc\checkmarki\checkmarkr\checkmarkc\checkmark$\checkmarkl\checkmark$\checkmark^\checkmark\\checkmarkc\checkmarki\checkmarkr\checkmarkc\checkmark$\checkmarke\checkmark$\checkmark^\checkmark\\checkmarkc\checkmarki\checkmarkr\checkmarkc\checkmark$\checkmark \checkmark$\checkmark^\checkmark\\checkmarkc\checkmarki\checkmarkr\checkmarkc\checkmark$\checkmarkc\checkmark$\checkmark^\checkmark\\checkmarkc\checkmarki\checkmarkr\checkmarkc\checkmark$\checkmarka\checkmark$\checkmark^\checkmark\\checkmarkc\checkmarki\checkmarkr\checkmarkc\checkmark$\checkmarkh\checkmark$\checkmark^\checkmark\\checkmarkc\checkmarki\checkmarkr\checkmarkc\checkmark$\checkmarki\checkmark$\checkmark^\checkmark\\checkmarkc\checkmarki\checkmarkr\checkmarkc\checkmark$\checkmarke\checkmark$\checkmark^\checkmark\\checkmarkc\checkmarki\checkmarkr\checkmarkc\checkmark$\checkmarkr\checkmark$\checkmark^\checkmark\\checkmarkc\checkmarki\checkmarkr\checkmarkc\checkmark$\checkmark \checkmark$\checkmark^\checkmark\\checkmarkc\checkmarki\checkmarkr\checkmarkc\checkmark$\checkmarkd\checkmark$\checkmark^\checkmark\\checkmarkc\checkmarki\checkmarkr\checkmarkc\checkmark$\checkmarke\checkmark$\checkmark^\checkmark\\checkmarkc\checkmarki\checkmarkr\checkmarkc\checkmark$\checkmarks\checkmark$\checkmark^\checkmark\\checkmarkc\checkmarki\checkmarkr\checkmarkc\checkmark$\checkmark \checkmark$\checkmark^\checkmark\\checkmarkc\checkmarki\checkmarkr\checkmarkc\checkmark$\checkmarkc\checkmark$\checkmark^\checkmark\\checkmarkc\checkmarki\checkmarkr\checkmarkc\checkmark$\checkmarkh\checkmark$\checkmark^\checkmark\\checkmarkc\checkmarki\checkmarkr\checkmarkc\checkmark$\checkmarka\checkmark$\checkmark^\checkmark\\checkmarkc\checkmarki\checkmarkr\checkmarkc\checkmark$\checkmarkr\checkmark$\checkmark^\checkmark\\checkmarkc\checkmarki\checkmarkr\checkmarkc\checkmark$\checkmarkg\checkmark$\checkmark^\checkmark\\checkmarkc\checkmarki\checkmarkr\checkmarkc\checkmark$\checkmarke\checkmark$\checkmark^\checkmark\\checkmarkc\checkmarki\checkmarkr\checkmarkc\checkmark$\checkmarks\checkmark$\checkmark^\checkmark\\checkmarkc\checkmarki\checkmarkr\checkmarkc\checkmark$\checkmark \checkmark$\checkmark^\checkmark\\checkmarkc\checkmarki\checkmarkr\checkmarkc\checkmark$\checkmark—\checkmark$\checkmark^\checkmark\\checkmarkc\checkmarki\checkmarkr\checkmarkc\checkmark$\checkmark \checkmark$\checkmark^\checkmark\\checkmarkc\checkmarki\checkmarkr\checkmarkc\checkmark$\checkmark$\checkmark$\checkmark^\checkmark\\checkmarkc\checkmarki\checkmarkr\checkmarkc\checkmark$\checkmark1\checkmark$\checkmark^\checkmark\\checkmarkc\checkmarki\checkmarkr\checkmarkc\checkmark$\checkmark1\checkmark$\checkmark^\checkmark\\checkmarkc\checkmarki\checkmarkr\checkmarkc\checkmark$\checkmark0\checkmark$\checkmark^\checkmark\\checkmarkc\checkmarki\checkmarkr\checkmarkc\checkmark$\checkmark$\checkmark$\checkmark^\checkmark\\checkmarkc\checkmarki\checkmarkr\checkmarkc\checkmark$\checkmark \checkmark$\checkmark^\checkmark\\checkmarkc\checkmarki\checkmarkr\checkmarkc\checkmark$\checkmarkH\checkmark$\checkmark^\checkmark\\checkmarkc\checkmarki\checkmarkr\checkmarkc\checkmark$\checkmarkB\checkmark$\checkmark^\checkmark\\checkmarkc\checkmarki\checkmarkr\checkmarkc\checkmark$\checkmark)\checkmark$\checkmark^\checkmark\\checkmarkc\checkmarki\checkmarkr\checkmarkc\checkmark$\checkmark.\checkmark$\checkmark^\checkmark\\checkmarkc\checkmarki\checkmarkr\checkmarkc\checkmark$\checkmark
\checkmark$\checkmark^\checkmark\\checkmarkc\checkmarki\checkmarkr\checkmarkc\checkmark$\checkmark
\checkmark$\checkmark^\checkmark\\checkmarkc\checkmarki\checkmarkr\checkmarkc\checkmark$\checkmarkL\checkmark$\checkmark^\checkmark\\checkmarkc\checkmarki\checkmarkr\checkmarkc\checkmark$\checkmarke\checkmark$\checkmark^\checkmark\\checkmarkc\checkmarki\checkmarkr\checkmarkc\checkmark$\checkmarks\checkmark$\checkmark^\checkmark\\checkmarkc\checkmarki\checkmarkr\checkmarkc\checkmark$\checkmark \checkmark$\checkmark^\checkmark\\checkmarkc\checkmarki\checkmarkr\checkmarkc\checkmark$\checkmarkp\checkmark$\checkmark^\checkmark\\checkmarkc\checkmarki\checkmarkr\checkmarkc\checkmark$\checkmarka\checkmark$\checkmark^\checkmark\\checkmarkc\checkmarki\checkmarkr\checkmarkc\checkmark$\checkmarkr\checkmark$\checkmark^\checkmark\\checkmarkc\checkmarki\checkmarkr\checkmarkc\checkmark$\checkmarka\checkmark$\checkmark^\checkmark\\checkmarkc\checkmarki\checkmarkr\checkmarkc\checkmark$\checkmarkm\checkmark$\checkmark^\checkmark\\checkmarkc\checkmarki\checkmarkr\checkmarkc\checkmark$\checkmarkè\checkmark$\checkmark^\checkmark\\checkmarkc\checkmarki\checkmarkr\checkmarkc\checkmark$\checkmarkt\checkmark$\checkmark^\checkmark\\checkmarkc\checkmarki\checkmarkr\checkmarkc\checkmark$\checkmarkr\checkmark$\checkmark^\checkmark\\checkmarkc\checkmarki\checkmarkr\checkmarkc\checkmark$\checkmarke\checkmark$\checkmark^\checkmark\\checkmarkc\checkmarki\checkmarkr\checkmarkc\checkmark$\checkmarks\checkmark$\checkmark^\checkmark\\checkmarkc\checkmarki\checkmarkr\checkmarkc\checkmark$\checkmark \checkmark$\checkmark^\checkmark\\checkmarkc\checkmarki\checkmarkr\checkmarkc\checkmark$\checkmarkd\checkmark$\checkmark^\checkmark\\checkmarkc\checkmarki\checkmarkr\checkmarkc\checkmark$\checkmarke\checkmark$\checkmark^\checkmark\\checkmarkc\checkmarki\checkmarkr\checkmarkc\checkmark$\checkmark \checkmark$\checkmark^\checkmark\\checkmarkc\checkmarki\checkmarkr\checkmarkc\checkmark$\checkmarkc\checkmark$\checkmark^\checkmark\\checkmarkc\checkmarki\checkmarkr\checkmarkc\checkmark$\checkmarko\checkmark$\checkmark^\checkmark\\checkmarkc\checkmarki\checkmarkr\checkmarkc\checkmark$\checkmarku\checkmark$\checkmark^\checkmark\\checkmarkc\checkmarki\checkmarkr\checkmarkc\checkmark$\checkmarkp\checkmark$\checkmark^\checkmark\\checkmarkc\checkmarki\checkmarkr\checkmarkc\checkmark$\checkmarke\checkmark$\checkmark^\checkmark\\checkmarkc\checkmarki\checkmarkr\checkmarkc\checkmark$\checkmark \checkmark$\checkmark^\checkmark\\checkmarkc\checkmarki\checkmarkr\checkmarkc\checkmark$\checkmarkr\checkmark$\checkmark^\checkmark\\checkmarkc\checkmarki\checkmarkr\checkmarkc\checkmark$\checkmarke\checkmark$\checkmark^\checkmark\\checkmarkc\checkmarki\checkmarkr\checkmarkc\checkmark$\checkmarkt\checkmark$\checkmark^\checkmark\\checkmarkc\checkmarki\checkmarkr\checkmarkc\checkmark$\checkmarke\checkmark$\checkmark^\checkmark\\checkmarkc\checkmarki\checkmarkr\checkmarkc\checkmark$\checkmarkn\checkmark$\checkmark^\checkmark\\checkmarkc\checkmarki\checkmarkr\checkmarkc\checkmark$\checkmarku\checkmark$\checkmark^\checkmark\\checkmarkc\checkmarki\checkmarkr\checkmarkc\checkmark$\checkmarks\checkmark$\checkmark^\checkmark\\checkmarkc\checkmarki\checkmarkr\checkmarkc\checkmark$\checkmark \checkmark$\checkmark^\checkmark\\checkmarkc\checkmarki\checkmarkr\checkmarkc\checkmark$\checkmark(\checkmark$\checkmark^\checkmark\\checkmarkc\checkmarki\checkmarkr\checkmarkc\checkmark$\checkmarki\checkmark$\checkmark^\checkmark\\checkmarkc\checkmarki\checkmarkr\checkmarkc\checkmark$\checkmarks\checkmark$\checkmark^\checkmark\\checkmarkc\checkmarki\checkmarkr\checkmarkc\checkmark$\checkmarks\checkmark$\checkmark^\checkmark\\checkmarkc\checkmarki\checkmarkr\checkmarkc\checkmark$\checkmarku\checkmark$\checkmark^\checkmark\\checkmarkc\checkmarki\checkmarkr\checkmarkc\checkmark$\checkmarks\checkmark$\checkmark^\checkmark\\checkmarkc\checkmarki\checkmarkr\checkmarkc\checkmark$\checkmark \checkmark$\checkmark^\checkmark\\checkmarkc\checkmarki\checkmarkr\checkmarkc\checkmark$\checkmarkd\checkmark$\checkmark^\checkmark\\checkmarkc\checkmarki\checkmarkr\checkmarkc\checkmark$\checkmarke\checkmark$\checkmark^\checkmark\\checkmarkc\checkmarki\checkmarkr\checkmarkc\checkmark$\checkmarks\checkmark$\checkmark^\checkmark\\checkmarkc\checkmarki\checkmarkr\checkmarkc\checkmark$\checkmark \checkmark$\checkmark^\checkmark\\checkmarkc\checkmarki\checkmarkr\checkmarkc\checkmark$\checkmarka\checkmark$\checkmark^\checkmark\\checkmarkc\checkmarki\checkmarkr\checkmarkc\checkmark$\checkmarkb\checkmark$\checkmark^\checkmark\\checkmarkc\checkmarki\checkmarkr\checkmarkc\checkmark$\checkmarka\checkmark$\checkmark^\checkmark\\checkmarkc\checkmarki\checkmarkr\checkmarkc\checkmark$\checkmarkq\checkmark$\checkmark^\checkmark\\checkmarkc\checkmarki\checkmarkr\checkmarkc\checkmark$\checkmarku\checkmark$\checkmark^\checkmark\\checkmarkc\checkmarki\checkmarkr\checkmarkc\checkmark$\checkmarke\checkmark$\checkmark^\checkmark\\checkmarkc\checkmarki\checkmarkr\checkmarkc\checkmark$\checkmarks\checkmark$\checkmark^\checkmark\\checkmarkc\checkmarki\checkmarkr\checkmarkc\checkmark$\checkmark \checkmark$\checkmark^\checkmark\\checkmarkc\checkmarki\checkmarkr\checkmarkc\checkmark$\checkmarkd\checkmark$\checkmark^\checkmark\\checkmarkc\checkmarki\checkmarkr\checkmarkc\checkmark$\checkmarke\checkmark$\checkmark^\checkmark\\checkmarkc\checkmarki\checkmarkr\checkmarkc\checkmark$\checkmark \checkmark$\checkmark^\checkmark\\checkmarkc\checkmarki\checkmarkr\checkmarkc\checkmark$\checkmarkf\checkmark$\checkmark^\checkmark\\checkmarkc\checkmarki\checkmarkr\checkmarkc\checkmark$\checkmarkr\checkmark$\checkmark^\checkmark\\checkmarkc\checkmarki\checkmarkr\checkmarkc\checkmark$\checkmarka\checkmark$\checkmark^\checkmark\\checkmarkc\checkmarki\checkmarkr\checkmarkc\checkmark$\checkmarki\checkmark$\checkmark^\checkmark\\checkmarkc\checkmarki\checkmarkr\checkmarkc\checkmark$\checkmarks\checkmark$\checkmark^\checkmark\\checkmarkc\checkmarki\checkmarkr\checkmarkc\checkmark$\checkmarka\checkmark$\checkmark^\checkmark\\checkmarkc\checkmarki\checkmarkr\checkmarkc\checkmark$\checkmarkg\checkmark$\checkmark^\checkmark\\checkmarkc\checkmarki\checkmarkr\checkmarkc\checkmark$\checkmarke\checkmark$\checkmark^\checkmark\\checkmarkc\checkmarki\checkmarkr\checkmarkc\checkmark$\checkmark \checkmark$\checkmark^\checkmark\\checkmarkc\checkmarki\checkmarkr\checkmarkc\checkmark$\checkmarkp\checkmark$\checkmark^\checkmark\\checkmarkc\checkmarki\checkmarkr\checkmarkc\checkmark$\checkmarko\checkmark$\checkmark^\checkmark\\checkmarkc\checkmarki\checkmarkr\checkmarkc\checkmark$\checkmarku\checkmark$\checkmark^\checkmark\\checkmarkc\checkmarki\checkmarkr\checkmarkc\checkmark$\checkmarkr\checkmark$\checkmark^\checkmark\\checkmarkc\checkmarki\checkmarkr\checkmarkc\checkmark$\checkmark \checkmark$\checkmark^\checkmark\\checkmarkc\checkmarki\checkmarkr\checkmarkc\checkmark$\checkmarkl\checkmark$\checkmark^\checkmark\\checkmarkc\checkmarki\checkmarkr\checkmarkc\checkmark$\checkmark'\checkmark$\checkmark^\checkmark\\checkmarkc\checkmarki\checkmarkr\checkmarkc\checkmark$\checkmarka\checkmark$\checkmark^\checkmark\\checkmarkc\checkmarki\checkmarkr\checkmarkc\checkmark$\checkmarkl\checkmark$\checkmark^\checkmark\\checkmarkc\checkmarki\checkmarkr\checkmarkc\checkmark$\checkmarku\checkmark$\checkmark^\checkmark\\checkmarkc\checkmarki\checkmarkr\checkmarkc\checkmark$\checkmarkm\checkmark$\checkmark^\checkmark\\checkmarkc\checkmarki\checkmarkr\checkmarkc\checkmark$\checkmarki\checkmark$\checkmark^\checkmark\\checkmarkc\checkmarki\checkmarkr\checkmarkc\checkmark$\checkmarkn\checkmark$\checkmark^\checkmark\\checkmarkc\checkmarki\checkmarkr\checkmarkc\checkmark$\checkmarki\checkmark$\checkmark^\checkmark\\checkmarkc\checkmarki\checkmarkr\checkmarkc\checkmark$\checkmarku\checkmark$\checkmark^\checkmark\\checkmarkc\checkmarki\checkmarkr\checkmarkc\checkmark$\checkmarkm\checkmark$\checkmark^\checkmark\\checkmarkc\checkmarki\checkmarkr\checkmarkc\checkmark$\checkmark \checkmark$\checkmark^\checkmark\\checkmarkc\checkmarki\checkmarkr\checkmarkc\checkmark$\checkmark—\checkmark$\checkmark^\checkmark\\checkmarkc\checkmarki\checkmarkr\checkmarkc\checkmark$\checkmark \checkmark$\checkmark^\checkmark\\checkmarkc\checkmarki\checkmarkr\checkmarkc\checkmark$\checkmarkS\checkmark$\checkmark^\checkmark\\checkmarkc\checkmarki\checkmarkr\checkmarkc\checkmark$\checkmarka\checkmark$\checkmark^\checkmark\\checkmarkc\checkmarki\checkmarkr\checkmarkc\checkmark$\checkmarkn\checkmark$\checkmark^\checkmark\\checkmarkc\checkmarki\checkmarkr\checkmarkc\checkmark$\checkmarkd\checkmark$\checkmark^\checkmark\\checkmarkc\checkmarki\checkmarkr\checkmarkc\checkmark$\checkmarkv\checkmark$\checkmark^\checkmark\\checkmarkc\checkmarki\checkmarkr\checkmarkc\checkmark$\checkmarki\checkmark$\checkmark^\checkmark\\checkmarkc\checkmarki\checkmarkr\checkmarkc\checkmark$\checkmarkk\checkmark$\checkmark^\checkmark\\checkmarkc\checkmarki\checkmarkr\checkmarkc\checkmark$\checkmark \checkmark$\checkmark^\checkmark\\checkmarkc\checkmarki\checkmarkr\checkmarkc\checkmark$\checkmarkC\checkmark$\checkmark^\checkmark\\checkmarkc\checkmarki\checkmarkr\checkmarkc\checkmark$\checkmarko\checkmark$\checkmark^\checkmark\\checkmarkc\checkmarki\checkmarkr\checkmarkc\checkmark$\checkmarkr\checkmark$\checkmark^\checkmark\\checkmarkc\checkmarki\checkmarkr\checkmarkc\checkmark$\checkmarko\checkmark$\checkmark^\checkmark\\checkmarkc\checkmarki\checkmarkr\checkmarkc\checkmark$\checkmarkm\checkmark$\checkmark^\checkmark\\checkmarkc\checkmarki\checkmarkr\checkmarkc\checkmark$\checkmarka\checkmark$\checkmark^\checkmark\\checkmarkc\checkmarki\checkmarkr\checkmarkc\checkmark$\checkmarkn\checkmark$\checkmark^\checkmark\\checkmarkc\checkmarki\checkmarkr\checkmarkc\checkmark$\checkmarkt\checkmark$\checkmark^\checkmark\\checkmarkc\checkmarki\checkmarkr\checkmarkc\checkmark$\checkmark \checkmark$\checkmark^\checkmark\\checkmarkc\checkmarki\checkmarkr\checkmarkc\checkmark$\checkmark[\checkmark$\checkmark^\checkmark\\checkmarkc\checkmarki\checkmarkr\checkmarkc\checkmark$\checkmarkR\checkmark$\checkmark^\checkmark\\checkmarkc\checkmarki\checkmarkr\checkmarkc\checkmark$\checkmarké\checkmark$\checkmark^\checkmark\\checkmarkc\checkmarki\checkmarkr\checkmarkc\checkmark$\checkmarkf\checkmark$\checkmark^\checkmark\\checkmarkc\checkmarki\checkmarkr\checkmarkc\checkmark$\checkmark.\checkmark$\checkmark^\checkmark\\checkmarkc\checkmarki\checkmarkr\checkmarkc\checkmark$\checkmark]\checkmark$\checkmark^\checkmark\\checkmarkc\checkmarki\checkmarkr\checkmarkc\checkmark$\checkmark)\checkmark$\checkmark^\checkmark\\checkmarkc\checkmarki\checkmarkr\checkmarkc\checkmark$\checkmark \checkmark$\checkmark^\checkmark\\checkmarkc\checkmarki\checkmarkr\checkmarkc\checkmark$\checkmarks\checkmark$\checkmark^\checkmark\\checkmarkc\checkmarki\checkmarkr\checkmarkc\checkmark$\checkmarko\checkmark$\checkmark^\checkmark\\checkmarkc\checkmarki\checkmarkr\checkmarkc\checkmark$\checkmarkn\checkmark$\checkmark^\checkmark\\checkmarkc\checkmarki\checkmarkr\checkmarkc\checkmark$\checkmarkt\checkmark$\checkmark^\checkmark\\checkmarkc\checkmarki\checkmarkr\checkmarkc\checkmark$\checkmark \checkmark$\checkmark^\checkmark\\checkmarkc\checkmarki\checkmarkr\checkmarkc\checkmark$\checkmark:\checkmark$\checkmark^\checkmark\\checkmarkc\checkmarki\checkmarkr\checkmarkc\checkmark$\checkmark
\checkmark$\checkmark^\checkmark\\checkmarkc\checkmarki\checkmarkr\checkmarkc\checkmark$\checkmark
\checkmark$\checkmark^\checkmark\\checkmarkc\checkmarki\checkmarkr\checkmarkc\checkmark$\checkmark\\checkmark$\checkmark^\checkmark\\checkmarkc\checkmarki\checkmarkr\checkmarkc\checkmark$\checkmarkb\checkmark$\checkmark^\checkmark\\checkmarkc\checkmarki\checkmarkr\checkmarkc\checkmark$\checkmarke\checkmark$\checkmark^\checkmark\\checkmarkc\checkmarki\checkmarkr\checkmarkc\checkmark$\checkmarkg\checkmark$\checkmark^\checkmark\\checkmarkc\checkmarki\checkmarkr\checkmarkc\checkmark$\checkmarki\checkmark$\checkmark^\checkmark\\checkmarkc\checkmarki\checkmarkr\checkmarkc\checkmark$\checkmarkn\checkmark$\checkmark^\checkmark\\checkmarkc\checkmarki\checkmarkr\checkmarkc\checkmark$\checkmark{\checkmark$\checkmark^\checkmark\\checkmarkc\checkmarki\checkmarkr\checkmarkc\checkmark$\checkmarkt\checkmark$\checkmark^\checkmark\\checkmarkc\checkmarki\checkmarkr\checkmarkc\checkmark$\checkmarka\checkmark$\checkmark^\checkmark\\checkmarkc\checkmarki\checkmarkr\checkmarkc\checkmark$\checkmarkb\checkmark$\checkmark^\checkmark\\checkmarkc\checkmarki\checkmarkr\checkmarkc\checkmark$\checkmarkl\checkmark$\checkmark^\checkmark\\checkmarkc\checkmarki\checkmarkr\checkmarkc\checkmark$\checkmarke\checkmark$\checkmark^\checkmark\\checkmarkc\checkmarki\checkmarkr\checkmarkc\checkmark$\checkmark}\checkmark$\checkmark^\checkmark\\checkmarkc\checkmarki\checkmarkr\checkmarkc\checkmark$\checkmark[\checkmark$\checkmark^\checkmark\\checkmarkc\checkmarki\checkmarkr\checkmarkc\checkmark$\checkmarkh\checkmark$\checkmark^\checkmark\\checkmarkc\checkmarki\checkmarkr\checkmarkc\checkmark$\checkmarkt\checkmark$\checkmark^\checkmark\\checkmarkc\checkmarki\checkmarkr\checkmarkc\checkmark$\checkmarkb\checkmark$\checkmark^\checkmark\\checkmarkc\checkmarki\checkmarkr\checkmarkc\checkmark$\checkmarkp\checkmark$\checkmark^\checkmark\\checkmarkc\checkmarki\checkmarkr\checkmarkc\checkmark$\checkmark]\checkmark$\checkmark^\checkmark\\checkmarkc\checkmarki\checkmarkr\checkmarkc\checkmark$\checkmark
\checkmark$\checkmark^\checkmark\\checkmarkc\checkmarki\checkmarkr\checkmarkc\checkmark$\checkmark\\checkmark$\checkmark^\checkmark\\checkmarkc\checkmarki\checkmarkr\checkmarkc\checkmark$\checkmarkc\checkmark$\checkmark^\checkmark\\checkmarkc\checkmarki\checkmarkr\checkmarkc\checkmark$\checkmarke\checkmark$\checkmark^\checkmark\\checkmarkc\checkmarki\checkmarkr\checkmarkc\checkmark$\checkmarkn\checkmark$\checkmark^\checkmark\\checkmarkc\checkmarki\checkmarkr\checkmarkc\checkmark$\checkmarkt\checkmark$\checkmark^\checkmark\\checkmarkc\checkmarki\checkmarkr\checkmarkc\checkmark$\checkmarke\checkmark$\checkmark^\checkmark\\checkmarkc\checkmarki\checkmarkr\checkmarkc\checkmark$\checkmarkr\checkmark$\checkmark^\checkmark\\checkmarkc\checkmarki\checkmarkr\checkmarkc\checkmark$\checkmarki\checkmark$\checkmark^\checkmark\\checkmarkc\checkmarki\checkmarkr\checkmarkc\checkmark$\checkmarkn\checkmark$\checkmark^\checkmark\\checkmarkc\checkmarki\checkmarkr\checkmarkc\checkmark$\checkmarkg\checkmark$\checkmark^\checkmark\\checkmarkc\checkmarki\checkmarkr\checkmarkc\checkmark$\checkmark
\checkmark$\checkmark^\checkmark\\checkmarkc\checkmarki\checkmarkr\checkmarkc\checkmark$\checkmark\\checkmark$\checkmark^\checkmark\\checkmarkc\checkmarki\checkmarkr\checkmarkc\checkmark$\checkmarkb\checkmark$\checkmark^\checkmark\\checkmarkc\checkmarki\checkmarkr\checkmarkc\checkmark$\checkmarke\checkmark$\checkmark^\checkmark\\checkmarkc\checkmarki\checkmarkr\checkmarkc\checkmark$\checkmarkg\checkmark$\checkmark^\checkmark\\checkmarkc\checkmarki\checkmarkr\checkmarkc\checkmark$\checkmarki\checkmark$\checkmark^\checkmark\\checkmarkc\checkmarki\checkmarkr\checkmarkc\checkmark$\checkmarkn\checkmark$\checkmark^\checkmark\\checkmarkc\checkmarki\checkmarkr\checkmarkc\checkmark$\checkmark{\checkmark$\checkmark^\checkmark\\checkmarkc\checkmarki\checkmarkr\checkmarkc\checkmark$\checkmarkt\checkmark$\checkmark^\checkmark\\checkmarkc\checkmarki\checkmarkr\checkmarkc\checkmark$\checkmarka\checkmark$\checkmark^\checkmark\\checkmarkc\checkmarki\checkmarkr\checkmarkc\checkmark$\checkmarkb\checkmark$\checkmark^\checkmark\\checkmarkc\checkmarki\checkmarkr\checkmarkc\checkmark$\checkmarku\checkmark$\checkmark^\checkmark\\checkmarkc\checkmarki\checkmarkr\checkmarkc\checkmark$\checkmarkl\checkmark$\checkmark^\checkmark\\checkmarkc\checkmarki\checkmarkr\checkmarkc\checkmark$\checkmarka\checkmark$\checkmark^\checkmark\\checkmarkc\checkmarki\checkmarkr\checkmarkc\checkmark$\checkmarkr\checkmark$\checkmark^\checkmark\\checkmarkc\checkmarki\checkmarkr\checkmarkc\checkmark$\checkmark}\checkmark$\checkmark^\checkmark\\checkmarkc\checkmarki\checkmarkr\checkmarkc\checkmark$\checkmark{\checkmark$\checkmark^\checkmark\\checkmarkc\checkmarki\checkmarkr\checkmarkc\checkmark$\checkmark|\checkmark$\checkmark^\checkmark\\checkmarkc\checkmarki\checkmarkr\checkmarkc\checkmark$\checkmarkl\checkmark$\checkmark^\checkmark\\checkmarkc\checkmarki\checkmarkr\checkmarkc\checkmark$\checkmark|\checkmark$\checkmark^\checkmark\\checkmarkc\checkmarki\checkmarkr\checkmarkc\checkmark$\checkmarkc\checkmark$\checkmark^\checkmark\\checkmarkc\checkmarki\checkmarkr\checkmarkc\checkmark$\checkmark|\checkmark$\checkmark^\checkmark\\checkmarkc\checkmarki\checkmarkr\checkmarkc\checkmark$\checkmarkl\checkmark$\checkmark^\checkmark\\checkmarkc\checkmarki\checkmarkr\checkmarkc\checkmark$\checkmark|\checkmark$\checkmark^\checkmark\\checkmarkc\checkmarki\checkmarkr\checkmarkc\checkmark$\checkmark}\checkmark$\checkmark^\checkmark\\checkmarkc\checkmarki\checkmarkr\checkmarkc\checkmark$\checkmark
\checkmark$\checkmark^\checkmark\\checkmarkc\checkmarki\checkmarkr\checkmarkc\checkmark$\checkmark\\checkmark$\checkmark^\checkmark\\checkmarkc\checkmarki\checkmarkr\checkmarkc\checkmark$\checkmarkh\checkmark$\checkmark^\checkmark\\checkmarkc\checkmarki\checkmarkr\checkmarkc\checkmark$\checkmarkl\checkmark$\checkmark^\checkmark\\checkmarkc\checkmarki\checkmarkr\checkmarkc\checkmark$\checkmarki\checkmark$\checkmark^\checkmark\\checkmarkc\checkmarki\checkmarkr\checkmarkc\checkmark$\checkmarkn\checkmark$\checkmark^\checkmark\\checkmarkc\checkmarki\checkmarkr\checkmarkc\checkmark$\checkmarke\checkmark$\checkmark^\checkmark\\checkmarkc\checkmarki\checkmarkr\checkmarkc\checkmark$\checkmark
\checkmark$\checkmark^\checkmark\\checkmarkc\checkmarki\checkmarkr\checkmarkc\checkmark$\checkmark\\checkmark$\checkmark^\checkmark\\checkmarkc\checkmarki\checkmarkr\checkmarkc\checkmark$\checkmarkt\checkmark$\checkmark^\checkmark\\checkmarkc\checkmarki\checkmarkr\checkmarkc\checkmark$\checkmarke\checkmark$\checkmark^\checkmark\\checkmarkc\checkmarki\checkmarkr\checkmarkc\checkmark$\checkmarkx\checkmark$\checkmark^\checkmark\\checkmarkc\checkmarki\checkmarkr\checkmarkc\checkmark$\checkmarkt\checkmark$\checkmark^\checkmark\\checkmarkc\checkmarki\checkmarkr\checkmarkc\checkmark$\checkmarkb\checkmark$\checkmark^\checkmark\\checkmarkc\checkmarki\checkmarkr\checkmarkc\checkmark$\checkmarkf\checkmark$\checkmark^\checkmark\\checkmarkc\checkmarki\checkmarkr\checkmarkc\checkmark$\checkmark{\checkmark$\checkmark^\checkmark\\checkmarkc\checkmarki\checkmarkr\checkmarkc\checkmark$\checkmarkP\checkmark$\checkmark^\checkmark\\checkmarkc\checkmarki\checkmarkr\checkmarkc\checkmark$\checkmarka\checkmark$\checkmark^\checkmark\\checkmarkc\checkmarki\checkmarkr\checkmarkc\checkmark$\checkmarkr\checkmark$\checkmark^\checkmark\\checkmarkc\checkmarki\checkmarkr\checkmarkc\checkmark$\checkmarka\checkmark$\checkmark^\checkmark\\checkmarkc\checkmarki\checkmarkr\checkmarkc\checkmark$\checkmarkm\checkmark$\checkmark^\checkmark\\checkmarkc\checkmarki\checkmarkr\checkmarkc\checkmark$\checkmarkè\checkmark$\checkmark^\checkmark\\checkmarkc\checkmarki\checkmarkr\checkmarkc\checkmark$\checkmarkt\checkmark$\checkmark^\checkmark\\checkmarkc\checkmarki\checkmarkr\checkmarkc\checkmark$\checkmarkr\checkmark$\checkmark^\checkmark\\checkmarkc\checkmarki\checkmarkr\checkmarkc\checkmark$\checkmarke\checkmark$\checkmark^\checkmark\\checkmarkc\checkmarki\checkmarkr\checkmarkc\checkmark$\checkmark}\checkmark$\checkmark^\checkmark\\checkmarkc\checkmarki\checkmarkr\checkmarkc\checkmark$\checkmark \checkmark$\checkmark^\checkmark\\checkmarkc\checkmarki\checkmarkr\checkmarkc\checkmark$\checkmark&\checkmark$\checkmark^\checkmark\\checkmarkc\checkmarki\checkmarkr\checkmarkc\checkmark$\checkmark \checkmark$\checkmark^\checkmark\\checkmarkc\checkmarki\checkmarkr\checkmarkc\checkmark$\checkmark\\checkmark$\checkmark^\checkmark\\checkmarkc\checkmarki\checkmarkr\checkmarkc\checkmark$\checkmarkt\checkmark$\checkmark^\checkmark\\checkmarkc\checkmarki\checkmarkr\checkmarkc\checkmark$\checkmarke\checkmark$\checkmark^\checkmark\\checkmarkc\checkmarki\checkmarkr\checkmarkc\checkmark$\checkmarkx\checkmark$\checkmark^\checkmark\\checkmarkc\checkmarki\checkmarkr\checkmarkc\checkmark$\checkmarkt\checkmark$\checkmark^\checkmark\\checkmarkc\checkmarki\checkmarkr\checkmarkc\checkmark$\checkmarkb\checkmark$\checkmark^\checkmark\\checkmarkc\checkmarki\checkmarkr\checkmarkc\checkmark$\checkmarkf\checkmark$\checkmark^\checkmark\\checkmarkc\checkmarki\checkmarkr\checkmarkc\checkmark$\checkmark{\checkmark$\checkmark^\checkmark\\checkmarkc\checkmarki\checkmarkr\checkmarkc\checkmark$\checkmarkV\checkmark$\checkmark^\checkmark\\checkmarkc\checkmarki\checkmarkr\checkmarkc\checkmark$\checkmarka\checkmark$\checkmark^\checkmark\\checkmarkc\checkmarki\checkmarkr\checkmarkc\checkmark$\checkmarkl\checkmark$\checkmark^\checkmark\\checkmarkc\checkmarki\checkmarkr\checkmarkc\checkmark$\checkmarke\checkmark$\checkmark^\checkmark\\checkmarkc\checkmarki\checkmarkr\checkmarkc\checkmark$\checkmarku\checkmark$\checkmark^\checkmark\\checkmarkc\checkmarki\checkmarkr\checkmarkc\checkmark$\checkmarkr\checkmark$\checkmark^\checkmark\\checkmarkc\checkmarki\checkmarkr\checkmarkc\checkmark$\checkmark}\checkmark$\checkmark^\checkmark\\checkmarkc\checkmarki\checkmarkr\checkmarkc\checkmark$\checkmark \checkmark$\checkmark^\checkmark\\checkmarkc\checkmarki\checkmarkr\checkmarkc\checkmark$\checkmark&\checkmark$\checkmark^\checkmark\\checkmarkc\checkmarki\checkmarkr\checkmarkc\checkmark$\checkmark \checkmark$\checkmark^\checkmark\\checkmarkc\checkmarki\checkmarkr\checkmarkc\checkmark$\checkmark\\checkmark$\checkmark^\checkmark\\checkmarkc\checkmarki\checkmarkr\checkmarkc\checkmark$\checkmarkt\checkmark$\checkmark^\checkmark\\checkmarkc\checkmarki\checkmarkr\checkmarkc\checkmark$\checkmarke\checkmark$\checkmark^\checkmark\\checkmarkc\checkmarki\checkmarkr\checkmarkc\checkmark$\checkmarkx\checkmark$\checkmark^\checkmark\\checkmarkc\checkmarki\checkmarkr\checkmarkc\checkmark$\checkmarkt\checkmark$\checkmark^\checkmark\\checkmarkc\checkmarki\checkmarkr\checkmarkc\checkmark$\checkmarkb\checkmark$\checkmark^\checkmark\\checkmarkc\checkmarki\checkmarkr\checkmarkc\checkmark$\checkmarkf\checkmark$\checkmark^\checkmark\\checkmarkc\checkmarki\checkmarkr\checkmarkc\checkmark$\checkmark{\checkmark$\checkmark^\checkmark\\checkmarkc\checkmarki\checkmarkr\checkmarkc\checkmark$\checkmarkJ\checkmark$\checkmark^\checkmark\\checkmarkc\checkmarki\checkmarkr\checkmarkc\checkmark$\checkmarku\checkmark$\checkmark^\checkmark\\checkmarkc\checkmarki\checkmarkr\checkmarkc\checkmark$\checkmarks\checkmark$\checkmark^\checkmark\\checkmarkc\checkmarki\checkmarkr\checkmarkc\checkmark$\checkmarkt\checkmark$\checkmark^\checkmark\\checkmarkc\checkmarki\checkmarkr\checkmarkc\checkmark$\checkmarki\checkmark$\checkmark^\checkmark\\checkmarkc\checkmarki\checkmarkr\checkmarkc\checkmark$\checkmarkf\checkmark$\checkmark^\checkmark\\checkmarkc\checkmarki\checkmarkr\checkmarkc\checkmark$\checkmarki\checkmark$\checkmark^\checkmark\\checkmarkc\checkmarki\checkmarkr\checkmarkc\checkmark$\checkmarkc\checkmark$\checkmark^\checkmark\\checkmarkc\checkmarki\checkmarkr\checkmarkc\checkmark$\checkmarka\checkmark$\checkmark^\checkmark\\checkmarkc\checkmarki\checkmarkr\checkmarkc\checkmark$\checkmarkt\checkmark$\checkmark^\checkmark\\checkmarkc\checkmarki\checkmarkr\checkmarkc\checkmark$\checkmarki\checkmark$\checkmark^\checkmark\\checkmarkc\checkmarki\checkmarkr\checkmarkc\checkmark$\checkmarko\checkmark$\checkmark^\checkmark\\checkmarkc\checkmarki\checkmarkr\checkmarkc\checkmark$\checkmarkn\checkmark$\checkmark^\checkmark\\checkmarkc\checkmarki\checkmarkr\checkmarkc\checkmark$\checkmark}\checkmark$\checkmark^\checkmark\\checkmarkc\checkmarki\checkmarkr\checkmarkc\checkmark$\checkmark \checkmark$\checkmark^\checkmark\\checkmarkc\checkmarki\checkmarkr\checkmarkc\checkmark$\checkmark\\checkmark$\checkmark^\checkmark\\checkmarkc\checkmarki\checkmarkr\checkmarkc\checkmark$\checkmark\\checkmark$\checkmark^\checkmark\\checkmarkc\checkmarki\checkmarkr\checkmarkc\checkmark$\checkmark \checkmark$\checkmark^\checkmark\\checkmarkc\checkmarki\checkmarkr\checkmarkc\checkmark$\checkmark\\checkmark$\checkmark^\checkmark\\checkmarkc\checkmarki\checkmarkr\checkmarkc\checkmark$\checkmarkh\checkmark$\checkmark^\checkmark\\checkmarkc\checkmarki\checkmarkr\checkmarkc\checkmark$\checkmarkl\checkmark$\checkmark^\checkmark\\checkmarkc\checkmarki\checkmarkr\checkmarkc\checkmark$\checkmarki\checkmark$\checkmark^\checkmark\\checkmarkc\checkmarki\checkmarkr\checkmarkc\checkmark$\checkmarkn\checkmark$\checkmark^\checkmark\\checkmarkc\checkmarki\checkmarkr\checkmarkc\checkmark$\checkmarke\checkmark$\checkmark^\checkmark\\checkmarkc\checkmarki\checkmarkr\checkmarkc\checkmark$\checkmark
\checkmark$\checkmark^\checkmark\\checkmarkc\checkmarki\checkmarkr\checkmarkc\checkmark$\checkmarkD\checkmark$\checkmark^\checkmark\\checkmarkc\checkmarki\checkmarkr\checkmarkc\checkmark$\checkmarki\checkmark$\checkmark^\checkmark\\checkmarkc\checkmarki\checkmarkr\checkmarkc\checkmark$\checkmarka\checkmark$\checkmark^\checkmark\\checkmarkc\checkmarki\checkmarkr\checkmarkc\checkmark$\checkmarkm\checkmark$\checkmark^\checkmark\\checkmarkc\checkmarki\checkmarkr\checkmarkc\checkmark$\checkmarkè\checkmark$\checkmark^\checkmark\\checkmarkc\checkmarki\checkmarkr\checkmarkc\checkmark$\checkmarkt\checkmark$\checkmark^\checkmark\\checkmarkc\checkmarki\checkmarkr\checkmarkc\checkmark$\checkmarkr\checkmark$\checkmark^\checkmark\\checkmarkc\checkmarki\checkmarkr\checkmarkc\checkmark$\checkmarke\checkmark$\checkmark^\checkmark\\checkmarkc\checkmarki\checkmarkr\checkmarkc\checkmark$\checkmark \checkmark$\checkmark^\checkmark\\checkmarkc\checkmarki\checkmarkr\checkmarkc\checkmark$\checkmarkd\checkmark$\checkmark^\checkmark\\checkmarkc\checkmarki\checkmarkr\checkmarkc\checkmark$\checkmarke\checkmark$\checkmark^\checkmark\\checkmarkc\checkmarki\checkmarkr\checkmarkc\checkmark$\checkmark \checkmark$\checkmark^\checkmark\\checkmarkc\checkmarki\checkmarkr\checkmarkc\checkmark$\checkmarkf\checkmark$\checkmark^\checkmark\\checkmarkc\checkmarki\checkmarkr\checkmarkc\checkmark$\checkmarkr\checkmark$\checkmark^\checkmark\\checkmarkc\checkmarki\checkmarkr\checkmarkc\checkmark$\checkmarka\checkmark$\checkmark^\checkmark\\checkmarkc\checkmarki\checkmarkr\checkmarkc\checkmark$\checkmarki\checkmark$\checkmark^\checkmark\\checkmarkc\checkmarki\checkmarkr\checkmarkc\checkmark$\checkmarks\checkmark$\checkmark^\checkmark\\checkmarkc\checkmarki\checkmarkr\checkmarkc\checkmark$\checkmarke\checkmark$\checkmark^\checkmark\\checkmarkc\checkmarki\checkmarkr\checkmarkc\checkmark$\checkmark \checkmark$\checkmark^\checkmark\\checkmarkc\checkmarki\checkmarkr\checkmarkc\checkmark$\checkmark(\checkmark$\checkmark^\checkmark\\checkmarkc\checkmarki\checkmarkr\checkmarkc\checkmark$\checkmark$\checkmark$\checkmark^\checkmark\\checkmarkc\checkmarki\checkmarkr\checkmarkc\checkmark$\checkmarkD\checkmark$\checkmark^\checkmark\\checkmarkc\checkmarki\checkmarkr\checkmarkc\checkmark$\checkmark_\checkmark$\checkmark^\checkmark\\checkmarkc\checkmarki\checkmarkr\checkmarkc\checkmark$\checkmarkc\checkmark$\checkmark^\checkmark\\checkmarkc\checkmarki\checkmarkr\checkmarkc\checkmark$\checkmark$\checkmark$\checkmark^\checkmark\\checkmarkc\checkmarki\checkmarkr\checkmarkc\checkmark$\checkmark)\checkmark$\checkmark^\checkmark\\checkmarkc\checkmarki\checkmarkr\checkmarkc\checkmark$\checkmark \checkmark$\checkmark^\checkmark\\checkmarkc\checkmarki\checkmarkr\checkmarkc\checkmark$\checkmark&\checkmark$\checkmark^\checkmark\\checkmarkc\checkmarki\checkmarkr\checkmarkc\checkmark$\checkmark \checkmark$\checkmark^\checkmark\\checkmarkc\checkmarki\checkmarkr\checkmarkc\checkmark$\checkmark$\checkmark$\checkmark^\checkmark\\checkmarkc\checkmarki\checkmarkr\checkmarkc\checkmark$\checkmark6\checkmark$\checkmark^\checkmark\\checkmarkc\checkmarki\checkmarkr\checkmarkc\checkmark$\checkmark$\checkmark$\checkmark^\checkmark\\checkmarkc\checkmarki\checkmarkr\checkmarkc\checkmark$\checkmark \checkmark$\checkmark^\checkmark\\checkmarkc\checkmarki\checkmarkr\checkmarkc\checkmark$\checkmarkm\checkmark$\checkmark^\checkmark\\checkmarkc\checkmarki\checkmarkr\checkmarkc\checkmark$\checkmarkm\checkmark$\checkmark^\checkmark\\checkmarkc\checkmarki\checkmarkr\checkmarkc\checkmark$\checkmark \checkmark$\checkmark^\checkmark\\checkmarkc\checkmarki\checkmarkr\checkmarkc\checkmark$\checkmark&\checkmark$\checkmark^\checkmark\\checkmarkc\checkmarki\checkmarkr\checkmarkc\checkmark$\checkmark \checkmark$\checkmark^\checkmark\\checkmarkc\checkmarki\checkmarkr\checkmarkc\checkmark$\checkmarkF\checkmark$\checkmark^\checkmark\\checkmarkc\checkmarki\checkmarkr\checkmarkc\checkmark$\checkmarkr\checkmark$\checkmark^\checkmark\\checkmarkc\checkmarki\checkmarkr\checkmarkc\checkmark$\checkmarka\checkmark$\checkmark^\checkmark\\checkmarkc\checkmarki\checkmarkr\checkmarkc\checkmark$\checkmarki\checkmark$\checkmark^\checkmark\\checkmarkc\checkmarki\checkmarkr\checkmarkc\checkmark$\checkmarks\checkmark$\checkmark^\checkmark\\checkmarkc\checkmarki\checkmarkr\checkmarkc\checkmark$\checkmarke\checkmark$\checkmark^\checkmark\\checkmarkc\checkmarki\checkmarkr\checkmarkc\checkmark$\checkmark \checkmark$\checkmark^\checkmark\\checkmarkc\checkmarki\checkmarkr\checkmarkc\checkmark$\checkmark2\checkmark$\checkmark^\checkmark\\checkmarkc\checkmarki\checkmarkr\checkmarkc\checkmark$\checkmark \checkmark$\checkmark^\checkmark\\checkmarkc\checkmarki\checkmarkr\checkmarkc\checkmark$\checkmarkd\checkmark$\checkmark^\checkmark\\checkmarkc\checkmarki\checkmarkr\checkmarkc\checkmark$\checkmarke\checkmark$\checkmark^\checkmark\\checkmarkc\checkmarki\checkmarkr\checkmarkc\checkmark$\checkmarkn\checkmark$\checkmark^\checkmark\\checkmarkc\checkmarki\checkmarkr\checkmarkc\checkmark$\checkmarkt\checkmark$\checkmark^\checkmark\\checkmarkc\checkmarki\checkmarkr\checkmarkc\checkmark$\checkmarks\checkmark$\checkmark^\checkmark\\checkmarkc\checkmarki\checkmarkr\checkmarkc\checkmark$\checkmark,\checkmark$\checkmark^\checkmark\\checkmarkc\checkmarki\checkmarkr\checkmarkc\checkmark$\checkmark \checkmark$\checkmark^\checkmark\\checkmarkc\checkmarki\checkmarkr\checkmarkc\checkmark$\checkmarku\checkmark$\checkmark^\checkmark\\checkmarkc\checkmarki\checkmarkr\checkmarkc\checkmark$\checkmarks\checkmark$\checkmark^\checkmark\\checkmarkc\checkmarki\checkmarkr\checkmarkc\checkmark$\checkmarka\checkmark$\checkmark^\checkmark\\checkmarkc\checkmarki\checkmarkr\checkmarkc\checkmark$\checkmarkg\checkmark$\checkmark^\checkmark\\checkmarkc\checkmarki\checkmarkr\checkmarkc\checkmark$\checkmarke\checkmark$\checkmark^\checkmark\\checkmarkc\checkmarki\checkmarkr\checkmarkc\checkmark$\checkmark \checkmark$\checkmark^\checkmark\\checkmarkc\checkmarki\checkmarkr\checkmarkc\checkmark$\checkmarkc\checkmark$\checkmark^\checkmark\\checkmarkc\checkmarki\checkmarkr\checkmarkc\checkmark$\checkmarko\checkmark$\checkmark^\checkmark\\checkmarkc\checkmarki\checkmarkr\checkmarkc\checkmark$\checkmarku\checkmark$\checkmark^\checkmark\\checkmarkc\checkmarki\checkmarkr\checkmarkc\checkmark$\checkmarkr\checkmark$\checkmark^\checkmark\\checkmarkc\checkmarki\checkmarkr\checkmarkc\checkmark$\checkmarka\checkmark$\checkmark^\checkmark\\checkmarkc\checkmarki\checkmarkr\checkmarkc\checkmark$\checkmarkn\checkmark$\checkmark^\checkmark\\checkmarkc\checkmarki\checkmarkr\checkmarkc\checkmark$\checkmarkt\checkmark$\checkmark^\checkmark\\checkmarkc\checkmarki\checkmarkr\checkmarkc\checkmark$\checkmark \checkmark$\checkmark^\checkmark\\checkmarkc\checkmarki\checkmarkr\checkmarkc\checkmark$\checkmark\\checkmark$\checkmark^\checkmark\\checkmarkc\checkmarki\checkmarkr\checkmarkc\checkmark$\checkmark\\checkmark$\checkmark^\checkmark\\checkmarkc\checkmarki\checkmarkr\checkmarkc\checkmark$\checkmark \checkmark$\checkmark^\checkmark\\checkmarkc\checkmarki\checkmarkr\checkmarkc\checkmark$\checkmark\\checkmark$\checkmark^\checkmark\\checkmarkc\checkmarki\checkmarkr\checkmarkc\checkmark$\checkmarkh\checkmark$\checkmark^\checkmark\\checkmarkc\checkmarki\checkmarkr\checkmarkc\checkmark$\checkmarkl\checkmark$\checkmark^\checkmark\\checkmarkc\checkmarki\checkmarkr\checkmarkc\checkmark$\checkmarki\checkmark$\checkmark^\checkmark\\checkmarkc\checkmarki\checkmarkr\checkmarkc\checkmark$\checkmarkn\checkmark$\checkmark^\checkmark\\checkmarkc\checkmarki\checkmarkr\checkmarkc\checkmark$\checkmarke\checkmark$\checkmark^\checkmark\\checkmarkc\checkmarki\checkmarkr\checkmarkc\checkmark$\checkmark
\checkmark$\checkmark^\checkmark\\checkmarkc\checkmarki\checkmarkr\checkmarkc\checkmark$\checkmarkN\checkmark$\checkmark^\checkmark\\checkmarkc\checkmarki\checkmarkr\checkmarkc\checkmark$\checkmarko\checkmark$\checkmark^\checkmark\\checkmarkc\checkmarki\checkmarkr\checkmarkc\checkmark$\checkmarkm\checkmark$\checkmark^\checkmark\\checkmarkc\checkmarki\checkmarkr\checkmarkc\checkmark$\checkmarkb\checkmark$\checkmark^\checkmark\\checkmarkc\checkmarki\checkmarkr\checkmarkc\checkmark$\checkmarkr\checkmark$\checkmark^\checkmark\\checkmarkc\checkmarki\checkmarkr\checkmarkc\checkmark$\checkmarke\checkmark$\checkmark^\checkmark\\checkmarkc\checkmarki\checkmarkr\checkmarkc\checkmark$\checkmark \checkmark$\checkmark^\checkmark\\checkmarkc\checkmarki\checkmarkr\checkmarkc\checkmark$\checkmarkd\checkmark$\checkmark^\checkmark\\checkmarkc\checkmarki\checkmarkr\checkmarkc\checkmark$\checkmarke\checkmark$\checkmark^\checkmark\\checkmarkc\checkmarki\checkmarkr\checkmarkc\checkmark$\checkmark \checkmark$\checkmark^\checkmark\\checkmarkc\checkmarki\checkmarkr\checkmarkc\checkmark$\checkmarkd\checkmark$\checkmark^\checkmark\\checkmarkc\checkmarki\checkmarkr\checkmarkc\checkmark$\checkmarke\checkmark$\checkmark^\checkmark\\checkmarkc\checkmarki\checkmarkr\checkmarkc\checkmark$\checkmarkn\checkmark$\checkmark^\checkmark\\checkmarkc\checkmarki\checkmarkr\checkmarkc\checkmark$\checkmarkt\checkmark$\checkmark^\checkmark\\checkmarkc\checkmarki\checkmarkr\checkmarkc\checkmark$\checkmarks\checkmark$\checkmark^\checkmark\\checkmarkc\checkmarki\checkmarkr\checkmarkc\checkmark$\checkmark \checkmark$\checkmark^\checkmark\\checkmarkc\checkmarki\checkmarkr\checkmarkc\checkmark$\checkmark(\checkmark$\checkmark^\checkmark\\checkmarkc\checkmarki\checkmarkr\checkmarkc\checkmark$\checkmark$\checkmark$\checkmark^\checkmark\\checkmarkc\checkmarki\checkmarkr\checkmarkc\checkmark$\checkmarkZ\checkmark$\checkmark^\checkmark\\checkmarkc\checkmarki\checkmarkr\checkmarkc\checkmark$\checkmark_\checkmark$\checkmark^\checkmark\\checkmarkc\checkmarki\checkmarkr\checkmarkc\checkmark$\checkmarkc\checkmark$\checkmark^\checkmark\\checkmarkc\checkmarki\checkmarkr\checkmarkc\checkmark$\checkmark$\checkmark$\checkmark^\checkmark\\checkmarkc\checkmarki\checkmarkr\checkmarkc\checkmark$\checkmark)\checkmark$\checkmark^\checkmark\\checkmarkc\checkmarki\checkmarkr\checkmarkc\checkmark$\checkmark \checkmark$\checkmark^\checkmark\\checkmarkc\checkmarki\checkmarkr\checkmarkc\checkmark$\checkmark&\checkmark$\checkmark^\checkmark\\checkmarkc\checkmarki\checkmarkr\checkmarkc\checkmark$\checkmark \checkmark$\checkmark^\checkmark\\checkmarkc\checkmarki\checkmarkr\checkmarkc\checkmark$\checkmark$\checkmark$\checkmark^\checkmark\\checkmarkc\checkmarki\checkmarkr\checkmarkc\checkmark$\checkmark2\checkmark$\checkmark^\checkmark\\checkmarkc\checkmarki\checkmarkr\checkmarkc\checkmark$\checkmark$\checkmark$\checkmark^\checkmark\\checkmarkc\checkmarki\checkmarkr\checkmarkc\checkmark$\checkmark \checkmark$\checkmark^\checkmark\\checkmarkc\checkmarki\checkmarkr\checkmarkc\checkmark$\checkmark&\checkmark$\checkmark^\checkmark\\checkmarkc\checkmarki\checkmarkr\checkmarkc\checkmark$\checkmark \checkmark$\checkmark^\checkmark\\checkmarkc\checkmarki\checkmarkr\checkmarkc\checkmark$\checkmarkF\checkmark$\checkmark^\checkmark\\checkmarkc\checkmarki\checkmarkr\checkmarkc\checkmark$\checkmarkr\checkmark$\checkmark^\checkmark\\checkmarkc\checkmarki\checkmarkr\checkmarkc\checkmark$\checkmarka\checkmark$\checkmark^\checkmark\\checkmarkc\checkmarki\checkmarkr\checkmarkc\checkmark$\checkmarki\checkmark$\checkmark^\checkmark\\checkmarkc\checkmarki\checkmarkr\checkmarkc\checkmark$\checkmarks\checkmark$\checkmark^\checkmark\\checkmarkc\checkmarki\checkmarkr\checkmarkc\checkmark$\checkmarke\checkmark$\checkmark^\checkmark\\checkmarkc\checkmarki\checkmarkr\checkmarkc\checkmark$\checkmark \checkmark$\checkmark^\checkmark\\checkmarkc\checkmarki\checkmarkr\checkmarkc\checkmark$\checkmarks\checkmark$\checkmark^\checkmark\\checkmarkc\checkmarki\checkmarkr\checkmarkc\checkmark$\checkmarkt\checkmark$\checkmark^\checkmark\\checkmarkc\checkmarki\checkmarkr\checkmarkc\checkmark$\checkmarka\checkmark$\checkmark^\checkmark\\checkmarkc\checkmarki\checkmarkr\checkmarkc\checkmark$\checkmarkn\checkmark$\checkmark^\checkmark\\checkmarkc\checkmarki\checkmarkr\checkmarkc\checkmark$\checkmarkd\checkmark$\checkmark^\checkmark\\checkmarkc\checkmarki\checkmarkr\checkmarkc\checkmark$\checkmarka\checkmark$\checkmark^\checkmark\\checkmarkc\checkmarki\checkmarkr\checkmarkc\checkmark$\checkmarkr\checkmark$\checkmark^\checkmark\\checkmarkc\checkmarki\checkmarkr\checkmarkc\checkmark$\checkmarkd\checkmark$\checkmark^\checkmark\\checkmarkc\checkmarki\checkmarkr\checkmarkc\checkmark$\checkmark \checkmark$\checkmark^\checkmark\\checkmarkc\checkmarki\checkmarkr\checkmarkc\checkmark$\checkmarka\checkmark$\checkmark^\checkmark\\checkmarkc\checkmarki\checkmarkr\checkmarkc\checkmark$\checkmarkl\checkmark$\checkmark^\checkmark\\checkmarkc\checkmarki\checkmarkr\checkmarkc\checkmark$\checkmarku\checkmark$\checkmark^\checkmark\\checkmarkc\checkmarki\checkmarkr\checkmarkc\checkmark$\checkmarkm\checkmark$\checkmark^\checkmark\\checkmarkc\checkmarki\checkmarkr\checkmarkc\checkmark$\checkmarki\checkmark$\checkmark^\checkmark\\checkmarkc\checkmarki\checkmarkr\checkmarkc\checkmark$\checkmarkn\checkmark$\checkmark^\checkmark\\checkmarkc\checkmarki\checkmarkr\checkmarkc\checkmark$\checkmarki\checkmark$\checkmark^\checkmark\\checkmarkc\checkmarki\checkmarkr\checkmarkc\checkmark$\checkmarku\checkmark$\checkmark^\checkmark\\checkmarkc\checkmarki\checkmarkr\checkmarkc\checkmark$\checkmarkm\checkmark$\checkmark^\checkmark\\checkmarkc\checkmarki\checkmarkr\checkmarkc\checkmark$\checkmark \checkmark$\checkmark^\checkmark\\checkmarkc\checkmarki\checkmarkr\checkmarkc\checkmark$\checkmark\\checkmark$\checkmark^\checkmark\\checkmarkc\checkmarki\checkmarkr\checkmarkc\checkmark$\checkmark\\checkmark$\checkmark^\checkmark\\checkmarkc\checkmarki\checkmarkr\checkmarkc\checkmark$\checkmark \checkmark$\checkmark^\checkmark\\checkmarkc\checkmarki\checkmarkr\checkmarkc\checkmark$\checkmark\\checkmark$\checkmark^\checkmark\\checkmarkc\checkmarki\checkmarkr\checkmarkc\checkmark$\checkmarkh\checkmark$\checkmark^\checkmark\\checkmarkc\checkmarki\checkmarkr\checkmarkc\checkmark$\checkmarkl\checkmark$\checkmark^\checkmark\\checkmarkc\checkmarki\checkmarkr\checkmarkc\checkmark$\checkmarki\checkmark$\checkmark^\checkmark\\checkmarkc\checkmarki\checkmarkr\checkmarkc\checkmark$\checkmarkn\checkmark$\checkmark^\checkmark\\checkmarkc\checkmarki\checkmarkr\checkmarkc\checkmark$\checkmarke\checkmark$\checkmark^\checkmark\\checkmarkc\checkmarki\checkmarkr\checkmarkc\checkmark$\checkmark
\checkmark$\checkmark^\checkmark\\checkmarkc\checkmarki\checkmarkr\checkmarkc\checkmark$\checkmarkP\checkmark$\checkmark^\checkmark\\checkmarkc\checkmarki\checkmarkr\checkmarkc\checkmark$\checkmarkr\checkmark$\checkmark^\checkmark\\checkmarkc\checkmarki\checkmarkr\checkmarkc\checkmark$\checkmarko\checkmark$\checkmark^\checkmark\\checkmarkc\checkmarki\checkmarkr\checkmarkc\checkmark$\checkmarkf\checkmark$\checkmark^\checkmark\\checkmarkc\checkmarki\checkmarkr\checkmarkc\checkmark$\checkmarko\checkmark$\checkmark^\checkmark\\checkmarkc\checkmarki\checkmarkr\checkmarkc\checkmark$\checkmarkn\checkmark$\checkmark^\checkmark\\checkmarkc\checkmarki\checkmarkr\checkmarkc\checkmark$\checkmarkd\checkmark$\checkmark^\checkmark\\checkmarkc\checkmarki\checkmarkr\checkmarkc\checkmark$\checkmarke\checkmark$\checkmark^\checkmark\\checkmarkc\checkmarki\checkmarkr\checkmarkc\checkmark$\checkmarku\checkmark$\checkmark^\checkmark\\checkmarkc\checkmarki\checkmarkr\checkmarkc\checkmark$\checkmarkr\checkmark$\checkmark^\checkmark\\checkmarkc\checkmarki\checkmarkr\checkmarkc\checkmark$\checkmark \checkmark$\checkmark^\checkmark\\checkmarkc\checkmarki\checkmarkr\checkmarkc\checkmark$\checkmarkr\checkmark$\checkmark^\checkmark\\checkmarkc\checkmarki\checkmarkr\checkmarkc\checkmark$\checkmarka\checkmark$\checkmark^\checkmark\\checkmarkc\checkmarki\checkmarkr\checkmarkc\checkmark$\checkmarkd\checkmark$\checkmark^\checkmark\\checkmarkc\checkmarki\checkmarkr\checkmarkc\checkmark$\checkmarki\checkmark$\checkmark^\checkmark\\checkmarkc\checkmarki\checkmarkr\checkmarkc\checkmark$\checkmarka\checkmark$\checkmark^\checkmark\\checkmarkc\checkmarki\checkmarkr\checkmarkc\checkmark$\checkmarkl\checkmark$\checkmark^\checkmark\\checkmarkc\checkmarki\checkmarkr\checkmarkc\checkmark$\checkmarke\checkmark$\checkmark^\checkmark\\checkmarkc\checkmarki\checkmarkr\checkmarkc\checkmark$\checkmark \checkmark$\checkmark^\checkmark\\checkmarkc\checkmarki\checkmarkr\checkmarkc\checkmark$\checkmark(\checkmark$\checkmark^\checkmark\\checkmarkc\checkmarki\checkmarkr\checkmarkc\checkmark$\checkmark$\checkmark$\checkmark^\checkmark\\checkmarkc\checkmarki\checkmarkr\checkmarkc\checkmark$\checkmarka\checkmark$\checkmark^\checkmark\\checkmarkc\checkmarki\checkmarkr\checkmarkc\checkmark$\checkmark_\checkmark$\checkmark^\checkmark\\checkmarkc\checkmarki\checkmarkr\checkmarkc\checkmark$\checkmarke\checkmark$\checkmark^\checkmark\\checkmarkc\checkmarki\checkmarkr\checkmarkc\checkmark$\checkmark$\checkmark$\checkmark^\checkmark\\checkmarkc\checkmarki\checkmarkr\checkmarkc\checkmark$\checkmark)\checkmark$\checkmark^\checkmark\\checkmarkc\checkmarki\checkmarkr\checkmarkc\checkmark$\checkmark \checkmark$\checkmark^\checkmark\\checkmarkc\checkmarki\checkmarkr\checkmarkc\checkmark$\checkmark&\checkmark$\checkmark^\checkmark\\checkmarkc\checkmarki\checkmarkr\checkmarkc\checkmark$\checkmark \checkmark$\checkmark^\checkmark\\checkmarkc\checkmarki\checkmarkr\checkmarkc\checkmark$\checkmark$\checkmark$\checkmark^\checkmark\\checkmarkc\checkmarki\checkmarkr\checkmarkc\checkmark$\checkmark6\checkmark$\checkmark^\checkmark\\checkmarkc\checkmarki\checkmarkr\checkmarkc\checkmark$\checkmark$\checkmark$\checkmark^\checkmark\\checkmarkc\checkmarki\checkmarkr\checkmarkc\checkmark$\checkmark \checkmark$\checkmark^\checkmark\\checkmarkc\checkmarki\checkmarkr\checkmarkc\checkmark$\checkmarkm\checkmark$\checkmark^\checkmark\\checkmarkc\checkmarki\checkmarkr\checkmarkc\checkmark$\checkmarkm\checkmark$\checkmark^\checkmark\\checkmarkc\checkmarki\checkmarkr\checkmarkc\checkmark$\checkmark \checkmark$\checkmark^\checkmark\\checkmarkc\checkmarki\checkmarkr\checkmarkc\checkmark$\checkmark&\checkmark$\checkmark^\checkmark\\checkmarkc\checkmarki\checkmarkr\checkmarkc\checkmark$\checkmark \checkmark$\checkmark^\checkmark\\checkmarkc\checkmarki\checkmarkr\checkmarkc\checkmark$\checkmarkP\checkmark$\checkmark^\checkmark\\checkmarkc\checkmarki\checkmarkr\checkmarkc\checkmark$\checkmarkl\checkmark$\checkmark^\checkmark\\checkmarkc\checkmarki\checkmarkr\checkmarkc\checkmark$\checkmarke\checkmark$\checkmark^\checkmark\\checkmarkc\checkmarki\checkmarkr\checkmarkc\checkmark$\checkmarki\checkmark$\checkmark^\checkmark\\checkmarkc\checkmarki\checkmarkr\checkmarkc\checkmark$\checkmarkn\checkmark$\checkmark^\checkmark\\checkmarkc\checkmarki\checkmarkr\checkmarkc\checkmark$\checkmark \checkmark$\checkmark^\checkmark\\checkmarkc\checkmarki\checkmarkr\checkmarkc\checkmark$\checkmarkb\checkmark$\checkmark^\checkmark\\checkmarkc\checkmarki\checkmarkr\checkmarkc\checkmark$\checkmarko\checkmark$\checkmark^\checkmark\\checkmarkc\checkmarki\checkmarkr\checkmarkc\checkmark$\checkmarkr\checkmark$\checkmark^\checkmark\\checkmarkc\checkmarki\checkmarkr\checkmarkc\checkmark$\checkmarkd\checkmark$\checkmark^\checkmark\\checkmarkc\checkmarki\checkmarkr\checkmarkc\checkmark$\checkmark \checkmark$\checkmark^\checkmark\\checkmarkc\checkmarki\checkmarkr\checkmarkc\checkmark$\checkmark(\checkmark$\checkmark^\checkmark\\checkmarkc\checkmarki\checkmarkr\checkmarkc\checkmark$\checkmarkp\checkmark$\checkmark^\checkmark\\checkmarkc\checkmarki\checkmarkr\checkmarkc\checkmark$\checkmarka\checkmark$\checkmark^\checkmark\\checkmarkc\checkmarki\checkmarkr\checkmarkc\checkmark$\checkmarks\checkmark$\checkmark^\checkmark\\checkmarkc\checkmarki\checkmarkr\checkmarkc\checkmark$\checkmarks\checkmark$\checkmark^\checkmark\\checkmarkc\checkmarki\checkmarkr\checkmarkc\checkmark$\checkmarke\checkmark$\checkmark^\checkmark\\checkmarkc\checkmarki\checkmarkr\checkmarkc\checkmark$\checkmark \checkmark$\checkmark^\checkmark\\checkmarkc\checkmarki\checkmarkr\checkmarkc\checkmark$\checkmarkm\checkmark$\checkmark^\checkmark\\checkmarkc\checkmarki\checkmarkr\checkmarkc\checkmark$\checkmarka\checkmark$\checkmark^\checkmark\\checkmarkc\checkmarki\checkmarkr\checkmarkc\checkmark$\checkmarkx\checkmark$\checkmark^\checkmark\\checkmarkc\checkmarki\checkmarkr\checkmarkc\checkmark$\checkmarki\checkmark$\checkmark^\checkmark\\checkmarkc\checkmarki\checkmarkr\checkmarkc\checkmark$\checkmarkm\checkmark$\checkmark^\checkmark\\checkmarkc\checkmarki\checkmarkr\checkmarkc\checkmark$\checkmarka\checkmark$\checkmark^\checkmark\\checkmarkc\checkmarki\checkmarkr\checkmarkc\checkmark$\checkmarkl\checkmark$\checkmark^\checkmark\\checkmarkc\checkmarki\checkmarkr\checkmarkc\checkmark$\checkmarke\checkmark$\checkmark^\checkmark\\checkmarkc\checkmarki\checkmarkr\checkmarkc\checkmark$\checkmark)\checkmark$\checkmark^\checkmark\\checkmarkc\checkmarki\checkmarkr\checkmarkc\checkmark$\checkmark \checkmark$\checkmark^\checkmark\\checkmarkc\checkmarki\checkmarkr\checkmarkc\checkmark$\checkmark\\checkmark$\checkmark^\checkmark\\checkmarkc\checkmarki\checkmarkr\checkmarkc\checkmark$\checkmark\\checkmark$\checkmark^\checkmark\\checkmarkc\checkmarki\checkmarkr\checkmarkc\checkmark$\checkmark \checkmark$\checkmark^\checkmark\\checkmarkc\checkmarki\checkmarkr\checkmarkc\checkmark$\checkmark\\checkmark$\checkmark^\checkmark\\checkmarkc\checkmarki\checkmarkr\checkmarkc\checkmark$\checkmarkh\checkmark$\checkmark^\checkmark\\checkmarkc\checkmarki\checkmarkr\checkmarkc\checkmark$\checkmarkl\checkmark$\checkmark^\checkmark\\checkmarkc\checkmarki\checkmarkr\checkmarkc\checkmark$\checkmarki\checkmark$\checkmark^\checkmark\\checkmarkc\checkmarki\checkmarkr\checkmarkc\checkmark$\checkmarkn\checkmark$\checkmark^\checkmark\\checkmarkc\checkmarki\checkmarkr\checkmarkc\checkmark$\checkmarke\checkmark$\checkmark^\checkmark\\checkmarkc\checkmarki\checkmarkr\checkmarkc\checkmark$\checkmark
\checkmark$\checkmark^\checkmark\\checkmarkc\checkmarki\checkmarkr\checkmarkc\checkmark$\checkmarkP\checkmark$\checkmark^\checkmark\\checkmarkc\checkmarki\checkmarkr\checkmarkc\checkmark$\checkmarkr\checkmark$\checkmark^\checkmark\\checkmarkc\checkmarki\checkmarkr\checkmarkc\checkmark$\checkmarko\checkmark$\checkmark^\checkmark\\checkmarkc\checkmarki\checkmarkr\checkmarkc\checkmark$\checkmarkf\checkmark$\checkmark^\checkmark\\checkmarkc\checkmarki\checkmarkr\checkmarkc\checkmark$\checkmarko\checkmark$\checkmark^\checkmark\\checkmarkc\checkmarki\checkmarkr\checkmarkc\checkmark$\checkmarkn\checkmark$\checkmark^\checkmark\\checkmarkc\checkmarki\checkmarkr\checkmarkc\checkmark$\checkmarkd\checkmark$\checkmark^\checkmark\\checkmarkc\checkmarki\checkmarkr\checkmarkc\checkmark$\checkmarke\checkmark$\checkmark^\checkmark\\checkmarkc\checkmarki\checkmarkr\checkmarkc\checkmark$\checkmarku\checkmark$\checkmark^\checkmark\\checkmarkc\checkmarki\checkmarkr\checkmarkc\checkmark$\checkmarkr\checkmark$\checkmark^\checkmark\\checkmarkc\checkmarki\checkmarkr\checkmarkc\checkmark$\checkmark \checkmark$\checkmark^\checkmark\\checkmarkc\checkmarki\checkmarkr\checkmarkc\checkmark$\checkmarka\checkmark$\checkmark^\checkmark\\checkmarkc\checkmarki\checkmarkr\checkmarkc\checkmark$\checkmarkx\checkmark$\checkmark^\checkmark\\checkmarkc\checkmarki\checkmarkr\checkmarkc\checkmark$\checkmarki\checkmark$\checkmark^\checkmark\\checkmarkc\checkmarki\checkmarkr\checkmarkc\checkmark$\checkmarka\checkmark$\checkmark^\checkmark\\checkmarkc\checkmarki\checkmarkr\checkmarkc\checkmark$\checkmarkl\checkmark$\checkmark^\checkmark\\checkmarkc\checkmarki\checkmarkr\checkmarkc\checkmark$\checkmarke\checkmark$\checkmark^\checkmark\\checkmarkc\checkmarki\checkmarkr\checkmarkc\checkmark$\checkmark \checkmark$\checkmark^\checkmark\\checkmarkc\checkmarki\checkmarkr\checkmarkc\checkmark$\checkmark(\checkmark$\checkmark^\checkmark\\checkmarkc\checkmarki\checkmarkr\checkmarkc\checkmark$\checkmark$\checkmark$\checkmark^\checkmark\\checkmarkc\checkmarki\checkmarkr\checkmarkc\checkmark$\checkmarka\checkmark$\checkmark^\checkmark\\checkmarkc\checkmarki\checkmarkr\checkmarkc\checkmark$\checkmark_\checkmark$\checkmark^\checkmark\\checkmarkc\checkmarki\checkmarkr\checkmarkc\checkmark$\checkmarkp\checkmark$\checkmark^\checkmark\\checkmarkc\checkmarki\checkmarkr\checkmarkc\checkmark$\checkmark$\checkmark$\checkmark^\checkmark\\checkmarkc\checkmarki\checkmarkr\checkmarkc\checkmark$\checkmark)\checkmark$\checkmark^\checkmark\\checkmarkc\checkmarki\checkmarkr\checkmarkc\checkmark$\checkmark \checkmark$\checkmark^\checkmark\\checkmarkc\checkmarki\checkmarkr\checkmarkc\checkmark$\checkmark&\checkmark$\checkmark^\checkmark\\checkmarkc\checkmarki\checkmarkr\checkmarkc\checkmark$\checkmark \checkmark$\checkmark^\checkmark\\checkmarkc\checkmarki\checkmarkr\checkmarkc\checkmark$\checkmark$\checkmark$\checkmark^\checkmark\\checkmarkc\checkmarki\checkmarkr\checkmarkc\checkmark$\checkmark1\checkmark$\checkmark^\checkmark\\checkmarkc\checkmarki\checkmarkr\checkmarkc\checkmark$\checkmark$\checkmark$\checkmark^\checkmark\\checkmarkc\checkmarki\checkmarkr\checkmarkc\checkmark$\checkmark \checkmark$\checkmark^\checkmark\\checkmarkc\checkmarki\checkmarkr\checkmarkc\checkmark$\checkmarkm\checkmark$\checkmark^\checkmark\\checkmarkc\checkmarki\checkmarkr\checkmarkc\checkmark$\checkmarkm\checkmark$\checkmark^\checkmark\\checkmarkc\checkmarki\checkmarkr\checkmarkc\checkmark$\checkmark \checkmark$\checkmark^\checkmark\\checkmarkc\checkmarki\checkmarkr\checkmarkc\checkmark$\checkmark&\checkmark$\checkmark^\checkmark\\checkmarkc\checkmarki\checkmarkr\checkmarkc\checkmark$\checkmark \checkmark$\checkmark^\checkmark\\checkmarkc\checkmarki\checkmarkr\checkmarkc\checkmark$\checkmarkP\checkmark$\checkmark^\checkmark\\checkmarkc\checkmarki\checkmarkr\checkmarkc\checkmark$\checkmarka\checkmark$\checkmark^\checkmark\\checkmarkc\checkmarki\checkmarkr\checkmarkc\checkmark$\checkmarks\checkmark$\checkmark^\checkmark\\checkmarkc\checkmarki\checkmarkr\checkmarkc\checkmark$\checkmarks\checkmark$\checkmark^\checkmark\\checkmarkc\checkmarki\checkmarkr\checkmarkc\checkmark$\checkmarke\checkmark$\checkmark^\checkmark\\checkmarkc\checkmarki\checkmarkr\checkmarkc\checkmark$\checkmark \checkmark$\checkmark^\checkmark\\checkmarkc\checkmarki\checkmarkr\checkmarkc\checkmark$\checkmarkd\checkmark$\checkmark^\checkmark\\checkmarkc\checkmarki\checkmarkr\checkmarkc\checkmark$\checkmarke\checkmark$\checkmark^\checkmark\\checkmarkc\checkmarki\checkmarkr\checkmarkc\checkmark$\checkmark \checkmark$\checkmark^\checkmark\\checkmarkc\checkmarki\checkmarkr\checkmarkc\checkmark$\checkmarkf\checkmark$\checkmark^\checkmark\\checkmarkc\checkmarki\checkmarkr\checkmarkc\checkmark$\checkmarki\checkmark$\checkmark^\checkmark\\checkmarkc\checkmarki\checkmarkr\checkmarkc\checkmark$\checkmarkn\checkmark$\checkmark^\checkmark\\checkmarkc\checkmarki\checkmarkr\checkmarkc\checkmark$\checkmarki\checkmark$\checkmark^\checkmark\\checkmarkc\checkmarki\checkmarkr\checkmarkc\checkmark$\checkmarkt\checkmark$\checkmark^\checkmark\\checkmarkc\checkmarki\checkmarkr\checkmarkc\checkmark$\checkmarki\checkmark$\checkmark^\checkmark\\checkmarkc\checkmarki\checkmarkr\checkmarkc\checkmark$\checkmarko\checkmark$\checkmark^\checkmark\\checkmarkc\checkmarki\checkmarkr\checkmarkc\checkmark$\checkmarkn\checkmark$\checkmark^\checkmark\\checkmarkc\checkmarki\checkmarkr\checkmarkc\checkmark$\checkmark/\checkmark$\checkmark^\checkmark\\checkmarkc\checkmarki\checkmarkr\checkmarkc\checkmark$\checkmarkd\checkmark$\checkmark^\checkmark\\checkmarkc\checkmarki\checkmarkr\checkmarkc\checkmark$\checkmarké\checkmark$\checkmark^\checkmark\\checkmarkc\checkmarki\checkmarkr\checkmarkc\checkmark$\checkmarkg\checkmark$\checkmark^\checkmark\\checkmarkc\checkmarki\checkmarkr\checkmarkc\checkmark$\checkmarkr\checkmark$\checkmark^\checkmark\\checkmarkc\checkmarki\checkmarkr\checkmarkc\checkmark$\checkmarko\checkmark$\checkmark^\checkmark\\checkmarkc\checkmarki\checkmarkr\checkmarkc\checkmark$\checkmarks\checkmark$\checkmark^\checkmark\\checkmarkc\checkmarki\checkmarkr\checkmarkc\checkmark$\checkmarks\checkmark$\checkmark^\checkmark\\checkmarkc\checkmarki\checkmarkr\checkmarkc\checkmark$\checkmarki\checkmark$\checkmark^\checkmark\\checkmarkc\checkmarki\checkmarkr\checkmarkc\checkmark$\checkmarks\checkmark$\checkmark^\checkmark\\checkmarkc\checkmarki\checkmarkr\checkmarkc\checkmark$\checkmarks\checkmark$\checkmark^\checkmark\\checkmarkc\checkmarki\checkmarkr\checkmarkc\checkmark$\checkmarka\checkmark$\checkmark^\checkmark\\checkmarkc\checkmarki\checkmarkr\checkmarkc\checkmark$\checkmarkg\checkmark$\checkmark^\checkmark\\checkmarkc\checkmarki\checkmarkr\checkmarkc\checkmark$\checkmarke\checkmark$\checkmark^\checkmark\\checkmarkc\checkmarki\checkmarkr\checkmarkc\checkmark$\checkmark \checkmark$\checkmark^\checkmark\\checkmarkc\checkmarki\checkmarkr\checkmarkc\checkmark$\checkmark\\checkmark$\checkmark^\checkmark\\checkmarkc\checkmarki\checkmarkr\checkmarkc\checkmark$\checkmark\\checkmark$\checkmark^\checkmark\\checkmarkc\checkmarki\checkmarkr\checkmarkc\checkmark$\checkmark \checkmark$\checkmark^\checkmark\\checkmarkc\checkmarki\checkmarkr\checkmarkc\checkmark$\checkmark\\checkmark$\checkmark^\checkmark\\checkmarkc\checkmarki\checkmarkr\checkmarkc\checkmark$\checkmarkh\checkmark$\checkmark^\checkmark\\checkmarkc\checkmarki\checkmarkr\checkmarkc\checkmark$\checkmarkl\checkmark$\checkmark^\checkmark\\checkmarkc\checkmarki\checkmarkr\checkmarkc\checkmark$\checkmarki\checkmark$\checkmark^\checkmark\\checkmarkc\checkmarki\checkmarkr\checkmarkc\checkmark$\checkmarkn\checkmark$\checkmark^\checkmark\\checkmarkc\checkmarki\checkmarkr\checkmarkc\checkmark$\checkmarke\checkmark$\checkmark^\checkmark\\checkmarkc\checkmarki\checkmarkr\checkmarkc\checkmark$\checkmark
\checkmark$\checkmark^\checkmark\\checkmarkc\checkmarki\checkmarkr\checkmarkc\checkmark$\checkmarkA\checkmark$\checkmark^\checkmark\\checkmarkc\checkmarki\checkmarkr\checkmarkc\checkmark$\checkmarkv\checkmark$\checkmark^\checkmark\\checkmarkc\checkmarki\checkmarkr\checkmarkc\checkmark$\checkmarka\checkmark$\checkmark^\checkmark\\checkmarkc\checkmarki\checkmarkr\checkmarkc\checkmark$\checkmarkn\checkmark$\checkmark^\checkmark\\checkmarkc\checkmarki\checkmarkr\checkmarkc\checkmark$\checkmarkc\checkmark$\checkmark^\checkmark\\checkmarkc\checkmarki\checkmarkr\checkmarkc\checkmark$\checkmarke\checkmark$\checkmark^\checkmark\\checkmarkc\checkmarki\checkmarkr\checkmarkc\checkmark$\checkmark \checkmark$\checkmark^\checkmark\\checkmarkc\checkmarki\checkmarkr\checkmarkc\checkmark$\checkmarkp\checkmark$\checkmark^\checkmark\\checkmarkc\checkmarki\checkmarkr\checkmarkc\checkmark$\checkmarka\checkmark$\checkmark^\checkmark\\checkmarkc\checkmarki\checkmarkr\checkmarkc\checkmark$\checkmarkr\checkmark$\checkmark^\checkmark\\checkmarkc\checkmarki\checkmarkr\checkmarkc\checkmark$\checkmark \checkmark$\checkmark^\checkmark\\checkmarkc\checkmarki\checkmarkr\checkmarkc\checkmark$\checkmarkd\checkmark$\checkmark^\checkmark\\checkmarkc\checkmarki\checkmarkr\checkmarkc\checkmark$\checkmarke\checkmark$\checkmark^\checkmark\\checkmarkc\checkmarki\checkmarkr\checkmarkc\checkmark$\checkmarkn\checkmark$\checkmark^\checkmark\\checkmarkc\checkmarki\checkmarkr\checkmarkc\checkmark$\checkmarkt\checkmark$\checkmark^\checkmark\\checkmarkc\checkmarki\checkmarkr\checkmarkc\checkmark$\checkmark \checkmark$\checkmark^\checkmark\\checkmarkc\checkmarki\checkmarkr\checkmarkc\checkmark$\checkmark(\checkmark$\checkmark^\checkmark\\checkmarkc\checkmarki\checkmarkr\checkmarkc\checkmark$\checkmark$\checkmark$\checkmark^\checkmark\\checkmarkc\checkmarki\checkmarkr\checkmarkc\checkmark$\checkmarkf\checkmark$\checkmark^\checkmark\\checkmarkc\checkmarki\checkmarkr\checkmarkc\checkmark$\checkmark_\checkmark$\checkmark^\checkmark\\checkmarkc\checkmarki\checkmarkr\checkmarkc\checkmark$\checkmarkz\checkmark$\checkmark^\checkmark\\checkmarkc\checkmarki\checkmarkr\checkmarkc\checkmark$\checkmark$\checkmark$\checkmark^\checkmark\\checkmarkc\checkmarki\checkmarkr\checkmarkc\checkmark$\checkmark)\checkmark$\checkmark^\checkmark\\checkmarkc\checkmarki\checkmarkr\checkmarkc\checkmark$\checkmark \checkmark$\checkmark^\checkmark\\checkmarkc\checkmarki\checkmarkr\checkmarkc\checkmark$\checkmark&\checkmark$\checkmark^\checkmark\\checkmarkc\checkmarki\checkmarkr\checkmarkc\checkmark$\checkmark \checkmark$\checkmark^\checkmark\\checkmarkc\checkmarki\checkmarkr\checkmarkc\checkmark$\checkmark$\checkmark$\checkmark^\checkmark\\checkmarkc\checkmarki\checkmarkr\checkmarkc\checkmark$\checkmark0\checkmark$\checkmark^\checkmark\\checkmarkc\checkmarki\checkmarkr\checkmarkc\checkmark$\checkmark{\checkmark$\checkmark^\checkmark\\checkmarkc\checkmarki\checkmarkr\checkmarkc\checkmark$\checkmark,\checkmark$\checkmark^\checkmark\\checkmarkc\checkmarki\checkmarkr\checkmarkc\checkmark$\checkmark}\checkmark$\checkmark^\checkmark\\checkmarkc\checkmarki\checkmarkr\checkmarkc\checkmark$\checkmark0\checkmark$\checkmark^\checkmark\\checkmarkc\checkmarki\checkmarkr\checkmarkc\checkmark$\checkmark5\checkmark$\checkmark^\checkmark\\checkmarkc\checkmarki\checkmarkr\checkmarkc\checkmark$\checkmark$\checkmark$\checkmark^\checkmark\\checkmarkc\checkmarki\checkmarkr\checkmarkc\checkmark$\checkmark \checkmark$\checkmark^\checkmark\\checkmarkc\checkmarki\checkmarkr\checkmarkc\checkmark$\checkmarkm\checkmark$\checkmark^\checkmark\\checkmarkc\checkmarki\checkmarkr\checkmarkc\checkmark$\checkmarkm\checkmark$\checkmark^\checkmark\\checkmarkc\checkmarki\checkmarkr\checkmarkc\checkmark$\checkmark/\checkmark$\checkmark^\checkmark\\checkmarkc\checkmarki\checkmarkr\checkmarkc\checkmark$\checkmarkd\checkmark$\checkmark^\checkmark\\checkmarkc\checkmarki\checkmarkr\checkmarkc\checkmark$\checkmarke\checkmark$\checkmark^\checkmark\\checkmarkc\checkmarki\checkmarkr\checkmarkc\checkmark$\checkmarkn\checkmark$\checkmark^\checkmark\\checkmarkc\checkmarki\checkmarkr\checkmarkc\checkmark$\checkmarkt\checkmark$\checkmark^\checkmark\\checkmarkc\checkmarki\checkmarkr\checkmarkc\checkmark$\checkmark \checkmark$\checkmark^\checkmark\\checkmarkc\checkmarki\checkmarkr\checkmarkc\checkmark$\checkmark&\checkmark$\checkmark^\checkmark\\checkmarkc\checkmarki\checkmarkr\checkmarkc\checkmark$\checkmark \checkmark$\checkmark^\checkmark\\checkmarkc\checkmarki\checkmarkr\checkmarkc\checkmark$\checkmarkS\checkmark$\checkmark^\checkmark\\checkmarkc\checkmarki\checkmarkr\checkmarkc\checkmark$\checkmarkt\checkmark$\checkmark^\checkmark\\checkmarkc\checkmarki\checkmarkr\checkmarkc\checkmark$\checkmarka\checkmark$\checkmark^\checkmark\\checkmarkc\checkmarki\checkmarkr\checkmarkc\checkmark$\checkmarkn\checkmark$\checkmark^\checkmark\\checkmarkc\checkmarki\checkmarkr\checkmarkc\checkmark$\checkmarkd\checkmark$\checkmark^\checkmark\\checkmarkc\checkmarki\checkmarkr\checkmarkc\checkmark$\checkmarka\checkmark$\checkmark^\checkmark\\checkmarkc\checkmarki\checkmarkr\checkmarkc\checkmark$\checkmarkr\checkmark$\checkmark^\checkmark\\checkmarkc\checkmarki\checkmarkr\checkmarkc\checkmark$\checkmarkd\checkmark$\checkmark^\checkmark\\checkmarkc\checkmarki\checkmarkr\checkmarkc\checkmark$\checkmark \checkmark$\checkmark^\checkmark\\checkmarkc\checkmarki\checkmarkr\checkmarkc\checkmark$\checkmarkp\checkmark$\checkmark^\checkmark\\checkmarkc\checkmarki\checkmarkr\checkmarkc\checkmark$\checkmarko\checkmark$\checkmark^\checkmark\\checkmarkc\checkmarki\checkmarkr\checkmarkc\checkmark$\checkmarku\checkmark$\checkmark^\checkmark\\checkmarkc\checkmarki\checkmarkr\checkmarkc\checkmark$\checkmarkr\checkmark$\checkmark^\checkmark\\checkmarkc\checkmarki\checkmarkr\checkmarkc\checkmark$\checkmark \checkmark$\checkmark^\checkmark\\checkmarkc\checkmarki\checkmarkr\checkmarkc\checkmark$\checkmarkA\checkmark$\checkmark^\checkmark\\checkmarkc\checkmarki\checkmarkr\checkmarkc\checkmark$\checkmarkl\checkmark$\checkmark^\checkmark\\checkmarkc\checkmarki\checkmarkr\checkmarkc\checkmark$\checkmark \checkmark$\checkmark^\checkmark\\checkmarkc\checkmarki\checkmarkr\checkmarkc\checkmark$\checkmark\\checkmark$\checkmark^\checkmark\\checkmarkc\checkmarki\checkmarkr\checkmarkc\checkmark$\checkmark\\checkmark$\checkmark^\checkmark\\checkmarkc\checkmarki\checkmarkr\checkmarkc\checkmark$\checkmark \checkmark$\checkmark^\checkmark\\checkmarkc\checkmarki\checkmarkr\checkmarkc\checkmark$\checkmark\\checkmark$\checkmark^\checkmark\\checkmarkc\checkmarki\checkmarkr\checkmarkc\checkmark$\checkmarkh\checkmark$\checkmark^\checkmark\\checkmarkc\checkmarki\checkmarkr\checkmarkc\checkmark$\checkmarkl\checkmark$\checkmark^\checkmark\\checkmarkc\checkmarki\checkmarkr\checkmarkc\checkmark$\checkmarki\checkmark$\checkmark^\checkmark\\checkmarkc\checkmarki\checkmarkr\checkmarkc\checkmark$\checkmarkn\checkmark$\checkmark^\checkmark\\checkmarkc\checkmarki\checkmarkr\checkmarkc\checkmark$\checkmarke\checkmark$\checkmark^\checkmark\\checkmarkc\checkmarki\checkmarkr\checkmarkc\checkmark$\checkmark
\checkmark$\checkmark^\checkmark\\checkmarkc\checkmarki\checkmarkr\checkmarkc\checkmark$\checkmarkV\checkmark$\checkmark^\checkmark\\checkmarkc\checkmarki\checkmarkr\checkmarkc\checkmark$\checkmarki\checkmark$\checkmark^\checkmark\\checkmarkc\checkmarki\checkmarkr\checkmarkc\checkmark$\checkmarkt\checkmark$\checkmark^\checkmark\\checkmarkc\checkmarki\checkmarkr\checkmarkc\checkmark$\checkmarke\checkmark$\checkmark^\checkmark\\checkmarkc\checkmarki\checkmarkr\checkmarkc\checkmark$\checkmarks\checkmark$\checkmark^\checkmark\\checkmarkc\checkmarki\checkmarkr\checkmarkc\checkmark$\checkmarks\checkmark$\checkmark^\checkmark\\checkmarkc\checkmarki\checkmarkr\checkmarkc\checkmark$\checkmarke\checkmark$\checkmark^\checkmark\\checkmarkc\checkmarki\checkmarkr\checkmarkc\checkmark$\checkmark \checkmark$\checkmark^\checkmark\\checkmarkc\checkmarki\checkmarkr\checkmarkc\checkmark$\checkmarkd\checkmark$\checkmark^\checkmark\\checkmarkc\checkmarki\checkmarkr\checkmarkc\checkmark$\checkmarke\checkmark$\checkmark^\checkmark\\checkmarkc\checkmarki\checkmarkr\checkmarkc\checkmark$\checkmark \checkmark$\checkmark^\checkmark\\checkmarkc\checkmarki\checkmarkr\checkmarkc\checkmark$\checkmarkc\checkmark$\checkmark^\checkmark\\checkmarkc\checkmarki\checkmarkr\checkmarkc\checkmark$\checkmarko\checkmark$\checkmark^\checkmark\\checkmarkc\checkmarki\checkmarkr\checkmarkc\checkmark$\checkmarku\checkmark$\checkmark^\checkmark\\checkmarkc\checkmarki\checkmarkr\checkmarkc\checkmark$\checkmarkp\checkmark$\checkmark^\checkmark\\checkmarkc\checkmarki\checkmarkr\checkmarkc\checkmark$\checkmarke\checkmark$\checkmark^\checkmark\\checkmarkc\checkmarki\checkmarkr\checkmarkc\checkmark$\checkmark \checkmark$\checkmark^\checkmark\\checkmarkc\checkmarki\checkmarkr\checkmarkc\checkmark$\checkmark(\checkmark$\checkmark^\checkmark\\checkmarkc\checkmarki\checkmarkr\checkmarkc\checkmark$\checkmark$\checkmark$\checkmark^\checkmark\\checkmarkc\checkmarki\checkmarkr\checkmarkc\checkmark$\checkmarkV\checkmark$\checkmark^\checkmark\\checkmarkc\checkmarki\checkmarkr\checkmarkc\checkmark$\checkmark_\checkmark$\checkmark^\checkmark\\checkmarkc\checkmarki\checkmarkr\checkmarkc\checkmark$\checkmarkc\checkmark$\checkmark^\checkmark\\checkmarkc\checkmarki\checkmarkr\checkmarkc\checkmark$\checkmark$\checkmark$\checkmark^\checkmark\\checkmarkc\checkmarki\checkmarkr\checkmarkc\checkmark$\checkmark)\checkmark$\checkmark^\checkmark\\checkmarkc\checkmarki\checkmarkr\checkmarkc\checkmark$\checkmark \checkmark$\checkmark^\checkmark\\checkmarkc\checkmarki\checkmarkr\checkmarkc\checkmark$\checkmark&\checkmark$\checkmark^\checkmark\\checkmarkc\checkmarki\checkmarkr\checkmarkc\checkmark$\checkmark \checkmark$\checkmark^\checkmark\\checkmarkc\checkmarki\checkmarkr\checkmarkc\checkmark$\checkmark$\checkmark$\checkmark^\checkmark\\checkmarkc\checkmarki\checkmarkr\checkmarkc\checkmark$\checkmark1\checkmark$\checkmark^\checkmark\\checkmarkc\checkmarki\checkmarkr\checkmarkc\checkmark$\checkmark5\checkmark$\checkmark^\checkmark\\checkmarkc\checkmarki\checkmarkr\checkmarkc\checkmark$\checkmark0\checkmark$\checkmark^\checkmark\\checkmarkc\checkmarki\checkmarkr\checkmarkc\checkmark$\checkmark$\checkmark$\checkmark^\checkmark\\checkmarkc\checkmarki\checkmarkr\checkmarkc\checkmark$\checkmark \checkmark$\checkmark^\checkmark\\checkmarkc\checkmarki\checkmarkr\checkmarkc\checkmark$\checkmarkm\checkmark$\checkmark^\checkmark\\checkmarkc\checkmarki\checkmarkr\checkmarkc\checkmark$\checkmark/\checkmark$\checkmark^\checkmark\\checkmarkc\checkmarki\checkmarkr\checkmarkc\checkmark$\checkmarkm\checkmark$\checkmark^\checkmark\\checkmarkc\checkmarki\checkmarkr\checkmarkc\checkmark$\checkmarki\checkmark$\checkmark^\checkmark\\checkmarkc\checkmarki\checkmarkr\checkmarkc\checkmark$\checkmarkn\checkmark$\checkmark^\checkmark\\checkmarkc\checkmarki\checkmarkr\checkmarkc\checkmark$\checkmark \checkmark$\checkmark^\checkmark\\checkmarkc\checkmarki\checkmarkr\checkmarkc\checkmark$\checkmark&\checkmark$\checkmark^\checkmark\\checkmarkc\checkmarki\checkmarkr\checkmarkc\checkmark$\checkmark \checkmark$\checkmark^\checkmark\\checkmarkc\checkmarki\checkmarkr\checkmarkc\checkmark$\checkmark$\checkmark$\checkmark^\checkmark\\checkmarkc\checkmarki\checkmarkr\checkmarkc\checkmark$\checkmark\\checkmark$\checkmark^\checkmark\\checkmarkc\checkmarki\checkmarkr\checkmarkc\checkmark$\checkmarko\checkmark$\checkmark^\checkmark\\checkmarkc\checkmarki\checkmarkr\checkmarkc\checkmark$\checkmarkm\checkmark$\checkmark^\checkmark\\checkmarkc\checkmarki\checkmarkr\checkmarkc\checkmark$\checkmarke\checkmark$\checkmark^\checkmark\\checkmarkc\checkmarki\checkmarkr\checkmarkc\checkmark$\checkmarkg\checkmark$\checkmark^\checkmark\\checkmarkc\checkmarki\checkmarkr\checkmarkc\checkmark$\checkmarka\checkmark$\checkmark^\checkmark\\checkmarkc\checkmarki\checkmarkr\checkmarkc\checkmark$\checkmark \checkmark$\checkmark^\checkmark\\checkmarkc\checkmarki\checkmarkr\checkmarkc\checkmark$\checkmark=\checkmark$\checkmark^\checkmark\\checkmarkc\checkmarki\checkmarkr\checkmarkc\checkmark$\checkmark \checkmark$\checkmark^\checkmark\\checkmarkc\checkmarki\checkmarkr\checkmarkc\checkmark$\checkmark7\checkmark$\checkmark^\checkmark\\checkmarkc\checkmarki\checkmarkr\checkmarkc\checkmark$\checkmark9\checkmark$\checkmark^\checkmark\\checkmarkc\checkmarki\checkmarkr\checkmarkc\checkmark$\checkmark5\checkmark$\checkmark^\checkmark\\checkmarkc\checkmarki\checkmarkr\checkmarkc\checkmark$\checkmark8\checkmark$\checkmark^\checkmark\\checkmarkc\checkmarki\checkmarkr\checkmarkc\checkmark$\checkmark$\checkmark$\checkmark^\checkmark\\checkmarkc\checkmarki\checkmarkr\checkmarkc\checkmark$\checkmark \checkmark$\checkmark^\checkmark\\checkmarkc\checkmarki\checkmarkr\checkmarkc\checkmark$\checkmarkt\checkmark$\checkmark^\checkmark\\checkmarkc\checkmarki\checkmarkr\checkmarkc\checkmark$\checkmarkr\checkmark$\checkmark^\checkmark\\checkmarkc\checkmarki\checkmarkr\checkmarkc\checkmark$\checkmark/\checkmark$\checkmark^\checkmark\\checkmarkc\checkmarki\checkmarkr\checkmarkc\checkmark$\checkmarkm\checkmark$\checkmark^\checkmark\\checkmarkc\checkmarki\checkmarkr\checkmarkc\checkmark$\checkmarki\checkmark$\checkmark^\checkmark\\checkmarkc\checkmarki\checkmarkr\checkmarkc\checkmark$\checkmarkn\checkmark$\checkmark^\checkmark\\checkmarkc\checkmarki\checkmarkr\checkmarkc\checkmark$\checkmark \checkmark$\checkmark^\checkmark\\checkmarkc\checkmarki\checkmarkr\checkmarkc\checkmark$\checkmark\\checkmark$\checkmark^\checkmark\\checkmarkc\checkmarki\checkmarkr\checkmarkc\checkmark$\checkmark\\checkmark$\checkmark^\checkmark\\checkmarkc\checkmarki\checkmarkr\checkmarkc\checkmark$\checkmark \checkmark$\checkmark^\checkmark\\checkmarkc\checkmarki\checkmarkr\checkmarkc\checkmark$\checkmark\\checkmark$\checkmark^\checkmark\\checkmarkc\checkmarki\checkmarkr\checkmarkc\checkmark$\checkmarkh\checkmark$\checkmark^\checkmark\\checkmarkc\checkmarki\checkmarkr\checkmarkc\checkmark$\checkmarkl\checkmark$\checkmark^\checkmark\\checkmarkc\checkmarki\checkmarkr\checkmarkc\checkmark$\checkmarki\checkmark$\checkmark^\checkmark\\checkmarkc\checkmarki\checkmarkr\checkmarkc\checkmark$\checkmarkn\checkmark$\checkmark^\checkmark\\checkmarkc\checkmarki\checkmarkr\checkmarkc\checkmark$\checkmarke\checkmark$\checkmark^\checkmark\\checkmarkc\checkmarki\checkmarkr\checkmarkc\checkmark$\checkmark
\checkmark$\checkmark^\checkmark\\checkmarkc\checkmarki\checkmarkr\checkmarkc\checkmark$\checkmarkP\checkmark$\checkmark^\checkmark\\checkmarkc\checkmarki\checkmarkr\checkmarkc\checkmark$\checkmarkr\checkmark$\checkmark^\checkmark\\checkmarkc\checkmarki\checkmarkr\checkmarkc\checkmark$\checkmarke\checkmark$\checkmark^\checkmark\\checkmarkc\checkmarki\checkmarkr\checkmarkc\checkmark$\checkmarks\checkmark$\checkmark^\checkmark\\checkmarkc\checkmarki\checkmarkr\checkmarkc\checkmark$\checkmarks\checkmark$\checkmark^\checkmark\\checkmarkc\checkmarki\checkmarkr\checkmarkc\checkmark$\checkmarki\checkmark$\checkmark^\checkmark\\checkmarkc\checkmarki\checkmarkr\checkmarkc\checkmark$\checkmarko\checkmark$\checkmark^\checkmark\\checkmarkc\checkmarki\checkmarkr\checkmarkc\checkmark$\checkmarkn\checkmark$\checkmark^\checkmark\\checkmarkc\checkmarki\checkmarkr\checkmarkc\checkmark$\checkmark \checkmark$\checkmark^\checkmark\\checkmarkc\checkmarki\checkmarkr\checkmarkc\checkmark$\checkmarks\checkmark$\checkmark^\checkmark\\checkmarkc\checkmarki\checkmarkr\checkmarkc\checkmark$\checkmarkp\checkmark$\checkmark^\checkmark\\checkmarkc\checkmarki\checkmarkr\checkmarkc\checkmark$\checkmarké\checkmark$\checkmark^\checkmark\\checkmarkc\checkmarki\checkmarkr\checkmarkc\checkmark$\checkmarkc\checkmark$\checkmark^\checkmark\\checkmarkc\checkmarki\checkmarkr\checkmarkc\checkmark$\checkmarki\checkmark$\checkmark^\checkmark\\checkmarkc\checkmarki\checkmarkr\checkmarkc\checkmark$\checkmarkf\checkmark$\checkmark^\checkmark\\checkmarkc\checkmarki\checkmarkr\checkmarkc\checkmark$\checkmarki\checkmark$\checkmark^\checkmark\\checkmarkc\checkmarki\checkmarkr\checkmarkc\checkmark$\checkmarkq\checkmark$\checkmark^\checkmark\\checkmarkc\checkmarki\checkmarkr\checkmarkc\checkmark$\checkmarku\checkmark$\checkmark^\checkmark\\checkmarkc\checkmarki\checkmarkr\checkmarkc\checkmark$\checkmarke\checkmark$\checkmark^\checkmark\\checkmarkc\checkmarki\checkmarkr\checkmarkc\checkmark$\checkmark \checkmark$\checkmark^\checkmark\\checkmarkc\checkmarki\checkmarkr\checkmarkc\checkmark$\checkmarkd\checkmark$\checkmark^\checkmark\\checkmarkc\checkmarki\checkmarkr\checkmarkc\checkmark$\checkmarke\checkmark$\checkmark^\checkmark\\checkmarkc\checkmarki\checkmarkr\checkmarkc\checkmark$\checkmark \checkmark$\checkmark^\checkmark\\checkmarkc\checkmarki\checkmarkr\checkmarkc\checkmark$\checkmarkc\checkmark$\checkmark^\checkmark\\checkmarkc\checkmarki\checkmarkr\checkmarkc\checkmark$\checkmarko\checkmark$\checkmark^\checkmark\\checkmarkc\checkmarki\checkmarkr\checkmarkc\checkmark$\checkmarku\checkmark$\checkmark^\checkmark\\checkmarkc\checkmarki\checkmarkr\checkmarkc\checkmark$\checkmarkp\checkmark$\checkmark^\checkmark\\checkmarkc\checkmarki\checkmarkr\checkmarkc\checkmark$\checkmarke\checkmark$\checkmark^\checkmark\\checkmarkc\checkmarki\checkmarkr\checkmarkc\checkmark$\checkmark \checkmark$\checkmark^\checkmark\\checkmarkc\checkmarki\checkmarkr\checkmarkc\checkmark$\checkmark(\checkmark$\checkmark^\checkmark\\checkmarkc\checkmarki\checkmarkr\checkmarkc\checkmark$\checkmark$\checkmark$\checkmark^\checkmark\\checkmarkc\checkmarki\checkmarkr\checkmarkc\checkmark$\checkmarkk\checkmark$\checkmark^\checkmark\\checkmarkc\checkmarki\checkmarkr\checkmarkc\checkmark$\checkmark_\checkmark$\checkmark^\checkmark\\checkmarkc\checkmarki\checkmarkr\checkmarkc\checkmark$\checkmarkc\checkmark$\checkmark^\checkmark\\checkmarkc\checkmarki\checkmarkr\checkmarkc\checkmark$\checkmark$\checkmark$\checkmark^\checkmark\\checkmarkc\checkmarki\checkmarkr\checkmarkc\checkmark$\checkmark)\checkmark$\checkmark^\checkmark\\checkmarkc\checkmarki\checkmarkr\checkmarkc\checkmark$\checkmark \checkmark$\checkmark^\checkmark\\checkmarkc\checkmarki\checkmarkr\checkmarkc\checkmark$\checkmark&\checkmark$\checkmark^\checkmark\\checkmarkc\checkmarki\checkmarkr\checkmarkc\checkmark$\checkmark \checkmark$\checkmark^\checkmark\\checkmarkc\checkmarki\checkmarkr\checkmarkc\checkmark$\checkmark$\checkmark$\checkmark^\checkmark\\checkmarkc\checkmarki\checkmarkr\checkmarkc\checkmark$\checkmark7\checkmark$\checkmark^\checkmark\\checkmarkc\checkmarki\checkmarkr\checkmarkc\checkmark$\checkmark0\checkmark$\checkmark^\checkmark\\checkmarkc\checkmarki\checkmarkr\checkmarkc\checkmark$\checkmark0\checkmark$\checkmark^\checkmark\\checkmarkc\checkmarki\checkmarkr\checkmarkc\checkmark$\checkmark$\checkmark$\checkmark^\checkmark\\checkmarkc\checkmarki\checkmarkr\checkmarkc\checkmark$\checkmark \checkmark$\checkmark^\checkmark\\checkmarkc\checkmarki\checkmarkr\checkmarkc\checkmark$\checkmarkN\checkmark$\checkmark^\checkmark\\checkmarkc\checkmarki\checkmarkr\checkmarkc\checkmark$\checkmark/\checkmark$\checkmark^\checkmark\\checkmarkc\checkmarki\checkmarkr\checkmarkc\checkmark$\checkmarkm\checkmark$\checkmark^\checkmark\\checkmarkc\checkmarki\checkmarkr\checkmarkc\checkmark$\checkmarkm\checkmark$\checkmark^\checkmark\\checkmarkc\checkmarki\checkmarkr\checkmarkc\checkmark$\checkmark²\checkmark$\checkmark^\checkmark\\checkmarkc\checkmarki\checkmarkr\checkmarkc\checkmark$\checkmark \checkmark$\checkmark^\checkmark\\checkmarkc\checkmarki\checkmarkr\checkmarkc\checkmark$\checkmark&\checkmark$\checkmark^\checkmark\\checkmarkc\checkmarki\checkmarkr\checkmarkc\checkmark$\checkmark \checkmark$\checkmark^\checkmark\\checkmarkc\checkmarki\checkmarkr\checkmarkc\checkmark$\checkmarkA\checkmark$\checkmark^\checkmark\\checkmarkc\checkmarki\checkmarkr\checkmarkc\checkmark$\checkmarkl\checkmark$\checkmark^\checkmark\\checkmarkc\checkmarki\checkmarkr\checkmarkc\checkmark$\checkmarku\checkmark$\checkmark^\checkmark\\checkmarkc\checkmarki\checkmarkr\checkmarkc\checkmark$\checkmarkm\checkmark$\checkmark^\checkmark\\checkmarkc\checkmarki\checkmarkr\checkmarkc\checkmark$\checkmarki\checkmark$\checkmark^\checkmark\\checkmarkc\checkmarki\checkmarkr\checkmarkc\checkmark$\checkmarkn\checkmark$\checkmark^\checkmark\\checkmarkc\checkmarki\checkmarkr\checkmarkc\checkmark$\checkmarki\checkmark$\checkmark^\checkmark\\checkmarkc\checkmarki\checkmarkr\checkmarkc\checkmark$\checkmarku\checkmark$\checkmark^\checkmark\\checkmarkc\checkmarki\checkmarkr\checkmarkc\checkmark$\checkmarkm\checkmark$\checkmark^\checkmark\\checkmarkc\checkmarki\checkmarkr\checkmarkc\checkmark$\checkmark \checkmark$\checkmark^\checkmark\\checkmarkc\checkmarki\checkmarkr\checkmarkc\checkmark$\checkmarkA\checkmark$\checkmark^\checkmark\\checkmarkc\checkmarki\checkmarkr\checkmarkc\checkmark$\checkmarkW\checkmark$\checkmark^\checkmark\\checkmarkc\checkmarki\checkmarkr\checkmarkc\checkmark$\checkmark-\checkmark$\checkmark^\checkmark\\checkmarkc\checkmarki\checkmarkr\checkmarkc\checkmark$\checkmark2\checkmark$\checkmark^\checkmark\\checkmarkc\checkmarki\checkmarkr\checkmarkc\checkmark$\checkmark0\checkmark$\checkmark^\checkmark\\checkmarkc\checkmarki\checkmarkr\checkmarkc\checkmark$\checkmark1\checkmark$\checkmark^\checkmark\\checkmarkc\checkmarki\checkmarkr\checkmarkc\checkmark$\checkmark7\checkmark$\checkmark^\checkmark\\checkmarkc\checkmarki\checkmarkr\checkmarkc\checkmark$\checkmark \checkmark$\checkmark^\checkmark\\checkmarkc\checkmarki\checkmarkr\checkmarkc\checkmark$\checkmark\\checkmark$\checkmark^\checkmark\\checkmarkc\checkmarki\checkmarkr\checkmarkc\checkmark$\checkmark\\checkmark$\checkmark^\checkmark\\checkmarkc\checkmarki\checkmarkr\checkmarkc\checkmark$\checkmark \checkmark$\checkmark^\checkmark\\checkmarkc\checkmarki\checkmarkr\checkmarkc\checkmark$\checkmark\\checkmark$\checkmark^\checkmark\\checkmarkc\checkmarki\checkmarkr\checkmarkc\checkmark$\checkmarkh\checkmark$\checkmark^\checkmark\\checkmarkc\checkmarki\checkmarkr\checkmarkc\checkmark$\checkmarkl\checkmark$\checkmark^\checkmark\\checkmarkc\checkmarki\checkmarkr\checkmarkc\checkmark$\checkmarki\checkmark$\checkmark^\checkmark\\checkmarkc\checkmarki\checkmarkr\checkmarkc\checkmark$\checkmarkn\checkmark$\checkmark^\checkmark\\checkmarkc\checkmarki\checkmarkr\checkmarkc\checkmark$\checkmarke\checkmark$\checkmark^\checkmark\\checkmarkc\checkmarki\checkmarkr\checkmarkc\checkmark$\checkmark
\checkmark$\checkmark^\checkmark\\checkmarkc\checkmarki\checkmarkr\checkmarkc\checkmark$\checkmark\\checkmark$\checkmark^\checkmark\\checkmarkc\checkmarki\checkmarkr\checkmarkc\checkmark$\checkmarke\checkmark$\checkmark^\checkmark\\checkmarkc\checkmarki\checkmarkr\checkmarkc\checkmark$\checkmarkn\checkmark$\checkmark^\checkmark\\checkmarkc\checkmarki\checkmarkr\checkmarkc\checkmark$\checkmarkd\checkmark$\checkmark^\checkmark\\checkmarkc\checkmarki\checkmarkr\checkmarkc\checkmark$\checkmark{\checkmark$\checkmark^\checkmark\\checkmarkc\checkmarki\checkmarkr\checkmarkc\checkmark$\checkmarkt\checkmark$\checkmark^\checkmark\\checkmarkc\checkmarki\checkmarkr\checkmarkc\checkmark$\checkmarka\checkmark$\checkmark^\checkmark\\checkmarkc\checkmarki\checkmarkr\checkmarkc\checkmark$\checkmarkb\checkmark$\checkmark^\checkmark\\checkmarkc\checkmarki\checkmarkr\checkmarkc\checkmark$\checkmarku\checkmark$\checkmark^\checkmark\\checkmarkc\checkmarki\checkmarkr\checkmarkc\checkmark$\checkmarkl\checkmark$\checkmark^\checkmark\\checkmarkc\checkmarki\checkmarkr\checkmarkc\checkmark$\checkmarka\checkmark$\checkmark^\checkmark\\checkmarkc\checkmarki\checkmarkr\checkmarkc\checkmark$\checkmarkr\checkmark$\checkmark^\checkmark\\checkmarkc\checkmarki\checkmarkr\checkmarkc\checkmark$\checkmark}\checkmark$\checkmark^\checkmark\\checkmarkc\checkmarki\checkmarkr\checkmarkc\checkmark$\checkmark
\checkmark$\checkmark^\checkmark\\checkmarkc\checkmarki\checkmarkr\checkmarkc\checkmark$\checkmark\\checkmark$\checkmark^\checkmark\\checkmarkc\checkmarki\checkmarkr\checkmarkc\checkmark$\checkmarkc\checkmark$\checkmark^\checkmark\\checkmarkc\checkmarki\checkmarkr\checkmarkc\checkmark$\checkmarka\checkmark$\checkmark^\checkmark\\checkmarkc\checkmarki\checkmarkr\checkmarkc\checkmark$\checkmarkp\checkmark$\checkmark^\checkmark\\checkmarkc\checkmarki\checkmarkr\checkmarkc\checkmark$\checkmarkt\checkmark$\checkmark^\checkmark\\checkmarkc\checkmarki\checkmarkr\checkmarkc\checkmark$\checkmarki\checkmark$\checkmark^\checkmark\\checkmarkc\checkmarki\checkmarkr\checkmarkc\checkmark$\checkmarko\checkmark$\checkmark^\checkmark\\checkmarkc\checkmarki\checkmarkr\checkmarkc\checkmark$\checkmarkn\checkmark$\checkmark^\checkmark\\checkmarkc\checkmarki\checkmarkr\checkmarkc\checkmark$\checkmark{\checkmark$\checkmark^\checkmark\\checkmarkc\checkmarki\checkmarkr\checkmarkc\checkmark$\checkmarkP\checkmark$\checkmark^\checkmark\\checkmarkc\checkmarki\checkmarkr\checkmarkc\checkmark$\checkmarka\checkmark$\checkmark^\checkmark\\checkmarkc\checkmarki\checkmarkr\checkmarkc\checkmark$\checkmarkr\checkmark$\checkmark^\checkmark\\checkmarkc\checkmarki\checkmarkr\checkmarkc\checkmark$\checkmarka\checkmark$\checkmark^\checkmark\\checkmarkc\checkmarki\checkmarkr\checkmarkc\checkmark$\checkmarkm\checkmark$\checkmark^\checkmark\\checkmarkc\checkmarki\checkmarkr\checkmarkc\checkmark$\checkmarkè\checkmark$\checkmark^\checkmark\\checkmarkc\checkmarki\checkmarkr\checkmarkc\checkmark$\checkmarkt\checkmark$\checkmark^\checkmark\\checkmarkc\checkmarki\checkmarkr\checkmarkc\checkmark$\checkmarkr\checkmark$\checkmark^\checkmark\\checkmarkc\checkmarki\checkmarkr\checkmarkc\checkmark$\checkmarke\checkmark$\checkmark^\checkmark\\checkmarkc\checkmarki\checkmarkr\checkmarkc\checkmark$\checkmarks\checkmark$\checkmark^\checkmark\\checkmarkc\checkmarki\checkmarkr\checkmarkc\checkmark$\checkmark \checkmark$\checkmark^\checkmark\\checkmarkc\checkmarki\checkmarkr\checkmarkc\checkmark$\checkmarkd\checkmark$\checkmark^\checkmark\\checkmarkc\checkmarki\checkmarkr\checkmarkc\checkmark$\checkmarke\checkmark$\checkmark^\checkmark\\checkmarkc\checkmarki\checkmarkr\checkmarkc\checkmark$\checkmark \checkmark$\checkmark^\checkmark\\checkmarkc\checkmarki\checkmarkr\checkmarkc\checkmark$\checkmarkc\checkmark$\checkmark^\checkmark\\checkmarkc\checkmarki\checkmarkr\checkmarkc\checkmark$\checkmarko\checkmark$\checkmark^\checkmark\\checkmarkc\checkmarki\checkmarkr\checkmarkc\checkmark$\checkmarku\checkmark$\checkmark^\checkmark\\checkmarkc\checkmarki\checkmarkr\checkmarkc\checkmark$\checkmarkp\checkmark$\checkmark^\checkmark\\checkmarkc\checkmarki\checkmarkr\checkmarkc\checkmark$\checkmarke\checkmark$\checkmark^\checkmark\\checkmarkc\checkmarki\checkmarkr\checkmarkc\checkmark$\checkmark \checkmark$\checkmark^\checkmark\\checkmarkc\checkmarki\checkmarkr\checkmarkc\checkmark$\checkmarka\checkmark$\checkmark^\checkmark\\checkmarkc\checkmarki\checkmarkr\checkmarkc\checkmark$\checkmarkd\checkmark$\checkmark^\checkmark\\checkmarkc\checkmarki\checkmarkr\checkmarkc\checkmark$\checkmarko\checkmark$\checkmark^\checkmark\\checkmarkc\checkmarki\checkmarkr\checkmarkc\checkmark$\checkmarkp\checkmark$\checkmark^\checkmark\\checkmarkc\checkmarki\checkmarkr\checkmarkc\checkmark$\checkmarkt\checkmark$\checkmark^\checkmark\\checkmarkc\checkmarki\checkmarkr\checkmarkc\checkmark$\checkmarké\checkmark$\checkmark^\checkmark\\checkmarkc\checkmarki\checkmarkr\checkmarkc\checkmark$\checkmarks\checkmark$\checkmark^\checkmark\\checkmarkc\checkmarki\checkmarkr\checkmarkc\checkmark$\checkmark \checkmark$\checkmark^\checkmark\\checkmarkc\checkmarki\checkmarkr\checkmarkc\checkmark$\checkmarkp\checkmark$\checkmark^\checkmark\\checkmarkc\checkmarki\checkmarkr\checkmarkc\checkmark$\checkmarko\checkmark$\checkmark^\checkmark\\checkmarkc\checkmarki\checkmarkr\checkmarkc\checkmark$\checkmarku\checkmark$\checkmark^\checkmark\\checkmarkc\checkmarki\checkmarkr\checkmarkc\checkmark$\checkmarkr\checkmark$\checkmark^\checkmark\\checkmarkc\checkmarki\checkmarkr\checkmarkc\checkmark$\checkmark \checkmark$\checkmark^\checkmark\\checkmarkc\checkmarki\checkmarkr\checkmarkc\checkmark$\checkmarkl\checkmark$\checkmark^\checkmark\\checkmarkc\checkmarki\checkmarkr\checkmarkc\checkmark$\checkmarke\checkmark$\checkmark^\checkmark\\checkmarkc\checkmarki\checkmarkr\checkmarkc\checkmark$\checkmark \checkmark$\checkmark^\checkmark\\checkmarkc\checkmarki\checkmarkr\checkmarkc\checkmark$\checkmarkc\checkmark$\checkmark^\checkmark\\checkmarkc\checkmarki\checkmarkr\checkmarkc\checkmark$\checkmarka\checkmark$\checkmark^\checkmark\\checkmarkc\checkmarki\checkmarkr\checkmarkc\checkmark$\checkmarks\checkmark$\checkmark^\checkmark\\checkmarkc\checkmarki\checkmarkr\checkmarkc\checkmark$\checkmark \checkmark$\checkmark^\checkmark\\checkmarkc\checkmarki\checkmarkr\checkmarkc\checkmark$\checkmarkl\checkmark$\checkmark^\checkmark\\checkmarkc\checkmarki\checkmarkr\checkmarkc\checkmark$\checkmarke\checkmark$\checkmark^\checkmark\\checkmarkc\checkmarki\checkmarkr\checkmarkc\checkmark$\checkmark \checkmark$\checkmark^\checkmark\\checkmarkc\checkmarki\checkmarkr\checkmarkc\checkmark$\checkmarkp\checkmark$\checkmark^\checkmark\\checkmarkc\checkmarki\checkmarkr\checkmarkc\checkmark$\checkmarkl\checkmark$\checkmark^\checkmark\\checkmarkc\checkmarki\checkmarkr\checkmarkc\checkmark$\checkmarku\checkmark$\checkmark^\checkmark\\checkmarkc\checkmarki\checkmarkr\checkmarkc\checkmark$\checkmarks\checkmark$\checkmark^\checkmark\\checkmarkc\checkmarki\checkmarkr\checkmarkc\checkmark$\checkmark \checkmark$\checkmark^\checkmark\\checkmarkc\checkmarki\checkmarkr\checkmarkc\checkmark$\checkmarkd\checkmark$\checkmark^\checkmark\\checkmarkc\checkmarki\checkmarkr\checkmarkc\checkmark$\checkmarké\checkmark$\checkmark^\checkmark\\checkmarkc\checkmarki\checkmarkr\checkmarkc\checkmark$\checkmarkf\checkmark$\checkmark^\checkmark\\checkmarkc\checkmarki\checkmarkr\checkmarkc\checkmark$\checkmarka\checkmark$\checkmark^\checkmark\\checkmarkc\checkmarki\checkmarkr\checkmarkc\checkmark$\checkmarkv\checkmark$\checkmark^\checkmark\\checkmarkc\checkmarki\checkmarkr\checkmarkc\checkmark$\checkmarko\checkmark$\checkmark^\checkmark\\checkmarkc\checkmarki\checkmarkr\checkmarkc\checkmark$\checkmarkr\checkmark$\checkmark^\checkmark\\checkmarkc\checkmarki\checkmarkr\checkmarkc\checkmark$\checkmarka\checkmark$\checkmark^\checkmark\\checkmarkc\checkmarki\checkmarkr\checkmarkc\checkmark$\checkmarkb\checkmark$\checkmark^\checkmark\\checkmarkc\checkmarki\checkmarkr\checkmarkc\checkmark$\checkmarkl\checkmark$\checkmark^\checkmark\\checkmarkc\checkmarki\checkmarkr\checkmarkc\checkmark$\checkmarke\checkmark$\checkmark^\checkmark\\checkmarkc\checkmarki\checkmarkr\checkmarkc\checkmark$\checkmark \checkmark$\checkmark^\checkmark\\checkmarkc\checkmarki\checkmarkr\checkmarkc\checkmark$\checkmark(\checkmark$\checkmark^\checkmark\\checkmarkc\checkmarki\checkmarkr\checkmarkc\checkmark$\checkmarkA\checkmark$\checkmark^\checkmark\\checkmarkc\checkmarki\checkmarkr\checkmarkc\checkmark$\checkmarkl\checkmark$\checkmark^\checkmark\\checkmarkc\checkmarki\checkmarkr\checkmarkc\checkmark$\checkmark \checkmark$\checkmark^\checkmark\\checkmarkc\checkmarki\checkmarkr\checkmarkc\checkmark$\checkmarkA\checkmark$\checkmark^\checkmark\\checkmarkc\checkmarki\checkmarkr\checkmarkc\checkmark$\checkmarkW\checkmark$\checkmark^\checkmark\\checkmarkc\checkmarki\checkmarkr\checkmarkc\checkmark$\checkmark-\checkmark$\checkmark^\checkmark\\checkmarkc\checkmarki\checkmarkr\checkmarkc\checkmark$\checkmark2\checkmark$\checkmark^\checkmark\\checkmarkc\checkmarki\checkmarkr\checkmarkc\checkmark$\checkmark0\checkmark$\checkmark^\checkmark\\checkmarkc\checkmarki\checkmarkr\checkmarkc\checkmark$\checkmark1\checkmark$\checkmark^\checkmark\\checkmarkc\checkmarki\checkmarkr\checkmarkc\checkmark$\checkmark7\checkmark$\checkmark^\checkmark\\checkmarkc\checkmarki\checkmarkr\checkmarkc\checkmark$\checkmarkA\checkmark$\checkmark^\checkmark\\checkmarkc\checkmarki\checkmarkr\checkmarkc\checkmark$\checkmark)\checkmark$\checkmark^\checkmark\\checkmarkc\checkmarki\checkmarkr\checkmarkc\checkmark$\checkmark}\checkmark$\checkmark^\checkmark\\checkmarkc\checkmarki\checkmarkr\checkmarkc\checkmark$\checkmark
\checkmark$\checkmark^\checkmark\\checkmarkc\checkmarki\checkmarkr\checkmarkc\checkmark$\checkmark\\checkmark$\checkmark^\checkmark\\checkmarkc\checkmarki\checkmarkr\checkmarkc\checkmark$\checkmarkl\checkmark$\checkmark^\checkmark\\checkmarkc\checkmarki\checkmarkr\checkmarkc\checkmark$\checkmarka\checkmark$\checkmark^\checkmark\\checkmarkc\checkmarki\checkmarkr\checkmarkc\checkmark$\checkmarkb\checkmark$\checkmark^\checkmark\\checkmarkc\checkmarki\checkmarkr\checkmarkc\checkmark$\checkmarke\checkmark$\checkmark^\checkmark\\checkmarkc\checkmarki\checkmarkr\checkmarkc\checkmark$\checkmarkl\checkmark$\checkmark^\checkmark\\checkmarkc\checkmarki\checkmarkr\checkmarkc\checkmark$\checkmark{\checkmark$\checkmark^\checkmark\\checkmarkc\checkmarki\checkmarkr\checkmarkc\checkmark$\checkmarkt\checkmark$\checkmark^\checkmark\\checkmarkc\checkmarki\checkmarkr\checkmarkc\checkmark$\checkmarka\checkmark$\checkmark^\checkmark\\checkmarkc\checkmarki\checkmarkr\checkmarkc\checkmark$\checkmarkb\checkmark$\checkmark^\checkmark\\checkmarkc\checkmarki\checkmarkr\checkmarkc\checkmark$\checkmark:\checkmark$\checkmark^\checkmark\\checkmarkc\checkmarki\checkmarkr\checkmarkc\checkmark$\checkmarkp\checkmark$\checkmark^\checkmark\\checkmarkc\checkmarki\checkmarkr\checkmarkc\checkmark$\checkmarka\checkmark$\checkmark^\checkmark\\checkmarkc\checkmarki\checkmarkr\checkmarkc\checkmark$\checkmarkr\checkmark$\checkmark^\checkmark\\checkmarkc\checkmarki\checkmarkr\checkmarkc\checkmark$\checkmarka\checkmark$\checkmark^\checkmark\\checkmarkc\checkmarki\checkmarkr\checkmarkc\checkmark$\checkmarkm\checkmark$\checkmark^\checkmark\\checkmarkc\checkmarki\checkmarkr\checkmarkc\checkmark$\checkmarks\checkmark$\checkmark^\checkmark\\checkmarkc\checkmarki\checkmarkr\checkmarkc\checkmark$\checkmark_\checkmark$\checkmark^\checkmark\\checkmarkc\checkmarki\checkmarkr\checkmarkc\checkmark$\checkmarkc\checkmark$\checkmark^\checkmark\\checkmarkc\checkmarki\checkmarkr\checkmarkc\checkmark$\checkmarko\checkmark$\checkmark^\checkmark\\checkmarkc\checkmarki\checkmarkr\checkmarkc\checkmark$\checkmarku\checkmark$\checkmark^\checkmark\\checkmarkc\checkmarki\checkmarkr\checkmarkc\checkmark$\checkmarkp\checkmark$\checkmark^\checkmark\\checkmarkc\checkmarki\checkmarkr\checkmarkc\checkmark$\checkmarke\checkmark$\checkmark^\checkmark\\checkmarkc\checkmarki\checkmarkr\checkmarkc\checkmark$\checkmark}\checkmark$\checkmark^\checkmark\\checkmarkc\checkmarki\checkmarkr\checkmarkc\checkmark$\checkmark
\checkmark$\checkmark^\checkmark\\checkmarkc\checkmarki\checkmarkr\checkmarkc\checkmark$\checkmark\\checkmark$\checkmark^\checkmark\\checkmarkc\checkmarki\checkmarkr\checkmarkc\checkmark$\checkmarke\checkmark$\checkmark^\checkmark\\checkmarkc\checkmarki\checkmarkr\checkmarkc\checkmark$\checkmarkn\checkmark$\checkmark^\checkmark\\checkmarkc\checkmarki\checkmarkr\checkmarkc\checkmark$\checkmarkd\checkmark$\checkmark^\checkmark\\checkmarkc\checkmarki\checkmarkr\checkmarkc\checkmark$\checkmark{\checkmark$\checkmark^\checkmark\\checkmarkc\checkmarki\checkmarkr\checkmarkc\checkmark$\checkmarkt\checkmark$\checkmark^\checkmark\\checkmarkc\checkmarki\checkmarkr\checkmarkc\checkmark$\checkmarka\checkmark$\checkmark^\checkmark\\checkmarkc\checkmarki\checkmarkr\checkmarkc\checkmark$\checkmarkb\checkmark$\checkmark^\checkmark\\checkmarkc\checkmarki\checkmarkr\checkmarkc\checkmark$\checkmarkl\checkmark$\checkmark^\checkmark\\checkmarkc\checkmarki\checkmarkr\checkmarkc\checkmark$\checkmarke\checkmark$\checkmark^\checkmark\\checkmarkc\checkmarki\checkmarkr\checkmarkc\checkmark$\checkmark}\checkmark$\checkmark^\checkmark\\checkmarkc\checkmarki\checkmarkr\checkmarkc\checkmark$\checkmark
\checkmark$\checkmark^\checkmark\\checkmarkc\checkmarki\checkmarkr\checkmarkc\checkmark$\checkmark
\checkmark$\checkmark^\checkmark\\checkmarkc\checkmarki\checkmarkr\checkmarkc\checkmark$\checkmark\\checkmark$\checkmark^\checkmark\\checkmarkc\checkmarki\checkmarkr\checkmarkc\checkmark$\checkmarks\checkmark$\checkmark^\checkmark\\checkmarkc\checkmarki\checkmarkr\checkmarkc\checkmark$\checkmarku\checkmark$\checkmark^\checkmark\\checkmarkc\checkmarki\checkmarkr\checkmarkc\checkmark$\checkmarkb\checkmark$\checkmark^\checkmark\\checkmarkc\checkmarki\checkmarkr\checkmarkc\checkmark$\checkmarks\checkmark$\checkmark^\checkmark\\checkmarkc\checkmarki\checkmarkr\checkmarkc\checkmark$\checkmarke\checkmark$\checkmark^\checkmark\\checkmarkc\checkmarki\checkmarkr\checkmarkc\checkmark$\checkmarkc\checkmark$\checkmark^\checkmark\\checkmarkc\checkmarki\checkmarkr\checkmarkc\checkmark$\checkmarkt\checkmark$\checkmark^\checkmark\\checkmarkc\checkmarki\checkmarkr\checkmarkc\checkmark$\checkmarki\checkmark$\checkmark^\checkmark\\checkmarkc\checkmarki\checkmarkr\checkmarkc\checkmark$\checkmarko\checkmark$\checkmark^\checkmark\\checkmarkc\checkmarki\checkmarkr\checkmarkc\checkmark$\checkmarkn\checkmark$\checkmark^\checkmark\\checkmarkc\checkmarki\checkmarkr\checkmarkc\checkmark$\checkmark{\checkmark$\checkmark^\checkmark\\checkmarkc\checkmarki\checkmarkr\checkmarkc\checkmark$\checkmarkC\checkmark$\checkmark^\checkmark\\checkmarkc\checkmarki\checkmarkr\checkmarkc\checkmark$\checkmarka\checkmark$\checkmark^\checkmark\\checkmarkc\checkmarki\checkmarkr\checkmarkc\checkmark$\checkmarkl\checkmark$\checkmark^\checkmark\\checkmarkc\checkmarki\checkmarkr\checkmarkc\checkmark$\checkmarkc\checkmark$\checkmark^\checkmark\\checkmarkc\checkmarki\checkmarkr\checkmarkc\checkmark$\checkmarku\checkmark$\checkmark^\checkmark\\checkmarkc\checkmarki\checkmarkr\checkmarkc\checkmark$\checkmarkl\checkmark$\checkmark^\checkmark\\checkmarkc\checkmarki\checkmarkr\checkmarkc\checkmark$\checkmark \checkmark$\checkmark^\checkmark\\checkmarkc\checkmarki\checkmarkr\checkmarkc\checkmark$\checkmarkd\checkmark$\checkmark^\checkmark\\checkmarkc\checkmarki\checkmarkr\checkmarkc\checkmark$\checkmarke\checkmark$\checkmark^\checkmark\\checkmarkc\checkmarki\checkmarkr\checkmarkc\checkmark$\checkmarks\checkmark$\checkmark^\checkmark\\checkmarkc\checkmarki\checkmarkr\checkmarkc\checkmark$\checkmark \checkmark$\checkmark^\checkmark\\checkmarkc\checkmarki\checkmarkr\checkmarkc\checkmark$\checkmarkc\checkmark$\checkmark^\checkmark\\checkmarkc\checkmarki\checkmarkr\checkmarkc\checkmark$\checkmarko\checkmark$\checkmark^\checkmark\\checkmarkc\checkmarki\checkmarkr\checkmarkc\checkmark$\checkmarkm\checkmark$\checkmark^\checkmark\\checkmarkc\checkmarki\checkmarkr\checkmarkc\checkmark$\checkmarkp\checkmark$\checkmark^\checkmark\\checkmarkc\checkmarki\checkmarkr\checkmarkc\checkmark$\checkmarko\checkmark$\checkmark^\checkmark\\checkmarkc\checkmarki\checkmarkr\checkmarkc\checkmark$\checkmarks\checkmark$\checkmark^\checkmark\\checkmarkc\checkmarki\checkmarkr\checkmarkc\checkmark$\checkmarka\checkmark$\checkmark^\checkmark\\checkmarkc\checkmarki\checkmarkr\checkmarkc\checkmark$\checkmarkn\checkmark$\checkmark^\checkmark\\checkmarkc\checkmarki\checkmarkr\checkmarkc\checkmark$\checkmarkt\checkmark$\checkmark^\checkmark\\checkmarkc\checkmarki\checkmarkr\checkmarkc\checkmark$\checkmarke\checkmark$\checkmark^\checkmark\\checkmarkc\checkmarki\checkmarkr\checkmarkc\checkmark$\checkmarks\checkmark$\checkmark^\checkmark\\checkmarkc\checkmarki\checkmarkr\checkmarkc\checkmark$\checkmark \checkmark$\checkmark^\checkmark\\checkmarkc\checkmarki\checkmarkr\checkmarkc\checkmark$\checkmarkd\checkmark$\checkmark^\checkmark\\checkmarkc\checkmarki\checkmarkr\checkmarkc\checkmark$\checkmarke\checkmark$\checkmark^\checkmark\\checkmarkc\checkmarki\checkmarkr\checkmarkc\checkmark$\checkmark \checkmark$\checkmark^\checkmark\\checkmarkc\checkmarki\checkmarkr\checkmarkc\checkmark$\checkmarkl\checkmark$\checkmark^\checkmark\\checkmarkc\checkmarki\checkmarkr\checkmarkc\checkmark$\checkmarka\checkmark$\checkmark^\checkmark\\checkmarkc\checkmarki\checkmarkr\checkmarkc\checkmark$\checkmark \checkmark$\checkmark^\checkmark\\checkmarkc\checkmarki\checkmarkr\checkmarkc\checkmark$\checkmarkf\checkmark$\checkmark^\checkmark\\checkmarkc\checkmarki\checkmarkr\checkmarkc\checkmark$\checkmarko\checkmark$\checkmark^\checkmark\\checkmarkc\checkmarki\checkmarkr\checkmarkc\checkmark$\checkmarkr\checkmark$\checkmark^\checkmark\\checkmarkc\checkmarki\checkmarkr\checkmarkc\checkmark$\checkmarkc\checkmark$\checkmark^\checkmark\\checkmarkc\checkmarki\checkmarkr\checkmarkc\checkmark$\checkmarke\checkmark$\checkmark^\checkmark\\checkmarkc\checkmarki\checkmarkr\checkmarkc\checkmark$\checkmark \checkmark$\checkmark^\checkmark\\checkmarkc\checkmarki\checkmarkr\checkmarkc\checkmark$\checkmarkd\checkmark$\checkmark^\checkmark\\checkmarkc\checkmarki\checkmarkr\checkmarkc\checkmark$\checkmarke\checkmark$\checkmark^\checkmark\\checkmarkc\checkmarki\checkmarkr\checkmarkc\checkmark$\checkmark \checkmark$\checkmark^\checkmark\\checkmarkc\checkmarki\checkmarkr\checkmarkc\checkmark$\checkmarkc\checkmark$\checkmark^\checkmark\\checkmarkc\checkmarki\checkmarkr\checkmarkc\checkmark$\checkmarko\checkmark$\checkmark^\checkmark\\checkmarkc\checkmarki\checkmarkr\checkmarkc\checkmark$\checkmarku\checkmark$\checkmark^\checkmark\\checkmarkc\checkmarki\checkmarkr\checkmarkc\checkmark$\checkmarkp\checkmark$\checkmark^\checkmark\\checkmarkc\checkmarki\checkmarkr\checkmarkc\checkmark$\checkmarke\checkmark$\checkmark^\checkmark\\checkmarkc\checkmarki\checkmarkr\checkmarkc\checkmark$\checkmark}\checkmark$\checkmark^\checkmark\\checkmarkc\checkmarki\checkmarkr\checkmarkc\checkmark$\checkmark
\checkmark$\checkmark^\checkmark\\checkmarkc\checkmarki\checkmarkr\checkmarkc\checkmark$\checkmark
\checkmark$\checkmark^\checkmark\\checkmarkc\checkmarki\checkmarkr\checkmarkc\checkmark$\checkmark\\checkmark$\checkmark^\checkmark\\checkmarkc\checkmarki\checkmarkr\checkmarkc\checkmark$\checkmarks\checkmark$\checkmark^\checkmark\\checkmarkc\checkmarki\checkmarkr\checkmarkc\checkmark$\checkmarku\checkmark$\checkmark^\checkmark\\checkmarkc\checkmarki\checkmarkr\checkmarkc\checkmark$\checkmarkb\checkmark$\checkmark^\checkmark\\checkmarkc\checkmarki\checkmarkr\checkmarkc\checkmark$\checkmarks\checkmark$\checkmark^\checkmark\\checkmarkc\checkmarki\checkmarkr\checkmarkc\checkmark$\checkmarku\checkmark$\checkmark^\checkmark\\checkmarkc\checkmarki\checkmarkr\checkmarkc\checkmark$\checkmarkb\checkmark$\checkmark^\checkmark\\checkmarkc\checkmarki\checkmarkr\checkmarkc\checkmark$\checkmarks\checkmark$\checkmark^\checkmark\\checkmarkc\checkmarki\checkmarkr\checkmarkc\checkmark$\checkmarke\checkmark$\checkmark^\checkmark\\checkmarkc\checkmarki\checkmarkr\checkmarkc\checkmark$\checkmarkc\checkmark$\checkmark^\checkmark\\checkmarkc\checkmarki\checkmarkr\checkmarkc\checkmark$\checkmarkt\checkmark$\checkmark^\checkmark\\checkmarkc\checkmarki\checkmarkr\checkmarkc\checkmark$\checkmarki\checkmark$\checkmark^\checkmark\\checkmarkc\checkmarki\checkmarkr\checkmarkc\checkmark$\checkmarko\checkmark$\checkmark^\checkmark\\checkmarkc\checkmarki\checkmarkr\checkmarkc\checkmark$\checkmarkn\checkmark$\checkmark^\checkmark\\checkmarkc\checkmarki\checkmarkr\checkmarkc\checkmark$\checkmark{\checkmark$\checkmark^\checkmark\\checkmarkc\checkmarki\checkmarkr\checkmarkc\checkmark$\checkmarkF\checkmark$\checkmark^\checkmark\\checkmarkc\checkmarki\checkmarkr\checkmarkc\checkmark$\checkmarko\checkmark$\checkmark^\checkmark\\checkmarkc\checkmarki\checkmarkr\checkmarkc\checkmark$\checkmarkr\checkmark$\checkmark^\checkmark\\checkmarkc\checkmarki\checkmarkr\checkmarkc\checkmark$\checkmarkc\checkmark$\checkmark^\checkmark\\checkmarkc\checkmarki\checkmarkr\checkmarkc\checkmark$\checkmarke\checkmark$\checkmark^\checkmark\\checkmarkc\checkmarki\checkmarkr\checkmarkc\checkmark$\checkmark \checkmark$\checkmark^\checkmark\\checkmarkc\checkmarki\checkmarkr\checkmarkc\checkmark$\checkmarkd\checkmark$\checkmark^\checkmark\\checkmarkc\checkmarki\checkmarkr\checkmarkc\checkmark$\checkmarke\checkmark$\checkmark^\checkmark\\checkmarkc\checkmarki\checkmarkr\checkmarkc\checkmark$\checkmark \checkmark$\checkmark^\checkmark\\checkmarkc\checkmarki\checkmarkr\checkmarkc\checkmark$\checkmarkc\checkmark$\checkmark^\checkmark\\checkmarkc\checkmarki\checkmarkr\checkmarkc\checkmark$\checkmarko\checkmark$\checkmark^\checkmark\\checkmarkc\checkmarki\checkmarkr\checkmarkc\checkmark$\checkmarku\checkmark$\checkmark^\checkmark\\checkmarkc\checkmarki\checkmarkr\checkmarkc\checkmark$\checkmarkp\checkmark$\checkmark^\checkmark\\checkmarkc\checkmarki\checkmarkr\checkmarkc\checkmark$\checkmarke\checkmark$\checkmark^\checkmark\\checkmarkc\checkmarki\checkmarkr\checkmarkc\checkmark$\checkmark \checkmark$\checkmark^\checkmark\\checkmarkc\checkmarki\checkmarkr\checkmarkc\checkmark$\checkmarkt\checkmark$\checkmark^\checkmark\\checkmarkc\checkmarki\checkmarkr\checkmarkc\checkmark$\checkmarka\checkmark$\checkmark^\checkmark\\checkmarkc\checkmarki\checkmarkr\checkmarkc\checkmark$\checkmarkn\checkmark$\checkmark^\checkmark\\checkmarkc\checkmarki\checkmarkr\checkmarkc\checkmark$\checkmarkg\checkmark$\checkmark^\checkmark\\checkmarkc\checkmarki\checkmarkr\checkmarkc\checkmark$\checkmarke\checkmark$\checkmark^\checkmark\\checkmarkc\checkmarki\checkmarkr\checkmarkc\checkmark$\checkmarkn\checkmark$\checkmark^\checkmark\\checkmarkc\checkmarki\checkmarkr\checkmarkc\checkmark$\checkmarkt\checkmark$\checkmark^\checkmark\\checkmarkc\checkmarki\checkmarkr\checkmarkc\checkmark$\checkmarki\checkmark$\checkmark^\checkmark\\checkmarkc\checkmarki\checkmarkr\checkmarkc\checkmark$\checkmarke\checkmark$\checkmark^\checkmark\\checkmarkc\checkmarki\checkmarkr\checkmarkc\checkmark$\checkmarkl\checkmark$\checkmark^\checkmark\\checkmarkc\checkmarki\checkmarkr\checkmarkc\checkmark$\checkmarkl\checkmark$\checkmark^\checkmark\\checkmarkc\checkmarki\checkmarkr\checkmarkc\checkmark$\checkmarke\checkmark$\checkmark^\checkmark\\checkmarkc\checkmarki\checkmarkr\checkmarkc\checkmark$\checkmark \checkmark$\checkmark^\checkmark\\checkmarkc\checkmarki\checkmarkr\checkmarkc\checkmark$\checkmark$\checkmark$\checkmark^\checkmark\\checkmarkc\checkmarki\checkmarkr\checkmarkc\checkmark$\checkmarkF\checkmark$\checkmark^\checkmark\\checkmarkc\checkmarki\checkmarkr\checkmarkc\checkmark$\checkmark_\checkmark$\checkmark^\checkmark\\checkmarkc\checkmarki\checkmarkr\checkmarkc\checkmark$\checkmarkc\checkmark$\checkmark^\checkmark\\checkmarkc\checkmarki\checkmarkr\checkmarkc\checkmark$\checkmark$\checkmark$\checkmark^\checkmark\\checkmarkc\checkmarki\checkmarkr\checkmarkc\checkmark$\checkmark \checkmark$\checkmark^\checkmark\\checkmarkc\checkmarki\checkmarkr\checkmarkc\checkmark$\checkmark(\checkmark$\checkmark^\checkmark\\checkmarkc\checkmarki\checkmarkr\checkmarkc\checkmark$\checkmarkf\checkmark$\checkmark^\checkmark\\checkmarkc\checkmarki\checkmarkr\checkmarkc\checkmark$\checkmarko\checkmark$\checkmark^\checkmark\\checkmarkc\checkmarki\checkmarkr\checkmarkc\checkmark$\checkmarkr\checkmark$\checkmark^\checkmark\\checkmarkc\checkmarki\checkmarkr\checkmarkc\checkmark$\checkmarkm\checkmark$\checkmark^\checkmark\\checkmarkc\checkmarki\checkmarkr\checkmarkc\checkmark$\checkmarku\checkmark$\checkmark^\checkmark\\checkmarkc\checkmarki\checkmarkr\checkmarkc\checkmark$\checkmarkl\checkmark$\checkmark^\checkmark\\checkmarkc\checkmarki\checkmarkr\checkmarkc\checkmark$\checkmarke\checkmark$\checkmark^\checkmark\\checkmarkc\checkmarki\checkmarkr\checkmarkc\checkmark$\checkmark \checkmark$\checkmark^\checkmark\\checkmarkc\checkmarki\checkmarkr\checkmarkc\checkmark$\checkmarkd\checkmark$\checkmark^\checkmark\\checkmarkc\checkmarki\checkmarkr\checkmarkc\checkmark$\checkmarke\checkmark$\checkmark^\checkmark\\checkmarkc\checkmarki\checkmarkr\checkmarkc\checkmark$\checkmark \checkmark$\checkmark^\checkmark\\checkmarkc\checkmarki\checkmarkr\checkmarkc\checkmark$\checkmarkK\checkmark$\checkmark^\checkmark\\checkmarkc\checkmarki\checkmarkr\checkmarkc\checkmark$\checkmarki\checkmark$\checkmark^\checkmark\\checkmarkc\checkmarki\checkmarkr\checkmarkc\checkmark$\checkmarke\checkmark$\checkmark^\checkmark\\checkmarkc\checkmarki\checkmarkr\checkmarkc\checkmark$\checkmarkn\checkmark$\checkmark^\checkmark\\checkmarkc\checkmarki\checkmarkr\checkmarkc\checkmark$\checkmarkz\checkmark$\checkmark^\checkmark\\checkmarkc\checkmarki\checkmarkr\checkmarkc\checkmark$\checkmarkl\checkmark$\checkmark^\checkmark\\checkmarkc\checkmarki\checkmarkr\checkmarkc\checkmark$\checkmarke\checkmark$\checkmark^\checkmark\\checkmarkc\checkmarki\checkmarkr\checkmarkc\checkmark$\checkmark)\checkmark$\checkmark^\checkmark\\checkmarkc\checkmarki\checkmarkr\checkmarkc\checkmark$\checkmark}\checkmark$\checkmark^\checkmark\\checkmarkc\checkmarki\checkmarkr\checkmarkc\checkmark$\checkmark
\checkmark$\checkmark^\checkmark\\checkmarkc\checkmarki\checkmarkr\checkmarkc\checkmark$\checkmarkL\checkmark$\checkmark^\checkmark\\checkmarkc\checkmarki\checkmarkr\checkmarkc\checkmark$\checkmarka\checkmark$\checkmark^\checkmark\\checkmarkc\checkmarki\checkmarkr\checkmarkc\checkmark$\checkmark \checkmark$\checkmark^\checkmark\\checkmarkc\checkmarki\checkmarkr\checkmarkc\checkmark$\checkmarkf\checkmark$\checkmark^\checkmark\\checkmarkc\checkmarki\checkmarkr\checkmarkc\checkmark$\checkmarko\checkmark$\checkmark^\checkmark\\checkmarkc\checkmarki\checkmarkr\checkmarkc\checkmark$\checkmarkr\checkmark$\checkmark^\checkmark\\checkmarkc\checkmarki\checkmarkr\checkmarkc\checkmark$\checkmarkc\checkmark$\checkmark^\checkmark\\checkmarkc\checkmarki\checkmarkr\checkmarkc\checkmark$\checkmarke\checkmark$\checkmark^\checkmark\\checkmarkc\checkmarki\checkmarkr\checkmarkc\checkmark$\checkmark \checkmark$\checkmark^\checkmark\\checkmarkc\checkmarki\checkmarkr\checkmarkc\checkmark$\checkmarkt\checkmark$\checkmark^\checkmark\\checkmarkc\checkmarki\checkmarkr\checkmarkc\checkmark$\checkmarka\checkmark$\checkmark^\checkmark\\checkmarkc\checkmarki\checkmarkr\checkmarkc\checkmark$\checkmarkn\checkmark$\checkmark^\checkmark\\checkmarkc\checkmarki\checkmarkr\checkmarkc\checkmark$\checkmarkg\checkmark$\checkmark^\checkmark\\checkmarkc\checkmarki\checkmarkr\checkmarkc\checkmark$\checkmarke\checkmark$\checkmark^\checkmark\\checkmarkc\checkmarki\checkmarkr\checkmarkc\checkmark$\checkmarkn\checkmark$\checkmark^\checkmark\\checkmarkc\checkmarki\checkmarkr\checkmarkc\checkmark$\checkmarkt\checkmark$\checkmark^\checkmark\\checkmarkc\checkmarki\checkmarkr\checkmarkc\checkmark$\checkmarki\checkmark$\checkmark^\checkmark\\checkmarkc\checkmarki\checkmarkr\checkmarkc\checkmark$\checkmarke\checkmark$\checkmark^\checkmark\\checkmarkc\checkmarki\checkmarkr\checkmarkc\checkmark$\checkmarkl\checkmark$\checkmark^\checkmark\\checkmarkc\checkmarki\checkmarkr\checkmarkc\checkmark$\checkmarkl\checkmark$\checkmark^\checkmark\\checkmarkc\checkmarki\checkmarkr\checkmarkc\checkmark$\checkmarke\checkmark$\checkmark^\checkmark\\checkmarkc\checkmarki\checkmarkr\checkmarkc\checkmark$\checkmark \checkmark$\checkmark^\checkmark\\checkmarkc\checkmarki\checkmarkr\checkmarkc\checkmark$\checkmarkp\checkmark$\checkmark^\checkmark\\checkmarkc\checkmarki\checkmarkr\checkmarkc\checkmark$\checkmarkr\checkmark$\checkmark^\checkmark\\checkmarkc\checkmarki\checkmarkr\checkmarkc\checkmark$\checkmarki\checkmark$\checkmark^\checkmark\\checkmarkc\checkmarki\checkmarkr\checkmarkc\checkmark$\checkmarkn\checkmark$\checkmark^\checkmark\\checkmarkc\checkmarki\checkmarkr\checkmarkc\checkmark$\checkmarkc\checkmark$\checkmark^\checkmark\\checkmarkc\checkmarki\checkmarkr\checkmarkc\checkmark$\checkmarki\checkmark$\checkmark^\checkmark\\checkmarkc\checkmarki\checkmarkr\checkmarkc\checkmark$\checkmarkp\checkmark$\checkmark^\checkmark\\checkmarkc\checkmarki\checkmarkr\checkmarkc\checkmark$\checkmarka\checkmark$\checkmark^\checkmark\\checkmarkc\checkmarki\checkmarkr\checkmarkc\checkmark$\checkmarkl\checkmark$\checkmark^\checkmark\\checkmarkc\checkmarki\checkmarkr\checkmarkc\checkmark$\checkmarke\checkmark$\checkmark^\checkmark\\checkmarkc\checkmarki\checkmarkr\checkmarkc\checkmark$\checkmark \checkmark$\checkmark^\checkmark\\checkmarkc\checkmarki\checkmarkr\checkmarkc\checkmark$\checkmarkd\checkmark$\checkmark^\checkmark\\checkmarkc\checkmarki\checkmarkr\checkmarkc\checkmark$\checkmark'\checkmark$\checkmark^\checkmark\\checkmarkc\checkmarki\checkmarkr\checkmarkc\checkmark$\checkmarku\checkmark$\checkmark^\checkmark\\checkmarkc\checkmarki\checkmarkr\checkmarkc\checkmark$\checkmarks\checkmark$\checkmark^\checkmark\\checkmarkc\checkmarki\checkmarkr\checkmarkc\checkmark$\checkmarki\checkmark$\checkmark^\checkmark\\checkmarkc\checkmarki\checkmarkr\checkmarkc\checkmark$\checkmarkn\checkmark$\checkmark^\checkmark\\checkmarkc\checkmarki\checkmarkr\checkmarkc\checkmark$\checkmarka\checkmark$\checkmark^\checkmark\\checkmarkc\checkmarki\checkmarkr\checkmarkc\checkmark$\checkmarkg\checkmark$\checkmark^\checkmark\\checkmarkc\checkmarki\checkmarkr\checkmarkc\checkmark$\checkmarke\checkmark$\checkmark^\checkmark\\checkmarkc\checkmarki\checkmarkr\checkmarkc\checkmark$\checkmark \checkmark$\checkmark^\checkmark\\checkmarkc\checkmarki\checkmarkr\checkmarkc\checkmark$\checkmarke\checkmark$\checkmark^\checkmark\\checkmarkc\checkmarki\checkmarkr\checkmarkc\checkmark$\checkmarks\checkmark$\checkmark^\checkmark\\checkmarkc\checkmarki\checkmarkr\checkmarkc\checkmark$\checkmarkt\checkmark$\checkmark^\checkmark\\checkmarkc\checkmarki\checkmarkr\checkmarkc\checkmark$\checkmark \checkmark$\checkmark^\checkmark\\checkmarkc\checkmarki\checkmarkr\checkmarkc\checkmark$\checkmarkc\checkmark$\checkmark^\checkmark\\checkmarkc\checkmarki\checkmarkr\checkmarkc\checkmark$\checkmarka\checkmark$\checkmark^\checkmark\\checkmarkc\checkmarki\checkmarkr\checkmarkc\checkmark$\checkmarkl\checkmark$\checkmark^\checkmark\\checkmarkc\checkmarki\checkmarkr\checkmarkc\checkmark$\checkmarkc\checkmark$\checkmark^\checkmark\\checkmarkc\checkmarki\checkmarkr\checkmarkc\checkmark$\checkmarku\checkmark$\checkmark^\checkmark\\checkmarkc\checkmarki\checkmarkr\checkmarkc\checkmark$\checkmarkl\checkmark$\checkmark^\checkmark\\checkmarkc\checkmarki\checkmarkr\checkmarkc\checkmark$\checkmarké\checkmark$\checkmark^\checkmark\\checkmarkc\checkmarki\checkmarkr\checkmarkc\checkmark$\checkmarke\checkmark$\checkmark^\checkmark\\checkmarkc\checkmarki\checkmarkr\checkmarkc\checkmark$\checkmark \checkmark$\checkmark^\checkmark\\checkmarkc\checkmarki\checkmarkr\checkmarkc\checkmark$\checkmarkp\checkmark$\checkmark^\checkmark\\checkmarkc\checkmarki\checkmarkr\checkmarkc\checkmark$\checkmarka\checkmark$\checkmark^\checkmark\\checkmarkc\checkmarki\checkmarkr\checkmarkc\checkmark$\checkmarkr\checkmark$\checkmark^\checkmark\\checkmarkc\checkmarki\checkmarkr\checkmarkc\checkmark$\checkmark \checkmark$\checkmark^\checkmark\\checkmarkc\checkmarki\checkmarkr\checkmarkc\checkmark$\checkmarkl\checkmark$\checkmark^\checkmark\\checkmarkc\checkmarki\checkmarkr\checkmarkc\checkmark$\checkmarka\checkmark$\checkmark^\checkmark\\checkmarkc\checkmarki\checkmarkr\checkmarkc\checkmark$\checkmark \checkmark$\checkmark^\checkmark\\checkmarkc\checkmarki\checkmarkr\checkmarkc\checkmark$\checkmarkf\checkmark$\checkmark^\checkmark\\checkmarkc\checkmarki\checkmarkr\checkmarkc\checkmark$\checkmarko\checkmark$\checkmark^\checkmark\\checkmarkc\checkmarki\checkmarkr\checkmarkc\checkmark$\checkmarkr\checkmark$\checkmark^\checkmark\\checkmarkc\checkmarki\checkmarkr\checkmarkc\checkmark$\checkmarkm\checkmark$\checkmark^\checkmark\\checkmarkc\checkmarki\checkmarkr\checkmarkc\checkmark$\checkmarku\checkmark$\checkmark^\checkmark\\checkmarkc\checkmarki\checkmarkr\checkmarkc\checkmark$\checkmarkl\checkmark$\checkmark^\checkmark\\checkmarkc\checkmarki\checkmarkr\checkmarkc\checkmark$\checkmarke\checkmark$\checkmark^\checkmark\\checkmarkc\checkmarki\checkmarkr\checkmarkc\checkmark$\checkmark \checkmark$\checkmark^\checkmark\\checkmarkc\checkmarki\checkmarkr\checkmarkc\checkmark$\checkmarke\checkmark$\checkmark^\checkmark\\checkmarkc\checkmarki\checkmarkr\checkmarkc\checkmark$\checkmarkm\checkmark$\checkmark^\checkmark\\checkmarkc\checkmarki\checkmarkr\checkmarkc\checkmark$\checkmarkp\checkmark$\checkmark^\checkmark\\checkmarkc\checkmarki\checkmarkr\checkmarkc\checkmark$\checkmarki\checkmark$\checkmark^\checkmark\\checkmarkc\checkmarki\checkmarkr\checkmarkc\checkmark$\checkmarkr\checkmark$\checkmark^\checkmark\\checkmarkc\checkmarki\checkmarkr\checkmarkc\checkmark$\checkmarki\checkmark$\checkmark^\checkmark\\checkmarkc\checkmarki\checkmarkr\checkmarkc\checkmark$\checkmarkq\checkmark$\checkmark^\checkmark\\checkmarkc\checkmarki\checkmarkr\checkmarkc\checkmark$\checkmarku\checkmark$\checkmark^\checkmark\\checkmarkc\checkmarki\checkmarkr\checkmarkc\checkmark$\checkmarke\checkmark$\checkmark^\checkmark\\checkmarkc\checkmarki\checkmarkr\checkmarkc\checkmark$\checkmark \checkmark$\checkmark^\checkmark\\checkmarkc\checkmarki\checkmarkr\checkmarkc\checkmark$\checkmarkd\checkmark$\checkmark^\checkmark\\checkmarkc\checkmarki\checkmarkr\checkmarkc\checkmark$\checkmarke\checkmark$\checkmark^\checkmark\\checkmarkc\checkmarki\checkmarkr\checkmarkc\checkmark$\checkmark \checkmark$\checkmark^\checkmark\\checkmarkc\checkmarki\checkmarkr\checkmarkc\checkmark$\checkmarkK\checkmark$\checkmark^\checkmark\\checkmarkc\checkmarki\checkmarkr\checkmarkc\checkmark$\checkmarki\checkmark$\checkmark^\checkmark\\checkmarkc\checkmarki\checkmarkr\checkmarkc\checkmark$\checkmarke\checkmark$\checkmark^\checkmark\\checkmarkc\checkmarki\checkmarkr\checkmarkc\checkmark$\checkmarkn\checkmark$\checkmark^\checkmark\\checkmarkc\checkmarki\checkmarkr\checkmarkc\checkmark$\checkmarkz\checkmark$\checkmark^\checkmark\\checkmarkc\checkmarki\checkmarkr\checkmarkc\checkmark$\checkmarkl\checkmark$\checkmark^\checkmark\\checkmarkc\checkmarki\checkmarkr\checkmarkc\checkmark$\checkmarke\checkmark$\checkmark^\checkmark\\checkmarkc\checkmarki\checkmarkr\checkmarkc\checkmark$\checkmark \checkmark$\checkmark^\checkmark\\checkmarkc\checkmarki\checkmarkr\checkmarkc\checkmark$\checkmark:\checkmark$\checkmark^\checkmark\\checkmarkc\checkmarki\checkmarkr\checkmarkc\checkmark$\checkmark
\checkmark$\checkmark^\checkmark\\checkmarkc\checkmarki\checkmarkr\checkmarkc\checkmark$\checkmark\\checkmark$\checkmark^\checkmark\\checkmarkc\checkmarki\checkmarkr\checkmarkc\checkmark$\checkmarkb\checkmark$\checkmark^\checkmark\\checkmarkc\checkmarki\checkmarkr\checkmarkc\checkmark$\checkmarke\checkmark$\checkmark^\checkmark\\checkmarkc\checkmarki\checkmarkr\checkmarkc\checkmark$\checkmarkg\checkmark$\checkmark^\checkmark\\checkmarkc\checkmarki\checkmarkr\checkmarkc\checkmark$\checkmarki\checkmark$\checkmark^\checkmark\\checkmarkc\checkmarki\checkmarkr\checkmarkc\checkmark$\checkmarkn\checkmark$\checkmark^\checkmark\\checkmarkc\checkmarki\checkmarkr\checkmarkc\checkmark$\checkmark{\checkmark$\checkmark^\checkmark\\checkmarkc\checkmarki\checkmarkr\checkmarkc\checkmark$\checkmarke\checkmark$\checkmark^\checkmark\\checkmarkc\checkmarki\checkmarkr\checkmarkc\checkmark$\checkmarkq\checkmark$\checkmark^\checkmark\\checkmarkc\checkmarki\checkmarkr\checkmarkc\checkmark$\checkmarku\checkmark$\checkmark^\checkmark\\checkmarkc\checkmarki\checkmarkr\checkmarkc\checkmark$\checkmarka\checkmark$\checkmark^\checkmark\\checkmarkc\checkmarki\checkmarkr\checkmarkc\checkmark$\checkmarkt\checkmark$\checkmark^\checkmark\\checkmarkc\checkmarki\checkmarkr\checkmarkc\checkmark$\checkmarki\checkmark$\checkmark^\checkmark\\checkmarkc\checkmarki\checkmarkr\checkmarkc\checkmark$\checkmarko\checkmark$\checkmark^\checkmark\\checkmarkc\checkmarki\checkmarkr\checkmarkc\checkmark$\checkmarkn\checkmark$\checkmark^\checkmark\\checkmarkc\checkmarki\checkmarkr\checkmarkc\checkmark$\checkmark}\checkmark$\checkmark^\checkmark\\checkmarkc\checkmarki\checkmarkr\checkmarkc\checkmark$\checkmark
\checkmark$\checkmark^\checkmark\\checkmarkc\checkmarki\checkmarkr\checkmarkc\checkmark$\checkmark \checkmark$\checkmark^\checkmark\\checkmarkc\checkmarki\checkmarkr\checkmarkc\checkmark$\checkmark \checkmark$\checkmark^\checkmark\\checkmarkc\checkmarki\checkmarkr\checkmarkc\checkmark$\checkmark \checkmark$\checkmark^\checkmark\\checkmarkc\checkmarki\checkmarkr\checkmarkc\checkmark$\checkmark \checkmark$\checkmark^\checkmark\\checkmarkc\checkmarki\checkmarkr\checkmarkc\checkmark$\checkmarkF\checkmark$\checkmark^\checkmark\\checkmarkc\checkmarki\checkmarkr\checkmarkc\checkmark$\checkmark_\checkmark$\checkmark^\checkmark\\checkmarkc\checkmarki\checkmarkr\checkmarkc\checkmark$\checkmarkc\checkmark$\checkmark^\checkmark\\checkmarkc\checkmarki\checkmarkr\checkmarkc\checkmark$\checkmark \checkmark$\checkmark^\checkmark\\checkmarkc\checkmarki\checkmarkr\checkmarkc\checkmark$\checkmark=\checkmark$\checkmark^\checkmark\\checkmarkc\checkmarki\checkmarkr\checkmarkc\checkmark$\checkmark \checkmark$\checkmark^\checkmark\\checkmarkc\checkmarki\checkmarkr\checkmarkc\checkmark$\checkmarkk\checkmark$\checkmark^\checkmark\\checkmarkc\checkmarki\checkmarkr\checkmarkc\checkmark$\checkmark_\checkmark$\checkmark^\checkmark\\checkmarkc\checkmarki\checkmarkr\checkmarkc\checkmark$\checkmarkc\checkmark$\checkmark^\checkmark\\checkmarkc\checkmarki\checkmarkr\checkmarkc\checkmark$\checkmark \checkmark$\checkmark^\checkmark\\checkmarkc\checkmarki\checkmarkr\checkmarkc\checkmark$\checkmark\\checkmark$\checkmark^\checkmark\\checkmarkc\checkmarki\checkmarkr\checkmarkc\checkmark$\checkmarkc\checkmark$\checkmark^\checkmark\\checkmarkc\checkmarki\checkmarkr\checkmarkc\checkmark$\checkmarkd\checkmark$\checkmark^\checkmark\\checkmarkc\checkmarki\checkmarkr\checkmarkc\checkmark$\checkmarko\checkmark$\checkmark^\checkmark\\checkmarkc\checkmarki\checkmarkr\checkmarkc\checkmark$\checkmarkt\checkmark$\checkmark^\checkmark\\checkmarkc\checkmarki\checkmarkr\checkmarkc\checkmark$\checkmark \checkmark$\checkmark^\checkmark\\checkmarkc\checkmarki\checkmarkr\checkmarkc\checkmark$\checkmarka\checkmark$\checkmark^\checkmark\\checkmarkc\checkmarki\checkmarkr\checkmarkc\checkmark$\checkmark_\checkmark$\checkmark^\checkmark\\checkmarkc\checkmarki\checkmarkr\checkmarkc\checkmark$\checkmarkp\checkmark$\checkmark^\checkmark\\checkmarkc\checkmarki\checkmarkr\checkmarkc\checkmark$\checkmark \checkmark$\checkmark^\checkmark\\checkmarkc\checkmarki\checkmarkr\checkmarkc\checkmark$\checkmark\\checkmark$\checkmark^\checkmark\\checkmarkc\checkmarki\checkmarkr\checkmarkc\checkmark$\checkmarkc\checkmark$\checkmark^\checkmark\\checkmarkc\checkmarki\checkmarkr\checkmarkc\checkmark$\checkmarkd\checkmark$\checkmark^\checkmark\\checkmarkc\checkmarki\checkmarkr\checkmarkc\checkmark$\checkmarko\checkmark$\checkmark^\checkmark\\checkmarkc\checkmarki\checkmarkr\checkmarkc\checkmark$\checkmarkt\checkmark$\checkmark^\checkmark\\checkmarkc\checkmarki\checkmarkr\checkmarkc\checkmark$\checkmark \checkmark$\checkmark^\checkmark\\checkmarkc\checkmarki\checkmarkr\checkmarkc\checkmark$\checkmarkf\checkmark$\checkmark^\checkmark\\checkmarkc\checkmarki\checkmarkr\checkmarkc\checkmark$\checkmark_\checkmark$\checkmark^\checkmark\\checkmarkc\checkmarki\checkmarkr\checkmarkc\checkmark$\checkmarkz\checkmark$\checkmark^\checkmark\\checkmarkc\checkmarki\checkmarkr\checkmarkc\checkmark$\checkmark \checkmark$\checkmark^\checkmark\\checkmarkc\checkmarki\checkmarkr\checkmarkc\checkmark$\checkmark\\checkmark$\checkmark^\checkmark\\checkmarkc\checkmarki\checkmarkr\checkmarkc\checkmark$\checkmarkc\checkmark$\checkmark^\checkmark\\checkmarkc\checkmarki\checkmarkr\checkmarkc\checkmark$\checkmarkd\checkmark$\checkmark^\checkmark\\checkmarkc\checkmarki\checkmarkr\checkmarkc\checkmark$\checkmarko\checkmark$\checkmark^\checkmark\\checkmarkc\checkmarki\checkmarkr\checkmarkc\checkmark$\checkmarkt\checkmark$\checkmark^\checkmark\\checkmarkc\checkmarki\checkmarkr\checkmarkc\checkmark$\checkmark \checkmark$\checkmark^\checkmark\\checkmarkc\checkmarki\checkmarkr\checkmarkc\checkmark$\checkmarkZ\checkmark$\checkmark^\checkmark\\checkmarkc\checkmarki\checkmarkr\checkmarkc\checkmark$\checkmark_\checkmark$\checkmark^\checkmark\\checkmarkc\checkmarki\checkmarkr\checkmarkc\checkmark$\checkmarkc\checkmark$\checkmark^\checkmark\\checkmarkc\checkmarki\checkmarkr\checkmarkc\checkmark$\checkmark \checkmark$\checkmark^\checkmark\\checkmarkc\checkmarki\checkmarkr\checkmarkc\checkmark$\checkmark=\checkmark$\checkmark^\checkmark\\checkmarkc\checkmarki\checkmarkr\checkmarkc\checkmark$\checkmark \checkmark$\checkmark^\checkmark\\checkmarkc\checkmarki\checkmarkr\checkmarkc\checkmark$\checkmark7\checkmark$\checkmark^\checkmark\\checkmarkc\checkmarki\checkmarkr\checkmarkc\checkmark$\checkmark0\checkmark$\checkmark^\checkmark\\checkmarkc\checkmarki\checkmarkr\checkmarkc\checkmark$\checkmark0\checkmark$\checkmark^\checkmark\\checkmarkc\checkmarki\checkmarkr\checkmarkc\checkmark$\checkmark \checkmark$\checkmark^\checkmark\\checkmarkc\checkmarki\checkmarkr\checkmarkc\checkmark$\checkmark\\checkmark$\checkmark^\checkmark\\checkmarkc\checkmarki\checkmarkr\checkmarkc\checkmark$\checkmarkt\checkmark$\checkmark^\checkmark\\checkmarkc\checkmarki\checkmarkr\checkmarkc\checkmark$\checkmarki\checkmark$\checkmark^\checkmark\\checkmarkc\checkmarki\checkmarkr\checkmarkc\checkmark$\checkmarkm\checkmark$\checkmark^\checkmark\\checkmarkc\checkmarki\checkmarkr\checkmarkc\checkmark$\checkmarke\checkmark$\checkmark^\checkmark\\checkmarkc\checkmarki\checkmarkr\checkmarkc\checkmark$\checkmarks\checkmark$\checkmark^\checkmark\\checkmarkc\checkmarki\checkmarkr\checkmarkc\checkmark$\checkmark \checkmark$\checkmark^\checkmark\\checkmarkc\checkmarki\checkmarkr\checkmarkc\checkmark$\checkmark1\checkmark$\checkmark^\checkmark\\checkmarkc\checkmarki\checkmarkr\checkmarkc\checkmark$\checkmark \checkmark$\checkmark^\checkmark\\checkmarkc\checkmarki\checkmarkr\checkmarkc\checkmark$\checkmark\\checkmark$\checkmark^\checkmark\\checkmarkc\checkmarki\checkmarkr\checkmarkc\checkmark$\checkmarkt\checkmark$\checkmark^\checkmark\\checkmarkc\checkmarki\checkmarkr\checkmarkc\checkmark$\checkmarki\checkmark$\checkmark^\checkmark\\checkmarkc\checkmarki\checkmarkr\checkmarkc\checkmark$\checkmarkm\checkmark$\checkmark^\checkmark\\checkmarkc\checkmarki\checkmarkr\checkmarkc\checkmark$\checkmarke\checkmark$\checkmark^\checkmark\\checkmarkc\checkmarki\checkmarkr\checkmarkc\checkmark$\checkmarks\checkmark$\checkmark^\checkmark\\checkmarkc\checkmarki\checkmarkr\checkmarkc\checkmark$\checkmark \checkmark$\checkmark^\checkmark\\checkmarkc\checkmarki\checkmarkr\checkmarkc\checkmark$\checkmark0\checkmark$\checkmark^\checkmark\\checkmarkc\checkmarki\checkmarkr\checkmarkc\checkmark$\checkmark{\checkmark$\checkmark^\checkmark\\checkmarkc\checkmarki\checkmarkr\checkmarkc\checkmark$\checkmark,\checkmark$\checkmark^\checkmark\\checkmarkc\checkmarki\checkmarkr\checkmarkc\checkmark$\checkmark}\checkmark$\checkmark^\checkmark\\checkmarkc\checkmarki\checkmarkr\checkmarkc\checkmark$\checkmark0\checkmark$\checkmark^\checkmark\\checkmarkc\checkmarki\checkmarkr\checkmarkc\checkmark$\checkmark5\checkmark$\checkmark^\checkmark\\checkmarkc\checkmarki\checkmarkr\checkmarkc\checkmark$\checkmark \checkmark$\checkmark^\checkmark\\checkmarkc\checkmarki\checkmarkr\checkmarkc\checkmark$\checkmark\\checkmark$\checkmark^\checkmark\\checkmarkc\checkmarki\checkmarkr\checkmarkc\checkmark$\checkmarkt\checkmark$\checkmark^\checkmark\\checkmarkc\checkmarki\checkmarkr\checkmarkc\checkmark$\checkmarki\checkmark$\checkmark^\checkmark\\checkmarkc\checkmarki\checkmarkr\checkmarkc\checkmark$\checkmarkm\checkmark$\checkmark^\checkmark\\checkmarkc\checkmarki\checkmarkr\checkmarkc\checkmark$\checkmarke\checkmark$\checkmark^\checkmark\\checkmarkc\checkmarki\checkmarkr\checkmarkc\checkmark$\checkmarks\checkmark$\checkmark^\checkmark\\checkmarkc\checkmarki\checkmarkr\checkmarkc\checkmark$\checkmark \checkmark$\checkmark^\checkmark\\checkmarkc\checkmarki\checkmarkr\checkmarkc\checkmark$\checkmark2\checkmark$\checkmark^\checkmark\\checkmarkc\checkmarki\checkmarkr\checkmarkc\checkmark$\checkmark \checkmark$\checkmark^\checkmark\\checkmarkc\checkmarki\checkmarkr\checkmarkc\checkmark$\checkmark=\checkmark$\checkmark^\checkmark\\checkmarkc\checkmarki\checkmarkr\checkmarkc\checkmark$\checkmark \checkmark$\checkmark^\checkmark\\checkmarkc\checkmarki\checkmarkr\checkmarkc\checkmark$\checkmark\\checkmark$\checkmark^\checkmark\\checkmarkc\checkmarki\checkmarkr\checkmarkc\checkmark$\checkmarkb\checkmark$\checkmark^\checkmark\\checkmarkc\checkmarki\checkmarkr\checkmarkc\checkmark$\checkmarko\checkmark$\checkmark^\checkmark\\checkmarkc\checkmarki\checkmarkr\checkmarkc\checkmark$\checkmarkx\checkmark$\checkmark^\checkmark\\checkmarkc\checkmarki\checkmarkr\checkmarkc\checkmark$\checkmarke\checkmark$\checkmark^\checkmark\\checkmarkc\checkmarki\checkmarkr\checkmarkc\checkmark$\checkmarkd\checkmark$\checkmark^\checkmark\\checkmarkc\checkmarki\checkmarkr\checkmarkc\checkmark$\checkmark{\checkmark$\checkmark^\checkmark\\checkmarkc\checkmarki\checkmarkr\checkmarkc\checkmark$\checkmark7\checkmark$\checkmark^\checkmark\\checkmarkc\checkmarki\checkmarkr\checkmarkc\checkmark$\checkmark0\checkmark$\checkmark^\checkmark\\checkmarkc\checkmarki\checkmarkr\checkmarkc\checkmark$\checkmark \checkmark$\checkmark^\checkmark\\checkmarkc\checkmarki\checkmarkr\checkmarkc\checkmark$\checkmark\\checkmark$\checkmark^\checkmark\\checkmarkc\checkmarki\checkmarkr\checkmarkc\checkmark$\checkmarkt\checkmark$\checkmark^\checkmark\\checkmarkc\checkmarki\checkmarkr\checkmarkc\checkmark$\checkmarke\checkmark$\checkmark^\checkmark\\checkmarkc\checkmarki\checkmarkr\checkmarkc\checkmark$\checkmarkx\checkmark$\checkmark^\checkmark\\checkmarkc\checkmarki\checkmarkr\checkmarkc\checkmark$\checkmarkt\checkmark$\checkmark^\checkmark\\checkmarkc\checkmarki\checkmarkr\checkmarkc\checkmark$\checkmark{\checkmark$\checkmark^\checkmark\\checkmarkc\checkmarki\checkmarkr\checkmarkc\checkmark$\checkmark \checkmark$\checkmark^\checkmark\\checkmarkc\checkmarki\checkmarkr\checkmarkc\checkmark$\checkmarkN\checkmark$\checkmark^\checkmark\\checkmarkc\checkmarki\checkmarkr\checkmarkc\checkmark$\checkmark}\checkmark$\checkmark^\checkmark\\checkmarkc\checkmarki\checkmarkr\checkmarkc\checkmark$\checkmark}\checkmark$\checkmark^\checkmark\\checkmarkc\checkmarki\checkmarkr\checkmarkc\checkmark$\checkmark
\checkmark$\checkmark^\checkmark\\checkmarkc\checkmarki\checkmarkr\checkmarkc\checkmark$\checkmark \checkmark$\checkmark^\checkmark\\checkmarkc\checkmarki\checkmarkr\checkmarkc\checkmark$\checkmark \checkmark$\checkmark^\checkmark\\checkmarkc\checkmarki\checkmarkr\checkmarkc\checkmark$\checkmark \checkmark$\checkmark^\checkmark\\checkmarkc\checkmarki\checkmarkr\checkmarkc\checkmark$\checkmark \checkmark$\checkmark^\checkmark\\checkmarkc\checkmarki\checkmarkr\checkmarkc\checkmark$\checkmark\\checkmark$\checkmark^\checkmark\\checkmarkc\checkmarki\checkmarkr\checkmarkc\checkmark$\checkmarkl\checkmark$\checkmark^\checkmark\\checkmarkc\checkmarki\checkmarkr\checkmarkc\checkmark$\checkmarka\checkmark$\checkmark^\checkmark\\checkmarkc\checkmarki\checkmarkr\checkmarkc\checkmark$\checkmarkb\checkmark$\checkmark^\checkmark\\checkmarkc\checkmarki\checkmarkr\checkmarkc\checkmark$\checkmarke\checkmark$\checkmark^\checkmark\\checkmarkc\checkmarki\checkmarkr\checkmarkc\checkmark$\checkmarkl\checkmark$\checkmark^\checkmark\\checkmarkc\checkmarki\checkmarkr\checkmarkc\checkmark$\checkmark{\checkmark$\checkmark^\checkmark\\checkmarkc\checkmarki\checkmarkr\checkmarkc\checkmark$\checkmarke\checkmark$\checkmark^\checkmark\\checkmarkc\checkmarki\checkmarkr\checkmarkc\checkmark$\checkmarkq\checkmark$\checkmark^\checkmark\\checkmarkc\checkmarki\checkmarkr\checkmarkc\checkmark$\checkmark:\checkmark$\checkmark^\checkmark\\checkmarkc\checkmarki\checkmarkr\checkmarkc\checkmark$\checkmarkf\checkmark$\checkmark^\checkmark\\checkmarkc\checkmarki\checkmarkr\checkmarkc\checkmark$\checkmarko\checkmark$\checkmark^\checkmark\\checkmarkc\checkmarki\checkmarkr\checkmarkc\checkmark$\checkmarkr\checkmark$\checkmark^\checkmark\\checkmarkc\checkmarki\checkmarkr\checkmarkc\checkmark$\checkmarkc\checkmark$\checkmark^\checkmark\\checkmarkc\checkmarki\checkmarkr\checkmarkc\checkmark$\checkmarke\checkmark$\checkmark^\checkmark\\checkmarkc\checkmarki\checkmarkr\checkmarkc\checkmark$\checkmark_\checkmark$\checkmark^\checkmark\\checkmarkc\checkmarki\checkmarkr\checkmarkc\checkmark$\checkmarkc\checkmark$\checkmark^\checkmark\\checkmarkc\checkmarki\checkmarkr\checkmarkc\checkmark$\checkmarko\checkmark$\checkmark^\checkmark\\checkmarkc\checkmarki\checkmarkr\checkmarkc\checkmark$\checkmarku\checkmark$\checkmark^\checkmark\\checkmarkc\checkmarki\checkmarkr\checkmarkc\checkmark$\checkmarkp\checkmark$\checkmark^\checkmark\\checkmarkc\checkmarki\checkmarkr\checkmarkc\checkmark$\checkmarke\checkmark$\checkmark^\checkmark\\checkmarkc\checkmarki\checkmarkr\checkmarkc\checkmark$\checkmark}\checkmark$\checkmark^\checkmark\\checkmarkc\checkmarki\checkmarkr\checkmarkc\checkmark$\checkmark
\checkmark$\checkmark^\checkmark\\checkmarkc\checkmarki\checkmarkr\checkmarkc\checkmark$\checkmark\\checkmark$\checkmark^\checkmark\\checkmarkc\checkmarki\checkmarkr\checkmarkc\checkmark$\checkmarke\checkmark$\checkmark^\checkmark\\checkmarkc\checkmarki\checkmarkr\checkmarkc\checkmark$\checkmarkn\checkmark$\checkmark^\checkmark\\checkmarkc\checkmarki\checkmarkr\checkmarkc\checkmark$\checkmarkd\checkmark$\checkmark^\checkmark\\checkmarkc\checkmarki\checkmarkr\checkmarkc\checkmark$\checkmark{\checkmark$\checkmark^\checkmark\\checkmarkc\checkmarki\checkmarkr\checkmarkc\checkmark$\checkmarke\checkmark$\checkmark^\checkmark\\checkmarkc\checkmarki\checkmarkr\checkmarkc\checkmark$\checkmarkq\checkmark$\checkmark^\checkmark\\checkmarkc\checkmarki\checkmarkr\checkmarkc\checkmark$\checkmarku\checkmark$\checkmark^\checkmark\\checkmarkc\checkmarki\checkmarkr\checkmarkc\checkmark$\checkmarka\checkmark$\checkmark^\checkmark\\checkmarkc\checkmarki\checkmarkr\checkmarkc\checkmark$\checkmarkt\checkmark$\checkmark^\checkmark\\checkmarkc\checkmarki\checkmarkr\checkmarkc\checkmark$\checkmarki\checkmark$\checkmark^\checkmark\\checkmarkc\checkmarki\checkmarkr\checkmarkc\checkmark$\checkmarko\checkmark$\checkmark^\checkmark\\checkmarkc\checkmarki\checkmarkr\checkmarkc\checkmark$\checkmarkn\checkmark$\checkmark^\checkmark\\checkmarkc\checkmarki\checkmarkr\checkmarkc\checkmark$\checkmark}\checkmark$\checkmark^\checkmark\\checkmarkc\checkmarki\checkmarkr\checkmarkc\checkmark$\checkmark
\checkmark$\checkmark^\checkmark\\checkmarkc\checkmarki\checkmarkr\checkmarkc\checkmark$\checkmark
\checkmark$\checkmark^\checkmark\\checkmarkc\checkmarki\checkmarkr\checkmarkc\checkmark$\checkmark\\checkmark$\checkmark^\checkmark\\checkmarkc\checkmarki\checkmarkr\checkmarkc\checkmark$\checkmarks\checkmark$\checkmark^\checkmark\\checkmarkc\checkmarki\checkmarkr\checkmarkc\checkmark$\checkmarku\checkmark$\checkmark^\checkmark\\checkmarkc\checkmarki\checkmarkr\checkmarkc\checkmark$\checkmarkb\checkmark$\checkmark^\checkmark\\checkmarkc\checkmarki\checkmarkr\checkmarkc\checkmark$\checkmarks\checkmark$\checkmark^\checkmark\\checkmarkc\checkmarki\checkmarkr\checkmarkc\checkmark$\checkmarku\checkmark$\checkmark^\checkmark\\checkmarkc\checkmarki\checkmarkr\checkmarkc\checkmark$\checkmarkb\checkmark$\checkmark^\checkmark\\checkmarkc\checkmarki\checkmarkr\checkmarkc\checkmark$\checkmarks\checkmark$\checkmark^\checkmark\\checkmarkc\checkmarki\checkmarkr\checkmarkc\checkmark$\checkmarke\checkmark$\checkmark^\checkmark\\checkmarkc\checkmarki\checkmarkr\checkmarkc\checkmark$\checkmarkc\checkmark$\checkmark^\checkmark\\checkmarkc\checkmarki\checkmarkr\checkmarkc\checkmark$\checkmarkt\checkmark$\checkmark^\checkmark\\checkmarkc\checkmarki\checkmarkr\checkmarkc\checkmark$\checkmarki\checkmark$\checkmark^\checkmark\\checkmarkc\checkmarki\checkmarkr\checkmarkc\checkmark$\checkmarko\checkmark$\checkmark^\checkmark\\checkmarkc\checkmarki\checkmarkr\checkmarkc\checkmark$\checkmarkn\checkmark$\checkmark^\checkmark\\checkmarkc\checkmarki\checkmarkr\checkmarkc\checkmark$\checkmark{\checkmark$\checkmark^\checkmark\\checkmarkc\checkmarki\checkmarkr\checkmarkc\checkmark$\checkmarkF\checkmark$\checkmark^\checkmark\\checkmarkc\checkmarki\checkmarkr\checkmarkc\checkmark$\checkmarko\checkmark$\checkmark^\checkmark\\checkmarkc\checkmarki\checkmarkr\checkmarkc\checkmark$\checkmarkr\checkmark$\checkmark^\checkmark\\checkmarkc\checkmarki\checkmarkr\checkmarkc\checkmark$\checkmarkc\checkmark$\checkmark^\checkmark\\checkmarkc\checkmarki\checkmarkr\checkmarkc\checkmark$\checkmarke\checkmark$\checkmark^\checkmark\\checkmarkc\checkmarki\checkmarkr\checkmarkc\checkmark$\checkmark \checkmark$\checkmark^\checkmark\\checkmarkc\checkmarki\checkmarkr\checkmarkc\checkmark$\checkmarkd\checkmark$\checkmark^\checkmark\\checkmarkc\checkmarki\checkmarkr\checkmarkc\checkmark$\checkmark'\checkmark$\checkmark^\checkmark\\checkmarkc\checkmarki\checkmarkr\checkmarkc\checkmark$\checkmarka\checkmark$\checkmark^\checkmark\\checkmarkc\checkmarki\checkmarkr\checkmarkc\checkmark$\checkmarkv\checkmark$\checkmark^\checkmark\\checkmarkc\checkmarki\checkmarkr\checkmarkc\checkmark$\checkmarka\checkmark$\checkmark^\checkmark\\checkmarkc\checkmarki\checkmarkr\checkmarkc\checkmark$\checkmarkn\checkmark$\checkmark^\checkmark\\checkmarkc\checkmarki\checkmarkr\checkmarkc\checkmark$\checkmarkc\checkmark$\checkmark^\checkmark\\checkmarkc\checkmarki\checkmarkr\checkmarkc\checkmark$\checkmarke\checkmark$\checkmark^\checkmark\\checkmarkc\checkmarki\checkmarkr\checkmarkc\checkmark$\checkmark \checkmark$\checkmark^\checkmark\\checkmarkc\checkmarki\checkmarkr\checkmarkc\checkmark$\checkmark$\checkmark$\checkmark^\checkmark\\checkmarkc\checkmarki\checkmarkr\checkmarkc\checkmark$\checkmarkF\checkmark$\checkmark^\checkmark\\checkmarkc\checkmarki\checkmarkr\checkmarkc\checkmark$\checkmark_\checkmark$\checkmark^\checkmark\\checkmarkc\checkmarki\checkmarkr\checkmarkc\checkmark$\checkmarkf\checkmark$\checkmark^\checkmark\\checkmarkc\checkmarki\checkmarkr\checkmarkc\checkmark$\checkmark$\checkmark$\checkmark^\checkmark\\checkmarkc\checkmarki\checkmarkr\checkmarkc\checkmark$\checkmark \checkmark$\checkmark^\checkmark\\checkmarkc\checkmarki\checkmarkr\checkmarkc\checkmark$\checkmarke\checkmark$\checkmark^\checkmark\\checkmarkc\checkmarki\checkmarkr\checkmarkc\checkmark$\checkmarkt\checkmark$\checkmark^\checkmark\\checkmarkc\checkmarki\checkmarkr\checkmarkc\checkmark$\checkmark \checkmark$\checkmark^\checkmark\\checkmarkc\checkmarki\checkmarkr\checkmarkc\checkmark$\checkmarkf\checkmark$\checkmark^\checkmark\\checkmarkc\checkmarki\checkmarkr\checkmarkc\checkmark$\checkmarko\checkmark$\checkmark^\checkmark\\checkmarkc\checkmarki\checkmarkr\checkmarkc\checkmark$\checkmarkr\checkmark$\checkmark^\checkmark\\checkmarkc\checkmarki\checkmarkr\checkmarkc\checkmark$\checkmarkc\checkmark$\checkmark^\checkmark\\checkmarkc\checkmarki\checkmarkr\checkmarkc\checkmark$\checkmarke\checkmark$\checkmark^\checkmark\\checkmarkc\checkmarki\checkmarkr\checkmarkc\checkmark$\checkmark \checkmark$\checkmark^\checkmark\\checkmarkc\checkmarki\checkmarkr\checkmarkc\checkmark$\checkmarkd\checkmark$\checkmark^\checkmark\\checkmarkc\checkmarki\checkmarkr\checkmarkc\checkmark$\checkmarke\checkmark$\checkmark^\checkmark\\checkmarkc\checkmarki\checkmarkr\checkmarkc\checkmark$\checkmark \checkmark$\checkmark^\checkmark\\checkmarkc\checkmarki\checkmarkr\checkmarkc\checkmark$\checkmarkp\checkmark$\checkmark^\checkmark\\checkmarkc\checkmarki\checkmarkr\checkmarkc\checkmark$\checkmarké\checkmark$\checkmark^\checkmark\\checkmarkc\checkmarki\checkmarkr\checkmarkc\checkmark$\checkmarkn\checkmark$\checkmark^\checkmark\\checkmarkc\checkmarki\checkmarkr\checkmarkc\checkmark$\checkmarké\checkmark$\checkmark^\checkmark\\checkmarkc\checkmarki\checkmarkr\checkmarkc\checkmark$\checkmarkt\checkmark$\checkmark^\checkmark\\checkmarkc\checkmarki\checkmarkr\checkmarkc\checkmark$\checkmarkr\checkmark$\checkmark^\checkmark\\checkmarkc\checkmarki\checkmarkr\checkmarkc\checkmark$\checkmarka\checkmark$\checkmark^\checkmark\\checkmarkc\checkmarki\checkmarkr\checkmarkc\checkmark$\checkmarkt\checkmark$\checkmark^\checkmark\\checkmarkc\checkmarki\checkmarkr\checkmarkc\checkmark$\checkmarki\checkmark$\checkmark^\checkmark\\checkmarkc\checkmarki\checkmarkr\checkmarkc\checkmark$\checkmarko\checkmark$\checkmark^\checkmark\\checkmarkc\checkmarki\checkmarkr\checkmarkc\checkmark$\checkmarkn\checkmark$\checkmark^\checkmark\\checkmarkc\checkmarki\checkmarkr\checkmarkc\checkmark$\checkmark \checkmark$\checkmark^\checkmark\\checkmarkc\checkmarki\checkmarkr\checkmarkc\checkmark$\checkmark$\checkmark$\checkmark^\checkmark\\checkmarkc\checkmarki\checkmarkr\checkmarkc\checkmark$\checkmarkF\checkmark$\checkmark^\checkmark\\checkmarkc\checkmarki\checkmarkr\checkmarkc\checkmark$\checkmark_\checkmark$\checkmark^\checkmark\\checkmarkc\checkmarki\checkmarkr\checkmarkc\checkmark$\checkmarkp\checkmark$\checkmark^\checkmark\\checkmarkc\checkmarki\checkmarkr\checkmarkc\checkmark$\checkmark$\checkmark$\checkmark^\checkmark\\checkmarkc\checkmarki\checkmarkr\checkmarkc\checkmark$\checkmark}\checkmark$\checkmark^\checkmark\\checkmarkc\checkmarki\checkmarkr\checkmarkc\checkmark$\checkmark
\checkmark$\checkmark^\checkmark\\checkmarkc\checkmarki\checkmarkr\checkmarkc\checkmark$\checkmarkD\checkmark$\checkmark^\checkmark\\checkmarkc\checkmarki\checkmarkr\checkmarkc\checkmark$\checkmark'\checkmark$\checkmark^\checkmark\\checkmarkc\checkmarki\checkmarkr\checkmarkc\checkmark$\checkmarka\checkmark$\checkmark^\checkmark\\checkmarkc\checkmarki\checkmarkr\checkmarkc\checkmark$\checkmarkp\checkmark$\checkmark^\checkmark\\checkmarkc\checkmarki\checkmarkr\checkmarkc\checkmark$\checkmarkr\checkmark$\checkmark^\checkmark\\checkmarkc\checkmarki\checkmarkr\checkmarkc\checkmark$\checkmarkè\checkmark$\checkmark^\checkmark\\checkmarkc\checkmarki\checkmarkr\checkmarkc\checkmark$\checkmarks\checkmark$\checkmark^\checkmark\\checkmarkc\checkmarki\checkmarkr\checkmarkc\checkmark$\checkmark \checkmark$\checkmark^\checkmark\\checkmarkc\checkmarki\checkmarkr\checkmarkc\checkmark$\checkmarkl\checkmark$\checkmark^\checkmark\\checkmarkc\checkmarki\checkmarkr\checkmarkc\checkmark$\checkmarke\checkmark$\checkmark^\checkmark\\checkmarkc\checkmarki\checkmarkr\checkmarkc\checkmark$\checkmarks\checkmark$\checkmark^\checkmark\\checkmarkc\checkmarki\checkmarkr\checkmarkc\checkmark$\checkmark \checkmark$\checkmark^\checkmark\\checkmarkc\checkmarki\checkmarkr\checkmarkc\checkmark$\checkmarkt\checkmark$\checkmark^\checkmark\\checkmarkc\checkmarki\checkmarkr\checkmarkc\checkmark$\checkmarka\checkmark$\checkmark^\checkmark\\checkmarkc\checkmarki\checkmarkr\checkmarkc\checkmark$\checkmarkb\checkmark$\checkmark^\checkmark\\checkmarkc\checkmarki\checkmarkr\checkmarkc\checkmark$\checkmarkl\checkmark$\checkmark^\checkmark\\checkmarkc\checkmarki\checkmarkr\checkmarkc\checkmark$\checkmarke\checkmark$\checkmark^\checkmark\\checkmarkc\checkmarki\checkmarkr\checkmarkc\checkmark$\checkmarks\checkmark$\checkmark^\checkmark\\checkmarkc\checkmarki\checkmarkr\checkmarkc\checkmark$\checkmark \checkmark$\checkmark^\checkmark\\checkmarkc\checkmarki\checkmarkr\checkmarkc\checkmark$\checkmarkd\checkmark$\checkmark^\checkmark\\checkmarkc\checkmarki\checkmarkr\checkmarkc\checkmark$\checkmark'\checkmark$\checkmark^\checkmark\\checkmarkc\checkmarki\checkmarkr\checkmarkc\checkmark$\checkmarku\checkmark$\checkmark^\checkmark\\checkmarkc\checkmarki\checkmarkr\checkmarkc\checkmark$\checkmarks\checkmark$\checkmark^\checkmark\\checkmarkc\checkmarki\checkmarkr\checkmarkc\checkmark$\checkmarki\checkmark$\checkmark^\checkmark\\checkmarkc\checkmarki\checkmarkr\checkmarkc\checkmark$\checkmarkn\checkmark$\checkmark^\checkmark\\checkmarkc\checkmarki\checkmarkr\checkmarkc\checkmark$\checkmarka\checkmark$\checkmark^\checkmark\\checkmarkc\checkmarki\checkmarkr\checkmarkc\checkmark$\checkmarkg\checkmark$\checkmark^\checkmark\\checkmarkc\checkmarki\checkmarkr\checkmarkc\checkmark$\checkmarke\checkmark$\checkmark^\checkmark\\checkmarkc\checkmarki\checkmarkr\checkmarkc\checkmark$\checkmark \checkmark$\checkmark^\checkmark\\checkmarkc\checkmarki\checkmarkr\checkmarkc\checkmark$\checkmark[\checkmark$\checkmark^\checkmark\\checkmarkc\checkmarki\checkmarkr\checkmarkc\checkmark$\checkmarkR\checkmark$\checkmark^\checkmark\\checkmarkc\checkmarki\checkmarkr\checkmarkc\checkmark$\checkmarké\checkmark$\checkmark^\checkmark\\checkmarkc\checkmarki\checkmarkr\checkmarkc\checkmark$\checkmarkf\checkmark$\checkmark^\checkmark\\checkmarkc\checkmarki\checkmarkr\checkmarkc\checkmark$\checkmark.\checkmark$\checkmark^\checkmark\\checkmarkc\checkmarki\checkmarkr\checkmarkc\checkmark$\checkmark]\checkmark$\checkmark^\checkmark\\checkmarkc\checkmarki\checkmarkr\checkmarkc\checkmark$\checkmark,\checkmark$\checkmark^\checkmark\\checkmarkc\checkmarki\checkmarkr\checkmarkc\checkmark$\checkmark \checkmark$\checkmark^\checkmark\\checkmarkc\checkmarki\checkmarkr\checkmarkc\checkmark$\checkmarkp\checkmark$\checkmark^\checkmark\\checkmarkc\checkmarki\checkmarkr\checkmarkc\checkmark$\checkmarko\checkmark$\checkmark^\checkmark\\checkmarkc\checkmarki\checkmarkr\checkmarkc\checkmark$\checkmarku\checkmark$\checkmark^\checkmark\\checkmarkc\checkmarki\checkmarkr\checkmarkc\checkmark$\checkmarkr\checkmark$\checkmark^\checkmark\\checkmarkc\checkmarki\checkmarkr\checkmarkc\checkmark$\checkmark \checkmark$\checkmark^\checkmark\\checkmarkc\checkmarki\checkmarkr\checkmarkc\checkmark$\checkmarkl\checkmark$\checkmark^\checkmark\\checkmarkc\checkmarki\checkmarkr\checkmarkc\checkmark$\checkmarke\checkmark$\checkmark^\checkmark\\checkmarkc\checkmarki\checkmarkr\checkmarkc\checkmark$\checkmark \checkmark$\checkmark^\checkmark\\checkmarkc\checkmarki\checkmarkr\checkmarkc\checkmark$\checkmarkf\checkmark$\checkmark^\checkmark\\checkmarkc\checkmarki\checkmarkr\checkmarkc\checkmark$\checkmarkr\checkmark$\checkmark^\checkmark\\checkmarkc\checkmarki\checkmarkr\checkmarkc\checkmark$\checkmarka\checkmark$\checkmark^\checkmark\\checkmarkc\checkmarki\checkmarkr\checkmarkc\checkmark$\checkmarki\checkmark$\checkmark^\checkmark\\checkmarkc\checkmarki\checkmarkr\checkmarkc\checkmark$\checkmarks\checkmark$\checkmark^\checkmark\\checkmarkc\checkmarki\checkmarkr\checkmarkc\checkmark$\checkmarka\checkmark$\checkmark^\checkmark\\checkmarkc\checkmarki\checkmarkr\checkmarkc\checkmark$\checkmarkg\checkmark$\checkmark^\checkmark\\checkmarkc\checkmarki\checkmarkr\checkmarkc\checkmark$\checkmarke\checkmark$\checkmark^\checkmark\\checkmarkc\checkmarki\checkmarkr\checkmarkc\checkmark$\checkmark \checkmark$\checkmark^\checkmark\\checkmarkc\checkmarki\checkmarkr\checkmarkc\checkmark$\checkmarkd\checkmark$\checkmark^\checkmark\\checkmarkc\checkmarki\checkmarkr\checkmarkc\checkmark$\checkmarke\checkmark$\checkmark^\checkmark\\checkmarkc\checkmarki\checkmarkr\checkmarkc\checkmark$\checkmark \checkmark$\checkmark^\checkmark\\checkmarkc\checkmarki\checkmarkr\checkmarkc\checkmark$\checkmarkl\checkmark$\checkmark^\checkmark\\checkmarkc\checkmarki\checkmarkr\checkmarkc\checkmark$\checkmark'\checkmark$\checkmark^\checkmark\\checkmarkc\checkmarki\checkmarkr\checkmarkc\checkmark$\checkmarka\checkmark$\checkmark^\checkmark\\checkmarkc\checkmarki\checkmarkr\checkmarkc\checkmark$\checkmarkl\checkmark$\checkmark^\checkmark\\checkmarkc\checkmarki\checkmarkr\checkmarkc\checkmark$\checkmarku\checkmark$\checkmark^\checkmark\\checkmarkc\checkmarki\checkmarkr\checkmarkc\checkmark$\checkmarkm\checkmark$\checkmark^\checkmark\\checkmarkc\checkmarki\checkmarkr\checkmarkc\checkmark$\checkmarki\checkmark$\checkmark^\checkmark\\checkmarkc\checkmarki\checkmarkr\checkmarkc\checkmark$\checkmarkn\checkmark$\checkmark^\checkmark\\checkmarkc\checkmarki\checkmarkr\checkmarkc\checkmark$\checkmarki\checkmark$\checkmark^\checkmark\\checkmarkc\checkmarki\checkmarkr\checkmarkc\checkmark$\checkmarku\checkmark$\checkmark^\checkmark\\checkmarkc\checkmarki\checkmarkr\checkmarkc\checkmark$\checkmarkm\checkmark$\checkmark^\checkmark\\checkmarkc\checkmarki\checkmarkr\checkmarkc\checkmark$\checkmark \checkmark$\checkmark^\checkmark\\checkmarkc\checkmarki\checkmarkr\checkmarkc\checkmark$\checkmark:\checkmark$\checkmark^\checkmark\\checkmarkc\checkmarki\checkmarkr\checkmarkc\checkmark$\checkmark
\checkmark$\checkmark^\checkmark\\checkmarkc\checkmarki\checkmarkr\checkmarkc\checkmark$\checkmark\\checkmark$\checkmark^\checkmark\\checkmarkc\checkmarki\checkmarkr\checkmarkc\checkmark$\checkmarkb\checkmark$\checkmark^\checkmark\\checkmarkc\checkmarki\checkmarkr\checkmarkc\checkmark$\checkmarke\checkmark$\checkmark^\checkmark\\checkmarkc\checkmarki\checkmarkr\checkmarkc\checkmark$\checkmarkg\checkmark$\checkmark^\checkmark\\checkmarkc\checkmarki\checkmarkr\checkmarkc\checkmark$\checkmarki\checkmark$\checkmark^\checkmark\\checkmarkc\checkmarki\checkmarkr\checkmarkc\checkmark$\checkmarkn\checkmark$\checkmark^\checkmark\\checkmarkc\checkmarki\checkmarkr\checkmarkc\checkmark$\checkmark{\checkmark$\checkmark^\checkmark\\checkmarkc\checkmarki\checkmarkr\checkmarkc\checkmark$\checkmarka\checkmark$\checkmark^\checkmark\\checkmarkc\checkmarki\checkmarkr\checkmarkc\checkmark$\checkmarkl\checkmark$\checkmark^\checkmark\\checkmarkc\checkmarki\checkmarkr\checkmarkc\checkmark$\checkmarki\checkmark$\checkmark^\checkmark\\checkmarkc\checkmarki\checkmarkr\checkmarkc\checkmark$\checkmarkg\checkmark$\checkmark^\checkmark\\checkmarkc\checkmarki\checkmarkr\checkmarkc\checkmark$\checkmarkn\checkmark$\checkmark^\checkmark\\checkmarkc\checkmarki\checkmarkr\checkmarkc\checkmark$\checkmark}\checkmark$\checkmark^\checkmark\\checkmarkc\checkmarki\checkmarkr\checkmarkc\checkmark$\checkmark
\checkmark$\checkmark^\checkmark\\checkmarkc\checkmarki\checkmarkr\checkmarkc\checkmark$\checkmark \checkmark$\checkmark^\checkmark\\checkmarkc\checkmarki\checkmarkr\checkmarkc\checkmark$\checkmark \checkmark$\checkmark^\checkmark\\checkmarkc\checkmarki\checkmarkr\checkmarkc\checkmark$\checkmark \checkmark$\checkmark^\checkmark\\checkmarkc\checkmarki\checkmarkr\checkmarkc\checkmark$\checkmark \checkmark$\checkmark^\checkmark\\checkmarkc\checkmarki\checkmarkr\checkmarkc\checkmark$\checkmarkF\checkmark$\checkmark^\checkmark\\checkmarkc\checkmarki\checkmarkr\checkmarkc\checkmark$\checkmark_\checkmark$\checkmark^\checkmark\\checkmarkc\checkmarki\checkmarkr\checkmarkc\checkmark$\checkmarkf\checkmark$\checkmark^\checkmark\\checkmarkc\checkmarki\checkmarkr\checkmarkc\checkmark$\checkmark \checkmark$\checkmark^\checkmark\\checkmarkc\checkmarki\checkmarkr\checkmarkc\checkmark$\checkmark&\checkmark$\checkmark^\checkmark\\checkmarkc\checkmarki\checkmarkr\checkmarkc\checkmark$\checkmark=\checkmark$\checkmark^\checkmark\\checkmarkc\checkmarki\checkmarkr\checkmarkc\checkmark$\checkmark \checkmark$\checkmark^\checkmark\\checkmarkc\checkmarki\checkmarkr\checkmarkc\checkmark$\checkmark0\checkmark$\checkmark^\checkmark\\checkmarkc\checkmarki\checkmarkr\checkmarkc\checkmark$\checkmark{\checkmark$\checkmark^\checkmark\\checkmarkc\checkmarki\checkmarkr\checkmarkc\checkmark$\checkmark,\checkmark$\checkmark^\checkmark\\checkmarkc\checkmarki\checkmarkr\checkmarkc\checkmark$\checkmark}\checkmark$\checkmark^\checkmark\\checkmarkc\checkmarki\checkmarkr\checkmarkc\checkmark$\checkmark3\checkmark$\checkmark^\checkmark\\checkmarkc\checkmarki\checkmarkr\checkmarkc\checkmark$\checkmark \checkmark$\checkmark^\checkmark\\checkmarkc\checkmarki\checkmarkr\checkmarkc\checkmark$\checkmark\\checkmark$\checkmark^\checkmark\\checkmarkc\checkmarki\checkmarkr\checkmarkc\checkmark$\checkmarkt\checkmark$\checkmark^\checkmark\\checkmarkc\checkmarki\checkmarkr\checkmarkc\checkmark$\checkmarki\checkmark$\checkmark^\checkmark\\checkmarkc\checkmarki\checkmarkr\checkmarkc\checkmark$\checkmarkm\checkmark$\checkmark^\checkmark\\checkmarkc\checkmarki\checkmarkr\checkmarkc\checkmark$\checkmarke\checkmark$\checkmark^\checkmark\\checkmarkc\checkmarki\checkmarkr\checkmarkc\checkmark$\checkmarks\checkmark$\checkmark^\checkmark\\checkmarkc\checkmarki\checkmarkr\checkmarkc\checkmark$\checkmark \checkmark$\checkmark^\checkmark\\checkmarkc\checkmarki\checkmarkr\checkmarkc\checkmark$\checkmarkF\checkmark$\checkmark^\checkmark\\checkmarkc\checkmarki\checkmarkr\checkmarkc\checkmark$\checkmark_\checkmark$\checkmark^\checkmark\\checkmarkc\checkmarki\checkmarkr\checkmarkc\checkmark$\checkmarkc\checkmark$\checkmark^\checkmark\\checkmarkc\checkmarki\checkmarkr\checkmarkc\checkmark$\checkmark \checkmark$\checkmark^\checkmark\\checkmarkc\checkmarki\checkmarkr\checkmarkc\checkmark$\checkmark=\checkmark$\checkmark^\checkmark\\checkmarkc\checkmarki\checkmarkr\checkmarkc\checkmark$\checkmark \checkmark$\checkmark^\checkmark\\checkmarkc\checkmarki\checkmarkr\checkmarkc\checkmark$\checkmark0\checkmark$\checkmark^\checkmark\\checkmarkc\checkmarki\checkmarkr\checkmarkc\checkmark$\checkmark{\checkmark$\checkmark^\checkmark\\checkmarkc\checkmarki\checkmarkr\checkmarkc\checkmark$\checkmark,\checkmark$\checkmark^\checkmark\\checkmarkc\checkmarki\checkmarkr\checkmarkc\checkmark$\checkmark}\checkmark$\checkmark^\checkmark\\checkmarkc\checkmarki\checkmarkr\checkmarkc\checkmark$\checkmark3\checkmark$\checkmark^\checkmark\\checkmarkc\checkmarki\checkmarkr\checkmarkc\checkmark$\checkmark \checkmark$\checkmark^\checkmark\\checkmarkc\checkmarki\checkmarkr\checkmarkc\checkmark$\checkmark\\checkmark$\checkmark^\checkmark\\checkmarkc\checkmarki\checkmarkr\checkmarkc\checkmark$\checkmarkt\checkmark$\checkmark^\checkmark\\checkmarkc\checkmarki\checkmarkr\checkmarkc\checkmark$\checkmarki\checkmark$\checkmark^\checkmark\\checkmarkc\checkmarki\checkmarkr\checkmarkc\checkmark$\checkmarkm\checkmark$\checkmark^\checkmark\\checkmarkc\checkmarki\checkmarkr\checkmarkc\checkmark$\checkmarke\checkmark$\checkmark^\checkmark\\checkmarkc\checkmarki\checkmarkr\checkmarkc\checkmark$\checkmarks\checkmark$\checkmark^\checkmark\\checkmarkc\checkmarki\checkmarkr\checkmarkc\checkmark$\checkmark \checkmark$\checkmark^\checkmark\\checkmarkc\checkmarki\checkmarkr\checkmarkc\checkmark$\checkmark7\checkmark$\checkmark^\checkmark\\checkmarkc\checkmarki\checkmarkr\checkmarkc\checkmark$\checkmark0\checkmark$\checkmark^\checkmark\\checkmarkc\checkmarki\checkmarkr\checkmarkc\checkmark$\checkmark \checkmark$\checkmark^\checkmark\\checkmarkc\checkmarki\checkmarkr\checkmarkc\checkmark$\checkmark=\checkmark$\checkmark^\checkmark\\checkmarkc\checkmarki\checkmarkr\checkmarkc\checkmark$\checkmark \checkmark$\checkmark^\checkmark\\checkmarkc\checkmarki\checkmarkr\checkmarkc\checkmark$\checkmark\\checkmark$\checkmark^\checkmark\\checkmarkc\checkmarki\checkmarkr\checkmarkc\checkmark$\checkmarkb\checkmark$\checkmark^\checkmark\\checkmarkc\checkmarki\checkmarkr\checkmarkc\checkmark$\checkmarko\checkmark$\checkmark^\checkmark\\checkmarkc\checkmarki\checkmarkr\checkmarkc\checkmark$\checkmarkx\checkmark$\checkmark^\checkmark\\checkmarkc\checkmarki\checkmarkr\checkmarkc\checkmark$\checkmarke\checkmark$\checkmark^\checkmark\\checkmarkc\checkmarki\checkmarkr\checkmarkc\checkmark$\checkmarkd\checkmark$\checkmark^\checkmark\\checkmarkc\checkmarki\checkmarkr\checkmarkc\checkmark$\checkmark{\checkmark$\checkmark^\checkmark\\checkmarkc\checkmarki\checkmarkr\checkmarkc\checkmark$\checkmark2\checkmark$\checkmark^\checkmark\\checkmarkc\checkmarki\checkmarkr\checkmarkc\checkmark$\checkmark1\checkmark$\checkmark^\checkmark\\checkmarkc\checkmarki\checkmarkr\checkmarkc\checkmark$\checkmark \checkmark$\checkmark^\checkmark\\checkmarkc\checkmarki\checkmarkr\checkmarkc\checkmark$\checkmark\\checkmark$\checkmark^\checkmark\\checkmarkc\checkmarki\checkmarkr\checkmarkc\checkmark$\checkmarkt\checkmark$\checkmark^\checkmark\\checkmarkc\checkmarki\checkmarkr\checkmarkc\checkmark$\checkmarke\checkmark$\checkmark^\checkmark\\checkmarkc\checkmarki\checkmarkr\checkmarkc\checkmark$\checkmarkx\checkmark$\checkmark^\checkmark\\checkmarkc\checkmarki\checkmarkr\checkmarkc\checkmark$\checkmarkt\checkmark$\checkmark^\checkmark\\checkmarkc\checkmarki\checkmarkr\checkmarkc\checkmark$\checkmark{\checkmark$\checkmark^\checkmark\\checkmarkc\checkmarki\checkmarkr\checkmarkc\checkmark$\checkmark \checkmark$\checkmark^\checkmark\\checkmarkc\checkmarki\checkmarkr\checkmarkc\checkmark$\checkmarkN\checkmark$\checkmark^\checkmark\\checkmarkc\checkmarki\checkmarkr\checkmarkc\checkmark$\checkmark}\checkmark$\checkmark^\checkmark\\checkmarkc\checkmarki\checkmarkr\checkmarkc\checkmark$\checkmark}\checkmark$\checkmark^\checkmark\\checkmarkc\checkmarki\checkmarkr\checkmarkc\checkmark$\checkmark \checkmark$\checkmark^\checkmark\\checkmarkc\checkmarki\checkmarkr\checkmarkc\checkmark$\checkmark\\checkmark$\checkmark^\checkmark\\checkmarkc\checkmarki\checkmarkr\checkmarkc\checkmark$\checkmarkl\checkmark$\checkmark^\checkmark\\checkmarkc\checkmarki\checkmarkr\checkmarkc\checkmark$\checkmarka\checkmark$\checkmark^\checkmark\\checkmarkc\checkmarki\checkmarkr\checkmarkc\checkmark$\checkmarkb\checkmark$\checkmark^\checkmark\\checkmarkc\checkmarki\checkmarkr\checkmarkc\checkmark$\checkmarke\checkmark$\checkmark^\checkmark\\checkmarkc\checkmarki\checkmarkr\checkmarkc\checkmark$\checkmarkl\checkmark$\checkmark^\checkmark\\checkmarkc\checkmarki\checkmarkr\checkmarkc\checkmark$\checkmark{\checkmark$\checkmark^\checkmark\\checkmarkc\checkmarki\checkmarkr\checkmarkc\checkmark$\checkmarke\checkmark$\checkmark^\checkmark\\checkmarkc\checkmarki\checkmarkr\checkmarkc\checkmark$\checkmarkq\checkmark$\checkmark^\checkmark\\checkmarkc\checkmarki\checkmarkr\checkmarkc\checkmark$\checkmark:\checkmark$\checkmark^\checkmark\\checkmarkc\checkmarki\checkmarkr\checkmarkc\checkmark$\checkmarkf\checkmark$\checkmark^\checkmark\\checkmarkc\checkmarki\checkmarkr\checkmarkc\checkmark$\checkmarko\checkmark$\checkmark^\checkmark\\checkmarkc\checkmarki\checkmarkr\checkmarkc\checkmark$\checkmarkr\checkmark$\checkmark^\checkmark\\checkmarkc\checkmarki\checkmarkr\checkmarkc\checkmark$\checkmarkc\checkmark$\checkmark^\checkmark\\checkmarkc\checkmarki\checkmarkr\checkmarkc\checkmark$\checkmarke\checkmark$\checkmark^\checkmark\\checkmarkc\checkmarki\checkmarkr\checkmarkc\checkmark$\checkmark_\checkmark$\checkmark^\checkmark\\checkmarkc\checkmarki\checkmarkr\checkmarkc\checkmark$\checkmarka\checkmark$\checkmark^\checkmark\\checkmarkc\checkmarki\checkmarkr\checkmarkc\checkmark$\checkmarkv\checkmark$\checkmark^\checkmark\\checkmarkc\checkmarki\checkmarkr\checkmarkc\checkmark$\checkmarka\checkmark$\checkmark^\checkmark\\checkmarkc\checkmarki\checkmarkr\checkmarkc\checkmark$\checkmarkn\checkmark$\checkmark^\checkmark\\checkmarkc\checkmarki\checkmarkr\checkmarkc\checkmark$\checkmarkc\checkmark$\checkmark^\checkmark\\checkmarkc\checkmarki\checkmarkr\checkmarkc\checkmark$\checkmarke\checkmark$\checkmark^\checkmark\\checkmarkc\checkmarki\checkmarkr\checkmarkc\checkmark$\checkmark}\checkmark$\checkmark^\checkmark\\checkmarkc\checkmarki\checkmarkr\checkmarkc\checkmark$\checkmark \checkmark$\checkmark^\checkmark\\checkmarkc\checkmarki\checkmarkr\checkmarkc\checkmark$\checkmark\\checkmark$\checkmark^\checkmark\\checkmarkc\checkmarki\checkmarkr\checkmarkc\checkmark$\checkmark\\checkmark$\checkmark^\checkmark\\checkmarkc\checkmarki\checkmarkr\checkmarkc\checkmark$\checkmark
\checkmark$\checkmark^\checkmark\\checkmarkc\checkmarki\checkmarkr\checkmarkc\checkmark$\checkmark \checkmark$\checkmark^\checkmark\\checkmarkc\checkmarki\checkmarkr\checkmarkc\checkmark$\checkmark \checkmark$\checkmark^\checkmark\\checkmarkc\checkmarki\checkmarkr\checkmarkc\checkmark$\checkmark \checkmark$\checkmark^\checkmark\\checkmarkc\checkmarki\checkmarkr\checkmarkc\checkmark$\checkmark \checkmark$\checkmark^\checkmark\\checkmarkc\checkmarki\checkmarkr\checkmarkc\checkmark$\checkmarkF\checkmark$\checkmark^\checkmark\\checkmarkc\checkmarki\checkmarkr\checkmarkc\checkmark$\checkmark_\checkmark$\checkmark^\checkmark\\checkmarkc\checkmarki\checkmarkr\checkmarkc\checkmark$\checkmarkp\checkmark$\checkmark^\checkmark\\checkmarkc\checkmarki\checkmarkr\checkmarkc\checkmark$\checkmark \checkmark$\checkmark^\checkmark\\checkmarkc\checkmarki\checkmarkr\checkmarkc\checkmark$\checkmark&\checkmark$\checkmark^\checkmark\\checkmarkc\checkmarki\checkmarkr\checkmarkc\checkmark$\checkmark=\checkmark$\checkmark^\checkmark\\checkmarkc\checkmarki\checkmarkr\checkmarkc\checkmark$\checkmark \checkmark$\checkmark^\checkmark\\checkmarkc\checkmarki\checkmarkr\checkmarkc\checkmark$\checkmark0\checkmark$\checkmark^\checkmark\\checkmarkc\checkmarki\checkmarkr\checkmarkc\checkmark$\checkmark{\checkmark$\checkmark^\checkmark\\checkmarkc\checkmarki\checkmarkr\checkmarkc\checkmark$\checkmark,\checkmark$\checkmark^\checkmark\\checkmarkc\checkmarki\checkmarkr\checkmarkc\checkmark$\checkmark}\checkmark$\checkmark^\checkmark\\checkmarkc\checkmarki\checkmarkr\checkmarkc\checkmark$\checkmark4\checkmark$\checkmark^\checkmark\\checkmarkc\checkmarki\checkmarkr\checkmarkc\checkmark$\checkmark \checkmark$\checkmark^\checkmark\\checkmarkc\checkmarki\checkmarkr\checkmarkc\checkmark$\checkmark\\checkmark$\checkmark^\checkmark\\checkmarkc\checkmarki\checkmarkr\checkmarkc\checkmark$\checkmarkt\checkmark$\checkmark^\checkmark\\checkmarkc\checkmarki\checkmarkr\checkmarkc\checkmark$\checkmarki\checkmark$\checkmark^\checkmark\\checkmarkc\checkmarki\checkmarkr\checkmarkc\checkmark$\checkmarkm\checkmark$\checkmark^\checkmark\\checkmarkc\checkmarki\checkmarkr\checkmarkc\checkmark$\checkmarke\checkmark$\checkmark^\checkmark\\checkmarkc\checkmarki\checkmarkr\checkmarkc\checkmark$\checkmarks\checkmark$\checkmark^\checkmark\\checkmarkc\checkmarki\checkmarkr\checkmarkc\checkmark$\checkmark \checkmark$\checkmark^\checkmark\\checkmarkc\checkmarki\checkmarkr\checkmarkc\checkmark$\checkmarkF\checkmark$\checkmark^\checkmark\\checkmarkc\checkmarki\checkmarkr\checkmarkc\checkmark$\checkmark_\checkmark$\checkmark^\checkmark\\checkmarkc\checkmarki\checkmarkr\checkmarkc\checkmark$\checkmarkc\checkmark$\checkmark^\checkmark\\checkmarkc\checkmarki\checkmarkr\checkmarkc\checkmark$\checkmark \checkmark$\checkmark^\checkmark\\checkmarkc\checkmarki\checkmarkr\checkmarkc\checkmark$\checkmark=\checkmark$\checkmark^\checkmark\\checkmarkc\checkmarki\checkmarkr\checkmarkc\checkmark$\checkmark \checkmark$\checkmark^\checkmark\\checkmarkc\checkmarki\checkmarkr\checkmarkc\checkmark$\checkmark0\checkmark$\checkmark^\checkmark\\checkmarkc\checkmarki\checkmarkr\checkmarkc\checkmark$\checkmark{\checkmark$\checkmark^\checkmark\\checkmarkc\checkmarki\checkmarkr\checkmarkc\checkmark$\checkmark,\checkmark$\checkmark^\checkmark\\checkmarkc\checkmarki\checkmarkr\checkmarkc\checkmark$\checkmark}\checkmark$\checkmark^\checkmark\\checkmarkc\checkmarki\checkmarkr\checkmarkc\checkmark$\checkmark4\checkmark$\checkmark^\checkmark\\checkmarkc\checkmarki\checkmarkr\checkmarkc\checkmark$\checkmark \checkmark$\checkmark^\checkmark\\checkmarkc\checkmarki\checkmarkr\checkmarkc\checkmark$\checkmark\\checkmark$\checkmark^\checkmark\\checkmarkc\checkmarki\checkmarkr\checkmarkc\checkmark$\checkmarkt\checkmark$\checkmark^\checkmark\\checkmarkc\checkmarki\checkmarkr\checkmarkc\checkmark$\checkmarki\checkmark$\checkmark^\checkmark\\checkmarkc\checkmarki\checkmarkr\checkmarkc\checkmark$\checkmarkm\checkmark$\checkmark^\checkmark\\checkmarkc\checkmarki\checkmarkr\checkmarkc\checkmark$\checkmarke\checkmark$\checkmark^\checkmark\\checkmarkc\checkmarki\checkmarkr\checkmarkc\checkmark$\checkmarks\checkmark$\checkmark^\checkmark\\checkmarkc\checkmarki\checkmarkr\checkmarkc\checkmark$\checkmark \checkmark$\checkmark^\checkmark\\checkmarkc\checkmarki\checkmarkr\checkmarkc\checkmark$\checkmark7\checkmark$\checkmark^\checkmark\\checkmarkc\checkmarki\checkmarkr\checkmarkc\checkmark$\checkmark0\checkmark$\checkmark^\checkmark\\checkmarkc\checkmarki\checkmarkr\checkmarkc\checkmark$\checkmark \checkmark$\checkmark^\checkmark\\checkmarkc\checkmarki\checkmarkr\checkmarkc\checkmark$\checkmark=\checkmark$\checkmark^\checkmark\\checkmarkc\checkmarki\checkmarkr\checkmarkc\checkmark$\checkmark \checkmark$\checkmark^\checkmark\\checkmarkc\checkmarki\checkmarkr\checkmarkc\checkmark$\checkmark\\checkmark$\checkmark^\checkmark\\checkmarkc\checkmarki\checkmarkr\checkmarkc\checkmark$\checkmarkb\checkmark$\checkmark^\checkmark\\checkmarkc\checkmarki\checkmarkr\checkmarkc\checkmark$\checkmarko\checkmark$\checkmark^\checkmark\\checkmarkc\checkmarki\checkmarkr\checkmarkc\checkmark$\checkmarkx\checkmark$\checkmark^\checkmark\\checkmarkc\checkmarki\checkmarkr\checkmarkc\checkmark$\checkmarke\checkmark$\checkmark^\checkmark\\checkmarkc\checkmarki\checkmarkr\checkmarkc\checkmark$\checkmarkd\checkmark$\checkmark^\checkmark\\checkmarkc\checkmarki\checkmarkr\checkmarkc\checkmark$\checkmark{\checkmark$\checkmark^\checkmark\\checkmarkc\checkmarki\checkmarkr\checkmarkc\checkmark$\checkmark2\checkmark$\checkmark^\checkmark\\checkmarkc\checkmarki\checkmarkr\checkmarkc\checkmark$\checkmark8\checkmark$\checkmark^\checkmark\\checkmarkc\checkmarki\checkmarkr\checkmarkc\checkmark$\checkmark \checkmark$\checkmark^\checkmark\\checkmarkc\checkmarki\checkmarkr\checkmarkc\checkmark$\checkmark\\checkmark$\checkmark^\checkmark\\checkmarkc\checkmarki\checkmarkr\checkmarkc\checkmark$\checkmarkt\checkmark$\checkmark^\checkmark\\checkmarkc\checkmarki\checkmarkr\checkmarkc\checkmark$\checkmarke\checkmark$\checkmark^\checkmark\\checkmarkc\checkmarki\checkmarkr\checkmarkc\checkmark$\checkmarkx\checkmark$\checkmark^\checkmark\\checkmarkc\checkmarki\checkmarkr\checkmarkc\checkmark$\checkmarkt\checkmark$\checkmark^\checkmark\\checkmarkc\checkmarki\checkmarkr\checkmarkc\checkmark$\checkmark{\checkmark$\checkmark^\checkmark\\checkmarkc\checkmarki\checkmarkr\checkmarkc\checkmark$\checkmark \checkmark$\checkmark^\checkmark\\checkmarkc\checkmarki\checkmarkr\checkmarkc\checkmark$\checkmarkN\checkmark$\checkmark^\checkmark\\checkmarkc\checkmarki\checkmarkr\checkmarkc\checkmark$\checkmark}\checkmark$\checkmark^\checkmark\\checkmarkc\checkmarki\checkmarkr\checkmarkc\checkmark$\checkmark}\checkmark$\checkmark^\checkmark\\checkmarkc\checkmarki\checkmarkr\checkmarkc\checkmark$\checkmark \checkmark$\checkmark^\checkmark\\checkmarkc\checkmarki\checkmarkr\checkmarkc\checkmark$\checkmark\\checkmark$\checkmark^\checkmark\\checkmarkc\checkmarki\checkmarkr\checkmarkc\checkmark$\checkmarkl\checkmark$\checkmark^\checkmark\\checkmarkc\checkmarki\checkmarkr\checkmarkc\checkmark$\checkmarka\checkmark$\checkmark^\checkmark\\checkmarkc\checkmarki\checkmarkr\checkmarkc\checkmark$\checkmarkb\checkmark$\checkmark^\checkmark\\checkmarkc\checkmarki\checkmarkr\checkmarkc\checkmark$\checkmarke\checkmark$\checkmark^\checkmark\\checkmarkc\checkmarki\checkmarkr\checkmarkc\checkmark$\checkmarkl\checkmark$\checkmark^\checkmark\\checkmarkc\checkmarki\checkmarkr\checkmarkc\checkmark$\checkmark{\checkmark$\checkmark^\checkmark\\checkmarkc\checkmarki\checkmarkr\checkmarkc\checkmark$\checkmarke\checkmark$\checkmark^\checkmark\\checkmarkc\checkmarki\checkmarkr\checkmarkc\checkmark$\checkmarkq\checkmark$\checkmark^\checkmark\\checkmarkc\checkmarki\checkmarkr\checkmarkc\checkmark$\checkmark:\checkmark$\checkmark^\checkmark\\checkmarkc\checkmarki\checkmarkr\checkmarkc\checkmark$\checkmarkf\checkmark$\checkmark^\checkmark\\checkmarkc\checkmarki\checkmarkr\checkmarkc\checkmark$\checkmarko\checkmark$\checkmark^\checkmark\\checkmarkc\checkmarki\checkmarkr\checkmarkc\checkmark$\checkmarkr\checkmark$\checkmark^\checkmark\\checkmarkc\checkmarki\checkmarkr\checkmarkc\checkmark$\checkmarkc\checkmark$\checkmark^\checkmark\\checkmarkc\checkmarki\checkmarkr\checkmarkc\checkmark$\checkmarke\checkmark$\checkmark^\checkmark\\checkmarkc\checkmarki\checkmarkr\checkmarkc\checkmark$\checkmark_\checkmark$\checkmark^\checkmark\\checkmarkc\checkmarki\checkmarkr\checkmarkc\checkmark$\checkmarkp\checkmark$\checkmark^\checkmark\\checkmarkc\checkmarki\checkmarkr\checkmarkc\checkmark$\checkmarke\checkmark$\checkmark^\checkmark\\checkmarkc\checkmarki\checkmarkr\checkmarkc\checkmark$\checkmarkn\checkmark$\checkmark^\checkmark\\checkmarkc\checkmarki\checkmarkr\checkmarkc\checkmark$\checkmarke\checkmark$\checkmark^\checkmark\\checkmarkc\checkmarki\checkmarkr\checkmarkc\checkmark$\checkmarkt\checkmark$\checkmark^\checkmark\\checkmarkc\checkmarki\checkmarkr\checkmarkc\checkmark$\checkmarkr\checkmark$\checkmark^\checkmark\\checkmarkc\checkmarki\checkmarkr\checkmarkc\checkmark$\checkmarka\checkmark$\checkmark^\checkmark\\checkmarkc\checkmarki\checkmarkr\checkmarkc\checkmark$\checkmarkt\checkmark$\checkmark^\checkmark\\checkmarkc\checkmarki\checkmarkr\checkmarkc\checkmark$\checkmarki\checkmark$\checkmark^\checkmark\\checkmarkc\checkmarki\checkmarkr\checkmarkc\checkmark$\checkmarko\checkmark$\checkmark^\checkmark\\checkmarkc\checkmarki\checkmarkr\checkmarkc\checkmark$\checkmarkn\checkmark$\checkmark^\checkmark\\checkmarkc\checkmarki\checkmarkr\checkmarkc\checkmark$\checkmark}\checkmark$\checkmark^\checkmark\\checkmarkc\checkmarki\checkmarkr\checkmarkc\checkmark$\checkmark
\checkmark$\checkmark^\checkmark\\checkmarkc\checkmarki\checkmarkr\checkmarkc\checkmark$\checkmark\\checkmark$\checkmark^\checkmark\\checkmarkc\checkmarki\checkmarkr\checkmarkc\checkmark$\checkmarke\checkmark$\checkmark^\checkmark\\checkmarkc\checkmarki\checkmarkr\checkmarkc\checkmark$\checkmarkn\checkmark$\checkmark^\checkmark\\checkmarkc\checkmarki\checkmarkr\checkmarkc\checkmark$\checkmarkd\checkmark$\checkmark^\checkmark\\checkmarkc\checkmarki\checkmarkr\checkmarkc\checkmark$\checkmark{\checkmark$\checkmark^\checkmark\\checkmarkc\checkmarki\checkmarkr\checkmarkc\checkmark$\checkmarka\checkmark$\checkmark^\checkmark\\checkmarkc\checkmarki\checkmarkr\checkmarkc\checkmark$\checkmarkl\checkmark$\checkmark^\checkmark\\checkmarkc\checkmarki\checkmarkr\checkmarkc\checkmark$\checkmarki\checkmark$\checkmark^\checkmark\\checkmarkc\checkmarki\checkmarkr\checkmarkc\checkmark$\checkmarkg\checkmark$\checkmark^\checkmark\\checkmarkc\checkmarki\checkmarkr\checkmarkc\checkmark$\checkmarkn\checkmark$\checkmark^\checkmark\\checkmarkc\checkmarki\checkmarkr\checkmarkc\checkmark$\checkmark}\checkmark$\checkmark^\checkmark\\checkmarkc\checkmarki\checkmarkr\checkmarkc\checkmark$\checkmark
\checkmark$\checkmark^\checkmark\\checkmarkc\checkmarki\checkmarkr\checkmarkc\checkmark$\checkmark
\checkmark$\checkmark^\checkmark\\checkmarkc\checkmarki\checkmarkr\checkmarkc\checkmark$\checkmark\\checkmark$\checkmark^\checkmark\\checkmarkc\checkmarki\checkmarkr\checkmarkc\checkmark$\checkmarks\checkmark$\checkmark^\checkmark\\checkmarkc\checkmarki\checkmarkr\checkmarkc\checkmark$\checkmarku\checkmark$\checkmark^\checkmark\\checkmarkc\checkmarki\checkmarkr\checkmarkc\checkmark$\checkmarkb\checkmark$\checkmark^\checkmark\\checkmarkc\checkmarki\checkmarkr\checkmarkc\checkmark$\checkmarks\checkmark$\checkmark^\checkmark\\checkmarkc\checkmarki\checkmarkr\checkmarkc\checkmark$\checkmarku\checkmark$\checkmark^\checkmark\\checkmarkc\checkmarki\checkmarkr\checkmarkc\checkmark$\checkmarkb\checkmark$\checkmark^\checkmark\\checkmarkc\checkmarki\checkmarkr\checkmarkc\checkmark$\checkmarks\checkmark$\checkmark^\checkmark\\checkmarkc\checkmarki\checkmarkr\checkmarkc\checkmark$\checkmarke\checkmark$\checkmark^\checkmark\\checkmarkc\checkmarki\checkmarkr\checkmarkc\checkmark$\checkmarkc\checkmark$\checkmark^\checkmark\\checkmarkc\checkmarki\checkmarkr\checkmarkc\checkmark$\checkmarkt\checkmark$\checkmark^\checkmark\\checkmarkc\checkmarki\checkmarkr\checkmarkc\checkmark$\checkmarki\checkmark$\checkmark^\checkmark\\checkmarkc\checkmarki\checkmarkr\checkmarkc\checkmark$\checkmarko\checkmark$\checkmark^\checkmark\\checkmarkc\checkmarki\checkmarkr\checkmarkc\checkmark$\checkmarkn\checkmark$\checkmark^\checkmark\\checkmarkc\checkmarki\checkmarkr\checkmarkc\checkmark$\checkmark{\checkmark$\checkmark^\checkmark\\checkmarkc\checkmarki\checkmarkr\checkmarkc\checkmark$\checkmarkR\checkmark$\checkmark^\checkmark\\checkmarkc\checkmarki\checkmarkr\checkmarkc\checkmark$\checkmarké\checkmark$\checkmark^\checkmark\\checkmarkc\checkmarki\checkmarkr\checkmarkc\checkmark$\checkmarks\checkmark$\checkmark^\checkmark\\checkmarkc\checkmarki\checkmarkr\checkmarkc\checkmark$\checkmarku\checkmark$\checkmark^\checkmark\\checkmarkc\checkmarki\checkmarkr\checkmarkc\checkmark$\checkmarkl\checkmark$\checkmark^\checkmark\\checkmarkc\checkmarki\checkmarkr\checkmarkc\checkmark$\checkmarkt\checkmark$\checkmark^\checkmark\\checkmarkc\checkmarki\checkmarkr\checkmarkc\checkmark$\checkmarka\checkmark$\checkmark^\checkmark\\checkmarkc\checkmarki\checkmarkr\checkmarkc\checkmark$\checkmarkn\checkmark$\checkmark^\checkmark\\checkmarkc\checkmarki\checkmarkr\checkmarkc\checkmark$\checkmarkt\checkmark$\checkmark^\checkmark\\checkmarkc\checkmarki\checkmarkr\checkmarkc\checkmark$\checkmarke\checkmark$\checkmark^\checkmark\\checkmarkc\checkmarki\checkmarkr\checkmarkc\checkmark$\checkmark \checkmark$\checkmark^\checkmark\\checkmarkc\checkmarki\checkmarkr\checkmarkc\checkmark$\checkmarkd\checkmark$\checkmark^\checkmark\\checkmarkc\checkmarki\checkmarkr\checkmarkc\checkmark$\checkmarke\checkmark$\checkmark^\checkmark\\checkmarkc\checkmarki\checkmarkr\checkmarkc\checkmark$\checkmarks\checkmark$\checkmark^\checkmark\\checkmarkc\checkmarki\checkmarkr\checkmarkc\checkmark$\checkmark \checkmark$\checkmark^\checkmark\\checkmarkc\checkmarki\checkmarkr\checkmarkc\checkmark$\checkmarke\checkmark$\checkmark^\checkmark\\checkmarkc\checkmarki\checkmarkr\checkmarkc\checkmark$\checkmarkf\checkmark$\checkmark^\checkmark\\checkmarkc\checkmarki\checkmarkr\checkmarkc\checkmark$\checkmarkf\checkmark$\checkmark^\checkmark\\checkmarkc\checkmarki\checkmarkr\checkmarkc\checkmark$\checkmarko\checkmark$\checkmark^\checkmark\\checkmarkc\checkmarki\checkmarkr\checkmarkc\checkmark$\checkmarkr\checkmark$\checkmark^\checkmark\\checkmarkc\checkmarki\checkmarkr\checkmarkc\checkmark$\checkmarkt\checkmark$\checkmark^\checkmark\\checkmarkc\checkmarki\checkmarkr\checkmarkc\checkmark$\checkmarks\checkmark$\checkmark^\checkmark\\checkmarkc\checkmarki\checkmarkr\checkmarkc\checkmark$\checkmark \checkmark$\checkmark^\checkmark\\checkmarkc\checkmarki\checkmarkr\checkmarkc\checkmark$\checkmarkd\checkmark$\checkmark^\checkmark\\checkmarkc\checkmarki\checkmarkr\checkmarkc\checkmark$\checkmarke\checkmark$\checkmark^\checkmark\\checkmarkc\checkmarki\checkmarkr\checkmarkc\checkmark$\checkmark \checkmark$\checkmark^\checkmark\\checkmarkc\checkmarki\checkmarkr\checkmarkc\checkmark$\checkmarkc\checkmark$\checkmark^\checkmark\\checkmarkc\checkmarki\checkmarkr\checkmarkc\checkmark$\checkmarko\checkmark$\checkmark^\checkmark\\checkmarkc\checkmarki\checkmarkr\checkmarkc\checkmark$\checkmarku\checkmark$\checkmark^\checkmark\\checkmarkc\checkmarki\checkmarkr\checkmarkc\checkmark$\checkmarkp\checkmark$\checkmark^\checkmark\\checkmarkc\checkmarki\checkmarkr\checkmarkc\checkmark$\checkmarke\checkmark$\checkmark^\checkmark\\checkmarkc\checkmarki\checkmarkr\checkmarkc\checkmark$\checkmark}\checkmark$\checkmark^\checkmark\\checkmarkc\checkmarki\checkmarkr\checkmarkc\checkmark$\checkmark
\checkmark$\checkmark^\checkmark\\checkmarkc\checkmarki\checkmarkr\checkmarkc\checkmark$\checkmark\\checkmark$\checkmark^\checkmark\\checkmarkc\checkmarki\checkmarkr\checkmarkc\checkmark$\checkmarkb\checkmark$\checkmark^\checkmark\\checkmarkc\checkmarki\checkmarkr\checkmarkc\checkmark$\checkmarke\checkmark$\checkmark^\checkmark\\checkmarkc\checkmarki\checkmarkr\checkmarkc\checkmark$\checkmarkg\checkmark$\checkmark^\checkmark\\checkmarkc\checkmarki\checkmarkr\checkmarkc\checkmark$\checkmarki\checkmark$\checkmark^\checkmark\\checkmarkc\checkmarki\checkmarkr\checkmarkc\checkmark$\checkmarkn\checkmark$\checkmark^\checkmark\\checkmarkc\checkmarki\checkmarkr\checkmarkc\checkmark$\checkmark{\checkmark$\checkmark^\checkmark\\checkmarkc\checkmarki\checkmarkr\checkmarkc\checkmark$\checkmarke\checkmark$\checkmark^\checkmark\\checkmarkc\checkmarki\checkmarkr\checkmarkc\checkmark$\checkmarkq\checkmark$\checkmark^\checkmark\\checkmarkc\checkmarki\checkmarkr\checkmarkc\checkmark$\checkmarku\checkmark$\checkmark^\checkmark\\checkmarkc\checkmarki\checkmarkr\checkmarkc\checkmark$\checkmarka\checkmark$\checkmark^\checkmark\\checkmarkc\checkmarki\checkmarkr\checkmarkc\checkmark$\checkmarkt\checkmark$\checkmark^\checkmark\\checkmarkc\checkmarki\checkmarkr\checkmarkc\checkmark$\checkmarki\checkmark$\checkmark^\checkmark\\checkmarkc\checkmarki\checkmarkr\checkmarkc\checkmark$\checkmarko\checkmark$\checkmark^\checkmark\\checkmarkc\checkmarki\checkmarkr\checkmarkc\checkmark$\checkmarkn\checkmark$\checkmark^\checkmark\\checkmarkc\checkmarki\checkmarkr\checkmarkc\checkmark$\checkmark}\checkmark$\checkmark^\checkmark\\checkmarkc\checkmarki\checkmarkr\checkmarkc\checkmark$\checkmark
\checkmark$\checkmark^\checkmark\\checkmarkc\checkmarki\checkmarkr\checkmarkc\checkmark$\checkmark \checkmark$\checkmark^\checkmark\\checkmarkc\checkmarki\checkmarkr\checkmarkc\checkmark$\checkmark \checkmark$\checkmark^\checkmark\\checkmarkc\checkmarki\checkmarkr\checkmarkc\checkmark$\checkmark \checkmark$\checkmark^\checkmark\\checkmarkc\checkmarki\checkmarkr\checkmarkc\checkmark$\checkmark \checkmark$\checkmark^\checkmark\\checkmarkc\checkmarki\checkmarkr\checkmarkc\checkmark$\checkmark\\checkmark$\checkmark^\checkmark\\checkmarkc\checkmarki\checkmarkr\checkmarkc\checkmark$\checkmark|\checkmark$\checkmark^\checkmark\\checkmarkc\checkmarki\checkmarkr\checkmarkc\checkmark$\checkmark\\checkmark$\checkmark^\checkmark\\checkmarkc\checkmarki\checkmarkr\checkmarkc\checkmark$\checkmarkv\checkmark$\checkmark^\checkmark\\checkmarkc\checkmarki\checkmarkr\checkmarkc\checkmark$\checkmarke\checkmark$\checkmark^\checkmark\\checkmarkc\checkmarki\checkmarkr\checkmarkc\checkmark$\checkmarkc\checkmark$\checkmark^\checkmark\\checkmarkc\checkmarki\checkmarkr\checkmarkc\checkmark$\checkmark{\checkmark$\checkmark^\checkmark\\checkmarkc\checkmarki\checkmarkr\checkmarkc\checkmark$\checkmarkR\checkmark$\checkmark^\checkmark\\checkmarkc\checkmarki\checkmarkr\checkmarkc\checkmark$\checkmark}\checkmark$\checkmark^\checkmark\\checkmarkc\checkmarki\checkmarkr\checkmarkc\checkmark$\checkmark\\checkmark$\checkmark^\checkmark\\checkmarkc\checkmarki\checkmarkr\checkmarkc\checkmark$\checkmark|\checkmark$\checkmark^\checkmark\\checkmarkc\checkmarki\checkmarkr\checkmarkc\checkmark$\checkmark \checkmark$\checkmark^\checkmark\\checkmarkc\checkmarki\checkmarkr\checkmarkc\checkmark$\checkmark=\checkmark$\checkmark^\checkmark\\checkmarkc\checkmarki\checkmarkr\checkmarkc\checkmark$\checkmark \checkmark$\checkmark^\checkmark\\checkmarkc\checkmarki\checkmarkr\checkmarkc\checkmark$\checkmark\\checkmark$\checkmark^\checkmark\\checkmarkc\checkmarki\checkmarkr\checkmarkc\checkmark$\checkmarks\checkmark$\checkmark^\checkmark\\checkmarkc\checkmarki\checkmarkr\checkmarkc\checkmark$\checkmarkq\checkmark$\checkmark^\checkmark\\checkmarkc\checkmarki\checkmarkr\checkmarkc\checkmark$\checkmarkr\checkmark$\checkmark^\checkmark\\checkmarkc\checkmarki\checkmarkr\checkmarkc\checkmark$\checkmarkt\checkmark$\checkmark^\checkmark\\checkmarkc\checkmarki\checkmarkr\checkmarkc\checkmark$\checkmark{\checkmark$\checkmark^\checkmark\\checkmarkc\checkmarki\checkmarkr\checkmarkc\checkmark$\checkmarkF\checkmark$\checkmark^\checkmark\\checkmarkc\checkmarki\checkmarkr\checkmarkc\checkmark$\checkmark_\checkmark$\checkmark^\checkmark\\checkmarkc\checkmarki\checkmarkr\checkmarkc\checkmark$\checkmarkc\checkmark$\checkmark^\checkmark\\checkmarkc\checkmarki\checkmarkr\checkmarkc\checkmark$\checkmark^\checkmark$\checkmark^\checkmark\\checkmarkc\checkmarki\checkmarkr\checkmarkc\checkmark$\checkmark2\checkmark$\checkmark^\checkmark\\checkmarkc\checkmarki\checkmarkr\checkmarkc\checkmark$\checkmark \checkmark$\checkmark^\checkmark\\checkmarkc\checkmarki\checkmarkr\checkmarkc\checkmark$\checkmark+\checkmark$\checkmark^\checkmark\\checkmarkc\checkmarki\checkmarkr\checkmarkc\checkmark$\checkmark \checkmark$\checkmark^\checkmark\\checkmarkc\checkmarki\checkmarkr\checkmarkc\checkmark$\checkmarkF\checkmark$\checkmark^\checkmark\\checkmarkc\checkmarki\checkmarkr\checkmarkc\checkmark$\checkmark_\checkmark$\checkmark^\checkmark\\checkmarkc\checkmarki\checkmarkr\checkmarkc\checkmark$\checkmarkf\checkmark$\checkmark^\checkmark\\checkmarkc\checkmarki\checkmarkr\checkmarkc\checkmark$\checkmark^\checkmark$\checkmark^\checkmark\\checkmarkc\checkmarki\checkmarkr\checkmarkc\checkmark$\checkmark2\checkmark$\checkmark^\checkmark\\checkmarkc\checkmarki\checkmarkr\checkmarkc\checkmark$\checkmark \checkmark$\checkmark^\checkmark\\checkmarkc\checkmarki\checkmarkr\checkmarkc\checkmark$\checkmark+\checkmark$\checkmark^\checkmark\\checkmarkc\checkmarki\checkmarkr\checkmarkc\checkmark$\checkmark \checkmark$\checkmark^\checkmark\\checkmarkc\checkmarki\checkmarkr\checkmarkc\checkmark$\checkmarkF\checkmark$\checkmark^\checkmark\\checkmarkc\checkmarki\checkmarkr\checkmarkc\checkmark$\checkmark_\checkmark$\checkmark^\checkmark\\checkmarkc\checkmarki\checkmarkr\checkmarkc\checkmark$\checkmarkp\checkmark$\checkmark^\checkmark\\checkmarkc\checkmarki\checkmarkr\checkmarkc\checkmark$\checkmark^\checkmark$\checkmark^\checkmark\\checkmarkc\checkmarki\checkmarkr\checkmarkc\checkmark$\checkmark2\checkmark$\checkmark^\checkmark\\checkmarkc\checkmarki\checkmarkr\checkmarkc\checkmark$\checkmark}\checkmark$\checkmark^\checkmark\\checkmarkc\checkmarki\checkmarkr\checkmarkc\checkmark$\checkmark \checkmark$\checkmark^\checkmark\\checkmarkc\checkmarki\checkmarkr\checkmarkc\checkmark$\checkmark=\checkmark$\checkmark^\checkmark\\checkmarkc\checkmarki\checkmarkr\checkmarkc\checkmark$\checkmark \checkmark$\checkmark^\checkmark\\checkmarkc\checkmarki\checkmarkr\checkmarkc\checkmark$\checkmark\\checkmark$\checkmark^\checkmark\\checkmarkc\checkmarki\checkmarkr\checkmarkc\checkmark$\checkmarks\checkmark$\checkmark^\checkmark\\checkmarkc\checkmarki\checkmarkr\checkmarkc\checkmark$\checkmarkq\checkmark$\checkmark^\checkmark\\checkmarkc\checkmarki\checkmarkr\checkmarkc\checkmark$\checkmarkr\checkmark$\checkmark^\checkmark\\checkmarkc\checkmarki\checkmarkr\checkmarkc\checkmark$\checkmarkt\checkmark$\checkmark^\checkmark\\checkmarkc\checkmarki\checkmarkr\checkmarkc\checkmark$\checkmark{\checkmark$\checkmark^\checkmark\\checkmarkc\checkmarki\checkmarkr\checkmarkc\checkmark$\checkmark7\checkmark$\checkmark^\checkmark\\checkmarkc\checkmarki\checkmarkr\checkmarkc\checkmark$\checkmark0\checkmark$\checkmark^\checkmark\\checkmarkc\checkmarki\checkmarkr\checkmarkc\checkmark$\checkmark^\checkmark$\checkmark^\checkmark\\checkmarkc\checkmarki\checkmarkr\checkmarkc\checkmark$\checkmark2\checkmark$\checkmark^\checkmark\\checkmarkc\checkmarki\checkmarkr\checkmarkc\checkmark$\checkmark \checkmark$\checkmark^\checkmark\\checkmarkc\checkmarki\checkmarkr\checkmarkc\checkmark$\checkmark+\checkmark$\checkmark^\checkmark\\checkmarkc\checkmarki\checkmarkr\checkmarkc\checkmark$\checkmark \checkmark$\checkmark^\checkmark\\checkmarkc\checkmarki\checkmarkr\checkmarkc\checkmark$\checkmark2\checkmark$\checkmark^\checkmark\\checkmarkc\checkmarki\checkmarkr\checkmarkc\checkmark$\checkmark1\checkmark$\checkmark^\checkmark\\checkmarkc\checkmarki\checkmarkr\checkmarkc\checkmark$\checkmark^\checkmark$\checkmark^\checkmark\\checkmarkc\checkmarki\checkmarkr\checkmarkc\checkmark$\checkmark2\checkmark$\checkmark^\checkmark\\checkmarkc\checkmarki\checkmarkr\checkmarkc\checkmark$\checkmark \checkmark$\checkmark^\checkmark\\checkmarkc\checkmarki\checkmarkr\checkmarkc\checkmark$\checkmark+\checkmark$\checkmark^\checkmark\\checkmarkc\checkmarki\checkmarkr\checkmarkc\checkmark$\checkmark \checkmark$\checkmark^\checkmark\\checkmarkc\checkmarki\checkmarkr\checkmarkc\checkmark$\checkmark2\checkmark$\checkmark^\checkmark\\checkmarkc\checkmarki\checkmarkr\checkmarkc\checkmark$\checkmark8\checkmark$\checkmark^\checkmark\\checkmarkc\checkmarki\checkmarkr\checkmarkc\checkmark$\checkmark^\checkmark$\checkmark^\checkmark\\checkmarkc\checkmarki\checkmarkr\checkmarkc\checkmark$\checkmark2\checkmark$\checkmark^\checkmark\\checkmarkc\checkmarki\checkmarkr\checkmarkc\checkmark$\checkmark}\checkmark$\checkmark^\checkmark\\checkmarkc\checkmarki\checkmarkr\checkmarkc\checkmark$\checkmark \checkmark$\checkmark^\checkmark\\checkmarkc\checkmarki\checkmarkr\checkmarkc\checkmark$\checkmark=\checkmark$\checkmark^\checkmark\\checkmarkc\checkmarki\checkmarkr\checkmarkc\checkmark$\checkmark \checkmark$\checkmark^\checkmark\\checkmarkc\checkmarki\checkmarkr\checkmarkc\checkmark$\checkmark\\checkmark$\checkmark^\checkmark\\checkmarkc\checkmarki\checkmarkr\checkmarkc\checkmark$\checkmarks\checkmark$\checkmark^\checkmark\\checkmarkc\checkmarki\checkmarkr\checkmarkc\checkmark$\checkmarkq\checkmark$\checkmark^\checkmark\\checkmarkc\checkmarki\checkmarkr\checkmarkc\checkmark$\checkmarkr\checkmark$\checkmark^\checkmark\\checkmarkc\checkmarki\checkmarkr\checkmarkc\checkmark$\checkmarkt\checkmark$\checkmark^\checkmark\\checkmarkc\checkmarki\checkmarkr\checkmarkc\checkmark$\checkmark{\checkmark$\checkmark^\checkmark\\checkmarkc\checkmarki\checkmarkr\checkmarkc\checkmark$\checkmark4\checkmark$\checkmark^\checkmark\\checkmarkc\checkmarki\checkmarkr\checkmarkc\checkmark$\checkmark9\checkmark$\checkmark^\checkmark\\checkmarkc\checkmarki\checkmarkr\checkmarkc\checkmark$\checkmark0\checkmark$\checkmark^\checkmark\\checkmarkc\checkmarki\checkmarkr\checkmarkc\checkmark$\checkmark0\checkmark$\checkmark^\checkmark\\checkmarkc\checkmarki\checkmarkr\checkmarkc\checkmark$\checkmark \checkmark$\checkmark^\checkmark\\checkmarkc\checkmarki\checkmarkr\checkmarkc\checkmark$\checkmark+\checkmark$\checkmark^\checkmark\\checkmarkc\checkmarki\checkmarkr\checkmarkc\checkmark$\checkmark \checkmark$\checkmark^\checkmark\\checkmarkc\checkmarki\checkmarkr\checkmarkc\checkmark$\checkmark4\checkmark$\checkmark^\checkmark\\checkmarkc\checkmarki\checkmarkr\checkmarkc\checkmark$\checkmark4\checkmark$\checkmark^\checkmark\\checkmarkc\checkmarki\checkmarkr\checkmarkc\checkmark$\checkmark1\checkmark$\checkmark^\checkmark\\checkmarkc\checkmarki\checkmarkr\checkmarkc\checkmark$\checkmark \checkmark$\checkmark^\checkmark\\checkmarkc\checkmarki\checkmarkr\checkmarkc\checkmark$\checkmark+\checkmark$\checkmark^\checkmark\\checkmarkc\checkmarki\checkmarkr\checkmarkc\checkmark$\checkmark \checkmark$\checkmark^\checkmark\\checkmarkc\checkmarki\checkmarkr\checkmarkc\checkmark$\checkmark7\checkmark$\checkmark^\checkmark\\checkmarkc\checkmarki\checkmarkr\checkmarkc\checkmark$\checkmark8\checkmark$\checkmark^\checkmark\\checkmarkc\checkmarki\checkmarkr\checkmarkc\checkmark$\checkmark4\checkmark$\checkmark^\checkmark\\checkmarkc\checkmarki\checkmarkr\checkmarkc\checkmark$\checkmark}\checkmark$\checkmark^\checkmark\\checkmarkc\checkmarki\checkmarkr\checkmarkc\checkmark$\checkmark \checkmark$\checkmark^\checkmark\\checkmarkc\checkmarki\checkmarkr\checkmarkc\checkmark$\checkmark\\checkmark$\checkmark^\checkmark\\checkmarkc\checkmarki\checkmarkr\checkmarkc\checkmark$\checkmarka\checkmark$\checkmark^\checkmark\\checkmarkc\checkmarki\checkmarkr\checkmarkc\checkmark$\checkmarkp\checkmark$\checkmark^\checkmark\\checkmarkc\checkmarki\checkmarkr\checkmarkc\checkmark$\checkmarkp\checkmark$\checkmark^\checkmark\\checkmarkc\checkmarki\checkmarkr\checkmarkc\checkmark$\checkmarkr\checkmark$\checkmark^\checkmark\\checkmarkc\checkmarki\checkmarkr\checkmarkc\checkmark$\checkmarko\checkmark$\checkmark^\checkmark\\checkmarkc\checkmarki\checkmarkr\checkmarkc\checkmark$\checkmarkx\checkmark$\checkmark^\checkmark\\checkmarkc\checkmarki\checkmarkr\checkmarkc\checkmark$\checkmark \checkmark$\checkmark^\checkmark\\checkmarkc\checkmarki\checkmarkr\checkmarkc\checkmark$\checkmark\\checkmark$\checkmark^\checkmark\\checkmarkc\checkmarki\checkmarkr\checkmarkc\checkmark$\checkmarkb\checkmark$\checkmark^\checkmark\\checkmarkc\checkmarki\checkmarkr\checkmarkc\checkmark$\checkmarko\checkmark$\checkmark^\checkmark\\checkmarkc\checkmarki\checkmarkr\checkmarkc\checkmark$\checkmarkx\checkmark$\checkmark^\checkmark\\checkmarkc\checkmarki\checkmarkr\checkmarkc\checkmark$\checkmarke\checkmark$\checkmark^\checkmark\\checkmarkc\checkmarki\checkmarkr\checkmarkc\checkmark$\checkmarkd\checkmark$\checkmark^\checkmark\\checkmarkc\checkmarki\checkmarkr\checkmarkc\checkmark$\checkmark{\checkmark$\checkmark^\checkmark\\checkmarkc\checkmarki\checkmarkr\checkmarkc\checkmark$\checkmark7\checkmark$\checkmark^\checkmark\\checkmarkc\checkmarki\checkmarkr\checkmarkc\checkmark$\checkmark8\checkmark$\checkmark^\checkmark\\checkmarkc\checkmarki\checkmarkr\checkmarkc\checkmark$\checkmark \checkmark$\checkmark^\checkmark\\checkmarkc\checkmarki\checkmarkr\checkmarkc\checkmark$\checkmark\\checkmark$\checkmark^\checkmark\\checkmarkc\checkmarki\checkmarkr\checkmarkc\checkmark$\checkmarkt\checkmark$\checkmark^\checkmark\\checkmarkc\checkmarki\checkmarkr\checkmarkc\checkmark$\checkmarke\checkmark$\checkmark^\checkmark\\checkmarkc\checkmarki\checkmarkr\checkmarkc\checkmark$\checkmarkx\checkmark$\checkmark^\checkmark\\checkmarkc\checkmarki\checkmarkr\checkmarkc\checkmark$\checkmarkt\checkmark$\checkmark^\checkmark\\checkmarkc\checkmarki\checkmarkr\checkmarkc\checkmark$\checkmark{\checkmark$\checkmark^\checkmark\\checkmarkc\checkmarki\checkmarkr\checkmarkc\checkmark$\checkmark \checkmark$\checkmark^\checkmark\\checkmarkc\checkmarki\checkmarkr\checkmarkc\checkmark$\checkmarkN\checkmark$\checkmark^\checkmark\\checkmarkc\checkmarki\checkmarkr\checkmarkc\checkmark$\checkmark}\checkmark$\checkmark^\checkmark\\checkmarkc\checkmarki\checkmarkr\checkmarkc\checkmark$\checkmark}\checkmark$\checkmark^\checkmark\\checkmarkc\checkmarki\checkmarkr\checkmarkc\checkmark$\checkmark
\checkmark$\checkmark^\checkmark\\checkmarkc\checkmarki\checkmarkr\checkmarkc\checkmark$\checkmark \checkmark$\checkmark^\checkmark\\checkmarkc\checkmarki\checkmarkr\checkmarkc\checkmark$\checkmark \checkmark$\checkmark^\checkmark\\checkmarkc\checkmarki\checkmarkr\checkmarkc\checkmark$\checkmark \checkmark$\checkmark^\checkmark\\checkmarkc\checkmarki\checkmarkr\checkmarkc\checkmark$\checkmark \checkmark$\checkmark^\checkmark\\checkmarkc\checkmarki\checkmarkr\checkmarkc\checkmark$\checkmark\\checkmark$\checkmark^\checkmark\\checkmarkc\checkmarki\checkmarkr\checkmarkc\checkmark$\checkmarkl\checkmark$\checkmark^\checkmark\\checkmarkc\checkmarki\checkmarkr\checkmarkc\checkmark$\checkmarka\checkmark$\checkmark^\checkmark\\checkmarkc\checkmarki\checkmarkr\checkmarkc\checkmark$\checkmarkb\checkmark$\checkmark^\checkmark\\checkmarkc\checkmarki\checkmarkr\checkmarkc\checkmark$\checkmarke\checkmark$\checkmark^\checkmark\\checkmarkc\checkmarki\checkmarkr\checkmarkc\checkmark$\checkmarkl\checkmark$\checkmark^\checkmark\\checkmarkc\checkmarki\checkmarkr\checkmarkc\checkmark$\checkmark{\checkmark$\checkmark^\checkmark\\checkmarkc\checkmarki\checkmarkr\checkmarkc\checkmark$\checkmarke\checkmark$\checkmark^\checkmark\\checkmarkc\checkmarki\checkmarkr\checkmarkc\checkmark$\checkmarkq\checkmark$\checkmark^\checkmark\\checkmarkc\checkmarki\checkmarkr\checkmarkc\checkmark$\checkmark:\checkmark$\checkmark^\checkmark\\checkmarkc\checkmarki\checkmarkr\checkmarkc\checkmark$\checkmarkr\checkmark$\checkmark^\checkmark\\checkmarkc\checkmarki\checkmarkr\checkmarkc\checkmark$\checkmarke\checkmark$\checkmark^\checkmark\\checkmarkc\checkmarki\checkmarkr\checkmarkc\checkmark$\checkmarks\checkmark$\checkmark^\checkmark\\checkmarkc\checkmarki\checkmarkr\checkmarkc\checkmark$\checkmarku\checkmark$\checkmark^\checkmark\\checkmarkc\checkmarki\checkmarkr\checkmarkc\checkmark$\checkmarkl\checkmark$\checkmark^\checkmark\\checkmarkc\checkmarki\checkmarkr\checkmarkc\checkmark$\checkmarkt\checkmark$\checkmark^\checkmark\\checkmarkc\checkmarki\checkmarkr\checkmarkc\checkmark$\checkmarka\checkmark$\checkmark^\checkmark\\checkmarkc\checkmarki\checkmarkr\checkmarkc\checkmark$\checkmarkn\checkmark$\checkmark^\checkmark\\checkmarkc\checkmarki\checkmarkr\checkmarkc\checkmark$\checkmarkt\checkmark$\checkmark^\checkmark\\checkmarkc\checkmarki\checkmarkr\checkmarkc\checkmark$\checkmarke\checkmark$\checkmark^\checkmark\\checkmarkc\checkmarki\checkmarkr\checkmarkc\checkmark$\checkmark}\checkmark$\checkmark^\checkmark\\checkmarkc\checkmarki\checkmarkr\checkmarkc\checkmark$\checkmark
\checkmark$\checkmark^\checkmark\\checkmarkc\checkmarki\checkmarkr\checkmarkc\checkmark$\checkmark\\checkmark$\checkmark^\checkmark\\checkmarkc\checkmarki\checkmarkr\checkmarkc\checkmark$\checkmarke\checkmark$\checkmark^\checkmark\\checkmarkc\checkmarki\checkmarkr\checkmarkc\checkmark$\checkmarkn\checkmark$\checkmark^\checkmark\\checkmarkc\checkmarki\checkmarkr\checkmarkc\checkmark$\checkmarkd\checkmark$\checkmark^\checkmark\\checkmarkc\checkmarki\checkmarkr\checkmarkc\checkmark$\checkmark{\checkmark$\checkmark^\checkmark\\checkmarkc\checkmarki\checkmarkr\checkmarkc\checkmark$\checkmarke\checkmark$\checkmark^\checkmark\\checkmarkc\checkmarki\checkmarkr\checkmarkc\checkmark$\checkmarkq\checkmark$\checkmark^\checkmark\\checkmarkc\checkmarki\checkmarkr\checkmarkc\checkmark$\checkmarku\checkmark$\checkmark^\checkmark\\checkmarkc\checkmarki\checkmarkr\checkmarkc\checkmark$\checkmarka\checkmark$\checkmark^\checkmark\\checkmarkc\checkmarki\checkmarkr\checkmarkc\checkmark$\checkmarkt\checkmark$\checkmark^\checkmark\\checkmarkc\checkmarki\checkmarkr\checkmarkc\checkmark$\checkmarki\checkmark$\checkmark^\checkmark\\checkmarkc\checkmarki\checkmarkr\checkmarkc\checkmark$\checkmarko\checkmark$\checkmark^\checkmark\\checkmarkc\checkmarki\checkmarkr\checkmarkc\checkmark$\checkmarkn\checkmark$\checkmark^\checkmark\\checkmarkc\checkmarki\checkmarkr\checkmarkc\checkmark$\checkmark}\checkmark$\checkmark^\checkmark\\checkmarkc\checkmarki\checkmarkr\checkmarkc\checkmark$\checkmark
\checkmark$\checkmark^\checkmark\\checkmarkc\checkmarki\checkmarkr\checkmarkc\checkmark$\checkmark
\checkmark$\checkmark^\checkmark\\checkmarkc\checkmarki\checkmarkr\checkmarkc\checkmark$\checkmark\\checkmark$\checkmark^\checkmark\\checkmarkc\checkmarki\checkmarkr\checkmarkc\checkmark$\checkmarkt\checkmark$\checkmark^\checkmark\\checkmarkc\checkmarki\checkmarkr\checkmarkc\checkmark$\checkmarke\checkmark$\checkmark^\checkmark\\checkmarkc\checkmarki\checkmarkr\checkmarkc\checkmark$\checkmarkx\checkmark$\checkmark^\checkmark\\checkmarkc\checkmarki\checkmarkr\checkmarkc\checkmark$\checkmarkt\checkmark$\checkmark^\checkmark\\checkmarkc\checkmarki\checkmarkr\checkmarkc\checkmark$\checkmarkb\checkmark$\checkmark^\checkmark\\checkmarkc\checkmarki\checkmarkr\checkmarkc\checkmark$\checkmarkf\checkmark$\checkmark^\checkmark\\checkmarkc\checkmarki\checkmarkr\checkmarkc\checkmark$\checkmark{\checkmark$\checkmark^\checkmark\\checkmarkc\checkmarki\checkmarkr\checkmarkc\checkmark$\checkmarkR\checkmark$\checkmark^\checkmark\\checkmarkc\checkmarki\checkmarkr\checkmarkc\checkmark$\checkmarke\checkmark$\checkmark^\checkmark\\checkmarkc\checkmarki\checkmarkr\checkmarkc\checkmark$\checkmarkm\checkmark$\checkmark^\checkmark\\checkmarkc\checkmarki\checkmarkr\checkmarkc\checkmark$\checkmarka\checkmark$\checkmark^\checkmark\\checkmarkc\checkmarki\checkmarkr\checkmarkc\checkmark$\checkmarkr\checkmark$\checkmark^\checkmark\\checkmarkc\checkmarki\checkmarkr\checkmarkc\checkmark$\checkmarkq\checkmark$\checkmark^\checkmark\\checkmarkc\checkmarki\checkmarkr\checkmarkc\checkmark$\checkmarku\checkmark$\checkmark^\checkmark\\checkmarkc\checkmarki\checkmarkr\checkmarkc\checkmark$\checkmarke\checkmark$\checkmark^\checkmark\\checkmarkc\checkmarki\checkmarkr\checkmarkc\checkmark$\checkmark \checkmark$\checkmark^\checkmark\\checkmarkc\checkmarki\checkmarkr\checkmarkc\checkmark$\checkmark:\checkmark$\checkmark^\checkmark\\checkmarkc\checkmarki\checkmarkr\checkmarkc\checkmark$\checkmark}\checkmark$\checkmark^\checkmark\\checkmarkc\checkmarki\checkmarkr\checkmarkc\checkmark$\checkmark \checkmark$\checkmark^\checkmark\\checkmarkc\checkmarki\checkmarkr\checkmarkc\checkmark$\checkmarkP\checkmark$\checkmark^\checkmark\\checkmarkc\checkmarki\checkmarkr\checkmarkc\checkmark$\checkmarka\checkmark$\checkmark^\checkmark\\checkmarkc\checkmarki\checkmarkr\checkmarkc\checkmark$\checkmarkr\checkmark$\checkmark^\checkmark\\checkmarkc\checkmarki\checkmarkr\checkmarkc\checkmark$\checkmark \checkmark$\checkmark^\checkmark\\checkmarkc\checkmarki\checkmarkr\checkmarkc\checkmark$\checkmarkc\checkmark$\checkmark^\checkmark\\checkmarkc\checkmarki\checkmarkr\checkmarkc\checkmark$\checkmarko\checkmark$\checkmark^\checkmark\\checkmarkc\checkmarki\checkmarkr\checkmarkc\checkmark$\checkmarkn\checkmark$\checkmark^\checkmark\\checkmarkc\checkmarki\checkmarkr\checkmarkc\checkmark$\checkmarkv\checkmark$\checkmark^\checkmark\\checkmarkc\checkmarki\checkmarkr\checkmarkc\checkmark$\checkmarke\checkmark$\checkmark^\checkmark\\checkmarkc\checkmarki\checkmarkr\checkmarkc\checkmark$\checkmarkn\checkmark$\checkmark^\checkmark\\checkmarkc\checkmarki\checkmarkr\checkmarkc\checkmark$\checkmarkt\checkmark$\checkmark^\checkmark\\checkmarkc\checkmarki\checkmarkr\checkmarkc\checkmark$\checkmarki\checkmark$\checkmark^\checkmark\\checkmarkc\checkmarki\checkmarkr\checkmarkc\checkmark$\checkmarko\checkmark$\checkmark^\checkmark\\checkmarkc\checkmarki\checkmarkr\checkmarkc\checkmark$\checkmarkn\checkmark$\checkmark^\checkmark\\checkmarkc\checkmarki\checkmarkr\checkmarkc\checkmark$\checkmark \checkmark$\checkmark^\checkmark\\checkmarkc\checkmarki\checkmarkr\checkmarkc\checkmark$\checkmarkd\checkmark$\checkmark^\checkmark\\checkmarkc\checkmarki\checkmarkr\checkmarkc\checkmark$\checkmarke\checkmark$\checkmark^\checkmark\\checkmarkc\checkmarki\checkmarkr\checkmarkc\checkmark$\checkmark \checkmark$\checkmark^\checkmark\\checkmarkc\checkmarki\checkmarkr\checkmarkc\checkmark$\checkmarks\checkmark$\checkmark^\checkmark\\checkmarkc\checkmarki\checkmarkr\checkmarkc\checkmark$\checkmarké\checkmark$\checkmark^\checkmark\\checkmarkc\checkmarki\checkmarkr\checkmarkc\checkmark$\checkmarkc\checkmark$\checkmark^\checkmark\\checkmarkc\checkmarki\checkmarkr\checkmarkc\checkmark$\checkmarku\checkmark$\checkmark^\checkmark\\checkmarkc\checkmarki\checkmarkr\checkmarkc\checkmark$\checkmarkr\checkmark$\checkmark^\checkmark\\checkmarkc\checkmarki\checkmarkr\checkmarkc\checkmark$\checkmarki\checkmark$\checkmark^\checkmark\\checkmarkc\checkmarki\checkmarkr\checkmarkc\checkmark$\checkmarkt\checkmark$\checkmark^\checkmark\\checkmarkc\checkmarki\checkmarkr\checkmarkc\checkmark$\checkmarké\checkmark$\checkmark^\checkmark\\checkmarkc\checkmarki\checkmarkr\checkmarkc\checkmark$\checkmark \checkmark$\checkmark^\checkmark\\checkmarkc\checkmarki\checkmarkr\checkmarkc\checkmark$\checkmark(\checkmark$\checkmark^\checkmark\\checkmarkc\checkmarki\checkmarkr\checkmarkc\checkmark$\checkmark$\checkmark$\checkmark^\checkmark\\checkmarkc\checkmarki\checkmarkr\checkmarkc\checkmark$\checkmarks\checkmark$\checkmark^\checkmark\\checkmarkc\checkmarki\checkmarkr\checkmarkc\checkmark$\checkmark \checkmark$\checkmark^\checkmark\\checkmarkc\checkmarki\checkmarkr\checkmarkc\checkmark$\checkmark=\checkmark$\checkmark^\checkmark\\checkmarkc\checkmarki\checkmarkr\checkmarkc\checkmark$\checkmark \checkmark$\checkmark^\checkmark\\checkmarkc\checkmarki\checkmarkr\checkmarkc\checkmark$\checkmark1\checkmark$\checkmark^\checkmark\\checkmarkc\checkmarki\checkmarkr\checkmarkc\checkmark$\checkmark{\checkmark$\checkmark^\checkmark\\checkmarkc\checkmarki\checkmarkr\checkmarkc\checkmark$\checkmark,\checkmark$\checkmark^\checkmark\\checkmarkc\checkmarki\checkmarkr\checkmarkc\checkmark$\checkmark}\checkmark$\checkmark^\checkmark\\checkmarkc\checkmarki\checkmarkr\checkmarkc\checkmark$\checkmark2\checkmark$\checkmark^\checkmark\\checkmarkc\checkmarki\checkmarkr\checkmarkc\checkmark$\checkmark8\checkmark$\checkmark^\checkmark\\checkmarkc\checkmarki\checkmarkr\checkmarkc\checkmark$\checkmark$\checkmark$\checkmark^\checkmark\\checkmarkc\checkmarki\checkmarkr\checkmarkc\checkmark$\checkmark)\checkmark$\checkmark^\checkmark\\checkmarkc\checkmarki\checkmarkr\checkmarkc\checkmark$\checkmark,\checkmark$\checkmark^\checkmark\\checkmarkc\checkmarki\checkmarkr\checkmarkc\checkmark$\checkmark \checkmark$\checkmark^\checkmark\\checkmarkc\checkmarki\checkmarkr\checkmarkc\checkmark$\checkmarkn\checkmark$\checkmark^\checkmark\\checkmarkc\checkmarki\checkmarkr\checkmarkc\checkmark$\checkmarko\checkmark$\checkmark^\checkmark\\checkmarkc\checkmarki\checkmarkr\checkmarkc\checkmark$\checkmarku\checkmark$\checkmark^\checkmark\\checkmarkc\checkmarki\checkmarkr\checkmarkc\checkmark$\checkmarks\checkmark$\checkmark^\checkmark\\checkmarkc\checkmarki\checkmarkr\checkmarkc\checkmark$\checkmark \checkmark$\checkmark^\checkmark\\checkmarkc\checkmarki\checkmarkr\checkmarkc\checkmark$\checkmarka\checkmark$\checkmark^\checkmark\\checkmarkc\checkmarki\checkmarkr\checkmarkc\checkmark$\checkmarkd\checkmark$\checkmark^\checkmark\\checkmarkc\checkmarki\checkmarkr\checkmarkc\checkmark$\checkmarko\checkmark$\checkmark^\checkmark\\checkmarkc\checkmarki\checkmarkr\checkmarkc\checkmark$\checkmarkp\checkmark$\checkmark^\checkmark\\checkmarkc\checkmarki\checkmarkr\checkmarkc\checkmark$\checkmarkt\checkmark$\checkmark^\checkmark\\checkmarkc\checkmarki\checkmarkr\checkmarkc\checkmark$\checkmarko\checkmark$\checkmark^\checkmark\\checkmarkc\checkmarki\checkmarkr\checkmarkc\checkmark$\checkmarkn\checkmark$\checkmark^\checkmark\\checkmarkc\checkmarki\checkmarkr\checkmarkc\checkmark$\checkmarks\checkmark$\checkmark^\checkmark\\checkmarkc\checkmarki\checkmarkr\checkmarkc\checkmark$\checkmark \checkmark$\checkmark^\checkmark\\checkmarkc\checkmarki\checkmarkr\checkmarkc\checkmark$\checkmark$\checkmark$\checkmark^\checkmark\\checkmarkc\checkmarki\checkmarkr\checkmarkc\checkmark$\checkmark\\checkmark$\checkmark^\checkmark\\checkmarkc\checkmarki\checkmarkr\checkmarkc\checkmark$\checkmark|\checkmark$\checkmark^\checkmark\\checkmarkc\checkmarki\checkmarkr\checkmarkc\checkmark$\checkmark\\checkmark$\checkmark^\checkmark\\checkmarkc\checkmarki\checkmarkr\checkmarkc\checkmark$\checkmarkv\checkmark$\checkmark^\checkmark\\checkmarkc\checkmarki\checkmarkr\checkmarkc\checkmark$\checkmarke\checkmark$\checkmark^\checkmark\\checkmarkc\checkmarki\checkmarkr\checkmarkc\checkmark$\checkmarkc\checkmark$\checkmark^\checkmark\\checkmarkc\checkmarki\checkmarkr\checkmarkc\checkmark$\checkmark{\checkmark$\checkmark^\checkmark\\checkmarkc\checkmarki\checkmarkr\checkmarkc\checkmark$\checkmarkR\checkmark$\checkmark^\checkmark\\checkmarkc\checkmarki\checkmarkr\checkmarkc\checkmark$\checkmark}\checkmark$\checkmark^\checkmark\\checkmarkc\checkmarki\checkmarkr\checkmarkc\checkmark$\checkmark\\checkmark$\checkmark^\checkmark\\checkmarkc\checkmarki\checkmarkr\checkmarkc\checkmark$\checkmark|\checkmark$\checkmark^\checkmark\\checkmarkc\checkmarki\checkmarkr\checkmarkc\checkmark$\checkmark \checkmark$\checkmark^\checkmark\\checkmarkc\checkmarki\checkmarkr\checkmarkc\checkmark$\checkmark=\checkmark$\checkmark^\checkmark\\checkmarkc\checkmarki\checkmarkr\checkmarkc\checkmark$\checkmark \checkmark$\checkmark^\checkmark\\checkmarkc\checkmarki\checkmarkr\checkmarkc\checkmark$\checkmark\\checkmark$\checkmark^\checkmark\\checkmarkc\checkmarki\checkmarkr\checkmarkc\checkmark$\checkmarkm\checkmark$\checkmark^\checkmark\\checkmarkc\checkmarki\checkmarkr\checkmarkc\checkmark$\checkmarka\checkmark$\checkmark^\checkmark\\checkmarkc\checkmarki\checkmarkr\checkmarkc\checkmark$\checkmarkt\checkmark$\checkmark^\checkmark\\checkmarkc\checkmarki\checkmarkr\checkmarkc\checkmark$\checkmarkh\checkmark$\checkmark^\checkmark\\checkmarkc\checkmarki\checkmarkr\checkmarkc\checkmark$\checkmarkb\checkmark$\checkmark^\checkmark\\checkmarkc\checkmarki\checkmarkr\checkmarkc\checkmark$\checkmarkf\checkmark$\checkmark^\checkmark\\checkmarkc\checkmarki\checkmarkr\checkmarkc\checkmark$\checkmark{\checkmark$\checkmark^\checkmark\\checkmarkc\checkmarki\checkmarkr\checkmarkc\checkmark$\checkmarkF\checkmark$\checkmark^\checkmark\\checkmarkc\checkmarki\checkmarkr\checkmarkc\checkmark$\checkmark_\checkmark$\checkmark^\checkmark\\checkmarkc\checkmarki\checkmarkr\checkmarkc\checkmark$\checkmark{\checkmark$\checkmark^\checkmark\\checkmarkc\checkmarki\checkmarkr\checkmarkc\checkmark$\checkmarkm\checkmark$\checkmark^\checkmark\\checkmarkc\checkmarki\checkmarkr\checkmarkc\checkmark$\checkmarka\checkmark$\checkmark^\checkmark\\checkmarkc\checkmarki\checkmarkr\checkmarkc\checkmark$\checkmarkx\checkmark$\checkmark^\checkmark\\checkmarkc\checkmarki\checkmarkr\checkmarkc\checkmark$\checkmark}\checkmark$\checkmark^\checkmark\\checkmarkc\checkmarki\checkmarkr\checkmarkc\checkmark$\checkmark \checkmark$\checkmark^\checkmark\\checkmarkc\checkmarki\checkmarkr\checkmarkc\checkmark$\checkmark=\checkmark$\checkmark^\checkmark\\checkmarkc\checkmarki\checkmarkr\checkmarkc\checkmark$\checkmark \checkmark$\checkmark^\checkmark\\checkmarkc\checkmarki\checkmarkr\checkmarkc\checkmark$\checkmark1\checkmark$\checkmark^\checkmark\\checkmarkc\checkmarki\checkmarkr\checkmarkc\checkmark$\checkmark0\checkmark$\checkmark^\checkmark\\checkmarkc\checkmarki\checkmarkr\checkmarkc\checkmark$\checkmark0\checkmark$\checkmark^\checkmark\\checkmarkc\checkmarki\checkmarkr\checkmarkc\checkmark$\checkmark \checkmark$\checkmark^\checkmark\\checkmarkc\checkmarki\checkmarkr\checkmarkc\checkmark$\checkmark\\checkmark$\checkmark^\checkmark\\checkmarkc\checkmarki\checkmarkr\checkmarkc\checkmark$\checkmarkt\checkmark$\checkmark^\checkmark\\checkmarkc\checkmarki\checkmarkr\checkmarkc\checkmark$\checkmarke\checkmark$\checkmark^\checkmark\\checkmarkc\checkmarki\checkmarkr\checkmarkc\checkmark$\checkmarkx\checkmark$\checkmark^\checkmark\\checkmarkc\checkmarki\checkmarkr\checkmarkc\checkmark$\checkmarkt\checkmark$\checkmark^\checkmark\\checkmarkc\checkmarki\checkmarkr\checkmarkc\checkmark$\checkmark{\checkmark$\checkmark^\checkmark\\checkmarkc\checkmarki\checkmarkr\checkmarkc\checkmark$\checkmark \checkmark$\checkmark^\checkmark\\checkmarkc\checkmarki\checkmarkr\checkmarkc\checkmark$\checkmarkN\checkmark$\checkmark^\checkmark\\checkmarkc\checkmarki\checkmarkr\checkmarkc\checkmark$\checkmark}\checkmark$\checkmark^\checkmark\\checkmarkc\checkmarki\checkmarkr\checkmarkc\checkmark$\checkmark}\checkmark$\checkmark^\checkmark\\checkmarkc\checkmarki\checkmarkr\checkmarkc\checkmark$\checkmark$\checkmark$\checkmark^\checkmark\\checkmarkc\checkmarki\checkmarkr\checkmarkc\checkmark$\checkmark \checkmark$\checkmark^\checkmark\\checkmarkc\checkmarki\checkmarkr\checkmarkc\checkmark$\checkmarkc\checkmark$\checkmark^\checkmark\\checkmarkc\checkmarki\checkmarkr\checkmarkc\checkmark$\checkmarko\checkmark$\checkmark^\checkmark\\checkmarkc\checkmarki\checkmarkr\checkmarkc\checkmark$\checkmarkm\checkmark$\checkmark^\checkmark\\checkmarkc\checkmarki\checkmarkr\checkmarkc\checkmark$\checkmarkm\checkmark$\checkmark^\checkmark\\checkmarkc\checkmarki\checkmarkr\checkmarkc\checkmark$\checkmarke\checkmark$\checkmark^\checkmark\\checkmarkc\checkmarki\checkmarkr\checkmarkc\checkmark$\checkmark \checkmark$\checkmark^\checkmark\\checkmarkc\checkmarki\checkmarkr\checkmarkc\checkmark$\checkmarkv\checkmark$\checkmark^\checkmark\\checkmarkc\checkmarki\checkmarkr\checkmarkc\checkmark$\checkmarka\checkmark$\checkmark^\checkmark\\checkmarkc\checkmarki\checkmarkr\checkmarkc\checkmark$\checkmarkl\checkmark$\checkmark^\checkmark\\checkmarkc\checkmarki\checkmarkr\checkmarkc\checkmark$\checkmarke\checkmark$\checkmark^\checkmark\\checkmarkc\checkmarki\checkmarkr\checkmarkc\checkmark$\checkmarku\checkmark$\checkmark^\checkmark\\checkmarkc\checkmarki\checkmarkr\checkmarkc\checkmark$\checkmarkr\checkmark$\checkmark^\checkmark\\checkmarkc\checkmarki\checkmarkr\checkmarkc\checkmark$\checkmark \checkmark$\checkmark^\checkmark\\checkmarkc\checkmarki\checkmarkr\checkmarkc\checkmark$\checkmarkd\checkmark$\checkmark^\checkmark\\checkmarkc\checkmarki\checkmarkr\checkmarkc\checkmark$\checkmarke\checkmark$\checkmark^\checkmark\\checkmarkc\checkmarki\checkmarkr\checkmarkc\checkmark$\checkmark \checkmark$\checkmark^\checkmark\\checkmarkc\checkmarki\checkmarkr\checkmarkc\checkmark$\checkmarkd\checkmark$\checkmark^\checkmark\\checkmarkc\checkmarki\checkmarkr\checkmarkc\checkmark$\checkmarki\checkmark$\checkmark^\checkmark\\checkmarkc\checkmarki\checkmarkr\checkmarkc\checkmark$\checkmarkm\checkmark$\checkmark^\checkmark\\checkmarkc\checkmarki\checkmarkr\checkmarkc\checkmark$\checkmarke\checkmark$\checkmark^\checkmark\\checkmarkc\checkmarki\checkmarkr\checkmarkc\checkmark$\checkmarkn\checkmark$\checkmark^\checkmark\\checkmarkc\checkmarki\checkmarkr\checkmarkc\checkmark$\checkmarks\checkmark$\checkmark^\checkmark\\checkmarkc\checkmarki\checkmarkr\checkmarkc\checkmark$\checkmarki\checkmark$\checkmark^\checkmark\\checkmarkc\checkmarki\checkmarkr\checkmarkc\checkmark$\checkmarko\checkmark$\checkmark^\checkmark\\checkmarkc\checkmarki\checkmarkr\checkmarkc\checkmark$\checkmarkn\checkmark$\checkmark^\checkmark\\checkmarkc\checkmarki\checkmarkr\checkmarkc\checkmark$\checkmarkn\checkmark$\checkmark^\checkmark\\checkmarkc\checkmarki\checkmarkr\checkmarkc\checkmark$\checkmarke\checkmark$\checkmark^\checkmark\\checkmarkc\checkmarki\checkmarkr\checkmarkc\checkmark$\checkmarkm\checkmark$\checkmark^\checkmark\\checkmarkc\checkmarki\checkmarkr\checkmarkc\checkmark$\checkmarke\checkmark$\checkmark^\checkmark\\checkmarkc\checkmarki\checkmarkr\checkmarkc\checkmark$\checkmarkn\checkmark$\checkmark^\checkmark\\checkmarkc\checkmarki\checkmarkr\checkmarkc\checkmark$\checkmarkt\checkmark$\checkmark^\checkmark\\checkmarkc\checkmarki\checkmarkr\checkmarkc\checkmark$\checkmark \checkmark$\checkmark^\checkmark\\checkmarkc\checkmarki\checkmarkr\checkmarkc\checkmark$\checkmarkp\checkmark$\checkmark^\checkmark\\checkmarkc\checkmarki\checkmarkr\checkmarkc\checkmark$\checkmarko\checkmark$\checkmark^\checkmark\\checkmarkc\checkmarki\checkmarkr\checkmarkc\checkmark$\checkmarku\checkmark$\checkmark^\checkmark\\checkmarkc\checkmarki\checkmarkr\checkmarkc\checkmark$\checkmarkr\checkmark$\checkmark^\checkmark\\checkmarkc\checkmarki\checkmarkr\checkmarkc\checkmark$\checkmark \checkmark$\checkmark^\checkmark\\checkmarkc\checkmarki\checkmarkr\checkmarkc\checkmark$\checkmarkt\checkmark$\checkmark^\checkmark\\checkmarkc\checkmarki\checkmarkr\checkmarkc\checkmark$\checkmarko\checkmark$\checkmark^\checkmark\\checkmarkc\checkmarki\checkmarkr\checkmarkc\checkmark$\checkmarku\checkmark$\checkmark^\checkmark\\checkmarkc\checkmarki\checkmarkr\checkmarkc\checkmark$\checkmarks\checkmark$\checkmark^\checkmark\\checkmarkc\checkmarki\checkmarkr\checkmarkc\checkmark$\checkmark \checkmark$\checkmark^\checkmark\\checkmarkc\checkmarki\checkmarkr\checkmarkc\checkmark$\checkmarkl\checkmark$\checkmark^\checkmark\\checkmarkc\checkmarki\checkmarkr\checkmarkc\checkmark$\checkmarke\checkmark$\checkmark^\checkmark\\checkmarkc\checkmarki\checkmarkr\checkmarkc\checkmark$\checkmarks\checkmark$\checkmark^\checkmark\\checkmarkc\checkmarki\checkmarkr\checkmarkc\checkmark$\checkmark \checkmark$\checkmark^\checkmark\\checkmarkc\checkmarki\checkmarkr\checkmarkc\checkmark$\checkmarkc\checkmark$\checkmark^\checkmark\\checkmarkc\checkmarki\checkmarkr\checkmarkc\checkmark$\checkmarka\checkmark$\checkmark^\checkmark\\checkmarkc\checkmarki\checkmarkr\checkmarkc\checkmark$\checkmarkl\checkmark$\checkmark^\checkmark\\checkmarkc\checkmarki\checkmarkr\checkmarkc\checkmark$\checkmarkc\checkmark$\checkmark^\checkmark\\checkmarkc\checkmarki\checkmarkr\checkmarkc\checkmark$\checkmarku\checkmark$\checkmark^\checkmark\\checkmarkc\checkmarki\checkmarkr\checkmarkc\checkmark$\checkmarkl\checkmark$\checkmark^\checkmark\\checkmarkc\checkmarki\checkmarkr\checkmarkc\checkmark$\checkmarks\checkmark$\checkmark^\checkmark\\checkmarkc\checkmarki\checkmarkr\checkmarkc\checkmark$\checkmark \checkmark$\checkmark^\checkmark\\checkmarkc\checkmarki\checkmarkr\checkmarkc\checkmark$\checkmarkq\checkmark$\checkmark^\checkmark\\checkmarkc\checkmarki\checkmarkr\checkmarkc\checkmark$\checkmarku\checkmark$\checkmark^\checkmark\\checkmarkc\checkmarki\checkmarkr\checkmarkc\checkmark$\checkmarki\checkmark$\checkmark^\checkmark\\checkmarkc\checkmarki\checkmarkr\checkmarkc\checkmark$\checkmark \checkmark$\checkmark^\checkmark\\checkmarkc\checkmarki\checkmarkr\checkmarkc\checkmark$\checkmarks\checkmark$\checkmark^\checkmark\\checkmarkc\checkmarki\checkmarkr\checkmarkc\checkmark$\checkmarku\checkmark$\checkmark^\checkmark\\checkmarkc\checkmarki\checkmarkr\checkmarkc\checkmark$\checkmarki\checkmark$\checkmark^\checkmark\\checkmarkc\checkmarki\checkmarkr\checkmarkc\checkmark$\checkmarkv\checkmark$\checkmark^\checkmark\\checkmarkc\checkmarki\checkmarkr\checkmarkc\checkmark$\checkmarke\checkmark$\checkmark^\checkmark\\checkmarkc\checkmarki\checkmarkr\checkmarkc\checkmark$\checkmarkn\checkmark$\checkmark^\checkmark\\checkmarkc\checkmarki\checkmarkr\checkmarkc\checkmark$\checkmarkt\checkmark$\checkmark^\checkmark\\checkmarkc\checkmarki\checkmarkr\checkmarkc\checkmark$\checkmark.\checkmark$\checkmark^\checkmark\\checkmarkc\checkmarki\checkmarkr\checkmarkc\checkmark$\checkmark
\checkmark$\checkmark^\checkmark\\checkmarkc\checkmarki\checkmarkr\checkmarkc\checkmark$\checkmark
\checkmark$\checkmark^\checkmark\\checkmarkc\checkmarki\checkmarkr\checkmarkc\checkmark$\checkmark\\checkmark$\checkmark^\checkmark\\checkmarkc\checkmarki\checkmarkr\checkmarkc\checkmark$\checkmarks\checkmark$\checkmark^\checkmark\\checkmarkc\checkmarki\checkmarkr\checkmarkc\checkmark$\checkmarku\checkmark$\checkmark^\checkmark\\checkmarkc\checkmarki\checkmarkr\checkmarkc\checkmark$\checkmarkb\checkmark$\checkmark^\checkmark\\checkmarkc\checkmarki\checkmarkr\checkmarkc\checkmark$\checkmarks\checkmark$\checkmark^\checkmark\\checkmarkc\checkmarki\checkmarkr\checkmarkc\checkmark$\checkmarke\checkmark$\checkmark^\checkmark\\checkmarkc\checkmarki\checkmarkr\checkmarkc\checkmark$\checkmarkc\checkmark$\checkmark^\checkmark\\checkmarkc\checkmarki\checkmarkr\checkmarkc\checkmark$\checkmarkt\checkmark$\checkmark^\checkmark\\checkmarkc\checkmarki\checkmarkr\checkmarkc\checkmark$\checkmarki\checkmark$\checkmark^\checkmark\\checkmarkc\checkmarki\checkmarkr\checkmarkc\checkmark$\checkmarko\checkmark$\checkmark^\checkmark\\checkmarkc\checkmarki\checkmarkr\checkmarkc\checkmark$\checkmarkn\checkmark$\checkmark^\checkmark\\checkmarkc\checkmarki\checkmarkr\checkmarkc\checkmark$\checkmark{\checkmark$\checkmark^\checkmark\\checkmarkc\checkmarki\checkmarkr\checkmarkc\checkmark$\checkmarkT\checkmark$\checkmark^\checkmark\\checkmarkc\checkmarki\checkmarkr\checkmarkc\checkmark$\checkmarko\checkmark$\checkmark^\checkmark\\checkmarkc\checkmarki\checkmarkr\checkmarkc\checkmark$\checkmarkr\checkmark$\checkmark^\checkmark\\checkmarkc\checkmarki\checkmarkr\checkmarkc\checkmark$\checkmarks\checkmark$\checkmark^\checkmark\\checkmarkc\checkmarki\checkmarkr\checkmarkc\checkmark$\checkmarke\checkmark$\checkmark^\checkmark\\checkmarkc\checkmarki\checkmarkr\checkmarkc\checkmark$\checkmarku\checkmark$\checkmark^\checkmark\\checkmarkc\checkmarki\checkmarkr\checkmarkc\checkmark$\checkmarkr\checkmark$\checkmark^\checkmark\\checkmarkc\checkmarki\checkmarkr\checkmarkc\checkmark$\checkmark \checkmark$\checkmark^\checkmark\\checkmarkc\checkmarki\checkmarkr\checkmarkc\checkmark$\checkmarkd\checkmark$\checkmark^\checkmark\\checkmarkc\checkmarki\checkmarkr\checkmarkc\checkmark$\checkmarke\checkmark$\checkmark^\checkmark\\checkmarkc\checkmarki\checkmarkr\checkmarkc\checkmark$\checkmarks\checkmark$\checkmark^\checkmark\\checkmarkc\checkmarki\checkmarkr\checkmarkc\checkmark$\checkmark \checkmark$\checkmark^\checkmark\\checkmarkc\checkmarki\checkmarkr\checkmarkc\checkmark$\checkmarka\checkmark$\checkmark^\checkmark\\checkmarkc\checkmarki\checkmarkr\checkmarkc\checkmark$\checkmarkc\checkmark$\checkmark^\checkmark\\checkmarkc\checkmarki\checkmarkr\checkmarkc\checkmark$\checkmarkt\checkmark$\checkmark^\checkmark\\checkmarkc\checkmarki\checkmarkr\checkmarkc\checkmark$\checkmarki\checkmark$\checkmark^\checkmark\\checkmarkc\checkmarki\checkmarkr\checkmarkc\checkmark$\checkmarko\checkmark$\checkmark^\checkmark\\checkmarkc\checkmarki\checkmarkr\checkmarkc\checkmark$\checkmarkn\checkmark$\checkmark^\checkmark\\checkmarkc\checkmarki\checkmarkr\checkmarkc\checkmark$\checkmarks\checkmark$\checkmark^\checkmark\\checkmarkc\checkmarki\checkmarkr\checkmarkc\checkmark$\checkmark \checkmark$\checkmark^\checkmark\\checkmarkc\checkmarki\checkmarkr\checkmarkc\checkmark$\checkmarkm\checkmark$\checkmark^\checkmark\\checkmarkc\checkmarki\checkmarkr\checkmarkc\checkmark$\checkmarké\checkmark$\checkmark^\checkmark\\checkmarkc\checkmarki\checkmarkr\checkmarkc\checkmark$\checkmarkc\checkmark$\checkmark^\checkmark\\checkmarkc\checkmarki\checkmarkr\checkmarkc\checkmark$\checkmarka\checkmark$\checkmark^\checkmark\\checkmarkc\checkmarki\checkmarkr\checkmarkc\checkmark$\checkmarkn\checkmark$\checkmark^\checkmark\\checkmarkc\checkmarki\checkmarkr\checkmarkc\checkmark$\checkmarki\checkmark$\checkmark^\checkmark\\checkmarkc\checkmarki\checkmarkr\checkmarkc\checkmark$\checkmarkq\checkmark$\checkmark^\checkmark\\checkmarkc\checkmarki\checkmarkr\checkmarkc\checkmark$\checkmarku\checkmark$\checkmark^\checkmark\\checkmarkc\checkmarki\checkmarkr\checkmarkc\checkmark$\checkmarke\checkmark$\checkmark^\checkmark\\checkmarkc\checkmarki\checkmarkr\checkmarkc\checkmark$\checkmarks\checkmark$\checkmark^\checkmark\\checkmarkc\checkmarki\checkmarkr\checkmarkc\checkmark$\checkmark}\checkmark$\checkmark^\checkmark\\checkmarkc\checkmarki\checkmarkr\checkmarkc\checkmark$\checkmark
\checkmark$\checkmark^\checkmark\\checkmarkc\checkmarki\checkmarkr\checkmarkc\checkmark$\checkmarkL\checkmark$\checkmark^\checkmark\\checkmarkc\checkmarki\checkmarkr\checkmarkc\checkmark$\checkmark'\checkmark$\checkmark^\checkmark\\checkmarkc\checkmarki\checkmarkr\checkmarkc\checkmark$\checkmarka\checkmark$\checkmark^\checkmark\\checkmarkc\checkmarki\checkmarkr\checkmarkc\checkmark$\checkmarkc\checkmark$\checkmark^\checkmark\\checkmarkc\checkmarki\checkmarkr\checkmarkc\checkmark$\checkmarkt\checkmark$\checkmark^\checkmark\\checkmarkc\checkmarki\checkmarkr\checkmarkc\checkmark$\checkmarki\checkmark$\checkmark^\checkmark\\checkmarkc\checkmarki\checkmarkr\checkmarkc\checkmark$\checkmarko\checkmark$\checkmark^\checkmark\\checkmarkc\checkmarki\checkmarkr\checkmarkc\checkmark$\checkmarkn\checkmark$\checkmark^\checkmark\\checkmarkc\checkmarki\checkmarkr\checkmarkc\checkmark$\checkmark \checkmark$\checkmark^\checkmark\\checkmarkc\checkmarki\checkmarkr\checkmarkc\checkmark$\checkmarkd\checkmark$\checkmark^\checkmark\\checkmarkc\checkmarki\checkmarkr\checkmarkc\checkmark$\checkmarke\checkmark$\checkmark^\checkmark\\checkmarkc\checkmarki\checkmarkr\checkmarkc\checkmark$\checkmark \checkmark$\checkmark^\checkmark\\checkmarkc\checkmarki\checkmarkr\checkmarkc\checkmark$\checkmarkl\checkmark$\checkmark^\checkmark\\checkmarkc\checkmarki\checkmarkr\checkmarkc\checkmark$\checkmark'\checkmark$\checkmark^\checkmark\\checkmarkc\checkmarki\checkmarkr\checkmarkc\checkmark$\checkmarko\checkmark$\checkmark^\checkmark\\checkmarkc\checkmarki\checkmarkr\checkmarkc\checkmark$\checkmarku\checkmark$\checkmark^\checkmark\\checkmarkc\checkmarki\checkmarkr\checkmarkc\checkmark$\checkmarkt\checkmark$\checkmark^\checkmark\\checkmarkc\checkmarki\checkmarkr\checkmarkc\checkmark$\checkmarki\checkmark$\checkmark^\checkmark\\checkmarkc\checkmarki\checkmarkr\checkmarkc\checkmark$\checkmarkl\checkmark$\checkmark^\checkmark\\checkmarkc\checkmarki\checkmarkr\checkmarkc\checkmark$\checkmark \checkmark$\checkmark^\checkmark\\checkmarkc\checkmarki\checkmarkr\checkmarkc\checkmark$\checkmarks\checkmark$\checkmark^\checkmark\\checkmarkc\checkmarki\checkmarkr\checkmarkc\checkmark$\checkmarku\checkmark$\checkmark^\checkmark\\checkmarkc\checkmarki\checkmarkr\checkmarkc\checkmark$\checkmarkr\checkmark$\checkmark^\checkmark\\checkmarkc\checkmarki\checkmarkr\checkmarkc\checkmark$\checkmark \checkmark$\checkmark^\checkmark\\checkmarkc\checkmarki\checkmarkr\checkmarkc\checkmark$\checkmarkl\checkmark$\checkmark^\checkmark\\checkmarkc\checkmarki\checkmarkr\checkmarkc\checkmark$\checkmarka\checkmark$\checkmark^\checkmark\\checkmarkc\checkmarki\checkmarkr\checkmarkc\checkmark$\checkmark \checkmark$\checkmark^\checkmark\\checkmarkc\checkmarki\checkmarkr\checkmarkc\checkmark$\checkmarkp\checkmark$\checkmark^\checkmark\\checkmarkc\checkmarki\checkmarkr\checkmarkc\checkmark$\checkmarki\checkmark$\checkmark^\checkmark\\checkmarkc\checkmarki\checkmarkr\checkmarkc\checkmark$\checkmarkè\checkmark$\checkmark^\checkmark\\checkmarkc\checkmarki\checkmarkr\checkmarkc\checkmark$\checkmarkc\checkmark$\checkmark^\checkmark\\checkmarkc\checkmarki\checkmarkr\checkmarkc\checkmark$\checkmarke\checkmark$\checkmark^\checkmark\\checkmarkc\checkmarki\checkmarkr\checkmarkc\checkmark$\checkmark \checkmark$\checkmark^\checkmark\\checkmarkc\checkmarki\checkmarkr\checkmarkc\checkmark$\checkmarka\checkmark$\checkmark^\checkmark\\checkmarkc\checkmarki\checkmarkr\checkmarkc\checkmark$\checkmarku\checkmark$\checkmark^\checkmark\\checkmarkc\checkmarki\checkmarkr\checkmarkc\checkmark$\checkmark \checkmark$\checkmark^\checkmark\\checkmarkc\checkmarki\checkmarkr\checkmarkc\checkmark$\checkmarkp\checkmark$\checkmark^\checkmark\\checkmarkc\checkmarki\checkmarkr\checkmarkc\checkmark$\checkmarko\checkmark$\checkmark^\checkmark\\checkmarkc\checkmarki\checkmarkr\checkmarkc\checkmark$\checkmarki\checkmark$\checkmark^\checkmark\\checkmarkc\checkmarki\checkmarkr\checkmarkc\checkmark$\checkmarkn\checkmark$\checkmark^\checkmark\\checkmarkc\checkmarki\checkmarkr\checkmarkc\checkmark$\checkmarkt\checkmark$\checkmark^\checkmark\\checkmarkc\checkmarki\checkmarkr\checkmarkc\checkmark$\checkmark \checkmark$\checkmark^\checkmark\\checkmarkc\checkmarki\checkmarkr\checkmarkc\checkmark$\checkmarkd\checkmark$\checkmark^\checkmark\\checkmarkc\checkmarki\checkmarkr\checkmarkc\checkmark$\checkmarke\checkmark$\checkmark^\checkmark\\checkmarkc\checkmarki\checkmarkr\checkmarkc\checkmark$\checkmark \checkmark$\checkmark^\checkmark\\checkmarkc\checkmarki\checkmarkr\checkmarkc\checkmark$\checkmarkc\checkmark$\checkmark^\checkmark\\checkmarkc\checkmarki\checkmarkr\checkmarkc\checkmark$\checkmarko\checkmark$\checkmark^\checkmark\\checkmarkc\checkmarki\checkmarkr\checkmarkc\checkmark$\checkmarkn\checkmark$\checkmark^\checkmark\\checkmarkc\checkmarki\checkmarkr\checkmarkc\checkmark$\checkmarkt\checkmark$\checkmark^\checkmark\\checkmarkc\checkmarki\checkmarkr\checkmarkc\checkmark$\checkmarka\checkmark$\checkmark^\checkmark\\checkmarkc\checkmarki\checkmarkr\checkmarkc\checkmark$\checkmarkc\checkmark$\checkmark^\checkmark\\checkmarkc\checkmarki\checkmarkr\checkmarkc\checkmark$\checkmarkt\checkmark$\checkmark^\checkmark\\checkmarkc\checkmarki\checkmarkr\checkmarkc\checkmark$\checkmark \checkmark$\checkmark^\checkmark\\checkmarkc\checkmarki\checkmarkr\checkmarkc\checkmark$\checkmark$\checkmark$\checkmark^\checkmark\\checkmarkc\checkmarki\checkmarkr\checkmarkc\checkmark$\checkmarkO\checkmark$\checkmark^\checkmark\\checkmarkc\checkmarki\checkmarkr\checkmarkc\checkmark$\checkmark$\checkmark$\checkmark^\checkmark\\checkmarkc\checkmarki\checkmarkr\checkmarkc\checkmark$\checkmark,\checkmark$\checkmark^\checkmark\\checkmarkc\checkmarki\checkmarkr\checkmarkc\checkmark$\checkmark \checkmark$\checkmark^\checkmark\\checkmarkc\checkmarki\checkmarkr\checkmarkc\checkmark$\checkmarkd\checkmark$\checkmark^\checkmark\\checkmarkc\checkmarki\checkmarkr\checkmarkc\checkmark$\checkmarka\checkmark$\checkmark^\checkmark\\checkmarkc\checkmarki\checkmarkr\checkmarkc\checkmark$\checkmarkn\checkmark$\checkmark^\checkmark\\checkmarkc\checkmarki\checkmarkr\checkmarkc\checkmark$\checkmarks\checkmark$\checkmark^\checkmark\\checkmarkc\checkmarki\checkmarkr\checkmarkc\checkmark$\checkmark \checkmark$\checkmark^\checkmark\\checkmarkc\checkmarki\checkmarkr\checkmarkc\checkmark$\checkmarkl\checkmark$\checkmark^\checkmark\\checkmarkc\checkmarki\checkmarkr\checkmarkc\checkmark$\checkmarke\checkmark$\checkmark^\checkmark\\checkmarkc\checkmarki\checkmarkr\checkmarkc\checkmark$\checkmark \checkmark$\checkmark^\checkmark\\checkmarkc\checkmarki\checkmarkr\checkmarkc\checkmark$\checkmarkr\checkmark$\checkmark^\checkmark\\checkmarkc\checkmarki\checkmarkr\checkmarkc\checkmark$\checkmarké\checkmark$\checkmark^\checkmark\\checkmarkc\checkmarki\checkmarkr\checkmarkc\checkmark$\checkmarkf\checkmark$\checkmark^\checkmark\\checkmarkc\checkmarki\checkmarkr\checkmarkc\checkmark$\checkmarké\checkmark$\checkmark^\checkmark\\checkmarkc\checkmarki\checkmarkr\checkmarkc\checkmark$\checkmarkr\checkmark$\checkmark^\checkmark\\checkmarkc\checkmarki\checkmarkr\checkmarkc\checkmark$\checkmarke\checkmark$\checkmark^\checkmark\\checkmarkc\checkmarki\checkmarkr\checkmarkc\checkmark$\checkmarkn\checkmark$\checkmark^\checkmark\\checkmarkc\checkmarki\checkmarkr\checkmarkc\checkmark$\checkmarkt\checkmark$\checkmark^\checkmark\\checkmarkc\checkmarki\checkmarkr\checkmarkc\checkmark$\checkmarki\checkmark$\checkmark^\checkmark\\checkmarkc\checkmarki\checkmarkr\checkmarkc\checkmark$\checkmarke\checkmark$\checkmark^\checkmark\\checkmarkc\checkmarki\checkmarkr\checkmarkc\checkmark$\checkmarkl\checkmark$\checkmark^\checkmark\\checkmarkc\checkmarki\checkmarkr\checkmarkc\checkmark$\checkmark \checkmark$\checkmark^\checkmark\\checkmarkc\checkmarki\checkmarkr\checkmarkc\checkmark$\checkmarkd\checkmark$\checkmark^\checkmark\\checkmarkc\checkmarki\checkmarkr\checkmarkc\checkmark$\checkmarke\checkmark$\checkmark^\checkmark\\checkmarkc\checkmarki\checkmarkr\checkmarkc\checkmark$\checkmark \checkmark$\checkmark^\checkmark\\checkmarkc\checkmarki\checkmarkr\checkmarkc\checkmark$\checkmarkl\checkmark$\checkmark^\checkmark\\checkmarkc\checkmarki\checkmarkr\checkmarkc\checkmark$\checkmarka\checkmark$\checkmark^\checkmark\\checkmarkc\checkmarki\checkmarkr\checkmarkc\checkmark$\checkmark \checkmark$\checkmark^\checkmark\\checkmarkc\checkmarki\checkmarkr\checkmarkc\checkmark$\checkmarkm\checkmark$\checkmark^\checkmark\\checkmarkc\checkmarki\checkmarkr\checkmarkc\checkmark$\checkmarka\checkmark$\checkmark^\checkmark\\checkmarkc\checkmarki\checkmarkr\checkmarkc\checkmark$\checkmarkc\checkmark$\checkmark^\checkmark\\checkmarkc\checkmarki\checkmarkr\checkmarkc\checkmark$\checkmarkh\checkmark$\checkmark^\checkmark\\checkmarkc\checkmarki\checkmarkr\checkmarkc\checkmark$\checkmarki\checkmark$\checkmark^\checkmark\\checkmarkc\checkmarki\checkmarkr\checkmarkc\checkmark$\checkmarkn\checkmark$\checkmark^\checkmark\\checkmarkc\checkmarki\checkmarkr\checkmarkc\checkmark$\checkmarke\checkmark$\checkmark^\checkmark\\checkmarkc\checkmarki\checkmarkr\checkmarkc\checkmark$\checkmark \checkmark$\checkmark^\checkmark\\checkmarkc\checkmarki\checkmarkr\checkmarkc\checkmark$\checkmark$\checkmark$\checkmark^\checkmark\\checkmarkc\checkmarki\checkmarkr\checkmarkc\checkmark$\checkmark\\checkmark$\checkmark^\checkmark\\checkmarkc\checkmarki\checkmarkr\checkmarkc\checkmark$\checkmarkm\checkmark$\checkmark^\checkmark\\checkmarkc\checkmarki\checkmarkr\checkmarkc\checkmark$\checkmarka\checkmark$\checkmark^\checkmark\\checkmarkc\checkmarki\checkmarkr\checkmarkc\checkmark$\checkmarkt\checkmark$\checkmark^\checkmark\\checkmarkc\checkmarki\checkmarkr\checkmarkc\checkmark$\checkmarkh\checkmark$\checkmark^\checkmark\\checkmarkc\checkmarki\checkmarkr\checkmarkc\checkmark$\checkmarkc\checkmark$\checkmark^\checkmark\\checkmarkc\checkmarki\checkmarkr\checkmarkc\checkmark$\checkmarka\checkmark$\checkmark^\checkmark\\checkmarkc\checkmarki\checkmarkr\checkmarkc\checkmark$\checkmarkl\checkmark$\checkmark^\checkmark\\checkmarkc\checkmarki\checkmarkr\checkmarkc\checkmark$\checkmark{\checkmark$\checkmark^\checkmark\\checkmarkc\checkmarki\checkmarkr\checkmarkc\checkmark$\checkmarkR\checkmark$\checkmark^\checkmark\\checkmarkc\checkmarki\checkmarkr\checkmarkc\checkmark$\checkmark}\checkmark$\checkmark^\checkmark\\checkmarkc\checkmarki\checkmarkr\checkmarkc\checkmark$\checkmark(\checkmark$\checkmark^\checkmark\\checkmarkc\checkmarki\checkmarkr\checkmarkc\checkmark$\checkmarkO\checkmark$\checkmark^\checkmark\\checkmarkc\checkmarki\checkmarkr\checkmarkc\checkmark$\checkmark,\checkmark$\checkmark^\checkmark\\checkmarkc\checkmarki\checkmarkr\checkmarkc\checkmark$\checkmark \checkmark$\checkmark^\checkmark\\checkmarkc\checkmarki\checkmarkr\checkmarkc\checkmark$\checkmark\\checkmark$\checkmark^\checkmark\\checkmarkc\checkmarki\checkmarkr\checkmarkc\checkmark$\checkmarkv\checkmark$\checkmark^\checkmark\\checkmarkc\checkmarki\checkmarkr\checkmarkc\checkmark$\checkmarke\checkmark$\checkmark^\checkmark\\checkmarkc\checkmarki\checkmarkr\checkmarkc\checkmark$\checkmarkc\checkmark$\checkmark^\checkmark\\checkmarkc\checkmarki\checkmarkr\checkmarkc\checkmark$\checkmark{\checkmark$\checkmark^\checkmark\\checkmarkc\checkmarki\checkmarkr\checkmarkc\checkmark$\checkmarkx\checkmark$\checkmark^\checkmark\\checkmarkc\checkmarki\checkmarkr\checkmarkc\checkmark$\checkmark}\checkmark$\checkmark^\checkmark\\checkmarkc\checkmarki\checkmarkr\checkmarkc\checkmark$\checkmark,\checkmark$\checkmark^\checkmark\\checkmarkc\checkmarki\checkmarkr\checkmarkc\checkmark$\checkmark \checkmark$\checkmark^\checkmark\\checkmarkc\checkmarki\checkmarkr\checkmarkc\checkmark$\checkmark\\checkmark$\checkmark^\checkmark\\checkmarkc\checkmarki\checkmarkr\checkmarkc\checkmark$\checkmarkv\checkmark$\checkmark^\checkmark\\checkmarkc\checkmarki\checkmarkr\checkmarkc\checkmark$\checkmarke\checkmark$\checkmark^\checkmark\\checkmarkc\checkmarki\checkmarkr\checkmarkc\checkmark$\checkmarkc\checkmark$\checkmark^\checkmark\\checkmarkc\checkmarki\checkmarkr\checkmarkc\checkmark$\checkmark{\checkmark$\checkmark^\checkmark\\checkmarkc\checkmarki\checkmarkr\checkmarkc\checkmark$\checkmarky\checkmark$\checkmark^\checkmark\\checkmarkc\checkmarki\checkmarkr\checkmarkc\checkmark$\checkmark}\checkmark$\checkmark^\checkmark\\checkmarkc\checkmarki\checkmarkr\checkmarkc\checkmark$\checkmark,\checkmark$\checkmark^\checkmark\\checkmarkc\checkmarki\checkmarkr\checkmarkc\checkmark$\checkmark \checkmark$\checkmark^\checkmark\\checkmarkc\checkmarki\checkmarkr\checkmarkc\checkmark$\checkmark\\checkmark$\checkmark^\checkmark\\checkmarkc\checkmarki\checkmarkr\checkmarkc\checkmark$\checkmarkv\checkmark$\checkmark^\checkmark\\checkmarkc\checkmarki\checkmarkr\checkmarkc\checkmark$\checkmarke\checkmark$\checkmark^\checkmark\\checkmarkc\checkmarki\checkmarkr\checkmarkc\checkmark$\checkmarkc\checkmark$\checkmark^\checkmark\\checkmarkc\checkmarki\checkmarkr\checkmarkc\checkmark$\checkmark{\checkmark$\checkmark^\checkmark\\checkmarkc\checkmarki\checkmarkr\checkmarkc\checkmark$\checkmarkz\checkmark$\checkmark^\checkmark\\checkmarkc\checkmarki\checkmarkr\checkmarkc\checkmark$\checkmark}\checkmark$\checkmark^\checkmark\\checkmarkc\checkmarki\checkmarkr\checkmarkc\checkmark$\checkmark)\checkmark$\checkmark^\checkmark\\checkmarkc\checkmarki\checkmarkr\checkmarkc\checkmark$\checkmark$\checkmark$\checkmark^\checkmark\\checkmarkc\checkmarki\checkmarkr\checkmarkc\checkmark$\checkmark,\checkmark$\checkmark^\checkmark\\checkmarkc\checkmarki\checkmarkr\checkmarkc\checkmark$\checkmark \checkmark$\checkmark^\checkmark\\checkmarkc\checkmarki\checkmarkr\checkmarkc\checkmark$\checkmarke\checkmark$\checkmark^\checkmark\\checkmarkc\checkmarki\checkmarkr\checkmarkc\checkmark$\checkmarks\checkmark$\checkmark^\checkmark\\checkmarkc\checkmarki\checkmarkr\checkmarkc\checkmark$\checkmarkt\checkmark$\checkmark^\checkmark\\checkmarkc\checkmarki\checkmarkr\checkmarkc\checkmark$\checkmark \checkmark$\checkmark^\checkmark\\checkmarkc\checkmarki\checkmarkr\checkmarkc\checkmark$\checkmarkm\checkmark$\checkmark^\checkmark\\checkmarkc\checkmarki\checkmarkr\checkmarkc\checkmark$\checkmarko\checkmark$\checkmark^\checkmark\\checkmarkc\checkmarki\checkmarkr\checkmarkc\checkmark$\checkmarkd\checkmark$\checkmark^\checkmark\\checkmarkc\checkmarki\checkmarkr\checkmarkc\checkmark$\checkmarké\checkmark$\checkmark^\checkmark\\checkmarkc\checkmarki\checkmarkr\checkmarkc\checkmark$\checkmarkl\checkmark$\checkmark^\checkmark\\checkmarkc\checkmarki\checkmarkr\checkmarkc\checkmark$\checkmarki\checkmark$\checkmark^\checkmark\\checkmarkc\checkmarki\checkmarkr\checkmarkc\checkmark$\checkmarks\checkmark$\checkmark^\checkmark\\checkmarkc\checkmarki\checkmarkr\checkmarkc\checkmark$\checkmarké\checkmark$\checkmark^\checkmark\\checkmarkc\checkmarki\checkmarkr\checkmarkc\checkmark$\checkmarke\checkmark$\checkmark^\checkmark\\checkmarkc\checkmarki\checkmarkr\checkmarkc\checkmark$\checkmark \checkmark$\checkmark^\checkmark\\checkmarkc\checkmarki\checkmarkr\checkmarkc\checkmark$\checkmarkp\checkmark$\checkmark^\checkmark\\checkmarkc\checkmarki\checkmarkr\checkmarkc\checkmark$\checkmarka\checkmark$\checkmark^\checkmark\\checkmarkc\checkmarki\checkmarkr\checkmarkc\checkmark$\checkmarkr\checkmark$\checkmark^\checkmark\\checkmarkc\checkmarki\checkmarkr\checkmarkc\checkmark$\checkmark \checkmark$\checkmark^\checkmark\\checkmarkc\checkmarki\checkmarkr\checkmarkc\checkmark$\checkmarkl\checkmark$\checkmark^\checkmark\\checkmarkc\checkmarki\checkmarkr\checkmarkc\checkmark$\checkmarke\checkmark$\checkmark^\checkmark\\checkmarkc\checkmarki\checkmarkr\checkmarkc\checkmark$\checkmark \checkmark$\checkmark^\checkmark\\checkmarkc\checkmarki\checkmarkr\checkmarkc\checkmark$\checkmarkt\checkmark$\checkmark^\checkmark\\checkmarkc\checkmarki\checkmarkr\checkmarkc\checkmark$\checkmarko\checkmark$\checkmark^\checkmark\\checkmarkc\checkmarki\checkmarkr\checkmarkc\checkmark$\checkmarkr\checkmark$\checkmark^\checkmark\\checkmarkc\checkmarki\checkmarkr\checkmarkc\checkmark$\checkmarks\checkmark$\checkmark^\checkmark\\checkmarkc\checkmarki\checkmarkr\checkmarkc\checkmark$\checkmarke\checkmark$\checkmark^\checkmark\\checkmarkc\checkmarki\checkmarkr\checkmarkc\checkmark$\checkmarku\checkmark$\checkmark^\checkmark\\checkmarkc\checkmarki\checkmarkr\checkmarkc\checkmark$\checkmarkr\checkmark$\checkmark^\checkmark\\checkmarkc\checkmarki\checkmarkr\checkmarkc\checkmark$\checkmark \checkmark$\checkmark^\checkmark\\checkmarkc\checkmarki\checkmarkr\checkmarkc\checkmark$\checkmarkd\checkmark$\checkmark^\checkmark\\checkmarkc\checkmarki\checkmarkr\checkmarkc\checkmark$\checkmark'\checkmark$\checkmark^\checkmark\\checkmarkc\checkmarki\checkmarkr\checkmarkc\checkmark$\checkmarka\checkmark$\checkmark^\checkmark\\checkmarkc\checkmarki\checkmarkr\checkmarkc\checkmark$\checkmarkc\checkmark$\checkmark^\checkmark\\checkmarkc\checkmarki\checkmarkr\checkmarkc\checkmark$\checkmarkt\checkmark$\checkmark^\checkmark\\checkmarkc\checkmarki\checkmarkr\checkmarkc\checkmark$\checkmarki\checkmark$\checkmark^\checkmark\\checkmarkc\checkmarki\checkmarkr\checkmarkc\checkmark$\checkmarko\checkmark$\checkmark^\checkmark\\checkmarkc\checkmarki\checkmarkr\checkmarkc\checkmark$\checkmarkn\checkmark$\checkmark^\checkmark\\checkmarkc\checkmarki\checkmarkr\checkmarkc\checkmark$\checkmark \checkmark$\checkmark^\checkmark\\checkmarkc\checkmarki\checkmarkr\checkmarkc\checkmark$\checkmarkm\checkmark$\checkmark^\checkmark\\checkmarkc\checkmarki\checkmarkr\checkmarkc\checkmark$\checkmarké\checkmark$\checkmark^\checkmark\\checkmarkc\checkmarki\checkmarkr\checkmarkc\checkmark$\checkmarkc\checkmark$\checkmark^\checkmark\\checkmarkc\checkmarki\checkmarkr\checkmarkc\checkmark$\checkmarka\checkmark$\checkmark^\checkmark\\checkmarkc\checkmarki\checkmarkr\checkmarkc\checkmark$\checkmarkn\checkmark$\checkmark^\checkmark\\checkmarkc\checkmarki\checkmarkr\checkmarkc\checkmark$\checkmarki\checkmark$\checkmark^\checkmark\\checkmarkc\checkmarki\checkmarkr\checkmarkc\checkmark$\checkmarkq\checkmark$\checkmark^\checkmark\\checkmarkc\checkmarki\checkmarkr\checkmarkc\checkmark$\checkmarku\checkmark$\checkmark^\checkmark\\checkmarkc\checkmarki\checkmarkr\checkmarkc\checkmark$\checkmarke\checkmark$\checkmark^\checkmark\\checkmarkc\checkmarki\checkmarkr\checkmarkc\checkmark$\checkmark \checkmark$\checkmark^\checkmark\\checkmarkc\checkmarki\checkmarkr\checkmarkc\checkmark$\checkmark:\checkmark$\checkmark^\checkmark\\checkmarkc\checkmarki\checkmarkr\checkmarkc\checkmark$\checkmark
\checkmark$\checkmark^\checkmark\\checkmarkc\checkmarki\checkmarkr\checkmarkc\checkmark$\checkmark\\checkmark$\checkmark^\checkmark\\checkmarkc\checkmarki\checkmarkr\checkmarkc\checkmark$\checkmarkb\checkmark$\checkmark^\checkmark\\checkmarkc\checkmarki\checkmarkr\checkmarkc\checkmark$\checkmarke\checkmark$\checkmark^\checkmark\\checkmarkc\checkmarki\checkmarkr\checkmarkc\checkmark$\checkmarkg\checkmark$\checkmark^\checkmark\\checkmarkc\checkmarki\checkmarkr\checkmarkc\checkmark$\checkmarki\checkmark$\checkmark^\checkmark\\checkmarkc\checkmarki\checkmarkr\checkmarkc\checkmark$\checkmarkn\checkmark$\checkmark^\checkmark\\checkmarkc\checkmarki\checkmarkr\checkmarkc\checkmark$\checkmark{\checkmark$\checkmark^\checkmark\\checkmarkc\checkmarki\checkmarkr\checkmarkc\checkmark$\checkmarke\checkmark$\checkmark^\checkmark\\checkmarkc\checkmarki\checkmarkr\checkmarkc\checkmark$\checkmarkq\checkmark$\checkmark^\checkmark\\checkmarkc\checkmarki\checkmarkr\checkmarkc\checkmark$\checkmarku\checkmark$\checkmark^\checkmark\\checkmarkc\checkmarki\checkmarkr\checkmarkc\checkmark$\checkmarka\checkmark$\checkmark^\checkmark\\checkmarkc\checkmarki\checkmarkr\checkmarkc\checkmark$\checkmarkt\checkmark$\checkmark^\checkmark\\checkmarkc\checkmarki\checkmarkr\checkmarkc\checkmark$\checkmarki\checkmark$\checkmark^\checkmark\\checkmarkc\checkmarki\checkmarkr\checkmarkc\checkmark$\checkmarko\checkmark$\checkmark^\checkmark\\checkmarkc\checkmarki\checkmarkr\checkmarkc\checkmark$\checkmarkn\checkmark$\checkmark^\checkmark\\checkmarkc\checkmarki\checkmarkr\checkmarkc\checkmark$\checkmark}\checkmark$\checkmark^\checkmark\\checkmarkc\checkmarki\checkmarkr\checkmarkc\checkmark$\checkmark
\checkmark$\checkmark^\checkmark\\checkmarkc\checkmarki\checkmarkr\checkmarkc\checkmark$\checkmark \checkmark$\checkmark^\checkmark\\checkmarkc\checkmarki\checkmarkr\checkmarkc\checkmark$\checkmark \checkmark$\checkmark^\checkmark\\checkmarkc\checkmarki\checkmarkr\checkmarkc\checkmark$\checkmark \checkmark$\checkmark^\checkmark\\checkmarkc\checkmarki\checkmarkr\checkmarkc\checkmark$\checkmark \checkmark$\checkmark^\checkmark\\checkmarkc\checkmarki\checkmarkr\checkmarkc\checkmark$\checkmark\\checkmark$\checkmark^\checkmark\\checkmarkc\checkmarki\checkmarkr\checkmarkc\checkmark$\checkmark{\checkmark$\checkmark^\checkmark\\checkmarkc\checkmarki\checkmarkr\checkmarkc\checkmark$\checkmarkT\checkmark$\checkmark^\checkmark\\checkmarkc\checkmarki\checkmarkr\checkmarkc\checkmark$\checkmark\\checkmark$\checkmark^\checkmark\\checkmarkc\checkmarki\checkmarkr\checkmarkc\checkmark$\checkmark}\checkmark$\checkmark^\checkmark\\checkmarkc\checkmarki\checkmarkr\checkmarkc\checkmark$\checkmark_\checkmark$\checkmark^\checkmark\\checkmarkc\checkmarki\checkmarkr\checkmarkc\checkmark$\checkmark{\checkmark$\checkmark^\checkmark\\checkmarkc\checkmarki\checkmarkr\checkmarkc\checkmark$\checkmark\\checkmark$\checkmark^\checkmark\\checkmarkc\checkmarki\checkmarkr\checkmarkc\checkmark$\checkmarkt\checkmark$\checkmark^\checkmark\\checkmarkc\checkmarki\checkmarkr\checkmarkc\checkmark$\checkmarke\checkmark$\checkmark^\checkmark\\checkmarkc\checkmarki\checkmarkr\checkmarkc\checkmark$\checkmarkx\checkmark$\checkmark^\checkmark\\checkmarkc\checkmarki\checkmarkr\checkmarkc\checkmark$\checkmarkt\checkmark$\checkmark^\checkmark\\checkmarkc\checkmarki\checkmarkr\checkmarkc\checkmark$\checkmark{\checkmark$\checkmark^\checkmark\\checkmarkc\checkmarki\checkmarkr\checkmarkc\checkmark$\checkmarko\checkmark$\checkmark^\checkmark\\checkmarkc\checkmarki\checkmarkr\checkmarkc\checkmark$\checkmarku\checkmark$\checkmark^\checkmark\\checkmarkc\checkmarki\checkmarkr\checkmarkc\checkmark$\checkmarkt\checkmark$\checkmark^\checkmark\\checkmarkc\checkmarki\checkmarkr\checkmarkc\checkmark$\checkmarki\checkmark$\checkmark^\checkmark\\checkmarkc\checkmarki\checkmarkr\checkmarkc\checkmark$\checkmarkl\checkmark$\checkmark^\checkmark\\checkmarkc\checkmarki\checkmarkr\checkmarkc\checkmark$\checkmark}\checkmark$\checkmark^\checkmark\\checkmarkc\checkmarki\checkmarkr\checkmarkc\checkmark$\checkmark \checkmark$\checkmark^\checkmark\\checkmarkc\checkmarki\checkmarkr\checkmarkc\checkmark$\checkmark\\checkmark$\checkmark^\checkmark\\checkmarkc\checkmarki\checkmarkr\checkmarkc\checkmark$\checkmarkt\checkmark$\checkmark^\checkmark\\checkmarkc\checkmarki\checkmarkr\checkmarkc\checkmark$\checkmarko\checkmark$\checkmark^\checkmark\\checkmarkc\checkmarki\checkmarkr\checkmarkc\checkmark$\checkmark \checkmark$\checkmark^\checkmark\\checkmarkc\checkmarki\checkmarkr\checkmarkc\checkmark$\checkmark\\checkmark$\checkmark^\checkmark\\checkmarkc\checkmarki\checkmarkr\checkmarkc\checkmark$\checkmarkt\checkmark$\checkmark^\checkmark\\checkmarkc\checkmarki\checkmarkr\checkmarkc\checkmark$\checkmarke\checkmark$\checkmark^\checkmark\\checkmarkc\checkmarki\checkmarkr\checkmarkc\checkmark$\checkmarkx\checkmark$\checkmark^\checkmark\\checkmarkc\checkmarki\checkmarkr\checkmarkc\checkmark$\checkmarkt\checkmark$\checkmark^\checkmark\\checkmarkc\checkmarki\checkmarkr\checkmarkc\checkmark$\checkmark{\checkmark$\checkmark^\checkmark\\checkmarkc\checkmarki\checkmarkr\checkmarkc\checkmark$\checkmarkp\checkmark$\checkmark^\checkmark\\checkmarkc\checkmarki\checkmarkr\checkmarkc\checkmark$\checkmarki\checkmark$\checkmark^\checkmark\\checkmarkc\checkmarki\checkmarkr\checkmarkc\checkmark$\checkmarkè\checkmark$\checkmark^\checkmark\\checkmarkc\checkmarki\checkmarkr\checkmarkc\checkmark$\checkmarkc\checkmark$\checkmark^\checkmark\\checkmarkc\checkmarki\checkmarkr\checkmarkc\checkmark$\checkmarke\checkmark$\checkmark^\checkmark\\checkmarkc\checkmarki\checkmarkr\checkmarkc\checkmark$\checkmark}\checkmark$\checkmark^\checkmark\\checkmarkc\checkmarki\checkmarkr\checkmarkc\checkmark$\checkmark}\checkmark$\checkmark^\checkmark\\checkmarkc\checkmarki\checkmarkr\checkmarkc\checkmark$\checkmark \checkmark$\checkmark^\checkmark\\checkmarkc\checkmarki\checkmarkr\checkmarkc\checkmark$\checkmark=\checkmark$\checkmark^\checkmark\\checkmarkc\checkmarki\checkmarkr\checkmarkc\checkmark$\checkmark
\checkmark$\checkmark^\checkmark\\checkmarkc\checkmarki\checkmarkr\checkmarkc\checkmark$\checkmark \checkmark$\checkmark^\checkmark\\checkmarkc\checkmarki\checkmarkr\checkmarkc\checkmark$\checkmark \checkmark$\checkmark^\checkmark\\checkmarkc\checkmarki\checkmarkr\checkmarkc\checkmark$\checkmark \checkmark$\checkmark^\checkmark\\checkmarkc\checkmarki\checkmarkr\checkmarkc\checkmark$\checkmark \checkmark$\checkmark^\checkmark\\checkmarkc\checkmarki\checkmarkr\checkmarkc\checkmark$\checkmark\\checkmark$\checkmark^\checkmark\\checkmarkc\checkmarki\checkmarkr\checkmarkc\checkmark$\checkmarkb\checkmark$\checkmark^\checkmark\\checkmarkc\checkmarki\checkmarkr\checkmarkc\checkmark$\checkmarke\checkmark$\checkmark^\checkmark\\checkmarkc\checkmarki\checkmarkr\checkmarkc\checkmark$\checkmarkg\checkmark$\checkmark^\checkmark\\checkmarkc\checkmarki\checkmarkr\checkmarkc\checkmark$\checkmarki\checkmark$\checkmark^\checkmark\\checkmarkc\checkmarki\checkmarkr\checkmarkc\checkmark$\checkmarkn\checkmark$\checkmark^\checkmark\\checkmarkc\checkmarki\checkmarkr\checkmarkc\checkmark$\checkmark{\checkmark$\checkmark^\checkmark\\checkmarkc\checkmarki\checkmarkr\checkmarkc\checkmark$\checkmarkB\checkmark$\checkmark^\checkmark\\checkmarkc\checkmarki\checkmarkr\checkmarkc\checkmark$\checkmarkm\checkmark$\checkmark^\checkmark\\checkmarkc\checkmarki\checkmarkr\checkmarkc\checkmark$\checkmarka\checkmark$\checkmark^\checkmark\\checkmarkc\checkmarki\checkmarkr\checkmarkc\checkmark$\checkmarkt\checkmark$\checkmark^\checkmark\\checkmarkc\checkmarki\checkmarkr\checkmarkc\checkmark$\checkmarkr\checkmark$\checkmark^\checkmark\\checkmarkc\checkmarki\checkmarkr\checkmarkc\checkmark$\checkmarki\checkmark$\checkmark^\checkmark\\checkmarkc\checkmarki\checkmarkr\checkmarkc\checkmark$\checkmarkx\checkmark$\checkmark^\checkmark\\checkmarkc\checkmarki\checkmarkr\checkmarkc\checkmark$\checkmark}\checkmark$\checkmark^\checkmark\\checkmarkc\checkmarki\checkmarkr\checkmarkc\checkmark$\checkmark
\checkmark$\checkmark^\checkmark\\checkmarkc\checkmarki\checkmarkr\checkmarkc\checkmark$\checkmark \checkmark$\checkmark^\checkmark\\checkmarkc\checkmarki\checkmarkr\checkmarkc\checkmark$\checkmark \checkmark$\checkmark^\checkmark\\checkmarkc\checkmarki\checkmarkr\checkmarkc\checkmark$\checkmark \checkmark$\checkmark^\checkmark\\checkmarkc\checkmarki\checkmarkr\checkmarkc\checkmark$\checkmark \checkmark$\checkmark^\checkmark\\checkmarkc\checkmarki\checkmarkr\checkmarkc\checkmark$\checkmark \checkmark$\checkmark^\checkmark\\checkmarkc\checkmarki\checkmarkr\checkmarkc\checkmark$\checkmark \checkmark$\checkmark^\checkmark\\checkmarkc\checkmarki\checkmarkr\checkmarkc\checkmark$\checkmark \checkmark$\checkmark^\checkmark\\checkmarkc\checkmarki\checkmarkr\checkmarkc\checkmark$\checkmark \checkmark$\checkmark^\checkmark\\checkmarkc\checkmarki\checkmarkr\checkmarkc\checkmark$\checkmark\\checkmark$\checkmark^\checkmark\\checkmarkc\checkmarki\checkmarkr\checkmarkc\checkmark$\checkmarkv\checkmark$\checkmark^\checkmark\\checkmarkc\checkmarki\checkmarkr\checkmarkc\checkmark$\checkmarke\checkmark$\checkmark^\checkmark\\checkmarkc\checkmarki\checkmarkr\checkmarkc\checkmark$\checkmarkc\checkmark$\checkmark^\checkmark\\checkmarkc\checkmarki\checkmarkr\checkmarkc\checkmark$\checkmark{\checkmark$\checkmark^\checkmark\\checkmarkc\checkmarki\checkmarkr\checkmarkc\checkmark$\checkmarkR\checkmark$\checkmark^\checkmark\\checkmarkc\checkmarki\checkmarkr\checkmarkc\checkmark$\checkmark}\checkmark$\checkmark^\checkmark\\checkmarkc\checkmarki\checkmarkr\checkmarkc\checkmark$\checkmark \checkmark$\checkmark^\checkmark\\checkmarkc\checkmarki\checkmarkr\checkmarkc\checkmark$\checkmark=\checkmark$\checkmark^\checkmark\\checkmarkc\checkmarki\checkmarkr\checkmarkc\checkmark$\checkmark \checkmark$\checkmark^\checkmark\\checkmarkc\checkmarki\checkmarkr\checkmarkc\checkmark$\checkmarkF\checkmark$\checkmark^\checkmark\\checkmarkc\checkmarki\checkmarkr\checkmarkc\checkmark$\checkmark_\checkmark$\checkmark^\checkmark\\checkmarkc\checkmarki\checkmarkr\checkmarkc\checkmark$\checkmarkc\checkmark$\checkmark^\checkmark\\checkmarkc\checkmarki\checkmarkr\checkmarkc\checkmark$\checkmark\\checkmark$\checkmark^\checkmark\\checkmarkc\checkmarki\checkmarkr\checkmarkc\checkmark$\checkmark,\checkmark$\checkmark^\checkmark\\checkmarkc\checkmarki\checkmarkr\checkmarkc\checkmark$\checkmark\\checkmark$\checkmark^\checkmark\\checkmarkc\checkmarki\checkmarkr\checkmarkc\checkmark$\checkmarkv\checkmark$\checkmark^\checkmark\\checkmarkc\checkmarki\checkmarkr\checkmarkc\checkmark$\checkmarke\checkmark$\checkmark^\checkmark\\checkmarkc\checkmarki\checkmarkr\checkmarkc\checkmark$\checkmarkc\checkmark$\checkmark^\checkmark\\checkmarkc\checkmarki\checkmarkr\checkmarkc\checkmark$\checkmark{\checkmark$\checkmark^\checkmark\\checkmarkc\checkmarki\checkmarkr\checkmarkc\checkmark$\checkmarkx\checkmark$\checkmark^\checkmark\\checkmarkc\checkmarki\checkmarkr\checkmarkc\checkmark$\checkmark}\checkmark$\checkmark^\checkmark\\checkmarkc\checkmarki\checkmarkr\checkmarkc\checkmark$\checkmark \checkmark$\checkmark^\checkmark\\checkmarkc\checkmarki\checkmarkr\checkmarkc\checkmark$\checkmark+\checkmark$\checkmark^\checkmark\\checkmarkc\checkmarki\checkmarkr\checkmarkc\checkmark$\checkmark \checkmark$\checkmark^\checkmark\\checkmarkc\checkmarki\checkmarkr\checkmarkc\checkmark$\checkmarkF\checkmark$\checkmark^\checkmark\\checkmarkc\checkmarki\checkmarkr\checkmarkc\checkmark$\checkmark_\checkmark$\checkmark^\checkmark\\checkmarkc\checkmarki\checkmarkr\checkmarkc\checkmark$\checkmarkf\checkmark$\checkmark^\checkmark\\checkmarkc\checkmarki\checkmarkr\checkmarkc\checkmark$\checkmark\\checkmark$\checkmark^\checkmark\\checkmarkc\checkmarki\checkmarkr\checkmarkc\checkmark$\checkmark,\checkmark$\checkmark^\checkmark\\checkmarkc\checkmarki\checkmarkr\checkmarkc\checkmark$\checkmark\\checkmark$\checkmark^\checkmark\\checkmarkc\checkmarki\checkmarkr\checkmarkc\checkmark$\checkmarkv\checkmark$\checkmark^\checkmark\\checkmarkc\checkmarki\checkmarkr\checkmarkc\checkmark$\checkmarke\checkmark$\checkmark^\checkmark\\checkmarkc\checkmarki\checkmarkr\checkmarkc\checkmark$\checkmarkc\checkmark$\checkmark^\checkmark\\checkmarkc\checkmarki\checkmarkr\checkmarkc\checkmark$\checkmark{\checkmark$\checkmark^\checkmark\\checkmarkc\checkmarki\checkmarkr\checkmarkc\checkmark$\checkmarky\checkmark$\checkmark^\checkmark\\checkmarkc\checkmarki\checkmarkr\checkmarkc\checkmark$\checkmark}\checkmark$\checkmark^\checkmark\\checkmarkc\checkmarki\checkmarkr\checkmarkc\checkmark$\checkmark \checkmark$\checkmark^\checkmark\\checkmarkc\checkmarki\checkmarkr\checkmarkc\checkmark$\checkmark+\checkmark$\checkmark^\checkmark\\checkmarkc\checkmarki\checkmarkr\checkmarkc\checkmark$\checkmark \checkmark$\checkmark^\checkmark\\checkmarkc\checkmarki\checkmarkr\checkmarkc\checkmark$\checkmarkF\checkmark$\checkmark^\checkmark\\checkmarkc\checkmarki\checkmarkr\checkmarkc\checkmark$\checkmark_\checkmark$\checkmark^\checkmark\\checkmarkc\checkmarki\checkmarkr\checkmarkc\checkmark$\checkmarkp\checkmark$\checkmark^\checkmark\\checkmarkc\checkmarki\checkmarkr\checkmarkc\checkmark$\checkmark\\checkmark$\checkmark^\checkmark\\checkmarkc\checkmarki\checkmarkr\checkmarkc\checkmark$\checkmark,\checkmark$\checkmark^\checkmark\\checkmarkc\checkmarki\checkmarkr\checkmarkc\checkmark$\checkmark\\checkmark$\checkmark^\checkmark\\checkmarkc\checkmarki\checkmarkr\checkmarkc\checkmark$\checkmarkv\checkmark$\checkmark^\checkmark\\checkmarkc\checkmarki\checkmarkr\checkmarkc\checkmark$\checkmarke\checkmark$\checkmark^\checkmark\\checkmarkc\checkmarki\checkmarkr\checkmarkc\checkmark$\checkmarkc\checkmark$\checkmark^\checkmark\\checkmarkc\checkmarki\checkmarkr\checkmarkc\checkmark$\checkmark{\checkmark$\checkmark^\checkmark\\checkmarkc\checkmarki\checkmarkr\checkmarkc\checkmark$\checkmarkz\checkmark$\checkmark^\checkmark\\checkmarkc\checkmarki\checkmarkr\checkmarkc\checkmark$\checkmark}\checkmark$\checkmark^\checkmark\\checkmarkc\checkmarki\checkmarkr\checkmarkc\checkmark$\checkmark \checkmark$\checkmark^\checkmark\\checkmarkc\checkmarki\checkmarkr\checkmarkc\checkmark$\checkmark\\checkmark$\checkmark^\checkmark\\checkmarkc\checkmarki\checkmarkr\checkmarkc\checkmark$\checkmark\\checkmark$\checkmark^\checkmark\\checkmarkc\checkmarki\checkmarkr\checkmarkc\checkmark$\checkmark
\checkmark$\checkmark^\checkmark\\checkmarkc\checkmarki\checkmarkr\checkmarkc\checkmark$\checkmark \checkmark$\checkmark^\checkmark\\checkmarkc\checkmarki\checkmarkr\checkmarkc\checkmark$\checkmark \checkmark$\checkmark^\checkmark\\checkmarkc\checkmarki\checkmarkr\checkmarkc\checkmark$\checkmark \checkmark$\checkmark^\checkmark\\checkmarkc\checkmarki\checkmarkr\checkmarkc\checkmark$\checkmark \checkmark$\checkmark^\checkmark\\checkmarkc\checkmarki\checkmarkr\checkmarkc\checkmark$\checkmark \checkmark$\checkmark^\checkmark\\checkmarkc\checkmarki\checkmarkr\checkmarkc\checkmark$\checkmark \checkmark$\checkmark^\checkmark\\checkmarkc\checkmarki\checkmarkr\checkmarkc\checkmark$\checkmark \checkmark$\checkmark^\checkmark\\checkmarkc\checkmarki\checkmarkr\checkmarkc\checkmark$\checkmark \checkmark$\checkmark^\checkmark\\checkmarkc\checkmarki\checkmarkr\checkmarkc\checkmark$\checkmark\\checkmark$\checkmark^\checkmark\\checkmarkc\checkmarki\checkmarkr\checkmarkc\checkmark$\checkmarkv\checkmark$\checkmark^\checkmark\\checkmarkc\checkmarki\checkmarkr\checkmarkc\checkmark$\checkmarke\checkmark$\checkmark^\checkmark\\checkmarkc\checkmarki\checkmarkr\checkmarkc\checkmark$\checkmarkc\checkmark$\checkmark^\checkmark\\checkmarkc\checkmarki\checkmarkr\checkmarkc\checkmark$\checkmark{\checkmark$\checkmark^\checkmark\\checkmarkc\checkmarki\checkmarkr\checkmarkc\checkmark$\checkmarkM\checkmark$\checkmark^\checkmark\\checkmarkc\checkmarki\checkmarkr\checkmarkc\checkmark$\checkmark}\checkmark$\checkmark^\checkmark\\checkmarkc\checkmarki\checkmarkr\checkmarkc\checkmark$\checkmark_\checkmark$\checkmark^\checkmark\\checkmarkc\checkmarki\checkmarkr\checkmarkc\checkmark$\checkmarkO\checkmark$\checkmark^\checkmark\\checkmarkc\checkmarki\checkmarkr\checkmarkc\checkmark$\checkmark \checkmark$\checkmark^\checkmark\\checkmarkc\checkmarki\checkmarkr\checkmarkc\checkmark$\checkmark\\checkmark$\checkmark^\checkmark\\checkmarkc\checkmarki\checkmarkr\checkmarkc\checkmark$\checkmarka\checkmark$\checkmark^\checkmark\\checkmarkc\checkmarki\checkmarkr\checkmarkc\checkmark$\checkmarkp\checkmark$\checkmark^\checkmark\\checkmarkc\checkmarki\checkmarkr\checkmarkc\checkmark$\checkmarkp\checkmark$\checkmark^\checkmark\\checkmarkc\checkmarki\checkmarkr\checkmarkc\checkmark$\checkmarkr\checkmark$\checkmark^\checkmark\\checkmarkc\checkmarki\checkmarkr\checkmarkc\checkmark$\checkmarko\checkmark$\checkmark^\checkmark\\checkmarkc\checkmarki\checkmarkr\checkmarkc\checkmark$\checkmarkx\checkmark$\checkmark^\checkmark\\checkmarkc\checkmarki\checkmarkr\checkmarkc\checkmark$\checkmark \checkmark$\checkmark^\checkmark\\checkmarkc\checkmarki\checkmarkr\checkmarkc\checkmark$\checkmark\\checkmark$\checkmark^\checkmark\\checkmarkc\checkmarki\checkmarkr\checkmarkc\checkmark$\checkmarkv\checkmark$\checkmark^\checkmark\\checkmarkc\checkmarki\checkmarkr\checkmarkc\checkmark$\checkmarke\checkmark$\checkmark^\checkmark\\checkmarkc\checkmarki\checkmarkr\checkmarkc\checkmark$\checkmarkc\checkmark$\checkmark^\checkmark\\checkmarkc\checkmarki\checkmarkr\checkmarkc\checkmark$\checkmark{\checkmark$\checkmark^\checkmark\\checkmarkc\checkmarki\checkmarkr\checkmarkc\checkmark$\checkmark0\checkmark$\checkmark^\checkmark\\checkmarkc\checkmarki\checkmarkr\checkmarkc\checkmark$\checkmark}\checkmark$\checkmark^\checkmark\\checkmarkc\checkmarki\checkmarkr\checkmarkc\checkmark$\checkmark \checkmark$\checkmark^\checkmark\\checkmarkc\checkmarki\checkmarkr\checkmarkc\checkmark$\checkmark\\checkmark$\checkmark^\checkmark\\checkmarkc\checkmarki\checkmarkr\checkmarkc\checkmark$\checkmarkq\checkmark$\checkmark^\checkmark\\checkmarkc\checkmarki\checkmarkr\checkmarkc\checkmark$\checkmarku\checkmark$\checkmark^\checkmark\\checkmarkc\checkmarki\checkmarkr\checkmarkc\checkmark$\checkmarka\checkmark$\checkmark^\checkmark\\checkmarkc\checkmarki\checkmarkr\checkmarkc\checkmark$\checkmarkd\checkmark$\checkmark^\checkmark\\checkmarkc\checkmarki\checkmarkr\checkmarkc\checkmark$\checkmark \checkmark$\checkmark^\checkmark\\checkmarkc\checkmarki\checkmarkr\checkmarkc\checkmark$\checkmark\\checkmark$\checkmark^\checkmark\\checkmarkc\checkmarki\checkmarkr\checkmarkc\checkmark$\checkmarkt\checkmark$\checkmark^\checkmark\\checkmarkc\checkmarki\checkmarkr\checkmarkc\checkmark$\checkmarke\checkmark$\checkmark^\checkmark\\checkmarkc\checkmarki\checkmarkr\checkmarkc\checkmark$\checkmarkx\checkmark$\checkmark^\checkmark\\checkmarkc\checkmarki\checkmarkr\checkmarkc\checkmark$\checkmarkt\checkmark$\checkmark^\checkmark\\checkmarkc\checkmarki\checkmarkr\checkmarkc\checkmark$\checkmark{\checkmark$\checkmark^\checkmark\\checkmarkc\checkmarki\checkmarkr\checkmarkc\checkmark$\checkmark(\checkmark$\checkmark^\checkmark\\checkmarkc\checkmarki\checkmarkr\checkmarkc\checkmark$\checkmarkm\checkmark$\checkmark^\checkmark\\checkmarkc\checkmarki\checkmarkr\checkmarkc\checkmark$\checkmarko\checkmark$\checkmark^\checkmark\\checkmarkc\checkmarki\checkmarkr\checkmarkc\checkmark$\checkmarkm\checkmark$\checkmark^\checkmark\\checkmarkc\checkmarki\checkmarkr\checkmarkc\checkmark$\checkmarke\checkmark$\checkmark^\checkmark\\checkmarkc\checkmarki\checkmarkr\checkmarkc\checkmark$\checkmarkn\checkmark$\checkmark^\checkmark\\checkmarkc\checkmarki\checkmarkr\checkmarkc\checkmark$\checkmarkt\checkmark$\checkmark^\checkmark\\checkmarkc\checkmarki\checkmarkr\checkmarkc\checkmark$\checkmark \checkmark$\checkmark^\checkmark\\checkmarkc\checkmarki\checkmarkr\checkmarkc\checkmark$\checkmarkd\checkmark$\checkmark^\checkmark\\checkmarkc\checkmarki\checkmarkr\checkmarkc\checkmark$\checkmark'\checkmark$\checkmark^\checkmark\\checkmarkc\checkmarki\checkmarkr\checkmarkc\checkmark$\checkmarke\checkmark$\checkmark^\checkmark\\checkmarkc\checkmarki\checkmarkr\checkmarkc\checkmark$\checkmarkn\checkmark$\checkmark^\checkmark\\checkmarkc\checkmarki\checkmarkr\checkmarkc\checkmark$\checkmarkc\checkmark$\checkmark^\checkmark\\checkmarkc\checkmarki\checkmarkr\checkmarkc\checkmark$\checkmarka\checkmark$\checkmark^\checkmark\\checkmarkc\checkmarki\checkmarkr\checkmarkc\checkmark$\checkmarks\checkmark$\checkmark^\checkmark\\checkmarkc\checkmarki\checkmarkr\checkmarkc\checkmark$\checkmarkt\checkmark$\checkmark^\checkmark\\checkmarkc\checkmarki\checkmarkr\checkmarkc\checkmark$\checkmarkr\checkmark$\checkmark^\checkmark\\checkmarkc\checkmarki\checkmarkr\checkmarkc\checkmark$\checkmarke\checkmark$\checkmark^\checkmark\\checkmarkc\checkmarki\checkmarkr\checkmarkc\checkmark$\checkmarkm\checkmark$\checkmark^\checkmark\\checkmarkc\checkmarki\checkmarkr\checkmarkc\checkmark$\checkmarke\checkmark$\checkmark^\checkmark\\checkmarkc\checkmarki\checkmarkr\checkmarkc\checkmark$\checkmarkn\checkmark$\checkmark^\checkmark\\checkmarkc\checkmarki\checkmarkr\checkmarkc\checkmark$\checkmarkt\checkmark$\checkmark^\checkmark\\checkmarkc\checkmarki\checkmarkr\checkmarkc\checkmark$\checkmark \checkmark$\checkmark^\checkmark\\checkmarkc\checkmarki\checkmarkr\checkmarkc\checkmark$\checkmarkn\checkmark$\checkmark^\checkmark\\checkmarkc\checkmarki\checkmarkr\checkmarkc\checkmark$\checkmarké\checkmark$\checkmark^\checkmark\\checkmarkc\checkmarki\checkmarkr\checkmarkc\checkmark$\checkmarkg\checkmark$\checkmark^\checkmark\\checkmarkc\checkmarki\checkmarkr\checkmarkc\checkmark$\checkmarkl\checkmark$\checkmark^\checkmark\\checkmarkc\checkmarki\checkmarkr\checkmarkc\checkmark$\checkmarki\checkmark$\checkmark^\checkmark\\checkmarkc\checkmarki\checkmarkr\checkmarkc\checkmark$\checkmarkg\checkmark$\checkmark^\checkmark\\checkmarkc\checkmarki\checkmarkr\checkmarkc\checkmark$\checkmarké\checkmark$\checkmark^\checkmark\\checkmarkc\checkmarki\checkmarkr\checkmarkc\checkmark$\checkmark)\checkmark$\checkmark^\checkmark\\checkmarkc\checkmarki\checkmarkr\checkmarkc\checkmark$\checkmark}\checkmark$\checkmark^\checkmark\\checkmarkc\checkmarki\checkmarkr\checkmarkc\checkmark$\checkmark
\checkmark$\checkmark^\checkmark\\checkmarkc\checkmarki\checkmarkr\checkmarkc\checkmark$\checkmark \checkmark$\checkmark^\checkmark\\checkmarkc\checkmarki\checkmarkr\checkmarkc\checkmark$\checkmark \checkmark$\checkmark^\checkmark\\checkmarkc\checkmarki\checkmarkr\checkmarkc\checkmark$\checkmark \checkmark$\checkmark^\checkmark\\checkmarkc\checkmarki\checkmarkr\checkmarkc\checkmark$\checkmark \checkmark$\checkmark^\checkmark\\checkmarkc\checkmarki\checkmarkr\checkmarkc\checkmark$\checkmark\\checkmark$\checkmark^\checkmark\\checkmarkc\checkmarki\checkmarkr\checkmarkc\checkmark$\checkmarke\checkmark$\checkmark^\checkmark\\checkmarkc\checkmarki\checkmarkr\checkmarkc\checkmark$\checkmarkn\checkmark$\checkmark^\checkmark\\checkmarkc\checkmarki\checkmarkr\checkmarkc\checkmark$\checkmarkd\checkmark$\checkmark^\checkmark\\checkmarkc\checkmarki\checkmarkr\checkmarkc\checkmark$\checkmark{\checkmark$\checkmark^\checkmark\\checkmarkc\checkmarki\checkmarkr\checkmarkc\checkmark$\checkmarkB\checkmark$\checkmark^\checkmark\\checkmarkc\checkmarki\checkmarkr\checkmarkc\checkmark$\checkmarkm\checkmark$\checkmark^\checkmark\\checkmarkc\checkmarki\checkmarkr\checkmarkc\checkmark$\checkmarka\checkmark$\checkmark^\checkmark\\checkmarkc\checkmarki\checkmarkr\checkmarkc\checkmark$\checkmarkt\checkmark$\checkmark^\checkmark\\checkmarkc\checkmarki\checkmarkr\checkmarkc\checkmark$\checkmarkr\checkmark$\checkmark^\checkmark\\checkmarkc\checkmarki\checkmarkr\checkmarkc\checkmark$\checkmarki\checkmark$\checkmark^\checkmark\\checkmarkc\checkmarki\checkmarkr\checkmarkc\checkmark$\checkmarkx\checkmark$\checkmark^\checkmark\\checkmarkc\checkmarki\checkmarkr\checkmarkc\checkmark$\checkmark}\checkmark$\checkmark^\checkmark\\checkmarkc\checkmarki\checkmarkr\checkmarkc\checkmark$\checkmark_\checkmark$\checkmark^\checkmark\\checkmarkc\checkmarki\checkmarkr\checkmarkc\checkmark$\checkmarkO\checkmark$\checkmark^\checkmark\\checkmarkc\checkmarki\checkmarkr\checkmarkc\checkmark$\checkmark
\checkmark$\checkmark^\checkmark\\checkmarkc\checkmarki\checkmarkr\checkmarkc\checkmark$\checkmark \checkmark$\checkmark^\checkmark\\checkmarkc\checkmarki\checkmarkr\checkmarkc\checkmark$\checkmark \checkmark$\checkmark^\checkmark\\checkmarkc\checkmarki\checkmarkr\checkmarkc\checkmark$\checkmark \checkmark$\checkmark^\checkmark\\checkmarkc\checkmarki\checkmarkr\checkmarkc\checkmark$\checkmark \checkmark$\checkmark^\checkmark\\checkmarkc\checkmarki\checkmarkr\checkmarkc\checkmark$\checkmark\\checkmark$\checkmark^\checkmark\\checkmarkc\checkmarki\checkmarkr\checkmarkc\checkmark$\checkmarkl\checkmark$\checkmark^\checkmark\\checkmarkc\checkmarki\checkmarkr\checkmarkc\checkmark$\checkmarka\checkmark$\checkmark^\checkmark\\checkmarkc\checkmarki\checkmarkr\checkmarkc\checkmark$\checkmarkb\checkmark$\checkmark^\checkmark\\checkmarkc\checkmarki\checkmarkr\checkmarkc\checkmark$\checkmarke\checkmark$\checkmark^\checkmark\\checkmarkc\checkmarki\checkmarkr\checkmarkc\checkmark$\checkmarkl\checkmark$\checkmark^\checkmark\\checkmarkc\checkmarki\checkmarkr\checkmarkc\checkmark$\checkmark{\checkmark$\checkmark^\checkmark\\checkmarkc\checkmarki\checkmarkr\checkmarkc\checkmark$\checkmarke\checkmark$\checkmark^\checkmark\\checkmarkc\checkmarki\checkmarkr\checkmarkc\checkmark$\checkmarkq\checkmark$\checkmark^\checkmark\\checkmarkc\checkmarki\checkmarkr\checkmarkc\checkmark$\checkmark:\checkmark$\checkmark^\checkmark\\checkmarkc\checkmarki\checkmarkr\checkmarkc\checkmark$\checkmarkt\checkmark$\checkmark^\checkmark\\checkmarkc\checkmarki\checkmarkr\checkmarkc\checkmark$\checkmarko\checkmark$\checkmark^\checkmark\\checkmarkc\checkmarki\checkmarkr\checkmarkc\checkmark$\checkmarkr\checkmark$\checkmark^\checkmark\\checkmarkc\checkmarki\checkmarkr\checkmarkc\checkmark$\checkmarks\checkmark$\checkmark^\checkmark\\checkmarkc\checkmarki\checkmarkr\checkmarkc\checkmark$\checkmarke\checkmark$\checkmark^\checkmark\\checkmarkc\checkmarki\checkmarkr\checkmarkc\checkmark$\checkmarku\checkmark$\checkmark^\checkmark\\checkmarkc\checkmarki\checkmarkr\checkmarkc\checkmark$\checkmarkr\checkmark$\checkmark^\checkmark\\checkmarkc\checkmarki\checkmarkr\checkmarkc\checkmark$\checkmark}\checkmark$\checkmark^\checkmark\\checkmarkc\checkmarki\checkmarkr\checkmarkc\checkmark$\checkmark
\checkmark$\checkmark^\checkmark\\checkmarkc\checkmarki\checkmarkr\checkmarkc\checkmark$\checkmark\\checkmark$\checkmark^\checkmark\\checkmarkc\checkmarki\checkmarkr\checkmarkc\checkmark$\checkmarke\checkmark$\checkmark^\checkmark\\checkmarkc\checkmarki\checkmarkr\checkmarkc\checkmark$\checkmarkn\checkmark$\checkmark^\checkmark\\checkmarkc\checkmarki\checkmarkr\checkmarkc\checkmark$\checkmarkd\checkmark$\checkmark^\checkmark\\checkmarkc\checkmarki\checkmarkr\checkmarkc\checkmark$\checkmark{\checkmark$\checkmark^\checkmark\\checkmarkc\checkmarki\checkmarkr\checkmarkc\checkmark$\checkmarke\checkmark$\checkmark^\checkmark\\checkmarkc\checkmarki\checkmarkr\checkmarkc\checkmark$\checkmarkq\checkmark$\checkmark^\checkmark\\checkmarkc\checkmarki\checkmarkr\checkmarkc\checkmark$\checkmarku\checkmark$\checkmark^\checkmark\\checkmarkc\checkmarki\checkmarkr\checkmarkc\checkmark$\checkmarka\checkmark$\checkmark^\checkmark\\checkmarkc\checkmarki\checkmarkr\checkmarkc\checkmark$\checkmarkt\checkmark$\checkmark^\checkmark\\checkmarkc\checkmarki\checkmarkr\checkmarkc\checkmark$\checkmarki\checkmark$\checkmark^\checkmark\\checkmarkc\checkmarki\checkmarkr\checkmarkc\checkmark$\checkmarko\checkmark$\checkmark^\checkmark\\checkmarkc\checkmarki\checkmarkr\checkmarkc\checkmark$\checkmarkn\checkmark$\checkmark^\checkmark\\checkmarkc\checkmarki\checkmarkr\checkmarkc\checkmark$\checkmark}\checkmark$\checkmark^\checkmark\\checkmarkc\checkmarki\checkmarkr\checkmarkc\checkmark$\checkmark
\checkmark$\checkmark^\checkmark\\checkmarkc\checkmarki\checkmarkr\checkmarkc\checkmark$\checkmark
\checkmark$\checkmark^\checkmark\\checkmarkc\checkmarki\checkmarkr\checkmarkc\checkmark$\checkmark\\checkmark$\checkmark^\checkmark\\checkmarkc\checkmarki\checkmarkr\checkmarkc\checkmark$\checkmarks\checkmark$\checkmark^\checkmark\\checkmarkc\checkmarki\checkmarkr\checkmarkc\checkmark$\checkmarku\checkmark$\checkmark^\checkmark\\checkmarkc\checkmarki\checkmarkr\checkmarkc\checkmark$\checkmarkb\checkmark$\checkmark^\checkmark\\checkmarkc\checkmarki\checkmarkr\checkmarkc\checkmark$\checkmarks\checkmark$\checkmark^\checkmark\\checkmarkc\checkmarki\checkmarkr\checkmarkc\checkmark$\checkmarke\checkmark$\checkmark^\checkmark\\checkmarkc\checkmarki\checkmarkr\checkmarkc\checkmark$\checkmarkc\checkmark$\checkmark^\checkmark\\checkmarkc\checkmarki\checkmarkr\checkmarkc\checkmark$\checkmarkt\checkmark$\checkmark^\checkmark\\checkmarkc\checkmarki\checkmarkr\checkmarkc\checkmark$\checkmarki\checkmark$\checkmark^\checkmark\\checkmarkc\checkmarki\checkmarkr\checkmarkc\checkmark$\checkmarko\checkmark$\checkmark^\checkmark\\checkmarkc\checkmarki\checkmarkr\checkmarkc\checkmark$\checkmarkn\checkmark$\checkmark^\checkmark\\checkmarkc\checkmarki\checkmarkr\checkmarkc\checkmark$\checkmark{\checkmark$\checkmark^\checkmark\\checkmarkc\checkmarki\checkmarkr\checkmarkc\checkmark$\checkmarkD\checkmark$\checkmark^\checkmark\\checkmarkc\checkmarki\checkmarkr\checkmarkc\checkmark$\checkmarki\checkmark$\checkmark^\checkmark\\checkmarkc\checkmarki\checkmarkr\checkmarkc\checkmark$\checkmarkm\checkmark$\checkmark^\checkmark\\checkmarkc\checkmarki\checkmarkr\checkmarkc\checkmark$\checkmarke\checkmark$\checkmark^\checkmark\\checkmarkc\checkmarki\checkmarkr\checkmarkc\checkmark$\checkmarkn\checkmark$\checkmark^\checkmark\\checkmarkc\checkmarki\checkmarkr\checkmarkc\checkmark$\checkmarks\checkmark$\checkmark^\checkmark\\checkmarkc\checkmarki\checkmarkr\checkmarkc\checkmark$\checkmarki\checkmark$\checkmark^\checkmark\\checkmarkc\checkmarki\checkmarkr\checkmarkc\checkmark$\checkmarko\checkmark$\checkmark^\checkmark\\checkmarkc\checkmarki\checkmarkr\checkmarkc\checkmark$\checkmarkn\checkmark$\checkmark^\checkmark\\checkmarkc\checkmarki\checkmarkr\checkmarkc\checkmark$\checkmarkn\checkmark$\checkmark^\checkmark\\checkmarkc\checkmarki\checkmarkr\checkmarkc\checkmark$\checkmarke\checkmark$\checkmark^\checkmark\\checkmarkc\checkmarki\checkmarkr\checkmarkc\checkmark$\checkmarkm\checkmark$\checkmark^\checkmark\\checkmarkc\checkmarki\checkmarkr\checkmarkc\checkmark$\checkmarke\checkmark$\checkmark^\checkmark\\checkmarkc\checkmarki\checkmarkr\checkmarkc\checkmark$\checkmarkn\checkmark$\checkmark^\checkmark\\checkmarkc\checkmarki\checkmarkr\checkmarkc\checkmark$\checkmarkt\checkmark$\checkmark^\checkmark\\checkmarkc\checkmarki\checkmarkr\checkmarkc\checkmark$\checkmark \checkmark$\checkmark^\checkmark\\checkmarkc\checkmarki\checkmarkr\checkmarkc\checkmark$\checkmarkd\checkmark$\checkmark^\checkmark\\checkmarkc\checkmarki\checkmarkr\checkmarkc\checkmark$\checkmarke\checkmark$\checkmark^\checkmark\\checkmarkc\checkmarki\checkmarkr\checkmarkc\checkmark$\checkmark \checkmark$\checkmark^\checkmark\\checkmarkc\checkmarki\checkmarkr\checkmarkc\checkmark$\checkmarkl\checkmark$\checkmark^\checkmark\\checkmarkc\checkmarki\checkmarkr\checkmarkc\checkmark$\checkmarka\checkmark$\checkmark^\checkmark\\checkmarkc\checkmarki\checkmarkr\checkmarkc\checkmark$\checkmark \checkmark$\checkmark^\checkmark\\checkmarkc\checkmarki\checkmarkr\checkmarkc\checkmark$\checkmarkp\checkmark$\checkmark^\checkmark\\checkmarkc\checkmarki\checkmarkr\checkmarkc\checkmark$\checkmarku\checkmark$\checkmark^\checkmark\\checkmarkc\checkmarki\checkmarkr\checkmarkc\checkmark$\checkmarki\checkmark$\checkmark^\checkmark\\checkmarkc\checkmarki\checkmarkr\checkmarkc\checkmark$\checkmarks\checkmark$\checkmark^\checkmark\\checkmarkc\checkmarki\checkmarkr\checkmarkc\checkmark$\checkmarks\checkmark$\checkmark^\checkmark\\checkmarkc\checkmarki\checkmarkr\checkmarkc\checkmark$\checkmarka\checkmark$\checkmark^\checkmark\\checkmarkc\checkmarki\checkmarkr\checkmarkc\checkmark$\checkmarkn\checkmark$\checkmark^\checkmark\\checkmarkc\checkmarki\checkmarkr\checkmarkc\checkmark$\checkmarkc\checkmark$\checkmark^\checkmark\\checkmarkc\checkmarki\checkmarkr\checkmarkc\checkmark$\checkmarke\checkmark$\checkmark^\checkmark\\checkmarkc\checkmarki\checkmarkr\checkmarkc\checkmark$\checkmark \checkmark$\checkmark^\checkmark\\checkmarkc\checkmarki\checkmarkr\checkmarkc\checkmark$\checkmarkd\checkmark$\checkmark^\checkmark\\checkmarkc\checkmarki\checkmarkr\checkmarkc\checkmark$\checkmarke\checkmark$\checkmark^\checkmark\\checkmarkc\checkmarki\checkmarkr\checkmarkc\checkmark$\checkmark \checkmark$\checkmark^\checkmark\\checkmarkc\checkmarki\checkmarkr\checkmarkc\checkmark$\checkmarkl\checkmark$\checkmark^\checkmark\\checkmarkc\checkmarki\checkmarkr\checkmarkc\checkmark$\checkmarka\checkmark$\checkmark^\checkmark\\checkmarkc\checkmarki\checkmarkr\checkmarkc\checkmark$\checkmark \checkmark$\checkmark^\checkmark\\checkmarkc\checkmarki\checkmarkr\checkmarkc\checkmark$\checkmarkb\checkmark$\checkmark^\checkmark\\checkmarkc\checkmarki\checkmarkr\checkmarkc\checkmark$\checkmarkr\checkmark$\checkmark^\checkmark\\checkmarkc\checkmarki\checkmarkr\checkmarkc\checkmark$\checkmarko\checkmark$\checkmark^\checkmark\\checkmarkc\checkmarki\checkmarkr\checkmarkc\checkmark$\checkmarkc\checkmark$\checkmark^\checkmark\\checkmarkc\checkmarki\checkmarkr\checkmarkc\checkmark$\checkmarkh\checkmark$\checkmark^\checkmark\\checkmarkc\checkmarki\checkmarkr\checkmarkc\checkmark$\checkmarke\checkmark$\checkmark^\checkmark\\checkmarkc\checkmarki\checkmarkr\checkmarkc\checkmark$\checkmark}\checkmark$\checkmark^\checkmark\\checkmarkc\checkmarki\checkmarkr\checkmarkc\checkmark$\checkmark
\checkmark$\checkmark^\checkmark\\checkmarkc\checkmarki\checkmarkr\checkmarkc\checkmark$\checkmarkL\checkmark$\checkmark^\checkmark\\checkmarkc\checkmarki\checkmarkr\checkmarkc\checkmark$\checkmarka\checkmark$\checkmark^\checkmark\\checkmarkc\checkmarki\checkmarkr\checkmarkc\checkmark$\checkmark \checkmark$\checkmark^\checkmark\\checkmarkc\checkmarki\checkmarkr\checkmarkc\checkmark$\checkmarkp\checkmark$\checkmark^\checkmark\\checkmarkc\checkmarki\checkmarkr\checkmarkc\checkmark$\checkmarku\checkmark$\checkmark^\checkmark\\checkmarkc\checkmarki\checkmarkr\checkmarkc\checkmark$\checkmarki\checkmark$\checkmark^\checkmark\\checkmarkc\checkmarki\checkmarkr\checkmarkc\checkmark$\checkmarks\checkmark$\checkmark^\checkmark\\checkmarkc\checkmarki\checkmarkr\checkmarkc\checkmark$\checkmarks\checkmark$\checkmark^\checkmark\\checkmarkc\checkmarki\checkmarkr\checkmarkc\checkmark$\checkmarka\checkmark$\checkmark^\checkmark\\checkmarkc\checkmarki\checkmarkr\checkmarkc\checkmark$\checkmarkn\checkmark$\checkmark^\checkmark\\checkmarkc\checkmarki\checkmarkr\checkmarkc\checkmark$\checkmarkc\checkmark$\checkmark^\checkmark\\checkmarkc\checkmarki\checkmarkr\checkmarkc\checkmark$\checkmarke\checkmark$\checkmark^\checkmark\\checkmarkc\checkmarki\checkmarkr\checkmarkc\checkmark$\checkmark \checkmark$\checkmark^\checkmark\\checkmarkc\checkmarki\checkmarkr\checkmarkc\checkmark$\checkmarkm\checkmark$\checkmark^\checkmark\\checkmarkc\checkmarki\checkmarkr\checkmarkc\checkmark$\checkmarké\checkmark$\checkmark^\checkmark\\checkmarkc\checkmarki\checkmarkr\checkmarkc\checkmark$\checkmarkc\checkmark$\checkmark^\checkmark\\checkmarkc\checkmarki\checkmarkr\checkmarkc\checkmark$\checkmarka\checkmark$\checkmark^\checkmark\\checkmarkc\checkmarki\checkmarkr\checkmarkc\checkmark$\checkmarkn\checkmark$\checkmark^\checkmark\\checkmarkc\checkmarki\checkmarkr\checkmarkc\checkmark$\checkmarki\checkmark$\checkmark^\checkmark\\checkmarkc\checkmarki\checkmarkr\checkmarkc\checkmark$\checkmarkq\checkmark$\checkmark^\checkmark\\checkmarkc\checkmarki\checkmarkr\checkmarkc\checkmark$\checkmarku\checkmark$\checkmark^\checkmark\\checkmarkc\checkmarki\checkmarkr\checkmarkc\checkmark$\checkmarke\checkmark$\checkmark^\checkmark\\checkmarkc\checkmarki\checkmarkr\checkmarkc\checkmark$\checkmark \checkmark$\checkmark^\checkmark\\checkmarkc\checkmarki\checkmarkr\checkmarkc\checkmark$\checkmarkd\checkmark$\checkmark^\checkmark\\checkmarkc\checkmarki\checkmarkr\checkmarkc\checkmark$\checkmarke\checkmark$\checkmark^\checkmark\\checkmarkc\checkmarki\checkmarkr\checkmarkc\checkmark$\checkmark \checkmark$\checkmark^\checkmark\\checkmarkc\checkmarki\checkmarkr\checkmarkc\checkmark$\checkmarkc\checkmark$\checkmark^\checkmark\\checkmarkc\checkmarki\checkmarkr\checkmarkc\checkmark$\checkmarko\checkmark$\checkmark^\checkmark\\checkmarkc\checkmarki\checkmarkr\checkmarkc\checkmark$\checkmarku\checkmark$\checkmark^\checkmark\\checkmarkc\checkmarki\checkmarkr\checkmarkc\checkmark$\checkmarkp\checkmark$\checkmark^\checkmark\\checkmarkc\checkmarki\checkmarkr\checkmarkc\checkmark$\checkmarke\checkmark$\checkmark^\checkmark\\checkmarkc\checkmarki\checkmarkr\checkmarkc\checkmark$\checkmark \checkmark$\checkmark^\checkmark\\checkmarkc\checkmarki\checkmarkr\checkmarkc\checkmark$\checkmarkr\checkmark$\checkmark^\checkmark\\checkmarkc\checkmarki\checkmarkr\checkmarkc\checkmark$\checkmarke\checkmark$\checkmark^\checkmark\\checkmarkc\checkmarki\checkmarkr\checkmarkc\checkmark$\checkmarkq\checkmark$\checkmark^\checkmark\\checkmarkc\checkmarki\checkmarkr\checkmarkc\checkmark$\checkmarku\checkmark$\checkmark^\checkmark\\checkmarkc\checkmarki\checkmarkr\checkmarkc\checkmark$\checkmarki\checkmark$\checkmark^\checkmark\\checkmarkc\checkmarki\checkmarkr\checkmarkc\checkmark$\checkmarks\checkmark$\checkmark^\checkmark\\checkmarkc\checkmarki\checkmarkr\checkmarkc\checkmark$\checkmarke\checkmark$\checkmark^\checkmark\\checkmarkc\checkmarki\checkmarkr\checkmarkc\checkmark$\checkmark \checkmark$\checkmark^\checkmark\\checkmarkc\checkmarki\checkmarkr\checkmarkc\checkmark$\checkmarkà\checkmark$\checkmark^\checkmark\\checkmarkc\checkmarki\checkmarkr\checkmarkc\checkmark$\checkmark \checkmark$\checkmark^\checkmark\\checkmarkc\checkmarki\checkmarkr\checkmarkc\checkmark$\checkmarkl\checkmark$\checkmark^\checkmark\\checkmarkc\checkmarki\checkmarkr\checkmarkc\checkmark$\checkmark'\checkmark$\checkmark^\checkmark\\checkmarkc\checkmarki\checkmarkr\checkmarkc\checkmark$\checkmarko\checkmark$\checkmark^\checkmark\\checkmarkc\checkmarki\checkmarkr\checkmarkc\checkmark$\checkmarku\checkmark$\checkmark^\checkmark\\checkmarkc\checkmarki\checkmarkr\checkmarkc\checkmark$\checkmarkt\checkmark$\checkmark^\checkmark\\checkmarkc\checkmarki\checkmarkr\checkmarkc\checkmark$\checkmarki\checkmark$\checkmark^\checkmark\\checkmarkc\checkmarki\checkmarkr\checkmarkc\checkmark$\checkmarkl\checkmark$\checkmark^\checkmark\\checkmarkc\checkmarki\checkmarkr\checkmarkc\checkmark$\checkmark,\checkmark$\checkmark^\checkmark\\checkmarkc\checkmarki\checkmarkr\checkmarkc\checkmark$\checkmark \checkmark$\checkmark^\checkmark\\checkmarkc\checkmarki\checkmarkr\checkmarkc\checkmark$\checkmarka\checkmark$\checkmark^\checkmark\\checkmarkc\checkmarki\checkmarkr\checkmarkc\checkmark$\checkmarkv\checkmark$\checkmark^\checkmark\\checkmarkc\checkmarki\checkmarkr\checkmarkc\checkmark$\checkmarke\checkmark$\checkmark^\checkmark\\checkmarkc\checkmarki\checkmarkr\checkmarkc\checkmark$\checkmarkc\checkmark$\checkmark^\checkmark\\checkmarkc\checkmarki\checkmarkr\checkmarkc\checkmark$\checkmark \checkmark$\checkmark^\checkmark\\checkmarkc\checkmarki\checkmarkr\checkmarkc\checkmark$\checkmark$\checkmark$\checkmark^\checkmark\\checkmarkc\checkmarki\checkmarkr\checkmarkc\checkmark$\checkmarkV\checkmark$\checkmark^\checkmark\\checkmarkc\checkmarki\checkmarkr\checkmarkc\checkmark$\checkmark_\checkmark$\checkmark^\checkmark\\checkmarkc\checkmarki\checkmarkr\checkmarkc\checkmark$\checkmarkc\checkmark$\checkmark^\checkmark\\checkmarkc\checkmarki\checkmarkr\checkmarkc\checkmark$\checkmark \checkmark$\checkmark^\checkmark\\checkmarkc\checkmarki\checkmarkr\checkmarkc\checkmark$\checkmark=\checkmark$\checkmark^\checkmark\\checkmarkc\checkmarki\checkmarkr\checkmarkc\checkmark$\checkmark \checkmark$\checkmark^\checkmark\\checkmarkc\checkmarki\checkmarkr\checkmarkc\checkmark$\checkmark1\checkmark$\checkmark^\checkmark\\checkmarkc\checkmarki\checkmarkr\checkmarkc\checkmark$\checkmark5\checkmark$\checkmark^\checkmark\\checkmarkc\checkmarki\checkmarkr\checkmarkc\checkmark$\checkmark0\checkmark$\checkmark^\checkmark\\checkmarkc\checkmarki\checkmarkr\checkmarkc\checkmark$\checkmark \checkmark$\checkmark^\checkmark\\checkmarkc\checkmarki\checkmarkr\checkmarkc\checkmark$\checkmark\\checkmark$\checkmark^\checkmark\\checkmarkc\checkmarki\checkmarkr\checkmarkc\checkmark$\checkmarkt\checkmark$\checkmark^\checkmark\\checkmarkc\checkmarki\checkmarkr\checkmarkc\checkmark$\checkmarke\checkmark$\checkmark^\checkmark\\checkmarkc\checkmarki\checkmarkr\checkmarkc\checkmark$\checkmarkx\checkmark$\checkmark^\checkmark\\checkmarkc\checkmarki\checkmarkr\checkmarkc\checkmark$\checkmarkt\checkmark$\checkmark^\checkmark\\checkmarkc\checkmarki\checkmarkr\checkmarkc\checkmark$\checkmark{\checkmark$\checkmark^\checkmark\\checkmarkc\checkmarki\checkmarkr\checkmarkc\checkmark$\checkmark \checkmark$\checkmark^\checkmark\\checkmarkc\checkmarki\checkmarkr\checkmarkc\checkmark$\checkmarkm\checkmark$\checkmark^\checkmark\\checkmarkc\checkmarki\checkmarkr\checkmarkc\checkmark$\checkmark/\checkmark$\checkmark^\checkmark\\checkmarkc\checkmarki\checkmarkr\checkmarkc\checkmark$\checkmarkm\checkmark$\checkmark^\checkmark\\checkmarkc\checkmarki\checkmarkr\checkmarkc\checkmark$\checkmarki\checkmark$\checkmark^\checkmark\\checkmarkc\checkmarki\checkmarkr\checkmarkc\checkmark$\checkmarkn\checkmark$\checkmark^\checkmark\\checkmarkc\checkmarki\checkmarkr\checkmarkc\checkmark$\checkmark}\checkmark$\checkmark^\checkmark\\checkmarkc\checkmarki\checkmarkr\checkmarkc\checkmark$\checkmark \checkmark$\checkmark^\checkmark\\checkmarkc\checkmarki\checkmarkr\checkmarkc\checkmark$\checkmark=\checkmark$\checkmark^\checkmark\\checkmarkc\checkmarki\checkmarkr\checkmarkc\checkmark$\checkmark \checkmark$\checkmark^\checkmark\\checkmarkc\checkmarki\checkmarkr\checkmarkc\checkmark$\checkmark2\checkmark$\checkmark^\checkmark\\checkmarkc\checkmarki\checkmarkr\checkmarkc\checkmark$\checkmark{\checkmark$\checkmark^\checkmark\\checkmarkc\checkmarki\checkmarkr\checkmarkc\checkmark$\checkmark,\checkmark$\checkmark^\checkmark\\checkmarkc\checkmarki\checkmarkr\checkmarkc\checkmark$\checkmark}\checkmark$\checkmark^\checkmark\\checkmarkc\checkmarki\checkmarkr\checkmarkc\checkmark$\checkmark5\checkmark$\checkmark^\checkmark\\checkmarkc\checkmarki\checkmarkr\checkmarkc\checkmark$\checkmark \checkmark$\checkmark^\checkmark\\checkmarkc\checkmarki\checkmarkr\checkmarkc\checkmark$\checkmark\\checkmark$\checkmark^\checkmark\\checkmarkc\checkmarki\checkmarkr\checkmarkc\checkmark$\checkmarkt\checkmark$\checkmark^\checkmark\\checkmarkc\checkmarki\checkmarkr\checkmarkc\checkmark$\checkmarke\checkmark$\checkmark^\checkmark\\checkmarkc\checkmarki\checkmarkr\checkmarkc\checkmark$\checkmarkx\checkmark$\checkmark^\checkmark\\checkmarkc\checkmarki\checkmarkr\checkmarkc\checkmark$\checkmarkt\checkmark$\checkmark^\checkmark\\checkmarkc\checkmarki\checkmarkr\checkmarkc\checkmark$\checkmark{\checkmark$\checkmark^\checkmark\\checkmarkc\checkmarki\checkmarkr\checkmarkc\checkmark$\checkmark \checkmark$\checkmark^\checkmark\\checkmarkc\checkmarki\checkmarkr\checkmarkc\checkmark$\checkmarkm\checkmark$\checkmark^\checkmark\\checkmarkc\checkmarki\checkmarkr\checkmarkc\checkmark$\checkmark/\checkmark$\checkmark^\checkmark\\checkmarkc\checkmarki\checkmarkr\checkmarkc\checkmark$\checkmarks\checkmark$\checkmark^\checkmark\\checkmarkc\checkmarki\checkmarkr\checkmarkc\checkmark$\checkmark}\checkmark$\checkmark^\checkmark\\checkmarkc\checkmarki\checkmarkr\checkmarkc\checkmark$\checkmark$\checkmark$\checkmark^\checkmark\\checkmarkc\checkmarki\checkmarkr\checkmarkc\checkmark$\checkmark \checkmark$\checkmark^\checkmark\\checkmarkc\checkmarki\checkmarkr\checkmarkc\checkmark$\checkmark:\checkmark$\checkmark^\checkmark\\checkmarkc\checkmarki\checkmarkr\checkmarkc\checkmark$\checkmark
\checkmark$\checkmark^\checkmark\\checkmarkc\checkmarki\checkmarkr\checkmarkc\checkmark$\checkmark\\checkmark$\checkmark^\checkmark\\checkmarkc\checkmarki\checkmarkr\checkmarkc\checkmark$\checkmarkb\checkmark$\checkmark^\checkmark\\checkmarkc\checkmarki\checkmarkr\checkmarkc\checkmark$\checkmarke\checkmark$\checkmark^\checkmark\\checkmarkc\checkmarki\checkmarkr\checkmarkc\checkmark$\checkmarkg\checkmark$\checkmark^\checkmark\\checkmarkc\checkmarki\checkmarkr\checkmarkc\checkmark$\checkmarki\checkmark$\checkmark^\checkmark\\checkmarkc\checkmarki\checkmarkr\checkmarkc\checkmark$\checkmarkn\checkmark$\checkmark^\checkmark\\checkmarkc\checkmarki\checkmarkr\checkmarkc\checkmark$\checkmark{\checkmark$\checkmark^\checkmark\\checkmarkc\checkmarki\checkmarkr\checkmarkc\checkmark$\checkmarke\checkmark$\checkmark^\checkmark\\checkmarkc\checkmarki\checkmarkr\checkmarkc\checkmark$\checkmarkq\checkmark$\checkmark^\checkmark\\checkmarkc\checkmarki\checkmarkr\checkmarkc\checkmark$\checkmarku\checkmark$\checkmark^\checkmark\\checkmarkc\checkmarki\checkmarkr\checkmarkc\checkmark$\checkmarka\checkmark$\checkmark^\checkmark\\checkmarkc\checkmarki\checkmarkr\checkmarkc\checkmark$\checkmarkt\checkmark$\checkmark^\checkmark\\checkmarkc\checkmarki\checkmarkr\checkmarkc\checkmark$\checkmarki\checkmark$\checkmark^\checkmark\\checkmarkc\checkmarki\checkmarkr\checkmarkc\checkmark$\checkmarko\checkmark$\checkmark^\checkmark\\checkmarkc\checkmarki\checkmarkr\checkmarkc\checkmark$\checkmarkn\checkmark$\checkmark^\checkmark\\checkmarkc\checkmarki\checkmarkr\checkmarkc\checkmark$\checkmark}\checkmark$\checkmark^\checkmark\\checkmarkc\checkmarki\checkmarkr\checkmarkc\checkmark$\checkmark
\checkmark$\checkmark^\checkmark\\checkmarkc\checkmarki\checkmarkr\checkmarkc\checkmark$\checkmark \checkmark$\checkmark^\checkmark\\checkmarkc\checkmarki\checkmarkr\checkmarkc\checkmark$\checkmark \checkmark$\checkmark^\checkmark\\checkmarkc\checkmarki\checkmarkr\checkmarkc\checkmark$\checkmark \checkmark$\checkmark^\checkmark\\checkmarkc\checkmarki\checkmarkr\checkmarkc\checkmark$\checkmark \checkmark$\checkmark^\checkmark\\checkmarkc\checkmarki\checkmarkr\checkmarkc\checkmark$\checkmarkP\checkmark$\checkmark^\checkmark\\checkmarkc\checkmarki\checkmarkr\checkmarkc\checkmark$\checkmark_\checkmark$\checkmark^\checkmark\\checkmarkc\checkmarki\checkmarkr\checkmarkc\checkmark$\checkmarkc\checkmark$\checkmark^\checkmark\\checkmarkc\checkmarki\checkmarkr\checkmarkc\checkmark$\checkmark \checkmark$\checkmark^\checkmark\\checkmarkc\checkmarki\checkmarkr\checkmarkc\checkmark$\checkmark=\checkmark$\checkmark^\checkmark\\checkmarkc\checkmarki\checkmarkr\checkmarkc\checkmark$\checkmark \checkmark$\checkmark^\checkmark\\checkmarkc\checkmarki\checkmarkr\checkmarkc\checkmark$\checkmarkF\checkmark$\checkmark^\checkmark\\checkmarkc\checkmarki\checkmarkr\checkmarkc\checkmark$\checkmark_\checkmark$\checkmark^\checkmark\\checkmarkc\checkmarki\checkmarkr\checkmarkc\checkmark$\checkmarkc\checkmark$\checkmark^\checkmark\\checkmarkc\checkmarki\checkmarkr\checkmarkc\checkmark$\checkmark \checkmark$\checkmark^\checkmark\\checkmarkc\checkmarki\checkmarkr\checkmarkc\checkmark$\checkmark\\checkmark$\checkmark^\checkmark\\checkmarkc\checkmarki\checkmarkr\checkmarkc\checkmark$\checkmarkt\checkmark$\checkmark^\checkmark\\checkmarkc\checkmarki\checkmarkr\checkmarkc\checkmark$\checkmarki\checkmark$\checkmark^\checkmark\\checkmarkc\checkmarki\checkmarkr\checkmarkc\checkmark$\checkmarkm\checkmark$\checkmark^\checkmark\\checkmarkc\checkmarki\checkmarkr\checkmarkc\checkmark$\checkmarke\checkmark$\checkmark^\checkmark\\checkmarkc\checkmarki\checkmarkr\checkmarkc\checkmark$\checkmarks\checkmark$\checkmark^\checkmark\\checkmarkc\checkmarki\checkmarkr\checkmarkc\checkmark$\checkmark \checkmark$\checkmark^\checkmark\\checkmarkc\checkmarki\checkmarkr\checkmarkc\checkmark$\checkmarkV\checkmark$\checkmark^\checkmark\\checkmarkc\checkmarki\checkmarkr\checkmarkc\checkmark$\checkmark_\checkmark$\checkmark^\checkmark\\checkmarkc\checkmarki\checkmarkr\checkmarkc\checkmark$\checkmarkc\checkmark$\checkmark^\checkmark\\checkmarkc\checkmarki\checkmarkr\checkmarkc\checkmark$\checkmark \checkmark$\checkmark^\checkmark\\checkmarkc\checkmarki\checkmarkr\checkmarkc\checkmark$\checkmark=\checkmark$\checkmark^\checkmark\\checkmarkc\checkmarki\checkmarkr\checkmarkc\checkmark$\checkmark \checkmark$\checkmark^\checkmark\\checkmarkc\checkmarki\checkmarkr\checkmarkc\checkmark$\checkmark7\checkmark$\checkmark^\checkmark\\checkmarkc\checkmarki\checkmarkr\checkmarkc\checkmark$\checkmark0\checkmark$\checkmark^\checkmark\\checkmarkc\checkmarki\checkmarkr\checkmarkc\checkmark$\checkmark \checkmark$\checkmark^\checkmark\\checkmarkc\checkmarki\checkmarkr\checkmarkc\checkmark$\checkmark\\checkmark$\checkmark^\checkmark\\checkmarkc\checkmarki\checkmarkr\checkmarkc\checkmark$\checkmarkt\checkmark$\checkmark^\checkmark\\checkmarkc\checkmarki\checkmarkr\checkmarkc\checkmark$\checkmarki\checkmark$\checkmark^\checkmark\\checkmarkc\checkmarki\checkmarkr\checkmarkc\checkmark$\checkmarkm\checkmark$\checkmark^\checkmark\\checkmarkc\checkmarki\checkmarkr\checkmarkc\checkmark$\checkmarke\checkmark$\checkmark^\checkmark\\checkmarkc\checkmarki\checkmarkr\checkmarkc\checkmark$\checkmarks\checkmark$\checkmark^\checkmark\\checkmarkc\checkmarki\checkmarkr\checkmarkc\checkmark$\checkmark \checkmark$\checkmark^\checkmark\\checkmarkc\checkmarki\checkmarkr\checkmarkc\checkmark$\checkmark2\checkmark$\checkmark^\checkmark\\checkmarkc\checkmarki\checkmarkr\checkmarkc\checkmark$\checkmark{\checkmark$\checkmark^\checkmark\\checkmarkc\checkmarki\checkmarkr\checkmarkc\checkmark$\checkmark,\checkmark$\checkmark^\checkmark\\checkmarkc\checkmarki\checkmarkr\checkmarkc\checkmark$\checkmark}\checkmark$\checkmark^\checkmark\\checkmarkc\checkmarki\checkmarkr\checkmarkc\checkmark$\checkmark5\checkmark$\checkmark^\checkmark\\checkmarkc\checkmarki\checkmarkr\checkmarkc\checkmark$\checkmark \checkmark$\checkmark^\checkmark\\checkmarkc\checkmarki\checkmarkr\checkmarkc\checkmark$\checkmark=\checkmark$\checkmark^\checkmark\\checkmarkc\checkmarki\checkmarkr\checkmarkc\checkmark$\checkmark \checkmark$\checkmark^\checkmark\\checkmarkc\checkmarki\checkmarkr\checkmarkc\checkmark$\checkmark\\checkmark$\checkmark^\checkmark\\checkmarkc\checkmarki\checkmarkr\checkmarkc\checkmark$\checkmarkb\checkmark$\checkmark^\checkmark\\checkmarkc\checkmarki\checkmarkr\checkmarkc\checkmark$\checkmarko\checkmark$\checkmark^\checkmark\\checkmarkc\checkmarki\checkmarkr\checkmarkc\checkmark$\checkmarkx\checkmark$\checkmark^\checkmark\\checkmarkc\checkmarki\checkmarkr\checkmarkc\checkmark$\checkmarke\checkmark$\checkmark^\checkmark\\checkmarkc\checkmarki\checkmarkr\checkmarkc\checkmark$\checkmarkd\checkmark$\checkmark^\checkmark\\checkmarkc\checkmarki\checkmarkr\checkmarkc\checkmark$\checkmark{\checkmark$\checkmark^\checkmark\\checkmarkc\checkmarki\checkmarkr\checkmarkc\checkmark$\checkmark1\checkmark$\checkmark^\checkmark\\checkmarkc\checkmarki\checkmarkr\checkmarkc\checkmark$\checkmark7\checkmark$\checkmark^\checkmark\\checkmarkc\checkmarki\checkmarkr\checkmarkc\checkmark$\checkmark5\checkmark$\checkmark^\checkmark\\checkmarkc\checkmarki\checkmarkr\checkmarkc\checkmark$\checkmark \checkmark$\checkmark^\checkmark\\checkmarkc\checkmarki\checkmarkr\checkmarkc\checkmark$\checkmark\\checkmark$\checkmark^\checkmark\\checkmarkc\checkmarki\checkmarkr\checkmarkc\checkmark$\checkmarkt\checkmark$\checkmark^\checkmark\\checkmarkc\checkmarki\checkmarkr\checkmarkc\checkmark$\checkmarke\checkmark$\checkmark^\checkmark\\checkmarkc\checkmarki\checkmarkr\checkmarkc\checkmark$\checkmarkx\checkmark$\checkmark^\checkmark\\checkmarkc\checkmarki\checkmarkr\checkmarkc\checkmark$\checkmarkt\checkmark$\checkmark^\checkmark\\checkmarkc\checkmarki\checkmarkr\checkmarkc\checkmark$\checkmark{\checkmark$\checkmark^\checkmark\\checkmarkc\checkmarki\checkmarkr\checkmarkc\checkmark$\checkmark \checkmark$\checkmark^\checkmark\\checkmarkc\checkmarki\checkmarkr\checkmarkc\checkmark$\checkmarkW\checkmark$\checkmark^\checkmark\\checkmarkc\checkmarki\checkmarkr\checkmarkc\checkmark$\checkmark}\checkmark$\checkmark^\checkmark\\checkmarkc\checkmarki\checkmarkr\checkmarkc\checkmark$\checkmark}\checkmark$\checkmark^\checkmark\\checkmarkc\checkmarki\checkmarkr\checkmarkc\checkmark$\checkmark
\checkmark$\checkmark^\checkmark\\checkmarkc\checkmarki\checkmarkr\checkmarkc\checkmark$\checkmark\\checkmark$\checkmark^\checkmark\\checkmarkc\checkmarki\checkmarkr\checkmarkc\checkmark$\checkmarke\checkmark$\checkmark^\checkmark\\checkmarkc\checkmarki\checkmarkr\checkmarkc\checkmark$\checkmarkn\checkmark$\checkmark^\checkmark\\checkmarkc\checkmarki\checkmarkr\checkmarkc\checkmark$\checkmarkd\checkmark$\checkmark^\checkmark\\checkmarkc\checkmarki\checkmarkr\checkmarkc\checkmark$\checkmark{\checkmark$\checkmark^\checkmark\\checkmarkc\checkmarki\checkmarkr\checkmarkc\checkmark$\checkmarke\checkmark$\checkmark^\checkmark\\checkmarkc\checkmarki\checkmarkr\checkmarkc\checkmark$\checkmarkq\checkmark$\checkmark^\checkmark\\checkmarkc\checkmarki\checkmarkr\checkmarkc\checkmark$\checkmarku\checkmark$\checkmark^\checkmark\\checkmarkc\checkmarki\checkmarkr\checkmarkc\checkmark$\checkmarka\checkmark$\checkmark^\checkmark\\checkmarkc\checkmarki\checkmarkr\checkmarkc\checkmark$\checkmarkt\checkmark$\checkmark^\checkmark\\checkmarkc\checkmarki\checkmarkr\checkmarkc\checkmark$\checkmarki\checkmark$\checkmark^\checkmark\\checkmarkc\checkmarki\checkmarkr\checkmarkc\checkmark$\checkmarko\checkmark$\checkmark^\checkmark\\checkmarkc\checkmarki\checkmarkr\checkmarkc\checkmark$\checkmarkn\checkmark$\checkmark^\checkmark\\checkmarkc\checkmarki\checkmarkr\checkmarkc\checkmark$\checkmark}\checkmark$\checkmark^\checkmark\\checkmarkc\checkmarki\checkmarkr\checkmarkc\checkmark$\checkmark
\checkmark$\checkmark^\checkmark\\checkmarkc\checkmarki\checkmarkr\checkmarkc\checkmark$\checkmark
\checkmark$\checkmark^\checkmark\\checkmarkc\checkmarki\checkmarkr\checkmarkc\checkmark$\checkmarkE\checkmark$\checkmark^\checkmark\\checkmarkc\checkmarki\checkmarkr\checkmarkc\checkmark$\checkmarkn\checkmark$\checkmark^\checkmark\\checkmarkc\checkmarki\checkmarkr\checkmarkc\checkmark$\checkmark \checkmark$\checkmark^\checkmark\\checkmarkc\checkmarki\checkmarkr\checkmarkc\checkmark$\checkmarkc\checkmark$\checkmark^\checkmark\\checkmarkc\checkmarki\checkmarkr\checkmarkc\checkmark$\checkmarko\checkmark$\checkmark^\checkmark\\checkmarkc\checkmarki\checkmarkr\checkmarkc\checkmark$\checkmarkn\checkmark$\checkmark^\checkmark\\checkmarkc\checkmarki\checkmarkr\checkmarkc\checkmark$\checkmarks\checkmark$\checkmark^\checkmark\\checkmarkc\checkmarki\checkmarkr\checkmarkc\checkmark$\checkmarki\checkmark$\checkmark^\checkmark\\checkmarkc\checkmarki\checkmarkr\checkmarkc\checkmark$\checkmarkd\checkmark$\checkmark^\checkmark\\checkmarkc\checkmarki\checkmarkr\checkmarkc\checkmark$\checkmarké\checkmark$\checkmark^\checkmark\\checkmarkc\checkmarki\checkmarkr\checkmarkc\checkmark$\checkmarkr\checkmark$\checkmark^\checkmark\\checkmarkc\checkmarki\checkmarkr\checkmarkc\checkmark$\checkmarka\checkmark$\checkmark^\checkmark\\checkmarkc\checkmarki\checkmarkr\checkmarkc\checkmark$\checkmarkn\checkmark$\checkmark^\checkmark\\checkmarkc\checkmarki\checkmarkr\checkmarkc\checkmark$\checkmarkt\checkmark$\checkmark^\checkmark\\checkmarkc\checkmarki\checkmarkr\checkmarkc\checkmark$\checkmark \checkmark$\checkmark^\checkmark\\checkmarkc\checkmarki\checkmarkr\checkmarkc\checkmark$\checkmarku\checkmark$\checkmark^\checkmark\\checkmarkc\checkmarki\checkmarkr\checkmarkc\checkmark$\checkmarkn\checkmark$\checkmark^\checkmark\\checkmarkc\checkmarki\checkmarkr\checkmarkc\checkmark$\checkmark \checkmark$\checkmark^\checkmark\\checkmarkc\checkmarki\checkmarkr\checkmarkc\checkmark$\checkmarkr\checkmark$\checkmark^\checkmark\\checkmarkc\checkmarki\checkmarkr\checkmarkc\checkmark$\checkmarke\checkmark$\checkmark^\checkmark\\checkmarkc\checkmarki\checkmarkr\checkmarkc\checkmark$\checkmarkn\checkmark$\checkmark^\checkmark\\checkmarkc\checkmarki\checkmarkr\checkmarkc\checkmark$\checkmarkd\checkmark$\checkmark^\checkmark\\checkmarkc\checkmarki\checkmarkr\checkmarkc\checkmark$\checkmarke\checkmark$\checkmark^\checkmark\\checkmarkc\checkmarki\checkmarkr\checkmarkc\checkmark$\checkmarkm\checkmark$\checkmark^\checkmark\\checkmarkc\checkmarki\checkmarkr\checkmarkc\checkmark$\checkmarke\checkmark$\checkmark^\checkmark\\checkmarkc\checkmarki\checkmarkr\checkmarkc\checkmark$\checkmarkn\checkmark$\checkmark^\checkmark\\checkmarkc\checkmarki\checkmarkr\checkmarkc\checkmark$\checkmarkt\checkmark$\checkmark^\checkmark\\checkmarkc\checkmarki\checkmarkr\checkmarkc\checkmark$\checkmark \checkmark$\checkmark^\checkmark\\checkmarkc\checkmarki\checkmarkr\checkmarkc\checkmark$\checkmarkd\checkmark$\checkmark^\checkmark\\checkmarkc\checkmarki\checkmarkr\checkmarkc\checkmark$\checkmarke\checkmark$\checkmark^\checkmark\\checkmarkc\checkmarki\checkmarkr\checkmarkc\checkmark$\checkmark \checkmark$\checkmark^\checkmark\\checkmarkc\checkmarki\checkmarkr\checkmarkc\checkmark$\checkmarkt\checkmark$\checkmark^\checkmark\\checkmarkc\checkmarki\checkmarkr\checkmarkc\checkmark$\checkmarkr\checkmark$\checkmark^\checkmark\\checkmarkc\checkmarki\checkmarkr\checkmarkc\checkmark$\checkmarka\checkmark$\checkmark^\checkmark\\checkmarkc\checkmarki\checkmarkr\checkmarkc\checkmark$\checkmarkn\checkmark$\checkmark^\checkmark\\checkmarkc\checkmarki\checkmarkr\checkmarkc\checkmark$\checkmarks\checkmark$\checkmark^\checkmark\\checkmarkc\checkmarki\checkmarkr\checkmarkc\checkmark$\checkmarkm\checkmark$\checkmark^\checkmark\\checkmarkc\checkmarki\checkmarkr\checkmarkc\checkmark$\checkmarki\checkmark$\checkmark^\checkmark\\checkmarkc\checkmarki\checkmarkr\checkmarkc\checkmark$\checkmarks\checkmark$\checkmark^\checkmark\\checkmarkc\checkmarki\checkmarkr\checkmarkc\checkmark$\checkmarks\checkmark$\checkmark^\checkmark\\checkmarkc\checkmarki\checkmarkr\checkmarkc\checkmark$\checkmarki\checkmark$\checkmark^\checkmark\\checkmarkc\checkmarki\checkmarkr\checkmarkc\checkmark$\checkmarko\checkmark$\checkmark^\checkmark\\checkmarkc\checkmarki\checkmarkr\checkmarkc\checkmark$\checkmarkn\checkmark$\checkmark^\checkmark\\checkmarkc\checkmarki\checkmarkr\checkmarkc\checkmark$\checkmark \checkmark$\checkmark^\checkmark\\checkmarkc\checkmarki\checkmarkr\checkmarkc\checkmark$\checkmarki\checkmark$\checkmark^\checkmark\\checkmarkc\checkmarki\checkmarkr\checkmarkc\checkmark$\checkmarkn\checkmark$\checkmark^\checkmark\\checkmarkc\checkmarki\checkmarkr\checkmarkc\checkmark$\checkmarkt\checkmark$\checkmark^\checkmark\\checkmarkc\checkmarki\checkmarkr\checkmarkc\checkmark$\checkmarke\checkmark$\checkmark^\checkmark\\checkmarkc\checkmarki\checkmarkr\checkmarkc\checkmark$\checkmarkr\checkmark$\checkmark^\checkmark\\checkmarkc\checkmarki\checkmarkr\checkmarkc\checkmark$\checkmarkn\checkmark$\checkmark^\checkmark\\checkmarkc\checkmarki\checkmarkr\checkmarkc\checkmark$\checkmarke\checkmark$\checkmark^\checkmark\\checkmarkc\checkmarki\checkmarkr\checkmarkc\checkmark$\checkmark \checkmark$\checkmark^\checkmark\\checkmarkc\checkmarki\checkmarkr\checkmarkc\checkmark$\checkmarkd\checkmark$\checkmark^\checkmark\\checkmarkc\checkmarki\checkmarkr\checkmarkc\checkmark$\checkmarke\checkmark$\checkmark^\checkmark\\checkmarkc\checkmarki\checkmarkr\checkmarkc\checkmark$\checkmark \checkmark$\checkmark^\checkmark\\checkmarkc\checkmarki\checkmarkr\checkmarkc\checkmark$\checkmarkl\checkmark$\checkmark^\checkmark\\checkmarkc\checkmarki\checkmarkr\checkmarkc\checkmark$\checkmark'\checkmark$\checkmark^\checkmark\\checkmarkc\checkmarki\checkmarkr\checkmarkc\checkmark$\checkmarké\checkmark$\checkmark^\checkmark\\checkmarkc\checkmarki\checkmarkr\checkmarkc\checkmark$\checkmarkl\checkmark$\checkmark^\checkmark\\checkmarkc\checkmarki\checkmarkr\checkmarkc\checkmark$\checkmarke\checkmark$\checkmark^\checkmark\\checkmarkc\checkmarki\checkmarkr\checkmarkc\checkmark$\checkmarkc\checkmark$\checkmark^\checkmark\\checkmarkc\checkmarki\checkmarkr\checkmarkc\checkmark$\checkmarkt\checkmark$\checkmark^\checkmark\\checkmarkc\checkmarki\checkmarkr\checkmarkc\checkmark$\checkmarkr\checkmark$\checkmark^\checkmark\\checkmarkc\checkmarki\checkmarkr\checkmarkc\checkmark$\checkmarko\checkmark$\checkmark^\checkmark\\checkmarkc\checkmarki\checkmarkr\checkmarkc\checkmark$\checkmarkb\checkmark$\checkmark^\checkmark\\checkmarkc\checkmarki\checkmarkr\checkmarkc\checkmark$\checkmarkr\checkmark$\checkmark^\checkmark\\checkmarkc\checkmarki\checkmarkr\checkmarkc\checkmark$\checkmarko\checkmark$\checkmark^\checkmark\\checkmarkc\checkmarki\checkmarkr\checkmarkc\checkmark$\checkmarkc\checkmark$\checkmark^\checkmark\\checkmarkc\checkmarki\checkmarkr\checkmarkc\checkmark$\checkmarkh\checkmark$\checkmark^\checkmark\\checkmarkc\checkmarki\checkmarkr\checkmarkc\checkmark$\checkmarke\checkmark$\checkmark^\checkmark\\checkmarkc\checkmarki\checkmarkr\checkmarkc\checkmark$\checkmark \checkmark$\checkmark^\checkmark\\checkmarkc\checkmarki\checkmarkr\checkmarkc\checkmark$\checkmark$\checkmark$\checkmark^\checkmark\\checkmarkc\checkmarki\checkmarkr\checkmarkc\checkmark$\checkmark\\checkmark$\checkmark^\checkmark\\checkmarkc\checkmarki\checkmarkr\checkmarkc\checkmark$\checkmarke\checkmark$\checkmark^\checkmark\\checkmarkc\checkmarki\checkmarkr\checkmarkc\checkmark$\checkmarkt\checkmark$\checkmark^\checkmark\\checkmarkc\checkmarki\checkmarkr\checkmarkc\checkmark$\checkmarka\checkmark$\checkmark^\checkmark\\checkmarkc\checkmarki\checkmarkr\checkmarkc\checkmark$\checkmark_\checkmark$\checkmark^\checkmark\\checkmarkc\checkmarki\checkmarkr\checkmarkc\checkmark$\checkmark{\checkmark$\checkmark^\checkmark\\checkmarkc\checkmarki\checkmarkr\checkmarkc\checkmark$\checkmarkb\checkmark$\checkmark^\checkmark\\checkmarkc\checkmarki\checkmarkr\checkmarkc\checkmark$\checkmarkr\checkmark$\checkmark^\checkmark\\checkmarkc\checkmarki\checkmarkr\checkmarkc\checkmark$\checkmarko\checkmark$\checkmark^\checkmark\\checkmarkc\checkmarki\checkmarkr\checkmarkc\checkmark$\checkmarkc\checkmark$\checkmark^\checkmark\\checkmarkc\checkmarki\checkmarkr\checkmarkc\checkmark$\checkmarkh\checkmark$\checkmark^\checkmark\\checkmarkc\checkmarki\checkmarkr\checkmarkc\checkmark$\checkmarke\checkmark$\checkmark^\checkmark\\checkmarkc\checkmarki\checkmarkr\checkmarkc\checkmark$\checkmark}\checkmark$\checkmark^\checkmark\\checkmarkc\checkmarki\checkmarkr\checkmarkc\checkmark$\checkmark \checkmark$\checkmark^\checkmark\\checkmarkc\checkmarki\checkmarkr\checkmarkc\checkmark$\checkmark=\checkmark$\checkmark^\checkmark\\checkmarkc\checkmarki\checkmarkr\checkmarkc\checkmark$\checkmark \checkmark$\checkmark^\checkmark\\checkmarkc\checkmarki\checkmarkr\checkmarkc\checkmark$\checkmark0\checkmark$\checkmark^\checkmark\\checkmarkc\checkmarki\checkmarkr\checkmarkc\checkmark$\checkmark{\checkmark$\checkmark^\checkmark\\checkmarkc\checkmarki\checkmarkr\checkmarkc\checkmark$\checkmark,\checkmark$\checkmark^\checkmark\\checkmarkc\checkmarki\checkmarkr\checkmarkc\checkmark$\checkmark}\checkmark$\checkmark^\checkmark\\checkmarkc\checkmarki\checkmarkr\checkmarkc\checkmark$\checkmark8\checkmark$\checkmark^\checkmark\\checkmarkc\checkmarki\checkmarkr\checkmarkc\checkmark$\checkmark0\checkmark$\checkmark^\checkmark\\checkmarkc\checkmarki\checkmarkr\checkmarkc\checkmark$\checkmark$\checkmark$\checkmark^\checkmark\\checkmarkc\checkmarki\checkmarkr\checkmarkc\checkmark$\checkmark \checkmark$\checkmark^\checkmark\\checkmarkc\checkmarki\checkmarkr\checkmarkc\checkmark$\checkmark:\checkmark$\checkmark^\checkmark\\checkmarkc\checkmarki\checkmarkr\checkmarkc\checkmark$\checkmark
\checkmark$\checkmark^\checkmark\\checkmarkc\checkmarki\checkmarkr\checkmarkc\checkmark$\checkmark\\checkmark$\checkmark^\checkmark\\checkmarkc\checkmarki\checkmarkr\checkmarkc\checkmark$\checkmarkb\checkmark$\checkmark^\checkmark\\checkmarkc\checkmarki\checkmarkr\checkmarkc\checkmark$\checkmarke\checkmark$\checkmark^\checkmark\\checkmarkc\checkmarki\checkmarkr\checkmarkc\checkmark$\checkmarkg\checkmark$\checkmark^\checkmark\\checkmarkc\checkmarki\checkmarkr\checkmarkc\checkmark$\checkmarki\checkmark$\checkmark^\checkmark\\checkmarkc\checkmarki\checkmarkr\checkmarkc\checkmark$\checkmarkn\checkmark$\checkmark^\checkmark\\checkmarkc\checkmarki\checkmarkr\checkmarkc\checkmark$\checkmark{\checkmark$\checkmark^\checkmark\\checkmarkc\checkmarki\checkmarkr\checkmarkc\checkmark$\checkmarke\checkmark$\checkmark^\checkmark\\checkmarkc\checkmarki\checkmarkr\checkmarkc\checkmark$\checkmarkq\checkmark$\checkmark^\checkmark\\checkmarkc\checkmarki\checkmarkr\checkmarkc\checkmark$\checkmarku\checkmark$\checkmark^\checkmark\\checkmarkc\checkmarki\checkmarkr\checkmarkc\checkmark$\checkmarka\checkmark$\checkmark^\checkmark\\checkmarkc\checkmarki\checkmarkr\checkmarkc\checkmark$\checkmarkt\checkmark$\checkmark^\checkmark\\checkmarkc\checkmarki\checkmarkr\checkmarkc\checkmark$\checkmarki\checkmark$\checkmark^\checkmark\\checkmarkc\checkmarki\checkmarkr\checkmarkc\checkmark$\checkmarko\checkmark$\checkmark^\checkmark\\checkmarkc\checkmarki\checkmarkr\checkmarkc\checkmark$\checkmarkn\checkmark$\checkmark^\checkmark\\checkmarkc\checkmarki\checkmarkr\checkmarkc\checkmark$\checkmark}\checkmark$\checkmark^\checkmark\\checkmarkc\checkmarki\checkmarkr\checkmarkc\checkmark$\checkmark
\checkmark$\checkmark^\checkmark\\checkmarkc\checkmarki\checkmarkr\checkmarkc\checkmark$\checkmark \checkmark$\checkmark^\checkmark\\checkmarkc\checkmarki\checkmarkr\checkmarkc\checkmark$\checkmark \checkmark$\checkmark^\checkmark\\checkmarkc\checkmarki\checkmarkr\checkmarkc\checkmark$\checkmark \checkmark$\checkmark^\checkmark\\checkmarkc\checkmarki\checkmarkr\checkmarkc\checkmark$\checkmark \checkmark$\checkmark^\checkmark\\checkmarkc\checkmarki\checkmarkr\checkmarkc\checkmark$\checkmarkP\checkmark$\checkmark^\checkmark\\checkmarkc\checkmarki\checkmarkr\checkmarkc\checkmark$\checkmark_\checkmark$\checkmark^\checkmark\\checkmarkc\checkmarki\checkmarkr\checkmarkc\checkmark$\checkmark{\checkmark$\checkmark^\checkmark\\checkmarkc\checkmarki\checkmarkr\checkmarkc\checkmark$\checkmarkm\checkmark$\checkmark^\checkmark\\checkmarkc\checkmarki\checkmarkr\checkmarkc\checkmark$\checkmarko\checkmark$\checkmark^\checkmark\\checkmarkc\checkmarki\checkmarkr\checkmarkc\checkmark$\checkmarkt\checkmark$\checkmark^\checkmark\\checkmarkc\checkmarki\checkmarkr\checkmarkc\checkmark$\checkmarke\checkmark$\checkmark^\checkmark\\checkmarkc\checkmarki\checkmarkr\checkmarkc\checkmark$\checkmarku\checkmark$\checkmark^\checkmark\\checkmarkc\checkmarki\checkmarkr\checkmarkc\checkmark$\checkmarkr\checkmark$\checkmark^\checkmark\\checkmarkc\checkmarki\checkmarkr\checkmarkc\checkmark$\checkmark}\checkmark$\checkmark^\checkmark\\checkmarkc\checkmarki\checkmarkr\checkmarkc\checkmark$\checkmark \checkmark$\checkmark^\checkmark\\checkmarkc\checkmarki\checkmarkr\checkmarkc\checkmark$\checkmark=\checkmark$\checkmark^\checkmark\\checkmarkc\checkmarki\checkmarkr\checkmarkc\checkmark$\checkmark \checkmark$\checkmark^\checkmark\\checkmarkc\checkmarki\checkmarkr\checkmarkc\checkmark$\checkmark\\checkmark$\checkmark^\checkmark\\checkmarkc\checkmarki\checkmarkr\checkmarkc\checkmark$\checkmarkf\checkmark$\checkmark^\checkmark\\checkmarkc\checkmarki\checkmarkr\checkmarkc\checkmark$\checkmarkr\checkmark$\checkmark^\checkmark\\checkmarkc\checkmarki\checkmarkr\checkmarkc\checkmark$\checkmarka\checkmark$\checkmark^\checkmark\\checkmarkc\checkmarki\checkmarkr\checkmarkc\checkmark$\checkmarkc\checkmark$\checkmark^\checkmark\\checkmarkc\checkmarki\checkmarkr\checkmarkc\checkmark$\checkmark{\checkmark$\checkmark^\checkmark\\checkmarkc\checkmarki\checkmarkr\checkmarkc\checkmark$\checkmarkP\checkmark$\checkmark^\checkmark\\checkmarkc\checkmarki\checkmarkr\checkmarkc\checkmark$\checkmark_\checkmark$\checkmark^\checkmark\\checkmarkc\checkmarki\checkmarkr\checkmarkc\checkmark$\checkmarkc\checkmark$\checkmark^\checkmark\\checkmarkc\checkmarki\checkmarkr\checkmarkc\checkmark$\checkmark}\checkmark$\checkmark^\checkmark\\checkmarkc\checkmarki\checkmarkr\checkmarkc\checkmark$\checkmark{\checkmark$\checkmark^\checkmark\\checkmarkc\checkmarki\checkmarkr\checkmarkc\checkmark$\checkmark\\checkmark$\checkmark^\checkmark\\checkmarkc\checkmarki\checkmarkr\checkmarkc\checkmark$\checkmarke\checkmark$\checkmark^\checkmark\\checkmarkc\checkmarki\checkmarkr\checkmarkc\checkmark$\checkmarkt\checkmark$\checkmark^\checkmark\\checkmarkc\checkmarki\checkmarkr\checkmarkc\checkmark$\checkmarka\checkmark$\checkmark^\checkmark\\checkmarkc\checkmarki\checkmarkr\checkmarkc\checkmark$\checkmark_\checkmark$\checkmark^\checkmark\\checkmarkc\checkmarki\checkmarkr\checkmarkc\checkmark$\checkmark{\checkmark$\checkmark^\checkmark\\checkmarkc\checkmarki\checkmarkr\checkmarkc\checkmark$\checkmarkb\checkmark$\checkmark^\checkmark\\checkmarkc\checkmarki\checkmarkr\checkmarkc\checkmark$\checkmarkr\checkmark$\checkmark^\checkmark\\checkmarkc\checkmarki\checkmarkr\checkmarkc\checkmark$\checkmarko\checkmark$\checkmark^\checkmark\\checkmarkc\checkmarki\checkmarkr\checkmarkc\checkmark$\checkmarkc\checkmark$\checkmark^\checkmark\\checkmarkc\checkmarki\checkmarkr\checkmarkc\checkmark$\checkmarkh\checkmark$\checkmark^\checkmark\\checkmarkc\checkmarki\checkmarkr\checkmarkc\checkmark$\checkmarke\checkmark$\checkmark^\checkmark\\checkmarkc\checkmarki\checkmarkr\checkmarkc\checkmark$\checkmark}\checkmark$\checkmark^\checkmark\\checkmarkc\checkmarki\checkmarkr\checkmarkc\checkmark$\checkmark}\checkmark$\checkmark^\checkmark\\checkmarkc\checkmarki\checkmarkr\checkmarkc\checkmark$\checkmark \checkmark$\checkmark^\checkmark\\checkmarkc\checkmarki\checkmarkr\checkmarkc\checkmark$\checkmark=\checkmark$\checkmark^\checkmark\\checkmarkc\checkmarki\checkmarkr\checkmarkc\checkmark$\checkmark \checkmark$\checkmark^\checkmark\\checkmarkc\checkmarki\checkmarkr\checkmarkc\checkmark$\checkmark\\checkmark$\checkmark^\checkmark\\checkmarkc\checkmarki\checkmarkr\checkmarkc\checkmark$\checkmarkf\checkmark$\checkmark^\checkmark\\checkmarkc\checkmarki\checkmarkr\checkmarkc\checkmark$\checkmarkr\checkmark$\checkmark^\checkmark\\checkmarkc\checkmarki\checkmarkr\checkmarkc\checkmark$\checkmarka\checkmark$\checkmark^\checkmark\\checkmarkc\checkmarki\checkmarkr\checkmarkc\checkmark$\checkmarkc\checkmark$\checkmark^\checkmark\\checkmarkc\checkmarki\checkmarkr\checkmarkc\checkmark$\checkmark{\checkmark$\checkmark^\checkmark\\checkmarkc\checkmarki\checkmarkr\checkmarkc\checkmark$\checkmark1\checkmark$\checkmark^\checkmark\\checkmarkc\checkmarki\checkmarkr\checkmarkc\checkmark$\checkmark7\checkmark$\checkmark^\checkmark\\checkmarkc\checkmarki\checkmarkr\checkmarkc\checkmark$\checkmark5\checkmark$\checkmark^\checkmark\\checkmarkc\checkmarki\checkmarkr\checkmarkc\checkmark$\checkmark}\checkmark$\checkmark^\checkmark\\checkmarkc\checkmarki\checkmarkr\checkmarkc\checkmark$\checkmark{\checkmark$\checkmark^\checkmark\\checkmarkc\checkmarki\checkmarkr\checkmarkc\checkmark$\checkmark0\checkmark$\checkmark^\checkmark\\checkmarkc\checkmarki\checkmarkr\checkmarkc\checkmark$\checkmark{\checkmark$\checkmark^\checkmark\\checkmarkc\checkmarki\checkmarkr\checkmarkc\checkmark$\checkmark,\checkmark$\checkmark^\checkmark\\checkmarkc\checkmarki\checkmarkr\checkmarkc\checkmark$\checkmark}\checkmark$\checkmark^\checkmark\\checkmarkc\checkmarki\checkmarkr\checkmarkc\checkmark$\checkmark8\checkmark$\checkmark^\checkmark\\checkmarkc\checkmarki\checkmarkr\checkmarkc\checkmark$\checkmark0\checkmark$\checkmark^\checkmark\\checkmarkc\checkmarki\checkmarkr\checkmarkc\checkmark$\checkmark}\checkmark$\checkmark^\checkmark\\checkmarkc\checkmarki\checkmarkr\checkmarkc\checkmark$\checkmark \checkmark$\checkmark^\checkmark\\checkmarkc\checkmarki\checkmarkr\checkmarkc\checkmark$\checkmark\\checkmark$\checkmark^\checkmark\\checkmarkc\checkmarki\checkmarkr\checkmarkc\checkmark$\checkmarka\checkmark$\checkmark^\checkmark\\checkmarkc\checkmarki\checkmarkr\checkmarkc\checkmark$\checkmarkp\checkmark$\checkmark^\checkmark\\checkmarkc\checkmarki\checkmarkr\checkmarkc\checkmark$\checkmarkp\checkmark$\checkmark^\checkmark\\checkmarkc\checkmarki\checkmarkr\checkmarkc\checkmark$\checkmarkr\checkmark$\checkmark^\checkmark\\checkmarkc\checkmarki\checkmarkr\checkmarkc\checkmark$\checkmarko\checkmark$\checkmark^\checkmark\\checkmarkc\checkmarki\checkmarkr\checkmarkc\checkmark$\checkmarkx\checkmark$\checkmark^\checkmark\\checkmarkc\checkmarki\checkmarkr\checkmarkc\checkmark$\checkmark \checkmark$\checkmark^\checkmark\\checkmarkc\checkmarki\checkmarkr\checkmarkc\checkmark$\checkmark\\checkmark$\checkmark^\checkmark\\checkmarkc\checkmarki\checkmarkr\checkmarkc\checkmark$\checkmarkb\checkmark$\checkmark^\checkmark\\checkmarkc\checkmarki\checkmarkr\checkmarkc\checkmark$\checkmarko\checkmark$\checkmark^\checkmark\\checkmarkc\checkmarki\checkmarkr\checkmarkc\checkmark$\checkmarkx\checkmark$\checkmark^\checkmark\\checkmarkc\checkmarki\checkmarkr\checkmarkc\checkmark$\checkmarke\checkmark$\checkmark^\checkmark\\checkmarkc\checkmarki\checkmarkr\checkmarkc\checkmark$\checkmarkd\checkmark$\checkmark^\checkmark\\checkmarkc\checkmarki\checkmarkr\checkmarkc\checkmark$\checkmark{\checkmark$\checkmark^\checkmark\\checkmarkc\checkmarki\checkmarkr\checkmarkc\checkmark$\checkmark2\checkmark$\checkmark^\checkmark\\checkmarkc\checkmarki\checkmarkr\checkmarkc\checkmark$\checkmark1\checkmark$\checkmark^\checkmark\\checkmarkc\checkmarki\checkmarkr\checkmarkc\checkmark$\checkmark9\checkmark$\checkmark^\checkmark\\checkmarkc\checkmarki\checkmarkr\checkmarkc\checkmark$\checkmark \checkmark$\checkmark^\checkmark\\checkmarkc\checkmarki\checkmarkr\checkmarkc\checkmark$\checkmark\\checkmark$\checkmark^\checkmark\\checkmarkc\checkmarki\checkmarkr\checkmarkc\checkmark$\checkmarkt\checkmark$\checkmark^\checkmark\\checkmarkc\checkmarki\checkmarkr\checkmarkc\checkmark$\checkmarke\checkmark$\checkmark^\checkmark\\checkmarkc\checkmarki\checkmarkr\checkmarkc\checkmark$\checkmarkx\checkmark$\checkmark^\checkmark\\checkmarkc\checkmarki\checkmarkr\checkmarkc\checkmark$\checkmarkt\checkmark$\checkmark^\checkmark\\checkmarkc\checkmarki\checkmarkr\checkmarkc\checkmark$\checkmark{\checkmark$\checkmark^\checkmark\\checkmarkc\checkmarki\checkmarkr\checkmarkc\checkmark$\checkmark \checkmark$\checkmark^\checkmark\\checkmarkc\checkmarki\checkmarkr\checkmarkc\checkmark$\checkmarkW\checkmark$\checkmark^\checkmark\\checkmarkc\checkmarki\checkmarkr\checkmarkc\checkmark$\checkmark}\checkmark$\checkmark^\checkmark\\checkmarkc\checkmarki\checkmarkr\checkmarkc\checkmark$\checkmark}\checkmark$\checkmark^\checkmark\\checkmarkc\checkmarki\checkmarkr\checkmarkc\checkmark$\checkmark
\checkmark$\checkmark^\checkmark\\checkmarkc\checkmarki\checkmarkr\checkmarkc\checkmark$\checkmark\\checkmark$\checkmark^\checkmark\\checkmarkc\checkmarki\checkmarkr\checkmarkc\checkmark$\checkmarke\checkmark$\checkmark^\checkmark\\checkmarkc\checkmarki\checkmarkr\checkmarkc\checkmark$\checkmarkn\checkmark$\checkmark^\checkmark\\checkmarkc\checkmarki\checkmarkr\checkmarkc\checkmark$\checkmarkd\checkmark$\checkmark^\checkmark\\checkmarkc\checkmarki\checkmarkr\checkmarkc\checkmark$\checkmark{\checkmark$\checkmark^\checkmark\\checkmarkc\checkmarki\checkmarkr\checkmarkc\checkmark$\checkmarke\checkmark$\checkmark^\checkmark\\checkmarkc\checkmarki\checkmarkr\checkmarkc\checkmark$\checkmarkq\checkmark$\checkmark^\checkmark\\checkmarkc\checkmarki\checkmarkr\checkmarkc\checkmark$\checkmarku\checkmark$\checkmark^\checkmark\\checkmarkc\checkmarki\checkmarkr\checkmarkc\checkmark$\checkmarka\checkmark$\checkmark^\checkmark\\checkmarkc\checkmarki\checkmarkr\checkmarkc\checkmark$\checkmarkt\checkmark$\checkmark^\checkmark\\checkmarkc\checkmarki\checkmarkr\checkmarkc\checkmark$\checkmarki\checkmark$\checkmark^\checkmark\\checkmarkc\checkmarki\checkmarkr\checkmarkc\checkmark$\checkmarko\checkmark$\checkmark^\checkmark\\checkmarkc\checkmarki\checkmarkr\checkmarkc\checkmark$\checkmarkn\checkmark$\checkmark^\checkmark\\checkmarkc\checkmarki\checkmarkr\checkmarkc\checkmark$\checkmark}\checkmark$\checkmark^\checkmark\\checkmarkc\checkmarki\checkmarkr\checkmarkc\checkmark$\checkmark
\checkmark$\checkmark^\checkmark\\checkmarkc\checkmarki\checkmarkr\checkmarkc\checkmark$\checkmark
\checkmark$\checkmark^\checkmark\\checkmarkc\checkmarki\checkmarkr\checkmarkc\checkmark$\checkmark\\checkmark$\checkmark^\checkmark\\checkmarkc\checkmarki\checkmarkr\checkmarkc\checkmark$\checkmarkt\checkmark$\checkmark^\checkmark\\checkmarkc\checkmarki\checkmarkr\checkmarkc\checkmark$\checkmarke\checkmark$\checkmark^\checkmark\\checkmarkc\checkmarki\checkmarkr\checkmarkc\checkmark$\checkmarkx\checkmark$\checkmark^\checkmark\\checkmarkc\checkmarki\checkmarkr\checkmarkc\checkmark$\checkmarkt\checkmark$\checkmark^\checkmark\\checkmarkc\checkmarki\checkmarkr\checkmarkc\checkmark$\checkmarkb\checkmark$\checkmark^\checkmark\\checkmarkc\checkmarki\checkmarkr\checkmarkc\checkmark$\checkmarkf\checkmark$\checkmark^\checkmark\\checkmarkc\checkmarki\checkmarkr\checkmarkc\checkmark$\checkmark{\checkmark$\checkmark^\checkmark\\checkmarkc\checkmarki\checkmarkr\checkmarkc\checkmark$\checkmarkC\checkmark$\checkmark^\checkmark\\checkmarkc\checkmarki\checkmarkr\checkmarkc\checkmark$\checkmarko\checkmark$\checkmark^\checkmark\\checkmarkc\checkmarki\checkmarkr\checkmarkc\checkmark$\checkmarkn\checkmark$\checkmark^\checkmark\\checkmarkc\checkmarki\checkmarkr\checkmarkc\checkmark$\checkmarkc\checkmark$\checkmark^\checkmark\\checkmarkc\checkmarki\checkmarkr\checkmarkc\checkmark$\checkmarkl\checkmark$\checkmark^\checkmark\\checkmarkc\checkmarki\checkmarkr\checkmarkc\checkmark$\checkmarku\checkmark$\checkmark^\checkmark\\checkmarkc\checkmarki\checkmarkr\checkmarkc\checkmark$\checkmarks\checkmark$\checkmark^\checkmark\\checkmarkc\checkmarki\checkmarkr\checkmarkc\checkmark$\checkmarki\checkmark$\checkmark^\checkmark\\checkmarkc\checkmarki\checkmarkr\checkmarkc\checkmark$\checkmarko\checkmark$\checkmark^\checkmark\\checkmarkc\checkmarki\checkmarkr\checkmarkc\checkmark$\checkmarkn\checkmark$\checkmark^\checkmark\\checkmarkc\checkmarki\checkmarkr\checkmarkc\checkmark$\checkmark \checkmark$\checkmark^\checkmark\\checkmarkc\checkmarki\checkmarkr\checkmarkc\checkmark$\checkmark:\checkmark$\checkmark^\checkmark\\checkmarkc\checkmarki\checkmarkr\checkmarkc\checkmark$\checkmark}\checkmark$\checkmark^\checkmark\\checkmarkc\checkmarki\checkmarkr\checkmarkc\checkmark$\checkmark \checkmark$\checkmark^\checkmark\\checkmarkc\checkmarki\checkmarkr\checkmarkc\checkmark$\checkmarkC\checkmark$\checkmark^\checkmark\\checkmarkc\checkmarki\checkmarkr\checkmarkc\checkmark$\checkmarke\checkmark$\checkmark^\checkmark\\checkmarkc\checkmarki\checkmarkr\checkmarkc\checkmark$\checkmark \checkmark$\checkmark^\checkmark\\checkmarkc\checkmarki\checkmarkr\checkmarkc\checkmark$\checkmarkd\checkmark$\checkmark^\checkmark\\checkmarkc\checkmarki\checkmarkr\checkmarkc\checkmark$\checkmarki\checkmark$\checkmark^\checkmark\\checkmarkc\checkmarki\checkmarkr\checkmarkc\checkmark$\checkmarkm\checkmark$\checkmark^\checkmark\\checkmarkc\checkmarki\checkmarkr\checkmarkc\checkmark$\checkmarke\checkmark$\checkmark^\checkmark\\checkmarkc\checkmarki\checkmarkr\checkmarkc\checkmark$\checkmarkn\checkmark$\checkmark^\checkmark\\checkmarkc\checkmarki\checkmarkr\checkmarkc\checkmark$\checkmarks\checkmark$\checkmark^\checkmark\\checkmarkc\checkmarki\checkmarkr\checkmarkc\checkmark$\checkmarki\checkmark$\checkmark^\checkmark\\checkmarkc\checkmarki\checkmarkr\checkmarkc\checkmark$\checkmarko\checkmark$\checkmark^\checkmark\\checkmarkc\checkmarki\checkmarkr\checkmarkc\checkmark$\checkmarkn\checkmark$\checkmark^\checkmark\\checkmarkc\checkmarki\checkmarkr\checkmarkc\checkmark$\checkmarkn\checkmark$\checkmark^\checkmark\\checkmarkc\checkmarki\checkmarkr\checkmarkc\checkmark$\checkmarke\checkmark$\checkmark^\checkmark\\checkmarkc\checkmarki\checkmarkr\checkmarkc\checkmark$\checkmarkm\checkmark$\checkmark^\checkmark\\checkmarkc\checkmarki\checkmarkr\checkmarkc\checkmark$\checkmarke\checkmark$\checkmark^\checkmark\\checkmarkc\checkmarki\checkmarkr\checkmarkc\checkmark$\checkmarkn\checkmark$\checkmark^\checkmark\\checkmarkc\checkmarki\checkmarkr\checkmarkc\checkmark$\checkmarkt\checkmark$\checkmark^\checkmark\\checkmarkc\checkmarki\checkmarkr\checkmarkc\checkmark$\checkmark \checkmark$\checkmark^\checkmark\\checkmarkc\checkmarki\checkmarkr\checkmarkc\checkmark$\checkmarkj\checkmark$\checkmark^\checkmark\\checkmarkc\checkmarki\checkmarkr\checkmarkc\checkmark$\checkmarku\checkmark$\checkmark^\checkmark\\checkmarkc\checkmarki\checkmarkr\checkmarkc\checkmark$\checkmarks\checkmark$\checkmark^\checkmark\\checkmarkc\checkmarki\checkmarkr\checkmarkc\checkmark$\checkmarkt\checkmark$\checkmark^\checkmark\\checkmarkc\checkmarki\checkmarkr\checkmarkc\checkmark$\checkmarki\checkmark$\checkmark^\checkmark\\checkmarkc\checkmarki\checkmarkr\checkmarkc\checkmark$\checkmarkf\checkmark$\checkmark^\checkmark\\checkmarkc\checkmarki\checkmarkr\checkmarkc\checkmark$\checkmarki\checkmark$\checkmark^\checkmark\\checkmarkc\checkmarki\checkmarkr\checkmarkc\checkmark$\checkmarke\checkmark$\checkmark^\checkmark\\checkmarkc\checkmarki\checkmarkr\checkmarkc\checkmark$\checkmark \checkmark$\checkmark^\checkmark\\checkmarkc\checkmarki\checkmarkr\checkmarkc\checkmark$\checkmarkl\checkmark$\checkmark^\checkmark\\checkmarkc\checkmarki\checkmarkr\checkmarkc\checkmark$\checkmarke\checkmark$\checkmark^\checkmark\\checkmarkc\checkmarki\checkmarkr\checkmarkc\checkmark$\checkmark \checkmark$\checkmark^\checkmark\\checkmarkc\checkmarki\checkmarkr\checkmarkc\checkmark$\checkmarkc\checkmark$\checkmark^\checkmark\\checkmarkc\checkmarki\checkmarkr\checkmarkc\checkmark$\checkmarkh\checkmark$\checkmark^\checkmark\\checkmarkc\checkmarki\checkmarkr\checkmarkc\checkmark$\checkmarko\checkmark$\checkmark^\checkmark\\checkmarkc\checkmarki\checkmarkr\checkmarkc\checkmark$\checkmarki\checkmark$\checkmark^\checkmark\\checkmarkc\checkmarki\checkmarkr\checkmarkc\checkmark$\checkmarkx\checkmark$\checkmark^\checkmark\\checkmarkc\checkmarki\checkmarkr\checkmarkc\checkmark$\checkmark \checkmark$\checkmark^\checkmark\\checkmarkc\checkmarki\checkmarkr\checkmarkc\checkmark$\checkmarkd\checkmark$\checkmark^\checkmark\\checkmarkc\checkmarki\checkmarkr\checkmarkc\checkmark$\checkmark'\checkmark$\checkmark^\checkmark\\checkmarkc\checkmarki\checkmarkr\checkmarkc\checkmark$\checkmarku\checkmark$\checkmark^\checkmark\\checkmarkc\checkmarki\checkmarkr\checkmarkc\checkmark$\checkmarkn\checkmark$\checkmark^\checkmark\\checkmarkc\checkmarki\checkmarkr\checkmarkc\checkmark$\checkmarke\checkmark$\checkmark^\checkmark\\checkmarkc\checkmarki\checkmarkr\checkmarkc\checkmark$\checkmark \checkmark$\checkmark^\checkmark\\checkmarkc\checkmarki\checkmarkr\checkmarkc\checkmark$\checkmarké\checkmark$\checkmark^\checkmark\\checkmarkc\checkmarki\checkmarkr\checkmarkc\checkmark$\checkmarkl\checkmark$\checkmark^\checkmark\\checkmarkc\checkmarki\checkmarkr\checkmarkc\checkmark$\checkmarke\checkmark$\checkmark^\checkmark\\checkmarkc\checkmarki\checkmarkr\checkmarkc\checkmark$\checkmarkc\checkmark$\checkmark^\checkmark\\checkmarkc\checkmarki\checkmarkr\checkmarkc\checkmark$\checkmarkt\checkmark$\checkmark^\checkmark\\checkmarkc\checkmarki\checkmarkr\checkmarkc\checkmark$\checkmarkr\checkmark$\checkmark^\checkmark\\checkmarkc\checkmarki\checkmarkr\checkmarkc\checkmark$\checkmarko\checkmark$\checkmark^\checkmark\\checkmarkc\checkmarki\checkmarkr\checkmarkc\checkmark$\checkmarkb\checkmark$\checkmark^\checkmark\\checkmarkc\checkmarki\checkmarkr\checkmarkc\checkmark$\checkmarkr\checkmark$\checkmark^\checkmark\\checkmarkc\checkmarki\checkmarkr\checkmarkc\checkmark$\checkmarko\checkmark$\checkmark^\checkmark\\checkmarkc\checkmarki\checkmarkr\checkmarkc\checkmark$\checkmarkc\checkmark$\checkmark^\checkmark\\checkmarkc\checkmarki\checkmarkr\checkmarkc\checkmark$\checkmarkh\checkmark$\checkmark^\checkmark\\checkmarkc\checkmarki\checkmarkr\checkmarkc\checkmark$\checkmarke\checkmark$\checkmark^\checkmark\\checkmarkc\checkmarki\checkmarkr\checkmarkc\checkmark$\checkmark \checkmark$\checkmark^\checkmark\\checkmarkc\checkmarki\checkmarkr\checkmarkc\checkmark$\checkmarkB\checkmark$\checkmark^\checkmark\\checkmarkc\checkmarki\checkmarkr\checkmarkc\checkmark$\checkmarkr\checkmark$\checkmark^\checkmark\\checkmarkc\checkmarki\checkmarkr\checkmarkc\checkmark$\checkmarku\checkmark$\checkmark^\checkmark\\checkmarkc\checkmarki\checkmarkr\checkmarkc\checkmark$\checkmarks\checkmark$\checkmark^\checkmark\\checkmarkc\checkmarki\checkmarkr\checkmarkc\checkmark$\checkmarkh\checkmark$\checkmark^\checkmark\\checkmarkc\checkmarki\checkmarkr\checkmarkc\checkmark$\checkmarkl\checkmark$\checkmark^\checkmark\\checkmarkc\checkmarki\checkmarkr\checkmarkc\checkmark$\checkmarke\checkmark$\checkmark^\checkmark\\checkmarkc\checkmarki\checkmarkr\checkmarkc\checkmark$\checkmarks\checkmark$\checkmark^\checkmark\\checkmarkc\checkmarki\checkmarkr\checkmarkc\checkmark$\checkmarks\checkmark$\checkmark^\checkmark\\checkmarkc\checkmarki\checkmarkr\checkmarkc\checkmark$\checkmark \checkmark$\checkmark^\checkmark\\checkmarkc\checkmarki\checkmarkr\checkmarkc\checkmark$\checkmarkD\checkmark$\checkmark^\checkmark\\checkmarkc\checkmarki\checkmarkr\checkmarkc\checkmark$\checkmarkC\checkmark$\checkmark^\checkmark\\checkmarkc\checkmarki\checkmarkr\checkmarkc\checkmark$\checkmark \checkmark$\checkmark^\checkmark\\checkmarkc\checkmarki\checkmarkr\checkmarkc\checkmark$\checkmarkd\checkmark$\checkmark^\checkmark\\checkmarkc\checkmarki\checkmarkr\checkmarkc\checkmark$\checkmarke\checkmark$\checkmark^\checkmark\\checkmarkc\checkmarki\checkmarkr\checkmarkc\checkmark$\checkmark \checkmark$\checkmark^\checkmark\\checkmarkc\checkmarki\checkmarkr\checkmarkc\checkmark$\checkmark\\checkmark$\checkmark^\checkmark\\checkmarkc\checkmarki\checkmarkr\checkmarkc\checkmark$\checkmarkt\checkmark$\checkmark^\checkmark\\checkmarkc\checkmarki\checkmarkr\checkmarkc\checkmark$\checkmarke\checkmark$\checkmark^\checkmark\\checkmarkc\checkmarki\checkmarkr\checkmarkc\checkmark$\checkmarkx\checkmark$\checkmark^\checkmark\\checkmarkc\checkmarki\checkmarkr\checkmarkc\checkmark$\checkmarkt\checkmark$\checkmark^\checkmark\\checkmarkc\checkmarki\checkmarkr\checkmarkc\checkmark$\checkmarkb\checkmark$\checkmark^\checkmark\\checkmarkc\checkmarki\checkmarkr\checkmarkc\checkmark$\checkmarkf\checkmark$\checkmark^\checkmark\\checkmarkc\checkmarki\checkmarkr\checkmarkc\checkmark$\checkmark{\checkmark$\checkmark^\checkmark\\checkmarkc\checkmarki\checkmarkr\checkmarkc\checkmark$\checkmark3\checkmark$\checkmark^\checkmark\\checkmarkc\checkmarki\checkmarkr\checkmarkc\checkmark$\checkmark0\checkmark$\checkmark^\checkmark\\checkmarkc\checkmarki\checkmarkr\checkmarkc\checkmark$\checkmark0\checkmark$\checkmark^\checkmark\\checkmarkc\checkmarki\checkmarkr\checkmarkc\checkmark$\checkmark \checkmark$\checkmark^\checkmark\\checkmarkc\checkmarki\checkmarkr\checkmarkc\checkmark$\checkmarkW\checkmark$\checkmark^\checkmark\\checkmarkc\checkmarki\checkmarkr\checkmarkc\checkmark$\checkmark \checkmark$\checkmark^\checkmark\\checkmarkc\checkmarki\checkmarkr\checkmarkc\checkmark$\checkmarkà\checkmark$\checkmark^\checkmark\\checkmarkc\checkmarki\checkmarkr\checkmarkc\checkmark$\checkmark \checkmark$\checkmark^\checkmark\\checkmarkc\checkmarki\checkmarkr\checkmarkc\checkmark$\checkmark5\checkmark$\checkmark^\checkmark\\checkmarkc\checkmarki\checkmarkr\checkmarkc\checkmark$\checkmark0\checkmark$\checkmark^\checkmark\\checkmarkc\checkmarki\checkmarkr\checkmarkc\checkmark$\checkmark0\checkmark$\checkmark^\checkmark\\checkmarkc\checkmarki\checkmarkr\checkmarkc\checkmark$\checkmark \checkmark$\checkmark^\checkmark\\checkmarkc\checkmarki\checkmarkr\checkmarkc\checkmark$\checkmarkW\checkmark$\checkmark^\checkmark\\checkmarkc\checkmarki\checkmarkr\checkmarkc\checkmark$\checkmark}\checkmark$\checkmark^\checkmark\\checkmarkc\checkmarki\checkmarkr\checkmarkc\checkmark$\checkmark,\checkmark$\checkmark^\checkmark\\checkmarkc\checkmarki\checkmarkr\checkmarkc\checkmark$\checkmark \checkmark$\checkmark^\checkmark\\checkmarkc\checkmarki\checkmarkr\checkmarkc\checkmark$\checkmarkq\checkmark$\checkmark^\checkmark\\checkmarkc\checkmarki\checkmarkr\checkmarkc\checkmark$\checkmarku\checkmark$\checkmark^\checkmark\\checkmarkc\checkmarki\checkmarkr\checkmarkc\checkmark$\checkmarki\checkmark$\checkmark^\checkmark\\checkmarkc\checkmarki\checkmarkr\checkmarkc\checkmark$\checkmark \checkmark$\checkmark^\checkmark\\checkmarkc\checkmarki\checkmarkr\checkmarkc\checkmark$\checkmarko\checkmark$\checkmark^\checkmark\\checkmarkc\checkmarki\checkmarkr\checkmarkc\checkmark$\checkmarkf\checkmark$\checkmark^\checkmark\\checkmarkc\checkmarki\checkmarkr\checkmarkc\checkmark$\checkmarkf\checkmark$\checkmark^\checkmark\\checkmarkc\checkmarki\checkmarkr\checkmarkc\checkmark$\checkmarkr\checkmark$\checkmark^\checkmark\\checkmarkc\checkmarki\checkmarkr\checkmarkc\checkmark$\checkmarke\checkmark$\checkmark^\checkmark\\checkmarkc\checkmarki\checkmarkr\checkmarkc\checkmark$\checkmark \checkmark$\checkmark^\checkmark\\checkmarkc\checkmarki\checkmarkr\checkmarkc\checkmark$\checkmarku\checkmark$\checkmark^\checkmark\\checkmarkc\checkmarki\checkmarkr\checkmarkc\checkmark$\checkmarkn\checkmark$\checkmark^\checkmark\\checkmarkc\checkmarki\checkmarkr\checkmarkc\checkmark$\checkmark \checkmark$\checkmark^\checkmark\\checkmarkc\checkmarki\checkmarkr\checkmarkc\checkmark$\checkmarkc\checkmark$\checkmark^\checkmark\\checkmarkc\checkmarki\checkmarkr\checkmarkc\checkmark$\checkmarko\checkmark$\checkmark^\checkmark\\checkmarkc\checkmarki\checkmarkr\checkmarkc\checkmark$\checkmarke\checkmark$\checkmark^\checkmark\\checkmarkc\checkmarki\checkmarkr\checkmarkc\checkmark$\checkmarkf\checkmark$\checkmark^\checkmark\\checkmarkc\checkmarki\checkmarkr\checkmarkc\checkmark$\checkmarkf\checkmark$\checkmark^\checkmark\\checkmarkc\checkmarki\checkmarkr\checkmarkc\checkmark$\checkmarki\checkmark$\checkmark^\checkmark\\checkmarkc\checkmarki\checkmarkr\checkmarkc\checkmark$\checkmarkc\checkmark$\checkmark^\checkmark\\checkmarkc\checkmarki\checkmarkr\checkmarkc\checkmark$\checkmarki\checkmark$\checkmark^\checkmark\\checkmarkc\checkmarki\checkmarkr\checkmarkc\checkmark$\checkmarke\checkmark$\checkmark^\checkmark\\checkmarkc\checkmarki\checkmarkr\checkmarkc\checkmark$\checkmarkn\checkmark$\checkmark^\checkmark\\checkmarkc\checkmarki\checkmarkr\checkmarkc\checkmark$\checkmarkt\checkmark$\checkmark^\checkmark\\checkmarkc\checkmarki\checkmarkr\checkmarkc\checkmark$\checkmark \checkmark$\checkmark^\checkmark\\checkmarkc\checkmarki\checkmarkr\checkmarkc\checkmark$\checkmarkd\checkmark$\checkmark^\checkmark\\checkmarkc\checkmarki\checkmarkr\checkmarkc\checkmark$\checkmarke\checkmark$\checkmark^\checkmark\\checkmarkc\checkmarki\checkmarkr\checkmarkc\checkmark$\checkmark \checkmark$\checkmark^\checkmark\\checkmarkc\checkmarki\checkmarkr\checkmarkc\checkmark$\checkmarks\checkmark$\checkmark^\checkmark\\checkmarkc\checkmarki\checkmarkr\checkmarkc\checkmark$\checkmarké\checkmark$\checkmark^\checkmark\\checkmarkc\checkmarki\checkmarkr\checkmarkc\checkmark$\checkmarkc\checkmark$\checkmark^\checkmark\\checkmarkc\checkmarki\checkmarkr\checkmarkc\checkmark$\checkmarku\checkmark$\checkmark^\checkmark\\checkmarkc\checkmarki\checkmarkr\checkmarkc\checkmark$\checkmarkr\checkmark$\checkmark^\checkmark\\checkmarkc\checkmarki\checkmarkr\checkmarkc\checkmark$\checkmarki\checkmark$\checkmark^\checkmark\\checkmarkc\checkmarki\checkmarkr\checkmarkc\checkmark$\checkmarkt\checkmark$\checkmark^\checkmark\\checkmarkc\checkmarki\checkmarkr\checkmarkc\checkmark$\checkmarké\checkmark$\checkmark^\checkmark\\checkmarkc\checkmarki\checkmarkr\checkmarkc\checkmark$\checkmark \checkmark$\checkmark^\checkmark\\checkmarkc\checkmarki\checkmarkr\checkmarkc\checkmark$\checkmarkd\checkmark$\checkmark^\checkmark\\checkmarkc\checkmarki\checkmarkr\checkmarkc\checkmark$\checkmarke\checkmark$\checkmark^\checkmark\\checkmarkc\checkmarki\checkmarkr\checkmarkc\checkmark$\checkmark \checkmark$\checkmark^\checkmark\\checkmarkc\checkmarki\checkmarkr\checkmarkc\checkmark$\checkmark$\checkmark$\checkmark^\checkmark\\checkmarkc\checkmarki\checkmarkr\checkmarkc\checkmark$\checkmarks\checkmark$\checkmark^\checkmark\\checkmarkc\checkmarki\checkmarkr\checkmarkc\checkmark$\checkmark \checkmark$\checkmark^\checkmark\\checkmarkc\checkmarki\checkmarkr\checkmarkc\checkmark$\checkmark=\checkmark$\checkmark^\checkmark\\checkmarkc\checkmarki\checkmarkr\checkmarkc\checkmark$\checkmark \checkmark$\checkmark^\checkmark\\checkmarkc\checkmarki\checkmarkr\checkmarkc\checkmark$\checkmark3\checkmark$\checkmark^\checkmark\\checkmarkc\checkmarki\checkmarkr\checkmarkc\checkmark$\checkmark0\checkmark$\checkmark^\checkmark\\checkmarkc\checkmarki\checkmarkr\checkmarkc\checkmark$\checkmark0\checkmark$\checkmark^\checkmark\\checkmarkc\checkmarki\checkmarkr\checkmarkc\checkmark$\checkmark/\checkmark$\checkmark^\checkmark\\checkmarkc\checkmarki\checkmarkr\checkmarkc\checkmark$\checkmark2\checkmark$\checkmark^\checkmark\\checkmarkc\checkmarki\checkmarkr\checkmarkc\checkmark$\checkmark1\checkmark$\checkmark^\checkmark\\checkmarkc\checkmarki\checkmarkr\checkmarkc\checkmark$\checkmark9\checkmark$\checkmark^\checkmark\\checkmarkc\checkmarki\checkmarkr\checkmarkc\checkmark$\checkmark \checkmark$\checkmark^\checkmark\\checkmarkc\checkmarki\checkmarkr\checkmarkc\checkmark$\checkmark\\checkmark$\checkmark^\checkmark\\checkmarkc\checkmarki\checkmarkr\checkmarkc\checkmark$\checkmarka\checkmark$\checkmark^\checkmark\\checkmarkc\checkmarki\checkmarkr\checkmarkc\checkmark$\checkmarkp\checkmark$\checkmark^\checkmark\\checkmarkc\checkmarki\checkmarkr\checkmarkc\checkmark$\checkmarkp\checkmark$\checkmark^\checkmark\\checkmarkc\checkmarki\checkmarkr\checkmarkc\checkmark$\checkmarkr\checkmark$\checkmark^\checkmark\\checkmarkc\checkmarki\checkmarkr\checkmarkc\checkmark$\checkmarko\checkmark$\checkmark^\checkmark\\checkmarkc\checkmarki\checkmarkr\checkmarkc\checkmark$\checkmarkx\checkmark$\checkmark^\checkmark\\checkmarkc\checkmarki\checkmarkr\checkmarkc\checkmark$\checkmark \checkmark$\checkmark^\checkmark\\checkmarkc\checkmarki\checkmarkr\checkmarkc\checkmark$\checkmark1\checkmark$\checkmark^\checkmark\\checkmarkc\checkmarki\checkmarkr\checkmarkc\checkmark$\checkmark{\checkmark$\checkmark^\checkmark\\checkmarkc\checkmarki\checkmarkr\checkmarkc\checkmark$\checkmark,\checkmark$\checkmark^\checkmark\\checkmarkc\checkmarki\checkmarkr\checkmarkc\checkmark$\checkmark}\checkmark$\checkmark^\checkmark\\checkmarkc\checkmarki\checkmarkr\checkmarkc\checkmark$\checkmark3\checkmark$\checkmark^\checkmark\\checkmarkc\checkmarki\checkmarkr\checkmarkc\checkmark$\checkmark7\checkmark$\checkmark^\checkmark\\checkmarkc\checkmarki\checkmarkr\checkmarkc\checkmark$\checkmark$\checkmark$\checkmark^\checkmark\\checkmarkc\checkmarki\checkmarkr\checkmarkc\checkmark$\checkmark \checkmark$\checkmark^\checkmark\\checkmarkc\checkmarki\checkmarkr\checkmarkc\checkmark$\checkmarks\checkmark$\checkmark^\checkmark\\checkmarkc\checkmarki\checkmarkr\checkmarkc\checkmark$\checkmarku\checkmark$\checkmark^\checkmark\\checkmarkc\checkmarki\checkmarkr\checkmarkc\checkmark$\checkmarkr\checkmark$\checkmark^\checkmark\\checkmarkc\checkmarki\checkmarkr\checkmarkc\checkmark$\checkmark \checkmark$\checkmark^\checkmark\\checkmarkc\checkmarki\checkmarkr\checkmarkc\checkmark$\checkmarkl\checkmark$\checkmark^\checkmark\\checkmarkc\checkmarki\checkmarkr\checkmarkc\checkmark$\checkmarka\checkmark$\checkmark^\checkmark\\checkmarkc\checkmarki\checkmarkr\checkmarkc\checkmark$\checkmark \checkmark$\checkmark^\checkmark\\checkmarkc\checkmarki\checkmarkr\checkmarkc\checkmark$\checkmarkp\checkmark$\checkmark^\checkmark\\checkmarkc\checkmarki\checkmarkr\checkmarkc\checkmark$\checkmarku\checkmark$\checkmark^\checkmark\\checkmarkc\checkmarki\checkmarkr\checkmarkc\checkmark$\checkmarki\checkmark$\checkmark^\checkmark\\checkmarkc\checkmarki\checkmarkr\checkmarkc\checkmark$\checkmarks\checkmark$\checkmark^\checkmark\\checkmarkc\checkmarki\checkmarkr\checkmarkc\checkmark$\checkmarks\checkmark$\checkmark^\checkmark\\checkmarkc\checkmarki\checkmarkr\checkmarkc\checkmark$\checkmarka\checkmark$\checkmark^\checkmark\\checkmarkc\checkmarki\checkmarkr\checkmarkc\checkmark$\checkmarkn\checkmark$\checkmark^\checkmark\\checkmarkc\checkmarki\checkmarkr\checkmarkc\checkmark$\checkmarkc\checkmark$\checkmark^\checkmark\\checkmarkc\checkmarki\checkmarkr\checkmarkc\checkmark$\checkmarke\checkmark$\checkmark^\checkmark\\checkmarkc\checkmarki\checkmarkr\checkmarkc\checkmark$\checkmark,\checkmark$\checkmark^\checkmark\\checkmarkc\checkmarki\checkmarkr\checkmarkc\checkmark$\checkmark \checkmark$\checkmark^\checkmark\\checkmarkc\checkmarki\checkmarkr\checkmarkc\checkmark$\checkmarkc\checkmark$\checkmark^\checkmark\\checkmarkc\checkmarki\checkmarkr\checkmarkc\checkmark$\checkmarko\checkmark$\checkmark^\checkmark\\checkmarkc\checkmarki\checkmarkr\checkmarkc\checkmark$\checkmarkn\checkmark$\checkmark^\checkmark\\checkmarkc\checkmarki\checkmarkr\checkmarkc\checkmark$\checkmarkf\checkmark$\checkmark^\checkmark\\checkmarkc\checkmarki\checkmarkr\checkmarkc\checkmark$\checkmarko\checkmark$\checkmark^\checkmark\\checkmarkc\checkmarki\checkmarkr\checkmarkc\checkmark$\checkmarkr\checkmark$\checkmark^\checkmark\\checkmarkc\checkmarki\checkmarkr\checkmarkc\checkmark$\checkmarkm\checkmark$\checkmark^\checkmark\\checkmarkc\checkmarki\checkmarkr\checkmarkc\checkmark$\checkmarke\checkmark$\checkmark^\checkmark\\checkmarkc\checkmarki\checkmarkr\checkmarkc\checkmark$\checkmark \checkmark$\checkmark^\checkmark\\checkmarkc\checkmarki\checkmarkr\checkmarkc\checkmark$\checkmarka\checkmark$\checkmark^\checkmark\\checkmarkc\checkmarki\checkmarkr\checkmarkc\checkmark$\checkmarku\checkmark$\checkmark^\checkmark\\checkmarkc\checkmarki\checkmarkr\checkmarkc\checkmark$\checkmarkx\checkmark$\checkmark^\checkmark\\checkmarkc\checkmarki\checkmarkr\checkmarkc\checkmark$\checkmark \checkmark$\checkmark^\checkmark\\checkmarkc\checkmarki\checkmarkr\checkmarkc\checkmark$\checkmarke\checkmark$\checkmark^\checkmark\\checkmarkc\checkmarki\checkmarkr\checkmarkc\checkmark$\checkmarkx\checkmark$\checkmark^\checkmark\\checkmarkc\checkmarki\checkmarkr\checkmarkc\checkmark$\checkmarki\checkmark$\checkmark^\checkmark\\checkmarkc\checkmarki\checkmarkr\checkmarkc\checkmark$\checkmarkg\checkmark$\checkmark^\checkmark\\checkmarkc\checkmarki\checkmarkr\checkmarkc\checkmark$\checkmarke\checkmark$\checkmark^\checkmark\\checkmarkc\checkmarki\checkmarkr\checkmarkc\checkmark$\checkmarkn\checkmark$\checkmark^\checkmark\\checkmarkc\checkmarki\checkmarkr\checkmarkc\checkmark$\checkmarkc\checkmark$\checkmark^\checkmark\\checkmarkc\checkmarki\checkmarkr\checkmarkc\checkmark$\checkmarke\checkmark$\checkmark^\checkmark\\checkmarkc\checkmarki\checkmarkr\checkmarkc\checkmark$\checkmarks\checkmark$\checkmark^\checkmark\\checkmarkc\checkmarki\checkmarkr\checkmarkc\checkmark$\checkmark.\checkmark$\checkmark^\checkmark\\checkmarkc\checkmarki\checkmarkr\checkmarkc\checkmark$\checkmark
\checkmark$\checkmark^\checkmark\\checkmarkc\checkmarki\checkmarkr\checkmarkc\checkmark$\checkmark
\checkmark$\checkmark^\checkmark\\checkmarkc\checkmarki\checkmarkr\checkmarkc\checkmark$\checkmark%\checkmark$\checkmark^\checkmark\\checkmarkc\checkmarki\checkmarkr\checkmarkc\checkmark$\checkmark \checkmark$\checkmark^\checkmark\\checkmarkc\checkmarki\checkmarkr\checkmarkc\checkmark$\checkmark=\checkmark$\checkmark^\checkmark\\checkmarkc\checkmarki\checkmarkr\checkmarkc\checkmark$\checkmark=\checkmark$\checkmark^\checkmark\\checkmarkc\checkmarki\checkmarkr\checkmarkc\checkmark$\checkmark=\checkmark$\checkmark^\checkmark\\checkmarkc\checkmarki\checkmarkr\checkmarkc\checkmark$\checkmark=\checkmark$\checkmark^\checkmark\\checkmarkc\checkmarki\checkmarkr\checkmarkc\checkmark$\checkmark=\checkmark$\checkmark^\checkmark\\checkmarkc\checkmarki\checkmarkr\checkmarkc\checkmark$\checkmark=\checkmark$\checkmark^\checkmark\\checkmarkc\checkmarki\checkmarkr\checkmarkc\checkmark$\checkmark=\checkmark$\checkmark^\checkmark\\checkmarkc\checkmarki\checkmarkr\checkmarkc\checkmark$\checkmark=\checkmark$\checkmark^\checkmark\\checkmarkc\checkmarki\checkmarkr\checkmarkc\checkmark$\checkmark=\checkmark$\checkmark^\checkmark\\checkmarkc\checkmarki\checkmarkr\checkmarkc\checkmark$\checkmark=\checkmark$\checkmark^\checkmark\\checkmarkc\checkmarki\checkmarkr\checkmarkc\checkmark$\checkmark=\checkmark$\checkmark^\checkmark\\checkmarkc\checkmarki\checkmarkr\checkmarkc\checkmark$\checkmark=\checkmark$\checkmark^\checkmark\\checkmarkc\checkmarki\checkmarkr\checkmarkc\checkmark$\checkmark=\checkmark$\checkmark^\checkmark\\checkmarkc\checkmarki\checkmarkr\checkmarkc\checkmark$\checkmark=\checkmark$\checkmark^\checkmark\\checkmarkc\checkmarki\checkmarkr\checkmarkc\checkmark$\checkmark=\checkmark$\checkmark^\checkmark\\checkmarkc\checkmarki\checkmarkr\checkmarkc\checkmark$\checkmark=\checkmark$\checkmark^\checkmark\\checkmarkc\checkmarki\checkmarkr\checkmarkc\checkmark$\checkmark=\checkmark$\checkmark^\checkmark\\checkmarkc\checkmarki\checkmarkr\checkmarkc\checkmark$\checkmark=\checkmark$\checkmark^\checkmark\\checkmarkc\checkmarki\checkmarkr\checkmarkc\checkmark$\checkmark=\checkmark$\checkmark^\checkmark\\checkmarkc\checkmarki\checkmarkr\checkmarkc\checkmark$\checkmark=\checkmark$\checkmark^\checkmark\\checkmarkc\checkmarki\checkmarkr\checkmarkc\checkmark$\checkmark=\checkmark$\checkmark^\checkmark\\checkmarkc\checkmarki\checkmarkr\checkmarkc\checkmark$\checkmark=\checkmark$\checkmark^\checkmark\\checkmarkc\checkmarki\checkmarkr\checkmarkc\checkmark$\checkmark=\checkmark$\checkmark^\checkmark\\checkmarkc\checkmarki\checkmarkr\checkmarkc\checkmark$\checkmark=\checkmark$\checkmark^\checkmark\\checkmarkc\checkmarki\checkmarkr\checkmarkc\checkmark$\checkmark=\checkmark$\checkmark^\checkmark\\checkmarkc\checkmarki\checkmarkr\checkmarkc\checkmark$\checkmark=\checkmark$\checkmark^\checkmark\\checkmarkc\checkmarki\checkmarkr\checkmarkc\checkmark$\checkmark=\checkmark$\checkmark^\checkmark\\checkmarkc\checkmarki\checkmarkr\checkmarkc\checkmark$\checkmark=\checkmark$\checkmark^\checkmark\\checkmarkc\checkmarki\checkmarkr\checkmarkc\checkmark$\checkmark=\checkmark$\checkmark^\checkmark\\checkmarkc\checkmarki\checkmarkr\checkmarkc\checkmark$\checkmark=\checkmark$\checkmark^\checkmark\\checkmarkc\checkmarki\checkmarkr\checkmarkc\checkmark$\checkmark=\checkmark$\checkmark^\checkmark\\checkmarkc\checkmarki\checkmarkr\checkmarkc\checkmark$\checkmark=\checkmark$\checkmark^\checkmark\\checkmarkc\checkmarki\checkmarkr\checkmarkc\checkmark$\checkmark=\checkmark$\checkmark^\checkmark\\checkmarkc\checkmarki\checkmarkr\checkmarkc\checkmark$\checkmark=\checkmark$\checkmark^\checkmark\\checkmarkc\checkmarki\checkmarkr\checkmarkc\checkmark$\checkmark=\checkmark$\checkmark^\checkmark\\checkmarkc\checkmarki\checkmarkr\checkmarkc\checkmark$\checkmark=\checkmark$\checkmark^\checkmark\\checkmarkc\checkmarki\checkmarkr\checkmarkc\checkmark$\checkmark=\checkmark$\checkmark^\checkmark\\checkmarkc\checkmarki\checkmarkr\checkmarkc\checkmark$\checkmark=\checkmark$\checkmark^\checkmark\\checkmarkc\checkmarki\checkmarkr\checkmarkc\checkmark$\checkmark=\checkmark$\checkmark^\checkmark\\checkmarkc\checkmarki\checkmarkr\checkmarkc\checkmark$\checkmark=\checkmark$\checkmark^\checkmark\\checkmarkc\checkmarki\checkmarkr\checkmarkc\checkmark$\checkmark=\checkmark$\checkmark^\checkmark\\checkmarkc\checkmarki\checkmarkr\checkmarkc\checkmark$\checkmark=\checkmark$\checkmark^\checkmark\\checkmarkc\checkmarki\checkmarkr\checkmarkc\checkmark$\checkmark=\checkmark$\checkmark^\checkmark\\checkmarkc\checkmarki\checkmarkr\checkmarkc\checkmark$\checkmark=\checkmark$\checkmark^\checkmark\\checkmarkc\checkmarki\checkmarkr\checkmarkc\checkmark$\checkmark=\checkmark$\checkmark^\checkmark\\checkmarkc\checkmarki\checkmarkr\checkmarkc\checkmark$\checkmark=\checkmark$\checkmark^\checkmark\\checkmarkc\checkmarki\checkmarkr\checkmarkc\checkmark$\checkmark=\checkmark$\checkmark^\checkmark\\checkmarkc\checkmarki\checkmarkr\checkmarkc\checkmark$\checkmark=\checkmark$\checkmark^\checkmark\\checkmarkc\checkmarki\checkmarkr\checkmarkc\checkmark$\checkmark=\checkmark$\checkmark^\checkmark\\checkmarkc\checkmarki\checkmarkr\checkmarkc\checkmark$\checkmark=\checkmark$\checkmark^\checkmark\\checkmarkc\checkmarki\checkmarkr\checkmarkc\checkmark$\checkmark=\checkmark$\checkmark^\checkmark\\checkmarkc\checkmarki\checkmarkr\checkmarkc\checkmark$\checkmark=\checkmark$\checkmark^\checkmark\\checkmarkc\checkmarki\checkmarkr\checkmarkc\checkmark$\checkmark=\checkmark$\checkmark^\checkmark\\checkmarkc\checkmarki\checkmarkr\checkmarkc\checkmark$\checkmark=\checkmark$\checkmark^\checkmark\\checkmarkc\checkmarki\checkmarkr\checkmarkc\checkmark$\checkmark=\checkmark$\checkmark^\checkmark\\checkmarkc\checkmarki\checkmarkr\checkmarkc\checkmark$\checkmark=\checkmark$\checkmark^\checkmark\\checkmarkc\checkmarki\checkmarkr\checkmarkc\checkmark$\checkmark=\checkmark$\checkmark^\checkmark\\checkmarkc\checkmarki\checkmarkr\checkmarkc\checkmark$\checkmark=\checkmark$\checkmark^\checkmark\\checkmarkc\checkmarki\checkmarkr\checkmarkc\checkmark$\checkmark=\checkmark$\checkmark^\checkmark\\checkmarkc\checkmarki\checkmarkr\checkmarkc\checkmark$\checkmark=\checkmark$\checkmark^\checkmark\\checkmarkc\checkmarki\checkmarkr\checkmarkc\checkmark$\checkmark=\checkmark$\checkmark^\checkmark\\checkmarkc\checkmarki\checkmarkr\checkmarkc\checkmark$\checkmark=\checkmark$\checkmark^\checkmark\\checkmarkc\checkmarki\checkmarkr\checkmarkc\checkmark$\checkmark
\checkmark$\checkmark^\checkmark\\checkmarkc\checkmarki\checkmarkr\checkmarkc\checkmark$\checkmark\\checkmark$\checkmark^\checkmark\\checkmarkc\checkmarki\checkmarkr\checkmarkc\checkmark$\checkmarks\checkmark$\checkmark^\checkmark\\checkmarkc\checkmarki\checkmarkr\checkmarkc\checkmark$\checkmarke\checkmark$\checkmark^\checkmark\\checkmarkc\checkmarki\checkmarkr\checkmarkc\checkmark$\checkmarkc\checkmark$\checkmark^\checkmark\\checkmarkc\checkmarki\checkmarkr\checkmarkc\checkmark$\checkmarkt\checkmark$\checkmark^\checkmark\\checkmarkc\checkmarki\checkmarkr\checkmarkc\checkmark$\checkmarki\checkmark$\checkmark^\checkmark\\checkmarkc\checkmarki\checkmarkr\checkmarkc\checkmark$\checkmarko\checkmark$\checkmark^\checkmark\\checkmarkc\checkmarki\checkmarkr\checkmarkc\checkmark$\checkmarkn\checkmark$\checkmark^\checkmark\\checkmarkc\checkmarki\checkmarkr\checkmarkc\checkmark$\checkmark{\checkmark$\checkmark^\checkmark\\checkmarkc\checkmarki\checkmarkr\checkmarkc\checkmark$\checkmarkÉ\checkmark$\checkmark^\checkmark\\checkmarkc\checkmarki\checkmarkr\checkmarkc\checkmark$\checkmarkt\checkmark$\checkmark^\checkmark\\checkmarkc\checkmarki\checkmarkr\checkmarkc\checkmark$\checkmarku\checkmark$\checkmark^\checkmark\\checkmarkc\checkmarki\checkmarkr\checkmarkc\checkmark$\checkmarkd\checkmark$\checkmark^\checkmark\\checkmarkc\checkmarki\checkmarkr\checkmarkc\checkmark$\checkmarke\checkmark$\checkmark^\checkmark\\checkmarkc\checkmarki\checkmarkr\checkmarkc\checkmark$\checkmark \checkmark$\checkmark^\checkmark\\checkmarkc\checkmarki\checkmarkr\checkmarkc\checkmark$\checkmarkA\checkmark$\checkmark^\checkmark\\checkmarkc\checkmarki\checkmarkr\checkmarkc\checkmark$\checkmarkn\checkmark$\checkmark^\checkmark\\checkmarkc\checkmarki\checkmarkr\checkmarkc\checkmark$\checkmarka\checkmark$\checkmark^\checkmark\\checkmarkc\checkmarki\checkmarkr\checkmarkc\checkmark$\checkmarkl\checkmark$\checkmark^\checkmark\\checkmarkc\checkmarki\checkmarkr\checkmarkc\checkmark$\checkmarky\checkmark$\checkmark^\checkmark\\checkmarkc\checkmarki\checkmarkr\checkmarkc\checkmark$\checkmarkt\checkmark$\checkmark^\checkmark\\checkmarkc\checkmarki\checkmarkr\checkmarkc\checkmark$\checkmarki\checkmark$\checkmark^\checkmark\\checkmarkc\checkmarki\checkmarkr\checkmarkc\checkmark$\checkmarkq\checkmark$\checkmark^\checkmark\\checkmarkc\checkmarki\checkmarkr\checkmarkc\checkmark$\checkmarku\checkmark$\checkmark^\checkmark\\checkmarkc\checkmarki\checkmarkr\checkmarkc\checkmark$\checkmarke\checkmark$\checkmark^\checkmark\\checkmarkc\checkmarki\checkmarkr\checkmarkc\checkmark$\checkmark \checkmark$\checkmark^\checkmark\\checkmarkc\checkmarki\checkmarkr\checkmarkc\checkmark$\checkmarkd\checkmark$\checkmark^\checkmark\\checkmarkc\checkmarki\checkmarkr\checkmarkc\checkmark$\checkmarku\checkmark$\checkmark^\checkmark\\checkmarkc\checkmarki\checkmarkr\checkmarkc\checkmark$\checkmark \checkmark$\checkmark^\checkmark\\checkmarkc\checkmarki\checkmarkr\checkmarkc\checkmark$\checkmarkC\checkmark$\checkmark^\checkmark\\checkmarkc\checkmarki\checkmarkr\checkmarkc\checkmark$\checkmarkh\checkmark$\checkmark^\checkmark\\checkmarkc\checkmarki\checkmarkr\checkmarkc\checkmark$\checkmarko\checkmark$\checkmark^\checkmark\\checkmarkc\checkmarki\checkmarkr\checkmarkc\checkmark$\checkmarki\checkmark$\checkmark^\checkmark\\checkmarkc\checkmarki\checkmarkr\checkmarkc\checkmark$\checkmarkx\checkmark$\checkmark^\checkmark\\checkmarkc\checkmarki\checkmarkr\checkmarkc\checkmark$\checkmark \checkmark$\checkmark^\checkmark\\checkmarkc\checkmarki\checkmarkr\checkmarkc\checkmark$\checkmarkd\checkmark$\checkmark^\checkmark\\checkmarkc\checkmarki\checkmarkr\checkmarkc\checkmark$\checkmarke\checkmark$\checkmark^\checkmark\\checkmarkc\checkmarki\checkmarkr\checkmarkc\checkmark$\checkmarks\checkmark$\checkmark^\checkmark\\checkmarkc\checkmarki\checkmarkr\checkmarkc\checkmark$\checkmark \checkmark$\checkmark^\checkmark\\checkmarkc\checkmarki\checkmarkr\checkmarkc\checkmark$\checkmarkO\checkmark$\checkmark^\checkmark\\checkmarkc\checkmarki\checkmarkr\checkmarkc\checkmark$\checkmarkr\checkmark$\checkmark^\checkmark\\checkmarkc\checkmarki\checkmarkr\checkmarkc\checkmark$\checkmarkg\checkmark$\checkmark^\checkmark\\checkmarkc\checkmarki\checkmarkr\checkmarkc\checkmark$\checkmarka\checkmark$\checkmark^\checkmark\\checkmarkc\checkmarki\checkmarkr\checkmarkc\checkmark$\checkmarkn\checkmark$\checkmark^\checkmark\\checkmarkc\checkmarki\checkmarkr\checkmarkc\checkmark$\checkmarke\checkmark$\checkmark^\checkmark\\checkmarkc\checkmarki\checkmarkr\checkmarkc\checkmark$\checkmarks\checkmark$\checkmark^\checkmark\\checkmarkc\checkmarki\checkmarkr\checkmarkc\checkmark$\checkmark}\checkmark$\checkmark^\checkmark\\checkmarkc\checkmarki\checkmarkr\checkmarkc\checkmark$\checkmark
\checkmark$\checkmark^\checkmark\\checkmarkc\checkmarki\checkmarkr\checkmarkc\checkmark$\checkmark
\checkmark$\checkmark^\checkmark\\checkmarkc\checkmarki\checkmarkr\checkmarkc\checkmark$\checkmark\\checkmark$\checkmark^\checkmark\\checkmarkc\checkmarki\checkmarkr\checkmarkc\checkmark$\checkmarks\checkmark$\checkmark^\checkmark\\checkmarkc\checkmarki\checkmarkr\checkmarkc\checkmark$\checkmarku\checkmark$\checkmark^\checkmark\\checkmarkc\checkmarki\checkmarkr\checkmarkc\checkmark$\checkmarkb\checkmark$\checkmark^\checkmark\\checkmarkc\checkmarki\checkmarkr\checkmarkc\checkmark$\checkmarks\checkmark$\checkmark^\checkmark\\checkmarkc\checkmarki\checkmarkr\checkmarkc\checkmark$\checkmarke\checkmark$\checkmark^\checkmark\\checkmarkc\checkmarki\checkmarkr\checkmarkc\checkmark$\checkmarkc\checkmark$\checkmark^\checkmark\\checkmarkc\checkmarki\checkmarkr\checkmarkc\checkmark$\checkmarkt\checkmark$\checkmark^\checkmark\\checkmarkc\checkmarki\checkmarkr\checkmarkc\checkmark$\checkmarki\checkmark$\checkmark^\checkmark\\checkmarkc\checkmarki\checkmarkr\checkmarkc\checkmark$\checkmarko\checkmark$\checkmark^\checkmark\\checkmarkc\checkmarki\checkmarkr\checkmarkc\checkmark$\checkmarkn\checkmark$\checkmark^\checkmark\\checkmarkc\checkmarki\checkmarkr\checkmarkc\checkmark$\checkmark{\checkmark$\checkmark^\checkmark\\checkmarkc\checkmarki\checkmarkr\checkmarkc\checkmark$\checkmarkC\checkmark$\checkmark^\checkmark\\checkmarkc\checkmarki\checkmarkr\checkmarkc\checkmark$\checkmarkh\checkmark$\checkmark^\checkmark\\checkmarkc\checkmarki\checkmarkr\checkmarkc\checkmark$\checkmarko\checkmark$\checkmark^\checkmark\\checkmarkc\checkmarki\checkmarkr\checkmarkc\checkmark$\checkmarki\checkmark$\checkmark^\checkmark\\checkmarkc\checkmarki\checkmarkr\checkmarkc\checkmark$\checkmarkx\checkmark$\checkmark^\checkmark\\checkmarkc\checkmarki\checkmarkr\checkmarkc\checkmark$\checkmark \checkmark$\checkmark^\checkmark\\checkmarkc\checkmarki\checkmarkr\checkmarkc\checkmark$\checkmarkd\checkmark$\checkmark^\checkmark\\checkmarkc\checkmarki\checkmarkr\checkmarkc\checkmark$\checkmarke\checkmark$\checkmark^\checkmark\\checkmarkc\checkmarki\checkmarkr\checkmarkc\checkmark$\checkmarks\checkmark$\checkmark^\checkmark\\checkmarkc\checkmarki\checkmarkr\checkmarkc\checkmark$\checkmark \checkmark$\checkmark^\checkmark\\checkmarkc\checkmarki\checkmarkr\checkmarkc\checkmark$\checkmarkg\checkmark$\checkmark^\checkmark\\checkmarkc\checkmarki\checkmarkr\checkmarkc\checkmark$\checkmarku\checkmark$\checkmark^\checkmark\\checkmarkc\checkmarki\checkmarkr\checkmarkc\checkmark$\checkmarki\checkmark$\checkmark^\checkmark\\checkmarkc\checkmarki\checkmarkr\checkmarkc\checkmark$\checkmarkd\checkmark$\checkmark^\checkmark\\checkmarkc\checkmarki\checkmarkr\checkmarkc\checkmark$\checkmarka\checkmark$\checkmark^\checkmark\\checkmarkc\checkmarki\checkmarkr\checkmarkc\checkmark$\checkmarkg\checkmark$\checkmark^\checkmark\\checkmarkc\checkmarki\checkmarkr\checkmarkc\checkmark$\checkmarke\checkmark$\checkmark^\checkmark\\checkmarkc\checkmarki\checkmarkr\checkmarkc\checkmark$\checkmarks\checkmark$\checkmark^\checkmark\\checkmarkc\checkmarki\checkmarkr\checkmarkc\checkmark$\checkmark \checkmark$\checkmark^\checkmark\\checkmarkc\checkmarki\checkmarkr\checkmarkc\checkmark$\checkmarkl\checkmark$\checkmark^\checkmark\\checkmarkc\checkmarki\checkmarkr\checkmarkc\checkmark$\checkmarki\checkmark$\checkmark^\checkmark\\checkmarkc\checkmarki\checkmarkr\checkmarkc\checkmark$\checkmarkn\checkmark$\checkmark^\checkmark\\checkmarkc\checkmarki\checkmarkr\checkmarkc\checkmark$\checkmarké\checkmark$\checkmark^\checkmark\\checkmarkc\checkmarki\checkmarkr\checkmarkc\checkmark$\checkmarka\checkmark$\checkmark^\checkmark\\checkmarkc\checkmarki\checkmarkr\checkmarkc\checkmark$\checkmarki\checkmark$\checkmark^\checkmark\\checkmarkc\checkmarki\checkmarkr\checkmarkc\checkmark$\checkmarkr\checkmark$\checkmark^\checkmark\\checkmarkc\checkmarki\checkmarkr\checkmarkc\checkmark$\checkmarke\checkmark$\checkmark^\checkmark\\checkmarkc\checkmarki\checkmarkr\checkmarkc\checkmark$\checkmarks\checkmark$\checkmark^\checkmark\\checkmarkc\checkmarki\checkmarkr\checkmarkc\checkmark$\checkmark}\checkmark$\checkmark^\checkmark\\checkmarkc\checkmarki\checkmarkr\checkmarkc\checkmark$\checkmark
\checkmark$\checkmark^\checkmark\\checkmarkc\checkmarki\checkmarkr\checkmarkc\checkmark$\checkmark
\checkmark$\checkmark^\checkmark\\checkmarkc\checkmarki\checkmarkr\checkmarkc\checkmark$\checkmark\\checkmark$\checkmark^\checkmark\\checkmarkc\checkmarki\checkmarkr\checkmarkc\checkmark$\checkmarks\checkmark$\checkmark^\checkmark\\checkmarkc\checkmarki\checkmarkr\checkmarkc\checkmark$\checkmarku\checkmark$\checkmark^\checkmark\\checkmarkc\checkmarki\checkmarkr\checkmarkc\checkmark$\checkmarkb\checkmark$\checkmark^\checkmark\\checkmarkc\checkmarki\checkmarkr\checkmarkc\checkmark$\checkmarks\checkmark$\checkmark^\checkmark\\checkmarkc\checkmarki\checkmarkr\checkmarkc\checkmark$\checkmarku\checkmark$\checkmark^\checkmark\\checkmarkc\checkmarki\checkmarkr\checkmarkc\checkmark$\checkmarkb\checkmark$\checkmark^\checkmark\\checkmarkc\checkmarki\checkmarkr\checkmarkc\checkmark$\checkmarks\checkmark$\checkmark^\checkmark\\checkmarkc\checkmarki\checkmarkr\checkmarkc\checkmark$\checkmarke\checkmark$\checkmark^\checkmark\\checkmarkc\checkmarki\checkmarkr\checkmarkc\checkmark$\checkmarkc\checkmark$\checkmark^\checkmark\\checkmarkc\checkmarki\checkmarkr\checkmarkc\checkmark$\checkmarkt\checkmark$\checkmark^\checkmark\\checkmarkc\checkmarki\checkmarkr\checkmarkc\checkmark$\checkmarki\checkmark$\checkmark^\checkmark\\checkmarkc\checkmarki\checkmarkr\checkmarkc\checkmark$\checkmarko\checkmark$\checkmark^\checkmark\\checkmarkc\checkmarki\checkmarkr\checkmarkc\checkmark$\checkmarkn\checkmark$\checkmark^\checkmark\\checkmarkc\checkmarki\checkmarkr\checkmarkc\checkmark$\checkmark{\checkmark$\checkmark^\checkmark\\checkmarkc\checkmarki\checkmarkr\checkmarkc\checkmark$\checkmarkC\checkmark$\checkmark^\checkmark\\checkmarkc\checkmarki\checkmarkr\checkmarkc\checkmark$\checkmarko\checkmark$\checkmark^\checkmark\\checkmarkc\checkmarki\checkmarkr\checkmarkc\checkmark$\checkmarkn\checkmark$\checkmark^\checkmark\\checkmarkc\checkmarki\checkmarkr\checkmarkc\checkmark$\checkmarkd\checkmark$\checkmark^\checkmark\\checkmarkc\checkmarki\checkmarkr\checkmarkc\checkmark$\checkmarki\checkmark$\checkmark^\checkmark\\checkmarkc\checkmarki\checkmarkr\checkmarkc\checkmark$\checkmarkt\checkmark$\checkmark^\checkmark\\checkmarkc\checkmarki\checkmarkr\checkmarkc\checkmark$\checkmarki\checkmark$\checkmark^\checkmark\\checkmarkc\checkmarki\checkmarkr\checkmarkc\checkmark$\checkmarko\checkmark$\checkmark^\checkmark\\checkmarkc\checkmarki\checkmarkr\checkmarkc\checkmark$\checkmarkn\checkmark$\checkmark^\checkmark\\checkmarkc\checkmarki\checkmarkr\checkmarkc\checkmark$\checkmarks\checkmark$\checkmark^\checkmark\\checkmarkc\checkmarki\checkmarkr\checkmarkc\checkmark$\checkmark \checkmark$\checkmark^\checkmark\\checkmarkc\checkmarki\checkmarkr\checkmarkc\checkmark$\checkmarks\checkmark$\checkmark^\checkmark\\checkmarkc\checkmarki\checkmarkr\checkmarkc\checkmark$\checkmarka\checkmark$\checkmark^\checkmark\\checkmarkc\checkmarki\checkmarkr\checkmarkc\checkmark$\checkmarkt\checkmark$\checkmark^\checkmark\\checkmarkc\checkmarki\checkmarkr\checkmarkc\checkmark$\checkmarki\checkmark$\checkmark^\checkmark\\checkmarkc\checkmarki\checkmarkr\checkmarkc\checkmark$\checkmarks\checkmark$\checkmark^\checkmark\\checkmarkc\checkmarki\checkmarkr\checkmarkc\checkmark$\checkmarkf\checkmark$\checkmark^\checkmark\\checkmarkc\checkmarki\checkmarkr\checkmarkc\checkmark$\checkmarka\checkmark$\checkmark^\checkmark\\checkmarkc\checkmarki\checkmarkr\checkmarkc\checkmark$\checkmarki\checkmark$\checkmark^\checkmark\\checkmarkc\checkmarki\checkmarkr\checkmarkc\checkmark$\checkmarks\checkmark$\checkmark^\checkmark\\checkmarkc\checkmarki\checkmarkr\checkmarkc\checkmark$\checkmarka\checkmark$\checkmark^\checkmark\\checkmarkc\checkmarki\checkmarkr\checkmarkc\checkmark$\checkmarkn\checkmark$\checkmark^\checkmark\\checkmarkc\checkmarki\checkmarkr\checkmarkc\checkmark$\checkmarkt\checkmark$\checkmark^\checkmark\\checkmarkc\checkmarki\checkmarkr\checkmarkc\checkmark$\checkmarke\checkmark$\checkmark^\checkmark\\checkmarkc\checkmarki\checkmarkr\checkmarkc\checkmark$\checkmarks\checkmark$\checkmark^\checkmark\\checkmarkc\checkmarki\checkmarkr\checkmarkc\checkmark$\checkmark}\checkmark$\checkmark^\checkmark\\checkmarkc\checkmarki\checkmarkr\checkmarkc\checkmark$\checkmark
\checkmark$\checkmark^\checkmark\\checkmarkc\checkmarki\checkmarkr\checkmarkc\checkmark$\checkmarkL\checkmark$\checkmark^\checkmark\\checkmarkc\checkmarki\checkmarkr\checkmarkc\checkmark$\checkmarke\checkmark$\checkmark^\checkmark\\checkmarkc\checkmarki\checkmarkr\checkmarkc\checkmark$\checkmarks\checkmark$\checkmark^\checkmark\\checkmarkc\checkmarki\checkmarkr\checkmarkc\checkmark$\checkmark \checkmark$\checkmark^\checkmark\\checkmarkc\checkmarki\checkmarkr\checkmarkc\checkmark$\checkmarkg\checkmark$\checkmark^\checkmark\\checkmarkc\checkmarki\checkmarkr\checkmarkc\checkmark$\checkmarku\checkmark$\checkmark^\checkmark\\checkmarkc\checkmarki\checkmarkr\checkmarkc\checkmark$\checkmarki\checkmark$\checkmark^\checkmark\\checkmarkc\checkmarki\checkmarkr\checkmarkc\checkmark$\checkmarkd\checkmark$\checkmark^\checkmark\\checkmarkc\checkmarki\checkmarkr\checkmarkc\checkmark$\checkmarka\checkmark$\checkmark^\checkmark\\checkmarkc\checkmarki\checkmarkr\checkmarkc\checkmark$\checkmarkg\checkmark$\checkmark^\checkmark\\checkmarkc\checkmarki\checkmarkr\checkmarkc\checkmark$\checkmarke\checkmark$\checkmark^\checkmark\\checkmarkc\checkmarki\checkmarkr\checkmarkc\checkmark$\checkmarks\checkmark$\checkmark^\checkmark\\checkmarkc\checkmarki\checkmarkr\checkmarkc\checkmark$\checkmark \checkmark$\checkmark^\checkmark\\checkmarkc\checkmarki\checkmarkr\checkmarkc\checkmark$\checkmarkl\checkmark$\checkmark^\checkmark\\checkmarkc\checkmarki\checkmarkr\checkmarkc\checkmark$\checkmarki\checkmark$\checkmark^\checkmark\\checkmarkc\checkmarki\checkmarkr\checkmarkc\checkmark$\checkmarkn\checkmark$\checkmark^\checkmark\\checkmarkc\checkmarki\checkmarkr\checkmarkc\checkmark$\checkmarké\checkmark$\checkmark^\checkmark\\checkmarkc\checkmarki\checkmarkr\checkmarkc\checkmark$\checkmarka\checkmark$\checkmark^\checkmark\\checkmarkc\checkmarki\checkmarkr\checkmarkc\checkmark$\checkmarki\checkmark$\checkmark^\checkmark\\checkmarkc\checkmarki\checkmarkr\checkmarkc\checkmark$\checkmarkr\checkmark$\checkmark^\checkmark\\checkmarkc\checkmarki\checkmarkr\checkmarkc\checkmark$\checkmarke\checkmark$\checkmark^\checkmark\\checkmarkc\checkmarki\checkmarkr\checkmarkc\checkmark$\checkmarks\checkmark$\checkmark^\checkmark\\checkmarkc\checkmarki\checkmarkr\checkmarkc\checkmark$\checkmark \checkmark$\checkmark^\checkmark\\checkmarkc\checkmarki\checkmarkr\checkmarkc\checkmark$\checkmarkd\checkmark$\checkmark^\checkmark\\checkmarkc\checkmarki\checkmarkr\checkmarkc\checkmark$\checkmarko\checkmark$\checkmark^\checkmark\\checkmarkc\checkmarki\checkmarkr\checkmarkc\checkmark$\checkmarki\checkmark$\checkmark^\checkmark\\checkmarkc\checkmarki\checkmarkr\checkmarkc\checkmark$\checkmarkv\checkmark$\checkmark^\checkmark\\checkmarkc\checkmarki\checkmarkr\checkmarkc\checkmark$\checkmarke\checkmark$\checkmark^\checkmark\\checkmarkc\checkmarki\checkmarkr\checkmarkc\checkmark$\checkmarkn\checkmark$\checkmark^\checkmark\\checkmarkc\checkmarki\checkmarkr\checkmarkc\checkmark$\checkmarkt\checkmark$\checkmark^\checkmark\\checkmarkc\checkmarki\checkmarkr\checkmarkc\checkmark$\checkmark \checkmark$\checkmark^\checkmark\\checkmarkc\checkmarki\checkmarkr\checkmarkc\checkmark$\checkmarkr\checkmark$\checkmark^\checkmark\\checkmarkc\checkmarki\checkmarkr\checkmarkc\checkmark$\checkmarke\checkmark$\checkmark^\checkmark\\checkmarkc\checkmarki\checkmarkr\checkmarkc\checkmark$\checkmarkm\checkmark$\checkmark^\checkmark\\checkmarkc\checkmarki\checkmarkr\checkmarkc\checkmark$\checkmarkp\checkmark$\checkmark^\checkmark\\checkmarkc\checkmarki\checkmarkr\checkmarkc\checkmark$\checkmarkl\checkmark$\checkmark^\checkmark\\checkmarkc\checkmarki\checkmarkr\checkmarkc\checkmark$\checkmarki\checkmark$\checkmark^\checkmark\\checkmarkc\checkmarki\checkmarkr\checkmarkc\checkmark$\checkmarkr\checkmark$\checkmark^\checkmark\\checkmarkc\checkmarki\checkmarkr\checkmarkc\checkmark$\checkmark \checkmark$\checkmark^\checkmark\\checkmarkc\checkmarki\checkmarkr\checkmarkc\checkmark$\checkmarkl\checkmark$\checkmark^\checkmark\\checkmarkc\checkmarki\checkmarkr\checkmarkc\checkmark$\checkmarke\checkmark$\checkmark^\checkmark\\checkmarkc\checkmarki\checkmarkr\checkmarkc\checkmark$\checkmarks\checkmark$\checkmark^\checkmark\\checkmarkc\checkmarki\checkmarkr\checkmarkc\checkmark$\checkmark \checkmark$\checkmark^\checkmark\\checkmarkc\checkmarki\checkmarkr\checkmarkc\checkmark$\checkmarkc\checkmark$\checkmark^\checkmark\\checkmarkc\checkmarki\checkmarkr\checkmarkc\checkmark$\checkmarko\checkmark$\checkmark^\checkmark\\checkmarkc\checkmarki\checkmarkr\checkmarkc\checkmark$\checkmarkn\checkmark$\checkmark^\checkmark\\checkmarkc\checkmarki\checkmarkr\checkmarkc\checkmark$\checkmarkd\checkmark$\checkmark^\checkmark\\checkmarkc\checkmarki\checkmarkr\checkmarkc\checkmark$\checkmarki\checkmark$\checkmark^\checkmark\\checkmarkc\checkmarki\checkmarkr\checkmarkc\checkmark$\checkmarkt\checkmark$\checkmark^\checkmark\\checkmarkc\checkmarki\checkmarkr\checkmarkc\checkmark$\checkmarki\checkmark$\checkmark^\checkmark\\checkmarkc\checkmarki\checkmarkr\checkmarkc\checkmark$\checkmarko\checkmark$\checkmark^\checkmark\\checkmarkc\checkmarki\checkmarkr\checkmarkc\checkmark$\checkmarkn\checkmark$\checkmark^\checkmark\\checkmarkc\checkmarki\checkmarkr\checkmarkc\checkmark$\checkmarks\checkmark$\checkmark^\checkmark\\checkmarkc\checkmarki\checkmarkr\checkmarkc\checkmark$\checkmark \checkmark$\checkmark^\checkmark\\checkmarkc\checkmarki\checkmarkr\checkmarkc\checkmark$\checkmarks\checkmark$\checkmark^\checkmark\\checkmarkc\checkmarki\checkmarkr\checkmarkc\checkmark$\checkmarku\checkmark$\checkmark^\checkmark\\checkmarkc\checkmarki\checkmarkr\checkmarkc\checkmark$\checkmarki\checkmark$\checkmark^\checkmark\\checkmarkc\checkmarki\checkmarkr\checkmarkc\checkmark$\checkmarkv\checkmark$\checkmark^\checkmark\\checkmarkc\checkmarki\checkmarkr\checkmarkc\checkmark$\checkmarka\checkmark$\checkmark^\checkmark\\checkmarkc\checkmarki\checkmarkr\checkmarkc\checkmark$\checkmarkn\checkmark$\checkmark^\checkmark\\checkmarkc\checkmarki\checkmarkr\checkmarkc\checkmark$\checkmarkt\checkmark$\checkmark^\checkmark\\checkmarkc\checkmarki\checkmarkr\checkmarkc\checkmark$\checkmarke\checkmark$\checkmark^\checkmark\\checkmarkc\checkmarki\checkmarkr\checkmarkc\checkmark$\checkmarks\checkmark$\checkmark^\checkmark\\checkmarkc\checkmarki\checkmarkr\checkmarkc\checkmark$\checkmark \checkmark$\checkmark^\checkmark\\checkmarkc\checkmarki\checkmarkr\checkmarkc\checkmark$\checkmark:\checkmark$\checkmark^\checkmark\\checkmarkc\checkmarki\checkmarkr\checkmarkc\checkmark$\checkmark
\checkmark$\checkmark^\checkmark\\checkmarkc\checkmarki\checkmarkr\checkmarkc\checkmark$\checkmark\\checkmark$\checkmark^\checkmark\\checkmarkc\checkmarki\checkmarkr\checkmarkc\checkmark$\checkmarkb\checkmark$\checkmark^\checkmark\\checkmarkc\checkmarki\checkmarkr\checkmarkc\checkmark$\checkmarke\checkmark$\checkmark^\checkmark\\checkmarkc\checkmarki\checkmarkr\checkmarkc\checkmark$\checkmarkg\checkmark$\checkmark^\checkmark\\checkmarkc\checkmarki\checkmarkr\checkmarkc\checkmark$\checkmarki\checkmark$\checkmark^\checkmark\\checkmarkc\checkmarki\checkmarkr\checkmarkc\checkmark$\checkmarkn\checkmark$\checkmark^\checkmark\\checkmarkc\checkmarki\checkmarkr\checkmarkc\checkmark$\checkmark{\checkmark$\checkmark^\checkmark\\checkmarkc\checkmarki\checkmarkr\checkmarkc\checkmark$\checkmarki\checkmark$\checkmark^\checkmark\\checkmarkc\checkmarki\checkmarkr\checkmarkc\checkmark$\checkmarkt\checkmark$\checkmark^\checkmark\\checkmarkc\checkmarki\checkmarkr\checkmarkc\checkmark$\checkmarke\checkmark$\checkmark^\checkmark\\checkmarkc\checkmarki\checkmarkr\checkmarkc\checkmark$\checkmarkm\checkmark$\checkmark^\checkmark\\checkmarkc\checkmarki\checkmarkr\checkmarkc\checkmark$\checkmarki\checkmark$\checkmark^\checkmark\\checkmarkc\checkmarki\checkmarkr\checkmarkc\checkmark$\checkmarkz\checkmark$\checkmark^\checkmark\\checkmarkc\checkmarki\checkmarkr\checkmarkc\checkmark$\checkmarke\checkmark$\checkmark^\checkmark\\checkmarkc\checkmarki\checkmarkr\checkmarkc\checkmark$\checkmark}\checkmark$\checkmark^\checkmark\\checkmarkc\checkmarki\checkmarkr\checkmarkc\checkmark$\checkmark
\checkmark$\checkmark^\checkmark\\checkmarkc\checkmarki\checkmarkr\checkmarkc\checkmark$\checkmark \checkmark$\checkmark^\checkmark\\checkmarkc\checkmarki\checkmarkr\checkmarkc\checkmark$\checkmark \checkmark$\checkmark^\checkmark\\checkmarkc\checkmarki\checkmarkr\checkmarkc\checkmark$\checkmark \checkmark$\checkmark^\checkmark\\checkmarkc\checkmarki\checkmarkr\checkmarkc\checkmark$\checkmark \checkmark$\checkmark^\checkmark\\checkmarkc\checkmarki\checkmarkr\checkmarkc\checkmark$\checkmark\\checkmark$\checkmark^\checkmark\\checkmarkc\checkmarki\checkmarkr\checkmarkc\checkmark$\checkmarki\checkmark$\checkmark^\checkmark\\checkmarkc\checkmarki\checkmarkr\checkmarkc\checkmark$\checkmarkt\checkmark$\checkmark^\checkmark\\checkmarkc\checkmarki\checkmarkr\checkmarkc\checkmark$\checkmarke\checkmark$\checkmark^\checkmark\\checkmarkc\checkmarki\checkmarkr\checkmarkc\checkmark$\checkmarkm\checkmark$\checkmark^\checkmark\\checkmarkc\checkmarki\checkmarkr\checkmarkc\checkmark$\checkmark \checkmark$\checkmark^\checkmark\\checkmarkc\checkmarki\checkmarkr\checkmarkc\checkmark$\checkmarkR\checkmark$\checkmark^\checkmark\\checkmarkc\checkmarki\checkmarkr\checkmarkc\checkmark$\checkmarké\checkmark$\checkmark^\checkmark\\checkmarkc\checkmarki\checkmarkr\checkmarkc\checkmark$\checkmarks\checkmark$\checkmark^\checkmark\\checkmarkc\checkmarki\checkmarkr\checkmarkc\checkmark$\checkmarki\checkmark$\checkmark^\checkmark\\checkmarkc\checkmarki\checkmarkr\checkmarkc\checkmark$\checkmarks\checkmark$\checkmark^\checkmark\\checkmarkc\checkmarki\checkmarkr\checkmarkc\checkmark$\checkmarkt\checkmark$\checkmark^\checkmark\\checkmarkc\checkmarki\checkmarkr\checkmarkc\checkmark$\checkmarke\checkmark$\checkmark^\checkmark\\checkmarkc\checkmarki\checkmarkr\checkmarkc\checkmark$\checkmarkr\checkmark$\checkmark^\checkmark\\checkmarkc\checkmarki\checkmarkr\checkmarkc\checkmark$\checkmark \checkmark$\checkmark^\checkmark\\checkmarkc\checkmarki\checkmarkr\checkmarkc\checkmark$\checkmarka\checkmark$\checkmark^\checkmark\\checkmarkc\checkmarki\checkmarkr\checkmarkc\checkmark$\checkmarku\checkmark$\checkmark^\checkmark\\checkmarkc\checkmarki\checkmarkr\checkmarkc\checkmark$\checkmarkx\checkmark$\checkmark^\checkmark\\checkmarkc\checkmarki\checkmarkr\checkmarkc\checkmark$\checkmark \checkmark$\checkmark^\checkmark\\checkmarkc\checkmarki\checkmarkr\checkmarkc\checkmark$\checkmarke\checkmark$\checkmark^\checkmark\\checkmarkc\checkmarki\checkmarkr\checkmarkc\checkmark$\checkmarkf\checkmark$\checkmark^\checkmark\\checkmarkc\checkmarki\checkmarkr\checkmarkc\checkmark$\checkmarkf\checkmark$\checkmark^\checkmark\\checkmarkc\checkmarki\checkmarkr\checkmarkc\checkmark$\checkmarko\checkmark$\checkmark^\checkmark\\checkmarkc\checkmarki\checkmarkr\checkmarkc\checkmark$\checkmarkr\checkmark$\checkmark^\checkmark\\checkmarkc\checkmarki\checkmarkr\checkmarkc\checkmark$\checkmarkt\checkmark$\checkmark^\checkmark\\checkmarkc\checkmarki\checkmarkr\checkmarkc\checkmark$\checkmarks\checkmark$\checkmark^\checkmark\\checkmarkc\checkmarki\checkmarkr\checkmarkc\checkmark$\checkmark \checkmark$\checkmark^\checkmark\\checkmarkc\checkmarki\checkmarkr\checkmarkc\checkmark$\checkmarkr\checkmark$\checkmark^\checkmark\\checkmarkc\checkmarki\checkmarkr\checkmarkc\checkmark$\checkmarka\checkmark$\checkmark^\checkmark\\checkmarkc\checkmarki\checkmarkr\checkmarkc\checkmark$\checkmarkd\checkmark$\checkmark^\checkmark\\checkmarkc\checkmarki\checkmarkr\checkmarkc\checkmark$\checkmarki\checkmark$\checkmark^\checkmark\\checkmarkc\checkmarki\checkmarkr\checkmarkc\checkmark$\checkmarka\checkmark$\checkmark^\checkmark\\checkmarkc\checkmarki\checkmarkr\checkmarkc\checkmark$\checkmarku\checkmark$\checkmark^\checkmark\\checkmarkc\checkmarki\checkmarkr\checkmarkc\checkmark$\checkmarkx\checkmark$\checkmark^\checkmark\\checkmarkc\checkmarki\checkmarkr\checkmarkc\checkmark$\checkmark \checkmark$\checkmark^\checkmark\\checkmarkc\checkmarki\checkmarkr\checkmarkc\checkmark$\checkmarke\checkmark$\checkmark^\checkmark\\checkmarkc\checkmarki\checkmarkr\checkmarkc\checkmark$\checkmarkt\checkmark$\checkmark^\checkmark\\checkmarkc\checkmarki\checkmarkr\checkmarkc\checkmark$\checkmark \checkmark$\checkmark^\checkmark\\checkmarkc\checkmarki\checkmarkr\checkmarkc\checkmark$\checkmarka\checkmark$\checkmark^\checkmark\\checkmarkc\checkmarki\checkmarkr\checkmarkc\checkmark$\checkmarku\checkmark$\checkmark^\checkmark\\checkmarkc\checkmarki\checkmarkr\checkmarkc\checkmark$\checkmarkx\checkmark$\checkmark^\checkmark\\checkmarkc\checkmarki\checkmarkr\checkmarkc\checkmark$\checkmark \checkmark$\checkmark^\checkmark\\checkmarkc\checkmarki\checkmarkr\checkmarkc\checkmark$\checkmarkm\checkmark$\checkmark^\checkmark\\checkmarkc\checkmarki\checkmarkr\checkmarkc\checkmark$\checkmarko\checkmark$\checkmark^\checkmark\\checkmarkc\checkmarki\checkmarkr\checkmarkc\checkmark$\checkmarkm\checkmark$\checkmark^\checkmark\\checkmarkc\checkmarki\checkmarkr\checkmarkc\checkmark$\checkmarke\checkmark$\checkmark^\checkmark\\checkmarkc\checkmarki\checkmarkr\checkmarkc\checkmark$\checkmarkn\checkmark$\checkmark^\checkmark\\checkmarkc\checkmarki\checkmarkr\checkmarkc\checkmark$\checkmarkt\checkmark$\checkmark^\checkmark\\checkmarkc\checkmarki\checkmarkr\checkmarkc\checkmark$\checkmarks\checkmark$\checkmark^\checkmark\\checkmarkc\checkmarki\checkmarkr\checkmarkc\checkmark$\checkmark \checkmark$\checkmark^\checkmark\\checkmarkc\checkmarki\checkmarkr\checkmarkc\checkmark$\checkmarkd\checkmark$\checkmark^\checkmark\\checkmarkc\checkmarki\checkmarkr\checkmarkc\checkmark$\checkmarke\checkmark$\checkmark^\checkmark\\checkmarkc\checkmarki\checkmarkr\checkmarkc\checkmark$\checkmark \checkmark$\checkmark^\checkmark\\checkmarkc\checkmarki\checkmarkr\checkmarkc\checkmark$\checkmarkr\checkmark$\checkmark^\checkmark\\checkmarkc\checkmarki\checkmarkr\checkmarkc\checkmark$\checkmarke\checkmark$\checkmark^\checkmark\\checkmarkc\checkmarki\checkmarkr\checkmarkc\checkmark$\checkmarkn\checkmark$\checkmark^\checkmark\\checkmarkc\checkmarki\checkmarkr\checkmarkc\checkmark$\checkmarkv\checkmark$\checkmark^\checkmark\\checkmarkc\checkmarki\checkmarkr\checkmarkc\checkmark$\checkmarke\checkmark$\checkmark^\checkmark\\checkmarkc\checkmarki\checkmarkr\checkmarkc\checkmark$\checkmarkr\checkmark$\checkmark^\checkmark\\checkmarkc\checkmarki\checkmarkr\checkmarkc\checkmark$\checkmarks\checkmark$\checkmark^\checkmark\\checkmarkc\checkmarki\checkmarkr\checkmarkc\checkmark$\checkmarke\checkmark$\checkmark^\checkmark\\checkmarkc\checkmarki\checkmarkr\checkmarkc\checkmark$\checkmarkm\checkmark$\checkmark^\checkmark\\checkmarkc\checkmarki\checkmarkr\checkmarkc\checkmark$\checkmarke\checkmark$\checkmark^\checkmark\\checkmarkc\checkmarki\checkmarkr\checkmarkc\checkmark$\checkmarkn\checkmark$\checkmark^\checkmark\\checkmarkc\checkmarki\checkmarkr\checkmarkc\checkmark$\checkmarkt\checkmark$\checkmark^\checkmark\\checkmarkc\checkmarki\checkmarkr\checkmarkc\checkmark$\checkmark \checkmark$\checkmark^\checkmark\\checkmarkc\checkmarki\checkmarkr\checkmarkc\checkmark$\checkmarke\checkmark$\checkmark^\checkmark\\checkmarkc\checkmarki\checkmarkr\checkmarkc\checkmark$\checkmarkn\checkmark$\checkmark^\checkmark\\checkmarkc\checkmarki\checkmarkr\checkmarkc\checkmark$\checkmark \checkmark$\checkmark^\checkmark\\checkmarkc\checkmarki\checkmarkr\checkmarkc\checkmark$\checkmarku\checkmark$\checkmark^\checkmark\\checkmarkc\checkmarki\checkmarkr\checkmarkc\checkmark$\checkmarks\checkmark$\checkmark^\checkmark\\checkmarkc\checkmarki\checkmarkr\checkmarkc\checkmark$\checkmarki\checkmark$\checkmark^\checkmark\\checkmarkc\checkmarki\checkmarkr\checkmarkc\checkmark$\checkmarkn\checkmark$\checkmark^\checkmark\\checkmarkc\checkmarki\checkmarkr\checkmarkc\checkmark$\checkmarka\checkmark$\checkmark^\checkmark\\checkmarkc\checkmarki\checkmarkr\checkmarkc\checkmark$\checkmarkg\checkmark$\checkmark^\checkmark\\checkmarkc\checkmarki\checkmarkr\checkmarkc\checkmark$\checkmarke\checkmark$\checkmark^\checkmark\\checkmarkc\checkmarki\checkmarkr\checkmarkc\checkmark$\checkmark \checkmark$\checkmark^\checkmark\\checkmarkc\checkmarki\checkmarkr\checkmarkc\checkmark$\checkmark;\checkmark$\checkmark^\checkmark\\checkmarkc\checkmarki\checkmarkr\checkmarkc\checkmark$\checkmark
\checkmark$\checkmark^\checkmark\\checkmarkc\checkmarki\checkmarkr\checkmarkc\checkmark$\checkmark \checkmark$\checkmark^\checkmark\\checkmarkc\checkmarki\checkmarkr\checkmarkc\checkmark$\checkmark \checkmark$\checkmark^\checkmark\\checkmarkc\checkmarki\checkmarkr\checkmarkc\checkmark$\checkmark \checkmark$\checkmark^\checkmark\\checkmarkc\checkmarki\checkmarkr\checkmarkc\checkmark$\checkmark \checkmark$\checkmark^\checkmark\\checkmarkc\checkmarki\checkmarkr\checkmarkc\checkmark$\checkmark\\checkmark$\checkmark^\checkmark\\checkmarkc\checkmarki\checkmarkr\checkmarkc\checkmark$\checkmarki\checkmark$\checkmark^\checkmark\\checkmarkc\checkmarki\checkmarkr\checkmarkc\checkmark$\checkmarkt\checkmark$\checkmark^\checkmark\\checkmarkc\checkmarki\checkmarkr\checkmarkc\checkmark$\checkmarke\checkmark$\checkmark^\checkmark\\checkmarkc\checkmarki\checkmarkr\checkmarkc\checkmark$\checkmarkm\checkmark$\checkmark^\checkmark\\checkmarkc\checkmarki\checkmarkr\checkmarkc\checkmark$\checkmark \checkmark$\checkmark^\checkmark\\checkmarkc\checkmarki\checkmarkr\checkmarkc\checkmark$\checkmarkO\checkmark$\checkmark^\checkmark\\checkmarkc\checkmarki\checkmarkr\checkmarkc\checkmark$\checkmarkf\checkmark$\checkmark^\checkmark\\checkmarkc\checkmarki\checkmarkr\checkmarkc\checkmark$\checkmarkf\checkmark$\checkmark^\checkmark\\checkmarkc\checkmarki\checkmarkr\checkmarkc\checkmark$\checkmarkr\checkmark$\checkmark^\checkmark\\checkmarkc\checkmarki\checkmarkr\checkmarkc\checkmark$\checkmarki\checkmark$\checkmark^\checkmark\\checkmarkc\checkmarki\checkmarkr\checkmarkc\checkmark$\checkmarkr\checkmark$\checkmark^\checkmark\\checkmarkc\checkmarki\checkmarkr\checkmarkc\checkmark$\checkmark \checkmark$\checkmark^\checkmark\\checkmarkc\checkmarki\checkmarkr\checkmarkc\checkmark$\checkmarku\checkmark$\checkmark^\checkmark\\checkmarkc\checkmarki\checkmarkr\checkmarkc\checkmark$\checkmarkn\checkmark$\checkmark^\checkmark\\checkmarkc\checkmarki\checkmarkr\checkmarkc\checkmark$\checkmarke\checkmark$\checkmark^\checkmark\\checkmarkc\checkmarki\checkmarkr\checkmarkc\checkmark$\checkmark \checkmark$\checkmark^\checkmark\\checkmarkc\checkmarki\checkmarkr\checkmarkc\checkmark$\checkmarkr\checkmark$\checkmark^\checkmark\\checkmarkc\checkmarki\checkmarkr\checkmarkc\checkmark$\checkmarki\checkmark$\checkmark^\checkmark\\checkmarkc\checkmarki\checkmarkr\checkmarkc\checkmark$\checkmarkg\checkmark$\checkmark^\checkmark\\checkmarkc\checkmarki\checkmarkr\checkmarkc\checkmark$\checkmarki\checkmark$\checkmark^\checkmark\\checkmarkc\checkmarki\checkmarkr\checkmarkc\checkmark$\checkmarkd\checkmark$\checkmark^\checkmark\\checkmarkc\checkmarki\checkmarkr\checkmarkc\checkmark$\checkmarki\checkmark$\checkmark^\checkmark\\checkmarkc\checkmarki\checkmarkr\checkmarkc\checkmark$\checkmarkt\checkmark$\checkmark^\checkmark\\checkmarkc\checkmarki\checkmarkr\checkmarkc\checkmark$\checkmarké\checkmark$\checkmark^\checkmark\\checkmarkc\checkmarki\checkmarkr\checkmarkc\checkmark$\checkmark \checkmark$\checkmark^\checkmark\\checkmarkc\checkmarki\checkmarkr\checkmarkc\checkmark$\checkmarks\checkmark$\checkmark^\checkmark\\checkmarkc\checkmarki\checkmarkr\checkmarkc\checkmark$\checkmarku\checkmark$\checkmark^\checkmark\\checkmarkc\checkmarki\checkmarkr\checkmarkc\checkmark$\checkmarkp\checkmark$\checkmark^\checkmark\\checkmarkc\checkmarki\checkmarkr\checkmarkc\checkmark$\checkmarké\checkmark$\checkmark^\checkmark\\checkmarkc\checkmarki\checkmarkr\checkmarkc\checkmark$\checkmarkr\checkmark$\checkmark^\checkmark\\checkmarkc\checkmarki\checkmarkr\checkmarkc\checkmark$\checkmarki\checkmark$\checkmark^\checkmark\\checkmarkc\checkmarki\checkmarkr\checkmarkc\checkmark$\checkmarke\checkmark$\checkmark^\checkmark\\checkmarkc\checkmarki\checkmarkr\checkmarkc\checkmark$\checkmarku\checkmark$\checkmark^\checkmark\\checkmarkc\checkmarki\checkmarkr\checkmarkc\checkmark$\checkmarkr\checkmark$\checkmark^\checkmark\\checkmarkc\checkmarki\checkmarkr\checkmarkc\checkmark$\checkmarke\checkmark$\checkmark^\checkmark\\checkmarkc\checkmarki\checkmarkr\checkmarkc\checkmark$\checkmark \checkmark$\checkmark^\checkmark\\checkmarkc\checkmarki\checkmarkr\checkmarkc\checkmark$\checkmarka\checkmark$\checkmark^\checkmark\\checkmarkc\checkmarki\checkmarkr\checkmarkc\checkmark$\checkmarku\checkmark$\checkmark^\checkmark\\checkmarkc\checkmarki\checkmarkr\checkmarkc\checkmark$\checkmarkx\checkmark$\checkmark^\checkmark\\checkmarkc\checkmarki\checkmarkr\checkmarkc\checkmark$\checkmark \checkmark$\checkmark^\checkmark\\checkmarkc\checkmarki\checkmarkr\checkmarkc\checkmark$\checkmarka\checkmark$\checkmark^\checkmark\\checkmarkc\checkmarki\checkmarkr\checkmarkc\checkmark$\checkmarkl\checkmark$\checkmark^\checkmark\\checkmarkc\checkmarki\checkmarkr\checkmarkc\checkmark$\checkmarkt\checkmark$\checkmark^\checkmark\\checkmarkc\checkmarki\checkmarkr\checkmarkc\checkmark$\checkmarke\checkmark$\checkmark^\checkmark\\checkmarkc\checkmarki\checkmarkr\checkmarkc\checkmark$\checkmarkr\checkmark$\checkmark^\checkmark\\checkmarkc\checkmarki\checkmarkr\checkmarkc\checkmark$\checkmarkn\checkmark$\checkmark^\checkmark\\checkmarkc\checkmarki\checkmarkr\checkmarkc\checkmark$\checkmarka\checkmark$\checkmark^\checkmark\\checkmarkc\checkmarki\checkmarkr\checkmarkc\checkmark$\checkmarkt\checkmark$\checkmark^\checkmark\\checkmarkc\checkmarki\checkmarkr\checkmarkc\checkmark$\checkmarki\checkmark$\checkmark^\checkmark\\checkmarkc\checkmarki\checkmarkr\checkmarkc\checkmark$\checkmarkv\checkmark$\checkmark^\checkmark\\checkmarkc\checkmarki\checkmarkr\checkmarkc\checkmark$\checkmarke\checkmark$\checkmark^\checkmark\\checkmarkc\checkmarki\checkmarkr\checkmarkc\checkmark$\checkmarks\checkmark$\checkmark^\checkmark\\checkmarkc\checkmarki\checkmarkr\checkmarkc\checkmark$\checkmark \checkmark$\checkmark^\checkmark\\checkmarkc\checkmarki\checkmarkr\checkmarkc\checkmark$\checkmark(\checkmark$\checkmark^\checkmark\\checkmarkc\checkmarki\checkmarkr\checkmarkc\checkmark$\checkmarkt\checkmark$\checkmark^\checkmark\\checkmarkc\checkmarki\checkmarkr\checkmarkc\checkmark$\checkmarki\checkmark$\checkmark^\checkmark\\checkmarkc\checkmarki\checkmarkr\checkmarkc\checkmark$\checkmarkg\checkmark$\checkmark^\checkmark\\checkmarkc\checkmarki\checkmarkr\checkmarkc\checkmark$\checkmarke\checkmark$\checkmark^\checkmark\\checkmarkc\checkmarki\checkmarkr\checkmarkc\checkmark$\checkmarks\checkmark$\checkmark^\checkmark\\checkmarkc\checkmarki\checkmarkr\checkmarkc\checkmark$\checkmark \checkmark$\checkmark^\checkmark\\checkmarkc\checkmarki\checkmarkr\checkmarkc\checkmark$\checkmarkl\checkmark$\checkmark^\checkmark\\checkmarkc\checkmarki\checkmarkr\checkmarkc\checkmark$\checkmarki\checkmark$\checkmark^\checkmark\\checkmarkc\checkmarki\checkmarkr\checkmarkc\checkmark$\checkmarks\checkmark$\checkmark^\checkmark\\checkmarkc\checkmarki\checkmarkr\checkmarkc\checkmark$\checkmarks\checkmark$\checkmark^\checkmark\\checkmarkc\checkmarki\checkmarkr\checkmarkc\checkmark$\checkmarke\checkmark$\checkmark^\checkmark\\checkmarkc\checkmarki\checkmarkr\checkmarkc\checkmark$\checkmarks\checkmark$\checkmark^\checkmark\\checkmarkc\checkmarki\checkmarkr\checkmarkc\checkmark$\checkmark,\checkmark$\checkmark^\checkmark\\checkmarkc\checkmarki\checkmarkr\checkmarkc\checkmark$\checkmark \checkmark$\checkmark^\checkmark\\checkmarkc\checkmarki\checkmarkr\checkmarkc\checkmark$\checkmarkV\checkmark$\checkmark^\checkmark\\checkmarkc\checkmarki\checkmarkr\checkmarkc\checkmark$\checkmark-\checkmark$\checkmark^\checkmark\\checkmarkc\checkmarki\checkmarkr\checkmarkc\checkmark$\checkmarkw\checkmark$\checkmark^\checkmark\\checkmarkc\checkmarki\checkmarkr\checkmarkc\checkmark$\checkmarkh\checkmark$\checkmark^\checkmark\\checkmarkc\checkmarki\checkmarkr\checkmarkc\checkmark$\checkmarke\checkmark$\checkmark^\checkmark\\checkmarkc\checkmarki\checkmarkr\checkmarkc\checkmark$\checkmarke\checkmark$\checkmark^\checkmark\\checkmarkc\checkmarki\checkmarkr\checkmarkc\checkmark$\checkmarkl\checkmark$\checkmark^\checkmark\\checkmarkc\checkmarki\checkmarkr\checkmarkc\checkmark$\checkmarks\checkmark$\checkmark^\checkmark\\checkmarkc\checkmarki\checkmarkr\checkmarkc\checkmark$\checkmark)\checkmark$\checkmark^\checkmark\\checkmarkc\checkmarki\checkmarkr\checkmarkc\checkmark$\checkmark \checkmark$\checkmark^\checkmark\\checkmarkc\checkmarki\checkmarkr\checkmarkc\checkmark$\checkmark;\checkmark$\checkmark^\checkmark\\checkmarkc\checkmarki\checkmarkr\checkmarkc\checkmark$\checkmark
\checkmark$\checkmark^\checkmark\\checkmarkc\checkmarki\checkmarkr\checkmarkc\checkmark$\checkmark \checkmark$\checkmark^\checkmark\\checkmarkc\checkmarki\checkmarkr\checkmarkc\checkmark$\checkmark \checkmark$\checkmark^\checkmark\\checkmarkc\checkmarki\checkmarkr\checkmarkc\checkmark$\checkmark \checkmark$\checkmark^\checkmark\\checkmarkc\checkmarki\checkmarkr\checkmarkc\checkmark$\checkmark \checkmark$\checkmark^\checkmark\\checkmarkc\checkmarki\checkmarkr\checkmarkc\checkmark$\checkmark\\checkmark$\checkmark^\checkmark\\checkmarkc\checkmarki\checkmarkr\checkmarkc\checkmark$\checkmarki\checkmark$\checkmark^\checkmark\\checkmarkc\checkmarki\checkmarkr\checkmarkc\checkmark$\checkmarkt\checkmark$\checkmark^\checkmark\\checkmarkc\checkmarki\checkmarkr\checkmarkc\checkmark$\checkmarke\checkmark$\checkmark^\checkmark\\checkmarkc\checkmarki\checkmarkr\checkmarkc\checkmark$\checkmarkm\checkmark$\checkmark^\checkmark\\checkmarkc\checkmarki\checkmarkr\checkmarkc\checkmark$\checkmark \checkmark$\checkmark^\checkmark\\checkmarkc\checkmarki\checkmarkr\checkmarkc\checkmark$\checkmarkP\checkmark$\checkmark^\checkmark\\checkmarkc\checkmarki\checkmarkr\checkmarkc\checkmark$\checkmarke\checkmark$\checkmark^\checkmark\\checkmarkc\checkmarki\checkmarkr\checkmarkc\checkmark$\checkmarkr\checkmark$\checkmark^\checkmark\\checkmarkc\checkmarki\checkmarkr\checkmarkc\checkmark$\checkmarkm\checkmark$\checkmark^\checkmark\\checkmarkc\checkmarki\checkmarkr\checkmarkc\checkmark$\checkmarke\checkmark$\checkmark^\checkmark\\checkmarkc\checkmarki\checkmarkr\checkmarkc\checkmark$\checkmarkt\checkmark$\checkmark^\checkmark\\checkmarkc\checkmarki\checkmarkr\checkmarkc\checkmark$\checkmarkt\checkmark$\checkmark^\checkmark\\checkmarkc\checkmarki\checkmarkr\checkmarkc\checkmark$\checkmarkr\checkmark$\checkmark^\checkmark\\checkmarkc\checkmarki\checkmarkr\checkmarkc\checkmark$\checkmarke\checkmark$\checkmark^\checkmark\\checkmarkc\checkmarki\checkmarkr\checkmarkc\checkmark$\checkmark \checkmark$\checkmark^\checkmark\\checkmarkc\checkmarki\checkmarkr\checkmarkc\checkmark$\checkmarku\checkmark$\checkmark^\checkmark\\checkmarkc\checkmarki\checkmarkr\checkmarkc\checkmark$\checkmarkn\checkmark$\checkmark^\checkmark\\checkmarkc\checkmarki\checkmarkr\checkmarkc\checkmark$\checkmark \checkmark$\checkmark^\checkmark\\checkmarkc\checkmarki\checkmarkr\checkmarkc\checkmark$\checkmarkd\checkmark$\checkmark^\checkmark\\checkmarkc\checkmarki\checkmarkr\checkmarkc\checkmark$\checkmarké\checkmark$\checkmark^\checkmark\\checkmarkc\checkmarki\checkmarkr\checkmarkc\checkmark$\checkmarkp\checkmark$\checkmark^\checkmark\\checkmarkc\checkmarki\checkmarkr\checkmarkc\checkmark$\checkmarkl\checkmark$\checkmark^\checkmark\\checkmarkc\checkmarki\checkmarkr\checkmarkc\checkmark$\checkmarka\checkmark$\checkmark^\checkmark\\checkmarkc\checkmarki\checkmarkr\checkmarkc\checkmark$\checkmarkc\checkmark$\checkmark^\checkmark\\checkmarkc\checkmarki\checkmarkr\checkmarkc\checkmark$\checkmarke\checkmark$\checkmark^\checkmark\\checkmarkc\checkmarki\checkmarkr\checkmarkc\checkmark$\checkmarkm\checkmark$\checkmark^\checkmark\\checkmarkc\checkmarki\checkmarkr\checkmarkc\checkmark$\checkmarke\checkmark$\checkmark^\checkmark\\checkmarkc\checkmarki\checkmarkr\checkmarkc\checkmark$\checkmarkn\checkmark$\checkmark^\checkmark\\checkmarkc\checkmarki\checkmarkr\checkmarkc\checkmark$\checkmarkt\checkmark$\checkmark^\checkmark\\checkmarkc\checkmarki\checkmarkr\checkmarkc\checkmark$\checkmark \checkmark$\checkmark^\checkmark\\checkmarkc\checkmarki\checkmarkr\checkmarkc\checkmark$\checkmarks\checkmark$\checkmark^\checkmark\\checkmarkc\checkmarki\checkmarkr\checkmarkc\checkmark$\checkmarka\checkmark$\checkmark^\checkmark\\checkmarkc\checkmarki\checkmarkr\checkmarkc\checkmark$\checkmarkn\checkmark$\checkmark^\checkmark\\checkmarkc\checkmarki\checkmarkr\checkmarkc\checkmark$\checkmarks\checkmark$\checkmark^\checkmark\\checkmarkc\checkmarki\checkmarkr\checkmarkc\checkmark$\checkmark \checkmark$\checkmark^\checkmark\\checkmarkc\checkmarki\checkmarkr\checkmarkc\checkmark$\checkmarkj\checkmark$\checkmark^\checkmark\\checkmarkc\checkmarki\checkmarkr\checkmarkc\checkmark$\checkmarke\checkmark$\checkmark^\checkmark\\checkmarkc\checkmarki\checkmarkr\checkmarkc\checkmark$\checkmarku\checkmark$\checkmark^\checkmark\\checkmarkc\checkmarki\checkmarkr\checkmarkc\checkmark$\checkmark \checkmark$\checkmark^\checkmark\\checkmarkc\checkmarki\checkmarkr\checkmarkc\checkmark$\checkmarke\checkmark$\checkmark^\checkmark\\checkmarkc\checkmarki\checkmarkr\checkmarkc\checkmark$\checkmarkt\checkmark$\checkmark^\checkmark\\checkmarkc\checkmarki\checkmarkr\checkmarkc\checkmark$\checkmark \checkmark$\checkmark^\checkmark\\checkmarkc\checkmarki\checkmarkr\checkmarkc\checkmark$\checkmarks\checkmark$\checkmark^\checkmark\\checkmarkc\checkmarki\checkmarkr\checkmarkc\checkmark$\checkmarka\checkmark$\checkmark^\checkmark\\checkmarkc\checkmarki\checkmarkr\checkmarkc\checkmark$\checkmarkn\checkmark$\checkmark^\checkmark\\checkmarkc\checkmarki\checkmarkr\checkmarkc\checkmark$\checkmarks\checkmark$\checkmark^\checkmark\\checkmarkc\checkmarki\checkmarkr\checkmarkc\checkmark$\checkmark \checkmark$\checkmark^\checkmark\\checkmarkc\checkmarki\checkmarkr\checkmarkc\checkmark$\checkmarkf\checkmark$\checkmark^\checkmark\\checkmarkc\checkmarki\checkmarkr\checkmarkc\checkmark$\checkmarkr\checkmark$\checkmark^\checkmark\\checkmarkc\checkmarki\checkmarkr\checkmarkc\checkmark$\checkmarko\checkmark$\checkmark^\checkmark\\checkmarkc\checkmarki\checkmarkr\checkmarkc\checkmark$\checkmarkt\checkmark$\checkmark^\checkmark\\checkmarkc\checkmarki\checkmarkr\checkmarkc\checkmark$\checkmarkt\checkmark$\checkmark^\checkmark\\checkmarkc\checkmarki\checkmarkr\checkmarkc\checkmark$\checkmarke\checkmark$\checkmark^\checkmark\\checkmarkc\checkmarki\checkmarkr\checkmarkc\checkmark$\checkmarkm\checkmark$\checkmark^\checkmark\\checkmarkc\checkmarki\checkmarkr\checkmarkc\checkmark$\checkmarke\checkmark$\checkmark^\checkmark\\checkmarkc\checkmarki\checkmarkr\checkmarkc\checkmark$\checkmarkn\checkmark$\checkmark^\checkmark\\checkmarkc\checkmarki\checkmarkr\checkmarkc\checkmark$\checkmarkt\checkmark$\checkmark^\checkmark\\checkmarkc\checkmarki\checkmarkr\checkmarkc\checkmark$\checkmark \checkmark$\checkmark^\checkmark\\checkmarkc\checkmarki\checkmarkr\checkmarkc\checkmark$\checkmarks\checkmark$\checkmark^\checkmark\\checkmarkc\checkmarki\checkmarkr\checkmarkc\checkmark$\checkmarke\checkmark$\checkmark^\checkmark\\checkmarkc\checkmarki\checkmarkr\checkmarkc\checkmark$\checkmarkc\checkmark$\checkmark^\checkmark\\checkmarkc\checkmarki\checkmarkr\checkmarkc\checkmark$\checkmark \checkmark$\checkmark^\checkmark\\checkmarkc\checkmarki\checkmarkr\checkmarkc\checkmark$\checkmark;\checkmark$\checkmark^\checkmark\\checkmarkc\checkmarki\checkmarkr\checkmarkc\checkmark$\checkmark
\checkmark$\checkmark^\checkmark\\checkmarkc\checkmarki\checkmarkr\checkmarkc\checkmark$\checkmark \checkmark$\checkmark^\checkmark\\checkmarkc\checkmarki\checkmarkr\checkmarkc\checkmark$\checkmark \checkmark$\checkmark^\checkmark\\checkmarkc\checkmarki\checkmarkr\checkmarkc\checkmark$\checkmark \checkmark$\checkmark^\checkmark\\checkmarkc\checkmarki\checkmarkr\checkmarkc\checkmark$\checkmark \checkmark$\checkmark^\checkmark\\checkmarkc\checkmarki\checkmarkr\checkmarkc\checkmark$\checkmark\\checkmark$\checkmark^\checkmark\\checkmarkc\checkmarki\checkmarkr\checkmarkc\checkmark$\checkmarki\checkmark$\checkmark^\checkmark\\checkmarkc\checkmarki\checkmarkr\checkmarkc\checkmark$\checkmarkt\checkmark$\checkmark^\checkmark\\checkmarkc\checkmarki\checkmarkr\checkmarkc\checkmark$\checkmarke\checkmark$\checkmark^\checkmark\\checkmarkc\checkmarki\checkmarkr\checkmarkc\checkmark$\checkmarkm\checkmark$\checkmark^\checkmark\\checkmarkc\checkmarki\checkmarkr\checkmarkc\checkmark$\checkmark \checkmark$\checkmark^\checkmark\\checkmarkc\checkmarki\checkmarkr\checkmarkc\checkmark$\checkmarkÊ\checkmark$\checkmark^\checkmark\\checkmarkc\checkmarki\checkmarkr\checkmarkc\checkmark$\checkmarkt\checkmark$\checkmark^\checkmark\\checkmarkc\checkmarki\checkmarkr\checkmarkc\checkmark$\checkmarkr\checkmark$\checkmark^\checkmark\\checkmarkc\checkmarki\checkmarkr\checkmarkc\checkmark$\checkmarke\checkmark$\checkmark^\checkmark\\checkmarkc\checkmarki\checkmarkr\checkmarkc\checkmark$\checkmark \checkmark$\checkmark^\checkmark\\checkmarkc\checkmarki\checkmarkr\checkmarkc\checkmark$\checkmarkn\checkmark$\checkmark^\checkmark\\checkmarkc\checkmarki\checkmarkr\checkmarkc\checkmark$\checkmarko\checkmark$\checkmark^\checkmark\\checkmarkc\checkmarki\checkmarkr\checkmarkc\checkmark$\checkmarkr\checkmark$\checkmark^\checkmark\\checkmarkc\checkmarki\checkmarkr\checkmarkc\checkmark$\checkmarkm\checkmark$\checkmark^\checkmark\\checkmarkc\checkmarki\checkmarkr\checkmarkc\checkmark$\checkmarka\checkmark$\checkmark^\checkmark\\checkmarkc\checkmarki\checkmarkr\checkmarkc\checkmark$\checkmarkl\checkmark$\checkmark^\checkmark\\checkmarkc\checkmarki\checkmarkr\checkmarkc\checkmark$\checkmarki\checkmark$\checkmark^\checkmark\\checkmarkc\checkmarki\checkmarkr\checkmarkc\checkmark$\checkmarks\checkmark$\checkmark^\checkmark\\checkmarkc\checkmarki\checkmarkr\checkmarkc\checkmark$\checkmarké\checkmark$\checkmark^\checkmark\\checkmarkc\checkmarki\checkmarkr\checkmarkc\checkmark$\checkmarks\checkmark$\checkmark^\checkmark\\checkmarkc\checkmarki\checkmarkr\checkmarkc\checkmark$\checkmark \checkmark$\checkmark^\checkmark\\checkmarkc\checkmarki\checkmarkr\checkmarkc\checkmark$\checkmarke\checkmark$\checkmark^\checkmark\\checkmarkc\checkmarki\checkmarkr\checkmarkc\checkmark$\checkmarkt\checkmark$\checkmark^\checkmark\\checkmarkc\checkmarki\checkmarkr\checkmarkc\checkmark$\checkmark \checkmark$\checkmark^\checkmark\\checkmarkc\checkmarki\checkmarkr\checkmarkc\checkmark$\checkmarkd\checkmark$\checkmark^\checkmark\\checkmarkc\checkmarki\checkmarkr\checkmarkc\checkmark$\checkmarki\checkmark$\checkmark^\checkmark\\checkmarkc\checkmarki\checkmarkr\checkmarkc\checkmark$\checkmarks\checkmark$\checkmark^\checkmark\\checkmarkc\checkmarki\checkmarkr\checkmarkc\checkmark$\checkmarkp\checkmark$\checkmark^\checkmark\\checkmarkc\checkmarki\checkmarkr\checkmarkc\checkmark$\checkmarko\checkmark$\checkmark^\checkmark\\checkmarkc\checkmarki\checkmarkr\checkmarkc\checkmark$\checkmarkn\checkmark$\checkmark^\checkmark\\checkmarkc\checkmarki\checkmarkr\checkmarkc\checkmark$\checkmarki\checkmark$\checkmark^\checkmark\\checkmarkc\checkmarki\checkmarkr\checkmarkc\checkmark$\checkmarkb\checkmark$\checkmark^\checkmark\\checkmarkc\checkmarki\checkmarkr\checkmarkc\checkmark$\checkmarkl\checkmark$\checkmark^\checkmark\\checkmarkc\checkmarki\checkmarkr\checkmarkc\checkmark$\checkmarke\checkmark$\checkmark^\checkmark\\checkmarkc\checkmarki\checkmarkr\checkmarkc\checkmark$\checkmarks\checkmark$\checkmark^\checkmark\\checkmarkc\checkmarki\checkmarkr\checkmarkc\checkmark$\checkmark \checkmark$\checkmark^\checkmark\\checkmarkc\checkmarki\checkmarkr\checkmarkc\checkmark$\checkmarka\checkmark$\checkmark^\checkmark\\checkmarkc\checkmarki\checkmarkr\checkmarkc\checkmark$\checkmarku\checkmark$\checkmark^\checkmark\\checkmarkc\checkmarki\checkmarkr\checkmarkc\checkmark$\checkmark \checkmark$\checkmark^\checkmark\\checkmarkc\checkmarki\checkmarkr\checkmarkc\checkmark$\checkmarkc\checkmark$\checkmark^\checkmark\\checkmarkc\checkmarki\checkmarkr\checkmarkc\checkmark$\checkmarka\checkmark$\checkmark^\checkmark\\checkmarkc\checkmarki\checkmarkr\checkmarkc\checkmark$\checkmarkt\checkmark$\checkmark^\checkmark\\checkmarkc\checkmarki\checkmarkr\checkmarkc\checkmark$\checkmarka\checkmark$\checkmark^\checkmark\\checkmarkc\checkmarki\checkmarkr\checkmarkc\checkmark$\checkmarkl\checkmark$\checkmark^\checkmark\\checkmarkc\checkmarki\checkmarkr\checkmarkc\checkmark$\checkmarko\checkmark$\checkmark^\checkmark\\checkmarkc\checkmarki\checkmarkr\checkmarkc\checkmark$\checkmarkg\checkmark$\checkmark^\checkmark\\checkmarkc\checkmarki\checkmarkr\checkmarkc\checkmark$\checkmarku\checkmark$\checkmark^\checkmark\\checkmarkc\checkmarki\checkmarkr\checkmarkc\checkmark$\checkmarke\checkmark$\checkmark^\checkmark\\checkmarkc\checkmarki\checkmarkr\checkmarkc\checkmark$\checkmark.\checkmark$\checkmark^\checkmark\\checkmarkc\checkmarki\checkmarkr\checkmarkc\checkmark$\checkmark
\checkmark$\checkmark^\checkmark\\checkmarkc\checkmarki\checkmarkr\checkmarkc\checkmark$\checkmark\\checkmark$\checkmark^\checkmark\\checkmarkc\checkmarki\checkmarkr\checkmarkc\checkmark$\checkmarke\checkmark$\checkmark^\checkmark\\checkmarkc\checkmarki\checkmarkr\checkmarkc\checkmark$\checkmarkn\checkmark$\checkmark^\checkmark\\checkmarkc\checkmarki\checkmarkr\checkmarkc\checkmark$\checkmarkd\checkmark$\checkmark^\checkmark\\checkmarkc\checkmarki\checkmarkr\checkmarkc\checkmark$\checkmark{\checkmark$\checkmark^\checkmark\\checkmarkc\checkmarki\checkmarkr\checkmarkc\checkmark$\checkmarki\checkmark$\checkmark^\checkmark\\checkmarkc\checkmarki\checkmarkr\checkmarkc\checkmark$\checkmarkt\checkmark$\checkmark^\checkmark\\checkmarkc\checkmarki\checkmarkr\checkmarkc\checkmark$\checkmarke\checkmark$\checkmark^\checkmark\\checkmarkc\checkmarki\checkmarkr\checkmarkc\checkmark$\checkmarkm\checkmark$\checkmark^\checkmark\\checkmarkc\checkmarki\checkmarkr\checkmarkc\checkmark$\checkmarki\checkmark$\checkmark^\checkmark\\checkmarkc\checkmarki\checkmarkr\checkmarkc\checkmark$\checkmarkz\checkmark$\checkmark^\checkmark\\checkmarkc\checkmarki\checkmarkr\checkmarkc\checkmark$\checkmarke\checkmark$\checkmark^\checkmark\\checkmarkc\checkmarki\checkmarkr\checkmarkc\checkmark$\checkmark}\checkmark$\checkmark^\checkmark\\checkmarkc\checkmarki\checkmarkr\checkmarkc\checkmark$\checkmark
\checkmark$\checkmark^\checkmark\\checkmarkc\checkmarki\checkmarkr\checkmarkc\checkmark$\checkmark
\checkmark$\checkmark^\checkmark\\checkmarkc\checkmarki\checkmarkr\checkmarkc\checkmark$\checkmark\\checkmark$\checkmark^\checkmark\\checkmarkc\checkmarki\checkmarkr\checkmarkc\checkmark$\checkmarks\checkmark$\checkmark^\checkmark\\checkmarkc\checkmarki\checkmarkr\checkmarkc\checkmark$\checkmarku\checkmark$\checkmark^\checkmark\\checkmarkc\checkmarki\checkmarkr\checkmarkc\checkmark$\checkmarkb\checkmark$\checkmark^\checkmark\\checkmarkc\checkmarki\checkmarkr\checkmarkc\checkmark$\checkmarks\checkmark$\checkmark^\checkmark\\checkmarkc\checkmarki\checkmarkr\checkmarkc\checkmark$\checkmarku\checkmark$\checkmark^\checkmark\\checkmarkc\checkmarki\checkmarkr\checkmarkc\checkmark$\checkmarkb\checkmark$\checkmark^\checkmark\\checkmarkc\checkmarki\checkmarkr\checkmarkc\checkmark$\checkmarks\checkmark$\checkmark^\checkmark\\checkmarkc\checkmarki\checkmarkr\checkmarkc\checkmark$\checkmarke\checkmark$\checkmark^\checkmark\\checkmarkc\checkmarki\checkmarkr\checkmarkc\checkmark$\checkmarkc\checkmark$\checkmark^\checkmark\\checkmarkc\checkmarki\checkmarkr\checkmarkc\checkmark$\checkmarkt\checkmark$\checkmark^\checkmark\\checkmarkc\checkmarki\checkmarkr\checkmarkc\checkmark$\checkmarki\checkmark$\checkmark^\checkmark\\checkmarkc\checkmarki\checkmarkr\checkmarkc\checkmark$\checkmarko\checkmark$\checkmark^\checkmark\\checkmarkc\checkmarki\checkmarkr\checkmarkc\checkmark$\checkmarkn\checkmark$\checkmark^\checkmark\\checkmarkc\checkmarki\checkmarkr\checkmarkc\checkmark$\checkmark{\checkmark$\checkmark^\checkmark\\checkmarkc\checkmarki\checkmarkr\checkmarkc\checkmark$\checkmarkC\checkmark$\checkmark^\checkmark\\checkmarkc\checkmarki\checkmarkr\checkmarkc\checkmark$\checkmarkh\checkmark$\checkmark^\checkmark\\checkmarkc\checkmarki\checkmarkr\checkmarkc\checkmark$\checkmarko\checkmark$\checkmark^\checkmark\\checkmarkc\checkmarki\checkmarkr\checkmarkc\checkmark$\checkmarki\checkmark$\checkmark^\checkmark\\checkmarkc\checkmarki\checkmarkr\checkmarkc\checkmark$\checkmarkx\checkmark$\checkmark^\checkmark\\checkmarkc\checkmarki\checkmarkr\checkmarkc\checkmark$\checkmark \checkmark$\checkmark^\checkmark\\checkmarkc\checkmarki\checkmarkr\checkmarkc\checkmark$\checkmarkc\checkmark$\checkmark^\checkmark\\checkmarkc\checkmarki\checkmarkr\checkmarkc\checkmark$\checkmarko\checkmark$\checkmark^\checkmark\\checkmarkc\checkmarki\checkmarkr\checkmarkc\checkmark$\checkmarkm\checkmark$\checkmark^\checkmark\\checkmarkc\checkmarki\checkmarkr\checkmarkc\checkmark$\checkmarkp\checkmark$\checkmark^\checkmark\\checkmarkc\checkmarki\checkmarkr\checkmarkc\checkmark$\checkmarka\checkmark$\checkmark^\checkmark\\checkmarkc\checkmarki\checkmarkr\checkmarkc\checkmark$\checkmarkr\checkmark$\checkmark^\checkmark\\checkmarkc\checkmarki\checkmarkr\checkmarkc\checkmark$\checkmarka\checkmark$\checkmark^\checkmark\\checkmarkc\checkmarki\checkmarkr\checkmarkc\checkmark$\checkmarkt\checkmark$\checkmark^\checkmark\\checkmarkc\checkmarki\checkmarkr\checkmarkc\checkmark$\checkmarki\checkmark$\checkmark^\checkmark\\checkmarkc\checkmarki\checkmarkr\checkmarkc\checkmark$\checkmarkf\checkmark$\checkmark^\checkmark\\checkmarkc\checkmarki\checkmarkr\checkmarkc\checkmark$\checkmark}\checkmark$\checkmark^\checkmark\\checkmarkc\checkmarki\checkmarkr\checkmarkc\checkmark$\checkmark
\checkmark$\checkmark^\checkmark\\checkmarkc\checkmarki\checkmarkr\checkmarkc\checkmark$\checkmark\\checkmark$\checkmark^\checkmark\\checkmarkc\checkmarki\checkmarkr\checkmarkc\checkmark$\checkmarkb\checkmark$\checkmark^\checkmark\\checkmarkc\checkmarki\checkmarkr\checkmarkc\checkmark$\checkmarke\checkmark$\checkmark^\checkmark\\checkmarkc\checkmarki\checkmarkr\checkmarkc\checkmark$\checkmarkg\checkmark$\checkmark^\checkmark\\checkmarkc\checkmarki\checkmarkr\checkmarkc\checkmark$\checkmarki\checkmark$\checkmark^\checkmark\\checkmarkc\checkmarki\checkmarkr\checkmarkc\checkmark$\checkmarkn\checkmark$\checkmark^\checkmark\\checkmarkc\checkmarki\checkmarkr\checkmarkc\checkmark$\checkmark{\checkmark$\checkmark^\checkmark\\checkmarkc\checkmarki\checkmarkr\checkmarkc\checkmark$\checkmarkt\checkmark$\checkmark^\checkmark\\checkmarkc\checkmarki\checkmarkr\checkmarkc\checkmark$\checkmarka\checkmark$\checkmark^\checkmark\\checkmarkc\checkmarki\checkmarkr\checkmarkc\checkmark$\checkmarkb\checkmark$\checkmark^\checkmark\\checkmarkc\checkmarki\checkmarkr\checkmarkc\checkmark$\checkmarkl\checkmark$\checkmark^\checkmark\\checkmarkc\checkmarki\checkmarkr\checkmarkc\checkmark$\checkmarke\checkmark$\checkmark^\checkmark\\checkmarkc\checkmarki\checkmarkr\checkmarkc\checkmark$\checkmark}\checkmark$\checkmark^\checkmark\\checkmarkc\checkmarki\checkmarkr\checkmarkc\checkmark$\checkmark[\checkmark$\checkmark^\checkmark\\checkmarkc\checkmarki\checkmarkr\checkmarkc\checkmark$\checkmarkh\checkmark$\checkmark^\checkmark\\checkmarkc\checkmarki\checkmarkr\checkmarkc\checkmark$\checkmarkt\checkmark$\checkmark^\checkmark\\checkmarkc\checkmarki\checkmarkr\checkmarkc\checkmark$\checkmarkb\checkmark$\checkmark^\checkmark\\checkmarkc\checkmarki\checkmarkr\checkmarkc\checkmark$\checkmarkp\checkmark$\checkmark^\checkmark\\checkmarkc\checkmarki\checkmarkr\checkmarkc\checkmark$\checkmark]\checkmark$\checkmark^\checkmark\\checkmarkc\checkmarki\checkmarkr\checkmarkc\checkmark$\checkmark
\checkmark$\checkmark^\checkmark\\checkmarkc\checkmarki\checkmarkr\checkmarkc\checkmark$\checkmark\\checkmark$\checkmark^\checkmark\\checkmarkc\checkmarki\checkmarkr\checkmarkc\checkmark$\checkmarkc\checkmark$\checkmark^\checkmark\\checkmarkc\checkmarki\checkmarkr\checkmarkc\checkmark$\checkmarke\checkmark$\checkmark^\checkmark\\checkmarkc\checkmarki\checkmarkr\checkmarkc\checkmark$\checkmarkn\checkmark$\checkmark^\checkmark\\checkmarkc\checkmarki\checkmarkr\checkmarkc\checkmark$\checkmarkt\checkmark$\checkmark^\checkmark\\checkmarkc\checkmarki\checkmarkr\checkmarkc\checkmark$\checkmarke\checkmark$\checkmark^\checkmark\\checkmarkc\checkmarki\checkmarkr\checkmarkc\checkmark$\checkmarkr\checkmark$\checkmark^\checkmark\\checkmarkc\checkmarki\checkmarkr\checkmarkc\checkmark$\checkmarki\checkmark$\checkmark^\checkmark\\checkmarkc\checkmarki\checkmarkr\checkmarkc\checkmark$\checkmarkn\checkmark$\checkmark^\checkmark\\checkmarkc\checkmarki\checkmarkr\checkmarkc\checkmark$\checkmarkg\checkmark$\checkmark^\checkmark\\checkmarkc\checkmarki\checkmarkr\checkmarkc\checkmark$\checkmark
\checkmark$\checkmark^\checkmark\\checkmarkc\checkmarki\checkmarkr\checkmarkc\checkmark$\checkmark\\checkmark$\checkmark^\checkmark\\checkmarkc\checkmarki\checkmarkr\checkmarkc\checkmark$\checkmarkb\checkmark$\checkmark^\checkmark\\checkmarkc\checkmarki\checkmarkr\checkmarkc\checkmark$\checkmarke\checkmark$\checkmark^\checkmark\\checkmarkc\checkmarki\checkmarkr\checkmarkc\checkmark$\checkmarkg\checkmark$\checkmark^\checkmark\\checkmarkc\checkmarki\checkmarkr\checkmarkc\checkmark$\checkmarki\checkmark$\checkmark^\checkmark\\checkmarkc\checkmarki\checkmarkr\checkmarkc\checkmark$\checkmarkn\checkmark$\checkmark^\checkmark\\checkmarkc\checkmarki\checkmarkr\checkmarkc\checkmark$\checkmark{\checkmark$\checkmark^\checkmark\\checkmarkc\checkmarki\checkmarkr\checkmarkc\checkmark$\checkmarkt\checkmark$\checkmark^\checkmark\\checkmarkc\checkmarki\checkmarkr\checkmarkc\checkmark$\checkmarka\checkmark$\checkmark^\checkmark\\checkmarkc\checkmarki\checkmarkr\checkmarkc\checkmark$\checkmarkb\checkmark$\checkmark^\checkmark\\checkmarkc\checkmarki\checkmarkr\checkmarkc\checkmark$\checkmarku\checkmark$\checkmark^\checkmark\\checkmarkc\checkmarki\checkmarkr\checkmarkc\checkmark$\checkmarkl\checkmark$\checkmark^\checkmark\\checkmarkc\checkmarki\checkmarkr\checkmarkc\checkmark$\checkmarka\checkmark$\checkmark^\checkmark\\checkmarkc\checkmarki\checkmarkr\checkmarkc\checkmark$\checkmarkr\checkmark$\checkmark^\checkmark\\checkmarkc\checkmarki\checkmarkr\checkmarkc\checkmark$\checkmark}\checkmark$\checkmark^\checkmark\\checkmarkc\checkmarki\checkmarkr\checkmarkc\checkmark$\checkmark{\checkmark$\checkmark^\checkmark\\checkmarkc\checkmarki\checkmarkr\checkmarkc\checkmark$\checkmark|\checkmark$\checkmark^\checkmark\\checkmarkc\checkmarki\checkmarkr\checkmarkc\checkmark$\checkmarkl\checkmark$\checkmark^\checkmark\\checkmarkc\checkmarki\checkmarkr\checkmarkc\checkmark$\checkmark|\checkmark$\checkmark^\checkmark\\checkmarkc\checkmarki\checkmarkr\checkmarkc\checkmark$\checkmarkc\checkmark$\checkmark^\checkmark\\checkmarkc\checkmarki\checkmarkr\checkmarkc\checkmark$\checkmark|\checkmark$\checkmark^\checkmark\\checkmarkc\checkmarki\checkmarkr\checkmarkc\checkmark$\checkmarkc\checkmark$\checkmark^\checkmark\\checkmarkc\checkmarki\checkmarkr\checkmarkc\checkmark$\checkmark|\checkmark$\checkmark^\checkmark\\checkmarkc\checkmarki\checkmarkr\checkmarkc\checkmark$\checkmarkc\checkmark$\checkmark^\checkmark\\checkmarkc\checkmarki\checkmarkr\checkmarkc\checkmark$\checkmark|\checkmark$\checkmark^\checkmark\\checkmarkc\checkmarki\checkmarkr\checkmarkc\checkmark$\checkmarkc\checkmark$\checkmark^\checkmark\\checkmarkc\checkmarki\checkmarkr\checkmarkc\checkmark$\checkmark|\checkmark$\checkmark^\checkmark\\checkmarkc\checkmarki\checkmarkr\checkmarkc\checkmark$\checkmark}\checkmark$\checkmark^\checkmark\\checkmarkc\checkmarki\checkmarkr\checkmarkc\checkmark$\checkmark
\checkmark$\checkmark^\checkmark\\checkmarkc\checkmarki\checkmarkr\checkmarkc\checkmark$\checkmark\\checkmark$\checkmark^\checkmark\\checkmarkc\checkmarki\checkmarkr\checkmarkc\checkmark$\checkmarkh\checkmark$\checkmark^\checkmark\\checkmarkc\checkmarki\checkmarkr\checkmarkc\checkmark$\checkmarkl\checkmark$\checkmark^\checkmark\\checkmarkc\checkmarki\checkmarkr\checkmarkc\checkmark$\checkmarki\checkmark$\checkmark^\checkmark\\checkmarkc\checkmarki\checkmarkr\checkmarkc\checkmark$\checkmarkn\checkmark$\checkmark^\checkmark\\checkmarkc\checkmarki\checkmarkr\checkmarkc\checkmark$\checkmarke\checkmark$\checkmark^\checkmark\\checkmarkc\checkmarki\checkmarkr\checkmarkc\checkmark$\checkmark
\checkmark$\checkmark^\checkmark\\checkmarkc\checkmarki\checkmarkr\checkmarkc\checkmark$\checkmark\\checkmark$\checkmark^\checkmark\\checkmarkc\checkmarki\checkmarkr\checkmarkc\checkmark$\checkmarkt\checkmark$\checkmark^\checkmark\\checkmarkc\checkmarki\checkmarkr\checkmarkc\checkmark$\checkmarke\checkmark$\checkmark^\checkmark\\checkmarkc\checkmarki\checkmarkr\checkmarkc\checkmark$\checkmarkx\checkmark$\checkmark^\checkmark\\checkmarkc\checkmarki\checkmarkr\checkmarkc\checkmark$\checkmarkt\checkmark$\checkmark^\checkmark\\checkmarkc\checkmarki\checkmarkr\checkmarkc\checkmark$\checkmarkb\checkmark$\checkmark^\checkmark\\checkmarkc\checkmarki\checkmarkr\checkmarkc\checkmark$\checkmarkf\checkmark$\checkmark^\checkmark\\checkmarkc\checkmarki\checkmarkr\checkmarkc\checkmark$\checkmark{\checkmark$\checkmark^\checkmark\\checkmarkc\checkmarki\checkmarkr\checkmarkc\checkmark$\checkmarkT\checkmark$\checkmark^\checkmark\\checkmarkc\checkmarki\checkmarkr\checkmarkc\checkmark$\checkmarky\checkmark$\checkmark^\checkmark\\checkmarkc\checkmarki\checkmarkr\checkmarkc\checkmark$\checkmarkp\checkmark$\checkmark^\checkmark\\checkmarkc\checkmarki\checkmarkr\checkmarkc\checkmark$\checkmarke\checkmark$\checkmark^\checkmark\\checkmarkc\checkmarki\checkmarkr\checkmarkc\checkmark$\checkmark \checkmark$\checkmark^\checkmark\\checkmarkc\checkmarki\checkmarkr\checkmarkc\checkmark$\checkmarkd\checkmark$\checkmark^\checkmark\\checkmarkc\checkmarki\checkmarkr\checkmarkc\checkmark$\checkmarke\checkmark$\checkmark^\checkmark\\checkmarkc\checkmarki\checkmarkr\checkmarkc\checkmark$\checkmark \checkmark$\checkmark^\checkmark\\checkmarkc\checkmarki\checkmarkr\checkmarkc\checkmark$\checkmarkg\checkmark$\checkmark^\checkmark\\checkmarkc\checkmarki\checkmarkr\checkmarkc\checkmark$\checkmarku\checkmark$\checkmark^\checkmark\\checkmarkc\checkmarki\checkmarkr\checkmarkc\checkmark$\checkmarki\checkmark$\checkmark^\checkmark\\checkmarkc\checkmarki\checkmarkr\checkmarkc\checkmark$\checkmarkd\checkmark$\checkmark^\checkmark\\checkmarkc\checkmarki\checkmarkr\checkmarkc\checkmark$\checkmarka\checkmark$\checkmark^\checkmark\\checkmarkc\checkmarki\checkmarkr\checkmarkc\checkmark$\checkmarkg\checkmark$\checkmark^\checkmark\\checkmarkc\checkmarki\checkmarkr\checkmarkc\checkmark$\checkmarke\checkmark$\checkmark^\checkmark\\checkmarkc\checkmarki\checkmarkr\checkmarkc\checkmark$\checkmark}\checkmark$\checkmark^\checkmark\\checkmarkc\checkmarki\checkmarkr\checkmarkc\checkmark$\checkmark \checkmark$\checkmark^\checkmark\\checkmarkc\checkmarki\checkmarkr\checkmarkc\checkmark$\checkmark&\checkmark$\checkmark^\checkmark\\checkmarkc\checkmarki\checkmarkr\checkmarkc\checkmark$\checkmark \checkmark$\checkmark^\checkmark\\checkmarkc\checkmarki\checkmarkr\checkmarkc\checkmark$\checkmark\\checkmark$\checkmark^\checkmark\\checkmarkc\checkmarki\checkmarkr\checkmarkc\checkmark$\checkmarkt\checkmark$\checkmark^\checkmark\\checkmarkc\checkmarki\checkmarkr\checkmarkc\checkmark$\checkmarke\checkmark$\checkmark^\checkmark\\checkmarkc\checkmarki\checkmarkr\checkmarkc\checkmark$\checkmarkx\checkmark$\checkmark^\checkmark\\checkmarkc\checkmarki\checkmarkr\checkmarkc\checkmark$\checkmarkt\checkmark$\checkmark^\checkmark\\checkmarkc\checkmarki\checkmarkr\checkmarkc\checkmark$\checkmarkb\checkmark$\checkmark^\checkmark\\checkmarkc\checkmarki\checkmarkr\checkmarkc\checkmark$\checkmarkf\checkmark$\checkmark^\checkmark\\checkmarkc\checkmarki\checkmarkr\checkmarkc\checkmark$\checkmark{\checkmark$\checkmark^\checkmark\\checkmarkc\checkmarki\checkmarkr\checkmarkc\checkmark$\checkmarkR\checkmark$\checkmark^\checkmark\\checkmarkc\checkmarki\checkmarkr\checkmarkc\checkmark$\checkmarki\checkmark$\checkmark^\checkmark\\checkmarkc\checkmarki\checkmarkr\checkmarkc\checkmark$\checkmarkg\checkmark$\checkmark^\checkmark\\checkmarkc\checkmarki\checkmarkr\checkmarkc\checkmark$\checkmarki\checkmark$\checkmark^\checkmark\\checkmarkc\checkmarki\checkmarkr\checkmarkc\checkmark$\checkmarkd\checkmark$\checkmark^\checkmark\\checkmarkc\checkmarki\checkmarkr\checkmarkc\checkmark$\checkmarki\checkmark$\checkmark^\checkmark\\checkmarkc\checkmarki\checkmarkr\checkmarkc\checkmark$\checkmarkt\checkmark$\checkmark^\checkmark\\checkmarkc\checkmarki\checkmarkr\checkmarkc\checkmark$\checkmarké\checkmark$\checkmark^\checkmark\\checkmarkc\checkmarki\checkmarkr\checkmarkc\checkmark$\checkmark}\checkmark$\checkmark^\checkmark\\checkmarkc\checkmarki\checkmarkr\checkmarkc\checkmark$\checkmark \checkmark$\checkmark^\checkmark\\checkmarkc\checkmarki\checkmarkr\checkmarkc\checkmark$\checkmark&\checkmark$\checkmark^\checkmark\\checkmarkc\checkmarki\checkmarkr\checkmarkc\checkmark$\checkmark \checkmark$\checkmark^\checkmark\\checkmarkc\checkmarki\checkmarkr\checkmarkc\checkmark$\checkmark\\checkmark$\checkmark^\checkmark\\checkmarkc\checkmarki\checkmarkr\checkmarkc\checkmark$\checkmarkt\checkmark$\checkmark^\checkmark\\checkmarkc\checkmarki\checkmarkr\checkmarkc\checkmark$\checkmarke\checkmark$\checkmark^\checkmark\\checkmarkc\checkmarki\checkmarkr\checkmarkc\checkmark$\checkmarkx\checkmark$\checkmark^\checkmark\\checkmarkc\checkmarki\checkmarkr\checkmarkc\checkmark$\checkmarkt\checkmark$\checkmark^\checkmark\\checkmarkc\checkmarki\checkmarkr\checkmarkc\checkmark$\checkmarkb\checkmark$\checkmark^\checkmark\\checkmarkc\checkmarki\checkmarkr\checkmarkc\checkmark$\checkmarkf\checkmark$\checkmark^\checkmark\\checkmarkc\checkmarki\checkmarkr\checkmarkc\checkmark$\checkmark{\checkmark$\checkmark^\checkmark\\checkmarkc\checkmarki\checkmarkr\checkmarkc\checkmark$\checkmarkA\checkmark$\checkmark^\checkmark\\checkmarkc\checkmarki\checkmarkr\checkmarkc\checkmark$\checkmarkn\checkmark$\checkmark^\checkmark\\checkmarkc\checkmarki\checkmarkr\checkmarkc\checkmark$\checkmarkt\checkmark$\checkmark^\checkmark\\checkmarkc\checkmarki\checkmarkr\checkmarkc\checkmark$\checkmarki\checkmark$\checkmark^\checkmark\\checkmarkc\checkmarki\checkmarkr\checkmarkc\checkmark$\checkmark-\checkmark$\checkmark^\checkmark\\checkmarkc\checkmarki\checkmarkr\checkmarkc\checkmark$\checkmarkb\checkmark$\checkmark^\checkmark\\checkmarkc\checkmarki\checkmarkr\checkmarkc\checkmark$\checkmarka\checkmark$\checkmark^\checkmark\\checkmarkc\checkmarki\checkmarkr\checkmarkc\checkmark$\checkmarkc\checkmark$\checkmark^\checkmark\\checkmarkc\checkmarki\checkmarkr\checkmarkc\checkmark$\checkmarkk\checkmark$\checkmark^\checkmark\\checkmarkc\checkmarki\checkmarkr\checkmarkc\checkmark$\checkmarkl\checkmark$\checkmark^\checkmark\\checkmarkc\checkmarki\checkmarkr\checkmarkc\checkmark$\checkmarka\checkmark$\checkmark^\checkmark\\checkmarkc\checkmarki\checkmarkr\checkmarkc\checkmark$\checkmarks\checkmark$\checkmark^\checkmark\\checkmarkc\checkmarki\checkmarkr\checkmarkc\checkmark$\checkmarkh\checkmark$\checkmark^\checkmark\\checkmarkc\checkmarki\checkmarkr\checkmarkc\checkmark$\checkmark}\checkmark$\checkmark^\checkmark\\checkmarkc\checkmarki\checkmarkr\checkmarkc\checkmark$\checkmark \checkmark$\checkmark^\checkmark\\checkmarkc\checkmarki\checkmarkr\checkmarkc\checkmark$\checkmark&\checkmark$\checkmark^\checkmark\\checkmarkc\checkmarki\checkmarkr\checkmarkc\checkmark$\checkmark \checkmark$\checkmark^\checkmark\\checkmarkc\checkmarki\checkmarkr\checkmarkc\checkmark$\checkmark\\checkmark$\checkmark^\checkmark\\checkmarkc\checkmarki\checkmarkr\checkmarkc\checkmark$\checkmarkt\checkmark$\checkmark^\checkmark\\checkmarkc\checkmarki\checkmarkr\checkmarkc\checkmark$\checkmarke\checkmark$\checkmark^\checkmark\\checkmarkc\checkmarki\checkmarkr\checkmarkc\checkmark$\checkmarkx\checkmark$\checkmark^\checkmark\\checkmarkc\checkmarki\checkmarkr\checkmarkc\checkmark$\checkmarkt\checkmark$\checkmark^\checkmark\\checkmarkc\checkmarki\checkmarkr\checkmarkc\checkmark$\checkmarkb\checkmark$\checkmark^\checkmark\\checkmarkc\checkmarki\checkmarkr\checkmarkc\checkmark$\checkmarkf\checkmark$\checkmark^\checkmark\\checkmarkc\checkmarki\checkmarkr\checkmarkc\checkmark$\checkmark{\checkmark$\checkmark^\checkmark\\checkmarkc\checkmarki\checkmarkr\checkmarkc\checkmark$\checkmarkC\checkmark$\checkmark^\checkmark\\checkmarkc\checkmarki\checkmarkr\checkmarkc\checkmark$\checkmarko\checkmark$\checkmark^\checkmark\\checkmarkc\checkmarki\checkmarkr\checkmarkc\checkmark$\checkmarkû\checkmark$\checkmark^\checkmark\\checkmarkc\checkmarki\checkmarkr\checkmarkc\checkmark$\checkmarkt\checkmark$\checkmark^\checkmark\\checkmarkc\checkmarki\checkmarkr\checkmarkc\checkmark$\checkmark}\checkmark$\checkmark^\checkmark\\checkmarkc\checkmarki\checkmarkr\checkmarkc\checkmark$\checkmark \checkmark$\checkmark^\checkmark\\checkmarkc\checkmarki\checkmarkr\checkmarkc\checkmark$\checkmark&\checkmark$\checkmark^\checkmark\\checkmarkc\checkmarki\checkmarkr\checkmarkc\checkmark$\checkmark \checkmark$\checkmark^\checkmark\\checkmarkc\checkmarki\checkmarkr\checkmarkc\checkmark$\checkmark\\checkmark$\checkmark^\checkmark\\checkmarkc\checkmarki\checkmarkr\checkmarkc\checkmark$\checkmarkt\checkmark$\checkmark^\checkmark\\checkmarkc\checkmarki\checkmarkr\checkmarkc\checkmark$\checkmarke\checkmark$\checkmark^\checkmark\\checkmarkc\checkmarki\checkmarkr\checkmarkc\checkmark$\checkmarkx\checkmark$\checkmark^\checkmark\\checkmarkc\checkmarki\checkmarkr\checkmarkc\checkmark$\checkmarkt\checkmark$\checkmark^\checkmark\\checkmarkc\checkmarki\checkmarkr\checkmarkc\checkmark$\checkmarkb\checkmark$\checkmark^\checkmark\\checkmarkc\checkmarki\checkmarkr\checkmarkc\checkmark$\checkmarkf\checkmark$\checkmark^\checkmark\\checkmarkc\checkmarki\checkmarkr\checkmarkc\checkmark$\checkmark{\checkmark$\checkmark^\checkmark\\checkmarkc\checkmarki\checkmarkr\checkmarkc\checkmark$\checkmarkD\checkmark$\checkmark^\checkmark\\checkmarkc\checkmarki\checkmarkr\checkmarkc\checkmark$\checkmarké\checkmark$\checkmark^\checkmark\\checkmarkc\checkmarki\checkmarkr\checkmarkc\checkmark$\checkmarkc\checkmark$\checkmark^\checkmark\\checkmarkc\checkmarki\checkmarkr\checkmarkc\checkmark$\checkmarki\checkmark$\checkmark^\checkmark\\checkmarkc\checkmarki\checkmarkr\checkmarkc\checkmark$\checkmarks\checkmark$\checkmark^\checkmark\\checkmarkc\checkmarki\checkmarkr\checkmarkc\checkmark$\checkmarki\checkmark$\checkmark^\checkmark\\checkmarkc\checkmarki\checkmarkr\checkmarkc\checkmark$\checkmarko\checkmark$\checkmark^\checkmark\\checkmarkc\checkmarki\checkmarkr\checkmarkc\checkmark$\checkmarkn\checkmark$\checkmark^\checkmark\\checkmarkc\checkmarki\checkmarkr\checkmarkc\checkmark$\checkmark}\checkmark$\checkmark^\checkmark\\checkmarkc\checkmarki\checkmarkr\checkmarkc\checkmark$\checkmark \checkmark$\checkmark^\checkmark\\checkmarkc\checkmarki\checkmarkr\checkmarkc\checkmark$\checkmark\\checkmark$\checkmark^\checkmark\\checkmarkc\checkmarki\checkmarkr\checkmarkc\checkmark$\checkmark\\checkmark$\checkmark^\checkmark\\checkmarkc\checkmarki\checkmarkr\checkmarkc\checkmark$\checkmark \checkmark$\checkmark^\checkmark\\checkmarkc\checkmarki\checkmarkr\checkmarkc\checkmark$\checkmark\\checkmark$\checkmark^\checkmark\\checkmarkc\checkmarki\checkmarkr\checkmarkc\checkmark$\checkmarkh\checkmark$\checkmark^\checkmark\\checkmarkc\checkmarki\checkmarkr\checkmarkc\checkmark$\checkmarkl\checkmark$\checkmark^\checkmark\\checkmarkc\checkmarki\checkmarkr\checkmarkc\checkmark$\checkmarki\checkmark$\checkmark^\checkmark\\checkmarkc\checkmarki\checkmarkr\checkmarkc\checkmark$\checkmarkn\checkmark$\checkmark^\checkmark\\checkmarkc\checkmarki\checkmarkr\checkmarkc\checkmark$\checkmarke\checkmark$\checkmark^\checkmark\\checkmarkc\checkmarki\checkmarkr\checkmarkc\checkmark$\checkmark
\checkmark$\checkmark^\checkmark\\checkmarkc\checkmarki\checkmarkr\checkmarkc\checkmark$\checkmarkT\checkmark$\checkmark^\checkmark\\checkmarkc\checkmarki\checkmarkr\checkmarkc\checkmark$\checkmarki\checkmark$\checkmark^\checkmark\\checkmarkc\checkmarki\checkmarkr\checkmarkc\checkmark$\checkmarkg\checkmark$\checkmark^\checkmark\\checkmarkc\checkmarki\checkmarkr\checkmarkc\checkmark$\checkmarke\checkmark$\checkmark^\checkmark\\checkmarkc\checkmarki\checkmarkr\checkmarkc\checkmark$\checkmark \checkmark$\checkmark^\checkmark\\checkmarkc\checkmarki\checkmarkr\checkmarkc\checkmark$\checkmarkl\checkmark$\checkmark^\checkmark\\checkmarkc\checkmarki\checkmarkr\checkmarkc\checkmark$\checkmarki\checkmark$\checkmark^\checkmark\\checkmarkc\checkmarki\checkmarkr\checkmarkc\checkmark$\checkmarks\checkmark$\checkmark^\checkmark\\checkmarkc\checkmarki\checkmarkr\checkmarkc\checkmark$\checkmarks\checkmark$\checkmark^\checkmark\\checkmarkc\checkmarki\checkmarkr\checkmarkc\checkmark$\checkmarke\checkmark$\checkmark^\checkmark\\checkmarkc\checkmarki\checkmarkr\checkmarkc\checkmark$\checkmark \checkmark$\checkmark^\checkmark\\checkmarkc\checkmarki\checkmarkr\checkmarkc\checkmark$\checkmark+\checkmark$\checkmark^\checkmark\\checkmarkc\checkmarki\checkmarkr\checkmarkc\checkmark$\checkmark \checkmark$\checkmark^\checkmark\\checkmarkc\checkmarki\checkmarkr\checkmarkc\checkmark$\checkmarkb\checkmark$\checkmark^\checkmark\\checkmarkc\checkmarki\checkmarkr\checkmarkc\checkmark$\checkmarka\checkmark$\checkmark^\checkmark\\checkmarkc\checkmarki\checkmarkr\checkmarkc\checkmark$\checkmarkg\checkmark$\checkmark^\checkmark\\checkmarkc\checkmarki\checkmarkr\checkmarkc\checkmark$\checkmarku\checkmark$\checkmark^\checkmark\\checkmarkc\checkmarki\checkmarkr\checkmarkc\checkmark$\checkmarke\checkmark$\checkmark^\checkmark\\checkmarkc\checkmarki\checkmarkr\checkmarkc\checkmark$\checkmark \checkmark$\checkmark^\checkmark\\checkmarkc\checkmarki\checkmarkr\checkmarkc\checkmark$\checkmarkS\checkmark$\checkmark^\checkmark\\checkmarkc\checkmarki\checkmarkr\checkmarkc\checkmark$\checkmarkC\checkmark$\checkmark^\checkmark\\checkmarkc\checkmarki\checkmarkr\checkmarkc\checkmark$\checkmarkS\checkmark$\checkmark^\checkmark\\checkmarkc\checkmarki\checkmarkr\checkmarkc\checkmark$\checkmark \checkmark$\checkmark^\checkmark\\checkmarkc\checkmarki\checkmarkr\checkmarkc\checkmark$\checkmark&\checkmark$\checkmark^\checkmark\\checkmarkc\checkmarki\checkmarkr\checkmarkc\checkmark$\checkmark \checkmark$\checkmark^\checkmark\\checkmarkc\checkmarki\checkmarkr\checkmarkc\checkmark$\checkmarkF\checkmark$\checkmark^\checkmark\\checkmarkc\checkmarki\checkmarkr\checkmarkc\checkmark$\checkmarka\checkmark$\checkmark^\checkmark\\checkmarkc\checkmarki\checkmarkr\checkmarkc\checkmark$\checkmarki\checkmark$\checkmark^\checkmark\\checkmarkc\checkmarki\checkmarkr\checkmarkc\checkmark$\checkmarkb\checkmark$\checkmark^\checkmark\\checkmarkc\checkmarki\checkmarkr\checkmarkc\checkmark$\checkmarkl\checkmark$\checkmark^\checkmark\\checkmarkc\checkmarki\checkmarkr\checkmarkc\checkmark$\checkmarke\checkmark$\checkmark^\checkmark\\checkmarkc\checkmarki\checkmarkr\checkmarkc\checkmark$\checkmark \checkmark$\checkmark^\checkmark\\checkmarkc\checkmarki\checkmarkr\checkmarkc\checkmark$\checkmark&\checkmark$\checkmark^\checkmark\\checkmarkc\checkmarki\checkmarkr\checkmarkc\checkmark$\checkmark \checkmark$\checkmark^\checkmark\\checkmarkc\checkmarki\checkmarkr\checkmarkc\checkmark$\checkmarkN\checkmark$\checkmark^\checkmark\\checkmarkc\checkmarki\checkmarkr\checkmarkc\checkmark$\checkmarko\checkmark$\checkmark^\checkmark\\checkmarkc\checkmarki\checkmarkr\checkmarkc\checkmark$\checkmarkn\checkmark$\checkmark^\checkmark\\checkmarkc\checkmarki\checkmarkr\checkmarkc\checkmark$\checkmark \checkmark$\checkmark^\checkmark\\checkmarkc\checkmarki\checkmarkr\checkmarkc\checkmark$\checkmark&\checkmark$\checkmark^\checkmark\\checkmarkc\checkmarki\checkmarkr\checkmarkc\checkmark$\checkmark \checkmark$\checkmark^\checkmark\\checkmarkc\checkmarki\checkmarkr\checkmarkc\checkmark$\checkmarkB\checkmark$\checkmark^\checkmark\\checkmarkc\checkmarki\checkmarkr\checkmarkc\checkmark$\checkmarka\checkmark$\checkmark^\checkmark\\checkmarkc\checkmarki\checkmarkr\checkmarkc\checkmark$\checkmarks\checkmark$\checkmark^\checkmark\\checkmarkc\checkmarki\checkmarkr\checkmarkc\checkmark$\checkmark \checkmark$\checkmark^\checkmark\\checkmarkc\checkmarki\checkmarkr\checkmarkc\checkmark$\checkmark&\checkmark$\checkmark^\checkmark\\checkmarkc\checkmarki\checkmarkr\checkmarkc\checkmark$\checkmark \checkmark$\checkmark^\checkmark\\checkmarkc\checkmarki\checkmarkr\checkmarkc\checkmark$\checkmarkR\checkmark$\checkmark^\checkmark\\checkmarkc\checkmarki\checkmarkr\checkmarkc\checkmark$\checkmarke\checkmark$\checkmark^\checkmark\\checkmarkc\checkmarki\checkmarkr\checkmarkc\checkmark$\checkmarkj\checkmark$\checkmark^\checkmark\\checkmarkc\checkmarki\checkmarkr\checkmarkc\checkmark$\checkmarke\checkmark$\checkmark^\checkmark\\checkmarkc\checkmarki\checkmarkr\checkmarkc\checkmark$\checkmarkt\checkmark$\checkmark^\checkmark\\checkmarkc\checkmarki\checkmarkr\checkmarkc\checkmark$\checkmarké\checkmark$\checkmark^\checkmark\\checkmarkc\checkmarki\checkmarkr\checkmarkc\checkmark$\checkmark \checkmark$\checkmark^\checkmark\\checkmarkc\checkmarki\checkmarkr\checkmarkc\checkmark$\checkmark\\checkmark$\checkmark^\checkmark\\checkmarkc\checkmarki\checkmarkr\checkmarkc\checkmark$\checkmark\\checkmark$\checkmark^\checkmark\\checkmarkc\checkmarki\checkmarkr\checkmarkc\checkmark$\checkmark \checkmark$\checkmark^\checkmark\\checkmarkc\checkmarki\checkmarkr\checkmarkc\checkmark$\checkmark\\checkmark$\checkmark^\checkmark\\checkmarkc\checkmarki\checkmarkr\checkmarkc\checkmark$\checkmarkh\checkmark$\checkmark^\checkmark\\checkmarkc\checkmarki\checkmarkr\checkmarkc\checkmark$\checkmarkl\checkmark$\checkmark^\checkmark\\checkmarkc\checkmarki\checkmarkr\checkmarkc\checkmark$\checkmarki\checkmark$\checkmark^\checkmark\\checkmarkc\checkmarki\checkmarkr\checkmarkc\checkmark$\checkmarkn\checkmark$\checkmark^\checkmark\\checkmarkc\checkmarki\checkmarkr\checkmarkc\checkmark$\checkmarke\checkmark$\checkmark^\checkmark\\checkmarkc\checkmarki\checkmarkr\checkmarkc\checkmark$\checkmark
\checkmark$\checkmark^\checkmark\\checkmarkc\checkmarki\checkmarkr\checkmarkc\checkmark$\checkmarkV\checkmark$\checkmark^\checkmark\\checkmarkc\checkmarki\checkmarkr\checkmarkc\checkmark$\checkmark-\checkmark$\checkmark^\checkmark\\checkmarkc\checkmarki\checkmarkr\checkmarkc\checkmark$\checkmarkW\checkmark$\checkmark^\checkmark\\checkmarkc\checkmarki\checkmarkr\checkmarkc\checkmark$\checkmarkh\checkmark$\checkmark^\checkmark\\checkmarkc\checkmarki\checkmarkr\checkmarkc\checkmark$\checkmarke\checkmark$\checkmark^\checkmark\\checkmarkc\checkmarki\checkmarkr\checkmarkc\checkmark$\checkmarke\checkmark$\checkmark^\checkmark\\checkmarkc\checkmarki\checkmarkr\checkmarkc\checkmark$\checkmarkl\checkmark$\checkmark^\checkmark\\checkmarkc\checkmarki\checkmarkr\checkmarkc\checkmark$\checkmark \checkmark$\checkmark^\checkmark\\checkmarkc\checkmarki\checkmarkr\checkmarkc\checkmark$\checkmarkp\checkmark$\checkmark^\checkmark\\checkmarkc\checkmarki\checkmarkr\checkmarkc\checkmark$\checkmarko\checkmark$\checkmark^\checkmark\\checkmarkc\checkmarki\checkmarkr\checkmarkc\checkmark$\checkmarkl\checkmark$\checkmark^\checkmark\\checkmarkc\checkmarki\checkmarkr\checkmarkc\checkmark$\checkmarky\checkmark$\checkmark^\checkmark\\checkmarkc\checkmarki\checkmarkr\checkmarkc\checkmark$\checkmarkc\checkmark$\checkmark^\checkmark\\checkmarkc\checkmarki\checkmarkr\checkmarkc\checkmark$\checkmarka\checkmark$\checkmark^\checkmark\\checkmarkc\checkmarki\checkmarkr\checkmarkc\checkmark$\checkmarkr\checkmark$\checkmark^\checkmark\\checkmarkc\checkmarki\checkmarkr\checkmarkc\checkmark$\checkmarkb\checkmark$\checkmark^\checkmark\\checkmarkc\checkmarki\checkmarkr\checkmarkc\checkmark$\checkmarko\checkmark$\checkmark^\checkmark\\checkmarkc\checkmarki\checkmarkr\checkmarkc\checkmark$\checkmarkn\checkmark$\checkmark^\checkmark\\checkmarkc\checkmarki\checkmarkr\checkmarkc\checkmark$\checkmarka\checkmark$\checkmark^\checkmark\\checkmarkc\checkmarki\checkmarkr\checkmarkc\checkmark$\checkmarkt\checkmark$\checkmark^\checkmark\\checkmarkc\checkmarki\checkmarkr\checkmarkc\checkmark$\checkmarke\checkmark$\checkmark^\checkmark\\checkmarkc\checkmarki\checkmarkr\checkmarkc\checkmark$\checkmark \checkmark$\checkmark^\checkmark\\checkmarkc\checkmarki\checkmarkr\checkmarkc\checkmark$\checkmark&\checkmark$\checkmark^\checkmark\\checkmarkc\checkmarki\checkmarkr\checkmarkc\checkmark$\checkmark \checkmark$\checkmark^\checkmark\\checkmarkc\checkmarki\checkmarkr\checkmarkc\checkmark$\checkmarkT\checkmark$\checkmark^\checkmark\\checkmarkc\checkmarki\checkmarkr\checkmarkc\checkmark$\checkmarkr\checkmark$\checkmark^\checkmark\\checkmarkc\checkmarki\checkmarkr\checkmarkc\checkmark$\checkmarkè\checkmark$\checkmark^\checkmark\\checkmarkc\checkmarki\checkmarkr\checkmarkc\checkmark$\checkmarks\checkmark$\checkmark^\checkmark\\checkmarkc\checkmarki\checkmarkr\checkmarkc\checkmark$\checkmark \checkmark$\checkmark^\checkmark\\checkmarkc\checkmarki\checkmarkr\checkmarkc\checkmark$\checkmarkf\checkmark$\checkmark^\checkmark\\checkmarkc\checkmarki\checkmarkr\checkmarkc\checkmark$\checkmarka\checkmark$\checkmark^\checkmark\\checkmarkc\checkmarki\checkmarkr\checkmarkc\checkmark$\checkmarki\checkmark$\checkmark^\checkmark\\checkmarkc\checkmarki\checkmarkr\checkmarkc\checkmark$\checkmarkb\checkmark$\checkmark^\checkmark\\checkmarkc\checkmarki\checkmarkr\checkmarkc\checkmark$\checkmarkl\checkmark$\checkmark^\checkmark\\checkmarkc\checkmarki\checkmarkr\checkmarkc\checkmark$\checkmarke\checkmark$\checkmark^\checkmark\\checkmarkc\checkmarki\checkmarkr\checkmarkc\checkmark$\checkmark \checkmark$\checkmark^\checkmark\\checkmarkc\checkmarki\checkmarkr\checkmarkc\checkmark$\checkmark&\checkmark$\checkmark^\checkmark\\checkmarkc\checkmarki\checkmarkr\checkmarkc\checkmark$\checkmark \checkmark$\checkmark^\checkmark\\checkmarkc\checkmarki\checkmarkr\checkmarkc\checkmark$\checkmarkN\checkmark$\checkmark^\checkmark\\checkmarkc\checkmarki\checkmarkr\checkmarkc\checkmark$\checkmarko\checkmark$\checkmark^\checkmark\\checkmarkc\checkmarki\checkmarkr\checkmarkc\checkmark$\checkmarkn\checkmark$\checkmark^\checkmark\\checkmarkc\checkmarki\checkmarkr\checkmarkc\checkmark$\checkmark \checkmark$\checkmark^\checkmark\\checkmarkc\checkmarki\checkmarkr\checkmarkc\checkmark$\checkmark&\checkmark$\checkmark^\checkmark\\checkmarkc\checkmarki\checkmarkr\checkmarkc\checkmark$\checkmark \checkmark$\checkmark^\checkmark\\checkmarkc\checkmarki\checkmarkr\checkmarkc\checkmark$\checkmarkT\checkmark$\checkmark^\checkmark\\checkmarkc\checkmarki\checkmarkr\checkmarkc\checkmark$\checkmarkr\checkmark$\checkmark^\checkmark\\checkmarkc\checkmarki\checkmarkr\checkmarkc\checkmark$\checkmarkè\checkmark$\checkmark^\checkmark\\checkmarkc\checkmarki\checkmarkr\checkmarkc\checkmark$\checkmarks\checkmark$\checkmark^\checkmark\\checkmarkc\checkmarki\checkmarkr\checkmarkc\checkmark$\checkmark \checkmark$\checkmark^\checkmark\\checkmarkc\checkmarki\checkmarkr\checkmarkc\checkmark$\checkmarkb\checkmark$\checkmark^\checkmark\\checkmarkc\checkmarki\checkmarkr\checkmarkc\checkmark$\checkmarka\checkmark$\checkmark^\checkmark\\checkmarkc\checkmarki\checkmarkr\checkmarkc\checkmark$\checkmarks\checkmark$\checkmark^\checkmark\\checkmarkc\checkmarki\checkmarkr\checkmarkc\checkmark$\checkmark \checkmark$\checkmark^\checkmark\\checkmarkc\checkmarki\checkmarkr\checkmarkc\checkmark$\checkmark&\checkmark$\checkmark^\checkmark\\checkmarkc\checkmarki\checkmarkr\checkmarkc\checkmark$\checkmark \checkmark$\checkmark^\checkmark\\checkmarkc\checkmarki\checkmarkr\checkmarkc\checkmark$\checkmarkR\checkmark$\checkmark^\checkmark\\checkmarkc\checkmarki\checkmarkr\checkmarkc\checkmark$\checkmarke\checkmark$\checkmark^\checkmark\\checkmarkc\checkmarki\checkmarkr\checkmarkc\checkmark$\checkmarkj\checkmark$\checkmark^\checkmark\\checkmarkc\checkmarki\checkmarkr\checkmarkc\checkmark$\checkmarke\checkmark$\checkmark^\checkmark\\checkmarkc\checkmarki\checkmarkr\checkmarkc\checkmark$\checkmarkt\checkmark$\checkmark^\checkmark\\checkmarkc\checkmarki\checkmarkr\checkmarkc\checkmark$\checkmarké\checkmark$\checkmark^\checkmark\\checkmarkc\checkmarki\checkmarkr\checkmarkc\checkmark$\checkmark \checkmark$\checkmark^\checkmark\\checkmarkc\checkmarki\checkmarkr\checkmarkc\checkmark$\checkmark\\checkmark$\checkmark^\checkmark\\checkmarkc\checkmarki\checkmarkr\checkmarkc\checkmark$\checkmark\\checkmark$\checkmark^\checkmark\\checkmarkc\checkmarki\checkmarkr\checkmarkc\checkmark$\checkmark \checkmark$\checkmark^\checkmark\\checkmarkc\checkmarki\checkmarkr\checkmarkc\checkmark$\checkmark\\checkmark$\checkmark^\checkmark\\checkmarkc\checkmarki\checkmarkr\checkmarkc\checkmark$\checkmarkh\checkmark$\checkmark^\checkmark\\checkmarkc\checkmarki\checkmarkr\checkmarkc\checkmark$\checkmarkl\checkmark$\checkmark^\checkmark\\checkmarkc\checkmarki\checkmarkr\checkmarkc\checkmark$\checkmarki\checkmark$\checkmark^\checkmark\\checkmarkc\checkmarki\checkmarkr\checkmarkc\checkmark$\checkmarkn\checkmark$\checkmark^\checkmark\\checkmarkc\checkmarki\checkmarkr\checkmarkc\checkmark$\checkmarke\checkmark$\checkmark^\checkmark\\checkmarkc\checkmarki\checkmarkr\checkmarkc\checkmark$\checkmark
\checkmark$\checkmark^\checkmark\\checkmarkc\checkmarki\checkmarkr\checkmarkc\checkmark$\checkmark\\checkmark$\checkmark^\checkmark\\checkmarkc\checkmarki\checkmarkr\checkmarkc\checkmark$\checkmarkt\checkmark$\checkmark^\checkmark\\checkmarkc\checkmarki\checkmarkr\checkmarkc\checkmark$\checkmarke\checkmark$\checkmark^\checkmark\\checkmarkc\checkmarki\checkmarkr\checkmarkc\checkmark$\checkmarkx\checkmark$\checkmark^\checkmark\\checkmarkc\checkmarki\checkmarkr\checkmarkc\checkmark$\checkmarkt\checkmark$\checkmark^\checkmark\\checkmarkc\checkmarki\checkmarkr\checkmarkc\checkmark$\checkmarkb\checkmark$\checkmark^\checkmark\\checkmarkc\checkmarki\checkmarkr\checkmarkc\checkmark$\checkmarkf\checkmark$\checkmark^\checkmark\\checkmarkc\checkmarki\checkmarkr\checkmarkc\checkmark$\checkmark{\checkmark$\checkmark^\checkmark\\checkmarkc\checkmarki\checkmarkr\checkmarkc\checkmark$\checkmarkR\checkmark$\checkmark^\checkmark\\checkmarkc\checkmarki\checkmarkr\checkmarkc\checkmark$\checkmarka\checkmark$\checkmark^\checkmark\\checkmarkc\checkmarki\checkmarkr\checkmarkc\checkmark$\checkmarki\checkmark$\checkmark^\checkmark\\checkmarkc\checkmarki\checkmarkr\checkmarkc\checkmark$\checkmarkl\checkmark$\checkmark^\checkmark\\checkmarkc\checkmarki\checkmarkr\checkmarkc\checkmark$\checkmark \checkmark$\checkmark^\checkmark\\checkmarkc\checkmarki\checkmarkr\checkmarkc\checkmark$\checkmarkH\checkmark$\checkmark^\checkmark\\checkmarkc\checkmarki\checkmarkr\checkmarkc\checkmark$\checkmarkG\checkmark$\checkmark^\checkmark\\checkmarkc\checkmarki\checkmarkr\checkmarkc\checkmark$\checkmarkR\checkmark$\checkmark^\checkmark\\checkmarkc\checkmarki\checkmarkr\checkmarkc\checkmark$\checkmark1\checkmark$\checkmark^\checkmark\\checkmarkc\checkmarki\checkmarkr\checkmarkc\checkmark$\checkmark5\checkmark$\checkmark^\checkmark\\checkmarkc\checkmarki\checkmarkr\checkmarkc\checkmark$\checkmark \checkmark$\checkmark^\checkmark\\checkmarkc\checkmarki\checkmarkr\checkmarkc\checkmark$\checkmark(\checkmark$\checkmark^\checkmark\\checkmarkc\checkmarki\checkmarkr\checkmarkc\checkmark$\checkmarkH\checkmark$\checkmark^\checkmark\\checkmarkc\checkmarki\checkmarkr\checkmarkc\checkmark$\checkmarkI\checkmark$\checkmark^\checkmark\\checkmarkc\checkmarki\checkmarkr\checkmarkc\checkmark$\checkmarkW\checkmark$\checkmark^\checkmark\\checkmarkc\checkmarki\checkmarkr\checkmarkc\checkmark$\checkmarkI\checkmark$\checkmark^\checkmark\\checkmarkc\checkmarki\checkmarkr\checkmarkc\checkmark$\checkmarkN\checkmark$\checkmark^\checkmark\\checkmarkc\checkmarki\checkmarkr\checkmarkc\checkmark$\checkmark)\checkmark$\checkmark^\checkmark\\checkmarkc\checkmarki\checkmarkr\checkmarkc\checkmark$\checkmark}\checkmark$\checkmark^\checkmark\\checkmarkc\checkmarki\checkmarkr\checkmarkc\checkmark$\checkmark \checkmark$\checkmark^\checkmark\\checkmarkc\checkmarki\checkmarkr\checkmarkc\checkmark$\checkmark&\checkmark$\checkmark^\checkmark\\checkmarkc\checkmarki\checkmarkr\checkmarkc\checkmark$\checkmark \checkmark$\checkmark^\checkmark\\checkmarkc\checkmarki\checkmarkr\checkmarkc\checkmark$\checkmark\\checkmark$\checkmark^\checkmark\\checkmarkc\checkmarki\checkmarkr\checkmarkc\checkmark$\checkmarkt\checkmark$\checkmark^\checkmark\\checkmarkc\checkmarki\checkmarkr\checkmarkc\checkmark$\checkmarke\checkmark$\checkmark^\checkmark\\checkmarkc\checkmarki\checkmarkr\checkmarkc\checkmark$\checkmarkx\checkmark$\checkmark^\checkmark\\checkmarkc\checkmarki\checkmarkr\checkmarkc\checkmark$\checkmarkt\checkmark$\checkmark^\checkmark\\checkmarkc\checkmarki\checkmarkr\checkmarkc\checkmark$\checkmarkb\checkmark$\checkmark^\checkmark\\checkmarkc\checkmarki\checkmarkr\checkmarkc\checkmark$\checkmarkf\checkmark$\checkmark^\checkmark\\checkmarkc\checkmarki\checkmarkr\checkmarkc\checkmark$\checkmark{\checkmark$\checkmark^\checkmark\\checkmarkc\checkmarki\checkmarkr\checkmarkc\checkmark$\checkmarkÉ\checkmark$\checkmark^\checkmark\\checkmarkc\checkmarki\checkmarkr\checkmarkc\checkmark$\checkmarkl\checkmark$\checkmark^\checkmark\\checkmarkc\checkmarki\checkmarkr\checkmarkc\checkmark$\checkmarke\checkmark$\checkmark^\checkmark\\checkmarkc\checkmarki\checkmarkr\checkmarkc\checkmark$\checkmarkv\checkmark$\checkmark^\checkmark\\checkmarkc\checkmarki\checkmarkr\checkmarkc\checkmark$\checkmarké\checkmark$\checkmark^\checkmark\\checkmarkc\checkmarki\checkmarkr\checkmarkc\checkmark$\checkmarke\checkmark$\checkmark^\checkmark\\checkmarkc\checkmarki\checkmarkr\checkmarkc\checkmark$\checkmark}\checkmark$\checkmark^\checkmark\\checkmarkc\checkmarki\checkmarkr\checkmarkc\checkmark$\checkmark \checkmark$\checkmark^\checkmark\\checkmarkc\checkmarki\checkmarkr\checkmarkc\checkmark$\checkmark&\checkmark$\checkmark^\checkmark\\checkmarkc\checkmarki\checkmarkr\checkmarkc\checkmark$\checkmark \checkmark$\checkmark^\checkmark\\checkmarkc\checkmarki\checkmarkr\checkmarkc\checkmark$\checkmark\\checkmark$\checkmark^\checkmark\\checkmarkc\checkmarki\checkmarkr\checkmarkc\checkmark$\checkmarkt\checkmark$\checkmark^\checkmark\\checkmarkc\checkmarki\checkmarkr\checkmarkc\checkmark$\checkmarke\checkmark$\checkmark^\checkmark\\checkmarkc\checkmarki\checkmarkr\checkmarkc\checkmark$\checkmarkx\checkmark$\checkmark^\checkmark\\checkmarkc\checkmarki\checkmarkr\checkmarkc\checkmark$\checkmarkt\checkmark$\checkmark^\checkmark\\checkmarkc\checkmarki\checkmarkr\checkmarkc\checkmark$\checkmarkb\checkmark$\checkmark^\checkmark\\checkmarkc\checkmarki\checkmarkr\checkmarkc\checkmark$\checkmarkf\checkmark$\checkmark^\checkmark\\checkmarkc\checkmarki\checkmarkr\checkmarkc\checkmark$\checkmark{\checkmark$\checkmark^\checkmark\\checkmarkc\checkmarki\checkmarkr\checkmarkc\checkmark$\checkmarkO\checkmark$\checkmark^\checkmark\\checkmarkc\checkmarki\checkmarkr\checkmarkc\checkmark$\checkmarku\checkmark$\checkmark^\checkmark\\checkmarkc\checkmarki\checkmarkr\checkmarkc\checkmark$\checkmarki\checkmark$\checkmark^\checkmark\\checkmarkc\checkmarki\checkmarkr\checkmarkc\checkmark$\checkmark}\checkmark$\checkmark^\checkmark\\checkmarkc\checkmarki\checkmarkr\checkmarkc\checkmark$\checkmark \checkmark$\checkmark^\checkmark\\checkmarkc\checkmarki\checkmarkr\checkmarkc\checkmark$\checkmark&\checkmark$\checkmark^\checkmark\\checkmarkc\checkmarki\checkmarkr\checkmarkc\checkmark$\checkmark \checkmark$\checkmark^\checkmark\\checkmarkc\checkmarki\checkmarkr\checkmarkc\checkmark$\checkmark\\checkmark$\checkmark^\checkmark\\checkmarkc\checkmarki\checkmarkr\checkmarkc\checkmark$\checkmarkt\checkmark$\checkmark^\checkmark\\checkmarkc\checkmarki\checkmarkr\checkmarkc\checkmark$\checkmarke\checkmark$\checkmark^\checkmark\\checkmarkc\checkmarki\checkmarkr\checkmarkc\checkmark$\checkmarkx\checkmark$\checkmark^\checkmark\\checkmarkc\checkmarki\checkmarkr\checkmarkc\checkmark$\checkmarkt\checkmark$\checkmark^\checkmark\\checkmarkc\checkmarki\checkmarkr\checkmarkc\checkmark$\checkmarkb\checkmark$\checkmark^\checkmark\\checkmarkc\checkmarki\checkmarkr\checkmarkc\checkmark$\checkmarkf\checkmark$\checkmark^\checkmark\\checkmarkc\checkmarki\checkmarkr\checkmarkc\checkmark$\checkmark{\checkmark$\checkmark^\checkmark\\checkmarkc\checkmarki\checkmarkr\checkmarkc\checkmark$\checkmarkM\checkmark$\checkmark^\checkmark\\checkmarkc\checkmarki\checkmarkr\checkmarkc\checkmark$\checkmarko\checkmark$\checkmark^\checkmark\\checkmarkc\checkmarki\checkmarkr\checkmarkc\checkmark$\checkmarky\checkmark$\checkmark^\checkmark\\checkmarkc\checkmarki\checkmarkr\checkmarkc\checkmark$\checkmarke\checkmark$\checkmark^\checkmark\\checkmarkc\checkmarki\checkmarkr\checkmarkc\checkmark$\checkmarkn\checkmark$\checkmark^\checkmark\\checkmarkc\checkmarki\checkmarkr\checkmarkc\checkmark$\checkmark}\checkmark$\checkmark^\checkmark\\checkmarkc\checkmarki\checkmarkr\checkmarkc\checkmark$\checkmark \checkmark$\checkmark^\checkmark\\checkmarkc\checkmarki\checkmarkr\checkmarkc\checkmark$\checkmark&\checkmark$\checkmark^\checkmark\\checkmarkc\checkmarki\checkmarkr\checkmarkc\checkmark$\checkmark \checkmark$\checkmark^\checkmark\\checkmarkc\checkmarki\checkmarkr\checkmarkc\checkmark$\checkmark\\checkmark$\checkmark^\checkmark\\checkmarkc\checkmarki\checkmarkr\checkmarkc\checkmark$\checkmarkt\checkmark$\checkmark^\checkmark\\checkmarkc\checkmarki\checkmarkr\checkmarkc\checkmark$\checkmarke\checkmark$\checkmark^\checkmark\\checkmarkc\checkmarki\checkmarkr\checkmarkc\checkmark$\checkmarkx\checkmark$\checkmark^\checkmark\\checkmarkc\checkmarki\checkmarkr\checkmarkc\checkmark$\checkmarkt\checkmark$\checkmark^\checkmark\\checkmarkc\checkmarki\checkmarkr\checkmarkc\checkmark$\checkmarkb\checkmark$\checkmark^\checkmark\\checkmarkc\checkmarki\checkmarkr\checkmarkc\checkmark$\checkmarkf\checkmark$\checkmark^\checkmark\\checkmarkc\checkmarki\checkmarkr\checkmarkc\checkmark$\checkmark{\checkmark$\checkmark^\checkmark\\checkmarkc\checkmarki\checkmarkr\checkmarkc\checkmark$\checkmarkR\checkmark$\checkmark^\checkmark\\checkmarkc\checkmarki\checkmarkr\checkmarkc\checkmark$\checkmarke\checkmark$\checkmark^\checkmark\\checkmarkc\checkmarki\checkmarkr\checkmarkc\checkmark$\checkmarkt\checkmark$\checkmark^\checkmark\\checkmarkc\checkmarki\checkmarkr\checkmarkc\checkmark$\checkmarke\checkmark$\checkmark^\checkmark\\checkmarkc\checkmarki\checkmarkr\checkmarkc\checkmark$\checkmarkn\checkmark$\checkmark^\checkmark\\checkmarkc\checkmarki\checkmarkr\checkmarkc\checkmark$\checkmarku\checkmark$\checkmark^\checkmark\\checkmarkc\checkmarki\checkmarkr\checkmarkc\checkmark$\checkmark}\checkmark$\checkmark^\checkmark\\checkmarkc\checkmarki\checkmarkr\checkmarkc\checkmark$\checkmark \checkmark$\checkmark^\checkmark\\checkmarkc\checkmarki\checkmarkr\checkmarkc\checkmark$\checkmark\\checkmark$\checkmark^\checkmark\\checkmarkc\checkmarki\checkmarkr\checkmarkc\checkmark$\checkmark\\checkmark$\checkmark^\checkmark\\checkmarkc\checkmarki\checkmarkr\checkmarkc\checkmark$\checkmark \checkmark$\checkmark^\checkmark\\checkmarkc\checkmarki\checkmarkr\checkmarkc\checkmark$\checkmark\\checkmark$\checkmark^\checkmark\\checkmarkc\checkmarki\checkmarkr\checkmarkc\checkmark$\checkmarkh\checkmark$\checkmark^\checkmark\\checkmarkc\checkmarki\checkmarkr\checkmarkc\checkmark$\checkmarkl\checkmark$\checkmark^\checkmark\\checkmarkc\checkmarki\checkmarkr\checkmarkc\checkmark$\checkmarki\checkmark$\checkmark^\checkmark\\checkmarkc\checkmarki\checkmarkr\checkmarkc\checkmark$\checkmarkn\checkmark$\checkmark^\checkmark\\checkmarkc\checkmarki\checkmarkr\checkmarkc\checkmark$\checkmarke\checkmark$\checkmark^\checkmark\\checkmarkc\checkmarki\checkmarkr\checkmarkc\checkmark$\checkmark
\checkmark$\checkmark^\checkmark\\checkmarkc\checkmarki\checkmarkr\checkmarkc\checkmark$\checkmarkR\checkmark$\checkmark^\checkmark\\checkmarkc\checkmarki\checkmarkr\checkmarkc\checkmark$\checkmarka\checkmark$\checkmark^\checkmark\\checkmarkc\checkmarki\checkmarkr\checkmarkc\checkmark$\checkmarki\checkmark$\checkmark^\checkmark\\checkmarkc\checkmarki\checkmarkr\checkmarkc\checkmark$\checkmarkl\checkmark$\checkmark^\checkmark\\checkmarkc\checkmarki\checkmarkr\checkmarkc\checkmark$\checkmark \checkmark$\checkmark^\checkmark\\checkmarkc\checkmarki\checkmarkr\checkmarkc\checkmark$\checkmarkH\checkmark$\checkmark^\checkmark\\checkmarkc\checkmarki\checkmarkr\checkmarkc\checkmark$\checkmarkG\checkmark$\checkmark^\checkmark\\checkmarkc\checkmarki\checkmarkr\checkmarkc\checkmark$\checkmarkH\checkmark$\checkmark^\checkmark\\checkmarkc\checkmarki\checkmarkr\checkmarkc\checkmark$\checkmark2\checkmark$\checkmark^\checkmark\\checkmarkc\checkmarki\checkmarkr\checkmarkc\checkmark$\checkmark0\checkmark$\checkmark^\checkmark\\checkmarkc\checkmarki\checkmarkr\checkmarkc\checkmark$\checkmark \checkmark$\checkmark^\checkmark\\checkmarkc\checkmarki\checkmarkr\checkmarkc\checkmark$\checkmark&\checkmark$\checkmark^\checkmark\\checkmarkc\checkmarki\checkmarkr\checkmarkc\checkmark$\checkmark \checkmark$\checkmark^\checkmark\\checkmarkc\checkmarki\checkmarkr\checkmarkc\checkmark$\checkmarkT\checkmark$\checkmark^\checkmark\\checkmarkc\checkmarki\checkmarkr\checkmarkc\checkmark$\checkmarkr\checkmark$\checkmark^\checkmark\\checkmarkc\checkmarki\checkmarkr\checkmarkc\checkmark$\checkmarkè\checkmark$\checkmark^\checkmark\\checkmarkc\checkmarki\checkmarkr\checkmarkc\checkmark$\checkmarks\checkmark$\checkmark^\checkmark\\checkmarkc\checkmarki\checkmarkr\checkmarkc\checkmark$\checkmark \checkmark$\checkmark^\checkmark\\checkmarkc\checkmarki\checkmarkr\checkmarkc\checkmark$\checkmarké\checkmark$\checkmark^\checkmark\\checkmarkc\checkmarki\checkmarkr\checkmarkc\checkmark$\checkmarkl\checkmark$\checkmark^\checkmark\\checkmarkc\checkmarki\checkmarkr\checkmarkc\checkmark$\checkmarke\checkmark$\checkmark^\checkmark\\checkmarkc\checkmarki\checkmarkr\checkmarkc\checkmark$\checkmarkv\checkmark$\checkmark^\checkmark\\checkmarkc\checkmarki\checkmarkr\checkmarkc\checkmark$\checkmarké\checkmark$\checkmark^\checkmark\\checkmarkc\checkmarki\checkmarkr\checkmarkc\checkmark$\checkmarke\checkmark$\checkmark^\checkmark\\checkmarkc\checkmarki\checkmarkr\checkmarkc\checkmark$\checkmark \checkmark$\checkmark^\checkmark\\checkmarkc\checkmarki\checkmarkr\checkmarkc\checkmark$\checkmark&\checkmark$\checkmark^\checkmark\\checkmarkc\checkmarki\checkmarkr\checkmarkc\checkmark$\checkmark \checkmark$\checkmark^\checkmark\\checkmarkc\checkmarki\checkmarkr\checkmarkc\checkmark$\checkmarkO\checkmark$\checkmark^\checkmark\\checkmarkc\checkmarki\checkmarkr\checkmarkc\checkmark$\checkmarku\checkmark$\checkmark^\checkmark\\checkmarkc\checkmarki\checkmarkr\checkmarkc\checkmark$\checkmarki\checkmark$\checkmark^\checkmark\\checkmarkc\checkmarki\checkmarkr\checkmarkc\checkmark$\checkmark \checkmark$\checkmark^\checkmark\\checkmarkc\checkmarki\checkmarkr\checkmarkc\checkmark$\checkmark&\checkmark$\checkmark^\checkmark\\checkmarkc\checkmarki\checkmarkr\checkmarkc\checkmark$\checkmark \checkmark$\checkmark^\checkmark\\checkmarkc\checkmarki\checkmarkr\checkmarkc\checkmark$\checkmarkÉ\checkmark$\checkmark^\checkmark\\checkmarkc\checkmarki\checkmarkr\checkmarkc\checkmark$\checkmarkl\checkmark$\checkmark^\checkmark\\checkmarkc\checkmarki\checkmarkr\checkmarkc\checkmark$\checkmarke\checkmark$\checkmark^\checkmark\\checkmarkc\checkmarki\checkmarkr\checkmarkc\checkmark$\checkmarkv\checkmark$\checkmark^\checkmark\\checkmarkc\checkmarki\checkmarkr\checkmarkc\checkmark$\checkmarké\checkmark$\checkmark^\checkmark\\checkmarkc\checkmarki\checkmarkr\checkmarkc\checkmark$\checkmark \checkmark$\checkmark^\checkmark\\checkmarkc\checkmarki\checkmarkr\checkmarkc\checkmark$\checkmark&\checkmark$\checkmark^\checkmark\\checkmarkc\checkmarki\checkmarkr\checkmarkc\checkmark$\checkmark \checkmark$\checkmark^\checkmark\\checkmarkc\checkmarki\checkmarkr\checkmarkc\checkmark$\checkmarkS\checkmark$\checkmark^\checkmark\\checkmarkc\checkmarki\checkmarkr\checkmarkc\checkmark$\checkmarku\checkmark$\checkmark^\checkmark\\checkmarkc\checkmarki\checkmarkr\checkmarkc\checkmark$\checkmarkr\checkmark$\checkmark^\checkmark\\checkmarkc\checkmarki\checkmarkr\checkmarkc\checkmark$\checkmarkd\checkmark$\checkmark^\checkmark\\checkmarkc\checkmarki\checkmarkr\checkmarkc\checkmark$\checkmarki\checkmark$\checkmark^\checkmark\\checkmarkc\checkmarki\checkmarkr\checkmarkc\checkmark$\checkmarkm\checkmark$\checkmark^\checkmark\\checkmarkc\checkmarki\checkmarkr\checkmarkc\checkmark$\checkmarke\checkmark$\checkmark^\checkmark\\checkmarkc\checkmarki\checkmarkr\checkmarkc\checkmark$\checkmarkn\checkmark$\checkmark^\checkmark\\checkmarkc\checkmarki\checkmarkr\checkmarkc\checkmark$\checkmarks\checkmark$\checkmark^\checkmark\\checkmarkc\checkmarki\checkmarkr\checkmarkc\checkmark$\checkmarki\checkmark$\checkmark^\checkmark\\checkmarkc\checkmarki\checkmarkr\checkmarkc\checkmark$\checkmarko\checkmark$\checkmark^\checkmark\\checkmarkc\checkmarki\checkmarkr\checkmarkc\checkmark$\checkmarkn\checkmark$\checkmark^\checkmark\\checkmarkc\checkmarki\checkmarkr\checkmarkc\checkmark$\checkmarkn\checkmark$\checkmark^\checkmark\\checkmarkc\checkmarki\checkmarkr\checkmarkc\checkmark$\checkmarké\checkmark$\checkmark^\checkmark\\checkmarkc\checkmarki\checkmarkr\checkmarkc\checkmark$\checkmark \checkmark$\checkmark^\checkmark\\checkmarkc\checkmarki\checkmarkr\checkmarkc\checkmark$\checkmark\\checkmark$\checkmark^\checkmark\\checkmarkc\checkmarki\checkmarkr\checkmarkc\checkmark$\checkmark\\checkmark$\checkmark^\checkmark\\checkmarkc\checkmarki\checkmarkr\checkmarkc\checkmark$\checkmark \checkmark$\checkmark^\checkmark\\checkmarkc\checkmarki\checkmarkr\checkmarkc\checkmark$\checkmark\\checkmark$\checkmark^\checkmark\\checkmarkc\checkmarki\checkmarkr\checkmarkc\checkmark$\checkmarkh\checkmark$\checkmark^\checkmark\\checkmarkc\checkmarki\checkmarkr\checkmarkc\checkmark$\checkmarkl\checkmark$\checkmark^\checkmark\\checkmarkc\checkmarki\checkmarkr\checkmarkc\checkmark$\checkmarki\checkmark$\checkmark^\checkmark\\checkmarkc\checkmarki\checkmarkr\checkmarkc\checkmark$\checkmarkn\checkmark$\checkmark^\checkmark\\checkmarkc\checkmarki\checkmarkr\checkmarkc\checkmark$\checkmarke\checkmark$\checkmark^\checkmark\\checkmarkc\checkmarki\checkmarkr\checkmarkc\checkmark$\checkmark
\checkmark$\checkmark^\checkmark\\checkmarkc\checkmarki\checkmarkr\checkmarkc\checkmark$\checkmark\\checkmark$\checkmark^\checkmark\\checkmarkc\checkmarki\checkmarkr\checkmarkc\checkmark$\checkmarke\checkmark$\checkmark^\checkmark\\checkmarkc\checkmarki\checkmarkr\checkmarkc\checkmark$\checkmarkn\checkmark$\checkmark^\checkmark\\checkmarkc\checkmarki\checkmarkr\checkmarkc\checkmark$\checkmarkd\checkmark$\checkmark^\checkmark\\checkmarkc\checkmarki\checkmarkr\checkmarkc\checkmark$\checkmark{\checkmark$\checkmark^\checkmark\\checkmarkc\checkmarki\checkmarkr\checkmarkc\checkmark$\checkmarkt\checkmark$\checkmark^\checkmark\\checkmarkc\checkmarki\checkmarkr\checkmarkc\checkmark$\checkmarka\checkmark$\checkmark^\checkmark\\checkmarkc\checkmarki\checkmarkr\checkmarkc\checkmark$\checkmarkb\checkmark$\checkmark^\checkmark\\checkmarkc\checkmarki\checkmarkr\checkmarkc\checkmark$\checkmarku\checkmark$\checkmark^\checkmark\\checkmarkc\checkmarki\checkmarkr\checkmarkc\checkmark$\checkmarkl\checkmark$\checkmark^\checkmark\\checkmarkc\checkmarki\checkmarkr\checkmarkc\checkmark$\checkmarka\checkmark$\checkmark^\checkmark\\checkmarkc\checkmarki\checkmarkr\checkmarkc\checkmark$\checkmarkr\checkmark$\checkmark^\checkmark\\checkmarkc\checkmarki\checkmarkr\checkmarkc\checkmark$\checkmark}\checkmark$\checkmark^\checkmark\\checkmarkc\checkmarki\checkmarkr\checkmarkc\checkmark$\checkmark
\checkmark$\checkmark^\checkmark\\checkmarkc\checkmarki\checkmarkr\checkmarkc\checkmark$\checkmark\\checkmark$\checkmark^\checkmark\\checkmarkc\checkmarki\checkmarkr\checkmarkc\checkmark$\checkmarkc\checkmark$\checkmark^\checkmark\\checkmarkc\checkmarki\checkmarkr\checkmarkc\checkmark$\checkmarka\checkmark$\checkmark^\checkmark\\checkmarkc\checkmarki\checkmarkr\checkmarkc\checkmark$\checkmarkp\checkmark$\checkmark^\checkmark\\checkmarkc\checkmarki\checkmarkr\checkmarkc\checkmark$\checkmarkt\checkmark$\checkmark^\checkmark\\checkmarkc\checkmarki\checkmarkr\checkmarkc\checkmark$\checkmarki\checkmark$\checkmark^\checkmark\\checkmarkc\checkmarki\checkmarkr\checkmarkc\checkmark$\checkmarko\checkmark$\checkmark^\checkmark\\checkmarkc\checkmarki\checkmarkr\checkmarkc\checkmark$\checkmarkn\checkmark$\checkmark^\checkmark\\checkmarkc\checkmarki\checkmarkr\checkmarkc\checkmark$\checkmark{\checkmark$\checkmark^\checkmark\\checkmarkc\checkmarki\checkmarkr\checkmarkc\checkmark$\checkmarkT\checkmark$\checkmark^\checkmark\\checkmarkc\checkmarki\checkmarkr\checkmarkc\checkmark$\checkmarka\checkmark$\checkmark^\checkmark\\checkmarkc\checkmarki\checkmarkr\checkmarkc\checkmark$\checkmarkb\checkmark$\checkmark^\checkmark\\checkmarkc\checkmarki\checkmarkr\checkmarkc\checkmark$\checkmarkl\checkmark$\checkmark^\checkmark\\checkmarkc\checkmarki\checkmarkr\checkmarkc\checkmark$\checkmarke\checkmark$\checkmark^\checkmark\\checkmarkc\checkmarki\checkmarkr\checkmarkc\checkmark$\checkmarka\checkmark$\checkmark^\checkmark\\checkmarkc\checkmarki\checkmarkr\checkmarkc\checkmark$\checkmarku\checkmark$\checkmark^\checkmark\\checkmarkc\checkmarki\checkmarkr\checkmarkc\checkmark$\checkmark \checkmark$\checkmark^\checkmark\\checkmarkc\checkmarki\checkmarkr\checkmarkc\checkmark$\checkmarkc\checkmark$\checkmark^\checkmark\\checkmarkc\checkmarki\checkmarkr\checkmarkc\checkmark$\checkmarko\checkmark$\checkmark^\checkmark\\checkmarkc\checkmarki\checkmarkr\checkmarkc\checkmark$\checkmarkm\checkmark$\checkmark^\checkmark\\checkmarkc\checkmarki\checkmarkr\checkmarkc\checkmark$\checkmarkp\checkmark$\checkmark^\checkmark\\checkmarkc\checkmarki\checkmarkr\checkmarkc\checkmark$\checkmarka\checkmark$\checkmark^\checkmark\\checkmarkc\checkmarki\checkmarkr\checkmarkc\checkmark$\checkmarkr\checkmark$\checkmark^\checkmark\\checkmarkc\checkmarki\checkmarkr\checkmarkc\checkmark$\checkmarka\checkmark$\checkmark^\checkmark\\checkmarkc\checkmarki\checkmarkr\checkmarkc\checkmark$\checkmarkt\checkmark$\checkmark^\checkmark\\checkmarkc\checkmarki\checkmarkr\checkmarkc\checkmark$\checkmarki\checkmark$\checkmark^\checkmark\\checkmarkc\checkmarki\checkmarkr\checkmarkc\checkmark$\checkmarkf\checkmark$\checkmark^\checkmark\\checkmarkc\checkmarki\checkmarkr\checkmarkc\checkmark$\checkmark \checkmark$\checkmark^\checkmark\\checkmarkc\checkmarki\checkmarkr\checkmarkc\checkmark$\checkmarkd\checkmark$\checkmark^\checkmark\\checkmarkc\checkmarki\checkmarkr\checkmarkc\checkmark$\checkmarke\checkmark$\checkmark^\checkmark\\checkmarkc\checkmarki\checkmarkr\checkmarkc\checkmark$\checkmarks\checkmark$\checkmark^\checkmark\\checkmarkc\checkmarki\checkmarkr\checkmarkc\checkmark$\checkmark \checkmark$\checkmark^\checkmark\\checkmarkc\checkmarki\checkmarkr\checkmarkc\checkmark$\checkmarks\checkmark$\checkmark^\checkmark\\checkmarkc\checkmarki\checkmarkr\checkmarkc\checkmark$\checkmarko\checkmark$\checkmark^\checkmark\\checkmarkc\checkmarki\checkmarkr\checkmarkc\checkmark$\checkmarkl\checkmark$\checkmark^\checkmark\\checkmarkc\checkmarki\checkmarkr\checkmarkc\checkmark$\checkmarku\checkmark$\checkmark^\checkmark\\checkmarkc\checkmarki\checkmarkr\checkmarkc\checkmark$\checkmarkt\checkmark$\checkmark^\checkmark\\checkmarkc\checkmarki\checkmarkr\checkmarkc\checkmark$\checkmarki\checkmark$\checkmark^\checkmark\\checkmarkc\checkmarki\checkmarkr\checkmarkc\checkmark$\checkmarko\checkmark$\checkmark^\checkmark\\checkmarkc\checkmarki\checkmarkr\checkmarkc\checkmark$\checkmarkn\checkmark$\checkmark^\checkmark\\checkmarkc\checkmarki\checkmarkr\checkmarkc\checkmark$\checkmarks\checkmark$\checkmark^\checkmark\\checkmarkc\checkmarki\checkmarkr\checkmarkc\checkmark$\checkmark \checkmark$\checkmark^\checkmark\\checkmarkc\checkmarki\checkmarkr\checkmarkc\checkmark$\checkmarkd\checkmark$\checkmark^\checkmark\\checkmarkc\checkmarki\checkmarkr\checkmarkc\checkmark$\checkmarke\checkmark$\checkmark^\checkmark\\checkmarkc\checkmarki\checkmarkr\checkmarkc\checkmark$\checkmark \checkmark$\checkmark^\checkmark\\checkmarkc\checkmarki\checkmarkr\checkmarkc\checkmark$\checkmarkg\checkmark$\checkmark^\checkmark\\checkmarkc\checkmarki\checkmarkr\checkmarkc\checkmark$\checkmarku\checkmark$\checkmark^\checkmark\\checkmarkc\checkmarki\checkmarkr\checkmarkc\checkmark$\checkmarki\checkmark$\checkmark^\checkmark\\checkmarkc\checkmarki\checkmarkr\checkmarkc\checkmark$\checkmarkd\checkmark$\checkmark^\checkmark\\checkmarkc\checkmarki\checkmarkr\checkmarkc\checkmark$\checkmarka\checkmark$\checkmark^\checkmark\\checkmarkc\checkmarki\checkmarkr\checkmarkc\checkmark$\checkmarkg\checkmark$\checkmark^\checkmark\\checkmarkc\checkmarki\checkmarkr\checkmarkc\checkmark$\checkmarke\checkmark$\checkmark^\checkmark\\checkmarkc\checkmarki\checkmarkr\checkmarkc\checkmark$\checkmark \checkmark$\checkmark^\checkmark\\checkmarkc\checkmarki\checkmarkr\checkmarkc\checkmark$\checkmarkl\checkmark$\checkmark^\checkmark\\checkmarkc\checkmarki\checkmarkr\checkmarkc\checkmark$\checkmarki\checkmark$\checkmark^\checkmark\\checkmarkc\checkmarki\checkmarkr\checkmarkc\checkmark$\checkmarkn\checkmark$\checkmark^\checkmark\\checkmarkc\checkmarki\checkmarkr\checkmarkc\checkmark$\checkmarké\checkmark$\checkmark^\checkmark\\checkmarkc\checkmarki\checkmarkr\checkmarkc\checkmark$\checkmarka\checkmark$\checkmark^\checkmark\\checkmarkc\checkmarki\checkmarkr\checkmarkc\checkmark$\checkmarki\checkmark$\checkmark^\checkmark\\checkmarkc\checkmarki\checkmarkr\checkmarkc\checkmark$\checkmarkr\checkmark$\checkmark^\checkmark\\checkmarkc\checkmarki\checkmarkr\checkmarkc\checkmark$\checkmarke\checkmark$\checkmark^\checkmark\\checkmarkc\checkmarki\checkmarkr\checkmarkc\checkmark$\checkmark}\checkmark$\checkmark^\checkmark\\checkmarkc\checkmarki\checkmarkr\checkmarkc\checkmark$\checkmark
\checkmark$\checkmark^\checkmark\\checkmarkc\checkmarki\checkmarkr\checkmarkc\checkmark$\checkmark\\checkmark$\checkmark^\checkmark\\checkmarkc\checkmarki\checkmarkr\checkmarkc\checkmark$\checkmarkl\checkmark$\checkmark^\checkmark\\checkmarkc\checkmarki\checkmarkr\checkmarkc\checkmark$\checkmarka\checkmark$\checkmark^\checkmark\\checkmarkc\checkmarki\checkmarkr\checkmarkc\checkmark$\checkmarkb\checkmark$\checkmark^\checkmark\\checkmarkc\checkmarki\checkmarkr\checkmarkc\checkmark$\checkmarke\checkmark$\checkmark^\checkmark\\checkmarkc\checkmarki\checkmarkr\checkmarkc\checkmark$\checkmarkl\checkmark$\checkmark^\checkmark\\checkmarkc\checkmarki\checkmarkr\checkmarkc\checkmark$\checkmark{\checkmark$\checkmark^\checkmark\\checkmarkc\checkmarki\checkmarkr\checkmarkc\checkmark$\checkmarkt\checkmark$\checkmark^\checkmark\\checkmarkc\checkmarki\checkmarkr\checkmarkc\checkmark$\checkmarka\checkmark$\checkmark^\checkmark\\checkmarkc\checkmarki\checkmarkr\checkmarkc\checkmark$\checkmarkb\checkmark$\checkmark^\checkmark\\checkmarkc\checkmarki\checkmarkr\checkmarkc\checkmark$\checkmark:\checkmark$\checkmark^\checkmark\\checkmarkc\checkmarki\checkmarkr\checkmarkc\checkmark$\checkmarkc\checkmark$\checkmark^\checkmark\\checkmarkc\checkmarki\checkmarkr\checkmarkc\checkmark$\checkmarkh\checkmark$\checkmark^\checkmark\\checkmarkc\checkmarki\checkmarkr\checkmarkc\checkmark$\checkmarko\checkmark$\checkmark^\checkmark\\checkmarkc\checkmarki\checkmarkr\checkmarkc\checkmark$\checkmarki\checkmark$\checkmark^\checkmark\\checkmarkc\checkmarki\checkmarkr\checkmarkc\checkmark$\checkmarkx\checkmark$\checkmark^\checkmark\\checkmarkc\checkmarki\checkmarkr\checkmarkc\checkmark$\checkmark_\checkmark$\checkmark^\checkmark\\checkmarkc\checkmarki\checkmarkr\checkmarkc\checkmark$\checkmarkg\checkmark$\checkmark^\checkmark\\checkmarkc\checkmarki\checkmarkr\checkmarkc\checkmark$\checkmarku\checkmark$\checkmark^\checkmark\\checkmarkc\checkmarki\checkmarkr\checkmarkc\checkmark$\checkmarki\checkmark$\checkmark^\checkmark\\checkmarkc\checkmarki\checkmarkr\checkmarkc\checkmark$\checkmarkd\checkmark$\checkmark^\checkmark\\checkmarkc\checkmarki\checkmarkr\checkmarkc\checkmark$\checkmarka\checkmark$\checkmark^\checkmark\\checkmarkc\checkmarki\checkmarkr\checkmarkc\checkmark$\checkmarkg\checkmark$\checkmark^\checkmark\\checkmarkc\checkmarki\checkmarkr\checkmarkc\checkmark$\checkmarke\checkmark$\checkmark^\checkmark\\checkmarkc\checkmarki\checkmarkr\checkmarkc\checkmark$\checkmark}\checkmark$\checkmark^\checkmark\\checkmarkc\checkmarki\checkmarkr\checkmarkc\checkmark$\checkmark
\checkmark$\checkmark^\checkmark\\checkmarkc\checkmarki\checkmarkr\checkmarkc\checkmark$\checkmark\\checkmark$\checkmark^\checkmark\\checkmarkc\checkmarki\checkmarkr\checkmarkc\checkmark$\checkmarke\checkmark$\checkmark^\checkmark\\checkmarkc\checkmarki\checkmarkr\checkmarkc\checkmark$\checkmarkn\checkmark$\checkmark^\checkmark\\checkmarkc\checkmarki\checkmarkr\checkmarkc\checkmark$\checkmarkd\checkmark$\checkmark^\checkmark\\checkmarkc\checkmarki\checkmarkr\checkmarkc\checkmark$\checkmark{\checkmark$\checkmark^\checkmark\\checkmarkc\checkmarki\checkmarkr\checkmarkc\checkmark$\checkmarkt\checkmark$\checkmark^\checkmark\\checkmarkc\checkmarki\checkmarkr\checkmarkc\checkmark$\checkmarka\checkmark$\checkmark^\checkmark\\checkmarkc\checkmarki\checkmarkr\checkmarkc\checkmark$\checkmarkb\checkmark$\checkmark^\checkmark\\checkmarkc\checkmarki\checkmarkr\checkmarkc\checkmark$\checkmarkl\checkmark$\checkmark^\checkmark\\checkmarkc\checkmarki\checkmarkr\checkmarkc\checkmark$\checkmarke\checkmark$\checkmark^\checkmark\\checkmarkc\checkmarki\checkmarkr\checkmarkc\checkmark$\checkmark}\checkmark$\checkmark^\checkmark\\checkmarkc\checkmarki\checkmarkr\checkmarkc\checkmark$\checkmark
\checkmark$\checkmark^\checkmark\\checkmarkc\checkmarki\checkmarkr\checkmarkc\checkmark$\checkmark
\checkmark$\checkmark^\checkmark\\checkmarkc\checkmarki\checkmarkr\checkmarkc\checkmark$\checkmark\\checkmark$\checkmark^\checkmark\\checkmarkc\checkmarki\checkmarkr\checkmarkc\checkmark$\checkmarkt\checkmark$\checkmark^\checkmark\\checkmarkc\checkmarki\checkmarkr\checkmarkc\checkmark$\checkmarke\checkmark$\checkmark^\checkmark\\checkmarkc\checkmarki\checkmarkr\checkmarkc\checkmark$\checkmarkx\checkmark$\checkmark^\checkmark\\checkmarkc\checkmarki\checkmarkr\checkmarkc\checkmark$\checkmarkt\checkmark$\checkmark^\checkmark\\checkmarkc\checkmarki\checkmarkr\checkmarkc\checkmark$\checkmarkb\checkmark$\checkmark^\checkmark\\checkmarkc\checkmarki\checkmarkr\checkmarkc\checkmark$\checkmarkf\checkmark$\checkmark^\checkmark\\checkmarkc\checkmarki\checkmarkr\checkmarkc\checkmark$\checkmark{\checkmark$\checkmark^\checkmark\\checkmarkc\checkmarki\checkmarkr\checkmarkc\checkmark$\checkmarkG\checkmark$\checkmark^\checkmark\\checkmarkc\checkmarki\checkmarkr\checkmarkc\checkmark$\checkmarku\checkmark$\checkmark^\checkmark\\checkmarkc\checkmarki\checkmarkr\checkmarkc\checkmark$\checkmarki\checkmark$\checkmark^\checkmark\\checkmarkc\checkmarki\checkmarkr\checkmarkc\checkmark$\checkmarkd\checkmark$\checkmark^\checkmark\\checkmarkc\checkmarki\checkmarkr\checkmarkc\checkmark$\checkmarka\checkmark$\checkmark^\checkmark\\checkmarkc\checkmarki\checkmarkr\checkmarkc\checkmark$\checkmarkg\checkmark$\checkmark^\checkmark\\checkmarkc\checkmarki\checkmarkr\checkmarkc\checkmark$\checkmarke\checkmark$\checkmark^\checkmark\\checkmarkc\checkmarki\checkmarkr\checkmarkc\checkmark$\checkmark \checkmark$\checkmark^\checkmark\\checkmarkc\checkmarki\checkmarkr\checkmarkc\checkmark$\checkmarkr\checkmark$\checkmark^\checkmark\\checkmarkc\checkmarki\checkmarkr\checkmarkc\checkmark$\checkmarke\checkmark$\checkmark^\checkmark\\checkmarkc\checkmarki\checkmarkr\checkmarkc\checkmark$\checkmarkt\checkmark$\checkmark^\checkmark\\checkmarkc\checkmarki\checkmarkr\checkmarkc\checkmark$\checkmarke\checkmark$\checkmark^\checkmark\\checkmarkc\checkmarki\checkmarkr\checkmarkc\checkmark$\checkmarkn\checkmark$\checkmark^\checkmark\\checkmarkc\checkmarki\checkmarkr\checkmarkc\checkmark$\checkmarku\checkmark$\checkmark^\checkmark\\checkmarkc\checkmarki\checkmarkr\checkmarkc\checkmark$\checkmark \checkmark$\checkmark^\checkmark\\checkmarkc\checkmarki\checkmarkr\checkmarkc\checkmark$\checkmark:\checkmark$\checkmark^\checkmark\\checkmarkc\checkmarki\checkmarkr\checkmarkc\checkmark$\checkmark \checkmark$\checkmark^\checkmark\\checkmarkc\checkmarki\checkmarkr\checkmarkc\checkmark$\checkmarkR\checkmark$\checkmark^\checkmark\\checkmarkc\checkmarki\checkmarkr\checkmarkc\checkmark$\checkmarka\checkmark$\checkmark^\checkmark\\checkmarkc\checkmarki\checkmarkr\checkmarkc\checkmark$\checkmarki\checkmark$\checkmark^\checkmark\\checkmarkc\checkmarki\checkmarkr\checkmarkc\checkmark$\checkmarkl\checkmark$\checkmark^\checkmark\\checkmarkc\checkmarki\checkmarkr\checkmarkc\checkmark$\checkmark \checkmark$\checkmark^\checkmark\\checkmarkc\checkmarki\checkmarkr\checkmarkc\checkmark$\checkmarkH\checkmark$\checkmark^\checkmark\\checkmarkc\checkmarki\checkmarkr\checkmarkc\checkmark$\checkmarkI\checkmark$\checkmark^\checkmark\\checkmarkc\checkmarki\checkmarkr\checkmarkc\checkmark$\checkmarkW\checkmark$\checkmark^\checkmark\\checkmarkc\checkmarki\checkmarkr\checkmarkc\checkmark$\checkmarkI\checkmark$\checkmark^\checkmark\\checkmarkc\checkmarki\checkmarkr\checkmarkc\checkmark$\checkmarkN\checkmark$\checkmark^\checkmark\\checkmarkc\checkmarki\checkmarkr\checkmarkc\checkmark$\checkmark \checkmark$\checkmark^\checkmark\\checkmarkc\checkmarki\checkmarkr\checkmarkc\checkmark$\checkmarkH\checkmark$\checkmark^\checkmark\\checkmarkc\checkmarki\checkmarkr\checkmarkc\checkmark$\checkmarkG\checkmark$\checkmark^\checkmark\\checkmarkc\checkmarki\checkmarkr\checkmarkc\checkmark$\checkmarkR\checkmark$\checkmark^\checkmark\\checkmarkc\checkmarki\checkmarkr\checkmarkc\checkmark$\checkmark1\checkmark$\checkmark^\checkmark\\checkmarkc\checkmarki\checkmarkr\checkmarkc\checkmark$\checkmark5\checkmark$\checkmark^\checkmark\\checkmarkc\checkmarki\checkmarkr\checkmarkc\checkmark$\checkmark}\checkmark$\checkmark^\checkmark\\checkmarkc\checkmarki\checkmarkr\checkmarkc\checkmark$\checkmark \checkmark$\checkmark^\checkmark\\checkmarkc\checkmarki\checkmarkr\checkmarkc\checkmark$\checkmark(\checkmark$\checkmark^\checkmark\\checkmarkc\checkmarki\checkmarkr\checkmarkc\checkmark$\checkmark4\checkmark$\checkmark^\checkmark\\checkmarkc\checkmarki\checkmarkr\checkmarkc\checkmark$\checkmark \checkmark$\checkmark^\checkmark\\checkmarkc\checkmarki\checkmarkr\checkmarkc\checkmark$\checkmarkr\checkmark$\checkmark^\checkmark\\checkmarkc\checkmarki\checkmarkr\checkmarkc\checkmark$\checkmarka\checkmark$\checkmark^\checkmark\\checkmarkc\checkmarki\checkmarkr\checkmarkc\checkmark$\checkmarkn\checkmark$\checkmark^\checkmark\\checkmarkc\checkmarki\checkmarkr\checkmarkc\checkmark$\checkmarkg\checkmark$\checkmark^\checkmark\\checkmarkc\checkmarki\checkmarkr\checkmarkc\checkmark$\checkmarké\checkmark$\checkmark^\checkmark\\checkmarkc\checkmarki\checkmarkr\checkmarkc\checkmark$\checkmarke\checkmark$\checkmark^\checkmark\\checkmarkc\checkmarki\checkmarkr\checkmarkc\checkmark$\checkmarks\checkmark$\checkmark^\checkmark\\checkmarkc\checkmarki\checkmarkr\checkmarkc\checkmark$\checkmark \checkmark$\checkmark^\checkmark\\checkmarkc\checkmarki\checkmarkr\checkmarkc\checkmark$\checkmarkd\checkmark$\checkmark^\checkmark\\checkmarkc\checkmarki\checkmarkr\checkmarkc\checkmark$\checkmarke\checkmark$\checkmark^\checkmark\\checkmarkc\checkmarki\checkmarkr\checkmarkc\checkmark$\checkmark \checkmark$\checkmark^\checkmark\\checkmarkc\checkmarki\checkmarkr\checkmarkc\checkmark$\checkmarkb\checkmark$\checkmark^\checkmark\\checkmarkc\checkmarki\checkmarkr\checkmarkc\checkmark$\checkmarki\checkmark$\checkmark^\checkmark\\checkmarkc\checkmarki\checkmarkr\checkmarkc\checkmark$\checkmarkl\checkmark$\checkmark^\checkmark\\checkmarkc\checkmarki\checkmarkr\checkmarkc\checkmark$\checkmarkl\checkmark$\checkmark^\checkmark\\checkmarkc\checkmarki\checkmarkr\checkmarkc\checkmark$\checkmarke\checkmark$\checkmark^\checkmark\\checkmarkc\checkmarki\checkmarkr\checkmarkc\checkmark$\checkmarks\checkmark$\checkmark^\checkmark\\checkmarkc\checkmarki\checkmarkr\checkmarkc\checkmark$\checkmark,\checkmark$\checkmark^\checkmark\\checkmarkc\checkmarki\checkmarkr\checkmarkc\checkmark$\checkmark \checkmark$\checkmark^\checkmark\\checkmarkc\checkmarki\checkmarkr\checkmarkc\checkmark$\checkmarkc\checkmark$\checkmark^\checkmark\\checkmarkc\checkmarki\checkmarkr\checkmarkc\checkmark$\checkmarko\checkmark$\checkmark^\checkmark\\checkmarkc\checkmarki\checkmarkr\checkmarkc\checkmark$\checkmarkn\checkmark$\checkmark^\checkmark\\checkmarkc\checkmarki\checkmarkr\checkmarkc\checkmark$\checkmarkt\checkmark$\checkmark^\checkmark\\checkmarkc\checkmarki\checkmarkr\checkmarkc\checkmark$\checkmarka\checkmark$\checkmark^\checkmark\\checkmarkc\checkmarki\checkmarkr\checkmarkc\checkmark$\checkmarkc\checkmark$\checkmark^\checkmark\\checkmarkc\checkmarki\checkmarkr\checkmarkc\checkmark$\checkmarkt\checkmark$\checkmark^\checkmark\\checkmarkc\checkmarki\checkmarkr\checkmarkc\checkmark$\checkmark \checkmark$\checkmark^\checkmark\\checkmarkc\checkmarki\checkmarkr\checkmarkc\checkmark$\checkmark4\checkmark$\checkmark^\checkmark\\checkmarkc\checkmarki\checkmarkr\checkmarkc\checkmark$\checkmark \checkmark$\checkmark^\checkmark\\checkmarkc\checkmarki\checkmarkr\checkmarkc\checkmark$\checkmarkp\checkmark$\checkmark^\checkmark\\checkmarkc\checkmarki\checkmarkr\checkmarkc\checkmark$\checkmarko\checkmark$\checkmark^\checkmark\\checkmarkc\checkmarki\checkmarkr\checkmarkc\checkmark$\checkmarki\checkmark$\checkmark^\checkmark\\checkmarkc\checkmarki\checkmarkr\checkmarkc\checkmark$\checkmarkn\checkmark$\checkmark^\checkmark\\checkmarkc\checkmarki\checkmarkr\checkmarkc\checkmark$\checkmarkt\checkmark$\checkmark^\checkmark\\checkmarkc\checkmarki\checkmarkr\checkmarkc\checkmark$\checkmarks\checkmark$\checkmark^\checkmark\\checkmarkc\checkmarki\checkmarkr\checkmarkc\checkmark$\checkmark)\checkmark$\checkmark^\checkmark\\checkmarkc\checkmarki\checkmarkr\checkmarkc\checkmark$\checkmark
\checkmark$\checkmark^\checkmark\\checkmarkc\checkmarki\checkmarkr\checkmarkc\checkmark$\checkmark\\checkmark$\checkmark^\checkmark\\checkmarkc\checkmarki\checkmarkr\checkmarkc\checkmark$\checkmarkb\checkmark$\checkmark^\checkmark\\checkmarkc\checkmarki\checkmarkr\checkmarkc\checkmark$\checkmarke\checkmark$\checkmark^\checkmark\\checkmarkc\checkmarki\checkmarkr\checkmarkc\checkmark$\checkmarkg\checkmark$\checkmark^\checkmark\\checkmarkc\checkmarki\checkmarkr\checkmarkc\checkmark$\checkmarki\checkmark$\checkmark^\checkmark\\checkmarkc\checkmarki\checkmarkr\checkmarkc\checkmark$\checkmarkn\checkmark$\checkmark^\checkmark\\checkmarkc\checkmarki\checkmarkr\checkmarkc\checkmark$\checkmark{\checkmark$\checkmark^\checkmark\\checkmarkc\checkmarki\checkmarkr\checkmarkc\checkmark$\checkmarki\checkmark$\checkmark^\checkmark\\checkmarkc\checkmarki\checkmarkr\checkmarkc\checkmark$\checkmarkt\checkmark$\checkmark^\checkmark\\checkmarkc\checkmarki\checkmarkr\checkmarkc\checkmark$\checkmarke\checkmark$\checkmark^\checkmark\\checkmarkc\checkmarki\checkmarkr\checkmarkc\checkmark$\checkmarkm\checkmark$\checkmark^\checkmark\\checkmarkc\checkmarki\checkmarkr\checkmarkc\checkmark$\checkmarki\checkmark$\checkmark^\checkmark\\checkmarkc\checkmarki\checkmarkr\checkmarkc\checkmark$\checkmarkz\checkmark$\checkmark^\checkmark\\checkmarkc\checkmarki\checkmarkr\checkmarkc\checkmark$\checkmarke\checkmark$\checkmark^\checkmark\\checkmarkc\checkmarki\checkmarkr\checkmarkc\checkmark$\checkmark}\checkmark$\checkmark^\checkmark\\checkmarkc\checkmarki\checkmarkr\checkmarkc\checkmark$\checkmark
\checkmark$\checkmark^\checkmark\\checkmarkc\checkmarki\checkmarkr\checkmarkc\checkmark$\checkmark \checkmark$\checkmark^\checkmark\\checkmarkc\checkmarki\checkmarkr\checkmarkc\checkmark$\checkmark \checkmark$\checkmark^\checkmark\\checkmarkc\checkmarki\checkmarkr\checkmarkc\checkmark$\checkmark \checkmark$\checkmark^\checkmark\\checkmarkc\checkmarki\checkmarkr\checkmarkc\checkmark$\checkmark \checkmark$\checkmark^\checkmark\\checkmarkc\checkmarki\checkmarkr\checkmarkc\checkmark$\checkmark\\checkmark$\checkmark^\checkmark\\checkmarkc\checkmarki\checkmarkr\checkmarkc\checkmark$\checkmarki\checkmark$\checkmark^\checkmark\\checkmarkc\checkmarki\checkmarkr\checkmarkc\checkmark$\checkmarkt\checkmark$\checkmark^\checkmark\\checkmarkc\checkmarki\checkmarkr\checkmarkc\checkmark$\checkmarke\checkmark$\checkmark^\checkmark\\checkmarkc\checkmarki\checkmarkr\checkmarkc\checkmark$\checkmarkm\checkmark$\checkmark^\checkmark\\checkmarkc\checkmarki\checkmarkr\checkmarkc\checkmark$\checkmark \checkmark$\checkmark^\checkmark\\checkmarkc\checkmarki\checkmarkr\checkmarkc\checkmark$\checkmarkC\checkmark$\checkmark^\checkmark\\checkmarkc\checkmarki\checkmarkr\checkmarkc\checkmark$\checkmarkh\checkmark$\checkmark^\checkmark\\checkmarkc\checkmarki\checkmarkr\checkmarkc\checkmark$\checkmarka\checkmark$\checkmark^\checkmark\\checkmarkc\checkmarki\checkmarkr\checkmarkc\checkmark$\checkmarkr\checkmark$\checkmark^\checkmark\\checkmarkc\checkmarki\checkmarkr\checkmarkc\checkmark$\checkmarkg\checkmark$\checkmark^\checkmark\\checkmarkc\checkmarki\checkmarkr\checkmarkc\checkmark$\checkmarke\checkmark$\checkmark^\checkmark\\checkmarkc\checkmarki\checkmarkr\checkmarkc\checkmark$\checkmark \checkmark$\checkmark^\checkmark\\checkmarkc\checkmarki\checkmarkr\checkmarkc\checkmark$\checkmarkd\checkmark$\checkmark^\checkmark\\checkmarkc\checkmarki\checkmarkr\checkmarkc\checkmark$\checkmarky\checkmark$\checkmark^\checkmark\\checkmarkc\checkmarki\checkmarkr\checkmarkc\checkmark$\checkmarkn\checkmark$\checkmark^\checkmark\\checkmarkc\checkmarki\checkmarkr\checkmarkc\checkmark$\checkmarka\checkmark$\checkmark^\checkmark\\checkmarkc\checkmarki\checkmarkr\checkmarkc\checkmark$\checkmarkm\checkmark$\checkmark^\checkmark\\checkmarkc\checkmarki\checkmarkr\checkmarkc\checkmark$\checkmarki\checkmark$\checkmark^\checkmark\\checkmarkc\checkmarki\checkmarkr\checkmarkc\checkmark$\checkmarkq\checkmark$\checkmark^\checkmark\\checkmarkc\checkmarki\checkmarkr\checkmarkc\checkmark$\checkmarku\checkmark$\checkmark^\checkmark\\checkmarkc\checkmarki\checkmarkr\checkmarkc\checkmark$\checkmarke\checkmark$\checkmark^\checkmark\\checkmarkc\checkmarki\checkmarkr\checkmarkc\checkmark$\checkmark \checkmark$\checkmark^\checkmark\\checkmarkc\checkmarki\checkmarkr\checkmarkc\checkmark$\checkmarkd\checkmark$\checkmark^\checkmark\\checkmarkc\checkmarki\checkmarkr\checkmarkc\checkmark$\checkmarke\checkmark$\checkmark^\checkmark\\checkmarkc\checkmarki\checkmarkr\checkmarkc\checkmark$\checkmark \checkmark$\checkmark^\checkmark\\checkmarkc\checkmarki\checkmarkr\checkmarkc\checkmark$\checkmarkb\checkmark$\checkmark^\checkmark\\checkmarkc\checkmarki\checkmarkr\checkmarkc\checkmark$\checkmarka\checkmark$\checkmark^\checkmark\\checkmarkc\checkmarki\checkmarkr\checkmarkc\checkmark$\checkmarks\checkmark$\checkmark^\checkmark\\checkmarkc\checkmarki\checkmarkr\checkmarkc\checkmark$\checkmarke\checkmark$\checkmark^\checkmark\\checkmarkc\checkmarki\checkmarkr\checkmarkc\checkmark$\checkmark \checkmark$\checkmark^\checkmark\\checkmarkc\checkmarki\checkmarkr\checkmarkc\checkmark$\checkmark:\checkmark$\checkmark^\checkmark\\checkmarkc\checkmarki\checkmarkr\checkmarkc\checkmark$\checkmark \checkmark$\checkmark^\checkmark\\checkmarkc\checkmarki\checkmarkr\checkmarkc\checkmark$\checkmark$\checkmark$\checkmark^\checkmark\\checkmarkc\checkmarki\checkmarkr\checkmarkc\checkmark$\checkmarkC\checkmark$\checkmark^\checkmark\\checkmarkc\checkmarki\checkmarkr\checkmarkc\checkmark$\checkmark_\checkmark$\checkmark^\checkmark\\checkmarkc\checkmarki\checkmarkr\checkmarkc\checkmark$\checkmark{\checkmark$\checkmark^\checkmark\\checkmarkc\checkmarki\checkmarkr\checkmarkc\checkmark$\checkmarkd\checkmark$\checkmark^\checkmark\\checkmarkc\checkmarki\checkmarkr\checkmarkc\checkmark$\checkmarky\checkmark$\checkmark^\checkmark\\checkmarkc\checkmarki\checkmarkr\checkmarkc\checkmark$\checkmarkn\checkmark$\checkmark^\checkmark\\checkmarkc\checkmarki\checkmarkr\checkmarkc\checkmark$\checkmark}\checkmark$\checkmark^\checkmark\\checkmarkc\checkmarki\checkmarkr\checkmarkc\checkmark$\checkmark \checkmark$\checkmark^\checkmark\\checkmarkc\checkmarki\checkmarkr\checkmarkc\checkmark$\checkmark=\checkmark$\checkmark^\checkmark\\checkmarkc\checkmarki\checkmarkr\checkmarkc\checkmark$\checkmark \checkmark$\checkmark^\checkmark\\checkmarkc\checkmarki\checkmarkr\checkmarkc\checkmark$\checkmark1\checkmark$\checkmark^\checkmark\\checkmarkc\checkmarki\checkmarkr\checkmarkc\checkmark$\checkmark1\checkmark$\checkmark^\checkmark\\checkmarkc\checkmarki\checkmarkr\checkmarkc\checkmark$\checkmark\\checkmark$\checkmark^\checkmark\\checkmarkc\checkmarki\checkmarkr\checkmarkc\checkmark$\checkmark,\checkmark$\checkmark^\checkmark\\checkmarkc\checkmarki\checkmarkr\checkmarkc\checkmark$\checkmark5\checkmark$\checkmark^\checkmark\\checkmarkc\checkmarki\checkmarkr\checkmarkc\checkmark$\checkmark0\checkmark$\checkmark^\checkmark\\checkmarkc\checkmarki\checkmarkr\checkmarkc\checkmark$\checkmark0\checkmark$\checkmark^\checkmark\\checkmarkc\checkmarki\checkmarkr\checkmarkc\checkmark$\checkmark$\checkmark$\checkmark^\checkmark\\checkmarkc\checkmarki\checkmarkr\checkmarkc\checkmark$\checkmark \checkmark$\checkmark^\checkmark\\checkmarkc\checkmarki\checkmarkr\checkmarkc\checkmark$\checkmarkN\checkmark$\checkmark^\checkmark\\checkmarkc\checkmarki\checkmarkr\checkmarkc\checkmark$\checkmark
\checkmark$\checkmark^\checkmark\\checkmarkc\checkmarki\checkmarkr\checkmarkc\checkmark$\checkmark \checkmark$\checkmark^\checkmark\\checkmarkc\checkmarki\checkmarkr\checkmarkc\checkmark$\checkmark \checkmark$\checkmark^\checkmark\\checkmarkc\checkmarki\checkmarkr\checkmarkc\checkmark$\checkmark \checkmark$\checkmark^\checkmark\\checkmarkc\checkmarki\checkmarkr\checkmarkc\checkmark$\checkmark \checkmark$\checkmark^\checkmark\\checkmarkc\checkmarki\checkmarkr\checkmarkc\checkmark$\checkmark\\checkmark$\checkmark^\checkmark\\checkmarkc\checkmarki\checkmarkr\checkmarkc\checkmark$\checkmarki\checkmark$\checkmark^\checkmark\\checkmarkc\checkmarki\checkmarkr\checkmarkc\checkmark$\checkmarkt\checkmark$\checkmark^\checkmark\\checkmarkc\checkmarki\checkmarkr\checkmarkc\checkmark$\checkmarke\checkmark$\checkmark^\checkmark\\checkmarkc\checkmarki\checkmarkr\checkmarkc\checkmark$\checkmarkm\checkmark$\checkmark^\checkmark\\checkmarkc\checkmarki\checkmarkr\checkmarkc\checkmark$\checkmark \checkmark$\checkmark^\checkmark\\checkmarkc\checkmarki\checkmarkr\checkmarkc\checkmark$\checkmarkC\checkmark$\checkmark^\checkmark\\checkmarkc\checkmarki\checkmarkr\checkmarkc\checkmark$\checkmarkh\checkmark$\checkmark^\checkmark\\checkmarkc\checkmarki\checkmarkr\checkmarkc\checkmark$\checkmarka\checkmark$\checkmark^\checkmark\\checkmarkc\checkmarki\checkmarkr\checkmarkc\checkmark$\checkmarkr\checkmark$\checkmark^\checkmark\\checkmarkc\checkmarki\checkmarkr\checkmarkc\checkmark$\checkmarkg\checkmark$\checkmark^\checkmark\\checkmarkc\checkmarki\checkmarkr\checkmarkc\checkmark$\checkmarke\checkmark$\checkmark^\checkmark\\checkmarkc\checkmarki\checkmarkr\checkmarkc\checkmark$\checkmark \checkmark$\checkmark^\checkmark\\checkmarkc\checkmarki\checkmarkr\checkmarkc\checkmark$\checkmarks\checkmark$\checkmark^\checkmark\\checkmarkc\checkmarki\checkmarkr\checkmarkc\checkmark$\checkmarkt\checkmark$\checkmark^\checkmark\\checkmarkc\checkmarki\checkmarkr\checkmarkc\checkmark$\checkmarka\checkmark$\checkmark^\checkmark\\checkmarkc\checkmarki\checkmarkr\checkmarkc\checkmark$\checkmarkt\checkmark$\checkmark^\checkmark\\checkmarkc\checkmarki\checkmarkr\checkmarkc\checkmark$\checkmarki\checkmark$\checkmark^\checkmark\\checkmarkc\checkmarki\checkmarkr\checkmarkc\checkmark$\checkmarkq\checkmark$\checkmark^\checkmark\\checkmarkc\checkmarki\checkmarkr\checkmarkc\checkmark$\checkmarku\checkmark$\checkmark^\checkmark\\checkmarkc\checkmarki\checkmarkr\checkmarkc\checkmark$\checkmarke\checkmark$\checkmark^\checkmark\\checkmarkc\checkmarki\checkmarkr\checkmarkc\checkmark$\checkmark \checkmark$\checkmark^\checkmark\\checkmarkc\checkmarki\checkmarkr\checkmarkc\checkmark$\checkmarkd\checkmark$\checkmark^\checkmark\\checkmarkc\checkmarki\checkmarkr\checkmarkc\checkmark$\checkmarke\checkmark$\checkmark^\checkmark\\checkmarkc\checkmarki\checkmarkr\checkmarkc\checkmark$\checkmark \checkmark$\checkmark^\checkmark\\checkmarkc\checkmarki\checkmarkr\checkmarkc\checkmark$\checkmarkb\checkmark$\checkmark^\checkmark\\checkmarkc\checkmarki\checkmarkr\checkmarkc\checkmark$\checkmarka\checkmark$\checkmark^\checkmark\\checkmarkc\checkmarki\checkmarkr\checkmarkc\checkmark$\checkmarks\checkmark$\checkmark^\checkmark\\checkmarkc\checkmarki\checkmarkr\checkmarkc\checkmark$\checkmarke\checkmark$\checkmark^\checkmark\\checkmarkc\checkmarki\checkmarkr\checkmarkc\checkmark$\checkmark \checkmark$\checkmark^\checkmark\\checkmarkc\checkmarki\checkmarkr\checkmarkc\checkmark$\checkmark:\checkmark$\checkmark^\checkmark\\checkmarkc\checkmarki\checkmarkr\checkmarkc\checkmark$\checkmark \checkmark$\checkmark^\checkmark\\checkmarkc\checkmarki\checkmarkr\checkmarkc\checkmark$\checkmark$\checkmark$\checkmark^\checkmark\\checkmarkc\checkmarki\checkmarkr\checkmarkc\checkmark$\checkmarkC\checkmark$\checkmark^\checkmark\\checkmarkc\checkmarki\checkmarkr\checkmarkc\checkmark$\checkmark_\checkmark$\checkmark^\checkmark\\checkmarkc\checkmarki\checkmarkr\checkmarkc\checkmark$\checkmark{\checkmark$\checkmark^\checkmark\\checkmarkc\checkmarki\checkmarkr\checkmarkc\checkmark$\checkmark0\checkmark$\checkmark^\checkmark\\checkmarkc\checkmarki\checkmarkr\checkmarkc\checkmark$\checkmark}\checkmark$\checkmark^\checkmark\\checkmarkc\checkmarki\checkmarkr\checkmarkc\checkmark$\checkmark \checkmark$\checkmark^\checkmark\\checkmarkc\checkmarki\checkmarkr\checkmarkc\checkmark$\checkmark=\checkmark$\checkmark^\checkmark\\checkmarkc\checkmarki\checkmarkr\checkmarkc\checkmark$\checkmark \checkmark$\checkmark^\checkmark\\checkmarkc\checkmarki\checkmarkr\checkmarkc\checkmark$\checkmark1\checkmark$\checkmark^\checkmark\\checkmarkc\checkmarki\checkmarkr\checkmarkc\checkmark$\checkmark4\checkmark$\checkmark^\checkmark\\checkmarkc\checkmarki\checkmarkr\checkmarkc\checkmark$\checkmark\\checkmark$\checkmark^\checkmark\\checkmarkc\checkmarki\checkmarkr\checkmarkc\checkmark$\checkmark,\checkmark$\checkmark^\checkmark\\checkmarkc\checkmarki\checkmarkr\checkmarkc\checkmark$\checkmark6\checkmark$\checkmark^\checkmark\\checkmarkc\checkmarki\checkmarkr\checkmarkc\checkmark$\checkmark0\checkmark$\checkmark^\checkmark\\checkmarkc\checkmarki\checkmarkr\checkmarkc\checkmark$\checkmark0\checkmark$\checkmark^\checkmark\\checkmarkc\checkmarki\checkmarkr\checkmarkc\checkmark$\checkmark$\checkmark$\checkmark^\checkmark\\checkmarkc\checkmarki\checkmarkr\checkmarkc\checkmark$\checkmark \checkmark$\checkmark^\checkmark\\checkmarkc\checkmarki\checkmarkr\checkmarkc\checkmark$\checkmarkN\checkmark$\checkmark^\checkmark\\checkmarkc\checkmarki\checkmarkr\checkmarkc\checkmark$\checkmark
\checkmark$\checkmark^\checkmark\\checkmarkc\checkmarki\checkmarkr\checkmarkc\checkmark$\checkmark \checkmark$\checkmark^\checkmark\\checkmarkc\checkmarki\checkmarkr\checkmarkc\checkmark$\checkmark \checkmark$\checkmark^\checkmark\\checkmarkc\checkmarki\checkmarkr\checkmarkc\checkmark$\checkmark \checkmark$\checkmark^\checkmark\\checkmarkc\checkmarki\checkmarkr\checkmarkc\checkmark$\checkmark \checkmark$\checkmark^\checkmark\\checkmarkc\checkmarki\checkmarkr\checkmarkc\checkmark$\checkmark\\checkmark$\checkmark^\checkmark\\checkmarkc\checkmarki\checkmarkr\checkmarkc\checkmark$\checkmarki\checkmark$\checkmark^\checkmark\\checkmarkc\checkmarki\checkmarkr\checkmarkc\checkmark$\checkmarkt\checkmark$\checkmark^\checkmark\\checkmarkc\checkmarki\checkmarkr\checkmarkc\checkmark$\checkmarke\checkmark$\checkmark^\checkmark\\checkmarkc\checkmarki\checkmarkr\checkmarkc\checkmark$\checkmarkm\checkmark$\checkmark^\checkmark\\checkmarkc\checkmarki\checkmarkr\checkmarkc\checkmark$\checkmark \checkmark$\checkmark^\checkmark\\checkmarkc\checkmarki\checkmarkr\checkmarkc\checkmark$\checkmarkC\checkmark$\checkmark^\checkmark\\checkmarkc\checkmarki\checkmarkr\checkmarkc\checkmark$\checkmarko\checkmark$\checkmark^\checkmark\\checkmarkc\checkmarki\checkmarkr\checkmarkc\checkmark$\checkmarku\checkmark$\checkmark^\checkmark\\checkmarkc\checkmarki\checkmarkr\checkmarkc\checkmark$\checkmarkp\checkmark$\checkmark^\checkmark\\checkmarkc\checkmarki\checkmarkr\checkmarkc\checkmark$\checkmarkl\checkmark$\checkmark^\checkmark\\checkmarkc\checkmarki\checkmarkr\checkmarkc\checkmark$\checkmarke\checkmark$\checkmark^\checkmark\\checkmarkc\checkmarki\checkmarkr\checkmarkc\checkmark$\checkmark \checkmark$\checkmark^\checkmark\\checkmarkc\checkmarki\checkmarkr\checkmarkc\checkmark$\checkmarkr\checkmark$\checkmark^\checkmark\\checkmarkc\checkmarki\checkmarkr\checkmarkc\checkmark$\checkmarké\checkmark$\checkmark^\checkmark\\checkmarkc\checkmarki\checkmarkr\checkmarkc\checkmark$\checkmarks\checkmark$\checkmark^\checkmark\\checkmarkc\checkmarki\checkmarkr\checkmarkc\checkmark$\checkmarki\checkmark$\checkmark^\checkmark\\checkmarkc\checkmarki\checkmarkr\checkmarkc\checkmark$\checkmarks\checkmark$\checkmark^\checkmark\\checkmarkc\checkmarki\checkmarkr\checkmarkc\checkmark$\checkmarkt\checkmark$\checkmark^\checkmark\\checkmarkc\checkmarki\checkmarkr\checkmarkc\checkmark$\checkmarka\checkmark$\checkmark^\checkmark\\checkmarkc\checkmarki\checkmarkr\checkmarkc\checkmark$\checkmarkn\checkmark$\checkmark^\checkmark\\checkmarkc\checkmarki\checkmarkr\checkmarkc\checkmark$\checkmarkt\checkmark$\checkmark^\checkmark\\checkmarkc\checkmarki\checkmarkr\checkmarkc\checkmark$\checkmark \checkmark$\checkmark^\checkmark\\checkmarkc\checkmarki\checkmarkr\checkmarkc\checkmark$\checkmarka\checkmark$\checkmark^\checkmark\\checkmarkc\checkmarki\checkmarkr\checkmarkc\checkmark$\checkmarkd\checkmark$\checkmark^\checkmark\\checkmarkc\checkmarki\checkmarkr\checkmarkc\checkmark$\checkmarkm\checkmark$\checkmark^\checkmark\\checkmarkc\checkmarki\checkmarkr\checkmarkc\checkmark$\checkmarki\checkmark$\checkmark^\checkmark\\checkmarkc\checkmarki\checkmarkr\checkmarkc\checkmark$\checkmarks\checkmark$\checkmark^\checkmark\\checkmarkc\checkmarki\checkmarkr\checkmarkc\checkmark$\checkmarks\checkmark$\checkmark^\checkmark\\checkmarkc\checkmarki\checkmarkr\checkmarkc\checkmark$\checkmarki\checkmark$\checkmark^\checkmark\\checkmarkc\checkmarki\checkmarkr\checkmarkc\checkmark$\checkmarkb\checkmark$\checkmark^\checkmark\\checkmarkc\checkmarki\checkmarkr\checkmarkc\checkmark$\checkmarkl\checkmark$\checkmark^\checkmark\\checkmarkc\checkmarki\checkmarkr\checkmarkc\checkmark$\checkmarke\checkmark$\checkmark^\checkmark\\checkmarkc\checkmarki\checkmarkr\checkmarkc\checkmark$\checkmark \checkmark$\checkmark^\checkmark\\checkmarkc\checkmarki\checkmarkr\checkmarkc\checkmark$\checkmark:\checkmark$\checkmark^\checkmark\\checkmarkc\checkmarki\checkmarkr\checkmarkc\checkmark$\checkmark \checkmark$\checkmark^\checkmark\\checkmarkc\checkmarki\checkmarkr\checkmarkc\checkmark$\checkmark$\checkmark$\checkmark^\checkmark\\checkmarkc\checkmarki\checkmarkr\checkmarkc\checkmark$\checkmarkM\checkmark$\checkmark^\checkmark\\checkmarkc\checkmarki\checkmarkr\checkmarkc\checkmark$\checkmark_\checkmark$\checkmark^\checkmark\\checkmarkc\checkmarki\checkmarkr\checkmarkc\checkmark$\checkmark{\checkmark$\checkmark^\checkmark\\checkmarkc\checkmarki\checkmarkr\checkmarkc\checkmark$\checkmarka\checkmark$\checkmark^\checkmark\\checkmarkc\checkmarki\checkmarkr\checkmarkc\checkmark$\checkmarkd\checkmark$\checkmark^\checkmark\\checkmarkc\checkmarki\checkmarkr\checkmarkc\checkmark$\checkmarkm\checkmark$\checkmark^\checkmark\\checkmarkc\checkmarki\checkmarkr\checkmarkc\checkmark$\checkmark}\checkmark$\checkmark^\checkmark\\checkmarkc\checkmarki\checkmarkr\checkmarkc\checkmark$\checkmark \checkmark$\checkmark^\checkmark\\checkmarkc\checkmarki\checkmarkr\checkmarkc\checkmark$\checkmark=\checkmark$\checkmark^\checkmark\\checkmarkc\checkmarki\checkmarkr\checkmarkc\checkmark$\checkmark \checkmark$\checkmark^\checkmark\\checkmarkc\checkmarki\checkmarkr\checkmarkc\checkmark$\checkmark9\checkmark$\checkmark^\checkmark\\checkmarkc\checkmarki\checkmarkr\checkmarkc\checkmark$\checkmark0\checkmark$\checkmark^\checkmark\\checkmarkc\checkmarki\checkmarkr\checkmarkc\checkmark$\checkmark$\checkmark$\checkmark^\checkmark\\checkmarkc\checkmarki\checkmarkr\checkmarkc\checkmark$\checkmark \checkmark$\checkmark^\checkmark\\checkmarkc\checkmarki\checkmarkr\checkmarkc\checkmark$\checkmarkN\checkmark$\checkmark^\checkmark\\checkmarkc\checkmarki\checkmarkr\checkmarkc\checkmark$\checkmark$\checkmark$\checkmark^\checkmark\\checkmarkc\checkmarki\checkmarkr\checkmarkc\checkmark$\checkmark\\checkmark$\checkmark^\checkmark\\checkmarkc\checkmarki\checkmarkr\checkmarkc\checkmark$\checkmarkc\checkmark$\checkmark^\checkmark\\checkmarkc\checkmarki\checkmarkr\checkmarkc\checkmark$\checkmarkd\checkmark$\checkmark^\checkmark\\checkmarkc\checkmarki\checkmarkr\checkmarkc\checkmark$\checkmarko\checkmark$\checkmark^\checkmark\\checkmarkc\checkmarki\checkmarkr\checkmarkc\checkmark$\checkmarkt\checkmark$\checkmark^\checkmark\\checkmarkc\checkmarki\checkmarkr\checkmarkc\checkmark$\checkmark$\checkmark$\checkmark^\checkmark\\checkmarkc\checkmarki\checkmarkr\checkmarkc\checkmark$\checkmarkm\checkmark$\checkmark^\checkmark\\checkmarkc\checkmarki\checkmarkr\checkmarkc\checkmark$\checkmark
\checkmark$\checkmark^\checkmark\\checkmarkc\checkmarki\checkmarkr\checkmarkc\checkmark$\checkmark\\checkmark$\checkmark^\checkmark\\checkmarkc\checkmarki\checkmarkr\checkmarkc\checkmark$\checkmarke\checkmark$\checkmark^\checkmark\\checkmarkc\checkmarki\checkmarkr\checkmarkc\checkmark$\checkmarkn\checkmark$\checkmark^\checkmark\\checkmarkc\checkmarki\checkmarkr\checkmarkc\checkmark$\checkmarkd\checkmark$\checkmark^\checkmark\\checkmarkc\checkmarki\checkmarkr\checkmarkc\checkmark$\checkmark{\checkmark$\checkmark^\checkmark\\checkmarkc\checkmarki\checkmarkr\checkmarkc\checkmark$\checkmarki\checkmark$\checkmark^\checkmark\\checkmarkc\checkmarki\checkmarkr\checkmarkc\checkmark$\checkmarkt\checkmark$\checkmark^\checkmark\\checkmarkc\checkmarki\checkmarkr\checkmarkc\checkmark$\checkmarke\checkmark$\checkmark^\checkmark\\checkmarkc\checkmarki\checkmarkr\checkmarkc\checkmark$\checkmarkm\checkmark$\checkmark^\checkmark\\checkmarkc\checkmarki\checkmarkr\checkmarkc\checkmark$\checkmarki\checkmark$\checkmark^\checkmark\\checkmarkc\checkmarki\checkmarkr\checkmarkc\checkmark$\checkmarkz\checkmark$\checkmark^\checkmark\\checkmarkc\checkmarki\checkmarkr\checkmarkc\checkmark$\checkmarke\checkmark$\checkmark^\checkmark\\checkmarkc\checkmarki\checkmarkr\checkmarkc\checkmark$\checkmark}\checkmark$\checkmark^\checkmark\\checkmarkc\checkmarki\checkmarkr\checkmarkc\checkmark$\checkmark
\checkmark$\checkmark^\checkmark\\checkmarkc\checkmarki\checkmarkr\checkmarkc\checkmark$\checkmark
\checkmark$\checkmark^\checkmark\\checkmarkc\checkmarki\checkmarkr\checkmarkc\checkmark$\checkmark\\checkmark$\checkmark^\checkmark\\checkmarkc\checkmarki\checkmarkr\checkmarkc\checkmark$\checkmarks\checkmark$\checkmark^\checkmark\\checkmarkc\checkmarki\checkmarkr\checkmarkc\checkmark$\checkmarku\checkmark$\checkmark^\checkmark\\checkmarkc\checkmarki\checkmarkr\checkmarkc\checkmark$\checkmarkb\checkmark$\checkmark^\checkmark\\checkmarkc\checkmarki\checkmarkr\checkmarkc\checkmark$\checkmarks\checkmark$\checkmark^\checkmark\\checkmarkc\checkmarki\checkmarkr\checkmarkc\checkmark$\checkmarke\checkmark$\checkmark^\checkmark\\checkmarkc\checkmarki\checkmarkr\checkmarkc\checkmark$\checkmarkc\checkmark$\checkmark^\checkmark\\checkmarkc\checkmarki\checkmarkr\checkmarkc\checkmark$\checkmarkt\checkmark$\checkmark^\checkmark\\checkmarkc\checkmarki\checkmarkr\checkmarkc\checkmark$\checkmarki\checkmark$\checkmark^\checkmark\\checkmarkc\checkmarki\checkmarkr\checkmarkc\checkmark$\checkmarko\checkmark$\checkmark^\checkmark\\checkmarkc\checkmarki\checkmarkr\checkmarkc\checkmark$\checkmarkn\checkmark$\checkmark^\checkmark\\checkmarkc\checkmarki\checkmarkr\checkmarkc\checkmark$\checkmark{\checkmark$\checkmark^\checkmark\\checkmarkc\checkmarki\checkmarkr\checkmarkc\checkmark$\checkmarkC\checkmark$\checkmark^\checkmark\\checkmarkc\checkmarki\checkmarkr\checkmarkc\checkmark$\checkmarkh\checkmark$\checkmark^\checkmark\\checkmarkc\checkmarki\checkmarkr\checkmarkc\checkmark$\checkmarko\checkmark$\checkmark^\checkmark\\checkmarkc\checkmarki\checkmarkr\checkmarkc\checkmark$\checkmarki\checkmark$\checkmark^\checkmark\\checkmarkc\checkmarki\checkmarkr\checkmarkc\checkmark$\checkmarkx\checkmark$\checkmark^\checkmark\\checkmarkc\checkmarki\checkmarkr\checkmarkc\checkmark$\checkmark \checkmark$\checkmark^\checkmark\\checkmarkc\checkmarki\checkmarkr\checkmarkc\checkmark$\checkmarkd\checkmark$\checkmark^\checkmark\\checkmarkc\checkmarki\checkmarkr\checkmarkc\checkmark$\checkmarke\checkmark$\checkmark^\checkmark\\checkmarkc\checkmarki\checkmarkr\checkmarkc\checkmark$\checkmarks\checkmark$\checkmark^\checkmark\\checkmarkc\checkmarki\checkmarkr\checkmarkc\checkmark$\checkmark \checkmark$\checkmark^\checkmark\\checkmarkc\checkmarki\checkmarkr\checkmarkc\checkmark$\checkmarkv\checkmark$\checkmark^\checkmark\\checkmarkc\checkmarki\checkmarkr\checkmarkc\checkmark$\checkmarki\checkmark$\checkmark^\checkmark\\checkmarkc\checkmarki\checkmarkr\checkmarkc\checkmark$\checkmarks\checkmark$\checkmark^\checkmark\\checkmarkc\checkmarki\checkmarkr\checkmarkc\checkmark$\checkmark \checkmark$\checkmark^\checkmark\\checkmarkc\checkmarki\checkmarkr\checkmarkc\checkmark$\checkmarkà\checkmark$\checkmark^\checkmark\\checkmarkc\checkmarki\checkmarkr\checkmarkc\checkmark$\checkmark \checkmark$\checkmark^\checkmark\\checkmarkc\checkmarki\checkmarkr\checkmarkc\checkmark$\checkmarkb\checkmark$\checkmark^\checkmark\\checkmarkc\checkmarki\checkmarkr\checkmarkc\checkmark$\checkmarki\checkmark$\checkmark^\checkmark\\checkmarkc\checkmarki\checkmarkr\checkmarkc\checkmark$\checkmarkl\checkmark$\checkmark^\checkmark\\checkmarkc\checkmarki\checkmarkr\checkmarkc\checkmark$\checkmarkl\checkmark$\checkmark^\checkmark\\checkmarkc\checkmarki\checkmarkr\checkmarkc\checkmark$\checkmarke\checkmark$\checkmark^\checkmark\\checkmarkc\checkmarki\checkmarkr\checkmarkc\checkmark$\checkmarks\checkmark$\checkmark^\checkmark\\checkmarkc\checkmarki\checkmarkr\checkmarkc\checkmark$\checkmark}\checkmark$\checkmark^\checkmark\\checkmarkc\checkmarki\checkmarkr\checkmarkc\checkmark$\checkmark
\checkmark$\checkmark^\checkmark\\checkmarkc\checkmarki\checkmarkr\checkmarkc\checkmark$\checkmark\\checkmark$\checkmark^\checkmark\\checkmarkc\checkmarki\checkmarkr\checkmarkc\checkmark$\checkmarks\checkmark$\checkmark^\checkmark\\checkmarkc\checkmarki\checkmarkr\checkmarkc\checkmark$\checkmarku\checkmark$\checkmark^\checkmark\\checkmarkc\checkmarki\checkmarkr\checkmarkc\checkmark$\checkmarkb\checkmark$\checkmark^\checkmark\\checkmarkc\checkmarki\checkmarkr\checkmarkc\checkmark$\checkmarks\checkmark$\checkmark^\checkmark\\checkmarkc\checkmarki\checkmarkr\checkmarkc\checkmark$\checkmarku\checkmark$\checkmark^\checkmark\\checkmarkc\checkmarki\checkmarkr\checkmarkc\checkmark$\checkmarkb\checkmark$\checkmark^\checkmark\\checkmarkc\checkmarki\checkmarkr\checkmarkc\checkmark$\checkmarks\checkmark$\checkmark^\checkmark\\checkmarkc\checkmarki\checkmarkr\checkmarkc\checkmark$\checkmarke\checkmark$\checkmark^\checkmark\\checkmarkc\checkmarki\checkmarkr\checkmarkc\checkmark$\checkmarkc\checkmark$\checkmark^\checkmark\\checkmarkc\checkmarki\checkmarkr\checkmarkc\checkmark$\checkmarkt\checkmark$\checkmark^\checkmark\\checkmarkc\checkmarki\checkmarkr\checkmarkc\checkmark$\checkmarki\checkmark$\checkmark^\checkmark\\checkmarkc\checkmarki\checkmarkr\checkmarkc\checkmark$\checkmarko\checkmark$\checkmark^\checkmark\\checkmarkc\checkmarki\checkmarkr\checkmarkc\checkmark$\checkmarkn\checkmark$\checkmark^\checkmark\\checkmarkc\checkmarki\checkmarkr\checkmarkc\checkmark$\checkmark{\checkmark$\checkmark^\checkmark\\checkmarkc\checkmarki\checkmarkr\checkmarkc\checkmark$\checkmarkC\checkmark$\checkmark^\checkmark\\checkmarkc\checkmarki\checkmarkr\checkmarkc\checkmark$\checkmarko\checkmark$\checkmark^\checkmark\\checkmarkc\checkmarki\checkmarkr\checkmarkc\checkmark$\checkmarkn\checkmark$\checkmark^\checkmark\\checkmarkc\checkmarki\checkmarkr\checkmarkc\checkmark$\checkmarkd\checkmark$\checkmark^\checkmark\\checkmarkc\checkmarki\checkmarkr\checkmarkc\checkmark$\checkmarki\checkmark$\checkmark^\checkmark\\checkmarkc\checkmarki\checkmarkr\checkmarkc\checkmark$\checkmarkt\checkmark$\checkmark^\checkmark\\checkmarkc\checkmarki\checkmarkr\checkmarkc\checkmark$\checkmarki\checkmark$\checkmark^\checkmark\\checkmarkc\checkmarki\checkmarkr\checkmarkc\checkmark$\checkmarko\checkmark$\checkmark^\checkmark\\checkmarkc\checkmarki\checkmarkr\checkmarkc\checkmark$\checkmarkn\checkmark$\checkmark^\checkmark\\checkmarkc\checkmarki\checkmarkr\checkmarkc\checkmark$\checkmarks\checkmark$\checkmark^\checkmark\\checkmarkc\checkmarki\checkmarkr\checkmarkc\checkmark$\checkmark \checkmark$\checkmark^\checkmark\\checkmarkc\checkmarki\checkmarkr\checkmarkc\checkmark$\checkmarks\checkmark$\checkmark^\checkmark\\checkmarkc\checkmarki\checkmarkr\checkmarkc\checkmark$\checkmarka\checkmark$\checkmark^\checkmark\\checkmarkc\checkmarki\checkmarkr\checkmarkc\checkmark$\checkmarkt\checkmark$\checkmark^\checkmark\\checkmarkc\checkmarki\checkmarkr\checkmarkc\checkmark$\checkmarki\checkmark$\checkmark^\checkmark\\checkmarkc\checkmarki\checkmarkr\checkmarkc\checkmark$\checkmarks\checkmark$\checkmark^\checkmark\\checkmarkc\checkmarki\checkmarkr\checkmarkc\checkmark$\checkmarkf\checkmark$\checkmark^\checkmark\\checkmarkc\checkmarki\checkmarkr\checkmarkc\checkmark$\checkmarka\checkmark$\checkmark^\checkmark\\checkmarkc\checkmarki\checkmarkr\checkmarkc\checkmark$\checkmarki\checkmark$\checkmark^\checkmark\\checkmarkc\checkmarki\checkmarkr\checkmarkc\checkmark$\checkmarks\checkmark$\checkmark^\checkmark\\checkmarkc\checkmarki\checkmarkr\checkmarkc\checkmark$\checkmarka\checkmark$\checkmark^\checkmark\\checkmarkc\checkmarki\checkmarkr\checkmarkc\checkmark$\checkmarkn\checkmark$\checkmark^\checkmark\\checkmarkc\checkmarki\checkmarkr\checkmarkc\checkmark$\checkmarkt\checkmark$\checkmark^\checkmark\\checkmarkc\checkmarki\checkmarkr\checkmarkc\checkmark$\checkmarke\checkmark$\checkmark^\checkmark\\checkmarkc\checkmarki\checkmarkr\checkmarkc\checkmark$\checkmarks\checkmark$\checkmark^\checkmark\\checkmarkc\checkmarki\checkmarkr\checkmarkc\checkmark$\checkmark}\checkmark$\checkmark^\checkmark\\checkmarkc\checkmarki\checkmarkr\checkmarkc\checkmark$\checkmark
\checkmark$\checkmark^\checkmark\\checkmarkc\checkmarki\checkmarkr\checkmarkc\checkmark$\checkmarkL\checkmark$\checkmark^\checkmark\\checkmarkc\checkmarki\checkmarkr\checkmarkc\checkmark$\checkmarka\checkmark$\checkmark^\checkmark\\checkmarkc\checkmarki\checkmarkr\checkmarkc\checkmark$\checkmark \checkmark$\checkmark^\checkmark\\checkmarkc\checkmarki\checkmarkr\checkmarkc\checkmark$\checkmarkv\checkmark$\checkmark^\checkmark\\checkmarkc\checkmarki\checkmarkr\checkmarkc\checkmark$\checkmarki\checkmark$\checkmark^\checkmark\\checkmarkc\checkmarki\checkmarkr\checkmarkc\checkmark$\checkmarks\checkmark$\checkmark^\checkmark\\checkmarkc\checkmarki\checkmarkr\checkmarkc\checkmark$\checkmark \checkmark$\checkmark^\checkmark\\checkmarkc\checkmarki\checkmarkr\checkmarkc\checkmark$\checkmarkà\checkmark$\checkmark^\checkmark\\checkmarkc\checkmarki\checkmarkr\checkmarkc\checkmark$\checkmark \checkmark$\checkmark^\checkmark\\checkmarkc\checkmarki\checkmarkr\checkmarkc\checkmark$\checkmarkb\checkmark$\checkmark^\checkmark\\checkmarkc\checkmarki\checkmarkr\checkmarkc\checkmark$\checkmarki\checkmark$\checkmark^\checkmark\\checkmarkc\checkmarki\checkmarkr\checkmarkc\checkmark$\checkmarkl\checkmark$\checkmark^\checkmark\\checkmarkc\checkmarki\checkmarkr\checkmarkc\checkmark$\checkmarkl\checkmark$\checkmark^\checkmark\\checkmarkc\checkmarki\checkmarkr\checkmarkc\checkmark$\checkmarke\checkmark$\checkmark^\checkmark\\checkmarkc\checkmarki\checkmarkr\checkmarkc\checkmark$\checkmarks\checkmark$\checkmark^\checkmark\\checkmarkc\checkmarki\checkmarkr\checkmarkc\checkmark$\checkmark \checkmark$\checkmark^\checkmark\\checkmarkc\checkmarki\checkmarkr\checkmarkc\checkmark$\checkmarkd\checkmark$\checkmark^\checkmark\\checkmarkc\checkmarki\checkmarkr\checkmarkc\checkmark$\checkmarko\checkmark$\checkmark^\checkmark\\checkmarkc\checkmarki\checkmarkr\checkmarkc\checkmark$\checkmarki\checkmark$\checkmark^\checkmark\\checkmarkc\checkmarki\checkmarkr\checkmarkc\checkmark$\checkmarkt\checkmark$\checkmark^\checkmark\\checkmarkc\checkmarki\checkmarkr\checkmarkc\checkmark$\checkmark \checkmark$\checkmark^\checkmark\\checkmarkc\checkmarki\checkmarkr\checkmarkc\checkmark$\checkmark:\checkmark$\checkmark^\checkmark\\checkmarkc\checkmarki\checkmarkr\checkmarkc\checkmark$\checkmark \checkmark$\checkmark^\checkmark\\checkmarkc\checkmarki\checkmarkr\checkmarkc\checkmark$\checkmarkt\checkmark$\checkmark^\checkmark\\checkmarkc\checkmarki\checkmarkr\checkmarkc\checkmark$\checkmarkr\checkmark$\checkmark^\checkmark\\checkmarkc\checkmarki\checkmarkr\checkmarkc\checkmark$\checkmarka\checkmark$\checkmark^\checkmark\\checkmarkc\checkmarki\checkmarkr\checkmarkc\checkmark$\checkmarkn\checkmark$\checkmark^\checkmark\\checkmarkc\checkmarki\checkmarkr\checkmarkc\checkmark$\checkmarks\checkmark$\checkmark^\checkmark\\checkmarkc\checkmarki\checkmarkr\checkmarkc\checkmark$\checkmarkf\checkmark$\checkmark^\checkmark\\checkmarkc\checkmarki\checkmarkr\checkmarkc\checkmark$\checkmarko\checkmark$\checkmark^\checkmark\\checkmarkc\checkmarki\checkmarkr\checkmarkc\checkmark$\checkmarkr\checkmark$\checkmark^\checkmark\\checkmarkc\checkmarki\checkmarkr\checkmarkc\checkmark$\checkmarkm\checkmark$\checkmark^\checkmark\\checkmarkc\checkmarki\checkmarkr\checkmarkc\checkmark$\checkmarke\checkmark$\checkmark^\checkmark\\checkmarkc\checkmarki\checkmarkr\checkmarkc\checkmark$\checkmarkr\checkmark$\checkmark^\checkmark\\checkmarkc\checkmarki\checkmarkr\checkmarkc\checkmark$\checkmark \checkmark$\checkmark^\checkmark\\checkmarkc\checkmarki\checkmarkr\checkmarkc\checkmark$\checkmarkl\checkmark$\checkmark^\checkmark\\checkmarkc\checkmarki\checkmarkr\checkmarkc\checkmark$\checkmarke\checkmark$\checkmark^\checkmark\\checkmarkc\checkmarki\checkmarkr\checkmarkc\checkmark$\checkmark \checkmark$\checkmark^\checkmark\\checkmarkc\checkmarki\checkmarkr\checkmarkc\checkmark$\checkmarkm\checkmark$\checkmark^\checkmark\\checkmarkc\checkmarki\checkmarkr\checkmarkc\checkmark$\checkmarko\checkmark$\checkmark^\checkmark\\checkmarkc\checkmarki\checkmarkr\checkmarkc\checkmark$\checkmarku\checkmark$\checkmark^\checkmark\\checkmarkc\checkmarki\checkmarkr\checkmarkc\checkmark$\checkmarkv\checkmark$\checkmark^\checkmark\\checkmarkc\checkmarki\checkmarkr\checkmarkc\checkmark$\checkmarke\checkmark$\checkmark^\checkmark\\checkmarkc\checkmarki\checkmarkr\checkmarkc\checkmark$\checkmarkm\checkmark$\checkmark^\checkmark\\checkmarkc\checkmarki\checkmarkr\checkmarkc\checkmark$\checkmarke\checkmark$\checkmark^\checkmark\\checkmarkc\checkmarki\checkmarkr\checkmarkc\checkmark$\checkmarkn\checkmark$\checkmark^\checkmark\\checkmarkc\checkmarki\checkmarkr\checkmarkc\checkmark$\checkmarkt\checkmark$\checkmark^\checkmark\\checkmarkc\checkmarki\checkmarkr\checkmarkc\checkmark$\checkmark \checkmark$\checkmark^\checkmark\\checkmarkc\checkmarki\checkmarkr\checkmarkc\checkmark$\checkmarkd\checkmark$\checkmark^\checkmark\\checkmarkc\checkmarki\checkmarkr\checkmarkc\checkmark$\checkmarke\checkmark$\checkmark^\checkmark\\checkmarkc\checkmarki\checkmarkr\checkmarkc\checkmark$\checkmark \checkmark$\checkmark^\checkmark\\checkmarkc\checkmarki\checkmarkr\checkmarkc\checkmark$\checkmarkr\checkmark$\checkmark^\checkmark\\checkmarkc\checkmarki\checkmarkr\checkmarkc\checkmark$\checkmarko\checkmark$\checkmark^\checkmark\\checkmarkc\checkmarki\checkmarkr\checkmarkc\checkmark$\checkmarkt\checkmark$\checkmark^\checkmark\\checkmarkc\checkmarki\checkmarkr\checkmarkc\checkmark$\checkmarka\checkmark$\checkmark^\checkmark\\checkmarkc\checkmarki\checkmarkr\checkmarkc\checkmark$\checkmarkt\checkmark$\checkmark^\checkmark\\checkmarkc\checkmarki\checkmarkr\checkmarkc\checkmark$\checkmarki\checkmark$\checkmark^\checkmark\\checkmarkc\checkmarki\checkmarkr\checkmarkc\checkmark$\checkmarko\checkmark$\checkmark^\checkmark\\checkmarkc\checkmarki\checkmarkr\checkmarkc\checkmark$\checkmarkn\checkmark$\checkmark^\checkmark\\checkmarkc\checkmarki\checkmarkr\checkmarkc\checkmark$\checkmark \checkmark$\checkmark^\checkmark\\checkmarkc\checkmarki\checkmarkr\checkmarkc\checkmark$\checkmarke\checkmark$\checkmark^\checkmark\\checkmarkc\checkmarki\checkmarkr\checkmarkc\checkmark$\checkmarkn\checkmark$\checkmark^\checkmark\\checkmarkc\checkmarki\checkmarkr\checkmarkc\checkmark$\checkmark \checkmark$\checkmark^\checkmark\\checkmarkc\checkmarki\checkmarkr\checkmarkc\checkmark$\checkmarkt\checkmark$\checkmark^\checkmark\\checkmarkc\checkmarki\checkmarkr\checkmarkc\checkmark$\checkmarkr\checkmark$\checkmark^\checkmark\\checkmarkc\checkmarki\checkmarkr\checkmarkc\checkmark$\checkmarka\checkmark$\checkmark^\checkmark\\checkmarkc\checkmarki\checkmarkr\checkmarkc\checkmark$\checkmarkn\checkmark$\checkmark^\checkmark\\checkmarkc\checkmarki\checkmarkr\checkmarkc\checkmark$\checkmarks\checkmark$\checkmark^\checkmark\\checkmarkc\checkmarki\checkmarkr\checkmarkc\checkmark$\checkmarkl\checkmark$\checkmark^\checkmark\\checkmarkc\checkmarki\checkmarkr\checkmarkc\checkmark$\checkmarka\checkmark$\checkmark^\checkmark\\checkmarkc\checkmarki\checkmarkr\checkmarkc\checkmark$\checkmarkt\checkmark$\checkmark^\checkmark\\checkmarkc\checkmarki\checkmarkr\checkmarkc\checkmark$\checkmarki\checkmark$\checkmark^\checkmark\\checkmarkc\checkmarki\checkmarkr\checkmarkc\checkmark$\checkmarko\checkmark$\checkmark^\checkmark\\checkmarkc\checkmarki\checkmarkr\checkmarkc\checkmark$\checkmarkn\checkmark$\checkmark^\checkmark\\checkmarkc\checkmarki\checkmarkr\checkmarkc\checkmark$\checkmark \checkmark$\checkmark^\checkmark\\checkmarkc\checkmarki\checkmarkr\checkmarkc\checkmark$\checkmarka\checkmark$\checkmark^\checkmark\\checkmarkc\checkmarki\checkmarkr\checkmarkc\checkmark$\checkmarkv\checkmark$\checkmark^\checkmark\\checkmarkc\checkmarki\checkmarkr\checkmarkc\checkmark$\checkmarke\checkmark$\checkmark^\checkmark\\checkmarkc\checkmarki\checkmarkr\checkmarkc\checkmark$\checkmarkc\checkmark$\checkmark^\checkmark\\checkmarkc\checkmarki\checkmarkr\checkmarkc\checkmark$\checkmark \checkmark$\checkmark^\checkmark\\checkmarkc\checkmarki\checkmarkr\checkmarkc\checkmark$\checkmarku\checkmark$\checkmark^\checkmark\\checkmarkc\checkmarki\checkmarkr\checkmarkc\checkmark$\checkmarkn\checkmark$\checkmark^\checkmark\\checkmarkc\checkmarki\checkmarkr\checkmarkc\checkmark$\checkmark \checkmark$\checkmark^\checkmark\\checkmarkc\checkmarki\checkmarkr\checkmarkc\checkmark$\checkmarkr\checkmark$\checkmark^\checkmark\\checkmarkc\checkmarki\checkmarkr\checkmarkc\checkmark$\checkmarke\checkmark$\checkmark^\checkmark\\checkmarkc\checkmarki\checkmarkr\checkmarkc\checkmark$\checkmarkn\checkmark$\checkmark^\checkmark\\checkmarkc\checkmarki\checkmarkr\checkmarkc\checkmark$\checkmarkd\checkmark$\checkmark^\checkmark\\checkmarkc\checkmarki\checkmarkr\checkmarkc\checkmark$\checkmarke\checkmark$\checkmark^\checkmark\\checkmarkc\checkmarki\checkmarkr\checkmarkc\checkmark$\checkmarkm\checkmark$\checkmark^\checkmark\\checkmarkc\checkmarki\checkmarkr\checkmarkc\checkmark$\checkmarke\checkmark$\checkmark^\checkmark\\checkmarkc\checkmarki\checkmarkr\checkmarkc\checkmark$\checkmarkn\checkmark$\checkmark^\checkmark\\checkmarkc\checkmarki\checkmarkr\checkmarkc\checkmark$\checkmarkt\checkmark$\checkmark^\checkmark\\checkmarkc\checkmarki\checkmarkr\checkmarkc\checkmark$\checkmark \checkmark$\checkmark^\checkmark\\checkmarkc\checkmarki\checkmarkr\checkmarkc\checkmark$\checkmarké\checkmark$\checkmark^\checkmark\\checkmarkc\checkmarki\checkmarkr\checkmarkc\checkmark$\checkmarkl\checkmark$\checkmark^\checkmark\\checkmarkc\checkmarki\checkmarkr\checkmarkc\checkmark$\checkmarke\checkmark$\checkmark^\checkmark\\checkmarkc\checkmarki\checkmarkr\checkmarkc\checkmark$\checkmarkv\checkmark$\checkmark^\checkmark\\checkmarkc\checkmarki\checkmarkr\checkmarkc\checkmark$\checkmarké\checkmark$\checkmark^\checkmark\\checkmarkc\checkmarki\checkmarkr\checkmarkc\checkmark$\checkmark,\checkmark$\checkmark^\checkmark\\checkmarkc\checkmarki\checkmarkr\checkmarkc\checkmark$\checkmark \checkmark$\checkmark^\checkmark\\checkmarkc\checkmarki\checkmarkr\checkmarkc\checkmark$\checkmarka\checkmark$\checkmark^\checkmark\\checkmarkc\checkmarki\checkmarkr\checkmarkc\checkmark$\checkmarkn\checkmark$\checkmark^\checkmark\\checkmarkc\checkmarki\checkmarkr\checkmarkc\checkmark$\checkmarkn\checkmark$\checkmark^\checkmark\\checkmarkc\checkmarki\checkmarkr\checkmarkc\checkmark$\checkmarku\checkmark$\checkmark^\checkmark\\checkmarkc\checkmarki\checkmarkr\checkmarkc\checkmark$\checkmarkl\checkmark$\checkmark^\checkmark\\checkmarkc\checkmarki\checkmarkr\checkmarkc\checkmark$\checkmarke\checkmark$\checkmark^\checkmark\\checkmarkc\checkmarki\checkmarkr\checkmarkc\checkmark$\checkmarkr\checkmark$\checkmark^\checkmark\\checkmarkc\checkmarki\checkmarkr\checkmarkc\checkmark$\checkmark \checkmark$\checkmark^\checkmark\\checkmarkc\checkmarki\checkmarkr\checkmarkc\checkmark$\checkmarkl\checkmark$\checkmark^\checkmark\\checkmarkc\checkmarki\checkmarkr\checkmarkc\checkmark$\checkmarke\checkmark$\checkmark^\checkmark\\checkmarkc\checkmarki\checkmarkr\checkmarkc\checkmark$\checkmark \checkmark$\checkmark^\checkmark\\checkmarkc\checkmarki\checkmarkr\checkmarkc\checkmark$\checkmarkj\checkmark$\checkmark^\checkmark\\checkmarkc\checkmarki\checkmarkr\checkmarkc\checkmark$\checkmarke\checkmark$\checkmark^\checkmark\\checkmarkc\checkmarki\checkmarkr\checkmarkc\checkmark$\checkmarku\checkmark$\checkmark^\checkmark\\checkmarkc\checkmarki\checkmarkr\checkmarkc\checkmark$\checkmark \checkmark$\checkmark^\checkmark\\checkmarkc\checkmarki\checkmarkr\checkmarkc\checkmark$\checkmarkd\checkmark$\checkmark^\checkmark\\checkmarkc\checkmarki\checkmarkr\checkmarkc\checkmark$\checkmark'\checkmark$\checkmark^\checkmark\\checkmarkc\checkmarki\checkmarkr\checkmarkc\checkmark$\checkmarki\checkmark$\checkmark^\checkmark\\checkmarkc\checkmarki\checkmarkr\checkmarkc\checkmark$\checkmarkn\checkmark$\checkmark^\checkmark\\checkmarkc\checkmarki\checkmarkr\checkmarkc\checkmark$\checkmarkv\checkmark$\checkmark^\checkmark\\checkmarkc\checkmarki\checkmarkr\checkmarkc\checkmark$\checkmarke\checkmark$\checkmark^\checkmark\\checkmarkc\checkmarki\checkmarkr\checkmarkc\checkmark$\checkmarkr\checkmark$\checkmark^\checkmark\\checkmarkc\checkmarki\checkmarkr\checkmarkc\checkmark$\checkmarks\checkmark$\checkmark^\checkmark\\checkmarkc\checkmarki\checkmarkr\checkmarkc\checkmark$\checkmarki\checkmark$\checkmark^\checkmark\\checkmarkc\checkmarki\checkmarkr\checkmarkc\checkmark$\checkmarko\checkmark$\checkmark^\checkmark\\checkmarkc\checkmarki\checkmarkr\checkmarkc\checkmark$\checkmarkn\checkmark$\checkmark^\checkmark\\checkmarkc\checkmarki\checkmarkr\checkmarkc\checkmark$\checkmark \checkmark$\checkmark^\checkmark\\checkmarkc\checkmarki\checkmarkr\checkmarkc\checkmark$\checkmark(\checkmark$\checkmark^\checkmark\\checkmarkc\checkmarki\checkmarkr\checkmarkc\checkmark$\checkmarkB\checkmark$\checkmark^\checkmark\\checkmarkc\checkmarki\checkmarkr\checkmarkc\checkmark$\checkmarka\checkmark$\checkmark^\checkmark\\checkmarkc\checkmarki\checkmarkr\checkmarkc\checkmark$\checkmarkc\checkmark$\checkmark^\checkmark\\checkmarkc\checkmarki\checkmarkr\checkmarkc\checkmark$\checkmarkk\checkmark$\checkmark^\checkmark\\checkmarkc\checkmarki\checkmarkr\checkmarkc\checkmark$\checkmarkl\checkmark$\checkmark^\checkmark\\checkmarkc\checkmarki\checkmarkr\checkmarkc\checkmark$\checkmarka\checkmark$\checkmark^\checkmark\\checkmarkc\checkmarki\checkmarkr\checkmarkc\checkmark$\checkmarks\checkmark$\checkmark^\checkmark\\checkmarkc\checkmarki\checkmarkr\checkmarkc\checkmark$\checkmarkh\checkmark$\checkmark^\checkmark\\checkmarkc\checkmarki\checkmarkr\checkmarkc\checkmark$\checkmark)\checkmark$\checkmark^\checkmark\\checkmarkc\checkmarki\checkmarkr\checkmarkc\checkmark$\checkmark,\checkmark$\checkmark^\checkmark\\checkmarkc\checkmarki\checkmarkr\checkmarkc\checkmark$\checkmark \checkmark$\checkmark^\checkmark\\checkmarkc\checkmarki\checkmarkr\checkmarkc\checkmark$\checkmarke\checkmark$\checkmark^\checkmark\\checkmarkc\checkmarki\checkmarkr\checkmarkc\checkmark$\checkmarkt\checkmark$\checkmark^\checkmark\\checkmarkc\checkmarki\checkmarkr\checkmarkc\checkmark$\checkmark \checkmark$\checkmark^\checkmark\\checkmarkc\checkmarki\checkmarkr\checkmarkc\checkmark$\checkmarkr\checkmark$\checkmark^\checkmark\\checkmarkc\checkmarki\checkmarkr\checkmarkc\checkmark$\checkmarké\checkmark$\checkmark^\checkmark\\checkmarkc\checkmarki\checkmarkr\checkmarkc\checkmark$\checkmarks\checkmark$\checkmark^\checkmark\\checkmarkc\checkmarki\checkmarkr\checkmarkc\checkmark$\checkmarki\checkmark$\checkmark^\checkmark\\checkmarkc\checkmarki\checkmarkr\checkmarkc\checkmark$\checkmarks\checkmark$\checkmark^\checkmark\\checkmarkc\checkmarki\checkmarkr\checkmarkc\checkmark$\checkmarkt\checkmark$\checkmark^\checkmark\\checkmarkc\checkmarki\checkmarkr\checkmarkc\checkmark$\checkmarke\checkmark$\checkmark^\checkmark\\checkmarkc\checkmarki\checkmarkr\checkmarkc\checkmark$\checkmarkr\checkmark$\checkmark^\checkmark\\checkmarkc\checkmarki\checkmarkr\checkmarkc\checkmark$\checkmark \checkmark$\checkmark^\checkmark\\checkmarkc\checkmarki\checkmarkr\checkmarkc\checkmark$\checkmarka\checkmark$\checkmark^\checkmark\\checkmarkc\checkmarki\checkmarkr\checkmarkc\checkmark$\checkmarku\checkmark$\checkmark^\checkmark\\checkmarkc\checkmarki\checkmarkr\checkmarkc\checkmark$\checkmark \checkmark$\checkmark^\checkmark\\checkmarkc\checkmarki\checkmarkr\checkmarkc\checkmark$\checkmarkf\checkmark$\checkmark^\checkmark\\checkmarkc\checkmarki\checkmarkr\checkmarkc\checkmark$\checkmarkl\checkmark$\checkmark^\checkmark\\checkmarkc\checkmarki\checkmarkr\checkmarkc\checkmark$\checkmarka\checkmark$\checkmark^\checkmark\\checkmarkc\checkmarki\checkmarkr\checkmarkc\checkmark$\checkmarkm\checkmark$\checkmark^\checkmark\\checkmarkc\checkmarki\checkmarkr\checkmarkc\checkmark$\checkmarkb\checkmark$\checkmark^\checkmark\\checkmarkc\checkmarki\checkmarkr\checkmarkc\checkmark$\checkmarka\checkmark$\checkmark^\checkmark\\checkmarkc\checkmarki\checkmarkr\checkmarkc\checkmark$\checkmarkg\checkmark$\checkmark^\checkmark\\checkmarkc\checkmarki\checkmarkr\checkmarkc\checkmark$\checkmarke\checkmark$\checkmark^\checkmark\\checkmarkc\checkmarki\checkmarkr\checkmarkc\checkmark$\checkmark \checkmark$\checkmark^\checkmark\\checkmarkc\checkmarki\checkmarkr\checkmarkc\checkmark$\checkmarks\checkmark$\checkmark^\checkmark\\checkmarkc\checkmarki\checkmarkr\checkmarkc\checkmark$\checkmarko\checkmark$\checkmark^\checkmark\\checkmarkc\checkmarki\checkmarkr\checkmarkc\checkmark$\checkmarku\checkmark$\checkmark^\checkmark\\checkmarkc\checkmarki\checkmarkr\checkmarkc\checkmark$\checkmarks\checkmark$\checkmark^\checkmark\\checkmarkc\checkmarki\checkmarkr\checkmarkc\checkmark$\checkmark \checkmark$\checkmark^\checkmark\\checkmarkc\checkmarki\checkmarkr\checkmarkc\checkmark$\checkmarkc\checkmark$\checkmark^\checkmark\\checkmarkc\checkmarki\checkmarkr\checkmarkc\checkmark$\checkmarkh\checkmark$\checkmark^\checkmark\\checkmarkc\checkmarki\checkmarkr\checkmarkc\checkmark$\checkmarka\checkmark$\checkmark^\checkmark\\checkmarkc\checkmarki\checkmarkr\checkmarkc\checkmark$\checkmarkr\checkmark$\checkmark^\checkmark\\checkmarkc\checkmarki\checkmarkr\checkmarkc\checkmark$\checkmarkg\checkmark$\checkmark^\checkmark\\checkmarkc\checkmarki\checkmarkr\checkmarkc\checkmark$\checkmarke\checkmark$\checkmark^\checkmark\\checkmarkc\checkmarki\checkmarkr\checkmarkc\checkmark$\checkmark \checkmark$\checkmark^\checkmark\\checkmarkc\checkmarki\checkmarkr\checkmarkc\checkmark$\checkmarka\checkmark$\checkmark^\checkmark\\checkmarkc\checkmarki\checkmarkr\checkmarkc\checkmark$\checkmarkx\checkmark$\checkmark^\checkmark\\checkmarkc\checkmarki\checkmarkr\checkmarkc\checkmark$\checkmarki\checkmark$\checkmark^\checkmark\\checkmarkc\checkmarki\checkmarkr\checkmarkc\checkmark$\checkmarka\checkmark$\checkmark^\checkmark\\checkmarkc\checkmarki\checkmarkr\checkmarkc\checkmark$\checkmarkl\checkmark$\checkmark^\checkmark\\checkmarkc\checkmarki\checkmarkr\checkmarkc\checkmark$\checkmarke\checkmark$\checkmark^\checkmark\\checkmarkc\checkmarki\checkmarkr\checkmarkc\checkmark$\checkmark.\checkmark$\checkmark^\checkmark\\checkmarkc\checkmarki\checkmarkr\checkmarkc\checkmark$\checkmark
\checkmark$\checkmark^\checkmark\\checkmarkc\checkmarki\checkmarkr\checkmarkc\checkmark$\checkmark
\checkmark$\checkmark^\checkmark\\checkmarkc\checkmarki\checkmarkr\checkmarkc\checkmark$\checkmark\\checkmark$\checkmark^\checkmark\\checkmarkc\checkmarki\checkmarkr\checkmarkc\checkmark$\checkmarks\checkmark$\checkmark^\checkmark\\checkmarkc\checkmarki\checkmarkr\checkmarkc\checkmark$\checkmarku\checkmark$\checkmark^\checkmark\\checkmarkc\checkmarki\checkmarkr\checkmarkc\checkmark$\checkmarkb\checkmark$\checkmark^\checkmark\\checkmarkc\checkmarki\checkmarkr\checkmarkc\checkmark$\checkmarks\checkmark$\checkmark^\checkmark\\checkmarkc\checkmarki\checkmarkr\checkmarkc\checkmark$\checkmarku\checkmark$\checkmark^\checkmark\\checkmarkc\checkmarki\checkmarkr\checkmarkc\checkmark$\checkmarkb\checkmark$\checkmark^\checkmark\\checkmarkc\checkmarki\checkmarkr\checkmarkc\checkmark$\checkmarks\checkmark$\checkmark^\checkmark\\checkmarkc\checkmarki\checkmarkr\checkmarkc\checkmark$\checkmarke\checkmark$\checkmark^\checkmark\\checkmarkc\checkmarki\checkmarkr\checkmarkc\checkmark$\checkmarkc\checkmark$\checkmark^\checkmark\\checkmarkc\checkmarki\checkmarkr\checkmarkc\checkmark$\checkmarkt\checkmark$\checkmark^\checkmark\\checkmarkc\checkmarki\checkmarkr\checkmarkc\checkmark$\checkmarki\checkmark$\checkmark^\checkmark\\checkmarkc\checkmarki\checkmarkr\checkmarkc\checkmark$\checkmarko\checkmark$\checkmark^\checkmark\\checkmarkc\checkmarki\checkmarkr\checkmarkc\checkmark$\checkmarkn\checkmark$\checkmark^\checkmark\\checkmarkc\checkmarki\checkmarkr\checkmarkc\checkmark$\checkmark{\checkmark$\checkmark^\checkmark\\checkmarkc\checkmarki\checkmarkr\checkmarkc\checkmark$\checkmarkC\checkmark$\checkmark^\checkmark\\checkmarkc\checkmarki\checkmarkr\checkmarkc\checkmark$\checkmarkh\checkmark$\checkmark^\checkmark\\checkmarkc\checkmarki\checkmarkr\checkmarkc\checkmark$\checkmarko\checkmark$\checkmark^\checkmark\\checkmarkc\checkmarki\checkmarkr\checkmarkc\checkmark$\checkmarki\checkmark$\checkmark^\checkmark\\checkmarkc\checkmarki\checkmarkr\checkmarkc\checkmark$\checkmarkx\checkmark$\checkmark^\checkmark\\checkmarkc\checkmarki\checkmarkr\checkmarkc\checkmark$\checkmark \checkmark$\checkmark^\checkmark\\checkmarkc\checkmarki\checkmarkr\checkmarkc\checkmark$\checkmarkc\checkmark$\checkmark^\checkmark\\checkmarkc\checkmarki\checkmarkr\checkmarkc\checkmark$\checkmarko\checkmark$\checkmark^\checkmark\\checkmarkc\checkmarki\checkmarkr\checkmarkc\checkmark$\checkmarkm\checkmark$\checkmark^\checkmark\\checkmarkc\checkmarki\checkmarkr\checkmarkc\checkmark$\checkmarkp\checkmark$\checkmark^\checkmark\\checkmarkc\checkmarki\checkmarkr\checkmarkc\checkmark$\checkmarka\checkmark$\checkmark^\checkmark\\checkmarkc\checkmarki\checkmarkr\checkmarkc\checkmark$\checkmarkr\checkmark$\checkmark^\checkmark\\checkmarkc\checkmarki\checkmarkr\checkmarkc\checkmark$\checkmarka\checkmark$\checkmark^\checkmark\\checkmarkc\checkmarki\checkmarkr\checkmarkc\checkmark$\checkmarkt\checkmark$\checkmark^\checkmark\\checkmarkc\checkmarki\checkmarkr\checkmarkc\checkmark$\checkmarki\checkmark$\checkmark^\checkmark\\checkmarkc\checkmarki\checkmarkr\checkmarkc\checkmark$\checkmarkf\checkmark$\checkmark^\checkmark\\checkmarkc\checkmarki\checkmarkr\checkmarkc\checkmark$\checkmark}\checkmark$\checkmark^\checkmark\\checkmarkc\checkmarki\checkmarkr\checkmarkc\checkmark$\checkmark
\checkmark$\checkmark^\checkmark\\checkmarkc\checkmarki\checkmarkr\checkmarkc\checkmark$\checkmark\\checkmark$\checkmark^\checkmark\\checkmarkc\checkmarki\checkmarkr\checkmarkc\checkmark$\checkmarkb\checkmark$\checkmark^\checkmark\\checkmarkc\checkmarki\checkmarkr\checkmarkc\checkmark$\checkmarke\checkmark$\checkmark^\checkmark\\checkmarkc\checkmarki\checkmarkr\checkmarkc\checkmark$\checkmarkg\checkmark$\checkmark^\checkmark\\checkmarkc\checkmarki\checkmarkr\checkmarkc\checkmark$\checkmarki\checkmark$\checkmark^\checkmark\\checkmarkc\checkmarki\checkmarkr\checkmarkc\checkmark$\checkmarkn\checkmark$\checkmark^\checkmark\\checkmarkc\checkmarki\checkmarkr\checkmarkc\checkmark$\checkmark{\checkmark$\checkmark^\checkmark\\checkmarkc\checkmarki\checkmarkr\checkmarkc\checkmark$\checkmarkt\checkmark$\checkmark^\checkmark\\checkmarkc\checkmarki\checkmarkr\checkmarkc\checkmark$\checkmarka\checkmark$\checkmark^\checkmark\\checkmarkc\checkmarki\checkmarkr\checkmarkc\checkmark$\checkmarkb\checkmark$\checkmark^\checkmark\\checkmarkc\checkmarki\checkmarkr\checkmarkc\checkmark$\checkmarkl\checkmark$\checkmark^\checkmark\\checkmarkc\checkmarki\checkmarkr\checkmarkc\checkmark$\checkmarke\checkmark$\checkmark^\checkmark\\checkmarkc\checkmarki\checkmarkr\checkmarkc\checkmark$\checkmark}\checkmark$\checkmark^\checkmark\\checkmarkc\checkmarki\checkmarkr\checkmarkc\checkmark$\checkmark[\checkmark$\checkmark^\checkmark\\checkmarkc\checkmarki\checkmarkr\checkmarkc\checkmark$\checkmarkh\checkmark$\checkmark^\checkmark\\checkmarkc\checkmarki\checkmarkr\checkmarkc\checkmark$\checkmarkt\checkmark$\checkmark^\checkmark\\checkmarkc\checkmarki\checkmarkr\checkmarkc\checkmark$\checkmarkb\checkmark$\checkmark^\checkmark\\checkmarkc\checkmarki\checkmarkr\checkmarkc\checkmark$\checkmarkp\checkmark$\checkmark^\checkmark\\checkmarkc\checkmarki\checkmarkr\checkmarkc\checkmark$\checkmark]\checkmark$\checkmark^\checkmark\\checkmarkc\checkmarki\checkmarkr\checkmarkc\checkmark$\checkmark
\checkmark$\checkmark^\checkmark\\checkmarkc\checkmarki\checkmarkr\checkmarkc\checkmark$\checkmark\\checkmark$\checkmark^\checkmark\\checkmarkc\checkmarki\checkmarkr\checkmarkc\checkmark$\checkmarkc\checkmark$\checkmark^\checkmark\\checkmarkc\checkmarki\checkmarkr\checkmarkc\checkmark$\checkmarke\checkmark$\checkmark^\checkmark\\checkmarkc\checkmarki\checkmarkr\checkmarkc\checkmark$\checkmarkn\checkmark$\checkmark^\checkmark\\checkmarkc\checkmarki\checkmarkr\checkmarkc\checkmark$\checkmarkt\checkmark$\checkmark^\checkmark\\checkmarkc\checkmarki\checkmarkr\checkmarkc\checkmark$\checkmarke\checkmark$\checkmark^\checkmark\\checkmarkc\checkmarki\checkmarkr\checkmarkc\checkmark$\checkmarkr\checkmark$\checkmark^\checkmark\\checkmarkc\checkmarki\checkmarkr\checkmarkc\checkmark$\checkmarki\checkmark$\checkmark^\checkmark\\checkmarkc\checkmarki\checkmarkr\checkmarkc\checkmark$\checkmarkn\checkmark$\checkmark^\checkmark\\checkmarkc\checkmarki\checkmarkr\checkmarkc\checkmark$\checkmarkg\checkmark$\checkmark^\checkmark\\checkmarkc\checkmarki\checkmarkr\checkmarkc\checkmark$\checkmark
\checkmark$\checkmark^\checkmark\\checkmarkc\checkmarki\checkmarkr\checkmarkc\checkmark$\checkmark\\checkmark$\checkmark^\checkmark\\checkmarkc\checkmarki\checkmarkr\checkmarkc\checkmark$\checkmarkb\checkmark$\checkmark^\checkmark\\checkmarkc\checkmarki\checkmarkr\checkmarkc\checkmark$\checkmarke\checkmark$\checkmark^\checkmark\\checkmarkc\checkmarki\checkmarkr\checkmarkc\checkmark$\checkmarkg\checkmark$\checkmark^\checkmark\\checkmarkc\checkmarki\checkmarkr\checkmarkc\checkmark$\checkmarki\checkmark$\checkmark^\checkmark\\checkmarkc\checkmarki\checkmarkr\checkmarkc\checkmark$\checkmarkn\checkmark$\checkmark^\checkmark\\checkmarkc\checkmarki\checkmarkr\checkmarkc\checkmark$\checkmark{\checkmark$\checkmark^\checkmark\\checkmarkc\checkmarki\checkmarkr\checkmarkc\checkmark$\checkmarkt\checkmark$\checkmark^\checkmark\\checkmarkc\checkmarki\checkmarkr\checkmarkc\checkmark$\checkmarka\checkmark$\checkmark^\checkmark\\checkmarkc\checkmarki\checkmarkr\checkmarkc\checkmark$\checkmarkb\checkmark$\checkmark^\checkmark\\checkmarkc\checkmarki\checkmarkr\checkmarkc\checkmark$\checkmarku\checkmark$\checkmark^\checkmark\\checkmarkc\checkmarki\checkmarkr\checkmarkc\checkmark$\checkmarkl\checkmark$\checkmark^\checkmark\\checkmarkc\checkmarki\checkmarkr\checkmarkc\checkmark$\checkmarka\checkmark$\checkmark^\checkmark\\checkmarkc\checkmarki\checkmarkr\checkmarkc\checkmark$\checkmarkr\checkmark$\checkmark^\checkmark\\checkmarkc\checkmarki\checkmarkr\checkmarkc\checkmark$\checkmark}\checkmark$\checkmark^\checkmark\\checkmarkc\checkmarki\checkmarkr\checkmarkc\checkmark$\checkmark{\checkmark$\checkmark^\checkmark\\checkmarkc\checkmarki\checkmarkr\checkmarkc\checkmark$\checkmark|\checkmark$\checkmark^\checkmark\\checkmarkc\checkmarki\checkmarkr\checkmarkc\checkmark$\checkmarkl\checkmark$\checkmark^\checkmark\\checkmarkc\checkmarki\checkmarkr\checkmarkc\checkmark$\checkmark|\checkmark$\checkmark^\checkmark\\checkmarkc\checkmarki\checkmarkr\checkmarkc\checkmark$\checkmarkc\checkmark$\checkmark^\checkmark\\checkmarkc\checkmarki\checkmarkr\checkmarkc\checkmark$\checkmark|\checkmark$\checkmark^\checkmark\\checkmarkc\checkmarki\checkmarkr\checkmarkc\checkmark$\checkmarkc\checkmark$\checkmark^\checkmark\\checkmarkc\checkmarki\checkmarkr\checkmarkc\checkmark$\checkmark|\checkmark$\checkmark^\checkmark\\checkmarkc\checkmarki\checkmarkr\checkmarkc\checkmark$\checkmarkc\checkmark$\checkmark^\checkmark\\checkmarkc\checkmarki\checkmarkr\checkmarkc\checkmark$\checkmark|\checkmark$\checkmark^\checkmark\\checkmarkc\checkmarki\checkmarkr\checkmarkc\checkmark$\checkmarkc\checkmark$\checkmark^\checkmark\\checkmarkc\checkmarki\checkmarkr\checkmarkc\checkmark$\checkmark|\checkmark$\checkmark^\checkmark\\checkmarkc\checkmarki\checkmarkr\checkmarkc\checkmark$\checkmark}\checkmark$\checkmark^\checkmark\\checkmarkc\checkmarki\checkmarkr\checkmarkc\checkmark$\checkmark
\checkmark$\checkmark^\checkmark\\checkmarkc\checkmarki\checkmarkr\checkmarkc\checkmark$\checkmark\\checkmark$\checkmark^\checkmark\\checkmarkc\checkmarki\checkmarkr\checkmarkc\checkmark$\checkmarkh\checkmark$\checkmark^\checkmark\\checkmarkc\checkmarki\checkmarkr\checkmarkc\checkmark$\checkmarkl\checkmark$\checkmark^\checkmark\\checkmarkc\checkmarki\checkmarkr\checkmarkc\checkmark$\checkmarki\checkmark$\checkmark^\checkmark\\checkmarkc\checkmarki\checkmarkr\checkmarkc\checkmark$\checkmarkn\checkmark$\checkmark^\checkmark\\checkmarkc\checkmarki\checkmarkr\checkmarkc\checkmark$\checkmarke\checkmark$\checkmark^\checkmark\\checkmarkc\checkmarki\checkmarkr\checkmarkc\checkmark$\checkmark
\checkmark$\checkmark^\checkmark\\checkmarkc\checkmarki\checkmarkr\checkmarkc\checkmark$\checkmark\\checkmark$\checkmark^\checkmark\\checkmarkc\checkmarki\checkmarkr\checkmarkc\checkmark$\checkmarkt\checkmark$\checkmark^\checkmark\\checkmarkc\checkmarki\checkmarkr\checkmarkc\checkmark$\checkmarke\checkmark$\checkmark^\checkmark\\checkmarkc\checkmarki\checkmarkr\checkmarkc\checkmark$\checkmarkx\checkmark$\checkmark^\checkmark\\checkmarkc\checkmarki\checkmarkr\checkmarkc\checkmark$\checkmarkt\checkmark$\checkmark^\checkmark\\checkmarkc\checkmarki\checkmarkr\checkmarkc\checkmark$\checkmarkb\checkmark$\checkmark^\checkmark\\checkmarkc\checkmarki\checkmarkr\checkmarkc\checkmark$\checkmarkf\checkmark$\checkmark^\checkmark\\checkmarkc\checkmarki\checkmarkr\checkmarkc\checkmark$\checkmark{\checkmark$\checkmark^\checkmark\\checkmarkc\checkmarki\checkmarkr\checkmarkc\checkmark$\checkmarkT\checkmark$\checkmark^\checkmark\\checkmarkc\checkmarki\checkmarkr\checkmarkc\checkmark$\checkmarky\checkmark$\checkmark^\checkmark\\checkmarkc\checkmarki\checkmarkr\checkmarkc\checkmark$\checkmarkp\checkmark$\checkmark^\checkmark\\checkmarkc\checkmarki\checkmarkr\checkmarkc\checkmark$\checkmarke\checkmark$\checkmark^\checkmark\\checkmarkc\checkmarki\checkmarkr\checkmarkc\checkmark$\checkmark}\checkmark$\checkmark^\checkmark\\checkmarkc\checkmarki\checkmarkr\checkmarkc\checkmark$\checkmark \checkmark$\checkmark^\checkmark\\checkmarkc\checkmarki\checkmarkr\checkmarkc\checkmark$\checkmark&\checkmark$\checkmark^\checkmark\\checkmarkc\checkmarki\checkmarkr\checkmarkc\checkmark$\checkmark \checkmark$\checkmark^\checkmark\\checkmarkc\checkmarki\checkmarkr\checkmarkc\checkmark$\checkmark\\checkmark$\checkmark^\checkmark\\checkmarkc\checkmarki\checkmarkr\checkmarkc\checkmark$\checkmarkt\checkmark$\checkmark^\checkmark\\checkmarkc\checkmarki\checkmarkr\checkmarkc\checkmark$\checkmarke\checkmark$\checkmark^\checkmark\\checkmarkc\checkmarki\checkmarkr\checkmarkc\checkmark$\checkmarkx\checkmark$\checkmark^\checkmark\\checkmarkc\checkmarki\checkmarkr\checkmarkc\checkmark$\checkmarkt\checkmark$\checkmark^\checkmark\\checkmarkc\checkmarki\checkmarkr\checkmarkc\checkmark$\checkmarkb\checkmark$\checkmark^\checkmark\\checkmarkc\checkmarki\checkmarkr\checkmarkc\checkmark$\checkmarkf\checkmark$\checkmark^\checkmark\\checkmarkc\checkmarki\checkmarkr\checkmarkc\checkmark$\checkmark{\checkmark$\checkmark^\checkmark\\checkmarkc\checkmarki\checkmarkr\checkmarkc\checkmark$\checkmarkR\checkmark$\checkmark^\checkmark\\checkmarkc\checkmarki\checkmarkr\checkmarkc\checkmark$\checkmarke\checkmark$\checkmark^\checkmark\\checkmarkc\checkmarki\checkmarkr\checkmarkc\checkmark$\checkmarkn\checkmark$\checkmark^\checkmark\\checkmarkc\checkmarki\checkmarkr\checkmarkc\checkmark$\checkmarkd\checkmark$\checkmark^\checkmark\\checkmarkc\checkmarki\checkmarkr\checkmarkc\checkmark$\checkmarke\checkmark$\checkmark^\checkmark\\checkmarkc\checkmarki\checkmarkr\checkmarkc\checkmark$\checkmarkm\checkmark$\checkmark^\checkmark\\checkmarkc\checkmarki\checkmarkr\checkmarkc\checkmark$\checkmarke\checkmark$\checkmark^\checkmark\\checkmarkc\checkmarki\checkmarkr\checkmarkc\checkmark$\checkmarkn\checkmark$\checkmark^\checkmark\\checkmarkc\checkmarki\checkmarkr\checkmarkc\checkmark$\checkmarkt\checkmark$\checkmark^\checkmark\\checkmarkc\checkmarki\checkmarkr\checkmarkc\checkmark$\checkmark \checkmark$\checkmark^\checkmark\\checkmarkc\checkmarki\checkmarkr\checkmarkc\checkmark$\checkmark$\checkmark$\checkmark^\checkmark\\checkmarkc\checkmarki\checkmarkr\checkmarkc\checkmark$\checkmark\\checkmark$\checkmark^\checkmark\\checkmarkc\checkmarki\checkmarkr\checkmarkc\checkmark$\checkmarke\checkmark$\checkmark^\checkmark\\checkmarkc\checkmarki\checkmarkr\checkmarkc\checkmark$\checkmarkt\checkmark$\checkmark^\checkmark\\checkmarkc\checkmarki\checkmarkr\checkmarkc\checkmark$\checkmarka\checkmark$\checkmark^\checkmark\\checkmarkc\checkmarki\checkmarkr\checkmarkc\checkmark$\checkmark$\checkmark$\checkmark^\checkmark\\checkmarkc\checkmarki\checkmarkr\checkmarkc\checkmark$\checkmark}\checkmark$\checkmark^\checkmark\\checkmarkc\checkmarki\checkmarkr\checkmarkc\checkmark$\checkmark \checkmark$\checkmark^\checkmark\\checkmarkc\checkmarki\checkmarkr\checkmarkc\checkmark$\checkmark&\checkmark$\checkmark^\checkmark\\checkmarkc\checkmarki\checkmarkr\checkmarkc\checkmark$\checkmark \checkmark$\checkmark^\checkmark\\checkmarkc\checkmarki\checkmarkr\checkmarkc\checkmark$\checkmark\\checkmark$\checkmark^\checkmark\\checkmarkc\checkmarki\checkmarkr\checkmarkc\checkmark$\checkmarkt\checkmark$\checkmark^\checkmark\\checkmarkc\checkmarki\checkmarkr\checkmarkc\checkmark$\checkmarke\checkmark$\checkmark^\checkmark\\checkmarkc\checkmarki\checkmarkr\checkmarkc\checkmark$\checkmarkx\checkmark$\checkmark^\checkmark\\checkmarkc\checkmarki\checkmarkr\checkmarkc\checkmark$\checkmarkt\checkmark$\checkmark^\checkmark\\checkmarkc\checkmarki\checkmarkr\checkmarkc\checkmark$\checkmarkb\checkmark$\checkmark^\checkmark\\checkmarkc\checkmarki\checkmarkr\checkmarkc\checkmark$\checkmarkf\checkmark$\checkmark^\checkmark\\checkmarkc\checkmarki\checkmarkr\checkmarkc\checkmark$\checkmark{\checkmark$\checkmark^\checkmark\\checkmarkc\checkmarki\checkmarkr\checkmarkc\checkmark$\checkmarkB\checkmark$\checkmark^\checkmark\\checkmarkc\checkmarki\checkmarkr\checkmarkc\checkmark$\checkmarka\checkmark$\checkmark^\checkmark\\checkmarkc\checkmarki\checkmarkr\checkmarkc\checkmark$\checkmarkc\checkmark$\checkmark^\checkmark\\checkmarkc\checkmarki\checkmarkr\checkmarkc\checkmark$\checkmarkk\checkmark$\checkmark^\checkmark\\checkmarkc\checkmarki\checkmarkr\checkmarkc\checkmark$\checkmarkl\checkmark$\checkmark^\checkmark\\checkmarkc\checkmarki\checkmarkr\checkmarkc\checkmark$\checkmarka\checkmark$\checkmark^\checkmark\\checkmarkc\checkmarki\checkmarkr\checkmarkc\checkmark$\checkmarks\checkmark$\checkmark^\checkmark\\checkmarkc\checkmarki\checkmarkr\checkmarkc\checkmark$\checkmarkh\checkmark$\checkmark^\checkmark\\checkmarkc\checkmarki\checkmarkr\checkmarkc\checkmark$\checkmark}\checkmark$\checkmark^\checkmark\\checkmarkc\checkmarki\checkmarkr\checkmarkc\checkmark$\checkmark \checkmark$\checkmark^\checkmark\\checkmarkc\checkmarki\checkmarkr\checkmarkc\checkmark$\checkmark&\checkmark$\checkmark^\checkmark\\checkmarkc\checkmarki\checkmarkr\checkmarkc\checkmark$\checkmark \checkmark$\checkmark^\checkmark\\checkmarkc\checkmarki\checkmarkr\checkmarkc\checkmark$\checkmark\\checkmark$\checkmark^\checkmark\\checkmarkc\checkmarki\checkmarkr\checkmarkc\checkmark$\checkmarkt\checkmark$\checkmark^\checkmark\\checkmarkc\checkmarki\checkmarkr\checkmarkc\checkmark$\checkmarke\checkmark$\checkmark^\checkmark\\checkmarkc\checkmarki\checkmarkr\checkmarkc\checkmark$\checkmarkx\checkmark$\checkmark^\checkmark\\checkmarkc\checkmarki\checkmarkr\checkmarkc\checkmark$\checkmarkt\checkmark$\checkmark^\checkmark\\checkmarkc\checkmarki\checkmarkr\checkmarkc\checkmark$\checkmarkb\checkmark$\checkmark^\checkmark\\checkmarkc\checkmarki\checkmarkr\checkmarkc\checkmark$\checkmarkf\checkmark$\checkmark^\checkmark\\checkmarkc\checkmarki\checkmarkr\checkmarkc\checkmark$\checkmark{\checkmark$\checkmark^\checkmark\\checkmarkc\checkmarki\checkmarkr\checkmarkc\checkmark$\checkmarkR\checkmark$\checkmark^\checkmark\\checkmarkc\checkmarki\checkmarkr\checkmarkc\checkmark$\checkmarké\checkmark$\checkmark^\checkmark\\checkmarkc\checkmarki\checkmarkr\checkmarkc\checkmark$\checkmarks\checkmark$\checkmark^\checkmark\\checkmarkc\checkmarki\checkmarkr\checkmarkc\checkmark$\checkmarki\checkmark$\checkmark^\checkmark\\checkmarkc\checkmarki\checkmarkr\checkmarkc\checkmark$\checkmarks\checkmark$\checkmark^\checkmark\\checkmarkc\checkmarki\checkmarkr\checkmarkc\checkmark$\checkmarkt\checkmark$\checkmark^\checkmark\\checkmarkc\checkmarki\checkmarkr\checkmarkc\checkmark$\checkmarka\checkmark$\checkmark^\checkmark\\checkmarkc\checkmarki\checkmarkr\checkmarkc\checkmark$\checkmarkn\checkmark$\checkmark^\checkmark\\checkmarkc\checkmarki\checkmarkr\checkmarkc\checkmark$\checkmarkc\checkmark$\checkmark^\checkmark\\checkmarkc\checkmarki\checkmarkr\checkmarkc\checkmark$\checkmarke\checkmark$\checkmark^\checkmark\\checkmarkc\checkmarki\checkmarkr\checkmarkc\checkmark$\checkmark \checkmark$\checkmark^\checkmark\\checkmarkc\checkmarki\checkmarkr\checkmarkc\checkmark$\checkmarkf\checkmark$\checkmark^\checkmark\\checkmarkc\checkmarki\checkmarkr\checkmarkc\checkmark$\checkmarkl\checkmark$\checkmark^\checkmark\\checkmarkc\checkmarki\checkmarkr\checkmarkc\checkmark$\checkmarka\checkmark$\checkmark^\checkmark\\checkmarkc\checkmarki\checkmarkr\checkmarkc\checkmark$\checkmarkm\checkmark$\checkmark^\checkmark\\checkmarkc\checkmarki\checkmarkr\checkmarkc\checkmark$\checkmarkb\checkmark$\checkmark^\checkmark\\checkmarkc\checkmarki\checkmarkr\checkmarkc\checkmark$\checkmarka\checkmark$\checkmark^\checkmark\\checkmarkc\checkmarki\checkmarkr\checkmarkc\checkmark$\checkmarkg\checkmark$\checkmark^\checkmark\\checkmarkc\checkmarki\checkmarkr\checkmarkc\checkmark$\checkmarke\checkmark$\checkmark^\checkmark\\checkmarkc\checkmarki\checkmarkr\checkmarkc\checkmark$\checkmark}\checkmark$\checkmark^\checkmark\\checkmarkc\checkmarki\checkmarkr\checkmarkc\checkmark$\checkmark \checkmark$\checkmark^\checkmark\\checkmarkc\checkmarki\checkmarkr\checkmarkc\checkmark$\checkmark&\checkmark$\checkmark^\checkmark\\checkmarkc\checkmarki\checkmarkr\checkmarkc\checkmark$\checkmark \checkmark$\checkmark^\checkmark\\checkmarkc\checkmarki\checkmarkr\checkmarkc\checkmark$\checkmark\\checkmark$\checkmark^\checkmark\\checkmarkc\checkmarki\checkmarkr\checkmarkc\checkmark$\checkmarkt\checkmark$\checkmark^\checkmark\\checkmarkc\checkmarki\checkmarkr\checkmarkc\checkmark$\checkmarke\checkmark$\checkmark^\checkmark\\checkmarkc\checkmarki\checkmarkr\checkmarkc\checkmark$\checkmarkx\checkmark$\checkmark^\checkmark\\checkmarkc\checkmarki\checkmarkr\checkmarkc\checkmark$\checkmarkt\checkmark$\checkmark^\checkmark\\checkmarkc\checkmarki\checkmarkr\checkmarkc\checkmark$\checkmarkb\checkmark$\checkmark^\checkmark\\checkmarkc\checkmarki\checkmarkr\checkmarkc\checkmark$\checkmarkf\checkmark$\checkmark^\checkmark\\checkmarkc\checkmarki\checkmarkr\checkmarkc\checkmark$\checkmark{\checkmark$\checkmark^\checkmark\\checkmarkc\checkmarki\checkmarkr\checkmarkc\checkmark$\checkmarkD\checkmark$\checkmark^\checkmark\\checkmarkc\checkmarki\checkmarkr\checkmarkc\checkmark$\checkmarké\checkmark$\checkmark^\checkmark\\checkmarkc\checkmarki\checkmarkr\checkmarkc\checkmark$\checkmarkc\checkmark$\checkmark^\checkmark\\checkmarkc\checkmarki\checkmarkr\checkmarkc\checkmark$\checkmarki\checkmark$\checkmark^\checkmark\\checkmarkc\checkmarki\checkmarkr\checkmarkc\checkmark$\checkmarks\checkmark$\checkmark^\checkmark\\checkmarkc\checkmarki\checkmarkr\checkmarkc\checkmark$\checkmarki\checkmark$\checkmark^\checkmark\\checkmarkc\checkmarki\checkmarkr\checkmarkc\checkmark$\checkmarko\checkmark$\checkmark^\checkmark\\checkmarkc\checkmarki\checkmarkr\checkmarkc\checkmark$\checkmarkn\checkmark$\checkmark^\checkmark\\checkmarkc\checkmarki\checkmarkr\checkmarkc\checkmark$\checkmark}\checkmark$\checkmark^\checkmark\\checkmarkc\checkmarki\checkmarkr\checkmarkc\checkmark$\checkmark \checkmark$\checkmark^\checkmark\\checkmarkc\checkmarki\checkmarkr\checkmarkc\checkmark$\checkmark\\checkmark$\checkmark^\checkmark\\checkmarkc\checkmarki\checkmarkr\checkmarkc\checkmark$\checkmark\\checkmark$\checkmark^\checkmark\\checkmarkc\checkmarki\checkmarkr\checkmarkc\checkmark$\checkmark \checkmark$\checkmark^\checkmark\\checkmarkc\checkmarki\checkmarkr\checkmarkc\checkmark$\checkmark\\checkmark$\checkmark^\checkmark\\checkmarkc\checkmarki\checkmarkr\checkmarkc\checkmark$\checkmarkh\checkmark$\checkmark^\checkmark\\checkmarkc\checkmarki\checkmarkr\checkmarkc\checkmark$\checkmarkl\checkmark$\checkmark^\checkmark\\checkmarkc\checkmarki\checkmarkr\checkmarkc\checkmark$\checkmarki\checkmark$\checkmark^\checkmark\\checkmarkc\checkmarki\checkmarkr\checkmarkc\checkmark$\checkmarkn\checkmark$\checkmark^\checkmark\\checkmarkc\checkmarki\checkmarkr\checkmarkc\checkmark$\checkmarke\checkmark$\checkmark^\checkmark\\checkmarkc\checkmarki\checkmarkr\checkmarkc\checkmark$\checkmark
\checkmark$\checkmark^\checkmark\\checkmarkc\checkmarki\checkmarkr\checkmarkc\checkmark$\checkmarkV\checkmark$\checkmark^\checkmark\\checkmarkc\checkmarki\checkmarkr\checkmarkc\checkmark$\checkmarki\checkmark$\checkmark^\checkmark\\checkmarkc\checkmarki\checkmarkr\checkmarkc\checkmark$\checkmarks\checkmark$\checkmark^\checkmark\\checkmarkc\checkmarki\checkmarkr\checkmarkc\checkmark$\checkmark \checkmark$\checkmark^\checkmark\\checkmarkc\checkmarki\checkmarkr\checkmarkc\checkmark$\checkmarkt\checkmark$\checkmark^\checkmark\\checkmarkc\checkmarki\checkmarkr\checkmarkc\checkmark$\checkmarkr\checkmark$\checkmark^\checkmark\\checkmarkc\checkmarki\checkmarkr\checkmarkc\checkmark$\checkmarka\checkmark$\checkmark^\checkmark\\checkmarkc\checkmarki\checkmarkr\checkmarkc\checkmark$\checkmarkp\checkmark$\checkmark^\checkmark\\checkmarkc\checkmarki\checkmarkr\checkmarkc\checkmark$\checkmarké\checkmark$\checkmark^\checkmark\\checkmarkc\checkmarki\checkmarkr\checkmarkc\checkmark$\checkmarkz\checkmark$\checkmark^\checkmark\\checkmarkc\checkmarki\checkmarkr\checkmarkc\checkmark$\checkmarko\checkmark$\checkmark^\checkmark\\checkmarkc\checkmarki\checkmarkr\checkmarkc\checkmark$\checkmarkï\checkmark$\checkmark^\checkmark\\checkmarkc\checkmarki\checkmarkr\checkmarkc\checkmark$\checkmarkd\checkmark$\checkmark^\checkmark\\checkmarkc\checkmarki\checkmarkr\checkmarkc\checkmark$\checkmarka\checkmark$\checkmark^\checkmark\\checkmarkc\checkmarki\checkmarkr\checkmarkc\checkmark$\checkmarkl\checkmark$\checkmark^\checkmark\\checkmarkc\checkmarki\checkmarkr\checkmarkc\checkmark$\checkmarke\checkmark$\checkmark^\checkmark\\checkmarkc\checkmarki\checkmarkr\checkmarkc\checkmark$\checkmark \checkmark$\checkmark^\checkmark\\checkmarkc\checkmarki\checkmarkr\checkmarkc\checkmark$\checkmarkT\checkmark$\checkmark^\checkmark\\checkmarkc\checkmarki\checkmarkr\checkmarkc\checkmark$\checkmarkr\checkmark$\checkmark^\checkmark\\checkmarkc\checkmarki\checkmarkr\checkmarkc\checkmark$\checkmark \checkmark$\checkmark^\checkmark\\checkmarkc\checkmarki\checkmarkr\checkmarkc\checkmark$\checkmark&\checkmark$\checkmark^\checkmark\\checkmarkc\checkmarki\checkmarkr\checkmarkc\checkmark$\checkmark \checkmark$\checkmark^\checkmark\\checkmarkc\checkmarki\checkmarkr\checkmarkc\checkmark$\checkmark$\checkmark$\checkmark^\checkmark\\checkmarkc\checkmarki\checkmarkr\checkmarkc\checkmark$\checkmark3\checkmark$\checkmark^\checkmark\\checkmarkc\checkmarki\checkmarkr\checkmarkc\checkmark$\checkmark0\checkmark$\checkmark^\checkmark\\checkmarkc\checkmarki\checkmarkr\checkmarkc\checkmark$\checkmark\\checkmark$\checkmark^\checkmark\\checkmarkc\checkmarki\checkmarkr\checkmarkc\checkmark$\checkmarkt\checkmark$\checkmark^\checkmark\\checkmarkc\checkmarki\checkmarkr\checkmarkc\checkmark$\checkmarke\checkmark$\checkmark^\checkmark\\checkmarkc\checkmarki\checkmarkr\checkmarkc\checkmark$\checkmarkx\checkmark$\checkmark^\checkmark\\checkmarkc\checkmarki\checkmarkr\checkmarkc\checkmark$\checkmarkt\checkmark$\checkmark^\checkmark\\checkmarkc\checkmarki\checkmarkr\checkmarkc\checkmark$\checkmark{\checkmark$\checkmark^\checkmark\\checkmarkc\checkmarki\checkmarkr\checkmarkc\checkmark$\checkmark–\checkmark$\checkmark^\checkmark\\checkmarkc\checkmarki\checkmarkr\checkmarkc\checkmark$\checkmark}\checkmark$\checkmark^\checkmark\\checkmarkc\checkmarki\checkmarkr\checkmarkc\checkmark$\checkmark5\checkmark$\checkmark^\checkmark\\checkmarkc\checkmarki\checkmarkr\checkmarkc\checkmark$\checkmark0\checkmark$\checkmark^\checkmark\\checkmarkc\checkmarki\checkmarkr\checkmarkc\checkmark$\checkmark\\checkmark$\checkmark^\checkmark\\checkmarkc\checkmarki\checkmarkr\checkmarkc\checkmark$\checkmark%\checkmark$\checkmark^\checkmark\\checkmarkc\checkmarki\checkmarkr\checkmarkc\checkmark$\checkmark$\checkmark$\checkmark^\checkmark\\checkmarkc\checkmarki\checkmarkr\checkmarkc\checkmark$\checkmark \checkmark$\checkmark^\checkmark\\checkmarkc\checkmarki\checkmarkr\checkmarkc\checkmark$\checkmark&\checkmark$\checkmark^\checkmark\\checkmarkc\checkmarki\checkmarkr\checkmarkc\checkmark$\checkmark \checkmark$\checkmark^\checkmark\\checkmarkc\checkmarki\checkmarkr\checkmarkc\checkmark$\checkmarkI\checkmark$\checkmark^\checkmark\\checkmarkc\checkmarki\checkmarkr\checkmarkc\checkmark$\checkmarkm\checkmark$\checkmark^\checkmark\\checkmarkc\checkmarki\checkmarkr\checkmarkc\checkmark$\checkmarkp\checkmark$\checkmark^\checkmark\\checkmarkc\checkmarki\checkmarkr\checkmarkc\checkmark$\checkmarko\checkmark$\checkmark^\checkmark\\checkmarkc\checkmarki\checkmarkr\checkmarkc\checkmark$\checkmarkr\checkmark$\checkmark^\checkmark\\checkmarkc\checkmarki\checkmarkr\checkmarkc\checkmark$\checkmarkt\checkmark$\checkmark^\checkmark\\checkmarkc\checkmarki\checkmarkr\checkmarkc\checkmark$\checkmarka\checkmark$\checkmark^\checkmark\\checkmarkc\checkmarki\checkmarkr\checkmarkc\checkmark$\checkmarkn\checkmark$\checkmark^\checkmark\\checkmarkc\checkmarki\checkmarkr\checkmarkc\checkmark$\checkmarkt\checkmark$\checkmark^\checkmark\\checkmarkc\checkmarki\checkmarkr\checkmarkc\checkmark$\checkmark \checkmark$\checkmark^\checkmark\\checkmarkc\checkmarki\checkmarkr\checkmarkc\checkmark$\checkmark&\checkmark$\checkmark^\checkmark\\checkmarkc\checkmarki\checkmarkr\checkmarkc\checkmark$\checkmark \checkmark$\checkmark^\checkmark\\checkmarkc\checkmarki\checkmarkr\checkmarkc\checkmark$\checkmarkF\checkmark$\checkmark^\checkmark\\checkmarkc\checkmarki\checkmarkr\checkmarkc\checkmark$\checkmarka\checkmark$\checkmark^\checkmark\\checkmarkc\checkmarki\checkmarkr\checkmarkc\checkmark$\checkmarki\checkmark$\checkmark^\checkmark\\checkmarkc\checkmarki\checkmarkr\checkmarkc\checkmark$\checkmarkb\checkmark$\checkmark^\checkmark\\checkmarkc\checkmarki\checkmarkr\checkmarkc\checkmark$\checkmarkl\checkmark$\checkmark^\checkmark\\checkmarkc\checkmarki\checkmarkr\checkmarkc\checkmark$\checkmarke\checkmark$\checkmark^\checkmark\\checkmarkc\checkmarki\checkmarkr\checkmarkc\checkmark$\checkmark \checkmark$\checkmark^\checkmark\\checkmarkc\checkmarki\checkmarkr\checkmarkc\checkmark$\checkmark&\checkmark$\checkmark^\checkmark\\checkmarkc\checkmarki\checkmarkr\checkmarkc\checkmark$\checkmark \checkmark$\checkmark^\checkmark\\checkmarkc\checkmarki\checkmarkr\checkmarkc\checkmark$\checkmarkR\checkmark$\checkmark^\checkmark\\checkmarkc\checkmarki\checkmarkr\checkmarkc\checkmark$\checkmarke\checkmark$\checkmark^\checkmark\\checkmarkc\checkmarki\checkmarkr\checkmarkc\checkmark$\checkmarkj\checkmark$\checkmark^\checkmark\\checkmarkc\checkmarki\checkmarkr\checkmarkc\checkmark$\checkmarke\checkmark$\checkmark^\checkmark\\checkmarkc\checkmarki\checkmarkr\checkmarkc\checkmark$\checkmarkt\checkmark$\checkmark^\checkmark\\checkmarkc\checkmarki\checkmarkr\checkmarkc\checkmark$\checkmarké\checkmark$\checkmark^\checkmark\\checkmarkc\checkmarki\checkmarkr\checkmarkc\checkmark$\checkmark \checkmark$\checkmark^\checkmark\\checkmarkc\checkmarki\checkmarkr\checkmarkc\checkmark$\checkmark\\checkmark$\checkmark^\checkmark\\checkmarkc\checkmarki\checkmarkr\checkmarkc\checkmark$\checkmark\\checkmark$\checkmark^\checkmark\\checkmarkc\checkmarki\checkmarkr\checkmarkc\checkmark$\checkmark \checkmark$\checkmark^\checkmark\\checkmarkc\checkmarki\checkmarkr\checkmarkc\checkmark$\checkmark\\checkmark$\checkmark^\checkmark\\checkmarkc\checkmarki\checkmarkr\checkmarkc\checkmark$\checkmarkh\checkmark$\checkmark^\checkmark\\checkmarkc\checkmarki\checkmarkr\checkmarkc\checkmark$\checkmarkl\checkmark$\checkmark^\checkmark\\checkmarkc\checkmarki\checkmarkr\checkmarkc\checkmark$\checkmarki\checkmark$\checkmark^\checkmark\\checkmarkc\checkmarki\checkmarkr\checkmarkc\checkmark$\checkmarkn\checkmark$\checkmark^\checkmark\\checkmarkc\checkmarki\checkmarkr\checkmarkc\checkmark$\checkmarke\checkmark$\checkmark^\checkmark\\checkmarkc\checkmarki\checkmarkr\checkmarkc\checkmark$\checkmark
\checkmark$\checkmark^\checkmark\\checkmarkc\checkmarki\checkmarkr\checkmarkc\checkmark$\checkmarkV\checkmark$\checkmark^\checkmark\\checkmarkc\checkmarki\checkmarkr\checkmarkc\checkmark$\checkmarki\checkmark$\checkmark^\checkmark\\checkmarkc\checkmarki\checkmarkr\checkmarkc\checkmark$\checkmarks\checkmark$\checkmark^\checkmark\\checkmarkc\checkmarki\checkmarkr\checkmarkc\checkmark$\checkmark \checkmark$\checkmark^\checkmark\\checkmarkc\checkmarki\checkmarkr\checkmarkc\checkmark$\checkmarkA\checkmark$\checkmark^\checkmark\\checkmarkc\checkmarki\checkmarkr\checkmarkc\checkmark$\checkmarkC\checkmark$\checkmark^\checkmark\\checkmarkc\checkmarki\checkmarkr\checkmarkc\checkmark$\checkmarkM\checkmark$\checkmark^\checkmark\\checkmarkc\checkmarki\checkmarkr\checkmarkc\checkmark$\checkmarkE\checkmark$\checkmark^\checkmark\\checkmarkc\checkmarki\checkmarkr\checkmarkc\checkmark$\checkmark \checkmark$\checkmark^\checkmark\\checkmarkc\checkmarki\checkmarkr\checkmarkc\checkmark$\checkmark&\checkmark$\checkmark^\checkmark\\checkmarkc\checkmarki\checkmarkr\checkmarkc\checkmark$\checkmark \checkmark$\checkmark^\checkmark\\checkmarkc\checkmarki\checkmarkr\checkmarkc\checkmark$\checkmark$\checkmark$\checkmark^\checkmark\\checkmarkc\checkmarki\checkmarkr\checkmarkc\checkmark$\checkmark4\checkmark$\checkmark^\checkmark\\checkmarkc\checkmarki\checkmarkr\checkmarkc\checkmark$\checkmark0\checkmark$\checkmark^\checkmark\\checkmarkc\checkmarki\checkmarkr\checkmarkc\checkmark$\checkmark\\checkmark$\checkmark^\checkmark\\checkmarkc\checkmarki\checkmarkr\checkmarkc\checkmark$\checkmarkt\checkmark$\checkmark^\checkmark\\checkmarkc\checkmarki\checkmarkr\checkmarkc\checkmark$\checkmarke\checkmark$\checkmark^\checkmark\\checkmarkc\checkmarki\checkmarkr\checkmarkc\checkmark$\checkmarkx\checkmark$\checkmark^\checkmark\\checkmarkc\checkmarki\checkmarkr\checkmarkc\checkmark$\checkmarkt\checkmark$\checkmark^\checkmark\\checkmarkc\checkmarki\checkmarkr\checkmarkc\checkmark$\checkmark{\checkmark$\checkmark^\checkmark\\checkmarkc\checkmarki\checkmarkr\checkmarkc\checkmark$\checkmark–\checkmark$\checkmark^\checkmark\\checkmarkc\checkmarki\checkmarkr\checkmarkc\checkmark$\checkmark}\checkmark$\checkmark^\checkmark\\checkmarkc\checkmarki\checkmarkr\checkmarkc\checkmark$\checkmark6\checkmark$\checkmark^\checkmark\\checkmarkc\checkmarki\checkmarkr\checkmarkc\checkmark$\checkmark0\checkmark$\checkmark^\checkmark\\checkmarkc\checkmarki\checkmarkr\checkmarkc\checkmark$\checkmark\\checkmark$\checkmark^\checkmark\\checkmarkc\checkmarki\checkmarkr\checkmarkc\checkmark$\checkmark%\checkmark$\checkmark^\checkmark\\checkmarkc\checkmarki\checkmarkr\checkmarkc\checkmark$\checkmark$\checkmark$\checkmark^\checkmark\\checkmarkc\checkmarki\checkmarkr\checkmarkc\checkmark$\checkmark \checkmark$\checkmark^\checkmark\\checkmarkc\checkmarki\checkmarkr\checkmarkc\checkmark$\checkmark&\checkmark$\checkmark^\checkmark\\checkmarkc\checkmarki\checkmarkr\checkmarkc\checkmark$\checkmark \checkmark$\checkmark^\checkmark\\checkmarkc\checkmarki\checkmarkr\checkmarkc\checkmark$\checkmarkM\checkmark$\checkmark^\checkmark\\checkmarkc\checkmarki\checkmarkr\checkmarkc\checkmark$\checkmarko\checkmark$\checkmark^\checkmark\\checkmarkc\checkmarki\checkmarkr\checkmarkc\checkmark$\checkmarky\checkmark$\checkmark^\checkmark\\checkmarkc\checkmarki\checkmarkr\checkmarkc\checkmark$\checkmarke\checkmark$\checkmark^\checkmark\\checkmarkc\checkmarki\checkmarkr\checkmarkc\checkmark$\checkmarkn\checkmark$\checkmark^\checkmark\\checkmarkc\checkmarki\checkmarkr\checkmarkc\checkmark$\checkmark \checkmark$\checkmark^\checkmark\\checkmarkc\checkmarki\checkmarkr\checkmarkc\checkmark$\checkmark&\checkmark$\checkmark^\checkmark\\checkmarkc\checkmarki\checkmarkr\checkmarkc\checkmark$\checkmark \checkmark$\checkmark^\checkmark\\checkmarkc\checkmarki\checkmarkr\checkmarkc\checkmark$\checkmarkM\checkmark$\checkmark^\checkmark\\checkmarkc\checkmarki\checkmarkr\checkmarkc\checkmark$\checkmarko\checkmark$\checkmark^\checkmark\\checkmarkc\checkmarki\checkmarkr\checkmarkc\checkmark$\checkmarky\checkmark$\checkmark^\checkmark\\checkmarkc\checkmarki\checkmarkr\checkmarkc\checkmark$\checkmarke\checkmark$\checkmark^\checkmark\\checkmarkc\checkmarki\checkmarkr\checkmarkc\checkmark$\checkmarkn\checkmark$\checkmark^\checkmark\\checkmarkc\checkmarki\checkmarkr\checkmarkc\checkmark$\checkmark \checkmark$\checkmark^\checkmark\\checkmarkc\checkmarki\checkmarkr\checkmarkc\checkmark$\checkmark&\checkmark$\checkmark^\checkmark\\checkmarkc\checkmarki\checkmarkr\checkmarkc\checkmark$\checkmark \checkmark$\checkmark^\checkmark\\checkmarkc\checkmarki\checkmarkr\checkmarkc\checkmark$\checkmarkR\checkmark$\checkmark^\checkmark\\checkmarkc\checkmarki\checkmarkr\checkmarkc\checkmark$\checkmarke\checkmark$\checkmark^\checkmark\\checkmarkc\checkmarki\checkmarkr\checkmarkc\checkmark$\checkmarkj\checkmark$\checkmark^\checkmark\\checkmarkc\checkmarki\checkmarkr\checkmarkc\checkmark$\checkmarke\checkmark$\checkmark^\checkmark\\checkmarkc\checkmarki\checkmarkr\checkmarkc\checkmark$\checkmarkt\checkmark$\checkmark^\checkmark\\checkmarkc\checkmarki\checkmarkr\checkmarkc\checkmark$\checkmarké\checkmark$\checkmark^\checkmark\\checkmarkc\checkmarki\checkmarkr\checkmarkc\checkmark$\checkmark \checkmark$\checkmark^\checkmark\\checkmarkc\checkmarki\checkmarkr\checkmarkc\checkmark$\checkmark\\checkmark$\checkmark^\checkmark\\checkmarkc\checkmarki\checkmarkr\checkmarkc\checkmark$\checkmark\\checkmark$\checkmark^\checkmark\\checkmarkc\checkmarki\checkmarkr\checkmarkc\checkmark$\checkmark \checkmark$\checkmark^\checkmark\\checkmarkc\checkmarki\checkmarkr\checkmarkc\checkmark$\checkmark\\checkmark$\checkmark^\checkmark\\checkmarkc\checkmarki\checkmarkr\checkmarkc\checkmark$\checkmarkh\checkmark$\checkmark^\checkmark\\checkmarkc\checkmarki\checkmarkr\checkmarkc\checkmark$\checkmarkl\checkmark$\checkmark^\checkmark\\checkmarkc\checkmarki\checkmarkr\checkmarkc\checkmark$\checkmarki\checkmark$\checkmark^\checkmark\\checkmarkc\checkmarki\checkmarkr\checkmarkc\checkmark$\checkmarkn\checkmark$\checkmark^\checkmark\\checkmarkc\checkmarki\checkmarkr\checkmarkc\checkmark$\checkmarke\checkmark$\checkmark^\checkmark\\checkmarkc\checkmarki\checkmarkr\checkmarkc\checkmark$\checkmark
\checkmark$\checkmark^\checkmark\\checkmarkc\checkmarki\checkmarkr\checkmarkc\checkmark$\checkmark\\checkmark$\checkmark^\checkmark\\checkmarkc\checkmarki\checkmarkr\checkmarkc\checkmark$\checkmarkt\checkmark$\checkmark^\checkmark\\checkmarkc\checkmarki\checkmarkr\checkmarkc\checkmark$\checkmarke\checkmark$\checkmark^\checkmark\\checkmarkc\checkmarki\checkmarkr\checkmarkc\checkmark$\checkmarkx\checkmark$\checkmark^\checkmark\\checkmarkc\checkmarki\checkmarkr\checkmarkc\checkmark$\checkmarkt\checkmark$\checkmark^\checkmark\\checkmarkc\checkmarki\checkmarkr\checkmarkc\checkmark$\checkmarkb\checkmark$\checkmark^\checkmark\\checkmarkc\checkmarki\checkmarkr\checkmarkc\checkmark$\checkmarkf\checkmark$\checkmark^\checkmark\\checkmarkc\checkmarki\checkmarkr\checkmarkc\checkmark$\checkmark{\checkmark$\checkmark^\checkmark\\checkmarkc\checkmarki\checkmarkr\checkmarkc\checkmark$\checkmarkV\checkmark$\checkmark^\checkmark\\checkmarkc\checkmarki\checkmarkr\checkmarkc\checkmark$\checkmarki\checkmark$\checkmark^\checkmark\\checkmarkc\checkmarki\checkmarkr\checkmarkc\checkmark$\checkmarks\checkmark$\checkmark^\checkmark\\checkmarkc\checkmarki\checkmarkr\checkmarkc\checkmark$\checkmark \checkmark$\checkmark^\checkmark\\checkmarkc\checkmarki\checkmarkr\checkmarkc\checkmark$\checkmarkà\checkmark$\checkmark^\checkmark\\checkmarkc\checkmarki\checkmarkr\checkmarkc\checkmark$\checkmark \checkmark$\checkmark^\checkmark\\checkmarkc\checkmarki\checkmarkr\checkmarkc\checkmark$\checkmarkb\checkmark$\checkmark^\checkmark\\checkmarkc\checkmarki\checkmarkr\checkmarkc\checkmark$\checkmarki\checkmark$\checkmark^\checkmark\\checkmarkc\checkmarki\checkmarkr\checkmarkc\checkmark$\checkmarkl\checkmark$\checkmark^\checkmark\\checkmarkc\checkmarki\checkmarkr\checkmarkc\checkmark$\checkmarkl\checkmark$\checkmark^\checkmark\\checkmarkc\checkmarki\checkmarkr\checkmarkc\checkmark$\checkmarke\checkmark$\checkmark^\checkmark\\checkmarkc\checkmarki\checkmarkr\checkmarkc\checkmark$\checkmarks\checkmark$\checkmark^\checkmark\\checkmarkc\checkmarki\checkmarkr\checkmarkc\checkmark$\checkmark \checkmark$\checkmark^\checkmark\\checkmarkc\checkmarki\checkmarkr\checkmarkc\checkmark$\checkmarkS\checkmark$\checkmark^\checkmark\\checkmarkc\checkmarki\checkmarkr\checkmarkc\checkmark$\checkmarkF\checkmark$\checkmark^\checkmark\\checkmarkc\checkmarki\checkmarkr\checkmarkc\checkmark$\checkmarkU\checkmark$\checkmark^\checkmark\\checkmarkc\checkmarki\checkmarkr\checkmarkc\checkmark$\checkmark1\checkmark$\checkmark^\checkmark\\checkmarkc\checkmarki\checkmarkr\checkmarkc\checkmark$\checkmark6\checkmark$\checkmark^\checkmark\\checkmarkc\checkmarki\checkmarkr\checkmarkc\checkmark$\checkmark0\checkmark$\checkmark^\checkmark\\checkmarkc\checkmarki\checkmarkr\checkmarkc\checkmark$\checkmark5\checkmark$\checkmark^\checkmark\\checkmarkc\checkmarki\checkmarkr\checkmarkc\checkmark$\checkmark}\checkmark$\checkmark^\checkmark\\checkmarkc\checkmarki\checkmarkr\checkmarkc\checkmark$\checkmark \checkmark$\checkmark^\checkmark\\checkmarkc\checkmarki\checkmarkr\checkmarkc\checkmark$\checkmark&\checkmark$\checkmark^\checkmark\\checkmarkc\checkmarki\checkmarkr\checkmarkc\checkmark$\checkmark \checkmark$\checkmark^\checkmark\\checkmarkc\checkmarki\checkmarkr\checkmarkc\checkmark$\checkmark$\checkmark$\checkmark^\checkmark\\checkmarkc\checkmarki\checkmarkr\checkmarkc\checkmark$\checkmark\\checkmark$\checkmark^\checkmark\\checkmarkc\checkmarki\checkmarkr\checkmarkc\checkmark$\checkmarkm\checkmark$\checkmark^\checkmark\\checkmarkc\checkmarki\checkmarkr\checkmarkc\checkmark$\checkmarka\checkmark$\checkmark^\checkmark\\checkmarkc\checkmarki\checkmarkr\checkmarkc\checkmark$\checkmarkt\checkmark$\checkmark^\checkmark\\checkmarkc\checkmarki\checkmarkr\checkmarkc\checkmark$\checkmarkh\checkmark$\checkmark^\checkmark\\checkmarkc\checkmarki\checkmarkr\checkmarkc\checkmark$\checkmarkb\checkmark$\checkmark^\checkmark\\checkmarkc\checkmarki\checkmarkr\checkmarkc\checkmark$\checkmarkf\checkmark$\checkmark^\checkmark\\checkmarkc\checkmarki\checkmarkr\checkmarkc\checkmark$\checkmark{\checkmark$\checkmark^\checkmark\\checkmarkc\checkmarki\checkmarkr\checkmarkc\checkmark$\checkmark>\checkmark$\checkmark^\checkmark\\checkmarkc\checkmarki\checkmarkr\checkmarkc\checkmark$\checkmark9\checkmark$\checkmark^\checkmark\\checkmarkc\checkmarki\checkmarkr\checkmarkc\checkmark$\checkmark0\checkmark$\checkmark^\checkmark\\checkmarkc\checkmarki\checkmarkr\checkmarkc\checkmark$\checkmark\\checkmark$\checkmark^\checkmark\\checkmarkc\checkmarki\checkmarkr\checkmarkc\checkmark$\checkmark%\checkmark$\checkmark^\checkmark\\checkmarkc\checkmarki\checkmarkr\checkmarkc\checkmark$\checkmark}\checkmark$\checkmark^\checkmark\\checkmarkc\checkmarki\checkmarkr\checkmarkc\checkmark$\checkmark$\checkmark$\checkmark^\checkmark\\checkmarkc\checkmarki\checkmarkr\checkmarkc\checkmark$\checkmark \checkmark$\checkmark^\checkmark\\checkmarkc\checkmarki\checkmarkr\checkmarkc\checkmark$\checkmark&\checkmark$\checkmark^\checkmark\\checkmarkc\checkmarki\checkmarkr\checkmarkc\checkmark$\checkmark \checkmark$\checkmark^\checkmark\\checkmarkc\checkmarki\checkmarkr\checkmarkc\checkmark$\checkmark\\checkmark$\checkmark^\checkmark\\checkmarkc\checkmarki\checkmarkr\checkmarkc\checkmark$\checkmarkt\checkmark$\checkmark^\checkmark\\checkmarkc\checkmarki\checkmarkr\checkmarkc\checkmark$\checkmarke\checkmark$\checkmark^\checkmark\\checkmarkc\checkmarki\checkmarkr\checkmarkc\checkmark$\checkmarkx\checkmark$\checkmark^\checkmark\\checkmarkc\checkmarki\checkmarkr\checkmarkc\checkmark$\checkmarkt\checkmark$\checkmark^\checkmark\\checkmarkc\checkmarki\checkmarkr\checkmarkc\checkmark$\checkmarkb\checkmark$\checkmark^\checkmark\\checkmarkc\checkmarki\checkmarkr\checkmarkc\checkmark$\checkmarkf\checkmark$\checkmark^\checkmark\\checkmarkc\checkmarki\checkmarkr\checkmarkc\checkmark$\checkmark{\checkmark$\checkmark^\checkmark\\checkmarkc\checkmarki\checkmarkr\checkmarkc\checkmark$\checkmarkQ\checkmark$\checkmark^\checkmark\\checkmarkc\checkmarki\checkmarkr\checkmarkc\checkmark$\checkmarku\checkmark$\checkmark^\checkmark\\checkmarkc\checkmarki\checkmarkr\checkmarkc\checkmark$\checkmarka\checkmark$\checkmark^\checkmark\\checkmarkc\checkmarki\checkmarkr\checkmarkc\checkmark$\checkmarks\checkmark$\checkmark^\checkmark\\checkmarkc\checkmarki\checkmarkr\checkmarkc\checkmark$\checkmarki\checkmark$\checkmark^\checkmark\\checkmarkc\checkmarki\checkmarkr\checkmarkc\checkmark$\checkmark \checkmark$\checkmark^\checkmark\\checkmarkc\checkmarki\checkmarkr\checkmarkc\checkmark$\checkmarkn\checkmark$\checkmark^\checkmark\\checkmarkc\checkmarki\checkmarkr\checkmarkc\checkmark$\checkmarku\checkmark$\checkmark^\checkmark\\checkmarkc\checkmarki\checkmarkr\checkmarkc\checkmark$\checkmarkl\checkmark$\checkmark^\checkmark\\checkmarkc\checkmarki\checkmarkr\checkmarkc\checkmark$\checkmark}\checkmark$\checkmark^\checkmark\\checkmarkc\checkmarki\checkmarkr\checkmarkc\checkmark$\checkmark \checkmark$\checkmark^\checkmark\\checkmarkc\checkmarki\checkmarkr\checkmarkc\checkmark$\checkmark&\checkmark$\checkmark^\checkmark\\checkmarkc\checkmarki\checkmarkr\checkmarkc\checkmark$\checkmark \checkmark$\checkmark^\checkmark\\checkmarkc\checkmarki\checkmarkr\checkmarkc\checkmark$\checkmark\\checkmark$\checkmark^\checkmark\\checkmarkc\checkmarki\checkmarkr\checkmarkc\checkmark$\checkmarkt\checkmark$\checkmark^\checkmark\\checkmarkc\checkmarki\checkmarkr\checkmarkc\checkmark$\checkmarke\checkmark$\checkmark^\checkmark\\checkmarkc\checkmarki\checkmarkr\checkmarkc\checkmark$\checkmarkx\checkmark$\checkmark^\checkmark\\checkmarkc\checkmarki\checkmarkr\checkmarkc\checkmark$\checkmarkt\checkmark$\checkmark^\checkmark\\checkmarkc\checkmarki\checkmarkr\checkmarkc\checkmark$\checkmarkb\checkmark$\checkmark^\checkmark\\checkmarkc\checkmarki\checkmarkr\checkmarkc\checkmark$\checkmarkf\checkmark$\checkmark^\checkmark\\checkmarkc\checkmarki\checkmarkr\checkmarkc\checkmark$\checkmark{\checkmark$\checkmark^\checkmark\\checkmarkc\checkmarki\checkmarkr\checkmarkc\checkmark$\checkmarkÉ\checkmark$\checkmark^\checkmark\\checkmarkc\checkmarki\checkmarkr\checkmarkc\checkmark$\checkmarkl\checkmark$\checkmark^\checkmark\\checkmarkc\checkmarki\checkmarkr\checkmarkc\checkmark$\checkmarke\checkmark$\checkmark^\checkmark\\checkmarkc\checkmarki\checkmarkr\checkmarkc\checkmark$\checkmarkv\checkmark$\checkmark^\checkmark\\checkmarkc\checkmarki\checkmarkr\checkmarkc\checkmark$\checkmarké\checkmark$\checkmark^\checkmark\\checkmarkc\checkmarki\checkmarkr\checkmarkc\checkmark$\checkmarke\checkmark$\checkmark^\checkmark\\checkmarkc\checkmarki\checkmarkr\checkmarkc\checkmark$\checkmark}\checkmark$\checkmark^\checkmark\\checkmarkc\checkmarki\checkmarkr\checkmarkc\checkmark$\checkmark \checkmark$\checkmark^\checkmark\\checkmarkc\checkmarki\checkmarkr\checkmarkc\checkmark$\checkmark&\checkmark$\checkmark^\checkmark\\checkmarkc\checkmarki\checkmarkr\checkmarkc\checkmark$\checkmark \checkmark$\checkmark^\checkmark\\checkmarkc\checkmarki\checkmarkr\checkmarkc\checkmark$\checkmark\\checkmark$\checkmark^\checkmark\\checkmarkc\checkmarki\checkmarkr\checkmarkc\checkmark$\checkmarkt\checkmark$\checkmark^\checkmark\\checkmarkc\checkmarki\checkmarkr\checkmarkc\checkmark$\checkmarke\checkmark$\checkmark^\checkmark\\checkmarkc\checkmarki\checkmarkr\checkmarkc\checkmark$\checkmarkx\checkmark$\checkmark^\checkmark\\checkmarkc\checkmarki\checkmarkr\checkmarkc\checkmark$\checkmarkt\checkmark$\checkmark^\checkmark\\checkmarkc\checkmarki\checkmarkr\checkmarkc\checkmark$\checkmarkb\checkmark$\checkmark^\checkmark\\checkmarkc\checkmarki\checkmarkr\checkmarkc\checkmark$\checkmarkf\checkmark$\checkmark^\checkmark\\checkmarkc\checkmarki\checkmarkr\checkmarkc\checkmark$\checkmark{\checkmark$\checkmark^\checkmark\\checkmarkc\checkmarki\checkmarkr\checkmarkc\checkmark$\checkmarkR\checkmark$\checkmark^\checkmark\\checkmarkc\checkmarki\checkmarkr\checkmarkc\checkmark$\checkmarke\checkmark$\checkmark^\checkmark\\checkmarkc\checkmarki\checkmarkr\checkmarkc\checkmark$\checkmarkt\checkmark$\checkmark^\checkmark\\checkmarkc\checkmarki\checkmarkr\checkmarkc\checkmark$\checkmarke\checkmark$\checkmark^\checkmark\\checkmarkc\checkmarki\checkmarkr\checkmarkc\checkmark$\checkmarkn\checkmark$\checkmark^\checkmark\\checkmarkc\checkmarki\checkmarkr\checkmarkc\checkmark$\checkmarku\checkmark$\checkmark^\checkmark\\checkmarkc\checkmarki\checkmarkr\checkmarkc\checkmark$\checkmark}\checkmark$\checkmark^\checkmark\\checkmarkc\checkmarki\checkmarkr\checkmarkc\checkmark$\checkmark \checkmark$\checkmark^\checkmark\\checkmarkc\checkmarki\checkmarkr\checkmarkc\checkmark$\checkmark\\checkmark$\checkmark^\checkmark\\checkmarkc\checkmarki\checkmarkr\checkmarkc\checkmark$\checkmark\\checkmark$\checkmark^\checkmark\\checkmarkc\checkmarki\checkmarkr\checkmarkc\checkmark$\checkmark \checkmark$\checkmark^\checkmark\\checkmarkc\checkmarki\checkmarkr\checkmarkc\checkmark$\checkmark\\checkmark$\checkmark^\checkmark\\checkmarkc\checkmarki\checkmarkr\checkmarkc\checkmark$\checkmarkh\checkmark$\checkmark^\checkmark\\checkmarkc\checkmarki\checkmarkr\checkmarkc\checkmark$\checkmarkl\checkmark$\checkmark^\checkmark\\checkmarkc\checkmarki\checkmarkr\checkmarkc\checkmark$\checkmarki\checkmark$\checkmark^\checkmark\\checkmarkc\checkmarki\checkmarkr\checkmarkc\checkmark$\checkmarkn\checkmark$\checkmark^\checkmark\\checkmarkc\checkmarki\checkmarkr\checkmarkc\checkmark$\checkmarke\checkmark$\checkmark^\checkmark\\checkmarkc\checkmarki\checkmarkr\checkmarkc\checkmark$\checkmark
\checkmark$\checkmark^\checkmark\\checkmarkc\checkmarki\checkmarkr\checkmarkc\checkmark$\checkmarkV\checkmark$\checkmark^\checkmark\\checkmarkc\checkmarki\checkmarkr\checkmarkc\checkmark$\checkmarki\checkmark$\checkmark^\checkmark\\checkmarkc\checkmarki\checkmarkr\checkmarkc\checkmark$\checkmarks\checkmark$\checkmark^\checkmark\\checkmarkc\checkmarki\checkmarkr\checkmarkc\checkmark$\checkmark \checkmark$\checkmark^\checkmark\\checkmarkc\checkmarki\checkmarkr\checkmarkc\checkmark$\checkmarkà\checkmark$\checkmark^\checkmark\\checkmarkc\checkmarki\checkmarkr\checkmarkc\checkmark$\checkmark \checkmark$\checkmark^\checkmark\\checkmarkc\checkmarki\checkmarkr\checkmarkc\checkmark$\checkmarkb\checkmark$\checkmark^\checkmark\\checkmarkc\checkmarki\checkmarkr\checkmarkc\checkmark$\checkmarki\checkmark$\checkmark^\checkmark\\checkmarkc\checkmarki\checkmarkr\checkmarkc\checkmark$\checkmarkl\checkmark$\checkmark^\checkmark\\checkmarkc\checkmarki\checkmarkr\checkmarkc\checkmark$\checkmarkl\checkmark$\checkmark^\checkmark\\checkmarkc\checkmarki\checkmarkr\checkmarkc\checkmark$\checkmarke\checkmark$\checkmark^\checkmark\\checkmarkc\checkmarki\checkmarkr\checkmarkc\checkmark$\checkmarks\checkmark$\checkmark^\checkmark\\checkmarkc\checkmarki\checkmarkr\checkmarkc\checkmark$\checkmark \checkmark$\checkmark^\checkmark\\checkmarkc\checkmarki\checkmarkr\checkmarkc\checkmark$\checkmarkS\checkmark$\checkmark^\checkmark\\checkmarkc\checkmarki\checkmarkr\checkmarkc\checkmark$\checkmarkF\checkmark$\checkmark^\checkmark\\checkmarkc\checkmarki\checkmarkr\checkmarkc\checkmark$\checkmarkU\checkmark$\checkmark^\checkmark\\checkmarkc\checkmarki\checkmarkr\checkmarkc\checkmark$\checkmark2\checkmark$\checkmark^\checkmark\\checkmarkc\checkmarki\checkmarkr\checkmarkc\checkmark$\checkmark0\checkmark$\checkmark^\checkmark\\checkmarkc\checkmarki\checkmarkr\checkmarkc\checkmark$\checkmark0\checkmark$\checkmark^\checkmark\\checkmarkc\checkmarki\checkmarkr\checkmarkc\checkmark$\checkmark5\checkmark$\checkmark^\checkmark\\checkmarkc\checkmarki\checkmarkr\checkmarkc\checkmark$\checkmark \checkmark$\checkmark^\checkmark\\checkmarkc\checkmarki\checkmarkr\checkmarkc\checkmark$\checkmark&\checkmark$\checkmark^\checkmark\\checkmarkc\checkmarki\checkmarkr\checkmarkc\checkmark$\checkmark \checkmark$\checkmark^\checkmark\\checkmarkc\checkmarki\checkmarkr\checkmarkc\checkmark$\checkmark$\checkmark$\checkmark^\checkmark\\checkmarkc\checkmarki\checkmarkr\checkmarkc\checkmark$\checkmark>\checkmark$\checkmark^\checkmark\\checkmarkc\checkmarki\checkmarkr\checkmarkc\checkmark$\checkmark9\checkmark$\checkmark^\checkmark\\checkmarkc\checkmarki\checkmarkr\checkmarkc\checkmark$\checkmark0\checkmark$\checkmark^\checkmark\\checkmarkc\checkmarki\checkmarkr\checkmarkc\checkmark$\checkmark\\checkmark$\checkmark^\checkmark\\checkmarkc\checkmarki\checkmarkr\checkmarkc\checkmark$\checkmark%\checkmark$\checkmark^\checkmark\\checkmarkc\checkmarki\checkmarkr\checkmarkc\checkmark$\checkmark$\checkmark$\checkmark^\checkmark\\checkmarkc\checkmarki\checkmarkr\checkmarkc\checkmark$\checkmark \checkmark$\checkmark^\checkmark\\checkmarkc\checkmarki\checkmarkr\checkmarkc\checkmark$\checkmark&\checkmark$\checkmark^\checkmark\\checkmarkc\checkmarki\checkmarkr\checkmarkc\checkmark$\checkmark \checkmark$\checkmark^\checkmark\\checkmarkc\checkmarki\checkmarkr\checkmarkc\checkmark$\checkmarkQ\checkmark$\checkmark^\checkmark\\checkmarkc\checkmarki\checkmarkr\checkmarkc\checkmark$\checkmarku\checkmark$\checkmark^\checkmark\\checkmarkc\checkmarki\checkmarkr\checkmarkc\checkmark$\checkmarka\checkmark$\checkmark^\checkmark\\checkmarkc\checkmarki\checkmarkr\checkmarkc\checkmark$\checkmarks\checkmark$\checkmark^\checkmark\\checkmarkc\checkmarki\checkmarkr\checkmarkc\checkmark$\checkmarki\checkmark$\checkmark^\checkmark\\checkmarkc\checkmarki\checkmarkr\checkmarkc\checkmark$\checkmark \checkmark$\checkmark^\checkmark\\checkmarkc\checkmarki\checkmarkr\checkmarkc\checkmark$\checkmarkn\checkmark$\checkmark^\checkmark\\checkmarkc\checkmarki\checkmarkr\checkmarkc\checkmark$\checkmarku\checkmark$\checkmark^\checkmark\\checkmarkc\checkmarki\checkmarkr\checkmarkc\checkmark$\checkmarkl\checkmark$\checkmark^\checkmark\\checkmarkc\checkmarki\checkmarkr\checkmarkc\checkmark$\checkmark \checkmark$\checkmark^\checkmark\\checkmarkc\checkmarki\checkmarkr\checkmarkc\checkmark$\checkmark&\checkmark$\checkmark^\checkmark\\checkmarkc\checkmarki\checkmarkr\checkmarkc\checkmark$\checkmark \checkmark$\checkmark^\checkmark\\checkmarkc\checkmarki\checkmarkr\checkmarkc\checkmark$\checkmarkT\checkmark$\checkmark^\checkmark\\checkmarkc\checkmarki\checkmarkr\checkmarkc\checkmark$\checkmarkr\checkmark$\checkmark^\checkmark\\checkmarkc\checkmarki\checkmarkr\checkmarkc\checkmark$\checkmarkè\checkmark$\checkmark^\checkmark\\checkmarkc\checkmarki\checkmarkr\checkmarkc\checkmark$\checkmarks\checkmark$\checkmark^\checkmark\\checkmarkc\checkmarki\checkmarkr\checkmarkc\checkmark$\checkmark \checkmark$\checkmark^\checkmark\\checkmarkc\checkmarki\checkmarkr\checkmarkc\checkmark$\checkmarké\checkmark$\checkmark^\checkmark\\checkmarkc\checkmarki\checkmarkr\checkmarkc\checkmark$\checkmarkl\checkmark$\checkmark^\checkmark\\checkmarkc\checkmarki\checkmarkr\checkmarkc\checkmark$\checkmarke\checkmark$\checkmark^\checkmark\\checkmarkc\checkmarki\checkmarkr\checkmarkc\checkmark$\checkmarkv\checkmark$\checkmark^\checkmark\\checkmarkc\checkmarki\checkmarkr\checkmarkc\checkmark$\checkmarké\checkmark$\checkmark^\checkmark\\checkmarkc\checkmarki\checkmarkr\checkmarkc\checkmark$\checkmarke\checkmark$\checkmark^\checkmark\\checkmarkc\checkmarki\checkmarkr\checkmarkc\checkmark$\checkmark \checkmark$\checkmark^\checkmark\\checkmarkc\checkmarki\checkmarkr\checkmarkc\checkmark$\checkmark&\checkmark$\checkmark^\checkmark\\checkmarkc\checkmarki\checkmarkr\checkmarkc\checkmark$\checkmark \checkmark$\checkmark^\checkmark\\checkmarkc\checkmarki\checkmarkr\checkmarkc\checkmark$\checkmarkS\checkmark$\checkmark^\checkmark\\checkmarkc\checkmarki\checkmarkr\checkmarkc\checkmark$\checkmarku\checkmark$\checkmark^\checkmark\\checkmarkc\checkmarki\checkmarkr\checkmarkc\checkmark$\checkmarkr\checkmark$\checkmark^\checkmark\\checkmarkc\checkmarki\checkmarkr\checkmarkc\checkmark$\checkmarkd\checkmark$\checkmark^\checkmark\\checkmarkc\checkmarki\checkmarkr\checkmarkc\checkmark$\checkmarki\checkmark$\checkmark^\checkmark\\checkmarkc\checkmarki\checkmarkr\checkmarkc\checkmark$\checkmarkm\checkmark$\checkmark^\checkmark\\checkmarkc\checkmarki\checkmarkr\checkmarkc\checkmark$\checkmarke\checkmark$\checkmark^\checkmark\\checkmarkc\checkmarki\checkmarkr\checkmarkc\checkmark$\checkmarkn\checkmark$\checkmark^\checkmark\\checkmarkc\checkmarki\checkmarkr\checkmarkc\checkmark$\checkmarks\checkmark$\checkmark^\checkmark\\checkmarkc\checkmarki\checkmarkr\checkmarkc\checkmark$\checkmarki\checkmark$\checkmark^\checkmark\\checkmarkc\checkmarki\checkmarkr\checkmarkc\checkmark$\checkmarko\checkmark$\checkmark^\checkmark\\checkmarkc\checkmarki\checkmarkr\checkmarkc\checkmark$\checkmarkn\checkmark$\checkmark^\checkmark\\checkmarkc\checkmarki\checkmarkr\checkmarkc\checkmark$\checkmarkn\checkmark$\checkmark^\checkmark\\checkmarkc\checkmarki\checkmarkr\checkmarkc\checkmark$\checkmarké\checkmark$\checkmark^\checkmark\\checkmarkc\checkmarki\checkmarkr\checkmarkc\checkmark$\checkmark \checkmark$\checkmark^\checkmark\\checkmarkc\checkmarki\checkmarkr\checkmarkc\checkmark$\checkmark\\checkmark$\checkmark^\checkmark\\checkmarkc\checkmarki\checkmarkr\checkmarkc\checkmark$\checkmark\\checkmark$\checkmark^\checkmark\\checkmarkc\checkmarki\checkmarkr\checkmarkc\checkmark$\checkmark \checkmark$\checkmark^\checkmark\\checkmarkc\checkmarki\checkmarkr\checkmarkc\checkmark$\checkmark\\checkmark$\checkmark^\checkmark\\checkmarkc\checkmarki\checkmarkr\checkmarkc\checkmark$\checkmarkh\checkmark$\checkmark^\checkmark\\checkmarkc\checkmarki\checkmarkr\checkmarkc\checkmark$\checkmarkl\checkmark$\checkmark^\checkmark\\checkmarkc\checkmarki\checkmarkr\checkmarkc\checkmark$\checkmarki\checkmark$\checkmark^\checkmark\\checkmarkc\checkmarki\checkmarkr\checkmarkc\checkmark$\checkmarkn\checkmark$\checkmark^\checkmark\\checkmarkc\checkmarki\checkmarkr\checkmarkc\checkmark$\checkmarke\checkmark$\checkmark^\checkmark\\checkmarkc\checkmarki\checkmarkr\checkmarkc\checkmark$\checkmark
\checkmark$\checkmark^\checkmark\\checkmarkc\checkmarki\checkmarkr\checkmarkc\checkmark$\checkmark\\checkmark$\checkmark^\checkmark\\checkmarkc\checkmarki\checkmarkr\checkmarkc\checkmark$\checkmarke\checkmark$\checkmark^\checkmark\\checkmarkc\checkmarki\checkmarkr\checkmarkc\checkmark$\checkmarkn\checkmark$\checkmark^\checkmark\\checkmarkc\checkmarki\checkmarkr\checkmarkc\checkmark$\checkmarkd\checkmark$\checkmark^\checkmark\\checkmarkc\checkmarki\checkmarkr\checkmarkc\checkmark$\checkmark{\checkmark$\checkmark^\checkmark\\checkmarkc\checkmarki\checkmarkr\checkmarkc\checkmark$\checkmarkt\checkmark$\checkmark^\checkmark\\checkmarkc\checkmarki\checkmarkr\checkmarkc\checkmark$\checkmarka\checkmark$\checkmark^\checkmark\\checkmarkc\checkmarki\checkmarkr\checkmarkc\checkmark$\checkmarkb\checkmark$\checkmark^\checkmark\\checkmarkc\checkmarki\checkmarkr\checkmarkc\checkmark$\checkmarku\checkmark$\checkmark^\checkmark\\checkmarkc\checkmarki\checkmarkr\checkmarkc\checkmark$\checkmarkl\checkmark$\checkmark^\checkmark\\checkmarkc\checkmarki\checkmarkr\checkmarkc\checkmark$\checkmarka\checkmark$\checkmark^\checkmark\\checkmarkc\checkmarki\checkmarkr\checkmarkc\checkmark$\checkmarkr\checkmark$\checkmark^\checkmark\\checkmarkc\checkmarki\checkmarkr\checkmarkc\checkmark$\checkmark}\checkmark$\checkmark^\checkmark\\checkmarkc\checkmarki\checkmarkr\checkmarkc\checkmark$\checkmark
\checkmark$\checkmark^\checkmark\\checkmarkc\checkmarki\checkmarkr\checkmarkc\checkmark$\checkmark\\checkmark$\checkmark^\checkmark\\checkmarkc\checkmarki\checkmarkr\checkmarkc\checkmark$\checkmarkc\checkmark$\checkmark^\checkmark\\checkmarkc\checkmarki\checkmarkr\checkmarkc\checkmark$\checkmarka\checkmark$\checkmark^\checkmark\\checkmarkc\checkmarki\checkmarkr\checkmarkc\checkmark$\checkmarkp\checkmark$\checkmark^\checkmark\\checkmarkc\checkmarki\checkmarkr\checkmarkc\checkmark$\checkmarkt\checkmark$\checkmark^\checkmark\\checkmarkc\checkmarki\checkmarkr\checkmarkc\checkmark$\checkmarki\checkmark$\checkmark^\checkmark\\checkmarkc\checkmarki\checkmarkr\checkmarkc\checkmark$\checkmarko\checkmark$\checkmark^\checkmark\\checkmarkc\checkmarki\checkmarkr\checkmarkc\checkmark$\checkmarkn\checkmark$\checkmark^\checkmark\\checkmarkc\checkmarki\checkmarkr\checkmarkc\checkmark$\checkmark{\checkmark$\checkmark^\checkmark\\checkmarkc\checkmarki\checkmarkr\checkmarkc\checkmark$\checkmarkT\checkmark$\checkmark^\checkmark\\checkmarkc\checkmarki\checkmarkr\checkmarkc\checkmark$\checkmarka\checkmark$\checkmark^\checkmark\\checkmarkc\checkmarki\checkmarkr\checkmarkc\checkmark$\checkmarkb\checkmark$\checkmark^\checkmark\\checkmarkc\checkmarki\checkmarkr\checkmarkc\checkmark$\checkmarkl\checkmark$\checkmark^\checkmark\\checkmarkc\checkmarki\checkmarkr\checkmarkc\checkmark$\checkmarke\checkmark$\checkmark^\checkmark\\checkmarkc\checkmarki\checkmarkr\checkmarkc\checkmark$\checkmarka\checkmark$\checkmark^\checkmark\\checkmarkc\checkmarki\checkmarkr\checkmarkc\checkmark$\checkmarku\checkmark$\checkmark^\checkmark\\checkmarkc\checkmarki\checkmarkr\checkmarkc\checkmark$\checkmark \checkmark$\checkmark^\checkmark\\checkmarkc\checkmarki\checkmarkr\checkmarkc\checkmark$\checkmarkc\checkmark$\checkmark^\checkmark\\checkmarkc\checkmarki\checkmarkr\checkmarkc\checkmark$\checkmarko\checkmark$\checkmark^\checkmark\\checkmarkc\checkmarki\checkmarkr\checkmarkc\checkmark$\checkmarkm\checkmark$\checkmark^\checkmark\\checkmarkc\checkmarki\checkmarkr\checkmarkc\checkmark$\checkmarkp\checkmark$\checkmark^\checkmark\\checkmarkc\checkmarki\checkmarkr\checkmarkc\checkmark$\checkmarka\checkmark$\checkmark^\checkmark\\checkmarkc\checkmarki\checkmarkr\checkmarkc\checkmark$\checkmarkr\checkmark$\checkmark^\checkmark\\checkmarkc\checkmarki\checkmarkr\checkmarkc\checkmark$\checkmarka\checkmark$\checkmark^\checkmark\\checkmarkc\checkmarki\checkmarkr\checkmarkc\checkmark$\checkmarkt\checkmark$\checkmark^\checkmark\\checkmarkc\checkmarki\checkmarkr\checkmarkc\checkmark$\checkmarki\checkmark$\checkmark^\checkmark\\checkmarkc\checkmarki\checkmarkr\checkmarkc\checkmark$\checkmarkf\checkmark$\checkmark^\checkmark\\checkmarkc\checkmarki\checkmarkr\checkmarkc\checkmark$\checkmark \checkmark$\checkmark^\checkmark\\checkmarkc\checkmarki\checkmarkr\checkmarkc\checkmark$\checkmarkd\checkmark$\checkmark^\checkmark\\checkmarkc\checkmarki\checkmarkr\checkmarkc\checkmark$\checkmarke\checkmark$\checkmark^\checkmark\\checkmarkc\checkmarki\checkmarkr\checkmarkc\checkmark$\checkmarks\checkmark$\checkmark^\checkmark\\checkmarkc\checkmarki\checkmarkr\checkmarkc\checkmark$\checkmark \checkmark$\checkmark^\checkmark\\checkmarkc\checkmarki\checkmarkr\checkmarkc\checkmark$\checkmarks\checkmark$\checkmark^\checkmark\\checkmarkc\checkmarki\checkmarkr\checkmarkc\checkmark$\checkmarko\checkmark$\checkmark^\checkmark\\checkmarkc\checkmarki\checkmarkr\checkmarkc\checkmark$\checkmarkl\checkmark$\checkmark^\checkmark\\checkmarkc\checkmarki\checkmarkr\checkmarkc\checkmark$\checkmarku\checkmark$\checkmark^\checkmark\\checkmarkc\checkmarki\checkmarkr\checkmarkc\checkmark$\checkmarkt\checkmark$\checkmark^\checkmark\\checkmarkc\checkmarki\checkmarkr\checkmarkc\checkmark$\checkmarki\checkmark$\checkmark^\checkmark\\checkmarkc\checkmarki\checkmarkr\checkmarkc\checkmark$\checkmarko\checkmark$\checkmark^\checkmark\\checkmarkc\checkmarki\checkmarkr\checkmarkc\checkmark$\checkmarkn\checkmark$\checkmark^\checkmark\\checkmarkc\checkmarki\checkmarkr\checkmarkc\checkmark$\checkmarks\checkmark$\checkmark^\checkmark\\checkmarkc\checkmarki\checkmarkr\checkmarkc\checkmark$\checkmark \checkmark$\checkmark^\checkmark\\checkmarkc\checkmarki\checkmarkr\checkmarkc\checkmark$\checkmarkd\checkmark$\checkmark^\checkmark\\checkmarkc\checkmarki\checkmarkr\checkmarkc\checkmark$\checkmarke\checkmark$\checkmark^\checkmark\\checkmarkc\checkmarki\checkmarkr\checkmarkc\checkmark$\checkmark \checkmark$\checkmark^\checkmark\\checkmarkc\checkmarki\checkmarkr\checkmarkc\checkmark$\checkmarkt\checkmark$\checkmark^\checkmark\\checkmarkc\checkmarki\checkmarkr\checkmarkc\checkmark$\checkmarkr\checkmark$\checkmark^\checkmark\\checkmarkc\checkmarki\checkmarkr\checkmarkc\checkmark$\checkmarka\checkmark$\checkmark^\checkmark\\checkmarkc\checkmarki\checkmarkr\checkmarkc\checkmark$\checkmarkn\checkmark$\checkmark^\checkmark\\checkmarkc\checkmarki\checkmarkr\checkmarkc\checkmark$\checkmarks\checkmark$\checkmark^\checkmark\\checkmarkc\checkmarki\checkmarkr\checkmarkc\checkmark$\checkmarkm\checkmark$\checkmark^\checkmark\\checkmarkc\checkmarki\checkmarkr\checkmarkc\checkmark$\checkmarki\checkmark$\checkmark^\checkmark\\checkmarkc\checkmarki\checkmarkr\checkmarkc\checkmark$\checkmarks\checkmark$\checkmark^\checkmark\\checkmarkc\checkmarki\checkmarkr\checkmarkc\checkmark$\checkmarks\checkmark$\checkmark^\checkmark\\checkmarkc\checkmarki\checkmarkr\checkmarkc\checkmark$\checkmarki\checkmark$\checkmark^\checkmark\\checkmarkc\checkmarki\checkmarkr\checkmarkc\checkmark$\checkmarko\checkmark$\checkmark^\checkmark\\checkmarkc\checkmarki\checkmarkr\checkmarkc\checkmark$\checkmarkn\checkmark$\checkmark^\checkmark\\checkmarkc\checkmarki\checkmarkr\checkmarkc\checkmark$\checkmark \checkmark$\checkmark^\checkmark\\checkmarkc\checkmarki\checkmarkr\checkmarkc\checkmark$\checkmarkl\checkmark$\checkmark^\checkmark\\checkmarkc\checkmarki\checkmarkr\checkmarkc\checkmark$\checkmarki\checkmark$\checkmark^\checkmark\\checkmarkc\checkmarki\checkmarkr\checkmarkc\checkmark$\checkmarkn\checkmark$\checkmark^\checkmark\\checkmarkc\checkmarki\checkmarkr\checkmarkc\checkmark$\checkmarké\checkmark$\checkmark^\checkmark\\checkmarkc\checkmarki\checkmarkr\checkmarkc\checkmark$\checkmarka\checkmark$\checkmark^\checkmark\\checkmarkc\checkmarki\checkmarkr\checkmarkc\checkmark$\checkmarki\checkmark$\checkmark^\checkmark\\checkmarkc\checkmarki\checkmarkr\checkmarkc\checkmark$\checkmarkr\checkmark$\checkmark^\checkmark\\checkmarkc\checkmarki\checkmarkr\checkmarkc\checkmark$\checkmarke\checkmark$\checkmark^\checkmark\\checkmarkc\checkmarki\checkmarkr\checkmarkc\checkmark$\checkmark}\checkmark$\checkmark^\checkmark\\checkmarkc\checkmarki\checkmarkr\checkmarkc\checkmark$\checkmark
\checkmark$\checkmark^\checkmark\\checkmarkc\checkmarki\checkmarkr\checkmarkc\checkmark$\checkmark\\checkmark$\checkmark^\checkmark\\checkmarkc\checkmarki\checkmarkr\checkmarkc\checkmark$\checkmarkl\checkmark$\checkmark^\checkmark\\checkmarkc\checkmarki\checkmarkr\checkmarkc\checkmark$\checkmarka\checkmark$\checkmark^\checkmark\\checkmarkc\checkmarki\checkmarkr\checkmarkc\checkmark$\checkmarkb\checkmark$\checkmark^\checkmark\\checkmarkc\checkmarki\checkmarkr\checkmarkc\checkmark$\checkmarke\checkmark$\checkmark^\checkmark\\checkmarkc\checkmarki\checkmarkr\checkmarkc\checkmark$\checkmarkl\checkmark$\checkmark^\checkmark\\checkmarkc\checkmarki\checkmarkr\checkmarkc\checkmark$\checkmark{\checkmark$\checkmark^\checkmark\\checkmarkc\checkmarki\checkmarkr\checkmarkc\checkmark$\checkmarkt\checkmark$\checkmark^\checkmark\\checkmarkc\checkmarki\checkmarkr\checkmarkc\checkmark$\checkmarka\checkmark$\checkmark^\checkmark\\checkmarkc\checkmarki\checkmarkr\checkmarkc\checkmark$\checkmarkb\checkmark$\checkmark^\checkmark\\checkmarkc\checkmarki\checkmarkr\checkmarkc\checkmark$\checkmark:\checkmark$\checkmark^\checkmark\\checkmarkc\checkmarki\checkmarkr\checkmarkc\checkmark$\checkmarkc\checkmark$\checkmark^\checkmark\\checkmarkc\checkmarki\checkmarkr\checkmarkc\checkmark$\checkmarkh\checkmark$\checkmark^\checkmark\\checkmarkc\checkmarki\checkmarkr\checkmarkc\checkmark$\checkmarko\checkmark$\checkmark^\checkmark\\checkmarkc\checkmarki\checkmarkr\checkmarkc\checkmark$\checkmarki\checkmark$\checkmark^\checkmark\\checkmarkc\checkmarki\checkmarkr\checkmarkc\checkmark$\checkmarkx\checkmark$\checkmark^\checkmark\\checkmarkc\checkmarki\checkmarkr\checkmarkc\checkmark$\checkmark_\checkmark$\checkmark^\checkmark\\checkmarkc\checkmarki\checkmarkr\checkmarkc\checkmark$\checkmarkv\checkmark$\checkmark^\checkmark\\checkmarkc\checkmarki\checkmarkr\checkmarkc\checkmark$\checkmarki\checkmark$\checkmark^\checkmark\\checkmarkc\checkmarki\checkmarkr\checkmarkc\checkmark$\checkmarks\checkmark$\checkmark^\checkmark\\checkmarkc\checkmarki\checkmarkr\checkmarkc\checkmark$\checkmark}\checkmark$\checkmark^\checkmark\\checkmarkc\checkmarki\checkmarkr\checkmarkc\checkmark$\checkmark
\checkmark$\checkmark^\checkmark\\checkmarkc\checkmarki\checkmarkr\checkmarkc\checkmark$\checkmark\\checkmark$\checkmark^\checkmark\\checkmarkc\checkmarki\checkmarkr\checkmarkc\checkmark$\checkmarke\checkmark$\checkmark^\checkmark\\checkmarkc\checkmarki\checkmarkr\checkmarkc\checkmark$\checkmarkn\checkmark$\checkmark^\checkmark\\checkmarkc\checkmarki\checkmarkr\checkmarkc\checkmark$\checkmarkd\checkmark$\checkmark^\checkmark\\checkmarkc\checkmarki\checkmarkr\checkmarkc\checkmark$\checkmark{\checkmark$\checkmark^\checkmark\\checkmarkc\checkmarki\checkmarkr\checkmarkc\checkmark$\checkmarkt\checkmark$\checkmark^\checkmark\\checkmarkc\checkmarki\checkmarkr\checkmarkc\checkmark$\checkmarka\checkmark$\checkmark^\checkmark\\checkmarkc\checkmarki\checkmarkr\checkmarkc\checkmark$\checkmarkb\checkmark$\checkmark^\checkmark\\checkmarkc\checkmarki\checkmarkr\checkmarkc\checkmark$\checkmarkl\checkmark$\checkmark^\checkmark\\checkmarkc\checkmarki\checkmarkr\checkmarkc\checkmark$\checkmarke\checkmark$\checkmark^\checkmark\\checkmarkc\checkmarki\checkmarkr\checkmarkc\checkmark$\checkmark}\checkmark$\checkmark^\checkmark\\checkmarkc\checkmarki\checkmarkr\checkmarkc\checkmark$\checkmark
\checkmark$\checkmark^\checkmark\\checkmarkc\checkmarki\checkmarkr\checkmarkc\checkmark$\checkmark
\checkmark$\checkmark^\checkmark\\checkmarkc\checkmarki\checkmarkr\checkmarkc\checkmark$\checkmark\\checkmark$\checkmark^\checkmark\\checkmarkc\checkmarki\checkmarkr\checkmarkc\checkmark$\checkmarkt\checkmark$\checkmark^\checkmark\\checkmarkc\checkmarki\checkmarkr\checkmarkc\checkmark$\checkmarke\checkmark$\checkmark^\checkmark\\checkmarkc\checkmarki\checkmarkr\checkmarkc\checkmark$\checkmarkx\checkmark$\checkmark^\checkmark\\checkmarkc\checkmarki\checkmarkr\checkmarkc\checkmark$\checkmarkt\checkmark$\checkmark^\checkmark\\checkmarkc\checkmarki\checkmarkr\checkmarkc\checkmark$\checkmarkb\checkmark$\checkmark^\checkmark\\checkmarkc\checkmarki\checkmarkr\checkmarkc\checkmark$\checkmarkf\checkmark$\checkmark^\checkmark\\checkmarkc\checkmarki\checkmarkr\checkmarkc\checkmark$\checkmark{\checkmark$\checkmark^\checkmark\\checkmarkc\checkmarki\checkmarkr\checkmarkc\checkmark$\checkmarkV\checkmark$\checkmark^\checkmark\\checkmarkc\checkmarki\checkmarkr\checkmarkc\checkmark$\checkmarki\checkmark$\checkmark^\checkmark\\checkmarkc\checkmarki\checkmarkr\checkmarkc\checkmark$\checkmarks\checkmark$\checkmark^\checkmark\\checkmarkc\checkmarki\checkmarkr\checkmarkc\checkmark$\checkmark \checkmark$\checkmark^\checkmark\\checkmarkc\checkmarki\checkmarkr\checkmarkc\checkmark$\checkmarkà\checkmark$\checkmark^\checkmark\\checkmarkc\checkmarki\checkmarkr\checkmarkc\checkmark$\checkmark \checkmark$\checkmark^\checkmark\\checkmarkc\checkmarki\checkmarkr\checkmarkc\checkmark$\checkmarkb\checkmark$\checkmark^\checkmark\\checkmarkc\checkmarki\checkmarkr\checkmarkc\checkmark$\checkmarki\checkmark$\checkmark^\checkmark\\checkmarkc\checkmarki\checkmarkr\checkmarkc\checkmark$\checkmarkl\checkmark$\checkmark^\checkmark\\checkmarkc\checkmarki\checkmarkr\checkmarkc\checkmark$\checkmarkl\checkmark$\checkmark^\checkmark\\checkmarkc\checkmarki\checkmarkr\checkmarkc\checkmark$\checkmarke\checkmark$\checkmark^\checkmark\\checkmarkc\checkmarki\checkmarkr\checkmarkc\checkmark$\checkmarks\checkmark$\checkmark^\checkmark\\checkmarkc\checkmarki\checkmarkr\checkmarkc\checkmark$\checkmark \checkmark$\checkmark^\checkmark\\checkmarkc\checkmarki\checkmarkr\checkmarkc\checkmark$\checkmarkr\checkmark$\checkmark^\checkmark\\checkmarkc\checkmarki\checkmarkr\checkmarkc\checkmark$\checkmarke\checkmark$\checkmark^\checkmark\\checkmarkc\checkmarki\checkmarkr\checkmarkc\checkmark$\checkmarkt\checkmark$\checkmark^\checkmark\\checkmarkc\checkmarki\checkmarkr\checkmarkc\checkmark$\checkmarke\checkmark$\checkmark^\checkmark\\checkmarkc\checkmarki\checkmarkr\checkmarkc\checkmark$\checkmarkn\checkmark$\checkmark^\checkmark\\checkmarkc\checkmarki\checkmarkr\checkmarkc\checkmark$\checkmarku\checkmark$\checkmark^\checkmark\\checkmarkc\checkmarki\checkmarkr\checkmarkc\checkmark$\checkmarke\checkmark$\checkmark^\checkmark\\checkmarkc\checkmarki\checkmarkr\checkmarkc\checkmark$\checkmark \checkmark$\checkmark^\checkmark\\checkmarkc\checkmarki\checkmarkr\checkmarkc\checkmark$\checkmark:\checkmark$\checkmark^\checkmark\\checkmarkc\checkmarki\checkmarkr\checkmarkc\checkmark$\checkmark \checkmark$\checkmark^\checkmark\\checkmarkc\checkmarki\checkmarkr\checkmarkc\checkmark$\checkmarkS\checkmark$\checkmark^\checkmark\\checkmarkc\checkmarki\checkmarkr\checkmarkc\checkmark$\checkmarkF\checkmark$\checkmark^\checkmark\\checkmarkc\checkmarki\checkmarkr\checkmarkc\checkmark$\checkmarkU\checkmark$\checkmark^\checkmark\\checkmarkc\checkmarki\checkmarkr\checkmarkc\checkmark$\checkmark1\checkmark$\checkmark^\checkmark\\checkmarkc\checkmarki\checkmarkr\checkmarkc\checkmark$\checkmark6\checkmark$\checkmark^\checkmark\\checkmarkc\checkmarki\checkmarkr\checkmarkc\checkmark$\checkmark0\checkmark$\checkmark^\checkmark\\checkmarkc\checkmarki\checkmarkr\checkmarkc\checkmark$\checkmark5\checkmark$\checkmark^\checkmark\\checkmarkc\checkmarki\checkmarkr\checkmarkc\checkmark$\checkmark}\checkmark$\checkmark^\checkmark\\checkmarkc\checkmarki\checkmarkr\checkmarkc\checkmark$\checkmark
\checkmark$\checkmark^\checkmark\\checkmarkc\checkmarki\checkmarkr\checkmarkc\checkmark$\checkmark\\checkmark$\checkmark^\checkmark\\checkmarkc\checkmarki\checkmarkr\checkmarkc\checkmark$\checkmarkb\checkmark$\checkmark^\checkmark\\checkmarkc\checkmarki\checkmarkr\checkmarkc\checkmark$\checkmarke\checkmark$\checkmark^\checkmark\\checkmarkc\checkmarki\checkmarkr\checkmarkc\checkmark$\checkmarkg\checkmark$\checkmark^\checkmark\\checkmarkc\checkmarki\checkmarkr\checkmarkc\checkmark$\checkmarki\checkmark$\checkmark^\checkmark\\checkmarkc\checkmarki\checkmarkr\checkmarkc\checkmark$\checkmarkn\checkmark$\checkmark^\checkmark\\checkmarkc\checkmarki\checkmarkr\checkmarkc\checkmark$\checkmark{\checkmark$\checkmark^\checkmark\\checkmarkc\checkmarki\checkmarkr\checkmarkc\checkmark$\checkmarki\checkmark$\checkmark^\checkmark\\checkmarkc\checkmarki\checkmarkr\checkmarkc\checkmark$\checkmarkt\checkmark$\checkmark^\checkmark\\checkmarkc\checkmarki\checkmarkr\checkmarkc\checkmark$\checkmarke\checkmark$\checkmark^\checkmark\\checkmarkc\checkmarki\checkmarkr\checkmarkc\checkmark$\checkmarkm\checkmark$\checkmark^\checkmark\\checkmarkc\checkmarki\checkmarkr\checkmarkc\checkmark$\checkmarki\checkmark$\checkmark^\checkmark\\checkmarkc\checkmarki\checkmarkr\checkmarkc\checkmark$\checkmarkz\checkmark$\checkmark^\checkmark\\checkmarkc\checkmarki\checkmarkr\checkmarkc\checkmark$\checkmarke\checkmark$\checkmark^\checkmark\\checkmarkc\checkmarki\checkmarkr\checkmarkc\checkmark$\checkmark}\checkmark$\checkmark^\checkmark\\checkmarkc\checkmarki\checkmarkr\checkmarkc\checkmark$\checkmark
\checkmark$\checkmark^\checkmark\\checkmarkc\checkmarki\checkmarkr\checkmarkc\checkmark$\checkmark \checkmark$\checkmark^\checkmark\\checkmarkc\checkmarki\checkmarkr\checkmarkc\checkmark$\checkmark \checkmark$\checkmark^\checkmark\\checkmarkc\checkmarki\checkmarkr\checkmarkc\checkmark$\checkmark \checkmark$\checkmark^\checkmark\\checkmarkc\checkmarki\checkmarkr\checkmarkc\checkmark$\checkmark \checkmark$\checkmark^\checkmark\\checkmarkc\checkmarki\checkmarkr\checkmarkc\checkmark$\checkmark\\checkmark$\checkmark^\checkmark\\checkmarkc\checkmarki\checkmarkr\checkmarkc\checkmark$\checkmarki\checkmark$\checkmark^\checkmark\\checkmarkc\checkmarki\checkmarkr\checkmarkc\checkmark$\checkmarkt\checkmark$\checkmark^\checkmark\\checkmarkc\checkmarki\checkmarkr\checkmarkc\checkmark$\checkmarke\checkmark$\checkmark^\checkmark\\checkmarkc\checkmarki\checkmarkr\checkmarkc\checkmark$\checkmarkm\checkmark$\checkmark^\checkmark\\checkmarkc\checkmarki\checkmarkr\checkmarkc\checkmark$\checkmark \checkmark$\checkmark^\checkmark\\checkmarkc\checkmarki\checkmarkr\checkmarkc\checkmark$\checkmarkD\checkmark$\checkmark^\checkmark\\checkmarkc\checkmarki\checkmarkr\checkmarkc\checkmark$\checkmarki\checkmark$\checkmark^\checkmark\\checkmarkc\checkmarki\checkmarkr\checkmarkc\checkmark$\checkmarka\checkmark$\checkmark^\checkmark\\checkmarkc\checkmarki\checkmarkr\checkmarkc\checkmark$\checkmarkm\checkmark$\checkmark^\checkmark\\checkmarkc\checkmarki\checkmarkr\checkmarkc\checkmark$\checkmarkè\checkmark$\checkmark^\checkmark\\checkmarkc\checkmarki\checkmarkr\checkmarkc\checkmark$\checkmarkt\checkmark$\checkmark^\checkmark\\checkmarkc\checkmarki\checkmarkr\checkmarkc\checkmark$\checkmarkr\checkmark$\checkmark^\checkmark\\checkmarkc\checkmarki\checkmarkr\checkmarkc\checkmark$\checkmarke\checkmark$\checkmark^\checkmark\\checkmarkc\checkmarki\checkmarkr\checkmarkc\checkmark$\checkmark \checkmark$\checkmark^\checkmark\\checkmarkc\checkmarki\checkmarkr\checkmarkc\checkmark$\checkmarkn\checkmark$\checkmark^\checkmark\\checkmarkc\checkmarki\checkmarkr\checkmarkc\checkmark$\checkmarko\checkmark$\checkmark^\checkmark\\checkmarkc\checkmarki\checkmarkr\checkmarkc\checkmark$\checkmarkm\checkmark$\checkmark^\checkmark\\checkmarkc\checkmarki\checkmarkr\checkmarkc\checkmark$\checkmarki\checkmark$\checkmark^\checkmark\\checkmarkc\checkmarki\checkmarkr\checkmarkc\checkmark$\checkmarkn\checkmark$\checkmark^\checkmark\\checkmarkc\checkmarki\checkmarkr\checkmarkc\checkmark$\checkmarka\checkmark$\checkmark^\checkmark\\checkmarkc\checkmarki\checkmarkr\checkmarkc\checkmark$\checkmarkl\checkmark$\checkmark^\checkmark\\checkmarkc\checkmarki\checkmarkr\checkmarkc\checkmark$\checkmark \checkmark$\checkmark^\checkmark\\checkmarkc\checkmarki\checkmarkr\checkmarkc\checkmark$\checkmark:\checkmark$\checkmark^\checkmark\\checkmarkc\checkmarki\checkmarkr\checkmarkc\checkmark$\checkmark \checkmark$\checkmark^\checkmark\\checkmarkc\checkmarki\checkmarkr\checkmarkc\checkmark$\checkmark$\checkmark$\checkmark^\checkmark\\checkmarkc\checkmarki\checkmarkr\checkmarkc\checkmark$\checkmarkD\checkmark$\checkmark^\checkmark\\checkmarkc\checkmarki\checkmarkr\checkmarkc\checkmark$\checkmark_\checkmark$\checkmark^\checkmark\\checkmarkc\checkmarki\checkmarkr\checkmarkc\checkmark$\checkmarkN\checkmark$\checkmark^\checkmark\\checkmarkc\checkmarki\checkmarkr\checkmarkc\checkmark$\checkmark \checkmark$\checkmark^\checkmark\\checkmarkc\checkmarki\checkmarkr\checkmarkc\checkmark$\checkmark=\checkmark$\checkmark^\checkmark\\checkmarkc\checkmarki\checkmarkr\checkmarkc\checkmark$\checkmark \checkmark$\checkmark^\checkmark\\checkmarkc\checkmarki\checkmarkr\checkmarkc\checkmark$\checkmark1\checkmark$\checkmark^\checkmark\\checkmarkc\checkmarki\checkmarkr\checkmarkc\checkmark$\checkmark6\checkmark$\checkmark^\checkmark\\checkmarkc\checkmarki\checkmarkr\checkmarkc\checkmark$\checkmark$\checkmark$\checkmark^\checkmark\\checkmarkc\checkmarki\checkmarkr\checkmarkc\checkmark$\checkmark \checkmark$\checkmark^\checkmark\\checkmarkc\checkmarki\checkmarkr\checkmarkc\checkmark$\checkmarkm\checkmark$\checkmark^\checkmark\\checkmarkc\checkmarki\checkmarkr\checkmarkc\checkmark$\checkmarkm\checkmark$\checkmark^\checkmark\\checkmarkc\checkmarki\checkmarkr\checkmarkc\checkmark$\checkmark
\checkmark$\checkmark^\checkmark\\checkmarkc\checkmarki\checkmarkr\checkmarkc\checkmark$\checkmark \checkmark$\checkmark^\checkmark\\checkmarkc\checkmarki\checkmarkr\checkmarkc\checkmark$\checkmark \checkmark$\checkmark^\checkmark\\checkmarkc\checkmarki\checkmarkr\checkmarkc\checkmark$\checkmark \checkmark$\checkmark^\checkmark\\checkmarkc\checkmarki\checkmarkr\checkmarkc\checkmark$\checkmark \checkmark$\checkmark^\checkmark\\checkmarkc\checkmarki\checkmarkr\checkmarkc\checkmark$\checkmark\\checkmark$\checkmark^\checkmark\\checkmarkc\checkmarki\checkmarkr\checkmarkc\checkmark$\checkmarki\checkmark$\checkmark^\checkmark\\checkmarkc\checkmarki\checkmarkr\checkmarkc\checkmark$\checkmarkt\checkmark$\checkmark^\checkmark\\checkmarkc\checkmarki\checkmarkr\checkmarkc\checkmark$\checkmarke\checkmark$\checkmark^\checkmark\\checkmarkc\checkmarki\checkmarkr\checkmarkc\checkmark$\checkmarkm\checkmark$\checkmark^\checkmark\\checkmarkc\checkmarki\checkmarkr\checkmarkc\checkmark$\checkmark \checkmark$\checkmark^\checkmark\\checkmarkc\checkmarki\checkmarkr\checkmarkc\checkmark$\checkmarkP\checkmark$\checkmark^\checkmark\\checkmarkc\checkmarki\checkmarkr\checkmarkc\checkmark$\checkmarka\checkmark$\checkmark^\checkmark\\checkmarkc\checkmarki\checkmarkr\checkmarkc\checkmark$\checkmarks\checkmark$\checkmark^\checkmark\\checkmarkc\checkmarki\checkmarkr\checkmarkc\checkmark$\checkmark \checkmark$\checkmark^\checkmark\\checkmarkc\checkmarki\checkmarkr\checkmarkc\checkmark$\checkmarkd\checkmark$\checkmark^\checkmark\\checkmarkc\checkmarki\checkmarkr\checkmarkc\checkmark$\checkmarke\checkmark$\checkmark^\checkmark\\checkmarkc\checkmarki\checkmarkr\checkmarkc\checkmark$\checkmark \checkmark$\checkmark^\checkmark\\checkmarkc\checkmarki\checkmarkr\checkmarkc\checkmark$\checkmarkl\checkmark$\checkmark^\checkmark\\checkmarkc\checkmarki\checkmarkr\checkmarkc\checkmark$\checkmarka\checkmark$\checkmark^\checkmark\\checkmarkc\checkmarki\checkmarkr\checkmarkc\checkmark$\checkmark \checkmark$\checkmark^\checkmark\\checkmarkc\checkmarki\checkmarkr\checkmarkc\checkmark$\checkmarkv\checkmark$\checkmark^\checkmark\\checkmarkc\checkmarki\checkmarkr\checkmarkc\checkmark$\checkmarki\checkmark$\checkmark^\checkmark\\checkmarkc\checkmarki\checkmarkr\checkmarkc\checkmark$\checkmarks\checkmark$\checkmark^\checkmark\\checkmarkc\checkmarki\checkmarkr\checkmarkc\checkmark$\checkmark \checkmark$\checkmark^\checkmark\\checkmarkc\checkmarki\checkmarkr\checkmarkc\checkmark$\checkmark:\checkmark$\checkmark^\checkmark\\checkmarkc\checkmarki\checkmarkr\checkmarkc\checkmark$\checkmark \checkmark$\checkmark^\checkmark\\checkmarkc\checkmarki\checkmarkr\checkmarkc\checkmark$\checkmark$\checkmark$\checkmark^\checkmark\\checkmarkc\checkmarki\checkmarkr\checkmarkc\checkmark$\checkmarkP\checkmark$\checkmark^\checkmark\\checkmarkc\checkmarki\checkmarkr\checkmarkc\checkmark$\checkmark_\checkmark$\checkmark^\checkmark\\checkmarkc\checkmarki\checkmarkr\checkmarkc\checkmark$\checkmarkh\checkmark$\checkmark^\checkmark\\checkmarkc\checkmarki\checkmarkr\checkmarkc\checkmark$\checkmark \checkmark$\checkmark^\checkmark\\checkmarkc\checkmarki\checkmarkr\checkmarkc\checkmark$\checkmark=\checkmark$\checkmark^\checkmark\\checkmarkc\checkmarki\checkmarkr\checkmarkc\checkmark$\checkmark \checkmark$\checkmark^\checkmark\\checkmarkc\checkmarki\checkmarkr\checkmarkc\checkmark$\checkmark5\checkmark$\checkmark^\checkmark\\checkmarkc\checkmarki\checkmarkr\checkmarkc\checkmark$\checkmark$\checkmark$\checkmark^\checkmark\\checkmarkc\checkmarki\checkmarkr\checkmarkc\checkmark$\checkmark \checkmark$\checkmark^\checkmark\\checkmarkc\checkmarki\checkmarkr\checkmarkc\checkmark$\checkmarkm\checkmark$\checkmark^\checkmark\\checkmarkc\checkmarki\checkmarkr\checkmarkc\checkmark$\checkmarkm\checkmark$\checkmark^\checkmark\\checkmarkc\checkmarki\checkmarkr\checkmarkc\checkmark$\checkmark/\checkmark$\checkmark^\checkmark\\checkmarkc\checkmarki\checkmarkr\checkmarkc\checkmark$\checkmarkt\checkmark$\checkmark^\checkmark\\checkmarkc\checkmarki\checkmarkr\checkmarkc\checkmark$\checkmarko\checkmark$\checkmark^\checkmark\\checkmarkc\checkmarki\checkmarkr\checkmarkc\checkmark$\checkmarku\checkmark$\checkmark^\checkmark\\checkmarkc\checkmarki\checkmarkr\checkmarkc\checkmark$\checkmarkr\checkmark$\checkmark^\checkmark\\checkmarkc\checkmarki\checkmarkr\checkmarkc\checkmark$\checkmark
\checkmark$\checkmark^\checkmark\\checkmarkc\checkmarki\checkmarkr\checkmarkc\checkmark$\checkmark \checkmark$\checkmark^\checkmark\\checkmarkc\checkmarki\checkmarkr\checkmarkc\checkmark$\checkmark \checkmark$\checkmark^\checkmark\\checkmarkc\checkmarki\checkmarkr\checkmarkc\checkmark$\checkmark \checkmark$\checkmark^\checkmark\\checkmarkc\checkmarki\checkmarkr\checkmarkc\checkmark$\checkmark \checkmark$\checkmark^\checkmark\\checkmarkc\checkmarki\checkmarkr\checkmarkc\checkmark$\checkmark\\checkmark$\checkmark^\checkmark\\checkmarkc\checkmarki\checkmarkr\checkmarkc\checkmark$\checkmarki\checkmark$\checkmark^\checkmark\\checkmarkc\checkmarki\checkmarkr\checkmarkc\checkmark$\checkmarkt\checkmark$\checkmark^\checkmark\\checkmarkc\checkmarki\checkmarkr\checkmarkc\checkmark$\checkmarke\checkmark$\checkmark^\checkmark\\checkmarkc\checkmarki\checkmarkr\checkmarkc\checkmark$\checkmarkm\checkmark$\checkmark^\checkmark\\checkmarkc\checkmarki\checkmarkr\checkmarkc\checkmark$\checkmark \checkmark$\checkmark^\checkmark\\checkmarkc\checkmarki\checkmarkr\checkmarkc\checkmark$\checkmarkC\checkmark$\checkmark^\checkmark\\checkmarkc\checkmarki\checkmarkr\checkmarkc\checkmark$\checkmarkh\checkmark$\checkmark^\checkmark\\checkmarkc\checkmarki\checkmarkr\checkmarkc\checkmark$\checkmarka\checkmark$\checkmark^\checkmark\\checkmarkc\checkmarki\checkmarkr\checkmarkc\checkmark$\checkmarkr\checkmark$\checkmark^\checkmark\\checkmarkc\checkmarki\checkmarkr\checkmarkc\checkmark$\checkmarkg\checkmark$\checkmark^\checkmark\\checkmarkc\checkmarki\checkmarkr\checkmarkc\checkmark$\checkmarke\checkmark$\checkmark^\checkmark\\checkmarkc\checkmarki\checkmarkr\checkmarkc\checkmark$\checkmark \checkmark$\checkmark^\checkmark\\checkmarkc\checkmarki\checkmarkr\checkmarkc\checkmark$\checkmarkd\checkmark$\checkmark^\checkmark\\checkmarkc\checkmarki\checkmarkr\checkmarkc\checkmark$\checkmarky\checkmark$\checkmark^\checkmark\\checkmarkc\checkmarki\checkmarkr\checkmarkc\checkmark$\checkmarkn\checkmark$\checkmark^\checkmark\\checkmarkc\checkmarki\checkmarkr\checkmarkc\checkmark$\checkmarka\checkmark$\checkmark^\checkmark\\checkmarkc\checkmarki\checkmarkr\checkmarkc\checkmark$\checkmarkm\checkmark$\checkmark^\checkmark\\checkmarkc\checkmarki\checkmarkr\checkmarkc\checkmark$\checkmarki\checkmark$\checkmark^\checkmark\\checkmarkc\checkmarki\checkmarkr\checkmarkc\checkmark$\checkmarkq\checkmark$\checkmark^\checkmark\\checkmarkc\checkmarki\checkmarkr\checkmarkc\checkmark$\checkmarku\checkmark$\checkmark^\checkmark\\checkmarkc\checkmarki\checkmarkr\checkmarkc\checkmark$\checkmarke\checkmark$\checkmark^\checkmark\\checkmarkc\checkmarki\checkmarkr\checkmarkc\checkmark$\checkmark \checkmark$\checkmark^\checkmark\\checkmarkc\checkmarki\checkmarkr\checkmarkc\checkmark$\checkmarkd\checkmark$\checkmark^\checkmark\\checkmarkc\checkmarki\checkmarkr\checkmarkc\checkmark$\checkmarke\checkmark$\checkmark^\checkmark\\checkmarkc\checkmarki\checkmarkr\checkmarkc\checkmark$\checkmark \checkmark$\checkmark^\checkmark\\checkmarkc\checkmarki\checkmarkr\checkmarkc\checkmark$\checkmarkb\checkmark$\checkmark^\checkmark\\checkmarkc\checkmarki\checkmarkr\checkmarkc\checkmark$\checkmarka\checkmark$\checkmark^\checkmark\\checkmarkc\checkmarki\checkmarkr\checkmarkc\checkmark$\checkmarks\checkmark$\checkmark^\checkmark\\checkmarkc\checkmarki\checkmarkr\checkmarkc\checkmark$\checkmarke\checkmark$\checkmark^\checkmark\\checkmarkc\checkmarki\checkmarkr\checkmarkc\checkmark$\checkmark \checkmark$\checkmark^\checkmark\\checkmarkc\checkmarki\checkmarkr\checkmarkc\checkmark$\checkmark:\checkmark$\checkmark^\checkmark\\checkmarkc\checkmarki\checkmarkr\checkmarkc\checkmark$\checkmark \checkmark$\checkmark^\checkmark\\checkmarkc\checkmarki\checkmarkr\checkmarkc\checkmark$\checkmark$\checkmark$\checkmark^\checkmark\\checkmarkc\checkmarki\checkmarkr\checkmarkc\checkmark$\checkmarkC\checkmark$\checkmark^\checkmark\\checkmarkc\checkmarki\checkmarkr\checkmarkc\checkmark$\checkmark_\checkmark$\checkmark^\checkmark\\checkmarkc\checkmarki\checkmarkr\checkmarkc\checkmark$\checkmark{\checkmark$\checkmark^\checkmark\\checkmarkc\checkmarki\checkmarkr\checkmarkc\checkmark$\checkmarkd\checkmark$\checkmark^\checkmark\\checkmarkc\checkmarki\checkmarkr\checkmarkc\checkmark$\checkmarky\checkmark$\checkmark^\checkmark\\checkmarkc\checkmarki\checkmarkr\checkmarkc\checkmark$\checkmarkn\checkmark$\checkmark^\checkmark\\checkmarkc\checkmarki\checkmarkr\checkmarkc\checkmark$\checkmark}\checkmark$\checkmark^\checkmark\\checkmarkc\checkmarki\checkmarkr\checkmarkc\checkmark$\checkmark \checkmark$\checkmark^\checkmark\\checkmarkc\checkmarki\checkmarkr\checkmarkc\checkmark$\checkmark=\checkmark$\checkmark^\checkmark\\checkmarkc\checkmarki\checkmarkr\checkmarkc\checkmark$\checkmark \checkmark$\checkmark^\checkmark\\checkmarkc\checkmarki\checkmarkr\checkmarkc\checkmark$\checkmark6\checkmark$\checkmark^\checkmark\\checkmarkc\checkmarki\checkmarkr\checkmarkc\checkmark$\checkmark\\checkmark$\checkmark^\checkmark\\checkmarkc\checkmarki\checkmarkr\checkmarkc\checkmark$\checkmark,\checkmark$\checkmark^\checkmark\\checkmarkc\checkmarki\checkmarkr\checkmarkc\checkmark$\checkmark8\checkmark$\checkmark^\checkmark\\checkmarkc\checkmarki\checkmarkr\checkmarkc\checkmark$\checkmark0\checkmark$\checkmark^\checkmark\\checkmarkc\checkmarki\checkmarkr\checkmarkc\checkmark$\checkmark0\checkmark$\checkmark^\checkmark\\checkmarkc\checkmarki\checkmarkr\checkmarkc\checkmark$\checkmark$\checkmark$\checkmark^\checkmark\\checkmarkc\checkmarki\checkmarkr\checkmarkc\checkmark$\checkmark \checkmark$\checkmark^\checkmark\\checkmarkc\checkmarki\checkmarkr\checkmarkc\checkmark$\checkmarkN\checkmark$\checkmark^\checkmark\\checkmarkc\checkmarki\checkmarkr\checkmarkc\checkmark$\checkmark
\checkmark$\checkmark^\checkmark\\checkmarkc\checkmarki\checkmarkr\checkmarkc\checkmark$\checkmark \checkmark$\checkmark^\checkmark\\checkmarkc\checkmarki\checkmarkr\checkmarkc\checkmark$\checkmark \checkmark$\checkmark^\checkmark\\checkmarkc\checkmarki\checkmarkr\checkmarkc\checkmark$\checkmark \checkmark$\checkmark^\checkmark\\checkmarkc\checkmarki\checkmarkr\checkmarkc\checkmark$\checkmark \checkmark$\checkmark^\checkmark\\checkmarkc\checkmarki\checkmarkr\checkmarkc\checkmark$\checkmark\\checkmark$\checkmark^\checkmark\\checkmarkc\checkmarki\checkmarkr\checkmarkc\checkmark$\checkmarki\checkmark$\checkmark^\checkmark\\checkmarkc\checkmarki\checkmarkr\checkmarkc\checkmark$\checkmarkt\checkmark$\checkmark^\checkmark\\checkmarkc\checkmarki\checkmarkr\checkmarkc\checkmark$\checkmarke\checkmark$\checkmark^\checkmark\\checkmarkc\checkmarki\checkmarkr\checkmarkc\checkmark$\checkmarkm\checkmark$\checkmark^\checkmark\\checkmarkc\checkmarki\checkmarkr\checkmarkc\checkmark$\checkmark \checkmark$\checkmark^\checkmark\\checkmarkc\checkmarki\checkmarkr\checkmarkc\checkmark$\checkmarkM\checkmark$\checkmark^\checkmark\\checkmarkc\checkmarki\checkmarkr\checkmarkc\checkmark$\checkmarka\checkmark$\checkmark^\checkmark\\checkmarkc\checkmarki\checkmarkr\checkmarkc\checkmark$\checkmarkt\checkmark$\checkmark^\checkmark\\checkmarkc\checkmarki\checkmarkr\checkmarkc\checkmark$\checkmarké\checkmark$\checkmark^\checkmark\\checkmarkc\checkmarki\checkmarkr\checkmarkc\checkmark$\checkmarkr\checkmark$\checkmark^\checkmark\\checkmarkc\checkmarki\checkmarkr\checkmarkc\checkmark$\checkmarki\checkmark$\checkmark^\checkmark\\checkmarkc\checkmarki\checkmarkr\checkmarkc\checkmark$\checkmarka\checkmark$\checkmark^\checkmark\\checkmarkc\checkmarki\checkmarkr\checkmarkc\checkmark$\checkmarku\checkmark$\checkmark^\checkmark\\checkmarkc\checkmarki\checkmarkr\checkmarkc\checkmark$\checkmark \checkmark$\checkmark^\checkmark\\checkmarkc\checkmarki\checkmarkr\checkmarkc\checkmark$\checkmark:\checkmark$\checkmark^\checkmark\\checkmarkc\checkmarki\checkmarkr\checkmarkc\checkmark$\checkmark \checkmark$\checkmark^\checkmark\\checkmarkc\checkmarki\checkmarkr\checkmarkc\checkmark$\checkmarkA\checkmark$\checkmark^\checkmark\\checkmarkc\checkmarki\checkmarkr\checkmarkc\checkmark$\checkmarkc\checkmark$\checkmark^\checkmark\\checkmarkc\checkmarki\checkmarkr\checkmarkc\checkmark$\checkmarki\checkmark$\checkmark^\checkmark\\checkmarkc\checkmarki\checkmarkr\checkmarkc\checkmark$\checkmarke\checkmark$\checkmark^\checkmark\\checkmarkc\checkmarki\checkmarkr\checkmarkc\checkmark$\checkmarkr\checkmark$\checkmark^\checkmark\\checkmarkc\checkmarki\checkmarkr\checkmarkc\checkmark$\checkmark \checkmark$\checkmark^\checkmark\\checkmarkc\checkmarki\checkmarkr\checkmarkc\checkmark$\checkmark5\checkmark$\checkmark^\checkmark\\checkmarkc\checkmarki\checkmarkr\checkmarkc\checkmark$\checkmark2\checkmark$\checkmark^\checkmark\\checkmarkc\checkmarki\checkmarkr\checkmarkc\checkmark$\checkmark1\checkmark$\checkmark^\checkmark\\checkmarkc\checkmarki\checkmarkr\checkmarkc\checkmark$\checkmark0\checkmark$\checkmark^\checkmark\\checkmarkc\checkmarki\checkmarkr\checkmarkc\checkmark$\checkmark0\checkmark$\checkmark^\checkmark\\checkmarkc\checkmarki\checkmarkr\checkmarkc\checkmark$\checkmark \checkmark$\checkmark^\checkmark\\checkmarkc\checkmarki\checkmarkr\checkmarkc\checkmark$\checkmark(\checkmark$\checkmark^\checkmark\\checkmarkc\checkmarki\checkmarkr\checkmarkc\checkmark$\checkmark$\checkmark$\checkmark^\checkmark\\checkmarkc\checkmarki\checkmarkr\checkmarkc\checkmark$\checkmarkR\checkmark$\checkmark^\checkmark\\checkmarkc\checkmarki\checkmarkr\checkmarkc\checkmark$\checkmark_\checkmark$\checkmark^\checkmark\\checkmarkc\checkmarki\checkmarkr\checkmarkc\checkmark$\checkmarkm\checkmark$\checkmark^\checkmark\\checkmarkc\checkmarki\checkmarkr\checkmarkc\checkmark$\checkmark \checkmark$\checkmark^\checkmark\\checkmarkc\checkmarki\checkmarkr\checkmarkc\checkmark$\checkmark=\checkmark$\checkmark^\checkmark\\checkmarkc\checkmarki\checkmarkr\checkmarkc\checkmark$\checkmark \checkmark$\checkmark^\checkmark\\checkmarkc\checkmarki\checkmarkr\checkmarkc\checkmark$\checkmark2\checkmark$\checkmark^\checkmark\\checkmarkc\checkmarki\checkmarkr\checkmarkc\checkmark$\checkmark\\checkmark$\checkmark^\checkmark\\checkmarkc\checkmarki\checkmarkr\checkmarkc\checkmark$\checkmark,\checkmark$\checkmark^\checkmark\\checkmarkc\checkmarki\checkmarkr\checkmarkc\checkmark$\checkmark0\checkmark$\checkmark^\checkmark\\checkmarkc\checkmarki\checkmarkr\checkmarkc\checkmark$\checkmark0\checkmark$\checkmark^\checkmark\\checkmarkc\checkmarki\checkmarkr\checkmarkc\checkmark$\checkmark0\checkmark$\checkmark^\checkmark\\checkmarkc\checkmarki\checkmarkr\checkmarkc\checkmark$\checkmark$\checkmark$\checkmark^\checkmark\\checkmarkc\checkmarki\checkmarkr\checkmarkc\checkmark$\checkmark \checkmark$\checkmark^\checkmark\\checkmarkc\checkmarki\checkmarkr\checkmarkc\checkmark$\checkmarkM\checkmark$\checkmark^\checkmark\\checkmarkc\checkmarki\checkmarkr\checkmarkc\checkmark$\checkmarkP\checkmark$\checkmark^\checkmark\\checkmarkc\checkmarki\checkmarkr\checkmarkc\checkmark$\checkmarka\checkmark$\checkmark^\checkmark\\checkmarkc\checkmarki\checkmarkr\checkmarkc\checkmark$\checkmark,\checkmark$\checkmark^\checkmark\\checkmarkc\checkmarki\checkmarkr\checkmarkc\checkmark$\checkmark \checkmark$\checkmark^\checkmark\\checkmarkc\checkmarki\checkmarkr\checkmarkc\checkmark$\checkmark$\checkmark$\checkmark^\checkmark\\checkmarkc\checkmarki\checkmarkr\checkmarkc\checkmark$\checkmarkH\checkmark$\checkmark^\checkmark\\checkmarkc\checkmarki\checkmarkr\checkmarkc\checkmark$\checkmarkR\checkmark$\checkmark^\checkmark\\checkmarkc\checkmarki\checkmarkr\checkmarkc\checkmark$\checkmarkC\checkmark$\checkmark^\checkmark\\checkmarkc\checkmarki\checkmarkr\checkmarkc\checkmark$\checkmark \checkmark$\checkmark^\checkmark\\checkmarkc\checkmarki\checkmarkr\checkmarkc\checkmark$\checkmark>\checkmark$\checkmark^\checkmark\\checkmarkc\checkmarki\checkmarkr\checkmarkc\checkmark$\checkmark \checkmark$\checkmark^\checkmark\\checkmarkc\checkmarki\checkmarkr\checkmarkc\checkmark$\checkmark5\checkmark$\checkmark^\checkmark\\checkmarkc\checkmarki\checkmarkr\checkmarkc\checkmark$\checkmark8\checkmark$\checkmark^\checkmark\\checkmarkc\checkmarki\checkmarkr\checkmarkc\checkmark$\checkmark$\checkmark$\checkmark^\checkmark\\checkmarkc\checkmarki\checkmarkr\checkmarkc\checkmark$\checkmark)\checkmark$\checkmark^\checkmark\\checkmarkc\checkmarki\checkmarkr\checkmarkc\checkmark$\checkmark
\checkmark$\checkmark^\checkmark\\checkmarkc\checkmarki\checkmarkr\checkmarkc\checkmark$\checkmark\\checkmark$\checkmark^\checkmark\\checkmarkc\checkmarki\checkmarkr\checkmarkc\checkmark$\checkmarke\checkmark$\checkmark^\checkmark\\checkmarkc\checkmarki\checkmarkr\checkmarkc\checkmark$\checkmarkn\checkmark$\checkmark^\checkmark\\checkmarkc\checkmarki\checkmarkr\checkmarkc\checkmark$\checkmarkd\checkmark$\checkmark^\checkmark\\checkmarkc\checkmarki\checkmarkr\checkmarkc\checkmark$\checkmark{\checkmark$\checkmark^\checkmark\\checkmarkc\checkmarki\checkmarkr\checkmarkc\checkmark$\checkmarki\checkmark$\checkmark^\checkmark\\checkmarkc\checkmarki\checkmarkr\checkmarkc\checkmark$\checkmarkt\checkmark$\checkmark^\checkmark\\checkmarkc\checkmarki\checkmarkr\checkmarkc\checkmark$\checkmarke\checkmark$\checkmark^\checkmark\\checkmarkc\checkmarki\checkmarkr\checkmarkc\checkmark$\checkmarkm\checkmark$\checkmark^\checkmark\\checkmarkc\checkmarki\checkmarkr\checkmarkc\checkmark$\checkmarki\checkmark$\checkmark^\checkmark\\checkmarkc\checkmarki\checkmarkr\checkmarkc\checkmark$\checkmarkz\checkmark$\checkmark^\checkmark\\checkmarkc\checkmarki\checkmarkr\checkmarkc\checkmark$\checkmarke\checkmark$\checkmark^\checkmark\\checkmarkc\checkmarki\checkmarkr\checkmarkc\checkmark$\checkmark}\checkmark$\checkmark^\checkmark\\checkmarkc\checkmarki\checkmarkr\checkmarkc\checkmark$\checkmark
\checkmark$\checkmark^\checkmark\\checkmarkc\checkmarki\checkmarkr\checkmarkc\checkmark$\checkmark
\checkmark$\checkmark^\checkmark\\checkmarkc\checkmarki\checkmarkr\checkmarkc\checkmark$\checkmark\\checkmark$\checkmark^\checkmark\\checkmarkc\checkmarki\checkmarkr\checkmarkc\checkmark$\checkmarks\checkmark$\checkmark^\checkmark\\checkmarkc\checkmarki\checkmarkr\checkmarkc\checkmark$\checkmarku\checkmark$\checkmark^\checkmark\\checkmarkc\checkmarki\checkmarkr\checkmarkc\checkmark$\checkmarkb\checkmark$\checkmark^\checkmark\\checkmarkc\checkmarki\checkmarkr\checkmarkc\checkmark$\checkmarks\checkmark$\checkmark^\checkmark\\checkmarkc\checkmarki\checkmarkr\checkmarkc\checkmark$\checkmarke\checkmark$\checkmark^\checkmark\\checkmarkc\checkmarki\checkmarkr\checkmarkc\checkmark$\checkmarkc\checkmark$\checkmark^\checkmark\\checkmarkc\checkmarki\checkmarkr\checkmarkc\checkmark$\checkmarkt\checkmark$\checkmark^\checkmark\\checkmarkc\checkmarki\checkmarkr\checkmarkc\checkmark$\checkmarki\checkmark$\checkmark^\checkmark\\checkmarkc\checkmarki\checkmarkr\checkmarkc\checkmark$\checkmarko\checkmark$\checkmark^\checkmark\\checkmarkc\checkmarki\checkmarkr\checkmarkc\checkmark$\checkmarkn\checkmark$\checkmark^\checkmark\\checkmarkc\checkmarki\checkmarkr\checkmarkc\checkmark$\checkmark{\checkmark$\checkmark^\checkmark\\checkmarkc\checkmarki\checkmarkr\checkmarkc\checkmark$\checkmarkC\checkmark$\checkmark^\checkmark\\checkmarkc\checkmarki\checkmarkr\checkmarkc\checkmark$\checkmarkh\checkmark$\checkmark^\checkmark\\checkmarkc\checkmarki\checkmarkr\checkmarkc\checkmark$\checkmarko\checkmark$\checkmark^\checkmark\\checkmarkc\checkmarki\checkmarkr\checkmarkc\checkmark$\checkmarki\checkmark$\checkmark^\checkmark\\checkmarkc\checkmarki\checkmarkr\checkmarkc\checkmark$\checkmarkx\checkmark$\checkmark^\checkmark\\checkmarkc\checkmarki\checkmarkr\checkmarkc\checkmark$\checkmark \checkmark$\checkmark^\checkmark\\checkmarkc\checkmarki\checkmarkr\checkmarkc\checkmark$\checkmarkd\checkmark$\checkmark^\checkmark\\checkmarkc\checkmarki\checkmarkr\checkmarkc\checkmark$\checkmarke\checkmark$\checkmark^\checkmark\\checkmarkc\checkmarki\checkmarkr\checkmarkc\checkmark$\checkmarks\checkmark$\checkmark^\checkmark\\checkmarkc\checkmarki\checkmarkr\checkmarkc\checkmark$\checkmark \checkmark$\checkmark^\checkmark\\checkmarkc\checkmarki\checkmarkr\checkmarkc\checkmark$\checkmarkm\checkmark$\checkmark^\checkmark\\checkmarkc\checkmarki\checkmarkr\checkmarkc\checkmark$\checkmarko\checkmark$\checkmark^\checkmark\\checkmarkc\checkmarki\checkmarkr\checkmarkc\checkmark$\checkmarkt\checkmark$\checkmark^\checkmark\\checkmarkc\checkmarki\checkmarkr\checkmarkc\checkmark$\checkmarke\checkmark$\checkmark^\checkmark\\checkmarkc\checkmarki\checkmarkr\checkmarkc\checkmark$\checkmarku\checkmark$\checkmark^\checkmark\\checkmarkc\checkmarki\checkmarkr\checkmarkc\checkmark$\checkmarkr\checkmark$\checkmark^\checkmark\\checkmarkc\checkmarki\checkmarkr\checkmarkc\checkmark$\checkmarks\checkmark$\checkmark^\checkmark\\checkmarkc\checkmarki\checkmarkr\checkmarkc\checkmark$\checkmark}\checkmark$\checkmark^\checkmark\\checkmarkc\checkmarki\checkmarkr\checkmarkc\checkmark$\checkmark
\checkmark$\checkmark^\checkmark\\checkmarkc\checkmarki\checkmarkr\checkmarkc\checkmark$\checkmark\\checkmark$\checkmark^\checkmark\\checkmarkc\checkmarki\checkmarkr\checkmarkc\checkmark$\checkmarks\checkmark$\checkmark^\checkmark\\checkmarkc\checkmarki\checkmarkr\checkmarkc\checkmark$\checkmarku\checkmark$\checkmark^\checkmark\\checkmarkc\checkmarki\checkmarkr\checkmarkc\checkmark$\checkmarkb\checkmark$\checkmark^\checkmark\\checkmarkc\checkmarki\checkmarkr\checkmarkc\checkmark$\checkmarks\checkmark$\checkmark^\checkmark\\checkmarkc\checkmarki\checkmarkr\checkmarkc\checkmark$\checkmarku\checkmark$\checkmark^\checkmark\\checkmarkc\checkmarki\checkmarkr\checkmarkc\checkmark$\checkmarkb\checkmark$\checkmark^\checkmark\\checkmarkc\checkmarki\checkmarkr\checkmarkc\checkmark$\checkmarks\checkmark$\checkmark^\checkmark\\checkmarkc\checkmarki\checkmarkr\checkmarkc\checkmark$\checkmarke\checkmark$\checkmark^\checkmark\\checkmarkc\checkmarki\checkmarkr\checkmarkc\checkmark$\checkmarkc\checkmark$\checkmark^\checkmark\\checkmarkc\checkmarki\checkmarkr\checkmarkc\checkmark$\checkmarkt\checkmark$\checkmark^\checkmark\\checkmarkc\checkmarki\checkmarkr\checkmarkc\checkmark$\checkmarki\checkmark$\checkmark^\checkmark\\checkmarkc\checkmarki\checkmarkr\checkmarkc\checkmark$\checkmarko\checkmark$\checkmark^\checkmark\\checkmarkc\checkmarki\checkmarkr\checkmarkc\checkmark$\checkmarkn\checkmark$\checkmark^\checkmark\\checkmarkc\checkmarki\checkmarkr\checkmarkc\checkmark$\checkmark{\checkmark$\checkmark^\checkmark\\checkmarkc\checkmarki\checkmarkr\checkmarkc\checkmark$\checkmarkC\checkmark$\checkmark^\checkmark\\checkmarkc\checkmarki\checkmarkr\checkmarkc\checkmark$\checkmarko\checkmark$\checkmark^\checkmark\\checkmarkc\checkmarki\checkmarkr\checkmarkc\checkmark$\checkmarkn\checkmark$\checkmark^\checkmark\\checkmarkc\checkmarki\checkmarkr\checkmarkc\checkmark$\checkmarkd\checkmark$\checkmark^\checkmark\\checkmarkc\checkmarki\checkmarkr\checkmarkc\checkmark$\checkmarki\checkmark$\checkmark^\checkmark\\checkmarkc\checkmarki\checkmarkr\checkmarkc\checkmark$\checkmarkt\checkmark$\checkmark^\checkmark\\checkmarkc\checkmarki\checkmarkr\checkmarkc\checkmark$\checkmarki\checkmark$\checkmark^\checkmark\\checkmarkc\checkmarki\checkmarkr\checkmarkc\checkmark$\checkmarko\checkmark$\checkmark^\checkmark\\checkmarkc\checkmarki\checkmarkr\checkmarkc\checkmark$\checkmarkn\checkmark$\checkmark^\checkmark\\checkmarkc\checkmarki\checkmarkr\checkmarkc\checkmark$\checkmarks\checkmark$\checkmark^\checkmark\\checkmarkc\checkmarki\checkmarkr\checkmarkc\checkmark$\checkmark \checkmark$\checkmark^\checkmark\\checkmarkc\checkmarki\checkmarkr\checkmarkc\checkmark$\checkmarks\checkmark$\checkmark^\checkmark\\checkmarkc\checkmarki\checkmarkr\checkmarkc\checkmark$\checkmarka\checkmark$\checkmark^\checkmark\\checkmarkc\checkmarki\checkmarkr\checkmarkc\checkmark$\checkmarkt\checkmark$\checkmark^\checkmark\\checkmarkc\checkmarki\checkmarkr\checkmarkc\checkmark$\checkmarki\checkmark$\checkmark^\checkmark\\checkmarkc\checkmarki\checkmarkr\checkmarkc\checkmark$\checkmarks\checkmark$\checkmark^\checkmark\\checkmarkc\checkmarki\checkmarkr\checkmarkc\checkmark$\checkmarkf\checkmark$\checkmark^\checkmark\\checkmarkc\checkmarki\checkmarkr\checkmarkc\checkmark$\checkmarka\checkmark$\checkmark^\checkmark\\checkmarkc\checkmarki\checkmarkr\checkmarkc\checkmark$\checkmarki\checkmark$\checkmark^\checkmark\\checkmarkc\checkmarki\checkmarkr\checkmarkc\checkmark$\checkmarks\checkmark$\checkmark^\checkmark\\checkmarkc\checkmarki\checkmarkr\checkmarkc\checkmark$\checkmarka\checkmark$\checkmark^\checkmark\\checkmarkc\checkmarki\checkmarkr\checkmarkc\checkmark$\checkmarkn\checkmark$\checkmark^\checkmark\\checkmarkc\checkmarki\checkmarkr\checkmarkc\checkmark$\checkmarkt\checkmark$\checkmark^\checkmark\\checkmarkc\checkmarki\checkmarkr\checkmarkc\checkmark$\checkmarke\checkmark$\checkmark^\checkmark\\checkmarkc\checkmarki\checkmarkr\checkmarkc\checkmark$\checkmarks\checkmark$\checkmark^\checkmark\\checkmarkc\checkmarki\checkmarkr\checkmarkc\checkmark$\checkmark}\checkmark$\checkmark^\checkmark\\checkmarkc\checkmarki\checkmarkr\checkmarkc\checkmark$\checkmark
\checkmark$\checkmark^\checkmark\\checkmarkc\checkmarki\checkmarkr\checkmarkc\checkmark$\checkmarkL\checkmark$\checkmark^\checkmark\\checkmarkc\checkmarki\checkmarkr\checkmarkc\checkmark$\checkmarke\checkmark$\checkmark^\checkmark\\checkmarkc\checkmarki\checkmarkr\checkmarkc\checkmark$\checkmarks\checkmark$\checkmark^\checkmark\\checkmarkc\checkmarki\checkmarkr\checkmarkc\checkmark$\checkmark \checkmark$\checkmark^\checkmark\\checkmarkc\checkmarki\checkmarkr\checkmarkc\checkmark$\checkmarkm\checkmark$\checkmark^\checkmark\\checkmarkc\checkmarki\checkmarkr\checkmarkc\checkmark$\checkmarko\checkmark$\checkmark^\checkmark\\checkmarkc\checkmarki\checkmarkr\checkmarkc\checkmark$\checkmarkt\checkmark$\checkmark^\checkmark\\checkmarkc\checkmarki\checkmarkr\checkmarkc\checkmark$\checkmarke\checkmark$\checkmark^\checkmark\\checkmarkc\checkmarki\checkmarkr\checkmarkc\checkmark$\checkmarku\checkmark$\checkmark^\checkmark\\checkmarkc\checkmarki\checkmarkr\checkmarkc\checkmark$\checkmarkr\checkmark$\checkmark^\checkmark\\checkmarkc\checkmarki\checkmarkr\checkmarkc\checkmark$\checkmarks\checkmark$\checkmark^\checkmark\\checkmarkc\checkmarki\checkmarkr\checkmarkc\checkmark$\checkmark \checkmark$\checkmark^\checkmark\\checkmarkc\checkmarki\checkmarkr\checkmarkc\checkmark$\checkmarkd\checkmark$\checkmark^\checkmark\\checkmarkc\checkmarki\checkmarkr\checkmarkc\checkmark$\checkmark'\checkmark$\checkmark^\checkmark\\checkmarkc\checkmarki\checkmarkr\checkmarkc\checkmark$\checkmarke\checkmark$\checkmark^\checkmark\\checkmarkc\checkmarki\checkmarkr\checkmarkc\checkmark$\checkmarkn\checkmark$\checkmark^\checkmark\\checkmarkc\checkmarki\checkmarkr\checkmarkc\checkmark$\checkmarkt\checkmark$\checkmark^\checkmark\\checkmarkc\checkmarki\checkmarkr\checkmarkc\checkmark$\checkmarkr\checkmark$\checkmark^\checkmark\\checkmarkc\checkmarki\checkmarkr\checkmarkc\checkmark$\checkmarka\checkmark$\checkmark^\checkmark\\checkmarkc\checkmarki\checkmarkr\checkmarkc\checkmark$\checkmarkî\checkmark$\checkmark^\checkmark\\checkmarkc\checkmarki\checkmarkr\checkmarkc\checkmark$\checkmarkn\checkmark$\checkmark^\checkmark\\checkmarkc\checkmarki\checkmarkr\checkmarkc\checkmark$\checkmarke\checkmark$\checkmark^\checkmark\\checkmarkc\checkmarki\checkmarkr\checkmarkc\checkmark$\checkmarkm\checkmark$\checkmark^\checkmark\\checkmarkc\checkmarki\checkmarkr\checkmarkc\checkmark$\checkmarke\checkmark$\checkmark^\checkmark\\checkmarkc\checkmarki\checkmarkr\checkmarkc\checkmark$\checkmarkn\checkmark$\checkmark^\checkmark\\checkmarkc\checkmarki\checkmarkr\checkmarkc\checkmark$\checkmarkt\checkmark$\checkmark^\checkmark\\checkmarkc\checkmarki\checkmarkr\checkmarkc\checkmark$\checkmark \checkmark$\checkmark^\checkmark\\checkmarkc\checkmarki\checkmarkr\checkmarkc\checkmark$\checkmarkd\checkmark$\checkmark^\checkmark\\checkmarkc\checkmarki\checkmarkr\checkmarkc\checkmark$\checkmarko\checkmark$\checkmark^\checkmark\\checkmarkc\checkmarki\checkmarkr\checkmarkc\checkmark$\checkmarki\checkmark$\checkmark^\checkmark\\checkmarkc\checkmarki\checkmarkr\checkmarkc\checkmark$\checkmarkv\checkmark$\checkmark^\checkmark\\checkmarkc\checkmarki\checkmarkr\checkmarkc\checkmark$\checkmarke\checkmark$\checkmark^\checkmark\\checkmarkc\checkmarki\checkmarkr\checkmarkc\checkmark$\checkmarkn\checkmark$\checkmark^\checkmark\\checkmarkc\checkmarki\checkmarkr\checkmarkc\checkmark$\checkmarkt\checkmark$\checkmark^\checkmark\\checkmarkc\checkmarki\checkmarkr\checkmarkc\checkmark$\checkmark \checkmark$\checkmark^\checkmark\\checkmarkc\checkmarki\checkmarkr\checkmarkc\checkmark$\checkmark:\checkmark$\checkmark^\checkmark\\checkmarkc\checkmarki\checkmarkr\checkmarkc\checkmark$\checkmark \checkmark$\checkmark^\checkmark\\checkmarkc\checkmarki\checkmarkr\checkmarkc\checkmark$\checkmarkf\checkmark$\checkmark^\checkmark\\checkmarkc\checkmarki\checkmarkr\checkmarkc\checkmark$\checkmarko\checkmark$\checkmark^\checkmark\\checkmarkc\checkmarki\checkmarkr\checkmarkc\checkmark$\checkmarku\checkmark$\checkmark^\checkmark\\checkmarkc\checkmarki\checkmarkr\checkmarkc\checkmark$\checkmarkr\checkmark$\checkmark^\checkmark\\checkmarkc\checkmarki\checkmarkr\checkmarkc\checkmark$\checkmarkn\checkmark$\checkmark^\checkmark\\checkmarkc\checkmarki\checkmarkr\checkmarkc\checkmark$\checkmarki\checkmark$\checkmark^\checkmark\\checkmarkc\checkmarki\checkmarkr\checkmarkc\checkmark$\checkmarkr\checkmark$\checkmark^\checkmark\\checkmarkc\checkmarki\checkmarkr\checkmarkc\checkmark$\checkmark \checkmark$\checkmark^\checkmark\\checkmarkc\checkmarki\checkmarkr\checkmarkc\checkmark$\checkmarku\checkmark$\checkmark^\checkmark\\checkmarkc\checkmarki\checkmarkr\checkmarkc\checkmark$\checkmarkn\checkmark$\checkmark^\checkmark\\checkmarkc\checkmarki\checkmarkr\checkmarkc\checkmark$\checkmark \checkmark$\checkmark^\checkmark\\checkmarkc\checkmarki\checkmarkr\checkmarkc\checkmark$\checkmarkc\checkmark$\checkmark^\checkmark\\checkmarkc\checkmarki\checkmarkr\checkmarkc\checkmark$\checkmarko\checkmark$\checkmark^\checkmark\\checkmarkc\checkmarki\checkmarkr\checkmarkc\checkmark$\checkmarku\checkmark$\checkmark^\checkmark\\checkmarkc\checkmarki\checkmarkr\checkmarkc\checkmark$\checkmarkp\checkmark$\checkmark^\checkmark\\checkmarkc\checkmarki\checkmarkr\checkmarkc\checkmark$\checkmarkl\checkmark$\checkmark^\checkmark\\checkmarkc\checkmarki\checkmarkr\checkmarkc\checkmark$\checkmarke\checkmark$\checkmark^\checkmark\\checkmarkc\checkmarki\checkmarkr\checkmarkc\checkmark$\checkmark \checkmark$\checkmark^\checkmark\\checkmarkc\checkmarki\checkmarkr\checkmarkc\checkmark$\checkmarks\checkmark$\checkmark^\checkmark\\checkmarkc\checkmarki\checkmarkr\checkmarkc\checkmark$\checkmarku\checkmark$\checkmark^\checkmark\\checkmarkc\checkmarki\checkmarkr\checkmarkc\checkmark$\checkmarkf\checkmark$\checkmark^\checkmark\\checkmarkc\checkmarki\checkmarkr\checkmarkc\checkmark$\checkmarkf\checkmark$\checkmark^\checkmark\\checkmarkc\checkmarki\checkmarkr\checkmarkc\checkmark$\checkmarki\checkmark$\checkmark^\checkmark\\checkmarkc\checkmarki\checkmarkr\checkmarkc\checkmark$\checkmarks\checkmark$\checkmark^\checkmark\\checkmarkc\checkmarki\checkmarkr\checkmarkc\checkmark$\checkmarka\checkmark$\checkmark^\checkmark\\checkmarkc\checkmarki\checkmarkr\checkmarkc\checkmark$\checkmarkn\checkmark$\checkmark^\checkmark\\checkmarkc\checkmarki\checkmarkr\checkmarkc\checkmark$\checkmarkt\checkmark$\checkmark^\checkmark\\checkmarkc\checkmarki\checkmarkr\checkmarkc\checkmark$\checkmark \checkmark$\checkmark^\checkmark\\checkmarkc\checkmarki\checkmarkr\checkmarkc\checkmark$\checkmarks\checkmark$\checkmark^\checkmark\\checkmarkc\checkmarki\checkmarkr\checkmarkc\checkmark$\checkmarka\checkmark$\checkmark^\checkmark\\checkmarkc\checkmarki\checkmarkr\checkmarkc\checkmark$\checkmarkn\checkmark$\checkmark^\checkmark\\checkmarkc\checkmarki\checkmarkr\checkmarkc\checkmark$\checkmarks\checkmark$\checkmark^\checkmark\\checkmarkc\checkmarki\checkmarkr\checkmarkc\checkmark$\checkmark \checkmark$\checkmark^\checkmark\\checkmarkc\checkmarki\checkmarkr\checkmarkc\checkmark$\checkmarkp\checkmark$\checkmark^\checkmark\\checkmarkc\checkmarki\checkmarkr\checkmarkc\checkmark$\checkmarke\checkmark$\checkmark^\checkmark\\checkmarkc\checkmarki\checkmarkr\checkmarkc\checkmark$\checkmarkr\checkmark$\checkmark^\checkmark\\checkmarkc\checkmarki\checkmarkr\checkmarkc\checkmark$\checkmarkd\checkmark$\checkmark^\checkmark\\checkmarkc\checkmarki\checkmarkr\checkmarkc\checkmark$\checkmarkr\checkmark$\checkmark^\checkmark\\checkmarkc\checkmarki\checkmarkr\checkmarkc\checkmark$\checkmarke\checkmark$\checkmark^\checkmark\\checkmarkc\checkmarki\checkmarkr\checkmarkc\checkmark$\checkmark \checkmark$\checkmark^\checkmark\\checkmarkc\checkmarki\checkmarkr\checkmarkc\checkmark$\checkmarkd\checkmark$\checkmark^\checkmark\\checkmarkc\checkmarki\checkmarkr\checkmarkc\checkmark$\checkmarke\checkmark$\checkmark^\checkmark\\checkmarkc\checkmarki\checkmarkr\checkmarkc\checkmark$\checkmark \checkmark$\checkmark^\checkmark\\checkmarkc\checkmarki\checkmarkr\checkmarkc\checkmark$\checkmarkp\checkmark$\checkmark^\checkmark\\checkmarkc\checkmarki\checkmarkr\checkmarkc\checkmark$\checkmarka\checkmark$\checkmark^\checkmark\\checkmarkc\checkmarki\checkmarkr\checkmarkc\checkmark$\checkmarks\checkmark$\checkmark^\checkmark\\checkmarkc\checkmarki\checkmarkr\checkmarkc\checkmark$\checkmark,\checkmark$\checkmark^\checkmark\\checkmarkc\checkmarki\checkmarkr\checkmarkc\checkmark$\checkmark \checkmark$\checkmark^\checkmark\\checkmarkc\checkmarki\checkmarkr\checkmarkc\checkmark$\checkmarkê\checkmark$\checkmark^\checkmark\\checkmarkc\checkmarki\checkmarkr\checkmarkc\checkmark$\checkmarkt\checkmark$\checkmark^\checkmark\\checkmarkc\checkmarki\checkmarkr\checkmarkc\checkmark$\checkmarkr\checkmark$\checkmark^\checkmark\\checkmarkc\checkmarki\checkmarkr\checkmarkc\checkmark$\checkmarke\checkmark$\checkmark^\checkmark\\checkmarkc\checkmarki\checkmarkr\checkmarkc\checkmark$\checkmark \checkmark$\checkmark^\checkmark\\checkmarkc\checkmarki\checkmarkr\checkmarkc\checkmark$\checkmarkc\checkmark$\checkmark^\checkmark\\checkmarkc\checkmarki\checkmarkr\checkmarkc\checkmark$\checkmarko\checkmark$\checkmark^\checkmark\\checkmarkc\checkmarki\checkmarkr\checkmarkc\checkmark$\checkmarkm\checkmark$\checkmark^\checkmark\\checkmarkc\checkmarki\checkmarkr\checkmarkc\checkmark$\checkmarkp\checkmark$\checkmark^\checkmark\\checkmarkc\checkmarki\checkmarkr\checkmarkc\checkmark$\checkmarka\checkmark$\checkmark^\checkmark\\checkmarkc\checkmarki\checkmarkr\checkmarkc\checkmark$\checkmarkt\checkmark$\checkmark^\checkmark\\checkmarkc\checkmarki\checkmarkr\checkmarkc\checkmark$\checkmarki\checkmark$\checkmark^\checkmark\\checkmarkc\checkmarki\checkmarkr\checkmarkc\checkmark$\checkmarkb\checkmark$\checkmark^\checkmark\\checkmarkc\checkmarki\checkmarkr\checkmarkc\checkmark$\checkmarkl\checkmark$\checkmark^\checkmark\\checkmarkc\checkmarki\checkmarkr\checkmarkc\checkmark$\checkmarke\checkmark$\checkmark^\checkmark\\checkmarkc\checkmarki\checkmarkr\checkmarkc\checkmark$\checkmarks\checkmark$\checkmark^\checkmark\\checkmarkc\checkmarki\checkmarkr\checkmarkc\checkmark$\checkmark \checkmark$\checkmark^\checkmark\\checkmarkc\checkmarki\checkmarkr\checkmarkc\checkmark$\checkmarka\checkmark$\checkmark^\checkmark\\checkmarkc\checkmarki\checkmarkr\checkmarkc\checkmark$\checkmarkv\checkmark$\checkmark^\checkmark\\checkmarkc\checkmarki\checkmarkr\checkmarkc\checkmark$\checkmarke\checkmark$\checkmark^\checkmark\\checkmarkc\checkmarki\checkmarkr\checkmarkc\checkmark$\checkmarkc\checkmark$\checkmark^\checkmark\\checkmarkc\checkmarki\checkmarkr\checkmarkc\checkmark$\checkmark \checkmark$\checkmark^\checkmark\\checkmarkc\checkmarki\checkmarkr\checkmarkc\checkmark$\checkmarkl\checkmark$\checkmark^\checkmark\\checkmarkc\checkmarki\checkmarkr\checkmarkc\checkmark$\checkmarke\checkmark$\checkmark^\checkmark\\checkmarkc\checkmarki\checkmarkr\checkmarkc\checkmark$\checkmark \checkmark$\checkmark^\checkmark\\checkmarkc\checkmarki\checkmarkr\checkmarkc\checkmark$\checkmarkc\checkmark$\checkmark^\checkmark\\checkmarkc\checkmarki\checkmarkr\checkmarkc\checkmark$\checkmarko\checkmark$\checkmark^\checkmark\\checkmarkc\checkmarki\checkmarkr\checkmarkc\checkmark$\checkmarkn\checkmark$\checkmark^\checkmark\\checkmarkc\checkmarki\checkmarkr\checkmarkc\checkmark$\checkmarkt\checkmark$\checkmark^\checkmark\\checkmarkc\checkmarki\checkmarkr\checkmarkc\checkmark$\checkmarkr\checkmark$\checkmark^\checkmark\\checkmarkc\checkmarki\checkmarkr\checkmarkc\checkmark$\checkmarkô\checkmark$\checkmark^\checkmark\\checkmarkc\checkmarki\checkmarkr\checkmarkc\checkmark$\checkmarkl\checkmark$\checkmark^\checkmark\\checkmarkc\checkmarki\checkmarkr\checkmarkc\checkmark$\checkmarke\checkmark$\checkmark^\checkmark\\checkmarkc\checkmarki\checkmarkr\checkmarkc\checkmark$\checkmarku\checkmark$\checkmark^\checkmark\\checkmarkc\checkmarki\checkmarkr\checkmarkc\checkmark$\checkmarkr\checkmark$\checkmark^\checkmark\\checkmarkc\checkmarki\checkmarkr\checkmarkc\checkmark$\checkmark \checkmark$\checkmark^\checkmark\\checkmarkc\checkmarki\checkmarkr\checkmarkc\checkmark$\checkmarkE\checkmark$\checkmark^\checkmark\\checkmarkc\checkmarki\checkmarkr\checkmarkc\checkmark$\checkmarkS\checkmark$\checkmark^\checkmark\\checkmarkc\checkmarki\checkmarkr\checkmarkc\checkmark$\checkmarkP\checkmark$\checkmark^\checkmark\\checkmarkc\checkmarki\checkmarkr\checkmarkc\checkmark$\checkmark3\checkmark$\checkmark^\checkmark\\checkmarkc\checkmarki\checkmarkr\checkmarkc\checkmark$\checkmark2\checkmark$\checkmark^\checkmark\\checkmarkc\checkmarki\checkmarkr\checkmarkc\checkmark$\checkmark \checkmark$\checkmark^\checkmark\\checkmarkc\checkmarki\checkmarkr\checkmarkc\checkmark$\checkmark(\checkmark$\checkmark^\checkmark\\checkmarkc\checkmarki\checkmarkr\checkmarkc\checkmark$\checkmarkd\checkmark$\checkmark^\checkmark\\checkmarkc\checkmarki\checkmarkr\checkmarkc\checkmark$\checkmarkr\checkmark$\checkmark^\checkmark\\checkmarkc\checkmarki\checkmarkr\checkmarkc\checkmark$\checkmarki\checkmark$\checkmark^\checkmark\\checkmarkc\checkmarki\checkmarkr\checkmarkc\checkmark$\checkmarkv\checkmark$\checkmark^\checkmark\\checkmarkc\checkmarki\checkmarkr\checkmarkc\checkmark$\checkmarke\checkmark$\checkmark^\checkmark\\checkmarkc\checkmarki\checkmarkr\checkmarkc\checkmark$\checkmarkr\checkmark$\checkmark^\checkmark\\checkmarkc\checkmarki\checkmarkr\checkmarkc\checkmark$\checkmark \checkmark$\checkmark^\checkmark\\checkmarkc\checkmarki\checkmarkr\checkmarkc\checkmark$\checkmarkA\checkmark$\checkmark^\checkmark\\checkmarkc\checkmarki\checkmarkr\checkmarkc\checkmark$\checkmark4\checkmark$\checkmark^\checkmark\\checkmarkc\checkmarki\checkmarkr\checkmarkc\checkmark$\checkmark9\checkmark$\checkmark^\checkmark\\checkmarkc\checkmarki\checkmarkr\checkmarkc\checkmark$\checkmark8\checkmark$\checkmark^\checkmark\\checkmarkc\checkmarki\checkmarkr\checkmarkc\checkmark$\checkmark8\checkmark$\checkmark^\checkmark\\checkmarkc\checkmarki\checkmarkr\checkmarkc\checkmark$\checkmark \checkmark$\checkmark^\checkmark\\checkmarkc\checkmarki\checkmarkr\checkmarkc\checkmark$\checkmark/\checkmark$\checkmark^\checkmark\\checkmarkc\checkmarki\checkmarkr\checkmarkc\checkmark$\checkmark \checkmark$\checkmark^\checkmark\\checkmarkc\checkmarki\checkmarkr\checkmarkc\checkmark$\checkmarkT\checkmark$\checkmark^\checkmark\\checkmarkc\checkmarki\checkmarkr\checkmarkc\checkmark$\checkmarkM\checkmark$\checkmark^\checkmark\\checkmarkc\checkmarki\checkmarkr\checkmarkc\checkmark$\checkmarkC\checkmark$\checkmark^\checkmark\\checkmarkc\checkmarki\checkmarkr\checkmarkc\checkmark$\checkmark2\checkmark$\checkmark^\checkmark\\checkmarkc\checkmarki\checkmarkr\checkmarkc\checkmark$\checkmark2\checkmark$\checkmark^\checkmark\\checkmarkc\checkmarki\checkmarkr\checkmarkc\checkmark$\checkmark0\checkmark$\checkmark^\checkmark\\checkmarkc\checkmarki\checkmarkr\checkmarkc\checkmark$\checkmark9\checkmark$\checkmark^\checkmark\\checkmarkc\checkmarki\checkmarkr\checkmarkc\checkmark$\checkmark)\checkmark$\checkmark^\checkmark\\checkmarkc\checkmarki\checkmarkr\checkmarkc\checkmark$\checkmark,\checkmark$\checkmark^\checkmark\\checkmarkc\checkmarki\checkmarkr\checkmarkc\checkmark$\checkmark \checkmark$\checkmark^\checkmark\\checkmarkc\checkmarki\checkmarkr\checkmarkc\checkmark$\checkmarke\checkmark$\checkmark^\checkmark\\checkmarkc\checkmarki\checkmarkr\checkmarkc\checkmark$\checkmarkt\checkmark$\checkmark^\checkmark\\checkmarkc\checkmarki\checkmarkr\checkmarkc\checkmark$\checkmark \checkmark$\checkmark^\checkmark\\checkmarkc\checkmarki\checkmarkr\checkmarkc\checkmark$\checkmarka\checkmark$\checkmark^\checkmark\\checkmarkc\checkmarki\checkmarkr\checkmarkc\checkmark$\checkmarkv\checkmark$\checkmark^\checkmark\\checkmarkc\checkmarki\checkmarkr\checkmarkc\checkmark$\checkmarko\checkmark$\checkmark^\checkmark\\checkmarkc\checkmarki\checkmarkr\checkmarkc\checkmark$\checkmarki\checkmark$\checkmark^\checkmark\\checkmarkc\checkmarki\checkmarkr\checkmarkc\checkmark$\checkmarkr\checkmark$\checkmark^\checkmark\\checkmarkc\checkmarki\checkmarkr\checkmarkc\checkmark$\checkmark \checkmark$\checkmark^\checkmark\\checkmarkc\checkmarki\checkmarkr\checkmarkc\checkmark$\checkmarku\checkmark$\checkmark^\checkmark\\checkmarkc\checkmarki\checkmarkr\checkmarkc\checkmark$\checkmarkn\checkmark$\checkmark^\checkmark\\checkmarkc\checkmarki\checkmarkr\checkmarkc\checkmark$\checkmark \checkmark$\checkmark^\checkmark\\checkmarkc\checkmarki\checkmarkr\checkmarkc\checkmark$\checkmarké\checkmark$\checkmark^\checkmark\\checkmarkc\checkmarki\checkmarkr\checkmarkc\checkmark$\checkmarkc\checkmark$\checkmark^\checkmark\\checkmarkc\checkmarki\checkmarkr\checkmarkc\checkmark$\checkmarkh\checkmark$\checkmark^\checkmark\\checkmarkc\checkmarki\checkmarkr\checkmarkc\checkmark$\checkmarka\checkmark$\checkmark^\checkmark\\checkmarkc\checkmarki\checkmarkr\checkmarkc\checkmark$\checkmarku\checkmark$\checkmark^\checkmark\\checkmarkc\checkmarki\checkmarkr\checkmarkc\checkmark$\checkmarkf\checkmark$\checkmark^\checkmark\\checkmarkc\checkmarki\checkmarkr\checkmarkc\checkmark$\checkmarkf\checkmark$\checkmark^\checkmark\\checkmarkc\checkmarki\checkmarkr\checkmarkc\checkmark$\checkmarke\checkmark$\checkmark^\checkmark\\checkmarkc\checkmarki\checkmarkr\checkmarkc\checkmark$\checkmarkm\checkmark$\checkmark^\checkmark\\checkmarkc\checkmarki\checkmarkr\checkmarkc\checkmark$\checkmarke\checkmark$\checkmark^\checkmark\\checkmarkc\checkmarki\checkmarkr\checkmarkc\checkmark$\checkmarkn\checkmark$\checkmark^\checkmark\\checkmarkc\checkmarki\checkmarkr\checkmarkc\checkmark$\checkmarkt\checkmark$\checkmark^\checkmark\\checkmarkc\checkmarki\checkmarkr\checkmarkc\checkmark$\checkmark \checkmark$\checkmark^\checkmark\\checkmarkc\checkmarki\checkmarkr\checkmarkc\checkmark$\checkmarkl\checkmark$\checkmark^\checkmark\\checkmarkc\checkmarki\checkmarkr\checkmarkc\checkmark$\checkmarki\checkmark$\checkmark^\checkmark\\checkmarkc\checkmarki\checkmarkr\checkmarkc\checkmark$\checkmarkm\checkmark$\checkmark^\checkmark\\checkmarkc\checkmarki\checkmarkr\checkmarkc\checkmark$\checkmarki\checkmark$\checkmark^\checkmark\\checkmarkc\checkmarki\checkmarkr\checkmarkc\checkmark$\checkmarkt\checkmark$\checkmark^\checkmark\\checkmarkc\checkmarki\checkmarkr\checkmarkc\checkmark$\checkmarké\checkmark$\checkmark^\checkmark\\checkmarkc\checkmarki\checkmarkr\checkmarkc\checkmark$\checkmark \checkmark$\checkmark^\checkmark\\checkmarkc\checkmarki\checkmarkr\checkmarkc\checkmark$\checkmarke\checkmark$\checkmark^\checkmark\\checkmarkc\checkmarki\checkmarkr\checkmarkc\checkmark$\checkmarkn\checkmark$\checkmark^\checkmark\\checkmarkc\checkmarki\checkmarkr\checkmarkc\checkmark$\checkmark \checkmark$\checkmark^\checkmark\\checkmarkc\checkmarki\checkmarkr\checkmarkc\checkmark$\checkmarks\checkmark$\checkmark^\checkmark\\checkmarkc\checkmarki\checkmarkr\checkmarkc\checkmark$\checkmarke\checkmark$\checkmark^\checkmark\\checkmarkc\checkmarki\checkmarkr\checkmarkc\checkmark$\checkmarkr\checkmark$\checkmark^\checkmark\\checkmarkc\checkmarki\checkmarkr\checkmarkc\checkmark$\checkmarkv\checkmark$\checkmark^\checkmark\\checkmarkc\checkmarki\checkmarkr\checkmarkc\checkmark$\checkmarki\checkmark$\checkmark^\checkmark\\checkmarkc\checkmarki\checkmarkr\checkmarkc\checkmark$\checkmarkc\checkmark$\checkmark^\checkmark\\checkmarkc\checkmarki\checkmarkr\checkmarkc\checkmark$\checkmarke\checkmark$\checkmark^\checkmark\\checkmarkc\checkmarki\checkmarkr\checkmarkc\checkmark$\checkmark \checkmark$\checkmark^\checkmark\\checkmarkc\checkmarki\checkmarkr\checkmarkc\checkmark$\checkmarkc\checkmark$\checkmark^\checkmark\\checkmarkc\checkmarki\checkmarkr\checkmarkc\checkmark$\checkmarko\checkmark$\checkmark^\checkmark\\checkmarkc\checkmarki\checkmarkr\checkmarkc\checkmark$\checkmarkn\checkmark$\checkmark^\checkmark\\checkmarkc\checkmarki\checkmarkr\checkmarkc\checkmark$\checkmarkt\checkmark$\checkmark^\checkmark\\checkmarkc\checkmarki\checkmarkr\checkmarkc\checkmark$\checkmarki\checkmark$\checkmark^\checkmark\\checkmarkc\checkmarki\checkmarkr\checkmarkc\checkmark$\checkmarkn\checkmark$\checkmark^\checkmark\\checkmarkc\checkmarki\checkmarkr\checkmarkc\checkmark$\checkmarku\checkmark$\checkmark^\checkmark\\checkmarkc\checkmarki\checkmarkr\checkmarkc\checkmark$\checkmark.\checkmark$\checkmark^\checkmark\\checkmarkc\checkmarki\checkmarkr\checkmarkc\checkmark$\checkmark
\checkmark$\checkmark^\checkmark\\checkmarkc\checkmarki\checkmarkr\checkmarkc\checkmark$\checkmark
\checkmark$\checkmark^\checkmark\\checkmarkc\checkmarki\checkmarkr\checkmarkc\checkmark$\checkmark\\checkmark$\checkmark^\checkmark\\checkmarkc\checkmarki\checkmarkr\checkmarkc\checkmark$\checkmarks\checkmark$\checkmark^\checkmark\\checkmarkc\checkmarki\checkmarkr\checkmarkc\checkmark$\checkmarku\checkmark$\checkmark^\checkmark\\checkmarkc\checkmarki\checkmarkr\checkmarkc\checkmark$\checkmarkb\checkmark$\checkmark^\checkmark\\checkmarkc\checkmarki\checkmarkr\checkmarkc\checkmark$\checkmarks\checkmark$\checkmark^\checkmark\\checkmarkc\checkmarki\checkmarkr\checkmarkc\checkmark$\checkmarku\checkmark$\checkmark^\checkmark\\checkmarkc\checkmarki\checkmarkr\checkmarkc\checkmark$\checkmarkb\checkmark$\checkmark^\checkmark\\checkmarkc\checkmarki\checkmarkr\checkmarkc\checkmark$\checkmarks\checkmark$\checkmark^\checkmark\\checkmarkc\checkmarki\checkmarkr\checkmarkc\checkmark$\checkmarke\checkmark$\checkmark^\checkmark\\checkmarkc\checkmarki\checkmarkr\checkmarkc\checkmark$\checkmarkc\checkmark$\checkmark^\checkmark\\checkmarkc\checkmarki\checkmarkr\checkmarkc\checkmark$\checkmarkt\checkmark$\checkmark^\checkmark\\checkmarkc\checkmarki\checkmarkr\checkmarkc\checkmark$\checkmarki\checkmark$\checkmark^\checkmark\\checkmarkc\checkmarki\checkmarkr\checkmarkc\checkmark$\checkmarko\checkmark$\checkmark^\checkmark\\checkmarkc\checkmarki\checkmarkr\checkmarkc\checkmark$\checkmarkn\checkmark$\checkmark^\checkmark\\checkmarkc\checkmarki\checkmarkr\checkmarkc\checkmark$\checkmark{\checkmark$\checkmark^\checkmark\\checkmarkc\checkmarki\checkmarkr\checkmarkc\checkmark$\checkmarkC\checkmark$\checkmark^\checkmark\\checkmarkc\checkmarki\checkmarkr\checkmarkc\checkmark$\checkmarkh\checkmark$\checkmark^\checkmark\\checkmarkc\checkmarki\checkmarkr\checkmarkc\checkmark$\checkmarko\checkmark$\checkmark^\checkmark\\checkmarkc\checkmarki\checkmarkr\checkmarkc\checkmark$\checkmarki\checkmark$\checkmark^\checkmark\\checkmarkc\checkmarki\checkmarkr\checkmarkc\checkmark$\checkmarkx\checkmark$\checkmark^\checkmark\\checkmarkc\checkmarki\checkmarkr\checkmarkc\checkmark$\checkmark \checkmark$\checkmark^\checkmark\\checkmarkc\checkmarki\checkmarkr\checkmarkc\checkmark$\checkmarkd\checkmark$\checkmark^\checkmark\\checkmarkc\checkmarki\checkmarkr\checkmarkc\checkmark$\checkmarké\checkmark$\checkmark^\checkmark\\checkmarkc\checkmarki\checkmarkr\checkmarkc\checkmark$\checkmarkf\checkmark$\checkmark^\checkmark\\checkmarkc\checkmarki\checkmarkr\checkmarkc\checkmark$\checkmarki\checkmark$\checkmark^\checkmark\\checkmarkc\checkmarki\checkmarkr\checkmarkc\checkmark$\checkmarkn\checkmark$\checkmark^\checkmark\\checkmarkc\checkmarki\checkmarkr\checkmarkc\checkmark$\checkmarki\checkmark$\checkmark^\checkmark\\checkmarkc\checkmarki\checkmarkr\checkmarkc\checkmark$\checkmarkt\checkmark$\checkmark^\checkmark\\checkmarkc\checkmarki\checkmarkr\checkmarkc\checkmark$\checkmarki\checkmark$\checkmark^\checkmark\\checkmarkc\checkmarki\checkmarkr\checkmarkc\checkmark$\checkmarkf\checkmark$\checkmark^\checkmark\\checkmarkc\checkmarki\checkmarkr\checkmarkc\checkmark$\checkmark}\checkmark$\checkmark^\checkmark\\checkmarkc\checkmarki\checkmarkr\checkmarkc\checkmark$\checkmark
\checkmark$\checkmark^\checkmark\\checkmarkc\checkmarki\checkmarkr\checkmarkc\checkmark$\checkmark\\checkmark$\checkmark^\checkmark\\checkmarkc\checkmarki\checkmarkr\checkmarkc\checkmark$\checkmarkb\checkmark$\checkmark^\checkmark\\checkmarkc\checkmarki\checkmarkr\checkmarkc\checkmark$\checkmarke\checkmark$\checkmark^\checkmark\\checkmarkc\checkmarki\checkmarkr\checkmarkc\checkmark$\checkmarkg\checkmark$\checkmark^\checkmark\\checkmarkc\checkmarki\checkmarkr\checkmarkc\checkmark$\checkmarki\checkmark$\checkmark^\checkmark\\checkmarkc\checkmarki\checkmarkr\checkmarkc\checkmark$\checkmarkn\checkmark$\checkmark^\checkmark\\checkmarkc\checkmarki\checkmarkr\checkmarkc\checkmark$\checkmark{\checkmark$\checkmark^\checkmark\\checkmarkc\checkmarki\checkmarkr\checkmarkc\checkmark$\checkmarkt\checkmark$\checkmark^\checkmark\\checkmarkc\checkmarki\checkmarkr\checkmarkc\checkmark$\checkmarka\checkmark$\checkmark^\checkmark\\checkmarkc\checkmarki\checkmarkr\checkmarkc\checkmark$\checkmarkb\checkmark$\checkmark^\checkmark\\checkmarkc\checkmarki\checkmarkr\checkmarkc\checkmark$\checkmarkl\checkmark$\checkmark^\checkmark\\checkmarkc\checkmarki\checkmarkr\checkmarkc\checkmark$\checkmarke\checkmark$\checkmark^\checkmark\\checkmarkc\checkmarki\checkmarkr\checkmarkc\checkmark$\checkmark}\checkmark$\checkmark^\checkmark\\checkmarkc\checkmarki\checkmarkr\checkmarkc\checkmark$\checkmark[\checkmark$\checkmark^\checkmark\\checkmarkc\checkmarki\checkmarkr\checkmarkc\checkmark$\checkmarkh\checkmark$\checkmark^\checkmark\\checkmarkc\checkmarki\checkmarkr\checkmarkc\checkmark$\checkmarkt\checkmark$\checkmark^\checkmark\\checkmarkc\checkmarki\checkmarkr\checkmarkc\checkmark$\checkmarkb\checkmark$\checkmark^\checkmark\\checkmarkc\checkmarki\checkmarkr\checkmarkc\checkmark$\checkmarkp\checkmark$\checkmark^\checkmark\\checkmarkc\checkmarki\checkmarkr\checkmarkc\checkmark$\checkmark]\checkmark$\checkmark^\checkmark\\checkmarkc\checkmarki\checkmarkr\checkmarkc\checkmark$\checkmark
\checkmark$\checkmark^\checkmark\\checkmarkc\checkmarki\checkmarkr\checkmarkc\checkmark$\checkmark\\checkmark$\checkmark^\checkmark\\checkmarkc\checkmarki\checkmarkr\checkmarkc\checkmark$\checkmarkc\checkmark$\checkmark^\checkmark\\checkmarkc\checkmarki\checkmarkr\checkmarkc\checkmark$\checkmarke\checkmark$\checkmark^\checkmark\\checkmarkc\checkmarki\checkmarkr\checkmarkc\checkmark$\checkmarkn\checkmark$\checkmark^\checkmark\\checkmarkc\checkmarki\checkmarkr\checkmarkc\checkmark$\checkmarkt\checkmark$\checkmark^\checkmark\\checkmarkc\checkmarki\checkmarkr\checkmarkc\checkmark$\checkmarke\checkmark$\checkmark^\checkmark\\checkmarkc\checkmarki\checkmarkr\checkmarkc\checkmark$\checkmarkr\checkmark$\checkmark^\checkmark\\checkmarkc\checkmarki\checkmarkr\checkmarkc\checkmark$\checkmarki\checkmark$\checkmark^\checkmark\\checkmarkc\checkmarki\checkmarkr\checkmarkc\checkmark$\checkmarkn\checkmark$\checkmark^\checkmark\\checkmarkc\checkmarki\checkmarkr\checkmarkc\checkmark$\checkmarkg\checkmark$\checkmark^\checkmark\\checkmarkc\checkmarki\checkmarkr\checkmarkc\checkmark$\checkmark
\checkmark$\checkmark^\checkmark\\checkmarkc\checkmarki\checkmarkr\checkmarkc\checkmark$\checkmark\\checkmark$\checkmark^\checkmark\\checkmarkc\checkmarki\checkmarkr\checkmarkc\checkmark$\checkmarkb\checkmark$\checkmark^\checkmark\\checkmarkc\checkmarki\checkmarkr\checkmarkc\checkmark$\checkmarke\checkmark$\checkmark^\checkmark\\checkmarkc\checkmarki\checkmarkr\checkmarkc\checkmark$\checkmarkg\checkmark$\checkmark^\checkmark\\checkmarkc\checkmarki\checkmarkr\checkmarkc\checkmark$\checkmarki\checkmark$\checkmark^\checkmark\\checkmarkc\checkmarki\checkmarkr\checkmarkc\checkmark$\checkmarkn\checkmark$\checkmark^\checkmark\\checkmarkc\checkmarki\checkmarkr\checkmarkc\checkmark$\checkmark{\checkmark$\checkmark^\checkmark\\checkmarkc\checkmarki\checkmarkr\checkmarkc\checkmark$\checkmarkt\checkmark$\checkmark^\checkmark\\checkmarkc\checkmarki\checkmarkr\checkmarkc\checkmark$\checkmarka\checkmark$\checkmark^\checkmark\\checkmarkc\checkmarki\checkmarkr\checkmarkc\checkmark$\checkmarkb\checkmark$\checkmark^\checkmark\\checkmarkc\checkmarki\checkmarkr\checkmarkc\checkmark$\checkmarku\checkmark$\checkmark^\checkmark\\checkmarkc\checkmarki\checkmarkr\checkmarkc\checkmark$\checkmarkl\checkmark$\checkmark^\checkmark\\checkmarkc\checkmarki\checkmarkr\checkmarkc\checkmark$\checkmarka\checkmark$\checkmark^\checkmark\\checkmarkc\checkmarki\checkmarkr\checkmarkc\checkmark$\checkmarkr\checkmark$\checkmark^\checkmark\\checkmarkc\checkmarki\checkmarkr\checkmarkc\checkmark$\checkmark}\checkmark$\checkmark^\checkmark\\checkmarkc\checkmarki\checkmarkr\checkmarkc\checkmark$\checkmark{\checkmark$\checkmark^\checkmark\\checkmarkc\checkmarki\checkmarkr\checkmarkc\checkmark$\checkmark|\checkmark$\checkmark^\checkmark\\checkmarkc\checkmarki\checkmarkr\checkmarkc\checkmark$\checkmarkl\checkmark$\checkmark^\checkmark\\checkmarkc\checkmarki\checkmarkr\checkmarkc\checkmark$\checkmark|\checkmark$\checkmark^\checkmark\\checkmarkc\checkmarki\checkmarkr\checkmarkc\checkmark$\checkmarkc\checkmark$\checkmark^\checkmark\\checkmarkc\checkmarki\checkmarkr\checkmarkc\checkmark$\checkmark|\checkmark$\checkmark^\checkmark\\checkmarkc\checkmarki\checkmarkr\checkmarkc\checkmark$\checkmarkc\checkmark$\checkmark^\checkmark\\checkmarkc\checkmarki\checkmarkr\checkmarkc\checkmark$\checkmark|\checkmark$\checkmark^\checkmark\\checkmarkc\checkmarki\checkmarkr\checkmarkc\checkmark$\checkmarkc\checkmark$\checkmark^\checkmark\\checkmarkc\checkmarki\checkmarkr\checkmarkc\checkmark$\checkmark|\checkmark$\checkmark^\checkmark\\checkmarkc\checkmarki\checkmarkr\checkmarkc\checkmark$\checkmarkl\checkmark$\checkmark^\checkmark\\checkmarkc\checkmarki\checkmarkr\checkmarkc\checkmark$\checkmark|\checkmark$\checkmark^\checkmark\\checkmarkc\checkmarki\checkmarkr\checkmarkc\checkmark$\checkmark}\checkmark$\checkmark^\checkmark\\checkmarkc\checkmarki\checkmarkr\checkmarkc\checkmark$\checkmark
\checkmark$\checkmark^\checkmark\\checkmarkc\checkmarki\checkmarkr\checkmarkc\checkmark$\checkmark\\checkmark$\checkmark^\checkmark\\checkmarkc\checkmarki\checkmarkr\checkmarkc\checkmark$\checkmarkh\checkmark$\checkmark^\checkmark\\checkmarkc\checkmarki\checkmarkr\checkmarkc\checkmark$\checkmarkl\checkmark$\checkmark^\checkmark\\checkmarkc\checkmarki\checkmarkr\checkmarkc\checkmark$\checkmarki\checkmark$\checkmark^\checkmark\\checkmarkc\checkmarki\checkmarkr\checkmarkc\checkmark$\checkmarkn\checkmark$\checkmark^\checkmark\\checkmarkc\checkmarki\checkmarkr\checkmarkc\checkmark$\checkmarke\checkmark$\checkmark^\checkmark\\checkmarkc\checkmarki\checkmarkr\checkmarkc\checkmark$\checkmark
\checkmark$\checkmark^\checkmark\\checkmarkc\checkmarki\checkmarkr\checkmarkc\checkmark$\checkmark\\checkmark$\checkmark^\checkmark\\checkmarkc\checkmarki\checkmarkr\checkmarkc\checkmark$\checkmarkt\checkmark$\checkmark^\checkmark\\checkmarkc\checkmarki\checkmarkr\checkmarkc\checkmark$\checkmarke\checkmark$\checkmark^\checkmark\\checkmarkc\checkmarki\checkmarkr\checkmarkc\checkmark$\checkmarkx\checkmark$\checkmark^\checkmark\\checkmarkc\checkmarki\checkmarkr\checkmarkc\checkmark$\checkmarkt\checkmark$\checkmark^\checkmark\\checkmarkc\checkmarki\checkmarkr\checkmarkc\checkmark$\checkmarkb\checkmark$\checkmark^\checkmark\\checkmarkc\checkmarki\checkmarkr\checkmarkc\checkmark$\checkmarkf\checkmark$\checkmark^\checkmark\\checkmarkc\checkmarki\checkmarkr\checkmarkc\checkmark$\checkmark{\checkmark$\checkmark^\checkmark\\checkmarkc\checkmarki\checkmarkr\checkmarkc\checkmark$\checkmarkA\checkmark$\checkmark^\checkmark\\checkmarkc\checkmarki\checkmarkr\checkmarkc\checkmark$\checkmarkx\checkmark$\checkmark^\checkmark\\checkmarkc\checkmarki\checkmarkr\checkmarkc\checkmark$\checkmarke\checkmark$\checkmark^\checkmark\\checkmarkc\checkmarki\checkmarkr\checkmarkc\checkmark$\checkmark}\checkmark$\checkmark^\checkmark\\checkmarkc\checkmarki\checkmarkr\checkmarkc\checkmark$\checkmark \checkmark$\checkmark^\checkmark\\checkmarkc\checkmarki\checkmarkr\checkmarkc\checkmark$\checkmark&\checkmark$\checkmark^\checkmark\\checkmarkc\checkmarki\checkmarkr\checkmarkc\checkmark$\checkmark \checkmark$\checkmark^\checkmark\\checkmarkc\checkmarki\checkmarkr\checkmarkc\checkmark$\checkmark\\checkmark$\checkmark^\checkmark\\checkmarkc\checkmarki\checkmarkr\checkmarkc\checkmark$\checkmarkt\checkmark$\checkmark^\checkmark\\checkmarkc\checkmarki\checkmarkr\checkmarkc\checkmark$\checkmarke\checkmark$\checkmark^\checkmark\\checkmarkc\checkmarki\checkmarkr\checkmarkc\checkmark$\checkmarkx\checkmark$\checkmark^\checkmark\\checkmarkc\checkmarki\checkmarkr\checkmarkc\checkmark$\checkmarkt\checkmark$\checkmark^\checkmark\\checkmarkc\checkmarki\checkmarkr\checkmarkc\checkmark$\checkmarkb\checkmark$\checkmark^\checkmark\\checkmarkc\checkmarki\checkmarkr\checkmarkc\checkmark$\checkmarkf\checkmark$\checkmark^\checkmark\\checkmarkc\checkmarki\checkmarkr\checkmarkc\checkmark$\checkmark{\checkmark$\checkmark^\checkmark\\checkmarkc\checkmarki\checkmarkr\checkmarkc\checkmark$\checkmarkM\checkmark$\checkmark^\checkmark\\checkmarkc\checkmarki\checkmarkr\checkmarkc\checkmark$\checkmarko\checkmark$\checkmark^\checkmark\\checkmarkc\checkmarki\checkmarkr\checkmarkc\checkmark$\checkmarkd\checkmark$\checkmark^\checkmark\\checkmarkc\checkmarki\checkmarkr\checkmarkc\checkmark$\checkmarkè\checkmark$\checkmark^\checkmark\\checkmarkc\checkmarki\checkmarkr\checkmarkc\checkmark$\checkmarkl\checkmark$\checkmark^\checkmark\\checkmarkc\checkmarki\checkmarkr\checkmarkc\checkmark$\checkmarke\checkmark$\checkmark^\checkmark\\checkmarkc\checkmarki\checkmarkr\checkmarkc\checkmark$\checkmark}\checkmark$\checkmark^\checkmark\\checkmarkc\checkmarki\checkmarkr\checkmarkc\checkmark$\checkmark \checkmark$\checkmark^\checkmark\\checkmarkc\checkmarki\checkmarkr\checkmarkc\checkmark$\checkmark&\checkmark$\checkmark^\checkmark\\checkmarkc\checkmarki\checkmarkr\checkmarkc\checkmark$\checkmark \checkmark$\checkmark^\checkmark\\checkmarkc\checkmarki\checkmarkr\checkmarkc\checkmark$\checkmark\\checkmark$\checkmark^\checkmark\\checkmarkc\checkmarki\checkmarkr\checkmarkc\checkmark$\checkmarkt\checkmark$\checkmark^\checkmark\\checkmarkc\checkmarki\checkmarkr\checkmarkc\checkmark$\checkmarke\checkmark$\checkmark^\checkmark\\checkmarkc\checkmarki\checkmarkr\checkmarkc\checkmark$\checkmarkx\checkmark$\checkmark^\checkmark\\checkmarkc\checkmarki\checkmarkr\checkmarkc\checkmark$\checkmarkt\checkmark$\checkmark^\checkmark\\checkmarkc\checkmarki\checkmarkr\checkmarkc\checkmark$\checkmarkb\checkmark$\checkmark^\checkmark\\checkmarkc\checkmarki\checkmarkr\checkmarkc\checkmark$\checkmarkf\checkmark$\checkmark^\checkmark\\checkmarkc\checkmarki\checkmarkr\checkmarkc\checkmark$\checkmark{\checkmark$\checkmark^\checkmark\\checkmarkc\checkmarki\checkmarkr\checkmarkc\checkmark$\checkmarkC\checkmark$\checkmark^\checkmark\\checkmarkc\checkmarki\checkmarkr\checkmarkc\checkmark$\checkmarko\checkmark$\checkmark^\checkmark\\checkmarkc\checkmarki\checkmarkr\checkmarkc\checkmark$\checkmarku\checkmark$\checkmark^\checkmark\\checkmarkc\checkmarki\checkmarkr\checkmarkc\checkmark$\checkmarkp\checkmark$\checkmark^\checkmark\\checkmarkc\checkmarki\checkmarkr\checkmarkc\checkmark$\checkmarkl\checkmark$\checkmark^\checkmark\\checkmarkc\checkmarki\checkmarkr\checkmarkc\checkmark$\checkmarke\checkmark$\checkmark^\checkmark\\checkmarkc\checkmarki\checkmarkr\checkmarkc\checkmark$\checkmark \checkmark$\checkmark^\checkmark\\checkmarkc\checkmarki\checkmarkr\checkmarkc\checkmark$\checkmarkm\checkmark$\checkmark^\checkmark\\checkmarkc\checkmarki\checkmarkr\checkmarkc\checkmark$\checkmarka\checkmark$\checkmark^\checkmark\\checkmarkc\checkmarki\checkmarkr\checkmarkc\checkmark$\checkmarki\checkmark$\checkmark^\checkmark\\checkmarkc\checkmarki\checkmarkr\checkmarkc\checkmark$\checkmarkn\checkmark$\checkmark^\checkmark\\checkmarkc\checkmarki\checkmarkr\checkmarkc\checkmark$\checkmarkt\checkmark$\checkmark^\checkmark\\checkmarkc\checkmarki\checkmarkr\checkmarkc\checkmark$\checkmarki\checkmark$\checkmark^\checkmark\\checkmarkc\checkmarki\checkmarkr\checkmarkc\checkmark$\checkmarke\checkmark$\checkmark^\checkmark\\checkmarkc\checkmarki\checkmarkr\checkmarkc\checkmark$\checkmarkn\checkmark$\checkmark^\checkmark\\checkmarkc\checkmarki\checkmarkr\checkmarkc\checkmark$\checkmark}\checkmark$\checkmark^\checkmark\\checkmarkc\checkmarki\checkmarkr\checkmarkc\checkmark$\checkmark \checkmark$\checkmark^\checkmark\\checkmarkc\checkmarki\checkmarkr\checkmarkc\checkmark$\checkmark&\checkmark$\checkmark^\checkmark\\checkmarkc\checkmarki\checkmarkr\checkmarkc\checkmark$\checkmark \checkmark$\checkmark^\checkmark\\checkmarkc\checkmarki\checkmarkr\checkmarkc\checkmark$\checkmark\\checkmark$\checkmark^\checkmark\\checkmarkc\checkmarki\checkmarkr\checkmarkc\checkmark$\checkmarkt\checkmark$\checkmark^\checkmark\\checkmarkc\checkmarki\checkmarkr\checkmarkc\checkmark$\checkmarke\checkmark$\checkmark^\checkmark\\checkmarkc\checkmarki\checkmarkr\checkmarkc\checkmark$\checkmarkx\checkmark$\checkmark^\checkmark\\checkmarkc\checkmarki\checkmarkr\checkmarkc\checkmark$\checkmarkt\checkmark$\checkmark^\checkmark\\checkmarkc\checkmarki\checkmarkr\checkmarkc\checkmark$\checkmarkb\checkmark$\checkmark^\checkmark\\checkmarkc\checkmarki\checkmarkr\checkmarkc\checkmark$\checkmarkf\checkmark$\checkmark^\checkmark\\checkmarkc\checkmarki\checkmarkr\checkmarkc\checkmark$\checkmark{\checkmark$\checkmark^\checkmark\\checkmarkc\checkmarki\checkmarkr\checkmarkc\checkmark$\checkmarkC\checkmark$\checkmark^\checkmark\\checkmarkc\checkmarki\checkmarkr\checkmarkc\checkmark$\checkmarko\checkmark$\checkmark^\checkmark\\checkmarkc\checkmarki\checkmarkr\checkmarkc\checkmark$\checkmarku\checkmark$\checkmark^\checkmark\\checkmarkc\checkmarki\checkmarkr\checkmarkc\checkmark$\checkmarkr\checkmark$\checkmark^\checkmark\\checkmarkc\checkmarki\checkmarkr\checkmarkc\checkmark$\checkmarka\checkmark$\checkmark^\checkmark\\checkmarkc\checkmarki\checkmarkr\checkmarkc\checkmark$\checkmarkn\checkmark$\checkmark^\checkmark\\checkmarkc\checkmarki\checkmarkr\checkmarkc\checkmark$\checkmarkt\checkmark$\checkmark^\checkmark\\checkmarkc\checkmarki\checkmarkr\checkmarkc\checkmark$\checkmark}\checkmark$\checkmark^\checkmark\\checkmarkc\checkmarki\checkmarkr\checkmarkc\checkmark$\checkmark \checkmark$\checkmark^\checkmark\\checkmarkc\checkmarki\checkmarkr\checkmarkc\checkmark$\checkmark&\checkmark$\checkmark^\checkmark\\checkmarkc\checkmarki\checkmarkr\checkmarkc\checkmark$\checkmark \checkmark$\checkmark^\checkmark\\checkmarkc\checkmarki\checkmarkr\checkmarkc\checkmark$\checkmark\\checkmark$\checkmark^\checkmark\\checkmarkc\checkmarki\checkmarkr\checkmarkc\checkmark$\checkmarkt\checkmark$\checkmark^\checkmark\\checkmarkc\checkmarki\checkmarkr\checkmarkc\checkmark$\checkmarke\checkmark$\checkmark^\checkmark\\checkmarkc\checkmarki\checkmarkr\checkmarkc\checkmark$\checkmarkx\checkmark$\checkmark^\checkmark\\checkmarkc\checkmarki\checkmarkr\checkmarkc\checkmark$\checkmarkt\checkmark$\checkmark^\checkmark\\checkmarkc\checkmarki\checkmarkr\checkmarkc\checkmark$\checkmarkb\checkmark$\checkmark^\checkmark\\checkmarkc\checkmarki\checkmarkr\checkmarkc\checkmark$\checkmarkf\checkmark$\checkmark^\checkmark\\checkmarkc\checkmarki\checkmarkr\checkmarkc\checkmark$\checkmark{\checkmark$\checkmark^\checkmark\\checkmarkc\checkmarki\checkmarkr\checkmarkc\checkmark$\checkmarkJ\checkmark$\checkmark^\checkmark\\checkmarkc\checkmarki\checkmarkr\checkmarkc\checkmark$\checkmarku\checkmark$\checkmark^\checkmark\\checkmarkc\checkmarki\checkmarkr\checkmarkc\checkmark$\checkmarks\checkmark$\checkmark^\checkmark\\checkmarkc\checkmarki\checkmarkr\checkmarkc\checkmark$\checkmarkt\checkmark$\checkmark^\checkmark\\checkmarkc\checkmarki\checkmarkr\checkmarkc\checkmark$\checkmarki\checkmark$\checkmark^\checkmark\\checkmarkc\checkmarki\checkmarkr\checkmarkc\checkmark$\checkmarkf\checkmark$\checkmark^\checkmark\\checkmarkc\checkmarki\checkmarkr\checkmarkc\checkmark$\checkmarki\checkmark$\checkmark^\checkmark\\checkmarkc\checkmarki\checkmarkr\checkmarkc\checkmark$\checkmarkc\checkmark$\checkmark^\checkmark\\checkmarkc\checkmarki\checkmarkr\checkmarkc\checkmark$\checkmarka\checkmark$\checkmark^\checkmark\\checkmarkc\checkmarki\checkmarkr\checkmarkc\checkmark$\checkmarkt\checkmark$\checkmark^\checkmark\\checkmarkc\checkmarki\checkmarkr\checkmarkc\checkmark$\checkmarki\checkmark$\checkmark^\checkmark\\checkmarkc\checkmarki\checkmarkr\checkmarkc\checkmark$\checkmarko\checkmark$\checkmark^\checkmark\\checkmarkc\checkmarki\checkmarkr\checkmarkc\checkmark$\checkmarkn\checkmark$\checkmark^\checkmark\\checkmarkc\checkmarki\checkmarkr\checkmarkc\checkmark$\checkmark}\checkmark$\checkmark^\checkmark\\checkmarkc\checkmarki\checkmarkr\checkmarkc\checkmark$\checkmark \checkmark$\checkmark^\checkmark\\checkmarkc\checkmarki\checkmarkr\checkmarkc\checkmark$\checkmark\\checkmark$\checkmark^\checkmark\\checkmarkc\checkmarki\checkmarkr\checkmarkc\checkmark$\checkmark\\checkmark$\checkmark^\checkmark\\checkmarkc\checkmarki\checkmarkr\checkmarkc\checkmark$\checkmark \checkmark$\checkmark^\checkmark\\checkmarkc\checkmarki\checkmarkr\checkmarkc\checkmark$\checkmark\\checkmark$\checkmark^\checkmark\\checkmarkc\checkmarki\checkmarkr\checkmarkc\checkmark$\checkmarkh\checkmark$\checkmark^\checkmark\\checkmarkc\checkmarki\checkmarkr\checkmarkc\checkmark$\checkmarkl\checkmark$\checkmark^\checkmark\\checkmarkc\checkmarki\checkmarkr\checkmarkc\checkmark$\checkmarki\checkmark$\checkmark^\checkmark\\checkmarkc\checkmarki\checkmarkr\checkmarkc\checkmark$\checkmarkn\checkmark$\checkmark^\checkmark\\checkmarkc\checkmarki\checkmarkr\checkmarkc\checkmark$\checkmarke\checkmark$\checkmark^\checkmark\\checkmarkc\checkmarki\checkmarkr\checkmarkc\checkmark$\checkmark
\checkmark$\checkmark^\checkmark\\checkmarkc\checkmarki\checkmarkr\checkmarkc\checkmark$\checkmarkX\checkmark$\checkmark^\checkmark\\checkmarkc\checkmarki\checkmarkr\checkmarkc\checkmark$\checkmark,\checkmark$\checkmark^\checkmark\\checkmarkc\checkmarki\checkmarkr\checkmarkc\checkmark$\checkmark \checkmark$\checkmark^\checkmark\\checkmarkc\checkmarki\checkmarkr\checkmarkc\checkmark$\checkmarkY\checkmark$\checkmark^\checkmark\\checkmarkc\checkmarki\checkmarkr\checkmarkc\checkmark$\checkmark,\checkmark$\checkmark^\checkmark\\checkmarkc\checkmarki\checkmarkr\checkmarkc\checkmark$\checkmark \checkmark$\checkmark^\checkmark\\checkmarkc\checkmarki\checkmarkr\checkmarkc\checkmark$\checkmarkZ\checkmark$\checkmark^\checkmark\\checkmarkc\checkmarki\checkmarkr\checkmarkc\checkmark$\checkmark \checkmark$\checkmark^\checkmark\\checkmarkc\checkmarki\checkmarkr\checkmarkc\checkmark$\checkmark&\checkmark$\checkmark^\checkmark\\checkmarkc\checkmarki\checkmarkr\checkmarkc\checkmark$\checkmark \checkmark$\checkmark^\checkmark\\checkmarkc\checkmarki\checkmarkr\checkmarkc\checkmark$\checkmarkN\checkmark$\checkmark^\checkmark\\checkmarkc\checkmarki\checkmarkr\checkmarkc\checkmark$\checkmarkE\checkmark$\checkmark^\checkmark\\checkmarkc\checkmarki\checkmarkr\checkmarkc\checkmark$\checkmarkM\checkmark$\checkmark^\checkmark\\checkmarkc\checkmarki\checkmarkr\checkmarkc\checkmark$\checkmarkA\checkmark$\checkmark^\checkmark\\checkmarkc\checkmarki\checkmarkr\checkmarkc\checkmark$\checkmark \checkmark$\checkmark^\checkmark\\checkmarkc\checkmarki\checkmarkr\checkmarkc\checkmark$\checkmark1\checkmark$\checkmark^\checkmark\\checkmarkc\checkmarki\checkmarkr\checkmarkc\checkmark$\checkmark7\checkmark$\checkmark^\checkmark\\checkmarkc\checkmarki\checkmarkr\checkmarkc\checkmark$\checkmark \checkmark$\checkmark^\checkmark\\checkmarkc\checkmarki\checkmarkr\checkmarkc\checkmark$\checkmark(\checkmark$\checkmark^\checkmark\\checkmarkc\checkmarki\checkmarkr\checkmarkc\checkmark$\checkmark1\checkmark$\checkmark^\checkmark\\checkmarkc\checkmarki\checkmarkr\checkmarkc\checkmark$\checkmark7\checkmark$\checkmark^\checkmark\\checkmarkc\checkmarki\checkmarkr\checkmarkc\checkmark$\checkmarkH\checkmark$\checkmark^\checkmark\\checkmarkc\checkmarki\checkmarkr\checkmarkc\checkmark$\checkmarkS\checkmark$\checkmark^\checkmark\\checkmarkc\checkmarki\checkmarkr\checkmarkc\checkmark$\checkmark4\checkmark$\checkmark^\checkmark\\checkmarkc\checkmarki\checkmarkr\checkmarkc\checkmark$\checkmark4\checkmark$\checkmark^\checkmark\\checkmarkc\checkmarki\checkmarkr\checkmarkc\checkmark$\checkmark0\checkmark$\checkmark^\checkmark\\checkmarkc\checkmarki\checkmarkr\checkmarkc\checkmark$\checkmark1\checkmark$\checkmark^\checkmark\\checkmarkc\checkmarki\checkmarkr\checkmarkc\checkmark$\checkmark)\checkmark$\checkmark^\checkmark\\checkmarkc\checkmarki\checkmarkr\checkmarkc\checkmark$\checkmark \checkmark$\checkmark^\checkmark\\checkmarkc\checkmarki\checkmarkr\checkmarkc\checkmark$\checkmark&\checkmark$\checkmark^\checkmark\\checkmarkc\checkmarki\checkmarkr\checkmarkc\checkmark$\checkmark \checkmark$\checkmark^\checkmark\\checkmarkc\checkmarki\checkmarkr\checkmarkc\checkmark$\checkmark$\checkmark$\checkmark^\checkmark\\checkmarkc\checkmarki\checkmarkr\checkmarkc\checkmark$\checkmark0\checkmark$\checkmark^\checkmark\\checkmarkc\checkmarki\checkmarkr\checkmarkc\checkmark$\checkmark{\checkmark$\checkmark^\checkmark\\checkmarkc\checkmarki\checkmarkr\checkmarkc\checkmark$\checkmark,\checkmark$\checkmark^\checkmark\\checkmarkc\checkmarki\checkmarkr\checkmarkc\checkmark$\checkmark}\checkmark$\checkmark^\checkmark\\checkmarkc\checkmarki\checkmarkr\checkmarkc\checkmark$\checkmark4\checkmark$\checkmark^\checkmark\\checkmarkc\checkmarki\checkmarkr\checkmarkc\checkmark$\checkmark2\checkmark$\checkmark^\checkmark\\checkmarkc\checkmarki\checkmarkr\checkmarkc\checkmark$\checkmark$\checkmark$\checkmark^\checkmark\\checkmarkc\checkmarki\checkmarkr\checkmarkc\checkmark$\checkmark \checkmark$\checkmark^\checkmark\\checkmarkc\checkmarki\checkmarkr\checkmarkc\checkmark$\checkmarkN\checkmark$\checkmark^\checkmark\\checkmarkc\checkmarki\checkmarkr\checkmarkc\checkmark$\checkmark$\checkmark$\checkmark^\checkmark\\checkmarkc\checkmarki\checkmarkr\checkmarkc\checkmark$\checkmark\\checkmark$\checkmark^\checkmark\\checkmarkc\checkmarki\checkmarkr\checkmarkc\checkmark$\checkmarkc\checkmark$\checkmark^\checkmark\\checkmarkc\checkmarki\checkmarkr\checkmarkc\checkmark$\checkmarkd\checkmark$\checkmark^\checkmark\\checkmarkc\checkmarki\checkmarkr\checkmarkc\checkmark$\checkmarko\checkmark$\checkmark^\checkmark\\checkmarkc\checkmarki\checkmarkr\checkmarkc\checkmark$\checkmarkt\checkmark$\checkmark^\checkmark\\checkmarkc\checkmarki\checkmarkr\checkmarkc\checkmark$\checkmark$\checkmark$\checkmark^\checkmark\\checkmarkc\checkmarki\checkmarkr\checkmarkc\checkmark$\checkmarkm\checkmark$\checkmark^\checkmark\\checkmarkc\checkmarki\checkmarkr\checkmarkc\checkmark$\checkmark \checkmark$\checkmark^\checkmark\\checkmarkc\checkmarki\checkmarkr\checkmarkc\checkmark$\checkmark&\checkmark$\checkmark^\checkmark\\checkmarkc\checkmarki\checkmarkr\checkmarkc\checkmark$\checkmark \checkmark$\checkmark^\checkmark\\checkmarkc\checkmarki\checkmarkr\checkmarkc\checkmark$\checkmark$\checkmark$\checkmark^\checkmark\\checkmarkc\checkmarki\checkmarkr\checkmarkc\checkmark$\checkmark1\checkmark$\checkmark^\checkmark\\checkmarkc\checkmarki\checkmarkr\checkmarkc\checkmark$\checkmark{\checkmark$\checkmark^\checkmark\\checkmarkc\checkmarki\checkmarkr\checkmarkc\checkmark$\checkmark,\checkmark$\checkmark^\checkmark\\checkmarkc\checkmarki\checkmarkr\checkmarkc\checkmark$\checkmark}\checkmark$\checkmark^\checkmark\\checkmarkc\checkmarki\checkmarkr\checkmarkc\checkmark$\checkmark7\checkmark$\checkmark^\checkmark\\checkmarkc\checkmarki\checkmarkr\checkmarkc\checkmark$\checkmark$\checkmark$\checkmark^\checkmark\\checkmarkc\checkmarki\checkmarkr\checkmarkc\checkmark$\checkmark \checkmark$\checkmark^\checkmark\\checkmarkc\checkmarki\checkmarkr\checkmarkc\checkmark$\checkmarkA\checkmark$\checkmark^\checkmark\\checkmarkc\checkmarki\checkmarkr\checkmarkc\checkmark$\checkmark \checkmark$\checkmark^\checkmark\\checkmarkc\checkmarki\checkmarkr\checkmarkc\checkmark$\checkmark&\checkmark$\checkmark^\checkmark\\checkmarkc\checkmarki\checkmarkr\checkmarkc\checkmark$\checkmark \checkmark$\checkmark^\checkmark\\checkmarkc\checkmarki\checkmarkr\checkmarkc\checkmark$\checkmarkT\checkmark$\checkmark^\checkmark\\checkmarkc\checkmarki\checkmarkr\checkmarkc\checkmark$\checkmarkr\checkmark$\checkmark^\checkmark\\checkmarkc\checkmarki\checkmarkr\checkmarkc\checkmark$\checkmarka\checkmark$\checkmark^\checkmark\\checkmarkc\checkmarki\checkmarkr\checkmarkc\checkmark$\checkmarkn\checkmark$\checkmark^\checkmark\\checkmarkc\checkmarki\checkmarkr\checkmarkc\checkmark$\checkmarks\checkmark$\checkmark^\checkmark\\checkmarkc\checkmarki\checkmarkr\checkmarkc\checkmark$\checkmarkl\checkmark$\checkmark^\checkmark\\checkmarkc\checkmarki\checkmarkr\checkmarkc\checkmark$\checkmarka\checkmark$\checkmark^\checkmark\\checkmarkc\checkmarki\checkmarkr\checkmarkc\checkmark$\checkmarkt\checkmark$\checkmark^\checkmark\\checkmarkc\checkmarki\checkmarkr\checkmarkc\checkmark$\checkmarki\checkmark$\checkmark^\checkmark\\checkmarkc\checkmarki\checkmarkr\checkmarkc\checkmark$\checkmarko\checkmark$\checkmark^\checkmark\\checkmarkc\checkmarki\checkmarkr\checkmarkc\checkmark$\checkmarkn\checkmark$\checkmark^\checkmark\\checkmarkc\checkmarki\checkmarkr\checkmarkc\checkmark$\checkmark \checkmark$\checkmark^\checkmark\\checkmarkc\checkmarki\checkmarkr\checkmarkc\checkmark$\checkmarkl\checkmark$\checkmark^\checkmark\\checkmarkc\checkmarki\checkmarkr\checkmarkc\checkmark$\checkmarké\checkmark$\checkmark^\checkmark\\checkmarkc\checkmarki\checkmarkr\checkmarkc\checkmark$\checkmarkg\checkmark$\checkmark^\checkmark\\checkmarkc\checkmarki\checkmarkr\checkmarkc\checkmark$\checkmarkè\checkmark$\checkmark^\checkmark\\checkmarkc\checkmarki\checkmarkr\checkmarkc\checkmark$\checkmarkr\checkmark$\checkmark^\checkmark\\checkmarkc\checkmarki\checkmarkr\checkmarkc\checkmark$\checkmarke\checkmark$\checkmark^\checkmark\\checkmarkc\checkmarki\checkmarkr\checkmarkc\checkmark$\checkmark \checkmark$\checkmark^\checkmark\\checkmarkc\checkmarki\checkmarkr\checkmarkc\checkmark$\checkmark\\checkmark$\checkmark^\checkmark\\checkmarkc\checkmarki\checkmarkr\checkmarkc\checkmark$\checkmark\\checkmark$\checkmark^\checkmark\\checkmarkc\checkmarki\checkmarkr\checkmarkc\checkmark$\checkmark \checkmark$\checkmark^\checkmark\\checkmarkc\checkmarki\checkmarkr\checkmarkc\checkmark$\checkmark\\checkmark$\checkmark^\checkmark\\checkmarkc\checkmarki\checkmarkr\checkmarkc\checkmark$\checkmarkh\checkmark$\checkmark^\checkmark\\checkmarkc\checkmarki\checkmarkr\checkmarkc\checkmark$\checkmarkl\checkmark$\checkmark^\checkmark\\checkmarkc\checkmarki\checkmarkr\checkmarkc\checkmark$\checkmarki\checkmark$\checkmark^\checkmark\\checkmarkc\checkmarki\checkmarkr\checkmarkc\checkmark$\checkmarkn\checkmark$\checkmark^\checkmark\\checkmarkc\checkmarki\checkmarkr\checkmarkc\checkmark$\checkmarke\checkmark$\checkmark^\checkmark\\checkmarkc\checkmarki\checkmarkr\checkmarkc\checkmark$\checkmark
\checkmark$\checkmark^\checkmark\\checkmarkc\checkmarki\checkmarkr\checkmarkc\checkmark$\checkmark\\checkmark$\checkmark^\checkmark\\checkmarkc\checkmarki\checkmarkr\checkmarkc\checkmark$\checkmarkt\checkmark$\checkmark^\checkmark\\checkmarkc\checkmarki\checkmarkr\checkmarkc\checkmark$\checkmarke\checkmark$\checkmark^\checkmark\\checkmarkc\checkmarki\checkmarkr\checkmarkc\checkmark$\checkmarkx\checkmark$\checkmark^\checkmark\\checkmarkc\checkmarki\checkmarkr\checkmarkc\checkmark$\checkmarkt\checkmark$\checkmark^\checkmark\\checkmarkc\checkmarki\checkmarkr\checkmarkc\checkmark$\checkmarkb\checkmark$\checkmark^\checkmark\\checkmarkc\checkmarki\checkmarkr\checkmarkc\checkmark$\checkmarkf\checkmark$\checkmark^\checkmark\\checkmarkc\checkmarki\checkmarkr\checkmarkc\checkmark$\checkmark{\checkmark$\checkmark^\checkmark\\checkmarkc\checkmarki\checkmarkr\checkmarkc\checkmark$\checkmarkA\checkmark$\checkmark^\checkmark\\checkmarkc\checkmarki\checkmarkr\checkmarkc\checkmark$\checkmark \checkmark$\checkmark^\checkmark\\checkmarkc\checkmarki\checkmarkr\checkmarkc\checkmark$\checkmark(\checkmark$\checkmark^\checkmark\\checkmarkc\checkmarki\checkmarkr\checkmarkc\checkmark$\checkmarkB\checkmark$\checkmark^\checkmark\\checkmarkc\checkmarki\checkmarkr\checkmarkc\checkmark$\checkmarke\checkmark$\checkmark^\checkmark\\checkmarkc\checkmarki\checkmarkr\checkmarkc\checkmark$\checkmarkr\checkmark$\checkmark^\checkmark\\checkmarkc\checkmarki\checkmarkr\checkmarkc\checkmark$\checkmarkc\checkmark$\checkmark^\checkmark\\checkmarkc\checkmarki\checkmarkr\checkmarkc\checkmark$\checkmarke\checkmark$\checkmark^\checkmark\\checkmarkc\checkmarki\checkmarkr\checkmarkc\checkmark$\checkmarka\checkmark$\checkmark^\checkmark\\checkmarkc\checkmarki\checkmarkr\checkmarkc\checkmark$\checkmarku\checkmark$\checkmark^\checkmark\\checkmarkc\checkmarki\checkmarkr\checkmarkc\checkmark$\checkmark)\checkmark$\checkmark^\checkmark\\checkmarkc\checkmarki\checkmarkr\checkmarkc\checkmark$\checkmark}\checkmark$\checkmark^\checkmark\\checkmarkc\checkmarki\checkmarkr\checkmarkc\checkmark$\checkmark \checkmark$\checkmark^\checkmark\\checkmarkc\checkmarki\checkmarkr\checkmarkc\checkmark$\checkmark&\checkmark$\checkmark^\checkmark\\checkmarkc\checkmarki\checkmarkr\checkmarkc\checkmark$\checkmark \checkmark$\checkmark^\checkmark\\checkmarkc\checkmarki\checkmarkr\checkmarkc\checkmark$\checkmark\\checkmark$\checkmark^\checkmark\\checkmarkc\checkmarki\checkmarkr\checkmarkc\checkmark$\checkmarkt\checkmark$\checkmark^\checkmark\\checkmarkc\checkmarki\checkmarkr\checkmarkc\checkmark$\checkmarke\checkmark$\checkmark^\checkmark\\checkmarkc\checkmarki\checkmarkr\checkmarkc\checkmark$\checkmarkx\checkmark$\checkmark^\checkmark\\checkmarkc\checkmarki\checkmarkr\checkmarkc\checkmark$\checkmarkt\checkmark$\checkmark^\checkmark\\checkmarkc\checkmarki\checkmarkr\checkmarkc\checkmark$\checkmarkb\checkmark$\checkmark^\checkmark\\checkmarkc\checkmarki\checkmarkr\checkmarkc\checkmark$\checkmarkf\checkmark$\checkmark^\checkmark\\checkmarkc\checkmarki\checkmarkr\checkmarkc\checkmark$\checkmark{\checkmark$\checkmark^\checkmark\\checkmarkc\checkmarki\checkmarkr\checkmarkc\checkmark$\checkmarkN\checkmark$\checkmark^\checkmark\\checkmarkc\checkmarki\checkmarkr\checkmarkc\checkmark$\checkmarkE\checkmark$\checkmark^\checkmark\\checkmarkc\checkmarki\checkmarkr\checkmarkc\checkmark$\checkmarkM\checkmark$\checkmark^\checkmark\\checkmarkc\checkmarki\checkmarkr\checkmarkc\checkmark$\checkmarkA\checkmark$\checkmark^\checkmark\\checkmarkc\checkmarki\checkmarkr\checkmarkc\checkmark$\checkmark \checkmark$\checkmark^\checkmark\\checkmarkc\checkmarki\checkmarkr\checkmarkc\checkmark$\checkmark2\checkmark$\checkmark^\checkmark\\checkmarkc\checkmarki\checkmarkr\checkmarkc\checkmark$\checkmark3\checkmark$\checkmark^\checkmark\\checkmarkc\checkmarki\checkmarkr\checkmarkc\checkmark$\checkmark \checkmark$\checkmark^\checkmark\\checkmarkc\checkmarki\checkmarkr\checkmarkc\checkmark$\checkmark(\checkmark$\checkmark^\checkmark\\checkmarkc\checkmarki\checkmarkr\checkmarkc\checkmark$\checkmark5\checkmark$\checkmark^\checkmark\\checkmarkc\checkmarki\checkmarkr\checkmarkc\checkmark$\checkmark7\checkmark$\checkmark^\checkmark\\checkmarkc\checkmarki\checkmarkr\checkmarkc\checkmark$\checkmarkH\checkmark$\checkmark^\checkmark\\checkmarkc\checkmarki\checkmarkr\checkmarkc\checkmark$\checkmarkS\checkmark$\checkmark^\checkmark\\checkmarkc\checkmarki\checkmarkr\checkmarkc\checkmark$\checkmark7\checkmark$\checkmark^\checkmark\\checkmarkc\checkmarki\checkmarkr\checkmarkc\checkmark$\checkmark6\checkmark$\checkmark^\checkmark\\checkmarkc\checkmarki\checkmarkr\checkmarkc\checkmark$\checkmark)\checkmark$\checkmark^\checkmark\\checkmarkc\checkmarki\checkmarkr\checkmarkc\checkmark$\checkmark}\checkmark$\checkmark^\checkmark\\checkmarkc\checkmarki\checkmarkr\checkmarkc\checkmark$\checkmark \checkmark$\checkmark^\checkmark\\checkmarkc\checkmarki\checkmarkr\checkmarkc\checkmark$\checkmark&\checkmark$\checkmark^\checkmark\\checkmarkc\checkmarki\checkmarkr\checkmarkc\checkmark$\checkmark \checkmark$\checkmark^\checkmark\\checkmarkc\checkmarki\checkmarkr\checkmarkc\checkmark$\checkmark$\checkmark$\checkmark^\checkmark\\checkmarkc\checkmarki\checkmarkr\checkmarkc\checkmark$\checkmark\\checkmark$\checkmark^\checkmark\\checkmarkc\checkmarki\checkmarkr\checkmarkc\checkmark$\checkmarkm\checkmark$\checkmark^\checkmark\\checkmarkc\checkmarki\checkmarkr\checkmarkc\checkmark$\checkmarka\checkmark$\checkmark^\checkmark\\checkmarkc\checkmarki\checkmarkr\checkmarkc\checkmark$\checkmarkt\checkmark$\checkmark^\checkmark\\checkmarkc\checkmarki\checkmarkr\checkmarkc\checkmark$\checkmarkh\checkmark$\checkmark^\checkmark\\checkmarkc\checkmarki\checkmarkr\checkmarkc\checkmark$\checkmarkb\checkmark$\checkmark^\checkmark\\checkmarkc\checkmarki\checkmarkr\checkmarkc\checkmark$\checkmarkf\checkmark$\checkmark^\checkmark\\checkmarkc\checkmarki\checkmarkr\checkmarkc\checkmark$\checkmark{\checkmark$\checkmark^\checkmark\\checkmarkc\checkmarki\checkmarkr\checkmarkc\checkmark$\checkmark1\checkmark$\checkmark^\checkmark\\checkmarkc\checkmarki\checkmarkr\checkmarkc\checkmark$\checkmark{\checkmark$\checkmark^\checkmark\\checkmarkc\checkmarki\checkmarkr\checkmarkc\checkmark$\checkmark,\checkmark$\checkmark^\checkmark\\checkmarkc\checkmarki\checkmarkr\checkmarkc\checkmark$\checkmark}\checkmark$\checkmark^\checkmark\\checkmarkc\checkmarki\checkmarkr\checkmarkc\checkmark$\checkmark2\checkmark$\checkmark^\checkmark\\checkmarkc\checkmarki\checkmarkr\checkmarkc\checkmark$\checkmark6\checkmark$\checkmark^\checkmark\\checkmarkc\checkmarki\checkmarkr\checkmarkc\checkmark$\checkmark}\checkmark$\checkmark^\checkmark\\checkmarkc\checkmarki\checkmarkr\checkmarkc\checkmark$\checkmark$\checkmark$\checkmark^\checkmark\\checkmarkc\checkmarki\checkmarkr\checkmarkc\checkmark$\checkmark \checkmark$\checkmark^\checkmark\\checkmarkc\checkmarki\checkmarkr\checkmarkc\checkmark$\checkmark\\checkmark$\checkmark^\checkmark\\checkmarkc\checkmarki\checkmarkr\checkmarkc\checkmark$\checkmarkt\checkmark$\checkmark^\checkmark\\checkmarkc\checkmarki\checkmarkr\checkmarkc\checkmark$\checkmarke\checkmark$\checkmark^\checkmark\\checkmarkc\checkmarki\checkmarkr\checkmarkc\checkmark$\checkmarkx\checkmark$\checkmark^\checkmark\\checkmarkc\checkmarki\checkmarkr\checkmarkc\checkmark$\checkmarkt\checkmark$\checkmark^\checkmark\\checkmarkc\checkmarki\checkmarkr\checkmarkc\checkmark$\checkmarkb\checkmark$\checkmark^\checkmark\\checkmarkc\checkmarki\checkmarkr\checkmarkc\checkmark$\checkmarkf\checkmark$\checkmark^\checkmark\\checkmarkc\checkmarki\checkmarkr\checkmarkc\checkmark$\checkmark{\checkmark$\checkmark^\checkmark\\checkmarkc\checkmarki\checkmarkr\checkmarkc\checkmark$\checkmarkN\checkmark$\checkmark^\checkmark\\checkmarkc\checkmarki\checkmarkr\checkmarkc\checkmark$\checkmark$\checkmark$\checkmark^\checkmark\\checkmarkc\checkmarki\checkmarkr\checkmarkc\checkmark$\checkmark\\checkmark$\checkmark^\checkmark\\checkmarkc\checkmarki\checkmarkr\checkmarkc\checkmark$\checkmarkc\checkmark$\checkmark^\checkmark\\checkmarkc\checkmarki\checkmarkr\checkmarkc\checkmark$\checkmarkd\checkmark$\checkmark^\checkmark\\checkmarkc\checkmarki\checkmarkr\checkmarkc\checkmark$\checkmarko\checkmark$\checkmark^\checkmark\\checkmarkc\checkmarki\checkmarkr\checkmarkc\checkmark$\checkmarkt\checkmark$\checkmark^\checkmark\\checkmarkc\checkmarki\checkmarkr\checkmarkc\checkmark$\checkmark$\checkmark$\checkmark^\checkmark\\checkmarkc\checkmarki\checkmarkr\checkmarkc\checkmark$\checkmarkm\checkmark$\checkmark^\checkmark\\checkmarkc\checkmarki\checkmarkr\checkmarkc\checkmark$\checkmark}\checkmark$\checkmark^\checkmark\\checkmarkc\checkmarki\checkmarkr\checkmarkc\checkmark$\checkmark \checkmark$\checkmark^\checkmark\\checkmarkc\checkmarki\checkmarkr\checkmarkc\checkmark$\checkmark&\checkmark$\checkmark^\checkmark\\checkmarkc\checkmarki\checkmarkr\checkmarkc\checkmark$\checkmark \checkmark$\checkmark^\checkmark\\checkmarkc\checkmarki\checkmarkr\checkmarkc\checkmark$\checkmark$\checkmark$\checkmark^\checkmark\\checkmarkc\checkmarki\checkmarkr\checkmarkc\checkmark$\checkmark3\checkmark$\checkmark^\checkmark\\checkmarkc\checkmarki\checkmarkr\checkmarkc\checkmark$\checkmark{\checkmark$\checkmark^\checkmark\\checkmarkc\checkmarki\checkmarkr\checkmarkc\checkmark$\checkmark,\checkmark$\checkmark^\checkmark\\checkmarkc\checkmarki\checkmarkr\checkmarkc\checkmark$\checkmark}\checkmark$\checkmark^\checkmark\\checkmarkc\checkmarki\checkmarkr\checkmarkc\checkmark$\checkmark0\checkmark$\checkmark^\checkmark\\checkmarkc\checkmarki\checkmarkr\checkmarkc\checkmark$\checkmark$\checkmark$\checkmark^\checkmark\\checkmarkc\checkmarki\checkmarkr\checkmarkc\checkmark$\checkmark \checkmark$\checkmark^\checkmark\\checkmarkc\checkmarki\checkmarkr\checkmarkc\checkmark$\checkmarkA\checkmark$\checkmark^\checkmark\\checkmarkc\checkmarki\checkmarkr\checkmarkc\checkmark$\checkmark \checkmark$\checkmark^\checkmark\\checkmarkc\checkmarki\checkmarkr\checkmarkc\checkmark$\checkmark&\checkmark$\checkmark^\checkmark\\checkmarkc\checkmarki\checkmarkr\checkmarkc\checkmark$\checkmark \checkmark$\checkmark^\checkmark\\checkmarkc\checkmarki\checkmarkr\checkmarkc\checkmark$\checkmarkM\checkmark$\checkmark^\checkmark\\checkmarkc\checkmarki\checkmarkr\checkmarkc\checkmark$\checkmarko\checkmark$\checkmark^\checkmark\\checkmarkc\checkmarki\checkmarkr\checkmarkc\checkmark$\checkmarkm\checkmark$\checkmark^\checkmark\\checkmarkc\checkmarki\checkmarkr\checkmarkc\checkmark$\checkmarke\checkmark$\checkmark^\checkmark\\checkmarkc\checkmarki\checkmarkr\checkmarkc\checkmark$\checkmarkn\checkmark$\checkmark^\checkmark\\checkmarkc\checkmarki\checkmarkr\checkmarkc\checkmark$\checkmarkt\checkmark$\checkmark^\checkmark\\checkmarkc\checkmarki\checkmarkr\checkmarkc\checkmark$\checkmark \checkmark$\checkmark^\checkmark\\checkmarkc\checkmarki\checkmarkr\checkmarkc\checkmark$\checkmarkd\checkmark$\checkmark^\checkmark\\checkmarkc\checkmarki\checkmarkr\checkmarkc\checkmark$\checkmarke\checkmark$\checkmark^\checkmark\\checkmarkc\checkmarki\checkmarkr\checkmarkc\checkmark$\checkmark \checkmark$\checkmark^\checkmark\\checkmarkc\checkmarki\checkmarkr\checkmarkc\checkmark$\checkmarkr\checkmark$\checkmark^\checkmark\\checkmarkc\checkmarki\checkmarkr\checkmarkc\checkmark$\checkmarke\checkmark$\checkmark^\checkmark\\checkmarkc\checkmarki\checkmarkr\checkmarkc\checkmark$\checkmarkn\checkmark$\checkmark^\checkmark\\checkmarkc\checkmarki\checkmarkr\checkmarkc\checkmark$\checkmarkv\checkmark$\checkmark^\checkmark\\checkmarkc\checkmarki\checkmarkr\checkmarkc\checkmark$\checkmarke\checkmark$\checkmark^\checkmark\\checkmarkc\checkmarki\checkmarkr\checkmarkc\checkmark$\checkmarkr\checkmark$\checkmark^\checkmark\\checkmarkc\checkmarki\checkmarkr\checkmarkc\checkmark$\checkmarks\checkmark$\checkmark^\checkmark\\checkmarkc\checkmarki\checkmarkr\checkmarkc\checkmark$\checkmarke\checkmark$\checkmark^\checkmark\\checkmarkc\checkmarki\checkmarkr\checkmarkc\checkmark$\checkmarkm\checkmark$\checkmark^\checkmark\\checkmarkc\checkmarki\checkmarkr\checkmarkc\checkmark$\checkmarke\checkmark$\checkmark^\checkmark\\checkmarkc\checkmarki\checkmarkr\checkmarkc\checkmark$\checkmarkn\checkmark$\checkmark^\checkmark\\checkmarkc\checkmarki\checkmarkr\checkmarkc\checkmark$\checkmarkt\checkmark$\checkmark^\checkmark\\checkmarkc\checkmarki\checkmarkr\checkmarkc\checkmark$\checkmark \checkmark$\checkmark^\checkmark\\checkmarkc\checkmarki\checkmarkr\checkmarkc\checkmark$\checkmark\\checkmark$\checkmark^\checkmark\\checkmarkc\checkmarki\checkmarkr\checkmarkc\checkmark$\checkmark\\checkmark$\checkmark^\checkmark\\checkmarkc\checkmarki\checkmarkr\checkmarkc\checkmark$\checkmark \checkmark$\checkmark^\checkmark\\checkmarkc\checkmarki\checkmarkr\checkmarkc\checkmark$\checkmark\\checkmark$\checkmark^\checkmark\\checkmarkc\checkmarki\checkmarkr\checkmarkc\checkmark$\checkmarkh\checkmark$\checkmark^\checkmark\\checkmarkc\checkmarki\checkmarkr\checkmarkc\checkmark$\checkmarkl\checkmark$\checkmark^\checkmark\\checkmarkc\checkmarki\checkmarkr\checkmarkc\checkmark$\checkmarki\checkmark$\checkmark^\checkmark\\checkmarkc\checkmarki\checkmarkr\checkmarkc\checkmark$\checkmarkn\checkmark$\checkmark^\checkmark\\checkmarkc\checkmarki\checkmarkr\checkmarkc\checkmark$\checkmarke\checkmark$\checkmark^\checkmark\\checkmarkc\checkmarki\checkmarkr\checkmarkc\checkmark$\checkmark
\checkmark$\checkmark^\checkmark\\checkmarkc\checkmarki\checkmarkr\checkmarkc\checkmark$\checkmarkB\checkmark$\checkmark^\checkmark\\checkmarkc\checkmarki\checkmarkr\checkmarkc\checkmark$\checkmark \checkmark$\checkmark^\checkmark\\checkmarkc\checkmarki\checkmarkr\checkmarkc\checkmark$\checkmark(\checkmark$\checkmark^\checkmark\\checkmarkc\checkmarki\checkmarkr\checkmarkc\checkmark$\checkmarkP\checkmark$\checkmark^\checkmark\\checkmarkc\checkmarki\checkmarkr\checkmarkc\checkmark$\checkmarkl\checkmark$\checkmark^\checkmark\\checkmarkc\checkmarki\checkmarkr\checkmarkc\checkmark$\checkmarka\checkmark$\checkmark^\checkmark\\checkmarkc\checkmarki\checkmarkr\checkmarkc\checkmark$\checkmarkt\checkmark$\checkmark^\checkmark\\checkmarkc\checkmarki\checkmarkr\checkmarkc\checkmark$\checkmarke\checkmark$\checkmark^\checkmark\\checkmarkc\checkmarki\checkmarkr\checkmarkc\checkmark$\checkmarka\checkmark$\checkmark^\checkmark\\checkmarkc\checkmarki\checkmarkr\checkmarkc\checkmark$\checkmarku\checkmark$\checkmark^\checkmark\\checkmarkc\checkmarki\checkmarkr\checkmarkc\checkmark$\checkmark)\checkmark$\checkmark^\checkmark\\checkmarkc\checkmarki\checkmarkr\checkmarkc\checkmark$\checkmark \checkmark$\checkmark^\checkmark\\checkmarkc\checkmarki\checkmarkr\checkmarkc\checkmark$\checkmark&\checkmark$\checkmark^\checkmark\\checkmarkc\checkmarki\checkmarkr\checkmarkc\checkmark$\checkmark \checkmark$\checkmark^\checkmark\\checkmarkc\checkmarki\checkmarkr\checkmarkc\checkmark$\checkmarkN\checkmark$\checkmark^\checkmark\\checkmarkc\checkmarki\checkmarkr\checkmarkc\checkmark$\checkmarkE\checkmark$\checkmark^\checkmark\\checkmarkc\checkmarki\checkmarkr\checkmarkc\checkmark$\checkmarkM\checkmark$\checkmark^\checkmark\\checkmarkc\checkmarki\checkmarkr\checkmarkc\checkmark$\checkmarkA\checkmark$\checkmark^\checkmark\\checkmarkc\checkmarki\checkmarkr\checkmarkc\checkmark$\checkmark \checkmark$\checkmark^\checkmark\\checkmarkc\checkmarki\checkmarkr\checkmarkc\checkmark$\checkmark1\checkmark$\checkmark^\checkmark\\checkmarkc\checkmarki\checkmarkr\checkmarkc\checkmark$\checkmark7\checkmark$\checkmark^\checkmark\\checkmarkc\checkmarki\checkmarkr\checkmarkc\checkmark$\checkmark \checkmark$\checkmark^\checkmark\\checkmarkc\checkmarki\checkmarkr\checkmarkc\checkmark$\checkmarkl\checkmark$\checkmark^\checkmark\\checkmarkc\checkmarki\checkmarkr\checkmarkc\checkmark$\checkmarko\checkmark$\checkmark^\checkmark\\checkmarkc\checkmarki\checkmarkr\checkmarkc\checkmark$\checkmarkn\checkmark$\checkmark^\checkmark\\checkmarkc\checkmarki\checkmarkr\checkmarkc\checkmark$\checkmarkg\checkmark$\checkmark^\checkmark\\checkmarkc\checkmarki\checkmarkr\checkmarkc\checkmark$\checkmark \checkmark$\checkmark^\checkmark\\checkmarkc\checkmarki\checkmarkr\checkmarkc\checkmark$\checkmark&\checkmark$\checkmark^\checkmark\\checkmarkc\checkmarki\checkmarkr\checkmarkc\checkmark$\checkmark \checkmark$\checkmark^\checkmark\\checkmarkc\checkmarki\checkmarkr\checkmarkc\checkmark$\checkmark$\checkmark$\checkmark^\checkmark\\checkmarkc\checkmarki\checkmarkr\checkmarkc\checkmark$\checkmark0\checkmark$\checkmark^\checkmark\\checkmarkc\checkmarki\checkmarkr\checkmarkc\checkmark$\checkmark{\checkmark$\checkmark^\checkmark\\checkmarkc\checkmarki\checkmarkr\checkmarkc\checkmark$\checkmark,\checkmark$\checkmark^\checkmark\\checkmarkc\checkmarki\checkmarkr\checkmarkc\checkmark$\checkmark}\checkmark$\checkmark^\checkmark\\checkmarkc\checkmarki\checkmarkr\checkmarkc\checkmark$\checkmark5\checkmark$\checkmark^\checkmark\\checkmarkc\checkmarki\checkmarkr\checkmarkc\checkmark$\checkmark5\checkmark$\checkmark^\checkmark\\checkmarkc\checkmarki\checkmarkr\checkmarkc\checkmark$\checkmark$\checkmark$\checkmark^\checkmark\\checkmarkc\checkmarki\checkmarkr\checkmarkc\checkmark$\checkmark \checkmark$\checkmark^\checkmark\\checkmarkc\checkmarki\checkmarkr\checkmarkc\checkmark$\checkmarkN\checkmark$\checkmark^\checkmark\\checkmarkc\checkmarki\checkmarkr\checkmarkc\checkmark$\checkmark$\checkmark$\checkmark^\checkmark\\checkmarkc\checkmarki\checkmarkr\checkmarkc\checkmark$\checkmark\\checkmark$\checkmark^\checkmark\\checkmarkc\checkmarki\checkmarkr\checkmarkc\checkmark$\checkmarkc\checkmark$\checkmark^\checkmark\\checkmarkc\checkmarki\checkmarkr\checkmarkc\checkmark$\checkmarkd\checkmark$\checkmark^\checkmark\\checkmarkc\checkmarki\checkmarkr\checkmarkc\checkmark$\checkmarko\checkmark$\checkmark^\checkmark\\checkmarkc\checkmarki\checkmarkr\checkmarkc\checkmark$\checkmarkt\checkmark$\checkmark^\checkmark\\checkmarkc\checkmarki\checkmarkr\checkmarkc\checkmark$\checkmark$\checkmark$\checkmark^\checkmark\\checkmarkc\checkmarki\checkmarkr\checkmarkc\checkmark$\checkmarkm\checkmark$\checkmark^\checkmark\\checkmarkc\checkmarki\checkmarkr\checkmarkc\checkmark$\checkmark \checkmark$\checkmark^\checkmark\\checkmarkc\checkmarki\checkmarkr\checkmarkc\checkmark$\checkmark&\checkmark$\checkmark^\checkmark\\checkmarkc\checkmarki\checkmarkr\checkmarkc\checkmark$\checkmark \checkmark$\checkmark^\checkmark\\checkmarkc\checkmarki\checkmarkr\checkmarkc\checkmark$\checkmark$\checkmark$\checkmark^\checkmark\\checkmarkc\checkmarki\checkmarkr\checkmarkc\checkmark$\checkmark2\checkmark$\checkmark^\checkmark\\checkmarkc\checkmarki\checkmarkr\checkmarkc\checkmark$\checkmark{\checkmark$\checkmark^\checkmark\\checkmarkc\checkmarki\checkmarkr\checkmarkc\checkmark$\checkmark,\checkmark$\checkmark^\checkmark\\checkmarkc\checkmarki\checkmarkr\checkmarkc\checkmark$\checkmark}\checkmark$\checkmark^\checkmark\\checkmarkc\checkmarki\checkmarkr\checkmarkc\checkmark$\checkmark0\checkmark$\checkmark^\checkmark\\checkmarkc\checkmarki\checkmarkr\checkmarkc\checkmark$\checkmark$\checkmark$\checkmark^\checkmark\\checkmarkc\checkmarki\checkmarkr\checkmarkc\checkmark$\checkmark \checkmark$\checkmark^\checkmark\\checkmarkc\checkmarki\checkmarkr\checkmarkc\checkmark$\checkmarkA\checkmark$\checkmark^\checkmark\\checkmarkc\checkmarki\checkmarkr\checkmarkc\checkmark$\checkmark \checkmark$\checkmark^\checkmark\\checkmarkc\checkmarki\checkmarkr\checkmarkc\checkmark$\checkmark&\checkmark$\checkmark^\checkmark\\checkmarkc\checkmarki\checkmarkr\checkmarkc\checkmark$\checkmark \checkmark$\checkmark^\checkmark\\checkmarkc\checkmarki\checkmarkr\checkmarkc\checkmark$\checkmarkI\checkmark$\checkmark^\checkmark\\checkmarkc\checkmarki\checkmarkr\checkmarkc\checkmark$\checkmarkn\checkmark$\checkmark^\checkmark\\checkmarkc\checkmarki\checkmarkr\checkmarkc\checkmark$\checkmarke\checkmark$\checkmark^\checkmark\\checkmarkc\checkmarki\checkmarkr\checkmarkc\checkmark$\checkmarkr\checkmark$\checkmark^\checkmark\\checkmarkc\checkmarki\checkmarkr\checkmarkc\checkmark$\checkmarkt\checkmark$\checkmark^\checkmark\\checkmarkc\checkmarki\checkmarkr\checkmarkc\checkmark$\checkmarki\checkmark$\checkmark^\checkmark\\checkmarkc\checkmarki\checkmarkr\checkmarkc\checkmark$\checkmarke\checkmark$\checkmark^\checkmark\\checkmarkc\checkmarki\checkmarkr\checkmarkc\checkmark$\checkmark \checkmark$\checkmark^\checkmark\\checkmarkc\checkmarki\checkmarkr\checkmarkc\checkmark$\checkmarkp\checkmark$\checkmark^\checkmark\\checkmarkc\checkmarki\checkmarkr\checkmarkc\checkmark$\checkmarkl\checkmark$\checkmark^\checkmark\\checkmarkc\checkmarki\checkmarkr\checkmarkc\checkmark$\checkmarka\checkmark$\checkmark^\checkmark\\checkmarkc\checkmarki\checkmarkr\checkmarkc\checkmark$\checkmarkt\checkmark$\checkmark^\checkmark\\checkmarkc\checkmarki\checkmarkr\checkmarkc\checkmark$\checkmarke\checkmark$\checkmark^\checkmark\\checkmarkc\checkmarki\checkmarkr\checkmarkc\checkmark$\checkmarka\checkmark$\checkmark^\checkmark\\checkmarkc\checkmarki\checkmarkr\checkmarkc\checkmark$\checkmarku\checkmark$\checkmark^\checkmark\\checkmarkc\checkmarki\checkmarkr\checkmarkc\checkmark$\checkmark \checkmark$\checkmark^\checkmark\\checkmarkc\checkmarki\checkmarkr\checkmarkc\checkmark$\checkmark\\checkmark$\checkmark^\checkmark\\checkmarkc\checkmarki\checkmarkr\checkmarkc\checkmark$\checkmark\\checkmark$\checkmark^\checkmark\\checkmarkc\checkmarki\checkmarkr\checkmarkc\checkmark$\checkmark \checkmark$\checkmark^\checkmark\\checkmarkc\checkmarki\checkmarkr\checkmarkc\checkmark$\checkmark\\checkmark$\checkmark^\checkmark\\checkmarkc\checkmarki\checkmarkr\checkmarkc\checkmark$\checkmarkh\checkmark$\checkmark^\checkmark\\checkmarkc\checkmarki\checkmarkr\checkmarkc\checkmark$\checkmarkl\checkmark$\checkmark^\checkmark\\checkmarkc\checkmarki\checkmarkr\checkmarkc\checkmark$\checkmarki\checkmark$\checkmark^\checkmark\\checkmarkc\checkmarki\checkmarkr\checkmarkc\checkmark$\checkmarkn\checkmark$\checkmark^\checkmark\\checkmarkc\checkmarki\checkmarkr\checkmarkc\checkmark$\checkmarke\checkmark$\checkmark^\checkmark\\checkmarkc\checkmarki\checkmarkr\checkmarkc\checkmark$\checkmark
\checkmark$\checkmark^\checkmark\\checkmarkc\checkmarki\checkmarkr\checkmarkc\checkmark$\checkmark\\checkmark$\checkmark^\checkmark\\checkmarkc\checkmarki\checkmarkr\checkmarkc\checkmark$\checkmarke\checkmark$\checkmark^\checkmark\\checkmarkc\checkmarki\checkmarkr\checkmarkc\checkmark$\checkmarkn\checkmark$\checkmark^\checkmark\\checkmarkc\checkmarki\checkmarkr\checkmarkc\checkmark$\checkmarkd\checkmark$\checkmark^\checkmark\\checkmarkc\checkmarki\checkmarkr\checkmarkc\checkmark$\checkmark{\checkmark$\checkmark^\checkmark\\checkmarkc\checkmarki\checkmarkr\checkmarkc\checkmark$\checkmarkt\checkmark$\checkmark^\checkmark\\checkmarkc\checkmarki\checkmarkr\checkmarkc\checkmark$\checkmarka\checkmark$\checkmark^\checkmark\\checkmarkc\checkmarki\checkmarkr\checkmarkc\checkmark$\checkmarkb\checkmark$\checkmark^\checkmark\\checkmarkc\checkmarki\checkmarkr\checkmarkc\checkmark$\checkmarku\checkmark$\checkmark^\checkmark\\checkmarkc\checkmarki\checkmarkr\checkmarkc\checkmark$\checkmarkl\checkmark$\checkmark^\checkmark\\checkmarkc\checkmarki\checkmarkr\checkmarkc\checkmark$\checkmarka\checkmark$\checkmark^\checkmark\\checkmarkc\checkmarki\checkmarkr\checkmarkc\checkmark$\checkmarkr\checkmark$\checkmark^\checkmark\\checkmarkc\checkmarki\checkmarkr\checkmarkc\checkmark$\checkmark}\checkmark$\checkmark^\checkmark\\checkmarkc\checkmarki\checkmarkr\checkmarkc\checkmark$\checkmark
\checkmark$\checkmark^\checkmark\\checkmarkc\checkmarki\checkmarkr\checkmarkc\checkmark$\checkmark\\checkmark$\checkmark^\checkmark\\checkmarkc\checkmarki\checkmarkr\checkmarkc\checkmark$\checkmarkc\checkmark$\checkmark^\checkmark\\checkmarkc\checkmarki\checkmarkr\checkmarkc\checkmark$\checkmarka\checkmark$\checkmark^\checkmark\\checkmarkc\checkmarki\checkmarkr\checkmarkc\checkmark$\checkmarkp\checkmark$\checkmark^\checkmark\\checkmarkc\checkmarki\checkmarkr\checkmarkc\checkmark$\checkmarkt\checkmark$\checkmark^\checkmark\\checkmarkc\checkmarki\checkmarkr\checkmarkc\checkmark$\checkmarki\checkmark$\checkmark^\checkmark\\checkmarkc\checkmarki\checkmarkr\checkmarkc\checkmark$\checkmarko\checkmark$\checkmark^\checkmark\\checkmarkc\checkmarki\checkmarkr\checkmarkc\checkmark$\checkmarkn\checkmark$\checkmark^\checkmark\\checkmarkc\checkmarki\checkmarkr\checkmarkc\checkmark$\checkmark{\checkmark$\checkmark^\checkmark\\checkmarkc\checkmarki\checkmarkr\checkmarkc\checkmark$\checkmarkT\checkmark$\checkmark^\checkmark\\checkmarkc\checkmarki\checkmarkr\checkmarkc\checkmark$\checkmarka\checkmark$\checkmark^\checkmark\\checkmarkc\checkmarki\checkmarkr\checkmarkc\checkmark$\checkmarkb\checkmark$\checkmark^\checkmark\\checkmarkc\checkmarki\checkmarkr\checkmarkc\checkmark$\checkmarkl\checkmark$\checkmark^\checkmark\\checkmarkc\checkmarki\checkmarkr\checkmarkc\checkmark$\checkmarke\checkmark$\checkmark^\checkmark\\checkmarkc\checkmarki\checkmarkr\checkmarkc\checkmark$\checkmarka\checkmark$\checkmark^\checkmark\\checkmarkc\checkmarki\checkmarkr\checkmarkc\checkmark$\checkmarku\checkmark$\checkmark^\checkmark\\checkmarkc\checkmarki\checkmarkr\checkmarkc\checkmark$\checkmark \checkmark$\checkmark^\checkmark\\checkmarkc\checkmarki\checkmarkr\checkmarkc\checkmark$\checkmarkd\checkmark$\checkmark^\checkmark\\checkmarkc\checkmarki\checkmarkr\checkmarkc\checkmark$\checkmarke\checkmark$\checkmark^\checkmark\\checkmarkc\checkmarki\checkmarkr\checkmarkc\checkmark$\checkmark \checkmark$\checkmark^\checkmark\\checkmarkc\checkmarki\checkmarkr\checkmarkc\checkmark$\checkmarkc\checkmark$\checkmark^\checkmark\\checkmarkc\checkmarki\checkmarkr\checkmarkc\checkmark$\checkmarkh\checkmark$\checkmark^\checkmark\\checkmarkc\checkmarki\checkmarkr\checkmarkc\checkmark$\checkmarko\checkmark$\checkmark^\checkmark\\checkmarkc\checkmarki\checkmarkr\checkmarkc\checkmark$\checkmarki\checkmark$\checkmark^\checkmark\\checkmarkc\checkmarki\checkmarkr\checkmarkc\checkmark$\checkmarkx\checkmark$\checkmark^\checkmark\\checkmarkc\checkmarki\checkmarkr\checkmarkc\checkmark$\checkmark \checkmark$\checkmark^\checkmark\\checkmarkc\checkmarki\checkmarkr\checkmarkc\checkmark$\checkmarkd\checkmark$\checkmark^\checkmark\\checkmarkc\checkmarki\checkmarkr\checkmarkc\checkmark$\checkmarke\checkmark$\checkmark^\checkmark\\checkmarkc\checkmarki\checkmarkr\checkmarkc\checkmark$\checkmarks\checkmark$\checkmark^\checkmark\\checkmarkc\checkmarki\checkmarkr\checkmarkc\checkmark$\checkmark \checkmark$\checkmark^\checkmark\\checkmarkc\checkmarki\checkmarkr\checkmarkc\checkmark$\checkmarkm\checkmark$\checkmark^\checkmark\\checkmarkc\checkmarki\checkmarkr\checkmarkc\checkmark$\checkmarko\checkmark$\checkmark^\checkmark\\checkmarkc\checkmarki\checkmarkr\checkmarkc\checkmark$\checkmarkt\checkmark$\checkmark^\checkmark\\checkmarkc\checkmarki\checkmarkr\checkmarkc\checkmark$\checkmarke\checkmark$\checkmark^\checkmark\\checkmarkc\checkmarki\checkmarkr\checkmarkc\checkmark$\checkmarku\checkmark$\checkmark^\checkmark\\checkmarkc\checkmarki\checkmarkr\checkmarkc\checkmark$\checkmarkr\checkmark$\checkmark^\checkmark\\checkmarkc\checkmarki\checkmarkr\checkmarkc\checkmark$\checkmarks\checkmark$\checkmark^\checkmark\\checkmarkc\checkmarki\checkmarkr\checkmarkc\checkmark$\checkmark \checkmark$\checkmark^\checkmark\\checkmarkc\checkmarki\checkmarkr\checkmarkc\checkmark$\checkmarkp\checkmark$\checkmark^\checkmark\\checkmarkc\checkmarki\checkmarkr\checkmarkc\checkmark$\checkmarka\checkmark$\checkmark^\checkmark\\checkmarkc\checkmarki\checkmarkr\checkmarkc\checkmark$\checkmarks\checkmark$\checkmark^\checkmark\\checkmarkc\checkmarki\checkmarkr\checkmarkc\checkmark$\checkmark \checkmark$\checkmark^\checkmark\\checkmarkc\checkmarki\checkmarkr\checkmarkc\checkmark$\checkmarkà\checkmark$\checkmark^\checkmark\\checkmarkc\checkmarki\checkmarkr\checkmarkc\checkmark$\checkmark \checkmark$\checkmark^\checkmark\\checkmarkc\checkmarki\checkmarkr\checkmarkc\checkmark$\checkmarkp\checkmark$\checkmark^\checkmark\\checkmarkc\checkmarki\checkmarkr\checkmarkc\checkmark$\checkmarka\checkmark$\checkmark^\checkmark\\checkmarkc\checkmarki\checkmarkr\checkmarkc\checkmark$\checkmarks\checkmark$\checkmark^\checkmark\\checkmarkc\checkmarki\checkmarkr\checkmarkc\checkmark$\checkmark \checkmark$\checkmark^\checkmark\\checkmarkc\checkmarki\checkmarkr\checkmarkc\checkmark$\checkmarkp\checkmark$\checkmark^\checkmark\\checkmarkc\checkmarki\checkmarkr\checkmarkc\checkmark$\checkmarko\checkmark$\checkmark^\checkmark\\checkmarkc\checkmarki\checkmarkr\checkmarkc\checkmark$\checkmarku\checkmark$\checkmark^\checkmark\\checkmarkc\checkmarki\checkmarkr\checkmarkc\checkmark$\checkmarkr\checkmark$\checkmark^\checkmark\\checkmarkc\checkmarki\checkmarkr\checkmarkc\checkmark$\checkmark \checkmark$\checkmark^\checkmark\\checkmarkc\checkmarki\checkmarkr\checkmarkc\checkmark$\checkmarkc\checkmark$\checkmark^\checkmark\\checkmarkc\checkmarki\checkmarkr\checkmarkc\checkmark$\checkmarkh\checkmark$\checkmark^\checkmark\\checkmarkc\checkmarki\checkmarkr\checkmarkc\checkmark$\checkmarka\checkmark$\checkmark^\checkmark\\checkmarkc\checkmarki\checkmarkr\checkmarkc\checkmark$\checkmarkq\checkmark$\checkmark^\checkmark\\checkmarkc\checkmarki\checkmarkr\checkmarkc\checkmark$\checkmarku\checkmark$\checkmark^\checkmark\\checkmarkc\checkmarki\checkmarkr\checkmarkc\checkmark$\checkmarke\checkmark$\checkmark^\checkmark\\checkmarkc\checkmarki\checkmarkr\checkmarkc\checkmark$\checkmark \checkmark$\checkmark^\checkmark\\checkmarkc\checkmarki\checkmarkr\checkmarkc\checkmark$\checkmarka\checkmark$\checkmark^\checkmark\\checkmarkc\checkmarki\checkmarkr\checkmarkc\checkmark$\checkmarkx\checkmark$\checkmark^\checkmark\\checkmarkc\checkmarki\checkmarkr\checkmarkc\checkmark$\checkmarke\checkmark$\checkmark^\checkmark\\checkmarkc\checkmarki\checkmarkr\checkmarkc\checkmark$\checkmark}\checkmark$\checkmark^\checkmark\\checkmarkc\checkmarki\checkmarkr\checkmarkc\checkmark$\checkmark
\checkmark$\checkmark^\checkmark\\checkmarkc\checkmarki\checkmarkr\checkmarkc\checkmark$\checkmark\\checkmark$\checkmark^\checkmark\\checkmarkc\checkmarki\checkmarkr\checkmarkc\checkmark$\checkmarkl\checkmark$\checkmark^\checkmark\\checkmarkc\checkmarki\checkmarkr\checkmarkc\checkmark$\checkmarka\checkmark$\checkmark^\checkmark\\checkmarkc\checkmarki\checkmarkr\checkmarkc\checkmark$\checkmarkb\checkmark$\checkmark^\checkmark\\checkmarkc\checkmarki\checkmarkr\checkmarkc\checkmark$\checkmarke\checkmark$\checkmark^\checkmark\\checkmarkc\checkmarki\checkmarkr\checkmarkc\checkmark$\checkmarkl\checkmark$\checkmark^\checkmark\\checkmarkc\checkmarki\checkmarkr\checkmarkc\checkmark$\checkmark{\checkmark$\checkmark^\checkmark\\checkmarkc\checkmarki\checkmarkr\checkmarkc\checkmark$\checkmarkt\checkmark$\checkmark^\checkmark\\checkmarkc\checkmarki\checkmarkr\checkmarkc\checkmark$\checkmarka\checkmark$\checkmark^\checkmark\\checkmarkc\checkmarki\checkmarkr\checkmarkc\checkmark$\checkmarkb\checkmark$\checkmark^\checkmark\\checkmarkc\checkmarki\checkmarkr\checkmarkc\checkmark$\checkmark:\checkmark$\checkmark^\checkmark\\checkmarkc\checkmarki\checkmarkr\checkmarkc\checkmark$\checkmarkc\checkmark$\checkmark^\checkmark\\checkmarkc\checkmarki\checkmarkr\checkmarkc\checkmark$\checkmarkh\checkmark$\checkmark^\checkmark\\checkmarkc\checkmarki\checkmarkr\checkmarkc\checkmark$\checkmarko\checkmark$\checkmark^\checkmark\\checkmarkc\checkmarki\checkmarkr\checkmarkc\checkmark$\checkmarki\checkmark$\checkmark^\checkmark\\checkmarkc\checkmarki\checkmarkr\checkmarkc\checkmark$\checkmarkx\checkmark$\checkmark^\checkmark\\checkmarkc\checkmarki\checkmarkr\checkmarkc\checkmark$\checkmark_\checkmark$\checkmark^\checkmark\\checkmarkc\checkmarki\checkmarkr\checkmarkc\checkmark$\checkmarkm\checkmark$\checkmark^\checkmark\\checkmarkc\checkmarki\checkmarkr\checkmarkc\checkmark$\checkmarko\checkmark$\checkmark^\checkmark\\checkmarkc\checkmarki\checkmarkr\checkmarkc\checkmark$\checkmarkt\checkmark$\checkmark^\checkmark\\checkmarkc\checkmarki\checkmarkr\checkmarkc\checkmark$\checkmarke\checkmark$\checkmark^\checkmark\\checkmarkc\checkmarki\checkmarkr\checkmarkc\checkmark$\checkmarku\checkmark$\checkmark^\checkmark\\checkmarkc\checkmarki\checkmarkr\checkmarkc\checkmark$\checkmarkr\checkmark$\checkmark^\checkmark\\checkmarkc\checkmarki\checkmarkr\checkmarkc\checkmark$\checkmarks\checkmark$\checkmark^\checkmark\\checkmarkc\checkmarki\checkmarkr\checkmarkc\checkmark$\checkmark}\checkmark$\checkmark^\checkmark\\checkmarkc\checkmarki\checkmarkr\checkmarkc\checkmark$\checkmark
\checkmark$\checkmark^\checkmark\\checkmarkc\checkmarki\checkmarkr\checkmarkc\checkmark$\checkmark\\checkmark$\checkmark^\checkmark\\checkmarkc\checkmarki\checkmarkr\checkmarkc\checkmark$\checkmarke\checkmark$\checkmark^\checkmark\\checkmarkc\checkmarki\checkmarkr\checkmarkc\checkmark$\checkmarkn\checkmark$\checkmark^\checkmark\\checkmarkc\checkmarki\checkmarkr\checkmarkc\checkmark$\checkmarkd\checkmark$\checkmark^\checkmark\\checkmarkc\checkmarki\checkmarkr\checkmarkc\checkmark$\checkmark{\checkmark$\checkmark^\checkmark\\checkmarkc\checkmarki\checkmarkr\checkmarkc\checkmark$\checkmarkt\checkmark$\checkmark^\checkmark\\checkmarkc\checkmarki\checkmarkr\checkmarkc\checkmark$\checkmarka\checkmark$\checkmark^\checkmark\\checkmarkc\checkmarki\checkmarkr\checkmarkc\checkmark$\checkmarkb\checkmark$\checkmark^\checkmark\\checkmarkc\checkmarki\checkmarkr\checkmarkc\checkmark$\checkmarkl\checkmark$\checkmark^\checkmark\\checkmarkc\checkmarki\checkmarkr\checkmarkc\checkmark$\checkmarke\checkmark$\checkmark^\checkmark\\checkmarkc\checkmarki\checkmarkr\checkmarkc\checkmark$\checkmark}\checkmark$\checkmark^\checkmark\\checkmarkc\checkmarki\checkmarkr\checkmarkc\checkmark$\checkmark
\checkmark$\checkmark^\checkmark\\checkmarkc\checkmarki\checkmarkr\checkmarkc\checkmark$\checkmark
\checkmark$\checkmark^\checkmark\\checkmarkc\checkmarki\checkmarkr\checkmarkc\checkmark$\checkmark%\checkmark$\checkmark^\checkmark\\checkmarkc\checkmarki\checkmarkr\checkmarkc\checkmark$\checkmark \checkmark$\checkmark^\checkmark\\checkmarkc\checkmarki\checkmarkr\checkmarkc\checkmark$\checkmark=\checkmark$\checkmark^\checkmark\\checkmarkc\checkmarki\checkmarkr\checkmarkc\checkmark$\checkmark=\checkmark$\checkmark^\checkmark\\checkmarkc\checkmarki\checkmarkr\checkmarkc\checkmark$\checkmark=\checkmark$\checkmark^\checkmark\\checkmarkc\checkmarki\checkmarkr\checkmarkc\checkmark$\checkmark=\checkmark$\checkmark^\checkmark\\checkmarkc\checkmarki\checkmarkr\checkmarkc\checkmark$\checkmark=\checkmark$\checkmark^\checkmark\\checkmarkc\checkmarki\checkmarkr\checkmarkc\checkmark$\checkmark=\checkmark$\checkmark^\checkmark\\checkmarkc\checkmarki\checkmarkr\checkmarkc\checkmark$\checkmark=\checkmark$\checkmark^\checkmark\\checkmarkc\checkmarki\checkmarkr\checkmarkc\checkmark$\checkmark=\checkmark$\checkmark^\checkmark\\checkmarkc\checkmarki\checkmarkr\checkmarkc\checkmark$\checkmark=\checkmark$\checkmark^\checkmark\\checkmarkc\checkmarki\checkmarkr\checkmarkc\checkmark$\checkmark=\checkmark$\checkmark^\checkmark\\checkmarkc\checkmarki\checkmarkr\checkmarkc\checkmark$\checkmark=\checkmark$\checkmark^\checkmark\\checkmarkc\checkmarki\checkmarkr\checkmarkc\checkmark$\checkmark=\checkmark$\checkmark^\checkmark\\checkmarkc\checkmarki\checkmarkr\checkmarkc\checkmark$\checkmark=\checkmark$\checkmark^\checkmark\\checkmarkc\checkmarki\checkmarkr\checkmarkc\checkmark$\checkmark=\checkmark$\checkmark^\checkmark\\checkmarkc\checkmarki\checkmarkr\checkmarkc\checkmark$\checkmark=\checkmark$\checkmark^\checkmark\\checkmarkc\checkmarki\checkmarkr\checkmarkc\checkmark$\checkmark=\checkmark$\checkmark^\checkmark\\checkmarkc\checkmarki\checkmarkr\checkmarkc\checkmark$\checkmark=\checkmark$\checkmark^\checkmark\\checkmarkc\checkmarki\checkmarkr\checkmarkc\checkmark$\checkmark=\checkmark$\checkmark^\checkmark\\checkmarkc\checkmarki\checkmarkr\checkmarkc\checkmark$\checkmark=\checkmark$\checkmark^\checkmark\\checkmarkc\checkmarki\checkmarkr\checkmarkc\checkmark$\checkmark=\checkmark$\checkmark^\checkmark\\checkmarkc\checkmarki\checkmarkr\checkmarkc\checkmark$\checkmark=\checkmark$\checkmark^\checkmark\\checkmarkc\checkmarki\checkmarkr\checkmarkc\checkmark$\checkmark=\checkmark$\checkmark^\checkmark\\checkmarkc\checkmarki\checkmarkr\checkmarkc\checkmark$\checkmark=\checkmark$\checkmark^\checkmark\\checkmarkc\checkmarki\checkmarkr\checkmarkc\checkmark$\checkmark=\checkmark$\checkmark^\checkmark\\checkmarkc\checkmarki\checkmarkr\checkmarkc\checkmark$\checkmark=\checkmark$\checkmark^\checkmark\\checkmarkc\checkmarki\checkmarkr\checkmarkc\checkmark$\checkmark=\checkmark$\checkmark^\checkmark\\checkmarkc\checkmarki\checkmarkr\checkmarkc\checkmark$\checkmark=\checkmark$\checkmark^\checkmark\\checkmarkc\checkmarki\checkmarkr\checkmarkc\checkmark$\checkmark=\checkmark$\checkmark^\checkmark\\checkmarkc\checkmarki\checkmarkr\checkmarkc\checkmark$\checkmark=\checkmark$\checkmark^\checkmark\\checkmarkc\checkmarki\checkmarkr\checkmarkc\checkmark$\checkmark=\checkmark$\checkmark^\checkmark\\checkmarkc\checkmarki\checkmarkr\checkmarkc\checkmark$\checkmark=\checkmark$\checkmark^\checkmark\\checkmarkc\checkmarki\checkmarkr\checkmarkc\checkmark$\checkmark=\checkmark$\checkmark^\checkmark\\checkmarkc\checkmarki\checkmarkr\checkmarkc\checkmark$\checkmark=\checkmark$\checkmark^\checkmark\\checkmarkc\checkmarki\checkmarkr\checkmarkc\checkmark$\checkmark=\checkmark$\checkmark^\checkmark\\checkmarkc\checkmarki\checkmarkr\checkmarkc\checkmark$\checkmark=\checkmark$\checkmark^\checkmark\\checkmarkc\checkmarki\checkmarkr\checkmarkc\checkmark$\checkmark=\checkmark$\checkmark^\checkmark\\checkmarkc\checkmarki\checkmarkr\checkmarkc\checkmark$\checkmark=\checkmark$\checkmark^\checkmark\\checkmarkc\checkmarki\checkmarkr\checkmarkc\checkmark$\checkmark=\checkmark$\checkmark^\checkmark\\checkmarkc\checkmarki\checkmarkr\checkmarkc\checkmark$\checkmark=\checkmark$\checkmark^\checkmark\\checkmarkc\checkmarki\checkmarkr\checkmarkc\checkmark$\checkmark=\checkmark$\checkmark^\checkmark\\checkmarkc\checkmarki\checkmarkr\checkmarkc\checkmark$\checkmark=\checkmark$\checkmark^\checkmark\\checkmarkc\checkmarki\checkmarkr\checkmarkc\checkmark$\checkmark=\checkmark$\checkmark^\checkmark\\checkmarkc\checkmarki\checkmarkr\checkmarkc\checkmark$\checkmark=\checkmark$\checkmark^\checkmark\\checkmarkc\checkmarki\checkmarkr\checkmarkc\checkmark$\checkmark=\checkmark$\checkmark^\checkmark\\checkmarkc\checkmarki\checkmarkr\checkmarkc\checkmark$\checkmark=\checkmark$\checkmark^\checkmark\\checkmarkc\checkmarki\checkmarkr\checkmarkc\checkmark$\checkmark=\checkmark$\checkmark^\checkmark\\checkmarkc\checkmarki\checkmarkr\checkmarkc\checkmark$\checkmark=\checkmark$\checkmark^\checkmark\\checkmarkc\checkmarki\checkmarkr\checkmarkc\checkmark$\checkmark=\checkmark$\checkmark^\checkmark\\checkmarkc\checkmarki\checkmarkr\checkmarkc\checkmark$\checkmark=\checkmark$\checkmark^\checkmark\\checkmarkc\checkmarki\checkmarkr\checkmarkc\checkmark$\checkmark=\checkmark$\checkmark^\checkmark\\checkmarkc\checkmarki\checkmarkr\checkmarkc\checkmark$\checkmark=\checkmark$\checkmark^\checkmark\\checkmarkc\checkmarki\checkmarkr\checkmarkc\checkmark$\checkmark=\checkmark$\checkmark^\checkmark\\checkmarkc\checkmarki\checkmarkr\checkmarkc\checkmark$\checkmark=\checkmark$\checkmark^\checkmark\\checkmarkc\checkmarki\checkmarkr\checkmarkc\checkmark$\checkmark=\checkmark$\checkmark^\checkmark\\checkmarkc\checkmarki\checkmarkr\checkmarkc\checkmark$\checkmark=\checkmark$\checkmark^\checkmark\\checkmarkc\checkmarki\checkmarkr\checkmarkc\checkmark$\checkmark=\checkmark$\checkmark^\checkmark\\checkmarkc\checkmarki\checkmarkr\checkmarkc\checkmark$\checkmark=\checkmark$\checkmark^\checkmark\\checkmarkc\checkmarki\checkmarkr\checkmarkc\checkmark$\checkmark=\checkmark$\checkmark^\checkmark\\checkmarkc\checkmarki\checkmarkr\checkmarkc\checkmark$\checkmark=\checkmark$\checkmark^\checkmark\\checkmarkc\checkmarki\checkmarkr\checkmarkc\checkmark$\checkmark=\checkmark$\checkmark^\checkmark\\checkmarkc\checkmarki\checkmarkr\checkmarkc\checkmark$\checkmark=\checkmark$\checkmark^\checkmark\\checkmarkc\checkmarki\checkmarkr\checkmarkc\checkmark$\checkmark=\checkmark$\checkmark^\checkmark\\checkmarkc\checkmarki\checkmarkr\checkmarkc\checkmark$\checkmark
\checkmark$\checkmark^\checkmark\\checkmarkc\checkmarki\checkmarkr\checkmarkc\checkmark$\checkmark\\checkmark$\checkmark^\checkmark\\checkmarkc\checkmarki\checkmarkr\checkmarkc\checkmark$\checkmarks\checkmark$\checkmark^\checkmark\\checkmarkc\checkmarki\checkmarkr\checkmarkc\checkmark$\checkmarke\checkmark$\checkmark^\checkmark\\checkmarkc\checkmarki\checkmarkr\checkmarkc\checkmark$\checkmarkc\checkmark$\checkmark^\checkmark\\checkmarkc\checkmarki\checkmarkr\checkmarkc\checkmark$\checkmarkt\checkmark$\checkmark^\checkmark\\checkmarkc\checkmarki\checkmarkr\checkmarkc\checkmark$\checkmarki\checkmark$\checkmark^\checkmark\\checkmarkc\checkmarki\checkmarkr\checkmarkc\checkmark$\checkmarko\checkmark$\checkmark^\checkmark\\checkmarkc\checkmarki\checkmarkr\checkmarkc\checkmark$\checkmarkn\checkmark$\checkmark^\checkmark\\checkmarkc\checkmarki\checkmarkr\checkmarkc\checkmark$\checkmark{\checkmark$\checkmark^\checkmark\\checkmarkc\checkmarki\checkmarkr\checkmarkc\checkmark$\checkmarkD\checkmark$\checkmark^\checkmark\\checkmarkc\checkmarki\checkmarkr\checkmarkc\checkmark$\checkmarki\checkmark$\checkmark^\checkmark\\checkmarkc\checkmarki\checkmarkr\checkmarkc\checkmark$\checkmarkm\checkmark$\checkmark^\checkmark\\checkmarkc\checkmarki\checkmarkr\checkmarkc\checkmark$\checkmarke\checkmark$\checkmark^\checkmark\\checkmarkc\checkmarki\checkmarkr\checkmarkc\checkmark$\checkmarkn\checkmark$\checkmark^\checkmark\\checkmarkc\checkmarki\checkmarkr\checkmarkc\checkmark$\checkmarks\checkmark$\checkmark^\checkmark\\checkmarkc\checkmarki\checkmarkr\checkmarkc\checkmark$\checkmarki\checkmark$\checkmark^\checkmark\\checkmarkc\checkmarki\checkmarkr\checkmarkc\checkmark$\checkmarko\checkmark$\checkmark^\checkmark\\checkmarkc\checkmarki\checkmarkr\checkmarkc\checkmark$\checkmarkn\checkmark$\checkmark^\checkmark\\checkmarkc\checkmarki\checkmarkr\checkmarkc\checkmark$\checkmarkn\checkmark$\checkmark^\checkmark\\checkmarkc\checkmarki\checkmarkr\checkmarkc\checkmark$\checkmarke\checkmark$\checkmark^\checkmark\\checkmarkc\checkmarki\checkmarkr\checkmarkc\checkmark$\checkmarkm\checkmark$\checkmark^\checkmark\\checkmarkc\checkmarki\checkmarkr\checkmarkc\checkmark$\checkmarke\checkmark$\checkmark^\checkmark\\checkmarkc\checkmarki\checkmarkr\checkmarkc\checkmark$\checkmarkn\checkmark$\checkmark^\checkmark\\checkmarkc\checkmarki\checkmarkr\checkmarkc\checkmark$\checkmarkt\checkmark$\checkmark^\checkmark\\checkmarkc\checkmarki\checkmarkr\checkmarkc\checkmark$\checkmark \checkmark$\checkmark^\checkmark\\checkmarkc\checkmarki\checkmarkr\checkmarkc\checkmark$\checkmarkd\checkmark$\checkmark^\checkmark\\checkmarkc\checkmarki\checkmarkr\checkmarkc\checkmark$\checkmarke\checkmark$\checkmark^\checkmark\\checkmarkc\checkmarki\checkmarkr\checkmarkc\checkmark$\checkmarks\checkmark$\checkmark^\checkmark\\checkmarkc\checkmarki\checkmarkr\checkmarkc\checkmark$\checkmark \checkmark$\checkmark^\checkmark\\checkmarkc\checkmarki\checkmarkr\checkmarkc\checkmark$\checkmarkO\checkmark$\checkmark^\checkmark\\checkmarkc\checkmarki\checkmarkr\checkmarkc\checkmark$\checkmarkr\checkmark$\checkmark^\checkmark\\checkmarkc\checkmarki\checkmarkr\checkmarkc\checkmark$\checkmarkg\checkmark$\checkmark^\checkmark\\checkmarkc\checkmarki\checkmarkr\checkmarkc\checkmark$\checkmarka\checkmark$\checkmark^\checkmark\\checkmarkc\checkmarki\checkmarkr\checkmarkc\checkmark$\checkmarkn\checkmark$\checkmark^\checkmark\\checkmarkc\checkmarki\checkmarkr\checkmarkc\checkmark$\checkmarke\checkmark$\checkmark^\checkmark\\checkmarkc\checkmarki\checkmarkr\checkmarkc\checkmark$\checkmarks\checkmark$\checkmark^\checkmark\\checkmarkc\checkmarki\checkmarkr\checkmarkc\checkmark$\checkmark \checkmark$\checkmark^\checkmark\\checkmarkc\checkmarki\checkmarkr\checkmarkc\checkmark$\checkmarkd\checkmark$\checkmark^\checkmark\\checkmarkc\checkmarki\checkmarkr\checkmarkc\checkmark$\checkmarke\checkmark$\checkmark^\checkmark\\checkmarkc\checkmarki\checkmarkr\checkmarkc\checkmark$\checkmark \checkmark$\checkmark^\checkmark\\checkmarkc\checkmarki\checkmarkr\checkmarkc\checkmark$\checkmarkT\checkmark$\checkmark^\checkmark\\checkmarkc\checkmarki\checkmarkr\checkmarkc\checkmark$\checkmarkr\checkmark$\checkmark^\checkmark\\checkmarkc\checkmarki\checkmarkr\checkmarkc\checkmark$\checkmarka\checkmark$\checkmark^\checkmark\\checkmarkc\checkmarki\checkmarkr\checkmarkc\checkmark$\checkmarkn\checkmark$\checkmark^\checkmark\\checkmarkc\checkmarki\checkmarkr\checkmarkc\checkmark$\checkmarks\checkmark$\checkmark^\checkmark\\checkmarkc\checkmarki\checkmarkr\checkmarkc\checkmark$\checkmarkl\checkmark$\checkmark^\checkmark\\checkmarkc\checkmarki\checkmarkr\checkmarkc\checkmark$\checkmarka\checkmark$\checkmark^\checkmark\\checkmarkc\checkmarki\checkmarkr\checkmarkc\checkmark$\checkmarkt\checkmark$\checkmark^\checkmark\\checkmarkc\checkmarki\checkmarkr\checkmarkc\checkmark$\checkmarki\checkmark$\checkmark^\checkmark\\checkmarkc\checkmarki\checkmarkr\checkmarkc\checkmark$\checkmarko\checkmark$\checkmark^\checkmark\\checkmarkc\checkmarki\checkmarkr\checkmarkc\checkmark$\checkmarkn\checkmark$\checkmark^\checkmark\\checkmarkc\checkmarki\checkmarkr\checkmarkc\checkmark$\checkmark \checkmark$\checkmark^\checkmark\\checkmarkc\checkmarki\checkmarkr\checkmarkc\checkmark$\checkmark(\checkmark$\checkmark^\checkmark\\checkmarkc\checkmarki\checkmarkr\checkmarkc\checkmark$\checkmarkA\checkmark$\checkmark^\checkmark\\checkmarkc\checkmarki\checkmarkr\checkmarkc\checkmark$\checkmarkx\checkmark$\checkmark^\checkmark\\checkmarkc\checkmarki\checkmarkr\checkmarkc\checkmark$\checkmarke\checkmark$\checkmark^\checkmark\\checkmarkc\checkmarki\checkmarkr\checkmarkc\checkmark$\checkmarks\checkmark$\checkmark^\checkmark\\checkmarkc\checkmarki\checkmarkr\checkmarkc\checkmark$\checkmark \checkmark$\checkmark^\checkmark\\checkmarkc\checkmarki\checkmarkr\checkmarkc\checkmark$\checkmarkX\checkmark$\checkmark^\checkmark\\checkmarkc\checkmarki\checkmarkr\checkmarkc\checkmark$\checkmark,\checkmark$\checkmark^\checkmark\\checkmarkc\checkmarki\checkmarkr\checkmarkc\checkmark$\checkmark \checkmark$\checkmark^\checkmark\\checkmarkc\checkmarki\checkmarkr\checkmarkc\checkmark$\checkmarkY\checkmark$\checkmark^\checkmark\\checkmarkc\checkmarki\checkmarkr\checkmarkc\checkmark$\checkmark,\checkmark$\checkmark^\checkmark\\checkmarkc\checkmarki\checkmarkr\checkmarkc\checkmark$\checkmark \checkmark$\checkmark^\checkmark\\checkmarkc\checkmarki\checkmarkr\checkmarkc\checkmark$\checkmarkZ\checkmark$\checkmark^\checkmark\\checkmarkc\checkmarki\checkmarkr\checkmarkc\checkmark$\checkmark)\checkmark$\checkmark^\checkmark\\checkmarkc\checkmarki\checkmarkr\checkmarkc\checkmark$\checkmark}\checkmark$\checkmark^\checkmark\\checkmarkc\checkmarki\checkmarkr\checkmarkc\checkmark$\checkmark
\checkmark$\checkmark^\checkmark\\checkmarkc\checkmarki\checkmarkr\checkmarkc\checkmark$\checkmark
\checkmark$\checkmark^\checkmark\\checkmarkc\checkmarki\checkmarkr\checkmarkc\checkmark$\checkmark\\checkmark$\checkmark^\checkmark\\checkmarkc\checkmarki\checkmarkr\checkmarkc\checkmark$\checkmarks\checkmark$\checkmark^\checkmark\\checkmarkc\checkmarki\checkmarkr\checkmarkc\checkmark$\checkmarku\checkmark$\checkmark^\checkmark\\checkmarkc\checkmarki\checkmarkr\checkmarkc\checkmark$\checkmarkb\checkmark$\checkmark^\checkmark\\checkmarkc\checkmarki\checkmarkr\checkmarkc\checkmark$\checkmarks\checkmark$\checkmark^\checkmark\\checkmarkc\checkmarki\checkmarkr\checkmarkc\checkmark$\checkmarke\checkmark$\checkmark^\checkmark\\checkmarkc\checkmarki\checkmarkr\checkmarkc\checkmark$\checkmarkc\checkmark$\checkmark^\checkmark\\checkmarkc\checkmarki\checkmarkr\checkmarkc\checkmark$\checkmarkt\checkmark$\checkmark^\checkmark\\checkmarkc\checkmarki\checkmarkr\checkmarkc\checkmark$\checkmarki\checkmark$\checkmark^\checkmark\\checkmarkc\checkmarki\checkmarkr\checkmarkc\checkmark$\checkmarko\checkmark$\checkmark^\checkmark\\checkmarkc\checkmarki\checkmarkr\checkmarkc\checkmark$\checkmarkn\checkmark$\checkmark^\checkmark\\checkmarkc\checkmarki\checkmarkr\checkmarkc\checkmark$\checkmark{\checkmark$\checkmark^\checkmark\\checkmarkc\checkmarki\checkmarkr\checkmarkc\checkmark$\checkmarkI\checkmark$\checkmark^\checkmark\\checkmarkc\checkmarki\checkmarkr\checkmarkc\checkmark$\checkmarkn\checkmark$\checkmark^\checkmark\\checkmarkc\checkmarki\checkmarkr\checkmarkc\checkmark$\checkmarkt\checkmark$\checkmark^\checkmark\\checkmarkc\checkmarki\checkmarkr\checkmarkc\checkmark$\checkmarkr\checkmark$\checkmark^\checkmark\\checkmarkc\checkmarki\checkmarkr\checkmarkc\checkmark$\checkmarko\checkmark$\checkmark^\checkmark\\checkmarkc\checkmarki\checkmarkr\checkmarkc\checkmark$\checkmarkd\checkmark$\checkmark^\checkmark\\checkmarkc\checkmarki\checkmarkr\checkmarkc\checkmark$\checkmarku\checkmark$\checkmark^\checkmark\\checkmarkc\checkmarki\checkmarkr\checkmarkc\checkmark$\checkmarkc\checkmark$\checkmark^\checkmark\\checkmarkc\checkmarki\checkmarkr\checkmarkc\checkmark$\checkmarkt\checkmark$\checkmark^\checkmark\\checkmarkc\checkmarki\checkmarkr\checkmarkc\checkmark$\checkmarki\checkmark$\checkmark^\checkmark\\checkmarkc\checkmarki\checkmarkr\checkmarkc\checkmark$\checkmarko\checkmark$\checkmark^\checkmark\\checkmarkc\checkmarki\checkmarkr\checkmarkc\checkmark$\checkmarkn\checkmark$\checkmark^\checkmark\\checkmarkc\checkmarki\checkmarkr\checkmarkc\checkmark$\checkmark}\checkmark$\checkmark^\checkmark\\checkmarkc\checkmarki\checkmarkr\checkmarkc\checkmark$\checkmark
\checkmark$\checkmark^\checkmark\\checkmarkc\checkmarki\checkmarkr\checkmarkc\checkmark$\checkmarkL\checkmark$\checkmark^\checkmark\\checkmarkc\checkmarki\checkmarkr\checkmarkc\checkmark$\checkmarke\checkmark$\checkmark^\checkmark\\checkmarkc\checkmarki\checkmarkr\checkmarkc\checkmark$\checkmark \checkmark$\checkmark^\checkmark\\checkmarkc\checkmarki\checkmarkr\checkmarkc\checkmark$\checkmarkd\checkmark$\checkmark^\checkmark\\checkmarkc\checkmarki\checkmarkr\checkmarkc\checkmark$\checkmarki\checkmark$\checkmark^\checkmark\\checkmarkc\checkmarki\checkmarkr\checkmarkc\checkmark$\checkmarkm\checkmark$\checkmark^\checkmark\\checkmarkc\checkmarki\checkmarkr\checkmarkc\checkmark$\checkmarke\checkmark$\checkmark^\checkmark\\checkmarkc\checkmarki\checkmarkr\checkmarkc\checkmark$\checkmarkn\checkmark$\checkmark^\checkmark\\checkmarkc\checkmarki\checkmarkr\checkmarkc\checkmark$\checkmarks\checkmark$\checkmark^\checkmark\\checkmarkc\checkmarki\checkmarkr\checkmarkc\checkmark$\checkmarki\checkmark$\checkmark^\checkmark\\checkmarkc\checkmarki\checkmarkr\checkmarkc\checkmark$\checkmarko\checkmark$\checkmark^\checkmark\\checkmarkc\checkmarki\checkmarkr\checkmarkc\checkmark$\checkmarkn\checkmark$\checkmark^\checkmark\\checkmarkc\checkmarki\checkmarkr\checkmarkc\checkmark$\checkmarkn\checkmark$\checkmark^\checkmark\\checkmarkc\checkmarki\checkmarkr\checkmarkc\checkmark$\checkmarke\checkmark$\checkmark^\checkmark\\checkmarkc\checkmarki\checkmarkr\checkmarkc\checkmark$\checkmarkm\checkmark$\checkmark^\checkmark\\checkmarkc\checkmarki\checkmarkr\checkmarkc\checkmark$\checkmarke\checkmark$\checkmark^\checkmark\\checkmarkc\checkmarki\checkmarkr\checkmarkc\checkmark$\checkmarkn\checkmark$\checkmark^\checkmark\\checkmarkc\checkmarki\checkmarkr\checkmarkc\checkmark$\checkmarkt\checkmark$\checkmark^\checkmark\\checkmarkc\checkmarki\checkmarkr\checkmarkc\checkmark$\checkmark \checkmark$\checkmark^\checkmark\\checkmarkc\checkmarki\checkmarkr\checkmarkc\checkmark$\checkmarkd\checkmark$\checkmark^\checkmark\\checkmarkc\checkmarki\checkmarkr\checkmarkc\checkmark$\checkmarke\checkmark$\checkmark^\checkmark\\checkmarkc\checkmarki\checkmarkr\checkmarkc\checkmark$\checkmark \checkmark$\checkmark^\checkmark\\checkmarkc\checkmarki\checkmarkr\checkmarkc\checkmark$\checkmarkl\checkmark$\checkmark^\checkmark\\checkmarkc\checkmarki\checkmarkr\checkmarkc\checkmark$\checkmarka\checkmark$\checkmark^\checkmark\\checkmarkc\checkmarki\checkmarkr\checkmarkc\checkmark$\checkmark \checkmark$\checkmark^\checkmark\\checkmarkc\checkmarki\checkmarkr\checkmarkc\checkmark$\checkmarkv\checkmark$\checkmark^\checkmark\\checkmarkc\checkmarki\checkmarkr\checkmarkc\checkmark$\checkmarki\checkmark$\checkmark^\checkmark\\checkmarkc\checkmarki\checkmarkr\checkmarkc\checkmark$\checkmarks\checkmark$\checkmark^\checkmark\\checkmarkc\checkmarki\checkmarkr\checkmarkc\checkmark$\checkmark \checkmark$\checkmark^\checkmark\\checkmarkc\checkmarki\checkmarkr\checkmarkc\checkmark$\checkmarkà\checkmark$\checkmark^\checkmark\\checkmarkc\checkmarki\checkmarkr\checkmarkc\checkmark$\checkmark \checkmark$\checkmark^\checkmark\\checkmarkc\checkmarki\checkmarkr\checkmarkc\checkmark$\checkmarkb\checkmark$\checkmark^\checkmark\\checkmarkc\checkmarki\checkmarkr\checkmarkc\checkmark$\checkmarki\checkmark$\checkmark^\checkmark\\checkmarkc\checkmarki\checkmarkr\checkmarkc\checkmark$\checkmarkl\checkmark$\checkmark^\checkmark\\checkmarkc\checkmarki\checkmarkr\checkmarkc\checkmark$\checkmarkl\checkmark$\checkmark^\checkmark\\checkmarkc\checkmarki\checkmarkr\checkmarkc\checkmark$\checkmarke\checkmark$\checkmark^\checkmark\\checkmarkc\checkmarki\checkmarkr\checkmarkc\checkmark$\checkmarks\checkmark$\checkmark^\checkmark\\checkmarkc\checkmarki\checkmarkr\checkmarkc\checkmark$\checkmark \checkmark$\checkmark^\checkmark\\checkmarkc\checkmarki\checkmarkr\checkmarkc\checkmark$\checkmarkS\checkmark$\checkmark^\checkmark\\checkmarkc\checkmarki\checkmarkr\checkmarkc\checkmark$\checkmarkF\checkmark$\checkmark^\checkmark\\checkmarkc\checkmarki\checkmarkr\checkmarkc\checkmark$\checkmarkU\checkmark$\checkmark^\checkmark\\checkmarkc\checkmarki\checkmarkr\checkmarkc\checkmark$\checkmark1\checkmark$\checkmark^\checkmark\\checkmarkc\checkmarki\checkmarkr\checkmarkc\checkmark$\checkmark6\checkmark$\checkmark^\checkmark\\checkmarkc\checkmarki\checkmarkr\checkmarkc\checkmark$\checkmark0\checkmark$\checkmark^\checkmark\\checkmarkc\checkmarki\checkmarkr\checkmarkc\checkmark$\checkmark5\checkmark$\checkmark^\checkmark\\checkmarkc\checkmarki\checkmarkr\checkmarkc\checkmark$\checkmark \checkmark$\checkmark^\checkmark\\checkmarkc\checkmarki\checkmarkr\checkmarkc\checkmark$\checkmarks\checkmark$\checkmark^\checkmark\\checkmarkc\checkmarki\checkmarkr\checkmarkc\checkmark$\checkmark'\checkmark$\checkmark^\checkmark\\checkmarkc\checkmarki\checkmarkr\checkmarkc\checkmark$\checkmarke\checkmark$\checkmark^\checkmark\\checkmarkc\checkmarki\checkmarkr\checkmarkc\checkmark$\checkmarkf\checkmark$\checkmark^\checkmark\\checkmarkc\checkmarki\checkmarkr\checkmarkc\checkmark$\checkmarkf\checkmark$\checkmark^\checkmark\\checkmarkc\checkmarki\checkmarkr\checkmarkc\checkmark$\checkmarke\checkmark$\checkmark^\checkmark\\checkmarkc\checkmarki\checkmarkr\checkmarkc\checkmark$\checkmarkc\checkmark$\checkmark^\checkmark\\checkmarkc\checkmarki\checkmarkr\checkmarkc\checkmark$\checkmarkt\checkmark$\checkmark^\checkmark\\checkmarkc\checkmarki\checkmarkr\checkmarkc\checkmark$\checkmarku\checkmark$\checkmark^\checkmark\\checkmarkc\checkmarki\checkmarkr\checkmarkc\checkmark$\checkmarke\checkmark$\checkmark^\checkmark\\checkmarkc\checkmarki\checkmarkr\checkmarkc\checkmark$\checkmark \checkmark$\checkmark^\checkmark\\checkmarkc\checkmarki\checkmarkr\checkmarkc\checkmark$\checkmarkp\checkmark$\checkmark^\checkmark\\checkmarkc\checkmarki\checkmarkr\checkmarkc\checkmark$\checkmarko\checkmark$\checkmark^\checkmark\\checkmarkc\checkmarki\checkmarkr\checkmarkc\checkmark$\checkmarku\checkmark$\checkmark^\checkmark\\checkmarkc\checkmarki\checkmarkr\checkmarkc\checkmark$\checkmarkr\checkmark$\checkmark^\checkmark\\checkmarkc\checkmarki\checkmarkr\checkmarkc\checkmark$\checkmark \checkmark$\checkmark^\checkmark\\checkmarkc\checkmarki\checkmarkr\checkmarkc\checkmark$\checkmarkl\checkmark$\checkmark^\checkmark\\checkmarkc\checkmarki\checkmarkr\checkmarkc\checkmark$\checkmark'\checkmark$\checkmark^\checkmark\\checkmarkc\checkmarki\checkmarkr\checkmarkc\checkmark$\checkmarka\checkmark$\checkmark^\checkmark\\checkmarkc\checkmarki\checkmarkr\checkmarkc\checkmark$\checkmarkx\checkmark$\checkmark^\checkmark\\checkmarkc\checkmarki\checkmarkr\checkmarkc\checkmark$\checkmarke\checkmark$\checkmark^\checkmark\\checkmarkc\checkmarki\checkmarkr\checkmarkc\checkmark$\checkmark \checkmark$\checkmark^\checkmark\\checkmarkc\checkmarki\checkmarkr\checkmarkc\checkmark$\checkmarkX\checkmark$\checkmark^\checkmark\\checkmarkc\checkmarki\checkmarkr\checkmarkc\checkmark$\checkmark \checkmark$\checkmark^\checkmark\\checkmarkc\checkmarki\checkmarkr\checkmarkc\checkmark$\checkmark(\checkmark$\checkmark^\checkmark\\checkmarkc\checkmarki\checkmarkr\checkmarkc\checkmark$\checkmarkl\checkmark$\checkmark^\checkmark\\checkmarkc\checkmarki\checkmarkr\checkmarkc\checkmark$\checkmarke\checkmark$\checkmark^\checkmark\\checkmarkc\checkmarki\checkmarkr\checkmarkc\checkmark$\checkmark \checkmark$\checkmark^\checkmark\\checkmarkc\checkmarki\checkmarkr\checkmarkc\checkmark$\checkmarkp\checkmark$\checkmark^\checkmark\\checkmarkc\checkmarki\checkmarkr\checkmarkc\checkmark$\checkmarkl\checkmark$\checkmark^\checkmark\\checkmarkc\checkmarki\checkmarkr\checkmarkc\checkmark$\checkmarku\checkmark$\checkmark^\checkmark\\checkmarkc\checkmarki\checkmarkr\checkmarkc\checkmark$\checkmarks\checkmark$\checkmark^\checkmark\\checkmarkc\checkmarki\checkmarkr\checkmarkc\checkmark$\checkmark \checkmark$\checkmark^\checkmark\\checkmarkc\checkmarki\checkmarkr\checkmarkc\checkmark$\checkmarkc\checkmark$\checkmark^\checkmark\\checkmarkc\checkmarki\checkmarkr\checkmarkc\checkmark$\checkmarkh\checkmark$\checkmark^\checkmark\\checkmarkc\checkmarki\checkmarkr\checkmarkc\checkmark$\checkmarka\checkmark$\checkmark^\checkmark\\checkmarkc\checkmarki\checkmarkr\checkmarkc\checkmark$\checkmarkr\checkmark$\checkmark^\checkmark\\checkmarkc\checkmarki\checkmarkr\checkmarkc\checkmark$\checkmarkg\checkmark$\checkmark^\checkmark\\checkmarkc\checkmarki\checkmarkr\checkmarkc\checkmark$\checkmarké\checkmark$\checkmark^\checkmark\\checkmarkc\checkmarki\checkmarkr\checkmarkc\checkmark$\checkmark,\checkmark$\checkmark^\checkmark\\checkmarkc\checkmarki\checkmarkr\checkmarkc\checkmark$\checkmark \checkmark$\checkmark^\checkmark\\checkmarkc\checkmarki\checkmarkr\checkmarkc\checkmark$\checkmarkp\checkmark$\checkmark^\checkmark\\checkmarkc\checkmarki\checkmarkr\checkmarkc\checkmark$\checkmarko\checkmark$\checkmark^\checkmark\\checkmarkc\checkmarki\checkmarkr\checkmarkc\checkmark$\checkmarkr\checkmark$\checkmark^\checkmark\\checkmarkc\checkmarki\checkmarkr\checkmarkc\checkmark$\checkmarkt\checkmark$\checkmark^\checkmark\\checkmarkc\checkmarki\checkmarkr\checkmarkc\checkmark$\checkmarka\checkmark$\checkmark^\checkmark\\checkmarkc\checkmarki\checkmarkr\checkmarkc\checkmark$\checkmarkn\checkmark$\checkmark^\checkmark\\checkmarkc\checkmarki\checkmarkr\checkmarkc\checkmark$\checkmarkt\checkmark$\checkmark^\checkmark\\checkmarkc\checkmarki\checkmarkr\checkmarkc\checkmark$\checkmark \checkmark$\checkmark^\checkmark\\checkmarkc\checkmarki\checkmarkr\checkmarkc\checkmark$\checkmarkl\checkmark$\checkmark^\checkmark\\checkmarkc\checkmarki\checkmarkr\checkmarkc\checkmark$\checkmarka\checkmark$\checkmark^\checkmark\\checkmarkc\checkmarki\checkmarkr\checkmarkc\checkmark$\checkmark \checkmark$\checkmark^\checkmark\\checkmarkc\checkmarki\checkmarkr\checkmarkc\checkmark$\checkmarkb\checkmark$\checkmark^\checkmark\\checkmarkc\checkmarki\checkmarkr\checkmarkc\checkmark$\checkmarkr\checkmark$\checkmark^\checkmark\\checkmarkc\checkmarki\checkmarkr\checkmarkc\checkmark$\checkmarko\checkmark$\checkmark^\checkmark\\checkmarkc\checkmarki\checkmarkr\checkmarkc\checkmark$\checkmarkc\checkmark$\checkmark^\checkmark\\checkmarkc\checkmarki\checkmarkr\checkmarkc\checkmark$\checkmarkh\checkmark$\checkmark^\checkmark\\checkmarkc\checkmarki\checkmarkr\checkmarkc\checkmark$\checkmarke\checkmark$\checkmark^\checkmark\\checkmarkc\checkmarki\checkmarkr\checkmarkc\checkmark$\checkmark)\checkmark$\checkmark^\checkmark\\checkmarkc\checkmarki\checkmarkr\checkmarkc\checkmark$\checkmark,\checkmark$\checkmark^\checkmark\\checkmarkc\checkmarki\checkmarkr\checkmarkc\checkmark$\checkmark \checkmark$\checkmark^\checkmark\\checkmarkc\checkmarki\checkmarkr\checkmarkc\checkmark$\checkmarkp\checkmark$\checkmark^\checkmark\\checkmarkc\checkmarki\checkmarkr\checkmarkc\checkmark$\checkmarko\checkmark$\checkmark^\checkmark\\checkmarkc\checkmarki\checkmarkr\checkmarkc\checkmark$\checkmarku\checkmark$\checkmark^\checkmark\\checkmarkc\checkmarki\checkmarkr\checkmarkc\checkmark$\checkmarkr\checkmark$\checkmark^\checkmark\\checkmarkc\checkmarki\checkmarkr\checkmarkc\checkmark$\checkmark \checkmark$\checkmark^\checkmark\\checkmarkc\checkmarki\checkmarkr\checkmarkc\checkmark$\checkmarku\checkmark$\checkmark^\checkmark\\checkmarkc\checkmarki\checkmarkr\checkmarkc\checkmark$\checkmarkn\checkmark$\checkmark^\checkmark\\checkmarkc\checkmarki\checkmarkr\checkmarkc\checkmark$\checkmarke\checkmark$\checkmark^\checkmark\\checkmarkc\checkmarki\checkmarkr\checkmarkc\checkmark$\checkmark \checkmark$\checkmark^\checkmark\\checkmarkc\checkmarki\checkmarkr\checkmarkc\checkmark$\checkmarkl\checkmark$\checkmark^\checkmark\\checkmarkc\checkmarki\checkmarkr\checkmarkc\checkmark$\checkmarko\checkmark$\checkmark^\checkmark\\checkmarkc\checkmarki\checkmarkr\checkmarkc\checkmark$\checkmarkn\checkmark$\checkmark^\checkmark\\checkmarkc\checkmarki\checkmarkr\checkmarkc\checkmark$\checkmarkg\checkmark$\checkmark^\checkmark\\checkmarkc\checkmarki\checkmarkr\checkmarkc\checkmark$\checkmarku\checkmark$\checkmark^\checkmark\\checkmarkc\checkmarki\checkmarkr\checkmarkc\checkmark$\checkmarke\checkmark$\checkmark^\checkmark\\checkmarkc\checkmarki\checkmarkr\checkmarkc\checkmark$\checkmarku\checkmark$\checkmark^\checkmark\\checkmarkc\checkmarki\checkmarkr\checkmarkc\checkmark$\checkmarkr\checkmark$\checkmark^\checkmark\\checkmarkc\checkmarki\checkmarkr\checkmarkc\checkmark$\checkmark \checkmark$\checkmark^\checkmark\\checkmarkc\checkmarki\checkmarkr\checkmarkc\checkmark$\checkmarkd\checkmark$\checkmark^\checkmark\\checkmarkc\checkmarki\checkmarkr\checkmarkc\checkmark$\checkmarke\checkmark$\checkmark^\checkmark\\checkmarkc\checkmarki\checkmarkr\checkmarkc\checkmark$\checkmark \checkmark$\checkmark^\checkmark\\checkmarkc\checkmarki\checkmarkr\checkmarkc\checkmark$\checkmarkv\checkmark$\checkmark^\checkmark\\checkmarkc\checkmarki\checkmarkr\checkmarkc\checkmark$\checkmarki\checkmark$\checkmark^\checkmark\\checkmarkc\checkmarki\checkmarkr\checkmarkc\checkmark$\checkmarks\checkmark$\checkmark^\checkmark\\checkmarkc\checkmarki\checkmarkr\checkmarkc\checkmark$\checkmark \checkmark$\checkmark^\checkmark\\checkmarkc\checkmarki\checkmarkr\checkmarkc\checkmark$\checkmark$\checkmark$\checkmark^\checkmark\\checkmarkc\checkmarki\checkmarkr\checkmarkc\checkmark$\checkmarkL\checkmark$\checkmark^\checkmark\\checkmarkc\checkmarki\checkmarkr\checkmarkc\checkmark$\checkmark_\checkmark$\checkmark^\checkmark\\checkmarkc\checkmarki\checkmarkr\checkmarkc\checkmark$\checkmark{\checkmark$\checkmark^\checkmark\\checkmarkc\checkmarki\checkmarkr\checkmarkc\checkmark$\checkmarkv\checkmark$\checkmark^\checkmark\\checkmarkc\checkmarki\checkmarkr\checkmarkc\checkmark$\checkmarki\checkmark$\checkmark^\checkmark\\checkmarkc\checkmarki\checkmarkr\checkmarkc\checkmark$\checkmarks\checkmark$\checkmark^\checkmark\\checkmarkc\checkmarki\checkmarkr\checkmarkc\checkmark$\checkmark}\checkmark$\checkmark^\checkmark\\checkmarkc\checkmarki\checkmarkr\checkmarkc\checkmark$\checkmark \checkmark$\checkmark^\checkmark\\checkmarkc\checkmarki\checkmarkr\checkmarkc\checkmark$\checkmark=\checkmark$\checkmark^\checkmark\\checkmarkc\checkmarki\checkmarkr\checkmarkc\checkmark$\checkmark \checkmark$\checkmark^\checkmark\\checkmarkc\checkmarki\checkmarkr\checkmarkc\checkmark$\checkmark4\checkmark$\checkmark^\checkmark\\checkmarkc\checkmarki\checkmarkr\checkmarkc\checkmark$\checkmark0\checkmark$\checkmark^\checkmark\\checkmarkc\checkmarki\checkmarkr\checkmarkc\checkmark$\checkmark0\checkmark$\checkmark^\checkmark\\checkmarkc\checkmarki\checkmarkr\checkmarkc\checkmark$\checkmark$\checkmark$\checkmark^\checkmark\\checkmarkc\checkmarki\checkmarkr\checkmarkc\checkmark$\checkmark \checkmark$\checkmark^\checkmark\\checkmarkc\checkmarki\checkmarkr\checkmarkc\checkmark$\checkmarkm\checkmark$\checkmark^\checkmark\\checkmarkc\checkmarki\checkmarkr\checkmarkc\checkmark$\checkmarkm\checkmark$\checkmark^\checkmark\\checkmarkc\checkmarki\checkmarkr\checkmarkc\checkmark$\checkmark.\checkmark$\checkmark^\checkmark\\checkmarkc\checkmarki\checkmarkr\checkmarkc\checkmark$\checkmark \checkmark$\checkmark^\checkmark\\checkmarkc\checkmarki\checkmarkr\checkmarkc\checkmark$\checkmarkL\checkmark$\checkmark^\checkmark\\checkmarkc\checkmarki\checkmarkr\checkmarkc\checkmark$\checkmarke\checkmark$\checkmark^\checkmark\\checkmarkc\checkmarki\checkmarkr\checkmarkc\checkmark$\checkmarks\checkmark$\checkmark^\checkmark\\checkmarkc\checkmarki\checkmarkr\checkmarkc\checkmark$\checkmark \checkmark$\checkmark^\checkmark\\checkmarkc\checkmarki\checkmarkr\checkmarkc\checkmark$\checkmarkh\checkmark$\checkmark^\checkmark\\checkmarkc\checkmarki\checkmarkr\checkmarkc\checkmark$\checkmarky\checkmark$\checkmark^\checkmark\\checkmarkc\checkmarki\checkmarkr\checkmarkc\checkmark$\checkmarkp\checkmark$\checkmark^\checkmark\\checkmarkc\checkmarki\checkmarkr\checkmarkc\checkmark$\checkmarko\checkmark$\checkmark^\checkmark\\checkmarkc\checkmarki\checkmarkr\checkmarkc\checkmark$\checkmarkt\checkmark$\checkmark^\checkmark\\checkmarkc\checkmarki\checkmarkr\checkmarkc\checkmark$\checkmarkh\checkmark$\checkmark^\checkmark\\checkmarkc\checkmarki\checkmarkr\checkmarkc\checkmark$\checkmarkè\checkmark$\checkmark^\checkmark\\checkmarkc\checkmarki\checkmarkr\checkmarkc\checkmark$\checkmarks\checkmark$\checkmark^\checkmark\\checkmarkc\checkmarki\checkmarkr\checkmarkc\checkmark$\checkmarke\checkmark$\checkmark^\checkmark\\checkmarkc\checkmarki\checkmarkr\checkmarkc\checkmark$\checkmarks\checkmark$\checkmark^\checkmark\\checkmarkc\checkmarki\checkmarkr\checkmarkc\checkmark$\checkmark \checkmark$\checkmark^\checkmark\\checkmarkc\checkmarki\checkmarkr\checkmarkc\checkmark$\checkmarks\checkmark$\checkmark^\checkmark\\checkmarkc\checkmarki\checkmarkr\checkmarkc\checkmark$\checkmarko\checkmark$\checkmark^\checkmark\\checkmarkc\checkmarki\checkmarkr\checkmarkc\checkmark$\checkmarkn\checkmark$\checkmark^\checkmark\\checkmarkc\checkmarki\checkmarkr\checkmarkc\checkmark$\checkmarkt\checkmark$\checkmark^\checkmark\\checkmarkc\checkmarki\checkmarkr\checkmarkc\checkmark$\checkmark \checkmark$\checkmark^\checkmark\\checkmarkc\checkmarki\checkmarkr\checkmarkc\checkmark$\checkmark:\checkmark$\checkmark^\checkmark\\checkmarkc\checkmarki\checkmarkr\checkmarkc\checkmark$\checkmark
\checkmark$\checkmark^\checkmark\\checkmarkc\checkmarki\checkmarkr\checkmarkc\checkmark$\checkmark\\checkmark$\checkmark^\checkmark\\checkmarkc\checkmarki\checkmarkr\checkmarkc\checkmark$\checkmarkb\checkmark$\checkmark^\checkmark\\checkmarkc\checkmarki\checkmarkr\checkmarkc\checkmark$\checkmarke\checkmark$\checkmark^\checkmark\\checkmarkc\checkmarki\checkmarkr\checkmarkc\checkmark$\checkmarkg\checkmark$\checkmark^\checkmark\\checkmarkc\checkmarki\checkmarkr\checkmarkc\checkmark$\checkmarki\checkmark$\checkmark^\checkmark\\checkmarkc\checkmarki\checkmarkr\checkmarkc\checkmark$\checkmarkn\checkmark$\checkmark^\checkmark\\checkmarkc\checkmarki\checkmarkr\checkmarkc\checkmark$\checkmark{\checkmark$\checkmark^\checkmark\\checkmarkc\checkmarki\checkmarkr\checkmarkc\checkmark$\checkmarki\checkmark$\checkmark^\checkmark\\checkmarkc\checkmarki\checkmarkr\checkmarkc\checkmark$\checkmarkt\checkmark$\checkmark^\checkmark\\checkmarkc\checkmarki\checkmarkr\checkmarkc\checkmark$\checkmarke\checkmark$\checkmark^\checkmark\\checkmarkc\checkmarki\checkmarkr\checkmarkc\checkmark$\checkmarkm\checkmark$\checkmark^\checkmark\\checkmarkc\checkmarki\checkmarkr\checkmarkc\checkmark$\checkmarki\checkmark$\checkmark^\checkmark\\checkmarkc\checkmarki\checkmarkr\checkmarkc\checkmark$\checkmarkz\checkmark$\checkmark^\checkmark\\checkmarkc\checkmarki\checkmarkr\checkmarkc\checkmark$\checkmarke\checkmark$\checkmark^\checkmark\\checkmarkc\checkmarki\checkmarkr\checkmarkc\checkmark$\checkmark}\checkmark$\checkmark^\checkmark\\checkmarkc\checkmarki\checkmarkr\checkmarkc\checkmark$\checkmark
\checkmark$\checkmark^\checkmark\\checkmarkc\checkmarki\checkmarkr\checkmarkc\checkmark$\checkmark \checkmark$\checkmark^\checkmark\\checkmarkc\checkmarki\checkmarkr\checkmarkc\checkmark$\checkmark \checkmark$\checkmark^\checkmark\\checkmarkc\checkmarki\checkmarkr\checkmarkc\checkmark$\checkmark \checkmark$\checkmark^\checkmark\\checkmarkc\checkmarki\checkmarkr\checkmarkc\checkmark$\checkmark \checkmark$\checkmark^\checkmark\\checkmarkc\checkmarki\checkmarkr\checkmarkc\checkmark$\checkmark\\checkmark$\checkmark^\checkmark\\checkmarkc\checkmarki\checkmarkr\checkmarkc\checkmark$\checkmarki\checkmark$\checkmark^\checkmark\\checkmarkc\checkmarki\checkmarkr\checkmarkc\checkmark$\checkmarkt\checkmark$\checkmark^\checkmark\\checkmarkc\checkmarki\checkmarkr\checkmarkc\checkmark$\checkmarke\checkmark$\checkmark^\checkmark\\checkmarkc\checkmarki\checkmarkr\checkmarkc\checkmark$\checkmarkm\checkmark$\checkmark^\checkmark\\checkmarkc\checkmarki\checkmarkr\checkmarkc\checkmark$\checkmark \checkmark$\checkmark^\checkmark\\checkmarkc\checkmarki\checkmarkr\checkmarkc\checkmark$\checkmarkL\checkmark$\checkmark^\checkmark\\checkmarkc\checkmarki\checkmarkr\checkmarkc\checkmark$\checkmarke\checkmark$\checkmark^\checkmark\\checkmarkc\checkmarki\checkmarkr\checkmarkc\checkmark$\checkmark \checkmark$\checkmark^\checkmark\\checkmarkc\checkmarki\checkmarkr\checkmarkc\checkmark$\checkmarkc\checkmark$\checkmark^\checkmark\\checkmarkc\checkmarki\checkmarkr\checkmarkc\checkmark$\checkmarko\checkmark$\checkmark^\checkmark\\checkmarkc\checkmarki\checkmarkr\checkmarkc\checkmark$\checkmarkn\checkmark$\checkmark^\checkmark\\checkmarkc\checkmarki\checkmarkr\checkmarkc\checkmark$\checkmarkt\checkmark$\checkmark^\checkmark\\checkmarkc\checkmarki\checkmarkr\checkmarkc\checkmark$\checkmarka\checkmark$\checkmark^\checkmark\\checkmarkc\checkmarki\checkmarkr\checkmarkc\checkmark$\checkmarkc\checkmark$\checkmark^\checkmark\\checkmarkc\checkmarki\checkmarkr\checkmarkc\checkmark$\checkmarkt\checkmark$\checkmark^\checkmark\\checkmarkc\checkmarki\checkmarkr\checkmarkc\checkmark$\checkmark \checkmark$\checkmark^\checkmark\\checkmarkc\checkmarki\checkmarkr\checkmarkc\checkmark$\checkmarke\checkmark$\checkmark^\checkmark\\checkmarkc\checkmarki\checkmarkr\checkmarkc\checkmark$\checkmarkn\checkmark$\checkmark^\checkmark\\checkmarkc\checkmarki\checkmarkr\checkmarkc\checkmark$\checkmarkt\checkmark$\checkmark^\checkmark\\checkmarkc\checkmarki\checkmarkr\checkmarkc\checkmark$\checkmarkr\checkmark$\checkmark^\checkmark\\checkmarkc\checkmarki\checkmarkr\checkmarkc\checkmark$\checkmarke\checkmark$\checkmark^\checkmark\\checkmarkc\checkmarki\checkmarkr\checkmarkc\checkmark$\checkmark \checkmark$\checkmark^\checkmark\\checkmarkc\checkmarki\checkmarkr\checkmarkc\checkmark$\checkmarkn\checkmark$\checkmark^\checkmark\\checkmarkc\checkmarki\checkmarkr\checkmarkc\checkmark$\checkmarko\checkmark$\checkmark^\checkmark\\checkmarkc\checkmarki\checkmarkr\checkmarkc\checkmark$\checkmarki\checkmark$\checkmark^\checkmark\\checkmarkc\checkmarki\checkmarkr\checkmarkc\checkmark$\checkmarkx\checkmark$\checkmark^\checkmark\\checkmarkc\checkmarki\checkmarkr\checkmarkc\checkmark$\checkmark \checkmark$\checkmark^\checkmark\\checkmarkc\checkmarki\checkmarkr\checkmarkc\checkmark$\checkmarke\checkmark$\checkmark^\checkmark\\checkmarkc\checkmarki\checkmarkr\checkmarkc\checkmark$\checkmarkt\checkmark$\checkmark^\checkmark\\checkmarkc\checkmarki\checkmarkr\checkmarkc\checkmark$\checkmark \checkmark$\checkmark^\checkmark\\checkmarkc\checkmarki\checkmarkr\checkmarkc\checkmark$\checkmarkv\checkmark$\checkmark^\checkmark\\checkmarkc\checkmarki\checkmarkr\checkmarkc\checkmark$\checkmarki\checkmark$\checkmark^\checkmark\\checkmarkc\checkmarki\checkmarkr\checkmarkc\checkmark$\checkmarks\checkmark$\checkmark^\checkmark\\checkmarkc\checkmarki\checkmarkr\checkmarkc\checkmark$\checkmark \checkmark$\checkmark^\checkmark\\checkmarkc\checkmarki\checkmarkr\checkmarkc\checkmark$\checkmarke\checkmark$\checkmark^\checkmark\\checkmarkc\checkmarki\checkmarkr\checkmarkc\checkmark$\checkmarks\checkmark$\checkmark^\checkmark\\checkmarkc\checkmarki\checkmarkr\checkmarkc\checkmark$\checkmarkt\checkmark$\checkmark^\checkmark\\checkmarkc\checkmarki\checkmarkr\checkmarkc\checkmark$\checkmark \checkmark$\checkmark^\checkmark\\checkmarkc\checkmarki\checkmarkr\checkmarkc\checkmark$\checkmarks\checkmark$\checkmark^\checkmark\\checkmarkc\checkmarki\checkmarkr\checkmarkc\checkmark$\checkmarku\checkmark$\checkmark^\checkmark\\checkmarkc\checkmarki\checkmarkr\checkmarkc\checkmark$\checkmarkp\checkmark$\checkmark^\checkmark\\checkmarkc\checkmarki\checkmarkr\checkmarkc\checkmark$\checkmarkp\checkmark$\checkmark^\checkmark\\checkmarkc\checkmarki\checkmarkr\checkmarkc\checkmark$\checkmarko\checkmark$\checkmark^\checkmark\\checkmarkc\checkmarki\checkmarkr\checkmarkc\checkmark$\checkmarks\checkmark$\checkmark^\checkmark\\checkmarkc\checkmarki\checkmarkr\checkmarkc\checkmark$\checkmarké\checkmark$\checkmark^\checkmark\\checkmarkc\checkmarki\checkmarkr\checkmarkc\checkmark$\checkmark \checkmark$\checkmark^\checkmark\\checkmarkc\checkmarki\checkmarkr\checkmarkc\checkmark$\checkmarks\checkmark$\checkmark^\checkmark\\checkmarkc\checkmarki\checkmarkr\checkmarkc\checkmark$\checkmarka\checkmark$\checkmark^\checkmark\\checkmarkc\checkmarki\checkmarkr\checkmarkc\checkmark$\checkmarkn\checkmark$\checkmark^\checkmark\\checkmarkc\checkmarki\checkmarkr\checkmarkc\checkmark$\checkmarks\checkmark$\checkmark^\checkmark\\checkmarkc\checkmarki\checkmarkr\checkmarkc\checkmark$\checkmark \checkmark$\checkmark^\checkmark\\checkmarkc\checkmarki\checkmarkr\checkmarkc\checkmark$\checkmarkf\checkmark$\checkmark^\checkmark\\checkmarkc\checkmarki\checkmarkr\checkmarkc\checkmark$\checkmarkr\checkmark$\checkmark^\checkmark\\checkmarkc\checkmarki\checkmarkr\checkmarkc\checkmark$\checkmarko\checkmark$\checkmark^\checkmark\\checkmarkc\checkmarki\checkmarkr\checkmarkc\checkmark$\checkmarkt\checkmark$\checkmark^\checkmark\\checkmarkc\checkmarki\checkmarkr\checkmarkc\checkmark$\checkmarkt\checkmark$\checkmark^\checkmark\\checkmarkc\checkmarki\checkmarkr\checkmarkc\checkmark$\checkmarke\checkmark$\checkmark^\checkmark\\checkmarkc\checkmarki\checkmarkr\checkmarkc\checkmark$\checkmarkm\checkmark$\checkmark^\checkmark\\checkmarkc\checkmarki\checkmarkr\checkmarkc\checkmark$\checkmarke\checkmark$\checkmark^\checkmark\\checkmarkc\checkmarki\checkmarkr\checkmarkc\checkmark$\checkmarkn\checkmark$\checkmark^\checkmark\\checkmarkc\checkmarki\checkmarkr\checkmarkc\checkmark$\checkmarkt\checkmark$\checkmark^\checkmark\\checkmarkc\checkmarki\checkmarkr\checkmarkc\checkmark$\checkmark \checkmark$\checkmark^\checkmark\\checkmarkc\checkmarki\checkmarkr\checkmarkc\checkmark$\checkmark(\checkmark$\checkmark^\checkmark\\checkmarkc\checkmarki\checkmarkr\checkmarkc\checkmark$\checkmarkr\checkmark$\checkmark^\checkmark\\checkmarkc\checkmarki\checkmarkr\checkmarkc\checkmark$\checkmarke\checkmark$\checkmark^\checkmark\\checkmarkc\checkmarki\checkmarkr\checkmarkc\checkmark$\checkmarkn\checkmark$\checkmark^\checkmark\\checkmarkc\checkmarki\checkmarkr\checkmarkc\checkmark$\checkmarkd\checkmark$\checkmark^\checkmark\\checkmarkc\checkmarki\checkmarkr\checkmarkc\checkmark$\checkmarke\checkmark$\checkmark^\checkmark\\checkmarkc\checkmarki\checkmarkr\checkmarkc\checkmark$\checkmarkm\checkmark$\checkmark^\checkmark\\checkmarkc\checkmarki\checkmarkr\checkmarkc\checkmark$\checkmarke\checkmark$\checkmark^\checkmark\\checkmarkc\checkmarki\checkmarkr\checkmarkc\checkmark$\checkmarkn\checkmark$\checkmark^\checkmark\\checkmarkc\checkmarki\checkmarkr\checkmarkc\checkmark$\checkmarkt\checkmark$\checkmark^\checkmark\\checkmarkc\checkmarki\checkmarkr\checkmarkc\checkmark$\checkmark \checkmark$\checkmark^\checkmark\\checkmarkc\checkmarki\checkmarkr\checkmarkc\checkmark$\checkmark$\checkmark$\checkmark^\checkmark\\checkmarkc\checkmarki\checkmarkr\checkmarkc\checkmark$\checkmark\\checkmark$\checkmark^\checkmark\\checkmarkc\checkmarki\checkmarkr\checkmarkc\checkmark$\checkmarke\checkmark$\checkmark^\checkmark\\checkmarkc\checkmarki\checkmarkr\checkmarkc\checkmark$\checkmarkt\checkmark$\checkmark^\checkmark\\checkmarkc\checkmarki\checkmarkr\checkmarkc\checkmark$\checkmarka\checkmark$\checkmark^\checkmark\\checkmarkc\checkmarki\checkmarkr\checkmarkc\checkmark$\checkmark \checkmark$\checkmark^\checkmark\\checkmarkc\checkmarki\checkmarkr\checkmarkc\checkmark$\checkmark=\checkmark$\checkmark^\checkmark\\checkmarkc\checkmarki\checkmarkr\checkmarkc\checkmark$\checkmark \checkmark$\checkmark^\checkmark\\checkmarkc\checkmarki\checkmarkr\checkmarkc\checkmark$\checkmark0\checkmark$\checkmark^\checkmark\\checkmarkc\checkmarki\checkmarkr\checkmarkc\checkmark$\checkmark{\checkmark$\checkmark^\checkmark\\checkmarkc\checkmarki\checkmarkr\checkmarkc\checkmark$\checkmark,\checkmark$\checkmark^\checkmark\\checkmarkc\checkmarki\checkmarkr\checkmarkc\checkmark$\checkmark}\checkmark$\checkmark^\checkmark\\checkmarkc\checkmarki\checkmarkr\checkmarkc\checkmark$\checkmark9\checkmark$\checkmark^\checkmark\\checkmarkc\checkmarki\checkmarkr\checkmarkc\checkmark$\checkmark0\checkmark$\checkmark^\checkmark\\checkmarkc\checkmarki\checkmarkr\checkmarkc\checkmark$\checkmark$\checkmark$\checkmark^\checkmark\\checkmarkc\checkmarki\checkmarkr\checkmarkc\checkmark$\checkmark)\checkmark$\checkmark^\checkmark\\checkmarkc\checkmarki\checkmarkr\checkmarkc\checkmark$\checkmark \checkmark$\checkmark^\checkmark\\checkmarkc\checkmarki\checkmarkr\checkmarkc\checkmark$\checkmark;\checkmark$\checkmark^\checkmark\\checkmarkc\checkmarki\checkmarkr\checkmarkc\checkmark$\checkmark
\checkmark$\checkmark^\checkmark\\checkmarkc\checkmarki\checkmarkr\checkmarkc\checkmark$\checkmark \checkmark$\checkmark^\checkmark\\checkmarkc\checkmarki\checkmarkr\checkmarkc\checkmark$\checkmark \checkmark$\checkmark^\checkmark\\checkmarkc\checkmarki\checkmarkr\checkmarkc\checkmark$\checkmark \checkmark$\checkmark^\checkmark\\checkmarkc\checkmarki\checkmarkr\checkmarkc\checkmark$\checkmark \checkmark$\checkmark^\checkmark\\checkmarkc\checkmarki\checkmarkr\checkmarkc\checkmark$\checkmark\\checkmark$\checkmark^\checkmark\\checkmarkc\checkmarki\checkmarkr\checkmarkc\checkmark$\checkmarki\checkmark$\checkmark^\checkmark\\checkmarkc\checkmarki\checkmarkr\checkmarkc\checkmark$\checkmarkt\checkmark$\checkmark^\checkmark\\checkmarkc\checkmarki\checkmarkr\checkmarkc\checkmark$\checkmarke\checkmark$\checkmark^\checkmark\\checkmarkc\checkmarki\checkmarkr\checkmarkc\checkmark$\checkmarkm\checkmark$\checkmark^\checkmark\\checkmarkc\checkmarki\checkmarkr\checkmarkc\checkmark$\checkmark \checkmark$\checkmark^\checkmark\\checkmarkc\checkmarki\checkmarkr\checkmarkc\checkmark$\checkmarkL\checkmark$\checkmark^\checkmark\\checkmarkc\checkmarki\checkmarkr\checkmarkc\checkmark$\checkmarka\checkmark$\checkmark^\checkmark\\checkmarkc\checkmarki\checkmarkr\checkmarkc\checkmark$\checkmark \checkmark$\checkmark^\checkmark\\checkmarkc\checkmarki\checkmarkr\checkmarkc\checkmark$\checkmarkm\checkmark$\checkmark^\checkmark\\checkmarkc\checkmarki\checkmarkr\checkmarkc\checkmark$\checkmarka\checkmark$\checkmark^\checkmark\\checkmarkc\checkmarki\checkmarkr\checkmarkc\checkmark$\checkmarks\checkmark$\checkmark^\checkmark\\checkmarkc\checkmarki\checkmarkr\checkmarkc\checkmark$\checkmarks\checkmark$\checkmark^\checkmark\\checkmarkc\checkmarki\checkmarkr\checkmarkc\checkmark$\checkmarke\checkmark$\checkmark^\checkmark\\checkmarkc\checkmarki\checkmarkr\checkmarkc\checkmark$\checkmark \checkmark$\checkmark^\checkmark\\checkmarkc\checkmarki\checkmarkr\checkmarkc\checkmark$\checkmarkm\checkmark$\checkmark^\checkmark\\checkmarkc\checkmarki\checkmarkr\checkmarkc\checkmark$\checkmarko\checkmark$\checkmark^\checkmark\\checkmarkc\checkmarki\checkmarkr\checkmarkc\checkmark$\checkmarkb\checkmark$\checkmark^\checkmark\\checkmarkc\checkmarki\checkmarkr\checkmarkc\checkmark$\checkmarki\checkmark$\checkmark^\checkmark\\checkmarkc\checkmarki\checkmarkr\checkmarkc\checkmark$\checkmarkl\checkmark$\checkmark^\checkmark\\checkmarkc\checkmarki\checkmarkr\checkmarkc\checkmark$\checkmarke\checkmark$\checkmark^\checkmark\\checkmarkc\checkmarki\checkmarkr\checkmarkc\checkmark$\checkmark \checkmark$\checkmark^\checkmark\\checkmarkc\checkmarki\checkmarkr\checkmarkc\checkmark$\checkmark(\checkmark$\checkmark^\checkmark\\checkmarkc\checkmarki\checkmarkr\checkmarkc\checkmark$\checkmarkb\checkmark$\checkmark^\checkmark\\checkmarkc\checkmarki\checkmarkr\checkmarkc\checkmark$\checkmarkr\checkmark$\checkmark^\checkmark\\checkmarkc\checkmarki\checkmarkr\checkmarkc\checkmark$\checkmarko\checkmark$\checkmark^\checkmark\\checkmarkc\checkmarki\checkmarkr\checkmarkc\checkmark$\checkmarkc\checkmark$\checkmark^\checkmark\\checkmarkc\checkmarki\checkmarkr\checkmarkc\checkmark$\checkmarkh\checkmark$\checkmark^\checkmark\\checkmarkc\checkmarki\checkmarkr\checkmarkc\checkmark$\checkmarke\checkmark$\checkmark^\checkmark\\checkmarkc\checkmarki\checkmarkr\checkmarkc\checkmark$\checkmark \checkmark$\checkmark^\checkmark\\checkmarkc\checkmarki\checkmarkr\checkmarkc\checkmark$\checkmark+\checkmark$\checkmark^\checkmark\\checkmarkc\checkmarki\checkmarkr\checkmarkc\checkmark$\checkmark \checkmark$\checkmark^\checkmark\\checkmarkc\checkmarki\checkmarkr\checkmarkc\checkmark$\checkmarkc\checkmark$\checkmark^\checkmark\\checkmarkc\checkmarki\checkmarkr\checkmarkc\checkmark$\checkmarkh\checkmark$\checkmark^\checkmark\\checkmarkc\checkmarki\checkmarkr\checkmarkc\checkmark$\checkmarka\checkmark$\checkmark^\checkmark\\checkmarkc\checkmarki\checkmarkr\checkmarkc\checkmark$\checkmarkr\checkmark$\checkmark^\checkmark\\checkmarkc\checkmarki\checkmarkr\checkmarkc\checkmark$\checkmarki\checkmark$\checkmark^\checkmark\\checkmarkc\checkmarki\checkmarkr\checkmarkc\checkmark$\checkmarko\checkmark$\checkmark^\checkmark\\checkmarkc\checkmarki\checkmarkr\checkmarkc\checkmark$\checkmarkt\checkmark$\checkmark^\checkmark\\checkmarkc\checkmarki\checkmarkr\checkmarkc\checkmark$\checkmark \checkmark$\checkmark^\checkmark\\checkmarkc\checkmarki\checkmarkr\checkmarkc\checkmark$\checkmarkX\checkmark$\checkmark^\checkmark\\checkmarkc\checkmarki\checkmarkr\checkmarkc\checkmark$\checkmark)\checkmark$\checkmark^\checkmark\\checkmarkc\checkmarki\checkmarkr\checkmarkc\checkmark$\checkmark \checkmark$\checkmark^\checkmark\\checkmarkc\checkmarki\checkmarkr\checkmarkc\checkmark$\checkmarke\checkmark$\checkmark^\checkmark\\checkmarkc\checkmarki\checkmarkr\checkmarkc\checkmark$\checkmarks\checkmark$\checkmark^\checkmark\\checkmarkc\checkmarki\checkmarkr\checkmarkc\checkmark$\checkmarkt\checkmark$\checkmark^\checkmark\\checkmarkc\checkmarki\checkmarkr\checkmarkc\checkmark$\checkmark \checkmark$\checkmark^\checkmark\\checkmarkc\checkmarki\checkmarkr\checkmarkc\checkmark$\checkmark$\checkmark$\checkmark^\checkmark\\checkmarkc\checkmarki\checkmarkr\checkmarkc\checkmark$\checkmarkM\checkmark$\checkmark^\checkmark\\checkmarkc\checkmarki\checkmarkr\checkmarkc\checkmark$\checkmark \checkmark$\checkmark^\checkmark\\checkmarkc\checkmarki\checkmarkr\checkmarkc\checkmark$\checkmark=\checkmark$\checkmark^\checkmark\\checkmarkc\checkmarki\checkmarkr\checkmarkc\checkmark$\checkmark \checkmark$\checkmark^\checkmark\\checkmarkc\checkmarki\checkmarkr\checkmarkc\checkmark$\checkmark1\checkmark$\checkmark^\checkmark\\checkmarkc\checkmarki\checkmarkr\checkmarkc\checkmark$\checkmark0\checkmark$\checkmark^\checkmark\\checkmarkc\checkmarki\checkmarkr\checkmarkc\checkmark$\checkmark$\checkmark$\checkmark^\checkmark\\checkmarkc\checkmarki\checkmarkr\checkmarkc\checkmark$\checkmark \checkmark$\checkmark^\checkmark\\checkmarkc\checkmarki\checkmarkr\checkmarkc\checkmark$\checkmarkk\checkmark$\checkmark^\checkmark\\checkmarkc\checkmarki\checkmarkr\checkmarkc\checkmark$\checkmarkg\checkmark$\checkmark^\checkmark\\checkmarkc\checkmarki\checkmarkr\checkmarkc\checkmark$\checkmark \checkmark$\checkmark^\checkmark\\checkmarkc\checkmarki\checkmarkr\checkmarkc\checkmark$\checkmark;\checkmark$\checkmark^\checkmark\\checkmarkc\checkmarki\checkmarkr\checkmarkc\checkmark$\checkmark
\checkmark$\checkmark^\checkmark\\checkmarkc\checkmarki\checkmarkr\checkmarkc\checkmark$\checkmark \checkmark$\checkmark^\checkmark\\checkmarkc\checkmarki\checkmarkr\checkmarkc\checkmark$\checkmark \checkmark$\checkmark^\checkmark\\checkmarkc\checkmarki\checkmarkr\checkmarkc\checkmark$\checkmark \checkmark$\checkmark^\checkmark\\checkmarkc\checkmarki\checkmarkr\checkmarkc\checkmark$\checkmark \checkmark$\checkmark^\checkmark\\checkmarkc\checkmarki\checkmarkr\checkmarkc\checkmark$\checkmark\\checkmark$\checkmark^\checkmark\\checkmarkc\checkmarki\checkmarkr\checkmarkc\checkmark$\checkmarki\checkmark$\checkmark^\checkmark\\checkmarkc\checkmarki\checkmarkr\checkmarkc\checkmark$\checkmarkt\checkmark$\checkmark^\checkmark\\checkmarkc\checkmarki\checkmarkr\checkmarkc\checkmark$\checkmarke\checkmark$\checkmark^\checkmark\\checkmarkc\checkmarki\checkmarkr\checkmarkc\checkmark$\checkmarkm\checkmark$\checkmark^\checkmark\\checkmarkc\checkmarki\checkmarkr\checkmarkc\checkmark$\checkmark \checkmark$\checkmark^\checkmark\\checkmarkc\checkmarki\checkmarkr\checkmarkc\checkmark$\checkmarkL\checkmark$\checkmark^\checkmark\\checkmarkc\checkmarki\checkmarkr\checkmarkc\checkmark$\checkmarka\checkmark$\checkmark^\checkmark\\checkmarkc\checkmarki\checkmarkr\checkmarkc\checkmark$\checkmark \checkmark$\checkmark^\checkmark\\checkmarkc\checkmarki\checkmarkr\checkmarkc\checkmark$\checkmarkv\checkmark$\checkmark^\checkmark\\checkmarkc\checkmarki\checkmarkr\checkmarkc\checkmark$\checkmarki\checkmark$\checkmark^\checkmark\\checkmarkc\checkmarki\checkmarkr\checkmarkc\checkmark$\checkmarks\checkmark$\checkmark^\checkmark\\checkmarkc\checkmarki\checkmarkr\checkmarkc\checkmark$\checkmark \checkmark$\checkmark^\checkmark\\checkmarkc\checkmarki\checkmarkr\checkmarkc\checkmark$\checkmarke\checkmark$\checkmark^\checkmark\\checkmarkc\checkmarki\checkmarkr\checkmarkc\checkmark$\checkmarks\checkmark$\checkmark^\checkmark\\checkmarkc\checkmarki\checkmarkr\checkmarkc\checkmark$\checkmarkt\checkmark$\checkmark^\checkmark\\checkmarkc\checkmarki\checkmarkr\checkmarkc\checkmark$\checkmark \checkmark$\checkmark^\checkmark\\checkmarkc\checkmarki\checkmarkr\checkmarkc\checkmark$\checkmarkm\checkmark$\checkmark^\checkmark\\checkmarkc\checkmarki\checkmarkr\checkmarkc\checkmark$\checkmarko\checkmark$\checkmark^\checkmark\\checkmarkc\checkmarki\checkmarkr\checkmarkc\checkmark$\checkmarkn\checkmark$\checkmark^\checkmark\\checkmarkc\checkmarki\checkmarkr\checkmarkc\checkmark$\checkmarkt\checkmark$\checkmark^\checkmark\\checkmarkc\checkmarki\checkmarkr\checkmarkc\checkmark$\checkmarké\checkmark$\checkmark^\checkmark\\checkmarkc\checkmarki\checkmarkr\checkmarkc\checkmark$\checkmarke\checkmark$\checkmark^\checkmark\\checkmarkc\checkmarki\checkmarkr\checkmarkc\checkmark$\checkmark \checkmark$\checkmark^\checkmark\\checkmarkc\checkmarki\checkmarkr\checkmarkc\checkmark$\checkmarke\checkmark$\checkmark^\checkmark\\checkmarkc\checkmarki\checkmarkr\checkmarkc\checkmark$\checkmarkn\checkmark$\checkmark^\checkmark\\checkmarkc\checkmarki\checkmarkr\checkmarkc\checkmark$\checkmark \checkmark$\checkmark^\checkmark\\checkmarkc\checkmarki\checkmarkr\checkmarkc\checkmark$\checkmarkc\checkmark$\checkmark^\checkmark\\checkmarkc\checkmarki\checkmarkr\checkmarkc\checkmark$\checkmarko\checkmark$\checkmark^\checkmark\\checkmarkc\checkmarki\checkmarkr\checkmarkc\checkmark$\checkmarkn\checkmark$\checkmark^\checkmark\\checkmarkc\checkmarki\checkmarkr\checkmarkc\checkmark$\checkmarkf\checkmark$\checkmark^\checkmark\\checkmarkc\checkmarki\checkmarkr\checkmarkc\checkmark$\checkmarki\checkmark$\checkmark^\checkmark\\checkmarkc\checkmarki\checkmarkr\checkmarkc\checkmark$\checkmarkg\checkmark$\checkmark^\checkmark\\checkmarkc\checkmarki\checkmarkr\checkmarkc\checkmark$\checkmarku\checkmark$\checkmark^\checkmark\\checkmarkc\checkmarki\checkmarkr\checkmarkc\checkmark$\checkmarkr\checkmark$\checkmark^\checkmark\\checkmarkc\checkmarki\checkmarkr\checkmarkc\checkmark$\checkmarka\checkmark$\checkmark^\checkmark\\checkmarkc\checkmarki\checkmarkr\checkmarkc\checkmark$\checkmarkt\checkmark$\checkmark^\checkmark\\checkmarkc\checkmarki\checkmarkr\checkmarkc\checkmark$\checkmarki\checkmark$\checkmark^\checkmark\\checkmarkc\checkmarki\checkmarkr\checkmarkc\checkmark$\checkmarko\checkmark$\checkmark^\checkmark\\checkmarkc\checkmarki\checkmarkr\checkmarkc\checkmark$\checkmarkn\checkmark$\checkmark^\checkmark\\checkmarkc\checkmarki\checkmarkr\checkmarkc\checkmark$\checkmark \checkmark$\checkmark^\checkmark\\checkmarkc\checkmarki\checkmarkr\checkmarkc\checkmark$\checkmark«\checkmark$\checkmark^\checkmark\\checkmarkc\checkmarki\checkmarkr\checkmarkc\checkmark$\checkmarke\checkmark$\checkmark^\checkmark\\checkmarkc\checkmarki\checkmarkr\checkmarkc\checkmark$\checkmarkn\checkmark$\checkmark^\checkmark\\checkmarkc\checkmarki\checkmarkr\checkmarkc\checkmark$\checkmarkc\checkmark$\checkmark^\checkmark\\checkmarkc\checkmarki\checkmarkr\checkmarkc\checkmark$\checkmarka\checkmark$\checkmark^\checkmark\\checkmarkc\checkmarki\checkmarkr\checkmarkc\checkmark$\checkmarks\checkmark$\checkmark^\checkmark\\checkmarkc\checkmarki\checkmarkr\checkmarkc\checkmark$\checkmarkt\checkmark$\checkmark^\checkmark\\checkmarkc\checkmarki\checkmarkr\checkmarkc\checkmark$\checkmarkr\checkmark$\checkmark^\checkmark\\checkmarkc\checkmarki\checkmarkr\checkmarkc\checkmark$\checkmarké\checkmark$\checkmark^\checkmark\\checkmarkc\checkmarki\checkmarkr\checkmarkc\checkmark$\checkmark-\checkmark$\checkmark^\checkmark\\checkmarkc\checkmarki\checkmarkr\checkmarkc\checkmark$\checkmarkl\checkmark$\checkmark^\checkmark\\checkmarkc\checkmarki\checkmarkr\checkmarkc\checkmark$\checkmarki\checkmark$\checkmark^\checkmark\\checkmarkc\checkmarki\checkmarkr\checkmarkc\checkmark$\checkmarkb\checkmark$\checkmark^\checkmark\\checkmarkc\checkmarki\checkmarkr\checkmarkc\checkmark$\checkmarkr\checkmark$\checkmark^\checkmark\\checkmarkc\checkmarki\checkmarkr\checkmarkc\checkmark$\checkmarke\checkmark$\checkmark^\checkmark\\checkmarkc\checkmarki\checkmarkr\checkmarkc\checkmark$\checkmark»\checkmark$\checkmark^\checkmark\\checkmarkc\checkmarki\checkmarkr\checkmarkc\checkmark$\checkmark \checkmark$\checkmark^\checkmark\\checkmarkc\checkmarki\checkmarkr\checkmarkc\checkmark$\checkmark(\checkmark$\checkmark^\checkmark\\checkmarkc\checkmarki\checkmarkr\checkmarkc\checkmark$\checkmark$\checkmark$\checkmark^\checkmark\\checkmarkc\checkmarki\checkmarkr\checkmarkc\checkmark$\checkmarkk\checkmark$\checkmark^\checkmark\\checkmarkc\checkmarki\checkmarkr\checkmarkc\checkmark$\checkmark_\checkmark$\checkmark^\checkmark\\checkmarkc\checkmarki\checkmarkr\checkmarkc\checkmark$\checkmarkL\checkmark$\checkmark^\checkmark\\checkmarkc\checkmarki\checkmarkr\checkmarkc\checkmark$\checkmark \checkmark$\checkmark^\checkmark\\checkmarkc\checkmarki\checkmarkr\checkmarkc\checkmark$\checkmark=\checkmark$\checkmark^\checkmark\\checkmarkc\checkmarki\checkmarkr\checkmarkc\checkmark$\checkmark \checkmark$\checkmark^\checkmark\\checkmarkc\checkmarki\checkmarkr\checkmarkc\checkmark$\checkmark0\checkmark$\checkmark^\checkmark\\checkmarkc\checkmarki\checkmarkr\checkmarkc\checkmark$\checkmark{\checkmark$\checkmark^\checkmark\\checkmarkc\checkmarki\checkmarkr\checkmarkc\checkmark$\checkmark,\checkmark$\checkmark^\checkmark\\checkmarkc\checkmarki\checkmarkr\checkmarkc\checkmark$\checkmark}\checkmark$\checkmark^\checkmark\\checkmarkc\checkmarki\checkmarkr\checkmarkc\checkmark$\checkmark7\checkmark$\checkmark^\checkmark\\checkmarkc\checkmarki\checkmarkr\checkmarkc\checkmark$\checkmark0\checkmark$\checkmark^\checkmark\\checkmarkc\checkmarki\checkmarkr\checkmarkc\checkmark$\checkmark7\checkmark$\checkmark^\checkmark\\checkmarkc\checkmarki\checkmarkr\checkmarkc\checkmark$\checkmark$\checkmark$\checkmark^\checkmark\\checkmarkc\checkmarki\checkmarkr\checkmarkc\checkmark$\checkmark)\checkmark$\checkmark^\checkmark\\checkmarkc\checkmarki\checkmarkr\checkmarkc\checkmark$\checkmark \checkmark$\checkmark^\checkmark\\checkmarkc\checkmarki\checkmarkr\checkmarkc\checkmark$\checkmarkp\checkmark$\checkmark^\checkmark\\checkmarkc\checkmarki\checkmarkr\checkmarkc\checkmark$\checkmarko\checkmark$\checkmark^\checkmark\\checkmarkc\checkmarki\checkmarkr\checkmarkc\checkmark$\checkmarku\checkmark$\checkmark^\checkmark\\checkmarkc\checkmarki\checkmarkr\checkmarkc\checkmark$\checkmarkr\checkmark$\checkmark^\checkmark\\checkmarkc\checkmarki\checkmarkr\checkmarkc\checkmark$\checkmark \checkmark$\checkmark^\checkmark\\checkmarkc\checkmarki\checkmarkr\checkmarkc\checkmark$\checkmarkl\checkmark$\checkmark^\checkmark\\checkmarkc\checkmarki\checkmarkr\checkmarkc\checkmark$\checkmarke\checkmark$\checkmark^\checkmark\\checkmarkc\checkmarki\checkmarkr\checkmarkc\checkmark$\checkmark \checkmark$\checkmark^\checkmark\\checkmarkc\checkmarki\checkmarkr\checkmarkc\checkmark$\checkmarkf\checkmark$\checkmark^\checkmark\\checkmarkc\checkmarki\checkmarkr\checkmarkc\checkmark$\checkmarkl\checkmark$\checkmark^\checkmark\\checkmarkc\checkmarki\checkmarkr\checkmarkc\checkmark$\checkmarka\checkmark$\checkmark^\checkmark\\checkmarkc\checkmarki\checkmarkr\checkmarkc\checkmark$\checkmarkm\checkmark$\checkmark^\checkmark\\checkmarkc\checkmarki\checkmarkr\checkmarkc\checkmark$\checkmarkb\checkmark$\checkmark^\checkmark\\checkmarkc\checkmarki\checkmarkr\checkmarkc\checkmark$\checkmarka\checkmark$\checkmark^\checkmark\\checkmarkc\checkmarki\checkmarkr\checkmarkc\checkmark$\checkmarkg\checkmark$\checkmark^\checkmark\\checkmarkc\checkmarki\checkmarkr\checkmarkc\checkmark$\checkmarke\checkmark$\checkmark^\checkmark\\checkmarkc\checkmarki\checkmarkr\checkmarkc\checkmark$\checkmark.\checkmark$\checkmark^\checkmark\\checkmarkc\checkmarki\checkmarkr\checkmarkc\checkmark$\checkmark
\checkmark$\checkmark^\checkmark\\checkmarkc\checkmarki\checkmarkr\checkmarkc\checkmark$\checkmark\\checkmark$\checkmark^\checkmark\\checkmarkc\checkmarki\checkmarkr\checkmarkc\checkmark$\checkmarke\checkmark$\checkmark^\checkmark\\checkmarkc\checkmarki\checkmarkr\checkmarkc\checkmark$\checkmarkn\checkmark$\checkmark^\checkmark\\checkmarkc\checkmarki\checkmarkr\checkmarkc\checkmark$\checkmarkd\checkmark$\checkmark^\checkmark\\checkmarkc\checkmarki\checkmarkr\checkmarkc\checkmark$\checkmark{\checkmark$\checkmark^\checkmark\\checkmarkc\checkmarki\checkmarkr\checkmarkc\checkmark$\checkmarki\checkmark$\checkmark^\checkmark\\checkmarkc\checkmarki\checkmarkr\checkmarkc\checkmark$\checkmarkt\checkmark$\checkmark^\checkmark\\checkmarkc\checkmarki\checkmarkr\checkmarkc\checkmark$\checkmarke\checkmark$\checkmark^\checkmark\\checkmarkc\checkmarki\checkmarkr\checkmarkc\checkmark$\checkmarkm\checkmark$\checkmark^\checkmark\\checkmarkc\checkmarki\checkmarkr\checkmarkc\checkmark$\checkmarki\checkmark$\checkmark^\checkmark\\checkmarkc\checkmarki\checkmarkr\checkmarkc\checkmark$\checkmarkz\checkmark$\checkmark^\checkmark\\checkmarkc\checkmarki\checkmarkr\checkmarkc\checkmark$\checkmarke\checkmark$\checkmark^\checkmark\\checkmarkc\checkmarki\checkmarkr\checkmarkc\checkmark$\checkmark}\checkmark$\checkmark^\checkmark\\checkmarkc\checkmarki\checkmarkr\checkmarkc\checkmark$\checkmark
\checkmark$\checkmark^\checkmark\\checkmarkc\checkmarki\checkmarkr\checkmarkc\checkmark$\checkmark
\checkmark$\checkmark^\checkmark\\checkmarkc\checkmarki\checkmarkr\checkmarkc\checkmark$\checkmark\\checkmark$\checkmark^\checkmark\\checkmarkc\checkmarki\checkmarkr\checkmarkc\checkmark$\checkmarks\checkmark$\checkmark^\checkmark\\checkmarkc\checkmarki\checkmarkr\checkmarkc\checkmark$\checkmarku\checkmark$\checkmark^\checkmark\\checkmarkc\checkmarki\checkmarkr\checkmarkc\checkmark$\checkmarkb\checkmark$\checkmark^\checkmark\\checkmarkc\checkmarki\checkmarkr\checkmarkc\checkmark$\checkmarks\checkmark$\checkmark^\checkmark\\checkmarkc\checkmarki\checkmarkr\checkmarkc\checkmark$\checkmarke\checkmark$\checkmark^\checkmark\\checkmarkc\checkmarki\checkmarkr\checkmarkc\checkmark$\checkmarkc\checkmark$\checkmark^\checkmark\\checkmarkc\checkmarki\checkmarkr\checkmarkc\checkmark$\checkmarkt\checkmark$\checkmark^\checkmark\\checkmarkc\checkmarki\checkmarkr\checkmarkc\checkmark$\checkmarki\checkmark$\checkmark^\checkmark\\checkmarkc\checkmarki\checkmarkr\checkmarkc\checkmark$\checkmarko\checkmark$\checkmark^\checkmark\\checkmarkc\checkmarki\checkmarkr\checkmarkc\checkmark$\checkmarkn\checkmark$\checkmark^\checkmark\\checkmarkc\checkmarki\checkmarkr\checkmarkc\checkmark$\checkmark{\checkmark$\checkmark^\checkmark\\checkmarkc\checkmarki\checkmarkr\checkmarkc\checkmark$\checkmarkD\checkmark$\checkmark^\checkmark\\checkmarkc\checkmarki\checkmarkr\checkmarkc\checkmark$\checkmarké\checkmark$\checkmark^\checkmark\\checkmarkc\checkmarki\checkmarkr\checkmarkc\checkmark$\checkmarkt\checkmark$\checkmark^\checkmark\\checkmarkc\checkmarki\checkmarkr\checkmarkc\checkmark$\checkmarke\checkmark$\checkmark^\checkmark\\checkmarkc\checkmarki\checkmarkr\checkmarkc\checkmark$\checkmarkr\checkmark$\checkmark^\checkmark\\checkmarkc\checkmarki\checkmarkr\checkmarkc\checkmark$\checkmarkm\checkmark$\checkmark^\checkmark\\checkmarkc\checkmarki\checkmarkr\checkmarkc\checkmark$\checkmarki\checkmark$\checkmark^\checkmark\\checkmarkc\checkmarki\checkmarkr\checkmarkc\checkmark$\checkmarkn\checkmark$\checkmark^\checkmark\\checkmarkc\checkmarki\checkmarkr\checkmarkc\checkmark$\checkmarka\checkmark$\checkmark^\checkmark\\checkmarkc\checkmarki\checkmarkr\checkmarkc\checkmark$\checkmarkt\checkmark$\checkmark^\checkmark\\checkmarkc\checkmarki\checkmarkr\checkmarkc\checkmark$\checkmarki\checkmark$\checkmark^\checkmark\\checkmarkc\checkmarki\checkmarkr\checkmarkc\checkmark$\checkmarko\checkmark$\checkmark^\checkmark\\checkmarkc\checkmarki\checkmarkr\checkmarkc\checkmark$\checkmarkn\checkmark$\checkmark^\checkmark\\checkmarkc\checkmarki\checkmarkr\checkmarkc\checkmark$\checkmark \checkmark$\checkmark^\checkmark\\checkmarkc\checkmarki\checkmarkr\checkmarkc\checkmark$\checkmarkd\checkmark$\checkmark^\checkmark\\checkmarkc\checkmarki\checkmarkr\checkmarkc\checkmark$\checkmarke\checkmark$\checkmark^\checkmark\\checkmarkc\checkmarki\checkmarkr\checkmarkc\checkmark$\checkmark \checkmark$\checkmark^\checkmark\\checkmarkc\checkmarki\checkmarkr\checkmarkc\checkmark$\checkmarkl\checkmark$\checkmark^\checkmark\\checkmarkc\checkmarki\checkmarkr\checkmarkc\checkmark$\checkmarka\checkmark$\checkmark^\checkmark\\checkmarkc\checkmarki\checkmarkr\checkmarkc\checkmark$\checkmark \checkmark$\checkmark^\checkmark\\checkmarkc\checkmarki\checkmarkr\checkmarkc\checkmark$\checkmarkf\checkmark$\checkmark^\checkmark\\checkmarkc\checkmarki\checkmarkr\checkmarkc\checkmark$\checkmarko\checkmark$\checkmark^\checkmark\\checkmarkc\checkmarki\checkmarkr\checkmarkc\checkmark$\checkmarkr\checkmark$\checkmark^\checkmark\\checkmarkc\checkmarki\checkmarkr\checkmarkc\checkmark$\checkmarkc\checkmark$\checkmark^\checkmark\\checkmarkc\checkmarki\checkmarkr\checkmarkc\checkmark$\checkmarke\checkmark$\checkmark^\checkmark\\checkmarkc\checkmarki\checkmarkr\checkmarkc\checkmark$\checkmark \checkmark$\checkmark^\checkmark\\checkmarkc\checkmarki\checkmarkr\checkmarkc\checkmark$\checkmarka\checkmark$\checkmark^\checkmark\\checkmarkc\checkmarki\checkmarkr\checkmarkc\checkmark$\checkmarkx\checkmark$\checkmark^\checkmark\\checkmarkc\checkmarki\checkmarkr\checkmarkc\checkmark$\checkmarki\checkmark$\checkmark^\checkmark\\checkmarkc\checkmarki\checkmarkr\checkmarkc\checkmark$\checkmarka\checkmark$\checkmark^\checkmark\\checkmarkc\checkmarki\checkmarkr\checkmarkc\checkmark$\checkmarkl\checkmark$\checkmark^\checkmark\\checkmarkc\checkmarki\checkmarkr\checkmarkc\checkmark$\checkmarke\checkmark$\checkmark^\checkmark\\checkmarkc\checkmarki\checkmarkr\checkmarkc\checkmark$\checkmark \checkmark$\checkmark^\checkmark\\checkmarkc\checkmarki\checkmarkr\checkmarkc\checkmark$\checkmarks\checkmark$\checkmark^\checkmark\\checkmarkc\checkmarki\checkmarkr\checkmarkc\checkmark$\checkmarku\checkmark$\checkmark^\checkmark\\checkmarkc\checkmarki\checkmarkr\checkmarkc\checkmark$\checkmarkr\checkmark$\checkmark^\checkmark\\checkmarkc\checkmarki\checkmarkr\checkmarkc\checkmark$\checkmark \checkmark$\checkmark^\checkmark\\checkmarkc\checkmarki\checkmarkr\checkmarkc\checkmark$\checkmarkl\checkmark$\checkmark^\checkmark\\checkmarkc\checkmarki\checkmarkr\checkmarkc\checkmark$\checkmarka\checkmark$\checkmark^\checkmark\\checkmarkc\checkmarki\checkmarkr\checkmarkc\checkmark$\checkmark \checkmark$\checkmark^\checkmark\\checkmarkc\checkmarki\checkmarkr\checkmarkc\checkmark$\checkmarkv\checkmark$\checkmark^\checkmark\\checkmarkc\checkmarki\checkmarkr\checkmarkc\checkmark$\checkmarki\checkmark$\checkmark^\checkmark\\checkmarkc\checkmarki\checkmarkr\checkmarkc\checkmark$\checkmarks\checkmark$\checkmark^\checkmark\\checkmarkc\checkmarki\checkmarkr\checkmarkc\checkmark$\checkmark}\checkmark$\checkmark^\checkmark\\checkmarkc\checkmarki\checkmarkr\checkmarkc\checkmark$\checkmark
\checkmark$\checkmark^\checkmark\\checkmarkc\checkmarki\checkmarkr\checkmarkc\checkmark$\checkmarkL\checkmark$\checkmark^\checkmark\\checkmarkc\checkmarki\checkmarkr\checkmarkc\checkmark$\checkmarka\checkmark$\checkmark^\checkmark\\checkmarkc\checkmarki\checkmarkr\checkmarkc\checkmark$\checkmark \checkmark$\checkmark^\checkmark\\checkmarkc\checkmarki\checkmarkr\checkmarkc\checkmark$\checkmarkf\checkmark$\checkmark^\checkmark\\checkmarkc\checkmarki\checkmarkr\checkmarkc\checkmark$\checkmarko\checkmark$\checkmark^\checkmark\\checkmarkc\checkmarki\checkmarkr\checkmarkc\checkmark$\checkmarkr\checkmark$\checkmark^\checkmark\\checkmarkc\checkmarki\checkmarkr\checkmarkc\checkmark$\checkmarkc\checkmark$\checkmark^\checkmark\\checkmarkc\checkmarki\checkmarkr\checkmarkc\checkmark$\checkmarke\checkmark$\checkmark^\checkmark\\checkmarkc\checkmarki\checkmarkr\checkmarkc\checkmark$\checkmark \checkmark$\checkmark^\checkmark\\checkmarkc\checkmarki\checkmarkr\checkmarkc\checkmark$\checkmarka\checkmark$\checkmark^\checkmark\\checkmarkc\checkmarki\checkmarkr\checkmarkc\checkmark$\checkmarkx\checkmark$\checkmark^\checkmark\\checkmarkc\checkmarki\checkmarkr\checkmarkc\checkmark$\checkmarki\checkmark$\checkmark^\checkmark\\checkmarkc\checkmarki\checkmarkr\checkmarkc\checkmark$\checkmarka\checkmark$\checkmark^\checkmark\\checkmarkc\checkmarki\checkmarkr\checkmarkc\checkmark$\checkmarkl\checkmark$\checkmark^\checkmark\\checkmarkc\checkmarki\checkmarkr\checkmarkc\checkmark$\checkmarke\checkmark$\checkmark^\checkmark\\checkmarkc\checkmarki\checkmarkr\checkmarkc\checkmark$\checkmark \checkmark$\checkmark^\checkmark\\checkmarkc\checkmarki\checkmarkr\checkmarkc\checkmark$\checkmark$\checkmark$\checkmark^\checkmark\\checkmarkc\checkmarki\checkmarkr\checkmarkc\checkmark$\checkmarkF\checkmark$\checkmark^\checkmark\\checkmarkc\checkmarki\checkmarkr\checkmarkc\checkmark$\checkmark_\checkmark$\checkmark^\checkmark\\checkmarkc\checkmarki\checkmarkr\checkmarkc\checkmark$\checkmark{\checkmark$\checkmark^\checkmark\\checkmarkc\checkmarki\checkmarkr\checkmarkc\checkmark$\checkmarka\checkmark$\checkmark^\checkmark\\checkmarkc\checkmarki\checkmarkr\checkmarkc\checkmark$\checkmarkx\checkmark$\checkmark^\checkmark\\checkmarkc\checkmarki\checkmarkr\checkmarkc\checkmark$\checkmark}\checkmark$\checkmark^\checkmark\\checkmarkc\checkmarki\checkmarkr\checkmarkc\checkmark$\checkmark$\checkmark$\checkmark^\checkmark\\checkmarkc\checkmarki\checkmarkr\checkmarkc\checkmark$\checkmark \checkmark$\checkmark^\checkmark\\checkmarkc\checkmarki\checkmarkr\checkmarkc\checkmark$\checkmarkr\checkmark$\checkmark^\checkmark\\checkmarkc\checkmarki\checkmarkr\checkmarkc\checkmark$\checkmarké\checkmark$\checkmark^\checkmark\\checkmarkc\checkmarki\checkmarkr\checkmarkc\checkmark$\checkmarks\checkmark$\checkmark^\checkmark\\checkmarkc\checkmarki\checkmarkr\checkmarkc\checkmark$\checkmarku\checkmark$\checkmark^\checkmark\\checkmarkc\checkmarki\checkmarkr\checkmarkc\checkmark$\checkmarkl\checkmark$\checkmark^\checkmark\\checkmarkc\checkmarki\checkmarkr\checkmarkc\checkmark$\checkmarkt\checkmark$\checkmark^\checkmark\\checkmarkc\checkmarki\checkmarkr\checkmarkc\checkmark$\checkmarke\checkmark$\checkmark^\checkmark\\checkmarkc\checkmarki\checkmarkr\checkmarkc\checkmark$\checkmark \checkmark$\checkmark^\checkmark\\checkmarkc\checkmarki\checkmarkr\checkmarkc\checkmark$\checkmarkd\checkmark$\checkmark^\checkmark\\checkmarkc\checkmarki\checkmarkr\checkmarkc\checkmark$\checkmarke\checkmark$\checkmark^\checkmark\\checkmarkc\checkmarki\checkmarkr\checkmarkc\checkmark$\checkmark \checkmark$\checkmark^\checkmark\\checkmarkc\checkmarki\checkmarkr\checkmarkc\checkmark$\checkmarkl\checkmark$\checkmark^\checkmark\\checkmarkc\checkmarki\checkmarkr\checkmarkc\checkmark$\checkmarka\checkmark$\checkmark^\checkmark\\checkmarkc\checkmarki\checkmarkr\checkmarkc\checkmark$\checkmark \checkmark$\checkmark^\checkmark\\checkmarkc\checkmarki\checkmarkr\checkmarkc\checkmark$\checkmarks\checkmark$\checkmark^\checkmark\\checkmarkc\checkmarki\checkmarkr\checkmarkc\checkmark$\checkmarku\checkmark$\checkmark^\checkmark\\checkmarkc\checkmarki\checkmarkr\checkmarkc\checkmark$\checkmarkp\checkmark$\checkmark^\checkmark\\checkmarkc\checkmarki\checkmarkr\checkmarkc\checkmark$\checkmarke\checkmark$\checkmark^\checkmark\\checkmarkc\checkmarki\checkmarkr\checkmarkc\checkmark$\checkmarkr\checkmark$\checkmark^\checkmark\\checkmarkc\checkmarki\checkmarkr\checkmarkc\checkmark$\checkmarkp\checkmark$\checkmark^\checkmark\\checkmarkc\checkmarki\checkmarkr\checkmarkc\checkmark$\checkmarko\checkmark$\checkmark^\checkmark\\checkmarkc\checkmarki\checkmarkr\checkmarkc\checkmark$\checkmarks\checkmark$\checkmark^\checkmark\\checkmarkc\checkmarki\checkmarkr\checkmarkc\checkmark$\checkmarki\checkmark$\checkmark^\checkmark\\checkmarkc\checkmarki\checkmarkr\checkmarkc\checkmark$\checkmarkt\checkmark$\checkmark^\checkmark\\checkmarkc\checkmarki\checkmarkr\checkmarkc\checkmark$\checkmarki\checkmark$\checkmark^\checkmark\\checkmarkc\checkmarki\checkmarkr\checkmarkc\checkmark$\checkmarko\checkmark$\checkmark^\checkmark\\checkmarkc\checkmarki\checkmarkr\checkmarkc\checkmark$\checkmarkn\checkmark$\checkmark^\checkmark\\checkmarkc\checkmarki\checkmarkr\checkmarkc\checkmark$\checkmark \checkmark$\checkmark^\checkmark\\checkmarkc\checkmarki\checkmarkr\checkmarkc\checkmark$\checkmarkd\checkmark$\checkmark^\checkmark\\checkmarkc\checkmarki\checkmarkr\checkmarkc\checkmark$\checkmarke\checkmark$\checkmark^\checkmark\\checkmarkc\checkmarki\checkmarkr\checkmarkc\checkmark$\checkmark \checkmark$\checkmark^\checkmark\\checkmarkc\checkmarki\checkmarkr\checkmarkc\checkmark$\checkmarkl\checkmark$\checkmark^\checkmark\\checkmarkc\checkmarki\checkmarkr\checkmarkc\checkmark$\checkmark'\checkmark$\checkmark^\checkmark\\checkmarkc\checkmarki\checkmarkr\checkmarkc\checkmark$\checkmarke\checkmark$\checkmark^\checkmark\\checkmarkc\checkmarki\checkmarkr\checkmarkc\checkmark$\checkmarkf\checkmark$\checkmark^\checkmark\\checkmarkc\checkmarki\checkmarkr\checkmarkc\checkmark$\checkmarkf\checkmark$\checkmark^\checkmark\\checkmarkc\checkmarki\checkmarkr\checkmarkc\checkmark$\checkmarko\checkmark$\checkmark^\checkmark\\checkmarkc\checkmarki\checkmarkr\checkmarkc\checkmark$\checkmarkr\checkmark$\checkmark^\checkmark\\checkmarkc\checkmarki\checkmarkr\checkmarkc\checkmark$\checkmarkt\checkmark$\checkmark^\checkmark\\checkmarkc\checkmarki\checkmarkr\checkmarkc\checkmark$\checkmark \checkmark$\checkmark^\checkmark\\checkmarkc\checkmarki\checkmarkr\checkmarkc\checkmark$\checkmarkd\checkmark$\checkmark^\checkmark\\checkmarkc\checkmarki\checkmarkr\checkmarkc\checkmark$\checkmark'\checkmark$\checkmark^\checkmark\\checkmarkc\checkmarki\checkmarkr\checkmarkc\checkmark$\checkmarku\checkmark$\checkmark^\checkmark\\checkmarkc\checkmarki\checkmarkr\checkmarkc\checkmark$\checkmarks\checkmark$\checkmark^\checkmark\\checkmarkc\checkmarki\checkmarkr\checkmarkc\checkmark$\checkmarki\checkmark$\checkmark^\checkmark\\checkmarkc\checkmarki\checkmarkr\checkmarkc\checkmark$\checkmarkn\checkmark$\checkmark^\checkmark\\checkmarkc\checkmarki\checkmarkr\checkmarkc\checkmark$\checkmarka\checkmark$\checkmark^\checkmark\\checkmarkc\checkmarki\checkmarkr\checkmarkc\checkmark$\checkmarkg\checkmark$\checkmark^\checkmark\\checkmarkc\checkmarki\checkmarkr\checkmarkc\checkmark$\checkmarke\checkmark$\checkmark^\checkmark\\checkmarkc\checkmarki\checkmarkr\checkmarkc\checkmark$\checkmark \checkmark$\checkmark^\checkmark\\checkmarkc\checkmarki\checkmarkr\checkmarkc\checkmark$\checkmarke\checkmark$\checkmark^\checkmark\\checkmarkc\checkmarki\checkmarkr\checkmarkc\checkmark$\checkmarkt\checkmark$\checkmark^\checkmark\\checkmarkc\checkmarki\checkmarkr\checkmarkc\checkmark$\checkmark \checkmark$\checkmark^\checkmark\\checkmarkc\checkmarki\checkmarkr\checkmarkc\checkmark$\checkmarkd\checkmark$\checkmark^\checkmark\\checkmarkc\checkmarki\checkmarkr\checkmarkc\checkmark$\checkmarke\checkmark$\checkmark^\checkmark\\checkmarkc\checkmarki\checkmarkr\checkmarkc\checkmark$\checkmarks\checkmark$\checkmark^\checkmark\\checkmarkc\checkmarki\checkmarkr\checkmarkc\checkmark$\checkmark \checkmark$\checkmark^\checkmark\\checkmarkc\checkmarki\checkmarkr\checkmarkc\checkmark$\checkmarkf\checkmark$\checkmark^\checkmark\\checkmarkc\checkmarki\checkmarkr\checkmarkc\checkmark$\checkmarko\checkmark$\checkmark^\checkmark\\checkmarkc\checkmarki\checkmarkr\checkmarkc\checkmark$\checkmarkr\checkmark$\checkmark^\checkmark\\checkmarkc\checkmarki\checkmarkr\checkmarkc\checkmark$\checkmarkc\checkmark$\checkmark^\checkmark\\checkmarkc\checkmarki\checkmarkr\checkmarkc\checkmark$\checkmarke\checkmark$\checkmark^\checkmark\\checkmarkc\checkmarki\checkmarkr\checkmarkc\checkmark$\checkmarks\checkmark$\checkmark^\checkmark\\checkmarkc\checkmarki\checkmarkr\checkmarkc\checkmark$\checkmark \checkmark$\checkmark^\checkmark\\checkmarkc\checkmarki\checkmarkr\checkmarkc\checkmark$\checkmarkd\checkmark$\checkmark^\checkmark\\checkmarkc\checkmarki\checkmarkr\checkmarkc\checkmark$\checkmark'\checkmark$\checkmark^\checkmark\\checkmarkc\checkmarki\checkmarkr\checkmarkc\checkmark$\checkmarki\checkmark$\checkmark^\checkmark\\checkmarkc\checkmarki\checkmarkr\checkmarkc\checkmark$\checkmarkn\checkmark$\checkmark^\checkmark\\checkmarkc\checkmarki\checkmarkr\checkmarkc\checkmark$\checkmarke\checkmark$\checkmark^\checkmark\\checkmarkc\checkmarki\checkmarkr\checkmarkc\checkmark$\checkmarkr\checkmark$\checkmark^\checkmark\\checkmarkc\checkmarki\checkmarkr\checkmarkc\checkmark$\checkmarkt\checkmark$\checkmark^\checkmark\\checkmarkc\checkmarki\checkmarkr\checkmarkc\checkmark$\checkmarki\checkmark$\checkmark^\checkmark\\checkmarkc\checkmarki\checkmarkr\checkmarkc\checkmark$\checkmarke\checkmark$\checkmark^\checkmark\\checkmarkc\checkmarki\checkmarkr\checkmarkc\checkmark$\checkmark \checkmark$\checkmark^\checkmark\\checkmarkc\checkmarki\checkmarkr\checkmarkc\checkmark$\checkmarkl\checkmark$\checkmark^\checkmark\\checkmarkc\checkmarki\checkmarkr\checkmarkc\checkmark$\checkmarko\checkmark$\checkmark^\checkmark\\checkmarkc\checkmarki\checkmarkr\checkmarkc\checkmark$\checkmarkr\checkmark$\checkmark^\checkmark\\checkmarkc\checkmarki\checkmarkr\checkmarkc\checkmark$\checkmarks\checkmark$\checkmark^\checkmark\\checkmarkc\checkmarki\checkmarkr\checkmarkc\checkmark$\checkmark \checkmark$\checkmark^\checkmark\\checkmarkc\checkmarki\checkmarkr\checkmarkc\checkmark$\checkmarkd\checkmark$\checkmark^\checkmark\\checkmarkc\checkmarki\checkmarkr\checkmarkc\checkmark$\checkmarke\checkmark$\checkmark^\checkmark\\checkmarkc\checkmarki\checkmarkr\checkmarkc\checkmark$\checkmarks\checkmark$\checkmark^\checkmark\\checkmarkc\checkmarki\checkmarkr\checkmarkc\checkmark$\checkmark \checkmark$\checkmark^\checkmark\\checkmarkc\checkmarki\checkmarkr\checkmarkc\checkmark$\checkmarka\checkmark$\checkmark^\checkmark\\checkmarkc\checkmarki\checkmarkr\checkmarkc\checkmark$\checkmarkc\checkmark$\checkmark^\checkmark\\checkmarkc\checkmarki\checkmarkr\checkmarkc\checkmark$\checkmarkc\checkmark$\checkmark^\checkmark\\checkmarkc\checkmarki\checkmarkr\checkmarkc\checkmark$\checkmarké\checkmark$\checkmark^\checkmark\\checkmarkc\checkmarki\checkmarkr\checkmarkc\checkmark$\checkmarkl\checkmark$\checkmark^\checkmark\\checkmarkc\checkmarki\checkmarkr\checkmarkc\checkmark$\checkmarké\checkmark$\checkmark^\checkmark\\checkmarkc\checkmarki\checkmarkr\checkmarkc\checkmark$\checkmarkr\checkmark$\checkmark^\checkmark\\checkmarkc\checkmarki\checkmarkr\checkmarkc\checkmark$\checkmarka\checkmark$\checkmark^\checkmark\\checkmarkc\checkmarki\checkmarkr\checkmarkc\checkmark$\checkmarkt\checkmark$\checkmark^\checkmark\\checkmarkc\checkmarki\checkmarkr\checkmarkc\checkmark$\checkmarki\checkmark$\checkmark^\checkmark\\checkmarkc\checkmarki\checkmarkr\checkmarkc\checkmark$\checkmarko\checkmark$\checkmark^\checkmark\\checkmarkc\checkmarki\checkmarkr\checkmarkc\checkmark$\checkmarkn\checkmark$\checkmark^\checkmark\\checkmarkc\checkmarki\checkmarkr\checkmarkc\checkmark$\checkmarks\checkmark$\checkmark^\checkmark\\checkmarkc\checkmarki\checkmarkr\checkmarkc\checkmark$\checkmark.\checkmark$\checkmark^\checkmark\\checkmarkc\checkmarki\checkmarkr\checkmarkc\checkmark$\checkmark
\checkmark$\checkmark^\checkmark\\checkmarkc\checkmarki\checkmarkr\checkmarkc\checkmark$\checkmark
\checkmark$\checkmark^\checkmark\\checkmarkc\checkmarki\checkmarkr\checkmarkc\checkmark$\checkmark\\checkmark$\checkmark^\checkmark\\checkmarkc\checkmarki\checkmarkr\checkmarkc\checkmark$\checkmarkb\checkmark$\checkmark^\checkmark\\checkmarkc\checkmarki\checkmarkr\checkmarkc\checkmark$\checkmarke\checkmark$\checkmark^\checkmark\\checkmarkc\checkmarki\checkmarkr\checkmarkc\checkmark$\checkmarkg\checkmark$\checkmark^\checkmark\\checkmarkc\checkmarki\checkmarkr\checkmarkc\checkmark$\checkmarki\checkmark$\checkmark^\checkmark\\checkmarkc\checkmarki\checkmarkr\checkmarkc\checkmark$\checkmarkn\checkmark$\checkmark^\checkmark\\checkmarkc\checkmarki\checkmarkr\checkmarkc\checkmark$\checkmark{\checkmark$\checkmark^\checkmark\\checkmarkc\checkmarki\checkmarkr\checkmarkc\checkmark$\checkmarki\checkmark$\checkmark^\checkmark\\checkmarkc\checkmarki\checkmarkr\checkmarkc\checkmark$\checkmarkt\checkmark$\checkmark^\checkmark\\checkmarkc\checkmarki\checkmarkr\checkmarkc\checkmark$\checkmarke\checkmark$\checkmark^\checkmark\\checkmarkc\checkmarki\checkmarkr\checkmarkc\checkmark$\checkmarkm\checkmark$\checkmark^\checkmark\\checkmarkc\checkmarki\checkmarkr\checkmarkc\checkmark$\checkmarki\checkmark$\checkmark^\checkmark\\checkmarkc\checkmarki\checkmarkr\checkmarkc\checkmark$\checkmarkz\checkmark$\checkmark^\checkmark\\checkmarkc\checkmarki\checkmarkr\checkmarkc\checkmark$\checkmarke\checkmark$\checkmark^\checkmark\\checkmarkc\checkmarki\checkmarkr\checkmarkc\checkmark$\checkmark}\checkmark$\checkmark^\checkmark\\checkmarkc\checkmarki\checkmarkr\checkmarkc\checkmark$\checkmark
\checkmark$\checkmark^\checkmark\\checkmarkc\checkmarki\checkmarkr\checkmarkc\checkmark$\checkmark \checkmark$\checkmark^\checkmark\\checkmarkc\checkmarki\checkmarkr\checkmarkc\checkmark$\checkmark \checkmark$\checkmark^\checkmark\\checkmarkc\checkmarki\checkmarkr\checkmarkc\checkmark$\checkmark \checkmark$\checkmark^\checkmark\\checkmarkc\checkmarki\checkmarkr\checkmarkc\checkmark$\checkmark \checkmark$\checkmark^\checkmark\\checkmarkc\checkmarki\checkmarkr\checkmarkc\checkmark$\checkmark\\checkmark$\checkmark^\checkmark\\checkmarkc\checkmarki\checkmarkr\checkmarkc\checkmark$\checkmarki\checkmark$\checkmark^\checkmark\\checkmarkc\checkmarki\checkmarkr\checkmarkc\checkmark$\checkmarkt\checkmark$\checkmark^\checkmark\\checkmarkc\checkmarki\checkmarkr\checkmarkc\checkmark$\checkmarke\checkmark$\checkmark^\checkmark\\checkmarkc\checkmarki\checkmarkr\checkmarkc\checkmark$\checkmarkm\checkmark$\checkmark^\checkmark\\checkmarkc\checkmarki\checkmarkr\checkmarkc\checkmark$\checkmark \checkmark$\checkmark^\checkmark\\checkmarkc\checkmarki\checkmarkr\checkmarkc\checkmark$\checkmarkE\checkmark$\checkmark^\checkmark\\checkmarkc\checkmarki\checkmarkr\checkmarkc\checkmark$\checkmarkf\checkmark$\checkmark^\checkmark\\checkmarkc\checkmarki\checkmarkr\checkmarkc\checkmark$\checkmarkf\checkmark$\checkmark^\checkmark\\checkmarkc\checkmarki\checkmarkr\checkmarkc\checkmark$\checkmarko\checkmark$\checkmark^\checkmark\\checkmarkc\checkmarki\checkmarkr\checkmarkc\checkmark$\checkmarkr\checkmark$\checkmark^\checkmark\\checkmarkc\checkmarki\checkmarkr\checkmarkc\checkmark$\checkmarkt\checkmark$\checkmark^\checkmark\\checkmarkc\checkmarki\checkmarkr\checkmarkc\checkmark$\checkmark \checkmark$\checkmark^\checkmark\\checkmarkc\checkmarki\checkmarkr\checkmarkc\checkmark$\checkmarkd\checkmark$\checkmark^\checkmark\\checkmarkc\checkmarki\checkmarkr\checkmarkc\checkmark$\checkmarke\checkmark$\checkmark^\checkmark\\checkmarkc\checkmarki\checkmarkr\checkmarkc\checkmark$\checkmark \checkmark$\checkmark^\checkmark\\checkmarkc\checkmarki\checkmarkr\checkmarkc\checkmark$\checkmarkc\checkmark$\checkmark^\checkmark\\checkmarkc\checkmarki\checkmarkr\checkmarkc\checkmark$\checkmarko\checkmark$\checkmark^\checkmark\\checkmarkc\checkmarki\checkmarkr\checkmarkc\checkmark$\checkmarku\checkmark$\checkmark^\checkmark\\checkmarkc\checkmarki\checkmarkr\checkmarkc\checkmark$\checkmarkp\checkmark$\checkmark^\checkmark\\checkmarkc\checkmarki\checkmarkr\checkmarkc\checkmark$\checkmarke\checkmark$\checkmark^\checkmark\\checkmarkc\checkmarki\checkmarkr\checkmarkc\checkmark$\checkmark \checkmark$\checkmark^\checkmark\\checkmarkc\checkmarki\checkmarkr\checkmarkc\checkmark$\checkmarkt\checkmark$\checkmark^\checkmark\\checkmarkc\checkmarki\checkmarkr\checkmarkc\checkmark$\checkmarkr\checkmark$\checkmark^\checkmark\\checkmarkc\checkmarki\checkmarkr\checkmarkc\checkmark$\checkmarka\checkmark$\checkmark^\checkmark\\checkmarkc\checkmarki\checkmarkr\checkmarkc\checkmark$\checkmarkn\checkmark$\checkmark^\checkmark\\checkmarkc\checkmarki\checkmarkr\checkmarkc\checkmark$\checkmarks\checkmark$\checkmark^\checkmark\\checkmarkc\checkmarki\checkmarkr\checkmarkc\checkmark$\checkmarkm\checkmark$\checkmark^\checkmark\\checkmarkc\checkmarki\checkmarkr\checkmarkc\checkmark$\checkmarki\checkmark$\checkmark^\checkmark\\checkmarkc\checkmarki\checkmarkr\checkmarkc\checkmark$\checkmarks\checkmark$\checkmark^\checkmark\\checkmarkc\checkmarki\checkmarkr\checkmarkc\checkmark$\checkmark \checkmark$\checkmark^\checkmark\\checkmarkc\checkmarki\checkmarkr\checkmarkc\checkmark$\checkmarka\checkmark$\checkmark^\checkmark\\checkmarkc\checkmarki\checkmarkr\checkmarkc\checkmark$\checkmarkx\checkmark$\checkmark^\checkmark\\checkmarkc\checkmarki\checkmarkr\checkmarkc\checkmark$\checkmarki\checkmark$\checkmark^\checkmark\\checkmarkc\checkmarki\checkmarkr\checkmarkc\checkmark$\checkmarka\checkmark$\checkmark^\checkmark\\checkmarkc\checkmarki\checkmarkr\checkmarkc\checkmark$\checkmarkl\checkmark$\checkmark^\checkmark\\checkmarkc\checkmarki\checkmarkr\checkmarkc\checkmark$\checkmarke\checkmark$\checkmark^\checkmark\\checkmarkc\checkmarki\checkmarkr\checkmarkc\checkmark$\checkmarkm\checkmark$\checkmark^\checkmark\\checkmarkc\checkmarki\checkmarkr\checkmarkc\checkmark$\checkmarke\checkmark$\checkmark^\checkmark\\checkmarkc\checkmarki\checkmarkr\checkmarkc\checkmark$\checkmarkn\checkmark$\checkmark^\checkmark\\checkmarkc\checkmarki\checkmarkr\checkmarkc\checkmark$\checkmarkt\checkmark$\checkmark^\checkmark\\checkmarkc\checkmarki\checkmarkr\checkmarkc\checkmark$\checkmark \checkmark$\checkmark^\checkmark\\checkmarkc\checkmarki\checkmarkr\checkmarkc\checkmark$\checkmark:\checkmark$\checkmark^\checkmark\\checkmarkc\checkmarki\checkmarkr\checkmarkc\checkmark$\checkmark \checkmark$\checkmark^\checkmark\\checkmarkc\checkmarki\checkmarkr\checkmarkc\checkmark$\checkmark$\checkmark$\checkmark^\checkmark\\checkmarkc\checkmarki\checkmarkr\checkmarkc\checkmark$\checkmarkF\checkmark$\checkmark^\checkmark\\checkmarkc\checkmarki\checkmarkr\checkmarkc\checkmark$\checkmark_\checkmark$\checkmark^\checkmark\\checkmarkc\checkmarki\checkmarkr\checkmarkc\checkmark$\checkmark{\checkmark$\checkmark^\checkmark\\checkmarkc\checkmarki\checkmarkr\checkmarkc\checkmark$\checkmarku\checkmark$\checkmark^\checkmark\\checkmarkc\checkmarki\checkmarkr\checkmarkc\checkmark$\checkmarks\checkmark$\checkmark^\checkmark\\checkmarkc\checkmarki\checkmarkr\checkmarkc\checkmark$\checkmarki\checkmark$\checkmark^\checkmark\\checkmarkc\checkmarki\checkmarkr\checkmarkc\checkmark$\checkmarkn\checkmark$\checkmark^\checkmark\\checkmarkc\checkmarki\checkmarkr\checkmarkc\checkmark$\checkmarka\checkmark$\checkmark^\checkmark\\checkmarkc\checkmarki\checkmarkr\checkmarkc\checkmark$\checkmarkg\checkmark$\checkmark^\checkmark\\checkmarkc\checkmarki\checkmarkr\checkmarkc\checkmark$\checkmarke\checkmark$\checkmark^\checkmark\\checkmarkc\checkmarki\checkmarkr\checkmarkc\checkmark$\checkmark}\checkmark$\checkmark^\checkmark\\checkmarkc\checkmarki\checkmarkr\checkmarkc\checkmark$\checkmark \checkmark$\checkmark^\checkmark\\checkmarkc\checkmarki\checkmarkr\checkmarkc\checkmark$\checkmark=\checkmark$\checkmark^\checkmark\\checkmarkc\checkmarki\checkmarkr\checkmarkc\checkmark$\checkmark \checkmark$\checkmark^\checkmark\\checkmarkc\checkmarki\checkmarkr\checkmarkc\checkmark$\checkmarkF\checkmark$\checkmark^\checkmark\\checkmarkc\checkmarki\checkmarkr\checkmarkc\checkmark$\checkmark_\checkmark$\checkmark^\checkmark\\checkmarkc\checkmarki\checkmarkr\checkmarkc\checkmark$\checkmark{\checkmark$\checkmark^\checkmark\\checkmarkc\checkmarki\checkmarkr\checkmarkc\checkmark$\checkmarkm\checkmark$\checkmark^\checkmark\\checkmarkc\checkmarki\checkmarkr\checkmarkc\checkmark$\checkmarka\checkmark$\checkmark^\checkmark\\checkmarkc\checkmarki\checkmarkr\checkmarkc\checkmark$\checkmarkx\checkmark$\checkmark^\checkmark\\checkmarkc\checkmarki\checkmarkr\checkmarkc\checkmark$\checkmark}\checkmark$\checkmark^\checkmark\\checkmarkc\checkmarki\checkmarkr\checkmarkc\checkmark$\checkmark \checkmark$\checkmark^\checkmark\\checkmarkc\checkmarki\checkmarkr\checkmarkc\checkmark$\checkmark=\checkmark$\checkmark^\checkmark\\checkmarkc\checkmarki\checkmarkr\checkmarkc\checkmark$\checkmark \checkmark$\checkmark^\checkmark\\checkmarkc\checkmarki\checkmarkr\checkmarkc\checkmark$\checkmark1\checkmark$\checkmark^\checkmark\\checkmarkc\checkmarki\checkmarkr\checkmarkc\checkmark$\checkmark0\checkmark$\checkmark^\checkmark\\checkmarkc\checkmarki\checkmarkr\checkmarkc\checkmark$\checkmark0\checkmark$\checkmark^\checkmark\\checkmarkc\checkmarki\checkmarkr\checkmarkc\checkmark$\checkmark$\checkmark$\checkmark^\checkmark\\checkmarkc\checkmarki\checkmarkr\checkmarkc\checkmark$\checkmark \checkmark$\checkmark^\checkmark\\checkmarkc\checkmarki\checkmarkr\checkmarkc\checkmark$\checkmarkN\checkmark$\checkmark^\checkmark\\checkmarkc\checkmarki\checkmarkr\checkmarkc\checkmark$\checkmark
\checkmark$\checkmark^\checkmark\\checkmarkc\checkmarki\checkmarkr\checkmarkc\checkmark$\checkmark \checkmark$\checkmark^\checkmark\\checkmarkc\checkmarki\checkmarkr\checkmarkc\checkmark$\checkmark \checkmark$\checkmark^\checkmark\\checkmarkc\checkmarki\checkmarkr\checkmarkc\checkmark$\checkmark \checkmark$\checkmark^\checkmark\\checkmarkc\checkmarki\checkmarkr\checkmarkc\checkmark$\checkmark \checkmark$\checkmark^\checkmark\\checkmarkc\checkmarki\checkmarkr\checkmarkc\checkmark$\checkmark\\checkmark$\checkmark^\checkmark\\checkmarkc\checkmarki\checkmarkr\checkmarkc\checkmark$\checkmarki\checkmark$\checkmark^\checkmark\\checkmarkc\checkmarki\checkmarkr\checkmarkc\checkmark$\checkmarkt\checkmark$\checkmark^\checkmark\\checkmarkc\checkmarki\checkmarkr\checkmarkc\checkmark$\checkmarke\checkmark$\checkmark^\checkmark\\checkmarkc\checkmarki\checkmarkr\checkmarkc\checkmark$\checkmarkm\checkmark$\checkmark^\checkmark\\checkmarkc\checkmarki\checkmarkr\checkmarkc\checkmark$\checkmark \checkmark$\checkmark^\checkmark\\checkmarkc\checkmarki\checkmarkr\checkmarkc\checkmark$\checkmarkF\checkmark$\checkmark^\checkmark\\checkmarkc\checkmarki\checkmarkr\checkmarkc\checkmark$\checkmarko\checkmark$\checkmark^\checkmark\\checkmarkc\checkmarki\checkmarkr\checkmarkc\checkmark$\checkmarkr\checkmark$\checkmark^\checkmark\\checkmarkc\checkmarki\checkmarkr\checkmarkc\checkmark$\checkmarkc\checkmark$\checkmark^\checkmark\\checkmarkc\checkmarki\checkmarkr\checkmarkc\checkmark$\checkmarke\checkmark$\checkmark^\checkmark\\checkmarkc\checkmarki\checkmarkr\checkmarkc\checkmark$\checkmark \checkmark$\checkmark^\checkmark\\checkmarkc\checkmarki\checkmarkr\checkmarkc\checkmark$\checkmarkd\checkmark$\checkmark^\checkmark\\checkmarkc\checkmarki\checkmarkr\checkmarkc\checkmark$\checkmarke\checkmark$\checkmark^\checkmark\\checkmarkc\checkmarki\checkmarkr\checkmarkc\checkmark$\checkmark \checkmark$\checkmark^\checkmark\\checkmarkc\checkmarki\checkmarkr\checkmarkc\checkmark$\checkmarkf\checkmark$\checkmark^\checkmark\\checkmarkc\checkmarki\checkmarkr\checkmarkc\checkmark$\checkmarkr\checkmark$\checkmark^\checkmark\\checkmarkc\checkmarki\checkmarkr\checkmarkc\checkmark$\checkmarko\checkmark$\checkmark^\checkmark\\checkmarkc\checkmarki\checkmarkr\checkmarkc\checkmark$\checkmarkt\checkmark$\checkmark^\checkmark\\checkmarkc\checkmarki\checkmarkr\checkmarkc\checkmark$\checkmarkt\checkmark$\checkmark^\checkmark\\checkmarkc\checkmarki\checkmarkr\checkmarkc\checkmark$\checkmarke\checkmark$\checkmark^\checkmark\\checkmarkc\checkmarki\checkmarkr\checkmarkc\checkmark$\checkmarkm\checkmark$\checkmark^\checkmark\\checkmarkc\checkmarki\checkmarkr\checkmarkc\checkmark$\checkmarke\checkmark$\checkmark^\checkmark\\checkmarkc\checkmarki\checkmarkr\checkmarkc\checkmark$\checkmarkn\checkmark$\checkmark^\checkmark\\checkmarkc\checkmarki\checkmarkr\checkmarkc\checkmark$\checkmarkt\checkmark$\checkmark^\checkmark\\checkmarkc\checkmarki\checkmarkr\checkmarkc\checkmark$\checkmark \checkmark$\checkmark^\checkmark\\checkmarkc\checkmarki\checkmarkr\checkmarkc\checkmark$\checkmarks\checkmark$\checkmark^\checkmark\\checkmarkc\checkmarki\checkmarkr\checkmarkc\checkmark$\checkmarku\checkmark$\checkmark^\checkmark\\checkmarkc\checkmarki\checkmarkr\checkmarkc\checkmark$\checkmarkr\checkmark$\checkmark^\checkmark\\checkmarkc\checkmarki\checkmarkr\checkmarkc\checkmark$\checkmark \checkmark$\checkmark^\checkmark\\checkmarkc\checkmarki\checkmarkr\checkmarkc\checkmark$\checkmarkl\checkmark$\checkmark^\checkmark\\checkmarkc\checkmarki\checkmarkr\checkmarkc\checkmark$\checkmarke\checkmark$\checkmark^\checkmark\\checkmarkc\checkmarki\checkmarkr\checkmarkc\checkmark$\checkmarks\checkmark$\checkmark^\checkmark\\checkmarkc\checkmarki\checkmarkr\checkmarkc\checkmark$\checkmark \checkmark$\checkmark^\checkmark\\checkmarkc\checkmarki\checkmarkr\checkmarkc\checkmark$\checkmarkg\checkmark$\checkmark^\checkmark\\checkmarkc\checkmarki\checkmarkr\checkmarkc\checkmark$\checkmarku\checkmark$\checkmark^\checkmark\\checkmarkc\checkmarki\checkmarkr\checkmarkc\checkmark$\checkmarki\checkmark$\checkmark^\checkmark\\checkmarkc\checkmarki\checkmarkr\checkmarkc\checkmark$\checkmarkd\checkmark$\checkmark^\checkmark\\checkmarkc\checkmarki\checkmarkr\checkmarkc\checkmark$\checkmarka\checkmark$\checkmark^\checkmark\\checkmarkc\checkmarki\checkmarkr\checkmarkc\checkmark$\checkmarkg\checkmark$\checkmark^\checkmark\\checkmarkc\checkmarki\checkmarkr\checkmarkc\checkmark$\checkmarke\checkmark$\checkmark^\checkmark\\checkmarkc\checkmarki\checkmarkr\checkmarkc\checkmark$\checkmarks\checkmark$\checkmark^\checkmark\\checkmarkc\checkmarki\checkmarkr\checkmarkc\checkmark$\checkmark \checkmark$\checkmark^\checkmark\\checkmarkc\checkmarki\checkmarkr\checkmarkc\checkmark$\checkmark:\checkmark$\checkmark^\checkmark\\checkmarkc\checkmarki\checkmarkr\checkmarkc\checkmark$\checkmark \checkmark$\checkmark^\checkmark\\checkmarkc\checkmarki\checkmarkr\checkmarkc\checkmark$\checkmark$\checkmark$\checkmark^\checkmark\\checkmarkc\checkmarki\checkmarkr\checkmarkc\checkmark$\checkmarkF\checkmark$\checkmark^\checkmark\\checkmarkc\checkmarki\checkmarkr\checkmarkc\checkmark$\checkmark_\checkmark$\checkmark^\checkmark\\checkmarkc\checkmarki\checkmarkr\checkmarkc\checkmark$\checkmark{\checkmark$\checkmark^\checkmark\\checkmarkc\checkmarki\checkmarkr\checkmarkc\checkmark$\checkmarkf\checkmark$\checkmark^\checkmark\\checkmarkc\checkmarki\checkmarkr\checkmarkc\checkmark$\checkmarkr\checkmark$\checkmark^\checkmark\\checkmarkc\checkmarki\checkmarkr\checkmarkc\checkmark$\checkmarko\checkmark$\checkmark^\checkmark\\checkmarkc\checkmarki\checkmarkr\checkmarkc\checkmark$\checkmarkt\checkmark$\checkmark^\checkmark\\checkmarkc\checkmarki\checkmarkr\checkmarkc\checkmark$\checkmarkt\checkmark$\checkmark^\checkmark\\checkmarkc\checkmarki\checkmarkr\checkmarkc\checkmark$\checkmark}\checkmark$\checkmark^\checkmark\\checkmarkc\checkmarki\checkmarkr\checkmarkc\checkmark$\checkmark \checkmark$\checkmark^\checkmark\\checkmarkc\checkmarki\checkmarkr\checkmarkc\checkmark$\checkmark=\checkmark$\checkmark^\checkmark\\checkmarkc\checkmarki\checkmarkr\checkmarkc\checkmark$\checkmark \checkmark$\checkmark^\checkmark\\checkmarkc\checkmarki\checkmarkr\checkmarkc\checkmark$\checkmark\\checkmark$\checkmark^\checkmark\\checkmarkc\checkmarki\checkmarkr\checkmarkc\checkmark$\checkmarkm\checkmark$\checkmark^\checkmark\\checkmarkc\checkmarki\checkmarkr\checkmarkc\checkmark$\checkmarku\checkmark$\checkmark^\checkmark\\checkmarkc\checkmarki\checkmarkr\checkmarkc\checkmark$\checkmark \checkmark$\checkmark^\checkmark\\checkmarkc\checkmarki\checkmarkr\checkmarkc\checkmark$\checkmark\\checkmark$\checkmark^\checkmark\\checkmarkc\checkmarki\checkmarkr\checkmarkc\checkmark$\checkmarkt\checkmark$\checkmark^\checkmark\\checkmarkc\checkmarki\checkmarkr\checkmarkc\checkmark$\checkmarki\checkmark$\checkmark^\checkmark\\checkmarkc\checkmarki\checkmarkr\checkmarkc\checkmark$\checkmarkm\checkmark$\checkmark^\checkmark\\checkmarkc\checkmarki\checkmarkr\checkmarkc\checkmark$\checkmarke\checkmark$\checkmark^\checkmark\\checkmarkc\checkmarki\checkmarkr\checkmarkc\checkmark$\checkmarks\checkmark$\checkmark^\checkmark\\checkmarkc\checkmarki\checkmarkr\checkmarkc\checkmark$\checkmark \checkmark$\checkmark^\checkmark\\checkmarkc\checkmarki\checkmarkr\checkmarkc\checkmark$\checkmarkM\checkmark$\checkmark^\checkmark\\checkmarkc\checkmarki\checkmarkr\checkmarkc\checkmark$\checkmark \checkmark$\checkmark^\checkmark\\checkmarkc\checkmarki\checkmarkr\checkmarkc\checkmark$\checkmark\\checkmark$\checkmark^\checkmark\\checkmarkc\checkmarki\checkmarkr\checkmarkc\checkmark$\checkmarkt\checkmark$\checkmark^\checkmark\\checkmarkc\checkmarki\checkmarkr\checkmarkc\checkmark$\checkmarki\checkmark$\checkmark^\checkmark\\checkmarkc\checkmarki\checkmarkr\checkmarkc\checkmark$\checkmarkm\checkmark$\checkmark^\checkmark\\checkmarkc\checkmarki\checkmarkr\checkmarkc\checkmark$\checkmarke\checkmark$\checkmark^\checkmark\\checkmarkc\checkmarki\checkmarkr\checkmarkc\checkmark$\checkmarks\checkmark$\checkmark^\checkmark\\checkmarkc\checkmarki\checkmarkr\checkmarkc\checkmark$\checkmark \checkmark$\checkmark^\checkmark\\checkmarkc\checkmarki\checkmarkr\checkmarkc\checkmark$\checkmarkg\checkmark$\checkmark^\checkmark\\checkmarkc\checkmarki\checkmarkr\checkmarkc\checkmark$\checkmark \checkmark$\checkmark^\checkmark\\checkmarkc\checkmarki\checkmarkr\checkmarkc\checkmark$\checkmark=\checkmark$\checkmark^\checkmark\\checkmarkc\checkmarki\checkmarkr\checkmarkc\checkmark$\checkmark \checkmark$\checkmark^\checkmark\\checkmarkc\checkmarki\checkmarkr\checkmarkc\checkmark$\checkmark0\checkmark$\checkmark^\checkmark\\checkmarkc\checkmarki\checkmarkr\checkmarkc\checkmark$\checkmark{\checkmark$\checkmark^\checkmark\\checkmarkc\checkmarki\checkmarkr\checkmarkc\checkmark$\checkmark,\checkmark$\checkmark^\checkmark\\checkmarkc\checkmarki\checkmarkr\checkmarkc\checkmark$\checkmark}\checkmark$\checkmark^\checkmark\\checkmarkc\checkmarki\checkmarkr\checkmarkc\checkmark$\checkmark0\checkmark$\checkmark^\checkmark\\checkmarkc\checkmarki\checkmarkr\checkmarkc\checkmark$\checkmark0\checkmark$\checkmark^\checkmark\\checkmarkc\checkmarki\checkmarkr\checkmarkc\checkmark$\checkmark5\checkmark$\checkmark^\checkmark\\checkmarkc\checkmarki\checkmarkr\checkmarkc\checkmark$\checkmark \checkmark$\checkmark^\checkmark\\checkmarkc\checkmarki\checkmarkr\checkmarkc\checkmark$\checkmark\\checkmark$\checkmark^\checkmark\\checkmarkc\checkmarki\checkmarkr\checkmarkc\checkmark$\checkmarkt\checkmark$\checkmark^\checkmark\\checkmarkc\checkmarki\checkmarkr\checkmarkc\checkmark$\checkmarki\checkmark$\checkmark^\checkmark\\checkmarkc\checkmarki\checkmarkr\checkmarkc\checkmark$\checkmarkm\checkmark$\checkmark^\checkmark\\checkmarkc\checkmarki\checkmarkr\checkmarkc\checkmark$\checkmarke\checkmark$\checkmark^\checkmark\\checkmarkc\checkmarki\checkmarkr\checkmarkc\checkmark$\checkmarks\checkmark$\checkmark^\checkmark\\checkmarkc\checkmarki\checkmarkr\checkmarkc\checkmark$\checkmark \checkmark$\checkmark^\checkmark\\checkmarkc\checkmarki\checkmarkr\checkmarkc\checkmark$\checkmark1\checkmark$\checkmark^\checkmark\\checkmarkc\checkmarki\checkmarkr\checkmarkc\checkmark$\checkmark0\checkmark$\checkmark^\checkmark\\checkmarkc\checkmarki\checkmarkr\checkmarkc\checkmark$\checkmark \checkmark$\checkmark^\checkmark\\checkmarkc\checkmarki\checkmarkr\checkmarkc\checkmark$\checkmark\\checkmark$\checkmark^\checkmark\\checkmarkc\checkmarki\checkmarkr\checkmarkc\checkmark$\checkmarkt\checkmark$\checkmark^\checkmark\\checkmarkc\checkmarki\checkmarkr\checkmarkc\checkmark$\checkmarki\checkmark$\checkmark^\checkmark\\checkmarkc\checkmarki\checkmarkr\checkmarkc\checkmark$\checkmarkm\checkmark$\checkmark^\checkmark\\checkmarkc\checkmarki\checkmarkr\checkmarkc\checkmark$\checkmarke\checkmark$\checkmark^\checkmark\\checkmarkc\checkmarki\checkmarkr\checkmarkc\checkmark$\checkmarks\checkmark$\checkmark^\checkmark\\checkmarkc\checkmarki\checkmarkr\checkmarkc\checkmark$\checkmark \checkmark$\checkmark^\checkmark\\checkmarkc\checkmarki\checkmarkr\checkmarkc\checkmark$\checkmark9\checkmark$\checkmark^\checkmark\\checkmarkc\checkmarki\checkmarkr\checkmarkc\checkmark$\checkmark{\checkmark$\checkmark^\checkmark\\checkmarkc\checkmarki\checkmarkr\checkmarkc\checkmark$\checkmark,\checkmark$\checkmark^\checkmark\\checkmarkc\checkmarki\checkmarkr\checkmarkc\checkmark$\checkmark}\checkmark$\checkmark^\checkmark\\checkmarkc\checkmarki\checkmarkr\checkmarkc\checkmark$\checkmark8\checkmark$\checkmark^\checkmark\\checkmarkc\checkmarki\checkmarkr\checkmarkc\checkmark$\checkmark1\checkmark$\checkmark^\checkmark\\checkmarkc\checkmarki\checkmarkr\checkmarkc\checkmark$\checkmark \checkmark$\checkmark^\checkmark\\checkmarkc\checkmarki\checkmarkr\checkmarkc\checkmark$\checkmark=\checkmark$\checkmark^\checkmark\\checkmarkc\checkmarki\checkmarkr\checkmarkc\checkmark$\checkmark \checkmark$\checkmark^\checkmark\\checkmarkc\checkmarki\checkmarkr\checkmarkc\checkmark$\checkmark0\checkmark$\checkmark^\checkmark\\checkmarkc\checkmarki\checkmarkr\checkmarkc\checkmark$\checkmark{\checkmark$\checkmark^\checkmark\\checkmarkc\checkmarki\checkmarkr\checkmarkc\checkmark$\checkmark,\checkmark$\checkmark^\checkmark\\checkmarkc\checkmarki\checkmarkr\checkmarkc\checkmark$\checkmark}\checkmark$\checkmark^\checkmark\\checkmarkc\checkmarki\checkmarkr\checkmarkc\checkmark$\checkmark4\checkmark$\checkmark^\checkmark\\checkmarkc\checkmarki\checkmarkr\checkmarkc\checkmark$\checkmark9\checkmark$\checkmark^\checkmark\\checkmarkc\checkmarki\checkmarkr\checkmarkc\checkmark$\checkmark$\checkmark$\checkmark^\checkmark\\checkmarkc\checkmarki\checkmarkr\checkmarkc\checkmark$\checkmark \checkmark$\checkmark^\checkmark\\checkmarkc\checkmarki\checkmarkr\checkmarkc\checkmark$\checkmarkN\checkmark$\checkmark^\checkmark\\checkmarkc\checkmarki\checkmarkr\checkmarkc\checkmark$\checkmark \checkmark$\checkmark^\checkmark\\checkmarkc\checkmarki\checkmarkr\checkmarkc\checkmark$\checkmark$\checkmark$\checkmark^\checkmark\\checkmarkc\checkmarki\checkmarkr\checkmarkc\checkmark$\checkmark\\checkmark$\checkmark^\checkmark\\checkmarkc\checkmarki\checkmarkr\checkmarkc\checkmark$\checkmarka\checkmark$\checkmark^\checkmark\\checkmarkc\checkmarki\checkmarkr\checkmarkc\checkmark$\checkmarkp\checkmark$\checkmark^\checkmark\\checkmarkc\checkmarki\checkmarkr\checkmarkc\checkmark$\checkmarkp\checkmark$\checkmark^\checkmark\\checkmarkc\checkmarki\checkmarkr\checkmarkc\checkmark$\checkmarkr\checkmark$\checkmark^\checkmark\\checkmarkc\checkmarki\checkmarkr\checkmarkc\checkmark$\checkmarko\checkmark$\checkmark^\checkmark\\checkmarkc\checkmarki\checkmarkr\checkmarkc\checkmark$\checkmarkx\checkmark$\checkmark^\checkmark\\checkmarkc\checkmarki\checkmarkr\checkmarkc\checkmark$\checkmark \checkmark$\checkmark^\checkmark\\checkmarkc\checkmarki\checkmarkr\checkmarkc\checkmark$\checkmark0\checkmark$\checkmark^\checkmark\\checkmarkc\checkmarki\checkmarkr\checkmarkc\checkmark$\checkmark{\checkmark$\checkmark^\checkmark\\checkmarkc\checkmarki\checkmarkr\checkmarkc\checkmark$\checkmark,\checkmark$\checkmark^\checkmark\\checkmarkc\checkmarki\checkmarkr\checkmarkc\checkmark$\checkmark}\checkmark$\checkmark^\checkmark\\checkmarkc\checkmarki\checkmarkr\checkmarkc\checkmark$\checkmark5\checkmark$\checkmark^\checkmark\\checkmarkc\checkmarki\checkmarkr\checkmarkc\checkmark$\checkmark$\checkmark$\checkmark^\checkmark\\checkmarkc\checkmarki\checkmarkr\checkmarkc\checkmark$\checkmark \checkmark$\checkmark^\checkmark\\checkmarkc\checkmarki\checkmarkr\checkmarkc\checkmark$\checkmarkN\checkmark$\checkmark^\checkmark\\checkmarkc\checkmarki\checkmarkr\checkmarkc\checkmark$\checkmark \checkmark$\checkmark^\checkmark\\checkmarkc\checkmarki\checkmarkr\checkmarkc\checkmark$\checkmark(\checkmark$\checkmark^\checkmark\\checkmarkc\checkmarki\checkmarkr\checkmarkc\checkmark$\checkmarkn\checkmark$\checkmark^\checkmark\\checkmarkc\checkmarki\checkmarkr\checkmarkc\checkmark$\checkmarké\checkmark$\checkmark^\checkmark\\checkmarkc\checkmarki\checkmarkr\checkmarkc\checkmark$\checkmarkg\checkmark$\checkmark^\checkmark\\checkmarkc\checkmarki\checkmarkr\checkmarkc\checkmark$\checkmarkl\checkmark$\checkmark^\checkmark\\checkmarkc\checkmarki\checkmarkr\checkmarkc\checkmark$\checkmarki\checkmark$\checkmark^\checkmark\\checkmarkc\checkmarki\checkmarkr\checkmarkc\checkmark$\checkmarkg\checkmark$\checkmark^\checkmark\\checkmarkc\checkmarki\checkmarkr\checkmarkc\checkmark$\checkmarke\checkmark$\checkmark^\checkmark\\checkmarkc\checkmarki\checkmarkr\checkmarkc\checkmark$\checkmarka\checkmark$\checkmark^\checkmark\\checkmarkc\checkmarki\checkmarkr\checkmarkc\checkmark$\checkmarkb\checkmark$\checkmark^\checkmark\\checkmarkc\checkmarki\checkmarkr\checkmarkc\checkmark$\checkmarkl\checkmark$\checkmark^\checkmark\\checkmarkc\checkmarki\checkmarkr\checkmarkc\checkmark$\checkmarke\checkmark$\checkmark^\checkmark\\checkmarkc\checkmarki\checkmarkr\checkmarkc\checkmark$\checkmark)\checkmark$\checkmark^\checkmark\\checkmarkc\checkmarki\checkmarkr\checkmarkc\checkmark$\checkmark
\checkmark$\checkmark^\checkmark\\checkmarkc\checkmarki\checkmarkr\checkmarkc\checkmark$\checkmark \checkmark$\checkmark^\checkmark\\checkmarkc\checkmarki\checkmarkr\checkmarkc\checkmark$\checkmark \checkmark$\checkmark^\checkmark\\checkmarkc\checkmarki\checkmarkr\checkmarkc\checkmark$\checkmark \checkmark$\checkmark^\checkmark\\checkmarkc\checkmarki\checkmarkr\checkmarkc\checkmark$\checkmark \checkmark$\checkmark^\checkmark\\checkmarkc\checkmarki\checkmarkr\checkmarkc\checkmark$\checkmark\\checkmark$\checkmark^\checkmark\\checkmarkc\checkmarki\checkmarkr\checkmarkc\checkmark$\checkmarki\checkmark$\checkmark^\checkmark\\checkmarkc\checkmarki\checkmarkr\checkmarkc\checkmark$\checkmarkt\checkmark$\checkmark^\checkmark\\checkmarkc\checkmarki\checkmarkr\checkmarkc\checkmark$\checkmarke\checkmark$\checkmark^\checkmark\\checkmarkc\checkmarki\checkmarkr\checkmarkc\checkmark$\checkmarkm\checkmark$\checkmark^\checkmark\\checkmarkc\checkmarki\checkmarkr\checkmarkc\checkmark$\checkmark \checkmark$\checkmark^\checkmark\\checkmarkc\checkmarki\checkmarkr\checkmarkc\checkmark$\checkmarkF\checkmark$\checkmark^\checkmark\\checkmarkc\checkmarki\checkmarkr\checkmarkc\checkmark$\checkmarko\checkmark$\checkmark^\checkmark\\checkmarkc\checkmarki\checkmarkr\checkmarkc\checkmark$\checkmarkr\checkmark$\checkmark^\checkmark\\checkmarkc\checkmarki\checkmarkr\checkmarkc\checkmark$\checkmarkc\checkmark$\checkmark^\checkmark\\checkmarkc\checkmarki\checkmarkr\checkmarkc\checkmark$\checkmarke\checkmark$\checkmark^\checkmark\\checkmarkc\checkmarki\checkmarkr\checkmarkc\checkmark$\checkmark \checkmark$\checkmark^\checkmark\\checkmarkc\checkmarki\checkmarkr\checkmarkc\checkmark$\checkmarkd\checkmark$\checkmark^\checkmark\\checkmarkc\checkmarki\checkmarkr\checkmarkc\checkmark$\checkmark'\checkmark$\checkmark^\checkmark\\checkmarkc\checkmarki\checkmarkr\checkmarkc\checkmark$\checkmarki\checkmark$\checkmark^\checkmark\\checkmarkc\checkmarki\checkmarkr\checkmarkc\checkmark$\checkmarkn\checkmark$\checkmark^\checkmark\\checkmarkc\checkmarki\checkmarkr\checkmarkc\checkmark$\checkmarke\checkmark$\checkmark^\checkmark\\checkmarkc\checkmarki\checkmarkr\checkmarkc\checkmark$\checkmarkr\checkmark$\checkmark^\checkmark\\checkmarkc\checkmarki\checkmarkr\checkmarkc\checkmark$\checkmarkt\checkmark$\checkmark^\checkmark\\checkmarkc\checkmarki\checkmarkr\checkmarkc\checkmark$\checkmarki\checkmark$\checkmark^\checkmark\\checkmarkc\checkmarki\checkmarkr\checkmarkc\checkmark$\checkmarke\checkmark$\checkmark^\checkmark\\checkmarkc\checkmarki\checkmarkr\checkmarkc\checkmark$\checkmark \checkmark$\checkmark^\checkmark\\checkmarkc\checkmarki\checkmarkr\checkmarkc\checkmark$\checkmarkl\checkmark$\checkmark^\checkmark\\checkmarkc\checkmarki\checkmarkr\checkmarkc\checkmark$\checkmarko\checkmark$\checkmark^\checkmark\\checkmarkc\checkmarki\checkmarkr\checkmarkc\checkmark$\checkmarkr\checkmark$\checkmark^\checkmark\\checkmarkc\checkmarki\checkmarkr\checkmarkc\checkmark$\checkmarks\checkmark$\checkmark^\checkmark\\checkmarkc\checkmarki\checkmarkr\checkmarkc\checkmark$\checkmark \checkmark$\checkmark^\checkmark\\checkmarkc\checkmarki\checkmarkr\checkmarkc\checkmark$\checkmarkd\checkmark$\checkmark^\checkmark\\checkmarkc\checkmarki\checkmarkr\checkmarkc\checkmark$\checkmarke\checkmark$\checkmark^\checkmark\\checkmarkc\checkmarki\checkmarkr\checkmarkc\checkmark$\checkmark \checkmark$\checkmark^\checkmark\\checkmarkc\checkmarki\checkmarkr\checkmarkc\checkmark$\checkmarkl\checkmark$\checkmark^\checkmark\\checkmarkc\checkmarki\checkmarkr\checkmarkc\checkmark$\checkmark'\checkmark$\checkmark^\checkmark\\checkmarkc\checkmarki\checkmarkr\checkmarkc\checkmark$\checkmarka\checkmark$\checkmark^\checkmark\\checkmarkc\checkmarki\checkmarkr\checkmarkc\checkmark$\checkmarkc\checkmark$\checkmark^\checkmark\\checkmarkc\checkmarki\checkmarkr\checkmarkc\checkmark$\checkmarkc\checkmark$\checkmark^\checkmark\\checkmarkc\checkmarki\checkmarkr\checkmarkc\checkmark$\checkmarké\checkmark$\checkmark^\checkmark\\checkmarkc\checkmarki\checkmarkr\checkmarkc\checkmark$\checkmarkl\checkmark$\checkmark^\checkmark\\checkmarkc\checkmarki\checkmarkr\checkmarkc\checkmark$\checkmarké\checkmark$\checkmark^\checkmark\\checkmarkc\checkmarki\checkmarkr\checkmarkc\checkmark$\checkmarkr\checkmark$\checkmark^\checkmark\\checkmarkc\checkmarki\checkmarkr\checkmarkc\checkmark$\checkmarka\checkmark$\checkmark^\checkmark\\checkmarkc\checkmarki\checkmarkr\checkmarkc\checkmark$\checkmarkt\checkmark$\checkmark^\checkmark\\checkmarkc\checkmarki\checkmarkr\checkmarkc\checkmark$\checkmarki\checkmark$\checkmark^\checkmark\\checkmarkc\checkmarki\checkmarkr\checkmarkc\checkmark$\checkmarko\checkmark$\checkmark^\checkmark\\checkmarkc\checkmarki\checkmarkr\checkmarkc\checkmark$\checkmarkn\checkmark$\checkmark^\checkmark\\checkmarkc\checkmarki\checkmarkr\checkmarkc\checkmark$\checkmark \checkmark$\checkmark^\checkmark\\checkmarkc\checkmarki\checkmarkr\checkmarkc\checkmark$\checkmark(\checkmark$\checkmark^\checkmark\\checkmarkc\checkmarki\checkmarkr\checkmarkc\checkmark$\checkmark$\checkmark$\checkmark^\checkmark\\checkmarkc\checkmarki\checkmarkr\checkmarkc\checkmark$\checkmarka\checkmark$\checkmark^\checkmark\\checkmarkc\checkmarki\checkmarkr\checkmarkc\checkmark$\checkmark_\checkmark$\checkmark^\checkmark\\checkmarkc\checkmarki\checkmarkr\checkmarkc\checkmark$\checkmark{\checkmark$\checkmark^\checkmark\\checkmarkc\checkmarki\checkmarkr\checkmarkc\checkmark$\checkmarkm\checkmark$\checkmark^\checkmark\\checkmarkc\checkmarki\checkmarkr\checkmarkc\checkmark$\checkmarka\checkmark$\checkmark^\checkmark\\checkmarkc\checkmarki\checkmarkr\checkmarkc\checkmark$\checkmarkx\checkmark$\checkmark^\checkmark\\checkmarkc\checkmarki\checkmarkr\checkmarkc\checkmark$\checkmark}\checkmark$\checkmark^\checkmark\\checkmarkc\checkmarki\checkmarkr\checkmarkc\checkmark$\checkmark \checkmark$\checkmark^\checkmark\\checkmarkc\checkmarki\checkmarkr\checkmarkc\checkmark$\checkmark=\checkmark$\checkmark^\checkmark\\checkmarkc\checkmarki\checkmarkr\checkmarkc\checkmark$\checkmark \checkmark$\checkmark^\checkmark\\checkmarkc\checkmarki\checkmarkr\checkmarkc\checkmark$\checkmark0\checkmark$\checkmark^\checkmark\\checkmarkc\checkmarki\checkmarkr\checkmarkc\checkmark$\checkmark{\checkmark$\checkmark^\checkmark\\checkmarkc\checkmarki\checkmarkr\checkmarkc\checkmark$\checkmark,\checkmark$\checkmark^\checkmark\\checkmarkc\checkmarki\checkmarkr\checkmarkc\checkmark$\checkmark}\checkmark$\checkmark^\checkmark\\checkmarkc\checkmarki\checkmarkr\checkmarkc\checkmark$\checkmark5\checkmark$\checkmark^\checkmark\\checkmarkc\checkmarki\checkmarkr\checkmarkc\checkmark$\checkmark$\checkmark$\checkmark^\checkmark\\checkmarkc\checkmarki\checkmarkr\checkmarkc\checkmark$\checkmark \checkmark$\checkmark^\checkmark\\checkmarkc\checkmarki\checkmarkr\checkmarkc\checkmark$\checkmarkm\checkmark$\checkmark^\checkmark\\checkmarkc\checkmarki\checkmarkr\checkmarkc\checkmark$\checkmark/\checkmark$\checkmark^\checkmark\\checkmarkc\checkmarki\checkmarkr\checkmarkc\checkmark$\checkmarks\checkmark$\checkmark^\checkmark\\checkmarkc\checkmarki\checkmarkr\checkmarkc\checkmark$\checkmark²\checkmark$\checkmark^\checkmark\\checkmarkc\checkmarki\checkmarkr\checkmarkc\checkmark$\checkmark)\checkmark$\checkmark^\checkmark\\checkmarkc\checkmarki\checkmarkr\checkmarkc\checkmark$\checkmark \checkmark$\checkmark^\checkmark\\checkmarkc\checkmarki\checkmarkr\checkmarkc\checkmark$\checkmark:\checkmark$\checkmark^\checkmark\\checkmarkc\checkmarki\checkmarkr\checkmarkc\checkmark$\checkmark \checkmark$\checkmark^\checkmark\\checkmarkc\checkmarki\checkmarkr\checkmarkc\checkmark$\checkmark$\checkmark$\checkmark^\checkmark\\checkmarkc\checkmarki\checkmarkr\checkmarkc\checkmark$\checkmarkF\checkmark$\checkmark^\checkmark\\checkmarkc\checkmarki\checkmarkr\checkmarkc\checkmark$\checkmark_\checkmark$\checkmark^\checkmark\\checkmarkc\checkmarki\checkmarkr\checkmarkc\checkmark$\checkmark{\checkmark$\checkmark^\checkmark\\checkmarkc\checkmarki\checkmarkr\checkmarkc\checkmark$\checkmarki\checkmark$\checkmark^\checkmark\\checkmarkc\checkmarki\checkmarkr\checkmarkc\checkmark$\checkmarkn\checkmark$\checkmark^\checkmark\\checkmarkc\checkmarki\checkmarkr\checkmarkc\checkmark$\checkmarke\checkmark$\checkmark^\checkmark\\checkmarkc\checkmarki\checkmarkr\checkmarkc\checkmark$\checkmarkr\checkmark$\checkmark^\checkmark\\checkmarkc\checkmarki\checkmarkr\checkmarkc\checkmark$\checkmarkt\checkmark$\checkmark^\checkmark\\checkmarkc\checkmarki\checkmarkr\checkmarkc\checkmark$\checkmarki\checkmark$\checkmark^\checkmark\\checkmarkc\checkmarki\checkmarkr\checkmarkc\checkmark$\checkmarke\checkmark$\checkmark^\checkmark\\checkmarkc\checkmarki\checkmarkr\checkmarkc\checkmark$\checkmark}\checkmark$\checkmark^\checkmark\\checkmarkc\checkmarki\checkmarkr\checkmarkc\checkmark$\checkmark \checkmark$\checkmark^\checkmark\\checkmarkc\checkmarki\checkmarkr\checkmarkc\checkmark$\checkmark=\checkmark$\checkmark^\checkmark\\checkmarkc\checkmarki\checkmarkr\checkmarkc\checkmark$\checkmark \checkmark$\checkmark^\checkmark\\checkmarkc\checkmarki\checkmarkr\checkmarkc\checkmark$\checkmarkM\checkmark$\checkmark^\checkmark\\checkmarkc\checkmarki\checkmarkr\checkmarkc\checkmark$\checkmark \checkmark$\checkmark^\checkmark\\checkmarkc\checkmarki\checkmarkr\checkmarkc\checkmark$\checkmark\\checkmark$\checkmark^\checkmark\\checkmarkc\checkmarki\checkmarkr\checkmarkc\checkmark$\checkmarkt\checkmark$\checkmark^\checkmark\\checkmarkc\checkmarki\checkmarkr\checkmarkc\checkmark$\checkmarki\checkmark$\checkmark^\checkmark\\checkmarkc\checkmarki\checkmarkr\checkmarkc\checkmark$\checkmarkm\checkmark$\checkmark^\checkmark\\checkmarkc\checkmarki\checkmarkr\checkmarkc\checkmark$\checkmarke\checkmark$\checkmark^\checkmark\\checkmarkc\checkmarki\checkmarkr\checkmarkc\checkmark$\checkmarks\checkmark$\checkmark^\checkmark\\checkmarkc\checkmarki\checkmarkr\checkmarkc\checkmark$\checkmark \checkmark$\checkmark^\checkmark\\checkmarkc\checkmarki\checkmarkr\checkmarkc\checkmark$\checkmarka\checkmark$\checkmark^\checkmark\\checkmarkc\checkmarki\checkmarkr\checkmarkc\checkmark$\checkmark \checkmark$\checkmark^\checkmark\\checkmarkc\checkmarki\checkmarkr\checkmarkc\checkmark$\checkmark=\checkmark$\checkmark^\checkmark\\checkmarkc\checkmarki\checkmarkr\checkmarkc\checkmark$\checkmark \checkmark$\checkmark^\checkmark\\checkmarkc\checkmarki\checkmarkr\checkmarkc\checkmark$\checkmark1\checkmark$\checkmark^\checkmark\\checkmarkc\checkmarki\checkmarkr\checkmarkc\checkmark$\checkmark0\checkmark$\checkmark^\checkmark\\checkmarkc\checkmarki\checkmarkr\checkmarkc\checkmark$\checkmark \checkmark$\checkmark^\checkmark\\checkmarkc\checkmarki\checkmarkr\checkmarkc\checkmark$\checkmark\\checkmark$\checkmark^\checkmark\\checkmarkc\checkmarki\checkmarkr\checkmarkc\checkmark$\checkmarkt\checkmark$\checkmark^\checkmark\\checkmarkc\checkmarki\checkmarkr\checkmarkc\checkmark$\checkmarki\checkmark$\checkmark^\checkmark\\checkmarkc\checkmarki\checkmarkr\checkmarkc\checkmark$\checkmarkm\checkmark$\checkmark^\checkmark\\checkmarkc\checkmarki\checkmarkr\checkmarkc\checkmark$\checkmarke\checkmark$\checkmark^\checkmark\\checkmarkc\checkmarki\checkmarkr\checkmarkc\checkmark$\checkmarks\checkmark$\checkmark^\checkmark\\checkmarkc\checkmarki\checkmarkr\checkmarkc\checkmark$\checkmark \checkmark$\checkmark^\checkmark\\checkmarkc\checkmarki\checkmarkr\checkmarkc\checkmark$\checkmark0\checkmark$\checkmark^\checkmark\\checkmarkc\checkmarki\checkmarkr\checkmarkc\checkmark$\checkmark{\checkmark$\checkmark^\checkmark\\checkmarkc\checkmarki\checkmarkr\checkmarkc\checkmark$\checkmark,\checkmark$\checkmark^\checkmark\\checkmarkc\checkmarki\checkmarkr\checkmarkc\checkmark$\checkmark}\checkmark$\checkmark^\checkmark\\checkmarkc\checkmarki\checkmarkr\checkmarkc\checkmark$\checkmark5\checkmark$\checkmark^\checkmark\\checkmarkc\checkmarki\checkmarkr\checkmarkc\checkmark$\checkmark \checkmark$\checkmark^\checkmark\\checkmarkc\checkmarki\checkmarkr\checkmarkc\checkmark$\checkmark=\checkmark$\checkmark^\checkmark\\checkmarkc\checkmarki\checkmarkr\checkmarkc\checkmark$\checkmark \checkmark$\checkmark^\checkmark\\checkmarkc\checkmarki\checkmarkr\checkmarkc\checkmark$\checkmark5\checkmark$\checkmark^\checkmark\\checkmarkc\checkmarki\checkmarkr\checkmarkc\checkmark$\checkmark$\checkmark$\checkmark^\checkmark\\checkmarkc\checkmarki\checkmarkr\checkmarkc\checkmark$\checkmark \checkmark$\checkmark^\checkmark\\checkmarkc\checkmarki\checkmarkr\checkmarkc\checkmark$\checkmarkN\checkmark$\checkmark^\checkmark\\checkmarkc\checkmarki\checkmarkr\checkmarkc\checkmark$\checkmark
\checkmark$\checkmark^\checkmark\\checkmarkc\checkmarki\checkmarkr\checkmarkc\checkmark$\checkmark\\checkmark$\checkmark^\checkmark\\checkmarkc\checkmarki\checkmarkr\checkmarkc\checkmark$\checkmarke\checkmark$\checkmark^\checkmark\\checkmarkc\checkmarki\checkmarkr\checkmarkc\checkmark$\checkmarkn\checkmark$\checkmark^\checkmark\\checkmarkc\checkmarki\checkmarkr\checkmarkc\checkmark$\checkmarkd\checkmark$\checkmark^\checkmark\\checkmarkc\checkmarki\checkmarkr\checkmarkc\checkmark$\checkmark{\checkmark$\checkmark^\checkmark\\checkmarkc\checkmarki\checkmarkr\checkmarkc\checkmark$\checkmarki\checkmark$\checkmark^\checkmark\\checkmarkc\checkmarki\checkmarkr\checkmarkc\checkmark$\checkmarkt\checkmark$\checkmark^\checkmark\\checkmarkc\checkmarki\checkmarkr\checkmarkc\checkmark$\checkmarke\checkmark$\checkmark^\checkmark\\checkmarkc\checkmarki\checkmarkr\checkmarkc\checkmark$\checkmarkm\checkmark$\checkmark^\checkmark\\checkmarkc\checkmarki\checkmarkr\checkmarkc\checkmark$\checkmarki\checkmark$\checkmark^\checkmark\\checkmarkc\checkmarki\checkmarkr\checkmarkc\checkmark$\checkmarkz\checkmark$\checkmark^\checkmark\\checkmarkc\checkmarki\checkmarkr\checkmarkc\checkmark$\checkmarke\checkmark$\checkmark^\checkmark\\checkmarkc\checkmarki\checkmarkr\checkmarkc\checkmark$\checkmark}\checkmark$\checkmark^\checkmark\\checkmarkc\checkmarki\checkmarkr\checkmarkc\checkmark$\checkmark
\checkmark$\checkmark^\checkmark\\checkmarkc\checkmarki\checkmarkr\checkmarkc\checkmark$\checkmark
\checkmark$\checkmark^\checkmark\\checkmarkc\checkmarki\checkmarkr\checkmarkc\checkmark$\checkmarkL\checkmark$\checkmark^\checkmark\\checkmarkc\checkmarki\checkmarkr\checkmarkc\checkmark$\checkmarka\checkmark$\checkmark^\checkmark\\checkmarkc\checkmarki\checkmarkr\checkmarkc\checkmark$\checkmark \checkmark$\checkmark^\checkmark\\checkmarkc\checkmarki\checkmarkr\checkmarkc\checkmark$\checkmarkf\checkmark$\checkmark^\checkmark\\checkmarkc\checkmarki\checkmarkr\checkmarkc\checkmark$\checkmarko\checkmark$\checkmark^\checkmark\\checkmarkc\checkmarki\checkmarkr\checkmarkc\checkmark$\checkmarkr\checkmark$\checkmark^\checkmark\\checkmarkc\checkmarki\checkmarkr\checkmarkc\checkmark$\checkmarkc\checkmark$\checkmark^\checkmark\\checkmarkc\checkmarki\checkmarkr\checkmarkc\checkmark$\checkmarke\checkmark$\checkmark^\checkmark\\checkmarkc\checkmarki\checkmarkr\checkmarkc\checkmark$\checkmark \checkmark$\checkmark^\checkmark\\checkmarkc\checkmarki\checkmarkr\checkmarkc\checkmark$\checkmarka\checkmark$\checkmark^\checkmark\\checkmarkc\checkmarki\checkmarkr\checkmarkc\checkmark$\checkmarkx\checkmark$\checkmark^\checkmark\\checkmarkc\checkmarki\checkmarkr\checkmarkc\checkmark$\checkmarki\checkmark$\checkmark^\checkmark\\checkmarkc\checkmarki\checkmarkr\checkmarkc\checkmark$\checkmarka\checkmark$\checkmark^\checkmark\\checkmarkc\checkmarki\checkmarkr\checkmarkc\checkmark$\checkmarkl\checkmark$\checkmark^\checkmark\\checkmarkc\checkmarki\checkmarkr\checkmarkc\checkmark$\checkmarke\checkmark$\checkmark^\checkmark\\checkmarkc\checkmarki\checkmarkr\checkmarkc\checkmark$\checkmark \checkmark$\checkmark^\checkmark\\checkmarkc\checkmarki\checkmarkr\checkmarkc\checkmark$\checkmarkt\checkmark$\checkmark^\checkmark\\checkmarkc\checkmarki\checkmarkr\checkmarkc\checkmark$\checkmarko\checkmark$\checkmark^\checkmark\\checkmarkc\checkmarki\checkmarkr\checkmarkc\checkmark$\checkmarkt\checkmark$\checkmark^\checkmark\\checkmarkc\checkmarki\checkmarkr\checkmarkc\checkmark$\checkmarka\checkmark$\checkmark^\checkmark\\checkmarkc\checkmarki\checkmarkr\checkmarkc\checkmark$\checkmarkl\checkmark$\checkmark^\checkmark\\checkmarkc\checkmarki\checkmarkr\checkmarkc\checkmark$\checkmarke\checkmark$\checkmark^\checkmark\\checkmarkc\checkmarki\checkmarkr\checkmarkc\checkmark$\checkmark \checkmark$\checkmark^\checkmark\\checkmarkc\checkmarki\checkmarkr\checkmarkc\checkmark$\checkmarkd\checkmark$\checkmark^\checkmark\\checkmarkc\checkmarki\checkmarkr\checkmarkc\checkmark$\checkmarke\checkmark$\checkmark^\checkmark\\checkmarkc\checkmarki\checkmarkr\checkmarkc\checkmark$\checkmark \checkmark$\checkmark^\checkmark\\checkmarkc\checkmarki\checkmarkr\checkmarkc\checkmark$\checkmarkd\checkmark$\checkmark^\checkmark\\checkmarkc\checkmarki\checkmarkr\checkmarkc\checkmark$\checkmarki\checkmark$\checkmark^\checkmark\\checkmarkc\checkmarki\checkmarkr\checkmarkc\checkmark$\checkmarkm\checkmark$\checkmark^\checkmark\\checkmarkc\checkmarki\checkmarkr\checkmarkc\checkmark$\checkmarke\checkmark$\checkmark^\checkmark\\checkmarkc\checkmarki\checkmarkr\checkmarkc\checkmark$\checkmarkn\checkmark$\checkmark^\checkmark\\checkmarkc\checkmarki\checkmarkr\checkmarkc\checkmark$\checkmarks\checkmark$\checkmark^\checkmark\\checkmarkc\checkmarki\checkmarkr\checkmarkc\checkmark$\checkmarki\checkmark$\checkmark^\checkmark\\checkmarkc\checkmarki\checkmarkr\checkmarkc\checkmark$\checkmarko\checkmark$\checkmark^\checkmark\\checkmarkc\checkmarki\checkmarkr\checkmarkc\checkmark$\checkmarkn\checkmark$\checkmark^\checkmark\\checkmarkc\checkmarki\checkmarkr\checkmarkc\checkmark$\checkmarkn\checkmark$\checkmark^\checkmark\\checkmarkc\checkmarki\checkmarkr\checkmarkc\checkmark$\checkmarke\checkmark$\checkmark^\checkmark\\checkmarkc\checkmarki\checkmarkr\checkmarkc\checkmark$\checkmarkm\checkmark$\checkmark^\checkmark\\checkmarkc\checkmarki\checkmarkr\checkmarkc\checkmark$\checkmarke\checkmark$\checkmark^\checkmark\\checkmarkc\checkmarki\checkmarkr\checkmarkc\checkmark$\checkmarkn\checkmark$\checkmark^\checkmark\\checkmarkc\checkmarki\checkmarkr\checkmarkc\checkmark$\checkmarkt\checkmark$\checkmark^\checkmark\\checkmarkc\checkmarki\checkmarkr\checkmarkc\checkmark$\checkmark \checkmark$\checkmark^\checkmark\\checkmarkc\checkmarki\checkmarkr\checkmarkc\checkmark$\checkmarke\checkmark$\checkmark^\checkmark\\checkmarkc\checkmarki\checkmarkr\checkmarkc\checkmark$\checkmarks\checkmark$\checkmark^\checkmark\\checkmarkc\checkmarki\checkmarkr\checkmarkc\checkmark$\checkmarkt\checkmark$\checkmark^\checkmark\\checkmarkc\checkmarki\checkmarkr\checkmarkc\checkmark$\checkmark \checkmark$\checkmark^\checkmark\\checkmarkc\checkmarki\checkmarkr\checkmarkc\checkmark$\checkmark:\checkmark$\checkmark^\checkmark\\checkmarkc\checkmarki\checkmarkr\checkmarkc\checkmark$\checkmark
\checkmark$\checkmark^\checkmark\\checkmarkc\checkmarki\checkmarkr\checkmarkc\checkmark$\checkmark\\checkmark$\checkmark^\checkmark\\checkmarkc\checkmarki\checkmarkr\checkmarkc\checkmark$\checkmarkb\checkmark$\checkmark^\checkmark\\checkmarkc\checkmarki\checkmarkr\checkmarkc\checkmark$\checkmarke\checkmark$\checkmark^\checkmark\\checkmarkc\checkmarki\checkmarkr\checkmarkc\checkmark$\checkmarkg\checkmark$\checkmark^\checkmark\\checkmarkc\checkmarki\checkmarkr\checkmarkc\checkmark$\checkmarki\checkmark$\checkmark^\checkmark\\checkmarkc\checkmarki\checkmarkr\checkmarkc\checkmark$\checkmarkn\checkmark$\checkmark^\checkmark\\checkmarkc\checkmarki\checkmarkr\checkmarkc\checkmark$\checkmark{\checkmark$\checkmark^\checkmark\\checkmarkc\checkmarki\checkmarkr\checkmarkc\checkmark$\checkmarke\checkmark$\checkmark^\checkmark\\checkmarkc\checkmarki\checkmarkr\checkmarkc\checkmark$\checkmarkq\checkmark$\checkmark^\checkmark\\checkmarkc\checkmarki\checkmarkr\checkmarkc\checkmark$\checkmarku\checkmark$\checkmark^\checkmark\\checkmarkc\checkmarki\checkmarkr\checkmarkc\checkmark$\checkmarka\checkmark$\checkmark^\checkmark\\checkmarkc\checkmarki\checkmarkr\checkmarkc\checkmark$\checkmarkt\checkmark$\checkmark^\checkmark\\checkmarkc\checkmarki\checkmarkr\checkmarkc\checkmark$\checkmarki\checkmark$\checkmark^\checkmark\\checkmarkc\checkmarki\checkmarkr\checkmarkc\checkmark$\checkmarko\checkmark$\checkmark^\checkmark\\checkmarkc\checkmarki\checkmarkr\checkmarkc\checkmark$\checkmarkn\checkmark$\checkmark^\checkmark\\checkmarkc\checkmarki\checkmarkr\checkmarkc\checkmark$\checkmark}\checkmark$\checkmark^\checkmark\\checkmarkc\checkmarki\checkmarkr\checkmarkc\checkmark$\checkmark
\checkmark$\checkmark^\checkmark\\checkmarkc\checkmarki\checkmarkr\checkmarkc\checkmark$\checkmark \checkmark$\checkmark^\checkmark\\checkmarkc\checkmarki\checkmarkr\checkmarkc\checkmark$\checkmark \checkmark$\checkmark^\checkmark\\checkmarkc\checkmarki\checkmarkr\checkmarkc\checkmark$\checkmark \checkmark$\checkmark^\checkmark\\checkmarkc\checkmarki\checkmarkr\checkmarkc\checkmark$\checkmark \checkmark$\checkmark^\checkmark\\checkmarkc\checkmarki\checkmarkr\checkmarkc\checkmark$\checkmarkF\checkmark$\checkmark^\checkmark\\checkmarkc\checkmarki\checkmarkr\checkmarkc\checkmark$\checkmark_\checkmark$\checkmark^\checkmark\\checkmarkc\checkmarki\checkmarkr\checkmarkc\checkmark$\checkmark{\checkmark$\checkmark^\checkmark\\checkmarkc\checkmarki\checkmarkr\checkmarkc\checkmark$\checkmarka\checkmark$\checkmark^\checkmark\\checkmarkc\checkmarki\checkmarkr\checkmarkc\checkmark$\checkmarkx\checkmark$\checkmark^\checkmark\\checkmarkc\checkmarki\checkmarkr\checkmarkc\checkmark$\checkmark}\checkmark$\checkmark^\checkmark\\checkmarkc\checkmarki\checkmarkr\checkmarkc\checkmark$\checkmark \checkmark$\checkmark^\checkmark\\checkmarkc\checkmarki\checkmarkr\checkmarkc\checkmark$\checkmark=\checkmark$\checkmark^\checkmark\\checkmarkc\checkmarki\checkmarkr\checkmarkc\checkmark$\checkmark \checkmark$\checkmark^\checkmark\\checkmarkc\checkmarki\checkmarkr\checkmarkc\checkmark$\checkmarkF\checkmark$\checkmark^\checkmark\\checkmarkc\checkmarki\checkmarkr\checkmarkc\checkmark$\checkmark_\checkmark$\checkmark^\checkmark\\checkmarkc\checkmarki\checkmarkr\checkmarkc\checkmark$\checkmark{\checkmark$\checkmark^\checkmark\\checkmarkc\checkmarki\checkmarkr\checkmarkc\checkmark$\checkmarku\checkmark$\checkmark^\checkmark\\checkmarkc\checkmarki\checkmarkr\checkmarkc\checkmark$\checkmarks\checkmark$\checkmark^\checkmark\\checkmarkc\checkmarki\checkmarkr\checkmarkc\checkmark$\checkmarki\checkmark$\checkmark^\checkmark\\checkmarkc\checkmarki\checkmarkr\checkmarkc\checkmark$\checkmarkn\checkmark$\checkmark^\checkmark\\checkmarkc\checkmarki\checkmarkr\checkmarkc\checkmark$\checkmarka\checkmark$\checkmark^\checkmark\\checkmarkc\checkmarki\checkmarkr\checkmarkc\checkmark$\checkmarkg\checkmark$\checkmark^\checkmark\\checkmarkc\checkmarki\checkmarkr\checkmarkc\checkmark$\checkmarke\checkmark$\checkmark^\checkmark\\checkmarkc\checkmarki\checkmarkr\checkmarkc\checkmark$\checkmark}\checkmark$\checkmark^\checkmark\\checkmarkc\checkmarki\checkmarkr\checkmarkc\checkmark$\checkmark \checkmark$\checkmark^\checkmark\\checkmarkc\checkmarki\checkmarkr\checkmarkc\checkmark$\checkmark+\checkmark$\checkmark^\checkmark\\checkmarkc\checkmarki\checkmarkr\checkmarkc\checkmark$\checkmark \checkmark$\checkmark^\checkmark\\checkmarkc\checkmarki\checkmarkr\checkmarkc\checkmark$\checkmarkF\checkmark$\checkmark^\checkmark\\checkmarkc\checkmarki\checkmarkr\checkmarkc\checkmark$\checkmark_\checkmark$\checkmark^\checkmark\\checkmarkc\checkmarki\checkmarkr\checkmarkc\checkmark$\checkmark{\checkmark$\checkmark^\checkmark\\checkmarkc\checkmarki\checkmarkr\checkmarkc\checkmark$\checkmarkf\checkmark$\checkmark^\checkmark\\checkmarkc\checkmarki\checkmarkr\checkmarkc\checkmark$\checkmarkr\checkmark$\checkmark^\checkmark\\checkmarkc\checkmarki\checkmarkr\checkmarkc\checkmark$\checkmarko\checkmark$\checkmark^\checkmark\\checkmarkc\checkmarki\checkmarkr\checkmarkc\checkmark$\checkmarkt\checkmark$\checkmark^\checkmark\\checkmarkc\checkmarki\checkmarkr\checkmarkc\checkmark$\checkmarkt\checkmark$\checkmark^\checkmark\\checkmarkc\checkmarki\checkmarkr\checkmarkc\checkmark$\checkmark}\checkmark$\checkmark^\checkmark\\checkmarkc\checkmarki\checkmarkr\checkmarkc\checkmark$\checkmark \checkmark$\checkmark^\checkmark\\checkmarkc\checkmarki\checkmarkr\checkmarkc\checkmark$\checkmark+\checkmark$\checkmark^\checkmark\\checkmarkc\checkmarki\checkmarkr\checkmarkc\checkmark$\checkmark \checkmark$\checkmark^\checkmark\\checkmarkc\checkmarki\checkmarkr\checkmarkc\checkmark$\checkmarkF\checkmark$\checkmark^\checkmark\\checkmarkc\checkmarki\checkmarkr\checkmarkc\checkmark$\checkmark_\checkmark$\checkmark^\checkmark\\checkmarkc\checkmarki\checkmarkr\checkmarkc\checkmark$\checkmark{\checkmark$\checkmark^\checkmark\\checkmarkc\checkmarki\checkmarkr\checkmarkc\checkmark$\checkmarki\checkmark$\checkmark^\checkmark\\checkmarkc\checkmarki\checkmarkr\checkmarkc\checkmark$\checkmarkn\checkmark$\checkmark^\checkmark\\checkmarkc\checkmarki\checkmarkr\checkmarkc\checkmark$\checkmarke\checkmark$\checkmark^\checkmark\\checkmarkc\checkmarki\checkmarkr\checkmarkc\checkmark$\checkmarkr\checkmark$\checkmark^\checkmark\\checkmarkc\checkmarki\checkmarkr\checkmarkc\checkmark$\checkmarkt\checkmark$\checkmark^\checkmark\\checkmarkc\checkmarki\checkmarkr\checkmarkc\checkmark$\checkmarki\checkmark$\checkmark^\checkmark\\checkmarkc\checkmarki\checkmarkr\checkmarkc\checkmark$\checkmarke\checkmark$\checkmark^\checkmark\\checkmarkc\checkmarki\checkmarkr\checkmarkc\checkmark$\checkmark}\checkmark$\checkmark^\checkmark\\checkmarkc\checkmarki\checkmarkr\checkmarkc\checkmark$\checkmark \checkmark$\checkmark^\checkmark\\checkmarkc\checkmarki\checkmarkr\checkmarkc\checkmark$\checkmark=\checkmark$\checkmark^\checkmark\\checkmarkc\checkmarki\checkmarkr\checkmarkc\checkmark$\checkmark \checkmark$\checkmark^\checkmark\\checkmarkc\checkmarki\checkmarkr\checkmarkc\checkmark$\checkmark1\checkmark$\checkmark^\checkmark\\checkmarkc\checkmarki\checkmarkr\checkmarkc\checkmark$\checkmark0\checkmark$\checkmark^\checkmark\\checkmarkc\checkmarki\checkmarkr\checkmarkc\checkmark$\checkmark0\checkmark$\checkmark^\checkmark\\checkmarkc\checkmarki\checkmarkr\checkmarkc\checkmark$\checkmark \checkmark$\checkmark^\checkmark\\checkmarkc\checkmarki\checkmarkr\checkmarkc\checkmark$\checkmark+\checkmark$\checkmark^\checkmark\\checkmarkc\checkmarki\checkmarkr\checkmarkc\checkmark$\checkmark \checkmark$\checkmark^\checkmark\\checkmarkc\checkmarki\checkmarkr\checkmarkc\checkmark$\checkmark0\checkmark$\checkmark^\checkmark\\checkmarkc\checkmarki\checkmarkr\checkmarkc\checkmark$\checkmark{\checkmark$\checkmark^\checkmark\\checkmarkc\checkmarki\checkmarkr\checkmarkc\checkmark$\checkmark,\checkmark$\checkmark^\checkmark\\checkmarkc\checkmarki\checkmarkr\checkmarkc\checkmark$\checkmark}\checkmark$\checkmark^\checkmark\\checkmarkc\checkmarki\checkmarkr\checkmarkc\checkmark$\checkmark5\checkmark$\checkmark^\checkmark\\checkmarkc\checkmarki\checkmarkr\checkmarkc\checkmark$\checkmark \checkmark$\checkmark^\checkmark\\checkmarkc\checkmarki\checkmarkr\checkmarkc\checkmark$\checkmark+\checkmark$\checkmark^\checkmark\\checkmarkc\checkmarki\checkmarkr\checkmarkc\checkmark$\checkmark \checkmark$\checkmark^\checkmark\\checkmarkc\checkmarki\checkmarkr\checkmarkc\checkmark$\checkmark5\checkmark$\checkmark^\checkmark\\checkmarkc\checkmarki\checkmarkr\checkmarkc\checkmark$\checkmark \checkmark$\checkmark^\checkmark\\checkmarkc\checkmarki\checkmarkr\checkmarkc\checkmark$\checkmark=\checkmark$\checkmark^\checkmark\\checkmarkc\checkmarki\checkmarkr\checkmarkc\checkmark$\checkmark \checkmark$\checkmark^\checkmark\\checkmarkc\checkmarki\checkmarkr\checkmarkc\checkmark$\checkmark\\checkmark$\checkmark^\checkmark\\checkmarkc\checkmarki\checkmarkr\checkmarkc\checkmark$\checkmarkb\checkmark$\checkmark^\checkmark\\checkmarkc\checkmarki\checkmarkr\checkmarkc\checkmark$\checkmarko\checkmark$\checkmark^\checkmark\\checkmarkc\checkmarki\checkmarkr\checkmarkc\checkmark$\checkmarkx\checkmark$\checkmark^\checkmark\\checkmarkc\checkmarki\checkmarkr\checkmarkc\checkmark$\checkmarke\checkmark$\checkmark^\checkmark\\checkmarkc\checkmarki\checkmarkr\checkmarkc\checkmark$\checkmarkd\checkmark$\checkmark^\checkmark\\checkmarkc\checkmarki\checkmarkr\checkmarkc\checkmark$\checkmark{\checkmark$\checkmark^\checkmark\\checkmarkc\checkmarki\checkmarkr\checkmarkc\checkmark$\checkmark1\checkmark$\checkmark^\checkmark\\checkmarkc\checkmarki\checkmarkr\checkmarkc\checkmark$\checkmark0\checkmark$\checkmark^\checkmark\\checkmarkc\checkmarki\checkmarkr\checkmarkc\checkmark$\checkmark5\checkmark$\checkmark^\checkmark\\checkmarkc\checkmarki\checkmarkr\checkmarkc\checkmark$\checkmark{\checkmark$\checkmark^\checkmark\\checkmarkc\checkmarki\checkmarkr\checkmarkc\checkmark$\checkmark,\checkmark$\checkmark^\checkmark\\checkmarkc\checkmarki\checkmarkr\checkmarkc\checkmark$\checkmark}\checkmark$\checkmark^\checkmark\\checkmarkc\checkmarki\checkmarkr\checkmarkc\checkmark$\checkmark5\checkmark$\checkmark^\checkmark\\checkmarkc\checkmarki\checkmarkr\checkmarkc\checkmark$\checkmark \checkmark$\checkmark^\checkmark\\checkmarkc\checkmarki\checkmarkr\checkmarkc\checkmark$\checkmark\\checkmark$\checkmark^\checkmark\\checkmarkc\checkmarki\checkmarkr\checkmarkc\checkmark$\checkmarkt\checkmark$\checkmark^\checkmark\\checkmarkc\checkmarki\checkmarkr\checkmarkc\checkmark$\checkmarke\checkmark$\checkmark^\checkmark\\checkmarkc\checkmarki\checkmarkr\checkmarkc\checkmark$\checkmarkx\checkmark$\checkmark^\checkmark\\checkmarkc\checkmarki\checkmarkr\checkmarkc\checkmark$\checkmarkt\checkmark$\checkmark^\checkmark\\checkmarkc\checkmarki\checkmarkr\checkmarkc\checkmark$\checkmark{\checkmark$\checkmark^\checkmark\\checkmarkc\checkmarki\checkmarkr\checkmarkc\checkmark$\checkmark \checkmark$\checkmark^\checkmark\\checkmarkc\checkmarki\checkmarkr\checkmarkc\checkmark$\checkmarkN\checkmark$\checkmark^\checkmark\\checkmarkc\checkmarki\checkmarkr\checkmarkc\checkmark$\checkmark}\checkmark$\checkmark^\checkmark\\checkmarkc\checkmarki\checkmarkr\checkmarkc\checkmark$\checkmark}\checkmark$\checkmark^\checkmark\\checkmarkc\checkmarki\checkmarkr\checkmarkc\checkmark$\checkmark
\checkmark$\checkmark^\checkmark\\checkmarkc\checkmarki\checkmarkr\checkmarkc\checkmark$\checkmark \checkmark$\checkmark^\checkmark\\checkmarkc\checkmarki\checkmarkr\checkmarkc\checkmark$\checkmark \checkmark$\checkmark^\checkmark\\checkmarkc\checkmarki\checkmarkr\checkmarkc\checkmark$\checkmark \checkmark$\checkmark^\checkmark\\checkmarkc\checkmarki\checkmarkr\checkmarkc\checkmark$\checkmark \checkmark$\checkmark^\checkmark\\checkmarkc\checkmarki\checkmarkr\checkmarkc\checkmark$\checkmark\\checkmark$\checkmark^\checkmark\\checkmarkc\checkmarki\checkmarkr\checkmarkc\checkmark$\checkmarkl\checkmark$\checkmark^\checkmark\\checkmarkc\checkmarki\checkmarkr\checkmarkc\checkmark$\checkmarka\checkmark$\checkmark^\checkmark\\checkmarkc\checkmarki\checkmarkr\checkmarkc\checkmark$\checkmarkb\checkmark$\checkmark^\checkmark\\checkmarkc\checkmarki\checkmarkr\checkmarkc\checkmark$\checkmarke\checkmark$\checkmark^\checkmark\\checkmarkc\checkmarki\checkmarkr\checkmarkc\checkmark$\checkmarkl\checkmark$\checkmark^\checkmark\\checkmarkc\checkmarki\checkmarkr\checkmarkc\checkmark$\checkmark{\checkmark$\checkmark^\checkmark\\checkmarkc\checkmarki\checkmarkr\checkmarkc\checkmark$\checkmarke\checkmark$\checkmark^\checkmark\\checkmarkc\checkmarki\checkmarkr\checkmarkc\checkmark$\checkmarkq\checkmark$\checkmark^\checkmark\\checkmarkc\checkmarki\checkmarkr\checkmarkc\checkmark$\checkmark:\checkmark$\checkmark^\checkmark\\checkmarkc\checkmarki\checkmarkr\checkmarkc\checkmark$\checkmarkf\checkmark$\checkmark^\checkmark\\checkmarkc\checkmarki\checkmarkr\checkmarkc\checkmark$\checkmarko\checkmark$\checkmark^\checkmark\\checkmarkc\checkmarki\checkmarkr\checkmarkc\checkmark$\checkmarkr\checkmark$\checkmark^\checkmark\\checkmarkc\checkmarki\checkmarkr\checkmarkc\checkmark$\checkmarkc\checkmark$\checkmark^\checkmark\\checkmarkc\checkmarki\checkmarkr\checkmarkc\checkmark$\checkmarke\checkmark$\checkmark^\checkmark\\checkmarkc\checkmarki\checkmarkr\checkmarkc\checkmark$\checkmark_\checkmark$\checkmark^\checkmark\\checkmarkc\checkmarki\checkmarkr\checkmarkc\checkmark$\checkmarka\checkmark$\checkmark^\checkmark\\checkmarkc\checkmarki\checkmarkr\checkmarkc\checkmark$\checkmarkx\checkmark$\checkmark^\checkmark\\checkmarkc\checkmarki\checkmarkr\checkmarkc\checkmark$\checkmarki\checkmark$\checkmark^\checkmark\\checkmarkc\checkmarki\checkmarkr\checkmarkc\checkmark$\checkmarka\checkmark$\checkmark^\checkmark\\checkmarkc\checkmarki\checkmarkr\checkmarkc\checkmark$\checkmarkl\checkmark$\checkmark^\checkmark\\checkmarkc\checkmarki\checkmarkr\checkmarkc\checkmark$\checkmarke\checkmark$\checkmark^\checkmark\\checkmarkc\checkmarki\checkmarkr\checkmarkc\checkmark$\checkmark}\checkmark$\checkmark^\checkmark\\checkmarkc\checkmarki\checkmarkr\checkmarkc\checkmark$\checkmark
\checkmark$\checkmark^\checkmark\\checkmarkc\checkmarki\checkmarkr\checkmarkc\checkmark$\checkmark\\checkmark$\checkmark^\checkmark\\checkmarkc\checkmarki\checkmarkr\checkmarkc\checkmark$\checkmarke\checkmark$\checkmark^\checkmark\\checkmarkc\checkmarki\checkmarkr\checkmarkc\checkmark$\checkmarkn\checkmark$\checkmark^\checkmark\\checkmarkc\checkmarki\checkmarkr\checkmarkc\checkmark$\checkmarkd\checkmark$\checkmark^\checkmark\\checkmarkc\checkmarki\checkmarkr\checkmarkc\checkmark$\checkmark{\checkmark$\checkmark^\checkmark\\checkmarkc\checkmarki\checkmarkr\checkmarkc\checkmark$\checkmarke\checkmark$\checkmark^\checkmark\\checkmarkc\checkmarki\checkmarkr\checkmarkc\checkmark$\checkmarkq\checkmark$\checkmark^\checkmark\\checkmarkc\checkmarki\checkmarkr\checkmarkc\checkmark$\checkmarku\checkmark$\checkmark^\checkmark\\checkmarkc\checkmarki\checkmarkr\checkmarkc\checkmark$\checkmarka\checkmark$\checkmark^\checkmark\\checkmarkc\checkmarki\checkmarkr\checkmarkc\checkmark$\checkmarkt\checkmark$\checkmark^\checkmark\\checkmarkc\checkmarki\checkmarkr\checkmarkc\checkmark$\checkmarki\checkmark$\checkmark^\checkmark\\checkmarkc\checkmarki\checkmarkr\checkmarkc\checkmark$\checkmarko\checkmark$\checkmark^\checkmark\\checkmarkc\checkmarki\checkmarkr\checkmarkc\checkmark$\checkmarkn\checkmark$\checkmark^\checkmark\\checkmarkc\checkmarki\checkmarkr\checkmarkc\checkmark$\checkmark}\checkmark$\checkmark^\checkmark\\checkmarkc\checkmarki\checkmarkr\checkmarkc\checkmark$\checkmark
\checkmark$\checkmark^\checkmark\\checkmarkc\checkmarki\checkmarkr\checkmarkc\checkmark$\checkmark
\checkmark$\checkmark^\checkmark\\checkmarkc\checkmarki\checkmarkr\checkmarkc\checkmark$\checkmark\\checkmark$\checkmark^\checkmark\\checkmarkc\checkmarki\checkmarkr\checkmarkc\checkmark$\checkmarkt\checkmark$\checkmark^\checkmark\\checkmarkc\checkmarki\checkmarkr\checkmarkc\checkmark$\checkmarke\checkmark$\checkmark^\checkmark\\checkmarkc\checkmarki\checkmarkr\checkmarkc\checkmark$\checkmarkx\checkmark$\checkmark^\checkmark\\checkmarkc\checkmarki\checkmarkr\checkmarkc\checkmark$\checkmarkt\checkmark$\checkmark^\checkmark\\checkmarkc\checkmarki\checkmarkr\checkmarkc\checkmark$\checkmarkb\checkmark$\checkmark^\checkmark\\checkmarkc\checkmarki\checkmarkr\checkmarkc\checkmark$\checkmarkf\checkmark$\checkmark^\checkmark\\checkmarkc\checkmarki\checkmarkr\checkmarkc\checkmark$\checkmark{\checkmark$\checkmark^\checkmark\\checkmarkc\checkmarki\checkmarkr\checkmarkc\checkmark$\checkmarkN\checkmark$\checkmark^\checkmark\\checkmarkc\checkmarki\checkmarkr\checkmarkc\checkmark$\checkmarko\checkmark$\checkmark^\checkmark\\checkmarkc\checkmarki\checkmarkr\checkmarkc\checkmark$\checkmarkm\checkmark$\checkmark^\checkmark\\checkmarkc\checkmarki\checkmarkr\checkmarkc\checkmark$\checkmarke\checkmark$\checkmark^\checkmark\\checkmarkc\checkmarki\checkmarkr\checkmarkc\checkmark$\checkmarkn\checkmark$\checkmark^\checkmark\\checkmarkc\checkmarki\checkmarkr\checkmarkc\checkmark$\checkmarkc\checkmark$\checkmark^\checkmark\\checkmarkc\checkmarki\checkmarkr\checkmarkc\checkmark$\checkmarkl\checkmark$\checkmark^\checkmark\\checkmarkc\checkmarki\checkmarkr\checkmarkc\checkmark$\checkmarka\checkmark$\checkmark^\checkmark\\checkmarkc\checkmarki\checkmarkr\checkmarkc\checkmark$\checkmarkt\checkmark$\checkmark^\checkmark\\checkmarkc\checkmarki\checkmarkr\checkmarkc\checkmark$\checkmarku\checkmark$\checkmark^\checkmark\\checkmarkc\checkmarki\checkmarkr\checkmarkc\checkmark$\checkmarkr\checkmark$\checkmark^\checkmark\\checkmarkc\checkmarki\checkmarkr\checkmarkc\checkmark$\checkmarke\checkmark$\checkmark^\checkmark\\checkmarkc\checkmarki\checkmarkr\checkmarkc\checkmark$\checkmark \checkmark$\checkmark^\checkmark\\checkmarkc\checkmarki\checkmarkr\checkmarkc\checkmark$\checkmarkd\checkmark$\checkmark^\checkmark\\checkmarkc\checkmarki\checkmarkr\checkmarkc\checkmark$\checkmarke\checkmark$\checkmark^\checkmark\\checkmarkc\checkmarki\checkmarkr\checkmarkc\checkmark$\checkmarks\checkmark$\checkmark^\checkmark\\checkmarkc\checkmarki\checkmarkr\checkmarkc\checkmark$\checkmark \checkmark$\checkmark^\checkmark\\checkmarkc\checkmarki\checkmarkr\checkmarkc\checkmark$\checkmarke\checkmark$\checkmark^\checkmark\\checkmarkc\checkmarki\checkmarkr\checkmarkc\checkmark$\checkmarkf\checkmark$\checkmark^\checkmark\\checkmarkc\checkmarki\checkmarkr\checkmarkc\checkmark$\checkmarkf\checkmark$\checkmark^\checkmark\\checkmarkc\checkmarki\checkmarkr\checkmarkc\checkmark$\checkmarko\checkmark$\checkmark^\checkmark\\checkmarkc\checkmarki\checkmarkr\checkmarkc\checkmark$\checkmarkr\checkmark$\checkmark^\checkmark\\checkmarkc\checkmarki\checkmarkr\checkmarkc\checkmark$\checkmarkt\checkmark$\checkmark^\checkmark\\checkmarkc\checkmarki\checkmarkr\checkmarkc\checkmark$\checkmarks\checkmark$\checkmark^\checkmark\\checkmarkc\checkmarki\checkmarkr\checkmarkc\checkmark$\checkmark \checkmark$\checkmark^\checkmark\\checkmarkc\checkmarki\checkmarkr\checkmarkc\checkmark$\checkmark:\checkmark$\checkmark^\checkmark\\checkmarkc\checkmarki\checkmarkr\checkmarkc\checkmark$\checkmark}\checkmark$\checkmark^\checkmark\\checkmarkc\checkmarki\checkmarkr\checkmarkc\checkmark$\checkmark
\checkmark$\checkmark^\checkmark\\checkmarkc\checkmarki\checkmarkr\checkmarkc\checkmark$\checkmark\\checkmark$\checkmark^\checkmark\\checkmarkc\checkmarki\checkmarkr\checkmarkc\checkmark$\checkmarkb\checkmark$\checkmark^\checkmark\\checkmarkc\checkmarki\checkmarkr\checkmarkc\checkmark$\checkmarke\checkmark$\checkmark^\checkmark\\checkmarkc\checkmarki\checkmarkr\checkmarkc\checkmark$\checkmarkg\checkmark$\checkmark^\checkmark\\checkmarkc\checkmarki\checkmarkr\checkmarkc\checkmark$\checkmarki\checkmark$\checkmark^\checkmark\\checkmarkc\checkmarki\checkmarkr\checkmarkc\checkmark$\checkmarkn\checkmark$\checkmark^\checkmark\\checkmarkc\checkmarki\checkmarkr\checkmarkc\checkmark$\checkmark{\checkmark$\checkmark^\checkmark\\checkmarkc\checkmarki\checkmarkr\checkmarkc\checkmark$\checkmarkt\checkmark$\checkmark^\checkmark\\checkmarkc\checkmarki\checkmarkr\checkmarkc\checkmark$\checkmarka\checkmark$\checkmark^\checkmark\\checkmarkc\checkmarki\checkmarkr\checkmarkc\checkmark$\checkmarkb\checkmark$\checkmark^\checkmark\\checkmarkc\checkmarki\checkmarkr\checkmarkc\checkmark$\checkmarkl\checkmark$\checkmark^\checkmark\\checkmarkc\checkmarki\checkmarkr\checkmarkc\checkmark$\checkmarke\checkmark$\checkmark^\checkmark\\checkmarkc\checkmarki\checkmarkr\checkmarkc\checkmark$\checkmark}\checkmark$\checkmark^\checkmark\\checkmarkc\checkmarki\checkmarkr\checkmarkc\checkmark$\checkmark[\checkmark$\checkmark^\checkmark\\checkmarkc\checkmarki\checkmarkr\checkmarkc\checkmark$\checkmarkh\checkmark$\checkmark^\checkmark\\checkmarkc\checkmarki\checkmarkr\checkmarkc\checkmark$\checkmarkt\checkmark$\checkmark^\checkmark\\checkmarkc\checkmarki\checkmarkr\checkmarkc\checkmark$\checkmarkb\checkmark$\checkmark^\checkmark\\checkmarkc\checkmarki\checkmarkr\checkmarkc\checkmark$\checkmarkp\checkmark$\checkmark^\checkmark\\checkmarkc\checkmarki\checkmarkr\checkmarkc\checkmark$\checkmark]\checkmark$\checkmark^\checkmark\\checkmarkc\checkmarki\checkmarkr\checkmarkc\checkmark$\checkmark
\checkmark$\checkmark^\checkmark\\checkmarkc\checkmarki\checkmarkr\checkmarkc\checkmark$\checkmark\\checkmark$\checkmark^\checkmark\\checkmarkc\checkmarki\checkmarkr\checkmarkc\checkmark$\checkmarkc\checkmark$\checkmark^\checkmark\\checkmarkc\checkmarki\checkmarkr\checkmarkc\checkmark$\checkmarke\checkmark$\checkmark^\checkmark\\checkmarkc\checkmarki\checkmarkr\checkmarkc\checkmark$\checkmarkn\checkmark$\checkmark^\checkmark\\checkmarkc\checkmarki\checkmarkr\checkmarkc\checkmark$\checkmarkt\checkmark$\checkmark^\checkmark\\checkmarkc\checkmarki\checkmarkr\checkmarkc\checkmark$\checkmarke\checkmark$\checkmark^\checkmark\\checkmarkc\checkmarki\checkmarkr\checkmarkc\checkmark$\checkmarkr\checkmark$\checkmark^\checkmark\\checkmarkc\checkmarki\checkmarkr\checkmarkc\checkmark$\checkmarki\checkmark$\checkmark^\checkmark\\checkmarkc\checkmarki\checkmarkr\checkmarkc\checkmark$\checkmarkn\checkmark$\checkmark^\checkmark\\checkmarkc\checkmarki\checkmarkr\checkmarkc\checkmark$\checkmarkg\checkmark$\checkmark^\checkmark\\checkmarkc\checkmarki\checkmarkr\checkmarkc\checkmark$\checkmark
\checkmark$\checkmark^\checkmark\\checkmarkc\checkmarki\checkmarkr\checkmarkc\checkmark$\checkmark\\checkmark$\checkmark^\checkmark\\checkmarkc\checkmarki\checkmarkr\checkmarkc\checkmark$\checkmarkb\checkmark$\checkmark^\checkmark\\checkmarkc\checkmarki\checkmarkr\checkmarkc\checkmark$\checkmarke\checkmark$\checkmark^\checkmark\\checkmarkc\checkmarki\checkmarkr\checkmarkc\checkmark$\checkmarkg\checkmark$\checkmark^\checkmark\\checkmarkc\checkmarki\checkmarkr\checkmarkc\checkmark$\checkmarki\checkmark$\checkmark^\checkmark\\checkmarkc\checkmarki\checkmarkr\checkmarkc\checkmark$\checkmarkn\checkmark$\checkmark^\checkmark\\checkmarkc\checkmarki\checkmarkr\checkmarkc\checkmark$\checkmark{\checkmark$\checkmark^\checkmark\\checkmarkc\checkmarki\checkmarkr\checkmarkc\checkmark$\checkmarkt\checkmark$\checkmark^\checkmark\\checkmarkc\checkmarki\checkmarkr\checkmarkc\checkmark$\checkmarka\checkmark$\checkmark^\checkmark\\checkmarkc\checkmarki\checkmarkr\checkmarkc\checkmark$\checkmarkb\checkmark$\checkmark^\checkmark\\checkmarkc\checkmarki\checkmarkr\checkmarkc\checkmark$\checkmarku\checkmark$\checkmark^\checkmark\\checkmarkc\checkmarki\checkmarkr\checkmarkc\checkmark$\checkmarkl\checkmark$\checkmark^\checkmark\\checkmarkc\checkmarki\checkmarkr\checkmarkc\checkmark$\checkmarka\checkmark$\checkmark^\checkmark\\checkmarkc\checkmarki\checkmarkr\checkmarkc\checkmark$\checkmarkr\checkmark$\checkmark^\checkmark\\checkmarkc\checkmarki\checkmarkr\checkmarkc\checkmark$\checkmark}\checkmark$\checkmark^\checkmark\\checkmarkc\checkmarki\checkmarkr\checkmarkc\checkmark$\checkmark{\checkmark$\checkmark^\checkmark\\checkmarkc\checkmarki\checkmarkr\checkmarkc\checkmark$\checkmark|\checkmark$\checkmark^\checkmark\\checkmarkc\checkmarki\checkmarkr\checkmarkc\checkmark$\checkmarkc\checkmark$\checkmark^\checkmark\\checkmarkc\checkmarki\checkmarkr\checkmarkc\checkmark$\checkmark|\checkmark$\checkmark^\checkmark\\checkmarkc\checkmarki\checkmarkr\checkmarkc\checkmark$\checkmarkl\checkmark$\checkmark^\checkmark\\checkmarkc\checkmarki\checkmarkr\checkmarkc\checkmark$\checkmark|\checkmark$\checkmark^\checkmark\\checkmarkc\checkmarki\checkmarkr\checkmarkc\checkmark$\checkmarkc\checkmark$\checkmark^\checkmark\\checkmarkc\checkmarki\checkmarkr\checkmarkc\checkmark$\checkmark|\checkmark$\checkmark^\checkmark\\checkmarkc\checkmarki\checkmarkr\checkmarkc\checkmark$\checkmark}\checkmark$\checkmark^\checkmark\\checkmarkc\checkmarki\checkmarkr\checkmarkc\checkmark$\checkmark
\checkmark$\checkmark^\checkmark\\checkmarkc\checkmarki\checkmarkr\checkmarkc\checkmark$\checkmark\\checkmark$\checkmark^\checkmark\\checkmarkc\checkmarki\checkmarkr\checkmarkc\checkmark$\checkmarkh\checkmark$\checkmark^\checkmark\\checkmarkc\checkmarki\checkmarkr\checkmarkc\checkmark$\checkmarkl\checkmark$\checkmark^\checkmark\\checkmarkc\checkmarki\checkmarkr\checkmarkc\checkmark$\checkmarki\checkmark$\checkmark^\checkmark\\checkmarkc\checkmarki\checkmarkr\checkmarkc\checkmark$\checkmarkn\checkmark$\checkmark^\checkmark\\checkmarkc\checkmarki\checkmarkr\checkmarkc\checkmark$\checkmarke\checkmark$\checkmark^\checkmark\\checkmarkc\checkmarki\checkmarkr\checkmarkc\checkmark$\checkmark
\checkmark$\checkmark^\checkmark\\checkmarkc\checkmarki\checkmarkr\checkmarkc\checkmark$\checkmark\\checkmark$\checkmark^\checkmark\\checkmarkc\checkmarki\checkmarkr\checkmarkc\checkmark$\checkmarkt\checkmark$\checkmark^\checkmark\\checkmarkc\checkmarki\checkmarkr\checkmarkc\checkmark$\checkmarke\checkmark$\checkmark^\checkmark\\checkmarkc\checkmarki\checkmarkr\checkmarkc\checkmark$\checkmarkx\checkmark$\checkmark^\checkmark\\checkmarkc\checkmarki\checkmarkr\checkmarkc\checkmark$\checkmarkt\checkmark$\checkmark^\checkmark\\checkmarkc\checkmarki\checkmarkr\checkmarkc\checkmark$\checkmarkb\checkmark$\checkmark^\checkmark\\checkmarkc\checkmarki\checkmarkr\checkmarkc\checkmark$\checkmarkf\checkmark$\checkmark^\checkmark\\checkmarkc\checkmarki\checkmarkr\checkmarkc\checkmark$\checkmark{\checkmark$\checkmark^\checkmark\\checkmarkc\checkmarki\checkmarkr\checkmarkc\checkmark$\checkmarkS\checkmark$\checkmark^\checkmark\\checkmarkc\checkmarki\checkmarkr\checkmarkc\checkmark$\checkmarky\checkmark$\checkmark^\checkmark\\checkmarkc\checkmarki\checkmarkr\checkmarkc\checkmark$\checkmarkm\checkmark$\checkmark^\checkmark\\checkmarkc\checkmarki\checkmarkr\checkmarkc\checkmark$\checkmarkb\checkmark$\checkmark^\checkmark\\checkmarkc\checkmarki\checkmarkr\checkmarkc\checkmark$\checkmarko\checkmark$\checkmark^\checkmark\\checkmarkc\checkmarki\checkmarkr\checkmarkc\checkmark$\checkmarkl\checkmark$\checkmark^\checkmark\\checkmarkc\checkmarki\checkmarkr\checkmarkc\checkmark$\checkmarke\checkmark$\checkmark^\checkmark\\checkmarkc\checkmarki\checkmarkr\checkmarkc\checkmark$\checkmark}\checkmark$\checkmark^\checkmark\\checkmarkc\checkmarki\checkmarkr\checkmarkc\checkmark$\checkmark \checkmark$\checkmark^\checkmark\\checkmarkc\checkmarki\checkmarkr\checkmarkc\checkmark$\checkmark&\checkmark$\checkmark^\checkmark\\checkmarkc\checkmarki\checkmarkr\checkmarkc\checkmark$\checkmark \checkmark$\checkmark^\checkmark\\checkmarkc\checkmarki\checkmarkr\checkmarkc\checkmark$\checkmark\\checkmark$\checkmark^\checkmark\\checkmarkc\checkmarki\checkmarkr\checkmarkc\checkmark$\checkmarkt\checkmark$\checkmark^\checkmark\\checkmarkc\checkmarki\checkmarkr\checkmarkc\checkmark$\checkmarke\checkmark$\checkmark^\checkmark\\checkmarkc\checkmarki\checkmarkr\checkmarkc\checkmark$\checkmarkx\checkmark$\checkmark^\checkmark\\checkmarkc\checkmarki\checkmarkr\checkmarkc\checkmark$\checkmarkt\checkmark$\checkmark^\checkmark\\checkmarkc\checkmarki\checkmarkr\checkmarkc\checkmark$\checkmarkb\checkmark$\checkmark^\checkmark\\checkmarkc\checkmarki\checkmarkr\checkmarkc\checkmark$\checkmarkf\checkmark$\checkmark^\checkmark\\checkmarkc\checkmarki\checkmarkr\checkmarkc\checkmark$\checkmark{\checkmark$\checkmark^\checkmark\\checkmarkc\checkmarki\checkmarkr\checkmarkc\checkmark$\checkmarkD\checkmark$\checkmark^\checkmark\\checkmarkc\checkmarki\checkmarkr\checkmarkc\checkmark$\checkmarké\checkmark$\checkmark^\checkmark\\checkmarkc\checkmarki\checkmarkr\checkmarkc\checkmark$\checkmarks\checkmark$\checkmark^\checkmark\\checkmarkc\checkmarki\checkmarkr\checkmarkc\checkmark$\checkmarki\checkmark$\checkmark^\checkmark\\checkmarkc\checkmarki\checkmarkr\checkmarkc\checkmark$\checkmarkg\checkmark$\checkmark^\checkmark\\checkmarkc\checkmarki\checkmarkr\checkmarkc\checkmark$\checkmarkn\checkmark$\checkmark^\checkmark\\checkmarkc\checkmarki\checkmarkr\checkmarkc\checkmark$\checkmarka\checkmark$\checkmark^\checkmark\\checkmarkc\checkmarki\checkmarkr\checkmarkc\checkmark$\checkmarkt\checkmark$\checkmark^\checkmark\\checkmarkc\checkmarki\checkmarkr\checkmarkc\checkmark$\checkmarki\checkmark$\checkmark^\checkmark\\checkmarkc\checkmarki\checkmarkr\checkmarkc\checkmark$\checkmarko\checkmark$\checkmark^\checkmark\\checkmarkc\checkmarki\checkmarkr\checkmarkc\checkmark$\checkmarkn\checkmark$\checkmark^\checkmark\\checkmarkc\checkmarki\checkmarkr\checkmarkc\checkmark$\checkmark}\checkmark$\checkmark^\checkmark\\checkmarkc\checkmarki\checkmarkr\checkmarkc\checkmark$\checkmark \checkmark$\checkmark^\checkmark\\checkmarkc\checkmarki\checkmarkr\checkmarkc\checkmark$\checkmark&\checkmark$\checkmark^\checkmark\\checkmarkc\checkmarki\checkmarkr\checkmarkc\checkmark$\checkmark \checkmark$\checkmark^\checkmark\\checkmarkc\checkmarki\checkmarkr\checkmarkc\checkmark$\checkmark\\checkmark$\checkmark^\checkmark\\checkmarkc\checkmarki\checkmarkr\checkmarkc\checkmark$\checkmarkt\checkmark$\checkmark^\checkmark\\checkmarkc\checkmarki\checkmarkr\checkmarkc\checkmark$\checkmarke\checkmark$\checkmark^\checkmark\\checkmarkc\checkmarki\checkmarkr\checkmarkc\checkmark$\checkmarkx\checkmark$\checkmark^\checkmark\\checkmarkc\checkmarki\checkmarkr\checkmarkc\checkmark$\checkmarkt\checkmark$\checkmark^\checkmark\\checkmarkc\checkmarki\checkmarkr\checkmarkc\checkmark$\checkmarkb\checkmark$\checkmark^\checkmark\\checkmarkc\checkmarki\checkmarkr\checkmarkc\checkmark$\checkmarkf\checkmark$\checkmark^\checkmark\\checkmarkc\checkmarki\checkmarkr\checkmarkc\checkmark$\checkmark{\checkmark$\checkmark^\checkmark\\checkmarkc\checkmarki\checkmarkr\checkmarkc\checkmark$\checkmarkU\checkmark$\checkmark^\checkmark\\checkmarkc\checkmarki\checkmarkr\checkmarkc\checkmark$\checkmarkn\checkmark$\checkmark^\checkmark\\checkmarkc\checkmarki\checkmarkr\checkmarkc\checkmark$\checkmarki\checkmark$\checkmark^\checkmark\\checkmarkc\checkmarki\checkmarkr\checkmarkc\checkmark$\checkmarkt\checkmark$\checkmark^\checkmark\\checkmarkc\checkmarki\checkmarkr\checkmarkc\checkmark$\checkmarké\checkmark$\checkmark^\checkmark\\checkmarkc\checkmarki\checkmarkr\checkmarkc\checkmark$\checkmark}\checkmark$\checkmark^\checkmark\\checkmarkc\checkmarki\checkmarkr\checkmarkc\checkmark$\checkmark \checkmark$\checkmark^\checkmark\\checkmarkc\checkmarki\checkmarkr\checkmarkc\checkmark$\checkmark\\checkmark$\checkmark^\checkmark\\checkmarkc\checkmarki\checkmarkr\checkmarkc\checkmark$\checkmark\\checkmark$\checkmark^\checkmark\\checkmarkc\checkmarki\checkmarkr\checkmarkc\checkmark$\checkmark \checkmark$\checkmark^\checkmark\\checkmarkc\checkmarki\checkmarkr\checkmarkc\checkmark$\checkmark\\checkmark$\checkmark^\checkmark\\checkmarkc\checkmarki\checkmarkr\checkmarkc\checkmark$\checkmarkh\checkmark$\checkmark^\checkmark\\checkmarkc\checkmarki\checkmarkr\checkmarkc\checkmark$\checkmarkl\checkmark$\checkmark^\checkmark\\checkmarkc\checkmarki\checkmarkr\checkmarkc\checkmark$\checkmarki\checkmark$\checkmark^\checkmark\\checkmarkc\checkmarki\checkmarkr\checkmarkc\checkmark$\checkmarkn\checkmark$\checkmark^\checkmark\\checkmarkc\checkmarki\checkmarkr\checkmarkc\checkmark$\checkmarke\checkmark$\checkmark^\checkmark\\checkmarkc\checkmarki\checkmarkr\checkmarkc\checkmark$\checkmark
\checkmark$\checkmark^\checkmark\\checkmarkc\checkmarki\checkmarkr\checkmarkc\checkmark$\checkmark$\checkmark$\checkmark^\checkmark\\checkmarkc\checkmarki\checkmarkr\checkmarkc\checkmark$\checkmarkF\checkmark$\checkmark^\checkmark\\checkmarkc\checkmarki\checkmarkr\checkmarkc\checkmark$\checkmark_\checkmark$\checkmark^\checkmark\\checkmarkc\checkmarki\checkmarkr\checkmarkc\checkmark$\checkmark{\checkmark$\checkmark^\checkmark\\checkmarkc\checkmarki\checkmarkr\checkmarkc\checkmark$\checkmarka\checkmark$\checkmark^\checkmark\\checkmarkc\checkmarki\checkmarkr\checkmarkc\checkmark$\checkmarkx\checkmark$\checkmark^\checkmark\\checkmarkc\checkmarki\checkmarkr\checkmarkc\checkmark$\checkmark}\checkmark$\checkmark^\checkmark\\checkmarkc\checkmarki\checkmarkr\checkmarkc\checkmark$\checkmark$\checkmark$\checkmark^\checkmark\\checkmarkc\checkmarki\checkmarkr\checkmarkc\checkmark$\checkmark \checkmark$\checkmark^\checkmark\\checkmarkc\checkmarki\checkmarkr\checkmarkc\checkmark$\checkmark&\checkmark$\checkmark^\checkmark\\checkmarkc\checkmarki\checkmarkr\checkmarkc\checkmark$\checkmark \checkmark$\checkmark^\checkmark\\checkmarkc\checkmarki\checkmarkr\checkmarkc\checkmark$\checkmarkF\checkmark$\checkmark^\checkmark\\checkmarkc\checkmarki\checkmarkr\checkmarkc\checkmark$\checkmarko\checkmark$\checkmark^\checkmark\\checkmarkc\checkmarki\checkmarkr\checkmarkc\checkmark$\checkmarkr\checkmark$\checkmark^\checkmark\\checkmarkc\checkmarki\checkmarkr\checkmarkc\checkmark$\checkmarkc\checkmark$\checkmark^\checkmark\\checkmarkc\checkmarki\checkmarkr\checkmarkc\checkmark$\checkmarke\checkmark$\checkmark^\checkmark\\checkmarkc\checkmarki\checkmarkr\checkmarkc\checkmark$\checkmark \checkmark$\checkmark^\checkmark\\checkmarkc\checkmarki\checkmarkr\checkmarkc\checkmark$\checkmarka\checkmark$\checkmark^\checkmark\\checkmarkc\checkmarki\checkmarkr\checkmarkc\checkmark$\checkmarkx\checkmark$\checkmark^\checkmark\\checkmarkc\checkmarki\checkmarkr\checkmarkc\checkmark$\checkmarki\checkmark$\checkmark^\checkmark\\checkmarkc\checkmarki\checkmarkr\checkmarkc\checkmark$\checkmarka\checkmark$\checkmark^\checkmark\\checkmarkc\checkmarki\checkmarkr\checkmarkc\checkmark$\checkmarkl\checkmark$\checkmark^\checkmark\\checkmarkc\checkmarki\checkmarkr\checkmarkc\checkmark$\checkmarke\checkmark$\checkmark^\checkmark\\checkmarkc\checkmarki\checkmarkr\checkmarkc\checkmark$\checkmark \checkmark$\checkmark^\checkmark\\checkmarkc\checkmarki\checkmarkr\checkmarkc\checkmark$\checkmarkt\checkmark$\checkmark^\checkmark\\checkmarkc\checkmarki\checkmarkr\checkmarkc\checkmark$\checkmarko\checkmark$\checkmark^\checkmark\\checkmarkc\checkmarki\checkmarkr\checkmarkc\checkmark$\checkmarkt\checkmark$\checkmark^\checkmark\\checkmarkc\checkmarki\checkmarkr\checkmarkc\checkmark$\checkmarka\checkmark$\checkmark^\checkmark\\checkmarkc\checkmarki\checkmarkr\checkmarkc\checkmark$\checkmarkl\checkmark$\checkmark^\checkmark\\checkmarkc\checkmarki\checkmarkr\checkmarkc\checkmark$\checkmarke\checkmark$\checkmark^\checkmark\\checkmarkc\checkmarki\checkmarkr\checkmarkc\checkmark$\checkmark \checkmark$\checkmark^\checkmark\\checkmarkc\checkmarki\checkmarkr\checkmarkc\checkmark$\checkmarks\checkmark$\checkmark^\checkmark\\checkmarkc\checkmarki\checkmarkr\checkmarkc\checkmark$\checkmarku\checkmark$\checkmark^\checkmark\\checkmarkc\checkmarki\checkmarkr\checkmarkc\checkmark$\checkmarkr\checkmark$\checkmark^\checkmark\\checkmarkc\checkmarki\checkmarkr\checkmarkc\checkmark$\checkmark \checkmark$\checkmark^\checkmark\\checkmarkc\checkmarki\checkmarkr\checkmarkc\checkmark$\checkmarkl\checkmark$\checkmark^\checkmark\\checkmarkc\checkmarki\checkmarkr\checkmarkc\checkmark$\checkmarka\checkmark$\checkmark^\checkmark\\checkmarkc\checkmarki\checkmarkr\checkmarkc\checkmark$\checkmark \checkmark$\checkmark^\checkmark\\checkmarkc\checkmarki\checkmarkr\checkmarkc\checkmark$\checkmarkv\checkmark$\checkmark^\checkmark\\checkmarkc\checkmarki\checkmarkr\checkmarkc\checkmark$\checkmarki\checkmark$\checkmark^\checkmark\\checkmarkc\checkmarki\checkmarkr\checkmarkc\checkmark$\checkmarks\checkmark$\checkmark^\checkmark\\checkmarkc\checkmarki\checkmarkr\checkmarkc\checkmark$\checkmark \checkmark$\checkmark^\checkmark\\checkmarkc\checkmarki\checkmarkr\checkmarkc\checkmark$\checkmark&\checkmark$\checkmark^\checkmark\\checkmarkc\checkmarki\checkmarkr\checkmarkc\checkmark$\checkmark \checkmark$\checkmark^\checkmark\\checkmarkc\checkmarki\checkmarkr\checkmarkc\checkmark$\checkmarkN\checkmark$\checkmark^\checkmark\\checkmarkc\checkmarki\checkmarkr\checkmarkc\checkmark$\checkmark \checkmark$\checkmark^\checkmark\\checkmarkc\checkmarki\checkmarkr\checkmarkc\checkmark$\checkmark\\checkmark$\checkmark^\checkmark\\checkmarkc\checkmarki\checkmarkr\checkmarkc\checkmark$\checkmark\\checkmark$\checkmark^\checkmark\\checkmarkc\checkmarki\checkmarkr\checkmarkc\checkmark$\checkmark \checkmark$\checkmark^\checkmark\\checkmarkc\checkmarki\checkmarkr\checkmarkc\checkmark$\checkmark\\checkmark$\checkmark^\checkmark\\checkmarkc\checkmarki\checkmarkr\checkmarkc\checkmark$\checkmarkh\checkmark$\checkmark^\checkmark\\checkmarkc\checkmarki\checkmarkr\checkmarkc\checkmark$\checkmarkl\checkmark$\checkmark^\checkmark\\checkmarkc\checkmarki\checkmarkr\checkmarkc\checkmark$\checkmarki\checkmark$\checkmark^\checkmark\\checkmarkc\checkmarki\checkmarkr\checkmarkc\checkmark$\checkmarkn\checkmark$\checkmark^\checkmark\\checkmarkc\checkmarki\checkmarkr\checkmarkc\checkmark$\checkmarke\checkmark$\checkmark^\checkmark\\checkmarkc\checkmarki\checkmarkr\checkmarkc\checkmark$\checkmark
\checkmark$\checkmark^\checkmark\\checkmarkc\checkmarki\checkmarkr\checkmarkc\checkmark$\checkmark$\checkmark$\checkmark^\checkmark\\checkmarkc\checkmarki\checkmarkr\checkmarkc\checkmark$\checkmarkF\checkmark$\checkmark^\checkmark\\checkmarkc\checkmarki\checkmarkr\checkmarkc\checkmark$\checkmark_\checkmark$\checkmark^\checkmark\\checkmarkc\checkmarki\checkmarkr\checkmarkc\checkmark$\checkmark{\checkmark$\checkmark^\checkmark\\checkmarkc\checkmarki\checkmarkr\checkmarkc\checkmark$\checkmarku\checkmark$\checkmark^\checkmark\\checkmarkc\checkmarki\checkmarkr\checkmarkc\checkmark$\checkmarks\checkmark$\checkmark^\checkmark\\checkmarkc\checkmarki\checkmarkr\checkmarkc\checkmark$\checkmarki\checkmark$\checkmark^\checkmark\\checkmarkc\checkmarki\checkmarkr\checkmarkc\checkmark$\checkmarkn\checkmark$\checkmark^\checkmark\\checkmarkc\checkmarki\checkmarkr\checkmarkc\checkmark$\checkmarka\checkmark$\checkmark^\checkmark\\checkmarkc\checkmarki\checkmarkr\checkmarkc\checkmark$\checkmarkg\checkmark$\checkmark^\checkmark\\checkmarkc\checkmarki\checkmarkr\checkmarkc\checkmark$\checkmarke\checkmark$\checkmark^\checkmark\\checkmarkc\checkmarki\checkmarkr\checkmarkc\checkmark$\checkmark}\checkmark$\checkmark^\checkmark\\checkmarkc\checkmarki\checkmarkr\checkmarkc\checkmark$\checkmark$\checkmark$\checkmark^\checkmark\\checkmarkc\checkmarki\checkmarkr\checkmarkc\checkmark$\checkmark \checkmark$\checkmark^\checkmark\\checkmarkc\checkmarki\checkmarkr\checkmarkc\checkmark$\checkmark&\checkmark$\checkmark^\checkmark\\checkmarkc\checkmarki\checkmarkr\checkmarkc\checkmark$\checkmark \checkmark$\checkmark^\checkmark\\checkmarkc\checkmarki\checkmarkr\checkmarkc\checkmark$\checkmarkE\checkmark$\checkmark^\checkmark\\checkmarkc\checkmarki\checkmarkr\checkmarkc\checkmark$\checkmarkf\checkmark$\checkmark^\checkmark\\checkmarkc\checkmarki\checkmarkr\checkmarkc\checkmark$\checkmarkf\checkmark$\checkmark^\checkmark\\checkmarkc\checkmarki\checkmarkr\checkmarkc\checkmark$\checkmarko\checkmark$\checkmark^\checkmark\\checkmarkc\checkmarki\checkmarkr\checkmarkc\checkmark$\checkmarkr\checkmark$\checkmark^\checkmark\\checkmarkc\checkmarki\checkmarkr\checkmarkc\checkmark$\checkmarkt\checkmark$\checkmark^\checkmark\\checkmarkc\checkmarki\checkmarkr\checkmarkc\checkmark$\checkmark \checkmark$\checkmark^\checkmark\\checkmarkc\checkmarki\checkmarkr\checkmarkc\checkmark$\checkmarkd\checkmark$\checkmark^\checkmark\\checkmarkc\checkmarki\checkmarkr\checkmarkc\checkmark$\checkmarke\checkmark$\checkmark^\checkmark\\checkmarkc\checkmarki\checkmarkr\checkmarkc\checkmark$\checkmark \checkmark$\checkmark^\checkmark\\checkmarkc\checkmarki\checkmarkr\checkmarkc\checkmark$\checkmarkc\checkmark$\checkmark^\checkmark\\checkmarkc\checkmarki\checkmarkr\checkmarkc\checkmark$\checkmarko\checkmark$\checkmark^\checkmark\\checkmarkc\checkmarki\checkmarkr\checkmarkc\checkmark$\checkmarku\checkmark$\checkmark^\checkmark\\checkmarkc\checkmarki\checkmarkr\checkmarkc\checkmark$\checkmarkp\checkmark$\checkmark^\checkmark\\checkmarkc\checkmarki\checkmarkr\checkmarkc\checkmark$\checkmarke\checkmark$\checkmark^\checkmark\\checkmarkc\checkmarki\checkmarkr\checkmarkc\checkmark$\checkmark \checkmark$\checkmark^\checkmark\\checkmarkc\checkmarki\checkmarkr\checkmarkc\checkmark$\checkmarkt\checkmark$\checkmark^\checkmark\\checkmarkc\checkmarki\checkmarkr\checkmarkc\checkmark$\checkmarkr\checkmark$\checkmark^\checkmark\\checkmarkc\checkmarki\checkmarkr\checkmarkc\checkmark$\checkmarka\checkmark$\checkmark^\checkmark\\checkmarkc\checkmarki\checkmarkr\checkmarkc\checkmark$\checkmarkn\checkmark$\checkmark^\checkmark\\checkmarkc\checkmarki\checkmarkr\checkmarkc\checkmark$\checkmarks\checkmark$\checkmark^\checkmark\\checkmarkc\checkmarki\checkmarkr\checkmarkc\checkmark$\checkmarkm\checkmark$\checkmark^\checkmark\\checkmarkc\checkmarki\checkmarkr\checkmarkc\checkmark$\checkmarki\checkmark$\checkmark^\checkmark\\checkmarkc\checkmarki\checkmarkr\checkmarkc\checkmark$\checkmarks\checkmark$\checkmark^\checkmark\\checkmarkc\checkmarki\checkmarkr\checkmarkc\checkmark$\checkmark \checkmark$\checkmark^\checkmark\\checkmarkc\checkmarki\checkmarkr\checkmarkc\checkmark$\checkmarkà\checkmark$\checkmark^\checkmark\\checkmarkc\checkmarki\checkmarkr\checkmarkc\checkmark$\checkmark \checkmark$\checkmark^\checkmark\\checkmarkc\checkmarki\checkmarkr\checkmarkc\checkmark$\checkmarkl\checkmark$\checkmark^\checkmark\\checkmarkc\checkmarki\checkmarkr\checkmarkc\checkmark$\checkmark'\checkmark$\checkmark^\checkmark\\checkmarkc\checkmarki\checkmarkr\checkmarkc\checkmark$\checkmarka\checkmark$\checkmark^\checkmark\\checkmarkc\checkmarki\checkmarkr\checkmarkc\checkmark$\checkmarkx\checkmark$\checkmark^\checkmark\\checkmarkc\checkmarki\checkmarkr\checkmarkc\checkmark$\checkmarke\checkmark$\checkmark^\checkmark\\checkmarkc\checkmarki\checkmarkr\checkmarkc\checkmark$\checkmark \checkmark$\checkmark^\checkmark\\checkmarkc\checkmarki\checkmarkr\checkmarkc\checkmark$\checkmark&\checkmark$\checkmark^\checkmark\\checkmarkc\checkmarki\checkmarkr\checkmarkc\checkmark$\checkmark \checkmark$\checkmark^\checkmark\\checkmarkc\checkmarki\checkmarkr\checkmarkc\checkmark$\checkmarkN\checkmark$\checkmark^\checkmark\\checkmarkc\checkmarki\checkmarkr\checkmarkc\checkmark$\checkmark \checkmark$\checkmark^\checkmark\\checkmarkc\checkmarki\checkmarkr\checkmarkc\checkmark$\checkmark\\checkmark$\checkmark^\checkmark\\checkmarkc\checkmarki\checkmarkr\checkmarkc\checkmark$\checkmark\\checkmark$\checkmark^\checkmark\\checkmarkc\checkmarki\checkmarkr\checkmarkc\checkmark$\checkmark \checkmark$\checkmark^\checkmark\\checkmarkc\checkmarki\checkmarkr\checkmarkc\checkmark$\checkmark\\checkmark$\checkmark^\checkmark\\checkmarkc\checkmarki\checkmarkr\checkmarkc\checkmark$\checkmarkh\checkmark$\checkmark^\checkmark\\checkmarkc\checkmarki\checkmarkr\checkmarkc\checkmark$\checkmarkl\checkmark$\checkmark^\checkmark\\checkmarkc\checkmarki\checkmarkr\checkmarkc\checkmark$\checkmarki\checkmark$\checkmark^\checkmark\\checkmarkc\checkmarki\checkmarkr\checkmarkc\checkmark$\checkmarkn\checkmark$\checkmark^\checkmark\\checkmarkc\checkmarki\checkmarkr\checkmarkc\checkmark$\checkmarke\checkmark$\checkmark^\checkmark\\checkmarkc\checkmarki\checkmarkr\checkmarkc\checkmark$\checkmark
\checkmark$\checkmark^\checkmark\\checkmarkc\checkmarki\checkmarkr\checkmarkc\checkmark$\checkmark$\checkmark$\checkmark^\checkmark\\checkmarkc\checkmarki\checkmarkr\checkmarkc\checkmark$\checkmarkF\checkmark$\checkmark^\checkmark\\checkmarkc\checkmarki\checkmarkr\checkmarkc\checkmark$\checkmark_\checkmark$\checkmark^\checkmark\\checkmarkc\checkmarki\checkmarkr\checkmarkc\checkmark$\checkmark{\checkmark$\checkmark^\checkmark\\checkmarkc\checkmarki\checkmarkr\checkmarkc\checkmark$\checkmarkf\checkmark$\checkmark^\checkmark\\checkmarkc\checkmarki\checkmarkr\checkmarkc\checkmark$\checkmarkr\checkmark$\checkmark^\checkmark\\checkmarkc\checkmarki\checkmarkr\checkmarkc\checkmark$\checkmarko\checkmark$\checkmark^\checkmark\\checkmarkc\checkmarki\checkmarkr\checkmarkc\checkmark$\checkmarkt\checkmark$\checkmark^\checkmark\\checkmarkc\checkmarki\checkmarkr\checkmarkc\checkmark$\checkmarkt\checkmark$\checkmark^\checkmark\\checkmarkc\checkmarki\checkmarkr\checkmarkc\checkmark$\checkmark}\checkmark$\checkmark^\checkmark\\checkmarkc\checkmarki\checkmarkr\checkmarkc\checkmark$\checkmark$\checkmark$\checkmark^\checkmark\\checkmarkc\checkmarki\checkmarkr\checkmarkc\checkmark$\checkmark \checkmark$\checkmark^\checkmark\\checkmarkc\checkmarki\checkmarkr\checkmarkc\checkmark$\checkmark&\checkmark$\checkmark^\checkmark\\checkmarkc\checkmarki\checkmarkr\checkmarkc\checkmark$\checkmark \checkmark$\checkmark^\checkmark\\checkmarkc\checkmarki\checkmarkr\checkmarkc\checkmark$\checkmarkF\checkmark$\checkmark^\checkmark\\checkmarkc\checkmarki\checkmarkr\checkmarkc\checkmark$\checkmarko\checkmark$\checkmark^\checkmark\\checkmarkc\checkmarki\checkmarkr\checkmarkc\checkmark$\checkmarkr\checkmark$\checkmark^\checkmark\\checkmarkc\checkmarki\checkmarkr\checkmarkc\checkmark$\checkmarkc\checkmark$\checkmark^\checkmark\\checkmarkc\checkmarki\checkmarkr\checkmarkc\checkmark$\checkmarke\checkmark$\checkmark^\checkmark\\checkmarkc\checkmarki\checkmarkr\checkmarkc\checkmark$\checkmark \checkmark$\checkmark^\checkmark\\checkmarkc\checkmarki\checkmarkr\checkmarkc\checkmark$\checkmarkd\checkmark$\checkmark^\checkmark\\checkmarkc\checkmarki\checkmarkr\checkmarkc\checkmark$\checkmarke\checkmark$\checkmark^\checkmark\\checkmarkc\checkmarki\checkmarkr\checkmarkc\checkmark$\checkmark \checkmark$\checkmark^\checkmark\\checkmarkc\checkmarki\checkmarkr\checkmarkc\checkmark$\checkmarkf\checkmark$\checkmark^\checkmark\\checkmarkc\checkmarki\checkmarkr\checkmarkc\checkmark$\checkmarkr\checkmark$\checkmark^\checkmark\\checkmarkc\checkmarki\checkmarkr\checkmarkc\checkmark$\checkmarko\checkmark$\checkmark^\checkmark\\checkmarkc\checkmarki\checkmarkr\checkmarkc\checkmark$\checkmarkt\checkmark$\checkmark^\checkmark\\checkmarkc\checkmarki\checkmarkr\checkmarkc\checkmark$\checkmarkt\checkmark$\checkmark^\checkmark\\checkmarkc\checkmarki\checkmarkr\checkmarkc\checkmark$\checkmarke\checkmark$\checkmark^\checkmark\\checkmarkc\checkmarki\checkmarkr\checkmarkc\checkmark$\checkmarkm\checkmark$\checkmark^\checkmark\\checkmarkc\checkmarki\checkmarkr\checkmarkc\checkmark$\checkmarke\checkmark$\checkmark^\checkmark\\checkmarkc\checkmarki\checkmarkr\checkmarkc\checkmark$\checkmarkn\checkmark$\checkmark^\checkmark\\checkmarkc\checkmarki\checkmarkr\checkmarkc\checkmark$\checkmarkt\checkmark$\checkmark^\checkmark\\checkmarkc\checkmarki\checkmarkr\checkmarkc\checkmark$\checkmark \checkmark$\checkmark^\checkmark\\checkmarkc\checkmarki\checkmarkr\checkmarkc\checkmark$\checkmarkd\checkmark$\checkmark^\checkmark\\checkmarkc\checkmarki\checkmarkr\checkmarkc\checkmark$\checkmarke\checkmark$\checkmark^\checkmark\\checkmarkc\checkmarki\checkmarkr\checkmarkc\checkmark$\checkmarks\checkmark$\checkmark^\checkmark\\checkmarkc\checkmarki\checkmarkr\checkmarkc\checkmark$\checkmark \checkmark$\checkmark^\checkmark\\checkmarkc\checkmarki\checkmarkr\checkmarkc\checkmark$\checkmarkg\checkmark$\checkmark^\checkmark\\checkmarkc\checkmarki\checkmarkr\checkmarkc\checkmark$\checkmarku\checkmark$\checkmark^\checkmark\\checkmarkc\checkmarki\checkmarkr\checkmarkc\checkmark$\checkmarki\checkmark$\checkmark^\checkmark\\checkmarkc\checkmarki\checkmarkr\checkmarkc\checkmark$\checkmarkd\checkmark$\checkmark^\checkmark\\checkmarkc\checkmarki\checkmarkr\checkmarkc\checkmark$\checkmarka\checkmark$\checkmark^\checkmark\\checkmarkc\checkmarki\checkmarkr\checkmarkc\checkmark$\checkmarkg\checkmark$\checkmark^\checkmark\\checkmarkc\checkmarki\checkmarkr\checkmarkc\checkmark$\checkmarke\checkmark$\checkmark^\checkmark\\checkmarkc\checkmarki\checkmarkr\checkmarkc\checkmark$\checkmarks\checkmark$\checkmark^\checkmark\\checkmarkc\checkmarki\checkmarkr\checkmarkc\checkmark$\checkmark \checkmark$\checkmark^\checkmark\\checkmarkc\checkmarki\checkmarkr\checkmarkc\checkmark$\checkmark&\checkmark$\checkmark^\checkmark\\checkmarkc\checkmarki\checkmarkr\checkmarkc\checkmark$\checkmark \checkmark$\checkmark^\checkmark\\checkmarkc\checkmarki\checkmarkr\checkmarkc\checkmark$\checkmarkN\checkmark$\checkmark^\checkmark\\checkmarkc\checkmarki\checkmarkr\checkmarkc\checkmark$\checkmark \checkmark$\checkmark^\checkmark\\checkmarkc\checkmarki\checkmarkr\checkmarkc\checkmark$\checkmark\\checkmark$\checkmark^\checkmark\\checkmarkc\checkmarki\checkmarkr\checkmarkc\checkmark$\checkmark\\checkmark$\checkmark^\checkmark\\checkmarkc\checkmarki\checkmarkr\checkmarkc\checkmark$\checkmark \checkmark$\checkmark^\checkmark\\checkmarkc\checkmarki\checkmarkr\checkmarkc\checkmark$\checkmark\\checkmark$\checkmark^\checkmark\\checkmarkc\checkmarki\checkmarkr\checkmarkc\checkmark$\checkmarkh\checkmark$\checkmark^\checkmark\\checkmarkc\checkmarki\checkmarkr\checkmarkc\checkmark$\checkmarkl\checkmark$\checkmark^\checkmark\\checkmarkc\checkmarki\checkmarkr\checkmarkc\checkmark$\checkmarki\checkmark$\checkmark^\checkmark\\checkmarkc\checkmarki\checkmarkr\checkmarkc\checkmark$\checkmarkn\checkmark$\checkmark^\checkmark\\checkmarkc\checkmarki\checkmarkr\checkmarkc\checkmark$\checkmarke\checkmark$\checkmark^\checkmark\\checkmarkc\checkmarki\checkmarkr\checkmarkc\checkmark$\checkmark
\checkmark$\checkmark^\checkmark\\checkmarkc\checkmarki\checkmarkr\checkmarkc\checkmark$\checkmark$\checkmark$\checkmark^\checkmark\\checkmarkc\checkmarki\checkmarkr\checkmarkc\checkmark$\checkmarkF\checkmark$\checkmark^\checkmark\\checkmarkc\checkmarki\checkmarkr\checkmarkc\checkmark$\checkmark_\checkmark$\checkmark^\checkmark\\checkmarkc\checkmarki\checkmarkr\checkmarkc\checkmark$\checkmark{\checkmark$\checkmark^\checkmark\\checkmarkc\checkmarki\checkmarkr\checkmarkc\checkmark$\checkmarki\checkmark$\checkmark^\checkmark\\checkmarkc\checkmarki\checkmarkr\checkmarkc\checkmark$\checkmarkn\checkmark$\checkmark^\checkmark\\checkmarkc\checkmarki\checkmarkr\checkmarkc\checkmark$\checkmarke\checkmark$\checkmark^\checkmark\\checkmarkc\checkmarki\checkmarkr\checkmarkc\checkmark$\checkmarkr\checkmark$\checkmark^\checkmark\\checkmarkc\checkmarki\checkmarkr\checkmarkc\checkmark$\checkmarkt\checkmark$\checkmark^\checkmark\\checkmarkc\checkmarki\checkmarkr\checkmarkc\checkmark$\checkmarki\checkmark$\checkmark^\checkmark\\checkmarkc\checkmarki\checkmarkr\checkmarkc\checkmark$\checkmarke\checkmark$\checkmark^\checkmark\\checkmarkc\checkmarki\checkmarkr\checkmarkc\checkmark$\checkmark}\checkmark$\checkmark^\checkmark\\checkmarkc\checkmarki\checkmarkr\checkmarkc\checkmark$\checkmark$\checkmark$\checkmark^\checkmark\\checkmarkc\checkmarki\checkmarkr\checkmarkc\checkmark$\checkmark \checkmark$\checkmark^\checkmark\\checkmarkc\checkmarki\checkmarkr\checkmarkc\checkmark$\checkmark&\checkmark$\checkmark^\checkmark\\checkmarkc\checkmarki\checkmarkr\checkmarkc\checkmark$\checkmark \checkmark$\checkmark^\checkmark\\checkmarkc\checkmarki\checkmarkr\checkmarkc\checkmark$\checkmarkF\checkmark$\checkmark^\checkmark\\checkmarkc\checkmarki\checkmarkr\checkmarkc\checkmark$\checkmarko\checkmark$\checkmark^\checkmark\\checkmarkc\checkmarki\checkmarkr\checkmarkc\checkmark$\checkmarkr\checkmark$\checkmark^\checkmark\\checkmarkc\checkmarki\checkmarkr\checkmarkc\checkmark$\checkmarkc\checkmark$\checkmark^\checkmark\\checkmarkc\checkmarki\checkmarkr\checkmarkc\checkmark$\checkmarke\checkmark$\checkmark^\checkmark\\checkmarkc\checkmarki\checkmarkr\checkmarkc\checkmark$\checkmark \checkmark$\checkmark^\checkmark\\checkmarkc\checkmarki\checkmarkr\checkmarkc\checkmark$\checkmarkd\checkmark$\checkmark^\checkmark\\checkmarkc\checkmarki\checkmarkr\checkmarkc\checkmark$\checkmark'\checkmark$\checkmark^\checkmark\\checkmarkc\checkmarki\checkmarkr\checkmarkc\checkmark$\checkmarki\checkmark$\checkmark^\checkmark\\checkmarkc\checkmarki\checkmarkr\checkmarkc\checkmark$\checkmarkn\checkmark$\checkmark^\checkmark\\checkmarkc\checkmarki\checkmarkr\checkmarkc\checkmark$\checkmarke\checkmark$\checkmark^\checkmark\\checkmarkc\checkmarki\checkmarkr\checkmarkc\checkmark$\checkmarkr\checkmark$\checkmark^\checkmark\\checkmarkc\checkmarki\checkmarkr\checkmarkc\checkmark$\checkmarkt\checkmark$\checkmark^\checkmark\\checkmarkc\checkmarki\checkmarkr\checkmarkc\checkmark$\checkmarki\checkmark$\checkmark^\checkmark\\checkmarkc\checkmarki\checkmarkr\checkmarkc\checkmark$\checkmarke\checkmark$\checkmark^\checkmark\\checkmarkc\checkmarki\checkmarkr\checkmarkc\checkmark$\checkmark \checkmark$\checkmark^\checkmark\\checkmarkc\checkmarki\checkmarkr\checkmarkc\checkmark$\checkmarkl\checkmark$\checkmark^\checkmark\\checkmarkc\checkmarki\checkmarkr\checkmarkc\checkmark$\checkmarko\checkmark$\checkmark^\checkmark\\checkmarkc\checkmarki\checkmarkr\checkmarkc\checkmark$\checkmarkr\checkmark$\checkmark^\checkmark\\checkmarkc\checkmarki\checkmarkr\checkmarkc\checkmark$\checkmarks\checkmark$\checkmark^\checkmark\\checkmarkc\checkmarki\checkmarkr\checkmarkc\checkmark$\checkmark \checkmark$\checkmark^\checkmark\\checkmarkc\checkmarki\checkmarkr\checkmarkc\checkmark$\checkmarkd\checkmark$\checkmark^\checkmark\\checkmarkc\checkmarki\checkmarkr\checkmarkc\checkmark$\checkmarke\checkmark$\checkmark^\checkmark\\checkmarkc\checkmarki\checkmarkr\checkmarkc\checkmark$\checkmark \checkmark$\checkmark^\checkmark\\checkmarkc\checkmarki\checkmarkr\checkmarkc\checkmark$\checkmarkl\checkmark$\checkmark^\checkmark\\checkmarkc\checkmarki\checkmarkr\checkmarkc\checkmark$\checkmark'\checkmark$\checkmark^\checkmark\\checkmarkc\checkmarki\checkmarkr\checkmarkc\checkmark$\checkmarka\checkmark$\checkmark^\checkmark\\checkmarkc\checkmarki\checkmarkr\checkmarkc\checkmark$\checkmarkc\checkmark$\checkmark^\checkmark\\checkmarkc\checkmarki\checkmarkr\checkmarkc\checkmark$\checkmarkc\checkmark$\checkmark^\checkmark\\checkmarkc\checkmarki\checkmarkr\checkmarkc\checkmark$\checkmarké\checkmark$\checkmark^\checkmark\\checkmarkc\checkmarki\checkmarkr\checkmarkc\checkmark$\checkmarkl\checkmark$\checkmark^\checkmark\\checkmarkc\checkmarki\checkmarkr\checkmarkc\checkmark$\checkmarké\checkmark$\checkmark^\checkmark\\checkmarkc\checkmarki\checkmarkr\checkmarkc\checkmark$\checkmarkr\checkmark$\checkmark^\checkmark\\checkmarkc\checkmarki\checkmarkr\checkmarkc\checkmark$\checkmarka\checkmark$\checkmark^\checkmark\\checkmarkc\checkmarki\checkmarkr\checkmarkc\checkmark$\checkmarkt\checkmark$\checkmark^\checkmark\\checkmarkc\checkmarki\checkmarkr\checkmarkc\checkmark$\checkmarki\checkmark$\checkmark^\checkmark\\checkmarkc\checkmarki\checkmarkr\checkmarkc\checkmark$\checkmarko\checkmark$\checkmark^\checkmark\\checkmarkc\checkmarki\checkmarkr\checkmarkc\checkmark$\checkmarkn\checkmark$\checkmark^\checkmark\\checkmarkc\checkmarki\checkmarkr\checkmarkc\checkmark$\checkmark \checkmark$\checkmark^\checkmark\\checkmarkc\checkmarki\checkmarkr\checkmarkc\checkmark$\checkmark&\checkmark$\checkmark^\checkmark\\checkmarkc\checkmarki\checkmarkr\checkmarkc\checkmark$\checkmark \checkmark$\checkmark^\checkmark\\checkmarkc\checkmarki\checkmarkr\checkmarkc\checkmark$\checkmarkN\checkmark$\checkmark^\checkmark\\checkmarkc\checkmarki\checkmarkr\checkmarkc\checkmark$\checkmark \checkmark$\checkmark^\checkmark\\checkmarkc\checkmarki\checkmarkr\checkmarkc\checkmark$\checkmark\\checkmark$\checkmark^\checkmark\\checkmarkc\checkmarki\checkmarkr\checkmarkc\checkmark$\checkmark\\checkmark$\checkmark^\checkmark\\checkmarkc\checkmarki\checkmarkr\checkmarkc\checkmark$\checkmark \checkmark$\checkmark^\checkmark\\checkmarkc\checkmarki\checkmarkr\checkmarkc\checkmark$\checkmark\\checkmark$\checkmark^\checkmark\\checkmarkc\checkmarki\checkmarkr\checkmarkc\checkmark$\checkmarkh\checkmark$\checkmark^\checkmark\\checkmarkc\checkmarki\checkmarkr\checkmarkc\checkmark$\checkmarkl\checkmark$\checkmark^\checkmark\\checkmarkc\checkmarki\checkmarkr\checkmarkc\checkmark$\checkmarki\checkmark$\checkmark^\checkmark\\checkmarkc\checkmarki\checkmarkr\checkmarkc\checkmark$\checkmarkn\checkmark$\checkmark^\checkmark\\checkmarkc\checkmarki\checkmarkr\checkmarkc\checkmark$\checkmarke\checkmark$\checkmark^\checkmark\\checkmarkc\checkmarki\checkmarkr\checkmarkc\checkmark$\checkmark
\checkmark$\checkmark^\checkmark\\checkmarkc\checkmarki\checkmarkr\checkmarkc\checkmark$\checkmark$\checkmark$\checkmark^\checkmark\\checkmarkc\checkmarki\checkmarkr\checkmarkc\checkmark$\checkmarkF\checkmark$\checkmark^\checkmark\\checkmarkc\checkmarki\checkmarkr\checkmarkc\checkmark$\checkmark_\checkmark$\checkmark^\checkmark\\checkmarkc\checkmarki\checkmarkr\checkmarkc\checkmark$\checkmarkk\checkmark$\checkmark^\checkmark\\checkmarkc\checkmarki\checkmarkr\checkmarkc\checkmark$\checkmark$\checkmark$\checkmark^\checkmark\\checkmarkc\checkmarki\checkmarkr\checkmarkc\checkmark$\checkmark \checkmark$\checkmark^\checkmark\\checkmarkc\checkmarki\checkmarkr\checkmarkc\checkmark$\checkmark&\checkmark$\checkmark^\checkmark\\checkmarkc\checkmarki\checkmarkr\checkmarkc\checkmark$\checkmark \checkmark$\checkmark^\checkmark\\checkmarkc\checkmarki\checkmarkr\checkmarkc\checkmark$\checkmarkC\checkmark$\checkmark^\checkmark\\checkmarkc\checkmarki\checkmarkr\checkmarkc\checkmark$\checkmarkh\checkmark$\checkmark^\checkmark\\checkmarkc\checkmarki\checkmarkr\checkmarkc\checkmark$\checkmarka\checkmark$\checkmark^\checkmark\\checkmarkc\checkmarki\checkmarkr\checkmarkc\checkmark$\checkmarkr\checkmark$\checkmark^\checkmark\\checkmarkc\checkmarki\checkmarkr\checkmarkc\checkmark$\checkmarkg\checkmark$\checkmark^\checkmark\\checkmarkc\checkmarki\checkmarkr\checkmarkc\checkmark$\checkmarke\checkmark$\checkmark^\checkmark\\checkmarkc\checkmarki\checkmarkr\checkmarkc\checkmark$\checkmark \checkmark$\checkmark^\checkmark\\checkmarkc\checkmarki\checkmarkr\checkmarkc\checkmark$\checkmarkc\checkmark$\checkmark^\checkmark\\checkmarkc\checkmarki\checkmarkr\checkmarkc\checkmark$\checkmarkr\checkmark$\checkmark^\checkmark\\checkmarkc\checkmarki\checkmarkr\checkmarkc\checkmark$\checkmarki\checkmark$\checkmark^\checkmark\\checkmarkc\checkmarki\checkmarkr\checkmarkc\checkmark$\checkmarkt\checkmark$\checkmark^\checkmark\\checkmarkc\checkmarki\checkmarkr\checkmarkc\checkmark$\checkmarki\checkmark$\checkmark^\checkmark\\checkmarkc\checkmarki\checkmarkr\checkmarkc\checkmark$\checkmarkq\checkmark$\checkmark^\checkmark\\checkmarkc\checkmarki\checkmarkr\checkmarkc\checkmark$\checkmarku\checkmark$\checkmark^\checkmark\\checkmarkc\checkmarki\checkmarkr\checkmarkc\checkmark$\checkmarke\checkmark$\checkmark^\checkmark\\checkmarkc\checkmarki\checkmarkr\checkmarkc\checkmark$\checkmark \checkmark$\checkmark^\checkmark\\checkmarkc\checkmarki\checkmarkr\checkmarkc\checkmark$\checkmarkd\checkmark$\checkmark^\checkmark\\checkmarkc\checkmarki\checkmarkr\checkmarkc\checkmark$\checkmarke\checkmark$\checkmark^\checkmark\\checkmarkc\checkmarki\checkmarkr\checkmarkc\checkmark$\checkmark \checkmark$\checkmark^\checkmark\\checkmarkc\checkmarki\checkmarkr\checkmarkc\checkmark$\checkmarkf\checkmark$\checkmark^\checkmark\\checkmarkc\checkmarki\checkmarkr\checkmarkc\checkmark$\checkmarkl\checkmark$\checkmark^\checkmark\\checkmarkc\checkmarki\checkmarkr\checkmarkc\checkmark$\checkmarka\checkmark$\checkmark^\checkmark\\checkmarkc\checkmarki\checkmarkr\checkmarkc\checkmark$\checkmarkm\checkmark$\checkmark^\checkmark\\checkmarkc\checkmarki\checkmarkr\checkmarkc\checkmark$\checkmarkb\checkmark$\checkmark^\checkmark\\checkmarkc\checkmarki\checkmarkr\checkmarkc\checkmark$\checkmarka\checkmark$\checkmark^\checkmark\\checkmarkc\checkmarki\checkmarkr\checkmarkc\checkmark$\checkmarkg\checkmark$\checkmark^\checkmark\\checkmarkc\checkmarki\checkmarkr\checkmarkc\checkmark$\checkmarke\checkmark$\checkmark^\checkmark\\checkmarkc\checkmarki\checkmarkr\checkmarkc\checkmark$\checkmark \checkmark$\checkmark^\checkmark\\checkmarkc\checkmarki\checkmarkr\checkmarkc\checkmark$\checkmark(\checkmark$\checkmark^\checkmark\\checkmarkc\checkmarki\checkmarkr\checkmarkc\checkmark$\checkmarkE\checkmark$\checkmark^\checkmark\\checkmarkc\checkmarki\checkmarkr\checkmarkc\checkmark$\checkmarku\checkmark$\checkmark^\checkmark\\checkmarkc\checkmarki\checkmarkr\checkmarkc\checkmark$\checkmarkl\checkmark$\checkmark^\checkmark\\checkmarkc\checkmarki\checkmarkr\checkmarkc\checkmark$\checkmarke\checkmark$\checkmark^\checkmark\\checkmarkc\checkmarki\checkmarkr\checkmarkc\checkmark$\checkmarkr\checkmark$\checkmark^\checkmark\\checkmarkc\checkmarki\checkmarkr\checkmarkc\checkmark$\checkmark)\checkmark$\checkmark^\checkmark\\checkmarkc\checkmarki\checkmarkr\checkmarkc\checkmark$\checkmark \checkmark$\checkmark^\checkmark\\checkmarkc\checkmarki\checkmarkr\checkmarkc\checkmark$\checkmark&\checkmark$\checkmark^\checkmark\\checkmarkc\checkmarki\checkmarkr\checkmarkc\checkmark$\checkmark \checkmark$\checkmark^\checkmark\\checkmarkc\checkmarki\checkmarkr\checkmarkc\checkmark$\checkmarkN\checkmark$\checkmark^\checkmark\\checkmarkc\checkmarki\checkmarkr\checkmarkc\checkmark$\checkmark \checkmark$\checkmark^\checkmark\\checkmarkc\checkmarki\checkmarkr\checkmarkc\checkmark$\checkmark\\checkmark$\checkmark^\checkmark\\checkmarkc\checkmarki\checkmarkr\checkmarkc\checkmark$\checkmark\\checkmark$\checkmark^\checkmark\\checkmarkc\checkmarki\checkmarkr\checkmarkc\checkmark$\checkmark \checkmark$\checkmark^\checkmark\\checkmarkc\checkmarki\checkmarkr\checkmarkc\checkmark$\checkmark\\checkmark$\checkmark^\checkmark\\checkmarkc\checkmarki\checkmarkr\checkmarkc\checkmark$\checkmarkh\checkmark$\checkmark^\checkmark\\checkmarkc\checkmarki\checkmarkr\checkmarkc\checkmark$\checkmarkl\checkmark$\checkmark^\checkmark\\checkmarkc\checkmarki\checkmarkr\checkmarkc\checkmark$\checkmarki\checkmark$\checkmark^\checkmark\\checkmarkc\checkmarki\checkmarkr\checkmarkc\checkmark$\checkmarkn\checkmark$\checkmark^\checkmark\\checkmarkc\checkmarki\checkmarkr\checkmarkc\checkmark$\checkmarke\checkmark$\checkmark^\checkmark\\checkmarkc\checkmarki\checkmarkr\checkmarkc\checkmark$\checkmark
\checkmark$\checkmark^\checkmark\\checkmarkc\checkmarki\checkmarkr\checkmarkc\checkmark$\checkmark$\checkmark$\checkmark^\checkmark\\checkmarkc\checkmarki\checkmarkr\checkmarkc\checkmark$\checkmarkC\checkmark$\checkmark^\checkmark\\checkmarkc\checkmarki\checkmarkr\checkmarkc\checkmark$\checkmark_\checkmark$\checkmark^\checkmark\\checkmarkc\checkmarki\checkmarkr\checkmarkc\checkmark$\checkmarkr\checkmark$\checkmark^\checkmark\\checkmarkc\checkmarki\checkmarkr\checkmarkc\checkmark$\checkmark$\checkmark$\checkmark^\checkmark\\checkmarkc\checkmarki\checkmarkr\checkmarkc\checkmark$\checkmark \checkmark$\checkmark^\checkmark\\checkmarkc\checkmarki\checkmarkr\checkmarkc\checkmark$\checkmark&\checkmark$\checkmark^\checkmark\\checkmarkc\checkmarki\checkmarkr\checkmarkc\checkmark$\checkmark \checkmark$\checkmark^\checkmark\\checkmarkc\checkmarki\checkmarkr\checkmarkc\checkmark$\checkmarkC\checkmark$\checkmark^\checkmark\\checkmarkc\checkmarki\checkmarkr\checkmarkc\checkmark$\checkmarko\checkmark$\checkmark^\checkmark\\checkmarkc\checkmarki\checkmarkr\checkmarkc\checkmark$\checkmarku\checkmark$\checkmark^\checkmark\\checkmarkc\checkmarki\checkmarkr\checkmarkc\checkmark$\checkmarkp\checkmark$\checkmark^\checkmark\\checkmarkc\checkmarki\checkmarkr\checkmarkc\checkmark$\checkmarkl\checkmark$\checkmark^\checkmark\\checkmarkc\checkmarki\checkmarkr\checkmarkc\checkmark$\checkmarke\checkmark$\checkmark^\checkmark\\checkmarkc\checkmarki\checkmarkr\checkmarkc\checkmark$\checkmark \checkmark$\checkmark^\checkmark\\checkmarkc\checkmarki\checkmarkr\checkmarkc\checkmark$\checkmarkr\checkmark$\checkmark^\checkmark\\checkmarkc\checkmarki\checkmarkr\checkmarkc\checkmark$\checkmarké\checkmark$\checkmark^\checkmark\\checkmarkc\checkmarki\checkmarkr\checkmarkc\checkmark$\checkmarks\checkmark$\checkmark^\checkmark\\checkmarkc\checkmarki\checkmarkr\checkmarkc\checkmark$\checkmarki\checkmark$\checkmark^\checkmark\\checkmarkc\checkmarki\checkmarkr\checkmarkc\checkmark$\checkmarks\checkmark$\checkmark^\checkmark\\checkmarkc\checkmarki\checkmarkr\checkmarkc\checkmark$\checkmarkt\checkmark$\checkmark^\checkmark\\checkmarkc\checkmarki\checkmarkr\checkmarkc\checkmark$\checkmarka\checkmark$\checkmark^\checkmark\\checkmarkc\checkmarki\checkmarkr\checkmarkc\checkmark$\checkmarkn\checkmark$\checkmark^\checkmark\\checkmarkc\checkmarki\checkmarkr\checkmarkc\checkmark$\checkmarkt\checkmark$\checkmark^\checkmark\\checkmarkc\checkmarki\checkmarkr\checkmarkc\checkmark$\checkmark \checkmark$\checkmark^\checkmark\\checkmarkc\checkmarki\checkmarkr\checkmarkc\checkmark$\checkmarkc\checkmark$\checkmark^\checkmark\\checkmarkc\checkmarki\checkmarkr\checkmarkc\checkmark$\checkmarko\checkmark$\checkmark^\checkmark\\checkmarkc\checkmarki\checkmarkr\checkmarkc\checkmark$\checkmarkn\checkmark$\checkmark^\checkmark\\checkmarkc\checkmarki\checkmarkr\checkmarkc\checkmark$\checkmarkt\checkmark$\checkmark^\checkmark\\checkmarkc\checkmarki\checkmarkr\checkmarkc\checkmark$\checkmarki\checkmark$\checkmark^\checkmark\\checkmarkc\checkmarki\checkmarkr\checkmarkc\checkmark$\checkmarkn\checkmark$\checkmark^\checkmark\\checkmarkc\checkmarki\checkmarkr\checkmarkc\checkmark$\checkmarku\checkmark$\checkmark^\checkmark\\checkmarkc\checkmarki\checkmarkr\checkmarkc\checkmark$\checkmark \checkmark$\checkmark^\checkmark\\checkmarkc\checkmarki\checkmarkr\checkmarkc\checkmark$\checkmark&\checkmark$\checkmark^\checkmark\\checkmarkc\checkmarki\checkmarkr\checkmarkc\checkmark$\checkmark \checkmark$\checkmark^\checkmark\\checkmarkc\checkmarki\checkmarkr\checkmarkc\checkmark$\checkmarkN\checkmark$\checkmark^\checkmark\\checkmarkc\checkmarki\checkmarkr\checkmarkc\checkmark$\checkmark$\checkmark$\checkmark^\checkmark\\checkmarkc\checkmarki\checkmarkr\checkmarkc\checkmark$\checkmark\\checkmark$\checkmark^\checkmark\\checkmarkc\checkmarki\checkmarkr\checkmarkc\checkmark$\checkmarkc\checkmark$\checkmark^\checkmark\\checkmarkc\checkmarki\checkmarkr\checkmarkc\checkmark$\checkmarkd\checkmark$\checkmark^\checkmark\\checkmarkc\checkmarki\checkmarkr\checkmarkc\checkmark$\checkmarko\checkmark$\checkmark^\checkmark\\checkmarkc\checkmarki\checkmarkr\checkmarkc\checkmark$\checkmarkt\checkmark$\checkmark^\checkmark\\checkmarkc\checkmarki\checkmarkr\checkmarkc\checkmark$\checkmark$\checkmark$\checkmark^\checkmark\\checkmarkc\checkmarki\checkmarkr\checkmarkc\checkmark$\checkmarkm\checkmark$\checkmark^\checkmark\\checkmarkc\checkmarki\checkmarkr\checkmarkc\checkmark$\checkmark \checkmark$\checkmark^\checkmark\\checkmarkc\checkmarki\checkmarkr\checkmarkc\checkmark$\checkmark\\checkmark$\checkmark^\checkmark\\checkmarkc\checkmarki\checkmarkr\checkmarkc\checkmark$\checkmark\\checkmark$\checkmark^\checkmark\\checkmarkc\checkmarki\checkmarkr\checkmarkc\checkmark$\checkmark \checkmark$\checkmark^\checkmark\\checkmarkc\checkmarki\checkmarkr\checkmarkc\checkmark$\checkmark\\checkmark$\checkmark^\checkmark\\checkmarkc\checkmarki\checkmarkr\checkmarkc\checkmark$\checkmarkh\checkmark$\checkmark^\checkmark\\checkmarkc\checkmarki\checkmarkr\checkmarkc\checkmark$\checkmarkl\checkmark$\checkmark^\checkmark\\checkmarkc\checkmarki\checkmarkr\checkmarkc\checkmark$\checkmarki\checkmark$\checkmark^\checkmark\\checkmarkc\checkmarki\checkmarkr\checkmarkc\checkmark$\checkmarkn\checkmark$\checkmark^\checkmark\\checkmarkc\checkmarki\checkmarkr\checkmarkc\checkmark$\checkmarke\checkmark$\checkmark^\checkmark\\checkmarkc\checkmarki\checkmarkr\checkmarkc\checkmark$\checkmark
\checkmark$\checkmark^\checkmark\\checkmarkc\checkmarki\checkmarkr\checkmarkc\checkmark$\checkmark$\checkmark$\checkmark^\checkmark\\checkmarkc\checkmarki\checkmarkr\checkmarkc\checkmark$\checkmarkC\checkmark$\checkmark^\checkmark\\checkmarkc\checkmarki\checkmarkr\checkmarkc\checkmark$\checkmark_\checkmark$\checkmark^\checkmark\\checkmarkc\checkmarki\checkmarkr\checkmarkc\checkmark$\checkmarkJ\checkmark$\checkmark^\checkmark\\checkmarkc\checkmarki\checkmarkr\checkmarkc\checkmark$\checkmark$\checkmark$\checkmark^\checkmark\\checkmarkc\checkmarki\checkmarkr\checkmarkc\checkmark$\checkmark \checkmark$\checkmark^\checkmark\\checkmarkc\checkmarki\checkmarkr\checkmarkc\checkmark$\checkmark&\checkmark$\checkmark^\checkmark\\checkmarkc\checkmarki\checkmarkr\checkmarkc\checkmark$\checkmark \checkmark$\checkmark^\checkmark\\checkmarkc\checkmarki\checkmarkr\checkmarkc\checkmark$\checkmarkC\checkmark$\checkmark^\checkmark\\checkmarkc\checkmarki\checkmarkr\checkmarkc\checkmark$\checkmarko\checkmark$\checkmark^\checkmark\\checkmarkc\checkmarki\checkmarkr\checkmarkc\checkmark$\checkmarku\checkmark$\checkmark^\checkmark\\checkmarkc\checkmarki\checkmarkr\checkmarkc\checkmark$\checkmarkp\checkmark$\checkmark^\checkmark\\checkmarkc\checkmarki\checkmarkr\checkmarkc\checkmark$\checkmarkl\checkmark$\checkmark^\checkmark\\checkmarkc\checkmarki\checkmarkr\checkmarkc\checkmark$\checkmarke\checkmark$\checkmark^\checkmark\\checkmarkc\checkmarki\checkmarkr\checkmarkc\checkmark$\checkmark \checkmark$\checkmark^\checkmark\\checkmarkc\checkmarki\checkmarkr\checkmarkc\checkmark$\checkmarkd\checkmark$\checkmark^\checkmark\\checkmarkc\checkmarki\checkmarkr\checkmarkc\checkmark$\checkmark'\checkmark$\checkmark^\checkmark\\checkmarkc\checkmarki\checkmarkr\checkmarkc\checkmark$\checkmarki\checkmark$\checkmark^\checkmark\\checkmarkc\checkmarki\checkmarkr\checkmarkc\checkmark$\checkmarkn\checkmark$\checkmark^\checkmark\\checkmarkc\checkmarki\checkmarkr\checkmarkc\checkmark$\checkmarke\checkmark$\checkmark^\checkmark\\checkmarkc\checkmarki\checkmarkr\checkmarkc\checkmark$\checkmarkr\checkmark$\checkmark^\checkmark\\checkmarkc\checkmarki\checkmarkr\checkmarkc\checkmark$\checkmarkt\checkmark$\checkmark^\checkmark\\checkmarkc\checkmarki\checkmarkr\checkmarkc\checkmark$\checkmarki\checkmark$\checkmark^\checkmark\\checkmarkc\checkmarki\checkmarkr\checkmarkc\checkmark$\checkmarke\checkmark$\checkmark^\checkmark\\checkmarkc\checkmarki\checkmarkr\checkmarkc\checkmark$\checkmark \checkmark$\checkmark^\checkmark\\checkmarkc\checkmarki\checkmarkr\checkmarkc\checkmark$\checkmarkl\checkmark$\checkmark^\checkmark\\checkmarkc\checkmarki\checkmarkr\checkmarkc\checkmark$\checkmarko\checkmark$\checkmark^\checkmark\\checkmarkc\checkmarki\checkmarkr\checkmarkc\checkmark$\checkmarkr\checkmark$\checkmark^\checkmark\\checkmarkc\checkmarki\checkmarkr\checkmarkc\checkmark$\checkmarks\checkmark$\checkmark^\checkmark\\checkmarkc\checkmarki\checkmarkr\checkmarkc\checkmark$\checkmark \checkmark$\checkmark^\checkmark\\checkmarkc\checkmarki\checkmarkr\checkmarkc\checkmark$\checkmarkd\checkmark$\checkmark^\checkmark\\checkmarkc\checkmarki\checkmarkr\checkmarkc\checkmark$\checkmarke\checkmark$\checkmark^\checkmark\\checkmarkc\checkmarki\checkmarkr\checkmarkc\checkmark$\checkmarks\checkmark$\checkmark^\checkmark\\checkmarkc\checkmarki\checkmarkr\checkmarkc\checkmark$\checkmark \checkmark$\checkmark^\checkmark\\checkmarkc\checkmarki\checkmarkr\checkmarkc\checkmark$\checkmarka\checkmark$\checkmark^\checkmark\\checkmarkc\checkmarki\checkmarkr\checkmarkc\checkmark$\checkmarkc\checkmark$\checkmark^\checkmark\\checkmarkc\checkmarki\checkmarkr\checkmarkc\checkmark$\checkmarkc\checkmark$\checkmark^\checkmark\\checkmarkc\checkmarki\checkmarkr\checkmarkc\checkmark$\checkmarké\checkmark$\checkmark^\checkmark\\checkmarkc\checkmarki\checkmarkr\checkmarkc\checkmark$\checkmarkl\checkmark$\checkmark^\checkmark\\checkmarkc\checkmarki\checkmarkr\checkmarkc\checkmark$\checkmarké\checkmark$\checkmark^\checkmark\\checkmarkc\checkmarki\checkmarkr\checkmarkc\checkmark$\checkmarkr\checkmark$\checkmark^\checkmark\\checkmarkc\checkmarki\checkmarkr\checkmarkc\checkmark$\checkmarka\checkmark$\checkmark^\checkmark\\checkmarkc\checkmarki\checkmarkr\checkmarkc\checkmark$\checkmarkt\checkmark$\checkmark^\checkmark\\checkmarkc\checkmarki\checkmarkr\checkmarkc\checkmark$\checkmarki\checkmark$\checkmark^\checkmark\\checkmarkc\checkmarki\checkmarkr\checkmarkc\checkmark$\checkmarko\checkmark$\checkmark^\checkmark\\checkmarkc\checkmarki\checkmarkr\checkmarkc\checkmark$\checkmarkn\checkmark$\checkmark^\checkmark\\checkmarkc\checkmarki\checkmarkr\checkmarkc\checkmark$\checkmarks\checkmark$\checkmark^\checkmark\\checkmarkc\checkmarki\checkmarkr\checkmarkc\checkmark$\checkmark \checkmark$\checkmark^\checkmark\\checkmarkc\checkmarki\checkmarkr\checkmarkc\checkmark$\checkmark&\checkmark$\checkmark^\checkmark\\checkmarkc\checkmarki\checkmarkr\checkmarkc\checkmark$\checkmark \checkmark$\checkmark^\checkmark\\checkmarkc\checkmarki\checkmarkr\checkmarkc\checkmark$\checkmarkN\checkmark$\checkmark^\checkmark\\checkmarkc\checkmarki\checkmarkr\checkmarkc\checkmark$\checkmark$\checkmark$\checkmark^\checkmark\\checkmarkc\checkmarki\checkmarkr\checkmarkc\checkmark$\checkmark\\checkmark$\checkmark^\checkmark\\checkmarkc\checkmarki\checkmarkr\checkmarkc\checkmark$\checkmarkc\checkmark$\checkmark^\checkmark\\checkmarkc\checkmarki\checkmarkr\checkmarkc\checkmark$\checkmarkd\checkmark$\checkmark^\checkmark\\checkmarkc\checkmarki\checkmarkr\checkmarkc\checkmark$\checkmarko\checkmark$\checkmark^\checkmark\\checkmarkc\checkmarki\checkmarkr\checkmarkc\checkmark$\checkmarkt\checkmark$\checkmark^\checkmark\\checkmarkc\checkmarki\checkmarkr\checkmarkc\checkmark$\checkmark$\checkmark$\checkmark^\checkmark\\checkmarkc\checkmarki\checkmarkr\checkmarkc\checkmark$\checkmarkm\checkmark$\checkmark^\checkmark\\checkmarkc\checkmarki\checkmarkr\checkmarkc\checkmark$\checkmark \checkmark$\checkmark^\checkmark\\checkmarkc\checkmarki\checkmarkr\checkmarkc\checkmark$\checkmark\\checkmark$\checkmark^\checkmark\\checkmarkc\checkmarki\checkmarkr\checkmarkc\checkmark$\checkmark\\checkmark$\checkmark^\checkmark\\checkmarkc\checkmarki\checkmarkr\checkmarkc\checkmark$\checkmark \checkmark$\checkmark^\checkmark\\checkmarkc\checkmarki\checkmarkr\checkmarkc\checkmark$\checkmark\\checkmark$\checkmark^\checkmark\\checkmarkc\checkmarki\checkmarkr\checkmarkc\checkmark$\checkmarkh\checkmark$\checkmark^\checkmark\\checkmarkc\checkmarki\checkmarkr\checkmarkc\checkmark$\checkmarkl\checkmark$\checkmark^\checkmark\\checkmarkc\checkmarki\checkmarkr\checkmarkc\checkmark$\checkmarki\checkmark$\checkmark^\checkmark\\checkmarkc\checkmarki\checkmarkr\checkmarkc\checkmark$\checkmarkn\checkmark$\checkmark^\checkmark\\checkmarkc\checkmarki\checkmarkr\checkmarkc\checkmark$\checkmarke\checkmark$\checkmark^\checkmark\\checkmarkc\checkmarki\checkmarkr\checkmarkc\checkmark$\checkmark
\checkmark$\checkmark^\checkmark\\checkmarkc\checkmarki\checkmarkr\checkmarkc\checkmark$\checkmark\\checkmark$\checkmark^\checkmark\\checkmarkc\checkmarki\checkmarkr\checkmarkc\checkmark$\checkmarke\checkmark$\checkmark^\checkmark\\checkmarkc\checkmarki\checkmarkr\checkmarkc\checkmark$\checkmarkn\checkmark$\checkmark^\checkmark\\checkmarkc\checkmarki\checkmarkr\checkmarkc\checkmark$\checkmarkd\checkmark$\checkmark^\checkmark\\checkmarkc\checkmarki\checkmarkr\checkmarkc\checkmark$\checkmark{\checkmark$\checkmark^\checkmark\\checkmarkc\checkmarki\checkmarkr\checkmarkc\checkmark$\checkmarkt\checkmark$\checkmark^\checkmark\\checkmarkc\checkmarki\checkmarkr\checkmarkc\checkmark$\checkmarka\checkmark$\checkmark^\checkmark\\checkmarkc\checkmarki\checkmarkr\checkmarkc\checkmark$\checkmarkb\checkmark$\checkmark^\checkmark\\checkmarkc\checkmarki\checkmarkr\checkmarkc\checkmark$\checkmarku\checkmark$\checkmark^\checkmark\\checkmarkc\checkmarki\checkmarkr\checkmarkc\checkmark$\checkmarkl\checkmark$\checkmark^\checkmark\\checkmarkc\checkmarki\checkmarkr\checkmarkc\checkmark$\checkmarka\checkmark$\checkmark^\checkmark\\checkmarkc\checkmarki\checkmarkr\checkmarkc\checkmark$\checkmarkr\checkmark$\checkmark^\checkmark\\checkmarkc\checkmarki\checkmarkr\checkmarkc\checkmark$\checkmark}\checkmark$\checkmark^\checkmark\\checkmarkc\checkmarki\checkmarkr\checkmarkc\checkmark$\checkmark
\checkmark$\checkmark^\checkmark\\checkmarkc\checkmarki\checkmarkr\checkmarkc\checkmark$\checkmark\\checkmark$\checkmark^\checkmark\\checkmarkc\checkmarki\checkmarkr\checkmarkc\checkmark$\checkmarkc\checkmark$\checkmark^\checkmark\\checkmarkc\checkmarki\checkmarkr\checkmarkc\checkmark$\checkmarka\checkmark$\checkmark^\checkmark\\checkmarkc\checkmarki\checkmarkr\checkmarkc\checkmark$\checkmarkp\checkmark$\checkmark^\checkmark\\checkmarkc\checkmarki\checkmarkr\checkmarkc\checkmark$\checkmarkt\checkmark$\checkmark^\checkmark\\checkmarkc\checkmarki\checkmarkr\checkmarkc\checkmark$\checkmarki\checkmark$\checkmark^\checkmark\\checkmarkc\checkmarki\checkmarkr\checkmarkc\checkmark$\checkmarko\checkmark$\checkmark^\checkmark\\checkmarkc\checkmarki\checkmarkr\checkmarkc\checkmark$\checkmarkn\checkmark$\checkmark^\checkmark\\checkmarkc\checkmarki\checkmarkr\checkmarkc\checkmark$\checkmark{\checkmark$\checkmark^\checkmark\\checkmarkc\checkmarki\checkmarkr\checkmarkc\checkmark$\checkmarkN\checkmark$\checkmark^\checkmark\\checkmarkc\checkmarki\checkmarkr\checkmarkc\checkmark$\checkmarko\checkmark$\checkmark^\checkmark\\checkmarkc\checkmarki\checkmarkr\checkmarkc\checkmark$\checkmarkm\checkmark$\checkmark^\checkmark\\checkmarkc\checkmarki\checkmarkr\checkmarkc\checkmark$\checkmarke\checkmark$\checkmark^\checkmark\\checkmarkc\checkmarki\checkmarkr\checkmarkc\checkmark$\checkmarkn\checkmark$\checkmark^\checkmark\\checkmarkc\checkmarki\checkmarkr\checkmarkc\checkmark$\checkmarkc\checkmark$\checkmark^\checkmark\\checkmarkc\checkmarki\checkmarkr\checkmarkc\checkmark$\checkmarkl\checkmark$\checkmark^\checkmark\\checkmarkc\checkmarki\checkmarkr\checkmarkc\checkmark$\checkmarka\checkmark$\checkmark^\checkmark\\checkmarkc\checkmarki\checkmarkr\checkmarkc\checkmark$\checkmarkt\checkmark$\checkmark^\checkmark\\checkmarkc\checkmarki\checkmarkr\checkmarkc\checkmark$\checkmarku\checkmark$\checkmark^\checkmark\\checkmarkc\checkmarki\checkmarkr\checkmarkc\checkmark$\checkmarkr\checkmark$\checkmark^\checkmark\\checkmarkc\checkmarki\checkmarkr\checkmarkc\checkmark$\checkmarke\checkmark$\checkmark^\checkmark\\checkmarkc\checkmarki\checkmarkr\checkmarkc\checkmark$\checkmark \checkmark$\checkmark^\checkmark\\checkmarkc\checkmarki\checkmarkr\checkmarkc\checkmark$\checkmarkd\checkmark$\checkmark^\checkmark\\checkmarkc\checkmarki\checkmarkr\checkmarkc\checkmark$\checkmarke\checkmark$\checkmark^\checkmark\\checkmarkc\checkmarki\checkmarkr\checkmarkc\checkmark$\checkmarks\checkmark$\checkmark^\checkmark\\checkmarkc\checkmarki\checkmarkr\checkmarkc\checkmark$\checkmark \checkmark$\checkmark^\checkmark\\checkmarkc\checkmarki\checkmarkr\checkmarkc\checkmark$\checkmarke\checkmark$\checkmark^\checkmark\\checkmarkc\checkmarki\checkmarkr\checkmarkc\checkmark$\checkmarkf\checkmark$\checkmark^\checkmark\\checkmarkc\checkmarki\checkmarkr\checkmarkc\checkmark$\checkmarkf\checkmark$\checkmark^\checkmark\\checkmarkc\checkmarki\checkmarkr\checkmarkc\checkmark$\checkmarko\checkmark$\checkmark^\checkmark\\checkmarkc\checkmarki\checkmarkr\checkmarkc\checkmark$\checkmarkr\checkmark$\checkmark^\checkmark\\checkmarkc\checkmarki\checkmarkr\checkmarkc\checkmark$\checkmarkt\checkmark$\checkmark^\checkmark\\checkmarkc\checkmarki\checkmarkr\checkmarkc\checkmark$\checkmarks\checkmark$\checkmark^\checkmark\\checkmarkc\checkmarki\checkmarkr\checkmarkc\checkmark$\checkmark \checkmark$\checkmark^\checkmark\\checkmarkc\checkmarki\checkmarkr\checkmarkc\checkmark$\checkmarks\checkmark$\checkmark^\checkmark\\checkmarkc\checkmarki\checkmarkr\checkmarkc\checkmark$\checkmarku\checkmark$\checkmark^\checkmark\\checkmarkc\checkmarki\checkmarkr\checkmarkc\checkmark$\checkmarkr\checkmark$\checkmark^\checkmark\\checkmarkc\checkmarki\checkmarkr\checkmarkc\checkmark$\checkmark \checkmark$\checkmark^\checkmark\\checkmarkc\checkmarki\checkmarkr\checkmarkc\checkmark$\checkmarkl\checkmark$\checkmark^\checkmark\\checkmarkc\checkmarki\checkmarkr\checkmarkc\checkmark$\checkmarka\checkmark$\checkmark^\checkmark\\checkmarkc\checkmarki\checkmarkr\checkmarkc\checkmark$\checkmark \checkmark$\checkmark^\checkmark\\checkmarkc\checkmarki\checkmarkr\checkmarkc\checkmark$\checkmarkv\checkmark$\checkmark^\checkmark\\checkmarkc\checkmarki\checkmarkr\checkmarkc\checkmark$\checkmarki\checkmark$\checkmark^\checkmark\\checkmarkc\checkmarki\checkmarkr\checkmarkc\checkmark$\checkmarks\checkmark$\checkmark^\checkmark\\checkmarkc\checkmarki\checkmarkr\checkmarkc\checkmark$\checkmark \checkmark$\checkmark^\checkmark\\checkmarkc\checkmarki\checkmarkr\checkmarkc\checkmark$\checkmarkà\checkmark$\checkmark^\checkmark\\checkmarkc\checkmarki\checkmarkr\checkmarkc\checkmark$\checkmark \checkmark$\checkmark^\checkmark\\checkmarkc\checkmarki\checkmarkr\checkmarkc\checkmark$\checkmarkb\checkmark$\checkmark^\checkmark\\checkmarkc\checkmarki\checkmarkr\checkmarkc\checkmark$\checkmarki\checkmark$\checkmark^\checkmark\\checkmarkc\checkmarki\checkmarkr\checkmarkc\checkmark$\checkmarkl\checkmark$\checkmark^\checkmark\\checkmarkc\checkmarki\checkmarkr\checkmarkc\checkmark$\checkmarkl\checkmark$\checkmark^\checkmark\\checkmarkc\checkmarki\checkmarkr\checkmarkc\checkmark$\checkmarke\checkmark$\checkmark^\checkmark\\checkmarkc\checkmarki\checkmarkr\checkmarkc\checkmark$\checkmarks\checkmark$\checkmark^\checkmark\\checkmarkc\checkmarki\checkmarkr\checkmarkc\checkmark$\checkmark \checkmark$\checkmark^\checkmark\\checkmarkc\checkmarki\checkmarkr\checkmarkc\checkmark$\checkmark(\checkmark$\checkmark^\checkmark\\checkmarkc\checkmarki\checkmarkr\checkmarkc\checkmark$\checkmarkA\checkmark$\checkmark^\checkmark\\checkmarkc\checkmarki\checkmarkr\checkmarkc\checkmark$\checkmarkx\checkmark$\checkmark^\checkmark\\checkmarkc\checkmarki\checkmarkr\checkmarkc\checkmark$\checkmarke\checkmark$\checkmark^\checkmark\\checkmarkc\checkmarki\checkmarkr\checkmarkc\checkmark$\checkmark \checkmark$\checkmark^\checkmark\\checkmarkc\checkmarki\checkmarkr\checkmarkc\checkmark$\checkmarkX\checkmark$\checkmark^\checkmark\\checkmarkc\checkmarki\checkmarkr\checkmarkc\checkmark$\checkmark)\checkmark$\checkmark^\checkmark\\checkmarkc\checkmarki\checkmarkr\checkmarkc\checkmark$\checkmark}\checkmark$\checkmark^\checkmark\\checkmarkc\checkmarki\checkmarkr\checkmarkc\checkmark$\checkmark
\checkmark$\checkmark^\checkmark\\checkmarkc\checkmarki\checkmarkr\checkmarkc\checkmark$\checkmark\\checkmark$\checkmark^\checkmark\\checkmarkc\checkmarki\checkmarkr\checkmarkc\checkmark$\checkmarkl\checkmark$\checkmark^\checkmark\\checkmarkc\checkmarki\checkmarkr\checkmarkc\checkmark$\checkmarka\checkmark$\checkmark^\checkmark\\checkmarkc\checkmarki\checkmarkr\checkmarkc\checkmark$\checkmarkb\checkmark$\checkmark^\checkmark\\checkmarkc\checkmarki\checkmarkr\checkmarkc\checkmark$\checkmarke\checkmark$\checkmark^\checkmark\\checkmarkc\checkmarki\checkmarkr\checkmarkc\checkmark$\checkmarkl\checkmark$\checkmark^\checkmark\\checkmarkc\checkmarki\checkmarkr\checkmarkc\checkmark$\checkmark{\checkmark$\checkmark^\checkmark\\checkmarkc\checkmarki\checkmarkr\checkmarkc\checkmark$\checkmarkt\checkmark$\checkmark^\checkmark\\checkmarkc\checkmarki\checkmarkr\checkmarkc\checkmark$\checkmarka\checkmark$\checkmark^\checkmark\\checkmarkc\checkmarki\checkmarkr\checkmarkc\checkmark$\checkmarkb\checkmark$\checkmark^\checkmark\\checkmarkc\checkmarki\checkmarkr\checkmarkc\checkmark$\checkmark:\checkmark$\checkmark^\checkmark\\checkmarkc\checkmarki\checkmarkr\checkmarkc\checkmark$\checkmarkn\checkmark$\checkmark^\checkmark\\checkmarkc\checkmarki\checkmarkr\checkmarkc\checkmark$\checkmarko\checkmark$\checkmark^\checkmark\\checkmarkc\checkmarki\checkmarkr\checkmarkc\checkmark$\checkmarkm\checkmark$\checkmark^\checkmark\\checkmarkc\checkmarki\checkmarkr\checkmarkc\checkmark$\checkmarke\checkmark$\checkmark^\checkmark\\checkmarkc\checkmarki\checkmarkr\checkmarkc\checkmark$\checkmarkn\checkmark$\checkmark^\checkmark\\checkmarkc\checkmarki\checkmarkr\checkmarkc\checkmark$\checkmarkc\checkmark$\checkmark^\checkmark\\checkmarkc\checkmarki\checkmarkr\checkmarkc\checkmark$\checkmarkl\checkmark$\checkmark^\checkmark\\checkmarkc\checkmarki\checkmarkr\checkmarkc\checkmark$\checkmarka\checkmark$\checkmark^\checkmark\\checkmarkc\checkmarki\checkmarkr\checkmarkc\checkmark$\checkmarkt\checkmark$\checkmark^\checkmark\\checkmarkc\checkmarki\checkmarkr\checkmarkc\checkmark$\checkmarku\checkmark$\checkmark^\checkmark\\checkmarkc\checkmarki\checkmarkr\checkmarkc\checkmark$\checkmarkr\checkmark$\checkmark^\checkmark\\checkmarkc\checkmarki\checkmarkr\checkmarkc\checkmark$\checkmarke\checkmark$\checkmark^\checkmark\\checkmarkc\checkmarki\checkmarkr\checkmarkc\checkmark$\checkmark_\checkmark$\checkmark^\checkmark\\checkmarkc\checkmarki\checkmarkr\checkmarkc\checkmark$\checkmarkv\checkmark$\checkmark^\checkmark\\checkmarkc\checkmarki\checkmarkr\checkmarkc\checkmark$\checkmarki\checkmark$\checkmark^\checkmark\\checkmarkc\checkmarki\checkmarkr\checkmarkc\checkmark$\checkmarks\checkmark$\checkmark^\checkmark\\checkmarkc\checkmarki\checkmarkr\checkmarkc\checkmark$\checkmark}\checkmark$\checkmark^\checkmark\\checkmarkc\checkmarki\checkmarkr\checkmarkc\checkmark$\checkmark
\checkmark$\checkmark^\checkmark\\checkmarkc\checkmarki\checkmarkr\checkmarkc\checkmark$\checkmark\\checkmark$\checkmark^\checkmark\\checkmarkc\checkmarki\checkmarkr\checkmarkc\checkmark$\checkmarke\checkmark$\checkmark^\checkmark\\checkmarkc\checkmarki\checkmarkr\checkmarkc\checkmark$\checkmarkn\checkmark$\checkmark^\checkmark\\checkmarkc\checkmarki\checkmarkr\checkmarkc\checkmark$\checkmarkd\checkmark$\checkmark^\checkmark\\checkmarkc\checkmarki\checkmarkr\checkmarkc\checkmark$\checkmark{\checkmark$\checkmark^\checkmark\\checkmarkc\checkmarki\checkmarkr\checkmarkc\checkmark$\checkmarkt\checkmark$\checkmark^\checkmark\\checkmarkc\checkmarki\checkmarkr\checkmarkc\checkmark$\checkmarka\checkmark$\checkmark^\checkmark\\checkmarkc\checkmarki\checkmarkr\checkmarkc\checkmark$\checkmarkb\checkmark$\checkmark^\checkmark\\checkmarkc\checkmarki\checkmarkr\checkmarkc\checkmark$\checkmarkl\checkmark$\checkmark^\checkmark\\checkmarkc\checkmarki\checkmarkr\checkmarkc\checkmark$\checkmarke\checkmark$\checkmark^\checkmark\\checkmarkc\checkmarki\checkmarkr\checkmarkc\checkmark$\checkmark}\checkmark$\checkmark^\checkmark\\checkmarkc\checkmarki\checkmarkr\checkmarkc\checkmark$\checkmark
\checkmark$\checkmark^\checkmark\\checkmarkc\checkmarki\checkmarkr\checkmarkc\checkmark$\checkmark
\checkmark$\checkmark^\checkmark\\checkmarkc\checkmarki\checkmarkr\checkmarkc\checkmark$\checkmark\\checkmark$\checkmark^\checkmark\\checkmarkc\checkmarki\checkmarkr\checkmarkc\checkmark$\checkmarks\checkmark$\checkmark^\checkmark\\checkmarkc\checkmarki\checkmarkr\checkmarkc\checkmark$\checkmarku\checkmark$\checkmark^\checkmark\\checkmarkc\checkmarki\checkmarkr\checkmarkc\checkmark$\checkmarkb\checkmark$\checkmark^\checkmark\\checkmarkc\checkmarki\checkmarkr\checkmarkc\checkmark$\checkmarks\checkmark$\checkmark^\checkmark\\checkmarkc\checkmarki\checkmarkr\checkmarkc\checkmark$\checkmarke\checkmark$\checkmark^\checkmark\\checkmarkc\checkmarki\checkmarkr\checkmarkc\checkmark$\checkmarkc\checkmark$\checkmark^\checkmark\\checkmarkc\checkmarki\checkmarkr\checkmarkc\checkmark$\checkmarkt\checkmark$\checkmark^\checkmark\\checkmarkc\checkmarki\checkmarkr\checkmarkc\checkmark$\checkmarki\checkmark$\checkmark^\checkmark\\checkmarkc\checkmarki\checkmarkr\checkmarkc\checkmark$\checkmarko\checkmark$\checkmark^\checkmark\\checkmarkc\checkmarki\checkmarkr\checkmarkc\checkmark$\checkmarkn\checkmark$\checkmark^\checkmark\\checkmarkc\checkmarki\checkmarkr\checkmarkc\checkmark$\checkmark{\checkmark$\checkmark^\checkmark\\checkmarkc\checkmarki\checkmarkr\checkmarkc\checkmark$\checkmarkV\checkmark$\checkmark^\checkmark\\checkmarkc\checkmarki\checkmarkr\checkmarkc\checkmark$\checkmarké\checkmark$\checkmark^\checkmark\\checkmarkc\checkmarki\checkmarkr\checkmarkc\checkmark$\checkmarkr\checkmark$\checkmark^\checkmark\\checkmarkc\checkmarki\checkmarkr\checkmarkc\checkmark$\checkmarki\checkmark$\checkmark^\checkmark\\checkmarkc\checkmarki\checkmarkr\checkmarkc\checkmark$\checkmarkf\checkmark$\checkmark^\checkmark\\checkmarkc\checkmarki\checkmarkr\checkmarkc\checkmark$\checkmarki\checkmark$\checkmark^\checkmark\\checkmarkc\checkmarki\checkmarkr\checkmarkc\checkmark$\checkmarkc\checkmark$\checkmark^\checkmark\\checkmarkc\checkmarki\checkmarkr\checkmarkc\checkmark$\checkmarka\checkmark$\checkmark^\checkmark\\checkmarkc\checkmarki\checkmarkr\checkmarkc\checkmark$\checkmarkt\checkmark$\checkmark^\checkmark\\checkmarkc\checkmarki\checkmarkr\checkmarkc\checkmark$\checkmarki\checkmark$\checkmark^\checkmark\\checkmarkc\checkmarki\checkmarkr\checkmarkc\checkmark$\checkmarko\checkmark$\checkmark^\checkmark\\checkmarkc\checkmarki\checkmarkr\checkmarkc\checkmark$\checkmarkn\checkmark$\checkmark^\checkmark\\checkmarkc\checkmarki\checkmarkr\checkmarkc\checkmark$\checkmark \checkmark$\checkmark^\checkmark\\checkmarkc\checkmarki\checkmarkr\checkmarkc\checkmark$\checkmarka\checkmark$\checkmark^\checkmark\\checkmarkc\checkmarki\checkmarkr\checkmarkc\checkmark$\checkmarku\checkmark$\checkmark^\checkmark\\checkmarkc\checkmarki\checkmarkr\checkmarkc\checkmark$\checkmark \checkmark$\checkmark^\checkmark\\checkmarkc\checkmarki\checkmarkr\checkmarkc\checkmark$\checkmarkF\checkmark$\checkmark^\checkmark\\checkmarkc\checkmarki\checkmarkr\checkmarkc\checkmark$\checkmarkl\checkmark$\checkmark^\checkmark\\checkmarkc\checkmarki\checkmarkr\checkmarkc\checkmark$\checkmarka\checkmark$\checkmark^\checkmark\\checkmarkc\checkmarki\checkmarkr\checkmarkc\checkmark$\checkmarkm\checkmark$\checkmark^\checkmark\\checkmarkc\checkmarki\checkmarkr\checkmarkc\checkmark$\checkmarkb\checkmark$\checkmark^\checkmark\\checkmarkc\checkmarki\checkmarkr\checkmarkc\checkmark$\checkmarka\checkmark$\checkmark^\checkmark\\checkmarkc\checkmarki\checkmarkr\checkmarkc\checkmark$\checkmarkg\checkmark$\checkmark^\checkmark\\checkmarkc\checkmarki\checkmarkr\checkmarkc\checkmark$\checkmarke\checkmark$\checkmark^\checkmark\\checkmarkc\checkmarki\checkmarkr\checkmarkc\checkmark$\checkmark \checkmark$\checkmark^\checkmark\\checkmarkc\checkmarki\checkmarkr\checkmarkc\checkmark$\checkmark(\checkmark$\checkmark^\checkmark\\checkmarkc\checkmarki\checkmarkr\checkmarkc\checkmark$\checkmarkC\checkmark$\checkmark^\checkmark\\checkmarkc\checkmarki\checkmarkr\checkmarkc\checkmark$\checkmarkr\checkmark$\checkmark^\checkmark\\checkmarkc\checkmarki\checkmarkr\checkmarkc\checkmark$\checkmarki\checkmark$\checkmark^\checkmark\\checkmarkc\checkmarki\checkmarkr\checkmarkc\checkmark$\checkmarkt\checkmark$\checkmark^\checkmark\\checkmarkc\checkmarki\checkmarkr\checkmarkc\checkmark$\checkmarkè\checkmark$\checkmark^\checkmark\\checkmarkc\checkmarki\checkmarkr\checkmarkc\checkmark$\checkmarkr\checkmark$\checkmark^\checkmark\\checkmarkc\checkmarki\checkmarkr\checkmarkc\checkmark$\checkmarke\checkmark$\checkmark^\checkmark\\checkmarkc\checkmarki\checkmarkr\checkmarkc\checkmark$\checkmark \checkmark$\checkmark^\checkmark\\checkmarkc\checkmarki\checkmarkr\checkmarkc\checkmark$\checkmarkd\checkmark$\checkmark^\checkmark\\checkmarkc\checkmarki\checkmarkr\checkmarkc\checkmark$\checkmark'\checkmark$\checkmark^\checkmark\\checkmarkc\checkmarki\checkmarkr\checkmarkc\checkmark$\checkmarkE\checkmark$\checkmark^\checkmark\\checkmarkc\checkmarki\checkmarkr\checkmarkc\checkmark$\checkmarku\checkmark$\checkmark^\checkmark\\checkmarkc\checkmarki\checkmarkr\checkmarkc\checkmark$\checkmarkl\checkmark$\checkmark^\checkmark\\checkmarkc\checkmarki\checkmarkr\checkmarkc\checkmark$\checkmarke\checkmark$\checkmark^\checkmark\\checkmarkc\checkmarki\checkmarkr\checkmarkc\checkmark$\checkmarkr\checkmark$\checkmark^\checkmark\\checkmarkc\checkmarki\checkmarkr\checkmarkc\checkmark$\checkmark)\checkmark$\checkmark^\checkmark\\checkmarkc\checkmarki\checkmarkr\checkmarkc\checkmark$\checkmark}\checkmark$\checkmark^\checkmark\\checkmarkc\checkmarki\checkmarkr\checkmarkc\checkmark$\checkmark
\checkmark$\checkmark^\checkmark\\checkmarkc\checkmarki\checkmarkr\checkmarkc\checkmark$\checkmarkL\checkmark$\checkmark^\checkmark\\checkmarkc\checkmarki\checkmarkr\checkmarkc\checkmark$\checkmarka\checkmark$\checkmark^\checkmark\\checkmarkc\checkmarki\checkmarkr\checkmarkc\checkmark$\checkmark \checkmark$\checkmark^\checkmark\\checkmarkc\checkmarki\checkmarkr\checkmarkc\checkmark$\checkmarkv\checkmark$\checkmark^\checkmark\\checkmarkc\checkmarki\checkmarkr\checkmarkc\checkmark$\checkmarki\checkmark$\checkmark^\checkmark\\checkmarkc\checkmarki\checkmarkr\checkmarkc\checkmark$\checkmarks\checkmark$\checkmark^\checkmark\\checkmarkc\checkmarki\checkmarkr\checkmarkc\checkmark$\checkmark \checkmark$\checkmark^\checkmark\\checkmarkc\checkmarki\checkmarkr\checkmarkc\checkmark$\checkmarke\checkmark$\checkmark^\checkmark\\checkmarkc\checkmarki\checkmarkr\checkmarkc\checkmark$\checkmarks\checkmark$\checkmark^\checkmark\\checkmarkc\checkmarki\checkmarkr\checkmarkc\checkmark$\checkmarkt\checkmark$\checkmark^\checkmark\\checkmarkc\checkmarki\checkmarkr\checkmarkc\checkmark$\checkmark \checkmark$\checkmark^\checkmark\\checkmarkc\checkmarki\checkmarkr\checkmarkc\checkmark$\checkmarkm\checkmark$\checkmark^\checkmark\\checkmarkc\checkmarki\checkmarkr\checkmarkc\checkmark$\checkmarko\checkmark$\checkmark^\checkmark\\checkmarkc\checkmarki\checkmarkr\checkmarkc\checkmark$\checkmarkd\checkmark$\checkmark^\checkmark\\checkmarkc\checkmarki\checkmarkr\checkmarkc\checkmark$\checkmarké\checkmark$\checkmark^\checkmark\\checkmarkc\checkmarki\checkmarkr\checkmarkc\checkmark$\checkmarkl\checkmark$\checkmark^\checkmark\\checkmarkc\checkmarki\checkmarkr\checkmarkc\checkmark$\checkmarki\checkmark$\checkmark^\checkmark\\checkmarkc\checkmarki\checkmarkr\checkmarkc\checkmark$\checkmarks\checkmark$\checkmark^\checkmark\\checkmarkc\checkmarki\checkmarkr\checkmarkc\checkmark$\checkmarké\checkmark$\checkmark^\checkmark\\checkmarkc\checkmarki\checkmarkr\checkmarkc\checkmark$\checkmarke\checkmark$\checkmark^\checkmark\\checkmarkc\checkmarki\checkmarkr\checkmarkc\checkmark$\checkmark \checkmark$\checkmark^\checkmark\\checkmarkc\checkmarki\checkmarkr\checkmarkc\checkmark$\checkmarkc\checkmark$\checkmark^\checkmark\\checkmarkc\checkmarki\checkmarkr\checkmarkc\checkmark$\checkmarko\checkmark$\checkmark^\checkmark\\checkmarkc\checkmarki\checkmarkr\checkmarkc\checkmark$\checkmarkm\checkmark$\checkmark^\checkmark\\checkmarkc\checkmarki\checkmarkr\checkmarkc\checkmark$\checkmarkm\checkmark$\checkmark^\checkmark\\checkmarkc\checkmarki\checkmarkr\checkmarkc\checkmark$\checkmarke\checkmark$\checkmark^\checkmark\\checkmarkc\checkmarki\checkmarkr\checkmarkc\checkmark$\checkmark \checkmark$\checkmark^\checkmark\\checkmarkc\checkmarki\checkmarkr\checkmarkc\checkmark$\checkmarku\checkmark$\checkmark^\checkmark\\checkmarkc\checkmarki\checkmarkr\checkmarkc\checkmark$\checkmarkn\checkmark$\checkmark^\checkmark\\checkmarkc\checkmarki\checkmarkr\checkmarkc\checkmark$\checkmarke\checkmark$\checkmark^\checkmark\\checkmarkc\checkmarki\checkmarkr\checkmarkc\checkmark$\checkmark \checkmark$\checkmark^\checkmark\\checkmarkc\checkmarki\checkmarkr\checkmarkc\checkmark$\checkmarkp\checkmark$\checkmark^\checkmark\\checkmarkc\checkmarki\checkmarkr\checkmarkc\checkmark$\checkmarko\checkmark$\checkmark^\checkmark\\checkmarkc\checkmarki\checkmarkr\checkmarkc\checkmark$\checkmarku\checkmark$\checkmark^\checkmark\\checkmarkc\checkmarki\checkmarkr\checkmarkc\checkmark$\checkmarkt\checkmark$\checkmark^\checkmark\\checkmarkc\checkmarki\checkmarkr\checkmarkc\checkmark$\checkmarkr\checkmark$\checkmark^\checkmark\\checkmarkc\checkmarki\checkmarkr\checkmarkc\checkmark$\checkmarke\checkmark$\checkmark^\checkmark\\checkmarkc\checkmarki\checkmarkr\checkmarkc\checkmark$\checkmark \checkmark$\checkmark^\checkmark\\checkmarkc\checkmarki\checkmarkr\checkmarkc\checkmark$\checkmarkc\checkmark$\checkmark^\checkmark\\checkmarkc\checkmarki\checkmarkr\checkmarkc\checkmark$\checkmarky\checkmark$\checkmark^\checkmark\\checkmarkc\checkmarki\checkmarkr\checkmarkc\checkmark$\checkmarkl\checkmark$\checkmark^\checkmark\\checkmarkc\checkmarki\checkmarkr\checkmarkc\checkmark$\checkmarki\checkmark$\checkmark^\checkmark\\checkmarkc\checkmarki\checkmarkr\checkmarkc\checkmark$\checkmarkn\checkmark$\checkmark^\checkmark\\checkmarkc\checkmarki\checkmarkr\checkmarkc\checkmark$\checkmarkd\checkmark$\checkmark^\checkmark\\checkmarkc\checkmarki\checkmarkr\checkmarkc\checkmark$\checkmarkr\checkmark$\checkmark^\checkmark\\checkmarkc\checkmarki\checkmarkr\checkmarkc\checkmark$\checkmarki\checkmark$\checkmark^\checkmark\\checkmarkc\checkmarki\checkmarkr\checkmarkc\checkmark$\checkmarkq\checkmark$\checkmark^\checkmark\\checkmarkc\checkmarki\checkmarkr\checkmarkc\checkmark$\checkmarku\checkmark$\checkmark^\checkmark\\checkmarkc\checkmarki\checkmarkr\checkmarkc\checkmark$\checkmarke\checkmark$\checkmark^\checkmark\\checkmarkc\checkmarki\checkmarkr\checkmarkc\checkmark$\checkmark \checkmark$\checkmark^\checkmark\\checkmarkc\checkmarki\checkmarkr\checkmarkc\checkmark$\checkmarkd\checkmark$\checkmark^\checkmark\\checkmarkc\checkmarki\checkmarkr\checkmarkc\checkmark$\checkmarke\checkmark$\checkmark^\checkmark\\checkmarkc\checkmarki\checkmarkr\checkmarkc\checkmark$\checkmark \checkmark$\checkmark^\checkmark\\checkmarkc\checkmarki\checkmarkr\checkmarkc\checkmark$\checkmarkd\checkmark$\checkmark^\checkmark\\checkmarkc\checkmarki\checkmarkr\checkmarkc\checkmark$\checkmarki\checkmark$\checkmark^\checkmark\\checkmarkc\checkmarki\checkmarkr\checkmarkc\checkmark$\checkmarka\checkmark$\checkmark^\checkmark\\checkmarkc\checkmarki\checkmarkr\checkmarkc\checkmark$\checkmarkm\checkmark$\checkmark^\checkmark\\checkmarkc\checkmarki\checkmarkr\checkmarkc\checkmark$\checkmarkè\checkmark$\checkmark^\checkmark\\checkmarkc\checkmarki\checkmarkr\checkmarkc\checkmark$\checkmarkt\checkmark$\checkmark^\checkmark\\checkmarkc\checkmarki\checkmarkr\checkmarkc\checkmark$\checkmarkr\checkmark$\checkmark^\checkmark\\checkmarkc\checkmarki\checkmarkr\checkmarkc\checkmark$\checkmarke\checkmark$\checkmark^\checkmark\\checkmarkc\checkmarki\checkmarkr\checkmarkc\checkmark$\checkmark \checkmark$\checkmark^\checkmark\\checkmarkc\checkmarki\checkmarkr\checkmarkc\checkmark$\checkmark$\checkmark$\checkmark^\checkmark\\checkmarkc\checkmarki\checkmarkr\checkmarkc\checkmark$\checkmarkd\checkmark$\checkmark^\checkmark\\checkmarkc\checkmarki\checkmarkr\checkmarkc\checkmark$\checkmark_\checkmark$\checkmark^\checkmark\\checkmarkc\checkmarki\checkmarkr\checkmarkc\checkmark$\checkmark{\checkmark$\checkmark^\checkmark\\checkmarkc\checkmarki\checkmarkr\checkmarkc\checkmark$\checkmarkv\checkmark$\checkmark^\checkmark\\checkmarkc\checkmarki\checkmarkr\checkmarkc\checkmark$\checkmarki\checkmark$\checkmark^\checkmark\\checkmarkc\checkmarki\checkmarkr\checkmarkc\checkmark$\checkmarks\checkmark$\checkmark^\checkmark\\checkmarkc\checkmarki\checkmarkr\checkmarkc\checkmark$\checkmark}\checkmark$\checkmark^\checkmark\\checkmarkc\checkmarki\checkmarkr\checkmarkc\checkmark$\checkmark \checkmark$\checkmark^\checkmark\\checkmarkc\checkmarki\checkmarkr\checkmarkc\checkmark$\checkmark=\checkmark$\checkmark^\checkmark\\checkmarkc\checkmarki\checkmarkr\checkmarkc\checkmark$\checkmark \checkmark$\checkmark^\checkmark\\checkmarkc\checkmarki\checkmarkr\checkmarkc\checkmark$\checkmark1\checkmark$\checkmark^\checkmark\\checkmarkc\checkmarki\checkmarkr\checkmarkc\checkmark$\checkmark2\checkmark$\checkmark^\checkmark\\checkmarkc\checkmarki\checkmarkr\checkmarkc\checkmark$\checkmark$\checkmark$\checkmark^\checkmark\\checkmarkc\checkmarki\checkmarkr\checkmarkc\checkmark$\checkmark \checkmark$\checkmark^\checkmark\\checkmarkc\checkmarki\checkmarkr\checkmarkc\checkmark$\checkmarkm\checkmark$\checkmark^\checkmark\\checkmarkc\checkmarki\checkmarkr\checkmarkc\checkmark$\checkmarkm\checkmark$\checkmark^\checkmark\\checkmarkc\checkmarki\checkmarkr\checkmarkc\checkmark$\checkmark \checkmark$\checkmark^\checkmark\\checkmarkc\checkmarki\checkmarkr\checkmarkc\checkmark$\checkmark(\checkmark$\checkmark^\checkmark\\checkmarkc\checkmarki\checkmarkr\checkmarkc\checkmark$\checkmarkd\checkmark$\checkmark^\checkmark\\checkmarkc\checkmarki\checkmarkr\checkmarkc\checkmark$\checkmarki\checkmark$\checkmark^\checkmark\\checkmarkc\checkmarki\checkmarkr\checkmarkc\checkmark$\checkmarka\checkmark$\checkmark^\checkmark\\checkmarkc\checkmarki\checkmarkr\checkmarkc\checkmark$\checkmarkm\checkmark$\checkmark^\checkmark\\checkmarkc\checkmarki\checkmarkr\checkmarkc\checkmark$\checkmarkè\checkmark$\checkmark^\checkmark\\checkmarkc\checkmarki\checkmarkr\checkmarkc\checkmark$\checkmarkt\checkmark$\checkmark^\checkmark\\checkmarkc\checkmarki\checkmarkr\checkmarkc\checkmark$\checkmarkr\checkmark$\checkmark^\checkmark\\checkmarkc\checkmarki\checkmarkr\checkmarkc\checkmark$\checkmarke\checkmark$\checkmark^\checkmark\\checkmarkc\checkmarki\checkmarkr\checkmarkc\checkmark$\checkmark \checkmark$\checkmark^\checkmark\\checkmarkc\checkmarki\checkmarkr\checkmarkc\checkmark$\checkmarkd\checkmark$\checkmark^\checkmark\\checkmarkc\checkmarki\checkmarkr\checkmarkc\checkmark$\checkmarku\checkmark$\checkmark^\checkmark\\checkmarkc\checkmarki\checkmarkr\checkmarkc\checkmark$\checkmark \checkmark$\checkmark^\checkmark\\checkmarkc\checkmarki\checkmarkr\checkmarkc\checkmark$\checkmarkn\checkmark$\checkmark^\checkmark\\checkmarkc\checkmarki\checkmarkr\checkmarkc\checkmark$\checkmarko\checkmark$\checkmark^\checkmark\\checkmarkc\checkmarki\checkmarkr\checkmarkc\checkmark$\checkmarky\checkmark$\checkmark^\checkmark\\checkmarkc\checkmarki\checkmarkr\checkmarkc\checkmark$\checkmarka\checkmark$\checkmark^\checkmark\\checkmarkc\checkmarki\checkmarkr\checkmarkc\checkmark$\checkmarku\checkmark$\checkmark^\checkmark\\checkmarkc\checkmarki\checkmarkr\checkmarkc\checkmark$\checkmark \checkmark$\checkmark^\checkmark\\checkmarkc\checkmarki\checkmarkr\checkmarkc\checkmark$\checkmarkf\checkmark$\checkmark^\checkmark\\checkmarkc\checkmarki\checkmarkr\checkmarkc\checkmark$\checkmarki\checkmark$\checkmark^\checkmark\\checkmarkc\checkmarki\checkmarkr\checkmarkc\checkmark$\checkmarkl\checkmark$\checkmark^\checkmark\\checkmarkc\checkmarki\checkmarkr\checkmarkc\checkmark$\checkmarke\checkmark$\checkmark^\checkmark\\checkmarkc\checkmarki\checkmarkr\checkmarkc\checkmark$\checkmarkt\checkmark$\checkmark^\checkmark\\checkmarkc\checkmarki\checkmarkr\checkmarkc\checkmark$\checkmark \checkmark$\checkmark^\checkmark\\checkmarkc\checkmarki\checkmarkr\checkmarkc\checkmark$\checkmarkS\checkmark$\checkmark^\checkmark\\checkmarkc\checkmarki\checkmarkr\checkmarkc\checkmark$\checkmarkF\checkmark$\checkmark^\checkmark\\checkmarkc\checkmarki\checkmarkr\checkmarkc\checkmark$\checkmarkU\checkmark$\checkmark^\checkmark\\checkmarkc\checkmarki\checkmarkr\checkmarkc\checkmark$\checkmark1\checkmark$\checkmark^\checkmark\\checkmarkc\checkmarki\checkmarkr\checkmarkc\checkmark$\checkmark6\checkmark$\checkmark^\checkmark\\checkmarkc\checkmarki\checkmarkr\checkmarkc\checkmark$\checkmark0\checkmark$\checkmark^\checkmark\\checkmarkc\checkmarki\checkmarkr\checkmarkc\checkmark$\checkmark5\checkmark$\checkmark^\checkmark\\checkmarkc\checkmarki\checkmarkr\checkmarkc\checkmark$\checkmark)\checkmark$\checkmark^\checkmark\\checkmarkc\checkmarki\checkmarkr\checkmarkc\checkmark$\checkmark,\checkmark$\checkmark^\checkmark\\checkmarkc\checkmarki\checkmarkr\checkmarkc\checkmark$\checkmark \checkmark$\checkmark^\checkmark\\checkmarkc\checkmarki\checkmarkr\checkmarkc\checkmark$\checkmarks\checkmark$\checkmark^\checkmark\\checkmarkc\checkmarki\checkmarkr\checkmarkc\checkmark$\checkmarko\checkmark$\checkmark^\checkmark\\checkmarkc\checkmarki\checkmarkr\checkmarkc\checkmark$\checkmarku\checkmark$\checkmark^\checkmark\\checkmarkc\checkmarki\checkmarkr\checkmarkc\checkmark$\checkmarkm\checkmark$\checkmark^\checkmark\\checkmarkc\checkmarki\checkmarkr\checkmarkc\checkmark$\checkmarki\checkmark$\checkmark^\checkmark\\checkmarkc\checkmarki\checkmarkr\checkmarkc\checkmark$\checkmarks\checkmark$\checkmark^\checkmark\\checkmarkc\checkmarki\checkmarkr\checkmarkc\checkmark$\checkmarke\checkmark$\checkmark^\checkmark\\checkmarkc\checkmarki\checkmarkr\checkmarkc\checkmark$\checkmark \checkmark$\checkmark^\checkmark\\checkmarkc\checkmarki\checkmarkr\checkmarkc\checkmark$\checkmarkà\checkmark$\checkmark^\checkmark\\checkmarkc\checkmarki\checkmarkr\checkmarkc\checkmark$\checkmark \checkmark$\checkmark^\checkmark\\checkmarkc\checkmarki\checkmarkr\checkmarkc\checkmark$\checkmarku\checkmark$\checkmark^\checkmark\\checkmarkc\checkmarki\checkmarkr\checkmarkc\checkmark$\checkmarkn\checkmark$\checkmark^\checkmark\\checkmarkc\checkmarki\checkmarkr\checkmarkc\checkmark$\checkmarke\checkmark$\checkmark^\checkmark\\checkmarkc\checkmarki\checkmarkr\checkmarkc\checkmark$\checkmark \checkmark$\checkmark^\checkmark\\checkmarkc\checkmarki\checkmarkr\checkmarkc\checkmark$\checkmarkc\checkmark$\checkmark^\checkmark\\checkmarkc\checkmarki\checkmarkr\checkmarkc\checkmark$\checkmarko\checkmark$\checkmark^\checkmark\\checkmarkc\checkmarki\checkmarkr\checkmarkc\checkmark$\checkmarkm\checkmark$\checkmark^\checkmark\\checkmarkc\checkmarki\checkmarkr\checkmarkc\checkmark$\checkmarkp\checkmark$\checkmark^\checkmark\\checkmarkc\checkmarki\checkmarkr\checkmarkc\checkmark$\checkmarkr\checkmark$\checkmark^\checkmark\\checkmarkc\checkmarki\checkmarkr\checkmarkc\checkmark$\checkmarke\checkmark$\checkmark^\checkmark\\checkmarkc\checkmarki\checkmarkr\checkmarkc\checkmark$\checkmarks\checkmark$\checkmark^\checkmark\\checkmarkc\checkmarki\checkmarkr\checkmarkc\checkmark$\checkmarks\checkmark$\checkmark^\checkmark\\checkmarkc\checkmarki\checkmarkr\checkmarkc\checkmark$\checkmarki\checkmark$\checkmark^\checkmark\\checkmarkc\checkmarki\checkmarkr\checkmarkc\checkmark$\checkmarko\checkmark$\checkmark^\checkmark\\checkmarkc\checkmarki\checkmarkr\checkmarkc\checkmark$\checkmarkn\checkmark$\checkmark^\checkmark\\checkmarkc\checkmarki\checkmarkr\checkmarkc\checkmark$\checkmark \checkmark$\checkmark^\checkmark\\checkmarkc\checkmarki\checkmarkr\checkmarkc\checkmark$\checkmarka\checkmark$\checkmark^\checkmark\\checkmarkc\checkmarki\checkmarkr\checkmarkc\checkmark$\checkmarkx\checkmark$\checkmark^\checkmark\\checkmarkc\checkmarki\checkmarkr\checkmarkc\checkmark$\checkmarki\checkmark$\checkmark^\checkmark\\checkmarkc\checkmarki\checkmarkr\checkmarkc\checkmark$\checkmarka\checkmark$\checkmark^\checkmark\\checkmarkc\checkmarki\checkmarkr\checkmarkc\checkmark$\checkmarkl\checkmark$\checkmark^\checkmark\\checkmarkc\checkmarki\checkmarkr\checkmarkc\checkmark$\checkmarke\checkmark$\checkmark^\checkmark\\checkmarkc\checkmarki\checkmarkr\checkmarkc\checkmark$\checkmark.\checkmark$\checkmark^\checkmark\\checkmarkc\checkmarki\checkmarkr\checkmarkc\checkmark$\checkmark
\checkmark$\checkmark^\checkmark\\checkmarkc\checkmarki\checkmarkr\checkmarkc\checkmark$\checkmark
\checkmark$\checkmark^\checkmark\\checkmarkc\checkmarki\checkmarkr\checkmarkc\checkmark$\checkmarkL\checkmark$\checkmark^\checkmark\\checkmarkc\checkmarki\checkmarkr\checkmarkc\checkmark$\checkmarke\checkmark$\checkmark^\checkmark\\checkmarkc\checkmarki\checkmarkr\checkmarkc\checkmark$\checkmark \checkmark$\checkmark^\checkmark\\checkmarkc\checkmarki\checkmarkr\checkmarkc\checkmark$\checkmarkm\checkmark$\checkmark^\checkmark\\checkmarkc\checkmarki\checkmarkr\checkmarkc\checkmark$\checkmarko\checkmark$\checkmark^\checkmark\\checkmarkc\checkmarki\checkmarkr\checkmarkc\checkmark$\checkmarkm\checkmark$\checkmark^\checkmark\\checkmarkc\checkmarki\checkmarkr\checkmarkc\checkmark$\checkmarke\checkmark$\checkmark^\checkmark\\checkmarkc\checkmarki\checkmarkr\checkmarkc\checkmark$\checkmarkn\checkmark$\checkmark^\checkmark\\checkmarkc\checkmarki\checkmarkr\checkmarkc\checkmark$\checkmarkt\checkmark$\checkmark^\checkmark\\checkmarkc\checkmarki\checkmarkr\checkmarkc\checkmark$\checkmark \checkmark$\checkmark^\checkmark\\checkmarkc\checkmarki\checkmarkr\checkmarkc\checkmark$\checkmarkq\checkmark$\checkmark^\checkmark\\checkmarkc\checkmarki\checkmarkr\checkmarkc\checkmark$\checkmarku\checkmark$\checkmark^\checkmark\\checkmarkc\checkmarki\checkmarkr\checkmarkc\checkmark$\checkmarka\checkmark$\checkmark^\checkmark\\checkmarkc\checkmarki\checkmarkr\checkmarkc\checkmark$\checkmarkd\checkmark$\checkmark^\checkmark\\checkmarkc\checkmarki\checkmarkr\checkmarkc\checkmark$\checkmarkr\checkmark$\checkmark^\checkmark\\checkmarkc\checkmarki\checkmarkr\checkmarkc\checkmark$\checkmarka\checkmark$\checkmark^\checkmark\\checkmarkc\checkmarki\checkmarkr\checkmarkc\checkmark$\checkmarkt\checkmark$\checkmark^\checkmark\\checkmarkc\checkmarki\checkmarkr\checkmarkc\checkmark$\checkmarki\checkmark$\checkmark^\checkmark\\checkmarkc\checkmarki\checkmarkr\checkmarkc\checkmark$\checkmarkq\checkmark$\checkmark^\checkmark\\checkmarkc\checkmarki\checkmarkr\checkmarkc\checkmark$\checkmarku\checkmark$\checkmark^\checkmark\\checkmarkc\checkmarki\checkmarkr\checkmarkc\checkmark$\checkmarke\checkmark$\checkmark^\checkmark\\checkmarkc\checkmarki\checkmarkr\checkmarkc\checkmark$\checkmark \checkmark$\checkmark^\checkmark\\checkmarkc\checkmarki\checkmarkr\checkmarkc\checkmark$\checkmarkd\checkmark$\checkmark^\checkmark\\checkmarkc\checkmarki\checkmarkr\checkmarkc\checkmark$\checkmarke\checkmark$\checkmark^\checkmark\\checkmarkc\checkmarki\checkmarkr\checkmarkc\checkmark$\checkmark \checkmark$\checkmark^\checkmark\\checkmarkc\checkmarki\checkmarkr\checkmarkc\checkmark$\checkmarkl\checkmark$\checkmark^\checkmark\\checkmarkc\checkmarki\checkmarkr\checkmarkc\checkmark$\checkmarka\checkmark$\checkmark^\checkmark\\checkmarkc\checkmarki\checkmarkr\checkmarkc\checkmark$\checkmark \checkmark$\checkmark^\checkmark\\checkmarkc\checkmarki\checkmarkr\checkmarkc\checkmark$\checkmarks\checkmark$\checkmark^\checkmark\\checkmarkc\checkmarki\checkmarkr\checkmarkc\checkmark$\checkmarke\checkmark$\checkmark^\checkmark\\checkmarkc\checkmarki\checkmarkr\checkmarkc\checkmark$\checkmarkc\checkmark$\checkmark^\checkmark\\checkmarkc\checkmarki\checkmarkr\checkmarkc\checkmark$\checkmarkt\checkmark$\checkmark^\checkmark\\checkmarkc\checkmarki\checkmarkr\checkmarkc\checkmark$\checkmarki\checkmark$\checkmark^\checkmark\\checkmarkc\checkmarki\checkmarkr\checkmarkc\checkmark$\checkmarko\checkmark$\checkmark^\checkmark\\checkmarkc\checkmarki\checkmarkr\checkmarkc\checkmark$\checkmarkn\checkmark$\checkmark^\checkmark\\checkmarkc\checkmarki\checkmarkr\checkmarkc\checkmark$\checkmark \checkmark$\checkmark^\checkmark\\checkmarkc\checkmarki\checkmarkr\checkmarkc\checkmark$\checkmarkc\checkmark$\checkmark^\checkmark\\checkmarkc\checkmarki\checkmarkr\checkmarkc\checkmark$\checkmarki\checkmark$\checkmark^\checkmark\\checkmarkc\checkmarki\checkmarkr\checkmarkc\checkmark$\checkmarkr\checkmark$\checkmark^\checkmark\\checkmarkc\checkmarki\checkmarkr\checkmarkc\checkmark$\checkmarkc\checkmark$\checkmark^\checkmark\\checkmarkc\checkmarki\checkmarkr\checkmarkc\checkmark$\checkmarku\checkmark$\checkmark^\checkmark\\checkmarkc\checkmarki\checkmarkr\checkmarkc\checkmark$\checkmarkl\checkmark$\checkmark^\checkmark\\checkmarkc\checkmarki\checkmarkr\checkmarkc\checkmark$\checkmarka\checkmark$\checkmark^\checkmark\\checkmarkc\checkmarki\checkmarkr\checkmarkc\checkmark$\checkmarki\checkmark$\checkmark^\checkmark\\checkmarkc\checkmarki\checkmarkr\checkmarkc\checkmark$\checkmarkr\checkmark$\checkmark^\checkmark\\checkmarkc\checkmarki\checkmarkr\checkmarkc\checkmark$\checkmarke\checkmark$\checkmark^\checkmark\\checkmarkc\checkmarki\checkmarkr\checkmarkc\checkmark$\checkmark \checkmark$\checkmark^\checkmark\\checkmarkc\checkmarki\checkmarkr\checkmarkc\checkmark$\checkmark:\checkmark$\checkmark^\checkmark\\checkmarkc\checkmarki\checkmarkr\checkmarkc\checkmark$\checkmark
\checkmark$\checkmark^\checkmark\\checkmarkc\checkmarki\checkmarkr\checkmarkc\checkmark$\checkmark\\checkmark$\checkmark^\checkmark\\checkmarkc\checkmarki\checkmarkr\checkmarkc\checkmark$\checkmarkb\checkmark$\checkmark^\checkmark\\checkmarkc\checkmarki\checkmarkr\checkmarkc\checkmark$\checkmarke\checkmark$\checkmark^\checkmark\\checkmarkc\checkmarki\checkmarkr\checkmarkc\checkmark$\checkmarkg\checkmark$\checkmark^\checkmark\\checkmarkc\checkmarki\checkmarkr\checkmarkc\checkmark$\checkmarki\checkmark$\checkmark^\checkmark\\checkmarkc\checkmarki\checkmarkr\checkmarkc\checkmark$\checkmarkn\checkmark$\checkmark^\checkmark\\checkmarkc\checkmarki\checkmarkr\checkmarkc\checkmark$\checkmark{\checkmark$\checkmark^\checkmark\\checkmarkc\checkmarki\checkmarkr\checkmarkc\checkmark$\checkmarke\checkmark$\checkmark^\checkmark\\checkmarkc\checkmarki\checkmarkr\checkmarkc\checkmark$\checkmarkq\checkmark$\checkmark^\checkmark\\checkmarkc\checkmarki\checkmarkr\checkmarkc\checkmark$\checkmarku\checkmark$\checkmark^\checkmark\\checkmarkc\checkmarki\checkmarkr\checkmarkc\checkmark$\checkmarka\checkmark$\checkmark^\checkmark\\checkmarkc\checkmarki\checkmarkr\checkmarkc\checkmark$\checkmarkt\checkmark$\checkmark^\checkmark\\checkmarkc\checkmarki\checkmarkr\checkmarkc\checkmark$\checkmarki\checkmark$\checkmark^\checkmark\\checkmarkc\checkmarki\checkmarkr\checkmarkc\checkmark$\checkmarko\checkmark$\checkmark^\checkmark\\checkmarkc\checkmarki\checkmarkr\checkmarkc\checkmark$\checkmarkn\checkmark$\checkmark^\checkmark\\checkmarkc\checkmarki\checkmarkr\checkmarkc\checkmark$\checkmark}\checkmark$\checkmark^\checkmark\\checkmarkc\checkmarki\checkmarkr\checkmarkc\checkmark$\checkmark
\checkmark$\checkmark^\checkmark\\checkmarkc\checkmarki\checkmarkr\checkmarkc\checkmark$\checkmark \checkmark$\checkmark^\checkmark\\checkmarkc\checkmarki\checkmarkr\checkmarkc\checkmark$\checkmark \checkmark$\checkmark^\checkmark\\checkmarkc\checkmarki\checkmarkr\checkmarkc\checkmark$\checkmark \checkmark$\checkmark^\checkmark\\checkmarkc\checkmarki\checkmarkr\checkmarkc\checkmark$\checkmark \checkmark$\checkmark^\checkmark\\checkmarkc\checkmarki\checkmarkr\checkmarkc\checkmark$\checkmarkI\checkmark$\checkmark^\checkmark\\checkmarkc\checkmarki\checkmarkr\checkmarkc\checkmark$\checkmark_\checkmark$\checkmark^\checkmark\\checkmarkc\checkmarki\checkmarkr\checkmarkc\checkmark$\checkmark{\checkmark$\checkmark^\checkmark\\checkmarkc\checkmarki\checkmarkr\checkmarkc\checkmark$\checkmarky\checkmark$\checkmark^\checkmark\\checkmarkc\checkmarki\checkmarkr\checkmarkc\checkmark$\checkmarky\checkmark$\checkmark^\checkmark\\checkmarkc\checkmarki\checkmarkr\checkmarkc\checkmark$\checkmark}\checkmark$\checkmark^\checkmark\\checkmarkc\checkmarki\checkmarkr\checkmarkc\checkmark$\checkmark \checkmark$\checkmark^\checkmark\\checkmarkc\checkmarki\checkmarkr\checkmarkc\checkmark$\checkmark=\checkmark$\checkmark^\checkmark\\checkmarkc\checkmarki\checkmarkr\checkmarkc\checkmark$\checkmark \checkmark$\checkmark^\checkmark\\checkmarkc\checkmarki\checkmarkr\checkmarkc\checkmark$\checkmark\\checkmark$\checkmark^\checkmark\\checkmarkc\checkmarki\checkmarkr\checkmarkc\checkmark$\checkmarkf\checkmark$\checkmark^\checkmark\\checkmarkc\checkmarki\checkmarkr\checkmarkc\checkmark$\checkmarkr\checkmark$\checkmark^\checkmark\\checkmarkc\checkmarki\checkmarkr\checkmarkc\checkmark$\checkmarka\checkmark$\checkmark^\checkmark\\checkmarkc\checkmarki\checkmarkr\checkmarkc\checkmark$\checkmarkc\checkmark$\checkmark^\checkmark\\checkmarkc\checkmarki\checkmarkr\checkmarkc\checkmark$\checkmark{\checkmark$\checkmark^\checkmark\\checkmarkc\checkmarki\checkmarkr\checkmarkc\checkmark$\checkmark\\checkmark$\checkmark^\checkmark\\checkmarkc\checkmarki\checkmarkr\checkmarkc\checkmark$\checkmarkp\checkmark$\checkmark^\checkmark\\checkmarkc\checkmarki\checkmarkr\checkmarkc\checkmark$\checkmarki\checkmark$\checkmark^\checkmark\\checkmarkc\checkmarki\checkmarkr\checkmarkc\checkmark$\checkmark \checkmark$\checkmark^\checkmark\\checkmarkc\checkmarki\checkmarkr\checkmarkc\checkmark$\checkmark\\checkmark$\checkmark^\checkmark\\checkmarkc\checkmarki\checkmarkr\checkmarkc\checkmark$\checkmarkc\checkmark$\checkmark^\checkmark\\checkmarkc\checkmarki\checkmarkr\checkmarkc\checkmark$\checkmarkd\checkmark$\checkmark^\checkmark\\checkmarkc\checkmarki\checkmarkr\checkmarkc\checkmark$\checkmarko\checkmark$\checkmark^\checkmark\\checkmarkc\checkmarki\checkmarkr\checkmarkc\checkmark$\checkmarkt\checkmark$\checkmark^\checkmark\\checkmarkc\checkmarki\checkmarkr\checkmarkc\checkmark$\checkmark \checkmark$\checkmark^\checkmark\\checkmarkc\checkmarki\checkmarkr\checkmarkc\checkmark$\checkmarkd\checkmark$\checkmark^\checkmark\\checkmarkc\checkmarki\checkmarkr\checkmarkc\checkmark$\checkmark_\checkmark$\checkmark^\checkmark\\checkmarkc\checkmarki\checkmarkr\checkmarkc\checkmark$\checkmark{\checkmark$\checkmark^\checkmark\\checkmarkc\checkmarki\checkmarkr\checkmarkc\checkmark$\checkmarkv\checkmark$\checkmark^\checkmark\\checkmarkc\checkmarki\checkmarkr\checkmarkc\checkmark$\checkmarki\checkmark$\checkmark^\checkmark\\checkmarkc\checkmarki\checkmarkr\checkmarkc\checkmark$\checkmarks\checkmark$\checkmark^\checkmark\\checkmarkc\checkmarki\checkmarkr\checkmarkc\checkmark$\checkmark}\checkmark$\checkmark^\checkmark\\checkmarkc\checkmarki\checkmarkr\checkmarkc\checkmark$\checkmark^\checkmark$\checkmark^\checkmark\\checkmarkc\checkmarki\checkmarkr\checkmarkc\checkmark$\checkmark4\checkmark$\checkmark^\checkmark\\checkmarkc\checkmarki\checkmarkr\checkmarkc\checkmark$\checkmark}\checkmark$\checkmark^\checkmark\\checkmarkc\checkmarki\checkmarkr\checkmarkc\checkmark$\checkmark{\checkmark$\checkmark^\checkmark\\checkmarkc\checkmarki\checkmarkr\checkmarkc\checkmark$\checkmark6\checkmark$\checkmark^\checkmark\\checkmarkc\checkmarki\checkmarkr\checkmarkc\checkmark$\checkmark4\checkmark$\checkmark^\checkmark\\checkmarkc\checkmarki\checkmarkr\checkmarkc\checkmark$\checkmark}\checkmark$\checkmark^\checkmark\\checkmarkc\checkmarki\checkmarkr\checkmarkc\checkmark$\checkmark \checkmark$\checkmark^\checkmark\\checkmarkc\checkmarki\checkmarkr\checkmarkc\checkmark$\checkmark=\checkmark$\checkmark^\checkmark\\checkmarkc\checkmarki\checkmarkr\checkmarkc\checkmark$\checkmark \checkmark$\checkmark^\checkmark\\checkmarkc\checkmarki\checkmarkr\checkmarkc\checkmark$\checkmark\\checkmark$\checkmark^\checkmark\\checkmarkc\checkmarki\checkmarkr\checkmarkc\checkmark$\checkmarkf\checkmark$\checkmark^\checkmark\\checkmarkc\checkmarki\checkmarkr\checkmarkc\checkmark$\checkmarkr\checkmark$\checkmark^\checkmark\\checkmarkc\checkmarki\checkmarkr\checkmarkc\checkmark$\checkmarka\checkmark$\checkmark^\checkmark\\checkmarkc\checkmarki\checkmarkr\checkmarkc\checkmark$\checkmarkc\checkmark$\checkmark^\checkmark\\checkmarkc\checkmarki\checkmarkr\checkmarkc\checkmark$\checkmark{\checkmark$\checkmark^\checkmark\\checkmarkc\checkmarki\checkmarkr\checkmarkc\checkmark$\checkmark\\checkmark$\checkmark^\checkmark\\checkmarkc\checkmarki\checkmarkr\checkmarkc\checkmark$\checkmarkp\checkmark$\checkmark^\checkmark\\checkmarkc\checkmarki\checkmarkr\checkmarkc\checkmark$\checkmarki\checkmark$\checkmark^\checkmark\\checkmarkc\checkmarki\checkmarkr\checkmarkc\checkmark$\checkmark \checkmark$\checkmark^\checkmark\\checkmarkc\checkmarki\checkmarkr\checkmarkc\checkmark$\checkmark\\checkmark$\checkmark^\checkmark\\checkmarkc\checkmarki\checkmarkr\checkmarkc\checkmark$\checkmarkt\checkmark$\checkmark^\checkmark\\checkmarkc\checkmarki\checkmarkr\checkmarkc\checkmark$\checkmarki\checkmark$\checkmark^\checkmark\\checkmarkc\checkmarki\checkmarkr\checkmarkc\checkmark$\checkmarkm\checkmark$\checkmark^\checkmark\\checkmarkc\checkmarki\checkmarkr\checkmarkc\checkmark$\checkmarke\checkmark$\checkmark^\checkmark\\checkmarkc\checkmarki\checkmarkr\checkmarkc\checkmark$\checkmarks\checkmark$\checkmark^\checkmark\\checkmarkc\checkmarki\checkmarkr\checkmarkc\checkmark$\checkmark \checkmark$\checkmark^\checkmark\\checkmarkc\checkmarki\checkmarkr\checkmarkc\checkmark$\checkmark1\checkmark$\checkmark^\checkmark\\checkmarkc\checkmarki\checkmarkr\checkmarkc\checkmark$\checkmark2\checkmark$\checkmark^\checkmark\\checkmarkc\checkmarki\checkmarkr\checkmarkc\checkmark$\checkmark^\checkmark$\checkmark^\checkmark\\checkmarkc\checkmarki\checkmarkr\checkmarkc\checkmark$\checkmark4\checkmark$\checkmark^\checkmark\\checkmarkc\checkmarki\checkmarkr\checkmarkc\checkmark$\checkmark}\checkmark$\checkmark^\checkmark\\checkmarkc\checkmarki\checkmarkr\checkmarkc\checkmark$\checkmark{\checkmark$\checkmark^\checkmark\\checkmarkc\checkmarki\checkmarkr\checkmarkc\checkmark$\checkmark6\checkmark$\checkmark^\checkmark\\checkmarkc\checkmarki\checkmarkr\checkmarkc\checkmark$\checkmark4\checkmark$\checkmark^\checkmark\\checkmarkc\checkmarki\checkmarkr\checkmarkc\checkmark$\checkmark}\checkmark$\checkmark^\checkmark\\checkmarkc\checkmarki\checkmarkr\checkmarkc\checkmark$\checkmark \checkmark$\checkmark^\checkmark\\checkmarkc\checkmarki\checkmarkr\checkmarkc\checkmark$\checkmark=\checkmark$\checkmark^\checkmark\\checkmarkc\checkmarki\checkmarkr\checkmarkc\checkmark$\checkmark \checkmark$\checkmark^\checkmark\\checkmarkc\checkmarki\checkmarkr\checkmarkc\checkmark$\checkmark1\checkmark$\checkmark^\checkmark\\checkmarkc\checkmarki\checkmarkr\checkmarkc\checkmark$\checkmark\\checkmark$\checkmark^\checkmark\\checkmarkc\checkmarki\checkmarkr\checkmarkc\checkmark$\checkmark,\checkmark$\checkmark^\checkmark\\checkmarkc\checkmarki\checkmarkr\checkmarkc\checkmark$\checkmark0\checkmark$\checkmark^\checkmark\\checkmarkc\checkmarki\checkmarkr\checkmarkc\checkmark$\checkmark1\checkmark$\checkmark^\checkmark\\checkmarkc\checkmarki\checkmarkr\checkmarkc\checkmark$\checkmark7\checkmark$\checkmark^\checkmark\\checkmarkc\checkmarki\checkmarkr\checkmarkc\checkmark$\checkmark{\checkmark$\checkmark^\checkmark\\checkmarkc\checkmarki\checkmarkr\checkmarkc\checkmark$\checkmark,\checkmark$\checkmark^\checkmark\\checkmarkc\checkmarki\checkmarkr\checkmarkc\checkmark$\checkmark}\checkmark$\checkmark^\checkmark\\checkmarkc\checkmarki\checkmarkr\checkmarkc\checkmark$\checkmark9\checkmark$\checkmark^\checkmark\\checkmarkc\checkmarki\checkmarkr\checkmarkc\checkmark$\checkmark \checkmark$\checkmark^\checkmark\\checkmarkc\checkmarki\checkmarkr\checkmarkc\checkmark$\checkmark\\checkmark$\checkmark^\checkmark\\checkmarkc\checkmarki\checkmarkr\checkmarkc\checkmark$\checkmarkt\checkmark$\checkmark^\checkmark\\checkmarkc\checkmarki\checkmarkr\checkmarkc\checkmark$\checkmarke\checkmark$\checkmark^\checkmark\\checkmarkc\checkmarki\checkmarkr\checkmarkc\checkmark$\checkmarkx\checkmark$\checkmark^\checkmark\\checkmarkc\checkmarki\checkmarkr\checkmarkc\checkmark$\checkmarkt\checkmark$\checkmark^\checkmark\\checkmarkc\checkmarki\checkmarkr\checkmarkc\checkmark$\checkmark{\checkmark$\checkmark^\checkmark\\checkmarkc\checkmarki\checkmarkr\checkmarkc\checkmark$\checkmark \checkmark$\checkmark^\checkmark\\checkmarkc\checkmarki\checkmarkr\checkmarkc\checkmark$\checkmarkm\checkmark$\checkmark^\checkmark\\checkmarkc\checkmarki\checkmarkr\checkmarkc\checkmark$\checkmarkm\checkmark$\checkmark^\checkmark\\checkmarkc\checkmarki\checkmarkr\checkmarkc\checkmark$\checkmark}\checkmark$\checkmark^\checkmark\\checkmarkc\checkmarki\checkmarkr\checkmarkc\checkmark$\checkmark^\checkmark$\checkmark^\checkmark\\checkmarkc\checkmarki\checkmarkr\checkmarkc\checkmark$\checkmark4\checkmark$\checkmark^\checkmark\\checkmarkc\checkmarki\checkmarkr\checkmarkc\checkmark$\checkmark
\checkmark$\checkmark^\checkmark\\checkmarkc\checkmarki\checkmarkr\checkmarkc\checkmark$\checkmark \checkmark$\checkmark^\checkmark\\checkmarkc\checkmarki\checkmarkr\checkmarkc\checkmark$\checkmark \checkmark$\checkmark^\checkmark\\checkmarkc\checkmarki\checkmarkr\checkmarkc\checkmark$\checkmark \checkmark$\checkmark^\checkmark\\checkmarkc\checkmarki\checkmarkr\checkmarkc\checkmark$\checkmark \checkmark$\checkmark^\checkmark\\checkmarkc\checkmarki\checkmarkr\checkmarkc\checkmark$\checkmark\\checkmark$\checkmark^\checkmark\\checkmarkc\checkmarki\checkmarkr\checkmarkc\checkmark$\checkmarkl\checkmark$\checkmark^\checkmark\\checkmarkc\checkmarki\checkmarkr\checkmarkc\checkmark$\checkmarka\checkmark$\checkmark^\checkmark\\checkmarkc\checkmarki\checkmarkr\checkmarkc\checkmark$\checkmarkb\checkmark$\checkmark^\checkmark\\checkmarkc\checkmarki\checkmarkr\checkmarkc\checkmark$\checkmarke\checkmark$\checkmark^\checkmark\\checkmarkc\checkmarki\checkmarkr\checkmarkc\checkmark$\checkmarkl\checkmark$\checkmark^\checkmark\\checkmarkc\checkmarki\checkmarkr\checkmarkc\checkmark$\checkmark{\checkmark$\checkmark^\checkmark\\checkmarkc\checkmarki\checkmarkr\checkmarkc\checkmark$\checkmarke\checkmark$\checkmark^\checkmark\\checkmarkc\checkmarki\checkmarkr\checkmarkc\checkmark$\checkmarkq\checkmark$\checkmark^\checkmark\\checkmarkc\checkmarki\checkmarkr\checkmarkc\checkmark$\checkmark:\checkmark$\checkmark^\checkmark\\checkmarkc\checkmarki\checkmarkr\checkmarkc\checkmark$\checkmarki\checkmark$\checkmark^\checkmark\\checkmarkc\checkmarki\checkmarkr\checkmarkc\checkmark$\checkmarkn\checkmark$\checkmark^\checkmark\\checkmarkc\checkmarki\checkmarkr\checkmarkc\checkmark$\checkmarke\checkmark$\checkmark^\checkmark\\checkmarkc\checkmarki\checkmarkr\checkmarkc\checkmark$\checkmarkr\checkmark$\checkmark^\checkmark\\checkmarkc\checkmarki\checkmarkr\checkmarkc\checkmark$\checkmarkt\checkmark$\checkmark^\checkmark\\checkmarkc\checkmarki\checkmarkr\checkmarkc\checkmark$\checkmarki\checkmark$\checkmark^\checkmark\\checkmarkc\checkmarki\checkmarkr\checkmarkc\checkmark$\checkmarke\checkmark$\checkmark^\checkmark\\checkmarkc\checkmarki\checkmarkr\checkmarkc\checkmark$\checkmark_\checkmark$\checkmark^\checkmark\\checkmarkc\checkmarki\checkmarkr\checkmarkc\checkmark$\checkmarks\checkmark$\checkmark^\checkmark\\checkmarkc\checkmarki\checkmarkr\checkmarkc\checkmark$\checkmarke\checkmark$\checkmark^\checkmark\\checkmarkc\checkmarki\checkmarkr\checkmarkc\checkmark$\checkmarkc\checkmark$\checkmark^\checkmark\\checkmarkc\checkmarki\checkmarkr\checkmarkc\checkmark$\checkmarkt\checkmark$\checkmark^\checkmark\\checkmarkc\checkmarki\checkmarkr\checkmarkc\checkmark$\checkmarki\checkmark$\checkmark^\checkmark\\checkmarkc\checkmarki\checkmarkr\checkmarkc\checkmark$\checkmarko\checkmark$\checkmark^\checkmark\\checkmarkc\checkmarki\checkmarkr\checkmarkc\checkmark$\checkmarkn\checkmark$\checkmark^\checkmark\\checkmarkc\checkmarki\checkmarkr\checkmarkc\checkmark$\checkmark}\checkmark$\checkmark^\checkmark\\checkmarkc\checkmarki\checkmarkr\checkmarkc\checkmark$\checkmark
\checkmark$\checkmark^\checkmark\\checkmarkc\checkmarki\checkmarkr\checkmarkc\checkmark$\checkmark\\checkmark$\checkmark^\checkmark\\checkmarkc\checkmarki\checkmarkr\checkmarkc\checkmark$\checkmarke\checkmark$\checkmark^\checkmark\\checkmarkc\checkmarki\checkmarkr\checkmarkc\checkmark$\checkmarkn\checkmark$\checkmark^\checkmark\\checkmarkc\checkmarki\checkmarkr\checkmarkc\checkmark$\checkmarkd\checkmark$\checkmark^\checkmark\\checkmarkc\checkmarki\checkmarkr\checkmarkc\checkmark$\checkmark{\checkmark$\checkmark^\checkmark\\checkmarkc\checkmarki\checkmarkr\checkmarkc\checkmark$\checkmarke\checkmark$\checkmark^\checkmark\\checkmarkc\checkmarki\checkmarkr\checkmarkc\checkmark$\checkmarkq\checkmark$\checkmark^\checkmark\\checkmarkc\checkmarki\checkmarkr\checkmarkc\checkmark$\checkmarku\checkmark$\checkmark^\checkmark\\checkmarkc\checkmarki\checkmarkr\checkmarkc\checkmark$\checkmarka\checkmark$\checkmark^\checkmark\\checkmarkc\checkmarki\checkmarkr\checkmarkc\checkmark$\checkmarkt\checkmark$\checkmark^\checkmark\\checkmarkc\checkmarki\checkmarkr\checkmarkc\checkmark$\checkmarki\checkmark$\checkmark^\checkmark\\checkmarkc\checkmarki\checkmarkr\checkmarkc\checkmark$\checkmarko\checkmark$\checkmark^\checkmark\\checkmarkc\checkmarki\checkmarkr\checkmarkc\checkmark$\checkmarkn\checkmark$\checkmark^\checkmark\\checkmarkc\checkmarki\checkmarkr\checkmarkc\checkmark$\checkmark}\checkmark$\checkmark^\checkmark\\checkmarkc\checkmarki\checkmarkr\checkmarkc\checkmark$\checkmark
\checkmark$\checkmark^\checkmark\\checkmarkc\checkmarki\checkmarkr\checkmarkc\checkmark$\checkmark
\checkmark$\checkmark^\checkmark\\checkmarkc\checkmarki\checkmarkr\checkmarkc\checkmark$\checkmarkL\checkmark$\checkmark^\checkmark\\checkmarkc\checkmarki\checkmarkr\checkmarkc\checkmark$\checkmarka\checkmark$\checkmark^\checkmark\\checkmarkc\checkmarki\checkmarkr\checkmarkc\checkmark$\checkmark \checkmark$\checkmark^\checkmark\\checkmarkc\checkmarki\checkmarkr\checkmarkc\checkmark$\checkmarkc\checkmark$\checkmark^\checkmark\\checkmarkc\checkmarki\checkmarkr\checkmarkc\checkmark$\checkmarkh\checkmark$\checkmark^\checkmark\\checkmarkc\checkmarki\checkmarkr\checkmarkc\checkmark$\checkmarka\checkmark$\checkmark^\checkmark\\checkmarkc\checkmarki\checkmarkr\checkmarkc\checkmark$\checkmarkr\checkmark$\checkmark^\checkmark\\checkmarkc\checkmarki\checkmarkr\checkmarkc\checkmark$\checkmarkg\checkmark$\checkmark^\checkmark\\checkmarkc\checkmarki\checkmarkr\checkmarkc\checkmark$\checkmarke\checkmark$\checkmark^\checkmark\\checkmarkc\checkmarki\checkmarkr\checkmarkc\checkmark$\checkmark \checkmark$\checkmark^\checkmark\\checkmarkc\checkmarki\checkmarkr\checkmarkc\checkmark$\checkmarkc\checkmark$\checkmark^\checkmark\\checkmarkc\checkmarki\checkmarkr\checkmarkc\checkmark$\checkmarkr\checkmark$\checkmark^\checkmark\\checkmarkc\checkmarki\checkmarkr\checkmarkc\checkmark$\checkmarki\checkmark$\checkmark^\checkmark\\checkmarkc\checkmarki\checkmarkr\checkmarkc\checkmark$\checkmarkt\checkmark$\checkmark^\checkmark\\checkmarkc\checkmarki\checkmarkr\checkmarkc\checkmark$\checkmarki\checkmark$\checkmark^\checkmark\\checkmarkc\checkmarki\checkmarkr\checkmarkc\checkmark$\checkmarkq\checkmark$\checkmark^\checkmark\\checkmarkc\checkmarki\checkmarkr\checkmarkc\checkmark$\checkmarku\checkmark$\checkmark^\checkmark\\checkmarkc\checkmarki\checkmarkr\checkmarkc\checkmark$\checkmarke\checkmark$\checkmark^\checkmark\\checkmarkc\checkmarki\checkmarkr\checkmarkc\checkmark$\checkmark \checkmark$\checkmark^\checkmark\\checkmarkc\checkmarki\checkmarkr\checkmarkc\checkmark$\checkmarkd\checkmark$\checkmark^\checkmark\\checkmarkc\checkmarki\checkmarkr\checkmarkc\checkmark$\checkmarke\checkmark$\checkmark^\checkmark\\checkmarkc\checkmarki\checkmarkr\checkmarkc\checkmark$\checkmark \checkmark$\checkmark^\checkmark\\checkmarkc\checkmarki\checkmarkr\checkmarkc\checkmark$\checkmarkf\checkmark$\checkmark^\checkmark\\checkmarkc\checkmarki\checkmarkr\checkmarkc\checkmark$\checkmarkl\checkmark$\checkmark^\checkmark\\checkmarkc\checkmarki\checkmarkr\checkmarkc\checkmark$\checkmarka\checkmark$\checkmark^\checkmark\\checkmarkc\checkmarki\checkmarkr\checkmarkc\checkmark$\checkmarkm\checkmark$\checkmark^\checkmark\\checkmarkc\checkmarki\checkmarkr\checkmarkc\checkmark$\checkmarkb\checkmark$\checkmark^\checkmark\\checkmarkc\checkmarki\checkmarkr\checkmarkc\checkmark$\checkmarka\checkmark$\checkmark^\checkmark\\checkmarkc\checkmarki\checkmarkr\checkmarkc\checkmark$\checkmarkg\checkmark$\checkmark^\checkmark\\checkmarkc\checkmarki\checkmarkr\checkmarkc\checkmark$\checkmarke\checkmark$\checkmark^\checkmark\\checkmarkc\checkmarki\checkmarkr\checkmarkc\checkmark$\checkmark \checkmark$\checkmark^\checkmark\\checkmarkc\checkmarki\checkmarkr\checkmarkc\checkmark$\checkmark(\checkmark$\checkmark^\checkmark\\checkmarkc\checkmarki\checkmarkr\checkmarkc\checkmark$\checkmarkf\checkmark$\checkmark^\checkmark\\checkmarkc\checkmarki\checkmarkr\checkmarkc\checkmark$\checkmarko\checkmark$\checkmark^\checkmark\\checkmarkc\checkmarki\checkmarkr\checkmarkc\checkmark$\checkmarkr\checkmark$\checkmark^\checkmark\\checkmarkc\checkmarki\checkmarkr\checkmarkc\checkmark$\checkmarkm\checkmark$\checkmark^\checkmark\\checkmarkc\checkmarki\checkmarkr\checkmarkc\checkmark$\checkmarku\checkmark$\checkmark^\checkmark\\checkmarkc\checkmarki\checkmarkr\checkmarkc\checkmark$\checkmarkl\checkmark$\checkmark^\checkmark\\checkmarkc\checkmarki\checkmarkr\checkmarkc\checkmark$\checkmarke\checkmark$\checkmark^\checkmark\\checkmarkc\checkmarki\checkmarkr\checkmarkc\checkmark$\checkmark \checkmark$\checkmark^\checkmark\\checkmarkc\checkmarki\checkmarkr\checkmarkc\checkmark$\checkmarkd\checkmark$\checkmark^\checkmark\\checkmarkc\checkmarki\checkmarkr\checkmarkc\checkmark$\checkmark'\checkmark$\checkmark^\checkmark\\checkmarkc\checkmarki\checkmarkr\checkmarkc\checkmark$\checkmarkE\checkmark$\checkmark^\checkmark\\checkmarkc\checkmarki\checkmarkr\checkmarkc\checkmark$\checkmarku\checkmark$\checkmark^\checkmark\\checkmarkc\checkmarki\checkmarkr\checkmarkc\checkmark$\checkmarkl\checkmark$\checkmark^\checkmark\\checkmarkc\checkmarki\checkmarkr\checkmarkc\checkmark$\checkmarke\checkmark$\checkmark^\checkmark\\checkmarkc\checkmarki\checkmarkr\checkmarkc\checkmark$\checkmarkr\checkmark$\checkmark^\checkmark\\checkmarkc\checkmarki\checkmarkr\checkmarkc\checkmark$\checkmark)\checkmark$\checkmark^\checkmark\\checkmarkc\checkmarki\checkmarkr\checkmarkc\checkmark$\checkmark \checkmark$\checkmark^\checkmark\\checkmarkc\checkmarki\checkmarkr\checkmarkc\checkmark$\checkmark:\checkmark$\checkmark^\checkmark\\checkmarkc\checkmarki\checkmarkr\checkmarkc\checkmark$\checkmark
\checkmark$\checkmark^\checkmark\\checkmarkc\checkmarki\checkmarkr\checkmarkc\checkmark$\checkmark\\checkmark$\checkmark^\checkmark\\checkmarkc\checkmarki\checkmarkr\checkmarkc\checkmark$\checkmarkb\checkmark$\checkmark^\checkmark\\checkmarkc\checkmarki\checkmarkr\checkmarkc\checkmark$\checkmarke\checkmark$\checkmark^\checkmark\\checkmarkc\checkmarki\checkmarkr\checkmarkc\checkmark$\checkmarkg\checkmark$\checkmark^\checkmark\\checkmarkc\checkmarki\checkmarkr\checkmarkc\checkmark$\checkmarki\checkmark$\checkmark^\checkmark\\checkmarkc\checkmarki\checkmarkr\checkmarkc\checkmark$\checkmarkn\checkmark$\checkmark^\checkmark\\checkmarkc\checkmarki\checkmarkr\checkmarkc\checkmark$\checkmark{\checkmark$\checkmark^\checkmark\\checkmarkc\checkmarki\checkmarkr\checkmarkc\checkmark$\checkmarke\checkmark$\checkmark^\checkmark\\checkmarkc\checkmarki\checkmarkr\checkmarkc\checkmark$\checkmarkq\checkmark$\checkmark^\checkmark\\checkmarkc\checkmarki\checkmarkr\checkmarkc\checkmark$\checkmarku\checkmark$\checkmark^\checkmark\\checkmarkc\checkmarki\checkmarkr\checkmarkc\checkmark$\checkmarka\checkmark$\checkmark^\checkmark\\checkmarkc\checkmarki\checkmarkr\checkmarkc\checkmark$\checkmarkt\checkmark$\checkmark^\checkmark\\checkmarkc\checkmarki\checkmarkr\checkmarkc\checkmark$\checkmarki\checkmark$\checkmark^\checkmark\\checkmarkc\checkmarki\checkmarkr\checkmarkc\checkmark$\checkmarko\checkmark$\checkmark^\checkmark\\checkmarkc\checkmarki\checkmarkr\checkmarkc\checkmark$\checkmarkn\checkmark$\checkmark^\checkmark\\checkmarkc\checkmarki\checkmarkr\checkmarkc\checkmark$\checkmark}\checkmark$\checkmark^\checkmark\\checkmarkc\checkmarki\checkmarkr\checkmarkc\checkmark$\checkmark
\checkmark$\checkmark^\checkmark\\checkmarkc\checkmarki\checkmarkr\checkmarkc\checkmark$\checkmark \checkmark$\checkmark^\checkmark\\checkmarkc\checkmarki\checkmarkr\checkmarkc\checkmark$\checkmark \checkmark$\checkmark^\checkmark\\checkmarkc\checkmarki\checkmarkr\checkmarkc\checkmark$\checkmark \checkmark$\checkmark^\checkmark\\checkmarkc\checkmarki\checkmarkr\checkmarkc\checkmark$\checkmark \checkmark$\checkmark^\checkmark\\checkmarkc\checkmarki\checkmarkr\checkmarkc\checkmark$\checkmarkF\checkmark$\checkmark^\checkmark\\checkmarkc\checkmarki\checkmarkr\checkmarkc\checkmark$\checkmark_\checkmark$\checkmark^\checkmark\\checkmarkc\checkmarki\checkmarkr\checkmarkc\checkmark$\checkmarkk\checkmark$\checkmark^\checkmark\\checkmarkc\checkmarki\checkmarkr\checkmarkc\checkmark$\checkmark \checkmark$\checkmark^\checkmark\\checkmarkc\checkmarki\checkmarkr\checkmarkc\checkmark$\checkmark=\checkmark$\checkmark^\checkmark\\checkmarkc\checkmarki\checkmarkr\checkmarkc\checkmark$\checkmark \checkmark$\checkmark^\checkmark\\checkmarkc\checkmarki\checkmarkr\checkmarkc\checkmark$\checkmark\\checkmark$\checkmark^\checkmark\\checkmarkc\checkmarki\checkmarkr\checkmarkc\checkmark$\checkmarkf\checkmark$\checkmark^\checkmark\\checkmarkc\checkmarki\checkmarkr\checkmarkc\checkmark$\checkmarkr\checkmark$\checkmark^\checkmark\\checkmarkc\checkmarki\checkmarkr\checkmarkc\checkmark$\checkmarka\checkmark$\checkmark^\checkmark\\checkmarkc\checkmarki\checkmarkr\checkmarkc\checkmark$\checkmarkc\checkmark$\checkmark^\checkmark\\checkmarkc\checkmarki\checkmarkr\checkmarkc\checkmark$\checkmark{\checkmark$\checkmark^\checkmark\\checkmarkc\checkmarki\checkmarkr\checkmarkc\checkmark$\checkmark\\checkmark$\checkmark^\checkmark\\checkmarkc\checkmarki\checkmarkr\checkmarkc\checkmark$\checkmarkp\checkmark$\checkmark^\checkmark\\checkmarkc\checkmarki\checkmarkr\checkmarkc\checkmark$\checkmarki\checkmark$\checkmark^\checkmark\\checkmarkc\checkmarki\checkmarkr\checkmarkc\checkmark$\checkmark^\checkmark$\checkmark^\checkmark\\checkmarkc\checkmarki\checkmarkr\checkmarkc\checkmark$\checkmark2\checkmark$\checkmark^\checkmark\\checkmarkc\checkmarki\checkmarkr\checkmarkc\checkmark$\checkmark \checkmark$\checkmark^\checkmark\\checkmarkc\checkmarki\checkmarkr\checkmarkc\checkmark$\checkmark\\checkmark$\checkmark^\checkmark\\checkmarkc\checkmarki\checkmarkr\checkmarkc\checkmark$\checkmarkc\checkmark$\checkmark^\checkmark\\checkmarkc\checkmarki\checkmarkr\checkmarkc\checkmark$\checkmarkd\checkmark$\checkmark^\checkmark\\checkmarkc\checkmarki\checkmarkr\checkmarkc\checkmark$\checkmarko\checkmark$\checkmark^\checkmark\\checkmarkc\checkmarki\checkmarkr\checkmarkc\checkmark$\checkmarkt\checkmark$\checkmark^\checkmark\\checkmarkc\checkmarki\checkmarkr\checkmarkc\checkmark$\checkmark \checkmark$\checkmark^\checkmark\\checkmarkc\checkmarki\checkmarkr\checkmarkc\checkmark$\checkmarkE\checkmark$\checkmark^\checkmark\\checkmarkc\checkmarki\checkmarkr\checkmarkc\checkmark$\checkmark \checkmark$\checkmark^\checkmark\\checkmarkc\checkmarki\checkmarkr\checkmarkc\checkmark$\checkmark\\checkmark$\checkmark^\checkmark\\checkmarkc\checkmarki\checkmarkr\checkmarkc\checkmark$\checkmarkc\checkmark$\checkmark^\checkmark\\checkmarkc\checkmarki\checkmarkr\checkmarkc\checkmark$\checkmarkd\checkmark$\checkmark^\checkmark\\checkmarkc\checkmarki\checkmarkr\checkmarkc\checkmark$\checkmarko\checkmark$\checkmark^\checkmark\\checkmarkc\checkmarki\checkmarkr\checkmarkc\checkmark$\checkmarkt\checkmark$\checkmark^\checkmark\\checkmarkc\checkmarki\checkmarkr\checkmarkc\checkmark$\checkmark \checkmark$\checkmark^\checkmark\\checkmarkc\checkmarki\checkmarkr\checkmarkc\checkmark$\checkmarkI\checkmark$\checkmark^\checkmark\\checkmarkc\checkmarki\checkmarkr\checkmarkc\checkmark$\checkmark_\checkmark$\checkmark^\checkmark\\checkmarkc\checkmarki\checkmarkr\checkmarkc\checkmark$\checkmark{\checkmark$\checkmark^\checkmark\\checkmarkc\checkmarki\checkmarkr\checkmarkc\checkmark$\checkmarky\checkmark$\checkmark^\checkmark\\checkmarkc\checkmarki\checkmarkr\checkmarkc\checkmark$\checkmarky\checkmark$\checkmark^\checkmark\\checkmarkc\checkmarki\checkmarkr\checkmarkc\checkmark$\checkmark}\checkmark$\checkmark^\checkmark\\checkmarkc\checkmarki\checkmarkr\checkmarkc\checkmark$\checkmark}\checkmark$\checkmark^\checkmark\\checkmarkc\checkmarki\checkmarkr\checkmarkc\checkmark$\checkmark{\checkmark$\checkmark^\checkmark\\checkmarkc\checkmarki\checkmarkr\checkmarkc\checkmark$\checkmark(\checkmark$\checkmark^\checkmark\\checkmarkc\checkmarki\checkmarkr\checkmarkc\checkmark$\checkmarkk\checkmark$\checkmark^\checkmark\\checkmarkc\checkmarki\checkmarkr\checkmarkc\checkmark$\checkmark_\checkmark$\checkmark^\checkmark\\checkmarkc\checkmarki\checkmarkr\checkmarkc\checkmark$\checkmarkL\checkmark$\checkmark^\checkmark\\checkmarkc\checkmarki\checkmarkr\checkmarkc\checkmark$\checkmark \checkmark$\checkmark^\checkmark\\checkmarkc\checkmarki\checkmarkr\checkmarkc\checkmark$\checkmark\\checkmark$\checkmark^\checkmark\\checkmarkc\checkmarki\checkmarkr\checkmarkc\checkmark$\checkmarkc\checkmark$\checkmark^\checkmark\\checkmarkc\checkmarki\checkmarkr\checkmarkc\checkmark$\checkmarkd\checkmark$\checkmark^\checkmark\\checkmarkc\checkmarki\checkmarkr\checkmarkc\checkmark$\checkmarko\checkmark$\checkmark^\checkmark\\checkmarkc\checkmarki\checkmarkr\checkmarkc\checkmark$\checkmarkt\checkmark$\checkmark^\checkmark\\checkmarkc\checkmarki\checkmarkr\checkmarkc\checkmark$\checkmark \checkmark$\checkmark^\checkmark\\checkmarkc\checkmarki\checkmarkr\checkmarkc\checkmark$\checkmarkL\checkmark$\checkmark^\checkmark\\checkmarkc\checkmarki\checkmarkr\checkmarkc\checkmark$\checkmark_\checkmark$\checkmark^\checkmark\\checkmarkc\checkmarki\checkmarkr\checkmarkc\checkmark$\checkmarkf\checkmark$\checkmark^\checkmark\\checkmarkc\checkmarki\checkmarkr\checkmarkc\checkmark$\checkmark)\checkmark$\checkmark^\checkmark\\checkmarkc\checkmarki\checkmarkr\checkmarkc\checkmark$\checkmark^\checkmark$\checkmark^\checkmark\\checkmarkc\checkmarki\checkmarkr\checkmarkc\checkmark$\checkmark2\checkmark$\checkmark^\checkmark\\checkmarkc\checkmarki\checkmarkr\checkmarkc\checkmark$\checkmark}\checkmark$\checkmark^\checkmark\\checkmarkc\checkmarki\checkmarkr\checkmarkc\checkmark$\checkmark \checkmark$\checkmark^\checkmark\\checkmarkc\checkmarki\checkmarkr\checkmarkc\checkmark$\checkmark=\checkmark$\checkmark^\checkmark\\checkmarkc\checkmarki\checkmarkr\checkmarkc\checkmark$\checkmark \checkmark$\checkmark^\checkmark\\checkmarkc\checkmarki\checkmarkr\checkmarkc\checkmark$\checkmark\\checkmark$\checkmark^\checkmark\\checkmarkc\checkmarki\checkmarkr\checkmarkc\checkmark$\checkmarkf\checkmark$\checkmark^\checkmark\\checkmarkc\checkmarki\checkmarkr\checkmarkc\checkmark$\checkmarkr\checkmark$\checkmark^\checkmark\\checkmarkc\checkmarki\checkmarkr\checkmarkc\checkmark$\checkmarka\checkmark$\checkmark^\checkmark\\checkmarkc\checkmarki\checkmarkr\checkmarkc\checkmark$\checkmarkc\checkmark$\checkmark^\checkmark\\checkmarkc\checkmarki\checkmarkr\checkmarkc\checkmark$\checkmark{\checkmark$\checkmark^\checkmark\\checkmarkc\checkmarki\checkmarkr\checkmarkc\checkmark$\checkmark\\checkmark$\checkmark^\checkmark\\checkmarkc\checkmarki\checkmarkr\checkmarkc\checkmark$\checkmarkp\checkmark$\checkmark^\checkmark\\checkmarkc\checkmarki\checkmarkr\checkmarkc\checkmark$\checkmarki\checkmark$\checkmark^\checkmark\\checkmarkc\checkmarki\checkmarkr\checkmarkc\checkmark$\checkmark^\checkmark$\checkmark^\checkmark\\checkmarkc\checkmarki\checkmarkr\checkmarkc\checkmark$\checkmark2\checkmark$\checkmark^\checkmark\\checkmarkc\checkmarki\checkmarkr\checkmarkc\checkmark$\checkmark \checkmark$\checkmark^\checkmark\\checkmarkc\checkmarki\checkmarkr\checkmarkc\checkmark$\checkmark\\checkmark$\checkmark^\checkmark\\checkmarkc\checkmarki\checkmarkr\checkmarkc\checkmark$\checkmarkt\checkmark$\checkmark^\checkmark\\checkmarkc\checkmarki\checkmarkr\checkmarkc\checkmark$\checkmarki\checkmark$\checkmark^\checkmark\\checkmarkc\checkmarki\checkmarkr\checkmarkc\checkmark$\checkmarkm\checkmark$\checkmark^\checkmark\\checkmarkc\checkmarki\checkmarkr\checkmarkc\checkmark$\checkmarke\checkmark$\checkmark^\checkmark\\checkmarkc\checkmarki\checkmarkr\checkmarkc\checkmark$\checkmarks\checkmark$\checkmark^\checkmark\\checkmarkc\checkmarki\checkmarkr\checkmarkc\checkmark$\checkmark \checkmark$\checkmark^\checkmark\\checkmarkc\checkmarki\checkmarkr\checkmarkc\checkmark$\checkmark2\checkmark$\checkmark^\checkmark\\checkmarkc\checkmarki\checkmarkr\checkmarkc\checkmark$\checkmark1\checkmark$\checkmark^\checkmark\\checkmarkc\checkmarki\checkmarkr\checkmarkc\checkmark$\checkmark0\checkmark$\checkmark^\checkmark\\checkmarkc\checkmarki\checkmarkr\checkmarkc\checkmark$\checkmark\\checkmark$\checkmark^\checkmark\\checkmarkc\checkmarki\checkmarkr\checkmarkc\checkmark$\checkmark,\checkmark$\checkmark^\checkmark\\checkmarkc\checkmarki\checkmarkr\checkmarkc\checkmark$\checkmark0\checkmark$\checkmark^\checkmark\\checkmarkc\checkmarki\checkmarkr\checkmarkc\checkmark$\checkmark0\checkmark$\checkmark^\checkmark\\checkmarkc\checkmarki\checkmarkr\checkmarkc\checkmark$\checkmark0\checkmark$\checkmark^\checkmark\\checkmarkc\checkmarki\checkmarkr\checkmarkc\checkmark$\checkmark \checkmark$\checkmark^\checkmark\\checkmarkc\checkmarki\checkmarkr\checkmarkc\checkmark$\checkmark\\checkmark$\checkmark^\checkmark\\checkmarkc\checkmarki\checkmarkr\checkmarkc\checkmark$\checkmarkt\checkmark$\checkmark^\checkmark\\checkmarkc\checkmarki\checkmarkr\checkmarkc\checkmark$\checkmarki\checkmark$\checkmark^\checkmark\\checkmarkc\checkmarki\checkmarkr\checkmarkc\checkmark$\checkmarkm\checkmark$\checkmark^\checkmark\\checkmarkc\checkmarki\checkmarkr\checkmarkc\checkmark$\checkmarke\checkmark$\checkmark^\checkmark\\checkmarkc\checkmarki\checkmarkr\checkmarkc\checkmark$\checkmarks\checkmark$\checkmark^\checkmark\\checkmarkc\checkmarki\checkmarkr\checkmarkc\checkmark$\checkmark \checkmark$\checkmark^\checkmark\\checkmarkc\checkmarki\checkmarkr\checkmarkc\checkmark$\checkmark1\checkmark$\checkmark^\checkmark\\checkmarkc\checkmarki\checkmarkr\checkmarkc\checkmark$\checkmark\\checkmark$\checkmark^\checkmark\\checkmarkc\checkmarki\checkmarkr\checkmarkc\checkmark$\checkmark,\checkmark$\checkmark^\checkmark\\checkmarkc\checkmarki\checkmarkr\checkmarkc\checkmark$\checkmark0\checkmark$\checkmark^\checkmark\\checkmarkc\checkmarki\checkmarkr\checkmarkc\checkmark$\checkmark1\checkmark$\checkmark^\checkmark\\checkmarkc\checkmarki\checkmarkr\checkmarkc\checkmark$\checkmark7\checkmark$\checkmark^\checkmark\\checkmarkc\checkmarki\checkmarkr\checkmarkc\checkmark$\checkmark{\checkmark$\checkmark^\checkmark\\checkmarkc\checkmarki\checkmarkr\checkmarkc\checkmark$\checkmark,\checkmark$\checkmark^\checkmark\\checkmarkc\checkmarki\checkmarkr\checkmarkc\checkmark$\checkmark}\checkmark$\checkmark^\checkmark\\checkmarkc\checkmarki\checkmarkr\checkmarkc\checkmark$\checkmark9\checkmark$\checkmark^\checkmark\\checkmarkc\checkmarki\checkmarkr\checkmarkc\checkmark$\checkmark}\checkmark$\checkmark^\checkmark\\checkmarkc\checkmarki\checkmarkr\checkmarkc\checkmark$\checkmark{\checkmark$\checkmark^\checkmark\\checkmarkc\checkmarki\checkmarkr\checkmarkc\checkmark$\checkmark(\checkmark$\checkmark^\checkmark\\checkmarkc\checkmarki\checkmarkr\checkmarkc\checkmark$\checkmark0\checkmark$\checkmark^\checkmark\\checkmarkc\checkmarki\checkmarkr\checkmarkc\checkmark$\checkmark{\checkmark$\checkmark^\checkmark\\checkmarkc\checkmarki\checkmarkr\checkmarkc\checkmark$\checkmark,\checkmark$\checkmark^\checkmark\\checkmarkc\checkmarki\checkmarkr\checkmarkc\checkmark$\checkmark}\checkmark$\checkmark^\checkmark\\checkmarkc\checkmarki\checkmarkr\checkmarkc\checkmark$\checkmark7\checkmark$\checkmark^\checkmark\\checkmarkc\checkmarki\checkmarkr\checkmarkc\checkmark$\checkmark0\checkmark$\checkmark^\checkmark\\checkmarkc\checkmarki\checkmarkr\checkmarkc\checkmark$\checkmark7\checkmark$\checkmark^\checkmark\\checkmarkc\checkmarki\checkmarkr\checkmarkc\checkmark$\checkmark \checkmark$\checkmark^\checkmark\\checkmarkc\checkmarki\checkmarkr\checkmarkc\checkmark$\checkmark\\checkmark$\checkmark^\checkmark\\checkmarkc\checkmarki\checkmarkr\checkmarkc\checkmark$\checkmarkt\checkmark$\checkmark^\checkmark\\checkmarkc\checkmarki\checkmarkr\checkmarkc\checkmark$\checkmarki\checkmark$\checkmark^\checkmark\\checkmarkc\checkmarki\checkmarkr\checkmarkc\checkmark$\checkmarkm\checkmark$\checkmark^\checkmark\\checkmarkc\checkmarki\checkmarkr\checkmarkc\checkmark$\checkmarke\checkmark$\checkmark^\checkmark\\checkmarkc\checkmarki\checkmarkr\checkmarkc\checkmark$\checkmarks\checkmark$\checkmark^\checkmark\\checkmarkc\checkmarki\checkmarkr\checkmarkc\checkmark$\checkmark \checkmark$\checkmark^\checkmark\\checkmarkc\checkmarki\checkmarkr\checkmarkc\checkmark$\checkmark4\checkmark$\checkmark^\checkmark\\checkmarkc\checkmarki\checkmarkr\checkmarkc\checkmark$\checkmark0\checkmark$\checkmark^\checkmark\\checkmarkc\checkmarki\checkmarkr\checkmarkc\checkmark$\checkmark0\checkmark$\checkmark^\checkmark\\checkmarkc\checkmarki\checkmarkr\checkmarkc\checkmark$\checkmark)\checkmark$\checkmark^\checkmark\\checkmarkc\checkmarki\checkmarkr\checkmarkc\checkmark$\checkmark^\checkmark$\checkmark^\checkmark\\checkmarkc\checkmarki\checkmarkr\checkmarkc\checkmark$\checkmark2\checkmark$\checkmark^\checkmark\\checkmarkc\checkmarki\checkmarkr\checkmarkc\checkmark$\checkmark}\checkmark$\checkmark^\checkmark\\checkmarkc\checkmarki\checkmarkr\checkmarkc\checkmark$\checkmark \checkmark$\checkmark^\checkmark\\checkmarkc\checkmarki\checkmarkr\checkmarkc\checkmark$\checkmark=\checkmark$\checkmark^\checkmark\\checkmarkc\checkmarki\checkmarkr\checkmarkc\checkmark$\checkmark \checkmark$\checkmark^\checkmark\\checkmarkc\checkmarki\checkmarkr\checkmarkc\checkmark$\checkmark\\checkmark$\checkmark^\checkmark\\checkmarkc\checkmarki\checkmarkr\checkmarkc\checkmark$\checkmarkb\checkmark$\checkmark^\checkmark\\checkmarkc\checkmarki\checkmarkr\checkmarkc\checkmark$\checkmarko\checkmark$\checkmark^\checkmark\\checkmarkc\checkmarki\checkmarkr\checkmarkc\checkmark$\checkmarkx\checkmark$\checkmark^\checkmark\\checkmarkc\checkmarki\checkmarkr\checkmarkc\checkmark$\checkmarke\checkmark$\checkmark^\checkmark\\checkmarkc\checkmarki\checkmarkr\checkmarkc\checkmark$\checkmarkd\checkmark$\checkmark^\checkmark\\checkmarkc\checkmarki\checkmarkr\checkmarkc\checkmark$\checkmark{\checkmark$\checkmark^\checkmark\\checkmarkc\checkmarki\checkmarkr\checkmarkc\checkmark$\checkmark2\checkmark$\checkmark^\checkmark\\checkmarkc\checkmarki\checkmarkr\checkmarkc\checkmark$\checkmark6\checkmark$\checkmark^\checkmark\\checkmarkc\checkmarki\checkmarkr\checkmarkc\checkmark$\checkmark\\checkmark$\checkmark^\checkmark\\checkmarkc\checkmarki\checkmarkr\checkmarkc\checkmark$\checkmark,\checkmark$\checkmark^\checkmark\\checkmarkc\checkmarki\checkmarkr\checkmarkc\checkmark$\checkmark3\checkmark$\checkmark^\checkmark\\checkmarkc\checkmarki\checkmarkr\checkmarkc\checkmark$\checkmark9\checkmark$\checkmark^\checkmark\\checkmarkc\checkmarki\checkmarkr\checkmarkc\checkmark$\checkmark5\checkmark$\checkmark^\checkmark\\checkmarkc\checkmarki\checkmarkr\checkmarkc\checkmark$\checkmark \checkmark$\checkmark^\checkmark\\checkmarkc\checkmarki\checkmarkr\checkmarkc\checkmark$\checkmark\\checkmark$\checkmark^\checkmark\\checkmarkc\checkmarki\checkmarkr\checkmarkc\checkmark$\checkmarkt\checkmark$\checkmark^\checkmark\\checkmarkc\checkmarki\checkmarkr\checkmarkc\checkmark$\checkmarke\checkmark$\checkmark^\checkmark\\checkmarkc\checkmarki\checkmarkr\checkmarkc\checkmark$\checkmarkx\checkmark$\checkmark^\checkmark\\checkmarkc\checkmarki\checkmarkr\checkmarkc\checkmark$\checkmarkt\checkmark$\checkmark^\checkmark\\checkmarkc\checkmarki\checkmarkr\checkmarkc\checkmark$\checkmark{\checkmark$\checkmark^\checkmark\\checkmarkc\checkmarki\checkmarkr\checkmarkc\checkmark$\checkmark \checkmark$\checkmark^\checkmark\\checkmarkc\checkmarki\checkmarkr\checkmarkc\checkmark$\checkmarkN\checkmark$\checkmark^\checkmark\\checkmarkc\checkmarki\checkmarkr\checkmarkc\checkmark$\checkmark}\checkmark$\checkmark^\checkmark\\checkmarkc\checkmarki\checkmarkr\checkmarkc\checkmark$\checkmark}\checkmark$\checkmark^\checkmark\\checkmarkc\checkmarki\checkmarkr\checkmarkc\checkmark$\checkmark
\checkmark$\checkmark^\checkmark\\checkmarkc\checkmarki\checkmarkr\checkmarkc\checkmark$\checkmark \checkmark$\checkmark^\checkmark\\checkmarkc\checkmarki\checkmarkr\checkmarkc\checkmark$\checkmark \checkmark$\checkmark^\checkmark\\checkmarkc\checkmarki\checkmarkr\checkmarkc\checkmark$\checkmark \checkmark$\checkmark^\checkmark\\checkmarkc\checkmarki\checkmarkr\checkmarkc\checkmark$\checkmark \checkmark$\checkmark^\checkmark\\checkmarkc\checkmarki\checkmarkr\checkmarkc\checkmark$\checkmark\\checkmark$\checkmark^\checkmark\\checkmarkc\checkmarki\checkmarkr\checkmarkc\checkmark$\checkmarkl\checkmark$\checkmark^\checkmark\\checkmarkc\checkmarki\checkmarkr\checkmarkc\checkmark$\checkmarka\checkmark$\checkmark^\checkmark\\checkmarkc\checkmarki\checkmarkr\checkmarkc\checkmark$\checkmarkb\checkmark$\checkmark^\checkmark\\checkmarkc\checkmarki\checkmarkr\checkmarkc\checkmark$\checkmarke\checkmark$\checkmark^\checkmark\\checkmarkc\checkmarki\checkmarkr\checkmarkc\checkmark$\checkmarkl\checkmark$\checkmark^\checkmark\\checkmarkc\checkmarki\checkmarkr\checkmarkc\checkmark$\checkmark{\checkmark$\checkmark^\checkmark\\checkmarkc\checkmarki\checkmarkr\checkmarkc\checkmark$\checkmarke\checkmark$\checkmark^\checkmark\\checkmarkc\checkmarki\checkmarkr\checkmarkc\checkmark$\checkmarkq\checkmark$\checkmark^\checkmark\\checkmarkc\checkmarki\checkmarkr\checkmarkc\checkmark$\checkmark:\checkmark$\checkmark^\checkmark\\checkmarkc\checkmarki\checkmarkr\checkmarkc\checkmark$\checkmarke\checkmark$\checkmark^\checkmark\\checkmarkc\checkmarki\checkmarkr\checkmarkc\checkmark$\checkmarku\checkmark$\checkmark^\checkmark\\checkmarkc\checkmarki\checkmarkr\checkmarkc\checkmark$\checkmarkl\checkmark$\checkmark^\checkmark\\checkmarkc\checkmarki\checkmarkr\checkmarkc\checkmark$\checkmarke\checkmark$\checkmark^\checkmark\\checkmarkc\checkmarki\checkmarkr\checkmarkc\checkmark$\checkmarkr\checkmark$\checkmark^\checkmark\\checkmarkc\checkmarki\checkmarkr\checkmarkc\checkmark$\checkmark}\checkmark$\checkmark^\checkmark\\checkmarkc\checkmarki\checkmarkr\checkmarkc\checkmark$\checkmark
\checkmark$\checkmark^\checkmark\\checkmarkc\checkmarki\checkmarkr\checkmarkc\checkmark$\checkmark\\checkmark$\checkmark^\checkmark\\checkmarkc\checkmarki\checkmarkr\checkmarkc\checkmark$\checkmarke\checkmark$\checkmark^\checkmark\\checkmarkc\checkmarki\checkmarkr\checkmarkc\checkmark$\checkmarkn\checkmark$\checkmark^\checkmark\\checkmarkc\checkmarki\checkmarkr\checkmarkc\checkmark$\checkmarkd\checkmark$\checkmark^\checkmark\\checkmarkc\checkmarki\checkmarkr\checkmarkc\checkmark$\checkmark{\checkmark$\checkmark^\checkmark\\checkmarkc\checkmarki\checkmarkr\checkmarkc\checkmark$\checkmarke\checkmark$\checkmark^\checkmark\\checkmarkc\checkmarki\checkmarkr\checkmarkc\checkmark$\checkmarkq\checkmark$\checkmark^\checkmark\\checkmarkc\checkmarki\checkmarkr\checkmarkc\checkmark$\checkmarku\checkmark$\checkmark^\checkmark\\checkmarkc\checkmarki\checkmarkr\checkmarkc\checkmark$\checkmarka\checkmark$\checkmark^\checkmark\\checkmarkc\checkmarki\checkmarkr\checkmarkc\checkmark$\checkmarkt\checkmark$\checkmark^\checkmark\\checkmarkc\checkmarki\checkmarkr\checkmarkc\checkmark$\checkmarki\checkmark$\checkmark^\checkmark\\checkmarkc\checkmarki\checkmarkr\checkmarkc\checkmark$\checkmarko\checkmark$\checkmark^\checkmark\\checkmarkc\checkmarki\checkmarkr\checkmarkc\checkmark$\checkmarkn\checkmark$\checkmark^\checkmark\\checkmarkc\checkmarki\checkmarkr\checkmarkc\checkmark$\checkmark}\checkmark$\checkmark^\checkmark\\checkmarkc\checkmarki\checkmarkr\checkmarkc\checkmark$\checkmark
\checkmark$\checkmark^\checkmark\\checkmarkc\checkmarki\checkmarkr\checkmarkc\checkmark$\checkmark
\checkmark$\checkmark^\checkmark\\checkmarkc\checkmarki\checkmarkr\checkmarkc\checkmark$\checkmarkL\checkmark$\checkmark^\checkmark\\checkmarkc\checkmarki\checkmarkr\checkmarkc\checkmark$\checkmarke\checkmark$\checkmark^\checkmark\\checkmarkc\checkmarki\checkmarkr\checkmarkc\checkmark$\checkmark \checkmark$\checkmark^\checkmark\\checkmarkc\checkmarki\checkmarkr\checkmarkc\checkmark$\checkmarkc\checkmark$\checkmark^\checkmark\\checkmarkc\checkmarki\checkmarkr\checkmarkc\checkmark$\checkmarko\checkmark$\checkmark^\checkmark\\checkmarkc\checkmarki\checkmarkr\checkmarkc\checkmark$\checkmarke\checkmark$\checkmark^\checkmark\\checkmarkc\checkmarki\checkmarkr\checkmarkc\checkmark$\checkmarkf\checkmark$\checkmark^\checkmark\\checkmarkc\checkmarki\checkmarkr\checkmarkc\checkmark$\checkmarkf\checkmark$\checkmark^\checkmark\\checkmarkc\checkmarki\checkmarkr\checkmarkc\checkmark$\checkmarki\checkmark$\checkmark^\checkmark\\checkmarkc\checkmarki\checkmarkr\checkmarkc\checkmark$\checkmarkc\checkmark$\checkmark^\checkmark\\checkmarkc\checkmarki\checkmarkr\checkmarkc\checkmark$\checkmarki\checkmark$\checkmark^\checkmark\\checkmarkc\checkmarki\checkmarkr\checkmarkc\checkmark$\checkmarke\checkmark$\checkmark^\checkmark\\checkmarkc\checkmarki\checkmarkr\checkmarkc\checkmark$\checkmarkn\checkmark$\checkmark^\checkmark\\checkmarkc\checkmarki\checkmarkr\checkmarkc\checkmark$\checkmarkt\checkmark$\checkmark^\checkmark\\checkmarkc\checkmarki\checkmarkr\checkmarkc\checkmark$\checkmark \checkmark$\checkmark^\checkmark\\checkmarkc\checkmarki\checkmarkr\checkmarkc\checkmark$\checkmarkd\checkmark$\checkmark^\checkmark\\checkmarkc\checkmarki\checkmarkr\checkmarkc\checkmark$\checkmarke\checkmark$\checkmark^\checkmark\\checkmarkc\checkmarki\checkmarkr\checkmarkc\checkmark$\checkmark \checkmark$\checkmark^\checkmark\\checkmarkc\checkmarki\checkmarkr\checkmarkc\checkmark$\checkmarks\checkmark$\checkmark^\checkmark\\checkmarkc\checkmarki\checkmarkr\checkmarkc\checkmark$\checkmarké\checkmark$\checkmark^\checkmark\\checkmarkc\checkmarki\checkmarkr\checkmarkc\checkmark$\checkmarkc\checkmark$\checkmark^\checkmark\\checkmarkc\checkmarki\checkmarkr\checkmarkc\checkmark$\checkmarku\checkmark$\checkmark^\checkmark\\checkmarkc\checkmarki\checkmarkr\checkmarkc\checkmark$\checkmarkr\checkmark$\checkmark^\checkmark\\checkmarkc\checkmarki\checkmarkr\checkmarkc\checkmark$\checkmarki\checkmark$\checkmark^\checkmark\\checkmarkc\checkmarki\checkmarkr\checkmarkc\checkmark$\checkmarkt\checkmark$\checkmark^\checkmark\\checkmarkc\checkmarki\checkmarkr\checkmarkc\checkmark$\checkmarké\checkmark$\checkmark^\checkmark\\checkmarkc\checkmarki\checkmarkr\checkmarkc\checkmark$\checkmark \checkmark$\checkmark^\checkmark\\checkmarkc\checkmarki\checkmarkr\checkmarkc\checkmark$\checkmarka\checkmark$\checkmark^\checkmark\\checkmarkc\checkmarki\checkmarkr\checkmarkc\checkmark$\checkmarku\checkmark$\checkmark^\checkmark\\checkmarkc\checkmarki\checkmarkr\checkmarkc\checkmark$\checkmark \checkmark$\checkmark^\checkmark\\checkmarkc\checkmarki\checkmarkr\checkmarkc\checkmark$\checkmarkf\checkmark$\checkmark^\checkmark\\checkmarkc\checkmarki\checkmarkr\checkmarkc\checkmark$\checkmarkl\checkmark$\checkmark^\checkmark\\checkmarkc\checkmarki\checkmarkr\checkmarkc\checkmark$\checkmarka\checkmark$\checkmark^\checkmark\\checkmarkc\checkmarki\checkmarkr\checkmarkc\checkmark$\checkmarkm\checkmark$\checkmark^\checkmark\\checkmarkc\checkmarki\checkmarkr\checkmarkc\checkmark$\checkmarkb\checkmark$\checkmark^\checkmark\\checkmarkc\checkmarki\checkmarkr\checkmarkc\checkmark$\checkmarka\checkmark$\checkmark^\checkmark\\checkmarkc\checkmarki\checkmarkr\checkmarkc\checkmark$\checkmarkg\checkmark$\checkmark^\checkmark\\checkmarkc\checkmarki\checkmarkr\checkmarkc\checkmark$\checkmarke\checkmark$\checkmark^\checkmark\\checkmarkc\checkmarki\checkmarkr\checkmarkc\checkmark$\checkmark \checkmark$\checkmark^\checkmark\\checkmarkc\checkmarki\checkmarkr\checkmarkc\checkmark$\checkmarke\checkmark$\checkmark^\checkmark\\checkmarkc\checkmarki\checkmarkr\checkmarkc\checkmark$\checkmarks\checkmark$\checkmark^\checkmark\\checkmarkc\checkmarki\checkmarkr\checkmarkc\checkmark$\checkmarkt\checkmark$\checkmark^\checkmark\\checkmarkc\checkmarki\checkmarkr\checkmarkc\checkmark$\checkmark \checkmark$\checkmark^\checkmark\\checkmarkc\checkmarki\checkmarkr\checkmarkc\checkmark$\checkmark:\checkmark$\checkmark^\checkmark\\checkmarkc\checkmarki\checkmarkr\checkmarkc\checkmark$\checkmark
\checkmark$\checkmark^\checkmark\\checkmarkc\checkmarki\checkmarkr\checkmarkc\checkmark$\checkmark\\checkmark$\checkmark^\checkmark\\checkmarkc\checkmarki\checkmarkr\checkmarkc\checkmark$\checkmarkb\checkmark$\checkmark^\checkmark\\checkmarkc\checkmarki\checkmarkr\checkmarkc\checkmark$\checkmarke\checkmark$\checkmark^\checkmark\\checkmarkc\checkmarki\checkmarkr\checkmarkc\checkmark$\checkmarkg\checkmark$\checkmark^\checkmark\\checkmarkc\checkmarki\checkmarkr\checkmarkc\checkmark$\checkmarki\checkmark$\checkmark^\checkmark\\checkmarkc\checkmarki\checkmarkr\checkmarkc\checkmark$\checkmarkn\checkmark$\checkmark^\checkmark\\checkmarkc\checkmarki\checkmarkr\checkmarkc\checkmark$\checkmark{\checkmark$\checkmark^\checkmark\\checkmarkc\checkmarki\checkmarkr\checkmarkc\checkmark$\checkmarke\checkmark$\checkmark^\checkmark\\checkmarkc\checkmarki\checkmarkr\checkmarkc\checkmark$\checkmarkq\checkmark$\checkmark^\checkmark\\checkmarkc\checkmarki\checkmarkr\checkmarkc\checkmark$\checkmarku\checkmark$\checkmark^\checkmark\\checkmarkc\checkmarki\checkmarkr\checkmarkc\checkmark$\checkmarka\checkmark$\checkmark^\checkmark\\checkmarkc\checkmarki\checkmarkr\checkmarkc\checkmark$\checkmarkt\checkmark$\checkmark^\checkmark\\checkmarkc\checkmarki\checkmarkr\checkmarkc\checkmark$\checkmarki\checkmark$\checkmark^\checkmark\\checkmarkc\checkmarki\checkmarkr\checkmarkc\checkmark$\checkmarko\checkmark$\checkmark^\checkmark\\checkmarkc\checkmarki\checkmarkr\checkmarkc\checkmark$\checkmarkn\checkmark$\checkmark^\checkmark\\checkmarkc\checkmarki\checkmarkr\checkmarkc\checkmark$\checkmark}\checkmark$\checkmark^\checkmark\\checkmarkc\checkmarki\checkmarkr\checkmarkc\checkmark$\checkmark
\checkmark$\checkmark^\checkmark\\checkmarkc\checkmarki\checkmarkr\checkmarkc\checkmark$\checkmark \checkmark$\checkmark^\checkmark\\checkmarkc\checkmarki\checkmarkr\checkmarkc\checkmark$\checkmark \checkmark$\checkmark^\checkmark\\checkmarkc\checkmarki\checkmarkr\checkmarkc\checkmark$\checkmark \checkmark$\checkmark^\checkmark\\checkmarkc\checkmarki\checkmarkr\checkmarkc\checkmark$\checkmark \checkmark$\checkmark^\checkmark\\checkmarkc\checkmarki\checkmarkr\checkmarkc\checkmark$\checkmarks\checkmark$\checkmark^\checkmark\\checkmarkc\checkmarki\checkmarkr\checkmarkc\checkmark$\checkmark_\checkmark$\checkmark^\checkmark\\checkmarkc\checkmarki\checkmarkr\checkmarkc\checkmark$\checkmark{\checkmark$\checkmark^\checkmark\\checkmarkc\checkmarki\checkmarkr\checkmarkc\checkmark$\checkmarkf\checkmark$\checkmark^\checkmark\\checkmarkc\checkmarki\checkmarkr\checkmarkc\checkmark$\checkmarkl\checkmark$\checkmark^\checkmark\\checkmarkc\checkmarki\checkmarkr\checkmarkc\checkmark$\checkmarka\checkmark$\checkmark^\checkmark\\checkmarkc\checkmarki\checkmarkr\checkmarkc\checkmark$\checkmarkm\checkmark$\checkmark^\checkmark\\checkmarkc\checkmarki\checkmarkr\checkmarkc\checkmark$\checkmarkb\checkmark$\checkmark^\checkmark\\checkmarkc\checkmarki\checkmarkr\checkmarkc\checkmark$\checkmarka\checkmark$\checkmark^\checkmark\\checkmarkc\checkmarki\checkmarkr\checkmarkc\checkmark$\checkmarkg\checkmark$\checkmark^\checkmark\\checkmarkc\checkmarki\checkmarkr\checkmarkc\checkmark$\checkmarke\checkmark$\checkmark^\checkmark\\checkmarkc\checkmarki\checkmarkr\checkmarkc\checkmark$\checkmark}\checkmark$\checkmark^\checkmark\\checkmarkc\checkmarki\checkmarkr\checkmarkc\checkmark$\checkmark \checkmark$\checkmark^\checkmark\\checkmarkc\checkmarki\checkmarkr\checkmarkc\checkmark$\checkmark=\checkmark$\checkmark^\checkmark\\checkmarkc\checkmarki\checkmarkr\checkmarkc\checkmark$\checkmark \checkmark$\checkmark^\checkmark\\checkmarkc\checkmarki\checkmarkr\checkmarkc\checkmark$\checkmark\\checkmark$\checkmark^\checkmark\\checkmarkc\checkmarki\checkmarkr\checkmarkc\checkmark$\checkmarkf\checkmark$\checkmark^\checkmark\\checkmarkc\checkmarki\checkmarkr\checkmarkc\checkmark$\checkmarkr\checkmark$\checkmark^\checkmark\\checkmarkc\checkmarki\checkmarkr\checkmarkc\checkmark$\checkmarka\checkmark$\checkmark^\checkmark\\checkmarkc\checkmarki\checkmarkr\checkmarkc\checkmark$\checkmarkc\checkmark$\checkmark^\checkmark\\checkmarkc\checkmarki\checkmarkr\checkmarkc\checkmark$\checkmark{\checkmark$\checkmark^\checkmark\\checkmarkc\checkmarki\checkmarkr\checkmarkc\checkmark$\checkmarkF\checkmark$\checkmark^\checkmark\\checkmarkc\checkmarki\checkmarkr\checkmarkc\checkmark$\checkmark_\checkmark$\checkmark^\checkmark\\checkmarkc\checkmarki\checkmarkr\checkmarkc\checkmark$\checkmarkk\checkmark$\checkmark^\checkmark\\checkmarkc\checkmarki\checkmarkr\checkmarkc\checkmark$\checkmark}\checkmark$\checkmark^\checkmark\\checkmarkc\checkmarki\checkmarkr\checkmarkc\checkmark$\checkmark{\checkmark$\checkmark^\checkmark\\checkmarkc\checkmarki\checkmarkr\checkmarkc\checkmark$\checkmarkF\checkmark$\checkmark^\checkmark\\checkmarkc\checkmarki\checkmarkr\checkmarkc\checkmark$\checkmark_\checkmark$\checkmark^\checkmark\\checkmarkc\checkmarki\checkmarkr\checkmarkc\checkmark$\checkmark{\checkmark$\checkmark^\checkmark\\checkmarkc\checkmarki\checkmarkr\checkmarkc\checkmark$\checkmarka\checkmark$\checkmark^\checkmark\\checkmarkc\checkmarki\checkmarkr\checkmarkc\checkmark$\checkmarkx\checkmark$\checkmark^\checkmark\\checkmarkc\checkmarki\checkmarkr\checkmarkc\checkmark$\checkmark}\checkmark$\checkmark^\checkmark\\checkmarkc\checkmarki\checkmarkr\checkmarkc\checkmark$\checkmark}\checkmark$\checkmark^\checkmark\\checkmarkc\checkmarki\checkmarkr\checkmarkc\checkmark$\checkmark \checkmark$\checkmark^\checkmark\\checkmarkc\checkmarki\checkmarkr\checkmarkc\checkmark$\checkmark=\checkmark$\checkmark^\checkmark\\checkmarkc\checkmarki\checkmarkr\checkmarkc\checkmark$\checkmark \checkmark$\checkmark^\checkmark\\checkmarkc\checkmarki\checkmarkr\checkmarkc\checkmark$\checkmark\\checkmark$\checkmark^\checkmark\\checkmarkc\checkmarki\checkmarkr\checkmarkc\checkmark$\checkmarkf\checkmark$\checkmark^\checkmark\\checkmarkc\checkmarki\checkmarkr\checkmarkc\checkmark$\checkmarkr\checkmark$\checkmark^\checkmark\\checkmarkc\checkmarki\checkmarkr\checkmarkc\checkmark$\checkmarka\checkmark$\checkmark^\checkmark\\checkmarkc\checkmarki\checkmarkr\checkmarkc\checkmark$\checkmarkc\checkmark$\checkmark^\checkmark\\checkmarkc\checkmarki\checkmarkr\checkmarkc\checkmark$\checkmark{\checkmark$\checkmark^\checkmark\\checkmarkc\checkmarki\checkmarkr\checkmarkc\checkmark$\checkmark2\checkmark$\checkmark^\checkmark\\checkmarkc\checkmarki\checkmarkr\checkmarkc\checkmark$\checkmark6\checkmark$\checkmark^\checkmark\\checkmarkc\checkmarki\checkmarkr\checkmarkc\checkmark$\checkmark\\checkmark$\checkmark^\checkmark\\checkmarkc\checkmarki\checkmarkr\checkmarkc\checkmark$\checkmark,\checkmark$\checkmark^\checkmark\\checkmarkc\checkmarki\checkmarkr\checkmarkc\checkmark$\checkmark3\checkmark$\checkmark^\checkmark\\checkmarkc\checkmarki\checkmarkr\checkmarkc\checkmark$\checkmark9\checkmark$\checkmark^\checkmark\\checkmarkc\checkmarki\checkmarkr\checkmarkc\checkmark$\checkmark5\checkmark$\checkmark^\checkmark\\checkmarkc\checkmarki\checkmarkr\checkmarkc\checkmark$\checkmark}\checkmark$\checkmark^\checkmark\\checkmarkc\checkmarki\checkmarkr\checkmarkc\checkmark$\checkmark{\checkmark$\checkmark^\checkmark\\checkmarkc\checkmarki\checkmarkr\checkmarkc\checkmark$\checkmark1\checkmark$\checkmark^\checkmark\\checkmarkc\checkmarki\checkmarkr\checkmarkc\checkmark$\checkmark0\checkmark$\checkmark^\checkmark\\checkmarkc\checkmarki\checkmarkr\checkmarkc\checkmark$\checkmark5\checkmark$\checkmark^\checkmark\\checkmarkc\checkmarki\checkmarkr\checkmarkc\checkmark$\checkmark{\checkmark$\checkmark^\checkmark\\checkmarkc\checkmarki\checkmarkr\checkmarkc\checkmark$\checkmark,\checkmark$\checkmark^\checkmark\\checkmarkc\checkmarki\checkmarkr\checkmarkc\checkmark$\checkmark}\checkmark$\checkmark^\checkmark\\checkmarkc\checkmarki\checkmarkr\checkmarkc\checkmark$\checkmark5\checkmark$\checkmark^\checkmark\\checkmarkc\checkmarki\checkmarkr\checkmarkc\checkmark$\checkmark}\checkmark$\checkmark^\checkmark\\checkmarkc\checkmarki\checkmarkr\checkmarkc\checkmark$\checkmark \checkmark$\checkmark^\checkmark\\checkmarkc\checkmarki\checkmarkr\checkmarkc\checkmark$\checkmark\\checkmark$\checkmark^\checkmark\\checkmarkc\checkmarki\checkmarkr\checkmarkc\checkmark$\checkmarka\checkmark$\checkmark^\checkmark\\checkmarkc\checkmarki\checkmarkr\checkmarkc\checkmark$\checkmarkp\checkmark$\checkmark^\checkmark\\checkmarkc\checkmarki\checkmarkr\checkmarkc\checkmark$\checkmarkp\checkmark$\checkmark^\checkmark\\checkmarkc\checkmarki\checkmarkr\checkmarkc\checkmark$\checkmarkr\checkmark$\checkmark^\checkmark\\checkmarkc\checkmarki\checkmarkr\checkmarkc\checkmark$\checkmarko\checkmark$\checkmark^\checkmark\\checkmarkc\checkmarki\checkmarkr\checkmarkc\checkmark$\checkmarkx\checkmark$\checkmark^\checkmark\\checkmarkc\checkmarki\checkmarkr\checkmarkc\checkmark$\checkmark \checkmark$\checkmark^\checkmark\\checkmarkc\checkmarki\checkmarkr\checkmarkc\checkmark$\checkmark\\checkmark$\checkmark^\checkmark\\checkmarkc\checkmarki\checkmarkr\checkmarkc\checkmark$\checkmarkb\checkmark$\checkmark^\checkmark\\checkmarkc\checkmarki\checkmarkr\checkmarkc\checkmark$\checkmarko\checkmark$\checkmark^\checkmark\\checkmarkc\checkmarki\checkmarkr\checkmarkc\checkmark$\checkmarkx\checkmark$\checkmark^\checkmark\\checkmarkc\checkmarki\checkmarkr\checkmarkc\checkmark$\checkmarke\checkmark$\checkmark^\checkmark\\checkmarkc\checkmarki\checkmarkr\checkmarkc\checkmark$\checkmarkd\checkmark$\checkmark^\checkmark\\checkmarkc\checkmarki\checkmarkr\checkmarkc\checkmark$\checkmark{\checkmark$\checkmark^\checkmark\\checkmarkc\checkmarki\checkmarkr\checkmarkc\checkmark$\checkmark2\checkmark$\checkmark^\checkmark\\checkmarkc\checkmarki\checkmarkr\checkmarkc\checkmark$\checkmark5\checkmark$\checkmark^\checkmark\\checkmarkc\checkmarki\checkmarkr\checkmarkc\checkmark$\checkmark0\checkmark$\checkmark^\checkmark\\checkmarkc\checkmarki\checkmarkr\checkmarkc\checkmark$\checkmark \checkmark$\checkmark^\checkmark\\checkmarkc\checkmarki\checkmarkr\checkmarkc\checkmark$\checkmark\\checkmark$\checkmark^\checkmark\\checkmarkc\checkmarki\checkmarkr\checkmarkc\checkmark$\checkmarkg\checkmark$\checkmark^\checkmark\\checkmarkc\checkmarki\checkmarkr\checkmarkc\checkmark$\checkmarkg\checkmark$\checkmark^\checkmark\\checkmarkc\checkmarki\checkmarkr\checkmarkc\checkmark$\checkmark \checkmark$\checkmark^\checkmark\\checkmarkc\checkmarki\checkmarkr\checkmarkc\checkmark$\checkmark1\checkmark$\checkmark^\checkmark\\checkmarkc\checkmarki\checkmarkr\checkmarkc\checkmark$\checkmark{\checkmark$\checkmark^\checkmark\\checkmarkc\checkmarki\checkmarkr\checkmarkc\checkmark$\checkmark,\checkmark$\checkmark^\checkmark\\checkmarkc\checkmarki\checkmarkr\checkmarkc\checkmark$\checkmark}\checkmark$\checkmark^\checkmark\\checkmarkc\checkmarki\checkmarkr\checkmarkc\checkmark$\checkmark5\checkmark$\checkmark^\checkmark\\checkmarkc\checkmarki\checkmarkr\checkmarkc\checkmark$\checkmark}\checkmark$\checkmark^\checkmark\\checkmarkc\checkmarki\checkmarkr\checkmarkc\checkmark$\checkmark
\checkmark$\checkmark^\checkmark\\checkmarkc\checkmarki\checkmarkr\checkmarkc\checkmark$\checkmark\\checkmark$\checkmark^\checkmark\\checkmarkc\checkmarki\checkmarkr\checkmarkc\checkmark$\checkmarke\checkmark$\checkmark^\checkmark\\checkmarkc\checkmarki\checkmarkr\checkmarkc\checkmark$\checkmarkn\checkmark$\checkmark^\checkmark\\checkmarkc\checkmarki\checkmarkr\checkmarkc\checkmark$\checkmarkd\checkmark$\checkmark^\checkmark\\checkmarkc\checkmarki\checkmarkr\checkmarkc\checkmark$\checkmark{\checkmark$\checkmark^\checkmark\\checkmarkc\checkmarki\checkmarkr\checkmarkc\checkmark$\checkmarke\checkmark$\checkmark^\checkmark\\checkmarkc\checkmarki\checkmarkr\checkmarkc\checkmark$\checkmarkq\checkmark$\checkmark^\checkmark\\checkmarkc\checkmarki\checkmarkr\checkmarkc\checkmark$\checkmarku\checkmark$\checkmark^\checkmark\\checkmarkc\checkmarki\checkmarkr\checkmarkc\checkmark$\checkmarka\checkmark$\checkmark^\checkmark\\checkmarkc\checkmarki\checkmarkr\checkmarkc\checkmark$\checkmarkt\checkmark$\checkmark^\checkmark\\checkmarkc\checkmarki\checkmarkr\checkmarkc\checkmark$\checkmarki\checkmark$\checkmark^\checkmark\\checkmarkc\checkmarki\checkmarkr\checkmarkc\checkmark$\checkmarko\checkmark$\checkmark^\checkmark\\checkmarkc\checkmarki\checkmarkr\checkmarkc\checkmark$\checkmarkn\checkmark$\checkmark^\checkmark\\checkmarkc\checkmarki\checkmarkr\checkmarkc\checkmark$\checkmark}\checkmark$\checkmark^\checkmark\\checkmarkc\checkmarki\checkmarkr\checkmarkc\checkmark$\checkmark
\checkmark$\checkmark^\checkmark\\checkmarkc\checkmarki\checkmarkr\checkmarkc\checkmark$\checkmark
\checkmark$\checkmark^\checkmark\\checkmarkc\checkmarki\checkmarkr\checkmarkc\checkmark$\checkmark\\checkmark$\checkmark^\checkmark\\checkmarkc\checkmarki\checkmarkr\checkmarkc\checkmark$\checkmarkt\checkmark$\checkmark^\checkmark\\checkmarkc\checkmarki\checkmarkr\checkmarkc\checkmark$\checkmarke\checkmark$\checkmark^\checkmark\\checkmarkc\checkmarki\checkmarkr\checkmarkc\checkmark$\checkmarkx\checkmark$\checkmark^\checkmark\\checkmarkc\checkmarki\checkmarkr\checkmarkc\checkmark$\checkmarkt\checkmark$\checkmark^\checkmark\\checkmarkc\checkmarki\checkmarkr\checkmarkc\checkmark$\checkmarkb\checkmark$\checkmark^\checkmark\\checkmarkc\checkmarki\checkmarkr\checkmarkc\checkmark$\checkmarkf\checkmark$\checkmark^\checkmark\\checkmarkc\checkmarki\checkmarkr\checkmarkc\checkmark$\checkmark{\checkmark$\checkmark^\checkmark\\checkmarkc\checkmarki\checkmarkr\checkmarkc\checkmark$\checkmarkV\checkmark$\checkmark^\checkmark\\checkmarkc\checkmarki\checkmarkr\checkmarkc\checkmark$\checkmarké\checkmark$\checkmark^\checkmark\\checkmarkc\checkmarki\checkmarkr\checkmarkc\checkmark$\checkmarkr\checkmark$\checkmark^\checkmark\\checkmarkc\checkmarki\checkmarkr\checkmarkc\checkmark$\checkmarki\checkmark$\checkmark^\checkmark\\checkmarkc\checkmarki\checkmarkr\checkmarkc\checkmark$\checkmarkf\checkmark$\checkmark^\checkmark\\checkmarkc\checkmarki\checkmarkr\checkmarkc\checkmark$\checkmarki\checkmark$\checkmark^\checkmark\\checkmarkc\checkmarki\checkmarkr\checkmarkc\checkmark$\checkmarkc\checkmark$\checkmark^\checkmark\\checkmarkc\checkmarki\checkmarkr\checkmarkc\checkmark$\checkmarka\checkmark$\checkmark^\checkmark\\checkmarkc\checkmarki\checkmarkr\checkmarkc\checkmark$\checkmarkt\checkmark$\checkmark^\checkmark\\checkmarkc\checkmarki\checkmarkr\checkmarkc\checkmark$\checkmarki\checkmark$\checkmark^\checkmark\\checkmarkc\checkmarki\checkmarkr\checkmarkc\checkmark$\checkmarko\checkmark$\checkmark^\checkmark\\checkmarkc\checkmarki\checkmarkr\checkmarkc\checkmark$\checkmarkn\checkmark$\checkmark^\checkmark\\checkmarkc\checkmarki\checkmarkr\checkmarkc\checkmark$\checkmark \checkmark$\checkmark^\checkmark\\checkmarkc\checkmarki\checkmarkr\checkmarkc\checkmark$\checkmark:\checkmark$\checkmark^\checkmark\\checkmarkc\checkmarki\checkmarkr\checkmarkc\checkmark$\checkmark}\checkmark$\checkmark^\checkmark\\checkmarkc\checkmarki\checkmarkr\checkmarkc\checkmark$\checkmark \checkmark$\checkmark^\checkmark\\checkmarkc\checkmarki\checkmarkr\checkmarkc\checkmark$\checkmark$\checkmark$\checkmark^\checkmark\\checkmarkc\checkmarki\checkmarkr\checkmarkc\checkmark$\checkmarks\checkmark$\checkmark^\checkmark\\checkmarkc\checkmarki\checkmarkr\checkmarkc\checkmark$\checkmark_\checkmark$\checkmark^\checkmark\\checkmarkc\checkmarki\checkmarkr\checkmarkc\checkmark$\checkmark{\checkmark$\checkmark^\checkmark\\checkmarkc\checkmarki\checkmarkr\checkmarkc\checkmark$\checkmarkf\checkmark$\checkmark^\checkmark\\checkmarkc\checkmarki\checkmarkr\checkmarkc\checkmark$\checkmarkl\checkmark$\checkmark^\checkmark\\checkmarkc\checkmarki\checkmarkr\checkmarkc\checkmark$\checkmarka\checkmark$\checkmark^\checkmark\\checkmarkc\checkmarki\checkmarkr\checkmarkc\checkmark$\checkmarkm\checkmark$\checkmark^\checkmark\\checkmarkc\checkmarki\checkmarkr\checkmarkc\checkmark$\checkmarkb\checkmark$\checkmark^\checkmark\\checkmarkc\checkmarki\checkmarkr\checkmarkc\checkmark$\checkmarka\checkmark$\checkmark^\checkmark\\checkmarkc\checkmarki\checkmarkr\checkmarkc\checkmark$\checkmarkg\checkmark$\checkmark^\checkmark\\checkmarkc\checkmarki\checkmarkr\checkmarkc\checkmark$\checkmarke\checkmark$\checkmark^\checkmark\\checkmarkc\checkmarki\checkmarkr\checkmarkc\checkmark$\checkmark}\checkmark$\checkmark^\checkmark\\checkmarkc\checkmarki\checkmarkr\checkmarkc\checkmark$\checkmark \checkmark$\checkmark^\checkmark\\checkmarkc\checkmarki\checkmarkr\checkmarkc\checkmark$\checkmark=\checkmark$\checkmark^\checkmark\\checkmarkc\checkmarki\checkmarkr\checkmarkc\checkmark$\checkmark \checkmark$\checkmark^\checkmark\\checkmarkc\checkmarki\checkmarkr\checkmarkc\checkmark$\checkmark2\checkmark$\checkmark^\checkmark\\checkmarkc\checkmarki\checkmarkr\checkmarkc\checkmark$\checkmark5\checkmark$\checkmark^\checkmark\\checkmarkc\checkmarki\checkmarkr\checkmarkc\checkmark$\checkmark0\checkmark$\checkmark^\checkmark\\checkmarkc\checkmarki\checkmarkr\checkmarkc\checkmark$\checkmark \checkmark$\checkmark^\checkmark\\checkmarkc\checkmarki\checkmarkr\checkmarkc\checkmark$\checkmark\\checkmark$\checkmark^\checkmark\\checkmarkc\checkmarki\checkmarkr\checkmarkc\checkmark$\checkmarkg\checkmark$\checkmark^\checkmark\\checkmarkc\checkmarki\checkmarkr\checkmarkc\checkmark$\checkmarkg\checkmark$\checkmark^\checkmark\\checkmarkc\checkmarki\checkmarkr\checkmarkc\checkmark$\checkmark \checkmark$\checkmark^\checkmark\\checkmarkc\checkmarki\checkmarkr\checkmarkc\checkmark$\checkmarks\checkmark$\checkmark^\checkmark\\checkmarkc\checkmarki\checkmarkr\checkmarkc\checkmark$\checkmark_\checkmark$\checkmark^\checkmark\\checkmarkc\checkmarki\checkmarkr\checkmarkc\checkmark$\checkmark{\checkmark$\checkmark^\checkmark\\checkmarkc\checkmarki\checkmarkr\checkmarkc\checkmark$\checkmarkm\checkmark$\checkmark^\checkmark\\checkmarkc\checkmarki\checkmarkr\checkmarkc\checkmark$\checkmarki\checkmark$\checkmark^\checkmark\\checkmarkc\checkmarki\checkmarkr\checkmarkc\checkmark$\checkmarkn\checkmark$\checkmark^\checkmark\\checkmarkc\checkmarki\checkmarkr\checkmarkc\checkmark$\checkmark}\checkmark$\checkmark^\checkmark\\checkmarkc\checkmarki\checkmarkr\checkmarkc\checkmark$\checkmark \checkmark$\checkmark^\checkmark\\checkmarkc\checkmarki\checkmarkr\checkmarkc\checkmark$\checkmark=\checkmark$\checkmark^\checkmark\\checkmarkc\checkmarki\checkmarkr\checkmarkc\checkmark$\checkmark \checkmark$\checkmark^\checkmark\\checkmarkc\checkmarki\checkmarkr\checkmarkc\checkmark$\checkmark1\checkmark$\checkmark^\checkmark\\checkmarkc\checkmarki\checkmarkr\checkmarkc\checkmark$\checkmark{\checkmark$\checkmark^\checkmark\\checkmarkc\checkmarki\checkmarkr\checkmarkc\checkmark$\checkmark,\checkmark$\checkmark^\checkmark\\checkmarkc\checkmarki\checkmarkr\checkmarkc\checkmark$\checkmark}\checkmark$\checkmark^\checkmark\\checkmarkc\checkmarki\checkmarkr\checkmarkc\checkmark$\checkmark5\checkmark$\checkmark^\checkmark\\checkmarkc\checkmarki\checkmarkr\checkmarkc\checkmark$\checkmark$\checkmark$\checkmark^\checkmark\\checkmarkc\checkmarki\checkmarkr\checkmarkc\checkmark$\checkmark.\checkmark$\checkmark^\checkmark\\checkmarkc\checkmarki\checkmarkr\checkmarkc\checkmark$\checkmark \checkmark$\checkmark^\checkmark\\checkmarkc\checkmarki\checkmarkr\checkmarkc\checkmark$\checkmarkL\checkmark$\checkmark^\checkmark\\checkmarkc\checkmarki\checkmarkr\checkmarkc\checkmark$\checkmarka\checkmark$\checkmark^\checkmark\\checkmarkc\checkmarki\checkmarkr\checkmarkc\checkmark$\checkmark \checkmark$\checkmark^\checkmark\\checkmarkc\checkmarki\checkmarkr\checkmarkc\checkmark$\checkmarkv\checkmark$\checkmark^\checkmark\\checkmarkc\checkmarki\checkmarkr\checkmarkc\checkmark$\checkmarki\checkmark$\checkmark^\checkmark\\checkmarkc\checkmarki\checkmarkr\checkmarkc\checkmark$\checkmarks\checkmark$\checkmark^\checkmark\\checkmarkc\checkmarki\checkmarkr\checkmarkc\checkmark$\checkmark \checkmark$\checkmark^\checkmark\\checkmarkc\checkmarki\checkmarkr\checkmarkc\checkmark$\checkmarkn\checkmark$\checkmark^\checkmark\\checkmarkc\checkmarki\checkmarkr\checkmarkc\checkmark$\checkmarke\checkmark$\checkmark^\checkmark\\checkmarkc\checkmarki\checkmarkr\checkmarkc\checkmark$\checkmark \checkmark$\checkmark^\checkmark\\checkmarkc\checkmarki\checkmarkr\checkmarkc\checkmark$\checkmarkp\checkmark$\checkmark^\checkmark\\checkmarkc\checkmarki\checkmarkr\checkmarkc\checkmark$\checkmarkr\checkmark$\checkmark^\checkmark\\checkmarkc\checkmarki\checkmarkr\checkmarkc\checkmark$\checkmarké\checkmark$\checkmark^\checkmark\\checkmarkc\checkmarki\checkmarkr\checkmarkc\checkmark$\checkmarks\checkmark$\checkmark^\checkmark\\checkmarkc\checkmarki\checkmarkr\checkmarkc\checkmark$\checkmarke\checkmark$\checkmark^\checkmark\\checkmarkc\checkmarki\checkmarkr\checkmarkc\checkmark$\checkmarkn\checkmark$\checkmark^\checkmark\\checkmarkc\checkmarki\checkmarkr\checkmarkc\checkmark$\checkmarkt\checkmark$\checkmark^\checkmark\\checkmarkc\checkmarki\checkmarkr\checkmarkc\checkmark$\checkmarke\checkmark$\checkmark^\checkmark\\checkmarkc\checkmarki\checkmarkr\checkmarkc\checkmark$\checkmark \checkmark$\checkmark^\checkmark\\checkmarkc\checkmarki\checkmarkr\checkmarkc\checkmark$\checkmarka\checkmark$\checkmark^\checkmark\\checkmarkc\checkmarki\checkmarkr\checkmarkc\checkmark$\checkmarku\checkmark$\checkmark^\checkmark\\checkmarkc\checkmarki\checkmarkr\checkmarkc\checkmark$\checkmarkc\checkmark$\checkmark^\checkmark\\checkmarkc\checkmarki\checkmarkr\checkmarkc\checkmark$\checkmarku\checkmark$\checkmark^\checkmark\\checkmarkc\checkmarki\checkmarkr\checkmarkc\checkmark$\checkmarkn\checkmark$\checkmark^\checkmark\\checkmarkc\checkmarki\checkmarkr\checkmarkc\checkmark$\checkmark \checkmark$\checkmark^\checkmark\\checkmarkc\checkmarki\checkmarkr\checkmarkc\checkmark$\checkmarkr\checkmark$\checkmark^\checkmark\\checkmarkc\checkmarki\checkmarkr\checkmarkc\checkmark$\checkmarki\checkmark$\checkmark^\checkmark\\checkmarkc\checkmarki\checkmarkr\checkmarkc\checkmark$\checkmarks\checkmark$\checkmark^\checkmark\\checkmarkc\checkmarki\checkmarkr\checkmarkc\checkmark$\checkmarkq\checkmark$\checkmark^\checkmark\\checkmarkc\checkmarki\checkmarkr\checkmarkc\checkmark$\checkmarku\checkmark$\checkmark^\checkmark\\checkmarkc\checkmarki\checkmarkr\checkmarkc\checkmark$\checkmarke\checkmark$\checkmark^\checkmark\\checkmarkc\checkmarki\checkmarkr\checkmarkc\checkmark$\checkmark \checkmark$\checkmark^\checkmark\\checkmarkc\checkmarki\checkmarkr\checkmarkc\checkmark$\checkmarkd\checkmark$\checkmark^\checkmark\\checkmarkc\checkmarki\checkmarkr\checkmarkc\checkmark$\checkmarke\checkmark$\checkmark^\checkmark\\checkmarkc\checkmarki\checkmarkr\checkmarkc\checkmark$\checkmark \checkmark$\checkmark^\checkmark\\checkmarkc\checkmarki\checkmarkr\checkmarkc\checkmark$\checkmarkf\checkmark$\checkmark^\checkmark\\checkmarkc\checkmarki\checkmarkr\checkmarkc\checkmark$\checkmarkl\checkmark$\checkmark^\checkmark\\checkmarkc\checkmarki\checkmarkr\checkmarkc\checkmark$\checkmarka\checkmark$\checkmark^\checkmark\\checkmarkc\checkmarki\checkmarkr\checkmarkc\checkmark$\checkmarkm\checkmark$\checkmark^\checkmark\\checkmarkc\checkmarki\checkmarkr\checkmarkc\checkmark$\checkmarkb\checkmark$\checkmark^\checkmark\\checkmarkc\checkmarki\checkmarkr\checkmarkc\checkmark$\checkmarka\checkmark$\checkmark^\checkmark\\checkmarkc\checkmarki\checkmarkr\checkmarkc\checkmark$\checkmarkg\checkmark$\checkmark^\checkmark\\checkmarkc\checkmarki\checkmarkr\checkmarkc\checkmark$\checkmarke\checkmark$\checkmark^\checkmark\\checkmarkc\checkmarki\checkmarkr\checkmarkc\checkmark$\checkmark.\checkmark$\checkmark^\checkmark\\checkmarkc\checkmarki\checkmarkr\checkmarkc\checkmark$\checkmark \checkmark$\checkmark^\checkmark\\checkmarkc\checkmarki\checkmarkr\checkmarkc\checkmark$\checkmarkL\checkmark$\checkmark^\checkmark\\checkmarkc\checkmarki\checkmarkr\checkmarkc\checkmark$\checkmarka\checkmark$\checkmark^\checkmark\\checkmarkc\checkmarki\checkmarkr\checkmarkc\checkmark$\checkmark \checkmark$\checkmark^\checkmark\\checkmarkc\checkmarki\checkmarkr\checkmarkc\checkmark$\checkmarkc\checkmark$\checkmark^\checkmark\\checkmarkc\checkmarki\checkmarkr\checkmarkc\checkmark$\checkmarko\checkmark$\checkmark^\checkmark\\checkmarkc\checkmarki\checkmarkr\checkmarkc\checkmark$\checkmarkn\checkmark$\checkmark^\checkmark\\checkmarkc\checkmarki\checkmarkr\checkmarkc\checkmark$\checkmarkd\checkmark$\checkmark^\checkmark\\checkmarkc\checkmarki\checkmarkr\checkmarkc\checkmark$\checkmarki\checkmark$\checkmark^\checkmark\\checkmarkc\checkmarki\checkmarkr\checkmarkc\checkmark$\checkmarkt\checkmark$\checkmark^\checkmark\\checkmarkc\checkmarki\checkmarkr\checkmarkc\checkmark$\checkmarki\checkmark$\checkmark^\checkmark\\checkmarkc\checkmarki\checkmarkr\checkmarkc\checkmark$\checkmarko\checkmark$\checkmark^\checkmark\\checkmarkc\checkmarki\checkmarkr\checkmarkc\checkmark$\checkmarkn\checkmark$\checkmark^\checkmark\\checkmarkc\checkmarki\checkmarkr\checkmarkc\checkmark$\checkmark \checkmark$\checkmark^\checkmark\\checkmarkc\checkmarki\checkmarkr\checkmarkc\checkmark$\checkmarke\checkmark$\checkmark^\checkmark\\checkmarkc\checkmarki\checkmarkr\checkmarkc\checkmark$\checkmarks\checkmark$\checkmark^\checkmark\\checkmarkc\checkmarki\checkmarkr\checkmarkc\checkmark$\checkmarkt\checkmark$\checkmark^\checkmark\\checkmarkc\checkmarki\checkmarkr\checkmarkc\checkmark$\checkmark \checkmark$\checkmark^\checkmark\\checkmarkc\checkmarki\checkmarkr\checkmarkc\checkmark$\checkmarkl\checkmark$\checkmark^\checkmark\\checkmarkc\checkmarki\checkmarkr\checkmarkc\checkmark$\checkmarka\checkmark$\checkmark^\checkmark\\checkmarkc\checkmarki\checkmarkr\checkmarkc\checkmark$\checkmarkr\checkmark$\checkmark^\checkmark\\checkmarkc\checkmarki\checkmarkr\checkmarkc\checkmark$\checkmarkg\checkmark$\checkmark^\checkmark\\checkmarkc\checkmarki\checkmarkr\checkmarkc\checkmark$\checkmarke\checkmark$\checkmark^\checkmark\\checkmarkc\checkmarki\checkmarkr\checkmarkc\checkmark$\checkmarkm\checkmark$\checkmark^\checkmark\\checkmarkc\checkmarki\checkmarkr\checkmarkc\checkmark$\checkmarke\checkmark$\checkmark^\checkmark\\checkmarkc\checkmarki\checkmarkr\checkmarkc\checkmark$\checkmarkn\checkmark$\checkmark^\checkmark\\checkmarkc\checkmarki\checkmarkr\checkmarkc\checkmark$\checkmarkt\checkmark$\checkmark^\checkmark\\checkmarkc\checkmarki\checkmarkr\checkmarkc\checkmark$\checkmark \checkmark$\checkmark^\checkmark\\checkmarkc\checkmarki\checkmarkr\checkmarkc\checkmark$\checkmarks\checkmark$\checkmark^\checkmark\\checkmarkc\checkmarki\checkmarkr\checkmarkc\checkmark$\checkmarka\checkmark$\checkmark^\checkmark\\checkmarkc\checkmarki\checkmarkr\checkmarkc\checkmark$\checkmarkt\checkmark$\checkmark^\checkmark\\checkmarkc\checkmarki\checkmarkr\checkmarkc\checkmark$\checkmarki\checkmark$\checkmark^\checkmark\\checkmarkc\checkmarki\checkmarkr\checkmarkc\checkmark$\checkmarks\checkmark$\checkmark^\checkmark\\checkmarkc\checkmarki\checkmarkr\checkmarkc\checkmark$\checkmarkf\checkmark$\checkmark^\checkmark\\checkmarkc\checkmarki\checkmarkr\checkmarkc\checkmark$\checkmarka\checkmark$\checkmark^\checkmark\\checkmarkc\checkmarki\checkmarkr\checkmarkc\checkmark$\checkmarki\checkmark$\checkmark^\checkmark\\checkmarkc\checkmarki\checkmarkr\checkmarkc\checkmark$\checkmarkt\checkmark$\checkmark^\checkmark\\checkmarkc\checkmarki\checkmarkr\checkmarkc\checkmark$\checkmarke\checkmark$\checkmark^\checkmark\\checkmarkc\checkmarki\checkmarkr\checkmarkc\checkmark$\checkmark.\checkmark$\checkmark^\checkmark\\checkmarkc\checkmarki\checkmarkr\checkmarkc\checkmark$\checkmark
\checkmark$\checkmark^\checkmark\\checkmarkc\checkmarki\checkmarkr\checkmarkc\checkmark$\checkmark
\checkmark$\checkmark^\checkmark\\checkmarkc\checkmarki\checkmarkr\checkmarkc\checkmark$\checkmark\\checkmark$\checkmark^\checkmark\\checkmarkc\checkmarki\checkmarkr\checkmarkc\checkmark$\checkmarks\checkmark$\checkmark^\checkmark\\checkmarkc\checkmarki\checkmarkr\checkmarkc\checkmark$\checkmarku\checkmark$\checkmark^\checkmark\\checkmarkc\checkmarki\checkmarkr\checkmarkc\checkmark$\checkmarkb\checkmark$\checkmark^\checkmark\\checkmarkc\checkmarki\checkmarkr\checkmarkc\checkmark$\checkmarks\checkmark$\checkmark^\checkmark\\checkmarkc\checkmarki\checkmarkr\checkmarkc\checkmark$\checkmarke\checkmark$\checkmark^\checkmark\\checkmarkc\checkmarki\checkmarkr\checkmarkc\checkmark$\checkmarkc\checkmark$\checkmark^\checkmark\\checkmarkc\checkmarki\checkmarkr\checkmarkc\checkmark$\checkmarkt\checkmark$\checkmark^\checkmark\\checkmarkc\checkmarki\checkmarkr\checkmarkc\checkmark$\checkmarki\checkmark$\checkmark^\checkmark\\checkmarkc\checkmarki\checkmarkr\checkmarkc\checkmark$\checkmarko\checkmark$\checkmark^\checkmark\\checkmarkc\checkmarki\checkmarkr\checkmarkc\checkmark$\checkmarkn\checkmark$\checkmark^\checkmark\\checkmarkc\checkmarki\checkmarkr\checkmarkc\checkmark$\checkmark{\checkmark$\checkmark^\checkmark\\checkmarkc\checkmarki\checkmarkr\checkmarkc\checkmark$\checkmarkD\checkmark$\checkmark^\checkmark\\checkmarkc\checkmarki\checkmarkr\checkmarkc\checkmark$\checkmarké\checkmark$\checkmark^\checkmark\\checkmarkc\checkmarki\checkmarkr\checkmarkc\checkmark$\checkmarkt\checkmark$\checkmark^\checkmark\\checkmarkc\checkmarki\checkmarkr\checkmarkc\checkmark$\checkmarke\checkmark$\checkmark^\checkmark\\checkmarkc\checkmarki\checkmarkr\checkmarkc\checkmark$\checkmarkr\checkmark$\checkmark^\checkmark\\checkmarkc\checkmarki\checkmarkr\checkmarkc\checkmark$\checkmarkm\checkmark$\checkmark^\checkmark\\checkmarkc\checkmarki\checkmarkr\checkmarkc\checkmark$\checkmarki\checkmark$\checkmark^\checkmark\\checkmarkc\checkmarki\checkmarkr\checkmarkc\checkmark$\checkmarkn\checkmark$\checkmark^\checkmark\\checkmarkc\checkmarki\checkmarkr\checkmarkc\checkmark$\checkmarka\checkmark$\checkmark^\checkmark\\checkmarkc\checkmarki\checkmarkr\checkmarkc\checkmark$\checkmarkt\checkmark$\checkmark^\checkmark\\checkmarkc\checkmarki\checkmarkr\checkmarkc\checkmark$\checkmarki\checkmark$\checkmark^\checkmark\\checkmarkc\checkmarki\checkmarkr\checkmarkc\checkmark$\checkmarko\checkmark$\checkmark^\checkmark\\checkmarkc\checkmarki\checkmarkr\checkmarkc\checkmark$\checkmarkn\checkmark$\checkmark^\checkmark\\checkmarkc\checkmarki\checkmarkr\checkmarkc\checkmark$\checkmark \checkmark$\checkmark^\checkmark\\checkmarkc\checkmarki\checkmarkr\checkmarkc\checkmark$\checkmarkd\checkmark$\checkmark^\checkmark\\checkmarkc\checkmarki\checkmarkr\checkmarkc\checkmark$\checkmarku\checkmark$\checkmark^\checkmark\\checkmarkc\checkmarki\checkmarkr\checkmarkc\checkmark$\checkmark \checkmark$\checkmark^\checkmark\\checkmarkc\checkmarki\checkmarkr\checkmarkc\checkmark$\checkmarkc\checkmark$\checkmark^\checkmark\\checkmarkc\checkmarki\checkmarkr\checkmarkc\checkmark$\checkmarko\checkmark$\checkmark^\checkmark\\checkmarkc\checkmarki\checkmarkr\checkmarkc\checkmark$\checkmarku\checkmark$\checkmark^\checkmark\\checkmarkc\checkmarki\checkmarkr\checkmarkc\checkmark$\checkmarkp\checkmark$\checkmark^\checkmark\\checkmarkc\checkmarki\checkmarkr\checkmarkc\checkmark$\checkmarkl\checkmark$\checkmark^\checkmark\\checkmarkc\checkmarki\checkmarkr\checkmarkc\checkmark$\checkmarke\checkmark$\checkmark^\checkmark\\checkmarkc\checkmarki\checkmarkr\checkmarkc\checkmark$\checkmark \checkmark$\checkmark^\checkmark\\checkmarkc\checkmarki\checkmarkr\checkmarkc\checkmark$\checkmarkm\checkmark$\checkmark^\checkmark\\checkmarkc\checkmarki\checkmarkr\checkmarkc\checkmark$\checkmarko\checkmark$\checkmark^\checkmark\\checkmarkc\checkmarki\checkmarkr\checkmarkc\checkmark$\checkmarkt\checkmark$\checkmark^\checkmark\\checkmarkc\checkmarki\checkmarkr\checkmarkc\checkmark$\checkmarke\checkmark$\checkmark^\checkmark\\checkmarkc\checkmarki\checkmarkr\checkmarkc\checkmark$\checkmarku\checkmark$\checkmark^\checkmark\\checkmarkc\checkmarki\checkmarkr\checkmarkc\checkmark$\checkmarkr\checkmark$\checkmark^\checkmark\\checkmarkc\checkmarki\checkmarkr\checkmarkc\checkmark$\checkmark \checkmark$\checkmark^\checkmark\\checkmarkc\checkmarki\checkmarkr\checkmarkc\checkmark$\checkmarkr\checkmark$\checkmark^\checkmark\\checkmarkc\checkmarki\checkmarkr\checkmarkc\checkmark$\checkmarké\checkmark$\checkmark^\checkmark\\checkmarkc\checkmarki\checkmarkr\checkmarkc\checkmark$\checkmarks\checkmark$\checkmark^\checkmark\\checkmarkc\checkmarki\checkmarkr\checkmarkc\checkmark$\checkmarki\checkmark$\checkmark^\checkmark\\checkmarkc\checkmarki\checkmarkr\checkmarkc\checkmark$\checkmarks\checkmark$\checkmark^\checkmark\\checkmarkc\checkmarki\checkmarkr\checkmarkc\checkmark$\checkmarkt\checkmark$\checkmark^\checkmark\\checkmarkc\checkmarki\checkmarkr\checkmarkc\checkmark$\checkmarka\checkmark$\checkmark^\checkmark\\checkmarkc\checkmarki\checkmarkr\checkmarkc\checkmark$\checkmarkn\checkmark$\checkmark^\checkmark\\checkmarkc\checkmarki\checkmarkr\checkmarkc\checkmark$\checkmarkt\checkmark$\checkmark^\checkmark\\checkmarkc\checkmarki\checkmarkr\checkmarkc\checkmark$\checkmark \checkmark$\checkmark^\checkmark\\checkmarkc\checkmarki\checkmarkr\checkmarkc\checkmark$\checkmarkc\checkmark$\checkmark^\checkmark\\checkmarkc\checkmarki\checkmarkr\checkmarkc\checkmark$\checkmarko\checkmark$\checkmark^\checkmark\\checkmarkc\checkmarki\checkmarkr\checkmarkc\checkmark$\checkmarkn\checkmark$\checkmark^\checkmark\\checkmarkc\checkmarki\checkmarkr\checkmarkc\checkmark$\checkmarkt\checkmark$\checkmark^\checkmark\\checkmarkc\checkmarki\checkmarkr\checkmarkc\checkmark$\checkmarki\checkmark$\checkmark^\checkmark\\checkmarkc\checkmarki\checkmarkr\checkmarkc\checkmark$\checkmarkn\checkmark$\checkmark^\checkmark\\checkmarkc\checkmarki\checkmarkr\checkmarkc\checkmark$\checkmarku\checkmark$\checkmark^\checkmark\\checkmarkc\checkmarki\checkmarkr\checkmarkc\checkmark$\checkmark}\checkmark$\checkmark^\checkmark\\checkmarkc\checkmarki\checkmarkr\checkmarkc\checkmark$\checkmark
\checkmark$\checkmark^\checkmark\\checkmarkc\checkmarki\checkmarkr\checkmarkc\checkmark$\checkmarkL\checkmark$\checkmark^\checkmark\\checkmarkc\checkmarki\checkmarkr\checkmarkc\checkmark$\checkmark'\checkmark$\checkmark^\checkmark\\checkmarkc\checkmarki\checkmarkr\checkmarkc\checkmark$\checkmarké\checkmark$\checkmark^\checkmark\\checkmarkc\checkmarki\checkmarkr\checkmarkc\checkmark$\checkmarkq\checkmark$\checkmark^\checkmark\\checkmarkc\checkmarki\checkmarkr\checkmarkc\checkmark$\checkmarku\checkmark$\checkmark^\checkmark\\checkmarkc\checkmarki\checkmarkr\checkmarkc\checkmark$\checkmarka\checkmark$\checkmark^\checkmark\\checkmarkc\checkmarki\checkmarkr\checkmarkc\checkmark$\checkmarkt\checkmark$\checkmark^\checkmark\\checkmarkc\checkmarki\checkmarkr\checkmarkc\checkmark$\checkmarki\checkmark$\checkmark^\checkmark\\checkmarkc\checkmarki\checkmarkr\checkmarkc\checkmark$\checkmarko\checkmark$\checkmark^\checkmark\\checkmarkc\checkmarki\checkmarkr\checkmarkc\checkmark$\checkmarkn\checkmark$\checkmark^\checkmark\\checkmarkc\checkmarki\checkmarkr\checkmarkc\checkmark$\checkmark \checkmark$\checkmark^\checkmark\\checkmarkc\checkmarki\checkmarkr\checkmarkc\checkmark$\checkmarkd\checkmark$\checkmark^\checkmark\\checkmarkc\checkmarki\checkmarkr\checkmarkc\checkmark$\checkmark'\checkmark$\checkmark^\checkmark\\checkmarkc\checkmarki\checkmarkr\checkmarkc\checkmark$\checkmarka\checkmark$\checkmark^\checkmark\\checkmarkc\checkmarki\checkmarkr\checkmarkc\checkmark$\checkmarkc\checkmark$\checkmark^\checkmark\\checkmarkc\checkmarki\checkmarkr\checkmarkc\checkmark$\checkmarkc\checkmark$\checkmark^\checkmark\\checkmarkc\checkmarki\checkmarkr\checkmarkc\checkmark$\checkmarko\checkmark$\checkmark^\checkmark\\checkmarkc\checkmarki\checkmarkr\checkmarkc\checkmark$\checkmarku\checkmark$\checkmark^\checkmark\\checkmarkc\checkmarki\checkmarkr\checkmarkc\checkmark$\checkmarkp\checkmark$\checkmark^\checkmark\\checkmarkc\checkmarki\checkmarkr\checkmarkc\checkmark$\checkmarkl\checkmark$\checkmark^\checkmark\\checkmarkc\checkmarki\checkmarkr\checkmarkc\checkmark$\checkmarke\checkmark$\checkmark^\checkmark\\checkmarkc\checkmarki\checkmarkr\checkmarkc\checkmark$\checkmarkm\checkmark$\checkmark^\checkmark\\checkmarkc\checkmarki\checkmarkr\checkmarkc\checkmark$\checkmarke\checkmark$\checkmark^\checkmark\\checkmarkc\checkmarki\checkmarkr\checkmarkc\checkmark$\checkmarkn\checkmark$\checkmark^\checkmark\\checkmarkc\checkmarki\checkmarkr\checkmarkc\checkmark$\checkmarkt\checkmark$\checkmark^\checkmark\\checkmarkc\checkmarki\checkmarkr\checkmarkc\checkmark$\checkmark \checkmark$\checkmark^\checkmark\\checkmarkc\checkmarki\checkmarkr\checkmarkc\checkmark$\checkmarkm\checkmark$\checkmark^\checkmark\\checkmarkc\checkmarki\checkmarkr\checkmarkc\checkmark$\checkmarko\checkmark$\checkmark^\checkmark\\checkmarkc\checkmarki\checkmarkr\checkmarkc\checkmark$\checkmarkt\checkmark$\checkmark^\checkmark\\checkmarkc\checkmarki\checkmarkr\checkmarkc\checkmark$\checkmarke\checkmark$\checkmark^\checkmark\\checkmarkc\checkmarki\checkmarkr\checkmarkc\checkmark$\checkmarku\checkmark$\checkmark^\checkmark\\checkmarkc\checkmarki\checkmarkr\checkmarkc\checkmark$\checkmarkr\checkmark$\checkmark^\checkmark\\checkmarkc\checkmarki\checkmarkr\checkmarkc\checkmark$\checkmark-\checkmark$\checkmark^\checkmark\\checkmarkc\checkmarki\checkmarkr\checkmarkc\checkmark$\checkmarkv\checkmark$\checkmark^\checkmark\\checkmarkc\checkmarki\checkmarkr\checkmarkc\checkmark$\checkmarki\checkmark$\checkmark^\checkmark\\checkmarkc\checkmarki\checkmarkr\checkmarkc\checkmark$\checkmarks\checkmark$\checkmark^\checkmark\\checkmarkc\checkmarki\checkmarkr\checkmarkc\checkmark$\checkmark \checkmark$\checkmark^\checkmark\\checkmarkc\checkmarki\checkmarkr\checkmarkc\checkmark$\checkmarkà\checkmark$\checkmark^\checkmark\\checkmarkc\checkmarki\checkmarkr\checkmarkc\checkmark$\checkmark \checkmark$\checkmark^\checkmark\\checkmarkc\checkmarki\checkmarkr\checkmarkc\checkmark$\checkmarkb\checkmark$\checkmark^\checkmark\\checkmarkc\checkmarki\checkmarkr\checkmarkc\checkmark$\checkmarki\checkmark$\checkmark^\checkmark\\checkmarkc\checkmarki\checkmarkr\checkmarkc\checkmark$\checkmarkl\checkmark$\checkmark^\checkmark\\checkmarkc\checkmarki\checkmarkr\checkmarkc\checkmark$\checkmarkl\checkmark$\checkmark^\checkmark\\checkmarkc\checkmarki\checkmarkr\checkmarkc\checkmark$\checkmarke\checkmark$\checkmark^\checkmark\\checkmarkc\checkmarki\checkmarkr\checkmarkc\checkmark$\checkmarks\checkmark$\checkmark^\checkmark\\checkmarkc\checkmarki\checkmarkr\checkmarkc\checkmark$\checkmark \checkmark$\checkmark^\checkmark\\checkmarkc\checkmarki\checkmarkr\checkmarkc\checkmark$\checkmarkr\checkmark$\checkmark^\checkmark\\checkmarkc\checkmarki\checkmarkr\checkmarkc\checkmark$\checkmarke\checkmark$\checkmark^\checkmark\\checkmarkc\checkmarki\checkmarkr\checkmarkc\checkmark$\checkmarkl\checkmark$\checkmark^\checkmark\\checkmarkc\checkmarki\checkmarkr\checkmarkc\checkmark$\checkmarki\checkmark$\checkmark^\checkmark\\checkmarkc\checkmarki\checkmarkr\checkmarkc\checkmark$\checkmarke\checkmark$\checkmark^\checkmark\\checkmarkc\checkmarki\checkmarkr\checkmarkc\checkmark$\checkmark \checkmark$\checkmark^\checkmark\\checkmarkc\checkmarki\checkmarkr\checkmarkc\checkmark$\checkmarkl\checkmark$\checkmark^\checkmark\\checkmarkc\checkmarki\checkmarkr\checkmarkc\checkmark$\checkmarka\checkmark$\checkmark^\checkmark\\checkmarkc\checkmarki\checkmarkr\checkmarkc\checkmark$\checkmark \checkmark$\checkmark^\checkmark\\checkmarkc\checkmarki\checkmarkr\checkmarkc\checkmark$\checkmarkf\checkmark$\checkmark^\checkmark\\checkmarkc\checkmarki\checkmarkr\checkmarkc\checkmark$\checkmarko\checkmark$\checkmark^\checkmark\\checkmarkc\checkmarki\checkmarkr\checkmarkc\checkmark$\checkmarkr\checkmark$\checkmark^\checkmark\\checkmarkc\checkmarki\checkmarkr\checkmarkc\checkmark$\checkmarkc\checkmark$\checkmark^\checkmark\\checkmarkc\checkmarki\checkmarkr\checkmarkc\checkmark$\checkmarke\checkmark$\checkmark^\checkmark\\checkmarkc\checkmarki\checkmarkr\checkmarkc\checkmark$\checkmark \checkmark$\checkmark^\checkmark\\checkmarkc\checkmarki\checkmarkr\checkmarkc\checkmark$\checkmarka\checkmark$\checkmark^\checkmark\\checkmarkc\checkmarki\checkmarkr\checkmarkc\checkmark$\checkmarkx\checkmark$\checkmark^\checkmark\\checkmarkc\checkmarki\checkmarkr\checkmarkc\checkmark$\checkmarki\checkmark$\checkmark^\checkmark\\checkmarkc\checkmarki\checkmarkr\checkmarkc\checkmark$\checkmarka\checkmark$\checkmark^\checkmark\\checkmarkc\checkmarki\checkmarkr\checkmarkc\checkmark$\checkmarkl\checkmark$\checkmark^\checkmark\\checkmarkc\checkmarki\checkmarkr\checkmarkc\checkmark$\checkmarke\checkmark$\checkmark^\checkmark\\checkmarkc\checkmarki\checkmarkr\checkmarkc\checkmark$\checkmark \checkmark$\checkmark^\checkmark\\checkmarkc\checkmarki\checkmarkr\checkmarkc\checkmark$\checkmark$\checkmark$\checkmark^\checkmark\\checkmarkc\checkmarki\checkmarkr\checkmarkc\checkmark$\checkmarkF\checkmark$\checkmark^\checkmark\\checkmarkc\checkmarki\checkmarkr\checkmarkc\checkmark$\checkmark_\checkmark$\checkmark^\checkmark\\checkmarkc\checkmarki\checkmarkr\checkmarkc\checkmark$\checkmark{\checkmark$\checkmark^\checkmark\\checkmarkc\checkmarki\checkmarkr\checkmarkc\checkmark$\checkmarka\checkmark$\checkmark^\checkmark\\checkmarkc\checkmarki\checkmarkr\checkmarkc\checkmark$\checkmarkx\checkmark$\checkmark^\checkmark\\checkmarkc\checkmarki\checkmarkr\checkmarkc\checkmark$\checkmark}\checkmark$\checkmark^\checkmark\\checkmarkc\checkmarki\checkmarkr\checkmarkc\checkmark$\checkmark$\checkmark$\checkmark^\checkmark\\checkmarkc\checkmarki\checkmarkr\checkmarkc\checkmark$\checkmark \checkmark$\checkmark^\checkmark\\checkmarkc\checkmarki\checkmarkr\checkmarkc\checkmark$\checkmarka\checkmark$\checkmark^\checkmark\\checkmarkc\checkmarki\checkmarkr\checkmarkc\checkmark$\checkmarku\checkmark$\checkmark^\checkmark\\checkmarkc\checkmarki\checkmarkr\checkmarkc\checkmark$\checkmark \checkmark$\checkmark^\checkmark\\checkmarkc\checkmarki\checkmarkr\checkmarkc\checkmark$\checkmarkc\checkmark$\checkmark^\checkmark\\checkmarkc\checkmarki\checkmarkr\checkmarkc\checkmark$\checkmarko\checkmark$\checkmark^\checkmark\\checkmarkc\checkmarki\checkmarkr\checkmarkc\checkmark$\checkmarku\checkmark$\checkmark^\checkmark\\checkmarkc\checkmarki\checkmarkr\checkmarkc\checkmark$\checkmarkp\checkmark$\checkmark^\checkmark\\checkmarkc\checkmarki\checkmarkr\checkmarkc\checkmark$\checkmarkl\checkmark$\checkmark^\checkmark\\checkmarkc\checkmarki\checkmarkr\checkmarkc\checkmark$\checkmarke\checkmark$\checkmark^\checkmark\\checkmarkc\checkmarki\checkmarkr\checkmarkc\checkmark$\checkmark \checkmark$\checkmark^\checkmark\\checkmarkc\checkmarki\checkmarkr\checkmarkc\checkmark$\checkmarkm\checkmark$\checkmark^\checkmark\\checkmarkc\checkmarki\checkmarkr\checkmarkc\checkmark$\checkmarko\checkmark$\checkmark^\checkmark\\checkmarkc\checkmarki\checkmarkr\checkmarkc\checkmark$\checkmarkt\checkmark$\checkmark^\checkmark\\checkmarkc\checkmarki\checkmarkr\checkmarkc\checkmark$\checkmarke\checkmark$\checkmark^\checkmark\\checkmarkc\checkmarki\checkmarkr\checkmarkc\checkmark$\checkmarku\checkmark$\checkmark^\checkmark\\checkmarkc\checkmarki\checkmarkr\checkmarkc\checkmark$\checkmarkr\checkmark$\checkmark^\checkmark\\checkmarkc\checkmarki\checkmarkr\checkmarkc\checkmark$\checkmark \checkmark$\checkmark^\checkmark\\checkmarkc\checkmarki\checkmarkr\checkmarkc\checkmark$\checkmarkr\checkmark$\checkmark^\checkmark\\checkmarkc\checkmarki\checkmarkr\checkmarkc\checkmark$\checkmarké\checkmark$\checkmark^\checkmark\\checkmarkc\checkmarki\checkmarkr\checkmarkc\checkmark$\checkmarks\checkmark$\checkmark^\checkmark\\checkmarkc\checkmarki\checkmarkr\checkmarkc\checkmark$\checkmarki\checkmark$\checkmark^\checkmark\\checkmarkc\checkmarki\checkmarkr\checkmarkc\checkmark$\checkmarks\checkmark$\checkmark^\checkmark\\checkmarkc\checkmarki\checkmarkr\checkmarkc\checkmark$\checkmarkt\checkmark$\checkmark^\checkmark\\checkmarkc\checkmarki\checkmarkr\checkmarkc\checkmark$\checkmarka\checkmark$\checkmark^\checkmark\\checkmarkc\checkmarki\checkmarkr\checkmarkc\checkmark$\checkmarkn\checkmark$\checkmark^\checkmark\\checkmarkc\checkmarki\checkmarkr\checkmarkc\checkmark$\checkmarkt\checkmark$\checkmark^\checkmark\\checkmarkc\checkmarki\checkmarkr\checkmarkc\checkmark$\checkmark \checkmark$\checkmark^\checkmark\\checkmarkc\checkmarki\checkmarkr\checkmarkc\checkmark$\checkmark$\checkmark$\checkmark^\checkmark\\checkmarkc\checkmarki\checkmarkr\checkmarkc\checkmark$\checkmarkC\checkmark$\checkmark^\checkmark\\checkmarkc\checkmarki\checkmarkr\checkmarkc\checkmark$\checkmark_\checkmark$\checkmark^\checkmark\\checkmarkc\checkmarki\checkmarkr\checkmarkc\checkmark$\checkmarkr\checkmark$\checkmark^\checkmark\\checkmarkc\checkmarki\checkmarkr\checkmarkc\checkmark$\checkmark$\checkmark$\checkmark^\checkmark\\checkmarkc\checkmarki\checkmarkr\checkmarkc\checkmark$\checkmark \checkmark$\checkmark^\checkmark\\checkmarkc\checkmarki\checkmarkr\checkmarkc\checkmark$\checkmark:\checkmark$\checkmark^\checkmark\\checkmarkc\checkmarki\checkmarkr\checkmarkc\checkmark$\checkmark
\checkmark$\checkmark^\checkmark\\checkmarkc\checkmarki\checkmarkr\checkmarkc\checkmark$\checkmark\\checkmark$\checkmark^\checkmark\\checkmarkc\checkmarki\checkmarkr\checkmarkc\checkmark$\checkmarkb\checkmark$\checkmark^\checkmark\\checkmarkc\checkmarki\checkmarkr\checkmarkc\checkmark$\checkmarke\checkmark$\checkmark^\checkmark\\checkmarkc\checkmarki\checkmarkr\checkmarkc\checkmark$\checkmarkg\checkmark$\checkmark^\checkmark\\checkmarkc\checkmarki\checkmarkr\checkmarkc\checkmark$\checkmarki\checkmark$\checkmark^\checkmark\\checkmarkc\checkmarki\checkmarkr\checkmarkc\checkmark$\checkmarkn\checkmark$\checkmark^\checkmark\\checkmarkc\checkmarki\checkmarkr\checkmarkc\checkmark$\checkmark{\checkmark$\checkmark^\checkmark\\checkmarkc\checkmarki\checkmarkr\checkmarkc\checkmark$\checkmarke\checkmark$\checkmark^\checkmark\\checkmarkc\checkmarki\checkmarkr\checkmarkc\checkmark$\checkmarkq\checkmark$\checkmark^\checkmark\\checkmarkc\checkmarki\checkmarkr\checkmarkc\checkmark$\checkmarku\checkmark$\checkmark^\checkmark\\checkmarkc\checkmarki\checkmarkr\checkmarkc\checkmark$\checkmarka\checkmark$\checkmark^\checkmark\\checkmarkc\checkmarki\checkmarkr\checkmarkc\checkmark$\checkmarkt\checkmark$\checkmark^\checkmark\\checkmarkc\checkmarki\checkmarkr\checkmarkc\checkmark$\checkmarki\checkmark$\checkmark^\checkmark\\checkmarkc\checkmarki\checkmarkr\checkmarkc\checkmark$\checkmarko\checkmark$\checkmark^\checkmark\\checkmarkc\checkmarki\checkmarkr\checkmarkc\checkmark$\checkmarkn\checkmark$\checkmark^\checkmark\\checkmarkc\checkmarki\checkmarkr\checkmarkc\checkmark$\checkmark}\checkmark$\checkmark^\checkmark\\checkmarkc\checkmarki\checkmarkr\checkmarkc\checkmark$\checkmark
\checkmark$\checkmark^\checkmark\\checkmarkc\checkmarki\checkmarkr\checkmarkc\checkmark$\checkmark \checkmark$\checkmark^\checkmark\\checkmarkc\checkmarki\checkmarkr\checkmarkc\checkmark$\checkmark \checkmark$\checkmark^\checkmark\\checkmarkc\checkmarki\checkmarkr\checkmarkc\checkmark$\checkmark \checkmark$\checkmark^\checkmark\\checkmarkc\checkmarki\checkmarkr\checkmarkc\checkmark$\checkmark \checkmark$\checkmark^\checkmark\\checkmarkc\checkmarki\checkmarkr\checkmarkc\checkmark$\checkmarkC\checkmark$\checkmark^\checkmark\\checkmarkc\checkmarki\checkmarkr\checkmarkc\checkmark$\checkmark_\checkmark$\checkmark^\checkmark\\checkmarkc\checkmarki\checkmarkr\checkmarkc\checkmark$\checkmarkr\checkmark$\checkmark^\checkmark\\checkmarkc\checkmarki\checkmarkr\checkmarkc\checkmark$\checkmark \checkmark$\checkmark^\checkmark\\checkmarkc\checkmarki\checkmarkr\checkmarkc\checkmark$\checkmark=\checkmark$\checkmark^\checkmark\\checkmarkc\checkmarki\checkmarkr\checkmarkc\checkmark$\checkmark \checkmark$\checkmark^\checkmark\\checkmarkc\checkmarki\checkmarkr\checkmarkc\checkmark$\checkmark\\checkmark$\checkmark^\checkmark\\checkmarkc\checkmarki\checkmarkr\checkmarkc\checkmark$\checkmarkf\checkmark$\checkmark^\checkmark\\checkmarkc\checkmarki\checkmarkr\checkmarkc\checkmark$\checkmarkr\checkmark$\checkmark^\checkmark\\checkmarkc\checkmarki\checkmarkr\checkmarkc\checkmark$\checkmarka\checkmark$\checkmark^\checkmark\\checkmarkc\checkmarki\checkmarkr\checkmarkc\checkmark$\checkmarkc\checkmark$\checkmark^\checkmark\\checkmarkc\checkmarki\checkmarkr\checkmarkc\checkmark$\checkmark{\checkmark$\checkmark^\checkmark\\checkmarkc\checkmarki\checkmarkr\checkmarkc\checkmark$\checkmarkF\checkmark$\checkmark^\checkmark\\checkmarkc\checkmarki\checkmarkr\checkmarkc\checkmark$\checkmark_\checkmark$\checkmark^\checkmark\\checkmarkc\checkmarki\checkmarkr\checkmarkc\checkmark$\checkmark{\checkmark$\checkmark^\checkmark\\checkmarkc\checkmarki\checkmarkr\checkmarkc\checkmark$\checkmarka\checkmark$\checkmark^\checkmark\\checkmarkc\checkmarki\checkmarkr\checkmarkc\checkmark$\checkmarkx\checkmark$\checkmark^\checkmark\\checkmarkc\checkmarki\checkmarkr\checkmarkc\checkmark$\checkmark}\checkmark$\checkmark^\checkmark\\checkmarkc\checkmarki\checkmarkr\checkmarkc\checkmark$\checkmark \checkmark$\checkmark^\checkmark\\checkmarkc\checkmarki\checkmarkr\checkmarkc\checkmark$\checkmark\\checkmark$\checkmark^\checkmark\\checkmarkc\checkmarki\checkmarkr\checkmarkc\checkmark$\checkmarkc\checkmark$\checkmark^\checkmark\\checkmarkc\checkmarki\checkmarkr\checkmarkc\checkmark$\checkmarkd\checkmark$\checkmark^\checkmark\\checkmarkc\checkmarki\checkmarkr\checkmarkc\checkmark$\checkmarko\checkmark$\checkmark^\checkmark\\checkmarkc\checkmarki\checkmarkr\checkmarkc\checkmark$\checkmarkt\checkmark$\checkmark^\checkmark\\checkmarkc\checkmarki\checkmarkr\checkmarkc\checkmark$\checkmark \checkmark$\checkmark^\checkmark\\checkmarkc\checkmarki\checkmarkr\checkmarkc\checkmark$\checkmarkP\checkmark$\checkmark^\checkmark\\checkmarkc\checkmarki\checkmarkr\checkmarkc\checkmark$\checkmark_\checkmark$\checkmark^\checkmark\\checkmarkc\checkmarki\checkmarkr\checkmarkc\checkmark$\checkmarkh\checkmark$\checkmark^\checkmark\\checkmarkc\checkmarki\checkmarkr\checkmarkc\checkmark$\checkmark}\checkmark$\checkmark^\checkmark\\checkmarkc\checkmarki\checkmarkr\checkmarkc\checkmark$\checkmark{\checkmark$\checkmark^\checkmark\\checkmarkc\checkmarki\checkmarkr\checkmarkc\checkmark$\checkmark2\checkmark$\checkmark^\checkmark\\checkmarkc\checkmarki\checkmarkr\checkmarkc\checkmark$\checkmark\\checkmark$\checkmark^\checkmark\\checkmarkc\checkmarki\checkmarkr\checkmarkc\checkmark$\checkmarkp\checkmark$\checkmark^\checkmark\\checkmarkc\checkmarki\checkmarkr\checkmarkc\checkmark$\checkmarki\checkmark$\checkmark^\checkmark\\checkmarkc\checkmarki\checkmarkr\checkmarkc\checkmark$\checkmark \checkmark$\checkmark^\checkmark\\checkmarkc\checkmarki\checkmarkr\checkmarkc\checkmark$\checkmark\\checkmark$\checkmark^\checkmark\\checkmarkc\checkmarki\checkmarkr\checkmarkc\checkmark$\checkmarkc\checkmark$\checkmark^\checkmark\\checkmarkc\checkmarki\checkmarkr\checkmarkc\checkmark$\checkmarkd\checkmark$\checkmark^\checkmark\\checkmarkc\checkmarki\checkmarkr\checkmarkc\checkmark$\checkmarko\checkmark$\checkmark^\checkmark\\checkmarkc\checkmarki\checkmarkr\checkmarkc\checkmark$\checkmarkt\checkmark$\checkmark^\checkmark\\checkmarkc\checkmarki\checkmarkr\checkmarkc\checkmark$\checkmark \checkmark$\checkmark^\checkmark\\checkmarkc\checkmarki\checkmarkr\checkmarkc\checkmark$\checkmark\\checkmark$\checkmark^\checkmark\\checkmarkc\checkmarki\checkmarkr\checkmarkc\checkmark$\checkmarke\checkmark$\checkmark^\checkmark\\checkmarkc\checkmarki\checkmarkr\checkmarkc\checkmark$\checkmarkt\checkmark$\checkmark^\checkmark\\checkmarkc\checkmarki\checkmarkr\checkmarkc\checkmark$\checkmarka\checkmark$\checkmark^\checkmark\\checkmarkc\checkmarki\checkmarkr\checkmarkc\checkmark$\checkmark}\checkmark$\checkmark^\checkmark\\checkmarkc\checkmarki\checkmarkr\checkmarkc\checkmark$\checkmark \checkmark$\checkmark^\checkmark\\checkmarkc\checkmarki\checkmarkr\checkmarkc\checkmark$\checkmark=\checkmark$\checkmark^\checkmark\\checkmarkc\checkmarki\checkmarkr\checkmarkc\checkmark$\checkmark \checkmark$\checkmark^\checkmark\\checkmarkc\checkmarki\checkmarkr\checkmarkc\checkmark$\checkmark\\checkmark$\checkmark^\checkmark\\checkmarkc\checkmarki\checkmarkr\checkmarkc\checkmark$\checkmarkf\checkmark$\checkmark^\checkmark\\checkmarkc\checkmarki\checkmarkr\checkmarkc\checkmark$\checkmarkr\checkmark$\checkmark^\checkmark\\checkmarkc\checkmarki\checkmarkr\checkmarkc\checkmark$\checkmarka\checkmark$\checkmark^\checkmark\\checkmarkc\checkmarki\checkmarkr\checkmarkc\checkmark$\checkmarkc\checkmark$\checkmark^\checkmark\\checkmarkc\checkmarki\checkmarkr\checkmarkc\checkmark$\checkmark{\checkmark$\checkmark^\checkmark\\checkmarkc\checkmarki\checkmarkr\checkmarkc\checkmark$\checkmark1\checkmark$\checkmark^\checkmark\\checkmarkc\checkmarki\checkmarkr\checkmarkc\checkmark$\checkmark0\checkmark$\checkmark^\checkmark\\checkmarkc\checkmarki\checkmarkr\checkmarkc\checkmark$\checkmark5\checkmark$\checkmark^\checkmark\\checkmarkc\checkmarki\checkmarkr\checkmarkc\checkmark$\checkmark{\checkmark$\checkmark^\checkmark\\checkmarkc\checkmarki\checkmarkr\checkmarkc\checkmark$\checkmark,\checkmark$\checkmark^\checkmark\\checkmarkc\checkmarki\checkmarkr\checkmarkc\checkmark$\checkmark}\checkmark$\checkmark^\checkmark\\checkmarkc\checkmarki\checkmarkr\checkmarkc\checkmark$\checkmark5\checkmark$\checkmark^\checkmark\\checkmarkc\checkmarki\checkmarkr\checkmarkc\checkmark$\checkmark \checkmark$\checkmark^\checkmark\\checkmarkc\checkmarki\checkmarkr\checkmarkc\checkmark$\checkmark\\checkmark$\checkmark^\checkmark\\checkmarkc\checkmarki\checkmarkr\checkmarkc\checkmark$\checkmarkt\checkmark$\checkmark^\checkmark\\checkmarkc\checkmarki\checkmarkr\checkmarkc\checkmark$\checkmarki\checkmark$\checkmark^\checkmark\\checkmarkc\checkmarki\checkmarkr\checkmarkc\checkmark$\checkmarkm\checkmark$\checkmark^\checkmark\\checkmarkc\checkmarki\checkmarkr\checkmarkc\checkmark$\checkmarke\checkmark$\checkmark^\checkmark\\checkmarkc\checkmarki\checkmarkr\checkmarkc\checkmark$\checkmarks\checkmark$\checkmark^\checkmark\\checkmarkc\checkmarki\checkmarkr\checkmarkc\checkmark$\checkmark \checkmark$\checkmark^\checkmark\\checkmarkc\checkmarki\checkmarkr\checkmarkc\checkmark$\checkmark0\checkmark$\checkmark^\checkmark\\checkmarkc\checkmarki\checkmarkr\checkmarkc\checkmark$\checkmark{\checkmark$\checkmark^\checkmark\\checkmarkc\checkmarki\checkmarkr\checkmarkc\checkmark$\checkmark,\checkmark$\checkmark^\checkmark\\checkmarkc\checkmarki\checkmarkr\checkmarkc\checkmark$\checkmark}\checkmark$\checkmark^\checkmark\\checkmarkc\checkmarki\checkmarkr\checkmarkc\checkmark$\checkmark0\checkmark$\checkmark^\checkmark\\checkmarkc\checkmarki\checkmarkr\checkmarkc\checkmark$\checkmark0\checkmark$\checkmark^\checkmark\\checkmarkc\checkmarki\checkmarkr\checkmarkc\checkmark$\checkmark5\checkmark$\checkmark^\checkmark\\checkmarkc\checkmarki\checkmarkr\checkmarkc\checkmark$\checkmark}\checkmark$\checkmark^\checkmark\\checkmarkc\checkmarki\checkmarkr\checkmarkc\checkmark$\checkmark{\checkmark$\checkmark^\checkmark\\checkmarkc\checkmarki\checkmarkr\checkmarkc\checkmark$\checkmark2\checkmark$\checkmark^\checkmark\\checkmarkc\checkmarki\checkmarkr\checkmarkc\checkmark$\checkmark\\checkmark$\checkmark^\checkmark\\checkmarkc\checkmarki\checkmarkr\checkmarkc\checkmark$\checkmarkp\checkmark$\checkmark^\checkmark\\checkmarkc\checkmarki\checkmarkr\checkmarkc\checkmark$\checkmarki\checkmark$\checkmark^\checkmark\\checkmarkc\checkmarki\checkmarkr\checkmarkc\checkmark$\checkmark \checkmark$\checkmark^\checkmark\\checkmarkc\checkmarki\checkmarkr\checkmarkc\checkmark$\checkmark\\checkmark$\checkmark^\checkmark\\checkmarkc\checkmarki\checkmarkr\checkmarkc\checkmark$\checkmarkt\checkmark$\checkmark^\checkmark\\checkmarkc\checkmarki\checkmarkr\checkmarkc\checkmark$\checkmarki\checkmark$\checkmark^\checkmark\\checkmarkc\checkmarki\checkmarkr\checkmarkc\checkmark$\checkmarkm\checkmark$\checkmark^\checkmark\\checkmarkc\checkmarki\checkmarkr\checkmarkc\checkmark$\checkmarke\checkmark$\checkmark^\checkmark\\checkmarkc\checkmarki\checkmarkr\checkmarkc\checkmark$\checkmarks\checkmark$\checkmark^\checkmark\\checkmarkc\checkmarki\checkmarkr\checkmarkc\checkmark$\checkmark \checkmark$\checkmark^\checkmark\\checkmarkc\checkmarki\checkmarkr\checkmarkc\checkmark$\checkmark0\checkmark$\checkmark^\checkmark\\checkmarkc\checkmarki\checkmarkr\checkmarkc\checkmark$\checkmark{\checkmark$\checkmark^\checkmark\\checkmarkc\checkmarki\checkmarkr\checkmarkc\checkmark$\checkmark,\checkmark$\checkmark^\checkmark\\checkmarkc\checkmarki\checkmarkr\checkmarkc\checkmark$\checkmark}\checkmark$\checkmark^\checkmark\\checkmarkc\checkmarki\checkmarkr\checkmarkc\checkmark$\checkmark9\checkmark$\checkmark^\checkmark\\checkmarkc\checkmarki\checkmarkr\checkmarkc\checkmark$\checkmark0\checkmark$\checkmark^\checkmark\\checkmarkc\checkmarki\checkmarkr\checkmarkc\checkmark$\checkmark}\checkmark$\checkmark^\checkmark\\checkmarkc\checkmarki\checkmarkr\checkmarkc\checkmark$\checkmark \checkmark$\checkmark^\checkmark\\checkmarkc\checkmarki\checkmarkr\checkmarkc\checkmark$\checkmark=\checkmark$\checkmark^\checkmark\\checkmarkc\checkmarki\checkmarkr\checkmarkc\checkmark$\checkmark \checkmark$\checkmark^\checkmark\\checkmarkc\checkmarki\checkmarkr\checkmarkc\checkmark$\checkmark\\checkmark$\checkmark^\checkmark\\checkmarkc\checkmarki\checkmarkr\checkmarkc\checkmark$\checkmarkb\checkmark$\checkmark^\checkmark\\checkmarkc\checkmarki\checkmarkr\checkmarkc\checkmark$\checkmarko\checkmark$\checkmark^\checkmark\\checkmarkc\checkmarki\checkmarkr\checkmarkc\checkmark$\checkmarkx\checkmark$\checkmark^\checkmark\\checkmarkc\checkmarki\checkmarkr\checkmarkc\checkmark$\checkmarke\checkmark$\checkmark^\checkmark\\checkmarkc\checkmarki\checkmarkr\checkmarkc\checkmark$\checkmarkd\checkmark$\checkmark^\checkmark\\checkmarkc\checkmarki\checkmarkr\checkmarkc\checkmark$\checkmark{\checkmark$\checkmark^\checkmark\\checkmarkc\checkmarki\checkmarkr\checkmarkc\checkmark$\checkmark0\checkmark$\checkmark^\checkmark\\checkmarkc\checkmarki\checkmarkr\checkmarkc\checkmark$\checkmark{\checkmark$\checkmark^\checkmark\\checkmarkc\checkmarki\checkmarkr\checkmarkc\checkmark$\checkmark,\checkmark$\checkmark^\checkmark\\checkmarkc\checkmarki\checkmarkr\checkmarkc\checkmark$\checkmark}\checkmark$\checkmark^\checkmark\\checkmarkc\checkmarki\checkmarkr\checkmarkc\checkmark$\checkmark0\checkmark$\checkmark^\checkmark\\checkmarkc\checkmarki\checkmarkr\checkmarkc\checkmark$\checkmark9\checkmark$\checkmark^\checkmark\\checkmarkc\checkmarki\checkmarkr\checkmarkc\checkmark$\checkmark3\checkmark$\checkmark^\checkmark\\checkmarkc\checkmarki\checkmarkr\checkmarkc\checkmark$\checkmark \checkmark$\checkmark^\checkmark\\checkmarkc\checkmarki\checkmarkr\checkmarkc\checkmark$\checkmark\\checkmark$\checkmark^\checkmark\\checkmarkc\checkmarki\checkmarkr\checkmarkc\checkmark$\checkmarkt\checkmark$\checkmark^\checkmark\\checkmarkc\checkmarki\checkmarkr\checkmarkc\checkmark$\checkmarke\checkmark$\checkmark^\checkmark\\checkmarkc\checkmarki\checkmarkr\checkmarkc\checkmark$\checkmarkx\checkmark$\checkmark^\checkmark\\checkmarkc\checkmarki\checkmarkr\checkmarkc\checkmark$\checkmarkt\checkmark$\checkmark^\checkmark\\checkmarkc\checkmarki\checkmarkr\checkmarkc\checkmark$\checkmark{\checkmark$\checkmark^\checkmark\\checkmarkc\checkmarki\checkmarkr\checkmarkc\checkmark$\checkmark \checkmark$\checkmark^\checkmark\\checkmarkc\checkmarki\checkmarkr\checkmarkc\checkmark$\checkmarkN\checkmark$\checkmark^\checkmark\\checkmarkc\checkmarki\checkmarkr\checkmarkc\checkmark$\checkmark$\checkmark$\checkmark^\checkmark\\checkmarkc\checkmarki\checkmarkr\checkmarkc\checkmark$\checkmark\\checkmark$\checkmark^\checkmark\\checkmarkc\checkmarki\checkmarkr\checkmarkc\checkmark$\checkmarkc\checkmark$\checkmark^\checkmark\\checkmarkc\checkmarki\checkmarkr\checkmarkc\checkmark$\checkmarkd\checkmark$\checkmark^\checkmark\\checkmarkc\checkmarki\checkmarkr\checkmarkc\checkmark$\checkmarko\checkmark$\checkmark^\checkmark\\checkmarkc\checkmarki\checkmarkr\checkmarkc\checkmark$\checkmarkt\checkmark$\checkmark^\checkmark\\checkmarkc\checkmarki\checkmarkr\checkmarkc\checkmark$\checkmark$\checkmark$\checkmark^\checkmark\\checkmarkc\checkmarki\checkmarkr\checkmarkc\checkmark$\checkmarkm\checkmark$\checkmark^\checkmark\\checkmarkc\checkmarki\checkmarkr\checkmarkc\checkmark$\checkmark}\checkmark$\checkmark^\checkmark\\checkmarkc\checkmarki\checkmarkr\checkmarkc\checkmark$\checkmark}\checkmark$\checkmark^\checkmark\\checkmarkc\checkmarki\checkmarkr\checkmarkc\checkmark$\checkmark
\checkmark$\checkmark^\checkmark\\checkmarkc\checkmarki\checkmarkr\checkmarkc\checkmark$\checkmark \checkmark$\checkmark^\checkmark\\checkmarkc\checkmarki\checkmarkr\checkmarkc\checkmark$\checkmark \checkmark$\checkmark^\checkmark\\checkmarkc\checkmarki\checkmarkr\checkmarkc\checkmark$\checkmark \checkmark$\checkmark^\checkmark\\checkmarkc\checkmarki\checkmarkr\checkmarkc\checkmark$\checkmark \checkmark$\checkmark^\checkmark\\checkmarkc\checkmarki\checkmarkr\checkmarkc\checkmark$\checkmark\\checkmark$\checkmark^\checkmark\\checkmarkc\checkmarki\checkmarkr\checkmarkc\checkmark$\checkmarkl\checkmark$\checkmark^\checkmark\\checkmarkc\checkmarki\checkmarkr\checkmarkc\checkmark$\checkmarka\checkmark$\checkmark^\checkmark\\checkmarkc\checkmarki\checkmarkr\checkmarkc\checkmark$\checkmarkb\checkmark$\checkmark^\checkmark\\checkmarkc\checkmarki\checkmarkr\checkmarkc\checkmark$\checkmarke\checkmark$\checkmark^\checkmark\\checkmarkc\checkmarki\checkmarkr\checkmarkc\checkmark$\checkmarkl\checkmark$\checkmark^\checkmark\\checkmarkc\checkmarki\checkmarkr\checkmarkc\checkmark$\checkmark{\checkmark$\checkmark^\checkmark\\checkmarkc\checkmarki\checkmarkr\checkmarkc\checkmark$\checkmarke\checkmark$\checkmark^\checkmark\\checkmarkc\checkmarki\checkmarkr\checkmarkc\checkmark$\checkmarkq\checkmark$\checkmark^\checkmark\\checkmarkc\checkmarki\checkmarkr\checkmarkc\checkmark$\checkmark:\checkmark$\checkmark^\checkmark\\checkmarkc\checkmarki\checkmarkr\checkmarkc\checkmark$\checkmarkc\checkmark$\checkmark^\checkmark\\checkmarkc\checkmarki\checkmarkr\checkmarkc\checkmark$\checkmarko\checkmark$\checkmark^\checkmark\\checkmarkc\checkmarki\checkmarkr\checkmarkc\checkmark$\checkmarku\checkmark$\checkmark^\checkmark\\checkmarkc\checkmarki\checkmarkr\checkmarkc\checkmark$\checkmarkp\checkmark$\checkmark^\checkmark\\checkmarkc\checkmarki\checkmarkr\checkmarkc\checkmark$\checkmarkl\checkmark$\checkmark^\checkmark\\checkmarkc\checkmarki\checkmarkr\checkmarkc\checkmark$\checkmarke\checkmark$\checkmark^\checkmark\\checkmarkc\checkmarki\checkmarkr\checkmarkc\checkmark$\checkmark_\checkmark$\checkmark^\checkmark\\checkmarkc\checkmarki\checkmarkr\checkmarkc\checkmark$\checkmarkr\checkmark$\checkmark^\checkmark\\checkmarkc\checkmarki\checkmarkr\checkmarkc\checkmark$\checkmarke\checkmark$\checkmark^\checkmark\\checkmarkc\checkmarki\checkmarkr\checkmarkc\checkmark$\checkmarks\checkmark$\checkmark^\checkmark\\checkmarkc\checkmarki\checkmarkr\checkmarkc\checkmark$\checkmarki\checkmark$\checkmark^\checkmark\\checkmarkc\checkmarki\checkmarkr\checkmarkc\checkmark$\checkmarks\checkmark$\checkmark^\checkmark\\checkmarkc\checkmarki\checkmarkr\checkmarkc\checkmark$\checkmarkt\checkmark$\checkmark^\checkmark\\checkmarkc\checkmarki\checkmarkr\checkmarkc\checkmark$\checkmarka\checkmark$\checkmark^\checkmark\\checkmarkc\checkmarki\checkmarkr\checkmarkc\checkmark$\checkmarkn\checkmark$\checkmark^\checkmark\\checkmarkc\checkmarki\checkmarkr\checkmarkc\checkmark$\checkmarkt\checkmark$\checkmark^\checkmark\\checkmarkc\checkmarki\checkmarkr\checkmarkc\checkmark$\checkmark}\checkmark$\checkmark^\checkmark\\checkmarkc\checkmarki\checkmarkr\checkmarkc\checkmark$\checkmark
\checkmark$\checkmark^\checkmark\\checkmarkc\checkmarki\checkmarkr\checkmarkc\checkmark$\checkmark\\checkmark$\checkmark^\checkmark\\checkmarkc\checkmarki\checkmarkr\checkmarkc\checkmark$\checkmarke\checkmark$\checkmark^\checkmark\\checkmarkc\checkmarki\checkmarkr\checkmarkc\checkmark$\checkmarkn\checkmark$\checkmark^\checkmark\\checkmarkc\checkmarki\checkmarkr\checkmarkc\checkmark$\checkmarkd\checkmark$\checkmark^\checkmark\\checkmarkc\checkmarki\checkmarkr\checkmarkc\checkmark$\checkmark{\checkmark$\checkmark^\checkmark\\checkmarkc\checkmarki\checkmarkr\checkmarkc\checkmark$\checkmarke\checkmark$\checkmark^\checkmark\\checkmarkc\checkmarki\checkmarkr\checkmarkc\checkmark$\checkmarkq\checkmark$\checkmark^\checkmark\\checkmarkc\checkmarki\checkmarkr\checkmarkc\checkmark$\checkmarku\checkmark$\checkmark^\checkmark\\checkmarkc\checkmarki\checkmarkr\checkmarkc\checkmark$\checkmarka\checkmark$\checkmark^\checkmark\\checkmarkc\checkmarki\checkmarkr\checkmarkc\checkmark$\checkmarkt\checkmark$\checkmark^\checkmark\\checkmarkc\checkmarki\checkmarkr\checkmarkc\checkmark$\checkmarki\checkmark$\checkmark^\checkmark\\checkmarkc\checkmarki\checkmarkr\checkmarkc\checkmark$\checkmarko\checkmark$\checkmark^\checkmark\\checkmarkc\checkmarki\checkmarkr\checkmarkc\checkmark$\checkmarkn\checkmark$\checkmark^\checkmark\\checkmarkc\checkmarki\checkmarkr\checkmarkc\checkmark$\checkmark}\checkmark$\checkmark^\checkmark\\checkmarkc\checkmarki\checkmarkr\checkmarkc\checkmark$\checkmark
\checkmark$\checkmark^\checkmark\\checkmarkc\checkmarki\checkmarkr\checkmarkc\checkmark$\checkmark
\checkmark$\checkmark^\checkmark\\checkmarkc\checkmarki\checkmarkr\checkmarkc\checkmark$\checkmark\\checkmark$\checkmark^\checkmark\\checkmarkc\checkmarki\checkmarkr\checkmarkc\checkmark$\checkmarks\checkmark$\checkmark^\checkmark\\checkmarkc\checkmarki\checkmarkr\checkmarkc\checkmark$\checkmarku\checkmark$\checkmark^\checkmark\\checkmarkc\checkmarki\checkmarkr\checkmarkc\checkmark$\checkmarkb\checkmark$\checkmark^\checkmark\\checkmarkc\checkmarki\checkmarkr\checkmarkc\checkmark$\checkmarks\checkmark$\checkmark^\checkmark\\checkmarkc\checkmarki\checkmarkr\checkmarkc\checkmark$\checkmarke\checkmark$\checkmark^\checkmark\\checkmarkc\checkmarki\checkmarkr\checkmarkc\checkmark$\checkmarkc\checkmark$\checkmark^\checkmark\\checkmarkc\checkmarki\checkmarkr\checkmarkc\checkmark$\checkmarkt\checkmark$\checkmark^\checkmark\\checkmarkc\checkmarki\checkmarkr\checkmarkc\checkmark$\checkmarki\checkmark$\checkmark^\checkmark\\checkmarkc\checkmarki\checkmarkr\checkmarkc\checkmark$\checkmarko\checkmark$\checkmark^\checkmark\\checkmarkc\checkmarki\checkmarkr\checkmarkc\checkmark$\checkmarkn\checkmark$\checkmark^\checkmark\\checkmarkc\checkmarki\checkmarkr\checkmarkc\checkmark$\checkmark{\checkmark$\checkmark^\checkmark\\checkmarkc\checkmarki\checkmarkr\checkmarkc\checkmark$\checkmarkD\checkmark$\checkmark^\checkmark\\checkmarkc\checkmarki\checkmarkr\checkmarkc\checkmark$\checkmarké\checkmark$\checkmark^\checkmark\\checkmarkc\checkmarki\checkmarkr\checkmarkc\checkmark$\checkmarkt\checkmark$\checkmark^\checkmark\\checkmarkc\checkmarki\checkmarkr\checkmarkc\checkmark$\checkmarke\checkmark$\checkmark^\checkmark\\checkmarkc\checkmarki\checkmarkr\checkmarkc\checkmark$\checkmarkr\checkmark$\checkmark^\checkmark\\checkmarkc\checkmarki\checkmarkr\checkmarkc\checkmark$\checkmarkm\checkmark$\checkmark^\checkmark\\checkmarkc\checkmarki\checkmarkr\checkmarkc\checkmark$\checkmarki\checkmark$\checkmark^\checkmark\\checkmarkc\checkmarki\checkmarkr\checkmarkc\checkmark$\checkmarkn\checkmark$\checkmark^\checkmark\\checkmarkc\checkmarki\checkmarkr\checkmarkc\checkmark$\checkmarka\checkmark$\checkmark^\checkmark\\checkmarkc\checkmarki\checkmarkr\checkmarkc\checkmark$\checkmarkt\checkmark$\checkmark^\checkmark\\checkmarkc\checkmarki\checkmarkr\checkmarkc\checkmark$\checkmarki\checkmark$\checkmark^\checkmark\\checkmarkc\checkmarki\checkmarkr\checkmarkc\checkmark$\checkmarko\checkmark$\checkmark^\checkmark\\checkmarkc\checkmarki\checkmarkr\checkmarkc\checkmark$\checkmarkn\checkmark$\checkmark^\checkmark\\checkmarkc\checkmarki\checkmarkr\checkmarkc\checkmark$\checkmark \checkmark$\checkmark^\checkmark\\checkmarkc\checkmarki\checkmarkr\checkmarkc\checkmark$\checkmarkd\checkmark$\checkmark^\checkmark\\checkmarkc\checkmarki\checkmarkr\checkmarkc\checkmark$\checkmarku\checkmark$\checkmark^\checkmark\\checkmarkc\checkmarki\checkmarkr\checkmarkc\checkmark$\checkmark \checkmark$\checkmark^\checkmark\\checkmarkc\checkmarki\checkmarkr\checkmarkc\checkmark$\checkmarkc\checkmark$\checkmark^\checkmark\\checkmarkc\checkmarki\checkmarkr\checkmarkc\checkmark$\checkmarko\checkmark$\checkmark^\checkmark\\checkmarkc\checkmarki\checkmarkr\checkmarkc\checkmark$\checkmarku\checkmark$\checkmark^\checkmark\\checkmarkc\checkmarki\checkmarkr\checkmarkc\checkmark$\checkmarkp\checkmark$\checkmark^\checkmark\\checkmarkc\checkmarki\checkmarkr\checkmarkc\checkmark$\checkmarkl\checkmark$\checkmark^\checkmark\\checkmarkc\checkmarki\checkmarkr\checkmarkc\checkmark$\checkmarke\checkmark$\checkmark^\checkmark\\checkmarkc\checkmarki\checkmarkr\checkmarkc\checkmark$\checkmark \checkmark$\checkmark^\checkmark\\checkmarkc\checkmarki\checkmarkr\checkmarkc\checkmark$\checkmarkd\checkmark$\checkmark^\checkmark\\checkmarkc\checkmarki\checkmarkr\checkmarkc\checkmark$\checkmark'\checkmark$\checkmark^\checkmark\\checkmarkc\checkmarki\checkmarkr\checkmarkc\checkmark$\checkmarki\checkmark$\checkmark^\checkmark\\checkmarkc\checkmarki\checkmarkr\checkmarkc\checkmark$\checkmarkn\checkmark$\checkmark^\checkmark\\checkmarkc\checkmarki\checkmarkr\checkmarkc\checkmark$\checkmarke\checkmark$\checkmark^\checkmark\\checkmarkc\checkmarki\checkmarkr\checkmarkc\checkmark$\checkmarkr\checkmark$\checkmark^\checkmark\\checkmarkc\checkmarki\checkmarkr\checkmarkc\checkmark$\checkmarkt\checkmark$\checkmark^\checkmark\\checkmarkc\checkmarki\checkmarkr\checkmarkc\checkmark$\checkmarki\checkmark$\checkmark^\checkmark\\checkmarkc\checkmarki\checkmarkr\checkmarkc\checkmark$\checkmarke\checkmark$\checkmark^\checkmark\\checkmarkc\checkmarki\checkmarkr\checkmarkc\checkmark$\checkmark}\checkmark$\checkmark^\checkmark\\checkmarkc\checkmarki\checkmarkr\checkmarkc\checkmark$\checkmark
\checkmark$\checkmark^\checkmark\\checkmarkc\checkmarki\checkmarkr\checkmarkc\checkmark$\checkmarkL\checkmark$\checkmark^\checkmark\\checkmarkc\checkmarki\checkmarkr\checkmarkc\checkmark$\checkmark'\checkmark$\checkmark^\checkmark\\checkmarkc\checkmarki\checkmarkr\checkmarkc\checkmark$\checkmarki\checkmark$\checkmark^\checkmark\\checkmarkc\checkmarki\checkmarkr\checkmarkc\checkmark$\checkmarkn\checkmark$\checkmark^\checkmark\\checkmarkc\checkmarki\checkmarkr\checkmarkc\checkmark$\checkmarke\checkmark$\checkmark^\checkmark\\checkmarkc\checkmarki\checkmarkr\checkmarkc\checkmark$\checkmarkr\checkmark$\checkmark^\checkmark\\checkmarkc\checkmarki\checkmarkr\checkmarkc\checkmark$\checkmarkt\checkmark$\checkmark^\checkmark\\checkmarkc\checkmarki\checkmarkr\checkmarkc\checkmark$\checkmarki\checkmark$\checkmark^\checkmark\\checkmarkc\checkmarki\checkmarkr\checkmarkc\checkmark$\checkmarke\checkmark$\checkmark^\checkmark\\checkmarkc\checkmarki\checkmarkr\checkmarkc\checkmark$\checkmark \checkmark$\checkmark^\checkmark\\checkmarkc\checkmarki\checkmarkr\checkmarkc\checkmark$\checkmarkt\checkmark$\checkmark^\checkmark\\checkmarkc\checkmarki\checkmarkr\checkmarkc\checkmark$\checkmarko\checkmark$\checkmark^\checkmark\\checkmarkc\checkmarki\checkmarkr\checkmarkc\checkmark$\checkmarkt\checkmark$\checkmark^\checkmark\\checkmarkc\checkmarki\checkmarkr\checkmarkc\checkmark$\checkmarka\checkmark$\checkmark^\checkmark\\checkmarkc\checkmarki\checkmarkr\checkmarkc\checkmark$\checkmarkl\checkmark$\checkmark^\checkmark\\checkmarkc\checkmarki\checkmarkr\checkmarkc\checkmark$\checkmarke\checkmark$\checkmark^\checkmark\\checkmarkc\checkmarki\checkmarkr\checkmarkc\checkmark$\checkmark \checkmark$\checkmark^\checkmark\\checkmarkc\checkmarki\checkmarkr\checkmarkc\checkmark$\checkmarkr\checkmark$\checkmark^\checkmark\\checkmarkc\checkmarki\checkmarkr\checkmarkc\checkmark$\checkmarka\checkmark$\checkmark^\checkmark\\checkmarkc\checkmarki\checkmarkr\checkmarkc\checkmark$\checkmarkm\checkmark$\checkmark^\checkmark\\checkmarkc\checkmarki\checkmarkr\checkmarkc\checkmark$\checkmarke\checkmark$\checkmark^\checkmark\\checkmarkc\checkmarki\checkmarkr\checkmarkc\checkmark$\checkmarkn\checkmark$\checkmark^\checkmark\\checkmarkc\checkmarki\checkmarkr\checkmarkc\checkmark$\checkmarké\checkmark$\checkmark^\checkmark\\checkmarkc\checkmarki\checkmarkr\checkmarkc\checkmark$\checkmarke\checkmark$\checkmark^\checkmark\\checkmarkc\checkmarki\checkmarkr\checkmarkc\checkmark$\checkmark \checkmark$\checkmark^\checkmark\\checkmarkc\checkmarki\checkmarkr\checkmarkc\checkmark$\checkmarks\checkmark$\checkmark^\checkmark\\checkmarkc\checkmarki\checkmarkr\checkmarkc\checkmark$\checkmarku\checkmark$\checkmark^\checkmark\\checkmarkc\checkmarki\checkmarkr\checkmarkc\checkmark$\checkmarkr\checkmark$\checkmark^\checkmark\\checkmarkc\checkmarki\checkmarkr\checkmarkc\checkmark$\checkmark \checkmark$\checkmark^\checkmark\\checkmarkc\checkmarki\checkmarkr\checkmarkc\checkmark$\checkmarkl\checkmark$\checkmark^\checkmark\\checkmarkc\checkmarki\checkmarkr\checkmarkc\checkmark$\checkmark'\checkmark$\checkmark^\checkmark\\checkmarkc\checkmarki\checkmarkr\checkmarkc\checkmark$\checkmarka\checkmark$\checkmark^\checkmark\\checkmarkc\checkmarki\checkmarkr\checkmarkc\checkmark$\checkmarkx\checkmark$\checkmark^\checkmark\\checkmarkc\checkmarki\checkmarkr\checkmarkc\checkmark$\checkmarke\checkmark$\checkmark^\checkmark\\checkmarkc\checkmarki\checkmarkr\checkmarkc\checkmark$\checkmark \checkmark$\checkmark^\checkmark\\checkmarkc\checkmarki\checkmarkr\checkmarkc\checkmark$\checkmarkm\checkmark$\checkmark^\checkmark\\checkmarkc\checkmarki\checkmarkr\checkmarkc\checkmark$\checkmarko\checkmark$\checkmark^\checkmark\\checkmarkc\checkmarki\checkmarkr\checkmarkc\checkmark$\checkmarkt\checkmark$\checkmark^\checkmark\\checkmarkc\checkmarki\checkmarkr\checkmarkc\checkmark$\checkmarke\checkmark$\checkmark^\checkmark\\checkmarkc\checkmarki\checkmarkr\checkmarkc\checkmark$\checkmarku\checkmark$\checkmark^\checkmark\\checkmarkc\checkmarki\checkmarkr\checkmarkc\checkmark$\checkmarkr\checkmark$\checkmark^\checkmark\\checkmarkc\checkmarki\checkmarkr\checkmarkc\checkmark$\checkmark \checkmark$\checkmark^\checkmark\\checkmarkc\checkmarki\checkmarkr\checkmarkc\checkmark$\checkmarke\checkmark$\checkmark^\checkmark\\checkmarkc\checkmarki\checkmarkr\checkmarkc\checkmark$\checkmarks\checkmark$\checkmark^\checkmark\\checkmarkc\checkmarki\checkmarkr\checkmarkc\checkmark$\checkmarkt\checkmark$\checkmark^\checkmark\\checkmarkc\checkmarki\checkmarkr\checkmarkc\checkmark$\checkmark \checkmark$\checkmark^\checkmark\\checkmarkc\checkmarki\checkmarkr\checkmarkc\checkmark$\checkmarkl\checkmark$\checkmark^\checkmark\\checkmarkc\checkmarki\checkmarkr\checkmarkc\checkmark$\checkmarka\checkmark$\checkmark^\checkmark\\checkmarkc\checkmarki\checkmarkr\checkmarkc\checkmark$\checkmark \checkmark$\checkmark^\checkmark\\checkmarkc\checkmarki\checkmarkr\checkmarkc\checkmark$\checkmarks\checkmark$\checkmark^\checkmark\\checkmarkc\checkmarki\checkmarkr\checkmarkc\checkmark$\checkmarko\checkmark$\checkmark^\checkmark\\checkmarkc\checkmarki\checkmarkr\checkmarkc\checkmark$\checkmarkm\checkmark$\checkmark^\checkmark\\checkmarkc\checkmarki\checkmarkr\checkmarkc\checkmark$\checkmarkm\checkmark$\checkmark^\checkmark\\checkmarkc\checkmarki\checkmarkr\checkmarkc\checkmark$\checkmarke\checkmark$\checkmark^\checkmark\\checkmarkc\checkmarki\checkmarkr\checkmarkc\checkmark$\checkmark \checkmark$\checkmark^\checkmark\\checkmarkc\checkmarki\checkmarkr\checkmarkc\checkmark$\checkmarkd\checkmark$\checkmark^\checkmark\\checkmarkc\checkmarki\checkmarkr\checkmarkc\checkmark$\checkmarke\checkmark$\checkmark^\checkmark\\checkmarkc\checkmarki\checkmarkr\checkmarkc\checkmark$\checkmark \checkmark$\checkmark^\checkmark\\checkmarkc\checkmarki\checkmarkr\checkmarkc\checkmark$\checkmarkl\checkmark$\checkmark^\checkmark\\checkmarkc\checkmarki\checkmarkr\checkmarkc\checkmark$\checkmark'\checkmark$\checkmark^\checkmark\\checkmarkc\checkmarki\checkmarkr\checkmarkc\checkmark$\checkmarki\checkmark$\checkmark^\checkmark\\checkmarkc\checkmarki\checkmarkr\checkmarkc\checkmark$\checkmarkn\checkmark$\checkmark^\checkmark\\checkmarkc\checkmarki\checkmarkr\checkmarkc\checkmark$\checkmarke\checkmark$\checkmark^\checkmark\\checkmarkc\checkmarki\checkmarkr\checkmarkc\checkmark$\checkmarkr\checkmark$\checkmark^\checkmark\\checkmarkc\checkmarki\checkmarkr\checkmarkc\checkmark$\checkmarkt\checkmark$\checkmark^\checkmark\\checkmarkc\checkmarki\checkmarkr\checkmarkc\checkmark$\checkmarki\checkmark$\checkmark^\checkmark\\checkmarkc\checkmarki\checkmarkr\checkmarkc\checkmark$\checkmarke\checkmark$\checkmark^\checkmark\\checkmarkc\checkmarki\checkmarkr\checkmarkc\checkmark$\checkmark \checkmark$\checkmark^\checkmark\\checkmarkc\checkmarki\checkmarkr\checkmarkc\checkmark$\checkmarkd\checkmark$\checkmark^\checkmark\\checkmarkc\checkmarki\checkmarkr\checkmarkc\checkmark$\checkmarke\checkmark$\checkmark^\checkmark\\checkmarkc\checkmarki\checkmarkr\checkmarkc\checkmark$\checkmark \checkmark$\checkmark^\checkmark\\checkmarkc\checkmarki\checkmarkr\checkmarkc\checkmark$\checkmarkl\checkmark$\checkmark^\checkmark\\checkmarkc\checkmarki\checkmarkr\checkmarkc\checkmark$\checkmarka\checkmark$\checkmark^\checkmark\\checkmarkc\checkmarki\checkmarkr\checkmarkc\checkmark$\checkmark \checkmark$\checkmark^\checkmark\\checkmarkc\checkmarki\checkmarkr\checkmarkc\checkmark$\checkmarkv\checkmark$\checkmark^\checkmark\\checkmarkc\checkmarki\checkmarkr\checkmarkc\checkmark$\checkmarki\checkmark$\checkmark^\checkmark\\checkmarkc\checkmarki\checkmarkr\checkmarkc\checkmark$\checkmarks\checkmark$\checkmark^\checkmark\\checkmarkc\checkmarki\checkmarkr\checkmarkc\checkmark$\checkmark \checkmark$\checkmark^\checkmark\\checkmarkc\checkmarki\checkmarkr\checkmarkc\checkmark$\checkmarke\checkmark$\checkmark^\checkmark\\checkmarkc\checkmarki\checkmarkr\checkmarkc\checkmark$\checkmarkt\checkmark$\checkmark^\checkmark\\checkmarkc\checkmarki\checkmarkr\checkmarkc\checkmark$\checkmark \checkmark$\checkmark^\checkmark\\checkmarkc\checkmarki\checkmarkr\checkmarkc\checkmark$\checkmarkd\checkmark$\checkmark^\checkmark\\checkmarkc\checkmarki\checkmarkr\checkmarkc\checkmark$\checkmarke\checkmark$\checkmark^\checkmark\\checkmarkc\checkmarki\checkmarkr\checkmarkc\checkmark$\checkmark \checkmark$\checkmark^\checkmark\\checkmarkc\checkmarki\checkmarkr\checkmarkc\checkmark$\checkmarkl\checkmark$\checkmark^\checkmark\\checkmarkc\checkmarki\checkmarkr\checkmarkc\checkmark$\checkmark'\checkmark$\checkmark^\checkmark\\checkmarkc\checkmarki\checkmarkr\checkmarkc\checkmark$\checkmarki\checkmark$\checkmark^\checkmark\\checkmarkc\checkmarki\checkmarkr\checkmarkc\checkmark$\checkmarkn\checkmark$\checkmark^\checkmark\\checkmarkc\checkmarki\checkmarkr\checkmarkc\checkmark$\checkmarke\checkmark$\checkmark^\checkmark\\checkmarkc\checkmarki\checkmarkr\checkmarkc\checkmark$\checkmarkr\checkmark$\checkmark^\checkmark\\checkmarkc\checkmarki\checkmarkr\checkmarkc\checkmark$\checkmarkt\checkmark$\checkmark^\checkmark\\checkmarkc\checkmarki\checkmarkr\checkmarkc\checkmark$\checkmarki\checkmark$\checkmark^\checkmark\\checkmarkc\checkmarki\checkmarkr\checkmarkc\checkmark$\checkmarke\checkmark$\checkmark^\checkmark\\checkmarkc\checkmarki\checkmarkr\checkmarkc\checkmark$\checkmark \checkmark$\checkmark^\checkmark\\checkmarkc\checkmarki\checkmarkr\checkmarkc\checkmark$\checkmarkd\checkmark$\checkmark^\checkmark\\checkmarkc\checkmarki\checkmarkr\checkmarkc\checkmark$\checkmarke\checkmark$\checkmark^\checkmark\\checkmarkc\checkmarki\checkmarkr\checkmarkc\checkmark$\checkmark \checkmark$\checkmark^\checkmark\\checkmarkc\checkmarki\checkmarkr\checkmarkc\checkmark$\checkmarkl\checkmark$\checkmark^\checkmark\\checkmarkc\checkmarki\checkmarkr\checkmarkc\checkmark$\checkmarka\checkmark$\checkmark^\checkmark\\checkmarkc\checkmarki\checkmarkr\checkmarkc\checkmark$\checkmark \checkmark$\checkmark^\checkmark\\checkmarkc\checkmarki\checkmarkr\checkmarkc\checkmark$\checkmarkm\checkmark$\checkmark^\checkmark\\checkmarkc\checkmarki\checkmarkr\checkmarkc\checkmark$\checkmarka\checkmark$\checkmark^\checkmark\\checkmarkc\checkmarki\checkmarkr\checkmarkc\checkmark$\checkmarks\checkmark$\checkmark^\checkmark\\checkmarkc\checkmarki\checkmarkr\checkmarkc\checkmark$\checkmarks\checkmark$\checkmark^\checkmark\\checkmarkc\checkmarki\checkmarkr\checkmarkc\checkmark$\checkmarke\checkmark$\checkmark^\checkmark\\checkmarkc\checkmarki\checkmarkr\checkmarkc\checkmark$\checkmark \checkmark$\checkmark^\checkmark\\checkmarkc\checkmarki\checkmarkr\checkmarkc\checkmark$\checkmarke\checkmark$\checkmark^\checkmark\\checkmarkc\checkmarki\checkmarkr\checkmarkc\checkmark$\checkmarkn\checkmark$\checkmark^\checkmark\\checkmarkc\checkmarki\checkmarkr\checkmarkc\checkmark$\checkmark \checkmark$\checkmark^\checkmark\\checkmarkc\checkmarki\checkmarkr\checkmarkc\checkmark$\checkmarkt\checkmark$\checkmark^\checkmark\\checkmarkc\checkmarki\checkmarkr\checkmarkc\checkmark$\checkmarkr\checkmark$\checkmark^\checkmark\\checkmarkc\checkmarki\checkmarkr\checkmarkc\checkmark$\checkmarka\checkmark$\checkmark^\checkmark\\checkmarkc\checkmarki\checkmarkr\checkmarkc\checkmark$\checkmarkn\checkmark$\checkmark^\checkmark\\checkmarkc\checkmarki\checkmarkr\checkmarkc\checkmark$\checkmarks\checkmark$\checkmark^\checkmark\\checkmarkc\checkmarki\checkmarkr\checkmarkc\checkmark$\checkmarkl\checkmark$\checkmark^\checkmark\\checkmarkc\checkmarki\checkmarkr\checkmarkc\checkmark$\checkmarka\checkmark$\checkmark^\checkmark\\checkmarkc\checkmarki\checkmarkr\checkmarkc\checkmark$\checkmarkt\checkmark$\checkmark^\checkmark\\checkmarkc\checkmarki\checkmarkr\checkmarkc\checkmark$\checkmarki\checkmark$\checkmark^\checkmark\\checkmarkc\checkmarki\checkmarkr\checkmarkc\checkmark$\checkmarko\checkmark$\checkmark^\checkmark\\checkmarkc\checkmarki\checkmarkr\checkmarkc\checkmark$\checkmarkn\checkmark$\checkmark^\checkmark\\checkmarkc\checkmarki\checkmarkr\checkmarkc\checkmark$\checkmark.\checkmark$\checkmark^\checkmark\\checkmarkc\checkmarki\checkmarkr\checkmarkc\checkmark$\checkmark
\checkmark$\checkmark^\checkmark\\checkmarkc\checkmarki\checkmarkr\checkmarkc\checkmark$\checkmark
\checkmark$\checkmark^\checkmark\\checkmarkc\checkmarki\checkmarkr\checkmarkc\checkmark$\checkmark\\checkmark$\checkmark^\checkmark\\checkmarkc\checkmarki\checkmarkr\checkmarkc\checkmark$\checkmarks\checkmark$\checkmark^\checkmark\\checkmarkc\checkmarki\checkmarkr\checkmarkc\checkmark$\checkmarku\checkmark$\checkmark^\checkmark\\checkmarkc\checkmarki\checkmarkr\checkmarkc\checkmark$\checkmarkb\checkmark$\checkmark^\checkmark\\checkmarkc\checkmarki\checkmarkr\checkmarkc\checkmark$\checkmarks\checkmark$\checkmark^\checkmark\\checkmarkc\checkmarki\checkmarkr\checkmarkc\checkmark$\checkmarku\checkmark$\checkmark^\checkmark\\checkmarkc\checkmarki\checkmarkr\checkmarkc\checkmark$\checkmarkb\checkmark$\checkmark^\checkmark\\checkmarkc\checkmarki\checkmarkr\checkmarkc\checkmark$\checkmarks\checkmark$\checkmark^\checkmark\\checkmarkc\checkmarki\checkmarkr\checkmarkc\checkmark$\checkmarke\checkmark$\checkmark^\checkmark\\checkmarkc\checkmarki\checkmarkr\checkmarkc\checkmark$\checkmarkc\checkmark$\checkmark^\checkmark\\checkmarkc\checkmarki\checkmarkr\checkmarkc\checkmark$\checkmarkt\checkmark$\checkmark^\checkmark\\checkmarkc\checkmarki\checkmarkr\checkmarkc\checkmark$\checkmarki\checkmark$\checkmark^\checkmark\\checkmarkc\checkmarki\checkmarkr\checkmarkc\checkmark$\checkmarko\checkmark$\checkmark^\checkmark\\checkmarkc\checkmarki\checkmarkr\checkmarkc\checkmark$\checkmarkn\checkmark$\checkmark^\checkmark\\checkmarkc\checkmarki\checkmarkr\checkmarkc\checkmark$\checkmark{\checkmark$\checkmark^\checkmark\\checkmarkc\checkmarki\checkmarkr\checkmarkc\checkmark$\checkmarkI\checkmark$\checkmark^\checkmark\\checkmarkc\checkmarki\checkmarkr\checkmarkc\checkmark$\checkmarkn\checkmark$\checkmark^\checkmark\\checkmarkc\checkmarki\checkmarkr\checkmarkc\checkmark$\checkmarke\checkmark$\checkmark^\checkmark\\checkmarkc\checkmarki\checkmarkr\checkmarkc\checkmark$\checkmarkr\checkmark$\checkmark^\checkmark\\checkmarkc\checkmarki\checkmarkr\checkmarkc\checkmark$\checkmarkt\checkmark$\checkmark^\checkmark\\checkmarkc\checkmarki\checkmarkr\checkmarkc\checkmark$\checkmarki\checkmark$\checkmark^\checkmark\\checkmarkc\checkmarki\checkmarkr\checkmarkc\checkmark$\checkmarke\checkmark$\checkmark^\checkmark\\checkmarkc\checkmarki\checkmarkr\checkmarkc\checkmark$\checkmark \checkmark$\checkmark^\checkmark\\checkmarkc\checkmarki\checkmarkr\checkmarkc\checkmark$\checkmarkd\checkmark$\checkmark^\checkmark\\checkmarkc\checkmarki\checkmarkr\checkmarkc\checkmark$\checkmarke\checkmark$\checkmark^\checkmark\\checkmarkc\checkmarki\checkmarkr\checkmarkc\checkmark$\checkmark \checkmark$\checkmark^\checkmark\\checkmarkc\checkmarki\checkmarkr\checkmarkc\checkmark$\checkmarkl\checkmark$\checkmark^\checkmark\\checkmarkc\checkmarki\checkmarkr\checkmarkc\checkmark$\checkmarka\checkmark$\checkmark^\checkmark\\checkmarkc\checkmarki\checkmarkr\checkmarkc\checkmark$\checkmark \checkmark$\checkmark^\checkmark\\checkmarkc\checkmarki\checkmarkr\checkmarkc\checkmark$\checkmarkv\checkmark$\checkmark^\checkmark\\checkmarkc\checkmarki\checkmarkr\checkmarkc\checkmark$\checkmarki\checkmark$\checkmark^\checkmark\\checkmarkc\checkmarki\checkmarkr\checkmarkc\checkmark$\checkmarks\checkmark$\checkmark^\checkmark\\checkmarkc\checkmarki\checkmarkr\checkmarkc\checkmark$\checkmark \checkmark$\checkmark^\checkmark\\checkmarkc\checkmarki\checkmarkr\checkmarkc\checkmark$\checkmark(\checkmark$\checkmark^\checkmark\\checkmarkc\checkmarki\checkmarkr\checkmarkc\checkmark$\checkmark$\checkmark$\checkmark^\checkmark\\checkmarkc\checkmarki\checkmarkr\checkmarkc\checkmark$\checkmarkm\checkmark$\checkmark^\checkmark\\checkmarkc\checkmarki\checkmarkr\checkmarkc\checkmark$\checkmark_\checkmark$\checkmark^\checkmark\\checkmarkc\checkmarki\checkmarkr\checkmarkc\checkmark$\checkmark{\checkmark$\checkmark^\checkmark\\checkmarkc\checkmarki\checkmarkr\checkmarkc\checkmark$\checkmarkv\checkmark$\checkmark^\checkmark\\checkmarkc\checkmarki\checkmarkr\checkmarkc\checkmark$\checkmarki\checkmark$\checkmark^\checkmark\\checkmarkc\checkmarki\checkmarkr\checkmarkc\checkmark$\checkmarks\checkmark$\checkmark^\checkmark\\checkmarkc\checkmarki\checkmarkr\checkmarkc\checkmark$\checkmark}\checkmark$\checkmark^\checkmark\\checkmarkc\checkmarki\checkmarkr\checkmarkc\checkmark$\checkmark \checkmark$\checkmark^\checkmark\\checkmarkc\checkmarki\checkmarkr\checkmarkc\checkmark$\checkmark\\checkmark$\checkmark^\checkmark\\checkmarkc\checkmarki\checkmarkr\checkmarkc\checkmark$\checkmarka\checkmark$\checkmark^\checkmark\\checkmarkc\checkmarki\checkmarkr\checkmarkc\checkmark$\checkmarkp\checkmark$\checkmark^\checkmark\\checkmarkc\checkmarki\checkmarkr\checkmarkc\checkmark$\checkmarkp\checkmark$\checkmark^\checkmark\\checkmarkc\checkmarki\checkmarkr\checkmarkc\checkmark$\checkmarkr\checkmark$\checkmark^\checkmark\\checkmarkc\checkmarki\checkmarkr\checkmarkc\checkmark$\checkmarko\checkmark$\checkmark^\checkmark\\checkmarkc\checkmarki\checkmarkr\checkmarkc\checkmark$\checkmarkx\checkmark$\checkmark^\checkmark\\checkmarkc\checkmarki\checkmarkr\checkmarkc\checkmark$\checkmark \checkmark$\checkmark^\checkmark\\checkmarkc\checkmarki\checkmarkr\checkmarkc\checkmark$\checkmark0\checkmark$\checkmark^\checkmark\\checkmarkc\checkmarki\checkmarkr\checkmarkc\checkmark$\checkmark{\checkmark$\checkmark^\checkmark\\checkmarkc\checkmarki\checkmarkr\checkmarkc\checkmark$\checkmark,\checkmark$\checkmark^\checkmark\\checkmarkc\checkmarki\checkmarkr\checkmarkc\checkmark$\checkmark}\checkmark$\checkmark^\checkmark\\checkmarkc\checkmarki\checkmarkr\checkmarkc\checkmark$\checkmark6\checkmark$\checkmark^\checkmark\\checkmarkc\checkmarki\checkmarkr\checkmarkc\checkmark$\checkmark3\checkmark$\checkmark^\checkmark\\checkmarkc\checkmarki\checkmarkr\checkmarkc\checkmark$\checkmark$\checkmark$\checkmark^\checkmark\\checkmarkc\checkmarki\checkmarkr\checkmarkc\checkmark$\checkmark \checkmark$\checkmark^\checkmark\\checkmarkc\checkmarki\checkmarkr\checkmarkc\checkmark$\checkmarkk\checkmark$\checkmark^\checkmark\\checkmarkc\checkmarki\checkmarkr\checkmarkc\checkmark$\checkmarkg\checkmark$\checkmark^\checkmark\\checkmarkc\checkmarki\checkmarkr\checkmarkc\checkmark$\checkmark)\checkmark$\checkmark^\checkmark\\checkmarkc\checkmarki\checkmarkr\checkmarkc\checkmark$\checkmark}\checkmark$\checkmark^\checkmark\\checkmarkc\checkmarki\checkmarkr\checkmarkc\checkmark$\checkmark
\checkmark$\checkmark^\checkmark\\checkmarkc\checkmarki\checkmarkr\checkmarkc\checkmark$\checkmark\\checkmark$\checkmark^\checkmark\\checkmarkc\checkmarki\checkmarkr\checkmarkc\checkmark$\checkmarkb\checkmark$\checkmark^\checkmark\\checkmarkc\checkmarki\checkmarkr\checkmarkc\checkmark$\checkmarke\checkmark$\checkmark^\checkmark\\checkmarkc\checkmarki\checkmarkr\checkmarkc\checkmark$\checkmarkg\checkmark$\checkmark^\checkmark\\checkmarkc\checkmarki\checkmarkr\checkmarkc\checkmark$\checkmarki\checkmark$\checkmark^\checkmark\\checkmarkc\checkmarki\checkmarkr\checkmarkc\checkmark$\checkmarkn\checkmark$\checkmark^\checkmark\\checkmarkc\checkmarki\checkmarkr\checkmarkc\checkmark$\checkmark{\checkmark$\checkmark^\checkmark\\checkmarkc\checkmarki\checkmarkr\checkmarkc\checkmark$\checkmarke\checkmark$\checkmark^\checkmark\\checkmarkc\checkmarki\checkmarkr\checkmarkc\checkmark$\checkmarkq\checkmark$\checkmark^\checkmark\\checkmarkc\checkmarki\checkmarkr\checkmarkc\checkmark$\checkmarku\checkmark$\checkmark^\checkmark\\checkmarkc\checkmarki\checkmarkr\checkmarkc\checkmark$\checkmarka\checkmark$\checkmark^\checkmark\\checkmarkc\checkmarki\checkmarkr\checkmarkc\checkmark$\checkmarkt\checkmark$\checkmark^\checkmark\\checkmarkc\checkmarki\checkmarkr\checkmarkc\checkmark$\checkmarki\checkmark$\checkmark^\checkmark\\checkmarkc\checkmarki\checkmarkr\checkmarkc\checkmark$\checkmarko\checkmark$\checkmark^\checkmark\\checkmarkc\checkmarki\checkmarkr\checkmarkc\checkmark$\checkmarkn\checkmark$\checkmark^\checkmark\\checkmarkc\checkmarki\checkmarkr\checkmarkc\checkmark$\checkmark}\checkmark$\checkmark^\checkmark\\checkmarkc\checkmarki\checkmarkr\checkmarkc\checkmark$\checkmark
\checkmark$\checkmark^\checkmark\\checkmarkc\checkmarki\checkmarkr\checkmarkc\checkmark$\checkmark \checkmark$\checkmark^\checkmark\\checkmarkc\checkmarki\checkmarkr\checkmarkc\checkmark$\checkmark \checkmark$\checkmark^\checkmark\\checkmarkc\checkmarki\checkmarkr\checkmarkc\checkmark$\checkmark \checkmark$\checkmark^\checkmark\\checkmarkc\checkmarki\checkmarkr\checkmarkc\checkmark$\checkmark \checkmark$\checkmark^\checkmark\\checkmarkc\checkmarki\checkmarkr\checkmarkc\checkmark$\checkmarkJ\checkmark$\checkmark^\checkmark\\checkmarkc\checkmarki\checkmarkr\checkmarkc\checkmark$\checkmark_\checkmark$\checkmark^\checkmark\\checkmarkc\checkmarki\checkmarkr\checkmarkc\checkmark$\checkmark{\checkmark$\checkmark^\checkmark\\checkmarkc\checkmarki\checkmarkr\checkmarkc\checkmark$\checkmarkv\checkmark$\checkmark^\checkmark\\checkmarkc\checkmarki\checkmarkr\checkmarkc\checkmark$\checkmarki\checkmark$\checkmark^\checkmark\\checkmarkc\checkmarki\checkmarkr\checkmarkc\checkmark$\checkmarks\checkmark$\checkmark^\checkmark\\checkmarkc\checkmarki\checkmarkr\checkmarkc\checkmark$\checkmark}\checkmark$\checkmark^\checkmark\\checkmarkc\checkmarki\checkmarkr\checkmarkc\checkmark$\checkmark \checkmark$\checkmark^\checkmark\\checkmarkc\checkmarki\checkmarkr\checkmarkc\checkmark$\checkmark=\checkmark$\checkmark^\checkmark\\checkmarkc\checkmarki\checkmarkr\checkmarkc\checkmark$\checkmark \checkmark$\checkmark^\checkmark\\checkmarkc\checkmarki\checkmarkr\checkmarkc\checkmark$\checkmark\\checkmark$\checkmark^\checkmark\\checkmarkc\checkmarki\checkmarkr\checkmarkc\checkmark$\checkmarkf\checkmark$\checkmark^\checkmark\\checkmarkc\checkmarki\checkmarkr\checkmarkc\checkmark$\checkmarkr\checkmark$\checkmark^\checkmark\\checkmarkc\checkmarki\checkmarkr\checkmarkc\checkmark$\checkmarka\checkmark$\checkmark^\checkmark\\checkmarkc\checkmarki\checkmarkr\checkmarkc\checkmark$\checkmarkc\checkmark$\checkmark^\checkmark\\checkmarkc\checkmarki\checkmarkr\checkmarkc\checkmark$\checkmark{\checkmark$\checkmark^\checkmark\\checkmarkc\checkmarki\checkmarkr\checkmarkc\checkmark$\checkmark1\checkmark$\checkmark^\checkmark\\checkmarkc\checkmarki\checkmarkr\checkmarkc\checkmark$\checkmark}\checkmark$\checkmark^\checkmark\\checkmarkc\checkmarki\checkmarkr\checkmarkc\checkmark$\checkmark{\checkmark$\checkmark^\checkmark\\checkmarkc\checkmarki\checkmarkr\checkmarkc\checkmark$\checkmark2\checkmark$\checkmark^\checkmark\\checkmarkc\checkmarki\checkmarkr\checkmarkc\checkmark$\checkmark}\checkmark$\checkmark^\checkmark\\checkmarkc\checkmarki\checkmarkr\checkmarkc\checkmark$\checkmark \checkmark$\checkmark^\checkmark\\checkmarkc\checkmarki\checkmarkr\checkmarkc\checkmark$\checkmarkm\checkmark$\checkmark^\checkmark\\checkmarkc\checkmarki\checkmarkr\checkmarkc\checkmark$\checkmark_\checkmark$\checkmark^\checkmark\\checkmarkc\checkmarki\checkmarkr\checkmarkc\checkmark$\checkmark{\checkmark$\checkmark^\checkmark\\checkmarkc\checkmarki\checkmarkr\checkmarkc\checkmark$\checkmarkv\checkmark$\checkmark^\checkmark\\checkmarkc\checkmarki\checkmarkr\checkmarkc\checkmark$\checkmarki\checkmark$\checkmark^\checkmark\\checkmarkc\checkmarki\checkmarkr\checkmarkc\checkmark$\checkmarks\checkmark$\checkmark^\checkmark\\checkmarkc\checkmarki\checkmarkr\checkmarkc\checkmark$\checkmark}\checkmark$\checkmark^\checkmark\\checkmarkc\checkmarki\checkmarkr\checkmarkc\checkmark$\checkmark \checkmark$\checkmark^\checkmark\\checkmarkc\checkmarki\checkmarkr\checkmarkc\checkmark$\checkmark\\checkmark$\checkmark^\checkmark\\checkmarkc\checkmarki\checkmarkr\checkmarkc\checkmark$\checkmarkc\checkmark$\checkmark^\checkmark\\checkmarkc\checkmarki\checkmarkr\checkmarkc\checkmark$\checkmarkd\checkmark$\checkmark^\checkmark\\checkmarkc\checkmarki\checkmarkr\checkmarkc\checkmark$\checkmarko\checkmark$\checkmark^\checkmark\\checkmarkc\checkmarki\checkmarkr\checkmarkc\checkmark$\checkmarkt\checkmark$\checkmark^\checkmark\\checkmarkc\checkmarki\checkmarkr\checkmarkc\checkmark$\checkmark \checkmark$\checkmark^\checkmark\\checkmarkc\checkmarki\checkmarkr\checkmarkc\checkmark$\checkmarkr\checkmark$\checkmark^\checkmark\\checkmarkc\checkmarki\checkmarkr\checkmarkc\checkmark$\checkmark_\checkmark$\checkmark^\checkmark\\checkmarkc\checkmarki\checkmarkr\checkmarkc\checkmark$\checkmark{\checkmark$\checkmark^\checkmark\\checkmarkc\checkmarki\checkmarkr\checkmarkc\checkmark$\checkmarkv\checkmark$\checkmark^\checkmark\\checkmarkc\checkmarki\checkmarkr\checkmarkc\checkmark$\checkmarki\checkmark$\checkmark^\checkmark\\checkmarkc\checkmarki\checkmarkr\checkmarkc\checkmark$\checkmarks\checkmark$\checkmark^\checkmark\\checkmarkc\checkmarki\checkmarkr\checkmarkc\checkmark$\checkmark}\checkmark$\checkmark^\checkmark\\checkmarkc\checkmarki\checkmarkr\checkmarkc\checkmark$\checkmark^\checkmark$\checkmark^\checkmark\\checkmarkc\checkmarki\checkmarkr\checkmarkc\checkmark$\checkmark2\checkmark$\checkmark^\checkmark\\checkmarkc\checkmarki\checkmarkr\checkmarkc\checkmark$\checkmark \checkmark$\checkmark^\checkmark\\checkmarkc\checkmarki\checkmarkr\checkmarkc\checkmark$\checkmark=\checkmark$\checkmark^\checkmark\\checkmarkc\checkmarki\checkmarkr\checkmarkc\checkmark$\checkmark \checkmark$\checkmark^\checkmark\\checkmarkc\checkmarki\checkmarkr\checkmarkc\checkmark$\checkmark\\checkmark$\checkmark^\checkmark\\checkmarkc\checkmarki\checkmarkr\checkmarkc\checkmark$\checkmarkf\checkmark$\checkmark^\checkmark\\checkmarkc\checkmarki\checkmarkr\checkmarkc\checkmark$\checkmarkr\checkmark$\checkmark^\checkmark\\checkmarkc\checkmarki\checkmarkr\checkmarkc\checkmark$\checkmarka\checkmark$\checkmark^\checkmark\\checkmarkc\checkmarki\checkmarkr\checkmarkc\checkmark$\checkmarkc\checkmark$\checkmark^\checkmark\\checkmarkc\checkmarki\checkmarkr\checkmarkc\checkmark$\checkmark{\checkmark$\checkmark^\checkmark\\checkmarkc\checkmarki\checkmarkr\checkmarkc\checkmark$\checkmark1\checkmark$\checkmark^\checkmark\\checkmarkc\checkmarki\checkmarkr\checkmarkc\checkmark$\checkmark}\checkmark$\checkmark^\checkmark\\checkmarkc\checkmarki\checkmarkr\checkmarkc\checkmark$\checkmark{\checkmark$\checkmark^\checkmark\\checkmarkc\checkmarki\checkmarkr\checkmarkc\checkmark$\checkmark2\checkmark$\checkmark^\checkmark\\checkmarkc\checkmarki\checkmarkr\checkmarkc\checkmark$\checkmark}\checkmark$\checkmark^\checkmark\\checkmarkc\checkmarki\checkmarkr\checkmarkc\checkmark$\checkmark \checkmark$\checkmark^\checkmark\\checkmarkc\checkmarki\checkmarkr\checkmarkc\checkmark$\checkmark\\checkmark$\checkmark^\checkmark\\checkmarkc\checkmarki\checkmarkr\checkmarkc\checkmark$\checkmarkt\checkmark$\checkmark^\checkmark\\checkmarkc\checkmarki\checkmarkr\checkmarkc\checkmark$\checkmarki\checkmark$\checkmark^\checkmark\\checkmarkc\checkmarki\checkmarkr\checkmarkc\checkmark$\checkmarkm\checkmark$\checkmark^\checkmark\\checkmarkc\checkmarki\checkmarkr\checkmarkc\checkmark$\checkmarke\checkmark$\checkmark^\checkmark\\checkmarkc\checkmarki\checkmarkr\checkmarkc\checkmark$\checkmarks\checkmark$\checkmark^\checkmark\\checkmarkc\checkmarki\checkmarkr\checkmarkc\checkmark$\checkmark \checkmark$\checkmark^\checkmark\\checkmarkc\checkmarki\checkmarkr\checkmarkc\checkmark$\checkmark0\checkmark$\checkmark^\checkmark\\checkmarkc\checkmarki\checkmarkr\checkmarkc\checkmark$\checkmark{\checkmark$\checkmark^\checkmark\\checkmarkc\checkmarki\checkmarkr\checkmarkc\checkmark$\checkmark,\checkmark$\checkmark^\checkmark\\checkmarkc\checkmarki\checkmarkr\checkmarkc\checkmark$\checkmark}\checkmark$\checkmark^\checkmark\\checkmarkc\checkmarki\checkmarkr\checkmarkc\checkmark$\checkmark6\checkmark$\checkmark^\checkmark\\checkmarkc\checkmarki\checkmarkr\checkmarkc\checkmark$\checkmark3\checkmark$\checkmark^\checkmark\\checkmarkc\checkmarki\checkmarkr\checkmarkc\checkmark$\checkmark \checkmark$\checkmark^\checkmark\\checkmarkc\checkmarki\checkmarkr\checkmarkc\checkmark$\checkmark\\checkmark$\checkmark^\checkmark\\checkmarkc\checkmarki\checkmarkr\checkmarkc\checkmark$\checkmarkt\checkmark$\checkmark^\checkmark\\checkmarkc\checkmarki\checkmarkr\checkmarkc\checkmark$\checkmarki\checkmark$\checkmark^\checkmark\\checkmarkc\checkmarki\checkmarkr\checkmarkc\checkmark$\checkmarkm\checkmark$\checkmark^\checkmark\\checkmarkc\checkmarki\checkmarkr\checkmarkc\checkmark$\checkmarke\checkmark$\checkmark^\checkmark\\checkmarkc\checkmarki\checkmarkr\checkmarkc\checkmark$\checkmarks\checkmark$\checkmark^\checkmark\\checkmarkc\checkmarki\checkmarkr\checkmarkc\checkmark$\checkmark \checkmark$\checkmark^\checkmark\\checkmarkc\checkmarki\checkmarkr\checkmarkc\checkmark$\checkmark(\checkmark$\checkmark^\checkmark\\checkmarkc\checkmarki\checkmarkr\checkmarkc\checkmark$\checkmark0\checkmark$\checkmark^\checkmark\\checkmarkc\checkmarki\checkmarkr\checkmarkc\checkmark$\checkmark{\checkmark$\checkmark^\checkmark\\checkmarkc\checkmarki\checkmarkr\checkmarkc\checkmark$\checkmark,\checkmark$\checkmark^\checkmark\\checkmarkc\checkmarki\checkmarkr\checkmarkc\checkmark$\checkmark}\checkmark$\checkmark^\checkmark\\checkmarkc\checkmarki\checkmarkr\checkmarkc\checkmark$\checkmark0\checkmark$\checkmark^\checkmark\\checkmarkc\checkmarki\checkmarkr\checkmarkc\checkmark$\checkmark0\checkmark$\checkmark^\checkmark\\checkmarkc\checkmarki\checkmarkr\checkmarkc\checkmark$\checkmark8\checkmark$\checkmark^\checkmark\\checkmarkc\checkmarki\checkmarkr\checkmarkc\checkmark$\checkmark)\checkmark$\checkmark^\checkmark\\checkmarkc\checkmarki\checkmarkr\checkmarkc\checkmark$\checkmark^\checkmark$\checkmark^\checkmark\\checkmarkc\checkmarki\checkmarkr\checkmarkc\checkmark$\checkmark2\checkmark$\checkmark^\checkmark\\checkmarkc\checkmarki\checkmarkr\checkmarkc\checkmark$\checkmark \checkmark$\checkmark^\checkmark\\checkmarkc\checkmarki\checkmarkr\checkmarkc\checkmark$\checkmark=\checkmark$\checkmark^\checkmark\\checkmarkc\checkmarki\checkmarkr\checkmarkc\checkmark$\checkmark \checkmark$\checkmark^\checkmark\\checkmarkc\checkmarki\checkmarkr\checkmarkc\checkmark$\checkmark2\checkmark$\checkmark^\checkmark\\checkmarkc\checkmarki\checkmarkr\checkmarkc\checkmark$\checkmark{\checkmark$\checkmark^\checkmark\\checkmarkc\checkmarki\checkmarkr\checkmarkc\checkmark$\checkmark,\checkmark$\checkmark^\checkmark\\checkmarkc\checkmarki\checkmarkr\checkmarkc\checkmark$\checkmark}\checkmark$\checkmark^\checkmark\\checkmarkc\checkmarki\checkmarkr\checkmarkc\checkmark$\checkmark0\checkmark$\checkmark^\checkmark\\checkmarkc\checkmarki\checkmarkr\checkmarkc\checkmark$\checkmark1\checkmark$\checkmark^\checkmark\\checkmarkc\checkmarki\checkmarkr\checkmarkc\checkmark$\checkmark6\checkmark$\checkmark^\checkmark\\checkmarkc\checkmarki\checkmarkr\checkmarkc\checkmark$\checkmark \checkmark$\checkmark^\checkmark\\checkmarkc\checkmarki\checkmarkr\checkmarkc\checkmark$\checkmark\\checkmark$\checkmark^\checkmark\\checkmarkc\checkmarki\checkmarkr\checkmarkc\checkmark$\checkmarkt\checkmark$\checkmark^\checkmark\\checkmarkc\checkmarki\checkmarkr\checkmarkc\checkmark$\checkmarki\checkmark$\checkmark^\checkmark\\checkmarkc\checkmarki\checkmarkr\checkmarkc\checkmark$\checkmarkm\checkmark$\checkmark^\checkmark\\checkmarkc\checkmarki\checkmarkr\checkmarkc\checkmark$\checkmarke\checkmark$\checkmark^\checkmark\\checkmarkc\checkmarki\checkmarkr\checkmarkc\checkmark$\checkmarks\checkmark$\checkmark^\checkmark\\checkmarkc\checkmarki\checkmarkr\checkmarkc\checkmark$\checkmark \checkmark$\checkmark^\checkmark\\checkmarkc\checkmarki\checkmarkr\checkmarkc\checkmark$\checkmark1\checkmark$\checkmark^\checkmark\\checkmarkc\checkmarki\checkmarkr\checkmarkc\checkmark$\checkmark0\checkmark$\checkmark^\checkmark\\checkmarkc\checkmarki\checkmarkr\checkmarkc\checkmark$\checkmark^\checkmark$\checkmark^\checkmark\\checkmarkc\checkmarki\checkmarkr\checkmarkc\checkmark$\checkmark{\checkmark$\checkmark^\checkmark\\checkmarkc\checkmarki\checkmarkr\checkmarkc\checkmark$\checkmark-\checkmark$\checkmark^\checkmark\\checkmarkc\checkmarki\checkmarkr\checkmarkc\checkmark$\checkmark5\checkmark$\checkmark^\checkmark\\checkmarkc\checkmarki\checkmarkr\checkmarkc\checkmark$\checkmark}\checkmark$\checkmark^\checkmark\\checkmarkc\checkmarki\checkmarkr\checkmarkc\checkmark$\checkmark \checkmark$\checkmark^\checkmark\\checkmarkc\checkmarki\checkmarkr\checkmarkc\checkmark$\checkmark\\checkmark$\checkmark^\checkmark\\checkmarkc\checkmarki\checkmarkr\checkmarkc\checkmark$\checkmarkt\checkmark$\checkmark^\checkmark\\checkmarkc\checkmarki\checkmarkr\checkmarkc\checkmark$\checkmarke\checkmark$\checkmark^\checkmark\\checkmarkc\checkmarki\checkmarkr\checkmarkc\checkmark$\checkmarkx\checkmark$\checkmark^\checkmark\\checkmarkc\checkmarki\checkmarkr\checkmarkc\checkmark$\checkmarkt\checkmark$\checkmark^\checkmark\\checkmarkc\checkmarki\checkmarkr\checkmarkc\checkmark$\checkmark{\checkmark$\checkmark^\checkmark\\checkmarkc\checkmarki\checkmarkr\checkmarkc\checkmark$\checkmark \checkmark$\checkmark^\checkmark\\checkmarkc\checkmarki\checkmarkr\checkmarkc\checkmark$\checkmarkk\checkmark$\checkmark^\checkmark\\checkmarkc\checkmarki\checkmarkr\checkmarkc\checkmark$\checkmarkg\checkmark$\checkmark^\checkmark\\checkmarkc\checkmarki\checkmarkr\checkmarkc\checkmark$\checkmark$\checkmark$\checkmark^\checkmark\\checkmarkc\checkmarki\checkmarkr\checkmarkc\checkmark$\checkmark\\checkmark$\checkmark^\checkmark\\checkmarkc\checkmarki\checkmarkr\checkmarkc\checkmark$\checkmarkc\checkmark$\checkmark^\checkmark\\checkmarkc\checkmarki\checkmarkr\checkmarkc\checkmark$\checkmarkd\checkmark$\checkmark^\checkmark\\checkmarkc\checkmarki\checkmarkr\checkmarkc\checkmark$\checkmarko\checkmark$\checkmark^\checkmark\\checkmarkc\checkmarki\checkmarkr\checkmarkc\checkmark$\checkmarkt\checkmark$\checkmark^\checkmark\\checkmarkc\checkmarki\checkmarkr\checkmarkc\checkmark$\checkmark$\checkmark$\checkmark^\checkmark\\checkmarkc\checkmarki\checkmarkr\checkmarkc\checkmark$\checkmarkm\checkmark$\checkmark^\checkmark\\checkmarkc\checkmarki\checkmarkr\checkmarkc\checkmark$\checkmark}\checkmark$\checkmark^\checkmark\\checkmarkc\checkmarki\checkmarkr\checkmarkc\checkmark$\checkmark^\checkmark$\checkmark^\checkmark\\checkmarkc\checkmarki\checkmarkr\checkmarkc\checkmark$\checkmark2\checkmark$\checkmark^\checkmark\\checkmarkc\checkmarki\checkmarkr\checkmarkc\checkmark$\checkmark
\checkmark$\checkmark^\checkmark\\checkmarkc\checkmarki\checkmarkr\checkmarkc\checkmark$\checkmark\\checkmark$\checkmark^\checkmark\\checkmarkc\checkmarki\checkmarkr\checkmarkc\checkmark$\checkmarke\checkmark$\checkmark^\checkmark\\checkmarkc\checkmarki\checkmarkr\checkmarkc\checkmark$\checkmarkn\checkmark$\checkmark^\checkmark\\checkmarkc\checkmarki\checkmarkr\checkmarkc\checkmark$\checkmarkd\checkmark$\checkmark^\checkmark\\checkmarkc\checkmarki\checkmarkr\checkmarkc\checkmark$\checkmark{\checkmark$\checkmark^\checkmark\\checkmarkc\checkmarki\checkmarkr\checkmarkc\checkmark$\checkmarke\checkmark$\checkmark^\checkmark\\checkmarkc\checkmarki\checkmarkr\checkmarkc\checkmark$\checkmarkq\checkmark$\checkmark^\checkmark\\checkmarkc\checkmarki\checkmarkr\checkmarkc\checkmark$\checkmarku\checkmark$\checkmark^\checkmark\\checkmarkc\checkmarki\checkmarkr\checkmarkc\checkmark$\checkmarka\checkmark$\checkmark^\checkmark\\checkmarkc\checkmarki\checkmarkr\checkmarkc\checkmark$\checkmarkt\checkmark$\checkmark^\checkmark\\checkmarkc\checkmarki\checkmarkr\checkmarkc\checkmark$\checkmarki\checkmark$\checkmark^\checkmark\\checkmarkc\checkmarki\checkmarkr\checkmarkc\checkmark$\checkmarko\checkmark$\checkmark^\checkmark\\checkmarkc\checkmarki\checkmarkr\checkmarkc\checkmark$\checkmarkn\checkmark$\checkmark^\checkmark\\checkmarkc\checkmarki\checkmarkr\checkmarkc\checkmark$\checkmark}\checkmark$\checkmark^\checkmark\\checkmarkc\checkmarki\checkmarkr\checkmarkc\checkmark$\checkmark
\checkmark$\checkmark^\checkmark\\checkmarkc\checkmarki\checkmarkr\checkmarkc\checkmark$\checkmark
\checkmark$\checkmark^\checkmark\\checkmarkc\checkmarki\checkmarkr\checkmarkc\checkmark$\checkmark\\checkmark$\checkmark^\checkmark\\checkmarkc\checkmarki\checkmarkr\checkmarkc\checkmark$\checkmarks\checkmark$\checkmark^\checkmark\\checkmarkc\checkmarki\checkmarkr\checkmarkc\checkmark$\checkmarku\checkmark$\checkmark^\checkmark\\checkmarkc\checkmarki\checkmarkr\checkmarkc\checkmark$\checkmarkb\checkmark$\checkmark^\checkmark\\checkmarkc\checkmarki\checkmarkr\checkmarkc\checkmark$\checkmarks\checkmark$\checkmark^\checkmark\\checkmarkc\checkmarki\checkmarkr\checkmarkc\checkmark$\checkmarku\checkmark$\checkmark^\checkmark\\checkmarkc\checkmarki\checkmarkr\checkmarkc\checkmark$\checkmarkb\checkmark$\checkmark^\checkmark\\checkmarkc\checkmarki\checkmarkr\checkmarkc\checkmark$\checkmarks\checkmark$\checkmark^\checkmark\\checkmarkc\checkmarki\checkmarkr\checkmarkc\checkmark$\checkmarke\checkmark$\checkmark^\checkmark\\checkmarkc\checkmarki\checkmarkr\checkmarkc\checkmark$\checkmarkc\checkmark$\checkmark^\checkmark\\checkmarkc\checkmarki\checkmarkr\checkmarkc\checkmark$\checkmarkt\checkmark$\checkmark^\checkmark\\checkmarkc\checkmarki\checkmarkr\checkmarkc\checkmark$\checkmarki\checkmark$\checkmark^\checkmark\\checkmarkc\checkmarki\checkmarkr\checkmarkc\checkmark$\checkmarko\checkmark$\checkmark^\checkmark\\checkmarkc\checkmarki\checkmarkr\checkmarkc\checkmark$\checkmarkn\checkmark$\checkmark^\checkmark\\checkmarkc\checkmarki\checkmarkr\checkmarkc\checkmark$\checkmark{\checkmark$\checkmark^\checkmark\\checkmarkc\checkmarki\checkmarkr\checkmarkc\checkmark$\checkmarkI\checkmark$\checkmark^\checkmark\\checkmarkc\checkmarki\checkmarkr\checkmarkc\checkmark$\checkmarkn\checkmark$\checkmark^\checkmark\\checkmarkc\checkmarki\checkmarkr\checkmarkc\checkmark$\checkmarke\checkmark$\checkmark^\checkmark\\checkmarkc\checkmarki\checkmarkr\checkmarkc\checkmark$\checkmarkr\checkmark$\checkmark^\checkmark\\checkmarkc\checkmarki\checkmarkr\checkmarkc\checkmark$\checkmarkt\checkmark$\checkmark^\checkmark\\checkmarkc\checkmarki\checkmarkr\checkmarkc\checkmark$\checkmarki\checkmark$\checkmark^\checkmark\\checkmarkc\checkmarki\checkmarkr\checkmarkc\checkmark$\checkmarke\checkmark$\checkmark^\checkmark\\checkmarkc\checkmarki\checkmarkr\checkmarkc\checkmark$\checkmark \checkmark$\checkmark^\checkmark\\checkmarkc\checkmarki\checkmarkr\checkmarkc\checkmark$\checkmarkd\checkmark$\checkmark^\checkmark\\checkmarkc\checkmarki\checkmarkr\checkmarkc\checkmark$\checkmarke\checkmark$\checkmark^\checkmark\\checkmarkc\checkmarki\checkmarkr\checkmarkc\checkmark$\checkmark \checkmark$\checkmark^\checkmark\\checkmarkc\checkmarki\checkmarkr\checkmarkc\checkmark$\checkmarkl\checkmark$\checkmark^\checkmark\\checkmarkc\checkmarki\checkmarkr\checkmarkc\checkmark$\checkmarka\checkmark$\checkmark^\checkmark\\checkmarkc\checkmarki\checkmarkr\checkmarkc\checkmark$\checkmark \checkmark$\checkmark^\checkmark\\checkmarkc\checkmarki\checkmarkr\checkmarkc\checkmark$\checkmarkm\checkmark$\checkmark^\checkmark\\checkmarkc\checkmarki\checkmarkr\checkmarkc\checkmark$\checkmarka\checkmark$\checkmark^\checkmark\\checkmarkc\checkmarki\checkmarkr\checkmarkc\checkmark$\checkmarks\checkmark$\checkmark^\checkmark\\checkmarkc\checkmarki\checkmarkr\checkmarkc\checkmark$\checkmarks\checkmark$\checkmark^\checkmark\\checkmarkc\checkmarki\checkmarkr\checkmarkc\checkmark$\checkmarke\checkmark$\checkmark^\checkmark\\checkmarkc\checkmarki\checkmarkr\checkmarkc\checkmark$\checkmark \checkmark$\checkmark^\checkmark\\checkmarkc\checkmarki\checkmarkr\checkmarkc\checkmark$\checkmarkm\checkmark$\checkmark^\checkmark\\checkmarkc\checkmarki\checkmarkr\checkmarkc\checkmark$\checkmarko\checkmark$\checkmark^\checkmark\\checkmarkc\checkmarki\checkmarkr\checkmarkc\checkmark$\checkmarkb\checkmark$\checkmark^\checkmark\\checkmarkc\checkmarki\checkmarkr\checkmarkc\checkmark$\checkmarki\checkmark$\checkmark^\checkmark\\checkmarkc\checkmarki\checkmarkr\checkmarkc\checkmark$\checkmarkl\checkmark$\checkmark^\checkmark\\checkmarkc\checkmarki\checkmarkr\checkmarkc\checkmark$\checkmarke\checkmark$\checkmark^\checkmark\\checkmarkc\checkmarki\checkmarkr\checkmarkc\checkmark$\checkmark \checkmark$\checkmark^\checkmark\\checkmarkc\checkmarki\checkmarkr\checkmarkc\checkmark$\checkmarke\checkmark$\checkmark^\checkmark\\checkmarkc\checkmarki\checkmarkr\checkmarkc\checkmark$\checkmarkn\checkmark$\checkmark^\checkmark\\checkmarkc\checkmarki\checkmarkr\checkmarkc\checkmark$\checkmark \checkmark$\checkmark^\checkmark\\checkmarkc\checkmarki\checkmarkr\checkmarkc\checkmark$\checkmarkt\checkmark$\checkmark^\checkmark\\checkmarkc\checkmarki\checkmarkr\checkmarkc\checkmark$\checkmarkr\checkmark$\checkmark^\checkmark\\checkmarkc\checkmarki\checkmarkr\checkmarkc\checkmark$\checkmarka\checkmark$\checkmark^\checkmark\\checkmarkc\checkmarki\checkmarkr\checkmarkc\checkmark$\checkmarkn\checkmark$\checkmark^\checkmark\\checkmarkc\checkmarki\checkmarkr\checkmarkc\checkmark$\checkmarks\checkmark$\checkmark^\checkmark\\checkmarkc\checkmarki\checkmarkr\checkmarkc\checkmark$\checkmarkl\checkmark$\checkmark^\checkmark\\checkmarkc\checkmarki\checkmarkr\checkmarkc\checkmark$\checkmarka\checkmark$\checkmark^\checkmark\\checkmarkc\checkmarki\checkmarkr\checkmarkc\checkmark$\checkmarkt\checkmark$\checkmark^\checkmark\\checkmarkc\checkmarki\checkmarkr\checkmarkc\checkmark$\checkmarki\checkmark$\checkmark^\checkmark\\checkmarkc\checkmarki\checkmarkr\checkmarkc\checkmark$\checkmarko\checkmark$\checkmark^\checkmark\\checkmarkc\checkmarki\checkmarkr\checkmarkc\checkmark$\checkmarkn\checkmark$\checkmark^\checkmark\\checkmarkc\checkmarki\checkmarkr\checkmarkc\checkmark$\checkmark}\checkmark$\checkmark^\checkmark\\checkmarkc\checkmarki\checkmarkr\checkmarkc\checkmark$\checkmark
\checkmark$\checkmark^\checkmark\\checkmarkc\checkmarki\checkmarkr\checkmarkc\checkmark$\checkmark\\checkmark$\checkmark^\checkmark\\checkmarkc\checkmarki\checkmarkr\checkmarkc\checkmark$\checkmarkb\checkmark$\checkmark^\checkmark\\checkmarkc\checkmarki\checkmarkr\checkmarkc\checkmark$\checkmarke\checkmark$\checkmark^\checkmark\\checkmarkc\checkmarki\checkmarkr\checkmarkc\checkmark$\checkmarkg\checkmark$\checkmark^\checkmark\\checkmarkc\checkmarki\checkmarkr\checkmarkc\checkmark$\checkmarki\checkmark$\checkmark^\checkmark\\checkmarkc\checkmarki\checkmarkr\checkmarkc\checkmark$\checkmarkn\checkmark$\checkmark^\checkmark\\checkmarkc\checkmarki\checkmarkr\checkmarkc\checkmark$\checkmark{\checkmark$\checkmark^\checkmark\\checkmarkc\checkmarki\checkmarkr\checkmarkc\checkmark$\checkmarke\checkmark$\checkmark^\checkmark\\checkmarkc\checkmarki\checkmarkr\checkmarkc\checkmark$\checkmarkq\checkmark$\checkmark^\checkmark\\checkmarkc\checkmarki\checkmarkr\checkmarkc\checkmark$\checkmarku\checkmark$\checkmark^\checkmark\\checkmarkc\checkmarki\checkmarkr\checkmarkc\checkmark$\checkmarka\checkmark$\checkmark^\checkmark\\checkmarkc\checkmarki\checkmarkr\checkmarkc\checkmark$\checkmarkt\checkmark$\checkmark^\checkmark\\checkmarkc\checkmarki\checkmarkr\checkmarkc\checkmark$\checkmarki\checkmark$\checkmark^\checkmark\\checkmarkc\checkmarki\checkmarkr\checkmarkc\checkmark$\checkmarko\checkmark$\checkmark^\checkmark\\checkmarkc\checkmarki\checkmarkr\checkmarkc\checkmark$\checkmarkn\checkmark$\checkmark^\checkmark\\checkmarkc\checkmarki\checkmarkr\checkmarkc\checkmark$\checkmark}\checkmark$\checkmark^\checkmark\\checkmarkc\checkmarki\checkmarkr\checkmarkc\checkmark$\checkmark
\checkmark$\checkmark^\checkmark\\checkmarkc\checkmarki\checkmarkr\checkmarkc\checkmark$\checkmark \checkmark$\checkmark^\checkmark\\checkmarkc\checkmarki\checkmarkr\checkmarkc\checkmark$\checkmark \checkmark$\checkmark^\checkmark\\checkmarkc\checkmarki\checkmarkr\checkmarkc\checkmark$\checkmark \checkmark$\checkmark^\checkmark\\checkmarkc\checkmarki\checkmarkr\checkmarkc\checkmark$\checkmark \checkmark$\checkmark^\checkmark\\checkmarkc\checkmarki\checkmarkr\checkmarkc\checkmark$\checkmarkJ\checkmark$\checkmark^\checkmark\\checkmarkc\checkmarki\checkmarkr\checkmarkc\checkmark$\checkmark_\checkmark$\checkmark^\checkmark\\checkmarkc\checkmarki\checkmarkr\checkmarkc\checkmark$\checkmark{\checkmark$\checkmark^\checkmark\\checkmarkc\checkmarki\checkmarkr\checkmarkc\checkmark$\checkmarkt\checkmark$\checkmark^\checkmark\\checkmarkc\checkmarki\checkmarkr\checkmarkc\checkmark$\checkmarkr\checkmark$\checkmark^\checkmark\\checkmarkc\checkmarki\checkmarkr\checkmarkc\checkmark$\checkmarka\checkmark$\checkmark^\checkmark\\checkmarkc\checkmarki\checkmarkr\checkmarkc\checkmark$\checkmarkn\checkmark$\checkmark^\checkmark\\checkmarkc\checkmarki\checkmarkr\checkmarkc\checkmark$\checkmarks\checkmark$\checkmark^\checkmark\\checkmarkc\checkmarki\checkmarkr\checkmarkc\checkmark$\checkmark}\checkmark$\checkmark^\checkmark\\checkmarkc\checkmarki\checkmarkr\checkmarkc\checkmark$\checkmark \checkmark$\checkmark^\checkmark\\checkmarkc\checkmarki\checkmarkr\checkmarkc\checkmark$\checkmark=\checkmark$\checkmark^\checkmark\\checkmarkc\checkmarki\checkmarkr\checkmarkc\checkmark$\checkmark \checkmark$\checkmark^\checkmark\\checkmarkc\checkmarki\checkmarkr\checkmarkc\checkmark$\checkmarkM\checkmark$\checkmark^\checkmark\\checkmarkc\checkmarki\checkmarkr\checkmarkc\checkmark$\checkmark \checkmark$\checkmark^\checkmark\\checkmarkc\checkmarki\checkmarkr\checkmarkc\checkmark$\checkmark\\checkmark$\checkmark^\checkmark\\checkmarkc\checkmarki\checkmarkr\checkmarkc\checkmark$\checkmarkc\checkmark$\checkmark^\checkmark\\checkmarkc\checkmarki\checkmarkr\checkmarkc\checkmark$\checkmarkd\checkmark$\checkmark^\checkmark\\checkmarkc\checkmarki\checkmarkr\checkmarkc\checkmark$\checkmarko\checkmark$\checkmark^\checkmark\\checkmarkc\checkmarki\checkmarkr\checkmarkc\checkmark$\checkmarkt\checkmark$\checkmark^\checkmark\\checkmarkc\checkmarki\checkmarkr\checkmarkc\checkmark$\checkmark \checkmark$\checkmark^\checkmark\\checkmarkc\checkmarki\checkmarkr\checkmarkc\checkmark$\checkmark\\checkmark$\checkmark^\checkmark\\checkmarkc\checkmarki\checkmarkr\checkmarkc\checkmark$\checkmarkl\checkmark$\checkmark^\checkmark\\checkmarkc\checkmarki\checkmarkr\checkmarkc\checkmark$\checkmarke\checkmark$\checkmark^\checkmark\\checkmarkc\checkmarki\checkmarkr\checkmarkc\checkmark$\checkmarkf\checkmark$\checkmark^\checkmark\\checkmarkc\checkmarki\checkmarkr\checkmarkc\checkmark$\checkmarkt\checkmark$\checkmark^\checkmark\\checkmarkc\checkmarki\checkmarkr\checkmarkc\checkmark$\checkmark(\checkmark$\checkmark^\checkmark\\checkmarkc\checkmarki\checkmarkr\checkmarkc\checkmark$\checkmark\\checkmark$\checkmark^\checkmark\\checkmarkc\checkmarki\checkmarkr\checkmarkc\checkmark$\checkmarkf\checkmark$\checkmark^\checkmark\\checkmarkc\checkmarki\checkmarkr\checkmarkc\checkmark$\checkmarkr\checkmark$\checkmark^\checkmark\\checkmarkc\checkmarki\checkmarkr\checkmarkc\checkmark$\checkmarka\checkmark$\checkmark^\checkmark\\checkmarkc\checkmarki\checkmarkr\checkmarkc\checkmark$\checkmarkc\checkmark$\checkmark^\checkmark\\checkmarkc\checkmarki\checkmarkr\checkmarkc\checkmark$\checkmark{\checkmark$\checkmark^\checkmark\\checkmarkc\checkmarki\checkmarkr\checkmarkc\checkmark$\checkmarkP\checkmark$\checkmark^\checkmark\\checkmarkc\checkmarki\checkmarkr\checkmarkc\checkmark$\checkmark_\checkmark$\checkmark^\checkmark\\checkmarkc\checkmarki\checkmarkr\checkmarkc\checkmark$\checkmarkh\checkmark$\checkmark^\checkmark\\checkmarkc\checkmarki\checkmarkr\checkmarkc\checkmark$\checkmark}\checkmark$\checkmark^\checkmark\\checkmarkc\checkmarki\checkmarkr\checkmarkc\checkmark$\checkmark{\checkmark$\checkmark^\checkmark\\checkmarkc\checkmarki\checkmarkr\checkmarkc\checkmark$\checkmark2\checkmark$\checkmark^\checkmark\\checkmarkc\checkmarki\checkmarkr\checkmarkc\checkmark$\checkmark\\checkmark$\checkmark^\checkmark\\checkmarkc\checkmarki\checkmarkr\checkmarkc\checkmark$\checkmarkp\checkmark$\checkmark^\checkmark\\checkmarkc\checkmarki\checkmarkr\checkmarkc\checkmark$\checkmarki\checkmark$\checkmark^\checkmark\\checkmarkc\checkmarki\checkmarkr\checkmarkc\checkmark$\checkmark}\checkmark$\checkmark^\checkmark\\checkmarkc\checkmarki\checkmarkr\checkmarkc\checkmark$\checkmark\\checkmark$\checkmark^\checkmark\\checkmarkc\checkmarki\checkmarkr\checkmarkc\checkmark$\checkmarkr\checkmark$\checkmark^\checkmark\\checkmarkc\checkmarki\checkmarkr\checkmarkc\checkmark$\checkmarki\checkmark$\checkmark^\checkmark\\checkmarkc\checkmarki\checkmarkr\checkmarkc\checkmark$\checkmarkg\checkmark$\checkmark^\checkmark\\checkmarkc\checkmarki\checkmarkr\checkmarkc\checkmark$\checkmarkh\checkmark$\checkmark^\checkmark\\checkmarkc\checkmarki\checkmarkr\checkmarkc\checkmark$\checkmarkt\checkmark$\checkmark^\checkmark\\checkmarkc\checkmarki\checkmarkr\checkmarkc\checkmark$\checkmark)\checkmark$\checkmark^\checkmark\\checkmarkc\checkmarki\checkmarkr\checkmarkc\checkmark$\checkmark^\checkmark$\checkmark^\checkmark\\checkmarkc\checkmarki\checkmarkr\checkmarkc\checkmark$\checkmark2\checkmark$\checkmark^\checkmark\\checkmarkc\checkmarki\checkmarkr\checkmarkc\checkmark$\checkmark \checkmark$\checkmark^\checkmark\\checkmarkc\checkmarki\checkmarkr\checkmarkc\checkmark$\checkmark=\checkmark$\checkmark^\checkmark\\checkmarkc\checkmarki\checkmarkr\checkmarkc\checkmark$\checkmark \checkmark$\checkmark^\checkmark\\checkmarkc\checkmarki\checkmarkr\checkmarkc\checkmark$\checkmark1\checkmark$\checkmark^\checkmark\\checkmarkc\checkmarki\checkmarkr\checkmarkc\checkmark$\checkmark0\checkmark$\checkmark^\checkmark\\checkmarkc\checkmarki\checkmarkr\checkmarkc\checkmark$\checkmark \checkmark$\checkmark^\checkmark\\checkmarkc\checkmarki\checkmarkr\checkmarkc\checkmark$\checkmark\\checkmark$\checkmark^\checkmark\\checkmarkc\checkmarki\checkmarkr\checkmarkc\checkmark$\checkmarkt\checkmark$\checkmark^\checkmark\\checkmarkc\checkmarki\checkmarkr\checkmarkc\checkmark$\checkmarki\checkmark$\checkmark^\checkmark\\checkmarkc\checkmarki\checkmarkr\checkmarkc\checkmark$\checkmarkm\checkmark$\checkmark^\checkmark\\checkmarkc\checkmarki\checkmarkr\checkmarkc\checkmark$\checkmarke\checkmark$\checkmark^\checkmark\\checkmarkc\checkmarki\checkmarkr\checkmarkc\checkmark$\checkmarks\checkmark$\checkmark^\checkmark\\checkmarkc\checkmarki\checkmarkr\checkmarkc\checkmark$\checkmark \checkmark$\checkmark^\checkmark\\checkmarkc\checkmarki\checkmarkr\checkmarkc\checkmark$\checkmark\\checkmark$\checkmark^\checkmark\\checkmarkc\checkmarki\checkmarkr\checkmarkc\checkmark$\checkmarkl\checkmark$\checkmark^\checkmark\\checkmarkc\checkmarki\checkmarkr\checkmarkc\checkmark$\checkmarke\checkmark$\checkmark^\checkmark\\checkmarkc\checkmarki\checkmarkr\checkmarkc\checkmark$\checkmarkf\checkmark$\checkmark^\checkmark\\checkmarkc\checkmarki\checkmarkr\checkmarkc\checkmark$\checkmarkt\checkmark$\checkmark^\checkmark\\checkmarkc\checkmarki\checkmarkr\checkmarkc\checkmark$\checkmark(\checkmark$\checkmark^\checkmark\\checkmarkc\checkmarki\checkmarkr\checkmarkc\checkmark$\checkmark\\checkmark$\checkmark^\checkmark\\checkmarkc\checkmarki\checkmarkr\checkmarkc\checkmark$\checkmarkf\checkmark$\checkmark^\checkmark\\checkmarkc\checkmarki\checkmarkr\checkmarkc\checkmark$\checkmarkr\checkmark$\checkmark^\checkmark\\checkmarkc\checkmarki\checkmarkr\checkmarkc\checkmark$\checkmarka\checkmark$\checkmark^\checkmark\\checkmarkc\checkmarki\checkmarkr\checkmarkc\checkmark$\checkmarkc\checkmark$\checkmark^\checkmark\\checkmarkc\checkmarki\checkmarkr\checkmarkc\checkmark$\checkmark{\checkmark$\checkmark^\checkmark\\checkmarkc\checkmarki\checkmarkr\checkmarkc\checkmark$\checkmark0\checkmark$\checkmark^\checkmark\\checkmarkc\checkmarki\checkmarkr\checkmarkc\checkmark$\checkmark{\checkmark$\checkmark^\checkmark\\checkmarkc\checkmarki\checkmarkr\checkmarkc\checkmark$\checkmark,\checkmark$\checkmark^\checkmark\\checkmarkc\checkmarki\checkmarkr\checkmarkc\checkmark$\checkmark}\checkmark$\checkmark^\checkmark\\checkmarkc\checkmarki\checkmarkr\checkmarkc\checkmark$\checkmark0\checkmark$\checkmark^\checkmark\\checkmarkc\checkmarki\checkmarkr\checkmarkc\checkmark$\checkmark0\checkmark$\checkmark^\checkmark\\checkmarkc\checkmarki\checkmarkr\checkmarkc\checkmark$\checkmark5\checkmark$\checkmark^\checkmark\\checkmarkc\checkmarki\checkmarkr\checkmarkc\checkmark$\checkmark}\checkmark$\checkmark^\checkmark\\checkmarkc\checkmarki\checkmarkr\checkmarkc\checkmark$\checkmark{\checkmark$\checkmark^\checkmark\\checkmarkc\checkmarki\checkmarkr\checkmarkc\checkmark$\checkmark2\checkmark$\checkmark^\checkmark\\checkmarkc\checkmarki\checkmarkr\checkmarkc\checkmark$\checkmark\\checkmark$\checkmark^\checkmark\\checkmarkc\checkmarki\checkmarkr\checkmarkc\checkmark$\checkmarkp\checkmark$\checkmark^\checkmark\\checkmarkc\checkmarki\checkmarkr\checkmarkc\checkmark$\checkmarki\checkmark$\checkmark^\checkmark\\checkmarkc\checkmarki\checkmarkr\checkmarkc\checkmark$\checkmark}\checkmark$\checkmark^\checkmark\\checkmarkc\checkmarki\checkmarkr\checkmarkc\checkmark$\checkmark\\checkmark$\checkmark^\checkmark\\checkmarkc\checkmarki\checkmarkr\checkmarkc\checkmark$\checkmarkr\checkmark$\checkmark^\checkmark\\checkmarkc\checkmarki\checkmarkr\checkmarkc\checkmark$\checkmarki\checkmark$\checkmark^\checkmark\\checkmarkc\checkmarki\checkmarkr\checkmarkc\checkmark$\checkmarkg\checkmark$\checkmark^\checkmark\\checkmarkc\checkmarki\checkmarkr\checkmarkc\checkmark$\checkmarkh\checkmark$\checkmark^\checkmark\\checkmarkc\checkmarki\checkmarkr\checkmarkc\checkmark$\checkmarkt\checkmark$\checkmark^\checkmark\\checkmarkc\checkmarki\checkmarkr\checkmarkc\checkmark$\checkmark)\checkmark$\checkmark^\checkmark\\checkmarkc\checkmarki\checkmarkr\checkmarkc\checkmark$\checkmark^\checkmark$\checkmark^\checkmark\\checkmarkc\checkmarki\checkmarkr\checkmarkc\checkmark$\checkmark2\checkmark$\checkmark^\checkmark\\checkmarkc\checkmarki\checkmarkr\checkmarkc\checkmark$\checkmark \checkmark$\checkmark^\checkmark\\checkmarkc\checkmarki\checkmarkr\checkmarkc\checkmark$\checkmark=\checkmark$\checkmark^\checkmark\\checkmarkc\checkmarki\checkmarkr\checkmarkc\checkmark$\checkmark \checkmark$\checkmark^\checkmark\\checkmarkc\checkmarki\checkmarkr\checkmarkc\checkmark$\checkmark6\checkmark$\checkmark^\checkmark\\checkmarkc\checkmarki\checkmarkr\checkmarkc\checkmark$\checkmark{\checkmark$\checkmark^\checkmark\\checkmarkc\checkmarki\checkmarkr\checkmarkc\checkmark$\checkmark,\checkmark$\checkmark^\checkmark\\checkmarkc\checkmarki\checkmarkr\checkmarkc\checkmark$\checkmark}\checkmark$\checkmark^\checkmark\\checkmarkc\checkmarki\checkmarkr\checkmarkc\checkmark$\checkmark3\checkmark$\checkmark^\checkmark\\checkmarkc\checkmarki\checkmarkr\checkmarkc\checkmark$\checkmark3\checkmark$\checkmark^\checkmark\\checkmarkc\checkmarki\checkmarkr\checkmarkc\checkmark$\checkmark \checkmark$\checkmark^\checkmark\\checkmarkc\checkmarki\checkmarkr\checkmarkc\checkmark$\checkmark\\checkmark$\checkmark^\checkmark\\checkmarkc\checkmarki\checkmarkr\checkmarkc\checkmark$\checkmarkt\checkmark$\checkmark^\checkmark\\checkmarkc\checkmarki\checkmarkr\checkmarkc\checkmark$\checkmarki\checkmark$\checkmark^\checkmark\\checkmarkc\checkmarki\checkmarkr\checkmarkc\checkmark$\checkmarkm\checkmark$\checkmark^\checkmark\\checkmarkc\checkmarki\checkmarkr\checkmarkc\checkmark$\checkmarke\checkmark$\checkmark^\checkmark\\checkmarkc\checkmarki\checkmarkr\checkmarkc\checkmark$\checkmarks\checkmark$\checkmark^\checkmark\\checkmarkc\checkmarki\checkmarkr\checkmarkc\checkmark$\checkmark \checkmark$\checkmark^\checkmark\\checkmarkc\checkmarki\checkmarkr\checkmarkc\checkmark$\checkmark1\checkmark$\checkmark^\checkmark\\checkmarkc\checkmarki\checkmarkr\checkmarkc\checkmark$\checkmark0\checkmark$\checkmark^\checkmark\\checkmarkc\checkmarki\checkmarkr\checkmarkc\checkmark$\checkmark^\checkmark$\checkmark^\checkmark\\checkmarkc\checkmarki\checkmarkr\checkmarkc\checkmark$\checkmark{\checkmark$\checkmark^\checkmark\\checkmarkc\checkmarki\checkmarkr\checkmarkc\checkmark$\checkmark-\checkmark$\checkmark^\checkmark\\checkmarkc\checkmarki\checkmarkr\checkmarkc\checkmark$\checkmark6\checkmark$\checkmark^\checkmark\\checkmarkc\checkmarki\checkmarkr\checkmarkc\checkmark$\checkmark}\checkmark$\checkmark^\checkmark\\checkmarkc\checkmarki\checkmarkr\checkmarkc\checkmark$\checkmark \checkmark$\checkmark^\checkmark\\checkmarkc\checkmarki\checkmarkr\checkmarkc\checkmark$\checkmark\\checkmark$\checkmark^\checkmark\\checkmarkc\checkmarki\checkmarkr\checkmarkc\checkmark$\checkmarkt\checkmark$\checkmark^\checkmark\\checkmarkc\checkmarki\checkmarkr\checkmarkc\checkmark$\checkmarke\checkmark$\checkmark^\checkmark\\checkmarkc\checkmarki\checkmarkr\checkmarkc\checkmark$\checkmarkx\checkmark$\checkmark^\checkmark\\checkmarkc\checkmarki\checkmarkr\checkmarkc\checkmark$\checkmarkt\checkmark$\checkmark^\checkmark\\checkmarkc\checkmarki\checkmarkr\checkmarkc\checkmark$\checkmark{\checkmark$\checkmark^\checkmark\\checkmarkc\checkmarki\checkmarkr\checkmarkc\checkmark$\checkmark \checkmark$\checkmark^\checkmark\\checkmarkc\checkmarki\checkmarkr\checkmarkc\checkmark$\checkmarkk\checkmark$\checkmark^\checkmark\\checkmarkc\checkmarki\checkmarkr\checkmarkc\checkmark$\checkmarkg\checkmark$\checkmark^\checkmark\\checkmarkc\checkmarki\checkmarkr\checkmarkc\checkmark$\checkmark$\checkmark$\checkmark^\checkmark\\checkmarkc\checkmarki\checkmarkr\checkmarkc\checkmark$\checkmark\\checkmark$\checkmark^\checkmark\\checkmarkc\checkmarki\checkmarkr\checkmarkc\checkmark$\checkmarkc\checkmark$\checkmark^\checkmark\\checkmarkc\checkmarki\checkmarkr\checkmarkc\checkmark$\checkmarkd\checkmark$\checkmark^\checkmark\\checkmarkc\checkmarki\checkmarkr\checkmarkc\checkmark$\checkmarko\checkmark$\checkmark^\checkmark\\checkmarkc\checkmarki\checkmarkr\checkmarkc\checkmark$\checkmarkt\checkmark$\checkmark^\checkmark\\checkmarkc\checkmarki\checkmarkr\checkmarkc\checkmark$\checkmark$\checkmark$\checkmark^\checkmark\\checkmarkc\checkmarki\checkmarkr\checkmarkc\checkmark$\checkmarkm\checkmark$\checkmark^\checkmark\\checkmarkc\checkmarki\checkmarkr\checkmarkc\checkmark$\checkmark}\checkmark$\checkmark^\checkmark\\checkmarkc\checkmarki\checkmarkr\checkmarkc\checkmark$\checkmark^\checkmark$\checkmark^\checkmark\\checkmarkc\checkmarki\checkmarkr\checkmarkc\checkmark$\checkmark2\checkmark$\checkmark^\checkmark\\checkmarkc\checkmarki\checkmarkr\checkmarkc\checkmark$\checkmark
\checkmark$\checkmark^\checkmark\\checkmarkc\checkmarki\checkmarkr\checkmarkc\checkmark$\checkmark\\checkmark$\checkmark^\checkmark\\checkmarkc\checkmarki\checkmarkr\checkmarkc\checkmark$\checkmarke\checkmark$\checkmark^\checkmark\\checkmarkc\checkmarki\checkmarkr\checkmarkc\checkmark$\checkmarkn\checkmark$\checkmark^\checkmark\\checkmarkc\checkmarki\checkmarkr\checkmarkc\checkmark$\checkmarkd\checkmark$\checkmark^\checkmark\\checkmarkc\checkmarki\checkmarkr\checkmarkc\checkmark$\checkmark{\checkmark$\checkmark^\checkmark\\checkmarkc\checkmarki\checkmarkr\checkmarkc\checkmark$\checkmarke\checkmark$\checkmark^\checkmark\\checkmarkc\checkmarki\checkmarkr\checkmarkc\checkmark$\checkmarkq\checkmark$\checkmark^\checkmark\\checkmarkc\checkmarki\checkmarkr\checkmarkc\checkmark$\checkmarku\checkmark$\checkmark^\checkmark\\checkmarkc\checkmarki\checkmarkr\checkmarkc\checkmark$\checkmarka\checkmark$\checkmark^\checkmark\\checkmarkc\checkmarki\checkmarkr\checkmarkc\checkmark$\checkmarkt\checkmark$\checkmark^\checkmark\\checkmarkc\checkmarki\checkmarkr\checkmarkc\checkmark$\checkmarki\checkmark$\checkmark^\checkmark\\checkmarkc\checkmarki\checkmarkr\checkmarkc\checkmark$\checkmarko\checkmark$\checkmark^\checkmark\\checkmarkc\checkmarki\checkmarkr\checkmarkc\checkmark$\checkmarkn\checkmark$\checkmark^\checkmark\\checkmarkc\checkmarki\checkmarkr\checkmarkc\checkmark$\checkmark}\checkmark$\checkmark^\checkmark\\checkmarkc\checkmarki\checkmarkr\checkmarkc\checkmark$\checkmark
\checkmark$\checkmark^\checkmark\\checkmarkc\checkmarki\checkmarkr\checkmarkc\checkmark$\checkmark
\checkmark$\checkmark^\checkmark\\checkmarkc\checkmarki\checkmarkr\checkmarkc\checkmark$\checkmark\\checkmark$\checkmark^\checkmark\\checkmarkc\checkmarki\checkmarkr\checkmarkc\checkmark$\checkmarks\checkmark$\checkmark^\checkmark\\checkmarkc\checkmarki\checkmarkr\checkmarkc\checkmark$\checkmarku\checkmark$\checkmark^\checkmark\\checkmarkc\checkmarki\checkmarkr\checkmarkc\checkmark$\checkmarkb\checkmark$\checkmark^\checkmark\\checkmarkc\checkmarki\checkmarkr\checkmarkc\checkmark$\checkmarks\checkmark$\checkmark^\checkmark\\checkmarkc\checkmarki\checkmarkr\checkmarkc\checkmark$\checkmarku\checkmark$\checkmark^\checkmark\\checkmarkc\checkmarki\checkmarkr\checkmarkc\checkmark$\checkmarkb\checkmark$\checkmark^\checkmark\\checkmarkc\checkmarki\checkmarkr\checkmarkc\checkmark$\checkmarks\checkmark$\checkmark^\checkmark\\checkmarkc\checkmarki\checkmarkr\checkmarkc\checkmark$\checkmarke\checkmark$\checkmark^\checkmark\\checkmarkc\checkmarki\checkmarkr\checkmarkc\checkmark$\checkmarkc\checkmark$\checkmark^\checkmark\\checkmarkc\checkmarki\checkmarkr\checkmarkc\checkmark$\checkmarkt\checkmark$\checkmark^\checkmark\\checkmarkc\checkmarki\checkmarkr\checkmarkc\checkmark$\checkmarki\checkmark$\checkmark^\checkmark\\checkmarkc\checkmarki\checkmarkr\checkmarkc\checkmark$\checkmarko\checkmark$\checkmark^\checkmark\\checkmarkc\checkmarki\checkmarkr\checkmarkc\checkmark$\checkmarkn\checkmark$\checkmark^\checkmark\\checkmarkc\checkmarki\checkmarkr\checkmarkc\checkmark$\checkmark{\checkmark$\checkmark^\checkmark\\checkmarkc\checkmarki\checkmarkr\checkmarkc\checkmark$\checkmarkI\checkmark$\checkmark^\checkmark\\checkmarkc\checkmarki\checkmarkr\checkmarkc\checkmark$\checkmarkn\checkmark$\checkmark^\checkmark\\checkmarkc\checkmarki\checkmarkr\checkmarkc\checkmark$\checkmarke\checkmark$\checkmark^\checkmark\\checkmarkc\checkmarki\checkmarkr\checkmarkc\checkmark$\checkmarkr\checkmark$\checkmark^\checkmark\\checkmarkc\checkmarki\checkmarkr\checkmarkc\checkmark$\checkmarkt\checkmark$\checkmark^\checkmark\\checkmarkc\checkmarki\checkmarkr\checkmarkc\checkmark$\checkmarki\checkmark$\checkmark^\checkmark\\checkmarkc\checkmarki\checkmarkr\checkmarkc\checkmark$\checkmarke\checkmark$\checkmark^\checkmark\\checkmarkc\checkmarki\checkmarkr\checkmarkc\checkmark$\checkmark \checkmark$\checkmark^\checkmark\\checkmarkc\checkmarki\checkmarkr\checkmarkc\checkmark$\checkmarkt\checkmark$\checkmark^\checkmark\\checkmarkc\checkmarki\checkmarkr\checkmarkc\checkmark$\checkmarko\checkmark$\checkmark^\checkmark\\checkmarkc\checkmarki\checkmarkr\checkmarkc\checkmark$\checkmarkt\checkmark$\checkmark^\checkmark\\checkmarkc\checkmarki\checkmarkr\checkmarkc\checkmark$\checkmarka\checkmark$\checkmark^\checkmark\\checkmarkc\checkmarki\checkmarkr\checkmarkc\checkmark$\checkmarkl\checkmark$\checkmark^\checkmark\\checkmarkc\checkmarki\checkmarkr\checkmarkc\checkmark$\checkmarke\checkmark$\checkmark^\checkmark\\checkmarkc\checkmarki\checkmarkr\checkmarkc\checkmark$\checkmark}\checkmark$\checkmark^\checkmark\\checkmarkc\checkmarki\checkmarkr\checkmarkc\checkmark$\checkmark
\checkmark$\checkmark^\checkmark\\checkmarkc\checkmarki\checkmarkr\checkmarkc\checkmark$\checkmarkE\checkmark$\checkmark^\checkmark\\checkmarkc\checkmarki\checkmarkr\checkmarkc\checkmark$\checkmarkn\checkmark$\checkmark^\checkmark\\checkmarkc\checkmarki\checkmarkr\checkmarkc\checkmark$\checkmark \checkmark$\checkmark^\checkmark\\checkmarkc\checkmarki\checkmarkr\checkmarkc\checkmark$\checkmarke\checkmark$\checkmark^\checkmark\\checkmarkc\checkmarki\checkmarkr\checkmarkc\checkmark$\checkmarks\checkmark$\checkmark^\checkmark\\checkmarkc\checkmarki\checkmarkr\checkmarkc\checkmark$\checkmarkt\checkmark$\checkmark^\checkmark\\checkmarkc\checkmarki\checkmarkr\checkmarkc\checkmark$\checkmarki\checkmark$\checkmark^\checkmark\\checkmarkc\checkmarki\checkmarkr\checkmarkc\checkmark$\checkmarkm\checkmark$\checkmark^\checkmark\\checkmarkc\checkmarki\checkmarkr\checkmarkc\checkmark$\checkmarka\checkmark$\checkmark^\checkmark\\checkmarkc\checkmarki\checkmarkr\checkmarkc\checkmark$\checkmarkn\checkmark$\checkmark^\checkmark\\checkmarkc\checkmarki\checkmarkr\checkmarkc\checkmark$\checkmarkt\checkmark$\checkmark^\checkmark\\checkmarkc\checkmarki\checkmarkr\checkmarkc\checkmark$\checkmark \checkmark$\checkmark^\checkmark\\checkmarkc\checkmarki\checkmarkr\checkmarkc\checkmark$\checkmarkl\checkmark$\checkmark^\checkmark\\checkmarkc\checkmarki\checkmarkr\checkmarkc\checkmark$\checkmark'\checkmark$\checkmark^\checkmark\\checkmarkc\checkmarki\checkmarkr\checkmarkc\checkmark$\checkmarki\checkmark$\checkmark^\checkmark\\checkmarkc\checkmarki\checkmarkr\checkmarkc\checkmark$\checkmarkn\checkmark$\checkmark^\checkmark\\checkmarkc\checkmarki\checkmarkr\checkmarkc\checkmark$\checkmarke\checkmark$\checkmark^\checkmark\\checkmarkc\checkmarki\checkmarkr\checkmarkc\checkmark$\checkmarkr\checkmark$\checkmark^\checkmark\\checkmarkc\checkmarki\checkmarkr\checkmarkc\checkmark$\checkmarkt\checkmark$\checkmark^\checkmark\\checkmarkc\checkmarki\checkmarkr\checkmarkc\checkmark$\checkmarki\checkmark$\checkmark^\checkmark\\checkmarkc\checkmarki\checkmarkr\checkmarkc\checkmark$\checkmarke\checkmark$\checkmark^\checkmark\\checkmarkc\checkmarki\checkmarkr\checkmarkc\checkmark$\checkmark \checkmark$\checkmark^\checkmark\\checkmarkc\checkmarki\checkmarkr\checkmarkc\checkmark$\checkmarkd\checkmark$\checkmark^\checkmark\\checkmarkc\checkmarki\checkmarkr\checkmarkc\checkmark$\checkmarku\checkmark$\checkmark^\checkmark\\checkmarkc\checkmarki\checkmarkr\checkmarkc\checkmark$\checkmark \checkmark$\checkmark^\checkmark\\checkmarkc\checkmarki\checkmarkr\checkmarkc\checkmark$\checkmarkr\checkmark$\checkmark^\checkmark\\checkmarkc\checkmarki\checkmarkr\checkmarkc\checkmark$\checkmarko\checkmark$\checkmark^\checkmark\\checkmarkc\checkmarki\checkmarkr\checkmarkc\checkmark$\checkmarkt\checkmark$\checkmark^\checkmark\\checkmarkc\checkmarki\checkmarkr\checkmarkc\checkmark$\checkmarko\checkmark$\checkmark^\checkmark\\checkmarkc\checkmarki\checkmarkr\checkmarkc\checkmark$\checkmarkr\checkmark$\checkmark^\checkmark\\checkmarkc\checkmarki\checkmarkr\checkmarkc\checkmark$\checkmark \checkmark$\checkmark^\checkmark\\checkmarkc\checkmarki\checkmarkr\checkmarkc\checkmark$\checkmarkd\checkmark$\checkmark^\checkmark\\checkmarkc\checkmarki\checkmarkr\checkmarkc\checkmark$\checkmarku\checkmark$\checkmark^\checkmark\\checkmarkc\checkmarki\checkmarkr\checkmarkc\checkmark$\checkmark \checkmark$\checkmark^\checkmark\\checkmarkc\checkmarki\checkmarkr\checkmarkc\checkmark$\checkmarkm\checkmark$\checkmark^\checkmark\\checkmarkc\checkmarki\checkmarkr\checkmarkc\checkmark$\checkmarko\checkmark$\checkmark^\checkmark\\checkmarkc\checkmarki\checkmarkr\checkmarkc\checkmark$\checkmarkt\checkmark$\checkmark^\checkmark\\checkmarkc\checkmarki\checkmarkr\checkmarkc\checkmark$\checkmarke\checkmark$\checkmark^\checkmark\\checkmarkc\checkmarki\checkmarkr\checkmarkc\checkmark$\checkmarku\checkmark$\checkmark^\checkmark\\checkmarkc\checkmarki\checkmarkr\checkmarkc\checkmark$\checkmarkr\checkmark$\checkmark^\checkmark\\checkmarkc\checkmarki\checkmarkr\checkmarkc\checkmark$\checkmark \checkmark$\checkmark^\checkmark\\checkmarkc\checkmarki\checkmarkr\checkmarkc\checkmark$\checkmarkN\checkmark$\checkmark^\checkmark\\checkmarkc\checkmarki\checkmarkr\checkmarkc\checkmark$\checkmarkE\checkmark$\checkmark^\checkmark\\checkmarkc\checkmarki\checkmarkr\checkmarkc\checkmark$\checkmarkM\checkmark$\checkmark^\checkmark\\checkmarkc\checkmarki\checkmarkr\checkmarkc\checkmark$\checkmarkA\checkmark$\checkmark^\checkmark\\checkmarkc\checkmarki\checkmarkr\checkmarkc\checkmark$\checkmark \checkmark$\checkmark^\checkmark\\checkmarkc\checkmarki\checkmarkr\checkmarkc\checkmark$\checkmark1\checkmark$\checkmark^\checkmark\\checkmarkc\checkmarki\checkmarkr\checkmarkc\checkmark$\checkmark7\checkmark$\checkmark^\checkmark\\checkmarkc\checkmarki\checkmarkr\checkmarkc\checkmark$\checkmark \checkmark$\checkmark^\checkmark\\checkmarkc\checkmarki\checkmarkr\checkmarkc\checkmark$\checkmark:\checkmark$\checkmark^\checkmark\\checkmarkc\checkmarki\checkmarkr\checkmarkc\checkmark$\checkmark \checkmark$\checkmark^\checkmark\\checkmarkc\checkmarki\checkmarkr\checkmarkc\checkmark$\checkmark$\checkmark$\checkmark^\checkmark\\checkmarkc\checkmarki\checkmarkr\checkmarkc\checkmark$\checkmarkJ\checkmark$\checkmark^\checkmark\\checkmarkc\checkmarki\checkmarkr\checkmarkc\checkmark$\checkmark_\checkmark$\checkmark^\checkmark\\checkmarkc\checkmarki\checkmarkr\checkmarkc\checkmark$\checkmark{\checkmark$\checkmark^\checkmark\\checkmarkc\checkmarki\checkmarkr\checkmarkc\checkmark$\checkmarkm\checkmark$\checkmark^\checkmark\\checkmarkc\checkmarki\checkmarkr\checkmarkc\checkmark$\checkmarko\checkmark$\checkmark^\checkmark\\checkmarkc\checkmarki\checkmarkr\checkmarkc\checkmark$\checkmarkt\checkmark$\checkmark^\checkmark\\checkmarkc\checkmarki\checkmarkr\checkmarkc\checkmark$\checkmark}\checkmark$\checkmark^\checkmark\\checkmarkc\checkmarki\checkmarkr\checkmarkc\checkmark$\checkmark \checkmark$\checkmark^\checkmark\\checkmarkc\checkmarki\checkmarkr\checkmarkc\checkmark$\checkmark\\checkmark$\checkmark^\checkmark\\checkmarkc\checkmarki\checkmarkr\checkmarkc\checkmark$\checkmarka\checkmark$\checkmark^\checkmark\\checkmarkc\checkmarki\checkmarkr\checkmarkc\checkmark$\checkmarkp\checkmark$\checkmark^\checkmark\\checkmarkc\checkmarki\checkmarkr\checkmarkc\checkmark$\checkmarkp\checkmark$\checkmark^\checkmark\\checkmarkc\checkmarki\checkmarkr\checkmarkc\checkmark$\checkmarkr\checkmark$\checkmark^\checkmark\\checkmarkc\checkmarki\checkmarkr\checkmarkc\checkmark$\checkmarko\checkmark$\checkmark^\checkmark\\checkmarkc\checkmarki\checkmarkr\checkmarkc\checkmark$\checkmarkx\checkmark$\checkmark^\checkmark\\checkmarkc\checkmarki\checkmarkr\checkmarkc\checkmark$\checkmark \checkmark$\checkmark^\checkmark\\checkmarkc\checkmarki\checkmarkr\checkmarkc\checkmark$\checkmark5\checkmark$\checkmark^\checkmark\\checkmarkc\checkmarki\checkmarkr\checkmarkc\checkmark$\checkmark{\checkmark$\checkmark^\checkmark\\checkmarkc\checkmarki\checkmarkr\checkmarkc\checkmark$\checkmark,\checkmark$\checkmark^\checkmark\\checkmarkc\checkmarki\checkmarkr\checkmarkc\checkmark$\checkmark}\checkmark$\checkmark^\checkmark\\checkmarkc\checkmarki\checkmarkr\checkmarkc\checkmark$\checkmark0\checkmark$\checkmark^\checkmark\\checkmarkc\checkmarki\checkmarkr\checkmarkc\checkmark$\checkmark \checkmark$\checkmark^\checkmark\\checkmarkc\checkmarki\checkmarkr\checkmarkc\checkmark$\checkmark\\checkmark$\checkmark^\checkmark\\checkmarkc\checkmarki\checkmarkr\checkmarkc\checkmark$\checkmarkt\checkmark$\checkmark^\checkmark\\checkmarkc\checkmarki\checkmarkr\checkmarkc\checkmark$\checkmarki\checkmark$\checkmark^\checkmark\\checkmarkc\checkmarki\checkmarkr\checkmarkc\checkmark$\checkmarkm\checkmark$\checkmark^\checkmark\\checkmarkc\checkmarki\checkmarkr\checkmarkc\checkmark$\checkmarke\checkmark$\checkmark^\checkmark\\checkmarkc\checkmarki\checkmarkr\checkmarkc\checkmark$\checkmarks\checkmark$\checkmark^\checkmark\\checkmarkc\checkmarki\checkmarkr\checkmarkc\checkmark$\checkmark \checkmark$\checkmark^\checkmark\\checkmarkc\checkmarki\checkmarkr\checkmarkc\checkmark$\checkmark1\checkmark$\checkmark^\checkmark\\checkmarkc\checkmarki\checkmarkr\checkmarkc\checkmark$\checkmark0\checkmark$\checkmark^\checkmark\\checkmarkc\checkmarki\checkmarkr\checkmarkc\checkmark$\checkmark^\checkmark$\checkmark^\checkmark\\checkmarkc\checkmarki\checkmarkr\checkmarkc\checkmark$\checkmark{\checkmark$\checkmark^\checkmark\\checkmarkc\checkmarki\checkmarkr\checkmarkc\checkmark$\checkmark-\checkmark$\checkmark^\checkmark\\checkmarkc\checkmarki\checkmarkr\checkmarkc\checkmark$\checkmark6\checkmark$\checkmark^\checkmark\\checkmarkc\checkmarki\checkmarkr\checkmarkc\checkmark$\checkmark}\checkmark$\checkmark^\checkmark\\checkmarkc\checkmarki\checkmarkr\checkmarkc\checkmark$\checkmark$\checkmark$\checkmark^\checkmark\\checkmarkc\checkmarki\checkmarkr\checkmarkc\checkmark$\checkmark \checkmark$\checkmark^\checkmark\\checkmarkc\checkmarki\checkmarkr\checkmarkc\checkmark$\checkmarkk\checkmark$\checkmark^\checkmark\\checkmarkc\checkmarki\checkmarkr\checkmarkc\checkmark$\checkmarkg\checkmark$\checkmark^\checkmark\\checkmarkc\checkmarki\checkmarkr\checkmarkc\checkmark$\checkmark$\checkmark$\checkmark^\checkmark\\checkmarkc\checkmarki\checkmarkr\checkmarkc\checkmark$\checkmark\\checkmark$\checkmark^\checkmark\\checkmarkc\checkmarki\checkmarkr\checkmarkc\checkmark$\checkmarkc\checkmark$\checkmark^\checkmark\\checkmarkc\checkmarki\checkmarkr\checkmarkc\checkmark$\checkmarkd\checkmark$\checkmark^\checkmark\\checkmarkc\checkmarki\checkmarkr\checkmarkc\checkmark$\checkmarko\checkmark$\checkmark^\checkmark\\checkmarkc\checkmarki\checkmarkr\checkmarkc\checkmark$\checkmarkt\checkmark$\checkmark^\checkmark\\checkmarkc\checkmarki\checkmarkr\checkmarkc\checkmark$\checkmark$\checkmark$\checkmark^\checkmark\\checkmarkc\checkmarki\checkmarkr\checkmarkc\checkmark$\checkmarkm\checkmark$\checkmark^\checkmark\\checkmarkc\checkmarki\checkmarkr\checkmarkc\checkmark$\checkmark²\checkmark$\checkmark^\checkmark\\checkmarkc\checkmarki\checkmarkr\checkmarkc\checkmark$\checkmark \checkmark$\checkmark^\checkmark\\checkmarkc\checkmarki\checkmarkr\checkmarkc\checkmark$\checkmark:\checkmark$\checkmark^\checkmark\\checkmarkc\checkmarki\checkmarkr\checkmarkc\checkmark$\checkmark
\checkmark$\checkmark^\checkmark\\checkmarkc\checkmarki\checkmarkr\checkmarkc\checkmark$\checkmark\\checkmark$\checkmark^\checkmark\\checkmarkc\checkmarki\checkmarkr\checkmarkc\checkmark$\checkmarkb\checkmark$\checkmark^\checkmark\\checkmarkc\checkmarki\checkmarkr\checkmarkc\checkmark$\checkmarke\checkmark$\checkmark^\checkmark\\checkmarkc\checkmarki\checkmarkr\checkmarkc\checkmark$\checkmarkg\checkmark$\checkmark^\checkmark\\checkmarkc\checkmarki\checkmarkr\checkmarkc\checkmark$\checkmarki\checkmark$\checkmark^\checkmark\\checkmarkc\checkmarki\checkmarkr\checkmarkc\checkmark$\checkmarkn\checkmark$\checkmark^\checkmark\\checkmarkc\checkmarki\checkmarkr\checkmarkc\checkmark$\checkmark{\checkmark$\checkmark^\checkmark\\checkmarkc\checkmarki\checkmarkr\checkmarkc\checkmark$\checkmarke\checkmark$\checkmark^\checkmark\\checkmarkc\checkmarki\checkmarkr\checkmarkc\checkmark$\checkmarkq\checkmark$\checkmark^\checkmark\\checkmarkc\checkmarki\checkmarkr\checkmarkc\checkmark$\checkmarku\checkmark$\checkmark^\checkmark\\checkmarkc\checkmarki\checkmarkr\checkmarkc\checkmark$\checkmarka\checkmark$\checkmark^\checkmark\\checkmarkc\checkmarki\checkmarkr\checkmarkc\checkmark$\checkmarkt\checkmark$\checkmark^\checkmark\\checkmarkc\checkmarki\checkmarkr\checkmarkc\checkmark$\checkmarki\checkmark$\checkmark^\checkmark\\checkmarkc\checkmarki\checkmarkr\checkmarkc\checkmark$\checkmarko\checkmark$\checkmark^\checkmark\\checkmarkc\checkmarki\checkmarkr\checkmarkc\checkmark$\checkmarkn\checkmark$\checkmark^\checkmark\\checkmarkc\checkmarki\checkmarkr\checkmarkc\checkmark$\checkmark}\checkmark$\checkmark^\checkmark\\checkmarkc\checkmarki\checkmarkr\checkmarkc\checkmark$\checkmark
\checkmark$\checkmark^\checkmark\\checkmarkc\checkmarki\checkmarkr\checkmarkc\checkmark$\checkmark \checkmark$\checkmark^\checkmark\\checkmarkc\checkmarki\checkmarkr\checkmarkc\checkmark$\checkmark \checkmark$\checkmark^\checkmark\\checkmarkc\checkmarki\checkmarkr\checkmarkc\checkmark$\checkmark \checkmark$\checkmark^\checkmark\\checkmarkc\checkmarki\checkmarkr\checkmarkc\checkmark$\checkmark \checkmark$\checkmark^\checkmark\\checkmarkc\checkmarki\checkmarkr\checkmarkc\checkmark$\checkmarkJ\checkmark$\checkmark^\checkmark\\checkmarkc\checkmarki\checkmarkr\checkmarkc\checkmark$\checkmark_\checkmark$\checkmark^\checkmark\\checkmarkc\checkmarki\checkmarkr\checkmarkc\checkmark$\checkmark{\checkmark$\checkmark^\checkmark\\checkmarkc\checkmarki\checkmarkr\checkmarkc\checkmark$\checkmarkt\checkmark$\checkmark^\checkmark\\checkmarkc\checkmarki\checkmarkr\checkmarkc\checkmark$\checkmarko\checkmark$\checkmark^\checkmark\\checkmarkc\checkmarki\checkmarkr\checkmarkc\checkmark$\checkmarkt\checkmark$\checkmark^\checkmark\\checkmarkc\checkmarki\checkmarkr\checkmarkc\checkmark$\checkmark}\checkmark$\checkmark^\checkmark\\checkmarkc\checkmarki\checkmarkr\checkmarkc\checkmark$\checkmark \checkmark$\checkmark^\checkmark\\checkmarkc\checkmarki\checkmarkr\checkmarkc\checkmark$\checkmark=\checkmark$\checkmark^\checkmark\\checkmarkc\checkmarki\checkmarkr\checkmarkc\checkmark$\checkmark \checkmark$\checkmark^\checkmark\\checkmarkc\checkmarki\checkmarkr\checkmarkc\checkmark$\checkmarkJ\checkmark$\checkmark^\checkmark\\checkmarkc\checkmarki\checkmarkr\checkmarkc\checkmark$\checkmark_\checkmark$\checkmark^\checkmark\\checkmarkc\checkmarki\checkmarkr\checkmarkc\checkmark$\checkmark{\checkmark$\checkmark^\checkmark\\checkmarkc\checkmarki\checkmarkr\checkmarkc\checkmark$\checkmarkv\checkmark$\checkmark^\checkmark\\checkmarkc\checkmarki\checkmarkr\checkmarkc\checkmark$\checkmarki\checkmark$\checkmark^\checkmark\\checkmarkc\checkmarki\checkmarkr\checkmarkc\checkmark$\checkmarks\checkmark$\checkmark^\checkmark\\checkmarkc\checkmarki\checkmarkr\checkmarkc\checkmark$\checkmark}\checkmark$\checkmark^\checkmark\\checkmarkc\checkmarki\checkmarkr\checkmarkc\checkmark$\checkmark \checkmark$\checkmark^\checkmark\\checkmarkc\checkmarki\checkmarkr\checkmarkc\checkmark$\checkmark+\checkmark$\checkmark^\checkmark\\checkmarkc\checkmarki\checkmarkr\checkmarkc\checkmark$\checkmark \checkmark$\checkmark^\checkmark\\checkmarkc\checkmarki\checkmarkr\checkmarkc\checkmark$\checkmarkJ\checkmark$\checkmark^\checkmark\\checkmarkc\checkmarki\checkmarkr\checkmarkc\checkmark$\checkmark_\checkmark$\checkmark^\checkmark\\checkmarkc\checkmarki\checkmarkr\checkmarkc\checkmark$\checkmark{\checkmark$\checkmark^\checkmark\\checkmarkc\checkmarki\checkmarkr\checkmarkc\checkmark$\checkmarkt\checkmark$\checkmark^\checkmark\\checkmarkc\checkmarki\checkmarkr\checkmarkc\checkmark$\checkmarkr\checkmark$\checkmark^\checkmark\\checkmarkc\checkmarki\checkmarkr\checkmarkc\checkmark$\checkmarka\checkmark$\checkmark^\checkmark\\checkmarkc\checkmarki\checkmarkr\checkmarkc\checkmark$\checkmarkn\checkmark$\checkmark^\checkmark\\checkmarkc\checkmarki\checkmarkr\checkmarkc\checkmark$\checkmarks\checkmark$\checkmark^\checkmark\\checkmarkc\checkmarki\checkmarkr\checkmarkc\checkmark$\checkmark}\checkmark$\checkmark^\checkmark\\checkmarkc\checkmarki\checkmarkr\checkmarkc\checkmark$\checkmark \checkmark$\checkmark^\checkmark\\checkmarkc\checkmarki\checkmarkr\checkmarkc\checkmark$\checkmark+\checkmark$\checkmark^\checkmark\\checkmarkc\checkmarki\checkmarkr\checkmarkc\checkmark$\checkmark \checkmark$\checkmark^\checkmark\\checkmarkc\checkmarki\checkmarkr\checkmarkc\checkmark$\checkmarkJ\checkmark$\checkmark^\checkmark\\checkmarkc\checkmarki\checkmarkr\checkmarkc\checkmark$\checkmark_\checkmark$\checkmark^\checkmark\\checkmarkc\checkmarki\checkmarkr\checkmarkc\checkmark$\checkmark{\checkmark$\checkmark^\checkmark\\checkmarkc\checkmarki\checkmarkr\checkmarkc\checkmark$\checkmarkm\checkmark$\checkmark^\checkmark\\checkmarkc\checkmarki\checkmarkr\checkmarkc\checkmark$\checkmarko\checkmark$\checkmark^\checkmark\\checkmarkc\checkmarki\checkmarkr\checkmarkc\checkmark$\checkmarkt\checkmark$\checkmark^\checkmark\\checkmarkc\checkmarki\checkmarkr\checkmarkc\checkmark$\checkmark}\checkmark$\checkmark^\checkmark\\checkmarkc\checkmarki\checkmarkr\checkmarkc\checkmark$\checkmark \checkmark$\checkmark^\checkmark\\checkmarkc\checkmarki\checkmarkr\checkmarkc\checkmark$\checkmark=\checkmark$\checkmark^\checkmark\\checkmarkc\checkmarki\checkmarkr\checkmarkc\checkmark$\checkmark \checkmark$\checkmark^\checkmark\\checkmarkc\checkmarki\checkmarkr\checkmarkc\checkmark$\checkmark2\checkmark$\checkmark^\checkmark\\checkmarkc\checkmarki\checkmarkr\checkmarkc\checkmark$\checkmark{\checkmark$\checkmark^\checkmark\\checkmarkc\checkmarki\checkmarkr\checkmarkc\checkmark$\checkmark,\checkmark$\checkmark^\checkmark\\checkmarkc\checkmarki\checkmarkr\checkmarkc\checkmark$\checkmark}\checkmark$\checkmark^\checkmark\\checkmarkc\checkmarki\checkmarkr\checkmarkc\checkmark$\checkmark0\checkmark$\checkmark^\checkmark\\checkmarkc\checkmarki\checkmarkr\checkmarkc\checkmark$\checkmark1\checkmark$\checkmark^\checkmark\\checkmarkc\checkmarki\checkmarkr\checkmarkc\checkmark$\checkmark6\checkmark$\checkmark^\checkmark\\checkmarkc\checkmarki\checkmarkr\checkmarkc\checkmark$\checkmark\\checkmark$\checkmark^\checkmark\\checkmarkc\checkmarki\checkmarkr\checkmarkc\checkmark$\checkmarkt\checkmark$\checkmark^\checkmark\\checkmarkc\checkmarki\checkmarkr\checkmarkc\checkmark$\checkmarki\checkmark$\checkmark^\checkmark\\checkmarkc\checkmarki\checkmarkr\checkmarkc\checkmark$\checkmarkm\checkmark$\checkmark^\checkmark\\checkmarkc\checkmarki\checkmarkr\checkmarkc\checkmark$\checkmarke\checkmark$\checkmark^\checkmark\\checkmarkc\checkmarki\checkmarkr\checkmarkc\checkmark$\checkmarks\checkmark$\checkmark^\checkmark\\checkmarkc\checkmarki\checkmarkr\checkmarkc\checkmark$\checkmark1\checkmark$\checkmark^\checkmark\\checkmarkc\checkmarki\checkmarkr\checkmarkc\checkmark$\checkmark0\checkmark$\checkmark^\checkmark\\checkmarkc\checkmarki\checkmarkr\checkmarkc\checkmark$\checkmark^\checkmark$\checkmark^\checkmark\\checkmarkc\checkmarki\checkmarkr\checkmarkc\checkmark$\checkmark{\checkmark$\checkmark^\checkmark\\checkmarkc\checkmarki\checkmarkr\checkmarkc\checkmark$\checkmark-\checkmark$\checkmark^\checkmark\\checkmarkc\checkmarki\checkmarkr\checkmarkc\checkmark$\checkmark5\checkmark$\checkmark^\checkmark\\checkmarkc\checkmarki\checkmarkr\checkmarkc\checkmark$\checkmark}\checkmark$\checkmark^\checkmark\\checkmarkc\checkmarki\checkmarkr\checkmarkc\checkmark$\checkmark \checkmark$\checkmark^\checkmark\\checkmarkc\checkmarki\checkmarkr\checkmarkc\checkmark$\checkmark+\checkmark$\checkmark^\checkmark\\checkmarkc\checkmarki\checkmarkr\checkmarkc\checkmark$\checkmark \checkmark$\checkmark^\checkmark\\checkmarkc\checkmarki\checkmarkr\checkmarkc\checkmark$\checkmark6\checkmark$\checkmark^\checkmark\\checkmarkc\checkmarki\checkmarkr\checkmarkc\checkmark$\checkmark{\checkmark$\checkmark^\checkmark\\checkmarkc\checkmarki\checkmarkr\checkmarkc\checkmark$\checkmark,\checkmark$\checkmark^\checkmark\\checkmarkc\checkmarki\checkmarkr\checkmarkc\checkmark$\checkmark}\checkmark$\checkmark^\checkmark\\checkmarkc\checkmarki\checkmarkr\checkmarkc\checkmark$\checkmark3\checkmark$\checkmark^\checkmark\\checkmarkc\checkmarki\checkmarkr\checkmarkc\checkmark$\checkmark3\checkmark$\checkmark^\checkmark\\checkmarkc\checkmarki\checkmarkr\checkmarkc\checkmark$\checkmark\\checkmark$\checkmark^\checkmark\\checkmarkc\checkmarki\checkmarkr\checkmarkc\checkmark$\checkmarkt\checkmark$\checkmark^\checkmark\\checkmarkc\checkmarki\checkmarkr\checkmarkc\checkmark$\checkmarki\checkmark$\checkmark^\checkmark\\checkmarkc\checkmarki\checkmarkr\checkmarkc\checkmark$\checkmarkm\checkmark$\checkmark^\checkmark\\checkmarkc\checkmarki\checkmarkr\checkmarkc\checkmark$\checkmarke\checkmark$\checkmark^\checkmark\\checkmarkc\checkmarki\checkmarkr\checkmarkc\checkmark$\checkmarks\checkmark$\checkmark^\checkmark\\checkmarkc\checkmarki\checkmarkr\checkmarkc\checkmark$\checkmark1\checkmark$\checkmark^\checkmark\\checkmarkc\checkmarki\checkmarkr\checkmarkc\checkmark$\checkmark0\checkmark$\checkmark^\checkmark\\checkmarkc\checkmarki\checkmarkr\checkmarkc\checkmark$\checkmark^\checkmark$\checkmark^\checkmark\\checkmarkc\checkmarki\checkmarkr\checkmarkc\checkmark$\checkmark{\checkmark$\checkmark^\checkmark\\checkmarkc\checkmarki\checkmarkr\checkmarkc\checkmark$\checkmark-\checkmark$\checkmark^\checkmark\\checkmarkc\checkmarki\checkmarkr\checkmarkc\checkmark$\checkmark6\checkmark$\checkmark^\checkmark\\checkmarkc\checkmarki\checkmarkr\checkmarkc\checkmark$\checkmark}\checkmark$\checkmark^\checkmark\\checkmarkc\checkmarki\checkmarkr\checkmarkc\checkmark$\checkmark \checkmark$\checkmark^\checkmark\\checkmarkc\checkmarki\checkmarkr\checkmarkc\checkmark$\checkmark+\checkmark$\checkmark^\checkmark\\checkmarkc\checkmarki\checkmarkr\checkmarkc\checkmark$\checkmark \checkmark$\checkmark^\checkmark\\checkmarkc\checkmarki\checkmarkr\checkmarkc\checkmark$\checkmark5\checkmark$\checkmark^\checkmark\\checkmarkc\checkmarki\checkmarkr\checkmarkc\checkmark$\checkmark{\checkmark$\checkmark^\checkmark\\checkmarkc\checkmarki\checkmarkr\checkmarkc\checkmark$\checkmark,\checkmark$\checkmark^\checkmark\\checkmarkc\checkmarki\checkmarkr\checkmarkc\checkmark$\checkmark}\checkmark$\checkmark^\checkmark\\checkmarkc\checkmarki\checkmarkr\checkmarkc\checkmark$\checkmark0\checkmark$\checkmark^\checkmark\\checkmarkc\checkmarki\checkmarkr\checkmarkc\checkmark$\checkmark\\checkmark$\checkmark^\checkmark\\checkmarkc\checkmarki\checkmarkr\checkmarkc\checkmark$\checkmarkt\checkmark$\checkmark^\checkmark\\checkmarkc\checkmarki\checkmarkr\checkmarkc\checkmark$\checkmarki\checkmark$\checkmark^\checkmark\\checkmarkc\checkmarki\checkmarkr\checkmarkc\checkmark$\checkmarkm\checkmark$\checkmark^\checkmark\\checkmarkc\checkmarki\checkmarkr\checkmarkc\checkmark$\checkmarke\checkmark$\checkmark^\checkmark\\checkmarkc\checkmarki\checkmarkr\checkmarkc\checkmark$\checkmarks\checkmark$\checkmark^\checkmark\\checkmarkc\checkmarki\checkmarkr\checkmarkc\checkmark$\checkmark1\checkmark$\checkmark^\checkmark\\checkmarkc\checkmarki\checkmarkr\checkmarkc\checkmark$\checkmark0\checkmark$\checkmark^\checkmark\\checkmarkc\checkmarki\checkmarkr\checkmarkc\checkmark$\checkmark^\checkmark$\checkmark^\checkmark\\checkmarkc\checkmarki\checkmarkr\checkmarkc\checkmark$\checkmark{\checkmark$\checkmark^\checkmark\\checkmarkc\checkmarki\checkmarkr\checkmarkc\checkmark$\checkmark-\checkmark$\checkmark^\checkmark\\checkmarkc\checkmarki\checkmarkr\checkmarkc\checkmark$\checkmark6\checkmark$\checkmark^\checkmark\\checkmarkc\checkmarki\checkmarkr\checkmarkc\checkmark$\checkmark}\checkmark$\checkmark^\checkmark\\checkmarkc\checkmarki\checkmarkr\checkmarkc\checkmark$\checkmark \checkmark$\checkmark^\checkmark\\checkmarkc\checkmarki\checkmarkr\checkmarkc\checkmark$\checkmark=\checkmark$\checkmark^\checkmark\\checkmarkc\checkmarki\checkmarkr\checkmarkc\checkmark$\checkmark \checkmark$\checkmark^\checkmark\\checkmarkc\checkmarki\checkmarkr\checkmarkc\checkmark$\checkmark\\checkmark$\checkmark^\checkmark\\checkmarkc\checkmarki\checkmarkr\checkmarkc\checkmark$\checkmarkb\checkmark$\checkmark^\checkmark\\checkmarkc\checkmarki\checkmarkr\checkmarkc\checkmark$\checkmarko\checkmark$\checkmark^\checkmark\\checkmarkc\checkmarki\checkmarkr\checkmarkc\checkmark$\checkmarkx\checkmark$\checkmark^\checkmark\\checkmarkc\checkmarki\checkmarkr\checkmarkc\checkmark$\checkmarke\checkmark$\checkmark^\checkmark\\checkmarkc\checkmarki\checkmarkr\checkmarkc\checkmark$\checkmarkd\checkmark$\checkmark^\checkmark\\checkmarkc\checkmarki\checkmarkr\checkmarkc\checkmark$\checkmark{\checkmark$\checkmark^\checkmark\\checkmarkc\checkmarki\checkmarkr\checkmarkc\checkmark$\checkmark3\checkmark$\checkmark^\checkmark\\checkmarkc\checkmarki\checkmarkr\checkmarkc\checkmark$\checkmark{\checkmark$\checkmark^\checkmark\\checkmarkc\checkmarki\checkmarkr\checkmarkc\checkmark$\checkmark,\checkmark$\checkmark^\checkmark\\checkmarkc\checkmarki\checkmarkr\checkmarkc\checkmark$\checkmark}\checkmark$\checkmark^\checkmark\\checkmarkc\checkmarki\checkmarkr\checkmarkc\checkmark$\checkmark1\checkmark$\checkmark^\checkmark\\checkmarkc\checkmarki\checkmarkr\checkmarkc\checkmark$\checkmark5\checkmark$\checkmark^\checkmark\\checkmarkc\checkmarki\checkmarkr\checkmarkc\checkmark$\checkmark \checkmark$\checkmark^\checkmark\\checkmarkc\checkmarki\checkmarkr\checkmarkc\checkmark$\checkmark\\checkmark$\checkmark^\checkmark\\checkmarkc\checkmarki\checkmarkr\checkmarkc\checkmark$\checkmarkt\checkmark$\checkmark^\checkmark\\checkmarkc\checkmarki\checkmarkr\checkmarkc\checkmark$\checkmarki\checkmark$\checkmark^\checkmark\\checkmarkc\checkmarki\checkmarkr\checkmarkc\checkmark$\checkmarkm\checkmark$\checkmark^\checkmark\\checkmarkc\checkmarki\checkmarkr\checkmarkc\checkmark$\checkmarke\checkmark$\checkmark^\checkmark\\checkmarkc\checkmarki\checkmarkr\checkmarkc\checkmark$\checkmarks\checkmark$\checkmark^\checkmark\\checkmarkc\checkmarki\checkmarkr\checkmarkc\checkmark$\checkmark \checkmark$\checkmark^\checkmark\\checkmarkc\checkmarki\checkmarkr\checkmarkc\checkmark$\checkmark1\checkmark$\checkmark^\checkmark\\checkmarkc\checkmarki\checkmarkr\checkmarkc\checkmark$\checkmark0\checkmark$\checkmark^\checkmark\\checkmarkc\checkmarki\checkmarkr\checkmarkc\checkmark$\checkmark^\checkmark$\checkmark^\checkmark\\checkmarkc\checkmarki\checkmarkr\checkmarkc\checkmark$\checkmark{\checkmark$\checkmark^\checkmark\\checkmarkc\checkmarki\checkmarkr\checkmarkc\checkmark$\checkmark-\checkmark$\checkmark^\checkmark\\checkmarkc\checkmarki\checkmarkr\checkmarkc\checkmark$\checkmark5\checkmark$\checkmark^\checkmark\\checkmarkc\checkmarki\checkmarkr\checkmarkc\checkmark$\checkmark}\checkmark$\checkmark^\checkmark\\checkmarkc\checkmarki\checkmarkr\checkmarkc\checkmark$\checkmark \checkmark$\checkmark^\checkmark\\checkmarkc\checkmarki\checkmarkr\checkmarkc\checkmark$\checkmark\\checkmark$\checkmark^\checkmark\\checkmarkc\checkmarki\checkmarkr\checkmarkc\checkmark$\checkmarkt\checkmark$\checkmark^\checkmark\\checkmarkc\checkmarki\checkmarkr\checkmarkc\checkmark$\checkmarke\checkmark$\checkmark^\checkmark\\checkmarkc\checkmarki\checkmarkr\checkmarkc\checkmark$\checkmarkx\checkmark$\checkmark^\checkmark\\checkmarkc\checkmarki\checkmarkr\checkmarkc\checkmark$\checkmarkt\checkmark$\checkmark^\checkmark\\checkmarkc\checkmarki\checkmarkr\checkmarkc\checkmark$\checkmark{\checkmark$\checkmark^\checkmark\\checkmarkc\checkmarki\checkmarkr\checkmarkc\checkmark$\checkmark \checkmark$\checkmark^\checkmark\\checkmarkc\checkmarki\checkmarkr\checkmarkc\checkmark$\checkmarkk\checkmark$\checkmark^\checkmark\\checkmarkc\checkmarki\checkmarkr\checkmarkc\checkmark$\checkmarkg\checkmark$\checkmark^\checkmark\\checkmarkc\checkmarki\checkmarkr\checkmarkc\checkmark$\checkmark$\checkmark$\checkmark^\checkmark\\checkmarkc\checkmarki\checkmarkr\checkmarkc\checkmark$\checkmark\\checkmark$\checkmark^\checkmark\\checkmarkc\checkmarki\checkmarkr\checkmarkc\checkmark$\checkmarkc\checkmark$\checkmark^\checkmark\\checkmarkc\checkmarki\checkmarkr\checkmarkc\checkmark$\checkmarkd\checkmark$\checkmark^\checkmark\\checkmarkc\checkmarki\checkmarkr\checkmarkc\checkmark$\checkmarko\checkmark$\checkmark^\checkmark\\checkmarkc\checkmarki\checkmarkr\checkmarkc\checkmark$\checkmarkt\checkmark$\checkmark^\checkmark\\checkmarkc\checkmarki\checkmarkr\checkmarkc\checkmark$\checkmark$\checkmark$\checkmark^\checkmark\\checkmarkc\checkmarki\checkmarkr\checkmarkc\checkmark$\checkmarkm\checkmark$\checkmark^\checkmark\\checkmarkc\checkmarki\checkmarkr\checkmarkc\checkmark$\checkmark}\checkmark$\checkmark^\checkmark\\checkmarkc\checkmarki\checkmarkr\checkmarkc\checkmark$\checkmark^\checkmark$\checkmark^\checkmark\\checkmarkc\checkmarki\checkmarkr\checkmarkc\checkmark$\checkmark2\checkmark$\checkmark^\checkmark\\checkmarkc\checkmarki\checkmarkr\checkmarkc\checkmark$\checkmark}\checkmark$\checkmark^\checkmark\\checkmarkc\checkmarki\checkmarkr\checkmarkc\checkmark$\checkmark
\checkmark$\checkmark^\checkmark\\checkmarkc\checkmarki\checkmarkr\checkmarkc\checkmark$\checkmark\\checkmark$\checkmark^\checkmark\\checkmarkc\checkmarki\checkmarkr\checkmarkc\checkmark$\checkmarke\checkmark$\checkmark^\checkmark\\checkmarkc\checkmarki\checkmarkr\checkmarkc\checkmark$\checkmarkn\checkmark$\checkmark^\checkmark\\checkmarkc\checkmarki\checkmarkr\checkmarkc\checkmark$\checkmarkd\checkmark$\checkmark^\checkmark\\checkmarkc\checkmarki\checkmarkr\checkmarkc\checkmark$\checkmark{\checkmark$\checkmark^\checkmark\\checkmarkc\checkmarki\checkmarkr\checkmarkc\checkmark$\checkmarke\checkmark$\checkmark^\checkmark\\checkmarkc\checkmarki\checkmarkr\checkmarkc\checkmark$\checkmarkq\checkmark$\checkmark^\checkmark\\checkmarkc\checkmarki\checkmarkr\checkmarkc\checkmark$\checkmarku\checkmark$\checkmark^\checkmark\\checkmarkc\checkmarki\checkmarkr\checkmarkc\checkmark$\checkmarka\checkmark$\checkmark^\checkmark\\checkmarkc\checkmarki\checkmarkr\checkmarkc\checkmark$\checkmarkt\checkmark$\checkmark^\checkmark\\checkmarkc\checkmarki\checkmarkr\checkmarkc\checkmark$\checkmarki\checkmark$\checkmark^\checkmark\\checkmarkc\checkmarki\checkmarkr\checkmarkc\checkmark$\checkmarko\checkmark$\checkmark^\checkmark\\checkmarkc\checkmarki\checkmarkr\checkmarkc\checkmark$\checkmarkn\checkmark$\checkmark^\checkmark\\checkmarkc\checkmarki\checkmarkr\checkmarkc\checkmark$\checkmark}\checkmark$\checkmark^\checkmark\\checkmarkc\checkmarki\checkmarkr\checkmarkc\checkmark$\checkmark
\checkmark$\checkmark^\checkmark\\checkmarkc\checkmarki\checkmarkr\checkmarkc\checkmark$\checkmark
\checkmark$\checkmark^\checkmark\\checkmarkc\checkmarki\checkmarkr\checkmarkc\checkmark$\checkmarkP\checkmark$\checkmark^\checkmark\\checkmarkc\checkmarki\checkmarkr\checkmarkc\checkmark$\checkmarko\checkmark$\checkmark^\checkmark\\checkmarkc\checkmarki\checkmarkr\checkmarkc\checkmark$\checkmarku\checkmark$\checkmark^\checkmark\\checkmarkc\checkmarki\checkmarkr\checkmarkc\checkmark$\checkmarkr\checkmark$\checkmark^\checkmark\\checkmarkc\checkmarki\checkmarkr\checkmarkc\checkmark$\checkmark \checkmark$\checkmark^\checkmark\\checkmarkc\checkmarki\checkmarkr\checkmarkc\checkmark$\checkmarku\checkmark$\checkmark^\checkmark\\checkmarkc\checkmarki\checkmarkr\checkmarkc\checkmark$\checkmarkn\checkmark$\checkmark^\checkmark\\checkmarkc\checkmarki\checkmarkr\checkmarkc\checkmark$\checkmarke\checkmark$\checkmark^\checkmark\\checkmarkc\checkmarki\checkmarkr\checkmarkc\checkmark$\checkmark \checkmark$\checkmark^\checkmark\\checkmarkc\checkmarki\checkmarkr\checkmarkc\checkmark$\checkmarka\checkmark$\checkmark^\checkmark\\checkmarkc\checkmarki\checkmarkr\checkmarkc\checkmark$\checkmarkc\checkmark$\checkmark^\checkmark\\checkmarkc\checkmarki\checkmarkr\checkmarkc\checkmark$\checkmarkc\checkmark$\checkmark^\checkmark\\checkmarkc\checkmarki\checkmarkr\checkmarkc\checkmark$\checkmarké\checkmark$\checkmark^\checkmark\\checkmarkc\checkmarki\checkmarkr\checkmarkc\checkmark$\checkmarkl\checkmark$\checkmark^\checkmark\\checkmarkc\checkmarki\checkmarkr\checkmarkc\checkmark$\checkmarké\checkmark$\checkmark^\checkmark\\checkmarkc\checkmarki\checkmarkr\checkmarkc\checkmark$\checkmarkr\checkmark$\checkmark^\checkmark\\checkmarkc\checkmarki\checkmarkr\checkmarkc\checkmark$\checkmarka\checkmark$\checkmark^\checkmark\\checkmarkc\checkmarki\checkmarkr\checkmarkc\checkmark$\checkmarkt\checkmark$\checkmark^\checkmark\\checkmarkc\checkmarki\checkmarkr\checkmarkc\checkmark$\checkmarki\checkmark$\checkmark^\checkmark\\checkmarkc\checkmarki\checkmarkr\checkmarkc\checkmark$\checkmarko\checkmark$\checkmark^\checkmark\\checkmarkc\checkmarki\checkmarkr\checkmarkc\checkmark$\checkmarkn\checkmark$\checkmark^\checkmark\\checkmarkc\checkmarki\checkmarkr\checkmarkc\checkmark$\checkmark \checkmark$\checkmark^\checkmark\\checkmarkc\checkmarki\checkmarkr\checkmarkc\checkmark$\checkmarkl\checkmark$\checkmark^\checkmark\\checkmarkc\checkmarki\checkmarkr\checkmarkc\checkmark$\checkmarki\checkmark$\checkmark^\checkmark\\checkmarkc\checkmarki\checkmarkr\checkmarkc\checkmark$\checkmarkn\checkmark$\checkmark^\checkmark\\checkmarkc\checkmarki\checkmarkr\checkmarkc\checkmark$\checkmarké\checkmark$\checkmark^\checkmark\\checkmarkc\checkmarki\checkmarkr\checkmarkc\checkmark$\checkmarka\checkmark$\checkmark^\checkmark\\checkmarkc\checkmarki\checkmarkr\checkmarkc\checkmark$\checkmarki\checkmark$\checkmark^\checkmark\\checkmarkc\checkmarki\checkmarkr\checkmarkc\checkmark$\checkmarkr\checkmark$\checkmark^\checkmark\\checkmarkc\checkmarki\checkmarkr\checkmarkc\checkmark$\checkmarke\checkmark$\checkmark^\checkmark\\checkmarkc\checkmarki\checkmarkr\checkmarkc\checkmark$\checkmark \checkmark$\checkmark^\checkmark\\checkmarkc\checkmarki\checkmarkr\checkmarkc\checkmark$\checkmark$\checkmark$\checkmark^\checkmark\\checkmarkc\checkmarki\checkmarkr\checkmarkc\checkmark$\checkmarka\checkmark$\checkmark^\checkmark\\checkmarkc\checkmarki\checkmarkr\checkmarkc\checkmark$\checkmark_\checkmark$\checkmark^\checkmark\\checkmarkc\checkmarki\checkmarkr\checkmarkc\checkmark$\checkmark{\checkmark$\checkmark^\checkmark\\checkmarkc\checkmarki\checkmarkr\checkmarkc\checkmark$\checkmarkm\checkmark$\checkmark^\checkmark\\checkmarkc\checkmarki\checkmarkr\checkmarkc\checkmark$\checkmarka\checkmark$\checkmark^\checkmark\\checkmarkc\checkmarki\checkmarkr\checkmarkc\checkmark$\checkmarkx\checkmark$\checkmark^\checkmark\\checkmarkc\checkmarki\checkmarkr\checkmarkc\checkmark$\checkmark}\checkmark$\checkmark^\checkmark\\checkmarkc\checkmarki\checkmarkr\checkmarkc\checkmark$\checkmark \checkmark$\checkmark^\checkmark\\checkmarkc\checkmarki\checkmarkr\checkmarkc\checkmark$\checkmark=\checkmark$\checkmark^\checkmark\\checkmarkc\checkmarki\checkmarkr\checkmarkc\checkmark$\checkmark \checkmark$\checkmark^\checkmark\\checkmarkc\checkmarki\checkmarkr\checkmarkc\checkmark$\checkmark1\checkmark$\checkmark^\checkmark\\checkmarkc\checkmarki\checkmarkr\checkmarkc\checkmark$\checkmark$\checkmark$\checkmark^\checkmark\\checkmarkc\checkmarki\checkmarkr\checkmarkc\checkmark$\checkmark \checkmark$\checkmark^\checkmark\\checkmarkc\checkmarki\checkmarkr\checkmarkc\checkmark$\checkmarkm\checkmark$\checkmark^\checkmark\\checkmarkc\checkmarki\checkmarkr\checkmarkc\checkmark$\checkmark/\checkmark$\checkmark^\checkmark\\checkmarkc\checkmarki\checkmarkr\checkmarkc\checkmark$\checkmarks\checkmark$\checkmark^\checkmark\\checkmarkc\checkmarki\checkmarkr\checkmarkc\checkmark$\checkmark²\checkmark$\checkmark^\checkmark\\checkmarkc\checkmarki\checkmarkr\checkmarkc\checkmark$\checkmark,\checkmark$\checkmark^\checkmark\\checkmarkc\checkmarki\checkmarkr\checkmarkc\checkmark$\checkmark \checkmark$\checkmark^\checkmark\\checkmarkc\checkmarki\checkmarkr\checkmarkc\checkmark$\checkmarkl\checkmark$\checkmark^\checkmark\\checkmarkc\checkmarki\checkmarkr\checkmarkc\checkmark$\checkmark'\checkmark$\checkmark^\checkmark\\checkmarkc\checkmarki\checkmarkr\checkmarkc\checkmark$\checkmarka\checkmark$\checkmark^\checkmark\\checkmarkc\checkmarki\checkmarkr\checkmarkc\checkmark$\checkmarkc\checkmark$\checkmark^\checkmark\\checkmarkc\checkmarki\checkmarkr\checkmarkc\checkmark$\checkmarkc\checkmark$\checkmark^\checkmark\\checkmarkc\checkmarki\checkmarkr\checkmarkc\checkmark$\checkmarké\checkmark$\checkmark^\checkmark\\checkmarkc\checkmarki\checkmarkr\checkmarkc\checkmark$\checkmarkl\checkmark$\checkmark^\checkmark\\checkmarkc\checkmarki\checkmarkr\checkmarkc\checkmark$\checkmarké\checkmark$\checkmark^\checkmark\\checkmarkc\checkmarki\checkmarkr\checkmarkc\checkmark$\checkmarkr\checkmark$\checkmark^\checkmark\\checkmarkc\checkmarki\checkmarkr\checkmarkc\checkmark$\checkmarka\checkmark$\checkmark^\checkmark\\checkmarkc\checkmarki\checkmarkr\checkmarkc\checkmark$\checkmarkt\checkmark$\checkmark^\checkmark\\checkmarkc\checkmarki\checkmarkr\checkmarkc\checkmark$\checkmarki\checkmark$\checkmark^\checkmark\\checkmarkc\checkmarki\checkmarkr\checkmarkc\checkmark$\checkmarko\checkmark$\checkmark^\checkmark\\checkmarkc\checkmarki\checkmarkr\checkmarkc\checkmark$\checkmarkn\checkmark$\checkmark^\checkmark\\checkmarkc\checkmarki\checkmarkr\checkmarkc\checkmark$\checkmark \checkmark$\checkmark^\checkmark\\checkmarkc\checkmarki\checkmarkr\checkmarkc\checkmark$\checkmarka\checkmark$\checkmark^\checkmark\\checkmarkc\checkmarki\checkmarkr\checkmarkc\checkmark$\checkmarkn\checkmark$\checkmark^\checkmark\\checkmarkc\checkmarki\checkmarkr\checkmarkc\checkmark$\checkmarkg\checkmark$\checkmark^\checkmark\\checkmarkc\checkmarki\checkmarkr\checkmarkc\checkmark$\checkmarku\checkmark$\checkmark^\checkmark\\checkmarkc\checkmarki\checkmarkr\checkmarkc\checkmark$\checkmarkl\checkmark$\checkmark^\checkmark\\checkmarkc\checkmarki\checkmarkr\checkmarkc\checkmark$\checkmarka\checkmark$\checkmark^\checkmark\\checkmarkc\checkmarki\checkmarkr\checkmarkc\checkmark$\checkmarki\checkmark$\checkmark^\checkmark\\checkmarkc\checkmarki\checkmarkr\checkmarkc\checkmark$\checkmarkr\checkmark$\checkmark^\checkmark\\checkmarkc\checkmarki\checkmarkr\checkmarkc\checkmark$\checkmarke\checkmark$\checkmark^\checkmark\\checkmarkc\checkmarki\checkmarkr\checkmarkc\checkmark$\checkmark \checkmark$\checkmark^\checkmark\\checkmarkc\checkmarki\checkmarkr\checkmarkc\checkmark$\checkmarke\checkmark$\checkmark^\checkmark\\checkmarkc\checkmarki\checkmarkr\checkmarkc\checkmark$\checkmarks\checkmark$\checkmark^\checkmark\\checkmarkc\checkmarki\checkmarkr\checkmarkc\checkmark$\checkmarkt\checkmark$\checkmark^\checkmark\\checkmarkc\checkmarki\checkmarkr\checkmarkc\checkmark$\checkmark \checkmark$\checkmark^\checkmark\\checkmarkc\checkmarki\checkmarkr\checkmarkc\checkmark$\checkmark:\checkmark$\checkmark^\checkmark\\checkmarkc\checkmarki\checkmarkr\checkmarkc\checkmark$\checkmark
\checkmark$\checkmark^\checkmark\\checkmarkc\checkmarki\checkmarkr\checkmarkc\checkmark$\checkmark\\checkmark$\checkmark^\checkmark\\checkmarkc\checkmarki\checkmarkr\checkmarkc\checkmark$\checkmarkb\checkmark$\checkmark^\checkmark\\checkmarkc\checkmarki\checkmarkr\checkmarkc\checkmark$\checkmarke\checkmark$\checkmark^\checkmark\\checkmarkc\checkmarki\checkmarkr\checkmarkc\checkmark$\checkmarkg\checkmark$\checkmark^\checkmark\\checkmarkc\checkmarki\checkmarkr\checkmarkc\checkmark$\checkmarki\checkmark$\checkmark^\checkmark\\checkmarkc\checkmarki\checkmarkr\checkmarkc\checkmark$\checkmarkn\checkmark$\checkmark^\checkmark\\checkmarkc\checkmarki\checkmarkr\checkmarkc\checkmark$\checkmark{\checkmark$\checkmark^\checkmark\\checkmarkc\checkmarki\checkmarkr\checkmarkc\checkmark$\checkmarke\checkmark$\checkmark^\checkmark\\checkmarkc\checkmarki\checkmarkr\checkmarkc\checkmark$\checkmarkq\checkmark$\checkmark^\checkmark\\checkmarkc\checkmarki\checkmarkr\checkmarkc\checkmark$\checkmarku\checkmark$\checkmark^\checkmark\\checkmarkc\checkmarki\checkmarkr\checkmarkc\checkmark$\checkmarka\checkmark$\checkmark^\checkmark\\checkmarkc\checkmarki\checkmarkr\checkmarkc\checkmark$\checkmarkt\checkmark$\checkmark^\checkmark\\checkmarkc\checkmarki\checkmarkr\checkmarkc\checkmark$\checkmarki\checkmark$\checkmark^\checkmark\\checkmarkc\checkmarki\checkmarkr\checkmarkc\checkmark$\checkmarko\checkmark$\checkmark^\checkmark\\checkmarkc\checkmarki\checkmarkr\checkmarkc\checkmark$\checkmarkn\checkmark$\checkmark^\checkmark\\checkmarkc\checkmarki\checkmarkr\checkmarkc\checkmark$\checkmark}\checkmark$\checkmark^\checkmark\\checkmarkc\checkmarki\checkmarkr\checkmarkc\checkmark$\checkmark
\checkmark$\checkmark^\checkmark\\checkmarkc\checkmarki\checkmarkr\checkmarkc\checkmark$\checkmark \checkmark$\checkmark^\checkmark\\checkmarkc\checkmarki\checkmarkr\checkmarkc\checkmark$\checkmark \checkmark$\checkmark^\checkmark\\checkmarkc\checkmarki\checkmarkr\checkmarkc\checkmark$\checkmark \checkmark$\checkmark^\checkmark\\checkmarkc\checkmarki\checkmarkr\checkmarkc\checkmark$\checkmark \checkmark$\checkmark^\checkmark\\checkmarkc\checkmarki\checkmarkr\checkmarkc\checkmark$\checkmark\\checkmark$\checkmark^\checkmark\\checkmarkc\checkmarki\checkmarkr\checkmarkc\checkmark$\checkmarka\checkmark$\checkmark^\checkmark\\checkmarkc\checkmarki\checkmarkr\checkmarkc\checkmark$\checkmarkl\checkmark$\checkmark^\checkmark\\checkmarkc\checkmarki\checkmarkr\checkmarkc\checkmark$\checkmarkp\checkmark$\checkmark^\checkmark\\checkmarkc\checkmarki\checkmarkr\checkmarkc\checkmark$\checkmarkh\checkmark$\checkmark^\checkmark\\checkmarkc\checkmarki\checkmarkr\checkmarkc\checkmark$\checkmarka\checkmark$\checkmark^\checkmark\\checkmarkc\checkmarki\checkmarkr\checkmarkc\checkmark$\checkmark \checkmark$\checkmark^\checkmark\\checkmarkc\checkmarki\checkmarkr\checkmarkc\checkmark$\checkmark=\checkmark$\checkmark^\checkmark\\checkmarkc\checkmarki\checkmarkr\checkmarkc\checkmark$\checkmark \checkmark$\checkmark^\checkmark\\checkmarkc\checkmarki\checkmarkr\checkmarkc\checkmark$\checkmarka\checkmark$\checkmark^\checkmark\\checkmarkc\checkmarki\checkmarkr\checkmarkc\checkmark$\checkmark_\checkmark$\checkmark^\checkmark\\checkmarkc\checkmarki\checkmarkr\checkmarkc\checkmark$\checkmark{\checkmark$\checkmark^\checkmark\\checkmarkc\checkmarki\checkmarkr\checkmarkc\checkmark$\checkmarkm\checkmark$\checkmark^\checkmark\\checkmarkc\checkmarki\checkmarkr\checkmarkc\checkmark$\checkmarka\checkmark$\checkmark^\checkmark\\checkmarkc\checkmarki\checkmarkr\checkmarkc\checkmark$\checkmarkx\checkmark$\checkmark^\checkmark\\checkmarkc\checkmarki\checkmarkr\checkmarkc\checkmark$\checkmark}\checkmark$\checkmark^\checkmark\\checkmarkc\checkmarki\checkmarkr\checkmarkc\checkmark$\checkmark \checkmark$\checkmark^\checkmark\\checkmarkc\checkmarki\checkmarkr\checkmarkc\checkmark$\checkmark\\checkmark$\checkmark^\checkmark\\checkmarkc\checkmarki\checkmarkr\checkmarkc\checkmark$\checkmarkt\checkmark$\checkmark^\checkmark\\checkmarkc\checkmarki\checkmarkr\checkmarkc\checkmark$\checkmarki\checkmark$\checkmark^\checkmark\\checkmarkc\checkmarki\checkmarkr\checkmarkc\checkmark$\checkmarkm\checkmark$\checkmark^\checkmark\\checkmarkc\checkmarki\checkmarkr\checkmarkc\checkmark$\checkmarke\checkmark$\checkmark^\checkmark\\checkmarkc\checkmarki\checkmarkr\checkmarkc\checkmark$\checkmarks\checkmark$\checkmark^\checkmark\\checkmarkc\checkmarki\checkmarkr\checkmarkc\checkmark$\checkmark \checkmark$\checkmark^\checkmark\\checkmarkc\checkmarki\checkmarkr\checkmarkc\checkmark$\checkmark\\checkmark$\checkmark^\checkmark\\checkmarkc\checkmarki\checkmarkr\checkmarkc\checkmark$\checkmarkf\checkmark$\checkmark^\checkmark\\checkmarkc\checkmarki\checkmarkr\checkmarkc\checkmark$\checkmarkr\checkmark$\checkmark^\checkmark\\checkmarkc\checkmarki\checkmarkr\checkmarkc\checkmark$\checkmarka\checkmark$\checkmark^\checkmark\\checkmarkc\checkmarki\checkmarkr\checkmarkc\checkmark$\checkmarkc\checkmark$\checkmark^\checkmark\\checkmarkc\checkmarki\checkmarkr\checkmarkc\checkmark$\checkmark{\checkmark$\checkmark^\checkmark\\checkmarkc\checkmarki\checkmarkr\checkmarkc\checkmark$\checkmark2\checkmark$\checkmark^\checkmark\\checkmarkc\checkmarki\checkmarkr\checkmarkc\checkmark$\checkmark\\checkmark$\checkmark^\checkmark\\checkmarkc\checkmarki\checkmarkr\checkmarkc\checkmark$\checkmarkp\checkmark$\checkmark^\checkmark\\checkmarkc\checkmarki\checkmarkr\checkmarkc\checkmark$\checkmarki\checkmark$\checkmark^\checkmark\\checkmarkc\checkmarki\checkmarkr\checkmarkc\checkmark$\checkmark}\checkmark$\checkmark^\checkmark\\checkmarkc\checkmarki\checkmarkr\checkmarkc\checkmark$\checkmark{\checkmark$\checkmark^\checkmark\\checkmarkc\checkmarki\checkmarkr\checkmarkc\checkmark$\checkmarkP\checkmark$\checkmark^\checkmark\\checkmarkc\checkmarki\checkmarkr\checkmarkc\checkmark$\checkmark_\checkmark$\checkmark^\checkmark\\checkmarkc\checkmarki\checkmarkr\checkmarkc\checkmark$\checkmarkh\checkmark$\checkmark^\checkmark\\checkmarkc\checkmarki\checkmarkr\checkmarkc\checkmark$\checkmark}\checkmark$\checkmark^\checkmark\\checkmarkc\checkmarki\checkmarkr\checkmarkc\checkmark$\checkmark \checkmark$\checkmark^\checkmark\\checkmarkc\checkmarki\checkmarkr\checkmarkc\checkmark$\checkmark=\checkmark$\checkmark^\checkmark\\checkmarkc\checkmarki\checkmarkr\checkmarkc\checkmark$\checkmark \checkmark$\checkmark^\checkmark\\checkmarkc\checkmarki\checkmarkr\checkmarkc\checkmark$\checkmark1\checkmark$\checkmark^\checkmark\\checkmarkc\checkmarki\checkmarkr\checkmarkc\checkmark$\checkmark \checkmark$\checkmark^\checkmark\\checkmarkc\checkmarki\checkmarkr\checkmarkc\checkmark$\checkmark\\checkmark$\checkmark^\checkmark\\checkmarkc\checkmarki\checkmarkr\checkmarkc\checkmark$\checkmarkt\checkmark$\checkmark^\checkmark\\checkmarkc\checkmarki\checkmarkr\checkmarkc\checkmark$\checkmarki\checkmark$\checkmark^\checkmark\\checkmarkc\checkmarki\checkmarkr\checkmarkc\checkmark$\checkmarkm\checkmark$\checkmark^\checkmark\\checkmarkc\checkmarki\checkmarkr\checkmarkc\checkmark$\checkmarke\checkmark$\checkmark^\checkmark\\checkmarkc\checkmarki\checkmarkr\checkmarkc\checkmark$\checkmarks\checkmark$\checkmark^\checkmark\\checkmarkc\checkmarki\checkmarkr\checkmarkc\checkmark$\checkmark \checkmark$\checkmark^\checkmark\\checkmarkc\checkmarki\checkmarkr\checkmarkc\checkmark$\checkmark\\checkmark$\checkmark^\checkmark\\checkmarkc\checkmarki\checkmarkr\checkmarkc\checkmark$\checkmarkf\checkmark$\checkmark^\checkmark\\checkmarkc\checkmarki\checkmarkr\checkmarkc\checkmark$\checkmarkr\checkmark$\checkmark^\checkmark\\checkmarkc\checkmarki\checkmarkr\checkmarkc\checkmark$\checkmarka\checkmark$\checkmark^\checkmark\\checkmarkc\checkmarki\checkmarkr\checkmarkc\checkmark$\checkmarkc\checkmark$\checkmark^\checkmark\\checkmarkc\checkmarki\checkmarkr\checkmarkc\checkmark$\checkmark{\checkmark$\checkmark^\checkmark\\checkmarkc\checkmarki\checkmarkr\checkmarkc\checkmark$\checkmark2\checkmark$\checkmark^\checkmark\\checkmarkc\checkmarki\checkmarkr\checkmarkc\checkmark$\checkmark\\checkmark$\checkmark^\checkmark\\checkmarkc\checkmarki\checkmarkr\checkmarkc\checkmark$\checkmarkp\checkmark$\checkmark^\checkmark\\checkmarkc\checkmarki\checkmarkr\checkmarkc\checkmark$\checkmarki\checkmark$\checkmark^\checkmark\\checkmarkc\checkmarki\checkmarkr\checkmarkc\checkmark$\checkmark}\checkmark$\checkmark^\checkmark\\checkmarkc\checkmarki\checkmarkr\checkmarkc\checkmark$\checkmark{\checkmark$\checkmark^\checkmark\\checkmarkc\checkmarki\checkmarkr\checkmarkc\checkmark$\checkmark0\checkmark$\checkmark^\checkmark\\checkmarkc\checkmarki\checkmarkr\checkmarkc\checkmark$\checkmark{\checkmark$\checkmark^\checkmark\\checkmarkc\checkmarki\checkmarkr\checkmarkc\checkmark$\checkmark,\checkmark$\checkmark^\checkmark\\checkmarkc\checkmarki\checkmarkr\checkmarkc\checkmark$\checkmark}\checkmark$\checkmark^\checkmark\\checkmarkc\checkmarki\checkmarkr\checkmarkc\checkmark$\checkmark0\checkmark$\checkmark^\checkmark\\checkmarkc\checkmarki\checkmarkr\checkmarkc\checkmark$\checkmark0\checkmark$\checkmark^\checkmark\\checkmarkc\checkmarki\checkmarkr\checkmarkc\checkmark$\checkmark5\checkmark$\checkmark^\checkmark\\checkmarkc\checkmarki\checkmarkr\checkmarkc\checkmark$\checkmark}\checkmark$\checkmark^\checkmark\\checkmarkc\checkmarki\checkmarkr\checkmarkc\checkmark$\checkmark \checkmark$\checkmark^\checkmark\\checkmarkc\checkmarki\checkmarkr\checkmarkc\checkmark$\checkmark=\checkmark$\checkmark^\checkmark\\checkmarkc\checkmarki\checkmarkr\checkmarkc\checkmark$\checkmark \checkmark$\checkmark^\checkmark\\checkmarkc\checkmarki\checkmarkr\checkmarkc\checkmark$\checkmark1\checkmark$\checkmark^\checkmark\\checkmarkc\checkmarki\checkmarkr\checkmarkc\checkmark$\checkmark\\checkmark$\checkmark^\checkmark\\checkmarkc\checkmarki\checkmarkr\checkmarkc\checkmark$\checkmark,\checkmark$\checkmark^\checkmark\\checkmarkc\checkmarki\checkmarkr\checkmarkc\checkmark$\checkmark2\checkmark$\checkmark^\checkmark\\checkmarkc\checkmarki\checkmarkr\checkmarkc\checkmark$\checkmark5\checkmark$\checkmark^\checkmark\\checkmarkc\checkmarki\checkmarkr\checkmarkc\checkmark$\checkmark6\checkmark$\checkmark^\checkmark\\checkmarkc\checkmarki\checkmarkr\checkmarkc\checkmark$\checkmark \checkmark$\checkmark^\checkmark\\checkmarkc\checkmarki\checkmarkr\checkmarkc\checkmark$\checkmark\\checkmark$\checkmark^\checkmark\\checkmarkc\checkmarki\checkmarkr\checkmarkc\checkmark$\checkmarkt\checkmark$\checkmark^\checkmark\\checkmarkc\checkmarki\checkmarkr\checkmarkc\checkmark$\checkmarke\checkmark$\checkmark^\checkmark\\checkmarkc\checkmarki\checkmarkr\checkmarkc\checkmark$\checkmarkx\checkmark$\checkmark^\checkmark\\checkmarkc\checkmarki\checkmarkr\checkmarkc\checkmark$\checkmarkt\checkmark$\checkmark^\checkmark\\checkmarkc\checkmarki\checkmarkr\checkmarkc\checkmark$\checkmark{\checkmark$\checkmark^\checkmark\\checkmarkc\checkmarki\checkmarkr\checkmarkc\checkmark$\checkmark \checkmark$\checkmark^\checkmark\\checkmarkc\checkmarki\checkmarkr\checkmarkc\checkmark$\checkmarkr\checkmark$\checkmark^\checkmark\\checkmarkc\checkmarki\checkmarkr\checkmarkc\checkmark$\checkmarka\checkmark$\checkmark^\checkmark\\checkmarkc\checkmarki\checkmarkr\checkmarkc\checkmark$\checkmarkd\checkmark$\checkmark^\checkmark\\checkmarkc\checkmarki\checkmarkr\checkmarkc\checkmark$\checkmark/\checkmark$\checkmark^\checkmark\\checkmarkc\checkmarki\checkmarkr\checkmarkc\checkmark$\checkmarks\checkmark$\checkmark^\checkmark\\checkmarkc\checkmarki\checkmarkr\checkmarkc\checkmark$\checkmark}\checkmark$\checkmark^\checkmark\\checkmarkc\checkmarki\checkmarkr\checkmarkc\checkmark$\checkmark^\checkmark$\checkmark^\checkmark\\checkmarkc\checkmarki\checkmarkr\checkmarkc\checkmark$\checkmark2\checkmark$\checkmark^\checkmark\\checkmarkc\checkmarki\checkmarkr\checkmarkc\checkmark$\checkmark
\checkmark$\checkmark^\checkmark\\checkmarkc\checkmarki\checkmarkr\checkmarkc\checkmark$\checkmark\\checkmark$\checkmark^\checkmark\\checkmarkc\checkmarki\checkmarkr\checkmarkc\checkmark$\checkmarke\checkmark$\checkmark^\checkmark\\checkmarkc\checkmarki\checkmarkr\checkmarkc\checkmark$\checkmarkn\checkmark$\checkmark^\checkmark\\checkmarkc\checkmarki\checkmarkr\checkmarkc\checkmark$\checkmarkd\checkmark$\checkmark^\checkmark\\checkmarkc\checkmarki\checkmarkr\checkmarkc\checkmark$\checkmark{\checkmark$\checkmark^\checkmark\\checkmarkc\checkmarki\checkmarkr\checkmarkc\checkmark$\checkmarke\checkmark$\checkmark^\checkmark\\checkmarkc\checkmarki\checkmarkr\checkmarkc\checkmark$\checkmarkq\checkmark$\checkmark^\checkmark\\checkmarkc\checkmarki\checkmarkr\checkmarkc\checkmark$\checkmarku\checkmark$\checkmark^\checkmark\\checkmarkc\checkmarki\checkmarkr\checkmarkc\checkmark$\checkmarka\checkmark$\checkmark^\checkmark\\checkmarkc\checkmarki\checkmarkr\checkmarkc\checkmark$\checkmarkt\checkmark$\checkmark^\checkmark\\checkmarkc\checkmarki\checkmarkr\checkmarkc\checkmark$\checkmarki\checkmark$\checkmark^\checkmark\\checkmarkc\checkmarki\checkmarkr\checkmarkc\checkmark$\checkmarko\checkmark$\checkmark^\checkmark\\checkmarkc\checkmarki\checkmarkr\checkmarkc\checkmark$\checkmarkn\checkmark$\checkmark^\checkmark\\checkmarkc\checkmarki\checkmarkr\checkmarkc\checkmark$\checkmark}\checkmark$\checkmark^\checkmark\\checkmarkc\checkmarki\checkmarkr\checkmarkc\checkmark$\checkmark
\checkmark$\checkmark^\checkmark\\checkmarkc\checkmarki\checkmarkr\checkmarkc\checkmark$\checkmark
\checkmark$\checkmark^\checkmark\\checkmarkc\checkmarki\checkmarkr\checkmarkc\checkmark$\checkmarkL\checkmark$\checkmark^\checkmark\\checkmarkc\checkmarki\checkmarkr\checkmarkc\checkmark$\checkmarke\checkmark$\checkmark^\checkmark\\checkmarkc\checkmarki\checkmarkr\checkmarkc\checkmark$\checkmark \checkmark$\checkmark^\checkmark\\checkmarkc\checkmarki\checkmarkr\checkmarkc\checkmark$\checkmarkc\checkmark$\checkmark^\checkmark\\checkmarkc\checkmarki\checkmarkr\checkmarkc\checkmark$\checkmarko\checkmark$\checkmark^\checkmark\\checkmarkc\checkmarki\checkmarkr\checkmarkc\checkmark$\checkmarku\checkmark$\checkmark^\checkmark\\checkmarkc\checkmarki\checkmarkr\checkmarkc\checkmark$\checkmarkp\checkmark$\checkmark^\checkmark\\checkmarkc\checkmarki\checkmarkr\checkmarkc\checkmark$\checkmarkl\checkmark$\checkmark^\checkmark\\checkmarkc\checkmarki\checkmarkr\checkmarkc\checkmark$\checkmarke\checkmark$\checkmark^\checkmark\\checkmarkc\checkmarki\checkmarkr\checkmarkc\checkmark$\checkmark \checkmark$\checkmark^\checkmark\\checkmarkc\checkmarki\checkmarkr\checkmarkc\checkmark$\checkmarkd\checkmark$\checkmark^\checkmark\\checkmarkc\checkmarki\checkmarkr\checkmarkc\checkmark$\checkmark'\checkmark$\checkmark^\checkmark\\checkmarkc\checkmarki\checkmarkr\checkmarkc\checkmark$\checkmarka\checkmark$\checkmark^\checkmark\\checkmarkc\checkmarki\checkmarkr\checkmarkc\checkmark$\checkmarkc\checkmark$\checkmark^\checkmark\\checkmarkc\checkmarki\checkmarkr\checkmarkc\checkmark$\checkmarkc\checkmark$\checkmark^\checkmark\\checkmarkc\checkmarki\checkmarkr\checkmarkc\checkmark$\checkmarké\checkmark$\checkmark^\checkmark\\checkmarkc\checkmarki\checkmarkr\checkmarkc\checkmark$\checkmarkl\checkmark$\checkmark^\checkmark\\checkmarkc\checkmarki\checkmarkr\checkmarkc\checkmark$\checkmarké\checkmark$\checkmark^\checkmark\\checkmarkc\checkmarki\checkmarkr\checkmarkc\checkmark$\checkmarkr\checkmark$\checkmark^\checkmark\\checkmarkc\checkmarki\checkmarkr\checkmarkc\checkmark$\checkmarka\checkmark$\checkmark^\checkmark\\checkmarkc\checkmarki\checkmarkr\checkmarkc\checkmark$\checkmarkt\checkmark$\checkmark^\checkmark\\checkmarkc\checkmarki\checkmarkr\checkmarkc\checkmark$\checkmarki\checkmark$\checkmark^\checkmark\\checkmarkc\checkmarki\checkmarkr\checkmarkc\checkmark$\checkmarko\checkmark$\checkmark^\checkmark\\checkmarkc\checkmarki\checkmarkr\checkmarkc\checkmark$\checkmarkn\checkmark$\checkmark^\checkmark\\checkmarkc\checkmarki\checkmarkr\checkmarkc\checkmark$\checkmark \checkmark$\checkmark^\checkmark\\checkmarkc\checkmarki\checkmarkr\checkmarkc\checkmark$\checkmark(\checkmark$\checkmark^\checkmark\\checkmarkc\checkmarki\checkmarkr\checkmarkc\checkmark$\checkmarkP\checkmark$\checkmark^\checkmark\\checkmarkc\checkmarki\checkmarkr\checkmarkc\checkmark$\checkmarkr\checkmark$\checkmark^\checkmark\\checkmarkc\checkmarki\checkmarkr\checkmarkc\checkmark$\checkmarki\checkmark$\checkmark^\checkmark\\checkmarkc\checkmarki\checkmarkr\checkmarkc\checkmark$\checkmarkn\checkmark$\checkmark^\checkmark\\checkmarkc\checkmarki\checkmarkr\checkmarkc\checkmark$\checkmarkc\checkmark$\checkmark^\checkmark\\checkmarkc\checkmarki\checkmarkr\checkmarkc\checkmark$\checkmarki\checkmark$\checkmark^\checkmark\\checkmarkc\checkmarki\checkmarkr\checkmarkc\checkmark$\checkmarkp\checkmark$\checkmark^\checkmark\\checkmarkc\checkmarki\checkmarkr\checkmarkc\checkmark$\checkmarke\checkmark$\checkmark^\checkmark\\checkmarkc\checkmarki\checkmarkr\checkmarkc\checkmark$\checkmark \checkmark$\checkmark^\checkmark\\checkmarkc\checkmarki\checkmarkr\checkmarkc\checkmark$\checkmarkF\checkmark$\checkmark^\checkmark\\checkmarkc\checkmarki\checkmarkr\checkmarkc\checkmark$\checkmarko\checkmark$\checkmark^\checkmark\\checkmarkc\checkmarki\checkmarkr\checkmarkc\checkmark$\checkmarkn\checkmark$\checkmark^\checkmark\\checkmarkc\checkmarki\checkmarkr\checkmarkc\checkmark$\checkmarkd\checkmark$\checkmark^\checkmark\\checkmarkc\checkmarki\checkmarkr\checkmarkc\checkmark$\checkmarka\checkmark$\checkmark^\checkmark\\checkmarkc\checkmarki\checkmarkr\checkmarkc\checkmark$\checkmarkm\checkmark$\checkmark^\checkmark\\checkmarkc\checkmarki\checkmarkr\checkmarkc\checkmark$\checkmarke\checkmark$\checkmark^\checkmark\\checkmarkc\checkmarki\checkmarkr\checkmarkc\checkmark$\checkmarkn\checkmark$\checkmark^\checkmark\\checkmarkc\checkmarki\checkmarkr\checkmarkc\checkmark$\checkmarkt\checkmark$\checkmark^\checkmark\\checkmarkc\checkmarki\checkmarkr\checkmarkc\checkmark$\checkmarka\checkmark$\checkmark^\checkmark\\checkmarkc\checkmarki\checkmarkr\checkmarkc\checkmark$\checkmarkl\checkmark$\checkmark^\checkmark\\checkmarkc\checkmarki\checkmarkr\checkmarkc\checkmark$\checkmark \checkmark$\checkmark^\checkmark\\checkmarkc\checkmarki\checkmarkr\checkmarkc\checkmark$\checkmarkd\checkmark$\checkmark^\checkmark\\checkmarkc\checkmarki\checkmarkr\checkmarkc\checkmark$\checkmarke\checkmark$\checkmark^\checkmark\\checkmarkc\checkmarki\checkmarkr\checkmarkc\checkmark$\checkmark \checkmark$\checkmark^\checkmark\\checkmarkc\checkmarki\checkmarkr\checkmarkc\checkmark$\checkmarkl\checkmark$\checkmark^\checkmark\\checkmarkc\checkmarki\checkmarkr\checkmarkc\checkmark$\checkmarka\checkmark$\checkmark^\checkmark\\checkmarkc\checkmarki\checkmarkr\checkmarkc\checkmark$\checkmark \checkmark$\checkmark^\checkmark\\checkmarkc\checkmarki\checkmarkr\checkmarkc\checkmark$\checkmarkD\checkmark$\checkmark^\checkmark\\checkmarkc\checkmarki\checkmarkr\checkmarkc\checkmark$\checkmarky\checkmark$\checkmark^\checkmark\\checkmarkc\checkmarki\checkmarkr\checkmarkc\checkmark$\checkmarkn\checkmark$\checkmark^\checkmark\\checkmarkc\checkmarki\checkmarkr\checkmarkc\checkmark$\checkmarka\checkmark$\checkmark^\checkmark\\checkmarkc\checkmarki\checkmarkr\checkmarkc\checkmark$\checkmarkm\checkmark$\checkmark^\checkmark\\checkmarkc\checkmarki\checkmarkr\checkmarkc\checkmark$\checkmarki\checkmark$\checkmark^\checkmark\\checkmarkc\checkmarki\checkmarkr\checkmarkc\checkmark$\checkmarkq\checkmark$\checkmark^\checkmark\\checkmarkc\checkmarki\checkmarkr\checkmarkc\checkmark$\checkmarku\checkmark$\checkmark^\checkmark\\checkmarkc\checkmarki\checkmarkr\checkmarkc\checkmark$\checkmarke\checkmark$\checkmark^\checkmark\\checkmarkc\checkmarki\checkmarkr\checkmarkc\checkmark$\checkmark)\checkmark$\checkmark^\checkmark\\checkmarkc\checkmarki\checkmarkr\checkmarkc\checkmark$\checkmark \checkmark$\checkmark^\checkmark\\checkmarkc\checkmarki\checkmarkr\checkmarkc\checkmark$\checkmark:\checkmark$\checkmark^\checkmark\\checkmarkc\checkmarki\checkmarkr\checkmarkc\checkmark$\checkmark
\checkmark$\checkmark^\checkmark\\checkmarkc\checkmarki\checkmarkr\checkmarkc\checkmark$\checkmark\\checkmark$\checkmark^\checkmark\\checkmarkc\checkmarki\checkmarkr\checkmarkc\checkmark$\checkmarkb\checkmark$\checkmark^\checkmark\\checkmarkc\checkmarki\checkmarkr\checkmarkc\checkmark$\checkmarke\checkmark$\checkmark^\checkmark\\checkmarkc\checkmarki\checkmarkr\checkmarkc\checkmark$\checkmarkg\checkmark$\checkmark^\checkmark\\checkmarkc\checkmarki\checkmarkr\checkmarkc\checkmark$\checkmarki\checkmark$\checkmark^\checkmark\\checkmarkc\checkmarki\checkmarkr\checkmarkc\checkmark$\checkmarkn\checkmark$\checkmark^\checkmark\\checkmarkc\checkmarki\checkmarkr\checkmarkc\checkmark$\checkmark{\checkmark$\checkmark^\checkmark\\checkmarkc\checkmarki\checkmarkr\checkmarkc\checkmark$\checkmarke\checkmark$\checkmark^\checkmark\\checkmarkc\checkmarki\checkmarkr\checkmarkc\checkmark$\checkmarkq\checkmark$\checkmark^\checkmark\\checkmarkc\checkmarki\checkmarkr\checkmarkc\checkmark$\checkmarku\checkmark$\checkmark^\checkmark\\checkmarkc\checkmarki\checkmarkr\checkmarkc\checkmark$\checkmarka\checkmark$\checkmark^\checkmark\\checkmarkc\checkmarki\checkmarkr\checkmarkc\checkmark$\checkmarkt\checkmark$\checkmark^\checkmark\\checkmarkc\checkmarki\checkmarkr\checkmarkc\checkmark$\checkmarki\checkmark$\checkmark^\checkmark\\checkmarkc\checkmarki\checkmarkr\checkmarkc\checkmark$\checkmarko\checkmark$\checkmark^\checkmark\\checkmarkc\checkmarki\checkmarkr\checkmarkc\checkmark$\checkmarkn\checkmark$\checkmark^\checkmark\\checkmarkc\checkmarki\checkmarkr\checkmarkc\checkmark$\checkmark}\checkmark$\checkmark^\checkmark\\checkmarkc\checkmarki\checkmarkr\checkmarkc\checkmark$\checkmark
\checkmark$\checkmark^\checkmark\\checkmarkc\checkmarki\checkmarkr\checkmarkc\checkmark$\checkmark \checkmark$\checkmark^\checkmark\\checkmarkc\checkmarki\checkmarkr\checkmarkc\checkmark$\checkmark \checkmark$\checkmark^\checkmark\\checkmarkc\checkmarki\checkmarkr\checkmarkc\checkmark$\checkmark \checkmark$\checkmark^\checkmark\\checkmarkc\checkmarki\checkmarkr\checkmarkc\checkmark$\checkmark \checkmark$\checkmark^\checkmark\\checkmarkc\checkmarki\checkmarkr\checkmarkc\checkmark$\checkmarkC\checkmark$\checkmark^\checkmark\\checkmarkc\checkmarki\checkmarkr\checkmarkc\checkmark$\checkmark_\checkmark$\checkmark^\checkmark\\checkmarkc\checkmarki\checkmarkr\checkmarkc\checkmark$\checkmarkJ\checkmark$\checkmark^\checkmark\\checkmarkc\checkmarki\checkmarkr\checkmarkc\checkmark$\checkmark \checkmark$\checkmark^\checkmark\\checkmarkc\checkmarki\checkmarkr\checkmarkc\checkmark$\checkmark=\checkmark$\checkmark^\checkmark\\checkmarkc\checkmarki\checkmarkr\checkmarkc\checkmark$\checkmark \checkmark$\checkmark^\checkmark\\checkmarkc\checkmarki\checkmarkr\checkmarkc\checkmark$\checkmarkJ\checkmark$\checkmark^\checkmark\\checkmarkc\checkmarki\checkmarkr\checkmarkc\checkmark$\checkmark_\checkmark$\checkmark^\checkmark\\checkmarkc\checkmarki\checkmarkr\checkmarkc\checkmark$\checkmark{\checkmark$\checkmark^\checkmark\\checkmarkc\checkmarki\checkmarkr\checkmarkc\checkmark$\checkmarkt\checkmark$\checkmark^\checkmark\\checkmarkc\checkmarki\checkmarkr\checkmarkc\checkmark$\checkmarko\checkmark$\checkmark^\checkmark\\checkmarkc\checkmarki\checkmarkr\checkmarkc\checkmark$\checkmarkt\checkmark$\checkmark^\checkmark\\checkmarkc\checkmarki\checkmarkr\checkmarkc\checkmark$\checkmark}\checkmark$\checkmark^\checkmark\\checkmarkc\checkmarki\checkmarkr\checkmarkc\checkmark$\checkmark \checkmark$\checkmark^\checkmark\\checkmarkc\checkmarki\checkmarkr\checkmarkc\checkmark$\checkmark\\checkmark$\checkmark^\checkmark\\checkmarkc\checkmarki\checkmarkr\checkmarkc\checkmark$\checkmarkc\checkmark$\checkmark^\checkmark\\checkmarkc\checkmarki\checkmarkr\checkmarkc\checkmark$\checkmarkd\checkmark$\checkmark^\checkmark\\checkmarkc\checkmarki\checkmarkr\checkmarkc\checkmark$\checkmarko\checkmark$\checkmark^\checkmark\\checkmarkc\checkmarki\checkmarkr\checkmarkc\checkmark$\checkmarkt\checkmark$\checkmark^\checkmark\\checkmarkc\checkmarki\checkmarkr\checkmarkc\checkmark$\checkmark \checkmark$\checkmark^\checkmark\\checkmarkc\checkmarki\checkmarkr\checkmarkc\checkmark$\checkmark\\checkmark$\checkmark^\checkmark\\checkmarkc\checkmarki\checkmarkr\checkmarkc\checkmark$\checkmarka\checkmark$\checkmark^\checkmark\\checkmarkc\checkmarki\checkmarkr\checkmarkc\checkmark$\checkmarkl\checkmark$\checkmark^\checkmark\\checkmarkc\checkmarki\checkmarkr\checkmarkc\checkmark$\checkmarkp\checkmark$\checkmark^\checkmark\\checkmarkc\checkmarki\checkmarkr\checkmarkc\checkmark$\checkmarkh\checkmark$\checkmark^\checkmark\\checkmarkc\checkmarki\checkmarkr\checkmarkc\checkmark$\checkmarka\checkmark$\checkmark^\checkmark\\checkmarkc\checkmarki\checkmarkr\checkmarkc\checkmark$\checkmark \checkmark$\checkmark^\checkmark\\checkmarkc\checkmarki\checkmarkr\checkmarkc\checkmark$\checkmark=\checkmark$\checkmark^\checkmark\\checkmarkc\checkmarki\checkmarkr\checkmarkc\checkmark$\checkmark \checkmark$\checkmark^\checkmark\\checkmarkc\checkmarki\checkmarkr\checkmarkc\checkmark$\checkmark3\checkmark$\checkmark^\checkmark\\checkmarkc\checkmarki\checkmarkr\checkmarkc\checkmark$\checkmark{\checkmark$\checkmark^\checkmark\\checkmarkc\checkmarki\checkmarkr\checkmarkc\checkmark$\checkmark,\checkmark$\checkmark^\checkmark\\checkmarkc\checkmarki\checkmarkr\checkmarkc\checkmark$\checkmark}\checkmark$\checkmark^\checkmark\\checkmarkc\checkmarki\checkmarkr\checkmarkc\checkmark$\checkmark1\checkmark$\checkmark^\checkmark\\checkmarkc\checkmarki\checkmarkr\checkmarkc\checkmark$\checkmark5\checkmark$\checkmark^\checkmark\\checkmarkc\checkmarki\checkmarkr\checkmarkc\checkmark$\checkmark \checkmark$\checkmark^\checkmark\\checkmarkc\checkmarki\checkmarkr\checkmarkc\checkmark$\checkmark\\checkmark$\checkmark^\checkmark\\checkmarkc\checkmarki\checkmarkr\checkmarkc\checkmark$\checkmarkt\checkmark$\checkmark^\checkmark\\checkmarkc\checkmarki\checkmarkr\checkmarkc\checkmark$\checkmarki\checkmark$\checkmark^\checkmark\\checkmarkc\checkmarki\checkmarkr\checkmarkc\checkmark$\checkmarkm\checkmark$\checkmark^\checkmark\\checkmarkc\checkmarki\checkmarkr\checkmarkc\checkmark$\checkmarke\checkmark$\checkmark^\checkmark\\checkmarkc\checkmarki\checkmarkr\checkmarkc\checkmark$\checkmarks\checkmark$\checkmark^\checkmark\\checkmarkc\checkmarki\checkmarkr\checkmarkc\checkmark$\checkmark \checkmark$\checkmark^\checkmark\\checkmarkc\checkmarki\checkmarkr\checkmarkc\checkmark$\checkmark1\checkmark$\checkmark^\checkmark\\checkmarkc\checkmarki\checkmarkr\checkmarkc\checkmark$\checkmark0\checkmark$\checkmark^\checkmark\\checkmarkc\checkmarki\checkmarkr\checkmarkc\checkmark$\checkmark^\checkmark$\checkmark^\checkmark\\checkmarkc\checkmarki\checkmarkr\checkmarkc\checkmark$\checkmark{\checkmark$\checkmark^\checkmark\\checkmarkc\checkmarki\checkmarkr\checkmarkc\checkmark$\checkmark-\checkmark$\checkmark^\checkmark\\checkmarkc\checkmarki\checkmarkr\checkmarkc\checkmark$\checkmark5\checkmark$\checkmark^\checkmark\\checkmarkc\checkmarki\checkmarkr\checkmarkc\checkmark$\checkmark}\checkmark$\checkmark^\checkmark\\checkmarkc\checkmarki\checkmarkr\checkmarkc\checkmark$\checkmark \checkmark$\checkmark^\checkmark\\checkmarkc\checkmarki\checkmarkr\checkmarkc\checkmark$\checkmark\\checkmark$\checkmark^\checkmark\\checkmarkc\checkmarki\checkmarkr\checkmarkc\checkmark$\checkmarkt\checkmark$\checkmark^\checkmark\\checkmarkc\checkmarki\checkmarkr\checkmarkc\checkmark$\checkmarki\checkmark$\checkmark^\checkmark\\checkmarkc\checkmarki\checkmarkr\checkmarkc\checkmark$\checkmarkm\checkmark$\checkmark^\checkmark\\checkmarkc\checkmarki\checkmarkr\checkmarkc\checkmark$\checkmarke\checkmark$\checkmark^\checkmark\\checkmarkc\checkmarki\checkmarkr\checkmarkc\checkmark$\checkmarks\checkmark$\checkmark^\checkmark\\checkmarkc\checkmarki\checkmarkr\checkmarkc\checkmark$\checkmark \checkmark$\checkmark^\checkmark\\checkmarkc\checkmarki\checkmarkr\checkmarkc\checkmark$\checkmark1\checkmark$\checkmark^\checkmark\\checkmarkc\checkmarki\checkmarkr\checkmarkc\checkmark$\checkmark\\checkmark$\checkmark^\checkmark\\checkmarkc\checkmarki\checkmarkr\checkmarkc\checkmark$\checkmark,\checkmark$\checkmark^\checkmark\\checkmarkc\checkmarki\checkmarkr\checkmarkc\checkmark$\checkmark2\checkmark$\checkmark^\checkmark\\checkmarkc\checkmarki\checkmarkr\checkmarkc\checkmark$\checkmark5\checkmark$\checkmark^\checkmark\\checkmarkc\checkmarki\checkmarkr\checkmarkc\checkmark$\checkmark6\checkmark$\checkmark^\checkmark\\checkmarkc\checkmarki\checkmarkr\checkmarkc\checkmark$\checkmark \checkmark$\checkmark^\checkmark\\checkmarkc\checkmarki\checkmarkr\checkmarkc\checkmark$\checkmark=\checkmark$\checkmark^\checkmark\\checkmarkc\checkmarki\checkmarkr\checkmarkc\checkmark$\checkmark \checkmark$\checkmark^\checkmark\\checkmarkc\checkmarki\checkmarkr\checkmarkc\checkmark$\checkmark\\checkmark$\checkmark^\checkmark\\checkmarkc\checkmarki\checkmarkr\checkmarkc\checkmark$\checkmarkb\checkmark$\checkmark^\checkmark\\checkmarkc\checkmarki\checkmarkr\checkmarkc\checkmark$\checkmarko\checkmark$\checkmark^\checkmark\\checkmarkc\checkmarki\checkmarkr\checkmarkc\checkmark$\checkmarkx\checkmark$\checkmark^\checkmark\\checkmarkc\checkmarki\checkmarkr\checkmarkc\checkmark$\checkmarke\checkmark$\checkmark^\checkmark\\checkmarkc\checkmarki\checkmarkr\checkmarkc\checkmark$\checkmarkd\checkmark$\checkmark^\checkmark\\checkmarkc\checkmarki\checkmarkr\checkmarkc\checkmark$\checkmark{\checkmark$\checkmark^\checkmark\\checkmarkc\checkmarki\checkmarkr\checkmarkc\checkmark$\checkmark0\checkmark$\checkmark^\checkmark\\checkmarkc\checkmarki\checkmarkr\checkmarkc\checkmark$\checkmark{\checkmark$\checkmark^\checkmark\\checkmarkc\checkmarki\checkmarkr\checkmarkc\checkmark$\checkmark,\checkmark$\checkmark^\checkmark\\checkmarkc\checkmarki\checkmarkr\checkmarkc\checkmark$\checkmark}\checkmark$\checkmark^\checkmark\\checkmarkc\checkmarki\checkmarkr\checkmarkc\checkmark$\checkmark0\checkmark$\checkmark^\checkmark\\checkmarkc\checkmarki\checkmarkr\checkmarkc\checkmark$\checkmark4\checkmark$\checkmark^\checkmark\\checkmarkc\checkmarki\checkmarkr\checkmarkc\checkmark$\checkmark0\checkmark$\checkmark^\checkmark\\checkmarkc\checkmarki\checkmarkr\checkmarkc\checkmark$\checkmark \checkmark$\checkmark^\checkmark\\checkmarkc\checkmarki\checkmarkr\checkmarkc\checkmark$\checkmark\\checkmark$\checkmark^\checkmark\\checkmarkc\checkmarki\checkmarkr\checkmarkc\checkmark$\checkmarkt\checkmark$\checkmark^\checkmark\\checkmarkc\checkmarki\checkmarkr\checkmarkc\checkmark$\checkmarke\checkmark$\checkmark^\checkmark\\checkmarkc\checkmarki\checkmarkr\checkmarkc\checkmark$\checkmarkx\checkmark$\checkmark^\checkmark\\checkmarkc\checkmarki\checkmarkr\checkmarkc\checkmark$\checkmarkt\checkmark$\checkmark^\checkmark\\checkmarkc\checkmarki\checkmarkr\checkmarkc\checkmark$\checkmark{\checkmark$\checkmark^\checkmark\\checkmarkc\checkmarki\checkmarkr\checkmarkc\checkmark$\checkmark \checkmark$\checkmark^\checkmark\\checkmarkc\checkmarki\checkmarkr\checkmarkc\checkmark$\checkmarkN\checkmark$\checkmark^\checkmark\\checkmarkc\checkmarki\checkmarkr\checkmarkc\checkmark$\checkmark$\checkmark$\checkmark^\checkmark\\checkmarkc\checkmarki\checkmarkr\checkmarkc\checkmark$\checkmark\\checkmark$\checkmark^\checkmark\\checkmarkc\checkmarki\checkmarkr\checkmarkc\checkmark$\checkmarkc\checkmark$\checkmark^\checkmark\\checkmarkc\checkmarki\checkmarkr\checkmarkc\checkmark$\checkmarkd\checkmark$\checkmark^\checkmark\\checkmarkc\checkmarki\checkmarkr\checkmarkc\checkmark$\checkmarko\checkmark$\checkmark^\checkmark\\checkmarkc\checkmarki\checkmarkr\checkmarkc\checkmark$\checkmarkt\checkmark$\checkmark^\checkmark\\checkmarkc\checkmarki\checkmarkr\checkmarkc\checkmark$\checkmark$\checkmark$\checkmark^\checkmark\\checkmarkc\checkmarki\checkmarkr\checkmarkc\checkmark$\checkmarkm\checkmark$\checkmark^\checkmark\\checkmarkc\checkmarki\checkmarkr\checkmarkc\checkmark$\checkmark}\checkmark$\checkmark^\checkmark\\checkmarkc\checkmarki\checkmarkr\checkmarkc\checkmark$\checkmark}\checkmark$\checkmark^\checkmark\\checkmarkc\checkmarki\checkmarkr\checkmarkc\checkmark$\checkmark
\checkmark$\checkmark^\checkmark\\checkmarkc\checkmarki\checkmarkr\checkmarkc\checkmark$\checkmark \checkmark$\checkmark^\checkmark\\checkmarkc\checkmarki\checkmarkr\checkmarkc\checkmark$\checkmark \checkmark$\checkmark^\checkmark\\checkmarkc\checkmarki\checkmarkr\checkmarkc\checkmark$\checkmark \checkmark$\checkmark^\checkmark\\checkmarkc\checkmarki\checkmarkr\checkmarkc\checkmark$\checkmark \checkmark$\checkmark^\checkmark\\checkmarkc\checkmarki\checkmarkr\checkmarkc\checkmark$\checkmark\\checkmark$\checkmark^\checkmark\\checkmarkc\checkmarki\checkmarkr\checkmarkc\checkmark$\checkmarkl\checkmark$\checkmark^\checkmark\\checkmarkc\checkmarki\checkmarkr\checkmarkc\checkmark$\checkmarka\checkmark$\checkmark^\checkmark\\checkmarkc\checkmarki\checkmarkr\checkmarkc\checkmark$\checkmarkb\checkmark$\checkmark^\checkmark\\checkmarkc\checkmarki\checkmarkr\checkmarkc\checkmark$\checkmarke\checkmark$\checkmark^\checkmark\\checkmarkc\checkmarki\checkmarkr\checkmarkc\checkmark$\checkmarkl\checkmark$\checkmark^\checkmark\\checkmarkc\checkmarki\checkmarkr\checkmarkc\checkmark$\checkmark{\checkmark$\checkmark^\checkmark\\checkmarkc\checkmarki\checkmarkr\checkmarkc\checkmark$\checkmarke\checkmark$\checkmark^\checkmark\\checkmarkc\checkmarki\checkmarkr\checkmarkc\checkmark$\checkmarkq\checkmark$\checkmark^\checkmark\\checkmarkc\checkmarki\checkmarkr\checkmarkc\checkmark$\checkmark:\checkmark$\checkmark^\checkmark\\checkmarkc\checkmarki\checkmarkr\checkmarkc\checkmark$\checkmarkc\checkmark$\checkmark^\checkmark\\checkmarkc\checkmarki\checkmarkr\checkmarkc\checkmark$\checkmarko\checkmark$\checkmark^\checkmark\\checkmarkc\checkmarki\checkmarkr\checkmarkc\checkmark$\checkmarku\checkmark$\checkmark^\checkmark\\checkmarkc\checkmarki\checkmarkr\checkmarkc\checkmark$\checkmarkp\checkmark$\checkmark^\checkmark\\checkmarkc\checkmarki\checkmarkr\checkmarkc\checkmark$\checkmarkl\checkmark$\checkmark^\checkmark\\checkmarkc\checkmarki\checkmarkr\checkmarkc\checkmark$\checkmarke\checkmark$\checkmark^\checkmark\\checkmarkc\checkmarki\checkmarkr\checkmarkc\checkmark$\checkmark_\checkmark$\checkmark^\checkmark\\checkmarkc\checkmarki\checkmarkr\checkmarkc\checkmark$\checkmarki\checkmark$\checkmark^\checkmark\\checkmarkc\checkmarki\checkmarkr\checkmarkc\checkmark$\checkmarkn\checkmark$\checkmark^\checkmark\\checkmarkc\checkmarki\checkmarkr\checkmarkc\checkmark$\checkmarke\checkmark$\checkmark^\checkmark\\checkmarkc\checkmarki\checkmarkr\checkmarkc\checkmark$\checkmarkr\checkmark$\checkmark^\checkmark\\checkmarkc\checkmarki\checkmarkr\checkmarkc\checkmark$\checkmarkt\checkmark$\checkmark^\checkmark\\checkmarkc\checkmarki\checkmarkr\checkmarkc\checkmark$\checkmarki\checkmark$\checkmark^\checkmark\\checkmarkc\checkmarki\checkmarkr\checkmarkc\checkmark$\checkmarke\checkmark$\checkmark^\checkmark\\checkmarkc\checkmarki\checkmarkr\checkmarkc\checkmark$\checkmark}\checkmark$\checkmark^\checkmark\\checkmarkc\checkmarki\checkmarkr\checkmarkc\checkmark$\checkmark
\checkmark$\checkmark^\checkmark\\checkmarkc\checkmarki\checkmarkr\checkmarkc\checkmark$\checkmark\\checkmark$\checkmark^\checkmark\\checkmarkc\checkmarki\checkmarkr\checkmarkc\checkmark$\checkmarke\checkmark$\checkmark^\checkmark\\checkmarkc\checkmarki\checkmarkr\checkmarkc\checkmark$\checkmarkn\checkmark$\checkmark^\checkmark\\checkmarkc\checkmarki\checkmarkr\checkmarkc\checkmark$\checkmarkd\checkmark$\checkmark^\checkmark\\checkmarkc\checkmarki\checkmarkr\checkmarkc\checkmark$\checkmark{\checkmark$\checkmark^\checkmark\\checkmarkc\checkmarki\checkmarkr\checkmarkc\checkmark$\checkmarke\checkmark$\checkmark^\checkmark\\checkmarkc\checkmarki\checkmarkr\checkmarkc\checkmark$\checkmarkq\checkmark$\checkmark^\checkmark\\checkmarkc\checkmarki\checkmarkr\checkmarkc\checkmark$\checkmarku\checkmark$\checkmark^\checkmark\\checkmarkc\checkmarki\checkmarkr\checkmarkc\checkmark$\checkmarka\checkmark$\checkmark^\checkmark\\checkmarkc\checkmarki\checkmarkr\checkmarkc\checkmark$\checkmarkt\checkmark$\checkmark^\checkmark\\checkmarkc\checkmarki\checkmarkr\checkmarkc\checkmark$\checkmarki\checkmark$\checkmark^\checkmark\\checkmarkc\checkmarki\checkmarkr\checkmarkc\checkmark$\checkmarko\checkmark$\checkmark^\checkmark\\checkmarkc\checkmarki\checkmarkr\checkmarkc\checkmark$\checkmarkn\checkmark$\checkmark^\checkmark\\checkmarkc\checkmarki\checkmarkr\checkmarkc\checkmark$\checkmark}\checkmark$\checkmark^\checkmark\\checkmarkc\checkmarki\checkmarkr\checkmarkc\checkmark$\checkmark
\checkmark$\checkmark^\checkmark\\checkmarkc\checkmarki\checkmarkr\checkmarkc\checkmark$\checkmark
\checkmark$\checkmark^\checkmark\\checkmarkc\checkmarki\checkmarkr\checkmarkc\checkmark$\checkmark\\checkmark$\checkmark^\checkmark\\checkmarkc\checkmarki\checkmarkr\checkmarkc\checkmark$\checkmarks\checkmark$\checkmark^\checkmark\\checkmarkc\checkmarki\checkmarkr\checkmarkc\checkmark$\checkmarku\checkmark$\checkmark^\checkmark\\checkmarkc\checkmarki\checkmarkr\checkmarkc\checkmark$\checkmarkb\checkmark$\checkmark^\checkmark\\checkmarkc\checkmarki\checkmarkr\checkmarkc\checkmark$\checkmarks\checkmark$\checkmark^\checkmark\\checkmarkc\checkmarki\checkmarkr\checkmarkc\checkmark$\checkmarke\checkmark$\checkmark^\checkmark\\checkmarkc\checkmarki\checkmarkr\checkmarkc\checkmark$\checkmarkc\checkmark$\checkmark^\checkmark\\checkmarkc\checkmarki\checkmarkr\checkmarkc\checkmark$\checkmarkt\checkmark$\checkmark^\checkmark\\checkmarkc\checkmarki\checkmarkr\checkmarkc\checkmark$\checkmarki\checkmark$\checkmark^\checkmark\\checkmarkc\checkmarki\checkmarkr\checkmarkc\checkmark$\checkmarko\checkmark$\checkmark^\checkmark\\checkmarkc\checkmarki\checkmarkr\checkmarkc\checkmark$\checkmarkn\checkmark$\checkmark^\checkmark\\checkmarkc\checkmarki\checkmarkr\checkmarkc\checkmark$\checkmark{\checkmark$\checkmark^\checkmark\\checkmarkc\checkmarki\checkmarkr\checkmarkc\checkmark$\checkmarkC\checkmark$\checkmark^\checkmark\\checkmarkc\checkmarki\checkmarkr\checkmarkc\checkmark$\checkmarko\checkmark$\checkmark^\checkmark\\checkmarkc\checkmarki\checkmarkr\checkmarkc\checkmark$\checkmarku\checkmark$\checkmark^\checkmark\\checkmarkc\checkmarki\checkmarkr\checkmarkc\checkmark$\checkmarkp\checkmark$\checkmark^\checkmark\\checkmarkc\checkmarki\checkmarkr\checkmarkc\checkmark$\checkmarkl\checkmark$\checkmark^\checkmark\\checkmarkc\checkmarki\checkmarkr\checkmarkc\checkmark$\checkmarke\checkmark$\checkmark^\checkmark\\checkmarkc\checkmarki\checkmarkr\checkmarkc\checkmark$\checkmark \checkmark$\checkmark^\checkmark\\checkmarkc\checkmarki\checkmarkr\checkmarkc\checkmark$\checkmarkd\checkmark$\checkmark^\checkmark\\checkmarkc\checkmarki\checkmarkr\checkmarkc\checkmark$\checkmarke\checkmark$\checkmark^\checkmark\\checkmarkc\checkmarki\checkmarkr\checkmarkc\checkmark$\checkmark \checkmark$\checkmark^\checkmark\\checkmarkc\checkmarki\checkmarkr\checkmarkc\checkmark$\checkmarkm\checkmark$\checkmark^\checkmark\\checkmarkc\checkmarki\checkmarkr\checkmarkc\checkmark$\checkmarka\checkmark$\checkmark^\checkmark\\checkmarkc\checkmarki\checkmarkr\checkmarkc\checkmark$\checkmarki\checkmark$\checkmark^\checkmark\\checkmarkc\checkmarki\checkmarkr\checkmarkc\checkmark$\checkmarkn\checkmark$\checkmark^\checkmark\\checkmarkc\checkmarki\checkmarkr\checkmarkc\checkmark$\checkmarkt\checkmark$\checkmark^\checkmark\\checkmarkc\checkmarki\checkmarkr\checkmarkc\checkmark$\checkmarki\checkmark$\checkmark^\checkmark\\checkmarkc\checkmarki\checkmarkr\checkmarkc\checkmark$\checkmarke\checkmark$\checkmark^\checkmark\\checkmarkc\checkmarki\checkmarkr\checkmarkc\checkmark$\checkmarkn\checkmark$\checkmark^\checkmark\\checkmarkc\checkmarki\checkmarkr\checkmarkc\checkmark$\checkmark \checkmark$\checkmark^\checkmark\\checkmarkc\checkmarki\checkmarkr\checkmarkc\checkmark$\checkmarkm\checkmark$\checkmark^\checkmark\\checkmarkc\checkmarki\checkmarkr\checkmarkc\checkmark$\checkmarki\checkmark$\checkmark^\checkmark\\checkmarkc\checkmarki\checkmarkr\checkmarkc\checkmark$\checkmarkn\checkmark$\checkmark^\checkmark\\checkmarkc\checkmarki\checkmarkr\checkmarkc\checkmark$\checkmarki\checkmark$\checkmark^\checkmark\\checkmarkc\checkmarki\checkmarkr\checkmarkc\checkmark$\checkmarkm\checkmark$\checkmark^\checkmark\\checkmarkc\checkmarki\checkmarkr\checkmarkc\checkmark$\checkmarka\checkmark$\checkmark^\checkmark\\checkmarkc\checkmarki\checkmarkr\checkmarkc\checkmark$\checkmarkl\checkmark$\checkmark^\checkmark\\checkmarkc\checkmarki\checkmarkr\checkmarkc\checkmark$\checkmark \checkmark$\checkmark^\checkmark\\checkmarkc\checkmarki\checkmarkr\checkmarkc\checkmark$\checkmarkr\checkmark$\checkmark^\checkmark\\checkmarkc\checkmarki\checkmarkr\checkmarkc\checkmark$\checkmarke\checkmark$\checkmark^\checkmark\\checkmarkc\checkmarki\checkmarkr\checkmarkc\checkmark$\checkmarkq\checkmark$\checkmark^\checkmark\\checkmarkc\checkmarki\checkmarkr\checkmarkc\checkmark$\checkmarku\checkmark$\checkmark^\checkmark\\checkmarkc\checkmarki\checkmarkr\checkmarkc\checkmark$\checkmarki\checkmark$\checkmark^\checkmark\\checkmarkc\checkmarki\checkmarkr\checkmarkc\checkmark$\checkmarks\checkmark$\checkmark^\checkmark\\checkmarkc\checkmarki\checkmarkr\checkmarkc\checkmark$\checkmark}\checkmark$\checkmark^\checkmark\\checkmarkc\checkmarki\checkmarkr\checkmarkc\checkmark$\checkmark
\checkmark$\checkmark^\checkmark\\checkmarkc\checkmarki\checkmarkr\checkmarkc\checkmark$\checkmarkL\checkmark$\checkmark^\checkmark\\checkmarkc\checkmarki\checkmarkr\checkmarkc\checkmark$\checkmarke\checkmark$\checkmark^\checkmark\\checkmarkc\checkmarki\checkmarkr\checkmarkc\checkmark$\checkmark \checkmark$\checkmark^\checkmark\\checkmarkc\checkmarki\checkmarkr\checkmarkc\checkmark$\checkmarkc\checkmark$\checkmark^\checkmark\\checkmarkc\checkmarki\checkmarkr\checkmarkc\checkmark$\checkmarko\checkmark$\checkmark^\checkmark\\checkmarkc\checkmarki\checkmarkr\checkmarkc\checkmark$\checkmarku\checkmark$\checkmark^\checkmark\\checkmarkc\checkmarki\checkmarkr\checkmarkc\checkmark$\checkmarkp\checkmark$\checkmark^\checkmark\\checkmarkc\checkmarki\checkmarkr\checkmarkc\checkmark$\checkmarkl\checkmark$\checkmark^\checkmark\\checkmarkc\checkmarki\checkmarkr\checkmarkc\checkmark$\checkmarke\checkmark$\checkmark^\checkmark\\checkmarkc\checkmarki\checkmarkr\checkmarkc\checkmark$\checkmark \checkmark$\checkmark^\checkmark\\checkmarkc\checkmarki\checkmarkr\checkmarkc\checkmark$\checkmarkn\checkmark$\checkmark^\checkmark\\checkmarkc\checkmarki\checkmarkr\checkmarkc\checkmark$\checkmarko\checkmark$\checkmark^\checkmark\\checkmarkc\checkmarki\checkmarkr\checkmarkc\checkmark$\checkmarkm\checkmark$\checkmark^\checkmark\\checkmarkc\checkmarki\checkmarkr\checkmarkc\checkmark$\checkmarki\checkmark$\checkmark^\checkmark\\checkmarkc\checkmarki\checkmarkr\checkmarkc\checkmark$\checkmarkn\checkmark$\checkmark^\checkmark\\checkmarkc\checkmarki\checkmarkr\checkmarkc\checkmark$\checkmarka\checkmark$\checkmark^\checkmark\\checkmarkc\checkmarki\checkmarkr\checkmarkc\checkmark$\checkmarkl\checkmark$\checkmark^\checkmark\\checkmarkc\checkmarki\checkmarkr\checkmarkc\checkmark$\checkmark \checkmark$\checkmark^\checkmark\\checkmarkc\checkmarki\checkmarkr\checkmarkc\checkmark$\checkmarkt\checkmark$\checkmark^\checkmark\\checkmarkc\checkmarki\checkmarkr\checkmarkc\checkmark$\checkmarko\checkmark$\checkmark^\checkmark\\checkmarkc\checkmarki\checkmarkr\checkmarkc\checkmark$\checkmarkt\checkmark$\checkmark^\checkmark\\checkmarkc\checkmarki\checkmarkr\checkmarkc\checkmark$\checkmarka\checkmark$\checkmark^\checkmark\\checkmarkc\checkmarki\checkmarkr\checkmarkc\checkmark$\checkmarkl\checkmark$\checkmark^\checkmark\\checkmarkc\checkmarki\checkmarkr\checkmarkc\checkmark$\checkmark \checkmark$\checkmark^\checkmark\\checkmarkc\checkmarki\checkmarkr\checkmarkc\checkmark$\checkmarke\checkmark$\checkmark^\checkmark\\checkmarkc\checkmarki\checkmarkr\checkmarkc\checkmark$\checkmarkn\checkmark$\checkmark^\checkmark\\checkmarkc\checkmarki\checkmarkr\checkmarkc\checkmark$\checkmark \checkmark$\checkmark^\checkmark\\checkmarkc\checkmarki\checkmarkr\checkmarkc\checkmark$\checkmarkr\checkmark$\checkmark^\checkmark\\checkmarkc\checkmarki\checkmarkr\checkmarkc\checkmark$\checkmarké\checkmark$\checkmark^\checkmark\\checkmarkc\checkmarki\checkmarkr\checkmarkc\checkmark$\checkmarkg\checkmark$\checkmark^\checkmark\\checkmarkc\checkmarki\checkmarkr\checkmarkc\checkmark$\checkmarki\checkmark$\checkmark^\checkmark\\checkmarkc\checkmarki\checkmarkr\checkmarkc\checkmark$\checkmarkm\checkmark$\checkmark^\checkmark\\checkmarkc\checkmarki\checkmarkr\checkmarkc\checkmark$\checkmarke\checkmark$\checkmark^\checkmark\\checkmarkc\checkmarki\checkmarkr\checkmarkc\checkmark$\checkmark \checkmark$\checkmark^\checkmark\\checkmarkc\checkmarki\checkmarkr\checkmarkc\checkmark$\checkmarkd\checkmark$\checkmark^\checkmark\\checkmarkc\checkmarki\checkmarkr\checkmarkc\checkmark$\checkmarky\checkmark$\checkmark^\checkmark\\checkmarkc\checkmarki\checkmarkr\checkmarkc\checkmark$\checkmarkn\checkmark$\checkmark^\checkmark\\checkmarkc\checkmarki\checkmarkr\checkmarkc\checkmark$\checkmarka\checkmark$\checkmark^\checkmark\\checkmarkc\checkmarki\checkmarkr\checkmarkc\checkmark$\checkmarkm\checkmark$\checkmark^\checkmark\\checkmarkc\checkmarki\checkmarkr\checkmarkc\checkmark$\checkmarki\checkmark$\checkmark^\checkmark\\checkmarkc\checkmarki\checkmarkr\checkmarkc\checkmark$\checkmarkq\checkmark$\checkmark^\checkmark\\checkmarkc\checkmarki\checkmarkr\checkmarkc\checkmark$\checkmarku\checkmark$\checkmark^\checkmark\\checkmarkc\checkmarki\checkmarkr\checkmarkc\checkmark$\checkmarke\checkmark$\checkmark^\checkmark\\checkmarkc\checkmarki\checkmarkr\checkmarkc\checkmark$\checkmark \checkmark$\checkmark^\checkmark\\checkmarkc\checkmarki\checkmarkr\checkmarkc\checkmark$\checkmark:\checkmark$\checkmark^\checkmark\\checkmarkc\checkmarki\checkmarkr\checkmarkc\checkmark$\checkmark
\checkmark$\checkmark^\checkmark\\checkmarkc\checkmarki\checkmarkr\checkmarkc\checkmark$\checkmark\\checkmark$\checkmark^\checkmark\\checkmarkc\checkmarki\checkmarkr\checkmarkc\checkmark$\checkmarkb\checkmark$\checkmark^\checkmark\\checkmarkc\checkmarki\checkmarkr\checkmarkc\checkmark$\checkmarke\checkmark$\checkmark^\checkmark\\checkmarkc\checkmarki\checkmarkr\checkmarkc\checkmark$\checkmarkg\checkmark$\checkmark^\checkmark\\checkmarkc\checkmarki\checkmarkr\checkmarkc\checkmark$\checkmarki\checkmark$\checkmark^\checkmark\\checkmarkc\checkmarki\checkmarkr\checkmarkc\checkmark$\checkmarkn\checkmark$\checkmark^\checkmark\\checkmarkc\checkmarki\checkmarkr\checkmarkc\checkmark$\checkmark{\checkmark$\checkmark^\checkmark\\checkmarkc\checkmarki\checkmarkr\checkmarkc\checkmark$\checkmarke\checkmark$\checkmark^\checkmark\\checkmarkc\checkmarki\checkmarkr\checkmarkc\checkmark$\checkmarkq\checkmark$\checkmark^\checkmark\\checkmarkc\checkmarki\checkmarkr\checkmarkc\checkmark$\checkmarku\checkmark$\checkmark^\checkmark\\checkmarkc\checkmarki\checkmarkr\checkmarkc\checkmark$\checkmarka\checkmark$\checkmark^\checkmark\\checkmarkc\checkmarki\checkmarkr\checkmarkc\checkmark$\checkmarkt\checkmark$\checkmark^\checkmark\\checkmarkc\checkmarki\checkmarkr\checkmarkc\checkmark$\checkmarki\checkmark$\checkmark^\checkmark\\checkmarkc\checkmarki\checkmarkr\checkmarkc\checkmark$\checkmarko\checkmark$\checkmark^\checkmark\\checkmarkc\checkmarki\checkmarkr\checkmarkc\checkmark$\checkmarkn\checkmark$\checkmark^\checkmark\\checkmarkc\checkmarki\checkmarkr\checkmarkc\checkmark$\checkmark}\checkmark$\checkmark^\checkmark\\checkmarkc\checkmarki\checkmarkr\checkmarkc\checkmark$\checkmark
\checkmark$\checkmark^\checkmark\\checkmarkc\checkmarki\checkmarkr\checkmarkc\checkmark$\checkmark \checkmark$\checkmark^\checkmark\\checkmarkc\checkmarki\checkmarkr\checkmarkc\checkmark$\checkmark \checkmark$\checkmark^\checkmark\\checkmarkc\checkmarki\checkmarkr\checkmarkc\checkmark$\checkmark \checkmark$\checkmark^\checkmark\\checkmarkc\checkmarki\checkmarkr\checkmarkc\checkmark$\checkmark \checkmark$\checkmark^\checkmark\\checkmarkc\checkmarki\checkmarkr\checkmarkc\checkmark$\checkmarkC\checkmark$\checkmark^\checkmark\\checkmarkc\checkmarki\checkmarkr\checkmarkc\checkmark$\checkmark_\checkmark$\checkmark^\checkmark\\checkmarkc\checkmarki\checkmarkr\checkmarkc\checkmark$\checkmark{\checkmark$\checkmark^\checkmark\\checkmarkc\checkmarki\checkmarkr\checkmarkc\checkmark$\checkmarkn\checkmark$\checkmark^\checkmark\\checkmarkc\checkmarki\checkmarkr\checkmarkc\checkmark$\checkmarko\checkmark$\checkmark^\checkmark\\checkmarkc\checkmarki\checkmarkr\checkmarkc\checkmark$\checkmarkm\checkmark$\checkmark^\checkmark\\checkmarkc\checkmarki\checkmarkr\checkmarkc\checkmark$\checkmark}\checkmark$\checkmark^\checkmark\\checkmarkc\checkmarki\checkmarkr\checkmarkc\checkmark$\checkmark \checkmark$\checkmark^\checkmark\\checkmarkc\checkmarki\checkmarkr\checkmarkc\checkmark$\checkmark=\checkmark$\checkmark^\checkmark\\checkmarkc\checkmarki\checkmarkr\checkmarkc\checkmark$\checkmark \checkmark$\checkmark^\checkmark\\checkmarkc\checkmarki\checkmarkr\checkmarkc\checkmark$\checkmarkC\checkmark$\checkmark^\checkmark\\checkmarkc\checkmarki\checkmarkr\checkmarkc\checkmark$\checkmark_\checkmark$\checkmark^\checkmark\\checkmarkc\checkmarki\checkmarkr\checkmarkc\checkmark$\checkmarkr\checkmark$\checkmark^\checkmark\\checkmarkc\checkmarki\checkmarkr\checkmarkc\checkmark$\checkmark \checkmark$\checkmark^\checkmark\\checkmarkc\checkmarki\checkmarkr\checkmarkc\checkmark$\checkmark+\checkmark$\checkmark^\checkmark\\checkmarkc\checkmarki\checkmarkr\checkmarkc\checkmark$\checkmark \checkmark$\checkmark^\checkmark\\checkmarkc\checkmarki\checkmarkr\checkmarkc\checkmark$\checkmarkC\checkmark$\checkmark^\checkmark\\checkmarkc\checkmarki\checkmarkr\checkmarkc\checkmark$\checkmark_\checkmark$\checkmark^\checkmark\\checkmarkc\checkmarki\checkmarkr\checkmarkc\checkmark$\checkmarkJ\checkmark$\checkmark^\checkmark\\checkmarkc\checkmarki\checkmarkr\checkmarkc\checkmark$\checkmark \checkmark$\checkmark^\checkmark\\checkmarkc\checkmarki\checkmarkr\checkmarkc\checkmark$\checkmark=\checkmark$\checkmark^\checkmark\\checkmarkc\checkmarki\checkmarkr\checkmarkc\checkmark$\checkmark \checkmark$\checkmark^\checkmark\\checkmarkc\checkmarki\checkmarkr\checkmarkc\checkmark$\checkmark0\checkmark$\checkmark^\checkmark\\checkmarkc\checkmarki\checkmarkr\checkmarkc\checkmark$\checkmark{\checkmark$\checkmark^\checkmark\\checkmarkc\checkmarki\checkmarkr\checkmarkc\checkmark$\checkmark,\checkmark$\checkmark^\checkmark\\checkmarkc\checkmarki\checkmarkr\checkmarkc\checkmark$\checkmark}\checkmark$\checkmark^\checkmark\\checkmarkc\checkmarki\checkmarkr\checkmarkc\checkmark$\checkmark0\checkmark$\checkmark^\checkmark\\checkmarkc\checkmarki\checkmarkr\checkmarkc\checkmark$\checkmark9\checkmark$\checkmark^\checkmark\\checkmarkc\checkmarki\checkmarkr\checkmarkc\checkmark$\checkmark3\checkmark$\checkmark^\checkmark\\checkmarkc\checkmarki\checkmarkr\checkmarkc\checkmark$\checkmark \checkmark$\checkmark^\checkmark\\checkmarkc\checkmarki\checkmarkr\checkmarkc\checkmark$\checkmark+\checkmark$\checkmark^\checkmark\\checkmarkc\checkmarki\checkmarkr\checkmarkc\checkmark$\checkmark \checkmark$\checkmark^\checkmark\\checkmarkc\checkmarki\checkmarkr\checkmarkc\checkmark$\checkmark0\checkmark$\checkmark^\checkmark\\checkmarkc\checkmarki\checkmarkr\checkmarkc\checkmark$\checkmark{\checkmark$\checkmark^\checkmark\\checkmarkc\checkmarki\checkmarkr\checkmarkc\checkmark$\checkmark,\checkmark$\checkmark^\checkmark\\checkmarkc\checkmarki\checkmarkr\checkmarkc\checkmark$\checkmark}\checkmark$\checkmark^\checkmark\\checkmarkc\checkmarki\checkmarkr\checkmarkc\checkmark$\checkmark0\checkmark$\checkmark^\checkmark\\checkmarkc\checkmarki\checkmarkr\checkmarkc\checkmark$\checkmark4\checkmark$\checkmark^\checkmark\\checkmarkc\checkmarki\checkmarkr\checkmarkc\checkmark$\checkmark0\checkmark$\checkmark^\checkmark\\checkmarkc\checkmarki\checkmarkr\checkmarkc\checkmark$\checkmark \checkmark$\checkmark^\checkmark\\checkmarkc\checkmarki\checkmarkr\checkmarkc\checkmark$\checkmark=\checkmark$\checkmark^\checkmark\\checkmarkc\checkmarki\checkmarkr\checkmarkc\checkmark$\checkmark \checkmark$\checkmark^\checkmark\\checkmarkc\checkmarki\checkmarkr\checkmarkc\checkmark$\checkmark0\checkmark$\checkmark^\checkmark\\checkmarkc\checkmarki\checkmarkr\checkmarkc\checkmark$\checkmark{\checkmark$\checkmark^\checkmark\\checkmarkc\checkmarki\checkmarkr\checkmarkc\checkmark$\checkmark,\checkmark$\checkmark^\checkmark\\checkmarkc\checkmarki\checkmarkr\checkmarkc\checkmark$\checkmark}\checkmark$\checkmark^\checkmark\\checkmarkc\checkmarki\checkmarkr\checkmarkc\checkmark$\checkmark1\checkmark$\checkmark^\checkmark\\checkmarkc\checkmarki\checkmarkr\checkmarkc\checkmark$\checkmark3\checkmark$\checkmark^\checkmark\\checkmarkc\checkmarki\checkmarkr\checkmarkc\checkmark$\checkmark3\checkmark$\checkmark^\checkmark\\checkmarkc\checkmarki\checkmarkr\checkmarkc\checkmark$\checkmark \checkmark$\checkmark^\checkmark\\checkmarkc\checkmarki\checkmarkr\checkmarkc\checkmark$\checkmark\\checkmark$\checkmark^\checkmark\\checkmarkc\checkmarki\checkmarkr\checkmarkc\checkmark$\checkmarkt\checkmark$\checkmark^\checkmark\\checkmarkc\checkmarki\checkmarkr\checkmarkc\checkmark$\checkmarke\checkmark$\checkmark^\checkmark\\checkmarkc\checkmarki\checkmarkr\checkmarkc\checkmark$\checkmarkx\checkmark$\checkmark^\checkmark\\checkmarkc\checkmarki\checkmarkr\checkmarkc\checkmark$\checkmarkt\checkmark$\checkmark^\checkmark\\checkmarkc\checkmarki\checkmarkr\checkmarkc\checkmark$\checkmark{\checkmark$\checkmark^\checkmark\\checkmarkc\checkmarki\checkmarkr\checkmarkc\checkmark$\checkmark \checkmark$\checkmark^\checkmark\\checkmarkc\checkmarki\checkmarkr\checkmarkc\checkmark$\checkmarkN\checkmark$\checkmark^\checkmark\\checkmarkc\checkmarki\checkmarkr\checkmarkc\checkmark$\checkmark$\checkmark$\checkmark^\checkmark\\checkmarkc\checkmarki\checkmarkr\checkmarkc\checkmark$\checkmark\\checkmark$\checkmark^\checkmark\\checkmarkc\checkmarki\checkmarkr\checkmarkc\checkmark$\checkmarkc\checkmark$\checkmark^\checkmark\\checkmarkc\checkmarki\checkmarkr\checkmarkc\checkmark$\checkmarkd\checkmark$\checkmark^\checkmark\\checkmarkc\checkmarki\checkmarkr\checkmarkc\checkmark$\checkmarko\checkmark$\checkmark^\checkmark\\checkmarkc\checkmarki\checkmarkr\checkmarkc\checkmark$\checkmarkt\checkmark$\checkmark^\checkmark\\checkmarkc\checkmarki\checkmarkr\checkmarkc\checkmark$\checkmark$\checkmark$\checkmark^\checkmark\\checkmarkc\checkmarki\checkmarkr\checkmarkc\checkmark$\checkmarkm\checkmark$\checkmark^\checkmark\\checkmarkc\checkmarki\checkmarkr\checkmarkc\checkmark$\checkmark}\checkmark$\checkmark^\checkmark\\checkmarkc\checkmarki\checkmarkr\checkmarkc\checkmark$\checkmark
\checkmark$\checkmark^\checkmark\\checkmarkc\checkmarki\checkmarkr\checkmarkc\checkmark$\checkmark\\checkmark$\checkmark^\checkmark\\checkmarkc\checkmarki\checkmarkr\checkmarkc\checkmark$\checkmarke\checkmark$\checkmark^\checkmark\\checkmarkc\checkmarki\checkmarkr\checkmarkc\checkmark$\checkmarkn\checkmark$\checkmark^\checkmark\\checkmarkc\checkmarki\checkmarkr\checkmarkc\checkmark$\checkmarkd\checkmark$\checkmark^\checkmark\\checkmarkc\checkmarki\checkmarkr\checkmarkc\checkmark$\checkmark{\checkmark$\checkmark^\checkmark\\checkmarkc\checkmarki\checkmarkr\checkmarkc\checkmark$\checkmarke\checkmark$\checkmark^\checkmark\\checkmarkc\checkmarki\checkmarkr\checkmarkc\checkmark$\checkmarkq\checkmark$\checkmark^\checkmark\\checkmarkc\checkmarki\checkmarkr\checkmarkc\checkmark$\checkmarku\checkmark$\checkmark^\checkmark\\checkmarkc\checkmarki\checkmarkr\checkmarkc\checkmark$\checkmarka\checkmark$\checkmark^\checkmark\\checkmarkc\checkmarki\checkmarkr\checkmarkc\checkmark$\checkmarkt\checkmark$\checkmark^\checkmark\\checkmarkc\checkmarki\checkmarkr\checkmarkc\checkmark$\checkmarki\checkmark$\checkmark^\checkmark\\checkmarkc\checkmarki\checkmarkr\checkmarkc\checkmark$\checkmarko\checkmark$\checkmark^\checkmark\\checkmarkc\checkmarki\checkmarkr\checkmarkc\checkmark$\checkmarkn\checkmark$\checkmark^\checkmark\\checkmarkc\checkmarki\checkmarkr\checkmarkc\checkmark$\checkmark}\checkmark$\checkmark^\checkmark\\checkmarkc\checkmarki\checkmarkr\checkmarkc\checkmark$\checkmark
\checkmark$\checkmark^\checkmark\\checkmarkc\checkmarki\checkmarkr\checkmarkc\checkmark$\checkmark
\checkmark$\checkmark^\checkmark\\checkmarkc\checkmarki\checkmarkr\checkmarkc\checkmark$\checkmarkE\checkmark$\checkmark^\checkmark\\checkmarkc\checkmarki\checkmarkr\checkmarkc\checkmark$\checkmarkn\checkmark$\checkmark^\checkmark\\checkmarkc\checkmarki\checkmarkr\checkmarkc\checkmark$\checkmark \checkmark$\checkmark^\checkmark\\checkmarkc\checkmarki\checkmarkr\checkmarkc\checkmark$\checkmarka\checkmark$\checkmark^\checkmark\\checkmarkc\checkmarki\checkmarkr\checkmarkc\checkmark$\checkmarkp\checkmark$\checkmark^\checkmark\\checkmarkc\checkmarki\checkmarkr\checkmarkc\checkmark$\checkmarkp\checkmark$\checkmark^\checkmark\\checkmarkc\checkmarki\checkmarkr\checkmarkc\checkmark$\checkmarkl\checkmark$\checkmark^\checkmark\\checkmarkc\checkmarki\checkmarkr\checkmarkc\checkmark$\checkmarki\checkmark$\checkmark^\checkmark\\checkmarkc\checkmarki\checkmarkr\checkmarkc\checkmark$\checkmarkq\checkmark$\checkmark^\checkmark\\checkmarkc\checkmarki\checkmarkr\checkmarkc\checkmark$\checkmarku\checkmark$\checkmark^\checkmark\\checkmarkc\checkmarki\checkmarkr\checkmarkc\checkmark$\checkmarka\checkmark$\checkmark^\checkmark\\checkmarkc\checkmarki\checkmarkr\checkmarkc\checkmark$\checkmarkn\checkmark$\checkmark^\checkmark\\checkmarkc\checkmarki\checkmarkr\checkmarkc\checkmark$\checkmarkt\checkmark$\checkmark^\checkmark\\checkmarkc\checkmarki\checkmarkr\checkmarkc\checkmark$\checkmark \checkmark$\checkmark^\checkmark\\checkmarkc\checkmarki\checkmarkr\checkmarkc\checkmark$\checkmarku\checkmark$\checkmark^\checkmark\\checkmarkc\checkmarki\checkmarkr\checkmarkc\checkmark$\checkmarkn\checkmark$\checkmark^\checkmark\\checkmarkc\checkmarki\checkmarkr\checkmarkc\checkmark$\checkmark \checkmark$\checkmark^\checkmark\\checkmarkc\checkmarki\checkmarkr\checkmarkc\checkmark$\checkmarkc\checkmark$\checkmark^\checkmark\\checkmarkc\checkmarki\checkmarkr\checkmarkc\checkmark$\checkmarko\checkmark$\checkmark^\checkmark\\checkmarkc\checkmarki\checkmarkr\checkmarkc\checkmark$\checkmarke\checkmark$\checkmark^\checkmark\\checkmarkc\checkmarki\checkmarkr\checkmarkc\checkmark$\checkmarkf\checkmark$\checkmark^\checkmark\\checkmarkc\checkmarki\checkmarkr\checkmarkc\checkmark$\checkmarkf\checkmark$\checkmark^\checkmark\\checkmarkc\checkmarki\checkmarkr\checkmarkc\checkmark$\checkmarki\checkmark$\checkmark^\checkmark\\checkmarkc\checkmarki\checkmarkr\checkmarkc\checkmark$\checkmarkc\checkmark$\checkmark^\checkmark\\checkmarkc\checkmarki\checkmarkr\checkmarkc\checkmark$\checkmarki\checkmark$\checkmark^\checkmark\\checkmarkc\checkmarki\checkmarkr\checkmarkc\checkmark$\checkmarke\checkmark$\checkmark^\checkmark\\checkmarkc\checkmarki\checkmarkr\checkmarkc\checkmark$\checkmarkn\checkmark$\checkmark^\checkmark\\checkmarkc\checkmarki\checkmarkr\checkmarkc\checkmark$\checkmarkt\checkmark$\checkmark^\checkmark\\checkmarkc\checkmarki\checkmarkr\checkmarkc\checkmark$\checkmark \checkmark$\checkmark^\checkmark\\checkmarkc\checkmarki\checkmarkr\checkmarkc\checkmark$\checkmarkd\checkmark$\checkmark^\checkmark\\checkmarkc\checkmarki\checkmarkr\checkmarkc\checkmark$\checkmarke\checkmark$\checkmark^\checkmark\\checkmarkc\checkmarki\checkmarkr\checkmarkc\checkmark$\checkmark \checkmark$\checkmark^\checkmark\\checkmarkc\checkmarki\checkmarkr\checkmarkc\checkmark$\checkmarks\checkmark$\checkmark^\checkmark\\checkmarkc\checkmarki\checkmarkr\checkmarkc\checkmark$\checkmarké\checkmark$\checkmark^\checkmark\\checkmarkc\checkmarki\checkmarkr\checkmarkc\checkmark$\checkmarkc\checkmark$\checkmark^\checkmark\\checkmarkc\checkmarki\checkmarkr\checkmarkc\checkmark$\checkmarku\checkmark$\checkmark^\checkmark\\checkmarkc\checkmarki\checkmarkr\checkmarkc\checkmark$\checkmarkr\checkmark$\checkmark^\checkmark\\checkmarkc\checkmarki\checkmarkr\checkmarkc\checkmark$\checkmarki\checkmark$\checkmark^\checkmark\\checkmarkc\checkmarki\checkmarkr\checkmarkc\checkmark$\checkmarkt\checkmark$\checkmark^\checkmark\\checkmarkc\checkmarki\checkmarkr\checkmarkc\checkmark$\checkmarké\checkmark$\checkmark^\checkmark\\checkmarkc\checkmarki\checkmarkr\checkmarkc\checkmark$\checkmark \checkmark$\checkmark^\checkmark\\checkmarkc\checkmarki\checkmarkr\checkmarkc\checkmark$\checkmarkd\checkmark$\checkmark^\checkmark\\checkmarkc\checkmarki\checkmarkr\checkmarkc\checkmark$\checkmarky\checkmark$\checkmark^\checkmark\\checkmarkc\checkmarki\checkmarkr\checkmarkc\checkmark$\checkmarkn\checkmark$\checkmark^\checkmark\\checkmarkc\checkmarki\checkmarkr\checkmarkc\checkmark$\checkmarka\checkmark$\checkmark^\checkmark\\checkmarkc\checkmarki\checkmarkr\checkmarkc\checkmark$\checkmarkm\checkmark$\checkmark^\checkmark\\checkmarkc\checkmarki\checkmarkr\checkmarkc\checkmark$\checkmarki\checkmark$\checkmark^\checkmark\\checkmarkc\checkmarki\checkmarkr\checkmarkc\checkmark$\checkmarkq\checkmark$\checkmark^\checkmark\\checkmarkc\checkmarki\checkmarkr\checkmarkc\checkmark$\checkmarku\checkmark$\checkmark^\checkmark\\checkmarkc\checkmarki\checkmarkr\checkmarkc\checkmark$\checkmarke\checkmark$\checkmark^\checkmark\\checkmarkc\checkmarki\checkmarkr\checkmarkc\checkmark$\checkmark \checkmark$\checkmark^\checkmark\\checkmarkc\checkmarki\checkmarkr\checkmarkc\checkmark$\checkmarkd\checkmark$\checkmark^\checkmark\\checkmarkc\checkmarki\checkmarkr\checkmarkc\checkmark$\checkmarke\checkmark$\checkmark^\checkmark\\checkmarkc\checkmarki\checkmarkr\checkmarkc\checkmark$\checkmark \checkmark$\checkmark^\checkmark\\checkmarkc\checkmarki\checkmarkr\checkmarkc\checkmark$\checkmark$\checkmark$\checkmark^\checkmark\\checkmarkc\checkmarki\checkmarkr\checkmarkc\checkmark$\checkmarks\checkmark$\checkmark^\checkmark\\checkmarkc\checkmarki\checkmarkr\checkmarkc\checkmark$\checkmark_\checkmark$\checkmark^\checkmark\\checkmarkc\checkmarki\checkmarkr\checkmarkc\checkmark$\checkmark{\checkmark$\checkmark^\checkmark\\checkmarkc\checkmarki\checkmarkr\checkmarkc\checkmark$\checkmarkd\checkmark$\checkmark^\checkmark\\checkmarkc\checkmarki\checkmarkr\checkmarkc\checkmark$\checkmarky\checkmark$\checkmark^\checkmark\\checkmarkc\checkmarki\checkmarkr\checkmarkc\checkmark$\checkmarkn\checkmark$\checkmark^\checkmark\\checkmarkc\checkmarki\checkmarkr\checkmarkc\checkmark$\checkmark}\checkmark$\checkmark^\checkmark\\checkmarkc\checkmarki\checkmarkr\checkmarkc\checkmark$\checkmark \checkmark$\checkmark^\checkmark\\checkmarkc\checkmarki\checkmarkr\checkmarkc\checkmark$\checkmark=\checkmark$\checkmark^\checkmark\\checkmarkc\checkmarki\checkmarkr\checkmarkc\checkmark$\checkmark \checkmark$\checkmark^\checkmark\\checkmarkc\checkmarki\checkmarkr\checkmarkc\checkmark$\checkmark2\checkmark$\checkmark^\checkmark\\checkmarkc\checkmarki\checkmarkr\checkmarkc\checkmark$\checkmark{\checkmark$\checkmark^\checkmark\\checkmarkc\checkmarki\checkmarkr\checkmarkc\checkmark$\checkmark,\checkmark$\checkmark^\checkmark\\checkmarkc\checkmarki\checkmarkr\checkmarkc\checkmark$\checkmark}\checkmark$\checkmark^\checkmark\\checkmarkc\checkmarki\checkmarkr\checkmarkc\checkmark$\checkmark5\checkmark$\checkmark^\checkmark\\checkmarkc\checkmarki\checkmarkr\checkmarkc\checkmark$\checkmark$\checkmark$\checkmark^\checkmark\\checkmarkc\checkmarki\checkmarkr\checkmarkc\checkmark$\checkmark \checkmark$\checkmark^\checkmark\\checkmarkc\checkmarki\checkmarkr\checkmarkc\checkmark$\checkmark(\checkmark$\checkmark^\checkmark\\checkmarkc\checkmarki\checkmarkr\checkmarkc\checkmark$\checkmarkp\checkmark$\checkmark^\checkmark\\checkmarkc\checkmarki\checkmarkr\checkmarkc\checkmark$\checkmarko\checkmark$\checkmark^\checkmark\\checkmarkc\checkmarki\checkmarkr\checkmarkc\checkmark$\checkmarku\checkmark$\checkmark^\checkmark\\checkmarkc\checkmarki\checkmarkr\checkmarkc\checkmark$\checkmarkr\checkmark$\checkmark^\checkmark\\checkmarkc\checkmarki\checkmarkr\checkmarkc\checkmark$\checkmark \checkmark$\checkmark^\checkmark\\checkmarkc\checkmarki\checkmarkr\checkmarkc\checkmark$\checkmarkp\checkmark$\checkmark^\checkmark\\checkmarkc\checkmarki\checkmarkr\checkmarkc\checkmark$\checkmarkr\checkmark$\checkmark^\checkmark\\checkmarkc\checkmarki\checkmarkr\checkmarkc\checkmark$\checkmarké\checkmark$\checkmark^\checkmark\\checkmarkc\checkmarki\checkmarkr\checkmarkc\checkmark$\checkmarkv\checkmark$\checkmark^\checkmark\\checkmarkc\checkmarki\checkmarkr\checkmarkc\checkmark$\checkmarke\checkmark$\checkmark^\checkmark\\checkmarkc\checkmarki\checkmarkr\checkmarkc\checkmark$\checkmarkn\checkmark$\checkmark^\checkmark\\checkmarkc\checkmarki\checkmarkr\checkmarkc\checkmark$\checkmarki\checkmark$\checkmark^\checkmark\\checkmarkc\checkmarki\checkmarkr\checkmarkc\checkmark$\checkmarkr\checkmark$\checkmark^\checkmark\\checkmarkc\checkmarki\checkmarkr\checkmarkc\checkmark$\checkmark \checkmark$\checkmark^\checkmark\\checkmarkc\checkmarki\checkmarkr\checkmarkc\checkmark$\checkmarkl\checkmark$\checkmark^\checkmark\\checkmarkc\checkmarki\checkmarkr\checkmarkc\checkmark$\checkmarke\checkmark$\checkmark^\checkmark\\checkmarkc\checkmarki\checkmarkr\checkmarkc\checkmark$\checkmarks\checkmark$\checkmark^\checkmark\\checkmarkc\checkmarki\checkmarkr\checkmarkc\checkmark$\checkmark \checkmark$\checkmark^\checkmark\\checkmarkc\checkmarki\checkmarkr\checkmarkc\checkmark$\checkmarkp\checkmark$\checkmark^\checkmark\\checkmarkc\checkmarki\checkmarkr\checkmarkc\checkmark$\checkmarke\checkmark$\checkmark^\checkmark\\checkmarkc\checkmarki\checkmarkr\checkmarkc\checkmark$\checkmarkr\checkmark$\checkmark^\checkmark\\checkmarkc\checkmarki\checkmarkr\checkmarkc\checkmark$\checkmarkt\checkmark$\checkmark^\checkmark\\checkmarkc\checkmarki\checkmarkr\checkmarkc\checkmark$\checkmarke\checkmark$\checkmark^\checkmark\\checkmarkc\checkmarki\checkmarkr\checkmarkc\checkmark$\checkmarks\checkmark$\checkmark^\checkmark\\checkmarkc\checkmarki\checkmarkr\checkmarkc\checkmark$\checkmark \checkmark$\checkmark^\checkmark\\checkmarkc\checkmarki\checkmarkr\checkmarkc\checkmark$\checkmarkd\checkmark$\checkmark^\checkmark\\checkmarkc\checkmarki\checkmarkr\checkmarkc\checkmark$\checkmarke\checkmark$\checkmark^\checkmark\\checkmarkc\checkmarki\checkmarkr\checkmarkc\checkmark$\checkmark \checkmark$\checkmark^\checkmark\\checkmarkc\checkmarki\checkmarkr\checkmarkc\checkmark$\checkmarkp\checkmark$\checkmark^\checkmark\\checkmarkc\checkmarki\checkmarkr\checkmarkc\checkmark$\checkmarka\checkmark$\checkmark^\checkmark\\checkmarkc\checkmarki\checkmarkr\checkmarkc\checkmark$\checkmarks\checkmark$\checkmark^\checkmark\\checkmarkc\checkmarki\checkmarkr\checkmarkc\checkmark$\checkmark)\checkmark$\checkmark^\checkmark\\checkmarkc\checkmarki\checkmarkr\checkmarkc\checkmark$\checkmark \checkmark$\checkmark^\checkmark\\checkmarkc\checkmarki\checkmarkr\checkmarkc\checkmark$\checkmark:\checkmark$\checkmark^\checkmark\\checkmarkc\checkmarki\checkmarkr\checkmarkc\checkmark$\checkmark
\checkmark$\checkmark^\checkmark\\checkmarkc\checkmarki\checkmarkr\checkmarkc\checkmark$\checkmark\\checkmark$\checkmark^\checkmark\\checkmarkc\checkmarki\checkmarkr\checkmarkc\checkmark$\checkmarkb\checkmark$\checkmark^\checkmark\\checkmarkc\checkmarki\checkmarkr\checkmarkc\checkmark$\checkmarke\checkmark$\checkmark^\checkmark\\checkmarkc\checkmarki\checkmarkr\checkmarkc\checkmark$\checkmarkg\checkmark$\checkmark^\checkmark\\checkmarkc\checkmarki\checkmarkr\checkmarkc\checkmark$\checkmarki\checkmark$\checkmark^\checkmark\\checkmarkc\checkmarki\checkmarkr\checkmarkc\checkmark$\checkmarkn\checkmark$\checkmark^\checkmark\\checkmarkc\checkmarki\checkmarkr\checkmarkc\checkmark$\checkmark{\checkmark$\checkmark^\checkmark\\checkmarkc\checkmarki\checkmarkr\checkmarkc\checkmark$\checkmarke\checkmark$\checkmark^\checkmark\\checkmarkc\checkmarki\checkmarkr\checkmarkc\checkmark$\checkmarkq\checkmark$\checkmark^\checkmark\\checkmarkc\checkmarki\checkmarkr\checkmarkc\checkmark$\checkmarku\checkmark$\checkmark^\checkmark\\checkmarkc\checkmarki\checkmarkr\checkmarkc\checkmark$\checkmarka\checkmark$\checkmark^\checkmark\\checkmarkc\checkmarki\checkmarkr\checkmarkc\checkmark$\checkmarkt\checkmark$\checkmark^\checkmark\\checkmarkc\checkmarki\checkmarkr\checkmarkc\checkmark$\checkmarki\checkmark$\checkmark^\checkmark\\checkmarkc\checkmarki\checkmarkr\checkmarkc\checkmark$\checkmarko\checkmark$\checkmark^\checkmark\\checkmarkc\checkmarki\checkmarkr\checkmarkc\checkmark$\checkmarkn\checkmark$\checkmark^\checkmark\\checkmarkc\checkmarki\checkmarkr\checkmarkc\checkmark$\checkmark}\checkmark$\checkmark^\checkmark\\checkmarkc\checkmarki\checkmarkr\checkmarkc\checkmark$\checkmark
\checkmark$\checkmark^\checkmark\\checkmarkc\checkmarki\checkmarkr\checkmarkc\checkmark$\checkmark \checkmark$\checkmark^\checkmark\\checkmarkc\checkmarki\checkmarkr\checkmarkc\checkmark$\checkmark \checkmark$\checkmark^\checkmark\\checkmarkc\checkmarki\checkmarkr\checkmarkc\checkmark$\checkmark \checkmark$\checkmark^\checkmark\\checkmarkc\checkmarki\checkmarkr\checkmarkc\checkmark$\checkmark \checkmark$\checkmark^\checkmark\\checkmarkc\checkmarki\checkmarkr\checkmarkc\checkmark$\checkmarkC\checkmark$\checkmark^\checkmark\\checkmarkc\checkmarki\checkmarkr\checkmarkc\checkmark$\checkmark_\checkmark$\checkmark^\checkmark\\checkmarkc\checkmarki\checkmarkr\checkmarkc\checkmark$\checkmark{\checkmark$\checkmark^\checkmark\\checkmarkc\checkmarki\checkmarkr\checkmarkc\checkmark$\checkmarkm\checkmark$\checkmark^\checkmark\\checkmarkc\checkmarki\checkmarkr\checkmarkc\checkmark$\checkmarka\checkmark$\checkmark^\checkmark\\checkmarkc\checkmarki\checkmarkr\checkmarkc\checkmark$\checkmarki\checkmark$\checkmark^\checkmark\\checkmarkc\checkmarki\checkmarkr\checkmarkc\checkmark$\checkmarkn\checkmark$\checkmark^\checkmark\\checkmarkc\checkmarki\checkmarkr\checkmarkc\checkmark$\checkmarkt\checkmark$\checkmark^\checkmark\\checkmarkc\checkmarki\checkmarkr\checkmarkc\checkmark$\checkmarki\checkmark$\checkmark^\checkmark\\checkmarkc\checkmarki\checkmarkr\checkmarkc\checkmark$\checkmarke\checkmark$\checkmark^\checkmark\\checkmarkc\checkmarki\checkmarkr\checkmarkc\checkmark$\checkmarkn\checkmark$\checkmark^\checkmark\\checkmarkc\checkmarki\checkmarkr\checkmarkc\checkmark$\checkmark\\checkmark$\checkmark^\checkmark\\checkmarkc\checkmarki\checkmarkr\checkmarkc\checkmark$\checkmark_\checkmark$\checkmark^\checkmark\\checkmarkc\checkmarki\checkmarkr\checkmarkc\checkmark$\checkmarkr\checkmark$\checkmark^\checkmark\\checkmarkc\checkmarki\checkmarkr\checkmarkc\checkmark$\checkmarke\checkmark$\checkmark^\checkmark\\checkmarkc\checkmarki\checkmarkr\checkmarkc\checkmark$\checkmarkq\checkmark$\checkmark^\checkmark\\checkmarkc\checkmarki\checkmarkr\checkmarkc\checkmark$\checkmark}\checkmark$\checkmark^\checkmark\\checkmarkc\checkmarki\checkmarkr\checkmarkc\checkmark$\checkmark \checkmark$\checkmark^\checkmark\\checkmarkc\checkmarki\checkmarkr\checkmarkc\checkmark$\checkmark=\checkmark$\checkmark^\checkmark\\checkmarkc\checkmarki\checkmarkr\checkmarkc\checkmark$\checkmark \checkmark$\checkmark^\checkmark\\checkmarkc\checkmarki\checkmarkr\checkmarkc\checkmark$\checkmarkC\checkmark$\checkmark^\checkmark\\checkmarkc\checkmarki\checkmarkr\checkmarkc\checkmark$\checkmark_\checkmark$\checkmark^\checkmark\\checkmarkc\checkmarki\checkmarkr\checkmarkc\checkmark$\checkmark{\checkmark$\checkmark^\checkmark\\checkmarkc\checkmarki\checkmarkr\checkmarkc\checkmark$\checkmarkn\checkmark$\checkmark^\checkmark\\checkmarkc\checkmarki\checkmarkr\checkmarkc\checkmark$\checkmarko\checkmark$\checkmark^\checkmark\\checkmarkc\checkmarki\checkmarkr\checkmarkc\checkmark$\checkmarkm\checkmark$\checkmark^\checkmark\\checkmarkc\checkmarki\checkmarkr\checkmarkc\checkmark$\checkmark}\checkmark$\checkmark^\checkmark\\checkmarkc\checkmarki\checkmarkr\checkmarkc\checkmark$\checkmark \checkmark$\checkmark^\checkmark\\checkmarkc\checkmarki\checkmarkr\checkmarkc\checkmark$\checkmark\\checkmark$\checkmark^\checkmark\\checkmarkc\checkmarki\checkmarkr\checkmarkc\checkmark$\checkmarkt\checkmark$\checkmark^\checkmark\\checkmarkc\checkmarki\checkmarkr\checkmarkc\checkmark$\checkmarki\checkmark$\checkmark^\checkmark\\checkmarkc\checkmarki\checkmarkr\checkmarkc\checkmark$\checkmarkm\checkmark$\checkmark^\checkmark\\checkmarkc\checkmarki\checkmarkr\checkmarkc\checkmark$\checkmarke\checkmark$\checkmark^\checkmark\\checkmarkc\checkmarki\checkmarkr\checkmarkc\checkmark$\checkmarks\checkmark$\checkmark^\checkmark\\checkmarkc\checkmarki\checkmarkr\checkmarkc\checkmark$\checkmark \checkmark$\checkmark^\checkmark\\checkmarkc\checkmarki\checkmarkr\checkmarkc\checkmark$\checkmarks\checkmark$\checkmark^\checkmark\\checkmarkc\checkmarki\checkmarkr\checkmarkc\checkmark$\checkmark_\checkmark$\checkmark^\checkmark\\checkmarkc\checkmarki\checkmarkr\checkmarkc\checkmark$\checkmark{\checkmark$\checkmark^\checkmark\\checkmarkc\checkmarki\checkmarkr\checkmarkc\checkmark$\checkmarkd\checkmark$\checkmark^\checkmark\\checkmarkc\checkmarki\checkmarkr\checkmarkc\checkmark$\checkmarky\checkmark$\checkmark^\checkmark\\checkmarkc\checkmarki\checkmarkr\checkmarkc\checkmark$\checkmarkn\checkmark$\checkmark^\checkmark\\checkmarkc\checkmarki\checkmarkr\checkmarkc\checkmark$\checkmark}\checkmark$\checkmark^\checkmark\\checkmarkc\checkmarki\checkmarkr\checkmarkc\checkmark$\checkmark \checkmark$\checkmark^\checkmark\\checkmarkc\checkmarki\checkmarkr\checkmarkc\checkmark$\checkmark=\checkmark$\checkmark^\checkmark\\checkmarkc\checkmarki\checkmarkr\checkmarkc\checkmark$\checkmark \checkmark$\checkmark^\checkmark\\checkmarkc\checkmarki\checkmarkr\checkmarkc\checkmark$\checkmark0\checkmark$\checkmark^\checkmark\\checkmarkc\checkmarki\checkmarkr\checkmarkc\checkmark$\checkmark{\checkmark$\checkmark^\checkmark\\checkmarkc\checkmarki\checkmarkr\checkmarkc\checkmark$\checkmark,\checkmark$\checkmark^\checkmark\\checkmarkc\checkmarki\checkmarkr\checkmarkc\checkmark$\checkmark}\checkmark$\checkmark^\checkmark\\checkmarkc\checkmarki\checkmarkr\checkmarkc\checkmark$\checkmark1\checkmark$\checkmark^\checkmark\\checkmarkc\checkmarki\checkmarkr\checkmarkc\checkmark$\checkmark3\checkmark$\checkmark^\checkmark\\checkmarkc\checkmarki\checkmarkr\checkmarkc\checkmark$\checkmark3\checkmark$\checkmark^\checkmark\\checkmarkc\checkmarki\checkmarkr\checkmarkc\checkmark$\checkmark \checkmark$\checkmark^\checkmark\\checkmarkc\checkmarki\checkmarkr\checkmarkc\checkmark$\checkmark\\checkmark$\checkmark^\checkmark\\checkmarkc\checkmarki\checkmarkr\checkmarkc\checkmark$\checkmarkt\checkmark$\checkmark^\checkmark\\checkmarkc\checkmarki\checkmarkr\checkmarkc\checkmark$\checkmarki\checkmark$\checkmark^\checkmark\\checkmarkc\checkmarki\checkmarkr\checkmarkc\checkmark$\checkmarkm\checkmark$\checkmark^\checkmark\\checkmarkc\checkmarki\checkmarkr\checkmarkc\checkmark$\checkmarke\checkmark$\checkmark^\checkmark\\checkmarkc\checkmarki\checkmarkr\checkmarkc\checkmark$\checkmarks\checkmark$\checkmark^\checkmark\\checkmarkc\checkmarki\checkmarkr\checkmarkc\checkmark$\checkmark \checkmark$\checkmark^\checkmark\\checkmarkc\checkmarki\checkmarkr\checkmarkc\checkmark$\checkmark2\checkmark$\checkmark^\checkmark\\checkmarkc\checkmarki\checkmarkr\checkmarkc\checkmark$\checkmark{\checkmark$\checkmark^\checkmark\\checkmarkc\checkmarki\checkmarkr\checkmarkc\checkmark$\checkmark,\checkmark$\checkmark^\checkmark\\checkmarkc\checkmarki\checkmarkr\checkmarkc\checkmark$\checkmark}\checkmark$\checkmark^\checkmark\\checkmarkc\checkmarki\checkmarkr\checkmarkc\checkmark$\checkmark5\checkmark$\checkmark^\checkmark\\checkmarkc\checkmarki\checkmarkr\checkmarkc\checkmark$\checkmark \checkmark$\checkmark^\checkmark\\checkmarkc\checkmarki\checkmarkr\checkmarkc\checkmark$\checkmark=\checkmark$\checkmark^\checkmark\\checkmarkc\checkmarki\checkmarkr\checkmarkc\checkmark$\checkmark \checkmark$\checkmark^\checkmark\\checkmarkc\checkmarki\checkmarkr\checkmarkc\checkmark$\checkmark\\checkmark$\checkmark^\checkmark\\checkmarkc\checkmarki\checkmarkr\checkmarkc\checkmark$\checkmarkb\checkmark$\checkmark^\checkmark\\checkmarkc\checkmarki\checkmarkr\checkmarkc\checkmark$\checkmarko\checkmark$\checkmark^\checkmark\\checkmarkc\checkmarki\checkmarkr\checkmarkc\checkmark$\checkmarkx\checkmark$\checkmark^\checkmark\\checkmarkc\checkmarki\checkmarkr\checkmarkc\checkmark$\checkmarke\checkmark$\checkmark^\checkmark\\checkmarkc\checkmarki\checkmarkr\checkmarkc\checkmark$\checkmarkd\checkmark$\checkmark^\checkmark\\checkmarkc\checkmarki\checkmarkr\checkmarkc\checkmark$\checkmark{\checkmark$\checkmark^\checkmark\\checkmarkc\checkmarki\checkmarkr\checkmarkc\checkmark$\checkmark0\checkmark$\checkmark^\checkmark\\checkmarkc\checkmarki\checkmarkr\checkmarkc\checkmark$\checkmark{\checkmark$\checkmark^\checkmark\\checkmarkc\checkmarki\checkmarkr\checkmarkc\checkmark$\checkmark,\checkmark$\checkmark^\checkmark\\checkmarkc\checkmarki\checkmarkr\checkmarkc\checkmark$\checkmark}\checkmark$\checkmark^\checkmark\\checkmarkc\checkmarki\checkmarkr\checkmarkc\checkmark$\checkmark3\checkmark$\checkmark^\checkmark\\checkmarkc\checkmarki\checkmarkr\checkmarkc\checkmark$\checkmark3\checkmark$\checkmark^\checkmark\\checkmarkc\checkmarki\checkmarkr\checkmarkc\checkmark$\checkmark3\checkmark$\checkmark^\checkmark\\checkmarkc\checkmarki\checkmarkr\checkmarkc\checkmark$\checkmark \checkmark$\checkmark^\checkmark\\checkmarkc\checkmarki\checkmarkr\checkmarkc\checkmark$\checkmark\\checkmark$\checkmark^\checkmark\\checkmarkc\checkmarki\checkmarkr\checkmarkc\checkmark$\checkmarkt\checkmark$\checkmark^\checkmark\\checkmarkc\checkmarki\checkmarkr\checkmarkc\checkmark$\checkmarke\checkmark$\checkmark^\checkmark\\checkmarkc\checkmarki\checkmarkr\checkmarkc\checkmark$\checkmarkx\checkmark$\checkmark^\checkmark\\checkmarkc\checkmarki\checkmarkr\checkmarkc\checkmark$\checkmarkt\checkmark$\checkmark^\checkmark\\checkmarkc\checkmarki\checkmarkr\checkmarkc\checkmark$\checkmark{\checkmark$\checkmark^\checkmark\\checkmarkc\checkmarki\checkmarkr\checkmarkc\checkmark$\checkmark \checkmark$\checkmark^\checkmark\\checkmarkc\checkmarki\checkmarkr\checkmarkc\checkmark$\checkmarkN\checkmark$\checkmark^\checkmark\\checkmarkc\checkmarki\checkmarkr\checkmarkc\checkmark$\checkmark$\checkmark$\checkmark^\checkmark\\checkmarkc\checkmarki\checkmarkr\checkmarkc\checkmark$\checkmark\\checkmark$\checkmark^\checkmark\\checkmarkc\checkmarki\checkmarkr\checkmarkc\checkmark$\checkmarkc\checkmark$\checkmark^\checkmark\\checkmarkc\checkmarki\checkmarkr\checkmarkc\checkmark$\checkmarkd\checkmark$\checkmark^\checkmark\\checkmarkc\checkmarki\checkmarkr\checkmarkc\checkmark$\checkmarko\checkmark$\checkmark^\checkmark\\checkmarkc\checkmarki\checkmarkr\checkmarkc\checkmark$\checkmarkt\checkmark$\checkmark^\checkmark\\checkmarkc\checkmarki\checkmarkr\checkmarkc\checkmark$\checkmark$\checkmark$\checkmark^\checkmark\\checkmarkc\checkmarki\checkmarkr\checkmarkc\checkmark$\checkmarkm\checkmark$\checkmark^\checkmark\\checkmarkc\checkmarki\checkmarkr\checkmarkc\checkmark$\checkmark}\checkmark$\checkmark^\checkmark\\checkmarkc\checkmarki\checkmarkr\checkmarkc\checkmark$\checkmark}\checkmark$\checkmark^\checkmark\\checkmarkc\checkmarki\checkmarkr\checkmarkc\checkmark$\checkmark
\checkmark$\checkmark^\checkmark\\checkmarkc\checkmarki\checkmarkr\checkmarkc\checkmark$\checkmark \checkmark$\checkmark^\checkmark\\checkmarkc\checkmarki\checkmarkr\checkmarkc\checkmark$\checkmark \checkmark$\checkmark^\checkmark\\checkmarkc\checkmarki\checkmarkr\checkmarkc\checkmark$\checkmark \checkmark$\checkmark^\checkmark\\checkmarkc\checkmarki\checkmarkr\checkmarkc\checkmark$\checkmark \checkmark$\checkmark^\checkmark\\checkmarkc\checkmarki\checkmarkr\checkmarkc\checkmark$\checkmark\\checkmark$\checkmark^\checkmark\\checkmarkc\checkmarki\checkmarkr\checkmarkc\checkmark$\checkmarkl\checkmark$\checkmark^\checkmark\\checkmarkc\checkmarki\checkmarkr\checkmarkc\checkmark$\checkmarka\checkmark$\checkmark^\checkmark\\checkmarkc\checkmarki\checkmarkr\checkmarkc\checkmark$\checkmarkb\checkmark$\checkmark^\checkmark\\checkmarkc\checkmarki\checkmarkr\checkmarkc\checkmark$\checkmarke\checkmark$\checkmark^\checkmark\\checkmarkc\checkmarki\checkmarkr\checkmarkc\checkmark$\checkmarkl\checkmark$\checkmark^\checkmark\\checkmarkc\checkmarki\checkmarkr\checkmarkc\checkmark$\checkmark{\checkmark$\checkmark^\checkmark\\checkmarkc\checkmarki\checkmarkr\checkmarkc\checkmark$\checkmarke\checkmark$\checkmark^\checkmark\\checkmarkc\checkmarki\checkmarkr\checkmarkc\checkmark$\checkmarkq\checkmark$\checkmark^\checkmark\\checkmarkc\checkmarki\checkmarkr\checkmarkc\checkmark$\checkmark:\checkmark$\checkmark^\checkmark\\checkmarkc\checkmarki\checkmarkr\checkmarkc\checkmark$\checkmarkc\checkmark$\checkmark^\checkmark\\checkmarkc\checkmarki\checkmarkr\checkmarkc\checkmark$\checkmarko\checkmark$\checkmark^\checkmark\\checkmarkc\checkmarki\checkmarkr\checkmarkc\checkmark$\checkmarku\checkmark$\checkmark^\checkmark\\checkmarkc\checkmarki\checkmarkr\checkmarkc\checkmark$\checkmarkp\checkmark$\checkmark^\checkmark\\checkmarkc\checkmarki\checkmarkr\checkmarkc\checkmark$\checkmarkl\checkmark$\checkmark^\checkmark\\checkmarkc\checkmarki\checkmarkr\checkmarkc\checkmark$\checkmarke\checkmark$\checkmark^\checkmark\\checkmarkc\checkmarki\checkmarkr\checkmarkc\checkmark$\checkmark_\checkmark$\checkmark^\checkmark\\checkmarkc\checkmarki\checkmarkr\checkmarkc\checkmark$\checkmarkm\checkmark$\checkmark^\checkmark\\checkmarkc\checkmarki\checkmarkr\checkmarkc\checkmark$\checkmarka\checkmark$\checkmark^\checkmark\\checkmarkc\checkmarki\checkmarkr\checkmarkc\checkmark$\checkmarki\checkmark$\checkmark^\checkmark\\checkmarkc\checkmarki\checkmarkr\checkmarkc\checkmark$\checkmarkn\checkmark$\checkmark^\checkmark\\checkmarkc\checkmarki\checkmarkr\checkmarkc\checkmark$\checkmarkt\checkmark$\checkmark^\checkmark\\checkmarkc\checkmarki\checkmarkr\checkmarkc\checkmark$\checkmarki\checkmark$\checkmark^\checkmark\\checkmarkc\checkmarki\checkmarkr\checkmarkc\checkmark$\checkmarke\checkmark$\checkmark^\checkmark\\checkmarkc\checkmarki\checkmarkr\checkmarkc\checkmark$\checkmarkn\checkmark$\checkmark^\checkmark\\checkmarkc\checkmarki\checkmarkr\checkmarkc\checkmark$\checkmark}\checkmark$\checkmark^\checkmark\\checkmarkc\checkmarki\checkmarkr\checkmarkc\checkmark$\checkmark
\checkmark$\checkmark^\checkmark\\checkmarkc\checkmarki\checkmarkr\checkmarkc\checkmark$\checkmark\\checkmark$\checkmark^\checkmark\\checkmarkc\checkmarki\checkmarkr\checkmarkc\checkmark$\checkmarke\checkmark$\checkmark^\checkmark\\checkmarkc\checkmarki\checkmarkr\checkmarkc\checkmark$\checkmarkn\checkmark$\checkmark^\checkmark\\checkmarkc\checkmarki\checkmarkr\checkmarkc\checkmark$\checkmarkd\checkmark$\checkmark^\checkmark\\checkmarkc\checkmarki\checkmarkr\checkmarkc\checkmark$\checkmark{\checkmark$\checkmark^\checkmark\\checkmarkc\checkmarki\checkmarkr\checkmarkc\checkmark$\checkmarke\checkmark$\checkmark^\checkmark\\checkmarkc\checkmarki\checkmarkr\checkmarkc\checkmark$\checkmarkq\checkmark$\checkmark^\checkmark\\checkmarkc\checkmarki\checkmarkr\checkmarkc\checkmark$\checkmarku\checkmark$\checkmark^\checkmark\\checkmarkc\checkmarki\checkmarkr\checkmarkc\checkmark$\checkmarka\checkmark$\checkmark^\checkmark\\checkmarkc\checkmarki\checkmarkr\checkmarkc\checkmark$\checkmarkt\checkmark$\checkmark^\checkmark\\checkmarkc\checkmarki\checkmarkr\checkmarkc\checkmark$\checkmarki\checkmark$\checkmark^\checkmark\\checkmarkc\checkmarki\checkmarkr\checkmarkc\checkmark$\checkmarko\checkmark$\checkmark^\checkmark\\checkmarkc\checkmarki\checkmarkr\checkmarkc\checkmark$\checkmarkn\checkmark$\checkmark^\checkmark\\checkmarkc\checkmarki\checkmarkr\checkmarkc\checkmark$\checkmark}\checkmark$\checkmark^\checkmark\\checkmarkc\checkmarki\checkmarkr\checkmarkc\checkmark$\checkmark
\checkmark$\checkmark^\checkmark\\checkmarkc\checkmarki\checkmarkr\checkmarkc\checkmark$\checkmark
\checkmark$\checkmark^\checkmark\\checkmarkc\checkmarki\checkmarkr\checkmarkc\checkmark$\checkmark\\checkmark$\checkmark^\checkmark\\checkmarkc\checkmarki\checkmarkr\checkmarkc\checkmark$\checkmarkt\checkmark$\checkmark^\checkmark\\checkmarkc\checkmarki\checkmarkr\checkmarkc\checkmark$\checkmarke\checkmark$\checkmark^\checkmark\\checkmarkc\checkmarki\checkmarkr\checkmarkc\checkmark$\checkmarkx\checkmark$\checkmark^\checkmark\\checkmarkc\checkmarki\checkmarkr\checkmarkc\checkmark$\checkmarkt\checkmark$\checkmark^\checkmark\\checkmarkc\checkmarki\checkmarkr\checkmarkc\checkmark$\checkmarkb\checkmark$\checkmark^\checkmark\\checkmarkc\checkmarki\checkmarkr\checkmarkc\checkmark$\checkmarkf\checkmark$\checkmark^\checkmark\\checkmarkc\checkmarki\checkmarkr\checkmarkc\checkmark$\checkmark{\checkmark$\checkmark^\checkmark\\checkmarkc\checkmarki\checkmarkr\checkmarkc\checkmark$\checkmarkV\checkmark$\checkmark^\checkmark\\checkmarkc\checkmarki\checkmarkr\checkmarkc\checkmark$\checkmarké\checkmark$\checkmark^\checkmark\\checkmarkc\checkmarki\checkmarkr\checkmarkc\checkmark$\checkmarkr\checkmark$\checkmark^\checkmark\\checkmarkc\checkmarki\checkmarkr\checkmarkc\checkmark$\checkmarki\checkmark$\checkmark^\checkmark\\checkmarkc\checkmarki\checkmarkr\checkmarkc\checkmark$\checkmarkf\checkmark$\checkmark^\checkmark\\checkmarkc\checkmarki\checkmarkr\checkmarkc\checkmark$\checkmarki\checkmark$\checkmark^\checkmark\\checkmarkc\checkmarki\checkmarkr\checkmarkc\checkmark$\checkmarkc\checkmark$\checkmark^\checkmark\\checkmarkc\checkmarki\checkmarkr\checkmarkc\checkmark$\checkmarka\checkmark$\checkmark^\checkmark\\checkmarkc\checkmarki\checkmarkr\checkmarkc\checkmark$\checkmarkt\checkmark$\checkmark^\checkmark\\checkmarkc\checkmarki\checkmarkr\checkmarkc\checkmark$\checkmarki\checkmark$\checkmark^\checkmark\\checkmarkc\checkmarki\checkmarkr\checkmarkc\checkmark$\checkmarko\checkmark$\checkmark^\checkmark\\checkmarkc\checkmarki\checkmarkr\checkmarkc\checkmark$\checkmarkn\checkmark$\checkmark^\checkmark\\checkmarkc\checkmarki\checkmarkr\checkmarkc\checkmark$\checkmark \checkmark$\checkmark^\checkmark\\checkmarkc\checkmarki\checkmarkr\checkmarkc\checkmark$\checkmark:\checkmark$\checkmark^\checkmark\\checkmarkc\checkmarki\checkmarkr\checkmarkc\checkmark$\checkmark}\checkmark$\checkmark^\checkmark\\checkmarkc\checkmarki\checkmarkr\checkmarkc\checkmark$\checkmark \checkmark$\checkmark^\checkmark\\checkmarkc\checkmarki\checkmarkr\checkmarkc\checkmark$\checkmarkL\checkmark$\checkmark^\checkmark\\checkmarkc\checkmarki\checkmarkr\checkmarkc\checkmark$\checkmarke\checkmark$\checkmark^\checkmark\\checkmarkc\checkmarki\checkmarkr\checkmarkc\checkmark$\checkmark \checkmark$\checkmark^\checkmark\\checkmarkc\checkmarki\checkmarkr\checkmarkc\checkmark$\checkmarkm\checkmark$\checkmark^\checkmark\\checkmarkc\checkmarki\checkmarkr\checkmarkc\checkmark$\checkmarko\checkmark$\checkmark^\checkmark\\checkmarkc\checkmarki\checkmarkr\checkmarkc\checkmark$\checkmarkt\checkmark$\checkmark^\checkmark\\checkmarkc\checkmarki\checkmarkr\checkmarkc\checkmark$\checkmarke\checkmark$\checkmark^\checkmark\\checkmarkc\checkmarki\checkmarkr\checkmarkc\checkmark$\checkmarku\checkmark$\checkmark^\checkmark\\checkmarkc\checkmarki\checkmarkr\checkmarkc\checkmark$\checkmarkr\checkmark$\checkmark^\checkmark\\checkmarkc\checkmarki\checkmarkr\checkmarkc\checkmark$\checkmark \checkmark$\checkmark^\checkmark\\checkmarkc\checkmarki\checkmarkr\checkmarkc\checkmark$\checkmarkN\checkmark$\checkmark^\checkmark\\checkmarkc\checkmarki\checkmarkr\checkmarkc\checkmark$\checkmarkE\checkmark$\checkmark^\checkmark\\checkmarkc\checkmarki\checkmarkr\checkmarkc\checkmark$\checkmarkM\checkmark$\checkmark^\checkmark\\checkmarkc\checkmarki\checkmarkr\checkmarkc\checkmark$\checkmarkA\checkmark$\checkmark^\checkmark\\checkmarkc\checkmarki\checkmarkr\checkmarkc\checkmark$\checkmark \checkmark$\checkmark^\checkmark\\checkmarkc\checkmarki\checkmarkr\checkmarkc\checkmark$\checkmark1\checkmark$\checkmark^\checkmark\\checkmarkc\checkmarki\checkmarkr\checkmarkc\checkmark$\checkmark7\checkmark$\checkmark^\checkmark\\checkmarkc\checkmarki\checkmarkr\checkmarkc\checkmark$\checkmark \checkmark$\checkmark^\checkmark\\checkmarkc\checkmarki\checkmarkr\checkmarkc\checkmark$\checkmark(\checkmark$\checkmark^\checkmark\\checkmarkc\checkmarki\checkmarkr\checkmarkc\checkmark$\checkmark1\checkmark$\checkmark^\checkmark\\checkmarkc\checkmarki\checkmarkr\checkmarkc\checkmark$\checkmark7\checkmark$\checkmark^\checkmark\\checkmarkc\checkmarki\checkmarkr\checkmarkc\checkmark$\checkmarkH\checkmark$\checkmark^\checkmark\\checkmarkc\checkmarki\checkmarkr\checkmarkc\checkmark$\checkmarkS\checkmark$\checkmark^\checkmark\\checkmarkc\checkmarki\checkmarkr\checkmarkc\checkmark$\checkmark4\checkmark$\checkmark^\checkmark\\checkmarkc\checkmarki\checkmarkr\checkmarkc\checkmark$\checkmark4\checkmark$\checkmark^\checkmark\\checkmarkc\checkmarki\checkmarkr\checkmarkc\checkmark$\checkmark0\checkmark$\checkmark^\checkmark\\checkmarkc\checkmarki\checkmarkr\checkmarkc\checkmark$\checkmark1\checkmark$\checkmark^\checkmark\\checkmarkc\checkmarki\checkmarkr\checkmarkc\checkmark$\checkmark)\checkmark$\checkmark^\checkmark\\checkmarkc\checkmarki\checkmarkr\checkmarkc\checkmark$\checkmark \checkmark$\checkmark^\checkmark\\checkmarkc\checkmarki\checkmarkr\checkmarkc\checkmark$\checkmarkf\checkmark$\checkmark^\checkmark\\checkmarkc\checkmarki\checkmarkr\checkmarkc\checkmark$\checkmarko\checkmark$\checkmark^\checkmark\\checkmarkc\checkmarki\checkmarkr\checkmarkc\checkmark$\checkmarku\checkmark$\checkmark^\checkmark\\checkmarkc\checkmarki\checkmarkr\checkmarkc\checkmark$\checkmarkr\checkmark$\checkmark^\checkmark\\checkmarkc\checkmarki\checkmarkr\checkmarkc\checkmark$\checkmarkn\checkmark$\checkmark^\checkmark\\checkmarkc\checkmarki\checkmarkr\checkmarkc\checkmark$\checkmarki\checkmark$\checkmark^\checkmark\\checkmarkc\checkmarki\checkmarkr\checkmarkc\checkmark$\checkmarkt\checkmark$\checkmark^\checkmark\\checkmarkc\checkmarki\checkmarkr\checkmarkc\checkmark$\checkmark \checkmark$\checkmark^\checkmark\\checkmarkc\checkmarki\checkmarkr\checkmarkc\checkmark$\checkmark$\checkmark$\checkmark^\checkmark\\checkmarkc\checkmarki\checkmarkr\checkmarkc\checkmark$\checkmarkC\checkmark$\checkmark^\checkmark\\checkmarkc\checkmarki\checkmarkr\checkmarkc\checkmark$\checkmark_\checkmark$\checkmark^\checkmark\\checkmarkc\checkmarki\checkmarkr\checkmarkc\checkmark$\checkmark{\checkmark$\checkmark^\checkmark\\checkmarkc\checkmarki\checkmarkr\checkmarkc\checkmark$\checkmarkm\checkmark$\checkmark^\checkmark\\checkmarkc\checkmarki\checkmarkr\checkmarkc\checkmark$\checkmarka\checkmark$\checkmark^\checkmark\\checkmarkc\checkmarki\checkmarkr\checkmarkc\checkmark$\checkmarki\checkmark$\checkmark^\checkmark\\checkmarkc\checkmarki\checkmarkr\checkmarkc\checkmark$\checkmarkn\checkmark$\checkmark^\checkmark\\checkmarkc\checkmarki\checkmarkr\checkmarkc\checkmark$\checkmarkt\checkmark$\checkmark^\checkmark\\checkmarkc\checkmarki\checkmarkr\checkmarkc\checkmark$\checkmarki\checkmark$\checkmark^\checkmark\\checkmarkc\checkmarki\checkmarkr\checkmarkc\checkmark$\checkmarke\checkmark$\checkmark^\checkmark\\checkmarkc\checkmarki\checkmarkr\checkmarkc\checkmark$\checkmarkn\checkmark$\checkmark^\checkmark\\checkmarkc\checkmarki\checkmarkr\checkmarkc\checkmark$\checkmark}\checkmark$\checkmark^\checkmark\\checkmarkc\checkmarki\checkmarkr\checkmarkc\checkmark$\checkmark \checkmark$\checkmark^\checkmark\\checkmarkc\checkmarki\checkmarkr\checkmarkc\checkmark$\checkmark=\checkmark$\checkmark^\checkmark\\checkmarkc\checkmarki\checkmarkr\checkmarkc\checkmark$\checkmark \checkmark$\checkmark^\checkmark\\checkmarkc\checkmarki\checkmarkr\checkmarkc\checkmark$\checkmark0\checkmark$\checkmark^\checkmark\\checkmarkc\checkmarki\checkmarkr\checkmarkc\checkmark$\checkmark{\checkmark$\checkmark^\checkmark\\checkmarkc\checkmarki\checkmarkr\checkmarkc\checkmark$\checkmark,\checkmark$\checkmark^\checkmark\\checkmarkc\checkmarki\checkmarkr\checkmarkc\checkmark$\checkmark}\checkmark$\checkmark^\checkmark\\checkmarkc\checkmarki\checkmarkr\checkmarkc\checkmark$\checkmark4\checkmark$\checkmark^\checkmark\\checkmarkc\checkmarki\checkmarkr\checkmarkc\checkmark$\checkmark2\checkmark$\checkmark^\checkmark\\checkmarkc\checkmarki\checkmarkr\checkmarkc\checkmark$\checkmark$\checkmark$\checkmark^\checkmark\\checkmarkc\checkmarki\checkmarkr\checkmarkc\checkmark$\checkmark \checkmark$\checkmark^\checkmark\\checkmarkc\checkmarki\checkmarkr\checkmarkc\checkmark$\checkmarkN\checkmark$\checkmark^\checkmark\\checkmarkc\checkmarki\checkmarkr\checkmarkc\checkmark$\checkmark$\checkmark$\checkmark^\checkmark\\checkmarkc\checkmarki\checkmarkr\checkmarkc\checkmark$\checkmark\\checkmark$\checkmark^\checkmark\\checkmarkc\checkmarki\checkmarkr\checkmarkc\checkmark$\checkmarkc\checkmark$\checkmark^\checkmark\\checkmarkc\checkmarki\checkmarkr\checkmarkc\checkmark$\checkmarkd\checkmark$\checkmark^\checkmark\\checkmarkc\checkmarki\checkmarkr\checkmarkc\checkmark$\checkmarko\checkmark$\checkmark^\checkmark\\checkmarkc\checkmarki\checkmarkr\checkmarkc\checkmark$\checkmarkt\checkmark$\checkmark^\checkmark\\checkmarkc\checkmarki\checkmarkr\checkmarkc\checkmark$\checkmark$\checkmark$\checkmark^\checkmark\\checkmarkc\checkmarki\checkmarkr\checkmarkc\checkmark$\checkmarkm\checkmark$\checkmark^\checkmark\\checkmarkc\checkmarki\checkmarkr\checkmarkc\checkmark$\checkmark.\checkmark$\checkmark^\checkmark\\checkmarkc\checkmarki\checkmarkr\checkmarkc\checkmark$\checkmark
\checkmark$\checkmark^\checkmark\\checkmarkc\checkmarki\checkmarkr\checkmarkc\checkmark$\checkmark\\checkmark$\checkmark^\checkmark\\checkmarkc\checkmarki\checkmarkr\checkmarkc\checkmark$\checkmarkb\checkmark$\checkmark^\checkmark\\checkmarkc\checkmarki\checkmarkr\checkmarkc\checkmark$\checkmarke\checkmark$\checkmark^\checkmark\\checkmarkc\checkmarki\checkmarkr\checkmarkc\checkmark$\checkmarkg\checkmark$\checkmark^\checkmark\\checkmarkc\checkmarki\checkmarkr\checkmarkc\checkmark$\checkmarki\checkmark$\checkmark^\checkmark\\checkmarkc\checkmarki\checkmarkr\checkmarkc\checkmark$\checkmarkn\checkmark$\checkmark^\checkmark\\checkmarkc\checkmarki\checkmarkr\checkmarkc\checkmark$\checkmark{\checkmark$\checkmark^\checkmark\\checkmarkc\checkmarki\checkmarkr\checkmarkc\checkmark$\checkmarke\checkmark$\checkmark^\checkmark\\checkmarkc\checkmarki\checkmarkr\checkmarkc\checkmark$\checkmarkq\checkmark$\checkmark^\checkmark\\checkmarkc\checkmarki\checkmarkr\checkmarkc\checkmark$\checkmarku\checkmark$\checkmark^\checkmark\\checkmarkc\checkmarki\checkmarkr\checkmarkc\checkmark$\checkmarka\checkmark$\checkmark^\checkmark\\checkmarkc\checkmarki\checkmarkr\checkmarkc\checkmark$\checkmarkt\checkmark$\checkmark^\checkmark\\checkmarkc\checkmarki\checkmarkr\checkmarkc\checkmark$\checkmarki\checkmark$\checkmark^\checkmark\\checkmarkc\checkmarki\checkmarkr\checkmarkc\checkmark$\checkmarko\checkmark$\checkmark^\checkmark\\checkmarkc\checkmarki\checkmarkr\checkmarkc\checkmark$\checkmarkn\checkmark$\checkmark^\checkmark\\checkmarkc\checkmarki\checkmarkr\checkmarkc\checkmark$\checkmark}\checkmark$\checkmark^\checkmark\\checkmarkc\checkmarki\checkmarkr\checkmarkc\checkmark$\checkmark
\checkmark$\checkmark^\checkmark\\checkmarkc\checkmarki\checkmarkr\checkmarkc\checkmark$\checkmark \checkmark$\checkmark^\checkmark\\checkmarkc\checkmarki\checkmarkr\checkmarkc\checkmark$\checkmark \checkmark$\checkmark^\checkmark\\checkmarkc\checkmarki\checkmarkr\checkmarkc\checkmark$\checkmark \checkmark$\checkmark^\checkmark\\checkmarkc\checkmarki\checkmarkr\checkmarkc\checkmark$\checkmark \checkmark$\checkmark^\checkmark\\checkmarkc\checkmarki\checkmarkr\checkmarkc\checkmark$\checkmarks\checkmark$\checkmark^\checkmark\\checkmarkc\checkmarki\checkmarkr\checkmarkc\checkmark$\checkmark_\checkmark$\checkmark^\checkmark\\checkmarkc\checkmarki\checkmarkr\checkmarkc\checkmark$\checkmark{\checkmark$\checkmark^\checkmark\\checkmarkc\checkmarki\checkmarkr\checkmarkc\checkmark$\checkmarkm\checkmark$\checkmark^\checkmark\\checkmarkc\checkmarki\checkmarkr\checkmarkc\checkmark$\checkmarko\checkmark$\checkmark^\checkmark\\checkmarkc\checkmarki\checkmarkr\checkmarkc\checkmark$\checkmarkt\checkmark$\checkmark^\checkmark\\checkmarkc\checkmarki\checkmarkr\checkmarkc\checkmark$\checkmarke\checkmark$\checkmark^\checkmark\\checkmarkc\checkmarki\checkmarkr\checkmarkc\checkmark$\checkmarku\checkmark$\checkmark^\checkmark\\checkmarkc\checkmarki\checkmarkr\checkmarkc\checkmark$\checkmarkr\checkmark$\checkmark^\checkmark\\checkmarkc\checkmarki\checkmarkr\checkmarkc\checkmark$\checkmark}\checkmark$\checkmark^\checkmark\\checkmarkc\checkmarki\checkmarkr\checkmarkc\checkmark$\checkmark \checkmark$\checkmark^\checkmark\\checkmarkc\checkmarki\checkmarkr\checkmarkc\checkmark$\checkmark=\checkmark$\checkmark^\checkmark\\checkmarkc\checkmarki\checkmarkr\checkmarkc\checkmark$\checkmark \checkmark$\checkmark^\checkmark\\checkmarkc\checkmarki\checkmarkr\checkmarkc\checkmark$\checkmark\\checkmark$\checkmark^\checkmark\\checkmarkc\checkmarki\checkmarkr\checkmarkc\checkmark$\checkmarkf\checkmark$\checkmark^\checkmark\\checkmarkc\checkmarki\checkmarkr\checkmarkc\checkmark$\checkmarkr\checkmark$\checkmark^\checkmark\\checkmarkc\checkmarki\checkmarkr\checkmarkc\checkmark$\checkmarka\checkmark$\checkmark^\checkmark\\checkmarkc\checkmarki\checkmarkr\checkmarkc\checkmark$\checkmarkc\checkmark$\checkmark^\checkmark\\checkmarkc\checkmarki\checkmarkr\checkmarkc\checkmark$\checkmark{\checkmark$\checkmark^\checkmark\\checkmarkc\checkmarki\checkmarkr\checkmarkc\checkmark$\checkmarkC\checkmark$\checkmark^\checkmark\\checkmarkc\checkmarki\checkmarkr\checkmarkc\checkmark$\checkmark_\checkmark$\checkmark^\checkmark\\checkmarkc\checkmarki\checkmarkr\checkmarkc\checkmark$\checkmark{\checkmark$\checkmark^\checkmark\\checkmarkc\checkmarki\checkmarkr\checkmarkc\checkmark$\checkmarkm\checkmark$\checkmark^\checkmark\\checkmarkc\checkmarki\checkmarkr\checkmarkc\checkmark$\checkmarka\checkmark$\checkmark^\checkmark\\checkmarkc\checkmarki\checkmarkr\checkmarkc\checkmark$\checkmarki\checkmark$\checkmark^\checkmark\\checkmarkc\checkmarki\checkmarkr\checkmarkc\checkmark$\checkmarkn\checkmark$\checkmark^\checkmark\\checkmarkc\checkmarki\checkmarkr\checkmarkc\checkmark$\checkmarkt\checkmark$\checkmark^\checkmark\\checkmarkc\checkmarki\checkmarkr\checkmarkc\checkmark$\checkmarki\checkmark$\checkmark^\checkmark\\checkmarkc\checkmarki\checkmarkr\checkmarkc\checkmark$\checkmarke\checkmark$\checkmark^\checkmark\\checkmarkc\checkmarki\checkmarkr\checkmarkc\checkmark$\checkmarkn\checkmark$\checkmark^\checkmark\\checkmarkc\checkmarki\checkmarkr\checkmarkc\checkmark$\checkmark}\checkmark$\checkmark^\checkmark\\checkmarkc\checkmarki\checkmarkr\checkmarkc\checkmark$\checkmark}\checkmark$\checkmark^\checkmark\\checkmarkc\checkmarki\checkmarkr\checkmarkc\checkmark$\checkmark{\checkmark$\checkmark^\checkmark\\checkmarkc\checkmarki\checkmarkr\checkmarkc\checkmark$\checkmarkC\checkmark$\checkmark^\checkmark\\checkmarkc\checkmarki\checkmarkr\checkmarkc\checkmark$\checkmark_\checkmark$\checkmark^\checkmark\\checkmarkc\checkmarki\checkmarkr\checkmarkc\checkmark$\checkmark{\checkmark$\checkmark^\checkmark\\checkmarkc\checkmarki\checkmarkr\checkmarkc\checkmark$\checkmarkm\checkmark$\checkmark^\checkmark\\checkmarkc\checkmarki\checkmarkr\checkmarkc\checkmark$\checkmarka\checkmark$\checkmark^\checkmark\\checkmarkc\checkmarki\checkmarkr\checkmarkc\checkmark$\checkmarki\checkmark$\checkmark^\checkmark\\checkmarkc\checkmarki\checkmarkr\checkmarkc\checkmark$\checkmarkn\checkmark$\checkmark^\checkmark\\checkmarkc\checkmarki\checkmarkr\checkmarkc\checkmark$\checkmarkt\checkmark$\checkmark^\checkmark\\checkmarkc\checkmarki\checkmarkr\checkmarkc\checkmark$\checkmarki\checkmark$\checkmark^\checkmark\\checkmarkc\checkmarki\checkmarkr\checkmarkc\checkmark$\checkmarke\checkmark$\checkmark^\checkmark\\checkmarkc\checkmarki\checkmarkr\checkmarkc\checkmark$\checkmarkn\checkmark$\checkmark^\checkmark\\checkmarkc\checkmarki\checkmarkr\checkmarkc\checkmark$\checkmark\\checkmark$\checkmark^\checkmark\\checkmarkc\checkmarki\checkmarkr\checkmarkc\checkmark$\checkmark_\checkmark$\checkmark^\checkmark\\checkmarkc\checkmarki\checkmarkr\checkmarkc\checkmark$\checkmarkr\checkmark$\checkmark^\checkmark\\checkmarkc\checkmarki\checkmarkr\checkmarkc\checkmark$\checkmarke\checkmark$\checkmark^\checkmark\\checkmarkc\checkmarki\checkmarkr\checkmarkc\checkmark$\checkmarkq\checkmark$\checkmark^\checkmark\\checkmarkc\checkmarki\checkmarkr\checkmarkc\checkmark$\checkmark}\checkmark$\checkmark^\checkmark\\checkmarkc\checkmarki\checkmarkr\checkmarkc\checkmark$\checkmark}\checkmark$\checkmark^\checkmark\\checkmarkc\checkmarki\checkmarkr\checkmarkc\checkmark$\checkmark \checkmark$\checkmark^\checkmark\\checkmarkc\checkmarki\checkmarkr\checkmarkc\checkmark$\checkmark=\checkmark$\checkmark^\checkmark\\checkmarkc\checkmarki\checkmarkr\checkmarkc\checkmark$\checkmark \checkmark$\checkmark^\checkmark\\checkmarkc\checkmarki\checkmarkr\checkmarkc\checkmark$\checkmark\\checkmark$\checkmark^\checkmark\\checkmarkc\checkmarki\checkmarkr\checkmarkc\checkmark$\checkmarkf\checkmark$\checkmark^\checkmark\\checkmarkc\checkmarki\checkmarkr\checkmarkc\checkmark$\checkmarkr\checkmark$\checkmark^\checkmark\\checkmarkc\checkmarki\checkmarkr\checkmarkc\checkmark$\checkmarka\checkmark$\checkmark^\checkmark\\checkmarkc\checkmarki\checkmarkr\checkmarkc\checkmark$\checkmarkc\checkmark$\checkmark^\checkmark\\checkmarkc\checkmarki\checkmarkr\checkmarkc\checkmark$\checkmark{\checkmark$\checkmark^\checkmark\\checkmarkc\checkmarki\checkmarkr\checkmarkc\checkmark$\checkmark0\checkmark$\checkmark^\checkmark\\checkmarkc\checkmarki\checkmarkr\checkmarkc\checkmark$\checkmark{\checkmark$\checkmark^\checkmark\\checkmarkc\checkmarki\checkmarkr\checkmarkc\checkmark$\checkmark,\checkmark$\checkmark^\checkmark\\checkmarkc\checkmarki\checkmarkr\checkmarkc\checkmark$\checkmark}\checkmark$\checkmark^\checkmark\\checkmarkc\checkmarki\checkmarkr\checkmarkc\checkmark$\checkmark4\checkmark$\checkmark^\checkmark\\checkmarkc\checkmarki\checkmarkr\checkmarkc\checkmark$\checkmark2\checkmark$\checkmark^\checkmark\\checkmarkc\checkmarki\checkmarkr\checkmarkc\checkmark$\checkmark}\checkmark$\checkmark^\checkmark\\checkmarkc\checkmarki\checkmarkr\checkmarkc\checkmark$\checkmark{\checkmark$\checkmark^\checkmark\\checkmarkc\checkmarki\checkmarkr\checkmarkc\checkmark$\checkmark0\checkmark$\checkmark^\checkmark\\checkmarkc\checkmarki\checkmarkr\checkmarkc\checkmark$\checkmark{\checkmark$\checkmark^\checkmark\\checkmarkc\checkmarki\checkmarkr\checkmarkc\checkmark$\checkmark,\checkmark$\checkmark^\checkmark\\checkmarkc\checkmarki\checkmarkr\checkmarkc\checkmark$\checkmark}\checkmark$\checkmark^\checkmark\\checkmarkc\checkmarki\checkmarkr\checkmarkc\checkmark$\checkmark3\checkmark$\checkmark^\checkmark\\checkmarkc\checkmarki\checkmarkr\checkmarkc\checkmark$\checkmark3\checkmark$\checkmark^\checkmark\\checkmarkc\checkmarki\checkmarkr\checkmarkc\checkmark$\checkmark3\checkmark$\checkmark^\checkmark\\checkmarkc\checkmarki\checkmarkr\checkmarkc\checkmark$\checkmark}\checkmark$\checkmark^\checkmark\\checkmarkc\checkmarki\checkmarkr\checkmarkc\checkmark$\checkmark \checkmark$\checkmark^\checkmark\\checkmarkc\checkmarki\checkmarkr\checkmarkc\checkmark$\checkmark=\checkmark$\checkmark^\checkmark\\checkmarkc\checkmarki\checkmarkr\checkmarkc\checkmark$\checkmark \checkmark$\checkmark^\checkmark\\checkmarkc\checkmarki\checkmarkr\checkmarkc\checkmark$\checkmark\\checkmark$\checkmark^\checkmark\\checkmarkc\checkmarki\checkmarkr\checkmarkc\checkmark$\checkmarkb\checkmark$\checkmark^\checkmark\\checkmarkc\checkmarki\checkmarkr\checkmarkc\checkmark$\checkmarko\checkmark$\checkmark^\checkmark\\checkmarkc\checkmarki\checkmarkr\checkmarkc\checkmark$\checkmarkx\checkmark$\checkmark^\checkmark\\checkmarkc\checkmarki\checkmarkr\checkmarkc\checkmark$\checkmarke\checkmark$\checkmark^\checkmark\\checkmarkc\checkmarki\checkmarkr\checkmarkc\checkmark$\checkmarkd\checkmark$\checkmark^\checkmark\\checkmarkc\checkmarki\checkmarkr\checkmarkc\checkmark$\checkmark{\checkmark$\checkmark^\checkmark\\checkmarkc\checkmarki\checkmarkr\checkmarkc\checkmark$\checkmark1\checkmark$\checkmark^\checkmark\\checkmarkc\checkmarki\checkmarkr\checkmarkc\checkmark$\checkmark{\checkmark$\checkmark^\checkmark\\checkmarkc\checkmarki\checkmarkr\checkmarkc\checkmark$\checkmark,\checkmark$\checkmark^\checkmark\\checkmarkc\checkmarki\checkmarkr\checkmarkc\checkmark$\checkmark}\checkmark$\checkmark^\checkmark\\checkmarkc\checkmarki\checkmarkr\checkmarkc\checkmark$\checkmark2\checkmark$\checkmark^\checkmark\\checkmarkc\checkmarki\checkmarkr\checkmarkc\checkmark$\checkmark6\checkmark$\checkmark^\checkmark\\checkmarkc\checkmarki\checkmarkr\checkmarkc\checkmark$\checkmark \checkmark$\checkmark^\checkmark\\checkmarkc\checkmarki\checkmarkr\checkmarkc\checkmark$\checkmark>\checkmark$\checkmark^\checkmark\\checkmarkc\checkmarki\checkmarkr\checkmarkc\checkmark$\checkmark \checkmark$\checkmark^\checkmark\\checkmarkc\checkmarki\checkmarkr\checkmarkc\checkmark$\checkmark1\checkmark$\checkmark^\checkmark\\checkmarkc\checkmarki\checkmarkr\checkmarkc\checkmark$\checkmark{\checkmark$\checkmark^\checkmark\\checkmarkc\checkmarki\checkmarkr\checkmarkc\checkmark$\checkmark,\checkmark$\checkmark^\checkmark\\checkmarkc\checkmarki\checkmarkr\checkmarkc\checkmark$\checkmark}\checkmark$\checkmark^\checkmark\\checkmarkc\checkmarki\checkmarkr\checkmarkc\checkmark$\checkmark0\checkmark$\checkmark^\checkmark\\checkmarkc\checkmarki\checkmarkr\checkmarkc\checkmark$\checkmark}\checkmark$\checkmark^\checkmark\\checkmarkc\checkmarki\checkmarkr\checkmarkc\checkmark$\checkmark
\checkmark$\checkmark^\checkmark\\checkmarkc\checkmarki\checkmarkr\checkmarkc\checkmark$\checkmark\\checkmark$\checkmark^\checkmark\\checkmarkc\checkmarki\checkmarkr\checkmarkc\checkmark$\checkmarke\checkmark$\checkmark^\checkmark\\checkmarkc\checkmarki\checkmarkr\checkmarkc\checkmark$\checkmarkn\checkmark$\checkmark^\checkmark\\checkmarkc\checkmarki\checkmarkr\checkmarkc\checkmark$\checkmarkd\checkmark$\checkmark^\checkmark\\checkmarkc\checkmarki\checkmarkr\checkmarkc\checkmark$\checkmark{\checkmark$\checkmark^\checkmark\\checkmarkc\checkmarki\checkmarkr\checkmarkc\checkmark$\checkmarke\checkmark$\checkmark^\checkmark\\checkmarkc\checkmarki\checkmarkr\checkmarkc\checkmark$\checkmarkq\checkmark$\checkmark^\checkmark\\checkmarkc\checkmarki\checkmarkr\checkmarkc\checkmark$\checkmarku\checkmark$\checkmark^\checkmark\\checkmarkc\checkmarki\checkmarkr\checkmarkc\checkmark$\checkmarka\checkmark$\checkmark^\checkmark\\checkmarkc\checkmarki\checkmarkr\checkmarkc\checkmark$\checkmarkt\checkmark$\checkmark^\checkmark\\checkmarkc\checkmarki\checkmarkr\checkmarkc\checkmark$\checkmarki\checkmark$\checkmark^\checkmark\\checkmarkc\checkmarki\checkmarkr\checkmarkc\checkmark$\checkmarko\checkmark$\checkmark^\checkmark\\checkmarkc\checkmarki\checkmarkr\checkmarkc\checkmark$\checkmarkn\checkmark$\checkmark^\checkmark\\checkmarkc\checkmarki\checkmarkr\checkmarkc\checkmark$\checkmark}\checkmark$\checkmark^\checkmark\\checkmarkc\checkmarki\checkmarkr\checkmarkc\checkmark$\checkmark
\checkmark$\checkmark^\checkmark\\checkmarkc\checkmarki\checkmarkr\checkmarkc\checkmark$\checkmarkL\checkmark$\checkmark^\checkmark\\checkmarkc\checkmarki\checkmarkr\checkmarkc\checkmark$\checkmarke\checkmark$\checkmark^\checkmark\\checkmarkc\checkmarki\checkmarkr\checkmarkc\checkmark$\checkmark \checkmark$\checkmark^\checkmark\\checkmarkc\checkmarki\checkmarkr\checkmarkc\checkmark$\checkmarkm\checkmark$\checkmark^\checkmark\\checkmarkc\checkmarki\checkmarkr\checkmarkc\checkmark$\checkmarko\checkmark$\checkmark^\checkmark\\checkmarkc\checkmarki\checkmarkr\checkmarkc\checkmark$\checkmarkt\checkmark$\checkmark^\checkmark\\checkmarkc\checkmarki\checkmarkr\checkmarkc\checkmark$\checkmarke\checkmark$\checkmark^\checkmark\\checkmarkc\checkmarki\checkmarkr\checkmarkc\checkmark$\checkmarku\checkmark$\checkmark^\checkmark\\checkmarkc\checkmarki\checkmarkr\checkmarkc\checkmark$\checkmarkr\checkmark$\checkmark^\checkmark\\checkmarkc\checkmarki\checkmarkr\checkmarkc\checkmark$\checkmark \checkmark$\checkmark^\checkmark\\checkmarkc\checkmarki\checkmarkr\checkmarkc\checkmark$\checkmarkN\checkmark$\checkmark^\checkmark\\checkmarkc\checkmarki\checkmarkr\checkmarkc\checkmark$\checkmarkE\checkmark$\checkmark^\checkmark\\checkmarkc\checkmarki\checkmarkr\checkmarkc\checkmark$\checkmarkM\checkmark$\checkmark^\checkmark\\checkmarkc\checkmarki\checkmarkr\checkmarkc\checkmark$\checkmarkA\checkmark$\checkmark^\checkmark\\checkmarkc\checkmarki\checkmarkr\checkmarkc\checkmark$\checkmark \checkmark$\checkmark^\checkmark\\checkmarkc\checkmarki\checkmarkr\checkmarkc\checkmark$\checkmark1\checkmark$\checkmark^\checkmark\\checkmarkc\checkmarki\checkmarkr\checkmarkc\checkmark$\checkmark7\checkmark$\checkmark^\checkmark\\checkmarkc\checkmarki\checkmarkr\checkmarkc\checkmark$\checkmark \checkmark$\checkmark^\checkmark\\checkmarkc\checkmarki\checkmarkr\checkmarkc\checkmark$\checkmarks\checkmark$\checkmark^\checkmark\\checkmarkc\checkmarki\checkmarkr\checkmarkc\checkmark$\checkmarka\checkmark$\checkmark^\checkmark\\checkmarkc\checkmarki\checkmarkr\checkmarkc\checkmark$\checkmarkt\checkmark$\checkmark^\checkmark\\checkmarkc\checkmarki\checkmarkr\checkmarkc\checkmark$\checkmarki\checkmark$\checkmark^\checkmark\\checkmarkc\checkmarki\checkmarkr\checkmarkc\checkmark$\checkmarks\checkmark$\checkmark^\checkmark\\checkmarkc\checkmarki\checkmarkr\checkmarkc\checkmark$\checkmarkf\checkmark$\checkmark^\checkmark\\checkmarkc\checkmarki\checkmarkr\checkmarkc\checkmark$\checkmarka\checkmark$\checkmark^\checkmark\\checkmarkc\checkmarki\checkmarkr\checkmarkc\checkmark$\checkmarki\checkmark$\checkmark^\checkmark\\checkmarkc\checkmarki\checkmarkr\checkmarkc\checkmark$\checkmarkt\checkmark$\checkmark^\checkmark\\checkmarkc\checkmarki\checkmarkr\checkmarkc\checkmark$\checkmark \checkmark$\checkmark^\checkmark\\checkmarkc\checkmarki\checkmarkr\checkmarkc\checkmark$\checkmarkl\checkmark$\checkmark^\checkmark\\checkmarkc\checkmarki\checkmarkr\checkmarkc\checkmark$\checkmarka\checkmark$\checkmark^\checkmark\\checkmarkc\checkmarki\checkmarkr\checkmarkc\checkmark$\checkmark \checkmark$\checkmark^\checkmark\\checkmarkc\checkmarki\checkmarkr\checkmarkc\checkmark$\checkmarkc\checkmark$\checkmark^\checkmark\\checkmarkc\checkmarki\checkmarkr\checkmarkc\checkmark$\checkmarko\checkmark$\checkmark^\checkmark\\checkmarkc\checkmarki\checkmarkr\checkmarkc\checkmark$\checkmarkn\checkmark$\checkmark^\checkmark\\checkmarkc\checkmarki\checkmarkr\checkmarkc\checkmark$\checkmarkd\checkmark$\checkmark^\checkmark\\checkmarkc\checkmarki\checkmarkr\checkmarkc\checkmark$\checkmarki\checkmark$\checkmark^\checkmark\\checkmarkc\checkmarki\checkmarkr\checkmarkc\checkmark$\checkmarkt\checkmark$\checkmark^\checkmark\\checkmarkc\checkmarki\checkmarkr\checkmarkc\checkmark$\checkmarki\checkmark$\checkmark^\checkmark\\checkmarkc\checkmarki\checkmarkr\checkmarkc\checkmark$\checkmarko\checkmark$\checkmark^\checkmark\\checkmarkc\checkmarki\checkmarkr\checkmarkc\checkmark$\checkmarkn\checkmark$\checkmark^\checkmark\\checkmarkc\checkmarki\checkmarkr\checkmarkc\checkmark$\checkmark \checkmark$\checkmark^\checkmark\\checkmarkc\checkmarki\checkmarkr\checkmarkc\checkmark$\checkmarkd\checkmark$\checkmark^\checkmark\\checkmarkc\checkmarki\checkmarkr\checkmarkc\checkmark$\checkmarke\checkmark$\checkmark^\checkmark\\checkmarkc\checkmarki\checkmarkr\checkmarkc\checkmark$\checkmark \checkmark$\checkmark^\checkmark\\checkmarkc\checkmarki\checkmarkr\checkmarkc\checkmark$\checkmarkc\checkmark$\checkmark^\checkmark\\checkmarkc\checkmarki\checkmarkr\checkmarkc\checkmark$\checkmarko\checkmark$\checkmark^\checkmark\\checkmarkc\checkmarki\checkmarkr\checkmarkc\checkmark$\checkmarku\checkmark$\checkmark^\checkmark\\checkmarkc\checkmarki\checkmarkr\checkmarkc\checkmark$\checkmarkp\checkmark$\checkmark^\checkmark\\checkmarkc\checkmarki\checkmarkr\checkmarkc\checkmark$\checkmarkl\checkmark$\checkmark^\checkmark\\checkmarkc\checkmarki\checkmarkr\checkmarkc\checkmark$\checkmarke\checkmark$\checkmark^\checkmark\\checkmarkc\checkmarki\checkmarkr\checkmarkc\checkmark$\checkmark \checkmark$\checkmark^\checkmark\\checkmarkc\checkmarki\checkmarkr\checkmarkc\checkmark$\checkmarkp\checkmark$\checkmark^\checkmark\\checkmarkc\checkmarki\checkmarkr\checkmarkc\checkmark$\checkmarko\checkmark$\checkmark^\checkmark\\checkmarkc\checkmarki\checkmarkr\checkmarkc\checkmark$\checkmarku\checkmark$\checkmark^\checkmark\\checkmarkc\checkmarki\checkmarkr\checkmarkc\checkmark$\checkmarkr\checkmark$\checkmark^\checkmark\\checkmarkc\checkmarki\checkmarkr\checkmarkc\checkmark$\checkmark \checkmark$\checkmark^\checkmark\\checkmarkc\checkmarki\checkmarkr\checkmarkc\checkmark$\checkmarkl\checkmark$\checkmark^\checkmark\\checkmarkc\checkmarki\checkmarkr\checkmarkc\checkmark$\checkmarke\checkmark$\checkmark^\checkmark\\checkmarkc\checkmarki\checkmarkr\checkmarkc\checkmark$\checkmarks\checkmark$\checkmark^\checkmark\\checkmarkc\checkmarki\checkmarkr\checkmarkc\checkmark$\checkmark \checkmark$\checkmark^\checkmark\\checkmarkc\checkmarki\checkmarkr\checkmarkc\checkmark$\checkmarka\checkmark$\checkmark^\checkmark\\checkmarkc\checkmarki\checkmarkr\checkmarkc\checkmark$\checkmarkx\checkmark$\checkmark^\checkmark\\checkmarkc\checkmarki\checkmarkr\checkmarkc\checkmark$\checkmarke\checkmark$\checkmark^\checkmark\\checkmarkc\checkmarki\checkmarkr\checkmarkc\checkmark$\checkmarks\checkmark$\checkmark^\checkmark\\checkmarkc\checkmarki\checkmarkr\checkmarkc\checkmark$\checkmark \checkmark$\checkmark^\checkmark\\checkmarkc\checkmarki\checkmarkr\checkmarkc\checkmark$\checkmarkX\checkmark$\checkmark^\checkmark\\checkmarkc\checkmarki\checkmarkr\checkmarkc\checkmark$\checkmark,\checkmark$\checkmark^\checkmark\\checkmarkc\checkmarki\checkmarkr\checkmarkc\checkmark$\checkmark \checkmark$\checkmark^\checkmark\\checkmarkc\checkmarki\checkmarkr\checkmarkc\checkmark$\checkmarkY\checkmark$\checkmark^\checkmark\\checkmarkc\checkmarki\checkmarkr\checkmarkc\checkmark$\checkmark \checkmark$\checkmark^\checkmark\\checkmarkc\checkmarki\checkmarkr\checkmarkc\checkmark$\checkmarke\checkmark$\checkmark^\checkmark\\checkmarkc\checkmarki\checkmarkr\checkmarkc\checkmark$\checkmarkt\checkmark$\checkmark^\checkmark\\checkmarkc\checkmarki\checkmarkr\checkmarkc\checkmark$\checkmark \checkmark$\checkmark^\checkmark\\checkmarkc\checkmarki\checkmarkr\checkmarkc\checkmark$\checkmarkZ\checkmark$\checkmark^\checkmark\\checkmarkc\checkmarki\checkmarkr\checkmarkc\checkmark$\checkmark.\checkmark$\checkmark^\checkmark\\checkmarkc\checkmarki\checkmarkr\checkmarkc\checkmark$\checkmark
\checkmark$\checkmark^\checkmark\\checkmarkc\checkmarki\checkmarkr\checkmarkc\checkmark$\checkmark
\checkmark$\checkmark^\checkmark\\checkmarkc\checkmarki\checkmarkr\checkmarkc\checkmark$\checkmark%\checkmark$\checkmark^\checkmark\\checkmarkc\checkmarki\checkmarkr\checkmarkc\checkmark$\checkmark \checkmark$\checkmark^\checkmark\\checkmarkc\checkmarki\checkmarkr\checkmarkc\checkmark$\checkmark=\checkmark$\checkmark^\checkmark\\checkmarkc\checkmarki\checkmarkr\checkmarkc\checkmark$\checkmark=\checkmark$\checkmark^\checkmark\\checkmarkc\checkmarki\checkmarkr\checkmarkc\checkmark$\checkmark=\checkmark$\checkmark^\checkmark\\checkmarkc\checkmarki\checkmarkr\checkmarkc\checkmark$\checkmark=\checkmark$\checkmark^\checkmark\\checkmarkc\checkmarki\checkmarkr\checkmarkc\checkmark$\checkmark=\checkmark$\checkmark^\checkmark\\checkmarkc\checkmarki\checkmarkr\checkmarkc\checkmark$\checkmark=\checkmark$\checkmark^\checkmark\\checkmarkc\checkmarki\checkmarkr\checkmarkc\checkmark$\checkmark=\checkmark$\checkmark^\checkmark\\checkmarkc\checkmarki\checkmarkr\checkmarkc\checkmark$\checkmark=\checkmark$\checkmark^\checkmark\\checkmarkc\checkmarki\checkmarkr\checkmarkc\checkmark$\checkmark=\checkmark$\checkmark^\checkmark\\checkmarkc\checkmarki\checkmarkr\checkmarkc\checkmark$\checkmark=\checkmark$\checkmark^\checkmark\\checkmarkc\checkmarki\checkmarkr\checkmarkc\checkmark$\checkmark=\checkmark$\checkmark^\checkmark\\checkmarkc\checkmarki\checkmarkr\checkmarkc\checkmark$\checkmark=\checkmark$\checkmark^\checkmark\\checkmarkc\checkmarki\checkmarkr\checkmarkc\checkmark$\checkmark=\checkmark$\checkmark^\checkmark\\checkmarkc\checkmarki\checkmarkr\checkmarkc\checkmark$\checkmark=\checkmark$\checkmark^\checkmark\\checkmarkc\checkmarki\checkmarkr\checkmarkc\checkmark$\checkmark=\checkmark$\checkmark^\checkmark\\checkmarkc\checkmarki\checkmarkr\checkmarkc\checkmark$\checkmark=\checkmark$\checkmark^\checkmark\\checkmarkc\checkmarki\checkmarkr\checkmarkc\checkmark$\checkmark=\checkmark$\checkmark^\checkmark\\checkmarkc\checkmarki\checkmarkr\checkmarkc\checkmark$\checkmark=\checkmark$\checkmark^\checkmark\\checkmarkc\checkmarki\checkmarkr\checkmarkc\checkmark$\checkmark=\checkmark$\checkmark^\checkmark\\checkmarkc\checkmarki\checkmarkr\checkmarkc\checkmark$\checkmark=\checkmark$\checkmark^\checkmark\\checkmarkc\checkmarki\checkmarkr\checkmarkc\checkmark$\checkmark=\checkmark$\checkmark^\checkmark\\checkmarkc\checkmarki\checkmarkr\checkmarkc\checkmark$\checkmark=\checkmark$\checkmark^\checkmark\\checkmarkc\checkmarki\checkmarkr\checkmarkc\checkmark$\checkmark=\checkmark$\checkmark^\checkmark\\checkmarkc\checkmarki\checkmarkr\checkmarkc\checkmark$\checkmark=\checkmark$\checkmark^\checkmark\\checkmarkc\checkmarki\checkmarkr\checkmarkc\checkmark$\checkmark=\checkmark$\checkmark^\checkmark\\checkmarkc\checkmarki\checkmarkr\checkmarkc\checkmark$\checkmark=\checkmark$\checkmark^\checkmark\\checkmarkc\checkmarki\checkmarkr\checkmarkc\checkmark$\checkmark=\checkmark$\checkmark^\checkmark\\checkmarkc\checkmarki\checkmarkr\checkmarkc\checkmark$\checkmark=\checkmark$\checkmark^\checkmark\\checkmarkc\checkmarki\checkmarkr\checkmarkc\checkmark$\checkmark=\checkmark$\checkmark^\checkmark\\checkmarkc\checkmarki\checkmarkr\checkmarkc\checkmark$\checkmark=\checkmark$\checkmark^\checkmark\\checkmarkc\checkmarki\checkmarkr\checkmarkc\checkmark$\checkmark=\checkmark$\checkmark^\checkmark\\checkmarkc\checkmarki\checkmarkr\checkmarkc\checkmark$\checkmark=\checkmark$\checkmark^\checkmark\\checkmarkc\checkmarki\checkmarkr\checkmarkc\checkmark$\checkmark=\checkmark$\checkmark^\checkmark\\checkmarkc\checkmarki\checkmarkr\checkmarkc\checkmark$\checkmark=\checkmark$\checkmark^\checkmark\\checkmarkc\checkmarki\checkmarkr\checkmarkc\checkmark$\checkmark=\checkmark$\checkmark^\checkmark\\checkmarkc\checkmarki\checkmarkr\checkmarkc\checkmark$\checkmark=\checkmark$\checkmark^\checkmark\\checkmarkc\checkmarki\checkmarkr\checkmarkc\checkmark$\checkmark=\checkmark$\checkmark^\checkmark\\checkmarkc\checkmarki\checkmarkr\checkmarkc\checkmark$\checkmark=\checkmark$\checkmark^\checkmark\\checkmarkc\checkmarki\checkmarkr\checkmarkc\checkmark$\checkmark=\checkmark$\checkmark^\checkmark\\checkmarkc\checkmarki\checkmarkr\checkmarkc\checkmark$\checkmark=\checkmark$\checkmark^\checkmark\\checkmarkc\checkmarki\checkmarkr\checkmarkc\checkmark$\checkmark=\checkmark$\checkmark^\checkmark\\checkmarkc\checkmarki\checkmarkr\checkmarkc\checkmark$\checkmark=\checkmark$\checkmark^\checkmark\\checkmarkc\checkmarki\checkmarkr\checkmarkc\checkmark$\checkmark=\checkmark$\checkmark^\checkmark\\checkmarkc\checkmarki\checkmarkr\checkmarkc\checkmark$\checkmark=\checkmark$\checkmark^\checkmark\\checkmarkc\checkmarki\checkmarkr\checkmarkc\checkmark$\checkmark=\checkmark$\checkmark^\checkmark\\checkmarkc\checkmarki\checkmarkr\checkmarkc\checkmark$\checkmark=\checkmark$\checkmark^\checkmark\\checkmarkc\checkmarki\checkmarkr\checkmarkc\checkmark$\checkmark=\checkmark$\checkmark^\checkmark\\checkmarkc\checkmarki\checkmarkr\checkmarkc\checkmark$\checkmark=\checkmark$\checkmark^\checkmark\\checkmarkc\checkmarki\checkmarkr\checkmarkc\checkmark$\checkmark=\checkmark$\checkmark^\checkmark\\checkmarkc\checkmarki\checkmarkr\checkmarkc\checkmark$\checkmark=\checkmark$\checkmark^\checkmark\\checkmarkc\checkmarki\checkmarkr\checkmarkc\checkmark$\checkmark=\checkmark$\checkmark^\checkmark\\checkmarkc\checkmarki\checkmarkr\checkmarkc\checkmark$\checkmark=\checkmark$\checkmark^\checkmark\\checkmarkc\checkmarki\checkmarkr\checkmarkc\checkmark$\checkmark=\checkmark$\checkmark^\checkmark\\checkmarkc\checkmarki\checkmarkr\checkmarkc\checkmark$\checkmark=\checkmark$\checkmark^\checkmark\\checkmarkc\checkmarki\checkmarkr\checkmarkc\checkmark$\checkmark=\checkmark$\checkmark^\checkmark\\checkmarkc\checkmarki\checkmarkr\checkmarkc\checkmark$\checkmark=\checkmark$\checkmark^\checkmark\\checkmarkc\checkmarki\checkmarkr\checkmarkc\checkmark$\checkmark=\checkmark$\checkmark^\checkmark\\checkmarkc\checkmarki\checkmarkr\checkmarkc\checkmark$\checkmark=\checkmark$\checkmark^\checkmark\\checkmarkc\checkmarki\checkmarkr\checkmarkc\checkmark$\checkmark=\checkmark$\checkmark^\checkmark\\checkmarkc\checkmarki\checkmarkr\checkmarkc\checkmark$\checkmark=\checkmark$\checkmark^\checkmark\\checkmarkc\checkmarki\checkmarkr\checkmarkc\checkmark$\checkmark=\checkmark$\checkmark^\checkmark\\checkmarkc\checkmarki\checkmarkr\checkmarkc\checkmark$\checkmark=\checkmark$\checkmark^\checkmark\\checkmarkc\checkmarki\checkmarkr\checkmarkc\checkmark$\checkmark
\checkmark$\checkmark^\checkmark\\checkmarkc\checkmarki\checkmarkr\checkmarkc\checkmark$\checkmark\\checkmark$\checkmark^\checkmark\\checkmarkc\checkmarki\checkmarkr\checkmarkc\checkmark$\checkmarks\checkmark$\checkmark^\checkmark\\checkmarkc\checkmarki\checkmarkr\checkmarkc\checkmark$\checkmarke\checkmark$\checkmark^\checkmark\\checkmarkc\checkmarki\checkmarkr\checkmarkc\checkmark$\checkmarkc\checkmark$\checkmark^\checkmark\\checkmarkc\checkmarki\checkmarkr\checkmarkc\checkmark$\checkmarkt\checkmark$\checkmark^\checkmark\\checkmarkc\checkmarki\checkmarkr\checkmarkc\checkmark$\checkmarki\checkmark$\checkmark^\checkmark\\checkmarkc\checkmarki\checkmarkr\checkmarkc\checkmark$\checkmarko\checkmark$\checkmark^\checkmark\\checkmarkc\checkmarki\checkmarkr\checkmarkc\checkmark$\checkmarkn\checkmark$\checkmark^\checkmark\\checkmarkc\checkmarki\checkmarkr\checkmarkc\checkmark$\checkmark{\checkmark$\checkmark^\checkmark\\checkmarkc\checkmarki\checkmarkr\checkmarkc\checkmark$\checkmarkD\checkmark$\checkmark^\checkmark\\checkmarkc\checkmarki\checkmarkr\checkmarkc\checkmark$\checkmarki\checkmark$\checkmark^\checkmark\\checkmarkc\checkmarki\checkmarkr\checkmarkc\checkmark$\checkmarkm\checkmark$\checkmark^\checkmark\\checkmarkc\checkmarki\checkmarkr\checkmarkc\checkmark$\checkmarke\checkmark$\checkmark^\checkmark\\checkmarkc\checkmarki\checkmarkr\checkmarkc\checkmark$\checkmarkn\checkmark$\checkmark^\checkmark\\checkmarkc\checkmarki\checkmarkr\checkmarkc\checkmark$\checkmarks\checkmark$\checkmark^\checkmark\\checkmarkc\checkmarki\checkmarkr\checkmarkc\checkmark$\checkmarki\checkmark$\checkmark^\checkmark\\checkmarkc\checkmarki\checkmarkr\checkmarkc\checkmark$\checkmarko\checkmark$\checkmark^\checkmark\\checkmarkc\checkmarki\checkmarkr\checkmarkc\checkmark$\checkmarkn\checkmark$\checkmark^\checkmark\\checkmarkc\checkmarki\checkmarkr\checkmarkc\checkmark$\checkmarkn\checkmark$\checkmark^\checkmark\\checkmarkc\checkmarki\checkmarkr\checkmarkc\checkmark$\checkmarke\checkmark$\checkmark^\checkmark\\checkmarkc\checkmarki\checkmarkr\checkmarkc\checkmark$\checkmarkm\checkmark$\checkmark^\checkmark\\checkmarkc\checkmarki\checkmarkr\checkmarkc\checkmark$\checkmarke\checkmark$\checkmark^\checkmark\\checkmarkc\checkmarki\checkmarkr\checkmarkc\checkmark$\checkmarkn\checkmark$\checkmark^\checkmark\\checkmarkc\checkmarki\checkmarkr\checkmarkc\checkmark$\checkmarkt\checkmark$\checkmark^\checkmark\\checkmarkc\checkmarki\checkmarkr\checkmarkc\checkmark$\checkmark \checkmark$\checkmark^\checkmark\\checkmarkc\checkmarki\checkmarkr\checkmarkc\checkmark$\checkmarkd\checkmark$\checkmark^\checkmark\\checkmarkc\checkmarki\checkmarkr\checkmarkc\checkmark$\checkmarke\checkmark$\checkmark^\checkmark\\checkmarkc\checkmarki\checkmarkr\checkmarkc\checkmark$\checkmarks\checkmark$\checkmark^\checkmark\\checkmarkc\checkmarki\checkmarkr\checkmarkc\checkmark$\checkmark \checkmark$\checkmark^\checkmark\\checkmarkc\checkmarki\checkmarkr\checkmarkc\checkmark$\checkmarkO\checkmark$\checkmark^\checkmark\\checkmarkc\checkmarki\checkmarkr\checkmarkc\checkmark$\checkmarkr\checkmark$\checkmark^\checkmark\\checkmarkc\checkmarki\checkmarkr\checkmarkc\checkmark$\checkmarkg\checkmark$\checkmark^\checkmark\\checkmarkc\checkmarki\checkmarkr\checkmarkc\checkmark$\checkmarka\checkmark$\checkmark^\checkmark\\checkmarkc\checkmarki\checkmarkr\checkmarkc\checkmark$\checkmarkn\checkmark$\checkmark^\checkmark\\checkmarkc\checkmarki\checkmarkr\checkmarkc\checkmark$\checkmarke\checkmark$\checkmark^\checkmark\\checkmarkc\checkmarki\checkmarkr\checkmarkc\checkmark$\checkmarks\checkmark$\checkmark^\checkmark\\checkmarkc\checkmarki\checkmarkr\checkmarkc\checkmark$\checkmark \checkmark$\checkmark^\checkmark\\checkmarkc\checkmarki\checkmarkr\checkmarkc\checkmark$\checkmarkd\checkmark$\checkmark^\checkmark\\checkmarkc\checkmarki\checkmarkr\checkmarkc\checkmark$\checkmarke\checkmark$\checkmark^\checkmark\\checkmarkc\checkmarki\checkmarkr\checkmarkc\checkmark$\checkmark \checkmark$\checkmark^\checkmark\\checkmarkc\checkmarki\checkmarkr\checkmarkc\checkmark$\checkmarkR\checkmark$\checkmark^\checkmark\\checkmarkc\checkmarki\checkmarkr\checkmarkc\checkmark$\checkmarko\checkmark$\checkmark^\checkmark\\checkmarkc\checkmarki\checkmarkr\checkmarkc\checkmark$\checkmarkt\checkmark$\checkmark^\checkmark\\checkmarkc\checkmarki\checkmarkr\checkmarkc\checkmark$\checkmarka\checkmark$\checkmark^\checkmark\\checkmarkc\checkmarki\checkmarkr\checkmarkc\checkmark$\checkmarkt\checkmark$\checkmark^\checkmark\\checkmarkc\checkmarki\checkmarkr\checkmarkc\checkmark$\checkmarki\checkmark$\checkmark^\checkmark\\checkmarkc\checkmarki\checkmarkr\checkmarkc\checkmark$\checkmarko\checkmark$\checkmark^\checkmark\\checkmarkc\checkmarki\checkmarkr\checkmarkc\checkmark$\checkmarkn\checkmark$\checkmark^\checkmark\\checkmarkc\checkmarki\checkmarkr\checkmarkc\checkmark$\checkmark \checkmark$\checkmark^\checkmark\\checkmarkc\checkmarki\checkmarkr\checkmarkc\checkmark$\checkmark(\checkmark$\checkmark^\checkmark\\checkmarkc\checkmarki\checkmarkr\checkmarkc\checkmark$\checkmarkA\checkmark$\checkmark^\checkmark\\checkmarkc\checkmarki\checkmarkr\checkmarkc\checkmark$\checkmarkx\checkmark$\checkmark^\checkmark\\checkmarkc\checkmarki\checkmarkr\checkmarkc\checkmark$\checkmarke\checkmark$\checkmark^\checkmark\\checkmarkc\checkmarki\checkmarkr\checkmarkc\checkmark$\checkmarks\checkmark$\checkmark^\checkmark\\checkmarkc\checkmarki\checkmarkr\checkmarkc\checkmark$\checkmark \checkmark$\checkmark^\checkmark\\checkmarkc\checkmarki\checkmarkr\checkmarkc\checkmark$\checkmarkA\checkmark$\checkmark^\checkmark\\checkmarkc\checkmarki\checkmarkr\checkmarkc\checkmark$\checkmark \checkmark$\checkmark^\checkmark\\checkmarkc\checkmarki\checkmarkr\checkmarkc\checkmark$\checkmarke\checkmark$\checkmark^\checkmark\\checkmarkc\checkmarki\checkmarkr\checkmarkc\checkmark$\checkmarkt\checkmark$\checkmark^\checkmark\\checkmarkc\checkmarki\checkmarkr\checkmarkc\checkmark$\checkmark \checkmark$\checkmark^\checkmark\\checkmarkc\checkmarki\checkmarkr\checkmarkc\checkmark$\checkmarkB\checkmark$\checkmark^\checkmark\\checkmarkc\checkmarki\checkmarkr\checkmarkc\checkmark$\checkmark)\checkmark$\checkmark^\checkmark\\checkmarkc\checkmarki\checkmarkr\checkmarkc\checkmark$\checkmark}\checkmark$\checkmark^\checkmark\\checkmarkc\checkmarki\checkmarkr\checkmarkc\checkmark$\checkmark
\checkmark$\checkmark^\checkmark\\checkmarkc\checkmarki\checkmarkr\checkmarkc\checkmark$\checkmark
\checkmark$\checkmark^\checkmark\\checkmarkc\checkmarki\checkmarkr\checkmarkc\checkmark$\checkmark\\checkmark$\checkmark^\checkmark\\checkmarkc\checkmarki\checkmarkr\checkmarkc\checkmark$\checkmarks\checkmark$\checkmark^\checkmark\\checkmarkc\checkmarki\checkmarkr\checkmarkc\checkmark$\checkmarku\checkmark$\checkmark^\checkmark\\checkmarkc\checkmarki\checkmarkr\checkmarkc\checkmark$\checkmarkb\checkmark$\checkmark^\checkmark\\checkmarkc\checkmarki\checkmarkr\checkmarkc\checkmark$\checkmarks\checkmark$\checkmark^\checkmark\\checkmarkc\checkmarki\checkmarkr\checkmarkc\checkmark$\checkmarke\checkmark$\checkmark^\checkmark\\checkmarkc\checkmarki\checkmarkr\checkmarkc\checkmark$\checkmarkc\checkmark$\checkmark^\checkmark\\checkmarkc\checkmarki\checkmarkr\checkmarkc\checkmark$\checkmarkt\checkmark$\checkmark^\checkmark\\checkmarkc\checkmarki\checkmarkr\checkmarkc\checkmark$\checkmarki\checkmark$\checkmark^\checkmark\\checkmarkc\checkmarki\checkmarkr\checkmarkc\checkmark$\checkmarko\checkmark$\checkmark^\checkmark\\checkmarkc\checkmarki\checkmarkr\checkmarkc\checkmark$\checkmarkn\checkmark$\checkmark^\checkmark\\checkmarkc\checkmarki\checkmarkr\checkmarkc\checkmark$\checkmark{\checkmark$\checkmark^\checkmark\\checkmarkc\checkmarki\checkmarkr\checkmarkc\checkmark$\checkmarkP\checkmark$\checkmark^\checkmark\\checkmarkc\checkmarki\checkmarkr\checkmarkc\checkmark$\checkmarkr\checkmark$\checkmark^\checkmark\\checkmarkc\checkmarki\checkmarkr\checkmarkc\checkmark$\checkmarko\checkmark$\checkmark^\checkmark\\checkmarkc\checkmarki\checkmarkr\checkmarkc\checkmark$\checkmarkb\checkmark$\checkmark^\checkmark\\checkmarkc\checkmarki\checkmarkr\checkmarkc\checkmark$\checkmarkl\checkmark$\checkmark^\checkmark\\checkmarkc\checkmarki\checkmarkr\checkmarkc\checkmark$\checkmarké\checkmark$\checkmark^\checkmark\\checkmarkc\checkmarki\checkmarkr\checkmarkc\checkmark$\checkmarkm\checkmark$\checkmark^\checkmark\\checkmarkc\checkmarki\checkmarkr\checkmarkc\checkmark$\checkmarka\checkmark$\checkmark^\checkmark\\checkmarkc\checkmarki\checkmarkr\checkmarkc\checkmark$\checkmarkt\checkmark$\checkmark^\checkmark\\checkmarkc\checkmarki\checkmarkr\checkmarkc\checkmark$\checkmarki\checkmark$\checkmark^\checkmark\\checkmarkc\checkmarki\checkmarkr\checkmarkc\checkmark$\checkmarkq\checkmark$\checkmark^\checkmark\\checkmarkc\checkmarki\checkmarkr\checkmarkc\checkmark$\checkmarku\checkmark$\checkmark^\checkmark\\checkmarkc\checkmarki\checkmarkr\checkmarkc\checkmark$\checkmarke\checkmark$\checkmark^\checkmark\\checkmarkc\checkmarki\checkmarkr\checkmarkc\checkmark$\checkmark \checkmark$\checkmark^\checkmark\\checkmarkc\checkmarki\checkmarkr\checkmarkc\checkmark$\checkmarkd\checkmark$\checkmark^\checkmark\\checkmarkc\checkmarki\checkmarkr\checkmarkc\checkmark$\checkmarku\checkmark$\checkmark^\checkmark\\checkmarkc\checkmarki\checkmarkr\checkmarkc\checkmark$\checkmark \checkmark$\checkmark^\checkmark\\checkmarkc\checkmarki\checkmarkr\checkmarkc\checkmark$\checkmarkB\checkmark$\checkmark^\checkmark\\checkmarkc\checkmarki\checkmarkr\checkmarkc\checkmark$\checkmarka\checkmark$\checkmark^\checkmark\\checkmarkc\checkmarki\checkmarkr\checkmarkc\checkmark$\checkmarkc\checkmark$\checkmark^\checkmark\\checkmarkc\checkmarki\checkmarkr\checkmarkc\checkmark$\checkmarkk\checkmark$\checkmark^\checkmark\\checkmarkc\checkmarki\checkmarkr\checkmarkc\checkmark$\checkmarkl\checkmark$\checkmark^\checkmark\\checkmarkc\checkmarki\checkmarkr\checkmarkc\checkmark$\checkmarka\checkmark$\checkmark^\checkmark\\checkmarkc\checkmarki\checkmarkr\checkmarkc\checkmark$\checkmarks\checkmark$\checkmark^\checkmark\\checkmarkc\checkmarki\checkmarkr\checkmarkc\checkmark$\checkmarkh\checkmark$\checkmark^\checkmark\\checkmarkc\checkmarki\checkmarkr\checkmarkc\checkmark$\checkmark \checkmark$\checkmark^\checkmark\\checkmarkc\checkmarki\checkmarkr\checkmarkc\checkmark$\checkmarka\checkmark$\checkmark^\checkmark\\checkmarkc\checkmarki\checkmarkr\checkmarkc\checkmark$\checkmarkn\checkmark$\checkmark^\checkmark\\checkmarkc\checkmarki\checkmarkr\checkmarkc\checkmark$\checkmarkg\checkmark$\checkmark^\checkmark\\checkmarkc\checkmarki\checkmarkr\checkmarkc\checkmark$\checkmarku\checkmark$\checkmark^\checkmark\\checkmarkc\checkmarki\checkmarkr\checkmarkc\checkmark$\checkmarkl\checkmark$\checkmark^\checkmark\\checkmarkc\checkmarki\checkmarkr\checkmarkc\checkmark$\checkmarka\checkmark$\checkmark^\checkmark\\checkmarkc\checkmarki\checkmarkr\checkmarkc\checkmark$\checkmarki\checkmark$\checkmark^\checkmark\\checkmarkc\checkmarki\checkmarkr\checkmarkc\checkmark$\checkmarkr\checkmark$\checkmark^\checkmark\\checkmarkc\checkmarki\checkmarkr\checkmarkc\checkmark$\checkmarke\checkmark$\checkmark^\checkmark\\checkmarkc\checkmarki\checkmarkr\checkmarkc\checkmark$\checkmark}\checkmark$\checkmark^\checkmark\\checkmarkc\checkmarki\checkmarkr\checkmarkc\checkmark$\checkmark
\checkmark$\checkmark^\checkmark\\checkmarkc\checkmarki\checkmarkr\checkmarkc\checkmark$\checkmarkU\checkmark$\checkmark^\checkmark\\checkmarkc\checkmarki\checkmarkr\checkmarkc\checkmark$\checkmarkn\checkmark$\checkmark^\checkmark\\checkmarkc\checkmarki\checkmarkr\checkmarkc\checkmark$\checkmarke\checkmark$\checkmark^\checkmark\\checkmarkc\checkmarki\checkmarkr\checkmarkc\checkmark$\checkmark \checkmark$\checkmark^\checkmark\\checkmarkc\checkmarki\checkmarkr\checkmarkc\checkmark$\checkmarke\checkmark$\checkmark^\checkmark\\checkmarkc\checkmarki\checkmarkr\checkmarkc\checkmark$\checkmarkr\checkmark$\checkmark^\checkmark\\checkmarkc\checkmarki\checkmarkr\checkmarkc\checkmark$\checkmarkr\checkmark$\checkmark^\checkmark\\checkmarkc\checkmarki\checkmarkr\checkmarkc\checkmark$\checkmarke\checkmark$\checkmark^\checkmark\\checkmarkc\checkmarki\checkmarkr\checkmarkc\checkmark$\checkmarku\checkmark$\checkmark^\checkmark\\checkmarkc\checkmarki\checkmarkr\checkmarkc\checkmark$\checkmarkr\checkmark$\checkmark^\checkmark\\checkmarkc\checkmarki\checkmarkr\checkmarkc\checkmark$\checkmark \checkmark$\checkmark^\checkmark\\checkmarkc\checkmarki\checkmarkr\checkmarkc\checkmark$\checkmarka\checkmark$\checkmark^\checkmark\\checkmarkc\checkmarki\checkmarkr\checkmarkc\checkmark$\checkmarkn\checkmark$\checkmark^\checkmark\\checkmarkc\checkmarki\checkmarkr\checkmarkc\checkmark$\checkmarkg\checkmark$\checkmark^\checkmark\\checkmarkc\checkmarki\checkmarkr\checkmarkc\checkmark$\checkmarku\checkmark$\checkmark^\checkmark\\checkmarkc\checkmarki\checkmarkr\checkmarkc\checkmark$\checkmarkl\checkmark$\checkmark^\checkmark\\checkmarkc\checkmarki\checkmarkr\checkmarkc\checkmark$\checkmarka\checkmark$\checkmark^\checkmark\\checkmarkc\checkmarki\checkmarkr\checkmarkc\checkmark$\checkmarki\checkmark$\checkmark^\checkmark\\checkmarkc\checkmarki\checkmarkr\checkmarkc\checkmark$\checkmarkr\checkmark$\checkmark^\checkmark\\checkmarkc\checkmarki\checkmarkr\checkmarkc\checkmark$\checkmarke\checkmark$\checkmark^\checkmark\\checkmarkc\checkmarki\checkmarkr\checkmarkc\checkmark$\checkmark \checkmark$\checkmark^\checkmark\\checkmarkc\checkmarki\checkmarkr\checkmarkc\checkmark$\checkmark$\checkmark$\checkmark^\checkmark\\checkmarkc\checkmarki\checkmarkr\checkmarkc\checkmark$\checkmark\\checkmark$\checkmark^\checkmark\\checkmarkc\checkmarki\checkmarkr\checkmarkc\checkmark$\checkmarkD\checkmark$\checkmark^\checkmark\\checkmarkc\checkmarki\checkmarkr\checkmarkc\checkmark$\checkmarke\checkmark$\checkmark^\checkmark\\checkmarkc\checkmarki\checkmarkr\checkmarkc\checkmark$\checkmarkl\checkmark$\checkmark^\checkmark\\checkmarkc\checkmarki\checkmarkr\checkmarkc\checkmark$\checkmarkt\checkmark$\checkmark^\checkmark\\checkmarkc\checkmarki\checkmarkr\checkmarkc\checkmark$\checkmarka\checkmark$\checkmark^\checkmark\\checkmarkc\checkmarki\checkmarkr\checkmarkc\checkmark$\checkmark\\checkmark$\checkmark^\checkmark\\checkmarkc\checkmarki\checkmarkr\checkmarkc\checkmark$\checkmarkt\checkmark$\checkmark^\checkmark\\checkmarkc\checkmarki\checkmarkr\checkmarkc\checkmark$\checkmarkh\checkmark$\checkmark^\checkmark\\checkmarkc\checkmarki\checkmarkr\checkmarkc\checkmark$\checkmarke\checkmark$\checkmark^\checkmark\\checkmarkc\checkmarki\checkmarkr\checkmarkc\checkmark$\checkmarkt\checkmark$\checkmark^\checkmark\\checkmarkc\checkmarki\checkmarkr\checkmarkc\checkmark$\checkmarka\checkmark$\checkmark^\checkmark\\checkmarkc\checkmarki\checkmarkr\checkmarkc\checkmark$\checkmark$\checkmark$\checkmark^\checkmark\\checkmarkc\checkmarki\checkmarkr\checkmarkc\checkmark$\checkmark \checkmark$\checkmark^\checkmark\\checkmarkc\checkmarki\checkmarkr\checkmarkc\checkmark$\checkmarks\checkmark$\checkmark^\checkmark\\checkmarkc\checkmarki\checkmarkr\checkmarkc\checkmark$\checkmarku\checkmark$\checkmark^\checkmark\\checkmarkc\checkmarki\checkmarkr\checkmarkc\checkmark$\checkmarkr\checkmark$\checkmark^\checkmark\\checkmarkc\checkmarki\checkmarkr\checkmarkc\checkmark$\checkmark \checkmark$\checkmark^\checkmark\\checkmarkc\checkmarki\checkmarkr\checkmarkc\checkmark$\checkmarkl\checkmark$\checkmark^\checkmark\\checkmarkc\checkmarki\checkmarkr\checkmarkc\checkmark$\checkmark'\checkmark$\checkmark^\checkmark\\checkmarkc\checkmarki\checkmarkr\checkmarkc\checkmark$\checkmarkA\checkmark$\checkmark^\checkmark\\checkmarkc\checkmarki\checkmarkr\checkmarkc\checkmark$\checkmarkx\checkmark$\checkmark^\checkmark\\checkmarkc\checkmarki\checkmarkr\checkmarkc\checkmark$\checkmarke\checkmark$\checkmark^\checkmark\\checkmarkc\checkmarki\checkmarkr\checkmarkc\checkmark$\checkmark \checkmark$\checkmark^\checkmark\\checkmarkc\checkmarki\checkmarkr\checkmarkc\checkmark$\checkmarkA\checkmark$\checkmark^\checkmark\\checkmarkc\checkmarki\checkmarkr\checkmarkc\checkmark$\checkmark \checkmark$\checkmark^\checkmark\\checkmarkc\checkmarki\checkmarkr\checkmarkc\checkmark$\checkmarks\checkmark$\checkmark^\checkmark\\checkmarkc\checkmarki\checkmarkr\checkmarkc\checkmark$\checkmarke\checkmark$\checkmark^\checkmark\\checkmarkc\checkmarki\checkmarkr\checkmarkc\checkmark$\checkmark \checkmark$\checkmark^\checkmark\\checkmarkc\checkmarki\checkmarkr\checkmarkc\checkmark$\checkmarkt\checkmark$\checkmark^\checkmark\\checkmarkc\checkmarki\checkmarkr\checkmarkc\checkmark$\checkmarkr\checkmark$\checkmark^\checkmark\\checkmarkc\checkmarki\checkmarkr\checkmarkc\checkmark$\checkmarka\checkmark$\checkmark^\checkmark\\checkmarkc\checkmarki\checkmarkr\checkmarkc\checkmark$\checkmarkd\checkmark$\checkmark^\checkmark\\checkmarkc\checkmarki\checkmarkr\checkmarkc\checkmark$\checkmarku\checkmark$\checkmark^\checkmark\\checkmarkc\checkmarki\checkmarkr\checkmarkc\checkmark$\checkmarki\checkmark$\checkmark^\checkmark\\checkmarkc\checkmarki\checkmarkr\checkmarkc\checkmark$\checkmarkt\checkmark$\checkmark^\checkmark\\checkmarkc\checkmarki\checkmarkr\checkmarkc\checkmark$\checkmark \checkmark$\checkmark^\checkmark\\checkmarkc\checkmarki\checkmarkr\checkmarkc\checkmark$\checkmarke\checkmark$\checkmark^\checkmark\\checkmarkc\checkmarki\checkmarkr\checkmarkc\checkmark$\checkmarkn\checkmark$\checkmark^\checkmark\\checkmarkc\checkmarki\checkmarkr\checkmarkc\checkmark$\checkmark \checkmark$\checkmark^\checkmark\\checkmarkc\checkmarki\checkmarkr\checkmarkc\checkmark$\checkmarku\checkmark$\checkmark^\checkmark\\checkmarkc\checkmarki\checkmarkr\checkmarkc\checkmark$\checkmarkn\checkmark$\checkmark^\checkmark\\checkmarkc\checkmarki\checkmarkr\checkmarkc\checkmark$\checkmarke\checkmark$\checkmark^\checkmark\\checkmarkc\checkmarki\checkmarkr\checkmarkc\checkmark$\checkmark \checkmark$\checkmark^\checkmark\\checkmarkc\checkmarki\checkmarkr\checkmarkc\checkmark$\checkmarke\checkmark$\checkmark^\checkmark\\checkmarkc\checkmarki\checkmarkr\checkmarkc\checkmark$\checkmarkr\checkmark$\checkmark^\checkmark\\checkmarkc\checkmarki\checkmarkr\checkmarkc\checkmark$\checkmarkr\checkmark$\checkmark^\checkmark\\checkmarkc\checkmarki\checkmarkr\checkmarkc\checkmark$\checkmarke\checkmark$\checkmark^\checkmark\\checkmarkc\checkmarki\checkmarkr\checkmarkc\checkmark$\checkmarku\checkmark$\checkmark^\checkmark\\checkmarkc\checkmarki\checkmarkr\checkmarkc\checkmark$\checkmarkr\checkmark$\checkmark^\checkmark\\checkmarkc\checkmarki\checkmarkr\checkmarkc\checkmark$\checkmark \checkmark$\checkmark^\checkmark\\checkmarkc\checkmarki\checkmarkr\checkmarkc\checkmark$\checkmarkl\checkmark$\checkmark^\checkmark\\checkmarkc\checkmarki\checkmarkr\checkmarkc\checkmark$\checkmarki\checkmark$\checkmark^\checkmark\\checkmarkc\checkmarki\checkmarkr\checkmarkc\checkmark$\checkmarkn\checkmark$\checkmark^\checkmark\\checkmarkc\checkmarki\checkmarkr\checkmarkc\checkmark$\checkmarké\checkmark$\checkmark^\checkmark\\checkmarkc\checkmarki\checkmarkr\checkmarkc\checkmark$\checkmarka\checkmark$\checkmark^\checkmark\\checkmarkc\checkmarki\checkmarkr\checkmarkc\checkmark$\checkmarki\checkmark$\checkmark^\checkmark\\checkmarkc\checkmarki\checkmarkr\checkmarkc\checkmark$\checkmarkr\checkmark$\checkmark^\checkmark\\checkmarkc\checkmarki\checkmarkr\checkmarkc\checkmark$\checkmarke\checkmark$\checkmark^\checkmark\\checkmarkc\checkmarki\checkmarkr\checkmarkc\checkmark$\checkmark \checkmark$\checkmark^\checkmark\\checkmarkc\checkmarki\checkmarkr\checkmarkc\checkmark$\checkmarkt\checkmark$\checkmark^\checkmark\\checkmarkc\checkmarki\checkmarkr\checkmarkc\checkmark$\checkmarka\checkmark$\checkmark^\checkmark\\checkmarkc\checkmarki\checkmarkr\checkmarkc\checkmark$\checkmarkn\checkmark$\checkmark^\checkmark\\checkmarkc\checkmarki\checkmarkr\checkmarkc\checkmark$\checkmarkg\checkmark$\checkmark^\checkmark\\checkmarkc\checkmarki\checkmarkr\checkmarkc\checkmark$\checkmarke\checkmark$\checkmark^\checkmark\\checkmarkc\checkmarki\checkmarkr\checkmarkc\checkmark$\checkmarkn\checkmark$\checkmark^\checkmark\\checkmarkc\checkmarki\checkmarkr\checkmarkc\checkmark$\checkmarkt\checkmark$\checkmark^\checkmark\\checkmarkc\checkmarki\checkmarkr\checkmarkc\checkmark$\checkmarki\checkmark$\checkmark^\checkmark\\checkmarkc\checkmarki\checkmarkr\checkmarkc\checkmark$\checkmarke\checkmark$\checkmark^\checkmark\\checkmarkc\checkmarki\checkmarkr\checkmarkc\checkmark$\checkmarkl\checkmark$\checkmark^\checkmark\\checkmarkc\checkmarki\checkmarkr\checkmarkc\checkmark$\checkmarkl\checkmark$\checkmark^\checkmark\\checkmarkc\checkmarki\checkmarkr\checkmarkc\checkmark$\checkmarke\checkmark$\checkmark^\checkmark\\checkmarkc\checkmarki\checkmarkr\checkmarkc\checkmark$\checkmark \checkmark$\checkmark^\checkmark\\checkmarkc\checkmarki\checkmarkr\checkmarkc\checkmark$\checkmark$\checkmark$\checkmark^\checkmark\\checkmarkc\checkmarki\checkmarkr\checkmarkc\checkmark$\checkmark\\checkmark$\checkmark^\checkmark\\checkmarkc\checkmarki\checkmarkr\checkmarkc\checkmark$\checkmarkD\checkmark$\checkmark^\checkmark\\checkmarkc\checkmarki\checkmarkr\checkmarkc\checkmark$\checkmarke\checkmark$\checkmark^\checkmark\\checkmarkc\checkmarki\checkmarkr\checkmarkc\checkmark$\checkmarkl\checkmark$\checkmark^\checkmark\\checkmarkc\checkmarki\checkmarkr\checkmarkc\checkmark$\checkmarkt\checkmark$\checkmark^\checkmark\\checkmarkc\checkmarki\checkmarkr\checkmarkc\checkmark$\checkmarka\checkmark$\checkmark^\checkmark\\checkmarkc\checkmarki\checkmarkr\checkmarkc\checkmark$\checkmark \checkmark$\checkmark^\checkmark\\checkmarkc\checkmarki\checkmarkr\checkmarkc\checkmark$\checkmarkL\checkmark$\checkmark^\checkmark\\checkmarkc\checkmarki\checkmarkr\checkmarkc\checkmark$\checkmark$\checkmark$\checkmark^\checkmark\\checkmarkc\checkmarki\checkmarkr\checkmarkc\checkmark$\checkmark \checkmark$\checkmark^\checkmark\\checkmarkc\checkmarki\checkmarkr\checkmarkc\checkmark$\checkmarkà\checkmark$\checkmark^\checkmark\\checkmarkc\checkmarki\checkmarkr\checkmarkc\checkmark$\checkmark \checkmark$\checkmark^\checkmark\\checkmarkc\checkmarki\checkmarkr\checkmarkc\checkmark$\checkmarkl\checkmark$\checkmark^\checkmark\\checkmarkc\checkmarki\checkmarkr\checkmarkc\checkmark$\checkmark'\checkmark$\checkmark^\checkmark\\checkmarkc\checkmarki\checkmarkr\checkmarkc\checkmark$\checkmarke\checkmark$\checkmark^\checkmark\\checkmarkc\checkmarki\checkmarkr\checkmarkc\checkmark$\checkmarkx\checkmark$\checkmark^\checkmark\\checkmarkc\checkmarki\checkmarkr\checkmarkc\checkmark$\checkmarkt\checkmark$\checkmark^\checkmark\\checkmarkc\checkmarki\checkmarkr\checkmarkc\checkmark$\checkmarkr\checkmark$\checkmark^\checkmark\\checkmarkc\checkmarki\checkmarkr\checkmarkc\checkmark$\checkmarké\checkmark$\checkmark^\checkmark\\checkmarkc\checkmarki\checkmarkr\checkmarkc\checkmark$\checkmarkm\checkmark$\checkmark^\checkmark\\checkmarkc\checkmarki\checkmarkr\checkmarkc\checkmark$\checkmarki\checkmark$\checkmark^\checkmark\\checkmarkc\checkmarki\checkmarkr\checkmarkc\checkmark$\checkmarkt\checkmark$\checkmark^\checkmark\\checkmarkc\checkmarki\checkmarkr\checkmarkc\checkmark$\checkmarké\checkmark$\checkmark^\checkmark\\checkmarkc\checkmarki\checkmarkr\checkmarkc\checkmark$\checkmark \checkmark$\checkmark^\checkmark\\checkmarkc\checkmarki\checkmarkr\checkmarkc\checkmark$\checkmarkd\checkmark$\checkmark^\checkmark\\checkmarkc\checkmarki\checkmarkr\checkmarkc\checkmark$\checkmarku\checkmark$\checkmark^\checkmark\\checkmarkc\checkmarki\checkmarkr\checkmarkc\checkmark$\checkmark \checkmark$\checkmark^\checkmark\\checkmarkc\checkmarki\checkmarkr\checkmarkc\checkmark$\checkmarkp\checkmark$\checkmark^\checkmark\\checkmarkc\checkmarki\checkmarkr\checkmarkc\checkmark$\checkmarkl\checkmark$\checkmark^\checkmark\\checkmarkc\checkmarki\checkmarkr\checkmarkc\checkmark$\checkmarka\checkmark$\checkmark^\checkmark\\checkmarkc\checkmarki\checkmarkr\checkmarkc\checkmark$\checkmarkt\checkmark$\checkmark^\checkmark\\checkmarkc\checkmarki\checkmarkr\checkmarkc\checkmark$\checkmarke\checkmark$\checkmark^\checkmark\\checkmarkc\checkmarki\checkmarkr\checkmarkc\checkmark$\checkmarka\checkmark$\checkmark^\checkmark\\checkmarkc\checkmarki\checkmarkr\checkmarkc\checkmark$\checkmarku\checkmark$\checkmark^\checkmark\\checkmarkc\checkmarki\checkmarkr\checkmarkc\checkmark$\checkmark \checkmark$\checkmark^\checkmark\\checkmarkc\checkmarki\checkmarkr\checkmarkc\checkmark$\checkmark(\checkmark$\checkmark^\checkmark\\checkmarkc\checkmarki\checkmarkr\checkmarkc\checkmark$\checkmarkr\checkmark$\checkmark^\checkmark\\checkmarkc\checkmarki\checkmarkr\checkmarkc\checkmark$\checkmarka\checkmark$\checkmark^\checkmark\\checkmarkc\checkmarki\checkmarkr\checkmarkc\checkmark$\checkmarky\checkmark$\checkmark^\checkmark\\checkmarkc\checkmarki\checkmarkr\checkmarkc\checkmark$\checkmarko\checkmark$\checkmark^\checkmark\\checkmarkc\checkmarki\checkmarkr\checkmarkc\checkmark$\checkmarkn\checkmark$\checkmark^\checkmark\\checkmarkc\checkmarki\checkmarkr\checkmarkc\checkmark$\checkmark \checkmark$\checkmark^\checkmark\\checkmarkc\checkmarki\checkmarkr\checkmarkc\checkmark$\checkmark$\checkmark$\checkmark^\checkmark\\checkmarkc\checkmarki\checkmarkr\checkmarkc\checkmark$\checkmarkR\checkmark$\checkmark^\checkmark\\checkmarkc\checkmarki\checkmarkr\checkmarkc\checkmark$\checkmark_\checkmark$\checkmark^\checkmark\\checkmarkc\checkmarki\checkmarkr\checkmarkc\checkmark$\checkmark{\checkmark$\checkmark^\checkmark\\checkmarkc\checkmarki\checkmarkr\checkmarkc\checkmark$\checkmarkp\checkmark$\checkmark^\checkmark\\checkmarkc\checkmarki\checkmarkr\checkmarkc\checkmark$\checkmarkl\checkmark$\checkmark^\checkmark\\checkmarkc\checkmarki\checkmarkr\checkmarkc\checkmark$\checkmarka\checkmark$\checkmark^\checkmark\\checkmarkc\checkmarki\checkmarkr\checkmarkc\checkmark$\checkmarkt\checkmark$\checkmark^\checkmark\\checkmarkc\checkmarki\checkmarkr\checkmarkc\checkmark$\checkmarke\checkmark$\checkmark^\checkmark\\checkmarkc\checkmarki\checkmarkr\checkmarkc\checkmark$\checkmarka\checkmark$\checkmark^\checkmark\\checkmarkc\checkmarki\checkmarkr\checkmarkc\checkmark$\checkmarku\checkmark$\checkmark^\checkmark\\checkmarkc\checkmarki\checkmarkr\checkmarkc\checkmark$\checkmark}\checkmark$\checkmark^\checkmark\\checkmarkc\checkmarki\checkmarkr\checkmarkc\checkmark$\checkmark \checkmark$\checkmark^\checkmark\\checkmarkc\checkmarki\checkmarkr\checkmarkc\checkmark$\checkmark=\checkmark$\checkmark^\checkmark\\checkmarkc\checkmarki\checkmarkr\checkmarkc\checkmark$\checkmark \checkmark$\checkmark^\checkmark\\checkmarkc\checkmarki\checkmarkr\checkmarkc\checkmark$\checkmark1\checkmark$\checkmark^\checkmark\\checkmarkc\checkmarki\checkmarkr\checkmarkc\checkmark$\checkmark0\checkmark$\checkmark^\checkmark\\checkmarkc\checkmarki\checkmarkr\checkmarkc\checkmark$\checkmark0\checkmark$\checkmark^\checkmark\\checkmarkc\checkmarki\checkmarkr\checkmarkc\checkmark$\checkmark$\checkmark$\checkmark^\checkmark\\checkmarkc\checkmarki\checkmarkr\checkmarkc\checkmark$\checkmark \checkmark$\checkmark^\checkmark\\checkmarkc\checkmarki\checkmarkr\checkmarkc\checkmark$\checkmarkm\checkmark$\checkmark^\checkmark\\checkmarkc\checkmarki\checkmarkr\checkmarkc\checkmark$\checkmarkm\checkmark$\checkmark^\checkmark\\checkmarkc\checkmarki\checkmarkr\checkmarkc\checkmark$\checkmark)\checkmark$\checkmark^\checkmark\\checkmarkc\checkmarki\checkmarkr\checkmarkc\checkmark$\checkmark \checkmark$\checkmark^\checkmark\\checkmarkc\checkmarki\checkmarkr\checkmarkc\checkmark$\checkmark:\checkmark$\checkmark^\checkmark\\checkmarkc\checkmarki\checkmarkr\checkmarkc\checkmark$\checkmark
\checkmark$\checkmark^\checkmark\\checkmarkc\checkmarki\checkmarkr\checkmarkc\checkmark$\checkmark\\checkmark$\checkmark^\checkmark\\checkmarkc\checkmarki\checkmarkr\checkmarkc\checkmark$\checkmarkb\checkmark$\checkmark^\checkmark\\checkmarkc\checkmarki\checkmarkr\checkmarkc\checkmark$\checkmarke\checkmark$\checkmark^\checkmark\\checkmarkc\checkmarki\checkmarkr\checkmarkc\checkmark$\checkmarkg\checkmark$\checkmark^\checkmark\\checkmarkc\checkmarki\checkmarkr\checkmarkc\checkmark$\checkmarki\checkmark$\checkmark^\checkmark\\checkmarkc\checkmarki\checkmarkr\checkmarkc\checkmark$\checkmarkn\checkmark$\checkmark^\checkmark\\checkmarkc\checkmarki\checkmarkr\checkmarkc\checkmark$\checkmark{\checkmark$\checkmark^\checkmark\\checkmarkc\checkmarki\checkmarkr\checkmarkc\checkmark$\checkmarke\checkmark$\checkmark^\checkmark\\checkmarkc\checkmarki\checkmarkr\checkmarkc\checkmark$\checkmarkq\checkmark$\checkmark^\checkmark\\checkmarkc\checkmarki\checkmarkr\checkmarkc\checkmark$\checkmarku\checkmark$\checkmark^\checkmark\\checkmarkc\checkmarki\checkmarkr\checkmarkc\checkmark$\checkmarka\checkmark$\checkmark^\checkmark\\checkmarkc\checkmarki\checkmarkr\checkmarkc\checkmark$\checkmarkt\checkmark$\checkmark^\checkmark\\checkmarkc\checkmarki\checkmarkr\checkmarkc\checkmark$\checkmarki\checkmark$\checkmark^\checkmark\\checkmarkc\checkmarki\checkmarkr\checkmarkc\checkmark$\checkmarko\checkmark$\checkmark^\checkmark\\checkmarkc\checkmarki\checkmarkr\checkmarkc\checkmark$\checkmarkn\checkmark$\checkmark^\checkmark\\checkmarkc\checkmarki\checkmarkr\checkmarkc\checkmark$\checkmark}\checkmark$\checkmark^\checkmark\\checkmarkc\checkmarki\checkmarkr\checkmarkc\checkmark$\checkmark
\checkmark$\checkmark^\checkmark\\checkmarkc\checkmarki\checkmarkr\checkmarkc\checkmark$\checkmark \checkmark$\checkmark^\checkmark\\checkmarkc\checkmarki\checkmarkr\checkmarkc\checkmark$\checkmark \checkmark$\checkmark^\checkmark\\checkmarkc\checkmarki\checkmarkr\checkmarkc\checkmark$\checkmark \checkmark$\checkmark^\checkmark\\checkmarkc\checkmarki\checkmarkr\checkmarkc\checkmark$\checkmark \checkmark$\checkmark^\checkmark\\checkmarkc\checkmarki\checkmarkr\checkmarkc\checkmark$\checkmark\\checkmark$\checkmark^\checkmark\\checkmarkc\checkmarki\checkmarkr\checkmarkc\checkmark$\checkmarkD\checkmark$\checkmark^\checkmark\\checkmarkc\checkmarki\checkmarkr\checkmarkc\checkmark$\checkmarke\checkmark$\checkmark^\checkmark\\checkmarkc\checkmarki\checkmarkr\checkmarkc\checkmark$\checkmarkl\checkmark$\checkmark^\checkmark\\checkmarkc\checkmarki\checkmarkr\checkmarkc\checkmark$\checkmarkt\checkmark$\checkmark^\checkmark\\checkmarkc\checkmarki\checkmarkr\checkmarkc\checkmark$\checkmarka\checkmark$\checkmark^\checkmark\\checkmarkc\checkmarki\checkmarkr\checkmarkc\checkmark$\checkmark \checkmark$\checkmark^\checkmark\\checkmarkc\checkmarki\checkmarkr\checkmarkc\checkmark$\checkmarkL\checkmark$\checkmark^\checkmark\\checkmarkc\checkmarki\checkmarkr\checkmarkc\checkmark$\checkmark \checkmark$\checkmark^\checkmark\\checkmarkc\checkmarki\checkmarkr\checkmarkc\checkmark$\checkmark=\checkmark$\checkmark^\checkmark\\checkmarkc\checkmarki\checkmarkr\checkmarkc\checkmark$\checkmark \checkmark$\checkmark^\checkmark\\checkmarkc\checkmarki\checkmarkr\checkmarkc\checkmark$\checkmarkR\checkmark$\checkmark^\checkmark\\checkmarkc\checkmarki\checkmarkr\checkmarkc\checkmark$\checkmark_\checkmark$\checkmark^\checkmark\\checkmarkc\checkmarki\checkmarkr\checkmarkc\checkmark$\checkmark{\checkmark$\checkmark^\checkmark\\checkmarkc\checkmarki\checkmarkr\checkmarkc\checkmark$\checkmarkp\checkmark$\checkmark^\checkmark\\checkmarkc\checkmarki\checkmarkr\checkmarkc\checkmark$\checkmarkl\checkmark$\checkmark^\checkmark\\checkmarkc\checkmarki\checkmarkr\checkmarkc\checkmark$\checkmarka\checkmark$\checkmark^\checkmark\\checkmarkc\checkmarki\checkmarkr\checkmarkc\checkmark$\checkmarkt\checkmark$\checkmark^\checkmark\\checkmarkc\checkmarki\checkmarkr\checkmarkc\checkmark$\checkmarke\checkmark$\checkmark^\checkmark\\checkmarkc\checkmarki\checkmarkr\checkmarkc\checkmark$\checkmarka\checkmark$\checkmark^\checkmark\\checkmarkc\checkmarki\checkmarkr\checkmarkc\checkmark$\checkmarku\checkmark$\checkmark^\checkmark\\checkmarkc\checkmarki\checkmarkr\checkmarkc\checkmark$\checkmark}\checkmark$\checkmark^\checkmark\\checkmarkc\checkmarki\checkmarkr\checkmarkc\checkmark$\checkmark \checkmark$\checkmark^\checkmark\\checkmarkc\checkmarki\checkmarkr\checkmarkc\checkmark$\checkmark\\checkmark$\checkmark^\checkmark\\checkmarkc\checkmarki\checkmarkr\checkmarkc\checkmark$\checkmarkc\checkmark$\checkmark^\checkmark\\checkmarkc\checkmarki\checkmarkr\checkmarkc\checkmark$\checkmarkd\checkmark$\checkmark^\checkmark\\checkmarkc\checkmarki\checkmarkr\checkmarkc\checkmark$\checkmarko\checkmark$\checkmark^\checkmark\\checkmarkc\checkmarki\checkmarkr\checkmarkc\checkmark$\checkmarkt\checkmark$\checkmark^\checkmark\\checkmarkc\checkmarki\checkmarkr\checkmarkc\checkmark$\checkmark \checkmark$\checkmark^\checkmark\\checkmarkc\checkmarki\checkmarkr\checkmarkc\checkmark$\checkmark\\checkmark$\checkmark^\checkmark\\checkmarkc\checkmarki\checkmarkr\checkmarkc\checkmark$\checkmarkt\checkmark$\checkmark^\checkmark\\checkmarkc\checkmarki\checkmarkr\checkmarkc\checkmark$\checkmarka\checkmark$\checkmark^\checkmark\\checkmarkc\checkmarki\checkmarkr\checkmarkc\checkmark$\checkmarkn\checkmark$\checkmark^\checkmark\\checkmarkc\checkmarki\checkmarkr\checkmarkc\checkmark$\checkmark(\checkmark$\checkmark^\checkmark\\checkmarkc\checkmarki\checkmarkr\checkmarkc\checkmark$\checkmark\\checkmark$\checkmark^\checkmark\\checkmarkc\checkmarki\checkmarkr\checkmarkc\checkmark$\checkmarkD\checkmark$\checkmark^\checkmark\\checkmarkc\checkmarki\checkmarkr\checkmarkc\checkmark$\checkmarke\checkmark$\checkmark^\checkmark\\checkmarkc\checkmarki\checkmarkr\checkmarkc\checkmark$\checkmarkl\checkmark$\checkmark^\checkmark\\checkmarkc\checkmarki\checkmarkr\checkmarkc\checkmark$\checkmarkt\checkmark$\checkmark^\checkmark\\checkmarkc\checkmarki\checkmarkr\checkmarkc\checkmark$\checkmarka\checkmark$\checkmark^\checkmark\\checkmarkc\checkmarki\checkmarkr\checkmarkc\checkmark$\checkmark\\checkmark$\checkmark^\checkmark\\checkmarkc\checkmarki\checkmarkr\checkmarkc\checkmark$\checkmarkt\checkmark$\checkmark^\checkmark\\checkmarkc\checkmarki\checkmarkr\checkmarkc\checkmark$\checkmarkh\checkmark$\checkmark^\checkmark\\checkmarkc\checkmarki\checkmarkr\checkmarkc\checkmark$\checkmarke\checkmark$\checkmark^\checkmark\\checkmarkc\checkmarki\checkmarkr\checkmarkc\checkmark$\checkmarkt\checkmark$\checkmark^\checkmark\\checkmarkc\checkmarki\checkmarkr\checkmarkc\checkmark$\checkmarka\checkmark$\checkmark^\checkmark\\checkmarkc\checkmarki\checkmarkr\checkmarkc\checkmark$\checkmark)\checkmark$\checkmark^\checkmark\\checkmarkc\checkmarki\checkmarkr\checkmarkc\checkmark$\checkmark
\checkmark$\checkmark^\checkmark\\checkmarkc\checkmarki\checkmarkr\checkmarkc\checkmark$\checkmark\\checkmark$\checkmark^\checkmark\\checkmarkc\checkmarki\checkmarkr\checkmarkc\checkmark$\checkmarke\checkmark$\checkmark^\checkmark\\checkmarkc\checkmarki\checkmarkr\checkmarkc\checkmark$\checkmarkn\checkmark$\checkmark^\checkmark\\checkmarkc\checkmarki\checkmarkr\checkmarkc\checkmark$\checkmarkd\checkmark$\checkmark^\checkmark\\checkmarkc\checkmarki\checkmarkr\checkmarkc\checkmark$\checkmark{\checkmark$\checkmark^\checkmark\\checkmarkc\checkmarki\checkmarkr\checkmarkc\checkmark$\checkmarke\checkmark$\checkmark^\checkmark\\checkmarkc\checkmarki\checkmarkr\checkmarkc\checkmark$\checkmarkq\checkmark$\checkmark^\checkmark\\checkmarkc\checkmarki\checkmarkr\checkmarkc\checkmark$\checkmarku\checkmark$\checkmark^\checkmark\\checkmarkc\checkmarki\checkmarkr\checkmarkc\checkmark$\checkmarka\checkmark$\checkmark^\checkmark\\checkmarkc\checkmarki\checkmarkr\checkmarkc\checkmark$\checkmarkt\checkmark$\checkmark^\checkmark\\checkmarkc\checkmarki\checkmarkr\checkmarkc\checkmark$\checkmarki\checkmark$\checkmark^\checkmark\\checkmarkc\checkmarki\checkmarkr\checkmarkc\checkmark$\checkmarko\checkmark$\checkmark^\checkmark\\checkmarkc\checkmarki\checkmarkr\checkmarkc\checkmark$\checkmarkn\checkmark$\checkmark^\checkmark\\checkmarkc\checkmarki\checkmarkr\checkmarkc\checkmark$\checkmark}\checkmark$\checkmark^\checkmark\\checkmarkc\checkmarki\checkmarkr\checkmarkc\checkmark$\checkmark
\checkmark$\checkmark^\checkmark\\checkmarkc\checkmarki\checkmarkr\checkmarkc\checkmark$\checkmark
\checkmark$\checkmark^\checkmark\\checkmarkc\checkmarki\checkmarkr\checkmarkc\checkmark$\checkmarkP\checkmark$\checkmark^\checkmark\\checkmarkc\checkmarki\checkmarkr\checkmarkc\checkmark$\checkmarko\checkmark$\checkmark^\checkmark\\checkmarkc\checkmarki\checkmarkr\checkmarkc\checkmark$\checkmarku\checkmark$\checkmark^\checkmark\\checkmarkc\checkmarki\checkmarkr\checkmarkc\checkmark$\checkmarkr\checkmark$\checkmark^\checkmark\\checkmarkc\checkmarki\checkmarkr\checkmarkc\checkmark$\checkmark \checkmark$\checkmark^\checkmark\\checkmarkc\checkmarki\checkmarkr\checkmarkc\checkmark$\checkmarku\checkmark$\checkmark^\checkmark\\checkmarkc\checkmarki\checkmarkr\checkmarkc\checkmark$\checkmarkn\checkmark$\checkmark^\checkmark\\checkmarkc\checkmarki\checkmarkr\checkmarkc\checkmark$\checkmark \checkmark$\checkmark^\checkmark\\checkmarkc\checkmarki\checkmarkr\checkmarkc\checkmark$\checkmarkj\checkmark$\checkmark^\checkmark\\checkmarkc\checkmarki\checkmarkr\checkmarkc\checkmark$\checkmarke\checkmark$\checkmark^\checkmark\\checkmarkc\checkmarki\checkmarkr\checkmarkc\checkmark$\checkmarku\checkmark$\checkmark^\checkmark\\checkmarkc\checkmarki\checkmarkr\checkmarkc\checkmark$\checkmark \checkmark$\checkmark^\checkmark\\checkmarkc\checkmarki\checkmarkr\checkmarkc\checkmark$\checkmarkd\checkmark$\checkmark^\checkmark\\checkmarkc\checkmarki\checkmarkr\checkmarkc\checkmark$\checkmark'\checkmark$\checkmark^\checkmark\\checkmarkc\checkmarki\checkmarkr\checkmarkc\checkmark$\checkmarki\checkmark$\checkmark^\checkmark\\checkmarkc\checkmarki\checkmarkr\checkmarkc\checkmark$\checkmarkn\checkmark$\checkmark^\checkmark\\checkmarkc\checkmarki\checkmarkr\checkmarkc\checkmark$\checkmarkv\checkmark$\checkmark^\checkmark\\checkmarkc\checkmarki\checkmarkr\checkmarkc\checkmark$\checkmarke\checkmark$\checkmark^\checkmark\\checkmarkc\checkmarki\checkmarkr\checkmarkc\checkmark$\checkmarkr\checkmark$\checkmark^\checkmark\\checkmarkc\checkmarki\checkmarkr\checkmarkc\checkmark$\checkmarks\checkmark$\checkmark^\checkmark\\checkmarkc\checkmarki\checkmarkr\checkmarkc\checkmark$\checkmarki\checkmark$\checkmark^\checkmark\\checkmarkc\checkmarki\checkmarkr\checkmarkc\checkmark$\checkmarko\checkmark$\checkmark^\checkmark\\checkmarkc\checkmarki\checkmarkr\checkmarkc\checkmark$\checkmarkn\checkmark$\checkmark^\checkmark\\checkmarkc\checkmarki\checkmarkr\checkmarkc\checkmark$\checkmark \checkmark$\checkmark^\checkmark\\checkmarkc\checkmarki\checkmarkr\checkmarkc\checkmark$\checkmarkt\checkmark$\checkmark^\checkmark\\checkmarkc\checkmarki\checkmarkr\checkmarkc\checkmark$\checkmarky\checkmark$\checkmark^\checkmark\\checkmarkc\checkmarki\checkmarkr\checkmarkc\checkmark$\checkmarkp\checkmark$\checkmark^\checkmark\\checkmarkc\checkmarki\checkmarkr\checkmarkc\checkmark$\checkmarki\checkmark$\checkmark^\checkmark\\checkmarkc\checkmarki\checkmarkr\checkmarkc\checkmark$\checkmarkq\checkmark$\checkmark^\checkmark\\checkmarkc\checkmarki\checkmarkr\checkmarkc\checkmark$\checkmarku\checkmark$\checkmark^\checkmark\\checkmarkc\checkmarki\checkmarkr\checkmarkc\checkmark$\checkmarke\checkmark$\checkmark^\checkmark\\checkmarkc\checkmarki\checkmarkr\checkmarkc\checkmark$\checkmark \checkmark$\checkmark^\checkmark\\checkmarkc\checkmarki\checkmarkr\checkmarkc\checkmark$\checkmarkd\checkmark$\checkmark^\checkmark\\checkmarkc\checkmarki\checkmarkr\checkmarkc\checkmark$\checkmark'\checkmark$\checkmark^\checkmark\\checkmarkc\checkmarki\checkmarkr\checkmarkc\checkmark$\checkmarku\checkmark$\checkmark^\checkmark\\checkmarkc\checkmarki\checkmarkr\checkmarkc\checkmark$\checkmarkn\checkmark$\checkmark^\checkmark\\checkmarkc\checkmarki\checkmarkr\checkmarkc\checkmark$\checkmark \checkmark$\checkmark^\checkmark\\checkmarkc\checkmarki\checkmarkr\checkmarkc\checkmark$\checkmarka\checkmark$\checkmark^\checkmark\\checkmarkc\checkmarki\checkmarkr\checkmarkc\checkmark$\checkmarkc\checkmark$\checkmark^\checkmark\\checkmarkc\checkmarki\checkmarkr\checkmarkc\checkmark$\checkmarkc\checkmark$\checkmark^\checkmark\\checkmarkc\checkmarki\checkmarkr\checkmarkc\checkmark$\checkmarko\checkmark$\checkmark^\checkmark\\checkmarkc\checkmarki\checkmarkr\checkmarkc\checkmark$\checkmarku\checkmark$\checkmark^\checkmark\\checkmarkc\checkmarki\checkmarkr\checkmarkc\checkmark$\checkmarkp\checkmark$\checkmark^\checkmark\\checkmarkc\checkmarki\checkmarkr\checkmarkc\checkmark$\checkmarkl\checkmark$\checkmark^\checkmark\\checkmarkc\checkmarki\checkmarkr\checkmarkc\checkmark$\checkmarke\checkmark$\checkmark^\checkmark\\checkmarkc\checkmarki\checkmarkr\checkmarkc\checkmark$\checkmarkm\checkmark$\checkmark^\checkmark\\checkmarkc\checkmarki\checkmarkr\checkmarkc\checkmark$\checkmarke\checkmark$\checkmark^\checkmark\\checkmarkc\checkmarki\checkmarkr\checkmarkc\checkmark$\checkmarkn\checkmark$\checkmark^\checkmark\\checkmarkc\checkmarki\checkmarkr\checkmarkc\checkmark$\checkmarkt\checkmark$\checkmark^\checkmark\\checkmarkc\checkmarki\checkmarkr\checkmarkc\checkmark$\checkmark \checkmark$\checkmark^\checkmark\\checkmarkc\checkmarki\checkmarkr\checkmarkc\checkmark$\checkmarkd\checkmark$\checkmark^\checkmark\\checkmarkc\checkmarki\checkmarkr\checkmarkc\checkmark$\checkmarki\checkmark$\checkmark^\checkmark\\checkmarkc\checkmarki\checkmarkr\checkmarkc\checkmark$\checkmarkr\checkmark$\checkmark^\checkmark\\checkmarkc\checkmarki\checkmarkr\checkmarkc\checkmark$\checkmarke\checkmark$\checkmark^\checkmark\\checkmarkc\checkmarki\checkmarkr\checkmarkc\checkmark$\checkmarkc\checkmark$\checkmark^\checkmark\\checkmarkc\checkmarki\checkmarkr\checkmarkc\checkmark$\checkmarkt\checkmark$\checkmark^\checkmark\\checkmarkc\checkmarki\checkmarkr\checkmarkc\checkmark$\checkmark \checkmark$\checkmark^\checkmark\\checkmarkc\checkmarki\checkmarkr\checkmarkc\checkmark$\checkmarkd\checkmark$\checkmark^\checkmark\\checkmarkc\checkmarki\checkmarkr\checkmarkc\checkmark$\checkmarke\checkmark$\checkmark^\checkmark\\checkmarkc\checkmarki\checkmarkr\checkmarkc\checkmark$\checkmark \checkmark$\checkmark^\checkmark\\checkmarkc\checkmarki\checkmarkr\checkmarkc\checkmark$\checkmark$\checkmark$\checkmark^\checkmark\\checkmarkc\checkmarki\checkmarkr\checkmarkc\checkmark$\checkmark\\checkmark$\checkmark^\checkmark\\checkmarkc\checkmarki\checkmarkr\checkmarkc\checkmark$\checkmarkD\checkmark$\checkmark^\checkmark\\checkmarkc\checkmarki\checkmarkr\checkmarkc\checkmark$\checkmarke\checkmark$\checkmark^\checkmark\\checkmarkc\checkmarki\checkmarkr\checkmarkc\checkmark$\checkmarkl\checkmark$\checkmark^\checkmark\\checkmarkc\checkmarki\checkmarkr\checkmarkc\checkmark$\checkmarkt\checkmark$\checkmark^\checkmark\\checkmarkc\checkmarki\checkmarkr\checkmarkc\checkmark$\checkmarka\checkmark$\checkmark^\checkmark\\checkmarkc\checkmarki\checkmarkr\checkmarkc\checkmark$\checkmark\\checkmark$\checkmark^\checkmark\\checkmarkc\checkmarki\checkmarkr\checkmarkc\checkmark$\checkmarkt\checkmark$\checkmark^\checkmark\\checkmarkc\checkmarki\checkmarkr\checkmarkc\checkmark$\checkmarkh\checkmark$\checkmark^\checkmark\\checkmarkc\checkmarki\checkmarkr\checkmarkc\checkmark$\checkmarke\checkmark$\checkmark^\checkmark\\checkmarkc\checkmarki\checkmarkr\checkmarkc\checkmark$\checkmarkt\checkmark$\checkmark^\checkmark\\checkmarkc\checkmarki\checkmarkr\checkmarkc\checkmark$\checkmarka\checkmark$\checkmark^\checkmark\\checkmarkc\checkmarki\checkmarkr\checkmarkc\checkmark$\checkmark \checkmark$\checkmark^\checkmark\\checkmarkc\checkmarki\checkmarkr\checkmarkc\checkmark$\checkmark=\checkmark$\checkmark^\checkmark\\checkmarkc\checkmarki\checkmarkr\checkmarkc\checkmark$\checkmark \checkmark$\checkmark^\checkmark\\checkmarkc\checkmarki\checkmarkr\checkmarkc\checkmark$\checkmark1\checkmark$\checkmark^\checkmark\\checkmarkc\checkmarki\checkmarkr\checkmarkc\checkmark$\checkmark°\checkmark$\checkmark^\checkmark\\checkmarkc\checkmarki\checkmarkr\checkmarkc\checkmark$\checkmark$\checkmark$\checkmark^\checkmark\\checkmarkc\checkmarki\checkmarkr\checkmarkc\checkmark$\checkmark \checkmark$\checkmark^\checkmark\\checkmarkc\checkmarki\checkmarkr\checkmarkc\checkmark$\checkmark:\checkmark$\checkmark^\checkmark\\checkmarkc\checkmarki\checkmarkr\checkmarkc\checkmark$\checkmark
\checkmark$\checkmark^\checkmark\\checkmarkc\checkmarki\checkmarkr\checkmarkc\checkmark$\checkmark\\checkmark$\checkmark^\checkmark\\checkmarkc\checkmarki\checkmarkr\checkmarkc\checkmark$\checkmarkb\checkmark$\checkmark^\checkmark\\checkmarkc\checkmarki\checkmarkr\checkmarkc\checkmark$\checkmarke\checkmark$\checkmark^\checkmark\\checkmarkc\checkmarki\checkmarkr\checkmarkc\checkmark$\checkmarkg\checkmark$\checkmark^\checkmark\\checkmarkc\checkmarki\checkmarkr\checkmarkc\checkmark$\checkmarki\checkmark$\checkmark^\checkmark\\checkmarkc\checkmarki\checkmarkr\checkmarkc\checkmark$\checkmarkn\checkmark$\checkmark^\checkmark\\checkmarkc\checkmarki\checkmarkr\checkmarkc\checkmark$\checkmark{\checkmark$\checkmark^\checkmark\\checkmarkc\checkmarki\checkmarkr\checkmarkc\checkmark$\checkmarke\checkmark$\checkmark^\checkmark\\checkmarkc\checkmarki\checkmarkr\checkmarkc\checkmark$\checkmarkq\checkmark$\checkmark^\checkmark\\checkmarkc\checkmarki\checkmarkr\checkmarkc\checkmark$\checkmarku\checkmark$\checkmark^\checkmark\\checkmarkc\checkmarki\checkmarkr\checkmarkc\checkmark$\checkmarka\checkmark$\checkmark^\checkmark\\checkmarkc\checkmarki\checkmarkr\checkmarkc\checkmark$\checkmarkt\checkmark$\checkmark^\checkmark\\checkmarkc\checkmarki\checkmarkr\checkmarkc\checkmark$\checkmarki\checkmark$\checkmark^\checkmark\\checkmarkc\checkmarki\checkmarkr\checkmarkc\checkmark$\checkmarko\checkmark$\checkmark^\checkmark\\checkmarkc\checkmarki\checkmarkr\checkmarkc\checkmark$\checkmarkn\checkmark$\checkmark^\checkmark\\checkmarkc\checkmarki\checkmarkr\checkmarkc\checkmark$\checkmark}\checkmark$\checkmark^\checkmark\\checkmarkc\checkmarki\checkmarkr\checkmarkc\checkmark$\checkmark
\checkmark$\checkmark^\checkmark\\checkmarkc\checkmarki\checkmarkr\checkmarkc\checkmark$\checkmark \checkmark$\checkmark^\checkmark\\checkmarkc\checkmarki\checkmarkr\checkmarkc\checkmark$\checkmark \checkmark$\checkmark^\checkmark\\checkmarkc\checkmarki\checkmarkr\checkmarkc\checkmark$\checkmark \checkmark$\checkmark^\checkmark\\checkmarkc\checkmarki\checkmarkr\checkmarkc\checkmark$\checkmark \checkmark$\checkmark^\checkmark\\checkmarkc\checkmarki\checkmarkr\checkmarkc\checkmark$\checkmark\\checkmark$\checkmark^\checkmark\\checkmarkc\checkmarki\checkmarkr\checkmarkc\checkmark$\checkmarkD\checkmark$\checkmark^\checkmark\\checkmarkc\checkmarki\checkmarkr\checkmarkc\checkmark$\checkmarke\checkmark$\checkmark^\checkmark\\checkmarkc\checkmarki\checkmarkr\checkmarkc\checkmark$\checkmarkl\checkmark$\checkmark^\checkmark\\checkmarkc\checkmarki\checkmarkr\checkmarkc\checkmark$\checkmarkt\checkmark$\checkmark^\checkmark\\checkmarkc\checkmarki\checkmarkr\checkmarkc\checkmark$\checkmarka\checkmark$\checkmark^\checkmark\\checkmarkc\checkmarki\checkmarkr\checkmarkc\checkmark$\checkmark \checkmark$\checkmark^\checkmark\\checkmarkc\checkmarki\checkmarkr\checkmarkc\checkmark$\checkmarkL\checkmark$\checkmark^\checkmark\\checkmarkc\checkmarki\checkmarkr\checkmarkc\checkmark$\checkmark \checkmark$\checkmark^\checkmark\\checkmarkc\checkmarki\checkmarkr\checkmarkc\checkmark$\checkmark=\checkmark$\checkmark^\checkmark\\checkmarkc\checkmarki\checkmarkr\checkmarkc\checkmark$\checkmark \checkmark$\checkmark^\checkmark\\checkmarkc\checkmarki\checkmarkr\checkmarkc\checkmark$\checkmark1\checkmark$\checkmark^\checkmark\\checkmarkc\checkmarki\checkmarkr\checkmarkc\checkmark$\checkmark0\checkmark$\checkmark^\checkmark\\checkmarkc\checkmarki\checkmarkr\checkmarkc\checkmark$\checkmark0\checkmark$\checkmark^\checkmark\\checkmarkc\checkmarki\checkmarkr\checkmarkc\checkmark$\checkmark \checkmark$\checkmark^\checkmark\\checkmarkc\checkmarki\checkmarkr\checkmarkc\checkmark$\checkmark\\checkmark$\checkmark^\checkmark\\checkmarkc\checkmarki\checkmarkr\checkmarkc\checkmark$\checkmarkt\checkmark$\checkmark^\checkmark\\checkmarkc\checkmarki\checkmarkr\checkmarkc\checkmark$\checkmarki\checkmark$\checkmark^\checkmark\\checkmarkc\checkmarki\checkmarkr\checkmarkc\checkmark$\checkmarkm\checkmark$\checkmark^\checkmark\\checkmarkc\checkmarki\checkmarkr\checkmarkc\checkmark$\checkmarke\checkmark$\checkmark^\checkmark\\checkmarkc\checkmarki\checkmarkr\checkmarkc\checkmark$\checkmarks\checkmark$\checkmark^\checkmark\\checkmarkc\checkmarki\checkmarkr\checkmarkc\checkmark$\checkmark \checkmark$\checkmark^\checkmark\\checkmarkc\checkmarki\checkmarkr\checkmarkc\checkmark$\checkmark\\checkmark$\checkmark^\checkmark\\checkmarkc\checkmarki\checkmarkr\checkmarkc\checkmark$\checkmarkt\checkmark$\checkmark^\checkmark\\checkmarkc\checkmarki\checkmarkr\checkmarkc\checkmark$\checkmarka\checkmark$\checkmark^\checkmark\\checkmarkc\checkmarki\checkmarkr\checkmarkc\checkmark$\checkmarkn\checkmark$\checkmark^\checkmark\\checkmarkc\checkmarki\checkmarkr\checkmarkc\checkmark$\checkmark(\checkmark$\checkmark^\checkmark\\checkmarkc\checkmarki\checkmarkr\checkmarkc\checkmark$\checkmark1\checkmark$\checkmark^\checkmark\\checkmarkc\checkmarki\checkmarkr\checkmarkc\checkmark$\checkmark°\checkmark$\checkmark^\checkmark\\checkmarkc\checkmarki\checkmarkr\checkmarkc\checkmark$\checkmark)\checkmark$\checkmark^\checkmark\\checkmarkc\checkmarki\checkmarkr\checkmarkc\checkmark$\checkmark \checkmark$\checkmark^\checkmark\\checkmarkc\checkmarki\checkmarkr\checkmarkc\checkmark$\checkmark=\checkmark$\checkmark^\checkmark\\checkmarkc\checkmarki\checkmarkr\checkmarkc\checkmark$\checkmark \checkmark$\checkmark^\checkmark\\checkmarkc\checkmarki\checkmarkr\checkmarkc\checkmark$\checkmark\\checkmark$\checkmark^\checkmark\\checkmarkc\checkmarki\checkmarkr\checkmarkc\checkmark$\checkmarkb\checkmark$\checkmark^\checkmark\\checkmarkc\checkmarki\checkmarkr\checkmarkc\checkmark$\checkmarko\checkmark$\checkmark^\checkmark\\checkmarkc\checkmarki\checkmarkr\checkmarkc\checkmark$\checkmarkx\checkmark$\checkmark^\checkmark\\checkmarkc\checkmarki\checkmarkr\checkmarkc\checkmark$\checkmarke\checkmark$\checkmark^\checkmark\\checkmarkc\checkmarki\checkmarkr\checkmarkc\checkmark$\checkmarkd\checkmark$\checkmark^\checkmark\\checkmarkc\checkmarki\checkmarkr\checkmarkc\checkmark$\checkmark{\checkmark$\checkmark^\checkmark\\checkmarkc\checkmarki\checkmarkr\checkmarkc\checkmark$\checkmark1\checkmark$\checkmark^\checkmark\\checkmarkc\checkmarki\checkmarkr\checkmarkc\checkmark$\checkmark{\checkmark$\checkmark^\checkmark\\checkmarkc\checkmarki\checkmarkr\checkmarkc\checkmark$\checkmark,\checkmark$\checkmark^\checkmark\\checkmarkc\checkmarki\checkmarkr\checkmarkc\checkmark$\checkmark}\checkmark$\checkmark^\checkmark\\checkmarkc\checkmarki\checkmarkr\checkmarkc\checkmark$\checkmark7\checkmark$\checkmark^\checkmark\\checkmarkc\checkmarki\checkmarkr\checkmarkc\checkmark$\checkmark4\checkmark$\checkmark^\checkmark\\checkmarkc\checkmarki\checkmarkr\checkmarkc\checkmark$\checkmark \checkmark$\checkmark^\checkmark\\checkmarkc\checkmarki\checkmarkr\checkmarkc\checkmark$\checkmark\\checkmark$\checkmark^\checkmark\\checkmarkc\checkmarki\checkmarkr\checkmarkc\checkmark$\checkmarkt\checkmark$\checkmark^\checkmark\\checkmarkc\checkmarki\checkmarkr\checkmarkc\checkmark$\checkmarke\checkmark$\checkmark^\checkmark\\checkmarkc\checkmarki\checkmarkr\checkmarkc\checkmark$\checkmarkx\checkmark$\checkmark^\checkmark\\checkmarkc\checkmarki\checkmarkr\checkmarkc\checkmark$\checkmarkt\checkmark$\checkmark^\checkmark\\checkmarkc\checkmarki\checkmarkr\checkmarkc\checkmark$\checkmark{\checkmark$\checkmark^\checkmark\\checkmarkc\checkmarki\checkmarkr\checkmarkc\checkmark$\checkmark \checkmark$\checkmark^\checkmark\\checkmarkc\checkmarki\checkmarkr\checkmarkc\checkmark$\checkmarkm\checkmark$\checkmark^\checkmark\\checkmarkc\checkmarki\checkmarkr\checkmarkc\checkmark$\checkmarkm\checkmark$\checkmark^\checkmark\\checkmarkc\checkmarki\checkmarkr\checkmarkc\checkmark$\checkmark}\checkmark$\checkmark^\checkmark\\checkmarkc\checkmarki\checkmarkr\checkmarkc\checkmark$\checkmark}\checkmark$\checkmark^\checkmark\\checkmarkc\checkmarki\checkmarkr\checkmarkc\checkmark$\checkmark
\checkmark$\checkmark^\checkmark\\checkmarkc\checkmarki\checkmarkr\checkmarkc\checkmark$\checkmark\\checkmark$\checkmark^\checkmark\\checkmarkc\checkmarki\checkmarkr\checkmarkc\checkmark$\checkmarke\checkmark$\checkmark^\checkmark\\checkmarkc\checkmarki\checkmarkr\checkmarkc\checkmark$\checkmarkn\checkmark$\checkmark^\checkmark\\checkmarkc\checkmarki\checkmarkr\checkmarkc\checkmark$\checkmarkd\checkmark$\checkmark^\checkmark\\checkmarkc\checkmarki\checkmarkr\checkmarkc\checkmark$\checkmark{\checkmark$\checkmark^\checkmark\\checkmarkc\checkmarki\checkmarkr\checkmarkc\checkmark$\checkmarke\checkmark$\checkmark^\checkmark\\checkmarkc\checkmarki\checkmarkr\checkmarkc\checkmark$\checkmarkq\checkmark$\checkmark^\checkmark\\checkmarkc\checkmarki\checkmarkr\checkmarkc\checkmark$\checkmarku\checkmark$\checkmark^\checkmark\\checkmarkc\checkmarki\checkmarkr\checkmarkc\checkmark$\checkmarka\checkmark$\checkmark^\checkmark\\checkmarkc\checkmarki\checkmarkr\checkmarkc\checkmark$\checkmarkt\checkmark$\checkmark^\checkmark\\checkmarkc\checkmarki\checkmarkr\checkmarkc\checkmark$\checkmarki\checkmark$\checkmark^\checkmark\\checkmarkc\checkmarki\checkmarkr\checkmarkc\checkmark$\checkmarko\checkmark$\checkmark^\checkmark\\checkmarkc\checkmarki\checkmarkr\checkmarkc\checkmark$\checkmarkn\checkmark$\checkmark^\checkmark\\checkmarkc\checkmarki\checkmarkr\checkmarkc\checkmark$\checkmark}\checkmark$\checkmark^\checkmark\\checkmarkc\checkmarki\checkmarkr\checkmarkc\checkmark$\checkmark
\checkmark$\checkmark^\checkmark\\checkmarkc\checkmarki\checkmarkr\checkmarkc\checkmark$\checkmarkC\checkmark$\checkmark^\checkmark\\checkmarkc\checkmarki\checkmarkr\checkmarkc\checkmark$\checkmarke\checkmark$\checkmark^\checkmark\\checkmarkc\checkmarki\checkmarkr\checkmarkc\checkmark$\checkmarkt\checkmark$\checkmark^\checkmark\\checkmarkc\checkmarki\checkmarkr\checkmarkc\checkmark$\checkmarkt\checkmark$\checkmark^\checkmark\\checkmarkc\checkmarki\checkmarkr\checkmarkc\checkmark$\checkmarke\checkmark$\checkmark^\checkmark\\checkmarkc\checkmarki\checkmarkr\checkmarkc\checkmark$\checkmark \checkmark$\checkmark^\checkmark\\checkmarkc\checkmarki\checkmarkr\checkmarkc\checkmark$\checkmarke\checkmark$\checkmark^\checkmark\\checkmarkc\checkmarki\checkmarkr\checkmarkc\checkmark$\checkmarkr\checkmark$\checkmark^\checkmark\\checkmarkc\checkmarki\checkmarkr\checkmarkc\checkmark$\checkmarkr\checkmark$\checkmark^\checkmark\\checkmarkc\checkmarki\checkmarkr\checkmarkc\checkmark$\checkmarke\checkmark$\checkmark^\checkmark\\checkmarkc\checkmarki\checkmarkr\checkmarkc\checkmark$\checkmarku\checkmark$\checkmark^\checkmark\\checkmarkc\checkmarki\checkmarkr\checkmarkc\checkmark$\checkmarkr\checkmark$\checkmark^\checkmark\\checkmarkc\checkmarki\checkmarkr\checkmarkc\checkmark$\checkmark \checkmark$\checkmark^\checkmark\\checkmarkc\checkmarki\checkmarkr\checkmarkc\checkmark$\checkmarkd\checkmark$\checkmark^\checkmark\\checkmarkc\checkmarki\checkmarkr\checkmarkc\checkmark$\checkmarke\checkmark$\checkmark^\checkmark\\checkmarkc\checkmarki\checkmarkr\checkmarkc\checkmark$\checkmark \checkmark$\checkmark^\checkmark\\checkmarkc\checkmarki\checkmarkr\checkmarkc\checkmark$\checkmark\\checkmark$\checkmark^\checkmark\\checkmarkc\checkmarki\checkmarkr\checkmarkc\checkmark$\checkmarkt\checkmark$\checkmark^\checkmark\\checkmarkc\checkmarki\checkmarkr\checkmarkc\checkmark$\checkmarke\checkmark$\checkmark^\checkmark\\checkmarkc\checkmarki\checkmarkr\checkmarkc\checkmark$\checkmarkx\checkmark$\checkmark^\checkmark\\checkmarkc\checkmarki\checkmarkr\checkmarkc\checkmark$\checkmarkt\checkmark$\checkmark^\checkmark\\checkmarkc\checkmarki\checkmarkr\checkmarkc\checkmark$\checkmarkb\checkmark$\checkmark^\checkmark\\checkmarkc\checkmarki\checkmarkr\checkmarkc\checkmark$\checkmarkf\checkmark$\checkmark^\checkmark\\checkmarkc\checkmarki\checkmarkr\checkmarkc\checkmark$\checkmark{\checkmark$\checkmark^\checkmark\\checkmarkc\checkmarki\checkmarkr\checkmarkc\checkmark$\checkmark1\checkmark$\checkmark^\checkmark\\checkmarkc\checkmarki\checkmarkr\checkmarkc\checkmark$\checkmark,\checkmark$\checkmark^\checkmark\\checkmarkc\checkmarki\checkmarkr\checkmarkc\checkmark$\checkmark7\checkmark$\checkmark^\checkmark\\checkmarkc\checkmarki\checkmarkr\checkmarkc\checkmark$\checkmark4\checkmark$\checkmark^\checkmark\\checkmarkc\checkmarki\checkmarkr\checkmarkc\checkmark$\checkmark \checkmark$\checkmark^\checkmark\\checkmarkc\checkmarki\checkmarkr\checkmarkc\checkmark$\checkmarkm\checkmark$\checkmark^\checkmark\\checkmarkc\checkmarki\checkmarkr\checkmarkc\checkmark$\checkmarkm\checkmark$\checkmark^\checkmark\\checkmarkc\checkmarki\checkmarkr\checkmarkc\checkmark$\checkmark}\checkmark$\checkmark^\checkmark\\checkmarkc\checkmarki\checkmarkr\checkmarkc\checkmark$\checkmark \checkmark$\checkmark^\checkmark\\checkmarkc\checkmarki\checkmarkr\checkmarkc\checkmark$\checkmarke\checkmark$\checkmark^\checkmark\\checkmarkc\checkmarki\checkmarkr\checkmarkc\checkmark$\checkmarks\checkmark$\checkmark^\checkmark\\checkmarkc\checkmarki\checkmarkr\checkmarkc\checkmark$\checkmarkt\checkmark$\checkmark^\checkmark\\checkmarkc\checkmarki\checkmarkr\checkmarkc\checkmark$\checkmark \checkmark$\checkmark^\checkmark\\checkmarkc\checkmarki\checkmarkr\checkmarkc\checkmark$\checkmarkt\checkmark$\checkmark^\checkmark\\checkmarkc\checkmarki\checkmarkr\checkmarkc\checkmark$\checkmarko\checkmark$\checkmark^\checkmark\\checkmarkc\checkmarki\checkmarkr\checkmarkc\checkmark$\checkmarkt\checkmark$\checkmark^\checkmark\\checkmarkc\checkmarki\checkmarkr\checkmarkc\checkmark$\checkmarka\checkmark$\checkmark^\checkmark\\checkmarkc\checkmarki\checkmarkr\checkmarkc\checkmark$\checkmarkl\checkmark$\checkmark^\checkmark\\checkmarkc\checkmarki\checkmarkr\checkmarkc\checkmark$\checkmarke\checkmark$\checkmark^\checkmark\\checkmarkc\checkmarki\checkmarkr\checkmarkc\checkmark$\checkmarkm\checkmark$\checkmark^\checkmark\\checkmarkc\checkmarki\checkmarkr\checkmarkc\checkmark$\checkmarke\checkmark$\checkmark^\checkmark\\checkmarkc\checkmarki\checkmarkr\checkmarkc\checkmark$\checkmarkn\checkmark$\checkmark^\checkmark\\checkmarkc\checkmarki\checkmarkr\checkmarkc\checkmark$\checkmarkt\checkmark$\checkmark^\checkmark\\checkmarkc\checkmarki\checkmarkr\checkmarkc\checkmark$\checkmark \checkmark$\checkmark^\checkmark\\checkmarkc\checkmarki\checkmarkr\checkmarkc\checkmark$\checkmarki\checkmark$\checkmark^\checkmark\\checkmarkc\checkmarki\checkmarkr\checkmarkc\checkmark$\checkmarkn\checkmark$\checkmark^\checkmark\\checkmarkc\checkmarki\checkmarkr\checkmarkc\checkmark$\checkmarka\checkmark$\checkmark^\checkmark\\checkmarkc\checkmarki\checkmarkr\checkmarkc\checkmark$\checkmarkc\checkmark$\checkmark^\checkmark\\checkmarkc\checkmarki\checkmarkr\checkmarkc\checkmark$\checkmarkc\checkmark$\checkmark^\checkmark\\checkmarkc\checkmarki\checkmarkr\checkmarkc\checkmark$\checkmarke\checkmark$\checkmark^\checkmark\\checkmarkc\checkmarki\checkmarkr\checkmarkc\checkmark$\checkmarkp\checkmark$\checkmark^\checkmark\\checkmarkc\checkmarki\checkmarkr\checkmarkc\checkmark$\checkmarkt\checkmark$\checkmark^\checkmark\\checkmarkc\checkmarki\checkmarkr\checkmarkc\checkmark$\checkmarka\checkmark$\checkmark^\checkmark\\checkmarkc\checkmarki\checkmarkr\checkmarkc\checkmark$\checkmarkb\checkmark$\checkmark^\checkmark\\checkmarkc\checkmarki\checkmarkr\checkmarkc\checkmark$\checkmarkl\checkmark$\checkmark^\checkmark\\checkmarkc\checkmarki\checkmarkr\checkmarkc\checkmark$\checkmarke\checkmark$\checkmark^\checkmark\\checkmarkc\checkmarki\checkmarkr\checkmarkc\checkmark$\checkmark \checkmark$\checkmark^\checkmark\\checkmarkc\checkmarki\checkmarkr\checkmarkc\checkmark$\checkmarka\checkmark$\checkmark^\checkmark\\checkmarkc\checkmarki\checkmarkr\checkmarkc\checkmark$\checkmarku\checkmark$\checkmark^\checkmark\\checkmarkc\checkmarki\checkmarkr\checkmarkc\checkmark$\checkmark \checkmark$\checkmark^\checkmark\\checkmarkc\checkmarki\checkmarkr\checkmarkc\checkmark$\checkmarkr\checkmark$\checkmark^\checkmark\\checkmarkc\checkmarki\checkmarkr\checkmarkc\checkmark$\checkmarke\checkmark$\checkmark^\checkmark\\checkmarkc\checkmarki\checkmarkr\checkmarkc\checkmark$\checkmarkg\checkmark$\checkmark^\checkmark\\checkmarkc\checkmarki\checkmarkr\checkmarkc\checkmark$\checkmarka\checkmark$\checkmark^\checkmark\\checkmarkc\checkmarki\checkmarkr\checkmarkc\checkmark$\checkmarkr\checkmark$\checkmark^\checkmark\\checkmarkc\checkmarki\checkmarkr\checkmarkc\checkmark$\checkmarkd\checkmark$\checkmark^\checkmark\\checkmarkc\checkmarki\checkmarkr\checkmarkc\checkmark$\checkmark \checkmark$\checkmark^\checkmark\\checkmarkc\checkmarki\checkmarkr\checkmarkc\checkmark$\checkmarkd\checkmark$\checkmark^\checkmark\\checkmarkc\checkmarki\checkmarkr\checkmarkc\checkmark$\checkmarke\checkmark$\checkmark^\checkmark\\checkmarkc\checkmarki\checkmarkr\checkmarkc\checkmark$\checkmark \checkmark$\checkmark^\checkmark\\checkmarkc\checkmarki\checkmarkr\checkmarkc\checkmark$\checkmarkl\checkmark$\checkmark^\checkmark\\checkmarkc\checkmarki\checkmarkr\checkmarkc\checkmark$\checkmarka\checkmark$\checkmark^\checkmark\\checkmarkc\checkmarki\checkmarkr\checkmarkc\checkmark$\checkmark \checkmark$\checkmark^\checkmark\\checkmarkc\checkmarki\checkmarkr\checkmarkc\checkmark$\checkmarkp\checkmark$\checkmark^\checkmark\\checkmarkc\checkmarki\checkmarkr\checkmarkc\checkmark$\checkmarkr\checkmark$\checkmark^\checkmark\\checkmarkc\checkmarki\checkmarkr\checkmarkc\checkmark$\checkmarké\checkmark$\checkmark^\checkmark\\checkmarkc\checkmarki\checkmarkr\checkmarkc\checkmark$\checkmarkc\checkmark$\checkmark^\checkmark\\checkmarkc\checkmarki\checkmarkr\checkmarkc\checkmark$\checkmarki\checkmark$\checkmark^\checkmark\\checkmarkc\checkmarki\checkmarkr\checkmarkc\checkmark$\checkmarks\checkmark$\checkmark^\checkmark\\checkmarkc\checkmarki\checkmarkr\checkmarkc\checkmark$\checkmarki\checkmark$\checkmark^\checkmark\\checkmarkc\checkmarki\checkmarkr\checkmarkc\checkmark$\checkmarko\checkmark$\checkmark^\checkmark\\checkmarkc\checkmarki\checkmarkr\checkmarkc\checkmark$\checkmarkn\checkmark$\checkmark^\checkmark\\checkmarkc\checkmarki\checkmarkr\checkmarkc\checkmark$\checkmark \checkmark$\checkmark^\checkmark\\checkmarkc\checkmarki\checkmarkr\checkmarkc\checkmark$\checkmarkr\checkmark$\checkmark^\checkmark\\checkmarkc\checkmarki\checkmarkr\checkmarkc\checkmark$\checkmarke\checkmark$\checkmark^\checkmark\\checkmarkc\checkmarki\checkmarkr\checkmarkc\checkmark$\checkmarkq\checkmark$\checkmark^\checkmark\\checkmarkc\checkmarki\checkmarkr\checkmarkc\checkmark$\checkmarku\checkmark$\checkmark^\checkmark\\checkmarkc\checkmarki\checkmarkr\checkmarkc\checkmark$\checkmarki\checkmark$\checkmark^\checkmark\\checkmarkc\checkmarki\checkmarkr\checkmarkc\checkmark$\checkmarks\checkmark$\checkmark^\checkmark\\checkmarkc\checkmarki\checkmarkr\checkmarkc\checkmark$\checkmarke\checkmark$\checkmark^\checkmark\\checkmarkc\checkmarki\checkmarkr\checkmarkc\checkmark$\checkmark \checkmark$\checkmark^\checkmark\\checkmarkc\checkmarki\checkmarkr\checkmarkc\checkmark$\checkmarkd\checkmark$\checkmark^\checkmark\\checkmarkc\checkmarki\checkmarkr\checkmarkc\checkmark$\checkmarke\checkmark$\checkmark^\checkmark\\checkmarkc\checkmarki\checkmarkr\checkmarkc\checkmark$\checkmark \checkmark$\checkmark^\checkmark\\checkmarkc\checkmarki\checkmarkr\checkmarkc\checkmark$\checkmark$\checkmark$\checkmark^\checkmark\\checkmarkc\checkmarki\checkmarkr\checkmarkc\checkmark$\checkmark\\checkmark$\checkmark^\checkmark\\checkmarkc\checkmarki\checkmarkr\checkmarkc\checkmark$\checkmarkp\checkmark$\checkmark^\checkmark\\checkmarkc\checkmarki\checkmarkr\checkmarkc\checkmark$\checkmarkm\checkmark$\checkmark^\checkmark\\checkmarkc\checkmarki\checkmarkr\checkmarkc\checkmark$\checkmark \checkmark$\checkmark^\checkmark\\checkmarkc\checkmarki\checkmarkr\checkmarkc\checkmark$\checkmark0\checkmark$\checkmark^\checkmark\\checkmarkc\checkmarki\checkmarkr\checkmarkc\checkmark$\checkmark{\checkmark$\checkmark^\checkmark\\checkmarkc\checkmarki\checkmarkr\checkmarkc\checkmark$\checkmark,\checkmark$\checkmark^\checkmark\\checkmarkc\checkmarki\checkmarkr\checkmarkc\checkmark$\checkmark}\checkmark$\checkmark^\checkmark\\checkmarkc\checkmarki\checkmarkr\checkmarkc\checkmark$\checkmark0\checkmark$\checkmark^\checkmark\\checkmarkc\checkmarki\checkmarkr\checkmarkc\checkmark$\checkmark5\checkmark$\checkmark^\checkmark\\checkmarkc\checkmarki\checkmarkr\checkmarkc\checkmark$\checkmark$\checkmark$\checkmark^\checkmark\\checkmarkc\checkmarki\checkmarkr\checkmarkc\checkmark$\checkmark \checkmark$\checkmark^\checkmark\\checkmarkc\checkmarki\checkmarkr\checkmarkc\checkmark$\checkmarkm\checkmark$\checkmark^\checkmark\\checkmarkc\checkmarki\checkmarkr\checkmarkc\checkmark$\checkmarkm\checkmark$\checkmark^\checkmark\\checkmarkc\checkmarki\checkmarkr\checkmarkc\checkmark$\checkmark.\checkmark$\checkmark^\checkmark\\checkmarkc\checkmarki\checkmarkr\checkmarkc\checkmark$\checkmark \checkmark$\checkmark^\checkmark\\checkmarkc\checkmarki\checkmarkr\checkmarkc\checkmark$\checkmarkC\checkmark$\checkmark^\checkmark\\checkmarkc\checkmarki\checkmarkr\checkmarkc\checkmark$\checkmarke\checkmark$\checkmark^\checkmark\\checkmarkc\checkmarki\checkmarkr\checkmarkc\checkmark$\checkmarkl\checkmark$\checkmark^\checkmark\\checkmarkc\checkmarki\checkmarkr\checkmarkc\checkmark$\checkmarka\checkmark$\checkmark^\checkmark\\checkmarkc\checkmarki\checkmarkr\checkmarkc\checkmark$\checkmark \checkmark$\checkmark^\checkmark\\checkmarkc\checkmarki\checkmarkr\checkmarkc\checkmark$\checkmarkj\checkmark$\checkmark^\checkmark\\checkmarkc\checkmarki\checkmarkr\checkmarkc\checkmark$\checkmarku\checkmark$\checkmark^\checkmark\\checkmarkc\checkmarki\checkmarkr\checkmarkc\checkmark$\checkmarks\checkmark$\checkmark^\checkmark\\checkmarkc\checkmarki\checkmarkr\checkmarkc\checkmark$\checkmarkt\checkmark$\checkmark^\checkmark\\checkmarkc\checkmarki\checkmarkr\checkmarkc\checkmark$\checkmarki\checkmark$\checkmark^\checkmark\\checkmarkc\checkmarki\checkmarkr\checkmarkc\checkmark$\checkmarkf\checkmark$\checkmark^\checkmark\\checkmarkc\checkmarki\checkmarkr\checkmarkc\checkmark$\checkmarki\checkmark$\checkmark^\checkmark\\checkmarkc\checkmarki\checkmarkr\checkmarkc\checkmark$\checkmarke\checkmark$\checkmark^\checkmark\\checkmarkc\checkmarki\checkmarkr\checkmarkc\checkmark$\checkmark \checkmark$\checkmark^\checkmark\\checkmarkc\checkmarki\checkmarkr\checkmarkc\checkmark$\checkmarki\checkmark$\checkmark^\checkmark\\checkmarkc\checkmarki\checkmarkr\checkmarkc\checkmark$\checkmarkm\checkmark$\checkmark^\checkmark\\checkmarkc\checkmarki\checkmarkr\checkmarkc\checkmark$\checkmarkp\checkmark$\checkmark^\checkmark\\checkmarkc\checkmarki\checkmarkr\checkmarkc\checkmark$\checkmarké\checkmark$\checkmark^\checkmark\\checkmarkc\checkmarki\checkmarkr\checkmarkc\checkmark$\checkmarkr\checkmark$\checkmark^\checkmark\\checkmarkc\checkmarki\checkmarkr\checkmarkc\checkmark$\checkmarka\checkmark$\checkmark^\checkmark\\checkmarkc\checkmarki\checkmarkr\checkmarkc\checkmark$\checkmarkt\checkmark$\checkmark^\checkmark\\checkmarkc\checkmarki\checkmarkr\checkmarkc\checkmark$\checkmarki\checkmark$\checkmark^\checkmark\\checkmarkc\checkmarki\checkmarkr\checkmarkc\checkmark$\checkmarkv\checkmark$\checkmark^\checkmark\\checkmarkc\checkmarki\checkmarkr\checkmarkc\checkmark$\checkmarke\checkmark$\checkmark^\checkmark\\checkmarkc\checkmarki\checkmarkr\checkmarkc\checkmark$\checkmarkm\checkmark$\checkmark^\checkmark\\checkmarkc\checkmarki\checkmarkr\checkmarkc\checkmark$\checkmarke\checkmark$\checkmark^\checkmark\\checkmarkc\checkmarki\checkmarkr\checkmarkc\checkmark$\checkmarkn\checkmark$\checkmark^\checkmark\\checkmarkc\checkmarki\checkmarkr\checkmarkc\checkmark$\checkmarkt\checkmark$\checkmark^\checkmark\\checkmarkc\checkmarki\checkmarkr\checkmarkc\checkmark$\checkmark \checkmark$\checkmark^\checkmark\\checkmarkc\checkmarki\checkmarkr\checkmarkc\checkmark$\checkmarku\checkmark$\checkmark^\checkmark\\checkmarkc\checkmarki\checkmarkr\checkmarkc\checkmark$\checkmarkn\checkmark$\checkmark^\checkmark\\checkmarkc\checkmarki\checkmarkr\checkmarkc\checkmark$\checkmark \checkmark$\checkmark^\checkmark\\checkmarkc\checkmarki\checkmarkr\checkmarkc\checkmark$\checkmarkr\checkmark$\checkmark^\checkmark\\checkmarkc\checkmarki\checkmarkr\checkmarkc\checkmark$\checkmarké\checkmark$\checkmark^\checkmark\\checkmarkc\checkmarki\checkmarkr\checkmarkc\checkmark$\checkmarkd\checkmark$\checkmark^\checkmark\\checkmarkc\checkmarki\checkmarkr\checkmarkc\checkmark$\checkmarku\checkmark$\checkmark^\checkmark\\checkmarkc\checkmarki\checkmarkr\checkmarkc\checkmark$\checkmarkc\checkmark$\checkmark^\checkmark\\checkmarkc\checkmarki\checkmarkr\checkmarkc\checkmark$\checkmarkt\checkmark$\checkmark^\checkmark\\checkmarkc\checkmarki\checkmarkr\checkmarkc\checkmark$\checkmarke\checkmark$\checkmark^\checkmark\\checkmarkc\checkmarki\checkmarkr\checkmarkc\checkmark$\checkmarku\checkmark$\checkmark^\checkmark\\checkmarkc\checkmarki\checkmarkr\checkmarkc\checkmark$\checkmarkr\checkmark$\checkmark^\checkmark\\checkmarkc\checkmarki\checkmarkr\checkmarkc\checkmark$\checkmark \checkmark$\checkmark^\checkmark\\checkmarkc\checkmarki\checkmarkr\checkmarkc\checkmark$\checkmarkà\checkmark$\checkmark^\checkmark\\checkmarkc\checkmarki\checkmarkr\checkmarkc\checkmark$\checkmark \checkmark$\checkmark^\checkmark\\checkmarkc\checkmarki\checkmarkr\checkmarkc\checkmark$\checkmarkr\checkmark$\checkmark^\checkmark\\checkmarkc\checkmarki\checkmarkr\checkmarkc\checkmark$\checkmarka\checkmark$\checkmark^\checkmark\\checkmarkc\checkmarki\checkmarkr\checkmarkc\checkmark$\checkmarkp\checkmark$\checkmark^\checkmark\\checkmarkc\checkmarki\checkmarkr\checkmarkc\checkmark$\checkmarkp\checkmark$\checkmark^\checkmark\\checkmarkc\checkmarki\checkmarkr\checkmarkc\checkmark$\checkmarko\checkmark$\checkmark^\checkmark\\checkmarkc\checkmarki\checkmarkr\checkmarkc\checkmark$\checkmarkr\checkmark$\checkmark^\checkmark\\checkmarkc\checkmarki\checkmarkr\checkmarkc\checkmark$\checkmarkt\checkmark$\checkmark^\checkmark\\checkmarkc\checkmarki\checkmarkr\checkmarkc\checkmark$\checkmark \checkmark$\checkmark^\checkmark\\checkmarkc\checkmarki\checkmarkr\checkmarkc\checkmark$\checkmarks\checkmark$\checkmark^\checkmark\\checkmarkc\checkmarki\checkmarkr\checkmarkc\checkmark$\checkmarké\checkmark$\checkmark^\checkmark\\checkmarkc\checkmarki\checkmarkr\checkmarkc\checkmark$\checkmarkv\checkmark$\checkmark^\checkmark\\checkmarkc\checkmarki\checkmarkr\checkmarkc\checkmark$\checkmarkè\checkmark$\checkmark^\checkmark\\checkmarkc\checkmarki\checkmarkr\checkmarkc\checkmark$\checkmarkr\checkmark$\checkmark^\checkmark\\checkmarkc\checkmarki\checkmarkr\checkmarkc\checkmark$\checkmarke\checkmark$\checkmark^\checkmark\\checkmarkc\checkmarki\checkmarkr\checkmarkc\checkmark$\checkmark.\checkmark$\checkmark^\checkmark\\checkmarkc\checkmarki\checkmarkr\checkmarkc\checkmark$\checkmark
\checkmark$\checkmark^\checkmark\\checkmarkc\checkmarki\checkmarkr\checkmarkc\checkmark$\checkmark
\checkmark$\checkmark^\checkmark\\checkmarkc\checkmarki\checkmarkr\checkmarkc\checkmark$\checkmark\\checkmark$\checkmark^\checkmark\\checkmarkc\checkmarki\checkmarkr\checkmarkc\checkmark$\checkmarks\checkmark$\checkmark^\checkmark\\checkmarkc\checkmarki\checkmarkr\checkmarkc\checkmark$\checkmarku\checkmark$\checkmark^\checkmark\\checkmarkc\checkmarki\checkmarkr\checkmarkc\checkmark$\checkmarkb\checkmark$\checkmark^\checkmark\\checkmarkc\checkmarki\checkmarkr\checkmarkc\checkmark$\checkmarks\checkmark$\checkmark^\checkmark\\checkmarkc\checkmarki\checkmarkr\checkmarkc\checkmark$\checkmarke\checkmark$\checkmark^\checkmark\\checkmarkc\checkmarki\checkmarkr\checkmarkc\checkmark$\checkmarkc\checkmark$\checkmark^\checkmark\\checkmarkc\checkmarki\checkmarkr\checkmarkc\checkmark$\checkmarkt\checkmark$\checkmark^\checkmark\\checkmarkc\checkmarki\checkmarkr\checkmarkc\checkmark$\checkmarki\checkmark$\checkmark^\checkmark\\checkmarkc\checkmarki\checkmarkr\checkmarkc\checkmark$\checkmarko\checkmark$\checkmark^\checkmark\\checkmarkc\checkmarki\checkmarkr\checkmarkc\checkmark$\checkmarkn\checkmark$\checkmark^\checkmark\\checkmarkc\checkmarki\checkmarkr\checkmarkc\checkmark$\checkmark{\checkmark$\checkmark^\checkmark\\checkmarkc\checkmarki\checkmarkr\checkmarkc\checkmark$\checkmarkC\checkmark$\checkmark^\checkmark\\checkmarkc\checkmarki\checkmarkr\checkmarkc\checkmark$\checkmarka\checkmark$\checkmark^\checkmark\\checkmarkc\checkmarki\checkmarkr\checkmarkc\checkmark$\checkmarkl\checkmark$\checkmark^\checkmark\\checkmarkc\checkmarki\checkmarkr\checkmarkc\checkmark$\checkmarkc\checkmark$\checkmark^\checkmark\\checkmarkc\checkmarki\checkmarkr\checkmarkc\checkmark$\checkmarku\checkmark$\checkmark^\checkmark\\checkmarkc\checkmarki\checkmarkr\checkmarkc\checkmark$\checkmarkl\checkmark$\checkmark^\checkmark\\checkmarkc\checkmarki\checkmarkr\checkmarkc\checkmark$\checkmark \checkmark$\checkmark^\checkmark\\checkmarkc\checkmarki\checkmarkr\checkmarkc\checkmark$\checkmarkd\checkmark$\checkmark^\checkmark\\checkmarkc\checkmarki\checkmarkr\checkmarkc\checkmark$\checkmarke\checkmark$\checkmark^\checkmark\\checkmarkc\checkmarki\checkmarkr\checkmarkc\checkmark$\checkmark \checkmark$\checkmark^\checkmark\\checkmarkc\checkmarki\checkmarkr\checkmarkc\checkmark$\checkmarkl\checkmark$\checkmark^\checkmark\\checkmarkc\checkmarki\checkmarkr\checkmarkc\checkmark$\checkmarka\checkmark$\checkmark^\checkmark\\checkmarkc\checkmarki\checkmarkr\checkmarkc\checkmark$\checkmark \checkmark$\checkmark^\checkmark\\checkmarkc\checkmarki\checkmarkr\checkmarkc\checkmark$\checkmarkr\checkmark$\checkmark^\checkmark\\checkmarkc\checkmarki\checkmarkr\checkmarkc\checkmark$\checkmarké\checkmark$\checkmark^\checkmark\\checkmarkc\checkmarki\checkmarkr\checkmarkc\checkmark$\checkmarks\checkmark$\checkmark^\checkmark\\checkmarkc\checkmarki\checkmarkr\checkmarkc\checkmark$\checkmarko\checkmark$\checkmark^\checkmark\\checkmarkc\checkmarki\checkmarkr\checkmarkc\checkmark$\checkmarkl\checkmark$\checkmark^\checkmark\\checkmarkc\checkmarki\checkmarkr\checkmarkc\checkmark$\checkmarku\checkmark$\checkmark^\checkmark\\checkmarkc\checkmarki\checkmarkr\checkmarkc\checkmark$\checkmarkt\checkmark$\checkmark^\checkmark\\checkmarkc\checkmarki\checkmarkr\checkmarkc\checkmark$\checkmarki\checkmark$\checkmark^\checkmark\\checkmarkc\checkmarki\checkmarkr\checkmarkc\checkmark$\checkmarko\checkmark$\checkmark^\checkmark\\checkmarkc\checkmarki\checkmarkr\checkmarkc\checkmark$\checkmarkn\checkmark$\checkmark^\checkmark\\checkmarkc\checkmarki\checkmarkr\checkmarkc\checkmark$\checkmark \checkmark$\checkmark^\checkmark\\checkmarkc\checkmarki\checkmarkr\checkmarkc\checkmark$\checkmarka\checkmark$\checkmark^\checkmark\\checkmarkc\checkmarki\checkmarkr\checkmarkc\checkmark$\checkmarkn\checkmark$\checkmark^\checkmark\\checkmarkc\checkmarki\checkmarkr\checkmarkc\checkmark$\checkmarkg\checkmark$\checkmark^\checkmark\\checkmarkc\checkmarki\checkmarkr\checkmarkc\checkmark$\checkmarku\checkmark$\checkmark^\checkmark\\checkmarkc\checkmarki\checkmarkr\checkmarkc\checkmark$\checkmarkl\checkmark$\checkmark^\checkmark\\checkmarkc\checkmarki\checkmarkr\checkmarkc\checkmark$\checkmarka\checkmark$\checkmark^\checkmark\\checkmarkc\checkmarki\checkmarkr\checkmarkc\checkmark$\checkmarki\checkmark$\checkmark^\checkmark\\checkmarkc\checkmarki\checkmarkr\checkmarkc\checkmark$\checkmarkr\checkmark$\checkmark^\checkmark\\checkmarkc\checkmarki\checkmarkr\checkmarkc\checkmark$\checkmarke\checkmark$\checkmark^\checkmark\\checkmarkc\checkmarki\checkmarkr\checkmarkc\checkmark$\checkmark \checkmark$\checkmark^\checkmark\\checkmarkc\checkmarki\checkmarkr\checkmarkc\checkmark$\checkmarka\checkmark$\checkmark^\checkmark\\checkmarkc\checkmarki\checkmarkr\checkmarkc\checkmark$\checkmarkv\checkmark$\checkmark^\checkmark\\checkmarkc\checkmarki\checkmarkr\checkmarkc\checkmark$\checkmarke\checkmark$\checkmark^\checkmark\\checkmarkc\checkmarki\checkmarkr\checkmarkc\checkmark$\checkmarkc\checkmark$\checkmark^\checkmark\\checkmarkc\checkmarki\checkmarkr\checkmarkc\checkmark$\checkmark \checkmark$\checkmark^\checkmark\\checkmarkc\checkmarki\checkmarkr\checkmarkc\checkmark$\checkmarkr\checkmark$\checkmark^\checkmark\\checkmarkc\checkmarki\checkmarkr\checkmarkc\checkmark$\checkmarké\checkmark$\checkmark^\checkmark\\checkmarkc\checkmarki\checkmarkr\checkmarkc\checkmark$\checkmarkd\checkmark$\checkmark^\checkmark\\checkmarkc\checkmarki\checkmarkr\checkmarkc\checkmark$\checkmarku\checkmark$\checkmark^\checkmark\\checkmarkc\checkmarki\checkmarkr\checkmarkc\checkmark$\checkmarkc\checkmark$\checkmark^\checkmark\\checkmarkc\checkmarki\checkmarkr\checkmarkc\checkmark$\checkmarkt\checkmark$\checkmark^\checkmark\\checkmarkc\checkmarki\checkmarkr\checkmarkc\checkmark$\checkmarki\checkmark$\checkmark^\checkmark\\checkmarkc\checkmarki\checkmarkr\checkmarkc\checkmark$\checkmarko\checkmark$\checkmark^\checkmark\\checkmarkc\checkmarki\checkmarkr\checkmarkc\checkmark$\checkmarkn\checkmark$\checkmark^\checkmark\\checkmarkc\checkmarki\checkmarkr\checkmarkc\checkmark$\checkmark \checkmark$\checkmark^\checkmark\\checkmarkc\checkmarki\checkmarkr\checkmarkc\checkmark$\checkmark1\checkmark$\checkmark^\checkmark\\checkmarkc\checkmarki\checkmarkr\checkmarkc\checkmark$\checkmark:\checkmark$\checkmark^\checkmark\\checkmarkc\checkmarki\checkmarkr\checkmarkc\checkmark$\checkmark5\checkmark$\checkmark^\checkmark\\checkmarkc\checkmarki\checkmarkr\checkmarkc\checkmark$\checkmark0\checkmark$\checkmark^\checkmark\\checkmarkc\checkmarki\checkmarkr\checkmarkc\checkmark$\checkmark}\checkmark$\checkmark^\checkmark\\checkmarkc\checkmarki\checkmarkr\checkmarkc\checkmark$\checkmark
\checkmark$\checkmark^\checkmark\\checkmarkc\checkmarki\checkmarkr\checkmarkc\checkmark$\checkmarkL\checkmark$\checkmark^\checkmark\\checkmarkc\checkmarki\checkmarkr\checkmarkc\checkmark$\checkmarke\checkmark$\checkmark^\checkmark\\checkmarkc\checkmarki\checkmarkr\checkmarkc\checkmark$\checkmark \checkmark$\checkmark^\checkmark\\checkmarkc\checkmarki\checkmarkr\checkmarkc\checkmark$\checkmarkm\checkmark$\checkmark^\checkmark\\checkmarkc\checkmarki\checkmarkr\checkmarkc\checkmark$\checkmarko\checkmark$\checkmark^\checkmark\\checkmarkc\checkmarki\checkmarkr\checkmarkc\checkmark$\checkmarkt\checkmark$\checkmark^\checkmark\\checkmarkc\checkmarki\checkmarkr\checkmarkc\checkmark$\checkmarke\checkmark$\checkmark^\checkmark\\checkmarkc\checkmarki\checkmarkr\checkmarkc\checkmark$\checkmarku\checkmark$\checkmark^\checkmark\\checkmarkc\checkmarki\checkmarkr\checkmarkc\checkmark$\checkmarkr\checkmark$\checkmark^\checkmark\\checkmarkc\checkmarki\checkmarkr\checkmarkc\checkmark$\checkmark \checkmark$\checkmark^\checkmark\\checkmarkc\checkmarki\checkmarkr\checkmarkc\checkmark$\checkmarkp\checkmark$\checkmark^\checkmark\\checkmarkc\checkmarki\checkmarkr\checkmarkc\checkmark$\checkmarka\checkmark$\checkmark^\checkmark\\checkmarkc\checkmarki\checkmarkr\checkmarkc\checkmark$\checkmarks\checkmark$\checkmark^\checkmark\\checkmarkc\checkmarki\checkmarkr\checkmarkc\checkmark$\checkmark \checkmark$\checkmark^\checkmark\\checkmarkc\checkmarki\checkmarkr\checkmarkc\checkmark$\checkmarkà\checkmark$\checkmark^\checkmark\\checkmarkc\checkmarki\checkmarkr\checkmarkc\checkmark$\checkmark \checkmark$\checkmark^\checkmark\\checkmarkc\checkmarki\checkmarkr\checkmarkc\checkmark$\checkmarkp\checkmark$\checkmark^\checkmark\\checkmarkc\checkmarki\checkmarkr\checkmarkc\checkmark$\checkmarka\checkmark$\checkmark^\checkmark\\checkmarkc\checkmarki\checkmarkr\checkmarkc\checkmark$\checkmarks\checkmark$\checkmark^\checkmark\\checkmarkc\checkmarki\checkmarkr\checkmarkc\checkmark$\checkmark \checkmark$\checkmark^\checkmark\\checkmarkc\checkmarki\checkmarkr\checkmarkc\checkmark$\checkmarkp\checkmark$\checkmark^\checkmark\\checkmarkc\checkmarki\checkmarkr\checkmarkc\checkmark$\checkmarko\checkmark$\checkmark^\checkmark\\checkmarkc\checkmarki\checkmarkr\checkmarkc\checkmark$\checkmarks\checkmark$\checkmark^\checkmark\\checkmarkc\checkmarki\checkmarkr\checkmarkc\checkmark$\checkmarks\checkmark$\checkmark^\checkmark\\checkmarkc\checkmarki\checkmarkr\checkmarkc\checkmark$\checkmarkè\checkmark$\checkmark^\checkmark\\checkmarkc\checkmarki\checkmarkr\checkmarkc\checkmark$\checkmarkd\checkmark$\checkmark^\checkmark\\checkmarkc\checkmarki\checkmarkr\checkmarkc\checkmark$\checkmarke\checkmark$\checkmark^\checkmark\\checkmarkc\checkmarki\checkmarkr\checkmarkc\checkmark$\checkmark \checkmark$\checkmark^\checkmark\\checkmarkc\checkmarki\checkmarkr\checkmarkc\checkmark$\checkmarku\checkmark$\checkmark^\checkmark\\checkmarkc\checkmarki\checkmarkr\checkmarkc\checkmark$\checkmarkn\checkmark$\checkmark^\checkmark\\checkmarkc\checkmarki\checkmarkr\checkmarkc\checkmark$\checkmark \checkmark$\checkmark^\checkmark\\checkmarkc\checkmarki\checkmarkr\checkmarkc\checkmark$\checkmarkp\checkmark$\checkmark^\checkmark\\checkmarkc\checkmarki\checkmarkr\checkmarkc\checkmark$\checkmarka\checkmark$\checkmark^\checkmark\\checkmarkc\checkmarki\checkmarkr\checkmarkc\checkmark$\checkmarks\checkmark$\checkmark^\checkmark\\checkmarkc\checkmarki\checkmarkr\checkmarkc\checkmark$\checkmark \checkmark$\checkmark^\checkmark\\checkmarkc\checkmarki\checkmarkr\checkmarkc\checkmark$\checkmarka\checkmark$\checkmark^\checkmark\\checkmarkc\checkmarki\checkmarkr\checkmarkc\checkmark$\checkmarkn\checkmark$\checkmark^\checkmark\\checkmarkc\checkmarki\checkmarkr\checkmarkc\checkmark$\checkmarkg\checkmark$\checkmark^\checkmark\\checkmarkc\checkmarki\checkmarkr\checkmarkc\checkmark$\checkmarku\checkmark$\checkmark^\checkmark\\checkmarkc\checkmarki\checkmarkr\checkmarkc\checkmark$\checkmarkl\checkmark$\checkmark^\checkmark\\checkmarkc\checkmarki\checkmarkr\checkmarkc\checkmark$\checkmarka\checkmark$\checkmark^\checkmark\\checkmarkc\checkmarki\checkmarkr\checkmarkc\checkmark$\checkmarki\checkmark$\checkmark^\checkmark\\checkmarkc\checkmarki\checkmarkr\checkmarkc\checkmark$\checkmarkr\checkmark$\checkmark^\checkmark\\checkmarkc\checkmarki\checkmarkr\checkmarkc\checkmark$\checkmarke\checkmark$\checkmark^\checkmark\\checkmarkc\checkmarki\checkmarkr\checkmarkc\checkmark$\checkmark \checkmark$\checkmark^\checkmark\\checkmarkc\checkmarki\checkmarkr\checkmarkc\checkmark$\checkmarkd\checkmark$\checkmark^\checkmark\\checkmarkc\checkmarki\checkmarkr\checkmarkc\checkmark$\checkmarke\checkmark$\checkmark^\checkmark\\checkmarkc\checkmarki\checkmarkr\checkmarkc\checkmark$\checkmark \checkmark$\checkmark^\checkmark\\checkmarkc\checkmarki\checkmarkr\checkmarkc\checkmark$\checkmark$\checkmark$\checkmark^\checkmark\\checkmarkc\checkmarki\checkmarkr\checkmarkc\checkmark$\checkmark1\checkmark$\checkmark^\checkmark\\checkmarkc\checkmarki\checkmarkr\checkmarkc\checkmark$\checkmark{\checkmark$\checkmark^\checkmark\\checkmarkc\checkmarki\checkmarkr\checkmarkc\checkmark$\checkmark,\checkmark$\checkmark^\checkmark\\checkmarkc\checkmarki\checkmarkr\checkmarkc\checkmark$\checkmark}\checkmark$\checkmark^\checkmark\\checkmarkc\checkmarki\checkmarkr\checkmarkc\checkmark$\checkmark8\checkmark$\checkmark^\checkmark\\checkmarkc\checkmarki\checkmarkr\checkmarkc\checkmark$\checkmark°\checkmark$\checkmark^\checkmark\\checkmarkc\checkmarki\checkmarkr\checkmarkc\checkmark$\checkmark$\checkmark$\checkmark^\checkmark\\checkmarkc\checkmarki\checkmarkr\checkmarkc\checkmark$\checkmark/\checkmark$\checkmark^\checkmark\\checkmarkc\checkmarki\checkmarkr\checkmarkc\checkmark$\checkmarkp\checkmark$\checkmark^\checkmark\\checkmarkc\checkmarki\checkmarkr\checkmarkc\checkmark$\checkmarka\checkmark$\checkmark^\checkmark\\checkmarkc\checkmarki\checkmarkr\checkmarkc\checkmark$\checkmarks\checkmark$\checkmark^\checkmark\\checkmarkc\checkmarki\checkmarkr\checkmarkc\checkmark$\checkmark.\checkmark$\checkmark^\checkmark\\checkmarkc\checkmarki\checkmarkr\checkmarkc\checkmark$\checkmark \checkmark$\checkmark^\checkmark\\checkmarkc\checkmarki\checkmarkr\checkmarkc\checkmark$\checkmarkA\checkmark$\checkmark^\checkmark\\checkmarkc\checkmarki\checkmarkr\checkmarkc\checkmark$\checkmarkv\checkmark$\checkmark^\checkmark\\checkmarkc\checkmarki\checkmarkr\checkmarkc\checkmark$\checkmarke\checkmark$\checkmark^\checkmark\\checkmarkc\checkmarki\checkmarkr\checkmarkc\checkmark$\checkmarkc\checkmark$\checkmark^\checkmark\\checkmarkc\checkmarki\checkmarkr\checkmarkc\checkmark$\checkmark \checkmark$\checkmark^\checkmark\\checkmarkc\checkmarki\checkmarkr\checkmarkc\checkmark$\checkmarkl\checkmark$\checkmark^\checkmark\\checkmarkc\checkmarki\checkmarkr\checkmarkc\checkmark$\checkmarke\checkmark$\checkmark^\checkmark\\checkmarkc\checkmarki\checkmarkr\checkmarkc\checkmark$\checkmark \checkmark$\checkmark^\checkmark\\checkmarkc\checkmarki\checkmarkr\checkmarkc\checkmark$\checkmarkm\checkmark$\checkmark^\checkmark\\checkmarkc\checkmarki\checkmarkr\checkmarkc\checkmark$\checkmarki\checkmark$\checkmark^\checkmark\\checkmarkc\checkmarki\checkmarkr\checkmarkc\checkmark$\checkmarkc\checkmark$\checkmark^\checkmark\\checkmarkc\checkmarki\checkmarkr\checkmarkc\checkmark$\checkmarkr\checkmark$\checkmark^\checkmark\\checkmarkc\checkmarki\checkmarkr\checkmarkc\checkmark$\checkmarko\checkmark$\checkmark^\checkmark\\checkmarkc\checkmarki\checkmarkr\checkmarkc\checkmark$\checkmarks\checkmark$\checkmark^\checkmark\\checkmarkc\checkmarki\checkmarkr\checkmarkc\checkmark$\checkmarkt\checkmark$\checkmark^\checkmark\\checkmarkc\checkmarki\checkmarkr\checkmarkc\checkmark$\checkmarke\checkmark$\checkmark^\checkmark\\checkmarkc\checkmarki\checkmarkr\checkmarkc\checkmark$\checkmarkp\checkmark$\checkmark^\checkmark\\checkmarkc\checkmarki\checkmarkr\checkmarkc\checkmark$\checkmarkp\checkmark$\checkmark^\checkmark\\checkmarkc\checkmarki\checkmarkr\checkmarkc\checkmark$\checkmarki\checkmark$\checkmark^\checkmark\\checkmarkc\checkmarki\checkmarkr\checkmarkc\checkmark$\checkmarkn\checkmark$\checkmark^\checkmark\\checkmarkc\checkmarki\checkmarkr\checkmarkc\checkmark$\checkmarkg\checkmark$\checkmark^\checkmark\\checkmarkc\checkmarki\checkmarkr\checkmarkc\checkmark$\checkmark \checkmark$\checkmark^\checkmark\\checkmarkc\checkmarki\checkmarkr\checkmarkc\checkmark$\checkmark1\checkmark$\checkmark^\checkmark\\checkmarkc\checkmarki\checkmarkr\checkmarkc\checkmark$\checkmark/\checkmark$\checkmark^\checkmark\\checkmarkc\checkmarki\checkmarkr\checkmarkc\checkmark$\checkmark1\checkmark$\checkmark^\checkmark\\checkmarkc\checkmarki\checkmarkr\checkmarkc\checkmark$\checkmark6\checkmark$\checkmark^\checkmark\\checkmarkc\checkmarki\checkmarkr\checkmarkc\checkmark$\checkmark \checkmark$\checkmark^\checkmark\\checkmarkc\checkmarki\checkmarkr\checkmarkc\checkmark$\checkmarke\checkmark$\checkmark^\checkmark\\checkmarkc\checkmarki\checkmarkr\checkmarkc\checkmark$\checkmarkt\checkmark$\checkmark^\checkmark\\checkmarkc\checkmarki\checkmarkr\checkmarkc\checkmark$\checkmark \checkmark$\checkmark^\checkmark\\checkmarkc\checkmarki\checkmarkr\checkmarkc\checkmark$\checkmarkl\checkmark$\checkmark^\checkmark\\checkmarkc\checkmarki\checkmarkr\checkmarkc\checkmark$\checkmarka\checkmark$\checkmark^\checkmark\\checkmarkc\checkmarki\checkmarkr\checkmarkc\checkmark$\checkmark \checkmark$\checkmark^\checkmark\\checkmarkc\checkmarki\checkmarkr\checkmarkc\checkmark$\checkmarkr\checkmark$\checkmark^\checkmark\\checkmarkc\checkmarki\checkmarkr\checkmarkc\checkmark$\checkmarké\checkmark$\checkmark^\checkmark\\checkmarkc\checkmarki\checkmarkr\checkmarkc\checkmark$\checkmarkd\checkmark$\checkmark^\checkmark\\checkmarkc\checkmarki\checkmarkr\checkmarkc\checkmark$\checkmarku\checkmark$\checkmark^\checkmark\\checkmarkc\checkmarki\checkmarkr\checkmarkc\checkmark$\checkmarkc\checkmark$\checkmark^\checkmark\\checkmarkc\checkmarki\checkmarkr\checkmarkc\checkmark$\checkmarkt\checkmark$\checkmark^\checkmark\\checkmarkc\checkmarki\checkmarkr\checkmarkc\checkmark$\checkmarki\checkmark$\checkmark^\checkmark\\checkmarkc\checkmarki\checkmarkr\checkmarkc\checkmark$\checkmarko\checkmark$\checkmark^\checkmark\\checkmarkc\checkmarki\checkmarkr\checkmarkc\checkmark$\checkmarkn\checkmark$\checkmark^\checkmark\\checkmarkc\checkmarki\checkmarkr\checkmarkc\checkmark$\checkmark \checkmark$\checkmark^\checkmark\\checkmarkc\checkmarki\checkmarkr\checkmarkc\checkmark$\checkmarkd\checkmark$\checkmark^\checkmark\\checkmarkc\checkmarki\checkmarkr\checkmarkc\checkmark$\checkmark'\checkmark$\checkmark^\checkmark\\checkmarkc\checkmarki\checkmarkr\checkmarkc\checkmark$\checkmarke\checkmark$\checkmark^\checkmark\\checkmarkc\checkmarki\checkmarkr\checkmarkc\checkmark$\checkmarkn\checkmark$\checkmark^\checkmark\\checkmarkc\checkmarki\checkmarkr\checkmarkc\checkmark$\checkmarkg\checkmark$\checkmark^\checkmark\\checkmarkc\checkmarki\checkmarkr\checkmarkc\checkmark$\checkmarkr\checkmark$\checkmark^\checkmark\\checkmarkc\checkmarki\checkmarkr\checkmarkc\checkmark$\checkmarke\checkmark$\checkmark^\checkmark\\checkmarkc\checkmarki\checkmarkr\checkmarkc\checkmark$\checkmarkn\checkmark$\checkmark^\checkmark\\checkmarkc\checkmarki\checkmarkr\checkmarkc\checkmark$\checkmarka\checkmark$\checkmark^\checkmark\\checkmarkc\checkmarki\checkmarkr\checkmarkc\checkmark$\checkmarkg\checkmark$\checkmark^\checkmark\\checkmarkc\checkmarki\checkmarkr\checkmarkc\checkmark$\checkmarke\checkmark$\checkmark^\checkmark\\checkmarkc\checkmarki\checkmarkr\checkmarkc\checkmark$\checkmark \checkmark$\checkmark^\checkmark\\checkmarkc\checkmarki\checkmarkr\checkmarkc\checkmark$\checkmark$\checkmark$\checkmark^\checkmark\\checkmarkc\checkmarki\checkmarkr\checkmarkc\checkmark$\checkmarki\checkmark$\checkmark^\checkmark\\checkmarkc\checkmarki\checkmarkr\checkmarkc\checkmark$\checkmark \checkmark$\checkmark^\checkmark\\checkmarkc\checkmarki\checkmarkr\checkmarkc\checkmark$\checkmark=\checkmark$\checkmark^\checkmark\\checkmarkc\checkmarki\checkmarkr\checkmarkc\checkmark$\checkmark \checkmark$\checkmark^\checkmark\\checkmarkc\checkmarki\checkmarkr\checkmarkc\checkmark$\checkmark1\checkmark$\checkmark^\checkmark\\checkmarkc\checkmarki\checkmarkr\checkmarkc\checkmark$\checkmark:\checkmark$\checkmark^\checkmark\\checkmarkc\checkmarki\checkmarkr\checkmarkc\checkmark$\checkmark5\checkmark$\checkmark^\checkmark\\checkmarkc\checkmarki\checkmarkr\checkmarkc\checkmark$\checkmark0\checkmark$\checkmark^\checkmark\\checkmarkc\checkmarki\checkmarkr\checkmarkc\checkmark$\checkmark$\checkmark$\checkmark^\checkmark\\checkmarkc\checkmarki\checkmarkr\checkmarkc\checkmark$\checkmark \checkmark$\checkmark^\checkmark\\checkmarkc\checkmarki\checkmarkr\checkmarkc\checkmark$\checkmark:\checkmark$\checkmark^\checkmark\\checkmarkc\checkmarki\checkmarkr\checkmarkc\checkmark$\checkmark
\checkmark$\checkmark^\checkmark\\checkmarkc\checkmarki\checkmarkr\checkmarkc\checkmark$\checkmark\\checkmark$\checkmark^\checkmark\\checkmarkc\checkmarki\checkmarkr\checkmarkc\checkmark$\checkmarkb\checkmark$\checkmark^\checkmark\\checkmarkc\checkmarki\checkmarkr\checkmarkc\checkmark$\checkmarke\checkmark$\checkmark^\checkmark\\checkmarkc\checkmarki\checkmarkr\checkmarkc\checkmark$\checkmarkg\checkmark$\checkmark^\checkmark\\checkmarkc\checkmarki\checkmarkr\checkmarkc\checkmark$\checkmarki\checkmark$\checkmark^\checkmark\\checkmarkc\checkmarki\checkmarkr\checkmarkc\checkmark$\checkmarkn\checkmark$\checkmark^\checkmark\\checkmarkc\checkmarki\checkmarkr\checkmarkc\checkmark$\checkmark{\checkmark$\checkmark^\checkmark\\checkmarkc\checkmarki\checkmarkr\checkmarkc\checkmark$\checkmarke\checkmark$\checkmark^\checkmark\\checkmarkc\checkmarki\checkmarkr\checkmarkc\checkmark$\checkmarkq\checkmark$\checkmark^\checkmark\\checkmarkc\checkmarki\checkmarkr\checkmarkc\checkmark$\checkmarku\checkmark$\checkmark^\checkmark\\checkmarkc\checkmarki\checkmarkr\checkmarkc\checkmark$\checkmarka\checkmark$\checkmark^\checkmark\\checkmarkc\checkmarki\checkmarkr\checkmarkc\checkmark$\checkmarkt\checkmark$\checkmark^\checkmark\\checkmarkc\checkmarki\checkmarkr\checkmarkc\checkmark$\checkmarki\checkmark$\checkmark^\checkmark\\checkmarkc\checkmarki\checkmarkr\checkmarkc\checkmark$\checkmarko\checkmark$\checkmark^\checkmark\\checkmarkc\checkmarki\checkmarkr\checkmarkc\checkmark$\checkmarkn\checkmark$\checkmark^\checkmark\\checkmarkc\checkmarki\checkmarkr\checkmarkc\checkmark$\checkmark}\checkmark$\checkmark^\checkmark\\checkmarkc\checkmarki\checkmarkr\checkmarkc\checkmark$\checkmark
\checkmark$\checkmark^\checkmark\\checkmarkc\checkmarki\checkmarkr\checkmarkc\checkmark$\checkmark \checkmark$\checkmark^\checkmark\\checkmarkc\checkmarki\checkmarkr\checkmarkc\checkmark$\checkmark \checkmark$\checkmark^\checkmark\\checkmarkc\checkmarki\checkmarkr\checkmarkc\checkmark$\checkmark \checkmark$\checkmark^\checkmark\\checkmarkc\checkmarki\checkmarkr\checkmarkc\checkmark$\checkmark \checkmark$\checkmark^\checkmark\\checkmarkc\checkmarki\checkmarkr\checkmarkc\checkmark$\checkmark\\checkmark$\checkmark^\checkmark\\checkmarkc\checkmarki\checkmarkr\checkmarkc\checkmark$\checkmarkt\checkmark$\checkmark^\checkmark\\checkmarkc\checkmarki\checkmarkr\checkmarkc\checkmark$\checkmarkh\checkmark$\checkmark^\checkmark\\checkmarkc\checkmarki\checkmarkr\checkmarkc\checkmark$\checkmarke\checkmark$\checkmark^\checkmark\\checkmarkc\checkmarki\checkmarkr\checkmarkc\checkmark$\checkmarkt\checkmark$\checkmark^\checkmark\\checkmarkc\checkmarki\checkmarkr\checkmarkc\checkmark$\checkmarka\checkmark$\checkmark^\checkmark\\checkmarkc\checkmarki\checkmarkr\checkmarkc\checkmark$\checkmark_\checkmark$\checkmark^\checkmark\\checkmarkc\checkmarki\checkmarkr\checkmarkc\checkmark$\checkmark{\checkmark$\checkmark^\checkmark\\checkmarkc\checkmarki\checkmarkr\checkmarkc\checkmark$\checkmarki\checkmark$\checkmark^\checkmark\\checkmarkc\checkmarki\checkmarkr\checkmarkc\checkmark$\checkmarkm\checkmark$\checkmark^\checkmark\\checkmarkc\checkmarki\checkmarkr\checkmarkc\checkmark$\checkmarkp\checkmark$\checkmark^\checkmark\\checkmarkc\checkmarki\checkmarkr\checkmarkc\checkmark$\checkmark}\checkmark$\checkmark^\checkmark\\checkmarkc\checkmarki\checkmarkr\checkmarkc\checkmark$\checkmark \checkmark$\checkmark^\checkmark\\checkmarkc\checkmarki\checkmarkr\checkmarkc\checkmark$\checkmark=\checkmark$\checkmark^\checkmark\\checkmarkc\checkmarki\checkmarkr\checkmarkc\checkmark$\checkmark \checkmark$\checkmark^\checkmark\\checkmarkc\checkmarki\checkmarkr\checkmarkc\checkmark$\checkmark\\checkmark$\checkmark^\checkmark\\checkmarkc\checkmarki\checkmarkr\checkmarkc\checkmark$\checkmarkf\checkmark$\checkmark^\checkmark\\checkmarkc\checkmarki\checkmarkr\checkmarkc\checkmark$\checkmarkr\checkmark$\checkmark^\checkmark\\checkmarkc\checkmarki\checkmarkr\checkmarkc\checkmark$\checkmarka\checkmark$\checkmark^\checkmark\\checkmarkc\checkmarki\checkmarkr\checkmarkc\checkmark$\checkmarkc\checkmark$\checkmark^\checkmark\\checkmarkc\checkmarki\checkmarkr\checkmarkc\checkmark$\checkmark{\checkmark$\checkmark^\checkmark\\checkmarkc\checkmarki\checkmarkr\checkmarkc\checkmark$\checkmark1\checkmark$\checkmark^\checkmark\\checkmarkc\checkmarki\checkmarkr\checkmarkc\checkmark$\checkmark{\checkmark$\checkmark^\checkmark\\checkmarkc\checkmarki\checkmarkr\checkmarkc\checkmark$\checkmark,\checkmark$\checkmark^\checkmark\\checkmarkc\checkmarki\checkmarkr\checkmarkc\checkmark$\checkmark}\checkmark$\checkmark^\checkmark\\checkmarkc\checkmarki\checkmarkr\checkmarkc\checkmark$\checkmark8\checkmark$\checkmark^\checkmark\\checkmarkc\checkmarki\checkmarkr\checkmarkc\checkmark$\checkmark°\checkmark$\checkmark^\checkmark\\checkmarkc\checkmarki\checkmarkr\checkmarkc\checkmark$\checkmark}\checkmark$\checkmark^\checkmark\\checkmarkc\checkmarki\checkmarkr\checkmarkc\checkmark$\checkmark{\checkmark$\checkmark^\checkmark\\checkmarkc\checkmarki\checkmarkr\checkmarkc\checkmark$\checkmark1\checkmark$\checkmark^\checkmark\\checkmarkc\checkmarki\checkmarkr\checkmarkc\checkmark$\checkmark6\checkmark$\checkmark^\checkmark\\checkmarkc\checkmarki\checkmarkr\checkmarkc\checkmark$\checkmark \checkmark$\checkmark^\checkmark\\checkmarkc\checkmarki\checkmarkr\checkmarkc\checkmark$\checkmark\\checkmark$\checkmark^\checkmark\\checkmarkc\checkmarki\checkmarkr\checkmarkc\checkmark$\checkmarkt\checkmark$\checkmark^\checkmark\\checkmarkc\checkmarki\checkmarkr\checkmarkc\checkmark$\checkmarki\checkmark$\checkmark^\checkmark\\checkmarkc\checkmarki\checkmarkr\checkmarkc\checkmark$\checkmarkm\checkmark$\checkmark^\checkmark\\checkmarkc\checkmarki\checkmarkr\checkmarkc\checkmark$\checkmarke\checkmark$\checkmark^\checkmark\\checkmarkc\checkmarki\checkmarkr\checkmarkc\checkmark$\checkmarks\checkmark$\checkmark^\checkmark\\checkmarkc\checkmarki\checkmarkr\checkmarkc\checkmark$\checkmark \checkmark$\checkmark^\checkmark\\checkmarkc\checkmarki\checkmarkr\checkmarkc\checkmark$\checkmark5\checkmark$\checkmark^\checkmark\\checkmarkc\checkmarki\checkmarkr\checkmarkc\checkmark$\checkmark0\checkmark$\checkmark^\checkmark\\checkmarkc\checkmarki\checkmarkr\checkmarkc\checkmark$\checkmark}\checkmark$\checkmark^\checkmark\\checkmarkc\checkmarki\checkmarkr\checkmarkc\checkmark$\checkmark \checkmark$\checkmark^\checkmark\\checkmarkc\checkmarki\checkmarkr\checkmarkc\checkmark$\checkmark=\checkmark$\checkmark^\checkmark\\checkmarkc\checkmarki\checkmarkr\checkmarkc\checkmark$\checkmark \checkmark$\checkmark^\checkmark\\checkmarkc\checkmarki\checkmarkr\checkmarkc\checkmark$\checkmark\\checkmark$\checkmark^\checkmark\\checkmarkc\checkmarki\checkmarkr\checkmarkc\checkmark$\checkmarkb\checkmark$\checkmark^\checkmark\\checkmarkc\checkmarki\checkmarkr\checkmarkc\checkmark$\checkmarko\checkmark$\checkmark^\checkmark\\checkmarkc\checkmarki\checkmarkr\checkmarkc\checkmark$\checkmarkx\checkmark$\checkmark^\checkmark\\checkmarkc\checkmarki\checkmarkr\checkmarkc\checkmark$\checkmarke\checkmark$\checkmark^\checkmark\\checkmarkc\checkmarki\checkmarkr\checkmarkc\checkmark$\checkmarkd\checkmark$\checkmark^\checkmark\\checkmarkc\checkmarki\checkmarkr\checkmarkc\checkmark$\checkmark{\checkmark$\checkmark^\checkmark\\checkmarkc\checkmarki\checkmarkr\checkmarkc\checkmark$\checkmark0\checkmark$\checkmark^\checkmark\\checkmarkc\checkmarki\checkmarkr\checkmarkc\checkmark$\checkmark{\checkmark$\checkmark^\checkmark\\checkmarkc\checkmarki\checkmarkr\checkmarkc\checkmark$\checkmark,\checkmark$\checkmark^\checkmark\\checkmarkc\checkmarki\checkmarkr\checkmarkc\checkmark$\checkmark}\checkmark$\checkmark^\checkmark\\checkmarkc\checkmarki\checkmarkr\checkmarkc\checkmark$\checkmark0\checkmark$\checkmark^\checkmark\\checkmarkc\checkmarki\checkmarkr\checkmarkc\checkmark$\checkmark0\checkmark$\checkmark^\checkmark\\checkmarkc\checkmarki\checkmarkr\checkmarkc\checkmark$\checkmark2\checkmark$\checkmark^\checkmark\\checkmarkc\checkmarki\checkmarkr\checkmarkc\checkmark$\checkmark2\checkmark$\checkmark^\checkmark\\checkmarkc\checkmarki\checkmarkr\checkmarkc\checkmark$\checkmark5\checkmark$\checkmark^\checkmark\\checkmarkc\checkmarki\checkmarkr\checkmarkc\checkmark$\checkmark°\checkmark$\checkmark^\checkmark\\checkmarkc\checkmarki\checkmarkr\checkmarkc\checkmark$\checkmark\\checkmark$\checkmark^\checkmark\\checkmarkc\checkmarki\checkmarkr\checkmarkc\checkmark$\checkmarkt\checkmark$\checkmark^\checkmark\\checkmarkc\checkmarki\checkmarkr\checkmarkc\checkmark$\checkmarke\checkmark$\checkmark^\checkmark\\checkmarkc\checkmarki\checkmarkr\checkmarkc\checkmark$\checkmarkx\checkmark$\checkmark^\checkmark\\checkmarkc\checkmarki\checkmarkr\checkmarkc\checkmark$\checkmarkt\checkmark$\checkmark^\checkmark\\checkmarkc\checkmarki\checkmarkr\checkmarkc\checkmark$\checkmark{\checkmark$\checkmark^\checkmark\\checkmarkc\checkmarki\checkmarkr\checkmarkc\checkmark$\checkmark/\checkmark$\checkmark^\checkmark\\checkmarkc\checkmarki\checkmarkr\checkmarkc\checkmark$\checkmarki\checkmark$\checkmark^\checkmark\\checkmarkc\checkmarki\checkmarkr\checkmarkc\checkmark$\checkmarkm\checkmark$\checkmark^\checkmark\\checkmarkc\checkmarki\checkmarkr\checkmarkc\checkmark$\checkmarkp\checkmark$\checkmark^\checkmark\\checkmarkc\checkmarki\checkmarkr\checkmarkc\checkmark$\checkmarku\checkmark$\checkmark^\checkmark\\checkmarkc\checkmarki\checkmarkr\checkmarkc\checkmark$\checkmarkl\checkmark$\checkmark^\checkmark\\checkmarkc\checkmarki\checkmarkr\checkmarkc\checkmark$\checkmarks\checkmark$\checkmark^\checkmark\\checkmarkc\checkmarki\checkmarkr\checkmarkc\checkmark$\checkmarki\checkmark$\checkmark^\checkmark\\checkmarkc\checkmarki\checkmarkr\checkmarkc\checkmark$\checkmarko\checkmark$\checkmark^\checkmark\\checkmarkc\checkmarki\checkmarkr\checkmarkc\checkmark$\checkmarkn\checkmark$\checkmark^\checkmark\\checkmarkc\checkmarki\checkmarkr\checkmarkc\checkmark$\checkmark}\checkmark$\checkmark^\checkmark\\checkmarkc\checkmarki\checkmarkr\checkmarkc\checkmark$\checkmark}\checkmark$\checkmark^\checkmark\\checkmarkc\checkmarki\checkmarkr\checkmarkc\checkmark$\checkmark
\checkmark$\checkmark^\checkmark\\checkmarkc\checkmarki\checkmarkr\checkmarkc\checkmark$\checkmark\\checkmark$\checkmark^\checkmark\\checkmarkc\checkmarki\checkmarkr\checkmarkc\checkmark$\checkmarke\checkmark$\checkmark^\checkmark\\checkmarkc\checkmarki\checkmarkr\checkmarkc\checkmark$\checkmarkn\checkmark$\checkmark^\checkmark\\checkmarkc\checkmarki\checkmarkr\checkmarkc\checkmark$\checkmarkd\checkmark$\checkmark^\checkmark\\checkmarkc\checkmarki\checkmarkr\checkmarkc\checkmark$\checkmark{\checkmark$\checkmark^\checkmark\\checkmarkc\checkmarki\checkmarkr\checkmarkc\checkmark$\checkmarke\checkmark$\checkmark^\checkmark\\checkmarkc\checkmarki\checkmarkr\checkmarkc\checkmark$\checkmarkq\checkmark$\checkmark^\checkmark\\checkmarkc\checkmarki\checkmarkr\checkmarkc\checkmark$\checkmarku\checkmark$\checkmark^\checkmark\\checkmarkc\checkmarki\checkmarkr\checkmarkc\checkmark$\checkmarka\checkmark$\checkmark^\checkmark\\checkmarkc\checkmarki\checkmarkr\checkmarkc\checkmark$\checkmarkt\checkmark$\checkmark^\checkmark\\checkmarkc\checkmarki\checkmarkr\checkmarkc\checkmark$\checkmarki\checkmark$\checkmark^\checkmark\\checkmarkc\checkmarki\checkmarkr\checkmarkc\checkmark$\checkmarko\checkmark$\checkmark^\checkmark\\checkmarkc\checkmarki\checkmarkr\checkmarkc\checkmark$\checkmarkn\checkmark$\checkmark^\checkmark\\checkmarkc\checkmarki\checkmarkr\checkmarkc\checkmark$\checkmark}\checkmark$\checkmark^\checkmark\\checkmarkc\checkmarki\checkmarkr\checkmarkc\checkmark$\checkmark
\checkmark$\checkmark^\checkmark\\checkmarkc\checkmarki\checkmarkr\checkmarkc\checkmark$\checkmark
\checkmark$\checkmark^\checkmark\\checkmarkc\checkmarki\checkmarkr\checkmarkc\checkmark$\checkmarkL\checkmark$\checkmark^\checkmark\\checkmarkc\checkmarki\checkmarkr\checkmarkc\checkmark$\checkmark'\checkmark$\checkmark^\checkmark\\checkmarkc\checkmarki\checkmarkr\checkmarkc\checkmark$\checkmarke\checkmark$\checkmark^\checkmark\\checkmarkc\checkmarki\checkmarkr\checkmarkc\checkmark$\checkmarkr\checkmark$\checkmark^\checkmark\\checkmarkc\checkmarki\checkmarkr\checkmarkc\checkmark$\checkmarkr\checkmark$\checkmark^\checkmark\\checkmarkc\checkmarki\checkmarkr\checkmarkc\checkmark$\checkmarke\checkmark$\checkmark^\checkmark\\checkmarkc\checkmarki\checkmarkr\checkmarkc\checkmark$\checkmarku\checkmark$\checkmark^\checkmark\\checkmarkc\checkmarki\checkmarkr\checkmarkc\checkmark$\checkmarkr\checkmark$\checkmark^\checkmark\\checkmarkc\checkmarki\checkmarkr\checkmarkc\checkmark$\checkmark \checkmark$\checkmark^\checkmark\\checkmarkc\checkmarki\checkmarkr\checkmarkc\checkmark$\checkmarkl\checkmark$\checkmark^\checkmark\\checkmarkc\checkmarki\checkmarkr\checkmarkc\checkmark$\checkmarki\checkmark$\checkmark^\checkmark\\checkmarkc\checkmarki\checkmarkr\checkmarkc\checkmark$\checkmarkn\checkmark$\checkmark^\checkmark\\checkmarkc\checkmarki\checkmarkr\checkmarkc\checkmark$\checkmarké\checkmark$\checkmark^\checkmark\\checkmarkc\checkmarki\checkmarkr\checkmarkc\checkmark$\checkmarka\checkmark$\checkmark^\checkmark\\checkmarkc\checkmarki\checkmarkr\checkmarkc\checkmark$\checkmarki\checkmark$\checkmark^\checkmark\\checkmarkc\checkmarki\checkmarkr\checkmarkc\checkmark$\checkmarkr\checkmark$\checkmark^\checkmark\\checkmarkc\checkmarki\checkmarkr\checkmarkc\checkmark$\checkmarke\checkmark$\checkmark^\checkmark\\checkmarkc\checkmarki\checkmarkr\checkmarkc\checkmark$\checkmark \checkmark$\checkmark^\checkmark\\checkmarkc\checkmarki\checkmarkr\checkmarkc\checkmark$\checkmarkc\checkmark$\checkmark^\checkmark\\checkmarkc\checkmarki\checkmarkr\checkmarkc\checkmark$\checkmarko\checkmark$\checkmark^\checkmark\\checkmarkc\checkmarki\checkmarkr\checkmarkc\checkmark$\checkmarkr\checkmark$\checkmark^\checkmark\\checkmarkc\checkmarki\checkmarkr\checkmarkc\checkmark$\checkmarkr\checkmark$\checkmark^\checkmark\\checkmarkc\checkmarki\checkmarkr\checkmarkc\checkmark$\checkmarke\checkmark$\checkmark^\checkmark\\checkmarkc\checkmarki\checkmarkr\checkmarkc\checkmark$\checkmarks\checkmark$\checkmark^\checkmark\\checkmarkc\checkmarki\checkmarkr\checkmarkc\checkmark$\checkmarkp\checkmark$\checkmark^\checkmark\\checkmarkc\checkmarki\checkmarkr\checkmarkc\checkmark$\checkmarko\checkmark$\checkmark^\checkmark\\checkmarkc\checkmarki\checkmarkr\checkmarkc\checkmark$\checkmarkn\checkmark$\checkmark^\checkmark\\checkmarkc\checkmarki\checkmarkr\checkmarkc\checkmark$\checkmarkd\checkmark$\checkmark^\checkmark\\checkmarkc\checkmarki\checkmarkr\checkmarkc\checkmark$\checkmarka\checkmark$\checkmark^\checkmark\\checkmarkc\checkmarki\checkmarkr\checkmarkc\checkmark$\checkmarkn\checkmark$\checkmark^\checkmark\\checkmarkc\checkmarki\checkmarkr\checkmarkc\checkmark$\checkmarkt\checkmark$\checkmark^\checkmark\\checkmarkc\checkmarki\checkmarkr\checkmarkc\checkmark$\checkmarke\checkmark$\checkmark^\checkmark\\checkmarkc\checkmarki\checkmarkr\checkmarkc\checkmark$\checkmark \checkmark$\checkmark^\checkmark\\checkmarkc\checkmarki\checkmarkr\checkmarkc\checkmark$\checkmarke\checkmark$\checkmark^\checkmark\\checkmarkc\checkmarki\checkmarkr\checkmarkc\checkmark$\checkmarkn\checkmark$\checkmark^\checkmark\\checkmarkc\checkmarki\checkmarkr\checkmarkc\checkmark$\checkmark \checkmark$\checkmark^\checkmark\\checkmarkc\checkmarki\checkmarkr\checkmarkc\checkmark$\checkmarkb\checkmark$\checkmark^\checkmark\\checkmarkc\checkmarki\checkmarkr\checkmarkc\checkmark$\checkmarko\checkmark$\checkmark^\checkmark\\checkmarkc\checkmarki\checkmarkr\checkmarkc\checkmark$\checkmarku\checkmark$\checkmark^\checkmark\\checkmarkc\checkmarki\checkmarkr\checkmarkc\checkmark$\checkmarkt\checkmark$\checkmark^\checkmark\\checkmarkc\checkmarki\checkmarkr\checkmarkc\checkmark$\checkmark \checkmark$\checkmark^\checkmark\\checkmarkc\checkmarki\checkmarkr\checkmarkc\checkmark$\checkmarkd\checkmark$\checkmark^\checkmark\\checkmarkc\checkmarki\checkmarkr\checkmarkc\checkmark$\checkmarku\checkmark$\checkmark^\checkmark\\checkmarkc\checkmarki\checkmarkr\checkmarkc\checkmark$\checkmark \checkmark$\checkmark^\checkmark\\checkmarkc\checkmarki\checkmarkr\checkmarkc\checkmark$\checkmarkp\checkmark$\checkmark^\checkmark\\checkmarkc\checkmarki\checkmarkr\checkmarkc\checkmark$\checkmarkl\checkmark$\checkmark^\checkmark\\checkmarkc\checkmarki\checkmarkr\checkmarkc\checkmark$\checkmarka\checkmark$\checkmark^\checkmark\\checkmarkc\checkmarki\checkmarkr\checkmarkc\checkmark$\checkmarkt\checkmark$\checkmark^\checkmark\\checkmarkc\checkmarki\checkmarkr\checkmarkc\checkmark$\checkmarke\checkmark$\checkmark^\checkmark\\checkmarkc\checkmarki\checkmarkr\checkmarkc\checkmark$\checkmarka\checkmark$\checkmark^\checkmark\\checkmarkc\checkmarki\checkmarkr\checkmarkc\checkmark$\checkmarku\checkmark$\checkmark^\checkmark\\checkmarkc\checkmarki\checkmarkr\checkmarkc\checkmark$\checkmark \checkmark$\checkmark^\checkmark\\checkmarkc\checkmarki\checkmarkr\checkmarkc\checkmark$\checkmark:\checkmark$\checkmark^\checkmark\\checkmarkc\checkmarki\checkmarkr\checkmarkc\checkmark$\checkmark
\checkmark$\checkmark^\checkmark\\checkmarkc\checkmarki\checkmarkr\checkmarkc\checkmark$\checkmark\\checkmark$\checkmark^\checkmark\\checkmarkc\checkmarki\checkmarkr\checkmarkc\checkmark$\checkmarkb\checkmark$\checkmark^\checkmark\\checkmarkc\checkmarki\checkmarkr\checkmarkc\checkmark$\checkmarke\checkmark$\checkmark^\checkmark\\checkmarkc\checkmarki\checkmarkr\checkmarkc\checkmark$\checkmarkg\checkmark$\checkmark^\checkmark\\checkmarkc\checkmarki\checkmarkr\checkmarkc\checkmark$\checkmarki\checkmark$\checkmark^\checkmark\\checkmarkc\checkmarki\checkmarkr\checkmarkc\checkmark$\checkmarkn\checkmark$\checkmark^\checkmark\\checkmarkc\checkmarki\checkmarkr\checkmarkc\checkmark$\checkmark{\checkmark$\checkmark^\checkmark\\checkmarkc\checkmarki\checkmarkr\checkmarkc\checkmark$\checkmarke\checkmark$\checkmark^\checkmark\\checkmarkc\checkmarki\checkmarkr\checkmarkc\checkmark$\checkmarkq\checkmark$\checkmark^\checkmark\\checkmarkc\checkmarki\checkmarkr\checkmarkc\checkmark$\checkmarku\checkmark$\checkmark^\checkmark\\checkmarkc\checkmarki\checkmarkr\checkmarkc\checkmark$\checkmarka\checkmark$\checkmark^\checkmark\\checkmarkc\checkmarki\checkmarkr\checkmarkc\checkmark$\checkmarkt\checkmark$\checkmark^\checkmark\\checkmarkc\checkmarki\checkmarkr\checkmarkc\checkmark$\checkmarki\checkmark$\checkmark^\checkmark\\checkmarkc\checkmarki\checkmarkr\checkmarkc\checkmark$\checkmarko\checkmark$\checkmark^\checkmark\\checkmarkc\checkmarki\checkmarkr\checkmarkc\checkmark$\checkmarkn\checkmark$\checkmark^\checkmark\\checkmarkc\checkmarki\checkmarkr\checkmarkc\checkmark$\checkmark}\checkmark$\checkmark^\checkmark\\checkmarkc\checkmarki\checkmarkr\checkmarkc\checkmark$\checkmark
\checkmark$\checkmark^\checkmark\\checkmarkc\checkmarki\checkmarkr\checkmarkc\checkmark$\checkmark \checkmark$\checkmark^\checkmark\\checkmarkc\checkmarki\checkmarkr\checkmarkc\checkmark$\checkmark \checkmark$\checkmark^\checkmark\\checkmarkc\checkmarki\checkmarkr\checkmarkc\checkmark$\checkmark \checkmark$\checkmark^\checkmark\\checkmarkc\checkmarki\checkmarkr\checkmarkc\checkmark$\checkmark \checkmark$\checkmark^\checkmark\\checkmarkc\checkmarki\checkmarkr\checkmarkc\checkmark$\checkmark\\checkmark$\checkmark^\checkmark\\checkmarkc\checkmarki\checkmarkr\checkmarkc\checkmark$\checkmarkD\checkmark$\checkmark^\checkmark\\checkmarkc\checkmarki\checkmarkr\checkmarkc\checkmark$\checkmarke\checkmark$\checkmark^\checkmark\\checkmarkc\checkmarki\checkmarkr\checkmarkc\checkmark$\checkmarkl\checkmark$\checkmark^\checkmark\\checkmarkc\checkmarki\checkmarkr\checkmarkc\checkmark$\checkmarkt\checkmark$\checkmark^\checkmark\\checkmarkc\checkmarki\checkmarkr\checkmarkc\checkmark$\checkmarka\checkmark$\checkmark^\checkmark\\checkmarkc\checkmarki\checkmarkr\checkmarkc\checkmark$\checkmark \checkmark$\checkmark^\checkmark\\checkmarkc\checkmarki\checkmarkr\checkmarkc\checkmark$\checkmarkL\checkmark$\checkmark^\checkmark\\checkmarkc\checkmarki\checkmarkr\checkmarkc\checkmark$\checkmark_\checkmark$\checkmark^\checkmark\\checkmarkc\checkmarki\checkmarkr\checkmarkc\checkmark$\checkmark{\checkmark$\checkmark^\checkmark\\checkmarkc\checkmarki\checkmarkr\checkmarkc\checkmark$\checkmarkr\checkmark$\checkmark^\checkmark\\checkmarkc\checkmarki\checkmarkr\checkmarkc\checkmark$\checkmarke\checkmark$\checkmark^\checkmark\\checkmarkc\checkmarki\checkmarkr\checkmarkc\checkmark$\checkmarks\checkmark$\checkmark^\checkmark\\checkmarkc\checkmarki\checkmarkr\checkmarkc\checkmark$\checkmark}\checkmark$\checkmark^\checkmark\\checkmarkc\checkmarki\checkmarkr\checkmarkc\checkmark$\checkmark \checkmark$\checkmark^\checkmark\\checkmarkc\checkmarki\checkmarkr\checkmarkc\checkmark$\checkmark=\checkmark$\checkmark^\checkmark\\checkmarkc\checkmarki\checkmarkr\checkmarkc\checkmark$\checkmark \checkmark$\checkmark^\checkmark\\checkmarkc\checkmarki\checkmarkr\checkmarkc\checkmark$\checkmarkR\checkmark$\checkmark^\checkmark\\checkmarkc\checkmarki\checkmarkr\checkmarkc\checkmark$\checkmark_\checkmark$\checkmark^\checkmark\\checkmarkc\checkmarki\checkmarkr\checkmarkc\checkmark$\checkmark{\checkmark$\checkmark^\checkmark\\checkmarkc\checkmarki\checkmarkr\checkmarkc\checkmark$\checkmarkp\checkmark$\checkmark^\checkmark\\checkmarkc\checkmarki\checkmarkr\checkmarkc\checkmark$\checkmarkl\checkmark$\checkmark^\checkmark\\checkmarkc\checkmarki\checkmarkr\checkmarkc\checkmark$\checkmarka\checkmark$\checkmark^\checkmark\\checkmarkc\checkmarki\checkmarkr\checkmarkc\checkmark$\checkmarkt\checkmark$\checkmark^\checkmark\\checkmarkc\checkmarki\checkmarkr\checkmarkc\checkmark$\checkmarke\checkmark$\checkmark^\checkmark\\checkmarkc\checkmarki\checkmarkr\checkmarkc\checkmark$\checkmarka\checkmark$\checkmark^\checkmark\\checkmarkc\checkmarki\checkmarkr\checkmarkc\checkmark$\checkmarku\checkmark$\checkmark^\checkmark\\checkmarkc\checkmarki\checkmarkr\checkmarkc\checkmark$\checkmark}\checkmark$\checkmark^\checkmark\\checkmarkc\checkmarki\checkmarkr\checkmarkc\checkmark$\checkmark \checkmark$\checkmark^\checkmark\\checkmarkc\checkmarki\checkmarkr\checkmarkc\checkmark$\checkmark\\checkmark$\checkmark^\checkmark\\checkmarkc\checkmarki\checkmarkr\checkmarkc\checkmark$\checkmarkt\checkmark$\checkmark^\checkmark\\checkmarkc\checkmarki\checkmarkr\checkmarkc\checkmark$\checkmarki\checkmark$\checkmark^\checkmark\\checkmarkc\checkmarki\checkmarkr\checkmarkc\checkmark$\checkmarkm\checkmark$\checkmark^\checkmark\\checkmarkc\checkmarki\checkmarkr\checkmarkc\checkmark$\checkmarke\checkmark$\checkmark^\checkmark\\checkmarkc\checkmarki\checkmarkr\checkmarkc\checkmark$\checkmarks\checkmark$\checkmark^\checkmark\\checkmarkc\checkmarki\checkmarkr\checkmarkc\checkmark$\checkmark \checkmark$\checkmark^\checkmark\\checkmarkc\checkmarki\checkmarkr\checkmarkc\checkmark$\checkmark\\checkmark$\checkmark^\checkmark\\checkmarkc\checkmarki\checkmarkr\checkmarkc\checkmark$\checkmarkt\checkmark$\checkmark^\checkmark\\checkmarkc\checkmarki\checkmarkr\checkmarkc\checkmark$\checkmarka\checkmark$\checkmark^\checkmark\\checkmarkc\checkmarki\checkmarkr\checkmarkc\checkmark$\checkmarkn\checkmark$\checkmark^\checkmark\\checkmarkc\checkmarki\checkmarkr\checkmarkc\checkmark$\checkmark(\checkmark$\checkmark^\checkmark\\checkmarkc\checkmarki\checkmarkr\checkmarkc\checkmark$\checkmark0\checkmark$\checkmark^\checkmark\\checkmarkc\checkmarki\checkmarkr\checkmarkc\checkmark$\checkmark{\checkmark$\checkmark^\checkmark\\checkmarkc\checkmarki\checkmarkr\checkmarkc\checkmark$\checkmark,\checkmark$\checkmark^\checkmark\\checkmarkc\checkmarki\checkmarkr\checkmarkc\checkmark$\checkmark}\checkmark$\checkmark^\checkmark\\checkmarkc\checkmarki\checkmarkr\checkmarkc\checkmark$\checkmark0\checkmark$\checkmark^\checkmark\\checkmarkc\checkmarki\checkmarkr\checkmarkc\checkmark$\checkmark0\checkmark$\checkmark^\checkmark\\checkmarkc\checkmarki\checkmarkr\checkmarkc\checkmark$\checkmark2\checkmark$\checkmark^\checkmark\\checkmarkc\checkmarki\checkmarkr\checkmarkc\checkmark$\checkmark2\checkmark$\checkmark^\checkmark\\checkmarkc\checkmarki\checkmarkr\checkmarkc\checkmark$\checkmark5\checkmark$\checkmark^\checkmark\\checkmarkc\checkmarki\checkmarkr\checkmarkc\checkmark$\checkmark°\checkmark$\checkmark^\checkmark\\checkmarkc\checkmarki\checkmarkr\checkmarkc\checkmark$\checkmark)\checkmark$\checkmark^\checkmark\\checkmarkc\checkmarki\checkmarkr\checkmarkc\checkmark$\checkmark \checkmark$\checkmark^\checkmark\\checkmarkc\checkmarki\checkmarkr\checkmarkc\checkmark$\checkmark=\checkmark$\checkmark^\checkmark\\checkmarkc\checkmarki\checkmarkr\checkmarkc\checkmark$\checkmark \checkmark$\checkmark^\checkmark\\checkmarkc\checkmarki\checkmarkr\checkmarkc\checkmark$\checkmark1\checkmark$\checkmark^\checkmark\\checkmarkc\checkmarki\checkmarkr\checkmarkc\checkmark$\checkmark0\checkmark$\checkmark^\checkmark\\checkmarkc\checkmarki\checkmarkr\checkmarkc\checkmark$\checkmark0\checkmark$\checkmark^\checkmark\\checkmarkc\checkmarki\checkmarkr\checkmarkc\checkmark$\checkmark \checkmark$\checkmark^\checkmark\\checkmarkc\checkmarki\checkmarkr\checkmarkc\checkmark$\checkmark\\checkmark$\checkmark^\checkmark\\checkmarkc\checkmarki\checkmarkr\checkmarkc\checkmark$\checkmarkt\checkmark$\checkmark^\checkmark\\checkmarkc\checkmarki\checkmarkr\checkmarkc\checkmark$\checkmarki\checkmark$\checkmark^\checkmark\\checkmarkc\checkmarki\checkmarkr\checkmarkc\checkmark$\checkmarkm\checkmark$\checkmark^\checkmark\\checkmarkc\checkmarki\checkmarkr\checkmarkc\checkmark$\checkmarke\checkmark$\checkmark^\checkmark\\checkmarkc\checkmarki\checkmarkr\checkmarkc\checkmark$\checkmarks\checkmark$\checkmark^\checkmark\\checkmarkc\checkmarki\checkmarkr\checkmarkc\checkmark$\checkmark \checkmark$\checkmark^\checkmark\\checkmarkc\checkmarki\checkmarkr\checkmarkc\checkmark$\checkmark\\checkmark$\checkmark^\checkmark\\checkmarkc\checkmarki\checkmarkr\checkmarkc\checkmark$\checkmarkt\checkmark$\checkmark^\checkmark\\checkmarkc\checkmarki\checkmarkr\checkmarkc\checkmark$\checkmarka\checkmark$\checkmark^\checkmark\\checkmarkc\checkmarki\checkmarkr\checkmarkc\checkmark$\checkmarkn\checkmark$\checkmark^\checkmark\\checkmarkc\checkmarki\checkmarkr\checkmarkc\checkmark$\checkmark(\checkmark$\checkmark^\checkmark\\checkmarkc\checkmarki\checkmarkr\checkmarkc\checkmark$\checkmark0\checkmark$\checkmark^\checkmark\\checkmarkc\checkmarki\checkmarkr\checkmarkc\checkmark$\checkmark{\checkmark$\checkmark^\checkmark\\checkmarkc\checkmarki\checkmarkr\checkmarkc\checkmark$\checkmark,\checkmark$\checkmark^\checkmark\\checkmarkc\checkmarki\checkmarkr\checkmarkc\checkmark$\checkmark}\checkmark$\checkmark^\checkmark\\checkmarkc\checkmarki\checkmarkr\checkmarkc\checkmark$\checkmark0\checkmark$\checkmark^\checkmark\\checkmarkc\checkmarki\checkmarkr\checkmarkc\checkmark$\checkmark0\checkmark$\checkmark^\checkmark\\checkmarkc\checkmarki\checkmarkr\checkmarkc\checkmark$\checkmark2\checkmark$\checkmark^\checkmark\\checkmarkc\checkmarki\checkmarkr\checkmarkc\checkmark$\checkmark2\checkmark$\checkmark^\checkmark\\checkmarkc\checkmarki\checkmarkr\checkmarkc\checkmark$\checkmark5\checkmark$\checkmark^\checkmark\\checkmarkc\checkmarki\checkmarkr\checkmarkc\checkmark$\checkmark°\checkmark$\checkmark^\checkmark\\checkmarkc\checkmarki\checkmarkr\checkmarkc\checkmark$\checkmark)\checkmark$\checkmark^\checkmark\\checkmarkc\checkmarki\checkmarkr\checkmarkc\checkmark$\checkmark \checkmark$\checkmark^\checkmark\\checkmarkc\checkmarki\checkmarkr\checkmarkc\checkmark$\checkmark\\checkmark$\checkmark^\checkmark\\checkmarkc\checkmarki\checkmarkr\checkmarkc\checkmark$\checkmarka\checkmark$\checkmark^\checkmark\\checkmarkc\checkmarki\checkmarkr\checkmarkc\checkmark$\checkmarkp\checkmark$\checkmark^\checkmark\\checkmarkc\checkmarki\checkmarkr\checkmarkc\checkmark$\checkmarkp\checkmark$\checkmark^\checkmark\\checkmarkc\checkmarki\checkmarkr\checkmarkc\checkmark$\checkmarkr\checkmark$\checkmark^\checkmark\\checkmarkc\checkmarki\checkmarkr\checkmarkc\checkmark$\checkmarko\checkmark$\checkmark^\checkmark\\checkmarkc\checkmarki\checkmarkr\checkmarkc\checkmark$\checkmarkx\checkmark$\checkmark^\checkmark\\checkmarkc\checkmarki\checkmarkr\checkmarkc\checkmark$\checkmark \checkmark$\checkmark^\checkmark\\checkmarkc\checkmarki\checkmarkr\checkmarkc\checkmark$\checkmark\\checkmark$\checkmark^\checkmark\\checkmarkc\checkmarki\checkmarkr\checkmarkc\checkmark$\checkmarkb\checkmark$\checkmark^\checkmark\\checkmarkc\checkmarki\checkmarkr\checkmarkc\checkmark$\checkmarko\checkmark$\checkmark^\checkmark\\checkmarkc\checkmarki\checkmarkr\checkmarkc\checkmark$\checkmarkx\checkmark$\checkmark^\checkmark\\checkmarkc\checkmarki\checkmarkr\checkmarkc\checkmark$\checkmarke\checkmark$\checkmark^\checkmark\\checkmarkc\checkmarki\checkmarkr\checkmarkc\checkmark$\checkmarkd\checkmark$\checkmark^\checkmark\\checkmarkc\checkmarki\checkmarkr\checkmarkc\checkmark$\checkmark{\checkmark$\checkmark^\checkmark\\checkmarkc\checkmarki\checkmarkr\checkmarkc\checkmark$\checkmark0\checkmark$\checkmark^\checkmark\\checkmarkc\checkmarki\checkmarkr\checkmarkc\checkmark$\checkmark{\checkmark$\checkmark^\checkmark\\checkmarkc\checkmarki\checkmarkr\checkmarkc\checkmark$\checkmark,\checkmark$\checkmark^\checkmark\\checkmarkc\checkmarki\checkmarkr\checkmarkc\checkmark$\checkmark}\checkmark$\checkmark^\checkmark\\checkmarkc\checkmarki\checkmarkr\checkmarkc\checkmark$\checkmark0\checkmark$\checkmark^\checkmark\\checkmarkc\checkmarki\checkmarkr\checkmarkc\checkmark$\checkmark0\checkmark$\checkmark^\checkmark\\checkmarkc\checkmarki\checkmarkr\checkmarkc\checkmark$\checkmark4\checkmark$\checkmark^\checkmark\\checkmarkc\checkmarki\checkmarkr\checkmarkc\checkmark$\checkmark \checkmark$\checkmark^\checkmark\\checkmarkc\checkmarki\checkmarkr\checkmarkc\checkmark$\checkmark\\checkmark$\checkmark^\checkmark\\checkmarkc\checkmarki\checkmarkr\checkmarkc\checkmark$\checkmarkt\checkmark$\checkmark^\checkmark\\checkmarkc\checkmarki\checkmarkr\checkmarkc\checkmark$\checkmarke\checkmark$\checkmark^\checkmark\\checkmarkc\checkmarki\checkmarkr\checkmarkc\checkmark$\checkmarkx\checkmark$\checkmark^\checkmark\\checkmarkc\checkmarki\checkmarkr\checkmarkc\checkmark$\checkmarkt\checkmark$\checkmark^\checkmark\\checkmarkc\checkmarki\checkmarkr\checkmarkc\checkmark$\checkmark{\checkmark$\checkmark^\checkmark\\checkmarkc\checkmarki\checkmarkr\checkmarkc\checkmark$\checkmark \checkmark$\checkmark^\checkmark\\checkmarkc\checkmarki\checkmarkr\checkmarkc\checkmark$\checkmarkm\checkmark$\checkmark^\checkmark\\checkmarkc\checkmarki\checkmarkr\checkmarkc\checkmark$\checkmarkm\checkmark$\checkmark^\checkmark\\checkmarkc\checkmarki\checkmarkr\checkmarkc\checkmark$\checkmark}\checkmark$\checkmark^\checkmark\\checkmarkc\checkmarki\checkmarkr\checkmarkc\checkmark$\checkmark}\checkmark$\checkmark^\checkmark\\checkmarkc\checkmarki\checkmarkr\checkmarkc\checkmark$\checkmark
\checkmark$\checkmark^\checkmark\\checkmarkc\checkmarki\checkmarkr\checkmarkc\checkmark$\checkmark\\checkmark$\checkmark^\checkmark\\checkmarkc\checkmarki\checkmarkr\checkmarkc\checkmark$\checkmarke\checkmark$\checkmark^\checkmark\\checkmarkc\checkmarki\checkmarkr\checkmarkc\checkmark$\checkmarkn\checkmark$\checkmark^\checkmark\\checkmarkc\checkmarki\checkmarkr\checkmarkc\checkmark$\checkmarkd\checkmark$\checkmark^\checkmark\\checkmarkc\checkmarki\checkmarkr\checkmarkc\checkmark$\checkmark{\checkmark$\checkmark^\checkmark\\checkmarkc\checkmarki\checkmarkr\checkmarkc\checkmark$\checkmarke\checkmark$\checkmark^\checkmark\\checkmarkc\checkmarki\checkmarkr\checkmarkc\checkmark$\checkmarkq\checkmark$\checkmark^\checkmark\\checkmarkc\checkmarki\checkmarkr\checkmarkc\checkmark$\checkmarku\checkmark$\checkmark^\checkmark\\checkmarkc\checkmarki\checkmarkr\checkmarkc\checkmark$\checkmarka\checkmark$\checkmark^\checkmark\\checkmarkc\checkmarki\checkmarkr\checkmarkc\checkmark$\checkmarkt\checkmark$\checkmark^\checkmark\\checkmarkc\checkmarki\checkmarkr\checkmarkc\checkmark$\checkmarki\checkmark$\checkmark^\checkmark\\checkmarkc\checkmarki\checkmarkr\checkmarkc\checkmark$\checkmarko\checkmark$\checkmark^\checkmark\\checkmarkc\checkmarki\checkmarkr\checkmarkc\checkmark$\checkmarkn\checkmark$\checkmark^\checkmark\\checkmarkc\checkmarki\checkmarkr\checkmarkc\checkmark$\checkmark}\checkmark$\checkmark^\checkmark\\checkmarkc\checkmarki\checkmarkr\checkmarkc\checkmark$\checkmark
\checkmark$\checkmark^\checkmark\\checkmarkc\checkmarki\checkmarkr\checkmarkc\checkmark$\checkmarkC\checkmark$\checkmark^\checkmark\\checkmarkc\checkmarki\checkmarkr\checkmarkc\checkmark$\checkmarke\checkmark$\checkmark^\checkmark\\checkmarkc\checkmarki\checkmarkr\checkmarkc\checkmark$\checkmarkt\checkmark$\checkmark^\checkmark\\checkmarkc\checkmarki\checkmarkr\checkmarkc\checkmark$\checkmarkt\checkmark$\checkmark^\checkmark\\checkmarkc\checkmarki\checkmarkr\checkmarkc\checkmark$\checkmarke\checkmark$\checkmark^\checkmark\\checkmarkc\checkmarki\checkmarkr\checkmarkc\checkmark$\checkmark \checkmark$\checkmark^\checkmark\\checkmarkc\checkmarki\checkmarkr\checkmarkc\checkmark$\checkmarkr\checkmark$\checkmark^\checkmark\\checkmarkc\checkmarki\checkmarkr\checkmarkc\checkmark$\checkmarké\checkmark$\checkmark^\checkmark\\checkmarkc\checkmarki\checkmarkr\checkmarkc\checkmark$\checkmarks\checkmark$\checkmark^\checkmark\\checkmarkc\checkmarki\checkmarkr\checkmarkc\checkmark$\checkmarko\checkmark$\checkmark^\checkmark\\checkmarkc\checkmarki\checkmarkr\checkmarkc\checkmark$\checkmarkl\checkmark$\checkmark^\checkmark\\checkmarkc\checkmarki\checkmarkr\checkmarkc\checkmark$\checkmarku\checkmark$\checkmark^\checkmark\\checkmarkc\checkmarki\checkmarkr\checkmarkc\checkmark$\checkmarkt\checkmark$\checkmark^\checkmark\\checkmarkc\checkmarki\checkmarkr\checkmarkc\checkmark$\checkmarki\checkmark$\checkmark^\checkmark\\checkmarkc\checkmarki\checkmarkr\checkmarkc\checkmark$\checkmarko\checkmark$\checkmark^\checkmark\\checkmarkc\checkmarki\checkmarkr\checkmarkc\checkmark$\checkmarkn\checkmark$\checkmark^\checkmark\\checkmarkc\checkmarki\checkmarkr\checkmarkc\checkmark$\checkmark \checkmark$\checkmark^\checkmark\\checkmarkc\checkmarki\checkmarkr\checkmarkc\checkmark$\checkmarkd\checkmark$\checkmark^\checkmark\\checkmarkc\checkmarki\checkmarkr\checkmarkc\checkmark$\checkmarke\checkmark$\checkmark^\checkmark\\checkmarkc\checkmarki\checkmarkr\checkmarkc\checkmark$\checkmark \checkmark$\checkmark^\checkmark\\checkmarkc\checkmarki\checkmarkr\checkmarkc\checkmark$\checkmark\\checkmark$\checkmark^\checkmark\\checkmarkc\checkmarki\checkmarkr\checkmarkc\checkmark$\checkmarkt\checkmark$\checkmark^\checkmark\\checkmarkc\checkmarki\checkmarkr\checkmarkc\checkmark$\checkmarke\checkmark$\checkmark^\checkmark\\checkmarkc\checkmarki\checkmarkr\checkmarkc\checkmark$\checkmarkx\checkmark$\checkmark^\checkmark\\checkmarkc\checkmarki\checkmarkr\checkmarkc\checkmark$\checkmarkt\checkmark$\checkmark^\checkmark\\checkmarkc\checkmarki\checkmarkr\checkmarkc\checkmark$\checkmarkb\checkmark$\checkmark^\checkmark\\checkmarkc\checkmarki\checkmarkr\checkmarkc\checkmark$\checkmarkf\checkmark$\checkmark^\checkmark\\checkmarkc\checkmarki\checkmarkr\checkmarkc\checkmark$\checkmark{\checkmark$\checkmark^\checkmark\\checkmarkc\checkmarki\checkmarkr\checkmarkc\checkmark$\checkmark4\checkmark$\checkmark^\checkmark\\checkmarkc\checkmarki\checkmarkr\checkmarkc\checkmark$\checkmark \checkmark$\checkmark^\checkmark\\checkmarkc\checkmarki\checkmarkr\checkmarkc\checkmark$\checkmarkµ\checkmark$\checkmark^\checkmark\\checkmarkc\checkmarki\checkmarkr\checkmarkc\checkmark$\checkmarkm\checkmark$\checkmark^\checkmark\\checkmarkc\checkmarki\checkmarkr\checkmarkc\checkmark$\checkmark}\checkmark$\checkmark^\checkmark\\checkmarkc\checkmarki\checkmarkr\checkmarkc\checkmark$\checkmark \checkmark$\checkmark^\checkmark\\checkmarkc\checkmarki\checkmarkr\checkmarkc\checkmark$\checkmarke\checkmark$\checkmark^\checkmark\\checkmarkc\checkmarki\checkmarkr\checkmarkc\checkmark$\checkmarks\checkmark$\checkmark^\checkmark\\checkmarkc\checkmarki\checkmarkr\checkmarkc\checkmark$\checkmarkt\checkmark$\checkmark^\checkmark\\checkmarkc\checkmarki\checkmarkr\checkmarkc\checkmark$\checkmark \checkmark$\checkmark^\checkmark\\checkmarkc\checkmarki\checkmarkr\checkmarkc\checkmark$\checkmarkb\checkmark$\checkmark^\checkmark\\checkmarkc\checkmarki\checkmarkr\checkmarkc\checkmark$\checkmarki\checkmark$\checkmark^\checkmark\\checkmarkc\checkmarki\checkmarkr\checkmarkc\checkmark$\checkmarke\checkmark$\checkmark^\checkmark\\checkmarkc\checkmarki\checkmarkr\checkmarkc\checkmark$\checkmarkn\checkmark$\checkmark^\checkmark\\checkmarkc\checkmarki\checkmarkr\checkmarkc\checkmark$\checkmark \checkmark$\checkmark^\checkmark\\checkmarkc\checkmarki\checkmarkr\checkmarkc\checkmark$\checkmarki\checkmark$\checkmark^\checkmark\\checkmarkc\checkmarki\checkmarkr\checkmarkc\checkmark$\checkmarkn\checkmark$\checkmark^\checkmark\\checkmarkc\checkmarki\checkmarkr\checkmarkc\checkmark$\checkmarkf\checkmark$\checkmark^\checkmark\\checkmarkc\checkmarki\checkmarkr\checkmarkc\checkmark$\checkmarké\checkmark$\checkmark^\checkmark\\checkmarkc\checkmarki\checkmarkr\checkmarkc\checkmark$\checkmarkr\checkmark$\checkmark^\checkmark\\checkmarkc\checkmarki\checkmarkr\checkmarkc\checkmark$\checkmarki\checkmark$\checkmark^\checkmark\\checkmarkc\checkmarki\checkmarkr\checkmarkc\checkmark$\checkmarke\checkmark$\checkmark^\checkmark\\checkmarkc\checkmarki\checkmarkr\checkmarkc\checkmark$\checkmarku\checkmark$\checkmark^\checkmark\\checkmarkc\checkmarki\checkmarkr\checkmarkc\checkmark$\checkmarkr\checkmark$\checkmark^\checkmark\\checkmarkc\checkmarki\checkmarkr\checkmarkc\checkmark$\checkmarke\checkmark$\checkmark^\checkmark\\checkmarkc\checkmarki\checkmarkr\checkmarkc\checkmark$\checkmark \checkmark$\checkmark^\checkmark\\checkmarkc\checkmarki\checkmarkr\checkmarkc\checkmark$\checkmarkà\checkmark$\checkmark^\checkmark\\checkmarkc\checkmarki\checkmarkr\checkmarkc\checkmark$\checkmark \checkmark$\checkmark^\checkmark\\checkmarkc\checkmarki\checkmarkr\checkmarkc\checkmark$\checkmarkl\checkmark$\checkmark^\checkmark\\checkmarkc\checkmarki\checkmarkr\checkmarkc\checkmark$\checkmarka\checkmark$\checkmark^\checkmark\\checkmarkc\checkmarki\checkmarkr\checkmarkc\checkmark$\checkmark \checkmark$\checkmark^\checkmark\\checkmarkc\checkmarki\checkmarkr\checkmarkc\checkmark$\checkmarkt\checkmark$\checkmark^\checkmark\\checkmarkc\checkmarki\checkmarkr\checkmarkc\checkmark$\checkmarko\checkmark$\checkmark^\checkmark\\checkmarkc\checkmarki\checkmarkr\checkmarkc\checkmark$\checkmarkl\checkmark$\checkmark^\checkmark\\checkmarkc\checkmarki\checkmarkr\checkmarkc\checkmark$\checkmarké\checkmark$\checkmark^\checkmark\\checkmarkc\checkmarki\checkmarkr\checkmarkc\checkmark$\checkmarkr\checkmark$\checkmark^\checkmark\\checkmarkc\checkmarki\checkmarkr\checkmarkc\checkmark$\checkmarka\checkmark$\checkmark^\checkmark\\checkmarkc\checkmarki\checkmarkr\checkmarkc\checkmark$\checkmarkn\checkmark$\checkmark^\checkmark\\checkmarkc\checkmarki\checkmarkr\checkmarkc\checkmark$\checkmarkc\checkmark$\checkmark^\checkmark\\checkmarkc\checkmarki\checkmarkr\checkmarkc\checkmark$\checkmarke\checkmark$\checkmark^\checkmark\\checkmarkc\checkmarki\checkmarkr\checkmarkc\checkmark$\checkmark \checkmark$\checkmark^\checkmark\\checkmarkc\checkmarki\checkmarkr\checkmarkc\checkmark$\checkmarke\checkmark$\checkmark^\checkmark\\checkmarkc\checkmarki\checkmarkr\checkmarkc\checkmark$\checkmarkx\checkmark$\checkmark^\checkmark\\checkmarkc\checkmarki\checkmarkr\checkmarkc\checkmark$\checkmarki\checkmark$\checkmark^\checkmark\\checkmarkc\checkmarki\checkmarkr\checkmarkc\checkmark$\checkmarkg\checkmark$\checkmark^\checkmark\\checkmarkc\checkmarki\checkmarkr\checkmarkc\checkmark$\checkmarké\checkmark$\checkmark^\checkmark\\checkmarkc\checkmarki\checkmarkr\checkmarkc\checkmark$\checkmarke\checkmark$\checkmark^\checkmark\\checkmarkc\checkmarki\checkmarkr\checkmarkc\checkmark$\checkmark \checkmark$\checkmark^\checkmark\\checkmarkc\checkmarki\checkmarkr\checkmarkc\checkmark$\checkmarkd\checkmark$\checkmark^\checkmark\\checkmarkc\checkmarki\checkmarkr\checkmarkc\checkmark$\checkmarke\checkmark$\checkmark^\checkmark\\checkmarkc\checkmarki\checkmarkr\checkmarkc\checkmark$\checkmark \checkmark$\checkmark^\checkmark\\checkmarkc\checkmarki\checkmarkr\checkmarkc\checkmark$\checkmark$\checkmark$\checkmark^\checkmark\\checkmarkc\checkmarki\checkmarkr\checkmarkc\checkmark$\checkmark\\checkmark$\checkmark^\checkmark\\checkmarkc\checkmarki\checkmarkr\checkmarkc\checkmark$\checkmarkp\checkmark$\checkmark^\checkmark\\checkmarkc\checkmarki\checkmarkr\checkmarkc\checkmark$\checkmarkm\checkmark$\checkmark^\checkmark\\checkmarkc\checkmarki\checkmarkr\checkmarkc\checkmark$\checkmark \checkmark$\checkmark^\checkmark\\checkmarkc\checkmarki\checkmarkr\checkmarkc\checkmark$\checkmark5\checkmark$\checkmark^\checkmark\\checkmarkc\checkmarki\checkmarkr\checkmarkc\checkmark$\checkmark0\checkmark$\checkmark^\checkmark\\checkmarkc\checkmarki\checkmarkr\checkmarkc\checkmark$\checkmark$\checkmark$\checkmark^\checkmark\\checkmarkc\checkmarki\checkmarkr\checkmarkc\checkmark$\checkmark \checkmark$\checkmark^\checkmark\\checkmarkc\checkmarki\checkmarkr\checkmarkc\checkmark$\checkmarkµ\checkmark$\checkmark^\checkmark\\checkmarkc\checkmarki\checkmarkr\checkmarkc\checkmark$\checkmarkm\checkmark$\checkmark^\checkmark\\checkmarkc\checkmarki\checkmarkr\checkmarkc\checkmark$\checkmark.\checkmark$\checkmark^\checkmark\\checkmarkc\checkmarki\checkmarkr\checkmarkc\checkmark$\checkmark \checkmark$\checkmark^\checkmark\\checkmarkc\checkmarki\checkmarkr\checkmarkc\checkmark$\checkmarkL\checkmark$\checkmark^\checkmark\\checkmarkc\checkmarki\checkmarkr\checkmarkc\checkmark$\checkmarka\checkmark$\checkmark^\checkmark\\checkmarkc\checkmarki\checkmarkr\checkmarkc\checkmark$\checkmark \checkmark$\checkmark^\checkmark\\checkmarkc\checkmarki\checkmarkr\checkmarkc\checkmark$\checkmarkr\checkmark$\checkmark^\checkmark\\checkmarkc\checkmarki\checkmarkr\checkmarkc\checkmark$\checkmarké\checkmark$\checkmark^\checkmark\\checkmarkc\checkmarki\checkmarkr\checkmarkc\checkmark$\checkmarkd\checkmark$\checkmark^\checkmark\\checkmarkc\checkmarki\checkmarkr\checkmarkc\checkmark$\checkmarku\checkmark$\checkmark^\checkmark\\checkmarkc\checkmarki\checkmarkr\checkmarkc\checkmark$\checkmarkc\checkmark$\checkmark^\checkmark\\checkmarkc\checkmarki\checkmarkr\checkmarkc\checkmark$\checkmarkt\checkmark$\checkmark^\checkmark\\checkmarkc\checkmarki\checkmarkr\checkmarkc\checkmark$\checkmarki\checkmark$\checkmark^\checkmark\\checkmarkc\checkmarki\checkmarkr\checkmarkc\checkmark$\checkmarko\checkmark$\checkmark^\checkmark\\checkmarkc\checkmarki\checkmarkr\checkmarkc\checkmark$\checkmarkn\checkmark$\checkmark^\checkmark\\checkmarkc\checkmarki\checkmarkr\checkmarkc\checkmark$\checkmark \checkmark$\checkmark^\checkmark\\checkmarkc\checkmarki\checkmarkr\checkmarkc\checkmark$\checkmark1\checkmark$\checkmark^\checkmark\\checkmarkc\checkmarki\checkmarkr\checkmarkc\checkmark$\checkmark:\checkmark$\checkmark^\checkmark\\checkmarkc\checkmarki\checkmarkr\checkmarkc\checkmark$\checkmark5\checkmark$\checkmark^\checkmark\\checkmarkc\checkmarki\checkmarkr\checkmarkc\checkmark$\checkmark0\checkmark$\checkmark^\checkmark\\checkmarkc\checkmarki\checkmarkr\checkmarkc\checkmark$\checkmark \checkmark$\checkmark^\checkmark\\checkmarkc\checkmarki\checkmarkr\checkmarkc\checkmark$\checkmarke\checkmark$\checkmark^\checkmark\\checkmarkc\checkmarki\checkmarkr\checkmarkc\checkmark$\checkmarks\checkmark$\checkmark^\checkmark\\checkmarkc\checkmarki\checkmarkr\checkmarkc\checkmark$\checkmarkt\checkmark$\checkmark^\checkmark\\checkmarkc\checkmarki\checkmarkr\checkmarkc\checkmark$\checkmark \checkmark$\checkmark^\checkmark\\checkmarkc\checkmarki\checkmarkr\checkmarkc\checkmark$\checkmarkv\checkmark$\checkmark^\checkmark\\checkmarkc\checkmarki\checkmarkr\checkmarkc\checkmark$\checkmarka\checkmark$\checkmark^\checkmark\\checkmarkc\checkmarki\checkmarkr\checkmarkc\checkmark$\checkmarkl\checkmark$\checkmark^\checkmark\\checkmarkc\checkmarki\checkmarkr\checkmarkc\checkmark$\checkmarki\checkmark$\checkmark^\checkmark\\checkmarkc\checkmarki\checkmarkr\checkmarkc\checkmark$\checkmarkd\checkmark$\checkmark^\checkmark\\checkmarkc\checkmarki\checkmarkr\checkmarkc\checkmark$\checkmarké\checkmark$\checkmark^\checkmark\\checkmarkc\checkmarki\checkmarkr\checkmarkc\checkmark$\checkmarke\checkmark$\checkmark^\checkmark\\checkmarkc\checkmarki\checkmarkr\checkmarkc\checkmark$\checkmark.\checkmark$\checkmark^\checkmark\\checkmarkc\checkmarki\checkmarkr\checkmarkc\checkmark$\checkmark
\checkmark$\checkmark^\checkmark\\checkmarkc\checkmarki\checkmarkr\checkmarkc\checkmark$\checkmark
\checkmark$\checkmark^\checkmark\\checkmarkc\checkmarki\checkmarkr\checkmarkc\checkmark$\checkmark\\checkmark$\checkmark^\checkmark\\checkmarkc\checkmarki\checkmarkr\checkmarkc\checkmark$\checkmarks\checkmark$\checkmark^\checkmark\\checkmarkc\checkmarki\checkmarkr\checkmarkc\checkmark$\checkmarku\checkmark$\checkmark^\checkmark\\checkmarkc\checkmarki\checkmarkr\checkmarkc\checkmark$\checkmarkb\checkmark$\checkmark^\checkmark\\checkmarkc\checkmarki\checkmarkr\checkmarkc\checkmark$\checkmarks\checkmark$\checkmark^\checkmark\\checkmarkc\checkmarki\checkmarkr\checkmarkc\checkmark$\checkmarke\checkmark$\checkmark^\checkmark\\checkmarkc\checkmarki\checkmarkr\checkmarkc\checkmark$\checkmarkc\checkmark$\checkmark^\checkmark\\checkmarkc\checkmarki\checkmarkr\checkmarkc\checkmark$\checkmarkt\checkmark$\checkmark^\checkmark\\checkmarkc\checkmarki\checkmarkr\checkmarkc\checkmark$\checkmarki\checkmark$\checkmark^\checkmark\\checkmarkc\checkmarki\checkmarkr\checkmarkc\checkmark$\checkmarko\checkmark$\checkmark^\checkmark\\checkmarkc\checkmarki\checkmarkr\checkmarkc\checkmark$\checkmarkn\checkmark$\checkmark^\checkmark\\checkmarkc\checkmarki\checkmarkr\checkmarkc\checkmark$\checkmark{\checkmark$\checkmark^\checkmark\\checkmarkc\checkmarki\checkmarkr\checkmarkc\checkmark$\checkmarkD\checkmark$\checkmark^\checkmark\\checkmarkc\checkmarki\checkmarkr\checkmarkc\checkmark$\checkmarki\checkmark$\checkmark^\checkmark\\checkmarkc\checkmarki\checkmarkr\checkmarkc\checkmark$\checkmarkm\checkmark$\checkmark^\checkmark\\checkmarkc\checkmarki\checkmarkr\checkmarkc\checkmark$\checkmarke\checkmark$\checkmark^\checkmark\\checkmarkc\checkmarki\checkmarkr\checkmarkc\checkmark$\checkmarkn\checkmark$\checkmark^\checkmark\\checkmarkc\checkmarki\checkmarkr\checkmarkc\checkmark$\checkmarks\checkmark$\checkmark^\checkmark\\checkmarkc\checkmarki\checkmarkr\checkmarkc\checkmark$\checkmarki\checkmark$\checkmark^\checkmark\\checkmarkc\checkmarki\checkmarkr\checkmarkc\checkmark$\checkmarko\checkmark$\checkmark^\checkmark\\checkmarkc\checkmarki\checkmarkr\checkmarkc\checkmark$\checkmarkn\checkmark$\checkmark^\checkmark\\checkmarkc\checkmarki\checkmarkr\checkmarkc\checkmark$\checkmarkn\checkmark$\checkmark^\checkmark\\checkmarkc\checkmarki\checkmarkr\checkmarkc\checkmark$\checkmarke\checkmark$\checkmark^\checkmark\\checkmarkc\checkmarki\checkmarkr\checkmarkc\checkmark$\checkmarkm\checkmark$\checkmark^\checkmark\\checkmarkc\checkmarki\checkmarkr\checkmarkc\checkmark$\checkmarke\checkmark$\checkmark^\checkmark\\checkmarkc\checkmarki\checkmarkr\checkmarkc\checkmark$\checkmarkn\checkmark$\checkmark^\checkmark\\checkmarkc\checkmarki\checkmarkr\checkmarkc\checkmark$\checkmarkt\checkmark$\checkmark^\checkmark\\checkmarkc\checkmarki\checkmarkr\checkmarkc\checkmark$\checkmark \checkmark$\checkmark^\checkmark\\checkmarkc\checkmarki\checkmarkr\checkmarkc\checkmark$\checkmarkd\checkmark$\checkmark^\checkmark\\checkmarkc\checkmarki\checkmarkr\checkmarkc\checkmark$\checkmarke\checkmark$\checkmark^\checkmark\\checkmarkc\checkmarki\checkmarkr\checkmarkc\checkmark$\checkmark \checkmark$\checkmark^\checkmark\\checkmarkc\checkmarki\checkmarkr\checkmarkc\checkmark$\checkmarkl\checkmark$\checkmark^\checkmark\\checkmarkc\checkmarki\checkmarkr\checkmarkc\checkmark$\checkmarka\checkmark$\checkmark^\checkmark\\checkmarkc\checkmarki\checkmarkr\checkmarkc\checkmark$\checkmark \checkmark$\checkmark^\checkmark\\checkmarkc\checkmarki\checkmarkr\checkmarkc\checkmark$\checkmarkt\checkmark$\checkmark^\checkmark\\checkmarkc\checkmarki\checkmarkr\checkmarkc\checkmark$\checkmarkr\checkmark$\checkmark^\checkmark\\checkmarkc\checkmarki\checkmarkr\checkmarkc\checkmark$\checkmarka\checkmark$\checkmark^\checkmark\\checkmarkc\checkmarki\checkmarkr\checkmarkc\checkmark$\checkmarkn\checkmark$\checkmark^\checkmark\\checkmarkc\checkmarki\checkmarkr\checkmarkc\checkmark$\checkmarks\checkmark$\checkmark^\checkmark\\checkmarkc\checkmarki\checkmarkr\checkmarkc\checkmark$\checkmarkm\checkmark$\checkmark^\checkmark\\checkmarkc\checkmarki\checkmarkr\checkmarkc\checkmark$\checkmarki\checkmark$\checkmark^\checkmark\\checkmarkc\checkmarki\checkmarkr\checkmarkc\checkmark$\checkmarks\checkmark$\checkmark^\checkmark\\checkmarkc\checkmarki\checkmarkr\checkmarkc\checkmark$\checkmarks\checkmark$\checkmark^\checkmark\\checkmarkc\checkmarki\checkmarkr\checkmarkc\checkmark$\checkmarki\checkmark$\checkmark^\checkmark\\checkmarkc\checkmarki\checkmarkr\checkmarkc\checkmark$\checkmarko\checkmark$\checkmark^\checkmark\\checkmarkc\checkmarki\checkmarkr\checkmarkc\checkmark$\checkmarkn\checkmark$\checkmark^\checkmark\\checkmarkc\checkmarki\checkmarkr\checkmarkc\checkmark$\checkmark \checkmark$\checkmark^\checkmark\\checkmarkc\checkmarki\checkmarkr\checkmarkc\checkmark$\checkmarkà\checkmark$\checkmark^\checkmark\\checkmarkc\checkmarki\checkmarkr\checkmarkc\checkmark$\checkmark \checkmark$\checkmark^\checkmark\\checkmarkc\checkmarki\checkmarkr\checkmarkc\checkmark$\checkmarkc\checkmark$\checkmark^\checkmark\\checkmarkc\checkmarki\checkmarkr\checkmarkc\checkmark$\checkmarko\checkmark$\checkmark^\checkmark\\checkmarkc\checkmarki\checkmarkr\checkmarkc\checkmark$\checkmarku\checkmark$\checkmark^\checkmark\\checkmarkc\checkmarki\checkmarkr\checkmarkc\checkmark$\checkmarkr\checkmark$\checkmark^\checkmark\\checkmarkc\checkmarki\checkmarkr\checkmarkc\checkmark$\checkmarkr\checkmark$\checkmark^\checkmark\\checkmarkc\checkmarki\checkmarkr\checkmarkc\checkmark$\checkmarko\checkmark$\checkmark^\checkmark\\checkmarkc\checkmarki\checkmarkr\checkmarkc\checkmark$\checkmarki\checkmark$\checkmark^\checkmark\\checkmarkc\checkmarki\checkmarkr\checkmarkc\checkmark$\checkmarke\checkmark$\checkmark^\checkmark\\checkmarkc\checkmarki\checkmarkr\checkmarkc\checkmark$\checkmark \checkmark$\checkmark^\checkmark\\checkmarkc\checkmarki\checkmarkr\checkmarkc\checkmark$\checkmarkc\checkmark$\checkmark^\checkmark\\checkmarkc\checkmarki\checkmarkr\checkmarkc\checkmark$\checkmarkr\checkmark$\checkmark^\checkmark\\checkmarkc\checkmarki\checkmarkr\checkmarkc\checkmark$\checkmarka\checkmark$\checkmark^\checkmark\\checkmarkc\checkmarki\checkmarkr\checkmarkc\checkmark$\checkmarkn\checkmark$\checkmark^\checkmark\\checkmarkc\checkmarki\checkmarkr\checkmarkc\checkmark$\checkmarkt\checkmark$\checkmark^\checkmark\\checkmarkc\checkmarki\checkmarkr\checkmarkc\checkmark$\checkmarké\checkmark$\checkmark^\checkmark\\checkmarkc\checkmarki\checkmarkr\checkmarkc\checkmark$\checkmarke\checkmark$\checkmark^\checkmark\\checkmarkc\checkmarki\checkmarkr\checkmarkc\checkmark$\checkmark \checkmark$\checkmark^\checkmark\\checkmarkc\checkmarki\checkmarkr\checkmarkc\checkmark$\checkmarkH\checkmark$\checkmark^\checkmark\\checkmarkc\checkmarki\checkmarkr\checkmarkc\checkmark$\checkmarkT\checkmark$\checkmark^\checkmark\\checkmarkc\checkmarki\checkmarkr\checkmarkc\checkmark$\checkmarkD\checkmark$\checkmark^\checkmark\\checkmarkc\checkmarki\checkmarkr\checkmarkc\checkmark$\checkmark}\checkmark$\checkmark^\checkmark\\checkmarkc\checkmarki\checkmarkr\checkmarkc\checkmark$\checkmark
\checkmark$\checkmark^\checkmark\\checkmarkc\checkmarki\checkmarkr\checkmarkc\checkmark$\checkmarkL\checkmark$\checkmark^\checkmark\\checkmarkc\checkmarki\checkmarkr\checkmarkc\checkmark$\checkmarka\checkmark$\checkmark^\checkmark\\checkmarkc\checkmarki\checkmarkr\checkmarkc\checkmark$\checkmark \checkmark$\checkmark^\checkmark\\checkmarkc\checkmarki\checkmarkr\checkmarkc\checkmark$\checkmarkt\checkmark$\checkmark^\checkmark\\checkmarkc\checkmarki\checkmarkr\checkmarkc\checkmark$\checkmarkr\checkmark$\checkmark^\checkmark\\checkmarkc\checkmarki\checkmarkr\checkmarkc\checkmark$\checkmarka\checkmark$\checkmark^\checkmark\\checkmarkc\checkmarki\checkmarkr\checkmarkc\checkmark$\checkmarkn\checkmark$\checkmark^\checkmark\\checkmarkc\checkmarki\checkmarkr\checkmarkc\checkmark$\checkmarks\checkmark$\checkmark^\checkmark\\checkmarkc\checkmarki\checkmarkr\checkmarkc\checkmark$\checkmarkm\checkmark$\checkmark^\checkmark\\checkmarkc\checkmarki\checkmarkr\checkmarkc\checkmark$\checkmarki\checkmark$\checkmark^\checkmark\\checkmarkc\checkmarki\checkmarkr\checkmarkc\checkmark$\checkmarks\checkmark$\checkmark^\checkmark\\checkmarkc\checkmarki\checkmarkr\checkmarkc\checkmark$\checkmarks\checkmark$\checkmark^\checkmark\\checkmarkc\checkmarki\checkmarkr\checkmarkc\checkmark$\checkmarki\checkmark$\checkmark^\checkmark\\checkmarkc\checkmarki\checkmarkr\checkmarkc\checkmark$\checkmarko\checkmark$\checkmark^\checkmark\\checkmarkc\checkmarki\checkmarkr\checkmarkc\checkmark$\checkmarkn\checkmark$\checkmark^\checkmark\\checkmarkc\checkmarki\checkmarkr\checkmarkc\checkmark$\checkmark \checkmark$\checkmark^\checkmark\\checkmarkc\checkmarki\checkmarkr\checkmarkc\checkmark$\checkmarkp\checkmark$\checkmark^\checkmark\\checkmarkc\checkmarki\checkmarkr\checkmarkc\checkmark$\checkmarka\checkmark$\checkmark^\checkmark\\checkmarkc\checkmarki\checkmarkr\checkmarkc\checkmark$\checkmarkr\checkmark$\checkmark^\checkmark\\checkmarkc\checkmarki\checkmarkr\checkmarkc\checkmark$\checkmark \checkmark$\checkmark^\checkmark\\checkmarkc\checkmarki\checkmarkr\checkmarkc\checkmark$\checkmarkc\checkmark$\checkmark^\checkmark\\checkmarkc\checkmarki\checkmarkr\checkmarkc\checkmark$\checkmarko\checkmark$\checkmark^\checkmark\\checkmarkc\checkmarki\checkmarkr\checkmarkc\checkmark$\checkmarku\checkmark$\checkmark^\checkmark\\checkmarkc\checkmarki\checkmarkr\checkmarkc\checkmark$\checkmarkr\checkmark$\checkmark^\checkmark\\checkmarkc\checkmarki\checkmarkr\checkmarkc\checkmark$\checkmarkr\checkmark$\checkmark^\checkmark\\checkmarkc\checkmarki\checkmarkr\checkmarkc\checkmark$\checkmarko\checkmark$\checkmark^\checkmark\\checkmarkc\checkmarki\checkmarkr\checkmarkc\checkmark$\checkmarki\checkmark$\checkmark^\checkmark\\checkmarkc\checkmarki\checkmarkr\checkmarkc\checkmark$\checkmarke\checkmark$\checkmark^\checkmark\\checkmarkc\checkmarki\checkmarkr\checkmarkc\checkmark$\checkmark \checkmark$\checkmark^\checkmark\\checkmarkc\checkmarki\checkmarkr\checkmarkc\checkmark$\checkmarkc\checkmark$\checkmark^\checkmark\\checkmarkc\checkmarki\checkmarkr\checkmarkc\checkmark$\checkmarkr\checkmark$\checkmark^\checkmark\\checkmarkc\checkmarki\checkmarkr\checkmarkc\checkmark$\checkmarka\checkmark$\checkmark^\checkmark\\checkmarkc\checkmarki\checkmarkr\checkmarkc\checkmark$\checkmarkn\checkmark$\checkmark^\checkmark\\checkmarkc\checkmarki\checkmarkr\checkmarkc\checkmark$\checkmarkt\checkmark$\checkmark^\checkmark\\checkmarkc\checkmarki\checkmarkr\checkmarkc\checkmark$\checkmarké\checkmark$\checkmark^\checkmark\\checkmarkc\checkmarki\checkmarkr\checkmarkc\checkmark$\checkmarke\checkmark$\checkmark^\checkmark\\checkmarkc\checkmarki\checkmarkr\checkmarkc\checkmark$\checkmark \checkmark$\checkmark^\checkmark\\checkmarkc\checkmarki\checkmarkr\checkmarkc\checkmark$\checkmarkH\checkmark$\checkmark^\checkmark\\checkmarkc\checkmarki\checkmarkr\checkmarkc\checkmark$\checkmarkT\checkmark$\checkmark^\checkmark\\checkmarkc\checkmarki\checkmarkr\checkmarkc\checkmark$\checkmarkD\checkmark$\checkmark^\checkmark\\checkmarkc\checkmarki\checkmarkr\checkmarkc\checkmark$\checkmark \checkmark$\checkmark^\checkmark\\checkmarkc\checkmarki\checkmarkr\checkmarkc\checkmark$\checkmark5\checkmark$\checkmark^\checkmark\\checkmarkc\checkmarki\checkmarkr\checkmarkc\checkmark$\checkmarkM\checkmark$\checkmark^\checkmark\\checkmarkc\checkmarki\checkmarkr\checkmarkc\checkmark$\checkmark \checkmark$\checkmark^\checkmark\\checkmarkc\checkmarki\checkmarkr\checkmarkc\checkmark$\checkmarka\checkmark$\checkmark^\checkmark\\checkmarkc\checkmarki\checkmarkr\checkmarkc\checkmark$\checkmark \checkmark$\checkmark^\checkmark\\checkmarkc\checkmarki\checkmarkr\checkmarkc\checkmark$\checkmarké\checkmark$\checkmark^\checkmark\\checkmarkc\checkmarki\checkmarkr\checkmarkc\checkmark$\checkmarkt\checkmark$\checkmark^\checkmark\\checkmarkc\checkmarki\checkmarkr\checkmarkc\checkmark$\checkmarké\checkmark$\checkmark^\checkmark\\checkmarkc\checkmarki\checkmarkr\checkmarkc\checkmark$\checkmark \checkmark$\checkmark^\checkmark\\checkmarkc\checkmarki\checkmarkr\checkmarkc\checkmark$\checkmarkr\checkmark$\checkmark^\checkmark\\checkmarkc\checkmarki\checkmarkr\checkmarkc\checkmark$\checkmarke\checkmark$\checkmark^\checkmark\\checkmarkc\checkmarki\checkmarkr\checkmarkc\checkmark$\checkmarkt\checkmark$\checkmark^\checkmark\\checkmarkc\checkmarki\checkmarkr\checkmarkc\checkmark$\checkmarke\checkmark$\checkmark^\checkmark\\checkmarkc\checkmarki\checkmarkr\checkmarkc\checkmark$\checkmarkn\checkmark$\checkmark^\checkmark\\checkmarkc\checkmarki\checkmarkr\checkmarkc\checkmark$\checkmarku\checkmark$\checkmark^\checkmark\\checkmarkc\checkmarki\checkmarkr\checkmarkc\checkmark$\checkmarke\checkmark$\checkmark^\checkmark\\checkmarkc\checkmarki\checkmarkr\checkmarkc\checkmark$\checkmark \checkmark$\checkmark^\checkmark\\checkmarkc\checkmarki\checkmarkr\checkmarkc\checkmark$\checkmarkp\checkmark$\checkmark^\checkmark\\checkmarkc\checkmarki\checkmarkr\checkmarkc\checkmark$\checkmarko\checkmark$\checkmark^\checkmark\\checkmarkc\checkmarki\checkmarkr\checkmarkc\checkmark$\checkmarku\checkmark$\checkmark^\checkmark\\checkmarkc\checkmarki\checkmarkr\checkmarkc\checkmark$\checkmarkr\checkmark$\checkmark^\checkmark\\checkmarkc\checkmarki\checkmarkr\checkmarkc\checkmark$\checkmark \checkmark$\checkmark^\checkmark\\checkmarkc\checkmarki\checkmarkr\checkmarkc\checkmark$\checkmarkr\checkmark$\checkmark^\checkmark\\checkmarkc\checkmarki\checkmarkr\checkmarkc\checkmark$\checkmarké\checkmark$\checkmark^\checkmark\\checkmarkc\checkmarki\checkmarkr\checkmarkc\checkmark$\checkmarka\checkmark$\checkmark^\checkmark\\checkmarkc\checkmarki\checkmarkr\checkmarkc\checkmark$\checkmarkl\checkmark$\checkmark^\checkmark\\checkmarkc\checkmarki\checkmarkr\checkmarkc\checkmark$\checkmarki\checkmark$\checkmark^\checkmark\\checkmarkc\checkmarki\checkmarkr\checkmarkc\checkmark$\checkmarks\checkmark$\checkmark^\checkmark\\checkmarkc\checkmarki\checkmarkr\checkmarkc\checkmark$\checkmarke\checkmark$\checkmark^\checkmark\\checkmarkc\checkmarki\checkmarkr\checkmarkc\checkmark$\checkmarkr\checkmark$\checkmark^\checkmark\\checkmarkc\checkmarki\checkmarkr\checkmarkc\checkmark$\checkmark \checkmark$\checkmark^\checkmark\\checkmarkc\checkmarki\checkmarkr\checkmarkc\checkmark$\checkmarkl\checkmark$\checkmark^\checkmark\\checkmarkc\checkmarki\checkmarkr\checkmarkc\checkmark$\checkmarka\checkmark$\checkmark^\checkmark\\checkmarkc\checkmarki\checkmarkr\checkmarkc\checkmark$\checkmark \checkmark$\checkmark^\checkmark\\checkmarkc\checkmarki\checkmarkr\checkmarkc\checkmark$\checkmarkr\checkmark$\checkmark^\checkmark\\checkmarkc\checkmarki\checkmarkr\checkmarkc\checkmark$\checkmarké\checkmark$\checkmark^\checkmark\\checkmarkc\checkmarki\checkmarkr\checkmarkc\checkmark$\checkmarkd\checkmark$\checkmark^\checkmark\\checkmarkc\checkmarki\checkmarkr\checkmarkc\checkmark$\checkmarku\checkmark$\checkmark^\checkmark\\checkmarkc\checkmarki\checkmarkr\checkmarkc\checkmark$\checkmarkc\checkmark$\checkmark^\checkmark\\checkmarkc\checkmarki\checkmarkr\checkmarkc\checkmark$\checkmarkt\checkmark$\checkmark^\checkmark\\checkmarkc\checkmarki\checkmarkr\checkmarkc\checkmark$\checkmarki\checkmark$\checkmark^\checkmark\\checkmarkc\checkmarki\checkmarkr\checkmarkc\checkmark$\checkmarko\checkmark$\checkmark^\checkmark\\checkmarkc\checkmarki\checkmarkr\checkmarkc\checkmark$\checkmarkn\checkmark$\checkmark^\checkmark\\checkmarkc\checkmarki\checkmarkr\checkmarkc\checkmark$\checkmark \checkmark$\checkmark^\checkmark\\checkmarkc\checkmarki\checkmarkr\checkmarkc\checkmark$\checkmark1\checkmark$\checkmark^\checkmark\\checkmarkc\checkmarki\checkmarkr\checkmarkc\checkmark$\checkmark:\checkmark$\checkmark^\checkmark\\checkmarkc\checkmarki\checkmarkr\checkmarkc\checkmark$\checkmark5\checkmark$\checkmark^\checkmark\\checkmarkc\checkmarki\checkmarkr\checkmarkc\checkmark$\checkmark0\checkmark$\checkmark^\checkmark\\checkmarkc\checkmarki\checkmarkr\checkmarkc\checkmark$\checkmark \checkmark$\checkmark^\checkmark\\checkmarkc\checkmarki\checkmarkr\checkmarkc\checkmark$\checkmarke\checkmark$\checkmark^\checkmark\\checkmarkc\checkmarki\checkmarkr\checkmarkc\checkmark$\checkmarkn\checkmark$\checkmark^\checkmark\\checkmarkc\checkmarki\checkmarkr\checkmarkc\checkmark$\checkmark \checkmark$\checkmark^\checkmark\\checkmarkc\checkmarki\checkmarkr\checkmarkc\checkmark$\checkmarkd\checkmark$\checkmark^\checkmark\\checkmarkc\checkmarki\checkmarkr\checkmarkc\checkmark$\checkmarke\checkmark$\checkmark^\checkmark\\checkmarkc\checkmarki\checkmarkr\checkmarkc\checkmark$\checkmarku\checkmark$\checkmark^\checkmark\\checkmarkc\checkmarki\checkmarkr\checkmarkc\checkmark$\checkmarkx\checkmark$\checkmark^\checkmark\\checkmarkc\checkmarki\checkmarkr\checkmarkc\checkmark$\checkmark \checkmark$\checkmark^\checkmark\\checkmarkc\checkmarki\checkmarkr\checkmarkc\checkmark$\checkmarké\checkmark$\checkmark^\checkmark\\checkmarkc\checkmarki\checkmarkr\checkmarkc\checkmark$\checkmarkt\checkmark$\checkmark^\checkmark\\checkmarkc\checkmarki\checkmarkr\checkmarkc\checkmark$\checkmarka\checkmark$\checkmark^\checkmark\\checkmarkc\checkmarki\checkmarkr\checkmarkc\checkmark$\checkmarkg\checkmark$\checkmark^\checkmark\\checkmarkc\checkmarki\checkmarkr\checkmarkc\checkmark$\checkmarke\checkmark$\checkmark^\checkmark\\checkmarkc\checkmarki\checkmarkr\checkmarkc\checkmark$\checkmarks\checkmark$\checkmark^\checkmark\\checkmarkc\checkmarki\checkmarkr\checkmarkc\checkmark$\checkmark \checkmark$\checkmark^\checkmark\\checkmarkc\checkmarki\checkmarkr\checkmarkc\checkmark$\checkmark(\checkmark$\checkmark^\checkmark\\checkmarkc\checkmarki\checkmarkr\checkmarkc\checkmark$\checkmark1\checkmark$\checkmark^\checkmark\\checkmarkc\checkmarki\checkmarkr\checkmarkc\checkmark$\checkmark:\checkmark$\checkmark^\checkmark\\checkmarkc\checkmarki\checkmarkr\checkmarkc\checkmark$\checkmark1\checkmark$\checkmark^\checkmark\\checkmarkc\checkmarki\checkmarkr\checkmarkc\checkmark$\checkmark0\checkmark$\checkmark^\checkmark\\checkmarkc\checkmarki\checkmarkr\checkmarkc\checkmark$\checkmark \checkmark$\checkmark^\checkmark\\checkmarkc\checkmarki\checkmarkr\checkmarkc\checkmark$\checkmark+\checkmark$\checkmark^\checkmark\\checkmarkc\checkmarki\checkmarkr\checkmarkc\checkmark$\checkmark \checkmark$\checkmark^\checkmark\\checkmarkc\checkmarki\checkmarkr\checkmarkc\checkmark$\checkmark1\checkmark$\checkmark^\checkmark\\checkmarkc\checkmarki\checkmarkr\checkmarkc\checkmark$\checkmark:\checkmark$\checkmark^\checkmark\\checkmarkc\checkmarki\checkmarkr\checkmarkc\checkmark$\checkmark5\checkmark$\checkmark^\checkmark\\checkmarkc\checkmarki\checkmarkr\checkmarkc\checkmark$\checkmark \checkmark$\checkmark^\checkmark\\checkmarkc\checkmarki\checkmarkr\checkmarkc\checkmark$\checkmark=\checkmark$\checkmark^\checkmark\\checkmarkc\checkmarki\checkmarkr\checkmarkc\checkmark$\checkmark \checkmark$\checkmark^\checkmark\\checkmarkc\checkmarki\checkmarkr\checkmarkc\checkmark$\checkmark1\checkmark$\checkmark^\checkmark\\checkmarkc\checkmarki\checkmarkr\checkmarkc\checkmark$\checkmark:\checkmark$\checkmark^\checkmark\\checkmarkc\checkmarki\checkmarkr\checkmarkc\checkmark$\checkmark5\checkmark$\checkmark^\checkmark\\checkmarkc\checkmarki\checkmarkr\checkmarkc\checkmark$\checkmark0\checkmark$\checkmark^\checkmark\\checkmarkc\checkmarki\checkmarkr\checkmarkc\checkmark$\checkmark)\checkmark$\checkmark^\checkmark\\checkmarkc\checkmarki\checkmarkr\checkmarkc\checkmark$\checkmark.\checkmark$\checkmark^\checkmark\\checkmarkc\checkmarki\checkmarkr\checkmarkc\checkmark$\checkmark
\checkmark$\checkmark^\checkmark\\checkmarkc\checkmarki\checkmarkr\checkmarkc\checkmark$\checkmark
\checkmark$\checkmark^\checkmark\\checkmarkc\checkmarki\checkmarkr\checkmarkc\checkmark$\checkmark\\checkmark$\checkmark^\checkmark\\checkmarkc\checkmarki\checkmarkr\checkmarkc\checkmark$\checkmarks\checkmark$\checkmark^\checkmark\\checkmarkc\checkmarki\checkmarkr\checkmarkc\checkmark$\checkmarku\checkmark$\checkmark^\checkmark\\checkmarkc\checkmarki\checkmarkr\checkmarkc\checkmark$\checkmarkb\checkmark$\checkmark^\checkmark\\checkmarkc\checkmarki\checkmarkr\checkmarkc\checkmark$\checkmarks\checkmark$\checkmark^\checkmark\\checkmarkc\checkmarki\checkmarkr\checkmarkc\checkmark$\checkmarku\checkmark$\checkmark^\checkmark\\checkmarkc\checkmarki\checkmarkr\checkmarkc\checkmark$\checkmarkb\checkmark$\checkmark^\checkmark\\checkmarkc\checkmarki\checkmarkr\checkmarkc\checkmark$\checkmarks\checkmark$\checkmark^\checkmark\\checkmarkc\checkmarki\checkmarkr\checkmarkc\checkmark$\checkmarke\checkmark$\checkmark^\checkmark\\checkmarkc\checkmarki\checkmarkr\checkmarkc\checkmark$\checkmarkc\checkmark$\checkmark^\checkmark\\checkmarkc\checkmarki\checkmarkr\checkmarkc\checkmark$\checkmarkt\checkmark$\checkmark^\checkmark\\checkmarkc\checkmarki\checkmarkr\checkmarkc\checkmark$\checkmarki\checkmark$\checkmark^\checkmark\\checkmarkc\checkmarki\checkmarkr\checkmarkc\checkmark$\checkmarko\checkmark$\checkmark^\checkmark\\checkmarkc\checkmarki\checkmarkr\checkmarkc\checkmark$\checkmarkn\checkmark$\checkmark^\checkmark\\checkmarkc\checkmarki\checkmarkr\checkmarkc\checkmark$\checkmark{\checkmark$\checkmark^\checkmark\\checkmarkc\checkmarki\checkmarkr\checkmarkc\checkmark$\checkmarkÉ\checkmark$\checkmark^\checkmark\\checkmarkc\checkmarki\checkmarkr\checkmarkc\checkmark$\checkmarkt\checkmark$\checkmark^\checkmark\\checkmarkc\checkmarki\checkmarkr\checkmarkc\checkmark$\checkmarka\checkmark$\checkmark^\checkmark\\checkmarkc\checkmarki\checkmarkr\checkmarkc\checkmark$\checkmarkg\checkmark$\checkmark^\checkmark\\checkmarkc\checkmarki\checkmarkr\checkmarkc\checkmark$\checkmarke\checkmark$\checkmark^\checkmark\\checkmarkc\checkmarki\checkmarkr\checkmarkc\checkmark$\checkmark \checkmark$\checkmark^\checkmark\\checkmarkc\checkmarki\checkmarkr\checkmarkc\checkmark$\checkmark1\checkmark$\checkmark^\checkmark\\checkmarkc\checkmarki\checkmarkr\checkmarkc\checkmark$\checkmark \checkmark$\checkmark^\checkmark\\checkmarkc\checkmarki\checkmarkr\checkmarkc\checkmark$\checkmark:\checkmark$\checkmark^\checkmark\\checkmarkc\checkmarki\checkmarkr\checkmarkc\checkmark$\checkmark \checkmark$\checkmark^\checkmark\\checkmarkc\checkmarki\checkmarkr\checkmarkc\checkmark$\checkmarkR\checkmark$\checkmark^\checkmark\\checkmarkc\checkmarki\checkmarkr\checkmarkc\checkmark$\checkmarké\checkmark$\checkmark^\checkmark\\checkmarkc\checkmarki\checkmarkr\checkmarkc\checkmark$\checkmarkd\checkmark$\checkmark^\checkmark\\checkmarkc\checkmarki\checkmarkr\checkmarkc\checkmark$\checkmarku\checkmark$\checkmark^\checkmark\\checkmarkc\checkmarki\checkmarkr\checkmarkc\checkmark$\checkmarkc\checkmark$\checkmark^\checkmark\\checkmarkc\checkmarki\checkmarkr\checkmarkc\checkmark$\checkmarkt\checkmark$\checkmark^\checkmark\\checkmarkc\checkmarki\checkmarkr\checkmarkc\checkmark$\checkmarki\checkmark$\checkmark^\checkmark\\checkmarkc\checkmarki\checkmarkr\checkmarkc\checkmark$\checkmarko\checkmark$\checkmark^\checkmark\\checkmarkc\checkmarki\checkmarkr\checkmarkc\checkmark$\checkmarkn\checkmark$\checkmark^\checkmark\\checkmarkc\checkmarki\checkmarkr\checkmarkc\checkmark$\checkmark \checkmark$\checkmark^\checkmark\\checkmarkc\checkmarki\checkmarkr\checkmarkc\checkmark$\checkmark1\checkmark$\checkmark^\checkmark\\checkmarkc\checkmarki\checkmarkr\checkmarkc\checkmark$\checkmark:\checkmark$\checkmark^\checkmark\\checkmarkc\checkmarki\checkmarkr\checkmarkc\checkmark$\checkmark1\checkmark$\checkmark^\checkmark\\checkmarkc\checkmarki\checkmarkr\checkmarkc\checkmark$\checkmark0\checkmark$\checkmark^\checkmark\\checkmarkc\checkmarki\checkmarkr\checkmarkc\checkmark$\checkmark}\checkmark$\checkmark^\checkmark\\checkmarkc\checkmarki\checkmarkr\checkmarkc\checkmark$\checkmark
\checkmark$\checkmark^\checkmark\\checkmarkc\checkmarki\checkmarkr\checkmarkc\checkmark$\checkmark\\checkmark$\checkmark^\checkmark\\checkmarkc\checkmarki\checkmarkr\checkmarkc\checkmark$\checkmarkb\checkmark$\checkmark^\checkmark\\checkmarkc\checkmarki\checkmarkr\checkmarkc\checkmark$\checkmarke\checkmark$\checkmark^\checkmark\\checkmarkc\checkmarki\checkmarkr\checkmarkc\checkmark$\checkmarkg\checkmark$\checkmark^\checkmark\\checkmarkc\checkmarki\checkmarkr\checkmarkc\checkmark$\checkmarki\checkmark$\checkmark^\checkmark\\checkmarkc\checkmarki\checkmarkr\checkmarkc\checkmark$\checkmarkn\checkmark$\checkmark^\checkmark\\checkmarkc\checkmarki\checkmarkr\checkmarkc\checkmark$\checkmark{\checkmark$\checkmark^\checkmark\\checkmarkc\checkmarki\checkmarkr\checkmarkc\checkmark$\checkmarki\checkmark$\checkmark^\checkmark\\checkmarkc\checkmarki\checkmarkr\checkmarkc\checkmark$\checkmarkt\checkmark$\checkmark^\checkmark\\checkmarkc\checkmarki\checkmarkr\checkmarkc\checkmark$\checkmarke\checkmark$\checkmark^\checkmark\\checkmarkc\checkmarki\checkmarkr\checkmarkc\checkmark$\checkmarkm\checkmark$\checkmark^\checkmark\\checkmarkc\checkmarki\checkmarkr\checkmarkc\checkmark$\checkmarki\checkmark$\checkmark^\checkmark\\checkmarkc\checkmarki\checkmarkr\checkmarkc\checkmark$\checkmarkz\checkmark$\checkmark^\checkmark\\checkmarkc\checkmarki\checkmarkr\checkmarkc\checkmark$\checkmarke\checkmark$\checkmark^\checkmark\\checkmarkc\checkmarki\checkmarkr\checkmarkc\checkmark$\checkmark}\checkmark$\checkmark^\checkmark\\checkmarkc\checkmarki\checkmarkr\checkmarkc\checkmark$\checkmark
\checkmark$\checkmark^\checkmark\\checkmarkc\checkmarki\checkmarkr\checkmarkc\checkmark$\checkmark \checkmark$\checkmark^\checkmark\\checkmarkc\checkmarki\checkmarkr\checkmarkc\checkmark$\checkmark \checkmark$\checkmark^\checkmark\\checkmarkc\checkmarki\checkmarkr\checkmarkc\checkmark$\checkmark \checkmark$\checkmark^\checkmark\\checkmarkc\checkmarki\checkmarkr\checkmarkc\checkmark$\checkmark \checkmark$\checkmark^\checkmark\\checkmarkc\checkmarki\checkmarkr\checkmarkc\checkmark$\checkmark\\checkmark$\checkmark^\checkmark\\checkmarkc\checkmarki\checkmarkr\checkmarkc\checkmark$\checkmarki\checkmark$\checkmark^\checkmark\\checkmarkc\checkmarki\checkmarkr\checkmarkc\checkmark$\checkmarkt\checkmark$\checkmark^\checkmark\\checkmarkc\checkmarki\checkmarkr\checkmarkc\checkmark$\checkmarke\checkmark$\checkmark^\checkmark\\checkmarkc\checkmarki\checkmarkr\checkmarkc\checkmark$\checkmarkm\checkmark$\checkmark^\checkmark\\checkmarkc\checkmarki\checkmarkr\checkmarkc\checkmark$\checkmark \checkmark$\checkmark^\checkmark\\checkmarkc\checkmarki\checkmarkr\checkmarkc\checkmark$\checkmarkP\checkmark$\checkmark^\checkmark\\checkmarkc\checkmarki\checkmarkr\checkmarkc\checkmark$\checkmarko\checkmark$\checkmark^\checkmark\\checkmarkc\checkmarki\checkmarkr\checkmarkc\checkmark$\checkmarku\checkmark$\checkmark^\checkmark\\checkmarkc\checkmarki\checkmarkr\checkmarkc\checkmark$\checkmarkl\checkmark$\checkmark^\checkmark\\checkmarkc\checkmarki\checkmarkr\checkmarkc\checkmark$\checkmarki\checkmark$\checkmark^\checkmark\\checkmarkc\checkmarki\checkmarkr\checkmarkc\checkmark$\checkmarke\checkmark$\checkmark^\checkmark\\checkmarkc\checkmarki\checkmarkr\checkmarkc\checkmark$\checkmark \checkmark$\checkmark^\checkmark\\checkmarkc\checkmarki\checkmarkr\checkmarkc\checkmark$\checkmarkm\checkmark$\checkmark^\checkmark\\checkmarkc\checkmarki\checkmarkr\checkmarkc\checkmark$\checkmarke\checkmark$\checkmark^\checkmark\\checkmarkc\checkmarki\checkmarkr\checkmarkc\checkmark$\checkmarkn\checkmark$\checkmark^\checkmark\\checkmarkc\checkmarki\checkmarkr\checkmarkc\checkmark$\checkmarka\checkmark$\checkmark^\checkmark\\checkmarkc\checkmarki\checkmarkr\checkmarkc\checkmark$\checkmarkn\checkmark$\checkmark^\checkmark\\checkmarkc\checkmarki\checkmarkr\checkmarkc\checkmark$\checkmarkt\checkmark$\checkmark^\checkmark\\checkmarkc\checkmarki\checkmarkr\checkmarkc\checkmark$\checkmarke\checkmark$\checkmark^\checkmark\\checkmarkc\checkmarki\checkmarkr\checkmarkc\checkmark$\checkmark \checkmark$\checkmark^\checkmark\\checkmarkc\checkmarki\checkmarkr\checkmarkc\checkmark$\checkmark(\checkmark$\checkmark^\checkmark\\checkmarkc\checkmarki\checkmarkr\checkmarkc\checkmark$\checkmarkm\checkmark$\checkmark^\checkmark\\checkmarkc\checkmarki\checkmarkr\checkmarkc\checkmark$\checkmarko\checkmark$\checkmark^\checkmark\\checkmarkc\checkmarki\checkmarkr\checkmarkc\checkmark$\checkmarkt\checkmark$\checkmark^\checkmark\\checkmarkc\checkmarki\checkmarkr\checkmarkc\checkmark$\checkmarke\checkmark$\checkmark^\checkmark\\checkmarkc\checkmarki\checkmarkr\checkmarkc\checkmark$\checkmarku\checkmark$\checkmark^\checkmark\\checkmarkc\checkmarki\checkmarkr\checkmarkc\checkmark$\checkmarkr\checkmark$\checkmark^\checkmark\\checkmarkc\checkmarki\checkmarkr\checkmarkc\checkmark$\checkmark)\checkmark$\checkmark^\checkmark\\checkmarkc\checkmarki\checkmarkr\checkmarkc\checkmark$\checkmark \checkmark$\checkmark^\checkmark\\checkmarkc\checkmarki\checkmarkr\checkmarkc\checkmark$\checkmark:\checkmark$\checkmark^\checkmark\\checkmarkc\checkmarki\checkmarkr\checkmarkc\checkmark$\checkmark \checkmark$\checkmark^\checkmark\\checkmarkc\checkmarki\checkmarkr\checkmarkc\checkmark$\checkmark$\checkmark$\checkmark^\checkmark\\checkmarkc\checkmarki\checkmarkr\checkmarkc\checkmark$\checkmarkZ\checkmark$\checkmark^\checkmark\\checkmarkc\checkmarki\checkmarkr\checkmarkc\checkmark$\checkmark_\checkmark$\checkmark^\checkmark\\checkmarkc\checkmarki\checkmarkr\checkmarkc\checkmark$\checkmark1\checkmark$\checkmark^\checkmark\\checkmarkc\checkmarki\checkmarkr\checkmarkc\checkmark$\checkmark \checkmark$\checkmark^\checkmark\\checkmarkc\checkmarki\checkmarkr\checkmarkc\checkmark$\checkmark=\checkmark$\checkmark^\checkmark\\checkmarkc\checkmarki\checkmarkr\checkmarkc\checkmark$\checkmark \checkmark$\checkmark^\checkmark\\checkmarkc\checkmarki\checkmarkr\checkmarkc\checkmark$\checkmark1\checkmark$\checkmark^\checkmark\\checkmarkc\checkmarki\checkmarkr\checkmarkc\checkmark$\checkmark5\checkmark$\checkmark^\checkmark\\checkmarkc\checkmarki\checkmarkr\checkmarkc\checkmark$\checkmark$\checkmark$\checkmark^\checkmark\\checkmarkc\checkmarki\checkmarkr\checkmarkc\checkmark$\checkmark \checkmark$\checkmark^\checkmark\\checkmarkc\checkmarki\checkmarkr\checkmarkc\checkmark$\checkmarkd\checkmark$\checkmark^\checkmark\\checkmarkc\checkmarki\checkmarkr\checkmarkc\checkmark$\checkmarke\checkmark$\checkmark^\checkmark\\checkmarkc\checkmarki\checkmarkr\checkmarkc\checkmark$\checkmarkn\checkmark$\checkmark^\checkmark\\checkmarkc\checkmarki\checkmarkr\checkmarkc\checkmark$\checkmarkt\checkmark$\checkmark^\checkmark\\checkmarkc\checkmarki\checkmarkr\checkmarkc\checkmark$\checkmarks\checkmark$\checkmark^\checkmark\\checkmarkc\checkmarki\checkmarkr\checkmarkc\checkmark$\checkmark
\checkmark$\checkmark^\checkmark\\checkmarkc\checkmarki\checkmarkr\checkmarkc\checkmark$\checkmark \checkmark$\checkmark^\checkmark\\checkmarkc\checkmarki\checkmarkr\checkmarkc\checkmark$\checkmark \checkmark$\checkmark^\checkmark\\checkmarkc\checkmarki\checkmarkr\checkmarkc\checkmark$\checkmark \checkmark$\checkmark^\checkmark\\checkmarkc\checkmarki\checkmarkr\checkmarkc\checkmark$\checkmark \checkmark$\checkmark^\checkmark\\checkmarkc\checkmarki\checkmarkr\checkmarkc\checkmark$\checkmark\\checkmark$\checkmark^\checkmark\\checkmarkc\checkmarki\checkmarkr\checkmarkc\checkmark$\checkmarki\checkmark$\checkmark^\checkmark\\checkmarkc\checkmarki\checkmarkr\checkmarkc\checkmark$\checkmarkt\checkmark$\checkmark^\checkmark\\checkmarkc\checkmarki\checkmarkr\checkmarkc\checkmark$\checkmarke\checkmark$\checkmark^\checkmark\\checkmarkc\checkmarki\checkmarkr\checkmarkc\checkmark$\checkmarkm\checkmark$\checkmark^\checkmark\\checkmarkc\checkmarki\checkmarkr\checkmarkc\checkmark$\checkmark \checkmark$\checkmark^\checkmark\\checkmarkc\checkmarki\checkmarkr\checkmarkc\checkmark$\checkmarkP\checkmark$\checkmark^\checkmark\\checkmarkc\checkmarki\checkmarkr\checkmarkc\checkmark$\checkmarko\checkmark$\checkmark^\checkmark\\checkmarkc\checkmarki\checkmarkr\checkmarkc\checkmark$\checkmarku\checkmark$\checkmark^\checkmark\\checkmarkc\checkmarki\checkmarkr\checkmarkc\checkmark$\checkmarkl\checkmark$\checkmark^\checkmark\\checkmarkc\checkmarki\checkmarkr\checkmarkc\checkmark$\checkmarki\checkmark$\checkmark^\checkmark\\checkmarkc\checkmarki\checkmarkr\checkmarkc\checkmark$\checkmarke\checkmark$\checkmark^\checkmark\\checkmarkc\checkmarki\checkmarkr\checkmarkc\checkmark$\checkmark \checkmark$\checkmark^\checkmark\\checkmarkc\checkmarki\checkmarkr\checkmarkc\checkmark$\checkmarkm\checkmark$\checkmark^\checkmark\\checkmarkc\checkmarki\checkmarkr\checkmarkc\checkmark$\checkmarke\checkmark$\checkmark^\checkmark\\checkmarkc\checkmarki\checkmarkr\checkmarkc\checkmark$\checkmarkn\checkmark$\checkmark^\checkmark\\checkmarkc\checkmarki\checkmarkr\checkmarkc\checkmark$\checkmarké\checkmark$\checkmark^\checkmark\\checkmarkc\checkmarki\checkmarkr\checkmarkc\checkmark$\checkmarke\checkmark$\checkmark^\checkmark\\checkmarkc\checkmarki\checkmarkr\checkmarkc\checkmark$\checkmark \checkmark$\checkmark^\checkmark\\checkmarkc\checkmarki\checkmarkr\checkmarkc\checkmark$\checkmark:\checkmark$\checkmark^\checkmark\\checkmarkc\checkmarki\checkmarkr\checkmarkc\checkmark$\checkmark \checkmark$\checkmark^\checkmark\\checkmarkc\checkmarki\checkmarkr\checkmarkc\checkmark$\checkmark$\checkmark$\checkmark^\checkmark\\checkmarkc\checkmarki\checkmarkr\checkmarkc\checkmark$\checkmarkZ\checkmark$\checkmark^\checkmark\\checkmarkc\checkmarki\checkmarkr\checkmarkc\checkmark$\checkmark_\checkmark$\checkmark^\checkmark\\checkmarkc\checkmarki\checkmarkr\checkmarkc\checkmark$\checkmark2\checkmark$\checkmark^\checkmark\\checkmarkc\checkmarki\checkmarkr\checkmarkc\checkmark$\checkmark \checkmark$\checkmark^\checkmark\\checkmarkc\checkmarki\checkmarkr\checkmarkc\checkmark$\checkmark=\checkmark$\checkmark^\checkmark\\checkmarkc\checkmarki\checkmarkr\checkmarkc\checkmark$\checkmark \checkmark$\checkmark^\checkmark\\checkmarkc\checkmarki\checkmarkr\checkmarkc\checkmark$\checkmark1\checkmark$\checkmark^\checkmark\\checkmarkc\checkmarki\checkmarkr\checkmarkc\checkmark$\checkmark5\checkmark$\checkmark^\checkmark\\checkmarkc\checkmarki\checkmarkr\checkmarkc\checkmark$\checkmark0\checkmark$\checkmark^\checkmark\\checkmarkc\checkmarki\checkmarkr\checkmarkc\checkmark$\checkmark$\checkmark$\checkmark^\checkmark\\checkmarkc\checkmarki\checkmarkr\checkmarkc\checkmark$\checkmark \checkmark$\checkmark^\checkmark\\checkmarkc\checkmarki\checkmarkr\checkmarkc\checkmark$\checkmarkd\checkmark$\checkmark^\checkmark\\checkmarkc\checkmarki\checkmarkr\checkmarkc\checkmark$\checkmarke\checkmark$\checkmark^\checkmark\\checkmarkc\checkmarki\checkmarkr\checkmarkc\checkmark$\checkmarkn\checkmark$\checkmark^\checkmark\\checkmarkc\checkmarki\checkmarkr\checkmarkc\checkmark$\checkmarkt\checkmark$\checkmark^\checkmark\\checkmarkc\checkmarki\checkmarkr\checkmarkc\checkmark$\checkmarks\checkmark$\checkmark^\checkmark\\checkmarkc\checkmarki\checkmarkr\checkmarkc\checkmark$\checkmark
\checkmark$\checkmark^\checkmark\\checkmarkc\checkmarki\checkmarkr\checkmarkc\checkmark$\checkmark \checkmark$\checkmark^\checkmark\\checkmarkc\checkmarki\checkmarkr\checkmarkc\checkmark$\checkmark \checkmark$\checkmark^\checkmark\\checkmarkc\checkmarki\checkmarkr\checkmarkc\checkmark$\checkmark \checkmark$\checkmark^\checkmark\\checkmarkc\checkmarki\checkmarkr\checkmarkc\checkmark$\checkmark \checkmark$\checkmark^\checkmark\\checkmarkc\checkmarki\checkmarkr\checkmarkc\checkmark$\checkmark\\checkmark$\checkmark^\checkmark\\checkmarkc\checkmarki\checkmarkr\checkmarkc\checkmark$\checkmarki\checkmark$\checkmark^\checkmark\\checkmarkc\checkmarki\checkmarkr\checkmarkc\checkmark$\checkmarkt\checkmark$\checkmark^\checkmark\\checkmarkc\checkmarki\checkmarkr\checkmarkc\checkmark$\checkmarke\checkmark$\checkmark^\checkmark\\checkmarkc\checkmarki\checkmarkr\checkmarkc\checkmark$\checkmarkm\checkmark$\checkmark^\checkmark\\checkmarkc\checkmarki\checkmarkr\checkmarkc\checkmark$\checkmark \checkmark$\checkmark^\checkmark\\checkmarkc\checkmarki\checkmarkr\checkmarkc\checkmark$\checkmarkM\checkmark$\checkmark^\checkmark\\checkmarkc\checkmarki\checkmarkr\checkmarkc\checkmark$\checkmarko\checkmark$\checkmark^\checkmark\\checkmarkc\checkmarki\checkmarkr\checkmarkc\checkmark$\checkmarkd\checkmark$\checkmark^\checkmark\\checkmarkc\checkmarki\checkmarkr\checkmarkc\checkmark$\checkmarku\checkmark$\checkmark^\checkmark\\checkmarkc\checkmarki\checkmarkr\checkmarkc\checkmark$\checkmarkl\checkmark$\checkmark^\checkmark\\checkmarkc\checkmarki\checkmarkr\checkmarkc\checkmark$\checkmarke\checkmark$\checkmark^\checkmark\\checkmarkc\checkmarki\checkmarkr\checkmarkc\checkmark$\checkmark \checkmark$\checkmark^\checkmark\\checkmarkc\checkmarki\checkmarkr\checkmarkc\checkmark$\checkmarkd\checkmark$\checkmark^\checkmark\\checkmarkc\checkmarki\checkmarkr\checkmarkc\checkmark$\checkmarke\checkmark$\checkmark^\checkmark\\checkmarkc\checkmarki\checkmarkr\checkmarkc\checkmark$\checkmark \checkmark$\checkmark^\checkmark\\checkmarkc\checkmarki\checkmarkr\checkmarkc\checkmark$\checkmarkp\checkmark$\checkmark^\checkmark\\checkmarkc\checkmarki\checkmarkr\checkmarkc\checkmark$\checkmarka\checkmark$\checkmark^\checkmark\\checkmarkc\checkmarki\checkmarkr\checkmarkc\checkmark$\checkmarks\checkmark$\checkmark^\checkmark\\checkmarkc\checkmarki\checkmarkr\checkmarkc\checkmark$\checkmark \checkmark$\checkmark^\checkmark\\checkmarkc\checkmarki\checkmarkr\checkmarkc\checkmark$\checkmark$\checkmark$\checkmark^\checkmark\\checkmarkc\checkmarki\checkmarkr\checkmarkc\checkmark$\checkmarkm\checkmark$\checkmark^\checkmark\\checkmarkc\checkmarki\checkmarkr\checkmarkc\checkmark$\checkmark_\checkmark$\checkmark^\checkmark\\checkmarkc\checkmarki\checkmarkr\checkmarkc\checkmark$\checkmarkt\checkmark$\checkmark^\checkmark\\checkmarkc\checkmarki\checkmarkr\checkmarkc\checkmark$\checkmark \checkmark$\checkmark^\checkmark\\checkmarkc\checkmarki\checkmarkr\checkmarkc\checkmark$\checkmark=\checkmark$\checkmark^\checkmark\\checkmarkc\checkmarki\checkmarkr\checkmarkc\checkmark$\checkmark \checkmark$\checkmark^\checkmark\\checkmarkc\checkmarki\checkmarkr\checkmarkc\checkmark$\checkmark5\checkmark$\checkmark^\checkmark\\checkmarkc\checkmarki\checkmarkr\checkmarkc\checkmark$\checkmark$\checkmark$\checkmark^\checkmark\\checkmarkc\checkmarki\checkmarkr\checkmarkc\checkmark$\checkmark \checkmark$\checkmark^\checkmark\\checkmarkc\checkmarki\checkmarkr\checkmarkc\checkmark$\checkmarkm\checkmark$\checkmark^\checkmark\\checkmarkc\checkmarki\checkmarkr\checkmarkc\checkmark$\checkmarkm\checkmark$\checkmark^\checkmark\\checkmarkc\checkmarki\checkmarkr\checkmarkc\checkmark$\checkmark \checkmark$\checkmark^\checkmark\\checkmarkc\checkmarki\checkmarkr\checkmarkc\checkmark$\checkmark(\checkmark$\checkmark^\checkmark\\checkmarkc\checkmarki\checkmarkr\checkmarkc\checkmark$\checkmarkH\checkmark$\checkmark^\checkmark\\checkmarkc\checkmarki\checkmarkr\checkmarkc\checkmark$\checkmarkT\checkmark$\checkmark^\checkmark\\checkmarkc\checkmarki\checkmarkr\checkmarkc\checkmark$\checkmarkD\checkmark$\checkmark^\checkmark\\checkmarkc\checkmarki\checkmarkr\checkmarkc\checkmark$\checkmark \checkmark$\checkmark^\checkmark\\checkmarkc\checkmarki\checkmarkr\checkmarkc\checkmark$\checkmark5\checkmark$\checkmark^\checkmark\\checkmarkc\checkmarki\checkmarkr\checkmarkc\checkmark$\checkmarkM\checkmark$\checkmark^\checkmark\\checkmarkc\checkmarki\checkmarkr\checkmarkc\checkmark$\checkmark)\checkmark$\checkmark^\checkmark\\checkmarkc\checkmarki\checkmarkr\checkmarkc\checkmark$\checkmark,\checkmark$\checkmark^\checkmark\\checkmarkc\checkmarki\checkmarkr\checkmarkc\checkmark$\checkmark \checkmark$\checkmark^\checkmark\\checkmarkc\checkmarki\checkmarkr\checkmarkc\checkmark$\checkmarkr\checkmark$\checkmark^\checkmark\\checkmarkc\checkmarki\checkmarkr\checkmarkc\checkmark$\checkmarka\checkmark$\checkmark^\checkmark\\checkmarkc\checkmarki\checkmarkr\checkmarkc\checkmark$\checkmarkp\checkmark$\checkmark^\checkmark\\checkmarkc\checkmarki\checkmarkr\checkmarkc\checkmark$\checkmarkp\checkmark$\checkmark^\checkmark\\checkmarkc\checkmarki\checkmarkr\checkmarkc\checkmark$\checkmarko\checkmark$\checkmark^\checkmark\\checkmarkc\checkmarki\checkmarkr\checkmarkc\checkmark$\checkmarkr\checkmark$\checkmark^\checkmark\\checkmarkc\checkmarki\checkmarkr\checkmarkc\checkmark$\checkmarkt\checkmark$\checkmark^\checkmark\\checkmarkc\checkmarki\checkmarkr\checkmarkc\checkmark$\checkmark \checkmark$\checkmark^\checkmark\\checkmarkc\checkmarki\checkmarkr\checkmarkc\checkmark$\checkmark:\checkmark$\checkmark^\checkmark\\checkmarkc\checkmarki\checkmarkr\checkmarkc\checkmark$\checkmark \checkmark$\checkmark^\checkmark\\checkmarkc\checkmarki\checkmarkr\checkmarkc\checkmark$\checkmark$\checkmark$\checkmark^\checkmark\\checkmarkc\checkmarki\checkmarkr\checkmarkc\checkmark$\checkmarki\checkmark$\checkmark^\checkmark\\checkmarkc\checkmarki\checkmarkr\checkmarkc\checkmark$\checkmark_\checkmark$\checkmark^\checkmark\\checkmarkc\checkmarki\checkmarkr\checkmarkc\checkmark$\checkmark1\checkmark$\checkmark^\checkmark\\checkmarkc\checkmarki\checkmarkr\checkmarkc\checkmark$\checkmark \checkmark$\checkmark^\checkmark\\checkmarkc\checkmarki\checkmarkr\checkmarkc\checkmark$\checkmark=\checkmark$\checkmark^\checkmark\\checkmarkc\checkmarki\checkmarkr\checkmarkc\checkmark$\checkmark \checkmark$\checkmark^\checkmark\\checkmarkc\checkmarki\checkmarkr\checkmarkc\checkmark$\checkmarkZ\checkmark$\checkmark^\checkmark\\checkmarkc\checkmarki\checkmarkr\checkmarkc\checkmark$\checkmark_\checkmark$\checkmark^\checkmark\\checkmarkc\checkmarki\checkmarkr\checkmarkc\checkmark$\checkmark2\checkmark$\checkmark^\checkmark\\checkmarkc\checkmarki\checkmarkr\checkmarkc\checkmark$\checkmark/\checkmark$\checkmark^\checkmark\\checkmarkc\checkmarki\checkmarkr\checkmarkc\checkmark$\checkmarkZ\checkmark$\checkmark^\checkmark\\checkmarkc\checkmarki\checkmarkr\checkmarkc\checkmark$\checkmark_\checkmark$\checkmark^\checkmark\\checkmarkc\checkmarki\checkmarkr\checkmarkc\checkmark$\checkmark1\checkmark$\checkmark^\checkmark\\checkmarkc\checkmarki\checkmarkr\checkmarkc\checkmark$\checkmark \checkmark$\checkmark^\checkmark\\checkmarkc\checkmarki\checkmarkr\checkmarkc\checkmark$\checkmark=\checkmark$\checkmark^\checkmark\\checkmarkc\checkmarki\checkmarkr\checkmarkc\checkmark$\checkmark \checkmark$\checkmark^\checkmark\\checkmarkc\checkmarki\checkmarkr\checkmarkc\checkmark$\checkmark1\checkmark$\checkmark^\checkmark\\checkmarkc\checkmarki\checkmarkr\checkmarkc\checkmark$\checkmark5\checkmark$\checkmark^\checkmark\\checkmarkc\checkmarki\checkmarkr\checkmarkc\checkmark$\checkmark0\checkmark$\checkmark^\checkmark\\checkmarkc\checkmarki\checkmarkr\checkmarkc\checkmark$\checkmark/\checkmark$\checkmark^\checkmark\\checkmarkc\checkmarki\checkmarkr\checkmarkc\checkmark$\checkmark1\checkmark$\checkmark^\checkmark\\checkmarkc\checkmarki\checkmarkr\checkmarkc\checkmark$\checkmark5\checkmark$\checkmark^\checkmark\\checkmarkc\checkmarki\checkmarkr\checkmarkc\checkmark$\checkmark \checkmark$\checkmark^\checkmark\\checkmarkc\checkmarki\checkmarkr\checkmarkc\checkmark$\checkmark=\checkmark$\checkmark^\checkmark\\checkmarkc\checkmarki\checkmarkr\checkmarkc\checkmark$\checkmark \checkmark$\checkmark^\checkmark\\checkmarkc\checkmarki\checkmarkr\checkmarkc\checkmark$\checkmark1\checkmark$\checkmark^\checkmark\\checkmarkc\checkmarki\checkmarkr\checkmarkc\checkmark$\checkmark0\checkmark$\checkmark^\checkmark\\checkmarkc\checkmarki\checkmarkr\checkmarkc\checkmark$\checkmark$\checkmark$\checkmark^\checkmark\\checkmarkc\checkmarki\checkmarkr\checkmarkc\checkmark$\checkmark
\checkmark$\checkmark^\checkmark\\checkmarkc\checkmarki\checkmarkr\checkmarkc\checkmark$\checkmark\\checkmark$\checkmark^\checkmark\\checkmarkc\checkmarki\checkmarkr\checkmarkc\checkmark$\checkmarke\checkmark$\checkmark^\checkmark\\checkmarkc\checkmarki\checkmarkr\checkmarkc\checkmark$\checkmarkn\checkmark$\checkmark^\checkmark\\checkmarkc\checkmarki\checkmarkr\checkmarkc\checkmark$\checkmarkd\checkmark$\checkmark^\checkmark\\checkmarkc\checkmarki\checkmarkr\checkmarkc\checkmark$\checkmark{\checkmark$\checkmark^\checkmark\\checkmarkc\checkmarki\checkmarkr\checkmarkc\checkmark$\checkmarki\checkmark$\checkmark^\checkmark\\checkmarkc\checkmarki\checkmarkr\checkmarkc\checkmark$\checkmarkt\checkmark$\checkmark^\checkmark\\checkmarkc\checkmarki\checkmarkr\checkmarkc\checkmark$\checkmarke\checkmark$\checkmark^\checkmark\\checkmarkc\checkmarki\checkmarkr\checkmarkc\checkmark$\checkmarkm\checkmark$\checkmark^\checkmark\\checkmarkc\checkmarki\checkmarkr\checkmarkc\checkmark$\checkmarki\checkmark$\checkmark^\checkmark\\checkmarkc\checkmarki\checkmarkr\checkmarkc\checkmark$\checkmarkz\checkmark$\checkmark^\checkmark\\checkmarkc\checkmarki\checkmarkr\checkmarkc\checkmark$\checkmarke\checkmark$\checkmark^\checkmark\\checkmarkc\checkmarki\checkmarkr\checkmarkc\checkmark$\checkmark}\checkmark$\checkmark^\checkmark\\checkmarkc\checkmarki\checkmarkr\checkmarkc\checkmark$\checkmark
\checkmark$\checkmark^\checkmark\\checkmarkc\checkmarki\checkmarkr\checkmarkc\checkmark$\checkmark
\checkmark$\checkmark^\checkmark\\checkmarkc\checkmarki\checkmarkr\checkmarkc\checkmark$\checkmark\\checkmark$\checkmark^\checkmark\\checkmarkc\checkmarki\checkmarkr\checkmarkc\checkmark$\checkmarks\checkmark$\checkmark^\checkmark\\checkmarkc\checkmarki\checkmarkr\checkmarkc\checkmark$\checkmarku\checkmark$\checkmark^\checkmark\\checkmarkc\checkmarki\checkmarkr\checkmarkc\checkmark$\checkmarkb\checkmark$\checkmark^\checkmark\\checkmarkc\checkmarki\checkmarkr\checkmarkc\checkmark$\checkmarks\checkmark$\checkmark^\checkmark\\checkmarkc\checkmarki\checkmarkr\checkmarkc\checkmark$\checkmarku\checkmark$\checkmark^\checkmark\\checkmarkc\checkmarki\checkmarkr\checkmarkc\checkmark$\checkmarkb\checkmark$\checkmark^\checkmark\\checkmarkc\checkmarki\checkmarkr\checkmarkc\checkmark$\checkmarks\checkmark$\checkmark^\checkmark\\checkmarkc\checkmarki\checkmarkr\checkmarkc\checkmark$\checkmarke\checkmark$\checkmark^\checkmark\\checkmarkc\checkmarki\checkmarkr\checkmarkc\checkmark$\checkmarkc\checkmark$\checkmark^\checkmark\\checkmarkc\checkmarki\checkmarkr\checkmarkc\checkmark$\checkmarkt\checkmark$\checkmark^\checkmark\\checkmarkc\checkmarki\checkmarkr\checkmarkc\checkmark$\checkmarki\checkmark$\checkmark^\checkmark\\checkmarkc\checkmarki\checkmarkr\checkmarkc\checkmark$\checkmarko\checkmark$\checkmark^\checkmark\\checkmarkc\checkmarki\checkmarkr\checkmarkc\checkmark$\checkmarkn\checkmark$\checkmark^\checkmark\\checkmarkc\checkmarki\checkmarkr\checkmarkc\checkmark$\checkmark{\checkmark$\checkmark^\checkmark\\checkmarkc\checkmarki\checkmarkr\checkmarkc\checkmark$\checkmarkÉ\checkmark$\checkmark^\checkmark\\checkmarkc\checkmarki\checkmarkr\checkmarkc\checkmark$\checkmarkt\checkmark$\checkmark^\checkmark\\checkmarkc\checkmarki\checkmarkr\checkmarkc\checkmark$\checkmarka\checkmark$\checkmark^\checkmark\\checkmarkc\checkmarki\checkmarkr\checkmarkc\checkmark$\checkmarkg\checkmark$\checkmark^\checkmark\\checkmarkc\checkmarki\checkmarkr\checkmarkc\checkmark$\checkmarke\checkmark$\checkmark^\checkmark\\checkmarkc\checkmarki\checkmarkr\checkmarkc\checkmark$\checkmark \checkmark$\checkmark^\checkmark\\checkmarkc\checkmarki\checkmarkr\checkmarkc\checkmark$\checkmark2\checkmark$\checkmark^\checkmark\\checkmarkc\checkmarki\checkmarkr\checkmarkc\checkmark$\checkmark \checkmark$\checkmark^\checkmark\\checkmarkc\checkmarki\checkmarkr\checkmarkc\checkmark$\checkmark:\checkmark$\checkmark^\checkmark\\checkmarkc\checkmarki\checkmarkr\checkmarkc\checkmark$\checkmark \checkmark$\checkmark^\checkmark\\checkmarkc\checkmarki\checkmarkr\checkmarkc\checkmark$\checkmarkR\checkmark$\checkmark^\checkmark\\checkmarkc\checkmarki\checkmarkr\checkmarkc\checkmark$\checkmarké\checkmark$\checkmark^\checkmark\\checkmarkc\checkmarki\checkmarkr\checkmarkc\checkmark$\checkmarkd\checkmark$\checkmark^\checkmark\\checkmarkc\checkmarki\checkmarkr\checkmarkc\checkmark$\checkmarku\checkmark$\checkmark^\checkmark\\checkmarkc\checkmarki\checkmarkr\checkmarkc\checkmark$\checkmarkc\checkmark$\checkmark^\checkmark\\checkmarkc\checkmarki\checkmarkr\checkmarkc\checkmark$\checkmarkt\checkmark$\checkmark^\checkmark\\checkmarkc\checkmarki\checkmarkr\checkmarkc\checkmark$\checkmarki\checkmark$\checkmark^\checkmark\\checkmarkc\checkmarki\checkmarkr\checkmarkc\checkmark$\checkmarko\checkmark$\checkmark^\checkmark\\checkmarkc\checkmarki\checkmarkr\checkmarkc\checkmark$\checkmarkn\checkmark$\checkmark^\checkmark\\checkmarkc\checkmarki\checkmarkr\checkmarkc\checkmark$\checkmark \checkmark$\checkmark^\checkmark\\checkmarkc\checkmarki\checkmarkr\checkmarkc\checkmark$\checkmark1\checkmark$\checkmark^\checkmark\\checkmarkc\checkmarki\checkmarkr\checkmarkc\checkmark$\checkmark:\checkmark$\checkmark^\checkmark\\checkmarkc\checkmarki\checkmarkr\checkmarkc\checkmark$\checkmark5\checkmark$\checkmark^\checkmark\\checkmarkc\checkmarki\checkmarkr\checkmarkc\checkmark$\checkmark}\checkmark$\checkmark^\checkmark\\checkmarkc\checkmarki\checkmarkr\checkmarkc\checkmark$\checkmark
\checkmark$\checkmark^\checkmark\\checkmarkc\checkmarki\checkmarkr\checkmarkc\checkmark$\checkmark\\checkmark$\checkmark^\checkmark\\checkmarkc\checkmarki\checkmarkr\checkmarkc\checkmark$\checkmarkb\checkmark$\checkmark^\checkmark\\checkmarkc\checkmarki\checkmarkr\checkmarkc\checkmark$\checkmarke\checkmark$\checkmark^\checkmark\\checkmarkc\checkmarki\checkmarkr\checkmarkc\checkmark$\checkmarkg\checkmark$\checkmark^\checkmark\\checkmarkc\checkmarki\checkmarkr\checkmarkc\checkmark$\checkmarki\checkmark$\checkmark^\checkmark\\checkmarkc\checkmarki\checkmarkr\checkmarkc\checkmark$\checkmarkn\checkmark$\checkmark^\checkmark\\checkmarkc\checkmarki\checkmarkr\checkmarkc\checkmark$\checkmark{\checkmark$\checkmark^\checkmark\\checkmarkc\checkmarki\checkmarkr\checkmarkc\checkmark$\checkmarki\checkmark$\checkmark^\checkmark\\checkmarkc\checkmarki\checkmarkr\checkmarkc\checkmark$\checkmarkt\checkmark$\checkmark^\checkmark\\checkmarkc\checkmarki\checkmarkr\checkmarkc\checkmark$\checkmarke\checkmark$\checkmark^\checkmark\\checkmarkc\checkmarki\checkmarkr\checkmarkc\checkmark$\checkmarkm\checkmark$\checkmark^\checkmark\\checkmarkc\checkmarki\checkmarkr\checkmarkc\checkmark$\checkmarki\checkmark$\checkmark^\checkmark\\checkmarkc\checkmarki\checkmarkr\checkmarkc\checkmark$\checkmarkz\checkmark$\checkmark^\checkmark\\checkmarkc\checkmarki\checkmarkr\checkmarkc\checkmark$\checkmarke\checkmark$\checkmark^\checkmark\\checkmarkc\checkmarki\checkmarkr\checkmarkc\checkmark$\checkmark}\checkmark$\checkmark^\checkmark\\checkmarkc\checkmarki\checkmarkr\checkmarkc\checkmark$\checkmark
\checkmark$\checkmark^\checkmark\\checkmarkc\checkmarki\checkmarkr\checkmarkc\checkmark$\checkmark \checkmark$\checkmark^\checkmark\\checkmarkc\checkmarki\checkmarkr\checkmarkc\checkmark$\checkmark \checkmark$\checkmark^\checkmark\\checkmarkc\checkmarki\checkmarkr\checkmarkc\checkmark$\checkmark \checkmark$\checkmark^\checkmark\\checkmarkc\checkmarki\checkmarkr\checkmarkc\checkmark$\checkmark \checkmark$\checkmark^\checkmark\\checkmarkc\checkmarki\checkmarkr\checkmarkc\checkmark$\checkmark\\checkmark$\checkmark^\checkmark\\checkmarkc\checkmarki\checkmarkr\checkmarkc\checkmark$\checkmarki\checkmark$\checkmark^\checkmark\\checkmarkc\checkmarki\checkmarkr\checkmarkc\checkmark$\checkmarkt\checkmark$\checkmark^\checkmark\\checkmarkc\checkmarki\checkmarkr\checkmarkc\checkmark$\checkmarke\checkmark$\checkmark^\checkmark\\checkmarkc\checkmarki\checkmarkr\checkmarkc\checkmark$\checkmarkm\checkmark$\checkmark^\checkmark\\checkmarkc\checkmarki\checkmarkr\checkmarkc\checkmark$\checkmark \checkmark$\checkmark^\checkmark\\checkmarkc\checkmarki\checkmarkr\checkmarkc\checkmark$\checkmarkP\checkmark$\checkmark^\checkmark\\checkmarkc\checkmarki\checkmarkr\checkmarkc\checkmark$\checkmarko\checkmark$\checkmark^\checkmark\\checkmarkc\checkmarki\checkmarkr\checkmarkc\checkmark$\checkmarku\checkmark$\checkmark^\checkmark\\checkmarkc\checkmarki\checkmarkr\checkmarkc\checkmark$\checkmarkl\checkmark$\checkmark^\checkmark\\checkmarkc\checkmarki\checkmarkr\checkmarkc\checkmark$\checkmarki\checkmark$\checkmark^\checkmark\\checkmarkc\checkmarki\checkmarkr\checkmarkc\checkmark$\checkmarke\checkmark$\checkmark^\checkmark\\checkmarkc\checkmarki\checkmarkr\checkmarkc\checkmark$\checkmark \checkmark$\checkmark^\checkmark\\checkmarkc\checkmarki\checkmarkr\checkmarkc\checkmark$\checkmarkm\checkmark$\checkmark^\checkmark\\checkmarkc\checkmarki\checkmarkr\checkmarkc\checkmark$\checkmarke\checkmark$\checkmark^\checkmark\\checkmarkc\checkmarki\checkmarkr\checkmarkc\checkmark$\checkmarkn\checkmark$\checkmark^\checkmark\\checkmarkc\checkmarki\checkmarkr\checkmarkc\checkmark$\checkmarka\checkmark$\checkmark^\checkmark\\checkmarkc\checkmarki\checkmarkr\checkmarkc\checkmark$\checkmarkn\checkmark$\checkmark^\checkmark\\checkmarkc\checkmarki\checkmarkr\checkmarkc\checkmark$\checkmarkt\checkmark$\checkmark^\checkmark\\checkmarkc\checkmarki\checkmarkr\checkmarkc\checkmark$\checkmarke\checkmark$\checkmark^\checkmark\\checkmarkc\checkmarki\checkmarkr\checkmarkc\checkmark$\checkmark \checkmark$\checkmark^\checkmark\\checkmarkc\checkmarki\checkmarkr\checkmarkc\checkmark$\checkmark:\checkmark$\checkmark^\checkmark\\checkmarkc\checkmarki\checkmarkr\checkmarkc\checkmark$\checkmark \checkmark$\checkmark^\checkmark\\checkmarkc\checkmarki\checkmarkr\checkmarkc\checkmark$\checkmark$\checkmark$\checkmark^\checkmark\\checkmarkc\checkmarki\checkmarkr\checkmarkc\checkmark$\checkmarkZ\checkmark$\checkmark^\checkmark\\checkmarkc\checkmarki\checkmarkr\checkmarkc\checkmark$\checkmark_\checkmark$\checkmark^\checkmark\\checkmarkc\checkmarki\checkmarkr\checkmarkc\checkmark$\checkmark3\checkmark$\checkmark^\checkmark\\checkmarkc\checkmarki\checkmarkr\checkmarkc\checkmark$\checkmark \checkmark$\checkmark^\checkmark\\checkmarkc\checkmarki\checkmarkr\checkmarkc\checkmark$\checkmark=\checkmark$\checkmark^\checkmark\\checkmarkc\checkmarki\checkmarkr\checkmarkc\checkmark$\checkmark \checkmark$\checkmark^\checkmark\\checkmarkc\checkmarki\checkmarkr\checkmarkc\checkmark$\checkmark2\checkmark$\checkmark^\checkmark\\checkmarkc\checkmarki\checkmarkr\checkmarkc\checkmark$\checkmark0\checkmark$\checkmark^\checkmark\\checkmarkc\checkmarki\checkmarkr\checkmarkc\checkmark$\checkmark$\checkmark$\checkmark^\checkmark\\checkmarkc\checkmarki\checkmarkr\checkmarkc\checkmark$\checkmark \checkmark$\checkmark^\checkmark\\checkmarkc\checkmarki\checkmarkr\checkmarkc\checkmark$\checkmarkd\checkmark$\checkmark^\checkmark\\checkmarkc\checkmarki\checkmarkr\checkmarkc\checkmark$\checkmarke\checkmark$\checkmark^\checkmark\\checkmarkc\checkmarki\checkmarkr\checkmarkc\checkmark$\checkmarkn\checkmark$\checkmark^\checkmark\\checkmarkc\checkmarki\checkmarkr\checkmarkc\checkmark$\checkmarkt\checkmark$\checkmark^\checkmark\\checkmarkc\checkmarki\checkmarkr\checkmarkc\checkmark$\checkmarks\checkmark$\checkmark^\checkmark\\checkmarkc\checkmarki\checkmarkr\checkmarkc\checkmark$\checkmark
\checkmark$\checkmark^\checkmark\\checkmarkc\checkmarki\checkmarkr\checkmarkc\checkmark$\checkmark \checkmark$\checkmark^\checkmark\\checkmarkc\checkmarki\checkmarkr\checkmarkc\checkmark$\checkmark \checkmark$\checkmark^\checkmark\\checkmarkc\checkmarki\checkmarkr\checkmarkc\checkmark$\checkmark \checkmark$\checkmark^\checkmark\\checkmarkc\checkmarki\checkmarkr\checkmarkc\checkmark$\checkmark \checkmark$\checkmark^\checkmark\\checkmarkc\checkmarki\checkmarkr\checkmarkc\checkmark$\checkmark\\checkmark$\checkmark^\checkmark\\checkmarkc\checkmarki\checkmarkr\checkmarkc\checkmark$\checkmarki\checkmark$\checkmark^\checkmark\\checkmarkc\checkmarki\checkmarkr\checkmarkc\checkmark$\checkmarkt\checkmark$\checkmark^\checkmark\\checkmarkc\checkmarki\checkmarkr\checkmarkc\checkmark$\checkmarke\checkmark$\checkmark^\checkmark\\checkmarkc\checkmarki\checkmarkr\checkmarkc\checkmark$\checkmarkm\checkmark$\checkmark^\checkmark\\checkmarkc\checkmarki\checkmarkr\checkmarkc\checkmark$\checkmark \checkmark$\checkmark^\checkmark\\checkmarkc\checkmarki\checkmarkr\checkmarkc\checkmark$\checkmarkP\checkmark$\checkmark^\checkmark\\checkmarkc\checkmarki\checkmarkr\checkmarkc\checkmark$\checkmarko\checkmark$\checkmark^\checkmark\\checkmarkc\checkmarki\checkmarkr\checkmarkc\checkmark$\checkmarku\checkmark$\checkmark^\checkmark\\checkmarkc\checkmarki\checkmarkr\checkmarkc\checkmark$\checkmarkl\checkmark$\checkmark^\checkmark\\checkmarkc\checkmarki\checkmarkr\checkmarkc\checkmark$\checkmarki\checkmark$\checkmark^\checkmark\\checkmarkc\checkmarki\checkmarkr\checkmarkc\checkmark$\checkmarke\checkmark$\checkmark^\checkmark\\checkmarkc\checkmarki\checkmarkr\checkmarkc\checkmark$\checkmark \checkmark$\checkmark^\checkmark\\checkmarkc\checkmarki\checkmarkr\checkmarkc\checkmark$\checkmarkm\checkmark$\checkmark^\checkmark\\checkmarkc\checkmarki\checkmarkr\checkmarkc\checkmark$\checkmarke\checkmark$\checkmark^\checkmark\\checkmarkc\checkmarki\checkmarkr\checkmarkc\checkmark$\checkmarkn\checkmark$\checkmark^\checkmark\\checkmarkc\checkmarki\checkmarkr\checkmarkc\checkmark$\checkmarké\checkmark$\checkmark^\checkmark\\checkmarkc\checkmarki\checkmarkr\checkmarkc\checkmark$\checkmarke\checkmark$\checkmark^\checkmark\\checkmarkc\checkmarki\checkmarkr\checkmarkc\checkmark$\checkmark \checkmark$\checkmark^\checkmark\\checkmarkc\checkmarki\checkmarkr\checkmarkc\checkmark$\checkmark(\checkmark$\checkmark^\checkmark\\checkmarkc\checkmarki\checkmarkr\checkmarkc\checkmark$\checkmarka\checkmark$\checkmark^\checkmark\\checkmarkc\checkmarki\checkmarkr\checkmarkc\checkmark$\checkmarkr\checkmark$\checkmark^\checkmark\\checkmarkc\checkmarki\checkmarkr\checkmarkc\checkmark$\checkmarkb\checkmark$\checkmark^\checkmark\\checkmarkc\checkmarki\checkmarkr\checkmarkc\checkmark$\checkmarkr\checkmark$\checkmark^\checkmark\\checkmarkc\checkmarki\checkmarkr\checkmarkc\checkmark$\checkmarke\checkmark$\checkmark^\checkmark\\checkmarkc\checkmarki\checkmarkr\checkmarkc\checkmark$\checkmark \checkmark$\checkmark^\checkmark\\checkmarkc\checkmarki\checkmarkr\checkmarkc\checkmark$\checkmarkA\checkmark$\checkmark^\checkmark\\checkmarkc\checkmarki\checkmarkr\checkmarkc\checkmark$\checkmarkx\checkmark$\checkmark^\checkmark\\checkmarkc\checkmarki\checkmarkr\checkmarkc\checkmark$\checkmarke\checkmark$\checkmark^\checkmark\\checkmarkc\checkmarki\checkmarkr\checkmarkc\checkmark$\checkmark \checkmark$\checkmark^\checkmark\\checkmarkc\checkmarki\checkmarkr\checkmarkc\checkmark$\checkmarkA\checkmark$\checkmark^\checkmark\\checkmarkc\checkmarki\checkmarkr\checkmarkc\checkmark$\checkmark)\checkmark$\checkmark^\checkmark\\checkmarkc\checkmarki\checkmarkr\checkmarkc\checkmark$\checkmark \checkmark$\checkmark^\checkmark\\checkmarkc\checkmarki\checkmarkr\checkmarkc\checkmark$\checkmark:\checkmark$\checkmark^\checkmark\\checkmarkc\checkmarki\checkmarkr\checkmarkc\checkmark$\checkmark \checkmark$\checkmark^\checkmark\\checkmarkc\checkmarki\checkmarkr\checkmarkc\checkmark$\checkmark$\checkmark$\checkmark^\checkmark\\checkmarkc\checkmarki\checkmarkr\checkmarkc\checkmark$\checkmarkZ\checkmark$\checkmark^\checkmark\\checkmarkc\checkmarki\checkmarkr\checkmarkc\checkmark$\checkmark_\checkmark$\checkmark^\checkmark\\checkmarkc\checkmarki\checkmarkr\checkmarkc\checkmark$\checkmark4\checkmark$\checkmark^\checkmark\\checkmarkc\checkmarki\checkmarkr\checkmarkc\checkmark$\checkmark \checkmark$\checkmark^\checkmark\\checkmarkc\checkmarki\checkmarkr\checkmarkc\checkmark$\checkmark=\checkmark$\checkmark^\checkmark\\checkmarkc\checkmarki\checkmarkr\checkmarkc\checkmark$\checkmark \checkmark$\checkmark^\checkmark\\checkmarkc\checkmarki\checkmarkr\checkmarkc\checkmark$\checkmark1\checkmark$\checkmark^\checkmark\\checkmarkc\checkmarki\checkmarkr\checkmarkc\checkmark$\checkmark0\checkmark$\checkmark^\checkmark\\checkmarkc\checkmarki\checkmarkr\checkmarkc\checkmark$\checkmark0\checkmark$\checkmark^\checkmark\\checkmarkc\checkmarki\checkmarkr\checkmarkc\checkmark$\checkmark$\checkmark$\checkmark^\checkmark\\checkmarkc\checkmarki\checkmarkr\checkmarkc\checkmark$\checkmark \checkmark$\checkmark^\checkmark\\checkmarkc\checkmarki\checkmarkr\checkmarkc\checkmark$\checkmarkd\checkmark$\checkmark^\checkmark\\checkmarkc\checkmarki\checkmarkr\checkmarkc\checkmark$\checkmarke\checkmark$\checkmark^\checkmark\\checkmarkc\checkmarki\checkmarkr\checkmarkc\checkmark$\checkmarkn\checkmark$\checkmark^\checkmark\\checkmarkc\checkmarki\checkmarkr\checkmarkc\checkmark$\checkmarkt\checkmark$\checkmark^\checkmark\\checkmarkc\checkmarki\checkmarkr\checkmarkc\checkmark$\checkmarks\checkmark$\checkmark^\checkmark\\checkmarkc\checkmarki\checkmarkr\checkmarkc\checkmark$\checkmark
\checkmark$\checkmark^\checkmark\\checkmarkc\checkmarki\checkmarkr\checkmarkc\checkmark$\checkmark \checkmark$\checkmark^\checkmark\\checkmarkc\checkmarki\checkmarkr\checkmarkc\checkmark$\checkmark \checkmark$\checkmark^\checkmark\\checkmarkc\checkmarki\checkmarkr\checkmarkc\checkmark$\checkmark \checkmark$\checkmark^\checkmark\\checkmarkc\checkmarki\checkmarkr\checkmarkc\checkmark$\checkmark \checkmark$\checkmark^\checkmark\\checkmarkc\checkmarki\checkmarkr\checkmarkc\checkmark$\checkmark\\checkmark$\checkmark^\checkmark\\checkmarkc\checkmarki\checkmarkr\checkmarkc\checkmark$\checkmarki\checkmark$\checkmark^\checkmark\\checkmarkc\checkmarki\checkmarkr\checkmarkc\checkmark$\checkmarkt\checkmark$\checkmark^\checkmark\\checkmarkc\checkmarki\checkmarkr\checkmarkc\checkmark$\checkmarke\checkmark$\checkmark^\checkmark\\checkmarkc\checkmarki\checkmarkr\checkmarkc\checkmark$\checkmarkm\checkmark$\checkmark^\checkmark\\checkmarkc\checkmarki\checkmarkr\checkmarkc\checkmark$\checkmark \checkmark$\checkmark^\checkmark\\checkmarkc\checkmarki\checkmarkr\checkmarkc\checkmark$\checkmarkR\checkmark$\checkmark^\checkmark\\checkmarkc\checkmarki\checkmarkr\checkmarkc\checkmark$\checkmarka\checkmark$\checkmark^\checkmark\\checkmarkc\checkmarki\checkmarkr\checkmarkc\checkmark$\checkmarkp\checkmark$\checkmark^\checkmark\\checkmarkc\checkmarki\checkmarkr\checkmarkc\checkmark$\checkmarkp\checkmark$\checkmark^\checkmark\\checkmarkc\checkmarki\checkmarkr\checkmarkc\checkmark$\checkmarko\checkmark$\checkmark^\checkmark\\checkmarkc\checkmarki\checkmarkr\checkmarkc\checkmark$\checkmarkr\checkmark$\checkmark^\checkmark\\checkmarkc\checkmarki\checkmarkr\checkmarkc\checkmark$\checkmarkt\checkmark$\checkmark^\checkmark\\checkmarkc\checkmarki\checkmarkr\checkmarkc\checkmark$\checkmark \checkmark$\checkmark^\checkmark\\checkmarkc\checkmarki\checkmarkr\checkmarkc\checkmark$\checkmark:\checkmark$\checkmark^\checkmark\\checkmarkc\checkmarki\checkmarkr\checkmarkc\checkmark$\checkmark \checkmark$\checkmark^\checkmark\\checkmarkc\checkmarki\checkmarkr\checkmarkc\checkmark$\checkmark$\checkmark$\checkmark^\checkmark\\checkmarkc\checkmarki\checkmarkr\checkmarkc\checkmark$\checkmarki\checkmark$\checkmark^\checkmark\\checkmarkc\checkmarki\checkmarkr\checkmarkc\checkmark$\checkmark_\checkmark$\checkmark^\checkmark\\checkmarkc\checkmarki\checkmarkr\checkmarkc\checkmark$\checkmark2\checkmark$\checkmark^\checkmark\\checkmarkc\checkmarki\checkmarkr\checkmarkc\checkmark$\checkmark \checkmark$\checkmark^\checkmark\\checkmarkc\checkmarki\checkmarkr\checkmarkc\checkmark$\checkmark=\checkmark$\checkmark^\checkmark\\checkmarkc\checkmarki\checkmarkr\checkmarkc\checkmark$\checkmark \checkmark$\checkmark^\checkmark\\checkmarkc\checkmarki\checkmarkr\checkmarkc\checkmark$\checkmarkZ\checkmark$\checkmark^\checkmark\\checkmarkc\checkmarki\checkmarkr\checkmarkc\checkmark$\checkmark_\checkmark$\checkmark^\checkmark\\checkmarkc\checkmarki\checkmarkr\checkmarkc\checkmark$\checkmark4\checkmark$\checkmark^\checkmark\\checkmarkc\checkmarki\checkmarkr\checkmarkc\checkmark$\checkmark/\checkmark$\checkmark^\checkmark\\checkmarkc\checkmarki\checkmarkr\checkmarkc\checkmark$\checkmarkZ\checkmark$\checkmark^\checkmark\\checkmarkc\checkmarki\checkmarkr\checkmarkc\checkmark$\checkmark_\checkmark$\checkmark^\checkmark\\checkmarkc\checkmarki\checkmarkr\checkmarkc\checkmark$\checkmark3\checkmark$\checkmark^\checkmark\\checkmarkc\checkmarki\checkmarkr\checkmarkc\checkmark$\checkmark \checkmark$\checkmark^\checkmark\\checkmarkc\checkmarki\checkmarkr\checkmarkc\checkmark$\checkmark=\checkmark$\checkmark^\checkmark\\checkmarkc\checkmarki\checkmarkr\checkmarkc\checkmark$\checkmark \checkmark$\checkmark^\checkmark\\checkmarkc\checkmarki\checkmarkr\checkmarkc\checkmark$\checkmark1\checkmark$\checkmark^\checkmark\\checkmarkc\checkmarki\checkmarkr\checkmarkc\checkmark$\checkmark0\checkmark$\checkmark^\checkmark\\checkmarkc\checkmarki\checkmarkr\checkmarkc\checkmark$\checkmark0\checkmark$\checkmark^\checkmark\\checkmarkc\checkmarki\checkmarkr\checkmarkc\checkmark$\checkmark/\checkmark$\checkmark^\checkmark\\checkmarkc\checkmarki\checkmarkr\checkmarkc\checkmark$\checkmark2\checkmark$\checkmark^\checkmark\\checkmarkc\checkmarki\checkmarkr\checkmarkc\checkmark$\checkmark0\checkmark$\checkmark^\checkmark\\checkmarkc\checkmarki\checkmarkr\checkmarkc\checkmark$\checkmark \checkmark$\checkmark^\checkmark\\checkmarkc\checkmarki\checkmarkr\checkmarkc\checkmark$\checkmark=\checkmark$\checkmark^\checkmark\\checkmarkc\checkmarki\checkmarkr\checkmarkc\checkmark$\checkmark \checkmark$\checkmark^\checkmark\\checkmarkc\checkmarki\checkmarkr\checkmarkc\checkmark$\checkmark5\checkmark$\checkmark^\checkmark\\checkmarkc\checkmarki\checkmarkr\checkmarkc\checkmark$\checkmark$\checkmark$\checkmark^\checkmark\\checkmarkc\checkmarki\checkmarkr\checkmarkc\checkmark$\checkmark
\checkmark$\checkmark^\checkmark\\checkmarkc\checkmarki\checkmarkr\checkmarkc\checkmark$\checkmark\\checkmark$\checkmark^\checkmark\\checkmarkc\checkmarki\checkmarkr\checkmarkc\checkmark$\checkmarke\checkmark$\checkmark^\checkmark\\checkmarkc\checkmarki\checkmarkr\checkmarkc\checkmark$\checkmarkn\checkmark$\checkmark^\checkmark\\checkmarkc\checkmarki\checkmarkr\checkmarkc\checkmark$\checkmarkd\checkmark$\checkmark^\checkmark\\checkmarkc\checkmarki\checkmarkr\checkmarkc\checkmark$\checkmark{\checkmark$\checkmark^\checkmark\\checkmarkc\checkmarki\checkmarkr\checkmarkc\checkmark$\checkmarki\checkmark$\checkmark^\checkmark\\checkmarkc\checkmarki\checkmarkr\checkmarkc\checkmark$\checkmarkt\checkmark$\checkmark^\checkmark\\checkmarkc\checkmarki\checkmarkr\checkmarkc\checkmark$\checkmarke\checkmark$\checkmark^\checkmark\\checkmarkc\checkmarki\checkmarkr\checkmarkc\checkmark$\checkmarkm\checkmark$\checkmark^\checkmark\\checkmarkc\checkmarki\checkmarkr\checkmarkc\checkmark$\checkmarki\checkmark$\checkmark^\checkmark\\checkmarkc\checkmarki\checkmarkr\checkmarkc\checkmark$\checkmarkz\checkmark$\checkmark^\checkmark\\checkmarkc\checkmarki\checkmarkr\checkmarkc\checkmark$\checkmarke\checkmark$\checkmark^\checkmark\\checkmarkc\checkmarki\checkmarkr\checkmarkc\checkmark$\checkmark}\checkmark$\checkmark^\checkmark\\checkmarkc\checkmarki\checkmarkr\checkmarkc\checkmark$\checkmark
\checkmark$\checkmark^\checkmark\\checkmarkc\checkmarki\checkmarkr\checkmarkc\checkmark$\checkmark
\checkmark$\checkmark^\checkmark\\checkmarkc\checkmarki\checkmarkr\checkmarkc\checkmark$\checkmark\\checkmark$\checkmark^\checkmark\\checkmarkc\checkmarki\checkmarkr\checkmarkc\checkmark$\checkmarkt\checkmark$\checkmark^\checkmark\\checkmarkc\checkmarki\checkmarkr\checkmarkc\checkmark$\checkmarke\checkmark$\checkmark^\checkmark\\checkmarkc\checkmarki\checkmarkr\checkmarkc\checkmark$\checkmarkx\checkmark$\checkmark^\checkmark\\checkmarkc\checkmarki\checkmarkr\checkmarkc\checkmark$\checkmarkt\checkmark$\checkmark^\checkmark\\checkmarkc\checkmarki\checkmarkr\checkmarkc\checkmark$\checkmarkb\checkmark$\checkmark^\checkmark\\checkmarkc\checkmarki\checkmarkr\checkmarkc\checkmark$\checkmarkf\checkmark$\checkmark^\checkmark\\checkmarkc\checkmarki\checkmarkr\checkmarkc\checkmark$\checkmark{\checkmark$\checkmark^\checkmark\\checkmarkc\checkmarki\checkmarkr\checkmarkc\checkmark$\checkmarkR\checkmark$\checkmark^\checkmark\\checkmarkc\checkmarki\checkmarkr\checkmarkc\checkmark$\checkmarka\checkmark$\checkmark^\checkmark\\checkmarkc\checkmarki\checkmarkr\checkmarkc\checkmark$\checkmarkp\checkmark$\checkmark^\checkmark\\checkmarkc\checkmarki\checkmarkr\checkmarkc\checkmark$\checkmarkp\checkmark$\checkmark^\checkmark\\checkmarkc\checkmarki\checkmarkr\checkmarkc\checkmark$\checkmarko\checkmark$\checkmark^\checkmark\\checkmarkc\checkmarki\checkmarkr\checkmarkc\checkmark$\checkmarkr\checkmark$\checkmark^\checkmark\\checkmarkc\checkmarki\checkmarkr\checkmarkc\checkmark$\checkmarkt\checkmark$\checkmark^\checkmark\\checkmarkc\checkmarki\checkmarkr\checkmarkc\checkmark$\checkmark \checkmark$\checkmark^\checkmark\\checkmarkc\checkmarki\checkmarkr\checkmarkc\checkmark$\checkmarkt\checkmark$\checkmark^\checkmark\\checkmarkc\checkmarki\checkmarkr\checkmarkc\checkmark$\checkmarko\checkmark$\checkmark^\checkmark\\checkmarkc\checkmarki\checkmarkr\checkmarkc\checkmark$\checkmarkt\checkmark$\checkmark^\checkmark\\checkmarkc\checkmarki\checkmarkr\checkmarkc\checkmark$\checkmarka\checkmark$\checkmark^\checkmark\\checkmarkc\checkmarki\checkmarkr\checkmarkc\checkmark$\checkmarkl\checkmark$\checkmark^\checkmark\\checkmarkc\checkmarki\checkmarkr\checkmarkc\checkmark$\checkmark \checkmark$\checkmark^\checkmark\\checkmarkc\checkmarki\checkmarkr\checkmarkc\checkmark$\checkmark:\checkmark$\checkmark^\checkmark\\checkmarkc\checkmarki\checkmarkr\checkmarkc\checkmark$\checkmark}\checkmark$\checkmark^\checkmark\\checkmarkc\checkmarki\checkmarkr\checkmarkc\checkmark$\checkmark \checkmark$\checkmark^\checkmark\\checkmarkc\checkmarki\checkmarkr\checkmarkc\checkmark$\checkmark$\checkmark$\checkmark^\checkmark\\checkmarkc\checkmarki\checkmarkr\checkmarkc\checkmark$\checkmarki\checkmark$\checkmark^\checkmark\\checkmarkc\checkmarki\checkmarkr\checkmarkc\checkmark$\checkmark_\checkmark$\checkmark^\checkmark\\checkmarkc\checkmarki\checkmarkr\checkmarkc\checkmark$\checkmark{\checkmark$\checkmark^\checkmark\\checkmarkc\checkmarki\checkmarkr\checkmarkc\checkmark$\checkmarkt\checkmark$\checkmark^\checkmark\\checkmarkc\checkmarki\checkmarkr\checkmarkc\checkmark$\checkmarko\checkmark$\checkmark^\checkmark\\checkmarkc\checkmarki\checkmarkr\checkmarkc\checkmark$\checkmarkt\checkmark$\checkmark^\checkmark\\checkmarkc\checkmarki\checkmarkr\checkmarkc\checkmark$\checkmarka\checkmark$\checkmark^\checkmark\\checkmarkc\checkmarki\checkmarkr\checkmarkc\checkmark$\checkmarkl\checkmark$\checkmark^\checkmark\\checkmarkc\checkmarki\checkmarkr\checkmarkc\checkmark$\checkmark}\checkmark$\checkmark^\checkmark\\checkmarkc\checkmarki\checkmarkr\checkmarkc\checkmark$\checkmark \checkmark$\checkmark^\checkmark\\checkmarkc\checkmarki\checkmarkr\checkmarkc\checkmark$\checkmark=\checkmark$\checkmark^\checkmark\\checkmarkc\checkmarki\checkmarkr\checkmarkc\checkmark$\checkmark \checkmark$\checkmark^\checkmark\\checkmarkc\checkmarki\checkmarkr\checkmarkc\checkmark$\checkmarki\checkmark$\checkmark^\checkmark\\checkmarkc\checkmarki\checkmarkr\checkmarkc\checkmark$\checkmark_\checkmark$\checkmark^\checkmark\\checkmarkc\checkmarki\checkmarkr\checkmarkc\checkmark$\checkmark1\checkmark$\checkmark^\checkmark\\checkmarkc\checkmarki\checkmarkr\checkmarkc\checkmark$\checkmark \checkmark$\checkmark^\checkmark\\checkmarkc\checkmarki\checkmarkr\checkmarkc\checkmark$\checkmark\\checkmark$\checkmark^\checkmark\\checkmarkc\checkmarki\checkmarkr\checkmarkc\checkmark$\checkmarkt\checkmark$\checkmark^\checkmark\\checkmarkc\checkmarki\checkmarkr\checkmarkc\checkmark$\checkmarki\checkmark$\checkmark^\checkmark\\checkmarkc\checkmarki\checkmarkr\checkmarkc\checkmark$\checkmarkm\checkmark$\checkmark^\checkmark\\checkmarkc\checkmarki\checkmarkr\checkmarkc\checkmark$\checkmarke\checkmark$\checkmark^\checkmark\\checkmarkc\checkmarki\checkmarkr\checkmarkc\checkmark$\checkmarks\checkmark$\checkmark^\checkmark\\checkmarkc\checkmarki\checkmarkr\checkmarkc\checkmark$\checkmark \checkmark$\checkmark^\checkmark\\checkmarkc\checkmarki\checkmarkr\checkmarkc\checkmark$\checkmarki\checkmark$\checkmark^\checkmark\\checkmarkc\checkmarki\checkmarkr\checkmarkc\checkmark$\checkmark_\checkmark$\checkmark^\checkmark\\checkmarkc\checkmarki\checkmarkr\checkmarkc\checkmark$\checkmark2\checkmark$\checkmark^\checkmark\\checkmarkc\checkmarki\checkmarkr\checkmarkc\checkmark$\checkmark \checkmark$\checkmark^\checkmark\\checkmarkc\checkmarki\checkmarkr\checkmarkc\checkmark$\checkmark=\checkmark$\checkmark^\checkmark\\checkmarkc\checkmarki\checkmarkr\checkmarkc\checkmark$\checkmark \checkmark$\checkmark^\checkmark\\checkmarkc\checkmarki\checkmarkr\checkmarkc\checkmark$\checkmark1\checkmark$\checkmark^\checkmark\\checkmarkc\checkmarki\checkmarkr\checkmarkc\checkmark$\checkmark0\checkmark$\checkmark^\checkmark\\checkmarkc\checkmarki\checkmarkr\checkmarkc\checkmark$\checkmark \checkmark$\checkmark^\checkmark\\checkmarkc\checkmarki\checkmarkr\checkmarkc\checkmark$\checkmark\\checkmark$\checkmark^\checkmark\\checkmarkc\checkmarki\checkmarkr\checkmarkc\checkmark$\checkmarkt\checkmark$\checkmark^\checkmark\\checkmarkc\checkmarki\checkmarkr\checkmarkc\checkmark$\checkmarki\checkmark$\checkmark^\checkmark\\checkmarkc\checkmarki\checkmarkr\checkmarkc\checkmark$\checkmarkm\checkmark$\checkmark^\checkmark\\checkmarkc\checkmarki\checkmarkr\checkmarkc\checkmark$\checkmarke\checkmark$\checkmark^\checkmark\\checkmarkc\checkmarki\checkmarkr\checkmarkc\checkmark$\checkmarks\checkmark$\checkmark^\checkmark\\checkmarkc\checkmarki\checkmarkr\checkmarkc\checkmark$\checkmark \checkmark$\checkmark^\checkmark\\checkmarkc\checkmarki\checkmarkr\checkmarkc\checkmark$\checkmark5\checkmark$\checkmark^\checkmark\\checkmarkc\checkmarki\checkmarkr\checkmarkc\checkmark$\checkmark \checkmark$\checkmark^\checkmark\\checkmarkc\checkmarki\checkmarkr\checkmarkc\checkmark$\checkmark=\checkmark$\checkmark^\checkmark\\checkmarkc\checkmarki\checkmarkr\checkmarkc\checkmark$\checkmark \checkmark$\checkmark^\checkmark\\checkmarkc\checkmarki\checkmarkr\checkmarkc\checkmark$\checkmark\\checkmark$\checkmark^\checkmark\\checkmarkc\checkmarki\checkmarkr\checkmarkc\checkmark$\checkmarkb\checkmark$\checkmark^\checkmark\\checkmarkc\checkmarki\checkmarkr\checkmarkc\checkmark$\checkmarko\checkmark$\checkmark^\checkmark\\checkmarkc\checkmarki\checkmarkr\checkmarkc\checkmark$\checkmarkx\checkmark$\checkmark^\checkmark\\checkmarkc\checkmarki\checkmarkr\checkmarkc\checkmark$\checkmarke\checkmark$\checkmark^\checkmark\\checkmarkc\checkmarki\checkmarkr\checkmarkc\checkmark$\checkmarkd\checkmark$\checkmark^\checkmark\\checkmarkc\checkmarki\checkmarkr\checkmarkc\checkmark$\checkmark{\checkmark$\checkmark^\checkmark\\checkmarkc\checkmarki\checkmarkr\checkmarkc\checkmark$\checkmark5\checkmark$\checkmark^\checkmark\\checkmarkc\checkmarki\checkmarkr\checkmarkc\checkmark$\checkmark0\checkmark$\checkmark^\checkmark\\checkmarkc\checkmarki\checkmarkr\checkmarkc\checkmark$\checkmark}\checkmark$\checkmark^\checkmark\\checkmarkc\checkmarki\checkmarkr\checkmarkc\checkmark$\checkmark$\checkmark$\checkmark^\checkmark\\checkmarkc\checkmarki\checkmarkr\checkmarkc\checkmark$\checkmark \checkmark$\checkmark^\checkmark\\checkmarkc\checkmarki\checkmarkr\checkmarkc\checkmark$\checkmark✓\checkmark$\checkmark^\checkmark\\checkmarkc\checkmarki\checkmarkr\checkmarkc\checkmark$\checkmark
\checkmark$\checkmark^\checkmark\\checkmarkc\checkmarki\checkmarkr\checkmarkc\checkmark$\checkmark
\checkmark$\checkmark^\checkmark\\checkmarkc\checkmarki\checkmarkr\checkmarkc\checkmark$\checkmark\\checkmark$\checkmark^\checkmark\\checkmarkc\checkmarki\checkmarkr\checkmarkc\checkmark$\checkmarks\checkmark$\checkmark^\checkmark\\checkmarkc\checkmarki\checkmarkr\checkmarkc\checkmark$\checkmarku\checkmark$\checkmark^\checkmark\\checkmarkc\checkmarki\checkmarkr\checkmarkc\checkmark$\checkmarkb\checkmark$\checkmark^\checkmark\\checkmarkc\checkmarki\checkmarkr\checkmarkc\checkmark$\checkmarks\checkmark$\checkmark^\checkmark\\checkmarkc\checkmarki\checkmarkr\checkmarkc\checkmark$\checkmarku\checkmark$\checkmark^\checkmark\\checkmarkc\checkmarki\checkmarkr\checkmarkc\checkmark$\checkmarkb\checkmark$\checkmark^\checkmark\\checkmarkc\checkmarki\checkmarkr\checkmarkc\checkmark$\checkmarks\checkmark$\checkmark^\checkmark\\checkmarkc\checkmarki\checkmarkr\checkmarkc\checkmark$\checkmarke\checkmark$\checkmark^\checkmark\\checkmarkc\checkmarki\checkmarkr\checkmarkc\checkmark$\checkmarkc\checkmark$\checkmark^\checkmark\\checkmarkc\checkmarki\checkmarkr\checkmarkc\checkmark$\checkmarkt\checkmark$\checkmark^\checkmark\\checkmarkc\checkmarki\checkmarkr\checkmarkc\checkmark$\checkmarki\checkmark$\checkmark^\checkmark\\checkmarkc\checkmarki\checkmarkr\checkmarkc\checkmark$\checkmarko\checkmark$\checkmark^\checkmark\\checkmarkc\checkmarki\checkmarkr\checkmarkc\checkmark$\checkmarkn\checkmark$\checkmark^\checkmark\\checkmarkc\checkmarki\checkmarkr\checkmarkc\checkmark$\checkmark{\checkmark$\checkmark^\checkmark\\checkmarkc\checkmarki\checkmarkr\checkmarkc\checkmark$\checkmarkC\checkmark$\checkmark^\checkmark\\checkmarkc\checkmarki\checkmarkr\checkmarkc\checkmark$\checkmarka\checkmark$\checkmark^\checkmark\\checkmarkc\checkmarki\checkmarkr\checkmarkc\checkmark$\checkmarkl\checkmark$\checkmark^\checkmark\\checkmarkc\checkmarki\checkmarkr\checkmarkc\checkmark$\checkmarkc\checkmark$\checkmark^\checkmark\\checkmarkc\checkmarki\checkmarkr\checkmarkc\checkmark$\checkmarku\checkmark$\checkmark^\checkmark\\checkmarkc\checkmarki\checkmarkr\checkmarkc\checkmark$\checkmarkl\checkmark$\checkmark^\checkmark\\checkmarkc\checkmarki\checkmarkr\checkmarkc\checkmark$\checkmark \checkmark$\checkmark^\checkmark\\checkmarkc\checkmarki\checkmarkr\checkmarkc\checkmark$\checkmarkd\checkmark$\checkmark^\checkmark\\checkmarkc\checkmarki\checkmarkr\checkmarkc\checkmark$\checkmarke\checkmark$\checkmark^\checkmark\\checkmarkc\checkmarki\checkmarkr\checkmarkc\checkmark$\checkmark \checkmark$\checkmark^\checkmark\\checkmarkc\checkmarki\checkmarkr\checkmarkc\checkmark$\checkmarkl\checkmark$\checkmark^\checkmark\\checkmarkc\checkmarki\checkmarkr\checkmarkc\checkmark$\checkmarka\checkmark$\checkmark^\checkmark\\checkmarkc\checkmarki\checkmarkr\checkmarkc\checkmark$\checkmark \checkmark$\checkmark^\checkmark\\checkmarkc\checkmarki\checkmarkr\checkmarkc\checkmark$\checkmarkt\checkmark$\checkmark^\checkmark\\checkmarkc\checkmarki\checkmarkr\checkmarkc\checkmark$\checkmarke\checkmark$\checkmark^\checkmark\\checkmarkc\checkmarki\checkmarkr\checkmarkc\checkmark$\checkmarkn\checkmark$\checkmark^\checkmark\\checkmarkc\checkmarki\checkmarkr\checkmarkc\checkmark$\checkmarks\checkmark$\checkmark^\checkmark\\checkmarkc\checkmarki\checkmarkr\checkmarkc\checkmark$\checkmarki\checkmark$\checkmark^\checkmark\\checkmarkc\checkmarki\checkmarkr\checkmarkc\checkmark$\checkmarko\checkmark$\checkmark^\checkmark\\checkmarkc\checkmarki\checkmarkr\checkmarkc\checkmark$\checkmarkn\checkmark$\checkmark^\checkmark\\checkmarkc\checkmarki\checkmarkr\checkmarkc\checkmark$\checkmark \checkmark$\checkmark^\checkmark\\checkmarkc\checkmarki\checkmarkr\checkmarkc\checkmark$\checkmarkd\checkmark$\checkmark^\checkmark\\checkmarkc\checkmarki\checkmarkr\checkmarkc\checkmark$\checkmarka\checkmark$\checkmark^\checkmark\\checkmarkc\checkmarki\checkmarkr\checkmarkc\checkmark$\checkmarkn\checkmark$\checkmark^\checkmark\\checkmarkc\checkmarki\checkmarkr\checkmarkc\checkmark$\checkmarks\checkmark$\checkmark^\checkmark\\checkmarkc\checkmarki\checkmarkr\checkmarkc\checkmark$\checkmark \checkmark$\checkmark^\checkmark\\checkmarkc\checkmarki\checkmarkr\checkmarkc\checkmark$\checkmarkl\checkmark$\checkmark^\checkmark\\checkmarkc\checkmarki\checkmarkr\checkmarkc\checkmark$\checkmarka\checkmark$\checkmark^\checkmark\\checkmarkc\checkmarki\checkmarkr\checkmarkc\checkmark$\checkmark \checkmark$\checkmark^\checkmark\\checkmarkc\checkmarki\checkmarkr\checkmarkc\checkmark$\checkmarkc\checkmark$\checkmark^\checkmark\\checkmarkc\checkmarki\checkmarkr\checkmarkc\checkmark$\checkmarko\checkmark$\checkmark^\checkmark\\checkmarkc\checkmarki\checkmarkr\checkmarkc\checkmark$\checkmarku\checkmark$\checkmark^\checkmark\\checkmarkc\checkmarki\checkmarkr\checkmarkc\checkmark$\checkmarkr\checkmark$\checkmark^\checkmark\\checkmarkc\checkmarki\checkmarkr\checkmarkc\checkmark$\checkmarkr\checkmark$\checkmark^\checkmark\\checkmarkc\checkmarki\checkmarkr\checkmarkc\checkmark$\checkmarko\checkmark$\checkmark^\checkmark\\checkmarkc\checkmarki\checkmarkr\checkmarkc\checkmark$\checkmarki\checkmark$\checkmark^\checkmark\\checkmarkc\checkmarki\checkmarkr\checkmarkc\checkmark$\checkmarke\checkmark$\checkmark^\checkmark\\checkmarkc\checkmarki\checkmarkr\checkmarkc\checkmark$\checkmark \checkmark$\checkmark^\checkmark\\checkmarkc\checkmarki\checkmarkr\checkmarkc\checkmark$\checkmark(\checkmark$\checkmark^\checkmark\\checkmarkc\checkmarki\checkmarkr\checkmarkc\checkmark$\checkmarkÉ\checkmark$\checkmark^\checkmark\\checkmarkc\checkmarki\checkmarkr\checkmarkc\checkmark$\checkmarkt\checkmark$\checkmark^\checkmark\\checkmarkc\checkmarki\checkmarkr\checkmarkc\checkmark$\checkmarka\checkmark$\checkmark^\checkmark\\checkmarkc\checkmarki\checkmarkr\checkmarkc\checkmark$\checkmarkg\checkmark$\checkmark^\checkmark\\checkmarkc\checkmarki\checkmarkr\checkmarkc\checkmark$\checkmarke\checkmark$\checkmark^\checkmark\\checkmarkc\checkmarki\checkmarkr\checkmarkc\checkmark$\checkmark \checkmark$\checkmark^\checkmark\\checkmarkc\checkmarki\checkmarkr\checkmarkc\checkmark$\checkmark2\checkmark$\checkmark^\checkmark\\checkmarkc\checkmarki\checkmarkr\checkmarkc\checkmark$\checkmark)\checkmark$\checkmark^\checkmark\\checkmarkc\checkmarki\checkmarkr\checkmarkc\checkmark$\checkmark}\checkmark$\checkmark^\checkmark\\checkmarkc\checkmarki\checkmarkr\checkmarkc\checkmark$\checkmark
\checkmark$\checkmark^\checkmark\\checkmarkc\checkmarki\checkmarkr\checkmarkc\checkmark$\checkmarkL\checkmark$\checkmark^\checkmark\\checkmarkc\checkmarki\checkmarkr\checkmarkc\checkmark$\checkmarke\checkmark$\checkmark^\checkmark\\checkmarkc\checkmarki\checkmarkr\checkmarkc\checkmark$\checkmark \checkmark$\checkmark^\checkmark\\checkmarkc\checkmarki\checkmarkr\checkmarkc\checkmark$\checkmarkm\checkmark$\checkmark^\checkmark\\checkmarkc\checkmarki\checkmarkr\checkmarkc\checkmark$\checkmarko\checkmark$\checkmark^\checkmark\\checkmarkc\checkmarki\checkmarkr\checkmarkc\checkmark$\checkmarkm\checkmark$\checkmark^\checkmark\\checkmarkc\checkmarki\checkmarkr\checkmarkc\checkmark$\checkmarke\checkmark$\checkmark^\checkmark\\checkmarkc\checkmarki\checkmarkr\checkmarkc\checkmark$\checkmarkn\checkmark$\checkmark^\checkmark\\checkmarkc\checkmarki\checkmarkr\checkmarkc\checkmark$\checkmarkt\checkmark$\checkmark^\checkmark\\checkmarkc\checkmarki\checkmarkr\checkmarkc\checkmark$\checkmark \checkmark$\checkmark^\checkmark\\checkmarkc\checkmarki\checkmarkr\checkmarkc\checkmark$\checkmarkd\checkmark$\checkmark^\checkmark\\checkmarkc\checkmarki\checkmarkr\checkmarkc\checkmark$\checkmarke\checkmark$\checkmark^\checkmark\\checkmarkc\checkmarki\checkmarkr\checkmarkc\checkmark$\checkmark \checkmark$\checkmark^\checkmark\\checkmarkc\checkmarki\checkmarkr\checkmarkc\checkmark$\checkmarkr\checkmark$\checkmark^\checkmark\\checkmarkc\checkmarki\checkmarkr\checkmarkc\checkmark$\checkmarke\checkmark$\checkmark^\checkmark\\checkmarkc\checkmarki\checkmarkr\checkmarkc\checkmark$\checkmarkn\checkmark$\checkmark^\checkmark\\checkmarkc\checkmarki\checkmarkr\checkmarkc\checkmark$\checkmarkv\checkmark$\checkmark^\checkmark\\checkmarkc\checkmarki\checkmarkr\checkmarkc\checkmark$\checkmarke\checkmark$\checkmark^\checkmark\\checkmarkc\checkmarki\checkmarkr\checkmarkc\checkmark$\checkmarkr\checkmark$\checkmark^\checkmark\\checkmarkc\checkmarki\checkmarkr\checkmarkc\checkmark$\checkmarks\checkmark$\checkmark^\checkmark\\checkmarkc\checkmarki\checkmarkr\checkmarkc\checkmark$\checkmarke\checkmark$\checkmark^\checkmark\\checkmarkc\checkmarki\checkmarkr\checkmarkc\checkmark$\checkmarkm\checkmark$\checkmark^\checkmark\\checkmarkc\checkmarki\checkmarkr\checkmarkc\checkmark$\checkmarke\checkmark$\checkmark^\checkmark\\checkmarkc\checkmarki\checkmarkr\checkmarkc\checkmark$\checkmarkn\checkmark$\checkmark^\checkmark\\checkmarkc\checkmarki\checkmarkr\checkmarkc\checkmark$\checkmarkt\checkmark$\checkmark^\checkmark\\checkmarkc\checkmarki\checkmarkr\checkmarkc\checkmark$\checkmark \checkmark$\checkmark^\checkmark\\checkmarkc\checkmarki\checkmarkr\checkmarkc\checkmark$\checkmarkc\checkmark$\checkmark^\checkmark\\checkmarkc\checkmarki\checkmarkr\checkmarkc\checkmark$\checkmarka\checkmark$\checkmark^\checkmark\\checkmarkc\checkmarki\checkmarkr\checkmarkc\checkmark$\checkmarkl\checkmark$\checkmark^\checkmark\\checkmarkc\checkmarki\checkmarkr\checkmarkc\checkmark$\checkmarkc\checkmark$\checkmark^\checkmark\\checkmarkc\checkmarki\checkmarkr\checkmarkc\checkmark$\checkmarku\checkmark$\checkmark^\checkmark\\checkmarkc\checkmarki\checkmarkr\checkmarkc\checkmark$\checkmarkl\checkmark$\checkmark^\checkmark\\checkmarkc\checkmarki\checkmarkr\checkmarkc\checkmark$\checkmarké\checkmark$\checkmark^\checkmark\\checkmarkc\checkmarki\checkmarkr\checkmarkc\checkmark$\checkmark \checkmark$\checkmark^\checkmark\\checkmarkc\checkmarki\checkmarkr\checkmarkc\checkmark$\checkmarks\checkmark$\checkmark^\checkmark\\checkmarkc\checkmarki\checkmarkr\checkmarkc\checkmark$\checkmarke\checkmark$\checkmark^\checkmark\\checkmarkc\checkmarki\checkmarkr\checkmarkc\checkmark$\checkmarkc\checkmark$\checkmark^\checkmark\\checkmarkc\checkmarki\checkmarkr\checkmarkc\checkmark$\checkmarkt\checkmark$\checkmark^\checkmark\\checkmarkc\checkmarki\checkmarkr\checkmarkc\checkmark$\checkmarki\checkmark$\checkmark^\checkmark\\checkmarkc\checkmarki\checkmarkr\checkmarkc\checkmark$\checkmarko\checkmark$\checkmark^\checkmark\\checkmarkc\checkmarki\checkmarkr\checkmarkc\checkmark$\checkmarkn\checkmark$\checkmark^\checkmark\\checkmarkc\checkmarki\checkmarkr\checkmarkc\checkmark$\checkmark \checkmark$\checkmark^\checkmark\\checkmarkc\checkmarki\checkmarkr\checkmarkc\checkmark$\checkmarks\checkmark$\checkmark^\checkmark\\checkmarkc\checkmarki\checkmarkr\checkmarkc\checkmark$\checkmarku\checkmark$\checkmark^\checkmark\\checkmarkc\checkmarki\checkmarkr\checkmarkc\checkmark$\checkmarki\checkmark$\checkmark^\checkmark\\checkmarkc\checkmarki\checkmarkr\checkmarkc\checkmark$\checkmarkv\checkmark$\checkmark^\checkmark\\checkmarkc\checkmarki\checkmarkr\checkmarkc\checkmark$\checkmarka\checkmark$\checkmark^\checkmark\\checkmarkc\checkmarki\checkmarkr\checkmarkc\checkmark$\checkmarkn\checkmark$\checkmark^\checkmark\\checkmarkc\checkmarki\checkmarkr\checkmarkc\checkmark$\checkmarkt\checkmark$\checkmark^\checkmark\\checkmarkc\checkmarki\checkmarkr\checkmarkc\checkmark$\checkmarke\checkmark$\checkmark^\checkmark\\checkmarkc\checkmarki\checkmarkr\checkmarkc\checkmark$\checkmark \checkmark$\checkmark^\checkmark\\checkmarkc\checkmarki\checkmarkr\checkmarkc\checkmark$\checkmarke\checkmark$\checkmark^\checkmark\\checkmarkc\checkmarki\checkmarkr\checkmarkc\checkmark$\checkmarks\checkmark$\checkmark^\checkmark\\checkmarkc\checkmarki\checkmarkr\checkmarkc\checkmark$\checkmarkt\checkmark$\checkmark^\checkmark\\checkmarkc\checkmarki\checkmarkr\checkmarkc\checkmark$\checkmark \checkmark$\checkmark^\checkmark\\checkmarkc\checkmarki\checkmarkr\checkmarkc\checkmark$\checkmark$\checkmark$\checkmark^\checkmark\\checkmarkc\checkmarki\checkmarkr\checkmarkc\checkmark$\checkmarkM\checkmark$\checkmark^\checkmark\\checkmarkc\checkmarki\checkmarkr\checkmarkc\checkmark$\checkmark_\checkmark$\checkmark^\checkmark\\checkmarkc\checkmarki\checkmarkr\checkmarkc\checkmark$\checkmark{\checkmark$\checkmark^\checkmark\\checkmarkc\checkmarki\checkmarkr\checkmarkc\checkmark$\checkmarkr\checkmark$\checkmark^\checkmark\\checkmarkc\checkmarki\checkmarkr\checkmarkc\checkmark$\checkmarke\checkmark$\checkmark^\checkmark\\checkmarkc\checkmarki\checkmarkr\checkmarkc\checkmark$\checkmarkn\checkmark$\checkmark^\checkmark\\checkmarkc\checkmarki\checkmarkr\checkmarkc\checkmark$\checkmarkv\checkmark$\checkmark^\checkmark\\checkmarkc\checkmarki\checkmarkr\checkmarkc\checkmark$\checkmark}\checkmark$\checkmark^\checkmark\\checkmarkc\checkmarki\checkmarkr\checkmarkc\checkmark$\checkmark \checkmark$\checkmark^\checkmark\\checkmarkc\checkmarki\checkmarkr\checkmarkc\checkmark$\checkmark=\checkmark$\checkmark^\checkmark\\checkmarkc\checkmarki\checkmarkr\checkmarkc\checkmark$\checkmark \checkmark$\checkmark^\checkmark\\checkmarkc\checkmarki\checkmarkr\checkmarkc\checkmark$\checkmark2\checkmark$\checkmark^\checkmark\\checkmarkc\checkmarki\checkmarkr\checkmarkc\checkmark$\checkmark0\checkmark$\checkmark^\checkmark\\checkmarkc\checkmarki\checkmarkr\checkmarkc\checkmark$\checkmark$\checkmark$\checkmark^\checkmark\\checkmarkc\checkmarki\checkmarkr\checkmarkc\checkmark$\checkmark \checkmark$\checkmark^\checkmark\\checkmarkc\checkmarki\checkmarkr\checkmarkc\checkmark$\checkmarkN\checkmark$\checkmark^\checkmark\\checkmarkc\checkmarki\checkmarkr\checkmarkc\checkmark$\checkmark$\checkmark$\checkmark^\checkmark\\checkmarkc\checkmarki\checkmarkr\checkmarkc\checkmark$\checkmark\\checkmark$\checkmark^\checkmark\\checkmarkc\checkmarki\checkmarkr\checkmarkc\checkmark$\checkmarkc\checkmark$\checkmark^\checkmark\\checkmarkc\checkmarki\checkmarkr\checkmarkc\checkmark$\checkmarkd\checkmark$\checkmark^\checkmark\\checkmarkc\checkmarki\checkmarkr\checkmarkc\checkmark$\checkmarko\checkmark$\checkmark^\checkmark\\checkmarkc\checkmarki\checkmarkr\checkmarkc\checkmark$\checkmarkt\checkmark$\checkmark^\checkmark\\checkmarkc\checkmarki\checkmarkr\checkmarkc\checkmark$\checkmark$\checkmark$\checkmark^\checkmark\\checkmarkc\checkmarki\checkmarkr\checkmarkc\checkmark$\checkmarkm\checkmark$\checkmark^\checkmark\\checkmarkc\checkmarki\checkmarkr\checkmarkc\checkmark$\checkmark.\checkmark$\checkmark^\checkmark\\checkmarkc\checkmarki\checkmarkr\checkmarkc\checkmark$\checkmark \checkmark$\checkmark^\checkmark\\checkmarkc\checkmarki\checkmarkr\checkmarkc\checkmark$\checkmarkL\checkmark$\checkmark^\checkmark\\checkmarkc\checkmarki\checkmarkr\checkmarkc\checkmark$\checkmarka\checkmark$\checkmark^\checkmark\\checkmarkc\checkmarki\checkmarkr\checkmarkc\checkmark$\checkmark \checkmark$\checkmark^\checkmark\\checkmarkc\checkmarki\checkmarkr\checkmarkc\checkmark$\checkmarkt\checkmark$\checkmark^\checkmark\\checkmarkc\checkmarki\checkmarkr\checkmarkc\checkmark$\checkmarke\checkmark$\checkmark^\checkmark\\checkmarkc\checkmarki\checkmarkr\checkmarkc\checkmark$\checkmarkn\checkmark$\checkmark^\checkmark\\checkmarkc\checkmarki\checkmarkr\checkmarkc\checkmark$\checkmarks\checkmark$\checkmark^\checkmark\\checkmarkc\checkmarki\checkmarkr\checkmarkc\checkmark$\checkmarki\checkmark$\checkmark^\checkmark\\checkmarkc\checkmarki\checkmarkr\checkmarkc\checkmark$\checkmarko\checkmark$\checkmark^\checkmark\\checkmarkc\checkmarki\checkmarkr\checkmarkc\checkmark$\checkmarkn\checkmark$\checkmark^\checkmark\\checkmarkc\checkmarki\checkmarkr\checkmarkc\checkmark$\checkmark \checkmark$\checkmark^\checkmark\\checkmarkc\checkmarki\checkmarkr\checkmarkc\checkmark$\checkmarkt\checkmark$\checkmark^\checkmark\\checkmarkc\checkmarki\checkmarkr\checkmarkc\checkmark$\checkmarka\checkmark$\checkmark^\checkmark\\checkmarkc\checkmarki\checkmarkr\checkmarkc\checkmark$\checkmarkn\checkmark$\checkmark^\checkmark\\checkmarkc\checkmarki\checkmarkr\checkmarkc\checkmark$\checkmarkg\checkmark$\checkmark^\checkmark\\checkmarkc\checkmarki\checkmarkr\checkmarkc\checkmark$\checkmarke\checkmark$\checkmark^\checkmark\\checkmarkc\checkmarki\checkmarkr\checkmarkc\checkmark$\checkmarkn\checkmark$\checkmark^\checkmark\\checkmarkc\checkmarki\checkmarkr\checkmarkc\checkmark$\checkmarkt\checkmark$\checkmark^\checkmark\\checkmarkc\checkmarki\checkmarkr\checkmarkc\checkmark$\checkmarki\checkmark$\checkmark^\checkmark\\checkmarkc\checkmarki\checkmarkr\checkmarkc\checkmark$\checkmarke\checkmark$\checkmark^\checkmark\\checkmarkc\checkmarki\checkmarkr\checkmarkc\checkmark$\checkmarkl\checkmark$\checkmark^\checkmark\\checkmarkc\checkmarki\checkmarkr\checkmarkc\checkmark$\checkmarkl\checkmark$\checkmark^\checkmark\\checkmarkc\checkmarki\checkmarkr\checkmarkc\checkmark$\checkmarke\checkmark$\checkmark^\checkmark\\checkmarkc\checkmarki\checkmarkr\checkmarkc\checkmark$\checkmark \checkmark$\checkmark^\checkmark\\checkmarkc\checkmarki\checkmarkr\checkmarkc\checkmark$\checkmarkd\checkmark$\checkmark^\checkmark\\checkmarkc\checkmarki\checkmarkr\checkmarkc\checkmark$\checkmarka\checkmark$\checkmark^\checkmark\\checkmarkc\checkmarki\checkmarkr\checkmarkc\checkmark$\checkmarkn\checkmark$\checkmark^\checkmark\\checkmarkc\checkmarki\checkmarkr\checkmarkc\checkmark$\checkmarks\checkmark$\checkmark^\checkmark\\checkmarkc\checkmarki\checkmarkr\checkmarkc\checkmark$\checkmark \checkmark$\checkmark^\checkmark\\checkmarkc\checkmarki\checkmarkr\checkmarkc\checkmark$\checkmarkl\checkmark$\checkmark^\checkmark\\checkmarkc\checkmarki\checkmarkr\checkmarkc\checkmark$\checkmarka\checkmark$\checkmark^\checkmark\\checkmarkc\checkmarki\checkmarkr\checkmarkc\checkmark$\checkmark \checkmark$\checkmark^\checkmark\\checkmarkc\checkmarki\checkmarkr\checkmarkc\checkmark$\checkmarkc\checkmark$\checkmark^\checkmark\\checkmarkc\checkmarki\checkmarkr\checkmarkc\checkmark$\checkmarko\checkmark$\checkmark^\checkmark\\checkmarkc\checkmarki\checkmarkr\checkmarkc\checkmark$\checkmarku\checkmark$\checkmark^\checkmark\\checkmarkc\checkmarki\checkmarkr\checkmarkc\checkmark$\checkmarkr\checkmark$\checkmark^\checkmark\\checkmarkc\checkmarki\checkmarkr\checkmarkc\checkmark$\checkmarkr\checkmark$\checkmark^\checkmark\\checkmarkc\checkmarki\checkmarkr\checkmarkc\checkmark$\checkmarko\checkmark$\checkmark^\checkmark\\checkmarkc\checkmarki\checkmarkr\checkmarkc\checkmark$\checkmarki\checkmark$\checkmark^\checkmark\\checkmarkc\checkmarki\checkmarkr\checkmarkc\checkmark$\checkmarke\checkmark$\checkmark^\checkmark\\checkmarkc\checkmarki\checkmarkr\checkmarkc\checkmark$\checkmark \checkmark$\checkmark^\checkmark\\checkmarkc\checkmarki\checkmarkr\checkmarkc\checkmark$\checkmarkd\checkmark$\checkmark^\checkmark\\checkmarkc\checkmarki\checkmarkr\checkmarkc\checkmark$\checkmarke\checkmark$\checkmark^\checkmark\\checkmarkc\checkmarki\checkmarkr\checkmarkc\checkmark$\checkmark \checkmark$\checkmark^\checkmark\\checkmarkc\checkmarki\checkmarkr\checkmarkc\checkmark$\checkmarkl\checkmark$\checkmark^\checkmark\\checkmarkc\checkmarki\checkmarkr\checkmarkc\checkmark$\checkmark'\checkmark$\checkmark^\checkmark\\checkmarkc\checkmarki\checkmarkr\checkmarkc\checkmark$\checkmarké\checkmark$\checkmark^\checkmark\\checkmarkc\checkmarki\checkmarkr\checkmarkc\checkmark$\checkmarkt\checkmark$\checkmark^\checkmark\\checkmarkc\checkmarki\checkmarkr\checkmarkc\checkmark$\checkmarka\checkmark$\checkmark^\checkmark\\checkmarkc\checkmarki\checkmarkr\checkmarkc\checkmark$\checkmarkg\checkmark$\checkmark^\checkmark\\checkmarkc\checkmarki\checkmarkr\checkmarkc\checkmark$\checkmarke\checkmark$\checkmark^\checkmark\\checkmarkc\checkmarki\checkmarkr\checkmarkc\checkmark$\checkmark \checkmark$\checkmark^\checkmark\\checkmarkc\checkmarki\checkmarkr\checkmarkc\checkmark$\checkmarkf\checkmark$\checkmark^\checkmark\\checkmarkc\checkmarki\checkmarkr\checkmarkc\checkmark$\checkmarki\checkmark$\checkmark^\checkmark\\checkmarkc\checkmarki\checkmarkr\checkmarkc\checkmark$\checkmarkn\checkmark$\checkmark^\checkmark\\checkmarkc\checkmarki\checkmarkr\checkmarkc\checkmark$\checkmarka\checkmark$\checkmark^\checkmark\\checkmarkc\checkmarki\checkmarkr\checkmarkc\checkmark$\checkmarkl\checkmark$\checkmark^\checkmark\\checkmarkc\checkmarki\checkmarkr\checkmarkc\checkmark$\checkmark \checkmark$\checkmark^\checkmark\\checkmarkc\checkmarki\checkmarkr\checkmarkc\checkmark$\checkmark(\checkmark$\checkmark^\checkmark\\checkmarkc\checkmarki\checkmarkr\checkmarkc\checkmark$\checkmarkr\checkmark$\checkmark^\checkmark\\checkmarkc\checkmarki\checkmarkr\checkmarkc\checkmark$\checkmarka\checkmark$\checkmark^\checkmark\\checkmarkc\checkmarki\checkmarkr\checkmarkc\checkmark$\checkmarky\checkmark$\checkmark^\checkmark\\checkmarkc\checkmarki\checkmarkr\checkmarkc\checkmark$\checkmarko\checkmark$\checkmark^\checkmark\\checkmarkc\checkmarki\checkmarkr\checkmarkc\checkmark$\checkmarkn\checkmark$\checkmark^\checkmark\\checkmarkc\checkmarki\checkmarkr\checkmarkc\checkmark$\checkmark \checkmark$\checkmark^\checkmark\\checkmarkc\checkmarki\checkmarkr\checkmarkc\checkmark$\checkmarkp\checkmark$\checkmark^\checkmark\\checkmarkc\checkmarki\checkmarkr\checkmarkc\checkmark$\checkmarkr\checkmark$\checkmark^\checkmark\\checkmarkc\checkmarki\checkmarkr\checkmarkc\checkmark$\checkmarki\checkmark$\checkmark^\checkmark\\checkmarkc\checkmarki\checkmarkr\checkmarkc\checkmark$\checkmarkm\checkmark$\checkmark^\checkmark\\checkmarkc\checkmarki\checkmarkr\checkmarkc\checkmark$\checkmarki\checkmark$\checkmark^\checkmark\\checkmarkc\checkmarki\checkmarkr\checkmarkc\checkmark$\checkmarkt\checkmark$\checkmark^\checkmark\\checkmarkc\checkmarki\checkmarkr\checkmarkc\checkmark$\checkmarki\checkmark$\checkmark^\checkmark\\checkmarkc\checkmarki\checkmarkr\checkmarkc\checkmark$\checkmarkf\checkmark$\checkmark^\checkmark\\checkmarkc\checkmarki\checkmarkr\checkmarkc\checkmark$\checkmark \checkmark$\checkmark^\checkmark\\checkmarkc\checkmarki\checkmarkr\checkmarkc\checkmark$\checkmarkd\checkmark$\checkmark^\checkmark\\checkmarkc\checkmarki\checkmarkr\checkmarkc\checkmark$\checkmarke\checkmark$\checkmark^\checkmark\\checkmarkc\checkmarki\checkmarkr\checkmarkc\checkmark$\checkmark \checkmark$\checkmark^\checkmark\\checkmarkc\checkmarki\checkmarkr\checkmarkc\checkmark$\checkmarkl\checkmark$\checkmark^\checkmark\\checkmarkc\checkmarki\checkmarkr\checkmarkc\checkmark$\checkmarka\checkmark$\checkmark^\checkmark\\checkmarkc\checkmarki\checkmarkr\checkmarkc\checkmark$\checkmark \checkmark$\checkmark^\checkmark\\checkmarkc\checkmarki\checkmarkr\checkmarkc\checkmark$\checkmarkp\checkmark$\checkmark^\checkmark\\checkmarkc\checkmarki\checkmarkr\checkmarkc\checkmark$\checkmarko\checkmark$\checkmark^\checkmark\\checkmarkc\checkmarki\checkmarkr\checkmarkc\checkmark$\checkmarku\checkmark$\checkmark^\checkmark\\checkmarkc\checkmarki\checkmarkr\checkmarkc\checkmark$\checkmarkl\checkmark$\checkmark^\checkmark\\checkmarkc\checkmarki\checkmarkr\checkmarkc\checkmark$\checkmarki\checkmark$\checkmark^\checkmark\\checkmarkc\checkmarki\checkmarkr\checkmarkc\checkmark$\checkmarke\checkmark$\checkmark^\checkmark\\checkmarkc\checkmarki\checkmarkr\checkmarkc\checkmark$\checkmark \checkmark$\checkmark^\checkmark\\checkmarkc\checkmarki\checkmarkr\checkmarkc\checkmark$\checkmarkm\checkmark$\checkmark^\checkmark\\checkmarkc\checkmarki\checkmarkr\checkmarkc\checkmark$\checkmarke\checkmark$\checkmark^\checkmark\\checkmarkc\checkmarki\checkmarkr\checkmarkc\checkmark$\checkmarkn\checkmark$\checkmark^\checkmark\\checkmarkc\checkmarki\checkmarkr\checkmarkc\checkmark$\checkmarké\checkmark$\checkmark^\checkmark\\checkmarkc\checkmarki\checkmarkr\checkmarkc\checkmark$\checkmarke\checkmark$\checkmark^\checkmark\\checkmarkc\checkmarki\checkmarkr\checkmarkc\checkmark$\checkmark \checkmark$\checkmark^\checkmark\\checkmarkc\checkmarki\checkmarkr\checkmarkc\checkmark$\checkmark$\checkmark$\checkmark^\checkmark\\checkmarkc\checkmarki\checkmarkr\checkmarkc\checkmark$\checkmarkr\checkmark$\checkmark^\checkmark\\checkmarkc\checkmarki\checkmarkr\checkmarkc\checkmark$\checkmark_\checkmark$\checkmark^\checkmark\\checkmarkc\checkmarki\checkmarkr\checkmarkc\checkmark$\checkmark4\checkmark$\checkmark^\checkmark\\checkmarkc\checkmarki\checkmarkr\checkmarkc\checkmark$\checkmark \checkmark$\checkmark^\checkmark\\checkmarkc\checkmarki\checkmarkr\checkmarkc\checkmark$\checkmark=\checkmark$\checkmark^\checkmark\\checkmarkc\checkmarki\checkmarkr\checkmarkc\checkmark$\checkmark \checkmark$\checkmark^\checkmark\\checkmarkc\checkmarki\checkmarkr\checkmarkc\checkmark$\checkmarkZ\checkmark$\checkmark^\checkmark\\checkmarkc\checkmarki\checkmarkr\checkmarkc\checkmark$\checkmark_\checkmark$\checkmark^\checkmark\\checkmarkc\checkmarki\checkmarkr\checkmarkc\checkmark$\checkmark4\checkmark$\checkmark^\checkmark\\checkmarkc\checkmarki\checkmarkr\checkmarkc\checkmark$\checkmark \checkmark$\checkmark^\checkmark\\checkmarkc\checkmarki\checkmarkr\checkmarkc\checkmark$\checkmark\\checkmark$\checkmark^\checkmark\\checkmarkc\checkmarki\checkmarkr\checkmarkc\checkmark$\checkmarkc\checkmark$\checkmark^\checkmark\\checkmarkc\checkmarki\checkmarkr\checkmarkc\checkmark$\checkmarkd\checkmark$\checkmark^\checkmark\\checkmarkc\checkmarki\checkmarkr\checkmarkc\checkmark$\checkmarko\checkmark$\checkmark^\checkmark\\checkmarkc\checkmarki\checkmarkr\checkmarkc\checkmark$\checkmarkt\checkmark$\checkmark^\checkmark\\checkmarkc\checkmarki\checkmarkr\checkmarkc\checkmark$\checkmark \checkmark$\checkmark^\checkmark\\checkmarkc\checkmarki\checkmarkr\checkmarkc\checkmark$\checkmarkm\checkmark$\checkmark^\checkmark\\checkmarkc\checkmarki\checkmarkr\checkmarkc\checkmark$\checkmark_\checkmark$\checkmark^\checkmark\\checkmarkc\checkmarki\checkmarkr\checkmarkc\checkmark$\checkmarkt\checkmark$\checkmark^\checkmark\\checkmarkc\checkmarki\checkmarkr\checkmarkc\checkmark$\checkmark \checkmark$\checkmark^\checkmark\\checkmarkc\checkmarki\checkmarkr\checkmarkc\checkmark$\checkmark/\checkmark$\checkmark^\checkmark\\checkmarkc\checkmarki\checkmarkr\checkmarkc\checkmark$\checkmark \checkmark$\checkmark^\checkmark\\checkmarkc\checkmarki\checkmarkr\checkmarkc\checkmark$\checkmark(\checkmark$\checkmark^\checkmark\\checkmarkc\checkmarki\checkmarkr\checkmarkc\checkmark$\checkmark2\checkmark$\checkmark^\checkmark\\checkmarkc\checkmarki\checkmarkr\checkmarkc\checkmark$\checkmark\\checkmark$\checkmark^\checkmark\\checkmarkc\checkmarki\checkmarkr\checkmarkc\checkmark$\checkmarkp\checkmark$\checkmark^\checkmark\\checkmarkc\checkmarki\checkmarkr\checkmarkc\checkmark$\checkmarki\checkmark$\checkmark^\checkmark\\checkmarkc\checkmarki\checkmarkr\checkmarkc\checkmark$\checkmark)\checkmark$\checkmark^\checkmark\\checkmarkc\checkmarki\checkmarkr\checkmarkc\checkmark$\checkmark \checkmark$\checkmark^\checkmark\\checkmarkc\checkmarki\checkmarkr\checkmarkc\checkmark$\checkmark=\checkmark$\checkmark^\checkmark\\checkmarkc\checkmarki\checkmarkr\checkmarkc\checkmark$\checkmark \checkmark$\checkmark^\checkmark\\checkmarkc\checkmarki\checkmarkr\checkmarkc\checkmark$\checkmark1\checkmark$\checkmark^\checkmark\\checkmarkc\checkmarki\checkmarkr\checkmarkc\checkmark$\checkmark0\checkmark$\checkmark^\checkmark\\checkmarkc\checkmarki\checkmarkr\checkmarkc\checkmark$\checkmark0\checkmark$\checkmark^\checkmark\\checkmarkc\checkmarki\checkmarkr\checkmarkc\checkmark$\checkmark \checkmark$\checkmark^\checkmark\\checkmarkc\checkmarki\checkmarkr\checkmarkc\checkmark$\checkmark\\checkmark$\checkmark^\checkmark\\checkmarkc\checkmarki\checkmarkr\checkmarkc\checkmark$\checkmarkt\checkmark$\checkmark^\checkmark\\checkmarkc\checkmarki\checkmarkr\checkmarkc\checkmark$\checkmarki\checkmark$\checkmark^\checkmark\\checkmarkc\checkmarki\checkmarkr\checkmarkc\checkmark$\checkmarkm\checkmark$\checkmark^\checkmark\\checkmarkc\checkmarki\checkmarkr\checkmarkc\checkmark$\checkmarke\checkmark$\checkmark^\checkmark\\checkmarkc\checkmarki\checkmarkr\checkmarkc\checkmark$\checkmarks\checkmark$\checkmark^\checkmark\\checkmarkc\checkmarki\checkmarkr\checkmarkc\checkmark$\checkmark \checkmark$\checkmark^\checkmark\\checkmarkc\checkmarki\checkmarkr\checkmarkc\checkmark$\checkmark5\checkmark$\checkmark^\checkmark\\checkmarkc\checkmarki\checkmarkr\checkmarkc\checkmark$\checkmark \checkmark$\checkmark^\checkmark\\checkmarkc\checkmarki\checkmarkr\checkmarkc\checkmark$\checkmark/\checkmark$\checkmark^\checkmark\\checkmarkc\checkmarki\checkmarkr\checkmarkc\checkmark$\checkmark \checkmark$\checkmark^\checkmark\\checkmarkc\checkmarki\checkmarkr\checkmarkc\checkmark$\checkmark(\checkmark$\checkmark^\checkmark\\checkmarkc\checkmarki\checkmarkr\checkmarkc\checkmark$\checkmark2\checkmark$\checkmark^\checkmark\\checkmarkc\checkmarki\checkmarkr\checkmarkc\checkmark$\checkmark\\checkmark$\checkmark^\checkmark\\checkmarkc\checkmarki\checkmarkr\checkmarkc\checkmark$\checkmarkp\checkmark$\checkmark^\checkmark\\checkmarkc\checkmarki\checkmarkr\checkmarkc\checkmark$\checkmarki\checkmark$\checkmark^\checkmark\\checkmarkc\checkmarki\checkmarkr\checkmarkc\checkmark$\checkmark)\checkmark$\checkmark^\checkmark\\checkmarkc\checkmarki\checkmarkr\checkmarkc\checkmark$\checkmark \checkmark$\checkmark^\checkmark\\checkmarkc\checkmarki\checkmarkr\checkmarkc\checkmark$\checkmark\\checkmark$\checkmark^\checkmark\\checkmarkc\checkmarki\checkmarkr\checkmarkc\checkmark$\checkmarka\checkmark$\checkmark^\checkmark\\checkmarkc\checkmarki\checkmarkr\checkmarkc\checkmark$\checkmarkp\checkmark$\checkmark^\checkmark\\checkmarkc\checkmarki\checkmarkr\checkmarkc\checkmark$\checkmarkp\checkmark$\checkmark^\checkmark\\checkmarkc\checkmarki\checkmarkr\checkmarkc\checkmark$\checkmarkr\checkmark$\checkmark^\checkmark\\checkmarkc\checkmarki\checkmarkr\checkmarkc\checkmark$\checkmarko\checkmark$\checkmark^\checkmark\\checkmarkc\checkmarki\checkmarkr\checkmarkc\checkmark$\checkmarkx\checkmark$\checkmark^\checkmark\\checkmarkc\checkmarki\checkmarkr\checkmarkc\checkmark$\checkmark \checkmark$\checkmark^\checkmark\\checkmarkc\checkmarki\checkmarkr\checkmarkc\checkmark$\checkmark7\checkmark$\checkmark^\checkmark\\checkmarkc\checkmarki\checkmarkr\checkmarkc\checkmark$\checkmark9\checkmark$\checkmark^\checkmark\\checkmarkc\checkmarki\checkmarkr\checkmarkc\checkmark$\checkmark{\checkmark$\checkmark^\checkmark\\checkmarkc\checkmarki\checkmarkr\checkmarkc\checkmark$\checkmark,\checkmark$\checkmark^\checkmark\\checkmarkc\checkmarki\checkmarkr\checkmarkc\checkmark$\checkmark}\checkmark$\checkmark^\checkmark\\checkmarkc\checkmarki\checkmarkr\checkmarkc\checkmark$\checkmark6\checkmark$\checkmark^\checkmark\\checkmarkc\checkmarki\checkmarkr\checkmarkc\checkmark$\checkmark$\checkmark$\checkmark^\checkmark\\checkmarkc\checkmarki\checkmarkr\checkmarkc\checkmark$\checkmark \checkmark$\checkmark^\checkmark\\checkmarkc\checkmarki\checkmarkr\checkmarkc\checkmark$\checkmarkm\checkmark$\checkmark^\checkmark\\checkmarkc\checkmarki\checkmarkr\checkmarkc\checkmark$\checkmarkm\checkmark$\checkmark^\checkmark\\checkmarkc\checkmarki\checkmarkr\checkmarkc\checkmark$\checkmark \checkmark$\checkmark^\checkmark\\checkmarkc\checkmarki\checkmarkr\checkmarkc\checkmark$\checkmark$\checkmark$\checkmark^\checkmark\\checkmarkc\checkmarki\checkmarkr\checkmarkc\checkmark$\checkmark\\checkmark$\checkmark^\checkmark\\checkmarkc\checkmarki\checkmarkr\checkmarkc\checkmark$\checkmarka\checkmark$\checkmark^\checkmark\\checkmarkc\checkmarki\checkmarkr\checkmarkc\checkmark$\checkmarkp\checkmark$\checkmark^\checkmark\\checkmarkc\checkmarki\checkmarkr\checkmarkc\checkmark$\checkmarkp\checkmark$\checkmark^\checkmark\\checkmarkc\checkmarki\checkmarkr\checkmarkc\checkmark$\checkmarkr\checkmark$\checkmark^\checkmark\\checkmarkc\checkmarki\checkmarkr\checkmarkc\checkmark$\checkmarko\checkmark$\checkmark^\checkmark\\checkmarkc\checkmarki\checkmarkr\checkmarkc\checkmark$\checkmarkx\checkmark$\checkmark^\checkmark\\checkmarkc\checkmarki\checkmarkr\checkmarkc\checkmark$\checkmark \checkmark$\checkmark^\checkmark\\checkmarkc\checkmarki\checkmarkr\checkmarkc\checkmark$\checkmark0\checkmark$\checkmark^\checkmark\\checkmarkc\checkmarki\checkmarkr\checkmarkc\checkmark$\checkmark{\checkmark$\checkmark^\checkmark\\checkmarkc\checkmarki\checkmarkr\checkmarkc\checkmark$\checkmark,\checkmark$\checkmark^\checkmark\\checkmarkc\checkmarki\checkmarkr\checkmarkc\checkmark$\checkmark}\checkmark$\checkmark^\checkmark\\checkmarkc\checkmarki\checkmarkr\checkmarkc\checkmark$\checkmark0\checkmark$\checkmark^\checkmark\\checkmarkc\checkmarki\checkmarkr\checkmarkc\checkmark$\checkmark8\checkmark$\checkmark^\checkmark\\checkmarkc\checkmarki\checkmarkr\checkmarkc\checkmark$\checkmark$\checkmark$\checkmark^\checkmark\\checkmarkc\checkmarki\checkmarkr\checkmarkc\checkmark$\checkmark \checkmark$\checkmark^\checkmark\\checkmarkc\checkmarki\checkmarkr\checkmarkc\checkmark$\checkmarkm\checkmark$\checkmark^\checkmark\\checkmarkc\checkmarki\checkmarkr\checkmarkc\checkmark$\checkmark)\checkmark$\checkmark^\checkmark\\checkmarkc\checkmarki\checkmarkr\checkmarkc\checkmark$\checkmark \checkmark$\checkmark^\checkmark\\checkmarkc\checkmarki\checkmarkr\checkmarkc\checkmark$\checkmark:\checkmark$\checkmark^\checkmark\\checkmarkc\checkmarki\checkmarkr\checkmarkc\checkmark$\checkmark
\checkmark$\checkmark^\checkmark\\checkmarkc\checkmarki\checkmarkr\checkmarkc\checkmark$\checkmark\\checkmark$\checkmark^\checkmark\\checkmarkc\checkmarki\checkmarkr\checkmarkc\checkmark$\checkmarkb\checkmark$\checkmark^\checkmark\\checkmarkc\checkmarki\checkmarkr\checkmarkc\checkmark$\checkmarke\checkmark$\checkmark^\checkmark\\checkmarkc\checkmarki\checkmarkr\checkmarkc\checkmark$\checkmarkg\checkmark$\checkmark^\checkmark\\checkmarkc\checkmarki\checkmarkr\checkmarkc\checkmark$\checkmarki\checkmark$\checkmark^\checkmark\\checkmarkc\checkmarki\checkmarkr\checkmarkc\checkmark$\checkmarkn\checkmark$\checkmark^\checkmark\\checkmarkc\checkmarki\checkmarkr\checkmarkc\checkmark$\checkmark{\checkmark$\checkmark^\checkmark\\checkmarkc\checkmarki\checkmarkr\checkmarkc\checkmark$\checkmarke\checkmark$\checkmark^\checkmark\\checkmarkc\checkmarki\checkmarkr\checkmarkc\checkmark$\checkmarkq\checkmark$\checkmark^\checkmark\\checkmarkc\checkmarki\checkmarkr\checkmarkc\checkmark$\checkmarku\checkmark$\checkmark^\checkmark\\checkmarkc\checkmarki\checkmarkr\checkmarkc\checkmark$\checkmarka\checkmark$\checkmark^\checkmark\\checkmarkc\checkmarki\checkmarkr\checkmarkc\checkmark$\checkmarkt\checkmark$\checkmark^\checkmark\\checkmarkc\checkmarki\checkmarkr\checkmarkc\checkmark$\checkmarki\checkmark$\checkmark^\checkmark\\checkmarkc\checkmarki\checkmarkr\checkmarkc\checkmark$\checkmarko\checkmark$\checkmark^\checkmark\\checkmarkc\checkmarki\checkmarkr\checkmarkc\checkmark$\checkmarkn\checkmark$\checkmark^\checkmark\\checkmarkc\checkmarki\checkmarkr\checkmarkc\checkmark$\checkmark}\checkmark$\checkmark^\checkmark\\checkmarkc\checkmarki\checkmarkr\checkmarkc\checkmark$\checkmark
\checkmark$\checkmark^\checkmark\\checkmarkc\checkmarki\checkmarkr\checkmarkc\checkmark$\checkmark \checkmark$\checkmark^\checkmark\\checkmarkc\checkmarki\checkmarkr\checkmarkc\checkmark$\checkmark \checkmark$\checkmark^\checkmark\\checkmarkc\checkmarki\checkmarkr\checkmarkc\checkmark$\checkmark \checkmark$\checkmark^\checkmark\\checkmarkc\checkmarki\checkmarkr\checkmarkc\checkmark$\checkmark \checkmark$\checkmark^\checkmark\\checkmarkc\checkmarki\checkmarkr\checkmarkc\checkmark$\checkmarkF\checkmark$\checkmark^\checkmark\\checkmarkc\checkmarki\checkmarkr\checkmarkc\checkmark$\checkmark_\checkmark$\checkmark^\checkmark\\checkmarkc\checkmarki\checkmarkr\checkmarkc\checkmark$\checkmark{\checkmark$\checkmark^\checkmark\\checkmarkc\checkmarki\checkmarkr\checkmarkc\checkmark$\checkmarkc\checkmark$\checkmark^\checkmark\\checkmarkc\checkmarki\checkmarkr\checkmarkc\checkmark$\checkmarko\checkmark$\checkmark^\checkmark\\checkmarkc\checkmarki\checkmarkr\checkmarkc\checkmark$\checkmarku\checkmark$\checkmark^\checkmark\\checkmarkc\checkmarki\checkmarkr\checkmarkc\checkmark$\checkmarkr\checkmark$\checkmark^\checkmark\\checkmarkc\checkmarki\checkmarkr\checkmarkc\checkmark$\checkmarkr\checkmark$\checkmark^\checkmark\\checkmarkc\checkmarki\checkmarkr\checkmarkc\checkmark$\checkmarko\checkmark$\checkmark^\checkmark\\checkmarkc\checkmarki\checkmarkr\checkmarkc\checkmark$\checkmarki\checkmark$\checkmark^\checkmark\\checkmarkc\checkmarki\checkmarkr\checkmarkc\checkmark$\checkmarke\checkmark$\checkmark^\checkmark\\checkmarkc\checkmarki\checkmarkr\checkmarkc\checkmark$\checkmark}\checkmark$\checkmark^\checkmark\\checkmarkc\checkmarki\checkmarkr\checkmarkc\checkmark$\checkmark \checkmark$\checkmark^\checkmark\\checkmarkc\checkmarki\checkmarkr\checkmarkc\checkmark$\checkmark=\checkmark$\checkmark^\checkmark\\checkmarkc\checkmarki\checkmarkr\checkmarkc\checkmark$\checkmark \checkmark$\checkmark^\checkmark\\checkmarkc\checkmarki\checkmarkr\checkmarkc\checkmark$\checkmark\\checkmark$\checkmark^\checkmark\\checkmarkc\checkmarki\checkmarkr\checkmarkc\checkmark$\checkmarkf\checkmark$\checkmark^\checkmark\\checkmarkc\checkmarki\checkmarkr\checkmarkc\checkmark$\checkmarkr\checkmark$\checkmark^\checkmark\\checkmarkc\checkmarki\checkmarkr\checkmarkc\checkmark$\checkmarka\checkmark$\checkmark^\checkmark\\checkmarkc\checkmarki\checkmarkr\checkmarkc\checkmark$\checkmarkc\checkmark$\checkmark^\checkmark\\checkmarkc\checkmarki\checkmarkr\checkmarkc\checkmark$\checkmark{\checkmark$\checkmark^\checkmark\\checkmarkc\checkmarki\checkmarkr\checkmarkc\checkmark$\checkmarkM\checkmark$\checkmark^\checkmark\\checkmarkc\checkmarki\checkmarkr\checkmarkc\checkmark$\checkmark_\checkmark$\checkmark^\checkmark\\checkmarkc\checkmarki\checkmarkr\checkmarkc\checkmark$\checkmark{\checkmark$\checkmark^\checkmark\\checkmarkc\checkmarki\checkmarkr\checkmarkc\checkmark$\checkmarkr\checkmark$\checkmark^\checkmark\\checkmarkc\checkmarki\checkmarkr\checkmarkc\checkmark$\checkmarke\checkmark$\checkmark^\checkmark\\checkmarkc\checkmarki\checkmarkr\checkmarkc\checkmark$\checkmarkn\checkmark$\checkmark^\checkmark\\checkmarkc\checkmarki\checkmarkr\checkmarkc\checkmark$\checkmarkv\checkmark$\checkmark^\checkmark\\checkmarkc\checkmarki\checkmarkr\checkmarkc\checkmark$\checkmark}\checkmark$\checkmark^\checkmark\\checkmarkc\checkmarki\checkmarkr\checkmarkc\checkmark$\checkmark}\checkmark$\checkmark^\checkmark\\checkmarkc\checkmarki\checkmarkr\checkmarkc\checkmark$\checkmark{\checkmark$\checkmark^\checkmark\\checkmarkc\checkmarki\checkmarkr\checkmarkc\checkmark$\checkmarkr\checkmark$\checkmark^\checkmark\\checkmarkc\checkmarki\checkmarkr\checkmarkc\checkmark$\checkmark_\checkmark$\checkmark^\checkmark\\checkmarkc\checkmarki\checkmarkr\checkmarkc\checkmark$\checkmark4\checkmark$\checkmark^\checkmark\\checkmarkc\checkmarki\checkmarkr\checkmarkc\checkmark$\checkmark}\checkmark$\checkmark^\checkmark\\checkmarkc\checkmarki\checkmarkr\checkmarkc\checkmark$\checkmark \checkmark$\checkmark^\checkmark\\checkmarkc\checkmarki\checkmarkr\checkmarkc\checkmark$\checkmark=\checkmark$\checkmark^\checkmark\\checkmarkc\checkmarki\checkmarkr\checkmarkc\checkmark$\checkmark \checkmark$\checkmark^\checkmark\\checkmarkc\checkmarki\checkmarkr\checkmarkc\checkmark$\checkmark\\checkmark$\checkmark^\checkmark\\checkmarkc\checkmarki\checkmarkr\checkmarkc\checkmark$\checkmarkf\checkmark$\checkmark^\checkmark\\checkmarkc\checkmarki\checkmarkr\checkmarkc\checkmark$\checkmarkr\checkmark$\checkmark^\checkmark\\checkmarkc\checkmarki\checkmarkr\checkmarkc\checkmark$\checkmarka\checkmark$\checkmark^\checkmark\\checkmarkc\checkmarki\checkmarkr\checkmarkc\checkmark$\checkmarkc\checkmark$\checkmark^\checkmark\\checkmarkc\checkmarki\checkmarkr\checkmarkc\checkmark$\checkmark{\checkmark$\checkmark^\checkmark\\checkmarkc\checkmarki\checkmarkr\checkmarkc\checkmark$\checkmark2\checkmark$\checkmark^\checkmark\\checkmarkc\checkmarki\checkmarkr\checkmarkc\checkmark$\checkmark0\checkmark$\checkmark^\checkmark\\checkmarkc\checkmarki\checkmarkr\checkmarkc\checkmark$\checkmark}\checkmark$\checkmark^\checkmark\\checkmarkc\checkmarki\checkmarkr\checkmarkc\checkmark$\checkmark{\checkmark$\checkmark^\checkmark\\checkmarkc\checkmarki\checkmarkr\checkmarkc\checkmark$\checkmark0\checkmark$\checkmark^\checkmark\\checkmarkc\checkmarki\checkmarkr\checkmarkc\checkmark$\checkmark{\checkmark$\checkmark^\checkmark\\checkmarkc\checkmarki\checkmarkr\checkmarkc\checkmark$\checkmark,\checkmark$\checkmark^\checkmark\\checkmarkc\checkmarki\checkmarkr\checkmarkc\checkmark$\checkmark}\checkmark$\checkmark^\checkmark\\checkmarkc\checkmarki\checkmarkr\checkmarkc\checkmark$\checkmark0\checkmark$\checkmark^\checkmark\\checkmarkc\checkmarki\checkmarkr\checkmarkc\checkmark$\checkmark8\checkmark$\checkmark^\checkmark\\checkmarkc\checkmarki\checkmarkr\checkmarkc\checkmark$\checkmark}\checkmark$\checkmark^\checkmark\\checkmarkc\checkmarki\checkmarkr\checkmarkc\checkmark$\checkmark \checkmark$\checkmark^\checkmark\\checkmarkc\checkmarki\checkmarkr\checkmarkc\checkmark$\checkmark=\checkmark$\checkmark^\checkmark\\checkmarkc\checkmarki\checkmarkr\checkmarkc\checkmark$\checkmark \checkmark$\checkmark^\checkmark\\checkmarkc\checkmarki\checkmarkr\checkmarkc\checkmark$\checkmark\\checkmark$\checkmark^\checkmark\\checkmarkc\checkmarki\checkmarkr\checkmarkc\checkmark$\checkmarkb\checkmark$\checkmark^\checkmark\\checkmarkc\checkmarki\checkmarkr\checkmarkc\checkmark$\checkmarko\checkmark$\checkmark^\checkmark\\checkmarkc\checkmarki\checkmarkr\checkmarkc\checkmark$\checkmarkx\checkmark$\checkmark^\checkmark\\checkmarkc\checkmarki\checkmarkr\checkmarkc\checkmark$\checkmarke\checkmark$\checkmark^\checkmark\\checkmarkc\checkmarki\checkmarkr\checkmarkc\checkmark$\checkmarkd\checkmark$\checkmark^\checkmark\\checkmarkc\checkmarki\checkmarkr\checkmarkc\checkmark$\checkmark{\checkmark$\checkmark^\checkmark\\checkmarkc\checkmarki\checkmarkr\checkmarkc\checkmark$\checkmark2\checkmark$\checkmark^\checkmark\\checkmarkc\checkmarki\checkmarkr\checkmarkc\checkmark$\checkmark5\checkmark$\checkmark^\checkmark\\checkmarkc\checkmarki\checkmarkr\checkmarkc\checkmark$\checkmark0\checkmark$\checkmark^\checkmark\\checkmarkc\checkmarki\checkmarkr\checkmarkc\checkmark$\checkmark \checkmark$\checkmark^\checkmark\\checkmarkc\checkmarki\checkmarkr\checkmarkc\checkmark$\checkmark\\checkmark$\checkmark^\checkmark\\checkmarkc\checkmarki\checkmarkr\checkmarkc\checkmark$\checkmarkt\checkmark$\checkmark^\checkmark\\checkmarkc\checkmarki\checkmarkr\checkmarkc\checkmark$\checkmarke\checkmark$\checkmark^\checkmark\\checkmarkc\checkmarki\checkmarkr\checkmarkc\checkmark$\checkmarkx\checkmark$\checkmark^\checkmark\\checkmarkc\checkmarki\checkmarkr\checkmarkc\checkmark$\checkmarkt\checkmark$\checkmark^\checkmark\\checkmarkc\checkmarki\checkmarkr\checkmarkc\checkmark$\checkmark{\checkmark$\checkmark^\checkmark\\checkmarkc\checkmarki\checkmarkr\checkmarkc\checkmark$\checkmark \checkmark$\checkmark^\checkmark\\checkmarkc\checkmarki\checkmarkr\checkmarkc\checkmark$\checkmarkN\checkmark$\checkmark^\checkmark\\checkmarkc\checkmarki\checkmarkr\checkmarkc\checkmark$\checkmark}\checkmark$\checkmark^\checkmark\\checkmarkc\checkmarki\checkmarkr\checkmarkc\checkmark$\checkmark}\checkmark$\checkmark^\checkmark\\checkmarkc\checkmarki\checkmarkr\checkmarkc\checkmark$\checkmark
\checkmark$\checkmark^\checkmark\\checkmarkc\checkmarki\checkmarkr\checkmarkc\checkmark$\checkmark\\checkmark$\checkmark^\checkmark\\checkmarkc\checkmarki\checkmarkr\checkmarkc\checkmark$\checkmarke\checkmark$\checkmark^\checkmark\\checkmarkc\checkmarki\checkmarkr\checkmarkc\checkmark$\checkmarkn\checkmark$\checkmark^\checkmark\\checkmarkc\checkmarki\checkmarkr\checkmarkc\checkmark$\checkmarkd\checkmark$\checkmark^\checkmark\\checkmarkc\checkmarki\checkmarkr\checkmarkc\checkmark$\checkmark{\checkmark$\checkmark^\checkmark\\checkmarkc\checkmarki\checkmarkr\checkmarkc\checkmark$\checkmarke\checkmark$\checkmark^\checkmark\\checkmarkc\checkmarki\checkmarkr\checkmarkc\checkmark$\checkmarkq\checkmark$\checkmark^\checkmark\\checkmarkc\checkmarki\checkmarkr\checkmarkc\checkmark$\checkmarku\checkmark$\checkmark^\checkmark\\checkmarkc\checkmarki\checkmarkr\checkmarkc\checkmark$\checkmarka\checkmark$\checkmark^\checkmark\\checkmarkc\checkmarki\checkmarkr\checkmarkc\checkmark$\checkmarkt\checkmark$\checkmark^\checkmark\\checkmarkc\checkmarki\checkmarkr\checkmarkc\checkmark$\checkmarki\checkmark$\checkmark^\checkmark\\checkmarkc\checkmarki\checkmarkr\checkmarkc\checkmark$\checkmarko\checkmark$\checkmark^\checkmark\\checkmarkc\checkmarki\checkmarkr\checkmarkc\checkmark$\checkmarkn\checkmark$\checkmark^\checkmark\\checkmarkc\checkmarki\checkmarkr\checkmarkc\checkmark$\checkmark}\checkmark$\checkmark^\checkmark\\checkmarkc\checkmarki\checkmarkr\checkmarkc\checkmark$\checkmark
\checkmark$\checkmark^\checkmark\\checkmarkc\checkmarki\checkmarkr\checkmarkc\checkmark$\checkmark
\checkmark$\checkmark^\checkmark\\checkmarkc\checkmarki\checkmarkr\checkmarkc\checkmark$\checkmarkE\checkmark$\checkmark^\checkmark\\checkmarkc\checkmarki\checkmarkr\checkmarkc\checkmark$\checkmarkn\checkmark$\checkmark^\checkmark\\checkmarkc\checkmarki\checkmarkr\checkmarkc\checkmark$\checkmark \checkmark$\checkmark^\checkmark\\checkmarkc\checkmarki\checkmarkr\checkmarkc\checkmark$\checkmarkc\checkmark$\checkmark^\checkmark\\checkmarkc\checkmarki\checkmarkr\checkmarkc\checkmark$\checkmarko\checkmark$\checkmark^\checkmark\\checkmarkc\checkmarki\checkmarkr\checkmarkc\checkmark$\checkmarkn\checkmark$\checkmark^\checkmark\\checkmarkc\checkmarki\checkmarkr\checkmarkc\checkmark$\checkmarks\checkmark$\checkmark^\checkmark\\checkmarkc\checkmarki\checkmarkr\checkmarkc\checkmark$\checkmarki\checkmark$\checkmark^\checkmark\\checkmarkc\checkmarki\checkmarkr\checkmarkc\checkmark$\checkmarkd\checkmark$\checkmark^\checkmark\\checkmarkc\checkmarki\checkmarkr\checkmarkc\checkmark$\checkmarké\checkmark$\checkmark^\checkmark\\checkmarkc\checkmarki\checkmarkr\checkmarkc\checkmark$\checkmarkr\checkmark$\checkmark^\checkmark\\checkmarkc\checkmarki\checkmarkr\checkmarkc\checkmark$\checkmarka\checkmark$\checkmark^\checkmark\\checkmarkc\checkmarki\checkmarkr\checkmarkc\checkmark$\checkmarkn\checkmark$\checkmark^\checkmark\\checkmarkc\checkmarki\checkmarkr\checkmarkc\checkmark$\checkmarkt\checkmark$\checkmark^\checkmark\\checkmarkc\checkmarki\checkmarkr\checkmarkc\checkmark$\checkmark \checkmark$\checkmark^\checkmark\\checkmarkc\checkmarki\checkmarkr\checkmarkc\checkmark$\checkmarkl\checkmark$\checkmark^\checkmark\\checkmarkc\checkmarki\checkmarkr\checkmarkc\checkmark$\checkmarke\checkmark$\checkmark^\checkmark\\checkmarkc\checkmarki\checkmarkr\checkmarkc\checkmark$\checkmark \checkmark$\checkmark^\checkmark\\checkmarkc\checkmarki\checkmarkr\checkmarkc\checkmark$\checkmarkr\checkmark$\checkmark^\checkmark\\checkmarkc\checkmarki\checkmarkr\checkmarkc\checkmark$\checkmarka\checkmark$\checkmark^\checkmark\\checkmarkc\checkmarki\checkmarkr\checkmarkc\checkmark$\checkmarkt\checkmark$\checkmark^\checkmark\\checkmarkc\checkmarki\checkmarkr\checkmarkc\checkmark$\checkmarki\checkmark$\checkmark^\checkmark\\checkmarkc\checkmarki\checkmarkr\checkmarkc\checkmark$\checkmarko\checkmark$\checkmark^\checkmark\\checkmarkc\checkmarki\checkmarkr\checkmarkc\checkmark$\checkmark \checkmark$\checkmark^\checkmark\\checkmarkc\checkmarki\checkmarkr\checkmarkc\checkmark$\checkmarkd\checkmark$\checkmark^\checkmark\\checkmarkc\checkmarki\checkmarkr\checkmarkc\checkmark$\checkmarke\checkmark$\checkmark^\checkmark\\checkmarkc\checkmarki\checkmarkr\checkmarkc\checkmark$\checkmark \checkmark$\checkmark^\checkmark\\checkmarkc\checkmarki\checkmarkr\checkmarkc\checkmark$\checkmarkt\checkmark$\checkmark^\checkmark\\checkmarkc\checkmarki\checkmarkr\checkmarkc\checkmark$\checkmarke\checkmark$\checkmark^\checkmark\\checkmarkc\checkmarki\checkmarkr\checkmarkc\checkmark$\checkmarkn\checkmark$\checkmark^\checkmark\\checkmarkc\checkmarki\checkmarkr\checkmarkc\checkmark$\checkmarks\checkmark$\checkmark^\checkmark\\checkmarkc\checkmarki\checkmarkr\checkmarkc\checkmark$\checkmarki\checkmark$\checkmark^\checkmark\\checkmarkc\checkmarki\checkmarkr\checkmarkc\checkmark$\checkmarko\checkmark$\checkmark^\checkmark\\checkmarkc\checkmarki\checkmarkr\checkmarkc\checkmark$\checkmarkn\checkmark$\checkmark^\checkmark\\checkmarkc\checkmarki\checkmarkr\checkmarkc\checkmark$\checkmark \checkmark$\checkmark^\checkmark\\checkmarkc\checkmarki\checkmarkr\checkmarkc\checkmark$\checkmarkc\checkmark$\checkmark^\checkmark\\checkmarkc\checkmarki\checkmarkr\checkmarkc\checkmark$\checkmarko\checkmark$\checkmark^\checkmark\\checkmarkc\checkmarki\checkmarkr\checkmarkc\checkmark$\checkmarku\checkmark$\checkmark^\checkmark\\checkmarkc\checkmarki\checkmarkr\checkmarkc\checkmark$\checkmarkr\checkmark$\checkmark^\checkmark\\checkmarkc\checkmarki\checkmarkr\checkmarkc\checkmark$\checkmarkr\checkmark$\checkmark^\checkmark\\checkmarkc\checkmarki\checkmarkr\checkmarkc\checkmark$\checkmarko\checkmark$\checkmark^\checkmark\\checkmarkc\checkmarki\checkmarkr\checkmarkc\checkmark$\checkmarki\checkmark$\checkmark^\checkmark\\checkmarkc\checkmarki\checkmarkr\checkmarkc\checkmark$\checkmarke\checkmark$\checkmark^\checkmark\\checkmarkc\checkmarki\checkmarkr\checkmarkc\checkmark$\checkmark \checkmark$\checkmark^\checkmark\\checkmarkc\checkmarki\checkmarkr\checkmarkc\checkmark$\checkmark(\checkmark$\checkmark^\checkmark\\checkmarkc\checkmarki\checkmarkr\checkmarkc\checkmark$\checkmarkc\checkmark$\checkmark^\checkmark\\checkmarkc\checkmarki\checkmarkr\checkmarkc\checkmark$\checkmarkô\checkmark$\checkmark^\checkmark\\checkmarkc\checkmarki\checkmarkr\checkmarkc\checkmark$\checkmarkt\checkmark$\checkmark^\checkmark\\checkmarkc\checkmarki\checkmarkr\checkmarkc\checkmark$\checkmarké\checkmark$\checkmark^\checkmark\\checkmarkc\checkmarki\checkmarkr\checkmarkc\checkmark$\checkmark \checkmark$\checkmark^\checkmark\\checkmarkc\checkmarki\checkmarkr\checkmarkc\checkmark$\checkmarkt\checkmark$\checkmark^\checkmark\\checkmarkc\checkmarki\checkmarkr\checkmarkc\checkmark$\checkmarke\checkmark$\checkmark^\checkmark\\checkmarkc\checkmarki\checkmarkr\checkmarkc\checkmark$\checkmarkn\checkmark$\checkmark^\checkmark\\checkmarkc\checkmarki\checkmarkr\checkmarkc\checkmark$\checkmarkd\checkmark$\checkmark^\checkmark\\checkmarkc\checkmarki\checkmarkr\checkmarkc\checkmark$\checkmarka\checkmark$\checkmark^\checkmark\\checkmarkc\checkmarki\checkmarkr\checkmarkc\checkmark$\checkmarkn\checkmark$\checkmark^\checkmark\\checkmarkc\checkmarki\checkmarkr\checkmarkc\checkmark$\checkmarkt\checkmark$\checkmark^\checkmark\\checkmarkc\checkmarki\checkmarkr\checkmarkc\checkmark$\checkmark/\checkmark$\checkmark^\checkmark\\checkmarkc\checkmarki\checkmarkr\checkmarkc\checkmark$\checkmarkc\checkmark$\checkmark^\checkmark\\checkmarkc\checkmarki\checkmarkr\checkmarkc\checkmark$\checkmarkô\checkmark$\checkmark^\checkmark\\checkmarkc\checkmarki\checkmarkr\checkmarkc\checkmark$\checkmarkt\checkmark$\checkmark^\checkmark\\checkmarkc\checkmarki\checkmarkr\checkmarkc\checkmark$\checkmarké\checkmark$\checkmark^\checkmark\\checkmarkc\checkmarki\checkmarkr\checkmarkc\checkmark$\checkmark \checkmark$\checkmark^\checkmark\\checkmarkc\checkmarki\checkmarkr\checkmarkc\checkmark$\checkmarkm\checkmark$\checkmark^\checkmark\\checkmarkc\checkmarki\checkmarkr\checkmarkc\checkmark$\checkmarko\checkmark$\checkmark^\checkmark\\checkmarkc\checkmarki\checkmarkr\checkmarkc\checkmark$\checkmarku\checkmark$\checkmark^\checkmark\\checkmarkc\checkmarki\checkmarkr\checkmarkc\checkmark$\checkmark)\checkmark$\checkmark^\checkmark\\checkmarkc\checkmarki\checkmarkr\checkmarkc\checkmark$\checkmark \checkmark$\checkmark^\checkmark\\checkmarkc\checkmarki\checkmarkr\checkmarkc\checkmark$\checkmark$\checkmark$\checkmark^\checkmark\\checkmarkc\checkmarki\checkmarkr\checkmarkc\checkmark$\checkmark\\checkmark$\checkmark^\checkmark\\checkmarkc\checkmarki\checkmarkr\checkmarkc\checkmark$\checkmarka\checkmark$\checkmark^\checkmark\\checkmarkc\checkmarki\checkmarkr\checkmarkc\checkmark$\checkmarkp\checkmark$\checkmark^\checkmark\\checkmarkc\checkmarki\checkmarkr\checkmarkc\checkmark$\checkmarkp\checkmark$\checkmark^\checkmark\\checkmarkc\checkmarki\checkmarkr\checkmarkc\checkmark$\checkmarkr\checkmark$\checkmark^\checkmark\\checkmarkc\checkmarki\checkmarkr\checkmarkc\checkmark$\checkmarko\checkmark$\checkmark^\checkmark\\checkmarkc\checkmarki\checkmarkr\checkmarkc\checkmark$\checkmarkx\checkmark$\checkmark^\checkmark\\checkmarkc\checkmarki\checkmarkr\checkmarkc\checkmark$\checkmark \checkmark$\checkmark^\checkmark\\checkmarkc\checkmarki\checkmarkr\checkmarkc\checkmark$\checkmark3\checkmark$\checkmark^\checkmark\\checkmarkc\checkmarki\checkmarkr\checkmarkc\checkmark$\checkmark$\checkmark$\checkmark^\checkmark\\checkmarkc\checkmarki\checkmarkr\checkmarkc\checkmark$\checkmark \checkmark$\checkmark^\checkmark\\checkmarkc\checkmarki\checkmarkr\checkmarkc\checkmark$\checkmarkp\checkmark$\checkmark^\checkmark\\checkmarkc\checkmarki\checkmarkr\checkmarkc\checkmark$\checkmarko\checkmark$\checkmark^\checkmark\\checkmarkc\checkmarki\checkmarkr\checkmarkc\checkmark$\checkmarku\checkmark$\checkmark^\checkmark\\checkmarkc\checkmarki\checkmarkr\checkmarkc\checkmark$\checkmarkr\checkmark$\checkmark^\checkmark\\checkmarkc\checkmarki\checkmarkr\checkmarkc\checkmark$\checkmark \checkmark$\checkmark^\checkmark\\checkmarkc\checkmarki\checkmarkr\checkmarkc\checkmark$\checkmarkl\checkmark$\checkmark^\checkmark\\checkmarkc\checkmarki\checkmarkr\checkmarkc\checkmark$\checkmarke\checkmark$\checkmark^\checkmark\\checkmarkc\checkmarki\checkmarkr\checkmarkc\checkmark$\checkmarks\checkmark$\checkmark^\checkmark\\checkmarkc\checkmarki\checkmarkr\checkmarkc\checkmark$\checkmark \checkmark$\checkmark^\checkmark\\checkmarkc\checkmarki\checkmarkr\checkmarkc\checkmark$\checkmarkc\checkmark$\checkmark^\checkmark\\checkmarkc\checkmarki\checkmarkr\checkmarkc\checkmark$\checkmarko\checkmark$\checkmark^\checkmark\\checkmarkc\checkmarki\checkmarkr\checkmarkc\checkmark$\checkmarku\checkmark$\checkmark^\checkmark\\checkmarkc\checkmarki\checkmarkr\checkmarkc\checkmark$\checkmarkr\checkmark$\checkmark^\checkmark\\checkmarkc\checkmarki\checkmarkr\checkmarkc\checkmark$\checkmarkr\checkmark$\checkmark^\checkmark\\checkmarkc\checkmarki\checkmarkr\checkmarkc\checkmark$\checkmarko\checkmark$\checkmark^\checkmark\\checkmarkc\checkmarki\checkmarkr\checkmarkc\checkmark$\checkmarki\checkmark$\checkmark^\checkmark\\checkmarkc\checkmarki\checkmarkr\checkmarkc\checkmark$\checkmarke\checkmark$\checkmark^\checkmark\\checkmarkc\checkmarki\checkmarkr\checkmarkc\checkmark$\checkmarks\checkmark$\checkmark^\checkmark\\checkmarkc\checkmarki\checkmarkr\checkmarkc\checkmark$\checkmark \checkmark$\checkmark^\checkmark\\checkmarkc\checkmarki\checkmarkr\checkmarkc\checkmark$\checkmarkH\checkmark$\checkmark^\checkmark\\checkmarkc\checkmarki\checkmarkr\checkmarkc\checkmark$\checkmarkT\checkmark$\checkmark^\checkmark\\checkmarkc\checkmarki\checkmarkr\checkmarkc\checkmark$\checkmarkD\checkmark$\checkmark^\checkmark\\checkmarkc\checkmarki\checkmarkr\checkmarkc\checkmark$\checkmark \checkmark$\checkmark^\checkmark\\checkmarkc\checkmarki\checkmarkr\checkmarkc\checkmark$\checkmark:\checkmark$\checkmark^\checkmark\\checkmarkc\checkmarki\checkmarkr\checkmarkc\checkmark$\checkmark
\checkmark$\checkmark^\checkmark\\checkmarkc\checkmarki\checkmarkr\checkmarkc\checkmark$\checkmark\\checkmark$\checkmark^\checkmark\\checkmarkc\checkmarki\checkmarkr\checkmarkc\checkmark$\checkmarkb\checkmark$\checkmark^\checkmark\\checkmarkc\checkmarki\checkmarkr\checkmarkc\checkmark$\checkmarke\checkmark$\checkmark^\checkmark\\checkmarkc\checkmarki\checkmarkr\checkmarkc\checkmark$\checkmarkg\checkmark$\checkmark^\checkmark\\checkmarkc\checkmarki\checkmarkr\checkmarkc\checkmark$\checkmarki\checkmark$\checkmark^\checkmark\\checkmarkc\checkmarki\checkmarkr\checkmarkc\checkmark$\checkmarkn\checkmark$\checkmark^\checkmark\\checkmarkc\checkmarki\checkmarkr\checkmarkc\checkmark$\checkmark{\checkmark$\checkmark^\checkmark\\checkmarkc\checkmarki\checkmarkr\checkmarkc\checkmark$\checkmarke\checkmark$\checkmark^\checkmark\\checkmarkc\checkmarki\checkmarkr\checkmarkc\checkmark$\checkmarkq\checkmark$\checkmark^\checkmark\\checkmarkc\checkmarki\checkmarkr\checkmarkc\checkmark$\checkmarku\checkmark$\checkmark^\checkmark\\checkmarkc\checkmarki\checkmarkr\checkmarkc\checkmark$\checkmarka\checkmark$\checkmark^\checkmark\\checkmarkc\checkmarki\checkmarkr\checkmarkc\checkmark$\checkmarkt\checkmark$\checkmark^\checkmark\\checkmarkc\checkmarki\checkmarkr\checkmarkc\checkmark$\checkmarki\checkmark$\checkmark^\checkmark\\checkmarkc\checkmarki\checkmarkr\checkmarkc\checkmark$\checkmarko\checkmark$\checkmark^\checkmark\\checkmarkc\checkmarki\checkmarkr\checkmarkc\checkmark$\checkmarkn\checkmark$\checkmark^\checkmark\\checkmarkc\checkmarki\checkmarkr\checkmarkc\checkmark$\checkmark}\checkmark$\checkmark^\checkmark\\checkmarkc\checkmarki\checkmarkr\checkmarkc\checkmark$\checkmark
\checkmark$\checkmark^\checkmark\\checkmarkc\checkmarki\checkmarkr\checkmarkc\checkmark$\checkmark \checkmark$\checkmark^\checkmark\\checkmarkc\checkmarki\checkmarkr\checkmarkc\checkmark$\checkmark \checkmark$\checkmark^\checkmark\\checkmarkc\checkmarki\checkmarkr\checkmarkc\checkmark$\checkmark \checkmark$\checkmark^\checkmark\\checkmarkc\checkmarki\checkmarkr\checkmarkc\checkmark$\checkmark \checkmark$\checkmark^\checkmark\\checkmarkc\checkmarki\checkmarkr\checkmarkc\checkmark$\checkmarkF\checkmark$\checkmark^\checkmark\\checkmarkc\checkmarki\checkmarkr\checkmarkc\checkmark$\checkmark_\checkmark$\checkmark^\checkmark\\checkmarkc\checkmarki\checkmarkr\checkmarkc\checkmark$\checkmark{\checkmark$\checkmark^\checkmark\\checkmarkc\checkmarki\checkmarkr\checkmarkc\checkmark$\checkmarkr\checkmark$\checkmark^\checkmark\\checkmarkc\checkmarki\checkmarkr\checkmarkc\checkmark$\checkmarka\checkmark$\checkmark^\checkmark\\checkmarkc\checkmarki\checkmarkr\checkmarkc\checkmark$\checkmarkd\checkmark$\checkmark^\checkmark\\checkmarkc\checkmarki\checkmarkr\checkmarkc\checkmark$\checkmarki\checkmark$\checkmark^\checkmark\\checkmarkc\checkmarki\checkmarkr\checkmarkc\checkmark$\checkmarka\checkmark$\checkmark^\checkmark\\checkmarkc\checkmarki\checkmarkr\checkmarkc\checkmark$\checkmarkl\checkmark$\checkmark^\checkmark\\checkmarkc\checkmarki\checkmarkr\checkmarkc\checkmark$\checkmarke\checkmark$\checkmark^\checkmark\\checkmarkc\checkmarki\checkmarkr\checkmarkc\checkmark$\checkmark\\checkmark$\checkmark^\checkmark\\checkmarkc\checkmarki\checkmarkr\checkmarkc\checkmark$\checkmark_\checkmark$\checkmark^\checkmark\\checkmarkc\checkmarki\checkmarkr\checkmarkc\checkmark$\checkmarka\checkmark$\checkmark^\checkmark\\checkmarkc\checkmarki\checkmarkr\checkmarkc\checkmark$\checkmarkr\checkmark$\checkmark^\checkmark\\checkmarkc\checkmarki\checkmarkr\checkmarkc\checkmark$\checkmarkb\checkmark$\checkmark^\checkmark\\checkmarkc\checkmarki\checkmarkr\checkmarkc\checkmark$\checkmarkr\checkmark$\checkmark^\checkmark\\checkmarkc\checkmarki\checkmarkr\checkmarkc\checkmark$\checkmarke\checkmark$\checkmark^\checkmark\\checkmarkc\checkmarki\checkmarkr\checkmarkc\checkmark$\checkmark}\checkmark$\checkmark^\checkmark\\checkmarkc\checkmarki\checkmarkr\checkmarkc\checkmark$\checkmark \checkmark$\checkmark^\checkmark\\checkmarkc\checkmarki\checkmarkr\checkmarkc\checkmark$\checkmark=\checkmark$\checkmark^\checkmark\\checkmarkc\checkmarki\checkmarkr\checkmarkc\checkmark$\checkmark \checkmark$\checkmark^\checkmark\\checkmarkc\checkmarki\checkmarkr\checkmarkc\checkmark$\checkmarkF\checkmark$\checkmark^\checkmark\\checkmarkc\checkmarki\checkmarkr\checkmarkc\checkmark$\checkmark_\checkmark$\checkmark^\checkmark\\checkmarkc\checkmarki\checkmarkr\checkmarkc\checkmark$\checkmark{\checkmark$\checkmark^\checkmark\\checkmarkc\checkmarki\checkmarkr\checkmarkc\checkmark$\checkmarkc\checkmark$\checkmark^\checkmark\\checkmarkc\checkmarki\checkmarkr\checkmarkc\checkmark$\checkmarko\checkmark$\checkmark^\checkmark\\checkmarkc\checkmarki\checkmarkr\checkmarkc\checkmark$\checkmarku\checkmark$\checkmark^\checkmark\\checkmarkc\checkmarki\checkmarkr\checkmarkc\checkmark$\checkmarkr\checkmark$\checkmark^\checkmark\\checkmarkc\checkmarki\checkmarkr\checkmarkc\checkmark$\checkmarkr\checkmark$\checkmark^\checkmark\\checkmarkc\checkmarki\checkmarkr\checkmarkc\checkmark$\checkmarko\checkmark$\checkmark^\checkmark\\checkmarkc\checkmarki\checkmarkr\checkmarkc\checkmark$\checkmarki\checkmark$\checkmark^\checkmark\\checkmarkc\checkmarki\checkmarkr\checkmarkc\checkmark$\checkmarke\checkmark$\checkmark^\checkmark\\checkmarkc\checkmarki\checkmarkr\checkmarkc\checkmark$\checkmark}\checkmark$\checkmark^\checkmark\\checkmarkc\checkmarki\checkmarkr\checkmarkc\checkmark$\checkmark \checkmark$\checkmark^\checkmark\\checkmarkc\checkmarki\checkmarkr\checkmarkc\checkmark$\checkmark\\checkmark$\checkmark^\checkmark\\checkmarkc\checkmarki\checkmarkr\checkmarkc\checkmark$\checkmarkt\checkmark$\checkmark^\checkmark\\checkmarkc\checkmarki\checkmarkr\checkmarkc\checkmark$\checkmarki\checkmark$\checkmark^\checkmark\\checkmarkc\checkmarki\checkmarkr\checkmarkc\checkmark$\checkmarkm\checkmark$\checkmark^\checkmark\\checkmarkc\checkmarki\checkmarkr\checkmarkc\checkmark$\checkmarke\checkmark$\checkmark^\checkmark\\checkmarkc\checkmarki\checkmarkr\checkmarkc\checkmark$\checkmarks\checkmark$\checkmark^\checkmark\\checkmarkc\checkmarki\checkmarkr\checkmarkc\checkmark$\checkmark \checkmark$\checkmark^\checkmark\\checkmarkc\checkmarki\checkmarkr\checkmarkc\checkmark$\checkmark(\checkmark$\checkmark^\checkmark\\checkmarkc\checkmarki\checkmarkr\checkmarkc\checkmark$\checkmark1\checkmark$\checkmark^\checkmark\\checkmarkc\checkmarki\checkmarkr\checkmarkc\checkmark$\checkmark \checkmark$\checkmark^\checkmark\\checkmarkc\checkmarki\checkmarkr\checkmarkc\checkmark$\checkmark+\checkmark$\checkmark^\checkmark\\checkmarkc\checkmarki\checkmarkr\checkmarkc\checkmark$\checkmark \checkmark$\checkmark^\checkmark\\checkmarkc\checkmarki\checkmarkr\checkmarkc\checkmark$\checkmark1\checkmark$\checkmark^\checkmark\\checkmarkc\checkmarki\checkmarkr\checkmarkc\checkmark$\checkmark/\checkmark$\checkmark^\checkmark\\checkmarkc\checkmarki\checkmarkr\checkmarkc\checkmark$\checkmark3\checkmark$\checkmark^\checkmark\\checkmarkc\checkmarki\checkmarkr\checkmarkc\checkmark$\checkmark)\checkmark$\checkmark^\checkmark\\checkmarkc\checkmarki\checkmarkr\checkmarkc\checkmark$\checkmark \checkmark$\checkmark^\checkmark\\checkmarkc\checkmarki\checkmarkr\checkmarkc\checkmark$\checkmark\\checkmark$\checkmark^\checkmark\\checkmarkc\checkmarki\checkmarkr\checkmarkc\checkmark$\checkmarka\checkmark$\checkmark^\checkmark\\checkmarkc\checkmarki\checkmarkr\checkmarkc\checkmark$\checkmarkp\checkmark$\checkmark^\checkmark\\checkmarkc\checkmarki\checkmarkr\checkmarkc\checkmark$\checkmarkp\checkmark$\checkmark^\checkmark\\checkmarkc\checkmarki\checkmarkr\checkmarkc\checkmark$\checkmarkr\checkmark$\checkmark^\checkmark\\checkmarkc\checkmarki\checkmarkr\checkmarkc\checkmark$\checkmarko\checkmark$\checkmark^\checkmark\\checkmarkc\checkmarki\checkmarkr\checkmarkc\checkmark$\checkmarkx\checkmark$\checkmark^\checkmark\\checkmarkc\checkmarki\checkmarkr\checkmarkc\checkmark$\checkmark \checkmark$\checkmark^\checkmark\\checkmarkc\checkmarki\checkmarkr\checkmarkc\checkmark$\checkmark\\checkmark$\checkmark^\checkmark\\checkmarkc\checkmarki\checkmarkr\checkmarkc\checkmark$\checkmarkb\checkmark$\checkmark^\checkmark\\checkmarkc\checkmarki\checkmarkr\checkmarkc\checkmark$\checkmarko\checkmark$\checkmark^\checkmark\\checkmarkc\checkmarki\checkmarkr\checkmarkc\checkmark$\checkmarkx\checkmark$\checkmark^\checkmark\\checkmarkc\checkmarki\checkmarkr\checkmarkc\checkmark$\checkmarke\checkmark$\checkmark^\checkmark\\checkmarkc\checkmarki\checkmarkr\checkmarkc\checkmark$\checkmarkd\checkmark$\checkmark^\checkmark\\checkmarkc\checkmarki\checkmarkr\checkmarkc\checkmark$\checkmark{\checkmark$\checkmark^\checkmark\\checkmarkc\checkmarki\checkmarkr\checkmarkc\checkmark$\checkmark3\checkmark$\checkmark^\checkmark\\checkmarkc\checkmarki\checkmarkr\checkmarkc\checkmark$\checkmark3\checkmark$\checkmark^\checkmark\\checkmarkc\checkmarki\checkmarkr\checkmarkc\checkmark$\checkmark3\checkmark$\checkmark^\checkmark\\checkmarkc\checkmarki\checkmarkr\checkmarkc\checkmark$\checkmark \checkmark$\checkmark^\checkmark\\checkmarkc\checkmarki\checkmarkr\checkmarkc\checkmark$\checkmark\\checkmark$\checkmark^\checkmark\\checkmarkc\checkmarki\checkmarkr\checkmarkc\checkmark$\checkmarkt\checkmark$\checkmark^\checkmark\\checkmarkc\checkmarki\checkmarkr\checkmarkc\checkmark$\checkmarke\checkmark$\checkmark^\checkmark\\checkmarkc\checkmarki\checkmarkr\checkmarkc\checkmark$\checkmarkx\checkmark$\checkmark^\checkmark\\checkmarkc\checkmarki\checkmarkr\checkmarkc\checkmark$\checkmarkt\checkmark$\checkmark^\checkmark\\checkmarkc\checkmarki\checkmarkr\checkmarkc\checkmark$\checkmark{\checkmark$\checkmark^\checkmark\\checkmarkc\checkmarki\checkmarkr\checkmarkc\checkmark$\checkmark \checkmark$\checkmark^\checkmark\\checkmarkc\checkmarki\checkmarkr\checkmarkc\checkmark$\checkmarkN\checkmark$\checkmark^\checkmark\\checkmarkc\checkmarki\checkmarkr\checkmarkc\checkmark$\checkmark}\checkmark$\checkmark^\checkmark\\checkmarkc\checkmarki\checkmarkr\checkmarkc\checkmark$\checkmark}\checkmark$\checkmark^\checkmark\\checkmarkc\checkmarki\checkmarkr\checkmarkc\checkmark$\checkmark
\checkmark$\checkmark^\checkmark\\checkmarkc\checkmarki\checkmarkr\checkmarkc\checkmark$\checkmark\\checkmark$\checkmark^\checkmark\\checkmarkc\checkmarki\checkmarkr\checkmarkc\checkmark$\checkmarke\checkmark$\checkmark^\checkmark\\checkmarkc\checkmarki\checkmarkr\checkmarkc\checkmark$\checkmarkn\checkmark$\checkmark^\checkmark\\checkmarkc\checkmarki\checkmarkr\checkmarkc\checkmark$\checkmarkd\checkmark$\checkmark^\checkmark\\checkmarkc\checkmarki\checkmarkr\checkmarkc\checkmark$\checkmark{\checkmark$\checkmark^\checkmark\\checkmarkc\checkmarki\checkmarkr\checkmarkc\checkmark$\checkmarke\checkmark$\checkmark^\checkmark\\checkmarkc\checkmarki\checkmarkr\checkmarkc\checkmark$\checkmarkq\checkmark$\checkmark^\checkmark\\checkmarkc\checkmarki\checkmarkr\checkmarkc\checkmark$\checkmarku\checkmark$\checkmark^\checkmark\\checkmarkc\checkmarki\checkmarkr\checkmarkc\checkmark$\checkmarka\checkmark$\checkmark^\checkmark\\checkmarkc\checkmarki\checkmarkr\checkmarkc\checkmark$\checkmarkt\checkmark$\checkmark^\checkmark\\checkmarkc\checkmarki\checkmarkr\checkmarkc\checkmark$\checkmarki\checkmark$\checkmark^\checkmark\\checkmarkc\checkmarki\checkmarkr\checkmarkc\checkmark$\checkmarko\checkmark$\checkmark^\checkmark\\checkmarkc\checkmarki\checkmarkr\checkmarkc\checkmark$\checkmarkn\checkmark$\checkmark^\checkmark\\checkmarkc\checkmarki\checkmarkr\checkmarkc\checkmark$\checkmark}\checkmark$\checkmark^\checkmark\\checkmarkc\checkmarki\checkmarkr\checkmarkc\checkmark$\checkmark
\checkmark$\checkmark^\checkmark\\checkmarkc\checkmarki\checkmarkr\checkmarkc\checkmark$\checkmarkC\checkmark$\checkmark^\checkmark\\checkmarkc\checkmarki\checkmarkr\checkmarkc\checkmark$\checkmarke\checkmark$\checkmark^\checkmark\\checkmarkc\checkmarki\checkmarkr\checkmarkc\checkmark$\checkmarkt\checkmark$\checkmark^\checkmark\\checkmarkc\checkmarki\checkmarkr\checkmarkc\checkmark$\checkmarkt\checkmark$\checkmark^\checkmark\\checkmarkc\checkmarki\checkmarkr\checkmarkc\checkmark$\checkmarke\checkmark$\checkmark^\checkmark\\checkmarkc\checkmarki\checkmarkr\checkmarkc\checkmark$\checkmark \checkmark$\checkmark^\checkmark\\checkmarkc\checkmarki\checkmarkr\checkmarkc\checkmark$\checkmarkv\checkmark$\checkmark^\checkmark\\checkmarkc\checkmarki\checkmarkr\checkmarkc\checkmark$\checkmarka\checkmark$\checkmark^\checkmark\\checkmarkc\checkmarki\checkmarkr\checkmarkc\checkmark$\checkmarkl\checkmark$\checkmark^\checkmark\\checkmarkc\checkmarki\checkmarkr\checkmarkc\checkmark$\checkmarke\checkmark$\checkmark^\checkmark\\checkmarkc\checkmarki\checkmarkr\checkmarkc\checkmark$\checkmarku\checkmark$\checkmark^\checkmark\\checkmarkc\checkmarki\checkmarkr\checkmarkc\checkmark$\checkmarkr\checkmark$\checkmark^\checkmark\\checkmarkc\checkmarki\checkmarkr\checkmarkc\checkmark$\checkmark \checkmark$\checkmark^\checkmark\\checkmarkc\checkmarki\checkmarkr\checkmarkc\checkmark$\checkmarks\checkmark$\checkmark^\checkmark\\checkmarkc\checkmarki\checkmarkr\checkmarkc\checkmark$\checkmarke\checkmark$\checkmark^\checkmark\\checkmarkc\checkmarki\checkmarkr\checkmarkc\checkmark$\checkmarkr\checkmark$\checkmark^\checkmark\\checkmarkc\checkmarki\checkmarkr\checkmarkc\checkmark$\checkmarka\checkmark$\checkmark^\checkmark\\checkmarkc\checkmarki\checkmarkr\checkmarkc\checkmark$\checkmark \checkmark$\checkmark^\checkmark\\checkmarkc\checkmarki\checkmarkr\checkmarkc\checkmark$\checkmarku\checkmark$\checkmark^\checkmark\\checkmarkc\checkmarki\checkmarkr\checkmarkc\checkmark$\checkmarkt\checkmark$\checkmark^\checkmark\\checkmarkc\checkmarki\checkmarkr\checkmarkc\checkmark$\checkmarki\checkmark$\checkmark^\checkmark\\checkmarkc\checkmarki\checkmarkr\checkmarkc\checkmark$\checkmarkl\checkmark$\checkmark^\checkmark\\checkmarkc\checkmarki\checkmarkr\checkmarkc\checkmark$\checkmarki\checkmark$\checkmark^\checkmark\\checkmarkc\checkmarki\checkmarkr\checkmarkc\checkmark$\checkmarks\checkmark$\checkmark^\checkmark\\checkmarkc\checkmarki\checkmarkr\checkmarkc\checkmark$\checkmarké\checkmark$\checkmark^\checkmark\\checkmarkc\checkmarki\checkmarkr\checkmarkc\checkmark$\checkmarke\checkmark$\checkmark^\checkmark\\checkmarkc\checkmarki\checkmarkr\checkmarkc\checkmark$\checkmark \checkmark$\checkmark^\checkmark\\checkmarkc\checkmarki\checkmarkr\checkmarkc\checkmark$\checkmarkc\checkmark$\checkmark^\checkmark\\checkmarkc\checkmarki\checkmarkr\checkmarkc\checkmark$\checkmarko\checkmark$\checkmark^\checkmark\\checkmarkc\checkmarki\checkmarkr\checkmarkc\checkmark$\checkmarkm\checkmark$\checkmark^\checkmark\\checkmarkc\checkmarki\checkmarkr\checkmarkc\checkmark$\checkmarkm\checkmark$\checkmark^\checkmark\\checkmarkc\checkmarki\checkmarkr\checkmarkc\checkmark$\checkmarke\checkmark$\checkmark^\checkmark\\checkmarkc\checkmarki\checkmarkr\checkmarkc\checkmark$\checkmark \checkmark$\checkmark^\checkmark\\checkmarkc\checkmarki\checkmarkr\checkmarkc\checkmark$\checkmarkc\checkmark$\checkmark^\checkmark\\checkmarkc\checkmarki\checkmarkr\checkmarkc\checkmark$\checkmarkh\checkmark$\checkmark^\checkmark\\checkmarkc\checkmarki\checkmarkr\checkmarkc\checkmark$\checkmarka\checkmark$\checkmark^\checkmark\\checkmarkc\checkmarki\checkmarkr\checkmarkc\checkmark$\checkmarkr\checkmark$\checkmark^\checkmark\\checkmarkc\checkmarki\checkmarkr\checkmarkc\checkmark$\checkmarkg\checkmark$\checkmark^\checkmark\\checkmarkc\checkmarki\checkmarkr\checkmarkc\checkmark$\checkmarke\checkmark$\checkmark^\checkmark\\checkmarkc\checkmarki\checkmarkr\checkmarkc\checkmark$\checkmark \checkmark$\checkmark^\checkmark\\checkmarkc\checkmarki\checkmarkr\checkmarkc\checkmark$\checkmarkr\checkmark$\checkmark^\checkmark\\checkmarkc\checkmarki\checkmarkr\checkmarkc\checkmark$\checkmarka\checkmark$\checkmark^\checkmark\\checkmarkc\checkmarki\checkmarkr\checkmarkc\checkmark$\checkmarkd\checkmark$\checkmark^\checkmark\\checkmarkc\checkmarki\checkmarkr\checkmarkc\checkmark$\checkmarki\checkmark$\checkmark^\checkmark\\checkmarkc\checkmarki\checkmarkr\checkmarkc\checkmark$\checkmarka\checkmark$\checkmark^\checkmark\\checkmarkc\checkmarki\checkmarkr\checkmarkc\checkmark$\checkmarkl\checkmark$\checkmark^\checkmark\\checkmarkc\checkmarki\checkmarkr\checkmarkc\checkmark$\checkmarke\checkmark$\checkmark^\checkmark\\checkmarkc\checkmarki\checkmarkr\checkmarkc\checkmark$\checkmark \checkmark$\checkmark^\checkmark\\checkmarkc\checkmarki\checkmarkr\checkmarkc\checkmark$\checkmarkd\checkmark$\checkmark^\checkmark\\checkmarkc\checkmarki\checkmarkr\checkmarkc\checkmark$\checkmarka\checkmark$\checkmark^\checkmark\\checkmarkc\checkmarki\checkmarkr\checkmarkc\checkmark$\checkmarkn\checkmark$\checkmark^\checkmark\\checkmarkc\checkmarki\checkmarkr\checkmarkc\checkmark$\checkmarks\checkmark$\checkmark^\checkmark\\checkmarkc\checkmarki\checkmarkr\checkmarkc\checkmark$\checkmark \checkmark$\checkmark^\checkmark\\checkmarkc\checkmarki\checkmarkr\checkmarkc\checkmark$\checkmarkl\checkmark$\checkmark^\checkmark\\checkmarkc\checkmarki\checkmarkr\checkmarkc\checkmark$\checkmarke\checkmark$\checkmark^\checkmark\\checkmarkc\checkmarki\checkmarkr\checkmarkc\checkmark$\checkmark \checkmark$\checkmark^\checkmark\\checkmarkc\checkmarki\checkmarkr\checkmarkc\checkmark$\checkmarkd\checkmark$\checkmark^\checkmark\\checkmarkc\checkmarki\checkmarkr\checkmarkc\checkmark$\checkmarki\checkmark$\checkmark^\checkmark\\checkmarkc\checkmarki\checkmarkr\checkmarkc\checkmark$\checkmarkm\checkmark$\checkmark^\checkmark\\checkmarkc\checkmarki\checkmarkr\checkmarkc\checkmark$\checkmarke\checkmark$\checkmark^\checkmark\\checkmarkc\checkmarki\checkmarkr\checkmarkc\checkmark$\checkmarkn\checkmark$\checkmark^\checkmark\\checkmarkc\checkmarki\checkmarkr\checkmarkc\checkmark$\checkmarks\checkmark$\checkmark^\checkmark\\checkmarkc\checkmarki\checkmarkr\checkmarkc\checkmark$\checkmarki\checkmark$\checkmark^\checkmark\\checkmarkc\checkmarki\checkmarkr\checkmarkc\checkmark$\checkmarko\checkmark$\checkmark^\checkmark\\checkmarkc\checkmarki\checkmarkr\checkmarkc\checkmark$\checkmarkn\checkmark$\checkmark^\checkmark\\checkmarkc\checkmarki\checkmarkr\checkmarkc\checkmark$\checkmarkn\checkmark$\checkmark^\checkmark\\checkmarkc\checkmarki\checkmarkr\checkmarkc\checkmark$\checkmarke\checkmark$\checkmark^\checkmark\\checkmarkc\checkmarki\checkmarkr\checkmarkc\checkmark$\checkmarkm\checkmark$\checkmark^\checkmark\\checkmarkc\checkmarki\checkmarkr\checkmarkc\checkmark$\checkmarke\checkmark$\checkmark^\checkmark\\checkmarkc\checkmarki\checkmarkr\checkmarkc\checkmark$\checkmarkn\checkmark$\checkmark^\checkmark\\checkmarkc\checkmarki\checkmarkr\checkmarkc\checkmark$\checkmarkt\checkmark$\checkmark^\checkmark\\checkmarkc\checkmarki\checkmarkr\checkmarkc\checkmark$\checkmark \checkmark$\checkmark^\checkmark\\checkmarkc\checkmarki\checkmarkr\checkmarkc\checkmark$\checkmarkd\checkmark$\checkmark^\checkmark\\checkmarkc\checkmarki\checkmarkr\checkmarkc\checkmark$\checkmarke\checkmark$\checkmark^\checkmark\\checkmarkc\checkmarki\checkmarkr\checkmarkc\checkmark$\checkmark \checkmark$\checkmark^\checkmark\\checkmarkc\checkmarki\checkmarkr\checkmarkc\checkmark$\checkmarkl\checkmark$\checkmark^\checkmark\\checkmarkc\checkmarki\checkmarkr\checkmarkc\checkmark$\checkmark'\checkmark$\checkmark^\checkmark\\checkmarkc\checkmarki\checkmarkr\checkmarkc\checkmark$\checkmarka\checkmark$\checkmark^\checkmark\\checkmarkc\checkmarki\checkmarkr\checkmarkc\checkmark$\checkmarkr\checkmark$\checkmark^\checkmark\\checkmarkc\checkmarki\checkmarkr\checkmarkc\checkmark$\checkmarkb\checkmark$\checkmark^\checkmark\\checkmarkc\checkmarki\checkmarkr\checkmarkc\checkmark$\checkmarkr\checkmark$\checkmark^\checkmark\\checkmarkc\checkmarki\checkmarkr\checkmarkc\checkmark$\checkmarke\checkmark$\checkmark^\checkmark\\checkmarkc\checkmarki\checkmarkr\checkmarkc\checkmark$\checkmark \checkmark$\checkmark^\checkmark\\checkmarkc\checkmarki\checkmarkr\checkmarkc\checkmark$\checkmarkd\checkmark$\checkmark^\checkmark\\checkmarkc\checkmarki\checkmarkr\checkmarkc\checkmark$\checkmarke\checkmark$\checkmark^\checkmark\\checkmarkc\checkmarki\checkmarkr\checkmarkc\checkmark$\checkmark \checkmark$\checkmark^\checkmark\\checkmarkc\checkmarki\checkmarkr\checkmarkc\checkmark$\checkmarkp\checkmark$\checkmark^\checkmark\\checkmarkc\checkmarki\checkmarkr\checkmarkc\checkmark$\checkmarki\checkmark$\checkmark^\checkmark\\checkmarkc\checkmarki\checkmarkr\checkmarkc\checkmark$\checkmarkv\checkmark$\checkmark^\checkmark\\checkmarkc\checkmarki\checkmarkr\checkmarkc\checkmark$\checkmarko\checkmark$\checkmark^\checkmark\\checkmarkc\checkmarki\checkmarkr\checkmarkc\checkmark$\checkmarkt\checkmark$\checkmark^\checkmark\\checkmarkc\checkmarki\checkmarkr\checkmarkc\checkmark$\checkmarke\checkmark$\checkmark^\checkmark\\checkmarkc\checkmarki\checkmarkr\checkmarkc\checkmark$\checkmarkm\checkmark$\checkmark^\checkmark\\checkmarkc\checkmarki\checkmarkr\checkmarkc\checkmark$\checkmarke\checkmark$\checkmark^\checkmark\\checkmarkc\checkmarki\checkmarkr\checkmarkc\checkmark$\checkmarkn\checkmark$\checkmark^\checkmark\\checkmarkc\checkmarki\checkmarkr\checkmarkc\checkmark$\checkmarkt\checkmark$\checkmark^\checkmark\\checkmarkc\checkmarki\checkmarkr\checkmarkc\checkmark$\checkmark.\checkmark$\checkmark^\checkmark\\checkmarkc\checkmarki\checkmarkr\checkmarkc\checkmark$\checkmark
\checkmark$\checkmark^\checkmark\\checkmarkc\checkmarki\checkmarkr\checkmarkc\checkmark$\checkmark
\checkmark$\checkmark^\checkmark\\checkmarkc\checkmarki\checkmarkr\checkmarkc\checkmark$\checkmark%\checkmark$\checkmark^\checkmark\\checkmarkc\checkmarki\checkmarkr\checkmarkc\checkmark$\checkmark \checkmark$\checkmark^\checkmark\\checkmarkc\checkmarki\checkmarkr\checkmarkc\checkmark$\checkmark=\checkmark$\checkmark^\checkmark\\checkmarkc\checkmarki\checkmarkr\checkmarkc\checkmark$\checkmark=\checkmark$\checkmark^\checkmark\\checkmarkc\checkmarki\checkmarkr\checkmarkc\checkmark$\checkmark=\checkmark$\checkmark^\checkmark\\checkmarkc\checkmarki\checkmarkr\checkmarkc\checkmark$\checkmark=\checkmark$\checkmark^\checkmark\\checkmarkc\checkmarki\checkmarkr\checkmarkc\checkmark$\checkmark=\checkmark$\checkmark^\checkmark\\checkmarkc\checkmarki\checkmarkr\checkmarkc\checkmark$\checkmark=\checkmark$\checkmark^\checkmark\\checkmarkc\checkmarki\checkmarkr\checkmarkc\checkmark$\checkmark=\checkmark$\checkmark^\checkmark\\checkmarkc\checkmarki\checkmarkr\checkmarkc\checkmark$\checkmark=\checkmark$\checkmark^\checkmark\\checkmarkc\checkmarki\checkmarkr\checkmarkc\checkmark$\checkmark=\checkmark$\checkmark^\checkmark\\checkmarkc\checkmarki\checkmarkr\checkmarkc\checkmark$\checkmark=\checkmark$\checkmark^\checkmark\\checkmarkc\checkmarki\checkmarkr\checkmarkc\checkmark$\checkmark=\checkmark$\checkmark^\checkmark\\checkmarkc\checkmarki\checkmarkr\checkmarkc\checkmark$\checkmark=\checkmark$\checkmark^\checkmark\\checkmarkc\checkmarki\checkmarkr\checkmarkc\checkmark$\checkmark=\checkmark$\checkmark^\checkmark\\checkmarkc\checkmarki\checkmarkr\checkmarkc\checkmark$\checkmark=\checkmark$\checkmark^\checkmark\\checkmarkc\checkmarki\checkmarkr\checkmarkc\checkmark$\checkmark=\checkmark$\checkmark^\checkmark\\checkmarkc\checkmarki\checkmarkr\checkmarkc\checkmark$\checkmark=\checkmark$\checkmark^\checkmark\\checkmarkc\checkmarki\checkmarkr\checkmarkc\checkmark$\checkmark=\checkmark$\checkmark^\checkmark\\checkmarkc\checkmarki\checkmarkr\checkmarkc\checkmark$\checkmark=\checkmark$\checkmark^\checkmark\\checkmarkc\checkmarki\checkmarkr\checkmarkc\checkmark$\checkmark=\checkmark$\checkmark^\checkmark\\checkmarkc\checkmarki\checkmarkr\checkmarkc\checkmark$\checkmark=\checkmark$\checkmark^\checkmark\\checkmarkc\checkmarki\checkmarkr\checkmarkc\checkmark$\checkmark=\checkmark$\checkmark^\checkmark\\checkmarkc\checkmarki\checkmarkr\checkmarkc\checkmark$\checkmark=\checkmark$\checkmark^\checkmark\\checkmarkc\checkmarki\checkmarkr\checkmarkc\checkmark$\checkmark=\checkmark$\checkmark^\checkmark\\checkmarkc\checkmarki\checkmarkr\checkmarkc\checkmark$\checkmark=\checkmark$\checkmark^\checkmark\\checkmarkc\checkmarki\checkmarkr\checkmarkc\checkmark$\checkmark=\checkmark$\checkmark^\checkmark\\checkmarkc\checkmarki\checkmarkr\checkmarkc\checkmark$\checkmark=\checkmark$\checkmark^\checkmark\\checkmarkc\checkmarki\checkmarkr\checkmarkc\checkmark$\checkmark=\checkmark$\checkmark^\checkmark\\checkmarkc\checkmarki\checkmarkr\checkmarkc\checkmark$\checkmark=\checkmark$\checkmark^\checkmark\\checkmarkc\checkmarki\checkmarkr\checkmarkc\checkmark$\checkmark=\checkmark$\checkmark^\checkmark\\checkmarkc\checkmarki\checkmarkr\checkmarkc\checkmark$\checkmark=\checkmark$\checkmark^\checkmark\\checkmarkc\checkmarki\checkmarkr\checkmarkc\checkmark$\checkmark=\checkmark$\checkmark^\checkmark\\checkmarkc\checkmarki\checkmarkr\checkmarkc\checkmark$\checkmark=\checkmark$\checkmark^\checkmark\\checkmarkc\checkmarki\checkmarkr\checkmarkc\checkmark$\checkmark=\checkmark$\checkmark^\checkmark\\checkmarkc\checkmarki\checkmarkr\checkmarkc\checkmark$\checkmark=\checkmark$\checkmark^\checkmark\\checkmarkc\checkmarki\checkmarkr\checkmarkc\checkmark$\checkmark=\checkmark$\checkmark^\checkmark\\checkmarkc\checkmarki\checkmarkr\checkmarkc\checkmark$\checkmark=\checkmark$\checkmark^\checkmark\\checkmarkc\checkmarki\checkmarkr\checkmarkc\checkmark$\checkmark=\checkmark$\checkmark^\checkmark\\checkmarkc\checkmarki\checkmarkr\checkmarkc\checkmark$\checkmark=\checkmark$\checkmark^\checkmark\\checkmarkc\checkmarki\checkmarkr\checkmarkc\checkmark$\checkmark=\checkmark$\checkmark^\checkmark\\checkmarkc\checkmarki\checkmarkr\checkmarkc\checkmark$\checkmark=\checkmark$\checkmark^\checkmark\\checkmarkc\checkmarki\checkmarkr\checkmarkc\checkmark$\checkmark=\checkmark$\checkmark^\checkmark\\checkmarkc\checkmarki\checkmarkr\checkmarkc\checkmark$\checkmark=\checkmark$\checkmark^\checkmark\\checkmarkc\checkmarki\checkmarkr\checkmarkc\checkmark$\checkmark=\checkmark$\checkmark^\checkmark\\checkmarkc\checkmarki\checkmarkr\checkmarkc\checkmark$\checkmark=\checkmark$\checkmark^\checkmark\\checkmarkc\checkmarki\checkmarkr\checkmarkc\checkmark$\checkmark=\checkmark$\checkmark^\checkmark\\checkmarkc\checkmarki\checkmarkr\checkmarkc\checkmark$\checkmark=\checkmark$\checkmark^\checkmark\\checkmarkc\checkmarki\checkmarkr\checkmarkc\checkmark$\checkmark=\checkmark$\checkmark^\checkmark\\checkmarkc\checkmarki\checkmarkr\checkmarkc\checkmark$\checkmark=\checkmark$\checkmark^\checkmark\\checkmarkc\checkmarki\checkmarkr\checkmarkc\checkmark$\checkmark=\checkmark$\checkmark^\checkmark\\checkmarkc\checkmarki\checkmarkr\checkmarkc\checkmark$\checkmark=\checkmark$\checkmark^\checkmark\\checkmarkc\checkmarki\checkmarkr\checkmarkc\checkmark$\checkmark=\checkmark$\checkmark^\checkmark\\checkmarkc\checkmarki\checkmarkr\checkmarkc\checkmark$\checkmark=\checkmark$\checkmark^\checkmark\\checkmarkc\checkmarki\checkmarkr\checkmarkc\checkmark$\checkmark=\checkmark$\checkmark^\checkmark\\checkmarkc\checkmarki\checkmarkr\checkmarkc\checkmark$\checkmark=\checkmark$\checkmark^\checkmark\\checkmarkc\checkmarki\checkmarkr\checkmarkc\checkmark$\checkmark=\checkmark$\checkmark^\checkmark\\checkmarkc\checkmarki\checkmarkr\checkmarkc\checkmark$\checkmark=\checkmark$\checkmark^\checkmark\\checkmarkc\checkmarki\checkmarkr\checkmarkc\checkmark$\checkmark=\checkmark$\checkmark^\checkmark\\checkmarkc\checkmarki\checkmarkr\checkmarkc\checkmark$\checkmark=\checkmark$\checkmark^\checkmark\\checkmarkc\checkmarki\checkmarkr\checkmarkc\checkmark$\checkmark=\checkmark$\checkmark^\checkmark\\checkmarkc\checkmarki\checkmarkr\checkmarkc\checkmark$\checkmark=\checkmark$\checkmark^\checkmark\\checkmarkc\checkmarki\checkmarkr\checkmarkc\checkmark$\checkmark=\checkmark$\checkmark^\checkmark\\checkmarkc\checkmarki\checkmarkr\checkmarkc\checkmark$\checkmark=\checkmark$\checkmark^\checkmark\\checkmarkc\checkmarki\checkmarkr\checkmarkc\checkmark$\checkmark
\checkmark$\checkmark^\checkmark\\checkmarkc\checkmarki\checkmarkr\checkmarkc\checkmark$\checkmark\\checkmark$\checkmark^\checkmark\\checkmarkc\checkmarki\checkmarkr\checkmarkc\checkmark$\checkmarks\checkmark$\checkmark^\checkmark\\checkmarkc\checkmarki\checkmarkr\checkmarkc\checkmark$\checkmarke\checkmark$\checkmark^\checkmark\\checkmarkc\checkmarki\checkmarkr\checkmarkc\checkmark$\checkmarkc\checkmark$\checkmark^\checkmark\\checkmarkc\checkmarki\checkmarkr\checkmarkc\checkmark$\checkmarkt\checkmark$\checkmark^\checkmark\\checkmarkc\checkmarki\checkmarkr\checkmarkc\checkmark$\checkmarki\checkmark$\checkmark^\checkmark\\checkmarkc\checkmarki\checkmarkr\checkmarkc\checkmark$\checkmarko\checkmark$\checkmark^\checkmark\\checkmarkc\checkmarki\checkmarkr\checkmarkc\checkmark$\checkmarkn\checkmark$\checkmark^\checkmark\\checkmarkc\checkmarki\checkmarkr\checkmarkc\checkmark$\checkmark{\checkmark$\checkmark^\checkmark\\checkmarkc\checkmarki\checkmarkr\checkmarkc\checkmark$\checkmarkD\checkmark$\checkmark^\checkmark\\checkmarkc\checkmarki\checkmarkr\checkmarkc\checkmark$\checkmarki\checkmark$\checkmark^\checkmark\\checkmarkc\checkmarki\checkmarkr\checkmarkc\checkmark$\checkmarkm\checkmark$\checkmark^\checkmark\\checkmarkc\checkmarki\checkmarkr\checkmarkc\checkmark$\checkmarke\checkmark$\checkmark^\checkmark\\checkmarkc\checkmarki\checkmarkr\checkmarkc\checkmark$\checkmarkn\checkmark$\checkmark^\checkmark\\checkmarkc\checkmarki\checkmarkr\checkmarkc\checkmark$\checkmarks\checkmark$\checkmark^\checkmark\\checkmarkc\checkmarki\checkmarkr\checkmarkc\checkmark$\checkmarki\checkmark$\checkmark^\checkmark\\checkmarkc\checkmarki\checkmarkr\checkmarkc\checkmark$\checkmarko\checkmark$\checkmark^\checkmark\\checkmarkc\checkmarki\checkmarkr\checkmarkc\checkmark$\checkmarkn\checkmark$\checkmark^\checkmark\\checkmarkc\checkmarki\checkmarkr\checkmarkc\checkmark$\checkmarkn\checkmark$\checkmark^\checkmark\\checkmarkc\checkmarki\checkmarkr\checkmarkc\checkmark$\checkmarke\checkmark$\checkmark^\checkmark\\checkmarkc\checkmarki\checkmarkr\checkmarkc\checkmark$\checkmarkm\checkmark$\checkmark^\checkmark\\checkmarkc\checkmarki\checkmarkr\checkmarkc\checkmark$\checkmarke\checkmark$\checkmark^\checkmark\\checkmarkc\checkmarki\checkmarkr\checkmarkc\checkmark$\checkmarkn\checkmark$\checkmark^\checkmark\\checkmarkc\checkmarki\checkmarkr\checkmarkc\checkmark$\checkmarkt\checkmark$\checkmark^\checkmark\\checkmarkc\checkmarki\checkmarkr\checkmarkc\checkmark$\checkmark \checkmark$\checkmark^\checkmark\\checkmarkc\checkmarki\checkmarkr\checkmarkc\checkmark$\checkmarkd\checkmark$\checkmark^\checkmark\\checkmarkc\checkmarki\checkmarkr\checkmarkc\checkmark$\checkmarke\checkmark$\checkmark^\checkmark\\checkmarkc\checkmarki\checkmarkr\checkmarkc\checkmark$\checkmark \checkmark$\checkmark^\checkmark\\checkmarkc\checkmarki\checkmarkr\checkmarkc\checkmark$\checkmarkl\checkmark$\checkmark^\checkmark\\checkmarkc\checkmarki\checkmarkr\checkmarkc\checkmark$\checkmark'\checkmark$\checkmark^\checkmark\\checkmarkc\checkmarki\checkmarkr\checkmarkc\checkmark$\checkmarkA\checkmark$\checkmark^\checkmark\\checkmarkc\checkmarki\checkmarkr\checkmarkc\checkmark$\checkmarkr\checkmark$\checkmark^\checkmark\\checkmarkc\checkmarki\checkmarkr\checkmarkc\checkmark$\checkmarkb\checkmark$\checkmark^\checkmark\\checkmarkc\checkmarki\checkmarkr\checkmarkc\checkmark$\checkmarkr\checkmark$\checkmark^\checkmark\\checkmarkc\checkmarki\checkmarkr\checkmarkc\checkmark$\checkmarke\checkmark$\checkmark^\checkmark\\checkmarkc\checkmarki\checkmarkr\checkmarkc\checkmark$\checkmark \checkmark$\checkmark^\checkmark\\checkmarkc\checkmarki\checkmarkr\checkmarkc\checkmark$\checkmarkd\checkmark$\checkmark^\checkmark\\checkmarkc\checkmarki\checkmarkr\checkmarkc\checkmark$\checkmarke\checkmark$\checkmark^\checkmark\\checkmarkc\checkmarki\checkmarkr\checkmarkc\checkmark$\checkmark \checkmark$\checkmark^\checkmark\\checkmarkc\checkmarki\checkmarkr\checkmarkc\checkmark$\checkmarkP\checkmark$\checkmark^\checkmark\\checkmarkc\checkmarki\checkmarkr\checkmarkc\checkmark$\checkmarki\checkmark$\checkmark^\checkmark\\checkmarkc\checkmarki\checkmarkr\checkmarkc\checkmark$\checkmarkv\checkmark$\checkmark^\checkmark\\checkmarkc\checkmarki\checkmarkr\checkmarkc\checkmark$\checkmarko\checkmark$\checkmark^\checkmark\\checkmarkc\checkmarki\checkmarkr\checkmarkc\checkmark$\checkmarkt\checkmark$\checkmark^\checkmark\\checkmarkc\checkmarki\checkmarkr\checkmarkc\checkmark$\checkmarke\checkmark$\checkmark^\checkmark\\checkmarkc\checkmarki\checkmarkr\checkmarkc\checkmark$\checkmarkm\checkmark$\checkmark^\checkmark\\checkmarkc\checkmarki\checkmarkr\checkmarkc\checkmark$\checkmarke\checkmark$\checkmark^\checkmark\\checkmarkc\checkmarki\checkmarkr\checkmarkc\checkmark$\checkmarkn\checkmark$\checkmark^\checkmark\\checkmarkc\checkmarki\checkmarkr\checkmarkc\checkmark$\checkmarkt\checkmark$\checkmark^\checkmark\\checkmarkc\checkmarki\checkmarkr\checkmarkc\checkmark$\checkmark \checkmark$\checkmark^\checkmark\\checkmarkc\checkmarki\checkmarkr\checkmarkc\checkmark$\checkmark(\checkmark$\checkmark^\checkmark\\checkmarkc\checkmarki\checkmarkr\checkmarkc\checkmark$\checkmarkA\checkmark$\checkmark^\checkmark\\checkmarkc\checkmarki\checkmarkr\checkmarkc\checkmark$\checkmarkx\checkmark$\checkmark^\checkmark\\checkmarkc\checkmarki\checkmarkr\checkmarkc\checkmark$\checkmarke\checkmark$\checkmark^\checkmark\\checkmarkc\checkmarki\checkmarkr\checkmarkc\checkmark$\checkmark \checkmark$\checkmark^\checkmark\\checkmarkc\checkmarki\checkmarkr\checkmarkc\checkmark$\checkmarkA\checkmark$\checkmark^\checkmark\\checkmarkc\checkmarki\checkmarkr\checkmarkc\checkmark$\checkmark)\checkmark$\checkmark^\checkmark\\checkmarkc\checkmarki\checkmarkr\checkmarkc\checkmark$\checkmark}\checkmark$\checkmark^\checkmark\\checkmarkc\checkmarki\checkmarkr\checkmarkc\checkmark$\checkmark
\checkmark$\checkmark^\checkmark\\checkmarkc\checkmarki\checkmarkr\checkmarkc\checkmark$\checkmark
\checkmark$\checkmark^\checkmark\\checkmarkc\checkmarki\checkmarkr\checkmarkc\checkmark$\checkmark\\checkmark$\checkmark^\checkmark\\checkmarkc\checkmarki\checkmarkr\checkmarkc\checkmark$\checkmarks\checkmark$\checkmark^\checkmark\\checkmarkc\checkmarki\checkmarkr\checkmarkc\checkmark$\checkmarku\checkmark$\checkmark^\checkmark\\checkmarkc\checkmarki\checkmarkr\checkmarkc\checkmark$\checkmarkb\checkmark$\checkmark^\checkmark\\checkmarkc\checkmarki\checkmarkr\checkmarkc\checkmark$\checkmarks\checkmark$\checkmark^\checkmark\\checkmarkc\checkmarki\checkmarkr\checkmarkc\checkmark$\checkmarke\checkmark$\checkmark^\checkmark\\checkmarkc\checkmarki\checkmarkr\checkmarkc\checkmark$\checkmarkc\checkmark$\checkmark^\checkmark\\checkmarkc\checkmarki\checkmarkr\checkmarkc\checkmark$\checkmarkt\checkmark$\checkmark^\checkmark\\checkmarkc\checkmarki\checkmarkr\checkmarkc\checkmark$\checkmarki\checkmark$\checkmark^\checkmark\\checkmarkc\checkmarki\checkmarkr\checkmarkc\checkmark$\checkmarko\checkmark$\checkmark^\checkmark\\checkmarkc\checkmarki\checkmarkr\checkmarkc\checkmark$\checkmarkn\checkmark$\checkmark^\checkmark\\checkmarkc\checkmarki\checkmarkr\checkmarkc\checkmark$\checkmark{\checkmark$\checkmark^\checkmark\\checkmarkc\checkmarki\checkmarkr\checkmarkc\checkmark$\checkmarkI\checkmark$\checkmark^\checkmark\\checkmarkc\checkmarki\checkmarkr\checkmarkc\checkmark$\checkmarkn\checkmark$\checkmark^\checkmark\\checkmarkc\checkmarki\checkmarkr\checkmarkc\checkmark$\checkmarkt\checkmark$\checkmark^\checkmark\\checkmarkc\checkmarki\checkmarkr\checkmarkc\checkmark$\checkmarkr\checkmark$\checkmark^\checkmark\\checkmarkc\checkmarki\checkmarkr\checkmarkc\checkmark$\checkmarko\checkmark$\checkmark^\checkmark\\checkmarkc\checkmarki\checkmarkr\checkmarkc\checkmark$\checkmarkd\checkmark$\checkmark^\checkmark\\checkmarkc\checkmarki\checkmarkr\checkmarkc\checkmark$\checkmarku\checkmark$\checkmark^\checkmark\\checkmarkc\checkmarki\checkmarkr\checkmarkc\checkmark$\checkmarkc\checkmark$\checkmark^\checkmark\\checkmarkc\checkmarki\checkmarkr\checkmarkc\checkmark$\checkmarkt\checkmark$\checkmark^\checkmark\\checkmarkc\checkmarki\checkmarkr\checkmarkc\checkmark$\checkmarki\checkmark$\checkmark^\checkmark\\checkmarkc\checkmarki\checkmarkr\checkmarkc\checkmark$\checkmarko\checkmark$\checkmark^\checkmark\\checkmarkc\checkmarki\checkmarkr\checkmarkc\checkmark$\checkmarkn\checkmark$\checkmark^\checkmark\\checkmarkc\checkmarki\checkmarkr\checkmarkc\checkmark$\checkmark}\checkmark$\checkmark^\checkmark\\checkmarkc\checkmarki\checkmarkr\checkmarkc\checkmark$\checkmark
\checkmark$\checkmark^\checkmark\\checkmarkc\checkmarki\checkmarkr\checkmarkc\checkmark$\checkmarkL\checkmark$\checkmark^\checkmark\\checkmarkc\checkmarki\checkmarkr\checkmarkc\checkmark$\checkmark'\checkmark$\checkmark^\checkmark\\checkmarkc\checkmarki\checkmarkr\checkmarkc\checkmark$\checkmarka\checkmark$\checkmark^\checkmark\\checkmarkc\checkmarki\checkmarkr\checkmarkc\checkmark$\checkmarkr\checkmark$\checkmark^\checkmark\\checkmarkc\checkmarki\checkmarkr\checkmarkc\checkmark$\checkmarkb\checkmark$\checkmark^\checkmark\\checkmarkc\checkmarki\checkmarkr\checkmarkc\checkmark$\checkmarkr\checkmark$\checkmark^\checkmark\\checkmarkc\checkmarki\checkmarkr\checkmarkc\checkmark$\checkmarke\checkmark$\checkmark^\checkmark\\checkmarkc\checkmarki\checkmarkr\checkmarkc\checkmark$\checkmark \checkmark$\checkmark^\checkmark\\checkmarkc\checkmarki\checkmarkr\checkmarkc\checkmark$\checkmarkd\checkmark$\checkmark^\checkmark\\checkmarkc\checkmarki\checkmarkr\checkmarkc\checkmark$\checkmarke\checkmark$\checkmark^\checkmark\\checkmarkc\checkmarki\checkmarkr\checkmarkc\checkmark$\checkmark \checkmark$\checkmark^\checkmark\\checkmarkc\checkmarki\checkmarkr\checkmarkc\checkmark$\checkmarkp\checkmark$\checkmark^\checkmark\\checkmarkc\checkmarki\checkmarkr\checkmarkc\checkmark$\checkmarki\checkmark$\checkmark^\checkmark\\checkmarkc\checkmarki\checkmarkr\checkmarkc\checkmark$\checkmarkv\checkmark$\checkmark^\checkmark\\checkmarkc\checkmarki\checkmarkr\checkmarkc\checkmark$\checkmarko\checkmark$\checkmark^\checkmark\\checkmarkc\checkmarki\checkmarkr\checkmarkc\checkmark$\checkmarkt\checkmark$\checkmark^\checkmark\\checkmarkc\checkmarki\checkmarkr\checkmarkc\checkmark$\checkmarke\checkmark$\checkmark^\checkmark\\checkmarkc\checkmarki\checkmarkr\checkmarkc\checkmark$\checkmarkm\checkmark$\checkmark^\checkmark\\checkmarkc\checkmarki\checkmarkr\checkmarkc\checkmark$\checkmarke\checkmark$\checkmark^\checkmark\\checkmarkc\checkmarki\checkmarkr\checkmarkc\checkmark$\checkmarkn\checkmark$\checkmark^\checkmark\\checkmarkc\checkmarki\checkmarkr\checkmarkc\checkmark$\checkmarkt\checkmark$\checkmark^\checkmark\\checkmarkc\checkmarki\checkmarkr\checkmarkc\checkmark$\checkmark \checkmark$\checkmark^\checkmark\\checkmarkc\checkmarki\checkmarkr\checkmarkc\checkmark$\checkmarke\checkmark$\checkmark^\checkmark\\checkmarkc\checkmarki\checkmarkr\checkmarkc\checkmark$\checkmarks\checkmark$\checkmark^\checkmark\\checkmarkc\checkmarki\checkmarkr\checkmarkc\checkmark$\checkmarkt\checkmark$\checkmark^\checkmark\\checkmarkc\checkmarki\checkmarkr\checkmarkc\checkmark$\checkmark \checkmark$\checkmark^\checkmark\\checkmarkc\checkmarki\checkmarkr\checkmarkc\checkmark$\checkmarkl\checkmark$\checkmark^\checkmark\\checkmarkc\checkmarki\checkmarkr\checkmarkc\checkmark$\checkmark'\checkmark$\checkmark^\checkmark\\checkmarkc\checkmarki\checkmarkr\checkmarkc\checkmark$\checkmarko\checkmark$\checkmark^\checkmark\\checkmarkc\checkmarki\checkmarkr\checkmarkc\checkmark$\checkmarkr\checkmark$\checkmark^\checkmark\\checkmarkc\checkmarki\checkmarkr\checkmarkc\checkmark$\checkmarkg\checkmark$\checkmark^\checkmark\\checkmarkc\checkmarki\checkmarkr\checkmarkc\checkmark$\checkmarka\checkmark$\checkmark^\checkmark\\checkmarkc\checkmarki\checkmarkr\checkmarkc\checkmark$\checkmarkn\checkmark$\checkmark^\checkmark\\checkmarkc\checkmarki\checkmarkr\checkmarkc\checkmark$\checkmarke\checkmark$\checkmark^\checkmark\\checkmarkc\checkmarki\checkmarkr\checkmarkc\checkmark$\checkmark \checkmark$\checkmark^\checkmark\\checkmarkc\checkmarki\checkmarkr\checkmarkc\checkmark$\checkmarkl\checkmark$\checkmark^\checkmark\\checkmarkc\checkmarki\checkmarkr\checkmarkc\checkmark$\checkmarke\checkmark$\checkmark^\checkmark\\checkmarkc\checkmarki\checkmarkr\checkmarkc\checkmark$\checkmark \checkmark$\checkmark^\checkmark\\checkmarkc\checkmarki\checkmarkr\checkmarkc\checkmark$\checkmarkp\checkmark$\checkmark^\checkmark\\checkmarkc\checkmarki\checkmarkr\checkmarkc\checkmark$\checkmarkl\checkmark$\checkmark^\checkmark\\checkmarkc\checkmarki\checkmarkr\checkmarkc\checkmark$\checkmarku\checkmark$\checkmark^\checkmark\\checkmarkc\checkmarki\checkmarkr\checkmarkc\checkmark$\checkmarks\checkmark$\checkmark^\checkmark\\checkmarkc\checkmarki\checkmarkr\checkmarkc\checkmark$\checkmark \checkmark$\checkmark^\checkmark\\checkmarkc\checkmarki\checkmarkr\checkmarkc\checkmark$\checkmarks\checkmark$\checkmark^\checkmark\\checkmarkc\checkmarki\checkmarkr\checkmarkc\checkmark$\checkmarko\checkmark$\checkmark^\checkmark\\checkmarkc\checkmarki\checkmarkr\checkmarkc\checkmark$\checkmarkl\checkmark$\checkmark^\checkmark\\checkmarkc\checkmarki\checkmarkr\checkmarkc\checkmark$\checkmarkl\checkmark$\checkmark^\checkmark\\checkmarkc\checkmarki\checkmarkr\checkmarkc\checkmark$\checkmarki\checkmark$\checkmark^\checkmark\\checkmarkc\checkmarki\checkmarkr\checkmarkc\checkmark$\checkmarkc\checkmark$\checkmark^\checkmark\\checkmarkc\checkmarki\checkmarkr\checkmarkc\checkmark$\checkmarki\checkmark$\checkmark^\checkmark\\checkmarkc\checkmarki\checkmarkr\checkmarkc\checkmark$\checkmarkt\checkmark$\checkmark^\checkmark\\checkmarkc\checkmarki\checkmarkr\checkmarkc\checkmark$\checkmarké\checkmark$\checkmark^\checkmark\\checkmarkc\checkmarki\checkmarkr\checkmarkc\checkmark$\checkmark \checkmark$\checkmark^\checkmark\\checkmarkc\checkmarki\checkmarkr\checkmarkc\checkmark$\checkmarkd\checkmark$\checkmark^\checkmark\\checkmarkc\checkmarki\checkmarkr\checkmarkc\checkmark$\checkmarke\checkmark$\checkmark^\checkmark\\checkmarkc\checkmarki\checkmarkr\checkmarkc\checkmark$\checkmark \checkmark$\checkmark^\checkmark\\checkmarkc\checkmarki\checkmarkr\checkmarkc\checkmark$\checkmarkl\checkmark$\checkmark^\checkmark\\checkmarkc\checkmarki\checkmarkr\checkmarkc\checkmark$\checkmarka\checkmark$\checkmark^\checkmark\\checkmarkc\checkmarki\checkmarkr\checkmarkc\checkmark$\checkmark \checkmark$\checkmark^\checkmark\\checkmarkc\checkmarki\checkmarkr\checkmarkc\checkmark$\checkmarkm\checkmark$\checkmark^\checkmark\\checkmarkc\checkmarki\checkmarkr\checkmarkc\checkmark$\checkmarka\checkmark$\checkmark^\checkmark\\checkmarkc\checkmarki\checkmarkr\checkmarkc\checkmark$\checkmarkc\checkmark$\checkmark^\checkmark\\checkmarkc\checkmarki\checkmarkr\checkmarkc\checkmark$\checkmarkh\checkmark$\checkmark^\checkmark\\checkmarkc\checkmarki\checkmarkr\checkmarkc\checkmark$\checkmarki\checkmark$\checkmark^\checkmark\\checkmarkc\checkmarki\checkmarkr\checkmarkc\checkmark$\checkmarkn\checkmark$\checkmark^\checkmark\\checkmarkc\checkmarki\checkmarkr\checkmarkc\checkmark$\checkmarke\checkmark$\checkmark^\checkmark\\checkmarkc\checkmarki\checkmarkr\checkmarkc\checkmark$\checkmark.\checkmark$\checkmark^\checkmark\\checkmarkc\checkmarki\checkmarkr\checkmarkc\checkmark$\checkmark \checkmark$\checkmark^\checkmark\\checkmarkc\checkmarki\checkmarkr\checkmarkc\checkmark$\checkmarkI\checkmark$\checkmark^\checkmark\\checkmarkc\checkmarki\checkmarkr\checkmarkc\checkmark$\checkmarkl\checkmark$\checkmark^\checkmark\\checkmarkc\checkmarki\checkmarkr\checkmarkc\checkmark$\checkmark \checkmark$\checkmark^\checkmark\\checkmarkc\checkmarki\checkmarkr\checkmarkc\checkmark$\checkmarkd\checkmark$\checkmark^\checkmark\\checkmarkc\checkmarki\checkmarkr\checkmarkc\checkmark$\checkmarko\checkmark$\checkmark^\checkmark\\checkmarkc\checkmarki\checkmarkr\checkmarkc\checkmark$\checkmarki\checkmark$\checkmark^\checkmark\\checkmarkc\checkmarki\checkmarkr\checkmarkc\checkmark$\checkmarkt\checkmark$\checkmark^\checkmark\\checkmarkc\checkmarki\checkmarkr\checkmarkc\checkmark$\checkmark \checkmark$\checkmark^\checkmark\\checkmarkc\checkmarki\checkmarkr\checkmarkc\checkmark$\checkmarks\checkmark$\checkmark^\checkmark\\checkmarkc\checkmarki\checkmarkr\checkmarkc\checkmark$\checkmarki\checkmark$\checkmark^\checkmark\\checkmarkc\checkmarki\checkmarkr\checkmarkc\checkmark$\checkmarkm\checkmark$\checkmark^\checkmark\\checkmarkc\checkmarki\checkmarkr\checkmarkc\checkmark$\checkmarku\checkmark$\checkmark^\checkmark\\checkmarkc\checkmarki\checkmarkr\checkmarkc\checkmark$\checkmarkl\checkmark$\checkmark^\checkmark\\checkmarkc\checkmarki\checkmarkr\checkmarkc\checkmark$\checkmarkt\checkmark$\checkmark^\checkmark\\checkmarkc\checkmarki\checkmarkr\checkmarkc\checkmark$\checkmarka\checkmark$\checkmark^\checkmark\\checkmarkc\checkmarki\checkmarkr\checkmarkc\checkmark$\checkmarkn\checkmark$\checkmark^\checkmark\\checkmarkc\checkmarki\checkmarkr\checkmarkc\checkmark$\checkmarké\checkmark$\checkmark^\checkmark\\checkmarkc\checkmarki\checkmarkr\checkmarkc\checkmark$\checkmarkm\checkmark$\checkmark^\checkmark\\checkmarkc\checkmarki\checkmarkr\checkmarkc\checkmark$\checkmarke\checkmark$\checkmark^\checkmark\\checkmarkc\checkmarki\checkmarkr\checkmarkc\checkmark$\checkmarkn\checkmark$\checkmark^\checkmark\\checkmarkc\checkmarki\checkmarkr\checkmarkc\checkmark$\checkmarkt\checkmark$\checkmark^\checkmark\\checkmarkc\checkmarki\checkmarkr\checkmarkc\checkmark$\checkmark \checkmark$\checkmark^\checkmark\\checkmarkc\checkmarki\checkmarkr\checkmarkc\checkmark$\checkmark:\checkmark$\checkmark^\checkmark\\checkmarkc\checkmarki\checkmarkr\checkmarkc\checkmark$\checkmark
\checkmark$\checkmark^\checkmark\\checkmarkc\checkmarki\checkmarkr\checkmarkc\checkmark$\checkmark\\checkmark$\checkmark^\checkmark\\checkmarkc\checkmarki\checkmarkr\checkmarkc\checkmark$\checkmarkb\checkmark$\checkmark^\checkmark\\checkmarkc\checkmarki\checkmarkr\checkmarkc\checkmark$\checkmarke\checkmark$\checkmark^\checkmark\\checkmarkc\checkmarki\checkmarkr\checkmarkc\checkmark$\checkmarkg\checkmark$\checkmark^\checkmark\\checkmarkc\checkmarki\checkmarkr\checkmarkc\checkmark$\checkmarki\checkmark$\checkmark^\checkmark\\checkmarkc\checkmarki\checkmarkr\checkmarkc\checkmark$\checkmarkn\checkmark$\checkmark^\checkmark\\checkmarkc\checkmarki\checkmarkr\checkmarkc\checkmark$\checkmark{\checkmark$\checkmark^\checkmark\\checkmarkc\checkmarki\checkmarkr\checkmarkc\checkmark$\checkmarki\checkmark$\checkmark^\checkmark\\checkmarkc\checkmarki\checkmarkr\checkmarkc\checkmark$\checkmarkt\checkmark$\checkmark^\checkmark\\checkmarkc\checkmarki\checkmarkr\checkmarkc\checkmark$\checkmarke\checkmark$\checkmark^\checkmark\\checkmarkc\checkmarki\checkmarkr\checkmarkc\checkmark$\checkmarkm\checkmark$\checkmark^\checkmark\\checkmarkc\checkmarki\checkmarkr\checkmarkc\checkmark$\checkmarki\checkmark$\checkmark^\checkmark\\checkmarkc\checkmarki\checkmarkr\checkmarkc\checkmark$\checkmarkz\checkmark$\checkmark^\checkmark\\checkmarkc\checkmarki\checkmarkr\checkmarkc\checkmark$\checkmarke\checkmark$\checkmark^\checkmark\\checkmarkc\checkmarki\checkmarkr\checkmarkc\checkmark$\checkmark}\checkmark$\checkmark^\checkmark\\checkmarkc\checkmarki\checkmarkr\checkmarkc\checkmark$\checkmark
\checkmark$\checkmark^\checkmark\\checkmarkc\checkmarki\checkmarkr\checkmarkc\checkmark$\checkmark \checkmark$\checkmark^\checkmark\\checkmarkc\checkmarki\checkmarkr\checkmarkc\checkmark$\checkmark \checkmark$\checkmark^\checkmark\\checkmarkc\checkmarki\checkmarkr\checkmarkc\checkmark$\checkmark \checkmark$\checkmark^\checkmark\\checkmarkc\checkmarki\checkmarkr\checkmarkc\checkmark$\checkmark \checkmark$\checkmark^\checkmark\\checkmarkc\checkmarki\checkmarkr\checkmarkc\checkmark$\checkmark\\checkmark$\checkmark^\checkmark\\checkmarkc\checkmarki\checkmarkr\checkmarkc\checkmark$\checkmarki\checkmark$\checkmark^\checkmark\\checkmarkc\checkmarki\checkmarkr\checkmarkc\checkmark$\checkmarkt\checkmark$\checkmark^\checkmark\\checkmarkc\checkmarki\checkmarkr\checkmarkc\checkmark$\checkmarke\checkmark$\checkmark^\checkmark\\checkmarkc\checkmarki\checkmarkr\checkmarkc\checkmark$\checkmarkm\checkmark$\checkmark^\checkmark\\checkmarkc\checkmarki\checkmarkr\checkmarkc\checkmark$\checkmark \checkmark$\checkmark^\checkmark\\checkmarkc\checkmarki\checkmarkr\checkmarkc\checkmark$\checkmarkT\checkmark$\checkmark^\checkmark\\checkmarkc\checkmarki\checkmarkr\checkmarkc\checkmark$\checkmarkr\checkmark$\checkmark^\checkmark\\checkmarkc\checkmarki\checkmarkr\checkmarkc\checkmark$\checkmarka\checkmark$\checkmark^\checkmark\\checkmarkc\checkmarki\checkmarkr\checkmarkc\checkmark$\checkmarkn\checkmark$\checkmark^\checkmark\\checkmarkc\checkmarki\checkmarkr\checkmarkc\checkmark$\checkmarks\checkmark$\checkmark^\checkmark\\checkmarkc\checkmarki\checkmarkr\checkmarkc\checkmark$\checkmarkm\checkmark$\checkmark^\checkmark\\checkmarkc\checkmarki\checkmarkr\checkmarkc\checkmark$\checkmarke\checkmark$\checkmark^\checkmark\\checkmarkc\checkmarki\checkmarkr\checkmarkc\checkmark$\checkmarkt\checkmark$\checkmark^\checkmark\\checkmarkc\checkmarki\checkmarkr\checkmarkc\checkmark$\checkmarkt\checkmark$\checkmark^\checkmark\\checkmarkc\checkmarki\checkmarkr\checkmarkc\checkmark$\checkmarkr\checkmark$\checkmark^\checkmark\\checkmarkc\checkmarki\checkmarkr\checkmarkc\checkmark$\checkmarke\checkmark$\checkmark^\checkmark\\checkmarkc\checkmarki\checkmarkr\checkmarkc\checkmark$\checkmark \checkmark$\checkmark^\checkmark\\checkmarkc\checkmarki\checkmarkr\checkmarkc\checkmark$\checkmarkl\checkmark$\checkmark^\checkmark\\checkmarkc\checkmarki\checkmarkr\checkmarkc\checkmark$\checkmarke\checkmark$\checkmark^\checkmark\\checkmarkc\checkmarki\checkmarkr\checkmarkc\checkmark$\checkmark \checkmark$\checkmark^\checkmark\\checkmarkc\checkmarki\checkmarkr\checkmarkc\checkmark$\checkmarkc\checkmark$\checkmark^\checkmark\\checkmarkc\checkmarki\checkmarkr\checkmarkc\checkmark$\checkmarko\checkmark$\checkmark^\checkmark\\checkmarkc\checkmarki\checkmarkr\checkmarkc\checkmark$\checkmarku\checkmark$\checkmark^\checkmark\\checkmarkc\checkmarki\checkmarkr\checkmarkc\checkmark$\checkmarkp\checkmark$\checkmark^\checkmark\\checkmarkc\checkmarki\checkmarkr\checkmarkc\checkmark$\checkmarkl\checkmark$\checkmark^\checkmark\\checkmarkc\checkmarki\checkmarkr\checkmarkc\checkmark$\checkmarke\checkmark$\checkmark^\checkmark\\checkmarkc\checkmarki\checkmarkr\checkmarkc\checkmark$\checkmark \checkmark$\checkmark^\checkmark\\checkmarkc\checkmarki\checkmarkr\checkmarkc\checkmark$\checkmarkd\checkmark$\checkmark^\checkmark\\checkmarkc\checkmarki\checkmarkr\checkmarkc\checkmark$\checkmarke\checkmark$\checkmark^\checkmark\\checkmarkc\checkmarki\checkmarkr\checkmarkc\checkmark$\checkmark \checkmark$\checkmark^\checkmark\\checkmarkc\checkmarki\checkmarkr\checkmarkc\checkmark$\checkmarkm\checkmark$\checkmark^\checkmark\\checkmarkc\checkmarki\checkmarkr\checkmarkc\checkmark$\checkmarka\checkmark$\checkmark^\checkmark\\checkmarkc\checkmarki\checkmarkr\checkmarkc\checkmark$\checkmarki\checkmark$\checkmark^\checkmark\\checkmarkc\checkmarki\checkmarkr\checkmarkc\checkmark$\checkmarkn\checkmark$\checkmark^\checkmark\\checkmarkc\checkmarki\checkmarkr\checkmarkc\checkmark$\checkmarkt\checkmark$\checkmark^\checkmark\\checkmarkc\checkmarki\checkmarkr\checkmarkc\checkmark$\checkmarki\checkmark$\checkmark^\checkmark\\checkmarkc\checkmarki\checkmarkr\checkmarkc\checkmark$\checkmarke\checkmark$\checkmark^\checkmark\\checkmarkc\checkmarki\checkmarkr\checkmarkc\checkmark$\checkmarkn\checkmark$\checkmark^\checkmark\\checkmarkc\checkmarki\checkmarkr\checkmarkc\checkmark$\checkmark \checkmark$\checkmark^\checkmark\\checkmarkc\checkmarki\checkmarkr\checkmarkc\checkmark$\checkmarkd\checkmark$\checkmark^\checkmark\\checkmarkc\checkmarki\checkmarkr\checkmarkc\checkmark$\checkmarku\checkmark$\checkmark^\checkmark\\checkmarkc\checkmarki\checkmarkr\checkmarkc\checkmark$\checkmark \checkmark$\checkmark^\checkmark\\checkmarkc\checkmarki\checkmarkr\checkmarkc\checkmark$\checkmarkm\checkmark$\checkmark^\checkmark\\checkmarkc\checkmarki\checkmarkr\checkmarkc\checkmark$\checkmarko\checkmark$\checkmark^\checkmark\\checkmarkc\checkmarki\checkmarkr\checkmarkc\checkmark$\checkmarkt\checkmark$\checkmark^\checkmark\\checkmarkc\checkmarki\checkmarkr\checkmarkc\checkmark$\checkmarke\checkmark$\checkmark^\checkmark\\checkmarkc\checkmarki\checkmarkr\checkmarkc\checkmark$\checkmarku\checkmark$\checkmark^\checkmark\\checkmarkc\checkmarki\checkmarkr\checkmarkc\checkmark$\checkmarkr\checkmark$\checkmark^\checkmark\\checkmarkc\checkmarki\checkmarkr\checkmarkc\checkmark$\checkmark \checkmark$\checkmark^\checkmark\\checkmarkc\checkmarki\checkmarkr\checkmarkc\checkmark$\checkmark(\checkmark$\checkmark^\checkmark\\checkmarkc\checkmarki\checkmarkr\checkmarkc\checkmark$\checkmarkv\checkmark$\checkmark^\checkmark\\checkmarkc\checkmarki\checkmarkr\checkmarkc\checkmark$\checkmarki\checkmark$\checkmark^\checkmark\\checkmarkc\checkmarki\checkmarkr\checkmarkc\checkmark$\checkmarka\checkmark$\checkmark^\checkmark\\checkmarkc\checkmarki\checkmarkr\checkmarkc\checkmark$\checkmark \checkmark$\checkmark^\checkmark\\checkmarkc\checkmarki\checkmarkr\checkmarkc\checkmark$\checkmarkl\checkmark$\checkmark^\checkmark\\checkmarkc\checkmarki\checkmarkr\checkmarkc\checkmark$\checkmarka\checkmark$\checkmark^\checkmark\\checkmarkc\checkmarki\checkmarkr\checkmarkc\checkmark$\checkmark \checkmark$\checkmark^\checkmark\\checkmarkc\checkmarki\checkmarkr\checkmarkc\checkmark$\checkmarkc\checkmark$\checkmark^\checkmark\\checkmarkc\checkmarki\checkmarkr\checkmarkc\checkmark$\checkmarko\checkmark$\checkmark^\checkmark\\checkmarkc\checkmarki\checkmarkr\checkmarkc\checkmark$\checkmarku\checkmark$\checkmark^\checkmark\\checkmarkc\checkmarki\checkmarkr\checkmarkc\checkmark$\checkmarkr\checkmark$\checkmark^\checkmark\\checkmarkc\checkmarki\checkmarkr\checkmarkc\checkmark$\checkmarkr\checkmark$\checkmark^\checkmark\\checkmarkc\checkmarki\checkmarkr\checkmarkc\checkmark$\checkmarko\checkmark$\checkmark^\checkmark\\checkmarkc\checkmarki\checkmarkr\checkmarkc\checkmark$\checkmarki\checkmark$\checkmark^\checkmark\\checkmarkc\checkmarki\checkmarkr\checkmarkc\checkmark$\checkmarke\checkmark$\checkmark^\checkmark\\checkmarkc\checkmarki\checkmarkr\checkmarkc\checkmark$\checkmark)\checkmark$\checkmark^\checkmark\\checkmarkc\checkmarki\checkmarkr\checkmarkc\checkmark$\checkmark \checkmark$\checkmark^\checkmark\\checkmarkc\checkmarki\checkmarkr\checkmarkc\checkmark$\checkmark;\checkmark$\checkmark^\checkmark\\checkmarkc\checkmarki\checkmarkr\checkmarkc\checkmark$\checkmark
\checkmark$\checkmark^\checkmark\\checkmarkc\checkmarki\checkmarkr\checkmarkc\checkmark$\checkmark \checkmark$\checkmark^\checkmark\\checkmarkc\checkmarki\checkmarkr\checkmarkc\checkmark$\checkmark \checkmark$\checkmark^\checkmark\\checkmarkc\checkmarki\checkmarkr\checkmarkc\checkmark$\checkmark \checkmark$\checkmark^\checkmark\\checkmarkc\checkmarki\checkmarkr\checkmarkc\checkmark$\checkmark \checkmark$\checkmark^\checkmark\\checkmarkc\checkmarki\checkmarkr\checkmarkc\checkmark$\checkmark\\checkmark$\checkmark^\checkmark\\checkmarkc\checkmarki\checkmarkr\checkmarkc\checkmark$\checkmarki\checkmark$\checkmark^\checkmark\\checkmarkc\checkmarki\checkmarkr\checkmarkc\checkmark$\checkmarkt\checkmark$\checkmark^\checkmark\\checkmarkc\checkmarki\checkmarkr\checkmarkc\checkmark$\checkmarke\checkmark$\checkmark^\checkmark\\checkmarkc\checkmarki\checkmarkr\checkmarkc\checkmark$\checkmarkm\checkmark$\checkmark^\checkmark\\checkmarkc\checkmarki\checkmarkr\checkmarkc\checkmark$\checkmark \checkmark$\checkmark^\checkmark\\checkmarkc\checkmarki\checkmarkr\checkmarkc\checkmark$\checkmarkR\checkmark$\checkmark^\checkmark\\checkmarkc\checkmarki\checkmarkr\checkmarkc\checkmark$\checkmarké\checkmark$\checkmark^\checkmark\\checkmarkc\checkmarki\checkmarkr\checkmarkc\checkmark$\checkmarks\checkmark$\checkmark^\checkmark\\checkmarkc\checkmarki\checkmarkr\checkmarkc\checkmark$\checkmarki\checkmark$\checkmark^\checkmark\\checkmarkc\checkmarki\checkmarkr\checkmarkc\checkmark$\checkmarks\checkmark$\checkmark^\checkmark\\checkmarkc\checkmarki\checkmarkr\checkmarkc\checkmark$\checkmarkt\checkmark$\checkmark^\checkmark\\checkmarkc\checkmarki\checkmarkr\checkmarkc\checkmark$\checkmarke\checkmark$\checkmark^\checkmark\\checkmarkc\checkmarki\checkmarkr\checkmarkc\checkmark$\checkmarkr\checkmark$\checkmark^\checkmark\\checkmarkc\checkmarki\checkmarkr\checkmarkc\checkmark$\checkmark \checkmark$\checkmark^\checkmark\\checkmarkc\checkmarki\checkmarkr\checkmarkc\checkmark$\checkmarka\checkmark$\checkmark^\checkmark\\checkmarkc\checkmarki\checkmarkr\checkmarkc\checkmark$\checkmarku\checkmark$\checkmark^\checkmark\\checkmarkc\checkmarki\checkmarkr\checkmarkc\checkmark$\checkmark \checkmark$\checkmark^\checkmark\\checkmarkc\checkmarki\checkmarkr\checkmarkc\checkmark$\checkmarkm\checkmark$\checkmark^\checkmark\\checkmarkc\checkmarki\checkmarkr\checkmarkc\checkmark$\checkmarko\checkmark$\checkmark^\checkmark\\checkmarkc\checkmarki\checkmarkr\checkmarkc\checkmark$\checkmarkm\checkmark$\checkmark^\checkmark\\checkmarkc\checkmarki\checkmarkr\checkmarkc\checkmark$\checkmarke\checkmark$\checkmark^\checkmark\\checkmarkc\checkmarki\checkmarkr\checkmarkc\checkmark$\checkmarkn\checkmark$\checkmark^\checkmark\\checkmarkc\checkmarki\checkmarkr\checkmarkc\checkmark$\checkmarkt\checkmark$\checkmark^\checkmark\\checkmarkc\checkmarki\checkmarkr\checkmarkc\checkmark$\checkmark \checkmark$\checkmark^\checkmark\\checkmarkc\checkmarki\checkmarkr\checkmarkc\checkmark$\checkmarkd\checkmark$\checkmark^\checkmark\\checkmarkc\checkmarki\checkmarkr\checkmarkc\checkmark$\checkmarke\checkmark$\checkmark^\checkmark\\checkmarkc\checkmarki\checkmarkr\checkmarkc\checkmark$\checkmark \checkmark$\checkmark^\checkmark\\checkmarkc\checkmarki\checkmarkr\checkmarkc\checkmark$\checkmarkr\checkmark$\checkmark^\checkmark\\checkmarkc\checkmarki\checkmarkr\checkmarkc\checkmark$\checkmarke\checkmark$\checkmark^\checkmark\\checkmarkc\checkmarki\checkmarkr\checkmarkc\checkmark$\checkmarkn\checkmark$\checkmark^\checkmark\\checkmarkc\checkmarki\checkmarkr\checkmarkc\checkmark$\checkmarkv\checkmark$\checkmark^\checkmark\\checkmarkc\checkmarki\checkmarkr\checkmarkc\checkmark$\checkmarke\checkmark$\checkmark^\checkmark\\checkmarkc\checkmarki\checkmarkr\checkmarkc\checkmark$\checkmarkr\checkmark$\checkmark^\checkmark\\checkmarkc\checkmarki\checkmarkr\checkmarkc\checkmark$\checkmarks\checkmark$\checkmark^\checkmark\\checkmarkc\checkmarki\checkmarkr\checkmarkc\checkmark$\checkmarke\checkmark$\checkmark^\checkmark\\checkmarkc\checkmarki\checkmarkr\checkmarkc\checkmark$\checkmarkm\checkmark$\checkmark^\checkmark\\checkmarkc\checkmarki\checkmarkr\checkmarkc\checkmark$\checkmarke\checkmark$\checkmark^\checkmark\\checkmarkc\checkmarki\checkmarkr\checkmarkc\checkmark$\checkmarkn\checkmark$\checkmark^\checkmark\\checkmarkc\checkmarki\checkmarkr\checkmarkc\checkmark$\checkmarkt\checkmark$\checkmark^\checkmark\\checkmarkc\checkmarki\checkmarkr\checkmarkc\checkmark$\checkmark \checkmark$\checkmark^\checkmark\\checkmarkc\checkmarki\checkmarkr\checkmarkc\checkmark$\checkmarkg\checkmark$\checkmark^\checkmark\\checkmarkc\checkmarki\checkmarkr\checkmarkc\checkmark$\checkmarké\checkmark$\checkmark^\checkmark\\checkmarkc\checkmarki\checkmarkr\checkmarkc\checkmark$\checkmarkn\checkmark$\checkmark^\checkmark\\checkmarkc\checkmarki\checkmarkr\checkmarkc\checkmark$\checkmarké\checkmark$\checkmark^\checkmark\\checkmarkc\checkmarki\checkmarkr\checkmarkc\checkmark$\checkmarkr\checkmark$\checkmark^\checkmark\\checkmarkc\checkmarki\checkmarkr\checkmarkc\checkmark$\checkmarké\checkmark$\checkmark^\checkmark\\checkmarkc\checkmarki\checkmarkr\checkmarkc\checkmark$\checkmark \checkmark$\checkmark^\checkmark\\checkmarkc\checkmarki\checkmarkr\checkmarkc\checkmark$\checkmarkp\checkmark$\checkmark^\checkmark\\checkmarkc\checkmarki\checkmarkr\checkmarkc\checkmark$\checkmarka\checkmark$\checkmark^\checkmark\\checkmarkc\checkmarki\checkmarkr\checkmarkc\checkmark$\checkmarkr\checkmark$\checkmark^\checkmark\\checkmarkc\checkmarki\checkmarkr\checkmarkc\checkmark$\checkmark \checkmark$\checkmark^\checkmark\\checkmarkc\checkmarki\checkmarkr\checkmarkc\checkmark$\checkmarkl\checkmark$\checkmark^\checkmark\\checkmarkc\checkmarki\checkmarkr\checkmarkc\checkmark$\checkmarka\checkmark$\checkmark^\checkmark\\checkmarkc\checkmarki\checkmarkr\checkmarkc\checkmark$\checkmark \checkmark$\checkmark^\checkmark\\checkmarkc\checkmarki\checkmarkr\checkmarkc\checkmark$\checkmarkf\checkmark$\checkmark^\checkmark\\checkmarkc\checkmarki\checkmarkr\checkmarkc\checkmark$\checkmarko\checkmark$\checkmark^\checkmark\\checkmarkc\checkmarki\checkmarkr\checkmarkc\checkmark$\checkmarkr\checkmark$\checkmark^\checkmark\\checkmarkc\checkmarki\checkmarkr\checkmarkc\checkmark$\checkmarkc\checkmark$\checkmark^\checkmark\\checkmarkc\checkmarki\checkmarkr\checkmarkc\checkmark$\checkmarke\checkmark$\checkmark^\checkmark\\checkmarkc\checkmarki\checkmarkr\checkmarkc\checkmark$\checkmark \checkmark$\checkmark^\checkmark\\checkmarkc\checkmarki\checkmarkr\checkmarkc\checkmark$\checkmarkd\checkmark$\checkmark^\checkmark\\checkmarkc\checkmarki\checkmarkr\checkmarkc\checkmark$\checkmarke\checkmark$\checkmark^\checkmark\\checkmarkc\checkmarki\checkmarkr\checkmarkc\checkmark$\checkmark \checkmark$\checkmark^\checkmark\\checkmarkc\checkmarki\checkmarkr\checkmarkc\checkmark$\checkmarkc\checkmark$\checkmark^\checkmark\\checkmarkc\checkmarki\checkmarkr\checkmarkc\checkmark$\checkmarko\checkmark$\checkmark^\checkmark\\checkmarkc\checkmarki\checkmarkr\checkmarkc\checkmark$\checkmarku\checkmark$\checkmark^\checkmark\\checkmarkc\checkmarki\checkmarkr\checkmarkc\checkmark$\checkmarkp\checkmark$\checkmark^\checkmark\\checkmarkc\checkmarki\checkmarkr\checkmarkc\checkmark$\checkmarke\checkmark$\checkmark^\checkmark\\checkmarkc\checkmarki\checkmarkr\checkmarkc\checkmark$\checkmark \checkmark$\checkmark^\checkmark\\checkmarkc\checkmarki\checkmarkr\checkmarkc\checkmark$\checkmark$\checkmark$\checkmark^\checkmark\\checkmarkc\checkmarki\checkmarkr\checkmarkc\checkmark$\checkmarkF\checkmark$\checkmark^\checkmark\\checkmarkc\checkmarki\checkmarkr\checkmarkc\checkmark$\checkmark_\checkmark$\checkmark^\checkmark\\checkmarkc\checkmarki\checkmarkr\checkmarkc\checkmark$\checkmark{\checkmark$\checkmark^\checkmark\\checkmarkc\checkmarki\checkmarkr\checkmarkc\checkmark$\checkmarkm\checkmark$\checkmark^\checkmark\\checkmarkc\checkmarki\checkmarkr\checkmarkc\checkmark$\checkmarka\checkmark$\checkmark^\checkmark\\checkmarkc\checkmarki\checkmarkr\checkmarkc\checkmark$\checkmarkx\checkmark$\checkmark^\checkmark\\checkmarkc\checkmarki\checkmarkr\checkmarkc\checkmark$\checkmark}\checkmark$\checkmark^\checkmark\\checkmarkc\checkmarki\checkmarkr\checkmarkc\checkmark$\checkmark \checkmark$\checkmark^\checkmark\\checkmarkc\checkmarki\checkmarkr\checkmarkc\checkmark$\checkmark=\checkmark$\checkmark^\checkmark\\checkmarkc\checkmarki\checkmarkr\checkmarkc\checkmark$\checkmark \checkmark$\checkmark^\checkmark\\checkmarkc\checkmarki\checkmarkr\checkmarkc\checkmark$\checkmark1\checkmark$\checkmark^\checkmark\\checkmarkc\checkmarki\checkmarkr\checkmarkc\checkmark$\checkmark0\checkmark$\checkmark^\checkmark\\checkmarkc\checkmarki\checkmarkr\checkmarkc\checkmark$\checkmark0\checkmark$\checkmark^\checkmark\\checkmarkc\checkmarki\checkmarkr\checkmarkc\checkmark$\checkmark$\checkmark$\checkmark^\checkmark\\checkmarkc\checkmarki\checkmarkr\checkmarkc\checkmark$\checkmark \checkmark$\checkmark^\checkmark\\checkmarkc\checkmarki\checkmarkr\checkmarkc\checkmark$\checkmarkN\checkmark$\checkmark^\checkmark\\checkmarkc\checkmarki\checkmarkr\checkmarkc\checkmark$\checkmark \checkmark$\checkmark^\checkmark\\checkmarkc\checkmarki\checkmarkr\checkmarkc\checkmark$\checkmarkà\checkmark$\checkmark^\checkmark\\checkmarkc\checkmarki\checkmarkr\checkmarkc\checkmark$\checkmark \checkmark$\checkmark^\checkmark\\checkmarkc\checkmarki\checkmarkr\checkmarkc\checkmark$\checkmarkh\checkmark$\checkmark^\checkmark\\checkmarkc\checkmarki\checkmarkr\checkmarkc\checkmark$\checkmarka\checkmark$\checkmark^\checkmark\\checkmarkc\checkmarki\checkmarkr\checkmarkc\checkmark$\checkmarku\checkmark$\checkmark^\checkmark\\checkmarkc\checkmarki\checkmarkr\checkmarkc\checkmark$\checkmarkt\checkmark$\checkmark^\checkmark\\checkmarkc\checkmarki\checkmarkr\checkmarkc\checkmark$\checkmarke\checkmark$\checkmark^\checkmark\\checkmarkc\checkmarki\checkmarkr\checkmarkc\checkmark$\checkmarku\checkmark$\checkmark^\checkmark\\checkmarkc\checkmarki\checkmarkr\checkmarkc\checkmark$\checkmarkr\checkmark$\checkmark^\checkmark\\checkmarkc\checkmarki\checkmarkr\checkmarkc\checkmark$\checkmark \checkmark$\checkmark^\checkmark\\checkmarkc\checkmarki\checkmarkr\checkmarkc\checkmark$\checkmark$\checkmark$\checkmark^\checkmark\\checkmarkc\checkmarki\checkmarkr\checkmarkc\checkmark$\checkmarkH\checkmark$\checkmark^\checkmark\\checkmarkc\checkmarki\checkmarkr\checkmarkc\checkmark$\checkmark \checkmark$\checkmark^\checkmark\\checkmarkc\checkmarki\checkmarkr\checkmarkc\checkmark$\checkmark=\checkmark$\checkmark^\checkmark\\checkmarkc\checkmarki\checkmarkr\checkmarkc\checkmark$\checkmark \checkmark$\checkmark^\checkmark\\checkmarkc\checkmarki\checkmarkr\checkmarkc\checkmark$\checkmark2\checkmark$\checkmark^\checkmark\\checkmarkc\checkmarki\checkmarkr\checkmarkc\checkmark$\checkmark0\checkmark$\checkmark^\checkmark\\checkmarkc\checkmarki\checkmarkr\checkmarkc\checkmark$\checkmark0\checkmark$\checkmark^\checkmark\\checkmarkc\checkmarki\checkmarkr\checkmarkc\checkmark$\checkmark$\checkmark$\checkmark^\checkmark\\checkmarkc\checkmarki\checkmarkr\checkmarkc\checkmark$\checkmark \checkmark$\checkmark^\checkmark\\checkmarkc\checkmarki\checkmarkr\checkmarkc\checkmark$\checkmarkm\checkmark$\checkmark^\checkmark\\checkmarkc\checkmarki\checkmarkr\checkmarkc\checkmark$\checkmarkm\checkmark$\checkmark^\checkmark\\checkmarkc\checkmarki\checkmarkr\checkmarkc\checkmark$\checkmark \checkmark$\checkmark^\checkmark\\checkmarkc\checkmarki\checkmarkr\checkmarkc\checkmark$\checkmark;\checkmark$\checkmark^\checkmark\\checkmarkc\checkmarki\checkmarkr\checkmarkc\checkmark$\checkmark
\checkmark$\checkmark^\checkmark\\checkmarkc\checkmarki\checkmarkr\checkmarkc\checkmark$\checkmark \checkmark$\checkmark^\checkmark\\checkmarkc\checkmarki\checkmarkr\checkmarkc\checkmark$\checkmark \checkmark$\checkmark^\checkmark\\checkmarkc\checkmarki\checkmarkr\checkmarkc\checkmark$\checkmark \checkmark$\checkmark^\checkmark\\checkmarkc\checkmarki\checkmarkr\checkmarkc\checkmark$\checkmark \checkmark$\checkmark^\checkmark\\checkmarkc\checkmarki\checkmarkr\checkmarkc\checkmark$\checkmark\\checkmark$\checkmark^\checkmark\\checkmarkc\checkmarki\checkmarkr\checkmarkc\checkmark$\checkmarki\checkmark$\checkmark^\checkmark\\checkmarkc\checkmarki\checkmarkr\checkmarkc\checkmark$\checkmarkt\checkmark$\checkmark^\checkmark\\checkmarkc\checkmarki\checkmarkr\checkmarkc\checkmark$\checkmarke\checkmark$\checkmark^\checkmark\\checkmarkc\checkmarki\checkmarkr\checkmarkc\checkmark$\checkmarkm\checkmark$\checkmark^\checkmark\\checkmarkc\checkmarki\checkmarkr\checkmarkc\checkmark$\checkmark \checkmark$\checkmark^\checkmark\\checkmarkc\checkmarki\checkmarkr\checkmarkc\checkmark$\checkmarkS\checkmark$\checkmark^\checkmark\\checkmarkc\checkmarki\checkmarkr\checkmarkc\checkmark$\checkmarku\checkmark$\checkmark^\checkmark\\checkmarkc\checkmarki\checkmarkr\checkmarkc\checkmark$\checkmarkp\checkmark$\checkmark^\checkmark\\checkmarkc\checkmarki\checkmarkr\checkmarkc\checkmark$\checkmarkp\checkmark$\checkmark^\checkmark\\checkmarkc\checkmarki\checkmarkr\checkmarkc\checkmark$\checkmarko\checkmark$\checkmark^\checkmark\\checkmarkc\checkmarki\checkmarkr\checkmarkc\checkmark$\checkmarkr\checkmark$\checkmark^\checkmark\\checkmarkc\checkmarki\checkmarkr\checkmarkc\checkmark$\checkmarkt\checkmark$\checkmark^\checkmark\\checkmarkc\checkmarki\checkmarkr\checkmarkc\checkmark$\checkmarke\checkmark$\checkmark^\checkmark\\checkmarkc\checkmarki\checkmarkr\checkmarkc\checkmark$\checkmarkr\checkmark$\checkmark^\checkmark\\checkmarkc\checkmarki\checkmarkr\checkmarkc\checkmark$\checkmark \checkmark$\checkmark^\checkmark\\checkmarkc\checkmarki\checkmarkr\checkmarkc\checkmark$\checkmarkl\checkmark$\checkmark^\checkmark\\checkmarkc\checkmarki\checkmarkr\checkmarkc\checkmark$\checkmarke\checkmark$\checkmark^\checkmark\\checkmarkc\checkmarki\checkmarkr\checkmarkc\checkmark$\checkmark \checkmark$\checkmark^\checkmark\\checkmarkc\checkmarki\checkmarkr\checkmarkc\checkmark$\checkmarkp\checkmark$\checkmark^\checkmark\\checkmarkc\checkmarki\checkmarkr\checkmarkc\checkmark$\checkmarko\checkmark$\checkmark^\checkmark\\checkmarkc\checkmarki\checkmarkr\checkmarkc\checkmark$\checkmarki\checkmark$\checkmark^\checkmark\\checkmarkc\checkmarki\checkmarkr\checkmarkc\checkmark$\checkmarkd\checkmark$\checkmark^\checkmark\\checkmarkc\checkmarki\checkmarkr\checkmarkc\checkmark$\checkmarks\checkmark$\checkmark^\checkmark\\checkmarkc\checkmarki\checkmarkr\checkmarkc\checkmark$\checkmark \checkmark$\checkmark^\checkmark\\checkmarkc\checkmarki\checkmarkr\checkmarkc\checkmark$\checkmarkd\checkmark$\checkmark^\checkmark\\checkmarkc\checkmarki\checkmarkr\checkmarkc\checkmark$\checkmarku\checkmark$\checkmark^\checkmark\\checkmarkc\checkmarki\checkmarkr\checkmarkc\checkmark$\checkmark \checkmark$\checkmark^\checkmark\\checkmarkc\checkmarki\checkmarkr\checkmarkc\checkmark$\checkmarkp\checkmark$\checkmark^\checkmark\\checkmarkc\checkmarki\checkmarkr\checkmarkc\checkmark$\checkmarkl\checkmark$\checkmark^\checkmark\\checkmarkc\checkmarki\checkmarkr\checkmarkc\checkmark$\checkmarka\checkmark$\checkmark^\checkmark\\checkmarkc\checkmarki\checkmarkr\checkmarkc\checkmark$\checkmarkt\checkmark$\checkmark^\checkmark\\checkmarkc\checkmarki\checkmarkr\checkmarkc\checkmark$\checkmarke\checkmark$\checkmark^\checkmark\\checkmarkc\checkmarki\checkmarkr\checkmarkc\checkmark$\checkmarka\checkmark$\checkmark^\checkmark\\checkmarkc\checkmarki\checkmarkr\checkmarkc\checkmark$\checkmarku\checkmark$\checkmark^\checkmark\\checkmarkc\checkmarki\checkmarkr\checkmarkc\checkmark$\checkmark \checkmark$\checkmark^\checkmark\\checkmarkc\checkmarki\checkmarkr\checkmarkc\checkmark$\checkmarke\checkmark$\checkmark^\checkmark\\checkmarkc\checkmarki\checkmarkr\checkmarkc\checkmark$\checkmarkt\checkmark$\checkmark^\checkmark\\checkmarkc\checkmarki\checkmarkr\checkmarkc\checkmark$\checkmark \checkmark$\checkmark^\checkmark\\checkmarkc\checkmarki\checkmarkr\checkmarkc\checkmark$\checkmarkd\checkmark$\checkmark^\checkmark\\checkmarkc\checkmarki\checkmarkr\checkmarkc\checkmark$\checkmarke\checkmark$\checkmark^\checkmark\\checkmarkc\checkmarki\checkmarkr\checkmarkc\checkmark$\checkmark \checkmark$\checkmark^\checkmark\\checkmarkc\checkmarki\checkmarkr\checkmarkc\checkmark$\checkmarkl\checkmark$\checkmark^\checkmark\\checkmarkc\checkmarki\checkmarkr\checkmarkc\checkmark$\checkmarka\checkmark$\checkmark^\checkmark\\checkmarkc\checkmarki\checkmarkr\checkmarkc\checkmark$\checkmark \checkmark$\checkmark^\checkmark\\checkmarkc\checkmarki\checkmarkr\checkmarkc\checkmark$\checkmarkp\checkmark$\checkmark^\checkmark\\checkmarkc\checkmarki\checkmarkr\checkmarkc\checkmark$\checkmarki\checkmark$\checkmark^\checkmark\\checkmarkc\checkmarki\checkmarkr\checkmarkc\checkmark$\checkmarkè\checkmark$\checkmark^\checkmark\\checkmarkc\checkmarki\checkmarkr\checkmarkc\checkmark$\checkmarkc\checkmark$\checkmark^\checkmark\\checkmarkc\checkmarki\checkmarkr\checkmarkc\checkmark$\checkmarke\checkmark$\checkmark^\checkmark\\checkmarkc\checkmarki\checkmarkr\checkmarkc\checkmark$\checkmark.\checkmark$\checkmark^\checkmark\\checkmarkc\checkmarki\checkmarkr\checkmarkc\checkmark$\checkmark
\checkmark$\checkmark^\checkmark\\checkmarkc\checkmarki\checkmarkr\checkmarkc\checkmark$\checkmark\\checkmark$\checkmark^\checkmark\\checkmarkc\checkmarki\checkmarkr\checkmarkc\checkmark$\checkmarke\checkmark$\checkmark^\checkmark\\checkmarkc\checkmarki\checkmarkr\checkmarkc\checkmark$\checkmarkn\checkmark$\checkmark^\checkmark\\checkmarkc\checkmarki\checkmarkr\checkmarkc\checkmark$\checkmarkd\checkmark$\checkmark^\checkmark\\checkmarkc\checkmarki\checkmarkr\checkmarkc\checkmark$\checkmark{\checkmark$\checkmark^\checkmark\\checkmarkc\checkmarki\checkmarkr\checkmarkc\checkmark$\checkmarki\checkmark$\checkmark^\checkmark\\checkmarkc\checkmarki\checkmarkr\checkmarkc\checkmark$\checkmarkt\checkmark$\checkmark^\checkmark\\checkmarkc\checkmarki\checkmarkr\checkmarkc\checkmark$\checkmarke\checkmark$\checkmark^\checkmark\\checkmarkc\checkmarki\checkmarkr\checkmarkc\checkmark$\checkmarkm\checkmark$\checkmark^\checkmark\\checkmarkc\checkmarki\checkmarkr\checkmarkc\checkmark$\checkmarki\checkmark$\checkmark^\checkmark\\checkmarkc\checkmarki\checkmarkr\checkmarkc\checkmark$\checkmarkz\checkmark$\checkmark^\checkmark\\checkmarkc\checkmarki\checkmarkr\checkmarkc\checkmark$\checkmarke\checkmark$\checkmark^\checkmark\\checkmarkc\checkmarki\checkmarkr\checkmarkc\checkmark$\checkmark}\checkmark$\checkmark^\checkmark\\checkmarkc\checkmarki\checkmarkr\checkmarkc\checkmark$\checkmark
\checkmark$\checkmark^\checkmark\\checkmarkc\checkmarki\checkmarkr\checkmarkc\checkmark$\checkmark
\checkmark$\checkmark^\checkmark\\checkmarkc\checkmarki\checkmarkr\checkmarkc\checkmark$\checkmark\\checkmark$\checkmark^\checkmark\\checkmarkc\checkmarki\checkmarkr\checkmarkc\checkmark$\checkmarks\checkmark$\checkmark^\checkmark\\checkmarkc\checkmarki\checkmarkr\checkmarkc\checkmark$\checkmarku\checkmark$\checkmark^\checkmark\\checkmarkc\checkmarki\checkmarkr\checkmarkc\checkmark$\checkmarkb\checkmark$\checkmark^\checkmark\\checkmarkc\checkmarki\checkmarkr\checkmarkc\checkmark$\checkmarks\checkmark$\checkmark^\checkmark\\checkmarkc\checkmarki\checkmarkr\checkmarkc\checkmark$\checkmarke\checkmark$\checkmark^\checkmark\\checkmarkc\checkmarki\checkmarkr\checkmarkc\checkmark$\checkmarkc\checkmark$\checkmark^\checkmark\\checkmarkc\checkmarki\checkmarkr\checkmarkc\checkmark$\checkmarkt\checkmark$\checkmark^\checkmark\\checkmarkc\checkmarki\checkmarkr\checkmarkc\checkmark$\checkmarki\checkmark$\checkmark^\checkmark\\checkmarkc\checkmarki\checkmarkr\checkmarkc\checkmark$\checkmarko\checkmark$\checkmark^\checkmark\\checkmarkc\checkmarki\checkmarkr\checkmarkc\checkmark$\checkmarkn\checkmark$\checkmark^\checkmark\\checkmarkc\checkmarki\checkmarkr\checkmarkc\checkmark$\checkmark{\checkmark$\checkmark^\checkmark\\checkmarkc\checkmarki\checkmarkr\checkmarkc\checkmark$\checkmarkC\checkmark$\checkmark^\checkmark\\checkmarkc\checkmarki\checkmarkr\checkmarkc\checkmark$\checkmarkh\checkmark$\checkmark^\checkmark\\checkmarkc\checkmarki\checkmarkr\checkmarkc\checkmark$\checkmarko\checkmark$\checkmark^\checkmark\\checkmarkc\checkmarki\checkmarkr\checkmarkc\checkmark$\checkmarki\checkmark$\checkmark^\checkmark\\checkmarkc\checkmarki\checkmarkr\checkmarkc\checkmark$\checkmarkx\checkmark$\checkmark^\checkmark\\checkmarkc\checkmarki\checkmarkr\checkmarkc\checkmark$\checkmark \checkmark$\checkmark^\checkmark\\checkmarkc\checkmarki\checkmarkr\checkmarkc\checkmark$\checkmarkd\checkmark$\checkmark^\checkmark\\checkmarkc\checkmarki\checkmarkr\checkmarkc\checkmark$\checkmarku\checkmark$\checkmark^\checkmark\\checkmarkc\checkmarki\checkmarkr\checkmarkc\checkmark$\checkmark \checkmark$\checkmark^\checkmark\\checkmarkc\checkmarki\checkmarkr\checkmarkc\checkmark$\checkmarkm\checkmark$\checkmark^\checkmark\\checkmarkc\checkmarki\checkmarkr\checkmarkc\checkmark$\checkmarka\checkmark$\checkmark^\checkmark\\checkmarkc\checkmarki\checkmarkr\checkmarkc\checkmark$\checkmarkt\checkmark$\checkmark^\checkmark\\checkmarkc\checkmarki\checkmarkr\checkmarkc\checkmark$\checkmarké\checkmark$\checkmark^\checkmark\\checkmarkc\checkmarki\checkmarkr\checkmarkc\checkmark$\checkmarkr\checkmark$\checkmark^\checkmark\\checkmarkc\checkmarki\checkmarkr\checkmarkc\checkmark$\checkmarki\checkmark$\checkmark^\checkmark\\checkmarkc\checkmarki\checkmarkr\checkmarkc\checkmark$\checkmarka\checkmark$\checkmark^\checkmark\\checkmarkc\checkmarki\checkmarkr\checkmarkc\checkmark$\checkmarku\checkmark$\checkmark^\checkmark\\checkmarkc\checkmarki\checkmarkr\checkmarkc\checkmark$\checkmark}\checkmark$\checkmark^\checkmark\\checkmarkc\checkmarki\checkmarkr\checkmarkc\checkmark$\checkmark
\checkmark$\checkmark^\checkmark\\checkmarkc\checkmarki\checkmarkr\checkmarkc\checkmark$\checkmark\\checkmark$\checkmark^\checkmark\\checkmarkc\checkmarki\checkmarkr\checkmarkc\checkmark$\checkmarkb\checkmark$\checkmark^\checkmark\\checkmarkc\checkmarki\checkmarkr\checkmarkc\checkmark$\checkmarke\checkmark$\checkmark^\checkmark\\checkmarkc\checkmarki\checkmarkr\checkmarkc\checkmark$\checkmarkg\checkmark$\checkmark^\checkmark\\checkmarkc\checkmarki\checkmarkr\checkmarkc\checkmark$\checkmarki\checkmark$\checkmark^\checkmark\\checkmarkc\checkmarki\checkmarkr\checkmarkc\checkmark$\checkmarkn\checkmark$\checkmark^\checkmark\\checkmarkc\checkmarki\checkmarkr\checkmarkc\checkmark$\checkmark{\checkmark$\checkmark^\checkmark\\checkmarkc\checkmarki\checkmarkr\checkmarkc\checkmark$\checkmarkt\checkmark$\checkmark^\checkmark\\checkmarkc\checkmarki\checkmarkr\checkmarkc\checkmark$\checkmarka\checkmark$\checkmark^\checkmark\\checkmarkc\checkmarki\checkmarkr\checkmarkc\checkmark$\checkmarkb\checkmark$\checkmark^\checkmark\\checkmarkc\checkmarki\checkmarkr\checkmarkc\checkmark$\checkmarkl\checkmark$\checkmark^\checkmark\\checkmarkc\checkmarki\checkmarkr\checkmarkc\checkmark$\checkmarke\checkmark$\checkmark^\checkmark\\checkmarkc\checkmarki\checkmarkr\checkmarkc\checkmark$\checkmark}\checkmark$\checkmark^\checkmark\\checkmarkc\checkmarki\checkmarkr\checkmarkc\checkmark$\checkmark[\checkmark$\checkmark^\checkmark\\checkmarkc\checkmarki\checkmarkr\checkmarkc\checkmark$\checkmarkh\checkmark$\checkmark^\checkmark\\checkmarkc\checkmarki\checkmarkr\checkmarkc\checkmark$\checkmarkt\checkmark$\checkmark^\checkmark\\checkmarkc\checkmarki\checkmarkr\checkmarkc\checkmark$\checkmarkb\checkmark$\checkmark^\checkmark\\checkmarkc\checkmarki\checkmarkr\checkmarkc\checkmark$\checkmarkp\checkmark$\checkmark^\checkmark\\checkmarkc\checkmarki\checkmarkr\checkmarkc\checkmark$\checkmark]\checkmark$\checkmark^\checkmark\\checkmarkc\checkmarki\checkmarkr\checkmarkc\checkmark$\checkmark
\checkmark$\checkmark^\checkmark\\checkmarkc\checkmarki\checkmarkr\checkmarkc\checkmark$\checkmark\\checkmark$\checkmark^\checkmark\\checkmarkc\checkmarki\checkmarkr\checkmarkc\checkmark$\checkmarkc\checkmark$\checkmark^\checkmark\\checkmarkc\checkmarki\checkmarkr\checkmarkc\checkmark$\checkmarke\checkmark$\checkmark^\checkmark\\checkmarkc\checkmarki\checkmarkr\checkmarkc\checkmark$\checkmarkn\checkmark$\checkmark^\checkmark\\checkmarkc\checkmarki\checkmarkr\checkmarkc\checkmark$\checkmarkt\checkmark$\checkmark^\checkmark\\checkmarkc\checkmarki\checkmarkr\checkmarkc\checkmark$\checkmarke\checkmark$\checkmark^\checkmark\\checkmarkc\checkmarki\checkmarkr\checkmarkc\checkmark$\checkmarkr\checkmark$\checkmark^\checkmark\\checkmarkc\checkmarki\checkmarkr\checkmarkc\checkmark$\checkmarki\checkmark$\checkmark^\checkmark\\checkmarkc\checkmarki\checkmarkr\checkmarkc\checkmark$\checkmarkn\checkmark$\checkmark^\checkmark\\checkmarkc\checkmarki\checkmarkr\checkmarkc\checkmark$\checkmarkg\checkmark$\checkmark^\checkmark\\checkmarkc\checkmarki\checkmarkr\checkmarkc\checkmark$\checkmark
\checkmark$\checkmark^\checkmark\\checkmarkc\checkmarki\checkmarkr\checkmarkc\checkmark$\checkmark\\checkmark$\checkmark^\checkmark\\checkmarkc\checkmarki\checkmarkr\checkmarkc\checkmark$\checkmarkb\checkmark$\checkmark^\checkmark\\checkmarkc\checkmarki\checkmarkr\checkmarkc\checkmark$\checkmarke\checkmark$\checkmark^\checkmark\\checkmarkc\checkmarki\checkmarkr\checkmarkc\checkmark$\checkmarkg\checkmark$\checkmark^\checkmark\\checkmarkc\checkmarki\checkmarkr\checkmarkc\checkmark$\checkmarki\checkmark$\checkmark^\checkmark\\checkmarkc\checkmarki\checkmarkr\checkmarkc\checkmark$\checkmarkn\checkmark$\checkmark^\checkmark\\checkmarkc\checkmarki\checkmarkr\checkmarkc\checkmark$\checkmark{\checkmark$\checkmark^\checkmark\\checkmarkc\checkmarki\checkmarkr\checkmarkc\checkmark$\checkmarkt\checkmark$\checkmark^\checkmark\\checkmarkc\checkmarki\checkmarkr\checkmarkc\checkmark$\checkmarka\checkmark$\checkmark^\checkmark\\checkmarkc\checkmarki\checkmarkr\checkmarkc\checkmark$\checkmarkb\checkmark$\checkmark^\checkmark\\checkmarkc\checkmarki\checkmarkr\checkmarkc\checkmark$\checkmarku\checkmark$\checkmark^\checkmark\\checkmarkc\checkmarki\checkmarkr\checkmarkc\checkmark$\checkmarkl\checkmark$\checkmark^\checkmark\\checkmarkc\checkmarki\checkmarkr\checkmarkc\checkmark$\checkmarka\checkmark$\checkmark^\checkmark\\checkmarkc\checkmarki\checkmarkr\checkmarkc\checkmark$\checkmarkr\checkmark$\checkmark^\checkmark\\checkmarkc\checkmarki\checkmarkr\checkmarkc\checkmark$\checkmark}\checkmark$\checkmark^\checkmark\\checkmarkc\checkmarki\checkmarkr\checkmarkc\checkmark$\checkmark{\checkmark$\checkmark^\checkmark\\checkmarkc\checkmarki\checkmarkr\checkmarkc\checkmark$\checkmark|\checkmark$\checkmark^\checkmark\\checkmarkc\checkmarki\checkmarkr\checkmarkc\checkmark$\checkmarkl\checkmark$\checkmark^\checkmark\\checkmarkc\checkmarki\checkmarkr\checkmarkc\checkmark$\checkmark|\checkmark$\checkmark^\checkmark\\checkmarkc\checkmarki\checkmarkr\checkmarkc\checkmark$\checkmarkc\checkmark$\checkmark^\checkmark\\checkmarkc\checkmarki\checkmarkr\checkmarkc\checkmark$\checkmark|\checkmark$\checkmark^\checkmark\\checkmarkc\checkmarki\checkmarkr\checkmarkc\checkmark$\checkmarkc\checkmark$\checkmark^\checkmark\\checkmarkc\checkmarki\checkmarkr\checkmarkc\checkmark$\checkmark|\checkmark$\checkmark^\checkmark\\checkmarkc\checkmarki\checkmarkr\checkmarkc\checkmark$\checkmarkc\checkmark$\checkmark^\checkmark\\checkmarkc\checkmarki\checkmarkr\checkmarkc\checkmark$\checkmark|\checkmark$\checkmark^\checkmark\\checkmarkc\checkmarki\checkmarkr\checkmarkc\checkmark$\checkmarkc\checkmark$\checkmark^\checkmark\\checkmarkc\checkmarki\checkmarkr\checkmarkc\checkmark$\checkmark|\checkmark$\checkmark^\checkmark\\checkmarkc\checkmarki\checkmarkr\checkmarkc\checkmark$\checkmark}\checkmark$\checkmark^\checkmark\\checkmarkc\checkmarki\checkmarkr\checkmarkc\checkmark$\checkmark
\checkmark$\checkmark^\checkmark\\checkmarkc\checkmarki\checkmarkr\checkmarkc\checkmark$\checkmark\\checkmark$\checkmark^\checkmark\\checkmarkc\checkmarki\checkmarkr\checkmarkc\checkmark$\checkmarkh\checkmark$\checkmark^\checkmark\\checkmarkc\checkmarki\checkmarkr\checkmarkc\checkmark$\checkmarkl\checkmark$\checkmark^\checkmark\\checkmarkc\checkmarki\checkmarkr\checkmarkc\checkmark$\checkmarki\checkmark$\checkmark^\checkmark\\checkmarkc\checkmarki\checkmarkr\checkmarkc\checkmark$\checkmarkn\checkmark$\checkmark^\checkmark\\checkmarkc\checkmarki\checkmarkr\checkmarkc\checkmark$\checkmarke\checkmark$\checkmark^\checkmark\\checkmarkc\checkmarki\checkmarkr\checkmarkc\checkmark$\checkmark
\checkmark$\checkmark^\checkmark\\checkmarkc\checkmarki\checkmarkr\checkmarkc\checkmark$\checkmark\\checkmark$\checkmark^\checkmark\\checkmarkc\checkmarki\checkmarkr\checkmarkc\checkmark$\checkmarkt\checkmark$\checkmark^\checkmark\\checkmarkc\checkmarki\checkmarkr\checkmarkc\checkmark$\checkmarke\checkmark$\checkmark^\checkmark\\checkmarkc\checkmarki\checkmarkr\checkmarkc\checkmark$\checkmarkx\checkmark$\checkmark^\checkmark\\checkmarkc\checkmarki\checkmarkr\checkmarkc\checkmark$\checkmarkt\checkmark$\checkmark^\checkmark\\checkmarkc\checkmarki\checkmarkr\checkmarkc\checkmark$\checkmarkb\checkmark$\checkmark^\checkmark\\checkmarkc\checkmarki\checkmarkr\checkmarkc\checkmark$\checkmarkf\checkmark$\checkmark^\checkmark\\checkmarkc\checkmarki\checkmarkr\checkmarkc\checkmark$\checkmark{\checkmark$\checkmark^\checkmark\\checkmarkc\checkmarki\checkmarkr\checkmarkc\checkmark$\checkmarkM\checkmark$\checkmark^\checkmark\\checkmarkc\checkmarki\checkmarkr\checkmarkc\checkmark$\checkmarka\checkmark$\checkmark^\checkmark\\checkmarkc\checkmarki\checkmarkr\checkmarkc\checkmark$\checkmarkt\checkmark$\checkmark^\checkmark\\checkmarkc\checkmarki\checkmarkr\checkmarkc\checkmark$\checkmarké\checkmark$\checkmark^\checkmark\\checkmarkc\checkmarki\checkmarkr\checkmarkc\checkmark$\checkmarkr\checkmark$\checkmark^\checkmark\\checkmarkc\checkmarki\checkmarkr\checkmarkc\checkmark$\checkmarki\checkmark$\checkmark^\checkmark\\checkmarkc\checkmarki\checkmarkr\checkmarkc\checkmark$\checkmarka\checkmark$\checkmark^\checkmark\\checkmarkc\checkmarki\checkmarkr\checkmarkc\checkmark$\checkmarku\checkmark$\checkmark^\checkmark\\checkmarkc\checkmarki\checkmarkr\checkmarkc\checkmark$\checkmark}\checkmark$\checkmark^\checkmark\\checkmarkc\checkmarki\checkmarkr\checkmarkc\checkmark$\checkmark \checkmark$\checkmark^\checkmark\\checkmarkc\checkmarki\checkmarkr\checkmarkc\checkmark$\checkmark&\checkmark$\checkmark^\checkmark\\checkmarkc\checkmarki\checkmarkr\checkmarkc\checkmark$\checkmark \checkmark$\checkmark^\checkmark\\checkmarkc\checkmarki\checkmarkr\checkmarkc\checkmark$\checkmark$\checkmark$\checkmark^\checkmark\\checkmarkc\checkmarki\checkmarkr\checkmarkc\checkmark$\checkmarkR\checkmark$\checkmark^\checkmark\\checkmarkc\checkmarki\checkmarkr\checkmarkc\checkmark$\checkmark_\checkmark$\checkmark^\checkmark\\checkmarkc\checkmarki\checkmarkr\checkmarkc\checkmark$\checkmarke\checkmark$\checkmark^\checkmark\\checkmarkc\checkmarki\checkmarkr\checkmarkc\checkmark$\checkmark$\checkmark$\checkmark^\checkmark\\checkmarkc\checkmarki\checkmarkr\checkmarkc\checkmark$\checkmark \checkmark$\checkmark^\checkmark\\checkmarkc\checkmarki\checkmarkr\checkmarkc\checkmark$\checkmark(\checkmark$\checkmark^\checkmark\\checkmarkc\checkmarki\checkmarkr\checkmarkc\checkmark$\checkmarkM\checkmark$\checkmark^\checkmark\\checkmarkc\checkmarki\checkmarkr\checkmarkc\checkmark$\checkmarkP\checkmark$\checkmark^\checkmark\\checkmarkc\checkmarki\checkmarkr\checkmarkc\checkmark$\checkmarka\checkmark$\checkmark^\checkmark\\checkmarkc\checkmarki\checkmarkr\checkmarkc\checkmark$\checkmark)\checkmark$\checkmark^\checkmark\\checkmarkc\checkmarki\checkmarkr\checkmarkc\checkmark$\checkmark \checkmark$\checkmark^\checkmark\\checkmarkc\checkmarki\checkmarkr\checkmarkc\checkmark$\checkmark&\checkmark$\checkmark^\checkmark\\checkmarkc\checkmarki\checkmarkr\checkmarkc\checkmark$\checkmark \checkmark$\checkmark^\checkmark\\checkmarkc\checkmarki\checkmarkr\checkmarkc\checkmark$\checkmark$\checkmark$\checkmark^\checkmark\\checkmarkc\checkmarki\checkmarkr\checkmarkc\checkmark$\checkmarkR\checkmark$\checkmark^\checkmark\\checkmarkc\checkmarki\checkmarkr\checkmarkc\checkmark$\checkmark_\checkmark$\checkmark^\checkmark\\checkmarkc\checkmarki\checkmarkr\checkmarkc\checkmark$\checkmarkm\checkmark$\checkmark^\checkmark\\checkmarkc\checkmarki\checkmarkr\checkmarkc\checkmark$\checkmark$\checkmark$\checkmark^\checkmark\\checkmarkc\checkmarki\checkmarkr\checkmarkc\checkmark$\checkmark \checkmark$\checkmark^\checkmark\\checkmarkc\checkmarki\checkmarkr\checkmarkc\checkmark$\checkmark(\checkmark$\checkmark^\checkmark\\checkmarkc\checkmarki\checkmarkr\checkmarkc\checkmark$\checkmarkM\checkmark$\checkmark^\checkmark\\checkmarkc\checkmarki\checkmarkr\checkmarkc\checkmark$\checkmarkP\checkmark$\checkmark^\checkmark\\checkmarkc\checkmarki\checkmarkr\checkmarkc\checkmark$\checkmarka\checkmark$\checkmark^\checkmark\\checkmarkc\checkmarki\checkmarkr\checkmarkc\checkmark$\checkmark)\checkmark$\checkmark^\checkmark\\checkmarkc\checkmarki\checkmarkr\checkmarkc\checkmark$\checkmark \checkmark$\checkmark^\checkmark\\checkmarkc\checkmarki\checkmarkr\checkmarkc\checkmark$\checkmark&\checkmark$\checkmark^\checkmark\\checkmarkc\checkmarki\checkmarkr\checkmarkc\checkmark$\checkmark \checkmark$\checkmark^\checkmark\\checkmarkc\checkmarki\checkmarkr\checkmarkc\checkmark$\checkmark$\checkmark$\checkmark^\checkmark\\checkmarkc\checkmarki\checkmarkr\checkmarkc\checkmark$\checkmarkE\checkmark$\checkmark^\checkmark\\checkmarkc\checkmarki\checkmarkr\checkmarkc\checkmark$\checkmark$\checkmark$\checkmark^\checkmark\\checkmarkc\checkmarki\checkmarkr\checkmarkc\checkmark$\checkmark \checkmark$\checkmark^\checkmark\\checkmarkc\checkmarki\checkmarkr\checkmarkc\checkmark$\checkmark(\checkmark$\checkmark^\checkmark\\checkmarkc\checkmarki\checkmarkr\checkmarkc\checkmark$\checkmarkG\checkmark$\checkmark^\checkmark\\checkmarkc\checkmarki\checkmarkr\checkmarkc\checkmark$\checkmarkP\checkmark$\checkmark^\checkmark\\checkmarkc\checkmarki\checkmarkr\checkmarkc\checkmark$\checkmarka\checkmark$\checkmark^\checkmark\\checkmarkc\checkmarki\checkmarkr\checkmarkc\checkmark$\checkmark)\checkmark$\checkmark^\checkmark\\checkmarkc\checkmarki\checkmarkr\checkmarkc\checkmark$\checkmark \checkmark$\checkmark^\checkmark\\checkmarkc\checkmarki\checkmarkr\checkmarkc\checkmark$\checkmark&\checkmark$\checkmark^\checkmark\\checkmarkc\checkmarki\checkmarkr\checkmarkc\checkmark$\checkmark \checkmark$\checkmark^\checkmark\\checkmarkc\checkmarki\checkmarkr\checkmarkc\checkmark$\checkmark\\checkmark$\checkmark^\checkmark\\checkmarkc\checkmarki\checkmarkr\checkmarkc\checkmark$\checkmarkt\checkmark$\checkmark^\checkmark\\checkmarkc\checkmarki\checkmarkr\checkmarkc\checkmark$\checkmarke\checkmark$\checkmark^\checkmark\\checkmarkc\checkmarki\checkmarkr\checkmarkc\checkmark$\checkmarkx\checkmark$\checkmark^\checkmark\\checkmarkc\checkmarki\checkmarkr\checkmarkc\checkmark$\checkmarkt\checkmark$\checkmark^\checkmark\\checkmarkc\checkmarki\checkmarkr\checkmarkc\checkmark$\checkmarkb\checkmark$\checkmark^\checkmark\\checkmarkc\checkmarki\checkmarkr\checkmarkc\checkmark$\checkmarkf\checkmark$\checkmark^\checkmark\\checkmarkc\checkmarki\checkmarkr\checkmarkc\checkmark$\checkmark{\checkmark$\checkmark^\checkmark\\checkmarkc\checkmarki\checkmarkr\checkmarkc\checkmark$\checkmarkD\checkmark$\checkmark^\checkmark\\checkmarkc\checkmarki\checkmarkr\checkmarkc\checkmark$\checkmarké\checkmark$\checkmark^\checkmark\\checkmarkc\checkmarki\checkmarkr\checkmarkc\checkmark$\checkmarkc\checkmark$\checkmark^\checkmark\\checkmarkc\checkmarki\checkmarkr\checkmarkc\checkmark$\checkmarki\checkmark$\checkmark^\checkmark\\checkmarkc\checkmarki\checkmarkr\checkmarkc\checkmark$\checkmarks\checkmark$\checkmark^\checkmark\\checkmarkc\checkmarki\checkmarkr\checkmarkc\checkmark$\checkmarki\checkmark$\checkmark^\checkmark\\checkmarkc\checkmarki\checkmarkr\checkmarkc\checkmark$\checkmarko\checkmark$\checkmark^\checkmark\\checkmarkc\checkmarki\checkmarkr\checkmarkc\checkmark$\checkmarkn\checkmark$\checkmark^\checkmark\\checkmarkc\checkmarki\checkmarkr\checkmarkc\checkmark$\checkmark}\checkmark$\checkmark^\checkmark\\checkmarkc\checkmarki\checkmarkr\checkmarkc\checkmark$\checkmark \checkmark$\checkmark^\checkmark\\checkmarkc\checkmarki\checkmarkr\checkmarkc\checkmark$\checkmark\\checkmark$\checkmark^\checkmark\\checkmarkc\checkmarki\checkmarkr\checkmarkc\checkmark$\checkmark\\checkmark$\checkmark^\checkmark\\checkmarkc\checkmarki\checkmarkr\checkmarkc\checkmark$\checkmark \checkmark$\checkmark^\checkmark\\checkmarkc\checkmarki\checkmarkr\checkmarkc\checkmark$\checkmark\\checkmark$\checkmark^\checkmark\\checkmarkc\checkmarki\checkmarkr\checkmarkc\checkmark$\checkmarkh\checkmark$\checkmark^\checkmark\\checkmarkc\checkmarki\checkmarkr\checkmarkc\checkmark$\checkmarkl\checkmark$\checkmark^\checkmark\\checkmarkc\checkmarki\checkmarkr\checkmarkc\checkmark$\checkmarki\checkmark$\checkmark^\checkmark\\checkmarkc\checkmarki\checkmarkr\checkmarkc\checkmark$\checkmarkn\checkmark$\checkmark^\checkmark\\checkmarkc\checkmarki\checkmarkr\checkmarkc\checkmark$\checkmarke\checkmark$\checkmark^\checkmark\\checkmarkc\checkmarki\checkmarkr\checkmarkc\checkmark$\checkmark
\checkmark$\checkmark^\checkmark\\checkmarkc\checkmarki\checkmarkr\checkmarkc\checkmark$\checkmarkA\checkmark$\checkmark^\checkmark\\checkmarkc\checkmarki\checkmarkr\checkmarkc\checkmark$\checkmarkl\checkmark$\checkmark^\checkmark\\checkmarkc\checkmarki\checkmarkr\checkmarkc\checkmark$\checkmarku\checkmark$\checkmark^\checkmark\\checkmarkc\checkmarki\checkmarkr\checkmarkc\checkmark$\checkmarkm\checkmark$\checkmark^\checkmark\\checkmarkc\checkmarki\checkmarkr\checkmarkc\checkmark$\checkmarki\checkmark$\checkmark^\checkmark\\checkmarkc\checkmarki\checkmarkr\checkmarkc\checkmark$\checkmarkn\checkmark$\checkmark^\checkmark\\checkmarkc\checkmarki\checkmarkr\checkmarkc\checkmark$\checkmarki\checkmark$\checkmark^\checkmark\\checkmarkc\checkmarki\checkmarkr\checkmarkc\checkmark$\checkmarku\checkmark$\checkmark^\checkmark\\checkmarkc\checkmarki\checkmarkr\checkmarkc\checkmark$\checkmarkm\checkmark$\checkmark^\checkmark\\checkmarkc\checkmarki\checkmarkr\checkmarkc\checkmark$\checkmark \checkmark$\checkmark^\checkmark\\checkmarkc\checkmarki\checkmarkr\checkmarkc\checkmark$\checkmark6\checkmark$\checkmark^\checkmark\\checkmarkc\checkmarki\checkmarkr\checkmarkc\checkmark$\checkmark0\checkmark$\checkmark^\checkmark\\checkmarkc\checkmarki\checkmarkr\checkmarkc\checkmark$\checkmark6\checkmark$\checkmark^\checkmark\\checkmarkc\checkmarki\checkmarkr\checkmarkc\checkmark$\checkmark1\checkmark$\checkmark^\checkmark\\checkmarkc\checkmarki\checkmarkr\checkmarkc\checkmark$\checkmark-\checkmark$\checkmark^\checkmark\\checkmarkc\checkmarki\checkmarkr\checkmarkc\checkmark$\checkmarkT\checkmark$\checkmark^\checkmark\\checkmarkc\checkmarki\checkmarkr\checkmarkc\checkmark$\checkmark6\checkmark$\checkmark^\checkmark\\checkmarkc\checkmarki\checkmarkr\checkmarkc\checkmark$\checkmark \checkmark$\checkmark^\checkmark\\checkmarkc\checkmarki\checkmarkr\checkmarkc\checkmark$\checkmark&\checkmark$\checkmark^\checkmark\\checkmarkc\checkmarki\checkmarkr\checkmarkc\checkmark$\checkmark \checkmark$\checkmark^\checkmark\\checkmarkc\checkmarki\checkmarkr\checkmarkc\checkmark$\checkmark2\checkmark$\checkmark^\checkmark\\checkmarkc\checkmarki\checkmarkr\checkmarkc\checkmark$\checkmark7\checkmark$\checkmark^\checkmark\\checkmarkc\checkmarki\checkmarkr\checkmarkc\checkmark$\checkmark6\checkmark$\checkmark^\checkmark\\checkmarkc\checkmarki\checkmarkr\checkmarkc\checkmark$\checkmark \checkmark$\checkmark^\checkmark\\checkmarkc\checkmarki\checkmarkr\checkmarkc\checkmark$\checkmark&\checkmark$\checkmark^\checkmark\\checkmarkc\checkmarki\checkmarkr\checkmarkc\checkmark$\checkmark \checkmark$\checkmark^\checkmark\\checkmarkc\checkmarki\checkmarkr\checkmarkc\checkmark$\checkmark3\checkmark$\checkmark^\checkmark\\checkmarkc\checkmarki\checkmarkr\checkmarkc\checkmark$\checkmark1\checkmark$\checkmark^\checkmark\\checkmarkc\checkmarki\checkmarkr\checkmarkc\checkmark$\checkmark0\checkmark$\checkmark^\checkmark\\checkmarkc\checkmarki\checkmarkr\checkmarkc\checkmark$\checkmark \checkmark$\checkmark^\checkmark\\checkmarkc\checkmarki\checkmarkr\checkmarkc\checkmark$\checkmark&\checkmark$\checkmark^\checkmark\\checkmarkc\checkmarki\checkmarkr\checkmarkc\checkmark$\checkmark \checkmark$\checkmark^\checkmark\\checkmarkc\checkmarki\checkmarkr\checkmarkc\checkmark$\checkmark6\checkmark$\checkmark^\checkmark\\checkmarkc\checkmarki\checkmarkr\checkmarkc\checkmark$\checkmark9\checkmark$\checkmark^\checkmark\\checkmarkc\checkmarki\checkmarkr\checkmarkc\checkmark$\checkmark \checkmark$\checkmark^\checkmark\\checkmarkc\checkmarki\checkmarkr\checkmarkc\checkmark$\checkmark&\checkmark$\checkmark^\checkmark\\checkmarkc\checkmarki\checkmarkr\checkmarkc\checkmark$\checkmark \checkmark$\checkmark^\checkmark\\checkmarkc\checkmarki\checkmarkr\checkmarkc\checkmark$\checkmarkR\checkmark$\checkmark^\checkmark\\checkmarkc\checkmarki\checkmarkr\checkmarkc\checkmark$\checkmarke\checkmark$\checkmark^\checkmark\\checkmarkc\checkmarki\checkmarkr\checkmarkc\checkmark$\checkmarkj\checkmark$\checkmark^\checkmark\\checkmarkc\checkmarki\checkmarkr\checkmarkc\checkmark$\checkmarke\checkmark$\checkmark^\checkmark\\checkmarkc\checkmarki\checkmarkr\checkmarkc\checkmark$\checkmarkt\checkmark$\checkmark^\checkmark\\checkmarkc\checkmarki\checkmarkr\checkmarkc\checkmark$\checkmarké\checkmark$\checkmark^\checkmark\\checkmarkc\checkmarki\checkmarkr\checkmarkc\checkmark$\checkmark \checkmark$\checkmark^\checkmark\\checkmarkc\checkmarki\checkmarkr\checkmarkc\checkmark$\checkmark(\checkmark$\checkmark^\checkmark\\checkmarkc\checkmarki\checkmarkr\checkmarkc\checkmark$\checkmarkt\checkmark$\checkmark^\checkmark\\checkmarkc\checkmarki\checkmarkr\checkmarkc\checkmark$\checkmarkr\checkmark$\checkmark^\checkmark\\checkmarkc\checkmarki\checkmarkr\checkmarkc\checkmark$\checkmarko\checkmark$\checkmark^\checkmark\\checkmarkc\checkmarki\checkmarkr\checkmarkc\checkmark$\checkmarkp\checkmark$\checkmark^\checkmark\\checkmarkc\checkmarki\checkmarkr\checkmarkc\checkmark$\checkmark \checkmark$\checkmark^\checkmark\\checkmarkc\checkmarki\checkmarkr\checkmarkc\checkmark$\checkmarkm\checkmark$\checkmark^\checkmark\\checkmarkc\checkmarki\checkmarkr\checkmarkc\checkmark$\checkmarko\checkmark$\checkmark^\checkmark\\checkmarkc\checkmarki\checkmarkr\checkmarkc\checkmark$\checkmarku\checkmark$\checkmark^\checkmark\\checkmarkc\checkmarki\checkmarkr\checkmarkc\checkmark$\checkmark)\checkmark$\checkmark^\checkmark\\checkmarkc\checkmarki\checkmarkr\checkmarkc\checkmark$\checkmark \checkmark$\checkmark^\checkmark\\checkmarkc\checkmarki\checkmarkr\checkmarkc\checkmark$\checkmark\\checkmark$\checkmark^\checkmark\\checkmarkc\checkmarki\checkmarkr\checkmarkc\checkmark$\checkmark\\checkmark$\checkmark^\checkmark\\checkmarkc\checkmarki\checkmarkr\checkmarkc\checkmark$\checkmark \checkmark$\checkmark^\checkmark\\checkmarkc\checkmarki\checkmarkr\checkmarkc\checkmark$\checkmark\\checkmark$\checkmark^\checkmark\\checkmarkc\checkmarki\checkmarkr\checkmarkc\checkmark$\checkmarkh\checkmark$\checkmark^\checkmark\\checkmarkc\checkmarki\checkmarkr\checkmarkc\checkmark$\checkmarkl\checkmark$\checkmark^\checkmark\\checkmarkc\checkmarki\checkmarkr\checkmarkc\checkmark$\checkmarki\checkmark$\checkmark^\checkmark\\checkmarkc\checkmarki\checkmarkr\checkmarkc\checkmark$\checkmarkn\checkmark$\checkmark^\checkmark\\checkmarkc\checkmarki\checkmarkr\checkmarkc\checkmark$\checkmarke\checkmark$\checkmark^\checkmark\\checkmarkc\checkmarki\checkmarkr\checkmarkc\checkmark$\checkmark
\checkmark$\checkmark^\checkmark\\checkmarkc\checkmarki\checkmarkr\checkmarkc\checkmark$\checkmarkA\checkmark$\checkmark^\checkmark\\checkmarkc\checkmarki\checkmarkr\checkmarkc\checkmark$\checkmarkc\checkmark$\checkmark^\checkmark\\checkmarkc\checkmarki\checkmarkr\checkmarkc\checkmark$\checkmarki\checkmark$\checkmark^\checkmark\\checkmarkc\checkmarki\checkmarkr\checkmarkc\checkmark$\checkmarke\checkmark$\checkmark^\checkmark\\checkmarkc\checkmarki\checkmarkr\checkmarkc\checkmark$\checkmarkr\checkmark$\checkmark^\checkmark\\checkmarkc\checkmarki\checkmarkr\checkmarkc\checkmark$\checkmark \checkmark$\checkmark^\checkmark\\checkmarkc\checkmarki\checkmarkr\checkmarkc\checkmark$\checkmarkX\checkmark$\checkmark^\checkmark\\checkmarkc\checkmarki\checkmarkr\checkmarkc\checkmark$\checkmarkC\checkmark$\checkmark^\checkmark\\checkmarkc\checkmarki\checkmarkr\checkmarkc\checkmark$\checkmark3\checkmark$\checkmark^\checkmark\\checkmarkc\checkmarki\checkmarkr\checkmarkc\checkmark$\checkmark8\checkmark$\checkmark^\checkmark\\checkmarkc\checkmarki\checkmarkr\checkmarkc\checkmark$\checkmark \checkmark$\checkmark^\checkmark\\checkmarkc\checkmarki\checkmarkr\checkmarkc\checkmark$\checkmark&\checkmark$\checkmark^\checkmark\\checkmarkc\checkmarki\checkmarkr\checkmarkc\checkmark$\checkmark \checkmark$\checkmark^\checkmark\\checkmarkc\checkmarki\checkmarkr\checkmarkc\checkmark$\checkmark2\checkmark$\checkmark^\checkmark\\checkmarkc\checkmarki\checkmarkr\checkmarkc\checkmark$\checkmark3\checkmark$\checkmark^\checkmark\\checkmarkc\checkmarki\checkmarkr\checkmarkc\checkmark$\checkmark5\checkmark$\checkmark^\checkmark\\checkmarkc\checkmarki\checkmarkr\checkmarkc\checkmark$\checkmark \checkmark$\checkmark^\checkmark\\checkmarkc\checkmarki\checkmarkr\checkmarkc\checkmark$\checkmark&\checkmark$\checkmark^\checkmark\\checkmarkc\checkmarki\checkmarkr\checkmarkc\checkmark$\checkmark \checkmark$\checkmark^\checkmark\\checkmarkc\checkmarki\checkmarkr\checkmarkc\checkmark$\checkmark4\checkmark$\checkmark^\checkmark\\checkmarkc\checkmarki\checkmarkr\checkmarkc\checkmark$\checkmark6\checkmark$\checkmark^\checkmark\\checkmarkc\checkmarki\checkmarkr\checkmarkc\checkmark$\checkmark0\checkmark$\checkmark^\checkmark\\checkmarkc\checkmarki\checkmarkr\checkmarkc\checkmark$\checkmark \checkmark$\checkmark^\checkmark\\checkmarkc\checkmarki\checkmarkr\checkmarkc\checkmark$\checkmark&\checkmark$\checkmark^\checkmark\\checkmarkc\checkmarki\checkmarkr\checkmarkc\checkmark$\checkmark \checkmark$\checkmark^\checkmark\\checkmarkc\checkmarki\checkmarkr\checkmarkc\checkmark$\checkmark2\checkmark$\checkmark^\checkmark\\checkmarkc\checkmarki\checkmarkr\checkmarkc\checkmark$\checkmark1\checkmark$\checkmark^\checkmark\\checkmarkc\checkmarki\checkmarkr\checkmarkc\checkmark$\checkmark0\checkmark$\checkmark^\checkmark\\checkmarkc\checkmarki\checkmarkr\checkmarkc\checkmark$\checkmark \checkmark$\checkmark^\checkmark\\checkmarkc\checkmarki\checkmarkr\checkmarkc\checkmark$\checkmark&\checkmark$\checkmark^\checkmark\\checkmarkc\checkmarki\checkmarkr\checkmarkc\checkmark$\checkmark \checkmark$\checkmark^\checkmark\\checkmarkc\checkmarki\checkmarkr\checkmarkc\checkmark$\checkmarkR\checkmark$\checkmark^\checkmark\\checkmarkc\checkmarki\checkmarkr\checkmarkc\checkmark$\checkmarke\checkmark$\checkmark^\checkmark\\checkmarkc\checkmarki\checkmarkr\checkmarkc\checkmark$\checkmarkj\checkmark$\checkmark^\checkmark\\checkmarkc\checkmarki\checkmarkr\checkmarkc\checkmark$\checkmarke\checkmark$\checkmark^\checkmark\\checkmarkc\checkmarki\checkmarkr\checkmarkc\checkmark$\checkmarkt\checkmark$\checkmark^\checkmark\\checkmarkc\checkmarki\checkmarkr\checkmarkc\checkmark$\checkmarké\checkmark$\checkmark^\checkmark\\checkmarkc\checkmarki\checkmarkr\checkmarkc\checkmark$\checkmark \checkmark$\checkmark^\checkmark\\checkmarkc\checkmarki\checkmarkr\checkmarkc\checkmark$\checkmark(\checkmark$\checkmark^\checkmark\\checkmarkc\checkmarki\checkmarkr\checkmarkc\checkmark$\checkmarki\checkmark$\checkmark^\checkmark\\checkmarkc\checkmarki\checkmarkr\checkmarkc\checkmark$\checkmarkn\checkmark$\checkmark^\checkmark\\checkmarkc\checkmarki\checkmarkr\checkmarkc\checkmark$\checkmarks\checkmark$\checkmark^\checkmark\\checkmarkc\checkmarki\checkmarkr\checkmarkc\checkmark$\checkmarku\checkmark$\checkmark^\checkmark\\checkmarkc\checkmarki\checkmarkr\checkmarkc\checkmark$\checkmarkf\checkmark$\checkmark^\checkmark\\checkmarkc\checkmarki\checkmarkr\checkmarkc\checkmark$\checkmarkf\checkmark$\checkmark^\checkmark\\checkmarkc\checkmarki\checkmarkr\checkmarkc\checkmark$\checkmarki\checkmark$\checkmark^\checkmark\\checkmarkc\checkmarki\checkmarkr\checkmarkc\checkmark$\checkmarks\checkmark$\checkmark^\checkmark\\checkmarkc\checkmarki\checkmarkr\checkmarkc\checkmark$\checkmarka\checkmark$\checkmark^\checkmark\\checkmarkc\checkmarki\checkmarkr\checkmarkc\checkmark$\checkmarkn\checkmark$\checkmark^\checkmark\\checkmarkc\checkmarki\checkmarkr\checkmarkc\checkmark$\checkmarkt\checkmark$\checkmark^\checkmark\\checkmarkc\checkmarki\checkmarkr\checkmarkc\checkmark$\checkmark)\checkmark$\checkmark^\checkmark\\checkmarkc\checkmarki\checkmarkr\checkmarkc\checkmark$\checkmark \checkmark$\checkmark^\checkmark\\checkmarkc\checkmarki\checkmarkr\checkmarkc\checkmark$\checkmark\\checkmark$\checkmark^\checkmark\\checkmarkc\checkmarki\checkmarkr\checkmarkc\checkmark$\checkmark\\checkmark$\checkmark^\checkmark\\checkmarkc\checkmarki\checkmarkr\checkmarkc\checkmark$\checkmark \checkmark$\checkmark^\checkmark\\checkmarkc\checkmarki\checkmarkr\checkmarkc\checkmark$\checkmark\\checkmark$\checkmark^\checkmark\\checkmarkc\checkmarki\checkmarkr\checkmarkc\checkmark$\checkmarkh\checkmark$\checkmark^\checkmark\\checkmarkc\checkmarki\checkmarkr\checkmarkc\checkmark$\checkmarkl\checkmark$\checkmark^\checkmark\\checkmarkc\checkmarki\checkmarkr\checkmarkc\checkmark$\checkmarki\checkmark$\checkmark^\checkmark\\checkmarkc\checkmarki\checkmarkr\checkmarkc\checkmark$\checkmarkn\checkmark$\checkmark^\checkmark\\checkmarkc\checkmarki\checkmarkr\checkmarkc\checkmark$\checkmarke\checkmark$\checkmark^\checkmark\\checkmarkc\checkmarki\checkmarkr\checkmarkc\checkmark$\checkmark
\checkmark$\checkmark^\checkmark\\checkmarkc\checkmarki\checkmarkr\checkmarkc\checkmark$\checkmark\\checkmark$\checkmark^\checkmark\\checkmarkc\checkmarki\checkmarkr\checkmarkc\checkmark$\checkmarkt\checkmark$\checkmark^\checkmark\\checkmarkc\checkmarki\checkmarkr\checkmarkc\checkmark$\checkmarke\checkmark$\checkmark^\checkmark\\checkmarkc\checkmarki\checkmarkr\checkmarkc\checkmark$\checkmarkx\checkmark$\checkmark^\checkmark\\checkmarkc\checkmarki\checkmarkr\checkmarkc\checkmark$\checkmarkt\checkmark$\checkmark^\checkmark\\checkmarkc\checkmarki\checkmarkr\checkmarkc\checkmark$\checkmarkb\checkmark$\checkmark^\checkmark\\checkmarkc\checkmarki\checkmarkr\checkmarkc\checkmark$\checkmarkf\checkmark$\checkmark^\checkmark\\checkmarkc\checkmarki\checkmarkr\checkmarkc\checkmark$\checkmark{\checkmark$\checkmark^\checkmark\\checkmarkc\checkmarki\checkmarkr\checkmarkc\checkmark$\checkmarkA\checkmark$\checkmark^\checkmark\\checkmarkc\checkmarki\checkmarkr\checkmarkc\checkmark$\checkmarkc\checkmark$\checkmark^\checkmark\\checkmarkc\checkmarki\checkmarkr\checkmarkc\checkmark$\checkmarki\checkmark$\checkmark^\checkmark\\checkmarkc\checkmarki\checkmarkr\checkmarkc\checkmark$\checkmarke\checkmark$\checkmark^\checkmark\\checkmarkc\checkmarki\checkmarkr\checkmarkc\checkmark$\checkmarkr\checkmark$\checkmark^\checkmark\\checkmarkc\checkmarki\checkmarkr\checkmarkc\checkmark$\checkmark \checkmark$\checkmark^\checkmark\\checkmarkc\checkmarki\checkmarkr\checkmarkc\checkmark$\checkmarkC\checkmark$\checkmark^\checkmark\\checkmarkc\checkmarki\checkmarkr\checkmarkc\checkmark$\checkmark4\checkmark$\checkmark^\checkmark\\checkmarkc\checkmarki\checkmarkr\checkmarkc\checkmark$\checkmark5\checkmark$\checkmark^\checkmark\\checkmarkc\checkmarki\checkmarkr\checkmarkc\checkmark$\checkmark \checkmark$\checkmark^\checkmark\\checkmarkc\checkmarki\checkmarkr\checkmarkc\checkmark$\checkmark(\checkmark$\checkmark^\checkmark\\checkmarkc\checkmarki\checkmarkr\checkmarkc\checkmark$\checkmark1\checkmark$\checkmark^\checkmark\\checkmarkc\checkmarki\checkmarkr\checkmarkc\checkmark$\checkmark.\checkmark$\checkmark^\checkmark\\checkmarkc\checkmarki\checkmarkr\checkmarkc\checkmark$\checkmark0\checkmark$\checkmark^\checkmark\\checkmarkc\checkmarki\checkmarkr\checkmarkc\checkmark$\checkmark5\checkmark$\checkmark^\checkmark\\checkmarkc\checkmarki\checkmarkr\checkmarkc\checkmark$\checkmark0\checkmark$\checkmark^\checkmark\\checkmarkc\checkmarki\checkmarkr\checkmarkc\checkmark$\checkmark3\checkmark$\checkmark^\checkmark\\checkmarkc\checkmarki\checkmarkr\checkmarkc\checkmark$\checkmark)\checkmark$\checkmark^\checkmark\\checkmarkc\checkmarki\checkmarkr\checkmarkc\checkmark$\checkmark}\checkmark$\checkmark^\checkmark\\checkmarkc\checkmarki\checkmarkr\checkmarkc\checkmark$\checkmark \checkmark$\checkmark^\checkmark\\checkmarkc\checkmarki\checkmarkr\checkmarkc\checkmark$\checkmark&\checkmark$\checkmark^\checkmark\\checkmarkc\checkmarki\checkmarkr\checkmarkc\checkmark$\checkmark \checkmark$\checkmark^\checkmark\\checkmarkc\checkmarki\checkmarkr\checkmarkc\checkmark$\checkmark\\checkmark$\checkmark^\checkmark\\checkmarkc\checkmarki\checkmarkr\checkmarkc\checkmark$\checkmarkt\checkmark$\checkmark^\checkmark\\checkmarkc\checkmarki\checkmarkr\checkmarkc\checkmark$\checkmarke\checkmark$\checkmark^\checkmark\\checkmarkc\checkmarki\checkmarkr\checkmarkc\checkmark$\checkmarkx\checkmark$\checkmark^\checkmark\\checkmarkc\checkmarki\checkmarkr\checkmarkc\checkmark$\checkmarkt\checkmark$\checkmark^\checkmark\\checkmarkc\checkmarki\checkmarkr\checkmarkc\checkmark$\checkmarkb\checkmark$\checkmark^\checkmark\\checkmarkc\checkmarki\checkmarkr\checkmarkc\checkmark$\checkmarkf\checkmark$\checkmark^\checkmark\\checkmarkc\checkmarki\checkmarkr\checkmarkc\checkmark$\checkmark{\checkmark$\checkmark^\checkmark\\checkmarkc\checkmarki\checkmarkr\checkmarkc\checkmark$\checkmark3\checkmark$\checkmark^\checkmark\\checkmarkc\checkmarki\checkmarkr\checkmarkc\checkmark$\checkmark5\checkmark$\checkmark^\checkmark\\checkmarkc\checkmarki\checkmarkr\checkmarkc\checkmark$\checkmark0\checkmark$\checkmark^\checkmark\\checkmarkc\checkmarki\checkmarkr\checkmarkc\checkmark$\checkmark}\checkmark$\checkmark^\checkmark\\checkmarkc\checkmarki\checkmarkr\checkmarkc\checkmark$\checkmark \checkmark$\checkmark^\checkmark\\checkmarkc\checkmarki\checkmarkr\checkmarkc\checkmark$\checkmark&\checkmark$\checkmark^\checkmark\\checkmarkc\checkmarki\checkmarkr\checkmarkc\checkmark$\checkmark \checkmark$\checkmark^\checkmark\\checkmarkc\checkmarki\checkmarkr\checkmarkc\checkmark$\checkmark\\checkmark$\checkmark^\checkmark\\checkmarkc\checkmarki\checkmarkr\checkmarkc\checkmark$\checkmarkt\checkmark$\checkmark^\checkmark\\checkmarkc\checkmarki\checkmarkr\checkmarkc\checkmark$\checkmarke\checkmark$\checkmark^\checkmark\\checkmarkc\checkmarki\checkmarkr\checkmarkc\checkmark$\checkmarkx\checkmark$\checkmark^\checkmark\\checkmarkc\checkmarki\checkmarkr\checkmarkc\checkmark$\checkmarkt\checkmark$\checkmark^\checkmark\\checkmarkc\checkmarki\checkmarkr\checkmarkc\checkmark$\checkmarkb\checkmark$\checkmark^\checkmark\\checkmarkc\checkmarki\checkmarkr\checkmarkc\checkmark$\checkmarkf\checkmark$\checkmark^\checkmark\\checkmarkc\checkmarki\checkmarkr\checkmarkc\checkmark$\checkmark{\checkmark$\checkmark^\checkmark\\checkmarkc\checkmarki\checkmarkr\checkmarkc\checkmark$\checkmark6\checkmark$\checkmark^\checkmark\\checkmarkc\checkmarki\checkmarkr\checkmarkc\checkmark$\checkmark0\checkmark$\checkmark^\checkmark\\checkmarkc\checkmarki\checkmarkr\checkmarkc\checkmark$\checkmark0\checkmark$\checkmark^\checkmark\\checkmarkc\checkmarki\checkmarkr\checkmarkc\checkmark$\checkmark}\checkmark$\checkmark^\checkmark\\checkmarkc\checkmarki\checkmarkr\checkmarkc\checkmark$\checkmark \checkmark$\checkmark^\checkmark\\checkmarkc\checkmarki\checkmarkr\checkmarkc\checkmark$\checkmark&\checkmark$\checkmark^\checkmark\\checkmarkc\checkmarki\checkmarkr\checkmarkc\checkmark$\checkmark \checkmark$\checkmark^\checkmark\\checkmarkc\checkmarki\checkmarkr\checkmarkc\checkmark$\checkmark\\checkmark$\checkmark^\checkmark\\checkmarkc\checkmarki\checkmarkr\checkmarkc\checkmark$\checkmarkt\checkmark$\checkmark^\checkmark\\checkmarkc\checkmarki\checkmarkr\checkmarkc\checkmark$\checkmarke\checkmark$\checkmark^\checkmark\\checkmarkc\checkmarki\checkmarkr\checkmarkc\checkmark$\checkmarkx\checkmark$\checkmark^\checkmark\\checkmarkc\checkmarki\checkmarkr\checkmarkc\checkmark$\checkmarkt\checkmark$\checkmark^\checkmark\\checkmarkc\checkmarki\checkmarkr\checkmarkc\checkmark$\checkmarkb\checkmark$\checkmark^\checkmark\\checkmarkc\checkmarki\checkmarkr\checkmarkc\checkmark$\checkmarkf\checkmark$\checkmark^\checkmark\\checkmarkc\checkmarki\checkmarkr\checkmarkc\checkmark$\checkmark{\checkmark$\checkmark^\checkmark\\checkmarkc\checkmarki\checkmarkr\checkmarkc\checkmark$\checkmark2\checkmark$\checkmark^\checkmark\\checkmarkc\checkmarki\checkmarkr\checkmarkc\checkmark$\checkmark1\checkmark$\checkmark^\checkmark\\checkmarkc\checkmarki\checkmarkr\checkmarkc\checkmark$\checkmark0\checkmark$\checkmark^\checkmark\\checkmarkc\checkmarki\checkmarkr\checkmarkc\checkmark$\checkmark}\checkmark$\checkmark^\checkmark\\checkmarkc\checkmarki\checkmarkr\checkmarkc\checkmark$\checkmark \checkmark$\checkmark^\checkmark\\checkmarkc\checkmarki\checkmarkr\checkmarkc\checkmark$\checkmark&\checkmark$\checkmark^\checkmark\\checkmarkc\checkmarki\checkmarkr\checkmarkc\checkmark$\checkmark \checkmark$\checkmark^\checkmark\\checkmarkc\checkmarki\checkmarkr\checkmarkc\checkmark$\checkmark\\checkmark$\checkmark^\checkmark\\checkmarkc\checkmarki\checkmarkr\checkmarkc\checkmark$\checkmarkt\checkmark$\checkmark^\checkmark\\checkmarkc\checkmarki\checkmarkr\checkmarkc\checkmark$\checkmarke\checkmark$\checkmark^\checkmark\\checkmarkc\checkmarki\checkmarkr\checkmarkc\checkmark$\checkmarkx\checkmark$\checkmark^\checkmark\\checkmarkc\checkmarki\checkmarkr\checkmarkc\checkmark$\checkmarkt\checkmark$\checkmark^\checkmark\\checkmarkc\checkmarki\checkmarkr\checkmarkc\checkmark$\checkmarkb\checkmark$\checkmark^\checkmark\\checkmarkc\checkmarki\checkmarkr\checkmarkc\checkmark$\checkmarkf\checkmark$\checkmark^\checkmark\\checkmarkc\checkmarki\checkmarkr\checkmarkc\checkmark$\checkmark{\checkmark$\checkmark^\checkmark\\checkmarkc\checkmarki\checkmarkr\checkmarkc\checkmark$\checkmarkR\checkmark$\checkmark^\checkmark\\checkmarkc\checkmarki\checkmarkr\checkmarkc\checkmark$\checkmarke\checkmark$\checkmark^\checkmark\\checkmarkc\checkmarki\checkmarkr\checkmarkc\checkmark$\checkmarkt\checkmark$\checkmark^\checkmark\\checkmarkc\checkmarki\checkmarkr\checkmarkc\checkmark$\checkmarke\checkmark$\checkmark^\checkmark\\checkmarkc\checkmarki\checkmarkr\checkmarkc\checkmark$\checkmarkn\checkmark$\checkmark^\checkmark\\checkmarkc\checkmarki\checkmarkr\checkmarkc\checkmark$\checkmarku\checkmark$\checkmark^\checkmark\\checkmarkc\checkmarki\checkmarkr\checkmarkc\checkmark$\checkmark}\checkmark$\checkmark^\checkmark\\checkmarkc\checkmarki\checkmarkr\checkmarkc\checkmark$\checkmark \checkmark$\checkmark^\checkmark\\checkmarkc\checkmarki\checkmarkr\checkmarkc\checkmark$\checkmark\\checkmark$\checkmark^\checkmark\\checkmarkc\checkmarki\checkmarkr\checkmarkc\checkmark$\checkmark\\checkmark$\checkmark^\checkmark\\checkmarkc\checkmarki\checkmarkr\checkmarkc\checkmark$\checkmark \checkmark$\checkmark^\checkmark\\checkmarkc\checkmarki\checkmarkr\checkmarkc\checkmark$\checkmark\\checkmark$\checkmark^\checkmark\\checkmarkc\checkmarki\checkmarkr\checkmarkc\checkmark$\checkmarkh\checkmark$\checkmark^\checkmark\\checkmarkc\checkmarki\checkmarkr\checkmarkc\checkmark$\checkmarkl\checkmark$\checkmark^\checkmark\\checkmarkc\checkmarki\checkmarkr\checkmarkc\checkmark$\checkmarki\checkmark$\checkmark^\checkmark\\checkmarkc\checkmarki\checkmarkr\checkmarkc\checkmark$\checkmarkn\checkmark$\checkmark^\checkmark\\checkmarkc\checkmarki\checkmarkr\checkmarkc\checkmark$\checkmarke\checkmark$\checkmark^\checkmark\\checkmarkc\checkmarki\checkmarkr\checkmarkc\checkmark$\checkmark
\checkmark$\checkmark^\checkmark\\checkmarkc\checkmarki\checkmarkr\checkmarkc\checkmark$\checkmarkA\checkmark$\checkmark^\checkmark\\checkmarkc\checkmarki\checkmarkr\checkmarkc\checkmark$\checkmarkc\checkmark$\checkmark^\checkmark\\checkmarkc\checkmarki\checkmarkr\checkmarkc\checkmark$\checkmarki\checkmark$\checkmark^\checkmark\\checkmarkc\checkmarki\checkmarkr\checkmarkc\checkmark$\checkmarke\checkmark$\checkmark^\checkmark\\checkmarkc\checkmarki\checkmarkr\checkmarkc\checkmark$\checkmarkr\checkmark$\checkmark^\checkmark\\checkmarkc\checkmarki\checkmarkr\checkmarkc\checkmark$\checkmark \checkmark$\checkmark^\checkmark\\checkmarkc\checkmarki\checkmarkr\checkmarkc\checkmark$\checkmark4\checkmark$\checkmark^\checkmark\\checkmarkc\checkmarki\checkmarkr\checkmarkc\checkmark$\checkmark2\checkmark$\checkmark^\checkmark\\checkmarkc\checkmarki\checkmarkr\checkmarkc\checkmark$\checkmarkC\checkmark$\checkmark^\checkmark\\checkmarkc\checkmarki\checkmarkr\checkmarkc\checkmark$\checkmarkr\checkmark$\checkmark^\checkmark\\checkmarkc\checkmarki\checkmarkr\checkmarkc\checkmark$\checkmarkM\checkmark$\checkmark^\checkmark\\checkmarkc\checkmarki\checkmarkr\checkmarkc\checkmark$\checkmarko\checkmark$\checkmark^\checkmark\\checkmarkc\checkmarki\checkmarkr\checkmarkc\checkmark$\checkmark4\checkmark$\checkmark^\checkmark\\checkmarkc\checkmarki\checkmarkr\checkmarkc\checkmark$\checkmark \checkmark$\checkmark^\checkmark\\checkmarkc\checkmarki\checkmarkr\checkmarkc\checkmark$\checkmark&\checkmark$\checkmark^\checkmark\\checkmarkc\checkmarki\checkmarkr\checkmarkc\checkmark$\checkmark \checkmark$\checkmark^\checkmark\\checkmarkc\checkmarki\checkmarkr\checkmarkc\checkmark$\checkmark6\checkmark$\checkmark^\checkmark\\checkmarkc\checkmarki\checkmarkr\checkmarkc\checkmark$\checkmark5\checkmark$\checkmark^\checkmark\\checkmarkc\checkmarki\checkmarkr\checkmarkc\checkmark$\checkmark0\checkmark$\checkmark^\checkmark\\checkmarkc\checkmarki\checkmarkr\checkmarkc\checkmark$\checkmark \checkmark$\checkmark^\checkmark\\checkmarkc\checkmarki\checkmarkr\checkmarkc\checkmark$\checkmark&\checkmark$\checkmark^\checkmark\\checkmarkc\checkmarki\checkmarkr\checkmarkc\checkmark$\checkmark \checkmark$\checkmark^\checkmark\\checkmarkc\checkmarki\checkmarkr\checkmarkc\checkmark$\checkmark9\checkmark$\checkmark^\checkmark\\checkmarkc\checkmarki\checkmarkr\checkmarkc\checkmark$\checkmark0\checkmark$\checkmark^\checkmark\\checkmarkc\checkmarki\checkmarkr\checkmarkc\checkmark$\checkmark0\checkmark$\checkmark^\checkmark\\checkmarkc\checkmarki\checkmarkr\checkmarkc\checkmark$\checkmark \checkmark$\checkmark^\checkmark\\checkmarkc\checkmarki\checkmarkr\checkmarkc\checkmark$\checkmark&\checkmark$\checkmark^\checkmark\\checkmarkc\checkmarki\checkmarkr\checkmarkc\checkmark$\checkmark \checkmark$\checkmark^\checkmark\\checkmarkc\checkmarki\checkmarkr\checkmarkc\checkmark$\checkmark2\checkmark$\checkmark^\checkmark\\checkmarkc\checkmarki\checkmarkr\checkmarkc\checkmark$\checkmark1\checkmark$\checkmark^\checkmark\\checkmarkc\checkmarki\checkmarkr\checkmarkc\checkmark$\checkmark0\checkmark$\checkmark^\checkmark\\checkmarkc\checkmarki\checkmarkr\checkmarkc\checkmark$\checkmark \checkmark$\checkmark^\checkmark\\checkmarkc\checkmarki\checkmarkr\checkmarkc\checkmark$\checkmark&\checkmark$\checkmark^\checkmark\\checkmarkc\checkmarki\checkmarkr\checkmarkc\checkmark$\checkmark \checkmark$\checkmark^\checkmark\\checkmarkc\checkmarki\checkmarkr\checkmarkc\checkmark$\checkmarkS\checkmark$\checkmark^\checkmark\\checkmarkc\checkmarki\checkmarkr\checkmarkc\checkmark$\checkmarku\checkmark$\checkmark^\checkmark\\checkmarkc\checkmarki\checkmarkr\checkmarkc\checkmark$\checkmarkr\checkmark$\checkmark^\checkmark\\checkmarkc\checkmarki\checkmarkr\checkmarkc\checkmark$\checkmarkd\checkmark$\checkmark^\checkmark\\checkmarkc\checkmarki\checkmarkr\checkmarkc\checkmark$\checkmarki\checkmark$\checkmark^\checkmark\\checkmarkc\checkmarki\checkmarkr\checkmarkc\checkmark$\checkmarkm\checkmark$\checkmark^\checkmark\\checkmarkc\checkmarki\checkmarkr\checkmarkc\checkmark$\checkmarke\checkmark$\checkmark^\checkmark\\checkmarkc\checkmarki\checkmarkr\checkmarkc\checkmark$\checkmarkn\checkmark$\checkmark^\checkmark\\checkmarkc\checkmarki\checkmarkr\checkmarkc\checkmark$\checkmarks\checkmark$\checkmark^\checkmark\\checkmarkc\checkmarki\checkmarkr\checkmarkc\checkmark$\checkmarki\checkmark$\checkmark^\checkmark\\checkmarkc\checkmarki\checkmarkr\checkmarkc\checkmark$\checkmarko\checkmark$\checkmark^\checkmark\\checkmarkc\checkmarki\checkmarkr\checkmarkc\checkmark$\checkmarkn\checkmark$\checkmark^\checkmark\\checkmarkc\checkmarki\checkmarkr\checkmarkc\checkmark$\checkmarkn\checkmark$\checkmark^\checkmark\\checkmarkc\checkmarki\checkmarkr\checkmarkc\checkmark$\checkmarké\checkmark$\checkmark^\checkmark\\checkmarkc\checkmarki\checkmarkr\checkmarkc\checkmark$\checkmark \checkmark$\checkmark^\checkmark\\checkmarkc\checkmarki\checkmarkr\checkmarkc\checkmark$\checkmark\\checkmark$\checkmark^\checkmark\\checkmarkc\checkmarki\checkmarkr\checkmarkc\checkmark$\checkmark\\checkmark$\checkmark^\checkmark\\checkmarkc\checkmarki\checkmarkr\checkmarkc\checkmark$\checkmark \checkmark$\checkmark^\checkmark\\checkmarkc\checkmarki\checkmarkr\checkmarkc\checkmark$\checkmark\\checkmark$\checkmark^\checkmark\\checkmarkc\checkmarki\checkmarkr\checkmarkc\checkmark$\checkmarkh\checkmark$\checkmark^\checkmark\\checkmarkc\checkmarki\checkmarkr\checkmarkc\checkmark$\checkmarkl\checkmark$\checkmark^\checkmark\\checkmarkc\checkmarki\checkmarkr\checkmarkc\checkmark$\checkmarki\checkmark$\checkmark^\checkmark\\checkmarkc\checkmarki\checkmarkr\checkmarkc\checkmark$\checkmarkn\checkmark$\checkmark^\checkmark\\checkmarkc\checkmarki\checkmarkr\checkmarkc\checkmark$\checkmarke\checkmark$\checkmark^\checkmark\\checkmarkc\checkmarki\checkmarkr\checkmarkc\checkmark$\checkmark
\checkmark$\checkmark^\checkmark\\checkmarkc\checkmarki\checkmarkr\checkmarkc\checkmark$\checkmark\\checkmark$\checkmark^\checkmark\\checkmarkc\checkmarki\checkmarkr\checkmarkc\checkmark$\checkmarke\checkmark$\checkmark^\checkmark\\checkmarkc\checkmarki\checkmarkr\checkmarkc\checkmark$\checkmarkn\checkmark$\checkmark^\checkmark\\checkmarkc\checkmarki\checkmarkr\checkmarkc\checkmark$\checkmarkd\checkmark$\checkmark^\checkmark\\checkmarkc\checkmarki\checkmarkr\checkmarkc\checkmark$\checkmark{\checkmark$\checkmark^\checkmark\\checkmarkc\checkmarki\checkmarkr\checkmarkc\checkmark$\checkmarkt\checkmark$\checkmark^\checkmark\\checkmarkc\checkmarki\checkmarkr\checkmarkc\checkmark$\checkmarka\checkmark$\checkmark^\checkmark\\checkmarkc\checkmarki\checkmarkr\checkmarkc\checkmark$\checkmarkb\checkmark$\checkmark^\checkmark\\checkmarkc\checkmarki\checkmarkr\checkmarkc\checkmark$\checkmarku\checkmark$\checkmark^\checkmark\\checkmarkc\checkmarki\checkmarkr\checkmarkc\checkmark$\checkmarkl\checkmark$\checkmark^\checkmark\\checkmarkc\checkmarki\checkmarkr\checkmarkc\checkmark$\checkmarka\checkmark$\checkmark^\checkmark\\checkmarkc\checkmarki\checkmarkr\checkmarkc\checkmark$\checkmarkr\checkmark$\checkmark^\checkmark\\checkmarkc\checkmarki\checkmarkr\checkmarkc\checkmark$\checkmark}\checkmark$\checkmark^\checkmark\\checkmarkc\checkmarki\checkmarkr\checkmarkc\checkmark$\checkmark
\checkmark$\checkmark^\checkmark\\checkmarkc\checkmarki\checkmarkr\checkmarkc\checkmark$\checkmark\\checkmark$\checkmark^\checkmark\\checkmarkc\checkmarki\checkmarkr\checkmarkc\checkmark$\checkmarkc\checkmark$\checkmark^\checkmark\\checkmarkc\checkmarki\checkmarkr\checkmarkc\checkmark$\checkmarka\checkmark$\checkmark^\checkmark\\checkmarkc\checkmarki\checkmarkr\checkmarkc\checkmark$\checkmarkp\checkmark$\checkmark^\checkmark\\checkmarkc\checkmarki\checkmarkr\checkmarkc\checkmark$\checkmarkt\checkmark$\checkmark^\checkmark\\checkmarkc\checkmarki\checkmarkr\checkmarkc\checkmark$\checkmarki\checkmark$\checkmark^\checkmark\\checkmarkc\checkmarki\checkmarkr\checkmarkc\checkmark$\checkmarko\checkmark$\checkmark^\checkmark\\checkmarkc\checkmarki\checkmarkr\checkmarkc\checkmark$\checkmarkn\checkmark$\checkmark^\checkmark\\checkmarkc\checkmarki\checkmarkr\checkmarkc\checkmark$\checkmark{\checkmark$\checkmark^\checkmark\\checkmarkc\checkmarki\checkmarkr\checkmarkc\checkmark$\checkmarkT\checkmark$\checkmark^\checkmark\\checkmarkc\checkmarki\checkmarkr\checkmarkc\checkmark$\checkmarka\checkmark$\checkmark^\checkmark\\checkmarkc\checkmarki\checkmarkr\checkmarkc\checkmark$\checkmarkb\checkmark$\checkmark^\checkmark\\checkmarkc\checkmarki\checkmarkr\checkmarkc\checkmark$\checkmarkl\checkmark$\checkmark^\checkmark\\checkmarkc\checkmarki\checkmarkr\checkmarkc\checkmark$\checkmarke\checkmark$\checkmark^\checkmark\\checkmarkc\checkmarki\checkmarkr\checkmarkc\checkmark$\checkmarka\checkmark$\checkmark^\checkmark\\checkmarkc\checkmarki\checkmarkr\checkmarkc\checkmark$\checkmarku\checkmark$\checkmark^\checkmark\\checkmarkc\checkmarki\checkmarkr\checkmarkc\checkmark$\checkmark \checkmark$\checkmark^\checkmark\\checkmarkc\checkmarki\checkmarkr\checkmarkc\checkmark$\checkmarkd\checkmark$\checkmark^\checkmark\\checkmarkc\checkmarki\checkmarkr\checkmarkc\checkmark$\checkmarke\checkmark$\checkmark^\checkmark\\checkmarkc\checkmarki\checkmarkr\checkmarkc\checkmark$\checkmarks\checkmark$\checkmark^\checkmark\\checkmarkc\checkmarki\checkmarkr\checkmarkc\checkmark$\checkmark \checkmark$\checkmark^\checkmark\\checkmarkc\checkmarki\checkmarkr\checkmarkc\checkmark$\checkmarkp\checkmark$\checkmark^\checkmark\\checkmarkc\checkmarki\checkmarkr\checkmarkc\checkmark$\checkmarkr\checkmark$\checkmark^\checkmark\\checkmarkc\checkmarki\checkmarkr\checkmarkc\checkmark$\checkmarki\checkmark$\checkmark^\checkmark\\checkmarkc\checkmarki\checkmarkr\checkmarkc\checkmark$\checkmarkn\checkmark$\checkmark^\checkmark\\checkmarkc\checkmarki\checkmarkr\checkmarkc\checkmark$\checkmarkc\checkmark$\checkmark^\checkmark\\checkmarkc\checkmarki\checkmarkr\checkmarkc\checkmark$\checkmarki\checkmark$\checkmark^\checkmark\\checkmarkc\checkmarki\checkmarkr\checkmarkc\checkmark$\checkmarkp\checkmark$\checkmark^\checkmark\\checkmarkc\checkmarki\checkmarkr\checkmarkc\checkmark$\checkmarka\checkmark$\checkmark^\checkmark\\checkmarkc\checkmarki\checkmarkr\checkmarkc\checkmark$\checkmarku\checkmark$\checkmark^\checkmark\\checkmarkc\checkmarki\checkmarkr\checkmarkc\checkmark$\checkmarkx\checkmark$\checkmark^\checkmark\\checkmarkc\checkmarki\checkmarkr\checkmarkc\checkmark$\checkmark \checkmark$\checkmark^\checkmark\\checkmarkc\checkmarki\checkmarkr\checkmarkc\checkmark$\checkmarka\checkmark$\checkmark^\checkmark\\checkmarkc\checkmarki\checkmarkr\checkmarkc\checkmark$\checkmarkc\checkmark$\checkmark^\checkmark\\checkmarkc\checkmarki\checkmarkr\checkmarkc\checkmark$\checkmarki\checkmark$\checkmark^\checkmark\\checkmarkc\checkmarki\checkmarkr\checkmarkc\checkmark$\checkmarke\checkmark$\checkmark^\checkmark\\checkmarkc\checkmarki\checkmarkr\checkmarkc\checkmark$\checkmarkr\checkmark$\checkmark^\checkmark\\checkmarkc\checkmarki\checkmarkr\checkmarkc\checkmark$\checkmarks\checkmark$\checkmark^\checkmark\\checkmarkc\checkmarki\checkmarkr\checkmarkc\checkmark$\checkmark \checkmark$\checkmark^\checkmark\\checkmarkc\checkmarki\checkmarkr\checkmarkc\checkmark$\checkmarkp\checkmark$\checkmark^\checkmark\\checkmarkc\checkmarki\checkmarkr\checkmarkc\checkmark$\checkmarko\checkmark$\checkmark^\checkmark\\checkmarkc\checkmarki\checkmarkr\checkmarkc\checkmark$\checkmarku\checkmark$\checkmark^\checkmark\\checkmarkc\checkmarki\checkmarkr\checkmarkc\checkmark$\checkmarkr\checkmark$\checkmark^\checkmark\\checkmarkc\checkmarki\checkmarkr\checkmarkc\checkmark$\checkmark \checkmark$\checkmark^\checkmark\\checkmarkc\checkmarki\checkmarkr\checkmarkc\checkmark$\checkmarka\checkmark$\checkmark^\checkmark\\checkmarkc\checkmarki\checkmarkr\checkmarkc\checkmark$\checkmarkr\checkmark$\checkmark^\checkmark\\checkmarkc\checkmarki\checkmarkr\checkmarkc\checkmark$\checkmarkb\checkmark$\checkmark^\checkmark\\checkmarkc\checkmarki\checkmarkr\checkmarkc\checkmark$\checkmarkr\checkmark$\checkmark^\checkmark\\checkmarkc\checkmarki\checkmarkr\checkmarkc\checkmark$\checkmarke\checkmark$\checkmark^\checkmark\\checkmarkc\checkmarki\checkmarkr\checkmarkc\checkmark$\checkmark \checkmark$\checkmark^\checkmark\\checkmarkc\checkmarki\checkmarkr\checkmarkc\checkmark$\checkmark—\checkmark$\checkmark^\checkmark\\checkmarkc\checkmarki\checkmarkr\checkmarkc\checkmark$\checkmark \checkmark$\checkmark^\checkmark\\checkmarkc\checkmarki\checkmarkr\checkmarkc\checkmark$\checkmarkC\checkmark$\checkmark^\checkmark\\checkmarkc\checkmarki\checkmarkr\checkmarkc\checkmark$\checkmarkh\checkmark$\checkmark^\checkmark\\checkmarkc\checkmarki\checkmarkr\checkmarkc\checkmark$\checkmarko\checkmark$\checkmark^\checkmark\\checkmarkc\checkmarki\checkmarkr\checkmarkc\checkmark$\checkmarki\checkmark$\checkmark^\checkmark\\checkmarkc\checkmarki\checkmarkr\checkmarkc\checkmark$\checkmarkx\checkmark$\checkmark^\checkmark\\checkmarkc\checkmarki\checkmarkr\checkmarkc\checkmark$\checkmark \checkmark$\checkmark^\checkmark\\checkmarkc\checkmarki\checkmarkr\checkmarkc\checkmark$\checkmarkd\checkmark$\checkmark^\checkmark\\checkmarkc\checkmarki\checkmarkr\checkmarkc\checkmark$\checkmarku\checkmark$\checkmark^\checkmark\\checkmarkc\checkmarki\checkmarkr\checkmarkc\checkmark$\checkmark \checkmark$\checkmark^\checkmark\\checkmarkc\checkmarki\checkmarkr\checkmarkc\checkmark$\checkmarkm\checkmark$\checkmark^\checkmark\\checkmarkc\checkmarki\checkmarkr\checkmarkc\checkmark$\checkmarka\checkmark$\checkmark^\checkmark\\checkmarkc\checkmarki\checkmarkr\checkmarkc\checkmark$\checkmarkt\checkmark$\checkmark^\checkmark\\checkmarkc\checkmarki\checkmarkr\checkmarkc\checkmark$\checkmarké\checkmark$\checkmark^\checkmark\\checkmarkc\checkmarki\checkmarkr\checkmarkc\checkmark$\checkmarkr\checkmark$\checkmark^\checkmark\\checkmarkc\checkmarki\checkmarkr\checkmarkc\checkmark$\checkmarki\checkmark$\checkmark^\checkmark\\checkmarkc\checkmarki\checkmarkr\checkmarkc\checkmark$\checkmarka\checkmark$\checkmark^\checkmark\\checkmarkc\checkmarki\checkmarkr\checkmarkc\checkmark$\checkmarku\checkmark$\checkmark^\checkmark\\checkmarkc\checkmarki\checkmarkr\checkmarkc\checkmark$\checkmark}\checkmark$\checkmark^\checkmark\\checkmarkc\checkmarki\checkmarkr\checkmarkc\checkmark$\checkmark
\checkmark$\checkmark^\checkmark\\checkmarkc\checkmarki\checkmarkr\checkmarkc\checkmark$\checkmark\\checkmark$\checkmark^\checkmark\\checkmarkc\checkmarki\checkmarkr\checkmarkc\checkmark$\checkmarkl\checkmark$\checkmark^\checkmark\\checkmarkc\checkmarki\checkmarkr\checkmarkc\checkmark$\checkmarka\checkmark$\checkmark^\checkmark\\checkmarkc\checkmarki\checkmarkr\checkmarkc\checkmark$\checkmarkb\checkmark$\checkmark^\checkmark\\checkmarkc\checkmarki\checkmarkr\checkmarkc\checkmark$\checkmarke\checkmark$\checkmark^\checkmark\\checkmarkc\checkmarki\checkmarkr\checkmarkc\checkmark$\checkmarkl\checkmark$\checkmark^\checkmark\\checkmarkc\checkmarki\checkmarkr\checkmarkc\checkmark$\checkmark{\checkmark$\checkmark^\checkmark\\checkmarkc\checkmarki\checkmarkr\checkmarkc\checkmark$\checkmarkt\checkmark$\checkmark^\checkmark\\checkmarkc\checkmarki\checkmarkr\checkmarkc\checkmark$\checkmarka\checkmark$\checkmark^\checkmark\\checkmarkc\checkmarki\checkmarkr\checkmarkc\checkmark$\checkmarkb\checkmark$\checkmark^\checkmark\\checkmarkc\checkmarki\checkmarkr\checkmarkc\checkmark$\checkmark:\checkmark$\checkmark^\checkmark\\checkmarkc\checkmarki\checkmarkr\checkmarkc\checkmark$\checkmarkm\checkmark$\checkmark^\checkmark\\checkmarkc\checkmarki\checkmarkr\checkmarkc\checkmark$\checkmarka\checkmark$\checkmark^\checkmark\\checkmarkc\checkmarki\checkmarkr\checkmarkc\checkmark$\checkmarkt\checkmark$\checkmark^\checkmark\\checkmarkc\checkmarki\checkmarkr\checkmarkc\checkmark$\checkmarke\checkmark$\checkmark^\checkmark\\checkmarkc\checkmarki\checkmarkr\checkmarkc\checkmark$\checkmarkr\checkmark$\checkmark^\checkmark\\checkmarkc\checkmarki\checkmarkr\checkmarkc\checkmark$\checkmarki\checkmark$\checkmark^\checkmark\\checkmarkc\checkmarki\checkmarkr\checkmarkc\checkmark$\checkmarka\checkmark$\checkmark^\checkmark\\checkmarkc\checkmarki\checkmarkr\checkmarkc\checkmark$\checkmarku\checkmark$\checkmark^\checkmark\\checkmarkc\checkmarki\checkmarkr\checkmarkc\checkmark$\checkmarkx\checkmark$\checkmark^\checkmark\\checkmarkc\checkmarki\checkmarkr\checkmarkc\checkmark$\checkmark_\checkmark$\checkmark^\checkmark\\checkmarkc\checkmarki\checkmarkr\checkmarkc\checkmark$\checkmarka\checkmark$\checkmark^\checkmark\\checkmarkc\checkmarki\checkmarkr\checkmarkc\checkmark$\checkmarkr\checkmark$\checkmark^\checkmark\\checkmarkc\checkmarki\checkmarkr\checkmarkc\checkmark$\checkmarkb\checkmark$\checkmark^\checkmark\\checkmarkc\checkmarki\checkmarkr\checkmarkc\checkmark$\checkmarkr\checkmark$\checkmark^\checkmark\\checkmarkc\checkmarki\checkmarkr\checkmarkc\checkmark$\checkmarke\checkmark$\checkmark^\checkmark\\checkmarkc\checkmarki\checkmarkr\checkmarkc\checkmark$\checkmark}\checkmark$\checkmark^\checkmark\\checkmarkc\checkmarki\checkmarkr\checkmarkc\checkmark$\checkmark
\checkmark$\checkmark^\checkmark\\checkmarkc\checkmarki\checkmarkr\checkmarkc\checkmark$\checkmark\\checkmark$\checkmark^\checkmark\\checkmarkc\checkmarki\checkmarkr\checkmarkc\checkmark$\checkmarke\checkmark$\checkmark^\checkmark\\checkmarkc\checkmarki\checkmarkr\checkmarkc\checkmark$\checkmarkn\checkmark$\checkmark^\checkmark\\checkmarkc\checkmarki\checkmarkr\checkmarkc\checkmark$\checkmarkd\checkmark$\checkmark^\checkmark\\checkmarkc\checkmarki\checkmarkr\checkmarkc\checkmark$\checkmark{\checkmark$\checkmark^\checkmark\\checkmarkc\checkmarki\checkmarkr\checkmarkc\checkmark$\checkmarkt\checkmark$\checkmark^\checkmark\\checkmarkc\checkmarki\checkmarkr\checkmarkc\checkmark$\checkmarka\checkmark$\checkmark^\checkmark\\checkmarkc\checkmarki\checkmarkr\checkmarkc\checkmark$\checkmarkb\checkmark$\checkmark^\checkmark\\checkmarkc\checkmarki\checkmarkr\checkmarkc\checkmark$\checkmarkl\checkmark$\checkmark^\checkmark\\checkmarkc\checkmarki\checkmarkr\checkmarkc\checkmark$\checkmarke\checkmark$\checkmark^\checkmark\\checkmarkc\checkmarki\checkmarkr\checkmarkc\checkmark$\checkmark}\checkmark$\checkmark^\checkmark\\checkmarkc\checkmarki\checkmarkr\checkmarkc\checkmark$\checkmark
\checkmark$\checkmark^\checkmark\\checkmarkc\checkmarki\checkmarkr\checkmarkc\checkmark$\checkmark
\checkmark$\checkmark^\checkmark\\checkmarkc\checkmarki\checkmarkr\checkmarkc\checkmark$\checkmark\\checkmark$\checkmark^\checkmark\\checkmarkc\checkmarki\checkmarkr\checkmarkc\checkmark$\checkmarkt\checkmark$\checkmark^\checkmark\\checkmarkc\checkmarki\checkmarkr\checkmarkc\checkmark$\checkmarke\checkmark$\checkmark^\checkmark\\checkmarkc\checkmarki\checkmarkr\checkmarkc\checkmark$\checkmarkx\checkmark$\checkmark^\checkmark\\checkmarkc\checkmarki\checkmarkr\checkmarkc\checkmark$\checkmarkt\checkmark$\checkmark^\checkmark\\checkmarkc\checkmarki\checkmarkr\checkmarkc\checkmark$\checkmarkb\checkmark$\checkmark^\checkmark\\checkmarkc\checkmarki\checkmarkr\checkmarkc\checkmark$\checkmarkf\checkmark$\checkmark^\checkmark\\checkmarkc\checkmarki\checkmarkr\checkmarkc\checkmark$\checkmark{\checkmark$\checkmark^\checkmark\\checkmarkc\checkmarki\checkmarkr\checkmarkc\checkmark$\checkmarkM\checkmark$\checkmark^\checkmark\\checkmarkc\checkmarki\checkmarkr\checkmarkc\checkmark$\checkmarka\checkmark$\checkmark^\checkmark\\checkmarkc\checkmarki\checkmarkr\checkmarkc\checkmark$\checkmarkt\checkmark$\checkmark^\checkmark\\checkmarkc\checkmarki\checkmarkr\checkmarkc\checkmark$\checkmarké\checkmark$\checkmark^\checkmark\\checkmarkc\checkmarki\checkmarkr\checkmarkc\checkmark$\checkmarkr\checkmark$\checkmark^\checkmark\\checkmarkc\checkmarki\checkmarkr\checkmarkc\checkmark$\checkmarki\checkmark$\checkmark^\checkmark\\checkmarkc\checkmarki\checkmarkr\checkmarkc\checkmark$\checkmarka\checkmark$\checkmark^\checkmark\\checkmarkc\checkmarki\checkmarkr\checkmarkc\checkmark$\checkmarku\checkmark$\checkmark^\checkmark\\checkmarkc\checkmarki\checkmarkr\checkmarkc\checkmark$\checkmark \checkmark$\checkmark^\checkmark\\checkmarkc\checkmarki\checkmarkr\checkmarkc\checkmark$\checkmarkr\checkmark$\checkmark^\checkmark\\checkmarkc\checkmarki\checkmarkr\checkmarkc\checkmark$\checkmarke\checkmark$\checkmark^\checkmark\\checkmarkc\checkmarki\checkmarkr\checkmarkc\checkmark$\checkmarkt\checkmark$\checkmark^\checkmark\\checkmarkc\checkmarki\checkmarkr\checkmarkc\checkmark$\checkmarke\checkmark$\checkmark^\checkmark\\checkmarkc\checkmarki\checkmarkr\checkmarkc\checkmark$\checkmarkn\checkmark$\checkmark^\checkmark\\checkmarkc\checkmarki\checkmarkr\checkmarkc\checkmark$\checkmarku\checkmark$\checkmark^\checkmark\\checkmarkc\checkmarki\checkmarkr\checkmarkc\checkmark$\checkmark \checkmark$\checkmark^\checkmark\\checkmarkc\checkmarki\checkmarkr\checkmarkc\checkmark$\checkmark:\checkmark$\checkmark^\checkmark\\checkmarkc\checkmarki\checkmarkr\checkmarkc\checkmark$\checkmark \checkmark$\checkmark^\checkmark\\checkmarkc\checkmarki\checkmarkr\checkmarkc\checkmark$\checkmarkA\checkmark$\checkmark^\checkmark\\checkmarkc\checkmarki\checkmarkr\checkmarkc\checkmark$\checkmarkc\checkmark$\checkmark^\checkmark\\checkmarkc\checkmarki\checkmarkr\checkmarkc\checkmark$\checkmarki\checkmark$\checkmark^\checkmark\\checkmarkc\checkmarki\checkmarkr\checkmarkc\checkmark$\checkmarke\checkmark$\checkmark^\checkmark\\checkmarkc\checkmarki\checkmarkr\checkmarkc\checkmark$\checkmarkr\checkmark$\checkmark^\checkmark\\checkmarkc\checkmarki\checkmarkr\checkmarkc\checkmark$\checkmark \checkmark$\checkmark^\checkmark\\checkmarkc\checkmarki\checkmarkr\checkmarkc\checkmark$\checkmarkC\checkmark$\checkmark^\checkmark\\checkmarkc\checkmarki\checkmarkr\checkmarkc\checkmark$\checkmark4\checkmark$\checkmark^\checkmark\\checkmarkc\checkmarki\checkmarkr\checkmarkc\checkmark$\checkmark5\checkmark$\checkmark^\checkmark\\checkmarkc\checkmarki\checkmarkr\checkmarkc\checkmark$\checkmark \checkmark$\checkmark^\checkmark\\checkmarkc\checkmarki\checkmarkr\checkmarkc\checkmark$\checkmark(\checkmark$\checkmark^\checkmark\\checkmarkc\checkmarki\checkmarkr\checkmarkc\checkmark$\checkmarkn\checkmark$\checkmark^\checkmark\\checkmarkc\checkmarki\checkmarkr\checkmarkc\checkmark$\checkmarko\checkmark$\checkmark^\checkmark\\checkmarkc\checkmarki\checkmarkr\checkmarkc\checkmark$\checkmarkr\checkmark$\checkmark^\checkmark\\checkmarkc\checkmarki\checkmarkr\checkmarkc\checkmark$\checkmarkm\checkmark$\checkmark^\checkmark\\checkmarkc\checkmarki\checkmarkr\checkmarkc\checkmark$\checkmarka\checkmark$\checkmark^\checkmark\\checkmarkc\checkmarki\checkmarkr\checkmarkc\checkmark$\checkmarkl\checkmark$\checkmark^\checkmark\\checkmarkc\checkmarki\checkmarkr\checkmarkc\checkmark$\checkmarki\checkmark$\checkmark^\checkmark\\checkmarkc\checkmarki\checkmarkr\checkmarkc\checkmark$\checkmarks\checkmark$\checkmark^\checkmark\\checkmarkc\checkmarki\checkmarkr\checkmarkc\checkmark$\checkmarké\checkmark$\checkmark^\checkmark\\checkmarkc\checkmarki\checkmarkr\checkmarkc\checkmark$\checkmark)\checkmark$\checkmark^\checkmark\\checkmarkc\checkmarki\checkmarkr\checkmarkc\checkmark$\checkmark}\checkmark$\checkmark^\checkmark\\checkmarkc\checkmarki\checkmarkr\checkmarkc\checkmark$\checkmark \checkmark$\checkmark^\checkmark\\checkmarkc\checkmarki\checkmarkr\checkmarkc\checkmark$\checkmark:\checkmark$\checkmark^\checkmark\\checkmarkc\checkmarki\checkmarkr\checkmarkc\checkmark$\checkmark
\checkmark$\checkmark^\checkmark\\checkmarkc\checkmarki\checkmarkr\checkmarkc\checkmark$\checkmark\\checkmark$\checkmark^\checkmark\\checkmarkc\checkmarki\checkmarkr\checkmarkc\checkmark$\checkmarkb\checkmark$\checkmark^\checkmark\\checkmarkc\checkmarki\checkmarkr\checkmarkc\checkmark$\checkmarke\checkmark$\checkmark^\checkmark\\checkmarkc\checkmarki\checkmarkr\checkmarkc\checkmark$\checkmarkg\checkmark$\checkmark^\checkmark\\checkmarkc\checkmarki\checkmarkr\checkmarkc\checkmark$\checkmarki\checkmark$\checkmark^\checkmark\\checkmarkc\checkmarki\checkmarkr\checkmarkc\checkmark$\checkmarkn\checkmark$\checkmark^\checkmark\\checkmarkc\checkmarki\checkmarkr\checkmarkc\checkmark$\checkmark{\checkmark$\checkmark^\checkmark\\checkmarkc\checkmarki\checkmarkr\checkmarkc\checkmark$\checkmarki\checkmark$\checkmark^\checkmark\\checkmarkc\checkmarki\checkmarkr\checkmarkc\checkmark$\checkmarkt\checkmark$\checkmark^\checkmark\\checkmarkc\checkmarki\checkmarkr\checkmarkc\checkmark$\checkmarke\checkmark$\checkmark^\checkmark\\checkmarkc\checkmarki\checkmarkr\checkmarkc\checkmark$\checkmarkm\checkmark$\checkmark^\checkmark\\checkmarkc\checkmarki\checkmarkr\checkmarkc\checkmark$\checkmarki\checkmark$\checkmark^\checkmark\\checkmarkc\checkmarki\checkmarkr\checkmarkc\checkmark$\checkmarkz\checkmark$\checkmark^\checkmark\\checkmarkc\checkmarki\checkmarkr\checkmarkc\checkmark$\checkmarke\checkmark$\checkmark^\checkmark\\checkmarkc\checkmarki\checkmarkr\checkmarkc\checkmark$\checkmark}\checkmark$\checkmark^\checkmark\\checkmarkc\checkmarki\checkmarkr\checkmarkc\checkmark$\checkmark
\checkmark$\checkmark^\checkmark\\checkmarkc\checkmarki\checkmarkr\checkmarkc\checkmark$\checkmark \checkmark$\checkmark^\checkmark\\checkmarkc\checkmarki\checkmarkr\checkmarkc\checkmark$\checkmark \checkmark$\checkmark^\checkmark\\checkmarkc\checkmarki\checkmarkr\checkmarkc\checkmark$\checkmark \checkmark$\checkmark^\checkmark\\checkmarkc\checkmarki\checkmarkr\checkmarkc\checkmark$\checkmark \checkmark$\checkmark^\checkmark\\checkmarkc\checkmarki\checkmarkr\checkmarkc\checkmark$\checkmark\\checkmark$\checkmark^\checkmark\\checkmarkc\checkmarki\checkmarkr\checkmarkc\checkmark$\checkmarki\checkmark$\checkmark^\checkmark\\checkmarkc\checkmarki\checkmarkr\checkmarkc\checkmark$\checkmarkt\checkmark$\checkmark^\checkmark\\checkmarkc\checkmarki\checkmarkr\checkmarkc\checkmark$\checkmarke\checkmark$\checkmark^\checkmark\\checkmarkc\checkmarki\checkmarkr\checkmarkc\checkmark$\checkmarkm\checkmark$\checkmark^\checkmark\\checkmarkc\checkmarki\checkmarkr\checkmarkc\checkmark$\checkmark \checkmark$\checkmark^\checkmark\\checkmarkc\checkmarki\checkmarkr\checkmarkc\checkmark$\checkmarkL\checkmark$\checkmark^\checkmark\\checkmarkc\checkmarki\checkmarkr\checkmarkc\checkmark$\checkmarki\checkmark$\checkmark^\checkmark\\checkmarkc\checkmarki\checkmarkr\checkmarkc\checkmark$\checkmarkm\checkmark$\checkmark^\checkmark\\checkmarkc\checkmarki\checkmarkr\checkmarkc\checkmark$\checkmarki\checkmark$\checkmark^\checkmark\\checkmarkc\checkmarki\checkmarkr\checkmarkc\checkmark$\checkmarkt\checkmark$\checkmark^\checkmark\\checkmarkc\checkmarki\checkmarkr\checkmarkc\checkmark$\checkmarke\checkmark$\checkmark^\checkmark\\checkmarkc\checkmarki\checkmarkr\checkmarkc\checkmark$\checkmark \checkmark$\checkmark^\checkmark\\checkmarkc\checkmarki\checkmarkr\checkmarkc\checkmark$\checkmarkd\checkmark$\checkmark^\checkmark\\checkmarkc\checkmarki\checkmarkr\checkmarkc\checkmark$\checkmark'\checkmark$\checkmark^\checkmark\\checkmarkc\checkmarki\checkmarkr\checkmarkc\checkmark$\checkmarké\checkmark$\checkmark^\checkmark\\checkmarkc\checkmarki\checkmarkr\checkmarkc\checkmark$\checkmarkl\checkmark$\checkmark^\checkmark\\checkmarkc\checkmarki\checkmarkr\checkmarkc\checkmark$\checkmarka\checkmark$\checkmark^\checkmark\\checkmarkc\checkmarki\checkmarkr\checkmarkc\checkmark$\checkmarks\checkmark$\checkmark^\checkmark\\checkmarkc\checkmarki\checkmarkr\checkmarkc\checkmark$\checkmarkt\checkmark$\checkmark^\checkmark\\checkmarkc\checkmarki\checkmarkr\checkmarkc\checkmark$\checkmarki\checkmark$\checkmark^\checkmark\\checkmarkc\checkmarki\checkmarkr\checkmarkc\checkmark$\checkmarkc\checkmark$\checkmark^\checkmark\\checkmarkc\checkmarki\checkmarkr\checkmarkc\checkmark$\checkmarki\checkmark$\checkmark^\checkmark\\checkmarkc\checkmarki\checkmarkr\checkmarkc\checkmark$\checkmarkt\checkmark$\checkmark^\checkmark\\checkmarkc\checkmarki\checkmarkr\checkmarkc\checkmark$\checkmarké\checkmark$\checkmark^\checkmark\\checkmarkc\checkmarki\checkmarkr\checkmarkc\checkmark$\checkmark \checkmark$\checkmark^\checkmark\\checkmarkc\checkmarki\checkmarkr\checkmarkc\checkmark$\checkmark:\checkmark$\checkmark^\checkmark\\checkmarkc\checkmarki\checkmarkr\checkmarkc\checkmark$\checkmark \checkmark$\checkmark^\checkmark\\checkmarkc\checkmarki\checkmarkr\checkmarkc\checkmark$\checkmark$\checkmark$\checkmark^\checkmark\\checkmarkc\checkmarki\checkmarkr\checkmarkc\checkmark$\checkmarkR\checkmark$\checkmark^\checkmark\\checkmarkc\checkmarki\checkmarkr\checkmarkc\checkmark$\checkmark_\checkmark$\checkmark^\checkmark\\checkmarkc\checkmarki\checkmarkr\checkmarkc\checkmark$\checkmarke\checkmark$\checkmark^\checkmark\\checkmarkc\checkmarki\checkmarkr\checkmarkc\checkmark$\checkmark \checkmark$\checkmark^\checkmark\\checkmarkc\checkmarki\checkmarkr\checkmarkc\checkmark$\checkmark=\checkmark$\checkmark^\checkmark\\checkmarkc\checkmarki\checkmarkr\checkmarkc\checkmark$\checkmark \checkmark$\checkmark^\checkmark\\checkmarkc\checkmarki\checkmarkr\checkmarkc\checkmark$\checkmark3\checkmark$\checkmark^\checkmark\\checkmarkc\checkmarki\checkmarkr\checkmarkc\checkmark$\checkmark5\checkmark$\checkmark^\checkmark\\checkmarkc\checkmarki\checkmarkr\checkmarkc\checkmark$\checkmark0\checkmark$\checkmark^\checkmark\\checkmarkc\checkmarki\checkmarkr\checkmarkc\checkmark$\checkmark$\checkmark$\checkmark^\checkmark\\checkmarkc\checkmarki\checkmarkr\checkmarkc\checkmark$\checkmark \checkmark$\checkmark^\checkmark\\checkmarkc\checkmarki\checkmarkr\checkmarkc\checkmark$\checkmarkM\checkmark$\checkmark^\checkmark\\checkmarkc\checkmarki\checkmarkr\checkmarkc\checkmark$\checkmarkP\checkmark$\checkmark^\checkmark\\checkmarkc\checkmarki\checkmarkr\checkmarkc\checkmark$\checkmarka\checkmark$\checkmark^\checkmark\\checkmarkc\checkmarki\checkmarkr\checkmarkc\checkmark$\checkmark
\checkmark$\checkmark^\checkmark\\checkmarkc\checkmarki\checkmarkr\checkmarkc\checkmark$\checkmark \checkmark$\checkmark^\checkmark\\checkmarkc\checkmarki\checkmarkr\checkmarkc\checkmark$\checkmark \checkmark$\checkmark^\checkmark\\checkmarkc\checkmarki\checkmarkr\checkmarkc\checkmark$\checkmark \checkmark$\checkmark^\checkmark\\checkmarkc\checkmarki\checkmarkr\checkmarkc\checkmark$\checkmark \checkmark$\checkmark^\checkmark\\checkmarkc\checkmarki\checkmarkr\checkmarkc\checkmark$\checkmark\\checkmark$\checkmark^\checkmark\\checkmarkc\checkmarki\checkmarkr\checkmarkc\checkmark$\checkmarki\checkmark$\checkmark^\checkmark\\checkmarkc\checkmarki\checkmarkr\checkmarkc\checkmark$\checkmarkt\checkmark$\checkmark^\checkmark\\checkmarkc\checkmarki\checkmarkr\checkmarkc\checkmark$\checkmarke\checkmark$\checkmark^\checkmark\\checkmarkc\checkmarki\checkmarkr\checkmarkc\checkmark$\checkmarkm\checkmark$\checkmark^\checkmark\\checkmarkc\checkmarki\checkmarkr\checkmarkc\checkmark$\checkmark \checkmark$\checkmark^\checkmark\\checkmarkc\checkmarki\checkmarkr\checkmarkc\checkmark$\checkmarkR\checkmark$\checkmark^\checkmark\\checkmarkc\checkmarki\checkmarkr\checkmarkc\checkmark$\checkmarké\checkmark$\checkmark^\checkmark\\checkmarkc\checkmarki\checkmarkr\checkmarkc\checkmark$\checkmarks\checkmark$\checkmark^\checkmark\\checkmarkc\checkmarki\checkmarkr\checkmarkc\checkmark$\checkmarki\checkmark$\checkmark^\checkmark\\checkmarkc\checkmarki\checkmarkr\checkmarkc\checkmark$\checkmarks\checkmark$\checkmark^\checkmark\\checkmarkc\checkmarki\checkmarkr\checkmarkc\checkmark$\checkmarkt\checkmark$\checkmark^\checkmark\\checkmarkc\checkmarki\checkmarkr\checkmarkc\checkmark$\checkmarka\checkmark$\checkmark^\checkmark\\checkmarkc\checkmarki\checkmarkr\checkmarkc\checkmark$\checkmarkn\checkmark$\checkmark^\checkmark\\checkmarkc\checkmarki\checkmarkr\checkmarkc\checkmark$\checkmarkc\checkmark$\checkmark^\checkmark\\checkmarkc\checkmarki\checkmarkr\checkmarkc\checkmark$\checkmarke\checkmark$\checkmark^\checkmark\\checkmarkc\checkmarki\checkmarkr\checkmarkc\checkmark$\checkmark \checkmark$\checkmark^\checkmark\\checkmarkc\checkmarki\checkmarkr\checkmarkc\checkmark$\checkmarkà\checkmark$\checkmark^\checkmark\\checkmarkc\checkmarki\checkmarkr\checkmarkc\checkmark$\checkmark \checkmark$\checkmark^\checkmark\\checkmarkc\checkmarki\checkmarkr\checkmarkc\checkmark$\checkmarkl\checkmark$\checkmark^\checkmark\\checkmarkc\checkmarki\checkmarkr\checkmarkc\checkmark$\checkmarka\checkmark$\checkmark^\checkmark\\checkmarkc\checkmarki\checkmarkr\checkmarkc\checkmark$\checkmark \checkmark$\checkmark^\checkmark\\checkmarkc\checkmarki\checkmarkr\checkmarkc\checkmark$\checkmarkr\checkmark$\checkmark^\checkmark\\checkmarkc\checkmarki\checkmarkr\checkmarkc\checkmark$\checkmarku\checkmark$\checkmark^\checkmark\\checkmarkc\checkmarki\checkmarkr\checkmarkc\checkmark$\checkmarkp\checkmark$\checkmark^\checkmark\\checkmarkc\checkmarki\checkmarkr\checkmarkc\checkmark$\checkmarkt\checkmark$\checkmark^\checkmark\\checkmarkc\checkmarki\checkmarkr\checkmarkc\checkmark$\checkmarku\checkmark$\checkmark^\checkmark\\checkmarkc\checkmarki\checkmarkr\checkmarkc\checkmark$\checkmarkr\checkmark$\checkmark^\checkmark\\checkmarkc\checkmarki\checkmarkr\checkmarkc\checkmark$\checkmarke\checkmark$\checkmark^\checkmark\\checkmarkc\checkmarki\checkmarkr\checkmarkc\checkmark$\checkmark \checkmark$\checkmark^\checkmark\\checkmarkc\checkmarki\checkmarkr\checkmarkc\checkmark$\checkmark:\checkmark$\checkmark^\checkmark\\checkmarkc\checkmarki\checkmarkr\checkmarkc\checkmark$\checkmark \checkmark$\checkmark^\checkmark\\checkmarkc\checkmarki\checkmarkr\checkmarkc\checkmark$\checkmark$\checkmark$\checkmark^\checkmark\\checkmarkc\checkmarki\checkmarkr\checkmarkc\checkmark$\checkmarkR\checkmark$\checkmark^\checkmark\\checkmarkc\checkmarki\checkmarkr\checkmarkc\checkmark$\checkmark_\checkmark$\checkmark^\checkmark\\checkmarkc\checkmarki\checkmarkr\checkmarkc\checkmark$\checkmarkm\checkmark$\checkmark^\checkmark\\checkmarkc\checkmarki\checkmarkr\checkmarkc\checkmark$\checkmark \checkmark$\checkmark^\checkmark\\checkmarkc\checkmarki\checkmarkr\checkmarkc\checkmark$\checkmark=\checkmark$\checkmark^\checkmark\\checkmarkc\checkmarki\checkmarkr\checkmarkc\checkmark$\checkmark \checkmark$\checkmark^\checkmark\\checkmarkc\checkmarki\checkmarkr\checkmarkc\checkmark$\checkmark6\checkmark$\checkmark^\checkmark\\checkmarkc\checkmarki\checkmarkr\checkmarkc\checkmark$\checkmark0\checkmark$\checkmark^\checkmark\\checkmarkc\checkmarki\checkmarkr\checkmarkc\checkmark$\checkmark0\checkmark$\checkmark^\checkmark\\checkmarkc\checkmarki\checkmarkr\checkmarkc\checkmark$\checkmark$\checkmark$\checkmark^\checkmark\\checkmarkc\checkmarki\checkmarkr\checkmarkc\checkmark$\checkmark \checkmark$\checkmark^\checkmark\\checkmarkc\checkmarki\checkmarkr\checkmarkc\checkmark$\checkmarkM\checkmark$\checkmark^\checkmark\\checkmarkc\checkmarki\checkmarkr\checkmarkc\checkmark$\checkmarkP\checkmark$\checkmark^\checkmark\\checkmarkc\checkmarki\checkmarkr\checkmarkc\checkmark$\checkmarka\checkmark$\checkmark^\checkmark\\checkmarkc\checkmarki\checkmarkr\checkmarkc\checkmark$\checkmark
\checkmark$\checkmark^\checkmark\\checkmarkc\checkmarki\checkmarkr\checkmarkc\checkmark$\checkmark \checkmark$\checkmark^\checkmark\\checkmarkc\checkmarki\checkmarkr\checkmarkc\checkmark$\checkmark \checkmark$\checkmark^\checkmark\\checkmarkc\checkmarki\checkmarkr\checkmarkc\checkmark$\checkmark \checkmark$\checkmark^\checkmark\\checkmarkc\checkmarki\checkmarkr\checkmarkc\checkmark$\checkmark \checkmark$\checkmark^\checkmark\\checkmarkc\checkmarki\checkmarkr\checkmarkc\checkmark$\checkmark\\checkmark$\checkmark^\checkmark\\checkmarkc\checkmarki\checkmarkr\checkmarkc\checkmark$\checkmarki\checkmark$\checkmark^\checkmark\\checkmarkc\checkmarki\checkmarkr\checkmarkc\checkmark$\checkmarkt\checkmark$\checkmark^\checkmark\\checkmarkc\checkmarki\checkmarkr\checkmarkc\checkmark$\checkmarke\checkmark$\checkmark^\checkmark\\checkmarkc\checkmarki\checkmarkr\checkmarkc\checkmark$\checkmarkm\checkmark$\checkmark^\checkmark\\checkmarkc\checkmarki\checkmarkr\checkmarkc\checkmark$\checkmark \checkmark$\checkmark^\checkmark\\checkmarkc\checkmarki\checkmarkr\checkmarkc\checkmark$\checkmarkM\checkmark$\checkmark^\checkmark\\checkmarkc\checkmarki\checkmarkr\checkmarkc\checkmark$\checkmarko\checkmark$\checkmark^\checkmark\\checkmarkc\checkmarki\checkmarkr\checkmarkc\checkmark$\checkmarkd\checkmark$\checkmark^\checkmark\\checkmarkc\checkmarki\checkmarkr\checkmarkc\checkmark$\checkmarku\checkmark$\checkmark^\checkmark\\checkmarkc\checkmarki\checkmarkr\checkmarkc\checkmark$\checkmarkl\checkmark$\checkmark^\checkmark\\checkmarkc\checkmarki\checkmarkr\checkmarkc\checkmark$\checkmarke\checkmark$\checkmark^\checkmark\\checkmarkc\checkmarki\checkmarkr\checkmarkc\checkmark$\checkmark \checkmark$\checkmark^\checkmark\\checkmarkc\checkmarki\checkmarkr\checkmarkc\checkmark$\checkmarkd\checkmark$\checkmark^\checkmark\\checkmarkc\checkmarki\checkmarkr\checkmarkc\checkmark$\checkmarke\checkmark$\checkmark^\checkmark\\checkmarkc\checkmarki\checkmarkr\checkmarkc\checkmark$\checkmark \checkmark$\checkmark^\checkmark\\checkmarkc\checkmarki\checkmarkr\checkmarkc\checkmark$\checkmarkY\checkmark$\checkmark^\checkmark\\checkmarkc\checkmarki\checkmarkr\checkmarkc\checkmark$\checkmarko\checkmark$\checkmark^\checkmark\\checkmarkc\checkmarki\checkmarkr\checkmarkc\checkmark$\checkmarku\checkmark$\checkmark^\checkmark\\checkmarkc\checkmarki\checkmarkr\checkmarkc\checkmark$\checkmarkn\checkmark$\checkmark^\checkmark\\checkmarkc\checkmarki\checkmarkr\checkmarkc\checkmark$\checkmarkg\checkmark$\checkmark^\checkmark\\checkmarkc\checkmarki\checkmarkr\checkmarkc\checkmark$\checkmark \checkmark$\checkmark^\checkmark\\checkmarkc\checkmarki\checkmarkr\checkmarkc\checkmark$\checkmark:\checkmark$\checkmark^\checkmark\\checkmarkc\checkmarki\checkmarkr\checkmarkc\checkmark$\checkmark \checkmark$\checkmark^\checkmark\\checkmarkc\checkmarki\checkmarkr\checkmarkc\checkmark$\checkmark$\checkmark$\checkmark^\checkmark\\checkmarkc\checkmarki\checkmarkr\checkmarkc\checkmark$\checkmarkE\checkmark$\checkmark^\checkmark\\checkmarkc\checkmarki\checkmarkr\checkmarkc\checkmark$\checkmark \checkmark$\checkmark^\checkmark\\checkmarkc\checkmarki\checkmarkr\checkmarkc\checkmark$\checkmark=\checkmark$\checkmark^\checkmark\\checkmarkc\checkmarki\checkmarkr\checkmarkc\checkmark$\checkmark \checkmark$\checkmark^\checkmark\\checkmarkc\checkmarki\checkmarkr\checkmarkc\checkmark$\checkmark2\checkmark$\checkmark^\checkmark\\checkmarkc\checkmarki\checkmarkr\checkmarkc\checkmark$\checkmark1\checkmark$\checkmark^\checkmark\\checkmarkc\checkmarki\checkmarkr\checkmarkc\checkmark$\checkmark0\checkmark$\checkmark^\checkmark\\checkmarkc\checkmarki\checkmarkr\checkmarkc\checkmark$\checkmark\\checkmark$\checkmark^\checkmark\\checkmarkc\checkmarki\checkmarkr\checkmarkc\checkmark$\checkmark,\checkmark$\checkmark^\checkmark\\checkmarkc\checkmarki\checkmarkr\checkmarkc\checkmark$\checkmark0\checkmark$\checkmark^\checkmark\\checkmarkc\checkmarki\checkmarkr\checkmarkc\checkmark$\checkmark0\checkmark$\checkmark^\checkmark\\checkmarkc\checkmarki\checkmarkr\checkmarkc\checkmark$\checkmark0\checkmark$\checkmark^\checkmark\\checkmarkc\checkmarki\checkmarkr\checkmarkc\checkmark$\checkmark$\checkmark$\checkmark^\checkmark\\checkmarkc\checkmarki\checkmarkr\checkmarkc\checkmark$\checkmark \checkmark$\checkmark^\checkmark\\checkmarkc\checkmarki\checkmarkr\checkmarkc\checkmark$\checkmarkM\checkmark$\checkmark^\checkmark\\checkmarkc\checkmarki\checkmarkr\checkmarkc\checkmark$\checkmarkP\checkmark$\checkmark^\checkmark\\checkmarkc\checkmarki\checkmarkr\checkmarkc\checkmark$\checkmarka\checkmark$\checkmark^\checkmark\\checkmarkc\checkmarki\checkmarkr\checkmarkc\checkmark$\checkmark
\checkmark$\checkmark^\checkmark\\checkmarkc\checkmarki\checkmarkr\checkmarkc\checkmark$\checkmark \checkmark$\checkmark^\checkmark\\checkmarkc\checkmarki\checkmarkr\checkmarkc\checkmark$\checkmark \checkmark$\checkmark^\checkmark\\checkmarkc\checkmarki\checkmarkr\checkmarkc\checkmark$\checkmark \checkmark$\checkmark^\checkmark\\checkmarkc\checkmarki\checkmarkr\checkmarkc\checkmark$\checkmark \checkmark$\checkmark^\checkmark\\checkmarkc\checkmarki\checkmarkr\checkmarkc\checkmark$\checkmark\\checkmark$\checkmark^\checkmark\\checkmarkc\checkmarki\checkmarkr\checkmarkc\checkmark$\checkmarki\checkmark$\checkmark^\checkmark\\checkmarkc\checkmarki\checkmarkr\checkmarkc\checkmark$\checkmarkt\checkmark$\checkmark^\checkmark\\checkmarkc\checkmarki\checkmarkr\checkmarkc\checkmark$\checkmarke\checkmark$\checkmark^\checkmark\\checkmarkc\checkmarki\checkmarkr\checkmarkc\checkmark$\checkmarkm\checkmark$\checkmark^\checkmark\\checkmarkc\checkmarki\checkmarkr\checkmarkc\checkmark$\checkmark \checkmark$\checkmark^\checkmark\\checkmarkc\checkmarki\checkmarkr\checkmarkc\checkmark$\checkmarkL\checkmark$\checkmark^\checkmark\\checkmarkc\checkmarki\checkmarkr\checkmarkc\checkmark$\checkmarki\checkmark$\checkmark^\checkmark\\checkmarkc\checkmarki\checkmarkr\checkmarkc\checkmark$\checkmarkm\checkmark$\checkmark^\checkmark\\checkmarkc\checkmarki\checkmarkr\checkmarkc\checkmark$\checkmarki\checkmark$\checkmark^\checkmark\\checkmarkc\checkmarki\checkmarkr\checkmarkc\checkmark$\checkmarkt\checkmark$\checkmark^\checkmark\\checkmarkc\checkmarki\checkmarkr\checkmarkc\checkmark$\checkmarke\checkmark$\checkmark^\checkmark\\checkmarkc\checkmarki\checkmarkr\checkmarkc\checkmark$\checkmark \checkmark$\checkmark^\checkmark\\checkmarkc\checkmarki\checkmarkr\checkmarkc\checkmark$\checkmarkd\checkmark$\checkmark^\checkmark\\checkmarkc\checkmarki\checkmarkr\checkmarkc\checkmark$\checkmarke\checkmark$\checkmark^\checkmark\\checkmarkc\checkmarki\checkmarkr\checkmarkc\checkmark$\checkmark \checkmark$\checkmark^\checkmark\\checkmarkc\checkmarki\checkmarkr\checkmarkc\checkmark$\checkmarkf\checkmark$\checkmark^\checkmark\\checkmarkc\checkmarki\checkmarkr\checkmarkc\checkmark$\checkmarka\checkmark$\checkmark^\checkmark\\checkmarkc\checkmarki\checkmarkr\checkmarkc\checkmark$\checkmarkt\checkmark$\checkmark^\checkmark\\checkmarkc\checkmarki\checkmarkr\checkmarkc\checkmark$\checkmarki\checkmark$\checkmark^\checkmark\\checkmarkc\checkmarki\checkmarkr\checkmarkc\checkmark$\checkmarkg\checkmark$\checkmark^\checkmark\\checkmarkc\checkmarki\checkmarkr\checkmarkc\checkmark$\checkmarku\checkmark$\checkmark^\checkmark\\checkmarkc\checkmarki\checkmarkr\checkmarkc\checkmark$\checkmarke\checkmark$\checkmark^\checkmark\\checkmarkc\checkmarki\checkmarkr\checkmarkc\checkmark$\checkmark \checkmark$\checkmark^\checkmark\\checkmarkc\checkmarki\checkmarkr\checkmarkc\checkmark$\checkmarke\checkmark$\checkmark^\checkmark\\checkmarkc\checkmarki\checkmarkr\checkmarkc\checkmark$\checkmarkn\checkmark$\checkmark^\checkmark\\checkmarkc\checkmarki\checkmarkr\checkmarkc\checkmark$\checkmark \checkmark$\checkmark^\checkmark\\checkmarkc\checkmarki\checkmarkr\checkmarkc\checkmark$\checkmarkf\checkmark$\checkmark^\checkmark\\checkmarkc\checkmarki\checkmarkr\checkmarkc\checkmark$\checkmarkl\checkmark$\checkmark^\checkmark\\checkmarkc\checkmarki\checkmarkr\checkmarkc\checkmark$\checkmarke\checkmark$\checkmark^\checkmark\\checkmarkc\checkmarki\checkmarkr\checkmarkc\checkmark$\checkmarkx\checkmark$\checkmark^\checkmark\\checkmarkc\checkmarki\checkmarkr\checkmarkc\checkmark$\checkmarki\checkmark$\checkmark^\checkmark\\checkmarkc\checkmarki\checkmarkr\checkmarkc\checkmark$\checkmarko\checkmark$\checkmark^\checkmark\\checkmarkc\checkmarki\checkmarkr\checkmarkc\checkmark$\checkmarkn\checkmark$\checkmark^\checkmark\\checkmarkc\checkmarki\checkmarkr\checkmarkc\checkmark$\checkmark \checkmark$\checkmark^\checkmark\\checkmarkc\checkmarki\checkmarkr\checkmarkc\checkmark$\checkmarka\checkmark$\checkmark^\checkmark\\checkmarkc\checkmarki\checkmarkr\checkmarkc\checkmark$\checkmarkl\checkmark$\checkmark^\checkmark\\checkmarkc\checkmarki\checkmarkr\checkmarkc\checkmark$\checkmarkt\checkmark$\checkmark^\checkmark\\checkmarkc\checkmarki\checkmarkr\checkmarkc\checkmark$\checkmarke\checkmark$\checkmark^\checkmark\\checkmarkc\checkmarki\checkmarkr\checkmarkc\checkmark$\checkmarkr\checkmark$\checkmark^\checkmark\\checkmarkc\checkmarki\checkmarkr\checkmarkc\checkmark$\checkmarkn\checkmark$\checkmark^\checkmark\\checkmarkc\checkmarki\checkmarkr\checkmarkc\checkmark$\checkmarké\checkmark$\checkmark^\checkmark\\checkmarkc\checkmarki\checkmarkr\checkmarkc\checkmark$\checkmarke\checkmark$\checkmark^\checkmark\\checkmarkc\checkmarki\checkmarkr\checkmarkc\checkmark$\checkmark \checkmark$\checkmark^\checkmark\\checkmarkc\checkmarki\checkmarkr\checkmarkc\checkmark$\checkmark:\checkmark$\checkmark^\checkmark\\checkmarkc\checkmarki\checkmarkr\checkmarkc\checkmark$\checkmark \checkmark$\checkmark^\checkmark\\checkmarkc\checkmarki\checkmarkr\checkmarkc\checkmark$\checkmark$\checkmark$\checkmark^\checkmark\\checkmarkc\checkmarki\checkmarkr\checkmarkc\checkmark$\checkmark\\checkmark$\checkmark^\checkmark\\checkmarkc\checkmarki\checkmarkr\checkmarkc\checkmark$\checkmarks\checkmark$\checkmark^\checkmark\\checkmarkc\checkmarki\checkmarkr\checkmarkc\checkmark$\checkmarki\checkmark$\checkmark^\checkmark\\checkmarkc\checkmarki\checkmarkr\checkmarkc\checkmark$\checkmarkg\checkmark$\checkmark^\checkmark\\checkmarkc\checkmarki\checkmarkr\checkmarkc\checkmark$\checkmarkm\checkmark$\checkmark^\checkmark\\checkmarkc\checkmarki\checkmarkr\checkmarkc\checkmark$\checkmarka\checkmark$\checkmark^\checkmark\\checkmarkc\checkmarki\checkmarkr\checkmarkc\checkmark$\checkmark_\checkmark$\checkmark^\checkmark\\checkmarkc\checkmarki\checkmarkr\checkmarkc\checkmark$\checkmarkD\checkmark$\checkmark^\checkmark\\checkmarkc\checkmarki\checkmarkr\checkmarkc\checkmark$\checkmark \checkmark$\checkmark^\checkmark\\checkmarkc\checkmarki\checkmarkr\checkmarkc\checkmark$\checkmark=\checkmark$\checkmark^\checkmark\\checkmarkc\checkmarki\checkmarkr\checkmarkc\checkmark$\checkmark \checkmark$\checkmark^\checkmark\\checkmarkc\checkmarki\checkmarkr\checkmarkc\checkmark$\checkmark0\checkmark$\checkmark^\checkmark\\checkmarkc\checkmarki\checkmarkr\checkmarkc\checkmark$\checkmark{\checkmark$\checkmark^\checkmark\\checkmarkc\checkmarki\checkmarkr\checkmarkc\checkmark$\checkmark,\checkmark$\checkmark^\checkmark\\checkmarkc\checkmarki\checkmarkr\checkmarkc\checkmark$\checkmark}\checkmark$\checkmark^\checkmark\\checkmarkc\checkmarki\checkmarkr\checkmarkc\checkmark$\checkmark4\checkmark$\checkmark^\checkmark\\checkmarkc\checkmarki\checkmarkr\checkmarkc\checkmark$\checkmark3\checkmark$\checkmark^\checkmark\\checkmarkc\checkmarki\checkmarkr\checkmarkc\checkmark$\checkmark \checkmark$\checkmark^\checkmark\\checkmarkc\checkmarki\checkmarkr\checkmarkc\checkmark$\checkmark\\checkmark$\checkmark^\checkmark\\checkmarkc\checkmarki\checkmarkr\checkmarkc\checkmark$\checkmarkt\checkmark$\checkmark^\checkmark\\checkmarkc\checkmarki\checkmarkr\checkmarkc\checkmark$\checkmarki\checkmark$\checkmark^\checkmark\\checkmarkc\checkmarki\checkmarkr\checkmarkc\checkmark$\checkmarkm\checkmark$\checkmark^\checkmark\\checkmarkc\checkmarki\checkmarkr\checkmarkc\checkmark$\checkmarke\checkmark$\checkmark^\checkmark\\checkmarkc\checkmarki\checkmarkr\checkmarkc\checkmark$\checkmarks\checkmark$\checkmark^\checkmark\\checkmarkc\checkmarki\checkmarkr\checkmarkc\checkmark$\checkmark \checkmark$\checkmark^\checkmark\\checkmarkc\checkmarki\checkmarkr\checkmarkc\checkmark$\checkmarkR\checkmark$\checkmark^\checkmark\\checkmarkc\checkmarki\checkmarkr\checkmarkc\checkmark$\checkmark_\checkmark$\checkmark^\checkmark\\checkmarkc\checkmarki\checkmarkr\checkmarkc\checkmark$\checkmarkm\checkmark$\checkmark^\checkmark\\checkmarkc\checkmarki\checkmarkr\checkmarkc\checkmark$\checkmark \checkmark$\checkmark^\checkmark\\checkmarkc\checkmarki\checkmarkr\checkmarkc\checkmark$\checkmark=\checkmark$\checkmark^\checkmark\\checkmarkc\checkmarki\checkmarkr\checkmarkc\checkmark$\checkmark \checkmark$\checkmark^\checkmark\\checkmarkc\checkmarki\checkmarkr\checkmarkc\checkmark$\checkmark0\checkmark$\checkmark^\checkmark\\checkmarkc\checkmarki\checkmarkr\checkmarkc\checkmark$\checkmark{\checkmark$\checkmark^\checkmark\\checkmarkc\checkmarki\checkmarkr\checkmarkc\checkmark$\checkmark,\checkmark$\checkmark^\checkmark\\checkmarkc\checkmarki\checkmarkr\checkmarkc\checkmark$\checkmark}\checkmark$\checkmark^\checkmark\\checkmarkc\checkmarki\checkmarkr\checkmarkc\checkmark$\checkmark4\checkmark$\checkmark^\checkmark\\checkmarkc\checkmarki\checkmarkr\checkmarkc\checkmark$\checkmark3\checkmark$\checkmark^\checkmark\\checkmarkc\checkmarki\checkmarkr\checkmarkc\checkmark$\checkmark \checkmark$\checkmark^\checkmark\\checkmarkc\checkmarki\checkmarkr\checkmarkc\checkmark$\checkmark\\checkmark$\checkmark^\checkmark\\checkmarkc\checkmarki\checkmarkr\checkmarkc\checkmark$\checkmarkt\checkmark$\checkmark^\checkmark\\checkmarkc\checkmarki\checkmarkr\checkmarkc\checkmark$\checkmarki\checkmark$\checkmark^\checkmark\\checkmarkc\checkmarki\checkmarkr\checkmarkc\checkmark$\checkmarkm\checkmark$\checkmark^\checkmark\\checkmarkc\checkmarki\checkmarkr\checkmarkc\checkmark$\checkmarke\checkmark$\checkmark^\checkmark\\checkmarkc\checkmarki\checkmarkr\checkmarkc\checkmark$\checkmarks\checkmark$\checkmark^\checkmark\\checkmarkc\checkmarki\checkmarkr\checkmarkc\checkmark$\checkmark \checkmark$\checkmark^\checkmark\\checkmarkc\checkmarki\checkmarkr\checkmarkc\checkmark$\checkmark6\checkmark$\checkmark^\checkmark\\checkmarkc\checkmarki\checkmarkr\checkmarkc\checkmark$\checkmark0\checkmark$\checkmark^\checkmark\\checkmarkc\checkmarki\checkmarkr\checkmarkc\checkmark$\checkmark0\checkmark$\checkmark^\checkmark\\checkmarkc\checkmarki\checkmarkr\checkmarkc\checkmark$\checkmark \checkmark$\checkmark^\checkmark\\checkmarkc\checkmarki\checkmarkr\checkmarkc\checkmark$\checkmark=\checkmark$\checkmark^\checkmark\\checkmarkc\checkmarki\checkmarkr\checkmarkc\checkmark$\checkmark \checkmark$\checkmark^\checkmark\\checkmarkc\checkmarki\checkmarkr\checkmarkc\checkmark$\checkmark2\checkmark$\checkmark^\checkmark\\checkmarkc\checkmarki\checkmarkr\checkmarkc\checkmark$\checkmark5\checkmark$\checkmark^\checkmark\\checkmarkc\checkmarki\checkmarkr\checkmarkc\checkmark$\checkmark8\checkmark$\checkmark^\checkmark\\checkmarkc\checkmarki\checkmarkr\checkmarkc\checkmark$\checkmark$\checkmark$\checkmark^\checkmark\\checkmarkc\checkmarki\checkmarkr\checkmarkc\checkmark$\checkmark \checkmark$\checkmark^\checkmark\\checkmarkc\checkmarki\checkmarkr\checkmarkc\checkmark$\checkmarkM\checkmark$\checkmark^\checkmark\\checkmarkc\checkmarki\checkmarkr\checkmarkc\checkmark$\checkmarkP\checkmark$\checkmark^\checkmark\\checkmarkc\checkmarki\checkmarkr\checkmarkc\checkmark$\checkmarka\checkmark$\checkmark^\checkmark\\checkmarkc\checkmarki\checkmarkr\checkmarkc\checkmark$\checkmark
\checkmark$\checkmark^\checkmark\\checkmarkc\checkmarki\checkmarkr\checkmarkc\checkmark$\checkmark\\checkmark$\checkmark^\checkmark\\checkmarkc\checkmarki\checkmarkr\checkmarkc\checkmark$\checkmarke\checkmark$\checkmark^\checkmark\\checkmarkc\checkmarki\checkmarkr\checkmarkc\checkmark$\checkmarkn\checkmark$\checkmark^\checkmark\\checkmarkc\checkmarki\checkmarkr\checkmarkc\checkmark$\checkmarkd\checkmark$\checkmark^\checkmark\\checkmarkc\checkmarki\checkmarkr\checkmarkc\checkmark$\checkmark{\checkmark$\checkmark^\checkmark\\checkmarkc\checkmarki\checkmarkr\checkmarkc\checkmark$\checkmarki\checkmark$\checkmark^\checkmark\\checkmarkc\checkmarki\checkmarkr\checkmarkc\checkmark$\checkmarkt\checkmark$\checkmark^\checkmark\\checkmarkc\checkmarki\checkmarkr\checkmarkc\checkmark$\checkmarke\checkmark$\checkmark^\checkmark\\checkmarkc\checkmarki\checkmarkr\checkmarkc\checkmark$\checkmarkm\checkmark$\checkmark^\checkmark\\checkmarkc\checkmarki\checkmarkr\checkmarkc\checkmark$\checkmarki\checkmark$\checkmark^\checkmark\\checkmarkc\checkmarki\checkmarkr\checkmarkc\checkmark$\checkmarkz\checkmark$\checkmark^\checkmark\\checkmarkc\checkmarki\checkmarkr\checkmarkc\checkmark$\checkmarke\checkmark$\checkmark^\checkmark\\checkmarkc\checkmarki\checkmarkr\checkmarkc\checkmark$\checkmark}\checkmark$\checkmark^\checkmark\\checkmarkc\checkmarki\checkmarkr\checkmarkc\checkmark$\checkmark
\checkmark$\checkmark^\checkmark\\checkmarkc\checkmarki\checkmarkr\checkmarkc\checkmark$\checkmark
\checkmark$\checkmark^\checkmark\\checkmarkc\checkmarki\checkmarkr\checkmarkc\checkmark$\checkmark\\checkmark$\checkmark^\checkmark\\checkmarkc\checkmarki\checkmarkr\checkmarkc\checkmark$\checkmarks\checkmark$\checkmark^\checkmark\\checkmarkc\checkmarki\checkmarkr\checkmarkc\checkmark$\checkmarku\checkmark$\checkmark^\checkmark\\checkmarkc\checkmarki\checkmarkr\checkmarkc\checkmark$\checkmarkb\checkmark$\checkmark^\checkmark\\checkmarkc\checkmarki\checkmarkr\checkmarkc\checkmark$\checkmarks\checkmark$\checkmark^\checkmark\\checkmarkc\checkmarki\checkmarkr\checkmarkc\checkmark$\checkmarke\checkmark$\checkmark^\checkmark\\checkmarkc\checkmarki\checkmarkr\checkmarkc\checkmark$\checkmarkc\checkmark$\checkmark^\checkmark\\checkmarkc\checkmarki\checkmarkr\checkmarkc\checkmark$\checkmarkt\checkmark$\checkmark^\checkmark\\checkmarkc\checkmarki\checkmarkr\checkmarkc\checkmark$\checkmarki\checkmark$\checkmark^\checkmark\\checkmarkc\checkmarki\checkmarkr\checkmarkc\checkmark$\checkmarko\checkmark$\checkmark^\checkmark\\checkmarkc\checkmarki\checkmarkr\checkmarkc\checkmark$\checkmarkn\checkmark$\checkmark^\checkmark\\checkmarkc\checkmarki\checkmarkr\checkmarkc\checkmark$\checkmark{\checkmark$\checkmark^\checkmark\\checkmarkc\checkmarki\checkmarkr\checkmarkc\checkmark$\checkmarkC\checkmark$\checkmark^\checkmark\\checkmarkc\checkmarki\checkmarkr\checkmarkc\checkmark$\checkmarka\checkmark$\checkmark^\checkmark\\checkmarkc\checkmarki\checkmarkr\checkmarkc\checkmark$\checkmarkl\checkmark$\checkmark^\checkmark\\checkmarkc\checkmarki\checkmarkr\checkmarkc\checkmark$\checkmarkc\checkmark$\checkmark^\checkmark\\checkmarkc\checkmarki\checkmarkr\checkmarkc\checkmark$\checkmarku\checkmark$\checkmark^\checkmark\\checkmarkc\checkmarki\checkmarkr\checkmarkc\checkmark$\checkmarkl\checkmark$\checkmark^\checkmark\\checkmarkc\checkmarki\checkmarkr\checkmarkc\checkmark$\checkmark \checkmark$\checkmark^\checkmark\\checkmarkc\checkmarki\checkmarkr\checkmarkc\checkmark$\checkmarkd\checkmark$\checkmark^\checkmark\\checkmarkc\checkmarki\checkmarkr\checkmarkc\checkmark$\checkmarku\checkmark$\checkmark^\checkmark\\checkmarkc\checkmarki\checkmarkr\checkmarkc\checkmark$\checkmark \checkmark$\checkmark^\checkmark\\checkmarkc\checkmarki\checkmarkr\checkmarkc\checkmark$\checkmarkm\checkmark$\checkmark^\checkmark\\checkmarkc\checkmarki\checkmarkr\checkmarkc\checkmark$\checkmarko\checkmark$\checkmark^\checkmark\\checkmarkc\checkmarki\checkmarkr\checkmarkc\checkmark$\checkmarkm\checkmark$\checkmark^\checkmark\\checkmarkc\checkmarki\checkmarkr\checkmarkc\checkmark$\checkmarke\checkmark$\checkmark^\checkmark\\checkmarkc\checkmarki\checkmarkr\checkmarkc\checkmark$\checkmarkn\checkmark$\checkmark^\checkmark\\checkmarkc\checkmarki\checkmarkr\checkmarkc\checkmark$\checkmarkt\checkmark$\checkmark^\checkmark\\checkmarkc\checkmarki\checkmarkr\checkmarkc\checkmark$\checkmark \checkmark$\checkmark^\checkmark\\checkmarkc\checkmarki\checkmarkr\checkmarkc\checkmark$\checkmarkd\checkmark$\checkmark^\checkmark\\checkmarkc\checkmarki\checkmarkr\checkmarkc\checkmark$\checkmarke\checkmark$\checkmark^\checkmark\\checkmarkc\checkmarki\checkmarkr\checkmarkc\checkmark$\checkmark \checkmark$\checkmark^\checkmark\\checkmarkc\checkmarki\checkmarkr\checkmarkc\checkmark$\checkmarkr\checkmark$\checkmark^\checkmark\\checkmarkc\checkmarki\checkmarkr\checkmarkc\checkmark$\checkmarke\checkmark$\checkmark^\checkmark\\checkmarkc\checkmarki\checkmarkr\checkmarkc\checkmark$\checkmarkn\checkmark$\checkmark^\checkmark\\checkmarkc\checkmarki\checkmarkr\checkmarkc\checkmark$\checkmarkv\checkmark$\checkmark^\checkmark\\checkmarkc\checkmarki\checkmarkr\checkmarkc\checkmark$\checkmarke\checkmark$\checkmark^\checkmark\\checkmarkc\checkmarki\checkmarkr\checkmarkc\checkmark$\checkmarkr\checkmark$\checkmark^\checkmark\\checkmarkc\checkmarki\checkmarkr\checkmarkc\checkmark$\checkmarks\checkmark$\checkmark^\checkmark\\checkmarkc\checkmarki\checkmarkr\checkmarkc\checkmark$\checkmarke\checkmark$\checkmark^\checkmark\\checkmarkc\checkmarki\checkmarkr\checkmarkc\checkmark$\checkmarkm\checkmark$\checkmark^\checkmark\\checkmarkc\checkmarki\checkmarkr\checkmarkc\checkmark$\checkmarke\checkmark$\checkmark^\checkmark\\checkmarkc\checkmarki\checkmarkr\checkmarkc\checkmark$\checkmarkn\checkmark$\checkmark^\checkmark\\checkmarkc\checkmarki\checkmarkr\checkmarkc\checkmark$\checkmarkt\checkmark$\checkmark^\checkmark\\checkmarkc\checkmarki\checkmarkr\checkmarkc\checkmark$\checkmark}\checkmark$\checkmark^\checkmark\\checkmarkc\checkmarki\checkmarkr\checkmarkc\checkmark$\checkmark
\checkmark$\checkmark^\checkmark\\checkmarkc\checkmarki\checkmarkr\checkmarkc\checkmark$\checkmarkL\checkmark$\checkmark^\checkmark\\checkmarkc\checkmarki\checkmarkr\checkmarkc\checkmark$\checkmarka\checkmark$\checkmark^\checkmark\\checkmarkc\checkmarki\checkmarkr\checkmarkc\checkmark$\checkmark \checkmark$\checkmark^\checkmark\\checkmarkc\checkmarki\checkmarkr\checkmarkc\checkmark$\checkmarkf\checkmark$\checkmark^\checkmark\\checkmarkc\checkmarki\checkmarkr\checkmarkc\checkmark$\checkmarko\checkmark$\checkmark^\checkmark\\checkmarkc\checkmarki\checkmarkr\checkmarkc\checkmark$\checkmarkr\checkmark$\checkmark^\checkmark\\checkmarkc\checkmarki\checkmarkr\checkmarkc\checkmark$\checkmarkc\checkmark$\checkmark^\checkmark\\checkmarkc\checkmarki\checkmarkr\checkmarkc\checkmark$\checkmarke\checkmark$\checkmark^\checkmark\\checkmarkc\checkmarki\checkmarkr\checkmarkc\checkmark$\checkmark \checkmark$\checkmark^\checkmark\\checkmarkc\checkmarki\checkmarkr\checkmarkc\checkmark$\checkmarkd\checkmark$\checkmark^\checkmark\\checkmarkc\checkmarki\checkmarkr\checkmarkc\checkmark$\checkmarke\checkmark$\checkmark^\checkmark\\checkmarkc\checkmarki\checkmarkr\checkmarkc\checkmark$\checkmark \checkmark$\checkmark^\checkmark\\checkmarkc\checkmarki\checkmarkr\checkmarkc\checkmark$\checkmarkc\checkmark$\checkmark^\checkmark\\checkmarkc\checkmarki\checkmarkr\checkmarkc\checkmark$\checkmarko\checkmark$\checkmark^\checkmark\\checkmarkc\checkmarki\checkmarkr\checkmarkc\checkmark$\checkmarku\checkmark$\checkmark^\checkmark\\checkmarkc\checkmarki\checkmarkr\checkmarkc\checkmark$\checkmarkp\checkmark$\checkmark^\checkmark\\checkmarkc\checkmarki\checkmarkr\checkmarkc\checkmark$\checkmarke\checkmark$\checkmark^\checkmark\\checkmarkc\checkmarki\checkmarkr\checkmarkc\checkmark$\checkmark \checkmark$\checkmark^\checkmark\\checkmarkc\checkmarki\checkmarkr\checkmarkc\checkmark$\checkmark$\checkmark$\checkmark^\checkmark\\checkmarkc\checkmarki\checkmarkr\checkmarkc\checkmark$\checkmarkF\checkmark$\checkmark^\checkmark\\checkmarkc\checkmarki\checkmarkr\checkmarkc\checkmark$\checkmark_\checkmark$\checkmark^\checkmark\\checkmarkc\checkmarki\checkmarkr\checkmarkc\checkmark$\checkmark{\checkmark$\checkmark^\checkmark\\checkmarkc\checkmarki\checkmarkr\checkmarkc\checkmark$\checkmarkm\checkmark$\checkmark^\checkmark\\checkmarkc\checkmarki\checkmarkr\checkmarkc\checkmark$\checkmarka\checkmark$\checkmark^\checkmark\\checkmarkc\checkmarki\checkmarkr\checkmarkc\checkmark$\checkmarkx\checkmark$\checkmark^\checkmark\\checkmarkc\checkmarki\checkmarkr\checkmarkc\checkmark$\checkmark}\checkmark$\checkmark^\checkmark\\checkmarkc\checkmarki\checkmarkr\checkmarkc\checkmark$\checkmark \checkmark$\checkmark^\checkmark\\checkmarkc\checkmarki\checkmarkr\checkmarkc\checkmark$\checkmark=\checkmark$\checkmark^\checkmark\\checkmarkc\checkmarki\checkmarkr\checkmarkc\checkmark$\checkmark \checkmark$\checkmark^\checkmark\\checkmarkc\checkmarki\checkmarkr\checkmarkc\checkmark$\checkmark1\checkmark$\checkmark^\checkmark\\checkmarkc\checkmarki\checkmarkr\checkmarkc\checkmark$\checkmark0\checkmark$\checkmark^\checkmark\\checkmarkc\checkmarki\checkmarkr\checkmarkc\checkmark$\checkmark0\checkmark$\checkmark^\checkmark\\checkmarkc\checkmarki\checkmarkr\checkmarkc\checkmark$\checkmark$\checkmark$\checkmark^\checkmark\\checkmarkc\checkmarki\checkmarkr\checkmarkc\checkmark$\checkmark \checkmark$\checkmark^\checkmark\\checkmarkc\checkmarki\checkmarkr\checkmarkc\checkmark$\checkmarkN\checkmark$\checkmark^\checkmark\\checkmarkc\checkmarki\checkmarkr\checkmarkc\checkmark$\checkmark,\checkmark$\checkmark^\checkmark\\checkmarkc\checkmarki\checkmarkr\checkmarkc\checkmark$\checkmark \checkmark$\checkmark^\checkmark\\checkmarkc\checkmarki\checkmarkr\checkmarkc\checkmark$\checkmarka\checkmark$\checkmark^\checkmark\\checkmarkc\checkmarki\checkmarkr\checkmarkc\checkmark$\checkmarkp\checkmark$\checkmark^\checkmark\\checkmarkc\checkmarki\checkmarkr\checkmarkc\checkmark$\checkmarkp\checkmark$\checkmark^\checkmark\\checkmarkc\checkmarki\checkmarkr\checkmarkc\checkmark$\checkmarkl\checkmark$\checkmark^\checkmark\\checkmarkc\checkmarki\checkmarkr\checkmarkc\checkmark$\checkmarki\checkmark$\checkmark^\checkmark\\checkmarkc\checkmarki\checkmarkr\checkmarkc\checkmark$\checkmarkq\checkmark$\checkmark^\checkmark\\checkmarkc\checkmarki\checkmarkr\checkmarkc\checkmark$\checkmarku\checkmark$\checkmark^\checkmark\\checkmarkc\checkmarki\checkmarkr\checkmarkc\checkmark$\checkmarké\checkmark$\checkmark^\checkmark\\checkmarkc\checkmarki\checkmarkr\checkmarkc\checkmark$\checkmarke\checkmark$\checkmark^\checkmark\\checkmarkc\checkmarki\checkmarkr\checkmarkc\checkmark$\checkmark \checkmark$\checkmark^\checkmark\\checkmarkc\checkmarki\checkmarkr\checkmarkc\checkmark$\checkmarkà\checkmark$\checkmark^\checkmark\\checkmarkc\checkmarki\checkmarkr\checkmarkc\checkmark$\checkmark \checkmark$\checkmark^\checkmark\\checkmarkc\checkmarki\checkmarkr\checkmarkc\checkmark$\checkmarkl\checkmark$\checkmark^\checkmark\\checkmarkc\checkmarki\checkmarkr\checkmarkc\checkmark$\checkmarka\checkmark$\checkmark^\checkmark\\checkmarkc\checkmarki\checkmarkr\checkmarkc\checkmark$\checkmark \checkmark$\checkmark^\checkmark\\checkmarkc\checkmarki\checkmarkr\checkmarkc\checkmark$\checkmarkh\checkmark$\checkmark^\checkmark\\checkmarkc\checkmarki\checkmarkr\checkmarkc\checkmark$\checkmarka\checkmark$\checkmark^\checkmark\\checkmarkc\checkmarki\checkmarkr\checkmarkc\checkmark$\checkmarku\checkmark$\checkmark^\checkmark\\checkmarkc\checkmarki\checkmarkr\checkmarkc\checkmark$\checkmarkt\checkmark$\checkmark^\checkmark\\checkmarkc\checkmarki\checkmarkr\checkmarkc\checkmark$\checkmarke\checkmark$\checkmark^\checkmark\\checkmarkc\checkmarki\checkmarkr\checkmarkc\checkmark$\checkmarku\checkmark$\checkmark^\checkmark\\checkmarkc\checkmarki\checkmarkr\checkmarkc\checkmark$\checkmarkr\checkmark$\checkmark^\checkmark\\checkmarkc\checkmarki\checkmarkr\checkmarkc\checkmark$\checkmark \checkmark$\checkmark^\checkmark\\checkmarkc\checkmarki\checkmarkr\checkmarkc\checkmark$\checkmarkd\checkmark$\checkmark^\checkmark\\checkmarkc\checkmarki\checkmarkr\checkmarkc\checkmark$\checkmarku\checkmark$\checkmark^\checkmark\\checkmarkc\checkmarki\checkmarkr\checkmarkc\checkmark$\checkmark \checkmark$\checkmark^\checkmark\\checkmarkc\checkmarki\checkmarkr\checkmarkc\checkmark$\checkmarkp\checkmark$\checkmark^\checkmark\\checkmarkc\checkmarki\checkmarkr\checkmarkc\checkmark$\checkmarkl\checkmark$\checkmark^\checkmark\\checkmarkc\checkmarki\checkmarkr\checkmarkc\checkmark$\checkmarka\checkmark$\checkmark^\checkmark\\checkmarkc\checkmarki\checkmarkr\checkmarkc\checkmark$\checkmarkt\checkmark$\checkmark^\checkmark\\checkmarkc\checkmarki\checkmarkr\checkmarkc\checkmark$\checkmarke\checkmark$\checkmark^\checkmark\\checkmarkc\checkmarki\checkmarkr\checkmarkc\checkmark$\checkmarka\checkmark$\checkmark^\checkmark\\checkmarkc\checkmarki\checkmarkr\checkmarkc\checkmark$\checkmarku\checkmark$\checkmark^\checkmark\\checkmarkc\checkmarki\checkmarkr\checkmarkc\checkmark$\checkmark \checkmark$\checkmark^\checkmark\\checkmarkc\checkmarki\checkmarkr\checkmarkc\checkmark$\checkmark$\checkmark$\checkmark^\checkmark\\checkmarkc\checkmarki\checkmarkr\checkmarkc\checkmark$\checkmarkH\checkmark$\checkmark^\checkmark\\checkmarkc\checkmarki\checkmarkr\checkmarkc\checkmark$\checkmark \checkmark$\checkmark^\checkmark\\checkmarkc\checkmarki\checkmarkr\checkmarkc\checkmark$\checkmark=\checkmark$\checkmark^\checkmark\\checkmarkc\checkmarki\checkmarkr\checkmarkc\checkmark$\checkmark \checkmark$\checkmark^\checkmark\\checkmarkc\checkmarki\checkmarkr\checkmarkc\checkmark$\checkmark2\checkmark$\checkmark^\checkmark\\checkmarkc\checkmarki\checkmarkr\checkmarkc\checkmark$\checkmark0\checkmark$\checkmark^\checkmark\\checkmarkc\checkmarki\checkmarkr\checkmarkc\checkmark$\checkmark0\checkmark$\checkmark^\checkmark\\checkmarkc\checkmarki\checkmarkr\checkmarkc\checkmark$\checkmark$\checkmark$\checkmark^\checkmark\\checkmarkc\checkmarki\checkmarkr\checkmarkc\checkmark$\checkmark \checkmark$\checkmark^\checkmark\\checkmarkc\checkmarki\checkmarkr\checkmarkc\checkmark$\checkmarkm\checkmark$\checkmark^\checkmark\\checkmarkc\checkmarki\checkmarkr\checkmarkc\checkmark$\checkmarkm\checkmark$\checkmark^\checkmark\\checkmarkc\checkmarki\checkmarkr\checkmarkc\checkmark$\checkmark \checkmark$\checkmark^\checkmark\\checkmarkc\checkmarki\checkmarkr\checkmarkc\checkmark$\checkmarka\checkmark$\checkmark^\checkmark\\checkmarkc\checkmarki\checkmarkr\checkmarkc\checkmark$\checkmarku\checkmark$\checkmark^\checkmark\\checkmarkc\checkmarki\checkmarkr\checkmarkc\checkmark$\checkmark-\checkmark$\checkmark^\checkmark\\checkmarkc\checkmarki\checkmarkr\checkmarkc\checkmark$\checkmarkd\checkmark$\checkmark^\checkmark\\checkmarkc\checkmarki\checkmarkr\checkmarkc\checkmark$\checkmarke\checkmark$\checkmark^\checkmark\\checkmarkc\checkmarki\checkmarkr\checkmarkc\checkmark$\checkmarks\checkmark$\checkmark^\checkmark\\checkmarkc\checkmarki\checkmarkr\checkmarkc\checkmark$\checkmarks\checkmark$\checkmark^\checkmark\\checkmarkc\checkmarki\checkmarkr\checkmarkc\checkmark$\checkmarku\checkmark$\checkmark^\checkmark\\checkmarkc\checkmarki\checkmarkr\checkmarkc\checkmark$\checkmarks\checkmark$\checkmark^\checkmark\\checkmarkc\checkmarki\checkmarkr\checkmarkc\checkmark$\checkmark \checkmark$\checkmark^\checkmark\\checkmarkc\checkmarki\checkmarkr\checkmarkc\checkmark$\checkmarkd\checkmark$\checkmark^\checkmark\\checkmarkc\checkmarki\checkmarkr\checkmarkc\checkmark$\checkmarke\checkmark$\checkmark^\checkmark\\checkmarkc\checkmarki\checkmarkr\checkmarkc\checkmark$\checkmarks\checkmark$\checkmark^\checkmark\\checkmarkc\checkmarki\checkmarkr\checkmarkc\checkmark$\checkmark \checkmark$\checkmark^\checkmark\\checkmarkc\checkmarki\checkmarkr\checkmarkc\checkmark$\checkmarkr\checkmark$\checkmark^\checkmark\\checkmarkc\checkmarki\checkmarkr\checkmarkc\checkmark$\checkmarko\checkmark$\checkmark^\checkmark\\checkmarkc\checkmarki\checkmarkr\checkmarkc\checkmark$\checkmarku\checkmark$\checkmark^\checkmark\\checkmarkc\checkmarki\checkmarkr\checkmarkc\checkmark$\checkmarkl\checkmark$\checkmark^\checkmark\\checkmarkc\checkmarki\checkmarkr\checkmarkc\checkmark$\checkmarke\checkmark$\checkmark^\checkmark\\checkmarkc\checkmarki\checkmarkr\checkmarkc\checkmark$\checkmarkm\checkmark$\checkmark^\checkmark\\checkmarkc\checkmarki\checkmarkr\checkmarkc\checkmark$\checkmarke\checkmark$\checkmark^\checkmark\\checkmarkc\checkmarki\checkmarkr\checkmarkc\checkmark$\checkmarkn\checkmark$\checkmark^\checkmark\\checkmarkc\checkmarki\checkmarkr\checkmarkc\checkmark$\checkmarkt\checkmark$\checkmark^\checkmark\\checkmarkc\checkmarki\checkmarkr\checkmarkc\checkmark$\checkmarks\checkmark$\checkmark^\checkmark\\checkmarkc\checkmarki\checkmarkr\checkmarkc\checkmark$\checkmark,\checkmark$\checkmark^\checkmark\\checkmarkc\checkmarki\checkmarkr\checkmarkc\checkmark$\checkmark \checkmark$\checkmark^\checkmark\\checkmarkc\checkmarki\checkmarkr\checkmarkc\checkmark$\checkmarkc\checkmark$\checkmark^\checkmark\\checkmarkc\checkmarki\checkmarkr\checkmarkc\checkmark$\checkmarkr\checkmark$\checkmark^\checkmark\\checkmarkc\checkmarki\checkmarkr\checkmarkc\checkmark$\checkmarké\checkmark$\checkmark^\checkmark\\checkmarkc\checkmarki\checkmarkr\checkmarkc\checkmark$\checkmarke\checkmark$\checkmark^\checkmark\\checkmarkc\checkmarki\checkmarkr\checkmarkc\checkmark$\checkmark \checkmark$\checkmark^\checkmark\\checkmarkc\checkmarki\checkmarkr\checkmarkc\checkmark$\checkmarku\checkmark$\checkmark^\checkmark\\checkmarkc\checkmarki\checkmarkr\checkmarkc\checkmark$\checkmarkn\checkmark$\checkmark^\checkmark\\checkmarkc\checkmarki\checkmarkr\checkmarkc\checkmark$\checkmark \checkmark$\checkmark^\checkmark\\checkmarkc\checkmarki\checkmarkr\checkmarkc\checkmark$\checkmarkm\checkmark$\checkmark^\checkmark\\checkmarkc\checkmarki\checkmarkr\checkmarkc\checkmark$\checkmarko\checkmark$\checkmark^\checkmark\\checkmarkc\checkmarki\checkmarkr\checkmarkc\checkmark$\checkmarkm\checkmark$\checkmark^\checkmark\\checkmarkc\checkmarki\checkmarkr\checkmarkc\checkmark$\checkmarke\checkmark$\checkmark^\checkmark\\checkmarkc\checkmarki\checkmarkr\checkmarkc\checkmark$\checkmarkn\checkmark$\checkmark^\checkmark\\checkmarkc\checkmarki\checkmarkr\checkmarkc\checkmark$\checkmarkt\checkmark$\checkmark^\checkmark\\checkmarkc\checkmarki\checkmarkr\checkmarkc\checkmark$\checkmark \checkmark$\checkmark^\checkmark\\checkmarkc\checkmarki\checkmarkr\checkmarkc\checkmark$\checkmarkd\checkmark$\checkmark^\checkmark\\checkmarkc\checkmarki\checkmarkr\checkmarkc\checkmark$\checkmarke\checkmark$\checkmark^\checkmark\\checkmarkc\checkmarki\checkmarkr\checkmarkc\checkmark$\checkmark \checkmark$\checkmark^\checkmark\\checkmarkc\checkmarki\checkmarkr\checkmarkc\checkmark$\checkmarkr\checkmark$\checkmark^\checkmark\\checkmarkc\checkmarki\checkmarkr\checkmarkc\checkmark$\checkmarke\checkmark$\checkmark^\checkmark\\checkmarkc\checkmarki\checkmarkr\checkmarkc\checkmark$\checkmarkn\checkmark$\checkmark^\checkmark\\checkmarkc\checkmarki\checkmarkr\checkmarkc\checkmark$\checkmarkv\checkmark$\checkmark^\checkmark\\checkmarkc\checkmarki\checkmarkr\checkmarkc\checkmark$\checkmarke\checkmark$\checkmark^\checkmark\\checkmarkc\checkmarki\checkmarkr\checkmarkc\checkmark$\checkmarkr\checkmark$\checkmark^\checkmark\\checkmarkc\checkmarki\checkmarkr\checkmarkc\checkmark$\checkmarks\checkmark$\checkmark^\checkmark\\checkmarkc\checkmarki\checkmarkr\checkmarkc\checkmark$\checkmarke\checkmark$\checkmark^\checkmark\\checkmarkc\checkmarki\checkmarkr\checkmarkc\checkmark$\checkmarkm\checkmark$\checkmark^\checkmark\\checkmarkc\checkmarki\checkmarkr\checkmarkc\checkmark$\checkmarke\checkmark$\checkmark^\checkmark\\checkmarkc\checkmarki\checkmarkr\checkmarkc\checkmark$\checkmarkn\checkmark$\checkmark^\checkmark\\checkmarkc\checkmarki\checkmarkr\checkmarkc\checkmark$\checkmarkt\checkmark$\checkmark^\checkmark\\checkmarkc\checkmarki\checkmarkr\checkmarkc\checkmark$\checkmark \checkmark$\checkmark^\checkmark\\checkmarkc\checkmarki\checkmarkr\checkmarkc\checkmark$\checkmark:\checkmark$\checkmark^\checkmark\\checkmarkc\checkmarki\checkmarkr\checkmarkc\checkmark$\checkmark
\checkmark$\checkmark^\checkmark\\checkmarkc\checkmarki\checkmarkr\checkmarkc\checkmark$\checkmark\\checkmark$\checkmark^\checkmark\\checkmarkc\checkmarki\checkmarkr\checkmarkc\checkmark$\checkmarkb\checkmark$\checkmark^\checkmark\\checkmarkc\checkmarki\checkmarkr\checkmarkc\checkmark$\checkmarke\checkmark$\checkmark^\checkmark\\checkmarkc\checkmarki\checkmarkr\checkmarkc\checkmark$\checkmarkg\checkmark$\checkmark^\checkmark\\checkmarkc\checkmarki\checkmarkr\checkmarkc\checkmark$\checkmarki\checkmark$\checkmark^\checkmark\\checkmarkc\checkmarki\checkmarkr\checkmarkc\checkmark$\checkmarkn\checkmark$\checkmark^\checkmark\\checkmarkc\checkmarki\checkmarkr\checkmarkc\checkmark$\checkmark{\checkmark$\checkmark^\checkmark\\checkmarkc\checkmarki\checkmarkr\checkmarkc\checkmark$\checkmarke\checkmark$\checkmark^\checkmark\\checkmarkc\checkmarki\checkmarkr\checkmarkc\checkmark$\checkmarkq\checkmark$\checkmark^\checkmark\\checkmarkc\checkmarki\checkmarkr\checkmarkc\checkmark$\checkmarku\checkmark$\checkmark^\checkmark\\checkmarkc\checkmarki\checkmarkr\checkmarkc\checkmark$\checkmarka\checkmark$\checkmark^\checkmark\\checkmarkc\checkmarki\checkmarkr\checkmarkc\checkmark$\checkmarkt\checkmark$\checkmark^\checkmark\\checkmarkc\checkmarki\checkmarkr\checkmarkc\checkmark$\checkmarki\checkmark$\checkmark^\checkmark\\checkmarkc\checkmarki\checkmarkr\checkmarkc\checkmark$\checkmarko\checkmark$\checkmark^\checkmark\\checkmarkc\checkmarki\checkmarkr\checkmarkc\checkmark$\checkmarkn\checkmark$\checkmark^\checkmark\\checkmarkc\checkmarki\checkmarkr\checkmarkc\checkmark$\checkmark}\checkmark$\checkmark^\checkmark\\checkmarkc\checkmarki\checkmarkr\checkmarkc\checkmark$\checkmark
\checkmark$\checkmark^\checkmark\\checkmarkc\checkmarki\checkmarkr\checkmarkc\checkmark$\checkmark \checkmark$\checkmark^\checkmark\\checkmarkc\checkmarki\checkmarkr\checkmarkc\checkmark$\checkmark \checkmark$\checkmark^\checkmark\\checkmarkc\checkmarki\checkmarkr\checkmarkc\checkmark$\checkmark \checkmark$\checkmark^\checkmark\\checkmarkc\checkmarki\checkmarkr\checkmarkc\checkmark$\checkmark \checkmark$\checkmark^\checkmark\\checkmarkc\checkmarki\checkmarkr\checkmarkc\checkmark$\checkmarkM\checkmark$\checkmark^\checkmark\\checkmarkc\checkmarki\checkmarkr\checkmarkc\checkmark$\checkmark_\checkmark$\checkmark^\checkmark\\checkmarkc\checkmarki\checkmarkr\checkmarkc\checkmark$\checkmark{\checkmark$\checkmark^\checkmark\\checkmarkc\checkmarki\checkmarkr\checkmarkc\checkmark$\checkmarkr\checkmark$\checkmark^\checkmark\\checkmarkc\checkmarki\checkmarkr\checkmarkc\checkmark$\checkmarke\checkmark$\checkmark^\checkmark\\checkmarkc\checkmarki\checkmarkr\checkmarkc\checkmark$\checkmarkn\checkmark$\checkmark^\checkmark\\checkmarkc\checkmarki\checkmarkr\checkmarkc\checkmark$\checkmarkv\checkmark$\checkmark^\checkmark\\checkmarkc\checkmarki\checkmarkr\checkmarkc\checkmark$\checkmark}\checkmark$\checkmark^\checkmark\\checkmarkc\checkmarki\checkmarkr\checkmarkc\checkmark$\checkmark \checkmark$\checkmark^\checkmark\\checkmarkc\checkmarki\checkmarkr\checkmarkc\checkmark$\checkmark=\checkmark$\checkmark^\checkmark\\checkmarkc\checkmarki\checkmarkr\checkmarkc\checkmark$\checkmark \checkmark$\checkmark^\checkmark\\checkmarkc\checkmarki\checkmarkr\checkmarkc\checkmark$\checkmarkF\checkmark$\checkmark^\checkmark\\checkmarkc\checkmarki\checkmarkr\checkmarkc\checkmark$\checkmark_\checkmark$\checkmark^\checkmark\\checkmarkc\checkmarki\checkmarkr\checkmarkc\checkmark$\checkmark{\checkmark$\checkmark^\checkmark\\checkmarkc\checkmarki\checkmarkr\checkmarkc\checkmark$\checkmarkm\checkmark$\checkmark^\checkmark\\checkmarkc\checkmarki\checkmarkr\checkmarkc\checkmark$\checkmarka\checkmark$\checkmark^\checkmark\\checkmarkc\checkmarki\checkmarkr\checkmarkc\checkmark$\checkmarkx\checkmark$\checkmark^\checkmark\\checkmarkc\checkmarki\checkmarkr\checkmarkc\checkmark$\checkmark}\checkmark$\checkmark^\checkmark\\checkmarkc\checkmarki\checkmarkr\checkmarkc\checkmark$\checkmark \checkmark$\checkmark^\checkmark\\checkmarkc\checkmarki\checkmarkr\checkmarkc\checkmark$\checkmark\\checkmark$\checkmark^\checkmark\\checkmarkc\checkmarki\checkmarkr\checkmarkc\checkmark$\checkmarkt\checkmark$\checkmark^\checkmark\\checkmarkc\checkmarki\checkmarkr\checkmarkc\checkmark$\checkmarki\checkmark$\checkmark^\checkmark\\checkmarkc\checkmarki\checkmarkr\checkmarkc\checkmark$\checkmarkm\checkmark$\checkmark^\checkmark\\checkmarkc\checkmarki\checkmarkr\checkmarkc\checkmark$\checkmarke\checkmark$\checkmark^\checkmark\\checkmarkc\checkmarki\checkmarkr\checkmarkc\checkmark$\checkmarks\checkmark$\checkmark^\checkmark\\checkmarkc\checkmarki\checkmarkr\checkmarkc\checkmark$\checkmark \checkmark$\checkmark^\checkmark\\checkmarkc\checkmarki\checkmarkr\checkmarkc\checkmark$\checkmarkH\checkmark$\checkmark^\checkmark\\checkmarkc\checkmarki\checkmarkr\checkmarkc\checkmark$\checkmark \checkmark$\checkmark^\checkmark\\checkmarkc\checkmarki\checkmarkr\checkmarkc\checkmark$\checkmark=\checkmark$\checkmark^\checkmark\\checkmarkc\checkmarki\checkmarkr\checkmarkc\checkmark$\checkmark \checkmark$\checkmark^\checkmark\\checkmarkc\checkmarki\checkmarkr\checkmarkc\checkmark$\checkmark1\checkmark$\checkmark^\checkmark\\checkmarkc\checkmarki\checkmarkr\checkmarkc\checkmark$\checkmark0\checkmark$\checkmark^\checkmark\\checkmarkc\checkmarki\checkmarkr\checkmarkc\checkmark$\checkmark0\checkmark$\checkmark^\checkmark\\checkmarkc\checkmarki\checkmarkr\checkmarkc\checkmark$\checkmark \checkmark$\checkmark^\checkmark\\checkmarkc\checkmarki\checkmarkr\checkmarkc\checkmark$\checkmark\\checkmark$\checkmark^\checkmark\\checkmarkc\checkmarki\checkmarkr\checkmarkc\checkmark$\checkmarkt\checkmark$\checkmark^\checkmark\\checkmarkc\checkmarki\checkmarkr\checkmarkc\checkmark$\checkmarki\checkmark$\checkmark^\checkmark\\checkmarkc\checkmarki\checkmarkr\checkmarkc\checkmark$\checkmarkm\checkmark$\checkmark^\checkmark\\checkmarkc\checkmarki\checkmarkr\checkmarkc\checkmark$\checkmarke\checkmark$\checkmark^\checkmark\\checkmarkc\checkmarki\checkmarkr\checkmarkc\checkmark$\checkmarks\checkmark$\checkmark^\checkmark\\checkmarkc\checkmarki\checkmarkr\checkmarkc\checkmark$\checkmark \checkmark$\checkmark^\checkmark\\checkmarkc\checkmarki\checkmarkr\checkmarkc\checkmark$\checkmark0\checkmark$\checkmark^\checkmark\\checkmarkc\checkmarki\checkmarkr\checkmarkc\checkmark$\checkmark{\checkmark$\checkmark^\checkmark\\checkmarkc\checkmarki\checkmarkr\checkmarkc\checkmark$\checkmark,\checkmark$\checkmark^\checkmark\\checkmarkc\checkmarki\checkmarkr\checkmarkc\checkmark$\checkmark}\checkmark$\checkmark^\checkmark\\checkmarkc\checkmarki\checkmarkr\checkmarkc\checkmark$\checkmark2\checkmark$\checkmark^\checkmark\\checkmarkc\checkmarki\checkmarkr\checkmarkc\checkmark$\checkmark0\checkmark$\checkmark^\checkmark\\checkmarkc\checkmarki\checkmarkr\checkmarkc\checkmark$\checkmark0\checkmark$\checkmark^\checkmark\\checkmarkc\checkmarki\checkmarkr\checkmarkc\checkmark$\checkmark \checkmark$\checkmark^\checkmark\\checkmarkc\checkmarki\checkmarkr\checkmarkc\checkmark$\checkmark=\checkmark$\checkmark^\checkmark\\checkmarkc\checkmarki\checkmarkr\checkmarkc\checkmark$\checkmark \checkmark$\checkmark^\checkmark\\checkmarkc\checkmarki\checkmarkr\checkmarkc\checkmark$\checkmark\\checkmark$\checkmark^\checkmark\\checkmarkc\checkmarki\checkmarkr\checkmarkc\checkmark$\checkmarkb\checkmark$\checkmark^\checkmark\\checkmarkc\checkmarki\checkmarkr\checkmarkc\checkmark$\checkmarko\checkmark$\checkmark^\checkmark\\checkmarkc\checkmarki\checkmarkr\checkmarkc\checkmark$\checkmarkx\checkmark$\checkmark^\checkmark\\checkmarkc\checkmarki\checkmarkr\checkmarkc\checkmark$\checkmarke\checkmark$\checkmark^\checkmark\\checkmarkc\checkmarki\checkmarkr\checkmarkc\checkmark$\checkmarkd\checkmark$\checkmark^\checkmark\\checkmarkc\checkmarki\checkmarkr\checkmarkc\checkmark$\checkmark{\checkmark$\checkmark^\checkmark\\checkmarkc\checkmarki\checkmarkr\checkmarkc\checkmark$\checkmark2\checkmark$\checkmark^\checkmark\\checkmarkc\checkmarki\checkmarkr\checkmarkc\checkmark$\checkmark0\checkmark$\checkmark^\checkmark\\checkmarkc\checkmarki\checkmarkr\checkmarkc\checkmark$\checkmark \checkmark$\checkmark^\checkmark\\checkmarkc\checkmarki\checkmarkr\checkmarkc\checkmark$\checkmark\\checkmark$\checkmark^\checkmark\\checkmarkc\checkmarki\checkmarkr\checkmarkc\checkmark$\checkmarkt\checkmark$\checkmark^\checkmark\\checkmarkc\checkmarki\checkmarkr\checkmarkc\checkmark$\checkmarke\checkmark$\checkmark^\checkmark\\checkmarkc\checkmarki\checkmarkr\checkmarkc\checkmark$\checkmarkx\checkmark$\checkmark^\checkmark\\checkmarkc\checkmarki\checkmarkr\checkmarkc\checkmark$\checkmarkt\checkmark$\checkmark^\checkmark\\checkmarkc\checkmarki\checkmarkr\checkmarkc\checkmark$\checkmark{\checkmark$\checkmark^\checkmark\\checkmarkc\checkmarki\checkmarkr\checkmarkc\checkmark$\checkmark \checkmark$\checkmark^\checkmark\\checkmarkc\checkmarki\checkmarkr\checkmarkc\checkmark$\checkmarkN\checkmark$\checkmark^\checkmark\\checkmarkc\checkmarki\checkmarkr\checkmarkc\checkmark$\checkmark$\checkmark$\checkmark^\checkmark\\checkmarkc\checkmarki\checkmarkr\checkmarkc\checkmark$\checkmark\\checkmark$\checkmark^\checkmark\\checkmarkc\checkmarki\checkmarkr\checkmarkc\checkmark$\checkmarkc\checkmark$\checkmark^\checkmark\\checkmarkc\checkmarki\checkmarkr\checkmarkc\checkmark$\checkmarkd\checkmark$\checkmark^\checkmark\\checkmarkc\checkmarki\checkmarkr\checkmarkc\checkmark$\checkmarko\checkmark$\checkmark^\checkmark\\checkmarkc\checkmarki\checkmarkr\checkmarkc\checkmark$\checkmarkt\checkmark$\checkmark^\checkmark\\checkmarkc\checkmarki\checkmarkr\checkmarkc\checkmark$\checkmark$\checkmark$\checkmark^\checkmark\\checkmarkc\checkmarki\checkmarkr\checkmarkc\checkmark$\checkmarkm\checkmark$\checkmark^\checkmark\\checkmarkc\checkmarki\checkmarkr\checkmarkc\checkmark$\checkmark}\checkmark$\checkmark^\checkmark\\checkmarkc\checkmarki\checkmarkr\checkmarkc\checkmark$\checkmark}\checkmark$\checkmark^\checkmark\\checkmarkc\checkmarki\checkmarkr\checkmarkc\checkmark$\checkmark
\checkmark$\checkmark^\checkmark\\checkmarkc\checkmarki\checkmarkr\checkmarkc\checkmark$\checkmark \checkmark$\checkmark^\checkmark\\checkmarkc\checkmarki\checkmarkr\checkmarkc\checkmark$\checkmark \checkmark$\checkmark^\checkmark\\checkmarkc\checkmarki\checkmarkr\checkmarkc\checkmark$\checkmark \checkmark$\checkmark^\checkmark\\checkmarkc\checkmarki\checkmarkr\checkmarkc\checkmark$\checkmark \checkmark$\checkmark^\checkmark\\checkmarkc\checkmarki\checkmarkr\checkmarkc\checkmark$\checkmark\\checkmark$\checkmark^\checkmark\\checkmarkc\checkmarki\checkmarkr\checkmarkc\checkmark$\checkmarkl\checkmark$\checkmark^\checkmark\\checkmarkc\checkmarki\checkmarkr\checkmarkc\checkmark$\checkmarka\checkmark$\checkmark^\checkmark\\checkmarkc\checkmarki\checkmarkr\checkmarkc\checkmark$\checkmarkb\checkmark$\checkmark^\checkmark\\checkmarkc\checkmarki\checkmarkr\checkmarkc\checkmark$\checkmarke\checkmark$\checkmark^\checkmark\\checkmarkc\checkmarki\checkmarkr\checkmarkc\checkmark$\checkmarkl\checkmark$\checkmark^\checkmark\\checkmarkc\checkmarki\checkmarkr\checkmarkc\checkmark$\checkmark{\checkmark$\checkmark^\checkmark\\checkmarkc\checkmarki\checkmarkr\checkmarkc\checkmark$\checkmarke\checkmark$\checkmark^\checkmark\\checkmarkc\checkmarki\checkmarkr\checkmarkc\checkmark$\checkmarkq\checkmark$\checkmark^\checkmark\\checkmarkc\checkmarki\checkmarkr\checkmarkc\checkmark$\checkmark:\checkmark$\checkmark^\checkmark\\checkmarkc\checkmarki\checkmarkr\checkmarkc\checkmark$\checkmarkm\checkmark$\checkmark^\checkmark\\checkmarkc\checkmarki\checkmarkr\checkmarkc\checkmark$\checkmarko\checkmark$\checkmark^\checkmark\\checkmarkc\checkmarki\checkmarkr\checkmarkc\checkmark$\checkmarkm\checkmark$\checkmark^\checkmark\\checkmarkc\checkmarki\checkmarkr\checkmarkc\checkmark$\checkmarke\checkmark$\checkmark^\checkmark\\checkmarkc\checkmarki\checkmarkr\checkmarkc\checkmark$\checkmarkn\checkmark$\checkmark^\checkmark\\checkmarkc\checkmarki\checkmarkr\checkmarkc\checkmark$\checkmarkt\checkmark$\checkmark^\checkmark\\checkmarkc\checkmarki\checkmarkr\checkmarkc\checkmark$\checkmark_\checkmark$\checkmark^\checkmark\\checkmarkc\checkmarki\checkmarkr\checkmarkc\checkmark$\checkmarkr\checkmark$\checkmark^\checkmark\\checkmarkc\checkmarki\checkmarkr\checkmarkc\checkmark$\checkmarke\checkmark$\checkmark^\checkmark\\checkmarkc\checkmarki\checkmarkr\checkmarkc\checkmark$\checkmarkn\checkmark$\checkmark^\checkmark\\checkmarkc\checkmarki\checkmarkr\checkmarkc\checkmark$\checkmarkv\checkmark$\checkmark^\checkmark\\checkmarkc\checkmarki\checkmarkr\checkmarkc\checkmark$\checkmarke\checkmark$\checkmark^\checkmark\\checkmarkc\checkmarki\checkmarkr\checkmarkc\checkmark$\checkmarkr\checkmark$\checkmark^\checkmark\\checkmarkc\checkmarki\checkmarkr\checkmarkc\checkmark$\checkmarks\checkmark$\checkmark^\checkmark\\checkmarkc\checkmarki\checkmarkr\checkmarkc\checkmark$\checkmarke\checkmark$\checkmark^\checkmark\\checkmarkc\checkmarki\checkmarkr\checkmarkc\checkmark$\checkmarkm\checkmark$\checkmark^\checkmark\\checkmarkc\checkmarki\checkmarkr\checkmarkc\checkmark$\checkmarke\checkmark$\checkmark^\checkmark\\checkmarkc\checkmarki\checkmarkr\checkmarkc\checkmark$\checkmarkn\checkmark$\checkmark^\checkmark\\checkmarkc\checkmarki\checkmarkr\checkmarkc\checkmark$\checkmarkt\checkmark$\checkmark^\checkmark\\checkmarkc\checkmarki\checkmarkr\checkmarkc\checkmark$\checkmark}\checkmark$\checkmark^\checkmark\\checkmarkc\checkmarki\checkmarkr\checkmarkc\checkmark$\checkmark
\checkmark$\checkmark^\checkmark\\checkmarkc\checkmarki\checkmarkr\checkmarkc\checkmark$\checkmark\\checkmark$\checkmark^\checkmark\\checkmarkc\checkmarki\checkmarkr\checkmarkc\checkmark$\checkmarke\checkmark$\checkmark^\checkmark\\checkmarkc\checkmarki\checkmarkr\checkmarkc\checkmark$\checkmarkn\checkmark$\checkmark^\checkmark\\checkmarkc\checkmarki\checkmarkr\checkmarkc\checkmark$\checkmarkd\checkmark$\checkmark^\checkmark\\checkmarkc\checkmarki\checkmarkr\checkmarkc\checkmark$\checkmark{\checkmark$\checkmark^\checkmark\\checkmarkc\checkmarki\checkmarkr\checkmarkc\checkmark$\checkmarke\checkmark$\checkmark^\checkmark\\checkmarkc\checkmarki\checkmarkr\checkmarkc\checkmark$\checkmarkq\checkmark$\checkmark^\checkmark\\checkmarkc\checkmarki\checkmarkr\checkmarkc\checkmark$\checkmarku\checkmark$\checkmark^\checkmark\\checkmarkc\checkmarki\checkmarkr\checkmarkc\checkmark$\checkmarka\checkmark$\checkmark^\checkmark\\checkmarkc\checkmarki\checkmarkr\checkmarkc\checkmark$\checkmarkt\checkmark$\checkmark^\checkmark\\checkmarkc\checkmarki\checkmarkr\checkmarkc\checkmark$\checkmarki\checkmark$\checkmark^\checkmark\\checkmarkc\checkmarki\checkmarkr\checkmarkc\checkmark$\checkmarko\checkmark$\checkmark^\checkmark\\checkmarkc\checkmarki\checkmarkr\checkmarkc\checkmark$\checkmarkn\checkmark$\checkmark^\checkmark\\checkmarkc\checkmarki\checkmarkr\checkmarkc\checkmark$\checkmark}\checkmark$\checkmark^\checkmark\\checkmarkc\checkmarki\checkmarkr\checkmarkc\checkmark$\checkmark
\checkmark$\checkmark^\checkmark\\checkmarkc\checkmarki\checkmarkr\checkmarkc\checkmark$\checkmark
\checkmark$\checkmark^\checkmark\\checkmarkc\checkmarki\checkmarkr\checkmarkc\checkmark$\checkmark\\checkmark$\checkmark^\checkmark\\checkmarkc\checkmarki\checkmarkr\checkmarkc\checkmark$\checkmarkb\checkmark$\checkmark^\checkmark\\checkmarkc\checkmarki\checkmarkr\checkmarkc\checkmark$\checkmarke\checkmark$\checkmark^\checkmark\\checkmarkc\checkmarki\checkmarkr\checkmarkc\checkmark$\checkmarkg\checkmark$\checkmark^\checkmark\\checkmarkc\checkmarki\checkmarkr\checkmarkc\checkmark$\checkmarki\checkmark$\checkmark^\checkmark\\checkmarkc\checkmarki\checkmarkr\checkmarkc\checkmark$\checkmarkn\checkmark$\checkmark^\checkmark\\checkmarkc\checkmarki\checkmarkr\checkmarkc\checkmark$\checkmark{\checkmark$\checkmark^\checkmark\\checkmarkc\checkmarki\checkmarkr\checkmarkc\checkmark$\checkmarkf\checkmark$\checkmark^\checkmark\\checkmarkc\checkmarki\checkmarkr\checkmarkc\checkmark$\checkmarki\checkmark$\checkmark^\checkmark\\checkmarkc\checkmarki\checkmarkr\checkmarkc\checkmark$\checkmarkg\checkmark$\checkmark^\checkmark\\checkmarkc\checkmarki\checkmarkr\checkmarkc\checkmark$\checkmarku\checkmark$\checkmark^\checkmark\\checkmarkc\checkmarki\checkmarkr\checkmarkc\checkmark$\checkmarkr\checkmark$\checkmark^\checkmark\\checkmarkc\checkmarki\checkmarkr\checkmarkc\checkmark$\checkmarke\checkmark$\checkmark^\checkmark\\checkmarkc\checkmarki\checkmarkr\checkmarkc\checkmark$\checkmark}\checkmark$\checkmark^\checkmark\\checkmarkc\checkmarki\checkmarkr\checkmarkc\checkmark$\checkmark[\checkmark$\checkmark^\checkmark\\checkmarkc\checkmarki\checkmarkr\checkmarkc\checkmark$\checkmarkh\checkmark$\checkmark^\checkmark\\checkmarkc\checkmarki\checkmarkr\checkmarkc\checkmark$\checkmarkt\checkmark$\checkmark^\checkmark\\checkmarkc\checkmarki\checkmarkr\checkmarkc\checkmark$\checkmarkb\checkmark$\checkmark^\checkmark\\checkmarkc\checkmarki\checkmarkr\checkmarkc\checkmark$\checkmarkp\checkmark$\checkmark^\checkmark\\checkmarkc\checkmarki\checkmarkr\checkmarkc\checkmark$\checkmark]\checkmark$\checkmark^\checkmark\\checkmarkc\checkmarki\checkmarkr\checkmarkc\checkmark$\checkmark
\checkmark$\checkmark^\checkmark\\checkmarkc\checkmarki\checkmarkr\checkmarkc\checkmark$\checkmark \checkmark$\checkmark^\checkmark\\checkmarkc\checkmarki\checkmarkr\checkmarkc\checkmark$\checkmark \checkmark$\checkmark^\checkmark\\checkmarkc\checkmarki\checkmarkr\checkmarkc\checkmark$\checkmark \checkmark$\checkmark^\checkmark\\checkmarkc\checkmarki\checkmarkr\checkmarkc\checkmark$\checkmark \checkmark$\checkmark^\checkmark\\checkmarkc\checkmarki\checkmarkr\checkmarkc\checkmark$\checkmark\\checkmark$\checkmark^\checkmark\\checkmarkc\checkmarki\checkmarkr\checkmarkc\checkmark$\checkmarkc\checkmark$\checkmark^\checkmark\\checkmarkc\checkmarki\checkmarkr\checkmarkc\checkmark$\checkmarke\checkmark$\checkmark^\checkmark\\checkmarkc\checkmarki\checkmarkr\checkmarkc\checkmark$\checkmarkn\checkmark$\checkmark^\checkmark\\checkmarkc\checkmarki\checkmarkr\checkmarkc\checkmark$\checkmarkt\checkmark$\checkmark^\checkmark\\checkmarkc\checkmarki\checkmarkr\checkmarkc\checkmark$\checkmarke\checkmark$\checkmark^\checkmark\\checkmarkc\checkmarki\checkmarkr\checkmarkc\checkmark$\checkmarkr\checkmark$\checkmark^\checkmark\\checkmarkc\checkmarki\checkmarkr\checkmarkc\checkmark$\checkmarki\checkmark$\checkmark^\checkmark\\checkmarkc\checkmarki\checkmarkr\checkmarkc\checkmark$\checkmarkn\checkmark$\checkmark^\checkmark\\checkmarkc\checkmarki\checkmarkr\checkmarkc\checkmark$\checkmarkg\checkmark$\checkmark^\checkmark\\checkmarkc\checkmarki\checkmarkr\checkmarkc\checkmark$\checkmark
\checkmark$\checkmark^\checkmark\\checkmarkc\checkmarki\checkmarkr\checkmarkc\checkmark$\checkmark \checkmark$\checkmark^\checkmark\\checkmarkc\checkmarki\checkmarkr\checkmarkc\checkmark$\checkmark \checkmark$\checkmark^\checkmark\\checkmarkc\checkmarki\checkmarkr\checkmarkc\checkmark$\checkmark \checkmark$\checkmark^\checkmark\\checkmarkc\checkmarki\checkmarkr\checkmarkc\checkmark$\checkmark \checkmark$\checkmark^\checkmark\\checkmarkc\checkmarki\checkmarkr\checkmarkc\checkmark$\checkmark\\checkmark$\checkmark^\checkmark\\checkmarkc\checkmarki\checkmarkr\checkmarkc\checkmark$\checkmarkb\checkmark$\checkmark^\checkmark\\checkmarkc\checkmarki\checkmarkr\checkmarkc\checkmark$\checkmarke\checkmark$\checkmark^\checkmark\\checkmarkc\checkmarki\checkmarkr\checkmarkc\checkmark$\checkmarkg\checkmark$\checkmark^\checkmark\\checkmarkc\checkmarki\checkmarkr\checkmarkc\checkmark$\checkmarki\checkmark$\checkmark^\checkmark\\checkmarkc\checkmarki\checkmarkr\checkmarkc\checkmark$\checkmarkn\checkmark$\checkmark^\checkmark\\checkmarkc\checkmarki\checkmarkr\checkmarkc\checkmark$\checkmark{\checkmark$\checkmark^\checkmark\\checkmarkc\checkmarki\checkmarkr\checkmarkc\checkmark$\checkmarkt\checkmark$\checkmark^\checkmark\\checkmarkc\checkmarki\checkmarkr\checkmarkc\checkmark$\checkmarki\checkmark$\checkmark^\checkmark\\checkmarkc\checkmarki\checkmarkr\checkmarkc\checkmark$\checkmarkk\checkmark$\checkmark^\checkmark\\checkmarkc\checkmarki\checkmarkr\checkmarkc\checkmark$\checkmarkz\checkmark$\checkmark^\checkmark\\checkmarkc\checkmarki\checkmarkr\checkmarkc\checkmark$\checkmarkp\checkmark$\checkmark^\checkmark\\checkmarkc\checkmarki\checkmarkr\checkmarkc\checkmark$\checkmarki\checkmark$\checkmark^\checkmark\\checkmarkc\checkmarki\checkmarkr\checkmarkc\checkmark$\checkmarkc\checkmark$\checkmark^\checkmark\\checkmarkc\checkmarki\checkmarkr\checkmarkc\checkmark$\checkmarkt\checkmark$\checkmark^\checkmark\\checkmarkc\checkmarki\checkmarkr\checkmarkc\checkmark$\checkmarku\checkmark$\checkmark^\checkmark\\checkmarkc\checkmarki\checkmarkr\checkmarkc\checkmark$\checkmarkr\checkmark$\checkmark^\checkmark\\checkmarkc\checkmarki\checkmarkr\checkmarkc\checkmark$\checkmarke\checkmark$\checkmark^\checkmark\\checkmarkc\checkmarki\checkmarkr\checkmarkc\checkmark$\checkmark}\checkmark$\checkmark^\checkmark\\checkmarkc\checkmarki\checkmarkr\checkmarkc\checkmark$\checkmark[\checkmark$\checkmark^\checkmark\\checkmarkc\checkmarki\checkmarkr\checkmarkc\checkmark$\checkmark
\checkmark$\checkmark^\checkmark\\checkmarkc\checkmarki\checkmarkr\checkmarkc\checkmark$\checkmark \checkmark$\checkmark^\checkmark\\checkmarkc\checkmarki\checkmarkr\checkmarkc\checkmark$\checkmark \checkmark$\checkmark^\checkmark\\checkmarkc\checkmarki\checkmarkr\checkmarkc\checkmark$\checkmark \checkmark$\checkmark^\checkmark\\checkmarkc\checkmarki\checkmarkr\checkmarkc\checkmark$\checkmark \checkmark$\checkmark^\checkmark\\checkmarkc\checkmarki\checkmarkr\checkmarkc\checkmark$\checkmark \checkmark$\checkmark^\checkmark\\checkmarkc\checkmarki\checkmarkr\checkmarkc\checkmark$\checkmark \checkmark$\checkmark^\checkmark\\checkmarkc\checkmarki\checkmarkr\checkmarkc\checkmark$\checkmark \checkmark$\checkmark^\checkmark\\checkmarkc\checkmarki\checkmarkr\checkmarkc\checkmark$\checkmark \checkmark$\checkmark^\checkmark\\checkmarkc\checkmarki\checkmarkr\checkmarkc\checkmark$\checkmarkf\checkmark$\checkmark^\checkmark\\checkmarkc\checkmarki\checkmarkr\checkmarkc\checkmark$\checkmarko\checkmark$\checkmark^\checkmark\\checkmarkc\checkmarki\checkmarkr\checkmarkc\checkmark$\checkmarkr\checkmark$\checkmark^\checkmark\\checkmarkc\checkmarki\checkmarkr\checkmarkc\checkmark$\checkmarkc\checkmark$\checkmark^\checkmark\\checkmarkc\checkmarki\checkmarkr\checkmarkc\checkmark$\checkmarke\checkmark$\checkmark^\checkmark\\checkmarkc\checkmarki\checkmarkr\checkmarkc\checkmark$\checkmark/\checkmark$\checkmark^\checkmark\\checkmarkc\checkmarki\checkmarkr\checkmarkc\checkmark$\checkmark.\checkmark$\checkmark^\checkmark\\checkmarkc\checkmarki\checkmarkr\checkmarkc\checkmark$\checkmarks\checkmark$\checkmark^\checkmark\\checkmarkc\checkmarki\checkmarkr\checkmarkc\checkmark$\checkmarkt\checkmark$\checkmark^\checkmark\\checkmarkc\checkmarki\checkmarkr\checkmarkc\checkmark$\checkmarky\checkmark$\checkmark^\checkmark\\checkmarkc\checkmarki\checkmarkr\checkmarkc\checkmark$\checkmarkl\checkmark$\checkmark^\checkmark\\checkmarkc\checkmarki\checkmarkr\checkmarkc\checkmark$\checkmarke\checkmark$\checkmark^\checkmark\\checkmarkc\checkmarki\checkmarkr\checkmarkc\checkmark$\checkmark=\checkmark$\checkmark^\checkmark\\checkmarkc\checkmarki\checkmarkr\checkmarkc\checkmark$\checkmark{\checkmark$\checkmark^\checkmark\\checkmarkc\checkmarki\checkmarkr\checkmarkc\checkmark$\checkmark-\checkmark$\checkmark^\checkmark\\checkmarkc\checkmarki\checkmarkr\checkmarkc\checkmark$\checkmark{\checkmark$\checkmark^\checkmark\\checkmarkc\checkmarki\checkmarkr\checkmarkc\checkmark$\checkmarkS\checkmark$\checkmark^\checkmark\\checkmarkc\checkmarki\checkmarkr\checkmarkc\checkmark$\checkmarkt\checkmark$\checkmark^\checkmark\\checkmarkc\checkmarki\checkmarkr\checkmarkc\checkmark$\checkmarke\checkmark$\checkmark^\checkmark\\checkmarkc\checkmarki\checkmarkr\checkmarkc\checkmark$\checkmarka\checkmark$\checkmark^\checkmark\\checkmarkc\checkmarki\checkmarkr\checkmarkc\checkmark$\checkmarkl\checkmark$\checkmark^\checkmark\\checkmarkc\checkmarki\checkmarkr\checkmarkc\checkmark$\checkmarkt\checkmark$\checkmark^\checkmark\\checkmarkc\checkmarki\checkmarkr\checkmarkc\checkmark$\checkmarkh\checkmark$\checkmark^\checkmark\\checkmarkc\checkmarki\checkmarkr\checkmarkc\checkmark$\checkmark[\checkmark$\checkmark^\checkmark\\checkmarkc\checkmarki\checkmarkr\checkmarkc\checkmark$\checkmarks\checkmark$\checkmark^\checkmark\\checkmarkc\checkmarki\checkmarkr\checkmarkc\checkmark$\checkmarkc\checkmark$\checkmark^\checkmark\\checkmarkc\checkmarki\checkmarkr\checkmarkc\checkmark$\checkmarka\checkmark$\checkmark^\checkmark\\checkmarkc\checkmarki\checkmarkr\checkmarkc\checkmark$\checkmarkl\checkmark$\checkmark^\checkmark\\checkmarkc\checkmarki\checkmarkr\checkmarkc\checkmark$\checkmarke\checkmark$\checkmark^\checkmark\\checkmarkc\checkmarki\checkmarkr\checkmarkc\checkmark$\checkmark=\checkmark$\checkmark^\checkmark\\checkmarkc\checkmarki\checkmarkr\checkmarkc\checkmark$\checkmark1\checkmark$\checkmark^\checkmark\\checkmarkc\checkmarki\checkmarkr\checkmarkc\checkmark$\checkmark.\checkmark$\checkmark^\checkmark\\checkmarkc\checkmarki\checkmarkr\checkmarkc\checkmark$\checkmark2\checkmark$\checkmark^\checkmark\\checkmarkc\checkmarki\checkmarkr\checkmarkc\checkmark$\checkmark]\checkmark$\checkmark^\checkmark\\checkmarkc\checkmarki\checkmarkr\checkmarkc\checkmark$\checkmark}\checkmark$\checkmark^\checkmark\\checkmarkc\checkmarki\checkmarkr\checkmarkc\checkmark$\checkmark,\checkmark$\checkmark^\checkmark\\checkmarkc\checkmarki\checkmarkr\checkmarkc\checkmark$\checkmark \checkmark$\checkmark^\checkmark\\checkmarkc\checkmarki\checkmarkr\checkmarkc\checkmark$\checkmarkt\checkmark$\checkmark^\checkmark\\checkmarkc\checkmarki\checkmarkr\checkmarkc\checkmark$\checkmarkh\checkmark$\checkmark^\checkmark\\checkmarkc\checkmarki\checkmarkr\checkmarkc\checkmark$\checkmarki\checkmark$\checkmark^\checkmark\\checkmarkc\checkmarki\checkmarkr\checkmarkc\checkmark$\checkmarkc\checkmark$\checkmark^\checkmark\\checkmarkc\checkmarki\checkmarkr\checkmarkc\checkmark$\checkmarkk\checkmark$\checkmark^\checkmark\\checkmarkc\checkmarki\checkmarkr\checkmarkc\checkmark$\checkmark,\checkmark$\checkmark^\checkmark\\checkmarkc\checkmarki\checkmarkr\checkmarkc\checkmark$\checkmark \checkmark$\checkmark^\checkmark\\checkmarkc\checkmarki\checkmarkr\checkmarkc\checkmark$\checkmarkr\checkmark$\checkmark^\checkmark\\checkmarkc\checkmarki\checkmarkr\checkmarkc\checkmark$\checkmarke\checkmark$\checkmark^\checkmark\\checkmarkc\checkmarki\checkmarkr\checkmarkc\checkmark$\checkmarkd\checkmark$\checkmark^\checkmark\\checkmarkc\checkmarki\checkmarkr\checkmarkc\checkmark$\checkmark}\checkmark$\checkmark^\checkmark\\checkmarkc\checkmarki\checkmarkr\checkmarkc\checkmark$\checkmark,\checkmark$\checkmark^\checkmark\\checkmarkc\checkmarki\checkmarkr\checkmarkc\checkmark$\checkmark
\checkmark$\checkmark^\checkmark\\checkmarkc\checkmarki\checkmarkr\checkmarkc\checkmark$\checkmark \checkmark$\checkmark^\checkmark\\checkmarkc\checkmarki\checkmarkr\checkmarkc\checkmark$\checkmark \checkmark$\checkmark^\checkmark\\checkmarkc\checkmarki\checkmarkr\checkmarkc\checkmark$\checkmark \checkmark$\checkmark^\checkmark\\checkmarkc\checkmarki\checkmarkr\checkmarkc\checkmark$\checkmark \checkmark$\checkmark^\checkmark\\checkmarkc\checkmarki\checkmarkr\checkmarkc\checkmark$\checkmark \checkmark$\checkmark^\checkmark\\checkmarkc\checkmarki\checkmarkr\checkmarkc\checkmark$\checkmark \checkmark$\checkmark^\checkmark\\checkmarkc\checkmarki\checkmarkr\checkmarkc\checkmark$\checkmark \checkmark$\checkmark^\checkmark\\checkmarkc\checkmarki\checkmarkr\checkmarkc\checkmark$\checkmark \checkmark$\checkmark^\checkmark\\checkmarkc\checkmarki\checkmarkr\checkmarkc\checkmark$\checkmarka\checkmark$\checkmark^\checkmark\\checkmarkc\checkmarki\checkmarkr\checkmarkc\checkmark$\checkmarkr\checkmark$\checkmark^\checkmark\\checkmarkc\checkmarki\checkmarkr\checkmarkc\checkmark$\checkmarkm\checkmark$\checkmark^\checkmark\\checkmarkc\checkmarki\checkmarkr\checkmarkc\checkmark$\checkmark/\checkmark$\checkmark^\checkmark\\checkmarkc\checkmarki\checkmarkr\checkmarkc\checkmark$\checkmark.\checkmark$\checkmark^\checkmark\\checkmarkc\checkmarki\checkmarkr\checkmarkc\checkmark$\checkmarks\checkmark$\checkmark^\checkmark\\checkmarkc\checkmarki\checkmarkr\checkmarkc\checkmark$\checkmarkt\checkmark$\checkmark^\checkmark\\checkmarkc\checkmarki\checkmarkr\checkmarkc\checkmark$\checkmarky\checkmark$\checkmark^\checkmark\\checkmarkc\checkmarki\checkmarkr\checkmarkc\checkmark$\checkmarkl\checkmark$\checkmark^\checkmark\\checkmarkc\checkmarki\checkmarkr\checkmarkc\checkmark$\checkmarke\checkmark$\checkmark^\checkmark\\checkmarkc\checkmarki\checkmarkr\checkmarkc\checkmark$\checkmark=\checkmark$\checkmark^\checkmark\\checkmarkc\checkmarki\checkmarkr\checkmarkc\checkmark$\checkmark{\checkmark$\checkmark^\checkmark\\checkmarkc\checkmarki\checkmarkr\checkmarkc\checkmark$\checkmarkt\checkmark$\checkmark^\checkmark\\checkmarkc\checkmarki\checkmarkr\checkmarkc\checkmark$\checkmarkh\checkmark$\checkmark^\checkmark\\checkmarkc\checkmarki\checkmarkr\checkmarkc\checkmark$\checkmarki\checkmark$\checkmark^\checkmark\\checkmarkc\checkmarki\checkmarkr\checkmarkc\checkmark$\checkmarkc\checkmark$\checkmark^\checkmark\\checkmarkc\checkmarki\checkmarkr\checkmarkc\checkmark$\checkmarkk\checkmark$\checkmark^\checkmark\\checkmarkc\checkmarki\checkmarkr\checkmarkc\checkmark$\checkmark,\checkmark$\checkmark^\checkmark\\checkmarkc\checkmarki\checkmarkr\checkmarkc\checkmark$\checkmark \checkmark$\checkmark^\checkmark\\checkmarkc\checkmarki\checkmarkr\checkmarkc\checkmark$\checkmarkf\checkmark$\checkmark^\checkmark\\checkmarkc\checkmarki\checkmarkr\checkmarkc\checkmark$\checkmarki\checkmark$\checkmark^\checkmark\\checkmarkc\checkmarki\checkmarkr\checkmarkc\checkmark$\checkmarkl\checkmark$\checkmark^\checkmark\\checkmarkc\checkmarki\checkmarkr\checkmarkc\checkmark$\checkmarkl\checkmark$\checkmark^\checkmark\\checkmarkc\checkmarki\checkmarkr\checkmarkc\checkmark$\checkmark=\checkmark$\checkmark^\checkmark\\checkmarkc\checkmarki\checkmarkr\checkmarkc\checkmark$\checkmarkb\checkmark$\checkmark^\checkmark\\checkmarkc\checkmarki\checkmarkr\checkmarkc\checkmark$\checkmarkl\checkmark$\checkmark^\checkmark\\checkmarkc\checkmarki\checkmarkr\checkmarkc\checkmark$\checkmarku\checkmark$\checkmark^\checkmark\\checkmarkc\checkmarki\checkmarkr\checkmarkc\checkmark$\checkmarke\checkmark$\checkmark^\checkmark\\checkmarkc\checkmarki\checkmarkr\checkmarkc\checkmark$\checkmark!\checkmark$\checkmark^\checkmark\\checkmarkc\checkmarki\checkmarkr\checkmarkc\checkmark$\checkmark8\checkmark$\checkmark^\checkmark\\checkmarkc\checkmarki\checkmarkr\checkmarkc\checkmark$\checkmark,\checkmark$\checkmark^\checkmark\\checkmarkc\checkmarki\checkmarkr\checkmarkc\checkmark$\checkmark \checkmark$\checkmark^\checkmark\\checkmarkc\checkmarki\checkmarkr\checkmarkc\checkmark$\checkmarkd\checkmark$\checkmark^\checkmark\\checkmarkc\checkmarki\checkmarkr\checkmarkc\checkmark$\checkmarkr\checkmark$\checkmark^\checkmark\\checkmarkc\checkmarki\checkmarkr\checkmarkc\checkmark$\checkmarka\checkmark$\checkmark^\checkmark\\checkmarkc\checkmarki\checkmarkr\checkmarkc\checkmark$\checkmarkw\checkmark$\checkmark^\checkmark\\checkmarkc\checkmarki\checkmarkr\checkmarkc\checkmark$\checkmark=\checkmark$\checkmark^\checkmark\\checkmarkc\checkmarki\checkmarkr\checkmarkc\checkmark$\checkmarkb\checkmark$\checkmark^\checkmark\\checkmarkc\checkmarki\checkmarkr\checkmarkc\checkmark$\checkmarkl\checkmark$\checkmark^\checkmark\\checkmarkc\checkmarki\checkmarkr\checkmarkc\checkmark$\checkmarka\checkmark$\checkmark^\checkmark\\checkmarkc\checkmarki\checkmarkr\checkmarkc\checkmark$\checkmarkc\checkmark$\checkmark^\checkmark\\checkmarkc\checkmarki\checkmarkr\checkmarkc\checkmark$\checkmarkk\checkmark$\checkmark^\checkmark\\checkmarkc\checkmarki\checkmarkr\checkmarkc\checkmark$\checkmark}\checkmark$\checkmark^\checkmark\\checkmarkc\checkmarki\checkmarkr\checkmarkc\checkmark$\checkmark
\checkmark$\checkmark^\checkmark\\checkmarkc\checkmarki\checkmarkr\checkmarkc\checkmark$\checkmark \checkmark$\checkmark^\checkmark\\checkmarkc\checkmarki\checkmarkr\checkmarkc\checkmark$\checkmark \checkmark$\checkmark^\checkmark\\checkmarkc\checkmarki\checkmarkr\checkmarkc\checkmark$\checkmark \checkmark$\checkmark^\checkmark\\checkmarkc\checkmarki\checkmarkr\checkmarkc\checkmark$\checkmark \checkmark$\checkmark^\checkmark\\checkmarkc\checkmarki\checkmarkr\checkmarkc\checkmark$\checkmark]\checkmark$\checkmark^\checkmark\\checkmarkc\checkmarki\checkmarkr\checkmarkc\checkmark$\checkmark
\checkmark$\checkmark^\checkmark\\checkmarkc\checkmarki\checkmarkr\checkmarkc\checkmark$\checkmark \checkmark$\checkmark^\checkmark\\checkmarkc\checkmarki\checkmarkr\checkmarkc\checkmark$\checkmark \checkmark$\checkmark^\checkmark\\checkmarkc\checkmarki\checkmarkr\checkmarkc\checkmark$\checkmark \checkmark$\checkmark^\checkmark\\checkmarkc\checkmarki\checkmarkr\checkmarkc\checkmark$\checkmark \checkmark$\checkmark^\checkmark\\checkmarkc\checkmarki\checkmarkr\checkmarkc\checkmark$\checkmark \checkmark$\checkmark^\checkmark\\checkmarkc\checkmarki\checkmarkr\checkmarkc\checkmark$\checkmark \checkmark$\checkmark^\checkmark\\checkmarkc\checkmarki\checkmarkr\checkmarkc\checkmark$\checkmark \checkmark$\checkmark^\checkmark\\checkmarkc\checkmarki\checkmarkr\checkmarkc\checkmark$\checkmark \checkmark$\checkmark^\checkmark\\checkmarkc\checkmarki\checkmarkr\checkmarkc\checkmark$\checkmark\\checkmark$\checkmark^\checkmark\\checkmarkc\checkmarki\checkmarkr\checkmarkc\checkmark$\checkmarkd\checkmark$\checkmark^\checkmark\\checkmarkc\checkmarki\checkmarkr\checkmarkc\checkmark$\checkmarkr\checkmark$\checkmark^\checkmark\\checkmarkc\checkmarki\checkmarkr\checkmarkc\checkmark$\checkmarka\checkmark$\checkmark^\checkmark\\checkmarkc\checkmarki\checkmarkr\checkmarkc\checkmark$\checkmarkw\checkmark$\checkmark^\checkmark\\checkmarkc\checkmarki\checkmarkr\checkmarkc\checkmark$\checkmark[\checkmark$\checkmark^\checkmark\\checkmarkc\checkmarki\checkmarkr\checkmarkc\checkmark$\checkmarka\checkmark$\checkmark^\checkmark\\checkmarkc\checkmarki\checkmarkr\checkmarkc\checkmark$\checkmarkr\checkmark$\checkmark^\checkmark\\checkmarkc\checkmarki\checkmarkr\checkmarkc\checkmark$\checkmarkm\checkmark$\checkmark^\checkmark\\checkmarkc\checkmarki\checkmarkr\checkmarkc\checkmark$\checkmark]\checkmark$\checkmark^\checkmark\\checkmarkc\checkmarki\checkmarkr\checkmarkc\checkmark$\checkmark \checkmark$\checkmark^\checkmark\\checkmarkc\checkmarki\checkmarkr\checkmarkc\checkmark$\checkmark(\checkmark$\checkmark^\checkmark\\checkmarkc\checkmarki\checkmarkr\checkmarkc\checkmark$\checkmark-\checkmark$\checkmark^\checkmark\\checkmarkc\checkmarki\checkmarkr\checkmarkc\checkmark$\checkmark3\checkmark$\checkmark^\checkmark\\checkmarkc\checkmarki\checkmarkr\checkmarkc\checkmark$\checkmark,\checkmark$\checkmark^\checkmark\\checkmarkc\checkmarki\checkmarkr\checkmarkc\checkmark$\checkmark-\checkmark$\checkmark^\checkmark\\checkmarkc\checkmarki\checkmarkr\checkmarkc\checkmark$\checkmark0\checkmark$\checkmark^\checkmark\\checkmarkc\checkmarki\checkmarkr\checkmarkc\checkmark$\checkmark.\checkmark$\checkmark^\checkmark\\checkmarkc\checkmarki\checkmarkr\checkmarkc\checkmark$\checkmark4\checkmark$\checkmark^\checkmark\\checkmarkc\checkmarki\checkmarkr\checkmarkc\checkmark$\checkmark)\checkmark$\checkmark^\checkmark\\checkmarkc\checkmarki\checkmarkr\checkmarkc\checkmark$\checkmark \checkmark$\checkmark^\checkmark\\checkmarkc\checkmarki\checkmarkr\checkmarkc\checkmark$\checkmarkr\checkmark$\checkmark^\checkmark\\checkmarkc\checkmarki\checkmarkr\checkmarkc\checkmark$\checkmarke\checkmark$\checkmark^\checkmark\\checkmarkc\checkmarki\checkmarkr\checkmarkc\checkmark$\checkmarkc\checkmark$\checkmark^\checkmark\\checkmarkc\checkmarki\checkmarkr\checkmarkc\checkmark$\checkmarkt\checkmark$\checkmark^\checkmark\\checkmarkc\checkmarki\checkmarkr\checkmarkc\checkmark$\checkmarka\checkmark$\checkmark^\checkmark\\checkmarkc\checkmarki\checkmarkr\checkmarkc\checkmark$\checkmarkn\checkmark$\checkmark^\checkmark\\checkmarkc\checkmarki\checkmarkr\checkmarkc\checkmark$\checkmarkg\checkmark$\checkmark^\checkmark\\checkmarkc\checkmarki\checkmarkr\checkmarkc\checkmark$\checkmarkl\checkmark$\checkmark^\checkmark\\checkmarkc\checkmarki\checkmarkr\checkmarkc\checkmark$\checkmarke\checkmark$\checkmark^\checkmark\\checkmarkc\checkmarki\checkmarkr\checkmarkc\checkmark$\checkmark \checkmark$\checkmark^\checkmark\\checkmarkc\checkmarki\checkmarkr\checkmarkc\checkmark$\checkmark(\checkmark$\checkmark^\checkmark\\checkmarkc\checkmarki\checkmarkr\checkmarkc\checkmark$\checkmark3\checkmark$\checkmark^\checkmark\\checkmarkc\checkmarki\checkmarkr\checkmarkc\checkmark$\checkmark,\checkmark$\checkmark^\checkmark\\checkmarkc\checkmarki\checkmarkr\checkmarkc\checkmark$\checkmark0\checkmark$\checkmark^\checkmark\\checkmarkc\checkmarki\checkmarkr\checkmarkc\checkmark$\checkmark)\checkmark$\checkmark^\checkmark\\checkmarkc\checkmarki\checkmarkr\checkmarkc\checkmark$\checkmark;\checkmark$\checkmark^\checkmark\\checkmarkc\checkmarki\checkmarkr\checkmarkc\checkmark$\checkmark
\checkmark$\checkmark^\checkmark\\checkmarkc\checkmarki\checkmarkr\checkmarkc\checkmark$\checkmark \checkmark$\checkmark^\checkmark\\checkmarkc\checkmarki\checkmarkr\checkmarkc\checkmark$\checkmark \checkmark$\checkmark^\checkmark\\checkmarkc\checkmarki\checkmarkr\checkmarkc\checkmark$\checkmark \checkmark$\checkmark^\checkmark\\checkmarkc\checkmarki\checkmarkr\checkmarkc\checkmark$\checkmark \checkmark$\checkmark^\checkmark\\checkmarkc\checkmarki\checkmarkr\checkmarkc\checkmark$\checkmark \checkmark$\checkmark^\checkmark\\checkmarkc\checkmarki\checkmarkr\checkmarkc\checkmark$\checkmark \checkmark$\checkmark^\checkmark\\checkmarkc\checkmarki\checkmarkr\checkmarkc\checkmark$\checkmark \checkmark$\checkmark^\checkmark\\checkmarkc\checkmarki\checkmarkr\checkmarkc\checkmark$\checkmark \checkmark$\checkmark^\checkmark\\checkmarkc\checkmarki\checkmarkr\checkmarkc\checkmark$\checkmark\\checkmark$\checkmark^\checkmark\\checkmarkc\checkmarki\checkmarkr\checkmarkc\checkmark$\checkmarkn\checkmark$\checkmark^\checkmark\\checkmarkc\checkmarki\checkmarkr\checkmarkc\checkmark$\checkmarko\checkmark$\checkmark^\checkmark\\checkmarkc\checkmarki\checkmarkr\checkmarkc\checkmark$\checkmarkd\checkmark$\checkmark^\checkmark\\checkmarkc\checkmarki\checkmarkr\checkmarkc\checkmark$\checkmarke\checkmark$\checkmark^\checkmark\\checkmarkc\checkmarki\checkmarkr\checkmarkc\checkmark$\checkmark[\checkmark$\checkmark^\checkmark\\checkmarkc\checkmarki\checkmarkr\checkmarkc\checkmark$\checkmarkb\checkmark$\checkmark^\checkmark\\checkmarkc\checkmarki\checkmarkr\checkmarkc\checkmark$\checkmarke\checkmark$\checkmark^\checkmark\\checkmarkc\checkmarki\checkmarkr\checkmarkc\checkmark$\checkmarkl\checkmark$\checkmark^\checkmark\\checkmarkc\checkmarki\checkmarkr\checkmarkc\checkmark$\checkmarko\checkmark$\checkmark^\checkmark\\checkmarkc\checkmarki\checkmarkr\checkmarkc\checkmark$\checkmarkw\checkmark$\checkmark^\checkmark\\checkmarkc\checkmarki\checkmarkr\checkmarkc\checkmark$\checkmark]\checkmark$\checkmark^\checkmark\\checkmarkc\checkmarki\checkmarkr\checkmarkc\checkmark$\checkmark \checkmark$\checkmark^\checkmark\\checkmarkc\checkmarki\checkmarkr\checkmarkc\checkmark$\checkmarka\checkmark$\checkmark^\checkmark\\checkmarkc\checkmarki\checkmarkr\checkmarkc\checkmark$\checkmarkt\checkmark$\checkmark^\checkmark\\checkmarkc\checkmarki\checkmarkr\checkmarkc\checkmark$\checkmark \checkmark$\checkmark^\checkmark\\checkmarkc\checkmarki\checkmarkr\checkmarkc\checkmark$\checkmark(\checkmark$\checkmark^\checkmark\\checkmarkc\checkmarki\checkmarkr\checkmarkc\checkmark$\checkmark0\checkmark$\checkmark^\checkmark\\checkmarkc\checkmarki\checkmarkr\checkmarkc\checkmark$\checkmark,\checkmark$\checkmark^\checkmark\\checkmarkc\checkmarki\checkmarkr\checkmarkc\checkmark$\checkmark-\checkmark$\checkmark^\checkmark\\checkmarkc\checkmarki\checkmarkr\checkmarkc\checkmark$\checkmark0\checkmark$\checkmark^\checkmark\\checkmarkc\checkmarki\checkmarkr\checkmarkc\checkmark$\checkmark.\checkmark$\checkmark^\checkmark\\checkmarkc\checkmarki\checkmarkr\checkmarkc\checkmark$\checkmark4\checkmark$\checkmark^\checkmark\\checkmarkc\checkmarki\checkmarkr\checkmarkc\checkmark$\checkmark)\checkmark$\checkmark^\checkmark\\checkmarkc\checkmarki\checkmarkr\checkmarkc\checkmark$\checkmark \checkmark$\checkmark^\checkmark\\checkmarkc\checkmarki\checkmarkr\checkmarkc\checkmark$\checkmark{\checkmark$\checkmark^\checkmark\\checkmarkc\checkmarki\checkmarkr\checkmarkc\checkmark$\checkmark\\checkmark$\checkmark^\checkmark\\checkmarkc\checkmarki\checkmarkr\checkmarkc\checkmark$\checkmarks\checkmark$\checkmark^\checkmark\\checkmarkc\checkmarki\checkmarkr\checkmarkc\checkmark$\checkmarkm\checkmark$\checkmark^\checkmark\\checkmarkc\checkmarki\checkmarkr\checkmarkc\checkmark$\checkmarka\checkmark$\checkmark^\checkmark\\checkmarkc\checkmarki\checkmarkr\checkmarkc\checkmark$\checkmarkl\checkmark$\checkmark^\checkmark\\checkmarkc\checkmarki\checkmarkr\checkmarkc\checkmark$\checkmarkl\checkmark$\checkmark^\checkmark\\checkmarkc\checkmarki\checkmarkr\checkmarkc\checkmark$\checkmark \checkmark$\checkmark^\checkmark\\checkmarkc\checkmarki\checkmarkr\checkmarkc\checkmark$\checkmarkR\checkmark$\checkmark^\checkmark\\checkmarkc\checkmarki\checkmarkr\checkmarkc\checkmark$\checkmarko\checkmark$\checkmark^\checkmark\\checkmarkc\checkmarki\checkmarkr\checkmarkc\checkmark$\checkmarku\checkmark$\checkmark^\checkmark\\checkmarkc\checkmarki\checkmarkr\checkmarkc\checkmark$\checkmarkl\checkmark$\checkmark^\checkmark\\checkmarkc\checkmarki\checkmarkr\checkmarkc\checkmark$\checkmarke\checkmark$\checkmark^\checkmark\\checkmarkc\checkmarki\checkmarkr\checkmarkc\checkmark$\checkmarkm\checkmark$\checkmark^\checkmark\\checkmarkc\checkmarki\checkmarkr\checkmarkc\checkmark$\checkmarke\checkmark$\checkmark^\checkmark\\checkmarkc\checkmarki\checkmarkr\checkmarkc\checkmark$\checkmarkn\checkmark$\checkmark^\checkmark\\checkmarkc\checkmarki\checkmarkr\checkmarkc\checkmark$\checkmarkt\checkmark$\checkmark^\checkmark\\checkmarkc\checkmarki\checkmarkr\checkmarkc\checkmark$\checkmarks\checkmark$\checkmark^\checkmark\\checkmarkc\checkmarki\checkmarkr\checkmarkc\checkmark$\checkmark \checkmark$\checkmark^\checkmark\\checkmarkc\checkmarki\checkmarkr\checkmarkc\checkmark$\checkmark(\checkmark$\checkmark^\checkmark\\checkmarkc\checkmarki\checkmarkr\checkmarkc\checkmark$\checkmarkA\checkmark$\checkmark^\checkmark\\checkmarkc\checkmarki\checkmarkr\checkmarkc\checkmark$\checkmarkx\checkmark$\checkmark^\checkmark\\checkmarkc\checkmarki\checkmarkr\checkmarkc\checkmark$\checkmarke\checkmark$\checkmark^\checkmark\\checkmarkc\checkmarki\checkmarkr\checkmarkc\checkmark$\checkmark \checkmark$\checkmark^\checkmark\\checkmarkc\checkmarki\checkmarkr\checkmarkc\checkmark$\checkmarkA\checkmark$\checkmark^\checkmark\\checkmarkc\checkmarki\checkmarkr\checkmarkc\checkmark$\checkmark)\checkmark$\checkmark^\checkmark\\checkmarkc\checkmarki\checkmarkr\checkmarkc\checkmark$\checkmark}\checkmark$\checkmark^\checkmark\\checkmarkc\checkmarki\checkmarkr\checkmarkc\checkmark$\checkmark;\checkmark$\checkmark^\checkmark\\checkmarkc\checkmarki\checkmarkr\checkmarkc\checkmark$\checkmark
\checkmark$\checkmark^\checkmark\\checkmarkc\checkmarki\checkmarkr\checkmarkc\checkmark$\checkmark \checkmark$\checkmark^\checkmark\\checkmarkc\checkmarki\checkmarkr\checkmarkc\checkmark$\checkmark \checkmark$\checkmark^\checkmark\\checkmarkc\checkmarki\checkmarkr\checkmarkc\checkmark$\checkmark \checkmark$\checkmark^\checkmark\\checkmarkc\checkmarki\checkmarkr\checkmarkc\checkmark$\checkmark \checkmark$\checkmark^\checkmark\\checkmarkc\checkmarki\checkmarkr\checkmarkc\checkmark$\checkmark \checkmark$\checkmark^\checkmark\\checkmarkc\checkmarki\checkmarkr\checkmarkc\checkmark$\checkmark \checkmark$\checkmark^\checkmark\\checkmarkc\checkmarki\checkmarkr\checkmarkc\checkmark$\checkmark \checkmark$\checkmark^\checkmark\\checkmarkc\checkmarki\checkmarkr\checkmarkc\checkmark$\checkmark \checkmark$\checkmark^\checkmark\\checkmarkc\checkmarki\checkmarkr\checkmarkc\checkmark$\checkmark\\checkmark$\checkmark^\checkmark\\checkmarkc\checkmarki\checkmarkr\checkmarkc\checkmark$\checkmarkd\checkmark$\checkmark^\checkmark\\checkmarkc\checkmarki\checkmarkr\checkmarkc\checkmark$\checkmarkr\checkmark$\checkmark^\checkmark\\checkmarkc\checkmarki\checkmarkr\checkmarkc\checkmark$\checkmarka\checkmark$\checkmark^\checkmark\\checkmarkc\checkmarki\checkmarkr\checkmarkc\checkmark$\checkmarkw\checkmark$\checkmark^\checkmark\\checkmarkc\checkmarki\checkmarkr\checkmarkc\checkmark$\checkmark[\checkmark$\checkmark^\checkmark\\checkmarkc\checkmarki\checkmarkr\checkmarkc\checkmark$\checkmarka\checkmark$\checkmark^\checkmark\\checkmarkc\checkmarki\checkmarkr\checkmarkc\checkmark$\checkmarkr\checkmark$\checkmark^\checkmark\\checkmarkc\checkmarki\checkmarkr\checkmarkc\checkmark$\checkmarkm\checkmark$\checkmark^\checkmark\\checkmarkc\checkmarki\checkmarkr\checkmarkc\checkmark$\checkmark]\checkmark$\checkmark^\checkmark\\checkmarkc\checkmarki\checkmarkr\checkmarkc\checkmark$\checkmark \checkmark$\checkmark^\checkmark\\checkmarkc\checkmarki\checkmarkr\checkmarkc\checkmark$\checkmark(\checkmark$\checkmark^\checkmark\\checkmarkc\checkmarki\checkmarkr\checkmarkc\checkmark$\checkmark-\checkmark$\checkmark^\checkmark\\checkmarkc\checkmarki\checkmarkr\checkmarkc\checkmark$\checkmark0\checkmark$\checkmark^\checkmark\\checkmarkc\checkmarki\checkmarkr\checkmarkc\checkmark$\checkmark.\checkmark$\checkmark^\checkmark\\checkmarkc\checkmarki\checkmarkr\checkmarkc\checkmark$\checkmark4\checkmark$\checkmark^\checkmark\\checkmarkc\checkmarki\checkmarkr\checkmarkc\checkmark$\checkmark,\checkmark$\checkmark^\checkmark\\checkmarkc\checkmarki\checkmarkr\checkmarkc\checkmark$\checkmark0\checkmark$\checkmark^\checkmark\\checkmarkc\checkmarki\checkmarkr\checkmarkc\checkmark$\checkmark)\checkmark$\checkmark^\checkmark\\checkmarkc\checkmarki\checkmarkr\checkmarkc\checkmark$\checkmark \checkmark$\checkmark^\checkmark\\checkmarkc\checkmarki\checkmarkr\checkmarkc\checkmark$\checkmarkr\checkmark$\checkmark^\checkmark\\checkmarkc\checkmarki\checkmarkr\checkmarkc\checkmark$\checkmarke\checkmark$\checkmark^\checkmark\\checkmarkc\checkmarki\checkmarkr\checkmarkc\checkmark$\checkmarkc\checkmark$\checkmark^\checkmark\\checkmarkc\checkmarki\checkmarkr\checkmarkc\checkmark$\checkmarkt\checkmark$\checkmark^\checkmark\\checkmarkc\checkmarki\checkmarkr\checkmarkc\checkmark$\checkmarka\checkmark$\checkmark^\checkmark\\checkmarkc\checkmarki\checkmarkr\checkmarkc\checkmark$\checkmarkn\checkmark$\checkmark^\checkmark\\checkmarkc\checkmarki\checkmarkr\checkmarkc\checkmark$\checkmarkg\checkmark$\checkmark^\checkmark\\checkmarkc\checkmarki\checkmarkr\checkmarkc\checkmark$\checkmarkl\checkmark$\checkmark^\checkmark\\checkmarkc\checkmarki\checkmarkr\checkmarkc\checkmark$\checkmarke\checkmark$\checkmark^\checkmark\\checkmarkc\checkmarki\checkmarkr\checkmarkc\checkmark$\checkmark \checkmark$\checkmark^\checkmark\\checkmarkc\checkmarki\checkmarkr\checkmarkc\checkmark$\checkmark(\checkmark$\checkmark^\checkmark\\checkmarkc\checkmarki\checkmarkr\checkmarkc\checkmark$\checkmark0\checkmark$\checkmark^\checkmark\\checkmarkc\checkmarki\checkmarkr\checkmarkc\checkmark$\checkmark.\checkmark$\checkmark^\checkmark\\checkmarkc\checkmarki\checkmarkr\checkmarkc\checkmark$\checkmark4\checkmark$\checkmark^\checkmark\\checkmarkc\checkmarki\checkmarkr\checkmarkc\checkmark$\checkmark,\checkmark$\checkmark^\checkmark\\checkmarkc\checkmarki\checkmarkr\checkmarkc\checkmark$\checkmark3\checkmark$\checkmark^\checkmark\\checkmarkc\checkmarki\checkmarkr\checkmarkc\checkmark$\checkmark)\checkmark$\checkmark^\checkmark\\checkmarkc\checkmarki\checkmarkr\checkmarkc\checkmark$\checkmark;\checkmark$\checkmark^\checkmark\\checkmarkc\checkmarki\checkmarkr\checkmarkc\checkmark$\checkmark
\checkmark$\checkmark^\checkmark\\checkmarkc\checkmarki\checkmarkr\checkmarkc\checkmark$\checkmark \checkmark$\checkmark^\checkmark\\checkmarkc\checkmarki\checkmarkr\checkmarkc\checkmark$\checkmark \checkmark$\checkmark^\checkmark\\checkmarkc\checkmarki\checkmarkr\checkmarkc\checkmark$\checkmark \checkmark$\checkmark^\checkmark\\checkmarkc\checkmarki\checkmarkr\checkmarkc\checkmark$\checkmark \checkmark$\checkmark^\checkmark\\checkmarkc\checkmarki\checkmarkr\checkmarkc\checkmark$\checkmark \checkmark$\checkmark^\checkmark\\checkmarkc\checkmarki\checkmarkr\checkmarkc\checkmark$\checkmark \checkmark$\checkmark^\checkmark\\checkmarkc\checkmarki\checkmarkr\checkmarkc\checkmark$\checkmark \checkmark$\checkmark^\checkmark\\checkmarkc\checkmarki\checkmarkr\checkmarkc\checkmark$\checkmark \checkmark$\checkmark^\checkmark\\checkmarkc\checkmarki\checkmarkr\checkmarkc\checkmark$\checkmark\\checkmark$\checkmark^\checkmark\\checkmarkc\checkmarki\checkmarkr\checkmarkc\checkmark$\checkmarkn\checkmark$\checkmark^\checkmark\\checkmarkc\checkmarki\checkmarkr\checkmarkc\checkmark$\checkmarko\checkmark$\checkmark^\checkmark\\checkmarkc\checkmarki\checkmarkr\checkmarkc\checkmark$\checkmarkd\checkmark$\checkmark^\checkmark\\checkmarkc\checkmarki\checkmarkr\checkmarkc\checkmark$\checkmarke\checkmark$\checkmark^\checkmark\\checkmarkc\checkmarki\checkmarkr\checkmarkc\checkmark$\checkmark[\checkmark$\checkmark^\checkmark\\checkmarkc\checkmarki\checkmarkr\checkmarkc\checkmark$\checkmarkr\checkmark$\checkmark^\checkmark\\checkmarkc\checkmarki\checkmarkr\checkmarkc\checkmark$\checkmarki\checkmark$\checkmark^\checkmark\\checkmarkc\checkmarki\checkmarkr\checkmarkc\checkmark$\checkmarkg\checkmark$\checkmark^\checkmark\\checkmarkc\checkmarki\checkmarkr\checkmarkc\checkmark$\checkmarkh\checkmark$\checkmark^\checkmark\\checkmarkc\checkmarki\checkmarkr\checkmarkc\checkmark$\checkmarkt\checkmark$\checkmark^\checkmark\\checkmarkc\checkmarki\checkmarkr\checkmarkc\checkmark$\checkmark,\checkmark$\checkmark^\checkmark\\checkmarkc\checkmarki\checkmarkr\checkmarkc\checkmark$\checkmark \checkmark$\checkmark^\checkmark\\checkmarkc\checkmarki\checkmarkr\checkmarkc\checkmark$\checkmarkf\checkmark$\checkmark^\checkmark\\checkmarkc\checkmarki\checkmarkr\checkmarkc\checkmark$\checkmarko\checkmark$\checkmark^\checkmark\\checkmarkc\checkmarki\checkmarkr\checkmarkc\checkmark$\checkmarkn\checkmark$\checkmark^\checkmark\\checkmarkc\checkmarki\checkmarkr\checkmarkc\checkmark$\checkmarkt\checkmark$\checkmark^\checkmark\\checkmarkc\checkmarki\checkmarkr\checkmarkc\checkmark$\checkmark=\checkmark$\checkmark^\checkmark\\checkmarkc\checkmarki\checkmarkr\checkmarkc\checkmark$\checkmark\\checkmark$\checkmark^\checkmark\\checkmarkc\checkmarki\checkmarkr\checkmarkc\checkmark$\checkmarkf\checkmark$\checkmark^\checkmark\\checkmarkc\checkmarki\checkmarkr\checkmarkc\checkmark$\checkmarko\checkmark$\checkmark^\checkmark\\checkmarkc\checkmarki\checkmarkr\checkmarkc\checkmark$\checkmarko\checkmark$\checkmark^\checkmark\\checkmarkc\checkmarki\checkmarkr\checkmarkc\checkmark$\checkmarkt\checkmark$\checkmark^\checkmark\\checkmarkc\checkmarki\checkmarkr\checkmarkc\checkmark$\checkmarkn\checkmark$\checkmark^\checkmark\\checkmarkc\checkmarki\checkmarkr\checkmarkc\checkmark$\checkmarko\checkmark$\checkmark^\checkmark\\checkmarkc\checkmarki\checkmarkr\checkmarkc\checkmark$\checkmarkt\checkmark$\checkmark^\checkmark\\checkmarkc\checkmarki\checkmarkr\checkmarkc\checkmark$\checkmarke\checkmark$\checkmark^\checkmark\\checkmarkc\checkmarki\checkmarkr\checkmarkc\checkmark$\checkmarks\checkmark$\checkmark^\checkmark\\checkmarkc\checkmarki\checkmarkr\checkmarkc\checkmark$\checkmarki\checkmark$\checkmark^\checkmark\\checkmarkc\checkmarki\checkmarkr\checkmarkc\checkmark$\checkmarkz\checkmark$\checkmark^\checkmark\\checkmarkc\checkmarki\checkmarkr\checkmarkc\checkmark$\checkmarke\checkmark$\checkmark^\checkmark\\checkmarkc\checkmarki\checkmarkr\checkmarkc\checkmark$\checkmark]\checkmark$\checkmark^\checkmark\\checkmarkc\checkmarki\checkmarkr\checkmarkc\checkmark$\checkmark \checkmark$\checkmark^\checkmark\\checkmarkc\checkmarki\checkmarkr\checkmarkc\checkmark$\checkmarka\checkmark$\checkmark^\checkmark\\checkmarkc\checkmarki\checkmarkr\checkmarkc\checkmark$\checkmarkt\checkmark$\checkmark^\checkmark\\checkmarkc\checkmarki\checkmarkr\checkmarkc\checkmark$\checkmark \checkmark$\checkmark^\checkmark\\checkmarkc\checkmarki\checkmarkr\checkmarkc\checkmark$\checkmark(\checkmark$\checkmark^\checkmark\\checkmarkc\checkmarki\checkmarkr\checkmarkc\checkmark$\checkmark0\checkmark$\checkmark^\checkmark\\checkmarkc\checkmarki\checkmarkr\checkmarkc\checkmark$\checkmark.\checkmark$\checkmark^\checkmark\\checkmarkc\checkmarki\checkmarkr\checkmarkc\checkmark$\checkmark5\checkmark$\checkmark^\checkmark\\checkmarkc\checkmarki\checkmarkr\checkmarkc\checkmark$\checkmark,\checkmark$\checkmark^\checkmark\\checkmarkc\checkmarki\checkmarkr\checkmarkc\checkmark$\checkmark1\checkmark$\checkmark^\checkmark\\checkmarkc\checkmarki\checkmarkr\checkmarkc\checkmark$\checkmark.\checkmark$\checkmark^\checkmark\\checkmarkc\checkmarki\checkmarkr\checkmarkc\checkmark$\checkmark5\checkmark$\checkmark^\checkmark\\checkmarkc\checkmarki\checkmarkr\checkmarkc\checkmark$\checkmark)\checkmark$\checkmark^\checkmark\\checkmarkc\checkmarki\checkmarkr\checkmarkc\checkmark$\checkmark \checkmark$\checkmark^\checkmark\\checkmarkc\checkmarki\checkmarkr\checkmarkc\checkmark$\checkmark{\checkmark$\checkmark^\checkmark\\checkmarkc\checkmarki\checkmarkr\checkmarkc\checkmark$\checkmarkA\checkmark$\checkmark^\checkmark\\checkmarkc\checkmarki\checkmarkr\checkmarkc\checkmark$\checkmarkr\checkmark$\checkmark^\checkmark\\checkmarkc\checkmarki\checkmarkr\checkmarkc\checkmark$\checkmarkb\checkmark$\checkmark^\checkmark\\checkmarkc\checkmarki\checkmarkr\checkmarkc\checkmark$\checkmarkr\checkmark$\checkmark^\checkmark\\checkmarkc\checkmarki\checkmarkr\checkmarkc\checkmark$\checkmarke\checkmark$\checkmark^\checkmark\\checkmarkc\checkmarki\checkmarkr\checkmarkc\checkmark$\checkmark \checkmark$\checkmark^\checkmark\\checkmarkc\checkmarki\checkmarkr\checkmarkc\checkmark$\checkmarkp\checkmark$\checkmark^\checkmark\\checkmarkc\checkmarki\checkmarkr\checkmarkc\checkmark$\checkmarki\checkmark$\checkmark^\checkmark\\checkmarkc\checkmarki\checkmarkr\checkmarkc\checkmark$\checkmarkv\checkmark$\checkmark^\checkmark\\checkmarkc\checkmarki\checkmarkr\checkmarkc\checkmark$\checkmarko\checkmark$\checkmark^\checkmark\\checkmarkc\checkmarki\checkmarkr\checkmarkc\checkmark$\checkmarkt\checkmark$\checkmark^\checkmark\\checkmarkc\checkmarki\checkmarkr\checkmarkc\checkmark$\checkmark \checkmark$\checkmark^\checkmark\\checkmarkc\checkmarki\checkmarkr\checkmarkc\checkmark$\checkmark\\checkmark$\checkmark^\checkmark\\checkmarkc\checkmarki\checkmarkr\checkmarkc\checkmark$\checkmark\\checkmark$\checkmark^\checkmark\\checkmarkc\checkmarki\checkmarkr\checkmarkc\checkmark$\checkmark \checkmark$\checkmark^\checkmark\\checkmarkc\checkmarki\checkmarkr\checkmarkc\checkmark$\checkmark$\checkmark$\checkmark^\checkmark\\checkmarkc\checkmarki\checkmarkr\checkmarkc\checkmark$\checkmarkH\checkmark$\checkmark^\checkmark\\checkmarkc\checkmarki\checkmarkr\checkmarkc\checkmark$\checkmark \checkmark$\checkmark^\checkmark\\checkmarkc\checkmarki\checkmarkr\checkmarkc\checkmark$\checkmark=\checkmark$\checkmark^\checkmark\\checkmarkc\checkmarki\checkmarkr\checkmarkc\checkmark$\checkmark \checkmark$\checkmark^\checkmark\\checkmarkc\checkmarki\checkmarkr\checkmarkc\checkmark$\checkmark2\checkmark$\checkmark^\checkmark\\checkmarkc\checkmarki\checkmarkr\checkmarkc\checkmark$\checkmark0\checkmark$\checkmark^\checkmark\\checkmarkc\checkmarki\checkmarkr\checkmarkc\checkmark$\checkmark0\checkmark$\checkmark^\checkmark\\checkmarkc\checkmarki\checkmarkr\checkmarkc\checkmark$\checkmark$\checkmark$\checkmark^\checkmark\\checkmarkc\checkmarki\checkmarkr\checkmarkc\checkmark$\checkmark \checkmark$\checkmark^\checkmark\\checkmarkc\checkmarki\checkmarkr\checkmarkc\checkmark$\checkmarkm\checkmark$\checkmark^\checkmark\\checkmarkc\checkmarki\checkmarkr\checkmarkc\checkmark$\checkmarkm\checkmark$\checkmark^\checkmark\\checkmarkc\checkmarki\checkmarkr\checkmarkc\checkmark$\checkmark}\checkmark$\checkmark^\checkmark\\checkmarkc\checkmarki\checkmarkr\checkmarkc\checkmark$\checkmark;\checkmark$\checkmark^\checkmark\\checkmarkc\checkmarki\checkmarkr\checkmarkc\checkmark$\checkmark
\checkmark$\checkmark^\checkmark\\checkmarkc\checkmarki\checkmarkr\checkmarkc\checkmark$\checkmark \checkmark$\checkmark^\checkmark\\checkmarkc\checkmarki\checkmarkr\checkmarkc\checkmark$\checkmark \checkmark$\checkmark^\checkmark\\checkmarkc\checkmarki\checkmarkr\checkmarkc\checkmark$\checkmark \checkmark$\checkmark^\checkmark\\checkmarkc\checkmarki\checkmarkr\checkmarkc\checkmark$\checkmark \checkmark$\checkmark^\checkmark\\checkmarkc\checkmarki\checkmarkr\checkmarkc\checkmark$\checkmark \checkmark$\checkmark^\checkmark\\checkmarkc\checkmarki\checkmarkr\checkmarkc\checkmark$\checkmark \checkmark$\checkmark^\checkmark\\checkmarkc\checkmarki\checkmarkr\checkmarkc\checkmark$\checkmark \checkmark$\checkmark^\checkmark\\checkmarkc\checkmarki\checkmarkr\checkmarkc\checkmark$\checkmark \checkmark$\checkmark^\checkmark\\checkmarkc\checkmarki\checkmarkr\checkmarkc\checkmark$\checkmark\\checkmark$\checkmark^\checkmark\\checkmarkc\checkmarki\checkmarkr\checkmarkc\checkmark$\checkmarkd\checkmark$\checkmark^\checkmark\\checkmarkc\checkmarki\checkmarkr\checkmarkc\checkmark$\checkmarkr\checkmark$\checkmark^\checkmark\\checkmarkc\checkmarki\checkmarkr\checkmarkc\checkmark$\checkmarka\checkmark$\checkmark^\checkmark\\checkmarkc\checkmarki\checkmarkr\checkmarkc\checkmark$\checkmarkw\checkmark$\checkmark^\checkmark\\checkmarkc\checkmarki\checkmarkr\checkmarkc\checkmark$\checkmark[\checkmark$\checkmark^\checkmark\\checkmarkc\checkmarki\checkmarkr\checkmarkc\checkmark$\checkmarka\checkmark$\checkmark^\checkmark\\checkmarkc\checkmarki\checkmarkr\checkmarkc\checkmark$\checkmarkr\checkmark$\checkmark^\checkmark\\checkmarkc\checkmarki\checkmarkr\checkmarkc\checkmark$\checkmarkm\checkmark$\checkmark^\checkmark\\checkmarkc\checkmarki\checkmarkr\checkmarkc\checkmark$\checkmark,\checkmark$\checkmark^\checkmark\\checkmarkc\checkmarki\checkmarkr\checkmarkc\checkmark$\checkmark \checkmark$\checkmark^\checkmark\\checkmarkc\checkmarki\checkmarkr\checkmarkc\checkmark$\checkmarkf\checkmark$\checkmark^\checkmark\\checkmarkc\checkmarki\checkmarkr\checkmarkc\checkmark$\checkmarki\checkmark$\checkmark^\checkmark\\checkmarkc\checkmarki\checkmarkr\checkmarkc\checkmark$\checkmarkl\checkmark$\checkmark^\checkmark\\checkmarkc\checkmarki\checkmarkr\checkmarkc\checkmark$\checkmarkl\checkmark$\checkmark^\checkmark\\checkmarkc\checkmarki\checkmarkr\checkmarkc\checkmark$\checkmark=\checkmark$\checkmark^\checkmark\\checkmarkc\checkmarki\checkmarkr\checkmarkc\checkmark$\checkmarkg\checkmark$\checkmark^\checkmark\\checkmarkc\checkmarki\checkmarkr\checkmarkc\checkmark$\checkmarkr\checkmark$\checkmark^\checkmark\\checkmarkc\checkmarki\checkmarkr\checkmarkc\checkmark$\checkmarka\checkmark$\checkmark^\checkmark\\checkmarkc\checkmarki\checkmarkr\checkmarkc\checkmark$\checkmarky\checkmark$\checkmark^\checkmark\\checkmarkc\checkmarki\checkmarkr\checkmarkc\checkmark$\checkmark!\checkmark$\checkmark^\checkmark\\checkmarkc\checkmarki\checkmarkr\checkmarkc\checkmark$\checkmark2\checkmark$\checkmark^\checkmark\\checkmarkc\checkmarki\checkmarkr\checkmarkc\checkmark$\checkmark0\checkmark$\checkmark^\checkmark\\checkmarkc\checkmarki\checkmarkr\checkmarkc\checkmark$\checkmark]\checkmark$\checkmark^\checkmark\\checkmarkc\checkmarki\checkmarkr\checkmarkc\checkmark$\checkmark \checkmark$\checkmark^\checkmark\\checkmarkc\checkmarki\checkmarkr\checkmarkc\checkmark$\checkmark(\checkmark$\checkmark^\checkmark\\checkmarkc\checkmarki\checkmarkr\checkmarkc\checkmark$\checkmark-\checkmark$\checkmark^\checkmark\\checkmarkc\checkmarki\checkmarkr\checkmarkc\checkmark$\checkmark2\checkmark$\checkmark^\checkmark\\checkmarkc\checkmarki\checkmarkr\checkmarkc\checkmark$\checkmark,\checkmark$\checkmark^\checkmark\\checkmarkc\checkmarki\checkmarkr\checkmarkc\checkmark$\checkmark3\checkmark$\checkmark^\checkmark\\checkmarkc\checkmarki\checkmarkr\checkmarkc\checkmark$\checkmark)\checkmark$\checkmark^\checkmark\\checkmarkc\checkmarki\checkmarkr\checkmarkc\checkmark$\checkmark \checkmark$\checkmark^\checkmark\\checkmarkc\checkmarki\checkmarkr\checkmarkc\checkmark$\checkmarkr\checkmark$\checkmark^\checkmark\\checkmarkc\checkmarki\checkmarkr\checkmarkc\checkmark$\checkmarke\checkmark$\checkmark^\checkmark\\checkmarkc\checkmarki\checkmarkr\checkmarkc\checkmark$\checkmarkc\checkmark$\checkmark^\checkmark\\checkmarkc\checkmarki\checkmarkr\checkmarkc\checkmark$\checkmarkt\checkmark$\checkmark^\checkmark\\checkmarkc\checkmarki\checkmarkr\checkmarkc\checkmark$\checkmarka\checkmark$\checkmark^\checkmark\\checkmarkc\checkmarki\checkmarkr\checkmarkc\checkmark$\checkmarkn\checkmark$\checkmark^\checkmark\\checkmarkc\checkmarki\checkmarkr\checkmarkc\checkmark$\checkmarkg\checkmark$\checkmark^\checkmark\\checkmarkc\checkmarki\checkmarkr\checkmarkc\checkmark$\checkmarkl\checkmark$\checkmark^\checkmark\\checkmarkc\checkmarki\checkmarkr\checkmarkc\checkmark$\checkmarke\checkmark$\checkmark^\checkmark\\checkmarkc\checkmarki\checkmarkr\checkmarkc\checkmark$\checkmark \checkmark$\checkmark^\checkmark\\checkmarkc\checkmarki\checkmarkr\checkmarkc\checkmark$\checkmark(\checkmark$\checkmark^\checkmark\\checkmarkc\checkmarki\checkmarkr\checkmarkc\checkmark$\checkmark2\checkmark$\checkmark^\checkmark\\checkmarkc\checkmarki\checkmarkr\checkmarkc\checkmark$\checkmark,\checkmark$\checkmark^\checkmark\\checkmarkc\checkmarki\checkmarkr\checkmarkc\checkmark$\checkmark3\checkmark$\checkmark^\checkmark\\checkmarkc\checkmarki\checkmarkr\checkmarkc\checkmark$\checkmark.\checkmark$\checkmark^\checkmark\\checkmarkc\checkmarki\checkmarkr\checkmarkc\checkmark$\checkmark4\checkmark$\checkmark^\checkmark\\checkmarkc\checkmarki\checkmarkr\checkmarkc\checkmark$\checkmark)\checkmark$\checkmark^\checkmark\\checkmarkc\checkmarki\checkmarkr\checkmarkc\checkmark$\checkmark;\checkmark$\checkmark^\checkmark\\checkmarkc\checkmarki\checkmarkr\checkmarkc\checkmark$\checkmark
\checkmark$\checkmark^\checkmark\\checkmarkc\checkmarki\checkmarkr\checkmarkc\checkmark$\checkmark \checkmark$\checkmark^\checkmark\\checkmarkc\checkmarki\checkmarkr\checkmarkc\checkmark$\checkmark \checkmark$\checkmark^\checkmark\\checkmarkc\checkmarki\checkmarkr\checkmarkc\checkmark$\checkmark \checkmark$\checkmark^\checkmark\\checkmarkc\checkmarki\checkmarkr\checkmarkc\checkmark$\checkmark \checkmark$\checkmark^\checkmark\\checkmarkc\checkmarki\checkmarkr\checkmarkc\checkmark$\checkmark \checkmark$\checkmark^\checkmark\\checkmarkc\checkmarki\checkmarkr\checkmarkc\checkmark$\checkmark \checkmark$\checkmark^\checkmark\\checkmarkc\checkmarki\checkmarkr\checkmarkc\checkmark$\checkmark \checkmark$\checkmark^\checkmark\\checkmarkc\checkmarki\checkmarkr\checkmarkc\checkmark$\checkmark \checkmark$\checkmark^\checkmark\\checkmarkc\checkmarki\checkmarkr\checkmarkc\checkmark$\checkmark\\checkmark$\checkmark^\checkmark\\checkmarkc\checkmarki\checkmarkr\checkmarkc\checkmark$\checkmarkn\checkmark$\checkmark^\checkmark\\checkmarkc\checkmarki\checkmarkr\checkmarkc\checkmark$\checkmarko\checkmark$\checkmark^\checkmark\\checkmarkc\checkmarki\checkmarkr\checkmarkc\checkmark$\checkmarkd\checkmark$\checkmark^\checkmark\\checkmarkc\checkmarki\checkmarkr\checkmarkc\checkmark$\checkmarke\checkmark$\checkmark^\checkmark\\checkmarkc\checkmarki\checkmarkr\checkmarkc\checkmark$\checkmark \checkmark$\checkmark^\checkmark\\checkmarkc\checkmarki\checkmarkr\checkmarkc\checkmark$\checkmarka\checkmark$\checkmark^\checkmark\\checkmarkc\checkmarki\checkmarkr\checkmarkc\checkmark$\checkmarkt\checkmark$\checkmark^\checkmark\\checkmarkc\checkmarki\checkmarkr\checkmarkc\checkmark$\checkmark \checkmark$\checkmark^\checkmark\\checkmarkc\checkmarki\checkmarkr\checkmarkc\checkmark$\checkmark(\checkmark$\checkmark^\checkmark\\checkmarkc\checkmarki\checkmarkr\checkmarkc\checkmark$\checkmark0\checkmark$\checkmark^\checkmark\\checkmarkc\checkmarki\checkmarkr\checkmarkc\checkmark$\checkmark,\checkmark$\checkmark^\checkmark\\checkmarkc\checkmarki\checkmarkr\checkmarkc\checkmark$\checkmark3\checkmark$\checkmark^\checkmark\\checkmarkc\checkmarki\checkmarkr\checkmarkc\checkmark$\checkmark.\checkmark$\checkmark^\checkmark\\checkmarkc\checkmarki\checkmarkr\checkmarkc\checkmark$\checkmark2\checkmark$\checkmark^\checkmark\\checkmarkc\checkmarki\checkmarkr\checkmarkc\checkmark$\checkmark)\checkmark$\checkmark^\checkmark\\checkmarkc\checkmarki\checkmarkr\checkmarkc\checkmark$\checkmark \checkmark$\checkmark^\checkmark\\checkmarkc\checkmarki\checkmarkr\checkmarkc\checkmark$\checkmark{\checkmark$\checkmark^\checkmark\\checkmarkc\checkmarki\checkmarkr\checkmarkc\checkmark$\checkmark\\checkmark$\checkmark^\checkmark\\checkmarkc\checkmarki\checkmarkr\checkmarkc\checkmark$\checkmarks\checkmark$\checkmark^\checkmark\\checkmarkc\checkmarki\checkmarkr\checkmarkc\checkmark$\checkmarkm\checkmark$\checkmark^\checkmark\\checkmarkc\checkmarki\checkmarkr\checkmarkc\checkmark$\checkmarka\checkmark$\checkmark^\checkmark\\checkmarkc\checkmarki\checkmarkr\checkmarkc\checkmark$\checkmarkl\checkmark$\checkmark^\checkmark\\checkmarkc\checkmarki\checkmarkr\checkmarkc\checkmark$\checkmarkl\checkmark$\checkmark^\checkmark\\checkmarkc\checkmarki\checkmarkr\checkmarkc\checkmark$\checkmark \checkmark$\checkmark^\checkmark\\checkmarkc\checkmarki\checkmarkr\checkmarkc\checkmark$\checkmarkP\checkmark$\checkmark^\checkmark\\checkmarkc\checkmarki\checkmarkr\checkmarkc\checkmark$\checkmarkl\checkmark$\checkmark^\checkmark\\checkmarkc\checkmarki\checkmarkr\checkmarkc\checkmark$\checkmarka\checkmark$\checkmark^\checkmark\\checkmarkc\checkmarki\checkmarkr\checkmarkc\checkmark$\checkmarkt\checkmark$\checkmark^\checkmark\\checkmarkc\checkmarki\checkmarkr\checkmarkc\checkmark$\checkmarke\checkmark$\checkmark^\checkmark\\checkmarkc\checkmarki\checkmarkr\checkmarkc\checkmark$\checkmarka\checkmark$\checkmark^\checkmark\\checkmarkc\checkmarki\checkmarkr\checkmarkc\checkmark$\checkmarku\checkmark$\checkmark^\checkmark\\checkmarkc\checkmarki\checkmarkr\checkmarkc\checkmark$\checkmark \checkmark$\checkmark^\checkmark\\checkmarkc\checkmarki\checkmarkr\checkmarkc\checkmark$\checkmarkB\checkmark$\checkmark^\checkmark\\checkmarkc\checkmarki\checkmarkr\checkmarkc\checkmark$\checkmark \checkmark$\checkmark^\checkmark\\checkmarkc\checkmarki\checkmarkr\checkmarkc\checkmark$\checkmark+\checkmark$\checkmark^\checkmark\\checkmarkc\checkmarki\checkmarkr\checkmarkc\checkmark$\checkmark \checkmark$\checkmark^\checkmark\\checkmarkc\checkmarki\checkmarkr\checkmarkc\checkmark$\checkmarkP\checkmark$\checkmark^\checkmark\\checkmarkc\checkmarki\checkmarkr\checkmarkc\checkmark$\checkmarki\checkmark$\checkmark^\checkmark\\checkmarkc\checkmarki\checkmarkr\checkmarkc\checkmark$\checkmarkè\checkmark$\checkmark^\checkmark\\checkmarkc\checkmarki\checkmarkr\checkmarkc\checkmark$\checkmarkc\checkmark$\checkmark^\checkmark\\checkmarkc\checkmarki\checkmarkr\checkmarkc\checkmark$\checkmarke\checkmark$\checkmark^\checkmark\\checkmarkc\checkmarki\checkmarkr\checkmarkc\checkmark$\checkmark}\checkmark$\checkmark^\checkmark\\checkmarkc\checkmarki\checkmarkr\checkmarkc\checkmark$\checkmark;\checkmark$\checkmark^\checkmark\\checkmarkc\checkmarki\checkmarkr\checkmarkc\checkmark$\checkmark
\checkmark$\checkmark^\checkmark\\checkmarkc\checkmarki\checkmarkr\checkmarkc\checkmark$\checkmark \checkmark$\checkmark^\checkmark\\checkmarkc\checkmarki\checkmarkr\checkmarkc\checkmark$\checkmark \checkmark$\checkmark^\checkmark\\checkmarkc\checkmarki\checkmarkr\checkmarkc\checkmark$\checkmark \checkmark$\checkmark^\checkmark\\checkmarkc\checkmarki\checkmarkr\checkmarkc\checkmark$\checkmark \checkmark$\checkmark^\checkmark\\checkmarkc\checkmarki\checkmarkr\checkmarkc\checkmark$\checkmark \checkmark$\checkmark^\checkmark\\checkmarkc\checkmarki\checkmarkr\checkmarkc\checkmark$\checkmark \checkmark$\checkmark^\checkmark\\checkmarkc\checkmarki\checkmarkr\checkmarkc\checkmark$\checkmark \checkmark$\checkmark^\checkmark\\checkmarkc\checkmarki\checkmarkr\checkmarkc\checkmark$\checkmark \checkmark$\checkmark^\checkmark\\checkmarkc\checkmarki\checkmarkr\checkmarkc\checkmark$\checkmark\\checkmark$\checkmark^\checkmark\\checkmarkc\checkmarki\checkmarkr\checkmarkc\checkmark$\checkmarkd\checkmark$\checkmark^\checkmark\\checkmarkc\checkmarki\checkmarkr\checkmarkc\checkmark$\checkmarkr\checkmark$\checkmark^\checkmark\\checkmarkc\checkmarki\checkmarkr\checkmarkc\checkmark$\checkmarka\checkmark$\checkmark^\checkmark\\checkmarkc\checkmarki\checkmarkr\checkmarkc\checkmark$\checkmarkw\checkmark$\checkmark^\checkmark\\checkmarkc\checkmarki\checkmarkr\checkmarkc\checkmark$\checkmark[\checkmark$\checkmark^\checkmark\\checkmarkc\checkmarki\checkmarkr\checkmarkc\checkmark$\checkmarkf\checkmark$\checkmark^\checkmark\\checkmarkc\checkmarki\checkmarkr\checkmarkc\checkmark$\checkmarko\checkmark$\checkmark^\checkmark\\checkmarkc\checkmarki\checkmarkr\checkmarkc\checkmark$\checkmarkr\checkmark$\checkmark^\checkmark\\checkmarkc\checkmarki\checkmarkr\checkmarkc\checkmark$\checkmarkc\checkmark$\checkmark^\checkmark\\checkmarkc\checkmarki\checkmarkr\checkmarkc\checkmark$\checkmarke\checkmark$\checkmark^\checkmark\\checkmarkc\checkmarki\checkmarkr\checkmarkc\checkmark$\checkmark]\checkmark$\checkmark^\checkmark\\checkmarkc\checkmarki\checkmarkr\checkmarkc\checkmark$\checkmark \checkmark$\checkmark^\checkmark\\checkmarkc\checkmarki\checkmarkr\checkmarkc\checkmark$\checkmark(\checkmark$\checkmark^\checkmark\\checkmarkc\checkmarki\checkmarkr\checkmarkc\checkmark$\checkmark-\checkmark$\checkmark^\checkmark\\checkmarkc\checkmarki\checkmarkr\checkmarkc\checkmark$\checkmark3\checkmark$\checkmark^\checkmark\\checkmarkc\checkmarki\checkmarkr\checkmarkc\checkmark$\checkmark,\checkmark$\checkmark^\checkmark\\checkmarkc\checkmarki\checkmarkr\checkmarkc\checkmark$\checkmark3\checkmark$\checkmark^\checkmark\\checkmarkc\checkmarki\checkmarkr\checkmarkc\checkmark$\checkmark.\checkmark$\checkmark^\checkmark\\checkmarkc\checkmarki\checkmarkr\checkmarkc\checkmark$\checkmark4\checkmark$\checkmark^\checkmark\\checkmarkc\checkmarki\checkmarkr\checkmarkc\checkmark$\checkmark)\checkmark$\checkmark^\checkmark\\checkmarkc\checkmarki\checkmarkr\checkmarkc\checkmark$\checkmark \checkmark$\checkmark^\checkmark\\checkmarkc\checkmarki\checkmarkr\checkmarkc\checkmark$\checkmark-\checkmark$\checkmark^\checkmark\\checkmarkc\checkmarki\checkmarkr\checkmarkc\checkmark$\checkmark-\checkmark$\checkmark^\checkmark\\checkmarkc\checkmarki\checkmarkr\checkmarkc\checkmark$\checkmark \checkmark$\checkmark^\checkmark\\checkmarkc\checkmarki\checkmarkr\checkmarkc\checkmark$\checkmark(\checkmark$\checkmark^\checkmark\\checkmarkc\checkmarki\checkmarkr\checkmarkc\checkmark$\checkmark-\checkmark$\checkmark^\checkmark\\checkmarkc\checkmarki\checkmarkr\checkmarkc\checkmark$\checkmark1\checkmark$\checkmark^\checkmark\\checkmarkc\checkmarki\checkmarkr\checkmarkc\checkmark$\checkmark.\checkmark$\checkmark^\checkmark\\checkmarkc\checkmarki\checkmarkr\checkmarkc\checkmark$\checkmark5\checkmark$\checkmark^\checkmark\\checkmarkc\checkmarki\checkmarkr\checkmarkc\checkmark$\checkmark,\checkmark$\checkmark^\checkmark\\checkmarkc\checkmarki\checkmarkr\checkmarkc\checkmark$\checkmark3\checkmark$\checkmark^\checkmark\\checkmarkc\checkmarki\checkmarkr\checkmarkc\checkmark$\checkmark.\checkmark$\checkmark^\checkmark\\checkmarkc\checkmarki\checkmarkr\checkmarkc\checkmark$\checkmark4\checkmark$\checkmark^\checkmark\\checkmarkc\checkmarki\checkmarkr\checkmarkc\checkmark$\checkmark)\checkmark$\checkmark^\checkmark\\checkmarkc\checkmarki\checkmarkr\checkmarkc\checkmark$\checkmark \checkmark$\checkmark^\checkmark\\checkmarkc\checkmarki\checkmarkr\checkmarkc\checkmark$\checkmarkn\checkmark$\checkmark^\checkmark\\checkmarkc\checkmarki\checkmarkr\checkmarkc\checkmark$\checkmarko\checkmark$\checkmark^\checkmark\\checkmarkc\checkmarki\checkmarkr\checkmarkc\checkmark$\checkmarkd\checkmark$\checkmark^\checkmark\\checkmarkc\checkmarki\checkmarkr\checkmarkc\checkmark$\checkmarke\checkmark$\checkmark^\checkmark\\checkmarkc\checkmarki\checkmarkr\checkmarkc\checkmark$\checkmark[\checkmark$\checkmark^\checkmark\\checkmarkc\checkmarki\checkmarkr\checkmarkc\checkmark$\checkmarkm\checkmark$\checkmark^\checkmark\\checkmarkc\checkmarki\checkmarkr\checkmarkc\checkmark$\checkmarki\checkmark$\checkmark^\checkmark\\checkmarkc\checkmarki\checkmarkr\checkmarkc\checkmark$\checkmarkd\checkmark$\checkmark^\checkmark\\checkmarkc\checkmarki\checkmarkr\checkmarkc\checkmark$\checkmarkw\checkmark$\checkmark^\checkmark\\checkmarkc\checkmarki\checkmarkr\checkmarkc\checkmark$\checkmarka\checkmark$\checkmark^\checkmark\\checkmarkc\checkmarki\checkmarkr\checkmarkc\checkmark$\checkmarky\checkmark$\checkmark^\checkmark\\checkmarkc\checkmarki\checkmarkr\checkmarkc\checkmark$\checkmark,\checkmark$\checkmark^\checkmark\\checkmarkc\checkmarki\checkmarkr\checkmarkc\checkmark$\checkmarka\checkmark$\checkmark^\checkmark\\checkmarkc\checkmarki\checkmarkr\checkmarkc\checkmark$\checkmarkb\checkmark$\checkmark^\checkmark\\checkmarkc\checkmarki\checkmarkr\checkmarkc\checkmark$\checkmarko\checkmark$\checkmark^\checkmark\\checkmarkc\checkmarki\checkmarkr\checkmarkc\checkmark$\checkmarkv\checkmark$\checkmark^\checkmark\\checkmarkc\checkmarki\checkmarkr\checkmarkc\checkmark$\checkmarke\checkmark$\checkmark^\checkmark\\checkmarkc\checkmarki\checkmarkr\checkmarkc\checkmark$\checkmark]\checkmark$\checkmark^\checkmark\\checkmarkc\checkmarki\checkmarkr\checkmarkc\checkmark$\checkmark \checkmark$\checkmark^\checkmark\\checkmarkc\checkmarki\checkmarkr\checkmarkc\checkmark$\checkmark{\checkmark$\checkmark^\checkmark\\checkmarkc\checkmarki\checkmarkr\checkmarkc\checkmark$\checkmark$\checkmark$\checkmark^\checkmark\\checkmarkc\checkmarki\checkmarkr\checkmarkc\checkmark$\checkmarkF\checkmark$\checkmark^\checkmark\\checkmarkc\checkmarki\checkmarkr\checkmarkc\checkmark$\checkmark_\checkmark$\checkmark^\checkmark\\checkmarkc\checkmarki\checkmarkr\checkmarkc\checkmark$\checkmark{\checkmark$\checkmark^\checkmark\\checkmarkc\checkmarki\checkmarkr\checkmarkc\checkmark$\checkmarkm\checkmark$\checkmark^\checkmark\\checkmarkc\checkmarki\checkmarkr\checkmarkc\checkmark$\checkmarka\checkmark$\checkmark^\checkmark\\checkmarkc\checkmarki\checkmarkr\checkmarkc\checkmark$\checkmarkx\checkmark$\checkmark^\checkmark\\checkmarkc\checkmarki\checkmarkr\checkmarkc\checkmark$\checkmark}\checkmark$\checkmark^\checkmark\\checkmarkc\checkmarki\checkmarkr\checkmarkc\checkmark$\checkmark=\checkmark$\checkmark^\checkmark\\checkmarkc\checkmarki\checkmarkr\checkmarkc\checkmark$\checkmark1\checkmark$\checkmark^\checkmark\\checkmarkc\checkmarki\checkmarkr\checkmarkc\checkmark$\checkmark0\checkmark$\checkmark^\checkmark\\checkmarkc\checkmarki\checkmarkr\checkmarkc\checkmark$\checkmark0\checkmark$\checkmark^\checkmark\\checkmarkc\checkmarki\checkmarkr\checkmarkc\checkmark$\checkmark$\checkmark$\checkmark^\checkmark\\checkmarkc\checkmarki\checkmarkr\checkmarkc\checkmark$\checkmark \checkmark$\checkmark^\checkmark\\checkmarkc\checkmarki\checkmarkr\checkmarkc\checkmark$\checkmarkN\checkmark$\checkmark^\checkmark\\checkmarkc\checkmarki\checkmarkr\checkmarkc\checkmark$\checkmark}\checkmark$\checkmark^\checkmark\\checkmarkc\checkmarki\checkmarkr\checkmarkc\checkmark$\checkmark;\checkmark$\checkmark^\checkmark\\checkmarkc\checkmarki\checkmarkr\checkmarkc\checkmark$\checkmark
\checkmark$\checkmark^\checkmark\\checkmarkc\checkmarki\checkmarkr\checkmarkc\checkmark$\checkmark \checkmark$\checkmark^\checkmark\\checkmarkc\checkmarki\checkmarkr\checkmarkc\checkmark$\checkmark \checkmark$\checkmark^\checkmark\\checkmarkc\checkmarki\checkmarkr\checkmarkc\checkmark$\checkmark \checkmark$\checkmark^\checkmark\\checkmarkc\checkmarki\checkmarkr\checkmarkc\checkmark$\checkmark \checkmark$\checkmark^\checkmark\\checkmarkc\checkmarki\checkmarkr\checkmarkc\checkmark$\checkmark \checkmark$\checkmark^\checkmark\\checkmarkc\checkmarki\checkmarkr\checkmarkc\checkmark$\checkmark \checkmark$\checkmark^\checkmark\\checkmarkc\checkmarki\checkmarkr\checkmarkc\checkmark$\checkmark \checkmark$\checkmark^\checkmark\\checkmarkc\checkmarki\checkmarkr\checkmarkc\checkmark$\checkmark \checkmark$\checkmark^\checkmark\\checkmarkc\checkmarki\checkmarkr\checkmarkc\checkmark$\checkmark\\checkmark$\checkmark^\checkmark\\checkmarkc\checkmarki\checkmarkr\checkmarkc\checkmark$\checkmarkd\checkmark$\checkmark^\checkmark\\checkmarkc\checkmarki\checkmarkr\checkmarkc\checkmark$\checkmarkr\checkmark$\checkmark^\checkmark\\checkmarkc\checkmarki\checkmarkr\checkmarkc\checkmark$\checkmarka\checkmark$\checkmark^\checkmark\\checkmarkc\checkmarki\checkmarkr\checkmarkc\checkmark$\checkmarkw\checkmark$\checkmark^\checkmark\\checkmarkc\checkmarki\checkmarkr\checkmarkc\checkmark$\checkmark[\checkmark$\checkmark^\checkmark\\checkmarkc\checkmarki\checkmarkr\checkmarkc\checkmark$\checkmarkt\checkmark$\checkmark^\checkmark\\checkmarkc\checkmarki\checkmarkr\checkmarkc\checkmark$\checkmarkh\checkmark$\checkmark^\checkmark\\checkmarkc\checkmarki\checkmarkr\checkmarkc\checkmark$\checkmarki\checkmark$\checkmark^\checkmark\\checkmarkc\checkmarki\checkmarkr\checkmarkc\checkmark$\checkmarkc\checkmark$\checkmark^\checkmark\\checkmarkc\checkmarki\checkmarkr\checkmarkc\checkmark$\checkmarkk\checkmark$\checkmark^\checkmark\\checkmarkc\checkmarki\checkmarkr\checkmarkc\checkmark$\checkmark,\checkmark$\checkmark^\checkmark\\checkmarkc\checkmarki\checkmarkr\checkmarkc\checkmark$\checkmark \checkmark$\checkmark^\checkmark\\checkmarkc\checkmarki\checkmarkr\checkmarkc\checkmark$\checkmarkb\checkmark$\checkmark^\checkmark\\checkmarkc\checkmarki\checkmarkr\checkmarkc\checkmark$\checkmarkl\checkmark$\checkmark^\checkmark\\checkmarkc\checkmarki\checkmarkr\checkmarkc\checkmark$\checkmarku\checkmark$\checkmark^\checkmark\\checkmarkc\checkmarki\checkmarkr\checkmarkc\checkmark$\checkmarke\checkmark$\checkmark^\checkmark\\checkmarkc\checkmarki\checkmarkr\checkmarkc\checkmark$\checkmark,\checkmark$\checkmark^\checkmark\\checkmarkc\checkmarki\checkmarkr\checkmarkc\checkmark$\checkmark \checkmark$\checkmark^\checkmark\\checkmarkc\checkmarki\checkmarkr\checkmarkc\checkmark$\checkmark-\checkmark$\checkmark^\checkmark\\checkmarkc\checkmarki\checkmarkr\checkmarkc\checkmark$\checkmark{\checkmark$\checkmark^\checkmark\\checkmarkc\checkmarki\checkmarkr\checkmarkc\checkmark$\checkmarkS\checkmark$\checkmark^\checkmark\\checkmarkc\checkmarki\checkmarkr\checkmarkc\checkmark$\checkmarkt\checkmark$\checkmark^\checkmark\\checkmarkc\checkmarki\checkmarkr\checkmarkc\checkmark$\checkmarke\checkmark$\checkmark^\checkmark\\checkmarkc\checkmarki\checkmarkr\checkmarkc\checkmark$\checkmarka\checkmark$\checkmark^\checkmark\\checkmarkc\checkmarki\checkmarkr\checkmarkc\checkmark$\checkmarkl\checkmark$\checkmark^\checkmark\\checkmarkc\checkmarki\checkmarkr\checkmarkc\checkmark$\checkmarkt\checkmark$\checkmark^\checkmark\\checkmarkc\checkmarki\checkmarkr\checkmarkc\checkmark$\checkmarkh\checkmark$\checkmark^\checkmark\\checkmarkc\checkmarki\checkmarkr\checkmarkc\checkmark$\checkmark[\checkmark$\checkmark^\checkmark\\checkmarkc\checkmarki\checkmarkr\checkmarkc\checkmark$\checkmarks\checkmark$\checkmark^\checkmark\\checkmarkc\checkmarki\checkmarkr\checkmarkc\checkmark$\checkmarkc\checkmark$\checkmark^\checkmark\\checkmarkc\checkmarki\checkmarkr\checkmarkc\checkmark$\checkmarka\checkmark$\checkmark^\checkmark\\checkmarkc\checkmarki\checkmarkr\checkmarkc\checkmark$\checkmarkl\checkmark$\checkmark^\checkmark\\checkmarkc\checkmarki\checkmarkr\checkmarkc\checkmark$\checkmarke\checkmark$\checkmark^\checkmark\\checkmarkc\checkmarki\checkmarkr\checkmarkc\checkmark$\checkmark=\checkmark$\checkmark^\checkmark\\checkmarkc\checkmarki\checkmarkr\checkmarkc\checkmark$\checkmark1\checkmark$\checkmark^\checkmark\\checkmarkc\checkmarki\checkmarkr\checkmarkc\checkmark$\checkmark.\checkmark$\checkmark^\checkmark\\checkmarkc\checkmarki\checkmarkr\checkmarkc\checkmark$\checkmark2\checkmark$\checkmark^\checkmark\\checkmarkc\checkmarki\checkmarkr\checkmarkc\checkmark$\checkmark]\checkmark$\checkmark^\checkmark\\checkmarkc\checkmarki\checkmarkr\checkmarkc\checkmark$\checkmark}\checkmark$\checkmark^\checkmark\\checkmarkc\checkmarki\checkmarkr\checkmarkc\checkmark$\checkmark]\checkmark$\checkmark^\checkmark\\checkmarkc\checkmarki\checkmarkr\checkmarkc\checkmark$\checkmark \checkmark$\checkmark^\checkmark\\checkmarkc\checkmarki\checkmarkr\checkmarkc\checkmark$\checkmark(\checkmark$\checkmark^\checkmark\\checkmarkc\checkmarki\checkmarkr\checkmarkc\checkmark$\checkmark-\checkmark$\checkmark^\checkmark\\checkmarkc\checkmarki\checkmarkr\checkmarkc\checkmark$\checkmark1\checkmark$\checkmark^\checkmark\\checkmarkc\checkmarki\checkmarkr\checkmarkc\checkmark$\checkmark,\checkmark$\checkmark^\checkmark\\checkmarkc\checkmarki\checkmarkr\checkmarkc\checkmark$\checkmark0\checkmark$\checkmark^\checkmark\\checkmarkc\checkmarki\checkmarkr\checkmarkc\checkmark$\checkmark.\checkmark$\checkmark^\checkmark\\checkmarkc\checkmarki\checkmarkr\checkmarkc\checkmark$\checkmark5\checkmark$\checkmark^\checkmark\\checkmarkc\checkmarki\checkmarkr\checkmarkc\checkmark$\checkmark)\checkmark$\checkmark^\checkmark\\checkmarkc\checkmarki\checkmarkr\checkmarkc\checkmark$\checkmark \checkmark$\checkmark^\checkmark\\checkmarkc\checkmarki\checkmarkr\checkmarkc\checkmark$\checkmarka\checkmark$\checkmark^\checkmark\\checkmarkc\checkmarki\checkmarkr\checkmarkc\checkmark$\checkmarkr\checkmark$\checkmark^\checkmark\\checkmarkc\checkmarki\checkmarkr\checkmarkc\checkmark$\checkmarkc\checkmark$\checkmark^\checkmark\\checkmarkc\checkmarki\checkmarkr\checkmarkc\checkmark$\checkmark[\checkmark$\checkmark^\checkmark\\checkmarkc\checkmarki\checkmarkr\checkmarkc\checkmark$\checkmarks\checkmark$\checkmark^\checkmark\\checkmarkc\checkmarki\checkmarkr\checkmarkc\checkmark$\checkmarkt\checkmark$\checkmark^\checkmark\\checkmarkc\checkmarki\checkmarkr\checkmarkc\checkmark$\checkmarka\checkmark$\checkmark^\checkmark\\checkmarkc\checkmarki\checkmarkr\checkmarkc\checkmark$\checkmarkr\checkmark$\checkmark^\checkmark\\checkmarkc\checkmarki\checkmarkr\checkmarkc\checkmark$\checkmarkt\checkmark$\checkmark^\checkmark\\checkmarkc\checkmarki\checkmarkr\checkmarkc\checkmark$\checkmark \checkmark$\checkmark^\checkmark\\checkmarkc\checkmarki\checkmarkr\checkmarkc\checkmark$\checkmarka\checkmark$\checkmark^\checkmark\\checkmarkc\checkmarki\checkmarkr\checkmarkc\checkmark$\checkmarkn\checkmark$\checkmark^\checkmark\\checkmarkc\checkmarki\checkmarkr\checkmarkc\checkmark$\checkmarkg\checkmark$\checkmark^\checkmark\\checkmarkc\checkmarki\checkmarkr\checkmarkc\checkmark$\checkmarkl\checkmark$\checkmark^\checkmark\\checkmarkc\checkmarki\checkmarkr\checkmarkc\checkmark$\checkmarke\checkmark$\checkmark^\checkmark\\checkmarkc\checkmarki\checkmarkr\checkmarkc\checkmark$\checkmark=\checkmark$\checkmark^\checkmark\\checkmarkc\checkmarki\checkmarkr\checkmarkc\checkmark$\checkmark1\checkmark$\checkmark^\checkmark\\checkmarkc\checkmarki\checkmarkr\checkmarkc\checkmark$\checkmark5\checkmark$\checkmark^\checkmark\\checkmarkc\checkmarki\checkmarkr\checkmarkc\checkmark$\checkmark0\checkmark$\checkmark^\checkmark\\checkmarkc\checkmarki\checkmarkr\checkmarkc\checkmark$\checkmark,\checkmark$\checkmark^\checkmark\\checkmarkc\checkmarki\checkmarkr\checkmarkc\checkmark$\checkmarke\checkmark$\checkmark^\checkmark\\checkmarkc\checkmarki\checkmarkr\checkmarkc\checkmark$\checkmarkn\checkmark$\checkmark^\checkmark\\checkmarkc\checkmarki\checkmarkr\checkmarkc\checkmark$\checkmarkd\checkmark$\checkmark^\checkmark\\checkmarkc\checkmarki\checkmarkr\checkmarkc\checkmark$\checkmark \checkmark$\checkmark^\checkmark\\checkmarkc\checkmarki\checkmarkr\checkmarkc\checkmark$\checkmarka\checkmark$\checkmark^\checkmark\\checkmarkc\checkmarki\checkmarkr\checkmarkc\checkmark$\checkmarkn\checkmark$\checkmark^\checkmark\\checkmarkc\checkmarki\checkmarkr\checkmarkc\checkmark$\checkmarkg\checkmark$\checkmark^\checkmark\\checkmarkc\checkmarki\checkmarkr\checkmarkc\checkmark$\checkmarkl\checkmark$\checkmark^\checkmark\\checkmarkc\checkmarki\checkmarkr\checkmarkc\checkmark$\checkmarke\checkmark$\checkmark^\checkmark\\checkmarkc\checkmarki\checkmarkr\checkmarkc\checkmark$\checkmark=\checkmark$\checkmark^\checkmark\\checkmarkc\checkmarki\checkmarkr\checkmarkc\checkmark$\checkmark3\checkmark$\checkmark^\checkmark\\checkmarkc\checkmarki\checkmarkr\checkmarkc\checkmark$\checkmark0\checkmark$\checkmark^\checkmark\\checkmarkc\checkmarki\checkmarkr\checkmarkc\checkmark$\checkmark,\checkmark$\checkmark^\checkmark\\checkmarkc\checkmarki\checkmarkr\checkmarkc\checkmark$\checkmarkr\checkmark$\checkmark^\checkmark\\checkmarkc\checkmarki\checkmarkr\checkmarkc\checkmark$\checkmarka\checkmark$\checkmark^\checkmark\\checkmarkc\checkmarki\checkmarkr\checkmarkc\checkmark$\checkmarkd\checkmark$\checkmark^\checkmark\\checkmarkc\checkmarki\checkmarkr\checkmarkc\checkmark$\checkmarki\checkmark$\checkmark^\checkmark\\checkmarkc\checkmarki\checkmarkr\checkmarkc\checkmark$\checkmarku\checkmark$\checkmark^\checkmark\\checkmarkc\checkmarki\checkmarkr\checkmarkc\checkmark$\checkmarks\checkmark$\checkmark^\checkmark\\checkmarkc\checkmarki\checkmarkr\checkmarkc\checkmark$\checkmark=\checkmark$\checkmark^\checkmark\\checkmarkc\checkmarki\checkmarkr\checkmarkc\checkmark$\checkmark1\checkmark$\checkmark^\checkmark\\checkmarkc\checkmarki\checkmarkr\checkmarkc\checkmark$\checkmark.\checkmark$\checkmark^\checkmark\\checkmarkc\checkmarki\checkmarkr\checkmarkc\checkmark$\checkmark2\checkmark$\checkmark^\checkmark\\checkmarkc\checkmarki\checkmarkr\checkmarkc\checkmark$\checkmark]\checkmark$\checkmark^\checkmark\\checkmarkc\checkmarki\checkmarkr\checkmarkc\checkmark$\checkmark \checkmark$\checkmark^\checkmark\\checkmarkc\checkmarki\checkmarkr\checkmarkc\checkmark$\checkmarkn\checkmark$\checkmark^\checkmark\\checkmarkc\checkmarki\checkmarkr\checkmarkc\checkmark$\checkmarko\checkmark$\checkmark^\checkmark\\checkmarkc\checkmarki\checkmarkr\checkmarkc\checkmark$\checkmarkd\checkmark$\checkmark^\checkmark\\checkmarkc\checkmarki\checkmarkr\checkmarkc\checkmark$\checkmarke\checkmark$\checkmark^\checkmark\\checkmarkc\checkmarki\checkmarkr\checkmarkc\checkmark$\checkmark[\checkmark$\checkmark^\checkmark\\checkmarkc\checkmarki\checkmarkr\checkmarkc\checkmark$\checkmarkm\checkmark$\checkmark^\checkmark\\checkmarkc\checkmarki\checkmarkr\checkmarkc\checkmark$\checkmarki\checkmark$\checkmark^\checkmark\\checkmarkc\checkmarki\checkmarkr\checkmarkc\checkmark$\checkmarkd\checkmark$\checkmark^\checkmark\\checkmarkc\checkmarki\checkmarkr\checkmarkc\checkmark$\checkmarkw\checkmark$\checkmark^\checkmark\\checkmarkc\checkmarki\checkmarkr\checkmarkc\checkmark$\checkmarka\checkmark$\checkmark^\checkmark\\checkmarkc\checkmarki\checkmarkr\checkmarkc\checkmark$\checkmarky\checkmark$\checkmark^\checkmark\\checkmarkc\checkmarki\checkmarkr\checkmarkc\checkmark$\checkmark,\checkmark$\checkmark^\checkmark\\checkmarkc\checkmarki\checkmarkr\checkmarkc\checkmark$\checkmarka\checkmark$\checkmark^\checkmark\\checkmarkc\checkmarki\checkmarkr\checkmarkc\checkmark$\checkmarkb\checkmark$\checkmark^\checkmark\\checkmarkc\checkmarki\checkmarkr\checkmarkc\checkmark$\checkmarko\checkmark$\checkmark^\checkmark\\checkmarkc\checkmarki\checkmarkr\checkmarkc\checkmark$\checkmarkv\checkmark$\checkmark^\checkmark\\checkmarkc\checkmarki\checkmarkr\checkmarkc\checkmark$\checkmarke\checkmark$\checkmark^\checkmark\\checkmarkc\checkmarki\checkmarkr\checkmarkc\checkmark$\checkmark]\checkmark$\checkmark^\checkmark\\checkmarkc\checkmarki\checkmarkr\checkmarkc\checkmark$\checkmark \checkmark$\checkmark^\checkmark\\checkmarkc\checkmarki\checkmarkr\checkmarkc\checkmark$\checkmark{\checkmark$\checkmark^\checkmark\\checkmarkc\checkmarki\checkmarkr\checkmarkc\checkmark$\checkmark$\checkmark$\checkmark^\checkmark\\checkmarkc\checkmarki\checkmarkr\checkmarkc\checkmark$\checkmarkM\checkmark$\checkmark^\checkmark\\checkmarkc\checkmarki\checkmarkr\checkmarkc\checkmark$\checkmark_\checkmark$\checkmark^\checkmark\\checkmarkc\checkmarki\checkmarkr\checkmarkc\checkmark$\checkmark{\checkmark$\checkmark^\checkmark\\checkmarkc\checkmarki\checkmarkr\checkmarkc\checkmark$\checkmarkr\checkmark$\checkmark^\checkmark\\checkmarkc\checkmarki\checkmarkr\checkmarkc\checkmark$\checkmarke\checkmark$\checkmark^\checkmark\\checkmarkc\checkmarki\checkmarkr\checkmarkc\checkmark$\checkmarkn\checkmark$\checkmark^\checkmark\\checkmarkc\checkmarki\checkmarkr\checkmarkc\checkmark$\checkmarkv\checkmark$\checkmark^\checkmark\\checkmarkc\checkmarki\checkmarkr\checkmarkc\checkmark$\checkmark}\checkmark$\checkmark^\checkmark\\checkmarkc\checkmarki\checkmarkr\checkmarkc\checkmark$\checkmark$\checkmark$\checkmark^\checkmark\\checkmarkc\checkmarki\checkmarkr\checkmarkc\checkmark$\checkmark}\checkmark$\checkmark^\checkmark\\checkmarkc\checkmarki\checkmarkr\checkmarkc\checkmark$\checkmark;\checkmark$\checkmark^\checkmark\\checkmarkc\checkmarki\checkmarkr\checkmarkc\checkmark$\checkmark
\checkmark$\checkmark^\checkmark\\checkmarkc\checkmarki\checkmarkr\checkmarkc\checkmark$\checkmark \checkmark$\checkmark^\checkmark\\checkmarkc\checkmarki\checkmarkr\checkmarkc\checkmark$\checkmark \checkmark$\checkmark^\checkmark\\checkmarkc\checkmarki\checkmarkr\checkmarkc\checkmark$\checkmark \checkmark$\checkmark^\checkmark\\checkmarkc\checkmarki\checkmarkr\checkmarkc\checkmark$\checkmark \checkmark$\checkmark^\checkmark\\checkmarkc\checkmarki\checkmarkr\checkmarkc\checkmark$\checkmark \checkmark$\checkmark^\checkmark\\checkmarkc\checkmarki\checkmarkr\checkmarkc\checkmark$\checkmark \checkmark$\checkmark^\checkmark\\checkmarkc\checkmarki\checkmarkr\checkmarkc\checkmark$\checkmark \checkmark$\checkmark^\checkmark\\checkmarkc\checkmarki\checkmarkr\checkmarkc\checkmark$\checkmark \checkmark$\checkmark^\checkmark\\checkmarkc\checkmarki\checkmarkr\checkmarkc\checkmark$\checkmark\\checkmark$\checkmark^\checkmark\\checkmarkc\checkmarki\checkmarkr\checkmarkc\checkmark$\checkmarkd\checkmark$\checkmark^\checkmark\\checkmarkc\checkmarki\checkmarkr\checkmarkc\checkmark$\checkmarkr\checkmark$\checkmark^\checkmark\\checkmarkc\checkmarki\checkmarkr\checkmarkc\checkmark$\checkmarka\checkmark$\checkmark^\checkmark\\checkmarkc\checkmarki\checkmarkr\checkmarkc\checkmark$\checkmarkw\checkmark$\checkmark^\checkmark\\checkmarkc\checkmarki\checkmarkr\checkmarkc\checkmark$\checkmark[\checkmark$\checkmark^\checkmark\\checkmarkc\checkmarki\checkmarkr\checkmarkc\checkmark$\checkmark|\checkmark$\checkmark^\checkmark\\checkmarkc\checkmarki\checkmarkr\checkmarkc\checkmark$\checkmark<\checkmark$\checkmark^\checkmark\\checkmarkc\checkmarki\checkmarkr\checkmarkc\checkmark$\checkmark-\checkmark$\checkmark^\checkmark\\checkmarkc\checkmarki\checkmarkr\checkmarkc\checkmark$\checkmark>\checkmark$\checkmark^\checkmark\\checkmarkc\checkmarki\checkmarkr\checkmarkc\checkmark$\checkmark|\checkmark$\checkmark^\checkmark\\checkmarkc\checkmarki\checkmarkr\checkmarkc\checkmark$\checkmark]\checkmark$\checkmark^\checkmark\\checkmarkc\checkmarki\checkmarkr\checkmarkc\checkmark$\checkmark \checkmark$\checkmark^\checkmark\\checkmarkc\checkmarki\checkmarkr\checkmarkc\checkmark$\checkmark(\checkmark$\checkmark^\checkmark\\checkmarkc\checkmarki\checkmarkr\checkmarkc\checkmark$\checkmark0\checkmark$\checkmark^\checkmark\\checkmarkc\checkmarki\checkmarkr\checkmarkc\checkmark$\checkmark.\checkmark$\checkmark^\checkmark\\checkmarkc\checkmarki\checkmarkr\checkmarkc\checkmark$\checkmark7\checkmark$\checkmark^\checkmark\\checkmarkc\checkmarki\checkmarkr\checkmarkc\checkmark$\checkmark,\checkmark$\checkmark^\checkmark\\checkmarkc\checkmarki\checkmarkr\checkmarkc\checkmark$\checkmark0\checkmark$\checkmark^\checkmark\\checkmarkc\checkmarki\checkmarkr\checkmarkc\checkmark$\checkmark)\checkmark$\checkmark^\checkmark\\checkmarkc\checkmarki\checkmarkr\checkmarkc\checkmark$\checkmark \checkmark$\checkmark^\checkmark\\checkmarkc\checkmarki\checkmarkr\checkmarkc\checkmark$\checkmark-\checkmark$\checkmark^\checkmark\\checkmarkc\checkmarki\checkmarkr\checkmarkc\checkmark$\checkmark-\checkmark$\checkmark^\checkmark\\checkmarkc\checkmarki\checkmarkr\checkmarkc\checkmark$\checkmark \checkmark$\checkmark^\checkmark\\checkmarkc\checkmarki\checkmarkr\checkmarkc\checkmark$\checkmark(\checkmark$\checkmark^\checkmark\\checkmarkc\checkmarki\checkmarkr\checkmarkc\checkmark$\checkmark0\checkmark$\checkmark^\checkmark\\checkmarkc\checkmarki\checkmarkr\checkmarkc\checkmark$\checkmark.\checkmark$\checkmark^\checkmark\\checkmarkc\checkmarki\checkmarkr\checkmarkc\checkmark$\checkmark7\checkmark$\checkmark^\checkmark\\checkmarkc\checkmarki\checkmarkr\checkmarkc\checkmark$\checkmark,\checkmark$\checkmark^\checkmark\\checkmarkc\checkmarki\checkmarkr\checkmarkc\checkmark$\checkmark3\checkmark$\checkmark^\checkmark\\checkmarkc\checkmarki\checkmarkr\checkmarkc\checkmark$\checkmark)\checkmark$\checkmark^\checkmark\\checkmarkc\checkmarki\checkmarkr\checkmarkc\checkmark$\checkmark \checkmark$\checkmark^\checkmark\\checkmarkc\checkmarki\checkmarkr\checkmarkc\checkmark$\checkmarkn\checkmark$\checkmark^\checkmark\\checkmarkc\checkmarki\checkmarkr\checkmarkc\checkmark$\checkmarko\checkmark$\checkmark^\checkmark\\checkmarkc\checkmarki\checkmarkr\checkmarkc\checkmark$\checkmarkd\checkmark$\checkmark^\checkmark\\checkmarkc\checkmarki\checkmarkr\checkmarkc\checkmark$\checkmarke\checkmark$\checkmark^\checkmark\\checkmarkc\checkmarki\checkmarkr\checkmarkc\checkmark$\checkmark[\checkmark$\checkmark^\checkmark\\checkmarkc\checkmarki\checkmarkr\checkmarkc\checkmark$\checkmarkm\checkmark$\checkmark^\checkmark\\checkmarkc\checkmarki\checkmarkr\checkmarkc\checkmark$\checkmarki\checkmark$\checkmark^\checkmark\\checkmarkc\checkmarki\checkmarkr\checkmarkc\checkmark$\checkmarkd\checkmark$\checkmark^\checkmark\\checkmarkc\checkmarki\checkmarkr\checkmarkc\checkmark$\checkmarkw\checkmark$\checkmark^\checkmark\\checkmarkc\checkmarki\checkmarkr\checkmarkc\checkmark$\checkmarka\checkmark$\checkmark^\checkmark\\checkmarkc\checkmarki\checkmarkr\checkmarkc\checkmark$\checkmarky\checkmark$\checkmark^\checkmark\\checkmarkc\checkmarki\checkmarkr\checkmarkc\checkmark$\checkmark,\checkmark$\checkmark^\checkmark\\checkmarkc\checkmarki\checkmarkr\checkmarkc\checkmark$\checkmarkr\checkmark$\checkmark^\checkmark\\checkmarkc\checkmarki\checkmarkr\checkmarkc\checkmark$\checkmarki\checkmark$\checkmark^\checkmark\\checkmarkc\checkmarki\checkmarkr\checkmarkc\checkmark$\checkmarkg\checkmark$\checkmark^\checkmark\\checkmarkc\checkmarki\checkmarkr\checkmarkc\checkmark$\checkmarkh\checkmark$\checkmark^\checkmark\\checkmarkc\checkmarki\checkmarkr\checkmarkc\checkmark$\checkmarkt\checkmark$\checkmark^\checkmark\\checkmarkc\checkmarki\checkmarkr\checkmarkc\checkmark$\checkmark]\checkmark$\checkmark^\checkmark\\checkmarkc\checkmarki\checkmarkr\checkmarkc\checkmark$\checkmark \checkmark$\checkmark^\checkmark\\checkmarkc\checkmarki\checkmarkr\checkmarkc\checkmark$\checkmark{\checkmark$\checkmark^\checkmark\\checkmarkc\checkmarki\checkmarkr\checkmarkc\checkmark$\checkmark$\checkmark$\checkmark^\checkmark\\checkmarkc\checkmarki\checkmarkr\checkmarkc\checkmark$\checkmarkH\checkmark$\checkmark^\checkmark\\checkmarkc\checkmarki\checkmarkr\checkmarkc\checkmark$\checkmark=\checkmark$\checkmark^\checkmark\\checkmarkc\checkmarki\checkmarkr\checkmarkc\checkmark$\checkmark2\checkmark$\checkmark^\checkmark\\checkmarkc\checkmarki\checkmarkr\checkmarkc\checkmark$\checkmark0\checkmark$\checkmark^\checkmark\\checkmarkc\checkmarki\checkmarkr\checkmarkc\checkmark$\checkmark0\checkmark$\checkmark^\checkmark\\checkmarkc\checkmarki\checkmarkr\checkmarkc\checkmark$\checkmark$\checkmark$\checkmark^\checkmark\\checkmarkc\checkmarki\checkmarkr\checkmarkc\checkmark$\checkmark \checkmark$\checkmark^\checkmark\\checkmarkc\checkmarki\checkmarkr\checkmarkc\checkmark$\checkmarkm\checkmark$\checkmark^\checkmark\\checkmarkc\checkmarki\checkmarkr\checkmarkc\checkmark$\checkmarkm\checkmark$\checkmark^\checkmark\\checkmarkc\checkmarki\checkmarkr\checkmarkc\checkmark$\checkmark}\checkmark$\checkmark^\checkmark\\checkmarkc\checkmarki\checkmarkr\checkmarkc\checkmark$\checkmark;\checkmark$\checkmark^\checkmark\\checkmarkc\checkmarki\checkmarkr\checkmarkc\checkmark$\checkmark
\checkmark$\checkmark^\checkmark\\checkmarkc\checkmarki\checkmarkr\checkmarkc\checkmark$\checkmark \checkmark$\checkmark^\checkmark\\checkmarkc\checkmarki\checkmarkr\checkmarkc\checkmark$\checkmark \checkmark$\checkmark^\checkmark\\checkmarkc\checkmarki\checkmarkr\checkmarkc\checkmark$\checkmark \checkmark$\checkmark^\checkmark\\checkmarkc\checkmarki\checkmarkr\checkmarkc\checkmark$\checkmark \checkmark$\checkmark^\checkmark\\checkmarkc\checkmarki\checkmarkr\checkmarkc\checkmark$\checkmark\\checkmark$\checkmark^\checkmark\\checkmarkc\checkmarki\checkmarkr\checkmarkc\checkmark$\checkmarke\checkmark$\checkmark^\checkmark\\checkmarkc\checkmarki\checkmarkr\checkmarkc\checkmark$\checkmarkn\checkmark$\checkmark^\checkmark\\checkmarkc\checkmarki\checkmarkr\checkmarkc\checkmark$\checkmarkd\checkmark$\checkmark^\checkmark\\checkmarkc\checkmarki\checkmarkr\checkmarkc\checkmark$\checkmark{\checkmark$\checkmark^\checkmark\\checkmarkc\checkmarki\checkmarkr\checkmarkc\checkmark$\checkmarkt\checkmark$\checkmark^\checkmark\\checkmarkc\checkmarki\checkmarkr\checkmarkc\checkmark$\checkmarki\checkmark$\checkmark^\checkmark\\checkmarkc\checkmarki\checkmarkr\checkmarkc\checkmark$\checkmarkk\checkmark$\checkmark^\checkmark\\checkmarkc\checkmarki\checkmarkr\checkmarkc\checkmark$\checkmarkz\checkmark$\checkmark^\checkmark\\checkmarkc\checkmarki\checkmarkr\checkmarkc\checkmark$\checkmarkp\checkmark$\checkmark^\checkmark\\checkmarkc\checkmarki\checkmarkr\checkmarkc\checkmark$\checkmarki\checkmark$\checkmark^\checkmark\\checkmarkc\checkmarki\checkmarkr\checkmarkc\checkmark$\checkmarkc\checkmark$\checkmark^\checkmark\\checkmarkc\checkmarki\checkmarkr\checkmarkc\checkmark$\checkmarkt\checkmark$\checkmark^\checkmark\\checkmarkc\checkmarki\checkmarkr\checkmarkc\checkmark$\checkmarku\checkmark$\checkmark^\checkmark\\checkmarkc\checkmarki\checkmarkr\checkmarkc\checkmark$\checkmarkr\checkmark$\checkmark^\checkmark\\checkmarkc\checkmarki\checkmarkr\checkmarkc\checkmark$\checkmarke\checkmark$\checkmark^\checkmark\\checkmarkc\checkmarki\checkmarkr\checkmarkc\checkmark$\checkmark}\checkmark$\checkmark^\checkmark\\checkmarkc\checkmarki\checkmarkr\checkmarkc\checkmark$\checkmark
\checkmark$\checkmark^\checkmark\\checkmarkc\checkmarki\checkmarkr\checkmarkc\checkmark$\checkmark \checkmark$\checkmark^\checkmark\\checkmarkc\checkmarki\checkmarkr\checkmarkc\checkmark$\checkmark \checkmark$\checkmark^\checkmark\\checkmarkc\checkmarki\checkmarkr\checkmarkc\checkmark$\checkmark \checkmark$\checkmark^\checkmark\\checkmarkc\checkmarki\checkmarkr\checkmarkc\checkmark$\checkmark \checkmark$\checkmark^\checkmark\\checkmarkc\checkmarki\checkmarkr\checkmarkc\checkmark$\checkmark\\checkmark$\checkmark^\checkmark\\checkmarkc\checkmarki\checkmarkr\checkmarkc\checkmark$\checkmarkc\checkmark$\checkmark^\checkmark\\checkmarkc\checkmarki\checkmarkr\checkmarkc\checkmark$\checkmarka\checkmark$\checkmark^\checkmark\\checkmarkc\checkmarki\checkmarkr\checkmarkc\checkmark$\checkmarkp\checkmark$\checkmark^\checkmark\\checkmarkc\checkmarki\checkmarkr\checkmarkc\checkmark$\checkmarkt\checkmark$\checkmark^\checkmark\\checkmarkc\checkmarki\checkmarkr\checkmarkc\checkmark$\checkmarki\checkmark$\checkmark^\checkmark\\checkmarkc\checkmarki\checkmarkr\checkmarkc\checkmark$\checkmarko\checkmark$\checkmark^\checkmark\\checkmarkc\checkmarki\checkmarkr\checkmarkc\checkmark$\checkmarkn\checkmark$\checkmark^\checkmark\\checkmarkc\checkmarki\checkmarkr\checkmarkc\checkmark$\checkmark{\checkmark$\checkmark^\checkmark\\checkmarkc\checkmarki\checkmarkr\checkmarkc\checkmark$\checkmarkS\checkmark$\checkmark^\checkmark\\checkmarkc\checkmarki\checkmarkr\checkmarkc\checkmark$\checkmarkc\checkmark$\checkmark^\checkmark\\checkmarkc\checkmarki\checkmarkr\checkmarkc\checkmark$\checkmarkh\checkmark$\checkmark^\checkmark\\checkmarkc\checkmarki\checkmarkr\checkmarkc\checkmark$\checkmarké\checkmark$\checkmark^\checkmark\\checkmarkc\checkmarki\checkmarkr\checkmarkc\checkmark$\checkmarkm\checkmark$\checkmark^\checkmark\\checkmarkc\checkmarki\checkmarkr\checkmarkc\checkmark$\checkmarka\checkmark$\checkmark^\checkmark\\checkmarkc\checkmarki\checkmarkr\checkmarkc\checkmark$\checkmark \checkmark$\checkmark^\checkmark\\checkmarkc\checkmarki\checkmarkr\checkmarkc\checkmark$\checkmarkd\checkmark$\checkmark^\checkmark\\checkmarkc\checkmarki\checkmarkr\checkmarkc\checkmark$\checkmarku\checkmark$\checkmark^\checkmark\\checkmarkc\checkmarki\checkmarkr\checkmarkc\checkmark$\checkmark \checkmark$\checkmark^\checkmark\\checkmarkc\checkmarki\checkmarkr\checkmarkc\checkmark$\checkmarkm\checkmark$\checkmark^\checkmark\\checkmarkc\checkmarki\checkmarkr\checkmarkc\checkmark$\checkmarko\checkmark$\checkmark^\checkmark\\checkmarkc\checkmarki\checkmarkr\checkmarkc\checkmark$\checkmarkm\checkmark$\checkmark^\checkmark\\checkmarkc\checkmarki\checkmarkr\checkmarkc\checkmark$\checkmarke\checkmark$\checkmark^\checkmark\\checkmarkc\checkmarki\checkmarkr\checkmarkc\checkmark$\checkmarkn\checkmark$\checkmark^\checkmark\\checkmarkc\checkmarki\checkmarkr\checkmarkc\checkmark$\checkmarkt\checkmark$\checkmark^\checkmark\\checkmarkc\checkmarki\checkmarkr\checkmarkc\checkmark$\checkmark \checkmark$\checkmark^\checkmark\\checkmarkc\checkmarki\checkmarkr\checkmarkc\checkmark$\checkmarkd\checkmark$\checkmark^\checkmark\\checkmarkc\checkmarki\checkmarkr\checkmarkc\checkmark$\checkmarke\checkmark$\checkmark^\checkmark\\checkmarkc\checkmarki\checkmarkr\checkmarkc\checkmark$\checkmark \checkmark$\checkmark^\checkmark\\checkmarkc\checkmarki\checkmarkr\checkmarkc\checkmark$\checkmarkr\checkmark$\checkmark^\checkmark\\checkmarkc\checkmarki\checkmarkr\checkmarkc\checkmark$\checkmarke\checkmark$\checkmark^\checkmark\\checkmarkc\checkmarki\checkmarkr\checkmarkc\checkmark$\checkmarkn\checkmark$\checkmark^\checkmark\\checkmarkc\checkmarki\checkmarkr\checkmarkc\checkmark$\checkmarkv\checkmark$\checkmark^\checkmark\\checkmarkc\checkmarki\checkmarkr\checkmarkc\checkmark$\checkmarke\checkmark$\checkmark^\checkmark\\checkmarkc\checkmarki\checkmarkr\checkmarkc\checkmark$\checkmarkr\checkmark$\checkmark^\checkmark\\checkmarkc\checkmarki\checkmarkr\checkmarkc\checkmark$\checkmarks\checkmark$\checkmark^\checkmark\\checkmarkc\checkmarki\checkmarkr\checkmarkc\checkmark$\checkmarke\checkmark$\checkmark^\checkmark\\checkmarkc\checkmarki\checkmarkr\checkmarkc\checkmark$\checkmarkm\checkmark$\checkmark^\checkmark\\checkmarkc\checkmarki\checkmarkr\checkmarkc\checkmark$\checkmarke\checkmark$\checkmark^\checkmark\\checkmarkc\checkmarki\checkmarkr\checkmarkc\checkmark$\checkmarkn\checkmark$\checkmark^\checkmark\\checkmarkc\checkmarki\checkmarkr\checkmarkc\checkmark$\checkmarkt\checkmark$\checkmark^\checkmark\\checkmarkc\checkmarki\checkmarkr\checkmarkc\checkmark$\checkmark \checkmark$\checkmark^\checkmark\\checkmarkc\checkmarki\checkmarkr\checkmarkc\checkmark$\checkmarks\checkmark$\checkmark^\checkmark\\checkmarkc\checkmarki\checkmarkr\checkmarkc\checkmark$\checkmarku\checkmark$\checkmark^\checkmark\\checkmarkc\checkmarki\checkmarkr\checkmarkc\checkmark$\checkmarkr\checkmark$\checkmark^\checkmark\\checkmarkc\checkmarki\checkmarkr\checkmarkc\checkmark$\checkmark \checkmark$\checkmark^\checkmark\\checkmarkc\checkmarki\checkmarkr\checkmarkc\checkmark$\checkmarkl\checkmark$\checkmark^\checkmark\\checkmarkc\checkmarki\checkmarkr\checkmarkc\checkmark$\checkmark'\checkmark$\checkmark^\checkmark\\checkmarkc\checkmarki\checkmarkr\checkmarkc\checkmark$\checkmarka\checkmark$\checkmark^\checkmark\\checkmarkc\checkmarki\checkmarkr\checkmarkc\checkmark$\checkmarkr\checkmark$\checkmark^\checkmark\\checkmarkc\checkmarki\checkmarkr\checkmarkc\checkmark$\checkmarkb\checkmark$\checkmark^\checkmark\\checkmarkc\checkmarki\checkmarkr\checkmarkc\checkmark$\checkmarkr\checkmark$\checkmark^\checkmark\\checkmarkc\checkmarki\checkmarkr\checkmarkc\checkmark$\checkmarke\checkmark$\checkmark^\checkmark\\checkmarkc\checkmarki\checkmarkr\checkmarkc\checkmark$\checkmark \checkmark$\checkmark^\checkmark\\checkmarkc\checkmarki\checkmarkr\checkmarkc\checkmark$\checkmarkd\checkmark$\checkmark^\checkmark\\checkmarkc\checkmarki\checkmarkr\checkmarkc\checkmark$\checkmarke\checkmark$\checkmark^\checkmark\\checkmarkc\checkmarki\checkmarkr\checkmarkc\checkmark$\checkmark \checkmark$\checkmark^\checkmark\\checkmarkc\checkmarki\checkmarkr\checkmarkc\checkmark$\checkmarkp\checkmark$\checkmark^\checkmark\\checkmarkc\checkmarki\checkmarkr\checkmarkc\checkmark$\checkmarki\checkmark$\checkmark^\checkmark\\checkmarkc\checkmarki\checkmarkr\checkmarkc\checkmark$\checkmarkv\checkmark$\checkmark^\checkmark\\checkmarkc\checkmarki\checkmarkr\checkmarkc\checkmark$\checkmarko\checkmark$\checkmark^\checkmark\\checkmarkc\checkmarki\checkmarkr\checkmarkc\checkmark$\checkmarkt\checkmark$\checkmark^\checkmark\\checkmarkc\checkmarki\checkmarkr\checkmarkc\checkmark$\checkmarke\checkmark$\checkmark^\checkmark\\checkmarkc\checkmarki\checkmarkr\checkmarkc\checkmark$\checkmarkm\checkmark$\checkmark^\checkmark\\checkmarkc\checkmarki\checkmarkr\checkmarkc\checkmark$\checkmarke\checkmark$\checkmark^\checkmark\\checkmarkc\checkmarki\checkmarkr\checkmarkc\checkmark$\checkmarkn\checkmark$\checkmark^\checkmark\\checkmarkc\checkmarki\checkmarkr\checkmarkc\checkmark$\checkmarkt\checkmark$\checkmark^\checkmark\\checkmarkc\checkmarki\checkmarkr\checkmarkc\checkmark$\checkmark \checkmark$\checkmark^\checkmark\\checkmarkc\checkmarki\checkmarkr\checkmarkc\checkmark$\checkmark(\checkmark$\checkmark^\checkmark\\checkmarkc\checkmarki\checkmarkr\checkmarkc\checkmark$\checkmarkA\checkmark$\checkmark^\checkmark\\checkmarkc\checkmarki\checkmarkr\checkmarkc\checkmark$\checkmarkx\checkmark$\checkmark^\checkmark\\checkmarkc\checkmarki\checkmarkr\checkmarkc\checkmark$\checkmarke\checkmark$\checkmark^\checkmark\\checkmarkc\checkmarki\checkmarkr\checkmarkc\checkmark$\checkmark \checkmark$\checkmark^\checkmark\\checkmarkc\checkmarki\checkmarkr\checkmarkc\checkmark$\checkmarkA\checkmark$\checkmark^\checkmark\\checkmarkc\checkmarki\checkmarkr\checkmarkc\checkmark$\checkmark)\checkmark$\checkmark^\checkmark\\checkmarkc\checkmarki\checkmarkr\checkmarkc\checkmark$\checkmark}\checkmark$\checkmark^\checkmark\\checkmarkc\checkmarki\checkmarkr\checkmarkc\checkmark$\checkmark
\checkmark$\checkmark^\checkmark\\checkmarkc\checkmarki\checkmarkr\checkmarkc\checkmark$\checkmark \checkmark$\checkmark^\checkmark\\checkmarkc\checkmarki\checkmarkr\checkmarkc\checkmark$\checkmark \checkmark$\checkmark^\checkmark\\checkmarkc\checkmarki\checkmarkr\checkmarkc\checkmark$\checkmark \checkmark$\checkmark^\checkmark\\checkmarkc\checkmarki\checkmarkr\checkmarkc\checkmark$\checkmark \checkmark$\checkmark^\checkmark\\checkmarkc\checkmarki\checkmarkr\checkmarkc\checkmark$\checkmark\\checkmark$\checkmark^\checkmark\\checkmarkc\checkmarki\checkmarkr\checkmarkc\checkmark$\checkmarkl\checkmark$\checkmark^\checkmark\\checkmarkc\checkmarki\checkmarkr\checkmarkc\checkmark$\checkmarka\checkmark$\checkmark^\checkmark\\checkmarkc\checkmarki\checkmarkr\checkmarkc\checkmark$\checkmarkb\checkmark$\checkmark^\checkmark\\checkmarkc\checkmarki\checkmarkr\checkmarkc\checkmark$\checkmarke\checkmark$\checkmark^\checkmark\\checkmarkc\checkmarki\checkmarkr\checkmarkc\checkmark$\checkmarkl\checkmark$\checkmark^\checkmark\\checkmarkc\checkmarki\checkmarkr\checkmarkc\checkmark$\checkmark{\checkmark$\checkmark^\checkmark\\checkmarkc\checkmarki\checkmarkr\checkmarkc\checkmark$\checkmarkf\checkmark$\checkmark^\checkmark\\checkmarkc\checkmarki\checkmarkr\checkmarkc\checkmark$\checkmarki\checkmark$\checkmark^\checkmark\\checkmarkc\checkmarki\checkmarkr\checkmarkc\checkmark$\checkmarkg\checkmark$\checkmark^\checkmark\\checkmarkc\checkmarki\checkmarkr\checkmarkc\checkmark$\checkmark:\checkmark$\checkmark^\checkmark\\checkmarkc\checkmarki\checkmarkr\checkmarkc\checkmark$\checkmarkm\checkmark$\checkmark^\checkmark\\checkmarkc\checkmarki\checkmarkr\checkmarkc\checkmark$\checkmarko\checkmark$\checkmark^\checkmark\\checkmarkc\checkmarki\checkmarkr\checkmarkc\checkmark$\checkmarkm\checkmark$\checkmark^\checkmark\\checkmarkc\checkmarki\checkmarkr\checkmarkc\checkmark$\checkmarke\checkmark$\checkmark^\checkmark\\checkmarkc\checkmarki\checkmarkr\checkmarkc\checkmark$\checkmarkn\checkmark$\checkmark^\checkmark\\checkmarkc\checkmarki\checkmarkr\checkmarkc\checkmark$\checkmarkt\checkmark$\checkmark^\checkmark\\checkmarkc\checkmarki\checkmarkr\checkmarkc\checkmark$\checkmark_\checkmark$\checkmark^\checkmark\\checkmarkc\checkmarki\checkmarkr\checkmarkc\checkmark$\checkmarkr\checkmark$\checkmark^\checkmark\\checkmarkc\checkmarki\checkmarkr\checkmarkc\checkmark$\checkmarke\checkmark$\checkmark^\checkmark\\checkmarkc\checkmarki\checkmarkr\checkmarkc\checkmark$\checkmarkn\checkmark$\checkmark^\checkmark\\checkmarkc\checkmarki\checkmarkr\checkmarkc\checkmark$\checkmarkv\checkmark$\checkmark^\checkmark\\checkmarkc\checkmarki\checkmarkr\checkmarkc\checkmark$\checkmarke\checkmark$\checkmark^\checkmark\\checkmarkc\checkmarki\checkmarkr\checkmarkc\checkmark$\checkmarkr\checkmark$\checkmark^\checkmark\\checkmarkc\checkmarki\checkmarkr\checkmarkc\checkmark$\checkmarks\checkmark$\checkmark^\checkmark\\checkmarkc\checkmarki\checkmarkr\checkmarkc\checkmark$\checkmarke\checkmark$\checkmark^\checkmark\\checkmarkc\checkmarki\checkmarkr\checkmarkc\checkmark$\checkmarkm\checkmark$\checkmark^\checkmark\\checkmarkc\checkmarki\checkmarkr\checkmarkc\checkmark$\checkmarke\checkmark$\checkmark^\checkmark\\checkmarkc\checkmarki\checkmarkr\checkmarkc\checkmark$\checkmarkn\checkmark$\checkmark^\checkmark\\checkmarkc\checkmarki\checkmarkr\checkmarkc\checkmark$\checkmarkt\checkmark$\checkmark^\checkmark\\checkmarkc\checkmarki\checkmarkr\checkmarkc\checkmark$\checkmark}\checkmark$\checkmark^\checkmark\\checkmarkc\checkmarki\checkmarkr\checkmarkc\checkmark$\checkmark
\checkmark$\checkmark^\checkmark\\checkmarkc\checkmarki\checkmarkr\checkmarkc\checkmark$\checkmark\\checkmark$\checkmark^\checkmark\\checkmarkc\checkmarki\checkmarkr\checkmarkc\checkmark$\checkmarke\checkmark$\checkmark^\checkmark\\checkmarkc\checkmarki\checkmarkr\checkmarkc\checkmark$\checkmarkn\checkmark$\checkmark^\checkmark\\checkmarkc\checkmarki\checkmarkr\checkmarkc\checkmark$\checkmarkd\checkmark$\checkmark^\checkmark\\checkmarkc\checkmarki\checkmarkr\checkmarkc\checkmark$\checkmark{\checkmark$\checkmark^\checkmark\\checkmarkc\checkmarki\checkmarkr\checkmarkc\checkmark$\checkmarkf\checkmark$\checkmark^\checkmark\\checkmarkc\checkmarki\checkmarkr\checkmarkc\checkmark$\checkmarki\checkmark$\checkmark^\checkmark\\checkmarkc\checkmarki\checkmarkr\checkmarkc\checkmark$\checkmarkg\checkmark$\checkmark^\checkmark\\checkmarkc\checkmarki\checkmarkr\checkmarkc\checkmark$\checkmarku\checkmark$\checkmark^\checkmark\\checkmarkc\checkmarki\checkmarkr\checkmarkc\checkmark$\checkmarkr\checkmark$\checkmark^\checkmark\\checkmarkc\checkmarki\checkmarkr\checkmarkc\checkmark$\checkmarke\checkmark$\checkmark^\checkmark\\checkmarkc\checkmarki\checkmarkr\checkmarkc\checkmark$\checkmark}\checkmark$\checkmark^\checkmark\\checkmarkc\checkmarki\checkmarkr\checkmarkc\checkmark$\checkmark
\checkmark$\checkmark^\checkmark\\checkmarkc\checkmarki\checkmarkr\checkmarkc\checkmark$\checkmark
\checkmark$\checkmark^\checkmark\\checkmarkc\checkmarki\checkmarkr\checkmarkc\checkmark$\checkmark\\checkmark$\checkmark^\checkmark\\checkmarkc\checkmarki\checkmarkr\checkmarkc\checkmark$\checkmarks\checkmark$\checkmark^\checkmark\\checkmarkc\checkmarki\checkmarkr\checkmarkc\checkmark$\checkmarku\checkmark$\checkmark^\checkmark\\checkmarkc\checkmarki\checkmarkr\checkmarkc\checkmark$\checkmarkb\checkmark$\checkmark^\checkmark\\checkmarkc\checkmarki\checkmarkr\checkmarkc\checkmark$\checkmarks\checkmark$\checkmark^\checkmark\\checkmarkc\checkmarki\checkmarkr\checkmarkc\checkmark$\checkmarke\checkmark$\checkmark^\checkmark\\checkmarkc\checkmarki\checkmarkr\checkmarkc\checkmark$\checkmarkc\checkmark$\checkmark^\checkmark\\checkmarkc\checkmarki\checkmarkr\checkmarkc\checkmark$\checkmarkt\checkmark$\checkmark^\checkmark\\checkmarkc\checkmarki\checkmarkr\checkmarkc\checkmark$\checkmarki\checkmark$\checkmark^\checkmark\\checkmarkc\checkmarki\checkmarkr\checkmarkc\checkmark$\checkmarko\checkmark$\checkmark^\checkmark\\checkmarkc\checkmarki\checkmarkr\checkmarkc\checkmark$\checkmarkn\checkmark$\checkmark^\checkmark\\checkmarkc\checkmarki\checkmarkr\checkmarkc\checkmark$\checkmark{\checkmark$\checkmark^\checkmark\\checkmarkc\checkmarki\checkmarkr\checkmarkc\checkmark$\checkmarkM\checkmark$\checkmark^\checkmark\\checkmarkc\checkmarki\checkmarkr\checkmarkc\checkmark$\checkmarko\checkmark$\checkmark^\checkmark\\checkmarkc\checkmarki\checkmarkr\checkmarkc\checkmark$\checkmarkd\checkmark$\checkmark^\checkmark\\checkmarkc\checkmarki\checkmarkr\checkmarkc\checkmark$\checkmarkè\checkmark$\checkmark^\checkmark\\checkmarkc\checkmarki\checkmarkr\checkmarkc\checkmark$\checkmarkl\checkmark$\checkmark^\checkmark\\checkmarkc\checkmarki\checkmarkr\checkmarkc\checkmark$\checkmarke\checkmark$\checkmark^\checkmark\\checkmarkc\checkmarki\checkmarkr\checkmarkc\checkmark$\checkmark \checkmark$\checkmark^\checkmark\\checkmarkc\checkmarki\checkmarkr\checkmarkc\checkmark$\checkmarks\checkmark$\checkmark^\checkmark\\checkmarkc\checkmarki\checkmarkr\checkmarkc\checkmark$\checkmarkt\checkmark$\checkmark^\checkmark\\checkmarkc\checkmarki\checkmarkr\checkmarkc\checkmark$\checkmarka\checkmark$\checkmark^\checkmark\\checkmarkc\checkmarki\checkmarkr\checkmarkc\checkmark$\checkmarkt\checkmark$\checkmark^\checkmark\\checkmarkc\checkmarki\checkmarkr\checkmarkc\checkmark$\checkmarki\checkmark$\checkmark^\checkmark\\checkmarkc\checkmarki\checkmarkr\checkmarkc\checkmark$\checkmarkq\checkmark$\checkmark^\checkmark\\checkmarkc\checkmarki\checkmarkr\checkmarkc\checkmark$\checkmarku\checkmark$\checkmark^\checkmark\\checkmarkc\checkmarki\checkmarkr\checkmarkc\checkmark$\checkmarke\checkmark$\checkmark^\checkmark\\checkmarkc\checkmarki\checkmarkr\checkmarkc\checkmark$\checkmark \checkmark$\checkmark^\checkmark\\checkmarkc\checkmarki\checkmarkr\checkmarkc\checkmark$\checkmarkd\checkmark$\checkmark^\checkmark\\checkmarkc\checkmarki\checkmarkr\checkmarkc\checkmark$\checkmarke\checkmark$\checkmark^\checkmark\\checkmarkc\checkmarki\checkmarkr\checkmarkc\checkmark$\checkmark \checkmark$\checkmark^\checkmark\\checkmarkc\checkmarki\checkmarkr\checkmarkc\checkmark$\checkmarkl\checkmark$\checkmark^\checkmark\\checkmarkc\checkmarki\checkmarkr\checkmarkc\checkmark$\checkmark'\checkmark$\checkmark^\checkmark\\checkmarkc\checkmarki\checkmarkr\checkmarkc\checkmark$\checkmarka\checkmark$\checkmark^\checkmark\\checkmarkc\checkmarki\checkmarkr\checkmarkc\checkmark$\checkmarkr\checkmark$\checkmark^\checkmark\\checkmarkc\checkmarki\checkmarkr\checkmarkc\checkmark$\checkmarkb\checkmark$\checkmark^\checkmark\\checkmarkc\checkmarki\checkmarkr\checkmarkc\checkmark$\checkmarkr\checkmark$\checkmark^\checkmark\\checkmarkc\checkmarki\checkmarkr\checkmarkc\checkmark$\checkmarke\checkmark$\checkmark^\checkmark\\checkmarkc\checkmarki\checkmarkr\checkmarkc\checkmark$\checkmark}\checkmark$\checkmark^\checkmark\\checkmarkc\checkmarki\checkmarkr\checkmarkc\checkmark$\checkmark
\checkmark$\checkmark^\checkmark\\checkmarkc\checkmarki\checkmarkr\checkmarkc\checkmark$\checkmarkL\checkmark$\checkmark^\checkmark\\checkmarkc\checkmarki\checkmarkr\checkmarkc\checkmark$\checkmark'\checkmark$\checkmark^\checkmark\\checkmarkc\checkmarki\checkmarkr\checkmarkc\checkmark$\checkmarka\checkmark$\checkmark^\checkmark\\checkmarkc\checkmarki\checkmarkr\checkmarkc\checkmark$\checkmarkr\checkmark$\checkmark^\checkmark\\checkmarkc\checkmarki\checkmarkr\checkmarkc\checkmark$\checkmarkb\checkmark$\checkmark^\checkmark\\checkmarkc\checkmarki\checkmarkr\checkmarkc\checkmark$\checkmarkr\checkmark$\checkmark^\checkmark\\checkmarkc\checkmarki\checkmarkr\checkmarkc\checkmark$\checkmarke\checkmark$\checkmark^\checkmark\\checkmarkc\checkmarki\checkmarkr\checkmarkc\checkmark$\checkmark \checkmark$\checkmark^\checkmark\\checkmarkc\checkmarki\checkmarkr\checkmarkc\checkmark$\checkmarke\checkmark$\checkmark^\checkmark\\checkmarkc\checkmarki\checkmarkr\checkmarkc\checkmark$\checkmarks\checkmark$\checkmark^\checkmark\\checkmarkc\checkmarki\checkmarkr\checkmarkc\checkmark$\checkmarkt\checkmark$\checkmark^\checkmark\\checkmarkc\checkmarki\checkmarkr\checkmarkc\checkmark$\checkmark \checkmark$\checkmark^\checkmark\\checkmarkc\checkmarki\checkmarkr\checkmarkc\checkmark$\checkmarkm\checkmark$\checkmark^\checkmark\\checkmarkc\checkmarki\checkmarkr\checkmarkc\checkmark$\checkmarko\checkmark$\checkmark^\checkmark\\checkmarkc\checkmarki\checkmarkr\checkmarkc\checkmark$\checkmarkd\checkmark$\checkmark^\checkmark\\checkmarkc\checkmarki\checkmarkr\checkmarkc\checkmark$\checkmarké\checkmark$\checkmark^\checkmark\\checkmarkc\checkmarki\checkmarkr\checkmarkc\checkmark$\checkmarkl\checkmark$\checkmark^\checkmark\\checkmarkc\checkmarki\checkmarkr\checkmarkc\checkmark$\checkmarki\checkmark$\checkmark^\checkmark\\checkmarkc\checkmarki\checkmarkr\checkmarkc\checkmark$\checkmarks\checkmark$\checkmark^\checkmark\\checkmarkc\checkmarki\checkmarkr\checkmarkc\checkmark$\checkmarké\checkmark$\checkmark^\checkmark\\checkmarkc\checkmarki\checkmarkr\checkmarkc\checkmark$\checkmark \checkmark$\checkmark^\checkmark\\checkmarkc\checkmarki\checkmarkr\checkmarkc\checkmark$\checkmarkc\checkmark$\checkmark^\checkmark\\checkmarkc\checkmarki\checkmarkr\checkmarkc\checkmark$\checkmarko\checkmark$\checkmark^\checkmark\\checkmarkc\checkmarki\checkmarkr\checkmarkc\checkmark$\checkmarkm\checkmark$\checkmark^\checkmark\\checkmarkc\checkmarki\checkmarkr\checkmarkc\checkmark$\checkmarkm\checkmark$\checkmark^\checkmark\\checkmarkc\checkmarki\checkmarkr\checkmarkc\checkmark$\checkmarke\checkmark$\checkmark^\checkmark\\checkmarkc\checkmarki\checkmarkr\checkmarkc\checkmark$\checkmark \checkmark$\checkmark^\checkmark\\checkmarkc\checkmarki\checkmarkr\checkmarkc\checkmark$\checkmarku\checkmark$\checkmark^\checkmark\\checkmarkc\checkmarki\checkmarkr\checkmarkc\checkmark$\checkmarkn\checkmark$\checkmark^\checkmark\\checkmarkc\checkmarki\checkmarkr\checkmarkc\checkmark$\checkmarke\checkmark$\checkmark^\checkmark\\checkmarkc\checkmarki\checkmarkr\checkmarkc\checkmark$\checkmark \checkmark$\checkmark^\checkmark\\checkmarkc\checkmarki\checkmarkr\checkmarkc\checkmark$\checkmarkp\checkmark$\checkmark^\checkmark\\checkmarkc\checkmarki\checkmarkr\checkmarkc\checkmark$\checkmarko\checkmark$\checkmark^\checkmark\\checkmarkc\checkmarki\checkmarkr\checkmarkc\checkmark$\checkmarku\checkmark$\checkmark^\checkmark\\checkmarkc\checkmarki\checkmarkr\checkmarkc\checkmark$\checkmarkt\checkmark$\checkmark^\checkmark\\checkmarkc\checkmarki\checkmarkr\checkmarkc\checkmark$\checkmarkr\checkmark$\checkmark^\checkmark\\checkmarkc\checkmarki\checkmarkr\checkmarkc\checkmark$\checkmarke\checkmark$\checkmark^\checkmark\\checkmarkc\checkmarki\checkmarkr\checkmarkc\checkmark$\checkmark \checkmark$\checkmark^\checkmark\\checkmarkc\checkmarki\checkmarkr\checkmarkc\checkmark$\checkmarks\checkmark$\checkmark^\checkmark\\checkmarkc\checkmarki\checkmarkr\checkmarkc\checkmark$\checkmarku\checkmark$\checkmark^\checkmark\\checkmarkc\checkmarki\checkmarkr\checkmarkc\checkmark$\checkmarkr\checkmark$\checkmark^\checkmark\\checkmarkc\checkmarki\checkmarkr\checkmarkc\checkmark$\checkmark \checkmark$\checkmark^\checkmark\\checkmarkc\checkmarki\checkmarkr\checkmarkc\checkmark$\checkmarkd\checkmark$\checkmark^\checkmark\\checkmarkc\checkmarki\checkmarkr\checkmarkc\checkmark$\checkmarke\checkmark$\checkmark^\checkmark\\checkmarkc\checkmarki\checkmarkr\checkmarkc\checkmark$\checkmarku\checkmark$\checkmark^\checkmark\\checkmarkc\checkmarki\checkmarkr\checkmarkc\checkmark$\checkmarkx\checkmark$\checkmark^\checkmark\\checkmarkc\checkmarki\checkmarkr\checkmarkc\checkmark$\checkmark \checkmark$\checkmark^\checkmark\\checkmarkc\checkmarki\checkmarkr\checkmarkc\checkmark$\checkmarka\checkmark$\checkmark^\checkmark\\checkmarkc\checkmarki\checkmarkr\checkmarkc\checkmark$\checkmarkp\checkmark$\checkmark^\checkmark\\checkmarkc\checkmarki\checkmarkr\checkmarkc\checkmark$\checkmarkp\checkmark$\checkmark^\checkmark\\checkmarkc\checkmarki\checkmarkr\checkmarkc\checkmark$\checkmarku\checkmark$\checkmark^\checkmark\\checkmarkc\checkmarki\checkmarkr\checkmarkc\checkmark$\checkmarki\checkmark$\checkmark^\checkmark\\checkmarkc\checkmarki\checkmarkr\checkmarkc\checkmark$\checkmarks\checkmark$\checkmark^\checkmark\\checkmarkc\checkmarki\checkmarkr\checkmarkc\checkmark$\checkmark \checkmark$\checkmark^\checkmark\\checkmarkc\checkmarki\checkmarkr\checkmarkc\checkmark$\checkmark(\checkmark$\checkmark^\checkmark\\checkmarkc\checkmarki\checkmarkr\checkmarkc\checkmark$\checkmarkl\checkmark$\checkmark^\checkmark\\checkmarkc\checkmarki\checkmarkr\checkmarkc\checkmark$\checkmarke\checkmark$\checkmark^\checkmark\\checkmarkc\checkmarki\checkmarkr\checkmarkc\checkmark$\checkmarks\checkmark$\checkmark^\checkmark\\checkmarkc\checkmarki\checkmarkr\checkmarkc\checkmark$\checkmark \checkmark$\checkmark^\checkmark\\checkmarkc\checkmarki\checkmarkr\checkmarkc\checkmark$\checkmarkd\checkmark$\checkmark^\checkmark\\checkmarkc\checkmarki\checkmarkr\checkmarkc\checkmark$\checkmarke\checkmark$\checkmark^\checkmark\\checkmarkc\checkmarki\checkmarkr\checkmarkc\checkmark$\checkmarku\checkmark$\checkmark^\checkmark\\checkmarkc\checkmarki\checkmarkr\checkmarkc\checkmark$\checkmarkx\checkmark$\checkmark^\checkmark\\checkmarkc\checkmarki\checkmarkr\checkmarkc\checkmark$\checkmark \checkmark$\checkmark^\checkmark\\checkmarkc\checkmarki\checkmarkr\checkmarkc\checkmark$\checkmarkp\checkmark$\checkmark^\checkmark\\checkmarkc\checkmarki\checkmarkr\checkmarkc\checkmark$\checkmarka\checkmark$\checkmark^\checkmark\\checkmarkc\checkmarki\checkmarkr\checkmarkc\checkmark$\checkmarkl\checkmark$\checkmark^\checkmark\\checkmarkc\checkmarki\checkmarkr\checkmarkc\checkmark$\checkmarki\checkmark$\checkmark^\checkmark\\checkmarkc\checkmarki\checkmarkr\checkmarkc\checkmark$\checkmarke\checkmark$\checkmark^\checkmark\\checkmarkc\checkmarki\checkmarkr\checkmarkc\checkmark$\checkmarkr\checkmark$\checkmark^\checkmark\\checkmarkc\checkmarki\checkmarkr\checkmarkc\checkmark$\checkmarks\checkmark$\checkmark^\checkmark\\checkmarkc\checkmarki\checkmarkr\checkmarkc\checkmark$\checkmark \checkmark$\checkmark^\checkmark\\checkmarkc\checkmarki\checkmarkr\checkmarkc\checkmark$\checkmarkd\checkmark$\checkmark^\checkmark\\checkmarkc\checkmarki\checkmarkr\checkmarkc\checkmark$\checkmarke\checkmark$\checkmark^\checkmark\\checkmarkc\checkmarki\checkmarkr\checkmarkc\checkmark$\checkmark \checkmark$\checkmark^\checkmark\\checkmarkc\checkmarki\checkmarkr\checkmarkc\checkmark$\checkmarkr\checkmark$\checkmark^\checkmark\\checkmarkc\checkmarki\checkmarkr\checkmarkc\checkmark$\checkmarko\checkmark$\checkmark^\checkmark\\checkmarkc\checkmarki\checkmarkr\checkmarkc\checkmark$\checkmarku\checkmark$\checkmark^\checkmark\\checkmarkc\checkmarki\checkmarkr\checkmarkc\checkmark$\checkmarkl\checkmark$\checkmark^\checkmark\\checkmarkc\checkmarki\checkmarkr\checkmarkc\checkmark$\checkmarke\checkmark$\checkmark^\checkmark\\checkmarkc\checkmarki\checkmarkr\checkmarkc\checkmark$\checkmarkm\checkmark$\checkmark^\checkmark\\checkmarkc\checkmarki\checkmarkr\checkmarkc\checkmark$\checkmarke\checkmark$\checkmark^\checkmark\\checkmarkc\checkmarki\checkmarkr\checkmarkc\checkmark$\checkmarkn\checkmark$\checkmark^\checkmark\\checkmarkc\checkmarki\checkmarkr\checkmarkc\checkmark$\checkmarkt\checkmark$\checkmark^\checkmark\\checkmarkc\checkmarki\checkmarkr\checkmarkc\checkmark$\checkmark \checkmark$\checkmark^\checkmark\\checkmarkc\checkmarki\checkmarkr\checkmarkc\checkmark$\checkmarke\checkmark$\checkmark^\checkmark\\checkmarkc\checkmarki\checkmarkr\checkmarkc\checkmark$\checkmarks\checkmark$\checkmark^\checkmark\\checkmarkc\checkmarki\checkmarkr\checkmarkc\checkmark$\checkmarkp\checkmark$\checkmark^\checkmark\\checkmarkc\checkmarki\checkmarkr\checkmarkc\checkmark$\checkmarka\checkmark$\checkmark^\checkmark\\checkmarkc\checkmarki\checkmarkr\checkmarkc\checkmark$\checkmarkc\checkmark$\checkmark^\checkmark\\checkmarkc\checkmarki\checkmarkr\checkmarkc\checkmark$\checkmarké\checkmark$\checkmark^\checkmark\\checkmarkc\checkmarki\checkmarkr\checkmarkc\checkmark$\checkmarks\checkmark$\checkmark^\checkmark\\checkmarkc\checkmarki\checkmarkr\checkmarkc\checkmark$\checkmark \checkmark$\checkmark^\checkmark\\checkmarkc\checkmarki\checkmarkr\checkmarkc\checkmark$\checkmarkd\checkmark$\checkmark^\checkmark\\checkmarkc\checkmarki\checkmarkr\checkmarkc\checkmark$\checkmarke\checkmark$\checkmark^\checkmark\\checkmarkc\checkmarki\checkmarkr\checkmarkc\checkmark$\checkmark \checkmark$\checkmark^\checkmark\\checkmarkc\checkmarki\checkmarkr\checkmarkc\checkmark$\checkmark$\checkmark$\checkmark^\checkmark\\checkmarkc\checkmarki\checkmarkr\checkmarkc\checkmark$\checkmarkL\checkmark$\checkmark^\checkmark\\checkmarkc\checkmarki\checkmarkr\checkmarkc\checkmark$\checkmark_\checkmark$\checkmark^\checkmark\\checkmarkc\checkmarki\checkmarkr\checkmarkc\checkmark$\checkmark{\checkmark$\checkmark^\checkmark\\checkmarkc\checkmarki\checkmarkr\checkmarkc\checkmark$\checkmarkp\checkmark$\checkmark^\checkmark\\checkmarkc\checkmarki\checkmarkr\checkmarkc\checkmark$\checkmarka\checkmark$\checkmark^\checkmark\\checkmarkc\checkmarki\checkmarkr\checkmarkc\checkmark$\checkmarkl\checkmark$\checkmark^\checkmark\\checkmarkc\checkmarki\checkmarkr\checkmarkc\checkmark$\checkmarki\checkmark$\checkmark^\checkmark\\checkmarkc\checkmarki\checkmarkr\checkmarkc\checkmark$\checkmarke\checkmark$\checkmark^\checkmark\\checkmarkc\checkmarki\checkmarkr\checkmarkc\checkmark$\checkmarkr\checkmark$\checkmark^\checkmark\\checkmarkc\checkmarki\checkmarkr\checkmarkc\checkmark$\checkmarks\checkmark$\checkmark^\checkmark\\checkmarkc\checkmarki\checkmarkr\checkmarkc\checkmark$\checkmark}\checkmark$\checkmark^\checkmark\\checkmarkc\checkmarki\checkmarkr\checkmarkc\checkmark$\checkmark \checkmark$\checkmark^\checkmark\\checkmarkc\checkmarki\checkmarkr\checkmarkc\checkmark$\checkmark=\checkmark$\checkmark^\checkmark\\checkmarkc\checkmarki\checkmarkr\checkmarkc\checkmark$\checkmark \checkmark$\checkmark^\checkmark\\checkmarkc\checkmarki\checkmarkr\checkmarkc\checkmark$\checkmark1\checkmark$\checkmark^\checkmark\\checkmarkc\checkmarki\checkmarkr\checkmarkc\checkmark$\checkmark5\checkmark$\checkmark^\checkmark\\checkmarkc\checkmarki\checkmarkr\checkmarkc\checkmark$\checkmark0\checkmark$\checkmark^\checkmark\\checkmarkc\checkmarki\checkmarkr\checkmarkc\checkmark$\checkmark$\checkmark$\checkmark^\checkmark\\checkmarkc\checkmarki\checkmarkr\checkmarkc\checkmark$\checkmark \checkmark$\checkmark^\checkmark\\checkmarkc\checkmarki\checkmarkr\checkmarkc\checkmark$\checkmarkm\checkmark$\checkmark^\checkmark\\checkmarkc\checkmarki\checkmarkr\checkmarkc\checkmark$\checkmarkm\checkmark$\checkmark^\checkmark\\checkmarkc\checkmarki\checkmarkr\checkmarkc\checkmark$\checkmark)\checkmark$\checkmark^\checkmark\\checkmarkc\checkmarki\checkmarkr\checkmarkc\checkmark$\checkmark.\checkmark$\checkmark^\checkmark\\checkmarkc\checkmarki\checkmarkr\checkmarkc\checkmark$\checkmark \checkmark$\checkmark^\checkmark\\checkmarkc\checkmarki\checkmarkr\checkmarkc\checkmark$\checkmarkL\checkmark$\checkmark^\checkmark\\checkmarkc\checkmarki\checkmarkr\checkmarkc\checkmark$\checkmarke\checkmark$\checkmark^\checkmark\\checkmarkc\checkmarki\checkmarkr\checkmarkc\checkmark$\checkmarks\checkmark$\checkmark^\checkmark\\checkmarkc\checkmarki\checkmarkr\checkmarkc\checkmark$\checkmark \checkmark$\checkmark^\checkmark\\checkmarkc\checkmarki\checkmarkr\checkmarkc\checkmark$\checkmarkc\checkmark$\checkmark^\checkmark\\checkmarkc\checkmarki\checkmarkr\checkmarkc\checkmark$\checkmarkh\checkmark$\checkmark^\checkmark\\checkmarkc\checkmarki\checkmarkr\checkmarkc\checkmark$\checkmarka\checkmark$\checkmark^\checkmark\\checkmarkc\checkmarki\checkmarkr\checkmarkc\checkmark$\checkmarkr\checkmark$\checkmark^\checkmark\\checkmarkc\checkmarki\checkmarkr\checkmarkc\checkmark$\checkmarkg\checkmark$\checkmark^\checkmark\\checkmarkc\checkmarki\checkmarkr\checkmarkc\checkmark$\checkmarke\checkmark$\checkmark^\checkmark\\checkmarkc\checkmarki\checkmarkr\checkmarkc\checkmark$\checkmarks\checkmark$\checkmark^\checkmark\\checkmarkc\checkmarki\checkmarkr\checkmarkc\checkmark$\checkmark \checkmark$\checkmark^\checkmark\\checkmarkc\checkmarki\checkmarkr\checkmarkc\checkmark$\checkmarka\checkmark$\checkmark^\checkmark\\checkmarkc\checkmarki\checkmarkr\checkmarkc\checkmark$\checkmarkp\checkmark$\checkmark^\checkmark\\checkmarkc\checkmarki\checkmarkr\checkmarkc\checkmark$\checkmarkp\checkmark$\checkmark^\checkmark\\checkmarkc\checkmarki\checkmarkr\checkmarkc\checkmark$\checkmarkl\checkmark$\checkmark^\checkmark\\checkmarkc\checkmarki\checkmarkr\checkmarkc\checkmark$\checkmarki\checkmark$\checkmark^\checkmark\\checkmarkc\checkmarki\checkmarkr\checkmarkc\checkmark$\checkmarkq\checkmark$\checkmark^\checkmark\\checkmarkc\checkmarki\checkmarkr\checkmarkc\checkmark$\checkmarku\checkmark$\checkmark^\checkmark\\checkmarkc\checkmarki\checkmarkr\checkmarkc\checkmark$\checkmarké\checkmark$\checkmark^\checkmark\\checkmarkc\checkmarki\checkmarkr\checkmarkc\checkmark$\checkmarke\checkmark$\checkmark^\checkmark\\checkmarkc\checkmarki\checkmarkr\checkmarkc\checkmark$\checkmarks\checkmark$\checkmark^\checkmark\\checkmarkc\checkmarki\checkmarkr\checkmarkc\checkmark$\checkmark \checkmark$\checkmark^\checkmark\\checkmarkc\checkmarki\checkmarkr\checkmarkc\checkmark$\checkmarks\checkmark$\checkmark^\checkmark\\checkmarkc\checkmarki\checkmarkr\checkmarkc\checkmark$\checkmarko\checkmark$\checkmark^\checkmark\\checkmarkc\checkmarki\checkmarkr\checkmarkc\checkmark$\checkmarkn\checkmark$\checkmark^\checkmark\\checkmarkc\checkmarki\checkmarkr\checkmarkc\checkmark$\checkmarkt\checkmark$\checkmark^\checkmark\\checkmarkc\checkmarki\checkmarkr\checkmarkc\checkmark$\checkmark \checkmark$\checkmark^\checkmark\\checkmarkc\checkmarki\checkmarkr\checkmarkc\checkmark$\checkmarkl\checkmark$\checkmark^\checkmark\\checkmarkc\checkmarki\checkmarkr\checkmarkc\checkmark$\checkmarke\checkmark$\checkmark^\checkmark\\checkmarkc\checkmarki\checkmarkr\checkmarkc\checkmark$\checkmark \checkmark$\checkmark^\checkmark\\checkmarkc\checkmarki\checkmarkr\checkmarkc\checkmark$\checkmarkm\checkmark$\checkmark^\checkmark\\checkmarkc\checkmarki\checkmarkr\checkmarkc\checkmark$\checkmarko\checkmark$\checkmark^\checkmark\\checkmarkc\checkmarki\checkmarkr\checkmarkc\checkmark$\checkmarkm\checkmark$\checkmark^\checkmark\\checkmarkc\checkmarki\checkmarkr\checkmarkc\checkmark$\checkmarke\checkmark$\checkmark^\checkmark\\checkmarkc\checkmarki\checkmarkr\checkmarkc\checkmark$\checkmarkn\checkmark$\checkmark^\checkmark\\checkmarkc\checkmarki\checkmarkr\checkmarkc\checkmark$\checkmarkt\checkmark$\checkmark^\checkmark\\checkmarkc\checkmarki\checkmarkr\checkmarkc\checkmark$\checkmark \checkmark$\checkmark^\checkmark\\checkmarkc\checkmarki\checkmarkr\checkmarkc\checkmark$\checkmarkd\checkmark$\checkmark^\checkmark\\checkmarkc\checkmarki\checkmarkr\checkmarkc\checkmark$\checkmarke\checkmark$\checkmark^\checkmark\\checkmarkc\checkmarki\checkmarkr\checkmarkc\checkmark$\checkmark \checkmark$\checkmark^\checkmark\\checkmarkc\checkmarki\checkmarkr\checkmarkc\checkmark$\checkmarkr\checkmark$\checkmark^\checkmark\\checkmarkc\checkmarki\checkmarkr\checkmarkc\checkmark$\checkmarke\checkmark$\checkmark^\checkmark\\checkmarkc\checkmarki\checkmarkr\checkmarkc\checkmark$\checkmarkn\checkmark$\checkmark^\checkmark\\checkmarkc\checkmarki\checkmarkr\checkmarkc\checkmark$\checkmarkv\checkmark$\checkmark^\checkmark\\checkmarkc\checkmarki\checkmarkr\checkmarkc\checkmark$\checkmarke\checkmark$\checkmark^\checkmark\\checkmarkc\checkmarki\checkmarkr\checkmarkc\checkmark$\checkmarkr\checkmark$\checkmark^\checkmark\\checkmarkc\checkmarki\checkmarkr\checkmarkc\checkmark$\checkmarks\checkmark$\checkmark^\checkmark\\checkmarkc\checkmarki\checkmarkr\checkmarkc\checkmark$\checkmarke\checkmark$\checkmark^\checkmark\\checkmarkc\checkmarki\checkmarkr\checkmarkc\checkmark$\checkmarkm\checkmark$\checkmark^\checkmark\\checkmarkc\checkmarki\checkmarkr\checkmarkc\checkmark$\checkmarke\checkmark$\checkmark^\checkmark\\checkmarkc\checkmarki\checkmarkr\checkmarkc\checkmark$\checkmarkn\checkmark$\checkmark^\checkmark\\checkmarkc\checkmarki\checkmarkr\checkmarkc\checkmark$\checkmarkt\checkmark$\checkmark^\checkmark\\checkmarkc\checkmarki\checkmarkr\checkmarkc\checkmark$\checkmark \checkmark$\checkmark^\checkmark\\checkmarkc\checkmarki\checkmarkr\checkmarkc\checkmark$\checkmark$\checkmark$\checkmark^\checkmark\\checkmarkc\checkmarki\checkmarkr\checkmarkc\checkmark$\checkmarkM\checkmark$\checkmark^\checkmark\\checkmarkc\checkmarki\checkmarkr\checkmarkc\checkmark$\checkmark_\checkmark$\checkmark^\checkmark\\checkmarkc\checkmarki\checkmarkr\checkmarkc\checkmark$\checkmark{\checkmark$\checkmark^\checkmark\\checkmarkc\checkmarki\checkmarkr\checkmarkc\checkmark$\checkmarkr\checkmark$\checkmark^\checkmark\\checkmarkc\checkmarki\checkmarkr\checkmarkc\checkmark$\checkmarke\checkmark$\checkmark^\checkmark\\checkmarkc\checkmarki\checkmarkr\checkmarkc\checkmark$\checkmarkn\checkmark$\checkmark^\checkmark\\checkmarkc\checkmarki\checkmarkr\checkmarkc\checkmark$\checkmarkv\checkmark$\checkmark^\checkmark\\checkmarkc\checkmarki\checkmarkr\checkmarkc\checkmark$\checkmark}\checkmark$\checkmark^\checkmark\\checkmarkc\checkmarki\checkmarkr\checkmarkc\checkmark$\checkmark$\checkmark$\checkmark^\checkmark\\checkmarkc\checkmarki\checkmarkr\checkmarkc\checkmark$\checkmark \checkmark$\checkmark^\checkmark\\checkmarkc\checkmarki\checkmarkr\checkmarkc\checkmark$\checkmarke\checkmark$\checkmark^\checkmark\\checkmarkc\checkmarki\checkmarkr\checkmarkc\checkmark$\checkmarkt\checkmark$\checkmark^\checkmark\\checkmarkc\checkmarki\checkmarkr\checkmarkc\checkmark$\checkmark \checkmark$\checkmark^\checkmark\\checkmarkc\checkmarki\checkmarkr\checkmarkc\checkmark$\checkmarkl\checkmark$\checkmark^\checkmark\\checkmarkc\checkmarki\checkmarkr\checkmarkc\checkmark$\checkmarka\checkmark$\checkmark^\checkmark\\checkmarkc\checkmarki\checkmarkr\checkmarkc\checkmark$\checkmark \checkmark$\checkmark^\checkmark\\checkmarkc\checkmarki\checkmarkr\checkmarkc\checkmark$\checkmarkf\checkmark$\checkmark^\checkmark\\checkmarkc\checkmarki\checkmarkr\checkmarkc\checkmark$\checkmarko\checkmark$\checkmark^\checkmark\\checkmarkc\checkmarki\checkmarkr\checkmarkc\checkmark$\checkmarkr\checkmark$\checkmark^\checkmark\\checkmarkc\checkmarki\checkmarkr\checkmarkc\checkmark$\checkmarkc\checkmark$\checkmark^\checkmark\\checkmarkc\checkmarki\checkmarkr\checkmarkc\checkmark$\checkmarke\checkmark$\checkmark^\checkmark\\checkmarkc\checkmarki\checkmarkr\checkmarkc\checkmark$\checkmark \checkmark$\checkmark^\checkmark\\checkmarkc\checkmarki\checkmarkr\checkmarkc\checkmark$\checkmarkr\checkmark$\checkmark^\checkmark\\checkmarkc\checkmarki\checkmarkr\checkmarkc\checkmark$\checkmarka\checkmark$\checkmark^\checkmark\\checkmarkc\checkmarki\checkmarkr\checkmarkc\checkmark$\checkmarkd\checkmark$\checkmark^\checkmark\\checkmarkc\checkmarki\checkmarkr\checkmarkc\checkmark$\checkmarki\checkmark$\checkmark^\checkmark\\checkmarkc\checkmarki\checkmarkr\checkmarkc\checkmark$\checkmarka\checkmark$\checkmark^\checkmark\\checkmarkc\checkmarki\checkmarkr\checkmarkc\checkmark$\checkmarkl\checkmark$\checkmark^\checkmark\\checkmarkc\checkmarki\checkmarkr\checkmarkc\checkmark$\checkmarke\checkmark$\checkmark^\checkmark\\checkmarkc\checkmarki\checkmarkr\checkmarkc\checkmark$\checkmark \checkmark$\checkmark^\checkmark\\checkmarkc\checkmarki\checkmarkr\checkmarkc\checkmark$\checkmarkd\checkmark$\checkmark^\checkmark\\checkmarkc\checkmarki\checkmarkr\checkmarkc\checkmark$\checkmarke\checkmark$\checkmark^\checkmark\\checkmarkc\checkmarki\checkmarkr\checkmarkc\checkmark$\checkmark \checkmark$\checkmark^\checkmark\\checkmarkc\checkmarki\checkmarkr\checkmarkc\checkmark$\checkmarkl\checkmark$\checkmark^\checkmark\\checkmarkc\checkmarki\checkmarkr\checkmarkc\checkmark$\checkmarka\checkmark$\checkmark^\checkmark\\checkmarkc\checkmarki\checkmarkr\checkmarkc\checkmark$\checkmark \checkmark$\checkmark^\checkmark\\checkmarkc\checkmarki\checkmarkr\checkmarkc\checkmark$\checkmarkc\checkmark$\checkmark^\checkmark\\checkmarkc\checkmarki\checkmarkr\checkmarkc\checkmark$\checkmarko\checkmark$\checkmark^\checkmark\\checkmarkc\checkmarki\checkmarkr\checkmarkc\checkmark$\checkmarku\checkmark$\checkmark^\checkmark\\checkmarkc\checkmarki\checkmarkr\checkmarkc\checkmark$\checkmarkr\checkmark$\checkmark^\checkmark\\checkmarkc\checkmarki\checkmarkr\checkmarkc\checkmark$\checkmarkr\checkmark$\checkmark^\checkmark\\checkmarkc\checkmarki\checkmarkr\checkmarkc\checkmark$\checkmarko\checkmark$\checkmark^\checkmark\\checkmarkc\checkmarki\checkmarkr\checkmarkc\checkmark$\checkmarki\checkmark$\checkmark^\checkmark\\checkmarkc\checkmarki\checkmarkr\checkmarkc\checkmark$\checkmarke\checkmark$\checkmark^\checkmark\\checkmarkc\checkmarki\checkmarkr\checkmarkc\checkmark$\checkmark \checkmark$\checkmark^\checkmark\\checkmarkc\checkmarki\checkmarkr\checkmarkc\checkmark$\checkmark$\checkmark$\checkmark^\checkmark\\checkmarkc\checkmarki\checkmarkr\checkmarkc\checkmark$\checkmarkF\checkmark$\checkmark^\checkmark\\checkmarkc\checkmarki\checkmarkr\checkmarkc\checkmark$\checkmark_\checkmark$\checkmark^\checkmark\\checkmarkc\checkmarki\checkmarkr\checkmarkc\checkmark$\checkmark{\checkmark$\checkmark^\checkmark\\checkmarkc\checkmarki\checkmarkr\checkmarkc\checkmark$\checkmarkc\checkmark$\checkmark^\checkmark\\checkmarkc\checkmarki\checkmarkr\checkmarkc\checkmark$\checkmarko\checkmark$\checkmark^\checkmark\\checkmarkc\checkmarki\checkmarkr\checkmarkc\checkmark$\checkmarku\checkmark$\checkmark^\checkmark\\checkmarkc\checkmarki\checkmarkr\checkmarkc\checkmark$\checkmarkr\checkmark$\checkmark^\checkmark\\checkmarkc\checkmarki\checkmarkr\checkmarkc\checkmark$\checkmarkr\checkmark$\checkmark^\checkmark\\checkmarkc\checkmarki\checkmarkr\checkmarkc\checkmark$\checkmarko\checkmark$\checkmark^\checkmark\\checkmarkc\checkmarki\checkmarkr\checkmarkc\checkmark$\checkmarki\checkmark$\checkmark^\checkmark\\checkmarkc\checkmarki\checkmarkr\checkmarkc\checkmark$\checkmarke\checkmark$\checkmark^\checkmark\\checkmarkc\checkmarki\checkmarkr\checkmarkc\checkmark$\checkmark}\checkmark$\checkmark^\checkmark\\checkmarkc\checkmarki\checkmarkr\checkmarkc\checkmark$\checkmark \checkmark$\checkmark^\checkmark\\checkmarkc\checkmarki\checkmarkr\checkmarkc\checkmark$\checkmark=\checkmark$\checkmark^\checkmark\\checkmarkc\checkmarki\checkmarkr\checkmarkc\checkmark$\checkmark \checkmark$\checkmark^\checkmark\\checkmarkc\checkmarki\checkmarkr\checkmarkc\checkmark$\checkmark3\checkmark$\checkmark^\checkmark\\checkmarkc\checkmarki\checkmarkr\checkmarkc\checkmark$\checkmark3\checkmark$\checkmark^\checkmark\\checkmarkc\checkmarki\checkmarkr\checkmarkc\checkmark$\checkmark3\checkmark$\checkmark^\checkmark\\checkmarkc\checkmarki\checkmarkr\checkmarkc\checkmark$\checkmark$\checkmark$\checkmark^\checkmark\\checkmarkc\checkmarki\checkmarkr\checkmarkc\checkmark$\checkmark \checkmark$\checkmark^\checkmark\\checkmarkc\checkmarki\checkmarkr\checkmarkc\checkmark$\checkmarkN\checkmark$\checkmark^\checkmark\\checkmarkc\checkmarki\checkmarkr\checkmarkc\checkmark$\checkmark.\checkmark$\checkmark^\checkmark\\checkmarkc\checkmarki\checkmarkr\checkmarkc\checkmark$\checkmark
\checkmark$\checkmark^\checkmark\\checkmarkc\checkmarki\checkmarkr\checkmarkc\checkmark$\checkmark
\checkmark$\checkmark^\checkmark\\checkmarkc\checkmarki\checkmarkr\checkmarkc\checkmark$\checkmark\\checkmark$\checkmark^\checkmark\\checkmarkc\checkmarki\checkmarkr\checkmarkc\checkmark$\checkmarkt\checkmark$\checkmark^\checkmark\\checkmarkc\checkmarki\checkmarkr\checkmarkc\checkmark$\checkmarke\checkmark$\checkmark^\checkmark\\checkmarkc\checkmarki\checkmarkr\checkmarkc\checkmark$\checkmarkx\checkmark$\checkmark^\checkmark\\checkmarkc\checkmarki\checkmarkr\checkmarkc\checkmark$\checkmarkt\checkmark$\checkmark^\checkmark\\checkmarkc\checkmarki\checkmarkr\checkmarkc\checkmark$\checkmarkb\checkmark$\checkmark^\checkmark\\checkmarkc\checkmarki\checkmarkr\checkmarkc\checkmark$\checkmarkf\checkmark$\checkmark^\checkmark\\checkmarkc\checkmarki\checkmarkr\checkmarkc\checkmark$\checkmark{\checkmark$\checkmark^\checkmark\\checkmarkc\checkmarki\checkmarkr\checkmarkc\checkmark$\checkmarkR\checkmark$\checkmark^\checkmark\\checkmarkc\checkmarki\checkmarkr\checkmarkc\checkmark$\checkmarké\checkmark$\checkmark^\checkmark\\checkmarkc\checkmarki\checkmarkr\checkmarkc\checkmark$\checkmarka\checkmark$\checkmark^\checkmark\\checkmarkc\checkmarki\checkmarkr\checkmarkc\checkmark$\checkmarkc\checkmark$\checkmark^\checkmark\\checkmarkc\checkmarki\checkmarkr\checkmarkc\checkmark$\checkmarkt\checkmark$\checkmark^\checkmark\\checkmarkc\checkmarki\checkmarkr\checkmarkc\checkmark$\checkmarki\checkmark$\checkmark^\checkmark\\checkmarkc\checkmarki\checkmarkr\checkmarkc\checkmark$\checkmarko\checkmark$\checkmark^\checkmark\\checkmarkc\checkmarki\checkmarkr\checkmarkc\checkmark$\checkmarkn\checkmark$\checkmark^\checkmark\\checkmarkc\checkmarki\checkmarkr\checkmarkc\checkmark$\checkmarks\checkmark$\checkmark^\checkmark\\checkmarkc\checkmarki\checkmarkr\checkmarkc\checkmark$\checkmark \checkmark$\checkmark^\checkmark\\checkmarkc\checkmarki\checkmarkr\checkmarkc\checkmark$\checkmarka\checkmark$\checkmark^\checkmark\\checkmarkc\checkmarki\checkmarkr\checkmarkc\checkmark$\checkmarku\checkmark$\checkmark^\checkmark\\checkmarkc\checkmarki\checkmarkr\checkmarkc\checkmark$\checkmarkx\checkmark$\checkmark^\checkmark\\checkmarkc\checkmarki\checkmarkr\checkmarkc\checkmark$\checkmark \checkmark$\checkmark^\checkmark\\checkmarkc\checkmarki\checkmarkr\checkmarkc\checkmark$\checkmarka\checkmark$\checkmark^\checkmark\\checkmarkc\checkmarki\checkmarkr\checkmarkc\checkmark$\checkmarkp\checkmark$\checkmark^\checkmark\\checkmarkc\checkmarki\checkmarkr\checkmarkc\checkmark$\checkmarkp\checkmark$\checkmark^\checkmark\\checkmarkc\checkmarki\checkmarkr\checkmarkc\checkmark$\checkmarku\checkmark$\checkmark^\checkmark\\checkmarkc\checkmarki\checkmarkr\checkmarkc\checkmark$\checkmarki\checkmark$\checkmark^\checkmark\\checkmarkc\checkmarki\checkmarkr\checkmarkc\checkmark$\checkmarks\checkmark$\checkmark^\checkmark\\checkmarkc\checkmarki\checkmarkr\checkmarkc\checkmark$\checkmark \checkmark$\checkmark^\checkmark\\checkmarkc\checkmarki\checkmarkr\checkmarkc\checkmark$\checkmark:\checkmark$\checkmark^\checkmark\\checkmarkc\checkmarki\checkmarkr\checkmarkc\checkmark$\checkmark}\checkmark$\checkmark^\checkmark\\checkmarkc\checkmarki\checkmarkr\checkmarkc\checkmark$\checkmark
\checkmark$\checkmark^\checkmark\\checkmarkc\checkmarki\checkmarkr\checkmarkc\checkmark$\checkmark\\checkmark$\checkmark^\checkmark\\checkmarkc\checkmarki\checkmarkr\checkmarkc\checkmark$\checkmarkb\checkmark$\checkmark^\checkmark\\checkmarkc\checkmarki\checkmarkr\checkmarkc\checkmark$\checkmarke\checkmark$\checkmark^\checkmark\\checkmarkc\checkmarki\checkmarkr\checkmarkc\checkmark$\checkmarkg\checkmark$\checkmark^\checkmark\\checkmarkc\checkmarki\checkmarkr\checkmarkc\checkmark$\checkmarki\checkmark$\checkmark^\checkmark\\checkmarkc\checkmarki\checkmarkr\checkmarkc\checkmark$\checkmarkn\checkmark$\checkmark^\checkmark\\checkmarkc\checkmarki\checkmarkr\checkmarkc\checkmark$\checkmark{\checkmark$\checkmark^\checkmark\\checkmarkc\checkmarki\checkmarkr\checkmarkc\checkmark$\checkmarke\checkmark$\checkmark^\checkmark\\checkmarkc\checkmarki\checkmarkr\checkmarkc\checkmark$\checkmarkq\checkmark$\checkmark^\checkmark\\checkmarkc\checkmarki\checkmarkr\checkmarkc\checkmark$\checkmarku\checkmark$\checkmark^\checkmark\\checkmarkc\checkmarki\checkmarkr\checkmarkc\checkmark$\checkmarka\checkmark$\checkmark^\checkmark\\checkmarkc\checkmarki\checkmarkr\checkmarkc\checkmark$\checkmarkt\checkmark$\checkmark^\checkmark\\checkmarkc\checkmarki\checkmarkr\checkmarkc\checkmark$\checkmarki\checkmark$\checkmark^\checkmark\\checkmarkc\checkmarki\checkmarkr\checkmarkc\checkmark$\checkmarko\checkmark$\checkmark^\checkmark\\checkmarkc\checkmarki\checkmarkr\checkmarkc\checkmark$\checkmarkn\checkmark$\checkmark^\checkmark\\checkmarkc\checkmarki\checkmarkr\checkmarkc\checkmark$\checkmark}\checkmark$\checkmark^\checkmark\\checkmarkc\checkmarki\checkmarkr\checkmarkc\checkmark$\checkmark
\checkmark$\checkmark^\checkmark\\checkmarkc\checkmarki\checkmarkr\checkmarkc\checkmark$\checkmark \checkmark$\checkmark^\checkmark\\checkmarkc\checkmarki\checkmarkr\checkmarkc\checkmark$\checkmark \checkmark$\checkmark^\checkmark\\checkmarkc\checkmarki\checkmarkr\checkmarkc\checkmark$\checkmark \checkmark$\checkmark^\checkmark\\checkmarkc\checkmarki\checkmarkr\checkmarkc\checkmark$\checkmark \checkmark$\checkmark^\checkmark\\checkmarkc\checkmarki\checkmarkr\checkmarkc\checkmark$\checkmarkF\checkmark$\checkmark^\checkmark\\checkmarkc\checkmarki\checkmarkr\checkmarkc\checkmark$\checkmark_\checkmark$\checkmark^\checkmark\\checkmarkc\checkmarki\checkmarkr\checkmarkc\checkmark$\checkmark{\checkmark$\checkmark^\checkmark\\checkmarkc\checkmarki\checkmarkr\checkmarkc\checkmark$\checkmarkr\checkmark$\checkmark^\checkmark\\checkmarkc\checkmarki\checkmarkr\checkmarkc\checkmark$\checkmark\\checkmark$\checkmark^\checkmark\\checkmarkc\checkmarki\checkmarkr\checkmarkc\checkmark$\checkmark_\checkmark$\checkmark^\checkmark\\checkmarkc\checkmarki\checkmarkr\checkmarkc\checkmark$\checkmarkp\checkmark$\checkmark^\checkmark\\checkmarkc\checkmarki\checkmarkr\checkmarkc\checkmark$\checkmarka\checkmark$\checkmark^\checkmark\\checkmarkc\checkmarki\checkmarkr\checkmarkc\checkmark$\checkmarkl\checkmark$\checkmark^\checkmark\\checkmarkc\checkmarki\checkmarkr\checkmarkc\checkmark$\checkmarki\checkmark$\checkmark^\checkmark\\checkmarkc\checkmarki\checkmarkr\checkmarkc\checkmark$\checkmarke\checkmark$\checkmark^\checkmark\\checkmarkc\checkmarki\checkmarkr\checkmarkc\checkmark$\checkmarkr\checkmark$\checkmark^\checkmark\\checkmarkc\checkmarki\checkmarkr\checkmarkc\checkmark$\checkmark}\checkmark$\checkmark^\checkmark\\checkmarkc\checkmarki\checkmarkr\checkmarkc\checkmark$\checkmark \checkmark$\checkmark^\checkmark\\checkmarkc\checkmarki\checkmarkr\checkmarkc\checkmark$\checkmark=\checkmark$\checkmark^\checkmark\\checkmarkc\checkmarki\checkmarkr\checkmarkc\checkmark$\checkmark \checkmark$\checkmark^\checkmark\\checkmarkc\checkmarki\checkmarkr\checkmarkc\checkmark$\checkmark\\checkmark$\checkmark^\checkmark\\checkmarkc\checkmarki\checkmarkr\checkmarkc\checkmark$\checkmarkf\checkmark$\checkmark^\checkmark\\checkmarkc\checkmarki\checkmarkr\checkmarkc\checkmark$\checkmarkr\checkmark$\checkmark^\checkmark\\checkmarkc\checkmarki\checkmarkr\checkmarkc\checkmark$\checkmarka\checkmark$\checkmark^\checkmark\\checkmarkc\checkmarki\checkmarkr\checkmarkc\checkmark$\checkmarkc\checkmark$\checkmark^\checkmark\\checkmarkc\checkmarki\checkmarkr\checkmarkc\checkmark$\checkmark{\checkmark$\checkmark^\checkmark\\checkmarkc\checkmarki\checkmarkr\checkmarkc\checkmark$\checkmarkM\checkmark$\checkmark^\checkmark\\checkmarkc\checkmarki\checkmarkr\checkmarkc\checkmark$\checkmark_\checkmark$\checkmark^\checkmark\\checkmarkc\checkmarki\checkmarkr\checkmarkc\checkmark$\checkmark{\checkmark$\checkmark^\checkmark\\checkmarkc\checkmarki\checkmarkr\checkmarkc\checkmark$\checkmarkr\checkmark$\checkmark^\checkmark\\checkmarkc\checkmarki\checkmarkr\checkmarkc\checkmark$\checkmarke\checkmark$\checkmark^\checkmark\\checkmarkc\checkmarki\checkmarkr\checkmarkc\checkmark$\checkmarkn\checkmark$\checkmark^\checkmark\\checkmarkc\checkmarki\checkmarkr\checkmarkc\checkmark$\checkmarkv\checkmark$\checkmark^\checkmark\\checkmarkc\checkmarki\checkmarkr\checkmarkc\checkmark$\checkmark}\checkmark$\checkmark^\checkmark\\checkmarkc\checkmarki\checkmarkr\checkmarkc\checkmark$\checkmark}\checkmark$\checkmark^\checkmark\\checkmarkc\checkmarki\checkmarkr\checkmarkc\checkmark$\checkmark{\checkmark$\checkmark^\checkmark\\checkmarkc\checkmarki\checkmarkr\checkmarkc\checkmark$\checkmarkL\checkmark$\checkmark^\checkmark\\checkmarkc\checkmarki\checkmarkr\checkmarkc\checkmark$\checkmark_\checkmark$\checkmark^\checkmark\\checkmarkc\checkmarki\checkmarkr\checkmarkc\checkmark$\checkmark{\checkmark$\checkmark^\checkmark\\checkmarkc\checkmarki\checkmarkr\checkmarkc\checkmark$\checkmarkp\checkmark$\checkmark^\checkmark\\checkmarkc\checkmarki\checkmarkr\checkmarkc\checkmark$\checkmarka\checkmark$\checkmark^\checkmark\\checkmarkc\checkmarki\checkmarkr\checkmarkc\checkmark$\checkmarkl\checkmark$\checkmark^\checkmark\\checkmarkc\checkmarki\checkmarkr\checkmarkc\checkmark$\checkmarki\checkmark$\checkmark^\checkmark\\checkmarkc\checkmarki\checkmarkr\checkmarkc\checkmark$\checkmarke\checkmark$\checkmark^\checkmark\\checkmarkc\checkmarki\checkmarkr\checkmarkc\checkmark$\checkmarkr\checkmark$\checkmark^\checkmark\\checkmarkc\checkmarki\checkmarkr\checkmarkc\checkmark$\checkmarks\checkmark$\checkmark^\checkmark\\checkmarkc\checkmarki\checkmarkr\checkmarkc\checkmark$\checkmark}\checkmark$\checkmark^\checkmark\\checkmarkc\checkmarki\checkmarkr\checkmarkc\checkmark$\checkmark}\checkmark$\checkmark^\checkmark\\checkmarkc\checkmarki\checkmarkr\checkmarkc\checkmark$\checkmark \checkmark$\checkmark^\checkmark\\checkmarkc\checkmarki\checkmarkr\checkmarkc\checkmark$\checkmark+\checkmark$\checkmark^\checkmark\\checkmarkc\checkmarki\checkmarkr\checkmarkc\checkmark$\checkmark \checkmark$\checkmark^\checkmark\\checkmarkc\checkmarki\checkmarkr\checkmarkc\checkmark$\checkmark\\checkmark$\checkmark^\checkmark\\checkmarkc\checkmarki\checkmarkr\checkmarkc\checkmark$\checkmarkf\checkmark$\checkmark^\checkmark\\checkmarkc\checkmarki\checkmarkr\checkmarkc\checkmark$\checkmarkr\checkmark$\checkmark^\checkmark\\checkmarkc\checkmarki\checkmarkr\checkmarkc\checkmark$\checkmarka\checkmark$\checkmark^\checkmark\\checkmarkc\checkmarki\checkmarkr\checkmarkc\checkmark$\checkmarkc\checkmark$\checkmark^\checkmark\\checkmarkc\checkmarki\checkmarkr\checkmarkc\checkmark$\checkmark{\checkmark$\checkmark^\checkmark\\checkmarkc\checkmarki\checkmarkr\checkmarkc\checkmark$\checkmarkF\checkmark$\checkmark^\checkmark\\checkmarkc\checkmarki\checkmarkr\checkmarkc\checkmark$\checkmark_\checkmark$\checkmark^\checkmark\\checkmarkc\checkmarki\checkmarkr\checkmarkc\checkmark$\checkmark{\checkmark$\checkmark^\checkmark\\checkmarkc\checkmarki\checkmarkr\checkmarkc\checkmark$\checkmarkc\checkmark$\checkmark^\checkmark\\checkmarkc\checkmarki\checkmarkr\checkmarkc\checkmark$\checkmarko\checkmark$\checkmark^\checkmark\\checkmarkc\checkmarki\checkmarkr\checkmarkc\checkmark$\checkmarku\checkmark$\checkmark^\checkmark\\checkmarkc\checkmarki\checkmarkr\checkmarkc\checkmark$\checkmarkr\checkmark$\checkmark^\checkmark\\checkmarkc\checkmarki\checkmarkr\checkmarkc\checkmark$\checkmarkr\checkmark$\checkmark^\checkmark\\checkmarkc\checkmarki\checkmarkr\checkmarkc\checkmark$\checkmarko\checkmark$\checkmark^\checkmark\\checkmarkc\checkmarki\checkmarkr\checkmarkc\checkmark$\checkmarki\checkmark$\checkmark^\checkmark\\checkmarkc\checkmarki\checkmarkr\checkmarkc\checkmark$\checkmarke\checkmark$\checkmark^\checkmark\\checkmarkc\checkmarki\checkmarkr\checkmarkc\checkmark$\checkmark}\checkmark$\checkmark^\checkmark\\checkmarkc\checkmarki\checkmarkr\checkmarkc\checkmark$\checkmark}\checkmark$\checkmark^\checkmark\\checkmarkc\checkmarki\checkmarkr\checkmarkc\checkmark$\checkmark{\checkmark$\checkmark^\checkmark\\checkmarkc\checkmarki\checkmarkr\checkmarkc\checkmark$\checkmark2\checkmark$\checkmark^\checkmark\\checkmarkc\checkmarki\checkmarkr\checkmarkc\checkmark$\checkmark}\checkmark$\checkmark^\checkmark\\checkmarkc\checkmarki\checkmarkr\checkmarkc\checkmark$\checkmark \checkmark$\checkmark^\checkmark\\checkmarkc\checkmarki\checkmarkr\checkmarkc\checkmark$\checkmark=\checkmark$\checkmark^\checkmark\\checkmarkc\checkmarki\checkmarkr\checkmarkc\checkmark$\checkmark \checkmark$\checkmark^\checkmark\\checkmarkc\checkmarki\checkmarkr\checkmarkc\checkmark$\checkmark\\checkmark$\checkmark^\checkmark\\checkmarkc\checkmarki\checkmarkr\checkmarkc\checkmark$\checkmarkf\checkmark$\checkmark^\checkmark\\checkmarkc\checkmarki\checkmarkr\checkmarkc\checkmark$\checkmarkr\checkmark$\checkmark^\checkmark\\checkmarkc\checkmarki\checkmarkr\checkmarkc\checkmark$\checkmarka\checkmark$\checkmark^\checkmark\\checkmarkc\checkmarki\checkmarkr\checkmarkc\checkmark$\checkmarkc\checkmark$\checkmark^\checkmark\\checkmarkc\checkmarki\checkmarkr\checkmarkc\checkmark$\checkmark{\checkmark$\checkmark^\checkmark\\checkmarkc\checkmarki\checkmarkr\checkmarkc\checkmark$\checkmark2\checkmark$\checkmark^\checkmark\\checkmarkc\checkmarki\checkmarkr\checkmarkc\checkmark$\checkmark0\checkmark$\checkmark^\checkmark\\checkmarkc\checkmarki\checkmarkr\checkmarkc\checkmark$\checkmark}\checkmark$\checkmark^\checkmark\\checkmarkc\checkmarki\checkmarkr\checkmarkc\checkmark$\checkmark{\checkmark$\checkmark^\checkmark\\checkmarkc\checkmarki\checkmarkr\checkmarkc\checkmark$\checkmark0\checkmark$\checkmark^\checkmark\\checkmarkc\checkmarki\checkmarkr\checkmarkc\checkmark$\checkmark{\checkmark$\checkmark^\checkmark\\checkmarkc\checkmarki\checkmarkr\checkmarkc\checkmark$\checkmark,\checkmark$\checkmark^\checkmark\\checkmarkc\checkmarki\checkmarkr\checkmarkc\checkmark$\checkmark}\checkmark$\checkmark^\checkmark\\checkmarkc\checkmarki\checkmarkr\checkmarkc\checkmark$\checkmark1\checkmark$\checkmark^\checkmark\\checkmarkc\checkmarki\checkmarkr\checkmarkc\checkmark$\checkmark5\checkmark$\checkmark^\checkmark\\checkmarkc\checkmarki\checkmarkr\checkmarkc\checkmark$\checkmark0\checkmark$\checkmark^\checkmark\\checkmarkc\checkmarki\checkmarkr\checkmarkc\checkmark$\checkmark}\checkmark$\checkmark^\checkmark\\checkmarkc\checkmarki\checkmarkr\checkmarkc\checkmark$\checkmark \checkmark$\checkmark^\checkmark\\checkmarkc\checkmarki\checkmarkr\checkmarkc\checkmark$\checkmark+\checkmark$\checkmark^\checkmark\\checkmarkc\checkmarki\checkmarkr\checkmarkc\checkmark$\checkmark \checkmark$\checkmark^\checkmark\\checkmarkc\checkmarki\checkmarkr\checkmarkc\checkmark$\checkmark\\checkmark$\checkmark^\checkmark\\checkmarkc\checkmarki\checkmarkr\checkmarkc\checkmark$\checkmarkf\checkmark$\checkmark^\checkmark\\checkmarkc\checkmarki\checkmarkr\checkmarkc\checkmark$\checkmarkr\checkmark$\checkmark^\checkmark\\checkmarkc\checkmarki\checkmarkr\checkmarkc\checkmark$\checkmarka\checkmark$\checkmark^\checkmark\\checkmarkc\checkmarki\checkmarkr\checkmarkc\checkmark$\checkmarkc\checkmark$\checkmark^\checkmark\\checkmarkc\checkmarki\checkmarkr\checkmarkc\checkmark$\checkmark{\checkmark$\checkmark^\checkmark\\checkmarkc\checkmarki\checkmarkr\checkmarkc\checkmark$\checkmark3\checkmark$\checkmark^\checkmark\\checkmarkc\checkmarki\checkmarkr\checkmarkc\checkmark$\checkmark3\checkmark$\checkmark^\checkmark\\checkmarkc\checkmarki\checkmarkr\checkmarkc\checkmark$\checkmark3\checkmark$\checkmark^\checkmark\\checkmarkc\checkmarki\checkmarkr\checkmarkc\checkmark$\checkmark}\checkmark$\checkmark^\checkmark\\checkmarkc\checkmarki\checkmarkr\checkmarkc\checkmark$\checkmark{\checkmark$\checkmark^\checkmark\\checkmarkc\checkmarki\checkmarkr\checkmarkc\checkmark$\checkmark2\checkmark$\checkmark^\checkmark\\checkmarkc\checkmarki\checkmarkr\checkmarkc\checkmark$\checkmark}\checkmark$\checkmark^\checkmark\\checkmarkc\checkmarki\checkmarkr\checkmarkc\checkmark$\checkmark \checkmark$\checkmark^\checkmark\\checkmarkc\checkmarki\checkmarkr\checkmarkc\checkmark$\checkmark=\checkmark$\checkmark^\checkmark\\checkmarkc\checkmarki\checkmarkr\checkmarkc\checkmark$\checkmark \checkmark$\checkmark^\checkmark\\checkmarkc\checkmarki\checkmarkr\checkmarkc\checkmark$\checkmark1\checkmark$\checkmark^\checkmark\\checkmarkc\checkmarki\checkmarkr\checkmarkc\checkmark$\checkmark3\checkmark$\checkmark^\checkmark\\checkmarkc\checkmarki\checkmarkr\checkmarkc\checkmark$\checkmark3\checkmark$\checkmark^\checkmark\\checkmarkc\checkmarki\checkmarkr\checkmarkc\checkmark$\checkmark{\checkmark$\checkmark^\checkmark\\checkmarkc\checkmarki\checkmarkr\checkmarkc\checkmark$\checkmark,\checkmark$\checkmark^\checkmark\\checkmarkc\checkmarki\checkmarkr\checkmarkc\checkmark$\checkmark}\checkmark$\checkmark^\checkmark\\checkmarkc\checkmarki\checkmarkr\checkmarkc\checkmark$\checkmark3\checkmark$\checkmark^\checkmark\\checkmarkc\checkmarki\checkmarkr\checkmarkc\checkmark$\checkmark \checkmark$\checkmark^\checkmark\\checkmarkc\checkmarki\checkmarkr\checkmarkc\checkmark$\checkmark+\checkmark$\checkmark^\checkmark\\checkmarkc\checkmarki\checkmarkr\checkmarkc\checkmark$\checkmark \checkmark$\checkmark^\checkmark\\checkmarkc\checkmarki\checkmarkr\checkmarkc\checkmark$\checkmark1\checkmark$\checkmark^\checkmark\\checkmarkc\checkmarki\checkmarkr\checkmarkc\checkmark$\checkmark6\checkmark$\checkmark^\checkmark\\checkmarkc\checkmarki\checkmarkr\checkmarkc\checkmark$\checkmark6\checkmark$\checkmark^\checkmark\\checkmarkc\checkmarki\checkmarkr\checkmarkc\checkmark$\checkmark{\checkmark$\checkmark^\checkmark\\checkmarkc\checkmarki\checkmarkr\checkmarkc\checkmark$\checkmark,\checkmark$\checkmark^\checkmark\\checkmarkc\checkmarki\checkmarkr\checkmarkc\checkmark$\checkmark}\checkmark$\checkmark^\checkmark\\checkmarkc\checkmarki\checkmarkr\checkmarkc\checkmark$\checkmark5\checkmark$\checkmark^\checkmark\\checkmarkc\checkmarki\checkmarkr\checkmarkc\checkmark$\checkmark \checkmark$\checkmark^\checkmark\\checkmarkc\checkmarki\checkmarkr\checkmarkc\checkmark$\checkmark=\checkmark$\checkmark^\checkmark\\checkmarkc\checkmarki\checkmarkr\checkmarkc\checkmark$\checkmark \checkmark$\checkmark^\checkmark\\checkmarkc\checkmarki\checkmarkr\checkmarkc\checkmark$\checkmark\\checkmark$\checkmark^\checkmark\\checkmarkc\checkmarki\checkmarkr\checkmarkc\checkmark$\checkmarkb\checkmark$\checkmark^\checkmark\\checkmarkc\checkmarki\checkmarkr\checkmarkc\checkmark$\checkmarko\checkmark$\checkmark^\checkmark\\checkmarkc\checkmarki\checkmarkr\checkmarkc\checkmark$\checkmarkx\checkmark$\checkmark^\checkmark\\checkmarkc\checkmarki\checkmarkr\checkmarkc\checkmark$\checkmarke\checkmark$\checkmark^\checkmark\\checkmarkc\checkmarki\checkmarkr\checkmarkc\checkmark$\checkmarkd\checkmark$\checkmark^\checkmark\\checkmarkc\checkmarki\checkmarkr\checkmarkc\checkmark$\checkmark{\checkmark$\checkmark^\checkmark\\checkmarkc\checkmarki\checkmarkr\checkmarkc\checkmark$\checkmark2\checkmark$\checkmark^\checkmark\\checkmarkc\checkmarki\checkmarkr\checkmarkc\checkmark$\checkmark9\checkmark$\checkmark^\checkmark\\checkmarkc\checkmarki\checkmarkr\checkmarkc\checkmark$\checkmark9\checkmark$\checkmark^\checkmark\\checkmarkc\checkmarki\checkmarkr\checkmarkc\checkmark$\checkmark{\checkmark$\checkmark^\checkmark\\checkmarkc\checkmarki\checkmarkr\checkmarkc\checkmark$\checkmark,\checkmark$\checkmark^\checkmark\\checkmarkc\checkmarki\checkmarkr\checkmarkc\checkmark$\checkmark}\checkmark$\checkmark^\checkmark\\checkmarkc\checkmarki\checkmarkr\checkmarkc\checkmark$\checkmark8\checkmark$\checkmark^\checkmark\\checkmarkc\checkmarki\checkmarkr\checkmarkc\checkmark$\checkmark \checkmark$\checkmark^\checkmark\\checkmarkc\checkmarki\checkmarkr\checkmarkc\checkmark$\checkmark\\checkmark$\checkmark^\checkmark\\checkmarkc\checkmarki\checkmarkr\checkmarkc\checkmark$\checkmarka\checkmark$\checkmark^\checkmark\\checkmarkc\checkmarki\checkmarkr\checkmarkc\checkmark$\checkmarkp\checkmark$\checkmark^\checkmark\\checkmarkc\checkmarki\checkmarkr\checkmarkc\checkmark$\checkmarkp\checkmark$\checkmark^\checkmark\\checkmarkc\checkmarki\checkmarkr\checkmarkc\checkmark$\checkmarkr\checkmark$\checkmark^\checkmark\\checkmarkc\checkmarki\checkmarkr\checkmarkc\checkmark$\checkmarko\checkmark$\checkmark^\checkmark\\checkmarkc\checkmarki\checkmarkr\checkmarkc\checkmark$\checkmarkx\checkmark$\checkmark^\checkmark\\checkmarkc\checkmarki\checkmarkr\checkmarkc\checkmark$\checkmark \checkmark$\checkmark^\checkmark\\checkmarkc\checkmarki\checkmarkr\checkmarkc\checkmark$\checkmark3\checkmark$\checkmark^\checkmark\\checkmarkc\checkmarki\checkmarkr\checkmarkc\checkmark$\checkmark0\checkmark$\checkmark^\checkmark\\checkmarkc\checkmarki\checkmarkr\checkmarkc\checkmark$\checkmark0\checkmark$\checkmark^\checkmark\\checkmarkc\checkmarki\checkmarkr\checkmarkc\checkmark$\checkmark \checkmark$\checkmark^\checkmark\\checkmarkc\checkmarki\checkmarkr\checkmarkc\checkmark$\checkmark\\checkmark$\checkmark^\checkmark\\checkmarkc\checkmarki\checkmarkr\checkmarkc\checkmark$\checkmarkt\checkmark$\checkmark^\checkmark\\checkmarkc\checkmarki\checkmarkr\checkmarkc\checkmark$\checkmarke\checkmark$\checkmark^\checkmark\\checkmarkc\checkmarki\checkmarkr\checkmarkc\checkmark$\checkmarkx\checkmark$\checkmark^\checkmark\\checkmarkc\checkmarki\checkmarkr\checkmarkc\checkmark$\checkmarkt\checkmark$\checkmark^\checkmark\\checkmarkc\checkmarki\checkmarkr\checkmarkc\checkmark$\checkmark{\checkmark$\checkmark^\checkmark\\checkmarkc\checkmarki\checkmarkr\checkmarkc\checkmark$\checkmark \checkmark$\checkmark^\checkmark\\checkmarkc\checkmarki\checkmarkr\checkmarkc\checkmark$\checkmarkN\checkmark$\checkmark^\checkmark\\checkmarkc\checkmarki\checkmarkr\checkmarkc\checkmark$\checkmark}\checkmark$\checkmark^\checkmark\\checkmarkc\checkmarki\checkmarkr\checkmarkc\checkmark$\checkmark}\checkmark$\checkmark^\checkmark\\checkmarkc\checkmarki\checkmarkr\checkmarkc\checkmark$\checkmark
\checkmark$\checkmark^\checkmark\\checkmarkc\checkmarki\checkmarkr\checkmarkc\checkmark$\checkmark \checkmark$\checkmark^\checkmark\\checkmarkc\checkmarki\checkmarkr\checkmarkc\checkmark$\checkmark \checkmark$\checkmark^\checkmark\\checkmarkc\checkmarki\checkmarkr\checkmarkc\checkmark$\checkmark \checkmark$\checkmark^\checkmark\\checkmarkc\checkmarki\checkmarkr\checkmarkc\checkmark$\checkmark \checkmark$\checkmark^\checkmark\\checkmarkc\checkmarki\checkmarkr\checkmarkc\checkmark$\checkmark\\checkmark$\checkmark^\checkmark\\checkmarkc\checkmarki\checkmarkr\checkmarkc\checkmark$\checkmarkl\checkmark$\checkmark^\checkmark\\checkmarkc\checkmarki\checkmarkr\checkmarkc\checkmark$\checkmarka\checkmark$\checkmark^\checkmark\\checkmarkc\checkmarki\checkmarkr\checkmarkc\checkmark$\checkmarkb\checkmark$\checkmark^\checkmark\\checkmarkc\checkmarki\checkmarkr\checkmarkc\checkmark$\checkmarke\checkmark$\checkmark^\checkmark\\checkmarkc\checkmarki\checkmarkr\checkmarkc\checkmark$\checkmarkl\checkmark$\checkmark^\checkmark\\checkmarkc\checkmarki\checkmarkr\checkmarkc\checkmark$\checkmark{\checkmark$\checkmark^\checkmark\\checkmarkc\checkmarki\checkmarkr\checkmarkc\checkmark$\checkmarke\checkmark$\checkmark^\checkmark\\checkmarkc\checkmarki\checkmarkr\checkmarkc\checkmark$\checkmarkq\checkmark$\checkmark^\checkmark\\checkmarkc\checkmarki\checkmarkr\checkmarkc\checkmark$\checkmark:\checkmark$\checkmark^\checkmark\\checkmarkc\checkmarki\checkmarkr\checkmarkc\checkmark$\checkmarkr\checkmark$\checkmark^\checkmark\\checkmarkc\checkmarki\checkmarkr\checkmarkc\checkmark$\checkmarke\checkmark$\checkmark^\checkmark\\checkmarkc\checkmarki\checkmarkr\checkmarkc\checkmark$\checkmarka\checkmark$\checkmark^\checkmark\\checkmarkc\checkmarki\checkmarkr\checkmarkc\checkmark$\checkmarkc\checkmark$\checkmark^\checkmark\\checkmarkc\checkmarki\checkmarkr\checkmarkc\checkmark$\checkmarkt\checkmark$\checkmark^\checkmark\\checkmarkc\checkmarki\checkmarkr\checkmarkc\checkmark$\checkmarki\checkmark$\checkmark^\checkmark\\checkmarkc\checkmarki\checkmarkr\checkmarkc\checkmark$\checkmarko\checkmark$\checkmark^\checkmark\\checkmarkc\checkmarki\checkmarkr\checkmarkc\checkmark$\checkmarkn\checkmark$\checkmark^\checkmark\\checkmarkc\checkmarki\checkmarkr\checkmarkc\checkmark$\checkmark_\checkmark$\checkmark^\checkmark\\checkmarkc\checkmarki\checkmarkr\checkmarkc\checkmark$\checkmarkp\checkmark$\checkmark^\checkmark\\checkmarkc\checkmarki\checkmarkr\checkmarkc\checkmark$\checkmarka\checkmark$\checkmark^\checkmark\\checkmarkc\checkmarki\checkmarkr\checkmarkc\checkmark$\checkmarkl\checkmark$\checkmark^\checkmark\\checkmarkc\checkmarki\checkmarkr\checkmarkc\checkmark$\checkmarki\checkmark$\checkmark^\checkmark\\checkmarkc\checkmarki\checkmarkr\checkmarkc\checkmark$\checkmarke\checkmark$\checkmark^\checkmark\\checkmarkc\checkmarki\checkmarkr\checkmarkc\checkmark$\checkmarkr\checkmark$\checkmark^\checkmark\\checkmarkc\checkmarki\checkmarkr\checkmarkc\checkmark$\checkmark}\checkmark$\checkmark^\checkmark\\checkmarkc\checkmarki\checkmarkr\checkmarkc\checkmark$\checkmark
\checkmark$\checkmark^\checkmark\\checkmarkc\checkmarki\checkmarkr\checkmarkc\checkmark$\checkmark\\checkmark$\checkmark^\checkmark\\checkmarkc\checkmarki\checkmarkr\checkmarkc\checkmark$\checkmarke\checkmark$\checkmark^\checkmark\\checkmarkc\checkmarki\checkmarkr\checkmarkc\checkmark$\checkmarkn\checkmark$\checkmark^\checkmark\\checkmarkc\checkmarki\checkmarkr\checkmarkc\checkmark$\checkmarkd\checkmark$\checkmark^\checkmark\\checkmarkc\checkmarki\checkmarkr\checkmarkc\checkmark$\checkmark{\checkmark$\checkmark^\checkmark\\checkmarkc\checkmarki\checkmarkr\checkmarkc\checkmark$\checkmarke\checkmark$\checkmark^\checkmark\\checkmarkc\checkmarki\checkmarkr\checkmarkc\checkmark$\checkmarkq\checkmark$\checkmark^\checkmark\\checkmarkc\checkmarki\checkmarkr\checkmarkc\checkmark$\checkmarku\checkmark$\checkmark^\checkmark\\checkmarkc\checkmarki\checkmarkr\checkmarkc\checkmark$\checkmarka\checkmark$\checkmark^\checkmark\\checkmarkc\checkmarki\checkmarkr\checkmarkc\checkmark$\checkmarkt\checkmark$\checkmark^\checkmark\\checkmarkc\checkmarki\checkmarkr\checkmarkc\checkmark$\checkmarki\checkmark$\checkmark^\checkmark\\checkmarkc\checkmarki\checkmarkr\checkmarkc\checkmark$\checkmarko\checkmark$\checkmark^\checkmark\\checkmarkc\checkmarki\checkmarkr\checkmarkc\checkmark$\checkmarkn\checkmark$\checkmark^\checkmark\\checkmarkc\checkmarki\checkmarkr\checkmarkc\checkmark$\checkmark}\checkmark$\checkmark^\checkmark\\checkmarkc\checkmarki\checkmarkr\checkmarkc\checkmark$\checkmark
\checkmark$\checkmark^\checkmark\\checkmarkc\checkmarki\checkmarkr\checkmarkc\checkmark$\checkmark
\checkmark$\checkmark^\checkmark\\checkmarkc\checkmarki\checkmarkr\checkmarkc\checkmark$\checkmark\\checkmark$\checkmark^\checkmark\\checkmarkc\checkmarki\checkmarkr\checkmarkc\checkmark$\checkmarks\checkmark$\checkmark^\checkmark\\checkmarkc\checkmarki\checkmarkr\checkmarkc\checkmark$\checkmarku\checkmark$\checkmark^\checkmark\\checkmarkc\checkmarki\checkmarkr\checkmarkc\checkmark$\checkmarkb\checkmark$\checkmark^\checkmark\\checkmarkc\checkmarki\checkmarkr\checkmarkc\checkmark$\checkmarks\checkmark$\checkmark^\checkmark\\checkmarkc\checkmarki\checkmarkr\checkmarkc\checkmark$\checkmarke\checkmark$\checkmark^\checkmark\\checkmarkc\checkmarki\checkmarkr\checkmarkc\checkmark$\checkmarkc\checkmark$\checkmark^\checkmark\\checkmarkc\checkmarki\checkmarkr\checkmarkc\checkmark$\checkmarkt\checkmark$\checkmark^\checkmark\\checkmarkc\checkmarki\checkmarkr\checkmarkc\checkmark$\checkmarki\checkmark$\checkmark^\checkmark\\checkmarkc\checkmarki\checkmarkr\checkmarkc\checkmark$\checkmarko\checkmark$\checkmark^\checkmark\\checkmarkc\checkmarki\checkmarkr\checkmarkc\checkmark$\checkmarkn\checkmark$\checkmark^\checkmark\\checkmarkc\checkmarki\checkmarkr\checkmarkc\checkmark$\checkmark{\checkmark$\checkmark^\checkmark\\checkmarkc\checkmarki\checkmarkr\checkmarkc\checkmark$\checkmarkD\checkmark$\checkmark^\checkmark\\checkmarkc\checkmarki\checkmarkr\checkmarkc\checkmark$\checkmarki\checkmark$\checkmark^\checkmark\\checkmarkc\checkmarki\checkmarkr\checkmarkc\checkmark$\checkmarkm\checkmark$\checkmark^\checkmark\\checkmarkc\checkmarki\checkmarkr\checkmarkc\checkmark$\checkmarke\checkmark$\checkmark^\checkmark\\checkmarkc\checkmarki\checkmarkr\checkmarkc\checkmark$\checkmarkn\checkmark$\checkmark^\checkmark\\checkmarkc\checkmarki\checkmarkr\checkmarkc\checkmark$\checkmarks\checkmark$\checkmark^\checkmark\\checkmarkc\checkmarki\checkmarkr\checkmarkc\checkmark$\checkmarki\checkmark$\checkmark^\checkmark\\checkmarkc\checkmarki\checkmarkr\checkmarkc\checkmark$\checkmarko\checkmark$\checkmark^\checkmark\\checkmarkc\checkmarki\checkmarkr\checkmarkc\checkmark$\checkmarkn\checkmark$\checkmark^\checkmark\\checkmarkc\checkmarki\checkmarkr\checkmarkc\checkmark$\checkmarkn\checkmark$\checkmark^\checkmark\\checkmarkc\checkmarki\checkmarkr\checkmarkc\checkmark$\checkmarke\checkmark$\checkmark^\checkmark\\checkmarkc\checkmarki\checkmarkr\checkmarkc\checkmark$\checkmarkm\checkmark$\checkmark^\checkmark\\checkmarkc\checkmarki\checkmarkr\checkmarkc\checkmark$\checkmarke\checkmark$\checkmark^\checkmark\\checkmarkc\checkmarki\checkmarkr\checkmarkc\checkmark$\checkmarkn\checkmark$\checkmark^\checkmark\\checkmarkc\checkmarki\checkmarkr\checkmarkc\checkmark$\checkmarkt\checkmark$\checkmark^\checkmark\\checkmarkc\checkmarki\checkmarkr\checkmarkc\checkmark$\checkmark \checkmark$\checkmark^\checkmark\\checkmarkc\checkmarki\checkmarkr\checkmarkc\checkmark$\checkmarke\checkmark$\checkmark^\checkmark\\checkmarkc\checkmarki\checkmarkr\checkmarkc\checkmark$\checkmarkn\checkmark$\checkmark^\checkmark\\checkmarkc\checkmarki\checkmarkr\checkmarkc\checkmark$\checkmark \checkmark$\checkmark^\checkmark\\checkmarkc\checkmarki\checkmarkr\checkmarkc\checkmark$\checkmarks\checkmark$\checkmark^\checkmark\\checkmarkc\checkmarki\checkmarkr\checkmarkc\checkmark$\checkmarkt\checkmark$\checkmark^\checkmark\\checkmarkc\checkmarki\checkmarkr\checkmarkc\checkmark$\checkmarka\checkmark$\checkmark^\checkmark\\checkmarkc\checkmarki\checkmarkr\checkmarkc\checkmark$\checkmarkt\checkmark$\checkmark^\checkmark\\checkmarkc\checkmarki\checkmarkr\checkmarkc\checkmark$\checkmarki\checkmark$\checkmark^\checkmark\\checkmarkc\checkmarki\checkmarkr\checkmarkc\checkmark$\checkmarkq\checkmark$\checkmark^\checkmark\\checkmarkc\checkmarki\checkmarkr\checkmarkc\checkmark$\checkmarku\checkmark$\checkmark^\checkmark\\checkmarkc\checkmarki\checkmarkr\checkmarkc\checkmark$\checkmarke\checkmark$\checkmark^\checkmark\\checkmarkc\checkmarki\checkmarkr\checkmarkc\checkmark$\checkmark \checkmark$\checkmark^\checkmark\\checkmarkc\checkmarki\checkmarkr\checkmarkc\checkmark$\checkmarkd\checkmark$\checkmark^\checkmark\\checkmarkc\checkmarki\checkmarkr\checkmarkc\checkmark$\checkmarke\checkmark$\checkmark^\checkmark\\checkmarkc\checkmarki\checkmarkr\checkmarkc\checkmark$\checkmark \checkmark$\checkmark^\checkmark\\checkmarkc\checkmarki\checkmarkr\checkmarkc\checkmark$\checkmarkl\checkmark$\checkmark^\checkmark\\checkmarkc\checkmarki\checkmarkr\checkmarkc\checkmark$\checkmark'\checkmark$\checkmark^\checkmark\\checkmarkc\checkmarki\checkmarkr\checkmarkc\checkmark$\checkmarka\checkmark$\checkmark^\checkmark\\checkmarkc\checkmarki\checkmarkr\checkmarkc\checkmark$\checkmarkr\checkmark$\checkmark^\checkmark\\checkmarkc\checkmarki\checkmarkr\checkmarkc\checkmark$\checkmarkb\checkmark$\checkmark^\checkmark\\checkmarkc\checkmarki\checkmarkr\checkmarkc\checkmark$\checkmarkr\checkmark$\checkmark^\checkmark\\checkmarkc\checkmarki\checkmarkr\checkmarkc\checkmark$\checkmarke\checkmark$\checkmark^\checkmark\\checkmarkc\checkmarki\checkmarkr\checkmarkc\checkmark$\checkmark \checkmark$\checkmark^\checkmark\\checkmarkc\checkmarki\checkmarkr\checkmarkc\checkmark$\checkmark(\checkmark$\checkmark^\checkmark\\checkmarkc\checkmarki\checkmarkr\checkmarkc\checkmark$\checkmarkM\checkmark$\checkmark^\checkmark\\checkmarkc\checkmarki\checkmarkr\checkmarkc\checkmark$\checkmarké\checkmark$\checkmark^\checkmark\\checkmarkc\checkmarki\checkmarkr\checkmarkc\checkmark$\checkmarkt\checkmark$\checkmark^\checkmark\\checkmarkc\checkmarki\checkmarkr\checkmarkc\checkmark$\checkmarkh\checkmark$\checkmark^\checkmark\\checkmarkc\checkmarki\checkmarkr\checkmarkc\checkmark$\checkmarko\checkmark$\checkmark^\checkmark\\checkmarkc\checkmarki\checkmarkr\checkmarkc\checkmark$\checkmarkd\checkmark$\checkmark^\checkmark\\checkmarkc\checkmarki\checkmarkr\checkmarkc\checkmark$\checkmarke\checkmark$\checkmark^\checkmark\\checkmarkc\checkmarki\checkmarkr\checkmarkc\checkmark$\checkmark \checkmark$\checkmark^\checkmark\\checkmarkc\checkmarki\checkmarkr\checkmarkc\checkmark$\checkmarkÉ\checkmark$\checkmark^\checkmark\\checkmarkc\checkmarki\checkmarkr\checkmarkc\checkmark$\checkmarkn\checkmark$\checkmark^\checkmark\\checkmarkc\checkmarki\checkmarkr\checkmarkc\checkmark$\checkmarke\checkmark$\checkmark^\checkmark\\checkmarkc\checkmarki\checkmarkr\checkmarkc\checkmark$\checkmarkr\checkmark$\checkmark^\checkmark\\checkmarkc\checkmarki\checkmarkr\checkmarkc\checkmark$\checkmarkg\checkmark$\checkmark^\checkmark\\checkmarkc\checkmarki\checkmarkr\checkmarkc\checkmark$\checkmarké\checkmark$\checkmark^\checkmark\\checkmarkc\checkmarki\checkmarkr\checkmarkc\checkmark$\checkmarkt\checkmark$\checkmark^\checkmark\\checkmarkc\checkmarki\checkmarkr\checkmarkc\checkmark$\checkmarki\checkmark$\checkmark^\checkmark\\checkmarkc\checkmarki\checkmarkr\checkmarkc\checkmark$\checkmarkq\checkmark$\checkmark^\checkmark\\checkmarkc\checkmarki\checkmarkr\checkmarkc\checkmark$\checkmarku\checkmark$\checkmark^\checkmark\\checkmarkc\checkmarki\checkmarkr\checkmarkc\checkmark$\checkmarke\checkmark$\checkmark^\checkmark\\checkmarkc\checkmarki\checkmarkr\checkmarkc\checkmark$\checkmark)\checkmark$\checkmark^\checkmark\\checkmarkc\checkmarki\checkmarkr\checkmarkc\checkmark$\checkmark}\checkmark$\checkmark^\checkmark\\checkmarkc\checkmarki\checkmarkr\checkmarkc\checkmark$\checkmark
\checkmark$\checkmark^\checkmark\\checkmarkc\checkmarki\checkmarkr\checkmarkc\checkmark$\checkmarkL\checkmark$\checkmark^\checkmark\\checkmarkc\checkmarki\checkmarkr\checkmarkc\checkmark$\checkmarke\checkmark$\checkmark^\checkmark\\checkmarkc\checkmarki\checkmarkr\checkmarkc\checkmark$\checkmark \checkmark$\checkmark^\checkmark\\checkmarkc\checkmarki\checkmarkr\checkmarkc\checkmark$\checkmarkm\checkmark$\checkmark^\checkmark\\checkmarkc\checkmarki\checkmarkr\checkmarkc\checkmark$\checkmarko\checkmark$\checkmark^\checkmark\\checkmarkc\checkmarki\checkmarkr\checkmarkc\checkmark$\checkmarkm\checkmark$\checkmark^\checkmark\\checkmarkc\checkmarki\checkmarkr\checkmarkc\checkmark$\checkmarke\checkmark$\checkmark^\checkmark\\checkmarkc\checkmarki\checkmarkr\checkmarkc\checkmark$\checkmarkn\checkmark$\checkmark^\checkmark\\checkmarkc\checkmarki\checkmarkr\checkmarkc\checkmark$\checkmarkt\checkmark$\checkmark^\checkmark\\checkmarkc\checkmarki\checkmarkr\checkmarkc\checkmark$\checkmark \checkmark$\checkmark^\checkmark\\checkmarkc\checkmarki\checkmarkr\checkmarkc\checkmark$\checkmarkf\checkmark$\checkmark^\checkmark\\checkmarkc\checkmarki\checkmarkr\checkmarkc\checkmark$\checkmarkl\checkmark$\checkmark^\checkmark\\checkmarkc\checkmarki\checkmarkr\checkmarkc\checkmark$\checkmarké\checkmark$\checkmark^\checkmark\\checkmarkc\checkmarki\checkmarkr\checkmarkc\checkmark$\checkmarkc\checkmark$\checkmark^\checkmark\\checkmarkc\checkmarki\checkmarkr\checkmarkc\checkmark$\checkmarkh\checkmark$\checkmark^\checkmark\\checkmarkc\checkmarki\checkmarkr\checkmarkc\checkmark$\checkmarki\checkmark$\checkmark^\checkmark\\checkmarkc\checkmarki\checkmarkr\checkmarkc\checkmark$\checkmarks\checkmark$\checkmark^\checkmark\\checkmarkc\checkmarki\checkmarkr\checkmarkc\checkmark$\checkmarks\checkmark$\checkmark^\checkmark\\checkmarkc\checkmarki\checkmarkr\checkmarkc\checkmark$\checkmarka\checkmark$\checkmark^\checkmark\\checkmarkc\checkmarki\checkmarkr\checkmarkc\checkmark$\checkmarkn\checkmark$\checkmark^\checkmark\\checkmarkc\checkmarki\checkmarkr\checkmarkc\checkmark$\checkmarkt\checkmark$\checkmark^\checkmark\\checkmarkc\checkmarki\checkmarkr\checkmarkc\checkmark$\checkmark \checkmark$\checkmark^\checkmark\\checkmarkc\checkmarki\checkmarkr\checkmarkc\checkmark$\checkmarkm\checkmark$\checkmark^\checkmark\\checkmarkc\checkmarki\checkmarkr\checkmarkc\checkmark$\checkmarka\checkmark$\checkmark^\checkmark\\checkmarkc\checkmarki\checkmarkr\checkmarkc\checkmark$\checkmarkx\checkmark$\checkmark^\checkmark\\checkmarkc\checkmarki\checkmarkr\checkmarkc\checkmark$\checkmarki\checkmark$\checkmark^\checkmark\\checkmarkc\checkmarki\checkmarkr\checkmarkc\checkmark$\checkmarkm\checkmark$\checkmark^\checkmark\\checkmarkc\checkmarki\checkmarkr\checkmarkc\checkmark$\checkmarka\checkmark$\checkmark^\checkmark\\checkmarkc\checkmarki\checkmarkr\checkmarkc\checkmark$\checkmarkl\checkmark$\checkmark^\checkmark\\checkmarkc\checkmarki\checkmarkr\checkmarkc\checkmark$\checkmark \checkmark$\checkmark^\checkmark\\checkmarkc\checkmarki\checkmarkr\checkmarkc\checkmark$\checkmark$\checkmark$\checkmark^\checkmark\\checkmarkc\checkmarki\checkmarkr\checkmarkc\checkmark$\checkmarkM\checkmark$\checkmark^\checkmark\\checkmarkc\checkmarki\checkmarkr\checkmarkc\checkmark$\checkmark_\checkmark$\checkmark^\checkmark\\checkmarkc\checkmarki\checkmarkr\checkmarkc\checkmark$\checkmark{\checkmark$\checkmark^\checkmark\\checkmarkc\checkmarki\checkmarkr\checkmarkc\checkmark$\checkmarkf\checkmark$\checkmark^\checkmark\\checkmarkc\checkmarki\checkmarkr\checkmarkc\checkmark$\checkmark\\checkmark$\checkmark^\checkmark\\checkmarkc\checkmarki\checkmarkr\checkmarkc\checkmark$\checkmark_\checkmark$\checkmark^\checkmark\\checkmarkc\checkmarki\checkmarkr\checkmarkc\checkmark$\checkmarkm\checkmark$\checkmark^\checkmark\\checkmarkc\checkmarki\checkmarkr\checkmarkc\checkmark$\checkmarka\checkmark$\checkmark^\checkmark\\checkmarkc\checkmarki\checkmarkr\checkmarkc\checkmark$\checkmarkx\checkmark$\checkmark^\checkmark\\checkmarkc\checkmarki\checkmarkr\checkmarkc\checkmark$\checkmark}\checkmark$\checkmark^\checkmark\\checkmarkc\checkmarki\checkmarkr\checkmarkc\checkmark$\checkmark$\checkmark$\checkmark^\checkmark\\checkmarkc\checkmarki\checkmarkr\checkmarkc\checkmark$\checkmark \checkmark$\checkmark^\checkmark\\checkmarkc\checkmarki\checkmarkr\checkmarkc\checkmark$\checkmarks\checkmark$\checkmark^\checkmark\\checkmarkc\checkmarki\checkmarkr\checkmarkc\checkmark$\checkmark'\checkmark$\checkmark^\checkmark\\checkmarkc\checkmarki\checkmarkr\checkmarkc\checkmark$\checkmarké\checkmark$\checkmark^\checkmark\\checkmarkc\checkmarki\checkmarkr\checkmarkc\checkmark$\checkmarkt\checkmark$\checkmark^\checkmark\\checkmarkc\checkmarki\checkmarkr\checkmarkc\checkmark$\checkmarka\checkmark$\checkmark^\checkmark\\checkmarkc\checkmarki\checkmarkr\checkmarkc\checkmark$\checkmarkb\checkmark$\checkmark^\checkmark\\checkmarkc\checkmarki\checkmarkr\checkmarkc\checkmark$\checkmarkl\checkmark$\checkmark^\checkmark\\checkmarkc\checkmarki\checkmarkr\checkmarkc\checkmark$\checkmarki\checkmark$\checkmark^\checkmark\\checkmarkc\checkmarki\checkmarkr\checkmarkc\checkmark$\checkmarkt\checkmark$\checkmark^\checkmark\\checkmarkc\checkmarki\checkmarkr\checkmarkc\checkmark$\checkmark \checkmark$\checkmark^\checkmark\\checkmarkc\checkmarki\checkmarkr\checkmarkc\checkmark$\checkmarkà\checkmark$\checkmark^\checkmark\\checkmarkc\checkmarki\checkmarkr\checkmarkc\checkmark$\checkmark \checkmark$\checkmark^\checkmark\\checkmarkc\checkmarki\checkmarkr\checkmarkc\checkmark$\checkmarkm\checkmark$\checkmark^\checkmark\\checkmarkc\checkmarki\checkmarkr\checkmarkc\checkmark$\checkmarki\checkmark$\checkmark^\checkmark\\checkmarkc\checkmarki\checkmarkr\checkmarkc\checkmark$\checkmark-\checkmark$\checkmark^\checkmark\\checkmarkc\checkmarki\checkmarkr\checkmarkc\checkmark$\checkmarkp\checkmark$\checkmark^\checkmark\\checkmarkc\checkmarki\checkmarkr\checkmarkc\checkmark$\checkmarko\checkmark$\checkmark^\checkmark\\checkmarkc\checkmarki\checkmarkr\checkmarkc\checkmark$\checkmarkr\checkmark$\checkmark^\checkmark\\checkmarkc\checkmarki\checkmarkr\checkmarkc\checkmark$\checkmarkt\checkmark$\checkmark^\checkmark\\checkmarkc\checkmarki\checkmarkr\checkmarkc\checkmark$\checkmarké\checkmark$\checkmark^\checkmark\\checkmarkc\checkmarki\checkmarkr\checkmarkc\checkmark$\checkmarke\checkmark$\checkmark^\checkmark\\checkmarkc\checkmarki\checkmarkr\checkmarkc\checkmark$\checkmark.\checkmark$\checkmark^\checkmark\\checkmarkc\checkmarki\checkmarkr\checkmarkc\checkmark$\checkmark \checkmark$\checkmark^\checkmark\\checkmarkc\checkmarki\checkmarkr\checkmarkc\checkmark$\checkmarkP\checkmark$\checkmark^\checkmark\\checkmarkc\checkmarki\checkmarkr\checkmarkc\checkmark$\checkmarko\checkmark$\checkmark^\checkmark\\checkmarkc\checkmarki\checkmarkr\checkmarkc\checkmark$\checkmarku\checkmark$\checkmark^\checkmark\\checkmarkc\checkmarki\checkmarkr\checkmarkc\checkmark$\checkmarkr\checkmark$\checkmark^\checkmark\\checkmarkc\checkmarki\checkmarkr\checkmarkc\checkmark$\checkmark \checkmark$\checkmark^\checkmark\\checkmarkc\checkmarki\checkmarkr\checkmarkc\checkmark$\checkmarkn\checkmark$\checkmark^\checkmark\\checkmarkc\checkmarki\checkmarkr\checkmarkc\checkmark$\checkmarko\checkmark$\checkmark^\checkmark\\checkmarkc\checkmarki\checkmarkr\checkmarkc\checkmark$\checkmarkt\checkmark$\checkmark^\checkmark\\checkmarkc\checkmarki\checkmarkr\checkmarkc\checkmark$\checkmarkr\checkmark$\checkmark^\checkmark\\checkmarkc\checkmarki\checkmarkr\checkmarkc\checkmark$\checkmarke\checkmark$\checkmark^\checkmark\\checkmarkc\checkmarki\checkmarkr\checkmarkc\checkmark$\checkmark \checkmark$\checkmark^\checkmark\\checkmarkc\checkmarki\checkmarkr\checkmarkc\checkmark$\checkmarkm\checkmark$\checkmark^\checkmark\\checkmarkc\checkmarki\checkmarkr\checkmarkc\checkmark$\checkmarko\checkmark$\checkmark^\checkmark\\checkmarkc\checkmarki\checkmarkr\checkmarkc\checkmark$\checkmarkd\checkmark$\checkmark^\checkmark\\checkmarkc\checkmarki\checkmarkr\checkmarkc\checkmark$\checkmarkè\checkmark$\checkmark^\checkmark\\checkmarkc\checkmarki\checkmarkr\checkmarkc\checkmark$\checkmarkl\checkmark$\checkmark^\checkmark\\checkmarkc\checkmarki\checkmarkr\checkmarkc\checkmark$\checkmarke\checkmark$\checkmark^\checkmark\\checkmarkc\checkmarki\checkmarkr\checkmarkc\checkmark$\checkmark,\checkmark$\checkmark^\checkmark\\checkmarkc\checkmarki\checkmarkr\checkmarkc\checkmark$\checkmark \checkmark$\checkmark^\checkmark\\checkmarkc\checkmarki\checkmarkr\checkmarkc\checkmark$\checkmarki\checkmark$\checkmark^\checkmark\\checkmarkc\checkmarki\checkmarkr\checkmarkc\checkmark$\checkmarkl\checkmark$\checkmark^\checkmark\\checkmarkc\checkmarki\checkmarkr\checkmarkc\checkmark$\checkmark \checkmark$\checkmark^\checkmark\\checkmarkc\checkmarki\checkmarkr\checkmarkc\checkmark$\checkmarkc\checkmark$\checkmark^\checkmark\\checkmarkc\checkmarki\checkmarkr\checkmarkc\checkmark$\checkmarko\checkmark$\checkmark^\checkmark\\checkmarkc\checkmarki\checkmarkr\checkmarkc\checkmark$\checkmarkr\checkmark$\checkmark^\checkmark\\checkmarkc\checkmarki\checkmarkr\checkmarkc\checkmark$\checkmarkr\checkmark$\checkmark^\checkmark\\checkmarkc\checkmarki\checkmarkr\checkmarkc\checkmark$\checkmarke\checkmark$\checkmark^\checkmark\\checkmarkc\checkmarki\checkmarkr\checkmarkc\checkmark$\checkmarks\checkmark$\checkmark^\checkmark\\checkmarkc\checkmarki\checkmarkr\checkmarkc\checkmark$\checkmarkp\checkmark$\checkmark^\checkmark\\checkmarkc\checkmarki\checkmarkr\checkmarkc\checkmark$\checkmarko\checkmark$\checkmark^\checkmark\\checkmarkc\checkmarki\checkmarkr\checkmarkc\checkmark$\checkmarkn\checkmark$\checkmark^\checkmark\\checkmarkc\checkmarki\checkmarkr\checkmarkc\checkmark$\checkmarkd\checkmark$\checkmark^\checkmark\\checkmarkc\checkmarki\checkmarkr\checkmarkc\checkmark$\checkmark \checkmark$\checkmark^\checkmark\\checkmarkc\checkmarki\checkmarkr\checkmarkc\checkmark$\checkmarkà\checkmark$\checkmark^\checkmark\\checkmarkc\checkmarki\checkmarkr\checkmarkc\checkmark$\checkmark \checkmark$\checkmark^\checkmark\\checkmarkc\checkmarki\checkmarkr\checkmarkc\checkmark$\checkmark:\checkmark$\checkmark^\checkmark\\checkmarkc\checkmarki\checkmarkr\checkmarkc\checkmark$\checkmark
\checkmark$\checkmark^\checkmark\\checkmarkc\checkmarki\checkmarkr\checkmarkc\checkmark$\checkmark\\checkmark$\checkmark^\checkmark\\checkmarkc\checkmarki\checkmarkr\checkmarkc\checkmark$\checkmarkb\checkmark$\checkmark^\checkmark\\checkmarkc\checkmarki\checkmarkr\checkmarkc\checkmark$\checkmarke\checkmark$\checkmark^\checkmark\\checkmarkc\checkmarki\checkmarkr\checkmarkc\checkmark$\checkmarkg\checkmark$\checkmark^\checkmark\\checkmarkc\checkmarki\checkmarkr\checkmarkc\checkmark$\checkmarki\checkmark$\checkmark^\checkmark\\checkmarkc\checkmarki\checkmarkr\checkmarkc\checkmark$\checkmarkn\checkmark$\checkmark^\checkmark\\checkmarkc\checkmarki\checkmarkr\checkmarkc\checkmark$\checkmark{\checkmark$\checkmark^\checkmark\\checkmarkc\checkmarki\checkmarkr\checkmarkc\checkmark$\checkmarke\checkmark$\checkmark^\checkmark\\checkmarkc\checkmarki\checkmarkr\checkmarkc\checkmark$\checkmarkq\checkmark$\checkmark^\checkmark\\checkmarkc\checkmarki\checkmarkr\checkmarkc\checkmark$\checkmarku\checkmark$\checkmark^\checkmark\\checkmarkc\checkmarki\checkmarkr\checkmarkc\checkmark$\checkmarka\checkmark$\checkmark^\checkmark\\checkmarkc\checkmarki\checkmarkr\checkmarkc\checkmark$\checkmarkt\checkmark$\checkmark^\checkmark\\checkmarkc\checkmarki\checkmarkr\checkmarkc\checkmark$\checkmarki\checkmark$\checkmark^\checkmark\\checkmarkc\checkmarki\checkmarkr\checkmarkc\checkmark$\checkmarko\checkmark$\checkmark^\checkmark\\checkmarkc\checkmarki\checkmarkr\checkmarkc\checkmark$\checkmarkn\checkmark$\checkmark^\checkmark\\checkmarkc\checkmarki\checkmarkr\checkmarkc\checkmark$\checkmark}\checkmark$\checkmark^\checkmark\\checkmarkc\checkmarki\checkmarkr\checkmarkc\checkmark$\checkmark
\checkmark$\checkmark^\checkmark\\checkmarkc\checkmarki\checkmarkr\checkmarkc\checkmark$\checkmark \checkmark$\checkmark^\checkmark\\checkmarkc\checkmarki\checkmarkr\checkmarkc\checkmark$\checkmark \checkmark$\checkmark^\checkmark\\checkmarkc\checkmarki\checkmarkr\checkmarkc\checkmark$\checkmark \checkmark$\checkmark^\checkmark\\checkmarkc\checkmarki\checkmarkr\checkmarkc\checkmark$\checkmark \checkmark$\checkmark^\checkmark\\checkmarkc\checkmarki\checkmarkr\checkmarkc\checkmark$\checkmarkM\checkmark$\checkmark^\checkmark\\checkmarkc\checkmarki\checkmarkr\checkmarkc\checkmark$\checkmark_\checkmark$\checkmark^\checkmark\\checkmarkc\checkmarki\checkmarkr\checkmarkc\checkmark$\checkmark{\checkmark$\checkmark^\checkmark\\checkmarkc\checkmarki\checkmarkr\checkmarkc\checkmark$\checkmarkf\checkmark$\checkmark^\checkmark\\checkmarkc\checkmarki\checkmarkr\checkmarkc\checkmark$\checkmark\\checkmark$\checkmark^\checkmark\\checkmarkc\checkmarki\checkmarkr\checkmarkc\checkmark$\checkmark_\checkmark$\checkmark^\checkmark\\checkmarkc\checkmarki\checkmarkr\checkmarkc\checkmark$\checkmarkm\checkmark$\checkmark^\checkmark\\checkmarkc\checkmarki\checkmarkr\checkmarkc\checkmark$\checkmarka\checkmark$\checkmark^\checkmark\\checkmarkc\checkmarki\checkmarkr\checkmarkc\checkmark$\checkmarkx\checkmark$\checkmark^\checkmark\\checkmarkc\checkmarki\checkmarkr\checkmarkc\checkmark$\checkmark}\checkmark$\checkmark^\checkmark\\checkmarkc\checkmarki\checkmarkr\checkmarkc\checkmark$\checkmark \checkmark$\checkmark^\checkmark\\checkmarkc\checkmarki\checkmarkr\checkmarkc\checkmark$\checkmark=\checkmark$\checkmark^\checkmark\\checkmarkc\checkmarki\checkmarkr\checkmarkc\checkmark$\checkmark \checkmark$\checkmark^\checkmark\\checkmarkc\checkmarki\checkmarkr\checkmarkc\checkmark$\checkmarkF\checkmark$\checkmark^\checkmark\\checkmarkc\checkmarki\checkmarkr\checkmarkc\checkmark$\checkmark_\checkmark$\checkmark^\checkmark\\checkmarkc\checkmarki\checkmarkr\checkmarkc\checkmark$\checkmark{\checkmark$\checkmark^\checkmark\\checkmarkc\checkmarki\checkmarkr\checkmarkc\checkmark$\checkmarkr\checkmark$\checkmark^\checkmark\\checkmarkc\checkmarki\checkmarkr\checkmarkc\checkmark$\checkmark\\checkmark$\checkmark^\checkmark\\checkmarkc\checkmarki\checkmarkr\checkmarkc\checkmark$\checkmark_\checkmark$\checkmark^\checkmark\\checkmarkc\checkmarki\checkmarkr\checkmarkc\checkmark$\checkmarkp\checkmark$\checkmark^\checkmark\\checkmarkc\checkmarki\checkmarkr\checkmarkc\checkmark$\checkmarka\checkmark$\checkmark^\checkmark\\checkmarkc\checkmarki\checkmarkr\checkmarkc\checkmark$\checkmarkl\checkmark$\checkmark^\checkmark\\checkmarkc\checkmarki\checkmarkr\checkmarkc\checkmark$\checkmarki\checkmark$\checkmark^\checkmark\\checkmarkc\checkmarki\checkmarkr\checkmarkc\checkmark$\checkmarke\checkmark$\checkmark^\checkmark\\checkmarkc\checkmarki\checkmarkr\checkmarkc\checkmark$\checkmarkr\checkmark$\checkmark^\checkmark\\checkmarkc\checkmarki\checkmarkr\checkmarkc\checkmark$\checkmark}\checkmark$\checkmark^\checkmark\\checkmarkc\checkmarki\checkmarkr\checkmarkc\checkmark$\checkmark \checkmark$\checkmark^\checkmark\\checkmarkc\checkmarki\checkmarkr\checkmarkc\checkmark$\checkmark\\checkmark$\checkmark^\checkmark\\checkmarkc\checkmarki\checkmarkr\checkmarkc\checkmark$\checkmarkt\checkmark$\checkmark^\checkmark\\checkmarkc\checkmarki\checkmarkr\checkmarkc\checkmark$\checkmarki\checkmark$\checkmark^\checkmark\\checkmarkc\checkmarki\checkmarkr\checkmarkc\checkmark$\checkmarkm\checkmark$\checkmark^\checkmark\\checkmarkc\checkmarki\checkmarkr\checkmarkc\checkmark$\checkmarke\checkmark$\checkmark^\checkmark\\checkmarkc\checkmarki\checkmarkr\checkmarkc\checkmark$\checkmarks\checkmark$\checkmark^\checkmark\\checkmarkc\checkmarki\checkmarkr\checkmarkc\checkmark$\checkmark \checkmark$\checkmark^\checkmark\\checkmarkc\checkmarki\checkmarkr\checkmarkc\checkmark$\checkmark\\checkmark$\checkmark^\checkmark\\checkmarkc\checkmarki\checkmarkr\checkmarkc\checkmark$\checkmarkf\checkmark$\checkmark^\checkmark\\checkmarkc\checkmarki\checkmarkr\checkmarkc\checkmark$\checkmarkr\checkmark$\checkmark^\checkmark\\checkmarkc\checkmarki\checkmarkr\checkmarkc\checkmark$\checkmarka\checkmark$\checkmark^\checkmark\\checkmarkc\checkmarki\checkmarkr\checkmarkc\checkmark$\checkmarkc\checkmark$\checkmark^\checkmark\\checkmarkc\checkmarki\checkmarkr\checkmarkc\checkmark$\checkmark{\checkmark$\checkmark^\checkmark\\checkmarkc\checkmarki\checkmarkr\checkmarkc\checkmark$\checkmarkL\checkmark$\checkmark^\checkmark\\checkmarkc\checkmarki\checkmarkr\checkmarkc\checkmark$\checkmark_\checkmark$\checkmark^\checkmark\\checkmarkc\checkmarki\checkmarkr\checkmarkc\checkmark$\checkmark{\checkmark$\checkmark^\checkmark\\checkmarkc\checkmarki\checkmarkr\checkmarkc\checkmark$\checkmarkp\checkmark$\checkmark^\checkmark\\checkmarkc\checkmarki\checkmarkr\checkmarkc\checkmark$\checkmarka\checkmark$\checkmark^\checkmark\\checkmarkc\checkmarki\checkmarkr\checkmarkc\checkmark$\checkmarkl\checkmark$\checkmark^\checkmark\\checkmarkc\checkmarki\checkmarkr\checkmarkc\checkmark$\checkmarki\checkmark$\checkmark^\checkmark\\checkmarkc\checkmarki\checkmarkr\checkmarkc\checkmark$\checkmarke\checkmark$\checkmark^\checkmark\\checkmarkc\checkmarki\checkmarkr\checkmarkc\checkmark$\checkmarkr\checkmark$\checkmark^\checkmark\\checkmarkc\checkmarki\checkmarkr\checkmarkc\checkmark$\checkmarks\checkmark$\checkmark^\checkmark\\checkmarkc\checkmarki\checkmarkr\checkmarkc\checkmark$\checkmark}\checkmark$\checkmark^\checkmark\\checkmarkc\checkmarki\checkmarkr\checkmarkc\checkmark$\checkmark}\checkmark$\checkmark^\checkmark\\checkmarkc\checkmarki\checkmarkr\checkmarkc\checkmark$\checkmark{\checkmark$\checkmark^\checkmark\\checkmarkc\checkmarki\checkmarkr\checkmarkc\checkmark$\checkmark4\checkmark$\checkmark^\checkmark\\checkmarkc\checkmarki\checkmarkr\checkmarkc\checkmark$\checkmark}\checkmark$\checkmark^\checkmark\\checkmarkc\checkmarki\checkmarkr\checkmarkc\checkmark$\checkmark \checkmark$\checkmark^\checkmark\\checkmarkc\checkmarki\checkmarkr\checkmarkc\checkmark$\checkmark=\checkmark$\checkmark^\checkmark\\checkmarkc\checkmarki\checkmarkr\checkmarkc\checkmark$\checkmark \checkmark$\checkmark^\checkmark\\checkmarkc\checkmarki\checkmarkr\checkmarkc\checkmark$\checkmark3\checkmark$\checkmark^\checkmark\\checkmarkc\checkmarki\checkmarkr\checkmarkc\checkmark$\checkmark0\checkmark$\checkmark^\checkmark\\checkmarkc\checkmarki\checkmarkr\checkmarkc\checkmark$\checkmark0\checkmark$\checkmark^\checkmark\\checkmarkc\checkmarki\checkmarkr\checkmarkc\checkmark$\checkmark \checkmark$\checkmark^\checkmark\\checkmarkc\checkmarki\checkmarkr\checkmarkc\checkmark$\checkmark\\checkmark$\checkmark^\checkmark\\checkmarkc\checkmarki\checkmarkr\checkmarkc\checkmark$\checkmarkt\checkmark$\checkmark^\checkmark\\checkmarkc\checkmarki\checkmarkr\checkmarkc\checkmark$\checkmarki\checkmark$\checkmark^\checkmark\\checkmarkc\checkmarki\checkmarkr\checkmarkc\checkmark$\checkmarkm\checkmark$\checkmark^\checkmark\\checkmarkc\checkmarki\checkmarkr\checkmarkc\checkmark$\checkmarke\checkmark$\checkmark^\checkmark\\checkmarkc\checkmarki\checkmarkr\checkmarkc\checkmark$\checkmarks\checkmark$\checkmark^\checkmark\\checkmarkc\checkmarki\checkmarkr\checkmarkc\checkmark$\checkmark \checkmark$\checkmark^\checkmark\\checkmarkc\checkmarki\checkmarkr\checkmarkc\checkmark$\checkmark\\checkmark$\checkmark^\checkmark\\checkmarkc\checkmarki\checkmarkr\checkmarkc\checkmark$\checkmarkf\checkmark$\checkmark^\checkmark\\checkmarkc\checkmarki\checkmarkr\checkmarkc\checkmark$\checkmarkr\checkmark$\checkmark^\checkmark\\checkmarkc\checkmarki\checkmarkr\checkmarkc\checkmark$\checkmarka\checkmark$\checkmark^\checkmark\\checkmarkc\checkmarki\checkmarkr\checkmarkc\checkmark$\checkmarkc\checkmark$\checkmark^\checkmark\\checkmarkc\checkmarki\checkmarkr\checkmarkc\checkmark$\checkmark{\checkmark$\checkmark^\checkmark\\checkmarkc\checkmarki\checkmarkr\checkmarkc\checkmark$\checkmark0\checkmark$\checkmark^\checkmark\\checkmarkc\checkmarki\checkmarkr\checkmarkc\checkmark$\checkmark{\checkmark$\checkmark^\checkmark\\checkmarkc\checkmarki\checkmarkr\checkmarkc\checkmark$\checkmark,\checkmark$\checkmark^\checkmark\\checkmarkc\checkmarki\checkmarkr\checkmarkc\checkmark$\checkmark}\checkmark$\checkmark^\checkmark\\checkmarkc\checkmarki\checkmarkr\checkmarkc\checkmark$\checkmark1\checkmark$\checkmark^\checkmark\\checkmarkc\checkmarki\checkmarkr\checkmarkc\checkmark$\checkmark5\checkmark$\checkmark^\checkmark\\checkmarkc\checkmarki\checkmarkr\checkmarkc\checkmark$\checkmark0\checkmark$\checkmark^\checkmark\\checkmarkc\checkmarki\checkmarkr\checkmarkc\checkmark$\checkmark}\checkmark$\checkmark^\checkmark\\checkmarkc\checkmarki\checkmarkr\checkmarkc\checkmark$\checkmark{\checkmark$\checkmark^\checkmark\\checkmarkc\checkmarki\checkmarkr\checkmarkc\checkmark$\checkmark4\checkmark$\checkmark^\checkmark\\checkmarkc\checkmarki\checkmarkr\checkmarkc\checkmark$\checkmark}\checkmark$\checkmark^\checkmark\\checkmarkc\checkmarki\checkmarkr\checkmarkc\checkmark$\checkmark \checkmark$\checkmark^\checkmark\\checkmarkc\checkmarki\checkmarkr\checkmarkc\checkmark$\checkmark=\checkmark$\checkmark^\checkmark\\checkmarkc\checkmarki\checkmarkr\checkmarkc\checkmark$\checkmark \checkmark$\checkmark^\checkmark\\checkmarkc\checkmarki\checkmarkr\checkmarkc\checkmark$\checkmark\\checkmark$\checkmark^\checkmark\\checkmarkc\checkmarki\checkmarkr\checkmarkc\checkmark$\checkmarkb\checkmark$\checkmark^\checkmark\\checkmarkc\checkmarki\checkmarkr\checkmarkc\checkmark$\checkmarko\checkmark$\checkmark^\checkmark\\checkmarkc\checkmarki\checkmarkr\checkmarkc\checkmark$\checkmarkx\checkmark$\checkmark^\checkmark\\checkmarkc\checkmarki\checkmarkr\checkmarkc\checkmark$\checkmarke\checkmark$\checkmark^\checkmark\\checkmarkc\checkmarki\checkmarkr\checkmarkc\checkmark$\checkmarkd\checkmark$\checkmark^\checkmark\\checkmarkc\checkmarki\checkmarkr\checkmarkc\checkmark$\checkmark{\checkmark$\checkmark^\checkmark\\checkmarkc\checkmarki\checkmarkr\checkmarkc\checkmark$\checkmark1\checkmark$\checkmark^\checkmark\\checkmarkc\checkmarki\checkmarkr\checkmarkc\checkmark$\checkmark1\checkmark$\checkmark^\checkmark\\checkmarkc\checkmarki\checkmarkr\checkmarkc\checkmark$\checkmark{\checkmark$\checkmark^\checkmark\\checkmarkc\checkmarki\checkmarkr\checkmarkc\checkmark$\checkmark,\checkmark$\checkmark^\checkmark\\checkmarkc\checkmarki\checkmarkr\checkmarkc\checkmark$\checkmark}\checkmark$\checkmark^\checkmark\\checkmarkc\checkmarki\checkmarkr\checkmarkc\checkmark$\checkmark2\checkmark$\checkmark^\checkmark\\checkmarkc\checkmarki\checkmarkr\checkmarkc\checkmark$\checkmark5\checkmark$\checkmark^\checkmark\\checkmarkc\checkmarki\checkmarkr\checkmarkc\checkmark$\checkmark \checkmark$\checkmark^\checkmark\\checkmarkc\checkmarki\checkmarkr\checkmarkc\checkmark$\checkmark\\checkmark$\checkmark^\checkmark\\checkmarkc\checkmarki\checkmarkr\checkmarkc\checkmark$\checkmarkt\checkmark$\checkmark^\checkmark\\checkmarkc\checkmarki\checkmarkr\checkmarkc\checkmark$\checkmarke\checkmark$\checkmark^\checkmark\\checkmarkc\checkmarki\checkmarkr\checkmarkc\checkmark$\checkmarkx\checkmark$\checkmark^\checkmark\\checkmarkc\checkmarki\checkmarkr\checkmarkc\checkmark$\checkmarkt\checkmark$\checkmark^\checkmark\\checkmarkc\checkmarki\checkmarkr\checkmarkc\checkmark$\checkmark{\checkmark$\checkmark^\checkmark\\checkmarkc\checkmarki\checkmarkr\checkmarkc\checkmark$\checkmark \checkmark$\checkmark^\checkmark\\checkmarkc\checkmarki\checkmarkr\checkmarkc\checkmark$\checkmarkN\checkmark$\checkmark^\checkmark\\checkmarkc\checkmarki\checkmarkr\checkmarkc\checkmark$\checkmark$\checkmark$\checkmark^\checkmark\\checkmarkc\checkmarki\checkmarkr\checkmarkc\checkmark$\checkmark\\checkmark$\checkmark^\checkmark\\checkmarkc\checkmarki\checkmarkr\checkmarkc\checkmark$\checkmarkc\checkmark$\checkmark^\checkmark\\checkmarkc\checkmarki\checkmarkr\checkmarkc\checkmark$\checkmarkd\checkmark$\checkmark^\checkmark\\checkmarkc\checkmarki\checkmarkr\checkmarkc\checkmark$\checkmarko\checkmark$\checkmark^\checkmark\\checkmarkc\checkmarki\checkmarkr\checkmarkc\checkmark$\checkmarkt\checkmark$\checkmark^\checkmark\\checkmarkc\checkmarki\checkmarkr\checkmarkc\checkmark$\checkmark$\checkmark$\checkmark^\checkmark\\checkmarkc\checkmarki\checkmarkr\checkmarkc\checkmark$\checkmarkm\checkmark$\checkmark^\checkmark\\checkmarkc\checkmarki\checkmarkr\checkmarkc\checkmark$\checkmark}\checkmark$\checkmark^\checkmark\\checkmarkc\checkmarki\checkmarkr\checkmarkc\checkmark$\checkmark \checkmark$\checkmark^\checkmark\\checkmarkc\checkmarki\checkmarkr\checkmarkc\checkmark$\checkmark=\checkmark$\checkmark^\checkmark\\checkmarkc\checkmarki\checkmarkr\checkmarkc\checkmark$\checkmark \checkmark$\checkmark^\checkmark\\checkmarkc\checkmarki\checkmarkr\checkmarkc\checkmark$\checkmark1\checkmark$\checkmark^\checkmark\\checkmarkc\checkmarki\checkmarkr\checkmarkc\checkmark$\checkmark1\checkmark$\checkmark^\checkmark\\checkmarkc\checkmarki\checkmarkr\checkmarkc\checkmark$\checkmark\\checkmark$\checkmark^\checkmark\\checkmarkc\checkmarki\checkmarkr\checkmarkc\checkmark$\checkmark,\checkmark$\checkmark^\checkmark\\checkmarkc\checkmarki\checkmarkr\checkmarkc\checkmark$\checkmark2\checkmark$\checkmark^\checkmark\\checkmarkc\checkmarki\checkmarkr\checkmarkc\checkmark$\checkmark5\checkmark$\checkmark^\checkmark\\checkmarkc\checkmarki\checkmarkr\checkmarkc\checkmark$\checkmark0\checkmark$\checkmark^\checkmark\\checkmarkc\checkmarki\checkmarkr\checkmarkc\checkmark$\checkmark \checkmark$\checkmark^\checkmark\\checkmarkc\checkmarki\checkmarkr\checkmarkc\checkmark$\checkmark\\checkmark$\checkmark^\checkmark\\checkmarkc\checkmarki\checkmarkr\checkmarkc\checkmark$\checkmarkt\checkmark$\checkmark^\checkmark\\checkmarkc\checkmarki\checkmarkr\checkmarkc\checkmark$\checkmarke\checkmark$\checkmark^\checkmark\\checkmarkc\checkmarki\checkmarkr\checkmarkc\checkmark$\checkmarkx\checkmark$\checkmark^\checkmark\\checkmarkc\checkmarki\checkmarkr\checkmarkc\checkmark$\checkmarkt\checkmark$\checkmark^\checkmark\\checkmarkc\checkmarki\checkmarkr\checkmarkc\checkmark$\checkmark{\checkmark$\checkmark^\checkmark\\checkmarkc\checkmarki\checkmarkr\checkmarkc\checkmark$\checkmark \checkmark$\checkmark^\checkmark\\checkmarkc\checkmarki\checkmarkr\checkmarkc\checkmark$\checkmarkN\checkmark$\checkmark^\checkmark\\checkmarkc\checkmarki\checkmarkr\checkmarkc\checkmark$\checkmark$\checkmark$\checkmark^\checkmark\\checkmarkc\checkmarki\checkmarkr\checkmarkc\checkmark$\checkmark\\checkmark$\checkmark^\checkmark\\checkmarkc\checkmarki\checkmarkr\checkmarkc\checkmark$\checkmarkc\checkmark$\checkmark^\checkmark\\checkmarkc\checkmarki\checkmarkr\checkmarkc\checkmark$\checkmarkd\checkmark$\checkmark^\checkmark\\checkmarkc\checkmarki\checkmarkr\checkmarkc\checkmark$\checkmarko\checkmark$\checkmark^\checkmark\\checkmarkc\checkmarki\checkmarkr\checkmarkc\checkmark$\checkmarkt\checkmark$\checkmark^\checkmark\\checkmarkc\checkmarki\checkmarkr\checkmarkc\checkmark$\checkmark$\checkmark$\checkmark^\checkmark\\checkmarkc\checkmarki\checkmarkr\checkmarkc\checkmark$\checkmarkm\checkmark$\checkmark^\checkmark\\checkmarkc\checkmarki\checkmarkr\checkmarkc\checkmark$\checkmarkm\checkmark$\checkmark^\checkmark\\checkmarkc\checkmarki\checkmarkr\checkmarkc\checkmark$\checkmark}\checkmark$\checkmark^\checkmark\\checkmarkc\checkmarki\checkmarkr\checkmarkc\checkmark$\checkmark}\checkmark$\checkmark^\checkmark\\checkmarkc\checkmarki\checkmarkr\checkmarkc\checkmark$\checkmark
\checkmark$\checkmark^\checkmark\\checkmarkc\checkmarki\checkmarkr\checkmarkc\checkmark$\checkmark\\checkmark$\checkmark^\checkmark\\checkmarkc\checkmarki\checkmarkr\checkmarkc\checkmark$\checkmarke\checkmark$\checkmark^\checkmark\\checkmarkc\checkmarki\checkmarkr\checkmarkc\checkmark$\checkmarkn\checkmark$\checkmark^\checkmark\\checkmarkc\checkmarki\checkmarkr\checkmarkc\checkmark$\checkmarkd\checkmark$\checkmark^\checkmark\\checkmarkc\checkmarki\checkmarkr\checkmarkc\checkmark$\checkmark{\checkmark$\checkmark^\checkmark\\checkmarkc\checkmarki\checkmarkr\checkmarkc\checkmark$\checkmarke\checkmark$\checkmark^\checkmark\\checkmarkc\checkmarki\checkmarkr\checkmarkc\checkmark$\checkmarkq\checkmark$\checkmark^\checkmark\\checkmarkc\checkmarki\checkmarkr\checkmarkc\checkmark$\checkmarku\checkmark$\checkmark^\checkmark\\checkmarkc\checkmarki\checkmarkr\checkmarkc\checkmark$\checkmarka\checkmark$\checkmark^\checkmark\\checkmarkc\checkmarki\checkmarkr\checkmarkc\checkmark$\checkmarkt\checkmark$\checkmark^\checkmark\\checkmarkc\checkmarki\checkmarkr\checkmarkc\checkmark$\checkmarki\checkmark$\checkmark^\checkmark\\checkmarkc\checkmarki\checkmarkr\checkmarkc\checkmark$\checkmarko\checkmark$\checkmark^\checkmark\\checkmarkc\checkmarki\checkmarkr\checkmarkc\checkmark$\checkmarkn\checkmark$\checkmark^\checkmark\\checkmarkc\checkmarki\checkmarkr\checkmarkc\checkmark$\checkmark}\checkmark$\checkmark^\checkmark\\checkmarkc\checkmarki\checkmarkr\checkmarkc\checkmark$\checkmark
\checkmark$\checkmark^\checkmark\\checkmarkc\checkmarki\checkmarkr\checkmarkc\checkmark$\checkmark
\checkmark$\checkmark^\checkmark\\checkmarkc\checkmarki\checkmarkr\checkmarkc\checkmark$\checkmarkL\checkmark$\checkmark^\checkmark\\checkmarkc\checkmarki\checkmarkr\checkmarkc\checkmark$\checkmarka\checkmark$\checkmark^\checkmark\\checkmarkc\checkmarki\checkmarkr\checkmarkc\checkmark$\checkmark \checkmark$\checkmark^\checkmark\\checkmarkc\checkmarki\checkmarkr\checkmarkc\checkmark$\checkmarkc\checkmark$\checkmark^\checkmark\\checkmarkc\checkmarki\checkmarkr\checkmarkc\checkmark$\checkmarko\checkmark$\checkmark^\checkmark\\checkmarkc\checkmarki\checkmarkr\checkmarkc\checkmark$\checkmarkn\checkmark$\checkmark^\checkmark\\checkmarkc\checkmarki\checkmarkr\checkmarkc\checkmark$\checkmarkt\checkmark$\checkmark^\checkmark\\checkmarkc\checkmarki\checkmarkr\checkmarkc\checkmark$\checkmarkr\checkmark$\checkmark^\checkmark\\checkmarkc\checkmarki\checkmarkr\checkmarkc\checkmark$\checkmarka\checkmark$\checkmark^\checkmark\\checkmarkc\checkmarki\checkmarkr\checkmarkc\checkmark$\checkmarki\checkmark$\checkmark^\checkmark\\checkmarkc\checkmarki\checkmarkr\checkmarkc\checkmark$\checkmarkn\checkmark$\checkmark^\checkmark\\checkmarkc\checkmarki\checkmarkr\checkmarkc\checkmark$\checkmarkt\checkmark$\checkmark^\checkmark\\checkmarkc\checkmarki\checkmarkr\checkmarkc\checkmark$\checkmarke\checkmark$\checkmark^\checkmark\\checkmarkc\checkmarki\checkmarkr\checkmarkc\checkmark$\checkmark \checkmark$\checkmark^\checkmark\\checkmarkc\checkmarki\checkmarkr\checkmarkc\checkmark$\checkmarkd\checkmark$\checkmark^\checkmark\\checkmarkc\checkmarki\checkmarkr\checkmarkc\checkmark$\checkmarke\checkmark$\checkmark^\checkmark\\checkmarkc\checkmarki\checkmarkr\checkmarkc\checkmark$\checkmark \checkmark$\checkmark^\checkmark\\checkmarkc\checkmarki\checkmarkr\checkmarkc\checkmark$\checkmarkf\checkmark$\checkmark^\checkmark\\checkmarkc\checkmarki\checkmarkr\checkmarkc\checkmark$\checkmarkl\checkmark$\checkmark^\checkmark\\checkmarkc\checkmarki\checkmarkr\checkmarkc\checkmark$\checkmarke\checkmark$\checkmark^\checkmark\\checkmarkc\checkmarki\checkmarkr\checkmarkc\checkmark$\checkmarkx\checkmark$\checkmark^\checkmark\\checkmarkc\checkmarki\checkmarkr\checkmarkc\checkmark$\checkmarki\checkmark$\checkmark^\checkmark\\checkmarkc\checkmarki\checkmarkr\checkmarkc\checkmark$\checkmarko\checkmark$\checkmark^\checkmark\\checkmarkc\checkmarki\checkmarkr\checkmarkc\checkmark$\checkmarkn\checkmark$\checkmark^\checkmark\\checkmarkc\checkmarki\checkmarkr\checkmarkc\checkmark$\checkmark \checkmark$\checkmark^\checkmark\\checkmarkc\checkmarki\checkmarkr\checkmarkc\checkmark$\checkmarke\checkmark$\checkmark^\checkmark\\checkmarkc\checkmarki\checkmarkr\checkmarkc\checkmark$\checkmarkn\checkmark$\checkmark^\checkmark\\checkmarkc\checkmarki\checkmarkr\checkmarkc\checkmark$\checkmark \checkmark$\checkmark^\checkmark\\checkmarkc\checkmarki\checkmarkr\checkmarkc\checkmark$\checkmarku\checkmark$\checkmark^\checkmark\\checkmarkc\checkmarki\checkmarkr\checkmarkc\checkmark$\checkmarkn\checkmark$\checkmark^\checkmark\\checkmarkc\checkmarki\checkmarkr\checkmarkc\checkmark$\checkmark \checkmark$\checkmark^\checkmark\\checkmarkc\checkmarki\checkmarkr\checkmarkc\checkmark$\checkmarkp\checkmark$\checkmark^\checkmark\\checkmarkc\checkmarki\checkmarkr\checkmarkc\checkmark$\checkmarko\checkmark$\checkmark^\checkmark\\checkmarkc\checkmarki\checkmarkr\checkmarkc\checkmark$\checkmarki\checkmark$\checkmark^\checkmark\\checkmarkc\checkmarki\checkmarkr\checkmarkc\checkmark$\checkmarkn\checkmark$\checkmark^\checkmark\\checkmarkc\checkmarki\checkmarkr\checkmarkc\checkmark$\checkmarkt\checkmark$\checkmark^\checkmark\\checkmarkc\checkmarki\checkmarkr\checkmarkc\checkmark$\checkmark \checkmark$\checkmark^\checkmark\\checkmarkc\checkmarki\checkmarkr\checkmarkc\checkmark$\checkmarkd\checkmark$\checkmark^\checkmark\\checkmarkc\checkmarki\checkmarkr\checkmarkc\checkmark$\checkmarke\checkmark$\checkmark^\checkmark\\checkmarkc\checkmarki\checkmarkr\checkmarkc\checkmark$\checkmark \checkmark$\checkmark^\checkmark\\checkmarkc\checkmarki\checkmarkr\checkmarkc\checkmark$\checkmarks\checkmark$\checkmark^\checkmark\\checkmarkc\checkmarki\checkmarkr\checkmarkc\checkmark$\checkmarke\checkmark$\checkmark^\checkmark\\checkmarkc\checkmarki\checkmarkr\checkmarkc\checkmark$\checkmarkc\checkmark$\checkmark^\checkmark\\checkmarkc\checkmarki\checkmarkr\checkmarkc\checkmark$\checkmarkt\checkmark$\checkmark^\checkmark\\checkmarkc\checkmarki\checkmarkr\checkmarkc\checkmark$\checkmarki\checkmark$\checkmark^\checkmark\\checkmarkc\checkmarki\checkmarkr\checkmarkc\checkmark$\checkmarko\checkmark$\checkmark^\checkmark\\checkmarkc\checkmarki\checkmarkr\checkmarkc\checkmark$\checkmarkn\checkmark$\checkmark^\checkmark\\checkmarkc\checkmarki\checkmarkr\checkmarkc\checkmark$\checkmark \checkmark$\checkmark^\checkmark\\checkmarkc\checkmarki\checkmarkr\checkmarkc\checkmark$\checkmarkc\checkmark$\checkmark^\checkmark\\checkmarkc\checkmarki\checkmarkr\checkmarkc\checkmark$\checkmarki\checkmark$\checkmark^\checkmark\\checkmarkc\checkmarki\checkmarkr\checkmarkc\checkmark$\checkmarkr\checkmark$\checkmark^\checkmark\\checkmarkc\checkmarki\checkmarkr\checkmarkc\checkmark$\checkmarkc\checkmark$\checkmark^\checkmark\\checkmarkc\checkmarki\checkmarkr\checkmarkc\checkmark$\checkmarku\checkmark$\checkmark^\checkmark\\checkmarkc\checkmarki\checkmarkr\checkmarkc\checkmark$\checkmarkl\checkmark$\checkmark^\checkmark\\checkmarkc\checkmarki\checkmarkr\checkmarkc\checkmark$\checkmarka\checkmark$\checkmark^\checkmark\\checkmarkc\checkmarki\checkmarkr\checkmarkc\checkmark$\checkmarki\checkmark$\checkmark^\checkmark\\checkmarkc\checkmarki\checkmarkr\checkmarkc\checkmark$\checkmarkr\checkmark$\checkmark^\checkmark\\checkmarkc\checkmarki\checkmarkr\checkmarkc\checkmark$\checkmarke\checkmark$\checkmark^\checkmark\\checkmarkc\checkmarki\checkmarkr\checkmarkc\checkmark$\checkmark \checkmark$\checkmark^\checkmark\\checkmarkc\checkmarki\checkmarkr\checkmarkc\checkmark$\checkmarkd\checkmark$\checkmark^\checkmark\\checkmarkc\checkmarki\checkmarkr\checkmarkc\checkmark$\checkmarke\checkmark$\checkmark^\checkmark\\checkmarkc\checkmarki\checkmarkr\checkmarkc\checkmark$\checkmark \checkmark$\checkmark^\checkmark\\checkmarkc\checkmarki\checkmarkr\checkmarkc\checkmark$\checkmarkd\checkmark$\checkmark^\checkmark\\checkmarkc\checkmarki\checkmarkr\checkmarkc\checkmark$\checkmarki\checkmark$\checkmark^\checkmark\\checkmarkc\checkmarki\checkmarkr\checkmarkc\checkmark$\checkmarka\checkmark$\checkmark^\checkmark\\checkmarkc\checkmarki\checkmarkr\checkmarkc\checkmark$\checkmarkm\checkmark$\checkmark^\checkmark\\checkmarkc\checkmarki\checkmarkr\checkmarkc\checkmark$\checkmarkè\checkmark$\checkmark^\checkmark\\checkmarkc\checkmarki\checkmarkr\checkmarkc\checkmark$\checkmarkt\checkmark$\checkmark^\checkmark\\checkmarkc\checkmarki\checkmarkr\checkmarkc\checkmark$\checkmarkr\checkmark$\checkmark^\checkmark\\checkmarkc\checkmarki\checkmarkr\checkmarkc\checkmark$\checkmarke\checkmark$\checkmark^\checkmark\\checkmarkc\checkmarki\checkmarkr\checkmarkc\checkmark$\checkmark \checkmark$\checkmark^\checkmark\\checkmarkc\checkmarki\checkmarkr\checkmarkc\checkmark$\checkmark$\checkmark$\checkmark^\checkmark\\checkmarkc\checkmarki\checkmarkr\checkmarkc\checkmark$\checkmarkd\checkmark$\checkmark^\checkmark\\checkmarkc\checkmarki\checkmarkr\checkmarkc\checkmark$\checkmark$\checkmark$\checkmark^\checkmark\\checkmarkc\checkmarki\checkmarkr\checkmarkc\checkmark$\checkmark \checkmark$\checkmark^\checkmark\\checkmarkc\checkmarki\checkmarkr\checkmarkc\checkmark$\checkmarke\checkmark$\checkmark^\checkmark\\checkmarkc\checkmarki\checkmarkr\checkmarkc\checkmark$\checkmarks\checkmark$\checkmark^\checkmark\\checkmarkc\checkmarki\checkmarkr\checkmarkc\checkmark$\checkmarkt\checkmark$\checkmark^\checkmark\\checkmarkc\checkmarki\checkmarkr\checkmarkc\checkmark$\checkmark \checkmark$\checkmark^\checkmark\\checkmarkc\checkmarki\checkmarkr\checkmarkc\checkmark$\checkmark:\checkmark$\checkmark^\checkmark\\checkmarkc\checkmarki\checkmarkr\checkmarkc\checkmark$\checkmark
\checkmark$\checkmark^\checkmark\\checkmarkc\checkmarki\checkmarkr\checkmarkc\checkmark$\checkmark\\checkmark$\checkmark^\checkmark\\checkmarkc\checkmarki\checkmarkr\checkmarkc\checkmark$\checkmarkb\checkmark$\checkmark^\checkmark\\checkmarkc\checkmarki\checkmarkr\checkmarkc\checkmark$\checkmarke\checkmark$\checkmark^\checkmark\\checkmarkc\checkmarki\checkmarkr\checkmarkc\checkmark$\checkmarkg\checkmark$\checkmark^\checkmark\\checkmarkc\checkmarki\checkmarkr\checkmarkc\checkmark$\checkmarki\checkmark$\checkmark^\checkmark\\checkmarkc\checkmarki\checkmarkr\checkmarkc\checkmark$\checkmarkn\checkmark$\checkmark^\checkmark\\checkmarkc\checkmarki\checkmarkr\checkmarkc\checkmark$\checkmark{\checkmark$\checkmark^\checkmark\\checkmarkc\checkmarki\checkmarkr\checkmarkc\checkmark$\checkmarke\checkmark$\checkmark^\checkmark\\checkmarkc\checkmarki\checkmarkr\checkmarkc\checkmark$\checkmarkq\checkmark$\checkmark^\checkmark\\checkmarkc\checkmarki\checkmarkr\checkmarkc\checkmark$\checkmarku\checkmark$\checkmark^\checkmark\\checkmarkc\checkmarki\checkmarkr\checkmarkc\checkmark$\checkmarka\checkmark$\checkmark^\checkmark\\checkmarkc\checkmarki\checkmarkr\checkmarkc\checkmark$\checkmarkt\checkmark$\checkmark^\checkmark\\checkmarkc\checkmarki\checkmarkr\checkmarkc\checkmark$\checkmarki\checkmark$\checkmark^\checkmark\\checkmarkc\checkmarki\checkmarkr\checkmarkc\checkmark$\checkmarko\checkmark$\checkmark^\checkmark\\checkmarkc\checkmarki\checkmarkr\checkmarkc\checkmark$\checkmarkn\checkmark$\checkmark^\checkmark\\checkmarkc\checkmarki\checkmarkr\checkmarkc\checkmark$\checkmark}\checkmark$\checkmark^\checkmark\\checkmarkc\checkmarki\checkmarkr\checkmarkc\checkmark$\checkmark
\checkmark$\checkmark^\checkmark\\checkmarkc\checkmarki\checkmarkr\checkmarkc\checkmark$\checkmark \checkmark$\checkmark^\checkmark\\checkmarkc\checkmarki\checkmarkr\checkmarkc\checkmark$\checkmark \checkmark$\checkmark^\checkmark\\checkmarkc\checkmarki\checkmarkr\checkmarkc\checkmark$\checkmark \checkmark$\checkmark^\checkmark\\checkmarkc\checkmarki\checkmarkr\checkmarkc\checkmark$\checkmark \checkmark$\checkmark^\checkmark\\checkmarkc\checkmarki\checkmarkr\checkmarkc\checkmark$\checkmark\\checkmark$\checkmark^\checkmark\\checkmarkc\checkmarki\checkmarkr\checkmarkc\checkmark$\checkmarks\checkmark$\checkmark^\checkmark\\checkmarkc\checkmarki\checkmarkr\checkmarkc\checkmark$\checkmarki\checkmark$\checkmark^\checkmark\\checkmarkc\checkmarki\checkmarkr\checkmarkc\checkmark$\checkmarkg\checkmark$\checkmark^\checkmark\\checkmarkc\checkmarki\checkmarkr\checkmarkc\checkmark$\checkmarkm\checkmark$\checkmark^\checkmark\\checkmarkc\checkmarki\checkmarkr\checkmarkc\checkmark$\checkmarka\checkmark$\checkmark^\checkmark\\checkmarkc\checkmarki\checkmarkr\checkmarkc\checkmark$\checkmark_\checkmark$\checkmark^\checkmark\\checkmarkc\checkmarki\checkmarkr\checkmarkc\checkmark$\checkmark{\checkmark$\checkmark^\checkmark\\checkmarkc\checkmarki\checkmarkr\checkmarkc\checkmark$\checkmarkf\checkmark$\checkmark^\checkmark\\checkmarkc\checkmarki\checkmarkr\checkmarkc\checkmark$\checkmarkl\checkmark$\checkmark^\checkmark\\checkmarkc\checkmarki\checkmarkr\checkmarkc\checkmark$\checkmarke\checkmark$\checkmark^\checkmark\\checkmarkc\checkmarki\checkmarkr\checkmarkc\checkmark$\checkmarkx\checkmark$\checkmark^\checkmark\\checkmarkc\checkmarki\checkmarkr\checkmarkc\checkmark$\checkmark}\checkmark$\checkmark^\checkmark\\checkmarkc\checkmarki\checkmarkr\checkmarkc\checkmark$\checkmark \checkmark$\checkmark^\checkmark\\checkmarkc\checkmarki\checkmarkr\checkmarkc\checkmark$\checkmark=\checkmark$\checkmark^\checkmark\\checkmarkc\checkmarki\checkmarkr\checkmarkc\checkmark$\checkmark \checkmark$\checkmark^\checkmark\\checkmarkc\checkmarki\checkmarkr\checkmarkc\checkmark$\checkmark\\checkmark$\checkmark^\checkmark\\checkmarkc\checkmarki\checkmarkr\checkmarkc\checkmark$\checkmarkf\checkmark$\checkmark^\checkmark\\checkmarkc\checkmarki\checkmarkr\checkmarkc\checkmark$\checkmarkr\checkmark$\checkmark^\checkmark\\checkmarkc\checkmarki\checkmarkr\checkmarkc\checkmark$\checkmarka\checkmark$\checkmark^\checkmark\\checkmarkc\checkmarki\checkmarkr\checkmarkc\checkmark$\checkmarkc\checkmark$\checkmark^\checkmark\\checkmarkc\checkmarki\checkmarkr\checkmarkc\checkmark$\checkmark{\checkmark$\checkmark^\checkmark\\checkmarkc\checkmarki\checkmarkr\checkmarkc\checkmark$\checkmarkM\checkmark$\checkmark^\checkmark\\checkmarkc\checkmarki\checkmarkr\checkmarkc\checkmark$\checkmark_\checkmark$\checkmark^\checkmark\\checkmarkc\checkmarki\checkmarkr\checkmarkc\checkmark$\checkmark{\checkmark$\checkmark^\checkmark\\checkmarkc\checkmarki\checkmarkr\checkmarkc\checkmark$\checkmarkf\checkmark$\checkmark^\checkmark\\checkmarkc\checkmarki\checkmarkr\checkmarkc\checkmark$\checkmark\\checkmark$\checkmark^\checkmark\\checkmarkc\checkmarki\checkmarkr\checkmarkc\checkmark$\checkmark_\checkmark$\checkmark^\checkmark\\checkmarkc\checkmarki\checkmarkr\checkmarkc\checkmark$\checkmarkm\checkmark$\checkmark^\checkmark\\checkmarkc\checkmarki\checkmarkr\checkmarkc\checkmark$\checkmarka\checkmark$\checkmark^\checkmark\\checkmarkc\checkmarki\checkmarkr\checkmarkc\checkmark$\checkmarkx\checkmark$\checkmark^\checkmark\\checkmarkc\checkmarki\checkmarkr\checkmarkc\checkmark$\checkmark}\checkmark$\checkmark^\checkmark\\checkmarkc\checkmarki\checkmarkr\checkmarkc\checkmark$\checkmark}\checkmark$\checkmark^\checkmark\\checkmarkc\checkmarki\checkmarkr\checkmarkc\checkmark$\checkmark{\checkmark$\checkmark^\checkmark\\checkmarkc\checkmarki\checkmarkr\checkmarkc\checkmark$\checkmarkW\checkmark$\checkmark^\checkmark\\checkmarkc\checkmarki\checkmarkr\checkmarkc\checkmark$\checkmark_\checkmark$\checkmark^\checkmark\\checkmarkc\checkmarki\checkmarkr\checkmarkc\checkmark$\checkmarkf\checkmark$\checkmark^\checkmark\\checkmarkc\checkmarki\checkmarkr\checkmarkc\checkmark$\checkmark}\checkmark$\checkmark^\checkmark\\checkmarkc\checkmarki\checkmarkr\checkmarkc\checkmark$\checkmark \checkmark$\checkmark^\checkmark\\checkmarkc\checkmarki\checkmarkr\checkmarkc\checkmark$\checkmark=\checkmark$\checkmark^\checkmark\\checkmarkc\checkmarki\checkmarkr\checkmarkc\checkmark$\checkmark \checkmark$\checkmark^\checkmark\\checkmarkc\checkmarki\checkmarkr\checkmarkc\checkmark$\checkmark\\checkmark$\checkmark^\checkmark\\checkmarkc\checkmarki\checkmarkr\checkmarkc\checkmark$\checkmarkf\checkmark$\checkmark^\checkmark\\checkmarkc\checkmarki\checkmarkr\checkmarkc\checkmark$\checkmarkr\checkmark$\checkmark^\checkmark\\checkmarkc\checkmarki\checkmarkr\checkmarkc\checkmark$\checkmarka\checkmark$\checkmark^\checkmark\\checkmarkc\checkmarki\checkmarkr\checkmarkc\checkmark$\checkmarkc\checkmark$\checkmark^\checkmark\\checkmarkc\checkmarki\checkmarkr\checkmarkc\checkmark$\checkmark{\checkmark$\checkmark^\checkmark\\checkmarkc\checkmarki\checkmarkr\checkmarkc\checkmark$\checkmark3\checkmark$\checkmark^\checkmark\\checkmarkc\checkmarki\checkmarkr\checkmarkc\checkmark$\checkmark2\checkmark$\checkmark^\checkmark\\checkmarkc\checkmarki\checkmarkr\checkmarkc\checkmark$\checkmark \checkmark$\checkmark^\checkmark\\checkmarkc\checkmarki\checkmarkr\checkmarkc\checkmark$\checkmark\\checkmark$\checkmark^\checkmark\\checkmarkc\checkmarki\checkmarkr\checkmarkc\checkmark$\checkmarkt\checkmark$\checkmark^\checkmark\\checkmarkc\checkmarki\checkmarkr\checkmarkc\checkmark$\checkmarki\checkmark$\checkmark^\checkmark\\checkmarkc\checkmarki\checkmarkr\checkmarkc\checkmark$\checkmarkm\checkmark$\checkmark^\checkmark\\checkmarkc\checkmarki\checkmarkr\checkmarkc\checkmark$\checkmarke\checkmark$\checkmark^\checkmark\\checkmarkc\checkmarki\checkmarkr\checkmarkc\checkmark$\checkmarks\checkmark$\checkmark^\checkmark\\checkmarkc\checkmarki\checkmarkr\checkmarkc\checkmark$\checkmark \checkmark$\checkmark^\checkmark\\checkmarkc\checkmarki\checkmarkr\checkmarkc\checkmark$\checkmarkM\checkmark$\checkmark^\checkmark\\checkmarkc\checkmarki\checkmarkr\checkmarkc\checkmark$\checkmark_\checkmark$\checkmark^\checkmark\\checkmarkc\checkmarki\checkmarkr\checkmarkc\checkmark$\checkmark{\checkmark$\checkmark^\checkmark\\checkmarkc\checkmarki\checkmarkr\checkmarkc\checkmark$\checkmarkf\checkmark$\checkmark^\checkmark\\checkmarkc\checkmarki\checkmarkr\checkmarkc\checkmark$\checkmark\\checkmark$\checkmark^\checkmark\\checkmarkc\checkmarki\checkmarkr\checkmarkc\checkmark$\checkmark_\checkmark$\checkmark^\checkmark\\checkmarkc\checkmarki\checkmarkr\checkmarkc\checkmark$\checkmarkm\checkmark$\checkmark^\checkmark\\checkmarkc\checkmarki\checkmarkr\checkmarkc\checkmark$\checkmarka\checkmark$\checkmark^\checkmark\\checkmarkc\checkmarki\checkmarkr\checkmarkc\checkmark$\checkmarkx\checkmark$\checkmark^\checkmark\\checkmarkc\checkmarki\checkmarkr\checkmarkc\checkmark$\checkmark}\checkmark$\checkmark^\checkmark\\checkmarkc\checkmarki\checkmarkr\checkmarkc\checkmark$\checkmark}\checkmark$\checkmark^\checkmark\\checkmarkc\checkmarki\checkmarkr\checkmarkc\checkmark$\checkmark{\checkmark$\checkmark^\checkmark\\checkmarkc\checkmarki\checkmarkr\checkmarkc\checkmark$\checkmark\\checkmark$\checkmark^\checkmark\\checkmarkc\checkmarki\checkmarkr\checkmarkc\checkmark$\checkmarkp\checkmark$\checkmark^\checkmark\\checkmarkc\checkmarki\checkmarkr\checkmarkc\checkmark$\checkmarki\checkmark$\checkmark^\checkmark\\checkmarkc\checkmarki\checkmarkr\checkmarkc\checkmark$\checkmark \checkmark$\checkmark^\checkmark\\checkmarkc\checkmarki\checkmarkr\checkmarkc\checkmark$\checkmark\\checkmark$\checkmark^\checkmark\\checkmarkc\checkmarki\checkmarkr\checkmarkc\checkmark$\checkmarkt\checkmark$\checkmark^\checkmark\\checkmarkc\checkmarki\checkmarkr\checkmarkc\checkmark$\checkmarki\checkmark$\checkmark^\checkmark\\checkmarkc\checkmarki\checkmarkr\checkmarkc\checkmark$\checkmarkm\checkmark$\checkmark^\checkmark\\checkmarkc\checkmarki\checkmarkr\checkmarkc\checkmark$\checkmarke\checkmark$\checkmark^\checkmark\\checkmarkc\checkmarki\checkmarkr\checkmarkc\checkmark$\checkmarks\checkmark$\checkmark^\checkmark\\checkmarkc\checkmarki\checkmarkr\checkmarkc\checkmark$\checkmark \checkmark$\checkmark^\checkmark\\checkmarkc\checkmarki\checkmarkr\checkmarkc\checkmark$\checkmarkd\checkmark$\checkmark^\checkmark\\checkmarkc\checkmarki\checkmarkr\checkmarkc\checkmark$\checkmark^\checkmark$\checkmark^\checkmark\\checkmarkc\checkmarki\checkmarkr\checkmarkc\checkmark$\checkmark3\checkmark$\checkmark^\checkmark\\checkmarkc\checkmarki\checkmarkr\checkmarkc\checkmark$\checkmark}\checkmark$\checkmark^\checkmark\\checkmarkc\checkmarki\checkmarkr\checkmarkc\checkmark$\checkmark
\checkmark$\checkmark^\checkmark\\checkmarkc\checkmarki\checkmarkr\checkmarkc\checkmark$\checkmark\\checkmark$\checkmark^\checkmark\\checkmarkc\checkmarki\checkmarkr\checkmarkc\checkmark$\checkmarke\checkmark$\checkmark^\checkmark\\checkmarkc\checkmarki\checkmarkr\checkmarkc\checkmark$\checkmarkn\checkmark$\checkmark^\checkmark\\checkmarkc\checkmarki\checkmarkr\checkmarkc\checkmark$\checkmarkd\checkmark$\checkmark^\checkmark\\checkmarkc\checkmarki\checkmarkr\checkmarkc\checkmark$\checkmark{\checkmark$\checkmark^\checkmark\\checkmarkc\checkmarki\checkmarkr\checkmarkc\checkmark$\checkmarke\checkmark$\checkmark^\checkmark\\checkmarkc\checkmarki\checkmarkr\checkmarkc\checkmark$\checkmarkq\checkmark$\checkmark^\checkmark\\checkmarkc\checkmarki\checkmarkr\checkmarkc\checkmark$\checkmarku\checkmark$\checkmark^\checkmark\\checkmarkc\checkmarki\checkmarkr\checkmarkc\checkmark$\checkmarka\checkmark$\checkmark^\checkmark\\checkmarkc\checkmarki\checkmarkr\checkmarkc\checkmark$\checkmarkt\checkmark$\checkmark^\checkmark\\checkmarkc\checkmarki\checkmarkr\checkmarkc\checkmark$\checkmarki\checkmark$\checkmark^\checkmark\\checkmarkc\checkmarki\checkmarkr\checkmarkc\checkmark$\checkmarko\checkmark$\checkmark^\checkmark\\checkmarkc\checkmarki\checkmarkr\checkmarkc\checkmark$\checkmarkn\checkmark$\checkmark^\checkmark\\checkmarkc\checkmarki\checkmarkr\checkmarkc\checkmark$\checkmark}\checkmark$\checkmark^\checkmark\\checkmarkc\checkmarki\checkmarkr\checkmarkc\checkmark$\checkmark
\checkmark$\checkmark^\checkmark\\checkmarkc\checkmarki\checkmarkr\checkmarkc\checkmark$\checkmark
\checkmark$\checkmark^\checkmark\\checkmarkc\checkmarki\checkmarkr\checkmarkc\checkmark$\checkmarkE\checkmark$\checkmark^\checkmark\\checkmarkc\checkmarki\checkmarkr\checkmarkc\checkmark$\checkmarkn\checkmark$\checkmark^\checkmark\\checkmarkc\checkmarki\checkmarkr\checkmarkc\checkmark$\checkmark \checkmark$\checkmark^\checkmark\\checkmarkc\checkmarki\checkmarkr\checkmarkc\checkmark$\checkmarki\checkmark$\checkmark^\checkmark\\checkmarkc\checkmarki\checkmarkr\checkmarkc\checkmark$\checkmarkm\checkmark$\checkmark^\checkmark\\checkmarkc\checkmarki\checkmarkr\checkmarkc\checkmark$\checkmarkp\checkmark$\checkmark^\checkmark\\checkmarkc\checkmarki\checkmarkr\checkmarkc\checkmark$\checkmarko\checkmark$\checkmark^\checkmark\\checkmarkc\checkmarki\checkmarkr\checkmarkc\checkmark$\checkmarks\checkmark$\checkmark^\checkmark\\checkmarkc\checkmarki\checkmarkr\checkmarkc\checkmark$\checkmarka\checkmark$\checkmark^\checkmark\\checkmarkc\checkmarki\checkmarkr\checkmarkc\checkmark$\checkmarkn\checkmark$\checkmark^\checkmark\\checkmarkc\checkmarki\checkmarkr\checkmarkc\checkmark$\checkmarkt\checkmark$\checkmark^\checkmark\\checkmarkc\checkmarki\checkmarkr\checkmarkc\checkmark$\checkmark \checkmark$\checkmark^\checkmark\\checkmarkc\checkmarki\checkmarkr\checkmarkc\checkmark$\checkmarkl\checkmark$\checkmark^\checkmark\\checkmarkc\checkmarki\checkmarkr\checkmarkc\checkmark$\checkmarka\checkmark$\checkmark^\checkmark\\checkmarkc\checkmarki\checkmarkr\checkmarkc\checkmark$\checkmark \checkmark$\checkmark^\checkmark\\checkmarkc\checkmarki\checkmarkr\checkmarkc\checkmark$\checkmarkc\checkmark$\checkmark^\checkmark\\checkmarkc\checkmarki\checkmarkr\checkmarkc\checkmark$\checkmarko\checkmark$\checkmark^\checkmark\\checkmarkc\checkmarki\checkmarkr\checkmarkc\checkmark$\checkmarkn\checkmark$\checkmark^\checkmark\\checkmarkc\checkmarki\checkmarkr\checkmarkc\checkmark$\checkmarkd\checkmark$\checkmark^\checkmark\\checkmarkc\checkmarki\checkmarkr\checkmarkc\checkmark$\checkmarki\checkmark$\checkmark^\checkmark\\checkmarkc\checkmarki\checkmarkr\checkmarkc\checkmark$\checkmarkt\checkmark$\checkmark^\checkmark\\checkmarkc\checkmarki\checkmarkr\checkmarkc\checkmark$\checkmarki\checkmark$\checkmark^\checkmark\\checkmarkc\checkmarki\checkmarkr\checkmarkc\checkmark$\checkmarko\checkmark$\checkmark^\checkmark\\checkmarkc\checkmarki\checkmarkr\checkmarkc\checkmark$\checkmarkn\checkmark$\checkmark^\checkmark\\checkmarkc\checkmarki\checkmarkr\checkmarkc\checkmark$\checkmark \checkmark$\checkmark^\checkmark\\checkmarkc\checkmarki\checkmarkr\checkmarkc\checkmark$\checkmarkd\checkmark$\checkmark^\checkmark\\checkmarkc\checkmarki\checkmarkr\checkmarkc\checkmark$\checkmarke\checkmark$\checkmark^\checkmark\\checkmarkc\checkmarki\checkmarkr\checkmarkc\checkmark$\checkmark \checkmark$\checkmark^\checkmark\\checkmarkc\checkmarki\checkmarkr\checkmarkc\checkmark$\checkmarkr\checkmark$\checkmark^\checkmark\\checkmarkc\checkmarki\checkmarkr\checkmarkc\checkmark$\checkmarké\checkmark$\checkmark^\checkmark\\checkmarkc\checkmarki\checkmarkr\checkmarkc\checkmark$\checkmarks\checkmark$\checkmark^\checkmark\\checkmarkc\checkmarki\checkmarkr\checkmarkc\checkmark$\checkmarki\checkmark$\checkmark^\checkmark\\checkmarkc\checkmarki\checkmarkr\checkmarkc\checkmark$\checkmarks\checkmark$\checkmark^\checkmark\\checkmarkc\checkmarki\checkmarkr\checkmarkc\checkmark$\checkmarkt\checkmark$\checkmark^\checkmark\\checkmarkc\checkmarki\checkmarkr\checkmarkc\checkmark$\checkmarka\checkmark$\checkmark^\checkmark\\checkmarkc\checkmarki\checkmarkr\checkmarkc\checkmark$\checkmarkn\checkmark$\checkmark^\checkmark\\checkmarkc\checkmarki\checkmarkr\checkmarkc\checkmark$\checkmarkc\checkmark$\checkmark^\checkmark\\checkmarkc\checkmarki\checkmarkr\checkmarkc\checkmark$\checkmarke\checkmark$\checkmark^\checkmark\\checkmarkc\checkmarki\checkmarkr\checkmarkc\checkmark$\checkmark \checkmark$\checkmark^\checkmark\\checkmarkc\checkmarki\checkmarkr\checkmarkc\checkmark$\checkmarka\checkmark$\checkmark^\checkmark\\checkmarkc\checkmarki\checkmarkr\checkmarkc\checkmark$\checkmarkv\checkmark$\checkmark^\checkmark\\checkmarkc\checkmarki\checkmarkr\checkmarkc\checkmark$\checkmarke\checkmark$\checkmark^\checkmark\\checkmarkc\checkmarki\checkmarkr\checkmarkc\checkmark$\checkmarkc\checkmark$\checkmark^\checkmark\\checkmarkc\checkmarki\checkmarkr\checkmarkc\checkmark$\checkmark \checkmark$\checkmark^\checkmark\\checkmarkc\checkmarki\checkmarkr\checkmarkc\checkmark$\checkmark$\checkmark$\checkmark^\checkmark\\checkmarkc\checkmarki\checkmarkr\checkmarkc\checkmark$\checkmarks\checkmark$\checkmark^\checkmark\\checkmarkc\checkmarki\checkmarkr\checkmarkc\checkmark$\checkmark \checkmark$\checkmark^\checkmark\\checkmarkc\checkmarki\checkmarkr\checkmarkc\checkmark$\checkmark=\checkmark$\checkmark^\checkmark\\checkmarkc\checkmarki\checkmarkr\checkmarkc\checkmark$\checkmark \checkmark$\checkmark^\checkmark\\checkmarkc\checkmarki\checkmarkr\checkmarkc\checkmark$\checkmark2\checkmark$\checkmark^\checkmark\\checkmarkc\checkmarki\checkmarkr\checkmarkc\checkmark$\checkmark$\checkmark$\checkmark^\checkmark\\checkmarkc\checkmarki\checkmarkr\checkmarkc\checkmark$\checkmark \checkmark$\checkmark^\checkmark\\checkmarkc\checkmarki\checkmarkr\checkmarkc\checkmark$\checkmark(\checkmark$\checkmark^\checkmark\\checkmarkc\checkmarki\checkmarkr\checkmarkc\checkmark$\checkmarkc\checkmark$\checkmark^\checkmark\\checkmarkc\checkmarki\checkmarkr\checkmarkc\checkmark$\checkmarko\checkmark$\checkmark^\checkmark\\checkmarkc\checkmarki\checkmarkr\checkmarkc\checkmark$\checkmarke\checkmark$\checkmark^\checkmark\\checkmarkc\checkmarki\checkmarkr\checkmarkc\checkmark$\checkmarkf\checkmark$\checkmark^\checkmark\\checkmarkc\checkmarki\checkmarkr\checkmarkc\checkmark$\checkmarkf\checkmark$\checkmark^\checkmark\\checkmarkc\checkmarki\checkmarkr\checkmarkc\checkmark$\checkmarki\checkmark$\checkmark^\checkmark\\checkmarkc\checkmarki\checkmarkr\checkmarkc\checkmark$\checkmarkc\checkmark$\checkmark^\checkmark\\checkmarkc\checkmarki\checkmarkr\checkmarkc\checkmark$\checkmarki\checkmark$\checkmark^\checkmark\\checkmarkc\checkmarki\checkmarkr\checkmarkc\checkmark$\checkmarke\checkmark$\checkmark^\checkmark\\checkmarkc\checkmarki\checkmarkr\checkmarkc\checkmark$\checkmarkn\checkmark$\checkmark^\checkmark\\checkmarkc\checkmarki\checkmarkr\checkmarkc\checkmark$\checkmarkt\checkmark$\checkmark^\checkmark\\checkmarkc\checkmarki\checkmarkr\checkmarkc\checkmark$\checkmark \checkmark$\checkmark^\checkmark\\checkmarkc\checkmarki\checkmarkr\checkmarkc\checkmark$\checkmarkd\checkmark$\checkmark^\checkmark\\checkmarkc\checkmarki\checkmarkr\checkmarkc\checkmark$\checkmarke\checkmark$\checkmark^\checkmark\\checkmarkc\checkmarki\checkmarkr\checkmarkc\checkmark$\checkmark \checkmark$\checkmark^\checkmark\\checkmarkc\checkmarki\checkmarkr\checkmarkc\checkmark$\checkmarks\checkmark$\checkmark^\checkmark\\checkmarkc\checkmarki\checkmarkr\checkmarkc\checkmark$\checkmarké\checkmark$\checkmark^\checkmark\\checkmarkc\checkmarki\checkmarkr\checkmarkc\checkmark$\checkmarkc\checkmark$\checkmark^\checkmark\\checkmarkc\checkmarki\checkmarkr\checkmarkc\checkmark$\checkmarku\checkmark$\checkmark^\checkmark\\checkmarkc\checkmarki\checkmarkr\checkmarkc\checkmark$\checkmarkr\checkmark$\checkmark^\checkmark\\checkmarkc\checkmarki\checkmarkr\checkmarkc\checkmark$\checkmarki\checkmark$\checkmark^\checkmark\\checkmarkc\checkmarki\checkmarkr\checkmarkc\checkmark$\checkmarkt\checkmark$\checkmark^\checkmark\\checkmarkc\checkmarki\checkmarkr\checkmarkc\checkmark$\checkmarké\checkmark$\checkmark^\checkmark\\checkmarkc\checkmarki\checkmarkr\checkmarkc\checkmark$\checkmark \checkmark$\checkmark^\checkmark\\checkmarkc\checkmarki\checkmarkr\checkmarkc\checkmark$\checkmarke\checkmark$\checkmark^\checkmark\\checkmarkc\checkmarki\checkmarkr\checkmarkc\checkmark$\checkmarkn\checkmark$\checkmark^\checkmark\\checkmarkc\checkmarki\checkmarkr\checkmarkc\checkmark$\checkmark \checkmark$\checkmark^\checkmark\\checkmarkc\checkmarki\checkmarkr\checkmarkc\checkmark$\checkmarks\checkmark$\checkmark^\checkmark\\checkmarkc\checkmarki\checkmarkr\checkmarkc\checkmark$\checkmarkt\checkmark$\checkmark^\checkmark\\checkmarkc\checkmarki\checkmarkr\checkmarkc\checkmark$\checkmarka\checkmark$\checkmark^\checkmark\\checkmarkc\checkmarki\checkmarkr\checkmarkc\checkmark$\checkmarkt\checkmark$\checkmark^\checkmark\\checkmarkc\checkmarki\checkmarkr\checkmarkc\checkmark$\checkmarki\checkmark$\checkmark^\checkmark\\checkmarkc\checkmarki\checkmarkr\checkmarkc\checkmark$\checkmarkq\checkmark$\checkmark^\checkmark\\checkmarkc\checkmarki\checkmarkr\checkmarkc\checkmark$\checkmarku\checkmark$\checkmark^\checkmark\\checkmarkc\checkmarki\checkmarkr\checkmarkc\checkmark$\checkmarke\checkmark$\checkmark^\checkmark\\checkmarkc\checkmarki\checkmarkr\checkmarkc\checkmark$\checkmark \checkmark$\checkmark^\checkmark\\checkmarkc\checkmarki\checkmarkr\checkmarkc\checkmark$\checkmarka\checkmark$\checkmark^\checkmark\\checkmarkc\checkmarki\checkmarkr\checkmarkc\checkmark$\checkmarkv\checkmark$\checkmark^\checkmark\\checkmarkc\checkmarki\checkmarkr\checkmarkc\checkmark$\checkmarke\checkmark$\checkmark^\checkmark\\checkmarkc\checkmarki\checkmarkr\checkmarkc\checkmark$\checkmarkc\checkmark$\checkmark^\checkmark\\checkmarkc\checkmarki\checkmarkr\checkmarkc\checkmark$\checkmark \checkmark$\checkmark^\checkmark\\checkmarkc\checkmarki\checkmarkr\checkmarkc\checkmark$\checkmarkm\checkmark$\checkmark^\checkmark\\checkmarkc\checkmarki\checkmarkr\checkmarkc\checkmark$\checkmarko\checkmark$\checkmark^\checkmark\\checkmarkc\checkmarki\checkmarkr\checkmarkc\checkmark$\checkmarkm\checkmark$\checkmark^\checkmark\\checkmarkc\checkmarki\checkmarkr\checkmarkc\checkmark$\checkmarke\checkmark$\checkmark^\checkmark\\checkmarkc\checkmarki\checkmarkr\checkmarkc\checkmark$\checkmarkn\checkmark$\checkmark^\checkmark\\checkmarkc\checkmarki\checkmarkr\checkmarkc\checkmark$\checkmarkt\checkmark$\checkmark^\checkmark\\checkmarkc\checkmarki\checkmarkr\checkmarkc\checkmark$\checkmarks\checkmark$\checkmark^\checkmark\\checkmarkc\checkmarki\checkmarkr\checkmarkc\checkmark$\checkmark \checkmark$\checkmark^\checkmark\\checkmarkc\checkmarki\checkmarkr\checkmarkc\checkmark$\checkmarka\checkmark$\checkmark^\checkmark\\checkmarkc\checkmarki\checkmarkr\checkmarkc\checkmark$\checkmarkl\checkmark$\checkmark^\checkmark\\checkmarkc\checkmarki\checkmarkr\checkmarkc\checkmark$\checkmarkt\checkmark$\checkmark^\checkmark\\checkmarkc\checkmarki\checkmarkr\checkmarkc\checkmark$\checkmarke\checkmark$\checkmark^\checkmark\\checkmarkc\checkmarki\checkmarkr\checkmarkc\checkmark$\checkmarkr\checkmark$\checkmark^\checkmark\\checkmarkc\checkmarki\checkmarkr\checkmarkc\checkmark$\checkmarkn\checkmark$\checkmark^\checkmark\\checkmarkc\checkmarki\checkmarkr\checkmarkc\checkmark$\checkmarké\checkmark$\checkmark^\checkmark\\checkmarkc\checkmarki\checkmarkr\checkmarkc\checkmark$\checkmarks\checkmark$\checkmark^\checkmark\\checkmarkc\checkmarki\checkmarkr\checkmarkc\checkmark$\checkmark)\checkmark$\checkmark^\checkmark\\checkmarkc\checkmarki\checkmarkr\checkmarkc\checkmark$\checkmark \checkmark$\checkmark^\checkmark\\checkmarkc\checkmarki\checkmarkr\checkmarkc\checkmark$\checkmark:\checkmark$\checkmark^\checkmark\\checkmarkc\checkmarki\checkmarkr\checkmarkc\checkmark$\checkmark
\checkmark$\checkmark^\checkmark\\checkmarkc\checkmarki\checkmarkr\checkmarkc\checkmark$\checkmark\\checkmark$\checkmark^\checkmark\\checkmarkc\checkmarki\checkmarkr\checkmarkc\checkmark$\checkmarkb\checkmark$\checkmark^\checkmark\\checkmarkc\checkmarki\checkmarkr\checkmarkc\checkmark$\checkmarke\checkmark$\checkmark^\checkmark\\checkmarkc\checkmarki\checkmarkr\checkmarkc\checkmark$\checkmarkg\checkmark$\checkmark^\checkmark\\checkmarkc\checkmarki\checkmarkr\checkmarkc\checkmark$\checkmarki\checkmark$\checkmark^\checkmark\\checkmarkc\checkmarki\checkmarkr\checkmarkc\checkmark$\checkmarkn\checkmark$\checkmark^\checkmark\\checkmarkc\checkmarki\checkmarkr\checkmarkc\checkmark$\checkmark{\checkmark$\checkmark^\checkmark\\checkmarkc\checkmarki\checkmarkr\checkmarkc\checkmark$\checkmarke\checkmark$\checkmark^\checkmark\\checkmarkc\checkmarki\checkmarkr\checkmarkc\checkmark$\checkmarkq\checkmark$\checkmark^\checkmark\\checkmarkc\checkmarki\checkmarkr\checkmarkc\checkmark$\checkmarku\checkmark$\checkmark^\checkmark\\checkmarkc\checkmarki\checkmarkr\checkmarkc\checkmark$\checkmarka\checkmark$\checkmark^\checkmark\\checkmarkc\checkmarki\checkmarkr\checkmarkc\checkmark$\checkmarkt\checkmark$\checkmark^\checkmark\\checkmarkc\checkmarki\checkmarkr\checkmarkc\checkmark$\checkmarki\checkmark$\checkmark^\checkmark\\checkmarkc\checkmarki\checkmarkr\checkmarkc\checkmark$\checkmarko\checkmark$\checkmark^\checkmark\\checkmarkc\checkmarki\checkmarkr\checkmarkc\checkmark$\checkmarkn\checkmark$\checkmark^\checkmark\\checkmarkc\checkmarki\checkmarkr\checkmarkc\checkmark$\checkmark}\checkmark$\checkmark^\checkmark\\checkmarkc\checkmarki\checkmarkr\checkmarkc\checkmark$\checkmark
\checkmark$\checkmark^\checkmark\\checkmarkc\checkmarki\checkmarkr\checkmarkc\checkmark$\checkmark \checkmark$\checkmark^\checkmark\\checkmarkc\checkmarki\checkmarkr\checkmarkc\checkmark$\checkmark \checkmark$\checkmark^\checkmark\\checkmarkc\checkmarki\checkmarkr\checkmarkc\checkmark$\checkmark \checkmark$\checkmark^\checkmark\\checkmarkc\checkmarki\checkmarkr\checkmarkc\checkmark$\checkmark \checkmark$\checkmark^\checkmark\\checkmarkc\checkmarki\checkmarkr\checkmarkc\checkmark$\checkmark\\checkmark$\checkmark^\checkmark\\checkmarkc\checkmarki\checkmarkr\checkmarkc\checkmark$\checkmarks\checkmark$\checkmark^\checkmark\\checkmarkc\checkmarki\checkmarkr\checkmarkc\checkmark$\checkmarki\checkmark$\checkmark^\checkmark\\checkmarkc\checkmarki\checkmarkr\checkmarkc\checkmark$\checkmarkg\checkmark$\checkmark^\checkmark\\checkmarkc\checkmarki\checkmarkr\checkmarkc\checkmark$\checkmarkm\checkmark$\checkmark^\checkmark\\checkmarkc\checkmarki\checkmarkr\checkmarkc\checkmark$\checkmarka\checkmark$\checkmark^\checkmark\\checkmarkc\checkmarki\checkmarkr\checkmarkc\checkmark$\checkmark_\checkmark$\checkmark^\checkmark\\checkmarkc\checkmarki\checkmarkr\checkmarkc\checkmark$\checkmark{\checkmark$\checkmark^\checkmark\\checkmarkc\checkmarki\checkmarkr\checkmarkc\checkmark$\checkmarkf\checkmark$\checkmark^\checkmark\\checkmarkc\checkmarki\checkmarkr\checkmarkc\checkmark$\checkmarkl\checkmark$\checkmark^\checkmark\\checkmarkc\checkmarki\checkmarkr\checkmarkc\checkmark$\checkmarke\checkmark$\checkmark^\checkmark\\checkmarkc\checkmarki\checkmarkr\checkmarkc\checkmark$\checkmarkx\checkmark$\checkmark^\checkmark\\checkmarkc\checkmarki\checkmarkr\checkmarkc\checkmark$\checkmark}\checkmark$\checkmark^\checkmark\\checkmarkc\checkmarki\checkmarkr\checkmarkc\checkmark$\checkmark \checkmark$\checkmark^\checkmark\\checkmarkc\checkmarki\checkmarkr\checkmarkc\checkmark$\checkmark\\checkmark$\checkmark^\checkmark\\checkmarkc\checkmarki\checkmarkr\checkmarkc\checkmark$\checkmarkl\checkmark$\checkmark^\checkmark\\checkmarkc\checkmarki\checkmarkr\checkmarkc\checkmark$\checkmarke\checkmark$\checkmark^\checkmark\\checkmarkc\checkmarki\checkmarkr\checkmarkc\checkmark$\checkmarkq\checkmark$\checkmark^\checkmark\\checkmarkc\checkmarki\checkmarkr\checkmarkc\checkmark$\checkmark \checkmark$\checkmark^\checkmark\\checkmarkc\checkmarki\checkmarkr\checkmarkc\checkmark$\checkmark\\checkmark$\checkmark^\checkmark\\checkmarkc\checkmarki\checkmarkr\checkmarkc\checkmark$\checkmarkf\checkmark$\checkmark^\checkmark\\checkmarkc\checkmarki\checkmarkr\checkmarkc\checkmark$\checkmarkr\checkmark$\checkmark^\checkmark\\checkmarkc\checkmarki\checkmarkr\checkmarkc\checkmark$\checkmarka\checkmark$\checkmark^\checkmark\\checkmarkc\checkmarki\checkmarkr\checkmarkc\checkmark$\checkmarkc\checkmark$\checkmark^\checkmark\\checkmarkc\checkmarki\checkmarkr\checkmarkc\checkmark$\checkmark{\checkmark$\checkmark^\checkmark\\checkmarkc\checkmarki\checkmarkr\checkmarkc\checkmark$\checkmarkR\checkmark$\checkmark^\checkmark\\checkmarkc\checkmarki\checkmarkr\checkmarkc\checkmark$\checkmark_\checkmark$\checkmark^\checkmark\\checkmarkc\checkmarki\checkmarkr\checkmarkc\checkmark$\checkmarke\checkmark$\checkmark^\checkmark\\checkmarkc\checkmarki\checkmarkr\checkmarkc\checkmark$\checkmark}\checkmark$\checkmark^\checkmark\\checkmarkc\checkmarki\checkmarkr\checkmarkc\checkmark$\checkmark{\checkmark$\checkmark^\checkmark\\checkmarkc\checkmarki\checkmarkr\checkmarkc\checkmark$\checkmarks\checkmark$\checkmark^\checkmark\\checkmarkc\checkmarki\checkmarkr\checkmarkc\checkmark$\checkmark}\checkmark$\checkmark^\checkmark\\checkmarkc\checkmarki\checkmarkr\checkmarkc\checkmark$\checkmark \checkmark$\checkmark^\checkmark\\checkmarkc\checkmarki\checkmarkr\checkmarkc\checkmark$\checkmark=\checkmark$\checkmark^\checkmark\\checkmarkc\checkmarki\checkmarkr\checkmarkc\checkmark$\checkmark \checkmark$\checkmark^\checkmark\\checkmarkc\checkmarki\checkmarkr\checkmarkc\checkmark$\checkmark\\checkmark$\checkmark^\checkmark\\checkmarkc\checkmarki\checkmarkr\checkmarkc\checkmark$\checkmarkf\checkmark$\checkmark^\checkmark\\checkmarkc\checkmarki\checkmarkr\checkmarkc\checkmark$\checkmarkr\checkmark$\checkmark^\checkmark\\checkmarkc\checkmarki\checkmarkr\checkmarkc\checkmark$\checkmarka\checkmark$\checkmark^\checkmark\\checkmarkc\checkmarki\checkmarkr\checkmarkc\checkmark$\checkmarkc\checkmark$\checkmark^\checkmark\\checkmarkc\checkmarki\checkmarkr\checkmarkc\checkmark$\checkmark{\checkmark$\checkmark^\checkmark\\checkmarkc\checkmarki\checkmarkr\checkmarkc\checkmark$\checkmark3\checkmark$\checkmark^\checkmark\\checkmarkc\checkmarki\checkmarkr\checkmarkc\checkmark$\checkmark5\checkmark$\checkmark^\checkmark\\checkmarkc\checkmarki\checkmarkr\checkmarkc\checkmark$\checkmark0\checkmark$\checkmark^\checkmark\\checkmarkc\checkmarki\checkmarkr\checkmarkc\checkmark$\checkmark}\checkmark$\checkmark^\checkmark\\checkmarkc\checkmarki\checkmarkr\checkmarkc\checkmark$\checkmark{\checkmark$\checkmark^\checkmark\\checkmarkc\checkmarki\checkmarkr\checkmarkc\checkmark$\checkmark2\checkmark$\checkmark^\checkmark\\checkmarkc\checkmarki\checkmarkr\checkmarkc\checkmark$\checkmark}\checkmark$\checkmark^\checkmark\\checkmarkc\checkmarki\checkmarkr\checkmarkc\checkmark$\checkmark \checkmark$\checkmark^\checkmark\\checkmarkc\checkmarki\checkmarkr\checkmarkc\checkmark$\checkmark=\checkmark$\checkmark^\checkmark\\checkmarkc\checkmarki\checkmarkr\checkmarkc\checkmark$\checkmark \checkmark$\checkmark^\checkmark\\checkmarkc\checkmarki\checkmarkr\checkmarkc\checkmark$\checkmark1\checkmark$\checkmark^\checkmark\\checkmarkc\checkmarki\checkmarkr\checkmarkc\checkmark$\checkmark7\checkmark$\checkmark^\checkmark\\checkmarkc\checkmarki\checkmarkr\checkmarkc\checkmark$\checkmark5\checkmark$\checkmark^\checkmark\\checkmarkc\checkmarki\checkmarkr\checkmarkc\checkmark$\checkmark \checkmark$\checkmark^\checkmark\\checkmarkc\checkmarki\checkmarkr\checkmarkc\checkmark$\checkmark\\checkmark$\checkmark^\checkmark\\checkmarkc\checkmarki\checkmarkr\checkmarkc\checkmark$\checkmarkt\checkmark$\checkmark^\checkmark\\checkmarkc\checkmarki\checkmarkr\checkmarkc\checkmark$\checkmarke\checkmark$\checkmark^\checkmark\\checkmarkc\checkmarki\checkmarkr\checkmarkc\checkmark$\checkmarkx\checkmark$\checkmark^\checkmark\\checkmarkc\checkmarki\checkmarkr\checkmarkc\checkmark$\checkmarkt\checkmark$\checkmark^\checkmark\\checkmarkc\checkmarki\checkmarkr\checkmarkc\checkmark$\checkmark{\checkmark$\checkmark^\checkmark\\checkmarkc\checkmarki\checkmarkr\checkmarkc\checkmark$\checkmark \checkmark$\checkmark^\checkmark\\checkmarkc\checkmarki\checkmarkr\checkmarkc\checkmark$\checkmarkM\checkmark$\checkmark^\checkmark\\checkmarkc\checkmarki\checkmarkr\checkmarkc\checkmark$\checkmarkP\checkmark$\checkmark^\checkmark\\checkmarkc\checkmarki\checkmarkr\checkmarkc\checkmark$\checkmarka\checkmark$\checkmark^\checkmark\\checkmarkc\checkmarki\checkmarkr\checkmarkc\checkmark$\checkmark}\checkmark$\checkmark^\checkmark\\checkmarkc\checkmarki\checkmarkr\checkmarkc\checkmark$\checkmark
\checkmark$\checkmark^\checkmark\\checkmarkc\checkmarki\checkmarkr\checkmarkc\checkmark$\checkmark\\checkmark$\checkmark^\checkmark\\checkmarkc\checkmarki\checkmarkr\checkmarkc\checkmark$\checkmarke\checkmark$\checkmark^\checkmark\\checkmarkc\checkmarki\checkmarkr\checkmarkc\checkmark$\checkmarkn\checkmark$\checkmark^\checkmark\\checkmarkc\checkmarki\checkmarkr\checkmarkc\checkmark$\checkmarkd\checkmark$\checkmark^\checkmark\\checkmarkc\checkmarki\checkmarkr\checkmarkc\checkmark$\checkmark{\checkmark$\checkmark^\checkmark\\checkmarkc\checkmarki\checkmarkr\checkmarkc\checkmark$\checkmarke\checkmark$\checkmark^\checkmark\\checkmarkc\checkmarki\checkmarkr\checkmarkc\checkmark$\checkmarkq\checkmark$\checkmark^\checkmark\\checkmarkc\checkmarki\checkmarkr\checkmarkc\checkmark$\checkmarku\checkmark$\checkmark^\checkmark\\checkmarkc\checkmarki\checkmarkr\checkmarkc\checkmark$\checkmarka\checkmark$\checkmark^\checkmark\\checkmarkc\checkmarki\checkmarkr\checkmarkc\checkmark$\checkmarkt\checkmark$\checkmark^\checkmark\\checkmarkc\checkmarki\checkmarkr\checkmarkc\checkmark$\checkmarki\checkmark$\checkmark^\checkmark\\checkmarkc\checkmarki\checkmarkr\checkmarkc\checkmark$\checkmarko\checkmark$\checkmark^\checkmark\\checkmarkc\checkmarki\checkmarkr\checkmarkc\checkmark$\checkmarkn\checkmark$\checkmark^\checkmark\\checkmarkc\checkmarki\checkmarkr\checkmarkc\checkmark$\checkmark}\checkmark$\checkmark^\checkmark\\checkmarkc\checkmarki\checkmarkr\checkmarkc\checkmark$\checkmark
\checkmark$\checkmark^\checkmark\\checkmarkc\checkmarki\checkmarkr\checkmarkc\checkmark$\checkmark
\checkmark$\checkmark^\checkmark\\checkmarkc\checkmarki\checkmarkr\checkmarkc\checkmark$\checkmarkE\checkmark$\checkmark^\checkmark\\checkmarkc\checkmarki\checkmarkr\checkmarkc\checkmark$\checkmarkn\checkmark$\checkmark^\checkmark\\checkmarkc\checkmarki\checkmarkr\checkmarkc\checkmark$\checkmark \checkmark$\checkmark^\checkmark\\checkmarkc\checkmarki\checkmarkr\checkmarkc\checkmark$\checkmarki\checkmark$\checkmark^\checkmark\\checkmarkc\checkmarki\checkmarkr\checkmarkc\checkmark$\checkmarks\checkmark$\checkmark^\checkmark\\checkmarkc\checkmarki\checkmarkr\checkmarkc\checkmark$\checkmarko\checkmark$\checkmark^\checkmark\\checkmarkc\checkmarki\checkmarkr\checkmarkc\checkmark$\checkmarkl\checkmark$\checkmark^\checkmark\\checkmarkc\checkmarki\checkmarkr\checkmarkc\checkmark$\checkmarka\checkmark$\checkmark^\checkmark\\checkmarkc\checkmarki\checkmarkr\checkmarkc\checkmark$\checkmarkn\checkmark$\checkmark^\checkmark\\checkmarkc\checkmarki\checkmarkr\checkmarkc\checkmark$\checkmarkt\checkmark$\checkmark^\checkmark\\checkmarkc\checkmarki\checkmarkr\checkmarkc\checkmark$\checkmark \checkmark$\checkmark^\checkmark\\checkmarkc\checkmarki\checkmarkr\checkmarkc\checkmark$\checkmark$\checkmark$\checkmark^\checkmark\\checkmarkc\checkmarki\checkmarkr\checkmarkc\checkmark$\checkmarkd\checkmark$\checkmark^\checkmark\\checkmarkc\checkmarki\checkmarkr\checkmarkc\checkmark$\checkmark$\checkmark$\checkmark^\checkmark\\checkmarkc\checkmarki\checkmarkr\checkmarkc\checkmark$\checkmark \checkmark$\checkmark^\checkmark\\checkmarkc\checkmarki\checkmarkr\checkmarkc\checkmark$\checkmark:\checkmark$\checkmark^\checkmark\\checkmarkc\checkmarki\checkmarkr\checkmarkc\checkmark$\checkmark
\checkmark$\checkmark^\checkmark\\checkmarkc\checkmarki\checkmarkr\checkmarkc\checkmark$\checkmark\\checkmark$\checkmark^\checkmark\\checkmarkc\checkmarki\checkmarkr\checkmarkc\checkmark$\checkmarkb\checkmark$\checkmark^\checkmark\\checkmarkc\checkmarki\checkmarkr\checkmarkc\checkmark$\checkmarke\checkmark$\checkmark^\checkmark\\checkmarkc\checkmarki\checkmarkr\checkmarkc\checkmark$\checkmarkg\checkmark$\checkmark^\checkmark\\checkmarkc\checkmarki\checkmarkr\checkmarkc\checkmark$\checkmarki\checkmark$\checkmark^\checkmark\\checkmarkc\checkmarki\checkmarkr\checkmarkc\checkmark$\checkmarkn\checkmark$\checkmark^\checkmark\\checkmarkc\checkmarki\checkmarkr\checkmarkc\checkmark$\checkmark{\checkmark$\checkmark^\checkmark\\checkmarkc\checkmarki\checkmarkr\checkmarkc\checkmark$\checkmarke\checkmark$\checkmark^\checkmark\\checkmarkc\checkmarki\checkmarkr\checkmarkc\checkmark$\checkmarkq\checkmark$\checkmark^\checkmark\\checkmarkc\checkmarki\checkmarkr\checkmarkc\checkmark$\checkmarku\checkmark$\checkmark^\checkmark\\checkmarkc\checkmarki\checkmarkr\checkmarkc\checkmark$\checkmarka\checkmark$\checkmark^\checkmark\\checkmarkc\checkmarki\checkmarkr\checkmarkc\checkmark$\checkmarkt\checkmark$\checkmark^\checkmark\\checkmarkc\checkmarki\checkmarkr\checkmarkc\checkmark$\checkmarki\checkmark$\checkmark^\checkmark\\checkmarkc\checkmarki\checkmarkr\checkmarkc\checkmark$\checkmarko\checkmark$\checkmark^\checkmark\\checkmarkc\checkmarki\checkmarkr\checkmarkc\checkmark$\checkmarkn\checkmark$\checkmark^\checkmark\\checkmarkc\checkmarki\checkmarkr\checkmarkc\checkmark$\checkmark}\checkmark$\checkmark^\checkmark\\checkmarkc\checkmarki\checkmarkr\checkmarkc\checkmark$\checkmark
\checkmark$\checkmark^\checkmark\\checkmarkc\checkmarki\checkmarkr\checkmarkc\checkmark$\checkmark \checkmark$\checkmark^\checkmark\\checkmarkc\checkmarki\checkmarkr\checkmarkc\checkmark$\checkmark \checkmark$\checkmark^\checkmark\\checkmarkc\checkmarki\checkmarkr\checkmarkc\checkmark$\checkmark \checkmark$\checkmark^\checkmark\\checkmarkc\checkmarki\checkmarkr\checkmarkc\checkmark$\checkmark \checkmark$\checkmark^\checkmark\\checkmarkc\checkmarki\checkmarkr\checkmarkc\checkmark$\checkmarkd\checkmark$\checkmark^\checkmark\\checkmarkc\checkmarki\checkmarkr\checkmarkc\checkmark$\checkmark \checkmark$\checkmark^\checkmark\\checkmarkc\checkmarki\checkmarkr\checkmarkc\checkmark$\checkmark\\checkmark$\checkmark^\checkmark\\checkmarkc\checkmarki\checkmarkr\checkmarkc\checkmark$\checkmarkg\checkmark$\checkmark^\checkmark\\checkmarkc\checkmarki\checkmarkr\checkmarkc\checkmark$\checkmarke\checkmark$\checkmark^\checkmark\\checkmarkc\checkmarki\checkmarkr\checkmarkc\checkmark$\checkmarkq\checkmark$\checkmark^\checkmark\\checkmarkc\checkmarki\checkmarkr\checkmarkc\checkmark$\checkmark \checkmark$\checkmark^\checkmark\\checkmarkc\checkmarki\checkmarkr\checkmarkc\checkmark$\checkmark\\checkmark$\checkmark^\checkmark\\checkmarkc\checkmarki\checkmarkr\checkmarkc\checkmark$\checkmarks\checkmark$\checkmark^\checkmark\\checkmarkc\checkmarki\checkmarkr\checkmarkc\checkmark$\checkmarkq\checkmark$\checkmark^\checkmark\\checkmarkc\checkmarki\checkmarkr\checkmarkc\checkmark$\checkmarkr\checkmark$\checkmark^\checkmark\\checkmarkc\checkmarki\checkmarkr\checkmarkc\checkmark$\checkmarkt\checkmark$\checkmark^\checkmark\\checkmarkc\checkmarki\checkmarkr\checkmarkc\checkmark$\checkmark[\checkmark$\checkmark^\checkmark\\checkmarkc\checkmarki\checkmarkr\checkmarkc\checkmark$\checkmark3\checkmark$\checkmark^\checkmark\\checkmarkc\checkmarki\checkmarkr\checkmarkc\checkmark$\checkmark]\checkmark$\checkmark^\checkmark\\checkmarkc\checkmarki\checkmarkr\checkmarkc\checkmark$\checkmark{\checkmark$\checkmark^\checkmark\\checkmarkc\checkmarki\checkmarkr\checkmarkc\checkmark$\checkmark\\checkmark$\checkmark^\checkmark\\checkmarkc\checkmarki\checkmarkr\checkmarkc\checkmark$\checkmarkf\checkmark$\checkmark^\checkmark\\checkmarkc\checkmarki\checkmarkr\checkmarkc\checkmark$\checkmarkr\checkmark$\checkmark^\checkmark\\checkmarkc\checkmarki\checkmarkr\checkmarkc\checkmark$\checkmarka\checkmark$\checkmark^\checkmark\\checkmarkc\checkmarki\checkmarkr\checkmarkc\checkmark$\checkmarkc\checkmark$\checkmark^\checkmark\\checkmarkc\checkmarki\checkmarkr\checkmarkc\checkmark$\checkmark{\checkmark$\checkmark^\checkmark\\checkmarkc\checkmarki\checkmarkr\checkmarkc\checkmark$\checkmark3\checkmark$\checkmark^\checkmark\\checkmarkc\checkmarki\checkmarkr\checkmarkc\checkmark$\checkmark2\checkmark$\checkmark^\checkmark\\checkmarkc\checkmarki\checkmarkr\checkmarkc\checkmark$\checkmark \checkmark$\checkmark^\checkmark\\checkmarkc\checkmarki\checkmarkr\checkmarkc\checkmark$\checkmark\\checkmark$\checkmark^\checkmark\\checkmarkc\checkmarki\checkmarkr\checkmarkc\checkmark$\checkmarkt\checkmark$\checkmark^\checkmark\\checkmarkc\checkmarki\checkmarkr\checkmarkc\checkmark$\checkmarki\checkmark$\checkmark^\checkmark\\checkmarkc\checkmarki\checkmarkr\checkmarkc\checkmark$\checkmarkm\checkmark$\checkmark^\checkmark\\checkmarkc\checkmarki\checkmarkr\checkmarkc\checkmark$\checkmarke\checkmark$\checkmark^\checkmark\\checkmarkc\checkmarki\checkmarkr\checkmarkc\checkmark$\checkmarks\checkmark$\checkmark^\checkmark\\checkmarkc\checkmarki\checkmarkr\checkmarkc\checkmark$\checkmark \checkmark$\checkmark^\checkmark\\checkmarkc\checkmarki\checkmarkr\checkmarkc\checkmark$\checkmarkM\checkmark$\checkmark^\checkmark\\checkmarkc\checkmarki\checkmarkr\checkmarkc\checkmark$\checkmark_\checkmark$\checkmark^\checkmark\\checkmarkc\checkmarki\checkmarkr\checkmarkc\checkmark$\checkmark{\checkmark$\checkmark^\checkmark\\checkmarkc\checkmarki\checkmarkr\checkmarkc\checkmark$\checkmarkf\checkmark$\checkmark^\checkmark\\checkmarkc\checkmarki\checkmarkr\checkmarkc\checkmark$\checkmark\\checkmark$\checkmark^\checkmark\\checkmarkc\checkmarki\checkmarkr\checkmarkc\checkmark$\checkmark_\checkmark$\checkmark^\checkmark\\checkmarkc\checkmarki\checkmarkr\checkmarkc\checkmark$\checkmarkm\checkmark$\checkmark^\checkmark\\checkmarkc\checkmarki\checkmarkr\checkmarkc\checkmark$\checkmarka\checkmark$\checkmark^\checkmark\\checkmarkc\checkmarki\checkmarkr\checkmarkc\checkmark$\checkmarkx\checkmark$\checkmark^\checkmark\\checkmarkc\checkmarki\checkmarkr\checkmarkc\checkmark$\checkmark}\checkmark$\checkmark^\checkmark\\checkmarkc\checkmarki\checkmarkr\checkmarkc\checkmark$\checkmark}\checkmark$\checkmark^\checkmark\\checkmarkc\checkmarki\checkmarkr\checkmarkc\checkmark$\checkmark{\checkmark$\checkmark^\checkmark\\checkmarkc\checkmarki\checkmarkr\checkmarkc\checkmark$\checkmark\\checkmark$\checkmark^\checkmark\\checkmarkc\checkmarki\checkmarkr\checkmarkc\checkmark$\checkmarkp\checkmark$\checkmark^\checkmark\\checkmarkc\checkmarki\checkmarkr\checkmarkc\checkmark$\checkmarki\checkmark$\checkmark^\checkmark\\checkmarkc\checkmarki\checkmarkr\checkmarkc\checkmark$\checkmark \checkmark$\checkmark^\checkmark\\checkmarkc\checkmarki\checkmarkr\checkmarkc\checkmark$\checkmark\\checkmark$\checkmark^\checkmark\\checkmarkc\checkmarki\checkmarkr\checkmarkc\checkmark$\checkmarkt\checkmark$\checkmark^\checkmark\\checkmarkc\checkmarki\checkmarkr\checkmarkc\checkmark$\checkmarki\checkmark$\checkmark^\checkmark\\checkmarkc\checkmarki\checkmarkr\checkmarkc\checkmark$\checkmarkm\checkmark$\checkmark^\checkmark\\checkmarkc\checkmarki\checkmarkr\checkmarkc\checkmark$\checkmarke\checkmark$\checkmark^\checkmark\\checkmarkc\checkmarki\checkmarkr\checkmarkc\checkmark$\checkmarks\checkmark$\checkmark^\checkmark\\checkmarkc\checkmarki\checkmarkr\checkmarkc\checkmark$\checkmark \checkmark$\checkmark^\checkmark\\checkmarkc\checkmarki\checkmarkr\checkmarkc\checkmark$\checkmark\\checkmark$\checkmark^\checkmark\\checkmarkc\checkmarki\checkmarkr\checkmarkc\checkmark$\checkmarks\checkmark$\checkmark^\checkmark\\checkmarkc\checkmarki\checkmarkr\checkmarkc\checkmark$\checkmarki\checkmark$\checkmark^\checkmark\\checkmarkc\checkmarki\checkmarkr\checkmarkc\checkmark$\checkmarkg\checkmark$\checkmark^\checkmark\\checkmarkc\checkmarki\checkmarkr\checkmarkc\checkmark$\checkmarkm\checkmark$\checkmark^\checkmark\\checkmarkc\checkmarki\checkmarkr\checkmarkc\checkmark$\checkmarka\checkmark$\checkmark^\checkmark\\checkmarkc\checkmarki\checkmarkr\checkmarkc\checkmark$\checkmark_\checkmark$\checkmark^\checkmark\\checkmarkc\checkmarki\checkmarkr\checkmarkc\checkmark$\checkmark{\checkmark$\checkmark^\checkmark\\checkmarkc\checkmarki\checkmarkr\checkmarkc\checkmark$\checkmarka\checkmark$\checkmark^\checkmark\\checkmarkc\checkmarki\checkmarkr\checkmarkc\checkmark$\checkmarkd\checkmark$\checkmark^\checkmark\\checkmarkc\checkmarki\checkmarkr\checkmarkc\checkmark$\checkmarkm\checkmark$\checkmark^\checkmark\\checkmarkc\checkmarki\checkmarkr\checkmarkc\checkmark$\checkmark}\checkmark$\checkmark^\checkmark\\checkmarkc\checkmarki\checkmarkr\checkmarkc\checkmark$\checkmark}\checkmark$\checkmark^\checkmark\\checkmarkc\checkmarki\checkmarkr\checkmarkc\checkmark$\checkmark}\checkmark$\checkmark^\checkmark\\checkmarkc\checkmarki\checkmarkr\checkmarkc\checkmark$\checkmark \checkmark$\checkmark^\checkmark\\checkmarkc\checkmarki\checkmarkr\checkmarkc\checkmark$\checkmark=\checkmark$\checkmark^\checkmark\\checkmarkc\checkmarki\checkmarkr\checkmarkc\checkmark$\checkmark \checkmark$\checkmark^\checkmark\\checkmarkc\checkmarki\checkmarkr\checkmarkc\checkmark$\checkmark\\checkmark$\checkmark^\checkmark\\checkmarkc\checkmarki\checkmarkr\checkmarkc\checkmark$\checkmarks\checkmark$\checkmark^\checkmark\\checkmarkc\checkmarki\checkmarkr\checkmarkc\checkmark$\checkmarkq\checkmark$\checkmark^\checkmark\\checkmarkc\checkmarki\checkmarkr\checkmarkc\checkmark$\checkmarkr\checkmark$\checkmark^\checkmark\\checkmarkc\checkmarki\checkmarkr\checkmarkc\checkmark$\checkmarkt\checkmark$\checkmark^\checkmark\\checkmarkc\checkmarki\checkmarkr\checkmarkc\checkmark$\checkmark[\checkmark$\checkmark^\checkmark\\checkmarkc\checkmarki\checkmarkr\checkmarkc\checkmark$\checkmark3\checkmark$\checkmark^\checkmark\\checkmarkc\checkmarki\checkmarkr\checkmarkc\checkmark$\checkmark]\checkmark$\checkmark^\checkmark\\checkmarkc\checkmarki\checkmarkr\checkmarkc\checkmark$\checkmark{\checkmark$\checkmark^\checkmark\\checkmarkc\checkmarki\checkmarkr\checkmarkc\checkmark$\checkmark\\checkmark$\checkmark^\checkmark\\checkmarkc\checkmarki\checkmarkr\checkmarkc\checkmark$\checkmarkf\checkmark$\checkmark^\checkmark\\checkmarkc\checkmarki\checkmarkr\checkmarkc\checkmark$\checkmarkr\checkmark$\checkmark^\checkmark\\checkmarkc\checkmarki\checkmarkr\checkmarkc\checkmark$\checkmarka\checkmark$\checkmark^\checkmark\\checkmarkc\checkmarki\checkmarkr\checkmarkc\checkmark$\checkmarkc\checkmark$\checkmark^\checkmark\\checkmarkc\checkmarki\checkmarkr\checkmarkc\checkmark$\checkmark{\checkmark$\checkmark^\checkmark\\checkmarkc\checkmarki\checkmarkr\checkmarkc\checkmark$\checkmark3\checkmark$\checkmark^\checkmark\\checkmarkc\checkmarki\checkmarkr\checkmarkc\checkmark$\checkmark2\checkmark$\checkmark^\checkmark\\checkmarkc\checkmarki\checkmarkr\checkmarkc\checkmark$\checkmark \checkmark$\checkmark^\checkmark\\checkmarkc\checkmarki\checkmarkr\checkmarkc\checkmark$\checkmark\\checkmark$\checkmark^\checkmark\\checkmarkc\checkmarki\checkmarkr\checkmarkc\checkmark$\checkmarkt\checkmark$\checkmark^\checkmark\\checkmarkc\checkmarki\checkmarkr\checkmarkc\checkmark$\checkmarki\checkmark$\checkmark^\checkmark\\checkmarkc\checkmarki\checkmarkr\checkmarkc\checkmark$\checkmarkm\checkmark$\checkmark^\checkmark\\checkmarkc\checkmarki\checkmarkr\checkmarkc\checkmark$\checkmarke\checkmark$\checkmark^\checkmark\\checkmarkc\checkmarki\checkmarkr\checkmarkc\checkmark$\checkmarks\checkmark$\checkmark^\checkmark\\checkmarkc\checkmarki\checkmarkr\checkmarkc\checkmark$\checkmark \checkmark$\checkmark^\checkmark\\checkmarkc\checkmarki\checkmarkr\checkmarkc\checkmark$\checkmark1\checkmark$\checkmark^\checkmark\\checkmarkc\checkmarki\checkmarkr\checkmarkc\checkmark$\checkmark1\checkmark$\checkmark^\checkmark\\checkmarkc\checkmarki\checkmarkr\checkmarkc\checkmark$\checkmark\\checkmark$\checkmark^\checkmark\\checkmarkc\checkmarki\checkmarkr\checkmarkc\checkmark$\checkmark,\checkmark$\checkmark^\checkmark\\checkmarkc\checkmarki\checkmarkr\checkmarkc\checkmark$\checkmark2\checkmark$\checkmark^\checkmark\\checkmarkc\checkmarki\checkmarkr\checkmarkc\checkmark$\checkmark5\checkmark$\checkmark^\checkmark\\checkmarkc\checkmarki\checkmarkr\checkmarkc\checkmark$\checkmark0\checkmark$\checkmark^\checkmark\\checkmarkc\checkmarki\checkmarkr\checkmarkc\checkmark$\checkmark}\checkmark$\checkmark^\checkmark\\checkmarkc\checkmarki\checkmarkr\checkmarkc\checkmark$\checkmark{\checkmark$\checkmark^\checkmark\\checkmarkc\checkmarki\checkmarkr\checkmarkc\checkmark$\checkmark\\checkmark$\checkmark^\checkmark\\checkmarkc\checkmarki\checkmarkr\checkmarkc\checkmark$\checkmarkp\checkmark$\checkmark^\checkmark\\checkmarkc\checkmarki\checkmarkr\checkmarkc\checkmark$\checkmarki\checkmark$\checkmark^\checkmark\\checkmarkc\checkmarki\checkmarkr\checkmarkc\checkmark$\checkmark \checkmark$\checkmark^\checkmark\\checkmarkc\checkmarki\checkmarkr\checkmarkc\checkmark$\checkmark\\checkmark$\checkmark^\checkmark\\checkmarkc\checkmarki\checkmarkr\checkmarkc\checkmark$\checkmarkt\checkmark$\checkmark^\checkmark\\checkmarkc\checkmarki\checkmarkr\checkmarkc\checkmark$\checkmarki\checkmark$\checkmark^\checkmark\\checkmarkc\checkmarki\checkmarkr\checkmarkc\checkmark$\checkmarkm\checkmark$\checkmark^\checkmark\\checkmarkc\checkmarki\checkmarkr\checkmarkc\checkmark$\checkmarke\checkmark$\checkmark^\checkmark\\checkmarkc\checkmarki\checkmarkr\checkmarkc\checkmark$\checkmarks\checkmark$\checkmark^\checkmark\\checkmarkc\checkmarki\checkmarkr\checkmarkc\checkmark$\checkmark \checkmark$\checkmark^\checkmark\\checkmarkc\checkmarki\checkmarkr\checkmarkc\checkmark$\checkmark1\checkmark$\checkmark^\checkmark\\checkmarkc\checkmarki\checkmarkr\checkmarkc\checkmark$\checkmark7\checkmark$\checkmark^\checkmark\\checkmarkc\checkmarki\checkmarkr\checkmarkc\checkmark$\checkmark5\checkmark$\checkmark^\checkmark\\checkmarkc\checkmarki\checkmarkr\checkmarkc\checkmark$\checkmark}\checkmark$\checkmark^\checkmark\\checkmarkc\checkmarki\checkmarkr\checkmarkc\checkmark$\checkmark}\checkmark$\checkmark^\checkmark\\checkmarkc\checkmarki\checkmarkr\checkmarkc\checkmark$\checkmark \checkmark$\checkmark^\checkmark\\checkmarkc\checkmarki\checkmarkr\checkmarkc\checkmark$\checkmark=\checkmark$\checkmark^\checkmark\\checkmarkc\checkmarki\checkmarkr\checkmarkc\checkmark$\checkmark \checkmark$\checkmark^\checkmark\\checkmarkc\checkmarki\checkmarkr\checkmarkc\checkmark$\checkmark\\checkmark$\checkmark^\checkmark\\checkmarkc\checkmarki\checkmarkr\checkmarkc\checkmark$\checkmarks\checkmark$\checkmark^\checkmark\\checkmarkc\checkmarki\checkmarkr\checkmarkc\checkmark$\checkmarkq\checkmark$\checkmark^\checkmark\\checkmarkc\checkmarki\checkmarkr\checkmarkc\checkmark$\checkmarkr\checkmark$\checkmark^\checkmark\\checkmarkc\checkmarki\checkmarkr\checkmarkc\checkmark$\checkmarkt\checkmark$\checkmark^\checkmark\\checkmarkc\checkmarki\checkmarkr\checkmarkc\checkmark$\checkmark[\checkmark$\checkmark^\checkmark\\checkmarkc\checkmarki\checkmarkr\checkmarkc\checkmark$\checkmark3\checkmark$\checkmark^\checkmark\\checkmarkc\checkmarki\checkmarkr\checkmarkc\checkmark$\checkmark]\checkmark$\checkmark^\checkmark\\checkmarkc\checkmarki\checkmarkr\checkmarkc\checkmark$\checkmark{\checkmark$\checkmark^\checkmark\\checkmarkc\checkmarki\checkmarkr\checkmarkc\checkmark$\checkmark\\checkmark$\checkmark^\checkmark\\checkmarkc\checkmarki\checkmarkr\checkmarkc\checkmark$\checkmarkf\checkmark$\checkmark^\checkmark\\checkmarkc\checkmarki\checkmarkr\checkmarkc\checkmark$\checkmarkr\checkmark$\checkmark^\checkmark\\checkmarkc\checkmarki\checkmarkr\checkmarkc\checkmark$\checkmarka\checkmark$\checkmark^\checkmark\\checkmarkc\checkmarki\checkmarkr\checkmarkc\checkmark$\checkmarkc\checkmark$\checkmark^\checkmark\\checkmarkc\checkmarki\checkmarkr\checkmarkc\checkmark$\checkmark{\checkmark$\checkmark^\checkmark\\checkmarkc\checkmarki\checkmarkr\checkmarkc\checkmark$\checkmark3\checkmark$\checkmark^\checkmark\\checkmarkc\checkmarki\checkmarkr\checkmarkc\checkmark$\checkmark6\checkmark$\checkmark^\checkmark\\checkmarkc\checkmarki\checkmarkr\checkmarkc\checkmark$\checkmark0\checkmark$\checkmark^\checkmark\\checkmarkc\checkmarki\checkmarkr\checkmarkc\checkmark$\checkmark\\checkmark$\checkmark^\checkmark\\checkmarkc\checkmarki\checkmarkr\checkmarkc\checkmark$\checkmark,\checkmark$\checkmark^\checkmark\\checkmarkc\checkmarki\checkmarkr\checkmarkc\checkmark$\checkmark0\checkmark$\checkmark^\checkmark\\checkmarkc\checkmarki\checkmarkr\checkmarkc\checkmark$\checkmark0\checkmark$\checkmark^\checkmark\\checkmarkc\checkmarki\checkmarkr\checkmarkc\checkmark$\checkmark0\checkmark$\checkmark^\checkmark\\checkmarkc\checkmarki\checkmarkr\checkmarkc\checkmark$\checkmark}\checkmark$\checkmark^\checkmark\\checkmarkc\checkmarki\checkmarkr\checkmarkc\checkmark$\checkmark{\checkmark$\checkmark^\checkmark\\checkmarkc\checkmarki\checkmarkr\checkmarkc\checkmark$\checkmark5\checkmark$\checkmark^\checkmark\\checkmarkc\checkmarki\checkmarkr\checkmarkc\checkmark$\checkmark4\checkmark$\checkmark^\checkmark\\checkmarkc\checkmarki\checkmarkr\checkmarkc\checkmark$\checkmark9\checkmark$\checkmark^\checkmark\\checkmarkc\checkmarki\checkmarkr\checkmarkc\checkmark$\checkmark{\checkmark$\checkmark^\checkmark\\checkmarkc\checkmarki\checkmarkr\checkmarkc\checkmark$\checkmark,\checkmark$\checkmark^\checkmark\\checkmarkc\checkmarki\checkmarkr\checkmarkc\checkmark$\checkmark}\checkmark$\checkmark^\checkmark\\checkmarkc\checkmarki\checkmarkr\checkmarkc\checkmark$\checkmark8\checkmark$\checkmark^\checkmark\\checkmarkc\checkmarki\checkmarkr\checkmarkc\checkmark$\checkmark}\checkmark$\checkmark^\checkmark\\checkmarkc\checkmarki\checkmarkr\checkmarkc\checkmark$\checkmark}\checkmark$\checkmark^\checkmark\\checkmarkc\checkmarki\checkmarkr\checkmarkc\checkmark$\checkmark \checkmark$\checkmark^\checkmark\\checkmarkc\checkmarki\checkmarkr\checkmarkc\checkmark$\checkmark=\checkmark$\checkmark^\checkmark\\checkmarkc\checkmarki\checkmarkr\checkmarkc\checkmark$\checkmark \checkmark$\checkmark^\checkmark\\checkmarkc\checkmarki\checkmarkr\checkmarkc\checkmark$\checkmark\\checkmark$\checkmark^\checkmark\\checkmarkc\checkmarki\checkmarkr\checkmarkc\checkmark$\checkmarks\checkmark$\checkmark^\checkmark\\checkmarkc\checkmarki\checkmarkr\checkmarkc\checkmark$\checkmarkq\checkmark$\checkmark^\checkmark\\checkmarkc\checkmarki\checkmarkr\checkmarkc\checkmark$\checkmarkr\checkmark$\checkmark^\checkmark\\checkmarkc\checkmarki\checkmarkr\checkmarkc\checkmark$\checkmarkt\checkmark$\checkmark^\checkmark\\checkmarkc\checkmarki\checkmarkr\checkmarkc\checkmark$\checkmark[\checkmark$\checkmark^\checkmark\\checkmarkc\checkmarki\checkmarkr\checkmarkc\checkmark$\checkmark3\checkmark$\checkmark^\checkmark\\checkmarkc\checkmarki\checkmarkr\checkmarkc\checkmark$\checkmark]\checkmark$\checkmark^\checkmark\\checkmarkc\checkmarki\checkmarkr\checkmarkc\checkmark$\checkmark{\checkmark$\checkmark^\checkmark\\checkmarkc\checkmarki\checkmarkr\checkmarkc\checkmark$\checkmark6\checkmark$\checkmark^\checkmark\\checkmarkc\checkmarki\checkmarkr\checkmarkc\checkmark$\checkmark5\checkmark$\checkmark^\checkmark\\checkmarkc\checkmarki\checkmarkr\checkmarkc\checkmark$\checkmark4\checkmark$\checkmark^\checkmark\\checkmarkc\checkmarki\checkmarkr\checkmarkc\checkmark$\checkmark{\checkmark$\checkmark^\checkmark\\checkmarkc\checkmarki\checkmarkr\checkmarkc\checkmark$\checkmark,\checkmark$\checkmark^\checkmark\\checkmarkc\checkmarki\checkmarkr\checkmarkc\checkmark$\checkmark}\checkmark$\checkmark^\checkmark\\checkmarkc\checkmarki\checkmarkr\checkmarkc\checkmark$\checkmark8\checkmark$\checkmark^\checkmark\\checkmarkc\checkmarki\checkmarkr\checkmarkc\checkmark$\checkmark}\checkmark$\checkmark^\checkmark\\checkmarkc\checkmarki\checkmarkr\checkmarkc\checkmark$\checkmark \checkmark$\checkmark^\checkmark\\checkmarkc\checkmarki\checkmarkr\checkmarkc\checkmark$\checkmark\\checkmark$\checkmark^\checkmark\\checkmarkc\checkmarki\checkmarkr\checkmarkc\checkmark$\checkmarka\checkmark$\checkmark^\checkmark\\checkmarkc\checkmarki\checkmarkr\checkmarkc\checkmark$\checkmarkp\checkmark$\checkmark^\checkmark\\checkmarkc\checkmarki\checkmarkr\checkmarkc\checkmark$\checkmarkp\checkmark$\checkmark^\checkmark\\checkmarkc\checkmarki\checkmarkr\checkmarkc\checkmark$\checkmarkr\checkmark$\checkmark^\checkmark\\checkmarkc\checkmarki\checkmarkr\checkmarkc\checkmark$\checkmarko\checkmark$\checkmark^\checkmark\\checkmarkc\checkmarki\checkmarkr\checkmarkc\checkmark$\checkmarkx\checkmark$\checkmark^\checkmark\\checkmarkc\checkmarki\checkmarkr\checkmarkc\checkmark$\checkmark \checkmark$\checkmark^\checkmark\\checkmarkc\checkmarki\checkmarkr\checkmarkc\checkmark$\checkmark\\checkmark$\checkmark^\checkmark\\checkmarkc\checkmarki\checkmarkr\checkmarkc\checkmark$\checkmarkb\checkmark$\checkmark^\checkmark\\checkmarkc\checkmarki\checkmarkr\checkmarkc\checkmark$\checkmarko\checkmark$\checkmark^\checkmark\\checkmarkc\checkmarki\checkmarkr\checkmarkc\checkmark$\checkmarkx\checkmark$\checkmark^\checkmark\\checkmarkc\checkmarki\checkmarkr\checkmarkc\checkmark$\checkmarke\checkmark$\checkmark^\checkmark\\checkmarkc\checkmarki\checkmarkr\checkmarkc\checkmark$\checkmarkd\checkmark$\checkmark^\checkmark\\checkmarkc\checkmarki\checkmarkr\checkmarkc\checkmark$\checkmark{\checkmark$\checkmark^\checkmark\\checkmarkc\checkmarki\checkmarkr\checkmarkc\checkmark$\checkmark8\checkmark$\checkmark^\checkmark\\checkmarkc\checkmarki\checkmarkr\checkmarkc\checkmark$\checkmark{\checkmark$\checkmark^\checkmark\\checkmarkc\checkmarki\checkmarkr\checkmarkc\checkmark$\checkmark,\checkmark$\checkmark^\checkmark\\checkmarkc\checkmarki\checkmarkr\checkmarkc\checkmark$\checkmark}\checkmark$\checkmark^\checkmark\\checkmarkc\checkmarki\checkmarkr\checkmarkc\checkmark$\checkmark6\checkmark$\checkmark^\checkmark\\checkmarkc\checkmarki\checkmarkr\checkmarkc\checkmark$\checkmark7\checkmark$\checkmark^\checkmark\\checkmarkc\checkmarki\checkmarkr\checkmarkc\checkmark$\checkmark \checkmark$\checkmark^\checkmark\\checkmarkc\checkmarki\checkmarkr\checkmarkc\checkmark$\checkmark\\checkmark$\checkmark^\checkmark\\checkmarkc\checkmarki\checkmarkr\checkmarkc\checkmark$\checkmarkt\checkmark$\checkmark^\checkmark\\checkmarkc\checkmarki\checkmarkr\checkmarkc\checkmark$\checkmarke\checkmark$\checkmark^\checkmark\\checkmarkc\checkmarki\checkmarkr\checkmarkc\checkmark$\checkmarkx\checkmark$\checkmark^\checkmark\\checkmarkc\checkmarki\checkmarkr\checkmarkc\checkmark$\checkmarkt\checkmark$\checkmark^\checkmark\\checkmarkc\checkmarki\checkmarkr\checkmarkc\checkmark$\checkmark{\checkmark$\checkmark^\checkmark\\checkmarkc\checkmarki\checkmarkr\checkmarkc\checkmark$\checkmark \checkmark$\checkmark^\checkmark\\checkmarkc\checkmarki\checkmarkr\checkmarkc\checkmark$\checkmarkm\checkmark$\checkmark^\checkmark\\checkmarkc\checkmarki\checkmarkr\checkmarkc\checkmark$\checkmarkm\checkmark$\checkmark^\checkmark\\checkmarkc\checkmarki\checkmarkr\checkmarkc\checkmark$\checkmark}\checkmark$\checkmark^\checkmark\\checkmarkc\checkmarki\checkmarkr\checkmarkc\checkmark$\checkmark}\checkmark$\checkmark^\checkmark\\checkmarkc\checkmarki\checkmarkr\checkmarkc\checkmark$\checkmark
\checkmark$\checkmark^\checkmark\\checkmarkc\checkmarki\checkmarkr\checkmarkc\checkmark$\checkmark\\checkmark$\checkmark^\checkmark\\checkmarkc\checkmarki\checkmarkr\checkmarkc\checkmark$\checkmarke\checkmark$\checkmark^\checkmark\\checkmarkc\checkmarki\checkmarkr\checkmarkc\checkmark$\checkmarkn\checkmark$\checkmark^\checkmark\\checkmarkc\checkmarki\checkmarkr\checkmarkc\checkmark$\checkmarkd\checkmark$\checkmark^\checkmark\\checkmarkc\checkmarki\checkmarkr\checkmarkc\checkmark$\checkmark{\checkmark$\checkmark^\checkmark\\checkmarkc\checkmarki\checkmarkr\checkmarkc\checkmark$\checkmarke\checkmark$\checkmark^\checkmark\\checkmarkc\checkmarki\checkmarkr\checkmarkc\checkmark$\checkmarkq\checkmark$\checkmark^\checkmark\\checkmarkc\checkmarki\checkmarkr\checkmarkc\checkmark$\checkmarku\checkmark$\checkmark^\checkmark\\checkmarkc\checkmarki\checkmarkr\checkmarkc\checkmark$\checkmarka\checkmark$\checkmark^\checkmark\\checkmarkc\checkmarki\checkmarkr\checkmarkc\checkmark$\checkmarkt\checkmark$\checkmark^\checkmark\\checkmarkc\checkmarki\checkmarkr\checkmarkc\checkmark$\checkmarki\checkmark$\checkmark^\checkmark\\checkmarkc\checkmarki\checkmarkr\checkmarkc\checkmark$\checkmarko\checkmark$\checkmark^\checkmark\\checkmarkc\checkmarki\checkmarkr\checkmarkc\checkmark$\checkmarkn\checkmark$\checkmark^\checkmark\\checkmarkc\checkmarki\checkmarkr\checkmarkc\checkmark$\checkmark}\checkmark$\checkmark^\checkmark\\checkmarkc\checkmarki\checkmarkr\checkmarkc\checkmark$\checkmark
\checkmark$\checkmark^\checkmark\\checkmarkc\checkmarki\checkmarkr\checkmarkc\checkmark$\checkmark
\checkmark$\checkmark^\checkmark\\checkmarkc\checkmarki\checkmarkr\checkmarkc\checkmark$\checkmarkO\checkmark$\checkmark^\checkmark\\checkmarkc\checkmarki\checkmarkr\checkmarkc\checkmark$\checkmarkn\checkmark$\checkmark^\checkmark\\checkmarkc\checkmarki\checkmarkr\checkmarkc\checkmark$\checkmark \checkmark$\checkmark^\checkmark\\checkmarkc\checkmarki\checkmarkr\checkmarkc\checkmark$\checkmarkr\checkmark$\checkmark^\checkmark\\checkmarkc\checkmarki\checkmarkr\checkmarkc\checkmark$\checkmarke\checkmark$\checkmark^\checkmark\\checkmarkc\checkmarki\checkmarkr\checkmarkc\checkmark$\checkmarkt\checkmark$\checkmark^\checkmark\\checkmarkc\checkmarki\checkmarkr\checkmarkc\checkmark$\checkmarki\checkmark$\checkmark^\checkmark\\checkmarkc\checkmarki\checkmarkr\checkmarkc\checkmark$\checkmarke\checkmark$\checkmark^\checkmark\\checkmarkc\checkmarki\checkmarkr\checkmarkc\checkmark$\checkmarkn\checkmark$\checkmark^\checkmark\\checkmarkc\checkmarki\checkmarkr\checkmarkc\checkmark$\checkmarkt\checkmark$\checkmark^\checkmark\\checkmarkc\checkmarki\checkmarkr\checkmarkc\checkmark$\checkmark \checkmark$\checkmark^\checkmark\\checkmarkc\checkmarki\checkmarkr\checkmarkc\checkmark$\checkmarkl\checkmark$\checkmark^\checkmark\\checkmarkc\checkmarki\checkmarkr\checkmarkc\checkmark$\checkmarke\checkmark$\checkmark^\checkmark\\checkmarkc\checkmarki\checkmarkr\checkmarkc\checkmark$\checkmark \checkmark$\checkmark^\checkmark\\checkmarkc\checkmarki\checkmarkr\checkmarkc\checkmark$\checkmarkd\checkmark$\checkmark^\checkmark\\checkmarkc\checkmarki\checkmarkr\checkmarkc\checkmark$\checkmarki\checkmark$\checkmark^\checkmark\\checkmarkc\checkmarki\checkmarkr\checkmarkc\checkmark$\checkmarka\checkmark$\checkmark^\checkmark\\checkmarkc\checkmarki\checkmarkr\checkmarkc\checkmark$\checkmarkm\checkmark$\checkmark^\checkmark\\checkmarkc\checkmarki\checkmarkr\checkmarkc\checkmark$\checkmarkè\checkmark$\checkmark^\checkmark\\checkmarkc\checkmarki\checkmarkr\checkmarkc\checkmark$\checkmarkt\checkmark$\checkmark^\checkmark\\checkmarkc\checkmarki\checkmarkr\checkmarkc\checkmark$\checkmarkr\checkmark$\checkmark^\checkmark\\checkmarkc\checkmarki\checkmarkr\checkmarkc\checkmark$\checkmarke\checkmark$\checkmark^\checkmark\\checkmarkc\checkmarki\checkmarkr\checkmarkc\checkmark$\checkmark \checkmark$\checkmark^\checkmark\\checkmarkc\checkmarki\checkmarkr\checkmarkc\checkmark$\checkmarkn\checkmark$\checkmark^\checkmark\\checkmarkc\checkmarki\checkmarkr\checkmarkc\checkmark$\checkmarko\checkmark$\checkmark^\checkmark\\checkmarkc\checkmarki\checkmarkr\checkmarkc\checkmark$\checkmarkr\checkmark$\checkmark^\checkmark\\checkmarkc\checkmarki\checkmarkr\checkmarkc\checkmark$\checkmarkm\checkmark$\checkmark^\checkmark\\checkmarkc\checkmarki\checkmarkr\checkmarkc\checkmark$\checkmarka\checkmark$\checkmark^\checkmark\\checkmarkc\checkmarki\checkmarkr\checkmarkc\checkmark$\checkmarkl\checkmark$\checkmark^\checkmark\\checkmarkc\checkmarki\checkmarkr\checkmarkc\checkmark$\checkmarki\checkmark$\checkmark^\checkmark\\checkmarkc\checkmarki\checkmarkr\checkmarkc\checkmark$\checkmarks\checkmark$\checkmark^\checkmark\\checkmarkc\checkmarki\checkmarkr\checkmarkc\checkmark$\checkmarké\checkmark$\checkmark^\checkmark\\checkmarkc\checkmarki\checkmarkr\checkmarkc\checkmark$\checkmark \checkmark$\checkmark^\checkmark\\checkmarkc\checkmarki\checkmarkr\checkmarkc\checkmark$\checkmarks\checkmark$\checkmark^\checkmark\\checkmarkc\checkmarki\checkmarkr\checkmarkc\checkmark$\checkmarku\checkmark$\checkmark^\checkmark\\checkmarkc\checkmarki\checkmarkr\checkmarkc\checkmark$\checkmarki\checkmark$\checkmark^\checkmark\\checkmarkc\checkmarki\checkmarkr\checkmarkc\checkmark$\checkmarkv\checkmark$\checkmark^\checkmark\\checkmarkc\checkmarki\checkmarkr\checkmarkc\checkmark$\checkmarka\checkmark$\checkmark^\checkmark\\checkmarkc\checkmarki\checkmarkr\checkmarkc\checkmark$\checkmarkn\checkmark$\checkmark^\checkmark\\checkmarkc\checkmarki\checkmarkr\checkmarkc\checkmark$\checkmarkt\checkmark$\checkmark^\checkmark\\checkmarkc\checkmarki\checkmarkr\checkmarkc\checkmark$\checkmark \checkmark$\checkmark^\checkmark\\checkmarkc\checkmarki\checkmarkr\checkmarkc\checkmark$\checkmarkl\checkmark$\checkmark^\checkmark\\checkmarkc\checkmarki\checkmarkr\checkmarkc\checkmark$\checkmarka\checkmark$\checkmark^\checkmark\\checkmarkc\checkmarki\checkmarkr\checkmarkc\checkmark$\checkmark \checkmark$\checkmark^\checkmark\\checkmarkc\checkmarki\checkmarkr\checkmarkc\checkmark$\checkmarkd\checkmark$\checkmark^\checkmark\\checkmarkc\checkmarki\checkmarkr\checkmarkc\checkmark$\checkmarki\checkmark$\checkmark^\checkmark\\checkmarkc\checkmarki\checkmarkr\checkmarkc\checkmark$\checkmarkm\checkmark$\checkmark^\checkmark\\checkmarkc\checkmarki\checkmarkr\checkmarkc\checkmark$\checkmarke\checkmark$\checkmark^\checkmark\\checkmarkc\checkmarki\checkmarkr\checkmarkc\checkmark$\checkmarkn\checkmark$\checkmark^\checkmark\\checkmarkc\checkmarki\checkmarkr\checkmarkc\checkmark$\checkmarks\checkmark$\checkmark^\checkmark\\checkmarkc\checkmarki\checkmarkr\checkmarkc\checkmark$\checkmarki\checkmark$\checkmark^\checkmark\\checkmarkc\checkmarki\checkmarkr\checkmarkc\checkmark$\checkmarko\checkmark$\checkmark^\checkmark\\checkmarkc\checkmarki\checkmarkr\checkmarkc\checkmark$\checkmarkn\checkmark$\checkmark^\checkmark\\checkmarkc\checkmarki\checkmarkr\checkmarkc\checkmark$\checkmark \checkmark$\checkmark^\checkmark\\checkmarkc\checkmarki\checkmarkr\checkmarkc\checkmark$\checkmarkd\checkmark$\checkmark^\checkmark\\checkmarkc\checkmarki\checkmarkr\checkmarkc\checkmark$\checkmarki\checkmark$\checkmark^\checkmark\\checkmarkc\checkmarki\checkmarkr\checkmarkc\checkmark$\checkmarks\checkmark$\checkmark^\checkmark\\checkmarkc\checkmarki\checkmarkr\checkmarkc\checkmark$\checkmarkp\checkmark$\checkmark^\checkmark\\checkmarkc\checkmarki\checkmarkr\checkmarkc\checkmark$\checkmarko\checkmark$\checkmark^\checkmark\\checkmarkc\checkmarki\checkmarkr\checkmarkc\checkmark$\checkmarkn\checkmark$\checkmark^\checkmark\\checkmarkc\checkmarki\checkmarkr\checkmarkc\checkmark$\checkmarki\checkmark$\checkmark^\checkmark\\checkmarkc\checkmarki\checkmarkr\checkmarkc\checkmark$\checkmarkb\checkmark$\checkmark^\checkmark\\checkmarkc\checkmarki\checkmarkr\checkmarkc\checkmark$\checkmarkl\checkmark$\checkmark^\checkmark\\checkmarkc\checkmarki\checkmarkr\checkmarkc\checkmark$\checkmarke\checkmark$\checkmark^\checkmark\\checkmarkc\checkmarki\checkmarkr\checkmarkc\checkmark$\checkmark \checkmark$\checkmark^\checkmark\\checkmarkc\checkmarki\checkmarkr\checkmarkc\checkmark$\checkmark:\checkmark$\checkmark^\checkmark\\checkmarkc\checkmarki\checkmarkr\checkmarkc\checkmark$\checkmark \checkmark$\checkmark^\checkmark\\checkmarkc\checkmarki\checkmarkr\checkmarkc\checkmark$\checkmark$\checkmark$\checkmark^\checkmark\\checkmarkc\checkmarki\checkmarkr\checkmarkc\checkmark$\checkmarkd\checkmark$\checkmark^\checkmark\\checkmarkc\checkmarki\checkmarkr\checkmarkc\checkmark$\checkmark_\checkmark$\checkmark^\checkmark\\checkmarkc\checkmarki\checkmarkr\checkmarkc\checkmark$\checkmark{\checkmark$\checkmark^\checkmark\\checkmarkc\checkmarki\checkmarkr\checkmarkc\checkmark$\checkmarka\checkmark$\checkmark^\checkmark\\checkmarkc\checkmarki\checkmarkr\checkmarkc\checkmark$\checkmarkr\checkmark$\checkmark^\checkmark\\checkmarkc\checkmarki\checkmarkr\checkmarkc\checkmark$\checkmarkb\checkmark$\checkmark^\checkmark\\checkmarkc\checkmarki\checkmarkr\checkmarkc\checkmark$\checkmarkr\checkmark$\checkmark^\checkmark\\checkmarkc\checkmarki\checkmarkr\checkmarkc\checkmark$\checkmarke\checkmark$\checkmark^\checkmark\\checkmarkc\checkmarki\checkmarkr\checkmarkc\checkmark$\checkmark}\checkmark$\checkmark^\checkmark\\checkmarkc\checkmarki\checkmarkr\checkmarkc\checkmark$\checkmark \checkmark$\checkmark^\checkmark\\checkmarkc\checkmarki\checkmarkr\checkmarkc\checkmark$\checkmark=\checkmark$\checkmark^\checkmark\\checkmarkc\checkmarki\checkmarkr\checkmarkc\checkmark$\checkmark \checkmark$\checkmark^\checkmark\\checkmarkc\checkmarki\checkmarkr\checkmarkc\checkmark$\checkmark\\checkmark$\checkmark^\checkmark\\checkmarkc\checkmarki\checkmarkr\checkmarkc\checkmark$\checkmarkb\checkmark$\checkmark^\checkmark\\checkmarkc\checkmarki\checkmarkr\checkmarkc\checkmark$\checkmarko\checkmark$\checkmark^\checkmark\\checkmarkc\checkmarki\checkmarkr\checkmarkc\checkmark$\checkmarkx\checkmark$\checkmark^\checkmark\\checkmarkc\checkmarki\checkmarkr\checkmarkc\checkmark$\checkmarke\checkmark$\checkmark^\checkmark\\checkmarkc\checkmarki\checkmarkr\checkmarkc\checkmark$\checkmarkd\checkmark$\checkmark^\checkmark\\checkmarkc\checkmarki\checkmarkr\checkmarkc\checkmark$\checkmark{\checkmark$\checkmark^\checkmark\\checkmarkc\checkmarki\checkmarkr\checkmarkc\checkmark$\checkmark2\checkmark$\checkmark^\checkmark\\checkmarkc\checkmarki\checkmarkr\checkmarkc\checkmark$\checkmark0\checkmark$\checkmark^\checkmark\\checkmarkc\checkmarki\checkmarkr\checkmarkc\checkmark$\checkmark \checkmark$\checkmark^\checkmark\\checkmarkc\checkmarki\checkmarkr\checkmarkc\checkmark$\checkmark\\checkmark$\checkmark^\checkmark\\checkmarkc\checkmarki\checkmarkr\checkmarkc\checkmark$\checkmarkt\checkmark$\checkmark^\checkmark\\checkmarkc\checkmarki\checkmarkr\checkmarkc\checkmark$\checkmarke\checkmark$\checkmark^\checkmark\\checkmarkc\checkmarki\checkmarkr\checkmarkc\checkmark$\checkmarkx\checkmark$\checkmark^\checkmark\\checkmarkc\checkmarki\checkmarkr\checkmarkc\checkmark$\checkmarkt\checkmark$\checkmark^\checkmark\\checkmarkc\checkmarki\checkmarkr\checkmarkc\checkmark$\checkmark{\checkmark$\checkmark^\checkmark\\checkmarkc\checkmarki\checkmarkr\checkmarkc\checkmark$\checkmark \checkmark$\checkmark^\checkmark\\checkmarkc\checkmarki\checkmarkr\checkmarkc\checkmark$\checkmarkm\checkmark$\checkmark^\checkmark\\checkmarkc\checkmarki\checkmarkr\checkmarkc\checkmark$\checkmarkm\checkmark$\checkmark^\checkmark\\checkmarkc\checkmarki\checkmarkr\checkmarkc\checkmark$\checkmark}\checkmark$\checkmark^\checkmark\\checkmarkc\checkmarki\checkmarkr\checkmarkc\checkmark$\checkmark}\checkmark$\checkmark^\checkmark\\checkmarkc\checkmarki\checkmarkr\checkmarkc\checkmark$\checkmark$\checkmark$\checkmark^\checkmark\\checkmarkc\checkmarki\checkmarkr\checkmarkc\checkmark$\checkmark.\checkmark$\checkmark^\checkmark\\checkmarkc\checkmarki\checkmarkr\checkmarkc\checkmark$\checkmark
\checkmark$\checkmark^\checkmark\\checkmarkc\checkmarki\checkmarkr\checkmarkc\checkmark$\checkmark
\checkmark$\checkmark^\checkmark\\checkmarkc\checkmarki\checkmarkr\checkmarkc\checkmark$\checkmarkC\checkmark$\checkmark^\checkmark\\checkmarkc\checkmarki\checkmarkr\checkmarkc\checkmark$\checkmarke\checkmark$\checkmark^\checkmark\\checkmarkc\checkmarki\checkmarkr\checkmarkc\checkmark$\checkmark \checkmark$\checkmark^\checkmark\\checkmarkc\checkmarki\checkmarkr\checkmarkc\checkmark$\checkmarkd\checkmark$\checkmark^\checkmark\\checkmarkc\checkmarki\checkmarkr\checkmarkc\checkmark$\checkmarki\checkmark$\checkmark^\checkmark\\checkmarkc\checkmarki\checkmarkr\checkmarkc\checkmark$\checkmarka\checkmark$\checkmark^\checkmark\\checkmarkc\checkmarki\checkmarkr\checkmarkc\checkmark$\checkmarkm\checkmark$\checkmark^\checkmark\\checkmarkc\checkmarki\checkmarkr\checkmarkc\checkmark$\checkmarkè\checkmark$\checkmark^\checkmark\\checkmarkc\checkmarki\checkmarkr\checkmarkc\checkmark$\checkmarkt\checkmark$\checkmark^\checkmark\\checkmarkc\checkmarki\checkmarkr\checkmarkc\checkmark$\checkmarkr\checkmark$\checkmark^\checkmark\\checkmarkc\checkmarki\checkmarkr\checkmarkc\checkmark$\checkmarke\checkmark$\checkmark^\checkmark\\checkmarkc\checkmarki\checkmarkr\checkmarkc\checkmark$\checkmark \checkmark$\checkmark^\checkmark\\checkmarkc\checkmarki\checkmarkr\checkmarkc\checkmark$\checkmarkd\checkmark$\checkmark^\checkmark\\checkmarkc\checkmarki\checkmarkr\checkmarkc\checkmark$\checkmarko\checkmark$\checkmark^\checkmark\\checkmarkc\checkmarki\checkmarkr\checkmarkc\checkmark$\checkmarkn\checkmark$\checkmark^\checkmark\\checkmarkc\checkmarki\checkmarkr\checkmarkc\checkmark$\checkmarkn\checkmark$\checkmark^\checkmark\\checkmarkc\checkmarki\checkmarkr\checkmarkc\checkmark$\checkmarke\checkmark$\checkmark^\checkmark\\checkmarkc\checkmarki\checkmarkr\checkmarkc\checkmark$\checkmark \checkmark$\checkmark^\checkmark\\checkmarkc\checkmarki\checkmarkr\checkmarkc\checkmark$\checkmarku\checkmark$\checkmark^\checkmark\\checkmarkc\checkmarki\checkmarkr\checkmarkc\checkmark$\checkmarkn\checkmark$\checkmark^\checkmark\\checkmarkc\checkmarki\checkmarkr\checkmarkc\checkmark$\checkmark \checkmark$\checkmark^\checkmark\\checkmarkc\checkmarki\checkmarkr\checkmarkc\checkmark$\checkmarkc\checkmark$\checkmark^\checkmark\\checkmarkc\checkmarki\checkmarkr\checkmarkc\checkmark$\checkmarko\checkmark$\checkmark^\checkmark\\checkmarkc\checkmarki\checkmarkr\checkmarkc\checkmark$\checkmarke\checkmark$\checkmark^\checkmark\\checkmarkc\checkmarki\checkmarkr\checkmarkc\checkmark$\checkmarkf\checkmark$\checkmark^\checkmark\\checkmarkc\checkmarki\checkmarkr\checkmarkc\checkmark$\checkmarkf\checkmark$\checkmark^\checkmark\\checkmarkc\checkmarki\checkmarkr\checkmarkc\checkmark$\checkmarki\checkmark$\checkmark^\checkmark\\checkmarkc\checkmarki\checkmarkr\checkmarkc\checkmark$\checkmarkc\checkmark$\checkmark^\checkmark\\checkmarkc\checkmarki\checkmarkr\checkmarkc\checkmark$\checkmarki\checkmark$\checkmark^\checkmark\\checkmarkc\checkmarki\checkmarkr\checkmarkc\checkmark$\checkmarke\checkmark$\checkmark^\checkmark\\checkmarkc\checkmarki\checkmarkr\checkmarkc\checkmark$\checkmarkn\checkmark$\checkmark^\checkmark\\checkmarkc\checkmarki\checkmarkr\checkmarkc\checkmark$\checkmarkt\checkmark$\checkmark^\checkmark\\checkmarkc\checkmarki\checkmarkr\checkmarkc\checkmark$\checkmark \checkmark$\checkmark^\checkmark\\checkmarkc\checkmarki\checkmarkr\checkmarkc\checkmark$\checkmarkd\checkmark$\checkmark^\checkmark\\checkmarkc\checkmarki\checkmarkr\checkmarkc\checkmark$\checkmarke\checkmark$\checkmark^\checkmark\\checkmarkc\checkmarki\checkmarkr\checkmarkc\checkmark$\checkmark \checkmark$\checkmark^\checkmark\\checkmarkc\checkmarki\checkmarkr\checkmarkc\checkmark$\checkmarks\checkmark$\checkmark^\checkmark\\checkmarkc\checkmarki\checkmarkr\checkmarkc\checkmark$\checkmarké\checkmark$\checkmark^\checkmark\\checkmarkc\checkmarki\checkmarkr\checkmarkc\checkmark$\checkmarkc\checkmark$\checkmark^\checkmark\\checkmarkc\checkmarki\checkmarkr\checkmarkc\checkmark$\checkmarku\checkmark$\checkmark^\checkmark\\checkmarkc\checkmarki\checkmarkr\checkmarkc\checkmark$\checkmarkr\checkmark$\checkmark^\checkmark\\checkmarkc\checkmarki\checkmarkr\checkmarkc\checkmark$\checkmarki\checkmark$\checkmark^\checkmark\\checkmarkc\checkmarki\checkmarkr\checkmarkc\checkmark$\checkmarkt\checkmark$\checkmark^\checkmark\\checkmarkc\checkmarki\checkmarkr\checkmarkc\checkmark$\checkmarké\checkmark$\checkmark^\checkmark\\checkmarkc\checkmarki\checkmarkr\checkmarkc\checkmark$\checkmark \checkmark$\checkmark^\checkmark\\checkmarkc\checkmarki\checkmarkr\checkmarkc\checkmark$\checkmarkr\checkmark$\checkmark^\checkmark\\checkmarkc\checkmarki\checkmarkr\checkmarkc\checkmark$\checkmarké\checkmark$\checkmark^\checkmark\\checkmarkc\checkmarki\checkmarkr\checkmarkc\checkmark$\checkmarke\checkmark$\checkmark^\checkmark\\checkmarkc\checkmarki\checkmarkr\checkmarkc\checkmark$\checkmarkl\checkmark$\checkmark^\checkmark\\checkmarkc\checkmarki\checkmarkr\checkmarkc\checkmark$\checkmark \checkmark$\checkmark^\checkmark\\checkmarkc\checkmarki\checkmarkr\checkmarkc\checkmark$\checkmark:\checkmark$\checkmark^\checkmark\\checkmarkc\checkmarki\checkmarkr\checkmarkc\checkmark$\checkmark
\checkmark$\checkmark^\checkmark\\checkmarkc\checkmarki\checkmarkr\checkmarkc\checkmark$\checkmark\\checkmark$\checkmark^\checkmark\\checkmarkc\checkmarki\checkmarkr\checkmarkc\checkmark$\checkmarkb\checkmark$\checkmark^\checkmark\\checkmarkc\checkmarki\checkmarkr\checkmarkc\checkmark$\checkmarke\checkmark$\checkmark^\checkmark\\checkmarkc\checkmarki\checkmarkr\checkmarkc\checkmark$\checkmarkg\checkmark$\checkmark^\checkmark\\checkmarkc\checkmarki\checkmarkr\checkmarkc\checkmark$\checkmarki\checkmark$\checkmark^\checkmark\\checkmarkc\checkmarki\checkmarkr\checkmarkc\checkmark$\checkmarkn\checkmark$\checkmark^\checkmark\\checkmarkc\checkmarki\checkmarkr\checkmarkc\checkmark$\checkmark{\checkmark$\checkmark^\checkmark\\checkmarkc\checkmarki\checkmarkr\checkmarkc\checkmark$\checkmarke\checkmark$\checkmark^\checkmark\\checkmarkc\checkmarki\checkmarkr\checkmarkc\checkmark$\checkmarkq\checkmark$\checkmark^\checkmark\\checkmarkc\checkmarki\checkmarkr\checkmarkc\checkmark$\checkmarku\checkmark$\checkmark^\checkmark\\checkmarkc\checkmarki\checkmarkr\checkmarkc\checkmark$\checkmarka\checkmark$\checkmark^\checkmark\\checkmarkc\checkmarki\checkmarkr\checkmarkc\checkmark$\checkmarkt\checkmark$\checkmark^\checkmark\\checkmarkc\checkmarki\checkmarkr\checkmarkc\checkmark$\checkmarki\checkmark$\checkmark^\checkmark\\checkmarkc\checkmarki\checkmarkr\checkmarkc\checkmark$\checkmarko\checkmark$\checkmark^\checkmark\\checkmarkc\checkmarki\checkmarkr\checkmarkc\checkmark$\checkmarkn\checkmark$\checkmark^\checkmark\\checkmarkc\checkmarki\checkmarkr\checkmarkc\checkmark$\checkmark}\checkmark$\checkmark^\checkmark\\checkmarkc\checkmarki\checkmarkr\checkmarkc\checkmark$\checkmark
\checkmark$\checkmark^\checkmark\\checkmarkc\checkmarki\checkmarkr\checkmarkc\checkmark$\checkmark \checkmark$\checkmark^\checkmark\\checkmarkc\checkmarki\checkmarkr\checkmarkc\checkmark$\checkmark \checkmark$\checkmark^\checkmark\\checkmarkc\checkmarki\checkmarkr\checkmarkc\checkmark$\checkmark \checkmark$\checkmark^\checkmark\\checkmarkc\checkmarki\checkmarkr\checkmarkc\checkmark$\checkmark \checkmark$\checkmark^\checkmark\\checkmarkc\checkmarki\checkmarkr\checkmarkc\checkmark$\checkmarks\checkmark$\checkmark^\checkmark\\checkmarkc\checkmarki\checkmarkr\checkmarkc\checkmark$\checkmark_\checkmark$\checkmark^\checkmark\\checkmarkc\checkmarki\checkmarkr\checkmarkc\checkmark$\checkmark{\checkmark$\checkmark^\checkmark\\checkmarkc\checkmarki\checkmarkr\checkmarkc\checkmark$\checkmarkr\checkmark$\checkmark^\checkmark\\checkmarkc\checkmarki\checkmarkr\checkmarkc\checkmark$\checkmarke\checkmark$\checkmark^\checkmark\\checkmarkc\checkmarki\checkmarkr\checkmarkc\checkmark$\checkmarke\checkmark$\checkmark^\checkmark\\checkmarkc\checkmarki\checkmarkr\checkmarkc\checkmark$\checkmarkl\checkmark$\checkmark^\checkmark\\checkmarkc\checkmarki\checkmarkr\checkmarkc\checkmark$\checkmark}\checkmark$\checkmark^\checkmark\\checkmarkc\checkmarki\checkmarkr\checkmarkc\checkmark$\checkmark \checkmark$\checkmark^\checkmark\\checkmarkc\checkmarki\checkmarkr\checkmarkc\checkmark$\checkmark=\checkmark$\checkmark^\checkmark\\checkmarkc\checkmarki\checkmarkr\checkmarkc\checkmark$\checkmark \checkmark$\checkmark^\checkmark\\checkmarkc\checkmarki\checkmarkr\checkmarkc\checkmark$\checkmark\\checkmark$\checkmark^\checkmark\\checkmarkc\checkmarki\checkmarkr\checkmarkc\checkmark$\checkmarkf\checkmark$\checkmark^\checkmark\\checkmarkc\checkmarki\checkmarkr\checkmarkc\checkmark$\checkmarkr\checkmark$\checkmark^\checkmark\\checkmarkc\checkmarki\checkmarkr\checkmarkc\checkmark$\checkmarka\checkmark$\checkmark^\checkmark\\checkmarkc\checkmarki\checkmarkr\checkmarkc\checkmark$\checkmarkc\checkmark$\checkmark^\checkmark\\checkmarkc\checkmarki\checkmarkr\checkmarkc\checkmark$\checkmark{\checkmark$\checkmark^\checkmark\\checkmarkc\checkmarki\checkmarkr\checkmarkc\checkmark$\checkmarkR\checkmark$\checkmark^\checkmark\\checkmarkc\checkmarki\checkmarkr\checkmarkc\checkmark$\checkmark_\checkmark$\checkmark^\checkmark\\checkmarkc\checkmarki\checkmarkr\checkmarkc\checkmark$\checkmarke\checkmark$\checkmark^\checkmark\\checkmarkc\checkmarki\checkmarkr\checkmarkc\checkmark$\checkmark}\checkmark$\checkmark^\checkmark\\checkmarkc\checkmarki\checkmarkr\checkmarkc\checkmark$\checkmark{\checkmark$\checkmark^\checkmark\\checkmarkc\checkmarki\checkmarkr\checkmarkc\checkmark$\checkmark\\checkmark$\checkmark^\checkmark\\checkmarkc\checkmarki\checkmarkr\checkmarkc\checkmark$\checkmarks\checkmark$\checkmark^\checkmark\\checkmarkc\checkmarki\checkmarkr\checkmarkc\checkmark$\checkmarki\checkmark$\checkmark^\checkmark\\checkmarkc\checkmarki\checkmarkr\checkmarkc\checkmark$\checkmarkg\checkmark$\checkmark^\checkmark\\checkmarkc\checkmarki\checkmarkr\checkmarkc\checkmark$\checkmarkm\checkmark$\checkmark^\checkmark\\checkmarkc\checkmarki\checkmarkr\checkmarkc\checkmark$\checkmarka\checkmark$\checkmark^\checkmark\\checkmarkc\checkmarki\checkmarkr\checkmarkc\checkmark$\checkmark_\checkmark$\checkmark^\checkmark\\checkmarkc\checkmarki\checkmarkr\checkmarkc\checkmark$\checkmark{\checkmark$\checkmark^\checkmark\\checkmarkc\checkmarki\checkmarkr\checkmarkc\checkmark$\checkmarkf\checkmark$\checkmark^\checkmark\\checkmarkc\checkmarki\checkmarkr\checkmarkc\checkmark$\checkmarkl\checkmark$\checkmark^\checkmark\\checkmarkc\checkmarki\checkmarkr\checkmarkc\checkmark$\checkmarke\checkmark$\checkmark^\checkmark\\checkmarkc\checkmarki\checkmarkr\checkmarkc\checkmark$\checkmarkx\checkmark$\checkmark^\checkmark\\checkmarkc\checkmarki\checkmarkr\checkmarkc\checkmark$\checkmark,\checkmark$\checkmark^\checkmark\\checkmarkc\checkmarki\checkmarkr\checkmarkc\checkmark$\checkmarkr\checkmark$\checkmark^\checkmark\\checkmarkc\checkmarki\checkmarkr\checkmarkc\checkmark$\checkmarke\checkmark$\checkmark^\checkmark\\checkmarkc\checkmarki\checkmarkr\checkmarkc\checkmark$\checkmarke\checkmark$\checkmark^\checkmark\\checkmarkc\checkmarki\checkmarkr\checkmarkc\checkmark$\checkmarkl\checkmark$\checkmark^\checkmark\\checkmarkc\checkmarki\checkmarkr\checkmarkc\checkmark$\checkmark}\checkmark$\checkmark^\checkmark\\checkmarkc\checkmarki\checkmarkr\checkmarkc\checkmark$\checkmark}\checkmark$\checkmark^\checkmark\\checkmarkc\checkmarki\checkmarkr\checkmarkc\checkmark$\checkmark \checkmark$\checkmark^\checkmark\\checkmarkc\checkmarki\checkmarkr\checkmarkc\checkmark$\checkmark=\checkmark$\checkmark^\checkmark\\checkmarkc\checkmarki\checkmarkr\checkmarkc\checkmark$\checkmark \checkmark$\checkmark^\checkmark\\checkmarkc\checkmarki\checkmarkr\checkmarkc\checkmark$\checkmark\\checkmark$\checkmark^\checkmark\\checkmarkc\checkmarki\checkmarkr\checkmarkc\checkmark$\checkmarkf\checkmark$\checkmark^\checkmark\\checkmarkc\checkmarki\checkmarkr\checkmarkc\checkmark$\checkmarkr\checkmark$\checkmark^\checkmark\\checkmarkc\checkmarki\checkmarkr\checkmarkc\checkmark$\checkmarka\checkmark$\checkmark^\checkmark\\checkmarkc\checkmarki\checkmarkr\checkmarkc\checkmark$\checkmarkc\checkmark$\checkmark^\checkmark\\checkmarkc\checkmarki\checkmarkr\checkmarkc\checkmark$\checkmark{\checkmark$\checkmark^\checkmark\\checkmarkc\checkmarki\checkmarkr\checkmarkc\checkmark$\checkmark3\checkmark$\checkmark^\checkmark\\checkmarkc\checkmarki\checkmarkr\checkmarkc\checkmark$\checkmark5\checkmark$\checkmark^\checkmark\\checkmarkc\checkmarki\checkmarkr\checkmarkc\checkmark$\checkmark0\checkmark$\checkmark^\checkmark\\checkmarkc\checkmarki\checkmarkr\checkmarkc\checkmark$\checkmark}\checkmark$\checkmark^\checkmark\\checkmarkc\checkmarki\checkmarkr\checkmarkc\checkmark$\checkmark{\checkmark$\checkmark^\checkmark\\checkmarkc\checkmarki\checkmarkr\checkmarkc\checkmark$\checkmark\\checkmark$\checkmark^\checkmark\\checkmarkc\checkmarki\checkmarkr\checkmarkc\checkmark$\checkmarkf\checkmark$\checkmark^\checkmark\\checkmarkc\checkmarki\checkmarkr\checkmarkc\checkmark$\checkmarkr\checkmark$\checkmark^\checkmark\\checkmarkc\checkmarki\checkmarkr\checkmarkc\checkmark$\checkmarka\checkmark$\checkmark^\checkmark\\checkmarkc\checkmarki\checkmarkr\checkmarkc\checkmark$\checkmarkc\checkmark$\checkmark^\checkmark\\checkmarkc\checkmarki\checkmarkr\checkmarkc\checkmark$\checkmark{\checkmark$\checkmark^\checkmark\\checkmarkc\checkmarki\checkmarkr\checkmarkc\checkmark$\checkmark3\checkmark$\checkmark^\checkmark\\checkmarkc\checkmarki\checkmarkr\checkmarkc\checkmark$\checkmark2\checkmark$\checkmark^\checkmark\\checkmarkc\checkmarki\checkmarkr\checkmarkc\checkmark$\checkmark \checkmark$\checkmark^\checkmark\\checkmarkc\checkmarki\checkmarkr\checkmarkc\checkmark$\checkmark\\checkmark$\checkmark^\checkmark\\checkmarkc\checkmarki\checkmarkr\checkmarkc\checkmark$\checkmarkt\checkmark$\checkmark^\checkmark\\checkmarkc\checkmarki\checkmarkr\checkmarkc\checkmark$\checkmarki\checkmark$\checkmark^\checkmark\\checkmarkc\checkmarki\checkmarkr\checkmarkc\checkmark$\checkmarkm\checkmark$\checkmark^\checkmark\\checkmarkc\checkmarki\checkmarkr\checkmarkc\checkmark$\checkmarke\checkmark$\checkmark^\checkmark\\checkmarkc\checkmarki\checkmarkr\checkmarkc\checkmark$\checkmarks\checkmark$\checkmark^\checkmark\\checkmarkc\checkmarki\checkmarkr\checkmarkc\checkmark$\checkmark \checkmark$\checkmark^\checkmark\\checkmarkc\checkmarki\checkmarkr\checkmarkc\checkmark$\checkmark1\checkmark$\checkmark^\checkmark\\checkmarkc\checkmarki\checkmarkr\checkmarkc\checkmark$\checkmark1\checkmark$\checkmark^\checkmark\\checkmarkc\checkmarki\checkmarkr\checkmarkc\checkmark$\checkmark\\checkmark$\checkmark^\checkmark\\checkmarkc\checkmarki\checkmarkr\checkmarkc\checkmark$\checkmark,\checkmark$\checkmark^\checkmark\\checkmarkc\checkmarki\checkmarkr\checkmarkc\checkmark$\checkmark2\checkmark$\checkmark^\checkmark\\checkmarkc\checkmarki\checkmarkr\checkmarkc\checkmark$\checkmark5\checkmark$\checkmark^\checkmark\\checkmarkc\checkmarki\checkmarkr\checkmarkc\checkmark$\checkmark0\checkmark$\checkmark^\checkmark\\checkmarkc\checkmarki\checkmarkr\checkmarkc\checkmark$\checkmark}\checkmark$\checkmark^\checkmark\\checkmarkc\checkmarki\checkmarkr\checkmarkc\checkmark$\checkmark{\checkmark$\checkmark^\checkmark\\checkmarkc\checkmarki\checkmarkr\checkmarkc\checkmark$\checkmark\\checkmark$\checkmark^\checkmark\\checkmarkc\checkmarki\checkmarkr\checkmarkc\checkmark$\checkmarkp\checkmark$\checkmark^\checkmark\\checkmarkc\checkmarki\checkmarkr\checkmarkc\checkmark$\checkmarki\checkmark$\checkmark^\checkmark\\checkmarkc\checkmarki\checkmarkr\checkmarkc\checkmark$\checkmark \checkmark$\checkmark^\checkmark\\checkmarkc\checkmarki\checkmarkr\checkmarkc\checkmark$\checkmark\\checkmark$\checkmark^\checkmark\\checkmarkc\checkmarki\checkmarkr\checkmarkc\checkmark$\checkmarkt\checkmark$\checkmark^\checkmark\\checkmarkc\checkmarki\checkmarkr\checkmarkc\checkmark$\checkmarki\checkmark$\checkmark^\checkmark\\checkmarkc\checkmarki\checkmarkr\checkmarkc\checkmark$\checkmarkm\checkmark$\checkmark^\checkmark\\checkmarkc\checkmarki\checkmarkr\checkmarkc\checkmark$\checkmarke\checkmark$\checkmark^\checkmark\\checkmarkc\checkmarki\checkmarkr\checkmarkc\checkmark$\checkmarks\checkmark$\checkmark^\checkmark\\checkmarkc\checkmarki\checkmarkr\checkmarkc\checkmark$\checkmark \checkmark$\checkmark^\checkmark\\checkmarkc\checkmarki\checkmarkr\checkmarkc\checkmark$\checkmark2\checkmark$\checkmark^\checkmark\\checkmarkc\checkmarki\checkmarkr\checkmarkc\checkmark$\checkmark0\checkmark$\checkmark^\checkmark\\checkmarkc\checkmarki\checkmarkr\checkmarkc\checkmark$\checkmark^\checkmark$\checkmark^\checkmark\\checkmarkc\checkmarki\checkmarkr\checkmarkc\checkmark$\checkmark3\checkmark$\checkmark^\checkmark\\checkmarkc\checkmarki\checkmarkr\checkmarkc\checkmark$\checkmark}\checkmark$\checkmark^\checkmark\\checkmarkc\checkmarki\checkmarkr\checkmarkc\checkmark$\checkmark}\checkmark$\checkmark^\checkmark\\checkmarkc\checkmarki\checkmarkr\checkmarkc\checkmark$\checkmark \checkmark$\checkmark^\checkmark\\checkmarkc\checkmarki\checkmarkr\checkmarkc\checkmark$\checkmark=\checkmark$\checkmark^\checkmark\\checkmarkc\checkmarki\checkmarkr\checkmarkc\checkmark$\checkmark \checkmark$\checkmark^\checkmark\\checkmarkc\checkmarki\checkmarkr\checkmarkc\checkmark$\checkmark\\checkmark$\checkmark^\checkmark\\checkmarkc\checkmarki\checkmarkr\checkmarkc\checkmark$\checkmarkf\checkmark$\checkmark^\checkmark\\checkmarkc\checkmarki\checkmarkr\checkmarkc\checkmark$\checkmarkr\checkmark$\checkmark^\checkmark\\checkmarkc\checkmarki\checkmarkr\checkmarkc\checkmark$\checkmarka\checkmark$\checkmark^\checkmark\\checkmarkc\checkmarki\checkmarkr\checkmarkc\checkmark$\checkmarkc\checkmark$\checkmark^\checkmark\\checkmarkc\checkmarki\checkmarkr\checkmarkc\checkmark$\checkmark{\checkmark$\checkmark^\checkmark\\checkmarkc\checkmarki\checkmarkr\checkmarkc\checkmark$\checkmark3\checkmark$\checkmark^\checkmark\\checkmarkc\checkmarki\checkmarkr\checkmarkc\checkmark$\checkmark5\checkmark$\checkmark^\checkmark\\checkmarkc\checkmarki\checkmarkr\checkmarkc\checkmark$\checkmark0\checkmark$\checkmark^\checkmark\\checkmarkc\checkmarki\checkmarkr\checkmarkc\checkmark$\checkmark}\checkmark$\checkmark^\checkmark\\checkmarkc\checkmarki\checkmarkr\checkmarkc\checkmark$\checkmark{\checkmark$\checkmark^\checkmark\\checkmarkc\checkmarki\checkmarkr\checkmarkc\checkmark$\checkmark\\checkmark$\checkmark^\checkmark\\checkmarkc\checkmarki\checkmarkr\checkmarkc\checkmark$\checkmarkf\checkmark$\checkmark^\checkmark\\checkmarkc\checkmarki\checkmarkr\checkmarkc\checkmark$\checkmarkr\checkmark$\checkmark^\checkmark\\checkmarkc\checkmarki\checkmarkr\checkmarkc\checkmark$\checkmarka\checkmark$\checkmark^\checkmark\\checkmarkc\checkmarki\checkmarkr\checkmarkc\checkmark$\checkmarkc\checkmark$\checkmark^\checkmark\\checkmarkc\checkmarki\checkmarkr\checkmarkc\checkmark$\checkmark{\checkmark$\checkmark^\checkmark\\checkmarkc\checkmarki\checkmarkr\checkmarkc\checkmark$\checkmark3\checkmark$\checkmark^\checkmark\\checkmarkc\checkmarki\checkmarkr\checkmarkc\checkmark$\checkmark6\checkmark$\checkmark^\checkmark\\checkmarkc\checkmarki\checkmarkr\checkmarkc\checkmark$\checkmark0\checkmark$\checkmark^\checkmark\\checkmarkc\checkmarki\checkmarkr\checkmarkc\checkmark$\checkmark\\checkmark$\checkmark^\checkmark\\checkmarkc\checkmarki\checkmarkr\checkmarkc\checkmark$\checkmark,\checkmark$\checkmark^\checkmark\\checkmarkc\checkmarki\checkmarkr\checkmarkc\checkmark$\checkmark0\checkmark$\checkmark^\checkmark\\checkmarkc\checkmarki\checkmarkr\checkmarkc\checkmark$\checkmark0\checkmark$\checkmark^\checkmark\\checkmarkc\checkmarki\checkmarkr\checkmarkc\checkmark$\checkmark0\checkmark$\checkmark^\checkmark\\checkmarkc\checkmarki\checkmarkr\checkmarkc\checkmark$\checkmark}\checkmark$\checkmark^\checkmark\\checkmarkc\checkmarki\checkmarkr\checkmarkc\checkmark$\checkmark{\checkmark$\checkmark^\checkmark\\checkmarkc\checkmarki\checkmarkr\checkmarkc\checkmark$\checkmark2\checkmark$\checkmark^\checkmark\\checkmarkc\checkmarki\checkmarkr\checkmarkc\checkmark$\checkmark5\checkmark$\checkmark^\checkmark\\checkmarkc\checkmarki\checkmarkr\checkmarkc\checkmark$\checkmark\\checkmark$\checkmark^\checkmark\\checkmarkc\checkmarki\checkmarkr\checkmarkc\checkmark$\checkmark,\checkmark$\checkmark^\checkmark\\checkmarkc\checkmarki\checkmarkr\checkmarkc\checkmark$\checkmark1\checkmark$\checkmark^\checkmark\\checkmarkc\checkmarki\checkmarkr\checkmarkc\checkmark$\checkmark3\checkmark$\checkmark^\checkmark\\checkmarkc\checkmarki\checkmarkr\checkmarkc\checkmark$\checkmark3\checkmark$\checkmark^\checkmark\\checkmarkc\checkmarki\checkmarkr\checkmarkc\checkmark$\checkmark}\checkmark$\checkmark^\checkmark\\checkmarkc\checkmarki\checkmarkr\checkmarkc\checkmark$\checkmark}\checkmark$\checkmark^\checkmark\\checkmarkc\checkmarki\checkmarkr\checkmarkc\checkmark$\checkmark \checkmark$\checkmark^\checkmark\\checkmarkc\checkmarki\checkmarkr\checkmarkc\checkmark$\checkmark=\checkmark$\checkmark^\checkmark\\checkmarkc\checkmarki\checkmarkr\checkmarkc\checkmark$\checkmark \checkmark$\checkmark^\checkmark\\checkmarkc\checkmarki\checkmarkr\checkmarkc\checkmark$\checkmark\\checkmark$\checkmark^\checkmark\\checkmarkc\checkmarki\checkmarkr\checkmarkc\checkmark$\checkmarkf\checkmark$\checkmark^\checkmark\\checkmarkc\checkmarki\checkmarkr\checkmarkc\checkmark$\checkmarkr\checkmark$\checkmark^\checkmark\\checkmarkc\checkmarki\checkmarkr\checkmarkc\checkmark$\checkmarka\checkmark$\checkmark^\checkmark\\checkmarkc\checkmarki\checkmarkr\checkmarkc\checkmark$\checkmarkc\checkmark$\checkmark^\checkmark\\checkmarkc\checkmarki\checkmarkr\checkmarkc\checkmark$\checkmark{\checkmark$\checkmark^\checkmark\\checkmarkc\checkmarki\checkmarkr\checkmarkc\checkmark$\checkmark3\checkmark$\checkmark^\checkmark\\checkmarkc\checkmarki\checkmarkr\checkmarkc\checkmark$\checkmark5\checkmark$\checkmark^\checkmark\\checkmarkc\checkmarki\checkmarkr\checkmarkc\checkmark$\checkmark0\checkmark$\checkmark^\checkmark\\checkmarkc\checkmarki\checkmarkr\checkmarkc\checkmark$\checkmark}\checkmark$\checkmark^\checkmark\\checkmarkc\checkmarki\checkmarkr\checkmarkc\checkmark$\checkmark{\checkmark$\checkmark^\checkmark\\checkmarkc\checkmarki\checkmarkr\checkmarkc\checkmark$\checkmark1\checkmark$\checkmark^\checkmark\\checkmarkc\checkmarki\checkmarkr\checkmarkc\checkmark$\checkmark4\checkmark$\checkmark^\checkmark\\checkmarkc\checkmarki\checkmarkr\checkmarkc\checkmark$\checkmark{\checkmark$\checkmark^\checkmark\\checkmarkc\checkmarki\checkmarkr\checkmarkc\checkmark$\checkmark,\checkmark$\checkmark^\checkmark\\checkmarkc\checkmarki\checkmarkr\checkmarkc\checkmark$\checkmark}\checkmark$\checkmark^\checkmark\\checkmarkc\checkmarki\checkmarkr\checkmarkc\checkmark$\checkmark3\checkmark$\checkmark^\checkmark\\checkmarkc\checkmarki\checkmarkr\checkmarkc\checkmark$\checkmark2\checkmark$\checkmark^\checkmark\\checkmarkc\checkmarki\checkmarkr\checkmarkc\checkmark$\checkmark}\checkmark$\checkmark^\checkmark\\checkmarkc\checkmarki\checkmarkr\checkmarkc\checkmark$\checkmark \checkmark$\checkmark^\checkmark\\checkmarkc\checkmarki\checkmarkr\checkmarkc\checkmark$\checkmark=\checkmark$\checkmark^\checkmark\\checkmarkc\checkmarki\checkmarkr\checkmarkc\checkmark$\checkmark \checkmark$\checkmark^\checkmark\\checkmarkc\checkmarki\checkmarkr\checkmarkc\checkmark$\checkmark\\checkmark$\checkmark^\checkmark\\checkmarkc\checkmarki\checkmarkr\checkmarkc\checkmark$\checkmarkb\checkmark$\checkmark^\checkmark\\checkmarkc\checkmarki\checkmarkr\checkmarkc\checkmark$\checkmarko\checkmark$\checkmark^\checkmark\\checkmarkc\checkmarki\checkmarkr\checkmarkc\checkmark$\checkmarkx\checkmark$\checkmark^\checkmark\\checkmarkc\checkmarki\checkmarkr\checkmarkc\checkmark$\checkmarke\checkmark$\checkmark^\checkmark\\checkmarkc\checkmarki\checkmarkr\checkmarkc\checkmark$\checkmarkd\checkmark$\checkmark^\checkmark\\checkmarkc\checkmarki\checkmarkr\checkmarkc\checkmark$\checkmark{\checkmark$\checkmark^\checkmark\\checkmarkc\checkmarki\checkmarkr\checkmarkc\checkmark$\checkmark2\checkmark$\checkmark^\checkmark\\checkmarkc\checkmarki\checkmarkr\checkmarkc\checkmark$\checkmark4\checkmark$\checkmark^\checkmark\\checkmarkc\checkmarki\checkmarkr\checkmarkc\checkmark$\checkmark{\checkmark$\checkmark^\checkmark\\checkmarkc\checkmarki\checkmarkr\checkmarkc\checkmark$\checkmark,\checkmark$\checkmark^\checkmark\\checkmarkc\checkmarki\checkmarkr\checkmarkc\checkmark$\checkmark}\checkmark$\checkmark^\checkmark\\checkmarkc\checkmarki\checkmarkr\checkmarkc\checkmark$\checkmark4\checkmark$\checkmark^\checkmark\\checkmarkc\checkmarki\checkmarkr\checkmarkc\checkmark$\checkmark \checkmark$\checkmark^\checkmark\\checkmarkc\checkmarki\checkmarkr\checkmarkc\checkmark$\checkmark\\checkmark$\checkmark^\checkmark\\checkmarkc\checkmarki\checkmarkr\checkmarkc\checkmark$\checkmarkg\checkmark$\checkmark^\checkmark\\checkmarkc\checkmarki\checkmarkr\checkmarkc\checkmark$\checkmarkg\checkmark$\checkmark^\checkmark\\checkmarkc\checkmarki\checkmarkr\checkmarkc\checkmark$\checkmark \checkmark$\checkmark^\checkmark\\checkmarkc\checkmarki\checkmarkr\checkmarkc\checkmark$\checkmark1\checkmark$\checkmark^\checkmark\\checkmarkc\checkmarki\checkmarkr\checkmarkc\checkmark$\checkmark{\checkmark$\checkmark^\checkmark\\checkmarkc\checkmarki\checkmarkr\checkmarkc\checkmark$\checkmark,\checkmark$\checkmark^\checkmark\\checkmarkc\checkmarki\checkmarkr\checkmarkc\checkmark$\checkmark}\checkmark$\checkmark^\checkmark\\checkmarkc\checkmarki\checkmarkr\checkmarkc\checkmark$\checkmark5\checkmark$\checkmark^\checkmark\\checkmarkc\checkmarki\checkmarkr\checkmarkc\checkmark$\checkmark}\checkmark$\checkmark^\checkmark\\checkmarkc\checkmarki\checkmarkr\checkmarkc\checkmark$\checkmark
\checkmark$\checkmark^\checkmark\\checkmarkc\checkmarki\checkmarkr\checkmarkc\checkmark$\checkmark\\checkmark$\checkmark^\checkmark\\checkmarkc\checkmarki\checkmarkr\checkmarkc\checkmark$\checkmarke\checkmark$\checkmark^\checkmark\\checkmarkc\checkmarki\checkmarkr\checkmarkc\checkmark$\checkmarkn\checkmark$\checkmark^\checkmark\\checkmarkc\checkmarki\checkmarkr\checkmarkc\checkmark$\checkmarkd\checkmark$\checkmark^\checkmark\\checkmarkc\checkmarki\checkmarkr\checkmarkc\checkmark$\checkmark{\checkmark$\checkmark^\checkmark\\checkmarkc\checkmarki\checkmarkr\checkmarkc\checkmark$\checkmarke\checkmark$\checkmark^\checkmark\\checkmarkc\checkmarki\checkmarkr\checkmarkc\checkmark$\checkmarkq\checkmark$\checkmark^\checkmark\\checkmarkc\checkmarki\checkmarkr\checkmarkc\checkmark$\checkmarku\checkmark$\checkmark^\checkmark\\checkmarkc\checkmarki\checkmarkr\checkmarkc\checkmark$\checkmarka\checkmark$\checkmark^\checkmark\\checkmarkc\checkmarki\checkmarkr\checkmarkc\checkmark$\checkmarkt\checkmark$\checkmark^\checkmark\\checkmarkc\checkmarki\checkmarkr\checkmarkc\checkmark$\checkmarki\checkmark$\checkmark^\checkmark\\checkmarkc\checkmarki\checkmarkr\checkmarkc\checkmark$\checkmarko\checkmark$\checkmark^\checkmark\\checkmarkc\checkmarki\checkmarkr\checkmarkc\checkmark$\checkmarkn\checkmark$\checkmark^\checkmark\\checkmarkc\checkmarki\checkmarkr\checkmarkc\checkmark$\checkmark}\checkmark$\checkmark^\checkmark\\checkmarkc\checkmarki\checkmarkr\checkmarkc\checkmark$\checkmark
\checkmark$\checkmark^\checkmark\\checkmarkc\checkmarki\checkmarkr\checkmarkc\checkmark$\checkmark
\checkmark$\checkmark^\checkmark\\checkmarkc\checkmarki\checkmarkr\checkmarkc\checkmark$\checkmark\\checkmark$\checkmark^\checkmark\\checkmarkc\checkmarki\checkmarkr\checkmarkc\checkmark$\checkmarkt\checkmark$\checkmark^\checkmark\\checkmarkc\checkmarki\checkmarkr\checkmarkc\checkmark$\checkmarke\checkmark$\checkmark^\checkmark\\checkmarkc\checkmarki\checkmarkr\checkmarkc\checkmark$\checkmarkx\checkmark$\checkmark^\checkmark\\checkmarkc\checkmarki\checkmarkr\checkmarkc\checkmark$\checkmarkt\checkmark$\checkmark^\checkmark\\checkmarkc\checkmarki\checkmarkr\checkmarkc\checkmark$\checkmarkb\checkmark$\checkmark^\checkmark\\checkmarkc\checkmarki\checkmarkr\checkmarkc\checkmark$\checkmarkf\checkmark$\checkmark^\checkmark\\checkmarkc\checkmarki\checkmarkr\checkmarkc\checkmark$\checkmark{\checkmark$\checkmark^\checkmark\\checkmarkc\checkmarki\checkmarkr\checkmarkc\checkmark$\checkmarkV\checkmark$\checkmark^\checkmark\\checkmarkc\checkmarki\checkmarkr\checkmarkc\checkmark$\checkmarké\checkmark$\checkmark^\checkmark\\checkmarkc\checkmarki\checkmarkr\checkmarkc\checkmark$\checkmarkr\checkmark$\checkmark^\checkmark\\checkmarkc\checkmarki\checkmarkr\checkmarkc\checkmark$\checkmarki\checkmark$\checkmark^\checkmark\\checkmarkc\checkmarki\checkmarkr\checkmarkc\checkmark$\checkmarkf\checkmark$\checkmark^\checkmark\\checkmarkc\checkmarki\checkmarkr\checkmarkc\checkmark$\checkmarki\checkmark$\checkmark^\checkmark\\checkmarkc\checkmarki\checkmarkr\checkmarkc\checkmark$\checkmarkc\checkmark$\checkmark^\checkmark\\checkmarkc\checkmarki\checkmarkr\checkmarkc\checkmark$\checkmarka\checkmark$\checkmark^\checkmark\\checkmarkc\checkmarki\checkmarkr\checkmarkc\checkmark$\checkmarkt\checkmark$\checkmark^\checkmark\\checkmarkc\checkmarki\checkmarkr\checkmarkc\checkmark$\checkmarki\checkmark$\checkmark^\checkmark\\checkmarkc\checkmarki\checkmarkr\checkmarkc\checkmark$\checkmarko\checkmark$\checkmark^\checkmark\\checkmarkc\checkmarki\checkmarkr\checkmarkc\checkmark$\checkmarkn\checkmark$\checkmark^\checkmark\\checkmarkc\checkmarki\checkmarkr\checkmarkc\checkmark$\checkmark \checkmark$\checkmark^\checkmark\\checkmarkc\checkmarki\checkmarkr\checkmarkc\checkmark$\checkmark:\checkmark$\checkmark^\checkmark\\checkmarkc\checkmarki\checkmarkr\checkmarkc\checkmark$\checkmark}\checkmark$\checkmark^\checkmark\\checkmarkc\checkmarki\checkmarkr\checkmarkc\checkmark$\checkmark \checkmark$\checkmark^\checkmark\\checkmarkc\checkmarki\checkmarkr\checkmarkc\checkmark$\checkmarkL\checkmark$\checkmark^\checkmark\\checkmarkc\checkmarki\checkmarkr\checkmarkc\checkmark$\checkmarke\checkmark$\checkmark^\checkmark\\checkmarkc\checkmarki\checkmarkr\checkmarkc\checkmark$\checkmark \checkmark$\checkmark^\checkmark\\checkmarkc\checkmarki\checkmarkr\checkmarkc\checkmark$\checkmarkd\checkmark$\checkmark^\checkmark\\checkmarkc\checkmarki\checkmarkr\checkmarkc\checkmark$\checkmarki\checkmark$\checkmark^\checkmark\\checkmarkc\checkmarki\checkmarkr\checkmarkc\checkmark$\checkmarka\checkmark$\checkmark^\checkmark\\checkmarkc\checkmarki\checkmarkr\checkmarkc\checkmark$\checkmarkm\checkmark$\checkmark^\checkmark\\checkmarkc\checkmarki\checkmarkr\checkmarkc\checkmark$\checkmarkè\checkmark$\checkmark^\checkmark\\checkmarkc\checkmarki\checkmarkr\checkmarkc\checkmark$\checkmarkt\checkmark$\checkmark^\checkmark\\checkmarkc\checkmarki\checkmarkr\checkmarkc\checkmark$\checkmarkr\checkmark$\checkmark^\checkmark\\checkmarkc\checkmarki\checkmarkr\checkmarkc\checkmark$\checkmarke\checkmark$\checkmark^\checkmark\\checkmarkc\checkmarki\checkmarkr\checkmarkc\checkmark$\checkmark \checkmark$\checkmark^\checkmark\\checkmarkc\checkmarki\checkmarkr\checkmarkc\checkmark$\checkmark$\checkmark$\checkmark^\checkmark\\checkmarkc\checkmarki\checkmarkr\checkmarkc\checkmark$\checkmarkd\checkmark$\checkmark^\checkmark\\checkmarkc\checkmarki\checkmarkr\checkmarkc\checkmark$\checkmark \checkmark$\checkmark^\checkmark\\checkmarkc\checkmarki\checkmarkr\checkmarkc\checkmark$\checkmark=\checkmark$\checkmark^\checkmark\\checkmarkc\checkmarki\checkmarkr\checkmarkc\checkmark$\checkmark \checkmark$\checkmark^\checkmark\\checkmarkc\checkmarki\checkmarkr\checkmarkc\checkmark$\checkmark2\checkmark$\checkmark^\checkmark\\checkmarkc\checkmarki\checkmarkr\checkmarkc\checkmark$\checkmark0\checkmark$\checkmark^\checkmark\\checkmarkc\checkmarki\checkmarkr\checkmarkc\checkmark$\checkmark$\checkmark$\checkmark^\checkmark\\checkmarkc\checkmarki\checkmarkr\checkmarkc\checkmark$\checkmark \checkmark$\checkmark^\checkmark\\checkmarkc\checkmarki\checkmarkr\checkmarkc\checkmark$\checkmarkm\checkmark$\checkmark^\checkmark\\checkmarkc\checkmarki\checkmarkr\checkmarkc\checkmark$\checkmarkm\checkmark$\checkmark^\checkmark\\checkmarkc\checkmarki\checkmarkr\checkmarkc\checkmark$\checkmark \checkmark$\checkmark^\checkmark\\checkmarkc\checkmarki\checkmarkr\checkmarkc\checkmark$\checkmarke\checkmark$\checkmark^\checkmark\\checkmarkc\checkmarki\checkmarkr\checkmarkc\checkmark$\checkmarks\checkmark$\checkmark^\checkmark\\checkmarkc\checkmarki\checkmarkr\checkmarkc\checkmark$\checkmarkt\checkmark$\checkmark^\checkmark\\checkmarkc\checkmarki\checkmarkr\checkmarkc\checkmark$\checkmark \checkmark$\checkmark^\checkmark\\checkmarkc\checkmarki\checkmarkr\checkmarkc\checkmark$\checkmarka\checkmark$\checkmark^\checkmark\\checkmarkc\checkmarki\checkmarkr\checkmarkc\checkmark$\checkmarkm\checkmark$\checkmark^\checkmark\\checkmarkc\checkmarki\checkmarkr\checkmarkc\checkmark$\checkmarkp\checkmark$\checkmark^\checkmark\\checkmarkc\checkmarki\checkmarkr\checkmarkc\checkmark$\checkmarkl\checkmark$\checkmark^\checkmark\\checkmarkc\checkmarki\checkmarkr\checkmarkc\checkmark$\checkmarke\checkmark$\checkmark^\checkmark\\checkmarkc\checkmarki\checkmarkr\checkmarkc\checkmark$\checkmarkm\checkmark$\checkmark^\checkmark\\checkmarkc\checkmarki\checkmarkr\checkmarkc\checkmark$\checkmarke\checkmark$\checkmark^\checkmark\\checkmarkc\checkmarki\checkmarkr\checkmarkc\checkmark$\checkmarkn\checkmark$\checkmark^\checkmark\\checkmarkc\checkmarki\checkmarkr\checkmarkc\checkmark$\checkmarkt\checkmark$\checkmark^\checkmark\\checkmarkc\checkmarki\checkmarkr\checkmarkc\checkmark$\checkmark \checkmark$\checkmark^\checkmark\\checkmarkc\checkmarki\checkmarkr\checkmarkc\checkmark$\checkmarks\checkmark$\checkmark^\checkmark\\checkmarkc\checkmarki\checkmarkr\checkmarkc\checkmark$\checkmarku\checkmark$\checkmark^\checkmark\\checkmarkc\checkmarki\checkmarkr\checkmarkc\checkmark$\checkmarkf\checkmark$\checkmark^\checkmark\\checkmarkc\checkmarki\checkmarkr\checkmarkc\checkmark$\checkmarkf\checkmark$\checkmark^\checkmark\\checkmarkc\checkmarki\checkmarkr\checkmarkc\checkmark$\checkmarki\checkmark$\checkmark^\checkmark\\checkmarkc\checkmarki\checkmarkr\checkmarkc\checkmark$\checkmarks\checkmark$\checkmark^\checkmark\\checkmarkc\checkmarki\checkmarkr\checkmarkc\checkmark$\checkmarka\checkmark$\checkmark^\checkmark\\checkmarkc\checkmarki\checkmarkr\checkmarkc\checkmark$\checkmarkn\checkmark$\checkmark^\checkmark\\checkmarkc\checkmarki\checkmarkr\checkmarkc\checkmark$\checkmarkt\checkmark$\checkmark^\checkmark\\checkmarkc\checkmarki\checkmarkr\checkmarkc\checkmark$\checkmark \checkmark$\checkmark^\checkmark\\checkmarkc\checkmarki\checkmarkr\checkmarkc\checkmark$\checkmarke\checkmark$\checkmark^\checkmark\\checkmarkc\checkmarki\checkmarkr\checkmarkc\checkmark$\checkmarkt\checkmark$\checkmark^\checkmark\\checkmarkc\checkmarki\checkmarkr\checkmarkc\checkmark$\checkmark \checkmark$\checkmark^\checkmark\\checkmarkc\checkmarki\checkmarkr\checkmarkc\checkmark$\checkmarkc\checkmark$\checkmark^\checkmark\\checkmarkc\checkmarki\checkmarkr\checkmarkc\checkmark$\checkmarko\checkmark$\checkmark^\checkmark\\checkmarkc\checkmarki\checkmarkr\checkmarkc\checkmark$\checkmarkn\checkmark$\checkmark^\checkmark\\checkmarkc\checkmarki\checkmarkr\checkmarkc\checkmark$\checkmarkf\checkmark$\checkmark^\checkmark\\checkmarkc\checkmarki\checkmarkr\checkmarkc\checkmark$\checkmarko\checkmark$\checkmark^\checkmark\\checkmarkc\checkmarki\checkmarkr\checkmarkc\checkmark$\checkmarkr\checkmark$\checkmark^\checkmark\\checkmarkc\checkmarki\checkmarkr\checkmarkc\checkmark$\checkmarkm\checkmark$\checkmark^\checkmark\\checkmarkc\checkmarki\checkmarkr\checkmarkc\checkmark$\checkmarke\checkmark$\checkmark^\checkmark\\checkmarkc\checkmarki\checkmarkr\checkmarkc\checkmark$\checkmark.\checkmark$\checkmark^\checkmark\\checkmarkc\checkmarki\checkmarkr\checkmarkc\checkmark$\checkmark
\checkmark$\checkmark^\checkmark\\checkmarkc\checkmarki\checkmarkr\checkmarkc\checkmark$\checkmark
\checkmark$\checkmark^\checkmark\\checkmarkc\checkmarki\checkmarkr\checkmarkc\checkmark$\checkmark%\checkmark$\checkmark^\checkmark\\checkmarkc\checkmarki\checkmarkr\checkmarkc\checkmark$\checkmark \checkmark$\checkmark^\checkmark\\checkmarkc\checkmarki\checkmarkr\checkmarkc\checkmark$\checkmark=\checkmark$\checkmark^\checkmark\\checkmarkc\checkmarki\checkmarkr\checkmarkc\checkmark$\checkmark=\checkmark$\checkmark^\checkmark\\checkmarkc\checkmarki\checkmarkr\checkmarkc\checkmark$\checkmark=\checkmark$\checkmark^\checkmark\\checkmarkc\checkmarki\checkmarkr\checkmarkc\checkmark$\checkmark=\checkmark$\checkmark^\checkmark\\checkmarkc\checkmarki\checkmarkr\checkmarkc\checkmark$\checkmark=\checkmark$\checkmark^\checkmark\\checkmarkc\checkmarki\checkmarkr\checkmarkc\checkmark$\checkmark=\checkmark$\checkmark^\checkmark\\checkmarkc\checkmarki\checkmarkr\checkmarkc\checkmark$\checkmark=\checkmark$\checkmark^\checkmark\\checkmarkc\checkmarki\checkmarkr\checkmarkc\checkmark$\checkmark=\checkmark$\checkmark^\checkmark\\checkmarkc\checkmarki\checkmarkr\checkmarkc\checkmark$\checkmark=\checkmark$\checkmark^\checkmark\\checkmarkc\checkmarki\checkmarkr\checkmarkc\checkmark$\checkmark=\checkmark$\checkmark^\checkmark\\checkmarkc\checkmarki\checkmarkr\checkmarkc\checkmark$\checkmark=\checkmark$\checkmark^\checkmark\\checkmarkc\checkmarki\checkmarkr\checkmarkc\checkmark$\checkmark=\checkmark$\checkmark^\checkmark\\checkmarkc\checkmarki\checkmarkr\checkmarkc\checkmark$\checkmark=\checkmark$\checkmark^\checkmark\\checkmarkc\checkmarki\checkmarkr\checkmarkc\checkmark$\checkmark=\checkmark$\checkmark^\checkmark\\checkmarkc\checkmarki\checkmarkr\checkmarkc\checkmark$\checkmark=\checkmark$\checkmark^\checkmark\\checkmarkc\checkmarki\checkmarkr\checkmarkc\checkmark$\checkmark=\checkmark$\checkmark^\checkmark\\checkmarkc\checkmarki\checkmarkr\checkmarkc\checkmark$\checkmark=\checkmark$\checkmark^\checkmark\\checkmarkc\checkmarki\checkmarkr\checkmarkc\checkmark$\checkmark=\checkmark$\checkmark^\checkmark\\checkmarkc\checkmarki\checkmarkr\checkmarkc\checkmark$\checkmark=\checkmark$\checkmark^\checkmark\\checkmarkc\checkmarki\checkmarkr\checkmarkc\checkmark$\checkmark=\checkmark$\checkmark^\checkmark\\checkmarkc\checkmarki\checkmarkr\checkmarkc\checkmark$\checkmark=\checkmark$\checkmark^\checkmark\\checkmarkc\checkmarki\checkmarkr\checkmarkc\checkmark$\checkmark=\checkmark$\checkmark^\checkmark\\checkmarkc\checkmarki\checkmarkr\checkmarkc\checkmark$\checkmark=\checkmark$\checkmark^\checkmark\\checkmarkc\checkmarki\checkmarkr\checkmarkc\checkmark$\checkmark=\checkmark$\checkmark^\checkmark\\checkmarkc\checkmarki\checkmarkr\checkmarkc\checkmark$\checkmark=\checkmark$\checkmark^\checkmark\\checkmarkc\checkmarki\checkmarkr\checkmarkc\checkmark$\checkmark=\checkmark$\checkmark^\checkmark\\checkmarkc\checkmarki\checkmarkr\checkmarkc\checkmark$\checkmark=\checkmark$\checkmark^\checkmark\\checkmarkc\checkmarki\checkmarkr\checkmarkc\checkmark$\checkmark=\checkmark$\checkmark^\checkmark\\checkmarkc\checkmarki\checkmarkr\checkmarkc\checkmark$\checkmark=\checkmark$\checkmark^\checkmark\\checkmarkc\checkmarki\checkmarkr\checkmarkc\checkmark$\checkmark=\checkmark$\checkmark^\checkmark\\checkmarkc\checkmarki\checkmarkr\checkmarkc\checkmark$\checkmark=\checkmark$\checkmark^\checkmark\\checkmarkc\checkmarki\checkmarkr\checkmarkc\checkmark$\checkmark=\checkmark$\checkmark^\checkmark\\checkmarkc\checkmarki\checkmarkr\checkmarkc\checkmark$\checkmark=\checkmark$\checkmark^\checkmark\\checkmarkc\checkmarki\checkmarkr\checkmarkc\checkmark$\checkmark=\checkmark$\checkmark^\checkmark\\checkmarkc\checkmarki\checkmarkr\checkmarkc\checkmark$\checkmark=\checkmark$\checkmark^\checkmark\\checkmarkc\checkmarki\checkmarkr\checkmarkc\checkmark$\checkmark=\checkmark$\checkmark^\checkmark\\checkmarkc\checkmarki\checkmarkr\checkmarkc\checkmark$\checkmark=\checkmark$\checkmark^\checkmark\\checkmarkc\checkmarki\checkmarkr\checkmarkc\checkmark$\checkmark=\checkmark$\checkmark^\checkmark\\checkmarkc\checkmarki\checkmarkr\checkmarkc\checkmark$\checkmark=\checkmark$\checkmark^\checkmark\\checkmarkc\checkmarki\checkmarkr\checkmarkc\checkmark$\checkmark=\checkmark$\checkmark^\checkmark\\checkmarkc\checkmarki\checkmarkr\checkmarkc\checkmark$\checkmark=\checkmark$\checkmark^\checkmark\\checkmarkc\checkmarki\checkmarkr\checkmarkc\checkmark$\checkmark=\checkmark$\checkmark^\checkmark\\checkmarkc\checkmarki\checkmarkr\checkmarkc\checkmark$\checkmark=\checkmark$\checkmark^\checkmark\\checkmarkc\checkmarki\checkmarkr\checkmarkc\checkmark$\checkmark=\checkmark$\checkmark^\checkmark\\checkmarkc\checkmarki\checkmarkr\checkmarkc\checkmark$\checkmark=\checkmark$\checkmark^\checkmark\\checkmarkc\checkmarki\checkmarkr\checkmarkc\checkmark$\checkmark=\checkmark$\checkmark^\checkmark\\checkmarkc\checkmarki\checkmarkr\checkmarkc\checkmark$\checkmark=\checkmark$\checkmark^\checkmark\\checkmarkc\checkmarki\checkmarkr\checkmarkc\checkmark$\checkmark=\checkmark$\checkmark^\checkmark\\checkmarkc\checkmarki\checkmarkr\checkmarkc\checkmark$\checkmark=\checkmark$\checkmark^\checkmark\\checkmarkc\checkmarki\checkmarkr\checkmarkc\checkmark$\checkmark=\checkmark$\checkmark^\checkmark\\checkmarkc\checkmarki\checkmarkr\checkmarkc\checkmark$\checkmark=\checkmark$\checkmark^\checkmark\\checkmarkc\checkmarki\checkmarkr\checkmarkc\checkmark$\checkmark=\checkmark$\checkmark^\checkmark\\checkmarkc\checkmarki\checkmarkr\checkmarkc\checkmark$\checkmark=\checkmark$\checkmark^\checkmark\\checkmarkc\checkmarki\checkmarkr\checkmarkc\checkmark$\checkmark=\checkmark$\checkmark^\checkmark\\checkmarkc\checkmarki\checkmarkr\checkmarkc\checkmark$\checkmark=\checkmark$\checkmark^\checkmark\\checkmarkc\checkmarki\checkmarkr\checkmarkc\checkmark$\checkmark=\checkmark$\checkmark^\checkmark\\checkmarkc\checkmarki\checkmarkr\checkmarkc\checkmark$\checkmark=\checkmark$\checkmark^\checkmark\\checkmarkc\checkmarki\checkmarkr\checkmarkc\checkmark$\checkmark=\checkmark$\checkmark^\checkmark\\checkmarkc\checkmarki\checkmarkr\checkmarkc\checkmark$\checkmark=\checkmark$\checkmark^\checkmark\\checkmarkc\checkmarki\checkmarkr\checkmarkc\checkmark$\checkmark=\checkmark$\checkmark^\checkmark\\checkmarkc\checkmarki\checkmarkr\checkmarkc\checkmark$\checkmark=\checkmark$\checkmark^\checkmark\\checkmarkc\checkmarki\checkmarkr\checkmarkc\checkmark$\checkmark=\checkmark$\checkmark^\checkmark\\checkmarkc\checkmarki\checkmarkr\checkmarkc\checkmark$\checkmark
\checkmark$\checkmark^\checkmark\\checkmarkc\checkmarki\checkmarkr\checkmarkc\checkmark$\checkmark\\checkmark$\checkmark^\checkmark\\checkmarkc\checkmarki\checkmarkr\checkmarkc\checkmark$\checkmarks\checkmark$\checkmark^\checkmark\\checkmarkc\checkmarki\checkmarkr\checkmarkc\checkmark$\checkmarke\checkmark$\checkmark^\checkmark\\checkmarkc\checkmarki\checkmarkr\checkmarkc\checkmark$\checkmarkc\checkmark$\checkmark^\checkmark\\checkmarkc\checkmarki\checkmarkr\checkmarkc\checkmark$\checkmarkt\checkmark$\checkmark^\checkmark\\checkmarkc\checkmarki\checkmarkr\checkmarkc\checkmark$\checkmarki\checkmark$\checkmark^\checkmark\\checkmarkc\checkmarki\checkmarkr\checkmarkc\checkmark$\checkmarko\checkmark$\checkmark^\checkmark\\checkmarkc\checkmarki\checkmarkr\checkmarkc\checkmark$\checkmarkn\checkmark$\checkmark^\checkmark\\checkmarkc\checkmarki\checkmarkr\checkmarkc\checkmark$\checkmark{\checkmark$\checkmark^\checkmark\\checkmarkc\checkmarki\checkmarkr\checkmarkc\checkmark$\checkmarkD\checkmark$\checkmark^\checkmark\\checkmarkc\checkmarki\checkmarkr\checkmarkc\checkmark$\checkmarki\checkmark$\checkmark^\checkmark\\checkmarkc\checkmarki\checkmarkr\checkmarkc\checkmark$\checkmarkm\checkmark$\checkmark^\checkmark\\checkmarkc\checkmarki\checkmarkr\checkmarkc\checkmark$\checkmarke\checkmark$\checkmark^\checkmark\\checkmarkc\checkmarki\checkmarkr\checkmarkc\checkmark$\checkmarkn\checkmark$\checkmark^\checkmark\\checkmarkc\checkmarki\checkmarkr\checkmarkc\checkmark$\checkmarks\checkmark$\checkmark^\checkmark\\checkmarkc\checkmarki\checkmarkr\checkmarkc\checkmark$\checkmarki\checkmark$\checkmark^\checkmark\\checkmarkc\checkmarki\checkmarkr\checkmarkc\checkmark$\checkmarko\checkmark$\checkmark^\checkmark\\checkmarkc\checkmarki\checkmarkr\checkmarkc\checkmark$\checkmarkn\checkmark$\checkmark^\checkmark\\checkmarkc\checkmarki\checkmarkr\checkmarkc\checkmark$\checkmarkn\checkmark$\checkmark^\checkmark\\checkmarkc\checkmarki\checkmarkr\checkmarkc\checkmark$\checkmarke\checkmark$\checkmark^\checkmark\\checkmarkc\checkmarki\checkmarkr\checkmarkc\checkmark$\checkmarkm\checkmark$\checkmark^\checkmark\\checkmarkc\checkmarki\checkmarkr\checkmarkc\checkmark$\checkmarke\checkmark$\checkmark^\checkmark\\checkmarkc\checkmarki\checkmarkr\checkmarkc\checkmark$\checkmarkn\checkmark$\checkmark^\checkmark\\checkmarkc\checkmarki\checkmarkr\checkmarkc\checkmark$\checkmarkt\checkmark$\checkmark^\checkmark\\checkmarkc\checkmarki\checkmarkr\checkmarkc\checkmark$\checkmark \checkmark$\checkmark^\checkmark\\checkmarkc\checkmarki\checkmarkr\checkmarkc\checkmark$\checkmarkd\checkmark$\checkmark^\checkmark\\checkmarkc\checkmarki\checkmarkr\checkmarkc\checkmark$\checkmarke\checkmark$\checkmark^\checkmark\\checkmarkc\checkmarki\checkmarkr\checkmarkc\checkmark$\checkmark \checkmark$\checkmark^\checkmark\\checkmarkc\checkmarki\checkmarkr\checkmarkc\checkmark$\checkmarkl\checkmark$\checkmark^\checkmark\\checkmarkc\checkmarki\checkmarkr\checkmarkc\checkmark$\checkmark'\checkmark$\checkmark^\checkmark\\checkmarkc\checkmarki\checkmarkr\checkmarkc\checkmark$\checkmarkI\checkmark$\checkmark^\checkmark\\checkmarkc\checkmarki\checkmarkr\checkmarkc\checkmark$\checkmarkn\checkmark$\checkmark^\checkmark\\checkmarkc\checkmarki\checkmarkr\checkmarkc\checkmark$\checkmarkf\checkmark$\checkmark^\checkmark\\checkmarkc\checkmarki\checkmarkr\checkmarkc\checkmark$\checkmarkr\checkmark$\checkmark^\checkmark\\checkmarkc\checkmarki\checkmarkr\checkmarkc\checkmark$\checkmarka\checkmark$\checkmark^\checkmark\\checkmarkc\checkmarki\checkmarkr\checkmarkc\checkmark$\checkmarks\checkmark$\checkmark^\checkmark\\checkmarkc\checkmarki\checkmarkr\checkmarkc\checkmark$\checkmarkt\checkmark$\checkmark^\checkmark\\checkmarkc\checkmarki\checkmarkr\checkmarkc\checkmark$\checkmarkr\checkmark$\checkmark^\checkmark\\checkmarkc\checkmarki\checkmarkr\checkmarkc\checkmark$\checkmarku\checkmark$\checkmark^\checkmark\\checkmarkc\checkmarki\checkmarkr\checkmarkc\checkmark$\checkmarkc\checkmark$\checkmark^\checkmark\\checkmarkc\checkmarki\checkmarkr\checkmarkc\checkmark$\checkmarkt\checkmark$\checkmark^\checkmark\\checkmarkc\checkmarki\checkmarkr\checkmarkc\checkmark$\checkmarku\checkmark$\checkmark^\checkmark\\checkmarkc\checkmarki\checkmarkr\checkmarkc\checkmark$\checkmarkr\checkmark$\checkmark^\checkmark\\checkmarkc\checkmarki\checkmarkr\checkmarkc\checkmark$\checkmarke\checkmark$\checkmark^\checkmark\\checkmarkc\checkmarki\checkmarkr\checkmarkc\checkmark$\checkmark \checkmark$\checkmark^\checkmark\\checkmarkc\checkmarki\checkmarkr\checkmarkc\checkmark$\checkmark(\checkmark$\checkmark^\checkmark\\checkmarkc\checkmarki\checkmarkr\checkmarkc\checkmark$\checkmarkP\checkmark$\checkmark^\checkmark\\checkmarkc\checkmarki\checkmarkr\checkmarkc\checkmark$\checkmarko\checkmark$\checkmark^\checkmark\\checkmarkc\checkmarki\checkmarkr\checkmarkc\checkmark$\checkmarkr\checkmark$\checkmark^\checkmark\\checkmarkc\checkmarki\checkmarkr\checkmarkc\checkmark$\checkmarkt\checkmark$\checkmark^\checkmark\\checkmarkc\checkmarki\checkmarkr\checkmarkc\checkmark$\checkmarki\checkmark$\checkmark^\checkmark\\checkmarkc\checkmarki\checkmarkr\checkmarkc\checkmark$\checkmarkq\checkmark$\checkmark^\checkmark\\checkmarkc\checkmarki\checkmarkr\checkmarkc\checkmark$\checkmarku\checkmark$\checkmark^\checkmark\\checkmarkc\checkmarki\checkmarkr\checkmarkc\checkmark$\checkmarke\checkmark$\checkmark^\checkmark\\checkmarkc\checkmarki\checkmarkr\checkmarkc\checkmark$\checkmark)\checkmark$\checkmark^\checkmark\\checkmarkc\checkmarki\checkmarkr\checkmarkc\checkmark$\checkmark}\checkmark$\checkmark^\checkmark\\checkmarkc\checkmarki\checkmarkr\checkmarkc\checkmark$\checkmark
\checkmark$\checkmark^\checkmark\\checkmarkc\checkmarki\checkmarkr\checkmarkc\checkmark$\checkmark
\checkmark$\checkmark^\checkmark\\checkmarkc\checkmarki\checkmarkr\checkmarkc\checkmark$\checkmark\\checkmark$\checkmark^\checkmark\\checkmarkc\checkmarki\checkmarkr\checkmarkc\checkmark$\checkmarks\checkmark$\checkmark^\checkmark\\checkmarkc\checkmarki\checkmarkr\checkmarkc\checkmark$\checkmarku\checkmark$\checkmark^\checkmark\\checkmarkc\checkmarki\checkmarkr\checkmarkc\checkmark$\checkmarkb\checkmark$\checkmark^\checkmark\\checkmarkc\checkmarki\checkmarkr\checkmarkc\checkmark$\checkmarks\checkmark$\checkmark^\checkmark\\checkmarkc\checkmarki\checkmarkr\checkmarkc\checkmark$\checkmarke\checkmark$\checkmark^\checkmark\\checkmarkc\checkmarki\checkmarkr\checkmarkc\checkmark$\checkmarkc\checkmark$\checkmark^\checkmark\\checkmarkc\checkmarki\checkmarkr\checkmarkc\checkmark$\checkmarkt\checkmark$\checkmark^\checkmark\\checkmarkc\checkmarki\checkmarkr\checkmarkc\checkmark$\checkmarki\checkmark$\checkmark^\checkmark\\checkmarkc\checkmarki\checkmarkr\checkmarkc\checkmark$\checkmarko\checkmark$\checkmark^\checkmark\\checkmarkc\checkmarki\checkmarkr\checkmarkc\checkmark$\checkmarkn\checkmark$\checkmark^\checkmark\\checkmarkc\checkmarki\checkmarkr\checkmarkc\checkmark$\checkmark{\checkmark$\checkmark^\checkmark\\checkmarkc\checkmarki\checkmarkr\checkmarkc\checkmark$\checkmarkC\checkmark$\checkmark^\checkmark\\checkmarkc\checkmarki\checkmarkr\checkmarkc\checkmark$\checkmarkh\checkmark$\checkmark^\checkmark\\checkmarkc\checkmarki\checkmarkr\checkmarkc\checkmark$\checkmarko\checkmark$\checkmark^\checkmark\\checkmarkc\checkmarki\checkmarkr\checkmarkc\checkmark$\checkmarki\checkmark$\checkmark^\checkmark\\checkmarkc\checkmarki\checkmarkr\checkmarkc\checkmark$\checkmarkx\checkmark$\checkmark^\checkmark\\checkmarkc\checkmarki\checkmarkr\checkmarkc\checkmark$\checkmark \checkmark$\checkmark^\checkmark\\checkmarkc\checkmarki\checkmarkr\checkmarkc\checkmark$\checkmarkd\checkmark$\checkmark^\checkmark\\checkmarkc\checkmarki\checkmarkr\checkmarkc\checkmark$\checkmarku\checkmark$\checkmark^\checkmark\\checkmarkc\checkmarki\checkmarkr\checkmarkc\checkmark$\checkmark \checkmark$\checkmark^\checkmark\\checkmarkc\checkmarki\checkmarkr\checkmarkc\checkmark$\checkmarkm\checkmark$\checkmark^\checkmark\\checkmarkc\checkmarki\checkmarkr\checkmarkc\checkmark$\checkmarka\checkmark$\checkmark^\checkmark\\checkmarkc\checkmarki\checkmarkr\checkmarkc\checkmark$\checkmarkt\checkmark$\checkmark^\checkmark\\checkmarkc\checkmarki\checkmarkr\checkmarkc\checkmark$\checkmarké\checkmark$\checkmark^\checkmark\\checkmarkc\checkmarki\checkmarkr\checkmarkc\checkmark$\checkmarkr\checkmark$\checkmark^\checkmark\\checkmarkc\checkmarki\checkmarkr\checkmarkc\checkmark$\checkmarki\checkmark$\checkmark^\checkmark\\checkmarkc\checkmarki\checkmarkr\checkmarkc\checkmark$\checkmarka\checkmark$\checkmark^\checkmark\\checkmarkc\checkmarki\checkmarkr\checkmarkc\checkmark$\checkmarku\checkmark$\checkmark^\checkmark\\checkmarkc\checkmarki\checkmarkr\checkmarkc\checkmark$\checkmark \checkmark$\checkmark^\checkmark\\checkmarkc\checkmarki\checkmarkr\checkmarkc\checkmark$\checkmark—\checkmark$\checkmark^\checkmark\\checkmarkc\checkmarki\checkmarkr\checkmarkc\checkmark$\checkmark \checkmark$\checkmark^\checkmark\\checkmarkc\checkmarki\checkmarkr\checkmarkc\checkmark$\checkmarkP\checkmark$\checkmark^\checkmark\\checkmarkc\checkmarki\checkmarkr\checkmarkc\checkmark$\checkmarkr\checkmark$\checkmark^\checkmark\\checkmarkc\checkmarki\checkmarkr\checkmarkc\checkmark$\checkmarko\checkmark$\checkmark^\checkmark\\checkmarkc\checkmarki\checkmarkr\checkmarkc\checkmark$\checkmarkf\checkmark$\checkmark^\checkmark\\checkmarkc\checkmarki\checkmarkr\checkmarkc\checkmark$\checkmarki\checkmark$\checkmark^\checkmark\\checkmarkc\checkmarki\checkmarkr\checkmarkc\checkmark$\checkmarkl\checkmark$\checkmark^\checkmark\\checkmarkc\checkmarki\checkmarkr\checkmarkc\checkmark$\checkmarké\checkmark$\checkmark^\checkmark\\checkmarkc\checkmarki\checkmarkr\checkmarkc\checkmark$\checkmark \checkmark$\checkmark^\checkmark\\checkmarkc\checkmarki\checkmarkr\checkmarkc\checkmark$\checkmarkA\checkmark$\checkmark^\checkmark\\checkmarkc\checkmarki\checkmarkr\checkmarkc\checkmark$\checkmarkl\checkmark$\checkmark^\checkmark\\checkmarkc\checkmarki\checkmarkr\checkmarkc\checkmark$\checkmarku\checkmark$\checkmark^\checkmark\\checkmarkc\checkmarki\checkmarkr\checkmarkc\checkmark$\checkmarkm\checkmark$\checkmark^\checkmark\\checkmarkc\checkmarki\checkmarkr\checkmarkc\checkmark$\checkmarki\checkmark$\checkmark^\checkmark\\checkmarkc\checkmarki\checkmarkr\checkmarkc\checkmark$\checkmarkn\checkmark$\checkmark^\checkmark\\checkmarkc\checkmarki\checkmarkr\checkmarkc\checkmark$\checkmarki\checkmark$\checkmark^\checkmark\\checkmarkc\checkmarki\checkmarkr\checkmarkc\checkmark$\checkmarku\checkmark$\checkmark^\checkmark\\checkmarkc\checkmarki\checkmarkr\checkmarkc\checkmark$\checkmarkm\checkmark$\checkmark^\checkmark\\checkmarkc\checkmarki\checkmarkr\checkmarkc\checkmark$\checkmark}\checkmark$\checkmark^\checkmark\\checkmarkc\checkmarki\checkmarkr\checkmarkc\checkmark$\checkmark
\checkmark$\checkmark^\checkmark\\checkmarkc\checkmarki\checkmarkr\checkmarkc\checkmark$\checkmarkL\checkmark$\checkmark^\checkmark\\checkmarkc\checkmarki\checkmarkr\checkmarkc\checkmark$\checkmark'\checkmark$\checkmark^\checkmark\\checkmarkc\checkmarki\checkmarkr\checkmarkc\checkmark$\checkmarko\checkmark$\checkmark^\checkmark\\checkmarkc\checkmarki\checkmarkr\checkmarkc\checkmark$\checkmarks\checkmark$\checkmark^\checkmark\\checkmarkc\checkmarki\checkmarkr\checkmarkc\checkmark$\checkmarks\checkmark$\checkmark^\checkmark\\checkmarkc\checkmarki\checkmarkr\checkmarkc\checkmark$\checkmarka\checkmark$\checkmark^\checkmark\\checkmarkc\checkmarki\checkmarkr\checkmarkc\checkmark$\checkmarkt\checkmark$\checkmark^\checkmark\\checkmarkc\checkmarki\checkmarkr\checkmarkc\checkmark$\checkmarku\checkmark$\checkmark^\checkmark\\checkmarkc\checkmarki\checkmarkr\checkmarkc\checkmark$\checkmarkr\checkmark$\checkmark^\checkmark\\checkmarkc\checkmarki\checkmarkr\checkmarkc\checkmark$\checkmarke\checkmark$\checkmark^\checkmark\\checkmarkc\checkmarki\checkmarkr\checkmarkc\checkmark$\checkmark \checkmark$\checkmark^\checkmark\\checkmarkc\checkmarki\checkmarkr\checkmarkc\checkmark$\checkmarke\checkmark$\checkmark^\checkmark\\checkmarkc\checkmarki\checkmarkr\checkmarkc\checkmark$\checkmarks\checkmark$\checkmark^\checkmark\\checkmarkc\checkmarki\checkmarkr\checkmarkc\checkmark$\checkmarkt\checkmark$\checkmark^\checkmark\\checkmarkc\checkmarki\checkmarkr\checkmarkc\checkmark$\checkmark \checkmark$\checkmark^\checkmark\\checkmarkc\checkmarki\checkmarkr\checkmarkc\checkmark$\checkmarkr\checkmark$\checkmark^\checkmark\\checkmarkc\checkmarki\checkmarkr\checkmarkc\checkmark$\checkmarké\checkmark$\checkmark^\checkmark\\checkmarkc\checkmarki\checkmarkr\checkmarkc\checkmark$\checkmarka\checkmark$\checkmark^\checkmark\\checkmarkc\checkmarki\checkmarkr\checkmarkc\checkmark$\checkmarkl\checkmark$\checkmark^\checkmark\\checkmarkc\checkmarki\checkmarkr\checkmarkc\checkmark$\checkmarki\checkmark$\checkmark^\checkmark\\checkmarkc\checkmarki\checkmarkr\checkmarkc\checkmark$\checkmarks\checkmark$\checkmark^\checkmark\\checkmarkc\checkmarki\checkmarkr\checkmarkc\checkmark$\checkmarké\checkmark$\checkmark^\checkmark\\checkmarkc\checkmarki\checkmarkr\checkmarkc\checkmark$\checkmarke\checkmark$\checkmark^\checkmark\\checkmarkc\checkmarki\checkmarkr\checkmarkc\checkmark$\checkmark \checkmark$\checkmark^\checkmark\\checkmarkc\checkmarki\checkmarkr\checkmarkc\checkmark$\checkmarke\checkmark$\checkmark^\checkmark\\checkmarkc\checkmarki\checkmarkr\checkmarkc\checkmark$\checkmarkn\checkmark$\checkmark^\checkmark\\checkmarkc\checkmarki\checkmarkr\checkmarkc\checkmark$\checkmark \checkmark$\checkmark^\checkmark\\checkmarkc\checkmarki\checkmarkr\checkmarkc\checkmark$\checkmarkp\checkmark$\checkmark^\checkmark\\checkmarkc\checkmarki\checkmarkr\checkmarkc\checkmark$\checkmarkr\checkmark$\checkmark^\checkmark\\checkmarkc\checkmarki\checkmarkr\checkmarkc\checkmark$\checkmarko\checkmark$\checkmark^\checkmark\\checkmarkc\checkmarki\checkmarkr\checkmarkc\checkmark$\checkmarkf\checkmark$\checkmark^\checkmark\\checkmarkc\checkmarki\checkmarkr\checkmarkc\checkmark$\checkmarki\checkmark$\checkmark^\checkmark\\checkmarkc\checkmarki\checkmarkr\checkmarkc\checkmark$\checkmarkl\checkmark$\checkmark^\checkmark\\checkmarkc\checkmarki\checkmarkr\checkmarkc\checkmark$\checkmarké\checkmark$\checkmark^\checkmark\\checkmarkc\checkmarki\checkmarkr\checkmarkc\checkmark$\checkmarks\checkmark$\checkmark^\checkmark\\checkmarkc\checkmarki\checkmarkr\checkmarkc\checkmark$\checkmark \checkmark$\checkmark^\checkmark\\checkmarkc\checkmarki\checkmarkr\checkmarkc\checkmark$\checkmarkm\checkmark$\checkmark^\checkmark\\checkmarkc\checkmarki\checkmarkr\checkmarkc\checkmark$\checkmarko\checkmark$\checkmark^\checkmark\\checkmarkc\checkmarki\checkmarkr\checkmarkc\checkmark$\checkmarkd\checkmark$\checkmark^\checkmark\\checkmarkc\checkmarki\checkmarkr\checkmarkc\checkmark$\checkmarku\checkmark$\checkmark^\checkmark\\checkmarkc\checkmarki\checkmarkr\checkmarkc\checkmark$\checkmarkl\checkmark$\checkmark^\checkmark\\checkmarkc\checkmarki\checkmarkr\checkmarkc\checkmark$\checkmarka\checkmark$\checkmark^\checkmark\\checkmarkc\checkmarki\checkmarkr\checkmarkc\checkmark$\checkmarki\checkmark$\checkmark^\checkmark\\checkmarkc\checkmarki\checkmarkr\checkmarkc\checkmark$\checkmarkr\checkmark$\checkmark^\checkmark\\checkmarkc\checkmarki\checkmarkr\checkmarkc\checkmark$\checkmarke\checkmark$\checkmark^\checkmark\\checkmarkc\checkmarki\checkmarkr\checkmarkc\checkmark$\checkmarks\checkmark$\checkmark^\checkmark\\checkmarkc\checkmarki\checkmarkr\checkmarkc\checkmark$\checkmark \checkmark$\checkmark^\checkmark\\checkmarkc\checkmarki\checkmarkr\checkmarkc\checkmark$\checkmarka\checkmark$\checkmark^\checkmark\\checkmarkc\checkmarki\checkmarkr\checkmarkc\checkmark$\checkmarkl\checkmark$\checkmark^\checkmark\\checkmarkc\checkmarki\checkmarkr\checkmarkc\checkmark$\checkmarku\checkmark$\checkmark^\checkmark\\checkmarkc\checkmarki\checkmarkr\checkmarkc\checkmark$\checkmarkm\checkmark$\checkmark^\checkmark\\checkmarkc\checkmarki\checkmarkr\checkmarkc\checkmark$\checkmarki\checkmark$\checkmark^\checkmark\\checkmarkc\checkmarki\checkmarkr\checkmarkc\checkmark$\checkmarkn\checkmark$\checkmark^\checkmark\\checkmarkc\checkmarki\checkmarkr\checkmarkc\checkmark$\checkmarki\checkmark$\checkmark^\checkmark\\checkmarkc\checkmarki\checkmarkr\checkmarkc\checkmark$\checkmarku\checkmark$\checkmark^\checkmark\\checkmarkc\checkmarki\checkmarkr\checkmarkc\checkmark$\checkmarkm\checkmark$\checkmark^\checkmark\\checkmarkc\checkmarki\checkmarkr\checkmarkc\checkmark$\checkmark \checkmark$\checkmark^\checkmark\\checkmarkc\checkmarki\checkmarkr\checkmarkc\checkmark$\checkmarkt\checkmark$\checkmark^\checkmark\\checkmarkc\checkmarki\checkmarkr\checkmarkc\checkmark$\checkmarky\checkmark$\checkmark^\checkmark\\checkmarkc\checkmarki\checkmarkr\checkmarkc\checkmark$\checkmarkp\checkmark$\checkmark^\checkmark\\checkmarkc\checkmarki\checkmarkr\checkmarkc\checkmark$\checkmarke\checkmark$\checkmark^\checkmark\\checkmarkc\checkmarki\checkmarkr\checkmarkc\checkmark$\checkmark \checkmark$\checkmark^\checkmark\\checkmarkc\checkmarki\checkmarkr\checkmarkc\checkmark$\checkmark\\checkmark$\checkmark^\checkmark\\checkmarkc\checkmarki\checkmarkr\checkmarkc\checkmark$\checkmarkt\checkmark$\checkmark^\checkmark\\checkmarkc\checkmarki\checkmarkr\checkmarkc\checkmark$\checkmarke\checkmark$\checkmark^\checkmark\\checkmarkc\checkmarki\checkmarkr\checkmarkc\checkmark$\checkmarkx\checkmark$\checkmark^\checkmark\\checkmarkc\checkmarki\checkmarkr\checkmarkc\checkmark$\checkmarkt\checkmark$\checkmark^\checkmark\\checkmarkc\checkmarki\checkmarkr\checkmarkc\checkmark$\checkmarkb\checkmark$\checkmark^\checkmark\\checkmarkc\checkmarki\checkmarkr\checkmarkc\checkmark$\checkmarkf\checkmark$\checkmark^\checkmark\\checkmarkc\checkmarki\checkmarkr\checkmarkc\checkmark$\checkmark{\checkmark$\checkmark^\checkmark\\checkmarkc\checkmarki\checkmarkr\checkmarkc\checkmark$\checkmarkV\checkmark$\checkmark^\checkmark\\checkmarkc\checkmarki\checkmarkr\checkmarkc\checkmark$\checkmark-\checkmark$\checkmark^\checkmark\\checkmarkc\checkmarki\checkmarkr\checkmarkc\checkmark$\checkmarkS\checkmark$\checkmark^\checkmark\\checkmarkc\checkmarki\checkmarkr\checkmarkc\checkmark$\checkmarkl\checkmark$\checkmark^\checkmark\\checkmarkc\checkmarki\checkmarkr\checkmarkc\checkmark$\checkmarko\checkmark$\checkmark^\checkmark\\checkmarkc\checkmarki\checkmarkr\checkmarkc\checkmark$\checkmarkt\checkmark$\checkmark^\checkmark\\checkmarkc\checkmarki\checkmarkr\checkmarkc\checkmark$\checkmark \checkmark$\checkmark^\checkmark\\checkmarkc\checkmarki\checkmarkr\checkmarkc\checkmark$\checkmark2\checkmark$\checkmark^\checkmark\\checkmarkc\checkmarki\checkmarkr\checkmarkc\checkmark$\checkmark0\checkmark$\checkmark^\checkmark\\checkmarkc\checkmarki\checkmarkr\checkmarkc\checkmark$\checkmark4\checkmark$\checkmark^\checkmark\\checkmarkc\checkmarki\checkmarkr\checkmarkc\checkmark$\checkmark0\checkmark$\checkmark^\checkmark\\checkmarkc\checkmarki\checkmarkr\checkmarkc\checkmark$\checkmark \checkmark$\checkmark^\checkmark\\checkmarkc\checkmarki\checkmarkr\checkmarkc\checkmark$\checkmarke\checkmark$\checkmark^\checkmark\\checkmarkc\checkmarki\checkmarkr\checkmarkc\checkmark$\checkmarkt\checkmark$\checkmark^\checkmark\\checkmarkc\checkmarki\checkmarkr\checkmarkc\checkmark$\checkmark \checkmark$\checkmark^\checkmark\\checkmarkc\checkmarki\checkmarkr\checkmarkc\checkmark$\checkmark2\checkmark$\checkmark^\checkmark\\checkmarkc\checkmarki\checkmarkr\checkmarkc\checkmark$\checkmark0\checkmark$\checkmark^\checkmark\\checkmarkc\checkmarki\checkmarkr\checkmarkc\checkmark$\checkmark8\checkmark$\checkmark^\checkmark\\checkmarkc\checkmarki\checkmarkr\checkmarkc\checkmark$\checkmark0\checkmark$\checkmark^\checkmark\\checkmarkc\checkmarki\checkmarkr\checkmarkc\checkmark$\checkmark}\checkmark$\checkmark^\checkmark\\checkmarkc\checkmarki\checkmarkr\checkmarkc\checkmark$\checkmark \checkmark$\checkmark^\checkmark\\checkmarkc\checkmarki\checkmarkr\checkmarkc\checkmark$\checkmark(\checkmark$\checkmark^\checkmark\\checkmarkc\checkmarki\checkmarkr\checkmarkc\checkmark$\checkmarkA\checkmark$\checkmark^\checkmark\\checkmarkc\checkmarki\checkmarkr\checkmarkc\checkmark$\checkmarkl\checkmark$\checkmark^\checkmark\\checkmarkc\checkmarki\checkmarkr\checkmarkc\checkmark$\checkmarku\checkmark$\checkmark^\checkmark\\checkmarkc\checkmarki\checkmarkr\checkmarkc\checkmark$\checkmarkm\checkmark$\checkmark^\checkmark\\checkmarkc\checkmarki\checkmarkr\checkmarkc\checkmark$\checkmarki\checkmark$\checkmark^\checkmark\\checkmarkc\checkmarki\checkmarkr\checkmarkc\checkmark$\checkmarkn\checkmark$\checkmark^\checkmark\\checkmarkc\checkmarki\checkmarkr\checkmarkc\checkmark$\checkmarki\checkmark$\checkmark^\checkmark\\checkmarkc\checkmarki\checkmarkr\checkmarkc\checkmark$\checkmarku\checkmark$\checkmark^\checkmark\\checkmarkc\checkmarki\checkmarkr\checkmarkc\checkmark$\checkmarkm\checkmark$\checkmark^\checkmark\\checkmarkc\checkmarki\checkmarkr\checkmarkc\checkmark$\checkmark \checkmark$\checkmark^\checkmark\\checkmarkc\checkmarki\checkmarkr\checkmarkc\checkmark$\checkmark6\checkmark$\checkmark^\checkmark\\checkmarkc\checkmarki\checkmarkr\checkmarkc\checkmark$\checkmark0\checkmark$\checkmark^\checkmark\\checkmarkc\checkmarki\checkmarkr\checkmarkc\checkmark$\checkmark6\checkmark$\checkmark^\checkmark\\checkmarkc\checkmarki\checkmarkr\checkmarkc\checkmark$\checkmark1\checkmark$\checkmark^\checkmark\\checkmarkc\checkmarki\checkmarkr\checkmarkc\checkmark$\checkmark-\checkmark$\checkmark^\checkmark\\checkmarkc\checkmarki\checkmarkr\checkmarkc\checkmark$\checkmarkT\checkmark$\checkmark^\checkmark\\checkmarkc\checkmarki\checkmarkr\checkmarkc\checkmark$\checkmark6\checkmark$\checkmark^\checkmark\\checkmarkc\checkmarki\checkmarkr\checkmarkc\checkmark$\checkmark)\checkmark$\checkmark^\checkmark\\checkmarkc\checkmarki\checkmarkr\checkmarkc\checkmark$\checkmark \checkmark$\checkmark^\checkmark\\checkmarkc\checkmarki\checkmarkr\checkmarkc\checkmark$\checkmarkd\checkmark$\checkmark^\checkmark\\checkmarkc\checkmarki\checkmarkr\checkmarkc\checkmark$\checkmarko\checkmark$\checkmark^\checkmark\\checkmarkc\checkmarki\checkmarkr\checkmarkc\checkmark$\checkmarkn\checkmark$\checkmark^\checkmark\\checkmarkc\checkmarki\checkmarkr\checkmarkc\checkmark$\checkmarkt\checkmark$\checkmark^\checkmark\\checkmarkc\checkmarki\checkmarkr\checkmarkc\checkmark$\checkmark \checkmark$\checkmark^\checkmark\\checkmarkc\checkmarki\checkmarkr\checkmarkc\checkmark$\checkmarkl\checkmark$\checkmark^\checkmark\\checkmarkc\checkmarki\checkmarkr\checkmarkc\checkmark$\checkmarke\checkmark$\checkmark^\checkmark\\checkmarkc\checkmarki\checkmarkr\checkmarkc\checkmark$\checkmarks\checkmark$\checkmark^\checkmark\\checkmarkc\checkmarki\checkmarkr\checkmarkc\checkmark$\checkmark \checkmark$\checkmark^\checkmark\\checkmarkc\checkmarki\checkmarkr\checkmarkc\checkmark$\checkmarkc\checkmark$\checkmark^\checkmark\\checkmarkc\checkmarki\checkmarkr\checkmarkc\checkmark$\checkmarka\checkmark$\checkmark^\checkmark\\checkmarkc\checkmarki\checkmarkr\checkmarkc\checkmark$\checkmarkr\checkmark$\checkmark^\checkmark\\checkmarkc\checkmarki\checkmarkr\checkmarkc\checkmark$\checkmarka\checkmark$\checkmark^\checkmark\\checkmarkc\checkmarki\checkmarkr\checkmarkc\checkmark$\checkmarkc\checkmark$\checkmark^\checkmark\\checkmarkc\checkmarki\checkmarkr\checkmarkc\checkmark$\checkmarkt\checkmark$\checkmark^\checkmark\\checkmarkc\checkmarki\checkmarkr\checkmarkc\checkmark$\checkmarké\checkmark$\checkmark^\checkmark\\checkmarkc\checkmarki\checkmarkr\checkmarkc\checkmark$\checkmarkr\checkmark$\checkmark^\checkmark\\checkmarkc\checkmarki\checkmarkr\checkmarkc\checkmark$\checkmarki\checkmark$\checkmark^\checkmark\\checkmarkc\checkmarki\checkmarkr\checkmarkc\checkmark$\checkmarks\checkmark$\checkmark^\checkmark\\checkmarkc\checkmarki\checkmarkr\checkmarkc\checkmark$\checkmarkt\checkmark$\checkmark^\checkmark\\checkmarkc\checkmarki\checkmarkr\checkmarkc\checkmark$\checkmarki\checkmark$\checkmark^\checkmark\\checkmarkc\checkmarki\checkmarkr\checkmarkc\checkmark$\checkmarkq\checkmark$\checkmark^\checkmark\\checkmarkc\checkmarki\checkmarkr\checkmarkc\checkmark$\checkmarku\checkmark$\checkmark^\checkmark\\checkmarkc\checkmarki\checkmarkr\checkmarkc\checkmark$\checkmarke\checkmark$\checkmark^\checkmark\\checkmarkc\checkmarki\checkmarkr\checkmarkc\checkmark$\checkmarks\checkmark$\checkmark^\checkmark\\checkmarkc\checkmarki\checkmarkr\checkmarkc\checkmark$\checkmark \checkmark$\checkmark^\checkmark\\checkmarkc\checkmarki\checkmarkr\checkmarkc\checkmark$\checkmarks\checkmark$\checkmark^\checkmark\\checkmarkc\checkmarki\checkmarkr\checkmarkc\checkmark$\checkmarko\checkmark$\checkmark^\checkmark\\checkmarkc\checkmarki\checkmarkr\checkmarkc\checkmark$\checkmarkn\checkmark$\checkmark^\checkmark\\checkmarkc\checkmarki\checkmarkr\checkmarkc\checkmark$\checkmarkt\checkmark$\checkmark^\checkmark\\checkmarkc\checkmarki\checkmarkr\checkmarkc\checkmark$\checkmark \checkmark$\checkmark^\checkmark\\checkmarkc\checkmarki\checkmarkr\checkmarkc\checkmark$\checkmark:\checkmark$\checkmark^\checkmark\\checkmarkc\checkmarki\checkmarkr\checkmarkc\checkmark$\checkmark
\checkmark$\checkmark^\checkmark\\checkmarkc\checkmarki\checkmarkr\checkmarkc\checkmark$\checkmark\\checkmark$\checkmark^\checkmark\\checkmarkc\checkmarki\checkmarkr\checkmarkc\checkmark$\checkmarkb\checkmark$\checkmark^\checkmark\\checkmarkc\checkmarki\checkmarkr\checkmarkc\checkmark$\checkmarke\checkmark$\checkmark^\checkmark\\checkmarkc\checkmarki\checkmarkr\checkmarkc\checkmark$\checkmarkg\checkmark$\checkmark^\checkmark\\checkmarkc\checkmarki\checkmarkr\checkmarkc\checkmark$\checkmarki\checkmark$\checkmark^\checkmark\\checkmarkc\checkmarki\checkmarkr\checkmarkc\checkmark$\checkmarkn\checkmark$\checkmark^\checkmark\\checkmarkc\checkmarki\checkmarkr\checkmarkc\checkmark$\checkmark{\checkmark$\checkmark^\checkmark\\checkmarkc\checkmarki\checkmarkr\checkmarkc\checkmark$\checkmarki\checkmark$\checkmark^\checkmark\\checkmarkc\checkmarki\checkmarkr\checkmarkc\checkmark$\checkmarkt\checkmark$\checkmark^\checkmark\\checkmarkc\checkmarki\checkmarkr\checkmarkc\checkmark$\checkmarke\checkmark$\checkmark^\checkmark\\checkmarkc\checkmarki\checkmarkr\checkmarkc\checkmark$\checkmarkm\checkmark$\checkmark^\checkmark\\checkmarkc\checkmarki\checkmarkr\checkmarkc\checkmark$\checkmarki\checkmark$\checkmark^\checkmark\\checkmarkc\checkmarki\checkmarkr\checkmarkc\checkmark$\checkmarkz\checkmark$\checkmark^\checkmark\\checkmarkc\checkmarki\checkmarkr\checkmarkc\checkmark$\checkmarke\checkmark$\checkmark^\checkmark\\checkmarkc\checkmarki\checkmarkr\checkmarkc\checkmark$\checkmark}\checkmark$\checkmark^\checkmark\\checkmarkc\checkmarki\checkmarkr\checkmarkc\checkmark$\checkmark
\checkmark$\checkmark^\checkmark\\checkmarkc\checkmarki\checkmarkr\checkmarkc\checkmark$\checkmark \checkmark$\checkmark^\checkmark\\checkmarkc\checkmarki\checkmarkr\checkmarkc\checkmark$\checkmark \checkmark$\checkmark^\checkmark\\checkmarkc\checkmarki\checkmarkr\checkmarkc\checkmark$\checkmark \checkmark$\checkmark^\checkmark\\checkmarkc\checkmarki\checkmarkr\checkmarkc\checkmark$\checkmark \checkmark$\checkmark^\checkmark\\checkmarkc\checkmarki\checkmarkr\checkmarkc\checkmark$\checkmark\\checkmark$\checkmark^\checkmark\\checkmarkc\checkmarki\checkmarkr\checkmarkc\checkmark$\checkmarki\checkmark$\checkmark^\checkmark\\checkmarkc\checkmarki\checkmarkr\checkmarkc\checkmark$\checkmarkt\checkmark$\checkmark^\checkmark\\checkmarkc\checkmarki\checkmarkr\checkmarkc\checkmark$\checkmarke\checkmark$\checkmark^\checkmark\\checkmarkc\checkmarki\checkmarkr\checkmarkc\checkmark$\checkmarkm\checkmark$\checkmark^\checkmark\\checkmarkc\checkmarki\checkmarkr\checkmarkc\checkmark$\checkmark \checkmark$\checkmark^\checkmark\\checkmarkc\checkmarki\checkmarkr\checkmarkc\checkmark$\checkmark$\checkmark$\checkmark^\checkmark\\checkmarkc\checkmarki\checkmarkr\checkmarkc\checkmark$\checkmarkR\checkmark$\checkmark^\checkmark\\checkmarkc\checkmarki\checkmarkr\checkmarkc\checkmark$\checkmark_\checkmark$\checkmark^\checkmark\\checkmarkc\checkmarki\checkmarkr\checkmarkc\checkmark$\checkmarke\checkmark$\checkmark^\checkmark\\checkmarkc\checkmarki\checkmarkr\checkmarkc\checkmark$\checkmark \checkmark$\checkmark^\checkmark\\checkmarkc\checkmarki\checkmarkr\checkmarkc\checkmark$\checkmark=\checkmark$\checkmark^\checkmark\\checkmarkc\checkmarki\checkmarkr\checkmarkc\checkmark$\checkmark \checkmark$\checkmark^\checkmark\\checkmarkc\checkmarki\checkmarkr\checkmarkc\checkmark$\checkmark2\checkmark$\checkmark^\checkmark\\checkmarkc\checkmarki\checkmarkr\checkmarkc\checkmark$\checkmark7\checkmark$\checkmark^\checkmark\\checkmarkc\checkmarki\checkmarkr\checkmarkc\checkmark$\checkmark6\checkmark$\checkmark^\checkmark\\checkmarkc\checkmarki\checkmarkr\checkmarkc\checkmark$\checkmark$\checkmark$\checkmark^\checkmark\\checkmarkc\checkmarki\checkmarkr\checkmarkc\checkmark$\checkmark \checkmark$\checkmark^\checkmark\\checkmarkc\checkmarki\checkmarkr\checkmarkc\checkmark$\checkmarkM\checkmark$\checkmark^\checkmark\\checkmarkc\checkmarki\checkmarkr\checkmarkc\checkmark$\checkmarkP\checkmark$\checkmark^\checkmark\\checkmarkc\checkmarki\checkmarkr\checkmarkc\checkmark$\checkmarka\checkmark$\checkmark^\checkmark\\checkmarkc\checkmarki\checkmarkr\checkmarkc\checkmark$\checkmark,\checkmark$\checkmark^\checkmark\\checkmarkc\checkmarki\checkmarkr\checkmarkc\checkmark$\checkmark \checkmark$\checkmark^\checkmark\\checkmarkc\checkmarki\checkmarkr\checkmarkc\checkmark$\checkmark$\checkmark$\checkmark^\checkmark\\checkmarkc\checkmarki\checkmarkr\checkmarkc\checkmark$\checkmarkE\checkmark$\checkmark^\checkmark\\checkmarkc\checkmarki\checkmarkr\checkmarkc\checkmark$\checkmark \checkmark$\checkmark^\checkmark\\checkmarkc\checkmarki\checkmarkr\checkmarkc\checkmark$\checkmark=\checkmark$\checkmark^\checkmark\\checkmarkc\checkmarki\checkmarkr\checkmarkc\checkmark$\checkmark \checkmark$\checkmark^\checkmark\\checkmarkc\checkmarki\checkmarkr\checkmarkc\checkmark$\checkmark6\checkmark$\checkmark^\checkmark\\checkmarkc\checkmarki\checkmarkr\checkmarkc\checkmark$\checkmark8\checkmark$\checkmark^\checkmark\\checkmarkc\checkmarki\checkmarkr\checkmarkc\checkmark$\checkmark\\checkmark$\checkmark^\checkmark\\checkmarkc\checkmarki\checkmarkr\checkmarkc\checkmark$\checkmark,\checkmark$\checkmark^\checkmark\\checkmarkc\checkmarki\checkmarkr\checkmarkc\checkmark$\checkmark9\checkmark$\checkmark^\checkmark\\checkmarkc\checkmarki\checkmarkr\checkmarkc\checkmark$\checkmark0\checkmark$\checkmark^\checkmark\\checkmarkc\checkmarki\checkmarkr\checkmarkc\checkmark$\checkmark0\checkmark$\checkmark^\checkmark\\checkmarkc\checkmarki\checkmarkr\checkmarkc\checkmark$\checkmark$\checkmark$\checkmark^\checkmark\\checkmarkc\checkmarki\checkmarkr\checkmarkc\checkmark$\checkmark \checkmark$\checkmark^\checkmark\\checkmarkc\checkmarki\checkmarkr\checkmarkc\checkmark$\checkmarkM\checkmark$\checkmark^\checkmark\\checkmarkc\checkmarki\checkmarkr\checkmarkc\checkmark$\checkmarkP\checkmark$\checkmark^\checkmark\\checkmarkc\checkmarki\checkmarkr\checkmarkc\checkmark$\checkmarka\checkmark$\checkmark^\checkmark\\checkmarkc\checkmarki\checkmarkr\checkmarkc\checkmark$\checkmark,\checkmark$\checkmark^\checkmark\\checkmarkc\checkmarki\checkmarkr\checkmarkc\checkmark$\checkmark \checkmark$\checkmark^\checkmark\\checkmarkc\checkmarki\checkmarkr\checkmarkc\checkmark$\checkmark$\checkmark$\checkmark^\checkmark\\checkmarkc\checkmarki\checkmarkr\checkmarkc\checkmark$\checkmark\\checkmark$\checkmark^\checkmark\\checkmarkc\checkmarki\checkmarkr\checkmarkc\checkmark$\checkmarkr\checkmark$\checkmark^\checkmark\\checkmarkc\checkmarki\checkmarkr\checkmarkc\checkmark$\checkmarkh\checkmark$\checkmark^\checkmark\\checkmarkc\checkmarki\checkmarkr\checkmarkc\checkmark$\checkmarko\checkmark$\checkmark^\checkmark\\checkmarkc\checkmarki\checkmarkr\checkmarkc\checkmark$\checkmark \checkmark$\checkmark^\checkmark\\checkmarkc\checkmarki\checkmarkr\checkmarkc\checkmark$\checkmark=\checkmark$\checkmark^\checkmark\\checkmarkc\checkmarki\checkmarkr\checkmarkc\checkmark$\checkmark \checkmark$\checkmark^\checkmark\\checkmarkc\checkmarki\checkmarkr\checkmarkc\checkmark$\checkmark2\checkmark$\checkmark^\checkmark\\checkmarkc\checkmarki\checkmarkr\checkmarkc\checkmark$\checkmark\\checkmark$\checkmark^\checkmark\\checkmarkc\checkmarki\checkmarkr\checkmarkc\checkmark$\checkmark,\checkmark$\checkmark^\checkmark\\checkmarkc\checkmarki\checkmarkr\checkmarkc\checkmark$\checkmark7\checkmark$\checkmark^\checkmark\\checkmarkc\checkmarki\checkmarkr\checkmarkc\checkmark$\checkmark0\checkmark$\checkmark^\checkmark\\checkmarkc\checkmarki\checkmarkr\checkmarkc\checkmark$\checkmark0\checkmark$\checkmark^\checkmark\\checkmarkc\checkmarki\checkmarkr\checkmarkc\checkmark$\checkmark$\checkmark$\checkmark^\checkmark\\checkmarkc\checkmarki\checkmarkr\checkmarkc\checkmark$\checkmark \checkmark$\checkmark^\checkmark\\checkmarkc\checkmarki\checkmarkr\checkmarkc\checkmark$\checkmarkk\checkmark$\checkmark^\checkmark\\checkmarkc\checkmarki\checkmarkr\checkmarkc\checkmark$\checkmarkg\checkmark$\checkmark^\checkmark\\checkmarkc\checkmarki\checkmarkr\checkmarkc\checkmark$\checkmark/\checkmark$\checkmark^\checkmark\\checkmarkc\checkmarki\checkmarkr\checkmarkc\checkmark$\checkmarkm\checkmark$\checkmark^\checkmark\\checkmarkc\checkmarki\checkmarkr\checkmarkc\checkmark$\checkmark³\checkmark$\checkmark^\checkmark\\checkmarkc\checkmarki\checkmarkr\checkmarkc\checkmark$\checkmark
\checkmark$\checkmark^\checkmark\\checkmarkc\checkmarki\checkmarkr\checkmarkc\checkmark$\checkmark \checkmark$\checkmark^\checkmark\\checkmarkc\checkmarki\checkmarkr\checkmarkc\checkmark$\checkmark \checkmark$\checkmark^\checkmark\\checkmarkc\checkmarki\checkmarkr\checkmarkc\checkmark$\checkmark \checkmark$\checkmark^\checkmark\\checkmarkc\checkmarki\checkmarkr\checkmarkc\checkmark$\checkmark \checkmark$\checkmark^\checkmark\\checkmarkc\checkmarki\checkmarkr\checkmarkc\checkmark$\checkmark\\checkmark$\checkmark^\checkmark\\checkmarkc\checkmarki\checkmarkr\checkmarkc\checkmark$\checkmarki\checkmark$\checkmark^\checkmark\\checkmarkc\checkmarki\checkmarkr\checkmarkc\checkmark$\checkmarkt\checkmark$\checkmark^\checkmark\\checkmarkc\checkmarki\checkmarkr\checkmarkc\checkmark$\checkmarke\checkmark$\checkmark^\checkmark\\checkmarkc\checkmarki\checkmarkr\checkmarkc\checkmark$\checkmarkm\checkmark$\checkmark^\checkmark\\checkmarkc\checkmarki\checkmarkr\checkmarkc\checkmark$\checkmark \checkmark$\checkmark^\checkmark\\checkmarkc\checkmarki\checkmarkr\checkmarkc\checkmark$\checkmarkM\checkmark$\checkmark^\checkmark\\checkmarkc\checkmarki\checkmarkr\checkmarkc\checkmark$\checkmarko\checkmark$\checkmark^\checkmark\\checkmarkc\checkmarki\checkmarkr\checkmarkc\checkmark$\checkmarkm\checkmark$\checkmark^\checkmark\\checkmarkc\checkmarki\checkmarkr\checkmarkc\checkmark$\checkmarke\checkmark$\checkmark^\checkmark\\checkmarkc\checkmarki\checkmarkr\checkmarkc\checkmark$\checkmarkn\checkmark$\checkmark^\checkmark\\checkmarkc\checkmarki\checkmarkr\checkmarkc\checkmark$\checkmarkt\checkmark$\checkmark^\checkmark\\checkmarkc\checkmarki\checkmarkr\checkmarkc\checkmark$\checkmark \checkmark$\checkmark^\checkmark\\checkmarkc\checkmarki\checkmarkr\checkmarkc\checkmark$\checkmarkq\checkmark$\checkmark^\checkmark\\checkmarkc\checkmarki\checkmarkr\checkmarkc\checkmark$\checkmarku\checkmark$\checkmark^\checkmark\\checkmarkc\checkmarki\checkmarkr\checkmarkc\checkmark$\checkmarka\checkmark$\checkmark^\checkmark\\checkmarkc\checkmarki\checkmarkr\checkmarkc\checkmark$\checkmarkd\checkmark$\checkmark^\checkmark\\checkmarkc\checkmarki\checkmarkr\checkmarkc\checkmark$\checkmarkr\checkmark$\checkmark^\checkmark\\checkmarkc\checkmarki\checkmarkr\checkmarkc\checkmark$\checkmarka\checkmark$\checkmark^\checkmark\\checkmarkc\checkmarki\checkmarkr\checkmarkc\checkmark$\checkmarkt\checkmark$\checkmark^\checkmark\\checkmarkc\checkmarki\checkmarkr\checkmarkc\checkmark$\checkmarki\checkmark$\checkmark^\checkmark\\checkmarkc\checkmarki\checkmarkr\checkmarkc\checkmark$\checkmarkq\checkmark$\checkmark^\checkmark\\checkmarkc\checkmarki\checkmarkr\checkmarkc\checkmark$\checkmarku\checkmark$\checkmark^\checkmark\\checkmarkc\checkmarki\checkmarkr\checkmarkc\checkmark$\checkmarke\checkmark$\checkmark^\checkmark\\checkmarkc\checkmarki\checkmarkr\checkmarkc\checkmark$\checkmark \checkmark$\checkmark^\checkmark\\checkmarkc\checkmarki\checkmarkr\checkmarkc\checkmark$\checkmarkd\checkmark$\checkmark^\checkmark\\checkmarkc\checkmarki\checkmarkr\checkmarkc\checkmark$\checkmarku\checkmark$\checkmark^\checkmark\\checkmarkc\checkmarki\checkmarkr\checkmarkc\checkmark$\checkmark \checkmark$\checkmark^\checkmark\\checkmarkc\checkmarki\checkmarkr\checkmarkc\checkmark$\checkmarkp\checkmark$\checkmark^\checkmark\\checkmarkc\checkmarki\checkmarkr\checkmarkc\checkmark$\checkmarkr\checkmark$\checkmark^\checkmark\\checkmarkc\checkmarki\checkmarkr\checkmarkc\checkmark$\checkmarko\checkmark$\checkmark^\checkmark\\checkmarkc\checkmarki\checkmarkr\checkmarkc\checkmark$\checkmarkf\checkmark$\checkmark^\checkmark\\checkmarkc\checkmarki\checkmarkr\checkmarkc\checkmark$\checkmarki\checkmark$\checkmark^\checkmark\\checkmarkc\checkmarki\checkmarkr\checkmarkc\checkmark$\checkmarkl\checkmark$\checkmark^\checkmark\\checkmarkc\checkmarki\checkmarkr\checkmarkc\checkmark$\checkmarké\checkmark$\checkmark^\checkmark\\checkmarkc\checkmarki\checkmarkr\checkmarkc\checkmark$\checkmark \checkmark$\checkmark^\checkmark\\checkmarkc\checkmarki\checkmarkr\checkmarkc\checkmark$\checkmark2\checkmark$\checkmark^\checkmark\\checkmarkc\checkmarki\checkmarkr\checkmarkc\checkmark$\checkmark0\checkmark$\checkmark^\checkmark\\checkmarkc\checkmarki\checkmarkr\checkmarkc\checkmark$\checkmark8\checkmark$\checkmark^\checkmark\\checkmarkc\checkmarki\checkmarkr\checkmarkc\checkmark$\checkmark0\checkmark$\checkmark^\checkmark\\checkmarkc\checkmarki\checkmarkr\checkmarkc\checkmark$\checkmark \checkmark$\checkmark^\checkmark\\checkmarkc\checkmarki\checkmarkr\checkmarkc\checkmark$\checkmark:\checkmark$\checkmark^\checkmark\\checkmarkc\checkmarki\checkmarkr\checkmarkc\checkmark$\checkmark \checkmark$\checkmark^\checkmark\\checkmarkc\checkmarki\checkmarkr\checkmarkc\checkmark$\checkmark$\checkmark$\checkmark^\checkmark\\checkmarkc\checkmarki\checkmarkr\checkmarkc\checkmark$\checkmarkI\checkmark$\checkmark^\checkmark\\checkmarkc\checkmarki\checkmarkr\checkmarkc\checkmark$\checkmark_\checkmark$\checkmark^\checkmark\\checkmarkc\checkmarki\checkmarkr\checkmarkc\checkmark$\checkmark{\checkmark$\checkmark^\checkmark\\checkmarkc\checkmarki\checkmarkr\checkmarkc\checkmark$\checkmarky\checkmark$\checkmark^\checkmark\\checkmarkc\checkmarki\checkmarkr\checkmarkc\checkmark$\checkmarky\checkmark$\checkmark^\checkmark\\checkmarkc\checkmarki\checkmarkr\checkmarkc\checkmark$\checkmark\\checkmark$\checkmark^\checkmark\\checkmarkc\checkmarki\checkmarkr\checkmarkc\checkmark$\checkmark_\checkmark$\checkmark^\checkmark\\checkmarkc\checkmarki\checkmarkr\checkmarkc\checkmark$\checkmark2\checkmark$\checkmark^\checkmark\\checkmarkc\checkmarki\checkmarkr\checkmarkc\checkmark$\checkmark0\checkmark$\checkmark^\checkmark\\checkmarkc\checkmarki\checkmarkr\checkmarkc\checkmark$\checkmark8\checkmark$\checkmark^\checkmark\\checkmarkc\checkmarki\checkmarkr\checkmarkc\checkmark$\checkmark0\checkmark$\checkmark^\checkmark\\checkmarkc\checkmarki\checkmarkr\checkmarkc\checkmark$\checkmark}\checkmark$\checkmark^\checkmark\\checkmarkc\checkmarki\checkmarkr\checkmarkc\checkmark$\checkmark \checkmark$\checkmark^\checkmark\\checkmarkc\checkmarki\checkmarkr\checkmarkc\checkmark$\checkmark\\checkmark$\checkmark^\checkmark\\checkmarkc\checkmarki\checkmarkr\checkmarkc\checkmark$\checkmarka\checkmark$\checkmark^\checkmark\\checkmarkc\checkmarki\checkmarkr\checkmarkc\checkmark$\checkmarkp\checkmark$\checkmark^\checkmark\\checkmarkc\checkmarki\checkmarkr\checkmarkc\checkmark$\checkmarkp\checkmark$\checkmark^\checkmark\\checkmarkc\checkmarki\checkmarkr\checkmarkc\checkmark$\checkmarkr\checkmark$\checkmark^\checkmark\\checkmarkc\checkmarki\checkmarkr\checkmarkc\checkmark$\checkmarko\checkmark$\checkmark^\checkmark\\checkmarkc\checkmarki\checkmarkr\checkmarkc\checkmark$\checkmarkx\checkmark$\checkmark^\checkmark\\checkmarkc\checkmarki\checkmarkr\checkmarkc\checkmark$\checkmark \checkmark$\checkmark^\checkmark\\checkmarkc\checkmarki\checkmarkr\checkmarkc\checkmark$\checkmark6\checkmark$\checkmark^\checkmark\\checkmarkc\checkmarki\checkmarkr\checkmarkc\checkmark$\checkmark5\checkmark$\checkmark^\checkmark\\checkmarkc\checkmarki\checkmarkr\checkmarkc\checkmark$\checkmark\\checkmark$\checkmark^\checkmark\\checkmarkc\checkmarki\checkmarkr\checkmarkc\checkmark$\checkmark,\checkmark$\checkmark^\checkmark\\checkmarkc\checkmarki\checkmarkr\checkmarkc\checkmark$\checkmark0\checkmark$\checkmark^\checkmark\\checkmarkc\checkmarki\checkmarkr\checkmarkc\checkmark$\checkmark0\checkmark$\checkmark^\checkmark\\checkmarkc\checkmarki\checkmarkr\checkmarkc\checkmark$\checkmark0\checkmark$\checkmark^\checkmark\\checkmarkc\checkmarki\checkmarkr\checkmarkc\checkmark$\checkmark$\checkmark$\checkmark^\checkmark\\checkmarkc\checkmarki\checkmarkr\checkmarkc\checkmark$\checkmark \checkmark$\checkmark^\checkmark\\checkmarkc\checkmarki\checkmarkr\checkmarkc\checkmark$\checkmarkm\checkmark$\checkmark^\checkmark\\checkmarkc\checkmarki\checkmarkr\checkmarkc\checkmark$\checkmarkm\checkmark$\checkmark^\checkmark\\checkmarkc\checkmarki\checkmarkr\checkmarkc\checkmark$\checkmark⁴\checkmark$\checkmark^\checkmark\\checkmarkc\checkmarki\checkmarkr\checkmarkc\checkmark$\checkmark \checkmark$\checkmark^\checkmark\\checkmarkc\checkmarki\checkmarkr\checkmarkc\checkmark$\checkmark(\checkmark$\checkmark^\checkmark\\checkmarkc\checkmarki\checkmarkr\checkmarkc\checkmark$\checkmarkd\checkmark$\checkmark^\checkmark\\checkmarkc\checkmarki\checkmarkr\checkmarkc\checkmark$\checkmarko\checkmark$\checkmark^\checkmark\\checkmarkc\checkmarki\checkmarkr\checkmarkc\checkmark$\checkmarkn\checkmark$\checkmark^\checkmark\\checkmarkc\checkmarki\checkmarkr\checkmarkc\checkmark$\checkmarkn\checkmark$\checkmark^\checkmark\\checkmarkc\checkmarki\checkmarkr\checkmarkc\checkmark$\checkmarké\checkmark$\checkmark^\checkmark\\checkmarkc\checkmarki\checkmarkr\checkmarkc\checkmark$\checkmarke\checkmark$\checkmark^\checkmark\\checkmarkc\checkmarki\checkmarkr\checkmarkc\checkmark$\checkmarks\checkmark$\checkmark^\checkmark\\checkmarkc\checkmarki\checkmarkr\checkmarkc\checkmark$\checkmark \checkmark$\checkmark^\checkmark\\checkmarkc\checkmarki\checkmarkr\checkmarkc\checkmark$\checkmarkc\checkmark$\checkmark^\checkmark\\checkmarkc\checkmarki\checkmarkr\checkmarkc\checkmark$\checkmarko\checkmark$\checkmark^\checkmark\\checkmarkc\checkmarki\checkmarkr\checkmarkc\checkmark$\checkmarkn\checkmark$\checkmark^\checkmark\\checkmarkc\checkmarki\checkmarkr\checkmarkc\checkmark$\checkmarks\checkmark$\checkmark^\checkmark\\checkmarkc\checkmarki\checkmarkr\checkmarkc\checkmark$\checkmarkt\checkmark$\checkmark^\checkmark\\checkmarkc\checkmarki\checkmarkr\checkmarkc\checkmark$\checkmarkr\checkmark$\checkmark^\checkmark\\checkmarkc\checkmarki\checkmarkr\checkmarkc\checkmark$\checkmarku\checkmark$\checkmark^\checkmark\\checkmarkc\checkmarki\checkmarkr\checkmarkc\checkmark$\checkmarkc\checkmark$\checkmark^\checkmark\\checkmarkc\checkmarki\checkmarkr\checkmarkc\checkmark$\checkmarkt\checkmark$\checkmark^\checkmark\\checkmarkc\checkmarki\checkmarkr\checkmarkc\checkmark$\checkmarke\checkmark$\checkmark^\checkmark\\checkmarkc\checkmarki\checkmarkr\checkmarkc\checkmark$\checkmarku\checkmark$\checkmark^\checkmark\\checkmarkc\checkmarki\checkmarkr\checkmarkc\checkmark$\checkmarkr\checkmark$\checkmark^\checkmark\\checkmarkc\checkmarki\checkmarkr\checkmarkc\checkmark$\checkmark)\checkmark$\checkmark^\checkmark\\checkmarkc\checkmarki\checkmarkr\checkmarkc\checkmark$\checkmark
\checkmark$\checkmark^\checkmark\\checkmarkc\checkmarki\checkmarkr\checkmarkc\checkmark$\checkmark\\checkmark$\checkmark^\checkmark\\checkmarkc\checkmarki\checkmarkr\checkmarkc\checkmark$\checkmarke\checkmark$\checkmark^\checkmark\\checkmarkc\checkmarki\checkmarkr\checkmarkc\checkmark$\checkmarkn\checkmark$\checkmark^\checkmark\\checkmarkc\checkmarki\checkmarkr\checkmarkc\checkmark$\checkmarkd\checkmark$\checkmark^\checkmark\\checkmarkc\checkmarki\checkmarkr\checkmarkc\checkmark$\checkmark{\checkmark$\checkmark^\checkmark\\checkmarkc\checkmarki\checkmarkr\checkmarkc\checkmark$\checkmarki\checkmark$\checkmark^\checkmark\\checkmarkc\checkmarki\checkmarkr\checkmarkc\checkmark$\checkmarkt\checkmark$\checkmark^\checkmark\\checkmarkc\checkmarki\checkmarkr\checkmarkc\checkmark$\checkmarke\checkmark$\checkmark^\checkmark\\checkmarkc\checkmarki\checkmarkr\checkmarkc\checkmark$\checkmarkm\checkmark$\checkmark^\checkmark\\checkmarkc\checkmarki\checkmarkr\checkmarkc\checkmark$\checkmarki\checkmark$\checkmark^\checkmark\\checkmarkc\checkmarki\checkmarkr\checkmarkc\checkmark$\checkmarkz\checkmark$\checkmark^\checkmark\\checkmarkc\checkmarki\checkmarkr\checkmarkc\checkmark$\checkmarke\checkmark$\checkmark^\checkmark\\checkmarkc\checkmarki\checkmarkr\checkmarkc\checkmark$\checkmark}\checkmark$\checkmark^\checkmark\\checkmarkc\checkmarki\checkmarkr\checkmarkc\checkmark$\checkmark
\checkmark$\checkmark^\checkmark\\checkmarkc\checkmarki\checkmarkr\checkmarkc\checkmark$\checkmark
\checkmark$\checkmark^\checkmark\\checkmarkc\checkmarki\checkmarkr\checkmarkc\checkmark$\checkmark\\checkmark$\checkmark^\checkmark\\checkmarkc\checkmarki\checkmarkr\checkmarkc\checkmark$\checkmarks\checkmark$\checkmark^\checkmark\\checkmarkc\checkmarki\checkmarkr\checkmarkc\checkmark$\checkmarku\checkmark$\checkmark^\checkmark\\checkmarkc\checkmarki\checkmarkr\checkmarkc\checkmark$\checkmarkb\checkmark$\checkmark^\checkmark\\checkmarkc\checkmarki\checkmarkr\checkmarkc\checkmark$\checkmarks\checkmark$\checkmark^\checkmark\\checkmarkc\checkmarki\checkmarkr\checkmarkc\checkmark$\checkmarke\checkmark$\checkmark^\checkmark\\checkmarkc\checkmarki\checkmarkr\checkmarkc\checkmark$\checkmarkc\checkmark$\checkmark^\checkmark\\checkmarkc\checkmarki\checkmarkr\checkmarkc\checkmark$\checkmarkt\checkmark$\checkmark^\checkmark\\checkmarkc\checkmarki\checkmarkr\checkmarkc\checkmark$\checkmarki\checkmark$\checkmark^\checkmark\\checkmarkc\checkmarki\checkmarkr\checkmarkc\checkmark$\checkmarko\checkmark$\checkmark^\checkmark\\checkmarkc\checkmarki\checkmarkr\checkmarkc\checkmark$\checkmarkn\checkmark$\checkmark^\checkmark\\checkmarkc\checkmarki\checkmarkr\checkmarkc\checkmark$\checkmark{\checkmark$\checkmark^\checkmark\\checkmarkc\checkmarki\checkmarkr\checkmarkc\checkmark$\checkmarkV\checkmark$\checkmark^\checkmark\\checkmarkc\checkmarki\checkmarkr\checkmarkc\checkmark$\checkmarké\checkmark$\checkmark^\checkmark\\checkmarkc\checkmarki\checkmarkr\checkmarkc\checkmark$\checkmarkr\checkmark$\checkmark^\checkmark\\checkmarkc\checkmarki\checkmarkr\checkmarkc\checkmark$\checkmarki\checkmark$\checkmark^\checkmark\\checkmarkc\checkmarki\checkmarkr\checkmarkc\checkmark$\checkmarkf\checkmark$\checkmark^\checkmark\\checkmarkc\checkmarki\checkmarkr\checkmarkc\checkmark$\checkmarki\checkmark$\checkmark^\checkmark\\checkmarkc\checkmarki\checkmarkr\checkmarkc\checkmark$\checkmarkc\checkmark$\checkmark^\checkmark\\checkmarkc\checkmarki\checkmarkr\checkmarkc\checkmark$\checkmarka\checkmark$\checkmark^\checkmark\\checkmarkc\checkmarki\checkmarkr\checkmarkc\checkmark$\checkmarkt\checkmark$\checkmark^\checkmark\\checkmarkc\checkmarki\checkmarkr\checkmarkc\checkmark$\checkmarki\checkmark$\checkmark^\checkmark\\checkmarkc\checkmarki\checkmarkr\checkmarkc\checkmark$\checkmarko\checkmark$\checkmark^\checkmark\\checkmarkc\checkmarki\checkmarkr\checkmarkc\checkmark$\checkmarkn\checkmark$\checkmark^\checkmark\\checkmarkc\checkmarki\checkmarkr\checkmarkc\checkmark$\checkmark \checkmark$\checkmark^\checkmark\\checkmarkc\checkmarki\checkmarkr\checkmarkc\checkmark$\checkmarkà\checkmark$\checkmark^\checkmark\\checkmarkc\checkmarki\checkmarkr\checkmarkc\checkmark$\checkmark \checkmark$\checkmark^\checkmark\\checkmarkc\checkmarki\checkmarkr\checkmarkc\checkmark$\checkmarkl\checkmark$\checkmark^\checkmark\\checkmarkc\checkmarki\checkmarkr\checkmarkc\checkmark$\checkmarka\checkmark$\checkmark^\checkmark\\checkmarkc\checkmarki\checkmarkr\checkmarkc\checkmark$\checkmark \checkmark$\checkmark^\checkmark\\checkmarkc\checkmarki\checkmarkr\checkmarkc\checkmark$\checkmarkD\checkmark$\checkmark^\checkmark\\checkmarkc\checkmarki\checkmarkr\checkmarkc\checkmark$\checkmarké\checkmark$\checkmark^\checkmark\\checkmarkc\checkmarki\checkmarkr\checkmarkc\checkmark$\checkmarkf\checkmark$\checkmark^\checkmark\\checkmarkc\checkmarki\checkmarkr\checkmarkc\checkmark$\checkmarkl\checkmark$\checkmark^\checkmark\\checkmarkc\checkmarki\checkmarkr\checkmarkc\checkmark$\checkmarke\checkmark$\checkmark^\checkmark\\checkmarkc\checkmarki\checkmarkr\checkmarkc\checkmark$\checkmarkx\checkmark$\checkmark^\checkmark\\checkmarkc\checkmarki\checkmarkr\checkmarkc\checkmark$\checkmarki\checkmark$\checkmark^\checkmark\\checkmarkc\checkmarki\checkmarkr\checkmarkc\checkmark$\checkmarko\checkmark$\checkmark^\checkmark\\checkmarkc\checkmarki\checkmarkr\checkmarkc\checkmark$\checkmarkn\checkmark$\checkmark^\checkmark\\checkmarkc\checkmarki\checkmarkr\checkmarkc\checkmark$\checkmark \checkmark$\checkmark^\checkmark\\checkmarkc\checkmarki\checkmarkr\checkmarkc\checkmark$\checkmarkd\checkmark$\checkmark^\checkmark\\checkmarkc\checkmarki\checkmarkr\checkmarkc\checkmark$\checkmarku\checkmark$\checkmark^\checkmark\\checkmarkc\checkmarki\checkmarkr\checkmarkc\checkmark$\checkmark \checkmark$\checkmark^\checkmark\\checkmarkc\checkmarki\checkmarkr\checkmarkc\checkmark$\checkmarkP\checkmark$\checkmark^\checkmark\\checkmarkc\checkmarki\checkmarkr\checkmarkc\checkmark$\checkmarko\checkmark$\checkmark^\checkmark\\checkmarkc\checkmarki\checkmarkr\checkmarkc\checkmark$\checkmarkr\checkmark$\checkmark^\checkmark\\checkmarkc\checkmarki\checkmarkr\checkmarkc\checkmark$\checkmarkt\checkmark$\checkmark^\checkmark\\checkmarkc\checkmarki\checkmarkr\checkmarkc\checkmark$\checkmarki\checkmark$\checkmark^\checkmark\\checkmarkc\checkmarki\checkmarkr\checkmarkc\checkmark$\checkmarkq\checkmark$\checkmark^\checkmark\\checkmarkc\checkmarki\checkmarkr\checkmarkc\checkmark$\checkmarku\checkmark$\checkmark^\checkmark\\checkmarkc\checkmarki\checkmarkr\checkmarkc\checkmark$\checkmarke\checkmark$\checkmark^\checkmark\\checkmarkc\checkmarki\checkmarkr\checkmarkc\checkmark$\checkmark \checkmark$\checkmark^\checkmark\\checkmarkc\checkmarki\checkmarkr\checkmarkc\checkmark$\checkmark(\checkmark$\checkmark^\checkmark\\checkmarkc\checkmarki\checkmarkr\checkmarkc\checkmark$\checkmarkA\checkmark$\checkmark^\checkmark\\checkmarkc\checkmarki\checkmarkr\checkmarkc\checkmark$\checkmarkx\checkmark$\checkmark^\checkmark\\checkmarkc\checkmarki\checkmarkr\checkmarkc\checkmark$\checkmarke\checkmark$\checkmark^\checkmark\\checkmarkc\checkmarki\checkmarkr\checkmarkc\checkmark$\checkmark \checkmark$\checkmark^\checkmark\\checkmarkc\checkmarki\checkmarkr\checkmarkc\checkmark$\checkmarkX\checkmark$\checkmark^\checkmark\\checkmarkc\checkmarki\checkmarkr\checkmarkc\checkmark$\checkmark)\checkmark$\checkmark^\checkmark\\checkmarkc\checkmarki\checkmarkr\checkmarkc\checkmark$\checkmark}\checkmark$\checkmark^\checkmark\\checkmarkc\checkmarki\checkmarkr\checkmarkc\checkmark$\checkmark
\checkmark$\checkmark^\checkmark\\checkmarkc\checkmarki\checkmarkr\checkmarkc\checkmark$\checkmarkL\checkmark$\checkmark^\checkmark\\checkmarkc\checkmarki\checkmarkr\checkmarkc\checkmark$\checkmarke\checkmark$\checkmark^\checkmark\\checkmarkc\checkmarki\checkmarkr\checkmarkc\checkmark$\checkmark \checkmark$\checkmark^\checkmark\\checkmarkc\checkmarki\checkmarkr\checkmarkc\checkmark$\checkmarkp\checkmark$\checkmark^\checkmark\\checkmarkc\checkmarki\checkmarkr\checkmarkc\checkmark$\checkmarko\checkmark$\checkmark^\checkmark\\checkmarkc\checkmarki\checkmarkr\checkmarkc\checkmark$\checkmarkr\checkmark$\checkmark^\checkmark\\checkmarkc\checkmarki\checkmarkr\checkmarkc\checkmark$\checkmarkt\checkmark$\checkmark^\checkmark\\checkmarkc\checkmarki\checkmarkr\checkmarkc\checkmark$\checkmarki\checkmark$\checkmark^\checkmark\\checkmarkc\checkmarki\checkmarkr\checkmarkc\checkmark$\checkmarkq\checkmark$\checkmark^\checkmark\\checkmarkc\checkmarki\checkmarkr\checkmarkc\checkmark$\checkmarku\checkmark$\checkmark^\checkmark\\checkmarkc\checkmarki\checkmarkr\checkmarkc\checkmark$\checkmarke\checkmark$\checkmark^\checkmark\\checkmarkc\checkmarki\checkmarkr\checkmarkc\checkmark$\checkmark \checkmark$\checkmark^\checkmark\\checkmarkc\checkmarki\checkmarkr\checkmarkc\checkmark$\checkmarke\checkmark$\checkmark^\checkmark\\checkmarkc\checkmarki\checkmarkr\checkmarkc\checkmark$\checkmarkn\checkmark$\checkmark^\checkmark\\checkmarkc\checkmarki\checkmarkr\checkmarkc\checkmark$\checkmark \checkmark$\checkmark^\checkmark\\checkmarkc\checkmarki\checkmarkr\checkmarkc\checkmark$\checkmarkX\checkmark$\checkmark^\checkmark\\checkmarkc\checkmarki\checkmarkr\checkmarkc\checkmark$\checkmark \checkmark$\checkmark^\checkmark\\checkmarkc\checkmarki\checkmarkr\checkmarkc\checkmark$\checkmark(\checkmark$\checkmark^\checkmark\\checkmarkc\checkmarki\checkmarkr\checkmarkc\checkmark$\checkmarkl\checkmark$\checkmark^\checkmark\\checkmarkc\checkmarki\checkmarkr\checkmarkc\checkmark$\checkmarko\checkmark$\checkmark^\checkmark\\checkmarkc\checkmarki\checkmarkr\checkmarkc\checkmark$\checkmarkn\checkmark$\checkmark^\checkmark\\checkmarkc\checkmarki\checkmarkr\checkmarkc\checkmark$\checkmarkg\checkmark$\checkmark^\checkmark\\checkmarkc\checkmarki\checkmarkr\checkmarkc\checkmark$\checkmarku\checkmark$\checkmark^\checkmark\\checkmarkc\checkmarki\checkmarkr\checkmarkc\checkmark$\checkmarke\checkmark$\checkmark^\checkmark\\checkmarkc\checkmarki\checkmarkr\checkmarkc\checkmark$\checkmarku\checkmark$\checkmark^\checkmark\\checkmarkc\checkmarki\checkmarkr\checkmarkc\checkmark$\checkmarkr\checkmark$\checkmark^\checkmark\\checkmarkc\checkmarki\checkmarkr\checkmarkc\checkmark$\checkmark \checkmark$\checkmark^\checkmark\\checkmarkc\checkmarki\checkmarkr\checkmarkc\checkmark$\checkmark$\checkmark$\checkmark^\checkmark\\checkmarkc\checkmarki\checkmarkr\checkmarkc\checkmark$\checkmarkL\checkmark$\checkmark^\checkmark\\checkmarkc\checkmarki\checkmarkr\checkmarkc\checkmark$\checkmark \checkmark$\checkmark^\checkmark\\checkmarkc\checkmarki\checkmarkr\checkmarkc\checkmark$\checkmark=\checkmark$\checkmark^\checkmark\\checkmarkc\checkmarki\checkmarkr\checkmarkc\checkmark$\checkmark \checkmark$\checkmark^\checkmark\\checkmarkc\checkmarki\checkmarkr\checkmarkc\checkmark$\checkmark4\checkmark$\checkmark^\checkmark\\checkmarkc\checkmarki\checkmarkr\checkmarkc\checkmark$\checkmark0\checkmark$\checkmark^\checkmark\\checkmarkc\checkmarki\checkmarkr\checkmarkc\checkmark$\checkmark0\checkmark$\checkmark^\checkmark\\checkmarkc\checkmarki\checkmarkr\checkmarkc\checkmark$\checkmark$\checkmark$\checkmark^\checkmark\\checkmarkc\checkmarki\checkmarkr\checkmarkc\checkmark$\checkmark \checkmark$\checkmark^\checkmark\\checkmarkc\checkmarki\checkmarkr\checkmarkc\checkmark$\checkmarkm\checkmark$\checkmark^\checkmark\\checkmarkc\checkmarki\checkmarkr\checkmarkc\checkmark$\checkmarkm\checkmark$\checkmark^\checkmark\\checkmarkc\checkmarki\checkmarkr\checkmarkc\checkmark$\checkmark)\checkmark$\checkmark^\checkmark\\checkmarkc\checkmarki\checkmarkr\checkmarkc\checkmark$\checkmark \checkmark$\checkmark^\checkmark\\checkmarkc\checkmarki\checkmarkr\checkmarkc\checkmark$\checkmarke\checkmark$\checkmark^\checkmark\\checkmarkc\checkmarki\checkmarkr\checkmarkc\checkmark$\checkmarks\checkmark$\checkmark^\checkmark\\checkmarkc\checkmarki\checkmarkr\checkmarkc\checkmark$\checkmarkt\checkmark$\checkmark^\checkmark\\checkmarkc\checkmarki\checkmarkr\checkmarkc\checkmark$\checkmark \checkmark$\checkmark^\checkmark\\checkmarkc\checkmarki\checkmarkr\checkmarkc\checkmark$\checkmarka\checkmark$\checkmark^\checkmark\\checkmarkc\checkmarki\checkmarkr\checkmarkc\checkmark$\checkmarks\checkmark$\checkmark^\checkmark\\checkmarkc\checkmarki\checkmarkr\checkmarkc\checkmark$\checkmarks\checkmark$\checkmark^\checkmark\\checkmarkc\checkmarki\checkmarkr\checkmarkc\checkmark$\checkmarki\checkmark$\checkmark^\checkmark\\checkmarkc\checkmarki\checkmarkr\checkmarkc\checkmark$\checkmarkm\checkmark$\checkmark^\checkmark\\checkmarkc\checkmarki\checkmarkr\checkmarkc\checkmark$\checkmarki\checkmark$\checkmark^\checkmark\\checkmarkc\checkmarki\checkmarkr\checkmarkc\checkmark$\checkmarkl\checkmark$\checkmark^\checkmark\\checkmarkc\checkmarki\checkmarkr\checkmarkc\checkmark$\checkmarké\checkmark$\checkmark^\checkmark\\checkmarkc\checkmarki\checkmarkr\checkmarkc\checkmark$\checkmark \checkmark$\checkmark^\checkmark\\checkmarkc\checkmarki\checkmarkr\checkmarkc\checkmark$\checkmarkà\checkmark$\checkmark^\checkmark\\checkmarkc\checkmarki\checkmarkr\checkmarkc\checkmark$\checkmark \checkmark$\checkmark^\checkmark\\checkmarkc\checkmarki\checkmarkr\checkmarkc\checkmark$\checkmarku\checkmark$\checkmark^\checkmark\\checkmarkc\checkmarki\checkmarkr\checkmarkc\checkmark$\checkmarkn\checkmark$\checkmark^\checkmark\\checkmarkc\checkmarki\checkmarkr\checkmarkc\checkmark$\checkmarke\checkmark$\checkmark^\checkmark\\checkmarkc\checkmarki\checkmarkr\checkmarkc\checkmark$\checkmark \checkmark$\checkmark^\checkmark\\checkmarkc\checkmarki\checkmarkr\checkmarkc\checkmark$\checkmarkp\checkmark$\checkmark^\checkmark\\checkmarkc\checkmarki\checkmarkr\checkmarkc\checkmark$\checkmarko\checkmark$\checkmark^\checkmark\\checkmarkc\checkmarki\checkmarkr\checkmarkc\checkmark$\checkmarku\checkmark$\checkmark^\checkmark\\checkmarkc\checkmarki\checkmarkr\checkmarkc\checkmark$\checkmarkt\checkmark$\checkmark^\checkmark\\checkmarkc\checkmarki\checkmarkr\checkmarkc\checkmark$\checkmarkr\checkmark$\checkmark^\checkmark\\checkmarkc\checkmarki\checkmarkr\checkmarkc\checkmark$\checkmarke\checkmark$\checkmark^\checkmark\\checkmarkc\checkmarki\checkmarkr\checkmarkc\checkmark$\checkmark \checkmark$\checkmark^\checkmark\\checkmarkc\checkmarki\checkmarkr\checkmarkc\checkmark$\checkmarks\checkmark$\checkmark^\checkmark\\checkmarkc\checkmarki\checkmarkr\checkmarkc\checkmark$\checkmarku\checkmark$\checkmark^\checkmark\\checkmarkc\checkmarki\checkmarkr\checkmarkc\checkmark$\checkmarkr\checkmark$\checkmark^\checkmark\\checkmarkc\checkmarki\checkmarkr\checkmarkc\checkmark$\checkmark \checkmark$\checkmark^\checkmark\\checkmarkc\checkmarki\checkmarkr\checkmarkc\checkmark$\checkmarkd\checkmark$\checkmark^\checkmark\\checkmarkc\checkmarki\checkmarkr\checkmarkc\checkmark$\checkmarke\checkmark$\checkmark^\checkmark\\checkmarkc\checkmarki\checkmarkr\checkmarkc\checkmark$\checkmarku\checkmark$\checkmark^\checkmark\\checkmarkc\checkmarki\checkmarkr\checkmarkc\checkmark$\checkmarkx\checkmark$\checkmark^\checkmark\\checkmarkc\checkmarki\checkmarkr\checkmarkc\checkmark$\checkmark \checkmark$\checkmark^\checkmark\\checkmarkc\checkmarki\checkmarkr\checkmarkc\checkmark$\checkmarka\checkmark$\checkmark^\checkmark\\checkmarkc\checkmarki\checkmarkr\checkmarkc\checkmark$\checkmarkp\checkmark$\checkmark^\checkmark\\checkmarkc\checkmarki\checkmarkr\checkmarkc\checkmark$\checkmarkp\checkmark$\checkmark^\checkmark\\checkmarkc\checkmarki\checkmarkr\checkmarkc\checkmark$\checkmarku\checkmark$\checkmark^\checkmark\\checkmarkc\checkmarki\checkmarkr\checkmarkc\checkmark$\checkmarki\checkmark$\checkmark^\checkmark\\checkmarkc\checkmarki\checkmarkr\checkmarkc\checkmark$\checkmarks\checkmark$\checkmark^\checkmark\\checkmarkc\checkmarki\checkmarkr\checkmarkc\checkmark$\checkmark \checkmark$\checkmark^\checkmark\\checkmarkc\checkmarki\checkmarkr\checkmarkc\checkmark$\checkmarks\checkmark$\checkmark^\checkmark\\checkmarkc\checkmarki\checkmarkr\checkmarkc\checkmark$\checkmarki\checkmark$\checkmark^\checkmark\\checkmarkc\checkmarki\checkmarkr\checkmarkc\checkmark$\checkmarkm\checkmark$\checkmark^\checkmark\\checkmarkc\checkmarki\checkmarkr\checkmarkc\checkmark$\checkmarkp\checkmark$\checkmark^\checkmark\\checkmarkc\checkmarki\checkmarkr\checkmarkc\checkmark$\checkmarkl\checkmark$\checkmark^\checkmark\\checkmarkc\checkmarki\checkmarkr\checkmarkc\checkmark$\checkmarke\checkmark$\checkmark^\checkmark\\checkmarkc\checkmarki\checkmarkr\checkmarkc\checkmark$\checkmarks\checkmark$\checkmark^\checkmark\\checkmarkc\checkmarki\checkmarkr\checkmarkc\checkmark$\checkmark \checkmark$\checkmark^\checkmark\\checkmarkc\checkmarki\checkmarkr\checkmarkc\checkmark$\checkmarks\checkmark$\checkmark^\checkmark\\checkmarkc\checkmarki\checkmarkr\checkmarkc\checkmark$\checkmarko\checkmark$\checkmark^\checkmark\\checkmarkc\checkmarki\checkmarkr\checkmarkc\checkmark$\checkmarku\checkmark$\checkmark^\checkmark\\checkmarkc\checkmarki\checkmarkr\checkmarkc\checkmark$\checkmarkm\checkmark$\checkmark^\checkmark\\checkmarkc\checkmarki\checkmarkr\checkmarkc\checkmark$\checkmarki\checkmark$\checkmark^\checkmark\\checkmarkc\checkmarki\checkmarkr\checkmarkc\checkmark$\checkmarks\checkmark$\checkmark^\checkmark\\checkmarkc\checkmarki\checkmarkr\checkmarkc\checkmark$\checkmarke\checkmark$\checkmark^\checkmark\\checkmarkc\checkmarki\checkmarkr\checkmarkc\checkmark$\checkmark \checkmark$\checkmark^\checkmark\\checkmarkc\checkmarki\checkmarkr\checkmarkc\checkmark$\checkmarkà\checkmark$\checkmark^\checkmark\\checkmarkc\checkmarki\checkmarkr\checkmarkc\checkmark$\checkmark \checkmark$\checkmark^\checkmark\\checkmarkc\checkmarki\checkmarkr\checkmarkc\checkmark$\checkmarku\checkmark$\checkmark^\checkmark\\checkmarkc\checkmarki\checkmarkr\checkmarkc\checkmark$\checkmarkn\checkmark$\checkmark^\checkmark\\checkmarkc\checkmarki\checkmarkr\checkmarkc\checkmark$\checkmarke\checkmark$\checkmark^\checkmark\\checkmarkc\checkmarki\checkmarkr\checkmarkc\checkmark$\checkmark \checkmark$\checkmark^\checkmark\\checkmarkc\checkmarki\checkmarkr\checkmarkc\checkmark$\checkmarkc\checkmark$\checkmark^\checkmark\\checkmarkc\checkmarki\checkmarkr\checkmarkc\checkmark$\checkmarkh\checkmark$\checkmark^\checkmark\\checkmarkc\checkmarki\checkmarkr\checkmarkc\checkmark$\checkmarka\checkmark$\checkmark^\checkmark\\checkmarkc\checkmarki\checkmarkr\checkmarkc\checkmark$\checkmarkr\checkmark$\checkmark^\checkmark\\checkmarkc\checkmarki\checkmarkr\checkmarkc\checkmark$\checkmarkg\checkmark$\checkmark^\checkmark\\checkmarkc\checkmarki\checkmarkr\checkmarkc\checkmark$\checkmarke\checkmark$\checkmark^\checkmark\\checkmarkc\checkmarki\checkmarkr\checkmarkc\checkmark$\checkmark \checkmark$\checkmark^\checkmark\\checkmarkc\checkmarki\checkmarkr\checkmarkc\checkmark$\checkmarkc\checkmark$\checkmark^\checkmark\\checkmarkc\checkmarki\checkmarkr\checkmarkc\checkmark$\checkmarke\checkmark$\checkmark^\checkmark\\checkmarkc\checkmarki\checkmarkr\checkmarkc\checkmark$\checkmarkn\checkmark$\checkmark^\checkmark\\checkmarkc\checkmarki\checkmarkr\checkmarkc\checkmark$\checkmarkt\checkmark$\checkmark^\checkmark\\checkmarkc\checkmarki\checkmarkr\checkmarkc\checkmark$\checkmarkr\checkmark$\checkmark^\checkmark\\checkmarkc\checkmarki\checkmarkr\checkmarkc\checkmark$\checkmarka\checkmark$\checkmark^\checkmark\\checkmarkc\checkmarki\checkmarkr\checkmarkc\checkmark$\checkmarkl\checkmark$\checkmark^\checkmark\\checkmarkc\checkmarki\checkmarkr\checkmarkc\checkmark$\checkmarke\checkmark$\checkmark^\checkmark\\checkmarkc\checkmarki\checkmarkr\checkmarkc\checkmark$\checkmark \checkmark$\checkmark^\checkmark\\checkmarkc\checkmarki\checkmarkr\checkmarkc\checkmark$\checkmark$\checkmark$\checkmark^\checkmark\\checkmarkc\checkmarki\checkmarkr\checkmarkc\checkmark$\checkmarkF\checkmark$\checkmark^\checkmark\\checkmarkc\checkmarki\checkmarkr\checkmarkc\checkmark$\checkmark_\checkmark$\checkmark^\checkmark\\checkmarkc\checkmarki\checkmarkr\checkmarkc\checkmark$\checkmark{\checkmark$\checkmark^\checkmark\\checkmarkc\checkmarki\checkmarkr\checkmarkc\checkmark$\checkmarkm\checkmark$\checkmark^\checkmark\\checkmarkc\checkmarki\checkmarkr\checkmarkc\checkmark$\checkmarka\checkmark$\checkmark^\checkmark\\checkmarkc\checkmarki\checkmarkr\checkmarkc\checkmark$\checkmarkx\checkmark$\checkmark^\checkmark\\checkmarkc\checkmarki\checkmarkr\checkmarkc\checkmark$\checkmark}\checkmark$\checkmark^\checkmark\\checkmarkc\checkmarki\checkmarkr\checkmarkc\checkmark$\checkmark \checkmark$\checkmark^\checkmark\\checkmarkc\checkmarki\checkmarkr\checkmarkc\checkmark$\checkmark=\checkmark$\checkmark^\checkmark\\checkmarkc\checkmarki\checkmarkr\checkmarkc\checkmark$\checkmark \checkmark$\checkmark^\checkmark\\checkmarkc\checkmarki\checkmarkr\checkmarkc\checkmark$\checkmark1\checkmark$\checkmark^\checkmark\\checkmarkc\checkmarki\checkmarkr\checkmarkc\checkmark$\checkmark4\checkmark$\checkmark^\checkmark\\checkmarkc\checkmarki\checkmarkr\checkmarkc\checkmark$\checkmark0\checkmark$\checkmark^\checkmark\\checkmarkc\checkmarki\checkmarkr\checkmarkc\checkmark$\checkmark$\checkmark$\checkmark^\checkmark\\checkmarkc\checkmarki\checkmarkr\checkmarkc\checkmark$\checkmark \checkmark$\checkmark^\checkmark\\checkmarkc\checkmarki\checkmarkr\checkmarkc\checkmark$\checkmarkN\checkmark$\checkmark^\checkmark\\checkmarkc\checkmarki\checkmarkr\checkmarkc\checkmark$\checkmark \checkmark$\checkmark^\checkmark\\checkmarkc\checkmarki\checkmarkr\checkmarkc\checkmark$\checkmark(\checkmark$\checkmark^\checkmark\\checkmarkc\checkmarki\checkmarkr\checkmarkc\checkmark$\checkmarkf\checkmark$\checkmark^\checkmark\\checkmarkc\checkmarki\checkmarkr\checkmarkc\checkmark$\checkmarko\checkmark$\checkmark^\checkmark\\checkmarkc\checkmarki\checkmarkr\checkmarkc\checkmark$\checkmarkr\checkmark$\checkmark^\checkmark\\checkmarkc\checkmarki\checkmarkr\checkmarkc\checkmark$\checkmarkc\checkmark$\checkmark^\checkmark\\checkmarkc\checkmarki\checkmarkr\checkmarkc\checkmark$\checkmarke\checkmark$\checkmark^\checkmark\\checkmarkc\checkmarki\checkmarkr\checkmarkc\checkmark$\checkmark \checkmark$\checkmark^\checkmark\\checkmarkc\checkmarki\checkmarkr\checkmarkc\checkmark$\checkmarkd\checkmark$\checkmark^\checkmark\\checkmarkc\checkmarki\checkmarkr\checkmarkc\checkmark$\checkmarke\checkmark$\checkmark^\checkmark\\checkmarkc\checkmarki\checkmarkr\checkmarkc\checkmark$\checkmark \checkmark$\checkmark^\checkmark\\checkmarkc\checkmarki\checkmarkr\checkmarkc\checkmark$\checkmarkc\checkmark$\checkmark^\checkmark\\checkmarkc\checkmarki\checkmarkr\checkmarkc\checkmark$\checkmarko\checkmark$\checkmark^\checkmark\\checkmarkc\checkmarki\checkmarkr\checkmarkc\checkmark$\checkmarku\checkmark$\checkmark^\checkmark\\checkmarkc\checkmarki\checkmarkr\checkmarkc\checkmark$\checkmarkp\checkmark$\checkmark^\checkmark\\checkmarkc\checkmarki\checkmarkr\checkmarkc\checkmark$\checkmarke\checkmark$\checkmark^\checkmark\\checkmarkc\checkmarki\checkmarkr\checkmarkc\checkmark$\checkmark \checkmark$\checkmark^\checkmark\\checkmarkc\checkmarki\checkmarkr\checkmarkc\checkmark$\checkmark+\checkmark$\checkmark^\checkmark\\checkmarkc\checkmarki\checkmarkr\checkmarkc\checkmark$\checkmark \checkmark$\checkmark^\checkmark\\checkmarkc\checkmarki\checkmarkr\checkmarkc\checkmark$\checkmarkp\checkmark$\checkmark^\checkmark\\checkmarkc\checkmarki\checkmarkr\checkmarkc\checkmark$\checkmarko\checkmark$\checkmark^\checkmark\\checkmarkc\checkmarki\checkmarkr\checkmarkc\checkmark$\checkmarki\checkmark$\checkmark^\checkmark\\checkmarkc\checkmarki\checkmarkr\checkmarkc\checkmark$\checkmarkd\checkmark$\checkmark^\checkmark\\checkmarkc\checkmarki\checkmarkr\checkmarkc\checkmark$\checkmarks\checkmark$\checkmark^\checkmark\\checkmarkc\checkmarki\checkmarkr\checkmarkc\checkmark$\checkmark \checkmark$\checkmark^\checkmark\\checkmarkc\checkmarki\checkmarkr\checkmarkc\checkmark$\checkmarkd\checkmark$\checkmark^\checkmark\\checkmarkc\checkmarki\checkmarkr\checkmarkc\checkmark$\checkmarke\checkmark$\checkmark^\checkmark\\checkmarkc\checkmarki\checkmarkr\checkmarkc\checkmark$\checkmark \checkmark$\checkmark^\checkmark\\checkmarkc\checkmarki\checkmarkr\checkmarkc\checkmark$\checkmarkl\checkmark$\checkmark^\checkmark\\checkmarkc\checkmarki\checkmarkr\checkmarkc\checkmark$\checkmarka\checkmark$\checkmark^\checkmark\\checkmarkc\checkmarki\checkmarkr\checkmarkc\checkmark$\checkmark \checkmark$\checkmark^\checkmark\\checkmarkc\checkmarki\checkmarkr\checkmarkc\checkmark$\checkmarkb\checkmark$\checkmark^\checkmark\\checkmarkc\checkmarki\checkmarkr\checkmarkc\checkmark$\checkmarkr\checkmark$\checkmark^\checkmark\\checkmarkc\checkmarki\checkmarkr\checkmarkc\checkmark$\checkmarko\checkmark$\checkmark^\checkmark\\checkmarkc\checkmarki\checkmarkr\checkmarkc\checkmark$\checkmarkc\checkmark$\checkmark^\checkmark\\checkmarkc\checkmarki\checkmarkr\checkmarkc\checkmark$\checkmarkh\checkmark$\checkmark^\checkmark\\checkmarkc\checkmarki\checkmarkr\checkmarkc\checkmark$\checkmarke\checkmark$\checkmark^\checkmark\\checkmarkc\checkmarki\checkmarkr\checkmarkc\checkmark$\checkmark)\checkmark$\checkmark^\checkmark\\checkmarkc\checkmarki\checkmarkr\checkmarkc\checkmark$\checkmark.\checkmark$\checkmark^\checkmark\\checkmarkc\checkmarki\checkmarkr\checkmarkc\checkmark$\checkmark
\checkmark$\checkmark^\checkmark\\checkmarkc\checkmarki\checkmarkr\checkmarkc\checkmark$\checkmark
\checkmark$\checkmark^\checkmark\\checkmarkc\checkmarki\checkmarkr\checkmarkc\checkmark$\checkmarkL\checkmark$\checkmark^\checkmark\\checkmarkc\checkmarki\checkmarkr\checkmarkc\checkmark$\checkmarka\checkmark$\checkmark^\checkmark\\checkmarkc\checkmarki\checkmarkr\checkmarkc\checkmark$\checkmark \checkmark$\checkmark^\checkmark\\checkmarkc\checkmarki\checkmarkr\checkmarkc\checkmark$\checkmarkf\checkmark$\checkmark^\checkmark\\checkmarkc\checkmarki\checkmarkr\checkmarkc\checkmark$\checkmarkl\checkmark$\checkmark^\checkmark\\checkmarkc\checkmarki\checkmarkr\checkmarkc\checkmark$\checkmarkè\checkmark$\checkmark^\checkmark\\checkmarkc\checkmarki\checkmarkr\checkmarkc\checkmark$\checkmarkc\checkmark$\checkmark^\checkmark\\checkmarkc\checkmarki\checkmarkr\checkmarkc\checkmark$\checkmarkh\checkmark$\checkmark^\checkmark\\checkmarkc\checkmarki\checkmarkr\checkmarkc\checkmark$\checkmarke\checkmark$\checkmark^\checkmark\\checkmarkc\checkmarki\checkmarkr\checkmarkc\checkmark$\checkmark \checkmark$\checkmark^\checkmark\\checkmarkc\checkmarki\checkmarkr\checkmarkc\checkmark$\checkmarkm\checkmark$\checkmark^\checkmark\\checkmarkc\checkmarki\checkmarkr\checkmarkc\checkmark$\checkmarka\checkmark$\checkmark^\checkmark\\checkmarkc\checkmarki\checkmarkr\checkmarkc\checkmark$\checkmarkx\checkmark$\checkmark^\checkmark\\checkmarkc\checkmarki\checkmarkr\checkmarkc\checkmark$\checkmarki\checkmark$\checkmark^\checkmark\\checkmarkc\checkmarki\checkmarkr\checkmarkc\checkmark$\checkmarkm\checkmark$\checkmark^\checkmark\\checkmarkc\checkmarki\checkmarkr\checkmarkc\checkmark$\checkmarka\checkmark$\checkmark^\checkmark\\checkmarkc\checkmarki\checkmarkr\checkmarkc\checkmark$\checkmarkl\checkmark$\checkmark^\checkmark\\checkmarkc\checkmarki\checkmarkr\checkmarkc\checkmark$\checkmarke\checkmark$\checkmark^\checkmark\\checkmarkc\checkmarki\checkmarkr\checkmarkc\checkmark$\checkmark \checkmark$\checkmark^\checkmark\\checkmarkc\checkmarki\checkmarkr\checkmarkc\checkmark$\checkmarka\checkmark$\checkmark^\checkmark\\checkmarkc\checkmarki\checkmarkr\checkmarkc\checkmark$\checkmarku\checkmark$\checkmark^\checkmark\\checkmarkc\checkmarki\checkmarkr\checkmarkc\checkmark$\checkmark \checkmark$\checkmark^\checkmark\\checkmarkc\checkmarki\checkmarkr\checkmarkc\checkmark$\checkmarkc\checkmark$\checkmark^\checkmark\\checkmarkc\checkmarki\checkmarkr\checkmarkc\checkmark$\checkmarke\checkmark$\checkmark^\checkmark\\checkmarkc\checkmarki\checkmarkr\checkmarkc\checkmark$\checkmarkn\checkmark$\checkmark^\checkmark\\checkmarkc\checkmarki\checkmarkr\checkmarkc\checkmark$\checkmarkt\checkmark$\checkmark^\checkmark\\checkmarkc\checkmarki\checkmarkr\checkmarkc\checkmark$\checkmarkr\checkmark$\checkmark^\checkmark\\checkmarkc\checkmarki\checkmarkr\checkmarkc\checkmark$\checkmarke\checkmark$\checkmark^\checkmark\\checkmarkc\checkmarki\checkmarkr\checkmarkc\checkmark$\checkmark \checkmark$\checkmark^\checkmark\\checkmarkc\checkmarki\checkmarkr\checkmarkc\checkmark$\checkmarkd\checkmark$\checkmark^\checkmark\\checkmarkc\checkmarki\checkmarkr\checkmarkc\checkmark$\checkmarke\checkmark$\checkmark^\checkmark\\checkmarkc\checkmarki\checkmarkr\checkmarkc\checkmark$\checkmark \checkmark$\checkmark^\checkmark\\checkmarkc\checkmarki\checkmarkr\checkmarkc\checkmark$\checkmarkl\checkmark$\checkmark^\checkmark\\checkmarkc\checkmarki\checkmarkr\checkmarkc\checkmark$\checkmarka\checkmark$\checkmark^\checkmark\\checkmarkc\checkmarki\checkmarkr\checkmarkc\checkmark$\checkmark \checkmark$\checkmark^\checkmark\\checkmarkc\checkmarki\checkmarkr\checkmarkc\checkmark$\checkmarkp\checkmark$\checkmark^\checkmark\\checkmarkc\checkmarki\checkmarkr\checkmarkc\checkmark$\checkmarko\checkmark$\checkmark^\checkmark\\checkmarkc\checkmarki\checkmarkr\checkmarkc\checkmark$\checkmarku\checkmark$\checkmark^\checkmark\\checkmarkc\checkmarki\checkmarkr\checkmarkc\checkmark$\checkmarkt\checkmark$\checkmark^\checkmark\\checkmarkc\checkmarki\checkmarkr\checkmarkc\checkmark$\checkmarkr\checkmark$\checkmark^\checkmark\\checkmarkc\checkmarki\checkmarkr\checkmarkc\checkmark$\checkmarke\checkmark$\checkmark^\checkmark\\checkmarkc\checkmarki\checkmarkr\checkmarkc\checkmark$\checkmark \checkmark$\checkmark^\checkmark\\checkmarkc\checkmarki\checkmarkr\checkmarkc\checkmark$\checkmark(\checkmark$\checkmark^\checkmark\\checkmarkc\checkmarki\checkmarkr\checkmarkc\checkmark$\checkmarkt\checkmark$\checkmark^\checkmark\\checkmarkc\checkmarki\checkmarkr\checkmarkc\checkmark$\checkmarkh\checkmark$\checkmark^\checkmark\\checkmarkc\checkmarki\checkmarkr\checkmarkc\checkmark$\checkmarké\checkmark$\checkmark^\checkmark\\checkmarkc\checkmarki\checkmarkr\checkmarkc\checkmark$\checkmarko\checkmark$\checkmark^\checkmark\\checkmarkc\checkmarki\checkmarkr\checkmarkc\checkmark$\checkmarkr\checkmark$\checkmark^\checkmark\\checkmarkc\checkmarki\checkmarkr\checkmarkc\checkmark$\checkmarkè\checkmark$\checkmark^\checkmark\\checkmarkc\checkmarki\checkmarkr\checkmarkc\checkmark$\checkmarkm\checkmark$\checkmark^\checkmark\\checkmarkc\checkmarki\checkmarkr\checkmarkc\checkmark$\checkmarke\checkmark$\checkmark^\checkmark\\checkmarkc\checkmarki\checkmarkr\checkmarkc\checkmark$\checkmark \checkmark$\checkmark^\checkmark\\checkmarkc\checkmarki\checkmarkr\checkmarkc\checkmark$\checkmarkd\checkmark$\checkmark^\checkmark\\checkmarkc\checkmarki\checkmarkr\checkmarkc\checkmark$\checkmarke\checkmark$\checkmark^\checkmark\\checkmarkc\checkmarki\checkmarkr\checkmarkc\checkmark$\checkmark \checkmark$\checkmark^\checkmark\\checkmarkc\checkmarki\checkmarkr\checkmarkc\checkmark$\checkmarkl\checkmark$\checkmark^\checkmark\\checkmarkc\checkmarki\checkmarkr\checkmarkc\checkmark$\checkmarka\checkmark$\checkmark^\checkmark\\checkmarkc\checkmarki\checkmarkr\checkmarkc\checkmark$\checkmark \checkmark$\checkmark^\checkmark\\checkmarkc\checkmarki\checkmarkr\checkmarkc\checkmark$\checkmarkd\checkmark$\checkmark^\checkmark\\checkmarkc\checkmarki\checkmarkr\checkmarkc\checkmark$\checkmarké\checkmark$\checkmark^\checkmark\\checkmarkc\checkmarki\checkmarkr\checkmarkc\checkmark$\checkmarkf\checkmark$\checkmark^\checkmark\\checkmarkc\checkmarki\checkmarkr\checkmarkc\checkmark$\checkmarkl\checkmark$\checkmark^\checkmark\\checkmarkc\checkmarki\checkmarkr\checkmarkc\checkmark$\checkmarke\checkmark$\checkmark^\checkmark\\checkmarkc\checkmarki\checkmarkr\checkmarkc\checkmark$\checkmarkx\checkmark$\checkmark^\checkmark\\checkmarkc\checkmarki\checkmarkr\checkmarkc\checkmark$\checkmarki\checkmark$\checkmark^\checkmark\\checkmarkc\checkmarki\checkmarkr\checkmarkc\checkmark$\checkmarko\checkmark$\checkmark^\checkmark\\checkmarkc\checkmarki\checkmarkr\checkmarkc\checkmark$\checkmarkn\checkmark$\checkmark^\checkmark\\checkmarkc\checkmarki\checkmarkr\checkmarkc\checkmark$\checkmark)\checkmark$\checkmark^\checkmark\\checkmarkc\checkmarki\checkmarkr\checkmarkc\checkmark$\checkmark \checkmark$\checkmark^\checkmark\\checkmarkc\checkmarki\checkmarkr\checkmarkc\checkmark$\checkmark:\checkmark$\checkmark^\checkmark\\checkmarkc\checkmarki\checkmarkr\checkmarkc\checkmark$\checkmark
\checkmark$\checkmark^\checkmark\\checkmarkc\checkmarki\checkmarkr\checkmarkc\checkmark$\checkmark\\checkmark$\checkmark^\checkmark\\checkmarkc\checkmarki\checkmarkr\checkmarkc\checkmark$\checkmarkb\checkmark$\checkmark^\checkmark\\checkmarkc\checkmarki\checkmarkr\checkmarkc\checkmark$\checkmarke\checkmark$\checkmark^\checkmark\\checkmarkc\checkmarki\checkmarkr\checkmarkc\checkmark$\checkmarkg\checkmark$\checkmark^\checkmark\\checkmarkc\checkmarki\checkmarkr\checkmarkc\checkmark$\checkmarki\checkmark$\checkmark^\checkmark\\checkmarkc\checkmarki\checkmarkr\checkmarkc\checkmark$\checkmarkn\checkmark$\checkmark^\checkmark\\checkmarkc\checkmarki\checkmarkr\checkmarkc\checkmark$\checkmark{\checkmark$\checkmark^\checkmark\\checkmarkc\checkmarki\checkmarkr\checkmarkc\checkmark$\checkmarke\checkmark$\checkmark^\checkmark\\checkmarkc\checkmarki\checkmarkr\checkmarkc\checkmark$\checkmarkq\checkmark$\checkmark^\checkmark\\checkmarkc\checkmarki\checkmarkr\checkmarkc\checkmark$\checkmarku\checkmark$\checkmark^\checkmark\\checkmarkc\checkmarki\checkmarkr\checkmarkc\checkmark$\checkmarka\checkmark$\checkmark^\checkmark\\checkmarkc\checkmarki\checkmarkr\checkmarkc\checkmark$\checkmarkt\checkmark$\checkmark^\checkmark\\checkmarkc\checkmarki\checkmarkr\checkmarkc\checkmark$\checkmarki\checkmark$\checkmark^\checkmark\\checkmarkc\checkmarki\checkmarkr\checkmarkc\checkmark$\checkmarko\checkmark$\checkmark^\checkmark\\checkmarkc\checkmarki\checkmarkr\checkmarkc\checkmark$\checkmarkn\checkmark$\checkmark^\checkmark\\checkmarkc\checkmarki\checkmarkr\checkmarkc\checkmark$\checkmark}\checkmark$\checkmark^\checkmark\\checkmarkc\checkmarki\checkmarkr\checkmarkc\checkmark$\checkmark
\checkmark$\checkmark^\checkmark\\checkmarkc\checkmarki\checkmarkr\checkmarkc\checkmark$\checkmark \checkmark$\checkmark^\checkmark\\checkmarkc\checkmarki\checkmarkr\checkmarkc\checkmark$\checkmark \checkmark$\checkmark^\checkmark\\checkmarkc\checkmarki\checkmarkr\checkmarkc\checkmark$\checkmark \checkmark$\checkmark^\checkmark\\checkmarkc\checkmarki\checkmarkr\checkmarkc\checkmark$\checkmark \checkmark$\checkmark^\checkmark\\checkmarkc\checkmarki\checkmarkr\checkmarkc\checkmark$\checkmarkf\checkmark$\checkmark^\checkmark\\checkmarkc\checkmarki\checkmarkr\checkmarkc\checkmark$\checkmark_\checkmark$\checkmark^\checkmark\\checkmarkc\checkmarki\checkmarkr\checkmarkc\checkmark$\checkmark{\checkmark$\checkmark^\checkmark\\checkmarkc\checkmarki\checkmarkr\checkmarkc\checkmark$\checkmarkm\checkmark$\checkmark^\checkmark\\checkmarkc\checkmarki\checkmarkr\checkmarkc\checkmark$\checkmarka\checkmark$\checkmark^\checkmark\\checkmarkc\checkmarki\checkmarkr\checkmarkc\checkmark$\checkmarkx\checkmark$\checkmark^\checkmark\\checkmarkc\checkmarki\checkmarkr\checkmarkc\checkmark$\checkmark}\checkmark$\checkmark^\checkmark\\checkmarkc\checkmarki\checkmarkr\checkmarkc\checkmark$\checkmark \checkmark$\checkmark^\checkmark\\checkmarkc\checkmarki\checkmarkr\checkmarkc\checkmark$\checkmark=\checkmark$\checkmark^\checkmark\\checkmarkc\checkmarki\checkmarkr\checkmarkc\checkmark$\checkmark \checkmark$\checkmark^\checkmark\\checkmarkc\checkmarki\checkmarkr\checkmarkc\checkmark$\checkmark\\checkmark$\checkmark^\checkmark\\checkmarkc\checkmarki\checkmarkr\checkmarkc\checkmark$\checkmarkf\checkmark$\checkmark^\checkmark\\checkmarkc\checkmarki\checkmarkr\checkmarkc\checkmark$\checkmarkr\checkmark$\checkmark^\checkmark\\checkmarkc\checkmarki\checkmarkr\checkmarkc\checkmark$\checkmarka\checkmark$\checkmark^\checkmark\\checkmarkc\checkmarki\checkmarkr\checkmarkc\checkmark$\checkmarkc\checkmark$\checkmark^\checkmark\\checkmarkc\checkmarki\checkmarkr\checkmarkc\checkmark$\checkmark{\checkmark$\checkmark^\checkmark\\checkmarkc\checkmarki\checkmarkr\checkmarkc\checkmark$\checkmarkF\checkmark$\checkmark^\checkmark\\checkmarkc\checkmarki\checkmarkr\checkmarkc\checkmark$\checkmark_\checkmark$\checkmark^\checkmark\\checkmarkc\checkmarki\checkmarkr\checkmarkc\checkmark$\checkmark{\checkmark$\checkmark^\checkmark\\checkmarkc\checkmarki\checkmarkr\checkmarkc\checkmark$\checkmarkm\checkmark$\checkmark^\checkmark\\checkmarkc\checkmarki\checkmarkr\checkmarkc\checkmark$\checkmarka\checkmark$\checkmark^\checkmark\\checkmarkc\checkmarki\checkmarkr\checkmarkc\checkmark$\checkmarkx\checkmark$\checkmark^\checkmark\\checkmarkc\checkmarki\checkmarkr\checkmarkc\checkmark$\checkmark}\checkmark$\checkmark^\checkmark\\checkmarkc\checkmarki\checkmarkr\checkmarkc\checkmark$\checkmark \checkmark$\checkmark^\checkmark\\checkmarkc\checkmarki\checkmarkr\checkmarkc\checkmark$\checkmark\\checkmark$\checkmark^\checkmark\\checkmarkc\checkmarki\checkmarkr\checkmarkc\checkmark$\checkmarkc\checkmark$\checkmark^\checkmark\\checkmarkc\checkmarki\checkmarkr\checkmarkc\checkmark$\checkmarkd\checkmark$\checkmark^\checkmark\\checkmarkc\checkmarki\checkmarkr\checkmarkc\checkmark$\checkmarko\checkmark$\checkmark^\checkmark\\checkmarkc\checkmarki\checkmarkr\checkmarkc\checkmark$\checkmarkt\checkmark$\checkmark^\checkmark\\checkmarkc\checkmarki\checkmarkr\checkmarkc\checkmark$\checkmark \checkmark$\checkmark^\checkmark\\checkmarkc\checkmarki\checkmarkr\checkmarkc\checkmark$\checkmarkL\checkmark$\checkmark^\checkmark\\checkmarkc\checkmarki\checkmarkr\checkmarkc\checkmark$\checkmark^\checkmark$\checkmark^\checkmark\\checkmarkc\checkmarki\checkmarkr\checkmarkc\checkmark$\checkmark3\checkmark$\checkmark^\checkmark\\checkmarkc\checkmarki\checkmarkr\checkmarkc\checkmark$\checkmark}\checkmark$\checkmark^\checkmark\\checkmarkc\checkmarki\checkmarkr\checkmarkc\checkmark$\checkmark{\checkmark$\checkmark^\checkmark\\checkmarkc\checkmarki\checkmarkr\checkmarkc\checkmark$\checkmark4\checkmark$\checkmark^\checkmark\\checkmarkc\checkmarki\checkmarkr\checkmarkc\checkmark$\checkmark8\checkmark$\checkmark^\checkmark\\checkmarkc\checkmarki\checkmarkr\checkmarkc\checkmark$\checkmark \checkmark$\checkmark^\checkmark\\checkmarkc\checkmarki\checkmarkr\checkmarkc\checkmark$\checkmark\\checkmark$\checkmark^\checkmark\\checkmarkc\checkmarki\checkmarkr\checkmarkc\checkmark$\checkmarkc\checkmark$\checkmark^\checkmark\\checkmarkc\checkmarki\checkmarkr\checkmarkc\checkmark$\checkmarkd\checkmark$\checkmark^\checkmark\\checkmarkc\checkmarki\checkmarkr\checkmarkc\checkmark$\checkmarko\checkmark$\checkmark^\checkmark\\checkmarkc\checkmarki\checkmarkr\checkmarkc\checkmark$\checkmarkt\checkmark$\checkmark^\checkmark\\checkmarkc\checkmarki\checkmarkr\checkmarkc\checkmark$\checkmark \checkmark$\checkmark^\checkmark\\checkmarkc\checkmarki\checkmarkr\checkmarkc\checkmark$\checkmarkE\checkmark$\checkmark^\checkmark\\checkmarkc\checkmarki\checkmarkr\checkmarkc\checkmark$\checkmark \checkmark$\checkmark^\checkmark\\checkmarkc\checkmarki\checkmarkr\checkmarkc\checkmark$\checkmark\\checkmark$\checkmark^\checkmark\\checkmarkc\checkmarki\checkmarkr\checkmarkc\checkmark$\checkmarkc\checkmark$\checkmark^\checkmark\\checkmarkc\checkmarki\checkmarkr\checkmarkc\checkmark$\checkmarkd\checkmark$\checkmark^\checkmark\\checkmarkc\checkmarki\checkmarkr\checkmarkc\checkmark$\checkmarko\checkmark$\checkmark^\checkmark\\checkmarkc\checkmarki\checkmarkr\checkmarkc\checkmark$\checkmarkt\checkmark$\checkmark^\checkmark\\checkmarkc\checkmarki\checkmarkr\checkmarkc\checkmark$\checkmark \checkmark$\checkmark^\checkmark\\checkmarkc\checkmarki\checkmarkr\checkmarkc\checkmark$\checkmarkI\checkmark$\checkmark^\checkmark\\checkmarkc\checkmarki\checkmarkr\checkmarkc\checkmark$\checkmark_\checkmark$\checkmark^\checkmark\\checkmarkc\checkmarki\checkmarkr\checkmarkc\checkmark$\checkmark{\checkmark$\checkmark^\checkmark\\checkmarkc\checkmarki\checkmarkr\checkmarkc\checkmark$\checkmarky\checkmark$\checkmark^\checkmark\\checkmarkc\checkmarki\checkmarkr\checkmarkc\checkmark$\checkmarky\checkmark$\checkmark^\checkmark\\checkmarkc\checkmarki\checkmarkr\checkmarkc\checkmark$\checkmark}\checkmark$\checkmark^\checkmark\\checkmarkc\checkmarki\checkmarkr\checkmarkc\checkmark$\checkmark}\checkmark$\checkmark^\checkmark\\checkmarkc\checkmarki\checkmarkr\checkmarkc\checkmark$\checkmark \checkmark$\checkmark^\checkmark\\checkmarkc\checkmarki\checkmarkr\checkmarkc\checkmark$\checkmark=\checkmark$\checkmark^\checkmark\\checkmarkc\checkmarki\checkmarkr\checkmarkc\checkmark$\checkmark \checkmark$\checkmark^\checkmark\\checkmarkc\checkmarki\checkmarkr\checkmarkc\checkmark$\checkmark\\checkmark$\checkmark^\checkmark\\checkmarkc\checkmarki\checkmarkr\checkmarkc\checkmark$\checkmarkf\checkmark$\checkmark^\checkmark\\checkmarkc\checkmarki\checkmarkr\checkmarkc\checkmark$\checkmarkr\checkmark$\checkmark^\checkmark\\checkmarkc\checkmarki\checkmarkr\checkmarkc\checkmark$\checkmarka\checkmark$\checkmark^\checkmark\\checkmarkc\checkmarki\checkmarkr\checkmarkc\checkmark$\checkmarkc\checkmark$\checkmark^\checkmark\\checkmarkc\checkmarki\checkmarkr\checkmarkc\checkmark$\checkmark{\checkmark$\checkmark^\checkmark\\checkmarkc\checkmarki\checkmarkr\checkmarkc\checkmark$\checkmark1\checkmark$\checkmark^\checkmark\\checkmarkc\checkmarki\checkmarkr\checkmarkc\checkmark$\checkmark4\checkmark$\checkmark^\checkmark\\checkmarkc\checkmarki\checkmarkr\checkmarkc\checkmark$\checkmark0\checkmark$\checkmark^\checkmark\\checkmarkc\checkmarki\checkmarkr\checkmarkc\checkmark$\checkmark \checkmark$\checkmark^\checkmark\\checkmarkc\checkmarki\checkmarkr\checkmarkc\checkmark$\checkmark\\checkmark$\checkmark^\checkmark\\checkmarkc\checkmarki\checkmarkr\checkmarkc\checkmark$\checkmarkt\checkmark$\checkmark^\checkmark\\checkmarkc\checkmarki\checkmarkr\checkmarkc\checkmark$\checkmarki\checkmark$\checkmark^\checkmark\\checkmarkc\checkmarki\checkmarkr\checkmarkc\checkmark$\checkmarkm\checkmark$\checkmark^\checkmark\\checkmarkc\checkmarki\checkmarkr\checkmarkc\checkmark$\checkmarke\checkmark$\checkmark^\checkmark\\checkmarkc\checkmarki\checkmarkr\checkmarkc\checkmark$\checkmarks\checkmark$\checkmark^\checkmark\\checkmarkc\checkmarki\checkmarkr\checkmarkc\checkmark$\checkmark \checkmark$\checkmark^\checkmark\\checkmarkc\checkmarki\checkmarkr\checkmarkc\checkmark$\checkmark4\checkmark$\checkmark^\checkmark\\checkmarkc\checkmarki\checkmarkr\checkmarkc\checkmark$\checkmark0\checkmark$\checkmark^\checkmark\\checkmarkc\checkmarki\checkmarkr\checkmarkc\checkmark$\checkmark0\checkmark$\checkmark^\checkmark\\checkmarkc\checkmarki\checkmarkr\checkmarkc\checkmark$\checkmark^\checkmark$\checkmark^\checkmark\\checkmarkc\checkmarki\checkmarkr\checkmarkc\checkmark$\checkmark3\checkmark$\checkmark^\checkmark\\checkmarkc\checkmarki\checkmarkr\checkmarkc\checkmark$\checkmark}\checkmark$\checkmark^\checkmark\\checkmarkc\checkmarki\checkmarkr\checkmarkc\checkmark$\checkmark{\checkmark$\checkmark^\checkmark\\checkmarkc\checkmarki\checkmarkr\checkmarkc\checkmark$\checkmark4\checkmark$\checkmark^\checkmark\\checkmarkc\checkmarki\checkmarkr\checkmarkc\checkmark$\checkmark8\checkmark$\checkmark^\checkmark\\checkmarkc\checkmarki\checkmarkr\checkmarkc\checkmark$\checkmark \checkmark$\checkmark^\checkmark\\checkmarkc\checkmarki\checkmarkr\checkmarkc\checkmark$\checkmark\\checkmark$\checkmark^\checkmark\\checkmarkc\checkmarki\checkmarkr\checkmarkc\checkmark$\checkmarkt\checkmark$\checkmark^\checkmark\\checkmarkc\checkmarki\checkmarkr\checkmarkc\checkmark$\checkmarki\checkmark$\checkmark^\checkmark\\checkmarkc\checkmarki\checkmarkr\checkmarkc\checkmark$\checkmarkm\checkmark$\checkmark^\checkmark\\checkmarkc\checkmarki\checkmarkr\checkmarkc\checkmark$\checkmarke\checkmark$\checkmark^\checkmark\\checkmarkc\checkmarki\checkmarkr\checkmarkc\checkmark$\checkmarks\checkmark$\checkmark^\checkmark\\checkmarkc\checkmarki\checkmarkr\checkmarkc\checkmark$\checkmark \checkmark$\checkmark^\checkmark\\checkmarkc\checkmarki\checkmarkr\checkmarkc\checkmark$\checkmark6\checkmark$\checkmark^\checkmark\\checkmarkc\checkmarki\checkmarkr\checkmarkc\checkmark$\checkmark8\checkmark$\checkmark^\checkmark\\checkmarkc\checkmarki\checkmarkr\checkmarkc\checkmark$\checkmark\\checkmark$\checkmark^\checkmark\\checkmarkc\checkmarki\checkmarkr\checkmarkc\checkmark$\checkmark,\checkmark$\checkmark^\checkmark\\checkmarkc\checkmarki\checkmarkr\checkmarkc\checkmark$\checkmark9\checkmark$\checkmark^\checkmark\\checkmarkc\checkmarki\checkmarkr\checkmarkc\checkmark$\checkmark0\checkmark$\checkmark^\checkmark\\checkmarkc\checkmarki\checkmarkr\checkmarkc\checkmark$\checkmark0\checkmark$\checkmark^\checkmark\\checkmarkc\checkmarki\checkmarkr\checkmarkc\checkmark$\checkmark \checkmark$\checkmark^\checkmark\\checkmarkc\checkmarki\checkmarkr\checkmarkc\checkmark$\checkmark\\checkmark$\checkmark^\checkmark\\checkmarkc\checkmarki\checkmarkr\checkmarkc\checkmark$\checkmarkt\checkmark$\checkmark^\checkmark\\checkmarkc\checkmarki\checkmarkr\checkmarkc\checkmark$\checkmarki\checkmark$\checkmark^\checkmark\\checkmarkc\checkmarki\checkmarkr\checkmarkc\checkmark$\checkmarkm\checkmark$\checkmark^\checkmark\\checkmarkc\checkmarki\checkmarkr\checkmarkc\checkmark$\checkmarke\checkmark$\checkmark^\checkmark\\checkmarkc\checkmarki\checkmarkr\checkmarkc\checkmark$\checkmarks\checkmark$\checkmark^\checkmark\\checkmarkc\checkmarki\checkmarkr\checkmarkc\checkmark$\checkmark \checkmark$\checkmark^\checkmark\\checkmarkc\checkmarki\checkmarkr\checkmarkc\checkmark$\checkmark6\checkmark$\checkmark^\checkmark\\checkmarkc\checkmarki\checkmarkr\checkmarkc\checkmark$\checkmark5\checkmark$\checkmark^\checkmark\\checkmarkc\checkmarki\checkmarkr\checkmarkc\checkmark$\checkmark\\checkmark$\checkmark^\checkmark\\checkmarkc\checkmarki\checkmarkr\checkmarkc\checkmark$\checkmark,\checkmark$\checkmark^\checkmark\\checkmarkc\checkmarki\checkmarkr\checkmarkc\checkmark$\checkmark0\checkmark$\checkmark^\checkmark\\checkmarkc\checkmarki\checkmarkr\checkmarkc\checkmark$\checkmark0\checkmark$\checkmark^\checkmark\\checkmarkc\checkmarki\checkmarkr\checkmarkc\checkmark$\checkmark0\checkmark$\checkmark^\checkmark\\checkmarkc\checkmarki\checkmarkr\checkmarkc\checkmark$\checkmark}\checkmark$\checkmark^\checkmark\\checkmarkc\checkmarki\checkmarkr\checkmarkc\checkmark$\checkmark \checkmark$\checkmark^\checkmark\\checkmarkc\checkmarki\checkmarkr\checkmarkc\checkmark$\checkmark=\checkmark$\checkmark^\checkmark\\checkmarkc\checkmarki\checkmarkr\checkmarkc\checkmark$\checkmark \checkmark$\checkmark^\checkmark\\checkmarkc\checkmarki\checkmarkr\checkmarkc\checkmark$\checkmark\\checkmark$\checkmark^\checkmark\\checkmarkc\checkmarki\checkmarkr\checkmarkc\checkmark$\checkmarkf\checkmark$\checkmark^\checkmark\\checkmarkc\checkmarki\checkmarkr\checkmarkc\checkmark$\checkmarkr\checkmark$\checkmark^\checkmark\\checkmarkc\checkmarki\checkmarkr\checkmarkc\checkmark$\checkmarka\checkmark$\checkmark^\checkmark\\checkmarkc\checkmarki\checkmarkr\checkmarkc\checkmark$\checkmarkc\checkmark$\checkmark^\checkmark\\checkmarkc\checkmarki\checkmarkr\checkmarkc\checkmark$\checkmark{\checkmark$\checkmark^\checkmark\\checkmarkc\checkmarki\checkmarkr\checkmarkc\checkmark$\checkmark1\checkmark$\checkmark^\checkmark\\checkmarkc\checkmarki\checkmarkr\checkmarkc\checkmark$\checkmark4\checkmark$\checkmark^\checkmark\\checkmarkc\checkmarki\checkmarkr\checkmarkc\checkmark$\checkmark0\checkmark$\checkmark^\checkmark\\checkmarkc\checkmarki\checkmarkr\checkmarkc\checkmark$\checkmark \checkmark$\checkmark^\checkmark\\checkmarkc\checkmarki\checkmarkr\checkmarkc\checkmark$\checkmark\\checkmark$\checkmark^\checkmark\\checkmarkc\checkmarki\checkmarkr\checkmarkc\checkmark$\checkmarkt\checkmark$\checkmark^\checkmark\\checkmarkc\checkmarki\checkmarkr\checkmarkc\checkmark$\checkmarki\checkmark$\checkmark^\checkmark\\checkmarkc\checkmarki\checkmarkr\checkmarkc\checkmark$\checkmarkm\checkmark$\checkmark^\checkmark\\checkmarkc\checkmarki\checkmarkr\checkmarkc\checkmark$\checkmarke\checkmark$\checkmark^\checkmark\\checkmarkc\checkmarki\checkmarkr\checkmarkc\checkmark$\checkmarks\checkmark$\checkmark^\checkmark\\checkmarkc\checkmarki\checkmarkr\checkmarkc\checkmark$\checkmark \checkmark$\checkmark^\checkmark\\checkmarkc\checkmarki\checkmarkr\checkmarkc\checkmark$\checkmark6\checkmark$\checkmark^\checkmark\\checkmarkc\checkmarki\checkmarkr\checkmarkc\checkmark$\checkmark4\checkmark$\checkmark^\checkmark\\checkmarkc\checkmarki\checkmarkr\checkmarkc\checkmark$\checkmark \checkmark$\checkmark^\checkmark\\checkmarkc\checkmarki\checkmarkr\checkmarkc\checkmark$\checkmark\\checkmark$\checkmark^\checkmark\\checkmarkc\checkmarki\checkmarkr\checkmarkc\checkmark$\checkmarkt\checkmark$\checkmark^\checkmark\\checkmarkc\checkmarki\checkmarkr\checkmarkc\checkmark$\checkmarki\checkmark$\checkmark^\checkmark\\checkmarkc\checkmarki\checkmarkr\checkmarkc\checkmark$\checkmarkm\checkmark$\checkmark^\checkmark\\checkmarkc\checkmarki\checkmarkr\checkmarkc\checkmark$\checkmarke\checkmark$\checkmark^\checkmark\\checkmarkc\checkmarki\checkmarkr\checkmarkc\checkmark$\checkmarks\checkmark$\checkmark^\checkmark\\checkmarkc\checkmarki\checkmarkr\checkmarkc\checkmark$\checkmark \checkmark$\checkmark^\checkmark\\checkmarkc\checkmarki\checkmarkr\checkmarkc\checkmark$\checkmark1\checkmark$\checkmark^\checkmark\\checkmarkc\checkmarki\checkmarkr\checkmarkc\checkmark$\checkmark0\checkmark$\checkmark^\checkmark\\checkmarkc\checkmarki\checkmarkr\checkmarkc\checkmark$\checkmark^\checkmark$\checkmark^\checkmark\\checkmarkc\checkmarki\checkmarkr\checkmarkc\checkmark$\checkmark6\checkmark$\checkmark^\checkmark\\checkmarkc\checkmarki\checkmarkr\checkmarkc\checkmark$\checkmark}\checkmark$\checkmark^\checkmark\\checkmarkc\checkmarki\checkmarkr\checkmarkc\checkmark$\checkmark{\checkmark$\checkmark^\checkmark\\checkmarkc\checkmarki\checkmarkr\checkmarkc\checkmark$\checkmark2\checkmark$\checkmark^\checkmark\\checkmarkc\checkmarki\checkmarkr\checkmarkc\checkmark$\checkmark1\checkmark$\checkmark^\checkmark\\checkmarkc\checkmarki\checkmarkr\checkmarkc\checkmark$\checkmark5\checkmark$\checkmark^\checkmark\\checkmarkc\checkmarki\checkmarkr\checkmarkc\checkmark$\checkmark\\checkmark$\checkmark^\checkmark\\checkmarkc\checkmarki\checkmarkr\checkmarkc\checkmark$\checkmark,\checkmark$\checkmark^\checkmark\\checkmarkc\checkmarki\checkmarkr\checkmarkc\checkmark$\checkmark4\checkmark$\checkmark^\checkmark\\checkmarkc\checkmarki\checkmarkr\checkmarkc\checkmark$\checkmark2\checkmark$\checkmark^\checkmark\\checkmarkc\checkmarki\checkmarkr\checkmarkc\checkmark$\checkmark4\checkmark$\checkmark^\checkmark\\checkmarkc\checkmarki\checkmarkr\checkmarkc\checkmark$\checkmark \checkmark$\checkmark^\checkmark\\checkmarkc\checkmarki\checkmarkr\checkmarkc\checkmark$\checkmark\\checkmark$\checkmark^\checkmark\\checkmarkc\checkmarki\checkmarkr\checkmarkc\checkmark$\checkmarkt\checkmark$\checkmark^\checkmark\\checkmarkc\checkmarki\checkmarkr\checkmarkc\checkmark$\checkmarki\checkmark$\checkmark^\checkmark\\checkmarkc\checkmarki\checkmarkr\checkmarkc\checkmark$\checkmarkm\checkmark$\checkmark^\checkmark\\checkmarkc\checkmarki\checkmarkr\checkmarkc\checkmark$\checkmarke\checkmark$\checkmark^\checkmark\\checkmarkc\checkmarki\checkmarkr\checkmarkc\checkmark$\checkmarks\checkmark$\checkmark^\checkmark\\checkmarkc\checkmarki\checkmarkr\checkmarkc\checkmark$\checkmark \checkmark$\checkmark^\checkmark\\checkmarkc\checkmarki\checkmarkr\checkmarkc\checkmark$\checkmark1\checkmark$\checkmark^\checkmark\\checkmarkc\checkmarki\checkmarkr\checkmarkc\checkmark$\checkmark0\checkmark$\checkmark^\checkmark\\checkmarkc\checkmarki\checkmarkr\checkmarkc\checkmark$\checkmark^\checkmark$\checkmark^\checkmark\\checkmarkc\checkmarki\checkmarkr\checkmarkc\checkmark$\checkmark6\checkmark$\checkmark^\checkmark\\checkmarkc\checkmarki\checkmarkr\checkmarkc\checkmark$\checkmark}\checkmark$\checkmark^\checkmark\\checkmarkc\checkmarki\checkmarkr\checkmarkc\checkmark$\checkmark \checkmark$\checkmark^\checkmark\\checkmarkc\checkmarki\checkmarkr\checkmarkc\checkmark$\checkmark=\checkmark$\checkmark^\checkmark\\checkmarkc\checkmarki\checkmarkr\checkmarkc\checkmark$\checkmark \checkmark$\checkmark^\checkmark\\checkmarkc\checkmarki\checkmarkr\checkmarkc\checkmark$\checkmark\\checkmark$\checkmark^\checkmark\\checkmarkc\checkmarki\checkmarkr\checkmarkc\checkmark$\checkmarkb\checkmark$\checkmark^\checkmark\\checkmarkc\checkmarki\checkmarkr\checkmarkc\checkmark$\checkmarko\checkmark$\checkmark^\checkmark\\checkmarkc\checkmarki\checkmarkr\checkmarkc\checkmark$\checkmarkx\checkmark$\checkmark^\checkmark\\checkmarkc\checkmarki\checkmarkr\checkmarkc\checkmark$\checkmarke\checkmark$\checkmark^\checkmark\\checkmarkc\checkmarki\checkmarkr\checkmarkc\checkmark$\checkmarkd\checkmark$\checkmark^\checkmark\\checkmarkc\checkmarki\checkmarkr\checkmarkc\checkmark$\checkmark{\checkmark$\checkmark^\checkmark\\checkmarkc\checkmarki\checkmarkr\checkmarkc\checkmark$\checkmark0\checkmark$\checkmark^\checkmark\\checkmarkc\checkmarki\checkmarkr\checkmarkc\checkmark$\checkmark{\checkmark$\checkmark^\checkmark\\checkmarkc\checkmarki\checkmarkr\checkmarkc\checkmark$\checkmark,\checkmark$\checkmark^\checkmark\\checkmarkc\checkmarki\checkmarkr\checkmarkc\checkmark$\checkmark}\checkmark$\checkmark^\checkmark\\checkmarkc\checkmarki\checkmarkr\checkmarkc\checkmark$\checkmark0\checkmark$\checkmark^\checkmark\\checkmarkc\checkmarki\checkmarkr\checkmarkc\checkmark$\checkmark4\checkmark$\checkmark^\checkmark\\checkmarkc\checkmarki\checkmarkr\checkmarkc\checkmark$\checkmark2\checkmark$\checkmark^\checkmark\\checkmarkc\checkmarki\checkmarkr\checkmarkc\checkmark$\checkmark \checkmark$\checkmark^\checkmark\\checkmarkc\checkmarki\checkmarkr\checkmarkc\checkmark$\checkmark\\checkmark$\checkmark^\checkmark\\checkmarkc\checkmarki\checkmarkr\checkmarkc\checkmark$\checkmarkt\checkmark$\checkmark^\checkmark\\checkmarkc\checkmarki\checkmarkr\checkmarkc\checkmark$\checkmarke\checkmark$\checkmark^\checkmark\\checkmarkc\checkmarki\checkmarkr\checkmarkc\checkmark$\checkmarkx\checkmark$\checkmark^\checkmark\\checkmarkc\checkmarki\checkmarkr\checkmarkc\checkmark$\checkmarkt\checkmark$\checkmark^\checkmark\\checkmarkc\checkmarki\checkmarkr\checkmarkc\checkmark$\checkmark{\checkmark$\checkmark^\checkmark\\checkmarkc\checkmarki\checkmarkr\checkmarkc\checkmark$\checkmark \checkmark$\checkmark^\checkmark\\checkmarkc\checkmarki\checkmarkr\checkmarkc\checkmark$\checkmarkm\checkmark$\checkmark^\checkmark\\checkmarkc\checkmarki\checkmarkr\checkmarkc\checkmark$\checkmarkm\checkmark$\checkmark^\checkmark\\checkmarkc\checkmarki\checkmarkr\checkmarkc\checkmark$\checkmark}\checkmark$\checkmark^\checkmark\\checkmarkc\checkmarki\checkmarkr\checkmarkc\checkmark$\checkmark}\checkmark$\checkmark^\checkmark\\checkmarkc\checkmarki\checkmarkr\checkmarkc\checkmark$\checkmark
\checkmark$\checkmark^\checkmark\\checkmarkc\checkmarki\checkmarkr\checkmarkc\checkmark$\checkmark \checkmark$\checkmark^\checkmark\\checkmarkc\checkmarki\checkmarkr\checkmarkc\checkmark$\checkmark \checkmark$\checkmark^\checkmark\\checkmarkc\checkmarki\checkmarkr\checkmarkc\checkmark$\checkmark \checkmark$\checkmark^\checkmark\\checkmarkc\checkmarki\checkmarkr\checkmarkc\checkmark$\checkmark \checkmark$\checkmark^\checkmark\\checkmarkc\checkmarki\checkmarkr\checkmarkc\checkmark$\checkmark\\checkmark$\checkmark^\checkmark\\checkmarkc\checkmarki\checkmarkr\checkmarkc\checkmark$\checkmarkl\checkmark$\checkmark^\checkmark\\checkmarkc\checkmarki\checkmarkr\checkmarkc\checkmark$\checkmarka\checkmark$\checkmark^\checkmark\\checkmarkc\checkmarki\checkmarkr\checkmarkc\checkmark$\checkmarkb\checkmark$\checkmark^\checkmark\\checkmarkc\checkmarki\checkmarkr\checkmarkc\checkmark$\checkmarke\checkmark$\checkmark^\checkmark\\checkmarkc\checkmarki\checkmarkr\checkmarkc\checkmark$\checkmarkl\checkmark$\checkmark^\checkmark\\checkmarkc\checkmarki\checkmarkr\checkmarkc\checkmark$\checkmark{\checkmark$\checkmark^\checkmark\\checkmarkc\checkmarki\checkmarkr\checkmarkc\checkmark$\checkmarke\checkmark$\checkmark^\checkmark\\checkmarkc\checkmarki\checkmarkr\checkmarkc\checkmark$\checkmarkq\checkmark$\checkmark^\checkmark\\checkmarkc\checkmarki\checkmarkr\checkmarkc\checkmark$\checkmark:\checkmark$\checkmark^\checkmark\\checkmarkc\checkmarki\checkmarkr\checkmarkc\checkmark$\checkmarkf\checkmark$\checkmark^\checkmark\\checkmarkc\checkmarki\checkmarkr\checkmarkc\checkmark$\checkmarkl\checkmark$\checkmark^\checkmark\\checkmarkc\checkmarki\checkmarkr\checkmarkc\checkmark$\checkmarke\checkmark$\checkmark^\checkmark\\checkmarkc\checkmarki\checkmarkr\checkmarkc\checkmark$\checkmarkc\checkmark$\checkmark^\checkmark\\checkmarkc\checkmarki\checkmarkr\checkmarkc\checkmark$\checkmarkh\checkmark$\checkmark^\checkmark\\checkmarkc\checkmarki\checkmarkr\checkmarkc\checkmark$\checkmarke\checkmark$\checkmark^\checkmark\\checkmarkc\checkmarki\checkmarkr\checkmarkc\checkmark$\checkmark}\checkmark$\checkmark^\checkmark\\checkmarkc\checkmarki\checkmarkr\checkmarkc\checkmark$\checkmark
\checkmark$\checkmark^\checkmark\\checkmarkc\checkmarki\checkmarkr\checkmarkc\checkmark$\checkmark\\checkmark$\checkmark^\checkmark\\checkmarkc\checkmarki\checkmarkr\checkmarkc\checkmark$\checkmarke\checkmark$\checkmark^\checkmark\\checkmarkc\checkmarki\checkmarkr\checkmarkc\checkmark$\checkmarkn\checkmark$\checkmark^\checkmark\\checkmarkc\checkmarki\checkmarkr\checkmarkc\checkmark$\checkmarkd\checkmark$\checkmark^\checkmark\\checkmarkc\checkmarki\checkmarkr\checkmarkc\checkmark$\checkmark{\checkmark$\checkmark^\checkmark\\checkmarkc\checkmarki\checkmarkr\checkmarkc\checkmark$\checkmarke\checkmark$\checkmark^\checkmark\\checkmarkc\checkmarki\checkmarkr\checkmarkc\checkmark$\checkmarkq\checkmark$\checkmark^\checkmark\\checkmarkc\checkmarki\checkmarkr\checkmarkc\checkmark$\checkmarku\checkmark$\checkmark^\checkmark\\checkmarkc\checkmarki\checkmarkr\checkmarkc\checkmark$\checkmarka\checkmark$\checkmark^\checkmark\\checkmarkc\checkmarki\checkmarkr\checkmarkc\checkmark$\checkmarkt\checkmark$\checkmark^\checkmark\\checkmarkc\checkmarki\checkmarkr\checkmarkc\checkmark$\checkmarki\checkmark$\checkmark^\checkmark\\checkmarkc\checkmarki\checkmarkr\checkmarkc\checkmark$\checkmarko\checkmark$\checkmark^\checkmark\\checkmarkc\checkmarki\checkmarkr\checkmarkc\checkmark$\checkmarkn\checkmark$\checkmark^\checkmark\\checkmarkc\checkmarki\checkmarkr\checkmarkc\checkmark$\checkmark}\checkmark$\checkmark^\checkmark\\checkmarkc\checkmarki\checkmarkr\checkmarkc\checkmark$\checkmark
\checkmark$\checkmark^\checkmark\\checkmarkc\checkmarki\checkmarkr\checkmarkc\checkmark$\checkmark
\checkmark$\checkmark^\checkmark\\checkmarkc\checkmarki\checkmarkr\checkmarkc\checkmark$\checkmark\\checkmark$\checkmark^\checkmark\\checkmarkc\checkmarki\checkmarkr\checkmarkc\checkmark$\checkmarkt\checkmark$\checkmark^\checkmark\\checkmarkc\checkmarki\checkmarkr\checkmarkc\checkmark$\checkmarke\checkmark$\checkmark^\checkmark\\checkmarkc\checkmarki\checkmarkr\checkmarkc\checkmark$\checkmarkx\checkmark$\checkmark^\checkmark\\checkmarkc\checkmarki\checkmarkr\checkmarkc\checkmark$\checkmarkt\checkmark$\checkmark^\checkmark\\checkmarkc\checkmarki\checkmarkr\checkmarkc\checkmark$\checkmarkb\checkmark$\checkmark^\checkmark\\checkmarkc\checkmarki\checkmarkr\checkmarkc\checkmark$\checkmarkf\checkmark$\checkmark^\checkmark\\checkmarkc\checkmarki\checkmarkr\checkmarkc\checkmark$\checkmark{\checkmark$\checkmark^\checkmark\\checkmarkc\checkmarki\checkmarkr\checkmarkc\checkmark$\checkmarkV\checkmark$\checkmark^\checkmark\\checkmarkc\checkmarki\checkmarkr\checkmarkc\checkmark$\checkmarké\checkmark$\checkmark^\checkmark\\checkmarkc\checkmarki\checkmarkr\checkmarkc\checkmark$\checkmarkr\checkmark$\checkmark^\checkmark\\checkmarkc\checkmarki\checkmarkr\checkmarkc\checkmark$\checkmarki\checkmark$\checkmark^\checkmark\\checkmarkc\checkmarki\checkmarkr\checkmarkc\checkmark$\checkmarkf\checkmark$\checkmark^\checkmark\\checkmarkc\checkmarki\checkmarkr\checkmarkc\checkmark$\checkmarki\checkmark$\checkmark^\checkmark\\checkmarkc\checkmarki\checkmarkr\checkmarkc\checkmark$\checkmarkc\checkmark$\checkmark^\checkmark\\checkmarkc\checkmarki\checkmarkr\checkmarkc\checkmark$\checkmarka\checkmark$\checkmark^\checkmark\\checkmarkc\checkmarki\checkmarkr\checkmarkc\checkmark$\checkmarkt\checkmark$\checkmark^\checkmark\\checkmarkc\checkmarki\checkmarkr\checkmarkc\checkmark$\checkmarki\checkmark$\checkmark^\checkmark\\checkmarkc\checkmarki\checkmarkr\checkmarkc\checkmark$\checkmarko\checkmark$\checkmark^\checkmark\\checkmarkc\checkmarki\checkmarkr\checkmarkc\checkmark$\checkmarkn\checkmark$\checkmark^\checkmark\\checkmarkc\checkmarki\checkmarkr\checkmarkc\checkmark$\checkmark \checkmark$\checkmark^\checkmark\\checkmarkc\checkmarki\checkmarkr\checkmarkc\checkmark$\checkmark:\checkmark$\checkmark^\checkmark\\checkmarkc\checkmarki\checkmarkr\checkmarkc\checkmark$\checkmark}\checkmark$\checkmark^\checkmark\\checkmarkc\checkmarki\checkmarkr\checkmarkc\checkmark$\checkmark \checkmark$\checkmark^\checkmark\\checkmarkc\checkmarki\checkmarkr\checkmarkc\checkmark$\checkmark$\checkmark$\checkmark^\checkmark\\checkmarkc\checkmarki\checkmarkr\checkmarkc\checkmark$\checkmarkf\checkmark$\checkmark^\checkmark\\checkmarkc\checkmarki\checkmarkr\checkmarkc\checkmark$\checkmark_\checkmark$\checkmark^\checkmark\\checkmarkc\checkmarki\checkmarkr\checkmarkc\checkmark$\checkmark{\checkmark$\checkmark^\checkmark\\checkmarkc\checkmarki\checkmarkr\checkmarkc\checkmark$\checkmarkm\checkmark$\checkmark^\checkmark\\checkmarkc\checkmarki\checkmarkr\checkmarkc\checkmark$\checkmarka\checkmark$\checkmark^\checkmark\\checkmarkc\checkmarki\checkmarkr\checkmarkc\checkmark$\checkmarkx\checkmark$\checkmark^\checkmark\\checkmarkc\checkmarki\checkmarkr\checkmarkc\checkmark$\checkmark}\checkmark$\checkmark^\checkmark\\checkmarkc\checkmarki\checkmarkr\checkmarkc\checkmark$\checkmark \checkmark$\checkmark^\checkmark\\checkmarkc\checkmarki\checkmarkr\checkmarkc\checkmark$\checkmark=\checkmark$\checkmark^\checkmark\\checkmarkc\checkmarki\checkmarkr\checkmarkc\checkmark$\checkmark \checkmark$\checkmark^\checkmark\\checkmarkc\checkmarki\checkmarkr\checkmarkc\checkmark$\checkmark0\checkmark$\checkmark^\checkmark\\checkmarkc\checkmarki\checkmarkr\checkmarkc\checkmark$\checkmark{\checkmark$\checkmark^\checkmark\\checkmarkc\checkmarki\checkmarkr\checkmarkc\checkmark$\checkmark,\checkmark$\checkmark^\checkmark\\checkmarkc\checkmarki\checkmarkr\checkmarkc\checkmark$\checkmark}\checkmark$\checkmark^\checkmark\\checkmarkc\checkmarki\checkmarkr\checkmarkc\checkmark$\checkmark0\checkmark$\checkmark^\checkmark\\checkmarkc\checkmarki\checkmarkr\checkmarkc\checkmark$\checkmark4\checkmark$\checkmark^\checkmark\\checkmarkc\checkmarki\checkmarkr\checkmarkc\checkmark$\checkmark2\checkmark$\checkmark^\checkmark\\checkmarkc\checkmarki\checkmarkr\checkmarkc\checkmark$\checkmark$\checkmark$\checkmark^\checkmark\\checkmarkc\checkmarki\checkmarkr\checkmarkc\checkmark$\checkmark \checkmark$\checkmark^\checkmark\\checkmarkc\checkmarki\checkmarkr\checkmarkc\checkmark$\checkmarkm\checkmark$\checkmark^\checkmark\\checkmarkc\checkmarki\checkmarkr\checkmarkc\checkmark$\checkmarkm\checkmark$\checkmark^\checkmark\\checkmarkc\checkmarki\checkmarkr\checkmarkc\checkmark$\checkmark \checkmark$\checkmark^\checkmark\\checkmarkc\checkmarki\checkmarkr\checkmarkc\checkmark$\checkmark$\checkmark$\checkmark^\checkmark\\checkmarkc\checkmarki\checkmarkr\checkmarkc\checkmark$\checkmark\\checkmark$\checkmark^\checkmark\\checkmarkc\checkmarki\checkmarkr\checkmarkc\checkmark$\checkmarkl\checkmark$\checkmark^\checkmark\\checkmarkc\checkmarki\checkmarkr\checkmarkc\checkmark$\checkmarkl\checkmark$\checkmark^\checkmark\\checkmarkc\checkmarki\checkmarkr\checkmarkc\checkmark$\checkmark \checkmark$\checkmark^\checkmark\\checkmarkc\checkmarki\checkmarkr\checkmarkc\checkmark$\checkmark0\checkmark$\checkmark^\checkmark\\checkmarkc\checkmarki\checkmarkr\checkmarkc\checkmark$\checkmark{\checkmark$\checkmark^\checkmark\\checkmarkc\checkmarki\checkmarkr\checkmarkc\checkmark$\checkmark,\checkmark$\checkmark^\checkmark\\checkmarkc\checkmarki\checkmarkr\checkmarkc\checkmark$\checkmark}\checkmark$\checkmark^\checkmark\\checkmarkc\checkmarki\checkmarkr\checkmarkc\checkmark$\checkmark1\checkmark$\checkmark^\checkmark\\checkmarkc\checkmarki\checkmarkr\checkmarkc\checkmark$\checkmark$\checkmark$\checkmark^\checkmark\\checkmarkc\checkmarki\checkmarkr\checkmarkc\checkmark$\checkmark \checkmark$\checkmark^\checkmark\\checkmarkc\checkmarki\checkmarkr\checkmarkc\checkmark$\checkmarkm\checkmark$\checkmark^\checkmark\\checkmarkc\checkmarki\checkmarkr\checkmarkc\checkmark$\checkmarkm\checkmark$\checkmark^\checkmark\\checkmarkc\checkmarki\checkmarkr\checkmarkc\checkmark$\checkmark \checkmark$\checkmark^\checkmark\\checkmarkc\checkmarki\checkmarkr\checkmarkc\checkmark$\checkmark(\checkmark$\checkmark^\checkmark\\checkmarkc\checkmarki\checkmarkr\checkmarkc\checkmark$\checkmarkt\checkmark$\checkmark^\checkmark\\checkmarkc\checkmarki\checkmarkr\checkmarkc\checkmark$\checkmarko\checkmark$\checkmark^\checkmark\\checkmarkc\checkmarki\checkmarkr\checkmarkc\checkmark$\checkmarkl\checkmark$\checkmark^\checkmark\\checkmarkc\checkmarki\checkmarkr\checkmarkc\checkmark$\checkmarké\checkmark$\checkmark^\checkmark\\checkmarkc\checkmarki\checkmarkr\checkmarkc\checkmark$\checkmarkr\checkmark$\checkmark^\checkmark\\checkmarkc\checkmarki\checkmarkr\checkmarkc\checkmark$\checkmarka\checkmark$\checkmark^\checkmark\\checkmarkc\checkmarki\checkmarkr\checkmarkc\checkmark$\checkmarkn\checkmark$\checkmark^\checkmark\\checkmarkc\checkmarki\checkmarkr\checkmarkc\checkmark$\checkmarkc\checkmark$\checkmark^\checkmark\\checkmarkc\checkmarki\checkmarkr\checkmarkc\checkmark$\checkmarke\checkmark$\checkmark^\checkmark\\checkmarkc\checkmarki\checkmarkr\checkmarkc\checkmark$\checkmark \checkmark$\checkmark^\checkmark\\checkmarkc\checkmarki\checkmarkr\checkmarkc\checkmark$\checkmarke\checkmark$\checkmark^\checkmark\\checkmarkc\checkmarki\checkmarkr\checkmarkc\checkmark$\checkmarkx\checkmark$\checkmark^\checkmark\\checkmarkc\checkmarki\checkmarkr\checkmarkc\checkmark$\checkmarki\checkmark$\checkmark^\checkmark\\checkmarkc\checkmarki\checkmarkr\checkmarkc\checkmark$\checkmarkg\checkmark$\checkmark^\checkmark\\checkmarkc\checkmarki\checkmarkr\checkmarkc\checkmark$\checkmarké\checkmark$\checkmark^\checkmark\\checkmarkc\checkmarki\checkmarkr\checkmarkc\checkmark$\checkmarke\checkmark$\checkmark^\checkmark\\checkmarkc\checkmarki\checkmarkr\checkmarkc\checkmark$\checkmark)\checkmark$\checkmark^\checkmark\\checkmarkc\checkmarki\checkmarkr\checkmarkc\checkmark$\checkmark.\checkmark$\checkmark^\checkmark\\checkmarkc\checkmarki\checkmarkr\checkmarkc\checkmark$\checkmark \checkmark$\checkmark^\checkmark\\checkmarkc\checkmarki\checkmarkr\checkmarkc\checkmark$\checkmarkL\checkmark$\checkmark^\checkmark\\checkmarkc\checkmarki\checkmarkr\checkmarkc\checkmark$\checkmarka\checkmark$\checkmark^\checkmark\\checkmarkc\checkmarki\checkmarkr\checkmarkc\checkmark$\checkmark \checkmark$\checkmark^\checkmark\\checkmarkc\checkmarki\checkmarkr\checkmarkc\checkmark$\checkmarkc\checkmark$\checkmark^\checkmark\\checkmarkc\checkmarki\checkmarkr\checkmarkc\checkmark$\checkmarko\checkmark$\checkmark^\checkmark\\checkmarkc\checkmarki\checkmarkr\checkmarkc\checkmark$\checkmarkn\checkmark$\checkmark^\checkmark\\checkmarkc\checkmarki\checkmarkr\checkmarkc\checkmark$\checkmarkd\checkmark$\checkmark^\checkmark\\checkmarkc\checkmarki\checkmarkr\checkmarkc\checkmark$\checkmarki\checkmark$\checkmark^\checkmark\\checkmarkc\checkmarki\checkmarkr\checkmarkc\checkmark$\checkmarkt\checkmark$\checkmark^\checkmark\\checkmarkc\checkmarki\checkmarkr\checkmarkc\checkmark$\checkmarki\checkmark$\checkmark^\checkmark\\checkmarkc\checkmarki\checkmarkr\checkmarkc\checkmark$\checkmarko\checkmark$\checkmark^\checkmark\\checkmarkc\checkmarki\checkmarkr\checkmarkc\checkmark$\checkmarkn\checkmark$\checkmark^\checkmark\\checkmarkc\checkmarki\checkmarkr\checkmarkc\checkmark$\checkmark \checkmark$\checkmark^\checkmark\\checkmarkc\checkmarki\checkmarkr\checkmarkc\checkmark$\checkmarkd\checkmark$\checkmark^\checkmark\\checkmarkc\checkmarki\checkmarkr\checkmarkc\checkmark$\checkmarke\checkmark$\checkmark^\checkmark\\checkmarkc\checkmarki\checkmarkr\checkmarkc\checkmark$\checkmark \checkmark$\checkmark^\checkmark\\checkmarkc\checkmarki\checkmarkr\checkmarkc\checkmark$\checkmarkd\checkmark$\checkmark^\checkmark\\checkmarkc\checkmarki\checkmarkr\checkmarkc\checkmark$\checkmarké\checkmark$\checkmark^\checkmark\\checkmarkc\checkmarki\checkmarkr\checkmarkc\checkmark$\checkmarkf\checkmark$\checkmark^\checkmark\\checkmarkc\checkmarki\checkmarkr\checkmarkc\checkmark$\checkmarkl\checkmark$\checkmark^\checkmark\\checkmarkc\checkmarki\checkmarkr\checkmarkc\checkmark$\checkmarke\checkmark$\checkmark^\checkmark\\checkmarkc\checkmarki\checkmarkr\checkmarkc\checkmark$\checkmarkx\checkmark$\checkmark^\checkmark\\checkmarkc\checkmarki\checkmarkr\checkmarkc\checkmark$\checkmarki\checkmark$\checkmark^\checkmark\\checkmarkc\checkmarki\checkmarkr\checkmarkc\checkmark$\checkmarko\checkmark$\checkmark^\checkmark\\checkmarkc\checkmarki\checkmarkr\checkmarkc\checkmark$\checkmarkn\checkmark$\checkmark^\checkmark\\checkmarkc\checkmarki\checkmarkr\checkmarkc\checkmark$\checkmark \checkmark$\checkmark^\checkmark\\checkmarkc\checkmarki\checkmarkr\checkmarkc\checkmark$\checkmarke\checkmark$\checkmark^\checkmark\\checkmarkc\checkmarki\checkmarkr\checkmarkc\checkmark$\checkmarks\checkmark$\checkmark^\checkmark\\checkmarkc\checkmarki\checkmarkr\checkmarkc\checkmark$\checkmarkt\checkmark$\checkmark^\checkmark\\checkmarkc\checkmarki\checkmarkr\checkmarkc\checkmark$\checkmark \checkmark$\checkmark^\checkmark\\checkmarkc\checkmarki\checkmarkr\checkmarkc\checkmark$\checkmarks\checkmark$\checkmark^\checkmark\\checkmarkc\checkmarki\checkmarkr\checkmarkc\checkmark$\checkmarka\checkmark$\checkmark^\checkmark\\checkmarkc\checkmarki\checkmarkr\checkmarkc\checkmark$\checkmarkt\checkmark$\checkmark^\checkmark\\checkmarkc\checkmarki\checkmarkr\checkmarkc\checkmark$\checkmarki\checkmark$\checkmark^\checkmark\\checkmarkc\checkmarki\checkmarkr\checkmarkc\checkmark$\checkmarks\checkmark$\checkmark^\checkmark\\checkmarkc\checkmarki\checkmarkr\checkmarkc\checkmark$\checkmarkf\checkmark$\checkmark^\checkmark\\checkmarkc\checkmarki\checkmarkr\checkmarkc\checkmark$\checkmarka\checkmark$\checkmark^\checkmark\\checkmarkc\checkmarki\checkmarkr\checkmarkc\checkmark$\checkmarki\checkmark$\checkmark^\checkmark\\checkmarkc\checkmarki\checkmarkr\checkmarkc\checkmark$\checkmarkt\checkmark$\checkmark^\checkmark\\checkmarkc\checkmarki\checkmarkr\checkmarkc\checkmark$\checkmarke\checkmark$\checkmark^\checkmark\\checkmarkc\checkmarki\checkmarkr\checkmarkc\checkmark$\checkmark.\checkmark$\checkmark^\checkmark\\checkmarkc\checkmarki\checkmarkr\checkmarkc\checkmark$\checkmark
\checkmark$\checkmark^\checkmark\\checkmarkc\checkmarki\checkmarkr\checkmarkc\checkmark$\checkmark
\checkmark$\checkmark^\checkmark\\checkmarkc\checkmarki\checkmarkr\checkmarkc\checkmark$\checkmark\\checkmark$\checkmark^\checkmark\\checkmarkc\checkmarki\checkmarkr\checkmarkc\checkmark$\checkmarks\checkmark$\checkmark^\checkmark\\checkmarkc\checkmarki\checkmarkr\checkmarkc\checkmark$\checkmarku\checkmark$\checkmark^\checkmark\\checkmarkc\checkmarki\checkmarkr\checkmarkc\checkmark$\checkmarkb\checkmark$\checkmark^\checkmark\\checkmarkc\checkmarki\checkmarkr\checkmarkc\checkmark$\checkmarks\checkmark$\checkmark^\checkmark\\checkmarkc\checkmarki\checkmarkr\checkmarkc\checkmark$\checkmarke\checkmark$\checkmark^\checkmark\\checkmarkc\checkmarki\checkmarkr\checkmarkc\checkmark$\checkmarkc\checkmark$\checkmark^\checkmark\\checkmarkc\checkmarki\checkmarkr\checkmarkc\checkmark$\checkmarkt\checkmark$\checkmark^\checkmark\\checkmarkc\checkmarki\checkmarkr\checkmarkc\checkmark$\checkmarki\checkmark$\checkmark^\checkmark\\checkmarkc\checkmarki\checkmarkr\checkmarkc\checkmark$\checkmarko\checkmark$\checkmark^\checkmark\\checkmarkc\checkmarki\checkmarkr\checkmarkc\checkmark$\checkmarkn\checkmark$\checkmark^\checkmark\\checkmarkc\checkmarki\checkmarkr\checkmarkc\checkmark$\checkmark{\checkmark$\checkmark^\checkmark\\checkmarkc\checkmarki\checkmarkr\checkmarkc\checkmark$\checkmarkV\checkmark$\checkmark^\checkmark\\checkmarkc\checkmarki\checkmarkr\checkmarkc\checkmark$\checkmarké\checkmark$\checkmark^\checkmark\\checkmarkc\checkmarki\checkmarkr\checkmarkc\checkmark$\checkmarkr\checkmark$\checkmark^\checkmark\\checkmarkc\checkmarki\checkmarkr\checkmarkc\checkmark$\checkmarki\checkmark$\checkmark^\checkmark\\checkmarkc\checkmarki\checkmarkr\checkmarkc\checkmark$\checkmarkf\checkmark$\checkmark^\checkmark\\checkmarkc\checkmarki\checkmarkr\checkmarkc\checkmark$\checkmarki\checkmark$\checkmark^\checkmark\\checkmarkc\checkmarki\checkmarkr\checkmarkc\checkmark$\checkmarkc\checkmark$\checkmark^\checkmark\\checkmarkc\checkmarki\checkmarkr\checkmarkc\checkmark$\checkmarka\checkmark$\checkmark^\checkmark\\checkmarkc\checkmarki\checkmarkr\checkmarkc\checkmark$\checkmarkt\checkmark$\checkmark^\checkmark\\checkmarkc\checkmarki\checkmarkr\checkmarkc\checkmark$\checkmarki\checkmark$\checkmark^\checkmark\\checkmarkc\checkmarki\checkmarkr\checkmarkc\checkmark$\checkmarko\checkmark$\checkmark^\checkmark\\checkmarkc\checkmarki\checkmarkr\checkmarkc\checkmark$\checkmarkn\checkmark$\checkmark^\checkmark\\checkmarkc\checkmarki\checkmarkr\checkmarkc\checkmark$\checkmark \checkmark$\checkmark^\checkmark\\checkmarkc\checkmarki\checkmarkr\checkmarkc\checkmark$\checkmarka\checkmark$\checkmark^\checkmark\\checkmarkc\checkmarki\checkmarkr\checkmarkc\checkmark$\checkmarku\checkmark$\checkmark^\checkmark\\checkmarkc\checkmarki\checkmarkr\checkmarkc\checkmark$\checkmark \checkmark$\checkmark^\checkmark\\checkmarkc\checkmarki\checkmarkr\checkmarkc\checkmark$\checkmarkC\checkmark$\checkmark^\checkmark\\checkmarkc\checkmarki\checkmarkr\checkmarkc\checkmark$\checkmarkr\checkmark$\checkmark^\checkmark\\checkmarkc\checkmarki\checkmarkr\checkmarkc\checkmark$\checkmarki\checkmark$\checkmark^\checkmark\\checkmarkc\checkmarki\checkmarkr\checkmarkc\checkmark$\checkmarkt\checkmark$\checkmark^\checkmark\\checkmarkc\checkmarki\checkmarkr\checkmarkc\checkmark$\checkmarkè\checkmark$\checkmark^\checkmark\\checkmarkc\checkmarki\checkmarkr\checkmarkc\checkmark$\checkmarkr\checkmark$\checkmark^\checkmark\\checkmarkc\checkmarki\checkmarkr\checkmarkc\checkmark$\checkmarke\checkmark$\checkmark^\checkmark\\checkmarkc\checkmarki\checkmarkr\checkmarkc\checkmark$\checkmark \checkmark$\checkmark^\checkmark\\checkmarkc\checkmarki\checkmarkr\checkmarkc\checkmark$\checkmarkd\checkmark$\checkmark^\checkmark\\checkmarkc\checkmarki\checkmarkr\checkmarkc\checkmark$\checkmarke\checkmark$\checkmark^\checkmark\\checkmarkc\checkmarki\checkmarkr\checkmarkc\checkmark$\checkmark \checkmark$\checkmark^\checkmark\\checkmarkc\checkmarki\checkmarkr\checkmarkc\checkmark$\checkmarkV\checkmark$\checkmark^\checkmark\\checkmarkc\checkmarki\checkmarkr\checkmarkc\checkmark$\checkmarko\checkmark$\checkmark^\checkmark\\checkmarkc\checkmarki\checkmarkr\checkmarkc\checkmark$\checkmarkn\checkmark$\checkmark^\checkmark\\checkmarkc\checkmarki\checkmarkr\checkmarkc\checkmark$\checkmark \checkmark$\checkmark^\checkmark\\checkmarkc\checkmarki\checkmarkr\checkmarkc\checkmark$\checkmarkM\checkmark$\checkmark^\checkmark\\checkmarkc\checkmarki\checkmarkr\checkmarkc\checkmark$\checkmarki\checkmark$\checkmark^\checkmark\\checkmarkc\checkmarki\checkmarkr\checkmarkc\checkmark$\checkmarks\checkmark$\checkmark^\checkmark\\checkmarkc\checkmarki\checkmarkr\checkmarkc\checkmark$\checkmarke\checkmark$\checkmark^\checkmark\\checkmarkc\checkmarki\checkmarkr\checkmarkc\checkmark$\checkmarks\checkmark$\checkmark^\checkmark\\checkmarkc\checkmarki\checkmarkr\checkmarkc\checkmark$\checkmark}\checkmark$\checkmark^\checkmark\\checkmarkc\checkmarki\checkmarkr\checkmarkc\checkmark$\checkmark
\checkmark$\checkmark^\checkmark\\checkmarkc\checkmarki\checkmarkr\checkmarkc\checkmark$\checkmarkL\checkmark$\checkmark^\checkmark\\checkmarkc\checkmarki\checkmarkr\checkmarkc\checkmark$\checkmarka\checkmark$\checkmark^\checkmark\\checkmarkc\checkmarki\checkmarkr\checkmarkc\checkmark$\checkmark \checkmark$\checkmark^\checkmark\\checkmarkc\checkmarki\checkmarkr\checkmarkc\checkmark$\checkmarkc\checkmark$\checkmark^\checkmark\\checkmarkc\checkmarki\checkmarkr\checkmarkc\checkmark$\checkmarko\checkmark$\checkmark^\checkmark\\checkmarkc\checkmarki\checkmarkr\checkmarkc\checkmark$\checkmarkn\checkmark$\checkmark^\checkmark\\checkmarkc\checkmarki\checkmarkr\checkmarkc\checkmark$\checkmarkt\checkmark$\checkmark^\checkmark\\checkmarkc\checkmarki\checkmarkr\checkmarkc\checkmark$\checkmarkr\checkmark$\checkmark^\checkmark\\checkmarkc\checkmarki\checkmarkr\checkmarkc\checkmark$\checkmarka\checkmark$\checkmark^\checkmark\\checkmarkc\checkmarki\checkmarkr\checkmarkc\checkmark$\checkmarki\checkmark$\checkmark^\checkmark\\checkmarkc\checkmarki\checkmarkr\checkmarkc\checkmark$\checkmarkn\checkmark$\checkmark^\checkmark\\checkmarkc\checkmarki\checkmarkr\checkmarkc\checkmark$\checkmarkt\checkmark$\checkmark^\checkmark\\checkmarkc\checkmarki\checkmarkr\checkmarkc\checkmark$\checkmarke\checkmark$\checkmark^\checkmark\\checkmarkc\checkmarki\checkmarkr\checkmarkc\checkmark$\checkmark \checkmark$\checkmark^\checkmark\\checkmarkc\checkmarki\checkmarkr\checkmarkc\checkmark$\checkmarkd\checkmark$\checkmark^\checkmark\\checkmarkc\checkmarki\checkmarkr\checkmarkc\checkmark$\checkmarke\checkmark$\checkmark^\checkmark\\checkmarkc\checkmarki\checkmarkr\checkmarkc\checkmark$\checkmark \checkmark$\checkmark^\checkmark\\checkmarkc\checkmarki\checkmarkr\checkmarkc\checkmark$\checkmarkf\checkmark$\checkmark^\checkmark\\checkmarkc\checkmarki\checkmarkr\checkmarkc\checkmark$\checkmarkl\checkmark$\checkmark^\checkmark\\checkmarkc\checkmarki\checkmarkr\checkmarkc\checkmark$\checkmarke\checkmark$\checkmark^\checkmark\\checkmarkc\checkmarki\checkmarkr\checkmarkc\checkmark$\checkmarkx\checkmark$\checkmark^\checkmark\\checkmarkc\checkmarki\checkmarkr\checkmarkc\checkmark$\checkmarki\checkmark$\checkmark^\checkmark\\checkmarkc\checkmarki\checkmarkr\checkmarkc\checkmark$\checkmarko\checkmark$\checkmark^\checkmark\\checkmarkc\checkmarki\checkmarkr\checkmarkc\checkmark$\checkmarkn\checkmark$\checkmark^\checkmark\\checkmarkc\checkmarki\checkmarkr\checkmarkc\checkmark$\checkmark \checkmark$\checkmark^\checkmark\\checkmarkc\checkmarki\checkmarkr\checkmarkc\checkmark$\checkmarkm\checkmark$\checkmark^\checkmark\\checkmarkc\checkmarki\checkmarkr\checkmarkc\checkmark$\checkmarka\checkmark$\checkmark^\checkmark\\checkmarkc\checkmarki\checkmarkr\checkmarkc\checkmark$\checkmarkx\checkmark$\checkmark^\checkmark\\checkmarkc\checkmarki\checkmarkr\checkmarkc\checkmark$\checkmarki\checkmark$\checkmark^\checkmark\\checkmarkc\checkmarki\checkmarkr\checkmarkc\checkmark$\checkmarkm\checkmark$\checkmark^\checkmark\\checkmarkc\checkmarki\checkmarkr\checkmarkc\checkmark$\checkmarka\checkmark$\checkmark^\checkmark\\checkmarkc\checkmarki\checkmarkr\checkmarkc\checkmark$\checkmarkl\checkmark$\checkmark^\checkmark\\checkmarkc\checkmarki\checkmarkr\checkmarkc\checkmark$\checkmarke\checkmark$\checkmark^\checkmark\\checkmarkc\checkmarki\checkmarkr\checkmarkc\checkmark$\checkmark \checkmark$\checkmark^\checkmark\\checkmarkc\checkmarki\checkmarkr\checkmarkc\checkmark$\checkmarka\checkmark$\checkmark^\checkmark\\checkmarkc\checkmarki\checkmarkr\checkmarkc\checkmark$\checkmarku\checkmark$\checkmark^\checkmark\\checkmarkc\checkmarki\checkmarkr\checkmarkc\checkmark$\checkmark \checkmark$\checkmark^\checkmark\\checkmarkc\checkmarki\checkmarkr\checkmarkc\checkmark$\checkmarkc\checkmark$\checkmark^\checkmark\\checkmarkc\checkmarki\checkmarkr\checkmarkc\checkmark$\checkmarke\checkmark$\checkmark^\checkmark\\checkmarkc\checkmarki\checkmarkr\checkmarkc\checkmark$\checkmarkn\checkmark$\checkmark^\checkmark\\checkmarkc\checkmarki\checkmarkr\checkmarkc\checkmark$\checkmarkt\checkmark$\checkmark^\checkmark\\checkmarkc\checkmarki\checkmarkr\checkmarkc\checkmark$\checkmarkr\checkmark$\checkmark^\checkmark\\checkmarkc\checkmarki\checkmarkr\checkmarkc\checkmark$\checkmarke\checkmark$\checkmark^\checkmark\\checkmarkc\checkmarki\checkmarkr\checkmarkc\checkmark$\checkmark \checkmark$\checkmark^\checkmark\\checkmarkc\checkmarki\checkmarkr\checkmarkc\checkmark$\checkmarkd\checkmark$\checkmark^\checkmark\\checkmarkc\checkmarki\checkmarkr\checkmarkc\checkmark$\checkmarku\checkmark$\checkmark^\checkmark\\checkmarkc\checkmarki\checkmarkr\checkmarkc\checkmark$\checkmark \checkmark$\checkmark^\checkmark\\checkmarkc\checkmarki\checkmarkr\checkmarkc\checkmark$\checkmarkp\checkmark$\checkmark^\checkmark\\checkmarkc\checkmarki\checkmarkr\checkmarkc\checkmark$\checkmarko\checkmark$\checkmark^\checkmark\\checkmarkc\checkmarki\checkmarkr\checkmarkc\checkmark$\checkmarkr\checkmark$\checkmark^\checkmark\\checkmarkc\checkmarki\checkmarkr\checkmarkc\checkmark$\checkmarkt\checkmark$\checkmark^\checkmark\\checkmarkc\checkmarki\checkmarkr\checkmarkc\checkmark$\checkmarki\checkmark$\checkmark^\checkmark\\checkmarkc\checkmarki\checkmarkr\checkmarkc\checkmark$\checkmarkq\checkmark$\checkmark^\checkmark\\checkmarkc\checkmarki\checkmarkr\checkmarkc\checkmark$\checkmarku\checkmark$\checkmark^\checkmark\\checkmarkc\checkmarki\checkmarkr\checkmarkc\checkmark$\checkmarke\checkmark$\checkmark^\checkmark\\checkmarkc\checkmarki\checkmarkr\checkmarkc\checkmark$\checkmark \checkmark$\checkmark^\checkmark\\checkmarkc\checkmarki\checkmarkr\checkmarkc\checkmark$\checkmark:\checkmark$\checkmark^\checkmark\\checkmarkc\checkmarki\checkmarkr\checkmarkc\checkmark$\checkmark
\checkmark$\checkmark^\checkmark\\checkmarkc\checkmarki\checkmarkr\checkmarkc\checkmark$\checkmark\\checkmark$\checkmark^\checkmark\\checkmarkc\checkmarki\checkmarkr\checkmarkc\checkmark$\checkmarkb\checkmark$\checkmark^\checkmark\\checkmarkc\checkmarki\checkmarkr\checkmarkc\checkmark$\checkmarke\checkmark$\checkmark^\checkmark\\checkmarkc\checkmarki\checkmarkr\checkmarkc\checkmark$\checkmarkg\checkmark$\checkmark^\checkmark\\checkmarkc\checkmarki\checkmarkr\checkmarkc\checkmark$\checkmarki\checkmark$\checkmark^\checkmark\\checkmarkc\checkmarki\checkmarkr\checkmarkc\checkmark$\checkmarkn\checkmark$\checkmark^\checkmark\\checkmarkc\checkmarki\checkmarkr\checkmarkc\checkmark$\checkmark{\checkmark$\checkmark^\checkmark\\checkmarkc\checkmarki\checkmarkr\checkmarkc\checkmark$\checkmarke\checkmark$\checkmark^\checkmark\\checkmarkc\checkmarki\checkmarkr\checkmarkc\checkmark$\checkmarkq\checkmark$\checkmark^\checkmark\\checkmarkc\checkmarki\checkmarkr\checkmarkc\checkmark$\checkmarku\checkmark$\checkmark^\checkmark\\checkmarkc\checkmarki\checkmarkr\checkmarkc\checkmark$\checkmarka\checkmark$\checkmark^\checkmark\\checkmarkc\checkmarki\checkmarkr\checkmarkc\checkmark$\checkmarkt\checkmark$\checkmark^\checkmark\\checkmarkc\checkmarki\checkmarkr\checkmarkc\checkmark$\checkmarki\checkmark$\checkmark^\checkmark\\checkmarkc\checkmarki\checkmarkr\checkmarkc\checkmark$\checkmarko\checkmark$\checkmark^\checkmark\\checkmarkc\checkmarki\checkmarkr\checkmarkc\checkmark$\checkmarkn\checkmark$\checkmark^\checkmark\\checkmarkc\checkmarki\checkmarkr\checkmarkc\checkmark$\checkmark}\checkmark$\checkmark^\checkmark\\checkmarkc\checkmarki\checkmarkr\checkmarkc\checkmark$\checkmark
\checkmark$\checkmark^\checkmark\\checkmarkc\checkmarki\checkmarkr\checkmarkc\checkmark$\checkmark \checkmark$\checkmark^\checkmark\\checkmarkc\checkmarki\checkmarkr\checkmarkc\checkmark$\checkmark \checkmark$\checkmark^\checkmark\\checkmarkc\checkmarki\checkmarkr\checkmarkc\checkmark$\checkmark \checkmark$\checkmark^\checkmark\\checkmarkc\checkmarki\checkmarkr\checkmarkc\checkmark$\checkmark \checkmark$\checkmark^\checkmark\\checkmarkc\checkmarki\checkmarkr\checkmarkc\checkmark$\checkmark\\checkmark$\checkmark^\checkmark\\checkmarkc\checkmarki\checkmarkr\checkmarkc\checkmark$\checkmarks\checkmark$\checkmark^\checkmark\\checkmarkc\checkmarki\checkmarkr\checkmarkc\checkmark$\checkmarki\checkmark$\checkmark^\checkmark\\checkmarkc\checkmarki\checkmarkr\checkmarkc\checkmark$\checkmarkg\checkmark$\checkmark^\checkmark\\checkmarkc\checkmarki\checkmarkr\checkmarkc\checkmark$\checkmarkm\checkmark$\checkmark^\checkmark\\checkmarkc\checkmarki\checkmarkr\checkmarkc\checkmark$\checkmarka\checkmark$\checkmark^\checkmark\\checkmarkc\checkmarki\checkmarkr\checkmarkc\checkmark$\checkmark_\checkmark$\checkmark^\checkmark\\checkmarkc\checkmarki\checkmarkr\checkmarkc\checkmark$\checkmark{\checkmark$\checkmark^\checkmark\\checkmarkc\checkmarki\checkmarkr\checkmarkc\checkmark$\checkmarkf\checkmark$\checkmark^\checkmark\\checkmarkc\checkmarki\checkmarkr\checkmarkc\checkmark$\checkmarkl\checkmark$\checkmark^\checkmark\\checkmarkc\checkmarki\checkmarkr\checkmarkc\checkmark$\checkmarke\checkmark$\checkmark^\checkmark\\checkmarkc\checkmarki\checkmarkr\checkmarkc\checkmark$\checkmarkx\checkmark$\checkmark^\checkmark\\checkmarkc\checkmarki\checkmarkr\checkmarkc\checkmark$\checkmark}\checkmark$\checkmark^\checkmark\\checkmarkc\checkmarki\checkmarkr\checkmarkc\checkmark$\checkmark \checkmark$\checkmark^\checkmark\\checkmarkc\checkmarki\checkmarkr\checkmarkc\checkmark$\checkmark=\checkmark$\checkmark^\checkmark\\checkmarkc\checkmarki\checkmarkr\checkmarkc\checkmark$\checkmark \checkmark$\checkmark^\checkmark\\checkmarkc\checkmarki\checkmarkr\checkmarkc\checkmark$\checkmark\\checkmark$\checkmark^\checkmark\\checkmarkc\checkmarki\checkmarkr\checkmarkc\checkmark$\checkmarkf\checkmark$\checkmark^\checkmark\\checkmarkc\checkmarki\checkmarkr\checkmarkc\checkmark$\checkmarkr\checkmark$\checkmark^\checkmark\\checkmarkc\checkmarki\checkmarkr\checkmarkc\checkmark$\checkmarka\checkmark$\checkmark^\checkmark\\checkmarkc\checkmarki\checkmarkr\checkmarkc\checkmark$\checkmarkc\checkmark$\checkmark^\checkmark\\checkmarkc\checkmarki\checkmarkr\checkmarkc\checkmark$\checkmark{\checkmark$\checkmark^\checkmark\\checkmarkc\checkmarki\checkmarkr\checkmarkc\checkmark$\checkmarkM\checkmark$\checkmark^\checkmark\\checkmarkc\checkmarki\checkmarkr\checkmarkc\checkmark$\checkmark_\checkmark$\checkmark^\checkmark\\checkmarkc\checkmarki\checkmarkr\checkmarkc\checkmark$\checkmark{\checkmark$\checkmark^\checkmark\\checkmarkc\checkmarki\checkmarkr\checkmarkc\checkmark$\checkmarkm\checkmark$\checkmark^\checkmark\\checkmarkc\checkmarki\checkmarkr\checkmarkc\checkmark$\checkmarka\checkmark$\checkmark^\checkmark\\checkmarkc\checkmarki\checkmarkr\checkmarkc\checkmark$\checkmarkx\checkmark$\checkmark^\checkmark\\checkmarkc\checkmarki\checkmarkr\checkmarkc\checkmark$\checkmark}\checkmark$\checkmark^\checkmark\\checkmarkc\checkmarki\checkmarkr\checkmarkc\checkmark$\checkmark \checkmark$\checkmark^\checkmark\\checkmarkc\checkmarki\checkmarkr\checkmarkc\checkmark$\checkmark\\checkmark$\checkmark^\checkmark\\checkmarkc\checkmarki\checkmarkr\checkmarkc\checkmark$\checkmarkc\checkmark$\checkmark^\checkmark\\checkmarkc\checkmarki\checkmarkr\checkmarkc\checkmark$\checkmarkd\checkmark$\checkmark^\checkmark\\checkmarkc\checkmarki\checkmarkr\checkmarkc\checkmark$\checkmarko\checkmark$\checkmark^\checkmark\\checkmarkc\checkmarki\checkmarkr\checkmarkc\checkmark$\checkmarkt\checkmark$\checkmark^\checkmark\\checkmarkc\checkmarki\checkmarkr\checkmarkc\checkmark$\checkmark \checkmark$\checkmark^\checkmark\\checkmarkc\checkmarki\checkmarkr\checkmarkc\checkmark$\checkmarkv\checkmark$\checkmark^\checkmark\\checkmarkc\checkmarki\checkmarkr\checkmarkc\checkmark$\checkmark}\checkmark$\checkmark^\checkmark\\checkmarkc\checkmarki\checkmarkr\checkmarkc\checkmark$\checkmark{\checkmark$\checkmark^\checkmark\\checkmarkc\checkmarki\checkmarkr\checkmarkc\checkmark$\checkmarkI\checkmark$\checkmark^\checkmark\\checkmarkc\checkmarki\checkmarkr\checkmarkc\checkmark$\checkmark_\checkmark$\checkmark^\checkmark\\checkmarkc\checkmarki\checkmarkr\checkmarkc\checkmark$\checkmark{\checkmark$\checkmark^\checkmark\\checkmarkc\checkmarki\checkmarkr\checkmarkc\checkmark$\checkmarky\checkmark$\checkmark^\checkmark\\checkmarkc\checkmarki\checkmarkr\checkmarkc\checkmark$\checkmarky\checkmark$\checkmark^\checkmark\\checkmarkc\checkmarki\checkmarkr\checkmarkc\checkmark$\checkmark}\checkmark$\checkmark^\checkmark\\checkmarkc\checkmarki\checkmarkr\checkmarkc\checkmark$\checkmark}\checkmark$\checkmark^\checkmark\\checkmarkc\checkmarki\checkmarkr\checkmarkc\checkmark$\checkmark \checkmark$\checkmark^\checkmark\\checkmarkc\checkmarki\checkmarkr\checkmarkc\checkmark$\checkmark=\checkmark$\checkmark^\checkmark\\checkmarkc\checkmarki\checkmarkr\checkmarkc\checkmark$\checkmark \checkmark$\checkmark^\checkmark\\checkmarkc\checkmarki\checkmarkr\checkmarkc\checkmark$\checkmark\\checkmark$\checkmark^\checkmark\\checkmarkc\checkmarki\checkmarkr\checkmarkc\checkmark$\checkmarkf\checkmark$\checkmark^\checkmark\\checkmarkc\checkmarki\checkmarkr\checkmarkc\checkmark$\checkmarkr\checkmark$\checkmark^\checkmark\\checkmarkc\checkmarki\checkmarkr\checkmarkc\checkmark$\checkmarka\checkmark$\checkmark^\checkmark\\checkmarkc\checkmarki\checkmarkr\checkmarkc\checkmark$\checkmarkc\checkmark$\checkmark^\checkmark\\checkmarkc\checkmarki\checkmarkr\checkmarkc\checkmark$\checkmark{\checkmark$\checkmark^\checkmark\\checkmarkc\checkmarki\checkmarkr\checkmarkc\checkmark$\checkmark(\checkmark$\checkmark^\checkmark\\checkmarkc\checkmarki\checkmarkr\checkmarkc\checkmark$\checkmarkF\checkmark$\checkmark^\checkmark\\checkmarkc\checkmarki\checkmarkr\checkmarkc\checkmark$\checkmark_\checkmark$\checkmark^\checkmark\\checkmarkc\checkmarki\checkmarkr\checkmarkc\checkmark$\checkmark{\checkmark$\checkmark^\checkmark\\checkmarkc\checkmarki\checkmarkr\checkmarkc\checkmark$\checkmarkm\checkmark$\checkmark^\checkmark\\checkmarkc\checkmarki\checkmarkr\checkmarkc\checkmark$\checkmarka\checkmark$\checkmark^\checkmark\\checkmarkc\checkmarki\checkmarkr\checkmarkc\checkmark$\checkmarkx\checkmark$\checkmark^\checkmark\\checkmarkc\checkmarki\checkmarkr\checkmarkc\checkmark$\checkmark}\checkmark$\checkmark^\checkmark\\checkmarkc\checkmarki\checkmarkr\checkmarkc\checkmark$\checkmark \checkmark$\checkmark^\checkmark\\checkmarkc\checkmarki\checkmarkr\checkmarkc\checkmark$\checkmark\\checkmark$\checkmark^\checkmark\\checkmarkc\checkmarki\checkmarkr\checkmarkc\checkmark$\checkmarkt\checkmark$\checkmark^\checkmark\\checkmarkc\checkmarki\checkmarkr\checkmarkc\checkmark$\checkmarki\checkmark$\checkmark^\checkmark\\checkmarkc\checkmarki\checkmarkr\checkmarkc\checkmark$\checkmarkm\checkmark$\checkmark^\checkmark\\checkmarkc\checkmarki\checkmarkr\checkmarkc\checkmark$\checkmarke\checkmark$\checkmark^\checkmark\\checkmarkc\checkmarki\checkmarkr\checkmarkc\checkmark$\checkmarks\checkmark$\checkmark^\checkmark\\checkmarkc\checkmarki\checkmarkr\checkmarkc\checkmark$\checkmark \checkmark$\checkmark^\checkmark\\checkmarkc\checkmarki\checkmarkr\checkmarkc\checkmark$\checkmarkL\checkmark$\checkmark^\checkmark\\checkmarkc\checkmarki\checkmarkr\checkmarkc\checkmark$\checkmark/\checkmark$\checkmark^\checkmark\\checkmarkc\checkmarki\checkmarkr\checkmarkc\checkmark$\checkmark4\checkmark$\checkmark^\checkmark\\checkmarkc\checkmarki\checkmarkr\checkmarkc\checkmark$\checkmark)\checkmark$\checkmark^\checkmark\\checkmarkc\checkmarki\checkmarkr\checkmarkc\checkmark$\checkmark \checkmark$\checkmark^\checkmark\\checkmarkc\checkmarki\checkmarkr\checkmarkc\checkmark$\checkmark\\checkmark$\checkmark^\checkmark\\checkmarkc\checkmarki\checkmarkr\checkmarkc\checkmark$\checkmarkc\checkmark$\checkmark^\checkmark\\checkmarkc\checkmarki\checkmarkr\checkmarkc\checkmark$\checkmarkd\checkmark$\checkmark^\checkmark\\checkmarkc\checkmarki\checkmarkr\checkmarkc\checkmark$\checkmarko\checkmark$\checkmark^\checkmark\\checkmarkc\checkmarki\checkmarkr\checkmarkc\checkmark$\checkmarkt\checkmark$\checkmark^\checkmark\\checkmarkc\checkmarki\checkmarkr\checkmarkc\checkmark$\checkmark \checkmark$\checkmark^\checkmark\\checkmarkc\checkmarki\checkmarkr\checkmarkc\checkmark$\checkmark(\checkmark$\checkmark^\checkmark\\checkmarkc\checkmarki\checkmarkr\checkmarkc\checkmark$\checkmarkh\checkmark$\checkmark^\checkmark\\checkmarkc\checkmarki\checkmarkr\checkmarkc\checkmark$\checkmark/\checkmark$\checkmark^\checkmark\\checkmarkc\checkmarki\checkmarkr\checkmarkc\checkmark$\checkmark2\checkmark$\checkmark^\checkmark\\checkmarkc\checkmarki\checkmarkr\checkmarkc\checkmark$\checkmark)\checkmark$\checkmark^\checkmark\\checkmarkc\checkmarki\checkmarkr\checkmarkc\checkmark$\checkmark}\checkmark$\checkmark^\checkmark\\checkmarkc\checkmarki\checkmarkr\checkmarkc\checkmark$\checkmark{\checkmark$\checkmark^\checkmark\\checkmarkc\checkmarki\checkmarkr\checkmarkc\checkmark$\checkmarkI\checkmark$\checkmark^\checkmark\\checkmarkc\checkmarki\checkmarkr\checkmarkc\checkmark$\checkmark_\checkmark$\checkmark^\checkmark\\checkmarkc\checkmarki\checkmarkr\checkmarkc\checkmark$\checkmark{\checkmark$\checkmark^\checkmark\\checkmarkc\checkmarki\checkmarkr\checkmarkc\checkmark$\checkmarky\checkmark$\checkmark^\checkmark\\checkmarkc\checkmarki\checkmarkr\checkmarkc\checkmark$\checkmarky\checkmark$\checkmark^\checkmark\\checkmarkc\checkmarki\checkmarkr\checkmarkc\checkmark$\checkmark}\checkmark$\checkmark^\checkmark\\checkmarkc\checkmarki\checkmarkr\checkmarkc\checkmark$\checkmark}\checkmark$\checkmark^\checkmark\\checkmarkc\checkmarki\checkmarkr\checkmarkc\checkmark$\checkmark \checkmark$\checkmark^\checkmark\\checkmarkc\checkmarki\checkmarkr\checkmarkc\checkmark$\checkmark=\checkmark$\checkmark^\checkmark\\checkmarkc\checkmarki\checkmarkr\checkmarkc\checkmark$\checkmark \checkmark$\checkmark^\checkmark\\checkmarkc\checkmarki\checkmarkr\checkmarkc\checkmark$\checkmark\\checkmark$\checkmark^\checkmark\\checkmarkc\checkmarki\checkmarkr\checkmarkc\checkmark$\checkmarkf\checkmark$\checkmark^\checkmark\\checkmarkc\checkmarki\checkmarkr\checkmarkc\checkmark$\checkmarkr\checkmark$\checkmark^\checkmark\\checkmarkc\checkmarki\checkmarkr\checkmarkc\checkmark$\checkmarka\checkmark$\checkmark^\checkmark\\checkmarkc\checkmarki\checkmarkr\checkmarkc\checkmark$\checkmarkc\checkmark$\checkmark^\checkmark\\checkmarkc\checkmarki\checkmarkr\checkmarkc\checkmark$\checkmark{\checkmark$\checkmark^\checkmark\\checkmarkc\checkmarki\checkmarkr\checkmarkc\checkmark$\checkmark(\checkmark$\checkmark^\checkmark\\checkmarkc\checkmarki\checkmarkr\checkmarkc\checkmark$\checkmark1\checkmark$\checkmark^\checkmark\\checkmarkc\checkmarki\checkmarkr\checkmarkc\checkmark$\checkmark4\checkmark$\checkmark^\checkmark\\checkmarkc\checkmarki\checkmarkr\checkmarkc\checkmark$\checkmark0\checkmark$\checkmark^\checkmark\\checkmarkc\checkmarki\checkmarkr\checkmarkc\checkmark$\checkmark \checkmark$\checkmark^\checkmark\\checkmarkc\checkmarki\checkmarkr\checkmarkc\checkmark$\checkmark\\checkmark$\checkmark^\checkmark\\checkmarkc\checkmarki\checkmarkr\checkmarkc\checkmark$\checkmarkt\checkmark$\checkmark^\checkmark\\checkmarkc\checkmarki\checkmarkr\checkmarkc\checkmark$\checkmarki\checkmark$\checkmark^\checkmark\\checkmarkc\checkmarki\checkmarkr\checkmarkc\checkmark$\checkmarkm\checkmark$\checkmark^\checkmark\\checkmarkc\checkmarki\checkmarkr\checkmarkc\checkmark$\checkmarke\checkmark$\checkmark^\checkmark\\checkmarkc\checkmarki\checkmarkr\checkmarkc\checkmark$\checkmarks\checkmark$\checkmark^\checkmark\\checkmarkc\checkmarki\checkmarkr\checkmarkc\checkmark$\checkmark \checkmark$\checkmark^\checkmark\\checkmarkc\checkmarki\checkmarkr\checkmarkc\checkmark$\checkmark1\checkmark$\checkmark^\checkmark\\checkmarkc\checkmarki\checkmarkr\checkmarkc\checkmark$\checkmark0\checkmark$\checkmark^\checkmark\\checkmarkc\checkmarki\checkmarkr\checkmarkc\checkmark$\checkmark0\checkmark$\checkmark^\checkmark\\checkmarkc\checkmarki\checkmarkr\checkmarkc\checkmark$\checkmark)\checkmark$\checkmark^\checkmark\\checkmarkc\checkmarki\checkmarkr\checkmarkc\checkmark$\checkmark \checkmark$\checkmark^\checkmark\\checkmarkc\checkmarki\checkmarkr\checkmarkc\checkmark$\checkmark\\checkmark$\checkmark^\checkmark\\checkmarkc\checkmarki\checkmarkr\checkmarkc\checkmark$\checkmarkt\checkmark$\checkmark^\checkmark\\checkmarkc\checkmarki\checkmarkr\checkmarkc\checkmark$\checkmarki\checkmark$\checkmark^\checkmark\\checkmarkc\checkmarki\checkmarkr\checkmarkc\checkmark$\checkmarkm\checkmark$\checkmark^\checkmark\\checkmarkc\checkmarki\checkmarkr\checkmarkc\checkmark$\checkmarke\checkmark$\checkmark^\checkmark\\checkmarkc\checkmarki\checkmarkr\checkmarkc\checkmark$\checkmarks\checkmark$\checkmark^\checkmark\\checkmarkc\checkmarki\checkmarkr\checkmarkc\checkmark$\checkmark \checkmark$\checkmark^\checkmark\\checkmarkc\checkmarki\checkmarkr\checkmarkc\checkmark$\checkmark4\checkmark$\checkmark^\checkmark\\checkmarkc\checkmarki\checkmarkr\checkmarkc\checkmark$\checkmark0\checkmark$\checkmark^\checkmark\\checkmarkc\checkmarki\checkmarkr\checkmarkc\checkmark$\checkmark}\checkmark$\checkmark^\checkmark\\checkmarkc\checkmarki\checkmarkr\checkmarkc\checkmark$\checkmark{\checkmark$\checkmark^\checkmark\\checkmarkc\checkmarki\checkmarkr\checkmarkc\checkmark$\checkmark6\checkmark$\checkmark^\checkmark\\checkmarkc\checkmarki\checkmarkr\checkmarkc\checkmark$\checkmark5\checkmark$\checkmark^\checkmark\\checkmarkc\checkmarki\checkmarkr\checkmarkc\checkmark$\checkmark\\checkmark$\checkmark^\checkmark\\checkmarkc\checkmarki\checkmarkr\checkmarkc\checkmark$\checkmark,\checkmark$\checkmark^\checkmark\\checkmarkc\checkmarki\checkmarkr\checkmarkc\checkmark$\checkmark0\checkmark$\checkmark^\checkmark\\checkmarkc\checkmarki\checkmarkr\checkmarkc\checkmark$\checkmark0\checkmark$\checkmark^\checkmark\\checkmarkc\checkmarki\checkmarkr\checkmarkc\checkmark$\checkmark0\checkmark$\checkmark^\checkmark\\checkmarkc\checkmarki\checkmarkr\checkmarkc\checkmark$\checkmark}\checkmark$\checkmark^\checkmark\\checkmarkc\checkmarki\checkmarkr\checkmarkc\checkmark$\checkmark \checkmark$\checkmark^\checkmark\\checkmarkc\checkmarki\checkmarkr\checkmarkc\checkmark$\checkmark=\checkmark$\checkmark^\checkmark\\checkmarkc\checkmarki\checkmarkr\checkmarkc\checkmark$\checkmark \checkmark$\checkmark^\checkmark\\checkmarkc\checkmarki\checkmarkr\checkmarkc\checkmark$\checkmark\\checkmark$\checkmark^\checkmark\\checkmarkc\checkmarki\checkmarkr\checkmarkc\checkmark$\checkmarkf\checkmark$\checkmark^\checkmark\\checkmarkc\checkmarki\checkmarkr\checkmarkc\checkmark$\checkmarkr\checkmark$\checkmark^\checkmark\\checkmarkc\checkmarki\checkmarkr\checkmarkc\checkmark$\checkmarka\checkmark$\checkmark^\checkmark\\checkmarkc\checkmarki\checkmarkr\checkmarkc\checkmark$\checkmarkc\checkmark$\checkmark^\checkmark\\checkmarkc\checkmarki\checkmarkr\checkmarkc\checkmark$\checkmark{\checkmark$\checkmark^\checkmark\\checkmarkc\checkmarki\checkmarkr\checkmarkc\checkmark$\checkmark5\checkmark$\checkmark^\checkmark\\checkmarkc\checkmarki\checkmarkr\checkmarkc\checkmark$\checkmark6\checkmark$\checkmark^\checkmark\\checkmarkc\checkmarki\checkmarkr\checkmarkc\checkmark$\checkmark0\checkmark$\checkmark^\checkmark\\checkmarkc\checkmarki\checkmarkr\checkmarkc\checkmark$\checkmark\\checkmark$\checkmark^\checkmark\\checkmarkc\checkmarki\checkmarkr\checkmarkc\checkmark$\checkmark,\checkmark$\checkmark^\checkmark\\checkmarkc\checkmarki\checkmarkr\checkmarkc\checkmark$\checkmark0\checkmark$\checkmark^\checkmark\\checkmarkc\checkmarki\checkmarkr\checkmarkc\checkmark$\checkmark0\checkmark$\checkmark^\checkmark\\checkmarkc\checkmarki\checkmarkr\checkmarkc\checkmark$\checkmark0\checkmark$\checkmark^\checkmark\\checkmarkc\checkmarki\checkmarkr\checkmarkc\checkmark$\checkmark}\checkmark$\checkmark^\checkmark\\checkmarkc\checkmarki\checkmarkr\checkmarkc\checkmark$\checkmark{\checkmark$\checkmark^\checkmark\\checkmarkc\checkmarki\checkmarkr\checkmarkc\checkmark$\checkmark6\checkmark$\checkmark^\checkmark\\checkmarkc\checkmarki\checkmarkr\checkmarkc\checkmark$\checkmark5\checkmark$\checkmark^\checkmark\\checkmarkc\checkmarki\checkmarkr\checkmarkc\checkmark$\checkmark\\checkmark$\checkmark^\checkmark\\checkmarkc\checkmarki\checkmarkr\checkmarkc\checkmark$\checkmark,\checkmark$\checkmark^\checkmark\\checkmarkc\checkmarki\checkmarkr\checkmarkc\checkmark$\checkmark0\checkmark$\checkmark^\checkmark\\checkmarkc\checkmarki\checkmarkr\checkmarkc\checkmark$\checkmark0\checkmark$\checkmark^\checkmark\\checkmarkc\checkmarki\checkmarkr\checkmarkc\checkmark$\checkmark0\checkmark$\checkmark^\checkmark\\checkmarkc\checkmarki\checkmarkr\checkmarkc\checkmark$\checkmark}\checkmark$\checkmark^\checkmark\\checkmarkc\checkmarki\checkmarkr\checkmarkc\checkmark$\checkmark \checkmark$\checkmark^\checkmark\\checkmarkc\checkmarki\checkmarkr\checkmarkc\checkmark$\checkmark\\checkmark$\checkmark^\checkmark\\checkmarkc\checkmarki\checkmarkr\checkmarkc\checkmark$\checkmarka\checkmark$\checkmark^\checkmark\\checkmarkc\checkmarki\checkmarkr\checkmarkc\checkmark$\checkmarkp\checkmark$\checkmark^\checkmark\\checkmarkc\checkmarki\checkmarkr\checkmarkc\checkmark$\checkmarkp\checkmark$\checkmark^\checkmark\\checkmarkc\checkmarki\checkmarkr\checkmarkc\checkmark$\checkmarkr\checkmark$\checkmark^\checkmark\\checkmarkc\checkmarki\checkmarkr\checkmarkc\checkmark$\checkmarko\checkmark$\checkmark^\checkmark\\checkmarkc\checkmarki\checkmarkr\checkmarkc\checkmark$\checkmarkx\checkmark$\checkmark^\checkmark\\checkmarkc\checkmarki\checkmarkr\checkmarkc\checkmark$\checkmark \checkmark$\checkmark^\checkmark\\checkmarkc\checkmarki\checkmarkr\checkmarkc\checkmark$\checkmark8\checkmark$\checkmark^\checkmark\\checkmarkc\checkmarki\checkmarkr\checkmarkc\checkmark$\checkmark{\checkmark$\checkmark^\checkmark\\checkmarkc\checkmarki\checkmarkr\checkmarkc\checkmark$\checkmark,\checkmark$\checkmark^\checkmark\\checkmarkc\checkmarki\checkmarkr\checkmarkc\checkmark$\checkmark}\checkmark$\checkmark^\checkmark\\checkmarkc\checkmarki\checkmarkr\checkmarkc\checkmark$\checkmark6\checkmark$\checkmark^\checkmark\\checkmarkc\checkmarki\checkmarkr\checkmarkc\checkmark$\checkmark \checkmark$\checkmark^\checkmark\\checkmarkc\checkmarki\checkmarkr\checkmarkc\checkmark$\checkmark\\checkmark$\checkmark^\checkmark\\checkmarkc\checkmarki\checkmarkr\checkmarkc\checkmark$\checkmarkt\checkmark$\checkmark^\checkmark\\checkmarkc\checkmarki\checkmarkr\checkmarkc\checkmark$\checkmarke\checkmark$\checkmark^\checkmark\\checkmarkc\checkmarki\checkmarkr\checkmarkc\checkmark$\checkmarkx\checkmark$\checkmark^\checkmark\\checkmarkc\checkmarki\checkmarkr\checkmarkc\checkmark$\checkmarkt\checkmark$\checkmark^\checkmark\\checkmarkc\checkmarki\checkmarkr\checkmarkc\checkmark$\checkmark{\checkmark$\checkmark^\checkmark\\checkmarkc\checkmarki\checkmarkr\checkmarkc\checkmark$\checkmark \checkmark$\checkmark^\checkmark\\checkmarkc\checkmarki\checkmarkr\checkmarkc\checkmark$\checkmarkM\checkmark$\checkmark^\checkmark\\checkmarkc\checkmarki\checkmarkr\checkmarkc\checkmark$\checkmarkP\checkmark$\checkmark^\checkmark\\checkmarkc\checkmarki\checkmarkr\checkmarkc\checkmark$\checkmarka\checkmark$\checkmark^\checkmark\\checkmarkc\checkmarki\checkmarkr\checkmarkc\checkmark$\checkmark}\checkmark$\checkmark^\checkmark\\checkmarkc\checkmarki\checkmarkr\checkmarkc\checkmark$\checkmark
\checkmark$\checkmark^\checkmark\\checkmarkc\checkmarki\checkmarkr\checkmarkc\checkmark$\checkmark\\checkmark$\checkmark^\checkmark\\checkmarkc\checkmarki\checkmarkr\checkmarkc\checkmark$\checkmarke\checkmark$\checkmark^\checkmark\\checkmarkc\checkmarki\checkmarkr\checkmarkc\checkmark$\checkmarkn\checkmark$\checkmark^\checkmark\\checkmarkc\checkmarki\checkmarkr\checkmarkc\checkmark$\checkmarkd\checkmark$\checkmark^\checkmark\\checkmarkc\checkmarki\checkmarkr\checkmarkc\checkmark$\checkmark{\checkmark$\checkmark^\checkmark\\checkmarkc\checkmarki\checkmarkr\checkmarkc\checkmark$\checkmarke\checkmark$\checkmark^\checkmark\\checkmarkc\checkmarki\checkmarkr\checkmarkc\checkmark$\checkmarkq\checkmark$\checkmark^\checkmark\\checkmarkc\checkmarki\checkmarkr\checkmarkc\checkmark$\checkmarku\checkmark$\checkmark^\checkmark\\checkmarkc\checkmarki\checkmarkr\checkmarkc\checkmark$\checkmarka\checkmark$\checkmark^\checkmark\\checkmarkc\checkmarki\checkmarkr\checkmarkc\checkmark$\checkmarkt\checkmark$\checkmark^\checkmark\\checkmarkc\checkmarki\checkmarkr\checkmarkc\checkmark$\checkmarki\checkmark$\checkmark^\checkmark\\checkmarkc\checkmarki\checkmarkr\checkmarkc\checkmark$\checkmarko\checkmark$\checkmark^\checkmark\\checkmarkc\checkmarki\checkmarkr\checkmarkc\checkmark$\checkmarkn\checkmark$\checkmark^\checkmark\\checkmarkc\checkmarki\checkmarkr\checkmarkc\checkmark$\checkmark}\checkmark$\checkmark^\checkmark\\checkmarkc\checkmarki\checkmarkr\checkmarkc\checkmark$\checkmark
\checkmark$\checkmark^\checkmark\\checkmarkc\checkmarki\checkmarkr\checkmarkc\checkmark$\checkmark
\checkmark$\checkmark^\checkmark\\checkmarkc\checkmarki\checkmarkr\checkmarkc\checkmark$\checkmarkL\checkmark$\checkmark^\checkmark\\checkmarkc\checkmarki\checkmarkr\checkmarkc\checkmark$\checkmarka\checkmark$\checkmark^\checkmark\\checkmarkc\checkmarki\checkmarkr\checkmarkc\checkmark$\checkmark \checkmark$\checkmark^\checkmark\\checkmarkc\checkmarki\checkmarkr\checkmarkc\checkmark$\checkmarkc\checkmark$\checkmark^\checkmark\\checkmarkc\checkmarki\checkmarkr\checkmarkc\checkmark$\checkmarko\checkmark$\checkmark^\checkmark\\checkmarkc\checkmarki\checkmarkr\checkmarkc\checkmark$\checkmarkn\checkmark$\checkmark^\checkmark\\checkmarkc\checkmarki\checkmarkr\checkmarkc\checkmark$\checkmarkt\checkmark$\checkmark^\checkmark\\checkmarkc\checkmarki\checkmarkr\checkmarkc\checkmark$\checkmarkr\checkmark$\checkmark^\checkmark\\checkmarkc\checkmarki\checkmarkr\checkmarkc\checkmark$\checkmarka\checkmark$\checkmark^\checkmark\\checkmarkc\checkmarki\checkmarkr\checkmarkc\checkmark$\checkmarki\checkmark$\checkmark^\checkmark\\checkmarkc\checkmarki\checkmarkr\checkmarkc\checkmark$\checkmarkn\checkmark$\checkmark^\checkmark\\checkmarkc\checkmarki\checkmarkr\checkmarkc\checkmark$\checkmarkt\checkmark$\checkmark^\checkmark\\checkmarkc\checkmarki\checkmarkr\checkmarkc\checkmark$\checkmarke\checkmark$\checkmark^\checkmark\\checkmarkc\checkmarki\checkmarkr\checkmarkc\checkmark$\checkmark \checkmark$\checkmark^\checkmark\\checkmarkc\checkmarki\checkmarkr\checkmarkc\checkmark$\checkmarkd\checkmark$\checkmark^\checkmark\\checkmarkc\checkmarki\checkmarkr\checkmarkc\checkmark$\checkmarke\checkmark$\checkmark^\checkmark\\checkmarkc\checkmarki\checkmarkr\checkmarkc\checkmark$\checkmark \checkmark$\checkmark^\checkmark\\checkmarkc\checkmarki\checkmarkr\checkmarkc\checkmark$\checkmarkc\checkmark$\checkmark^\checkmark\\checkmarkc\checkmarki\checkmarkr\checkmarkc\checkmark$\checkmarki\checkmark$\checkmark^\checkmark\\checkmarkc\checkmarki\checkmarkr\checkmarkc\checkmark$\checkmarks\checkmark$\checkmark^\checkmark\\checkmarkc\checkmarki\checkmarkr\checkmarkc\checkmark$\checkmarka\checkmark$\checkmark^\checkmark\\checkmarkc\checkmarki\checkmarkr\checkmarkc\checkmark$\checkmarki\checkmark$\checkmark^\checkmark\\checkmarkc\checkmarki\checkmarkr\checkmarkc\checkmark$\checkmarkl\checkmark$\checkmark^\checkmark\\checkmarkc\checkmarki\checkmarkr\checkmarkc\checkmark$\checkmarkl\checkmark$\checkmark^\checkmark\\checkmarkc\checkmarki\checkmarkr\checkmarkc\checkmark$\checkmarke\checkmark$\checkmark^\checkmark\\checkmarkc\checkmarki\checkmarkr\checkmarkc\checkmark$\checkmarkm\checkmark$\checkmark^\checkmark\\checkmarkc\checkmarki\checkmarkr\checkmarkc\checkmark$\checkmarke\checkmark$\checkmark^\checkmark\\checkmarkc\checkmarki\checkmarkr\checkmarkc\checkmark$\checkmarkn\checkmark$\checkmark^\checkmark\\checkmarkc\checkmarki\checkmarkr\checkmarkc\checkmark$\checkmarkt\checkmark$\checkmark^\checkmark\\checkmarkc\checkmarki\checkmarkr\checkmarkc\checkmark$\checkmark \checkmark$\checkmark^\checkmark\\checkmarkc\checkmarki\checkmarkr\checkmarkc\checkmark$\checkmark$\checkmark$\checkmark^\checkmark\\checkmarkc\checkmarki\checkmarkr\checkmarkc\checkmark$\checkmark\\checkmark$\checkmark^\checkmark\\checkmarkc\checkmarki\checkmarkr\checkmarkc\checkmark$\checkmarkt\checkmark$\checkmark^\checkmark\\checkmarkc\checkmarki\checkmarkr\checkmarkc\checkmark$\checkmarka\checkmark$\checkmark^\checkmark\\checkmarkc\checkmarki\checkmarkr\checkmarkc\checkmark$\checkmarku\checkmark$\checkmark^\checkmark\\checkmarkc\checkmarki\checkmarkr\checkmarkc\checkmark$\checkmark$\checkmark$\checkmark^\checkmark\\checkmarkc\checkmarki\checkmarkr\checkmarkc\checkmark$\checkmark \checkmark$\checkmark^\checkmark\\checkmarkc\checkmarki\checkmarkr\checkmarkc\checkmark$\checkmarke\checkmark$\checkmark^\checkmark\\checkmarkc\checkmarki\checkmarkr\checkmarkc\checkmark$\checkmarks\checkmark$\checkmark^\checkmark\\checkmarkc\checkmarki\checkmarkr\checkmarkc\checkmark$\checkmarkt\checkmark$\checkmark^\checkmark\\checkmarkc\checkmarki\checkmarkr\checkmarkc\checkmark$\checkmark \checkmark$\checkmark^\checkmark\\checkmarkc\checkmarki\checkmarkr\checkmarkc\checkmark$\checkmarkn\checkmark$\checkmark^\checkmark\\checkmarkc\checkmarki\checkmarkr\checkmarkc\checkmark$\checkmarké\checkmark$\checkmark^\checkmark\\checkmarkc\checkmarki\checkmarkr\checkmarkc\checkmark$\checkmarkg\checkmark$\checkmark^\checkmark\\checkmarkc\checkmarki\checkmarkr\checkmarkc\checkmark$\checkmarkl\checkmark$\checkmark^\checkmark\\checkmarkc\checkmarki\checkmarkr\checkmarkc\checkmark$\checkmarki\checkmark$\checkmark^\checkmark\\checkmarkc\checkmarki\checkmarkr\checkmarkc\checkmark$\checkmarkg\checkmark$\checkmark^\checkmark\\checkmarkc\checkmarki\checkmarkr\checkmarkc\checkmark$\checkmarke\checkmark$\checkmark^\checkmark\\checkmarkc\checkmarki\checkmarkr\checkmarkc\checkmark$\checkmarka\checkmark$\checkmark^\checkmark\\checkmarkc\checkmarki\checkmarkr\checkmarkc\checkmark$\checkmarkb\checkmark$\checkmark^\checkmark\\checkmarkc\checkmarki\checkmarkr\checkmarkc\checkmark$\checkmarkl\checkmark$\checkmark^\checkmark\\checkmarkc\checkmarki\checkmarkr\checkmarkc\checkmark$\checkmarke\checkmark$\checkmark^\checkmark\\checkmarkc\checkmarki\checkmarkr\checkmarkc\checkmark$\checkmark \checkmark$\checkmark^\checkmark\\checkmarkc\checkmarki\checkmarkr\checkmarkc\checkmark$\checkmarkp\checkmark$\checkmark^\checkmark\\checkmarkc\checkmarki\checkmarkr\checkmarkc\checkmark$\checkmarka\checkmark$\checkmark^\checkmark\\checkmarkc\checkmarki\checkmarkr\checkmarkc\checkmark$\checkmarkr\checkmark$\checkmark^\checkmark\\checkmarkc\checkmarki\checkmarkr\checkmarkc\checkmark$\checkmark \checkmark$\checkmark^\checkmark\\checkmarkc\checkmarki\checkmarkr\checkmarkc\checkmark$\checkmarkr\checkmark$\checkmark^\checkmark\\checkmarkc\checkmarki\checkmarkr\checkmarkc\checkmark$\checkmarka\checkmark$\checkmark^\checkmark\\checkmarkc\checkmarki\checkmarkr\checkmarkc\checkmark$\checkmarkp\checkmark$\checkmark^\checkmark\\checkmarkc\checkmarki\checkmarkr\checkmarkc\checkmark$\checkmarkp\checkmark$\checkmark^\checkmark\\checkmarkc\checkmarki\checkmarkr\checkmarkc\checkmark$\checkmarko\checkmark$\checkmark^\checkmark\\checkmarkc\checkmarki\checkmarkr\checkmarkc\checkmark$\checkmarkr\checkmark$\checkmark^\checkmark\\checkmarkc\checkmarki\checkmarkr\checkmarkc\checkmark$\checkmarkt\checkmark$\checkmark^\checkmark\\checkmarkc\checkmarki\checkmarkr\checkmarkc\checkmark$\checkmark \checkmark$\checkmark^\checkmark\\checkmarkc\checkmarki\checkmarkr\checkmarkc\checkmark$\checkmarkà\checkmark$\checkmark^\checkmark\\checkmarkc\checkmarki\checkmarkr\checkmarkc\checkmark$\checkmark \checkmark$\checkmark^\checkmark\\checkmarkc\checkmarki\checkmarkr\checkmarkc\checkmark$\checkmark$\checkmark$\checkmark^\checkmark\\checkmarkc\checkmarki\checkmarkr\checkmarkc\checkmark$\checkmark\\checkmark$\checkmark^\checkmark\\checkmarkc\checkmarki\checkmarkr\checkmarkc\checkmark$\checkmarks\checkmark$\checkmark^\checkmark\\checkmarkc\checkmarki\checkmarkr\checkmarkc\checkmark$\checkmarki\checkmark$\checkmark^\checkmark\\checkmarkc\checkmarki\checkmarkr\checkmarkc\checkmark$\checkmarkg\checkmark$\checkmark^\checkmark\\checkmarkc\checkmarki\checkmarkr\checkmarkc\checkmark$\checkmarkm\checkmark$\checkmark^\checkmark\\checkmarkc\checkmarki\checkmarkr\checkmarkc\checkmark$\checkmarka\checkmark$\checkmark^\checkmark\\checkmarkc\checkmarki\checkmarkr\checkmarkc\checkmark$\checkmark$\checkmark$\checkmark^\checkmark\\checkmarkc\checkmarki\checkmarkr\checkmarkc\checkmark$\checkmark \checkmark$\checkmark^\checkmark\\checkmarkc\checkmarki\checkmarkr\checkmarkc\checkmark$\checkmarkp\checkmark$\checkmark^\checkmark\\checkmarkc\checkmarki\checkmarkr\checkmarkc\checkmark$\checkmarko\checkmark$\checkmark^\checkmark\\checkmarkc\checkmarki\checkmarkr\checkmarkc\checkmark$\checkmarku\checkmark$\checkmark^\checkmark\\checkmarkc\checkmarki\checkmarkr\checkmarkc\checkmark$\checkmarkr\checkmark$\checkmark^\checkmark\\checkmarkc\checkmarki\checkmarkr\checkmarkc\checkmark$\checkmark \checkmark$\checkmark^\checkmark\\checkmarkc\checkmarki\checkmarkr\checkmarkc\checkmark$\checkmarku\checkmark$\checkmark^\checkmark\\checkmarkc\checkmarki\checkmarkr\checkmarkc\checkmark$\checkmarkn\checkmark$\checkmark^\checkmark\\checkmarkc\checkmarki\checkmarkr\checkmarkc\checkmark$\checkmarke\checkmark$\checkmark^\checkmark\\checkmarkc\checkmarki\checkmarkr\checkmarkc\checkmark$\checkmark \checkmark$\checkmark^\checkmark\\checkmarkc\checkmarki\checkmarkr\checkmarkc\checkmark$\checkmarkp\checkmark$\checkmark^\checkmark\\checkmarkc\checkmarki\checkmarkr\checkmarkc\checkmark$\checkmarko\checkmark$\checkmark^\checkmark\\checkmarkc\checkmarki\checkmarkr\checkmarkc\checkmark$\checkmarku\checkmark$\checkmark^\checkmark\\checkmarkc\checkmarki\checkmarkr\checkmarkc\checkmark$\checkmarkt\checkmark$\checkmark^\checkmark\\checkmarkc\checkmarki\checkmarkr\checkmarkc\checkmark$\checkmarkr\checkmark$\checkmark^\checkmark\\checkmarkc\checkmarki\checkmarkr\checkmarkc\checkmark$\checkmarke\checkmark$\checkmark^\checkmark\\checkmarkc\checkmarki\checkmarkr\checkmarkc\checkmark$\checkmark \checkmark$\checkmark^\checkmark\\checkmarkc\checkmarki\checkmarkr\checkmarkc\checkmark$\checkmarks\checkmark$\checkmark^\checkmark\\checkmarkc\checkmarki\checkmarkr\checkmarkc\checkmark$\checkmarkt\checkmark$\checkmark^\checkmark\\checkmarkc\checkmarki\checkmarkr\checkmarkc\checkmark$\checkmarka\checkmark$\checkmark^\checkmark\\checkmarkc\checkmarki\checkmarkr\checkmarkc\checkmark$\checkmarkn\checkmark$\checkmark^\checkmark\\checkmarkc\checkmarki\checkmarkr\checkmarkc\checkmark$\checkmarkd\checkmark$\checkmark^\checkmark\\checkmarkc\checkmarki\checkmarkr\checkmarkc\checkmark$\checkmarka\checkmark$\checkmark^\checkmark\\checkmarkc\checkmarki\checkmarkr\checkmarkc\checkmark$\checkmarkr\checkmark$\checkmark^\checkmark\\checkmarkc\checkmarki\checkmarkr\checkmarkc\checkmark$\checkmarkd\checkmark$\checkmark^\checkmark\\checkmarkc\checkmarki\checkmarkr\checkmarkc\checkmark$\checkmark.\checkmark$\checkmark^\checkmark\\checkmarkc\checkmarki\checkmarkr\checkmarkc\checkmark$\checkmark \checkmark$\checkmark^\checkmark\\checkmarkc\checkmarki\checkmarkr\checkmarkc\checkmark$\checkmarkL\checkmark$\checkmark^\checkmark\\checkmarkc\checkmarki\checkmarkr\checkmarkc\checkmark$\checkmarke\checkmark$\checkmark^\checkmark\\checkmarkc\checkmarki\checkmarkr\checkmarkc\checkmark$\checkmark \checkmark$\checkmark^\checkmark\\checkmarkc\checkmarki\checkmarkr\checkmarkc\checkmark$\checkmarkc\checkmark$\checkmark^\checkmark\\checkmarkc\checkmarki\checkmarkr\checkmarkc\checkmark$\checkmarkr\checkmark$\checkmark^\checkmark\\checkmarkc\checkmarki\checkmarkr\checkmarkc\checkmark$\checkmarki\checkmark$\checkmark^\checkmark\\checkmarkc\checkmarki\checkmarkr\checkmarkc\checkmark$\checkmarkt\checkmark$\checkmark^\checkmark\\checkmarkc\checkmarki\checkmarkr\checkmarkc\checkmark$\checkmarkè\checkmark$\checkmark^\checkmark\\checkmarkc\checkmarki\checkmarkr\checkmarkc\checkmark$\checkmarkr\checkmark$\checkmark^\checkmark\\checkmarkc\checkmarki\checkmarkr\checkmarkc\checkmark$\checkmarke\checkmark$\checkmark^\checkmark\\checkmarkc\checkmarki\checkmarkr\checkmarkc\checkmark$\checkmark \checkmark$\checkmark^\checkmark\\checkmarkc\checkmarki\checkmarkr\checkmarkc\checkmark$\checkmarkd\checkmark$\checkmark^\checkmark\\checkmarkc\checkmarki\checkmarkr\checkmarkc\checkmark$\checkmarke\checkmark$\checkmark^\checkmark\\checkmarkc\checkmarki\checkmarkr\checkmarkc\checkmark$\checkmark \checkmark$\checkmark^\checkmark\\checkmarkc\checkmarki\checkmarkr\checkmarkc\checkmark$\checkmarkV\checkmark$\checkmark^\checkmark\\checkmarkc\checkmarki\checkmarkr\checkmarkc\checkmark$\checkmarko\checkmark$\checkmark^\checkmark\\checkmarkc\checkmarki\checkmarkr\checkmarkc\checkmark$\checkmarkn\checkmark$\checkmark^\checkmark\\checkmarkc\checkmarki\checkmarkr\checkmarkc\checkmark$\checkmark \checkmark$\checkmark^\checkmark\\checkmarkc\checkmarki\checkmarkr\checkmarkc\checkmark$\checkmarkM\checkmark$\checkmark^\checkmark\\checkmarkc\checkmarki\checkmarkr\checkmarkc\checkmark$\checkmarki\checkmark$\checkmark^\checkmark\\checkmarkc\checkmarki\checkmarkr\checkmarkc\checkmark$\checkmarks\checkmark$\checkmark^\checkmark\\checkmarkc\checkmarki\checkmarkr\checkmarkc\checkmark$\checkmarke\checkmark$\checkmark^\checkmark\\checkmarkc\checkmarki\checkmarkr\checkmarkc\checkmark$\checkmarks\checkmark$\checkmark^\checkmark\\checkmarkc\checkmarki\checkmarkr\checkmarkc\checkmark$\checkmark \checkmark$\checkmark^\checkmark\\checkmarkc\checkmarki\checkmarkr\checkmarkc\checkmark$\checkmarkd\checkmark$\checkmark^\checkmark\\checkmarkc\checkmarki\checkmarkr\checkmarkc\checkmark$\checkmarko\checkmark$\checkmark^\checkmark\\checkmarkc\checkmarki\checkmarkr\checkmarkc\checkmark$\checkmarkn\checkmark$\checkmark^\checkmark\\checkmarkc\checkmarki\checkmarkr\checkmarkc\checkmark$\checkmarkn\checkmark$\checkmark^\checkmark\\checkmarkc\checkmarki\checkmarkr\checkmarkc\checkmark$\checkmarke\checkmark$\checkmark^\checkmark\\checkmarkc\checkmarki\checkmarkr\checkmarkc\checkmark$\checkmark \checkmark$\checkmark^\checkmark\\checkmarkc\checkmarki\checkmarkr\checkmarkc\checkmark$\checkmark:\checkmark$\checkmark^\checkmark\\checkmarkc\checkmarki\checkmarkr\checkmarkc\checkmark$\checkmark
\checkmark$\checkmark^\checkmark\\checkmarkc\checkmarki\checkmarkr\checkmarkc\checkmark$\checkmark\\checkmark$\checkmark^\checkmark\\checkmarkc\checkmarki\checkmarkr\checkmarkc\checkmark$\checkmarkb\checkmark$\checkmark^\checkmark\\checkmarkc\checkmarki\checkmarkr\checkmarkc\checkmark$\checkmarke\checkmark$\checkmark^\checkmark\\checkmarkc\checkmarki\checkmarkr\checkmarkc\checkmark$\checkmarkg\checkmark$\checkmark^\checkmark\\checkmarkc\checkmarki\checkmarkr\checkmarkc\checkmark$\checkmarki\checkmark$\checkmark^\checkmark\\checkmarkc\checkmarki\checkmarkr\checkmarkc\checkmark$\checkmarkn\checkmark$\checkmark^\checkmark\\checkmarkc\checkmarki\checkmarkr\checkmarkc\checkmark$\checkmark{\checkmark$\checkmark^\checkmark\\checkmarkc\checkmarki\checkmarkr\checkmarkc\checkmark$\checkmarke\checkmark$\checkmark^\checkmark\\checkmarkc\checkmarki\checkmarkr\checkmarkc\checkmark$\checkmarkq\checkmark$\checkmark^\checkmark\\checkmarkc\checkmarki\checkmarkr\checkmarkc\checkmark$\checkmarku\checkmark$\checkmark^\checkmark\\checkmarkc\checkmarki\checkmarkr\checkmarkc\checkmark$\checkmarka\checkmark$\checkmark^\checkmark\\checkmarkc\checkmarki\checkmarkr\checkmarkc\checkmark$\checkmarkt\checkmark$\checkmark^\checkmark\\checkmarkc\checkmarki\checkmarkr\checkmarkc\checkmark$\checkmarki\checkmark$\checkmark^\checkmark\\checkmarkc\checkmarki\checkmarkr\checkmarkc\checkmark$\checkmarko\checkmark$\checkmark^\checkmark\\checkmarkc\checkmarki\checkmarkr\checkmarkc\checkmark$\checkmarkn\checkmark$\checkmark^\checkmark\\checkmarkc\checkmarki\checkmarkr\checkmarkc\checkmark$\checkmark}\checkmark$\checkmark^\checkmark\\checkmarkc\checkmarki\checkmarkr\checkmarkc\checkmark$\checkmark
\checkmark$\checkmark^\checkmark\\checkmarkc\checkmarki\checkmarkr\checkmarkc\checkmark$\checkmark \checkmark$\checkmark^\checkmark\\checkmarkc\checkmarki\checkmarkr\checkmarkc\checkmark$\checkmark \checkmark$\checkmark^\checkmark\\checkmarkc\checkmarki\checkmarkr\checkmarkc\checkmark$\checkmark \checkmark$\checkmark^\checkmark\\checkmarkc\checkmarki\checkmarkr\checkmarkc\checkmark$\checkmark \checkmark$\checkmark^\checkmark\\checkmarkc\checkmarki\checkmarkr\checkmarkc\checkmark$\checkmark\\checkmark$\checkmark^\checkmark\\checkmarkc\checkmarki\checkmarkr\checkmarkc\checkmark$\checkmarks\checkmark$\checkmark^\checkmark\\checkmarkc\checkmarki\checkmarkr\checkmarkc\checkmark$\checkmarki\checkmark$\checkmark^\checkmark\\checkmarkc\checkmarki\checkmarkr\checkmarkc\checkmark$\checkmarkg\checkmark$\checkmark^\checkmark\\checkmarkc\checkmarki\checkmarkr\checkmarkc\checkmark$\checkmarkm\checkmark$\checkmark^\checkmark\\checkmarkc\checkmarki\checkmarkr\checkmarkc\checkmark$\checkmarka\checkmark$\checkmark^\checkmark\\checkmarkc\checkmarki\checkmarkr\checkmarkc\checkmark$\checkmark_\checkmark$\checkmark^\checkmark\\checkmarkc\checkmarki\checkmarkr\checkmarkc\checkmark$\checkmark{\checkmark$\checkmark^\checkmark\\checkmarkc\checkmarki\checkmarkr\checkmarkc\checkmark$\checkmarke\checkmark$\checkmark^\checkmark\\checkmarkc\checkmarki\checkmarkr\checkmarkc\checkmark$\checkmarkq\checkmark$\checkmark^\checkmark\\checkmarkc\checkmarki\checkmarkr\checkmarkc\checkmark$\checkmark}\checkmark$\checkmark^\checkmark\\checkmarkc\checkmarki\checkmarkr\checkmarkc\checkmark$\checkmark \checkmark$\checkmark^\checkmark\\checkmarkc\checkmarki\checkmarkr\checkmarkc\checkmark$\checkmark=\checkmark$\checkmark^\checkmark\\checkmarkc\checkmarki\checkmarkr\checkmarkc\checkmark$\checkmark \checkmark$\checkmark^\checkmark\\checkmarkc\checkmarki\checkmarkr\checkmarkc\checkmark$\checkmark\\checkmark$\checkmark^\checkmark\\checkmarkc\checkmarki\checkmarkr\checkmarkc\checkmark$\checkmarks\checkmark$\checkmark^\checkmark\\checkmarkc\checkmarki\checkmarkr\checkmarkc\checkmark$\checkmarkq\checkmark$\checkmark^\checkmark\\checkmarkc\checkmarki\checkmarkr\checkmarkc\checkmark$\checkmarkr\checkmark$\checkmark^\checkmark\\checkmarkc\checkmarki\checkmarkr\checkmarkc\checkmark$\checkmarkt\checkmark$\checkmark^\checkmark\\checkmarkc\checkmarki\checkmarkr\checkmarkc\checkmark$\checkmark{\checkmark$\checkmark^\checkmark\\checkmarkc\checkmarki\checkmarkr\checkmarkc\checkmark$\checkmark\\checkmark$\checkmark^\checkmark\\checkmarkc\checkmarki\checkmarkr\checkmarkc\checkmark$\checkmarks\checkmark$\checkmark^\checkmark\\checkmarkc\checkmarki\checkmarkr\checkmarkc\checkmark$\checkmarki\checkmark$\checkmark^\checkmark\\checkmarkc\checkmarki\checkmarkr\checkmarkc\checkmark$\checkmarkg\checkmark$\checkmark^\checkmark\\checkmarkc\checkmarki\checkmarkr\checkmarkc\checkmark$\checkmarkm\checkmark$\checkmark^\checkmark\\checkmarkc\checkmarki\checkmarkr\checkmarkc\checkmark$\checkmarka\checkmark$\checkmark^\checkmark\\checkmarkc\checkmarki\checkmarkr\checkmarkc\checkmark$\checkmark^\checkmark$\checkmark^\checkmark\\checkmarkc\checkmarki\checkmarkr\checkmarkc\checkmark$\checkmark2\checkmark$\checkmark^\checkmark\\checkmarkc\checkmarki\checkmarkr\checkmarkc\checkmark$\checkmark \checkmark$\checkmark^\checkmark\\checkmarkc\checkmarki\checkmarkr\checkmarkc\checkmark$\checkmark+\checkmark$\checkmark^\checkmark\\checkmarkc\checkmarki\checkmarkr\checkmarkc\checkmark$\checkmark \checkmark$\checkmark^\checkmark\\checkmarkc\checkmarki\checkmarkr\checkmarkc\checkmark$\checkmark3\checkmark$\checkmark^\checkmark\\checkmarkc\checkmarki\checkmarkr\checkmarkc\checkmark$\checkmark\\checkmark$\checkmark^\checkmark\\checkmarkc\checkmarki\checkmarkr\checkmarkc\checkmark$\checkmarkt\checkmark$\checkmark^\checkmark\\checkmarkc\checkmarki\checkmarkr\checkmarkc\checkmark$\checkmarka\checkmark$\checkmark^\checkmark\\checkmarkc\checkmarki\checkmarkr\checkmarkc\checkmark$\checkmarku\checkmark$\checkmark^\checkmark\\checkmarkc\checkmarki\checkmarkr\checkmarkc\checkmark$\checkmark^\checkmark$\checkmark^\checkmark\\checkmarkc\checkmarki\checkmarkr\checkmarkc\checkmark$\checkmark2\checkmark$\checkmark^\checkmark\\checkmarkc\checkmarki\checkmarkr\checkmarkc\checkmark$\checkmark}\checkmark$\checkmark^\checkmark\\checkmarkc\checkmarki\checkmarkr\checkmarkc\checkmark$\checkmark \checkmark$\checkmark^\checkmark\\checkmarkc\checkmarki\checkmarkr\checkmarkc\checkmark$\checkmark\\checkmark$\checkmark^\checkmark\\checkmarkc\checkmarki\checkmarkr\checkmarkc\checkmark$\checkmarka\checkmark$\checkmark^\checkmark\\checkmarkc\checkmarki\checkmarkr\checkmarkc\checkmark$\checkmarkp\checkmark$\checkmark^\checkmark\\checkmarkc\checkmarki\checkmarkr\checkmarkc\checkmark$\checkmarkp\checkmark$\checkmark^\checkmark\\checkmarkc\checkmarki\checkmarkr\checkmarkc\checkmark$\checkmarkr\checkmark$\checkmark^\checkmark\\checkmarkc\checkmarki\checkmarkr\checkmarkc\checkmark$\checkmarko\checkmark$\checkmark^\checkmark\\checkmarkc\checkmarki\checkmarkr\checkmarkc\checkmark$\checkmarkx\checkmark$\checkmark^\checkmark\\checkmarkc\checkmarki\checkmarkr\checkmarkc\checkmark$\checkmark \checkmark$\checkmark^\checkmark\\checkmarkc\checkmarki\checkmarkr\checkmarkc\checkmark$\checkmark\\checkmark$\checkmark^\checkmark\\checkmarkc\checkmarki\checkmarkr\checkmarkc\checkmark$\checkmarks\checkmark$\checkmark^\checkmark\\checkmarkc\checkmarki\checkmarkr\checkmarkc\checkmark$\checkmarki\checkmark$\checkmark^\checkmark\\checkmarkc\checkmarki\checkmarkr\checkmarkc\checkmark$\checkmarkg\checkmark$\checkmark^\checkmark\\checkmarkc\checkmarki\checkmarkr\checkmarkc\checkmark$\checkmarkm\checkmark$\checkmark^\checkmark\\checkmarkc\checkmarki\checkmarkr\checkmarkc\checkmark$\checkmarka\checkmark$\checkmark^\checkmark\\checkmarkc\checkmarki\checkmarkr\checkmarkc\checkmark$\checkmark_\checkmark$\checkmark^\checkmark\\checkmarkc\checkmarki\checkmarkr\checkmarkc\checkmark$\checkmark{\checkmark$\checkmark^\checkmark\\checkmarkc\checkmarki\checkmarkr\checkmarkc\checkmark$\checkmarkf\checkmark$\checkmark^\checkmark\\checkmarkc\checkmarki\checkmarkr\checkmarkc\checkmark$\checkmarkl\checkmark$\checkmark^\checkmark\\checkmarkc\checkmarki\checkmarkr\checkmarkc\checkmark$\checkmarke\checkmark$\checkmark^\checkmark\\checkmarkc\checkmarki\checkmarkr\checkmarkc\checkmark$\checkmarkx\checkmark$\checkmark^\checkmark\\checkmarkc\checkmarki\checkmarkr\checkmarkc\checkmark$\checkmark}\checkmark$\checkmark^\checkmark\\checkmarkc\checkmarki\checkmarkr\checkmarkc\checkmark$\checkmark \checkmark$\checkmark^\checkmark\\checkmarkc\checkmarki\checkmarkr\checkmarkc\checkmark$\checkmark=\checkmark$\checkmark^\checkmark\\checkmarkc\checkmarki\checkmarkr\checkmarkc\checkmark$\checkmark \checkmark$\checkmark^\checkmark\\checkmarkc\checkmarki\checkmarkr\checkmarkc\checkmark$\checkmark8\checkmark$\checkmark^\checkmark\\checkmarkc\checkmarki\checkmarkr\checkmarkc\checkmark$\checkmark{\checkmark$\checkmark^\checkmark\\checkmarkc\checkmarki\checkmarkr\checkmarkc\checkmark$\checkmark,\checkmark$\checkmark^\checkmark\\checkmarkc\checkmarki\checkmarkr\checkmarkc\checkmark$\checkmark}\checkmark$\checkmark^\checkmark\\checkmarkc\checkmarki\checkmarkr\checkmarkc\checkmark$\checkmark6\checkmark$\checkmark^\checkmark\\checkmarkc\checkmarki\checkmarkr\checkmarkc\checkmark$\checkmark \checkmark$\checkmark^\checkmark\\checkmarkc\checkmarki\checkmarkr\checkmarkc\checkmark$\checkmark\\checkmark$\checkmark^\checkmark\\checkmarkc\checkmarki\checkmarkr\checkmarkc\checkmark$\checkmarkt\checkmark$\checkmark^\checkmark\\checkmarkc\checkmarki\checkmarkr\checkmarkc\checkmark$\checkmarke\checkmark$\checkmark^\checkmark\\checkmarkc\checkmarki\checkmarkr\checkmarkc\checkmark$\checkmarkx\checkmark$\checkmark^\checkmark\\checkmarkc\checkmarki\checkmarkr\checkmarkc\checkmark$\checkmarkt\checkmark$\checkmark^\checkmark\\checkmarkc\checkmarki\checkmarkr\checkmarkc\checkmark$\checkmark{\checkmark$\checkmark^\checkmark\\checkmarkc\checkmarki\checkmarkr\checkmarkc\checkmark$\checkmark \checkmark$\checkmark^\checkmark\\checkmarkc\checkmarki\checkmarkr\checkmarkc\checkmark$\checkmarkM\checkmark$\checkmark^\checkmark\\checkmarkc\checkmarki\checkmarkr\checkmarkc\checkmark$\checkmarkP\checkmark$\checkmark^\checkmark\\checkmarkc\checkmarki\checkmarkr\checkmarkc\checkmark$\checkmarka\checkmark$\checkmark^\checkmark\\checkmarkc\checkmarki\checkmarkr\checkmarkc\checkmark$\checkmark}\checkmark$\checkmark^\checkmark\\checkmarkc\checkmarki\checkmarkr\checkmarkc\checkmark$\checkmark
\checkmark$\checkmark^\checkmark\\checkmarkc\checkmarki\checkmarkr\checkmarkc\checkmark$\checkmark\\checkmark$\checkmark^\checkmark\\checkmarkc\checkmarki\checkmarkr\checkmarkc\checkmark$\checkmarke\checkmark$\checkmark^\checkmark\\checkmarkc\checkmarki\checkmarkr\checkmarkc\checkmark$\checkmarkn\checkmark$\checkmark^\checkmark\\checkmarkc\checkmarki\checkmarkr\checkmarkc\checkmark$\checkmarkd\checkmark$\checkmark^\checkmark\\checkmarkc\checkmarki\checkmarkr\checkmarkc\checkmark$\checkmark{\checkmark$\checkmark^\checkmark\\checkmarkc\checkmarki\checkmarkr\checkmarkc\checkmark$\checkmarke\checkmark$\checkmark^\checkmark\\checkmarkc\checkmarki\checkmarkr\checkmarkc\checkmark$\checkmarkq\checkmark$\checkmark^\checkmark\\checkmarkc\checkmarki\checkmarkr\checkmarkc\checkmark$\checkmarku\checkmark$\checkmark^\checkmark\\checkmarkc\checkmarki\checkmarkr\checkmarkc\checkmark$\checkmarka\checkmark$\checkmark^\checkmark\\checkmarkc\checkmarki\checkmarkr\checkmarkc\checkmark$\checkmarkt\checkmark$\checkmark^\checkmark\\checkmarkc\checkmarki\checkmarkr\checkmarkc\checkmark$\checkmarki\checkmark$\checkmark^\checkmark\\checkmarkc\checkmarki\checkmarkr\checkmarkc\checkmark$\checkmarko\checkmark$\checkmark^\checkmark\\checkmarkc\checkmarki\checkmarkr\checkmarkc\checkmark$\checkmarkn\checkmark$\checkmark^\checkmark\\checkmarkc\checkmarki\checkmarkr\checkmarkc\checkmark$\checkmark}\checkmark$\checkmark^\checkmark\\checkmarkc\checkmarki\checkmarkr\checkmarkc\checkmark$\checkmark
\checkmark$\checkmark^\checkmark\\checkmarkc\checkmarki\checkmarkr\checkmarkc\checkmark$\checkmark
\checkmark$\checkmark^\checkmark\\checkmarkc\checkmarki\checkmarkr\checkmarkc\checkmark$\checkmark\\checkmark$\checkmark^\checkmark\\checkmarkc\checkmarki\checkmarkr\checkmarkc\checkmark$\checkmarkb\checkmark$\checkmark^\checkmark\\checkmarkc\checkmarki\checkmarkr\checkmarkc\checkmark$\checkmarke\checkmark$\checkmark^\checkmark\\checkmarkc\checkmarki\checkmarkr\checkmarkc\checkmark$\checkmarkg\checkmark$\checkmark^\checkmark\\checkmarkc\checkmarki\checkmarkr\checkmarkc\checkmark$\checkmarki\checkmark$\checkmark^\checkmark\\checkmarkc\checkmarki\checkmarkr\checkmarkc\checkmark$\checkmarkn\checkmark$\checkmark^\checkmark\\checkmarkc\checkmarki\checkmarkr\checkmarkc\checkmark$\checkmark{\checkmark$\checkmark^\checkmark\\checkmarkc\checkmarki\checkmarkr\checkmarkc\checkmark$\checkmarke\checkmark$\checkmark^\checkmark\\checkmarkc\checkmarki\checkmarkr\checkmarkc\checkmark$\checkmarkq\checkmark$\checkmark^\checkmark\\checkmarkc\checkmarki\checkmarkr\checkmarkc\checkmark$\checkmarku\checkmark$\checkmark^\checkmark\\checkmarkc\checkmarki\checkmarkr\checkmarkc\checkmark$\checkmarka\checkmark$\checkmark^\checkmark\\checkmarkc\checkmarki\checkmarkr\checkmarkc\checkmark$\checkmarkt\checkmark$\checkmark^\checkmark\\checkmarkc\checkmarki\checkmarkr\checkmarkc\checkmark$\checkmarki\checkmark$\checkmark^\checkmark\\checkmarkc\checkmarki\checkmarkr\checkmarkc\checkmark$\checkmarko\checkmark$\checkmark^\checkmark\\checkmarkc\checkmarki\checkmarkr\checkmarkc\checkmark$\checkmarkn\checkmark$\checkmark^\checkmark\\checkmarkc\checkmarki\checkmarkr\checkmarkc\checkmark$\checkmark}\checkmark$\checkmark^\checkmark\\checkmarkc\checkmarki\checkmarkr\checkmarkc\checkmark$\checkmark
\checkmark$\checkmark^\checkmark\\checkmarkc\checkmarki\checkmarkr\checkmarkc\checkmark$\checkmark \checkmark$\checkmark^\checkmark\\checkmarkc\checkmarki\checkmarkr\checkmarkc\checkmark$\checkmark \checkmark$\checkmark^\checkmark\\checkmarkc\checkmarki\checkmarkr\checkmarkc\checkmark$\checkmark \checkmark$\checkmark^\checkmark\\checkmarkc\checkmarki\checkmarkr\checkmarkc\checkmark$\checkmark \checkmark$\checkmark^\checkmark\\checkmarkc\checkmarki\checkmarkr\checkmarkc\checkmark$\checkmarks\checkmark$\checkmark^\checkmark\\checkmarkc\checkmarki\checkmarkr\checkmarkc\checkmark$\checkmark_\checkmark$\checkmark^\checkmark\\checkmarkc\checkmarki\checkmarkr\checkmarkc\checkmark$\checkmark{\checkmark$\checkmark^\checkmark\\checkmarkc\checkmarki\checkmarkr\checkmarkc\checkmark$\checkmarkV\checkmark$\checkmark^\checkmark\\checkmarkc\checkmarki\checkmarkr\checkmarkc\checkmark$\checkmarkM\checkmark$\checkmark^\checkmark\\checkmarkc\checkmarki\checkmarkr\checkmarkc\checkmark$\checkmark}\checkmark$\checkmark^\checkmark\\checkmarkc\checkmarki\checkmarkr\checkmarkc\checkmark$\checkmark \checkmark$\checkmark^\checkmark\\checkmarkc\checkmarki\checkmarkr\checkmarkc\checkmark$\checkmark=\checkmark$\checkmark^\checkmark\\checkmarkc\checkmarki\checkmarkr\checkmarkc\checkmark$\checkmark \checkmark$\checkmark^\checkmark\\checkmarkc\checkmarki\checkmarkr\checkmarkc\checkmark$\checkmark\\checkmark$\checkmark^\checkmark\\checkmarkc\checkmarki\checkmarkr\checkmarkc\checkmark$\checkmarkf\checkmark$\checkmark^\checkmark\\checkmarkc\checkmarki\checkmarkr\checkmarkc\checkmark$\checkmarkr\checkmark$\checkmark^\checkmark\\checkmarkc\checkmarki\checkmarkr\checkmarkc\checkmark$\checkmarka\checkmark$\checkmark^\checkmark\\checkmarkc\checkmarki\checkmarkr\checkmarkc\checkmark$\checkmarkc\checkmark$\checkmark^\checkmark\\checkmarkc\checkmarki\checkmarkr\checkmarkc\checkmark$\checkmark{\checkmark$\checkmark^\checkmark\\checkmarkc\checkmarki\checkmarkr\checkmarkc\checkmark$\checkmarkR\checkmark$\checkmark^\checkmark\\checkmarkc\checkmarki\checkmarkr\checkmarkc\checkmark$\checkmark_\checkmark$\checkmark^\checkmark\\checkmarkc\checkmarki\checkmarkr\checkmarkc\checkmark$\checkmarke\checkmark$\checkmark^\checkmark\\checkmarkc\checkmarki\checkmarkr\checkmarkc\checkmark$\checkmark}\checkmark$\checkmark^\checkmark\\checkmarkc\checkmarki\checkmarkr\checkmarkc\checkmark$\checkmark{\checkmark$\checkmark^\checkmark\\checkmarkc\checkmarki\checkmarkr\checkmarkc\checkmark$\checkmark\\checkmark$\checkmark^\checkmark\\checkmarkc\checkmarki\checkmarkr\checkmarkc\checkmark$\checkmarks\checkmark$\checkmark^\checkmark\\checkmarkc\checkmarki\checkmarkr\checkmarkc\checkmark$\checkmarki\checkmark$\checkmark^\checkmark\\checkmarkc\checkmarki\checkmarkr\checkmarkc\checkmark$\checkmarkg\checkmark$\checkmark^\checkmark\\checkmarkc\checkmarki\checkmarkr\checkmarkc\checkmark$\checkmarkm\checkmark$\checkmark^\checkmark\\checkmarkc\checkmarki\checkmarkr\checkmarkc\checkmark$\checkmarka\checkmark$\checkmark^\checkmark\\checkmarkc\checkmarki\checkmarkr\checkmarkc\checkmark$\checkmark_\checkmark$\checkmark^\checkmark\\checkmarkc\checkmarki\checkmarkr\checkmarkc\checkmark$\checkmark{\checkmark$\checkmark^\checkmark\\checkmarkc\checkmarki\checkmarkr\checkmarkc\checkmark$\checkmarke\checkmark$\checkmark^\checkmark\\checkmarkc\checkmarki\checkmarkr\checkmarkc\checkmark$\checkmarkq\checkmark$\checkmark^\checkmark\\checkmarkc\checkmarki\checkmarkr\checkmarkc\checkmark$\checkmark}\checkmark$\checkmark^\checkmark\\checkmarkc\checkmarki\checkmarkr\checkmarkc\checkmark$\checkmark}\checkmark$\checkmark^\checkmark\\checkmarkc\checkmarki\checkmarkr\checkmarkc\checkmark$\checkmark \checkmark$\checkmark^\checkmark\\checkmarkc\checkmarki\checkmarkr\checkmarkc\checkmark$\checkmark=\checkmark$\checkmark^\checkmark\\checkmarkc\checkmarki\checkmarkr\checkmarkc\checkmark$\checkmark \checkmark$\checkmark^\checkmark\\checkmarkc\checkmarki\checkmarkr\checkmarkc\checkmark$\checkmark\\checkmark$\checkmark^\checkmark\\checkmarkc\checkmarki\checkmarkr\checkmarkc\checkmark$\checkmarkf\checkmark$\checkmark^\checkmark\\checkmarkc\checkmarki\checkmarkr\checkmarkc\checkmark$\checkmarkr\checkmark$\checkmark^\checkmark\\checkmarkc\checkmarki\checkmarkr\checkmarkc\checkmark$\checkmarka\checkmark$\checkmark^\checkmark\\checkmarkc\checkmarki\checkmarkr\checkmarkc\checkmark$\checkmarkc\checkmark$\checkmark^\checkmark\\checkmarkc\checkmarki\checkmarkr\checkmarkc\checkmark$\checkmark{\checkmark$\checkmark^\checkmark\\checkmarkc\checkmarki\checkmarkr\checkmarkc\checkmark$\checkmark2\checkmark$\checkmark^\checkmark\\checkmarkc\checkmarki\checkmarkr\checkmarkc\checkmark$\checkmark7\checkmark$\checkmark^\checkmark\\checkmarkc\checkmarki\checkmarkr\checkmarkc\checkmark$\checkmark6\checkmark$\checkmark^\checkmark\\checkmarkc\checkmarki\checkmarkr\checkmarkc\checkmark$\checkmark}\checkmark$\checkmark^\checkmark\\checkmarkc\checkmarki\checkmarkr\checkmarkc\checkmark$\checkmark{\checkmark$\checkmark^\checkmark\\checkmarkc\checkmarki\checkmarkr\checkmarkc\checkmark$\checkmark8\checkmark$\checkmark^\checkmark\\checkmarkc\checkmarki\checkmarkr\checkmarkc\checkmark$\checkmark{\checkmark$\checkmark^\checkmark\\checkmarkc\checkmarki\checkmarkr\checkmarkc\checkmark$\checkmark,\checkmark$\checkmark^\checkmark\\checkmarkc\checkmarki\checkmarkr\checkmarkc\checkmark$\checkmark}\checkmark$\checkmark^\checkmark\\checkmarkc\checkmarki\checkmarkr\checkmarkc\checkmark$\checkmark6\checkmark$\checkmark^\checkmark\\checkmarkc\checkmarki\checkmarkr\checkmarkc\checkmark$\checkmark}\checkmark$\checkmark^\checkmark\\checkmarkc\checkmarki\checkmarkr\checkmarkc\checkmark$\checkmark \checkmark$\checkmark^\checkmark\\checkmarkc\checkmarki\checkmarkr\checkmarkc\checkmark$\checkmark=\checkmark$\checkmark^\checkmark\\checkmarkc\checkmarki\checkmarkr\checkmarkc\checkmark$\checkmark \checkmark$\checkmark^\checkmark\\checkmarkc\checkmarki\checkmarkr\checkmarkc\checkmark$\checkmark\\checkmark$\checkmark^\checkmark\\checkmarkc\checkmarki\checkmarkr\checkmarkc\checkmark$\checkmarkb\checkmark$\checkmark^\checkmark\\checkmarkc\checkmarki\checkmarkr\checkmarkc\checkmark$\checkmarko\checkmark$\checkmark^\checkmark\\checkmarkc\checkmarki\checkmarkr\checkmarkc\checkmark$\checkmarkx\checkmark$\checkmark^\checkmark\\checkmarkc\checkmarki\checkmarkr\checkmarkc\checkmark$\checkmarke\checkmark$\checkmark^\checkmark\\checkmarkc\checkmarki\checkmarkr\checkmarkc\checkmark$\checkmarkd\checkmark$\checkmark^\checkmark\\checkmarkc\checkmarki\checkmarkr\checkmarkc\checkmark$\checkmark{\checkmark$\checkmark^\checkmark\\checkmarkc\checkmarki\checkmarkr\checkmarkc\checkmark$\checkmark3\checkmark$\checkmark^\checkmark\\checkmarkc\checkmarki\checkmarkr\checkmarkc\checkmark$\checkmark2\checkmark$\checkmark^\checkmark\\checkmarkc\checkmarki\checkmarkr\checkmarkc\checkmark$\checkmark{\checkmark$\checkmark^\checkmark\\checkmarkc\checkmarki\checkmarkr\checkmarkc\checkmark$\checkmark,\checkmark$\checkmark^\checkmark\\checkmarkc\checkmarki\checkmarkr\checkmarkc\checkmark$\checkmark}\checkmark$\checkmark^\checkmark\\checkmarkc\checkmarki\checkmarkr\checkmarkc\checkmark$\checkmark1\checkmark$\checkmark^\checkmark\\checkmarkc\checkmarki\checkmarkr\checkmarkc\checkmark$\checkmark \checkmark$\checkmark^\checkmark\\checkmarkc\checkmarki\checkmarkr\checkmarkc\checkmark$\checkmark\\checkmark$\checkmark^\checkmark\\checkmarkc\checkmarki\checkmarkr\checkmarkc\checkmark$\checkmarkg\checkmark$\checkmark^\checkmark\\checkmarkc\checkmarki\checkmarkr\checkmarkc\checkmark$\checkmarkg\checkmark$\checkmark^\checkmark\\checkmarkc\checkmarki\checkmarkr\checkmarkc\checkmark$\checkmark \checkmark$\checkmark^\checkmark\\checkmarkc\checkmarki\checkmarkr\checkmarkc\checkmark$\checkmark1\checkmark$\checkmark^\checkmark\\checkmarkc\checkmarki\checkmarkr\checkmarkc\checkmark$\checkmark{\checkmark$\checkmark^\checkmark\\checkmarkc\checkmarki\checkmarkr\checkmarkc\checkmark$\checkmark,\checkmark$\checkmark^\checkmark\\checkmarkc\checkmarki\checkmarkr\checkmarkc\checkmark$\checkmark}\checkmark$\checkmark^\checkmark\\checkmarkc\checkmarki\checkmarkr\checkmarkc\checkmark$\checkmark5\checkmark$\checkmark^\checkmark\\checkmarkc\checkmarki\checkmarkr\checkmarkc\checkmark$\checkmark}\checkmark$\checkmark^\checkmark\\checkmarkc\checkmarki\checkmarkr\checkmarkc\checkmark$\checkmark
\checkmark$\checkmark^\checkmark\\checkmarkc\checkmarki\checkmarkr\checkmarkc\checkmark$\checkmark\\checkmark$\checkmark^\checkmark\\checkmarkc\checkmarki\checkmarkr\checkmarkc\checkmark$\checkmarke\checkmark$\checkmark^\checkmark\\checkmarkc\checkmarki\checkmarkr\checkmarkc\checkmark$\checkmarkn\checkmark$\checkmark^\checkmark\\checkmarkc\checkmarki\checkmarkr\checkmarkc\checkmark$\checkmarkd\checkmark$\checkmark^\checkmark\\checkmarkc\checkmarki\checkmarkr\checkmarkc\checkmark$\checkmark{\checkmark$\checkmark^\checkmark\\checkmarkc\checkmarki\checkmarkr\checkmarkc\checkmark$\checkmarke\checkmark$\checkmark^\checkmark\\checkmarkc\checkmarki\checkmarkr\checkmarkc\checkmark$\checkmarkq\checkmark$\checkmark^\checkmark\\checkmarkc\checkmarki\checkmarkr\checkmarkc\checkmark$\checkmarku\checkmark$\checkmark^\checkmark\\checkmarkc\checkmarki\checkmarkr\checkmarkc\checkmark$\checkmarka\checkmark$\checkmark^\checkmark\\checkmarkc\checkmarki\checkmarkr\checkmarkc\checkmark$\checkmarkt\checkmark$\checkmark^\checkmark\\checkmarkc\checkmarki\checkmarkr\checkmarkc\checkmark$\checkmarki\checkmark$\checkmark^\checkmark\\checkmarkc\checkmarki\checkmarkr\checkmarkc\checkmark$\checkmarko\checkmark$\checkmark^\checkmark\\checkmarkc\checkmarki\checkmarkr\checkmarkc\checkmark$\checkmarkn\checkmark$\checkmark^\checkmark\\checkmarkc\checkmarki\checkmarkr\checkmarkc\checkmark$\checkmark}\checkmark$\checkmark^\checkmark\\checkmarkc\checkmarki\checkmarkr\checkmarkc\checkmark$\checkmark
\checkmark$\checkmark^\checkmark\\checkmarkc\checkmarki\checkmarkr\checkmarkc\checkmark$\checkmark
\checkmark$\checkmark^\checkmark\\checkmarkc\checkmarki\checkmarkr\checkmarkc\checkmark$\checkmark\\checkmark$\checkmark^\checkmark\\checkmarkc\checkmarki\checkmarkr\checkmarkc\checkmark$\checkmarkt\checkmark$\checkmark^\checkmark\\checkmarkc\checkmarki\checkmarkr\checkmarkc\checkmark$\checkmarke\checkmark$\checkmark^\checkmark\\checkmarkc\checkmarki\checkmarkr\checkmarkc\checkmark$\checkmarkx\checkmark$\checkmark^\checkmark\\checkmarkc\checkmarki\checkmarkr\checkmarkc\checkmark$\checkmarkt\checkmark$\checkmark^\checkmark\\checkmarkc\checkmarki\checkmarkr\checkmarkc\checkmark$\checkmarkb\checkmark$\checkmark^\checkmark\\checkmarkc\checkmarki\checkmarkr\checkmarkc\checkmark$\checkmarkf\checkmark$\checkmark^\checkmark\\checkmarkc\checkmarki\checkmarkr\checkmarkc\checkmark$\checkmark{\checkmark$\checkmark^\checkmark\\checkmarkc\checkmarki\checkmarkr\checkmarkc\checkmark$\checkmarkV\checkmark$\checkmark^\checkmark\\checkmarkc\checkmarki\checkmarkr\checkmarkc\checkmark$\checkmarké\checkmark$\checkmark^\checkmark\\checkmarkc\checkmarki\checkmarkr\checkmarkc\checkmark$\checkmarkr\checkmark$\checkmark^\checkmark\\checkmarkc\checkmarki\checkmarkr\checkmarkc\checkmark$\checkmarki\checkmark$\checkmark^\checkmark\\checkmarkc\checkmarki\checkmarkr\checkmarkc\checkmark$\checkmarkf\checkmark$\checkmark^\checkmark\\checkmarkc\checkmarki\checkmarkr\checkmarkc\checkmark$\checkmarki\checkmark$\checkmark^\checkmark\\checkmarkc\checkmarki\checkmarkr\checkmarkc\checkmark$\checkmarkc\checkmark$\checkmark^\checkmark\\checkmarkc\checkmarki\checkmarkr\checkmarkc\checkmark$\checkmarka\checkmark$\checkmark^\checkmark\\checkmarkc\checkmarki\checkmarkr\checkmarkc\checkmark$\checkmarkt\checkmark$\checkmark^\checkmark\\checkmarkc\checkmarki\checkmarkr\checkmarkc\checkmark$\checkmarki\checkmark$\checkmark^\checkmark\\checkmarkc\checkmarki\checkmarkr\checkmarkc\checkmark$\checkmarko\checkmark$\checkmark^\checkmark\\checkmarkc\checkmarki\checkmarkr\checkmarkc\checkmark$\checkmarkn\checkmark$\checkmark^\checkmark\\checkmarkc\checkmarki\checkmarkr\checkmarkc\checkmark$\checkmark \checkmark$\checkmark^\checkmark\\checkmarkc\checkmarki\checkmarkr\checkmarkc\checkmark$\checkmark:\checkmark$\checkmark^\checkmark\\checkmarkc\checkmarki\checkmarkr\checkmarkc\checkmark$\checkmark}\checkmark$\checkmark^\checkmark\\checkmarkc\checkmarki\checkmarkr\checkmarkc\checkmark$\checkmark \checkmark$\checkmark^\checkmark\\checkmarkc\checkmarki\checkmarkr\checkmarkc\checkmark$\checkmarkC\checkmark$\checkmark^\checkmark\\checkmarkc\checkmarki\checkmarkr\checkmarkc\checkmark$\checkmarko\checkmark$\checkmark^\checkmark\\checkmarkc\checkmarki\checkmarkr\checkmarkc\checkmark$\checkmarke\checkmark$\checkmark^\checkmark\\checkmarkc\checkmarki\checkmarkr\checkmarkc\checkmark$\checkmarkf\checkmark$\checkmark^\checkmark\\checkmarkc\checkmarki\checkmarkr\checkmarkc\checkmark$\checkmarkf\checkmark$\checkmark^\checkmark\\checkmarkc\checkmarki\checkmarkr\checkmarkc\checkmark$\checkmarki\checkmark$\checkmark^\checkmark\\checkmarkc\checkmarki\checkmarkr\checkmarkc\checkmark$\checkmarkc\checkmark$\checkmark^\checkmark\\checkmarkc\checkmarki\checkmarkr\checkmarkc\checkmark$\checkmarki\checkmark$\checkmark^\checkmark\\checkmarkc\checkmarki\checkmarkr\checkmarkc\checkmark$\checkmarke\checkmark$\checkmark^\checkmark\\checkmarkc\checkmarki\checkmarkr\checkmarkc\checkmark$\checkmarkn\checkmark$\checkmark^\checkmark\\checkmarkc\checkmarki\checkmarkr\checkmarkc\checkmark$\checkmarkt\checkmark$\checkmark^\checkmark\\checkmarkc\checkmarki\checkmarkr\checkmarkc\checkmark$\checkmark \checkmark$\checkmark^\checkmark\\checkmarkc\checkmarki\checkmarkr\checkmarkc\checkmark$\checkmarkd\checkmark$\checkmark^\checkmark\\checkmarkc\checkmarki\checkmarkr\checkmarkc\checkmark$\checkmarke\checkmark$\checkmark^\checkmark\\checkmarkc\checkmarki\checkmarkr\checkmarkc\checkmark$\checkmark \checkmark$\checkmark^\checkmark\\checkmarkc\checkmarki\checkmarkr\checkmarkc\checkmark$\checkmarks\checkmark$\checkmark^\checkmark\\checkmarkc\checkmarki\checkmarkr\checkmarkc\checkmark$\checkmarké\checkmark$\checkmark^\checkmark\\checkmarkc\checkmarki\checkmarkr\checkmarkc\checkmark$\checkmarkc\checkmark$\checkmark^\checkmark\\checkmarkc\checkmarki\checkmarkr\checkmarkc\checkmark$\checkmarku\checkmark$\checkmark^\checkmark\\checkmarkc\checkmarki\checkmarkr\checkmarkc\checkmark$\checkmarkr\checkmark$\checkmark^\checkmark\\checkmarkc\checkmarki\checkmarkr\checkmarkc\checkmark$\checkmarki\checkmark$\checkmark^\checkmark\\checkmarkc\checkmarki\checkmarkr\checkmarkc\checkmark$\checkmarkt\checkmark$\checkmark^\checkmark\\checkmarkc\checkmarki\checkmarkr\checkmarkc\checkmark$\checkmarké\checkmark$\checkmark^\checkmark\\checkmarkc\checkmarki\checkmarkr\checkmarkc\checkmark$\checkmark \checkmark$\checkmark^\checkmark\\checkmarkc\checkmarki\checkmarkr\checkmarkc\checkmark$\checkmark$\checkmark$\checkmark^\checkmark\\checkmarkc\checkmarki\checkmarkr\checkmarkc\checkmark$\checkmark=\checkmark$\checkmark^\checkmark\\checkmarkc\checkmarki\checkmarkr\checkmarkc\checkmark$\checkmark \checkmark$\checkmark^\checkmark\\checkmarkc\checkmarki\checkmarkr\checkmarkc\checkmark$\checkmark3\checkmark$\checkmark^\checkmark\\checkmarkc\checkmarki\checkmarkr\checkmarkc\checkmark$\checkmark2\checkmark$\checkmark^\checkmark\\checkmarkc\checkmarki\checkmarkr\checkmarkc\checkmark$\checkmark{\checkmark$\checkmark^\checkmark\\checkmarkc\checkmarki\checkmarkr\checkmarkc\checkmark$\checkmark,\checkmark$\checkmark^\checkmark\\checkmarkc\checkmarki\checkmarkr\checkmarkc\checkmark$\checkmark}\checkmark$\checkmark^\checkmark\\checkmarkc\checkmarki\checkmarkr\checkmarkc\checkmark$\checkmark1\checkmark$\checkmark^\checkmark\\checkmarkc\checkmarki\checkmarkr\checkmarkc\checkmark$\checkmark \checkmark$\checkmark^\checkmark\\checkmarkc\checkmarki\checkmarkr\checkmarkc\checkmark$\checkmark\\checkmark$\checkmark^\checkmark\\checkmarkc\checkmarki\checkmarkr\checkmarkc\checkmark$\checkmarkg\checkmark$\checkmark^\checkmark\\checkmarkc\checkmarki\checkmarkr\checkmarkc\checkmark$\checkmarkg\checkmark$\checkmark^\checkmark\\checkmarkc\checkmarki\checkmarkr\checkmarkc\checkmark$\checkmark \checkmark$\checkmark^\checkmark\\checkmarkc\checkmarki\checkmarkr\checkmarkc\checkmark$\checkmark1\checkmark$\checkmark^\checkmark\\checkmarkc\checkmarki\checkmarkr\checkmarkc\checkmark$\checkmark{\checkmark$\checkmark^\checkmark\\checkmarkc\checkmarki\checkmarkr\checkmarkc\checkmark$\checkmark,\checkmark$\checkmark^\checkmark\\checkmarkc\checkmarki\checkmarkr\checkmarkc\checkmark$\checkmark}\checkmark$\checkmark^\checkmark\\checkmarkc\checkmarki\checkmarkr\checkmarkc\checkmark$\checkmark5\checkmark$\checkmark^\checkmark\\checkmarkc\checkmarki\checkmarkr\checkmarkc\checkmark$\checkmark$\checkmark$\checkmark^\checkmark\\checkmarkc\checkmarki\checkmarkr\checkmarkc\checkmark$\checkmark.\checkmark$\checkmark^\checkmark\\checkmarkc\checkmarki\checkmarkr\checkmarkc\checkmark$\checkmark \checkmark$\checkmark^\checkmark\\checkmarkc\checkmarki\checkmarkr\checkmarkc\checkmark$\checkmarkL\checkmark$\checkmark^\checkmark\\checkmarkc\checkmarki\checkmarkr\checkmarkc\checkmark$\checkmarka\checkmark$\checkmark^\checkmark\\checkmarkc\checkmarki\checkmarkr\checkmarkc\checkmark$\checkmark \checkmark$\checkmark^\checkmark\\checkmarkc\checkmarki\checkmarkr\checkmarkc\checkmark$\checkmarks\checkmark$\checkmark^\checkmark\\checkmarkc\checkmarki\checkmarkr\checkmarkc\checkmark$\checkmarkt\checkmark$\checkmark^\checkmark\\checkmarkc\checkmarki\checkmarkr\checkmarkc\checkmark$\checkmarkr\checkmark$\checkmark^\checkmark\\checkmarkc\checkmarki\checkmarkr\checkmarkc\checkmark$\checkmarku\checkmark$\checkmark^\checkmark\\checkmarkc\checkmarki\checkmarkr\checkmarkc\checkmark$\checkmarkc\checkmark$\checkmark^\checkmark\\checkmarkc\checkmarki\checkmarkr\checkmarkc\checkmark$\checkmarkt\checkmark$\checkmark^\checkmark\\checkmarkc\checkmarki\checkmarkr\checkmarkc\checkmark$\checkmarku\checkmark$\checkmark^\checkmark\\checkmarkc\checkmarki\checkmarkr\checkmarkc\checkmark$\checkmarkr\checkmark$\checkmark^\checkmark\\checkmarkc\checkmarki\checkmarkr\checkmarkc\checkmark$\checkmarke\checkmark$\checkmark^\checkmark\\checkmarkc\checkmarki\checkmarkr\checkmarkc\checkmark$\checkmark \checkmark$\checkmark^\checkmark\\checkmarkc\checkmarki\checkmarkr\checkmarkc\checkmark$\checkmarke\checkmark$\checkmark^\checkmark\\checkmarkc\checkmarki\checkmarkr\checkmarkc\checkmark$\checkmarks\checkmark$\checkmark^\checkmark\\checkmarkc\checkmarki\checkmarkr\checkmarkc\checkmark$\checkmarkt\checkmark$\checkmark^\checkmark\\checkmarkc\checkmarki\checkmarkr\checkmarkc\checkmark$\checkmark \checkmark$\checkmark^\checkmark\\checkmarkc\checkmarki\checkmarkr\checkmarkc\checkmark$\checkmark\\checkmark$\checkmark^\checkmark\\checkmarkc\checkmarki\checkmarkr\checkmarkc\checkmark$\checkmarkt\checkmark$\checkmark^\checkmark\\checkmarkc\checkmarki\checkmarkr\checkmarkc\checkmark$\checkmarke\checkmark$\checkmark^\checkmark\\checkmarkc\checkmarki\checkmarkr\checkmarkc\checkmark$\checkmarkx\checkmark$\checkmark^\checkmark\\checkmarkc\checkmarki\checkmarkr\checkmarkc\checkmark$\checkmarkt\checkmark$\checkmark^\checkmark\\checkmarkc\checkmarki\checkmarkr\checkmarkc\checkmark$\checkmarkb\checkmark$\checkmark^\checkmark\\checkmarkc\checkmarki\checkmarkr\checkmarkc\checkmark$\checkmarkf\checkmark$\checkmark^\checkmark\\checkmarkc\checkmarki\checkmarkr\checkmarkc\checkmark$\checkmark{\checkmark$\checkmark^\checkmark\\checkmarkc\checkmarki\checkmarkr\checkmarkc\checkmark$\checkmarkt\checkmark$\checkmark^\checkmark\\checkmarkc\checkmarki\checkmarkr\checkmarkc\checkmark$\checkmarkr\checkmark$\checkmark^\checkmark\\checkmarkc\checkmarki\checkmarkr\checkmarkc\checkmark$\checkmarkè\checkmark$\checkmark^\checkmark\\checkmarkc\checkmarki\checkmarkr\checkmarkc\checkmark$\checkmarks\checkmark$\checkmark^\checkmark\\checkmarkc\checkmarki\checkmarkr\checkmarkc\checkmark$\checkmark \checkmark$\checkmark^\checkmark\\checkmarkc\checkmarki\checkmarkr\checkmarkc\checkmark$\checkmarkl\checkmark$\checkmark^\checkmark\\checkmarkc\checkmarki\checkmarkr\checkmarkc\checkmark$\checkmarka\checkmark$\checkmark^\checkmark\\checkmarkc\checkmarki\checkmarkr\checkmarkc\checkmark$\checkmarkr\checkmark$\checkmark^\checkmark\\checkmarkc\checkmarki\checkmarkr\checkmarkc\checkmark$\checkmarkg\checkmark$\checkmark^\checkmark\\checkmarkc\checkmarki\checkmarkr\checkmarkc\checkmark$\checkmarke\checkmark$\checkmark^\checkmark\\checkmarkc\checkmarki\checkmarkr\checkmarkc\checkmark$\checkmarkm\checkmark$\checkmark^\checkmark\\checkmarkc\checkmarki\checkmarkr\checkmarkc\checkmark$\checkmarke\checkmark$\checkmark^\checkmark\\checkmarkc\checkmarki\checkmarkr\checkmarkc\checkmark$\checkmarkn\checkmark$\checkmark^\checkmark\\checkmarkc\checkmarki\checkmarkr\checkmarkc\checkmark$\checkmarkt\checkmark$\checkmark^\checkmark\\checkmarkc\checkmarki\checkmarkr\checkmarkc\checkmark$\checkmark \checkmark$\checkmark^\checkmark\\checkmarkc\checkmarki\checkmarkr\checkmarkc\checkmark$\checkmarks\checkmark$\checkmark^\checkmark\\checkmarkc\checkmarki\checkmarkr\checkmarkc\checkmark$\checkmarké\checkmark$\checkmark^\checkmark\\checkmarkc\checkmarki\checkmarkr\checkmarkc\checkmark$\checkmarkc\checkmark$\checkmark^\checkmark\\checkmarkc\checkmarki\checkmarkr\checkmarkc\checkmark$\checkmarku\checkmark$\checkmark^\checkmark\\checkmarkc\checkmarki\checkmarkr\checkmarkc\checkmark$\checkmarkr\checkmark$\checkmark^\checkmark\\checkmarkc\checkmarki\checkmarkr\checkmarkc\checkmark$\checkmarki\checkmark$\checkmark^\checkmark\\checkmarkc\checkmarki\checkmarkr\checkmarkc\checkmark$\checkmarks\checkmark$\checkmark^\checkmark\\checkmarkc\checkmarki\checkmarkr\checkmarkc\checkmark$\checkmarké\checkmark$\checkmark^\checkmark\\checkmarkc\checkmarki\checkmarkr\checkmarkc\checkmark$\checkmarke\checkmark$\checkmark^\checkmark\\checkmarkc\checkmarki\checkmarkr\checkmarkc\checkmark$\checkmark}\checkmark$\checkmark^\checkmark\\checkmarkc\checkmarki\checkmarkr\checkmarkc\checkmark$\checkmark.\checkmark$\checkmark^\checkmark\\checkmarkc\checkmarki\checkmarkr\checkmarkc\checkmark$\checkmark
\checkmark$\checkmark^\checkmark\\checkmarkc\checkmarki\checkmarkr\checkmarkc\checkmark$\checkmark
\checkmark$\checkmark^\checkmark\\checkmarkc\checkmarki\checkmarkr\checkmarkc\checkmark$\checkmark%\checkmark$\checkmark^\checkmark\\checkmarkc\checkmarki\checkmarkr\checkmarkc\checkmark$\checkmark \checkmark$\checkmark^\checkmark\\checkmarkc\checkmarki\checkmarkr\checkmarkc\checkmark$\checkmark=\checkmark$\checkmark^\checkmark\\checkmarkc\checkmarki\checkmarkr\checkmarkc\checkmark$\checkmark=\checkmark$\checkmark^\checkmark\\checkmarkc\checkmarki\checkmarkr\checkmarkc\checkmark$\checkmark=\checkmark$\checkmark^\checkmark\\checkmarkc\checkmarki\checkmarkr\checkmarkc\checkmark$\checkmark=\checkmark$\checkmark^\checkmark\\checkmarkc\checkmarki\checkmarkr\checkmarkc\checkmark$\checkmark=\checkmark$\checkmark^\checkmark\\checkmarkc\checkmarki\checkmarkr\checkmarkc\checkmark$\checkmark=\checkmark$\checkmark^\checkmark\\checkmarkc\checkmarki\checkmarkr\checkmarkc\checkmark$\checkmark=\checkmark$\checkmark^\checkmark\\checkmarkc\checkmarki\checkmarkr\checkmarkc\checkmark$\checkmark=\checkmark$\checkmark^\checkmark\\checkmarkc\checkmarki\checkmarkr\checkmarkc\checkmark$\checkmark=\checkmark$\checkmark^\checkmark\\checkmarkc\checkmarki\checkmarkr\checkmarkc\checkmark$\checkmark=\checkmark$\checkmark^\checkmark\\checkmarkc\checkmarki\checkmarkr\checkmarkc\checkmark$\checkmark=\checkmark$\checkmark^\checkmark\\checkmarkc\checkmarki\checkmarkr\checkmarkc\checkmark$\checkmark=\checkmark$\checkmark^\checkmark\\checkmarkc\checkmarki\checkmarkr\checkmarkc\checkmark$\checkmark=\checkmark$\checkmark^\checkmark\\checkmarkc\checkmarki\checkmarkr\checkmarkc\checkmark$\checkmark=\checkmark$\checkmark^\checkmark\\checkmarkc\checkmarki\checkmarkr\checkmarkc\checkmark$\checkmark=\checkmark$\checkmark^\checkmark\\checkmarkc\checkmarki\checkmarkr\checkmarkc\checkmark$\checkmark=\checkmark$\checkmark^\checkmark\\checkmarkc\checkmarki\checkmarkr\checkmarkc\checkmark$\checkmark=\checkmark$\checkmark^\checkmark\\checkmarkc\checkmarki\checkmarkr\checkmarkc\checkmark$\checkmark=\checkmark$\checkmark^\checkmark\\checkmarkc\checkmarki\checkmarkr\checkmarkc\checkmark$\checkmark=\checkmark$\checkmark^\checkmark\\checkmarkc\checkmarki\checkmarkr\checkmarkc\checkmark$\checkmark=\checkmark$\checkmark^\checkmark\\checkmarkc\checkmarki\checkmarkr\checkmarkc\checkmark$\checkmark=\checkmark$\checkmark^\checkmark\\checkmarkc\checkmarki\checkmarkr\checkmarkc\checkmark$\checkmark=\checkmark$\checkmark^\checkmark\\checkmarkc\checkmarki\checkmarkr\checkmarkc\checkmark$\checkmark=\checkmark$\checkmark^\checkmark\\checkmarkc\checkmarki\checkmarkr\checkmarkc\checkmark$\checkmark=\checkmark$\checkmark^\checkmark\\checkmarkc\checkmarki\checkmarkr\checkmarkc\checkmark$\checkmark=\checkmark$\checkmark^\checkmark\\checkmarkc\checkmarki\checkmarkr\checkmarkc\checkmark$\checkmark=\checkmark$\checkmark^\checkmark\\checkmarkc\checkmarki\checkmarkr\checkmarkc\checkmark$\checkmark=\checkmark$\checkmark^\checkmark\\checkmarkc\checkmarki\checkmarkr\checkmarkc\checkmark$\checkmark=\checkmark$\checkmark^\checkmark\\checkmarkc\checkmarki\checkmarkr\checkmarkc\checkmark$\checkmark=\checkmark$\checkmark^\checkmark\\checkmarkc\checkmarki\checkmarkr\checkmarkc\checkmark$\checkmark=\checkmark$\checkmark^\checkmark\\checkmarkc\checkmarki\checkmarkr\checkmarkc\checkmark$\checkmark=\checkmark$\checkmark^\checkmark\\checkmarkc\checkmarki\checkmarkr\checkmarkc\checkmark$\checkmark=\checkmark$\checkmark^\checkmark\\checkmarkc\checkmarki\checkmarkr\checkmarkc\checkmark$\checkmark=\checkmark$\checkmark^\checkmark\\checkmarkc\checkmarki\checkmarkr\checkmarkc\checkmark$\checkmark=\checkmark$\checkmark^\checkmark\\checkmarkc\checkmarki\checkmarkr\checkmarkc\checkmark$\checkmark=\checkmark$\checkmark^\checkmark\\checkmarkc\checkmarki\checkmarkr\checkmarkc\checkmark$\checkmark=\checkmark$\checkmark^\checkmark\\checkmarkc\checkmarki\checkmarkr\checkmarkc\checkmark$\checkmark=\checkmark$\checkmark^\checkmark\\checkmarkc\checkmarki\checkmarkr\checkmarkc\checkmark$\checkmark=\checkmark$\checkmark^\checkmark\\checkmarkc\checkmarki\checkmarkr\checkmarkc\checkmark$\checkmark=\checkmark$\checkmark^\checkmark\\checkmarkc\checkmarki\checkmarkr\checkmarkc\checkmark$\checkmark=\checkmark$\checkmark^\checkmark\\checkmarkc\checkmarki\checkmarkr\checkmarkc\checkmark$\checkmark=\checkmark$\checkmark^\checkmark\\checkmarkc\checkmarki\checkmarkr\checkmarkc\checkmark$\checkmark=\checkmark$\checkmark^\checkmark\\checkmarkc\checkmarki\checkmarkr\checkmarkc\checkmark$\checkmark=\checkmark$\checkmark^\checkmark\\checkmarkc\checkmarki\checkmarkr\checkmarkc\checkmark$\checkmark=\checkmark$\checkmark^\checkmark\\checkmarkc\checkmarki\checkmarkr\checkmarkc\checkmark$\checkmark=\checkmark$\checkmark^\checkmark\\checkmarkc\checkmarki\checkmarkr\checkmarkc\checkmark$\checkmark=\checkmark$\checkmark^\checkmark\\checkmarkc\checkmarki\checkmarkr\checkmarkc\checkmark$\checkmark=\checkmark$\checkmark^\checkmark\\checkmarkc\checkmarki\checkmarkr\checkmarkc\checkmark$\checkmark=\checkmark$\checkmark^\checkmark\\checkmarkc\checkmarki\checkmarkr\checkmarkc\checkmark$\checkmark=\checkmark$\checkmark^\checkmark\\checkmarkc\checkmarki\checkmarkr\checkmarkc\checkmark$\checkmark=\checkmark$\checkmark^\checkmark\\checkmarkc\checkmarki\checkmarkr\checkmarkc\checkmark$\checkmark=\checkmark$\checkmark^\checkmark\\checkmarkc\checkmarki\checkmarkr\checkmarkc\checkmark$\checkmark=\checkmark$\checkmark^\checkmark\\checkmarkc\checkmarki\checkmarkr\checkmarkc\checkmark$\checkmark=\checkmark$\checkmark^\checkmark\\checkmarkc\checkmarki\checkmarkr\checkmarkc\checkmark$\checkmark=\checkmark$\checkmark^\checkmark\\checkmarkc\checkmarki\checkmarkr\checkmarkc\checkmark$\checkmark=\checkmark$\checkmark^\checkmark\\checkmarkc\checkmarki\checkmarkr\checkmarkc\checkmark$\checkmark=\checkmark$\checkmark^\checkmark\\checkmarkc\checkmarki\checkmarkr\checkmarkc\checkmark$\checkmark=\checkmark$\checkmark^\checkmark\\checkmarkc\checkmarki\checkmarkr\checkmarkc\checkmark$\checkmark=\checkmark$\checkmark^\checkmark\\checkmarkc\checkmarki\checkmarkr\checkmarkc\checkmark$\checkmark=\checkmark$\checkmark^\checkmark\\checkmarkc\checkmarki\checkmarkr\checkmarkc\checkmark$\checkmark=\checkmark$\checkmark^\checkmark\\checkmarkc\checkmarki\checkmarkr\checkmarkc\checkmark$\checkmark=\checkmark$\checkmark^\checkmark\\checkmarkc\checkmarki\checkmarkr\checkmarkc\checkmark$\checkmark=\checkmark$\checkmark^\checkmark\\checkmarkc\checkmarki\checkmarkr\checkmarkc\checkmark$\checkmark
\checkmark$\checkmark^\checkmark\\checkmarkc\checkmarki\checkmarkr\checkmarkc\checkmark$\checkmark\\checkmark$\checkmark^\checkmark\\checkmarkc\checkmarki\checkmarkr\checkmarkc\checkmark$\checkmarks\checkmark$\checkmark^\checkmark\\checkmarkc\checkmarki\checkmarkr\checkmarkc\checkmark$\checkmarke\checkmark$\checkmark^\checkmark\\checkmarkc\checkmarki\checkmarkr\checkmarkc\checkmark$\checkmarkc\checkmark$\checkmark^\checkmark\\checkmarkc\checkmarki\checkmarkr\checkmarkc\checkmark$\checkmarkt\checkmark$\checkmark^\checkmark\\checkmarkc\checkmarki\checkmarkr\checkmarkc\checkmark$\checkmarki\checkmark$\checkmark^\checkmark\\checkmarkc\checkmarki\checkmarkr\checkmarkc\checkmark$\checkmarko\checkmark$\checkmark^\checkmark\\checkmarkc\checkmarki\checkmarkr\checkmarkc\checkmark$\checkmarkn\checkmark$\checkmark^\checkmark\\checkmarkc\checkmarki\checkmarkr\checkmarkc\checkmark$\checkmark{\checkmark$\checkmark^\checkmark\\checkmarkc\checkmarki\checkmarkr\checkmarkc\checkmark$\checkmarkD\checkmark$\checkmark^\checkmark\\checkmarkc\checkmarki\checkmarkr\checkmarkc\checkmark$\checkmarku\checkmark$\checkmark^\checkmark\\checkmarkc\checkmarki\checkmarkr\checkmarkc\checkmark$\checkmarkr\checkmark$\checkmark^\checkmark\\checkmarkc\checkmarki\checkmarkr\checkmarkc\checkmark$\checkmarké\checkmark$\checkmark^\checkmark\\checkmarkc\checkmarki\checkmarkr\checkmarkc\checkmark$\checkmarke\checkmark$\checkmark^\checkmark\\checkmarkc\checkmarki\checkmarkr\checkmarkc\checkmark$\checkmark \checkmark$\checkmark^\checkmark\\checkmarkc\checkmarki\checkmarkr\checkmarkc\checkmark$\checkmarkd\checkmark$\checkmark^\checkmark\\checkmarkc\checkmarki\checkmarkr\checkmarkc\checkmark$\checkmarke\checkmark$\checkmark^\checkmark\\checkmarkc\checkmarki\checkmarkr\checkmarkc\checkmark$\checkmark \checkmark$\checkmark^\checkmark\\checkmarkc\checkmarki\checkmarkr\checkmarkc\checkmark$\checkmarkV\checkmark$\checkmark^\checkmark\\checkmarkc\checkmarki\checkmarkr\checkmarkc\checkmark$\checkmarki\checkmark$\checkmark^\checkmark\\checkmarkc\checkmarki\checkmarkr\checkmarkc\checkmark$\checkmarke\checkmark$\checkmark^\checkmark\\checkmarkc\checkmarki\checkmarkr\checkmarkc\checkmark$\checkmark \checkmark$\checkmark^\checkmark\\checkmarkc\checkmarki\checkmarkr\checkmarkc\checkmark$\checkmarkd\checkmark$\checkmark^\checkmark\\checkmarkc\checkmarki\checkmarkr\checkmarkc\checkmark$\checkmarke\checkmark$\checkmark^\checkmark\\checkmarkc\checkmarki\checkmarkr\checkmarkc\checkmark$\checkmarks\checkmark$\checkmark^\checkmark\\checkmarkc\checkmarki\checkmarkr\checkmarkc\checkmark$\checkmark \checkmark$\checkmark^\checkmark\\checkmarkc\checkmarki\checkmarkr\checkmarkc\checkmark$\checkmarkG\checkmark$\checkmark^\checkmark\\checkmarkc\checkmarki\checkmarkr\checkmarkc\checkmark$\checkmarku\checkmark$\checkmark^\checkmark\\checkmarkc\checkmarki\checkmarkr\checkmarkc\checkmark$\checkmarki\checkmark$\checkmark^\checkmark\\checkmarkc\checkmarki\checkmarkr\checkmarkc\checkmark$\checkmarkd\checkmark$\checkmark^\checkmark\\checkmarkc\checkmarki\checkmarkr\checkmarkc\checkmark$\checkmarka\checkmark$\checkmark^\checkmark\\checkmarkc\checkmarki\checkmarkr\checkmarkc\checkmark$\checkmarkg\checkmark$\checkmark^\checkmark\\checkmarkc\checkmarki\checkmarkr\checkmarkc\checkmark$\checkmarke\checkmark$\checkmark^\checkmark\\checkmarkc\checkmarki\checkmarkr\checkmarkc\checkmark$\checkmarks\checkmark$\checkmark^\checkmark\\checkmarkc\checkmarki\checkmarkr\checkmarkc\checkmark$\checkmark \checkmark$\checkmark^\checkmark\\checkmarkc\checkmarki\checkmarkr\checkmarkc\checkmark$\checkmarkL\checkmark$\checkmark^\checkmark\\checkmarkc\checkmarki\checkmarkr\checkmarkc\checkmark$\checkmarki\checkmark$\checkmark^\checkmark\\checkmarkc\checkmarki\checkmarkr\checkmarkc\checkmark$\checkmarkn\checkmark$\checkmark^\checkmark\\checkmarkc\checkmarki\checkmarkr\checkmarkc\checkmark$\checkmarké\checkmark$\checkmark^\checkmark\\checkmarkc\checkmarki\checkmarkr\checkmarkc\checkmark$\checkmarka\checkmark$\checkmark^\checkmark\\checkmarkc\checkmarki\checkmarkr\checkmarkc\checkmark$\checkmarki\checkmark$\checkmark^\checkmark\\checkmarkc\checkmarki\checkmarkr\checkmarkc\checkmark$\checkmarkr\checkmark$\checkmark^\checkmark\\checkmarkc\checkmarki\checkmarkr\checkmarkc\checkmark$\checkmarke\checkmark$\checkmark^\checkmark\\checkmarkc\checkmarki\checkmarkr\checkmarkc\checkmark$\checkmarks\checkmark$\checkmark^\checkmark\\checkmarkc\checkmarki\checkmarkr\checkmarkc\checkmark$\checkmark \checkmark$\checkmark^\checkmark\\checkmarkc\checkmarki\checkmarkr\checkmarkc\checkmark$\checkmark(\checkmark$\checkmark^\checkmark\\checkmarkc\checkmarki\checkmarkr\checkmarkc\checkmark$\checkmarkN\checkmark$\checkmark^\checkmark\\checkmarkc\checkmarki\checkmarkr\checkmarkc\checkmark$\checkmarko\checkmark$\checkmark^\checkmark\\checkmarkc\checkmarki\checkmarkr\checkmarkc\checkmark$\checkmarkr\checkmark$\checkmark^\checkmark\\checkmarkc\checkmarki\checkmarkr\checkmarkc\checkmark$\checkmarkm\checkmark$\checkmark^\checkmark\\checkmarkc\checkmarki\checkmarkr\checkmarkc\checkmark$\checkmarke\checkmark$\checkmark^\checkmark\\checkmarkc\checkmarki\checkmarkr\checkmarkc\checkmark$\checkmark \checkmark$\checkmark^\checkmark\\checkmarkc\checkmarki\checkmarkr\checkmarkc\checkmark$\checkmarkI\checkmark$\checkmark^\checkmark\\checkmarkc\checkmarki\checkmarkr\checkmarkc\checkmark$\checkmarkS\checkmark$\checkmark^\checkmark\\checkmarkc\checkmarki\checkmarkr\checkmarkc\checkmark$\checkmarkO\checkmark$\checkmark^\checkmark\\checkmarkc\checkmarki\checkmarkr\checkmarkc\checkmark$\checkmark \checkmark$\checkmark^\checkmark\\checkmarkc\checkmarki\checkmarkr\checkmarkc\checkmark$\checkmark2\checkmark$\checkmark^\checkmark\\checkmarkc\checkmarki\checkmarkr\checkmarkc\checkmark$\checkmark8\checkmark$\checkmark^\checkmark\\checkmarkc\checkmarki\checkmarkr\checkmarkc\checkmark$\checkmark1\checkmark$\checkmark^\checkmark\\checkmarkc\checkmarki\checkmarkr\checkmarkc\checkmark$\checkmark)\checkmark$\checkmark^\checkmark\\checkmarkc\checkmarki\checkmarkr\checkmarkc\checkmark$\checkmark}\checkmark$\checkmark^\checkmark\\checkmarkc\checkmarki\checkmarkr\checkmarkc\checkmark$\checkmark
\checkmark$\checkmark^\checkmark\\checkmarkc\checkmarki\checkmarkr\checkmarkc\checkmark$\checkmark
\checkmark$\checkmark^\checkmark\\checkmarkc\checkmarki\checkmarkr\checkmarkc\checkmark$\checkmark\\checkmark$\checkmark^\checkmark\\checkmarkc\checkmarki\checkmarkr\checkmarkc\checkmark$\checkmarks\checkmark$\checkmark^\checkmark\\checkmarkc\checkmarki\checkmarkr\checkmarkc\checkmark$\checkmarku\checkmark$\checkmark^\checkmark\\checkmarkc\checkmarki\checkmarkr\checkmarkc\checkmark$\checkmarkb\checkmark$\checkmark^\checkmark\\checkmarkc\checkmarki\checkmarkr\checkmarkc\checkmark$\checkmarks\checkmark$\checkmark^\checkmark\\checkmarkc\checkmarki\checkmarkr\checkmarkc\checkmark$\checkmarke\checkmark$\checkmark^\checkmark\\checkmarkc\checkmarki\checkmarkr\checkmarkc\checkmark$\checkmarkc\checkmark$\checkmark^\checkmark\\checkmarkc\checkmarki\checkmarkr\checkmarkc\checkmark$\checkmarkt\checkmark$\checkmark^\checkmark\\checkmarkc\checkmarki\checkmarkr\checkmarkc\checkmark$\checkmarki\checkmark$\checkmark^\checkmark\\checkmarkc\checkmarki\checkmarkr\checkmarkc\checkmark$\checkmarko\checkmark$\checkmark^\checkmark\\checkmarkc\checkmarki\checkmarkr\checkmarkc\checkmark$\checkmarkn\checkmark$\checkmark^\checkmark\\checkmarkc\checkmarki\checkmarkr\checkmarkc\checkmark$\checkmark{\checkmark$\checkmark^\checkmark\\checkmarkc\checkmarki\checkmarkr\checkmarkc\checkmark$\checkmarkC\checkmark$\checkmark^\checkmark\\checkmarkc\checkmarki\checkmarkr\checkmarkc\checkmark$\checkmarkh\checkmark$\checkmark^\checkmark\\checkmarkc\checkmarki\checkmarkr\checkmarkc\checkmark$\checkmarka\checkmark$\checkmark^\checkmark\\checkmarkc\checkmarki\checkmarkr\checkmarkc\checkmark$\checkmarkr\checkmark$\checkmark^\checkmark\\checkmarkc\checkmarki\checkmarkr\checkmarkc\checkmark$\checkmarkg\checkmark$\checkmark^\checkmark\\checkmarkc\checkmarki\checkmarkr\checkmarkc\checkmark$\checkmarke\checkmark$\checkmark^\checkmark\\checkmarkc\checkmarki\checkmarkr\checkmarkc\checkmark$\checkmarks\checkmark$\checkmark^\checkmark\\checkmarkc\checkmarki\checkmarkr\checkmarkc\checkmark$\checkmark \checkmark$\checkmark^\checkmark\\checkmarkc\checkmarki\checkmarkr\checkmarkc\checkmark$\checkmarka\checkmark$\checkmark^\checkmark\\checkmarkc\checkmarki\checkmarkr\checkmarkc\checkmark$\checkmarkp\checkmark$\checkmark^\checkmark\\checkmarkc\checkmarki\checkmarkr\checkmarkc\checkmark$\checkmarkp\checkmark$\checkmark^\checkmark\\checkmarkc\checkmarki\checkmarkr\checkmarkc\checkmark$\checkmarkl\checkmark$\checkmark^\checkmark\\checkmarkc\checkmarki\checkmarkr\checkmarkc\checkmark$\checkmarki\checkmark$\checkmark^\checkmark\\checkmarkc\checkmarki\checkmarkr\checkmarkc\checkmark$\checkmarkq\checkmark$\checkmark^\checkmark\\checkmarkc\checkmarki\checkmarkr\checkmarkc\checkmark$\checkmarku\checkmark$\checkmark^\checkmark\\checkmarkc\checkmarki\checkmarkr\checkmarkc\checkmark$\checkmarké\checkmark$\checkmark^\checkmark\\checkmarkc\checkmarki\checkmarkr\checkmarkc\checkmark$\checkmarke\checkmark$\checkmark^\checkmark\\checkmarkc\checkmarki\checkmarkr\checkmarkc\checkmark$\checkmarks\checkmark$\checkmark^\checkmark\\checkmarkc\checkmarki\checkmarkr\checkmarkc\checkmark$\checkmark \checkmark$\checkmark^\checkmark\\checkmarkc\checkmarki\checkmarkr\checkmarkc\checkmark$\checkmarks\checkmark$\checkmark^\checkmark\\checkmarkc\checkmarki\checkmarkr\checkmarkc\checkmark$\checkmarku\checkmark$\checkmark^\checkmark\\checkmarkc\checkmarki\checkmarkr\checkmarkc\checkmark$\checkmarkr\checkmark$\checkmark^\checkmark\\checkmarkc\checkmarki\checkmarkr\checkmarkc\checkmark$\checkmark \checkmark$\checkmark^\checkmark\\checkmarkc\checkmarki\checkmarkr\checkmarkc\checkmark$\checkmarkl\checkmark$\checkmark^\checkmark\\checkmarkc\checkmarki\checkmarkr\checkmarkc\checkmark$\checkmarke\checkmark$\checkmark^\checkmark\\checkmarkc\checkmarki\checkmarkr\checkmarkc\checkmark$\checkmarks\checkmark$\checkmark^\checkmark\\checkmarkc\checkmarki\checkmarkr\checkmarkc\checkmark$\checkmark \checkmark$\checkmark^\checkmark\\checkmarkc\checkmarki\checkmarkr\checkmarkc\checkmark$\checkmarkp\checkmark$\checkmark^\checkmark\\checkmarkc\checkmarki\checkmarkr\checkmarkc\checkmark$\checkmarka\checkmark$\checkmark^\checkmark\\checkmarkc\checkmarki\checkmarkr\checkmarkc\checkmark$\checkmarkt\checkmark$\checkmark^\checkmark\\checkmarkc\checkmarki\checkmarkr\checkmarkc\checkmark$\checkmarki\checkmark$\checkmark^\checkmark\\checkmarkc\checkmarki\checkmarkr\checkmarkc\checkmark$\checkmarkn\checkmark$\checkmark^\checkmark\\checkmarkc\checkmarki\checkmarkr\checkmarkc\checkmark$\checkmarks\checkmark$\checkmark^\checkmark\\checkmarkc\checkmarki\checkmarkr\checkmarkc\checkmark$\checkmark \checkmark$\checkmark^\checkmark\\checkmarkc\checkmarki\checkmarkr\checkmarkc\checkmark$\checkmarkH\checkmark$\checkmark^\checkmark\\checkmarkc\checkmarki\checkmarkr\checkmarkc\checkmark$\checkmarkG\checkmark$\checkmark^\checkmark\\checkmarkc\checkmarki\checkmarkr\checkmarkc\checkmark$\checkmarkR\checkmark$\checkmark^\checkmark\\checkmarkc\checkmarki\checkmarkr\checkmarkc\checkmark$\checkmark1\checkmark$\checkmark^\checkmark\\checkmarkc\checkmarki\checkmarkr\checkmarkc\checkmark$\checkmark5\checkmark$\checkmark^\checkmark\\checkmarkc\checkmarki\checkmarkr\checkmarkc\checkmark$\checkmark}\checkmark$\checkmark^\checkmark\\checkmarkc\checkmarki\checkmarkr\checkmarkc\checkmark$\checkmark
\checkmark$\checkmark^\checkmark\\checkmarkc\checkmarki\checkmarkr\checkmarkc\checkmark$\checkmarkL\checkmark$\checkmark^\checkmark\\checkmarkc\checkmarki\checkmarkr\checkmarkc\checkmark$\checkmarka\checkmark$\checkmark^\checkmark\\checkmarkc\checkmarki\checkmarkr\checkmarkc\checkmark$\checkmark \checkmark$\checkmark^\checkmark\\checkmarkc\checkmarki\checkmarkr\checkmarkc\checkmark$\checkmarkc\checkmark$\checkmark^\checkmark\\checkmarkc\checkmarki\checkmarkr\checkmarkc\checkmark$\checkmarkh\checkmark$\checkmark^\checkmark\\checkmarkc\checkmarki\checkmarkr\checkmarkc\checkmark$\checkmarka\checkmark$\checkmark^\checkmark\\checkmarkc\checkmarki\checkmarkr\checkmarkc\checkmark$\checkmarkr\checkmark$\checkmark^\checkmark\\checkmarkc\checkmarki\checkmarkr\checkmarkc\checkmark$\checkmarkg\checkmark$\checkmark^\checkmark\\checkmarkc\checkmarki\checkmarkr\checkmarkc\checkmark$\checkmarke\checkmark$\checkmark^\checkmark\\checkmarkc\checkmarki\checkmarkr\checkmarkc\checkmark$\checkmark \checkmark$\checkmark^\checkmark\\checkmarkc\checkmarki\checkmarkr\checkmarkc\checkmark$\checkmarké\checkmark$\checkmark^\checkmark\\checkmarkc\checkmarki\checkmarkr\checkmarkc\checkmark$\checkmarkq\checkmark$\checkmark^\checkmark\\checkmarkc\checkmarki\checkmarkr\checkmarkc\checkmark$\checkmarku\checkmark$\checkmark^\checkmark\\checkmarkc\checkmarki\checkmarkr\checkmarkc\checkmark$\checkmarki\checkmark$\checkmark^\checkmark\\checkmarkc\checkmarki\checkmarkr\checkmarkc\checkmark$\checkmarkv\checkmark$\checkmark^\checkmark\\checkmarkc\checkmarki\checkmarkr\checkmarkc\checkmark$\checkmarka\checkmark$\checkmark^\checkmark\\checkmarkc\checkmarki\checkmarkr\checkmarkc\checkmark$\checkmarkl\checkmark$\checkmark^\checkmark\\checkmarkc\checkmarki\checkmarkr\checkmarkc\checkmark$\checkmarke\checkmark$\checkmark^\checkmark\\checkmarkc\checkmarki\checkmarkr\checkmarkc\checkmark$\checkmarkn\checkmark$\checkmark^\checkmark\\checkmarkc\checkmarki\checkmarkr\checkmarkc\checkmark$\checkmarkt\checkmark$\checkmark^\checkmark\\checkmarkc\checkmarki\checkmarkr\checkmarkc\checkmark$\checkmarke\checkmark$\checkmark^\checkmark\\checkmarkc\checkmarki\checkmarkr\checkmarkc\checkmark$\checkmark \checkmark$\checkmark^\checkmark\\checkmarkc\checkmarki\checkmarkr\checkmarkc\checkmark$\checkmark$\checkmark$\checkmark^\checkmark\\checkmarkc\checkmarki\checkmarkr\checkmarkc\checkmark$\checkmarkP\checkmark$\checkmark^\checkmark\\checkmarkc\checkmarki\checkmarkr\checkmarkc\checkmark$\checkmark_\checkmark$\checkmark^\checkmark\\checkmarkc\checkmarki\checkmarkr\checkmarkc\checkmark$\checkmark{\checkmark$\checkmark^\checkmark\\checkmarkc\checkmarki\checkmarkr\checkmarkc\checkmark$\checkmarke\checkmark$\checkmark^\checkmark\\checkmarkc\checkmarki\checkmarkr\checkmarkc\checkmark$\checkmarkq\checkmark$\checkmark^\checkmark\\checkmarkc\checkmarki\checkmarkr\checkmarkc\checkmark$\checkmarku\checkmark$\checkmark^\checkmark\\checkmarkc\checkmarki\checkmarkr\checkmarkc\checkmark$\checkmarki\checkmark$\checkmark^\checkmark\\checkmarkc\checkmarki\checkmarkr\checkmarkc\checkmark$\checkmarkv\checkmark$\checkmark^\checkmark\\checkmarkc\checkmarki\checkmarkr\checkmarkc\checkmark$\checkmark}\checkmark$\checkmark^\checkmark\\checkmarkc\checkmarki\checkmarkr\checkmarkc\checkmark$\checkmark$\checkmark$\checkmark^\checkmark\\checkmarkc\checkmarki\checkmarkr\checkmarkc\checkmark$\checkmark \checkmark$\checkmark^\checkmark\\checkmarkc\checkmarki\checkmarkr\checkmarkc\checkmark$\checkmarks\checkmark$\checkmark^\checkmark\\checkmarkc\checkmarki\checkmarkr\checkmarkc\checkmark$\checkmarku\checkmark$\checkmark^\checkmark\\checkmarkc\checkmarki\checkmarkr\checkmarkc\checkmark$\checkmarkr\checkmark$\checkmark^\checkmark\\checkmarkc\checkmarki\checkmarkr\checkmarkc\checkmark$\checkmark \checkmark$\checkmark^\checkmark\\checkmarkc\checkmarki\checkmarkr\checkmarkc\checkmark$\checkmarkl\checkmark$\checkmark^\checkmark\\checkmarkc\checkmarki\checkmarkr\checkmarkc\checkmark$\checkmarke\checkmark$\checkmark^\checkmark\\checkmarkc\checkmarki\checkmarkr\checkmarkc\checkmark$\checkmark \checkmark$\checkmark^\checkmark\\checkmarkc\checkmarki\checkmarkr\checkmarkc\checkmark$\checkmarkp\checkmark$\checkmark^\checkmark\\checkmarkc\checkmarki\checkmarkr\checkmarkc\checkmark$\checkmarka\checkmark$\checkmark^\checkmark\\checkmarkc\checkmarki\checkmarkr\checkmarkc\checkmark$\checkmarkt\checkmark$\checkmark^\checkmark\\checkmarkc\checkmarki\checkmarkr\checkmarkc\checkmark$\checkmarki\checkmark$\checkmark^\checkmark\\checkmarkc\checkmarki\checkmarkr\checkmarkc\checkmark$\checkmarkn\checkmark$\checkmark^\checkmark\\checkmarkc\checkmarki\checkmarkr\checkmarkc\checkmark$\checkmark \checkmark$\checkmark^\checkmark\\checkmarkc\checkmarki\checkmarkr\checkmarkc\checkmark$\checkmarkl\checkmark$\checkmark^\checkmark\\checkmarkc\checkmarki\checkmarkr\checkmarkc\checkmark$\checkmarke\checkmark$\checkmark^\checkmark\\checkmarkc\checkmarki\checkmarkr\checkmarkc\checkmark$\checkmark \checkmark$\checkmark^\checkmark\\checkmarkc\checkmarki\checkmarkr\checkmarkc\checkmark$\checkmarkp\checkmark$\checkmark^\checkmark\\checkmarkc\checkmarki\checkmarkr\checkmarkc\checkmark$\checkmarkl\checkmark$\checkmark^\checkmark\\checkmarkc\checkmarki\checkmarkr\checkmarkc\checkmark$\checkmarku\checkmark$\checkmark^\checkmark\\checkmarkc\checkmarki\checkmarkr\checkmarkc\checkmark$\checkmarks\checkmark$\checkmark^\checkmark\\checkmarkc\checkmarki\checkmarkr\checkmarkc\checkmark$\checkmark \checkmark$\checkmark^\checkmark\\checkmarkc\checkmarki\checkmarkr\checkmarkc\checkmark$\checkmarkc\checkmark$\checkmark^\checkmark\\checkmarkc\checkmarki\checkmarkr\checkmarkc\checkmark$\checkmarkh\checkmark$\checkmark^\checkmark\\checkmarkc\checkmarki\checkmarkr\checkmarkc\checkmark$\checkmarka\checkmark$\checkmark^\checkmark\\checkmarkc\checkmarki\checkmarkr\checkmarkc\checkmark$\checkmarkr\checkmark$\checkmark^\checkmark\\checkmarkc\checkmarki\checkmarkr\checkmarkc\checkmark$\checkmarkg\checkmark$\checkmark^\checkmark\\checkmarkc\checkmarki\checkmarkr\checkmarkc\checkmark$\checkmarké\checkmark$\checkmark^\checkmark\\checkmarkc\checkmarki\checkmarkr\checkmarkc\checkmark$\checkmark \checkmark$\checkmark^\checkmark\\checkmarkc\checkmarki\checkmarkr\checkmarkc\checkmark$\checkmark(\checkmark$\checkmark^\checkmark\\checkmarkc\checkmarki\checkmarkr\checkmarkc\checkmark$\checkmarkc\checkmark$\checkmark^\checkmark\\checkmarkc\checkmarki\checkmarkr\checkmarkc\checkmark$\checkmarkh\checkmark$\checkmark^\checkmark\\checkmarkc\checkmarki\checkmarkr\checkmarkc\checkmark$\checkmarka\checkmark$\checkmark^\checkmark\\checkmarkc\checkmarki\checkmarkr\checkmarkc\checkmark$\checkmarkr\checkmark$\checkmark^\checkmark\\checkmarkc\checkmarki\checkmarkr\checkmarkc\checkmark$\checkmarki\checkmark$\checkmark^\checkmark\\checkmarkc\checkmarki\checkmarkr\checkmarkc\checkmark$\checkmarko\checkmark$\checkmark^\checkmark\\checkmarkc\checkmarki\checkmarkr\checkmarkc\checkmark$\checkmarkt\checkmark$\checkmark^\checkmark\\checkmarkc\checkmarki\checkmarkr\checkmarkc\checkmark$\checkmark \checkmark$\checkmark^\checkmark\\checkmarkc\checkmarki\checkmarkr\checkmarkc\checkmark$\checkmarkY\checkmark$\checkmark^\checkmark\\checkmarkc\checkmarki\checkmarkr\checkmarkc\checkmark$\checkmark,\checkmark$\checkmark^\checkmark\\checkmarkc\checkmarki\checkmarkr\checkmarkc\checkmark$\checkmark \checkmark$\checkmark^\checkmark\\checkmarkc\checkmarki\checkmarkr\checkmarkc\checkmark$\checkmarkp\checkmark$\checkmark^\checkmark\\checkmarkc\checkmarki\checkmarkr\checkmarkc\checkmark$\checkmarko\checkmark$\checkmark^\checkmark\\checkmarkc\checkmarki\checkmarkr\checkmarkc\checkmark$\checkmarkr\checkmark$\checkmark^\checkmark\\checkmarkc\checkmarki\checkmarkr\checkmarkc\checkmark$\checkmarkt\checkmark$\checkmark^\checkmark\\checkmarkc\checkmarki\checkmarkr\checkmarkc\checkmark$\checkmarka\checkmark$\checkmark^\checkmark\\checkmarkc\checkmarki\checkmarkr\checkmarkc\checkmark$\checkmarkn\checkmark$\checkmark^\checkmark\\checkmarkc\checkmarki\checkmarkr\checkmarkc\checkmark$\checkmarkt\checkmark$\checkmark^\checkmark\\checkmarkc\checkmarki\checkmarkr\checkmarkc\checkmark$\checkmark \checkmark$\checkmark^\checkmark\\checkmarkc\checkmarki\checkmarkr\checkmarkc\checkmark$\checkmarkl\checkmark$\checkmark^\checkmark\\checkmarkc\checkmarki\checkmarkr\checkmarkc\checkmark$\checkmarka\checkmark$\checkmark^\checkmark\\checkmarkc\checkmarki\checkmarkr\checkmarkc\checkmark$\checkmark \checkmark$\checkmark^\checkmark\\checkmarkc\checkmarki\checkmarkr\checkmarkc\checkmark$\checkmarkt\checkmark$\checkmark^\checkmark\\checkmarkc\checkmarki\checkmarkr\checkmarkc\checkmark$\checkmarka\checkmark$\checkmark^\checkmark\\checkmarkc\checkmarki\checkmarkr\checkmarkc\checkmark$\checkmarkb\checkmark$\checkmark^\checkmark\\checkmarkc\checkmarki\checkmarkr\checkmarkc\checkmark$\checkmarkl\checkmark$\checkmark^\checkmark\\checkmarkc\checkmarki\checkmarkr\checkmarkc\checkmark$\checkmarke\checkmark$\checkmark^\checkmark\\checkmarkc\checkmarki\checkmarkr\checkmarkc\checkmark$\checkmark \checkmark$\checkmark^\checkmark\\checkmarkc\checkmarki\checkmarkr\checkmarkc\checkmark$\checkmarkT\checkmark$\checkmark^\checkmark\\checkmarkc\checkmarki\checkmarkr\checkmarkc\checkmark$\checkmarkr\checkmark$\checkmark^\checkmark\\checkmarkc\checkmarki\checkmarkr\checkmarkc\checkmark$\checkmarku\checkmark$\checkmark^\checkmark\\checkmarkc\checkmarki\checkmarkr\checkmarkc\checkmark$\checkmarkn\checkmark$\checkmark^\checkmark\\checkmarkc\checkmarki\checkmarkr\checkmarkc\checkmark$\checkmarkn\checkmark$\checkmark^\checkmark\\checkmarkc\checkmarki\checkmarkr\checkmarkc\checkmark$\checkmarki\checkmark$\checkmark^\checkmark\\checkmarkc\checkmarki\checkmarkr\checkmarkc\checkmark$\checkmarko\checkmark$\checkmark^\checkmark\\checkmarkc\checkmarki\checkmarkr\checkmarkc\checkmark$\checkmarkn\checkmark$\checkmark^\checkmark\\checkmarkc\checkmarki\checkmarkr\checkmarkc\checkmark$\checkmark)\checkmark$\checkmark^\checkmark\\checkmarkc\checkmarki\checkmarkr\checkmarkc\checkmark$\checkmark \checkmark$\checkmark^\checkmark\\checkmarkc\checkmarki\checkmarkr\checkmarkc\checkmark$\checkmark:\checkmark$\checkmark^\checkmark\\checkmarkc\checkmarki\checkmarkr\checkmarkc\checkmark$\checkmark
\checkmark$\checkmark^\checkmark\\checkmarkc\checkmarki\checkmarkr\checkmarkc\checkmark$\checkmark\\checkmark$\checkmark^\checkmark\\checkmarkc\checkmarki\checkmarkr\checkmarkc\checkmark$\checkmarkb\checkmark$\checkmark^\checkmark\\checkmarkc\checkmarki\checkmarkr\checkmarkc\checkmark$\checkmarke\checkmark$\checkmark^\checkmark\\checkmarkc\checkmarki\checkmarkr\checkmarkc\checkmark$\checkmarkg\checkmark$\checkmark^\checkmark\\checkmarkc\checkmarki\checkmarkr\checkmarkc\checkmark$\checkmarki\checkmark$\checkmark^\checkmark\\checkmarkc\checkmarki\checkmarkr\checkmarkc\checkmark$\checkmarkn\checkmark$\checkmark^\checkmark\\checkmarkc\checkmarki\checkmarkr\checkmarkc\checkmark$\checkmark{\checkmark$\checkmark^\checkmark\\checkmarkc\checkmarki\checkmarkr\checkmarkc\checkmark$\checkmarki\checkmark$\checkmark^\checkmark\\checkmarkc\checkmarki\checkmarkr\checkmarkc\checkmark$\checkmarkt\checkmark$\checkmark^\checkmark\\checkmarkc\checkmarki\checkmarkr\checkmarkc\checkmark$\checkmarke\checkmark$\checkmark^\checkmark\\checkmarkc\checkmarki\checkmarkr\checkmarkc\checkmark$\checkmarkm\checkmark$\checkmark^\checkmark\\checkmarkc\checkmarki\checkmarkr\checkmarkc\checkmark$\checkmarki\checkmark$\checkmark^\checkmark\\checkmarkc\checkmarki\checkmarkr\checkmarkc\checkmark$\checkmarkz\checkmark$\checkmark^\checkmark\\checkmarkc\checkmarki\checkmarkr\checkmarkc\checkmark$\checkmarke\checkmark$\checkmark^\checkmark\\checkmarkc\checkmarki\checkmarkr\checkmarkc\checkmark$\checkmark}\checkmark$\checkmark^\checkmark\\checkmarkc\checkmarki\checkmarkr\checkmarkc\checkmark$\checkmark
\checkmark$\checkmark^\checkmark\\checkmarkc\checkmarki\checkmarkr\checkmarkc\checkmark$\checkmark \checkmark$\checkmark^\checkmark\\checkmarkc\checkmarki\checkmarkr\checkmarkc\checkmark$\checkmark \checkmark$\checkmark^\checkmark\\checkmarkc\checkmarki\checkmarkr\checkmarkc\checkmark$\checkmark \checkmark$\checkmark^\checkmark\\checkmarkc\checkmarki\checkmarkr\checkmarkc\checkmark$\checkmark \checkmark$\checkmark^\checkmark\\checkmarkc\checkmarki\checkmarkr\checkmarkc\checkmark$\checkmark\\checkmark$\checkmark^\checkmark\\checkmarkc\checkmarki\checkmarkr\checkmarkc\checkmark$\checkmarki\checkmark$\checkmark^\checkmark\\checkmarkc\checkmarki\checkmarkr\checkmarkc\checkmark$\checkmarkt\checkmark$\checkmark^\checkmark\\checkmarkc\checkmarki\checkmarkr\checkmarkc\checkmark$\checkmarke\checkmark$\checkmark^\checkmark\\checkmarkc\checkmarki\checkmarkr\checkmarkc\checkmark$\checkmarkm\checkmark$\checkmark^\checkmark\\checkmarkc\checkmarki\checkmarkr\checkmarkc\checkmark$\checkmark \checkmark$\checkmark^\checkmark\\checkmarkc\checkmarki\checkmarkr\checkmarkc\checkmark$\checkmarkC\checkmark$\checkmark^\checkmark\\checkmarkc\checkmarki\checkmarkr\checkmarkc\checkmark$\checkmarkh\checkmark$\checkmark^\checkmark\\checkmarkc\checkmarki\checkmarkr\checkmarkc\checkmark$\checkmarka\checkmark$\checkmark^\checkmark\\checkmarkc\checkmarki\checkmarkr\checkmarkc\checkmark$\checkmarkr\checkmark$\checkmark^\checkmark\\checkmarkc\checkmarki\checkmarkr\checkmarkc\checkmark$\checkmarkg\checkmark$\checkmark^\checkmark\\checkmarkc\checkmarki\checkmarkr\checkmarkc\checkmark$\checkmarke\checkmark$\checkmark^\checkmark\\checkmarkc\checkmarki\checkmarkr\checkmarkc\checkmark$\checkmark \checkmark$\checkmark^\checkmark\\checkmarkc\checkmarki\checkmarkr\checkmarkc\checkmark$\checkmarks\checkmark$\checkmark^\checkmark\\checkmarkc\checkmarki\checkmarkr\checkmarkc\checkmark$\checkmarkt\checkmark$\checkmark^\checkmark\\checkmarkc\checkmarki\checkmarkr\checkmarkc\checkmark$\checkmarka\checkmark$\checkmark^\checkmark\\checkmarkc\checkmarki\checkmarkr\checkmarkc\checkmark$\checkmarkt\checkmark$\checkmark^\checkmark\\checkmarkc\checkmarki\checkmarkr\checkmarkc\checkmark$\checkmarki\checkmark$\checkmark^\checkmark\\checkmarkc\checkmarki\checkmarkr\checkmarkc\checkmark$\checkmarkq\checkmark$\checkmark^\checkmark\\checkmarkc\checkmarki\checkmarkr\checkmarkc\checkmark$\checkmarku\checkmark$\checkmark^\checkmark\\checkmarkc\checkmarki\checkmarkr\checkmarkc\checkmark$\checkmarke\checkmark$\checkmark^\checkmark\\checkmarkc\checkmarki\checkmarkr\checkmarkc\checkmark$\checkmark \checkmark$\checkmark^\checkmark\\checkmarkc\checkmarki\checkmarkr\checkmarkc\checkmark$\checkmark(\checkmark$\checkmark^\checkmark\\checkmarkc\checkmarki\checkmarkr\checkmarkc\checkmark$\checkmarkp\checkmark$\checkmark^\checkmark\\checkmarkc\checkmarki\checkmarkr\checkmarkc\checkmark$\checkmarko\checkmark$\checkmark^\checkmark\\checkmarkc\checkmarki\checkmarkr\checkmarkc\checkmark$\checkmarki\checkmark$\checkmark^\checkmark\\checkmarkc\checkmarki\checkmarkr\checkmarkc\checkmark$\checkmarkd\checkmark$\checkmark^\checkmark\\checkmarkc\checkmarki\checkmarkr\checkmarkc\checkmark$\checkmarks\checkmark$\checkmark^\checkmark\\checkmarkc\checkmarki\checkmarkr\checkmarkc\checkmark$\checkmark \checkmark$\checkmark^\checkmark\\checkmarkc\checkmarki\checkmarkr\checkmarkc\checkmark$\checkmarkt\checkmark$\checkmark^\checkmark\\checkmarkc\checkmarki\checkmarkr\checkmarkc\checkmark$\checkmarka\checkmark$\checkmark^\checkmark\\checkmarkc\checkmarki\checkmarkr\checkmarkc\checkmark$\checkmarkb\checkmark$\checkmark^\checkmark\\checkmarkc\checkmarki\checkmarkr\checkmarkc\checkmark$\checkmarkl\checkmark$\checkmark^\checkmark\\checkmarkc\checkmarki\checkmarkr\checkmarkc\checkmark$\checkmarke\checkmark$\checkmark^\checkmark\\checkmarkc\checkmarki\checkmarkr\checkmarkc\checkmark$\checkmark \checkmark$\checkmark^\checkmark\\checkmarkc\checkmarki\checkmarkr\checkmarkc\checkmark$\checkmark+\checkmark$\checkmark^\checkmark\\checkmarkc\checkmarki\checkmarkr\checkmarkc\checkmark$\checkmark \checkmark$\checkmark^\checkmark\\checkmarkc\checkmarki\checkmarkr\checkmarkc\checkmark$\checkmarkp\checkmark$\checkmark^\checkmark\\checkmarkc\checkmarki\checkmarkr\checkmarkc\checkmark$\checkmarki\checkmark$\checkmark^\checkmark\\checkmarkc\checkmarki\checkmarkr\checkmarkc\checkmark$\checkmarkè\checkmark$\checkmark^\checkmark\\checkmarkc\checkmarki\checkmarkr\checkmarkc\checkmark$\checkmarkc\checkmark$\checkmark^\checkmark\\checkmarkc\checkmarki\checkmarkr\checkmarkc\checkmark$\checkmarke\checkmark$\checkmark^\checkmark\\checkmarkc\checkmarki\checkmarkr\checkmarkc\checkmark$\checkmark)\checkmark$\checkmark^\checkmark\\checkmarkc\checkmarki\checkmarkr\checkmarkc\checkmark$\checkmark \checkmark$\checkmark^\checkmark\\checkmarkc\checkmarki\checkmarkr\checkmarkc\checkmark$\checkmark:\checkmark$\checkmark^\checkmark\\checkmarkc\checkmarki\checkmarkr\checkmarkc\checkmark$\checkmark \checkmark$\checkmark^\checkmark\\checkmarkc\checkmarki\checkmarkr\checkmarkc\checkmark$\checkmark$\checkmark$\checkmark^\checkmark\\checkmarkc\checkmarki\checkmarkr\checkmarkc\checkmark$\checkmarkP\checkmark$\checkmark^\checkmark\\checkmarkc\checkmarki\checkmarkr\checkmarkc\checkmark$\checkmark_\checkmark$\checkmark^\checkmark\\checkmarkc\checkmarki\checkmarkr\checkmarkc\checkmark$\checkmark{\checkmark$\checkmark^\checkmark\\checkmarkc\checkmarki\checkmarkr\checkmarkc\checkmark$\checkmarks\checkmark$\checkmark^\checkmark\\checkmarkc\checkmarki\checkmarkr\checkmarkc\checkmark$\checkmarkt\checkmark$\checkmark^\checkmark\\checkmarkc\checkmarki\checkmarkr\checkmarkc\checkmark$\checkmarka\checkmark$\checkmark^\checkmark\\checkmarkc\checkmarki\checkmarkr\checkmarkc\checkmark$\checkmarkt\checkmark$\checkmark^\checkmark\\checkmarkc\checkmarki\checkmarkr\checkmarkc\checkmark$\checkmark}\checkmark$\checkmark^\checkmark\\checkmarkc\checkmarki\checkmarkr\checkmarkc\checkmark$\checkmark \checkmark$\checkmark^\checkmark\\checkmarkc\checkmarki\checkmarkr\checkmarkc\checkmark$\checkmark=\checkmark$\checkmark^\checkmark\\checkmarkc\checkmarki\checkmarkr\checkmarkc\checkmark$\checkmark \checkmark$\checkmark^\checkmark\\checkmarkc\checkmarki\checkmarkr\checkmarkc\checkmark$\checkmark(\checkmark$\checkmark^\checkmark\\checkmarkc\checkmarki\checkmarkr\checkmarkc\checkmark$\checkmarkM\checkmark$\checkmark^\checkmark\\checkmarkc\checkmarki\checkmarkr\checkmarkc\checkmark$\checkmark_\checkmark$\checkmark^\checkmark\\checkmarkc\checkmarki\checkmarkr\checkmarkc\checkmark$\checkmark{\checkmark$\checkmark^\checkmark\\checkmarkc\checkmarki\checkmarkr\checkmarkc\checkmark$\checkmarkt\checkmark$\checkmark^\checkmark\\checkmarkc\checkmarki\checkmarkr\checkmarkc\checkmark$\checkmarka\checkmark$\checkmark^\checkmark\\checkmarkc\checkmarki\checkmarkr\checkmarkc\checkmark$\checkmarkb\checkmark$\checkmark^\checkmark\\checkmarkc\checkmarki\checkmarkr\checkmarkc\checkmark$\checkmarkl\checkmark$\checkmark^\checkmark\\checkmarkc\checkmarki\checkmarkr\checkmarkc\checkmark$\checkmarke\checkmark$\checkmark^\checkmark\\checkmarkc\checkmarki\checkmarkr\checkmarkc\checkmark$\checkmark}\checkmark$\checkmark^\checkmark\\checkmarkc\checkmarki\checkmarkr\checkmarkc\checkmark$\checkmark \checkmark$\checkmark^\checkmark\\checkmarkc\checkmarki\checkmarkr\checkmarkc\checkmark$\checkmark+\checkmark$\checkmark^\checkmark\\checkmarkc\checkmarki\checkmarkr\checkmarkc\checkmark$\checkmark \checkmark$\checkmark^\checkmark\\checkmarkc\checkmarki\checkmarkr\checkmarkc\checkmark$\checkmarkM\checkmark$\checkmark^\checkmark\\checkmarkc\checkmarki\checkmarkr\checkmarkc\checkmark$\checkmark_\checkmark$\checkmark^\checkmark\\checkmarkc\checkmarki\checkmarkr\checkmarkc\checkmark$\checkmark{\checkmark$\checkmark^\checkmark\\checkmarkc\checkmarki\checkmarkr\checkmarkc\checkmark$\checkmarkp\checkmark$\checkmark^\checkmark\\checkmarkc\checkmarki\checkmarkr\checkmarkc\checkmark$\checkmarki\checkmark$\checkmark^\checkmark\\checkmarkc\checkmarki\checkmarkr\checkmarkc\checkmark$\checkmarke\checkmark$\checkmark^\checkmark\\checkmarkc\checkmarki\checkmarkr\checkmarkc\checkmark$\checkmarkc\checkmark$\checkmark^\checkmark\\checkmarkc\checkmarki\checkmarkr\checkmarkc\checkmark$\checkmarke\checkmark$\checkmark^\checkmark\\checkmarkc\checkmarki\checkmarkr\checkmarkc\checkmark$\checkmark}\checkmark$\checkmark^\checkmark\\checkmarkc\checkmarki\checkmarkr\checkmarkc\checkmark$\checkmark)\checkmark$\checkmark^\checkmark\\checkmarkc\checkmarki\checkmarkr\checkmarkc\checkmark$\checkmark \checkmark$\checkmark^\checkmark\\checkmarkc\checkmarki\checkmarkr\checkmarkc\checkmark$\checkmark\\checkmark$\checkmark^\checkmark\\checkmarkc\checkmarki\checkmarkr\checkmarkc\checkmark$\checkmarkt\checkmark$\checkmark^\checkmark\\checkmarkc\checkmarki\checkmarkr\checkmarkc\checkmark$\checkmarki\checkmark$\checkmark^\checkmark\\checkmarkc\checkmarki\checkmarkr\checkmarkc\checkmark$\checkmarkm\checkmark$\checkmark^\checkmark\\checkmarkc\checkmarki\checkmarkr\checkmarkc\checkmark$\checkmarke\checkmark$\checkmark^\checkmark\\checkmarkc\checkmarki\checkmarkr\checkmarkc\checkmark$\checkmarks\checkmark$\checkmark^\checkmark\\checkmarkc\checkmarki\checkmarkr\checkmarkc\checkmark$\checkmark \checkmark$\checkmark^\checkmark\\checkmarkc\checkmarki\checkmarkr\checkmarkc\checkmark$\checkmarkg\checkmark$\checkmark^\checkmark\\checkmarkc\checkmarki\checkmarkr\checkmarkc\checkmark$\checkmark \checkmark$\checkmark^\checkmark\\checkmarkc\checkmarki\checkmarkr\checkmarkc\checkmark$\checkmark=\checkmark$\checkmark^\checkmark\\checkmarkc\checkmarki\checkmarkr\checkmarkc\checkmark$\checkmark \checkmark$\checkmark^\checkmark\\checkmarkc\checkmarki\checkmarkr\checkmarkc\checkmark$\checkmark(\checkmark$\checkmark^\checkmark\\checkmarkc\checkmarki\checkmarkr\checkmarkc\checkmark$\checkmark3\checkmark$\checkmark^\checkmark\\checkmarkc\checkmarki\checkmarkr\checkmarkc\checkmark$\checkmark \checkmark$\checkmark^\checkmark\\checkmarkc\checkmarki\checkmarkr\checkmarkc\checkmark$\checkmark+\checkmark$\checkmark^\checkmark\\checkmarkc\checkmarki\checkmarkr\checkmarkc\checkmark$\checkmark \checkmark$\checkmark^\checkmark\\checkmarkc\checkmarki\checkmarkr\checkmarkc\checkmark$\checkmark1\checkmark$\checkmark^\checkmark\\checkmarkc\checkmarki\checkmarkr\checkmarkc\checkmark$\checkmark)\checkmark$\checkmark^\checkmark\\checkmarkc\checkmarki\checkmarkr\checkmarkc\checkmark$\checkmark \checkmark$\checkmark^\checkmark\\checkmarkc\checkmarki\checkmarkr\checkmarkc\checkmark$\checkmark\\checkmark$\checkmark^\checkmark\\checkmarkc\checkmarki\checkmarkr\checkmarkc\checkmark$\checkmarkt\checkmark$\checkmark^\checkmark\\checkmarkc\checkmarki\checkmarkr\checkmarkc\checkmark$\checkmarki\checkmark$\checkmark^\checkmark\\checkmarkc\checkmarki\checkmarkr\checkmarkc\checkmark$\checkmarkm\checkmark$\checkmark^\checkmark\\checkmarkc\checkmarki\checkmarkr\checkmarkc\checkmark$\checkmarke\checkmark$\checkmark^\checkmark\\checkmarkc\checkmarki\checkmarkr\checkmarkc\checkmark$\checkmarks\checkmark$\checkmark^\checkmark\\checkmarkc\checkmarki\checkmarkr\checkmarkc\checkmark$\checkmark \checkmark$\checkmark^\checkmark\\checkmarkc\checkmarki\checkmarkr\checkmarkc\checkmark$\checkmark9\checkmark$\checkmark^\checkmark\\checkmarkc\checkmarki\checkmarkr\checkmarkc\checkmark$\checkmark{\checkmark$\checkmark^\checkmark\\checkmarkc\checkmarki\checkmarkr\checkmarkc\checkmark$\checkmark,\checkmark$\checkmark^\checkmark\\checkmarkc\checkmarki\checkmarkr\checkmarkc\checkmark$\checkmark}\checkmark$\checkmark^\checkmark\\checkmarkc\checkmarki\checkmarkr\checkmarkc\checkmark$\checkmark8\checkmark$\checkmark^\checkmark\\checkmarkc\checkmarki\checkmarkr\checkmarkc\checkmark$\checkmark1\checkmark$\checkmark^\checkmark\\checkmarkc\checkmarki\checkmarkr\checkmarkc\checkmark$\checkmark \checkmark$\checkmark^\checkmark\\checkmarkc\checkmarki\checkmarkr\checkmarkc\checkmark$\checkmark=\checkmark$\checkmark^\checkmark\\checkmarkc\checkmarki\checkmarkr\checkmarkc\checkmark$\checkmark \checkmark$\checkmark^\checkmark\\checkmarkc\checkmarki\checkmarkr\checkmarkc\checkmark$\checkmark3\checkmark$\checkmark^\checkmark\\checkmarkc\checkmarki\checkmarkr\checkmarkc\checkmark$\checkmark9\checkmark$\checkmark^\checkmark\\checkmarkc\checkmarki\checkmarkr\checkmarkc\checkmark$\checkmark{\checkmark$\checkmark^\checkmark\\checkmarkc\checkmarki\checkmarkr\checkmarkc\checkmark$\checkmark,\checkmark$\checkmark^\checkmark\\checkmarkc\checkmarki\checkmarkr\checkmarkc\checkmark$\checkmark}\checkmark$\checkmark^\checkmark\\checkmarkc\checkmarki\checkmarkr\checkmarkc\checkmark$\checkmark2\checkmark$\checkmark^\checkmark\\checkmarkc\checkmarki\checkmarkr\checkmarkc\checkmark$\checkmark4\checkmark$\checkmark^\checkmark\\checkmarkc\checkmarki\checkmarkr\checkmarkc\checkmark$\checkmark$\checkmark$\checkmark^\checkmark\\checkmarkc\checkmarki\checkmarkr\checkmarkc\checkmark$\checkmark \checkmark$\checkmark^\checkmark\\checkmarkc\checkmarki\checkmarkr\checkmarkc\checkmark$\checkmarkN\checkmark$\checkmark^\checkmark\\checkmarkc\checkmarki\checkmarkr\checkmarkc\checkmark$\checkmark
\checkmark$\checkmark^\checkmark\\checkmarkc\checkmarki\checkmarkr\checkmarkc\checkmark$\checkmark \checkmark$\checkmark^\checkmark\\checkmarkc\checkmarki\checkmarkr\checkmarkc\checkmark$\checkmark \checkmark$\checkmark^\checkmark\\checkmarkc\checkmarki\checkmarkr\checkmarkc\checkmark$\checkmark \checkmark$\checkmark^\checkmark\\checkmarkc\checkmarki\checkmarkr\checkmarkc\checkmark$\checkmark \checkmark$\checkmark^\checkmark\\checkmarkc\checkmarki\checkmarkr\checkmarkc\checkmark$\checkmark\\checkmark$\checkmark^\checkmark\\checkmarkc\checkmarki\checkmarkr\checkmarkc\checkmark$\checkmarki\checkmark$\checkmark^\checkmark\\checkmarkc\checkmarki\checkmarkr\checkmarkc\checkmark$\checkmarkt\checkmark$\checkmark^\checkmark\\checkmarkc\checkmarki\checkmarkr\checkmarkc\checkmark$\checkmarke\checkmark$\checkmark^\checkmark\\checkmarkc\checkmarki\checkmarkr\checkmarkc\checkmark$\checkmarkm\checkmark$\checkmark^\checkmark\\checkmarkc\checkmarki\checkmarkr\checkmarkc\checkmark$\checkmark \checkmark$\checkmark^\checkmark\\checkmarkc\checkmarki\checkmarkr\checkmarkc\checkmark$\checkmarkC\checkmark$\checkmark^\checkmark\\checkmarkc\checkmarki\checkmarkr\checkmarkc\checkmark$\checkmarkh\checkmark$\checkmark^\checkmark\\checkmarkc\checkmarki\checkmarkr\checkmarkc\checkmark$\checkmarka\checkmark$\checkmark^\checkmark\\checkmarkc\checkmarki\checkmarkr\checkmarkc\checkmark$\checkmarkr\checkmark$\checkmark^\checkmark\\checkmarkc\checkmarki\checkmarkr\checkmarkc\checkmark$\checkmarkg\checkmark$\checkmark^\checkmark\\checkmarkc\checkmarki\checkmarkr\checkmarkc\checkmark$\checkmarke\checkmark$\checkmark^\checkmark\\checkmarkc\checkmarki\checkmarkr\checkmarkc\checkmark$\checkmark \checkmark$\checkmark^\checkmark\\checkmarkc\checkmarki\checkmarkr\checkmarkc\checkmark$\checkmarkd\checkmark$\checkmark^\checkmark\\checkmarkc\checkmarki\checkmarkr\checkmarkc\checkmark$\checkmarky\checkmark$\checkmark^\checkmark\\checkmarkc\checkmarki\checkmarkr\checkmarkc\checkmark$\checkmarkn\checkmark$\checkmark^\checkmark\\checkmarkc\checkmarki\checkmarkr\checkmarkc\checkmark$\checkmarka\checkmark$\checkmark^\checkmark\\checkmarkc\checkmarki\checkmarkr\checkmarkc\checkmark$\checkmarkm\checkmark$\checkmark^\checkmark\\checkmarkc\checkmarki\checkmarkr\checkmarkc\checkmark$\checkmarki\checkmark$\checkmark^\checkmark\\checkmarkc\checkmarki\checkmarkr\checkmarkc\checkmark$\checkmarkq\checkmark$\checkmark^\checkmark\\checkmarkc\checkmarki\checkmarkr\checkmarkc\checkmark$\checkmarku\checkmark$\checkmark^\checkmark\\checkmarkc\checkmarki\checkmarkr\checkmarkc\checkmark$\checkmarke\checkmark$\checkmark^\checkmark\\checkmarkc\checkmarki\checkmarkr\checkmarkc\checkmark$\checkmark \checkmark$\checkmark^\checkmark\\checkmarkc\checkmarki\checkmarkr\checkmarkc\checkmark$\checkmark(\checkmark$\checkmark^\checkmark\\checkmarkc\checkmarki\checkmarkr\checkmarkc\checkmark$\checkmarkf\checkmark$\checkmark^\checkmark\\checkmarkc\checkmarki\checkmarkr\checkmarkc\checkmark$\checkmarko\checkmark$\checkmark^\checkmark\\checkmarkc\checkmarki\checkmarkr\checkmarkc\checkmark$\checkmarkr\checkmark$\checkmark^\checkmark\\checkmarkc\checkmarki\checkmarkr\checkmarkc\checkmark$\checkmarkc\checkmark$\checkmark^\checkmark\\checkmarkc\checkmarki\checkmarkr\checkmarkc\checkmark$\checkmarke\checkmark$\checkmark^\checkmark\\checkmarkc\checkmarki\checkmarkr\checkmarkc\checkmark$\checkmarks\checkmark$\checkmark^\checkmark\\checkmarkc\checkmarki\checkmarkr\checkmarkc\checkmark$\checkmark \checkmark$\checkmark^\checkmark\\checkmarkc\checkmarki\checkmarkr\checkmarkc\checkmark$\checkmarkd\checkmark$\checkmark^\checkmark\\checkmarkc\checkmarki\checkmarkr\checkmarkc\checkmark$\checkmarke\checkmark$\checkmark^\checkmark\\checkmarkc\checkmarki\checkmarkr\checkmarkc\checkmark$\checkmark \checkmark$\checkmark^\checkmark\\checkmarkc\checkmarki\checkmarkr\checkmarkc\checkmark$\checkmarkc\checkmark$\checkmark^\checkmark\\checkmarkc\checkmarki\checkmarkr\checkmarkc\checkmark$\checkmarko\checkmark$\checkmark^\checkmark\\checkmarkc\checkmarki\checkmarkr\checkmarkc\checkmark$\checkmarku\checkmark$\checkmark^\checkmark\\checkmarkc\checkmarki\checkmarkr\checkmarkc\checkmark$\checkmarkp\checkmark$\checkmark^\checkmark\\checkmarkc\checkmarki\checkmarkr\checkmarkc\checkmark$\checkmarke\checkmark$\checkmark^\checkmark\\checkmarkc\checkmarki\checkmarkr\checkmarkc\checkmark$\checkmark)\checkmark$\checkmark^\checkmark\\checkmarkc\checkmarki\checkmarkr\checkmarkc\checkmark$\checkmark \checkmark$\checkmark^\checkmark\\checkmarkc\checkmarki\checkmarkr\checkmarkc\checkmark$\checkmark:\checkmark$\checkmark^\checkmark\\checkmarkc\checkmarki\checkmarkr\checkmarkc\checkmark$\checkmark \checkmark$\checkmark^\checkmark\\checkmarkc\checkmarki\checkmarkr\checkmarkc\checkmark$\checkmark$\checkmark$\checkmark^\checkmark\\checkmarkc\checkmarki\checkmarkr\checkmarkc\checkmark$\checkmarkF\checkmark$\checkmark^\checkmark\\checkmarkc\checkmarki\checkmarkr\checkmarkc\checkmark$\checkmark_\checkmark$\checkmark^\checkmark\\checkmarkc\checkmarki\checkmarkr\checkmarkc\checkmark$\checkmark{\checkmark$\checkmark^\checkmark\\checkmarkc\checkmarki\checkmarkr\checkmarkc\checkmark$\checkmarkc\checkmark$\checkmark^\checkmark\\checkmarkc\checkmarki\checkmarkr\checkmarkc\checkmark$\checkmarko\checkmark$\checkmark^\checkmark\\checkmarkc\checkmarki\checkmarkr\checkmarkc\checkmark$\checkmarku\checkmark$\checkmark^\checkmark\\checkmarkc\checkmarki\checkmarkr\checkmarkc\checkmark$\checkmarkp\checkmark$\checkmark^\checkmark\\checkmarkc\checkmarki\checkmarkr\checkmarkc\checkmark$\checkmarke\checkmark$\checkmark^\checkmark\\checkmarkc\checkmarki\checkmarkr\checkmarkc\checkmark$\checkmark}\checkmark$\checkmark^\checkmark\\checkmarkc\checkmarki\checkmarkr\checkmarkc\checkmark$\checkmark \checkmark$\checkmark^\checkmark\\checkmarkc\checkmarki\checkmarkr\checkmarkc\checkmark$\checkmark=\checkmark$\checkmark^\checkmark\\checkmarkc\checkmarki\checkmarkr\checkmarkc\checkmark$\checkmark \checkmark$\checkmark^\checkmark\\checkmarkc\checkmarki\checkmarkr\checkmarkc\checkmark$\checkmark1\checkmark$\checkmark^\checkmark\\checkmarkc\checkmarki\checkmarkr\checkmarkc\checkmark$\checkmark0\checkmark$\checkmark^\checkmark\\checkmarkc\checkmarki\checkmarkr\checkmarkc\checkmark$\checkmark0\checkmark$\checkmark^\checkmark\\checkmarkc\checkmarki\checkmarkr\checkmarkc\checkmark$\checkmark$\checkmark$\checkmark^\checkmark\\checkmarkc\checkmarki\checkmarkr\checkmarkc\checkmark$\checkmark \checkmark$\checkmark^\checkmark\\checkmarkc\checkmarki\checkmarkr\checkmarkc\checkmark$\checkmarkN\checkmark$\checkmark^\checkmark\\checkmarkc\checkmarki\checkmarkr\checkmarkc\checkmark$\checkmark
\checkmark$\checkmark^\checkmark\\checkmarkc\checkmarki\checkmarkr\checkmarkc\checkmark$\checkmark \checkmark$\checkmark^\checkmark\\checkmarkc\checkmarki\checkmarkr\checkmarkc\checkmark$\checkmark \checkmark$\checkmark^\checkmark\\checkmarkc\checkmarki\checkmarkr\checkmarkc\checkmark$\checkmark \checkmark$\checkmark^\checkmark\\checkmarkc\checkmarki\checkmarkr\checkmarkc\checkmark$\checkmark \checkmark$\checkmark^\checkmark\\checkmarkc\checkmarki\checkmarkr\checkmarkc\checkmark$\checkmark\\checkmark$\checkmark^\checkmark\\checkmarkc\checkmarki\checkmarkr\checkmarkc\checkmark$\checkmarki\checkmark$\checkmark^\checkmark\\checkmarkc\checkmarki\checkmarkr\checkmarkc\checkmark$\checkmarkt\checkmark$\checkmark^\checkmark\\checkmarkc\checkmarki\checkmarkr\checkmarkc\checkmark$\checkmarke\checkmark$\checkmark^\checkmark\\checkmarkc\checkmarki\checkmarkr\checkmarkc\checkmark$\checkmarkm\checkmark$\checkmark^\checkmark\\checkmarkc\checkmarki\checkmarkr\checkmarkc\checkmark$\checkmark \checkmark$\checkmark^\checkmark\\checkmarkc\checkmarki\checkmarkr\checkmarkc\checkmark$\checkmarkC\checkmark$\checkmark^\checkmark\\checkmarkc\checkmarki\checkmarkr\checkmarkc\checkmark$\checkmarkh\checkmark$\checkmark^\checkmark\\checkmarkc\checkmarki\checkmarkr\checkmarkc\checkmark$\checkmarka\checkmark$\checkmark^\checkmark\\checkmarkc\checkmarki\checkmarkr\checkmarkc\checkmark$\checkmarkr\checkmark$\checkmark^\checkmark\\checkmarkc\checkmarki\checkmarkr\checkmarkc\checkmark$\checkmarkg\checkmark$\checkmark^\checkmark\\checkmarkc\checkmarki\checkmarkr\checkmarkc\checkmark$\checkmarke\checkmark$\checkmark^\checkmark\\checkmarkc\checkmarki\checkmarkr\checkmarkc\checkmark$\checkmark \checkmark$\checkmark^\checkmark\\checkmarkc\checkmarki\checkmarkr\checkmarkc\checkmark$\checkmarké\checkmark$\checkmark^\checkmark\\checkmarkc\checkmarki\checkmarkr\checkmarkc\checkmark$\checkmarkq\checkmark$\checkmark^\checkmark\\checkmarkc\checkmarki\checkmarkr\checkmarkc\checkmark$\checkmarku\checkmark$\checkmark^\checkmark\\checkmarkc\checkmarki\checkmarkr\checkmarkc\checkmark$\checkmarki\checkmark$\checkmark^\checkmark\\checkmarkc\checkmarki\checkmarkr\checkmarkc\checkmark$\checkmarkv\checkmark$\checkmark^\checkmark\\checkmarkc\checkmarki\checkmarkr\checkmarkc\checkmark$\checkmarka\checkmark$\checkmark^\checkmark\\checkmarkc\checkmarki\checkmarkr\checkmarkc\checkmark$\checkmarkl\checkmark$\checkmark^\checkmark\\checkmarkc\checkmarki\checkmarkr\checkmarkc\checkmark$\checkmarke\checkmark$\checkmark^\checkmark\\checkmarkc\checkmarki\checkmarkr\checkmarkc\checkmark$\checkmarkn\checkmark$\checkmark^\checkmark\\checkmarkc\checkmarki\checkmarkr\checkmarkc\checkmark$\checkmarkt\checkmark$\checkmark^\checkmark\\checkmarkc\checkmarki\checkmarkr\checkmarkc\checkmark$\checkmarke\checkmark$\checkmark^\checkmark\\checkmarkc\checkmarki\checkmarkr\checkmarkc\checkmark$\checkmark \checkmark$\checkmark^\checkmark\\checkmarkc\checkmarki\checkmarkr\checkmarkc\checkmark$\checkmark:\checkmark$\checkmark^\checkmark\\checkmarkc\checkmarki\checkmarkr\checkmarkc\checkmark$\checkmark
\checkmark$\checkmark^\checkmark\\checkmarkc\checkmarki\checkmarkr\checkmarkc\checkmark$\checkmark\\checkmark$\checkmark^\checkmark\\checkmarkc\checkmarki\checkmarkr\checkmarkc\checkmark$\checkmarke\checkmark$\checkmark^\checkmark\\checkmarkc\checkmarki\checkmarkr\checkmarkc\checkmark$\checkmarkn\checkmark$\checkmark^\checkmark\\checkmarkc\checkmarki\checkmarkr\checkmarkc\checkmark$\checkmarkd\checkmark$\checkmark^\checkmark\\checkmarkc\checkmarki\checkmarkr\checkmarkc\checkmark$\checkmark{\checkmark$\checkmark^\checkmark\\checkmarkc\checkmarki\checkmarkr\checkmarkc\checkmark$\checkmarki\checkmark$\checkmark^\checkmark\\checkmarkc\checkmarki\checkmarkr\checkmarkc\checkmark$\checkmarkt\checkmark$\checkmark^\checkmark\\checkmarkc\checkmarki\checkmarkr\checkmarkc\checkmark$\checkmarke\checkmark$\checkmark^\checkmark\\checkmarkc\checkmarki\checkmarkr\checkmarkc\checkmark$\checkmarkm\checkmark$\checkmark^\checkmark\\checkmarkc\checkmarki\checkmarkr\checkmarkc\checkmark$\checkmarki\checkmark$\checkmark^\checkmark\\checkmarkc\checkmarki\checkmarkr\checkmarkc\checkmark$\checkmarkz\checkmark$\checkmark^\checkmark\\checkmarkc\checkmarki\checkmarkr\checkmarkc\checkmark$\checkmarke\checkmark$\checkmark^\checkmark\\checkmarkc\checkmarki\checkmarkr\checkmarkc\checkmark$\checkmark}\checkmark$\checkmark^\checkmark\\checkmarkc\checkmarki\checkmarkr\checkmarkc\checkmark$\checkmark
\checkmark$\checkmark^\checkmark\\checkmarkc\checkmarki\checkmarkr\checkmarkc\checkmark$\checkmark\\checkmark$\checkmark^\checkmark\\checkmarkc\checkmarki\checkmarkr\checkmarkc\checkmark$\checkmarkb\checkmark$\checkmark^\checkmark\\checkmarkc\checkmarki\checkmarkr\checkmarkc\checkmark$\checkmarke\checkmark$\checkmark^\checkmark\\checkmarkc\checkmarki\checkmarkr\checkmarkc\checkmark$\checkmarkg\checkmark$\checkmark^\checkmark\\checkmarkc\checkmarki\checkmarkr\checkmarkc\checkmark$\checkmarki\checkmark$\checkmark^\checkmark\\checkmarkc\checkmarki\checkmarkr\checkmarkc\checkmark$\checkmarkn\checkmark$\checkmark^\checkmark\\checkmarkc\checkmarki\checkmarkr\checkmarkc\checkmark$\checkmark{\checkmark$\checkmark^\checkmark\\checkmarkc\checkmarki\checkmarkr\checkmarkc\checkmark$\checkmarke\checkmark$\checkmark^\checkmark\\checkmarkc\checkmarki\checkmarkr\checkmarkc\checkmark$\checkmarkq\checkmark$\checkmark^\checkmark\\checkmarkc\checkmarki\checkmarkr\checkmarkc\checkmark$\checkmarku\checkmark$\checkmark^\checkmark\\checkmarkc\checkmarki\checkmarkr\checkmarkc\checkmark$\checkmarka\checkmark$\checkmark^\checkmark\\checkmarkc\checkmarki\checkmarkr\checkmarkc\checkmark$\checkmarkt\checkmark$\checkmark^\checkmark\\checkmarkc\checkmarki\checkmarkr\checkmarkc\checkmark$\checkmarki\checkmark$\checkmark^\checkmark\\checkmarkc\checkmarki\checkmarkr\checkmarkc\checkmark$\checkmarko\checkmark$\checkmark^\checkmark\\checkmarkc\checkmarki\checkmarkr\checkmarkc\checkmark$\checkmarkn\checkmark$\checkmark^\checkmark\\checkmarkc\checkmarki\checkmarkr\checkmarkc\checkmark$\checkmark}\checkmark$\checkmark^\checkmark\\checkmarkc\checkmarki\checkmarkr\checkmarkc\checkmark$\checkmark
\checkmark$\checkmark^\checkmark\\checkmarkc\checkmarki\checkmarkr\checkmarkc\checkmark$\checkmark \checkmark$\checkmark^\checkmark\\checkmarkc\checkmarki\checkmarkr\checkmarkc\checkmark$\checkmark \checkmark$\checkmark^\checkmark\\checkmarkc\checkmarki\checkmarkr\checkmarkc\checkmark$\checkmark \checkmark$\checkmark^\checkmark\\checkmarkc\checkmarki\checkmarkr\checkmarkc\checkmark$\checkmark \checkmark$\checkmark^\checkmark\\checkmarkc\checkmarki\checkmarkr\checkmarkc\checkmark$\checkmarkP\checkmark$\checkmark^\checkmark\\checkmarkc\checkmarki\checkmarkr\checkmarkc\checkmark$\checkmark_\checkmark$\checkmark^\checkmark\\checkmarkc\checkmarki\checkmarkr\checkmarkc\checkmark$\checkmark{\checkmark$\checkmark^\checkmark\\checkmarkc\checkmarki\checkmarkr\checkmarkc\checkmark$\checkmarke\checkmark$\checkmark^\checkmark\\checkmarkc\checkmarki\checkmarkr\checkmarkc\checkmark$\checkmarkq\checkmark$\checkmark^\checkmark\\checkmarkc\checkmarki\checkmarkr\checkmarkc\checkmark$\checkmarku\checkmark$\checkmark^\checkmark\\checkmarkc\checkmarki\checkmarkr\checkmarkc\checkmark$\checkmarki\checkmark$\checkmark^\checkmark\\checkmarkc\checkmarki\checkmarkr\checkmarkc\checkmark$\checkmarkv\checkmark$\checkmark^\checkmark\\checkmarkc\checkmarki\checkmarkr\checkmarkc\checkmark$\checkmark}\checkmark$\checkmark^\checkmark\\checkmarkc\checkmarki\checkmarkr\checkmarkc\checkmark$\checkmark \checkmark$\checkmark^\checkmark\\checkmarkc\checkmarki\checkmarkr\checkmarkc\checkmark$\checkmark=\checkmark$\checkmark^\checkmark\\checkmarkc\checkmarki\checkmarkr\checkmarkc\checkmark$\checkmark \checkmark$\checkmark^\checkmark\\checkmarkc\checkmarki\checkmarkr\checkmarkc\checkmark$\checkmark\\checkmark$\checkmark^\checkmark\\checkmarkc\checkmarki\checkmarkr\checkmarkc\checkmark$\checkmarks\checkmark$\checkmark^\checkmark\\checkmarkc\checkmarki\checkmarkr\checkmarkc\checkmark$\checkmarkq\checkmark$\checkmark^\checkmark\\checkmarkc\checkmarki\checkmarkr\checkmarkc\checkmark$\checkmarkr\checkmark$\checkmark^\checkmark\\checkmarkc\checkmarki\checkmarkr\checkmarkc\checkmark$\checkmarkt\checkmark$\checkmark^\checkmark\\checkmarkc\checkmarki\checkmarkr\checkmarkc\checkmark$\checkmark{\checkmark$\checkmark^\checkmark\\checkmarkc\checkmarki\checkmarkr\checkmarkc\checkmark$\checkmarkP\checkmark$\checkmark^\checkmark\\checkmarkc\checkmarki\checkmarkr\checkmarkc\checkmark$\checkmark_\checkmark$\checkmark^\checkmark\\checkmarkc\checkmarki\checkmarkr\checkmarkc\checkmark$\checkmark{\checkmark$\checkmark^\checkmark\\checkmarkc\checkmarki\checkmarkr\checkmarkc\checkmark$\checkmarks\checkmark$\checkmark^\checkmark\\checkmarkc\checkmarki\checkmarkr\checkmarkc\checkmark$\checkmarkt\checkmark$\checkmark^\checkmark\\checkmarkc\checkmarki\checkmarkr\checkmarkc\checkmark$\checkmarka\checkmark$\checkmark^\checkmark\\checkmarkc\checkmarki\checkmarkr\checkmarkc\checkmark$\checkmarkt\checkmark$\checkmark^\checkmark\\checkmarkc\checkmarki\checkmarkr\checkmarkc\checkmark$\checkmark}\checkmark$\checkmark^\checkmark\\checkmarkc\checkmarki\checkmarkr\checkmarkc\checkmark$\checkmark^\checkmark$\checkmark^\checkmark\\checkmarkc\checkmarki\checkmarkr\checkmarkc\checkmark$\checkmark2\checkmark$\checkmark^\checkmark\\checkmarkc\checkmarki\checkmarkr\checkmarkc\checkmark$\checkmark \checkmark$\checkmark^\checkmark\\checkmarkc\checkmarki\checkmarkr\checkmarkc\checkmark$\checkmark+\checkmark$\checkmark^\checkmark\\checkmarkc\checkmarki\checkmarkr\checkmarkc\checkmark$\checkmark \checkmark$\checkmark^\checkmark\\checkmarkc\checkmarki\checkmarkr\checkmarkc\checkmark$\checkmarkF\checkmark$\checkmark^\checkmark\\checkmarkc\checkmarki\checkmarkr\checkmarkc\checkmark$\checkmark_\checkmark$\checkmark^\checkmark\\checkmarkc\checkmarki\checkmarkr\checkmarkc\checkmark$\checkmark{\checkmark$\checkmark^\checkmark\\checkmarkc\checkmarki\checkmarkr\checkmarkc\checkmark$\checkmarkc\checkmark$\checkmark^\checkmark\\checkmarkc\checkmarki\checkmarkr\checkmarkc\checkmark$\checkmarko\checkmark$\checkmark^\checkmark\\checkmarkc\checkmarki\checkmarkr\checkmarkc\checkmark$\checkmarku\checkmark$\checkmark^\checkmark\\checkmarkc\checkmarki\checkmarkr\checkmarkc\checkmark$\checkmarkp\checkmark$\checkmark^\checkmark\\checkmarkc\checkmarki\checkmarkr\checkmarkc\checkmark$\checkmarke\checkmark$\checkmark^\checkmark\\checkmarkc\checkmarki\checkmarkr\checkmarkc\checkmark$\checkmark}\checkmark$\checkmark^\checkmark\\checkmarkc\checkmarki\checkmarkr\checkmarkc\checkmark$\checkmark^\checkmark$\checkmark^\checkmark\\checkmarkc\checkmarki\checkmarkr\checkmarkc\checkmark$\checkmark2\checkmark$\checkmark^\checkmark\\checkmarkc\checkmarki\checkmarkr\checkmarkc\checkmark$\checkmark}\checkmark$\checkmark^\checkmark\\checkmarkc\checkmarki\checkmarkr\checkmarkc\checkmark$\checkmark \checkmark$\checkmark^\checkmark\\checkmarkc\checkmarki\checkmarkr\checkmarkc\checkmark$\checkmark=\checkmark$\checkmark^\checkmark\\checkmarkc\checkmarki\checkmarkr\checkmarkc\checkmark$\checkmark \checkmark$\checkmark^\checkmark\\checkmarkc\checkmarki\checkmarkr\checkmarkc\checkmark$\checkmark\\checkmark$\checkmark^\checkmark\\checkmarkc\checkmarki\checkmarkr\checkmarkc\checkmark$\checkmarks\checkmark$\checkmark^\checkmark\\checkmarkc\checkmarki\checkmarkr\checkmarkc\checkmark$\checkmarkq\checkmark$\checkmark^\checkmark\\checkmarkc\checkmarki\checkmarkr\checkmarkc\checkmark$\checkmarkr\checkmark$\checkmark^\checkmark\\checkmarkc\checkmarki\checkmarkr\checkmarkc\checkmark$\checkmarkt\checkmark$\checkmark^\checkmark\\checkmarkc\checkmarki\checkmarkr\checkmarkc\checkmark$\checkmark{\checkmark$\checkmark^\checkmark\\checkmarkc\checkmarki\checkmarkr\checkmarkc\checkmark$\checkmark3\checkmark$\checkmark^\checkmark\\checkmarkc\checkmarki\checkmarkr\checkmarkc\checkmark$\checkmark9\checkmark$\checkmark^\checkmark\\checkmarkc\checkmarki\checkmarkr\checkmarkc\checkmark$\checkmark{\checkmark$\checkmark^\checkmark\\checkmarkc\checkmarki\checkmarkr\checkmarkc\checkmark$\checkmark,\checkmark$\checkmark^\checkmark\\checkmarkc\checkmarki\checkmarkr\checkmarkc\checkmark$\checkmark}\checkmark$\checkmark^\checkmark\\checkmarkc\checkmarki\checkmarkr\checkmarkc\checkmark$\checkmark2\checkmark$\checkmark^\checkmark\\checkmarkc\checkmarki\checkmarkr\checkmarkc\checkmark$\checkmark4\checkmark$\checkmark^\checkmark\\checkmarkc\checkmarki\checkmarkr\checkmarkc\checkmark$\checkmark^\checkmark$\checkmark^\checkmark\\checkmarkc\checkmarki\checkmarkr\checkmarkc\checkmark$\checkmark2\checkmark$\checkmark^\checkmark\\checkmarkc\checkmarki\checkmarkr\checkmarkc\checkmark$\checkmark \checkmark$\checkmark^\checkmark\\checkmarkc\checkmarki\checkmarkr\checkmarkc\checkmark$\checkmark+\checkmark$\checkmark^\checkmark\\checkmarkc\checkmarki\checkmarkr\checkmarkc\checkmark$\checkmark \checkmark$\checkmark^\checkmark\\checkmarkc\checkmarki\checkmarkr\checkmarkc\checkmark$\checkmark1\checkmark$\checkmark^\checkmark\\checkmarkc\checkmarki\checkmarkr\checkmarkc\checkmark$\checkmark0\checkmark$\checkmark^\checkmark\\checkmarkc\checkmarki\checkmarkr\checkmarkc\checkmark$\checkmark0\checkmark$\checkmark^\checkmark\\checkmarkc\checkmarki\checkmarkr\checkmarkc\checkmark$\checkmark^\checkmark$\checkmark^\checkmark\\checkmarkc\checkmarki\checkmarkr\checkmarkc\checkmark$\checkmark2\checkmark$\checkmark^\checkmark\\checkmarkc\checkmarki\checkmarkr\checkmarkc\checkmark$\checkmark}\checkmark$\checkmark^\checkmark\\checkmarkc\checkmarki\checkmarkr\checkmarkc\checkmark$\checkmark \checkmark$\checkmark^\checkmark\\checkmarkc\checkmarki\checkmarkr\checkmarkc\checkmark$\checkmark=\checkmark$\checkmark^\checkmark\\checkmarkc\checkmarki\checkmarkr\checkmarkc\checkmark$\checkmark \checkmark$\checkmark^\checkmark\\checkmarkc\checkmarki\checkmarkr\checkmarkc\checkmark$\checkmark\\checkmark$\checkmark^\checkmark\\checkmarkc\checkmarki\checkmarkr\checkmarkc\checkmark$\checkmarks\checkmark$\checkmark^\checkmark\\checkmarkc\checkmarki\checkmarkr\checkmarkc\checkmark$\checkmarkq\checkmark$\checkmark^\checkmark\\checkmarkc\checkmarki\checkmarkr\checkmarkc\checkmark$\checkmarkr\checkmark$\checkmark^\checkmark\\checkmarkc\checkmarki\checkmarkr\checkmarkc\checkmark$\checkmarkt\checkmark$\checkmark^\checkmark\\checkmarkc\checkmarki\checkmarkr\checkmarkc\checkmark$\checkmark{\checkmark$\checkmark^\checkmark\\checkmarkc\checkmarki\checkmarkr\checkmarkc\checkmark$\checkmark1\checkmark$\checkmark^\checkmark\\checkmarkc\checkmarki\checkmarkr\checkmarkc\checkmark$\checkmark\\checkmark$\checkmark^\checkmark\\checkmarkc\checkmarki\checkmarkr\checkmarkc\checkmark$\checkmark,\checkmark$\checkmark^\checkmark\\checkmarkc\checkmarki\checkmarkr\checkmarkc\checkmark$\checkmark5\checkmark$\checkmark^\checkmark\\checkmarkc\checkmarki\checkmarkr\checkmarkc\checkmark$\checkmark3\checkmark$\checkmark^\checkmark\\checkmarkc\checkmarki\checkmarkr\checkmarkc\checkmark$\checkmark9\checkmark$\checkmark^\checkmark\\checkmarkc\checkmarki\checkmarkr\checkmarkc\checkmark$\checkmark{\checkmark$\checkmark^\checkmark\\checkmarkc\checkmarki\checkmarkr\checkmarkc\checkmark$\checkmark,\checkmark$\checkmark^\checkmark\\checkmarkc\checkmarki\checkmarkr\checkmarkc\checkmark$\checkmark}\checkmark$\checkmark^\checkmark\\checkmarkc\checkmarki\checkmarkr\checkmarkc\checkmark$\checkmark8\checkmark$\checkmark^\checkmark\\checkmarkc\checkmarki\checkmarkr\checkmarkc\checkmark$\checkmark \checkmark$\checkmark^\checkmark\\checkmarkc\checkmarki\checkmarkr\checkmarkc\checkmark$\checkmark+\checkmark$\checkmark^\checkmark\\checkmarkc\checkmarki\checkmarkr\checkmarkc\checkmark$\checkmark \checkmark$\checkmark^\checkmark\\checkmarkc\checkmarki\checkmarkr\checkmarkc\checkmark$\checkmark1\checkmark$\checkmark^\checkmark\\checkmarkc\checkmarki\checkmarkr\checkmarkc\checkmark$\checkmark0\checkmark$\checkmark^\checkmark\\checkmarkc\checkmarki\checkmarkr\checkmarkc\checkmark$\checkmark\\checkmark$\checkmark^\checkmark\\checkmarkc\checkmarki\checkmarkr\checkmarkc\checkmark$\checkmark,\checkmark$\checkmark^\checkmark\\checkmarkc\checkmarki\checkmarkr\checkmarkc\checkmark$\checkmark0\checkmark$\checkmark^\checkmark\\checkmarkc\checkmarki\checkmarkr\checkmarkc\checkmark$\checkmark0\checkmark$\checkmark^\checkmark\\checkmarkc\checkmarki\checkmarkr\checkmarkc\checkmark$\checkmark0\checkmark$\checkmark^\checkmark\\checkmarkc\checkmarki\checkmarkr\checkmarkc\checkmark$\checkmark}\checkmark$\checkmark^\checkmark\\checkmarkc\checkmarki\checkmarkr\checkmarkc\checkmark$\checkmark \checkmark$\checkmark^\checkmark\\checkmarkc\checkmarki\checkmarkr\checkmarkc\checkmark$\checkmark=\checkmark$\checkmark^\checkmark\\checkmarkc\checkmarki\checkmarkr\checkmarkc\checkmark$\checkmark \checkmark$\checkmark^\checkmark\\checkmarkc\checkmarki\checkmarkr\checkmarkc\checkmark$\checkmark\\checkmark$\checkmark^\checkmark\\checkmarkc\checkmarki\checkmarkr\checkmarkc\checkmark$\checkmarkb\checkmark$\checkmark^\checkmark\\checkmarkc\checkmarki\checkmarkr\checkmarkc\checkmark$\checkmarko\checkmark$\checkmark^\checkmark\\checkmarkc\checkmarki\checkmarkr\checkmarkc\checkmark$\checkmarkx\checkmark$\checkmark^\checkmark\\checkmarkc\checkmarki\checkmarkr\checkmarkc\checkmark$\checkmarke\checkmark$\checkmark^\checkmark\\checkmarkc\checkmarki\checkmarkr\checkmarkc\checkmark$\checkmarkd\checkmark$\checkmark^\checkmark\\checkmarkc\checkmarki\checkmarkr\checkmarkc\checkmark$\checkmark{\checkmark$\checkmark^\checkmark\\checkmarkc\checkmarki\checkmarkr\checkmarkc\checkmark$\checkmark1\checkmark$\checkmark^\checkmark\\checkmarkc\checkmarki\checkmarkr\checkmarkc\checkmark$\checkmark0\checkmark$\checkmark^\checkmark\\checkmarkc\checkmarki\checkmarkr\checkmarkc\checkmark$\checkmark7\checkmark$\checkmark^\checkmark\\checkmarkc\checkmarki\checkmarkr\checkmarkc\checkmark$\checkmark{\checkmark$\checkmark^\checkmark\\checkmarkc\checkmarki\checkmarkr\checkmarkc\checkmark$\checkmark,\checkmark$\checkmark^\checkmark\\checkmarkc\checkmarki\checkmarkr\checkmarkc\checkmark$\checkmark}\checkmark$\checkmark^\checkmark\\checkmarkc\checkmarki\checkmarkr\checkmarkc\checkmark$\checkmark4\checkmark$\checkmark^\checkmark\\checkmarkc\checkmarki\checkmarkr\checkmarkc\checkmark$\checkmark \checkmark$\checkmark^\checkmark\\checkmarkc\checkmarki\checkmarkr\checkmarkc\checkmark$\checkmark\\checkmark$\checkmark^\checkmark\\checkmarkc\checkmarki\checkmarkr\checkmarkc\checkmark$\checkmarkt\checkmark$\checkmark^\checkmark\\checkmarkc\checkmarki\checkmarkr\checkmarkc\checkmark$\checkmarke\checkmark$\checkmark^\checkmark\\checkmarkc\checkmarki\checkmarkr\checkmarkc\checkmark$\checkmarkx\checkmark$\checkmark^\checkmark\\checkmarkc\checkmarki\checkmarkr\checkmarkc\checkmark$\checkmarkt\checkmark$\checkmark^\checkmark\\checkmarkc\checkmarki\checkmarkr\checkmarkc\checkmark$\checkmark{\checkmark$\checkmark^\checkmark\\checkmarkc\checkmarki\checkmarkr\checkmarkc\checkmark$\checkmark \checkmark$\checkmark^\checkmark\\checkmarkc\checkmarki\checkmarkr\checkmarkc\checkmark$\checkmarkN\checkmark$\checkmark^\checkmark\\checkmarkc\checkmarki\checkmarkr\checkmarkc\checkmark$\checkmark}\checkmark$\checkmark^\checkmark\\checkmarkc\checkmarki\checkmarkr\checkmarkc\checkmark$\checkmark}\checkmark$\checkmark^\checkmark\\checkmarkc\checkmarki\checkmarkr\checkmarkc\checkmark$\checkmark
\checkmark$\checkmark^\checkmark\\checkmarkc\checkmarki\checkmarkr\checkmarkc\checkmark$\checkmark \checkmark$\checkmark^\checkmark\\checkmarkc\checkmarki\checkmarkr\checkmarkc\checkmark$\checkmark \checkmark$\checkmark^\checkmark\\checkmarkc\checkmarki\checkmarkr\checkmarkc\checkmark$\checkmark \checkmark$\checkmark^\checkmark\\checkmarkc\checkmarki\checkmarkr\checkmarkc\checkmark$\checkmark \checkmark$\checkmark^\checkmark\\checkmarkc\checkmarki\checkmarkr\checkmarkc\checkmark$\checkmark\\checkmark$\checkmark^\checkmark\\checkmarkc\checkmarki\checkmarkr\checkmarkc\checkmark$\checkmarkl\checkmark$\checkmark^\checkmark\\checkmarkc\checkmarki\checkmarkr\checkmarkc\checkmark$\checkmarka\checkmark$\checkmark^\checkmark\\checkmarkc\checkmarki\checkmarkr\checkmarkc\checkmark$\checkmarkb\checkmark$\checkmark^\checkmark\\checkmarkc\checkmarki\checkmarkr\checkmarkc\checkmark$\checkmarke\checkmark$\checkmark^\checkmark\\checkmarkc\checkmarki\checkmarkr\checkmarkc\checkmark$\checkmarkl\checkmark$\checkmark^\checkmark\\checkmarkc\checkmarki\checkmarkr\checkmarkc\checkmark$\checkmark{\checkmark$\checkmark^\checkmark\\checkmarkc\checkmarki\checkmarkr\checkmarkc\checkmark$\checkmarke\checkmark$\checkmark^\checkmark\\checkmarkc\checkmarki\checkmarkr\checkmarkc\checkmark$\checkmarkq\checkmark$\checkmark^\checkmark\\checkmarkc\checkmarki\checkmarkr\checkmarkc\checkmark$\checkmark:\checkmark$\checkmark^\checkmark\\checkmarkc\checkmarki\checkmarkr\checkmarkc\checkmark$\checkmarkc\checkmark$\checkmark^\checkmark\\checkmarkc\checkmarki\checkmarkr\checkmarkc\checkmark$\checkmarkh\checkmark$\checkmark^\checkmark\\checkmarkc\checkmarki\checkmarkr\checkmarkc\checkmark$\checkmarka\checkmark$\checkmark^\checkmark\\checkmarkc\checkmarki\checkmarkr\checkmarkc\checkmark$\checkmarkr\checkmark$\checkmark^\checkmark\\checkmarkc\checkmarki\checkmarkr\checkmarkc\checkmark$\checkmarkg\checkmark$\checkmark^\checkmark\\checkmarkc\checkmarki\checkmarkr\checkmarkc\checkmark$\checkmarke\checkmark$\checkmark^\checkmark\\checkmarkc\checkmarki\checkmarkr\checkmarkc\checkmark$\checkmark_\checkmark$\checkmark^\checkmark\\checkmarkc\checkmarki\checkmarkr\checkmarkc\checkmark$\checkmarke\checkmark$\checkmark^\checkmark\\checkmarkc\checkmarki\checkmarkr\checkmarkc\checkmark$\checkmarkq\checkmark$\checkmark^\checkmark\\checkmarkc\checkmarki\checkmarkr\checkmarkc\checkmark$\checkmarku\checkmark$\checkmark^\checkmark\\checkmarkc\checkmarki\checkmarkr\checkmarkc\checkmark$\checkmarki\checkmark$\checkmark^\checkmark\\checkmarkc\checkmarki\checkmarkr\checkmarkc\checkmark$\checkmarkv\checkmark$\checkmark^\checkmark\\checkmarkc\checkmarki\checkmarkr\checkmarkc\checkmark$\checkmark}\checkmark$\checkmark^\checkmark\\checkmarkc\checkmarki\checkmarkr\checkmarkc\checkmark$\checkmark
\checkmark$\checkmark^\checkmark\\checkmarkc\checkmarki\checkmarkr\checkmarkc\checkmark$\checkmark\\checkmark$\checkmark^\checkmark\\checkmarkc\checkmarki\checkmarkr\checkmarkc\checkmark$\checkmarke\checkmark$\checkmark^\checkmark\\checkmarkc\checkmarki\checkmarkr\checkmarkc\checkmark$\checkmarkn\checkmark$\checkmark^\checkmark\\checkmarkc\checkmarki\checkmarkr\checkmarkc\checkmark$\checkmarkd\checkmark$\checkmark^\checkmark\\checkmarkc\checkmarki\checkmarkr\checkmarkc\checkmark$\checkmark{\checkmark$\checkmark^\checkmark\\checkmarkc\checkmarki\checkmarkr\checkmarkc\checkmark$\checkmarke\checkmark$\checkmark^\checkmark\\checkmarkc\checkmarki\checkmarkr\checkmarkc\checkmark$\checkmarkq\checkmark$\checkmark^\checkmark\\checkmarkc\checkmarki\checkmarkr\checkmarkc\checkmark$\checkmarku\checkmark$\checkmark^\checkmark\\checkmarkc\checkmarki\checkmarkr\checkmarkc\checkmark$\checkmarka\checkmark$\checkmark^\checkmark\\checkmarkc\checkmarki\checkmarkr\checkmarkc\checkmark$\checkmarkt\checkmark$\checkmark^\checkmark\\checkmarkc\checkmarki\checkmarkr\checkmarkc\checkmark$\checkmarki\checkmark$\checkmark^\checkmark\\checkmarkc\checkmarki\checkmarkr\checkmarkc\checkmark$\checkmarko\checkmark$\checkmark^\checkmark\\checkmarkc\checkmarki\checkmarkr\checkmarkc\checkmark$\checkmarkn\checkmark$\checkmark^\checkmark\\checkmarkc\checkmarki\checkmarkr\checkmarkc\checkmark$\checkmark}\checkmark$\checkmark^\checkmark\\checkmarkc\checkmarki\checkmarkr\checkmarkc\checkmark$\checkmark
\checkmark$\checkmark^\checkmark\\checkmarkc\checkmarki\checkmarkr\checkmarkc\checkmark$\checkmark
\checkmark$\checkmark^\checkmark\\checkmarkc\checkmarki\checkmarkr\checkmarkc\checkmark$\checkmark\\checkmark$\checkmark^\checkmark\\checkmarkc\checkmarki\checkmarkr\checkmarkc\checkmark$\checkmarks\checkmark$\checkmark^\checkmark\\checkmarkc\checkmarki\checkmarkr\checkmarkc\checkmark$\checkmarku\checkmark$\checkmark^\checkmark\\checkmarkc\checkmarki\checkmarkr\checkmarkc\checkmark$\checkmarkb\checkmark$\checkmark^\checkmark\\checkmarkc\checkmarki\checkmarkr\checkmarkc\checkmark$\checkmarks\checkmark$\checkmark^\checkmark\\checkmarkc\checkmarki\checkmarkr\checkmarkc\checkmark$\checkmarke\checkmark$\checkmark^\checkmark\\checkmarkc\checkmarki\checkmarkr\checkmarkc\checkmark$\checkmarkc\checkmark$\checkmark^\checkmark\\checkmarkc\checkmarki\checkmarkr\checkmarkc\checkmark$\checkmarkt\checkmark$\checkmark^\checkmark\\checkmarkc\checkmarki\checkmarkr\checkmarkc\checkmark$\checkmarki\checkmark$\checkmark^\checkmark\\checkmarkc\checkmarki\checkmarkr\checkmarkc\checkmark$\checkmarko\checkmark$\checkmark^\checkmark\\checkmarkc\checkmarki\checkmarkr\checkmarkc\checkmark$\checkmarkn\checkmark$\checkmark^\checkmark\\checkmarkc\checkmarki\checkmarkr\checkmarkc\checkmark$\checkmark{\checkmark$\checkmark^\checkmark\\checkmarkc\checkmarki\checkmarkr\checkmarkc\checkmark$\checkmarkC\checkmark$\checkmark^\checkmark\\checkmarkc\checkmarki\checkmarkr\checkmarkc\checkmark$\checkmarka\checkmark$\checkmark^\checkmark\\checkmarkc\checkmarki\checkmarkr\checkmarkc\checkmark$\checkmarkl\checkmark$\checkmark^\checkmark\\checkmarkc\checkmarki\checkmarkr\checkmarkc\checkmark$\checkmarkc\checkmark$\checkmark^\checkmark\\checkmarkc\checkmarki\checkmarkr\checkmarkc\checkmark$\checkmarku\checkmark$\checkmark^\checkmark\\checkmarkc\checkmarki\checkmarkr\checkmarkc\checkmark$\checkmarkl\checkmark$\checkmark^\checkmark\\checkmarkc\checkmarki\checkmarkr\checkmarkc\checkmark$\checkmark \checkmark$\checkmark^\checkmark\\checkmarkc\checkmarki\checkmarkr\checkmarkc\checkmark$\checkmarkd\checkmark$\checkmark^\checkmark\\checkmarkc\checkmarki\checkmarkr\checkmarkc\checkmark$\checkmarke\checkmark$\checkmark^\checkmark\\checkmarkc\checkmarki\checkmarkr\checkmarkc\checkmark$\checkmark \checkmark$\checkmark^\checkmark\\checkmarkc\checkmarki\checkmarkr\checkmarkc\checkmark$\checkmarkl\checkmark$\checkmark^\checkmark\\checkmarkc\checkmarki\checkmarkr\checkmarkc\checkmark$\checkmarka\checkmark$\checkmark^\checkmark\\checkmarkc\checkmarki\checkmarkr\checkmarkc\checkmark$\checkmark \checkmark$\checkmark^\checkmark\\checkmarkc\checkmarki\checkmarkr\checkmarkc\checkmark$\checkmarkd\checkmark$\checkmark^\checkmark\\checkmarkc\checkmarki\checkmarkr\checkmarkc\checkmark$\checkmarku\checkmark$\checkmark^\checkmark\\checkmarkc\checkmarki\checkmarkr\checkmarkc\checkmark$\checkmarkr\checkmark$\checkmark^\checkmark\\checkmarkc\checkmarki\checkmarkr\checkmarkc\checkmark$\checkmarké\checkmark$\checkmark^\checkmark\\checkmarkc\checkmarki\checkmarkr\checkmarkc\checkmark$\checkmarke\checkmark$\checkmark^\checkmark\\checkmarkc\checkmarki\checkmarkr\checkmarkc\checkmark$\checkmark \checkmark$\checkmark^\checkmark\\checkmarkc\checkmarki\checkmarkr\checkmarkc\checkmark$\checkmarkd\checkmark$\checkmark^\checkmark\\checkmarkc\checkmarki\checkmarkr\checkmarkc\checkmark$\checkmarke\checkmark$\checkmark^\checkmark\\checkmarkc\checkmarki\checkmarkr\checkmarkc\checkmark$\checkmark \checkmark$\checkmark^\checkmark\\checkmarkc\checkmarki\checkmarkr\checkmarkc\checkmark$\checkmarkv\checkmark$\checkmark^\checkmark\\checkmarkc\checkmarki\checkmarkr\checkmarkc\checkmark$\checkmarki\checkmark$\checkmark^\checkmark\\checkmarkc\checkmarki\checkmarkr\checkmarkc\checkmark$\checkmarke\checkmark$\checkmark^\checkmark\\checkmarkc\checkmarki\checkmarkr\checkmarkc\checkmark$\checkmark \checkmark$\checkmark^\checkmark\\checkmarkc\checkmarki\checkmarkr\checkmarkc\checkmark$\checkmarkn\checkmark$\checkmark^\checkmark\\checkmarkc\checkmarki\checkmarkr\checkmarkc\checkmark$\checkmarko\checkmark$\checkmark^\checkmark\\checkmarkc\checkmarki\checkmarkr\checkmarkc\checkmark$\checkmarkm\checkmark$\checkmark^\checkmark\\checkmarkc\checkmarki\checkmarkr\checkmarkc\checkmark$\checkmarki\checkmark$\checkmark^\checkmark\\checkmarkc\checkmarki\checkmarkr\checkmarkc\checkmark$\checkmarkn\checkmark$\checkmark^\checkmark\\checkmarkc\checkmarki\checkmarkr\checkmarkc\checkmark$\checkmarka\checkmark$\checkmark^\checkmark\\checkmarkc\checkmarki\checkmarkr\checkmarkc\checkmark$\checkmarkl\checkmark$\checkmark^\checkmark\\checkmarkc\checkmarki\checkmarkr\checkmarkc\checkmark$\checkmarke\checkmark$\checkmark^\checkmark\\checkmarkc\checkmarki\checkmarkr\checkmarkc\checkmark$\checkmark \checkmark$\checkmark^\checkmark\\checkmarkc\checkmarki\checkmarkr\checkmarkc\checkmark$\checkmarkL\checkmark$\checkmark^\checkmark\\checkmarkc\checkmarki\checkmarkr\checkmarkc\checkmark$\checkmark$\checkmark$\checkmark^\checkmark\\checkmarkc\checkmarki\checkmarkr\checkmarkc\checkmark$\checkmark_\checkmark$\checkmark^\checkmark\\checkmarkc\checkmarki\checkmarkr\checkmarkc\checkmark$\checkmark{\checkmark$\checkmark^\checkmark\\checkmarkc\checkmarki\checkmarkr\checkmarkc\checkmark$\checkmark1\checkmark$\checkmark^\checkmark\\checkmarkc\checkmarki\checkmarkr\checkmarkc\checkmark$\checkmark0\checkmark$\checkmark^\checkmark\\checkmarkc\checkmarki\checkmarkr\checkmarkc\checkmark$\checkmark}\checkmark$\checkmark^\checkmark\\checkmarkc\checkmarki\checkmarkr\checkmarkc\checkmark$\checkmark$\checkmark$\checkmark^\checkmark\\checkmarkc\checkmarki\checkmarkr\checkmarkc\checkmark$\checkmark}\checkmark$\checkmark^\checkmark\\checkmarkc\checkmarki\checkmarkr\checkmarkc\checkmark$\checkmark
\checkmark$\checkmark^\checkmark\\checkmarkc\checkmarki\checkmarkr\checkmarkc\checkmark$\checkmarkD\checkmark$\checkmark^\checkmark\\checkmarkc\checkmarki\checkmarkr\checkmarkc\checkmark$\checkmark'\checkmark$\checkmark^\checkmark\\checkmarkc\checkmarki\checkmarkr\checkmarkc\checkmark$\checkmarka\checkmark$\checkmark^\checkmark\\checkmarkc\checkmarki\checkmarkr\checkmarkc\checkmark$\checkmarkp\checkmark$\checkmark^\checkmark\\checkmarkc\checkmarki\checkmarkr\checkmarkc\checkmark$\checkmarkr\checkmark$\checkmark^\checkmark\\checkmarkc\checkmarki\checkmarkr\checkmarkc\checkmark$\checkmarkè\checkmark$\checkmark^\checkmark\\checkmarkc\checkmarki\checkmarkr\checkmarkc\checkmark$\checkmarks\checkmark$\checkmark^\checkmark\\checkmarkc\checkmarki\checkmarkr\checkmarkc\checkmark$\checkmark \checkmark$\checkmark^\checkmark\\checkmarkc\checkmarki\checkmarkr\checkmarkc\checkmark$\checkmarkl\checkmark$\checkmark^\checkmark\\checkmarkc\checkmarki\checkmarkr\checkmarkc\checkmark$\checkmarka\checkmark$\checkmark^\checkmark\\checkmarkc\checkmarki\checkmarkr\checkmarkc\checkmark$\checkmark \checkmark$\checkmark^\checkmark\\checkmarkc\checkmarki\checkmarkr\checkmarkc\checkmark$\checkmarkn\checkmark$\checkmark^\checkmark\\checkmarkc\checkmarki\checkmarkr\checkmarkc\checkmark$\checkmarko\checkmark$\checkmark^\checkmark\\checkmarkc\checkmarki\checkmarkr\checkmarkc\checkmark$\checkmarkr\checkmark$\checkmark^\checkmark\\checkmarkc\checkmarki\checkmarkr\checkmarkc\checkmark$\checkmarkm\checkmark$\checkmark^\checkmark\\checkmarkc\checkmarki\checkmarkr\checkmarkc\checkmark$\checkmarke\checkmark$\checkmark^\checkmark\\checkmarkc\checkmarki\checkmarkr\checkmarkc\checkmark$\checkmark \checkmark$\checkmark^\checkmark\\checkmarkc\checkmarki\checkmarkr\checkmarkc\checkmark$\checkmarkI\checkmark$\checkmark^\checkmark\\checkmarkc\checkmarki\checkmarkr\checkmarkc\checkmark$\checkmarkS\checkmark$\checkmark^\checkmark\\checkmarkc\checkmarki\checkmarkr\checkmarkc\checkmark$\checkmarkO\checkmark$\checkmark^\checkmark\\checkmarkc\checkmarki\checkmarkr\checkmarkc\checkmark$\checkmark \checkmark$\checkmark^\checkmark\\checkmarkc\checkmarki\checkmarkr\checkmarkc\checkmark$\checkmark2\checkmark$\checkmark^\checkmark\\checkmarkc\checkmarki\checkmarkr\checkmarkc\checkmark$\checkmark8\checkmark$\checkmark^\checkmark\\checkmarkc\checkmarki\checkmarkr\checkmarkc\checkmark$\checkmark1\checkmark$\checkmark^\checkmark\\checkmarkc\checkmarki\checkmarkr\checkmarkc\checkmark$\checkmark,\checkmark$\checkmark^\checkmark\\checkmarkc\checkmarki\checkmarkr\checkmarkc\checkmark$\checkmark \checkmark$\checkmark^\checkmark\\checkmarkc\checkmarki\checkmarkr\checkmarkc\checkmark$\checkmarka\checkmark$\checkmark^\checkmark\\checkmarkc\checkmarki\checkmarkr\checkmarkc\checkmark$\checkmarkv\checkmark$\checkmark^\checkmark\\checkmarkc\checkmarki\checkmarkr\checkmarkc\checkmark$\checkmarke\checkmark$\checkmark^\checkmark\\checkmarkc\checkmarki\checkmarkr\checkmarkc\checkmark$\checkmarkc\checkmark$\checkmark^\checkmark\\checkmarkc\checkmarki\checkmarkr\checkmarkc\checkmark$\checkmark \checkmark$\checkmark^\checkmark\\checkmarkc\checkmarki\checkmarkr\checkmarkc\checkmark$\checkmarku\checkmark$\checkmark^\checkmark\\checkmarkc\checkmarki\checkmarkr\checkmarkc\checkmark$\checkmarkn\checkmark$\checkmark^\checkmark\\checkmarkc\checkmarki\checkmarkr\checkmarkc\checkmark$\checkmark \checkmark$\checkmark^\checkmark\\checkmarkc\checkmarki\checkmarkr\checkmarkc\checkmark$\checkmarke\checkmark$\checkmark^\checkmark\\checkmarkc\checkmarki\checkmarkr\checkmarkc\checkmark$\checkmarkx\checkmark$\checkmark^\checkmark\\checkmarkc\checkmarki\checkmarkr\checkmarkc\checkmark$\checkmarkp\checkmark$\checkmark^\checkmark\\checkmarkc\checkmarki\checkmarkr\checkmarkc\checkmark$\checkmarko\checkmark$\checkmark^\checkmark\\checkmarkc\checkmarki\checkmarkr\checkmarkc\checkmark$\checkmarks\checkmark$\checkmark^\checkmark\\checkmarkc\checkmarki\checkmarkr\checkmarkc\checkmark$\checkmarka\checkmark$\checkmark^\checkmark\\checkmarkc\checkmarki\checkmarkr\checkmarkc\checkmark$\checkmarkn\checkmark$\checkmark^\checkmark\\checkmarkc\checkmarki\checkmarkr\checkmarkc\checkmark$\checkmarkt\checkmark$\checkmark^\checkmark\\checkmarkc\checkmarki\checkmarkr\checkmarkc\checkmark$\checkmark \checkmark$\checkmark^\checkmark\\checkmarkc\checkmarki\checkmarkr\checkmarkc\checkmark$\checkmark$\checkmark$\checkmark^\checkmark\\checkmarkc\checkmarki\checkmarkr\checkmarkc\checkmark$\checkmarkp\checkmark$\checkmark^\checkmark\\checkmarkc\checkmarki\checkmarkr\checkmarkc\checkmark$\checkmark \checkmark$\checkmark^\checkmark\\checkmarkc\checkmarki\checkmarkr\checkmarkc\checkmark$\checkmark=\checkmark$\checkmark^\checkmark\\checkmarkc\checkmarki\checkmarkr\checkmarkc\checkmark$\checkmark \checkmark$\checkmark^\checkmark\\checkmarkc\checkmarki\checkmarkr\checkmarkc\checkmark$\checkmark3\checkmark$\checkmark^\checkmark\\checkmarkc\checkmarki\checkmarkr\checkmarkc\checkmark$\checkmark$\checkmark$\checkmark^\checkmark\\checkmarkc\checkmarki\checkmarkr\checkmarkc\checkmark$\checkmark \checkmark$\checkmark^\checkmark\\checkmarkc\checkmarki\checkmarkr\checkmarkc\checkmark$\checkmarkp\checkmark$\checkmark^\checkmark\\checkmarkc\checkmarki\checkmarkr\checkmarkc\checkmark$\checkmarko\checkmark$\checkmark^\checkmark\\checkmarkc\checkmarki\checkmarkr\checkmarkc\checkmark$\checkmarku\checkmark$\checkmark^\checkmark\\checkmarkc\checkmarki\checkmarkr\checkmarkc\checkmark$\checkmarkr\checkmark$\checkmark^\checkmark\\checkmarkc\checkmarki\checkmarkr\checkmarkc\checkmark$\checkmark \checkmark$\checkmark^\checkmark\\checkmarkc\checkmarki\checkmarkr\checkmarkc\checkmark$\checkmarkl\checkmark$\checkmark^\checkmark\\checkmarkc\checkmarki\checkmarkr\checkmarkc\checkmark$\checkmarke\checkmark$\checkmark^\checkmark\\checkmarkc\checkmarki\checkmarkr\checkmarkc\checkmark$\checkmarks\checkmark$\checkmark^\checkmark\\checkmarkc\checkmarki\checkmarkr\checkmarkc\checkmark$\checkmark \checkmark$\checkmark^\checkmark\\checkmarkc\checkmarki\checkmarkr\checkmarkc\checkmark$\checkmarkr\checkmark$\checkmark^\checkmark\\checkmarkc\checkmarki\checkmarkr\checkmarkc\checkmark$\checkmarko\checkmark$\checkmark^\checkmark\\checkmarkc\checkmarki\checkmarkr\checkmarkc\checkmark$\checkmarku\checkmark$\checkmark^\checkmark\\checkmarkc\checkmarki\checkmarkr\checkmarkc\checkmark$\checkmarkl\checkmark$\checkmark^\checkmark\\checkmarkc\checkmarki\checkmarkr\checkmarkc\checkmark$\checkmarke\checkmark$\checkmark^\checkmark\\checkmarkc\checkmarki\checkmarkr\checkmarkc\checkmark$\checkmarkm\checkmark$\checkmark^\checkmark\\checkmarkc\checkmarki\checkmarkr\checkmarkc\checkmark$\checkmarke\checkmark$\checkmark^\checkmark\\checkmarkc\checkmarki\checkmarkr\checkmarkc\checkmark$\checkmarkn\checkmark$\checkmark^\checkmark\\checkmarkc\checkmarki\checkmarkr\checkmarkc\checkmark$\checkmarkt\checkmark$\checkmark^\checkmark\\checkmarkc\checkmarki\checkmarkr\checkmarkc\checkmark$\checkmarks\checkmark$\checkmark^\checkmark\\checkmarkc\checkmarki\checkmarkr\checkmarkc\checkmark$\checkmark \checkmark$\checkmark^\checkmark\\checkmarkc\checkmarki\checkmarkr\checkmarkc\checkmark$\checkmarkà\checkmark$\checkmark^\checkmark\\checkmarkc\checkmarki\checkmarkr\checkmarkc\checkmark$\checkmark \checkmark$\checkmark^\checkmark\\checkmarkc\checkmarki\checkmarkr\checkmarkc\checkmark$\checkmarkb\checkmark$\checkmark^\checkmark\\checkmarkc\checkmarki\checkmarkr\checkmarkc\checkmark$\checkmarki\checkmark$\checkmark^\checkmark\\checkmarkc\checkmarki\checkmarkr\checkmarkc\checkmark$\checkmarkl\checkmark$\checkmark^\checkmark\\checkmarkc\checkmarki\checkmarkr\checkmarkc\checkmark$\checkmarkl\checkmark$\checkmark^\checkmark\\checkmarkc\checkmarki\checkmarkr\checkmarkc\checkmark$\checkmarke\checkmark$\checkmark^\checkmark\\checkmarkc\checkmarki\checkmarkr\checkmarkc\checkmark$\checkmarks\checkmark$\checkmark^\checkmark\\checkmarkc\checkmarki\checkmarkr\checkmarkc\checkmark$\checkmark \checkmark$\checkmark^\checkmark\\checkmarkc\checkmarki\checkmarkr\checkmarkc\checkmark$\checkmark:\checkmark$\checkmark^\checkmark\\checkmarkc\checkmarki\checkmarkr\checkmarkc\checkmark$\checkmark
\checkmark$\checkmark^\checkmark\\checkmarkc\checkmarki\checkmarkr\checkmarkc\checkmark$\checkmark\\checkmark$\checkmark^\checkmark\\checkmarkc\checkmarki\checkmarkr\checkmarkc\checkmark$\checkmarkb\checkmark$\checkmark^\checkmark\\checkmarkc\checkmarki\checkmarkr\checkmarkc\checkmark$\checkmarke\checkmark$\checkmark^\checkmark\\checkmarkc\checkmarki\checkmarkr\checkmarkc\checkmark$\checkmarkg\checkmark$\checkmark^\checkmark\\checkmarkc\checkmarki\checkmarkr\checkmarkc\checkmark$\checkmarki\checkmark$\checkmark^\checkmark\\checkmarkc\checkmarki\checkmarkr\checkmarkc\checkmark$\checkmarkn\checkmark$\checkmark^\checkmark\\checkmarkc\checkmarki\checkmarkr\checkmarkc\checkmark$\checkmark{\checkmark$\checkmark^\checkmark\\checkmarkc\checkmarki\checkmarkr\checkmarkc\checkmark$\checkmarke\checkmark$\checkmark^\checkmark\\checkmarkc\checkmarki\checkmarkr\checkmarkc\checkmark$\checkmarkq\checkmark$\checkmark^\checkmark\\checkmarkc\checkmarki\checkmarkr\checkmarkc\checkmark$\checkmarku\checkmark$\checkmark^\checkmark\\checkmarkc\checkmarki\checkmarkr\checkmarkc\checkmark$\checkmarka\checkmark$\checkmark^\checkmark\\checkmarkc\checkmarki\checkmarkr\checkmarkc\checkmark$\checkmarkt\checkmark$\checkmark^\checkmark\\checkmarkc\checkmarki\checkmarkr\checkmarkc\checkmark$\checkmarki\checkmark$\checkmark^\checkmark\\checkmarkc\checkmarki\checkmarkr\checkmarkc\checkmark$\checkmarko\checkmark$\checkmark^\checkmark\\checkmarkc\checkmarki\checkmarkr\checkmarkc\checkmark$\checkmarkn\checkmark$\checkmark^\checkmark\\checkmarkc\checkmarki\checkmarkr\checkmarkc\checkmark$\checkmark}\checkmark$\checkmark^\checkmark\\checkmarkc\checkmarki\checkmarkr\checkmarkc\checkmark$\checkmark
\checkmark$\checkmark^\checkmark\\checkmarkc\checkmarki\checkmarkr\checkmarkc\checkmark$\checkmark \checkmark$\checkmark^\checkmark\\checkmarkc\checkmarki\checkmarkr\checkmarkc\checkmark$\checkmark \checkmark$\checkmark^\checkmark\\checkmarkc\checkmarki\checkmarkr\checkmarkc\checkmark$\checkmark \checkmark$\checkmark^\checkmark\\checkmarkc\checkmarki\checkmarkr\checkmarkc\checkmark$\checkmark \checkmark$\checkmark^\checkmark\\checkmarkc\checkmarki\checkmarkr\checkmarkc\checkmark$\checkmarkL\checkmark$\checkmark^\checkmark\\checkmarkc\checkmarki\checkmarkr\checkmarkc\checkmark$\checkmark_\checkmark$\checkmark^\checkmark\\checkmarkc\checkmarki\checkmarkr\checkmarkc\checkmark$\checkmark{\checkmark$\checkmark^\checkmark\\checkmarkc\checkmarki\checkmarkr\checkmarkc\checkmark$\checkmark1\checkmark$\checkmark^\checkmark\\checkmarkc\checkmarki\checkmarkr\checkmarkc\checkmark$\checkmark0\checkmark$\checkmark^\checkmark\\checkmarkc\checkmarki\checkmarkr\checkmarkc\checkmark$\checkmark}\checkmark$\checkmark^\checkmark\\checkmarkc\checkmarki\checkmarkr\checkmarkc\checkmark$\checkmark \checkmark$\checkmark^\checkmark\\checkmarkc\checkmarki\checkmarkr\checkmarkc\checkmark$\checkmark=\checkmark$\checkmark^\checkmark\\checkmarkc\checkmarki\checkmarkr\checkmarkc\checkmark$\checkmark \checkmark$\checkmark^\checkmark\\checkmarkc\checkmarki\checkmarkr\checkmarkc\checkmark$\checkmark\\checkmark$\checkmark^\checkmark\\checkmarkc\checkmarki\checkmarkr\checkmarkc\checkmark$\checkmarkl\checkmark$\checkmark^\checkmark\\checkmarkc\checkmarki\checkmarkr\checkmarkc\checkmark$\checkmarke\checkmark$\checkmark^\checkmark\\checkmarkc\checkmarki\checkmarkr\checkmarkc\checkmark$\checkmarkf\checkmark$\checkmark^\checkmark\\checkmarkc\checkmarki\checkmarkr\checkmarkc\checkmark$\checkmarkt\checkmark$\checkmark^\checkmark\\checkmarkc\checkmarki\checkmarkr\checkmarkc\checkmark$\checkmark(\checkmark$\checkmark^\checkmark\\checkmarkc\checkmarki\checkmarkr\checkmarkc\checkmark$\checkmark\\checkmark$\checkmark^\checkmark\\checkmarkc\checkmarki\checkmarkr\checkmarkc\checkmark$\checkmarkf\checkmark$\checkmark^\checkmark\\checkmarkc\checkmarki\checkmarkr\checkmarkc\checkmark$\checkmarkr\checkmark$\checkmark^\checkmark\\checkmarkc\checkmarki\checkmarkr\checkmarkc\checkmark$\checkmarka\checkmark$\checkmark^\checkmark\\checkmarkc\checkmarki\checkmarkr\checkmarkc\checkmark$\checkmarkc\checkmark$\checkmark^\checkmark\\checkmarkc\checkmarki\checkmarkr\checkmarkc\checkmark$\checkmark{\checkmark$\checkmark^\checkmark\\checkmarkc\checkmarki\checkmarkr\checkmarkc\checkmark$\checkmarkC\checkmark$\checkmark^\checkmark\\checkmarkc\checkmarki\checkmarkr\checkmarkc\checkmark$\checkmark_\checkmark$\checkmark^\checkmark\\checkmarkc\checkmarki\checkmarkr\checkmarkc\checkmark$\checkmark{\checkmark$\checkmark^\checkmark\\checkmarkc\checkmarki\checkmarkr\checkmarkc\checkmark$\checkmarkd\checkmark$\checkmark^\checkmark\\checkmarkc\checkmarki\checkmarkr\checkmarkc\checkmark$\checkmarky\checkmark$\checkmark^\checkmark\\checkmarkc\checkmarki\checkmarkr\checkmarkc\checkmark$\checkmarkn\checkmark$\checkmark^\checkmark\\checkmarkc\checkmarki\checkmarkr\checkmarkc\checkmark$\checkmark}\checkmark$\checkmark^\checkmark\\checkmarkc\checkmarki\checkmarkr\checkmarkc\checkmark$\checkmark}\checkmark$\checkmark^\checkmark\\checkmarkc\checkmarki\checkmarkr\checkmarkc\checkmark$\checkmark{\checkmark$\checkmark^\checkmark\\checkmarkc\checkmarki\checkmarkr\checkmarkc\checkmark$\checkmarkP\checkmark$\checkmark^\checkmark\\checkmarkc\checkmarki\checkmarkr\checkmarkc\checkmark$\checkmark_\checkmark$\checkmark^\checkmark\\checkmarkc\checkmarki\checkmarkr\checkmarkc\checkmark$\checkmark{\checkmark$\checkmark^\checkmark\\checkmarkc\checkmarki\checkmarkr\checkmarkc\checkmark$\checkmarke\checkmark$\checkmark^\checkmark\\checkmarkc\checkmarki\checkmarkr\checkmarkc\checkmark$\checkmarkq\checkmark$\checkmark^\checkmark\\checkmarkc\checkmarki\checkmarkr\checkmarkc\checkmark$\checkmarku\checkmark$\checkmark^\checkmark\\checkmarkc\checkmarki\checkmarkr\checkmarkc\checkmark$\checkmarki\checkmark$\checkmark^\checkmark\\checkmarkc\checkmarki\checkmarkr\checkmarkc\checkmark$\checkmarkv\checkmark$\checkmark^\checkmark\\checkmarkc\checkmarki\checkmarkr\checkmarkc\checkmark$\checkmark}\checkmark$\checkmark^\checkmark\\checkmarkc\checkmarki\checkmarkr\checkmarkc\checkmark$\checkmark}\checkmark$\checkmark^\checkmark\\checkmarkc\checkmarki\checkmarkr\checkmarkc\checkmark$\checkmark\\checkmark$\checkmark^\checkmark\\checkmarkc\checkmarki\checkmarkr\checkmarkc\checkmark$\checkmarkr\checkmark$\checkmark^\checkmark\\checkmarkc\checkmarki\checkmarkr\checkmarkc\checkmark$\checkmarki\checkmark$\checkmark^\checkmark\\checkmarkc\checkmarki\checkmarkr\checkmarkc\checkmark$\checkmarkg\checkmark$\checkmark^\checkmark\\checkmarkc\checkmarki\checkmarkr\checkmarkc\checkmark$\checkmarkh\checkmark$\checkmark^\checkmark\\checkmarkc\checkmarki\checkmarkr\checkmarkc\checkmark$\checkmarkt\checkmark$\checkmark^\checkmark\\checkmarkc\checkmarki\checkmarkr\checkmarkc\checkmark$\checkmark)\checkmark$\checkmark^\checkmark\\checkmarkc\checkmarki\checkmarkr\checkmarkc\checkmark$\checkmark^\checkmark$\checkmark^\checkmark\\checkmarkc\checkmarki\checkmarkr\checkmarkc\checkmark$\checkmark3\checkmark$\checkmark^\checkmark\\checkmarkc\checkmarki\checkmarkr\checkmarkc\checkmark$\checkmark \checkmark$\checkmark^\checkmark\\checkmarkc\checkmarki\checkmarkr\checkmarkc\checkmark$\checkmark\\checkmark$\checkmark^\checkmark\\checkmarkc\checkmarki\checkmarkr\checkmarkc\checkmark$\checkmarkt\checkmark$\checkmark^\checkmark\\checkmarkc\checkmarki\checkmarkr\checkmarkc\checkmark$\checkmarki\checkmark$\checkmark^\checkmark\\checkmarkc\checkmarki\checkmarkr\checkmarkc\checkmark$\checkmarkm\checkmark$\checkmark^\checkmark\\checkmarkc\checkmarki\checkmarkr\checkmarkc\checkmark$\checkmarke\checkmark$\checkmark^\checkmark\\checkmarkc\checkmarki\checkmarkr\checkmarkc\checkmark$\checkmarks\checkmark$\checkmark^\checkmark\\checkmarkc\checkmarki\checkmarkr\checkmarkc\checkmark$\checkmark \checkmark$\checkmark^\checkmark\\checkmarkc\checkmarki\checkmarkr\checkmarkc\checkmark$\checkmark5\checkmark$\checkmark^\checkmark\\checkmarkc\checkmarki\checkmarkr\checkmarkc\checkmark$\checkmark0\checkmark$\checkmark^\checkmark\\checkmarkc\checkmarki\checkmarkr\checkmarkc\checkmark$\checkmark \checkmark$\checkmark^\checkmark\\checkmarkc\checkmarki\checkmarkr\checkmarkc\checkmark$\checkmark\\checkmark$\checkmark^\checkmark\\checkmarkc\checkmarki\checkmarkr\checkmarkc\checkmark$\checkmarkt\checkmark$\checkmark^\checkmark\\checkmarkc\checkmarki\checkmarkr\checkmarkc\checkmark$\checkmarke\checkmark$\checkmark^\checkmark\\checkmarkc\checkmarki\checkmarkr\checkmarkc\checkmark$\checkmarkx\checkmark$\checkmark^\checkmark\\checkmarkc\checkmarki\checkmarkr\checkmarkc\checkmark$\checkmarkt\checkmark$\checkmark^\checkmark\\checkmarkc\checkmarki\checkmarkr\checkmarkc\checkmark$\checkmark{\checkmark$\checkmark^\checkmark\\checkmarkc\checkmarki\checkmarkr\checkmarkc\checkmark$\checkmark \checkmark$\checkmark^\checkmark\\checkmarkc\checkmarki\checkmarkr\checkmarkc\checkmark$\checkmarkk\checkmark$\checkmark^\checkmark\\checkmarkc\checkmarki\checkmarkr\checkmarkc\checkmark$\checkmarkm\checkmark$\checkmark^\checkmark\\checkmarkc\checkmarki\checkmarkr\checkmarkc\checkmark$\checkmark}\checkmark$\checkmark^\checkmark\\checkmarkc\checkmarki\checkmarkr\checkmarkc\checkmark$\checkmark \checkmark$\checkmark^\checkmark\\checkmarkc\checkmarki\checkmarkr\checkmarkc\checkmark$\checkmark=\checkmark$\checkmark^\checkmark\\checkmarkc\checkmarki\checkmarkr\checkmarkc\checkmark$\checkmark \checkmark$\checkmark^\checkmark\\checkmarkc\checkmarki\checkmarkr\checkmarkc\checkmark$\checkmark\\checkmark$\checkmark^\checkmark\\checkmarkc\checkmarki\checkmarkr\checkmarkc\checkmark$\checkmarkl\checkmark$\checkmark^\checkmark\\checkmarkc\checkmarki\checkmarkr\checkmarkc\checkmark$\checkmarke\checkmark$\checkmark^\checkmark\\checkmarkc\checkmarki\checkmarkr\checkmarkc\checkmark$\checkmarkf\checkmark$\checkmark^\checkmark\\checkmarkc\checkmarki\checkmarkr\checkmarkc\checkmark$\checkmarkt\checkmark$\checkmark^\checkmark\\checkmarkc\checkmarki\checkmarkr\checkmarkc\checkmark$\checkmark(\checkmark$\checkmark^\checkmark\\checkmarkc\checkmarki\checkmarkr\checkmarkc\checkmark$\checkmark\\checkmark$\checkmark^\checkmark\\checkmarkc\checkmarki\checkmarkr\checkmarkc\checkmark$\checkmarkf\checkmark$\checkmark^\checkmark\\checkmarkc\checkmarki\checkmarkr\checkmarkc\checkmark$\checkmarkr\checkmark$\checkmark^\checkmark\\checkmarkc\checkmarki\checkmarkr\checkmarkc\checkmark$\checkmarka\checkmark$\checkmark^\checkmark\\checkmarkc\checkmarki\checkmarkr\checkmarkc\checkmark$\checkmarkc\checkmark$\checkmark^\checkmark\\checkmarkc\checkmarki\checkmarkr\checkmarkc\checkmark$\checkmark{\checkmark$\checkmark^\checkmark\\checkmarkc\checkmarki\checkmarkr\checkmarkc\checkmark$\checkmark1\checkmark$\checkmark^\checkmark\\checkmarkc\checkmarki\checkmarkr\checkmarkc\checkmark$\checkmark1\checkmark$\checkmark^\checkmark\\checkmarkc\checkmarki\checkmarkr\checkmarkc\checkmark$\checkmark\\checkmark$\checkmark^\checkmark\\checkmarkc\checkmarki\checkmarkr\checkmarkc\checkmark$\checkmark,\checkmark$\checkmark^\checkmark\\checkmarkc\checkmarki\checkmarkr\checkmarkc\checkmark$\checkmark5\checkmark$\checkmark^\checkmark\\checkmarkc\checkmarki\checkmarkr\checkmarkc\checkmark$\checkmark0\checkmark$\checkmark^\checkmark\\checkmarkc\checkmarki\checkmarkr\checkmarkc\checkmark$\checkmark0\checkmark$\checkmark^\checkmark\\checkmarkc\checkmarki\checkmarkr\checkmarkc\checkmark$\checkmark}\checkmark$\checkmark^\checkmark\\checkmarkc\checkmarki\checkmarkr\checkmarkc\checkmark$\checkmark{\checkmark$\checkmark^\checkmark\\checkmarkc\checkmarki\checkmarkr\checkmarkc\checkmark$\checkmark1\checkmark$\checkmark^\checkmark\\checkmarkc\checkmarki\checkmarkr\checkmarkc\checkmark$\checkmark0\checkmark$\checkmark^\checkmark\\checkmarkc\checkmarki\checkmarkr\checkmarkc\checkmark$\checkmark7\checkmark$\checkmark^\checkmark\\checkmarkc\checkmarki\checkmarkr\checkmarkc\checkmark$\checkmark{\checkmark$\checkmark^\checkmark\\checkmarkc\checkmarki\checkmarkr\checkmarkc\checkmark$\checkmark,\checkmark$\checkmark^\checkmark\\checkmarkc\checkmarki\checkmarkr\checkmarkc\checkmark$\checkmark}\checkmark$\checkmark^\checkmark\\checkmarkc\checkmarki\checkmarkr\checkmarkc\checkmark$\checkmark4\checkmark$\checkmark^\checkmark\\checkmarkc\checkmarki\checkmarkr\checkmarkc\checkmark$\checkmark}\checkmark$\checkmark^\checkmark\\checkmarkc\checkmarki\checkmarkr\checkmarkc\checkmark$\checkmark\\checkmark$\checkmark^\checkmark\\checkmarkc\checkmarki\checkmarkr\checkmarkc\checkmark$\checkmarkr\checkmark$\checkmark^\checkmark\\checkmarkc\checkmarki\checkmarkr\checkmarkc\checkmark$\checkmarki\checkmark$\checkmark^\checkmark\\checkmarkc\checkmarki\checkmarkr\checkmarkc\checkmark$\checkmarkg\checkmark$\checkmark^\checkmark\\checkmarkc\checkmarki\checkmarkr\checkmarkc\checkmark$\checkmarkh\checkmark$\checkmark^\checkmark\\checkmarkc\checkmarki\checkmarkr\checkmarkc\checkmark$\checkmarkt\checkmark$\checkmark^\checkmark\\checkmarkc\checkmarki\checkmarkr\checkmarkc\checkmark$\checkmark)\checkmark$\checkmark^\checkmark\\checkmarkc\checkmarki\checkmarkr\checkmarkc\checkmark$\checkmark^\checkmark$\checkmark^\checkmark\\checkmarkc\checkmarki\checkmarkr\checkmarkc\checkmark$\checkmark3\checkmark$\checkmark^\checkmark\\checkmarkc\checkmarki\checkmarkr\checkmarkc\checkmark$\checkmark \checkmark$\checkmark^\checkmark\\checkmarkc\checkmarki\checkmarkr\checkmarkc\checkmark$\checkmark\\checkmark$\checkmark^\checkmark\\checkmarkc\checkmarki\checkmarkr\checkmarkc\checkmark$\checkmarkt\checkmark$\checkmark^\checkmark\\checkmarkc\checkmarki\checkmarkr\checkmarkc\checkmark$\checkmarki\checkmark$\checkmark^\checkmark\\checkmarkc\checkmarki\checkmarkr\checkmarkc\checkmark$\checkmarkm\checkmark$\checkmark^\checkmark\\checkmarkc\checkmarki\checkmarkr\checkmarkc\checkmark$\checkmarke\checkmark$\checkmark^\checkmark\\checkmarkc\checkmarki\checkmarkr\checkmarkc\checkmark$\checkmarks\checkmark$\checkmark^\checkmark\\checkmarkc\checkmarki\checkmarkr\checkmarkc\checkmark$\checkmark \checkmark$\checkmark^\checkmark\\checkmarkc\checkmarki\checkmarkr\checkmarkc\checkmark$\checkmark5\checkmark$\checkmark^\checkmark\\checkmarkc\checkmarki\checkmarkr\checkmarkc\checkmark$\checkmark0\checkmark$\checkmark^\checkmark\\checkmarkc\checkmarki\checkmarkr\checkmarkc\checkmark$\checkmark \checkmark$\checkmark^\checkmark\\checkmarkc\checkmarki\checkmarkr\checkmarkc\checkmark$\checkmark=\checkmark$\checkmark^\checkmark\\checkmarkc\checkmarki\checkmarkr\checkmarkc\checkmark$\checkmark \checkmark$\checkmark^\checkmark\\checkmarkc\checkmarki\checkmarkr\checkmarkc\checkmark$\checkmark(\checkmark$\checkmark^\checkmark\\checkmarkc\checkmarki\checkmarkr\checkmarkc\checkmark$\checkmark1\checkmark$\checkmark^\checkmark\\checkmarkc\checkmarki\checkmarkr\checkmarkc\checkmark$\checkmark0\checkmark$\checkmark^\checkmark\\checkmarkc\checkmarki\checkmarkr\checkmarkc\checkmark$\checkmark7\checkmark$\checkmark^\checkmark\\checkmarkc\checkmarki\checkmarkr\checkmarkc\checkmark$\checkmark{\checkmark$\checkmark^\checkmark\\checkmarkc\checkmarki\checkmarkr\checkmarkc\checkmark$\checkmark,\checkmark$\checkmark^\checkmark\\checkmarkc\checkmarki\checkmarkr\checkmarkc\checkmark$\checkmark}\checkmark$\checkmark^\checkmark\\checkmarkc\checkmarki\checkmarkr\checkmarkc\checkmark$\checkmark0\checkmark$\checkmark^\checkmark\\checkmarkc\checkmarki\checkmarkr\checkmarkc\checkmark$\checkmark7\checkmark$\checkmark^\checkmark\\checkmarkc\checkmarki\checkmarkr\checkmarkc\checkmark$\checkmark)\checkmark$\checkmark^\checkmark\\checkmarkc\checkmarki\checkmarkr\checkmarkc\checkmark$\checkmark^\checkmark$\checkmark^\checkmark\\checkmarkc\checkmarki\checkmarkr\checkmarkc\checkmark$\checkmark3\checkmark$\checkmark^\checkmark\\checkmarkc\checkmarki\checkmarkr\checkmarkc\checkmark$\checkmark \checkmark$\checkmark^\checkmark\\checkmarkc\checkmarki\checkmarkr\checkmarkc\checkmark$\checkmark\\checkmark$\checkmark^\checkmark\\checkmarkc\checkmarki\checkmarkr\checkmarkc\checkmark$\checkmarkt\checkmark$\checkmark^\checkmark\\checkmarkc\checkmarki\checkmarkr\checkmarkc\checkmark$\checkmarki\checkmark$\checkmark^\checkmark\\checkmarkc\checkmarki\checkmarkr\checkmarkc\checkmark$\checkmarkm\checkmark$\checkmark^\checkmark\\checkmarkc\checkmarki\checkmarkr\checkmarkc\checkmark$\checkmarke\checkmark$\checkmark^\checkmark\\checkmarkc\checkmarki\checkmarkr\checkmarkc\checkmark$\checkmarks\checkmark$\checkmark^\checkmark\\checkmarkc\checkmarki\checkmarkr\checkmarkc\checkmark$\checkmark \checkmark$\checkmark^\checkmark\\checkmarkc\checkmarki\checkmarkr\checkmarkc\checkmark$\checkmark5\checkmark$\checkmark^\checkmark\\checkmarkc\checkmarki\checkmarkr\checkmarkc\checkmark$\checkmark0\checkmark$\checkmark^\checkmark\\checkmarkc\checkmarki\checkmarkr\checkmarkc\checkmark$\checkmark \checkmark$\checkmark^\checkmark\\checkmarkc\checkmarki\checkmarkr\checkmarkc\checkmark$\checkmark=\checkmark$\checkmark^\checkmark\\checkmarkc\checkmarki\checkmarkr\checkmarkc\checkmark$\checkmark \checkmark$\checkmark^\checkmark\\checkmarkc\checkmarki\checkmarkr\checkmarkc\checkmark$\checkmark1\checkmark$\checkmark^\checkmark\\checkmarkc\checkmarki\checkmarkr\checkmarkc\checkmark$\checkmark\\checkmark$\checkmark^\checkmark\\checkmarkc\checkmarki\checkmarkr\checkmarkc\checkmark$\checkmark,\checkmark$\checkmark^\checkmark\\checkmarkc\checkmarki\checkmarkr\checkmarkc\checkmark$\checkmark2\checkmark$\checkmark^\checkmark\\checkmarkc\checkmarki\checkmarkr\checkmarkc\checkmark$\checkmark2\checkmark$\checkmark^\checkmark\\checkmarkc\checkmarki\checkmarkr\checkmarkc\checkmark$\checkmark5\checkmark$\checkmark^\checkmark\\checkmarkc\checkmarki\checkmarkr\checkmarkc\checkmark$\checkmark\\checkmark$\checkmark^\checkmark\\checkmarkc\checkmarki\checkmarkr\checkmarkc\checkmark$\checkmark,\checkmark$\checkmark^\checkmark\\checkmarkc\checkmarki\checkmarkr\checkmarkc\checkmark$\checkmark5\checkmark$\checkmark^\checkmark\\checkmarkc\checkmarki\checkmarkr\checkmarkc\checkmark$\checkmark0\checkmark$\checkmark^\checkmark\\checkmarkc\checkmarki\checkmarkr\checkmarkc\checkmark$\checkmark0\checkmark$\checkmark^\checkmark\\checkmarkc\checkmarki\checkmarkr\checkmarkc\checkmark$\checkmark \checkmark$\checkmark^\checkmark\\checkmarkc\checkmarki\checkmarkr\checkmarkc\checkmark$\checkmark\\checkmark$\checkmark^\checkmark\\checkmarkc\checkmarki\checkmarkr\checkmarkc\checkmark$\checkmarkt\checkmark$\checkmark^\checkmark\\checkmarkc\checkmarki\checkmarkr\checkmarkc\checkmark$\checkmarki\checkmark$\checkmark^\checkmark\\checkmarkc\checkmarki\checkmarkr\checkmarkc\checkmark$\checkmarkm\checkmark$\checkmark^\checkmark\\checkmarkc\checkmarki\checkmarkr\checkmarkc\checkmark$\checkmarke\checkmark$\checkmark^\checkmark\\checkmarkc\checkmarki\checkmarkr\checkmarkc\checkmark$\checkmarks\checkmark$\checkmark^\checkmark\\checkmarkc\checkmarki\checkmarkr\checkmarkc\checkmark$\checkmark \checkmark$\checkmark^\checkmark\\checkmarkc\checkmarki\checkmarkr\checkmarkc\checkmark$\checkmark5\checkmark$\checkmark^\checkmark\\checkmarkc\checkmarki\checkmarkr\checkmarkc\checkmark$\checkmark0\checkmark$\checkmark^\checkmark\\checkmarkc\checkmarki\checkmarkr\checkmarkc\checkmark$\checkmark \checkmark$\checkmark^\checkmark\\checkmarkc\checkmarki\checkmarkr\checkmarkc\checkmark$\checkmark\\checkmark$\checkmark^\checkmark\\checkmarkc\checkmarki\checkmarkr\checkmarkc\checkmark$\checkmarka\checkmark$\checkmark^\checkmark\\checkmarkc\checkmarki\checkmarkr\checkmarkc\checkmark$\checkmarkp\checkmark$\checkmark^\checkmark\\checkmarkc\checkmarki\checkmarkr\checkmarkc\checkmark$\checkmarkp\checkmark$\checkmark^\checkmark\\checkmarkc\checkmarki\checkmarkr\checkmarkc\checkmark$\checkmarkr\checkmark$\checkmark^\checkmark\\checkmarkc\checkmarki\checkmarkr\checkmarkc\checkmark$\checkmarko\checkmark$\checkmark^\checkmark\\checkmarkc\checkmarki\checkmarkr\checkmarkc\checkmark$\checkmarkx\checkmark$\checkmark^\checkmark\\checkmarkc\checkmarki\checkmarkr\checkmarkc\checkmark$\checkmark \checkmark$\checkmark^\checkmark\\checkmarkc\checkmarki\checkmarkr\checkmarkc\checkmark$\checkmark\\checkmark$\checkmark^\checkmark\\checkmarkc\checkmarki\checkmarkr\checkmarkc\checkmark$\checkmarkb\checkmark$\checkmark^\checkmark\\checkmarkc\checkmarki\checkmarkr\checkmarkc\checkmark$\checkmarko\checkmark$\checkmark^\checkmark\\checkmarkc\checkmarki\checkmarkr\checkmarkc\checkmark$\checkmarkx\checkmark$\checkmark^\checkmark\\checkmarkc\checkmarki\checkmarkr\checkmarkc\checkmark$\checkmarke\checkmark$\checkmark^\checkmark\\checkmarkc\checkmarki\checkmarkr\checkmarkc\checkmark$\checkmarkd\checkmark$\checkmark^\checkmark\\checkmarkc\checkmarki\checkmarkr\checkmarkc\checkmark$\checkmark{\checkmark$\checkmark^\checkmark\\checkmarkc\checkmarki\checkmarkr\checkmarkc\checkmark$\checkmark6\checkmark$\checkmark^\checkmark\\checkmarkc\checkmarki\checkmarkr\checkmarkc\checkmark$\checkmark{\checkmark$\checkmark^\checkmark\\checkmarkc\checkmarki\checkmarkr\checkmarkc\checkmark$\checkmark,\checkmark$\checkmark^\checkmark\\checkmarkc\checkmarki\checkmarkr\checkmarkc\checkmark$\checkmark}\checkmark$\checkmark^\checkmark\\checkmarkc\checkmarki\checkmarkr\checkmarkc\checkmark$\checkmark1\checkmark$\checkmark^\checkmark\\checkmarkc\checkmarki\checkmarkr\checkmarkc\checkmark$\checkmark \checkmark$\checkmark^\checkmark\\checkmarkc\checkmarki\checkmarkr\checkmarkc\checkmark$\checkmark\\checkmark$\checkmark^\checkmark\\checkmarkc\checkmarki\checkmarkr\checkmarkc\checkmark$\checkmarkt\checkmark$\checkmark^\checkmark\\checkmarkc\checkmarki\checkmarkr\checkmarkc\checkmark$\checkmarki\checkmark$\checkmark^\checkmark\\checkmarkc\checkmarki\checkmarkr\checkmarkc\checkmark$\checkmarkm\checkmark$\checkmark^\checkmark\\checkmarkc\checkmarki\checkmarkr\checkmarkc\checkmark$\checkmarke\checkmark$\checkmark^\checkmark\\checkmarkc\checkmarki\checkmarkr\checkmarkc\checkmark$\checkmarks\checkmark$\checkmark^\checkmark\\checkmarkc\checkmarki\checkmarkr\checkmarkc\checkmark$\checkmark \checkmark$\checkmark^\checkmark\\checkmarkc\checkmarki\checkmarkr\checkmarkc\checkmark$\checkmark1\checkmark$\checkmark^\checkmark\\checkmarkc\checkmarki\checkmarkr\checkmarkc\checkmark$\checkmark0\checkmark$\checkmark^\checkmark\\checkmarkc\checkmarki\checkmarkr\checkmarkc\checkmark$\checkmark^\checkmark$\checkmark^\checkmark\\checkmarkc\checkmarki\checkmarkr\checkmarkc\checkmark$\checkmark7\checkmark$\checkmark^\checkmark\\checkmarkc\checkmarki\checkmarkr\checkmarkc\checkmark$\checkmark \checkmark$\checkmark^\checkmark\\checkmarkc\checkmarki\checkmarkr\checkmarkc\checkmark$\checkmark\\checkmark$\checkmark^\checkmark\\checkmarkc\checkmarki\checkmarkr\checkmarkc\checkmark$\checkmarkt\checkmark$\checkmark^\checkmark\\checkmarkc\checkmarki\checkmarkr\checkmarkc\checkmark$\checkmarke\checkmark$\checkmark^\checkmark\\checkmarkc\checkmarki\checkmarkr\checkmarkc\checkmark$\checkmarkx\checkmark$\checkmark^\checkmark\\checkmarkc\checkmarki\checkmarkr\checkmarkc\checkmark$\checkmarkt\checkmark$\checkmark^\checkmark\\checkmarkc\checkmarki\checkmarkr\checkmarkc\checkmark$\checkmark{\checkmark$\checkmark^\checkmark\\checkmarkc\checkmarki\checkmarkr\checkmarkc\checkmark$\checkmark \checkmark$\checkmark^\checkmark\\checkmarkc\checkmarki\checkmarkr\checkmarkc\checkmark$\checkmarkk\checkmark$\checkmark^\checkmark\\checkmarkc\checkmarki\checkmarkr\checkmarkc\checkmark$\checkmarkm\checkmark$\checkmark^\checkmark\\checkmarkc\checkmarki\checkmarkr\checkmarkc\checkmark$\checkmark}\checkmark$\checkmark^\checkmark\\checkmarkc\checkmarki\checkmarkr\checkmarkc\checkmark$\checkmark}\checkmark$\checkmark^\checkmark\\checkmarkc\checkmarki\checkmarkr\checkmarkc\checkmark$\checkmark
\checkmark$\checkmark^\checkmark\\checkmarkc\checkmarki\checkmarkr\checkmarkc\checkmark$\checkmark \checkmark$\checkmark^\checkmark\\checkmarkc\checkmarki\checkmarkr\checkmarkc\checkmark$\checkmark \checkmark$\checkmark^\checkmark\\checkmarkc\checkmarki\checkmarkr\checkmarkc\checkmark$\checkmark \checkmark$\checkmark^\checkmark\\checkmarkc\checkmarki\checkmarkr\checkmarkc\checkmark$\checkmark \checkmark$\checkmark^\checkmark\\checkmarkc\checkmarki\checkmarkr\checkmarkc\checkmark$\checkmark\\checkmark$\checkmark^\checkmark\\checkmarkc\checkmarki\checkmarkr\checkmarkc\checkmark$\checkmarkl\checkmark$\checkmark^\checkmark\\checkmarkc\checkmarki\checkmarkr\checkmarkc\checkmark$\checkmarka\checkmark$\checkmark^\checkmark\\checkmarkc\checkmarki\checkmarkr\checkmarkc\checkmark$\checkmarkb\checkmark$\checkmark^\checkmark\\checkmarkc\checkmarki\checkmarkr\checkmarkc\checkmark$\checkmarke\checkmark$\checkmark^\checkmark\\checkmarkc\checkmarki\checkmarkr\checkmarkc\checkmark$\checkmarkl\checkmark$\checkmark^\checkmark\\checkmarkc\checkmarki\checkmarkr\checkmarkc\checkmark$\checkmark{\checkmark$\checkmark^\checkmark\\checkmarkc\checkmarki\checkmarkr\checkmarkc\checkmark$\checkmarke\checkmark$\checkmark^\checkmark\\checkmarkc\checkmarki\checkmarkr\checkmarkc\checkmark$\checkmarkq\checkmark$\checkmark^\checkmark\\checkmarkc\checkmarki\checkmarkr\checkmarkc\checkmark$\checkmark:\checkmark$\checkmark^\checkmark\\checkmarkc\checkmarki\checkmarkr\checkmarkc\checkmark$\checkmarkl\checkmark$\checkmark^\checkmark\\checkmarkc\checkmarki\checkmarkr\checkmarkc\checkmark$\checkmark1\checkmark$\checkmark^\checkmark\\checkmarkc\checkmarki\checkmarkr\checkmarkc\checkmark$\checkmark0\checkmark$\checkmark^\checkmark\\checkmarkc\checkmarki\checkmarkr\checkmarkc\checkmark$\checkmark}\checkmark$\checkmark^\checkmark\\checkmarkc\checkmarki\checkmarkr\checkmarkc\checkmark$\checkmark
\checkmark$\checkmark^\checkmark\\checkmarkc\checkmarki\checkmarkr\checkmarkc\checkmark$\checkmark\\checkmark$\checkmark^\checkmark\\checkmarkc\checkmarki\checkmarkr\checkmarkc\checkmark$\checkmarke\checkmark$\checkmark^\checkmark\\checkmarkc\checkmarki\checkmarkr\checkmarkc\checkmark$\checkmarkn\checkmark$\checkmark^\checkmark\\checkmarkc\checkmarki\checkmarkr\checkmarkc\checkmark$\checkmarkd\checkmark$\checkmark^\checkmark\\checkmarkc\checkmarki\checkmarkr\checkmarkc\checkmark$\checkmark{\checkmark$\checkmark^\checkmark\\checkmarkc\checkmarki\checkmarkr\checkmarkc\checkmark$\checkmarke\checkmark$\checkmark^\checkmark\\checkmarkc\checkmarki\checkmarkr\checkmarkc\checkmark$\checkmarkq\checkmark$\checkmark^\checkmark\\checkmarkc\checkmarki\checkmarkr\checkmarkc\checkmark$\checkmarku\checkmark$\checkmark^\checkmark\\checkmarkc\checkmarki\checkmarkr\checkmarkc\checkmark$\checkmarka\checkmark$\checkmark^\checkmark\\checkmarkc\checkmarki\checkmarkr\checkmarkc\checkmark$\checkmarkt\checkmark$\checkmark^\checkmark\\checkmarkc\checkmarki\checkmarkr\checkmarkc\checkmark$\checkmarki\checkmark$\checkmark^\checkmark\\checkmarkc\checkmarki\checkmarkr\checkmarkc\checkmark$\checkmarko\checkmark$\checkmark^\checkmark\\checkmarkc\checkmarki\checkmarkr\checkmarkc\checkmark$\checkmarkn\checkmark$\checkmark^\checkmark\\checkmarkc\checkmarki\checkmarkr\checkmarkc\checkmark$\checkmark}\checkmark$\checkmark^\checkmark\\checkmarkc\checkmarki\checkmarkr\checkmarkc\checkmark$\checkmark
\checkmark$\checkmark^\checkmark\\checkmarkc\checkmarki\checkmarkr\checkmarkc\checkmark$\checkmark
\checkmark$\checkmark^\checkmark\\checkmarkc\checkmarki\checkmarkr\checkmarkc\checkmark$\checkmark\\checkmark$\checkmark^\checkmark\\checkmarkc\checkmarki\checkmarkr\checkmarkc\checkmark$\checkmarks\checkmark$\checkmark^\checkmark\\checkmarkc\checkmarki\checkmarkr\checkmarkc\checkmark$\checkmarku\checkmark$\checkmark^\checkmark\\checkmarkc\checkmarki\checkmarkr\checkmarkc\checkmark$\checkmarkb\checkmark$\checkmark^\checkmark\\checkmarkc\checkmarki\checkmarkr\checkmarkc\checkmark$\checkmarks\checkmark$\checkmark^\checkmark\\checkmarkc\checkmarki\checkmarkr\checkmarkc\checkmark$\checkmarke\checkmark$\checkmark^\checkmark\\checkmarkc\checkmarki\checkmarkr\checkmarkc\checkmark$\checkmarkc\checkmark$\checkmark^\checkmark\\checkmarkc\checkmarki\checkmarkr\checkmarkc\checkmark$\checkmarkt\checkmark$\checkmark^\checkmark\\checkmarkc\checkmarki\checkmarkr\checkmarkc\checkmark$\checkmarki\checkmark$\checkmark^\checkmark\\checkmarkc\checkmarki\checkmarkr\checkmarkc\checkmark$\checkmarko\checkmark$\checkmark^\checkmark\\checkmarkc\checkmarki\checkmarkr\checkmarkc\checkmark$\checkmarkn\checkmark$\checkmark^\checkmark\\checkmarkc\checkmarki\checkmarkr\checkmarkc\checkmark$\checkmark{\checkmark$\checkmark^\checkmark\\checkmarkc\checkmarki\checkmarkr\checkmarkc\checkmark$\checkmarkD\checkmark$\checkmark^\checkmark\\checkmarkc\checkmarki\checkmarkr\checkmarkc\checkmark$\checkmarku\checkmark$\checkmark^\checkmark\\checkmarkc\checkmarki\checkmarkr\checkmarkc\checkmark$\checkmarkr\checkmark$\checkmark^\checkmark\\checkmarkc\checkmarki\checkmarkr\checkmarkc\checkmark$\checkmarké\checkmark$\checkmark^\checkmark\\checkmarkc\checkmarki\checkmarkr\checkmarkc\checkmark$\checkmarke\checkmark$\checkmark^\checkmark\\checkmarkc\checkmarki\checkmarkr\checkmarkc\checkmark$\checkmark \checkmark$\checkmark^\checkmark\\checkmarkc\checkmarki\checkmarkr\checkmarkc\checkmark$\checkmarkd\checkmark$\checkmark^\checkmark\\checkmarkc\checkmarki\checkmarkr\checkmarkc\checkmark$\checkmarke\checkmark$\checkmark^\checkmark\\checkmarkc\checkmarki\checkmarkr\checkmarkc\checkmark$\checkmark \checkmark$\checkmark^\checkmark\\checkmarkc\checkmarki\checkmarkr\checkmarkc\checkmark$\checkmarkv\checkmark$\checkmark^\checkmark\\checkmarkc\checkmarki\checkmarkr\checkmarkc\checkmark$\checkmarki\checkmark$\checkmark^\checkmark\\checkmarkc\checkmarki\checkmarkr\checkmarkc\checkmark$\checkmarke\checkmark$\checkmark^\checkmark\\checkmarkc\checkmarki\checkmarkr\checkmarkc\checkmark$\checkmark \checkmark$\checkmark^\checkmark\\checkmarkc\checkmarki\checkmarkr\checkmarkc\checkmark$\checkmarke\checkmark$\checkmark^\checkmark\\checkmarkc\checkmarki\checkmarkr\checkmarkc\checkmark$\checkmarkn\checkmark$\checkmark^\checkmark\\checkmarkc\checkmarki\checkmarkr\checkmarkc\checkmark$\checkmark \checkmark$\checkmark^\checkmark\\checkmarkc\checkmarki\checkmarkr\checkmarkc\checkmark$\checkmarkh\checkmark$\checkmark^\checkmark\\checkmarkc\checkmarki\checkmarkr\checkmarkc\checkmark$\checkmarke\checkmark$\checkmark^\checkmark\\checkmarkc\checkmarki\checkmarkr\checkmarkc\checkmark$\checkmarku\checkmark$\checkmark^\checkmark\\checkmarkc\checkmarki\checkmarkr\checkmarkc\checkmark$\checkmarkr\checkmark$\checkmark^\checkmark\\checkmarkc\checkmarki\checkmarkr\checkmarkc\checkmark$\checkmarke\checkmark$\checkmark^\checkmark\\checkmarkc\checkmarki\checkmarkr\checkmarkc\checkmark$\checkmarks\checkmark$\checkmark^\checkmark\\checkmarkc\checkmarki\checkmarkr\checkmarkc\checkmark$\checkmark \checkmark$\checkmark^\checkmark\\checkmarkc\checkmarki\checkmarkr\checkmarkc\checkmark$\checkmarkd\checkmark$\checkmark^\checkmark\\checkmarkc\checkmarki\checkmarkr\checkmarkc\checkmark$\checkmarke\checkmark$\checkmark^\checkmark\\checkmarkc\checkmarki\checkmarkr\checkmarkc\checkmark$\checkmark \checkmark$\checkmark^\checkmark\\checkmarkc\checkmarki\checkmarkr\checkmarkc\checkmark$\checkmarks\checkmark$\checkmark^\checkmark\\checkmarkc\checkmarki\checkmarkr\checkmarkc\checkmark$\checkmarke\checkmark$\checkmark^\checkmark\\checkmarkc\checkmarki\checkmarkr\checkmarkc\checkmark$\checkmarkr\checkmark$\checkmark^\checkmark\\checkmarkc\checkmarki\checkmarkr\checkmarkc\checkmark$\checkmarkv\checkmark$\checkmark^\checkmark\\checkmarkc\checkmarki\checkmarkr\checkmarkc\checkmark$\checkmarki\checkmark$\checkmark^\checkmark\\checkmarkc\checkmarki\checkmarkr\checkmarkc\checkmark$\checkmarkc\checkmark$\checkmark^\checkmark\\checkmarkc\checkmarki\checkmarkr\checkmarkc\checkmark$\checkmarke\checkmark$\checkmark^\checkmark\\checkmarkc\checkmarki\checkmarkr\checkmarkc\checkmark$\checkmark}\checkmark$\checkmark^\checkmark\\checkmarkc\checkmarki\checkmarkr\checkmarkc\checkmark$\checkmark
\checkmark$\checkmark^\checkmark\\checkmarkc\checkmarki\checkmarkr\checkmarkc\checkmark$\checkmarkE\checkmark$\checkmark^\checkmark\\checkmarkc\checkmarki\checkmarkr\checkmarkc\checkmark$\checkmarkn\checkmark$\checkmark^\checkmark\\checkmarkc\checkmarki\checkmarkr\checkmarkc\checkmark$\checkmark \checkmark$\checkmark^\checkmark\\checkmarkc\checkmarki\checkmarkr\checkmarkc\checkmark$\checkmarks\checkmark$\checkmark^\checkmark\\checkmarkc\checkmarki\checkmarkr\checkmarkc\checkmark$\checkmarku\checkmark$\checkmark^\checkmark\\checkmarkc\checkmarki\checkmarkr\checkmarkc\checkmark$\checkmarkp\checkmark$\checkmark^\checkmark\\checkmarkc\checkmarki\checkmarkr\checkmarkc\checkmark$\checkmarkp\checkmark$\checkmark^\checkmark\\checkmarkc\checkmarki\checkmarkr\checkmarkc\checkmark$\checkmarko\checkmark$\checkmark^\checkmark\\checkmarkc\checkmarki\checkmarkr\checkmarkc\checkmark$\checkmarks\checkmark$\checkmark^\checkmark\\checkmarkc\checkmarki\checkmarkr\checkmarkc\checkmark$\checkmarka\checkmark$\checkmark^\checkmark\\checkmarkc\checkmarki\checkmarkr\checkmarkc\checkmark$\checkmarkn\checkmark$\checkmark^\checkmark\\checkmarkc\checkmarki\checkmarkr\checkmarkc\checkmark$\checkmarkt\checkmark$\checkmark^\checkmark\\checkmarkc\checkmarki\checkmarkr\checkmarkc\checkmark$\checkmark \checkmark$\checkmark^\checkmark\\checkmarkc\checkmarki\checkmarkr\checkmarkc\checkmark$\checkmarku\checkmark$\checkmark^\checkmark\\checkmarkc\checkmarki\checkmarkr\checkmarkc\checkmark$\checkmarkn\checkmark$\checkmark^\checkmark\\checkmarkc\checkmarki\checkmarkr\checkmarkc\checkmark$\checkmarke\checkmark$\checkmark^\checkmark\\checkmarkc\checkmarki\checkmarkr\checkmarkc\checkmark$\checkmark \checkmark$\checkmark^\checkmark\\checkmarkc\checkmarki\checkmarkr\checkmarkc\checkmark$\checkmarku\checkmark$\checkmark^\checkmark\\checkmarkc\checkmarki\checkmarkr\checkmarkc\checkmark$\checkmarkt\checkmark$\checkmark^\checkmark\\checkmarkc\checkmarki\checkmarkr\checkmarkc\checkmark$\checkmarki\checkmark$\checkmark^\checkmark\\checkmarkc\checkmarki\checkmarkr\checkmarkc\checkmark$\checkmarkl\checkmark$\checkmark^\checkmark\\checkmarkc\checkmarki\checkmarkr\checkmarkc\checkmark$\checkmarki\checkmark$\checkmark^\checkmark\\checkmarkc\checkmarki\checkmarkr\checkmarkc\checkmark$\checkmarks\checkmark$\checkmark^\checkmark\\checkmarkc\checkmarki\checkmarkr\checkmarkc\checkmark$\checkmarka\checkmark$\checkmark^\checkmark\\checkmarkc\checkmarki\checkmarkr\checkmarkc\checkmark$\checkmarkt\checkmark$\checkmark^\checkmark\\checkmarkc\checkmarki\checkmarkr\checkmarkc\checkmark$\checkmarki\checkmark$\checkmark^\checkmark\\checkmarkc\checkmarki\checkmarkr\checkmarkc\checkmark$\checkmarko\checkmark$\checkmark^\checkmark\\checkmarkc\checkmarki\checkmarkr\checkmarkc\checkmark$\checkmarkn\checkmark$\checkmark^\checkmark\\checkmarkc\checkmarki\checkmarkr\checkmarkc\checkmark$\checkmark \checkmark$\checkmark^\checkmark\\checkmarkc\checkmarki\checkmarkr\checkmarkc\checkmark$\checkmarki\checkmark$\checkmark^\checkmark\\checkmarkc\checkmarki\checkmarkr\checkmarkc\checkmark$\checkmarkn\checkmark$\checkmark^\checkmark\\checkmarkc\checkmarki\checkmarkr\checkmarkc\checkmark$\checkmarkt\checkmark$\checkmark^\checkmark\\checkmarkc\checkmarki\checkmarkr\checkmarkc\checkmark$\checkmarke\checkmark$\checkmark^\checkmark\\checkmarkc\checkmarki\checkmarkr\checkmarkc\checkmark$\checkmarkn\checkmark$\checkmark^\checkmark\\checkmarkc\checkmarki\checkmarkr\checkmarkc\checkmark$\checkmarks\checkmark$\checkmark^\checkmark\\checkmarkc\checkmarki\checkmarkr\checkmarkc\checkmark$\checkmarki\checkmark$\checkmark^\checkmark\\checkmarkc\checkmarki\checkmarkr\checkmarkc\checkmark$\checkmarkv\checkmark$\checkmark^\checkmark\\checkmarkc\checkmarki\checkmarkr\checkmarkc\checkmark$\checkmarke\checkmark$\checkmark^\checkmark\\checkmarkc\checkmarki\checkmarkr\checkmarkc\checkmark$\checkmark \checkmark$\checkmark^\checkmark\\checkmarkc\checkmarki\checkmarkr\checkmarkc\checkmark$\checkmarkd\checkmark$\checkmark^\checkmark\\checkmarkc\checkmarki\checkmarkr\checkmarkc\checkmark$\checkmarke\checkmark$\checkmark^\checkmark\\checkmarkc\checkmarki\checkmarkr\checkmarkc\checkmark$\checkmark \checkmark$\checkmark^\checkmark\\checkmarkc\checkmarki\checkmarkr\checkmarkc\checkmark$\checkmark$\checkmark$\checkmark^\checkmark\\checkmarkc\checkmarki\checkmarkr\checkmarkc\checkmark$\checkmarkv\checkmark$\checkmark^\checkmark\\checkmarkc\checkmarki\checkmarkr\checkmarkc\checkmark$\checkmark_\checkmark$\checkmark^\checkmark\\checkmarkc\checkmarki\checkmarkr\checkmarkc\checkmark$\checkmark{\checkmark$\checkmark^\checkmark\\checkmarkc\checkmarki\checkmarkr\checkmarkc\checkmark$\checkmarkm\checkmark$\checkmark^\checkmark\\checkmarkc\checkmarki\checkmarkr\checkmarkc\checkmark$\checkmarko\checkmark$\checkmark^\checkmark\\checkmarkc\checkmarki\checkmarkr\checkmarkc\checkmark$\checkmarky\checkmark$\checkmark^\checkmark\\checkmarkc\checkmarki\checkmarkr\checkmarkc\checkmark$\checkmark}\checkmark$\checkmark^\checkmark\\checkmarkc\checkmarki\checkmarkr\checkmarkc\checkmark$\checkmark \checkmark$\checkmark^\checkmark\\checkmarkc\checkmarki\checkmarkr\checkmarkc\checkmark$\checkmark=\checkmark$\checkmark^\checkmark\\checkmarkc\checkmarki\checkmarkr\checkmarkc\checkmark$\checkmark \checkmark$\checkmark^\checkmark\\checkmarkc\checkmarki\checkmarkr\checkmarkc\checkmark$\checkmark1\checkmark$\checkmark^\checkmark\\checkmarkc\checkmarki\checkmarkr\checkmarkc\checkmark$\checkmark\\checkmark$\checkmark^\checkmark\\checkmarkc\checkmarki\checkmarkr\checkmarkc\checkmark$\checkmark,\checkmark$\checkmark^\checkmark\\checkmarkc\checkmarki\checkmarkr\checkmarkc\checkmark$\checkmark0\checkmark$\checkmark^\checkmark\\checkmarkc\checkmarki\checkmarkr\checkmarkc\checkmark$\checkmark0\checkmark$\checkmark^\checkmark\\checkmarkc\checkmarki\checkmarkr\checkmarkc\checkmark$\checkmark0\checkmark$\checkmark^\checkmark\\checkmarkc\checkmarki\checkmarkr\checkmarkc\checkmark$\checkmark$\checkmark$\checkmark^\checkmark\\checkmarkc\checkmarki\checkmarkr\checkmarkc\checkmark$\checkmark \checkmark$\checkmark^\checkmark\\checkmarkc\checkmarki\checkmarkr\checkmarkc\checkmark$\checkmarkm\checkmark$\checkmark^\checkmark\\checkmarkc\checkmarki\checkmarkr\checkmarkc\checkmark$\checkmarkm\checkmark$\checkmark^\checkmark\\checkmarkc\checkmarki\checkmarkr\checkmarkc\checkmark$\checkmark/\checkmark$\checkmark^\checkmark\\checkmarkc\checkmarki\checkmarkr\checkmarkc\checkmark$\checkmarkm\checkmark$\checkmark^\checkmark\\checkmarkc\checkmarki\checkmarkr\checkmarkc\checkmark$\checkmarki\checkmark$\checkmark^\checkmark\\checkmarkc\checkmarki\checkmarkr\checkmarkc\checkmark$\checkmarkn\checkmark$\checkmark^\checkmark\\checkmarkc\checkmarki\checkmarkr\checkmarkc\checkmark$\checkmark \checkmark$\checkmark^\checkmark\\checkmarkc\checkmarki\checkmarkr\checkmarkc\checkmark$\checkmark=\checkmark$\checkmark^\checkmark\\checkmarkc\checkmarki\checkmarkr\checkmarkc\checkmark$\checkmark \checkmark$\checkmark^\checkmark\\checkmarkc\checkmarki\checkmarkr\checkmarkc\checkmark$\checkmark$\checkmark$\checkmark^\checkmark\\checkmarkc\checkmarki\checkmarkr\checkmarkc\checkmark$\checkmark0\checkmark$\checkmark^\checkmark\\checkmarkc\checkmarki\checkmarkr\checkmarkc\checkmark$\checkmark{\checkmark$\checkmark^\checkmark\\checkmarkc\checkmarki\checkmarkr\checkmarkc\checkmark$\checkmark,\checkmark$\checkmark^\checkmark\\checkmarkc\checkmarki\checkmarkr\checkmarkc\checkmark$\checkmark}\checkmark$\checkmark^\checkmark\\checkmarkc\checkmarki\checkmarkr\checkmarkc\checkmark$\checkmark0\checkmark$\checkmark^\checkmark\\checkmarkc\checkmarki\checkmarkr\checkmarkc\checkmark$\checkmark6\checkmark$\checkmark^\checkmark\\checkmarkc\checkmarki\checkmarkr\checkmarkc\checkmark$\checkmark$\checkmark$\checkmark^\checkmark\\checkmarkc\checkmarki\checkmarkr\checkmarkc\checkmark$\checkmark \checkmark$\checkmark^\checkmark\\checkmarkc\checkmarki\checkmarkr\checkmarkc\checkmark$\checkmarkk\checkmark$\checkmark^\checkmark\\checkmarkc\checkmarki\checkmarkr\checkmarkc\checkmark$\checkmarkm\checkmark$\checkmark^\checkmark\\checkmarkc\checkmarki\checkmarkr\checkmarkc\checkmark$\checkmark/\checkmark$\checkmark^\checkmark\\checkmarkc\checkmarki\checkmarkr\checkmarkc\checkmark$\checkmarkh\checkmark$\checkmark^\checkmark\\checkmarkc\checkmarki\checkmarkr\checkmarkc\checkmark$\checkmark,\checkmark$\checkmark^\checkmark\\checkmarkc\checkmarki\checkmarkr\checkmarkc\checkmark$\checkmark \checkmark$\checkmark^\checkmark\\checkmarkc\checkmarki\checkmarkr\checkmarkc\checkmark$\checkmarkp\checkmark$\checkmark^\checkmark\\checkmarkc\checkmarki\checkmarkr\checkmarkc\checkmark$\checkmarke\checkmark$\checkmark^\checkmark\\checkmarkc\checkmarki\checkmarkr\checkmarkc\checkmark$\checkmarkn\checkmark$\checkmark^\checkmark\\checkmarkc\checkmarki\checkmarkr\checkmarkc\checkmark$\checkmarkd\checkmark$\checkmark^\checkmark\\checkmarkc\checkmarki\checkmarkr\checkmarkc\checkmark$\checkmarka\checkmark$\checkmark^\checkmark\\checkmarkc\checkmarki\checkmarkr\checkmarkc\checkmark$\checkmarkn\checkmark$\checkmark^\checkmark\\checkmarkc\checkmarki\checkmarkr\checkmarkc\checkmark$\checkmarkt\checkmark$\checkmark^\checkmark\\checkmarkc\checkmarki\checkmarkr\checkmarkc\checkmark$\checkmark \checkmark$\checkmark^\checkmark\\checkmarkc\checkmarki\checkmarkr\checkmarkc\checkmark$\checkmark8\checkmark$\checkmark^\checkmark\\checkmarkc\checkmarki\checkmarkr\checkmarkc\checkmark$\checkmark \checkmark$\checkmark^\checkmark\\checkmarkc\checkmarki\checkmarkr\checkmarkc\checkmark$\checkmarkh\checkmark$\checkmark^\checkmark\\checkmarkc\checkmarki\checkmarkr\checkmarkc\checkmark$\checkmarke\checkmark$\checkmark^\checkmark\\checkmarkc\checkmarki\checkmarkr\checkmarkc\checkmark$\checkmarku\checkmark$\checkmark^\checkmark\\checkmarkc\checkmarki\checkmarkr\checkmarkc\checkmark$\checkmarkr\checkmark$\checkmark^\checkmark\\checkmarkc\checkmarki\checkmarkr\checkmarkc\checkmark$\checkmarke\checkmark$\checkmark^\checkmark\\checkmarkc\checkmarki\checkmarkr\checkmarkc\checkmark$\checkmarks\checkmark$\checkmark^\checkmark\\checkmarkc\checkmarki\checkmarkr\checkmarkc\checkmark$\checkmark \checkmark$\checkmark^\checkmark\\checkmarkc\checkmarki\checkmarkr\checkmarkc\checkmark$\checkmarkp\checkmark$\checkmark^\checkmark\\checkmarkc\checkmarki\checkmarkr\checkmarkc\checkmark$\checkmarka\checkmark$\checkmark^\checkmark\\checkmarkc\checkmarki\checkmarkr\checkmarkc\checkmark$\checkmarkr\checkmark$\checkmark^\checkmark\\checkmarkc\checkmarki\checkmarkr\checkmarkc\checkmark$\checkmark \checkmark$\checkmark^\checkmark\\checkmarkc\checkmarki\checkmarkr\checkmarkc\checkmark$\checkmarkj\checkmark$\checkmark^\checkmark\\checkmarkc\checkmarki\checkmarkr\checkmarkc\checkmark$\checkmarko\checkmark$\checkmark^\checkmark\\checkmarkc\checkmarki\checkmarkr\checkmarkc\checkmark$\checkmarku\checkmark$\checkmark^\checkmark\\checkmarkc\checkmarki\checkmarkr\checkmarkc\checkmark$\checkmarkr\checkmark$\checkmark^\checkmark\\checkmarkc\checkmarki\checkmarkr\checkmarkc\checkmark$\checkmark,\checkmark$\checkmark^\checkmark\\checkmarkc\checkmarki\checkmarkr\checkmarkc\checkmark$\checkmark \checkmark$\checkmark^\checkmark\\checkmarkc\checkmarki\checkmarkr\checkmarkc\checkmark$\checkmark2\checkmark$\checkmark^\checkmark\\checkmarkc\checkmarki\checkmarkr\checkmarkc\checkmark$\checkmark5\checkmark$\checkmark^\checkmark\\checkmarkc\checkmarki\checkmarkr\checkmarkc\checkmark$\checkmark0\checkmark$\checkmark^\checkmark\\checkmarkc\checkmarki\checkmarkr\checkmarkc\checkmark$\checkmark \checkmark$\checkmark^\checkmark\\checkmarkc\checkmarki\checkmarkr\checkmarkc\checkmark$\checkmarkj\checkmark$\checkmark^\checkmark\\checkmarkc\checkmarki\checkmarkr\checkmarkc\checkmark$\checkmarko\checkmark$\checkmark^\checkmark\\checkmarkc\checkmarki\checkmarkr\checkmarkc\checkmark$\checkmarku\checkmark$\checkmark^\checkmark\\checkmarkc\checkmarki\checkmarkr\checkmarkc\checkmark$\checkmarkr\checkmark$\checkmark^\checkmark\\checkmarkc\checkmarki\checkmarkr\checkmarkc\checkmark$\checkmarks\checkmark$\checkmark^\checkmark\\checkmarkc\checkmarki\checkmarkr\checkmarkc\checkmark$\checkmark \checkmark$\checkmark^\checkmark\\checkmarkc\checkmarki\checkmarkr\checkmarkc\checkmark$\checkmarkp\checkmark$\checkmark^\checkmark\\checkmarkc\checkmarki\checkmarkr\checkmarkc\checkmark$\checkmarka\checkmark$\checkmark^\checkmark\\checkmarkc\checkmarki\checkmarkr\checkmarkc\checkmark$\checkmarkr\checkmark$\checkmark^\checkmark\\checkmarkc\checkmarki\checkmarkr\checkmarkc\checkmark$\checkmark \checkmark$\checkmark^\checkmark\\checkmarkc\checkmarki\checkmarkr\checkmarkc\checkmark$\checkmarka\checkmark$\checkmark^\checkmark\\checkmarkc\checkmarki\checkmarkr\checkmarkc\checkmark$\checkmarkn\checkmark$\checkmark^\checkmark\\checkmarkc\checkmarki\checkmarkr\checkmarkc\checkmark$\checkmark \checkmark$\checkmark^\checkmark\\checkmarkc\checkmarki\checkmarkr\checkmarkc\checkmark$\checkmark:\checkmark$\checkmark^\checkmark\\checkmarkc\checkmarki\checkmarkr\checkmarkc\checkmark$\checkmark
\checkmark$\checkmark^\checkmark\\checkmarkc\checkmarki\checkmarkr\checkmarkc\checkmark$\checkmark\\checkmark$\checkmark^\checkmark\\checkmarkc\checkmarki\checkmarkr\checkmarkc\checkmark$\checkmarkb\checkmark$\checkmark^\checkmark\\checkmarkc\checkmarki\checkmarkr\checkmarkc\checkmark$\checkmarke\checkmark$\checkmark^\checkmark\\checkmarkc\checkmarki\checkmarkr\checkmarkc\checkmark$\checkmarkg\checkmark$\checkmark^\checkmark\\checkmarkc\checkmarki\checkmarkr\checkmarkc\checkmark$\checkmarki\checkmark$\checkmark^\checkmark\\checkmarkc\checkmarki\checkmarkr\checkmarkc\checkmark$\checkmarkn\checkmark$\checkmark^\checkmark\\checkmarkc\checkmarki\checkmarkr\checkmarkc\checkmark$\checkmark{\checkmark$\checkmark^\checkmark\\checkmarkc\checkmarki\checkmarkr\checkmarkc\checkmark$\checkmarke\checkmark$\checkmark^\checkmark\\checkmarkc\checkmarki\checkmarkr\checkmarkc\checkmark$\checkmarkq\checkmark$\checkmark^\checkmark\\checkmarkc\checkmarki\checkmarkr\checkmarkc\checkmark$\checkmarku\checkmark$\checkmark^\checkmark\\checkmarkc\checkmarki\checkmarkr\checkmarkc\checkmark$\checkmarka\checkmark$\checkmark^\checkmark\\checkmarkc\checkmarki\checkmarkr\checkmarkc\checkmark$\checkmarkt\checkmark$\checkmark^\checkmark\\checkmarkc\checkmarki\checkmarkr\checkmarkc\checkmark$\checkmarki\checkmark$\checkmark^\checkmark\\checkmarkc\checkmarki\checkmarkr\checkmarkc\checkmark$\checkmarko\checkmark$\checkmark^\checkmark\\checkmarkc\checkmarki\checkmarkr\checkmarkc\checkmark$\checkmarkn\checkmark$\checkmark^\checkmark\\checkmarkc\checkmarki\checkmarkr\checkmarkc\checkmark$\checkmark}\checkmark$\checkmark^\checkmark\\checkmarkc\checkmarki\checkmarkr\checkmarkc\checkmark$\checkmark
\checkmark$\checkmark^\checkmark\\checkmarkc\checkmarki\checkmarkr\checkmarkc\checkmark$\checkmark \checkmark$\checkmark^\checkmark\\checkmarkc\checkmarki\checkmarkr\checkmarkc\checkmark$\checkmark \checkmark$\checkmark^\checkmark\\checkmarkc\checkmarki\checkmarkr\checkmarkc\checkmark$\checkmark \checkmark$\checkmark^\checkmark\\checkmarkc\checkmarki\checkmarkr\checkmarkc\checkmark$\checkmark \checkmark$\checkmark^\checkmark\\checkmarkc\checkmarki\checkmarkr\checkmarkc\checkmark$\checkmarkL\checkmark$\checkmark^\checkmark\\checkmarkc\checkmarki\checkmarkr\checkmarkc\checkmark$\checkmark_\checkmark$\checkmark^\checkmark\\checkmarkc\checkmarki\checkmarkr\checkmarkc\checkmark$\checkmark{\checkmark$\checkmark^\checkmark\\checkmarkc\checkmarki\checkmarkr\checkmarkc\checkmark$\checkmark1\checkmark$\checkmark^\checkmark\\checkmarkc\checkmarki\checkmarkr\checkmarkc\checkmark$\checkmark0\checkmark$\checkmark^\checkmark\\checkmarkc\checkmarki\checkmarkr\checkmarkc\checkmark$\checkmarkh\checkmark$\checkmark^\checkmark\\checkmarkc\checkmarki\checkmarkr\checkmarkc\checkmark$\checkmark}\checkmark$\checkmark^\checkmark\\checkmarkc\checkmarki\checkmarkr\checkmarkc\checkmark$\checkmark \checkmark$\checkmark^\checkmark\\checkmarkc\checkmarki\checkmarkr\checkmarkc\checkmark$\checkmark=\checkmark$\checkmark^\checkmark\\checkmarkc\checkmarki\checkmarkr\checkmarkc\checkmark$\checkmark \checkmark$\checkmark^\checkmark\\checkmarkc\checkmarki\checkmarkr\checkmarkc\checkmark$\checkmark\\checkmark$\checkmark^\checkmark\\checkmarkc\checkmarki\checkmarkr\checkmarkc\checkmark$\checkmarkf\checkmark$\checkmark^\checkmark\\checkmarkc\checkmarki\checkmarkr\checkmarkc\checkmark$\checkmarkr\checkmark$\checkmark^\checkmark\\checkmarkc\checkmarki\checkmarkr\checkmarkc\checkmark$\checkmarka\checkmark$\checkmark^\checkmark\\checkmarkc\checkmarki\checkmarkr\checkmarkc\checkmark$\checkmarkc\checkmark$\checkmark^\checkmark\\checkmarkc\checkmarki\checkmarkr\checkmarkc\checkmark$\checkmark{\checkmark$\checkmark^\checkmark\\checkmarkc\checkmarki\checkmarkr\checkmarkc\checkmark$\checkmarkL\checkmark$\checkmark^\checkmark\\checkmarkc\checkmarki\checkmarkr\checkmarkc\checkmark$\checkmark_\checkmark$\checkmark^\checkmark\\checkmarkc\checkmarki\checkmarkr\checkmarkc\checkmark$\checkmark{\checkmark$\checkmark^\checkmark\\checkmarkc\checkmarki\checkmarkr\checkmarkc\checkmark$\checkmark1\checkmark$\checkmark^\checkmark\\checkmarkc\checkmarki\checkmarkr\checkmarkc\checkmark$\checkmark0\checkmark$\checkmark^\checkmark\\checkmarkc\checkmarki\checkmarkr\checkmarkc\checkmark$\checkmark}\checkmark$\checkmark^\checkmark\\checkmarkc\checkmarki\checkmarkr\checkmarkc\checkmark$\checkmark}\checkmark$\checkmark^\checkmark\\checkmarkc\checkmarki\checkmarkr\checkmarkc\checkmark$\checkmark{\checkmark$\checkmark^\checkmark\\checkmarkc\checkmarki\checkmarkr\checkmarkc\checkmark$\checkmarkv\checkmark$\checkmark^\checkmark\\checkmarkc\checkmarki\checkmarkr\checkmarkc\checkmark$\checkmark_\checkmark$\checkmark^\checkmark\\checkmarkc\checkmarki\checkmarkr\checkmarkc\checkmark$\checkmark{\checkmark$\checkmark^\checkmark\\checkmarkc\checkmarki\checkmarkr\checkmarkc\checkmark$\checkmarkm\checkmark$\checkmark^\checkmark\\checkmarkc\checkmarki\checkmarkr\checkmarkc\checkmark$\checkmarko\checkmark$\checkmark^\checkmark\\checkmarkc\checkmarki\checkmarkr\checkmarkc\checkmark$\checkmarky\checkmark$\checkmark^\checkmark\\checkmarkc\checkmarki\checkmarkr\checkmarkc\checkmark$\checkmark}\checkmark$\checkmark^\checkmark\\checkmarkc\checkmarki\checkmarkr\checkmarkc\checkmark$\checkmark \checkmark$\checkmark^\checkmark\\checkmarkc\checkmarki\checkmarkr\checkmarkc\checkmark$\checkmark\\checkmark$\checkmark^\checkmark\\checkmarkc\checkmarki\checkmarkr\checkmarkc\checkmark$\checkmarkt\checkmark$\checkmark^\checkmark\\checkmarkc\checkmarki\checkmarkr\checkmarkc\checkmark$\checkmarki\checkmark$\checkmark^\checkmark\\checkmarkc\checkmarki\checkmarkr\checkmarkc\checkmark$\checkmarkm\checkmark$\checkmark^\checkmark\\checkmarkc\checkmarki\checkmarkr\checkmarkc\checkmark$\checkmarke\checkmark$\checkmark^\checkmark\\checkmarkc\checkmarki\checkmarkr\checkmarkc\checkmark$\checkmarks\checkmark$\checkmark^\checkmark\\checkmarkc\checkmarki\checkmarkr\checkmarkc\checkmark$\checkmark \checkmark$\checkmark^\checkmark\\checkmarkc\checkmarki\checkmarkr\checkmarkc\checkmark$\checkmarkT\checkmark$\checkmark^\checkmark\\checkmarkc\checkmarki\checkmarkr\checkmarkc\checkmark$\checkmark_\checkmark$\checkmark^\checkmark\\checkmarkc\checkmarki\checkmarkr\checkmarkc\checkmark$\checkmark{\checkmark$\checkmark^\checkmark\\checkmarkc\checkmarki\checkmarkr\checkmarkc\checkmark$\checkmarkj\checkmark$\checkmark^\checkmark\\checkmarkc\checkmarki\checkmarkr\checkmarkc\checkmark$\checkmarko\checkmark$\checkmark^\checkmark\\checkmarkc\checkmarki\checkmarkr\checkmarkc\checkmark$\checkmarku\checkmark$\checkmark^\checkmark\\checkmarkc\checkmarki\checkmarkr\checkmarkc\checkmark$\checkmarkr\checkmark$\checkmark^\checkmark\\checkmarkc\checkmarki\checkmarkr\checkmarkc\checkmark$\checkmark}\checkmark$\checkmark^\checkmark\\checkmarkc\checkmarki\checkmarkr\checkmarkc\checkmark$\checkmark \checkmark$\checkmark^\checkmark\\checkmarkc\checkmarki\checkmarkr\checkmarkc\checkmark$\checkmark\\checkmark$\checkmark^\checkmark\\checkmarkc\checkmarki\checkmarkr\checkmarkc\checkmark$\checkmarkt\checkmark$\checkmark^\checkmark\\checkmarkc\checkmarki\checkmarkr\checkmarkc\checkmark$\checkmarki\checkmark$\checkmark^\checkmark\\checkmarkc\checkmarki\checkmarkr\checkmarkc\checkmark$\checkmarkm\checkmark$\checkmark^\checkmark\\checkmarkc\checkmarki\checkmarkr\checkmarkc\checkmark$\checkmarke\checkmark$\checkmark^\checkmark\\checkmarkc\checkmarki\checkmarkr\checkmarkc\checkmark$\checkmarks\checkmark$\checkmark^\checkmark\\checkmarkc\checkmarki\checkmarkr\checkmarkc\checkmark$\checkmark \checkmark$\checkmark^\checkmark\\checkmarkc\checkmarki\checkmarkr\checkmarkc\checkmark$\checkmarkT\checkmark$\checkmark^\checkmark\\checkmarkc\checkmarki\checkmarkr\checkmarkc\checkmark$\checkmark_\checkmark$\checkmark^\checkmark\\checkmarkc\checkmarki\checkmarkr\checkmarkc\checkmark$\checkmark{\checkmark$\checkmark^\checkmark\\checkmarkc\checkmarki\checkmarkr\checkmarkc\checkmark$\checkmarka\checkmark$\checkmark^\checkmark\\checkmarkc\checkmarki\checkmarkr\checkmarkc\checkmark$\checkmarkn\checkmark$\checkmark^\checkmark\\checkmarkc\checkmarki\checkmarkr\checkmarkc\checkmark$\checkmark}\checkmark$\checkmark^\checkmark\\checkmarkc\checkmarki\checkmarkr\checkmarkc\checkmark$\checkmark}\checkmark$\checkmark^\checkmark\\checkmarkc\checkmarki\checkmarkr\checkmarkc\checkmark$\checkmark \checkmark$\checkmark^\checkmark\\checkmarkc\checkmarki\checkmarkr\checkmarkc\checkmark$\checkmark=\checkmark$\checkmark^\checkmark\\checkmarkc\checkmarki\checkmarkr\checkmarkc\checkmark$\checkmark \checkmark$\checkmark^\checkmark\\checkmarkc\checkmarki\checkmarkr\checkmarkc\checkmark$\checkmark\\checkmark$\checkmark^\checkmark\\checkmarkc\checkmarki\checkmarkr\checkmarkc\checkmark$\checkmarkf\checkmark$\checkmark^\checkmark\\checkmarkc\checkmarki\checkmarkr\checkmarkc\checkmark$\checkmarkr\checkmark$\checkmark^\checkmark\\checkmarkc\checkmarki\checkmarkr\checkmarkc\checkmark$\checkmarka\checkmark$\checkmark^\checkmark\\checkmarkc\checkmarki\checkmarkr\checkmarkc\checkmark$\checkmarkc\checkmark$\checkmark^\checkmark\\checkmarkc\checkmarki\checkmarkr\checkmarkc\checkmark$\checkmark{\checkmark$\checkmark^\checkmark\\checkmarkc\checkmarki\checkmarkr\checkmarkc\checkmark$\checkmark6\checkmark$\checkmark^\checkmark\\checkmarkc\checkmarki\checkmarkr\checkmarkc\checkmark$\checkmark{\checkmark$\checkmark^\checkmark\\checkmarkc\checkmarki\checkmarkr\checkmarkc\checkmark$\checkmark,\checkmark$\checkmark^\checkmark\\checkmarkc\checkmarki\checkmarkr\checkmarkc\checkmark$\checkmark}\checkmark$\checkmark^\checkmark\\checkmarkc\checkmarki\checkmarkr\checkmarkc\checkmark$\checkmark1\checkmark$\checkmark^\checkmark\\checkmarkc\checkmarki\checkmarkr\checkmarkc\checkmark$\checkmark \checkmark$\checkmark^\checkmark\\checkmarkc\checkmarki\checkmarkr\checkmarkc\checkmark$\checkmark\\checkmark$\checkmark^\checkmark\\checkmarkc\checkmarki\checkmarkr\checkmarkc\checkmark$\checkmarkt\checkmark$\checkmark^\checkmark\\checkmarkc\checkmarki\checkmarkr\checkmarkc\checkmark$\checkmarki\checkmark$\checkmark^\checkmark\\checkmarkc\checkmarki\checkmarkr\checkmarkc\checkmark$\checkmarkm\checkmark$\checkmark^\checkmark\\checkmarkc\checkmarki\checkmarkr\checkmarkc\checkmark$\checkmarke\checkmark$\checkmark^\checkmark\\checkmarkc\checkmarki\checkmarkr\checkmarkc\checkmark$\checkmarks\checkmark$\checkmark^\checkmark\\checkmarkc\checkmarki\checkmarkr\checkmarkc\checkmark$\checkmark \checkmark$\checkmark^\checkmark\\checkmarkc\checkmarki\checkmarkr\checkmarkc\checkmark$\checkmark1\checkmark$\checkmark^\checkmark\\checkmarkc\checkmarki\checkmarkr\checkmarkc\checkmark$\checkmark0\checkmark$\checkmark^\checkmark\\checkmarkc\checkmarki\checkmarkr\checkmarkc\checkmark$\checkmark^\checkmark$\checkmark^\checkmark\\checkmarkc\checkmarki\checkmarkr\checkmarkc\checkmark$\checkmark7\checkmark$\checkmark^\checkmark\\checkmarkc\checkmarki\checkmarkr\checkmarkc\checkmark$\checkmark}\checkmark$\checkmark^\checkmark\\checkmarkc\checkmarki\checkmarkr\checkmarkc\checkmark$\checkmark{\checkmark$\checkmark^\checkmark\\checkmarkc\checkmarki\checkmarkr\checkmarkc\checkmark$\checkmark0\checkmark$\checkmark^\checkmark\\checkmarkc\checkmarki\checkmarkr\checkmarkc\checkmark$\checkmark{\checkmark$\checkmark^\checkmark\\checkmarkc\checkmarki\checkmarkr\checkmarkc\checkmark$\checkmark,\checkmark$\checkmark^\checkmark\\checkmarkc\checkmarki\checkmarkr\checkmarkc\checkmark$\checkmark}\checkmark$\checkmark^\checkmark\\checkmarkc\checkmarki\checkmarkr\checkmarkc\checkmark$\checkmark0\checkmark$\checkmark^\checkmark\\checkmarkc\checkmarki\checkmarkr\checkmarkc\checkmark$\checkmark6\checkmark$\checkmark^\checkmark\\checkmarkc\checkmarki\checkmarkr\checkmarkc\checkmark$\checkmark \checkmark$\checkmark^\checkmark\\checkmarkc\checkmarki\checkmarkr\checkmarkc\checkmark$\checkmark\\checkmark$\checkmark^\checkmark\\checkmarkc\checkmarki\checkmarkr\checkmarkc\checkmark$\checkmarkt\checkmark$\checkmark^\checkmark\\checkmarkc\checkmarki\checkmarkr\checkmarkc\checkmark$\checkmarki\checkmark$\checkmark^\checkmark\\checkmarkc\checkmarki\checkmarkr\checkmarkc\checkmark$\checkmarkm\checkmark$\checkmark^\checkmark\\checkmarkc\checkmarki\checkmarkr\checkmarkc\checkmark$\checkmarke\checkmark$\checkmark^\checkmark\\checkmarkc\checkmarki\checkmarkr\checkmarkc\checkmark$\checkmarks\checkmark$\checkmark^\checkmark\\checkmarkc\checkmarki\checkmarkr\checkmarkc\checkmark$\checkmark \checkmark$\checkmark^\checkmark\\checkmarkc\checkmarki\checkmarkr\checkmarkc\checkmark$\checkmark8\checkmark$\checkmark^\checkmark\\checkmarkc\checkmarki\checkmarkr\checkmarkc\checkmark$\checkmark \checkmark$\checkmark^\checkmark\\checkmarkc\checkmarki\checkmarkr\checkmarkc\checkmark$\checkmark\\checkmark$\checkmark^\checkmark\\checkmarkc\checkmarki\checkmarkr\checkmarkc\checkmark$\checkmarkt\checkmark$\checkmark^\checkmark\\checkmarkc\checkmarki\checkmarkr\checkmarkc\checkmark$\checkmarki\checkmark$\checkmark^\checkmark\\checkmarkc\checkmarki\checkmarkr\checkmarkc\checkmark$\checkmarkm\checkmark$\checkmark^\checkmark\\checkmarkc\checkmarki\checkmarkr\checkmarkc\checkmark$\checkmarke\checkmark$\checkmark^\checkmark\\checkmarkc\checkmarki\checkmarkr\checkmarkc\checkmark$\checkmarks\checkmark$\checkmark^\checkmark\\checkmarkc\checkmarki\checkmarkr\checkmarkc\checkmark$\checkmark \checkmark$\checkmark^\checkmark\\checkmarkc\checkmarki\checkmarkr\checkmarkc\checkmark$\checkmark2\checkmark$\checkmark^\checkmark\\checkmarkc\checkmarki\checkmarkr\checkmarkc\checkmark$\checkmark5\checkmark$\checkmark^\checkmark\\checkmarkc\checkmarki\checkmarkr\checkmarkc\checkmark$\checkmark0\checkmark$\checkmark^\checkmark\\checkmarkc\checkmarki\checkmarkr\checkmarkc\checkmark$\checkmark}\checkmark$\checkmark^\checkmark\\checkmarkc\checkmarki\checkmarkr\checkmarkc\checkmark$\checkmark \checkmark$\checkmark^\checkmark\\checkmarkc\checkmarki\checkmarkr\checkmarkc\checkmark$\checkmark\\checkmark$\checkmark^\checkmark\\checkmarkc\checkmarki\checkmarkr\checkmarkc\checkmark$\checkmarka\checkmark$\checkmark^\checkmark\\checkmarkc\checkmarki\checkmarkr\checkmarkc\checkmark$\checkmarkp\checkmark$\checkmark^\checkmark\\checkmarkc\checkmarki\checkmarkr\checkmarkc\checkmark$\checkmarkp\checkmark$\checkmark^\checkmark\\checkmarkc\checkmarki\checkmarkr\checkmarkc\checkmark$\checkmarkr\checkmark$\checkmark^\checkmark\\checkmarkc\checkmarki\checkmarkr\checkmarkc\checkmark$\checkmarko\checkmark$\checkmark^\checkmark\\checkmarkc\checkmarki\checkmarkr\checkmarkc\checkmark$\checkmarkx\checkmark$\checkmark^\checkmark\\checkmarkc\checkmarki\checkmarkr\checkmarkc\checkmark$\checkmark \checkmark$\checkmark^\checkmark\\checkmarkc\checkmarki\checkmarkr\checkmarkc\checkmark$\checkmark\\checkmark$\checkmark^\checkmark\\checkmarkc\checkmarki\checkmarkr\checkmarkc\checkmark$\checkmarkb\checkmark$\checkmark^\checkmark\\checkmarkc\checkmarki\checkmarkr\checkmarkc\checkmark$\checkmarko\checkmark$\checkmark^\checkmark\\checkmarkc\checkmarki\checkmarkr\checkmarkc\checkmark$\checkmarkx\checkmark$\checkmark^\checkmark\\checkmarkc\checkmarki\checkmarkr\checkmarkc\checkmark$\checkmarke\checkmark$\checkmark^\checkmark\\checkmarkc\checkmarki\checkmarkr\checkmarkc\checkmark$\checkmarkd\checkmark$\checkmark^\checkmark\\checkmarkc\checkmarki\checkmarkr\checkmarkc\checkmark$\checkmark{\checkmark$\checkmark^\checkmark\\checkmarkc\checkmarki\checkmarkr\checkmarkc\checkmark$\checkmark5\checkmark$\checkmark^\checkmark\\checkmarkc\checkmarki\checkmarkr\checkmarkc\checkmark$\checkmark0\checkmark$\checkmark^\checkmark\\checkmarkc\checkmarki\checkmarkr\checkmarkc\checkmark$\checkmark8\checkmark$\checkmark^\checkmark\\checkmarkc\checkmarki\checkmarkr\checkmarkc\checkmark$\checkmark\\checkmark$\checkmark^\checkmark\\checkmarkc\checkmarki\checkmarkr\checkmarkc\checkmark$\checkmark,\checkmark$\checkmark^\checkmark\\checkmarkc\checkmarki\checkmarkr\checkmarkc\checkmark$\checkmark3\checkmark$\checkmark^\checkmark\\checkmarkc\checkmarki\checkmarkr\checkmarkc\checkmark$\checkmark3\checkmark$\checkmark^\checkmark\\checkmarkc\checkmarki\checkmarkr\checkmarkc\checkmark$\checkmark3\checkmark$\checkmark^\checkmark\\checkmarkc\checkmarki\checkmarkr\checkmarkc\checkmark$\checkmark \checkmark$\checkmark^\checkmark\\checkmarkc\checkmarki\checkmarkr\checkmarkc\checkmark$\checkmark\\checkmark$\checkmark^\checkmark\\checkmarkc\checkmarki\checkmarkr\checkmarkc\checkmark$\checkmarkt\checkmark$\checkmark^\checkmark\\checkmarkc\checkmarki\checkmarkr\checkmarkc\checkmark$\checkmarke\checkmark$\checkmark^\checkmark\\checkmarkc\checkmarki\checkmarkr\checkmarkc\checkmark$\checkmarkx\checkmark$\checkmark^\checkmark\\checkmarkc\checkmarki\checkmarkr\checkmarkc\checkmark$\checkmarkt\checkmark$\checkmark^\checkmark\\checkmarkc\checkmarki\checkmarkr\checkmarkc\checkmark$\checkmark{\checkmark$\checkmark^\checkmark\\checkmarkc\checkmarki\checkmarkr\checkmarkc\checkmark$\checkmark \checkmark$\checkmark^\checkmark\\checkmarkc\checkmarki\checkmarkr\checkmarkc\checkmark$\checkmarkh\checkmark$\checkmark^\checkmark\\checkmarkc\checkmarki\checkmarkr\checkmarkc\checkmark$\checkmarke\checkmark$\checkmark^\checkmark\\checkmarkc\checkmarki\checkmarkr\checkmarkc\checkmark$\checkmarku\checkmark$\checkmark^\checkmark\\checkmarkc\checkmarki\checkmarkr\checkmarkc\checkmark$\checkmarkr\checkmark$\checkmark^\checkmark\\checkmarkc\checkmarki\checkmarkr\checkmarkc\checkmark$\checkmarke\checkmark$\checkmark^\checkmark\\checkmarkc\checkmarki\checkmarkr\checkmarkc\checkmark$\checkmarks\checkmark$\checkmark^\checkmark\\checkmarkc\checkmarki\checkmarkr\checkmarkc\checkmark$\checkmark}\checkmark$\checkmark^\checkmark\\checkmarkc\checkmarki\checkmarkr\checkmarkc\checkmark$\checkmark \checkmark$\checkmark^\checkmark\\checkmarkc\checkmarki\checkmarkr\checkmarkc\checkmark$\checkmark\\checkmark$\checkmark^\checkmark\\checkmarkc\checkmarki\checkmarkr\checkmarkc\checkmark$\checkmarka\checkmark$\checkmark^\checkmark\\checkmarkc\checkmarki\checkmarkr\checkmarkc\checkmark$\checkmarkp\checkmark$\checkmark^\checkmark\\checkmarkc\checkmarki\checkmarkr\checkmarkc\checkmark$\checkmarkp\checkmark$\checkmark^\checkmark\\checkmarkc\checkmarki\checkmarkr\checkmarkc\checkmark$\checkmarkr\checkmark$\checkmark^\checkmark\\checkmarkc\checkmarki\checkmarkr\checkmarkc\checkmark$\checkmarko\checkmark$\checkmark^\checkmark\\checkmarkc\checkmarki\checkmarkr\checkmarkc\checkmark$\checkmarkx\checkmark$\checkmark^\checkmark\\checkmarkc\checkmarki\checkmarkr\checkmarkc\checkmark$\checkmark \checkmark$\checkmark^\checkmark\\checkmarkc\checkmarki\checkmarkr\checkmarkc\checkmark$\checkmark2\checkmark$\checkmark^\checkmark\\checkmarkc\checkmarki\checkmarkr\checkmarkc\checkmark$\checkmark5\checkmark$\checkmark^\checkmark\\checkmarkc\checkmarki\checkmarkr\checkmarkc\checkmark$\checkmark4\checkmark$\checkmark^\checkmark\\checkmarkc\checkmarki\checkmarkr\checkmarkc\checkmark$\checkmark \checkmark$\checkmark^\checkmark\\checkmarkc\checkmarki\checkmarkr\checkmarkc\checkmark$\checkmark\\checkmark$\checkmark^\checkmark\\checkmarkc\checkmarki\checkmarkr\checkmarkc\checkmark$\checkmarkt\checkmark$\checkmark^\checkmark\\checkmarkc\checkmarki\checkmarkr\checkmarkc\checkmark$\checkmarke\checkmark$\checkmark^\checkmark\\checkmarkc\checkmarki\checkmarkr\checkmarkc\checkmark$\checkmarkx\checkmark$\checkmark^\checkmark\\checkmarkc\checkmarki\checkmarkr\checkmarkc\checkmark$\checkmarkt\checkmark$\checkmark^\checkmark\\checkmarkc\checkmarki\checkmarkr\checkmarkc\checkmark$\checkmark{\checkmark$\checkmark^\checkmark\\checkmarkc\checkmarki\checkmarkr\checkmarkc\checkmark$\checkmark \checkmark$\checkmark^\checkmark\\checkmarkc\checkmarki\checkmarkr\checkmarkc\checkmark$\checkmarka\checkmark$\checkmark^\checkmark\\checkmarkc\checkmarki\checkmarkr\checkmarkc\checkmark$\checkmarkn\checkmark$\checkmark^\checkmark\\checkmarkc\checkmarki\checkmarkr\checkmarkc\checkmark$\checkmarks\checkmark$\checkmark^\checkmark\\checkmarkc\checkmarki\checkmarkr\checkmarkc\checkmark$\checkmark}\checkmark$\checkmark^\checkmark\\checkmarkc\checkmarki\checkmarkr\checkmarkc\checkmark$\checkmark}\checkmark$\checkmark^\checkmark\\checkmarkc\checkmarki\checkmarkr\checkmarkc\checkmark$\checkmark
\checkmark$\checkmark^\checkmark\\checkmarkc\checkmarki\checkmarkr\checkmarkc\checkmark$\checkmark \checkmark$\checkmark^\checkmark\\checkmarkc\checkmarki\checkmarkr\checkmarkc\checkmark$\checkmark \checkmark$\checkmark^\checkmark\\checkmarkc\checkmarki\checkmarkr\checkmarkc\checkmark$\checkmark \checkmark$\checkmark^\checkmark\\checkmarkc\checkmarki\checkmarkr\checkmarkc\checkmark$\checkmark \checkmark$\checkmark^\checkmark\\checkmarkc\checkmarki\checkmarkr\checkmarkc\checkmark$\checkmark\\checkmark$\checkmark^\checkmark\\checkmarkc\checkmarki\checkmarkr\checkmarkc\checkmark$\checkmarkl\checkmark$\checkmark^\checkmark\\checkmarkc\checkmarki\checkmarkr\checkmarkc\checkmark$\checkmarka\checkmark$\checkmark^\checkmark\\checkmarkc\checkmarki\checkmarkr\checkmarkc\checkmark$\checkmarkb\checkmark$\checkmark^\checkmark\\checkmarkc\checkmarki\checkmarkr\checkmarkc\checkmark$\checkmarke\checkmark$\checkmark^\checkmark\\checkmarkc\checkmarki\checkmarkr\checkmarkc\checkmark$\checkmarkl\checkmark$\checkmark^\checkmark\\checkmarkc\checkmarki\checkmarkr\checkmarkc\checkmark$\checkmark{\checkmark$\checkmark^\checkmark\\checkmarkc\checkmarki\checkmarkr\checkmarkc\checkmark$\checkmarke\checkmark$\checkmark^\checkmark\\checkmarkc\checkmarki\checkmarkr\checkmarkc\checkmark$\checkmarkq\checkmark$\checkmark^\checkmark\\checkmarkc\checkmarki\checkmarkr\checkmarkc\checkmark$\checkmark:\checkmark$\checkmark^\checkmark\\checkmarkc\checkmarki\checkmarkr\checkmarkc\checkmark$\checkmarkl\checkmark$\checkmark^\checkmark\\checkmarkc\checkmarki\checkmarkr\checkmarkc\checkmark$\checkmark1\checkmark$\checkmark^\checkmark\\checkmarkc\checkmarki\checkmarkr\checkmarkc\checkmark$\checkmark0\checkmark$\checkmark^\checkmark\\checkmarkc\checkmarki\checkmarkr\checkmarkc\checkmark$\checkmarkh\checkmark$\checkmark^\checkmark\\checkmarkc\checkmarki\checkmarkr\checkmarkc\checkmark$\checkmark}\checkmark$\checkmark^\checkmark\\checkmarkc\checkmarki\checkmarkr\checkmarkc\checkmark$\checkmark
\checkmark$\checkmark^\checkmark\\checkmarkc\checkmarki\checkmarkr\checkmarkc\checkmark$\checkmark\\checkmark$\checkmark^\checkmark\\checkmarkc\checkmarki\checkmarkr\checkmarkc\checkmark$\checkmarke\checkmark$\checkmark^\checkmark\\checkmarkc\checkmarki\checkmarkr\checkmarkc\checkmark$\checkmarkn\checkmark$\checkmark^\checkmark\\checkmarkc\checkmarki\checkmarkr\checkmarkc\checkmark$\checkmarkd\checkmark$\checkmark^\checkmark\\checkmarkc\checkmarki\checkmarkr\checkmarkc\checkmark$\checkmark{\checkmark$\checkmark^\checkmark\\checkmarkc\checkmarki\checkmarkr\checkmarkc\checkmark$\checkmarke\checkmark$\checkmark^\checkmark\\checkmarkc\checkmarki\checkmarkr\checkmarkc\checkmark$\checkmarkq\checkmark$\checkmark^\checkmark\\checkmarkc\checkmarki\checkmarkr\checkmarkc\checkmark$\checkmarku\checkmark$\checkmark^\checkmark\\checkmarkc\checkmarki\checkmarkr\checkmarkc\checkmark$\checkmarka\checkmark$\checkmark^\checkmark\\checkmarkc\checkmarki\checkmarkr\checkmarkc\checkmark$\checkmarkt\checkmark$\checkmark^\checkmark\\checkmarkc\checkmarki\checkmarkr\checkmarkc\checkmark$\checkmarki\checkmark$\checkmark^\checkmark\\checkmarkc\checkmarki\checkmarkr\checkmarkc\checkmark$\checkmarko\checkmark$\checkmark^\checkmark\\checkmarkc\checkmarki\checkmarkr\checkmarkc\checkmark$\checkmarkn\checkmark$\checkmark^\checkmark\\checkmarkc\checkmarki\checkmarkr\checkmarkc\checkmark$\checkmark}\checkmark$\checkmark^\checkmark\\checkmarkc\checkmarki\checkmarkr\checkmarkc\checkmark$\checkmark
\checkmark$\checkmark^\checkmark\\checkmarkc\checkmarki\checkmarkr\checkmarkc\checkmark$\checkmark
\checkmark$\checkmark^\checkmark\\checkmarkc\checkmarki\checkmarkr\checkmarkc\checkmark$\checkmark\\checkmark$\checkmark^\checkmark\\checkmarkc\checkmarki\checkmarkr\checkmarkc\checkmark$\checkmarkt\checkmark$\checkmark^\checkmark\\checkmarkc\checkmarki\checkmarkr\checkmarkc\checkmark$\checkmarke\checkmark$\checkmark^\checkmark\\checkmarkc\checkmarki\checkmarkr\checkmarkc\checkmark$\checkmarkx\checkmark$\checkmark^\checkmark\\checkmarkc\checkmarki\checkmarkr\checkmarkc\checkmark$\checkmarkt\checkmark$\checkmark^\checkmark\\checkmarkc\checkmarki\checkmarkr\checkmarkc\checkmark$\checkmarkb\checkmark$\checkmark^\checkmark\\checkmarkc\checkmarki\checkmarkr\checkmarkc\checkmark$\checkmarkf\checkmark$\checkmark^\checkmark\\checkmarkc\checkmarki\checkmarkr\checkmarkc\checkmark$\checkmark{\checkmark$\checkmark^\checkmark\\checkmarkc\checkmarki\checkmarkr\checkmarkc\checkmark$\checkmarkC\checkmark$\checkmark^\checkmark\\checkmarkc\checkmarki\checkmarkr\checkmarkc\checkmark$\checkmarko\checkmark$\checkmark^\checkmark\\checkmarkc\checkmarki\checkmarkr\checkmarkc\checkmark$\checkmarkn\checkmark$\checkmark^\checkmark\\checkmarkc\checkmarki\checkmarkr\checkmarkc\checkmark$\checkmarkc\checkmark$\checkmark^\checkmark\\checkmarkc\checkmarki\checkmarkr\checkmarkc\checkmark$\checkmarkl\checkmark$\checkmark^\checkmark\\checkmarkc\checkmarki\checkmarkr\checkmarkc\checkmark$\checkmarku\checkmark$\checkmark^\checkmark\\checkmarkc\checkmarki\checkmarkr\checkmarkc\checkmark$\checkmarks\checkmark$\checkmark^\checkmark\\checkmarkc\checkmarki\checkmarkr\checkmarkc\checkmark$\checkmarki\checkmark$\checkmark^\checkmark\\checkmarkc\checkmarki\checkmarkr\checkmarkc\checkmark$\checkmarko\checkmark$\checkmark^\checkmark\\checkmarkc\checkmarki\checkmarkr\checkmarkc\checkmark$\checkmarkn\checkmark$\checkmark^\checkmark\\checkmarkc\checkmarki\checkmarkr\checkmarkc\checkmark$\checkmark \checkmark$\checkmark^\checkmark\\checkmarkc\checkmarki\checkmarkr\checkmarkc\checkmark$\checkmark:\checkmark$\checkmark^\checkmark\\checkmarkc\checkmarki\checkmarkr\checkmarkc\checkmark$\checkmark}\checkmark$\checkmark^\checkmark\\checkmarkc\checkmarki\checkmarkr\checkmarkc\checkmark$\checkmark \checkmark$\checkmark^\checkmark\\checkmarkc\checkmarki\checkmarkr\checkmarkc\checkmark$\checkmarkL\checkmark$\checkmark^\checkmark\\checkmarkc\checkmarki\checkmarkr\checkmarkc\checkmark$\checkmarka\checkmark$\checkmark^\checkmark\\checkmarkc\checkmarki\checkmarkr\checkmarkc\checkmark$\checkmark \checkmark$\checkmark^\checkmark\\checkmarkc\checkmarki\checkmarkr\checkmarkc\checkmark$\checkmarkd\checkmark$\checkmark^\checkmark\\checkmarkc\checkmarki\checkmarkr\checkmarkc\checkmark$\checkmarku\checkmark$\checkmark^\checkmark\\checkmarkc\checkmarki\checkmarkr\checkmarkc\checkmark$\checkmarkr\checkmark$\checkmark^\checkmark\\checkmarkc\checkmarki\checkmarkr\checkmarkc\checkmark$\checkmarké\checkmark$\checkmark^\checkmark\\checkmarkc\checkmarki\checkmarkr\checkmarkc\checkmark$\checkmarke\checkmark$\checkmark^\checkmark\\checkmarkc\checkmarki\checkmarkr\checkmarkc\checkmark$\checkmark \checkmark$\checkmark^\checkmark\\checkmarkc\checkmarki\checkmarkr\checkmarkc\checkmark$\checkmarkd\checkmark$\checkmark^\checkmark\\checkmarkc\checkmarki\checkmarkr\checkmarkc\checkmark$\checkmarke\checkmark$\checkmark^\checkmark\\checkmarkc\checkmarki\checkmarkr\checkmarkc\checkmark$\checkmark \checkmark$\checkmark^\checkmark\\checkmarkc\checkmarki\checkmarkr\checkmarkc\checkmark$\checkmarkv\checkmark$\checkmark^\checkmark\\checkmarkc\checkmarki\checkmarkr\checkmarkc\checkmark$\checkmarki\checkmark$\checkmark^\checkmark\\checkmarkc\checkmarki\checkmarkr\checkmarkc\checkmark$\checkmarke\checkmark$\checkmark^\checkmark\\checkmarkc\checkmarki\checkmarkr\checkmarkc\checkmark$\checkmark \checkmark$\checkmark^\checkmark\\checkmarkc\checkmarki\checkmarkr\checkmarkc\checkmark$\checkmarkc\checkmark$\checkmark^\checkmark\\checkmarkc\checkmarki\checkmarkr\checkmarkc\checkmark$\checkmarka\checkmark$\checkmark^\checkmark\\checkmarkc\checkmarki\checkmarkr\checkmarkc\checkmark$\checkmarkl\checkmark$\checkmark^\checkmark\\checkmarkc\checkmarki\checkmarkr\checkmarkc\checkmark$\checkmarkc\checkmark$\checkmark^\checkmark\\checkmarkc\checkmarki\checkmarkr\checkmarkc\checkmark$\checkmarku\checkmark$\checkmark^\checkmark\\checkmarkc\checkmarki\checkmarkr\checkmarkc\checkmark$\checkmarkl\checkmark$\checkmark^\checkmark\\checkmarkc\checkmarki\checkmarkr\checkmarkc\checkmark$\checkmarké\checkmark$\checkmark^\checkmark\\checkmarkc\checkmarki\checkmarkr\checkmarkc\checkmark$\checkmarke\checkmark$\checkmark^\checkmark\\checkmarkc\checkmarki\checkmarkr\checkmarkc\checkmark$\checkmark \checkmark$\checkmark^\checkmark\\checkmarkc\checkmarki\checkmarkr\checkmarkc\checkmark$\checkmarkd\checkmark$\checkmark^\checkmark\\checkmarkc\checkmarki\checkmarkr\checkmarkc\checkmark$\checkmarke\checkmark$\checkmark^\checkmark\\checkmarkc\checkmarki\checkmarkr\checkmarkc\checkmark$\checkmarks\checkmark$\checkmark^\checkmark\\checkmarkc\checkmarki\checkmarkr\checkmarkc\checkmark$\checkmark \checkmark$\checkmark^\checkmark\\checkmarkc\checkmarki\checkmarkr\checkmarkc\checkmark$\checkmarkg\checkmark$\checkmark^\checkmark\\checkmarkc\checkmarki\checkmarkr\checkmarkc\checkmark$\checkmarku\checkmark$\checkmark^\checkmark\\checkmarkc\checkmarki\checkmarkr\checkmarkc\checkmark$\checkmarki\checkmark$\checkmark^\checkmark\\checkmarkc\checkmarki\checkmarkr\checkmarkc\checkmark$\checkmarkd\checkmark$\checkmark^\checkmark\\checkmarkc\checkmarki\checkmarkr\checkmarkc\checkmark$\checkmarka\checkmark$\checkmark^\checkmark\\checkmarkc\checkmarki\checkmarkr\checkmarkc\checkmark$\checkmarkg\checkmark$\checkmark^\checkmark\\checkmarkc\checkmarki\checkmarkr\checkmarkc\checkmark$\checkmarke\checkmark$\checkmark^\checkmark\\checkmarkc\checkmarki\checkmarkr\checkmarkc\checkmark$\checkmarks\checkmark$\checkmark^\checkmark\\checkmarkc\checkmarki\checkmarkr\checkmarkc\checkmark$\checkmark \checkmark$\checkmark^\checkmark\\checkmarkc\checkmarki\checkmarkr\checkmarkc\checkmark$\checkmarkH\checkmark$\checkmark^\checkmark\\checkmarkc\checkmarki\checkmarkr\checkmarkc\checkmark$\checkmarkI\checkmark$\checkmark^\checkmark\\checkmarkc\checkmarki\checkmarkr\checkmarkc\checkmark$\checkmarkW\checkmark$\checkmark^\checkmark\\checkmarkc\checkmarki\checkmarkr\checkmarkc\checkmark$\checkmarkI\checkmark$\checkmark^\checkmark\\checkmarkc\checkmarki\checkmarkr\checkmarkc\checkmark$\checkmarkN\checkmark$\checkmark^\checkmark\\checkmarkc\checkmarki\checkmarkr\checkmarkc\checkmark$\checkmark \checkmark$\checkmark^\checkmark\\checkmarkc\checkmarki\checkmarkr\checkmarkc\checkmark$\checkmarkH\checkmark$\checkmark^\checkmark\\checkmarkc\checkmarki\checkmarkr\checkmarkc\checkmark$\checkmarkG\checkmark$\checkmark^\checkmark\\checkmarkc\checkmarki\checkmarkr\checkmarkc\checkmark$\checkmarkR\checkmark$\checkmark^\checkmark\\checkmarkc\checkmarki\checkmarkr\checkmarkc\checkmark$\checkmark1\checkmark$\checkmark^\checkmark\\checkmarkc\checkmarki\checkmarkr\checkmarkc\checkmark$\checkmark5\checkmark$\checkmark^\checkmark\\checkmarkc\checkmarki\checkmarkr\checkmarkc\checkmark$\checkmark \checkmark$\checkmark^\checkmark\\checkmarkc\checkmarki\checkmarkr\checkmarkc\checkmark$\checkmarkd\checkmark$\checkmark^\checkmark\\checkmarkc\checkmarki\checkmarkr\checkmarkc\checkmark$\checkmarké\checkmark$\checkmark^\checkmark\\checkmarkc\checkmarki\checkmarkr\checkmarkc\checkmark$\checkmarkp\checkmark$\checkmark^\checkmark\\checkmarkc\checkmarki\checkmarkr\checkmarkc\checkmark$\checkmarka\checkmark$\checkmark^\checkmark\\checkmarkc\checkmarki\checkmarkr\checkmarkc\checkmark$\checkmarks\checkmark$\checkmark^\checkmark\\checkmarkc\checkmarki\checkmarkr\checkmarkc\checkmark$\checkmarks\checkmark$\checkmark^\checkmark\\checkmarkc\checkmarki\checkmarkr\checkmarkc\checkmark$\checkmarke\checkmark$\checkmark^\checkmark\\checkmarkc\checkmarki\checkmarkr\checkmarkc\checkmark$\checkmark \checkmark$\checkmark^\checkmark\\checkmarkc\checkmarki\checkmarkr\checkmarkc\checkmark$\checkmarkl\checkmark$\checkmark^\checkmark\\checkmarkc\checkmarki\checkmarkr\checkmarkc\checkmark$\checkmarka\checkmark$\checkmark^\checkmark\\checkmarkc\checkmarki\checkmarkr\checkmarkc\checkmark$\checkmarkr\checkmark$\checkmark^\checkmark\\checkmarkc\checkmarki\checkmarkr\checkmarkc\checkmark$\checkmarkg\checkmark$\checkmark^\checkmark\\checkmarkc\checkmarki\checkmarkr\checkmarkc\checkmark$\checkmarke\checkmark$\checkmark^\checkmark\\checkmarkc\checkmarki\checkmarkr\checkmarkc\checkmark$\checkmarkm\checkmark$\checkmark^\checkmark\\checkmarkc\checkmarki\checkmarkr\checkmarkc\checkmark$\checkmarke\checkmark$\checkmark^\checkmark\\checkmarkc\checkmarki\checkmarkr\checkmarkc\checkmark$\checkmarkn\checkmark$\checkmark^\checkmark\\checkmarkc\checkmarki\checkmarkr\checkmarkc\checkmark$\checkmarkt\checkmark$\checkmark^\checkmark\\checkmarkc\checkmarki\checkmarkr\checkmarkc\checkmark$\checkmark \checkmark$\checkmark^\checkmark\\checkmarkc\checkmarki\checkmarkr\checkmarkc\checkmark$\checkmarkt\checkmark$\checkmark^\checkmark\\checkmarkc\checkmarki\checkmarkr\checkmarkc\checkmark$\checkmarko\checkmark$\checkmark^\checkmark\\checkmarkc\checkmarki\checkmarkr\checkmarkc\checkmark$\checkmarku\checkmark$\checkmark^\checkmark\\checkmarkc\checkmarki\checkmarkr\checkmarkc\checkmark$\checkmarkt\checkmark$\checkmark^\checkmark\\checkmarkc\checkmarki\checkmarkr\checkmarkc\checkmark$\checkmarke\checkmark$\checkmark^\checkmark\\checkmarkc\checkmarki\checkmarkr\checkmarkc\checkmark$\checkmark \checkmark$\checkmark^\checkmark\\checkmarkc\checkmarki\checkmarkr\checkmarkc\checkmark$\checkmarku\checkmark$\checkmark^\checkmark\\checkmarkc\checkmarki\checkmarkr\checkmarkc\checkmark$\checkmarkt\checkmark$\checkmark^\checkmark\\checkmarkc\checkmarki\checkmarkr\checkmarkc\checkmark$\checkmarki\checkmark$\checkmark^\checkmark\\checkmarkc\checkmarki\checkmarkr\checkmarkc\checkmark$\checkmarkl\checkmark$\checkmark^\checkmark\\checkmarkc\checkmarki\checkmarkr\checkmarkc\checkmark$\checkmarki\checkmark$\checkmark^\checkmark\\checkmarkc\checkmarki\checkmarkr\checkmarkc\checkmark$\checkmarks\checkmark$\checkmark^\checkmark\\checkmarkc\checkmarki\checkmarkr\checkmarkc\checkmark$\checkmarka\checkmark$\checkmark^\checkmark\\checkmarkc\checkmarki\checkmarkr\checkmarkc\checkmark$\checkmarkt\checkmark$\checkmark^\checkmark\\checkmarkc\checkmarki\checkmarkr\checkmarkc\checkmark$\checkmarki\checkmark$\checkmark^\checkmark\\checkmarkc\checkmarki\checkmarkr\checkmarkc\checkmark$\checkmarko\checkmark$\checkmark^\checkmark\\checkmarkc\checkmarki\checkmarkr\checkmarkc\checkmark$\checkmarkn\checkmark$\checkmark^\checkmark\\checkmarkc\checkmarki\checkmarkr\checkmarkc\checkmark$\checkmark \checkmark$\checkmark^\checkmark\\checkmarkc\checkmarki\checkmarkr\checkmarkc\checkmark$\checkmarkr\checkmark$\checkmark^\checkmark\\checkmarkc\checkmarki\checkmarkr\checkmarkc\checkmark$\checkmarké\checkmark$\checkmark^\checkmark\\checkmarkc\checkmarki\checkmarkr\checkmarkc\checkmark$\checkmarka\checkmark$\checkmark^\checkmark\\checkmarkc\checkmarki\checkmarkr\checkmarkc\checkmark$\checkmarkl\checkmark$\checkmark^\checkmark\\checkmarkc\checkmarki\checkmarkr\checkmarkc\checkmark$\checkmarki\checkmark$\checkmark^\checkmark\\checkmarkc\checkmarki\checkmarkr\checkmarkc\checkmark$\checkmarks\checkmark$\checkmark^\checkmark\\checkmarkc\checkmarki\checkmarkr\checkmarkc\checkmark$\checkmarkt\checkmark$\checkmark^\checkmark\\checkmarkc\checkmarki\checkmarkr\checkmarkc\checkmark$\checkmarke\checkmark$\checkmark^\checkmark\\checkmarkc\checkmarki\checkmarkr\checkmarkc\checkmark$\checkmark \checkmark$\checkmark^\checkmark\\checkmarkc\checkmarki\checkmarkr\checkmarkc\checkmark$\checkmarkd\checkmark$\checkmark^\checkmark\\checkmarkc\checkmarki\checkmarkr\checkmarkc\checkmark$\checkmark'\checkmark$\checkmark^\checkmark\\checkmarkc\checkmarki\checkmarkr\checkmarkc\checkmark$\checkmarku\checkmark$\checkmark^\checkmark\\checkmarkc\checkmarki\checkmarkr\checkmarkc\checkmark$\checkmarkn\checkmark$\checkmark^\checkmark\\checkmarkc\checkmarki\checkmarkr\checkmarkc\checkmark$\checkmarke\checkmark$\checkmark^\checkmark\\checkmarkc\checkmarki\checkmarkr\checkmarkc\checkmark$\checkmark \checkmark$\checkmark^\checkmark\\checkmarkc\checkmarki\checkmarkr\checkmarkc\checkmark$\checkmarkm\checkmark$\checkmark^\checkmark\\checkmarkc\checkmarki\checkmarkr\checkmarkc\checkmark$\checkmarka\checkmark$\checkmark^\checkmark\\checkmarkc\checkmarki\checkmarkr\checkmarkc\checkmark$\checkmarkc\checkmark$\checkmark^\checkmark\\checkmarkc\checkmarki\checkmarkr\checkmarkc\checkmark$\checkmarkh\checkmark$\checkmark^\checkmark\\checkmarkc\checkmarki\checkmarkr\checkmarkc\checkmark$\checkmarki\checkmark$\checkmark^\checkmark\\checkmarkc\checkmarki\checkmarkr\checkmarkc\checkmark$\checkmarkn\checkmark$\checkmark^\checkmark\\checkmarkc\checkmarki\checkmarkr\checkmarkc\checkmark$\checkmarke\checkmark$\checkmark^\checkmark\\checkmarkc\checkmarki\checkmarkr\checkmarkc\checkmark$\checkmark \checkmark$\checkmark^\checkmark\\checkmarkc\checkmarki\checkmarkr\checkmarkc\checkmark$\checkmarkd\checkmark$\checkmark^\checkmark\\checkmarkc\checkmarki\checkmarkr\checkmarkc\checkmark$\checkmarki\checkmark$\checkmark^\checkmark\\checkmarkc\checkmarki\checkmarkr\checkmarkc\checkmark$\checkmarkd\checkmark$\checkmark^\checkmark\\checkmarkc\checkmarki\checkmarkr\checkmarkc\checkmark$\checkmarka\checkmark$\checkmark^\checkmark\\checkmarkc\checkmarki\checkmarkr\checkmarkc\checkmark$\checkmarkc\checkmark$\checkmark^\checkmark\\checkmarkc\checkmarki\checkmarkr\checkmarkc\checkmark$\checkmarkt\checkmark$\checkmark^\checkmark\\checkmarkc\checkmarki\checkmarkr\checkmarkc\checkmark$\checkmarki\checkmark$\checkmark^\checkmark\\checkmarkc\checkmarki\checkmarkr\checkmarkc\checkmark$\checkmarkq\checkmark$\checkmark^\checkmark\\checkmarkc\checkmarki\checkmarkr\checkmarkc\checkmark$\checkmarku\checkmark$\checkmark^\checkmark\\checkmarkc\checkmarki\checkmarkr\checkmarkc\checkmark$\checkmarke\checkmark$\checkmark^\checkmark\\checkmarkc\checkmarki\checkmarkr\checkmarkc\checkmark$\checkmark,\checkmark$\checkmark^\checkmark\\checkmarkc\checkmarki\checkmarkr\checkmarkc\checkmark$\checkmark \checkmark$\checkmark^\checkmark\\checkmarkc\checkmarki\checkmarkr\checkmarkc\checkmark$\checkmarkv\checkmark$\checkmark^\checkmark\\checkmarkc\checkmarki\checkmarkr\checkmarkc\checkmark$\checkmarka\checkmark$\checkmark^\checkmark\\checkmarkc\checkmarki\checkmarkr\checkmarkc\checkmark$\checkmarkl\checkmark$\checkmark^\checkmark\\checkmarkc\checkmarki\checkmarkr\checkmarkc\checkmark$\checkmarki\checkmark$\checkmark^\checkmark\\checkmarkc\checkmarki\checkmarkr\checkmarkc\checkmark$\checkmarkd\checkmark$\checkmark^\checkmark\\checkmarkc\checkmarki\checkmarkr\checkmarkc\checkmark$\checkmarka\checkmark$\checkmark^\checkmark\\checkmarkc\checkmarki\checkmarkr\checkmarkc\checkmark$\checkmarkn\checkmark$\checkmark^\checkmark\\checkmarkc\checkmarki\checkmarkr\checkmarkc\checkmark$\checkmarkt\checkmark$\checkmark^\checkmark\\checkmarkc\checkmarki\checkmarkr\checkmarkc\checkmark$\checkmark \checkmark$\checkmark^\checkmark\\checkmarkc\checkmarki\checkmarkr\checkmarkc\checkmark$\checkmarka\checkmark$\checkmark^\checkmark\\checkmarkc\checkmarki\checkmarkr\checkmarkc\checkmark$\checkmarkm\checkmark$\checkmark^\checkmark\\checkmarkc\checkmarki\checkmarkr\checkmarkc\checkmark$\checkmarkp\checkmark$\checkmark^\checkmark\\checkmarkc\checkmarki\checkmarkr\checkmarkc\checkmark$\checkmarkl\checkmark$\checkmark^\checkmark\\checkmarkc\checkmarki\checkmarkr\checkmarkc\checkmark$\checkmarke\checkmark$\checkmark^\checkmark\\checkmarkc\checkmarki\checkmarkr\checkmarkc\checkmark$\checkmarkm\checkmark$\checkmark^\checkmark\\checkmarkc\checkmarki\checkmarkr\checkmarkc\checkmark$\checkmarke\checkmark$\checkmark^\checkmark\\checkmarkc\checkmarki\checkmarkr\checkmarkc\checkmark$\checkmarkn\checkmark$\checkmark^\checkmark\\checkmarkc\checkmarki\checkmarkr\checkmarkc\checkmark$\checkmarkt\checkmark$\checkmark^\checkmark\\checkmarkc\checkmarki\checkmarkr\checkmarkc\checkmark$\checkmark \checkmark$\checkmark^\checkmark\\checkmarkc\checkmarki\checkmarkr\checkmarkc\checkmark$\checkmarkc\checkmark$\checkmark^\checkmark\\checkmarkc\checkmarki\checkmarkr\checkmarkc\checkmark$\checkmarke\checkmark$\checkmark^\checkmark\\checkmarkc\checkmarki\checkmarkr\checkmarkc\checkmark$\checkmark \checkmark$\checkmark^\checkmark\\checkmarkc\checkmarki\checkmarkr\checkmarkc\checkmark$\checkmarkc\checkmark$\checkmark^\checkmark\\checkmarkc\checkmarki\checkmarkr\checkmarkc\checkmark$\checkmarkh\checkmark$\checkmark^\checkmark\\checkmarkc\checkmarki\checkmarkr\checkmarkc\checkmark$\checkmarko\checkmark$\checkmark^\checkmark\\checkmarkc\checkmarki\checkmarkr\checkmarkc\checkmark$\checkmarki\checkmark$\checkmark^\checkmark\\checkmarkc\checkmarki\checkmarkr\checkmarkc\checkmark$\checkmarkx\checkmark$\checkmark^\checkmark\\checkmarkc\checkmarki\checkmarkr\checkmarkc\checkmark$\checkmark \checkmark$\checkmark^\checkmark\\checkmarkc\checkmarki\checkmarkr\checkmarkc\checkmark$\checkmarkt\checkmark$\checkmark^\checkmark\\checkmarkc\checkmarki\checkmarkr\checkmarkc\checkmark$\checkmarke\checkmark$\checkmark^\checkmark\\checkmarkc\checkmarki\checkmarkr\checkmarkc\checkmark$\checkmarkc\checkmark$\checkmark^\checkmark\\checkmarkc\checkmarki\checkmarkr\checkmarkc\checkmark$\checkmarkh\checkmark$\checkmark^\checkmark\\checkmarkc\checkmarki\checkmarkr\checkmarkc\checkmark$\checkmarkn\checkmark$\checkmark^\checkmark\\checkmarkc\checkmarki\checkmarkr\checkmarkc\checkmark$\checkmarko\checkmark$\checkmark^\checkmark\\checkmarkc\checkmarki\checkmarkr\checkmarkc\checkmark$\checkmarkl\checkmark$\checkmark^\checkmark\\checkmarkc\checkmarki\checkmarkr\checkmarkc\checkmark$\checkmarko\checkmark$\checkmark^\checkmark\\checkmarkc\checkmarki\checkmarkr\checkmarkc\checkmark$\checkmarkg\checkmark$\checkmark^\checkmark\\checkmarkc\checkmarki\checkmarkr\checkmarkc\checkmark$\checkmarki\checkmark$\checkmark^\checkmark\\checkmarkc\checkmarki\checkmarkr\checkmarkc\checkmark$\checkmarkq\checkmark$\checkmark^\checkmark\\checkmarkc\checkmarki\checkmarkr\checkmarkc\checkmark$\checkmarku\checkmark$\checkmark^\checkmark\\checkmarkc\checkmarki\checkmarkr\checkmarkc\checkmark$\checkmarke\checkmark$\checkmark^\checkmark\\checkmarkc\checkmarki\checkmarkr\checkmarkc\checkmark$\checkmark.\checkmark$\checkmark^\checkmark\\checkmarkc\checkmarki\checkmarkr\checkmarkc\checkmark$\checkmark
\checkmark$\checkmark^\checkmark\\checkmarkc\checkmarki\checkmarkr\checkmarkc\checkmark$\checkmark
\checkmark$\checkmark^\checkmark\\checkmarkc\checkmarki\checkmarkr\checkmarkc\checkmark$\checkmark%\checkmark$\checkmark^\checkmark\\checkmarkc\checkmarki\checkmarkr\checkmarkc\checkmark$\checkmark \checkmark$\checkmark^\checkmark\\checkmarkc\checkmarki\checkmarkr\checkmarkc\checkmark$\checkmark=\checkmark$\checkmark^\checkmark\\checkmarkc\checkmarki\checkmarkr\checkmarkc\checkmark$\checkmark=\checkmark$\checkmark^\checkmark\\checkmarkc\checkmarki\checkmarkr\checkmarkc\checkmark$\checkmark=\checkmark$\checkmark^\checkmark\\checkmarkc\checkmarki\checkmarkr\checkmarkc\checkmark$\checkmark=\checkmark$\checkmark^\checkmark\\checkmarkc\checkmarki\checkmarkr\checkmarkc\checkmark$\checkmark=\checkmark$\checkmark^\checkmark\\checkmarkc\checkmarki\checkmarkr\checkmarkc\checkmark$\checkmark=\checkmark$\checkmark^\checkmark\\checkmarkc\checkmarki\checkmarkr\checkmarkc\checkmark$\checkmark=\checkmark$\checkmark^\checkmark\\checkmarkc\checkmarki\checkmarkr\checkmarkc\checkmark$\checkmark=\checkmark$\checkmark^\checkmark\\checkmarkc\checkmarki\checkmarkr\checkmarkc\checkmark$\checkmark=\checkmark$\checkmark^\checkmark\\checkmarkc\checkmarki\checkmarkr\checkmarkc\checkmark$\checkmark=\checkmark$\checkmark^\checkmark\\checkmarkc\checkmarki\checkmarkr\checkmarkc\checkmark$\checkmark=\checkmark$\checkmark^\checkmark\\checkmarkc\checkmarki\checkmarkr\checkmarkc\checkmark$\checkmark=\checkmark$\checkmark^\checkmark\\checkmarkc\checkmarki\checkmarkr\checkmarkc\checkmark$\checkmark=\checkmark$\checkmark^\checkmark\\checkmarkc\checkmarki\checkmarkr\checkmarkc\checkmark$\checkmark=\checkmark$\checkmark^\checkmark\\checkmarkc\checkmarki\checkmarkr\checkmarkc\checkmark$\checkmark=\checkmark$\checkmark^\checkmark\\checkmarkc\checkmarki\checkmarkr\checkmarkc\checkmark$\checkmark=\checkmark$\checkmark^\checkmark\\checkmarkc\checkmarki\checkmarkr\checkmarkc\checkmark$\checkmark=\checkmark$\checkmark^\checkmark\\checkmarkc\checkmarki\checkmarkr\checkmarkc\checkmark$\checkmark=\checkmark$\checkmark^\checkmark\\checkmarkc\checkmarki\checkmarkr\checkmarkc\checkmark$\checkmark=\checkmark$\checkmark^\checkmark\\checkmarkc\checkmarki\checkmarkr\checkmarkc\checkmark$\checkmark=\checkmark$\checkmark^\checkmark\\checkmarkc\checkmarki\checkmarkr\checkmarkc\checkmark$\checkmark=\checkmark$\checkmark^\checkmark\\checkmarkc\checkmarki\checkmarkr\checkmarkc\checkmark$\checkmark=\checkmark$\checkmark^\checkmark\\checkmarkc\checkmarki\checkmarkr\checkmarkc\checkmark$\checkmark=\checkmark$\checkmark^\checkmark\\checkmarkc\checkmarki\checkmarkr\checkmarkc\checkmark$\checkmark=\checkmark$\checkmark^\checkmark\\checkmarkc\checkmarki\checkmarkr\checkmarkc\checkmark$\checkmark=\checkmark$\checkmark^\checkmark\\checkmarkc\checkmarki\checkmarkr\checkmarkc\checkmark$\checkmark=\checkmark$\checkmark^\checkmark\\checkmarkc\checkmarki\checkmarkr\checkmarkc\checkmark$\checkmark=\checkmark$\checkmark^\checkmark\\checkmarkc\checkmarki\checkmarkr\checkmarkc\checkmark$\checkmark=\checkmark$\checkmark^\checkmark\\checkmarkc\checkmarki\checkmarkr\checkmarkc\checkmark$\checkmark=\checkmark$\checkmark^\checkmark\\checkmarkc\checkmarki\checkmarkr\checkmarkc\checkmark$\checkmark=\checkmark$\checkmark^\checkmark\\checkmarkc\checkmarki\checkmarkr\checkmarkc\checkmark$\checkmark=\checkmark$\checkmark^\checkmark\\checkmarkc\checkmarki\checkmarkr\checkmarkc\checkmark$\checkmark=\checkmark$\checkmark^\checkmark\\checkmarkc\checkmarki\checkmarkr\checkmarkc\checkmark$\checkmark=\checkmark$\checkmark^\checkmark\\checkmarkc\checkmarki\checkmarkr\checkmarkc\checkmark$\checkmark=\checkmark$\checkmark^\checkmark\\checkmarkc\checkmarki\checkmarkr\checkmarkc\checkmark$\checkmark=\checkmark$\checkmark^\checkmark\\checkmarkc\checkmarki\checkmarkr\checkmarkc\checkmark$\checkmark=\checkmark$\checkmark^\checkmark\\checkmarkc\checkmarki\checkmarkr\checkmarkc\checkmark$\checkmark=\checkmark$\checkmark^\checkmark\\checkmarkc\checkmarki\checkmarkr\checkmarkc\checkmark$\checkmark=\checkmark$\checkmark^\checkmark\\checkmarkc\checkmarki\checkmarkr\checkmarkc\checkmark$\checkmark=\checkmark$\checkmark^\checkmark\\checkmarkc\checkmarki\checkmarkr\checkmarkc\checkmark$\checkmark=\checkmark$\checkmark^\checkmark\\checkmarkc\checkmarki\checkmarkr\checkmarkc\checkmark$\checkmark=\checkmark$\checkmark^\checkmark\\checkmarkc\checkmarki\checkmarkr\checkmarkc\checkmark$\checkmark=\checkmark$\checkmark^\checkmark\\checkmarkc\checkmarki\checkmarkr\checkmarkc\checkmark$\checkmark=\checkmark$\checkmark^\checkmark\\checkmarkc\checkmarki\checkmarkr\checkmarkc\checkmark$\checkmark=\checkmark$\checkmark^\checkmark\\checkmarkc\checkmarki\checkmarkr\checkmarkc\checkmark$\checkmark=\checkmark$\checkmark^\checkmark\\checkmarkc\checkmarki\checkmarkr\checkmarkc\checkmark$\checkmark=\checkmark$\checkmark^\checkmark\\checkmarkc\checkmarki\checkmarkr\checkmarkc\checkmark$\checkmark=\checkmark$\checkmark^\checkmark\\checkmarkc\checkmarki\checkmarkr\checkmarkc\checkmark$\checkmark=\checkmark$\checkmark^\checkmark\\checkmarkc\checkmarki\checkmarkr\checkmarkc\checkmark$\checkmark=\checkmark$\checkmark^\checkmark\\checkmarkc\checkmarki\checkmarkr\checkmarkc\checkmark$\checkmark=\checkmark$\checkmark^\checkmark\\checkmarkc\checkmarki\checkmarkr\checkmarkc\checkmark$\checkmark=\checkmark$\checkmark^\checkmark\\checkmarkc\checkmarki\checkmarkr\checkmarkc\checkmark$\checkmark=\checkmark$\checkmark^\checkmark\\checkmarkc\checkmarki\checkmarkr\checkmarkc\checkmark$\checkmark=\checkmark$\checkmark^\checkmark\\checkmarkc\checkmarki\checkmarkr\checkmarkc\checkmark$\checkmark=\checkmark$\checkmark^\checkmark\\checkmarkc\checkmarki\checkmarkr\checkmarkc\checkmark$\checkmark=\checkmark$\checkmark^\checkmark\\checkmarkc\checkmarki\checkmarkr\checkmarkc\checkmark$\checkmark=\checkmark$\checkmark^\checkmark\\checkmarkc\checkmarki\checkmarkr\checkmarkc\checkmark$\checkmark=\checkmark$\checkmark^\checkmark\\checkmarkc\checkmarki\checkmarkr\checkmarkc\checkmark$\checkmark=\checkmark$\checkmark^\checkmark\\checkmarkc\checkmarki\checkmarkr\checkmarkc\checkmark$\checkmark=\checkmark$\checkmark^\checkmark\\checkmarkc\checkmarki\checkmarkr\checkmarkc\checkmark$\checkmark=\checkmark$\checkmark^\checkmark\\checkmarkc\checkmarki\checkmarkr\checkmarkc\checkmark$\checkmark=\checkmark$\checkmark^\checkmark\\checkmarkc\checkmarki\checkmarkr\checkmarkc\checkmark$\checkmark=\checkmark$\checkmark^\checkmark\\checkmarkc\checkmarki\checkmarkr\checkmarkc\checkmark$\checkmark
\checkmark$\checkmark^\checkmark\\checkmarkc\checkmarki\checkmarkr\checkmarkc\checkmark$\checkmark\\checkmark$\checkmark^\checkmark\\checkmarkc\checkmarki\checkmarkr\checkmarkc\checkmark$\checkmarks\checkmark$\checkmark^\checkmark\\checkmarkc\checkmarki\checkmarkr\checkmarkc\checkmark$\checkmarke\checkmark$\checkmark^\checkmark\\checkmarkc\checkmarki\checkmarkr\checkmarkc\checkmark$\checkmarkc\checkmark$\checkmark^\checkmark\\checkmarkc\checkmarki\checkmarkr\checkmarkc\checkmark$\checkmarkt\checkmark$\checkmark^\checkmark\\checkmarkc\checkmarki\checkmarkr\checkmarkc\checkmark$\checkmarki\checkmark$\checkmark^\checkmark\\checkmarkc\checkmarki\checkmarkr\checkmarkc\checkmark$\checkmarko\checkmark$\checkmark^\checkmark\\checkmarkc\checkmarki\checkmarkr\checkmarkc\checkmark$\checkmarkn\checkmark$\checkmark^\checkmark\\checkmarkc\checkmarki\checkmarkr\checkmarkc\checkmark$\checkmark{\checkmark$\checkmark^\checkmark\\checkmarkc\checkmarki\checkmarkr\checkmarkc\checkmark$\checkmarkV\checkmark$\checkmark^\checkmark\\checkmarkc\checkmarki\checkmarkr\checkmarkc\checkmark$\checkmarké\checkmark$\checkmark^\checkmark\\checkmarkc\checkmarki\checkmarkr\checkmarkc\checkmark$\checkmarkr\checkmark$\checkmark^\checkmark\\checkmarkc\checkmarki\checkmarkr\checkmarkc\checkmark$\checkmarki\checkmark$\checkmark^\checkmark\\checkmarkc\checkmarki\checkmarkr\checkmarkc\checkmark$\checkmarkf\checkmark$\checkmark^\checkmark\\checkmarkc\checkmarki\checkmarkr\checkmarkc\checkmark$\checkmarki\checkmark$\checkmark^\checkmark\\checkmarkc\checkmarki\checkmarkr\checkmarkc\checkmark$\checkmarkc\checkmark$\checkmark^\checkmark\\checkmarkc\checkmarki\checkmarkr\checkmarkc\checkmark$\checkmarka\checkmark$\checkmark^\checkmark\\checkmarkc\checkmarki\checkmarkr\checkmarkc\checkmark$\checkmarkt\checkmark$\checkmark^\checkmark\\checkmarkc\checkmarki\checkmarkr\checkmarkc\checkmark$\checkmarki\checkmark$\checkmark^\checkmark\\checkmarkc\checkmarki\checkmarkr\checkmarkc\checkmark$\checkmarko\checkmark$\checkmark^\checkmark\\checkmarkc\checkmarki\checkmarkr\checkmarkc\checkmark$\checkmarkn\checkmark$\checkmark^\checkmark\\checkmarkc\checkmarki\checkmarkr\checkmarkc\checkmark$\checkmark \checkmark$\checkmark^\checkmark\\checkmarkc\checkmarki\checkmarkr\checkmarkc\checkmark$\checkmarkF\checkmark$\checkmark^\checkmark\\checkmarkc\checkmarki\checkmarkr\checkmarkc\checkmark$\checkmarkr\checkmark$\checkmark^\checkmark\\checkmarkc\checkmarki\checkmarkr\checkmarkc\checkmark$\checkmarké\checkmark$\checkmark^\checkmark\\checkmarkc\checkmarki\checkmarkr\checkmarkc\checkmark$\checkmarkq\checkmark$\checkmark^\checkmark\\checkmarkc\checkmarki\checkmarkr\checkmarkc\checkmark$\checkmarku\checkmark$\checkmark^\checkmark\\checkmarkc\checkmarki\checkmarkr\checkmarkc\checkmark$\checkmarke\checkmark$\checkmark^\checkmark\\checkmarkc\checkmarki\checkmarkr\checkmarkc\checkmark$\checkmarkn\checkmark$\checkmark^\checkmark\\checkmarkc\checkmarki\checkmarkr\checkmarkc\checkmark$\checkmarkt\checkmark$\checkmark^\checkmark\\checkmarkc\checkmarki\checkmarkr\checkmarkc\checkmark$\checkmarki\checkmark$\checkmark^\checkmark\\checkmarkc\checkmarki\checkmarkr\checkmarkc\checkmark$\checkmarke\checkmark$\checkmark^\checkmark\\checkmarkc\checkmarki\checkmarkr\checkmarkc\checkmark$\checkmarkl\checkmark$\checkmark^\checkmark\\checkmarkc\checkmarki\checkmarkr\checkmarkc\checkmark$\checkmarkl\checkmark$\checkmark^\checkmark\\checkmarkc\checkmarki\checkmarkr\checkmarkc\checkmark$\checkmarke\checkmark$\checkmark^\checkmark\\checkmarkc\checkmarki\checkmarkr\checkmarkc\checkmark$\checkmark \checkmark$\checkmark^\checkmark\\checkmarkc\checkmarki\checkmarkr\checkmarkc\checkmark$\checkmark—\checkmark$\checkmark^\checkmark\\checkmarkc\checkmarki\checkmarkr\checkmarkc\checkmark$\checkmark \checkmark$\checkmark^\checkmark\\checkmarkc\checkmarki\checkmarkr\checkmarkc\checkmark$\checkmarkA\checkmark$\checkmark^\checkmark\\checkmarkc\checkmarki\checkmarkr\checkmarkc\checkmark$\checkmarkn\checkmark$\checkmark^\checkmark\\checkmarkc\checkmarki\checkmarkr\checkmarkc\checkmark$\checkmarka\checkmark$\checkmark^\checkmark\\checkmarkc\checkmarki\checkmarkr\checkmarkc\checkmark$\checkmarkl\checkmark$\checkmark^\checkmark\\checkmarkc\checkmarki\checkmarkr\checkmarkc\checkmark$\checkmarky\checkmark$\checkmark^\checkmark\\checkmarkc\checkmarki\checkmarkr\checkmarkc\checkmark$\checkmarks\checkmark$\checkmark^\checkmark\\checkmarkc\checkmarki\checkmarkr\checkmarkc\checkmark$\checkmarke\checkmark$\checkmark^\checkmark\\checkmarkc\checkmarki\checkmarkr\checkmarkc\checkmark$\checkmark \checkmark$\checkmark^\checkmark\\checkmarkc\checkmarki\checkmarkr\checkmarkc\checkmark$\checkmarkd\checkmark$\checkmark^\checkmark\\checkmarkc\checkmarki\checkmarkr\checkmarkc\checkmark$\checkmarke\checkmark$\checkmark^\checkmark\\checkmarkc\checkmarki\checkmarkr\checkmarkc\checkmark$\checkmark \checkmark$\checkmark^\checkmark\\checkmarkc\checkmarki\checkmarkr\checkmarkc\checkmark$\checkmarkR\checkmark$\checkmark^\checkmark\\checkmarkc\checkmarki\checkmarkr\checkmarkc\checkmark$\checkmarké\checkmark$\checkmark^\checkmark\\checkmarkc\checkmarki\checkmarkr\checkmarkc\checkmark$\checkmarks\checkmark$\checkmark^\checkmark\\checkmarkc\checkmarki\checkmarkr\checkmarkc\checkmark$\checkmarko\checkmark$\checkmark^\checkmark\\checkmarkc\checkmarki\checkmarkr\checkmarkc\checkmark$\checkmarkn\checkmark$\checkmark^\checkmark\\checkmarkc\checkmarki\checkmarkr\checkmarkc\checkmark$\checkmarka\checkmark$\checkmark^\checkmark\\checkmarkc\checkmarki\checkmarkr\checkmarkc\checkmark$\checkmarkn\checkmark$\checkmark^\checkmark\\checkmarkc\checkmarki\checkmarkr\checkmarkc\checkmark$\checkmarkc\checkmark$\checkmark^\checkmark\\checkmarkc\checkmarki\checkmarkr\checkmarkc\checkmark$\checkmarke\checkmark$\checkmark^\checkmark\\checkmarkc\checkmarki\checkmarkr\checkmarkc\checkmark$\checkmark}\checkmark$\checkmark^\checkmark\\checkmarkc\checkmarki\checkmarkr\checkmarkc\checkmark$\checkmark
\checkmark$\checkmark^\checkmark\\checkmarkc\checkmarki\checkmarkr\checkmarkc\checkmark$\checkmark
\checkmark$\checkmark^\checkmark\\checkmarkc\checkmarki\checkmarkr\checkmarkc\checkmark$\checkmarkL\checkmark$\checkmark^\checkmark\\checkmarkc\checkmarki\checkmarkr\checkmarkc\checkmark$\checkmarka\checkmark$\checkmark^\checkmark\\checkmarkc\checkmarki\checkmarkr\checkmarkc\checkmark$\checkmark \checkmark$\checkmark^\checkmark\\checkmarkc\checkmarki\checkmarkr\checkmarkc\checkmark$\checkmarkf\checkmark$\checkmark^\checkmark\\checkmarkc\checkmarki\checkmarkr\checkmarkc\checkmark$\checkmarkr\checkmark$\checkmark^\checkmark\\checkmarkc\checkmarki\checkmarkr\checkmarkc\checkmark$\checkmarké\checkmark$\checkmark^\checkmark\\checkmarkc\checkmarki\checkmarkr\checkmarkc\checkmark$\checkmarkq\checkmark$\checkmark^\checkmark\\checkmarkc\checkmarki\checkmarkr\checkmarkc\checkmark$\checkmarku\checkmark$\checkmark^\checkmark\\checkmarkc\checkmarki\checkmarkr\checkmarkc\checkmark$\checkmarke\checkmark$\checkmark^\checkmark\\checkmarkc\checkmarki\checkmarkr\checkmarkc\checkmark$\checkmarkn\checkmark$\checkmark^\checkmark\\checkmarkc\checkmarki\checkmarkr\checkmarkc\checkmark$\checkmarkc\checkmark$\checkmark^\checkmark\\checkmarkc\checkmarki\checkmarkr\checkmarkc\checkmark$\checkmarke\checkmark$\checkmark^\checkmark\\checkmarkc\checkmarki\checkmarkr\checkmarkc\checkmark$\checkmark \checkmark$\checkmark^\checkmark\\checkmarkc\checkmarki\checkmarkr\checkmarkc\checkmark$\checkmarkp\checkmark$\checkmark^\checkmark\\checkmarkc\checkmarki\checkmarkr\checkmarkc\checkmark$\checkmarkr\checkmark$\checkmark^\checkmark\\checkmarkc\checkmarki\checkmarkr\checkmarkc\checkmark$\checkmarko\checkmark$\checkmark^\checkmark\\checkmarkc\checkmarki\checkmarkr\checkmarkc\checkmark$\checkmarkp\checkmark$\checkmark^\checkmark\\checkmarkc\checkmarki\checkmarkr\checkmarkc\checkmark$\checkmarkr\checkmark$\checkmark^\checkmark\\checkmarkc\checkmarki\checkmarkr\checkmarkc\checkmark$\checkmarke\checkmark$\checkmark^\checkmark\\checkmarkc\checkmarki\checkmarkr\checkmarkc\checkmark$\checkmark \checkmark$\checkmark^\checkmark\\checkmarkc\checkmarki\checkmarkr\checkmarkc\checkmark$\checkmarkd\checkmark$\checkmark^\checkmark\\checkmarkc\checkmarki\checkmarkr\checkmarkc\checkmark$\checkmarke\checkmark$\checkmark^\checkmark\\checkmarkc\checkmarki\checkmarkr\checkmarkc\checkmark$\checkmark \checkmark$\checkmark^\checkmark\\checkmarkc\checkmarki\checkmarkr\checkmarkc\checkmark$\checkmarkl\checkmark$\checkmark^\checkmark\\checkmarkc\checkmarki\checkmarkr\checkmarkc\checkmark$\checkmarka\checkmark$\checkmark^\checkmark\\checkmarkc\checkmarki\checkmarkr\checkmarkc\checkmark$\checkmark \checkmark$\checkmark^\checkmark\\checkmarkc\checkmarki\checkmarkr\checkmarkc\checkmark$\checkmarks\checkmark$\checkmark^\checkmark\\checkmarkc\checkmarki\checkmarkr\checkmarkc\checkmark$\checkmarkt\checkmark$\checkmark^\checkmark\\checkmarkc\checkmarki\checkmarkr\checkmarkc\checkmark$\checkmarkr\checkmark$\checkmark^\checkmark\\checkmarkc\checkmarki\checkmarkr\checkmarkc\checkmark$\checkmarku\checkmark$\checkmark^\checkmark\\checkmarkc\checkmarki\checkmarkr\checkmarkc\checkmark$\checkmarkc\checkmark$\checkmark^\checkmark\\checkmarkc\checkmarki\checkmarkr\checkmarkc\checkmark$\checkmarkt\checkmark$\checkmark^\checkmark\\checkmarkc\checkmarki\checkmarkr\checkmarkc\checkmark$\checkmarku\checkmark$\checkmark^\checkmark\\checkmarkc\checkmarki\checkmarkr\checkmarkc\checkmark$\checkmarkr\checkmark$\checkmark^\checkmark\\checkmarkc\checkmarki\checkmarkr\checkmarkc\checkmark$\checkmarke\checkmark$\checkmark^\checkmark\\checkmarkc\checkmarki\checkmarkr\checkmarkc\checkmark$\checkmark \checkmark$\checkmark^\checkmark\\checkmarkc\checkmarki\checkmarkr\checkmarkc\checkmark$\checkmark(\checkmark$\checkmark^\checkmark\\checkmarkc\checkmarki\checkmarkr\checkmarkc\checkmark$\checkmarkp\checkmark$\checkmark^\checkmark\\checkmarkc\checkmarki\checkmarkr\checkmarkc\checkmark$\checkmarkr\checkmark$\checkmark^\checkmark\\checkmarkc\checkmarki\checkmarkr\checkmarkc\checkmark$\checkmarke\checkmark$\checkmark^\checkmark\\checkmarkc\checkmarki\checkmarkr\checkmarkc\checkmark$\checkmarkm\checkmark$\checkmark^\checkmark\\checkmarkc\checkmarki\checkmarkr\checkmarkc\checkmark$\checkmarki\checkmark$\checkmark^\checkmark\\checkmarkc\checkmarki\checkmarkr\checkmarkc\checkmark$\checkmarke\checkmark$\checkmark^\checkmark\\checkmarkc\checkmarki\checkmarkr\checkmarkc\checkmark$\checkmarkr\checkmark$\checkmark^\checkmark\\checkmarkc\checkmarki\checkmarkr\checkmarkc\checkmark$\checkmark \checkmark$\checkmark^\checkmark\\checkmarkc\checkmarki\checkmarkr\checkmarkc\checkmark$\checkmarkm\checkmark$\checkmark^\checkmark\\checkmarkc\checkmarki\checkmarkr\checkmarkc\checkmark$\checkmarko\checkmark$\checkmark^\checkmark\\checkmarkc\checkmarki\checkmarkr\checkmarkc\checkmark$\checkmarkd\checkmark$\checkmark^\checkmark\\checkmarkc\checkmarki\checkmarkr\checkmarkc\checkmark$\checkmarke\checkmark$\checkmark^\checkmark\\checkmarkc\checkmarki\checkmarkr\checkmarkc\checkmark$\checkmark \checkmark$\checkmark^\checkmark\\checkmarkc\checkmarki\checkmarkr\checkmarkc\checkmark$\checkmarkd\checkmark$\checkmark^\checkmark\\checkmarkc\checkmarki\checkmarkr\checkmarkc\checkmark$\checkmarke\checkmark$\checkmark^\checkmark\\checkmarkc\checkmarki\checkmarkr\checkmarkc\checkmark$\checkmark \checkmark$\checkmark^\checkmark\\checkmarkc\checkmarki\checkmarkr\checkmarkc\checkmark$\checkmarkv\checkmark$\checkmark^\checkmark\\checkmarkc\checkmarki\checkmarkr\checkmarkc\checkmark$\checkmarki\checkmark$\checkmark^\checkmark\\checkmarkc\checkmarki\checkmarkr\checkmarkc\checkmark$\checkmarkb\checkmark$\checkmark^\checkmark\\checkmarkc\checkmarki\checkmarkr\checkmarkc\checkmark$\checkmarkr\checkmark$\checkmark^\checkmark\\checkmarkc\checkmarki\checkmarkr\checkmarkc\checkmark$\checkmarka\checkmark$\checkmark^\checkmark\\checkmarkc\checkmarki\checkmarkr\checkmarkc\checkmark$\checkmarkt\checkmark$\checkmark^\checkmark\\checkmarkc\checkmarki\checkmarkr\checkmarkc\checkmark$\checkmarki\checkmark$\checkmark^\checkmark\\checkmarkc\checkmarki\checkmarkr\checkmarkc\checkmark$\checkmarko\checkmark$\checkmark^\checkmark\\checkmarkc\checkmarki\checkmarkr\checkmarkc\checkmark$\checkmarkn\checkmark$\checkmark^\checkmark\\checkmarkc\checkmarki\checkmarkr\checkmarkc\checkmark$\checkmark)\checkmark$\checkmark^\checkmark\\checkmarkc\checkmarki\checkmarkr\checkmarkc\checkmark$\checkmark \checkmark$\checkmark^\checkmark\\checkmarkc\checkmarki\checkmarkr\checkmarkc\checkmark$\checkmarkd\checkmark$\checkmark^\checkmark\\checkmarkc\checkmarki\checkmarkr\checkmarkc\checkmark$\checkmarko\checkmark$\checkmark^\checkmark\\checkmarkc\checkmarki\checkmarkr\checkmarkc\checkmark$\checkmarki\checkmark$\checkmark^\checkmark\\checkmarkc\checkmarki\checkmarkr\checkmarkc\checkmark$\checkmarkt\checkmark$\checkmark^\checkmark\\checkmarkc\checkmarki\checkmarkr\checkmarkc\checkmark$\checkmark \checkmark$\checkmark^\checkmark\\checkmarkc\checkmarki\checkmarkr\checkmarkc\checkmark$\checkmarkê\checkmark$\checkmark^\checkmark\\checkmarkc\checkmarki\checkmarkr\checkmarkc\checkmark$\checkmarkt\checkmark$\checkmark^\checkmark\\checkmarkc\checkmarki\checkmarkr\checkmarkc\checkmark$\checkmarkr\checkmark$\checkmark^\checkmark\\checkmarkc\checkmarki\checkmarkr\checkmarkc\checkmark$\checkmarke\checkmark$\checkmark^\checkmark\\checkmarkc\checkmarki\checkmarkr\checkmarkc\checkmark$\checkmark \checkmark$\checkmark^\checkmark\\checkmarkc\checkmarki\checkmarkr\checkmarkc\checkmark$\checkmarké\checkmark$\checkmark^\checkmark\\checkmarkc\checkmarki\checkmarkr\checkmarkc\checkmark$\checkmarkl\checkmark$\checkmark^\checkmark\\checkmarkc\checkmarki\checkmarkr\checkmarkc\checkmark$\checkmarko\checkmark$\checkmark^\checkmark\\checkmarkc\checkmarki\checkmarkr\checkmarkc\checkmark$\checkmarki\checkmark$\checkmark^\checkmark\\checkmarkc\checkmarki\checkmarkr\checkmarkc\checkmark$\checkmarkg\checkmark$\checkmark^\checkmark\\checkmarkc\checkmarki\checkmarkr\checkmarkc\checkmark$\checkmarkn\checkmark$\checkmark^\checkmark\\checkmarkc\checkmarki\checkmarkr\checkmarkc\checkmark$\checkmarké\checkmark$\checkmark^\checkmark\\checkmarkc\checkmarki\checkmarkr\checkmarkc\checkmark$\checkmarke\checkmark$\checkmark^\checkmark\\checkmarkc\checkmarki\checkmarkr\checkmarkc\checkmark$\checkmark \checkmark$\checkmark^\checkmark\\checkmarkc\checkmarki\checkmarkr\checkmarkc\checkmark$\checkmarkd\checkmark$\checkmark^\checkmark\\checkmarkc\checkmarki\checkmarkr\checkmarkc\checkmark$\checkmarke\checkmark$\checkmark^\checkmark\\checkmarkc\checkmarki\checkmarkr\checkmarkc\checkmark$\checkmarks\checkmark$\checkmark^\checkmark\\checkmarkc\checkmarki\checkmarkr\checkmarkc\checkmark$\checkmark \checkmark$\checkmark^\checkmark\\checkmarkc\checkmarki\checkmarkr\checkmarkc\checkmark$\checkmarkf\checkmark$\checkmark^\checkmark\\checkmarkc\checkmarki\checkmarkr\checkmarkc\checkmark$\checkmarkr\checkmark$\checkmark^\checkmark\\checkmarkc\checkmarki\checkmarkr\checkmarkc\checkmark$\checkmarké\checkmark$\checkmark^\checkmark\\checkmarkc\checkmarki\checkmarkr\checkmarkc\checkmark$\checkmarkq\checkmark$\checkmark^\checkmark\\checkmarkc\checkmarki\checkmarkr\checkmarkc\checkmark$\checkmarku\checkmark$\checkmark^\checkmark\\checkmarkc\checkmarki\checkmarkr\checkmarkc\checkmark$\checkmarke\checkmark$\checkmark^\checkmark\\checkmarkc\checkmarki\checkmarkr\checkmarkc\checkmark$\checkmarkn\checkmark$\checkmark^\checkmark\\checkmarkc\checkmarki\checkmarkr\checkmarkc\checkmark$\checkmarkc\checkmark$\checkmark^\checkmark\\checkmarkc\checkmarki\checkmarkr\checkmarkc\checkmark$\checkmarke\checkmark$\checkmark^\checkmark\\checkmarkc\checkmarki\checkmarkr\checkmarkc\checkmark$\checkmarks\checkmark$\checkmark^\checkmark\\checkmarkc\checkmarki\checkmarkr\checkmarkc\checkmark$\checkmark \checkmark$\checkmark^\checkmark\\checkmarkc\checkmarki\checkmarkr\checkmarkc\checkmark$\checkmarkd\checkmark$\checkmark^\checkmark\\checkmarkc\checkmarki\checkmarkr\checkmarkc\checkmark$\checkmark'\checkmark$\checkmark^\checkmark\\checkmarkc\checkmarki\checkmarkr\checkmarkc\checkmark$\checkmarke\checkmark$\checkmark^\checkmark\\checkmarkc\checkmarki\checkmarkr\checkmarkc\checkmark$\checkmarkx\checkmark$\checkmark^\checkmark\\checkmarkc\checkmarki\checkmarkr\checkmarkc\checkmark$\checkmarkc\checkmark$\checkmark^\checkmark\\checkmarkc\checkmarki\checkmarkr\checkmarkc\checkmark$\checkmarki\checkmark$\checkmark^\checkmark\\checkmarkc\checkmarki\checkmarkr\checkmarkc\checkmark$\checkmarkt\checkmark$\checkmark^\checkmark\\checkmarkc\checkmarki\checkmarkr\checkmarkc\checkmark$\checkmarka\checkmark$\checkmark^\checkmark\\checkmarkc\checkmarki\checkmarkr\checkmarkc\checkmark$\checkmarkt\checkmark$\checkmark^\checkmark\\checkmarkc\checkmarki\checkmarkr\checkmarkc\checkmark$\checkmarki\checkmark$\checkmark^\checkmark\\checkmarkc\checkmarki\checkmarkr\checkmarkc\checkmark$\checkmarko\checkmark$\checkmark^\checkmark\\checkmarkc\checkmarki\checkmarkr\checkmarkc\checkmark$\checkmarkn\checkmark$\checkmark^\checkmark\\checkmarkc\checkmarki\checkmarkr\checkmarkc\checkmark$\checkmark \checkmark$\checkmark^\checkmark\\checkmarkc\checkmarki\checkmarkr\checkmarkc\checkmark$\checkmarkd\checkmark$\checkmark^\checkmark\\checkmarkc\checkmarki\checkmarkr\checkmarkc\checkmark$\checkmarke\checkmark$\checkmark^\checkmark\\checkmarkc\checkmarki\checkmarkr\checkmarkc\checkmark$\checkmark \checkmark$\checkmark^\checkmark\\checkmarkc\checkmarki\checkmarkr\checkmarkc\checkmark$\checkmarkl\checkmark$\checkmark^\checkmark\\checkmarkc\checkmarki\checkmarkr\checkmarkc\checkmark$\checkmarka\checkmark$\checkmark^\checkmark\\checkmarkc\checkmarki\checkmarkr\checkmarkc\checkmark$\checkmark \checkmark$\checkmark^\checkmark\\checkmarkc\checkmarki\checkmarkr\checkmarkc\checkmark$\checkmarkb\checkmark$\checkmark^\checkmark\\checkmarkc\checkmarki\checkmarkr\checkmarkc\checkmark$\checkmarkr\checkmark$\checkmark^\checkmark\\checkmarkc\checkmarki\checkmarkr\checkmarkc\checkmark$\checkmarko\checkmark$\checkmark^\checkmark\\checkmarkc\checkmarki\checkmarkr\checkmarkc\checkmark$\checkmarkc\checkmark$\checkmark^\checkmark\\checkmarkc\checkmarki\checkmarkr\checkmarkc\checkmark$\checkmarkh\checkmark$\checkmark^\checkmark\\checkmarkc\checkmarki\checkmarkr\checkmarkc\checkmark$\checkmarke\checkmark$\checkmark^\checkmark\\checkmarkc\checkmarki\checkmarkr\checkmarkc\checkmark$\checkmark.\checkmark$\checkmark^\checkmark\\checkmarkc\checkmarki\checkmarkr\checkmarkc\checkmark$\checkmark \checkmark$\checkmark^\checkmark\\checkmarkc\checkmarki\checkmarkr\checkmarkc\checkmark$\checkmarkP\checkmark$\checkmark^\checkmark\\checkmarkc\checkmarki\checkmarkr\checkmarkc\checkmark$\checkmarko\checkmark$\checkmark^\checkmark\\checkmarkc\checkmarki\checkmarkr\checkmarkc\checkmark$\checkmarku\checkmark$\checkmark^\checkmark\\checkmarkc\checkmarki\checkmarkr\checkmarkc\checkmark$\checkmarkr\checkmark$\checkmark^\checkmark\\checkmarkc\checkmarki\checkmarkr\checkmarkc\checkmark$\checkmark \checkmark$\checkmark^\checkmark\\checkmarkc\checkmarki\checkmarkr\checkmarkc\checkmark$\checkmarku\checkmark$\checkmark^\checkmark\\checkmarkc\checkmarki\checkmarkr\checkmarkc\checkmark$\checkmarkn\checkmark$\checkmark^\checkmark\\checkmarkc\checkmarki\checkmarkr\checkmarkc\checkmark$\checkmarke\checkmark$\checkmark^\checkmark\\checkmarkc\checkmarki\checkmarkr\checkmarkc\checkmark$\checkmark \checkmark$\checkmark^\checkmark\\checkmarkc\checkmarki\checkmarkr\checkmarkc\checkmark$\checkmarkb\checkmark$\checkmark^\checkmark\\checkmarkc\checkmarki\checkmarkr\checkmarkc\checkmark$\checkmarkr\checkmark$\checkmark^\checkmark\\checkmarkc\checkmarki\checkmarkr\checkmarkc\checkmark$\checkmarko\checkmark$\checkmark^\checkmark\\checkmarkc\checkmarki\checkmarkr\checkmarkc\checkmark$\checkmarkc\checkmark$\checkmark^\checkmark\\checkmarkc\checkmarki\checkmarkr\checkmarkc\checkmark$\checkmarkh\checkmark$\checkmark^\checkmark\\checkmarkc\checkmarki\checkmarkr\checkmarkc\checkmark$\checkmarke\checkmark$\checkmark^\checkmark\\checkmarkc\checkmarki\checkmarkr\checkmarkc\checkmark$\checkmark \checkmark$\checkmark^\checkmark\\checkmarkc\checkmarki\checkmarkr\checkmarkc\checkmark$\checkmarkt\checkmark$\checkmark^\checkmark\\checkmarkc\checkmarki\checkmarkr\checkmarkc\checkmark$\checkmarko\checkmark$\checkmark^\checkmark\\checkmarkc\checkmarki\checkmarkr\checkmarkc\checkmark$\checkmarku\checkmark$\checkmark^\checkmark\\checkmarkc\checkmarki\checkmarkr\checkmarkc\checkmark$\checkmarkr\checkmark$\checkmark^\checkmark\\checkmarkc\checkmarki\checkmarkr\checkmarkc\checkmark$\checkmarkn\checkmark$\checkmark^\checkmark\\checkmarkc\checkmarki\checkmarkr\checkmarkc\checkmark$\checkmarka\checkmark$\checkmark^\checkmark\\checkmarkc\checkmarki\checkmarkr\checkmarkc\checkmark$\checkmarkn\checkmark$\checkmark^\checkmark\\checkmarkc\checkmarki\checkmarkr\checkmarkc\checkmark$\checkmarkt\checkmark$\checkmark^\checkmark\\checkmarkc\checkmarki\checkmarkr\checkmarkc\checkmark$\checkmark \checkmark$\checkmark^\checkmark\\checkmarkc\checkmarki\checkmarkr\checkmarkc\checkmark$\checkmarkà\checkmark$\checkmark^\checkmark\\checkmarkc\checkmarki\checkmarkr\checkmarkc\checkmark$\checkmark \checkmark$\checkmark^\checkmark\\checkmarkc\checkmarki\checkmarkr\checkmarkc\checkmark$\checkmark$\checkmark$\checkmark^\checkmark\\checkmarkc\checkmarki\checkmarkr\checkmarkc\checkmark$\checkmarkN\checkmark$\checkmark^\checkmark\\checkmarkc\checkmarki\checkmarkr\checkmarkc\checkmark$\checkmark \checkmark$\checkmark^\checkmark\\checkmarkc\checkmarki\checkmarkr\checkmarkc\checkmark$\checkmark=\checkmark$\checkmark^\checkmark\\checkmarkc\checkmarki\checkmarkr\checkmarkc\checkmark$\checkmark \checkmark$\checkmark^\checkmark\\checkmarkc\checkmarki\checkmarkr\checkmarkc\checkmark$\checkmark1\checkmark$\checkmark^\checkmark\\checkmarkc\checkmarki\checkmarkr\checkmarkc\checkmark$\checkmark8\checkmark$\checkmark^\checkmark\\checkmarkc\checkmarki\checkmarkr\checkmarkc\checkmark$\checkmark\\checkmark$\checkmark^\checkmark\\checkmarkc\checkmarki\checkmarkr\checkmarkc\checkmark$\checkmark,\checkmark$\checkmark^\checkmark\\checkmarkc\checkmarki\checkmarkr\checkmarkc\checkmark$\checkmark0\checkmark$\checkmark^\checkmark\\checkmarkc\checkmarki\checkmarkr\checkmarkc\checkmark$\checkmark0\checkmark$\checkmark^\checkmark\\checkmarkc\checkmarki\checkmarkr\checkmarkc\checkmark$\checkmark0\checkmark$\checkmark^\checkmark\\checkmarkc\checkmarki\checkmarkr\checkmarkc\checkmark$\checkmark$\checkmark$\checkmark^\checkmark\\checkmarkc\checkmarki\checkmarkr\checkmarkc\checkmark$\checkmark \checkmark$\checkmark^\checkmark\\checkmarkc\checkmarki\checkmarkr\checkmarkc\checkmark$\checkmarkt\checkmark$\checkmark^\checkmark\\checkmarkc\checkmarki\checkmarkr\checkmarkc\checkmark$\checkmarkr\checkmark$\checkmark^\checkmark\\checkmarkc\checkmarki\checkmarkr\checkmarkc\checkmark$\checkmark/\checkmark$\checkmark^\checkmark\\checkmarkc\checkmarki\checkmarkr\checkmarkc\checkmark$\checkmarkm\checkmark$\checkmark^\checkmark\\checkmarkc\checkmarki\checkmarkr\checkmarkc\checkmark$\checkmarki\checkmark$\checkmark^\checkmark\\checkmarkc\checkmarki\checkmarkr\checkmarkc\checkmark$\checkmarkn\checkmark$\checkmark^\checkmark\\checkmarkc\checkmarki\checkmarkr\checkmarkc\checkmark$\checkmark \checkmark$\checkmark^\checkmark\\checkmarkc\checkmarki\checkmarkr\checkmarkc\checkmark$\checkmarka\checkmark$\checkmark^\checkmark\\checkmarkc\checkmarki\checkmarkr\checkmarkc\checkmark$\checkmarkv\checkmark$\checkmark^\checkmark\\checkmarkc\checkmarki\checkmarkr\checkmarkc\checkmark$\checkmarke\checkmark$\checkmark^\checkmark\\checkmarkc\checkmarki\checkmarkr\checkmarkc\checkmark$\checkmarkc\checkmark$\checkmark^\checkmark\\checkmarkc\checkmarki\checkmarkr\checkmarkc\checkmark$\checkmark \checkmark$\checkmark^\checkmark\\checkmarkc\checkmarki\checkmarkr\checkmarkc\checkmark$\checkmarku\checkmark$\checkmark^\checkmark\\checkmarkc\checkmarki\checkmarkr\checkmarkc\checkmark$\checkmarkn\checkmark$\checkmark^\checkmark\\checkmarkc\checkmarki\checkmarkr\checkmarkc\checkmark$\checkmarke\checkmark$\checkmark^\checkmark\\checkmarkc\checkmarki\checkmarkr\checkmarkc\checkmark$\checkmark \checkmark$\checkmark^\checkmark\\checkmarkc\checkmarki\checkmarkr\checkmarkc\checkmark$\checkmarkf\checkmark$\checkmark^\checkmark\\checkmarkc\checkmarki\checkmarkr\checkmarkc\checkmark$\checkmarkr\checkmark$\checkmark^\checkmark\\checkmarkc\checkmarki\checkmarkr\checkmarkc\checkmark$\checkmarka\checkmark$\checkmark^\checkmark\\checkmarkc\checkmarki\checkmarkr\checkmarkc\checkmark$\checkmarki\checkmark$\checkmark^\checkmark\\checkmarkc\checkmarki\checkmarkr\checkmarkc\checkmark$\checkmarks\checkmark$\checkmark^\checkmark\\checkmarkc\checkmarki\checkmarkr\checkmarkc\checkmark$\checkmarke\checkmark$\checkmark^\checkmark\\checkmarkc\checkmarki\checkmarkr\checkmarkc\checkmark$\checkmark \checkmark$\checkmark^\checkmark\\checkmarkc\checkmarki\checkmarkr\checkmarkc\checkmark$\checkmarkà\checkmark$\checkmark^\checkmark\\checkmarkc\checkmarki\checkmarkr\checkmarkc\checkmark$\checkmark \checkmark$\checkmark^\checkmark\\checkmarkc\checkmarki\checkmarkr\checkmarkc\checkmark$\checkmark2\checkmark$\checkmark^\checkmark\\checkmarkc\checkmarki\checkmarkr\checkmarkc\checkmark$\checkmark \checkmark$\checkmark^\checkmark\\checkmarkc\checkmarki\checkmarkr\checkmarkc\checkmark$\checkmarkd\checkmark$\checkmark^\checkmark\\checkmarkc\checkmarki\checkmarkr\checkmarkc\checkmark$\checkmarke\checkmark$\checkmark^\checkmark\\checkmarkc\checkmarki\checkmarkr\checkmarkc\checkmark$\checkmarkn\checkmark$\checkmark^\checkmark\\checkmarkc\checkmarki\checkmarkr\checkmarkc\checkmark$\checkmarkt\checkmark$\checkmark^\checkmark\\checkmarkc\checkmarki\checkmarkr\checkmarkc\checkmark$\checkmarks\checkmark$\checkmark^\checkmark\\checkmarkc\checkmarki\checkmarkr\checkmarkc\checkmark$\checkmark \checkmark$\checkmark^\checkmark\\checkmarkc\checkmarki\checkmarkr\checkmarkc\checkmark$\checkmark:\checkmark$\checkmark^\checkmark\\checkmarkc\checkmarki\checkmarkr\checkmarkc\checkmark$\checkmark
\checkmark$\checkmark^\checkmark\\checkmarkc\checkmarki\checkmarkr\checkmarkc\checkmark$\checkmark\\checkmark$\checkmark^\checkmark\\checkmarkc\checkmarki\checkmarkr\checkmarkc\checkmark$\checkmarkb\checkmark$\checkmark^\checkmark\\checkmarkc\checkmarki\checkmarkr\checkmarkc\checkmark$\checkmarke\checkmark$\checkmark^\checkmark\\checkmarkc\checkmarki\checkmarkr\checkmarkc\checkmark$\checkmarkg\checkmark$\checkmark^\checkmark\\checkmarkc\checkmarki\checkmarkr\checkmarkc\checkmark$\checkmarki\checkmark$\checkmark^\checkmark\\checkmarkc\checkmarki\checkmarkr\checkmarkc\checkmark$\checkmarkn\checkmark$\checkmark^\checkmark\\checkmarkc\checkmarki\checkmarkr\checkmarkc\checkmark$\checkmark{\checkmark$\checkmark^\checkmark\\checkmarkc\checkmarki\checkmarkr\checkmarkc\checkmark$\checkmarke\checkmark$\checkmark^\checkmark\\checkmarkc\checkmarki\checkmarkr\checkmarkc\checkmark$\checkmarkq\checkmark$\checkmark^\checkmark\\checkmarkc\checkmarki\checkmarkr\checkmarkc\checkmark$\checkmarku\checkmark$\checkmark^\checkmark\\checkmarkc\checkmarki\checkmarkr\checkmarkc\checkmark$\checkmarka\checkmark$\checkmark^\checkmark\\checkmarkc\checkmarki\checkmarkr\checkmarkc\checkmark$\checkmarkt\checkmark$\checkmark^\checkmark\\checkmarkc\checkmarki\checkmarkr\checkmarkc\checkmark$\checkmarki\checkmark$\checkmark^\checkmark\\checkmarkc\checkmarki\checkmarkr\checkmarkc\checkmark$\checkmarko\checkmark$\checkmark^\checkmark\\checkmarkc\checkmarki\checkmarkr\checkmarkc\checkmark$\checkmarkn\checkmark$\checkmark^\checkmark\\checkmarkc\checkmarki\checkmarkr\checkmarkc\checkmark$\checkmark}\checkmark$\checkmark^\checkmark\\checkmarkc\checkmarki\checkmarkr\checkmarkc\checkmark$\checkmark
\checkmark$\checkmark^\checkmark\\checkmarkc\checkmarki\checkmarkr\checkmarkc\checkmark$\checkmark \checkmark$\checkmark^\checkmark\\checkmarkc\checkmarki\checkmarkr\checkmarkc\checkmark$\checkmark \checkmark$\checkmark^\checkmark\\checkmarkc\checkmarki\checkmarkr\checkmarkc\checkmark$\checkmark \checkmark$\checkmark^\checkmark\\checkmarkc\checkmarki\checkmarkr\checkmarkc\checkmark$\checkmark \checkmark$\checkmark^\checkmark\\checkmarkc\checkmarki\checkmarkr\checkmarkc\checkmark$\checkmarkf\checkmark$\checkmark^\checkmark\\checkmarkc\checkmarki\checkmarkr\checkmarkc\checkmark$\checkmark_\checkmark$\checkmark^\checkmark\\checkmarkc\checkmarki\checkmarkr\checkmarkc\checkmark$\checkmark{\checkmark$\checkmark^\checkmark\\checkmarkc\checkmarki\checkmarkr\checkmarkc\checkmark$\checkmarke\checkmark$\checkmark^\checkmark\\checkmarkc\checkmarki\checkmarkr\checkmarkc\checkmark$\checkmarkx\checkmark$\checkmark^\checkmark\\checkmarkc\checkmarki\checkmarkr\checkmarkc\checkmark$\checkmarkc\checkmark$\checkmark^\checkmark\\checkmarkc\checkmarki\checkmarkr\checkmarkc\checkmark$\checkmarki\checkmark$\checkmark^\checkmark\\checkmarkc\checkmarki\checkmarkr\checkmarkc\checkmark$\checkmarkt\checkmark$\checkmark^\checkmark\\checkmarkc\checkmarki\checkmarkr\checkmarkc\checkmark$\checkmarka\checkmark$\checkmark^\checkmark\\checkmarkc\checkmarki\checkmarkr\checkmarkc\checkmark$\checkmarkt\checkmark$\checkmark^\checkmark\\checkmarkc\checkmarki\checkmarkr\checkmarkc\checkmark$\checkmarki\checkmark$\checkmark^\checkmark\\checkmarkc\checkmarki\checkmarkr\checkmarkc\checkmark$\checkmarko\checkmark$\checkmark^\checkmark\\checkmarkc\checkmarki\checkmarkr\checkmarkc\checkmark$\checkmarkn\checkmark$\checkmark^\checkmark\\checkmarkc\checkmarki\checkmarkr\checkmarkc\checkmark$\checkmark}\checkmark$\checkmark^\checkmark\\checkmarkc\checkmarki\checkmarkr\checkmarkc\checkmark$\checkmark \checkmark$\checkmark^\checkmark\\checkmarkc\checkmarki\checkmarkr\checkmarkc\checkmark$\checkmark=\checkmark$\checkmark^\checkmark\\checkmarkc\checkmarki\checkmarkr\checkmarkc\checkmark$\checkmark \checkmark$\checkmark^\checkmark\\checkmarkc\checkmarki\checkmarkr\checkmarkc\checkmark$\checkmark\\checkmark$\checkmark^\checkmark\\checkmarkc\checkmarki\checkmarkr\checkmarkc\checkmark$\checkmarkf\checkmark$\checkmark^\checkmark\\checkmarkc\checkmarki\checkmarkr\checkmarkc\checkmark$\checkmarkr\checkmark$\checkmark^\checkmark\\checkmarkc\checkmarki\checkmarkr\checkmarkc\checkmark$\checkmarka\checkmark$\checkmark^\checkmark\\checkmarkc\checkmarki\checkmarkr\checkmarkc\checkmark$\checkmarkc\checkmark$\checkmark^\checkmark\\checkmarkc\checkmarki\checkmarkr\checkmarkc\checkmark$\checkmark{\checkmark$\checkmark^\checkmark\\checkmarkc\checkmarki\checkmarkr\checkmarkc\checkmark$\checkmarkN\checkmark$\checkmark^\checkmark\\checkmarkc\checkmarki\checkmarkr\checkmarkc\checkmark$\checkmark \checkmark$\checkmark^\checkmark\\checkmarkc\checkmarki\checkmarkr\checkmarkc\checkmark$\checkmark\\checkmark$\checkmark^\checkmark\\checkmarkc\checkmarki\checkmarkr\checkmarkc\checkmark$\checkmarkt\checkmark$\checkmark^\checkmark\\checkmarkc\checkmarki\checkmarkr\checkmarkc\checkmark$\checkmarki\checkmark$\checkmark^\checkmark\\checkmarkc\checkmarki\checkmarkr\checkmarkc\checkmark$\checkmarkm\checkmark$\checkmark^\checkmark\\checkmarkc\checkmarki\checkmarkr\checkmarkc\checkmark$\checkmarke\checkmark$\checkmark^\checkmark\\checkmarkc\checkmarki\checkmarkr\checkmarkc\checkmark$\checkmarks\checkmark$\checkmark^\checkmark\\checkmarkc\checkmarki\checkmarkr\checkmarkc\checkmark$\checkmark \checkmark$\checkmark^\checkmark\\checkmarkc\checkmarki\checkmarkr\checkmarkc\checkmark$\checkmarkZ\checkmark$\checkmark^\checkmark\\checkmarkc\checkmarki\checkmarkr\checkmarkc\checkmark$\checkmark_\checkmark$\checkmark^\checkmark\\checkmarkc\checkmarki\checkmarkr\checkmarkc\checkmark$\checkmarkc\checkmark$\checkmark^\checkmark\\checkmarkc\checkmarki\checkmarkr\checkmarkc\checkmark$\checkmark}\checkmark$\checkmark^\checkmark\\checkmarkc\checkmarki\checkmarkr\checkmarkc\checkmark$\checkmark{\checkmark$\checkmark^\checkmark\\checkmarkc\checkmarki\checkmarkr\checkmarkc\checkmark$\checkmark6\checkmark$\checkmark^\checkmark\\checkmarkc\checkmarki\checkmarkr\checkmarkc\checkmark$\checkmark0\checkmark$\checkmark^\checkmark\\checkmarkc\checkmarki\checkmarkr\checkmarkc\checkmark$\checkmark}\checkmark$\checkmark^\checkmark\\checkmarkc\checkmarki\checkmarkr\checkmarkc\checkmark$\checkmark \checkmark$\checkmark^\checkmark\\checkmarkc\checkmarki\checkmarkr\checkmarkc\checkmark$\checkmark=\checkmark$\checkmark^\checkmark\\checkmarkc\checkmarki\checkmarkr\checkmarkc\checkmark$\checkmark \checkmark$\checkmark^\checkmark\\checkmarkc\checkmarki\checkmarkr\checkmarkc\checkmark$\checkmark\\checkmark$\checkmark^\checkmark\\checkmarkc\checkmarki\checkmarkr\checkmarkc\checkmark$\checkmarkf\checkmark$\checkmark^\checkmark\\checkmarkc\checkmarki\checkmarkr\checkmarkc\checkmark$\checkmarkr\checkmark$\checkmark^\checkmark\\checkmarkc\checkmarki\checkmarkr\checkmarkc\checkmark$\checkmarka\checkmark$\checkmark^\checkmark\\checkmarkc\checkmarki\checkmarkr\checkmarkc\checkmark$\checkmarkc\checkmark$\checkmark^\checkmark\\checkmarkc\checkmarki\checkmarkr\checkmarkc\checkmark$\checkmark{\checkmark$\checkmark^\checkmark\\checkmarkc\checkmarki\checkmarkr\checkmarkc\checkmark$\checkmark1\checkmark$\checkmark^\checkmark\\checkmarkc\checkmarki\checkmarkr\checkmarkc\checkmark$\checkmark8\checkmark$\checkmark^\checkmark\\checkmarkc\checkmarki\checkmarkr\checkmarkc\checkmark$\checkmark\\checkmark$\checkmark^\checkmark\\checkmarkc\checkmarki\checkmarkr\checkmarkc\checkmark$\checkmark,\checkmark$\checkmark^\checkmark\\checkmarkc\checkmarki\checkmarkr\checkmarkc\checkmark$\checkmark0\checkmark$\checkmark^\checkmark\\checkmarkc\checkmarki\checkmarkr\checkmarkc\checkmark$\checkmark0\checkmark$\checkmark^\checkmark\\checkmarkc\checkmarki\checkmarkr\checkmarkc\checkmark$\checkmark0\checkmark$\checkmark^\checkmark\\checkmarkc\checkmarki\checkmarkr\checkmarkc\checkmark$\checkmark \checkmark$\checkmark^\checkmark\\checkmarkc\checkmarki\checkmarkr\checkmarkc\checkmark$\checkmark\\checkmark$\checkmark^\checkmark\\checkmarkc\checkmarki\checkmarkr\checkmarkc\checkmark$\checkmarkt\checkmark$\checkmark^\checkmark\\checkmarkc\checkmarki\checkmarkr\checkmarkc\checkmark$\checkmarki\checkmark$\checkmark^\checkmark\\checkmarkc\checkmarki\checkmarkr\checkmarkc\checkmark$\checkmarkm\checkmark$\checkmark^\checkmark\\checkmarkc\checkmarki\checkmarkr\checkmarkc\checkmark$\checkmarke\checkmark$\checkmark^\checkmark\\checkmarkc\checkmarki\checkmarkr\checkmarkc\checkmark$\checkmarks\checkmark$\checkmark^\checkmark\\checkmarkc\checkmarki\checkmarkr\checkmarkc\checkmark$\checkmark \checkmark$\checkmark^\checkmark\\checkmarkc\checkmarki\checkmarkr\checkmarkc\checkmark$\checkmark2\checkmark$\checkmark^\checkmark\\checkmarkc\checkmarki\checkmarkr\checkmarkc\checkmark$\checkmark}\checkmark$\checkmark^\checkmark\\checkmarkc\checkmarki\checkmarkr\checkmarkc\checkmark$\checkmark{\checkmark$\checkmark^\checkmark\\checkmarkc\checkmarki\checkmarkr\checkmarkc\checkmark$\checkmark6\checkmark$\checkmark^\checkmark\\checkmarkc\checkmarki\checkmarkr\checkmarkc\checkmark$\checkmark0\checkmark$\checkmark^\checkmark\\checkmarkc\checkmarki\checkmarkr\checkmarkc\checkmark$\checkmark}\checkmark$\checkmark^\checkmark\\checkmarkc\checkmarki\checkmarkr\checkmarkc\checkmark$\checkmark \checkmark$\checkmark^\checkmark\\checkmarkc\checkmarki\checkmarkr\checkmarkc\checkmark$\checkmark=\checkmark$\checkmark^\checkmark\\checkmarkc\checkmarki\checkmarkr\checkmarkc\checkmark$\checkmark \checkmark$\checkmark^\checkmark\\checkmarkc\checkmarki\checkmarkr\checkmarkc\checkmark$\checkmark\\checkmark$\checkmark^\checkmark\\checkmarkc\checkmarki\checkmarkr\checkmarkc\checkmark$\checkmarkb\checkmark$\checkmark^\checkmark\\checkmarkc\checkmarki\checkmarkr\checkmarkc\checkmark$\checkmarko\checkmark$\checkmark^\checkmark\\checkmarkc\checkmarki\checkmarkr\checkmarkc\checkmark$\checkmarkx\checkmark$\checkmark^\checkmark\\checkmarkc\checkmarki\checkmarkr\checkmarkc\checkmark$\checkmarke\checkmark$\checkmark^\checkmark\\checkmarkc\checkmarki\checkmarkr\checkmarkc\checkmark$\checkmarkd\checkmark$\checkmark^\checkmark\\checkmarkc\checkmarki\checkmarkr\checkmarkc\checkmark$\checkmark{\checkmark$\checkmark^\checkmark\\checkmarkc\checkmarki\checkmarkr\checkmarkc\checkmark$\checkmark6\checkmark$\checkmark^\checkmark\\checkmarkc\checkmarki\checkmarkr\checkmarkc\checkmark$\checkmark0\checkmark$\checkmark^\checkmark\\checkmarkc\checkmarki\checkmarkr\checkmarkc\checkmark$\checkmark0\checkmark$\checkmark^\checkmark\\checkmarkc\checkmarki\checkmarkr\checkmarkc\checkmark$\checkmark \checkmark$\checkmark^\checkmark\\checkmarkc\checkmarki\checkmarkr\checkmarkc\checkmark$\checkmark\\checkmark$\checkmark^\checkmark\\checkmarkc\checkmarki\checkmarkr\checkmarkc\checkmark$\checkmarkt\checkmark$\checkmark^\checkmark\\checkmarkc\checkmarki\checkmarkr\checkmarkc\checkmark$\checkmarke\checkmark$\checkmark^\checkmark\\checkmarkc\checkmarki\checkmarkr\checkmarkc\checkmark$\checkmarkx\checkmark$\checkmark^\checkmark\\checkmarkc\checkmarki\checkmarkr\checkmarkc\checkmark$\checkmarkt\checkmark$\checkmark^\checkmark\\checkmarkc\checkmarki\checkmarkr\checkmarkc\checkmark$\checkmark{\checkmark$\checkmark^\checkmark\\checkmarkc\checkmarki\checkmarkr\checkmarkc\checkmark$\checkmark \checkmark$\checkmark^\checkmark\\checkmarkc\checkmarki\checkmarkr\checkmarkc\checkmark$\checkmarkH\checkmark$\checkmark^\checkmark\\checkmarkc\checkmarki\checkmarkr\checkmarkc\checkmark$\checkmarkz\checkmark$\checkmark^\checkmark\\checkmarkc\checkmarki\checkmarkr\checkmarkc\checkmark$\checkmark}\checkmark$\checkmark^\checkmark\\checkmarkc\checkmarki\checkmarkr\checkmarkc\checkmark$\checkmark}\checkmark$\checkmark^\checkmark\\checkmarkc\checkmarki\checkmarkr\checkmarkc\checkmark$\checkmark
\checkmark$\checkmark^\checkmark\\checkmarkc\checkmarki\checkmarkr\checkmarkc\checkmark$\checkmark \checkmark$\checkmark^\checkmark\\checkmarkc\checkmarki\checkmarkr\checkmarkc\checkmark$\checkmark \checkmark$\checkmark^\checkmark\\checkmarkc\checkmarki\checkmarkr\checkmarkc\checkmark$\checkmark \checkmark$\checkmark^\checkmark\\checkmarkc\checkmarki\checkmarkr\checkmarkc\checkmark$\checkmark \checkmark$\checkmark^\checkmark\\checkmarkc\checkmarki\checkmarkr\checkmarkc\checkmark$\checkmark\\checkmark$\checkmark^\checkmark\\checkmarkc\checkmarki\checkmarkr\checkmarkc\checkmark$\checkmarkl\checkmark$\checkmark^\checkmark\\checkmarkc\checkmarki\checkmarkr\checkmarkc\checkmark$\checkmarka\checkmark$\checkmark^\checkmark\\checkmarkc\checkmarki\checkmarkr\checkmarkc\checkmark$\checkmarkb\checkmark$\checkmark^\checkmark\\checkmarkc\checkmarki\checkmarkr\checkmarkc\checkmark$\checkmarke\checkmark$\checkmark^\checkmark\\checkmarkc\checkmarki\checkmarkr\checkmarkc\checkmark$\checkmarkl\checkmark$\checkmark^\checkmark\\checkmarkc\checkmarki\checkmarkr\checkmarkc\checkmark$\checkmark{\checkmark$\checkmark^\checkmark\\checkmarkc\checkmarki\checkmarkr\checkmarkc\checkmark$\checkmarke\checkmark$\checkmark^\checkmark\\checkmarkc\checkmarki\checkmarkr\checkmarkc\checkmark$\checkmarkq\checkmark$\checkmark^\checkmark\\checkmarkc\checkmarki\checkmarkr\checkmarkc\checkmark$\checkmark:\checkmark$\checkmark^\checkmark\\checkmarkc\checkmarki\checkmarkr\checkmarkc\checkmark$\checkmarkf\checkmark$\checkmark^\checkmark\\checkmarkc\checkmarki\checkmarkr\checkmarkc\checkmark$\checkmarkr\checkmark$\checkmark^\checkmark\\checkmarkc\checkmarki\checkmarkr\checkmarkc\checkmark$\checkmarke\checkmark$\checkmark^\checkmark\\checkmarkc\checkmarki\checkmarkr\checkmarkc\checkmark$\checkmarkq\checkmark$\checkmark^\checkmark\\checkmarkc\checkmarki\checkmarkr\checkmarkc\checkmark$\checkmark_\checkmark$\checkmark^\checkmark\\checkmarkc\checkmarki\checkmarkr\checkmarkc\checkmark$\checkmarke\checkmark$\checkmark^\checkmark\\checkmarkc\checkmarki\checkmarkr\checkmarkc\checkmark$\checkmarkx\checkmark$\checkmark^\checkmark\\checkmarkc\checkmarki\checkmarkr\checkmarkc\checkmark$\checkmarkc\checkmark$\checkmark^\checkmark\\checkmarkc\checkmarki\checkmarkr\checkmarkc\checkmark$\checkmarki\checkmark$\checkmark^\checkmark\\checkmarkc\checkmarki\checkmarkr\checkmarkc\checkmark$\checkmarkt\checkmark$\checkmark^\checkmark\\checkmarkc\checkmarki\checkmarkr\checkmarkc\checkmark$\checkmarka\checkmark$\checkmark^\checkmark\\checkmarkc\checkmarki\checkmarkr\checkmarkc\checkmark$\checkmarkt\checkmark$\checkmark^\checkmark\\checkmarkc\checkmarki\checkmarkr\checkmarkc\checkmark$\checkmarki\checkmark$\checkmark^\checkmark\\checkmarkc\checkmarki\checkmarkr\checkmarkc\checkmark$\checkmarko\checkmark$\checkmark^\checkmark\\checkmarkc\checkmarki\checkmarkr\checkmarkc\checkmark$\checkmarkn\checkmark$\checkmark^\checkmark\\checkmarkc\checkmarki\checkmarkr\checkmarkc\checkmark$\checkmark}\checkmark$\checkmark^\checkmark\\checkmarkc\checkmarki\checkmarkr\checkmarkc\checkmark$\checkmark
\checkmark$\checkmark^\checkmark\\checkmarkc\checkmarki\checkmarkr\checkmarkc\checkmark$\checkmark\\checkmark$\checkmark^\checkmark\\checkmarkc\checkmarki\checkmarkr\checkmarkc\checkmark$\checkmarke\checkmark$\checkmark^\checkmark\\checkmarkc\checkmarki\checkmarkr\checkmarkc\checkmark$\checkmarkn\checkmark$\checkmark^\checkmark\\checkmarkc\checkmarki\checkmarkr\checkmarkc\checkmark$\checkmarkd\checkmark$\checkmark^\checkmark\\checkmarkc\checkmarki\checkmarkr\checkmarkc\checkmark$\checkmark{\checkmark$\checkmark^\checkmark\\checkmarkc\checkmarki\checkmarkr\checkmarkc\checkmark$\checkmarke\checkmark$\checkmark^\checkmark\\checkmarkc\checkmarki\checkmarkr\checkmarkc\checkmark$\checkmarkq\checkmark$\checkmark^\checkmark\\checkmarkc\checkmarki\checkmarkr\checkmarkc\checkmark$\checkmarku\checkmark$\checkmark^\checkmark\\checkmarkc\checkmarki\checkmarkr\checkmarkc\checkmark$\checkmarka\checkmark$\checkmark^\checkmark\\checkmarkc\checkmarki\checkmarkr\checkmarkc\checkmark$\checkmarkt\checkmark$\checkmark^\checkmark\\checkmarkc\checkmarki\checkmarkr\checkmarkc\checkmark$\checkmarki\checkmark$\checkmark^\checkmark\\checkmarkc\checkmarki\checkmarkr\checkmarkc\checkmark$\checkmarko\checkmark$\checkmark^\checkmark\\checkmarkc\checkmarki\checkmarkr\checkmarkc\checkmark$\checkmarkn\checkmark$\checkmark^\checkmark\\checkmarkc\checkmarki\checkmarkr\checkmarkc\checkmark$\checkmark}\checkmark$\checkmark^\checkmark\\checkmarkc\checkmarki\checkmarkr\checkmarkc\checkmark$\checkmark
\checkmark$\checkmark^\checkmark\\checkmarkc\checkmarki\checkmarkr\checkmarkc\checkmark$\checkmark
\checkmark$\checkmark^\checkmark\\checkmarkc\checkmarki\checkmarkr\checkmarkc\checkmark$\checkmarkL\checkmark$\checkmark^\checkmark\\checkmarkc\checkmarki\checkmarkr\checkmarkc\checkmark$\checkmarka\checkmark$\checkmark^\checkmark\\checkmarkc\checkmarki\checkmarkr\checkmarkc\checkmark$\checkmark \checkmark$\checkmark^\checkmark\\checkmarkc\checkmarki\checkmarkr\checkmarkc\checkmark$\checkmarkf\checkmark$\checkmark^\checkmark\\checkmarkc\checkmarki\checkmarkr\checkmarkc\checkmark$\checkmarkr\checkmark$\checkmark^\checkmark\\checkmarkc\checkmarki\checkmarkr\checkmarkc\checkmark$\checkmarké\checkmark$\checkmark^\checkmark\\checkmarkc\checkmarki\checkmarkr\checkmarkc\checkmark$\checkmarkq\checkmark$\checkmark^\checkmark\\checkmarkc\checkmarki\checkmarkr\checkmarkc\checkmark$\checkmarku\checkmark$\checkmark^\checkmark\\checkmarkc\checkmarki\checkmarkr\checkmarkc\checkmark$\checkmarke\checkmark$\checkmark^\checkmark\\checkmarkc\checkmarki\checkmarkr\checkmarkc\checkmark$\checkmarkn\checkmark$\checkmark^\checkmark\\checkmarkc\checkmarki\checkmarkr\checkmarkc\checkmark$\checkmarkc\checkmark$\checkmark^\checkmark\\checkmarkc\checkmarki\checkmarkr\checkmarkc\checkmark$\checkmarke\checkmark$\checkmark^\checkmark\\checkmarkc\checkmarki\checkmarkr\checkmarkc\checkmark$\checkmark \checkmark$\checkmark^\checkmark\\checkmarkc\checkmarki\checkmarkr\checkmarkc\checkmark$\checkmarkp\checkmark$\checkmark^\checkmark\\checkmarkc\checkmarki\checkmarkr\checkmarkc\checkmark$\checkmarkr\checkmark$\checkmark^\checkmark\\checkmarkc\checkmarki\checkmarkr\checkmarkc\checkmark$\checkmarko\checkmark$\checkmark^\checkmark\\checkmarkc\checkmarki\checkmarkr\checkmarkc\checkmark$\checkmarkp\checkmark$\checkmark^\checkmark\\checkmarkc\checkmarki\checkmarkr\checkmarkc\checkmark$\checkmarkr\checkmark$\checkmark^\checkmark\\checkmarkc\checkmarki\checkmarkr\checkmarkc\checkmark$\checkmarke\checkmark$\checkmark^\checkmark\\checkmarkc\checkmarki\checkmarkr\checkmarkc\checkmark$\checkmark \checkmark$\checkmark^\checkmark\\checkmarkc\checkmarki\checkmarkr\checkmarkc\checkmark$\checkmarkd\checkmark$\checkmark^\checkmark\\checkmarkc\checkmarki\checkmarkr\checkmarkc\checkmark$\checkmark'\checkmark$\checkmark^\checkmark\\checkmarkc\checkmarki\checkmarkr\checkmarkc\checkmark$\checkmarku\checkmark$\checkmark^\checkmark\\checkmarkc\checkmarki\checkmarkr\checkmarkc\checkmark$\checkmarkn\checkmark$\checkmark^\checkmark\\checkmarkc\checkmarki\checkmarkr\checkmarkc\checkmark$\checkmarke\checkmark$\checkmark^\checkmark\\checkmarkc\checkmarki\checkmarkr\checkmarkc\checkmark$\checkmark \checkmark$\checkmark^\checkmark\\checkmarkc\checkmarki\checkmarkr\checkmarkc\checkmark$\checkmarkp\checkmark$\checkmark^\checkmark\\checkmarkc\checkmarki\checkmarkr\checkmarkc\checkmark$\checkmarko\checkmark$\checkmark^\checkmark\\checkmarkc\checkmarki\checkmarkr\checkmarkc\checkmark$\checkmarku\checkmark$\checkmark^\checkmark\\checkmarkc\checkmarki\checkmarkr\checkmarkc\checkmark$\checkmarkt\checkmark$\checkmark^\checkmark\\checkmarkc\checkmarki\checkmarkr\checkmarkc\checkmark$\checkmarkr\checkmark$\checkmark^\checkmark\\checkmarkc\checkmarki\checkmarkr\checkmarkc\checkmark$\checkmarke\checkmark$\checkmark^\checkmark\\checkmarkc\checkmarki\checkmarkr\checkmarkc\checkmark$\checkmark \checkmark$\checkmark^\checkmark\\checkmarkc\checkmarki\checkmarkr\checkmarkc\checkmark$\checkmarke\checkmark$\checkmark^\checkmark\\checkmarkc\checkmarki\checkmarkr\checkmarkc\checkmark$\checkmarkn\checkmark$\checkmark^\checkmark\\checkmarkc\checkmarki\checkmarkr\checkmarkc\checkmark$\checkmarkc\checkmark$\checkmark^\checkmark\\checkmarkc\checkmarki\checkmarkr\checkmarkc\checkmark$\checkmarka\checkmark$\checkmark^\checkmark\\checkmarkc\checkmarki\checkmarkr\checkmarkc\checkmark$\checkmarks\checkmark$\checkmark^\checkmark\\checkmarkc\checkmarki\checkmarkr\checkmarkc\checkmark$\checkmarkt\checkmark$\checkmark^\checkmark\\checkmarkc\checkmarki\checkmarkr\checkmarkc\checkmark$\checkmarkr\checkmark$\checkmark^\checkmark\\checkmarkc\checkmarki\checkmarkr\checkmarkc\checkmark$\checkmarké\checkmark$\checkmark^\checkmark\\checkmarkc\checkmarki\checkmarkr\checkmarkc\checkmark$\checkmarke\checkmark$\checkmark^\checkmark\\checkmarkc\checkmarki\checkmarkr\checkmarkc\checkmark$\checkmark-\checkmark$\checkmark^\checkmark\\checkmarkc\checkmarki\checkmarkr\checkmarkc\checkmark$\checkmarkl\checkmark$\checkmark^\checkmark\\checkmarkc\checkmarki\checkmarkr\checkmarkc\checkmark$\checkmarki\checkmark$\checkmark^\checkmark\\checkmarkc\checkmarki\checkmarkr\checkmarkc\checkmark$\checkmarkb\checkmark$\checkmark^\checkmark\\checkmarkc\checkmarki\checkmarkr\checkmarkc\checkmark$\checkmarkr\checkmark$\checkmark^\checkmark\\checkmarkc\checkmarki\checkmarkr\checkmarkc\checkmark$\checkmarke\checkmark$\checkmark^\checkmark\\checkmarkc\checkmarki\checkmarkr\checkmarkc\checkmark$\checkmark \checkmark$\checkmark^\checkmark\\checkmarkc\checkmarki\checkmarkr\checkmarkc\checkmark$\checkmark(\checkmark$\checkmark^\checkmark\\checkmarkc\checkmarki\checkmarkr\checkmarkc\checkmark$\checkmarkp\checkmark$\checkmark^\checkmark\\checkmarkc\checkmarki\checkmarkr\checkmarkc\checkmark$\checkmarko\checkmark$\checkmark^\checkmark\\checkmarkc\checkmarki\checkmarkr\checkmarkc\checkmark$\checkmarkr\checkmark$\checkmark^\checkmark\\checkmarkc\checkmarki\checkmarkr\checkmarkc\checkmark$\checkmarkt\checkmark$\checkmark^\checkmark\\checkmarkc\checkmarki\checkmarkr\checkmarkc\checkmark$\checkmarki\checkmark$\checkmark^\checkmark\\checkmarkc\checkmarki\checkmarkr\checkmarkc\checkmark$\checkmarkq\checkmark$\checkmark^\checkmark\\checkmarkc\checkmarki\checkmarkr\checkmarkc\checkmark$\checkmarku\checkmark$\checkmark^\checkmark\\checkmarkc\checkmarki\checkmarkr\checkmarkc\checkmark$\checkmarke\checkmark$\checkmark^\checkmark\\checkmarkc\checkmarki\checkmarkr\checkmarkc\checkmark$\checkmark)\checkmark$\checkmark^\checkmark\\checkmarkc\checkmarki\checkmarkr\checkmarkc\checkmark$\checkmark \checkmark$\checkmark^\checkmark\\checkmarkc\checkmarki\checkmarkr\checkmarkc\checkmark$\checkmarke\checkmark$\checkmark^\checkmark\\checkmarkc\checkmarki\checkmarkr\checkmarkc\checkmark$\checkmarks\checkmark$\checkmark^\checkmark\\checkmarkc\checkmarki\checkmarkr\checkmarkc\checkmark$\checkmarkt\checkmark$\checkmark^\checkmark\\checkmarkc\checkmarki\checkmarkr\checkmarkc\checkmark$\checkmark \checkmark$\checkmark^\checkmark\\checkmarkc\checkmarki\checkmarkr\checkmarkc\checkmark$\checkmarka\checkmark$\checkmark^\checkmark\\checkmarkc\checkmarki\checkmarkr\checkmarkc\checkmark$\checkmarkp\checkmark$\checkmark^\checkmark\\checkmarkc\checkmarki\checkmarkr\checkmarkc\checkmark$\checkmarkp\checkmark$\checkmark^\checkmark\\checkmarkc\checkmarki\checkmarkr\checkmarkc\checkmark$\checkmarkr\checkmark$\checkmark^\checkmark\\checkmarkc\checkmarki\checkmarkr\checkmarkc\checkmark$\checkmarko\checkmark$\checkmark^\checkmark\\checkmarkc\checkmarki\checkmarkr\checkmarkc\checkmark$\checkmarkx\checkmark$\checkmark^\checkmark\\checkmarkc\checkmarki\checkmarkr\checkmarkc\checkmark$\checkmarki\checkmark$\checkmark^\checkmark\\checkmarkc\checkmarki\checkmarkr\checkmarkc\checkmark$\checkmarkm\checkmark$\checkmark^\checkmark\\checkmarkc\checkmarki\checkmarkr\checkmarkc\checkmark$\checkmarké\checkmark$\checkmark^\checkmark\\checkmarkc\checkmarki\checkmarkr\checkmarkc\checkmark$\checkmarke\checkmark$\checkmark^\checkmark\\checkmarkc\checkmarki\checkmarkr\checkmarkc\checkmark$\checkmark \checkmark$\checkmark^\checkmark\\checkmarkc\checkmarki\checkmarkr\checkmarkc\checkmark$\checkmarkp\checkmark$\checkmark^\checkmark\\checkmarkc\checkmarki\checkmarkr\checkmarkc\checkmark$\checkmarka\checkmark$\checkmark^\checkmark\\checkmarkc\checkmarki\checkmarkr\checkmarkc\checkmark$\checkmarkr\checkmark$\checkmark^\checkmark\\checkmarkc\checkmarki\checkmarkr\checkmarkc\checkmark$\checkmark \checkmark$\checkmark^\checkmark\\checkmarkc\checkmarki\checkmarkr\checkmarkc\checkmark$\checkmarkl\checkmark$\checkmark^\checkmark\\checkmarkc\checkmarki\checkmarkr\checkmarkc\checkmark$\checkmarka\checkmark$\checkmark^\checkmark\\checkmarkc\checkmarki\checkmarkr\checkmarkc\checkmark$\checkmark \checkmark$\checkmark^\checkmark\\checkmarkc\checkmarki\checkmarkr\checkmarkc\checkmark$\checkmarkm\checkmark$\checkmark^\checkmark\\checkmarkc\checkmarki\checkmarkr\checkmarkc\checkmark$\checkmarké\checkmark$\checkmark^\checkmark\\checkmarkc\checkmarki\checkmarkr\checkmarkc\checkmark$\checkmarkt\checkmark$\checkmark^\checkmark\\checkmarkc\checkmarki\checkmarkr\checkmarkc\checkmark$\checkmarkh\checkmark$\checkmark^\checkmark\\checkmarkc\checkmarki\checkmarkr\checkmarkc\checkmark$\checkmarko\checkmark$\checkmark^\checkmark\\checkmarkc\checkmarki\checkmarkr\checkmarkc\checkmark$\checkmarkd\checkmark$\checkmark^\checkmark\\checkmarkc\checkmarki\checkmarkr\checkmarkc\checkmark$\checkmarke\checkmark$\checkmark^\checkmark\\checkmarkc\checkmarki\checkmarkr\checkmarkc\checkmark$\checkmark \checkmark$\checkmark^\checkmark\\checkmarkc\checkmarki\checkmarkr\checkmarkc\checkmark$\checkmarkd\checkmark$\checkmark^\checkmark\\checkmarkc\checkmarki\checkmarkr\checkmarkc\checkmark$\checkmarke\checkmark$\checkmark^\checkmark\\checkmarkc\checkmarki\checkmarkr\checkmarkc\checkmark$\checkmark \checkmark$\checkmark^\checkmark\\checkmarkc\checkmarki\checkmarkr\checkmarkc\checkmark$\checkmarkR\checkmark$\checkmark^\checkmark\\checkmarkc\checkmarki\checkmarkr\checkmarkc\checkmark$\checkmarka\checkmark$\checkmark^\checkmark\\checkmarkc\checkmarki\checkmarkr\checkmarkc\checkmark$\checkmarky\checkmark$\checkmark^\checkmark\\checkmarkc\checkmarki\checkmarkr\checkmarkc\checkmark$\checkmarkl\checkmark$\checkmark^\checkmark\\checkmarkc\checkmarki\checkmarkr\checkmarkc\checkmark$\checkmarke\checkmark$\checkmark^\checkmark\\checkmarkc\checkmarki\checkmarkr\checkmarkc\checkmark$\checkmarki\checkmark$\checkmark^\checkmark\\checkmarkc\checkmarki\checkmarkr\checkmarkc\checkmark$\checkmarkg\checkmark$\checkmark^\checkmark\\checkmarkc\checkmarki\checkmarkr\checkmarkc\checkmark$\checkmarkh\checkmark$\checkmark^\checkmark\\checkmarkc\checkmarki\checkmarkr\checkmarkc\checkmark$\checkmark \checkmark$\checkmark^\checkmark\\checkmarkc\checkmarki\checkmarkr\checkmarkc\checkmark$\checkmark:\checkmark$\checkmark^\checkmark\\checkmarkc\checkmarki\checkmarkr\checkmarkc\checkmark$\checkmark
\checkmark$\checkmark^\checkmark\\checkmarkc\checkmarki\checkmarkr\checkmarkc\checkmark$\checkmark\\checkmark$\checkmark^\checkmark\\checkmarkc\checkmarki\checkmarkr\checkmarkc\checkmark$\checkmarkb\checkmark$\checkmark^\checkmark\\checkmarkc\checkmarki\checkmarkr\checkmarkc\checkmark$\checkmarke\checkmark$\checkmark^\checkmark\\checkmarkc\checkmarki\checkmarkr\checkmarkc\checkmark$\checkmarkg\checkmark$\checkmark^\checkmark\\checkmarkc\checkmarki\checkmarkr\checkmarkc\checkmark$\checkmarki\checkmark$\checkmark^\checkmark\\checkmarkc\checkmarki\checkmarkr\checkmarkc\checkmark$\checkmarkn\checkmark$\checkmark^\checkmark\\checkmarkc\checkmarki\checkmarkr\checkmarkc\checkmark$\checkmark{\checkmark$\checkmark^\checkmark\\checkmarkc\checkmarki\checkmarkr\checkmarkc\checkmark$\checkmarke\checkmark$\checkmark^\checkmark\\checkmarkc\checkmarki\checkmarkr\checkmarkc\checkmark$\checkmarkq\checkmark$\checkmark^\checkmark\\checkmarkc\checkmarki\checkmarkr\checkmarkc\checkmark$\checkmarku\checkmark$\checkmark^\checkmark\\checkmarkc\checkmarki\checkmarkr\checkmarkc\checkmark$\checkmarka\checkmark$\checkmark^\checkmark\\checkmarkc\checkmarki\checkmarkr\checkmarkc\checkmark$\checkmarkt\checkmark$\checkmark^\checkmark\\checkmarkc\checkmarki\checkmarkr\checkmarkc\checkmark$\checkmarki\checkmark$\checkmark^\checkmark\\checkmarkc\checkmarki\checkmarkr\checkmarkc\checkmark$\checkmarko\checkmark$\checkmark^\checkmark\\checkmarkc\checkmarki\checkmarkr\checkmarkc\checkmark$\checkmarkn\checkmark$\checkmark^\checkmark\\checkmarkc\checkmarki\checkmarkr\checkmarkc\checkmark$\checkmark}\checkmark$\checkmark^\checkmark\\checkmarkc\checkmarki\checkmarkr\checkmarkc\checkmark$\checkmark
\checkmark$\checkmark^\checkmark\\checkmarkc\checkmarki\checkmarkr\checkmarkc\checkmark$\checkmark \checkmark$\checkmark^\checkmark\\checkmarkc\checkmarki\checkmarkr\checkmarkc\checkmark$\checkmark \checkmark$\checkmark^\checkmark\\checkmarkc\checkmarki\checkmarkr\checkmarkc\checkmark$\checkmark \checkmark$\checkmark^\checkmark\\checkmarkc\checkmarki\checkmarkr\checkmarkc\checkmark$\checkmark \checkmark$\checkmark^\checkmark\\checkmarkc\checkmarki\checkmarkr\checkmarkc\checkmark$\checkmarkf\checkmark$\checkmark^\checkmark\\checkmarkc\checkmarki\checkmarkr\checkmarkc\checkmark$\checkmark_\checkmark$\checkmark^\checkmark\\checkmarkc\checkmarki\checkmarkr\checkmarkc\checkmark$\checkmark1\checkmark$\checkmark^\checkmark\\checkmarkc\checkmarki\checkmarkr\checkmarkc\checkmark$\checkmark \checkmark$\checkmark^\checkmark\\checkmarkc\checkmarki\checkmarkr\checkmarkc\checkmark$\checkmark=\checkmark$\checkmark^\checkmark\\checkmarkc\checkmarki\checkmarkr\checkmarkc\checkmark$\checkmark \checkmark$\checkmark^\checkmark\\checkmarkc\checkmarki\checkmarkr\checkmarkc\checkmark$\checkmark\\checkmark$\checkmark^\checkmark\\checkmarkc\checkmarki\checkmarkr\checkmarkc\checkmark$\checkmarkf\checkmark$\checkmark^\checkmark\\checkmarkc\checkmarki\checkmarkr\checkmarkc\checkmark$\checkmarkr\checkmark$\checkmark^\checkmark\\checkmarkc\checkmarki\checkmarkr\checkmarkc\checkmark$\checkmarka\checkmark$\checkmark^\checkmark\\checkmarkc\checkmarki\checkmarkr\checkmarkc\checkmark$\checkmarkc\checkmark$\checkmark^\checkmark\\checkmarkc\checkmarki\checkmarkr\checkmarkc\checkmark$\checkmark{\checkmark$\checkmark^\checkmark\\checkmarkc\checkmarki\checkmarkr\checkmarkc\checkmark$\checkmark1\checkmark$\checkmark^\checkmark\\checkmarkc\checkmarki\checkmarkr\checkmarkc\checkmark$\checkmark}\checkmark$\checkmark^\checkmark\\checkmarkc\checkmarki\checkmarkr\checkmarkc\checkmark$\checkmark{\checkmark$\checkmark^\checkmark\\checkmarkc\checkmarki\checkmarkr\checkmarkc\checkmark$\checkmark2\checkmark$\checkmark^\checkmark\\checkmarkc\checkmarki\checkmarkr\checkmarkc\checkmark$\checkmark\\checkmark$\checkmark^\checkmark\\checkmarkc\checkmarki\checkmarkr\checkmarkc\checkmark$\checkmarkp\checkmark$\checkmark^\checkmark\\checkmarkc\checkmarki\checkmarkr\checkmarkc\checkmark$\checkmarki\checkmark$\checkmark^\checkmark\\checkmarkc\checkmarki\checkmarkr\checkmarkc\checkmark$\checkmark}\checkmark$\checkmark^\checkmark\\checkmarkc\checkmarki\checkmarkr\checkmarkc\checkmark$\checkmark \checkmark$\checkmark^\checkmark\\checkmarkc\checkmarki\checkmarkr\checkmarkc\checkmark$\checkmark\\checkmark$\checkmark^\checkmark\\checkmarkc\checkmarki\checkmarkr\checkmarkc\checkmark$\checkmarks\checkmark$\checkmark^\checkmark\\checkmarkc\checkmarki\checkmarkr\checkmarkc\checkmark$\checkmarkq\checkmark$\checkmark^\checkmark\\checkmarkc\checkmarki\checkmarkr\checkmarkc\checkmark$\checkmarkr\checkmark$\checkmark^\checkmark\\checkmarkc\checkmarki\checkmarkr\checkmarkc\checkmark$\checkmarkt\checkmark$\checkmark^\checkmark\\checkmarkc\checkmarki\checkmarkr\checkmarkc\checkmark$\checkmark{\checkmark$\checkmark^\checkmark\\checkmarkc\checkmarki\checkmarkr\checkmarkc\checkmark$\checkmark\\checkmark$\checkmark^\checkmark\\checkmarkc\checkmarki\checkmarkr\checkmarkc\checkmark$\checkmarkf\checkmark$\checkmark^\checkmark\\checkmarkc\checkmarki\checkmarkr\checkmarkc\checkmark$\checkmarkr\checkmark$\checkmark^\checkmark\\checkmarkc\checkmarki\checkmarkr\checkmarkc\checkmark$\checkmarka\checkmark$\checkmark^\checkmark\\checkmarkc\checkmarki\checkmarkr\checkmarkc\checkmark$\checkmarkc\checkmark$\checkmark^\checkmark\\checkmarkc\checkmarki\checkmarkr\checkmarkc\checkmark$\checkmark{\checkmark$\checkmark^\checkmark\\checkmarkc\checkmarki\checkmarkr\checkmarkc\checkmark$\checkmarkk\checkmark$\checkmark^\checkmark\\checkmarkc\checkmarki\checkmarkr\checkmarkc\checkmark$\checkmark_\checkmark$\checkmark^\checkmark\\checkmarkc\checkmarki\checkmarkr\checkmarkc\checkmark$\checkmark{\checkmark$\checkmark^\checkmark\\checkmarkc\checkmarki\checkmarkr\checkmarkc\checkmark$\checkmarkp\checkmark$\checkmark^\checkmark\\checkmarkc\checkmarki\checkmarkr\checkmarkc\checkmark$\checkmarko\checkmark$\checkmark^\checkmark\\checkmarkc\checkmarki\checkmarkr\checkmarkc\checkmark$\checkmarkr\checkmark$\checkmark^\checkmark\\checkmarkc\checkmarki\checkmarkr\checkmarkc\checkmark$\checkmarkt\checkmark$\checkmark^\checkmark\\checkmarkc\checkmarki\checkmarkr\checkmarkc\checkmark$\checkmarki\checkmark$\checkmark^\checkmark\\checkmarkc\checkmarki\checkmarkr\checkmarkc\checkmark$\checkmarkq\checkmark$\checkmark^\checkmark\\checkmarkc\checkmarki\checkmarkr\checkmarkc\checkmark$\checkmarku\checkmark$\checkmark^\checkmark\\checkmarkc\checkmarki\checkmarkr\checkmarkc\checkmark$\checkmarke\checkmark$\checkmark^\checkmark\\checkmarkc\checkmarki\checkmarkr\checkmarkc\checkmark$\checkmark}\checkmark$\checkmark^\checkmark\\checkmarkc\checkmarki\checkmarkr\checkmarkc\checkmark$\checkmark}\checkmark$\checkmark^\checkmark\\checkmarkc\checkmarki\checkmarkr\checkmarkc\checkmark$\checkmark{\checkmark$\checkmark^\checkmark\\checkmarkc\checkmarki\checkmarkr\checkmarkc\checkmark$\checkmarkM\checkmark$\checkmark^\checkmark\\checkmarkc\checkmarki\checkmarkr\checkmarkc\checkmark$\checkmark_\checkmark$\checkmark^\checkmark\\checkmarkc\checkmarki\checkmarkr\checkmarkc\checkmark$\checkmark{\checkmark$\checkmark^\checkmark\\checkmarkc\checkmarki\checkmarkr\checkmarkc\checkmark$\checkmarkb\checkmark$\checkmark^\checkmark\\checkmarkc\checkmarki\checkmarkr\checkmarkc\checkmark$\checkmarkr\checkmark$\checkmark^\checkmark\\checkmarkc\checkmarki\checkmarkr\checkmarkc\checkmark$\checkmarko\checkmark$\checkmark^\checkmark\\checkmarkc\checkmarki\checkmarkr\checkmarkc\checkmark$\checkmarkc\checkmark$\checkmark^\checkmark\\checkmarkc\checkmarki\checkmarkr\checkmarkc\checkmark$\checkmarkh\checkmark$\checkmark^\checkmark\\checkmarkc\checkmarki\checkmarkr\checkmarkc\checkmark$\checkmarke\checkmark$\checkmark^\checkmark\\checkmarkc\checkmarki\checkmarkr\checkmarkc\checkmark$\checkmark}\checkmark$\checkmark^\checkmark\\checkmarkc\checkmarki\checkmarkr\checkmarkc\checkmark$\checkmark}\checkmark$\checkmark^\checkmark\\checkmarkc\checkmarki\checkmarkr\checkmarkc\checkmark$\checkmark}\checkmark$\checkmark^\checkmark\\checkmarkc\checkmarki\checkmarkr\checkmarkc\checkmark$\checkmark
\checkmark$\checkmark^\checkmark\\checkmarkc\checkmarki\checkmarkr\checkmarkc\checkmark$\checkmark\\checkmark$\checkmark^\checkmark\\checkmarkc\checkmarki\checkmarkr\checkmarkc\checkmark$\checkmarke\checkmark$\checkmark^\checkmark\\checkmarkc\checkmarki\checkmarkr\checkmarkc\checkmark$\checkmarkn\checkmark$\checkmark^\checkmark\\checkmarkc\checkmarki\checkmarkr\checkmarkc\checkmark$\checkmarkd\checkmark$\checkmark^\checkmark\\checkmarkc\checkmarki\checkmarkr\checkmarkc\checkmark$\checkmark{\checkmark$\checkmark^\checkmark\\checkmarkc\checkmarki\checkmarkr\checkmarkc\checkmark$\checkmarke\checkmark$\checkmark^\checkmark\\checkmarkc\checkmarki\checkmarkr\checkmarkc\checkmark$\checkmarkq\checkmark$\checkmark^\checkmark\\checkmarkc\checkmarki\checkmarkr\checkmarkc\checkmark$\checkmarku\checkmark$\checkmark^\checkmark\\checkmarkc\checkmarki\checkmarkr\checkmarkc\checkmark$\checkmarka\checkmark$\checkmark^\checkmark\\checkmarkc\checkmarki\checkmarkr\checkmarkc\checkmark$\checkmarkt\checkmark$\checkmark^\checkmark\\checkmarkc\checkmarki\checkmarkr\checkmarkc\checkmark$\checkmarki\checkmark$\checkmark^\checkmark\\checkmarkc\checkmarki\checkmarkr\checkmarkc\checkmark$\checkmarko\checkmark$\checkmark^\checkmark\\checkmarkc\checkmarki\checkmarkr\checkmarkc\checkmark$\checkmarkn\checkmark$\checkmark^\checkmark\\checkmarkc\checkmarki\checkmarkr\checkmarkc\checkmark$\checkmark}\checkmark$\checkmark^\checkmark\\checkmarkc\checkmarki\checkmarkr\checkmarkc\checkmark$\checkmark
\checkmark$\checkmark^\checkmark\\checkmarkc\checkmarki\checkmarkr\checkmarkc\checkmark$\checkmark
\checkmark$\checkmark^\checkmark\\checkmarkc\checkmarki\checkmarkr\checkmarkc\checkmark$\checkmarkA\checkmark$\checkmark^\checkmark\\checkmarkc\checkmarki\checkmarkr\checkmarkc\checkmark$\checkmarkv\checkmark$\checkmark^\checkmark\\checkmarkc\checkmarki\checkmarkr\checkmarkc\checkmark$\checkmarke\checkmark$\checkmark^\checkmark\\checkmarkc\checkmarki\checkmarkr\checkmarkc\checkmark$\checkmarkc\checkmark$\checkmark^\checkmark\\checkmarkc\checkmarki\checkmarkr\checkmarkc\checkmark$\checkmark \checkmark$\checkmark^\checkmark\\checkmarkc\checkmarki\checkmarkr\checkmarkc\checkmark$\checkmarkl\checkmark$\checkmark^\checkmark\\checkmarkc\checkmarki\checkmarkr\checkmarkc\checkmark$\checkmarka\checkmark$\checkmark^\checkmark\\checkmarkc\checkmarki\checkmarkr\checkmarkc\checkmark$\checkmark \checkmark$\checkmark^\checkmark\\checkmarkc\checkmarki\checkmarkr\checkmarkc\checkmark$\checkmarkr\checkmark$\checkmark^\checkmark\\checkmarkc\checkmarki\checkmarkr\checkmarkc\checkmark$\checkmarki\checkmark$\checkmark^\checkmark\\checkmarkc\checkmarki\checkmarkr\checkmarkc\checkmark$\checkmarkg\checkmark$\checkmark^\checkmark\\checkmarkc\checkmarki\checkmarkr\checkmarkc\checkmark$\checkmarki\checkmark$\checkmark^\checkmark\\checkmarkc\checkmarki\checkmarkr\checkmarkc\checkmark$\checkmarkd\checkmark$\checkmark^\checkmark\\checkmarkc\checkmarki\checkmarkr\checkmarkc\checkmark$\checkmarki\checkmark$\checkmark^\checkmark\\checkmarkc\checkmarki\checkmarkr\checkmarkc\checkmark$\checkmarkt\checkmark$\checkmark^\checkmark\\checkmarkc\checkmarki\checkmarkr\checkmarkc\checkmark$\checkmarké\checkmark$\checkmark^\checkmark\\checkmarkc\checkmarki\checkmarkr\checkmarkc\checkmark$\checkmark \checkmark$\checkmark^\checkmark\\checkmarkc\checkmarki\checkmarkr\checkmarkc\checkmark$\checkmarkd\checkmark$\checkmark^\checkmark\\checkmarkc\checkmarki\checkmarkr\checkmarkc\checkmark$\checkmarku\checkmark$\checkmark^\checkmark\\checkmarkc\checkmarki\checkmarkr\checkmarkc\checkmark$\checkmark \checkmark$\checkmark^\checkmark\\checkmarkc\checkmarki\checkmarkr\checkmarkc\checkmark$\checkmarkp\checkmark$\checkmark^\checkmark\\checkmarkc\checkmarki\checkmarkr\checkmarkc\checkmark$\checkmarko\checkmark$\checkmark^\checkmark\\checkmarkc\checkmarki\checkmarkr\checkmarkc\checkmark$\checkmarkr\checkmark$\checkmark^\checkmark\\checkmarkc\checkmarki\checkmarkr\checkmarkc\checkmark$\checkmarkt\checkmark$\checkmark^\checkmark\\checkmarkc\checkmarki\checkmarkr\checkmarkc\checkmark$\checkmarki\checkmark$\checkmark^\checkmark\\checkmarkc\checkmarki\checkmarkr\checkmarkc\checkmark$\checkmarkq\checkmark$\checkmark^\checkmark\\checkmarkc\checkmarki\checkmarkr\checkmarkc\checkmark$\checkmarku\checkmark$\checkmark^\checkmark\\checkmarkc\checkmarki\checkmarkr\checkmarkc\checkmark$\checkmarke\checkmark$\checkmark^\checkmark\\checkmarkc\checkmarki\checkmarkr\checkmarkc\checkmark$\checkmark \checkmark$\checkmark^\checkmark\\checkmarkc\checkmarki\checkmarkr\checkmarkc\checkmark$\checkmark$\checkmark$\checkmark^\checkmark\\checkmarkc\checkmarki\checkmarkr\checkmarkc\checkmark$\checkmarkk\checkmark$\checkmark^\checkmark\\checkmarkc\checkmarki\checkmarkr\checkmarkc\checkmark$\checkmark_\checkmark$\checkmark^\checkmark\\checkmarkc\checkmarki\checkmarkr\checkmarkc\checkmark$\checkmark{\checkmark$\checkmark^\checkmark\\checkmarkc\checkmarki\checkmarkr\checkmarkc\checkmark$\checkmarkp\checkmark$\checkmark^\checkmark\\checkmarkc\checkmarki\checkmarkr\checkmarkc\checkmark$\checkmarko\checkmark$\checkmark^\checkmark\\checkmarkc\checkmarki\checkmarkr\checkmarkc\checkmark$\checkmarkr\checkmark$\checkmark^\checkmark\\checkmarkc\checkmarki\checkmarkr\checkmarkc\checkmark$\checkmarkt\checkmark$\checkmark^\checkmark\\checkmarkc\checkmarki\checkmarkr\checkmarkc\checkmark$\checkmarki\checkmark$\checkmark^\checkmark\\checkmarkc\checkmarki\checkmarkr\checkmarkc\checkmark$\checkmarkq\checkmark$\checkmark^\checkmark\\checkmarkc\checkmarki\checkmarkr\checkmarkc\checkmark$\checkmarku\checkmark$\checkmark^\checkmark\\checkmarkc\checkmarki\checkmarkr\checkmarkc\checkmark$\checkmarke\checkmark$\checkmark^\checkmark\\checkmarkc\checkmarki\checkmarkr\checkmarkc\checkmark$\checkmark}\checkmark$\checkmark^\checkmark\\checkmarkc\checkmarki\checkmarkr\checkmarkc\checkmark$\checkmark \checkmark$\checkmark^\checkmark\\checkmarkc\checkmarki\checkmarkr\checkmarkc\checkmark$\checkmark=\checkmark$\checkmark^\checkmark\\checkmarkc\checkmarki\checkmarkr\checkmarkc\checkmark$\checkmark \checkmark$\checkmark^\checkmark\\checkmarkc\checkmarki\checkmarkr\checkmarkc\checkmark$\checkmark4\checkmark$\checkmark^\checkmark\\checkmarkc\checkmarki\checkmarkr\checkmarkc\checkmark$\checkmark8\checkmark$\checkmark^\checkmark\\checkmarkc\checkmarki\checkmarkr\checkmarkc\checkmark$\checkmarkE\checkmark$\checkmark^\checkmark\\checkmarkc\checkmarki\checkmarkr\checkmarkc\checkmark$\checkmarkI\checkmark$\checkmark^\checkmark\\checkmarkc\checkmarki\checkmarkr\checkmarkc\checkmark$\checkmark/\checkmark$\checkmark^\checkmark\\checkmarkc\checkmarki\checkmarkr\checkmarkc\checkmark$\checkmarkL\checkmark$\checkmark^\checkmark\\checkmarkc\checkmarki\checkmarkr\checkmarkc\checkmark$\checkmark^\checkmark$\checkmark^\checkmark\\checkmarkc\checkmarki\checkmarkr\checkmarkc\checkmark$\checkmark3\checkmark$\checkmark^\checkmark\\checkmarkc\checkmarki\checkmarkr\checkmarkc\checkmark$\checkmark \checkmark$\checkmark^\checkmark\\checkmarkc\checkmarki\checkmarkr\checkmarkc\checkmark$\checkmark=\checkmark$\checkmark^\checkmark\\checkmarkc\checkmarki\checkmarkr\checkmarkc\checkmark$\checkmark \checkmark$\checkmark^\checkmark\\checkmarkc\checkmarki\checkmarkr\checkmarkc\checkmark$\checkmark4\checkmark$\checkmark^\checkmark\\checkmarkc\checkmarki\checkmarkr\checkmarkc\checkmark$\checkmark8\checkmark$\checkmark^\checkmark\\checkmarkc\checkmarki\checkmarkr\checkmarkc\checkmark$\checkmark \checkmark$\checkmark^\checkmark\\checkmarkc\checkmarki\checkmarkr\checkmarkc\checkmark$\checkmark\\checkmark$\checkmark^\checkmark\\checkmarkc\checkmarki\checkmarkr\checkmarkc\checkmark$\checkmarkt\checkmark$\checkmark^\checkmark\\checkmarkc\checkmarki\checkmarkr\checkmarkc\checkmark$\checkmarki\checkmark$\checkmark^\checkmark\\checkmarkc\checkmarki\checkmarkr\checkmarkc\checkmark$\checkmarkm\checkmark$\checkmark^\checkmark\\checkmarkc\checkmarki\checkmarkr\checkmarkc\checkmark$\checkmarke\checkmark$\checkmark^\checkmark\\checkmarkc\checkmarki\checkmarkr\checkmarkc\checkmark$\checkmarks\checkmark$\checkmark^\checkmark\\checkmarkc\checkmarki\checkmarkr\checkmarkc\checkmark$\checkmark \checkmark$\checkmark^\checkmark\\checkmarkc\checkmarki\checkmarkr\checkmarkc\checkmark$\checkmark6\checkmark$\checkmark^\checkmark\\checkmarkc\checkmarki\checkmarkr\checkmarkc\checkmark$\checkmark8\checkmark$\checkmark^\checkmark\\checkmarkc\checkmarki\checkmarkr\checkmarkc\checkmark$\checkmark\\checkmark$\checkmark^\checkmark\\checkmarkc\checkmarki\checkmarkr\checkmarkc\checkmark$\checkmark,\checkmark$\checkmark^\checkmark\\checkmarkc\checkmarki\checkmarkr\checkmarkc\checkmark$\checkmark9\checkmark$\checkmark^\checkmark\\checkmarkc\checkmarki\checkmarkr\checkmarkc\checkmark$\checkmark0\checkmark$\checkmark^\checkmark\\checkmarkc\checkmarki\checkmarkr\checkmarkc\checkmark$\checkmark0\checkmark$\checkmark^\checkmark\\checkmarkc\checkmarki\checkmarkr\checkmarkc\checkmark$\checkmark \checkmark$\checkmark^\checkmark\\checkmarkc\checkmarki\checkmarkr\checkmarkc\checkmark$\checkmark\\checkmark$\checkmark^\checkmark\\checkmarkc\checkmarki\checkmarkr\checkmarkc\checkmark$\checkmarkt\checkmark$\checkmark^\checkmark\\checkmarkc\checkmarki\checkmarkr\checkmarkc\checkmark$\checkmarki\checkmark$\checkmark^\checkmark\\checkmarkc\checkmarki\checkmarkr\checkmarkc\checkmark$\checkmarkm\checkmark$\checkmark^\checkmark\\checkmarkc\checkmarki\checkmarkr\checkmarkc\checkmark$\checkmarke\checkmark$\checkmark^\checkmark\\checkmarkc\checkmarki\checkmarkr\checkmarkc\checkmark$\checkmarks\checkmark$\checkmark^\checkmark\\checkmarkc\checkmarki\checkmarkr\checkmarkc\checkmark$\checkmark \checkmark$\checkmark^\checkmark\\checkmarkc\checkmarki\checkmarkr\checkmarkc\checkmark$\checkmark6\checkmark$\checkmark^\checkmark\\checkmarkc\checkmarki\checkmarkr\checkmarkc\checkmark$\checkmark5\checkmark$\checkmark^\checkmark\\checkmarkc\checkmarki\checkmarkr\checkmarkc\checkmark$\checkmark\\checkmark$\checkmark^\checkmark\\checkmarkc\checkmarki\checkmarkr\checkmarkc\checkmark$\checkmark,\checkmark$\checkmark^\checkmark\\checkmarkc\checkmarki\checkmarkr\checkmarkc\checkmark$\checkmark0\checkmark$\checkmark^\checkmark\\checkmarkc\checkmarki\checkmarkr\checkmarkc\checkmark$\checkmark0\checkmark$\checkmark^\checkmark\\checkmarkc\checkmarki\checkmarkr\checkmarkc\checkmark$\checkmark0\checkmark$\checkmark^\checkmark\\checkmarkc\checkmarki\checkmarkr\checkmarkc\checkmark$\checkmark \checkmark$\checkmark^\checkmark\\checkmarkc\checkmarki\checkmarkr\checkmarkc\checkmark$\checkmark/\checkmark$\checkmark^\checkmark\\checkmarkc\checkmarki\checkmarkr\checkmarkc\checkmark$\checkmark \checkmark$\checkmark^\checkmark\\checkmarkc\checkmarki\checkmarkr\checkmarkc\checkmark$\checkmark4\checkmark$\checkmark^\checkmark\\checkmarkc\checkmarki\checkmarkr\checkmarkc\checkmark$\checkmark0\checkmark$\checkmark^\checkmark\\checkmarkc\checkmarki\checkmarkr\checkmarkc\checkmark$\checkmark0\checkmark$\checkmark^\checkmark\\checkmarkc\checkmarki\checkmarkr\checkmarkc\checkmark$\checkmark^\checkmark$\checkmark^\checkmark\\checkmarkc\checkmarki\checkmarkr\checkmarkc\checkmark$\checkmark3\checkmark$\checkmark^\checkmark\\checkmarkc\checkmarki\checkmarkr\checkmarkc\checkmark$\checkmark \checkmark$\checkmark^\checkmark\\checkmarkc\checkmarki\checkmarkr\checkmarkc\checkmark$\checkmark=\checkmark$\checkmark^\checkmark\\checkmarkc\checkmarki\checkmarkr\checkmarkc\checkmark$\checkmark \checkmark$\checkmark^\checkmark\\checkmarkc\checkmarki\checkmarkr\checkmarkc\checkmark$\checkmark3\checkmark$\checkmark^\checkmark\\checkmarkc\checkmarki\checkmarkr\checkmarkc\checkmark$\checkmark\\checkmark$\checkmark^\checkmark\\checkmarkc\checkmarki\checkmarkr\checkmarkc\checkmark$\checkmark,\checkmark$\checkmark^\checkmark\\checkmarkc\checkmarki\checkmarkr\checkmarkc\checkmark$\checkmark3\checkmark$\checkmark^\checkmark\\checkmarkc\checkmarki\checkmarkr\checkmarkc\checkmark$\checkmark3\checkmark$\checkmark^\checkmark\\checkmarkc\checkmarki\checkmarkr\checkmarkc\checkmark$\checkmark0\checkmark$\checkmark^\checkmark\\checkmarkc\checkmarki\checkmarkr\checkmarkc\checkmark$\checkmark$\checkmark$\checkmark^\checkmark\\checkmarkc\checkmarki\checkmarkr\checkmarkc\checkmark$\checkmark \checkmark$\checkmark^\checkmark\\checkmarkc\checkmarki\checkmarkr\checkmarkc\checkmark$\checkmarkN\checkmark$\checkmark^\checkmark\\checkmarkc\checkmarki\checkmarkr\checkmarkc\checkmark$\checkmark/\checkmark$\checkmark^\checkmark\\checkmarkc\checkmarki\checkmarkr\checkmarkc\checkmark$\checkmarkm\checkmark$\checkmark^\checkmark\\checkmarkc\checkmarki\checkmarkr\checkmarkc\checkmark$\checkmarkm\checkmark$\checkmark^\checkmark\\checkmarkc\checkmarki\checkmarkr\checkmarkc\checkmark$\checkmark \checkmark$\checkmark^\checkmark\\checkmarkc\checkmarki\checkmarkr\checkmarkc\checkmark$\checkmarke\checkmark$\checkmark^\checkmark\\checkmarkc\checkmarki\checkmarkr\checkmarkc\checkmark$\checkmarkt\checkmark$\checkmark^\checkmark\\checkmarkc\checkmarki\checkmarkr\checkmarkc\checkmark$\checkmark \checkmark$\checkmark^\checkmark\\checkmarkc\checkmarki\checkmarkr\checkmarkc\checkmark$\checkmark$\checkmark$\checkmark^\checkmark\\checkmarkc\checkmarki\checkmarkr\checkmarkc\checkmark$\checkmarkM\checkmark$\checkmark^\checkmark\\checkmarkc\checkmarki\checkmarkr\checkmarkc\checkmark$\checkmark_\checkmark$\checkmark^\checkmark\\checkmarkc\checkmarki\checkmarkr\checkmarkc\checkmark$\checkmark{\checkmark$\checkmark^\checkmark\\checkmarkc\checkmarki\checkmarkr\checkmarkc\checkmark$\checkmarkb\checkmark$\checkmark^\checkmark\\checkmarkc\checkmarki\checkmarkr\checkmarkc\checkmark$\checkmarkr\checkmark$\checkmark^\checkmark\\checkmarkc\checkmarki\checkmarkr\checkmarkc\checkmark$\checkmarko\checkmark$\checkmark^\checkmark\\checkmarkc\checkmarki\checkmarkr\checkmarkc\checkmark$\checkmarkc\checkmark$\checkmark^\checkmark\\checkmarkc\checkmarki\checkmarkr\checkmarkc\checkmark$\checkmarkh\checkmark$\checkmark^\checkmark\\checkmarkc\checkmarki\checkmarkr\checkmarkc\checkmark$\checkmarke\checkmark$\checkmark^\checkmark\\checkmarkc\checkmarki\checkmarkr\checkmarkc\checkmark$\checkmark}\checkmark$\checkmark^\checkmark\\checkmarkc\checkmarki\checkmarkr\checkmarkc\checkmark$\checkmark \checkmark$\checkmark^\checkmark\\checkmarkc\checkmarki\checkmarkr\checkmarkc\checkmark$\checkmark=\checkmark$\checkmark^\checkmark\\checkmarkc\checkmarki\checkmarkr\checkmarkc\checkmark$\checkmark \checkmark$\checkmark^\checkmark\\checkmarkc\checkmarki\checkmarkr\checkmarkc\checkmark$\checkmark1\checkmark$\checkmark^\checkmark\\checkmarkc\checkmarki\checkmarkr\checkmarkc\checkmark$\checkmark$\checkmark$\checkmark^\checkmark\\checkmarkc\checkmarki\checkmarkr\checkmarkc\checkmark$\checkmark \checkmark$\checkmark^\checkmark\\checkmarkc\checkmarki\checkmarkr\checkmarkc\checkmark$\checkmarkk\checkmark$\checkmark^\checkmark\\checkmarkc\checkmarki\checkmarkr\checkmarkc\checkmark$\checkmarkg\checkmark$\checkmark^\checkmark\\checkmarkc\checkmarki\checkmarkr\checkmarkc\checkmark$\checkmark \checkmark$\checkmark^\checkmark\\checkmarkc\checkmarki\checkmarkr\checkmarkc\checkmark$\checkmark:\checkmark$\checkmark^\checkmark\\checkmarkc\checkmarki\checkmarkr\checkmarkc\checkmark$\checkmark
\checkmark$\checkmark^\checkmark\\checkmarkc\checkmarki\checkmarkr\checkmarkc\checkmark$\checkmark\\checkmark$\checkmark^\checkmark\\checkmarkc\checkmarki\checkmarkr\checkmarkc\checkmark$\checkmarkb\checkmark$\checkmark^\checkmark\\checkmarkc\checkmarki\checkmarkr\checkmarkc\checkmark$\checkmarke\checkmark$\checkmark^\checkmark\\checkmarkc\checkmarki\checkmarkr\checkmarkc\checkmark$\checkmarkg\checkmark$\checkmark^\checkmark\\checkmarkc\checkmarki\checkmarkr\checkmarkc\checkmark$\checkmarki\checkmark$\checkmark^\checkmark\\checkmarkc\checkmarki\checkmarkr\checkmarkc\checkmark$\checkmarkn\checkmark$\checkmark^\checkmark\\checkmarkc\checkmarki\checkmarkr\checkmarkc\checkmark$\checkmark{\checkmark$\checkmark^\checkmark\\checkmarkc\checkmarki\checkmarkr\checkmarkc\checkmark$\checkmarke\checkmark$\checkmark^\checkmark\\checkmarkc\checkmarki\checkmarkr\checkmarkc\checkmark$\checkmarkq\checkmark$\checkmark^\checkmark\\checkmarkc\checkmarki\checkmarkr\checkmarkc\checkmark$\checkmarku\checkmark$\checkmark^\checkmark\\checkmarkc\checkmarki\checkmarkr\checkmarkc\checkmark$\checkmarka\checkmark$\checkmark^\checkmark\\checkmarkc\checkmarki\checkmarkr\checkmarkc\checkmark$\checkmarkt\checkmark$\checkmark^\checkmark\\checkmarkc\checkmarki\checkmarkr\checkmarkc\checkmark$\checkmarki\checkmark$\checkmark^\checkmark\\checkmarkc\checkmarki\checkmarkr\checkmarkc\checkmark$\checkmarko\checkmark$\checkmark^\checkmark\\checkmarkc\checkmarki\checkmarkr\checkmarkc\checkmark$\checkmarkn\checkmark$\checkmark^\checkmark\\checkmarkc\checkmarki\checkmarkr\checkmarkc\checkmark$\checkmark}\checkmark$\checkmark^\checkmark\\checkmarkc\checkmarki\checkmarkr\checkmarkc\checkmark$\checkmark
\checkmark$\checkmark^\checkmark\\checkmarkc\checkmarki\checkmarkr\checkmarkc\checkmark$\checkmark \checkmark$\checkmark^\checkmark\\checkmarkc\checkmarki\checkmarkr\checkmarkc\checkmark$\checkmark \checkmark$\checkmark^\checkmark\\checkmarkc\checkmarki\checkmarkr\checkmarkc\checkmark$\checkmark \checkmark$\checkmark^\checkmark\\checkmarkc\checkmarki\checkmarkr\checkmarkc\checkmark$\checkmark \checkmark$\checkmark^\checkmark\\checkmarkc\checkmarki\checkmarkr\checkmarkc\checkmark$\checkmarkf\checkmark$\checkmark^\checkmark\\checkmarkc\checkmarki\checkmarkr\checkmarkc\checkmark$\checkmark_\checkmark$\checkmark^\checkmark\\checkmarkc\checkmarki\checkmarkr\checkmarkc\checkmark$\checkmark1\checkmark$\checkmark^\checkmark\\checkmarkc\checkmarki\checkmarkr\checkmarkc\checkmark$\checkmark \checkmark$\checkmark^\checkmark\\checkmarkc\checkmarki\checkmarkr\checkmarkc\checkmark$\checkmark=\checkmark$\checkmark^\checkmark\\checkmarkc\checkmarki\checkmarkr\checkmarkc\checkmark$\checkmark \checkmark$\checkmark^\checkmark\\checkmarkc\checkmarki\checkmarkr\checkmarkc\checkmark$\checkmark\\checkmark$\checkmark^\checkmark\\checkmarkc\checkmarki\checkmarkr\checkmarkc\checkmark$\checkmarkf\checkmark$\checkmark^\checkmark\\checkmarkc\checkmarki\checkmarkr\checkmarkc\checkmark$\checkmarkr\checkmark$\checkmark^\checkmark\\checkmarkc\checkmarki\checkmarkr\checkmarkc\checkmark$\checkmarka\checkmark$\checkmark^\checkmark\\checkmarkc\checkmarki\checkmarkr\checkmarkc\checkmark$\checkmarkc\checkmark$\checkmark^\checkmark\\checkmarkc\checkmarki\checkmarkr\checkmarkc\checkmark$\checkmark{\checkmark$\checkmark^\checkmark\\checkmarkc\checkmarki\checkmarkr\checkmarkc\checkmark$\checkmark1\checkmark$\checkmark^\checkmark\\checkmarkc\checkmarki\checkmarkr\checkmarkc\checkmark$\checkmark}\checkmark$\checkmark^\checkmark\\checkmarkc\checkmarki\checkmarkr\checkmarkc\checkmark$\checkmark{\checkmark$\checkmark^\checkmark\\checkmarkc\checkmarki\checkmarkr\checkmarkc\checkmark$\checkmark2\checkmark$\checkmark^\checkmark\\checkmarkc\checkmarki\checkmarkr\checkmarkc\checkmark$\checkmark\\checkmark$\checkmark^\checkmark\\checkmarkc\checkmarki\checkmarkr\checkmarkc\checkmark$\checkmarkp\checkmark$\checkmark^\checkmark\\checkmarkc\checkmarki\checkmarkr\checkmarkc\checkmark$\checkmarki\checkmark$\checkmark^\checkmark\\checkmarkc\checkmarki\checkmarkr\checkmarkc\checkmark$\checkmark}\checkmark$\checkmark^\checkmark\\checkmarkc\checkmarki\checkmarkr\checkmarkc\checkmark$\checkmark \checkmark$\checkmark^\checkmark\\checkmarkc\checkmarki\checkmarkr\checkmarkc\checkmark$\checkmark\\checkmark$\checkmark^\checkmark\\checkmarkc\checkmarki\checkmarkr\checkmarkc\checkmark$\checkmarks\checkmark$\checkmark^\checkmark\\checkmarkc\checkmarki\checkmarkr\checkmarkc\checkmark$\checkmarkq\checkmark$\checkmark^\checkmark\\checkmarkc\checkmarki\checkmarkr\checkmarkc\checkmark$\checkmarkr\checkmark$\checkmark^\checkmark\\checkmarkc\checkmarki\checkmarkr\checkmarkc\checkmark$\checkmarkt\checkmark$\checkmark^\checkmark\\checkmarkc\checkmarki\checkmarkr\checkmarkc\checkmark$\checkmark{\checkmark$\checkmark^\checkmark\\checkmarkc\checkmarki\checkmarkr\checkmarkc\checkmark$\checkmark\\checkmark$\checkmark^\checkmark\\checkmarkc\checkmarki\checkmarkr\checkmarkc\checkmark$\checkmarkf\checkmark$\checkmark^\checkmark\\checkmarkc\checkmarki\checkmarkr\checkmarkc\checkmark$\checkmarkr\checkmark$\checkmark^\checkmark\\checkmarkc\checkmarki\checkmarkr\checkmarkc\checkmark$\checkmarka\checkmark$\checkmark^\checkmark\\checkmarkc\checkmarki\checkmarkr\checkmarkc\checkmark$\checkmarkc\checkmark$\checkmark^\checkmark\\checkmarkc\checkmarki\checkmarkr\checkmarkc\checkmark$\checkmark{\checkmark$\checkmark^\checkmark\\checkmarkc\checkmarki\checkmarkr\checkmarkc\checkmark$\checkmark3\checkmark$\checkmark^\checkmark\\checkmarkc\checkmarki\checkmarkr\checkmarkc\checkmark$\checkmark\\checkmark$\checkmark^\checkmark\\checkmarkc\checkmarki\checkmarkr\checkmarkc\checkmark$\checkmark,\checkmark$\checkmark^\checkmark\\checkmarkc\checkmarki\checkmarkr\checkmarkc\checkmark$\checkmark3\checkmark$\checkmark^\checkmark\\checkmarkc\checkmarki\checkmarkr\checkmarkc\checkmark$\checkmark3\checkmark$\checkmark^\checkmark\\checkmarkc\checkmarki\checkmarkr\checkmarkc\checkmark$\checkmark0\checkmark$\checkmark^\checkmark\\checkmarkc\checkmarki\checkmarkr\checkmarkc\checkmark$\checkmark\\checkmark$\checkmark^\checkmark\\checkmarkc\checkmarki\checkmarkr\checkmarkc\checkmark$\checkmark,\checkmark$\checkmark^\checkmark\\checkmarkc\checkmarki\checkmarkr\checkmarkc\checkmark$\checkmark0\checkmark$\checkmark^\checkmark\\checkmarkc\checkmarki\checkmarkr\checkmarkc\checkmark$\checkmark0\checkmark$\checkmark^\checkmark\\checkmarkc\checkmarki\checkmarkr\checkmarkc\checkmark$\checkmark0\checkmark$\checkmark^\checkmark\\checkmarkc\checkmarki\checkmarkr\checkmarkc\checkmark$\checkmark}\checkmark$\checkmark^\checkmark\\checkmarkc\checkmarki\checkmarkr\checkmarkc\checkmark$\checkmark{\checkmark$\checkmark^\checkmark\\checkmarkc\checkmarki\checkmarkr\checkmarkc\checkmark$\checkmark1\checkmark$\checkmark^\checkmark\\checkmarkc\checkmarki\checkmarkr\checkmarkc\checkmark$\checkmark}\checkmark$\checkmark^\checkmark\\checkmarkc\checkmarki\checkmarkr\checkmarkc\checkmark$\checkmark}\checkmark$\checkmark^\checkmark\\checkmarkc\checkmarki\checkmarkr\checkmarkc\checkmark$\checkmark \checkmark$\checkmark^\checkmark\\checkmarkc\checkmarki\checkmarkr\checkmarkc\checkmark$\checkmark=\checkmark$\checkmark^\checkmark\\checkmarkc\checkmarki\checkmarkr\checkmarkc\checkmark$\checkmark \checkmark$\checkmark^\checkmark\\checkmarkc\checkmarki\checkmarkr\checkmarkc\checkmark$\checkmark\\checkmark$\checkmark^\checkmark\\checkmarkc\checkmarki\checkmarkr\checkmarkc\checkmark$\checkmarkf\checkmark$\checkmark^\checkmark\\checkmarkc\checkmarki\checkmarkr\checkmarkc\checkmark$\checkmarkr\checkmark$\checkmark^\checkmark\\checkmarkc\checkmarki\checkmarkr\checkmarkc\checkmark$\checkmarka\checkmark$\checkmark^\checkmark\\checkmarkc\checkmarki\checkmarkr\checkmarkc\checkmark$\checkmarkc\checkmark$\checkmark^\checkmark\\checkmarkc\checkmarki\checkmarkr\checkmarkc\checkmark$\checkmark{\checkmark$\checkmark^\checkmark\\checkmarkc\checkmarki\checkmarkr\checkmarkc\checkmark$\checkmark1\checkmark$\checkmark^\checkmark\\checkmarkc\checkmarki\checkmarkr\checkmarkc\checkmark$\checkmark}\checkmark$\checkmark^\checkmark\\checkmarkc\checkmarki\checkmarkr\checkmarkc\checkmark$\checkmark{\checkmark$\checkmark^\checkmark\\checkmarkc\checkmarki\checkmarkr\checkmarkc\checkmark$\checkmark2\checkmark$\checkmark^\checkmark\\checkmarkc\checkmarki\checkmarkr\checkmarkc\checkmark$\checkmark\\checkmark$\checkmark^\checkmark\\checkmarkc\checkmarki\checkmarkr\checkmarkc\checkmark$\checkmarkp\checkmark$\checkmark^\checkmark\\checkmarkc\checkmarki\checkmarkr\checkmarkc\checkmark$\checkmarki\checkmark$\checkmark^\checkmark\\checkmarkc\checkmarki\checkmarkr\checkmarkc\checkmark$\checkmark}\checkmark$\checkmark^\checkmark\\checkmarkc\checkmarki\checkmarkr\checkmarkc\checkmark$\checkmark \checkmark$\checkmark^\checkmark\\checkmarkc\checkmarki\checkmarkr\checkmarkc\checkmark$\checkmark\\checkmark$\checkmark^\checkmark\\checkmarkc\checkmarki\checkmarkr\checkmarkc\checkmark$\checkmarkt\checkmark$\checkmark^\checkmark\\checkmarkc\checkmarki\checkmarkr\checkmarkc\checkmark$\checkmarki\checkmark$\checkmark^\checkmark\\checkmarkc\checkmarki\checkmarkr\checkmarkc\checkmark$\checkmarkm\checkmark$\checkmark^\checkmark\\checkmarkc\checkmarki\checkmarkr\checkmarkc\checkmark$\checkmarke\checkmark$\checkmark^\checkmark\\checkmarkc\checkmarki\checkmarkr\checkmarkc\checkmark$\checkmarks\checkmark$\checkmark^\checkmark\\checkmarkc\checkmarki\checkmarkr\checkmarkc\checkmark$\checkmark \checkmark$\checkmark^\checkmark\\checkmarkc\checkmarki\checkmarkr\checkmarkc\checkmark$\checkmark1\checkmark$\checkmark^\checkmark\\checkmarkc\checkmarki\checkmarkr\checkmarkc\checkmark$\checkmark\\checkmark$\checkmark^\checkmark\\checkmarkc\checkmarki\checkmarkr\checkmarkc\checkmark$\checkmark,\checkmark$\checkmark^\checkmark\\checkmarkc\checkmarki\checkmarkr\checkmarkc\checkmark$\checkmark8\checkmark$\checkmark^\checkmark\\checkmarkc\checkmarki\checkmarkr\checkmarkc\checkmark$\checkmark2\checkmark$\checkmark^\checkmark\\checkmarkc\checkmarki\checkmarkr\checkmarkc\checkmark$\checkmark5\checkmark$\checkmark^\checkmark\\checkmarkc\checkmarki\checkmarkr\checkmarkc\checkmark$\checkmark \checkmark$\checkmark^\checkmark\\checkmarkc\checkmarki\checkmarkr\checkmarkc\checkmark$\checkmark\\checkmark$\checkmark^\checkmark\\checkmarkc\checkmarki\checkmarkr\checkmarkc\checkmark$\checkmarka\checkmark$\checkmark^\checkmark\\checkmarkc\checkmarki\checkmarkr\checkmarkc\checkmark$\checkmarkp\checkmark$\checkmark^\checkmark\\checkmarkc\checkmarki\checkmarkr\checkmarkc\checkmark$\checkmarkp\checkmark$\checkmark^\checkmark\\checkmarkc\checkmarki\checkmarkr\checkmarkc\checkmark$\checkmarkr\checkmark$\checkmark^\checkmark\\checkmarkc\checkmarki\checkmarkr\checkmarkc\checkmark$\checkmarko\checkmark$\checkmark^\checkmark\\checkmarkc\checkmarki\checkmarkr\checkmarkc\checkmark$\checkmarkx\checkmark$\checkmark^\checkmark\\checkmarkc\checkmarki\checkmarkr\checkmarkc\checkmark$\checkmark \checkmark$\checkmark^\checkmark\\checkmarkc\checkmarki\checkmarkr\checkmarkc\checkmark$\checkmark\\checkmark$\checkmark^\checkmark\\checkmarkc\checkmarki\checkmarkr\checkmarkc\checkmark$\checkmarkb\checkmark$\checkmark^\checkmark\\checkmarkc\checkmarki\checkmarkr\checkmarkc\checkmark$\checkmarko\checkmark$\checkmark^\checkmark\\checkmarkc\checkmarki\checkmarkr\checkmarkc\checkmark$\checkmarkx\checkmark$\checkmark^\checkmark\\checkmarkc\checkmarki\checkmarkr\checkmarkc\checkmark$\checkmarke\checkmark$\checkmark^\checkmark\\checkmarkc\checkmarki\checkmarkr\checkmarkc\checkmark$\checkmarkd\checkmark$\checkmark^\checkmark\\checkmarkc\checkmarki\checkmarkr\checkmarkc\checkmark$\checkmark{\checkmark$\checkmark^\checkmark\\checkmarkc\checkmarki\checkmarkr\checkmarkc\checkmark$\checkmark2\checkmark$\checkmark^\checkmark\\checkmarkc\checkmarki\checkmarkr\checkmarkc\checkmark$\checkmark9\checkmark$\checkmark^\checkmark\\checkmarkc\checkmarki\checkmarkr\checkmarkc\checkmark$\checkmark0\checkmark$\checkmark^\checkmark\\checkmarkc\checkmarki\checkmarkr\checkmarkc\checkmark$\checkmark \checkmark$\checkmark^\checkmark\\checkmarkc\checkmarki\checkmarkr\checkmarkc\checkmark$\checkmark\\checkmark$\checkmark^\checkmark\\checkmarkc\checkmarki\checkmarkr\checkmarkc\checkmark$\checkmarkt\checkmark$\checkmark^\checkmark\\checkmarkc\checkmarki\checkmarkr\checkmarkc\checkmark$\checkmarke\checkmark$\checkmark^\checkmark\\checkmarkc\checkmarki\checkmarkr\checkmarkc\checkmark$\checkmarkx\checkmark$\checkmark^\checkmark\\checkmarkc\checkmarki\checkmarkr\checkmarkc\checkmark$\checkmarkt\checkmark$\checkmark^\checkmark\\checkmarkc\checkmarki\checkmarkr\checkmarkc\checkmark$\checkmark{\checkmark$\checkmark^\checkmark\\checkmarkc\checkmarki\checkmarkr\checkmarkc\checkmark$\checkmark \checkmark$\checkmark^\checkmark\\checkmarkc\checkmarki\checkmarkr\checkmarkc\checkmark$\checkmarkH\checkmark$\checkmark^\checkmark\\checkmarkc\checkmarki\checkmarkr\checkmarkc\checkmark$\checkmarkz\checkmark$\checkmark^\checkmark\\checkmarkc\checkmarki\checkmarkr\checkmarkc\checkmark$\checkmark}\checkmark$\checkmark^\checkmark\\checkmarkc\checkmarki\checkmarkr\checkmarkc\checkmark$\checkmark}\checkmark$\checkmark^\checkmark\\checkmarkc\checkmarki\checkmarkr\checkmarkc\checkmark$\checkmark
\checkmark$\checkmark^\checkmark\\checkmarkc\checkmarki\checkmarkr\checkmarkc\checkmark$\checkmark\\checkmark$\checkmark^\checkmark\\checkmarkc\checkmarki\checkmarkr\checkmarkc\checkmark$\checkmarke\checkmark$\checkmark^\checkmark\\checkmarkc\checkmarki\checkmarkr\checkmarkc\checkmark$\checkmarkn\checkmark$\checkmark^\checkmark\\checkmarkc\checkmarki\checkmarkr\checkmarkc\checkmark$\checkmarkd\checkmark$\checkmark^\checkmark\\checkmarkc\checkmarki\checkmarkr\checkmarkc\checkmark$\checkmark{\checkmark$\checkmark^\checkmark\\checkmarkc\checkmarki\checkmarkr\checkmarkc\checkmark$\checkmarke\checkmark$\checkmark^\checkmark\\checkmarkc\checkmarki\checkmarkr\checkmarkc\checkmark$\checkmarkq\checkmark$\checkmark^\checkmark\\checkmarkc\checkmarki\checkmarkr\checkmarkc\checkmark$\checkmarku\checkmark$\checkmark^\checkmark\\checkmarkc\checkmarki\checkmarkr\checkmarkc\checkmark$\checkmarka\checkmark$\checkmark^\checkmark\\checkmarkc\checkmarki\checkmarkr\checkmarkc\checkmark$\checkmarkt\checkmark$\checkmark^\checkmark\\checkmarkc\checkmarki\checkmarkr\checkmarkc\checkmark$\checkmarki\checkmark$\checkmark^\checkmark\\checkmarkc\checkmarki\checkmarkr\checkmarkc\checkmark$\checkmarko\checkmark$\checkmark^\checkmark\\checkmarkc\checkmarki\checkmarkr\checkmarkc\checkmark$\checkmarkn\checkmark$\checkmark^\checkmark\\checkmarkc\checkmarki\checkmarkr\checkmarkc\checkmark$\checkmark}\checkmark$\checkmark^\checkmark\\checkmarkc\checkmarki\checkmarkr\checkmarkc\checkmark$\checkmark
\checkmark$\checkmark^\checkmark\\checkmarkc\checkmarki\checkmarkr\checkmarkc\checkmark$\checkmark
\checkmark$\checkmark^\checkmark\\checkmarkc\checkmarki\checkmarkr\checkmarkc\checkmark$\checkmark\\checkmark$\checkmark^\checkmark\\checkmarkc\checkmarki\checkmarkr\checkmarkc\checkmark$\checkmarkt\checkmark$\checkmark^\checkmark\\checkmarkc\checkmarki\checkmarkr\checkmarkc\checkmark$\checkmarke\checkmark$\checkmark^\checkmark\\checkmarkc\checkmarki\checkmarkr\checkmarkc\checkmark$\checkmarkx\checkmark$\checkmark^\checkmark\\checkmarkc\checkmarki\checkmarkr\checkmarkc\checkmark$\checkmarkt\checkmark$\checkmark^\checkmark\\checkmarkc\checkmarki\checkmarkr\checkmarkc\checkmark$\checkmarkb\checkmark$\checkmark^\checkmark\\checkmarkc\checkmarki\checkmarkr\checkmarkc\checkmark$\checkmarkf\checkmark$\checkmark^\checkmark\\checkmarkc\checkmarki\checkmarkr\checkmarkc\checkmark$\checkmark{\checkmark$\checkmark^\checkmark\\checkmarkc\checkmarki\checkmarkr\checkmarkc\checkmark$\checkmarkV\checkmark$\checkmark^\checkmark\\checkmarkc\checkmarki\checkmarkr\checkmarkc\checkmark$\checkmarké\checkmark$\checkmark^\checkmark\\checkmarkc\checkmarki\checkmarkr\checkmarkc\checkmark$\checkmarkr\checkmark$\checkmark^\checkmark\\checkmarkc\checkmarki\checkmarkr\checkmarkc\checkmark$\checkmarki\checkmark$\checkmark^\checkmark\\checkmarkc\checkmarki\checkmarkr\checkmarkc\checkmark$\checkmarkf\checkmark$\checkmark^\checkmark\\checkmarkc\checkmarki\checkmarkr\checkmarkc\checkmark$\checkmarki\checkmark$\checkmark^\checkmark\\checkmarkc\checkmarki\checkmarkr\checkmarkc\checkmark$\checkmarkc\checkmark$\checkmark^\checkmark\\checkmarkc\checkmarki\checkmarkr\checkmarkc\checkmark$\checkmarka\checkmark$\checkmark^\checkmark\\checkmarkc\checkmarki\checkmarkr\checkmarkc\checkmark$\checkmarkt\checkmark$\checkmark^\checkmark\\checkmarkc\checkmarki\checkmarkr\checkmarkc\checkmark$\checkmarki\checkmark$\checkmark^\checkmark\\checkmarkc\checkmarki\checkmarkr\checkmarkc\checkmark$\checkmarko\checkmark$\checkmark^\checkmark\\checkmarkc\checkmarki\checkmarkr\checkmarkc\checkmark$\checkmarkn\checkmark$\checkmark^\checkmark\\checkmarkc\checkmarki\checkmarkr\checkmarkc\checkmark$\checkmark \checkmark$\checkmark^\checkmark\\checkmarkc\checkmarki\checkmarkr\checkmarkc\checkmark$\checkmarka\checkmark$\checkmark^\checkmark\\checkmarkc\checkmarki\checkmarkr\checkmarkc\checkmark$\checkmarkn\checkmark$\checkmark^\checkmark\\checkmarkc\checkmarki\checkmarkr\checkmarkc\checkmark$\checkmarkt\checkmark$\checkmark^\checkmark\\checkmarkc\checkmarki\checkmarkr\checkmarkc\checkmark$\checkmarki\checkmark$\checkmark^\checkmark\\checkmarkc\checkmarki\checkmarkr\checkmarkc\checkmark$\checkmark-\checkmark$\checkmark^\checkmark\\checkmarkc\checkmarki\checkmarkr\checkmarkc\checkmark$\checkmarkr\checkmark$\checkmark^\checkmark\\checkmarkc\checkmarki\checkmarkr\checkmarkc\checkmark$\checkmarké\checkmark$\checkmark^\checkmark\\checkmarkc\checkmarki\checkmarkr\checkmarkc\checkmark$\checkmarks\checkmark$\checkmark^\checkmark\\checkmarkc\checkmarki\checkmarkr\checkmarkc\checkmark$\checkmarko\checkmark$\checkmark^\checkmark\\checkmarkc\checkmarki\checkmarkr\checkmarkc\checkmark$\checkmarkn\checkmark$\checkmark^\checkmark\\checkmarkc\checkmarki\checkmarkr\checkmarkc\checkmark$\checkmarka\checkmark$\checkmark^\checkmark\\checkmarkc\checkmarki\checkmarkr\checkmarkc\checkmark$\checkmarkn\checkmark$\checkmark^\checkmark\\checkmarkc\checkmarki\checkmarkr\checkmarkc\checkmark$\checkmarkc\checkmark$\checkmark^\checkmark\\checkmarkc\checkmarki\checkmarkr\checkmarkc\checkmark$\checkmarke\checkmark$\checkmark^\checkmark\\checkmarkc\checkmarki\checkmarkr\checkmarkc\checkmark$\checkmark \checkmark$\checkmark^\checkmark\\checkmarkc\checkmarki\checkmarkr\checkmarkc\checkmark$\checkmark:\checkmark$\checkmark^\checkmark\\checkmarkc\checkmarki\checkmarkr\checkmarkc\checkmark$\checkmark}\checkmark$\checkmark^\checkmark\\checkmarkc\checkmarki\checkmarkr\checkmarkc\checkmark$\checkmark
\checkmark$\checkmark^\checkmark\\checkmarkc\checkmarki\checkmarkr\checkmarkc\checkmark$\checkmark\\checkmark$\checkmark^\checkmark\\checkmarkc\checkmarki\checkmarkr\checkmarkc\checkmark$\checkmarkb\checkmark$\checkmark^\checkmark\\checkmarkc\checkmarki\checkmarkr\checkmarkc\checkmark$\checkmarke\checkmark$\checkmark^\checkmark\\checkmarkc\checkmarki\checkmarkr\checkmarkc\checkmark$\checkmarkg\checkmark$\checkmark^\checkmark\\checkmarkc\checkmarki\checkmarkr\checkmarkc\checkmark$\checkmarki\checkmark$\checkmark^\checkmark\\checkmarkc\checkmarki\checkmarkr\checkmarkc\checkmark$\checkmarkn\checkmark$\checkmark^\checkmark\\checkmarkc\checkmarki\checkmarkr\checkmarkc\checkmark$\checkmark{\checkmark$\checkmark^\checkmark\\checkmarkc\checkmarki\checkmarkr\checkmarkc\checkmark$\checkmarke\checkmark$\checkmark^\checkmark\\checkmarkc\checkmarki\checkmarkr\checkmarkc\checkmark$\checkmarkq\checkmark$\checkmark^\checkmark\\checkmarkc\checkmarki\checkmarkr\checkmarkc\checkmark$\checkmarku\checkmark$\checkmark^\checkmark\\checkmarkc\checkmarki\checkmarkr\checkmarkc\checkmark$\checkmarka\checkmark$\checkmark^\checkmark\\checkmarkc\checkmarki\checkmarkr\checkmarkc\checkmark$\checkmarkt\checkmark$\checkmark^\checkmark\\checkmarkc\checkmarki\checkmarkr\checkmarkc\checkmark$\checkmarki\checkmark$\checkmark^\checkmark\\checkmarkc\checkmarki\checkmarkr\checkmarkc\checkmark$\checkmarko\checkmark$\checkmark^\checkmark\\checkmarkc\checkmarki\checkmarkr\checkmarkc\checkmark$\checkmarkn\checkmark$\checkmark^\checkmark\\checkmarkc\checkmarki\checkmarkr\checkmarkc\checkmark$\checkmark}\checkmark$\checkmark^\checkmark\\checkmarkc\checkmarki\checkmarkr\checkmarkc\checkmark$\checkmark
\checkmark$\checkmark^\checkmark\\checkmarkc\checkmarki\checkmarkr\checkmarkc\checkmark$\checkmark \checkmark$\checkmark^\checkmark\\checkmarkc\checkmarki\checkmarkr\checkmarkc\checkmark$\checkmark \checkmark$\checkmark^\checkmark\\checkmarkc\checkmarki\checkmarkr\checkmarkc\checkmark$\checkmark \checkmark$\checkmark^\checkmark\\checkmarkc\checkmarki\checkmarkr\checkmarkc\checkmark$\checkmark \checkmark$\checkmark^\checkmark\\checkmarkc\checkmarki\checkmarkr\checkmarkc\checkmark$\checkmark\\checkmark$\checkmark^\checkmark\\checkmarkc\checkmarki\checkmarkr\checkmarkc\checkmark$\checkmarkt\checkmark$\checkmark^\checkmark\\checkmarkc\checkmarki\checkmarkr\checkmarkc\checkmark$\checkmarke\checkmark$\checkmark^\checkmark\\checkmarkc\checkmarki\checkmarkr\checkmarkc\checkmark$\checkmarkx\checkmark$\checkmark^\checkmark\\checkmarkc\checkmarki\checkmarkr\checkmarkc\checkmark$\checkmarkt\checkmark$\checkmark^\checkmark\\checkmarkc\checkmarki\checkmarkr\checkmarkc\checkmark$\checkmark{\checkmark$\checkmark^\checkmark\\checkmarkc\checkmarki\checkmarkr\checkmarkc\checkmark$\checkmarkM\checkmark$\checkmark^\checkmark\\checkmarkc\checkmarki\checkmarkr\checkmarkc\checkmark$\checkmarka\checkmark$\checkmark^\checkmark\\checkmarkc\checkmarki\checkmarkr\checkmarkc\checkmark$\checkmarkr\checkmark$\checkmark^\checkmark\\checkmarkc\checkmarki\checkmarkr\checkmarkc\checkmark$\checkmarkg\checkmark$\checkmark^\checkmark\\checkmarkc\checkmarki\checkmarkr\checkmarkc\checkmark$\checkmarke\checkmark$\checkmark^\checkmark\\checkmarkc\checkmarki\checkmarkr\checkmarkc\checkmark$\checkmark \checkmark$\checkmark^\checkmark\\checkmarkc\checkmarki\checkmarkr\checkmarkc\checkmark$\checkmarkd\checkmark$\checkmark^\checkmark\\checkmarkc\checkmarki\checkmarkr\checkmarkc\checkmark$\checkmarke\checkmark$\checkmark^\checkmark\\checkmarkc\checkmarki\checkmarkr\checkmarkc\checkmark$\checkmark \checkmark$\checkmark^\checkmark\\checkmarkc\checkmarki\checkmarkr\checkmarkc\checkmark$\checkmarkf\checkmark$\checkmark^\checkmark\\checkmarkc\checkmarki\checkmarkr\checkmarkc\checkmark$\checkmarkr\checkmark$\checkmark^\checkmark\\checkmarkc\checkmarki\checkmarkr\checkmarkc\checkmark$\checkmarké\checkmark$\checkmark^\checkmark\\checkmarkc\checkmarki\checkmarkr\checkmarkc\checkmark$\checkmarkq\checkmark$\checkmark^\checkmark\\checkmarkc\checkmarki\checkmarkr\checkmarkc\checkmark$\checkmarku\checkmark$\checkmark^\checkmark\\checkmarkc\checkmarki\checkmarkr\checkmarkc\checkmark$\checkmarke\checkmark$\checkmark^\checkmark\\checkmarkc\checkmarki\checkmarkr\checkmarkc\checkmark$\checkmarkn\checkmark$\checkmark^\checkmark\\checkmarkc\checkmarki\checkmarkr\checkmarkc\checkmark$\checkmarkc\checkmark$\checkmark^\checkmark\\checkmarkc\checkmarki\checkmarkr\checkmarkc\checkmark$\checkmarke\checkmark$\checkmark^\checkmark\\checkmarkc\checkmarki\checkmarkr\checkmarkc\checkmark$\checkmark}\checkmark$\checkmark^\checkmark\\checkmarkc\checkmarki\checkmarkr\checkmarkc\checkmark$\checkmark \checkmark$\checkmark^\checkmark\\checkmarkc\checkmarki\checkmarkr\checkmarkc\checkmark$\checkmark=\checkmark$\checkmark^\checkmark\\checkmarkc\checkmarki\checkmarkr\checkmarkc\checkmark$\checkmark \checkmark$\checkmark^\checkmark\\checkmarkc\checkmarki\checkmarkr\checkmarkc\checkmark$\checkmark\\checkmark$\checkmark^\checkmark\\checkmarkc\checkmarki\checkmarkr\checkmarkc\checkmark$\checkmarkf\checkmark$\checkmark^\checkmark\\checkmarkc\checkmarki\checkmarkr\checkmarkc\checkmark$\checkmarkr\checkmark$\checkmark^\checkmark\\checkmarkc\checkmarki\checkmarkr\checkmarkc\checkmark$\checkmarka\checkmark$\checkmark^\checkmark\\checkmarkc\checkmarki\checkmarkr\checkmarkc\checkmark$\checkmarkc\checkmark$\checkmark^\checkmark\\checkmarkc\checkmarki\checkmarkr\checkmarkc\checkmark$\checkmark{\checkmark$\checkmark^\checkmark\\checkmarkc\checkmarki\checkmarkr\checkmarkc\checkmark$\checkmarkf\checkmark$\checkmark^\checkmark\\checkmarkc\checkmarki\checkmarkr\checkmarkc\checkmark$\checkmark_\checkmark$\checkmark^\checkmark\\checkmarkc\checkmarki\checkmarkr\checkmarkc\checkmark$\checkmark{\checkmark$\checkmark^\checkmark\\checkmarkc\checkmarki\checkmarkr\checkmarkc\checkmark$\checkmarke\checkmark$\checkmark^\checkmark\\checkmarkc\checkmarki\checkmarkr\checkmarkc\checkmark$\checkmarkx\checkmark$\checkmark^\checkmark\\checkmarkc\checkmarki\checkmarkr\checkmarkc\checkmark$\checkmarkc\checkmark$\checkmark^\checkmark\\checkmarkc\checkmarki\checkmarkr\checkmarkc\checkmark$\checkmarki\checkmark$\checkmark^\checkmark\\checkmarkc\checkmarki\checkmarkr\checkmarkc\checkmark$\checkmarkt\checkmark$\checkmark^\checkmark\\checkmarkc\checkmarki\checkmarkr\checkmarkc\checkmark$\checkmarka\checkmark$\checkmark^\checkmark\\checkmarkc\checkmarki\checkmarkr\checkmarkc\checkmark$\checkmarkt\checkmark$\checkmark^\checkmark\\checkmarkc\checkmarki\checkmarkr\checkmarkc\checkmark$\checkmarki\checkmark$\checkmark^\checkmark\\checkmarkc\checkmarki\checkmarkr\checkmarkc\checkmark$\checkmarko\checkmark$\checkmark^\checkmark\\checkmarkc\checkmarki\checkmarkr\checkmarkc\checkmark$\checkmarkn\checkmark$\checkmark^\checkmark\\checkmarkc\checkmarki\checkmarkr\checkmarkc\checkmark$\checkmark}\checkmark$\checkmark^\checkmark\\checkmarkc\checkmarki\checkmarkr\checkmarkc\checkmark$\checkmark}\checkmark$\checkmark^\checkmark\\checkmarkc\checkmarki\checkmarkr\checkmarkc\checkmark$\checkmark{\checkmark$\checkmark^\checkmark\\checkmarkc\checkmarki\checkmarkr\checkmarkc\checkmark$\checkmarkf\checkmark$\checkmark^\checkmark\\checkmarkc\checkmarki\checkmarkr\checkmarkc\checkmark$\checkmark_\checkmark$\checkmark^\checkmark\\checkmarkc\checkmarki\checkmarkr\checkmarkc\checkmark$\checkmark1\checkmark$\checkmark^\checkmark\\checkmarkc\checkmarki\checkmarkr\checkmarkc\checkmark$\checkmark}\checkmark$\checkmark^\checkmark\\checkmarkc\checkmarki\checkmarkr\checkmarkc\checkmark$\checkmark \checkmark$\checkmark^\checkmark\\checkmarkc\checkmarki\checkmarkr\checkmarkc\checkmark$\checkmark=\checkmark$\checkmark^\checkmark\\checkmarkc\checkmarki\checkmarkr\checkmarkc\checkmark$\checkmark \checkmark$\checkmark^\checkmark\\checkmarkc\checkmarki\checkmarkr\checkmarkc\checkmark$\checkmark\\checkmark$\checkmark^\checkmark\\checkmarkc\checkmarki\checkmarkr\checkmarkc\checkmark$\checkmarkf\checkmark$\checkmark^\checkmark\\checkmarkc\checkmarki\checkmarkr\checkmarkc\checkmark$\checkmarkr\checkmark$\checkmark^\checkmark\\checkmarkc\checkmarki\checkmarkr\checkmarkc\checkmark$\checkmarka\checkmark$\checkmark^\checkmark\\checkmarkc\checkmarki\checkmarkr\checkmarkc\checkmark$\checkmarkc\checkmark$\checkmark^\checkmark\\checkmarkc\checkmarki\checkmarkr\checkmarkc\checkmark$\checkmark{\checkmark$\checkmark^\checkmark\\checkmarkc\checkmarki\checkmarkr\checkmarkc\checkmark$\checkmark6\checkmark$\checkmark^\checkmark\\checkmarkc\checkmarki\checkmarkr\checkmarkc\checkmark$\checkmark0\checkmark$\checkmark^\checkmark\\checkmarkc\checkmarki\checkmarkr\checkmarkc\checkmark$\checkmark0\checkmark$\checkmark^\checkmark\\checkmarkc\checkmarki\checkmarkr\checkmarkc\checkmark$\checkmark}\checkmark$\checkmark^\checkmark\\checkmarkc\checkmarki\checkmarkr\checkmarkc\checkmark$\checkmark{\checkmark$\checkmark^\checkmark\\checkmarkc\checkmarki\checkmarkr\checkmarkc\checkmark$\checkmark2\checkmark$\checkmark^\checkmark\\checkmarkc\checkmarki\checkmarkr\checkmarkc\checkmark$\checkmark9\checkmark$\checkmark^\checkmark\\checkmarkc\checkmarki\checkmarkr\checkmarkc\checkmark$\checkmark0\checkmark$\checkmark^\checkmark\\checkmarkc\checkmarki\checkmarkr\checkmarkc\checkmark$\checkmark}\checkmark$\checkmark^\checkmark\\checkmarkc\checkmarki\checkmarkr\checkmarkc\checkmark$\checkmark \checkmark$\checkmark^\checkmark\\checkmarkc\checkmarki\checkmarkr\checkmarkc\checkmark$\checkmark\\checkmark$\checkmark^\checkmark\\checkmarkc\checkmarki\checkmarkr\checkmarkc\checkmark$\checkmarka\checkmark$\checkmark^\checkmark\\checkmarkc\checkmarki\checkmarkr\checkmarkc\checkmark$\checkmarkp\checkmark$\checkmark^\checkmark\\checkmarkc\checkmarki\checkmarkr\checkmarkc\checkmark$\checkmarkp\checkmark$\checkmark^\checkmark\\checkmarkc\checkmarki\checkmarkr\checkmarkc\checkmark$\checkmarkr\checkmark$\checkmark^\checkmark\\checkmarkc\checkmarki\checkmarkr\checkmarkc\checkmark$\checkmarko\checkmark$\checkmark^\checkmark\\checkmarkc\checkmarki\checkmarkr\checkmarkc\checkmark$\checkmarkx\checkmark$\checkmark^\checkmark\\checkmarkc\checkmarki\checkmarkr\checkmarkc\checkmark$\checkmark \checkmark$\checkmark^\checkmark\\checkmarkc\checkmarki\checkmarkr\checkmarkc\checkmark$\checkmark2\checkmark$\checkmark^\checkmark\\checkmarkc\checkmarki\checkmarkr\checkmarkc\checkmark$\checkmark{\checkmark$\checkmark^\checkmark\\checkmarkc\checkmarki\checkmarkr\checkmarkc\checkmark$\checkmark,\checkmark$\checkmark^\checkmark\\checkmarkc\checkmarki\checkmarkr\checkmarkc\checkmark$\checkmark}\checkmark$\checkmark^\checkmark\\checkmarkc\checkmarki\checkmarkr\checkmarkc\checkmark$\checkmark0\checkmark$\checkmark^\checkmark\\checkmarkc\checkmarki\checkmarkr\checkmarkc\checkmark$\checkmark7\checkmark$\checkmark^\checkmark\\checkmarkc\checkmarki\checkmarkr\checkmarkc\checkmark$\checkmark \checkmark$\checkmark^\checkmark\\checkmarkc\checkmarki\checkmarkr\checkmarkc\checkmark$\checkmark>\checkmark$\checkmark^\checkmark\\checkmarkc\checkmarki\checkmarkr\checkmarkc\checkmark$\checkmark \checkmark$\checkmark^\checkmark\\checkmarkc\checkmarki\checkmarkr\checkmarkc\checkmark$\checkmark1\checkmark$\checkmark^\checkmark\\checkmarkc\checkmarki\checkmarkr\checkmarkc\checkmark$\checkmark{\checkmark$\checkmark^\checkmark\\checkmarkc\checkmarki\checkmarkr\checkmarkc\checkmark$\checkmark,\checkmark$\checkmark^\checkmark\\checkmarkc\checkmarki\checkmarkr\checkmarkc\checkmark$\checkmark}\checkmark$\checkmark^\checkmark\\checkmarkc\checkmarki\checkmarkr\checkmarkc\checkmark$\checkmark5\checkmark$\checkmark^\checkmark\\checkmarkc\checkmarki\checkmarkr\checkmarkc\checkmark$\checkmark
\checkmark$\checkmark^\checkmark\\checkmarkc\checkmarki\checkmarkr\checkmarkc\checkmark$\checkmark\\checkmark$\checkmark^\checkmark\\checkmarkc\checkmarki\checkmarkr\checkmarkc\checkmark$\checkmarke\checkmark$\checkmark^\checkmark\\checkmarkc\checkmarki\checkmarkr\checkmarkc\checkmark$\checkmarkn\checkmark$\checkmark^\checkmark\\checkmarkc\checkmarki\checkmarkr\checkmarkc\checkmark$\checkmarkd\checkmark$\checkmark^\checkmark\\checkmarkc\checkmarki\checkmarkr\checkmarkc\checkmark$\checkmark{\checkmark$\checkmark^\checkmark\\checkmarkc\checkmarki\checkmarkr\checkmarkc\checkmark$\checkmarke\checkmark$\checkmark^\checkmark\\checkmarkc\checkmarki\checkmarkr\checkmarkc\checkmark$\checkmarkq\checkmark$\checkmark^\checkmark\\checkmarkc\checkmarki\checkmarkr\checkmarkc\checkmark$\checkmarku\checkmark$\checkmark^\checkmark\\checkmarkc\checkmarki\checkmarkr\checkmarkc\checkmark$\checkmarka\checkmark$\checkmark^\checkmark\\checkmarkc\checkmarki\checkmarkr\checkmarkc\checkmark$\checkmarkt\checkmark$\checkmark^\checkmark\\checkmarkc\checkmarki\checkmarkr\checkmarkc\checkmark$\checkmarki\checkmark$\checkmark^\checkmark\\checkmarkc\checkmarki\checkmarkr\checkmarkc\checkmark$\checkmarko\checkmark$\checkmark^\checkmark\\checkmarkc\checkmarki\checkmarkr\checkmarkc\checkmark$\checkmarkn\checkmark$\checkmark^\checkmark\\checkmarkc\checkmarki\checkmarkr\checkmarkc\checkmark$\checkmark}\checkmark$\checkmark^\checkmark\\checkmarkc\checkmarki\checkmarkr\checkmarkc\checkmark$\checkmark
\checkmark$\checkmark^\checkmark\\checkmarkc\checkmarki\checkmarkr\checkmarkc\checkmark$\checkmarkL\checkmark$\checkmark^\checkmark\\checkmarkc\checkmarki\checkmarkr\checkmarkc\checkmark$\checkmarke\checkmark$\checkmark^\checkmark\\checkmarkc\checkmarki\checkmarkr\checkmarkc\checkmark$\checkmark \checkmark$\checkmark^\checkmark\\checkmarkc\checkmarki\checkmarkr\checkmarkc\checkmark$\checkmarkr\checkmark$\checkmark^\checkmark\\checkmarkc\checkmarki\checkmarkr\checkmarkc\checkmark$\checkmarka\checkmark$\checkmark^\checkmark\\checkmarkc\checkmarki\checkmarkr\checkmarkc\checkmark$\checkmarkp\checkmark$\checkmark^\checkmark\\checkmarkc\checkmarki\checkmarkr\checkmarkc\checkmark$\checkmarkp\checkmark$\checkmark^\checkmark\\checkmarkc\checkmarki\checkmarkr\checkmarkc\checkmark$\checkmarko\checkmark$\checkmark^\checkmark\\checkmarkc\checkmarki\checkmarkr\checkmarkc\checkmark$\checkmarkr\checkmark$\checkmark^\checkmark\\checkmarkc\checkmarki\checkmarkr\checkmarkc\checkmark$\checkmarkt\checkmark$\checkmark^\checkmark\\checkmarkc\checkmarki\checkmarkr\checkmarkc\checkmark$\checkmark \checkmark$\checkmark^\checkmark\\checkmarkc\checkmarki\checkmarkr\checkmarkc\checkmark$\checkmark$\checkmark$\checkmark^\checkmark\\checkmarkc\checkmarki\checkmarkr\checkmarkc\checkmark$\checkmarkf\checkmark$\checkmark^\checkmark\\checkmarkc\checkmarki\checkmarkr\checkmarkc\checkmark$\checkmark_\checkmark$\checkmark^\checkmark\\checkmarkc\checkmarki\checkmarkr\checkmarkc\checkmark$\checkmark{\checkmark$\checkmark^\checkmark\\checkmarkc\checkmarki\checkmarkr\checkmarkc\checkmark$\checkmarke\checkmark$\checkmark^\checkmark\\checkmarkc\checkmarki\checkmarkr\checkmarkc\checkmark$\checkmarkx\checkmark$\checkmark^\checkmark\\checkmarkc\checkmarki\checkmarkr\checkmarkc\checkmark$\checkmarkc\checkmark$\checkmark^\checkmark\\checkmarkc\checkmarki\checkmarkr\checkmarkc\checkmark$\checkmarki\checkmark$\checkmark^\checkmark\\checkmarkc\checkmarki\checkmarkr\checkmarkc\checkmark$\checkmarkt\checkmark$\checkmark^\checkmark\\checkmarkc\checkmarki\checkmarkr\checkmarkc\checkmark$\checkmarka\checkmark$\checkmark^\checkmark\\checkmarkc\checkmarki\checkmarkr\checkmarkc\checkmark$\checkmarkt\checkmark$\checkmark^\checkmark\\checkmarkc\checkmarki\checkmarkr\checkmarkc\checkmark$\checkmarki\checkmark$\checkmark^\checkmark\\checkmarkc\checkmarki\checkmarkr\checkmarkc\checkmark$\checkmarko\checkmark$\checkmark^\checkmark\\checkmarkc\checkmarki\checkmarkr\checkmarkc\checkmark$\checkmarkn\checkmark$\checkmark^\checkmark\\checkmarkc\checkmarki\checkmarkr\checkmarkc\checkmark$\checkmark}\checkmark$\checkmark^\checkmark\\checkmarkc\checkmarki\checkmarkr\checkmarkc\checkmark$\checkmark/\checkmark$\checkmark^\checkmark\\checkmarkc\checkmarki\checkmarkr\checkmarkc\checkmark$\checkmarkf\checkmark$\checkmark^\checkmark\\checkmarkc\checkmarki\checkmarkr\checkmarkc\checkmark$\checkmark_\checkmark$\checkmark^\checkmark\\checkmarkc\checkmarki\checkmarkr\checkmarkc\checkmark$\checkmark1\checkmark$\checkmark^\checkmark\\checkmarkc\checkmarki\checkmarkr\checkmarkc\checkmark$\checkmark \checkmark$\checkmark^\checkmark\\checkmarkc\checkmarki\checkmarkr\checkmarkc\checkmark$\checkmark\\checkmark$\checkmark^\checkmark\\checkmarkc\checkmarki\checkmarkr\checkmarkc\checkmark$\checkmarka\checkmark$\checkmark^\checkmark\\checkmarkc\checkmarki\checkmarkr\checkmarkc\checkmark$\checkmarkp\checkmark$\checkmark^\checkmark\\checkmarkc\checkmarki\checkmarkr\checkmarkc\checkmark$\checkmarkp\checkmark$\checkmark^\checkmark\\checkmarkc\checkmarki\checkmarkr\checkmarkc\checkmark$\checkmarkr\checkmark$\checkmark^\checkmark\\checkmarkc\checkmarki\checkmarkr\checkmarkc\checkmark$\checkmarko\checkmark$\checkmark^\checkmark\\checkmarkc\checkmarki\checkmarkr\checkmarkc\checkmark$\checkmarkx\checkmark$\checkmark^\checkmark\\checkmarkc\checkmarki\checkmarkr\checkmarkc\checkmark$\checkmark \checkmark$\checkmark^\checkmark\\checkmarkc\checkmarki\checkmarkr\checkmarkc\checkmark$\checkmark2\checkmark$\checkmark^\checkmark\\checkmarkc\checkmarki\checkmarkr\checkmarkc\checkmark$\checkmark{\checkmark$\checkmark^\checkmark\\checkmarkc\checkmarki\checkmarkr\checkmarkc\checkmark$\checkmark,\checkmark$\checkmark^\checkmark\\checkmarkc\checkmarki\checkmarkr\checkmarkc\checkmark$\checkmark}\checkmark$\checkmark^\checkmark\\checkmarkc\checkmarki\checkmarkr\checkmarkc\checkmark$\checkmark0\checkmark$\checkmark^\checkmark\\checkmarkc\checkmarki\checkmarkr\checkmarkc\checkmark$\checkmark7\checkmark$\checkmark^\checkmark\\checkmarkc\checkmarki\checkmarkr\checkmarkc\checkmark$\checkmark$\checkmark$\checkmark^\checkmark\\checkmarkc\checkmarki\checkmarkr\checkmarkc\checkmark$\checkmark \checkmark$\checkmark^\checkmark\\checkmarkc\checkmarki\checkmarkr\checkmarkc\checkmark$\checkmarkg\checkmark$\checkmark^\checkmark\\checkmarkc\checkmarki\checkmarkr\checkmarkc\checkmark$\checkmarka\checkmark$\checkmark^\checkmark\\checkmarkc\checkmarki\checkmarkr\checkmarkc\checkmark$\checkmarkr\checkmark$\checkmark^\checkmark\\checkmarkc\checkmarki\checkmarkr\checkmarkc\checkmark$\checkmarka\checkmark$\checkmark^\checkmark\\checkmarkc\checkmarki\checkmarkr\checkmarkc\checkmark$\checkmarkn\checkmark$\checkmark^\checkmark\\checkmarkc\checkmarki\checkmarkr\checkmarkc\checkmark$\checkmarkt\checkmark$\checkmark^\checkmark\\checkmarkc\checkmarki\checkmarkr\checkmarkc\checkmark$\checkmarki\checkmark$\checkmark^\checkmark\\checkmarkc\checkmarki\checkmarkr\checkmarkc\checkmark$\checkmarkt\checkmark$\checkmark^\checkmark\\checkmarkc\checkmarki\checkmarkr\checkmarkc\checkmark$\checkmark \checkmark$\checkmark^\checkmark\\checkmarkc\checkmarki\checkmarkr\checkmarkc\checkmark$\checkmarkq\checkmark$\checkmark^\checkmark\\checkmarkc\checkmarki\checkmarkr\checkmarkc\checkmark$\checkmarku\checkmark$\checkmark^\checkmark\\checkmarkc\checkmarki\checkmarkr\checkmarkc\checkmark$\checkmark'\checkmark$\checkmark^\checkmark\\checkmarkc\checkmarki\checkmarkr\checkmarkc\checkmark$\checkmarka\checkmark$\checkmark^\checkmark\\checkmarkc\checkmarki\checkmarkr\checkmarkc\checkmark$\checkmarku\checkmark$\checkmark^\checkmark\\checkmarkc\checkmarki\checkmarkr\checkmarkc\checkmark$\checkmarkc\checkmark$\checkmark^\checkmark\\checkmarkc\checkmarki\checkmarkr\checkmarkc\checkmark$\checkmarku\checkmark$\checkmark^\checkmark\\checkmarkc\checkmarki\checkmarkr\checkmarkc\checkmark$\checkmarkn\checkmark$\checkmark^\checkmark\\checkmarkc\checkmarki\checkmarkr\checkmarkc\checkmark$\checkmark \checkmark$\checkmark^\checkmark\\checkmarkc\checkmarki\checkmarkr\checkmarkc\checkmark$\checkmarkp\checkmark$\checkmark^\checkmark\\checkmarkc\checkmarki\checkmarkr\checkmarkc\checkmark$\checkmarkh\checkmark$\checkmark^\checkmark\\checkmarkc\checkmarki\checkmarkr\checkmarkc\checkmark$\checkmarké\checkmark$\checkmark^\checkmark\\checkmarkc\checkmarki\checkmarkr\checkmarkc\checkmark$\checkmarkn\checkmark$\checkmark^\checkmark\\checkmarkc\checkmarki\checkmarkr\checkmarkc\checkmark$\checkmarko\checkmark$\checkmark^\checkmark\\checkmarkc\checkmarki\checkmarkr\checkmarkc\checkmark$\checkmarkm\checkmark$\checkmark^\checkmark\\checkmarkc\checkmarki\checkmarkr\checkmarkc\checkmark$\checkmarkè\checkmark$\checkmark^\checkmark\\checkmarkc\checkmarki\checkmarkr\checkmarkc\checkmark$\checkmarkn\checkmark$\checkmark^\checkmark\\checkmarkc\checkmarki\checkmarkr\checkmarkc\checkmark$\checkmarke\checkmark$\checkmark^\checkmark\\checkmarkc\checkmarki\checkmarkr\checkmarkc\checkmark$\checkmark \checkmark$\checkmark^\checkmark\\checkmarkc\checkmarki\checkmarkr\checkmarkc\checkmark$\checkmarkd\checkmark$\checkmark^\checkmark\\checkmarkc\checkmarki\checkmarkr\checkmarkc\checkmark$\checkmarke\checkmark$\checkmark^\checkmark\\checkmarkc\checkmarki\checkmarkr\checkmarkc\checkmark$\checkmark \checkmark$\checkmark^\checkmark\\checkmarkc\checkmarki\checkmarkr\checkmarkc\checkmark$\checkmarkr\checkmark$\checkmark^\checkmark\\checkmarkc\checkmarki\checkmarkr\checkmarkc\checkmark$\checkmarké\checkmark$\checkmark^\checkmark\\checkmarkc\checkmarki\checkmarkr\checkmarkc\checkmark$\checkmarks\checkmark$\checkmark^\checkmark\\checkmarkc\checkmarki\checkmarkr\checkmarkc\checkmark$\checkmarko\checkmark$\checkmark^\checkmark\\checkmarkc\checkmarki\checkmarkr\checkmarkc\checkmark$\checkmarkn\checkmark$\checkmark^\checkmark\\checkmarkc\checkmarki\checkmarkr\checkmarkc\checkmark$\checkmarka\checkmark$\checkmark^\checkmark\\checkmarkc\checkmarki\checkmarkr\checkmarkc\checkmark$\checkmarkn\checkmark$\checkmark^\checkmark\\checkmarkc\checkmarki\checkmarkr\checkmarkc\checkmark$\checkmarkc\checkmark$\checkmark^\checkmark\\checkmarkc\checkmarki\checkmarkr\checkmarkc\checkmark$\checkmarke\checkmark$\checkmark^\checkmark\\checkmarkc\checkmarki\checkmarkr\checkmarkc\checkmark$\checkmark \checkmark$\checkmark^\checkmark\\checkmarkc\checkmarki\checkmarkr\checkmarkc\checkmark$\checkmarkm\checkmark$\checkmark^\checkmark\\checkmarkc\checkmarki\checkmarkr\checkmarkc\checkmark$\checkmarké\checkmark$\checkmark^\checkmark\\checkmarkc\checkmarki\checkmarkr\checkmarkc\checkmark$\checkmarkc\checkmark$\checkmark^\checkmark\\checkmarkc\checkmarki\checkmarkr\checkmarkc\checkmark$\checkmarka\checkmark$\checkmark^\checkmark\\checkmarkc\checkmarki\checkmarkr\checkmarkc\checkmark$\checkmarkn\checkmark$\checkmark^\checkmark\\checkmarkc\checkmarki\checkmarkr\checkmarkc\checkmark$\checkmarki\checkmark$\checkmark^\checkmark\\checkmarkc\checkmarki\checkmarkr\checkmarkc\checkmark$\checkmarkq\checkmark$\checkmark^\checkmark\\checkmarkc\checkmarki\checkmarkr\checkmarkc\checkmark$\checkmarku\checkmark$\checkmark^\checkmark\\checkmarkc\checkmarki\checkmarkr\checkmarkc\checkmark$\checkmarke\checkmark$\checkmark^\checkmark\\checkmarkc\checkmarki\checkmarkr\checkmarkc\checkmark$\checkmark \checkmark$\checkmark^\checkmark\\checkmarkc\checkmarki\checkmarkr\checkmarkc\checkmark$\checkmarkn\checkmark$\checkmark^\checkmark\\checkmarkc\checkmarki\checkmarkr\checkmarkc\checkmark$\checkmarke\checkmark$\checkmark^\checkmark\\checkmarkc\checkmarki\checkmarkr\checkmarkc\checkmark$\checkmark \checkmark$\checkmark^\checkmark\\checkmarkc\checkmarki\checkmarkr\checkmarkc\checkmark$\checkmarks\checkmark$\checkmark^\checkmark\\checkmarkc\checkmarki\checkmarkr\checkmarkc\checkmark$\checkmarke\checkmark$\checkmark^\checkmark\\checkmarkc\checkmarki\checkmarkr\checkmarkc\checkmark$\checkmark \checkmark$\checkmark^\checkmark\\checkmarkc\checkmarki\checkmarkr\checkmarkc\checkmark$\checkmarkp\checkmark$\checkmark^\checkmark\\checkmarkc\checkmarki\checkmarkr\checkmarkc\checkmark$\checkmarkr\checkmark$\checkmark^\checkmark\\checkmarkc\checkmarki\checkmarkr\checkmarkc\checkmark$\checkmarko\checkmark$\checkmark^\checkmark\\checkmarkc\checkmarki\checkmarkr\checkmarkc\checkmark$\checkmarkd\checkmark$\checkmark^\checkmark\\checkmarkc\checkmarki\checkmarkr\checkmarkc\checkmark$\checkmarku\checkmark$\checkmark^\checkmark\\checkmarkc\checkmarki\checkmarkr\checkmarkc\checkmark$\checkmarki\checkmark$\checkmark^\checkmark\\checkmarkc\checkmarki\checkmarkr\checkmarkc\checkmark$\checkmarkr\checkmark$\checkmark^\checkmark\\checkmarkc\checkmarki\checkmarkr\checkmarkc\checkmark$\checkmarka\checkmark$\checkmark^\checkmark\\checkmarkc\checkmarki\checkmarkr\checkmarkc\checkmark$\checkmark \checkmark$\checkmark^\checkmark\\checkmarkc\checkmarki\checkmarkr\checkmarkc\checkmark$\checkmarke\checkmark$\checkmark^\checkmark\\checkmarkc\checkmarki\checkmarkr\checkmarkc\checkmark$\checkmarkn\checkmark$\checkmark^\checkmark\\checkmarkc\checkmarki\checkmarkr\checkmarkc\checkmark$\checkmark \checkmark$\checkmark^\checkmark\\checkmarkc\checkmarki\checkmarkr\checkmarkc\checkmark$\checkmarkr\checkmark$\checkmark^\checkmark\\checkmarkc\checkmarki\checkmarkr\checkmarkc\checkmark$\checkmarké\checkmark$\checkmark^\checkmark\\checkmarkc\checkmarki\checkmarkr\checkmarkc\checkmark$\checkmarkg\checkmark$\checkmark^\checkmark\\checkmarkc\checkmarki\checkmarkr\checkmarkc\checkmark$\checkmarki\checkmark$\checkmark^\checkmark\\checkmarkc\checkmarki\checkmarkr\checkmarkc\checkmark$\checkmarkm\checkmark$\checkmark^\checkmark\\checkmarkc\checkmarki\checkmarkr\checkmarkc\checkmark$\checkmarke\checkmark$\checkmark^\checkmark\\checkmarkc\checkmarki\checkmarkr\checkmarkc\checkmark$\checkmark \checkmark$\checkmark^\checkmark\\checkmarkc\checkmarki\checkmarkr\checkmarkc\checkmark$\checkmarkn\checkmark$\checkmark^\checkmark\\checkmarkc\checkmarki\checkmarkr\checkmarkc\checkmark$\checkmarko\checkmark$\checkmark^\checkmark\\checkmarkc\checkmarki\checkmarkr\checkmarkc\checkmark$\checkmarkm\checkmark$\checkmark^\checkmark\\checkmarkc\checkmarki\checkmarkr\checkmarkc\checkmark$\checkmarki\checkmark$\checkmark^\checkmark\\checkmarkc\checkmarki\checkmarkr\checkmarkc\checkmark$\checkmarkn\checkmark$\checkmark^\checkmark\\checkmarkc\checkmarki\checkmarkr\checkmarkc\checkmark$\checkmarka\checkmark$\checkmark^\checkmark\\checkmarkc\checkmarki\checkmarkr\checkmarkc\checkmark$\checkmarkl\checkmark$\checkmark^\checkmark\\checkmarkc\checkmarki\checkmarkr\checkmarkc\checkmark$\checkmark.\checkmark$\checkmark^\checkmark\\checkmarkc\checkmarki\checkmarkr\checkmarkc\checkmark$\checkmark
\checkmark$\checkmark^\checkmark\\checkmarkc\checkmarki\checkmarkr\checkmarkc\checkmark$\checkmark
\checkmark$\checkmark^\checkmark\\checkmarkc\checkmarki\checkmarkr\checkmarkc\checkmark$\checkmark%\checkmark$\checkmark^\checkmark\\checkmarkc\checkmarki\checkmarkr\checkmarkc\checkmark$\checkmark \checkmark$\checkmark^\checkmark\\checkmarkc\checkmarki\checkmarkr\checkmarkc\checkmark$\checkmark=\checkmark$\checkmark^\checkmark\\checkmarkc\checkmarki\checkmarkr\checkmarkc\checkmark$\checkmark=\checkmark$\checkmark^\checkmark\\checkmarkc\checkmarki\checkmarkr\checkmarkc\checkmark$\checkmark=\checkmark$\checkmark^\checkmark\\checkmarkc\checkmarki\checkmarkr\checkmarkc\checkmark$\checkmark=\checkmark$\checkmark^\checkmark\\checkmarkc\checkmarki\checkmarkr\checkmarkc\checkmark$\checkmark=\checkmark$\checkmark^\checkmark\\checkmarkc\checkmarki\checkmarkr\checkmarkc\checkmark$\checkmark=\checkmark$\checkmark^\checkmark\\checkmarkc\checkmarki\checkmarkr\checkmarkc\checkmark$\checkmark=\checkmark$\checkmark^\checkmark\\checkmarkc\checkmarki\checkmarkr\checkmarkc\checkmark$\checkmark=\checkmark$\checkmark^\checkmark\\checkmarkc\checkmarki\checkmarkr\checkmarkc\checkmark$\checkmark=\checkmark$\checkmark^\checkmark\\checkmarkc\checkmarki\checkmarkr\checkmarkc\checkmark$\checkmark=\checkmark$\checkmark^\checkmark\\checkmarkc\checkmarki\checkmarkr\checkmarkc\checkmark$\checkmark=\checkmark$\checkmark^\checkmark\\checkmarkc\checkmarki\checkmarkr\checkmarkc\checkmark$\checkmark=\checkmark$\checkmark^\checkmark\\checkmarkc\checkmarki\checkmarkr\checkmarkc\checkmark$\checkmark=\checkmark$\checkmark^\checkmark\\checkmarkc\checkmarki\checkmarkr\checkmarkc\checkmark$\checkmark=\checkmark$\checkmark^\checkmark\\checkmarkc\checkmarki\checkmarkr\checkmarkc\checkmark$\checkmark=\checkmark$\checkmark^\checkmark\\checkmarkc\checkmarki\checkmarkr\checkmarkc\checkmark$\checkmark=\checkmark$\checkmark^\checkmark\\checkmarkc\checkmarki\checkmarkr\checkmarkc\checkmark$\checkmark=\checkmark$\checkmark^\checkmark\\checkmarkc\checkmarki\checkmarkr\checkmarkc\checkmark$\checkmark=\checkmark$\checkmark^\checkmark\\checkmarkc\checkmarki\checkmarkr\checkmarkc\checkmark$\checkmark=\checkmark$\checkmark^\checkmark\\checkmarkc\checkmarki\checkmarkr\checkmarkc\checkmark$\checkmark=\checkmark$\checkmark^\checkmark\\checkmarkc\checkmarki\checkmarkr\checkmarkc\checkmark$\checkmark=\checkmark$\checkmark^\checkmark\\checkmarkc\checkmarki\checkmarkr\checkmarkc\checkmark$\checkmark=\checkmark$\checkmark^\checkmark\\checkmarkc\checkmarki\checkmarkr\checkmarkc\checkmark$\checkmark=\checkmark$\checkmark^\checkmark\\checkmarkc\checkmarki\checkmarkr\checkmarkc\checkmark$\checkmark=\checkmark$\checkmark^\checkmark\\checkmarkc\checkmarki\checkmarkr\checkmarkc\checkmark$\checkmark=\checkmark$\checkmark^\checkmark\\checkmarkc\checkmarki\checkmarkr\checkmarkc\checkmark$\checkmark=\checkmark$\checkmark^\checkmark\\checkmarkc\checkmarki\checkmarkr\checkmarkc\checkmark$\checkmark=\checkmark$\checkmark^\checkmark\\checkmarkc\checkmarki\checkmarkr\checkmarkc\checkmark$\checkmark=\checkmark$\checkmark^\checkmark\\checkmarkc\checkmarki\checkmarkr\checkmarkc\checkmark$\checkmark=\checkmark$\checkmark^\checkmark\\checkmarkc\checkmarki\checkmarkr\checkmarkc\checkmark$\checkmark=\checkmark$\checkmark^\checkmark\\checkmarkc\checkmarki\checkmarkr\checkmarkc\checkmark$\checkmark=\checkmark$\checkmark^\checkmark\\checkmarkc\checkmarki\checkmarkr\checkmarkc\checkmark$\checkmark=\checkmark$\checkmark^\checkmark\\checkmarkc\checkmarki\checkmarkr\checkmarkc\checkmark$\checkmark=\checkmark$\checkmark^\checkmark\\checkmarkc\checkmarki\checkmarkr\checkmarkc\checkmark$\checkmark=\checkmark$\checkmark^\checkmark\\checkmarkc\checkmarki\checkmarkr\checkmarkc\checkmark$\checkmark=\checkmark$\checkmark^\checkmark\\checkmarkc\checkmarki\checkmarkr\checkmarkc\checkmark$\checkmark=\checkmark$\checkmark^\checkmark\\checkmarkc\checkmarki\checkmarkr\checkmarkc\checkmark$\checkmark=\checkmark$\checkmark^\checkmark\\checkmarkc\checkmarki\checkmarkr\checkmarkc\checkmark$\checkmark=\checkmark$\checkmark^\checkmark\\checkmarkc\checkmarki\checkmarkr\checkmarkc\checkmark$\checkmark=\checkmark$\checkmark^\checkmark\\checkmarkc\checkmarki\checkmarkr\checkmarkc\checkmark$\checkmark=\checkmark$\checkmark^\checkmark\\checkmarkc\checkmarki\checkmarkr\checkmarkc\checkmark$\checkmark=\checkmark$\checkmark^\checkmark\\checkmarkc\checkmarki\checkmarkr\checkmarkc\checkmark$\checkmark=\checkmark$\checkmark^\checkmark\\checkmarkc\checkmarki\checkmarkr\checkmarkc\checkmark$\checkmark=\checkmark$\checkmark^\checkmark\\checkmarkc\checkmarki\checkmarkr\checkmarkc\checkmark$\checkmark=\checkmark$\checkmark^\checkmark\\checkmarkc\checkmarki\checkmarkr\checkmarkc\checkmark$\checkmark=\checkmark$\checkmark^\checkmark\\checkmarkc\checkmarki\checkmarkr\checkmarkc\checkmark$\checkmark=\checkmark$\checkmark^\checkmark\\checkmarkc\checkmarki\checkmarkr\checkmarkc\checkmark$\checkmark=\checkmark$\checkmark^\checkmark\\checkmarkc\checkmarki\checkmarkr\checkmarkc\checkmark$\checkmark=\checkmark$\checkmark^\checkmark\\checkmarkc\checkmarki\checkmarkr\checkmarkc\checkmark$\checkmark=\checkmark$\checkmark^\checkmark\\checkmarkc\checkmarki\checkmarkr\checkmarkc\checkmark$\checkmark=\checkmark$\checkmark^\checkmark\\checkmarkc\checkmarki\checkmarkr\checkmarkc\checkmark$\checkmark=\checkmark$\checkmark^\checkmark\\checkmarkc\checkmarki\checkmarkr\checkmarkc\checkmark$\checkmark=\checkmark$\checkmark^\checkmark\\checkmarkc\checkmarki\checkmarkr\checkmarkc\checkmark$\checkmark=\checkmark$\checkmark^\checkmark\\checkmarkc\checkmarki\checkmarkr\checkmarkc\checkmark$\checkmark=\checkmark$\checkmark^\checkmark\\checkmarkc\checkmarki\checkmarkr\checkmarkc\checkmark$\checkmark=\checkmark$\checkmark^\checkmark\\checkmarkc\checkmarki\checkmarkr\checkmarkc\checkmark$\checkmark=\checkmark$\checkmark^\checkmark\\checkmarkc\checkmarki\checkmarkr\checkmarkc\checkmark$\checkmark=\checkmark$\checkmark^\checkmark\\checkmarkc\checkmarki\checkmarkr\checkmarkc\checkmark$\checkmark=\checkmark$\checkmark^\checkmark\\checkmarkc\checkmarki\checkmarkr\checkmarkc\checkmark$\checkmark=\checkmark$\checkmark^\checkmark\\checkmarkc\checkmarki\checkmarkr\checkmarkc\checkmark$\checkmark=\checkmark$\checkmark^\checkmark\\checkmarkc\checkmarki\checkmarkr\checkmarkc\checkmark$\checkmark=\checkmark$\checkmark^\checkmark\\checkmarkc\checkmarki\checkmarkr\checkmarkc\checkmark$\checkmark=\checkmark$\checkmark^\checkmark\\checkmarkc\checkmarki\checkmarkr\checkmarkc\checkmark$\checkmark
\checkmark$\checkmark^\checkmark\\checkmarkc\checkmarki\checkmarkr\checkmarkc\checkmark$\checkmark\\checkmark$\checkmark^\checkmark\\checkmarkc\checkmarki\checkmarkr\checkmarkc\checkmark$\checkmarks\checkmark$\checkmark^\checkmark\\checkmarkc\checkmarki\checkmarkr\checkmarkc\checkmark$\checkmarke\checkmark$\checkmark^\checkmark\\checkmarkc\checkmarki\checkmarkr\checkmarkc\checkmark$\checkmarkc\checkmark$\checkmark^\checkmark\\checkmarkc\checkmarki\checkmarkr\checkmarkc\checkmark$\checkmarkt\checkmark$\checkmark^\checkmark\\checkmarkc\checkmarki\checkmarkr\checkmarkc\checkmark$\checkmarki\checkmark$\checkmark^\checkmark\\checkmarkc\checkmarki\checkmarkr\checkmarkc\checkmark$\checkmarko\checkmark$\checkmark^\checkmark\\checkmarkc\checkmarki\checkmarkr\checkmarkc\checkmark$\checkmarkn\checkmark$\checkmark^\checkmark\\checkmarkc\checkmarki\checkmarkr\checkmarkc\checkmark$\checkmark{\checkmark$\checkmark^\checkmark\\checkmarkc\checkmarki\checkmarkr\checkmarkc\checkmark$\checkmarkS\checkmark$\checkmark^\checkmark\\checkmarkc\checkmarki\checkmarkr\checkmarkc\checkmark$\checkmarky\checkmark$\checkmark^\checkmark\\checkmarkc\checkmarki\checkmarkr\checkmarkc\checkmark$\checkmarkn\checkmark$\checkmark^\checkmark\\checkmarkc\checkmarki\checkmarkr\checkmarkc\checkmark$\checkmarkt\checkmark$\checkmark^\checkmark\\checkmarkc\checkmarki\checkmarkr\checkmarkc\checkmark$\checkmarkh\checkmark$\checkmark^\checkmark\\checkmarkc\checkmarki\checkmarkr\checkmarkc\checkmark$\checkmarkè\checkmark$\checkmark^\checkmark\\checkmarkc\checkmarki\checkmarkr\checkmarkc\checkmark$\checkmarks\checkmark$\checkmark^\checkmark\\checkmarkc\checkmarki\checkmarkr\checkmarkc\checkmark$\checkmarke\checkmark$\checkmark^\checkmark\\checkmarkc\checkmarki\checkmarkr\checkmarkc\checkmark$\checkmark \checkmark$\checkmark^\checkmark\\checkmarkc\checkmarki\checkmarkr\checkmarkc\checkmark$\checkmarke\checkmark$\checkmark^\checkmark\\checkmarkc\checkmarki\checkmarkr\checkmarkc\checkmark$\checkmarkt\checkmark$\checkmark^\checkmark\\checkmarkc\checkmarki\checkmarkr\checkmarkc\checkmark$\checkmark \checkmark$\checkmark^\checkmark\\checkmarkc\checkmarki\checkmarkr\checkmarkc\checkmark$\checkmarkC\checkmark$\checkmark^\checkmark\\checkmarkc\checkmarki\checkmarkr\checkmarkc\checkmark$\checkmarko\checkmark$\checkmark^\checkmark\\checkmarkc\checkmarki\checkmarkr\checkmarkc\checkmark$\checkmarkn\checkmark$\checkmark^\checkmark\\checkmarkc\checkmarki\checkmarkr\checkmarkc\checkmark$\checkmarkc\checkmark$\checkmark^\checkmark\\checkmarkc\checkmarki\checkmarkr\checkmarkc\checkmark$\checkmarkl\checkmark$\checkmark^\checkmark\\checkmarkc\checkmarki\checkmarkr\checkmarkc\checkmark$\checkmarku\checkmark$\checkmark^\checkmark\\checkmarkc\checkmarki\checkmarkr\checkmarkc\checkmark$\checkmarks\checkmark$\checkmark^\checkmark\\checkmarkc\checkmarki\checkmarkr\checkmarkc\checkmark$\checkmarki\checkmark$\checkmark^\checkmark\\checkmarkc\checkmarki\checkmarkr\checkmarkc\checkmark$\checkmarko\checkmark$\checkmark^\checkmark\\checkmarkc\checkmarki\checkmarkr\checkmarkc\checkmark$\checkmarkn\checkmark$\checkmark^\checkmark\\checkmarkc\checkmarki\checkmarkr\checkmarkc\checkmark$\checkmark \checkmark$\checkmark^\checkmark\\checkmarkc\checkmarki\checkmarkr\checkmarkc\checkmark$\checkmarkd\checkmark$\checkmark^\checkmark\\checkmarkc\checkmarki\checkmarkr\checkmarkc\checkmark$\checkmarku\checkmark$\checkmark^\checkmark\\checkmarkc\checkmarki\checkmarkr\checkmarkc\checkmark$\checkmark \checkmark$\checkmark^\checkmark\\checkmarkc\checkmarki\checkmarkr\checkmarkc\checkmark$\checkmarkD\checkmark$\checkmark^\checkmark\\checkmarkc\checkmarki\checkmarkr\checkmarkc\checkmark$\checkmarki\checkmark$\checkmark^\checkmark\\checkmarkc\checkmarki\checkmarkr\checkmarkc\checkmark$\checkmarkm\checkmark$\checkmark^\checkmark\\checkmarkc\checkmarki\checkmarkr\checkmarkc\checkmark$\checkmarke\checkmark$\checkmark^\checkmark\\checkmarkc\checkmarki\checkmarkr\checkmarkc\checkmark$\checkmarkn\checkmark$\checkmark^\checkmark\\checkmarkc\checkmarki\checkmarkr\checkmarkc\checkmark$\checkmarks\checkmark$\checkmark^\checkmark\\checkmarkc\checkmarki\checkmarkr\checkmarkc\checkmark$\checkmarki\checkmark$\checkmark^\checkmark\\checkmarkc\checkmarki\checkmarkr\checkmarkc\checkmark$\checkmarko\checkmark$\checkmark^\checkmark\\checkmarkc\checkmarki\checkmarkr\checkmarkc\checkmark$\checkmarkn\checkmark$\checkmark^\checkmark\\checkmarkc\checkmarki\checkmarkr\checkmarkc\checkmark$\checkmarkn\checkmark$\checkmark^\checkmark\\checkmarkc\checkmarki\checkmarkr\checkmarkc\checkmark$\checkmarke\checkmark$\checkmark^\checkmark\\checkmarkc\checkmarki\checkmarkr\checkmarkc\checkmark$\checkmarkm\checkmark$\checkmark^\checkmark\\checkmarkc\checkmarki\checkmarkr\checkmarkc\checkmark$\checkmarke\checkmark$\checkmark^\checkmark\\checkmarkc\checkmarki\checkmarkr\checkmarkc\checkmark$\checkmarkn\checkmark$\checkmark^\checkmark\\checkmarkc\checkmarki\checkmarkr\checkmarkc\checkmark$\checkmarkt\checkmark$\checkmark^\checkmark\\checkmarkc\checkmarki\checkmarkr\checkmarkc\checkmark$\checkmark}\checkmark$\checkmark^\checkmark\\checkmarkc\checkmarki\checkmarkr\checkmarkc\checkmark$\checkmark
\checkmark$\checkmark^\checkmark\\checkmarkc\checkmarki\checkmarkr\checkmarkc\checkmark$\checkmark
\checkmark$\checkmark^\checkmark\\checkmarkc\checkmarki\checkmarkr\checkmarkc\checkmark$\checkmark\\checkmark$\checkmark^\checkmark\\checkmarkc\checkmarki\checkmarkr\checkmarkc\checkmark$\checkmarkb\checkmark$\checkmark^\checkmark\\checkmarkc\checkmarki\checkmarkr\checkmarkc\checkmark$\checkmarke\checkmark$\checkmark^\checkmark\\checkmarkc\checkmarki\checkmarkr\checkmarkc\checkmark$\checkmarkg\checkmark$\checkmark^\checkmark\\checkmarkc\checkmarki\checkmarkr\checkmarkc\checkmark$\checkmarki\checkmark$\checkmark^\checkmark\\checkmarkc\checkmarki\checkmarkr\checkmarkc\checkmark$\checkmarkn\checkmark$\checkmark^\checkmark\\checkmarkc\checkmarki\checkmarkr\checkmarkc\checkmark$\checkmark{\checkmark$\checkmark^\checkmark\\checkmarkc\checkmarki\checkmarkr\checkmarkc\checkmark$\checkmarkt\checkmark$\checkmark^\checkmark\\checkmarkc\checkmarki\checkmarkr\checkmarkc\checkmark$\checkmarka\checkmark$\checkmark^\checkmark\\checkmarkc\checkmarki\checkmarkr\checkmarkc\checkmark$\checkmarkb\checkmark$\checkmark^\checkmark\\checkmarkc\checkmarki\checkmarkr\checkmarkc\checkmark$\checkmarkl\checkmark$\checkmark^\checkmark\\checkmarkc\checkmarki\checkmarkr\checkmarkc\checkmark$\checkmarke\checkmark$\checkmark^\checkmark\\checkmarkc\checkmarki\checkmarkr\checkmarkc\checkmark$\checkmark}\checkmark$\checkmark^\checkmark\\checkmarkc\checkmarki\checkmarkr\checkmarkc\checkmark$\checkmark[\checkmark$\checkmark^\checkmark\\checkmarkc\checkmarki\checkmarkr\checkmarkc\checkmark$\checkmarkh\checkmark$\checkmark^\checkmark\\checkmarkc\checkmarki\checkmarkr\checkmarkc\checkmark$\checkmarkt\checkmark$\checkmark^\checkmark\\checkmarkc\checkmarki\checkmarkr\checkmarkc\checkmark$\checkmarkb\checkmark$\checkmark^\checkmark\\checkmarkc\checkmarki\checkmarkr\checkmarkc\checkmark$\checkmarkp\checkmark$\checkmark^\checkmark\\checkmarkc\checkmarki\checkmarkr\checkmarkc\checkmark$\checkmark]\checkmark$\checkmark^\checkmark\\checkmarkc\checkmarki\checkmarkr\checkmarkc\checkmark$\checkmark
\checkmark$\checkmark^\checkmark\\checkmarkc\checkmarki\checkmarkr\checkmarkc\checkmark$\checkmark\\checkmark$\checkmark^\checkmark\\checkmarkc\checkmarki\checkmarkr\checkmarkc\checkmark$\checkmarkc\checkmark$\checkmark^\checkmark\\checkmarkc\checkmarki\checkmarkr\checkmarkc\checkmark$\checkmarke\checkmark$\checkmark^\checkmark\\checkmarkc\checkmarki\checkmarkr\checkmarkc\checkmark$\checkmarkn\checkmark$\checkmark^\checkmark\\checkmarkc\checkmarki\checkmarkr\checkmarkc\checkmark$\checkmarkt\checkmark$\checkmark^\checkmark\\checkmarkc\checkmarki\checkmarkr\checkmarkc\checkmark$\checkmarke\checkmark$\checkmark^\checkmark\\checkmarkc\checkmarki\checkmarkr\checkmarkc\checkmark$\checkmarkr\checkmark$\checkmark^\checkmark\\checkmarkc\checkmarki\checkmarkr\checkmarkc\checkmark$\checkmarki\checkmark$\checkmark^\checkmark\\checkmarkc\checkmarki\checkmarkr\checkmarkc\checkmark$\checkmarkn\checkmark$\checkmark^\checkmark\\checkmarkc\checkmarki\checkmarkr\checkmarkc\checkmark$\checkmarkg\checkmark$\checkmark^\checkmark\\checkmarkc\checkmarki\checkmarkr\checkmarkc\checkmark$\checkmark
\checkmark$\checkmark^\checkmark\\checkmarkc\checkmarki\checkmarkr\checkmarkc\checkmark$\checkmark\\checkmark$\checkmark^\checkmark\\checkmarkc\checkmarki\checkmarkr\checkmarkc\checkmark$\checkmarkr\checkmark$\checkmark^\checkmark\\checkmarkc\checkmarki\checkmarkr\checkmarkc\checkmark$\checkmarke\checkmark$\checkmark^\checkmark\\checkmarkc\checkmarki\checkmarkr\checkmarkc\checkmark$\checkmarks\checkmark$\checkmark^\checkmark\\checkmarkc\checkmarki\checkmarkr\checkmarkc\checkmark$\checkmarki\checkmark$\checkmark^\checkmark\\checkmarkc\checkmarki\checkmarkr\checkmarkc\checkmark$\checkmarkz\checkmark$\checkmark^\checkmark\\checkmarkc\checkmarki\checkmarkr\checkmarkc\checkmark$\checkmarke\checkmark$\checkmark^\checkmark\\checkmarkc\checkmarki\checkmarkr\checkmarkc\checkmark$\checkmarkb\checkmark$\checkmark^\checkmark\\checkmarkc\checkmarki\checkmarkr\checkmarkc\checkmark$\checkmarko\checkmark$\checkmark^\checkmark\\checkmarkc\checkmarki\checkmarkr\checkmarkc\checkmark$\checkmarkx\checkmark$\checkmark^\checkmark\\checkmarkc\checkmarki\checkmarkr\checkmarkc\checkmark$\checkmark{\checkmark$\checkmark^\checkmark\\checkmarkc\checkmarki\checkmarkr\checkmarkc\checkmark$\checkmark\\checkmark$\checkmark^\checkmark\\checkmarkc\checkmarki\checkmarkr\checkmarkc\checkmark$\checkmarkt\checkmark$\checkmark^\checkmark\\checkmarkc\checkmarki\checkmarkr\checkmarkc\checkmark$\checkmarke\checkmark$\checkmark^\checkmark\\checkmarkc\checkmarki\checkmarkr\checkmarkc\checkmark$\checkmarkx\checkmark$\checkmark^\checkmark\\checkmarkc\checkmarki\checkmarkr\checkmarkc\checkmark$\checkmarkt\checkmark$\checkmark^\checkmark\\checkmarkc\checkmarki\checkmarkr\checkmarkc\checkmark$\checkmarkw\checkmark$\checkmark^\checkmark\\checkmarkc\checkmarki\checkmarkr\checkmarkc\checkmark$\checkmarki\checkmark$\checkmark^\checkmark\\checkmarkc\checkmarki\checkmarkr\checkmarkc\checkmark$\checkmarkd\checkmark$\checkmark^\checkmark\\checkmarkc\checkmarki\checkmarkr\checkmarkc\checkmark$\checkmarkt\checkmark$\checkmark^\checkmark\\checkmarkc\checkmarki\checkmarkr\checkmarkc\checkmark$\checkmarkh\checkmark$\checkmark^\checkmark\\checkmarkc\checkmarki\checkmarkr\checkmarkc\checkmark$\checkmark}\checkmark$\checkmark^\checkmark\\checkmarkc\checkmarki\checkmarkr\checkmarkc\checkmark$\checkmark{\checkmark$\checkmark^\checkmark\\checkmarkc\checkmarki\checkmarkr\checkmarkc\checkmark$\checkmark!\checkmark$\checkmark^\checkmark\\checkmarkc\checkmarki\checkmarkr\checkmarkc\checkmark$\checkmark}\checkmark$\checkmark^\checkmark\\checkmarkc\checkmarki\checkmarkr\checkmarkc\checkmark$\checkmark{\checkmark$\checkmark^\checkmark\\checkmarkc\checkmarki\checkmarkr\checkmarkc\checkmark$\checkmark
\checkmark$\checkmark^\checkmark\\checkmarkc\checkmarki\checkmarkr\checkmarkc\checkmark$\checkmark\\checkmark$\checkmark^\checkmark\\checkmarkc\checkmarki\checkmarkr\checkmarkc\checkmark$\checkmarkb\checkmark$\checkmark^\checkmark\\checkmarkc\checkmarki\checkmarkr\checkmarkc\checkmark$\checkmarke\checkmark$\checkmark^\checkmark\\checkmarkc\checkmarki\checkmarkr\checkmarkc\checkmark$\checkmarkg\checkmark$\checkmark^\checkmark\\checkmarkc\checkmarki\checkmarkr\checkmarkc\checkmark$\checkmarki\checkmark$\checkmark^\checkmark\\checkmarkc\checkmarki\checkmarkr\checkmarkc\checkmark$\checkmarkn\checkmark$\checkmark^\checkmark\\checkmarkc\checkmarki\checkmarkr\checkmarkc\checkmark$\checkmark{\checkmark$\checkmark^\checkmark\\checkmarkc\checkmarki\checkmarkr\checkmarkc\checkmark$\checkmarkt\checkmark$\checkmark^\checkmark\\checkmarkc\checkmarki\checkmarkr\checkmarkc\checkmark$\checkmarka\checkmark$\checkmark^\checkmark\\checkmarkc\checkmarki\checkmarkr\checkmarkc\checkmark$\checkmarkb\checkmark$\checkmark^\checkmark\\checkmarkc\checkmarki\checkmarkr\checkmarkc\checkmark$\checkmarku\checkmark$\checkmark^\checkmark\\checkmarkc\checkmarki\checkmarkr\checkmarkc\checkmark$\checkmarkl\checkmark$\checkmark^\checkmark\\checkmarkc\checkmarki\checkmarkr\checkmarkc\checkmark$\checkmarka\checkmark$\checkmark^\checkmark\\checkmarkc\checkmarki\checkmarkr\checkmarkc\checkmark$\checkmarkr\checkmark$\checkmark^\checkmark\\checkmarkc\checkmarki\checkmarkr\checkmarkc\checkmark$\checkmark}\checkmark$\checkmark^\checkmark\\checkmarkc\checkmarki\checkmarkr\checkmarkc\checkmark$\checkmark{\checkmark$\checkmark^\checkmark\\checkmarkc\checkmarki\checkmarkr\checkmarkc\checkmark$\checkmark|\checkmark$\checkmark^\checkmark\\checkmarkc\checkmarki\checkmarkr\checkmarkc\checkmark$\checkmarkl\checkmark$\checkmark^\checkmark\\checkmarkc\checkmarki\checkmarkr\checkmarkc\checkmark$\checkmark|\checkmark$\checkmark^\checkmark\\checkmarkc\checkmarki\checkmarkr\checkmarkc\checkmark$\checkmarkl\checkmark$\checkmark^\checkmark\\checkmarkc\checkmarki\checkmarkr\checkmarkc\checkmark$\checkmark|\checkmark$\checkmark^\checkmark\\checkmarkc\checkmarki\checkmarkr\checkmarkc\checkmark$\checkmarkc\checkmark$\checkmark^\checkmark\\checkmarkc\checkmarki\checkmarkr\checkmarkc\checkmark$\checkmark|\checkmark$\checkmark^\checkmark\\checkmarkc\checkmarki\checkmarkr\checkmarkc\checkmark$\checkmarkc\checkmark$\checkmark^\checkmark\\checkmarkc\checkmarki\checkmarkr\checkmarkc\checkmark$\checkmark|\checkmark$\checkmark^\checkmark\\checkmarkc\checkmarki\checkmarkr\checkmarkc\checkmark$\checkmarkc\checkmark$\checkmark^\checkmark\\checkmarkc\checkmarki\checkmarkr\checkmarkc\checkmark$\checkmark|\checkmark$\checkmark^\checkmark\\checkmarkc\checkmarki\checkmarkr\checkmarkc\checkmark$\checkmark}\checkmark$\checkmark^\checkmark\\checkmarkc\checkmarki\checkmarkr\checkmarkc\checkmark$\checkmark
\checkmark$\checkmark^\checkmark\\checkmarkc\checkmarki\checkmarkr\checkmarkc\checkmark$\checkmark\\checkmark$\checkmark^\checkmark\\checkmarkc\checkmarki\checkmarkr\checkmarkc\checkmark$\checkmarkh\checkmark$\checkmark^\checkmark\\checkmarkc\checkmarki\checkmarkr\checkmarkc\checkmark$\checkmarkl\checkmark$\checkmark^\checkmark\\checkmarkc\checkmarki\checkmarkr\checkmarkc\checkmark$\checkmarki\checkmark$\checkmark^\checkmark\\checkmarkc\checkmarki\checkmarkr\checkmarkc\checkmark$\checkmarkn\checkmark$\checkmark^\checkmark\\checkmarkc\checkmarki\checkmarkr\checkmarkc\checkmark$\checkmarke\checkmark$\checkmark^\checkmark\\checkmarkc\checkmarki\checkmarkr\checkmarkc\checkmark$\checkmark
\checkmark$\checkmark^\checkmark\\checkmarkc\checkmarki\checkmarkr\checkmarkc\checkmark$\checkmark\\checkmark$\checkmark^\checkmark\\checkmarkc\checkmarki\checkmarkr\checkmarkc\checkmark$\checkmarkt\checkmark$\checkmark^\checkmark\\checkmarkc\checkmarki\checkmarkr\checkmarkc\checkmark$\checkmarke\checkmark$\checkmark^\checkmark\\checkmarkc\checkmarki\checkmarkr\checkmarkc\checkmark$\checkmarkx\checkmark$\checkmark^\checkmark\\checkmarkc\checkmarki\checkmarkr\checkmarkc\checkmark$\checkmarkt\checkmark$\checkmark^\checkmark\\checkmarkc\checkmarki\checkmarkr\checkmarkc\checkmark$\checkmarkb\checkmark$\checkmark^\checkmark\\checkmarkc\checkmarki\checkmarkr\checkmarkc\checkmark$\checkmarkf\checkmark$\checkmark^\checkmark\\checkmarkc\checkmarki\checkmarkr\checkmarkc\checkmark$\checkmark{\checkmark$\checkmark^\checkmark\\checkmarkc\checkmarki\checkmarkr\checkmarkc\checkmark$\checkmarkO\checkmark$\checkmark^\checkmark\\checkmarkc\checkmarki\checkmarkr\checkmarkc\checkmark$\checkmarkr\checkmark$\checkmark^\checkmark\\checkmarkc\checkmarki\checkmarkr\checkmarkc\checkmark$\checkmarkg\checkmark$\checkmark^\checkmark\\checkmarkc\checkmarki\checkmarkr\checkmarkc\checkmark$\checkmarka\checkmark$\checkmark^\checkmark\\checkmarkc\checkmarki\checkmarkr\checkmarkc\checkmark$\checkmarkn\checkmark$\checkmark^\checkmark\\checkmarkc\checkmarki\checkmarkr\checkmarkc\checkmark$\checkmarke\checkmark$\checkmark^\checkmark\\checkmarkc\checkmarki\checkmarkr\checkmarkc\checkmark$\checkmark}\checkmark$\checkmark^\checkmark\\checkmarkc\checkmarki\checkmarkr\checkmarkc\checkmark$\checkmark \checkmark$\checkmark^\checkmark\\checkmarkc\checkmarki\checkmarkr\checkmarkc\checkmark$\checkmark&\checkmark$\checkmark^\checkmark\\checkmarkc\checkmarki\checkmarkr\checkmarkc\checkmark$\checkmark \checkmark$\checkmark^\checkmark\\checkmarkc\checkmarki\checkmarkr\checkmarkc\checkmark$\checkmark\\checkmark$\checkmark^\checkmark\\checkmarkc\checkmarki\checkmarkr\checkmarkc\checkmark$\checkmarkt\checkmark$\checkmark^\checkmark\\checkmarkc\checkmarki\checkmarkr\checkmarkc\checkmark$\checkmarke\checkmark$\checkmark^\checkmark\\checkmarkc\checkmarki\checkmarkr\checkmarkc\checkmark$\checkmarkx\checkmark$\checkmark^\checkmark\\checkmarkc\checkmarki\checkmarkr\checkmarkc\checkmark$\checkmarkt\checkmark$\checkmark^\checkmark\\checkmarkc\checkmarki\checkmarkr\checkmarkc\checkmark$\checkmarkb\checkmark$\checkmark^\checkmark\\checkmarkc\checkmarki\checkmarkr\checkmarkc\checkmark$\checkmarkf\checkmark$\checkmark^\checkmark\\checkmarkc\checkmarki\checkmarkr\checkmarkc\checkmark$\checkmark{\checkmark$\checkmark^\checkmark\\checkmarkc\checkmarki\checkmarkr\checkmarkc\checkmark$\checkmarkC\checkmark$\checkmark^\checkmark\\checkmarkc\checkmarki\checkmarkr\checkmarkc\checkmark$\checkmarkr\checkmark$\checkmark^\checkmark\\checkmarkc\checkmarki\checkmarkr\checkmarkc\checkmark$\checkmarki\checkmark$\checkmark^\checkmark\\checkmarkc\checkmarki\checkmarkr\checkmarkc\checkmark$\checkmarkt\checkmark$\checkmark^\checkmark\\checkmarkc\checkmarki\checkmarkr\checkmarkc\checkmark$\checkmarkè\checkmark$\checkmark^\checkmark\\checkmarkc\checkmarki\checkmarkr\checkmarkc\checkmark$\checkmarkr\checkmark$\checkmark^\checkmark\\checkmarkc\checkmarki\checkmarkr\checkmarkc\checkmark$\checkmarke\checkmark$\checkmark^\checkmark\\checkmarkc\checkmarki\checkmarkr\checkmarkc\checkmark$\checkmark \checkmark$\checkmark^\checkmark\\checkmarkc\checkmarki\checkmarkr\checkmarkc\checkmark$\checkmarkv\checkmark$\checkmark^\checkmark\\checkmarkc\checkmarki\checkmarkr\checkmarkc\checkmark$\checkmarké\checkmark$\checkmark^\checkmark\\checkmarkc\checkmarki\checkmarkr\checkmarkc\checkmark$\checkmarkr\checkmark$\checkmark^\checkmark\\checkmarkc\checkmarki\checkmarkr\checkmarkc\checkmark$\checkmarki\checkmark$\checkmark^\checkmark\\checkmarkc\checkmarki\checkmarkr\checkmarkc\checkmark$\checkmarkf\checkmark$\checkmark^\checkmark\\checkmarkc\checkmarki\checkmarkr\checkmarkc\checkmark$\checkmarki\checkmark$\checkmark^\checkmark\\checkmarkc\checkmarki\checkmarkr\checkmarkc\checkmark$\checkmarké\checkmark$\checkmark^\checkmark\\checkmarkc\checkmarki\checkmarkr\checkmarkc\checkmark$\checkmark}\checkmark$\checkmark^\checkmark\\checkmarkc\checkmarki\checkmarkr\checkmarkc\checkmark$\checkmark \checkmark$\checkmark^\checkmark\\checkmarkc\checkmarki\checkmarkr\checkmarkc\checkmark$\checkmark&\checkmark$\checkmark^\checkmark\\checkmarkc\checkmarki\checkmarkr\checkmarkc\checkmark$\checkmark \checkmark$\checkmark^\checkmark\\checkmarkc\checkmarki\checkmarkr\checkmarkc\checkmark$\checkmark\\checkmark$\checkmark^\checkmark\\checkmarkc\checkmarki\checkmarkr\checkmarkc\checkmark$\checkmarkt\checkmark$\checkmark^\checkmark\\checkmarkc\checkmarki\checkmarkr\checkmarkc\checkmark$\checkmarke\checkmark$\checkmark^\checkmark\\checkmarkc\checkmarki\checkmarkr\checkmarkc\checkmark$\checkmarkx\checkmark$\checkmark^\checkmark\\checkmarkc\checkmarki\checkmarkr\checkmarkc\checkmark$\checkmarkt\checkmark$\checkmark^\checkmark\\checkmarkc\checkmarki\checkmarkr\checkmarkc\checkmark$\checkmarkb\checkmark$\checkmark^\checkmark\\checkmarkc\checkmarki\checkmarkr\checkmarkc\checkmark$\checkmarkf\checkmark$\checkmark^\checkmark\\checkmarkc\checkmarki\checkmarkr\checkmarkc\checkmark$\checkmark{\checkmark$\checkmark^\checkmark\\checkmarkc\checkmarki\checkmarkr\checkmarkc\checkmark$\checkmarkV\checkmark$\checkmark^\checkmark\\checkmarkc\checkmarki\checkmarkr\checkmarkc\checkmark$\checkmarka\checkmark$\checkmark^\checkmark\\checkmarkc\checkmarki\checkmarkr\checkmarkc\checkmark$\checkmarkl\checkmark$\checkmark^\checkmark\\checkmarkc\checkmarki\checkmarkr\checkmarkc\checkmark$\checkmarke\checkmark$\checkmark^\checkmark\\checkmarkc\checkmarki\checkmarkr\checkmarkc\checkmark$\checkmarku\checkmark$\checkmark^\checkmark\\checkmarkc\checkmarki\checkmarkr\checkmarkc\checkmark$\checkmarkr\checkmark$\checkmark^\checkmark\\checkmarkc\checkmarki\checkmarkr\checkmarkc\checkmark$\checkmark \checkmark$\checkmark^\checkmark\\checkmarkc\checkmarki\checkmarkr\checkmarkc\checkmark$\checkmarkc\checkmark$\checkmark^\checkmark\\checkmarkc\checkmarki\checkmarkr\checkmarkc\checkmark$\checkmarka\checkmark$\checkmark^\checkmark\\checkmarkc\checkmarki\checkmarkr\checkmarkc\checkmark$\checkmarkl\checkmark$\checkmark^\checkmark\\checkmarkc\checkmarki\checkmarkr\checkmarkc\checkmark$\checkmarkc\checkmark$\checkmark^\checkmark\\checkmarkc\checkmarki\checkmarkr\checkmarkc\checkmark$\checkmarku\checkmark$\checkmark^\checkmark\\checkmarkc\checkmarki\checkmarkr\checkmarkc\checkmark$\checkmarkl\checkmark$\checkmark^\checkmark\\checkmarkc\checkmarki\checkmarkr\checkmarkc\checkmark$\checkmarké\checkmark$\checkmark^\checkmark\\checkmarkc\checkmarki\checkmarkr\checkmarkc\checkmark$\checkmarke\checkmark$\checkmark^\checkmark\\checkmarkc\checkmarki\checkmarkr\checkmarkc\checkmark$\checkmark}\checkmark$\checkmark^\checkmark\\checkmarkc\checkmarki\checkmarkr\checkmarkc\checkmark$\checkmark \checkmark$\checkmark^\checkmark\\checkmarkc\checkmarki\checkmarkr\checkmarkc\checkmark$\checkmark&\checkmark$\checkmark^\checkmark\\checkmarkc\checkmarki\checkmarkr\checkmarkc\checkmark$\checkmark \checkmark$\checkmark^\checkmark\\checkmarkc\checkmarki\checkmarkr\checkmarkc\checkmark$\checkmark\\checkmark$\checkmark^\checkmark\\checkmarkc\checkmarki\checkmarkr\checkmarkc\checkmark$\checkmarkt\checkmark$\checkmark^\checkmark\\checkmarkc\checkmarki\checkmarkr\checkmarkc\checkmark$\checkmarke\checkmark$\checkmark^\checkmark\\checkmarkc\checkmarki\checkmarkr\checkmarkc\checkmark$\checkmarkx\checkmark$\checkmark^\checkmark\\checkmarkc\checkmarki\checkmarkr\checkmarkc\checkmark$\checkmarkt\checkmark$\checkmark^\checkmark\\checkmarkc\checkmarki\checkmarkr\checkmarkc\checkmark$\checkmarkb\checkmark$\checkmark^\checkmark\\checkmarkc\checkmarki\checkmarkr\checkmarkc\checkmark$\checkmarkf\checkmark$\checkmark^\checkmark\\checkmarkc\checkmarki\checkmarkr\checkmarkc\checkmark$\checkmark{\checkmark$\checkmark^\checkmark\\checkmarkc\checkmarki\checkmarkr\checkmarkc\checkmark$\checkmarkL\checkmark$\checkmark^\checkmark\\checkmarkc\checkmarki\checkmarkr\checkmarkc\checkmark$\checkmarki\checkmark$\checkmark^\checkmark\\checkmarkc\checkmarki\checkmarkr\checkmarkc\checkmark$\checkmarkm\checkmark$\checkmark^\checkmark\\checkmarkc\checkmarki\checkmarkr\checkmarkc\checkmark$\checkmarki\checkmark$\checkmark^\checkmark\\checkmarkc\checkmarki\checkmarkr\checkmarkc\checkmark$\checkmarkt\checkmark$\checkmark^\checkmark\\checkmarkc\checkmarki\checkmarkr\checkmarkc\checkmark$\checkmarke\checkmark$\checkmark^\checkmark\\checkmarkc\checkmarki\checkmarkr\checkmarkc\checkmark$\checkmark \checkmark$\checkmark^\checkmark\\checkmarkc\checkmarki\checkmarkr\checkmarkc\checkmark$\checkmarka\checkmark$\checkmark^\checkmark\\checkmarkc\checkmarki\checkmarkr\checkmarkc\checkmark$\checkmarkd\checkmark$\checkmark^\checkmark\\checkmarkc\checkmarki\checkmarkr\checkmarkc\checkmark$\checkmarkm\checkmark$\checkmark^\checkmark\\checkmarkc\checkmarki\checkmarkr\checkmarkc\checkmark$\checkmarki\checkmark$\checkmark^\checkmark\\checkmarkc\checkmarki\checkmarkr\checkmarkc\checkmark$\checkmarks\checkmark$\checkmark^\checkmark\\checkmarkc\checkmarki\checkmarkr\checkmarkc\checkmark$\checkmarks\checkmark$\checkmark^\checkmark\\checkmarkc\checkmarki\checkmarkr\checkmarkc\checkmark$\checkmarki\checkmark$\checkmark^\checkmark\\checkmarkc\checkmarki\checkmarkr\checkmarkc\checkmark$\checkmarkb\checkmark$\checkmark^\checkmark\\checkmarkc\checkmarki\checkmarkr\checkmarkc\checkmark$\checkmarkl\checkmark$\checkmark^\checkmark\\checkmarkc\checkmarki\checkmarkr\checkmarkc\checkmark$\checkmarke\checkmark$\checkmark^\checkmark\\checkmarkc\checkmarki\checkmarkr\checkmarkc\checkmark$\checkmark}\checkmark$\checkmark^\checkmark\\checkmarkc\checkmarki\checkmarkr\checkmarkc\checkmark$\checkmark \checkmark$\checkmark^\checkmark\\checkmarkc\checkmarki\checkmarkr\checkmarkc\checkmark$\checkmark&\checkmark$\checkmark^\checkmark\\checkmarkc\checkmarki\checkmarkr\checkmarkc\checkmark$\checkmark \checkmark$\checkmark^\checkmark\\checkmarkc\checkmarki\checkmarkr\checkmarkc\checkmark$\checkmark\\checkmark$\checkmark^\checkmark\\checkmarkc\checkmarki\checkmarkr\checkmarkc\checkmark$\checkmarkt\checkmark$\checkmark^\checkmark\\checkmarkc\checkmarki\checkmarkr\checkmarkc\checkmark$\checkmarke\checkmark$\checkmark^\checkmark\\checkmarkc\checkmarki\checkmarkr\checkmarkc\checkmark$\checkmarkx\checkmark$\checkmark^\checkmark\\checkmarkc\checkmarki\checkmarkr\checkmarkc\checkmark$\checkmarkt\checkmark$\checkmark^\checkmark\\checkmarkc\checkmarki\checkmarkr\checkmarkc\checkmark$\checkmarkb\checkmark$\checkmark^\checkmark\\checkmarkc\checkmarki\checkmarkr\checkmarkc\checkmark$\checkmarkf\checkmark$\checkmark^\checkmark\\checkmarkc\checkmarki\checkmarkr\checkmarkc\checkmark$\checkmark{\checkmark$\checkmark^\checkmark\\checkmarkc\checkmarki\checkmarkr\checkmarkc\checkmark$\checkmark$\checkmark$\checkmark^\checkmark\\checkmarkc\checkmarki\checkmarkr\checkmarkc\checkmark$\checkmarks\checkmark$\checkmark^\checkmark\\checkmarkc\checkmarki\checkmarkr\checkmarkc\checkmark$\checkmark$\checkmark$\checkmark^\checkmark\\checkmarkc\checkmarki\checkmarkr\checkmarkc\checkmark$\checkmark \checkmark$\checkmark^\checkmark\\checkmarkc\checkmarki\checkmarkr\checkmarkc\checkmark$\checkmarkr\checkmark$\checkmark^\checkmark\\checkmarkc\checkmarki\checkmarkr\checkmarkc\checkmark$\checkmarké\checkmark$\checkmark^\checkmark\\checkmarkc\checkmarki\checkmarkr\checkmarkc\checkmark$\checkmarke\checkmark$\checkmark^\checkmark\\checkmarkc\checkmarki\checkmarkr\checkmarkc\checkmark$\checkmarkl\checkmark$\checkmark^\checkmark\\checkmarkc\checkmarki\checkmarkr\checkmarkc\checkmark$\checkmark}\checkmark$\checkmark^\checkmark\\checkmarkc\checkmarki\checkmarkr\checkmarkc\checkmark$\checkmark \checkmark$\checkmark^\checkmark\\checkmarkc\checkmarki\checkmarkr\checkmarkc\checkmark$\checkmark\\checkmark$\checkmark^\checkmark\\checkmarkc\checkmarki\checkmarkr\checkmarkc\checkmark$\checkmark\\checkmark$\checkmark^\checkmark\\checkmarkc\checkmarki\checkmarkr\checkmarkc\checkmark$\checkmark \checkmark$\checkmark^\checkmark\\checkmarkc\checkmarki\checkmarkr\checkmarkc\checkmark$\checkmark\\checkmark$\checkmark^\checkmark\\checkmarkc\checkmarki\checkmarkr\checkmarkc\checkmark$\checkmarkh\checkmark$\checkmark^\checkmark\\checkmarkc\checkmarki\checkmarkr\checkmarkc\checkmark$\checkmarkl\checkmark$\checkmark^\checkmark\\checkmarkc\checkmarki\checkmarkr\checkmarkc\checkmark$\checkmarki\checkmark$\checkmark^\checkmark\\checkmarkc\checkmarki\checkmarkr\checkmarkc\checkmark$\checkmarkn\checkmark$\checkmark^\checkmark\\checkmarkc\checkmarki\checkmarkr\checkmarkc\checkmark$\checkmarke\checkmark$\checkmark^\checkmark\\checkmarkc\checkmarki\checkmarkr\checkmarkc\checkmark$\checkmark
\checkmark$\checkmark^\checkmark\\checkmarkc\checkmarki\checkmarkr\checkmarkc\checkmark$\checkmarkV\checkmark$\checkmark^\checkmark\\checkmarkc\checkmarki\checkmarkr\checkmarkc\checkmark$\checkmarki\checkmark$\checkmark^\checkmark\\checkmarkc\checkmarki\checkmarkr\checkmarkc\checkmark$\checkmarks\checkmark$\checkmark^\checkmark\\checkmarkc\checkmarki\checkmarkr\checkmarkc\checkmark$\checkmark \checkmark$\checkmark^\checkmark\\checkmarkc\checkmarki\checkmarkr\checkmarkc\checkmark$\checkmarkà\checkmark$\checkmark^\checkmark\\checkmarkc\checkmarki\checkmarkr\checkmarkc\checkmark$\checkmark \checkmark$\checkmark^\checkmark\\checkmarkc\checkmarki\checkmarkr\checkmarkc\checkmark$\checkmarkb\checkmark$\checkmark^\checkmark\\checkmarkc\checkmarki\checkmarkr\checkmarkc\checkmark$\checkmarki\checkmark$\checkmark^\checkmark\\checkmarkc\checkmarki\checkmarkr\checkmarkc\checkmark$\checkmarkl\checkmark$\checkmark^\checkmark\\checkmarkc\checkmarki\checkmarkr\checkmarkc\checkmark$\checkmarkl\checkmark$\checkmark^\checkmark\\checkmarkc\checkmarki\checkmarkr\checkmarkc\checkmark$\checkmarke\checkmark$\checkmark^\checkmark\\checkmarkc\checkmarki\checkmarkr\checkmarkc\checkmark$\checkmarks\checkmark$\checkmark^\checkmark\\checkmarkc\checkmarki\checkmarkr\checkmarkc\checkmark$\checkmark \checkmark$\checkmark^\checkmark\\checkmarkc\checkmarki\checkmarkr\checkmarkc\checkmark$\checkmarkS\checkmark$\checkmark^\checkmark\\checkmarkc\checkmarki\checkmarkr\checkmarkc\checkmark$\checkmarkF\checkmark$\checkmark^\checkmark\\checkmarkc\checkmarki\checkmarkr\checkmarkc\checkmark$\checkmarkU\checkmark$\checkmark^\checkmark\\checkmarkc\checkmarki\checkmarkr\checkmarkc\checkmark$\checkmark1\checkmark$\checkmark^\checkmark\\checkmarkc\checkmarki\checkmarkr\checkmarkc\checkmark$\checkmark6\checkmark$\checkmark^\checkmark\\checkmarkc\checkmarki\checkmarkr\checkmarkc\checkmark$\checkmark0\checkmark$\checkmark^\checkmark\\checkmarkc\checkmarki\checkmarkr\checkmarkc\checkmark$\checkmark5\checkmark$\checkmark^\checkmark\\checkmarkc\checkmarki\checkmarkr\checkmarkc\checkmark$\checkmark \checkmark$\checkmark^\checkmark\\checkmarkc\checkmarki\checkmarkr\checkmarkc\checkmark$\checkmark&\checkmark$\checkmark^\checkmark\\checkmarkc\checkmarki\checkmarkr\checkmarkc\checkmark$\checkmark \checkmark$\checkmark^\checkmark\\checkmarkc\checkmarki\checkmarkr\checkmarkc\checkmark$\checkmarkF\checkmark$\checkmark^\checkmark\\checkmarkc\checkmarki\checkmarkr\checkmarkc\checkmark$\checkmarkl\checkmark$\checkmark^\checkmark\\checkmarkc\checkmarki\checkmarkr\checkmarkc\checkmark$\checkmarka\checkmark$\checkmark^\checkmark\\checkmarkc\checkmarki\checkmarkr\checkmarkc\checkmark$\checkmarkm\checkmark$\checkmark^\checkmark\\checkmarkc\checkmarki\checkmarkr\checkmarkc\checkmark$\checkmarkb\checkmark$\checkmark^\checkmark\\checkmarkc\checkmarki\checkmarkr\checkmarkc\checkmark$\checkmarka\checkmark$\checkmark^\checkmark\\checkmarkc\checkmarki\checkmarkr\checkmarkc\checkmark$\checkmarkg\checkmark$\checkmark^\checkmark\\checkmarkc\checkmarki\checkmarkr\checkmarkc\checkmark$\checkmarke\checkmark$\checkmark^\checkmark\\checkmarkc\checkmarki\checkmarkr\checkmarkc\checkmark$\checkmark \checkmark$\checkmark^\checkmark\\checkmarkc\checkmarki\checkmarkr\checkmarkc\checkmark$\checkmark(\checkmark$\checkmark^\checkmark\\checkmarkc\checkmarki\checkmarkr\checkmarkc\checkmark$\checkmarkE\checkmark$\checkmark^\checkmark\\checkmarkc\checkmarki\checkmarkr\checkmarkc\checkmark$\checkmarku\checkmark$\checkmark^\checkmark\\checkmarkc\checkmarki\checkmarkr\checkmarkc\checkmark$\checkmarkl\checkmark$\checkmark^\checkmark\\checkmarkc\checkmarki\checkmarkr\checkmarkc\checkmark$\checkmarke\checkmark$\checkmark^\checkmark\\checkmarkc\checkmarki\checkmarkr\checkmarkc\checkmark$\checkmarkr\checkmark$\checkmark^\checkmark\\checkmarkc\checkmarki\checkmarkr\checkmarkc\checkmark$\checkmark)\checkmark$\checkmark^\checkmark\\checkmarkc\checkmarki\checkmarkr\checkmarkc\checkmark$\checkmark \checkmark$\checkmark^\checkmark\\checkmarkc\checkmarki\checkmarkr\checkmarkc\checkmark$\checkmark&\checkmark$\checkmark^\checkmark\\checkmarkc\checkmarki\checkmarkr\checkmarkc\checkmark$\checkmark \checkmark$\checkmark^\checkmark\\checkmarkc\checkmarki\checkmarkr\checkmarkc\checkmark$\checkmark$\checkmark$\checkmark^\checkmark\\checkmarkc\checkmarki\checkmarkr\checkmarkc\checkmark$\checkmarkF\checkmark$\checkmark^\checkmark\\checkmarkc\checkmarki\checkmarkr\checkmarkc\checkmark$\checkmark_\checkmark$\checkmark^\checkmark\\checkmarkc\checkmarki\checkmarkr\checkmarkc\checkmark$\checkmark{\checkmark$\checkmark^\checkmark\\checkmarkc\checkmarki\checkmarkr\checkmarkc\checkmark$\checkmarka\checkmark$\checkmark^\checkmark\\checkmarkc\checkmarki\checkmarkr\checkmarkc\checkmark$\checkmarkx\checkmark$\checkmark^\checkmark\\checkmarkc\checkmarki\checkmarkr\checkmarkc\checkmark$\checkmark}\checkmark$\checkmark^\checkmark\\checkmarkc\checkmarki\checkmarkr\checkmarkc\checkmark$\checkmark \checkmark$\checkmark^\checkmark\\checkmarkc\checkmarki\checkmarkr\checkmarkc\checkmark$\checkmark=\checkmark$\checkmark^\checkmark\\checkmarkc\checkmarki\checkmarkr\checkmarkc\checkmark$\checkmark \checkmark$\checkmark^\checkmark\\checkmarkc\checkmarki\checkmarkr\checkmarkc\checkmark$\checkmark1\checkmark$\checkmark^\checkmark\\checkmarkc\checkmarki\checkmarkr\checkmarkc\checkmark$\checkmark0\checkmark$\checkmark^\checkmark\\checkmarkc\checkmarki\checkmarkr\checkmarkc\checkmark$\checkmark5\checkmark$\checkmark^\checkmark\\checkmarkc\checkmarki\checkmarkr\checkmarkc\checkmark$\checkmark{\checkmark$\checkmark^\checkmark\\checkmarkc\checkmarki\checkmarkr\checkmarkc\checkmark$\checkmark,\checkmark$\checkmark^\checkmark\\checkmarkc\checkmarki\checkmarkr\checkmarkc\checkmark$\checkmark}\checkmark$\checkmark^\checkmark\\checkmarkc\checkmarki\checkmarkr\checkmarkc\checkmark$\checkmark5\checkmark$\checkmark^\checkmark\\checkmarkc\checkmarki\checkmarkr\checkmarkc\checkmark$\checkmark$\checkmark$\checkmark^\checkmark\\checkmarkc\checkmarki\checkmarkr\checkmarkc\checkmark$\checkmark \checkmark$\checkmark^\checkmark\\checkmarkc\checkmarki\checkmarkr\checkmarkc\checkmark$\checkmarkN\checkmark$\checkmark^\checkmark\\checkmarkc\checkmarki\checkmarkr\checkmarkc\checkmark$\checkmark \checkmark$\checkmark^\checkmark\\checkmarkc\checkmarki\checkmarkr\checkmarkc\checkmark$\checkmark&\checkmark$\checkmark^\checkmark\\checkmarkc\checkmarki\checkmarkr\checkmarkc\checkmark$\checkmark \checkmark$\checkmark^\checkmark\\checkmarkc\checkmarki\checkmarkr\checkmarkc\checkmark$\checkmark$\checkmark$\checkmark^\checkmark\\checkmarkc\checkmarki\checkmarkr\checkmarkc\checkmark$\checkmarkF\checkmark$\checkmark^\checkmark\\checkmarkc\checkmarki\checkmarkr\checkmarkc\checkmark$\checkmark_\checkmark$\checkmark^\checkmark\\checkmarkc\checkmarki\checkmarkr\checkmarkc\checkmark$\checkmarkk\checkmark$\checkmark^\checkmark\\checkmarkc\checkmarki\checkmarkr\checkmarkc\checkmark$\checkmark \checkmark$\checkmark^\checkmark\\checkmarkc\checkmarki\checkmarkr\checkmarkc\checkmark$\checkmark=\checkmark$\checkmark^\checkmark\\checkmarkc\checkmarki\checkmarkr\checkmarkc\checkmark$\checkmark \checkmark$\checkmark^\checkmark\\checkmarkc\checkmarki\checkmarkr\checkmarkc\checkmark$\checkmark2\checkmark$\checkmark^\checkmark\\checkmarkc\checkmarki\checkmarkr\checkmarkc\checkmark$\checkmark6\checkmark$\checkmark^\checkmark\\checkmarkc\checkmarki\checkmarkr\checkmarkc\checkmark$\checkmark\\checkmark$\checkmark^\checkmark\\checkmarkc\checkmarki\checkmarkr\checkmarkc\checkmark$\checkmark,\checkmark$\checkmark^\checkmark\\checkmarkc\checkmarki\checkmarkr\checkmarkc\checkmark$\checkmark3\checkmark$\checkmark^\checkmark\\checkmarkc\checkmarki\checkmarkr\checkmarkc\checkmark$\checkmark9\checkmark$\checkmark^\checkmark\\checkmarkc\checkmarki\checkmarkr\checkmarkc\checkmark$\checkmark5\checkmark$\checkmark^\checkmark\\checkmarkc\checkmarki\checkmarkr\checkmarkc\checkmark$\checkmark$\checkmark$\checkmark^\checkmark\\checkmarkc\checkmarki\checkmarkr\checkmarkc\checkmark$\checkmark \checkmark$\checkmark^\checkmark\\checkmarkc\checkmarki\checkmarkr\checkmarkc\checkmark$\checkmarkN\checkmark$\checkmark^\checkmark\\checkmarkc\checkmarki\checkmarkr\checkmarkc\checkmark$\checkmark \checkmark$\checkmark^\checkmark\\checkmarkc\checkmarki\checkmarkr\checkmarkc\checkmark$\checkmark&\checkmark$\checkmark^\checkmark\\checkmarkc\checkmarki\checkmarkr\checkmarkc\checkmark$\checkmark \checkmark$\checkmark^\checkmark\\checkmarkc\checkmarki\checkmarkr\checkmarkc\checkmark$\checkmark$\checkmark$\checkmark^\checkmark\\checkmarkc\checkmarki\checkmarkr\checkmarkc\checkmark$\checkmark2\checkmark$\checkmark^\checkmark\\checkmarkc\checkmarki\checkmarkr\checkmarkc\checkmark$\checkmark5\checkmark$\checkmark^\checkmark\\checkmarkc\checkmarki\checkmarkr\checkmarkc\checkmark$\checkmark0\checkmark$\checkmark^\checkmark\\checkmarkc\checkmarki\checkmarkr\checkmarkc\checkmark$\checkmark$\checkmark$\checkmark^\checkmark\\checkmarkc\checkmarki\checkmarkr\checkmarkc\checkmark$\checkmark \checkmark$\checkmark^\checkmark\\checkmarkc\checkmarki\checkmarkr\checkmarkc\checkmark$\checkmark\\checkmark$\checkmark^\checkmark\\checkmarkc\checkmarki\checkmarkr\checkmarkc\checkmark$\checkmark\\checkmark$\checkmark^\checkmark\\checkmarkc\checkmarki\checkmarkr\checkmarkc\checkmark$\checkmark \checkmark$\checkmark^\checkmark\\checkmarkc\checkmarki\checkmarkr\checkmarkc\checkmark$\checkmark\\checkmark$\checkmark^\checkmark\\checkmarkc\checkmarki\checkmarkr\checkmarkc\checkmark$\checkmarkh\checkmark$\checkmark^\checkmark\\checkmarkc\checkmarki\checkmarkr\checkmarkc\checkmark$\checkmarkl\checkmark$\checkmark^\checkmark\\checkmarkc\checkmarki\checkmarkr\checkmarkc\checkmark$\checkmarki\checkmark$\checkmark^\checkmark\\checkmarkc\checkmarki\checkmarkr\checkmarkc\checkmark$\checkmarkn\checkmark$\checkmark^\checkmark\\checkmarkc\checkmarki\checkmarkr\checkmarkc\checkmark$\checkmarke\checkmark$\checkmark^\checkmark\\checkmarkc\checkmarki\checkmarkr\checkmarkc\checkmark$\checkmark
\checkmark$\checkmark^\checkmark\\checkmarkc\checkmarki\checkmarkr\checkmarkc\checkmark$\checkmarkM\checkmark$\checkmark^\checkmark\\checkmarkc\checkmarki\checkmarkr\checkmarkc\checkmark$\checkmarko\checkmark$\checkmark^\checkmark\\checkmarkc\checkmarki\checkmarkr\checkmarkc\checkmark$\checkmarkt\checkmark$\checkmark^\checkmark\\checkmarkc\checkmarki\checkmarkr\checkmarkc\checkmark$\checkmarke\checkmark$\checkmark^\checkmark\\checkmarkc\checkmarki\checkmarkr\checkmarkc\checkmark$\checkmarku\checkmark$\checkmark^\checkmark\\checkmarkc\checkmarki\checkmarkr\checkmarkc\checkmark$\checkmarkr\checkmark$\checkmark^\checkmark\\checkmarkc\checkmarki\checkmarkr\checkmarkc\checkmark$\checkmark \checkmark$\checkmark^\checkmark\\checkmarkc\checkmarki\checkmarkr\checkmarkc\checkmark$\checkmarkN\checkmark$\checkmark^\checkmark\\checkmarkc\checkmarki\checkmarkr\checkmarkc\checkmark$\checkmarkE\checkmark$\checkmark^\checkmark\\checkmarkc\checkmarki\checkmarkr\checkmarkc\checkmark$\checkmarkM\checkmark$\checkmark^\checkmark\\checkmarkc\checkmarki\checkmarkr\checkmarkc\checkmark$\checkmarkA\checkmark$\checkmark^\checkmark\\checkmarkc\checkmarki\checkmarkr\checkmarkc\checkmark$\checkmark \checkmark$\checkmark^\checkmark\\checkmarkc\checkmarki\checkmarkr\checkmarkc\checkmark$\checkmark1\checkmark$\checkmark^\checkmark\\checkmarkc\checkmarki\checkmarkr\checkmarkc\checkmark$\checkmark7\checkmark$\checkmark^\checkmark\\checkmarkc\checkmarki\checkmarkr\checkmarkc\checkmark$\checkmark \checkmark$\checkmark^\checkmark\\checkmarkc\checkmarki\checkmarkr\checkmarkc\checkmark$\checkmark(\checkmark$\checkmark^\checkmark\\checkmarkc\checkmarki\checkmarkr\checkmarkc\checkmark$\checkmarkX\checkmark$\checkmark^\checkmark\\checkmarkc\checkmarki\checkmarkr\checkmarkc\checkmark$\checkmark,\checkmark$\checkmark^\checkmark\\checkmarkc\checkmarki\checkmarkr\checkmarkc\checkmark$\checkmarkY\checkmark$\checkmark^\checkmark\\checkmarkc\checkmarki\checkmarkr\checkmarkc\checkmark$\checkmark,\checkmark$\checkmark^\checkmark\\checkmarkc\checkmarki\checkmarkr\checkmarkc\checkmark$\checkmarkZ\checkmark$\checkmark^\checkmark\\checkmarkc\checkmarki\checkmarkr\checkmarkc\checkmark$\checkmark)\checkmark$\checkmark^\checkmark\\checkmarkc\checkmarki\checkmarkr\checkmarkc\checkmark$\checkmark \checkmark$\checkmark^\checkmark\\checkmarkc\checkmarki\checkmarkr\checkmarkc\checkmark$\checkmark&\checkmark$\checkmark^\checkmark\\checkmarkc\checkmarki\checkmarkr\checkmarkc\checkmark$\checkmark \checkmark$\checkmark^\checkmark\\checkmarkc\checkmarki\checkmarkr\checkmarkc\checkmark$\checkmarkC\checkmark$\checkmark^\checkmark\\checkmarkc\checkmarki\checkmarkr\checkmarkc\checkmark$\checkmarko\checkmark$\checkmark^\checkmark\\checkmarkc\checkmarki\checkmarkr\checkmarkc\checkmark$\checkmarku\checkmark$\checkmark^\checkmark\\checkmarkc\checkmarki\checkmarkr\checkmarkc\checkmark$\checkmarkp\checkmark$\checkmark^\checkmark\\checkmarkc\checkmarki\checkmarkr\checkmarkc\checkmark$\checkmarkl\checkmark$\checkmark^\checkmark\\checkmarkc\checkmarki\checkmarkr\checkmarkc\checkmark$\checkmarke\checkmark$\checkmark^\checkmark\\checkmarkc\checkmarki\checkmarkr\checkmarkc\checkmark$\checkmark \checkmark$\checkmark^\checkmark\\checkmarkc\checkmarki\checkmarkr\checkmarkc\checkmark$\checkmarkd\checkmark$\checkmark^\checkmark\\checkmarkc\checkmarki\checkmarkr\checkmarkc\checkmark$\checkmarke\checkmark$\checkmark^\checkmark\\checkmarkc\checkmarki\checkmarkr\checkmarkc\checkmark$\checkmark \checkmark$\checkmark^\checkmark\\checkmarkc\checkmarki\checkmarkr\checkmarkc\checkmark$\checkmarkm\checkmark$\checkmark^\checkmark\\checkmarkc\checkmarki\checkmarkr\checkmarkc\checkmark$\checkmarka\checkmark$\checkmark^\checkmark\\checkmarkc\checkmarki\checkmarkr\checkmarkc\checkmark$\checkmarki\checkmark$\checkmark^\checkmark\\checkmarkc\checkmarki\checkmarkr\checkmarkc\checkmark$\checkmarkn\checkmark$\checkmark^\checkmark\\checkmarkc\checkmarki\checkmarkr\checkmarkc\checkmark$\checkmarkt\checkmark$\checkmark^\checkmark\\checkmarkc\checkmarki\checkmarkr\checkmarkc\checkmark$\checkmarki\checkmark$\checkmark^\checkmark\\checkmarkc\checkmarki\checkmarkr\checkmarkc\checkmark$\checkmarke\checkmark$\checkmark^\checkmark\\checkmarkc\checkmarki\checkmarkr\checkmarkc\checkmark$\checkmarkn\checkmark$\checkmark^\checkmark\\checkmarkc\checkmarki\checkmarkr\checkmarkc\checkmark$\checkmark \checkmark$\checkmark^\checkmark\\checkmarkc\checkmarki\checkmarkr\checkmarkc\checkmark$\checkmark&\checkmark$\checkmark^\checkmark\\checkmarkc\checkmarki\checkmarkr\checkmarkc\checkmark$\checkmark \checkmark$\checkmark^\checkmark\\checkmarkc\checkmarki\checkmarkr\checkmarkc\checkmark$\checkmark$\checkmark$\checkmark^\checkmark\\checkmarkc\checkmarki\checkmarkr\checkmarkc\checkmark$\checkmarkC\checkmark$\checkmark^\checkmark\\checkmarkc\checkmarki\checkmarkr\checkmarkc\checkmark$\checkmark_\checkmark$\checkmark^\checkmark\\checkmarkc\checkmarki\checkmarkr\checkmarkc\checkmark$\checkmark{\checkmark$\checkmark^\checkmark\\checkmarkc\checkmarki\checkmarkr\checkmarkc\checkmark$\checkmarkr\checkmark$\checkmark^\checkmark\\checkmarkc\checkmarki\checkmarkr\checkmarkc\checkmark$\checkmarke\checkmark$\checkmark^\checkmark\\checkmarkc\checkmarki\checkmarkr\checkmarkc\checkmark$\checkmarkq\checkmark$\checkmark^\checkmark\\checkmarkc\checkmarki\checkmarkr\checkmarkc\checkmark$\checkmark}\checkmark$\checkmark^\checkmark\\checkmarkc\checkmarki\checkmarkr\checkmarkc\checkmark$\checkmark \checkmark$\checkmark^\checkmark\\checkmarkc\checkmarki\checkmarkr\checkmarkc\checkmark$\checkmark=\checkmark$\checkmark^\checkmark\\checkmarkc\checkmarki\checkmarkr\checkmarkc\checkmark$\checkmark \checkmark$\checkmark^\checkmark\\checkmarkc\checkmarki\checkmarkr\checkmarkc\checkmark$\checkmark0\checkmark$\checkmark^\checkmark\\checkmarkc\checkmarki\checkmarkr\checkmarkc\checkmark$\checkmark{\checkmark$\checkmark^\checkmark\\checkmarkc\checkmarki\checkmarkr\checkmarkc\checkmark$\checkmark,\checkmark$\checkmark^\checkmark\\checkmarkc\checkmarki\checkmarkr\checkmarkc\checkmark$\checkmark}\checkmark$\checkmark^\checkmark\\checkmarkc\checkmarki\checkmarkr\checkmarkc\checkmark$\checkmark3\checkmark$\checkmark^\checkmark\\checkmarkc\checkmarki\checkmarkr\checkmarkc\checkmark$\checkmark3\checkmark$\checkmark^\checkmark\\checkmarkc\checkmarki\checkmarkr\checkmarkc\checkmark$\checkmark3\checkmark$\checkmark^\checkmark\\checkmarkc\checkmarki\checkmarkr\checkmarkc\checkmark$\checkmark$\checkmark$\checkmark^\checkmark\\checkmarkc\checkmarki\checkmarkr\checkmarkc\checkmark$\checkmark \checkmark$\checkmark^\checkmark\\checkmarkc\checkmarki\checkmarkr\checkmarkc\checkmark$\checkmarkN\checkmark$\checkmark^\checkmark\\checkmarkc\checkmarki\checkmarkr\checkmarkc\checkmark$\checkmark$\checkmark$\checkmark^\checkmark\\checkmarkc\checkmarki\checkmarkr\checkmarkc\checkmark$\checkmark\\checkmark$\checkmark^\checkmark\\checkmarkc\checkmarki\checkmarkr\checkmarkc\checkmark$\checkmarkc\checkmark$\checkmark^\checkmark\\checkmarkc\checkmarki\checkmarkr\checkmarkc\checkmark$\checkmarkd\checkmark$\checkmark^\checkmark\\checkmarkc\checkmarki\checkmarkr\checkmarkc\checkmark$\checkmarko\checkmark$\checkmark^\checkmark\\checkmarkc\checkmarki\checkmarkr\checkmarkc\checkmark$\checkmarkt\checkmark$\checkmark^\checkmark\\checkmarkc\checkmarki\checkmarkr\checkmarkc\checkmark$\checkmark$\checkmark$\checkmark^\checkmark\\checkmarkc\checkmarki\checkmarkr\checkmarkc\checkmark$\checkmarkm\checkmark$\checkmark^\checkmark\\checkmarkc\checkmarki\checkmarkr\checkmarkc\checkmark$\checkmark \checkmark$\checkmark^\checkmark\\checkmarkc\checkmarki\checkmarkr\checkmarkc\checkmark$\checkmark&\checkmark$\checkmark^\checkmark\\checkmarkc\checkmarki\checkmarkr\checkmarkc\checkmark$\checkmark \checkmark$\checkmark^\checkmark\\checkmarkc\checkmarki\checkmarkr\checkmarkc\checkmark$\checkmark$\checkmark$\checkmark^\checkmark\\checkmarkc\checkmarki\checkmarkr\checkmarkc\checkmark$\checkmarkC\checkmark$\checkmark^\checkmark\\checkmarkc\checkmarki\checkmarkr\checkmarkc\checkmark$\checkmark_\checkmark$\checkmark^\checkmark\\checkmarkc\checkmarki\checkmarkr\checkmarkc\checkmark$\checkmark{\checkmark$\checkmark^\checkmark\\checkmarkc\checkmarki\checkmarkr\checkmarkc\checkmark$\checkmarkm\checkmark$\checkmark^\checkmark\\checkmarkc\checkmarki\checkmarkr\checkmarkc\checkmark$\checkmarko\checkmark$\checkmark^\checkmark\\checkmarkc\checkmarki\checkmarkr\checkmarkc\checkmark$\checkmarkt\checkmark$\checkmark^\checkmark\\checkmarkc\checkmarki\checkmarkr\checkmarkc\checkmark$\checkmark}\checkmark$\checkmark^\checkmark\\checkmarkc\checkmarki\checkmarkr\checkmarkc\checkmark$\checkmark \checkmark$\checkmark^\checkmark\\checkmarkc\checkmarki\checkmarkr\checkmarkc\checkmark$\checkmark=\checkmark$\checkmark^\checkmark\\checkmarkc\checkmarki\checkmarkr\checkmarkc\checkmark$\checkmark \checkmark$\checkmark^\checkmark\\checkmarkc\checkmarki\checkmarkr\checkmarkc\checkmark$\checkmark0\checkmark$\checkmark^\checkmark\\checkmarkc\checkmarki\checkmarkr\checkmarkc\checkmark$\checkmark{\checkmark$\checkmark^\checkmark\\checkmarkc\checkmarki\checkmarkr\checkmarkc\checkmark$\checkmark,\checkmark$\checkmark^\checkmark\\checkmarkc\checkmarki\checkmarkr\checkmarkc\checkmark$\checkmark}\checkmark$\checkmark^\checkmark\\checkmarkc\checkmarki\checkmarkr\checkmarkc\checkmark$\checkmark4\checkmark$\checkmark^\checkmark\\checkmarkc\checkmarki\checkmarkr\checkmarkc\checkmark$\checkmark2\checkmark$\checkmark^\checkmark\\checkmarkc\checkmarki\checkmarkr\checkmarkc\checkmark$\checkmark$\checkmark$\checkmark^\checkmark\\checkmarkc\checkmarki\checkmarkr\checkmarkc\checkmark$\checkmark \checkmark$\checkmark^\checkmark\\checkmarkc\checkmarki\checkmarkr\checkmarkc\checkmark$\checkmarkN\checkmark$\checkmark^\checkmark\\checkmarkc\checkmarki\checkmarkr\checkmarkc\checkmark$\checkmark$\checkmark$\checkmark^\checkmark\\checkmarkc\checkmarki\checkmarkr\checkmarkc\checkmark$\checkmark\\checkmark$\checkmark^\checkmark\\checkmarkc\checkmarki\checkmarkr\checkmarkc\checkmark$\checkmarkc\checkmark$\checkmark^\checkmark\\checkmarkc\checkmarki\checkmarkr\checkmarkc\checkmark$\checkmarkd\checkmark$\checkmark^\checkmark\\checkmarkc\checkmarki\checkmarkr\checkmarkc\checkmark$\checkmarko\checkmark$\checkmark^\checkmark\\checkmarkc\checkmarki\checkmarkr\checkmarkc\checkmark$\checkmarkt\checkmark$\checkmark^\checkmark\\checkmarkc\checkmarki\checkmarkr\checkmarkc\checkmark$\checkmark$\checkmark$\checkmark^\checkmark\\checkmarkc\checkmarki\checkmarkr\checkmarkc\checkmark$\checkmarkm\checkmark$\checkmark^\checkmark\\checkmarkc\checkmarki\checkmarkr\checkmarkc\checkmark$\checkmark \checkmark$\checkmark^\checkmark\\checkmarkc\checkmarki\checkmarkr\checkmarkc\checkmark$\checkmark&\checkmark$\checkmark^\checkmark\\checkmarkc\checkmarki\checkmarkr\checkmarkc\checkmark$\checkmark \checkmark$\checkmark^\checkmark\\checkmarkc\checkmarki\checkmarkr\checkmarkc\checkmark$\checkmark$\checkmark$\checkmark^\checkmark\\checkmarkc\checkmarki\checkmarkr\checkmarkc\checkmark$\checkmark1\checkmark$\checkmark^\checkmark\\checkmarkc\checkmarki\checkmarkr\checkmarkc\checkmark$\checkmark{\checkmark$\checkmark^\checkmark\\checkmarkc\checkmarki\checkmarkr\checkmarkc\checkmark$\checkmark,\checkmark$\checkmark^\checkmark\\checkmarkc\checkmarki\checkmarkr\checkmarkc\checkmark$\checkmark}\checkmark$\checkmark^\checkmark\\checkmarkc\checkmarki\checkmarkr\checkmarkc\checkmark$\checkmark2\checkmark$\checkmark^\checkmark\\checkmarkc\checkmarki\checkmarkr\checkmarkc\checkmark$\checkmark6\checkmark$\checkmark^\checkmark\\checkmarkc\checkmarki\checkmarkr\checkmarkc\checkmark$\checkmark$\checkmark$\checkmark^\checkmark\\checkmarkc\checkmarki\checkmarkr\checkmarkc\checkmark$\checkmark \checkmark$\checkmark^\checkmark\\checkmarkc\checkmarki\checkmarkr\checkmarkc\checkmark$\checkmark\\checkmark$\checkmark^\checkmark\\checkmarkc\checkmarki\checkmarkr\checkmarkc\checkmark$\checkmark\\checkmark$\checkmark^\checkmark\\checkmarkc\checkmarki\checkmarkr\checkmarkc\checkmark$\checkmark \checkmark$\checkmark^\checkmark\\checkmarkc\checkmarki\checkmarkr\checkmarkc\checkmark$\checkmark\\checkmark$\checkmark^\checkmark\\checkmarkc\checkmarki\checkmarkr\checkmarkc\checkmark$\checkmarkh\checkmark$\checkmark^\checkmark\\checkmarkc\checkmarki\checkmarkr\checkmarkc\checkmark$\checkmarkl\checkmark$\checkmark^\checkmark\\checkmarkc\checkmarki\checkmarkr\checkmarkc\checkmark$\checkmarki\checkmark$\checkmark^\checkmark\\checkmarkc\checkmarki\checkmarkr\checkmarkc\checkmark$\checkmarkn\checkmark$\checkmark^\checkmark\\checkmarkc\checkmarki\checkmarkr\checkmarkc\checkmark$\checkmarke\checkmark$\checkmark^\checkmark\\checkmarkc\checkmarki\checkmarkr\checkmarkc\checkmark$\checkmark
\checkmark$\checkmark^\checkmark\\checkmarkc\checkmarki\checkmarkr\checkmarkc\checkmark$\checkmarkA\checkmark$\checkmark^\checkmark\\checkmarkc\checkmarki\checkmarkr\checkmarkc\checkmark$\checkmarkr\checkmark$\checkmark^\checkmark\\checkmarkc\checkmarki\checkmarkr\checkmarkc\checkmark$\checkmarkb\checkmark$\checkmark^\checkmark\\checkmarkc\checkmarki\checkmarkr\checkmarkc\checkmark$\checkmarkr\checkmark$\checkmark^\checkmark\\checkmarkc\checkmarki\checkmarkr\checkmarkc\checkmark$\checkmarke\checkmark$\checkmark^\checkmark\\checkmarkc\checkmarki\checkmarkr\checkmarkc\checkmark$\checkmark \checkmark$\checkmark^\checkmark\\checkmarkc\checkmarki\checkmarkr\checkmarkc\checkmark$\checkmarkp\checkmark$\checkmark^\checkmark\\checkmarkc\checkmarki\checkmarkr\checkmarkc\checkmark$\checkmarki\checkmark$\checkmark^\checkmark\\checkmarkc\checkmarki\checkmarkr\checkmarkc\checkmark$\checkmarkv\checkmark$\checkmark^\checkmark\\checkmarkc\checkmarki\checkmarkr\checkmarkc\checkmark$\checkmarko\checkmark$\checkmark^\checkmark\\checkmarkc\checkmarki\checkmarkr\checkmarkc\checkmark$\checkmarkt\checkmark$\checkmark^\checkmark\\checkmarkc\checkmarki\checkmarkr\checkmarkc\checkmark$\checkmark \checkmark$\checkmark^\checkmark\\checkmarkc\checkmarki\checkmarkr\checkmarkc\checkmark$\checkmark(\checkmark$\checkmark^\checkmark\\checkmarkc\checkmarki\checkmarkr\checkmarkc\checkmark$\checkmarkA\checkmark$\checkmark^\checkmark\\checkmarkc\checkmarki\checkmarkr\checkmarkc\checkmark$\checkmarkx\checkmark$\checkmark^\checkmark\\checkmarkc\checkmarki\checkmarkr\checkmarkc\checkmark$\checkmarke\checkmark$\checkmark^\checkmark\\checkmarkc\checkmarki\checkmarkr\checkmarkc\checkmark$\checkmark \checkmark$\checkmark^\checkmark\\checkmarkc\checkmarki\checkmarkr\checkmarkc\checkmark$\checkmarkA\checkmark$\checkmark^\checkmark\\checkmarkc\checkmarki\checkmarkr\checkmarkc\checkmark$\checkmark)\checkmark$\checkmark^\checkmark\\checkmarkc\checkmarki\checkmarkr\checkmarkc\checkmark$\checkmark \checkmark$\checkmark^\checkmark\\checkmarkc\checkmarki\checkmarkr\checkmarkc\checkmark$\checkmark&\checkmark$\checkmark^\checkmark\\checkmarkc\checkmarki\checkmarkr\checkmarkc\checkmark$\checkmark \checkmark$\checkmark^\checkmark\\checkmarkc\checkmarki\checkmarkr\checkmarkc\checkmark$\checkmarkV\checkmark$\checkmark^\checkmark\\checkmarkc\checkmarki\checkmarkr\checkmarkc\checkmark$\checkmarko\checkmark$\checkmark^\checkmark\\checkmarkc\checkmarki\checkmarkr\checkmarkc\checkmark$\checkmarkn\checkmark$\checkmark^\checkmark\\checkmarkc\checkmarki\checkmarkr\checkmarkc\checkmark$\checkmark \checkmark$\checkmark^\checkmark\\checkmarkc\checkmarki\checkmarkr\checkmarkc\checkmark$\checkmarkM\checkmark$\checkmark^\checkmark\\checkmarkc\checkmarki\checkmarkr\checkmarkc\checkmark$\checkmarki\checkmark$\checkmark^\checkmark\\checkmarkc\checkmarki\checkmarkr\checkmarkc\checkmark$\checkmarks\checkmark$\checkmark^\checkmark\\checkmarkc\checkmarki\checkmarkr\checkmarkc\checkmark$\checkmarke\checkmark$\checkmark^\checkmark\\checkmarkc\checkmarki\checkmarkr\checkmarkc\checkmark$\checkmarks\checkmark$\checkmark^\checkmark\\checkmarkc\checkmarki\checkmarkr\checkmarkc\checkmark$\checkmark \checkmark$\checkmark^\checkmark\\checkmarkc\checkmarki\checkmarkr\checkmarkc\checkmark$\checkmarke\checkmark$\checkmark^\checkmark\\checkmarkc\checkmarki\checkmarkr\checkmarkc\checkmark$\checkmarkn\checkmark$\checkmark^\checkmark\\checkmarkc\checkmarki\checkmarkr\checkmarkc\checkmark$\checkmark \checkmark$\checkmark^\checkmark\\checkmarkc\checkmarki\checkmarkr\checkmarkc\checkmark$\checkmarkf\checkmark$\checkmark^\checkmark\\checkmarkc\checkmarki\checkmarkr\checkmarkc\checkmark$\checkmarkl\checkmark$\checkmark^\checkmark\\checkmarkc\checkmarki\checkmarkr\checkmarkc\checkmark$\checkmarke\checkmark$\checkmark^\checkmark\\checkmarkc\checkmarki\checkmarkr\checkmarkc\checkmark$\checkmarkx\checkmark$\checkmark^\checkmark\\checkmarkc\checkmarki\checkmarkr\checkmarkc\checkmark$\checkmarki\checkmark$\checkmark^\checkmark\\checkmarkc\checkmarki\checkmarkr\checkmarkc\checkmark$\checkmarko\checkmark$\checkmark^\checkmark\\checkmarkc\checkmarki\checkmarkr\checkmarkc\checkmark$\checkmarkn\checkmark$\checkmark^\checkmark\\checkmarkc\checkmarki\checkmarkr\checkmarkc\checkmark$\checkmark \checkmark$\checkmark^\checkmark\\checkmarkc\checkmarki\checkmarkr\checkmarkc\checkmark$\checkmark&\checkmark$\checkmark^\checkmark\\checkmarkc\checkmarki\checkmarkr\checkmarkc\checkmark$\checkmark \checkmark$\checkmark^\checkmark\\checkmarkc\checkmarki\checkmarkr\checkmarkc\checkmark$\checkmark$\checkmark$\checkmark^\checkmark\\checkmarkc\checkmarki\checkmarkr\checkmarkc\checkmark$\checkmark\\checkmark$\checkmark^\checkmark\\checkmarkc\checkmarki\checkmarkr\checkmarkc\checkmark$\checkmarks\checkmark$\checkmark^\checkmark\\checkmarkc\checkmarki\checkmarkr\checkmarkc\checkmark$\checkmarki\checkmark$\checkmark^\checkmark\\checkmarkc\checkmarki\checkmarkr\checkmarkc\checkmark$\checkmarkg\checkmark$\checkmark^\checkmark\\checkmarkc\checkmarki\checkmarkr\checkmarkc\checkmark$\checkmarkm\checkmark$\checkmark^\checkmark\\checkmarkc\checkmarki\checkmarkr\checkmarkc\checkmark$\checkmarka\checkmark$\checkmark^\checkmark\\checkmarkc\checkmarki\checkmarkr\checkmarkc\checkmark$\checkmark_\checkmark$\checkmark^\checkmark\\checkmarkc\checkmarki\checkmarkr\checkmarkc\checkmark$\checkmark{\checkmark$\checkmark^\checkmark\\checkmarkc\checkmarki\checkmarkr\checkmarkc\checkmark$\checkmarke\checkmark$\checkmark^\checkmark\\checkmarkc\checkmarki\checkmarkr\checkmarkc\checkmark$\checkmarkq\checkmark$\checkmark^\checkmark\\checkmarkc\checkmarki\checkmarkr\checkmarkc\checkmark$\checkmark}\checkmark$\checkmark^\checkmark\\checkmarkc\checkmarki\checkmarkr\checkmarkc\checkmark$\checkmark \checkmark$\checkmark^\checkmark\\checkmarkc\checkmarki\checkmarkr\checkmarkc\checkmark$\checkmark=\checkmark$\checkmark^\checkmark\\checkmarkc\checkmarki\checkmarkr\checkmarkc\checkmark$\checkmark \checkmark$\checkmark^\checkmark\\checkmarkc\checkmarki\checkmarkr\checkmarkc\checkmark$\checkmark1\checkmark$\checkmark^\checkmark\\checkmarkc\checkmarki\checkmarkr\checkmarkc\checkmark$\checkmark4\checkmark$\checkmark^\checkmark\\checkmarkc\checkmarki\checkmarkr\checkmarkc\checkmark$\checkmark{\checkmark$\checkmark^\checkmark\\checkmarkc\checkmarki\checkmarkr\checkmarkc\checkmark$\checkmark,\checkmark$\checkmark^\checkmark\\checkmarkc\checkmarki\checkmarkr\checkmarkc\checkmark$\checkmark}\checkmark$\checkmark^\checkmark\\checkmarkc\checkmarki\checkmarkr\checkmarkc\checkmark$\checkmark3\checkmark$\checkmark^\checkmark\\checkmarkc\checkmarki\checkmarkr\checkmarkc\checkmark$\checkmark$\checkmark$\checkmark^\checkmark\\checkmarkc\checkmarki\checkmarkr\checkmarkc\checkmark$\checkmark \checkmark$\checkmark^\checkmark\\checkmarkc\checkmarki\checkmarkr\checkmarkc\checkmark$\checkmarkM\checkmark$\checkmark^\checkmark\\checkmarkc\checkmarki\checkmarkr\checkmarkc\checkmark$\checkmarkP\checkmark$\checkmark^\checkmark\\checkmarkc\checkmarki\checkmarkr\checkmarkc\checkmark$\checkmarka\checkmark$\checkmark^\checkmark\\checkmarkc\checkmarki\checkmarkr\checkmarkc\checkmark$\checkmark \checkmark$\checkmark^\checkmark\\checkmarkc\checkmarki\checkmarkr\checkmarkc\checkmark$\checkmark&\checkmark$\checkmark^\checkmark\\checkmarkc\checkmarki\checkmarkr\checkmarkc\checkmark$\checkmark \checkmark$\checkmark^\checkmark\\checkmarkc\checkmarki\checkmarkr\checkmarkc\checkmark$\checkmark$\checkmark$\checkmark^\checkmark\\checkmarkc\checkmarki\checkmarkr\checkmarkc\checkmark$\checkmarkR\checkmark$\checkmark^\checkmark\\checkmarkc\checkmarki\checkmarkr\checkmarkc\checkmark$\checkmark_\checkmark$\checkmark^\checkmark\\checkmarkc\checkmarki\checkmarkr\checkmarkc\checkmark$\checkmarke\checkmark$\checkmark^\checkmark\\checkmarkc\checkmarki\checkmarkr\checkmarkc\checkmark$\checkmark/\checkmark$\checkmark^\checkmark\\checkmarkc\checkmarki\checkmarkr\checkmarkc\checkmark$\checkmarks\checkmark$\checkmark^\checkmark\\checkmarkc\checkmarki\checkmarkr\checkmarkc\checkmark$\checkmark \checkmark$\checkmark^\checkmark\\checkmarkc\checkmarki\checkmarkr\checkmarkc\checkmark$\checkmark=\checkmark$\checkmark^\checkmark\\checkmarkc\checkmarki\checkmarkr\checkmarkc\checkmark$\checkmark \checkmark$\checkmark^\checkmark\\checkmarkc\checkmarki\checkmarkr\checkmarkc\checkmark$\checkmark1\checkmark$\checkmark^\checkmark\\checkmarkc\checkmarki\checkmarkr\checkmarkc\checkmark$\checkmark7\checkmark$\checkmark^\checkmark\\checkmarkc\checkmarki\checkmarkr\checkmarkc\checkmark$\checkmark5\checkmark$\checkmark^\checkmark\\checkmarkc\checkmarki\checkmarkr\checkmarkc\checkmark$\checkmark$\checkmark$\checkmark^\checkmark\\checkmarkc\checkmarki\checkmarkr\checkmarkc\checkmark$\checkmark \checkmark$\checkmark^\checkmark\\checkmarkc\checkmarki\checkmarkr\checkmarkc\checkmark$\checkmarkM\checkmark$\checkmark^\checkmark\\checkmarkc\checkmarki\checkmarkr\checkmarkc\checkmark$\checkmarkP\checkmark$\checkmark^\checkmark\\checkmarkc\checkmarki\checkmarkr\checkmarkc\checkmark$\checkmarka\checkmark$\checkmark^\checkmark\\checkmarkc\checkmarki\checkmarkr\checkmarkc\checkmark$\checkmark \checkmark$\checkmark^\checkmark\\checkmarkc\checkmarki\checkmarkr\checkmarkc\checkmark$\checkmark&\checkmark$\checkmark^\checkmark\\checkmarkc\checkmarki\checkmarkr\checkmarkc\checkmark$\checkmark \checkmark$\checkmark^\checkmark\\checkmarkc\checkmarki\checkmarkr\checkmarkc\checkmark$\checkmark$\checkmark$\checkmark^\checkmark\\checkmarkc\checkmarki\checkmarkr\checkmarkc\checkmark$\checkmark2\checkmark$\checkmark^\checkmark\\checkmarkc\checkmarki\checkmarkr\checkmarkc\checkmark$\checkmark4\checkmark$\checkmark^\checkmark\\checkmarkc\checkmarki\checkmarkr\checkmarkc\checkmark$\checkmark{\checkmark$\checkmark^\checkmark\\checkmarkc\checkmarki\checkmarkr\checkmarkc\checkmark$\checkmark,\checkmark$\checkmark^\checkmark\\checkmarkc\checkmarki\checkmarkr\checkmarkc\checkmark$\checkmark}\checkmark$\checkmark^\checkmark\\checkmarkc\checkmarki\checkmarkr\checkmarkc\checkmark$\checkmark4\checkmark$\checkmark^\checkmark\\checkmarkc\checkmarki\checkmarkr\checkmarkc\checkmark$\checkmark$\checkmark$\checkmark^\checkmark\\checkmarkc\checkmarki\checkmarkr\checkmarkc\checkmark$\checkmark \checkmark$\checkmark^\checkmark\\checkmarkc\checkmarki\checkmarkr\checkmarkc\checkmark$\checkmark\\checkmark$\checkmark^\checkmark\\checkmarkc\checkmarki\checkmarkr\checkmarkc\checkmark$\checkmark\\checkmark$\checkmark^\checkmark\\checkmarkc\checkmarki\checkmarkr\checkmarkc\checkmark$\checkmark \checkmark$\checkmark^\checkmark\\checkmarkc\checkmarki\checkmarkr\checkmarkc\checkmark$\checkmark\\checkmark$\checkmark^\checkmark\\checkmarkc\checkmarki\checkmarkr\checkmarkc\checkmark$\checkmarkh\checkmark$\checkmark^\checkmark\\checkmarkc\checkmarki\checkmarkr\checkmarkc\checkmark$\checkmarkl\checkmark$\checkmark^\checkmark\\checkmarkc\checkmarki\checkmarkr\checkmarkc\checkmark$\checkmarki\checkmark$\checkmark^\checkmark\\checkmarkc\checkmarki\checkmarkr\checkmarkc\checkmark$\checkmarkn\checkmark$\checkmark^\checkmark\\checkmarkc\checkmarki\checkmarkr\checkmarkc\checkmark$\checkmarke\checkmark$\checkmark^\checkmark\\checkmarkc\checkmarki\checkmarkr\checkmarkc\checkmark$\checkmark
\checkmark$\checkmark^\checkmark\\checkmarkc\checkmarki\checkmarkr\checkmarkc\checkmark$\checkmarkP\checkmark$\checkmark^\checkmark\\checkmarkc\checkmarki\checkmarkr\checkmarkc\checkmark$\checkmarko\checkmark$\checkmark^\checkmark\\checkmarkc\checkmarki\checkmarkr\checkmarkc\checkmark$\checkmarkr\checkmark$\checkmark^\checkmark\\checkmarkc\checkmarki\checkmarkr\checkmarkc\checkmark$\checkmarkt\checkmark$\checkmark^\checkmark\\checkmarkc\checkmarki\checkmarkr\checkmarkc\checkmark$\checkmarki\checkmark$\checkmark^\checkmark\\checkmarkc\checkmarki\checkmarkr\checkmarkc\checkmark$\checkmarkq\checkmark$\checkmark^\checkmark\\checkmarkc\checkmarki\checkmarkr\checkmarkc\checkmark$\checkmarku\checkmark$\checkmark^\checkmark\\checkmarkc\checkmarki\checkmarkr\checkmarkc\checkmark$\checkmarke\checkmark$\checkmark^\checkmark\\checkmarkc\checkmarki\checkmarkr\checkmarkc\checkmark$\checkmark \checkmark$\checkmark^\checkmark\\checkmarkc\checkmarki\checkmarkr\checkmarkc\checkmark$\checkmark(\checkmark$\checkmark^\checkmark\\checkmarkc\checkmarki\checkmarkr\checkmarkc\checkmark$\checkmarkA\checkmark$\checkmark^\checkmark\\checkmarkc\checkmarki\checkmarkr\checkmarkc\checkmark$\checkmarkx\checkmark$\checkmark^\checkmark\\checkmarkc\checkmarki\checkmarkr\checkmarkc\checkmark$\checkmarke\checkmark$\checkmark^\checkmark\\checkmarkc\checkmarki\checkmarkr\checkmarkc\checkmark$\checkmark \checkmark$\checkmark^\checkmark\\checkmarkc\checkmarki\checkmarkr\checkmarkc\checkmark$\checkmarkX\checkmark$\checkmark^\checkmark\\checkmarkc\checkmarki\checkmarkr\checkmarkc\checkmark$\checkmark)\checkmark$\checkmark^\checkmark\\checkmarkc\checkmarki\checkmarkr\checkmarkc\checkmark$\checkmark \checkmark$\checkmark^\checkmark\\checkmarkc\checkmarki\checkmarkr\checkmarkc\checkmark$\checkmark&\checkmark$\checkmark^\checkmark\\checkmarkc\checkmarki\checkmarkr\checkmarkc\checkmark$\checkmark \checkmark$\checkmark^\checkmark\\checkmarkc\checkmarki\checkmarkr\checkmarkc\checkmark$\checkmarkD\checkmark$\checkmark^\checkmark\\checkmarkc\checkmarki\checkmarkr\checkmarkc\checkmark$\checkmarké\checkmark$\checkmark^\checkmark\\checkmarkc\checkmarki\checkmarkr\checkmarkc\checkmark$\checkmarkf\checkmark$\checkmark^\checkmark\\checkmarkc\checkmarki\checkmarkr\checkmarkc\checkmark$\checkmarkl\checkmark$\checkmark^\checkmark\\checkmarkc\checkmarki\checkmarkr\checkmarkc\checkmark$\checkmarke\checkmark$\checkmark^\checkmark\\checkmarkc\checkmarki\checkmarkr\checkmarkc\checkmark$\checkmarkx\checkmark$\checkmark^\checkmark\\checkmarkc\checkmarki\checkmarkr\checkmarkc\checkmark$\checkmarki\checkmark$\checkmark^\checkmark\\checkmarkc\checkmarki\checkmarkr\checkmarkc\checkmark$\checkmarko\checkmark$\checkmark^\checkmark\\checkmarkc\checkmarki\checkmarkr\checkmarkc\checkmark$\checkmarkn\checkmark$\checkmark^\checkmark\\checkmarkc\checkmarki\checkmarkr\checkmarkc\checkmark$\checkmark \checkmark$\checkmark^\checkmark\\checkmarkc\checkmarki\checkmarkr\checkmarkc\checkmark$\checkmarks\checkmark$\checkmark^\checkmark\\checkmarkc\checkmarki\checkmarkr\checkmarkc\checkmark$\checkmarkt\checkmark$\checkmark^\checkmark\\checkmarkc\checkmarki\checkmarkr\checkmarkc\checkmark$\checkmarka\checkmark$\checkmark^\checkmark\\checkmarkc\checkmarki\checkmarkr\checkmarkc\checkmark$\checkmarkt\checkmark$\checkmark^\checkmark\\checkmarkc\checkmarki\checkmarkr\checkmarkc\checkmark$\checkmarki\checkmark$\checkmark^\checkmark\\checkmarkc\checkmarki\checkmarkr\checkmarkc\checkmark$\checkmarkq\checkmark$\checkmark^\checkmark\\checkmarkc\checkmarki\checkmarkr\checkmarkc\checkmark$\checkmarku\checkmark$\checkmark^\checkmark\\checkmarkc\checkmarki\checkmarkr\checkmarkc\checkmark$\checkmarke\checkmark$\checkmark^\checkmark\\checkmarkc\checkmarki\checkmarkr\checkmarkc\checkmark$\checkmark \checkmark$\checkmark^\checkmark\\checkmarkc\checkmarki\checkmarkr\checkmarkc\checkmark$\checkmark&\checkmark$\checkmark^\checkmark\\checkmarkc\checkmarki\checkmarkr\checkmarkc\checkmark$\checkmark \checkmark$\checkmark^\checkmark\\checkmarkc\checkmarki\checkmarkr\checkmarkc\checkmark$\checkmark$\checkmark$\checkmark^\checkmark\\checkmarkc\checkmarki\checkmarkr\checkmarkc\checkmark$\checkmarkf\checkmark$\checkmark^\checkmark\\checkmarkc\checkmarki\checkmarkr\checkmarkc\checkmark$\checkmark_\checkmark$\checkmark^\checkmark\\checkmarkc\checkmarki\checkmarkr\checkmarkc\checkmark$\checkmark{\checkmark$\checkmark^\checkmark\\checkmarkc\checkmarki\checkmarkr\checkmarkc\checkmark$\checkmarkm\checkmark$\checkmark^\checkmark\\checkmarkc\checkmarki\checkmarkr\checkmarkc\checkmark$\checkmarka\checkmark$\checkmark^\checkmark\\checkmarkc\checkmarki\checkmarkr\checkmarkc\checkmark$\checkmarkx\checkmark$\checkmark^\checkmark\\checkmarkc\checkmarki\checkmarkr\checkmarkc\checkmark$\checkmark}\checkmark$\checkmark^\checkmark\\checkmarkc\checkmarki\checkmarkr\checkmarkc\checkmark$\checkmark \checkmark$\checkmark^\checkmark\\checkmarkc\checkmarki\checkmarkr\checkmarkc\checkmark$\checkmark=\checkmark$\checkmark^\checkmark\\checkmarkc\checkmarki\checkmarkr\checkmarkc\checkmark$\checkmark \checkmark$\checkmark^\checkmark\\checkmarkc\checkmarki\checkmarkr\checkmarkc\checkmark$\checkmark0\checkmark$\checkmark^\checkmark\\checkmarkc\checkmarki\checkmarkr\checkmarkc\checkmark$\checkmark{\checkmark$\checkmark^\checkmark\\checkmarkc\checkmarki\checkmarkr\checkmarkc\checkmark$\checkmark,\checkmark$\checkmark^\checkmark\\checkmarkc\checkmarki\checkmarkr\checkmarkc\checkmark$\checkmark}\checkmark$\checkmark^\checkmark\\checkmarkc\checkmarki\checkmarkr\checkmarkc\checkmark$\checkmark0\checkmark$\checkmark^\checkmark\\checkmarkc\checkmarki\checkmarkr\checkmarkc\checkmark$\checkmark4\checkmark$\checkmark^\checkmark\\checkmarkc\checkmarki\checkmarkr\checkmarkc\checkmark$\checkmark2\checkmark$\checkmark^\checkmark\\checkmarkc\checkmarki\checkmarkr\checkmarkc\checkmark$\checkmark$\checkmark$\checkmark^\checkmark\\checkmarkc\checkmarki\checkmarkr\checkmarkc\checkmark$\checkmark \checkmark$\checkmark^\checkmark\\checkmarkc\checkmarki\checkmarkr\checkmarkc\checkmark$\checkmarkm\checkmark$\checkmark^\checkmark\\checkmarkc\checkmarki\checkmarkr\checkmarkc\checkmark$\checkmarkm\checkmark$\checkmark^\checkmark\\checkmarkc\checkmarki\checkmarkr\checkmarkc\checkmark$\checkmark \checkmark$\checkmark^\checkmark\\checkmarkc\checkmarki\checkmarkr\checkmarkc\checkmark$\checkmark&\checkmark$\checkmark^\checkmark\\checkmarkc\checkmarki\checkmarkr\checkmarkc\checkmark$\checkmark \checkmark$\checkmark^\checkmark\\checkmarkc\checkmarki\checkmarkr\checkmarkc\checkmark$\checkmark0\checkmark$\checkmark^\checkmark\\checkmarkc\checkmarki\checkmarkr\checkmarkc\checkmark$\checkmark{\checkmark$\checkmark^\checkmark\\checkmarkc\checkmarki\checkmarkr\checkmarkc\checkmark$\checkmark,\checkmark$\checkmark^\checkmark\\checkmarkc\checkmarki\checkmarkr\checkmarkc\checkmark$\checkmark}\checkmark$\checkmark^\checkmark\\checkmarkc\checkmarki\checkmarkr\checkmarkc\checkmark$\checkmark1\checkmark$\checkmark^\checkmark\\checkmarkc\checkmarki\checkmarkr\checkmarkc\checkmark$\checkmark \checkmark$\checkmark^\checkmark\\checkmarkc\checkmarki\checkmarkr\checkmarkc\checkmark$\checkmarkm\checkmark$\checkmark^\checkmark\\checkmarkc\checkmarki\checkmarkr\checkmarkc\checkmark$\checkmarkm\checkmark$\checkmark^\checkmark\\checkmarkc\checkmarki\checkmarkr\checkmarkc\checkmark$\checkmark \checkmark$\checkmark^\checkmark\\checkmarkc\checkmarki\checkmarkr\checkmarkc\checkmark$\checkmark&\checkmark$\checkmark^\checkmark\\checkmarkc\checkmarki\checkmarkr\checkmarkc\checkmark$\checkmark \checkmark$\checkmark^\checkmark\\checkmarkc\checkmarki\checkmarkr\checkmarkc\checkmark$\checkmark$\checkmark$\checkmark^\checkmark\\checkmarkc\checkmarki\checkmarkr\checkmarkc\checkmark$\checkmark2\checkmark$\checkmark^\checkmark\\checkmarkc\checkmarki\checkmarkr\checkmarkc\checkmark$\checkmark{\checkmark$\checkmark^\checkmark\\checkmarkc\checkmarki\checkmarkr\checkmarkc\checkmark$\checkmark,\checkmark$\checkmark^\checkmark\\checkmarkc\checkmarki\checkmarkr\checkmarkc\checkmark$\checkmark}\checkmark$\checkmark^\checkmark\\checkmarkc\checkmarki\checkmarkr\checkmarkc\checkmark$\checkmark4\checkmark$\checkmark^\checkmark\\checkmarkc\checkmarki\checkmarkr\checkmarkc\checkmark$\checkmark$\checkmark$\checkmark^\checkmark\\checkmarkc\checkmarki\checkmarkr\checkmarkc\checkmark$\checkmark \checkmark$\checkmark^\checkmark\\checkmarkc\checkmarki\checkmarkr\checkmarkc\checkmark$\checkmark\\checkmark$\checkmark^\checkmark\\checkmarkc\checkmarki\checkmarkr\checkmarkc\checkmark$\checkmark\\checkmark$\checkmark^\checkmark\\checkmarkc\checkmarki\checkmarkr\checkmarkc\checkmark$\checkmark \checkmark$\checkmark^\checkmark\\checkmarkc\checkmarki\checkmarkr\checkmarkc\checkmark$\checkmark\\checkmark$\checkmark^\checkmark\\checkmarkc\checkmarki\checkmarkr\checkmarkc\checkmark$\checkmarkh\checkmark$\checkmark^\checkmark\\checkmarkc\checkmarki\checkmarkr\checkmarkc\checkmark$\checkmarkl\checkmark$\checkmark^\checkmark\\checkmarkc\checkmarki\checkmarkr\checkmarkc\checkmark$\checkmarki\checkmark$\checkmark^\checkmark\\checkmarkc\checkmarki\checkmarkr\checkmarkc\checkmark$\checkmarkn\checkmark$\checkmark^\checkmark\\checkmarkc\checkmarki\checkmarkr\checkmarkc\checkmark$\checkmarke\checkmark$\checkmark^\checkmark\\checkmarkc\checkmarki\checkmarkr\checkmarkc\checkmark$\checkmark
\checkmark$\checkmark^\checkmark\\checkmarkc\checkmarki\checkmarkr\checkmarkc\checkmark$\checkmarkP\checkmark$\checkmark^\checkmark\\checkmarkc\checkmarki\checkmarkr\checkmarkc\checkmark$\checkmarko\checkmark$\checkmark^\checkmark\\checkmarkc\checkmarki\checkmarkr\checkmarkc\checkmark$\checkmarkr\checkmark$\checkmark^\checkmark\\checkmarkc\checkmarki\checkmarkr\checkmarkc\checkmark$\checkmarkt\checkmark$\checkmark^\checkmark\\checkmarkc\checkmarki\checkmarkr\checkmarkc\checkmark$\checkmarki\checkmark$\checkmark^\checkmark\\checkmarkc\checkmarki\checkmarkr\checkmarkc\checkmark$\checkmarkq\checkmark$\checkmark^\checkmark\\checkmarkc\checkmarki\checkmarkr\checkmarkc\checkmark$\checkmarku\checkmark$\checkmark^\checkmark\\checkmarkc\checkmarki\checkmarkr\checkmarkc\checkmark$\checkmarke\checkmark$\checkmark^\checkmark\\checkmarkc\checkmarki\checkmarkr\checkmarkc\checkmark$\checkmark \checkmark$\checkmark^\checkmark\\checkmarkc\checkmarki\checkmarkr\checkmarkc\checkmark$\checkmark(\checkmark$\checkmark^\checkmark\\checkmarkc\checkmarki\checkmarkr\checkmarkc\checkmark$\checkmarkA\checkmark$\checkmark^\checkmark\\checkmarkc\checkmarki\checkmarkr\checkmarkc\checkmark$\checkmarkx\checkmark$\checkmark^\checkmark\\checkmarkc\checkmarki\checkmarkr\checkmarkc\checkmark$\checkmarke\checkmark$\checkmark^\checkmark\\checkmarkc\checkmarki\checkmarkr\checkmarkc\checkmark$\checkmark \checkmark$\checkmark^\checkmark\\checkmarkc\checkmarki\checkmarkr\checkmarkc\checkmark$\checkmarkX\checkmark$\checkmark^\checkmark\\checkmarkc\checkmarki\checkmarkr\checkmarkc\checkmark$\checkmark)\checkmark$\checkmark^\checkmark\\checkmarkc\checkmarki\checkmarkr\checkmarkc\checkmark$\checkmark \checkmark$\checkmark^\checkmark\\checkmarkc\checkmarki\checkmarkr\checkmarkc\checkmark$\checkmark&\checkmark$\checkmark^\checkmark\\checkmarkc\checkmarki\checkmarkr\checkmarkc\checkmark$\checkmark \checkmark$\checkmark^\checkmark\\checkmarkc\checkmarki\checkmarkr\checkmarkc\checkmark$\checkmarkV\checkmark$\checkmark^\checkmark\\checkmarkc\checkmarki\checkmarkr\checkmarkc\checkmark$\checkmarko\checkmark$\checkmark^\checkmark\\checkmarkc\checkmarki\checkmarkr\checkmarkc\checkmark$\checkmarkn\checkmark$\checkmark^\checkmark\\checkmarkc\checkmarki\checkmarkr\checkmarkc\checkmark$\checkmark \checkmark$\checkmark^\checkmark\\checkmarkc\checkmarki\checkmarkr\checkmarkc\checkmark$\checkmarkM\checkmark$\checkmark^\checkmark\\checkmarkc\checkmarki\checkmarkr\checkmarkc\checkmark$\checkmarki\checkmark$\checkmark^\checkmark\\checkmarkc\checkmarki\checkmarkr\checkmarkc\checkmark$\checkmarks\checkmark$\checkmark^\checkmark\\checkmarkc\checkmarki\checkmarkr\checkmarkc\checkmark$\checkmarke\checkmark$\checkmark^\checkmark\\checkmarkc\checkmarki\checkmarkr\checkmarkc\checkmark$\checkmarks\checkmark$\checkmark^\checkmark\\checkmarkc\checkmarki\checkmarkr\checkmarkc\checkmark$\checkmark \checkmark$\checkmark^\checkmark\\checkmarkc\checkmarki\checkmarkr\checkmarkc\checkmark$\checkmark&\checkmark$\checkmark^\checkmark\\checkmarkc\checkmarki\checkmarkr\checkmarkc\checkmark$\checkmark \checkmark$\checkmark^\checkmark\\checkmarkc\checkmarki\checkmarkr\checkmarkc\checkmark$\checkmark$\checkmark$\checkmark^\checkmark\\checkmarkc\checkmarki\checkmarkr\checkmarkc\checkmark$\checkmark\\checkmark$\checkmark^\checkmark\\checkmarkc\checkmarki\checkmarkr\checkmarkc\checkmark$\checkmarks\checkmark$\checkmark^\checkmark\\checkmarkc\checkmarki\checkmarkr\checkmarkc\checkmark$\checkmarki\checkmark$\checkmark^\checkmark\\checkmarkc\checkmarki\checkmarkr\checkmarkc\checkmark$\checkmarkg\checkmark$\checkmark^\checkmark\\checkmarkc\checkmarki\checkmarkr\checkmarkc\checkmark$\checkmarkm\checkmark$\checkmark^\checkmark\\checkmarkc\checkmarki\checkmarkr\checkmarkc\checkmark$\checkmarka\checkmark$\checkmark^\checkmark\\checkmarkc\checkmarki\checkmarkr\checkmarkc\checkmark$\checkmark_\checkmark$\checkmark^\checkmark\\checkmarkc\checkmarki\checkmarkr\checkmarkc\checkmark$\checkmark{\checkmark$\checkmark^\checkmark\\checkmarkc\checkmarki\checkmarkr\checkmarkc\checkmark$\checkmarke\checkmark$\checkmark^\checkmark\\checkmarkc\checkmarki\checkmarkr\checkmarkc\checkmark$\checkmarkq\checkmark$\checkmark^\checkmark\\checkmarkc\checkmarki\checkmarkr\checkmarkc\checkmark$\checkmark}\checkmark$\checkmark^\checkmark\\checkmarkc\checkmarki\checkmarkr\checkmarkc\checkmark$\checkmark \checkmark$\checkmark^\checkmark\\checkmarkc\checkmarki\checkmarkr\checkmarkc\checkmark$\checkmark=\checkmark$\checkmark^\checkmark\\checkmarkc\checkmarki\checkmarkr\checkmarkc\checkmark$\checkmark \checkmark$\checkmark^\checkmark\\checkmarkc\checkmarki\checkmarkr\checkmarkc\checkmark$\checkmark8\checkmark$\checkmark^\checkmark\\checkmarkc\checkmarki\checkmarkr\checkmarkc\checkmark$\checkmark{\checkmark$\checkmark^\checkmark\\checkmarkc\checkmarki\checkmarkr\checkmarkc\checkmark$\checkmark,\checkmark$\checkmark^\checkmark\\checkmarkc\checkmarki\checkmarkr\checkmarkc\checkmark$\checkmark}\checkmark$\checkmark^\checkmark\\checkmarkc\checkmarki\checkmarkr\checkmarkc\checkmark$\checkmark6\checkmark$\checkmark^\checkmark\\checkmarkc\checkmarki\checkmarkr\checkmarkc\checkmark$\checkmark$\checkmark$\checkmark^\checkmark\\checkmarkc\checkmarki\checkmarkr\checkmarkc\checkmark$\checkmark \checkmark$\checkmark^\checkmark\\checkmarkc\checkmarki\checkmarkr\checkmarkc\checkmark$\checkmarkM\checkmark$\checkmark^\checkmark\\checkmarkc\checkmarki\checkmarkr\checkmarkc\checkmark$\checkmarkP\checkmark$\checkmark^\checkmark\\checkmarkc\checkmarki\checkmarkr\checkmarkc\checkmark$\checkmarka\checkmark$\checkmark^\checkmark\\checkmarkc\checkmarki\checkmarkr\checkmarkc\checkmark$\checkmark \checkmark$\checkmark^\checkmark\\checkmarkc\checkmarki\checkmarkr\checkmarkc\checkmark$\checkmark&\checkmark$\checkmark^\checkmark\\checkmarkc\checkmarki\checkmarkr\checkmarkc\checkmark$\checkmark \checkmark$\checkmark^\checkmark\\checkmarkc\checkmarki\checkmarkr\checkmarkc\checkmark$\checkmark$\checkmark$\checkmark^\checkmark\\checkmarkc\checkmarki\checkmarkr\checkmarkc\checkmark$\checkmarkR\checkmark$\checkmark^\checkmark\\checkmarkc\checkmarki\checkmarkr\checkmarkc\checkmark$\checkmark_\checkmark$\checkmark^\checkmark\\checkmarkc\checkmarki\checkmarkr\checkmarkc\checkmark$\checkmarke\checkmark$\checkmark^\checkmark\\checkmarkc\checkmarki\checkmarkr\checkmarkc\checkmark$\checkmark/\checkmark$\checkmark^\checkmark\\checkmarkc\checkmarki\checkmarkr\checkmarkc\checkmark$\checkmarks\checkmark$\checkmark^\checkmark\\checkmarkc\checkmarki\checkmarkr\checkmarkc\checkmark$\checkmark \checkmark$\checkmark^\checkmark\\checkmarkc\checkmarki\checkmarkr\checkmarkc\checkmark$\checkmark=\checkmark$\checkmark^\checkmark\\checkmarkc\checkmarki\checkmarkr\checkmarkc\checkmark$\checkmark \checkmark$\checkmark^\checkmark\\checkmarkc\checkmarki\checkmarkr\checkmarkc\checkmark$\checkmark9\checkmark$\checkmark^\checkmark\\checkmarkc\checkmarki\checkmarkr\checkmarkc\checkmark$\checkmark2\checkmark$\checkmark^\checkmark\\checkmarkc\checkmarki\checkmarkr\checkmarkc\checkmark$\checkmark$\checkmark$\checkmark^\checkmark\\checkmarkc\checkmarki\checkmarkr\checkmarkc\checkmark$\checkmark \checkmark$\checkmark^\checkmark\\checkmarkc\checkmarki\checkmarkr\checkmarkc\checkmark$\checkmarkM\checkmark$\checkmark^\checkmark\\checkmarkc\checkmarki\checkmarkr\checkmarkc\checkmark$\checkmarkP\checkmark$\checkmark^\checkmark\\checkmarkc\checkmarki\checkmarkr\checkmarkc\checkmark$\checkmarka\checkmark$\checkmark^\checkmark\\checkmarkc\checkmarki\checkmarkr\checkmarkc\checkmark$\checkmark \checkmark$\checkmark^\checkmark\\checkmarkc\checkmarki\checkmarkr\checkmarkc\checkmark$\checkmark&\checkmark$\checkmark^\checkmark\\checkmarkc\checkmarki\checkmarkr\checkmarkc\checkmark$\checkmark \checkmark$\checkmark^\checkmark\\checkmarkc\checkmarki\checkmarkr\checkmarkc\checkmark$\checkmark$\checkmark$\checkmark^\checkmark\\checkmarkc\checkmarki\checkmarkr\checkmarkc\checkmark$\checkmark3\checkmark$\checkmark^\checkmark\\checkmarkc\checkmarki\checkmarkr\checkmarkc\checkmark$\checkmark2\checkmark$\checkmark^\checkmark\\checkmarkc\checkmarki\checkmarkr\checkmarkc\checkmark$\checkmark{\checkmark$\checkmark^\checkmark\\checkmarkc\checkmarki\checkmarkr\checkmarkc\checkmark$\checkmark,\checkmark$\checkmark^\checkmark\\checkmarkc\checkmarki\checkmarkr\checkmarkc\checkmark$\checkmark}\checkmark$\checkmark^\checkmark\\checkmarkc\checkmarki\checkmarkr\checkmarkc\checkmark$\checkmark1\checkmark$\checkmark^\checkmark\\checkmarkc\checkmarki\checkmarkr\checkmarkc\checkmark$\checkmark$\checkmark$\checkmark^\checkmark\\checkmarkc\checkmarki\checkmarkr\checkmarkc\checkmark$\checkmark \checkmark$\checkmark^\checkmark\\checkmarkc\checkmarki\checkmarkr\checkmarkc\checkmark$\checkmark\\checkmark$\checkmark^\checkmark\\checkmarkc\checkmarki\checkmarkr\checkmarkc\checkmark$\checkmark\\checkmark$\checkmark^\checkmark\\checkmarkc\checkmarki\checkmarkr\checkmarkc\checkmark$\checkmark \checkmark$\checkmark^\checkmark\\checkmarkc\checkmarki\checkmarkr\checkmarkc\checkmark$\checkmark\\checkmark$\checkmark^\checkmark\\checkmarkc\checkmarki\checkmarkr\checkmarkc\checkmark$\checkmarkh\checkmark$\checkmark^\checkmark\\checkmarkc\checkmarki\checkmarkr\checkmarkc\checkmark$\checkmarkl\checkmark$\checkmark^\checkmark\\checkmarkc\checkmarki\checkmarkr\checkmarkc\checkmark$\checkmarki\checkmark$\checkmark^\checkmark\\checkmarkc\checkmarki\checkmarkr\checkmarkc\checkmark$\checkmarkn\checkmark$\checkmark^\checkmark\\checkmarkc\checkmarki\checkmarkr\checkmarkc\checkmark$\checkmarke\checkmark$\checkmark^\checkmark\\checkmarkc\checkmarki\checkmarkr\checkmarkc\checkmark$\checkmark
\checkmark$\checkmark^\checkmark\\checkmarkc\checkmarki\checkmarkr\checkmarkc\checkmark$\checkmarkG\checkmark$\checkmark^\checkmark\\checkmarkc\checkmarki\checkmarkr\checkmarkc\checkmark$\checkmarku\checkmark$\checkmark^\checkmark\\checkmarkc\checkmarki\checkmarkr\checkmarkc\checkmark$\checkmarki\checkmark$\checkmark^\checkmark\\checkmarkc\checkmarki\checkmarkr\checkmarkc\checkmark$\checkmarkd\checkmark$\checkmark^\checkmark\\checkmarkc\checkmarki\checkmarkr\checkmarkc\checkmark$\checkmarka\checkmark$\checkmark^\checkmark\\checkmarkc\checkmarki\checkmarkr\checkmarkc\checkmark$\checkmarkg\checkmark$\checkmark^\checkmark\\checkmarkc\checkmarki\checkmarkr\checkmarkc\checkmark$\checkmarke\checkmark$\checkmark^\checkmark\\checkmarkc\checkmarki\checkmarkr\checkmarkc\checkmark$\checkmark \checkmark$\checkmark^\checkmark\\checkmarkc\checkmarki\checkmarkr\checkmarkc\checkmark$\checkmarkH\checkmark$\checkmark^\checkmark\\checkmarkc\checkmarki\checkmarkr\checkmarkc\checkmark$\checkmarkG\checkmark$\checkmark^\checkmark\\checkmarkc\checkmarki\checkmarkr\checkmarkc\checkmark$\checkmarkR\checkmark$\checkmark^\checkmark\\checkmarkc\checkmarki\checkmarkr\checkmarkc\checkmark$\checkmark1\checkmark$\checkmark^\checkmark\\checkmarkc\checkmarki\checkmarkr\checkmarkc\checkmark$\checkmark5\checkmark$\checkmark^\checkmark\\checkmarkc\checkmarki\checkmarkr\checkmarkc\checkmark$\checkmark \checkmark$\checkmark^\checkmark\\checkmarkc\checkmarki\checkmarkr\checkmarkc\checkmark$\checkmark&\checkmark$\checkmark^\checkmark\\checkmarkc\checkmarki\checkmarkr\checkmarkc\checkmark$\checkmark \checkmark$\checkmark^\checkmark\\checkmarkc\checkmarki\checkmarkr\checkmarkc\checkmark$\checkmarkD\checkmark$\checkmark^\checkmark\\checkmarkc\checkmarki\checkmarkr\checkmarkc\checkmark$\checkmarku\checkmark$\checkmark^\checkmark\\checkmarkc\checkmarki\checkmarkr\checkmarkc\checkmark$\checkmarkr\checkmark$\checkmark^\checkmark\\checkmarkc\checkmarki\checkmarkr\checkmarkc\checkmark$\checkmarké\checkmark$\checkmark^\checkmark\\checkmarkc\checkmarki\checkmarkr\checkmarkc\checkmark$\checkmarke\checkmark$\checkmark^\checkmark\\checkmarkc\checkmarki\checkmarkr\checkmarkc\checkmark$\checkmark \checkmark$\checkmark^\checkmark\\checkmarkc\checkmarki\checkmarkr\checkmarkc\checkmark$\checkmarkd\checkmark$\checkmark^\checkmark\\checkmarkc\checkmarki\checkmarkr\checkmarkc\checkmark$\checkmarke\checkmark$\checkmark^\checkmark\\checkmarkc\checkmarki\checkmarkr\checkmarkc\checkmark$\checkmark \checkmark$\checkmark^\checkmark\\checkmarkc\checkmarki\checkmarkr\checkmarkc\checkmark$\checkmarkv\checkmark$\checkmark^\checkmark\\checkmarkc\checkmarki\checkmarkr\checkmarkc\checkmark$\checkmarki\checkmark$\checkmark^\checkmark\\checkmarkc\checkmarki\checkmarkr\checkmarkc\checkmark$\checkmarke\checkmark$\checkmark^\checkmark\\checkmarkc\checkmarki\checkmarkr\checkmarkc\checkmark$\checkmark \checkmark$\checkmark^\checkmark\\checkmarkc\checkmarki\checkmarkr\checkmarkc\checkmark$\checkmarkI\checkmark$\checkmark^\checkmark\\checkmarkc\checkmarki\checkmarkr\checkmarkc\checkmark$\checkmarkS\checkmark$\checkmark^\checkmark\\checkmarkc\checkmarki\checkmarkr\checkmarkc\checkmark$\checkmarkO\checkmark$\checkmark^\checkmark\\checkmarkc\checkmarki\checkmarkr\checkmarkc\checkmark$\checkmark \checkmark$\checkmark^\checkmark\\checkmarkc\checkmarki\checkmarkr\checkmarkc\checkmark$\checkmark2\checkmark$\checkmark^\checkmark\\checkmarkc\checkmarki\checkmarkr\checkmarkc\checkmark$\checkmark8\checkmark$\checkmark^\checkmark\\checkmarkc\checkmarki\checkmarkr\checkmarkc\checkmark$\checkmark1\checkmark$\checkmark^\checkmark\\checkmarkc\checkmarki\checkmarkr\checkmarkc\checkmark$\checkmark \checkmark$\checkmark^\checkmark\\checkmarkc\checkmarki\checkmarkr\checkmarkc\checkmark$\checkmark&\checkmark$\checkmark^\checkmark\\checkmarkc\checkmarki\checkmarkr\checkmarkc\checkmark$\checkmark \checkmark$\checkmark^\checkmark\\checkmarkc\checkmarki\checkmarkr\checkmarkc\checkmark$\checkmark$\checkmark$\checkmark^\checkmark\\checkmarkc\checkmarki\checkmarkr\checkmarkc\checkmark$\checkmarkL\checkmark$\checkmark^\checkmark\\checkmarkc\checkmarki\checkmarkr\checkmarkc\checkmark$\checkmark_\checkmark$\checkmark^\checkmark\\checkmarkc\checkmarki\checkmarkr\checkmarkc\checkmark$\checkmark{\checkmark$\checkmark^\checkmark\\checkmarkc\checkmarki\checkmarkr\checkmarkc\checkmark$\checkmark1\checkmark$\checkmark^\checkmark\\checkmarkc\checkmarki\checkmarkr\checkmarkc\checkmark$\checkmark0\checkmark$\checkmark^\checkmark\\checkmarkc\checkmarki\checkmarkr\checkmarkc\checkmark$\checkmark}\checkmark$\checkmark^\checkmark\\checkmarkc\checkmarki\checkmarkr\checkmarkc\checkmark$\checkmark \checkmark$\checkmark^\checkmark\\checkmarkc\checkmarki\checkmarkr\checkmarkc\checkmark$\checkmark=\checkmark$\checkmark^\checkmark\\checkmarkc\checkmarki\checkmarkr\checkmarkc\checkmark$\checkmark \checkmark$\checkmark^\checkmark\\checkmarkc\checkmarki\checkmarkr\checkmarkc\checkmark$\checkmark6\checkmark$\checkmark^\checkmark\\checkmarkc\checkmarki\checkmarkr\checkmarkc\checkmark$\checkmark{\checkmark$\checkmark^\checkmark\\checkmarkc\checkmarki\checkmarkr\checkmarkc\checkmark$\checkmark,\checkmark$\checkmark^\checkmark\\checkmarkc\checkmarki\checkmarkr\checkmarkc\checkmark$\checkmark}\checkmark$\checkmark^\checkmark\\checkmarkc\checkmarki\checkmarkr\checkmarkc\checkmark$\checkmark1\checkmark$\checkmark^\checkmark\\checkmarkc\checkmarki\checkmarkr\checkmarkc\checkmark$\checkmark\\checkmark$\checkmark^\checkmark\\checkmarkc\checkmarki\checkmarkr\checkmarkc\checkmark$\checkmarkt\checkmark$\checkmark^\checkmark\\checkmarkc\checkmarki\checkmarkr\checkmarkc\checkmark$\checkmarki\checkmark$\checkmark^\checkmark\\checkmarkc\checkmarki\checkmarkr\checkmarkc\checkmark$\checkmarkm\checkmark$\checkmark^\checkmark\\checkmarkc\checkmarki\checkmarkr\checkmarkc\checkmark$\checkmarke\checkmark$\checkmark^\checkmark\\checkmarkc\checkmarki\checkmarkr\checkmarkc\checkmark$\checkmarks\checkmark$\checkmark^\checkmark\\checkmarkc\checkmarki\checkmarkr\checkmarkc\checkmark$\checkmark1\checkmark$\checkmark^\checkmark\\checkmarkc\checkmarki\checkmarkr\checkmarkc\checkmark$\checkmark0\checkmark$\checkmark^\checkmark\\checkmarkc\checkmarki\checkmarkr\checkmarkc\checkmark$\checkmark^\checkmark$\checkmark^\checkmark\\checkmarkc\checkmarki\checkmarkr\checkmarkc\checkmark$\checkmark7\checkmark$\checkmark^\checkmark\\checkmarkc\checkmarki\checkmarkr\checkmarkc\checkmark$\checkmark$\checkmark$\checkmark^\checkmark\\checkmarkc\checkmarki\checkmarkr\checkmarkc\checkmark$\checkmark \checkmark$\checkmark^\checkmark\\checkmarkc\checkmarki\checkmarkr\checkmarkc\checkmark$\checkmarkk\checkmark$\checkmark^\checkmark\\checkmarkc\checkmarki\checkmarkr\checkmarkc\checkmark$\checkmarkm\checkmark$\checkmark^\checkmark\\checkmarkc\checkmarki\checkmarkr\checkmarkc\checkmark$\checkmark \checkmark$\checkmark^\checkmark\\checkmarkc\checkmarki\checkmarkr\checkmarkc\checkmark$\checkmark&\checkmark$\checkmark^\checkmark\\checkmarkc\checkmarki\checkmarkr\checkmarkc\checkmark$\checkmark \checkmark$\checkmark^\checkmark\\checkmarkc\checkmarki\checkmarkr\checkmarkc\checkmark$\checkmarkI\checkmark$\checkmark^\checkmark\\checkmarkc\checkmarki\checkmarkr\checkmarkc\checkmark$\checkmarkl\checkmark$\checkmark^\checkmark\\checkmarkc\checkmarki\checkmarkr\checkmarkc\checkmark$\checkmarkl\checkmark$\checkmark^\checkmark\\checkmarkc\checkmarki\checkmarkr\checkmarkc\checkmark$\checkmarki\checkmark$\checkmark^\checkmark\\checkmarkc\checkmarki\checkmarkr\checkmarkc\checkmark$\checkmarkm\checkmark$\checkmark^\checkmark\\checkmarkc\checkmarki\checkmarkr\checkmarkc\checkmark$\checkmarki\checkmark$\checkmark^\checkmark\\checkmarkc\checkmarki\checkmarkr\checkmarkc\checkmark$\checkmarkt\checkmark$\checkmark^\checkmark\\checkmarkc\checkmarki\checkmarkr\checkmarkc\checkmark$\checkmarké\checkmark$\checkmark^\checkmark\\checkmarkc\checkmarki\checkmarkr\checkmarkc\checkmark$\checkmarke\checkmark$\checkmark^\checkmark\\checkmarkc\checkmarki\checkmarkr\checkmarkc\checkmark$\checkmark \checkmark$\checkmark^\checkmark\\checkmarkc\checkmarki\checkmarkr\checkmarkc\checkmark$\checkmark&\checkmark$\checkmark^\checkmark\\checkmarkc\checkmarki\checkmarkr\checkmarkc\checkmark$\checkmark \checkmark$\checkmark^\checkmark\\checkmarkc\checkmarki\checkmarkr\checkmarkc\checkmark$\checkmark$\checkmark$\checkmark^\checkmark\\checkmarkc\checkmarki\checkmarkr\checkmarkc\checkmark$\checkmark\\checkmark$\checkmark^\checkmark\\checkmarkc\checkmarki\checkmarkr\checkmarkc\checkmark$\checkmarkg\checkmark$\checkmark^\checkmark\\checkmarkc\checkmarki\checkmarkr\checkmarkc\checkmark$\checkmarkg\checkmark$\checkmark^\checkmark\\checkmarkc\checkmarki\checkmarkr\checkmarkc\checkmark$\checkmark \checkmark$\checkmark^\checkmark\\checkmarkc\checkmarki\checkmarkr\checkmarkc\checkmark$\checkmark1\checkmark$\checkmark^\checkmark\\checkmarkc\checkmarki\checkmarkr\checkmarkc\checkmark$\checkmark$\checkmark$\checkmark^\checkmark\\checkmarkc\checkmarki\checkmarkr\checkmarkc\checkmark$\checkmark \checkmark$\checkmark^\checkmark\\checkmarkc\checkmarki\checkmarkr\checkmarkc\checkmark$\checkmark\\checkmark$\checkmark^\checkmark\\checkmarkc\checkmarki\checkmarkr\checkmarkc\checkmark$\checkmark\\checkmark$\checkmark^\checkmark\\checkmarkc\checkmarki\checkmarkr\checkmarkc\checkmark$\checkmark \checkmark$\checkmark^\checkmark\\checkmarkc\checkmarki\checkmarkr\checkmarkc\checkmark$\checkmark\\checkmark$\checkmark^\checkmark\\checkmarkc\checkmarki\checkmarkr\checkmarkc\checkmark$\checkmarkh\checkmark$\checkmark^\checkmark\\checkmarkc\checkmarki\checkmarkr\checkmarkc\checkmark$\checkmarkl\checkmark$\checkmark^\checkmark\\checkmarkc\checkmarki\checkmarkr\checkmarkc\checkmark$\checkmarki\checkmark$\checkmark^\checkmark\\checkmarkc\checkmarki\checkmarkr\checkmarkc\checkmark$\checkmarkn\checkmark$\checkmark^\checkmark\\checkmarkc\checkmarki\checkmarkr\checkmarkc\checkmark$\checkmarke\checkmark$\checkmark^\checkmark\\checkmarkc\checkmarki\checkmarkr\checkmarkc\checkmark$\checkmark
\checkmark$\checkmark^\checkmark\\checkmarkc\checkmarki\checkmarkr\checkmarkc\checkmark$\checkmarkS\checkmark$\checkmark^\checkmark\\checkmarkc\checkmarki\checkmarkr\checkmarkc\checkmark$\checkmarkt\checkmark$\checkmark^\checkmark\\checkmarkc\checkmarki\checkmarkr\checkmarkc\checkmark$\checkmarkr\checkmark$\checkmark^\checkmark\\checkmarkc\checkmarki\checkmarkr\checkmarkc\checkmark$\checkmarku\checkmark$\checkmark^\checkmark\\checkmarkc\checkmarki\checkmarkr\checkmarkc\checkmark$\checkmarkc\checkmark$\checkmark^\checkmark\\checkmarkc\checkmarki\checkmarkr\checkmarkc\checkmark$\checkmarkt\checkmark$\checkmark^\checkmark\\checkmarkc\checkmarki\checkmarkr\checkmarkc\checkmark$\checkmarku\checkmark$\checkmark^\checkmark\\checkmarkc\checkmarki\checkmarkr\checkmarkc\checkmark$\checkmarkr\checkmark$\checkmark^\checkmark\\checkmarkc\checkmarki\checkmarkr\checkmarkc\checkmark$\checkmarke\checkmark$\checkmark^\checkmark\\checkmarkc\checkmarki\checkmarkr\checkmarkc\checkmark$\checkmark \checkmark$\checkmark^\checkmark\\checkmarkc\checkmarki\checkmarkr\checkmarkc\checkmark$\checkmark(\checkmark$\checkmark^\checkmark\\checkmarkc\checkmarki\checkmarkr\checkmarkc\checkmark$\checkmarkA\checkmark$\checkmark^\checkmark\\checkmarkc\checkmarki\checkmarkr\checkmarkc\checkmark$\checkmarkn\checkmark$\checkmark^\checkmark\\checkmarkc\checkmarki\checkmarkr\checkmarkc\checkmark$\checkmarkt\checkmark$\checkmark^\checkmark\\checkmarkc\checkmarki\checkmarkr\checkmarkc\checkmark$\checkmarki\checkmark$\checkmark^\checkmark\\checkmarkc\checkmarki\checkmarkr\checkmarkc\checkmark$\checkmark-\checkmark$\checkmark^\checkmark\\checkmarkc\checkmarki\checkmarkr\checkmarkc\checkmark$\checkmarkr\checkmark$\checkmark^\checkmark\\checkmarkc\checkmarki\checkmarkr\checkmarkc\checkmark$\checkmarké\checkmark$\checkmark^\checkmark\\checkmarkc\checkmarki\checkmarkr\checkmarkc\checkmark$\checkmarks\checkmark$\checkmark^\checkmark\\checkmarkc\checkmarki\checkmarkr\checkmarkc\checkmark$\checkmarko\checkmark$\checkmark^\checkmark\\checkmarkc\checkmarki\checkmarkr\checkmarkc\checkmark$\checkmarkn\checkmark$\checkmark^\checkmark\\checkmarkc\checkmarki\checkmarkr\checkmarkc\checkmark$\checkmarka\checkmark$\checkmark^\checkmark\\checkmarkc\checkmarki\checkmarkr\checkmarkc\checkmark$\checkmarkn\checkmark$\checkmark^\checkmark\\checkmarkc\checkmarki\checkmarkr\checkmarkc\checkmark$\checkmarkc\checkmark$\checkmark^\checkmark\\checkmarkc\checkmarki\checkmarkr\checkmarkc\checkmark$\checkmarke\checkmark$\checkmark^\checkmark\\checkmarkc\checkmarki\checkmarkr\checkmarkc\checkmark$\checkmark)\checkmark$\checkmark^\checkmark\\checkmarkc\checkmarki\checkmarkr\checkmarkc\checkmark$\checkmark \checkmark$\checkmark^\checkmark\\checkmarkc\checkmarki\checkmarkr\checkmarkc\checkmark$\checkmark&\checkmark$\checkmark^\checkmark\\checkmarkc\checkmarki\checkmarkr\checkmarkc\checkmark$\checkmark \checkmark$\checkmark^\checkmark\\checkmarkc\checkmarki\checkmarkr\checkmarkc\checkmark$\checkmarkF\checkmark$\checkmark^\checkmark\\checkmarkc\checkmarki\checkmarkr\checkmarkc\checkmark$\checkmarkr\checkmark$\checkmark^\checkmark\\checkmarkc\checkmarki\checkmarkr\checkmarkc\checkmark$\checkmarké\checkmark$\checkmark^\checkmark\\checkmarkc\checkmarki\checkmarkr\checkmarkc\checkmark$\checkmarkq\checkmark$\checkmark^\checkmark\\checkmarkc\checkmarki\checkmarkr\checkmarkc\checkmark$\checkmarku\checkmark$\checkmark^\checkmark\\checkmarkc\checkmarki\checkmarkr\checkmarkc\checkmark$\checkmarke\checkmark$\checkmark^\checkmark\\checkmarkc\checkmarki\checkmarkr\checkmarkc\checkmark$\checkmarkn\checkmark$\checkmark^\checkmark\\checkmarkc\checkmarki\checkmarkr\checkmarkc\checkmark$\checkmarkc\checkmark$\checkmark^\checkmark\\checkmarkc\checkmarki\checkmarkr\checkmarkc\checkmark$\checkmarke\checkmark$\checkmark^\checkmark\\checkmarkc\checkmarki\checkmarkr\checkmarkc\checkmark$\checkmark \checkmark$\checkmark^\checkmark\\checkmarkc\checkmarki\checkmarkr\checkmarkc\checkmark$\checkmarkp\checkmark$\checkmark^\checkmark\\checkmarkc\checkmarki\checkmarkr\checkmarkc\checkmark$\checkmarkr\checkmark$\checkmark^\checkmark\\checkmarkc\checkmarki\checkmarkr\checkmarkc\checkmark$\checkmarko\checkmark$\checkmark^\checkmark\\checkmarkc\checkmarki\checkmarkr\checkmarkc\checkmark$\checkmarkp\checkmark$\checkmark^\checkmark\\checkmarkc\checkmarki\checkmarkr\checkmarkc\checkmark$\checkmarkr\checkmark$\checkmark^\checkmark\\checkmarkc\checkmarki\checkmarkr\checkmarkc\checkmark$\checkmarke\checkmark$\checkmark^\checkmark\\checkmarkc\checkmarki\checkmarkr\checkmarkc\checkmark$\checkmark \checkmark$\checkmark^\checkmark\\checkmarkc\checkmarki\checkmarkr\checkmarkc\checkmark$\checkmark&\checkmark$\checkmark^\checkmark\\checkmarkc\checkmarki\checkmarkr\checkmarkc\checkmark$\checkmark \checkmark$\checkmark^\checkmark\\checkmarkc\checkmarki\checkmarkr\checkmarkc\checkmark$\checkmark$\checkmark$\checkmark^\checkmark\\checkmarkc\checkmarki\checkmarkr\checkmarkc\checkmark$\checkmarkf\checkmark$\checkmark^\checkmark\\checkmarkc\checkmarki\checkmarkr\checkmarkc\checkmark$\checkmark_\checkmark$\checkmark^\checkmark\\checkmarkc\checkmarki\checkmarkr\checkmarkc\checkmark$\checkmark1\checkmark$\checkmark^\checkmark\\checkmarkc\checkmarki\checkmarkr\checkmarkc\checkmark$\checkmark \checkmark$\checkmark^\checkmark\\checkmarkc\checkmarki\checkmarkr\checkmarkc\checkmark$\checkmark=\checkmark$\checkmark^\checkmark\\checkmarkc\checkmarki\checkmarkr\checkmarkc\checkmark$\checkmark \checkmark$\checkmark^\checkmark\\checkmarkc\checkmarki\checkmarkr\checkmarkc\checkmark$\checkmark2\checkmark$\checkmark^\checkmark\\checkmarkc\checkmarki\checkmarkr\checkmarkc\checkmark$\checkmark9\checkmark$\checkmark^\checkmark\\checkmarkc\checkmarki\checkmarkr\checkmarkc\checkmark$\checkmark0\checkmark$\checkmark^\checkmark\\checkmarkc\checkmarki\checkmarkr\checkmarkc\checkmark$\checkmark$\checkmark$\checkmark^\checkmark\\checkmarkc\checkmarki\checkmarkr\checkmarkc\checkmark$\checkmark \checkmark$\checkmark^\checkmark\\checkmarkc\checkmarki\checkmarkr\checkmarkc\checkmark$\checkmarkH\checkmark$\checkmark^\checkmark\\checkmarkc\checkmarki\checkmarkr\checkmarkc\checkmark$\checkmarkz\checkmark$\checkmark^\checkmark\\checkmarkc\checkmarki\checkmarkr\checkmarkc\checkmark$\checkmark \checkmark$\checkmark^\checkmark\\checkmarkc\checkmarki\checkmarkr\checkmarkc\checkmark$\checkmark&\checkmark$\checkmark^\checkmark\\checkmarkc\checkmarki\checkmarkr\checkmarkc\checkmark$\checkmark \checkmark$\checkmark^\checkmark\\checkmarkc\checkmarki\checkmarkr\checkmarkc\checkmark$\checkmark$\checkmark$\checkmark^\checkmark\\checkmarkc\checkmarki\checkmarkr\checkmarkc\checkmark$\checkmarkf\checkmark$\checkmark^\checkmark\\checkmarkc\checkmarki\checkmarkr\checkmarkc\checkmark$\checkmark_\checkmark$\checkmark^\checkmark\\checkmarkc\checkmarki\checkmarkr\checkmarkc\checkmark$\checkmark{\checkmark$\checkmark^\checkmark\\checkmarkc\checkmarki\checkmarkr\checkmarkc\checkmark$\checkmarke\checkmark$\checkmark^\checkmark\\checkmarkc\checkmarki\checkmarkr\checkmarkc\checkmark$\checkmarkx\checkmark$\checkmark^\checkmark\\checkmarkc\checkmarki\checkmarkr\checkmarkc\checkmark$\checkmarkc\checkmark$\checkmark^\checkmark\\checkmarkc\checkmarki\checkmarkr\checkmarkc\checkmark$\checkmark}\checkmark$\checkmark^\checkmark\\checkmarkc\checkmarki\checkmarkr\checkmarkc\checkmark$\checkmark/\checkmark$\checkmark^\checkmark\\checkmarkc\checkmarki\checkmarkr\checkmarkc\checkmark$\checkmarkf\checkmark$\checkmark^\checkmark\\checkmarkc\checkmarki\checkmarkr\checkmarkc\checkmark$\checkmark_\checkmark$\checkmark^\checkmark\\checkmarkc\checkmarki\checkmarkr\checkmarkc\checkmark$\checkmark1\checkmark$\checkmark^\checkmark\\checkmarkc\checkmarki\checkmarkr\checkmarkc\checkmark$\checkmark \checkmark$\checkmark^\checkmark\\checkmarkc\checkmarki\checkmarkr\checkmarkc\checkmark$\checkmark<\checkmark$\checkmark^\checkmark\\checkmarkc\checkmarki\checkmarkr\checkmarkc\checkmark$\checkmark \checkmark$\checkmark^\checkmark\\checkmarkc\checkmarki\checkmarkr\checkmarkc\checkmark$\checkmark1\checkmark$\checkmark^\checkmark\\checkmarkc\checkmarki\checkmarkr\checkmarkc\checkmark$\checkmark{\checkmark$\checkmark^\checkmark\\checkmarkc\checkmarki\checkmarkr\checkmarkc\checkmark$\checkmark,\checkmark$\checkmark^\checkmark\\checkmarkc\checkmarki\checkmarkr\checkmarkc\checkmark$\checkmark}\checkmark$\checkmark^\checkmark\\checkmarkc\checkmarki\checkmarkr\checkmarkc\checkmark$\checkmark5\checkmark$\checkmark^\checkmark\\checkmarkc\checkmarki\checkmarkr\checkmarkc\checkmark$\checkmark$\checkmark$\checkmark^\checkmark\\checkmarkc\checkmarki\checkmarkr\checkmarkc\checkmark$\checkmark \checkmark$\checkmark^\checkmark\\checkmarkc\checkmarki\checkmarkr\checkmarkc\checkmark$\checkmark&\checkmark$\checkmark^\checkmark\\checkmarkc\checkmarki\checkmarkr\checkmarkc\checkmark$\checkmark \checkmark$\checkmark^\checkmark\\checkmarkc\checkmarki\checkmarkr\checkmarkc\checkmark$\checkmark$\checkmark$\checkmark^\checkmark\\checkmarkc\checkmarki\checkmarkr\checkmarkc\checkmark$\checkmark2\checkmark$\checkmark^\checkmark\\checkmarkc\checkmarki\checkmarkr\checkmarkc\checkmark$\checkmark{\checkmark$\checkmark^\checkmark\\checkmarkc\checkmarki\checkmarkr\checkmarkc\checkmark$\checkmark,\checkmark$\checkmark^\checkmark\\checkmarkc\checkmarki\checkmarkr\checkmarkc\checkmark$\checkmark}\checkmark$\checkmark^\checkmark\\checkmarkc\checkmarki\checkmarkr\checkmarkc\checkmark$\checkmark0\checkmark$\checkmark^\checkmark\\checkmarkc\checkmarki\checkmarkr\checkmarkc\checkmark$\checkmark7\checkmark$\checkmark^\checkmark\\checkmarkc\checkmarki\checkmarkr\checkmarkc\checkmark$\checkmark$\checkmark$\checkmark^\checkmark\\checkmarkc\checkmarki\checkmarkr\checkmarkc\checkmark$\checkmark \checkmark$\checkmark^\checkmark\\checkmarkc\checkmarki\checkmarkr\checkmarkc\checkmark$\checkmark\\checkmark$\checkmark^\checkmark\\checkmarkc\checkmarki\checkmarkr\checkmarkc\checkmark$\checkmark\\checkmark$\checkmark^\checkmark\\checkmarkc\checkmarki\checkmarkr\checkmarkc\checkmark$\checkmark \checkmark$\checkmark^\checkmark\\checkmarkc\checkmarki\checkmarkr\checkmarkc\checkmark$\checkmark\\checkmark$\checkmark^\checkmark\\checkmarkc\checkmarki\checkmarkr\checkmarkc\checkmark$\checkmarkh\checkmark$\checkmark^\checkmark\\checkmarkc\checkmarki\checkmarkr\checkmarkc\checkmark$\checkmarkl\checkmark$\checkmark^\checkmark\\checkmarkc\checkmarki\checkmarkr\checkmarkc\checkmark$\checkmarki\checkmark$\checkmark^\checkmark\\checkmarkc\checkmarki\checkmarkr\checkmarkc\checkmark$\checkmarkn\checkmark$\checkmark^\checkmark\\checkmarkc\checkmarki\checkmarkr\checkmarkc\checkmark$\checkmarke\checkmark$\checkmark^\checkmark\\checkmarkc\checkmarki\checkmarkr\checkmarkc\checkmark$\checkmark
\checkmark$\checkmark^\checkmark\\checkmarkc\checkmarki\checkmarkr\checkmarkc\checkmark$\checkmark\\checkmark$\checkmark^\checkmark\\checkmarkc\checkmarki\checkmarkr\checkmarkc\checkmark$\checkmarke\checkmark$\checkmark^\checkmark\\checkmarkc\checkmarki\checkmarkr\checkmarkc\checkmark$\checkmarkn\checkmark$\checkmark^\checkmark\\checkmarkc\checkmarki\checkmarkr\checkmarkc\checkmark$\checkmarkd\checkmark$\checkmark^\checkmark\\checkmarkc\checkmarki\checkmarkr\checkmarkc\checkmark$\checkmark{\checkmark$\checkmark^\checkmark\\checkmarkc\checkmarki\checkmarkr\checkmarkc\checkmark$\checkmarkt\checkmark$\checkmark^\checkmark\\checkmarkc\checkmarki\checkmarkr\checkmarkc\checkmark$\checkmarka\checkmark$\checkmark^\checkmark\\checkmarkc\checkmarki\checkmarkr\checkmarkc\checkmark$\checkmarkb\checkmark$\checkmark^\checkmark\\checkmarkc\checkmarki\checkmarkr\checkmarkc\checkmark$\checkmarku\checkmark$\checkmark^\checkmark\\checkmarkc\checkmarki\checkmarkr\checkmarkc\checkmark$\checkmarkl\checkmark$\checkmark^\checkmark\\checkmarkc\checkmarki\checkmarkr\checkmarkc\checkmark$\checkmarka\checkmark$\checkmark^\checkmark\\checkmarkc\checkmarki\checkmarkr\checkmarkc\checkmark$\checkmarkr\checkmark$\checkmark^\checkmark\\checkmarkc\checkmarki\checkmarkr\checkmarkc\checkmark$\checkmark}\checkmark$\checkmark^\checkmark\\checkmarkc\checkmarki\checkmarkr\checkmarkc\checkmark$\checkmark
\checkmark$\checkmark^\checkmark\\checkmarkc\checkmarki\checkmarkr\checkmarkc\checkmark$\checkmark}\checkmark$\checkmark^\checkmark\\checkmarkc\checkmarki\checkmarkr\checkmarkc\checkmark$\checkmark
\checkmark$\checkmark^\checkmark\\checkmarkc\checkmarki\checkmarkr\checkmarkc\checkmark$\checkmark\\checkmark$\checkmark^\checkmark\\checkmarkc\checkmarki\checkmarkr\checkmarkc\checkmark$\checkmarkc\checkmark$\checkmark^\checkmark\\checkmarkc\checkmarki\checkmarkr\checkmarkc\checkmark$\checkmarka\checkmark$\checkmark^\checkmark\\checkmarkc\checkmarki\checkmarkr\checkmarkc\checkmark$\checkmarkp\checkmark$\checkmark^\checkmark\\checkmarkc\checkmarki\checkmarkr\checkmarkc\checkmark$\checkmarkt\checkmark$\checkmark^\checkmark\\checkmarkc\checkmarki\checkmarkr\checkmarkc\checkmark$\checkmarki\checkmark$\checkmark^\checkmark\\checkmarkc\checkmarki\checkmarkr\checkmarkc\checkmark$\checkmarko\checkmark$\checkmark^\checkmark\\checkmarkc\checkmarki\checkmarkr\checkmarkc\checkmark$\checkmarkn\checkmark$\checkmark^\checkmark\\checkmarkc\checkmarki\checkmarkr\checkmarkc\checkmark$\checkmark{\checkmark$\checkmark^\checkmark\\checkmarkc\checkmarki\checkmarkr\checkmarkc\checkmark$\checkmarkT\checkmark$\checkmark^\checkmark\\checkmarkc\checkmarki\checkmarkr\checkmarkc\checkmark$\checkmarka\checkmark$\checkmark^\checkmark\\checkmarkc\checkmarki\checkmarkr\checkmarkc\checkmark$\checkmarkb\checkmark$\checkmark^\checkmark\\checkmarkc\checkmarki\checkmarkr\checkmarkc\checkmark$\checkmarkl\checkmark$\checkmark^\checkmark\\checkmarkc\checkmarki\checkmarkr\checkmarkc\checkmark$\checkmarke\checkmark$\checkmark^\checkmark\\checkmarkc\checkmarki\checkmarkr\checkmarkc\checkmark$\checkmarka\checkmark$\checkmark^\checkmark\\checkmarkc\checkmarki\checkmarkr\checkmarkc\checkmark$\checkmarku\checkmark$\checkmark^\checkmark\\checkmarkc\checkmarki\checkmarkr\checkmarkc\checkmark$\checkmark \checkmark$\checkmark^\checkmark\\checkmarkc\checkmarki\checkmarkr\checkmarkc\checkmark$\checkmarkd\checkmark$\checkmark^\checkmark\\checkmarkc\checkmarki\checkmarkr\checkmarkc\checkmark$\checkmarke\checkmark$\checkmark^\checkmark\\checkmarkc\checkmarki\checkmarkr\checkmarkc\checkmark$\checkmark \checkmark$\checkmark^\checkmark\\checkmarkc\checkmarki\checkmarkr\checkmarkc\checkmark$\checkmarks\checkmark$\checkmark^\checkmark\\checkmarkc\checkmarki\checkmarkr\checkmarkc\checkmark$\checkmarky\checkmark$\checkmark^\checkmark\\checkmarkc\checkmarki\checkmarkr\checkmarkc\checkmark$\checkmarkn\checkmark$\checkmark^\checkmark\\checkmarkc\checkmarki\checkmarkr\checkmarkc\checkmark$\checkmarkt\checkmark$\checkmark^\checkmark\\checkmarkc\checkmarki\checkmarkr\checkmarkc\checkmark$\checkmarkh\checkmark$\checkmark^\checkmark\\checkmarkc\checkmarki\checkmarkr\checkmarkc\checkmark$\checkmarkè\checkmark$\checkmark^\checkmark\\checkmarkc\checkmarki\checkmarkr\checkmarkc\checkmark$\checkmarks\checkmark$\checkmark^\checkmark\\checkmarkc\checkmarki\checkmarkr\checkmarkc\checkmark$\checkmarke\checkmark$\checkmark^\checkmark\\checkmarkc\checkmarki\checkmarkr\checkmarkc\checkmark$\checkmark \checkmark$\checkmark^\checkmark\\checkmarkc\checkmarki\checkmarkr\checkmarkc\checkmark$\checkmarkd\checkmark$\checkmark^\checkmark\\checkmarkc\checkmarki\checkmarkr\checkmarkc\checkmark$\checkmarke\checkmark$\checkmark^\checkmark\\checkmarkc\checkmarki\checkmarkr\checkmarkc\checkmark$\checkmarks\checkmark$\checkmark^\checkmark\\checkmarkc\checkmarki\checkmarkr\checkmarkc\checkmark$\checkmark \checkmark$\checkmark^\checkmark\\checkmarkc\checkmarki\checkmarkr\checkmarkc\checkmark$\checkmarkv\checkmark$\checkmark^\checkmark\\checkmarkc\checkmarki\checkmarkr\checkmarkc\checkmark$\checkmarké\checkmark$\checkmark^\checkmark\\checkmarkc\checkmarki\checkmarkr\checkmarkc\checkmark$\checkmarkr\checkmark$\checkmark^\checkmark\\checkmarkc\checkmarki\checkmarkr\checkmarkc\checkmark$\checkmarki\checkmark$\checkmark^\checkmark\\checkmarkc\checkmarki\checkmarkr\checkmarkc\checkmark$\checkmarkf\checkmark$\checkmark^\checkmark\\checkmarkc\checkmarki\checkmarkr\checkmarkc\checkmark$\checkmarki\checkmark$\checkmark^\checkmark\\checkmarkc\checkmarki\checkmarkr\checkmarkc\checkmark$\checkmarkc\checkmark$\checkmark^\checkmark\\checkmarkc\checkmarki\checkmarkr\checkmarkc\checkmark$\checkmarka\checkmark$\checkmark^\checkmark\\checkmarkc\checkmarki\checkmarkr\checkmarkc\checkmark$\checkmarkt\checkmark$\checkmark^\checkmark\\checkmarkc\checkmarki\checkmarkr\checkmarkc\checkmark$\checkmarki\checkmark$\checkmark^\checkmark\\checkmarkc\checkmarki\checkmarkr\checkmarkc\checkmark$\checkmarko\checkmark$\checkmark^\checkmark\\checkmarkc\checkmarki\checkmarkr\checkmarkc\checkmark$\checkmarkn\checkmark$\checkmark^\checkmark\\checkmarkc\checkmarki\checkmarkr\checkmarkc\checkmark$\checkmarks\checkmark$\checkmark^\checkmark\\checkmarkc\checkmarki\checkmarkr\checkmarkc\checkmark$\checkmark \checkmark$\checkmark^\checkmark\\checkmarkc\checkmarki\checkmarkr\checkmarkc\checkmark$\checkmarkd\checkmark$\checkmark^\checkmark\\checkmarkc\checkmarki\checkmarkr\checkmarkc\checkmark$\checkmarke\checkmark$\checkmark^\checkmark\\checkmarkc\checkmarki\checkmarkr\checkmarkc\checkmark$\checkmark \checkmark$\checkmark^\checkmark\\checkmarkc\checkmarki\checkmarkr\checkmarkc\checkmark$\checkmarkd\checkmark$\checkmark^\checkmark\\checkmarkc\checkmarki\checkmarkr\checkmarkc\checkmark$\checkmarki\checkmark$\checkmark^\checkmark\\checkmarkc\checkmarki\checkmarkr\checkmarkc\checkmark$\checkmarkm\checkmark$\checkmark^\checkmark\\checkmarkc\checkmarki\checkmarkr\checkmarkc\checkmark$\checkmarke\checkmark$\checkmark^\checkmark\\checkmarkc\checkmarki\checkmarkr\checkmarkc\checkmark$\checkmarkn\checkmark$\checkmark^\checkmark\\checkmarkc\checkmarki\checkmarkr\checkmarkc\checkmark$\checkmarks\checkmark$\checkmark^\checkmark\\checkmarkc\checkmarki\checkmarkr\checkmarkc\checkmark$\checkmarki\checkmark$\checkmark^\checkmark\\checkmarkc\checkmarki\checkmarkr\checkmarkc\checkmark$\checkmarko\checkmark$\checkmark^\checkmark\\checkmarkc\checkmarki\checkmarkr\checkmarkc\checkmark$\checkmarkn\checkmark$\checkmark^\checkmark\\checkmarkc\checkmarki\checkmarkr\checkmarkc\checkmark$\checkmarkn\checkmark$\checkmark^\checkmark\\checkmarkc\checkmarki\checkmarkr\checkmarkc\checkmark$\checkmarke\checkmark$\checkmark^\checkmark\\checkmarkc\checkmarki\checkmarkr\checkmarkc\checkmark$\checkmarkm\checkmark$\checkmark^\checkmark\\checkmarkc\checkmarki\checkmarkr\checkmarkc\checkmark$\checkmarke\checkmark$\checkmark^\checkmark\\checkmarkc\checkmarki\checkmarkr\checkmarkc\checkmark$\checkmarkn\checkmark$\checkmark^\checkmark\\checkmarkc\checkmarki\checkmarkr\checkmarkc\checkmark$\checkmarkt\checkmark$\checkmark^\checkmark\\checkmarkc\checkmarki\checkmarkr\checkmarkc\checkmark$\checkmark}\checkmark$\checkmark^\checkmark\\checkmarkc\checkmarki\checkmarkr\checkmarkc\checkmark$\checkmark
\checkmark$\checkmark^\checkmark\\checkmarkc\checkmarki\checkmarkr\checkmarkc\checkmark$\checkmark\\checkmark$\checkmark^\checkmark\\checkmarkc\checkmarki\checkmarkr\checkmarkc\checkmark$\checkmarkl\checkmark$\checkmark^\checkmark\\checkmarkc\checkmarki\checkmarkr\checkmarkc\checkmark$\checkmarka\checkmark$\checkmark^\checkmark\\checkmarkc\checkmarki\checkmarkr\checkmarkc\checkmark$\checkmarkb\checkmark$\checkmark^\checkmark\\checkmarkc\checkmarki\checkmarkr\checkmarkc\checkmark$\checkmarke\checkmark$\checkmark^\checkmark\\checkmarkc\checkmarki\checkmarkr\checkmarkc\checkmark$\checkmarkl\checkmark$\checkmark^\checkmark\\checkmarkc\checkmarki\checkmarkr\checkmarkc\checkmark$\checkmark{\checkmark$\checkmark^\checkmark\\checkmarkc\checkmarki\checkmarkr\checkmarkc\checkmark$\checkmarkt\checkmark$\checkmark^\checkmark\\checkmarkc\checkmarki\checkmarkr\checkmarkc\checkmark$\checkmarka\checkmark$\checkmark^\checkmark\\checkmarkc\checkmarki\checkmarkr\checkmarkc\checkmark$\checkmarkb\checkmark$\checkmark^\checkmark\\checkmarkc\checkmarki\checkmarkr\checkmarkc\checkmark$\checkmark:\checkmark$\checkmark^\checkmark\\checkmarkc\checkmarki\checkmarkr\checkmarkc\checkmark$\checkmarks\checkmark$\checkmark^\checkmark\\checkmarkc\checkmarki\checkmarkr\checkmarkc\checkmark$\checkmarky\checkmark$\checkmark^\checkmark\\checkmarkc\checkmarki\checkmarkr\checkmarkc\checkmark$\checkmarkn\checkmark$\checkmark^\checkmark\\checkmarkc\checkmarki\checkmarkr\checkmarkc\checkmark$\checkmarkt\checkmark$\checkmark^\checkmark\\checkmarkc\checkmarki\checkmarkr\checkmarkc\checkmark$\checkmarkh\checkmark$\checkmark^\checkmark\\checkmarkc\checkmarki\checkmarkr\checkmarkc\checkmark$\checkmarke\checkmark$\checkmark^\checkmark\\checkmarkc\checkmarki\checkmarkr\checkmarkc\checkmark$\checkmarks\checkmark$\checkmark^\checkmark\\checkmarkc\checkmarki\checkmarkr\checkmarkc\checkmark$\checkmarke\checkmark$\checkmark^\checkmark\\checkmarkc\checkmarki\checkmarkr\checkmarkc\checkmark$\checkmark_\checkmark$\checkmark^\checkmark\\checkmarkc\checkmarki\checkmarkr\checkmarkc\checkmark$\checkmarkr\checkmark$\checkmark^\checkmark\\checkmarkc\checkmarki\checkmarkr\checkmarkc\checkmark$\checkmarkd\checkmark$\checkmark^\checkmark\\checkmarkc\checkmarki\checkmarkr\checkmarkc\checkmark$\checkmarkm\checkmark$\checkmark^\checkmark\\checkmarkc\checkmarki\checkmarkr\checkmarkc\checkmark$\checkmark}\checkmark$\checkmark^\checkmark\\checkmarkc\checkmarki\checkmarkr\checkmarkc\checkmark$\checkmark
\checkmark$\checkmark^\checkmark\\checkmarkc\checkmarki\checkmarkr\checkmarkc\checkmark$\checkmark\\checkmark$\checkmark^\checkmark\\checkmarkc\checkmarki\checkmarkr\checkmarkc\checkmark$\checkmarke\checkmark$\checkmark^\checkmark\\checkmarkc\checkmarki\checkmarkr\checkmarkc\checkmark$\checkmarkn\checkmark$\checkmark^\checkmark\\checkmarkc\checkmarki\checkmarkr\checkmarkc\checkmark$\checkmarkd\checkmark$\checkmark^\checkmark\\checkmarkc\checkmarki\checkmarkr\checkmarkc\checkmark$\checkmark{\checkmark$\checkmark^\checkmark\\checkmarkc\checkmarki\checkmarkr\checkmarkc\checkmark$\checkmarkt\checkmark$\checkmark^\checkmark\\checkmarkc\checkmarki\checkmarkr\checkmarkc\checkmark$\checkmarka\checkmark$\checkmark^\checkmark\\checkmarkc\checkmarki\checkmarkr\checkmarkc\checkmark$\checkmarkb\checkmark$\checkmark^\checkmark\\checkmarkc\checkmarki\checkmarkr\checkmarkc\checkmark$\checkmarkl\checkmark$\checkmark^\checkmark\\checkmarkc\checkmarki\checkmarkr\checkmarkc\checkmark$\checkmarke\checkmark$\checkmark^\checkmark\\checkmarkc\checkmarki\checkmarkr\checkmarkc\checkmark$\checkmark}\checkmark$\checkmark^\checkmark\\checkmarkc\checkmarki\checkmarkr\checkmarkc\checkmark$\checkmark
\checkmark$\checkmark^\checkmark\\checkmarkc\checkmarki\checkmarkr\checkmarkc\checkmark$\checkmark
\checkmark$\checkmark^\checkmark\\checkmarkc\checkmarki\checkmarkr\checkmarkc\checkmark$\checkmarkL\checkmark$\checkmark^\checkmark\\checkmarkc\checkmarki\checkmarkr\checkmarkc\checkmark$\checkmarke\checkmark$\checkmark^\checkmark\\checkmarkc\checkmarki\checkmarkr\checkmarkc\checkmark$\checkmarks\checkmark$\checkmark^\checkmark\\checkmarkc\checkmarki\checkmarkr\checkmarkc\checkmark$\checkmark \checkmark$\checkmark^\checkmark\\checkmarkc\checkmarki\checkmarkr\checkmarkc\checkmark$\checkmarkc\checkmark$\checkmark^\checkmark\\checkmarkc\checkmarki\checkmarkr\checkmarkc\checkmark$\checkmarka\checkmark$\checkmark^\checkmark\\checkmarkc\checkmarki\checkmarkr\checkmarkc\checkmark$\checkmarkl\checkmark$\checkmark^\checkmark\\checkmarkc\checkmarki\checkmarkr\checkmarkc\checkmark$\checkmarkc\checkmark$\checkmark^\checkmark\\checkmarkc\checkmarki\checkmarkr\checkmarkc\checkmark$\checkmarku\checkmark$\checkmark^\checkmark\\checkmarkc\checkmarki\checkmarkr\checkmarkc\checkmark$\checkmarkl\checkmark$\checkmark^\checkmark\\checkmarkc\checkmarki\checkmarkr\checkmarkc\checkmark$\checkmarks\checkmark$\checkmark^\checkmark\\checkmarkc\checkmarki\checkmarkr\checkmarkc\checkmark$\checkmark \checkmark$\checkmark^\checkmark\\checkmarkc\checkmarki\checkmarkr\checkmarkc\checkmark$\checkmarkf\checkmark$\checkmark^\checkmark\\checkmarkc\checkmarki\checkmarkr\checkmarkc\checkmark$\checkmarko\checkmark$\checkmark^\checkmark\\checkmarkc\checkmarki\checkmarkr\checkmarkc\checkmark$\checkmarkr\checkmark$\checkmark^\checkmark\\checkmarkc\checkmarki\checkmarkr\checkmarkc\checkmark$\checkmarkm\checkmark$\checkmark^\checkmark\\checkmarkc\checkmarki\checkmarkr\checkmarkc\checkmark$\checkmarke\checkmark$\checkmark^\checkmark\\checkmarkc\checkmarki\checkmarkr\checkmarkc\checkmark$\checkmarkl\checkmark$\checkmark^\checkmark\\checkmarkc\checkmarki\checkmarkr\checkmarkc\checkmark$\checkmarks\checkmark$\checkmark^\checkmark\\checkmarkc\checkmarki\checkmarkr\checkmarkc\checkmark$\checkmark \checkmark$\checkmark^\checkmark\\checkmarkc\checkmarki\checkmarkr\checkmarkc\checkmark$\checkmarkd\checkmark$\checkmark^\checkmark\\checkmarkc\checkmarki\checkmarkr\checkmarkc\checkmark$\checkmarke\checkmark$\checkmark^\checkmark\\checkmarkc\checkmarki\checkmarkr\checkmarkc\checkmark$\checkmark \checkmark$\checkmark^\checkmark\\checkmarkc\checkmarki\checkmarkr\checkmarkc\checkmark$\checkmarkd\checkmark$\checkmark^\checkmark\\checkmarkc\checkmarki\checkmarkr\checkmarkc\checkmark$\checkmarki\checkmark$\checkmark^\checkmark\\checkmarkc\checkmarki\checkmarkr\checkmarkc\checkmark$\checkmarkm\checkmark$\checkmark^\checkmark\\checkmarkc\checkmarki\checkmarkr\checkmarkc\checkmark$\checkmarke\checkmark$\checkmark^\checkmark\\checkmarkc\checkmarki\checkmarkr\checkmarkc\checkmark$\checkmarkn\checkmark$\checkmark^\checkmark\\checkmarkc\checkmarki\checkmarkr\checkmarkc\checkmark$\checkmarks\checkmark$\checkmark^\checkmark\\checkmarkc\checkmarki\checkmarkr\checkmarkc\checkmark$\checkmarki\checkmark$\checkmark^\checkmark\\checkmarkc\checkmarki\checkmarkr\checkmarkc\checkmark$\checkmarko\checkmark$\checkmark^\checkmark\\checkmarkc\checkmarki\checkmarkr\checkmarkc\checkmark$\checkmarkn\checkmark$\checkmark^\checkmark\\checkmarkc\checkmarki\checkmarkr\checkmarkc\checkmark$\checkmarkn\checkmark$\checkmark^\checkmark\\checkmarkc\checkmarki\checkmarkr\checkmarkc\checkmark$\checkmarke\checkmark$\checkmark^\checkmark\\checkmarkc\checkmarki\checkmarkr\checkmarkc\checkmark$\checkmarkm\checkmark$\checkmark^\checkmark\\checkmarkc\checkmarki\checkmarkr\checkmarkc\checkmark$\checkmarke\checkmark$\checkmark^\checkmark\\checkmarkc\checkmarki\checkmarkr\checkmarkc\checkmark$\checkmarkn\checkmark$\checkmark^\checkmark\\checkmarkc\checkmarki\checkmarkr\checkmarkc\checkmark$\checkmarkt\checkmark$\checkmark^\checkmark\\checkmarkc\checkmarki\checkmarkr\checkmarkc\checkmark$\checkmark,\checkmark$\checkmark^\checkmark\\checkmarkc\checkmarki\checkmarkr\checkmarkc\checkmark$\checkmark \checkmark$\checkmark^\checkmark\\checkmarkc\checkmarki\checkmarkr\checkmarkc\checkmark$\checkmarkp\checkmark$\checkmark^\checkmark\\checkmarkc\checkmarki\checkmarkr\checkmarkc\checkmark$\checkmarkr\checkmark$\checkmark^\checkmark\\checkmarkc\checkmarki\checkmarkr\checkmarkc\checkmark$\checkmarké\checkmark$\checkmark^\checkmark\\checkmarkc\checkmarki\checkmarkr\checkmarkc\checkmark$\checkmarks\checkmark$\checkmark^\checkmark\\checkmarkc\checkmarki\checkmarkr\checkmarkc\checkmark$\checkmarke\checkmark$\checkmark^\checkmark\\checkmarkc\checkmarki\checkmarkr\checkmarkc\checkmark$\checkmarkn\checkmark$\checkmark^\checkmark\\checkmarkc\checkmarki\checkmarkr\checkmarkc\checkmark$\checkmarkt\checkmark$\checkmark^\checkmark\\checkmarkc\checkmarki\checkmarkr\checkmarkc\checkmark$\checkmarké\checkmark$\checkmark^\checkmark\\checkmarkc\checkmarki\checkmarkr\checkmarkc\checkmark$\checkmarks\checkmark$\checkmark^\checkmark\\checkmarkc\checkmarki\checkmarkr\checkmarkc\checkmark$\checkmark \checkmark$\checkmark^\checkmark\\checkmarkc\checkmarki\checkmarkr\checkmarkc\checkmark$\checkmarkd\checkmark$\checkmark^\checkmark\\checkmarkc\checkmarki\checkmarkr\checkmarkc\checkmark$\checkmarka\checkmark$\checkmark^\checkmark\\checkmarkc\checkmarki\checkmarkr\checkmarkc\checkmark$\checkmarkn\checkmark$\checkmark^\checkmark\\checkmarkc\checkmarki\checkmarkr\checkmarkc\checkmark$\checkmarks\checkmark$\checkmark^\checkmark\\checkmarkc\checkmarki\checkmarkr\checkmarkc\checkmark$\checkmark \checkmark$\checkmark^\checkmark\\checkmarkc\checkmarki\checkmarkr\checkmarkc\checkmark$\checkmarkc\checkmark$\checkmark^\checkmark\\checkmarkc\checkmarki\checkmarkr\checkmarkc\checkmark$\checkmarke\checkmark$\checkmark^\checkmark\\checkmarkc\checkmarki\checkmarkr\checkmarkc\checkmark$\checkmark \checkmark$\checkmark^\checkmark\\checkmarkc\checkmarki\checkmarkr\checkmarkc\checkmark$\checkmarkc\checkmark$\checkmark^\checkmark\\checkmarkc\checkmarki\checkmarkr\checkmarkc\checkmark$\checkmarkh\checkmark$\checkmark^\checkmark\\checkmarkc\checkmarki\checkmarkr\checkmarkc\checkmark$\checkmarka\checkmark$\checkmark^\checkmark\\checkmarkc\checkmarki\checkmarkr\checkmarkc\checkmark$\checkmarkp\checkmark$\checkmark^\checkmark\\checkmarkc\checkmarki\checkmarkr\checkmarkc\checkmark$\checkmarki\checkmark$\checkmark^\checkmark\\checkmarkc\checkmarki\checkmarkr\checkmarkc\checkmark$\checkmarkt\checkmark$\checkmark^\checkmark\\checkmarkc\checkmarki\checkmarkr\checkmarkc\checkmark$\checkmarkr\checkmark$\checkmark^\checkmark\\checkmarkc\checkmarki\checkmarkr\checkmarkc\checkmark$\checkmarke\checkmark$\checkmark^\checkmark\\checkmarkc\checkmarki\checkmarkr\checkmarkc\checkmark$\checkmark,\checkmark$\checkmark^\checkmark\\checkmarkc\checkmarki\checkmarkr\checkmarkc\checkmark$\checkmark \checkmark$\checkmark^\checkmark\\checkmarkc\checkmarki\checkmarkr\checkmarkc\checkmark$\checkmarkc\checkmark$\checkmark^\checkmark\\checkmarkc\checkmarki\checkmarkr\checkmarkc\checkmark$\checkmarko\checkmark$\checkmark^\checkmark\\checkmarkc\checkmarki\checkmarkr\checkmarkc\checkmark$\checkmarkn\checkmark$\checkmark^\checkmark\\checkmarkc\checkmarki\checkmarkr\checkmarkc\checkmark$\checkmarkf\checkmark$\checkmark^\checkmark\\checkmarkc\checkmarki\checkmarkr\checkmarkc\checkmark$\checkmarki\checkmark$\checkmark^\checkmark\\checkmarkc\checkmarki\checkmarkr\checkmarkc\checkmark$\checkmarkr\checkmark$\checkmark^\checkmark\\checkmarkc\checkmarki\checkmarkr\checkmarkc\checkmark$\checkmarkm\checkmark$\checkmark^\checkmark\\checkmarkc\checkmarki\checkmarkr\checkmarkc\checkmark$\checkmarke\checkmark$\checkmark^\checkmark\\checkmarkc\checkmarki\checkmarkr\checkmarkc\checkmark$\checkmarkn\checkmark$\checkmark^\checkmark\\checkmarkc\checkmarki\checkmarkr\checkmarkc\checkmark$\checkmarkt\checkmark$\checkmark^\checkmark\\checkmarkc\checkmarki\checkmarkr\checkmarkc\checkmark$\checkmark \checkmark$\checkmark^\checkmark\\checkmarkc\checkmarki\checkmarkr\checkmarkc\checkmark$\checkmarks\checkmark$\checkmark^\checkmark\\checkmarkc\checkmarki\checkmarkr\checkmarkc\checkmark$\checkmarka\checkmark$\checkmark^\checkmark\\checkmarkc\checkmarki\checkmarkr\checkmarkc\checkmark$\checkmarkn\checkmark$\checkmark^\checkmark\\checkmarkc\checkmarki\checkmarkr\checkmarkc\checkmark$\checkmarks\checkmark$\checkmark^\checkmark\\checkmarkc\checkmarki\checkmarkr\checkmarkc\checkmark$\checkmark \checkmark$\checkmark^\checkmark\\checkmarkc\checkmarki\checkmarkr\checkmarkc\checkmark$\checkmarka\checkmark$\checkmark^\checkmark\\checkmarkc\checkmarki\checkmarkr\checkmarkc\checkmark$\checkmarkm\checkmark$\checkmark^\checkmark\\checkmarkc\checkmarki\checkmarkr\checkmarkc\checkmark$\checkmarkb\checkmark$\checkmark^\checkmark\\checkmarkc\checkmarki\checkmarkr\checkmarkc\checkmark$\checkmarki\checkmark$\checkmark^\checkmark\\checkmarkc\checkmarki\checkmarkr\checkmarkc\checkmark$\checkmarkg\checkmark$\checkmark^\checkmark\\checkmarkc\checkmarki\checkmarkr\checkmarkc\checkmark$\checkmarku\checkmark$\checkmark^\checkmark\\checkmarkc\checkmarki\checkmarkr\checkmarkc\checkmark$\checkmarkï\checkmark$\checkmark^\checkmark\\checkmarkc\checkmarki\checkmarkr\checkmarkc\checkmark$\checkmarkt\checkmark$\checkmark^\checkmark\\checkmarkc\checkmarki\checkmarkr\checkmarkc\checkmark$\checkmarké\checkmark$\checkmark^\checkmark\\checkmarkc\checkmarki\checkmarkr\checkmarkc\checkmark$\checkmark \checkmark$\checkmark^\checkmark\\checkmarkc\checkmarki\checkmarkr\checkmarkc\checkmark$\checkmarkl\checkmark$\checkmark^\checkmark\\checkmarkc\checkmarki\checkmarkr\checkmarkc\checkmark$\checkmark'\checkmark$\checkmark^\checkmark\\checkmarkc\checkmarki\checkmarkr\checkmarkc\checkmark$\checkmarke\checkmark$\checkmark^\checkmark\\checkmarkc\checkmarki\checkmarkr\checkmarkc\checkmark$\checkmarkn\checkmark$\checkmark^\checkmark\\checkmarkc\checkmarki\checkmarkr\checkmarkc\checkmark$\checkmarks\checkmark$\checkmark^\checkmark\\checkmarkc\checkmarki\checkmarkr\checkmarkc\checkmark$\checkmarke\checkmark$\checkmark^\checkmark\\checkmarkc\checkmarki\checkmarkr\checkmarkc\checkmark$\checkmarkm\checkmark$\checkmark^\checkmark\\checkmarkc\checkmarki\checkmarkr\checkmarkc\checkmark$\checkmarkb\checkmark$\checkmark^\checkmark\\checkmarkc\checkmarki\checkmarkr\checkmarkc\checkmark$\checkmarkl\checkmark$\checkmark^\checkmark\\checkmarkc\checkmarki\checkmarkr\checkmarkc\checkmark$\checkmarke\checkmark$\checkmark^\checkmark\\checkmarkc\checkmarki\checkmarkr\checkmarkc\checkmark$\checkmark \checkmark$\checkmark^\checkmark\\checkmarkc\checkmarki\checkmarkr\checkmarkc\checkmark$\checkmarkd\checkmark$\checkmark^\checkmark\\checkmarkc\checkmarki\checkmarkr\checkmarkc\checkmark$\checkmarke\checkmark$\checkmark^\checkmark\\checkmarkc\checkmarki\checkmarkr\checkmarkc\checkmark$\checkmarks\checkmark$\checkmark^\checkmark\\checkmarkc\checkmarki\checkmarkr\checkmarkc\checkmark$\checkmark \checkmark$\checkmark^\checkmark\\checkmarkc\checkmarki\checkmarkr\checkmarkc\checkmark$\checkmarkc\checkmark$\checkmark^\checkmark\\checkmarkc\checkmarki\checkmarkr\checkmarkc\checkmark$\checkmarkh\checkmark$\checkmark^\checkmark\\checkmarkc\checkmarki\checkmarkr\checkmarkc\checkmark$\checkmarko\checkmark$\checkmark^\checkmark\\checkmarkc\checkmarki\checkmarkr\checkmarkc\checkmark$\checkmarki\checkmark$\checkmark^\checkmark\\checkmarkc\checkmarki\checkmarkr\checkmarkc\checkmark$\checkmarkx\checkmark$\checkmark^\checkmark\\checkmarkc\checkmarki\checkmarkr\checkmarkc\checkmark$\checkmark \checkmark$\checkmark^\checkmark\\checkmarkc\checkmarki\checkmarkr\checkmarkc\checkmark$\checkmarkt\checkmark$\checkmark^\checkmark\\checkmarkc\checkmarki\checkmarkr\checkmarkc\checkmark$\checkmarke\checkmark$\checkmark^\checkmark\\checkmarkc\checkmarki\checkmarkr\checkmarkc\checkmark$\checkmarkc\checkmark$\checkmark^\checkmark\\checkmarkc\checkmarki\checkmarkr\checkmarkc\checkmark$\checkmarkh\checkmark$\checkmark^\checkmark\\checkmarkc\checkmarki\checkmarkr\checkmarkc\checkmark$\checkmarkn\checkmark$\checkmark^\checkmark\\checkmarkc\checkmarki\checkmarkr\checkmarkc\checkmark$\checkmarko\checkmark$\checkmark^\checkmark\\checkmarkc\checkmarki\checkmarkr\checkmarkc\checkmark$\checkmarkl\checkmark$\checkmark^\checkmark\\checkmarkc\checkmarki\checkmarkr\checkmarkc\checkmark$\checkmarko\checkmark$\checkmark^\checkmark\\checkmarkc\checkmarki\checkmarkr\checkmarkc\checkmark$\checkmarkg\checkmark$\checkmark^\checkmark\\checkmarkc\checkmarki\checkmarkr\checkmarkc\checkmark$\checkmarki\checkmark$\checkmark^\checkmark\\checkmarkc\checkmarki\checkmarkr\checkmarkc\checkmark$\checkmarkq\checkmark$\checkmark^\checkmark\\checkmarkc\checkmarki\checkmarkr\checkmarkc\checkmark$\checkmarku\checkmark$\checkmark^\checkmark\\checkmarkc\checkmarki\checkmarkr\checkmarkc\checkmark$\checkmarke\checkmark$\checkmark^\checkmark\\checkmarkc\checkmarki\checkmarkr\checkmarkc\checkmark$\checkmarks\checkmark$\checkmark^\checkmark\\checkmarkc\checkmarki\checkmarkr\checkmarkc\checkmark$\checkmark \checkmark$\checkmark^\checkmark\\checkmarkc\checkmarki\checkmarkr\checkmarkc\checkmark$\checkmarki\checkmark$\checkmark^\checkmark\\checkmarkc\checkmarki\checkmarkr\checkmarkc\checkmark$\checkmarkn\checkmark$\checkmark^\checkmark\\checkmarkc\checkmarki\checkmarkr\checkmarkc\checkmark$\checkmarki\checkmark$\checkmark^\checkmark\\checkmarkc\checkmarki\checkmarkr\checkmarkc\checkmark$\checkmarkt\checkmark$\checkmark^\checkmark\\checkmarkc\checkmarki\checkmarkr\checkmarkc\checkmark$\checkmarki\checkmark$\checkmark^\checkmark\\checkmarkc\checkmarki\checkmarkr\checkmarkc\checkmark$\checkmarka\checkmark$\checkmark^\checkmark\\checkmarkc\checkmarki\checkmarkr\checkmarkc\checkmark$\checkmarku\checkmark$\checkmark^\checkmark\\checkmarkc\checkmarki\checkmarkr\checkmarkc\checkmark$\checkmarkx\checkmark$\checkmark^\checkmark\\checkmarkc\checkmarki\checkmarkr\checkmarkc\checkmark$\checkmark \checkmark$\checkmark^\checkmark\\checkmarkc\checkmarki\checkmarkr\checkmarkc\checkmark$\checkmarkd\checkmark$\checkmark^\checkmark\\checkmarkc\checkmarki\checkmarkr\checkmarkc\checkmark$\checkmarku\checkmark$\checkmark^\checkmark\\checkmarkc\checkmarki\checkmarkr\checkmarkc\checkmark$\checkmark \checkmark$\checkmark^\checkmark\\checkmarkc\checkmarki\checkmarkr\checkmarkc\checkmark$\checkmarkc\checkmark$\checkmark^\checkmark\\checkmarkc\checkmarki\checkmarkr\checkmarkc\checkmark$\checkmarka\checkmark$\checkmark^\checkmark\\checkmarkc\checkmarki\checkmarkr\checkmarkc\checkmark$\checkmarkh\checkmark$\checkmark^\checkmark\\checkmarkc\checkmarki\checkmarkr\checkmarkc\checkmark$\checkmarki\checkmark$\checkmark^\checkmark\\checkmarkc\checkmarki\checkmarkr\checkmarkc\checkmark$\checkmarke\checkmark$\checkmark^\checkmark\\checkmarkc\checkmarki\checkmarkr\checkmarkc\checkmark$\checkmarkr\checkmark$\checkmark^\checkmark\\checkmarkc\checkmarki\checkmarkr\checkmarkc\checkmark$\checkmark \checkmark$\checkmark^\checkmark\\checkmarkc\checkmarki\checkmarkr\checkmarkc\checkmark$\checkmarkd\checkmark$\checkmark^\checkmark\\checkmarkc\checkmarki\checkmarkr\checkmarkc\checkmark$\checkmarke\checkmark$\checkmark^\checkmark\\checkmarkc\checkmarki\checkmarkr\checkmarkc\checkmark$\checkmarks\checkmark$\checkmark^\checkmark\\checkmarkc\checkmarki\checkmarkr\checkmarkc\checkmark$\checkmark \checkmark$\checkmark^\checkmark\\checkmarkc\checkmarki\checkmarkr\checkmarkc\checkmark$\checkmarkc\checkmark$\checkmark^\checkmark\\checkmarkc\checkmarki\checkmarkr\checkmarkc\checkmark$\checkmarkh\checkmark$\checkmark^\checkmark\\checkmarkc\checkmarki\checkmarkr\checkmarkc\checkmark$\checkmarka\checkmark$\checkmark^\checkmark\\checkmarkc\checkmarki\checkmarkr\checkmarkc\checkmark$\checkmarkr\checkmark$\checkmark^\checkmark\\checkmarkc\checkmarki\checkmarkr\checkmarkc\checkmark$\checkmarkg\checkmark$\checkmark^\checkmark\\checkmarkc\checkmarki\checkmarkr\checkmarkc\checkmark$\checkmarke\checkmark$\checkmark^\checkmark\\checkmarkc\checkmarki\checkmarkr\checkmarkc\checkmark$\checkmarks\checkmark$\checkmark^\checkmark\\checkmarkc\checkmarki\checkmarkr\checkmarkc\checkmark$\checkmark.\checkmark$\checkmark^\checkmark\\checkmarkc\checkmarki\checkmarkr\checkmarkc\checkmark$\checkmark \checkmark$\checkmark^\checkmark\\checkmarkc\checkmarki\checkmarkr\checkmarkc\checkmark$\checkmarkT\checkmark$\checkmark^\checkmark\\checkmarkc\checkmarki\checkmarkr\checkmarkc\checkmark$\checkmarko\checkmark$\checkmark^\checkmark\\checkmarkc\checkmarki\checkmarkr\checkmarkc\checkmark$\checkmarku\checkmark$\checkmark^\checkmark\\checkmarkc\checkmarki\checkmarkr\checkmarkc\checkmark$\checkmarks\checkmark$\checkmark^\checkmark\\checkmarkc\checkmarki\checkmarkr\checkmarkc\checkmark$\checkmark \checkmark$\checkmark^\checkmark\\checkmarkc\checkmarki\checkmarkr\checkmarkc\checkmark$\checkmarkl\checkmark$\checkmark^\checkmark\\checkmarkc\checkmarki\checkmarkr\checkmarkc\checkmark$\checkmarke\checkmark$\checkmark^\checkmark\\checkmarkc\checkmarki\checkmarkr\checkmarkc\checkmark$\checkmarks\checkmark$\checkmark^\checkmark\\checkmarkc\checkmarki\checkmarkr\checkmarkc\checkmark$\checkmark \checkmark$\checkmark^\checkmark\\checkmarkc\checkmarki\checkmarkr\checkmarkc\checkmark$\checkmarkc\checkmark$\checkmark^\checkmark\\checkmarkc\checkmarki\checkmarkr\checkmarkc\checkmark$\checkmarko\checkmark$\checkmark^\checkmark\\checkmarkc\checkmarki\checkmarkr\checkmarkc\checkmark$\checkmarke\checkmark$\checkmark^\checkmark\\checkmarkc\checkmarki\checkmarkr\checkmarkc\checkmark$\checkmarkf\checkmark$\checkmark^\checkmark\\checkmarkc\checkmarki\checkmarkr\checkmarkc\checkmark$\checkmarkf\checkmark$\checkmark^\checkmark\\checkmarkc\checkmarki\checkmarkr\checkmarkc\checkmark$\checkmarki\checkmark$\checkmark^\checkmark\\checkmarkc\checkmarki\checkmarkr\checkmarkc\checkmark$\checkmarkc\checkmark$\checkmark^\checkmark\\checkmarkc\checkmarki\checkmarkr\checkmarkc\checkmark$\checkmarki\checkmark$\checkmark^\checkmark\\checkmarkc\checkmarki\checkmarkr\checkmarkc\checkmark$\checkmarke\checkmark$\checkmark^\checkmark\\checkmarkc\checkmarki\checkmarkr\checkmarkc\checkmark$\checkmarkn\checkmark$\checkmark^\checkmark\\checkmarkc\checkmarki\checkmarkr\checkmarkc\checkmark$\checkmarkt\checkmark$\checkmark^\checkmark\\checkmarkc\checkmarki\checkmarkr\checkmarkc\checkmark$\checkmarks\checkmark$\checkmark^\checkmark\\checkmarkc\checkmarki\checkmarkr\checkmarkc\checkmark$\checkmark \checkmark$\checkmark^\checkmark\\checkmarkc\checkmarki\checkmarkr\checkmarkc\checkmark$\checkmarkd\checkmark$\checkmark^\checkmark\\checkmarkc\checkmarki\checkmarkr\checkmarkc\checkmark$\checkmarke\checkmark$\checkmark^\checkmark\\checkmarkc\checkmarki\checkmarkr\checkmarkc\checkmark$\checkmark \checkmark$\checkmark^\checkmark\\checkmarkc\checkmarki\checkmarkr\checkmarkc\checkmark$\checkmarks\checkmark$\checkmark^\checkmark\\checkmarkc\checkmarki\checkmarkr\checkmarkc\checkmark$\checkmarké\checkmark$\checkmark^\checkmark\\checkmarkc\checkmarki\checkmarkr\checkmarkc\checkmark$\checkmarkc\checkmark$\checkmark^\checkmark\\checkmarkc\checkmarki\checkmarkr\checkmarkc\checkmark$\checkmarku\checkmark$\checkmark^\checkmark\\checkmarkc\checkmarki\checkmarkr\checkmarkc\checkmark$\checkmarkr\checkmark$\checkmark^\checkmark\\checkmarkc\checkmarki\checkmarkr\checkmarkc\checkmark$\checkmarki\checkmark$\checkmark^\checkmark\\checkmarkc\checkmarki\checkmarkr\checkmarkc\checkmark$\checkmarkt\checkmark$\checkmark^\checkmark\\checkmarkc\checkmarki\checkmarkr\checkmarkc\checkmark$\checkmarké\checkmark$\checkmark^\checkmark\\checkmarkc\checkmarki\checkmarkr\checkmarkc\checkmark$\checkmark \checkmark$\checkmark^\checkmark\\checkmarkc\checkmarki\checkmarkr\checkmarkc\checkmark$\checkmarkc\checkmark$\checkmark^\checkmark\\checkmarkc\checkmarki\checkmarkr\checkmarkc\checkmark$\checkmarka\checkmark$\checkmark^\checkmark\\checkmarkc\checkmarki\checkmarkr\checkmarkc\checkmark$\checkmarkl\checkmark$\checkmark^\checkmark\\checkmarkc\checkmarki\checkmarkr\checkmarkc\checkmark$\checkmarkc\checkmark$\checkmark^\checkmark\\checkmarkc\checkmarki\checkmarkr\checkmarkc\checkmark$\checkmarku\checkmark$\checkmark^\checkmark\\checkmarkc\checkmarki\checkmarkr\checkmarkc\checkmark$\checkmarkl\checkmark$\checkmark^\checkmark\\checkmarkc\checkmarki\checkmarkr\checkmarkc\checkmark$\checkmarké\checkmark$\checkmark^\checkmark\\checkmarkc\checkmarki\checkmarkr\checkmarkc\checkmark$\checkmarks\checkmark$\checkmark^\checkmark\\checkmarkc\checkmarki\checkmarkr\checkmarkc\checkmark$\checkmark \checkmark$\checkmark^\checkmark\\checkmarkc\checkmarki\checkmarkr\checkmarkc\checkmark$\checkmarks\checkmark$\checkmark^\checkmark\\checkmarkc\checkmarki\checkmarkr\checkmarkc\checkmark$\checkmarko\checkmark$\checkmark^\checkmark\\checkmarkc\checkmarki\checkmarkr\checkmarkc\checkmark$\checkmarkn\checkmark$\checkmark^\checkmark\\checkmarkc\checkmarki\checkmarkr\checkmarkc\checkmark$\checkmarkt\checkmark$\checkmark^\checkmark\\checkmarkc\checkmarki\checkmarkr\checkmarkc\checkmark$\checkmark \checkmark$\checkmark^\checkmark\\checkmarkc\checkmarki\checkmarkr\checkmarkc\checkmark$\checkmarks\checkmark$\checkmark^\checkmark\\checkmarkc\checkmarki\checkmarkr\checkmarkc\checkmark$\checkmarku\checkmark$\checkmark^\checkmark\\checkmarkc\checkmarki\checkmarkr\checkmarkc\checkmark$\checkmarkp\checkmark$\checkmark^\checkmark\\checkmarkc\checkmarki\checkmarkr\checkmarkc\checkmark$\checkmarké\checkmark$\checkmark^\checkmark\\checkmarkc\checkmarki\checkmarkr\checkmarkc\checkmark$\checkmarkr\checkmark$\checkmark^\checkmark\\checkmarkc\checkmarki\checkmarkr\checkmarkc\checkmark$\checkmarki\checkmark$\checkmark^\checkmark\\checkmarkc\checkmarki\checkmarkr\checkmarkc\checkmark$\checkmarke\checkmark$\checkmark^\checkmark\\checkmarkc\checkmarki\checkmarkr\checkmarkc\checkmark$\checkmarku\checkmark$\checkmark^\checkmark\\checkmarkc\checkmarki\checkmarkr\checkmarkc\checkmark$\checkmarkr\checkmark$\checkmark^\checkmark\\checkmarkc\checkmarki\checkmarkr\checkmarkc\checkmark$\checkmarks\checkmark$\checkmark^\checkmark\\checkmarkc\checkmarki\checkmarkr\checkmarkc\checkmark$\checkmark \checkmark$\checkmark^\checkmark\\checkmarkc\checkmarki\checkmarkr\checkmarkc\checkmark$\checkmarka\checkmark$\checkmark^\checkmark\\checkmarkc\checkmarki\checkmarkr\checkmarkc\checkmark$\checkmarku\checkmark$\checkmark^\checkmark\\checkmarkc\checkmarki\checkmarkr\checkmarkc\checkmark$\checkmarkx\checkmark$\checkmark^\checkmark\\checkmarkc\checkmarki\checkmarkr\checkmarkc\checkmark$\checkmark \checkmark$\checkmark^\checkmark\\checkmarkc\checkmarki\checkmarkr\checkmarkc\checkmark$\checkmarkm\checkmark$\checkmark^\checkmark\\checkmarkc\checkmarki\checkmarkr\checkmarkc\checkmark$\checkmarki\checkmark$\checkmark^\checkmark\\checkmarkc\checkmarki\checkmarkr\checkmarkc\checkmark$\checkmarkn\checkmark$\checkmark^\checkmark\\checkmarkc\checkmarki\checkmarkr\checkmarkc\checkmark$\checkmarki\checkmark$\checkmark^\checkmark\\checkmarkc\checkmarki\checkmarkr\checkmarkc\checkmark$\checkmarkm\checkmark$\checkmark^\checkmark\\checkmarkc\checkmarki\checkmarkr\checkmarkc\checkmark$\checkmarka\checkmark$\checkmark^\checkmark\\checkmarkc\checkmarki\checkmarkr\checkmarkc\checkmark$\checkmark \checkmark$\checkmark^\checkmark\\checkmarkc\checkmarki\checkmarkr\checkmarkc\checkmark$\checkmarki\checkmark$\checkmark^\checkmark\\checkmarkc\checkmarki\checkmarkr\checkmarkc\checkmark$\checkmarkm\checkmark$\checkmark^\checkmark\\checkmarkc\checkmarki\checkmarkr\checkmarkc\checkmark$\checkmarkp\checkmark$\checkmark^\checkmark\\checkmarkc\checkmarki\checkmarkr\checkmarkc\checkmark$\checkmarko\checkmark$\checkmark^\checkmark\\checkmarkc\checkmarki\checkmarkr\checkmarkc\checkmark$\checkmarks\checkmark$\checkmark^\checkmark\\checkmarkc\checkmarki\checkmarkr\checkmarkc\checkmark$\checkmarké\checkmark$\checkmark^\checkmark\\checkmarkc\checkmarki\checkmarkr\checkmarkc\checkmark$\checkmarks\checkmark$\checkmark^\checkmark\\checkmarkc\checkmarki\checkmarkr\checkmarkc\checkmark$\checkmark \checkmark$\checkmark^\checkmark\\checkmarkc\checkmarki\checkmarkr\checkmarkc\checkmark$\checkmarkp\checkmark$\checkmark^\checkmark\\checkmarkc\checkmarki\checkmarkr\checkmarkc\checkmark$\checkmarka\checkmark$\checkmark^\checkmark\\checkmarkc\checkmarki\checkmarkr\checkmarkc\checkmark$\checkmarkr\checkmark$\checkmark^\checkmark\\checkmarkc\checkmarki\checkmarkr\checkmarkc\checkmark$\checkmark \checkmark$\checkmark^\checkmark\\checkmarkc\checkmarki\checkmarkr\checkmarkc\checkmark$\checkmarkl\checkmark$\checkmark^\checkmark\\checkmarkc\checkmarki\checkmarkr\checkmarkc\checkmark$\checkmarke\checkmark$\checkmark^\checkmark\\checkmarkc\checkmarki\checkmarkr\checkmarkc\checkmark$\checkmarks\checkmark$\checkmark^\checkmark\\checkmarkc\checkmarki\checkmarkr\checkmarkc\checkmark$\checkmark \checkmark$\checkmark^\checkmark\\checkmarkc\checkmarki\checkmarkr\checkmarkc\checkmark$\checkmarkn\checkmark$\checkmark^\checkmark\\checkmarkc\checkmarki\checkmarkr\checkmarkc\checkmark$\checkmarko\checkmark$\checkmark^\checkmark\\checkmarkc\checkmarki\checkmarkr\checkmarkc\checkmark$\checkmarkr\checkmark$\checkmark^\checkmark\\checkmarkc\checkmarki\checkmarkr\checkmarkc\checkmark$\checkmarkm\checkmark$\checkmark^\checkmark\\checkmarkc\checkmarki\checkmarkr\checkmarkc\checkmark$\checkmarke\checkmark$\checkmark^\checkmark\\checkmarkc\checkmarki\checkmarkr\checkmarkc\checkmark$\checkmarks\checkmark$\checkmark^\checkmark\\checkmarkc\checkmarki\checkmarkr\checkmarkc\checkmark$\checkmark \checkmark$\checkmark^\checkmark\\checkmarkc\checkmarki\checkmarkr\checkmarkc\checkmark$\checkmarke\checkmark$\checkmark^\checkmark\\checkmarkc\checkmarki\checkmarkr\checkmarkc\checkmark$\checkmarkt\checkmark$\checkmark^\checkmark\\checkmarkc\checkmarki\checkmarkr\checkmarkc\checkmark$\checkmark \checkmark$\checkmark^\checkmark\\checkmarkc\checkmarki\checkmarkr\checkmarkc\checkmark$\checkmarkl\checkmark$\checkmark^\checkmark\\checkmarkc\checkmarki\checkmarkr\checkmarkc\checkmark$\checkmarke\checkmark$\checkmark^\checkmark\\checkmarkc\checkmarki\checkmarkr\checkmarkc\checkmark$\checkmarks\checkmark$\checkmark^\checkmark\\checkmarkc\checkmarki\checkmarkr\checkmarkc\checkmark$\checkmark \checkmark$\checkmark^\checkmark\\checkmarkc\checkmarki\checkmarkr\checkmarkc\checkmark$\checkmarkr\checkmark$\checkmark^\checkmark\\checkmarkc\checkmarki\checkmarkr\checkmarkc\checkmark$\checkmarkè\checkmark$\checkmark^\checkmark\\checkmarkc\checkmarki\checkmarkr\checkmarkc\checkmark$\checkmarkg\checkmark$\checkmark^\checkmark\\checkmarkc\checkmarki\checkmarkr\checkmarkc\checkmark$\checkmarkl\checkmark$\checkmark^\checkmark\\checkmarkc\checkmarki\checkmarkr\checkmarkc\checkmark$\checkmarke\checkmark$\checkmark^\checkmark\\checkmarkc\checkmarki\checkmarkr\checkmarkc\checkmark$\checkmarks\checkmark$\checkmark^\checkmark\\checkmarkc\checkmarki\checkmarkr\checkmarkc\checkmark$\checkmark \checkmark$\checkmark^\checkmark\\checkmarkc\checkmarki\checkmarkr\checkmarkc\checkmark$\checkmarkd\checkmark$\checkmark^\checkmark\\checkmarkc\checkmarki\checkmarkr\checkmarkc\checkmark$\checkmarke\checkmark$\checkmark^\checkmark\\checkmarkc\checkmarki\checkmarkr\checkmarkc\checkmark$\checkmark \checkmark$\checkmark^\checkmark\\checkmarkc\checkmarki\checkmarkr\checkmarkc\checkmark$\checkmarkl\checkmark$\checkmark^\checkmark\\checkmarkc\checkmarki\checkmarkr\checkmarkc\checkmark$\checkmark'\checkmark$\checkmark^\checkmark\\checkmarkc\checkmarki\checkmarkr\checkmarkc\checkmark$\checkmarka\checkmark$\checkmark^\checkmark\\checkmarkc\checkmarki\checkmarkr\checkmarkc\checkmark$\checkmarkr\checkmark$\checkmark^\checkmark\\checkmarkc\checkmarki\checkmarkr\checkmarkc\checkmark$\checkmarkt\checkmark$\checkmark^\checkmark\\checkmarkc\checkmarki\checkmarkr\checkmarkc\checkmark$\checkmark \checkmark$\checkmark^\checkmark\\checkmarkc\checkmarki\checkmarkr\checkmarkc\checkmark$\checkmarke\checkmark$\checkmark^\checkmark\\checkmarkc\checkmarki\checkmarkr\checkmarkc\checkmark$\checkmarkn\checkmark$\checkmark^\checkmark\\checkmarkc\checkmarki\checkmarkr\checkmarkc\checkmark$\checkmark \checkmark$\checkmark^\checkmark\\checkmarkc\checkmarki\checkmarkr\checkmarkc\checkmark$\checkmarkc\checkmark$\checkmark^\checkmark\\checkmarkc\checkmarki\checkmarkr\checkmarkc\checkmark$\checkmarko\checkmark$\checkmark^\checkmark\\checkmarkc\checkmarki\checkmarkr\checkmarkc\checkmark$\checkmarkn\checkmark$\checkmark^\checkmark\\checkmarkc\checkmarki\checkmarkr\checkmarkc\checkmark$\checkmarkc\checkmark$\checkmark^\checkmark\\checkmarkc\checkmarki\checkmarkr\checkmarkc\checkmark$\checkmarke\checkmark$\checkmark^\checkmark\\checkmarkc\checkmarki\checkmarkr\checkmarkc\checkmark$\checkmarkp\checkmark$\checkmark^\checkmark\\checkmarkc\checkmarki\checkmarkr\checkmarkc\checkmark$\checkmarkt\checkmark$\checkmark^\checkmark\\checkmarkc\checkmarki\checkmarkr\checkmarkc\checkmark$\checkmarki\checkmark$\checkmark^\checkmark\\checkmarkc\checkmarki\checkmarkr\checkmarkc\checkmark$\checkmarko\checkmark$\checkmark^\checkmark\\checkmarkc\checkmarki\checkmarkr\checkmarkc\checkmark$\checkmarkn\checkmark$\checkmark^\checkmark\\checkmarkc\checkmarki\checkmarkr\checkmarkc\checkmark$\checkmark \checkmark$\checkmark^\checkmark\\checkmarkc\checkmarki\checkmarkr\checkmarkc\checkmark$\checkmarkm\checkmark$\checkmark^\checkmark\\checkmarkc\checkmarki\checkmarkr\checkmarkc\checkmark$\checkmarké\checkmark$\checkmark^\checkmark\\checkmarkc\checkmarki\checkmarkr\checkmarkc\checkmark$\checkmarkc\checkmark$\checkmark^\checkmark\\checkmarkc\checkmarki\checkmarkr\checkmarkc\checkmark$\checkmarka\checkmark$\checkmark^\checkmark\\checkmarkc\checkmarki\checkmarkr\checkmarkc\checkmark$\checkmarkn\checkmark$\checkmark^\checkmark\\checkmarkc\checkmarki\checkmarkr\checkmarkc\checkmark$\checkmarki\checkmark$\checkmark^\checkmark\\checkmarkc\checkmarki\checkmarkr\checkmarkc\checkmark$\checkmarkq\checkmark$\checkmark^\checkmark\\checkmarkc\checkmarki\checkmarkr\checkmarkc\checkmark$\checkmarku\checkmark$\checkmark^\checkmark\\checkmarkc\checkmarki\checkmarkr\checkmarkc\checkmark$\checkmarke\checkmark$\checkmark^\checkmark\\checkmarkc\checkmarki\checkmarkr\checkmarkc\checkmark$\checkmark.\checkmark$\checkmark^\checkmark\\checkmarkc\checkmarki\checkmarkr\checkmarkc\checkmark$\checkmark
\checkmark$\checkmark^\checkmark\\checkmarkc\checkmarki\checkmarkr\checkmarkc\checkmark$\checkmark
\checkmark$\checkmark^\checkmark\\checkmarkc\checkmarki\checkmarkr\checkmarkc\checkmark$\checkmarkL\checkmark$\checkmark^\checkmark\\checkmarkc\checkmarki\checkmarkr\checkmarkc\checkmark$\checkmark'\checkmark$\checkmark^\checkmark\\checkmarkc\checkmarki\checkmarkr\checkmarkc\checkmark$\checkmarka\checkmark$\checkmark^\checkmark\\checkmarkc\checkmarki\checkmarkr\checkmarkc\checkmark$\checkmarkr\checkmark$\checkmark^\checkmark\\checkmarkc\checkmarki\checkmarkr\checkmarkc\checkmark$\checkmarkc\checkmark$\checkmark^\checkmark\\checkmarkc\checkmarki\checkmarkr\checkmarkc\checkmark$\checkmarkh\checkmark$\checkmark^\checkmark\\checkmarkc\checkmarki\checkmarkr\checkmarkc\checkmark$\checkmarki\checkmark$\checkmark^\checkmark\\checkmarkc\checkmarki\checkmarkr\checkmarkc\checkmark$\checkmarkt\checkmark$\checkmark^\checkmark\\checkmarkc\checkmarki\checkmarkr\checkmarkc\checkmark$\checkmarke\checkmark$\checkmark^\checkmark\\checkmarkc\checkmarki\checkmarkr\checkmarkc\checkmark$\checkmarkc\checkmark$\checkmark^\checkmark\\checkmarkc\checkmarki\checkmarkr\checkmarkc\checkmark$\checkmarkt\checkmark$\checkmark^\checkmark\\checkmarkc\checkmarki\checkmarkr\checkmarkc\checkmark$\checkmarku\checkmark$\checkmark^\checkmark\\checkmarkc\checkmarki\checkmarkr\checkmarkc\checkmark$\checkmarkr\checkmark$\checkmark^\checkmark\\checkmarkc\checkmarki\checkmarkr\checkmarkc\checkmark$\checkmarke\checkmark$\checkmark^\checkmark\\checkmarkc\checkmarki\checkmarkr\checkmarkc\checkmark$\checkmark \checkmark$\checkmark^\checkmark\\checkmarkc\checkmarki\checkmarkr\checkmarkc\checkmark$\checkmark«\checkmark$\checkmark^\checkmark\\checkmarkc\checkmarki\checkmarkr\checkmarkc\checkmark$\checkmarkT\checkmark$\checkmark^\checkmark\\checkmarkc\checkmarki\checkmarkr\checkmarkc\checkmark$\checkmarka\checkmark$\checkmark^\checkmark\\checkmarkc\checkmarki\checkmarkr\checkmarkc\checkmark$\checkmarkb\checkmark$\checkmark^\checkmark\\checkmarkc\checkmarki\checkmarkr\checkmarkc\checkmark$\checkmarkl\checkmark$\checkmark^\checkmark\\checkmarkc\checkmarki\checkmarkr\checkmarkc\checkmark$\checkmarke\checkmark$\checkmark^\checkmark\\checkmarkc\checkmarki\checkmarkr\checkmarkc\checkmark$\checkmark-\checkmark$\checkmark^\checkmark\\checkmarkc\checkmarki\checkmarkr\checkmarkc\checkmark$\checkmarkT\checkmark$\checkmark^\checkmark\\checkmarkc\checkmarki\checkmarkr\checkmarkc\checkmark$\checkmarka\checkmark$\checkmark^\checkmark\\checkmarkc\checkmarki\checkmarkr\checkmarkc\checkmark$\checkmarkb\checkmark$\checkmark^\checkmark\\checkmarkc\checkmarki\checkmarkr\checkmarkc\checkmark$\checkmarkl\checkmark$\checkmark^\checkmark\\checkmarkc\checkmarki\checkmarkr\checkmarkc\checkmark$\checkmarke\checkmark$\checkmark^\checkmark\\checkmarkc\checkmarki\checkmarkr\checkmarkc\checkmark$\checkmark»\checkmark$\checkmark^\checkmark\\checkmarkc\checkmarki\checkmarkr\checkmarkc\checkmark$\checkmark \checkmark$\checkmark^\checkmark\\checkmarkc\checkmarki\checkmarkr\checkmarkc\checkmark$\checkmark(\checkmark$\checkmark^\checkmark\\checkmarkc\checkmarki\checkmarkr\checkmarkc\checkmark$\checkmarkT\checkmark$\checkmark^\checkmark\\checkmarkc\checkmarki\checkmarkr\checkmarkc\checkmark$\checkmarkr\checkmark$\checkmark^\checkmark\\checkmarkc\checkmarki\checkmarkr\checkmarkc\checkmark$\checkmarku\checkmark$\checkmark^\checkmark\\checkmarkc\checkmarki\checkmarkr\checkmarkc\checkmark$\checkmarkn\checkmark$\checkmark^\checkmark\\checkmarkc\checkmarki\checkmarkr\checkmarkc\checkmark$\checkmarkn\checkmark$\checkmark^\checkmark\\checkmarkc\checkmarki\checkmarkr\checkmarkc\checkmark$\checkmarki\checkmark$\checkmark^\checkmark\\checkmarkc\checkmarki\checkmarkr\checkmarkc\checkmark$\checkmarko\checkmark$\checkmark^\checkmark\\checkmarkc\checkmarki\checkmarkr\checkmarkc\checkmark$\checkmarkn\checkmark$\checkmark^\checkmark\\checkmarkc\checkmarki\checkmarkr\checkmarkc\checkmark$\checkmark)\checkmark$\checkmark^\checkmark\\checkmarkc\checkmarki\checkmarkr\checkmarkc\checkmark$\checkmark \checkmark$\checkmark^\checkmark\\checkmarkc\checkmarki\checkmarkr\checkmarkc\checkmark$\checkmarkg\checkmark$\checkmark^\checkmark\\checkmarkc\checkmarki\checkmarkr\checkmarkc\checkmark$\checkmarké\checkmark$\checkmark^\checkmark\\checkmarkc\checkmarki\checkmarkr\checkmarkc\checkmark$\checkmarkn\checkmark$\checkmark^\checkmark\\checkmarkc\checkmarki\checkmarkr\checkmarkc\checkmark$\checkmarkè\checkmark$\checkmark^\checkmark\\checkmarkc\checkmarki\checkmarkr\checkmarkc\checkmark$\checkmarkr\checkmark$\checkmark^\checkmark\\checkmarkc\checkmarki\checkmarkr\checkmarkc\checkmark$\checkmarke\checkmark$\checkmark^\checkmark\\checkmarkc\checkmarki\checkmarkr\checkmarkc\checkmark$\checkmark \checkmark$\checkmark^\checkmark\\checkmarkc\checkmarki\checkmarkr\checkmarkc\checkmark$\checkmarkc\checkmark$\checkmark^\checkmark\\checkmarkc\checkmarki\checkmarkr\checkmarkc\checkmark$\checkmarke\checkmark$\checkmark^\checkmark\\checkmarkc\checkmarki\checkmarkr\checkmarkc\checkmark$\checkmarkr\checkmark$\checkmark^\checkmark\\checkmarkc\checkmarki\checkmarkr\checkmarkc\checkmark$\checkmarkt\checkmark$\checkmark^\checkmark\\checkmarkc\checkmarki\checkmarkr\checkmarkc\checkmark$\checkmarke\checkmark$\checkmark^\checkmark\\checkmarkc\checkmarki\checkmarkr\checkmarkc\checkmark$\checkmarks\checkmark$\checkmark^\checkmark\\checkmarkc\checkmarki\checkmarkr\checkmarkc\checkmark$\checkmark \checkmark$\checkmark^\checkmark\\checkmarkc\checkmarki\checkmarkr\checkmarkc\checkmark$\checkmarku\checkmark$\checkmark^\checkmark\\checkmarkc\checkmarki\checkmarkr\checkmarkc\checkmark$\checkmarkn\checkmark$\checkmark^\checkmark\\checkmarkc\checkmarki\checkmarkr\checkmarkc\checkmark$\checkmark \checkmark$\checkmark^\checkmark\\checkmarkc\checkmarki\checkmarkr\checkmarkc\checkmark$\checkmarkm\checkmark$\checkmark^\checkmark\\checkmarkc\checkmarki\checkmarkr\checkmarkc\checkmark$\checkmarko\checkmark$\checkmark^\checkmark\\checkmarkc\checkmarki\checkmarkr\checkmarkc\checkmark$\checkmarkm\checkmark$\checkmark^\checkmark\\checkmarkc\checkmarki\checkmarkr\checkmarkc\checkmark$\checkmarke\checkmark$\checkmark^\checkmark\\checkmarkc\checkmarki\checkmarkr\checkmarkc\checkmark$\checkmarkn\checkmark$\checkmark^\checkmark\\checkmarkc\checkmarki\checkmarkr\checkmarkc\checkmark$\checkmarkt\checkmark$\checkmark^\checkmark\\checkmarkc\checkmarki\checkmarkr\checkmarkc\checkmark$\checkmark \checkmark$\checkmark^\checkmark\\checkmarkc\checkmarki\checkmarkr\checkmarkc\checkmark$\checkmarkd\checkmark$\checkmark^\checkmark\\checkmarkc\checkmarki\checkmarkr\checkmarkc\checkmark$\checkmarke\checkmark$\checkmark^\checkmark\\checkmarkc\checkmarki\checkmarkr\checkmarkc\checkmark$\checkmark \checkmark$\checkmark^\checkmark\\checkmarkc\checkmarki\checkmarkr\checkmarkc\checkmark$\checkmarkr\checkmark$\checkmark^\checkmark\\checkmarkc\checkmarki\checkmarkr\checkmarkc\checkmark$\checkmarke\checkmark$\checkmark^\checkmark\\checkmarkc\checkmarki\checkmarkr\checkmarkc\checkmark$\checkmarkn\checkmark$\checkmark^\checkmark\\checkmarkc\checkmarki\checkmarkr\checkmarkc\checkmark$\checkmarkv\checkmark$\checkmark^\checkmark\\checkmarkc\checkmarki\checkmarkr\checkmarkc\checkmark$\checkmarke\checkmark$\checkmark^\checkmark\\checkmarkc\checkmarki\checkmarkr\checkmarkc\checkmark$\checkmarkr\checkmark$\checkmark^\checkmark\\checkmarkc\checkmarki\checkmarkr\checkmarkc\checkmark$\checkmarks\checkmark$\checkmark^\checkmark\\checkmarkc\checkmarki\checkmarkr\checkmarkc\checkmark$\checkmarke\checkmark$\checkmark^\checkmark\\checkmarkc\checkmarki\checkmarkr\checkmarkc\checkmark$\checkmarkm\checkmark$\checkmark^\checkmark\\checkmarkc\checkmarki\checkmarkr\checkmarkc\checkmark$\checkmarke\checkmark$\checkmark^\checkmark\\checkmarkc\checkmarki\checkmarkr\checkmarkc\checkmark$\checkmarkn\checkmark$\checkmark^\checkmark\\checkmarkc\checkmarki\checkmarkr\checkmarkc\checkmark$\checkmarkt\checkmark$\checkmark^\checkmark\\checkmarkc\checkmarki\checkmarkr\checkmarkc\checkmark$\checkmark \checkmark$\checkmark^\checkmark\\checkmarkc\checkmarki\checkmarkr\checkmarkc\checkmark$\checkmarkd\checkmark$\checkmark^\checkmark\\checkmarkc\checkmarki\checkmarkr\checkmarkc\checkmark$\checkmarke\checkmark$\checkmark^\checkmark\\checkmarkc\checkmarki\checkmarkr\checkmarkc\checkmark$\checkmark \checkmark$\checkmark^\checkmark\\checkmarkc\checkmarki\checkmarkr\checkmarkc\checkmark$\checkmark$\checkmark$\checkmark^\checkmark\\checkmarkc\checkmarki\checkmarkr\checkmarkc\checkmark$\checkmark2\checkmark$\checkmark^\checkmark\\checkmarkc\checkmarki\checkmarkr\checkmarkc\checkmark$\checkmark0\checkmark$\checkmark^\checkmark\\checkmarkc\checkmarki\checkmarkr\checkmarkc\checkmark$\checkmark$\checkmark$\checkmark^\checkmark\\checkmarkc\checkmarki\checkmarkr\checkmarkc\checkmark$\checkmark \checkmark$\checkmark^\checkmark\\checkmarkc\checkmarki\checkmarkr\checkmarkc\checkmark$\checkmarkN\checkmark$\checkmark^\checkmark\\checkmarkc\checkmarki\checkmarkr\checkmarkc\checkmark$\checkmark$\checkmark$\checkmark^\checkmark\\checkmarkc\checkmarki\checkmarkr\checkmarkc\checkmark$\checkmark\\checkmark$\checkmark^\checkmark\\checkmarkc\checkmarki\checkmarkr\checkmarkc\checkmark$\checkmarkc\checkmark$\checkmark^\checkmark\\checkmarkc\checkmarki\checkmarkr\checkmarkc\checkmark$\checkmarkd\checkmark$\checkmark^\checkmark\\checkmarkc\checkmarki\checkmarkr\checkmarkc\checkmark$\checkmarko\checkmark$\checkmark^\checkmark\\checkmarkc\checkmarki\checkmarkr\checkmarkc\checkmark$\checkmarkt\checkmark$\checkmark^\checkmark\\checkmarkc\checkmarki\checkmarkr\checkmarkc\checkmark$\checkmark$\checkmark$\checkmark^\checkmark\\checkmarkc\checkmarki\checkmarkr\checkmarkc\checkmark$\checkmarkm\checkmark$\checkmark^\checkmark\\checkmarkc\checkmarki\checkmarkr\checkmarkc\checkmark$\checkmark \checkmark$\checkmark^\checkmark\\checkmarkc\checkmarki\checkmarkr\checkmarkc\checkmark$\checkmarks\checkmark$\checkmark^\checkmark\\checkmarkc\checkmarki\checkmarkr\checkmarkc\checkmark$\checkmarku\checkmark$\checkmark^\checkmark\\checkmarkc\checkmarki\checkmarkr\checkmarkc\checkmark$\checkmarkr\checkmark$\checkmark^\checkmark\\checkmarkc\checkmarki\checkmarkr\checkmarkc\checkmark$\checkmark \checkmark$\checkmark^\checkmark\\checkmarkc\checkmarki\checkmarkr\checkmarkc\checkmark$\checkmarkl\checkmark$\checkmark^\checkmark\\checkmarkc\checkmarki\checkmarkr\checkmarkc\checkmark$\checkmarke\checkmark$\checkmark^\checkmark\\checkmarkc\checkmarki\checkmarkr\checkmarkc\checkmark$\checkmark \checkmark$\checkmark^\checkmark\\checkmarkc\checkmarki\checkmarkr\checkmarkc\checkmark$\checkmarkc\checkmark$\checkmark^\checkmark\\checkmarkc\checkmarki\checkmarkr\checkmarkc\checkmark$\checkmarkh\checkmark$\checkmark^\checkmark\\checkmarkc\checkmarki\checkmarkr\checkmarkc\checkmark$\checkmarka\checkmark$\checkmark^\checkmark\\checkmarkc\checkmarki\checkmarkr\checkmarkc\checkmark$\checkmarkr\checkmark$\checkmark^\checkmark\\checkmarkc\checkmarki\checkmarkr\checkmarkc\checkmark$\checkmarki\checkmark$\checkmark^\checkmark\\checkmarkc\checkmarki\checkmarkr\checkmarkc\checkmark$\checkmarko\checkmark$\checkmark^\checkmark\\checkmarkc\checkmarki\checkmarkr\checkmarkc\checkmark$\checkmarkt\checkmark$\checkmark^\checkmark\\checkmarkc\checkmarki\checkmarkr\checkmarkc\checkmark$\checkmark \checkmark$\checkmark^\checkmark\\checkmarkc\checkmarki\checkmarkr\checkmarkc\checkmark$\checkmarkY\checkmark$\checkmark^\checkmark\\checkmarkc\checkmarki\checkmarkr\checkmarkc\checkmark$\checkmark,\checkmark$\checkmark^\checkmark\\checkmarkc\checkmarki\checkmarkr\checkmarkc\checkmark$\checkmark \checkmark$\checkmark^\checkmark\\checkmarkc\checkmarki\checkmarkr\checkmarkc\checkmark$\checkmarkm\checkmark$\checkmark^\checkmark\\checkmarkc\checkmarki\checkmarkr\checkmarkc\checkmark$\checkmarka\checkmark$\checkmark^\checkmark\\checkmarkc\checkmarki\checkmarkr\checkmarkc\checkmark$\checkmarki\checkmark$\checkmark^\checkmark\\checkmarkc\checkmarki\checkmarkr\checkmarkc\checkmark$\checkmarks\checkmark$\checkmark^\checkmark\\checkmarkc\checkmarki\checkmarkr\checkmarkc\checkmark$\checkmark \checkmark$\checkmark^\checkmark\\checkmarkc\checkmarki\checkmarkr\checkmarkc\checkmark$\checkmarkl\checkmark$\checkmark^\checkmark\\checkmarkc\checkmarki\checkmarkr\checkmarkc\checkmark$\checkmark'\checkmark$\checkmark^\checkmark\\checkmarkc\checkmarki\checkmarkr\checkmarkc\checkmark$\checkmarku\checkmark$\checkmark^\checkmark\\checkmarkc\checkmarki\checkmarkr\checkmarkc\checkmark$\checkmarkt\checkmark$\checkmark^\checkmark\\checkmarkc\checkmarki\checkmarkr\checkmarkc\checkmark$\checkmarki\checkmark$\checkmark^\checkmark\\checkmarkc\checkmarki\checkmarkr\checkmarkc\checkmark$\checkmarkl\checkmark$\checkmark^\checkmark\\checkmarkc\checkmarki\checkmarkr\checkmarkc\checkmark$\checkmarki\checkmark$\checkmark^\checkmark\\checkmarkc\checkmarki\checkmarkr\checkmarkc\checkmark$\checkmarks\checkmark$\checkmark^\checkmark\\checkmarkc\checkmarki\checkmarkr\checkmarkc\checkmark$\checkmarka\checkmark$\checkmark^\checkmark\\checkmarkc\checkmarki\checkmarkr\checkmarkc\checkmark$\checkmarkt\checkmark$\checkmark^\checkmark\\checkmarkc\checkmarki\checkmarkr\checkmarkc\checkmark$\checkmarki\checkmark$\checkmark^\checkmark\\checkmarkc\checkmarki\checkmarkr\checkmarkc\checkmark$\checkmarko\checkmark$\checkmark^\checkmark\\checkmarkc\checkmarki\checkmarkr\checkmarkc\checkmark$\checkmarkn\checkmark$\checkmark^\checkmark\\checkmarkc\checkmarki\checkmarkr\checkmarkc\checkmark$\checkmark \checkmark$\checkmark^\checkmark\\checkmarkc\checkmarki\checkmarkr\checkmarkc\checkmark$\checkmarkd\checkmark$\checkmark^\checkmark\\checkmarkc\checkmarki\checkmarkr\checkmarkc\checkmark$\checkmarke\checkmark$\checkmark^\checkmark\\checkmarkc\checkmarki\checkmarkr\checkmarkc\checkmark$\checkmark \checkmark$\checkmark^\checkmark\\checkmarkc\checkmarki\checkmarkr\checkmarkc\checkmark$\checkmarkr\checkmark$\checkmark^\checkmark\\checkmarkc\checkmarki\checkmarkr\checkmarkc\checkmark$\checkmarka\checkmark$\checkmark^\checkmark\\checkmarkc\checkmarki\checkmarkr\checkmarkc\checkmark$\checkmarki\checkmark$\checkmark^\checkmark\\checkmarkc\checkmarki\checkmarkr\checkmarkc\checkmark$\checkmarkl\checkmark$\checkmark^\checkmark\\checkmarkc\checkmarki\checkmarkr\checkmarkc\checkmark$\checkmarks\checkmark$\checkmark^\checkmark\\checkmarkc\checkmarki\checkmarkr\checkmarkc\checkmark$\checkmark \checkmark$\checkmark^\checkmark\\checkmarkc\checkmarki\checkmarkr\checkmarkc\checkmark$\checkmarkp\checkmark$\checkmark^\checkmark\\checkmarkc\checkmarki\checkmarkr\checkmarkc\checkmark$\checkmarkr\checkmark$\checkmark^\checkmark\\checkmarkc\checkmarki\checkmarkr\checkmarkc\checkmark$\checkmarko\checkmark$\checkmark^\checkmark\\checkmarkc\checkmarki\checkmarkr\checkmarkc\checkmark$\checkmarkf\checkmark$\checkmark^\checkmark\\checkmarkc\checkmarki\checkmarkr\checkmarkc\checkmark$\checkmarki\checkmark$\checkmark^\checkmark\\checkmarkc\checkmarki\checkmarkr\checkmarkc\checkmark$\checkmarkl\checkmark$\checkmark^\checkmark\\checkmarkc\checkmarki\checkmarkr\checkmarkc\checkmark$\checkmarké\checkmark$\checkmark^\checkmark\\checkmarkc\checkmarki\checkmarkr\checkmarkc\checkmark$\checkmarks\checkmark$\checkmark^\checkmark\\checkmarkc\checkmarki\checkmarkr\checkmarkc\checkmark$\checkmark \checkmark$\checkmark^\checkmark\\checkmarkc\checkmarki\checkmarkr\checkmarkc\checkmark$\checkmarkH\checkmark$\checkmark^\checkmark\\checkmarkc\checkmarki\checkmarkr\checkmarkc\checkmark$\checkmarkI\checkmark$\checkmark^\checkmark\\checkmarkc\checkmarki\checkmarkr\checkmarkc\checkmark$\checkmarkW\checkmark$\checkmark^\checkmark\\checkmarkc\checkmarki\checkmarkr\checkmarkc\checkmark$\checkmarkI\checkmark$\checkmark^\checkmark\\checkmarkc\checkmarki\checkmarkr\checkmarkc\checkmark$\checkmarkN\checkmark$\checkmark^\checkmark\\checkmarkc\checkmarki\checkmarkr\checkmarkc\checkmark$\checkmark \checkmark$\checkmark^\checkmark\\checkmarkc\checkmarki\checkmarkr\checkmarkc\checkmark$\checkmarkH\checkmark$\checkmark^\checkmark\\checkmarkc\checkmarki\checkmarkr\checkmarkc\checkmark$\checkmarkG\checkmark$\checkmark^\checkmark\\checkmarkc\checkmarki\checkmarkr\checkmarkc\checkmark$\checkmarkR\checkmark$\checkmark^\checkmark\\checkmarkc\checkmarki\checkmarkr\checkmarkc\checkmark$\checkmark1\checkmark$\checkmark^\checkmark\\checkmarkc\checkmarki\checkmarkr\checkmarkc\checkmark$\checkmark5\checkmark$\checkmark^\checkmark\\checkmarkc\checkmarki\checkmarkr\checkmarkc\checkmark$\checkmark \checkmark$\checkmark^\checkmark\\checkmarkc\checkmarki\checkmarkr\checkmarkc\checkmark$\checkmark(\checkmark$\checkmark^\checkmark\\checkmarkc\checkmarki\checkmarkr\checkmarkc\checkmark$\checkmarkm\checkmark$\checkmark^\checkmark\\checkmarkc\checkmarki\checkmarkr\checkmarkc\checkmark$\checkmarko\checkmark$\checkmark^\checkmark\\checkmarkc\checkmarki\checkmarkr\checkmarkc\checkmark$\checkmarkm\checkmark$\checkmark^\checkmark\\checkmarkc\checkmarki\checkmarkr\checkmarkc\checkmark$\checkmarke\checkmark$\checkmark^\checkmark\\checkmarkc\checkmarki\checkmarkr\checkmarkc\checkmark$\checkmarkn\checkmark$\checkmark^\checkmark\\checkmarkc\checkmarki\checkmarkr\checkmarkc\checkmark$\checkmarkt\checkmark$\checkmark^\checkmark\\checkmarkc\checkmarki\checkmarkr\checkmarkc\checkmark$\checkmark \checkmark$\checkmark^\checkmark\\checkmarkc\checkmarki\checkmarkr\checkmarkc\checkmark$\checkmarka\checkmark$\checkmark^\checkmark\\checkmarkc\checkmarki\checkmarkr\checkmarkc\checkmark$\checkmarkd\checkmark$\checkmark^\checkmark\\checkmarkc\checkmarki\checkmarkr\checkmarkc\checkmark$\checkmarkm\checkmark$\checkmark^\checkmark\\checkmarkc\checkmarki\checkmarkr\checkmarkc\checkmark$\checkmarki\checkmark$\checkmark^\checkmark\\checkmarkc\checkmarki\checkmarkr\checkmarkc\checkmark$\checkmarks\checkmark$\checkmark^\checkmark\\checkmarkc\checkmarki\checkmarkr\checkmarkc\checkmark$\checkmarks\checkmark$\checkmark^\checkmark\\checkmarkc\checkmarki\checkmarkr\checkmarkc\checkmark$\checkmarki\checkmark$\checkmark^\checkmark\\checkmarkc\checkmarki\checkmarkr\checkmarkc\checkmark$\checkmarkb\checkmark$\checkmark^\checkmark\\checkmarkc\checkmarki\checkmarkr\checkmarkc\checkmark$\checkmarkl\checkmark$\checkmark^\checkmark\\checkmarkc\checkmarki\checkmarkr\checkmarkc\checkmark$\checkmarke\checkmark$\checkmark^\checkmark\\checkmarkc\checkmarki\checkmarkr\checkmarkc\checkmark$\checkmark \checkmark$\checkmark^\checkmark\\checkmarkc\checkmarki\checkmarkr\checkmarkc\checkmark$\checkmark$\checkmark$\checkmark^\checkmark\\checkmarkc\checkmarki\checkmarkr\checkmarkc\checkmark$\checkmark=\checkmark$\checkmark^\checkmark\\checkmarkc\checkmarki\checkmarkr\checkmarkc\checkmark$\checkmark \checkmark$\checkmark^\checkmark\\checkmarkc\checkmarki\checkmarkr\checkmarkc\checkmark$\checkmark9\checkmark$\checkmark^\checkmark\\checkmarkc\checkmarki\checkmarkr\checkmarkc\checkmark$\checkmark0\checkmark$\checkmark^\checkmark\\checkmarkc\checkmarki\checkmarkr\checkmarkc\checkmark$\checkmark$\checkmark$\checkmark^\checkmark\\checkmarkc\checkmarki\checkmarkr\checkmarkc\checkmark$\checkmark \checkmark$\checkmark^\checkmark\\checkmarkc\checkmarki\checkmarkr\checkmarkc\checkmark$\checkmarkN\checkmark$\checkmark^\checkmark\\checkmarkc\checkmarki\checkmarkr\checkmarkc\checkmark$\checkmark$\checkmark$\checkmark^\checkmark\\checkmarkc\checkmarki\checkmarkr\checkmarkc\checkmark$\checkmark\\checkmark$\checkmark^\checkmark\\checkmarkc\checkmarki\checkmarkr\checkmarkc\checkmark$\checkmarkc\checkmark$\checkmark^\checkmark\\checkmarkc\checkmarki\checkmarkr\checkmarkc\checkmark$\checkmarkd\checkmark$\checkmark^\checkmark\\checkmarkc\checkmarki\checkmarkr\checkmarkc\checkmark$\checkmarko\checkmark$\checkmark^\checkmark\\checkmarkc\checkmarki\checkmarkr\checkmarkc\checkmark$\checkmarkt\checkmark$\checkmark^\checkmark\\checkmarkc\checkmarki\checkmarkr\checkmarkc\checkmark$\checkmark$\checkmark$\checkmark^\checkmark\\checkmarkc\checkmarki\checkmarkr\checkmarkc\checkmark$\checkmarkm\checkmark$\checkmark^\checkmark\\checkmarkc\checkmarki\checkmarkr\checkmarkc\checkmark$\checkmark)\checkmark$\checkmark^\checkmark\\checkmarkc\checkmarki\checkmarkr\checkmarkc\checkmark$\checkmark \checkmark$\checkmark^\checkmark\\checkmarkc\checkmarki\checkmarkr\checkmarkc\checkmark$\checkmarka\checkmark$\checkmark^\checkmark\\checkmarkc\checkmarki\checkmarkr\checkmarkc\checkmark$\checkmarkv\checkmark$\checkmark^\checkmark\\checkmarkc\checkmarki\checkmarkr\checkmarkc\checkmark$\checkmarke\checkmark$\checkmark^\checkmark\\checkmarkc\checkmarki\checkmarkr\checkmarkc\checkmark$\checkmarkc\checkmark$\checkmark^\checkmark\\checkmarkc\checkmarki\checkmarkr\checkmarkc\checkmark$\checkmark \checkmark$\checkmark^\checkmark\\checkmarkc\checkmarki\checkmarkr\checkmarkc\checkmark$\checkmarku\checkmark$\checkmark^\checkmark\\checkmarkc\checkmarki\checkmarkr\checkmarkc\checkmark$\checkmarkn\checkmark$\checkmark^\checkmark\\checkmarkc\checkmarki\checkmarkr\checkmarkc\checkmark$\checkmark \checkmark$\checkmark^\checkmark\\checkmarkc\checkmarki\checkmarkr\checkmarkc\checkmark$\checkmarkc\checkmark$\checkmark^\checkmark\\checkmarkc\checkmarki\checkmarkr\checkmarkc\checkmark$\checkmarko\checkmark$\checkmark^\checkmark\\checkmarkc\checkmarki\checkmarkr\checkmarkc\checkmark$\checkmarke\checkmark$\checkmark^\checkmark\\checkmarkc\checkmarki\checkmarkr\checkmarkc\checkmark$\checkmarkf\checkmark$\checkmark^\checkmark\\checkmarkc\checkmarki\checkmarkr\checkmarkc\checkmark$\checkmarkf\checkmark$\checkmark^\checkmark\\checkmarkc\checkmarki\checkmarkr\checkmarkc\checkmark$\checkmarki\checkmark$\checkmark^\checkmark\\checkmarkc\checkmarki\checkmarkr\checkmarkc\checkmark$\checkmarkc\checkmark$\checkmark^\checkmark\\checkmarkc\checkmarki\checkmarkr\checkmarkc\checkmark$\checkmarki\checkmark$\checkmark^\checkmark\\checkmarkc\checkmarki\checkmarkr\checkmarkc\checkmark$\checkmarke\checkmark$\checkmark^\checkmark\\checkmarkc\checkmarki\checkmarkr\checkmarkc\checkmark$\checkmarkn\checkmark$\checkmark^\checkmark\\checkmarkc\checkmarki\checkmarkr\checkmarkc\checkmark$\checkmarkt\checkmark$\checkmark^\checkmark\\checkmarkc\checkmarki\checkmarkr\checkmarkc\checkmark$\checkmark \checkmark$\checkmark^\checkmark\\checkmarkc\checkmarki\checkmarkr\checkmarkc\checkmark$\checkmarkd\checkmark$\checkmark^\checkmark\\checkmarkc\checkmarki\checkmarkr\checkmarkc\checkmark$\checkmarke\checkmark$\checkmark^\checkmark\\checkmarkc\checkmarki\checkmarkr\checkmarkc\checkmark$\checkmark \checkmark$\checkmark^\checkmark\\checkmarkc\checkmarki\checkmarkr\checkmarkc\checkmark$\checkmarks\checkmark$\checkmark^\checkmark\\checkmarkc\checkmarki\checkmarkr\checkmarkc\checkmark$\checkmarké\checkmark$\checkmark^\checkmark\\checkmarkc\checkmarki\checkmarkr\checkmarkc\checkmark$\checkmarkc\checkmark$\checkmark^\checkmark\\checkmarkc\checkmarki\checkmarkr\checkmarkc\checkmark$\checkmarku\checkmark$\checkmark^\checkmark\\checkmarkc\checkmarki\checkmarkr\checkmarkc\checkmark$\checkmarkr\checkmark$\checkmark^\checkmark\\checkmarkc\checkmarki\checkmarkr\checkmarkc\checkmark$\checkmarki\checkmark$\checkmark^\checkmark\\checkmarkc\checkmarki\checkmarkr\checkmarkc\checkmark$\checkmarkt\checkmark$\checkmark^\checkmark\\checkmarkc\checkmarki\checkmarkr\checkmarkc\checkmark$\checkmarké\checkmark$\checkmark^\checkmark\\checkmarkc\checkmarki\checkmarkr\checkmarkc\checkmark$\checkmark \checkmark$\checkmark^\checkmark\\checkmarkc\checkmarki\checkmarkr\checkmarkc\checkmark$\checkmarkd\checkmark$\checkmark^\checkmark\\checkmarkc\checkmarki\checkmarkr\checkmarkc\checkmark$\checkmarke\checkmark$\checkmark^\checkmark\\checkmarkc\checkmarki\checkmarkr\checkmarkc\checkmark$\checkmark \checkmark$\checkmark^\checkmark\\checkmarkc\checkmarki\checkmarkr\checkmarkc\checkmark$\checkmark$\checkmark$\checkmark^\checkmark\\checkmarkc\checkmarki\checkmarkr\checkmarkc\checkmark$\checkmark9\checkmark$\checkmark^\checkmark\\checkmarkc\checkmarki\checkmarkr\checkmarkc\checkmark$\checkmark0\checkmark$\checkmark^\checkmark\\checkmarkc\checkmarki\checkmarkr\checkmarkc\checkmark$\checkmark/\checkmark$\checkmark^\checkmark\\checkmarkc\checkmarki\checkmarkr\checkmarkc\checkmark$\checkmark2\checkmark$\checkmark^\checkmark\\checkmarkc\checkmarki\checkmarkr\checkmarkc\checkmark$\checkmark0\checkmark$\checkmark^\checkmark\\checkmarkc\checkmarki\checkmarkr\checkmarkc\checkmark$\checkmark \checkmark$\checkmark^\checkmark\\checkmarkc\checkmarki\checkmarkr\checkmarkc\checkmark$\checkmark=\checkmark$\checkmark^\checkmark\\checkmarkc\checkmarki\checkmarkr\checkmarkc\checkmark$\checkmark \checkmark$\checkmark^\checkmark\\checkmarkc\checkmarki\checkmarkr\checkmarkc\checkmark$\checkmark4\checkmark$\checkmark^\checkmark\\checkmarkc\checkmarki\checkmarkr\checkmarkc\checkmark$\checkmark{\checkmark$\checkmark^\checkmark\\checkmarkc\checkmarki\checkmarkr\checkmarkc\checkmark$\checkmark,\checkmark$\checkmark^\checkmark\\checkmarkc\checkmarki\checkmarkr\checkmarkc\checkmark$\checkmark}\checkmark$\checkmark^\checkmark\\checkmarkc\checkmarki\checkmarkr\checkmarkc\checkmark$\checkmark5\checkmark$\checkmark^\checkmark\\checkmarkc\checkmarki\checkmarkr\checkmarkc\checkmark$\checkmark$\checkmark$\checkmark^\checkmark\\checkmarkc\checkmarki\checkmarkr\checkmarkc\checkmark$\checkmark \checkmark$\checkmark^\checkmark\\checkmarkc\checkmarki\checkmarkr\checkmarkc\checkmark$\checkmarka\checkmark$\checkmark^\checkmark\\checkmarkc\checkmarki\checkmarkr\checkmarkc\checkmark$\checkmarkb\checkmark$\checkmark^\checkmark\\checkmarkc\checkmarki\checkmarkr\checkmarkc\checkmark$\checkmarks\checkmark$\checkmark^\checkmark\\checkmarkc\checkmarki\checkmarkr\checkmarkc\checkmark$\checkmarko\checkmark$\checkmark^\checkmark\\checkmarkc\checkmarki\checkmarkr\checkmarkc\checkmark$\checkmarkr\checkmark$\checkmark^\checkmark\\checkmarkc\checkmarki\checkmarkr\checkmarkc\checkmark$\checkmarkb\checkmark$\checkmark^\checkmark\\checkmarkc\checkmarki\checkmarkr\checkmarkc\checkmark$\checkmarke\checkmark$\checkmark^\checkmark\\checkmarkc\checkmarki\checkmarkr\checkmarkc\checkmark$\checkmark \checkmark$\checkmark^\checkmark\\checkmarkc\checkmarki\checkmarkr\checkmarkc\checkmark$\checkmarkp\checkmark$\checkmark^\checkmark\\checkmarkc\checkmarki\checkmarkr\checkmarkc\checkmark$\checkmarka\checkmark$\checkmark^\checkmark\\checkmarkc\checkmarki\checkmarkr\checkmarkc\checkmark$\checkmarkr\checkmark$\checkmark^\checkmark\\checkmarkc\checkmarki\checkmarkr\checkmarkc\checkmark$\checkmarkf\checkmark$\checkmark^\checkmark\\checkmarkc\checkmarki\checkmarkr\checkmarkc\checkmark$\checkmarka\checkmark$\checkmark^\checkmark\\checkmarkc\checkmarki\checkmarkr\checkmarkc\checkmark$\checkmarki\checkmark$\checkmark^\checkmark\\checkmarkc\checkmarki\checkmarkr\checkmarkc\checkmark$\checkmarkt\checkmark$\checkmark^\checkmark\\checkmarkc\checkmarki\checkmarkr\checkmarkc\checkmark$\checkmarke\checkmark$\checkmark^\checkmark\\checkmarkc\checkmarki\checkmarkr\checkmarkc\checkmark$\checkmarkm\checkmark$\checkmark^\checkmark\\checkmarkc\checkmarki\checkmarkr\checkmarkc\checkmark$\checkmarke\checkmark$\checkmark^\checkmark\\checkmarkc\checkmarki\checkmarkr\checkmarkc\checkmark$\checkmarkn\checkmark$\checkmark^\checkmark\\checkmarkc\checkmarki\checkmarkr\checkmarkc\checkmark$\checkmarkt\checkmark$\checkmark^\checkmark\\checkmarkc\checkmarki\checkmarkr\checkmarkc\checkmark$\checkmark \checkmark$\checkmark^\checkmark\\checkmarkc\checkmarki\checkmarkr\checkmarkc\checkmark$\checkmarkc\checkmark$\checkmark^\checkmark\\checkmarkc\checkmarki\checkmarkr\checkmarkc\checkmark$\checkmarke\checkmark$\checkmark^\checkmark\\checkmarkc\checkmarki\checkmarkr\checkmarkc\checkmark$\checkmarkt\checkmark$\checkmark^\checkmark\\checkmarkc\checkmarki\checkmarkr\checkmarkc\checkmark$\checkmarkt\checkmark$\checkmark^\checkmark\\checkmarkc\checkmarki\checkmarkr\checkmarkc\checkmark$\checkmarke\checkmark$\checkmark^\checkmark\\checkmarkc\checkmarki\checkmarkr\checkmarkc\checkmark$\checkmark \checkmark$\checkmark^\checkmark\\checkmarkc\checkmarki\checkmarkr\checkmarkc\checkmark$\checkmarkc\checkmark$\checkmark^\checkmark\\checkmarkc\checkmarki\checkmarkr\checkmarkc\checkmark$\checkmarko\checkmark$\checkmark^\checkmark\\checkmarkc\checkmarki\checkmarkr\checkmarkc\checkmark$\checkmarkn\checkmark$\checkmark^\checkmark\\checkmarkc\checkmarki\checkmarkr\checkmarkc\checkmark$\checkmarkt\checkmark$\checkmark^\checkmark\\checkmarkc\checkmarki\checkmarkr\checkmarkc\checkmark$\checkmarkr\checkmark$\checkmark^\checkmark\\checkmarkc\checkmarki\checkmarkr\checkmarkc\checkmark$\checkmarka\checkmark$\checkmark^\checkmark\\checkmarkc\checkmarki\checkmarkr\checkmarkc\checkmark$\checkmarki\checkmark$\checkmark^\checkmark\\checkmarkc\checkmarki\checkmarkr\checkmarkc\checkmark$\checkmarkn\checkmark$\checkmark^\checkmark\\checkmarkc\checkmarki\checkmarkr\checkmarkc\checkmark$\checkmarkt\checkmark$\checkmark^\checkmark\\checkmarkc\checkmarki\checkmarkr\checkmarkc\checkmark$\checkmarke\checkmark$\checkmark^\checkmark\\checkmarkc\checkmarki\checkmarkr\checkmarkc\checkmark$\checkmark.\checkmark$\checkmark^\checkmark\\checkmarkc\checkmarki\checkmarkr\checkmarkc\checkmark$\checkmark
\checkmark$\checkmark^\checkmark\\checkmarkc\checkmarki\checkmarkr\checkmarkc\checkmark$\checkmark
\checkmark$\checkmark^\checkmark\\checkmarkc\checkmarki\checkmarkr\checkmarkc\checkmark$\checkmarkL\checkmark$\checkmark^\checkmark\\checkmarkc\checkmarki\checkmarkr\checkmarkc\checkmark$\checkmark'\checkmark$\checkmark^\checkmark\\checkmarkc\checkmarki\checkmarkr\checkmarkc\checkmark$\checkmarki\checkmark$\checkmark^\checkmark\\checkmarkc\checkmarki\checkmarkr\checkmarkc\checkmark$\checkmarkn\checkmark$\checkmark^\checkmark\\checkmarkc\checkmarki\checkmarkr\checkmarkc\checkmark$\checkmarke\checkmark$\checkmark^\checkmark\\checkmarkc\checkmarki\checkmarkr\checkmarkc\checkmark$\checkmarkr\checkmark$\checkmark^\checkmark\\checkmarkc\checkmarki\checkmarkr\checkmarkc\checkmark$\checkmarkt\checkmark$\checkmark^\checkmark\\checkmarkc\checkmarki\checkmarkr\checkmarkc\checkmark$\checkmarki\checkmark$\checkmark^\checkmark\\checkmarkc\checkmarki\checkmarkr\checkmarkc\checkmark$\checkmarke\checkmark$\checkmark^\checkmark\\checkmarkc\checkmarki\checkmarkr\checkmarkc\checkmark$\checkmark \checkmark$\checkmark^\checkmark\\checkmarkc\checkmarki\checkmarkr\checkmarkc\checkmark$\checkmarkg\checkmark$\checkmark^\checkmark\\checkmarkc\checkmarki\checkmarkr\checkmarkc\checkmark$\checkmarké\checkmark$\checkmark^\checkmark\\checkmarkc\checkmarki\checkmarkr\checkmarkc\checkmark$\checkmarkn\checkmark$\checkmark^\checkmark\\checkmarkc\checkmarki\checkmarkr\checkmarkc\checkmark$\checkmarké\checkmark$\checkmark^\checkmark\\checkmarkc\checkmarki\checkmarkr\checkmarkc\checkmark$\checkmarkr\checkmark$\checkmark^\checkmark\\checkmarkc\checkmarki\checkmarkr\checkmarkc\checkmark$\checkmarké\checkmark$\checkmark^\checkmark\\checkmarkc\checkmarki\checkmarkr\checkmarkc\checkmark$\checkmarke\checkmark$\checkmark^\checkmark\\checkmarkc\checkmarki\checkmarkr\checkmarkc\checkmark$\checkmark \checkmark$\checkmark^\checkmark\\checkmarkc\checkmarki\checkmarkr\checkmarkc\checkmark$\checkmarkp\checkmark$\checkmark^\checkmark\\checkmarkc\checkmarki\checkmarkr\checkmarkc\checkmark$\checkmarka\checkmark$\checkmark^\checkmark\\checkmarkc\checkmarki\checkmarkr\checkmarkc\checkmark$\checkmarkr\checkmark$\checkmark^\checkmark\\checkmarkc\checkmarki\checkmarkr\checkmarkc\checkmark$\checkmark \checkmark$\checkmark^\checkmark\\checkmarkc\checkmarki\checkmarkr\checkmarkc\checkmark$\checkmarkc\checkmark$\checkmark^\checkmark\\checkmarkc\checkmarki\checkmarkr\checkmarkc\checkmark$\checkmarke\checkmark$\checkmark^\checkmark\\checkmarkc\checkmarki\checkmarkr\checkmarkc\checkmark$\checkmarks\checkmark$\checkmark^\checkmark\\checkmarkc\checkmarki\checkmarkr\checkmarkc\checkmark$\checkmark \checkmark$\checkmark^\checkmark\\checkmarkc\checkmarki\checkmarkr\checkmarkc\checkmark$\checkmarkm\checkmark$\checkmark^\checkmark\\checkmarkc\checkmarki\checkmarkr\checkmarkc\checkmark$\checkmarké\checkmark$\checkmark^\checkmark\\checkmarkc\checkmarki\checkmarkr\checkmarkc\checkmark$\checkmarkc\checkmark$\checkmark^\checkmark\\checkmarkc\checkmarki\checkmarkr\checkmarkc\checkmark$\checkmarka\checkmark$\checkmark^\checkmark\\checkmarkc\checkmarki\checkmarkr\checkmarkc\checkmark$\checkmarkn\checkmark$\checkmark^\checkmark\\checkmarkc\checkmarki\checkmarkr\checkmarkc\checkmark$\checkmarki\checkmark$\checkmark^\checkmark\\checkmarkc\checkmarki\checkmarkr\checkmarkc\checkmark$\checkmarks\checkmark$\checkmark^\checkmark\\checkmarkc\checkmarki\checkmarkr\checkmarkc\checkmark$\checkmarkm\checkmark$\checkmark^\checkmark\\checkmarkc\checkmarki\checkmarkr\checkmarkc\checkmark$\checkmarke\checkmark$\checkmark^\checkmark\\checkmarkc\checkmarki\checkmarkr\checkmarkc\checkmark$\checkmarks\checkmark$\checkmark^\checkmark\\checkmarkc\checkmarki\checkmarkr\checkmarkc\checkmark$\checkmark \checkmark$\checkmark^\checkmark\\checkmarkc\checkmarki\checkmarkr\checkmarkc\checkmark$\checkmarkd\checkmark$\checkmark^\checkmark\\checkmarkc\checkmarki\checkmarkr\checkmarkc\checkmark$\checkmarke\checkmark$\checkmark^\checkmark\\checkmarkc\checkmarki\checkmarkr\checkmarkc\checkmark$\checkmark \checkmark$\checkmark^\checkmark\\checkmarkc\checkmarki\checkmarkr\checkmarkc\checkmark$\checkmarkp\checkmark$\checkmark^\checkmark\\checkmarkc\checkmarki\checkmarkr\checkmarkc\checkmark$\checkmarkr\checkmark$\checkmark^\checkmark\\checkmarkc\checkmarki\checkmarkr\checkmarkc\checkmark$\checkmarké\checkmark$\checkmark^\checkmark\\checkmarkc\checkmarki\checkmarkr\checkmarkc\checkmark$\checkmarkc\checkmark$\checkmark^\checkmark\\checkmarkc\checkmarki\checkmarkr\checkmarkc\checkmark$\checkmarki\checkmark$\checkmark^\checkmark\\checkmarkc\checkmarki\checkmarkr\checkmarkc\checkmark$\checkmarks\checkmark$\checkmark^\checkmark\\checkmarkc\checkmarki\checkmarkr\checkmarkc\checkmark$\checkmarki\checkmark$\checkmark^\checkmark\\checkmarkc\checkmarki\checkmarkr\checkmarkc\checkmark$\checkmarko\checkmark$\checkmark^\checkmark\\checkmarkc\checkmarki\checkmarkr\checkmarkc\checkmark$\checkmarkn\checkmark$\checkmark^\checkmark\\checkmarkc\checkmarki\checkmarkr\checkmarkc\checkmark$\checkmark \checkmark$\checkmark^\checkmark\\checkmarkc\checkmarki\checkmarkr\checkmarkc\checkmark$\checkmarki\checkmark$\checkmark^\checkmark\\checkmarkc\checkmarki\checkmarkr\checkmarkc\checkmark$\checkmarkm\checkmark$\checkmark^\checkmark\\checkmarkc\checkmarki\checkmarkr\checkmarkc\checkmark$\checkmarkp\checkmark$\checkmark^\checkmark\\checkmarkc\checkmarki\checkmarkr\checkmarkc\checkmark$\checkmarkl\checkmark$\checkmark^\checkmark\\checkmarkc\checkmarki\checkmarkr\checkmarkc\checkmark$\checkmarki\checkmark$\checkmark^\checkmark\\checkmarkc\checkmarki\checkmarkr\checkmarkc\checkmark$\checkmarkq\checkmark$\checkmark^\checkmark\\checkmarkc\checkmarki\checkmarkr\checkmarkc\checkmark$\checkmarku\checkmark$\checkmark^\checkmark\\checkmarkc\checkmarki\checkmarkr\checkmarkc\checkmark$\checkmarke\checkmark$\checkmark^\checkmark\\checkmarkc\checkmarki\checkmarkr\checkmarkc\checkmark$\checkmark \checkmark$\checkmark^\checkmark\\checkmarkc\checkmarki\checkmarkr\checkmarkc\checkmark$\checkmarkl\checkmark$\checkmark^\checkmark\\checkmarkc\checkmarki\checkmarkr\checkmarkc\checkmark$\checkmark'\checkmark$\checkmark^\checkmark\\checkmarkc\checkmarki\checkmarkr\checkmarkc\checkmark$\checkmarku\checkmark$\checkmark^\checkmark\\checkmarkc\checkmarki\checkmarkr\checkmarkc\checkmark$\checkmarkt\checkmark$\checkmark^\checkmark\\checkmarkc\checkmarki\checkmarkr\checkmarkc\checkmark$\checkmarki\checkmark$\checkmark^\checkmark\\checkmarkc\checkmarki\checkmarkr\checkmarkc\checkmark$\checkmarkl\checkmark$\checkmark^\checkmark\\checkmarkc\checkmarki\checkmarkr\checkmarkc\checkmark$\checkmarki\checkmark$\checkmark^\checkmark\\checkmarkc\checkmarki\checkmarkr\checkmarkc\checkmark$\checkmarks\checkmark$\checkmark^\checkmark\\checkmarkc\checkmarki\checkmarkr\checkmarkc\checkmark$\checkmarka\checkmark$\checkmark^\checkmark\\checkmarkc\checkmarki\checkmarkr\checkmarkc\checkmark$\checkmarkt\checkmark$\checkmark^\checkmark\\checkmarkc\checkmarki\checkmarkr\checkmarkc\checkmark$\checkmarki\checkmark$\checkmark^\checkmark\\checkmarkc\checkmarki\checkmarkr\checkmarkc\checkmark$\checkmarko\checkmark$\checkmark^\checkmark\\checkmarkc\checkmarki\checkmarkr\checkmarkc\checkmark$\checkmarkn\checkmark$\checkmark^\checkmark\\checkmarkc\checkmarki\checkmarkr\checkmarkc\checkmark$\checkmark \checkmark$\checkmark^\checkmark\\checkmarkc\checkmarki\checkmarkr\checkmarkc\checkmark$\checkmarkd\checkmark$\checkmark^\checkmark\\checkmarkc\checkmarki\checkmarkr\checkmarkc\checkmark$\checkmark'\checkmark$\checkmark^\checkmark\\checkmarkc\checkmarki\checkmarkr\checkmarkc\checkmark$\checkmarku\checkmark$\checkmark^\checkmark\\checkmarkc\checkmarki\checkmarkr\checkmarkc\checkmark$\checkmarkn\checkmark$\checkmark^\checkmark\\checkmarkc\checkmarki\checkmarkr\checkmarkc\checkmark$\checkmark \checkmark$\checkmark^\checkmark\\checkmarkc\checkmarki\checkmarkr\checkmarkc\checkmark$\checkmarks\checkmark$\checkmark^\checkmark\\checkmarkc\checkmarki\checkmarkr\checkmarkc\checkmark$\checkmarky\checkmark$\checkmark^\checkmark\\checkmarkc\checkmarki\checkmarkr\checkmarkc\checkmark$\checkmarks\checkmark$\checkmark^\checkmark\\checkmarkc\checkmarki\checkmarkr\checkmarkc\checkmark$\checkmarkt\checkmark$\checkmark^\checkmark\\checkmarkc\checkmarki\checkmarkr\checkmarkc\checkmark$\checkmarkè\checkmark$\checkmark^\checkmark\\checkmarkc\checkmarki\checkmarkr\checkmarkc\checkmark$\checkmarkm\checkmark$\checkmark^\checkmark\\checkmarkc\checkmarki\checkmarkr\checkmarkc\checkmark$\checkmarke\checkmark$\checkmark^\checkmark\\checkmarkc\checkmarki\checkmarkr\checkmarkc\checkmark$\checkmark \checkmark$\checkmark^\checkmark\\checkmarkc\checkmarki\checkmarkr\checkmarkc\checkmark$\checkmarkd\checkmark$\checkmark^\checkmark\\checkmarkc\checkmarki\checkmarkr\checkmarkc\checkmark$\checkmarke\checkmark$\checkmark^\checkmark\\checkmarkc\checkmarki\checkmarkr\checkmarkc\checkmark$\checkmark \checkmark$\checkmark^\checkmark\\checkmarkc\checkmarki\checkmarkr\checkmarkc\checkmark$\checkmarkc\checkmark$\checkmark^\checkmark\\checkmarkc\checkmarki\checkmarkr\checkmarkc\checkmark$\checkmarko\checkmark$\checkmark^\checkmark\\checkmarkc\checkmarki\checkmarkr\checkmarkc\checkmark$\checkmarkn\checkmark$\checkmark^\checkmark\\checkmarkc\checkmarki\checkmarkr\checkmarkc\checkmark$\checkmarkt\checkmark$\checkmark^\checkmark\\checkmarkc\checkmarki\checkmarkr\checkmarkc\checkmark$\checkmarkr\checkmark$\checkmark^\checkmark\\checkmarkc\checkmarki\checkmarkr\checkmarkc\checkmark$\checkmarkô\checkmark$\checkmark^\checkmark\\checkmarkc\checkmarki\checkmarkr\checkmarkc\checkmark$\checkmarkl\checkmark$\checkmark^\checkmark\\checkmarkc\checkmarki\checkmarkr\checkmarkc\checkmark$\checkmarke\checkmark$\checkmark^\checkmark\\checkmarkc\checkmarki\checkmarkr\checkmarkc\checkmark$\checkmark-\checkmark$\checkmark^\checkmark\\checkmarkc\checkmarki\checkmarkr\checkmarkc\checkmark$\checkmarkc\checkmark$\checkmark^\checkmark\\checkmarkc\checkmarki\checkmarkr\checkmarkc\checkmark$\checkmarko\checkmark$\checkmark^\checkmark\\checkmarkc\checkmarki\checkmarkr\checkmarkc\checkmark$\checkmarkm\checkmark$\checkmark^\checkmark\\checkmarkc\checkmarki\checkmarkr\checkmarkc\checkmark$\checkmarkm\checkmark$\checkmark^\checkmark\\checkmarkc\checkmarki\checkmarkr\checkmarkc\checkmark$\checkmarka\checkmark$\checkmark^\checkmark\\checkmarkc\checkmarki\checkmarkr\checkmarkc\checkmark$\checkmarkn\checkmark$\checkmark^\checkmark\\checkmarkc\checkmarki\checkmarkr\checkmarkc\checkmark$\checkmarkd\checkmark$\checkmark^\checkmark\\checkmarkc\checkmarki\checkmarkr\checkmarkc\checkmark$\checkmarke\checkmark$\checkmark^\checkmark\\checkmarkc\checkmarki\checkmarkr\checkmarkc\checkmark$\checkmark \checkmark$\checkmark^\checkmark\\checkmarkc\checkmarki\checkmarkr\checkmarkc\checkmark$\checkmarkr\checkmark$\checkmark^\checkmark\\checkmarkc\checkmarki\checkmarkr\checkmarkc\checkmark$\checkmarko\checkmark$\checkmark^\checkmark\\checkmarkc\checkmarki\checkmarkr\checkmarkc\checkmark$\checkmarkb\checkmark$\checkmark^\checkmark\\checkmarkc\checkmarki\checkmarkr\checkmarkc\checkmark$\checkmarku\checkmark$\checkmark^\checkmark\\checkmarkc\checkmarki\checkmarkr\checkmarkc\checkmark$\checkmarks\checkmark$\checkmark^\checkmark\\checkmarkc\checkmarki\checkmarkr\checkmarkc\checkmark$\checkmarkt\checkmark$\checkmark^\checkmark\\checkmarkc\checkmarki\checkmarkr\checkmarkc\checkmark$\checkmarke\checkmark$\checkmark^\checkmark\\checkmarkc\checkmarki\checkmarkr\checkmarkc\checkmark$\checkmark,\checkmark$\checkmark^\checkmark\\checkmarkc\checkmarki\checkmarkr\checkmarkc\checkmark$\checkmark \checkmark$\checkmark^\checkmark\\checkmarkc\checkmarki\checkmarkr\checkmarkc\checkmark$\checkmarkq\checkmark$\checkmark^\checkmark\\checkmarkc\checkmarki\checkmarkr\checkmarkc\checkmark$\checkmarku\checkmark$\checkmark^\checkmark\\checkmarkc\checkmarki\checkmarkr\checkmarkc\checkmark$\checkmarke\checkmark$\checkmark^\checkmark\\checkmarkc\checkmarki\checkmarkr\checkmarkc\checkmark$\checkmark \checkmark$\checkmark^\checkmark\\checkmarkc\checkmarki\checkmarkr\checkmarkc\checkmark$\checkmarkn\checkmark$\checkmark^\checkmark\\checkmarkc\checkmarki\checkmarkr\checkmarkc\checkmark$\checkmarko\checkmark$\checkmark^\checkmark\\checkmarkc\checkmarki\checkmarkr\checkmarkc\checkmark$\checkmarku\checkmark$\checkmark^\checkmark\\checkmarkc\checkmarki\checkmarkr\checkmarkc\checkmark$\checkmarks\checkmark$\checkmark^\checkmark\\checkmarkc\checkmarki\checkmarkr\checkmarkc\checkmark$\checkmark \checkmark$\checkmark^\checkmark\\checkmarkc\checkmarki\checkmarkr\checkmarkc\checkmark$\checkmarka\checkmark$\checkmark^\checkmark\\checkmarkc\checkmarki\checkmarkr\checkmarkc\checkmark$\checkmarkn\checkmark$\checkmark^\checkmark\\checkmarkc\checkmarki\checkmarkr\checkmarkc\checkmark$\checkmarka\checkmark$\checkmark^\checkmark\\checkmarkc\checkmarki\checkmarkr\checkmarkc\checkmark$\checkmarkl\checkmark$\checkmark^\checkmark\\checkmarkc\checkmarki\checkmarkr\checkmarkc\checkmark$\checkmarky\checkmark$\checkmark^\checkmark\\checkmarkc\checkmarki\checkmarkr\checkmarkc\checkmark$\checkmarks\checkmark$\checkmark^\checkmark\\checkmarkc\checkmarki\checkmarkr\checkmarkc\checkmark$\checkmarko\checkmark$\checkmark^\checkmark\\checkmarkc\checkmarki\checkmarkr\checkmarkc\checkmark$\checkmarkn\checkmark$\checkmark^\checkmark\\checkmarkc\checkmarki\checkmarkr\checkmarkc\checkmark$\checkmarks\checkmark$\checkmark^\checkmark\\checkmarkc\checkmarki\checkmarkr\checkmarkc\checkmark$\checkmark \checkmark$\checkmark^\checkmark\\checkmarkc\checkmarki\checkmarkr\checkmarkc\checkmark$\checkmarkd\checkmark$\checkmark^\checkmark\\checkmarkc\checkmarki\checkmarkr\checkmarkc\checkmark$\checkmarka\checkmark$\checkmark^\checkmark\\checkmarkc\checkmarki\checkmarkr\checkmarkc\checkmark$\checkmarkn\checkmark$\checkmark^\checkmark\\checkmarkc\checkmarki\checkmarkr\checkmarkc\checkmark$\checkmarks\checkmark$\checkmark^\checkmark\\checkmarkc\checkmarki\checkmarkr\checkmarkc\checkmark$\checkmark \checkmark$\checkmark^\checkmark\\checkmarkc\checkmarki\checkmarkr\checkmarkc\checkmark$\checkmarkl\checkmark$\checkmark^\checkmark\\checkmarkc\checkmarki\checkmarkr\checkmarkc\checkmark$\checkmarke\checkmark$\checkmark^\checkmark\\checkmarkc\checkmarki\checkmarkr\checkmarkc\checkmark$\checkmarks\checkmark$\checkmark^\checkmark\\checkmarkc\checkmarki\checkmarkr\checkmarkc\checkmark$\checkmark \checkmark$\checkmark^\checkmark\\checkmarkc\checkmarki\checkmarkr\checkmarkc\checkmark$\checkmarkc\checkmark$\checkmark^\checkmark\\checkmarkc\checkmarki\checkmarkr\checkmarkc\checkmark$\checkmarkh\checkmark$\checkmark^\checkmark\\checkmarkc\checkmarki\checkmarkr\checkmarkc\checkmark$\checkmarka\checkmark$\checkmark^\checkmark\\checkmarkc\checkmarki\checkmarkr\checkmarkc\checkmark$\checkmarkp\checkmark$\checkmark^\checkmark\\checkmarkc\checkmarki\checkmarkr\checkmarkc\checkmark$\checkmarki\checkmark$\checkmark^\checkmark\\checkmarkc\checkmarki\checkmarkr\checkmarkc\checkmark$\checkmarkt\checkmark$\checkmark^\checkmark\\checkmarkc\checkmarki\checkmarkr\checkmarkc\checkmark$\checkmarkr\checkmark$\checkmark^\checkmark\\checkmarkc\checkmarki\checkmarkr\checkmarkc\checkmark$\checkmarke\checkmark$\checkmark^\checkmark\\checkmarkc\checkmarki\checkmarkr\checkmarkc\checkmark$\checkmarks\checkmark$\checkmark^\checkmark\\checkmarkc\checkmarki\checkmarkr\checkmarkc\checkmark$\checkmark \checkmark$\checkmark^\checkmark\\checkmarkc\checkmarki\checkmarkr\checkmarkc\checkmark$\checkmarkd\checkmark$\checkmark^\checkmark\\checkmarkc\checkmarki\checkmarkr\checkmarkc\checkmark$\checkmark'\checkmark$\checkmark^\checkmark\\checkmarkc\checkmarki\checkmarkr\checkmarkc\checkmark$\checkmarka\checkmark$\checkmark^\checkmark\\checkmarkc\checkmarki\checkmarkr\checkmarkc\checkmark$\checkmarkr\checkmark$\checkmark^\checkmark\\checkmarkc\checkmarki\checkmarkr\checkmarkc\checkmark$\checkmarkc\checkmark$\checkmark^\checkmark\\checkmarkc\checkmarki\checkmarkr\checkmarkc\checkmark$\checkmarkh\checkmark$\checkmark^\checkmark\\checkmarkc\checkmarki\checkmarkr\checkmarkc\checkmark$\checkmarki\checkmark$\checkmark^\checkmark\\checkmarkc\checkmarki\checkmarkr\checkmarkc\checkmark$\checkmarkt\checkmark$\checkmark^\checkmark\\checkmarkc\checkmarki\checkmarkr\checkmarkc\checkmark$\checkmarke\checkmark$\checkmark^\checkmark\\checkmarkc\checkmarki\checkmarkr\checkmarkc\checkmark$\checkmarkc\checkmark$\checkmark^\checkmark\\checkmarkc\checkmarki\checkmarkr\checkmarkc\checkmark$\checkmarkt\checkmark$\checkmark^\checkmark\\checkmarkc\checkmarki\checkmarkr\checkmarkc\checkmark$\checkmarku\checkmark$\checkmark^\checkmark\\checkmarkc\checkmarki\checkmarkr\checkmarkc\checkmark$\checkmarkr\checkmark$\checkmark^\checkmark\\checkmarkc\checkmarki\checkmarkr\checkmarkc\checkmark$\checkmarke\checkmark$\checkmark^\checkmark\\checkmarkc\checkmarki\checkmarkr\checkmarkc\checkmark$\checkmark \checkmark$\checkmark^\checkmark\\checkmarkc\checkmarki\checkmarkr\checkmarkc\checkmark$\checkmarkm\checkmark$\checkmark^\checkmark\\checkmarkc\checkmarki\checkmarkr\checkmarkc\checkmark$\checkmarké\checkmark$\checkmark^\checkmark\\checkmarkc\checkmarki\checkmarkr\checkmarkc\checkmark$\checkmarkc\checkmark$\checkmark^\checkmark\\checkmarkc\checkmarki\checkmarkr\checkmarkc\checkmark$\checkmarka\checkmark$\checkmark^\checkmark\\checkmarkc\checkmarki\checkmarkr\checkmarkc\checkmark$\checkmarkt\checkmark$\checkmark^\checkmark\\checkmarkc\checkmarki\checkmarkr\checkmarkc\checkmark$\checkmarkr\checkmark$\checkmark^\checkmark\\checkmarkc\checkmarki\checkmarkr\checkmarkc\checkmark$\checkmarko\checkmark$\checkmark^\checkmark\\checkmarkc\checkmarki\checkmarkr\checkmarkc\checkmark$\checkmarkn\checkmark$\checkmark^\checkmark\\checkmarkc\checkmarki\checkmarkr\checkmarkc\checkmark$\checkmarki\checkmark$\checkmark^\checkmark\\checkmarkc\checkmarki\checkmarkr\checkmarkc\checkmark$\checkmarkq\checkmark$\checkmark^\checkmark\\checkmarkc\checkmarki\checkmarkr\checkmarkc\checkmark$\checkmarku\checkmark$\checkmark^\checkmark\\checkmarkc\checkmarki\checkmarkr\checkmarkc\checkmark$\checkmarke\checkmark$\checkmark^\checkmark\\checkmarkc\checkmarki\checkmarkr\checkmarkc\checkmark$\checkmark \checkmark$\checkmark^\checkmark\\checkmarkc\checkmarki\checkmarkr\checkmarkc\checkmark$\checkmarke\checkmark$\checkmark^\checkmark\\checkmarkc\checkmarki\checkmarkr\checkmarkc\checkmark$\checkmarkt\checkmark$\checkmark^\checkmark\\checkmarkc\checkmarki\checkmarkr\checkmarkc\checkmark$\checkmark \checkmark$\checkmark^\checkmark\\checkmarkc\checkmarki\checkmarkr\checkmarkc\checkmark$\checkmarkd\checkmark$\checkmark^\checkmark\\checkmarkc\checkmarki\checkmarkr\checkmarkc\checkmark$\checkmarke\checkmark$\checkmark^\checkmark\\checkmarkc\checkmarki\checkmarkr\checkmarkc\checkmark$\checkmark \checkmark$\checkmark^\checkmark\\checkmarkc\checkmarki\checkmarkr\checkmarkc\checkmark$\checkmarkd\checkmark$\checkmark^\checkmark\\checkmarkc\checkmarki\checkmarkr\checkmarkc\checkmark$\checkmarké\checkmark$\checkmark^\checkmark\\checkmarkc\checkmarki\checkmarkr\checkmarkc\checkmark$\checkmarkv\checkmark$\checkmark^\checkmark\\checkmarkc\checkmarki\checkmarkr\checkmarkc\checkmark$\checkmarke\checkmark$\checkmark^\checkmark\\checkmarkc\checkmarki\checkmarkr\checkmarkc\checkmark$\checkmarkl\checkmark$\checkmark^\checkmark\\checkmarkc\checkmarki\checkmarkr\checkmarkc\checkmark$\checkmarko\checkmark$\checkmark^\checkmark\\checkmarkc\checkmarki\checkmarkr\checkmarkc\checkmark$\checkmarkp\checkmark$\checkmark^\checkmark\\checkmarkc\checkmarki\checkmarkr\checkmarkc\checkmark$\checkmarkp\checkmark$\checkmark^\checkmark\\checkmarkc\checkmarki\checkmarkr\checkmarkc\checkmark$\checkmarke\checkmark$\checkmark^\checkmark\\checkmarkc\checkmarki\checkmarkr\checkmarkc\checkmark$\checkmarkm\checkmark$\checkmark^\checkmark\\checkmarkc\checkmarki\checkmarkr\checkmarkc\checkmark$\checkmarke\checkmark$\checkmark^\checkmark\\checkmarkc\checkmarki\checkmarkr\checkmarkc\checkmark$\checkmarkn\checkmark$\checkmark^\checkmark\\checkmarkc\checkmarki\checkmarkr\checkmarkc\checkmark$\checkmarkt\checkmark$\checkmark^\checkmark\\checkmarkc\checkmarki\checkmarkr\checkmarkc\checkmark$\checkmark \checkmark$\checkmark^\checkmark\\checkmarkc\checkmarki\checkmarkr\checkmarkc\checkmark$\checkmarkl\checkmark$\checkmark^\checkmark\\checkmarkc\checkmarki\checkmarkr\checkmarkc\checkmark$\checkmarko\checkmark$\checkmark^\checkmark\\checkmarkc\checkmarki\checkmarkr\checkmarkc\checkmark$\checkmarkg\checkmark$\checkmark^\checkmark\\checkmarkc\checkmarki\checkmarkr\checkmarkc\checkmark$\checkmarki\checkmark$\checkmark^\checkmark\\checkmarkc\checkmarki\checkmarkr\checkmarkc\checkmark$\checkmarkc\checkmark$\checkmark^\checkmark\\checkmarkc\checkmarki\checkmarkr\checkmarkc\checkmark$\checkmarki\checkmark$\checkmark^\checkmark\\checkmarkc\checkmarki\checkmarkr\checkmarkc\checkmark$\checkmarke\checkmark$\checkmark^\checkmark\\checkmarkc\checkmarki\checkmarkr\checkmarkc\checkmark$\checkmarkl\checkmark$\checkmark^\checkmark\\checkmarkc\checkmarki\checkmarkr\checkmarkc\checkmark$\checkmark \checkmark$\checkmark^\checkmark\\checkmarkc\checkmarki\checkmarkr\checkmarkc\checkmark$\checkmarkq\checkmark$\checkmark^\checkmark\\checkmarkc\checkmarki\checkmarkr\checkmarkc\checkmark$\checkmarku\checkmark$\checkmark^\checkmark\\checkmarkc\checkmarki\checkmarkr\checkmarkc\checkmark$\checkmarki\checkmark$\checkmark^\checkmark\\checkmarkc\checkmarki\checkmarkr\checkmarkc\checkmark$\checkmark \checkmark$\checkmark^\checkmark\\checkmarkc\checkmarki\checkmarkr\checkmarkc\checkmark$\checkmarks\checkmark$\checkmark^\checkmark\\checkmarkc\checkmarki\checkmarkr\checkmarkc\checkmark$\checkmarku\checkmark$\checkmark^\checkmark\\checkmarkc\checkmarki\checkmarkr\checkmarkc\checkmark$\checkmarki\checkmark$\checkmark^\checkmark\\checkmarkc\checkmarki\checkmarkr\checkmarkc\checkmark$\checkmarkv\checkmark$\checkmark^\checkmark\\checkmarkc\checkmarki\checkmarkr\checkmarkc\checkmark$\checkmarke\checkmark$\checkmark^\checkmark\\checkmarkc\checkmarki\checkmarkr\checkmarkc\checkmark$\checkmarkn\checkmark$\checkmark^\checkmark\\checkmarkc\checkmarki\checkmarkr\checkmarkc\checkmark$\checkmarkt\checkmark$\checkmark^\checkmark\\checkmarkc\checkmarki\checkmarkr\checkmarkc\checkmark$\checkmark.\checkmark$\checkmark^\checkmark\\checkmarkc\checkmarki\checkmarkr\checkmarkc\checkmark$\checkmark
          % Étude Mécanique et Dimensionnement
% !TEX root = ../main.tex
\chapter{Architecture Électronique et Ingénierie des Signaux}
\label{chap:mechatronics}

La conception mécanique, validée au chapitre précédent, ne peut fournir le summum de sa précision sans une électronique de commande hyper-réactive et immunisée contre les perturbations électromagnétiques. L'architecture matérielle d'une CNC 5 axes est fondamentalement différente d'un sytème cartésien simple, exigeant le pilotage asynchrone à haute fréquence de multiples onduleurs (Drivers) et la gestion de la sûreté des personnes \textit{via} des circuits câblés de coupure d'énergie.

% ========================================================
\section{L'Unité de Traitement Embarquée (Le "Cerveau")}

\subsection{Analyse Critique des Microcontrôleurs}
Le choix de l'unité de traitement est dicté par la fréquence de récurrence spatiale des interruptions temporelles (Timers).
Chacun des 5 moteurs nécessite 2 broches (GPIO) de contrôle direct (\textit{Step} et \textit{Direction}).
La fréquence de signal \textit{Step} (impulsions par seconde, Hz) requise pour faire tourner un moteur à la vitesse limite de $v_{max} = 3000 \text{ mm/min}$ (soit $50 \text{ mm/s}$) avec la vis à billes détaillée au chapitre précédent est :
\begin{equation}
    f_{step} = \frac{v_{max}}{\text{Résolution}} = \frac{50}{(5 / (200 \times 16))} = \frac{50}{0.0015625} \approx 32~000 \text{ Hz}
\end{equation}

* \textbf{Le problème des 8 bits (Arduino) :} Le microcontrôleur classique ATmega328P de 16 MHz a un cycle machine de $62.5 \text{ ns}$. Avec l'algorithmique cinématique complexe tournant dans la seule boucle `void loop()`, il engendre du \textit{Jitter} (gigue de phase temporelle) inévitable, entraînant le décrochage (calage) physique des moteurs.

* \textbf{La solution 32 bits, Modèle ESP32 :} L'Espressif ESP32-WROOM-32U a été sélectionné pour sa structure Xtensa Dual-Core à 240 MHz (cycle machine $\approx 4 \text{ ns}$). Outre sa puissance brute ($600 \text{ DMIPS}$), sa validation industrielle repose sur sa capacité à exécuter nativement un Système d'Exploitation Temps Réel : \textbf{FreeRTOS}.

\subsection{Allocation Mémoire et Tâches (FreeRTOS)}
L'intérêt majeur de l'ESP32 est la séparation asymétrique des tâches logicielles.
Grâce à l'API système d'Espressif (`xTaskCreatePinnedToCore`), l'architecture a été scindée physiquement :
\begin{itemize}
    \item \textbf{Sur le Core 0 (Tâches non-critiques temporellement) :} Gestion de la pile réseau (TCP/IP Wi-Fi), du serveur WebSockets, réception des trames JSON depuis l'interface Flutter, et parsing/calcul matriciel Denavit-Hartenberg.
    \item \textbf{Sur le Core 1 (Tâches hyper-déterministes) :} Gestion de l'\textit{Interrupt Service Routine} (ISR) d'un Timer Hardware dédié unilatéralement au balayage des pins STEP/DIR des cinq axes, lisant le buffer circulaire de trajectoire et pilotant les drivers sans la moindre instruction de débranchement concurrent.
\end{itemize}

% ========================================================
\section{Chaîne de Puissance et Drivers Moteurs}

L'ESP32 génère des signaux TTL logiques de $3.3 \text{ V}$ avec un courant maximal extractible d'environ $40 \text{ mA}$ par pin. Les moteurs NEMA 17 demandent des pointes dépassant les $2.5 \text{ A}$ (à $24 \text{ V}$). Cette différence flagrante de puissance d'un rapport de plus de $500$ impose d'intercaler de gros "muscles" électroniques : les Stepper Drivers autonomes complets (type DM556).

\begin{figure}[htbp]
    \centering
    \resizebox{0.95\textwidth}{!}{
    \begin{tikzpicture}[
        box/.style={draw, rounded corners, minimum height=1.5cm, text width=3.5cm, align=center, thick},
        line/.style={-{Stealth[scale=1.2]}, thick}
    ]
        % Composants
        \node[box, fill=orange!10] (pc) {PC / Flutter \\ (Via Wi-Fi UDP)};
        \node[box, fill=blue!10, right=2cm of pc] (esp32) {Microcontrôleur \\ ESP32 (3.3V)};
        \node[box, fill=yellow!10, right=2cm of esp32] (ls) {Level Shifter \\ (SN74AHCT125N)};
        \node[box, fill=red!10, right=2cm of ls] (driver) {Driver DM556 \\ (Optocoupleurs)};
        \node[box, fill=gray!20, below=1.5cm of driver] (motor) {Moteurs NEMA 17 \\ (Axes X,Y,Z,A,B)};
        \node[box, fill=green!10, below=1.5cm of ls] (psu) {Alimentation \\ 36V / 350W};

        % Connexions
        \draw[line, dashed, <->] (pc) -- (esp32) node[midway, above, font=\footnotesize] {G-Code (JSON)};
        \draw[line] (esp32) -- (ls) node[midway, above, font=\footnotesize] {Step/Dir (3.3V)};
        \draw[line] (ls) -- (driver) node[midway, above, font=\footnotesize] {Step/Dir (5V)};
        \draw[line] (driver) -- (motor) node[midway, right, font=\footnotesize] {Puissance hachée};
        \draw[line, draw=red, thick] (psu) -| (driver) node[near end, left, font=\footnotesize] {36V DC};

    \end{tikzpicture}
    }
    \caption{Architecture électronique globale et logique d'amplification des signaux}
    \label{fig:arch_elec}
\end{figure}

\subsection{Dimensionnement de l'Interface de Logique (Level Shifter)}
Un phénomène mécatronique destructeur de signaux est la chute de tension relative aux optocoupleurs asociaux intégrés dans les \textit{Drivers} industriels DM556.
L'optocoupleur (ex. puce 6N137) nécessite une tension de l'ordre de $V_f \approx 1.4 \text{ V}$ avec un courant passant $I_F = 10 \text{ mA}$ pour garantir le déclenchement de la DEL infrarouge interne en moins de $50 \text{ ns}$.
Or, la résistance de limitation d'entrée montée en usine sur le DM556 est de $270 \Omega$. 

Si l'ESP32 fournit une onde carrée à $3.3 \text{ V}$ (qui chute souvent à $3.0 \text{ V}$ en charge statique) :
\begin{equation}
    I_{\text{entree}} = \frac{V_{uP} - V_f}{R} = \frac{3.0 - 1.4}{270} \approx 5.9 \text{ mA}
\end{equation}
Ce courant de $5.9 \text{ mA}$ est suffisant "en théorie" pour allumer lentement l'optocoupleur, mais le bord montant de l'onde carrée (\textit{Rise Time} $t_r$) va se ramollir considérablement (passant à plus de $2 \mu s$), induisant des retards ou des impulsions manquées.
Afin d'obtenir une onde parfaite ($\mathbf{5 \text{ V}}$) garantissant $13.3 \text{ mA}$ par canal et un $t_r \approx 10 \text{ ns}$, nous avons intégré sur la carte mère (PCB personnalisé) un circuit adaptateur de niveau à portes logiques unidirectionnelles type \textbf{SN74AHCT125N}.

\subsection{Régime Micro-stepping du DM556}
Le driver assure également l'antiparasitage des ponts en H internes commandés en PWM ($20 \text{ kHz}$ hacheur de courant constant). Nous l'avons configuré sur un sous-découpage de l'onde électrique de $1/16$ pas.
Bien que la résolution mécanique s'accroît, nous l'utilisons spécifiquement pour gommer les harmoniques de résonance du moteur de basculement A, qui, s'il entre en résonance autour de sa fréquence propre (souvent $\approx 150 \text{ Hz}$ en NEMA 17), ferait vibrer tout le portique (effet \textit{chatter} de fraisage récalcitrant).

% ========================================================
\section{Bilan de Puissance et Dimensionnement des Alimentations}

Un principe inviolable de conception est la séparation stricte de l'énergie de puissance bruyante et de la source de commande pure.

\subsection{Dimensionnement de l'Alimentation Principale (PSU 36V)}
Pour garantir la dynamique de couples calculée au Chapitre 3, les drivers moteurs sont alimentés  $36 \text{ V DC}$.
Considérons l'appel de courant en crête (lorsque tous les axes opèrent en pleine accélération - 100\% d'usage - le pire cas).
* Courant Max par moteur NEMA 17 : $2 \text{ A}$.
* Nombre de moteurs : $n = 5$.
\begin{equation}
    I_{total\_RMS} \approx 0.707 \cdot n \cdot I_{max} \approx 0.707 \times 5 \times 2 \approx \mathbf{7.07 \text{ A}}
\end{equation}
La puissance nominale à dissiper de l'alimentation SMPS (Switch-Mode Power Supply) sera donc :
\begin{equation}
    P_{PSU} = V \times I \times K_{sécu} = 36 \times 7.07 \times 1.25 \approx \mathbf{318 \text{ Watts}}
\end{equation}
Une alimentation à découpage industrielle (ex. MeanWell $350\text{W} - 36\text{V}$) a été intégrée pour le réseau exclusif des Moteurs.

\subsection{Dimensionnement des Capacités de Découplage Logique (Réseau 5V / 3.3V)}
Le deuxième réseau est le bus $5 \text{ V}$ gérant la logique (ESP32 via le LDO Interne, et Level Shifters).
Lors de la transmission d'une trame Wi-Fi par l'ESP32, la puce Radio subit un pic radiofréquence absorbant subitement une grosse quantité de courant ($\Delta I \approx 500 \text{ mA}$ sur quelques microsecondes). Si la résistance propre de la piste en cuivre est mal gérée, la loi des nœuds de Kirchhoff entraîne une chute de la tension CPU ($\Delta U = R \times \Delta I$), menant instantanément à un "Brown-out Reset" (redémarrage d'urgence) catastrophique effaçant tout l'usinage.

On prévient celà en implémentant une \textit{Bulk Capacitance} (Condensateur Chute de tension) calculée ainsi :
\begin{equation}
    C_{bulk} = I_{pic} \times \frac{\Delta t}{\Delta V_{tolere}}
\end{equation}
Pour limiter la l'effondrement de tension CPU à seulement $100 \text{ mV}$ sous un burst temporel réseau d'une durée de $\Delta t = 20 \mu s$ en déchargeant $0.5 \text{ A}$ :
\begin{equation}
    C_{bulk} = 0.5 \times \frac{20 \cdot 10^{-6}}{0.1} = 100 \cdot 10^{-6} \text{ F} = \mathbf{100 \mu \text{F}}
\end{equation}
Un gros condensateur électrolytique \textit{Low-ESR} de $470 \mu \text{F}$ couplé en parallèle à un condensateur céramique ($100 \text{ nF}$) ont été brasés directement aux pattes $V_{in}$ et $GND$ de l'ESP32 sur la carte.

% ========================================================
\section{Ingénierie de Sûreté : L'Arrêt d'Urgence E-STOP et Relais}
L'Arrêt d'Urgence est le point névralgique homologant la sécurité. En mode 5 axes, le risque de collision entre la fraise de carbure rotative (broche à $20~000$ Tr/min) et le solide socle Trunnion basculant en axe A est élevé.

\begin{figure}[htbp]
    \centering
    \begin{tikzpicture}[
        node distance=2.5cm and 2.5cm,
        block/.style={rectangle, draw, rounded corners, fill=white, text width=3.5cm, align=center, thick},
        line/.style={draw, -{Stealth[scale=1.2]}, thick}
    ]
        \node[block, fill=red!20] (estop) {\textbf{Bouton E-STOP}\\(Contact NF)};
        \node[block, right=of estop, fill=gray!20] (relais) {\textbf{Contacteur}\\(Bobine AC KM1)};
        \node[block, fill=green!10, right=of relais] (psu) {\textbf{Réseau 36V}\\(Alim. Puissance)};
        \node[block, fill=blue!10, below=1.5cm of relais] (esp) {\textbf{Micro ESP32}\\(Interrupts ISR)};

        \draw[line] (estop) -- (relais) node[midway, above, font=\footnotesize] {Rupture maintien};
        \draw[line] (relais) -- (psu) node[midway, above, font=\footnotesize] {Coupure hard};
        \draw[line, dashed] (relais) -- (esp) node[midway, right, font=\footnotesize] {Alerte Stop};

    \end{tikzpicture}
    \caption{Logigramme de la chaîne matérielle de coupure d'urgence (Catégorie 0)}
    \label{fig:estop_logic}
\end{figure}

\subsection{Câblage en "Catégorie d'Arrêt 0" (Coupure non-conditionnelle)}
Appuyer sur l'E-STOp ne peut se contenter d'envoyer un signal logique de $0$ V à l'ESP32 (Le code pourrait bugger (Deadlock) et ignorer cet appel de fonction d'interruption).
Notre intégration observe un fonctionnement par relais "Hard-Cut".
L'appui physique sur le bouton d'arrêt (NC : Normalement Fermé) brise l'auto-maintien de la bobine d'un relais de puissance AC (Contacteur). 
Ce contacteur s'ouvre magnétiquement brisant instantanément :
\begin{enumerate}
    \item La liaison $220 \text{ V} / AC$ de l'alimentation PSU 36V, figeant brutalement les balourds d'inertie magnétiques des NEMA.
    \item La liaison logique $ENABLE$ de tous les pilotes de puissance DM556, relâchant l'attraction électromagnétique interne de la cage des rotors.
    \item Un pin d'interruption vers l'ESP32 afin qu'il suspende de lui-même la lecture des coordonnées Flutter.
\end{enumerate}
Le réenclenchement mécanique manuel du champignon rouge est obligatoirement requis, couplé à un appui sur le bouton "Reset/Resume" de l'armoire, interdisant de ce fait tout redémarrage inopiné intempestif préjudiciable physiquement.

% ========================================================
\section{Conclusion sur l'Étude Électronique}
L'infrastructure mécatronique mise en œuvre surcarte n'est pas qu'un amalgame de câbles : elle conditionne la vie ou la mort du signal algorithmique.  C’est seulement par l'optimisation des fronts temporels avec des Level Shifters, l'étanchéité absolue de l’isolement galvanique, et la provision d'un courant ininterrompu via condensateur, que nous autorisons le système d'exploitation de l'ESP32 à se consacrer intégralement à son immense tache conceptuelle : le calcul de la Cinématique Inverse pure. Le Chapitre 5 qui suit se plonge précisément au cœur de l'algèbre algorithmique de ce microcontrôleur dual-core et de son pendant Interface (l'IHM Flutter).     % Architecture Électronique et Mécatronique
% !TEX root = ../main.tex
\chapter{Réalisation, Assemblage, Tests et Conclusion}
\label{chap:tests_conclusion}

\section{Réalisation Mécanique et Électrique}

\subsection{Assemblage du Portique et du Plateau Trunnion}
L'assemblage a suivi un protocole métrologique strict pour garantir l'\textit{équerrage} (Squaring) parfait de la machine.
Le Plateau Trunnion (Axes A et B) constitue la phase critique. L'axe de rotation du plateau B doit croiser parfaitement (coplanarité) l'axe de basculement A. Tout décalage (offset $L_x$ induit) ruinerait la cinématique inverse. Le réducteur a été pré-contraint pour annuler le jeu d'inversion mécanique.

\subsection{Câblage et Intégration du Coffret Électrique}
Toute la "nappe nerveuse" (câbles moteurs et logiques) a été acheminée en dissociant la puissance (36V) et la logique (5V/3.3V). Chaque tresse de blindage a été mise à la terre en un point unique (Star Grounding) dans le boitier pour combattre les boucles de masse.

\section{Étalonnage Logiciel (Homing et Calibration)}

\subsection{Calibration des Steps par Millimètre}
La conversion Théorie/Réalité nécessite le calibrage des impulsions logiques.
Pour la vis à billes (pas de 5mm, $1/16$ microstepping) :
\begin{equation}
    Steps/mm = \frac{N_{pas/tour} \times Microstepping}{Pitch_{vis\_a\_billes}} = \frac{200 \times 16}{5} = 640 \text{ Steps/mm}
\end{equation}

\subsection{L'Étalonnage du Pivot Point (Offsets Cinématiques)}
C'est le sommet de la calibration 5 axes. La formule de Denavit-Hartenberg implémentée dépend de constantes d'excentricité ($L_y, L_z$). Celles-ci doivent être mesurées physiquement avec un Palpeur pour localiser mathématiquement le Centre Absolu de Rotation du Plateau B. Ces coordonnées ($TCP$) sont transmises au parser G-Code de l'application Flutter via Riverpod avant de lancer l'usinage.

\section{Tests d'usinage et Validation des Performances}

\subsection{Test 1 : Validation du contournage 3 Axes}
Gravure de circuits électroniques (PCB FR4) avec un foret "V-Bit" à $0.1 \text{ mm}$ de profondeur. Le résultat est net, sans "chatter", prouvant la forte rigidité verticale (Axe Z).

\subsection{Test 2 : Validation du 3+2 Axes}
Usiner les cinq faces visibles d'un cube sans démonter la pièce de l'étau. Le plateau B tourne de $90^\circ$, la broche usine la face 1. L'axe A bascule de $45^\circ$, la broche usine un chanfrein. Les trous sont parfaitement alignés. L'IHM Flutter n'a connu aucune désynchronisation depuis le stream Riverpod UART.

\subsection{Test 3 : Usinage simultané 5 Axes "Continu" (La Turbine)}
Fichier G-Code de plusieurs Mo contenant uniquement des micro-linéaires (`G1 X.. Y.. Z.. A.. B..`). L'algorithme Dual-Core a maintenu le Ring Buffer plein. La tête ne s'est pas arrêtée dans les virages. L'interface Flutter a restitué la visualisation transparente sur Android, confirmant la stabilité absolue de l'architecture distribuée.

\section*{Conclusion Générale et Perspectives}
\addcontentsline{toc}{section}{Conclusion Générale et Perspectives}

L'aboutissement de ce Projet de Fin d'Études marque la transition maîtrisée d'une conception théorique exigeante vers une matérialisation mécatronique fonctionnelle. L'objectif initial était ambitieux : démocratiser l'accès à une technologie d'industrie 4.0 lourde — l'usinage à 5 axes — par la conception locale d'une architecture globale "Desktop".

Le pari logiciel a été un grand succès : s'affranchissant des logiciels de commande obscurs ou limitatifs (Arduino 8-bits), la symbiose entre un Sender développé en \textbf{Flutter} (gérant son état via \textbf{Riverpod} de façon asynchrone) et un Firmware C++ sur l'\textbf{ESP32} (gérant le Buffer temporel et les complexes Matrices Trunnion de Denavit-Hartenberg) s'est avérée extrêmement fluide. Les tests valident cette architecture par la réalisation de parcours 5 axes sans saccades.

\textbf{Perspectives Evolutives :}
L'excellence industrielle appelle des améliorations continues. 
\begin{itemize}
    \item \textit{Mécanique :} Le remplacement des moteurs Open Loop par des Servomoteurs à codeurs (boucle fermée).
    \item \textit{Mécatronique :} L'ajout d'un Changeur d'Outil Automatique (ATC - Automatic Tool Changer).
    \item \textit{IHM / I.A. :} L'application Flutter pourrait monitorer la chute de courant d'un moteur pour extrapoler l'usure de l'outil de coupe, réalisant la promesse de la maintenance prédictive inhérente à l'Industrie 4.0.
\end{itemize}

Cette mini-fraiseuse CNC 5 axes n'est pas uniquement un équipement technique matériel, c'est un véritable laboratoire modulaire de R\&D local.
  % Développement Logiciel et Cinématique UML
% !TEX root = ../main.tex
\chapter{Réalisation, Assemblage, Tests et Qualification}
\label{chap:tests_conclusion}

La transition entre un jumeau numérique CAO ou une théorie de commande mathématique (aérienne) et une machine industrielle concrète ("au sol") est jonchée d'écueils physiques imperceptibles sur ordinateur. La flexion, les jeux thermiques et résiduels, ainsi que le bruit électromagnétique, exigent un processus d'étalonnage et de tests itératifs intensifs. Ce dernier chapitre solde la conception en confrontant la fraiseuse CNC 5 axes à l'épreuve de la métrologie et de l'usinage réel.

% ========================================================
\section{Processus d'Assemblage Métrologique}

Contrairement à une imprimante 3D (FDM) tolérant des désalignements structurels compensables dynamiquement par un capteur Z (type BLTouch), l'usinage par soustraction de matière 5 axes est incompensable si le bâti (la machine) n'est pas \textit{orthonormé} à l'origine.

\subsection{L'Équerrage du Portique (Squaring process)}
L'assemblage des V-Slots de la base Y et du portique X/Z s'est fait sur un plan de travail certifié (Marbre). 
\begin{enumerate}
    \item \textbf{Étalonnage X/Y :} Un jeu d'équerres mécaniques de précision classe 00 (DIN 875) a été positionné entre les traverses. Un serrage au couple via clé dynamométrique a été effectué pour garantir une perpendicularité spatiale $\perp_{XY} < 0.02 \text{ mm/m}$.
    \item \textbf{Calibration de la broche :} Un support à comparateur horizontal à cadran (Micromètre) fixé sur le plateau a permis d'usiner virtuellement "dans le vide". Le balayage manuel de la tête en Z confirme que l'outil est strictement tangentiel à l'axe Z (aucun faux-rond de fixation "Tramming" n'excédant $0.015 \text{ mm}$).
\end{enumerate}

\begin{figure}[htbp]
    \centering
    \begin{tikzpicture}[
        tool/.style={rectangle, draw, fill=gray!20, thick},
        comp/.style={circle, draw=black, fill=white, inner sep=0pt, minimum size=0.8cm, thick},
        line/.style={thick}
    ]
        % Base/Marbre
        \draw[thick, fill=blue!5] (-3,0) rectangle (3,-0.5);
        \node at (0,-0.25) {Marbre de métrologie (Base de référence)};
        
        % Equerre
        \draw[tool, fill=orange!20] (0.5,0) -- (0.5,2.5) -- (1,2.5) -- (1,0.5) -- (2.5,0.5) -- (2.5,0) -- cycle;
        \node[font=\footnotesize, align=center] at (1.8, 0.25) {Équerre\\DIN 875};
        
        % Broche Z
        \draw[tool, fill=gray!40] (-1.5,1) rectangle (-0.5,3);
        \node[font=\footnotesize, rotate=90] at (-1, 2) {Axe Z (Broche)};
        
        % Comparateur
        \draw[line] (-0.5,1.5) -- (0,1.5); % bras
        \node[comp] at (0.2, 1.5) {}; 
        \draw[-{Stealth[scale=1]}, thick, red] (0.4, 1.5) -- (0.5, 1.5); % palpeur
        \node[above, font=\scriptsize] at (0.2, 1.9) {Comparateur};
        
        % Mouvement
        \draw[<->, thick, blue, dashed] (-1.5,1.5) -- (-1.5,2.5) node[midway, left] {Balayage $\Delta Z$};
        
    \end{tikzpicture}
    \caption{Principe de mesure de la perpendicularité de la broche par balayage (Tramming)}
    \label{fig:squaring}
\end{figure}

\subsection{Câblage en Étoile (\textit{Star Grounding}) et Anti-Interférences}
Toute la "nappe nerveuse" (câbles codeurs, signaux \textit{Step/Dir}) a été acheminée dans des gaines (Drag Chains) de manière physiquement distante (séparation de 10 cm minimum) des câbles de puissance (36V AC du Spindle). L'innovation électrotechnique repose sur la \textit{mise à la terre en étoile} : la protection faradienne (\textit{Shield}) des câbles de la broche (Moteur VFD) n'est reliée à la terre ($GND$) que \textbf{d'un seul côté} (au niveau du coffret d'entrée), l'autre extrémité (côté tête d'usinage) étant laissée flottante. Cela rompt la boucle de courant harmonique (Ground Loop) qui parasitait l'ESP32.

% ========================================================
\section{Étalonnage Logiciel (Homing et Offsets Cinématiques)}

\subsection{La conversion logico-mécanique ($Steps/mm$)}
Le microcontrôleur ESP32 ne comprend que des impulsions (Ticks). Les valeurs saisies dans sa mémoire EEPROM établissent le rapport de transformation entre la rotation paramétrique du moteur et le millimètre ISO usiné.

Pour l'axe X (Vis à billes de pas $5\text{mm}$, Moteur $1.8^\circ=200\text{ Pas/tr}$, Microstepping du DM556 paramétré à $1/16$) :
\begin{equation}
    Steps/mm = \frac{N_{pas/tour} \times \mu{Step}}{P_h} = \frac{200 \times 16}{5} = \textbf{640 \text{ Impulsions/mm}}.
\end{equation}

Pour le plateau Trunnion rotatif Axe B (Réduction harmonique ratio $1:50$) :
\begin{equation}
    Steps/\text{Degré} = \frac{200 \times 16 \times 50}{360^\circ} = \textbf{444.444 \text{ Impulsions/Degré}}.
\end{equation}
L'exactitude "flottante" de ce paramétrage confirme à l'ESP32 l'extrême finesse de déclenchement d'une interpolation 5 axes continue.

\subsection{Palpage du \textit{Tool Center Point} (Offsets Trunnion)}
Le logiciel Riverpod (sous Flutter) nécessite l'intégration des variables globales de Denavit-Hartenberg $\mathbf{L_y}$ et $\mathbf{L_z}$ modélisées au Chapitre 5. 
Pour ce faire, une pige de centrage (Edge Finder mécanique à ressort) montée à 1000 Tr/min sur la broche a été abaissée lentement sur le bloc Trunnion. La désaxation du plateau lors d'une rotation B manuelle a permis de localiser "l'oeil du cyclone" (X, Y) et "l'apex" (Z) de l'axe de basculement A.
\begin{itemize}
    \item \textbf{Mesure retenue $L_{y}$ :} $34.50 \text{ mm}$ (décalage entre broche Z et centre de pivot A).
    \item \textbf{Mesure retenue $L_{z}$ :} $112.80 \text{ mm}$ (hauteur du plateau vis-à-vis du pivot).
\end{itemize}

% ========================================================
\section{Validation Métrologique par l'Usinage (Tests Pratiques)}

\subsection{Essai 1 : Contourage 2.5D (Rigidité et Répétabilité)}
\textbf{Protocole :} Usinage d'une poche rectangulaire (pocketing) et d'un alésage dans un bloc d'aluminium AW-2017A avec une fraise de $6\text{ mm}$ (Avance : $400 \text{ mm/min}$, Passe de $0.5 \text{ mm}$).
\textbf{Analyse des Résultats :} Le pied à coulisse au $1/50^{ème}$ constate une tolérance d'erreur de positionnement $\Delta_{moyen} = \pm 0.04 \text{ mm}$. La pièce n'affiche pas de trace d'arrachement par vibration (Broutage). Les vis sphériques à billes absorbent victorieusement les pics d'accélération inertielle du système sans flexion notable du portique (validant la RDM).

\subsection{Essai 2 : Usinage Positionnel "3+2 axes"}
\textbf{Protocole :} Opération multi-faces sur une pièce fixée une seule fois. Le plan G-Code pivote le plateau de $90^\circ$ (Axe B) pour le fraisage d'une gorge latérale, puis bascule de $45^\circ$ (Axe A) pour chanfreiner une arrête supérieure.
\textbf{Analyse des Résultats :} Tous les contours usinés se croisent ou s'alignent parfaitement. La transmission réseau UDP du G-Code depuis l'application tablette Flutter s'est maintenue stable durant les 27 minutes du cycle, confirmant la pertinence d'utiliser FreeRTOS (Dual Core).

\subsection{Essai 3 : Interpolation Simultanée 5 Axes "Continue"}
\textbf{Protocole :} Usinage surfacique complexe d'un dôme asymétrique et d'une pale d'hélice dans un bloc de résine dense (Cibatool). 
Les 5 moteurs tournent simultanément à des vitesses variées. L'outil "glisse" spatialement en pivotant avec la pièce pour rester toujours perpendiculaire (normale $I,J,K$) à la surface conique usinée.
\textbf{Analyse des Résultats :} Le Ring Buffer (Look-Ahead) du C++ a parfaitement lissé les micro-décélérations. Les drivers s'échauffent au sein de leur plage thermale tolérée ($< 55^\circ$C avec brassage de l'air de l'armoire). La fluidité spectaculaire des mouvements de la table Trunnion prouve définitivement la capacité de la matrice mathématique inversée à rattraper les dérives de basculement en temps réel.


\begin{figure}[htbp]
    \centering
    \begin{minipage}{0.31\textwidth}
        \centering
        \includegraphics[width=\linewidth, draft]{placeholder_test1_2D.png}
        \caption{Contourage 2.5D}
        \label{fig:test1}
    \end{minipage}\hfill
    \begin{minipage}{0.31\textwidth}
        \centering
        \includegraphics[width=\linewidth, draft]{placeholder_test2_3plus2.png}
        \caption{Chanfreins en 3+2}
        \label{fig:test2}
    \end{minipage}\hfill
    \begin{minipage}{0.31\textwidth}
        \centering
        \includegraphics[width=\linewidth, draft]{placeholder_test3_5axes.png}
        \caption{Dôme continu 5 axes}
        \label{fig:test3}
    \end{minipage}
\end{figure}

% ========================================================
\section{Analyse Technico-Économique (BOM)}

La viabilité de ce projet repose sur la démonstration qu'une technologie complexe (Usinage 5 axes), historiquement onéreuse, peut être reproduite à moindre coût sans sacrifier la précision industrielle.

\begin{table}[htbp]
    \centering
    \renewcommand{\arraystretch}{1.2}
    \begin{tabular}{|l|p{4cm}|c|c|r|}
    \hline
    \textbf{Catégorie} & \textbf{Désignation Composant} & \textbf{Qté} & \textbf{PU (FCFA)} & \textbf{P. Total (FCFA)} \\ \hline
    \rowcolor{gray!10} 
    \textbf{Structure Mécanique} & Rails HGR15 + Vis SFU1605 & 3 & 45 000 & 135 000 \\ \hline
    \rowcolor{gray!10}
    & V-Slot Alu 2040 \& 4040 & 4 & 10 000 & 40 000 \\ \hline
    \rowcolor{gray!10}
    & Broche ER11 500W & 1 & 65 000 & 65 000 \\ \hline
    \rowcolor{blue!5}
    \textbf{Électronique} & Moteurs NEMA 17 (1.7A) & 5 & 12 000 & 60 000 \\ \hline
    \rowcolor{blue!5}
    & Drivers DM556 & 5 & 9 000 & 45 000 \\ \hline
    \rowcolor{blue!5}
    & ESP32 + Isolation Optique & 1 & 18 000 & 18 000 \\ \hline
    \rowcolor{blue!5}
    & Alimentation 36V 350W & 1 & 25 000 & 25 000 \\ \hline
    \rowcolor{green!10}
    \textbf{Trunnion \& Transm.} & Réducteur Courroie (1:50) & 2 & 32 000 & 64 000 \\ \hline
    \rowcolor{green!10}
    & Roulements Croisés & 4 & 8 500 & 34 000 \\ \hline
    \rowcolor{yellow!10}
    \textbf{Divers} & Visserie, Câbles, E-STOP & - & 40 000 & 40 000 \\ \hline
    \hline
    \multicolumn{4}{|r|}{\textbf{COÛT GLOBAL ESTIMÉ DE REVIENT DU PROTOTYPE}} & \textbf{526 000 FCFA} \\ \hline
    \end{tabular}
    \caption{Nomenclature Financière (Bill of Materials) détaillée en Francs CFA}
    \label{tab:bom}
\end{table}

\textbf{Amortissement et Comparatif :} Une machine CNC 5 axes éducative équivalente importée (Ex: PocketNC V2) est facturée à plus de 4~000~000 FCFA (frais de port et douanes inclus). Notre solution souveraine, conçue intégralement à l'ENI-ABT pour un coût unitaire autour de $\sim$ 526 000 FCFA, représente une \textbf{économie de 86\%}. Ce projet devient ainsi un formidable levier industriel pour les PME locales (Prototypage plastique, bois, aluminium doux).

% ========================================================
\section{Sécurité Industrielle et Analyse des Risques (AMDEC)}

Toute conception de rang ingénieur exige la garantie de l'intégrité de l'opérateur et de l'outil machine. L’Analyse des Modes de Défaillance, de leurs Effets et de leur Criticité (AMDEC) justifie la présence des capteurs matériels et de la logique "Fail-Safe".

\begin{table}[htbp]
    \centering
    \resizebox{\textwidth}{!}{
    \renewcommand{\arraystretch}{1.3}
    \begin{tabular}{|p{3cm}|p{3.5cm}|p{3cm}|p{4cm}|p{3cm}|}
    \hline
    \rowcolor{macouleur!20}
    \textbf{Mode de défaillance} & \textbf{Causes Possibles} & \textbf{Effets sur Système} & \textbf{Mesure Corrective (Solution)} & \textbf{État ESP32} \\ \hline
    Accélération brusque & G-Code erroné (Avance \texttt{F} aberrante) & Cassure de la fraise & Parseur Flutter rejette le bloc avant l'envoi. & \texttt{IDLE} \\ \hline
    Hors limites physiques (Crash Axes) & Perte de pas, Offset Origine erroné & Torsion structurelle, Destruction Broche & Capteurs Fin de Course (Hard Limits) déclenchant une interruption. & \texttt{ALARM} \\ \hline
    Coupure de courant & Réseau instable & Outil coincé, Pièce perdue & Arrêt gravé (G-Code ligne sauvegardée dans EEPROM optionnelle). & - \\ \hline
    Intrusion dans l'espace de travail & Erreur de l'opérateur local & \textbf{Blessure grave} (Lame 10~000 Tr/min) & Coupure \textbf{hardware} du Relais via E-STOP + Signal Pull-UP. & \texttt{ALARM} (Stop) \\ \hline
    Surchauffe Moteur & Friction mécanique accrue & Perte magnétisme, Incendie & Abaissement manuel V-Ref + Ventilation du coffret actif. & \texttt{RUN} \\ \hline
    \end{tabular}
    }
    \caption{Tableau de synthèse AMDEC des sécurités embarquées de la machine}
    \label{tab:amdec}
\end{table}

% ========================================================
\section*{Conclusion Générale et Perspectives}
\addcontentsline{toc}{section}{Conclusion Générale et Perspectives}

L'aboutissement de ce Projet de Fin d'Études cristallise la transition aboutie d'une conception théorique exigeante vers un équipement mécatronique de précision d'industrie 4.0 fonctionnel. L'objectif initial – l'abstraction de la difficulté logicielle et du coût effarant d'une machine CNC 5 axes propriétaire pour les structures éducatives maliennes – s'est imposé comme le fil conducteur des choix d'ingénierie.

\textbf{Bilan Opérationnel et Technique :}
La démarche a validé successivement des ruptures technologiques vis-à-vis du prototypage conventionnel :
\begin{itemize}
    \item \textbf{Sur la mécanique (Chapitres 2 et 3) :} L'emploi systématique de composants industriels durcis (Guidages linéaires précontraints et Vis à billes roulées) a permis d'éradiquer la flexion mortelle induite par le lourd couple de renversement ($20 \text{ Nm}$) du centre de gravité mobile de la table Trunnion basculante.
    \item \textbf{Sur la mécatronique de supervision (Chapitres 4 et 5) :} Nous avons brisé le monopole des vieux microcontrôleurs 8-bits en démontrant l'efficacité retentissante du \textbf{processeur temps réel ESP32} couplé au système de \textit{Tasking} FreeRTOS. Le cerveau exécute avec véhémence des transformations de matrices homogènes ($x, y, z$ corrigés dynamiquement par $\cos/\sin$ selon A et B) afin de maintenir l'outil sur le point de contact TCP (Tool Center Point) nominal.
    \item \textbf{Sur l'interaction logicielle (Flutter / Riverpod) :} La conception \textit{From Scratch} d'un interpréteur asynchrone G-code et interface Homme-Machine (IHM) multi-plateformes, affranchie des saturations graphiques grâce au framework Riverpod de Google, offre au client final (l'opérateur) un portail de communication TCP/IP stable et esthétique, remplaçant la boite noire par un code auditable (Open-Source).
\end{itemize}

\textbf{Perspectives d'Innovation :}
Si ce PFE atteint ses objectifs cinématiques, ce laboratoire modulaire appelle naturellement une myriade d'évolutions incrémentales, tant sur le plan mécanique que dans le paradigme de l'Internet des Objets (IoT) :
\begin{enumerate}
    \item \textbf{Évolution de l'Atelier Logiciel (I.A. et Maintenance Prédictive) :} En branchant des résistances de type de Shunt sur l'alimentation 36V, l'ESP32 pourrait capter en temps réel les pics de courant consommés lors du fraisage de matière difficile, les numériser pour les remonter sur le réseau jusqu'à une \textit{Intelligence Artificielle} sur l'APP Flutter, laquelle prédirait l'usure prématurée d'une lame ou d'une avarie broche (Condition Monitoring).
    \item \textbf{Évolution de Servocommande :} Actuellement paramétrée en "Boucle Ouverte", la machine pourrait accueillir des servomoteurs "Closed-Loop" dotés d'encodeurs rotatifs absolus. Couplés au Firmware C++, ils compenseraient en temps-masqué une éventuelle perte de pas lors d'une collision imprévue, sauvant une pièce de plusieurs heures d'usinages.
    \item \textbf{Ergonomie "Industry 5.0" et Automatisation :} L'ajout d'une station logicielle générant la commutation pneumatique pour une broche dotée d'un "Automatic Tool Changer" (ATC) et d'un palpeur autonome de numérisation 3D sans contact.
\end{enumerate}

Le passage à l'usinage $5$ axes n'est plus, par ce projet, un tabou d'ingénierie exclusif. Cette mini-fraiseuse en constitue le socle tangible : une machine capable de former la prochaine génération de cadres concepteurs-usiniers à la matrice des contraintes cinématiques de l'industrie de pointe africaine.
        % Tests, Validation et Analyse Économique

% ===================== BACKMATTER =====================
% !TEX root = main.tex

% ========================================================
% BIBLIOGRAPHIE
% ========================================================
\newpage
\addcontentsline{toc}{chapter}{Bibliographie et Nétographie}
\begin{thebibliography}{99}

\bibitem{machining} 
\textbf{Smid, Peter.} \textit{CNC Programming Handbook, Third Edition}. Industrial Press Inc., 2007. \newline
(Référence canonique sur la structure et le parsing du G-Code ISO 6983).

\bibitem{robotics} 
\textbf{Craig, John J.} \textit{Introduction to Robotics: Mechanics and Control}. Pearson Prentice Hall, 2005. \newline
(Méthode de Denavit-Hartenberg et calculs matriciels de cinématique inverse).

\bibitem{rdm} 
\textbf{Timoshenko, Stephen P.} \textit{Résistance des matériaux}. Dunod. \newline
(Calcul des flèches sur les guides de translation HGR15).

\bibitem{esp32} 
\textbf{Espressif Systems.} \textit{ESP32 Technical Reference Manual}, 2023. \newline
\texttt{https://www.espressif.com/en/support/documents/technical-documents} \newline
(Architecture Dual-Core et FreeRTOS).

\bibitem{flutter} 
\textbf{Dev Google.} \textit{Flutter Architectural Overview}. \newline
\texttt{https://docs.flutter.dev/resources/architectural-overview} \newline
(Documentation sur le Render Tree, StateNotifier et le pattern Riverpod).

\bibitem{openbuilds}
\textbf{OpenBuilds Community.} \textit{Open Source Machine Design Archives}. \newline
\texttt{https://openbuilds.com/} \newline
(Inspiration pour les sous-assemblages de la géométrie V-Slot des profils aluminium).

\end{thebibliography}

% ========================================================
% ANNEXES
% ========================================================
\newpage
\appendix
\chapter{Annexes techniques}

\section{Extrait de l'algorithme "Look-Ahead" du Planner (C++)}
Cet annexe présente une portion simplifiée de l'accélération trapézoïdale programmée sur le Core 1 de l'ESP32 pour lisser les micro-segments du G-Code 5 axes.

\begin{tcolorbox}[colback=white, colframe=macouleur, title=Planner.cpp (Extrait)]
\begin{verbatim}
void plan_buffer_line(float x, float y, float z, float a, float b, float feed_rate) {
    // Calcul du vecteur déplacement 5D (Espace Machine)
    float dx = x - position[X_AXIS];
    float dy = y - position[Y_AXIS];
    float dz = z - position[Z_AXIS];
    float da = a - position[A_AXIS];
    float db = b - position[B_AXIS];

    float distance = sqrt(dx*dx + dy*dy + dz*dz + da*da + db*db);

    // Blocage si déplacement inférieur à la tolérance mécanique (0.01mm)
    if (distance < MINIMUM_PLANNER_SPEED) return;

    // Calcul de la vitesse de croisière cible (Target Speed)
    float nominal_speed = feed_rate;
    
    // Détection d'angle aigu (Jonction) entre l'ancien bloc et le nouveau
    float junction_cos_theta = calculate_junction_deviation(previous_vector,
                                                            current_vector);
                                                            
    // Ralentissement dynamique aux angles (Cornering Speed)
    float max_junction_speed = min(nominal_speed, 
                                   sqrt(JUNCTION_DEVIATION * acceleration 
                                   / (1.0 - junction_cos_theta)));

    // Insertion dans le Ring Buffer
    block_buffer[head].nominal_speed = nominal_speed;
    block_buffer[head].entry_speed = max_junction_speed;
    
    // Avance de l'index du Buffer Circulaire
    head = (head + 1) % BLOCK_BUFFER_SIZE;
}
\end{verbatim}
\end{tcolorbox}
                 % Bibliographie et Annexe Code (Look-Ahead)

% ===================== ANNEXES ========================
\newpage
\appendix
% !TEX root = ../main.tex
\chapter{Annexes Techniques et Documentation}
\label{annexe:documentation}

\section{A.1 — Code Source Commenté du Firmware ESP32}

\subsection{A.1.1 — Fichier Principal : \texttt{main.cpp}}

\begin{tcolorbox}[colback=gray!5, colframe=macouleur, title=\texttt{main.cpp} -- Point d'entrée du firmware ESP32]
\begin{verbatim}
/**
 * @file main.cpp
 * @brief Firmware CNC 5-axes pour ESP32 (configuration Trunnion Table)
 * @author Étudiant PFE - ENI-ABT, Génie Mécanique et Énergie
 * @date 2025
 * @hardware ESP32-WROOM-32U, DM556 x5, NEMA17 x5, MeanWell 36V/350W
 *
 * ARCHITECTURE DUAL-CORE:
 *   - Core 0: Réseau Wi-Fi (UDP), parsing JSON, cinématique inverse
 *   - Core 1: Hardware Timer ISR (Step/Dir), Ring Buffer, Planner
 */

#include <Arduino.h>
#include <WiFi.h>
#include <AsyncUDP.h>
#include "GCodeParser.h"       // Parser G-Code ISO 6983
#include "Planner5Axes.h"      // Planificateur Look-Ahead
#include "InvKinematics.h"     // Cinématique Inverse Trunnion
#include "StepGenerator.h"     // Générateur d'impulsions Step/Dir
#include "SafetyMonitor.h"     // Gestion E-STOP et fins de course

// ---- Configuration Réseau Wi-Fi ----
const char* WIFI_SSID = "CNC5AXES_AP";
const char* WIFI_PASS = "cnc12345";    // Réseau AP propre à la machine
const uint16_t UDP_PORT = 2222;

// ---- Instances Globales ----
AsyncUDP udp;
GCodeParser  parser;
Planner5Axes planner;
InvKinematics ik;
StepGenerator stepGen;
SafetyMonitor safety;

// Tâche réseau sur Core 0 (basse priorité)
void taskNetwork(void* pvParameters) {
    WiFi.softAP(WIFI_SSID, WIFI_PASS);
    udp.listen(UDP_PORT);
    
    udp.onPacket([](AsyncUDPPacket packet) {
        // Extraire la charge utile JSON
        String payload((char*)packet.data(), packet.length());
        
        // Valider le CRC-16 du paquet
        if (!validateCRC(payload)) return;
        
        // Parser le bloc G-Code
        GCodeBlock block = parser.parseBlock(payload);
        
        // Calcul cinématique inverse (TCP compensé)
        MotorPositions motors = ik.solve(
            block.x, block.y, block.z,
            block.a, block.b);
        
        // Insérer dans le Ring Buffer du Planner
        planner.bufferBlock(motors, block.feedRate);
        
        // Envoyer ACK avec position courante
        String ack = buildACKJson(ik.getCurrentPosition());
        packet.print(ack);
    });
    
    // Boucle infinie (la pile réseau gère l'événementiel)
    while(true) vTaskDelay(100 / portTICK_PERIOD_MS);
}

// Tâche temps-réel sur Core 1 (priorité maximale)
void taskRealTime(void* pvParameters) {
    // Configurer les Hardware Timers pour Step Generation
    stepGen.init();
    safety.init();
    
    while(true) {
        // Vérifier les fins de course et E-STOP
        if (safety.isEmergencyStop()) {
            stepGen.disableAllAxes();
            planner.clearBuffer();
            // Attendre l'acquittement de l'opérateur
            while(safety.isEmergencyStop()) vTaskDelay(1);
        }
        
        // Exécuter le bloc suivant s'il est disponible
        if (planner.hasBlock()) {
            MotionBlock block = planner.popBlock();
            stepGen.executeBlock(&block);  // Bloquant
        }
        
        vTaskDelay(1 / portTICK_PERIOD_MS); // Yield 1ms
    }
}

void setup() {
    Serial.begin(115200);
    
    // Démarrer les deux tâches sur leurs cœurs assignés
    xTaskCreatePinnedToCore(taskNetwork,  "NetTask",  8192, NULL, 1, NULL, 0);
    xTaskCreatePinnedToCore(taskRealTime, "RTTask",   4096, NULL, 5, NULL, 1);
    
    // Déléguer immédiatement à FreeRTOS Scheduler
}

void loop() { vTaskDelete(NULL); } // Loop non utilisée
\end{verbatim}
\end{tcolorbox}

\subsection{A.1.2 — Algorithme de Cinématique Inverse : \texttt{InvKinematics.cpp}}

\begin{tcolorbox}[colback=gray!5, colframe=macouleur, title=\texttt{InvKinematics.cpp} -- Calcul des compensations TCP Trunnion]
\begin{verbatim}
/**
 * @brief Résout la cinématique inverse de la table Trunnion.
 * Transforme les coordonnées pièce (x_w, y_w, z_w, A°, B°)
 * en positions machine corrigées (x_m, y_m, z_m).
 *
 * ÉQUATIONS MATHÉMATIQUES (Trunnion A+B):
 *   R_plateau = rayon du plateau = 100 mm
 *   H_berceau = hauteur du pivot sous la table = 80 mm
 *
 *   dZ = R * (1 - cos(A_rad))       // Décalage TCP en Z
 *   dY = R * sin(A_rad)              // Décalage TCP en Y
 *
 *   x_m = x_w   - dY * (-sin(B_rad)) // Projection X par rotation B
 *   y_m = y_w   - dY * cos(B_rad)    // Projection Y
 *   z_m = z_w   - dZ                 // Correction Z
 */
MotorPositions InvKinematics::solve(float x_w, float y_w, float z_w,
                                     float A_deg, float B_deg) {
    const float R = PLATEAU_RADIUS_MM;  // 100.0f mm
    
    // Conversion Degrés -> Radians
    float A_rad = A_deg * DEG_TO_RAD;
    float B_rad = B_deg * DEG_TO_RAD;
    
    // Calcul des offsets TCP (Table Trunnion A+B)
    float dZ_tcp = R * (1.0f - cosf(A_rad));
    float dY_tcp = R * sinf(A_rad);
    
    // Application de la rotation B sur les offsets (matrice 2x2)
    float dx_machine = dY_tcp * (-sinf(B_rad));
    float dy_machine = dY_tcp * cosf(B_rad);
    
    // Positions machine finales
    MotorPositions motors;
    motors.x = x_w - dx_machine;
    motors.y = y_w - dy_machine;
    motors.z = z_w - dZ_tcp;
    motors.a = A_deg;  // Angles moteurs A et B
    motors.b = B_deg;  // = angulations demandées
    
    return motors;
}
\end{verbatim}
\end{tcolorbox}

\subsection{A.1.3 — Interface Utilisateur Flutter : \texttt{dashboard\_page.dart}}

\begin{tcolorbox}[colback=gray!5, colframe=macouleur, title=\texttt{dashboard\_page.dart} -- Écran principal IHM Flutter]
\begin{verbatim}
/// Page principale de l'IHM - Tableau de bord de la CNC 5 Axes
/// Architecture : Riverpod v2 + Flutter 3.x
/// Widget Stateless consommant les états depuis les Providers
class DashboardPage extends ConsumerWidget {
  const DashboardPage({super.key});

  @override
  Widget build(BuildContext context, WidgetRef ref) {
    // Écoute réactive de l'état complet de la machine
    final machineState = ref.watch(machineStateProvider);
    final coordinates = ref.watch(coordinatesProvider);
    final connectionStatus = ref.watch(connectionStatusProvider);

    return Scaffold(
      backgroundColor: const Color(0xFF0A0F1E),
      body: Row(
        children: [
          // === PANNEAU GAUCHE : Coordonnées & Status ===
          SizedBox(
            width: 280,
            child: Column(
              children: [
                // En-tête Machine
                MachineHeaderWidget(
                  machineName: 'CNC-5AXES / ENI-ABT',
                  status: connectionStatus,
                ),
                // Affichage des 5 axes (monospace, grand)
                for (final axis in ['X', 'Y', 'Z', 'A', 'B'])
                  AxisCoordinateCard(
                    axisLabel: axis,
                    value: coordinates.getAxis(axis),
                    isRotary: axis == 'A' || axis == 'B',
                  ),
              ],
            ),
          ),
          
          // === PANNEAU CENTRE : Visualisateur 3D G-Code ===
          Expanded(
            child: GCodeViewer3D(
              gcodeBlocks: ref.watch(gcodeQueueProvider),
              currentLineIndex: machineState.currentLine,
            ),
          ),
          
          // === PANNEAU DROITE : Contrôles ===
          SizedBox(
            width: 300,
            child: Column(
              children: [
                // Boutons Action Principaux
                CycleControlButtons(
                  onStart: () => ref.read(
                    cycleControlProvider.notifier).startCycle(),
                  onPause: () => ref.read(
                    cycleControlProvider.notifier).feedHold(),
                  onEStop: () => ref.read(
                    safetyProvider.notifier).triggerEStop(),
                ),
                // Contrôle Jog Manuel
                const JogPadWidget(),
              ],
            ),
          ),
        ],
      ),
    );
  }
}
\end{verbatim}
\end{tcolorbox}

% ============================================
\section{A.2 — Fiches Techniques des Composants Principaux}

\subsection{A.2.1 — Extrait de la Fiche Technique ESP32-WROOM-32U (Espressif Systems)}

\begin{table}[htbp]
\centering
\renewcommand{\arraystretch}{1.3}
\begin{tabular}{|l|l|}
\hline
\rowcolor{macouleur!20}
\textbf{Paramètre} & \textbf{Valeur} \\
\hline
Processeur & Xtensa LX6 Dual-Core 32-bit \\
\hline
Fréquence CPU & 80 / 160 / 240 MHz (configurable) \\
\hline
SRAM interne & 520 KB \\
\hline
Flash embarquée & 4 MB (SPI Flash) \\
\hline
Wi-Fi & 802.11 b/g/n 2.4 GHz, 150 Mbps \\
\hline
Bluetooth & BT Classic 4.2 + BLE \\
\hline
GPIO disponibles & 34 broches (28 configurables) \\
\hline
Hardware Timers & 4 × 64 bits haute résolution \\
\hline
ADC & 18 canaux 12 bits (0--3.6V) \\
\hline
DAC & 2 canaux 8 bits \\
\hline
Tension d'alimentation & $3{,}0$ -- $3{,}6$ V (Typique $3{,}3$ V) \\
\hline
Courant de veille (Deep Sleep) & $10\,\mu A$ \\
\hline
Courant opération Wi-Fi (Pic Tx) & $500\,mA$ (Burst), $80\,mA$ nominal \\
\hline
Température opération & $-40°C$ -- $+85°C$ (Gamme Industrielle) \\
\hline
Dimensions module & $18 \times 25{,}5 \times 3{,}1$ mm \\
\hline
Certification & CE, FCC, TELEC, IC \\
\hline
\end{tabular}
\caption{Caractéristiques techniques clés de l'ESP32-WROOM-32U (Source : Espressif DS ESP32 Rev 3.3)}
\label{tab:esp32_specs}
\end{table}

\subsection{A.2.2 — Extrait de la Fiche Technique Driver DM556 (Leadshine)}

\begin{table}[htbp]
\centering
\renewcommand{\arraystretch}{1.3}
\begin{tabular}{|l|l|}
\hline
\rowcolor{macouleur!20}
\textbf{Paramètre} & \textbf{Valeur} \\
\hline
Courant moteur (pic) & 1{,}0 -- 5{,}6 A (réglable par DIP) \\
\hline
Tension d'alimentation & 18 -- 50 V DC \\
\hline
Micro-stepping & 400 / 800 / 1~600 / 3~200 / 6~400 / 12~800 / 25~600 pas/tour \\
\hline
Fréquence Step max & 200 kHz \\
\hline
Temps de réponse Step & $< 1\,\mu s$ (Isolation optique) \\
\hline
Niveau logique entrée & $3{,}5$ -- $5{,}5$ V (TTL 5V requis) \\
\hline
Courant d'entrée ($I_f$) & $10$ -- $14$ mA \\
\hline
Protection thermique & Oui ($75°C$ coupure auto) \\
\hline
Protection court-circuit & Oui \\
\hline
Protection surtension & Oui (Moteur back-EMF) \\
\hline
Rendement & $\approx 95\%$ \\
\hline
Dimensions & $118 \times 75.5 \times 34$ mm \\
\hline
\end{tabular}
\caption{Caractéristiques techniques clés du Driver DM556 (Source : Leadshine DS 2018)}
\label{tab:dm556_specs}
\end{table}

\subsection{A.2.3 — Caractéristiques du Moteur NEMA 17HS4401}

\begin{table}[htbp]
\centering
\renewcommand{\arraystretch}{1.3}
\begin{tabular}{|l|l|}
\hline
\rowcolor{macouleur!20}
\textbf{Paramètre} & \textbf{Valeur} \\
\hline
Pas angulaire (plein pas) & $1{,}8°$ (200 pas/tour) \\
\hline
Holding Torque & $0{,}45$ Nm \\
\hline
Courant de phase (nominal) & $1{,}7$ A \\
\hline
Résistance de phase & $1{,}5\,\Omega$ \\
\hline
Inductance de phase & $3{,}6$ mH \\
\hline
Inertie rotor & $68\,g.cm^2$ \\
\hline
Longueur (hors arbre) & 40 mm \\
\hline
Poids & 0{,}34 kg \\
\hline
Taille boîtier & $42{,}3 \times 42{,}3$ mm (NEMA 17) \\
\hline
Interface arbre & $\varnothing\, 5$ mm (simple méplat) \\
\hline
\end{tabular}
\caption{Caractéristiques techniques du moteur pas-à-pas NEMA 17HS4401}
\label{tab:nema17_specs}
\end{table}

% ============================================
\section{A.3 — Données Brutes des Tests de Métrologie (Axe Y -- 20 Mesures)}

\begin{table}[htbp]
\centering
\renewcommand{\arraystretch}{1.3}
\begin{tabular}{|c|c|c||c|c|c|}
\hline
\rowcolor{macouleur!20}
\textbf{N° Mesure} & \textbf{Valeur ($\mu m$)} & \textbf{Écart ($\mu m$)} & \textbf{N° Mesure} & \textbf{Valeur ($\mu m$)} & \textbf{Écart ($\mu m$)} \\
\hline
1 & $-26$ & $+3$ & 11 & $-28$ & $+1$ \\
2 & $-31$ & $-2$ & 12 & $-32$ & $-3$ \\
3 & $-29$ & $0$ & 13 & $-27$ & $+2$ \\
4 & $-28$ & $+1$ & 14 & $-30$ & $-1$ \\
5 & $-33$ & $-4$ & 15 & $-29$ & $0$ \\
6 & $-27$ & $+2$ & 16 & $-27$ & $+2$ \\
7 & $-29$ & $0$ & 17 & $-32$ & $-3$ \\
8 & $-30$ & $-1$ & 18 & $-28$ & $+1$ \\
9 & $-26$ & $+3$ & 19 & $-31$ & $-2$ \\
10 & $-29$ & $0$ & 20 & $-29$ & $0$ \\
\hline
\multicolumn{2}{|c|}{\textbf{Moyenne}} & $\boldsymbol{\bar{x} = -29}$ & \multicolumn{2}{c|}{\textbf{Déviation Std.}} & $\boldsymbol{\sigma = 1{,}9\,\mu m}$ \\
\hline
\end{tabular}
\caption{Données brutes des 20 mesures métrologie de l'axe Y (déplacement consigne = 50 mm)}
\label{tab:raw_data_y}
\end{table}

% ============================================
\section{A.4 — Manuel Opérateur (Guide de Démarrage Rapide)}

\subsection{A.4.1 — Procédure de Démarrage}
\begin{enumerate}
    \item \textbf{Vérifications avant mise sous tension :}
    \begin{itemize}
        \item S'assurer que la zone de travail est libre d'obstacles.
        \item Vérifier que le bouton E-STOP (champignon rouge) est \textbf{enfoncé}.
        \item Contrôler visuellement les connexions des câbles moteurs.
    \end{itemize}
    \item \textbf{Mise sous tension :} Activer l'interrupteur général du coffret. L'afficheur OLED affiche «BOOTING...» puis «IDLE — WIFI OK».
    \item \textbf{Connexion Wi-Fi :} Connecter l'ordinateur/tablette au réseau \textbf{CNC5AXES\_AP} (Mot de passe: \texttt{cnc12345}).
    \item \textbf{Lancement de l'IHM Flutter :} Ouvrir l'application. L'indicateur de connexion passe au vert.
    \item \textbf{Déverrouillage E-STOP :} Tourner le champignon rouge d'un quart de tour (sens anti-horaire). Cliquer sur «Déverrouiller (\$X)» dans l'IHM.
    \item \textbf{Homing (Référencement) :} Cliquer sur «Homing (\$H)» dans l'IHM. La machine va automatiquement referencer les 5 axes sur leurs fins de course. La LED verte «CONNECTÉ» s'allume.
    \item \textbf{Chargement du fichier G-Code :} Utiliser l'onglet «Éditeur» pour charger le fichier \texttt{.nc}. Vérifier le résultat de la validation.
    \item \textbf{Lancement de l'usinage :} Appuyer sur «LANCER L'USINAGE» (bouton vert). Surveiller l'avancement dans le visualisateur 3D.
\end{enumerate}

\subsection{A.4.2 — Procédure d'Arrêt d'Urgence}
En cas de comportement anormal (vibrations anormales, odeur de brûlé, collision imminente) :
\begin{enumerate}
    \item Appuyer immédiatement sur le champignon rouge E-STOP.
    \item Cela coupe l'alimentation 36V et désactive tous les moteurs.
    \item Ne pas relancer la machine avant d'avoir identifié et corrigé la cause.
    \item Un redémarrage complet (séquence depuis l'étape 1) est requis.
\end{enumerate}

\subsection{A.4.3 — Planning de Maintenance Préventive}

\begin{table}[htbp]
\centering
\renewcommand{\arraystretch}{1.3}
\begin{tabular}{|c|p{5cm}|l|p{4cm}|}
\hline
\rowcolor{macouleur!20}
\textbf{Fréquence} & \textbf{Opération} & \textbf{Composant} & \textbf{Produit/Outil} \\
\hline
Avant chaque utilisation & Inspection visuelle des câbles & Câbles moteurs & Visuel \\
\hline
Avant chaque utilisation & Serrage des guides & Vis de serrage HGR15 & Clé Allen 4mm \\
\hline
Hebdomadaire & Nettoyage des copeaux & Fraise, rails, plateau & Soufflette + chiffon \\
\hline
Mensuel & Graissage des rails HGR15 & Patins à billes & Graisse lithium EP2 \\
\hline
Mensuel & Graissage des écrous vis à billes & SFU1605 x3 & Huile machine SAE 10 \\
\hline
Trimestriel & Test E-STOP (Fonctionnel) & Contacteur KM1 & Manuel \\
\hline
Semestriel & Vérification des courroies (tension) & Entraînements A+B & Tensiomètre \\
\hline
Annuel & Recalibration complète (métrologie) & 5 axes & Comparateur $1\,\mu m$ \\
\hline
Aux 500 h & Remplacement courroies GT2 & Entraînements A+B & Courroie GT2-6mm neuve \\
\hline
Aux 1~000 h & Remplacement joints OLED & Afficheur OLED & SSD1306 neuf \\
\hline
\end{tabular}
\caption{Planning de maintenance préventive de la CNC 5 axes}
\label{tab:maintenance}
\end{table}
      % Annexes Techniques : Code, Datasheets, Données Brutes, Manuel

\end{document}